% This LaTeX document needs to be compiled with XeLaTeX.
\documentclass[10pt]{article}
\usepackage[utf8]{inputenc}
\usepackage{ucharclasses}
\usepackage{amsmath}
\usepackage{amsfonts}
\usepackage{amssymb}
\usepackage[version=4]{mhchem}
\usepackage{extpfeil}
\usepackage{stmaryrd}
\usepackage{graphicx}
\usepackage[export]{adjustbox}
\graphicspath{ {./images/} }
\usepackage{multirow}
\usepackage{fvextra, csquotes}
\usepackage{hyperref}
\hypersetup{colorlinks=true, linkcolor=blue, filecolor=magenta, urlcolor=cyan,}
\urlstyle{same}
\usepackage{mathrsfs}
\usepackage{bbold}
\usepackage[fallback]{xeCJK}
\usepackage{polyglossia}
\usepackage{fontspec}
\usepackage{newunicodechar}
\IfFontExistsTF{Noto Serif CJK JP}
{\setCJKmainfont{Noto Serif CJK JP}}
{\IfFontExistsTF{STSong}
  {\setCJKmainfont{STSong}}
  {\IfFontExistsTF{Droid Sans Fallback}
    {\setCJKmainfont{Droid Sans Fallback}}
    {\setCJKmainfont{SimSun}}
}}

\setmainlanguage{english}
\setotherlanguages{german, arabic, thai, tamil}
\IfFontExistsTF{Noto Naskh Arabic}
{\newfontfamily\arabicfont{Noto Naskh Arabic}}
{\IfFontExistsTF{Al Bayan}
  {\newfontfamily\arabicfont{Al Bayan}}
  {\IfFontExistsTF{FreeSerif}
    {\newfontfamily\arabicfont{FreeSerif}}
    {\IfFontExistsTF{DejaVu Sans}
      {\newfontfamily\arabicfont{DejaVu Sans}}
      {\IfFontExistsTF{Tahoma}
        {\newfontfamily\arabicfont{Tahoma}}
        {\newfontfamily\arabicfont{Arial Unicode MS}}
}}}}
\IfFontExistsTF{Noto Serif Thai}
{\newfontfamily\thaifont{Noto Serif Thai}}
{\IfFontExistsTF{Thonburi}
  {\newfontfamily\thaifont{Thonburi}}
  {\IfFontExistsTF{FreeSerif}
    {\newfontfamily\thaifont{FreeSerif}}
    {\IfFontExistsTF{Tahoma}
      {\newfontfamily\thaifont{Tahoma}}
      {\newfontfamily\thaifont{Arial Unicode MS}}
}}}
\IfFontExistsTF{Noto Serif Tamil}
{\newfontfamily\tamilfont{Noto Serif Tamil}}
{\IfFontExistsTF{Tamil MN}
  {\newfontfamily\tamilfont{Tamil MN}}
  {\IfFontExistsTF{Lohit Tamil}
    {\newfontfamily\tamilfont{Lohit Tamil}}
    {\IfFontExistsTF{FreeSerif}
      {\newfontfamily\tamilfont{FreeSerif}}
      {\newfontfamily\tamilfont{Latha}}
}}}
\IfFontExistsTF{CMU Serif}
{\newfontfamily\lgcfont{CMU Serif}}
{\IfFontExistsTF{DejaVu Sans}
  {\newfontfamily\lgcfont{DejaVu Sans}}
  {\newfontfamily\lgcfont{Georgia}}
}
\setDefaultTransitions{\lgcfont}{}
\setTransitionsFor{Arabic}{\arabicfont}{\lgcfont}
\setTransitionsFor{Thai}{\thaifont}{\lgcfont}
\setTransitionsFor{Tamil}{\tamilfont}{\lgcfont}

\title{TOPIC 2D State functions and exact differentials }

\author{\includegraphics[max width=\textwidth]{2025_02_27_d4b95d9718f0b9e2e6b9g-658}
}
\date{}


%New command to display footnote whose markers will always be hidden
\let\svthefootnote\thefootnote
\newcommand\blfootnotetext[1]{%
  \let\thefootnote\relax\footnote{#1}%
  \addtocounter{footnote}{-1}%
  \let\thefootnote\svthefootnote%
}

%Overriding the \footnotetext command to hide the marker if its value is `0`
\let\svfootnotetext\footnotetext
\renewcommand\footnotetext[2][?]{%
  \if\relax#1\relax%
    \ifnum\value{footnote}=0\blfootnotetext{#2}\else\svfootnotetext{#2}\fi%
  \else%
    \if?#1\ifnum\value{footnote}=0\blfootnotetext{#2}\else\svfootnotetext{#2}\fi%
    \else\svfootnotetext[#1]{#2}\fi%
  \fi
}

\def\AA{\mathring{\mathrm{A}}}

\newunicodechar{ı}{\ifmmode\imath\else{$\imath$}\fi}

\begin{document}
\maketitle
\chapter{FOCUS 1 The properties of gases}
A gas is a form of matter that fills whatever container it occupies. This Focus establishes the properties of gases that are used throughout the text.

\subsection*{1A The perfect gas}
This Topic is an account of an idealized version of a gas, a 'perfect gas', and shows how its equation of state may be assembled from the experimental observations summarized by Boyle's law, Charles's law, and Avogadro's principle.\\
1A. 1 Variables of state; 1A. 2 Equations of state

\subsection*{1B The kinetic model}
A central feature of physical chemistry is its role in building models of molecular behaviour that seek to explain observed phenomena. A prime example of this procedure is the development of a molecular model of a perfect gas in terms of a collection of molecules (or atoms) in ceaseless, essentially random motion. As well as accounting for the gas laws, this model can be used to predict the average speed at which molecules move in a gas, and its dependence on temperature. In combination with the Boltzmann distribution (see the text's Prologue), the model can also be used to predict the spread of molecular speeds and its dependence on molecular mass and temperature.

\subsection*{1C Real gases}
The perfect gas is a starting point for the discussion of properties of all gases, and its properties are invoked throughout thermodynamics. However, actual gases, 'real gases', have properties that differ from those of perfect gases, and it is necessary to be able to interpret these deviations and build the effects of molecular attractions and repulsions into the model. The discussion of real gases is another example of how initially primitive models in physical chemistry are elaborated to take into account more detailed observations.\\
1C. 1 Deviations from perfect behaviour; 1C. 2 The van der Waals equation

\subsection*{Web resources What is an application of this material?}
The perfect gas law and the kinetic theory can be applied to the study of phenomena confined to a reaction vessel or encompassing an entire planet or star. In Impact 1 the gas laws are used in the discussion of meteorological phenomena-the weather. Impact 2 examines how the kinetic model of gases has a surprising application: to the discussion of dense stellar media, such as the interior of the Sun.

\section*{TOPIC 1A The perfect gas}
\subsection*{Why do you need to know this material?}
Equations related to perfect gases provide the basis for the development of many relations in thermodynamics. The perfect gas law is also a good first approximation for accounting for the properties of real gases.

\subsection*{What is the key idea?}
The perfect gas law, which is based on a series of empirical observations, is a limiting law that is obeyed increasingly well as the pressure of a gas tends to zero.

\subsection*{What do you need to know already?}
You need to know how to handle quantities and units in calculations, as reviewed in The chemist's toolkit 1. You also need to be aware of the concepts of pressure, volume, amount of substance, and temperature, all reviewed in The chemist's toolkit 2.

The properties of gases were among the first to be established quantitatively (largely during the seventeenth and eighteenth centuries) when the technological requirements of travel in balloons stimulated their investigation. These properties set the stage for the development of the kinetic model of gases, as discussed in Topic 1B.

\subsection*{1A. 1 Variables of state}
The physical state of a sample of a substance, its physical condition, is defined by its physical properties. Two samples of the same substance that have the same physical properties are in the same state. The variables needed to specify the state of a system are the amount of substance it contains, $n$, the volume it occupies, $V$, the pressure, $p$, and the temperature, $T$.

\subsection*{(a) Pressure}
The origin of the force exerted by a gas is the incessant battering of the molecules on the walls of its container. The collisions are so numerous that they exert an effectively steady force, which is experienced as a steady pressure. The SI unit

Table 1A. 1 Pressure units*

\begin{center}
\begin{tabular}{lll}
\hline
Name & Symbol & Value \\
\hline
pascal & Pa & $1 \mathrm{~Pa}=1 \mathrm{~N} \mathrm{~m}^{-2}, 1 \mathrm{~kg} \mathrm{~m}^{-1} \mathrm{~s}^{-2}$ \\
bar & bar & $1 \mathrm{bar}=10^{5} \mathrm{~Pa}$ \\
atmosphere & atm & $1 \mathrm{~atm}=101.325 \mathrm{kPa}$ \\
torr & Torr & $1 \mathrm{Torr}=(101325 / 760) \mathrm{Pa}=133.32 \ldots \mathrm{~Pa}$ \\
millimetres of mercury & mmHg & $1 \mathrm{mmHg}=133.322 \ldots \mathrm{~Pa}$ \\
pounds per square inch & psi & $1 \mathrm{psi}=6.894757 \ldots \mathrm{kPa}$ \\
\hline
\end{tabular}
\end{center}

*Values in bold are exact.\\
of pressure, the pascal $\left(\mathrm{Pa}, 1 \mathrm{~Pa}=1 \mathrm{Nm}^{-2}\right)$, is introduced in The chemist's toolkit 1 . Several other units are still widely used (Table 1A.1). A pressure of 1 bar is the standard pressure for reporting data; it is denoted $p^{\ominus}$.

If two gases are in separate containers that share a common movable wall (Fig. 1A.1), the gas that has the higher pressure will tend to compress (reduce the volume of) the gas that has lower pressure. The pressure of the high-pressure gas will fall as it expands and that of the low-pressure gas will rise as it is compressed. There will come a stage when the two pressures are equal and the wall has no further tendency to move. This condition of equality of pressure on either side of a movable wall is a state of mechanical equilibrium between the two gases. The pressure of a gas is therefore an indication of whether a container that contains the gas will be in mechanical equilibrium with another gas with which it shares a movable wall.\\
\includegraphics[max width=\textwidth, center]{2025_02_27_d4b95d9718f0b9e2e6b9g-036}

Figure 1A. 1 When a region of high pressure is separated from a region of low pressure by a movable wall, the wall will be pushed into one region or the other, as in (a) and (c). However, if the two pressures are identical, the wall will not move (b). The latter condition is one of mechanical equilibrium between the two regions.

\subsection*{The chemist's toolkit 1}
The result of a measurement is a physical quantity that is reported as a numerical multiple of a unit:

$$
\text { physical quantity }=\text { numerical value } \times \text { unit }
$$

It follows that units may be treated like algebraic quantities and may be multiplied, divided, and cancelled. Thus, the expression (physical quantity)/unit is the numerical value (a dimensionless quantity) of the measurement in the specified units. For instance, the mass $m$ of an object could be reported as $m=2.5 \mathrm{~kg}$ or $m / \mathrm{kg}=2.5$. In this instance the unit of mass is 1 kg , but it is common to refer to the unit simply as kg (and likewise for other units). See Table A. 1 in the Resource section for a list of units.\\
Although it is good practice to use only SI units, there will be occasions where accepted practice is so deeply rooted that physical quantities are expressed using other, non-SI units. By international convention, all physical quantities are represented by oblique (sloping) letters (for instance, $m$ for mass); units are given in roman (upright) letters (for instance $m$ for metre).\\
Units may be modified by a prefix that denotes a factor of a power of 10 . Among the most common SI prefixes are those\\
listed in Table A. 2 in the Resource section. Examples of the use of these prefixes are:

$$
1 \mathrm{~nm}=10^{-9} \mathrm{~m} \quad 1 \mathrm{ps}=10^{-12} \mathrm{~s} \quad 1 \mu \mathrm{~mol}=10^{-6} \mathrm{~mol}
$$

Powers of units apply to the prefix as well as the unit they modify. For example, $1 \mathrm{~cm}^{3}=1(\mathrm{~cm})^{3}$, and $\left(10^{-2} \mathrm{~m}\right)^{3}=10^{-6} \mathrm{~m}^{3}$. Note that $1 \mathrm{~cm}^{3}$ does not mean $1 \mathrm{c}\left(\mathrm{m}^{3}\right)$. When carrying out numerical calculations, it is usually safest to write out the numerical value of an observable in scientific notation (as $n . n n n \times 10^{n}$ ).\\
There are seven SI base units, which are listed in Table A. 3 in the Resource section. All other physical quantities may be expressed as combinations of these base units. Molar concentration (more formally, but very rarely, amount of substance concentration) for example, which is an amount of substance divided by the volume it occupies, can be expressed using the derived units of $\mathrm{mol} \mathrm{dm}^{-3}$ as a combination of the base units for amount of substance and length. A number of these derived combinations of units have special names and symbols. For example, force is reported in the derived unit newton, $1 \mathrm{~N}=$ $1 \mathrm{~kg} \mathrm{~m} \mathrm{~s}^{-2}$ (see Table A. 4 in the Resource section).

The pressure exerted by the atmosphere is measured with a barometer. The original version of a barometer (which was invented by Torricelli, a student of Galileo) was an inverted tube of mercury sealed at the upper end. When the column of mercury is in mechanical equilibrium with the atmosphere, the pressure at its base is equal to that exerted by the atmosphere. It follows that the height of the mercury column is proportional to the external pressure.

The pressure of a sample of gas inside a container is measured by using a pressure gauge, which is a device with properties that respond to the pressure. For instance, a Bayard-Alpert pressure gauge is based on the ionization of the molecules present in the gas and the resulting current of ions is interpreted in terms of the pressure. In a capacitance manometer, the deflection of a diaphragm relative to a fixed electrode is monitored through its effect on the capacitance of the arrangement. Certain semiconductors also respond to pressure and are used as transducers in solid-state pressure gauges.

\subsection*{(b) Temperature}
The concept of temperature is introduced in The chemist's toolkit 2. In the early days of thermometry (and still in laboratory practice today), temperatures were related to the length of a column of liquid, and the difference in lengths shown when the thermometer was first in contact with melting ice and then with boiling water was divided into 100 steps called 'degrees', the lower point being labelled 0 . This procedure led\\
to the Celsius scale of temperature. In this text, temperatures on the Celsius scale are denoted $\theta$ (theta) and expressed in $d e$ grees Celsius $\left({ }^{\circ} \mathrm{C}\right)$. However, because different liquids expand to different extents, and do not always expand uniformly over a given range, thermometers constructed from different materials showed different numerical values of the temperature between their fixed points. The pressure of a gas, however, can be used to construct a perfect-gas temperature scale that is independent of the identity of the gas. The perfect-gas scale turns out to be identical to the thermodynamic temperature scale (Topic 3A), so the latter term is used from now on to avoid a proliferation of names.

On the thermodynamic temperature scale, temperatures are denoted $T$ and are normally reported in kelvins ( K ; not ${ }^{\circ} \mathrm{K}$ ). Thermodynamic and Celsius temperatures are related by the exact expression


\begin{equation*}
T / \mathrm{K}=\theta /{ }^{\circ} \mathrm{C}+273.15 \tag{1A.1}
\end{equation*}
 [definition]

This relation is the current definition of the Celsius scale in terms of the more fundamental Kelvin scale. It implies that a difference in temperature of $1^{\circ} \mathrm{C}$ is equivalent to a difference of 1 K .

\subsection*{Brief illustration 1A. 1}
To express $25.00^{\circ} \mathrm{C}$ as a temperature in kelvins, eqn 1 A .1 is used to write

$$
T / \mathrm{K}=\left(25.00^{\circ} \mathrm{C}\right) /{ }^{\circ} \mathrm{C}+273.15=25.00+273.15=298.15
$$

\subsection*{The chemist's toolkit 2}
The state of a bulk sample of matter is defined by specifying the values of various properties. Among them are:

The mass, $m$, a measure of the quantity of matter present (unit: kilogram, kg ).\\
The volume, $V$, a measure of the quantity of space the sample occupies (unit: cubic metre, $\mathrm{m}^{3}$ ).\\
The amount of substance, $n$, a measure of the number of specified entities (atoms, molecules, or formula units) present (unit: mole, mol).\\
The amount of substance, $n$ (colloquially, 'the number of moles'), is a measure of the number of specified entities present in the sample. 'Amount of substance' is the official name of the quantity; it is commonly simplified to 'chemical amount' or simply 'amount'. A mole is currently defined as the number of carbon atoms in exactly 12 g of carbon-12. (In 2011 the decision was taken to replace this definition, but the change has not yet, in 2018, been implemented.) The number of entities per mole is called Avogadro's constant, $N_{\mathrm{A}}$; the currently accepted value is $6.022 \times 10^{23} \mathrm{~mol}^{-1}$ (note that $N_{\mathrm{A}}$ is a constant with units, not a pure number).\\
The molar mass of a substance, $M$ (units: formally $\mathrm{kg} \mathrm{mol}^{-1}$ but commonly $\mathrm{g} \mathrm{mol}^{-1}$ ) is the mass per mole of its atoms, its molecules, or its formula units. The amount of substance of specified entities in a sample can readily be calculated from its mass, by noting that

$$
n=\frac{m}{M}
$$

Amount of substance

A note on good practice Be careful to distinguish atomic or molecular mass (the mass of a single atom or molecule; unit: kg ) from molar mass (the mass per mole of atoms or molecules; units: $\mathrm{kg} \mathrm{mol}^{-1}$ ). Relative molecular masses of atoms and molecules, $M_{\mathrm{r}}=m / m_{\mathrm{u}}$, where $m$ is the mass of the atom or molecule and $m_{\mathrm{u}}$ is the atomic mass constant (see inside front cover), are still widely called 'atomic weights' and 'molecular weights' even though they are dimensionless quantities and not weights ('weight' is the gravitational force exerted on an object).

A sample of matter may be subjected to a pressure, $p$ (unit: pascal, $\mathrm{Pa} ; 1 \mathrm{~Pa}=1 \mathrm{~kg} \mathrm{~m}^{-1} \mathrm{~s}^{-2}$ ), which is defined as the force, $F$, it is subjected to, divided by the area, $A$, to which that force is applied. Although the pascal is the SI unit of pressure, it is also common to express pressure in bar $\left(1 \mathrm{bar}=10^{5} \mathrm{~Pa}\right)$ or atmospheres $(1 \mathrm{~atm}=101325 \mathrm{~Pa}$ exactly), both of which correspond to typical atmospheric pressure. Because many physical properties depend on the pressure acting on a sample, it is appropriate to select a certain value of the pressure to report their values. The standard pressure for reporting physical quantities is currently defined as $p^{\ominus}=1$ bar exactly.\\
To specify the state of a sample fully it is also necessary to give its temperature, $T$. The temperature is formally a property that determines in which direction energy will flow as heat when two samples are placed in contact through thermally conducting walls: energy flows from the sample with the higher temperature to the sample with the lower temperature. The symbol $T$ is used to denote the thermodynamic temperature which is an absolute scale with $T=0$ as the lowest point. Temperatures above $T=0$ are then most commonly expressed by using the Kelvin scale, in which the gradations of temperature are expressed in kelvins (K). The Kelvin scale is currently defined by setting the triple point of water (the temperature at which ice, liquid water, and water vapour are in mutual equilibrium) at exactly 273.16 K (as for certain other units, a decision has been taken to revise this definition, but it has not yet, in 2018, been implemented). The freezing point of water (the melting point of ice) at 1 atm is then found experimentally to lie 0.01 K below the triple point, so the freezing point of water is 273.15 K .\\
Suppose a sample is divided into smaller samples. If a property of the original sample has a value that is equal to the sum of its values in all the smaller samples (as mass would), then it is said to be extensive. Mass and volume are extensive properties. If a property retains the same value as in the original sample for all the smaller samples (as temperature would), then it is said to be intensive. Temperature and pressure are intensive properties. Mass density, $\rho=m / V$, is also intensive because it would have the same value for all the smaller samples and the original sample. All molar properties, $X_{\mathrm{m}}=X / n$, are intensive, whereas $X$ and $n$ are both extensive.\\
$p=0$, regardless of the size of the units, such as bar or pascal). However, it is appropriate to write $0^{\circ} \mathrm{C}$ because the Celsius scale is not absolute.

\subsection*{1A. 2 Equations of state}
Although in principle the state of a pure substance is specified by giving the values of $n, V, p$, and $T$, it has been established experimentally that it is sufficient to specify only three of these variables since doing so fixes the value of the fourth variable.

That is, it is an experimental fact that each substance is described by an equation of state, an equation that interrelates these four variables.

The general form of an equation of state is


\begin{equation*}
p=f(T, V, n) \quad \text { General form of an equation of state } \tag{1A.2}
\end{equation*}


This equation states that if the values of $n, T$, and $V$ are known for a particular substance, then the pressure has a fixed value. Each substance is described by its own equation of state, but the explicit form of the equation is known in only a few special cases. One very important example is the equation of state of a 'perfect gas', which has the form $p=n R T / V$, where $R$ is a constant independent of the identity of the gas.

The equation of state of a perfect gas was established by combining a series of empirical laws.

\subsection*{(a) The empirical basis}
The following individual gas laws should be familiar:\\
Boyle's law: $\quad p V=$ constant, at constant $n, T$\\
Charles's law: $\quad V=$ constant $\times T$, at constant $n, p \quad$ (1A.3b)


\begin{equation*}
p=\text { constant } \times T, \text { at constant } n, V \tag{1A.3c}
\end{equation*}


Avogadro's principle:


\begin{equation*}
V=\text { constant } \times n \text { at constant } p, T \tag{1A.3d}
\end{equation*}


Boyle's and Charles's laws are examples of a limiting law, a law that is strictly true only in a certain limit, in this case $p \rightarrow 0$. For example, if it is found empirically that the volume of a substance fits an expression $V=a T+b p+c p^{2}$, then in the limit of $p \rightarrow 0, V=a T$. Many relations that are strictly true only at $p=0$ are nevertheless reasonably reliable at normal pressures ( $p \approx 1$ bar) and are used throughout chemistry.

Figure 1A. 2 depicts the variation of the pressure of a sample of gas as the volume is changed. Each of the curves in the\\
\includegraphics[max width=\textwidth, center]{2025_02_27_d4b95d9718f0b9e2e6b9g-039}

Figure 1A. 2 The pressure-volume dependence of a fixed amount of perfect gas at different temperatures. Each curve is a hyperbola ( $p V=$ constant) and is called an isotherm.\\
\includegraphics[max width=\textwidth, center]{2025_02_27_d4b95d9718f0b9e2e6b9g-039(1)}

Figure 1A. 3 Straight lines are obtained when the pressure of a perfect gas is plotted against $1 / V$ at constant temperature. These lines extrapolate to zero pressure at $1 / V=0$.\\
graph corresponds to a single temperature and hence is called an isotherm. According to Boyle's law, the isotherms of gases are hyperbolas (a curve obtained by plotting $y$ against $x$ with $x y=$ constant, or $y=\operatorname{constant} / x$ ). An alternative depiction, a plot of pressure against $1 /$ volume, is shown in Fig. 1A.3. The linear variation of volume with temperature summarized by Charles's law is illustrated in Fig. 1A.4. The lines in this illustration are examples of isobars, or lines showing the variation of properties at constant pressure. Figure 1A. 5 illustrates the linear variation of pressure with temperature. The lines in this diagram are isochores, or lines showing the variation of properties at constant volume.

A note on good practice To test the validity of a relation between two quantities, it is best to plot them in such a way that they should give a straight line, because deviations from a straight line are much easier to detect than deviations from a curve. The development of expressions that, when plotted, give a straight line is a very important and common procedure in physical chemistry.\\
\includegraphics[max width=\textwidth, center]{2025_02_27_d4b95d9718f0b9e2e6b9g-039(2)}

Figure 1A. 4 The variation of the volume of a fixed amount of a perfect gas with the temperature at constant pressure. Note that in each case the isobars extrapolate to zero volume at $T=0$, corresponding to $\theta=-273.15^{\circ} \mathrm{C}$.\\
\includegraphics[max width=\textwidth, center]{2025_02_27_d4b95d9718f0b9e2e6b9g-040(1)}

Figure 1A. 5 The pressure of a perfect gas also varies linearly with the temperature at constant volume, and extrapolates to zero at $T=0\left(-273.15^{\circ} \mathrm{C}\right)$.

The empirical observations summarized by eqn 1A. 3 can be combined into a single expression:

$$
p V=\text { constant } \times n T
$$

This expression is consistent with Boyle's law ( $p V=$ constant) when $n$ and $T$ are constant, with both forms of Charles's law $(p \propto T, V \propto T)$ when $n$ and either $V$ or $p$ are held constant, and with Avogadro's principle $(V \propto n)$ when $p$ and $T$ are constant. The constant of proportionality, which is found experimentally to be the same for all gases, is denoted $R$ and called the (molar) gas constant. The resulting expression

$$
p V=n R T
$$

Perfect gas law\\
(1A.4)\\
is the perfect gas law (or perfect gas equation of state). It is the approximate equation of state of any gas, and becomes increasingly exact as the pressure of the gas approaches zero. A gas that obeys eqn 1A. 4 exactly under all conditions is called a perfect gas (or ideal gas). A real gas, an actual gas, behaves more like a perfect gas the lower the pressure, and is described exactly by eqn 1 A .4 in the limit of $p \rightarrow 0$. The gas constant $R$ can be determined by evaluating $R=p V / n T$ for a gas in the limit of zero pressure (to guarantee that it is behaving perfectly).

A note on good practice Despite 'ideal gas' being the more common term, 'perfect gas' is preferable. As explained in Topic 5B, in an 'ideal mixture' of A and B , the $\mathrm{AA}, \mathrm{BB}$, and $A B$ interactions are all the same but not necessarily zero. In a perfect gas, not only are the interactions all the same, they are also zero.

The surface in Fig. 1A. 6 is a plot of the pressure of a fixed amount of perfect gas against its volume and thermodynamic temperature as given by eqn 1A.4. The surface depicts the only possible states of a perfect gas: the gas cannot exist in states that do not correspond to points on the surface. The graphs in Figs. 1A. 2 and 1A. 4 correspond to the sections through the surface (Fig. 1A.7).\\
\includegraphics[max width=\textwidth, center]{2025_02_27_d4b95d9718f0b9e2e6b9g-040}

Figure 1A.6 A region of the $p, V, T$ surface of a fixed amount of perfect gas. The points forming the surface represent the only states of the gas that can exist.\\
\includegraphics[max width=\textwidth, center]{2025_02_27_d4b95d9718f0b9e2e6b9g-040(2)}

Figure 1A. 7 Sections through the surface shown in Fig. 1A. 6 at constant temperature give the isotherms shown in Fig. 1A.2. Sections at constant pressure give the isobars shown in Fig. 1A.4. Sections at constant volume give the isochores shown in Fig. 1A.5.

\subsection*{Example 1A. 1 Using the perfect gas law}
In an industrial process, nitrogen gas is introduced into a vessel of constant volume at a pressure of 100 atm and a temperature of 300 K . The gas is then heated to 500 K . What pressure would the gas then exert, assuming that it behaved as a perfect gas?\\
Collect your thoughts The pressure is expected to be greater on account of the increase in temperature. The perfect gas law in the form $p V / n T=R$ implies that if the conditions are changed from one set of values to another, then because $p V / n T$ is equal to a constant, the two sets of values are related by the 'combined gas law'


\begin{equation*}
\frac{p_{1} V_{1}}{n_{1} T_{1}}=\frac{p_{2} V_{2}}{n_{2} T_{2}} \quad \text { Combined gas law } \tag{1A.5}
\end{equation*}


This expression is easily rearranged to give the unknown quantity (in this case $p_{2}$ ) in terms of the known. The known and unknown data are summarized as follows:

\begin{center}
\begin{tabular}{lllll}
\hline
 & $\boldsymbol{n}$ & $\boldsymbol{p}$ & $\boldsymbol{V}$ & $\boldsymbol{T}$ \\
\hline
Initial & Same & 100 atm & Same & 300 K \\
Final & Same & $?$ & Same & 500 K \\
\hline
\end{tabular}
\end{center}

The solution Cancellation of the volumes (because $V_{1}=V_{2}$ ) and amounts (because $n_{1}=n_{2}$ ) on each side of the combined gas law results in

$$
\frac{p_{1}}{T_{1}}=\frac{p_{2}}{T_{2}}
$$

which can be rearranged into

$$
p_{2}=\frac{T_{2}}{T_{1}} \times p_{1}
$$

Substitution of the data then gives

$$
p_{2}=\frac{500 \mathrm{~K}}{300 \mathrm{~K}} \times(100 \mathrm{~atm})=167 \mathrm{~atm}
$$

Self-test 1A. 1 What temperature would result in the same sample exerting a pressure of 300 atm ?

The perfect gas law is of the greatest importance in physical chemistry because it is used to derive a wide range of relations that are used throughout thermodynamics. However, it is also of considerable practical utility for calculating the properties of a gas under a variety of conditions. For instance, the molar volume, $V_{\mathrm{m}}=V / n$, of a perfect gas under the conditions called standard ambient temperature and pressure (SATP), which means 298.15 K and 1 bar (i.e. exactly $10^{5} \mathrm{~Pa}$ ), is easily calculated from $V_{\mathrm{m}}=R T / p$ to be $24.789 \mathrm{dm}^{3} \mathrm{~mol}^{-1}$. An earlier definition, standard temperature and pressure (STP), was $0^{\circ} \mathrm{C}$ and 1 atm ; at STP, the molar volume of a perfect gas is $22.414 \mathrm{dm}^{3} \mathrm{~mol}^{-1}$.

The molecular explanation of Boyle's law is that if a sample of gas is compressed to half its volume, then twice as many molecules strike the walls in a given period of time than before it was compressed. As a result, the average force exerted on the walls is doubled. Hence, when the volume is halved the pressure of the gas is doubled, and $p V$ is a constant. Boyle's law applies to all gases regardless of their chemical identity (provided the pressure is low) because at low pressures the average separation of molecules is so great that they exert no influence on one another and hence travel independently. The molecular explanation of Charles's law lies in the fact that raising the temperature of a gas increases the average speed of its molecules. The molecules collide with the walls more frequently and with greater impact. Therefore they exert a greater pressure on the walls of the container. For a quantitative account of these relations, see Topic 1B.

\subsection*{(b) Mixtures of gases}
When dealing with gaseous mixtures, it is often necessary to know the contribution that each component makes to the total pressure of the sample. The partial pressure, $p_{\mathrm{J}}$, of a gas J in a mixture (any gas, not just a perfect gas), is defined as

\[
\begin{array}{ll}
p_{\mathrm{J}}=x_{\mathrm{J}} p & \begin{array}{l}
\text { Partial pressure } \\
\text { [definition] }
\end{array} \tag{1A.6}
\end{array}
\]

where $x_{\mathrm{J}}$ is the mole fraction of the component J , the amount of J expressed as a fraction of the total amount of molecules, $n$, in the sample:

\[
x_{\mathrm{J}}=\frac{n_{\mathrm{J}}}{n} \quad n=n_{\mathrm{A}}+n_{\mathrm{B}}+\cdots \quad \begin{align*}
& \text { Mole fraction }  \tag{1A.7}\\
& \text { [definition] }
\end{align*}
\]

When no J molecules are present, $x_{\mathrm{J}}=0$; when only J molecules are present, $x_{\mathrm{J}}=1$. It follows from the definition of $x_{\mathrm{J}}$ that, whatever the composition of the mixture, $x_{\mathrm{A}}+x_{\mathrm{B}}+\cdots=1$ and therefore that the sum of the partial pressures is equal to the total pressure:


\begin{equation*}
p_{\mathrm{A}}+p_{\mathrm{B}}+\cdots=\left(x_{\mathrm{A}}+x_{\mathrm{B}}+\cdots\right) p=p \tag{1A.8}
\end{equation*}


This relation is true for both real and perfect gases.\\
When all the gases are perfect, the partial pressure as defined in eqn 1A. 6 is also the pressure that each gas would exert if it occupied the same container alone at the same temperature. The latter is the original meaning of 'partial pressure'. That identification was the basis of the original formulation of Dalton's law:

The pressure exerted by a mixture of gases is the sum of the pressures that each one would exert if it occupied the container alone.

Dalton's law\\
This law is valid only for mixtures of perfect gases, so it is not used to define partial pressure. Partial pressure is defined by eqn 1A.6, which is valid for all gases.

\subsection*{Example 1A. 2 Calculating partial pressures}
The mass percentage composition of dry air at sea level is approximately $\mathrm{N}_{2}: 75.5 ; \mathrm{O}_{2}: 23.2$; Ar: 1.3. What is the partial pressure of each component when the total pressure is 1.20 atm ?

Collect your thoughts Partial pressures are defined by eqn 1A.6. To use the equation, first calculate the mole fractions of the components, by using eqn 1A. 7 and the fact that the amount of atoms or molecules J of molar mass $M_{\mathrm{J}}$ in a sample of mass $m_{\mathrm{J}}$ is $n_{\mathrm{J}}=m_{\mathrm{J}} / M_{\mathrm{J}}$. The mole fractions are independent of the total mass of the sample, so choose the latter to be exactly 100 g (which makes the conversion from mass percentages very easy). Thus, the mass of $\mathrm{N}_{2}$ present is 75.5 per cent of 100 g , which is 75.5 g .

The solution The amounts of each type of atom or molecule present in 100 g of air are, in which the masses of $\mathrm{N}_{2}, \mathrm{O}_{2}$, and Ar are $75.5 \mathrm{~g}, 23.2 \mathrm{~g}$, and 1.3 g , respectively, are

$$
\begin{aligned}
& n\left(\mathrm{~N}_{2}\right)=\frac{75.5 \mathrm{~g}}{28.02 \mathrm{~g} \mathrm{~mol}^{-1}}=\frac{75.5}{28.02} \mathrm{~mol}=2.69 \mathrm{~mol} \\
& n\left(\mathrm{O}_{2}\right)=\frac{23.2 \mathrm{~g}}{32.00 \mathrm{~g} \mathrm{~mol}^{-1}}=\frac{23.2}{32.00} \mathrm{~mol}=0.725 \mathrm{~mol} \\
& n(\mathrm{Ar})=\frac{1.3 \mathrm{~g}}{39.95 \mathrm{~g} \mathrm{~mol}^{-1}}=\frac{1.3}{39.95} \mathrm{~mol}=0.033 \mathrm{~mol}
\end{aligned}
$$

The total is 3.45 mol . The mole fractions are obtained by dividing each of the above amounts by 3.45 mol and the partial pressures are then obtained by multiplying the mole fraction by the total pressure ( 1.20 atm ):

\begin{center}
\begin{tabular}{llll}
 & $\mathrm{N}_{2}$ & $\mathrm{O}_{2}$ & Ar \\
Mole fraction: & 0.780 & 0.210 & 0.0096 \\
Partial pressure/atm: & 0.936 & 0.252 & 0.012 \\
\end{tabular}
\end{center}

Self-test 1A. 2 When carbon dioxide is taken into account, the mass percentages are $75.52\left(\mathrm{~N}_{2}\right), 23.15\left(\mathrm{O}_{2}\right), 1.28(\mathrm{Ar})$, and $0.046\left(\mathrm{CO}_{2}\right)$. What are the partial pressures when the total pressure is 0.900 atm ?\\
\includegraphics[max width=\textwidth, center]{2025_02_27_d4b95d9718f0b9e2e6b9g-042}

\subsection*{Checklist of concepts}
$\square$ 1. The physical state of a sample of a substance, its physical condition, is defined by its physical properties.2. Mechanical equilibrium is the condition of equality of pressure on either side of a shared movable wall.3. An equation of state is an equation that interrelates the variables that define the state of a substance.4. Boyle's and Charles's laws are examples of a limiting law, a law that is strictly true only in a certain limit, in this case $p \rightarrow 0$.5. An isotherm is a line in a graph that corresponds to a single temperature.6. An isobar is a line in a graph that corresponds to a single pressure.7. An isochore is a line in a graph that corresponds to a single volume.8. A perfect gas is a gas that obeys the perfect gas law under all conditions.\\
9. Dalton's law states that the pressure exerted by a mixture of (perfect) gases is the sum of the pressures that each one would exert if it occupied the container alone.

\subsection*{Checklist of equations}
\begin{center}
\begin{tabular}{llll}
\hline
Property & Equation & Comment & Equation number \\
\hline
Relation between temperature scales & $T / \mathrm{K}=\theta /{ }^{\circ} \mathrm{C}+273.15$ & 273.15 is exact & 1 A .1 \\
Perfect gas law & $p V=n R T$ & Valid for real gases in the limit $p \rightarrow 0$ & 1 A .4 \\
Partial pressure & $p_{\mathrm{J}}=x_{\mathrm{J}} p$ & Valid for all gases & 1 A .6 \\
Mole fraction & $x_{\mathrm{J}}=n_{\mathrm{J}} / n$ & Definition & 1 A .7 \\
 & $n=n_{\mathrm{A}}+n_{\mathrm{B}}+\cdots$ &  &  \\
\hline
\end{tabular}
\end{center}

\section*{TOPIC 1B The kinetic model}
\subsection*{Why do you need to know this material?}
This material illustrates an important skill in science: the ability to extract quantitative information from a qualitative model. Moreover, the model is used in the discussion of the transport properties of gases (Topic 16A), reaction rates in gases (Topic 18A), and catalysis (Topic 19C).

\subsection*{What is the key idea?}
A gas consists of molecules of negligible size in ceaseless random motion and obeying the laws of classical mechanics in their collisions.

\subsection*{What do you need to know already?}
You need to be aware of Newton's second law of motion, that the acceleration of a body is proportional to the force acting on it, and the conservation of linear momentum (The chemist's toolkit 3).

\subsection*{The chemist's toolkit 3}
\subsection*{Momentum and force}
The speed, $v$, of a body is defined as the rate of change of position. The velocity, $v$, defines the direction of travel as well as the rate of motion, and particles travelling at the same speed but in different directions have different velocities. As shown in Sketch 1, the velocity can be depicted as an arrow in the direction of travel, its length being the speed $v$ and its components $v_{x}, v_{y}$, and $v_{z}$ along three perpendicular axes. These components have a sign: $v_{x}=+5 \mathrm{~m} \mathrm{~s}^{-1}$, for instance, indicates that a body is moving in the positive $x$-direction, whereas $v_{x}=$ $-5 \mathrm{~m} \mathrm{~s}^{-1}$ indicates that it is moving in the opposite direction. The length of the arrow (the speed) is related to the components by Pythagoras' theorem: $v^{2}=v_{x}^{2}+v_{y}^{2}+v_{z}^{2}$.\\
\includegraphics[max width=\textwidth, center]{2025_02_27_d4b95d9718f0b9e2e6b9g-043}

Sketch 1

In the kinetic theory of gases (which is sometimes called the kinetic-molecular theory, KMT) it is assumed that the only contribution to the energy of the gas is from the kinetic energies of the molecules. The kinetic model is one of the most re-markable-and arguably most beautiful-models in physical chemistry, for from a set of very slender assumptions, powerful quantitative conclusions can be reached.

\subsection*{1B. 1 The model}
The kinetic model is based on three assumptions:

\begin{enumerate}
  \item The gas consists of molecules of mass $m$ in ceaseless random motion obeying the laws of classical mechanics.
  \item The size of the molecules is negligible, in the sense that their diameters are much smaller than the average distance travelled between collisions; they are 'point-like'.
  \item The molecules interact only through brief elastic collisions.
\end{enumerate}

The concepts of classical mechanics are commonly expressed in terms of the linear momentum, $p$, which is defined as

$$
p=m v
$$

Linear momentum [definition]\\
Momentum also mirrors velocity in having a sense of direction; bodies of the same mass and moving at the same speed but in different directions have different linear momenta.\\
Acceleration, $\boldsymbol{a}$, is the rate of change of velocity. A body accelerates if its speed changes. A body also accelerates if its speed remains unchanged but its direction of motion changes. According to Newton's second law of motion, the acceleration of a body of mass $m$ is proportional to the force, $F$, acting on it:

$$
F=m a
$$

Force\\
Because $m v$ is the linear momentum and $\boldsymbol{a}$ is the rate of change of velocity, $m a$ is the rate of change of momentum. Therefore, an alternative statement of Newton's second law is that the force is equal to the rate of change of momentum. Newton's law indicates that the acceleration occurs in the same direction as the force acts. If, for an isolated system, no external force acts, then there is no acceleration. This statement is the law of conservation of momentum: that the momentum of a body is constant in the absence of a force acting on the body.

An elastic collision is a collision in which the total translational kinetic energy of the molecules is conserved.

\subsection*{(a) Pressure and molecular speeds}
From the very economical assumptions of the kinetic model, it is possible to derive an expression that relates the pressure and volume of a gas.

\subsection*{How is that done? 1B. 1 Using the kinetic model to derive an expression for the pressure of a gas}
Consider the arrangement in Fig. 1B.1, and then follow these steps.

\subsection*{Step 1 Set up the calculation of the change in momentum}
When a particle of mass $m$ that is travelling with a component of velocity $v_{x}$ parallel to the $x$-axis collides with the wall on the right and is reflected, its linear momentum changes from $m v_{x}$ before the collision to $-m v_{x}$ after the collision (when it is travelling in the opposite direction). The $x$-component of momentum therefore changes by $2 m v_{x}$ on each collision (the $y$ - and $z$-components are unchanged). Many molecules collide with the wall in an interval $\Delta t$, and the total change of momentum is the product of the change in momentum of each molecule multiplied by the number of molecules that reach the wall during the interval.

\subsection*{Step 2 Calculate the change in momentum}
Because a molecule with velocity component $v_{x}$ travels a distance $v_{x} \Delta t$ along the $x$-axis in an interval $\Delta t$, all the molecules within a distance $v_{x} \Delta t$ of the wall strike it if they are travelling towards it (Fig. 1B.2). It follows that if the wall has area $A$, then all the particles in a volume $A \times v_{x} \Delta t$ reach the wall (if they are travelling towards it). The number density of particles is $n N_{\mathrm{A}} / V$, where $n$ is the total amount of molecules in the container of volume $V$ and $N_{\mathrm{A}}$ is Avogadro's constant. It follows that the number of molecules in the volume $A v_{x} \Delta t$ is $\left(n N_{\mathrm{A}} / V\right) \times A v_{x} \Delta t$.\\
\includegraphics[max width=\textwidth, center]{2025_02_27_d4b95d9718f0b9e2e6b9g-044(1)}

Figure 1B. 1 The pressure of a gas arises from the impact of its molecules on the walls. In an elastic collision of a molecule with a wall perpendicular to the $x$-axis, the $x$-component of velocity is reversed but the $y$-and $z$-components are unchanged.\\
\includegraphics[max width=\textwidth, center]{2025_02_27_d4b95d9718f0b9e2e6b9g-044}

Figure 1B. 2 A molecule will reach the wall on the right within an interval of time $\Delta t$ if it is within a distance $v_{x} \Delta t$ of the wall and travelling to the right.

At any instant, half the particles are moving to the right and half are moving to the left. Therefore, the average number of collisions with the wall during the interval $\Delta t$ is $\frac{1}{2} n N_{A} A v_{x} \Delta t / V$. The total momentum change in that interval is the product of this number and the change $2 m v_{x}$ :

$$
\begin{aligned}
\text { Momentum change } & =\frac{n N_{\mathrm{A}} A v_{x} \Delta t}{2 V} \times 2 m v_{x} \\
& =\frac{n \overbrace{m N_{\mathrm{A}}}^{M} A v_{x}^{2} \Delta t}{V}=\frac{n M A v_{x}^{2} \Delta t}{V}
\end{aligned}
$$

\subsection*{Step 3 Calculate the force}
The rate of change of momentum, the change of momentum divided by the interval $\Delta t$ during which it occurs, is

$$
\text { Rate of change of momentum }=\frac{n M A v_{x}^{2}}{V}
$$

According to Newton's second law of motion this rate of change of momentum is equal to the force.

Step 4 Calculate the pressure\\
The pressure is this force $\left(n M A \nu_{x}^{2} / V\right)$ divided by the area ( $A$ ) on which the impacts occur. The areas cancel, leaving

$$
\text { Pressure }=\frac{n M v_{x}^{2}}{V}
$$

Not all the molecules travel with the same velocity, so the detected pressure, $p$, is the average (denoted $\langle\ldots\rangle$ ) of the quantity just calculated:

$$
p=\frac{n M\left\langle v_{x}^{2}\right\rangle}{V}
$$

The average values of $v_{x}^{2}, v_{y}^{2}$, and $v_{z}^{2}$ are all the same, and because $v^{2}=v_{x}^{2}+v_{y}^{2}+v_{z}^{2}$, it follows that $\left\langle v_{x}^{2}\right\rangle=\frac{1}{3}\left\langle v^{2}\right\rangle$.

At this stage it is useful to define the root-mean-square speed, $v_{\text {rms }}$, as the square root of the mean of the squares of the speeds, $v$, of the molecules. Therefore

\[
v_{\text {rms }}=\left\langle v^{2}\right\rangle^{1 / 2} \quad \begin{align*}
& \text { Root-mean-square speed }  \tag{1B.1}\\
& \text { [definition] }
\end{align*}
\]

The mean square speed in the expression for the pressure can therefore be written $\left\langle v_{x}^{2}\right\rangle=\frac{1}{3}\left\langle v^{2}\right\rangle=\frac{1}{3} v_{\text {rms }}^{2}$ to give

$$
\left.\begin{array}{l|l} 
\\
𠃌
\end{array} p V=\frac{1}{3} n M \nu_{\mathrm{rms}}^{2} \right\rvert\, \begin{aligned}
& \text { Relation between pressure and volume } \\
& \text { [KMT] }
\end{aligned}
$$

This equation is one of the key results of the kinetic model. If the root-mean-square speed of the molecules depends only on the temperature, then at constant temperature

$$
p V=\text { constant }
$$

which is the content of Boyle's law. The task now is to show that the right-hand side of eqn $1 B .2$ is equal to $n R T$.

\subsection*{(b) The Maxwell-Boltzmann distribution of speeds}
In a gas the speeds of individual molecules span a wide range, and the collisions in the gas ensure that their speeds are ceaselessly changing. Before a collision, a molecule may be travelling rapidly, but after a collision it may be accelerated to a higher speed, only to be slowed again by the next collision. To evaluate the root-mean-square speed it is necessary to know the fraction of molecules that have a given speed at any instant. The fraction of molecules that have speeds in the range $v$ to $v+\mathrm{d} v$ is proportional to the width of the range, and is written $f(v) \mathrm{d} v$, where $f(v)$ is called the distribution of speeds. An expression for this distribution can be found by recognizing that the energy of the molecules is entirely kinetic, and then using the Boltzmann distribution to describe how this energy is distributed over the molecules.

\subsection*{How is that done? 1B. 2 Deriving the distribution}
 of speedsThe starting point for this derivation is the Boltzmann distribution (see the text's Prologue).

Step 1 Write an expression for the distribution of the kinetic energy\\
The Boltzmann distribution implies that the fraction of molecules with velocity components $v_{x}, v_{y}$, and $v_{z}$ is proportional to an exponential function of their kinetic energy: $f(v)=K \mathrm{e}^{-\varepsilon / k T}$, where $K$ is a constant of proportionality. The kinetic energy is

$$
\varepsilon=\frac{1}{2} m v_{x}^{2}+\frac{1}{2} m v_{y}^{2}+\frac{1}{2} m v_{z}^{2}
$$

Therefore, use the relation $a^{x+y+z}=a^{x} a^{y} a^{z}$ to write

$$
f(v)=K \mathrm{e}^{-\left(m v_{x}^{2}+m v_{y}^{2}+m v_{z}^{2}\right) / 2 k T}=K \mathrm{e}^{-m v_{x}^{2} / 2 k T} \mathrm{e}^{-m v_{y}^{2} / 2 k T} \mathrm{e}^{-m v_{\Sigma}^{2} / 2 k T}
$$

The distribution factorizes into three terms as $f(v)=f\left(v_{x}\right) f\left(v_{y}\right) f\left(v_{z}\right)$ and $K=K_{x} K_{y} K_{z}$, with

$$
f\left(v_{x}\right)=K_{x} \mathrm{e}^{-m v_{x}^{2} / 2 k T}
$$

and likewise for the other two coordinates.\\
Step 2 Determine the constants $K_{x}, K_{y}$ and $K_{z}$\\
To determine the constant $K_{x}$, note that a molecule must have a velocity component somewhere in the range $-\infty<v_{x}<\infty$, so integration over the full range of possible values of $v_{x}$ must give a total probability of 1 :

$$
\int_{-\infty}^{\infty} f\left(v_{x}\right) \mathrm{d} v_{x}=1
$$

(See The chemist's toolkit 4 for the principles of integration.) Substitution of the expression for $f\left(v_{x}\right)$ then gives

$$
1=K_{x} \overbrace{\int_{-\infty}^{\infty} \mathrm{e}^{-m v_{x}^{2} / 2 k T} \mathrm{~d} v_{x}}^{\text {Integral G. } 1}=K_{x}\left(\frac{2 \pi k T}{m}\right)^{1 / 2}
$$

Therefore, $K_{x}=(m / 2 \pi k T)^{1 / 2}$ and


\begin{equation*}
f\left(v_{x}\right)=\left(\frac{m}{2 \pi k T}\right)^{1 / 2} \mathrm{e}^{-m v_{x}^{2} / 2 k T} \tag{1B.3}
\end{equation*}


The expressions for $f\left(v_{y}\right)$ and $f\left(v_{z}\right)$ are analogous.\\
Step 3 Write a preliminary expression for\\
$f\left(v_{x}\right) f\left(v_{y}\right) f\left(v_{z}\right) \mathrm{d} v_{x} \mathrm{~d} v_{y} \mathrm{~d} v_{z}$\\
The probability that a molecule has a velocity in the range $v_{x}$ to $v_{x}+\mathrm{d} v_{x}, v_{y}$ to $v_{y}+\mathrm{d} v_{y}, v_{z}$ to $v_{z}+\mathrm{d} v_{z}$, is

$$
\begin{aligned}
f\left(v_{x}\right) f\left(v_{y}\right) f\left(v_{z}\right) \mathrm{d} v_{x} \mathrm{~d} v_{y} \mathrm{~d} v_{z}= & \left(\frac{m}{2 \pi k T}\right)^{3 / 2} \overbrace{\mathrm{e}^{-m v_{x}^{2} / 2 k T} \mathrm{e}^{-m v_{y}^{2} / 2 k T} \mathrm{e}^{-m v_{z}^{2} / 2 k T}}^{\mathrm{e}^{-m\left(v_{2}^{2}+v_{2}^{2}+v_{2}^{2}\right) / 2 k T}} \\
& \times \mathrm{d} v_{x} \mathrm{~d} v_{y} \mathrm{~d} v_{z} \\
& =\left(\frac{m}{2 \pi k T}\right)^{3 / 2} \mathrm{e}^{-m v^{2} / 2 k T} \mathrm{~d} v_{x} \mathrm{~d} v_{y} \mathrm{~d} v_{z}
\end{aligned}
$$

where $v^{2}=v_{x}^{2}+v_{y}^{2}+v_{z}^{2}$.\\
Step 3 Calculate the probability that a molecule has a speed in the range $v$ to $v+\mathrm{d} v$\\
To evaluate the probability that a molecule has a speed in the range $v$ to $v+\mathrm{d} v$ regardless of direction, think of the three velocity components as defining three coordinates in 'velocity space', with the same properties as ordinary space except that the axes are labelled ( $v_{x}, v_{y}, v_{z}$ ) instead of $(x, y, z)$. Just as the volume element in ordinary space is $\mathrm{d} x \mathrm{~d} y \mathrm{~d} z$, so the volume element in velocity space is $\mathrm{d} v_{x} \mathrm{~d} v_{y} \mathrm{~d} v_{z}$. The sum of all the volume elements in ordinary space that lie at a distance $r$ from the centre is the volume of a spherical shell of radius $r$ and thickness $\mathrm{d} r$. That volume is the product of the surface area of the shell,\\
\includegraphics[max width=\textwidth, center]{2025_02_27_d4b95d9718f0b9e2e6b9g-046(1)}

Figure 1B. 3 To evaluate the probability that a molecule has a speed in the range $v$ to $v+\mathrm{d} v$, evaluate the total probability that the molecule will have a speed that is anywhere in a thin shell of radius $v=\left(v_{x}^{2}+v_{y}^{2}+v_{z}^{2}\right)^{1 / 2}$ and thickness $\mathrm{d} v$.\\
$4 \pi r^{2}$, and its thickness $\mathrm{d} r$, and is therefore $4 \pi r^{2} \mathrm{~d} r$. Similarly, the analogous volume in velocity space is the volume of a shell of radius $v$ and thickness $\mathrm{d} v$, namely $4 \pi \nu^{2} \mathrm{~d} v$ (Fig. 1B.3). Now, because $f\left(v_{x}\right) f\left(v_{y}\right) f\left(v_{z}\right)$, the term in blue in the last equation, depends only on $\nu^{2}$, and has the same value everywhere in a shell of radius $v$, the total probability of the molecules possessing a speed in the range $v$ to $v+\mathrm{d} v$ is the product of the term in blue and the volume of the shell of radius $v$ and thickness $\mathrm{d} v$. If this probability is written $f(v) \mathrm{d} v$, it follows that

$$
f(v) \mathrm{d} v=4 \pi v^{2} \mathrm{~d} v\left(\frac{m}{2 \pi k T}\right)^{3 / 2} \mathrm{e}^{-m v^{2} / 2 k T}
$$

\subsection*{The chemist's toolkit 4}
\subsection*{Integration}
Integration is concerned with the areas under curves. The integral of a function $f(x)$, which is denoted $\int f(x) \mathrm{d} x$ (the symbol $\int$ is an elongated S denoting a sum), between the two values $x=a$ and $x=b$ is defined by imagining the $x$-axis as divided into strips of width $\delta x$ and evaluating the following sum:

$$
\int_{a}^{b} f(x) \mathrm{d} x=\lim _{\delta x \rightarrow 0} \sum_{i} f\left(x_{i}\right) \delta x
$$

Integration [definition]

As can be appreciated from Sketch 1, the integral is the area under the curve between the limits $a$ and $b$. The function to be integrated is called the integrand. It is an astonishing mathematical fact that the integral of a function is the inverse of the differential of that function. In other words, if differentiation of $f$ is followed by integration of the resulting function, the result is the original function $f$ (to within a constant).\\
The integral in the preceding equation with the limits specified is called a definite integral. If it is written without the limits specified, it is called an indefinite integral. If the result of carrying out an indefinite integration is $g(x)+C$, where $C$ is a constant, the following procedure is used to evaluate the corresponding definite integral:\\
and $f(v)$ itself, after minor rearrangement, is

$$
f(v)=4 \pi\left(\frac{m}{2 \pi k T}\right)^{3 / 2} v^{2} \mathrm{e}^{-m v^{2} / 2 k T}
$$

Because $R=N_{\mathrm{A}} k$ (Table 1B.1), $m / k=m N_{\mathrm{A}} / R=M / R$, it follows that

$$
-\left|f(v)=4 \pi\left(\frac{M}{2 \pi R T}\right)^{3 / 2} v^{2} \mathrm{e}^{-M v^{2} / 2 R T}\right| \begin{aligned}
& \text { Maxwell-Boltzmann } \\
& \begin{array}{l}
\text { distribution } \\
{[\mathrm{KMT}]}
\end{array}
\end{aligned}
$$

The function $f(v)$ is called the Maxwell-Boltzmann distribution of speeds. Note that, in common with other distribution functions, $f(v)$ acquires physical significance only after it is multiplied by the range of speeds of interest.

Table 1B. 1 The (molar) gas constant*

\begin{center}
\begin{tabular}{ll}
\hline
$R$ &  \\
\hline
8.31447 & $\mathrm{~J} \mathrm{~K}^{-1} \mathrm{~mol}^{-1}$ \\
$8.20574 \times 10^{-2}$ & $\mathrm{dm}^{3} \mathrm{~atm} \mathrm{~K}^{-1} \mathrm{~mol}^{-1}$ \\
$8.31447 \times 10^{-2}$ & $\mathrm{dm}^{3} \mathrm{bar} \mathrm{K}^{-1} \mathrm{~mol}^{-1}$ \\
8.31447 & $\mathrm{~Pa} \mathrm{~m}^{3} \mathrm{~K}^{-1} \mathrm{~mol}^{-1}$ \\
62.364 & $\mathrm{dm}^{3} \mathrm{Torr} \mathrm{K}^{-1} \mathrm{~mol}^{-1}$ \\
1.98721 & $\mathrm{cal} \mathrm{K} \mathrm{mol}^{-1}$ \\
\hline
\end{tabular}
\end{center}

\begin{itemize}
  \item The gas constant is now defined as $R=N_{\mathrm{A}} k$, where $N_{\mathrm{A}}$ is Avogadro's constant and $k$ is Boltzmann's constant.
\end{itemize}

$$
\begin{aligned}
I=\int_{a}^{b} f(x) \mathrm{d} x & =\left.\{g(x)+C\}\right|_{a} ^{b}=\{g(b)+C\}-\{g(a)+C\} \\
& =g(b)-g(a)
\end{aligned}
$$

Note that the constant of integration disappears. The definite and indefinite integrals encountered in this text are listed in the Resource section. They may also be calculated by using mathematical software.\\
\includegraphics[max width=\textwidth, center]{2025_02_27_d4b95d9718f0b9e2e6b9g-046}

Sketch 1

The important features of the Maxwell-Boltzmann distribution are as follows (and are shown pictorially in Fig. 1B.4):

\begin{itemize}
  \item Equation 1B. 4 includes a decaying exponential function (more specifically, a Gaussian function). Its presence implies that the fraction of molecules with very high speeds is very small because $\mathrm{e}^{-x^{2}}$ becomes very small when $x$ is large.
  \item The factor $M / 2 R T$ multiplying $v^{2}$ in the exponent is large when the molar mass, $M$, is large, so the exponential factor goes most rapidly towards zero when $M$ is large. That is, heavy molecules are unlikely to be found with very high speeds.
  \item The opposite is true when the temperature, $T$, is high: then the factor $M / 2 R T$ in the exponent is small, so the exponential factor falls towards zero relatively slowly as $v$ increases. In other words, a greater fraction of the molecules can be expected to have high speeds at high temperatures than at low temperatures.
  \item A factor $v^{2}$ (the term before the e) multiplies the exponential. This factor goes to zero as $v$ goes to zero, so the fraction of molecules with very low speeds will also be very small whatever their mass.
  \item The remaining factors (the term in parentheses in eqn 1B. 4 and the $4 \pi$ ) simply ensure that, when the fractions are summed over the entire range of speeds from zero to infinity, the result is 1 .
\end{itemize}

\subsection*{(c) Mean values}
With the Maxwell-Boltzmann distribution in hand, it is possible to calculate the mean value of any power of the speed by evaluating the appropriate integral. For instance, to evaluate\\
\includegraphics[max width=\textwidth, center]{2025_02_27_d4b95d9718f0b9e2e6b9g-047}

Figure 1B. 4 The distribution of molecular speeds with temperature and molar mass. Note that the most probable speed (corresponding to the peak of the distribution) increases with temperature and with decreasing molar mass, and simultaneously the distribution becomes broader.\\
the fraction, $F$, of molecules with speeds in the range $v_{1}$ to $v_{2}$ evaluate the integral


\begin{equation*}
F\left(v_{1}, v_{2}\right)=\int_{v_{1}}^{v_{2}} f(v) \mathrm{d} v \tag{1B.5}
\end{equation*}


This integral is the area under the graph of $f$ as a function of $v$ and, except in special cases, has to be evaluated numerically by using mathematical software (Fig. 1B.5). The average value of $v^{n}$ is calculated as


\begin{equation*}
\left\langle v^{n}\right\rangle=\int_{0}^{\infty} v^{n} f(v) \mathrm{d} v \tag{1B.6}
\end{equation*}


In particular, integration with $n=2$ results in the mean square speed, $\left\langle v^{2}\right\rangle$, of the molecules at a temperature $T$ :

\[
\left\langle v^{2}\right\rangle=\frac{3 R T}{M} \quad \begin{align*}
& \text { Mean square speed }  \tag{1B.7}\\
& {[\text { KMMT] }}
\end{align*}
\]

It follows that the root-mean-square speed of the molecules of the gas is

\[
v_{\mathrm{rms}}=\left\langle v^{2}\right\rangle^{1 / 2}=\left(\frac{3 R T}{M}\right)^{1 / 2} \quad \begin{align*}
& \text { Root-mean-square speed }  \tag{1B.8}\\
& {[\mathrm{KMT}]}
\end{align*}
\]

which is proportional to the square root of the temperature and inversely proportional to the square root of the molar mass. That is, the higher the temperature, the higher the root-mean-square speed of the molecules, and, at a given temperature, heavy molecules travel more slowly than light molecules.

The important conclusion, however, is that when eqn 1B. 8 is substituted into eqn $1 B .2$, the result is $p V=n R T$, which is the equation of state of a perfect gas. This conclusion confirms that the kinetic model can be regarded as a model of a perfect gas.\\
\includegraphics[max width=\textwidth, center]{2025_02_27_d4b95d9718f0b9e2e6b9g-047(1)}

Figure 1B. 5 To calculate the probability that a molecule will have a speed in the range $v_{1}$ to $v_{2}$, integrate the distribution between those two limits; the integral is equal to the area under the curve between the limits, as shown shaded here.

\subsection*{Example 1B. 1 \\
 Calculating the mean speed of molecules}
 in a gasCalculate $v_{\text {rms }}$ and the mean speed, $v_{\text {mean }}$, of $\mathrm{N}_{2}$ molecules at $25^{\circ} \mathrm{C}$.

Collect your thoughts The root-mean-square speed is calculated from eqn $1 B .8$, with $M=28.02 \mathrm{~g} \mathrm{~mol}^{-1}$ (that is, $0.02802 \mathrm{~kg} \mathrm{~mol}^{-1}$ ) and $T=298 \mathrm{~K}$. The mean speed is obtained by evaluating the integral

$$
v_{\text {mean }}=\int_{0}^{\infty} v f(v) \mathrm{d} v
$$

with $f(v)$ given in eqn 1B.3. Use either mathematical software or the integrals listed in the Resource section and note that $1 \mathrm{~J}=1 \mathrm{~kg} \mathrm{~m}^{2} \mathrm{~s}^{-2}$.

The solution The root-mean-square speed is

$$
v_{\mathrm{rms}}=\left\{\frac{3 \times\left(8.3145 \mathrm{JK}^{-1} \mathrm{~mol}^{-1}\right) \times(298 \mathrm{~K})}{0.02802 \mathrm{~kg} \mathrm{~mol}^{-1}}\right\}^{1 / 2}=515 \mathrm{~ms}^{-1}
$$

The integral required for the calculation of $v_{\text {mean }}$ is

$$
\begin{aligned}
v_{\text {mean }} & =4 \pi\left(\frac{M}{2 \pi R T}\right)^{3 / 2} \overbrace{\int_{0}^{\infty} v^{3} \mathrm{e}^{-M v^{2} / 2 R T} \mathrm{~d} v}^{\text {Integral } 4} \\
& =4 \pi\left(\frac{M}{2 \pi R T}\right)^{3 / 2} \times \frac{1}{2}\left(\frac{2 R T}{M}\right)^{2}=\left(\frac{8 R T}{\pi M}\right)^{1 / 2}
\end{aligned}
$$

Substitution of the data then gives

$$
v_{\text {mean }}=\left(\frac{8 \times\left(8.3145 \mathrm{JK}^{-1} \mathrm{~mol}^{-1}\right) \times(298 \mathrm{~K})}{\pi \times\left(0.02802 \mathrm{~kg} \mathrm{~mol}^{-1}\right)}\right)^{1 / 2}=475 \mathrm{~m} \mathrm{~s}^{-1}
$$

Self-test 1B. 1 Confirm that eqn 1B. 7 follows from eqn 1B. 6 .

As shown in Example 1B.1, the Maxwell-Boltzmann distribution can be used to evaluate the mean speed, $v_{\text {mean }}$, of the molecules in a gas:


\begin{equation*}
v_{\text {mean }}=\left(\frac{8 R T}{\pi M}\right)^{1 / 2}=\left(\frac{8}{3 \pi}\right)^{1 / 2} v_{\text {rms }} \quad{\underset{[K M T]}{M e a n ~ s p e e d ~}}_{[ }^{[ } \tag{1B.9}
\end{equation*}


The most probable speed, $v_{\mathrm{mp}}$, can be identified from the location of the peak of the distribution by differentiating $f(v)$ with respect to $v$ and looking for the value of $v$ at which the derivative is zero (other than at $v=0$ and $v=\infty$; see Problem 1B.11):

\[
v_{\mathrm{mp}}=\left(\frac{2 R T}{M}\right)^{1 / 2}=\left(\frac{2}{3}\right)^{1 / 2} v_{\mathrm{rms}} \quad \begin{align*}
& \text { Most probable }  \tag{1B.10}\\
& \text { speed } \\
& {[\mathrm{KMT}]}
\end{align*}
\]

Figure 1B. 6 summarizes these results.\\
\includegraphics[max width=\textwidth, center]{2025_02_27_d4b95d9718f0b9e2e6b9g-048(1)}

Figure 1B. 6 A summary of the conclusions that can be deduced from the Maxwell distribution for molecules of molar mass $M$ at a temperature $T$ : $v_{\text {mp }}$ is the most probable speed, $v_{\text {mean }}$ is the mean speed, and $v_{\text {rms }}$ is the root-mean-square speed.

The mean relative speed, $v_{\text {rel }}$, the mean speed with which one molecule approaches another of the same kind, can also be calculated from the distribution:

\[
v_{\text {rel }}=2^{1 / 2} v_{\text {mean }} \quad \begin{align*}
& \text { Mean relative speed }  \tag{1B.11a}\\
& \text { [KMT, identical molecules] }]
\end{align*}
\]

This result is much harder to derive, but the diagram in Fig. 1B. 7 should help to show that it is plausible. For the relative mean speed of two dissimilar molecules of masses $m_{\mathrm{A}}$ and $m_{\mathrm{B}}$ :

\[
v_{\text {rel }}=\left(\frac{8 k T}{\pi \mu}\right)^{1 / 2} \quad \mu=\frac{m_{\mathrm{A}} m_{\mathrm{B}}}{m_{\mathrm{A}}+m_{\mathrm{B}}} \quad \begin{align*}
& \text { Mean relative }  \tag{1B.11b}\\
& \text { speed } \\
& \text { [perfect gas] }
\end{align*}
\]

\begin{center}
\includegraphics[max width=\textwidth]{2025_02_27_d4b95d9718f0b9e2e6b9g-048}
\end{center}

Figure 1B. 7 A simplified version of the argument to show that the mean relative speed of molecules in a gas is related to the mean speed. When the molecules are moving in the same direction, the mean relative speed is zero; it is $2 v$ when the molecules are approaching each other. A typical mean direction of approach is from the side, and the mean speed of approach is then $2^{1 / 2} v$. The last direction of approach is the most characteristic, so the mean speed of approach can be expected to be about $2^{1 / 2} v$. This value is confirmed by more detailed calculation.

\subsection*{Brief illustration 1B. 1}
As already seen (in Example 1B.1), the mean speed of $\mathrm{N}_{2}$ molecules at $25^{\circ} \mathrm{C}$ is $475 \mathrm{~m} \mathrm{~s}^{-1}$. It follows from eqn 1B.11a that their relative mean speed is

$$
v_{\mathrm{rel}}=2^{1 / 2} \times\left(475 \mathrm{~m} \mathrm{~s}^{-1}\right)=671 \mathrm{~ms}^{-1}
$$

\subsection*{18.2 Collisions}
The kinetic model can be used to develop the qualitative picture of a perfect gas, as a collection of ceaselessly moving, colliding molecules, into a quantitative, testable expression. In particular, it provides a way to calculate the average frequency with which molecular collisions occur and the average distance a molecule travels between collisions.

\subsection*{(a) The collision frequency}
Although the kinetic model assumes that the molecules are point-like, a 'hit' can be counted as occurring whenever the centres of two molecules come within a distance $d$ of each other, where $d$, the collision diameter, is of the order of the actual diameters of the molecules (for impenetrable hard spheres $d$ is the diameter). The kinetic model can be used to deduce the collision frequency, $z$, the number of collisions made by one molecule divided by the time interval during which the collisions are counted.

\subsection*{How is that done? 1B. 3 Using the kinetic model to derive}
 an expression for the collision frequencyConsider the positions of all the molecules except one to be frozen. Then note what happens as this one mobile molecule travels through the gas with a mean relative speed $v_{\text {rel }}$ for a time $\Delta t$. In doing so it sweeps out a 'collision tube' of crosssectional area $\sigma=\pi d^{2}$, length $v_{\text {rel }} \Delta t$ and therefore of volume $\sigma v_{\text {rel }} \Delta t$ (Fig. 1B.8). The number of stationary molecules with centres inside the collision tube is given by the volume $V$ of\\
\includegraphics[max width=\textwidth, center]{2025_02_27_d4b95d9718f0b9e2e6b9g-049(1)}

Figure 1B. 8 The basis of the calculation of the collision frequency in the kinetic theory of gases.

Table 1B. 2 Collision cross-sections*

\begin{center}
\begin{tabular}{ll}
\hline
 & $\sigma / \mathbf{n m}^{2}$ \\
\hline
$\mathrm{C}_{6} \mathrm{H}_{6}$ & 0.88 \\
$\mathrm{CO}_{2}$ & 0.52 \\
He & 0.21 \\
$\mathrm{~N}_{2}$ & 0.43 \\
\hline
\end{tabular}
\end{center}

\begin{itemize}
  \item More values are given in the Resource section.\\
the tube multiplied by the number density $\mathcal{N}=N / V$, where $N$ is the total number of molecules in the sample, and is $\mathcal{N} \sigma v_{\text {rel }} \Delta t$. The collision frequency $z$ is this number divided by $\Delta t$. It follows that\\
\includegraphics[max width=\textwidth, center]{2025_02_27_d4b95d9718f0b9e2e6b9g-049}
\end{itemize}

The parameter $\sigma$ is called the collision cross-section of the molecules. Some typical values are given in Table 1B.2.

An expression in terms of the pressure of the gas is obtained by using the perfect gas equation and $R=N_{\mathrm{A}} k$ to write the number density in terms of the pressure:

$$
\mathcal{N}=\frac{N}{V}=\frac{n N_{\mathrm{A}}}{V}=\frac{n N_{\mathrm{A}}}{n R T / p}=\frac{p N_{\mathrm{A}}}{R T}=\frac{p}{k T}
$$

Then


\begin{equation*}
z=\frac{\sigma v_{\mathrm{rel}} p}{k T} \quad \underset{\text { Collision frequency }}{ } \tag{1B.12b}
\end{equation*}


Equation 1B.12a shows that, at constant volume, the collision frequency increases with increasing temperature, because most molecules are moving faster. Equation 1B.12b shows that, at constant temperature, the collision frequency is proportional to the pressure. The greater the pressure, the greater the number density of molecules in the sample, and the rate at which they encounter one another is greater even though their average speed remains the same.

\subsection*{Brief illustration 1B. 2}
For an $\mathrm{N}_{2}$ molecule in a sample at $1.00 \mathrm{~atm}(101 \mathrm{kPa})$ and $25^{\circ} \mathrm{C}$, from Brief illustration 1B. $1 v_{\text {rel }}=671 \mathrm{~m} \mathrm{~s}^{-1}$. Therefore, from eqn 1B.12b, and taking $\sigma=0.45 \mathrm{~nm}^{2}$ (corresponding to $0.45 \times 10^{-18} \mathrm{~m}^{2}$ ) from Table 1B.2,

$$
\begin{aligned}
z & =\frac{\left(0.45 \times 10^{-18} \mathrm{~m}^{2}\right) \times\left(671 \mathrm{~m} \mathrm{~s}^{-1}\right) \times\left(1.01 \times 10^{5} \mathrm{~Pa}\right)}{\left(1.381 \times 10^{-23} \mathrm{JK}^{-1}\right) \times(298 \mathrm{~K})} \\
& =7.4 \times 10^{9} \mathrm{~s}^{-1}
\end{aligned}
$$

so a given molecule collides about $7 \times 10^{9}$ times each second. The timescale of events in gases is becoming clear.

\subsection*{(b) The mean free path}
The mean free path, $\lambda$ (lambda), is the average distance a molecule travels between collisions. If a molecule collides with a frequency $z$, it spends a time $1 / z$ in free flight between collisions, and therefore travels a distance $(1 / z) v_{\text {rel }}$. It follows that the mean free path is

\[
\lambda=\frac{v_{\mathrm{rel}}}{z} \quad \begin{align*}
& \text { Mean free path }  \tag{1B.13}\\
& {[\text { [KMT] }}
\end{align*}
\]

Substitution of the expression for $z$ from eqn 1B.12b gives

\[
\lambda=\frac{k T}{\sigma p} \quad l \begin{align*}
& \text { Mean free path }  \tag{1B.14}\\
& \text { [perfect gas] }
\end{align*}
\]

Doubling the pressure shortens the mean free path by a factor of 2 .

\subsection*{Brief illustration 1B. 3}
From Brief illustration 1B. $1 v_{\text {rel }}=671 \mathrm{~m} \mathrm{~s}^{-1}$ for $\mathrm{N}_{2}$ molecules at $25^{\circ} \mathrm{C}$, and from Brief illustration $1 \mathrm{~B} .2 z=7.4 \times 10^{9} \mathrm{~s}^{-1}$ when\\
the pressure is 1.00 atm . Under these circumstances, the mean free path of $\mathrm{N}_{2}$ molecules is

$$
\lambda=\frac{671 \mathrm{~m} \mathrm{~s}^{-1}}{7.4 \times 10^{9} \mathrm{~s}^{-1}}=9.1 \times 10^{-8} \mathrm{~m}
$$

or 91 nm , about $10^{3}$ molecular diameters.

Although the temperature appears in eqn 1B.14, in a sample of constant volume, the pressure is proportional to $T$, so $T / p$ remains constant when the temperature is increased. Therefore, the mean free path is independent of the temperature in a sample of gas provided the volume is constant. In a container of fixed volume the distance between collisions is determined by the number of molecules present in the given volume, not by the speed at which they travel.\\
In summary, a typical gas $\left(\mathrm{N}_{2}\right.$ or $\left.\mathrm{O}_{2}\right)$ at 1 atm and $25^{\circ} \mathrm{C}$ can be thought of as a collection of molecules travelling with a mean speed of about $500 \mathrm{~m} \mathrm{~s}^{-1}$. Each molecule makes a collision within about 1 ns , and between collisions it travels about $10^{3}$ molecular diameters.

\subsection*{Checklist of concepts}
$\square$ 1. The kinetic model of a gas considers only the contribution to the energy from the kinetic energies of the molecules.2. Important results from the model include expressions for the pressure and the root-mean-square speed.3. The Maxwell-Boltzmann distribution of speeds gives the fraction of molecules that have speeds in a specified range.\\
4. The collision frequency is the average number of collisions made by a molecule in an interval divided by the length of the interval.\\
5. The mean free path is the average distance a molecule travels between collisions.

\subsection*{Checklist of equations}
\begin{center}
\begin{tabular}{|c|c|c|c|}
\hline
Property & Equation & Comment & Equation number \\
\hline
Pressure of a perfect gas from the kinetic model & $p V=\frac{1}{3} n M \nu_{\text {rms }}^{2}$ & Kinetic model of a perfect gas & 1B. 2 \\
\hline
Maxwell-Boltzmann distribution of speeds & $f(v)=4 \pi(M / 2 \pi R T)^{3 / 2} v^{2} \mathrm{e}^{-M v^{2} / 2 R T}$ &  & 1B. 4 \\
\hline
Root-mean-square speed & $v_{\mathrm{rms}}=(3 R T / M)^{1 / 2}$ &  & 1B. 8 \\
\hline
Mean speed & $v_{\text {mean }}=(8 R T / \pi M)^{1 / 2}$ &  & 1B. 9 \\
\hline
Most probable speed & $v_{\text {mp }}=(2 R T / M)^{1 / 2}$ &  & 1B. 10 \\
\hline
Mean relative speed & \( \begin{aligned} v_{\mathrm{rel}} & =(8 k T / \pi \mu)^{1 / 2} \\ \mu & =m_{\mathrm{A}} m_{\mathrm{B}} /\left(m_{\mathrm{A}}+m_{\mathrm{B}}\right) \end{aligned} \) &  & 1B.11b \\
\hline
The collision frequency & $z=\sigma v_{\text {rel }} p / k T, \sigma=\pi d^{2}$ &  & 1B.12b \\
\hline
Mean free path & $\lambda=v_{\text {rel }} / z$ &  & 1B. 13 \\
\hline
\end{tabular}
\end{center}

\section*{TOPIC 1C Real gases}
\subsection*{> Why do you need to know this material?}
The properties of actual gases, so-called 'real gases', are different from those of a perfect gas. Moreover, the deviations from perfect behaviour give insight into the nature of the interactions between molecules.

\subsection*{What is the key idea?}
Attractions and repulsions between gas molecules account for modifications to the isotherms of a gas and account for critical behaviour.

\subsection*{What do you need to know already?}
This Topic builds on and extends the discussion of perfect gases in Topic 1A. The principal mathematical technique employed is the use of differentiation to identify a point of inflexion of a curve (The chemist's toolkit 5).\\
\includegraphics[max width=\textwidth, center]{2025_02_27_d4b95d9718f0b9e2e6b9g-051}

Figure 1C. 1 The dependence of the potential energy of two molecules on their internuclear separation. High positive potential energy (at very small separations) indicates that the interactions between them are strongly repulsive at these distances. At intermediate separations, attractive interactions dominate. At large separations (far to the right) the potential energy is zero and there is no interaction between the molecules.\\
ineffective when the molecules are far apart (well to the right in Fig. 1C.1). Intermolecular forces are also important when the temperature is so low that the molecules travel with such low mean speeds that they can be captured by one another.

The consequences of these interactions are shown by shapes of experimental isotherms (Fig. 1C.2). At low pressures, when the sample occupies a large volume, the molecules are so far apart for most of the time that the intermolecular forces play no significant role, and the gas behaves virtually perfectly. At moderate pressures, when the average separation of the molecules is only a few molecular diameters, the attractive forces dominate the repulsive forces. In this case, the gas can be expected to be more compressible than a perfect gas because the forces help to draw the molecules together. At high pressures, when the average separation of the molecules is small, the repulsive forces dominate and the gas can be expected to be less compressible because now the forces help to drive the molecules apart.

Consider what happens when a sample of gas initially in the state marked A in Fig. 1C.2b is compressed (its volume is reduced) at constant temperature by pushing in a piston. Near A, the pressure of the gas rises in approximate agreement with Boyle's law. Serious deviations from that law begin to appear when the volume has been reduced to $B$.

At C (which corresponds to about 60 atm for carbon dioxide), all similarity to perfect behaviour is lost, for suddenly the\\
\includegraphics[max width=\textwidth, center]{2025_02_27_d4b95d9718f0b9e2e6b9g-052(1)}

Figure 1C. 2 (a) Experimental isotherms of carbon dioxide at several temperatures. The 'critical isotherm', the isotherm at the critical temperature, is at $31.04^{\circ} \mathrm{C}$ (in blue). The critical point is marked with a star. (b) As explained in the text, the gas can condense only at and below the critical temperature as it is compressed along a horizontal line (such as CDE). The dotted black curve consists of points like $C$ and $E$ for all isotherms below the critical temperature.\\
piston slides in without any further rise in pressure: this stage is represented by the horizontal line CDE. Examination of the contents of the vessel shows that just to the left of $C$ a liquid appears, and there are two phases separated by a sharply defined surface. As the volume is decreased from $C$ through $D$ to $E$, the amount of liquid increases. There is no additional resistance to the piston because the gas can respond by condensing. The pressure corresponding to the line CDE, when both liquid and vapour are present in equilibrium, is called the vapour pressure of the liquid at the temperature of the experiment.

At E , the sample is entirely liquid and the piston rests on its surface. Any further reduction of volume requires the exertion of considerable pressure, as is indicated by the sharply rising line to the left of $E$. Even a small reduction of volume from $E$ to $F$ requires a great increase in pressure.

\subsection*{(a) The compression factor}
As a first step in understanding these observations it is useful to introduce the compression factor, $Z$, the ratio of the meas-\\
\includegraphics[max width=\textwidth, center]{2025_02_27_d4b95d9718f0b9e2e6b9g-052}

Figure 1C. 3 The variation of the compression factor, $Z$, with pressure for several gases at $0^{\circ} \mathrm{C}$. A perfect gas has $Z=1$ at all pressures. Notice that, although the curves approach 1 as $p \rightarrow 0$, they do so with different slopes.\\
ured molar volume of a gas, $V_{\mathrm{m}}=V / n$, to the molar volume of a perfect gas, $V_{\mathrm{m}}^{\circ}$, at the same pressure and temperature:

\[
Z=\frac{V_{\mathrm{m}}}{V_{\mathrm{m}}^{\mathrm{o}}} \quad l l \begin{align*}
& \text { Compression factor }  \tag{1C.1}\\
& \text { [definition] }
\end{align*}
\]

Because the molar volume of a perfect gas is equal to $R T / p$, an equivalent expression is $Z=p V_{\mathrm{m}} / R T$, which can be written as


\begin{equation*}
p V_{\mathrm{m}}=R T Z \tag{1C.2}
\end{equation*}


Because for a perfect gas $Z=1$ under all conditions, deviation of $Z$ from 1 is a measure of departure from perfect behaviour.

Some experimental values of $Z$ are plotted in Fig. 1C.3. At very low pressures, all the gases shown have $Z \approx 1$ and behave nearly perfectly. At high pressures, all the gases have $Z>1$, signifying that they have a larger molar volume than a perfect gas. Repulsive forces are now dominant. At intermediate pressures, most gases have $Z<1$, indicating that the attractive forces are reducing the molar volume relative to that of a perfect gas.

\subsection*{Brief illustration 1C. 1}
The molar volume of a perfect gas at 500 K and 100 bar is $V_{\mathrm{m}}^{\circ}=0.416 \mathrm{dm}^{3} \mathrm{~mol}^{-1}$. The molar volume of carbon dioxide under the same conditions is $V_{\mathrm{m}}=0.366 \mathrm{dm}^{3} \mathrm{~mol}^{-1}$. It follows that at 500 K

$$
Z=\frac{0.366 \mathrm{dm}^{3} \mathrm{~mol}^{-1}}{0.416 \mathrm{dm}^{3} \mathrm{~mol}^{-1}}=0.880
$$

The fact that $Z<1$ indicates that attractive forces dominate repulsive forces under these conditions.

\subsection*{(b) Virial coefficients}
At large molar volumes and high temperatures the real-gas isotherms do not differ greatly from perfect-gas isotherms.

Table 1C. 1 Second virial coefficients, $B /\left(\mathrm{cm}^{3} \mathrm{~mol}^{-1}\right)^{\star}$

\begin{center}
\begin{tabular}{lrc}
\hline
 & \multicolumn{2}{c}{Temperature} \\
\cline { 2 - 3 }
 & 273 K & $\mathbf{6 0 0} \mathrm{~K}$ \\
\hline
Ar & -21.7 & 11.9 \\
$\mathrm{CO}_{2}$ & -149.7 & -12.4 \\
$\mathrm{~N}_{2}$ & -10.5 & 21.7 \\
Xe & -153.7 & -19.6 \\
\hline
\end{tabular}
\end{center}

\begin{itemize}
  \item More values are given in the Resource section.
\end{itemize}

The small differences suggest that the perfect gas law $p V_{\mathrm{m}}=R T$ is in fact the first term in an expression of the form


\begin{equation*}
p V_{\mathrm{m}}=R T\left(1+B^{\prime} p+C^{\prime} p^{2}+\cdots\right) \tag{1C.3a}
\end{equation*}


This expression is an example of a common procedure in physical chemistry, in which a simple law that is known to be a good first approximation (in this case $p V_{\mathrm{m}}=R T$ ) is treated as the first term in a series in powers of a variable (in this case $p$ ). A more convenient expansion for many applications is

$$
p V_{\mathrm{m}}=R T\left(1+\frac{B}{V_{\mathrm{m}}}+\frac{C}{V_{\mathrm{m}}^{2}}+\cdots\right) \quad \text { Virial equation of state (1C.3b) }
$$

These two expressions are two versions of the virial equation of state. ${ }^{1}$ By comparing the expression with eqn 1C. 2 it is seen that the term in parentheses in eqn 1C.3b is just the compression factor, $Z$.

The coefficients $B, C, \ldots$, which depend on the temperature, are the second, third, ... virial coefficients (Table 1C.1); the first virial coefficient is 1 . The third virial coefficient, $C$, is usually less important than the second coefficient, $B$, in the sense that at typical molar volumes $C / V_{\mathrm{m}}^{2} \ll B / V_{\mathrm{m}}$. The values of the virial coefficients of a gas are determined from measurements of its compression factor.

\subsection*{Brief illustration 1C. 2}
To use eqn 1C.3b (up to the $B$ term) to calculate the pressure exerted at 100 K by $0.104 \mathrm{~mol}_{2}(\mathrm{~g})$ in a vessel of volume $0.225 \mathrm{dm}^{3}$, begin by calculating the molar volume:

$$
V_{\mathrm{m}}=\frac{V}{n_{\mathrm{O}_{2}}}=\frac{0.225 \mathrm{dm}^{3}}{0.104 \mathrm{~mol}}=2.16 \mathrm{dm}^{3} \mathrm{~mol}^{-1}=2.16 \times 10^{-3} \mathrm{~m}^{3} \mathrm{~mol}^{-1}
$$

Then, by using the value of $B$ found in Table 1C. 1 of the Resource section,

$$
p=\frac{R T}{V_{\mathrm{m}}}\left(1+\frac{B}{V_{\mathrm{m}}}\right)
$$

\footnotetext{${ }^{1}$ The name comes from the Latin word for force. The coefficients are sometimes denoted $B_{2}, B_{3}, \ldots$..
}
\$\$

\[
\begin{aligned}
& =\frac{\left(8.3145 \mathrm{Jmol}^{-1} \mathrm{~K}^{-1}\right) \times(100 \mathrm{~K})}{2.16 \times 10^{-3} \mathrm{~m}^{3} \mathrm{~mol}^{-1}}\left(1-\frac{1.975 \times 10^{-4} \mathrm{~m}^{3} \mathrm{~mol}^{-1}}{2.16 \times 10^{-3} \mathrm{~m}^{3} \mathrm{~mol}^{-1}}\right) \\
& =3.50 \times 10^{5} \mathrm{~Pa}, \text { or } 350 \mathrm{kPa}
\end{aligned}
\]

\$\$\\
where $1 \mathrm{~Pa}=1 \mathrm{Jm}^{-3}$. The perfect gas equation of state would give the calculated pressure as 385 kPa , or 10 per cent higher than the value calculated by using the virial equation of state. The difference is significant because under these conditions $B / V_{\mathrm{m}} \approx 0.1$ which is not negligible relative to 1 .

An important point is that although the equation of state of a real gas may coincide with the perfect gas law as $p \rightarrow 0$, not all its properties necessarily coincide with those of a perfect gas in that limit. Consider, for example, the value of $d Z / d p$, the slope of the graph of compression factor against pressure (see The chemist's toolkit 5 for a review of derivatives and differentiation). For a perfect gas $\mathrm{d} Z / \mathrm{d} p=0$ (because $Z=1$ at all pressures), but for a real gas from eqn 1C.3a


\begin{equation*}
\frac{\mathrm{d} Z}{\mathrm{~d} p}=B^{\prime}+2 p C^{\prime}+\cdots \rightarrow B^{\prime} \text { as } p \rightarrow 0 \tag{1C.4a}
\end{equation*}


However, $B^{\prime}$ is not necessarily zero, so the slope of $Z$ with respect to $p$ does not necessarily approach 0 (the perfect gas value), as can be seen in Fig. 1C.4. By a similar argument (see The chemist's toolkit 5 for evaluating derivatives of this kind),


\begin{equation*}
\frac{\mathrm{d} Z}{\mathrm{~d}\left(1 / V_{\mathrm{m}}\right)} \rightarrow B \text { as } V_{\mathrm{m}} \rightarrow \infty \tag{1C.4b}
\end{equation*}


Because the virial coefficients depend on the temperature, there may be a temperature at which $Z \rightarrow 1$ with zero slope at low pressure or high molar volume (as in Fig. 1C.4). At this temperature, which is called the Boyle temperature, $T_{B}$, the properties of the real gas do coincide with those of a per-\\
\includegraphics[max width=\textwidth, center]{2025_02_27_d4b95d9718f0b9e2e6b9g-053}

Figure 1C. 4 The compression factor, $Z$, approaches 1 at low pressures, but does so with different slopes. For a perfect gas, the slope is zero, but real gases may have either positive or negative slopes, and the slope may vary with temperature. At the Boyle temperature, the slope is zero at $p=0$ and the gas behaves perfectly over a wider range of conditions than at other temperatures.

\subsection*{The chemist's toolkit 5}
\subsection*{Differentiation}
Differentiation is concerned with the slopes of functions, such as the rate of change of a variable with time. The formal definition of the derivative, $\mathrm{d} f / \mathrm{d} x$, of a function $f(x)$ is

$$
\frac{\mathrm{d} f}{\mathrm{~d} x}=\lim _{\delta x \rightarrow 0} \frac{f(x+\delta x)-f(x)}{\delta x}
$$

First derivative [definition]

As shown in Sketch 1, the derivative can be interpreted as the slope of the tangent to the graph of $f(x)$ at a given value of $x$. A positive first derivative indicates that the function slopes upwards (as $x$ increases), and a negative first derivative indicates the opposite. It is sometimes convenient to denote the first derivative as $f^{\prime}(x)$. The second derivative, $\mathrm{d}^{2} f / \mathrm{d} x^{2}$, of a function is the derivative of the first derivative (here denoted $f^{\prime}$ ):

$$
\frac{\mathrm{d}^{2} f}{\mathrm{~d} x^{2}}=\lim _{\delta x \rightarrow 0} \frac{f^{\prime}(x+\delta x)-f^{\prime}(x)}{\delta x}
$$

Second derivative [definition]

It is sometimes convenient to denote the second derivative $f^{\prime \prime}$. As shown in Sketch 2, the second derivative of a function can be interpreted as an indication of the sharpness of the curvature of the function. A positive second derivative indicates that the function is $\cup$ shaped, and a negative second derivative indicates that it is $\cap$ shaped. The second derivative is zero at a point of inflection, where the first derivative changes sign.\\
The derivatives of some common functions are as follows:

$$
\begin{aligned}
& \frac{\mathrm{d}}{\mathrm{~d} x} x^{n}=n x^{n-1} \\
& \frac{\mathrm{~d}}{\mathrm{~d} x} \mathrm{e}^{a x}=a \mathrm{e}^{a x}
\end{aligned}
$$

\begin{center}
\includegraphics[max width=\textwidth]{2025_02_27_d4b95d9718f0b9e2e6b9g-054(1)}
\end{center}

Sketch 1\\
fect gas as $p \rightarrow 0$. According to eqn 1C.4a, $Z$ has zero slope as $p \rightarrow 0$ if $B^{\prime}=0$, so at the Boyle temperature $B^{\prime}=0$. It then follows from eqn 1 C .3 a that $p V_{\mathrm{m}} \approx R T_{\mathrm{B}}$ over a more extended range of pressures than at other temperatures because the first term after 1 (i.e. $B^{\prime} p$ ) in the virial equation is zero and $C^{\prime} p^{2}$ and higher terms are negligibly small. For helium $T_{\mathrm{B}}=22.64 \mathrm{~K}$; for air $T_{\mathrm{B}}=346.8 \mathrm{~K}$; more values are given in Table 1C.2.\\
\includegraphics[max width=\textwidth, center]{2025_02_27_d4b95d9718f0b9e2e6b9g-054}

It follows from the definition of the derivative that a variety of combinations of functions can be differentiated by using the following rules:

$$
\begin{aligned}
& \frac{\mathrm{d}}{\mathrm{~d} x}(u+v)=\frac{\mathrm{d} u}{\mathrm{~d} x}+\frac{\mathrm{d} v}{\mathrm{~d} x} \\
& \frac{\mathrm{~d}}{\mathrm{~d} x} u v=u \frac{\mathrm{~d} v}{\mathrm{~d} x}+v \frac{\mathrm{~d} u}{\mathrm{~d} x} \\
& \frac{\mathrm{~d}}{\mathrm{~d} x} \frac{u}{v}=\frac{1}{v} \frac{\mathrm{~d} u}{\mathrm{~d} x}-\frac{u}{v^{2}} \frac{\mathrm{~d} v}{\mathrm{~d} x}
\end{aligned}
$$

It is sometimes convenient to differentiate with respect to a function of $x$, rather than $x$ itself. For instance, suppose that

$$
f(x)=a+\frac{b}{x}+\frac{c}{x^{2}}
$$

where $a, b$, and $c$ are constants and you need to evaluate $\mathrm{d} f / \mathrm{d}(1 / x)$, rather than $\mathrm{d} f / \mathrm{d} x$. To begin, let $y=1 / x$. Then $f(y)=$ $a+b y+c y^{2}$ and

$$
\frac{\mathrm{d} f}{\mathrm{~d} y}=b+2 c y
$$

Because $y=1 / x$, it follows that

$$
\frac{\mathrm{d} f}{\mathrm{~d}(1 / x)}=b+\frac{2 c}{x}
$$

\subsection*{(c) Critical constants}
There is a temperature, called the critical temperature, $T_{c}$, which separates two regions of behaviour and plays a special role in the theory of the states of matter. An isotherm slightly below $T_{\mathrm{c}}$ behaves as already described: at a certain pressure, a liquid condenses from the gas and is distinguishable from it by

Table 1C. 2 Critical constants of gases*

\begin{center}
\begin{tabular}{lccrcc}
\hline
 & $p_{\mathrm{c}} / \mathrm{atm}$ & $V_{\mathrm{c}} /\left(\mathrm{cm}^{3} \mathrm{~mol}^{-1}\right)$ & $T_{c} / \mathrm{K}$ & $Z_{\mathrm{c}}$ & $T_{\mathrm{B}} / \mathrm{K}$ \\
\hline
Ar & 48.0 & 75.3 & 150.7 & 0.292 & 411.5 \\
$\mathrm{CO}_{2}$ & 72.9 & 94.0 & 304.2 & 0.274 & 714.8 \\
He & 2.26 & 57.8 & 5.2 & 0.305 & 22.64 \\
$\mathrm{O}_{2}$ & 50.14 & 78.0 & 154.8 & 0.308 & 405.9 \\
\hline
\end{tabular}
\end{center}

${ }^{*}$ More values are given in the Resource section.\\
the presence of a visible surface. If, however, the compression takes place at $T_{\mathrm{c}}$ itself, then a surface separating two phases does not appear and the volumes at each end of the horizontal part of the isotherm have merged to a single point, the critical point of the gas. The pressure and molar volume at the critical point are called the critical pressure, $p_{c}$, and critical molar volume, $V_{c}$, of the substance. Collectively, $p_{c}, V_{c}$, and $T_{c}$ are the critical constants of a substance (Table 1C.2).

At and above $T_{c}$, the sample has a single phase which occupies the entire volume of the container. Such a phase is, by definition, a gas. Hence, the liquid phase of a substance does not form above the critical temperature. The single phase that fills the entire volume when $T>T_{\mathrm{c}}$ may be much denser than considered typical for gases, and the name supercritical fluid is preferred.

\subsection*{Brief illustration 1C. 3}
The critical temperature of oxygen, 155 K , signifies that it is impossible to produce liquid oxygen by compression alone if its temperature is greater than 155 K . To liquefy oxygen the temperature must first be lowered to below 155 K , and then the gas compressed isothermally.

\subsection*{1c. 2 The van der Waals equation}
Conclusions may be drawn from the virial equations of state only by inserting specific values of the coefficients. It is often useful to have a broader, if less precise, view of all gases, such as that provided by an approximate equation of state.

\subsection*{(a) Formulation of the equation}
The equation introduced by J.D. van der Waals in 1873 is an excellent example of an expression that can be obtained by thinking scientifically about a mathematically complicated but physically simple problem; that is, it is a good example of 'model building'.

How is that done? 1C. 1 Deriving the van der Waals equation of state

The repulsive interaction between molecules is taken into account by supposing that it causes the molecules to behave as small but impenetrable spheres, so instead of moving in a volume $V$ they are restricted to a smaller volume $V-n b$, where $n b$ is approximately the total volume taken up by the molecules themselves. This argument suggests that the perfect gas law $p=n R T / V$ should be replaced by

$$
p=\frac{n R T}{V-n b}
$$

when repulsions are significant. To calculate the excluded volume, note that the closest distance of approach of two hardsphere molecules of radius $r$ (and volume $V_{\text {molecule }}=\frac{4}{3} \pi r^{3}$ ) is $2 r$, so the volume excluded is $\frac{4}{3} \pi(2 r)^{3}$, or $8 V_{\text {molecule }}$. The volume excluded per molecule is one-half this volume, or $4 V_{\text {molecule }}$, so $b \approx 4 V_{\text {molecule }} N_{\mathrm{A}}$.

The pressure depends on both the frequency of collisions with the walls and the force of each collision. Both the frequency of the collisions and their force are reduced by the attractive interaction, which acts with a strength proportional to the number of interacting molecules and therefore to the molar concentration, $n / V$, of molecules in the sample. Because both the frequency and the force of the collisions are reduced by the attractive interactions, the pressure is reduced in proportion to the square of this concentration. If the reduction of pressure is written as $a(n / V)^{2}$, where $a$ is a positive constant characteristic of each gas, the combined effect of the repulsive and attractive forces is the van der Waals equation:\\
$\left.-p=\frac{n R T}{V-n b}-a \frac{n^{2}}{V^{2}} \right\rvert\, \quad$ (1C.5a)\\
The constants $a$ and $b$ are called the van der Waals coefficients, with $a$ representing the strength of attractive interactions and $b$ that of the repulsive interactions between the molecules. They are characteristic of each gas and taken to be independent of the temperature (Table 1C.3). Although $a$ and $b$ are not precisely defined molecular properties, they correlate with physical properties that reflect the strength of intermolecular interactions, such as critical temperature, vapour pressure, and enthalpy of vaporization.

Table 1C. 3 van der Waals coefficients*

\begin{center}
\begin{tabular}{lll}
\hline
 & $\boldsymbol{a} /\left(\mathbf{a t m ~ d m}^{6} \mathrm{~mol}^{-2}\right)$ & $\boldsymbol{b} /\left(\mathbf{1 0}^{-\mathbf{2}} \mathrm{dm}^{3} \mathrm{~mol}^{-1}\right)$ \\
\hline
Ar & 1.337 & 3.20 \\
$\mathrm{CO}_{2}$ & 3.610 & 4.29 \\
He & 0.0341 & 2.38 \\
Xe & 4.137 & 5.16 \\
\hline
\end{tabular}
\end{center}

\footnotetext{\begin{itemize}
  \item More values are given in the Resource section.
\end{itemize}
}\subsection*{Brief illustration 1C. 4}
For benzene $a=18.57 \mathrm{~atm} \mathrm{dm}^{6} \mathrm{~mol}^{-2}\left(1.882 \mathrm{Pam}^{6} \mathrm{~mol}^{-2}\right)$ and $b=0.1193 \mathrm{dm}^{3} \mathrm{~mol}^{-1}\left(1.193 \times 10^{-4} \mathrm{~m}^{3} \mathrm{~mol}^{-1}\right)$; its normal boiling point is 353 K . Treated as a perfect gas at $T=400 \mathrm{~K}$ and $p=1.0 \mathrm{~atm}$, benzene vapour has a molar volume of $V_{\mathrm{m}}=$ $R T / p=33 \mathrm{dm}^{3} \mathrm{~mol}^{-1}$, so the criterion $V_{\mathrm{m}} \gg b$ for perfect gas behaviour is satisfied. It follows that $a / V_{\mathrm{m}}^{2} \approx 0.017 \mathrm{~atm}$, which is 1.7 per cent of 1.0 atm . Therefore, benzene vapour is expected to deviate only slightly from perfect gas behaviour at this temperature and pressure.

Equation 1C.5a is often written in terms of the molar volume $V_{\mathrm{m}}=V / n$ as


\begin{equation*}
p=\frac{R T}{V_{\mathrm{m}}-b}-\frac{a}{V_{\mathrm{m}}^{2}} \tag{1C.5b}
\end{equation*}


\subsection*{Example 1C. 1 Using the van der Waals equation to estimate a molar volume}
Estimate the molar volume of $\mathrm{CO}_{2}$ at 500 K and 100 atm by treating it as a van der Waals gas.

Collect your thoughts You need to find an expression for the molar volume by solving the van der Waals equation, eqn 1 C .5 b . To rearrange the equation into a suitable form, multiply both sides by $\left(V_{\mathrm{m}}-b\right) V_{\mathrm{m}}^{2}$, to obtain

$$
\left(V_{\mathrm{m}}-b\right) V_{\mathrm{m}}^{2} p=R T V_{\mathrm{m}}^{2}-\left(V_{\mathrm{m}}-b\right) a
$$

Then, after division by $p$, collect powers of $V_{\mathrm{m}}$ to obtain

$$
V_{\mathrm{m}}^{3}-\left(b+\frac{R T}{p}\right) V_{\mathrm{m}}^{2}+\left(\frac{a}{p}\right) V_{\mathrm{m}}+\frac{a b}{p}=0
$$

Although closed expressions for the roots of a cubic equation can be given, they are very complicated. Unless analytical solutions are essential, it is usually best to solve such equations with mathematical software; graphing calculators can also be used to help identify the acceptable root.

The solution According to Table 1C.3, $a=3.592 \mathrm{dm}^{6} \mathrm{~atm} \mathrm{~mol}^{-2}$ and $b=4.267 \times 10^{-2} \mathrm{dm}^{3} \mathrm{~mol}^{-1}$. Under the stated conditions, $R T / p=0.410 \mathrm{dm}^{3} \mathrm{~mol}^{-1}$. The coefficients in the equation for $V_{\mathrm{m}}$ are therefore

$$
\begin{aligned}
b+R T / p & =0.453 \mathrm{dm}^{3} \mathrm{~mol}^{-1} \\
a / p & =3.61 \times 10^{-2}\left(\mathrm{dm}^{3} \mathrm{~mol}^{-1}\right)^{2} \\
a b / p & =1.55 \times 10^{-3}\left(\mathrm{dm}^{3} \mathrm{~mol}^{-1}\right)^{3}
\end{aligned}
$$

Therefore, on writing $x=V_{\mathrm{m}} /\left(\mathrm{dm}^{3} \mathrm{~mol}^{-1}\right)$, the equation to solve is

$$
x^{3}-0.453 x^{2}+\left(3.61 \times 10^{-2}\right) x-\left(1.55 \times 10^{-3}\right)=0
$$

\begin{center}
\includegraphics[max width=\textwidth]{2025_02_27_d4b95d9718f0b9e2e6b9g-056(1)}
\end{center}

Figure 1C. 5 The graphical solution of the cubic equation for $V$ in Example 1C.1.

The acceptable root is $x=0.366$ (Fig. 1C.5), which implies that $V_{\mathrm{m}}=0.366 \mathrm{dm}^{3} \mathrm{~mol}^{-1}$. The molar volume of a perfect gas under these conditions is $0.410 \mathrm{dm}^{3} \mathrm{~mol}^{-1}$.

Self-test 1C. 1 Calculate the molar volume of argon at $100^{\circ} \mathrm{C}$ and 100 atm on the assumption that it is a van der Waals gas.\\
\includegraphics[max width=\textwidth, center]{2025_02_27_d4b95d9718f0b9e2e6b9g-056}

\subsection*{(b) The features of the equation}
To what extent does the van der Waals equation predict the behaviour of real gases? It is too optimistic to expect a single, simple expression to be the true equation of state of all substances, and accurate work on gases must resort to the virial equation, use tabulated values of the coefficients at various temperatures, and analyse the system numerically. The advantage of the van der Waals equation, however, is that it is analytical (that is, expressed symbolically) and allows some general conclusions about real gases to be drawn. When the equation fails another equation of state must be used (some are listed in Table 1C.4), yet another must be invented, or the virial equation is used.\\
The reliability of the equation can be judged by comparing the isotherms it predicts with the experimental isotherms in Fig. 1C.2. Some calculated isotherms are shown in Figs. 1C. 6 and 1C.7. Apart from the oscillations below the critical temperature, they do resemble experimental isotherms quite well. The oscillations, the van der Waals loops, are unrealistic because they suggest that under some conditions an increase of pressure results in an increase of volume. Therefore they are replaced by horizontal lines drawn so the loops define equal areas above and below the lines: this procedure is called the Maxwell construction (1). The van der Waals coefficients, such as those in Table 1C.3, are found by fitting the calculated curves to the experimental curves.\\
\includegraphics[max width=\textwidth, center]{2025_02_27_d4b95d9718f0b9e2e6b9g-056(2)}

Table 1C. 4 Selected equations of state\\
\includegraphics[max width=\textwidth, center]{2025_02_27_d4b95d9718f0b9e2e6b9g-057(2)}

\begin{itemize}
  \item Reduced variables are defined as $X_{\mathrm{r}}=X / X_{\mathrm{c}}$ with $X=p, V_{\mathrm{m}}$, and $T$. Equations of state are sometimes expressed in terms of the molar volume, $V_{\mathrm{m}}=V / n$.\\
\includegraphics[max width=\textwidth, center]{2025_02_27_d4b95d9718f0b9e2e6b9g-057(1)}
\end{itemize}

Figure 1C. 6 The surface of possible states allowed by the van der Waals equation. The curves drawn on the surface are isotherms, labelled with the value of $T / T_{c^{\prime}}$ and correspond to the isotherms in Fig. 1C.7.\\
\includegraphics[max width=\textwidth, center]{2025_02_27_d4b95d9718f0b9e2e6b9g-057}

Figure 1C. 7 Van der Waals isotherms at several values of $T / T_{c}$. The van der Waals loops are normally replaced by horizontal straight lines. The critical isotherm is the isotherm for $T / T_{c}=1$, and is shown in blue.

The principal features of the van der Waals equation can be summarized as follows.

\begin{enumerate}
  \item Perfect gas isotherms are obtained at high temperatures and large molar volumes.
\end{enumerate}

When the temperature is high, $R T$ may be so large that the first term in eqn 1C.5b greatly exceeds the second. Furthermore, if the molar volume is large in the sense $V_{\mathrm{m}} \gg b$, then the denominator $V_{\mathrm{m}}-b \approx V_{\mathrm{m}}$. Under these conditions, the equation reduces to $p=R T / V_{\mathrm{m}}$, the perfect gas equation.\\
2. Liquids and gases coexist when the attractive and repulsive effects are in balance.

The van der Waals loops occur when both terms in eqn 1C.5b have similar magnitudes. The first term arises from the kinetic energy of the molecules and their repulsive interactions; the second represents the effect of the attractive interactions.\\
3. The critical constants are related to the van der Waals coefficients.

For $T<T_{c}$, the calculated isotherms oscillate, and each one passes through a minimum followed by a maximum. These extrema converge as $T \rightarrow T_{c}$ and coincide at $T=T_{c}$; at the critical point the curve has a flat\\
\includegraphics[max width=\textwidth]{2025_02_27_d4b95d9718f0b9e2e6b9g-057(3)} inflexion (2). From the properties of curves, an inflexion of this type occurs when both the first and second derivatives are zero. Hence, the critical constants can be found by calculating these derivatives and setting them equal to zero at the critical point:

$$
\frac{\mathrm{d} p}{\mathrm{~d} V_{\mathrm{m}}}=-\frac{R T}{\left(V_{\mathrm{m}}-b\right)^{2}}+\frac{2 a}{V_{\mathrm{m}}^{3}}=0
$$

$$
\frac{\mathrm{d}^{2} p}{\mathrm{~d} V_{\mathrm{m}}^{2}}=\frac{2 R T}{\left(V_{\mathrm{m}}-b\right)^{3}}-\frac{6 a}{V_{\mathrm{m}}^{4}}=0
$$

The solutions of these two equations (and using eqn 1C.5b to calculate $p_{c}$ from $V_{c}$ and $T_{c}$; see Problem 1C.12) are


\begin{equation*}
V_{\mathrm{c}}=3 b \quad p_{\mathrm{c}}=\frac{a}{27 b^{2}} \quad T_{\mathrm{c}}=\frac{8 a}{27 b R} \tag{1C.6}
\end{equation*}


These relations provide an alternative route to the determination of $a$ and $b$ from the values of the critical constants. They can be tested by noting that the critical compression factor, $Z_{c}$, is predicted to be


\begin{equation*}
Z_{\mathrm{c}}=\frac{p_{\mathrm{c}} V_{\mathrm{c}}}{R T_{\mathrm{c}}}=\frac{3}{8} \tag{1C.7}
\end{equation*}


for all gases that are described by the van der Waals equation near the critical point. Table 1C. 2 shows that although $Z_{c}<\frac{3}{8}$ $=0.375$, it is approximately constant (at 0.3 ) and the discrepancy is reasonably small.

\subsection*{(c) The principle of corresponding states}
An important general technique in science for comparing the properties of objects is to choose a related fundamental property of the same kind and to set up a relative scale on that basis. The critical constants are characteristic properties of gases, so it may be that a scale can be set up by using them as yardsticks and to introduce the dimensionless reduced variables of a gas by dividing the actual variable by the corresponding critical constant:

$$
V_{\mathrm{r}}=\frac{V_{\mathrm{m}}}{V_{\mathrm{c}}} \quad p_{\mathrm{r}}=\frac{p}{p_{\mathrm{c}}} \quad T_{\mathrm{r}}=\frac{T}{T_{\mathrm{c}}} \quad \begin{aligned}
& \text { Reduced variables } \\
& \text { [definition] }
\end{aligned}
$$

(1C.8)

If the reduced pressure of a gas is given, its actual pressure is calculated by using $p=p_{\mathrm{r}} p_{\mathrm{c}}$, and likewise for the volume and temperature. Van der Waals, who first tried this procedure, hoped that gases confined to the same reduced volume, $V_{\mathrm{r}}$, at the same reduced temperature, $T_{\mathrm{r}}$, would exert the same reduced pressure, $p_{r}$. The hope was largely fulfilled (Fig. 1C.8). The illustration shows the dependence of the compression factor on the reduced pressure for a variety of gases at various reduced temperatures. The success of the procedure is strikingly clear: compare this graph with Fig. 1C.3, where similar data are plotted without using reduced variables.\\
The observation that real gases at the same reduced volume and reduced temperature exert the same reduced pressure is called the principle of corresponding states. The principle is only an approximation. It works best for gases composed of spherical molecules; it fails, sometimes badly, when the molecules are non-spherical or polar.\\
\includegraphics[max width=\textwidth, center]{2025_02_27_d4b95d9718f0b9e2e6b9g-058}

Figure 1C. 8 The compression factors of four of the gases shown in Fig. 1C. 3 plotted using reduced variables. The curves are labelled with the reduced temperature $T_{\mathrm{r}}=T / T_{\mathrm{c}}$. The use of reduced variables organizes the data on to single curves.

\subsection*{Brief illustration 1C. 5}
The critical constants of argon and carbon dioxide are given in Table 1C.2. Suppose argon is at 23 atm and 200 K , its reduced pressure and temperature are then

$$
p_{\mathrm{r}}=\frac{23 \mathrm{~atm}}{48.0 \mathrm{~atm}}=0.48 \quad T_{\mathrm{r}}=\frac{200 \mathrm{~K}}{150.7 \mathrm{~K}}=1.33
$$

For carbon dioxide to be in a corresponding state, its pressure and temperature would need to be

$$
p=0.48 \times(72.9 \mathrm{~atm})=35 \mathrm{~atm} \quad T=1.33 \times 304.2 \mathrm{~K}=405 \mathrm{~K}
$$

The van der Waals equation sheds some light on the principle. When eqn 1C.5b is expressed in terms of the reduced variables it becomes

$$
p_{\mathrm{r}} p_{\mathrm{c}}=\frac{R T_{\mathrm{r}} T_{\mathrm{c}}}{V_{\mathrm{r}} V_{\mathrm{c}}-b}-\frac{a}{V_{\mathrm{r}}^{2} V_{\mathrm{c}}^{2}}
$$

Now express the critical constants in terms of $a$ and $b$ by using eqn 1C.6:

$$
\frac{a p_{\mathrm{r}}}{27 b^{2}}=\frac{8 a T_{\mathrm{r}} / 27 b}{3 b V_{\mathrm{r}}-b}-\frac{a}{9 b^{2} V_{\mathrm{r}}^{2}}
$$

and, after multiplying both sides by $27 b^{2} / a$, reorganize it into


\begin{equation*}
p_{\mathrm{r}}=\frac{8 T_{\mathrm{r}}}{3 V_{\mathrm{r}}-1}-\frac{3}{V_{\mathrm{r}}^{2}} \tag{1C.9}
\end{equation*}


This equation has the same form as the original, but the coefficients $a$ and $b$, which differ from gas to gas, have disappeared. It follows that if the isotherms are plotted in terms of the reduced variables (as done in fact in Fig. 1C. 7 without drawing attention to the fact), then the same curves are obtained whatever the gas. This is precisely the content of the principle of corresponding states, so the van der Waals equation is compatible with it.

Looking for too much significance in this apparent triumph is mistaken, because other equations of state also accommodate the principle. In fact, any equation of state (such as those in Table 1C.4) with two parameters playing the roles of $a$ and $b$ can be manipulated into a reduced form. The observation that real gases obey the principle approximately amounts to\\
saying that the effects of the attractive and repulsive interactions can each be approximated in terms of a single parameter. The importance of the principle is then not so much its theoretical interpretation but the way that it enables the properties of a range of gases to be coordinated on to a single diagram (e.g. Fig. 1C. 8 instead of Fig. 1C.3).

\subsection*{Checklist of concepts}
1. The extent of deviations from perfect behaviour is summarized by introducing the compression factor.2. The virial equation is an empirical extension of the perfect gas equation that summarizes the behaviour of real gases over a range of conditions.3. The isotherms of a real gas introduce the concept of critical behaviour.4. A gas can be liquefied by pressure alone only if its temperature is at or below its critical temperature.$\square$ 5. The van der Waals equation is a model equation of state for a real gas expressed in terms of two parameters, one (a) representing molecular attractions and the other (b) representing molecular repulsions.\\
$\square$ 6. The van der Waals equation captures the general features of the behaviour of real gases, including their critical behaviour.\\
7. The properties of real gases are coordinated by expressing their equations of state in terms of reduced variables.

\subsection*{Checklist of equations}
\begin{center}
\begin{tabular}{llll}
\hline
Property & Equation & Comment & \begin{tabular}{c}
Equation \\
number \\
\end{tabular} \\
\hline
Compression factor & $Z=V_{\mathrm{m}} / V_{\mathrm{m}}^{\circ}$ & Definition & 1 C .1 \\
Virial equation of state & $p V_{\mathrm{m}}=R T\left(1+B / V_{\mathrm{m}}+C / V_{\mathrm{m}}^{2}+\cdots\right)$ & $B, C$ depend on temperature & $a$ parameterizes attractions, \\
van der Waals equation of state & $p=n R T /(V-n b)-a(n / V)^{2}$ & $b$ parameterizes repulsions & 1 C .3 b \\
Reduced variables & $X_{\mathrm{r}}=X / X_{\mathrm{c}}$ & $X=p, V_{\mathrm{m}}$, or $T$ & 1 C .5 a \\
\hline
\end{tabular}
\end{center}


\chapter{FOCUS 2 The First Law}
The release of energy can be used to provide heat when a fuel burns in a furnace, to produce mechanical work when a fuel burns in an engine, and to generate electrical work when a chemical reaction pumps electrons through a circuit. Chemical reactions can be harnessed to provide heat and work, liberate energy that is unused but which gives desired products, and drive the processes of life. Thermodynamics, the study of the transformations of energy, enables the discussion of all these matters quantitatively, allowing for useful predictions.

\subsection*{2A Internal energy}
This Topic examines the ways in which a system can exchange energy with its surroundings in terms of the work it may do or have done on it, or the heat that it may produce or absorb. These considerations lead to the definition of the 'internal energy', the total energy of a system, and the formulation of the 'First Law' of thermodynamics, which states that the internal energy of an isolated system is constant.\\
2A. 1 Work, heat, and energy; 2A. 2 The definition of internal energy;\\
2A. 3 Expansion work; 2A. 4 Heat transactions

\subsection*{2B Enthalpy}
The second major concept of the Focus is 'enthalpy', which is a very useful book-keeping property for keeping track of the heat output (or requirements) of physical processes and chemical reactions that take place at constant pressure. Experimentally, changes in internal energy or enthalpy may be measured by techniques known collectively as 'calorimetry'.\\
2B.1 The definition of enthalpy; 2B.2 The variation of enthalpy with temperature

\subsection*{2C Thermochemistry}
'Thermochemistry' is the study of heat transactions during chemical reactions. This Topic describes methods for the de-\\
termination of enthalpy changes associated with both physical and chemical changes.\\
2C. 1 Standard enthalpy changes; 2C. 2 Standard enthalpies of formation; 2C. 3 The temperature dependence of reaction enthalpies; 2C. 4 Experimental techniques

\subsection*{2D State functions and exact differentials}
The power of thermodynamics becomes apparent by establishing relations between different properties of a system. One very useful aspect of thermodynamics is that a property can be measured indirectly by measuring others and then combining their values. The relations derived in this Topic also apply to the discussion of the liquefaction of gases and to the relation between the heat capacities of a substance under different conditions.\\
2D. 1 Exact and inexact differentials; 2D. 2 Changes in internal energy; 2D. 3 Changes in enthalpy; 2D. 4 The Joule-Thomson effect

\subsection*{2E Adiabatic changes}
'Adiabatic' processes occur without transfer of energy as heat. This Topic describes reversible adiabatic changes involving perfect gases because they figure prominently in the presentation of thermodynamics.\\
2E. 1 The change in temperature; 2E. 2 The change in pressure

\subsection*{Web resource What is an application of this material?}
A major application of thermodynamics is to the assessment of fuels and their equivalent for organisms, food. Some thermochemical aspects of fuels and foods are described in Impact 3 on the website of this text.

\section*{TOPIC 2A Internal energy}
\subsection*{Why do you need to know this material?}
The First Law of thermodynamics is the foundation of the discussion of the role of energy in chemistry. Wherever the generation or use of energy in physical transformations or chemical reactions is of interest, lying in the background are the concepts introduced by the First Law.

\subsection*{What is the key idea?}
The total energy of an isolated system is constant.

\subsection*{> What do you need to know already?}
This Topic makes use of the discussion of the properties of gases (Topic 1A), particularly the perfect gas law. It builds on the definition of work given in The chemist's toolkit 6.

For the purposes of thermodynamics, the universe is divided into two parts, the system and its surroundings. The system is the part of the world of interest. It may be a reaction vessel, an engine, an electrochemical cell, a biological cell, and so on. The surroundings comprise the region outside the system and are where measurements are made. The type of system depends on the characteristics of the boundary that divides it from the surroundings (Fig. 2A.1). If matter can be transferred through the boundary between the system and its surroundings the system is classified as open. If matter cannot pass through the boundary the system is classified as closed. Both open and closed systems can exchange energy with their surroundings.\\
\includegraphics[max width=\textwidth, center]{2025_02_27_d4b95d9718f0b9e2e6b9g-066}

Figure 2A. 1 (a) An open system can exchange matter and energy with its surroundings. (b) A closed system can exchange energy with its surroundings, but it cannot exchange matter. (c) An isolated system can exchange neither energy nor matter with its surroundings.

For example, a closed system can expand and thereby raise a weight in the surroundings; a closed system may also transfer energy to the surroundings if they are at a lower temperature. An isolated system is a closed system that has neither mechanical nor thermal contact with its surroundings.

\subsection*{2A.1 Work, heat, and energy}
Although thermodynamics deals with observations on bulk systems, it is immeasurably enriched by understanding the molecular origins of these observations.

\subsection*{(a) Operational definitions}
The fundamental physical property in thermodynamics is work: work is done to achieve motion against an opposing force (The chemist's toolkit 6). A simple example is the process of raising a weight against the pull of gravity. A process does work if in principle it can be harnessed to raise a weight somewhere in the surroundings. An example of doing work is the expansion of a gas that pushes out a piston: the motion of the piston can in principle be used to raise a weight. Another example is a chemical reaction in a cell, which leads to an electric current that can drive a motor and be used to raise a weight.

The energy of a system is its capacity to do work (see The chemist's toolkit 6 for more detail). When work is done on an otherwise isolated system (for instance, by compressing a gas or winding a spring), the capacity of the system to do work is increased; in other words, the energy of the system is increased. When the system does work (when the piston moves out or the spring unwinds), the energy of the system is reduced and it can do less work than before.

Experiments have shown that the energy of a system may be changed by means other than work itself. When the energy of a system changes as a result of a temperature difference between the system and its surroundings the energy is said to be transferred as heat. When a heater is immersed in a beaker of water (the system), the capacity of the system to do work increases because hot water can be used to do more work than the same amount of cold water. Not all boundaries permit the transfer of energy even though there is a temperature difference between the system and its surroundings. Boundaries that do permit the transfer of energy as heat are called diathermic; those that do not are called adiabatic.

\subsection*{The chemist's toolkit 6}
\subsection*{Work and energy}
Work, $w$, is done when a body is moved against an opposing force. For an infinitesimal displacement through ds (a vector), the work done on the body is

$$
\begin{array}{ll}
\mathrm{d} w_{\text {body }}=-F \cdot \mathrm{~d} \boldsymbol{s} & \begin{array}{l}
\text { Work done on body } \\
\text { [definition] }
\end{array}
\end{array}
$$

where $F \cdot \mathrm{~d} s$ is the 'scalar product' of the vectors $F$ and $\mathrm{d} s$ :

$$
F \cdot \mathrm{~d} \boldsymbol{s}=F_{x} \mathrm{~d} x+F_{y} \mathrm{~d} y+F_{z} \mathrm{~d} z \quad \begin{aligned}
& \text { Scalar product } \\
& \text { [definition] }
\end{aligned}
$$

The energy lost as work by the system, $\mathrm{d} w$, is the negative of the work done on the body, so

$$
\mathrm{d} w=F \cdot \mathrm{~d} s
$$

Work done on system [definition]\\
For motion in one dimension, $\mathrm{d} w=F_{x} \mathrm{~d} x$, with $F_{x}<0$ (so $F_{x}=$ $\left.-\left|F_{x}\right|\right)$ if it opposed the motion. The total work done along a path is the integral of this expression, allowing for the possibility that $F$ changes in direction and magnitude at each point of the path. With force in newtons (N) and distance in metres, the units of work are joules ( J ), with

$$
1 \mathrm{~J}=1 \mathrm{Nm}=1 \mathrm{~kg} \mathrm{~m}^{2} \mathrm{~s}^{-2}
$$

Energy is the capacity to do work. The SI unit of energy is the same as that of work, namely the joule. The rate of supply of energy is called the power $(P)$, and is expressed in watts (W):

$$
1 \mathrm{~W}=1 \mathrm{~J} \mathrm{~s}^{-1}
$$

A particle may possess two kinds of energy, kinetic energy and potential energy. The kinetic energy, $E_{\mathrm{k}}$, of a body is the energy the body possesses as a result of its motion. For a body of mass $m$ travelling at a speed $v$,

$$
E_{\mathrm{k}}=\frac{1}{2} m v^{2}
$$

Kinetic energy [definition]\\
Because $p=m v$ (The chemist's toolkit 3 of Topic 1B), where $p$ is the magnitude of the linear momentum, it follows that

$$
E_{\mathrm{k}}=\frac{p^{2}}{2 m}
$$

Kinetic energy [definition]

The potential energy, $E_{\mathrm{p}}$, (and commonly $V$, but do not confuse that with the volume!) of a body is the energy it possesses as a result of its position. In the absence of losses, the potential energy of a stationary particle is equal to the work that had to be done on the body to bring it to its current location. Because $\mathrm{d} w_{\text {body }}=-F_{x} \mathrm{~d} x$, it follows that $\mathrm{d} E_{\mathrm{p}}=-F_{x} \mathrm{~d} x$ and therefore

$$
F_{x}=-\frac{\mathrm{d} E_{\mathrm{p}}}{\mathrm{~d} x}
$$

Potential energy [relation to force]

If $E_{\mathrm{p}}$ increases as $x$ increases, then $F_{x}$ is negative (directed towards negative $x$, Sketch 1). Thus, the steeper the gradient (the more strongly the potential energy depends on position), the greater is the force.\\
\includegraphics[max width=\textwidth, center]{2025_02_27_d4b95d9718f0b9e2e6b9g-067}

Sketch 1\\
No universal expression for the potential energy can be given because it depends on the type of force the body experiences. For a particle of mass $m$ at an altitude $h$ close to the surface of the Earth, the gravitational potential energy is

$$
E_{\mathrm{p}}(h)=E_{\mathrm{p}}(0)+m g h \quad \text { Gravitational potential energy }
$$

where $g$ is the acceleration of free fall ( $g$ depends on location, but its 'standard value' is close to $\left.9.81 \mathrm{~m} \mathrm{~s}^{-2}\right)$. The zero of potential energy is arbitrary. For a particle close to the surface of the Earth, it is common to set $E_{\mathrm{p}}(0)=0$.\\
The Coulomb potential energy of two electric charges, $Q_{1}$ and $Q_{2}$, separated by a distance $r$ is

$$
E_{\mathrm{p}}=\frac{Q_{1} Q_{2}}{4 \pi \varepsilon r}
$$

Coulomb potential energy\\
The quantity $\varepsilon$ (epsilon) is the permittivity; its value depends upon the nature of the medium between the charges. If the charges are separated by a vacuum, then the constant is known as the vacuum permittivity, $\varepsilon_{0}$ (epsilon zero), or the electric constant, which has the value $8.854 \times 10^{-12} \mathrm{~J}^{-1} \mathrm{C}^{2} \mathrm{~m}^{-1}$. The permittivity is greater for other media, such as air, water, or oil. It is commonly expressed as a multiple of the vacuum permittivity:

$$
\varepsilon=\varepsilon_{\mathrm{r}} \varepsilon_{0}
$$

Permittivity [definition]\\
with $\varepsilon_{\mathrm{r}}$ the dimensionless relative permittivity (formerly, the dielectric constant).\\
The total energy of a particle is the sum of its kinetic and potential energies:

$$
E=E_{\mathrm{k}}+E_{\mathrm{p}} \quad l \begin{array}{ll}
\text { Total energy } & \text { [definition] }
\end{array}
$$

Provided no external forces are acting on the body, its total energy is constant. This central statement of physics is known as the law of the conservation of energy. Potential and kinetic energy may be freely interchanged, but their sum remains constant in the absence of external influences.

An exothermic process is a process that releases energy as heat. For example, combustions are chemical reactions in which substances react with oxygen, normally with a flame. The combustion of methane gas, $\mathrm{CH}_{4}(\mathrm{~g})$, is written as:

$$
\mathrm{CH}_{4}(\mathrm{~g})+2 \mathrm{O}_{2}(\mathrm{~g}) \rightarrow \mathrm{CO}_{2}(\mathrm{~g})+2 \mathrm{H}_{2} \mathrm{O}(\mathrm{l})
$$

All combustions are exothermic. Although the temperature rises in the course of the combustion, given enough time, a system in a diathermic vessel returns to the temperature of its surroundings, so it is possible to speak of a combustion 'at $25^{\circ} \mathrm{C}$ ', for instance. If the combustion takes place in an adiabatic container, the energy released as heat remains inside the container and results in a permanent rise in temperature.

An endothermic process is a process in which energy is acquired as heat. An example of an endothermic process is the vaporization of water. To avoid a lot of awkward language, it is common to say that in an exothermic process energy is transferred 'as heat' to the surroundings and in an endothermic process energy is transferred 'as heat' from the surroundings into the system. However, it must never be forgotten that heat is a process (the transfer of energy as a result of a temperature difference), not an entity. An endothermic process in a diathermic container results in energy flowing into the system as heat to restore the temperature to that of the surroundings. An exothermic process in a similar diathermic container results in a release of energy as heat into the surroundings. When an endothermic process takes place in an adiabatic container, it results in a lowering of temperature of the system; an exothermic process results in a rise of temperature. These features are summarized in Fig. 2A.2.\\
\includegraphics[max width=\textwidth, center]{2025_02_27_d4b95d9718f0b9e2e6b9g-068(2)}

Figure 2A. 2 (a) When an endothermic process occurs in an adiabatic system, the temperature falls; (b) if the process is exothermic, then the temperature rises. (c) When an endothermic process occurs in a diathermic container, energy enters as heat from the surroundings (which remain at the same temperature), and the system remains at the same temperature. (d) If the process is exothermic, then energy leaves as heat, and the process is isothermal.

\subsection*{(b) The molecular interpretation of heat and work}
In molecular terms, heating is the transfer of energy that makes use of disorderly, apparently random, molecular motion in the surroundings. The disorderly motion of molecules is called thermal motion. The thermal motion of the molecules in the hot surroundings stimulates the molecules in the cooler system to move more vigorously and, as a result, the energy of the cooler system is increased. When a system heats its surroundings, molecules of the system stimulate the thermal motion of the molecules in the surroundings (Fig. 2A.3).

In contrast, work is the transfer of energy that makes use of organized motion in the surroundings (Fig. 2A.4). When a weight is raised or lowered, its atoms move in an organized way (up or down). The atoms in a spring move in an orderly way when it is wound; the electrons in an electric current\\
\includegraphics[max width=\textwidth, center]{2025_02_27_d4b95d9718f0b9e2e6b9g-068}

Figure 2A. 3 When energy is transferred to the surroundings as heat, the transfer stimulates random motion of the atoms in the surroundings. Transfer of energy from the surroundings to the system makes use of random motion (thermal motion) in the surroundings.\\
\includegraphics[max width=\textwidth, center]{2025_02_27_d4b95d9718f0b9e2e6b9g-068(1)}

Figure 2A. 4 When a system does work, it stimulates orderly motion in the surroundings. For instance, the atoms shown here may be part of a weight that is being raised. The ordered motion of the atoms in a falling weight does work on the system.\\
move in the same direction. When a system does work it causes atoms or electrons in its surroundings to move in an organized way. Likewise, when work is done on a system, molecules in the surroundings are used to transfer energy to it in an organized way, as the atoms in a weight are lowered or a current of electrons is passed.

The distinction between work and heat is made in the surroundings. The fact that a falling weight may stimulate thermal motion in the system is irrelevant to the distinction between heat and work: work is identified as energy transfer making use of the organized motion of atoms in the surroundings, and heat is identified as energy transfer making use of thermal motion in the surroundings. In the compression of a gas in an adiabatic enclosure, for instance, work is done on the system as the atoms of the compressing weight descend in an orderly way, but the effect of the incoming piston is to accelerate the gas molecules to higher average speeds. Because collisions between molecules quickly randomize their directions, the orderly motion of the atoms of the weight is in effect stimulating thermal motion in the gas. The weight is observed to fall, leading to the orderly descent of its atoms, and work is done even though it is stimulating thermal motion.

\subsection*{2A.2 The definition of internal energy}
In thermodynamics, the total energy of a system is called its internal energy, $U$. The internal energy is the total kinetic and potential energy of the constituents (the atoms, ions, or molecules) of the system. It does not include the kinetic energy arising from the motion of the system as a whole, such as its kinetic energy as it accompanies the Earth on its orbit round the Sun. That is, the internal energy is the energy 'internal' to the system. The change in internal energy is denoted by $\Delta U$ when a system changes from an initial state $i$ with internal energy $U_{\mathrm{i}}$ to a final state $f$ of internal energy $U_{\mathrm{f}}$ :


\begin{equation*}
\Delta U=U_{\mathrm{f}}-U_{\mathrm{i}} \tag{2A.1}
\end{equation*}


\subsection*{The chemist's toolkit 7 The equipartition theorem}
The Boltzmann distribution (see the Prologue) can be used to calculate the average energy associated with each mode of motion of an atom or molecule in a sample at a given temperature. However, when the temperature is so high that many energy levels are occupied, there is a much simpler way to find the average energy, through the equipartition theorem:

For a sample at thermal equilibrium the average value of each quadratic contribution to the energy is $\frac{1}{2} k T$.

A 'quadratic contribution' is a term that is proportional to the square of the momentum (as in the expression for the kinetic ener-

A convention used throughout thermodynamics is that $\Delta X=X_{\mathrm{f}}-X_{\mathrm{i}}$, where $X$ is a property (a 'state function') of the system.

The internal energy is a state function, a property with a value that depends only on the current state of the system and is independent of how that state has been prepared. In other words, internal energy is a function of the variables that determine the current state of the system. Changing any one of the state variables, such as the pressure, may result in a change in internal energy. That the internal energy is a state function has consequences of the greatest importance (Topic 2D).

The internal energy is an extensive property of a system (a property that depends on the amount of substance present; see The chemist's toolkit 2 in Topic 1A) and is measured in joules $\left(1 \mathrm{~J}=1 \mathrm{~kg} \mathrm{~m}^{2} \mathrm{~s}^{-2}\right)$. The molar internal energy, $U_{\mathrm{m}}$, is the internal energy divided by the amount of substance in a system, $U_{\mathrm{m}}=$ $U / n$; it is an intensive property (a property independent of the amount of substance) and is commonly reported in kilojoules per mole ( $\mathrm{kJ} \mathrm{mol}^{-1}$ ).

\subsection*{(a) Molecular interpretation of internal energy}
A molecule has a certain number of motional degrees of freedom, such as the ability to move through space (this motion is called 'translation'), rotate, or vibrate. Many physical and chemical properties depend on the energy associated with each of these modes of motion. For example, a chemical bond might break if a lot of energy becomes concentrated in it, for instance as vigorous vibration. The internal energy of a sample increases as the temperature is raised and states of higher energy become more highly populated.

The 'equipartition theorem' of classical mechanics, introduced in The chemist's toolkit 7, can be used to predict the contributions of each mode of motion of a molecule to the total energy of a collection of non-interacting molecules (that is, of a perfect gas, and providing quantum effects can be ignored).\\
gy, $E_{\mathrm{k}}=p^{2} / 2 m$; The chemist's toolkit 6) or the displacement from an equilibrium position (as for the potential energy of a harmonic oscillator, $E_{\mathrm{p}}=\frac{1}{2} k_{\mathrm{f}} x^{2}$ ). The theorem is a conclusion from classical mechanics and for quantized systems is applicable only when the separation between the energy levels is so small compared to $k T$ that many states are populated. Under normal conditions the equipartition theorem gives good estimates for the average energies associated with translation and rotation. However, the separation between vibrational and electronic states is typically much greater than for rotation or translation, and for these types of motion the equipartition theorem is unlikely to apply.

\subsection*{Brief illustration 2A. 1}
An atom in a gas can move in three dimensions, so its translational kinetic energy is the sum of three quadratic contributions:

$$
E_{\text {trans }}=\frac{1}{2} m v_{x}^{2}+\frac{1}{2} m v_{y}^{2}+\frac{1}{2} m v_{z}^{2}
$$

The equipartition theorem predicts that the average energy for each of these quadratic contributions is $\frac{1}{2} k T$. Thus, the average kinetic energy is $E_{\text {trans }}=3 \times \frac{1}{2} k T=\frac{3}{2} k T$. The molar translational energy is therefore $E_{\text {trans, m }}=\frac{3}{2} k T \times N_{\mathrm{A}}=\frac{3}{2} R T$. At $25^{\circ} \mathrm{C}, R T=$ $2.48 \mathrm{~kJ} \mathrm{~mol}^{-1}$, so the contribution of translation to the molar internal energy of a perfect gas is $3.72 \mathrm{~kJ} \mathrm{~mol}^{-1}$.

The contribution to the internal energy of a collection of perfect gas molecules is independent of the volume occupied by the molecules: there are no intermolecular interactions in a perfect gas, so the distance between the molecules has no effect on the energy. That is,

The internal energy of a perfect gas is independent of the volume it occupies.

The internal energy of interacting molecules in condensed phases also has a contribution from the potential energy of their interaction, but no simple expressions can be written down in general. Nevertheless, it remains true that as the temperature of a system is raised, the internal energy increases as the various modes of motion become more highly excited.

\subsection*{(b) The formulation of the First Law}
It has been found experimentally that the internal energy of a system may be changed either by doing work on the system or by heating it. Whereas it might be known how the energy transfer has occurred (if a weight has been raised or lowered in the surroundings, indicating transfer of energy by doing work, or if ice has melted in the surroundings, indicating transfer of energy as heat), the system is blind to the mode employed. That is,

Heat and work are equivalent ways of changing the internal energy of a system.

A system is like a bank: it accepts deposits in either currency (work or heat), but stores its reserves as internal energy. It is also found experimentally that if a system is isolated from its surroundings, meaning that it can exchange neither matter nor energy with its surroundings, then no change in internal energy takes place. This summary of observations is now known as the First Law of thermodynamics and is expressed as follows:

The internal energy of an isolated system is constant.

It is not possible to use a system to do work, leave it isolated, and then come back expecting to find it restored to its original state with the same capacity for doing work. The experimental evidence for this observation is that no 'perpetual motion machine', a machine that does work without consuming fuel or using some other source of energy, has ever been built.

These remarks may be expressed symbolically as follows. If $w$ is the work done on a system, $q$ is the energy transferred as heat to a system, and $\Delta U$ is the resulting change in internal energy, then


\begin{equation*}
\Delta U=q+w \quad \text { Mathematical statement of the First Law } \tag{2A.2}
\end{equation*}


Equation 2A. 2 summarizes the equivalence of heat and work for bringing about changes in the internal energy and the fact that the internal energy is constant in an isolated system (for which $q=0$ and $w=0$ ). It states that the change in internal energy of a closed system is equal to the energy that passes through its boundary as heat or work. Equation 2A. 2 employs the 'acquisitive convention', in which $w$ and $q$ are positive if energy is transferred to the system as work or heat and are negative if energy is lost from the system. ${ }^{1}$ In other words, the flow of energy as work or heat is viewed from the system's perspective.

\subsection*{Brief illustration 2A. 2}
If an electric motor produces 15 kJ of energy each second as mechanical work and loses 2 kJ as heat to the surroundings, then the change in the internal energy of the motor each second is $\Delta U=-2 \mathrm{~kJ}-15 \mathrm{~kJ}=-17 \mathrm{~kJ}$. Suppose that, when a spring is wound, 100 J of work is done on it but 15 J escapes to the surroundings as heat. The change in internal energy of the spring is $\Delta U=100 \mathrm{~J}-15 \mathrm{~J}=+85 \mathrm{~J}$.

A note on good practice Always include the sign of $\Delta U$ (and of $\Delta X$ in general), even if it is positive.

\subsection*{2A.3 Expansion work}
The way is opened to powerful methods of calculation by switching attention to infinitesimal changes in the variables that describe the state of the system (such as infinitesimal change in temperature) and infinitesimal changes in the internal energy $\mathrm{d} U$. Then, if the work done on a system is $\mathrm{d} w$ and the energy supplied to it as heat is $\mathrm{d} q$, in place of eqn 2 A .2 , it follows that


\begin{equation*}
\mathrm{d} U=\mathrm{d} q+\mathrm{d} w \tag{2A.3}
\end{equation*}


\footnotetext{${ }^{1}$ Many engineering texts adopt a different convention for work: $w>0$ if energy is used to do work in the surroundings.
}The ability to use this expression depends on being able to relate $\mathrm{d} q$ and $\mathrm{d} w$ to events taking place in the surroundings.

A good starting point is a discussion of expansion work, the work arising from a change in volume. This type of work includes the work done by a gas as it expands and drives back the atmosphere. Many chemical reactions result in the generation of gases (for instance, the thermal decomposition of calcium carbonate or the combustion of hydrocarbons), and the thermodynamic characteristics of the reaction depend on the work that must be done to make room for the gas it has produced. The term 'expansion work' also includes work associated with negative changes of volume, that is, compression.

\subsection*{(a) The general expression for work}
The calculation of expansion work starts from the definition in The chemist's toolkit 6 with the sign of the opposing force written explicitly:

\[
\begin{array}{ll}
\mathrm{d} w=-|F| \mathrm{d} z & \begin{array}{l}
\text { Work done } \\
\text { [definition] }
\end{array} \tag{2A.4}
\end{array}
\]

The negative sign implies that the internal energy of the system doing the work decreases when the system moves an object against an opposing force of magnitude $|F|$, and there are no other changes. That is, if $\mathrm{d} z$ is positive (motion to positive $z$ ), $\mathrm{d} w$ is negative, and the internal energy decreases ( $\mathrm{d} U$ in eqn 2 A .3 is negative provided that $\mathrm{d} q=0$ ).

Now consider the arrangement shown in Fig. 2A.5, in which one wall of a system is a massless, frictionless, rigid, perfectly fitting piston of area $A$. If the external pressure is $p_{\text {ex }}$, the magnitude of the force acting on the outer face of the piston is $|F|=$ $p_{\mathrm{ex}} A$. The work done when the system expands through a distance $\mathrm{d} z$ against an external pressure $p_{\mathrm{ex}}$, is $\mathrm{d} w=-p_{\mathrm{ex}} A \mathrm{~d} z$. The quantity $A \mathrm{~d} z$ is the change in volume, $\mathrm{d} V$, in the course of the expansion. Therefore, the work done when the system expands by $\mathrm{d} V$ against a pressure $p_{\mathrm{ex}}$ is

$$
\mathrm{d} w=-p_{\mathrm{ex}} \mathrm{~d} V
$$

Expansion work (2A.5a)\\
\includegraphics[max width=\textwidth, center]{2025_02_27_d4b95d9718f0b9e2e6b9g-071}

Figure 2A. 5 When a piston of area $A$ moves out through a distance $\mathrm{d} z$, it sweeps out a volume $\mathrm{d} V=A d z$. The external pressure $p_{\text {ex }}$ is equivalent to a weight pressing on the piston, and the magnitude of the force opposing expansion is $p_{\text {ex }} A$.

Table 2A. 1 Varieties of work*

\begin{center}
\begin{tabular}{|c|c|c|c|}
\hline
Type of work & $\mathrm{d} w$ & Comments & Units ${ }^{\dagger}$ \\
\hline
\multirow[t]{2}{*}{Expansion} & \multirow[t]{2}{*}{$-p_{\text {ex }} \mathrm{d} V$} & $p_{\text {ex }}$ is the external pressure & Pa \\
\hline
 &  & $\mathrm{d} V$ is the change in volume & $\mathrm{m}^{3}$ \\
\hline
\multirow[t]{2}{*}{Surface expansion} & \multirow[t]{2}{*}{$\gamma \mathrm{d} \sigma$} & $\gamma$ is the surface tension & $\mathrm{N} \mathrm{m}{ }^{-1}$ \\
\hline
 &  & $\mathrm{d} \sigma$ is the change in area & $\mathrm{m}^{2}$ \\
\hline
\multirow[t]{2}{*}{Extension} & \multirow[t]{2}{*}{$f \mathrm{~d} l$} & $f$ is the tension & N \\
\hline
 &  & $\mathrm{d} l$ is the change in length & m \\
\hline
\multirow[t]{4}{*}{Electrical} & \multirow[t]{2}{*}{$\phi \mathrm{d} Q$} & $\phi$ is the electric potential & V \\
\hline
 &  & $\mathrm{d} Q$ is the change in charge & C \\
\hline
 & \multirow[t]{2}{*}{$Q \mathrm{~d} \phi$} & $\mathrm{d} \phi$ is the potential difference & V \\
\hline
 &  & $Q$ is the charge transferred & C \\
\hline
\end{tabular}
\end{center}

\begin{itemize}
  \item In general, the work done on a system can be expressed in the form $\mathrm{d} w=-|F| \mathrm{d} z$, where $|F|$ is the magnitude of a 'generalized force' and $\mathrm{d} z$ is a 'generalized displacement'.\\
${ }^{\dagger}$ For work in joules (J). Note that $1 \mathrm{Nm}=1 \mathrm{~J}$ and $1 \mathrm{VC}=1 \mathrm{~J}$.
\end{itemize}

To obtain the total work done when the volume changes from an initial value $V_{\mathrm{i}}$ to a final value $V_{\mathrm{f}}$ it is necessary to integrate this expression between the initial and final volumes:


\begin{equation*}
w=-\int_{V_{\mathrm{i}}}^{V_{\mathrm{f}}} p_{\mathrm{ex}} \mathrm{~d} V \tag{2A.5b}
\end{equation*}


The force acting on the piston, $p_{\mathrm{ex}} A$, is equivalent to the force arising from a weight that is raised as the system expands. If the system is compressed instead, then the same weight is lowered in the surroundings and eqn 2 A .5 b can still be used, but now $V_{\mathrm{f}}<V_{\mathrm{i}}$. It is important to note that it is still the external pressure that determines the magnitude of the work. This somewhat perplexing conclusion seems to be inconsistent with the fact that the gas inside the container is opposing the compression. However, when a gas is compressed, the ability of the surroundings to do work is diminished to an extent determined by the weight that is lowered, and it is this energy that is transferred into the system.

Other types of work (e.g. electrical work), which are called either non-expansion work or additional work, have analogous expressions, with each one the product of an intensive factor (the pressure, for instance) and an extensive factor (such as a change in volume). Some are collected in Table 2A.1. The present discussion focuses on how the work associated with changing the volume, the expansion work, can be extracted from eqn 2A.5b.

\subsection*{(b) Expansion against constant pressure}
Suppose that the external pressure is constant throughout the expansion. For example, the piston might be pressed on by the atmosphere, which exerts the same pressure throughout the expansion. A chemical example of this condition is the expansion of a gas formed in a chemical reaction in a container\\
\includegraphics[max width=\textwidth, center]{2025_02_27_d4b95d9718f0b9e2e6b9g-072(1)}

Figure 2A. 6 The work done by a gas when it expands against a constant external pressure, $p_{\text {ex }}$ is equal to the shaded area in this example of an indicator diagram.\\
that can expand. Equation 2 A .5 b is then evaluated by taking the constant $p_{\text {ex }}$ outside the integral:

$$
w=-p_{\mathrm{ex}} \int_{V_{\mathrm{i}}}^{V_{\mathrm{f}}} \mathrm{~d} V=-p_{\mathrm{ex}}\left(V_{\mathrm{f}}-V_{\mathrm{i}}\right)
$$

Therefore, if the change in volume is written as $\Delta V=V_{\mathrm{f}}-V_{\mathrm{i}}$,

\[
w=-p_{\mathrm{ex}} \Delta V \quad \begin{align*}
& \text { Expansion work }  \tag{2A.6}\\
& \text { [constant external pressure] }
\end{align*}
\]

This result is illustrated graphically in Fig. 2A.6, which makes use of the fact that the magnitude of an integral can be interpreted as an area. The magnitude of $w$, denoted $|w|$, is equal to the area beneath the horizontal line at $p=$ $p_{\text {ex }}$ lying between the initial and final volumes. A $p, V$-graph used to illustrate expansion work is called an indicator diagram; James Watt first used one to indicate aspects of the operation of his steam engine.\\
Free expansion is expansion against zero opposing force. It occurs when $p_{\text {ex }}=0$. According to eqn 2A.6, in this case


\begin{equation*}
w=0 \tag{2A.7}
\end{equation*}


Work of free expansion\\
That is, no work is done when a system expands freely. Expansion of this kind occurs when a gas expands into a vacuum.

\subsection*{Example 2A. 1 Calculating the work of gas production}
Calculate the work done when 50 g of iron reacts with hydrochloric acid to produce $\mathrm{FeCl}_{2}(\mathrm{aq})$ and hydrogen in (a) a closed vessel of fixed volume, (b) an open beaker at $25^{\circ} \mathrm{C}$.

Collect your thoughts You need to judge the magnitude of the volume change and then to decide how the process occurs. If there is no change in volume, there is no expansion work however the process takes place. If the system expands against a constant external pressure, the work can be calculated from eqn 2A.6. A general feature of processes in which a condensed\\
phase changes into a gas is that you can usually neglect the volume of a condensed phase relative to the volume of the gas it forms.

The solution In (a) the volume cannot change, so no expansion work is done and $w=0$. In (b) the gas drives back the atmosphere and therefore $w=-p_{\mathrm{ex}} \Delta V$. The initial volume can be neglected because the final volume (after the production of gas) is so much larger and $\Delta V=V_{\mathrm{f}}-V_{\mathrm{i}} \approx V_{\mathrm{f}}=n R T / p_{\mathrm{ex}}$, where $n$ is the amount of $\mathrm{H}_{2}$ produced. Therefore,

$$
w=-p_{\mathrm{ex}} \Delta V \approx-p_{\mathrm{ex}} \times \frac{n R T}{p_{\mathrm{ex}}}=-n R T
$$

Because the reaction is $\mathrm{Fe}(\mathrm{s})+2 \mathrm{HCl}(\mathrm{aq}) \rightarrow \mathrm{FeCl}_{2}(\mathrm{aq})+\mathrm{H}_{2}(\mathrm{~g})$, $1 \mathrm{~mol} \mathrm{H}_{2}$ is generated when 1 mol Fe is consumed, and $n$ can be taken as the amount of Fe atoms that react. Because the molar mass of Fe is $55.85 \mathrm{~g} \mathrm{~mol}^{-1}$, it follows that

$$
\begin{aligned}
w & =-\frac{50 \mathrm{~g}}{55.85 \mathrm{~g} \mathrm{~mol}^{-1}} \times\left(8.3145 \mathrm{JK}^{-1} \mathrm{~mol}^{-1}\right) \times(298 \mathrm{~K}) \\
& \approx-2.2 \mathrm{~kJ}
\end{aligned}
$$

The system (the reaction mixture) does 2.2 kJ of work driving back the atmosphere.

Comment. The magnitude of the external pressure does not affect the final result: the lower the pressure, the larger is the volume occupied by the gas, so the effects cancel.

Self-test 2A.1 Calculate the expansion work done when 50 g of water is electrolysed under constant pressure at $25^{\circ} \mathrm{C}$.\\
\includegraphics[max width=\textwidth, center]{2025_02_27_d4b95d9718f0b9e2e6b9g-072}

\subsection*{(c) Reversible expansion}
A reversible change in thermodynamics is a change that can be reversed by an infinitesimal modification of a variable. The key word 'infinitesimal' sharpens the everyday meaning of the word 'reversible' as something that can change direction. One example of reversibility is the thermal equilibrium of two systems with the same temperature. The transfer of energy as heat between the two is reversible because, if the temperature of either system is lowered infinitesimally, then energy flows into the system with the lower temperature. If the temperature of either system at thermal equilibrium is raised infinitesimally, then energy flows out of the hotter system. There is obviously a very close relationship between reversibility and equilibrium: systems at equilibrium are poised to undergo reversible change.

Suppose a gas is confined by a piston and that the external pressure, $p_{\text {ex }}$, is set equal to the pressure, $p$, of the confined gas. Such a system is in mechanical equilibrium with its surroundings because an infinitesimal change in the external pressure in either direction causes changes in volume in\\
opposite directions. If the external pressure is reduced infinitesimally, the gas expands slightly. If the external pressure is increased infinitesimally, the gas contracts slightly. In either case the change is reversible in the thermodynamic sense. If, on the other hand, the external pressure is measurably greater than the internal pressure, then decreasing $p_{\text {ex }}$ infinitesimally will not decrease it below the pressure of the gas, so will not change the direction of the process. Such a system is not in mechanical equilibrium with its surroundings and the compression is thermodynamically irreversible.

To achieve reversible expansion $p_{\text {ex }}$ is set equal to $p$ at each stage of the expansion. In practice, this equalization could be achieved by gradually removing weights from the piston so that the downward force due to the weights always matches the changing upward force due to the pressure of the gas or by gradually adjusting the external pressure to match the pressure of the expanding gas. When $p_{\mathrm{ex}}=p$, eqn 2A.5a becomes


\begin{equation*}
\mathrm{d} w=-p_{\mathrm{ex}} \mathrm{~d} V=-p \mathrm{~d} V \quad \text { Reversible expansion work } \tag{2A.8a}
\end{equation*}


Although the pressure inside the system appears in this expression for the work, it does so only because $p_{\text {ex }}$ has been arranged to be equal to $p$ to ensure reversibility. The total work of reversible expansion from an initial volume $V_{\mathrm{i}}$ to a final volume $V_{\mathrm{f}}$ is therefore


\begin{equation*}
w=-\int_{V_{\mathrm{i}}}^{V_{\mathrm{f}}} p \mathrm{~d} V \tag{2A.8b}
\end{equation*}


The integral can be evaluated once it is known how the pressure of the confined gas depends on its volume. Equation 2 A .8 b is the link with the material covered in Focus 1 because, if the equation of state of the gas is known, $p$ can be expressed in terms of $V$ and the integral can be evaluated.

\subsection*{(d) Isothermal reversible expansion of a perfect gas}
Consider the isothermal reversible expansion of a perfect gas. The expansion is made isothermal by keeping the system in thermal contact with its unchanging surroundings (which may be a constant-temperature bath). Because the equation of state is $p V=n R T$, at each stage $p=n R T / V$, with $V$ the volume at that stage of the expansion. The temperature $T$ is constant in an isothermal expansion, so (together with $n$ and $R$ ) it may be taken outside the integral. It follows that the work of isothermal reversible expansion of a perfect gas from $V_{\mathrm{i}}$ to $V_{\mathrm{f}}$ at a temperature $T$ is

\[
w=-n R T \overbrace{\int_{V_{\mathrm{i}}}^{V_{\mathrm{i}}} \frac{\mathrm{~d} V}{V}}^{\text {Integral } \mathrm{A} .2}=-n R T \ln \frac{V_{\mathrm{f}}}{V_{\mathrm{i}}} \begin{align*}
& \begin{array}{l}
\text { revers of iblothe expanal } \\
\text { [perfect gas] }
\end{array} \tag{2A.9}
\end{align*}
\]

\subsection*{Brief Illustration 2A. 3}
When a sample of 1.00 mol Ar , regarded here as a perfect gas, undergoes an isothermal reversible expansion at $20.0^{\circ} \mathrm{C}$ from $10.0 \mathrm{dm}^{3}$ to $30.0 \mathrm{dm}^{3}$ the work done is

$$
\begin{aligned}
w & =-(1.00 \mathrm{~mol}) \times\left(8.3145 \mathrm{JK}^{-1} \mathrm{~mol}^{-1}\right) \times(293.2 \mathrm{~K}) \ln \frac{30.0 \mathrm{dm}^{3}}{10.0 \mathrm{dm}^{3}} \\
& =-2.68 \mathrm{~kJ}
\end{aligned}
$$

When the final volume is greater than the initial volume, as in an expansion, the logarithm in eqn 2A. 9 is positive and hence $w<0$. In this case, the system has done work on the surroundings and there is a corresponding negative contribution to its internal energy. (Note the cautious language: as seen later, there is a compensating influx of energy as heat, so overall the internal energy is constant for the isothermal expansion of a perfect gas.) The equations also show that more work is done for a given change of volume when the temperature is increased: at a higher temperature the greater pressure of the confined gas needs a higher opposing pressure to ensure reversibility and the work done is correspondingly greater.

The result of the calculation can be illustrated by an indicator diagram in which the magnitude of the work done is equal to the area under the isotherm $p=n R T / V$ (Fig. 2A.7). Superimposed on the diagram is the rectangular area obtained for irreversible expansion against constant external pressure fixed at the same final value as that reached in the reversible expansion. More work is obtained when the expansion is reversible (the area is greater) because matching the external pressure to the internal pressure at each stage of the process ensures that none of the pushing power of the system is wasted. It is not possible to obtain more work than that for\\
\includegraphics[max width=\textwidth, center]{2025_02_27_d4b95d9718f0b9e2e6b9g-073}

Figure 2A. 7 The work done by a perfect gas when it expands reversibly and isothermally is equal to the area under the isotherm $p=n R T / V$. The work done during the irreversible expansion against the same final pressure is equal to the rectangular area shown slightly darker. Note that the reversible work done is greater than the irreversible work done.\\
the reversible process because increasing the external pressure even infinitesimally at any stage results in compression. It can be inferred from this discussion that, because some pushing power is wasted when $p>p_{\text {ex }}$, the maximum work available from a system operating between specified initial and final states is obtained when the change takes place reversibly.

\subsection*{2A. 4 Heat transactions}
In general, the change in internal energy of a system is


\begin{equation*}
\mathrm{d} U=\mathrm{d} q+\mathrm{d} w_{\text {exp }}+\mathrm{d} w_{\mathrm{add}} \tag{2A.10}
\end{equation*}


where $\mathrm{d} w_{\text {add }}$ is work in addition ('add' for additional) to the expansion work, $\mathrm{d} w_{\text {exp }}$. For instance, $\mathrm{d} w_{\text {add }}$ might be the electrical work of driving a current of electrons through a circuit. A system kept at constant volume can do no expansion work, so in that case $\mathrm{d} w_{\text {exp }}=0$. If the system is also incapable of doing any other kind of work (if it is not, for instance, an electrochemical cell connected to an electric motor), then $\mathrm{d} w_{\text {add }}=0$ too. Under these circumstances:

\[
\mathrm{d} U=\mathrm{d} q \quad \begin{align*}
& \text { Heat transferred at }  \tag{2A.11a}\\
& \text { constant volume }
\end{align*}
\]

This relation can also be expressed as $\mathrm{d} U=\mathrm{d} q_{v}$, where the subscript implies the constraint of constant volume. For a measurable change between states $i$ and $f$ along a path at constant volume,

$$
\int_{\int_{i}^{f}}^{U_{f}-U_{i}} \underbrace{\mathrm{~d}_{\mathrm{i}}}=\overbrace{\int_{\mathrm{f}}^{\mathrm{f}} \mathrm{~d} q_{V}}^{q_{v}}
$$

which is summarized as


\begin{equation*}
\Delta U=q_{V} \tag{2A.11b}
\end{equation*}


Note that the integral over $\mathrm{d} q$ is not written as $\Delta q$ because $q$, unlike $U$, is not a state function. It follows from eqn 2A.11b that measuring the energy supplied as heat to a system at constant volume is equivalent to measuring the change in internal energy of the system.

\subsection*{(a) Calorimetry}
Calorimetry is the study of the transfer of energy as heat during a physical or chemical process. A calorimeter is a device for measuring energy transferred as heat. The most common device for measuring $q_{V}$ (and therefore $\Delta U$ ) is an adiabatic bomb calorimeter (Fig. 2A.8). The process to be studied-which may be a chemical reaction-is initiated inside a constant-volume container, the 'bomb'. The bomb is immersed in a stirred water bath, and the whole device is the calorimeter. The calorimeter is also immersed in an outer water bath. The water in the calorimeter and of the outer bath are both monitored and adjusted\\
\includegraphics[max width=\textwidth, center]{2025_02_27_d4b95d9718f0b9e2e6b9g-074}

Figure 2A.8 A constant-volume bomb calorimeter. The 'bomb' is the central vessel, which is strong enough to withstand high pressures. The calorimeter is the entire assembly shown here. To ensure adiabaticity, the calorimeter is immersed in a water bath with a temperature continuously readjusted to that of the calorimeter at each stage of the combustion.\\
to the same temperature. This arrangement ensures that there is no net loss of heat from the calorimeter to the surroundings (the bath) and hence that the calorimeter is adiabatic.

The change in temperature, $\Delta T$, of the calorimeter is proportional to the energy that the reaction releases or absorbs as heat. Therefore, $q_{V}$ and hence $\Delta U$ can be determined by measuring $\Delta T$. The conversion of $\Delta T$ to $q_{V}$ is best achieved by calibrating the calorimeter using a process of known output and determining the calorimeter constant, the constant $C$ in the relation


\begin{equation*}
q=C \Delta T \tag{2A.12}
\end{equation*}


The calorimeter constant may be measured electrically by passing a constant current, $I$, from a source of known potential difference, $\Delta \phi$, through a heater for a known period of time, $t$, for then (The chemist's toolkit 8)


\begin{equation*}
q=I t \Delta \phi \tag{2A.13}
\end{equation*}


\subsection*{Brief illustration 2A. 4}
If a current of 10.0 A from a 12 V supply is passed for 300 s , then from eqn 2A. 13 the energy supplied as heat is

$$
q=(10.0 \mathrm{~A}) \times(300 \mathrm{~s}) \times(12 \mathrm{~V})=3.6 \times 10^{4} \mathrm{~A} \mathrm{~V} \mathrm{~s}=36 \mathrm{~kJ}
$$

The result in joules is obtained by using $1 \mathrm{AVs}=1\left(\mathrm{C} \mathrm{s}^{-1}\right) \mathrm{Vs}=$ $1 \mathrm{CV}=1 \mathrm{~J}$. If the observed rise in temperature is 5.5 K , then the calorimeter constant is $C=(36 \mathrm{~kJ}) /(5.5 \mathrm{~K})=6.5 \mathrm{~kJ} \mathrm{~K}^{-1}$.

Alternatively, $C$ may be determined by burning a known mass of substance (benzoic acid is often used) that has a known heat output. With $C$ known, it is simple to interpret an observed temperature rise as a release of energy as heat.

\subsection*{The chemist's toolkit 8 Electrical charge, current, power, and energy}
Electrical charge, $Q$, is measured in coulombs, $C$. The fundamental charge, $e$, the magnitude of charge carried by a single electron or proton, is approximately $1.6 \times 10^{-19} \mathrm{C}$. The motion of charge gives rise to an electric current, $I$, measured in coulombs per second, or amperes, A , where $1 \mathrm{~A}=1 \mathrm{C} \mathrm{s}^{-1}$. If the electric charge is that of electrons (as it is for the current in a metal), then a current of 1 A represents the flow of $6 \times 10^{18}$ electrons ( $10 \mu \mathrm{~mol} \mathrm{e}^{-}$) per second.\\
When a current $I$ flows through a potential difference $\Delta \phi$ (measured in volts, V , with $1 \mathrm{~V}=1 \mathrm{~J} \mathrm{~A}^{-1}$ ), the power, $P$, is

$$
P=I \Delta \phi
$$

It follows that if a constant current flows for a period $t$ the energy supplied is

$$
E=P t=I t \Delta \phi
$$

Because $1 \mathrm{AVs}=1\left(\mathrm{C} \mathrm{s}^{-1}\right) \mathrm{Vs}=1 \mathrm{CV}=1 \mathrm{~J}$, the energy is obtained in joules with the current in amperes, the potential difference in volts, and the time in seconds. That energy may be supplied as either work (to drive a motor) or as heat (through a 'heater'). In the latter case

$$
q=I t \Delta \phi
$$

\subsection*{Brief illustration 2A. 5}
In Brief illustration 2A. 1 it is shown that the translational contribution to the molar internal energy of a perfect monatomic gas is $\frac{3}{2} R T$. Because this is the only contribution to the internal energy, $U_{\mathrm{m}}(T)=\frac{3}{2} R T$. It follows from eqn 2 A .14 that

$$
C_{V, \mathrm{~m}}=\frac{\partial}{\partial T}\left\{\frac{3}{2} R T\right\}=\frac{3}{2} R
$$

The numerical value is $12.47 \mathrm{~J} \mathrm{~K}^{-1} \mathrm{~mol}^{-1}$.

Heat capacities are extensive properties: 100 g of water, for instance, has 100 times the heat capacity of 1 g of water (and therefore requires 100 times the energy as heat to bring about the same rise in temperature). The molar heat capacity at constant volume, $C_{V, \mathrm{~m}}=C_{V} / n$, is the heat capacity per mole of substance, and is an intensive property (all molar quantities are intensive). For certain applications it is useful to know the\\
\includegraphics[max width=\textwidth, center]{2025_02_27_d4b95d9718f0b9e2e6b9g-075}

Figure 2A. 10 The internal energy of a system varies with volume and temperature, perhaps as shown here by the surface. The variation of the internal energy with temperature at one particular constant volume is illustrated by the curve drawn parallel to the temperature axis. The slope of this curve at any point is the partial derivative $(\partial U / \partial T)_{v}$.

\subsection*{The chemist's toolkit 9}
\subsection*{Partial derivatives}
A partial derivative of a function of more than one variable, such as $f(x, y)$, is the slope of the function with respect to one of the variables, all the other variables being held constant (Sketch 1). Although a partial derivative shows how a function changes when one variable changes, it may be used to determine how the function changes when more than one variable changes by an infinitesimal amount. Thus, if $f$ is a function of $x$ and $y$, then when $x$ and $y$ change by $\mathrm{d} x$ and $\mathrm{d} y$, respectively, $f$ changes by

$$
\mathrm{d} f=\left(\frac{\partial f}{\partial x}\right)_{y} \mathrm{~d} x+\left(\frac{\partial f}{\partial y}\right)_{x} \mathrm{~d} y
$$

where the symbol $\partial$ ('curly d') is used (instead of d) to denote a partial derivative and the subscript on the parentheses indicates which variable is being held constant.\\
\includegraphics[max width=\textwidth, center]{2025_02_27_d4b95d9718f0b9e2e6b9g-076}

Sketch 1\\
The quantity $\mathrm{d} f$ is also called the differential of $f$. Successive partial derivatives may be taken in any order:

$$
\left(\frac{\partial}{\partial y}\left(\frac{\partial f}{\partial x}\right)_{y}\right)_{x}=\left(\frac{\partial}{\partial x}\left(\frac{\partial f}{\partial y}\right)_{x}\right)_{y}
$$

For example, suppose that $f(x, y)=a x^{3} y+b y^{2}$ (the function plotted in Sketch 1) then

$$
\left(\frac{\partial f}{\partial x}\right)_{y}=3 a x^{2} y \quad\left(\frac{\partial f}{\partial y}\right)_{x}=a x^{3}+2 b y
$$

Then, when $x$ and $y$ undergo infinitesimal changes, $f$ changes by

$$
\mathrm{d} f=3 a x^{2} y \mathrm{~d} x+\left(a x^{3}+2 b y\right) \mathrm{d} y
$$

To verify that the order of taking the second partial derivative is irrelevant, form

$$
\begin{aligned}
& \left(\frac{\partial}{\partial y}\left(\frac{\partial f}{\partial x}\right)_{y}\right)_{x}=\left(\frac{\partial\left(3 a x^{2} y\right)}{\partial y}\right)_{x}=3 a x^{2} \\
& \left(\frac{\partial}{\partial x}\left(\frac{\partial f}{\partial y}\right)_{x}\right)_{y}=\left(\frac{\partial\left(a x^{3}+2 b y\right)}{\partial x}\right)_{y}=3 a x^{2}
\end{aligned}
$$

Now suppose that $z$ is a variable on which $x$ and $y$ depend (for example, $x, y$, and $z$ might correspond to $p, V$, and $T$ ). The following relations then apply:

Relation 1. When $x$ is changed at constant $z$ :

$$
\left(\frac{\partial f}{\partial x}\right)_{z}=\left(\frac{\partial f}{\partial x}\right)_{y}+\left(\frac{\partial f}{\partial y}\right)_{x}\left(\frac{\partial y}{\partial x}\right)_{z}
$$

Relation 2

$$
\left(\frac{\partial y}{\partial x}\right)_{z}=\frac{1}{(\partial x / \partial y)_{z}}
$$

\subsection*{Relation 3}
$$
\left(\frac{\partial x}{\partial y}\right)_{z}=-\left(\frac{\partial x}{\partial z}\right)_{y}\left(\frac{\partial z}{\partial y}\right)_{x}
$$

Combining Relations 2 and 3 results in the Euler chain relation:

$$
\left(\frac{\partial y}{\partial x}\right)_{z}\left(\frac{\partial x}{\partial z}\right)_{y}\left(\frac{\partial z}{\partial y}\right)_{x}=-1 \quad \text { Euler chain relation }
$$

specific heat capacity (more informally, the 'specific heat') of a substance, which is the heat capacity of the sample divided by its mass, usually in grams: $C_{V, s}=C_{V} / m$. The specific heat capacity of water at room temperature is close to $4.2 \mathrm{~J} \mathrm{~K}^{-1} \mathrm{~g}^{-1}$. In general, heat capacities depend on the temperature and decrease at low temperatures. However, over small ranges of temperature at and above room temperature, the variation is quite small and for approximate calculations heat capacities can be treated as almost independent of temperature.

The heat capacity is used to relate a change in internal energy to a change in temperature of a constant-volume system. It follows from eqn 2A. 14 that

$$
\mathrm{d} U=C_{V} \mathrm{~d} T \quad \begin{aligned}
& \text { Internal energy } \\
& \text { change on heating } \\
& \text { [constant volume] }
\end{aligned}
$$

That is, at constant volume, an infinitesimal change in temperature brings about an infinitesimal change in internal energy, and the constant of proportionality is $C_{V}$. If the heat capacity is independent of temperature over the range of temperatures of interest, then

$$
\Delta U=\int_{T_{1}}^{T_{2}} C_{V} \mathrm{~d} T=C_{V} \int_{T_{1}}^{T_{2}} \mathrm{~d} T=C_{V} \overbrace{\left(T_{2}-T_{1}\right)}^{\Delta T}
$$

A measurable change of temperature, $\Delta T$, brings about a measurable change in internal energy, $\Delta U$, with

$$
\Delta U=C_{V} \Delta T \quad \begin{aligned}
& \text { Internal energy } \\
& \text { change on heating } \\
& \text { [constant volume] }
\end{aligned}
$$

Because a change in internal energy can be identified with the heat supplied at constant volume (eqn 2A.11b), the last equation can also be written as


\begin{equation*}
q_{V}=C_{V} \Delta T \tag{2A.16}
\end{equation*}


This relation provides a simple way of measuring the heat capacity of a sample: a measured quantity of energy is transferred as heat to the sample (by electrical heating, for example) under constant volume conditions and the resulting increase in temperature is monitored. The ratio of the energy transferred as heat to the temperature rise it causes $\left(q_{v} / \Delta T\right)$ is\\
the constant-volume heat capacity of the sample. A large heat capacity implies that, for a given quantity of energy transferred as heat, there will be only a small increase in temperature (the sample has a large capacity for heat).

\subsection*{Brief illustration 2A. 6}
Suppose a 55 W electric heater immersed in a gas in a constantvolume adiabatic container was on for 120 s and it was found that the temperature of the gas rose by $5.0^{\circ} \mathrm{C}$ (an increase equivalent to 5.0 K$)$. The heat supplied is $(55 \mathrm{~W}) \times(120 \mathrm{~s})=$ 6.6 kJ (with $1 \mathrm{~J}=1 \mathrm{~W} \mathrm{~s}$ ), so the heat capacity of the sample is

$$
C_{V}=\frac{6.6 \mathrm{~kJ}}{5.0 \mathrm{~K}}=1.3 \mathrm{kJK}^{-1}
$$

\begin{enumerate}
  \setcounter{enumi}{7}
  \item The internal energy increases as the temperature is raised.
  \item The equipartition theorem can be used to estimate the contribution to the internal energy of each classically behaving mode of motion.
  \item The First Law states that the internal energy of an isolated system is constant.
  \item Free expansion (expansion against zero pressure) does no work.
  \item A reversible change is a change that can be reversed by an infinitesimal change in a variable.
  \item To achieve reversible expansion, the external pressure is matched at every stage to the pressure of the system.
  \item The energy transferred as heat at constant volume is equal to the change in internal energy of the system.
  \item Calorimetry is the measurement of heat transactions.
\end{enumerate}

\subsection*{Checklist of equations}
\begin{center}
\begin{tabular}{llll}
\hline
Property & Equation & Comment & Equation number \\
\hline
First Law of thermodynamics & $\Delta U=q+w$ & Convention &  \\
Work of expansion & $\mathrm{d} w=-p_{e x} \mathrm{~d} V$ &  & 2A.2 \\
Work of expansion against a constant external pressure & $w=-p_{e x} \Delta V$ & $p_{\text {ex }}=0$ for free expansion & 2A.5a \\
Reversible work of expansion of a gas & $w=-n R T \ln \left(V_{f} / V_{\mathrm{i}}\right)$ & Isothermal, perfect gas & 2A.9 \\
Internal energy change & $\Delta U=q_{V}$ & Constant volume, no other forms of work & 2A.11b \\
Electrical heating & $q=I t \Delta \phi$ &  & 2A.13 \\
Heat capacity at constant volume & $C_{V}=(\partial U / \partial T)_{V}$ & Definition & 2 A .14 \\
\hline
\end{tabular}
\end{center}

\section*{TOPIC 2B Enthalpy}
\subsection*{Why do you need to know this material?}
The concept of enthalpy is central to many thermodynamic discussions about processes, such as physical transformations and chemical reactions taking place under conditions of constant pressure.

\subsection*{What is the key idea?}
A change in enthalpy is equal to the energy transferred as heat at constant pressure.

\subsection*{What do you need to know already?}
This Topic makes use of the discussion of internal energy (Topic 2A) and draws on some aspects of perfect gases (Topic 1A).

The change in internal energy is not equal to the energy transferred as heat when the system is free to change its volume, such as when it is able to expand or contract under conditions of constant pressure. Under these circumstances some of the energy supplied as heat to the system is returned to the surroundings as expansion work (Fig. 2B.1), so $\mathrm{d} U$ is less than $\mathrm{d} q$. In this case the energy supplied as heat at constant pressure is equal to the change in another thermodynamic property of the system, the 'enthalpy'.\\
\includegraphics[max width=\textwidth, center]{2025_02_27_d4b95d9718f0b9e2e6b9g-078}

Figure 2B. 1 When a system is subjected to constant pressure and is free to change its volume, some of the energy supplied as heat may escape back into the surroundings as work. In such a case, the change in internal energy is smaller than the energy supplied as heat.

\subsection*{2B.1 The definition of enthalpy}
The enthalpy, $H$, is defined as


\begin{equation*}
H=U+p V \tag{2B.1}
\end{equation*}


Enthalpy [definition]\\
where $p$ is the pressure of the system and $V$ is its volume. Because $U, p$, and $V$ are all state functions, the enthalpy is a state function too. As is true of any state function, the change in enthalpy, $\Delta H$, between any pair of initial and final states is independent of the path between them.

\subsection*{(a) Enthalpy change and heat transfer}
An important consequence of the definition of enthalpy in eqn 2B. 1 is that it can be shown that the change in enthalpy is equal to the energy supplied as heat under conditions of constant pressure.

\subsection*{How is that done? 2B. 1 Deriving the relation between}
 enthalpy change and heat transfer at constant pressureIn a typical thermodynamic derivation, as here, a common way to proceed is to introduce successive definitions of the quantities of interest and then apply the appropriate constraints.

Step 1 Write an expression for $H+\mathrm{d} H$ in terms of the definition of $H$\\
For a general infinitesimal change in the state of the system, $U$ changes to $U+\mathrm{d} U, p$ changes to $p+\mathrm{d} p$, and $V$ changes to $V+\mathrm{d} V$, so from the definition in eqn 2B.1, $H$ changes by $\mathrm{d} H$ to

$$
\begin{aligned}
H+\mathrm{d} H & =(U+\mathrm{d} U)+(p+\mathrm{d} p)(V+\mathrm{d} V) \\
& =U+\mathrm{d} U+p V+p \mathrm{~d} V+V \mathrm{~d} p+\mathrm{d} p \mathrm{~d} V
\end{aligned}
$$

The last term is the product of two infinitesimally small quantities and can be neglected. Now recognize that $U+p V=H$ on the right (in blue), so

$$
H+\mathrm{d} H=H+\mathrm{d} U+p \mathrm{~d} V+V \mathrm{~d} p
$$

and hence

$$
\mathrm{d} H=\mathrm{d} U+p \mathrm{~d} V+V \mathrm{~d} p
$$

\subsection*{Step 2 Introduce the definition of $\mathrm{d} U$}
Because $\mathrm{d} U=\mathrm{d} q+\mathrm{d} w$ this expression becomes

$$
\mathrm{d} H=\mathrm{d} q+\mathrm{d} w+p \mathrm{~d} V+V \mathrm{~d} p
$$

\subsection*{Step 3 Apply the appropriate constraints}
If the system is in mechanical equilibrium with its surroundings at a pressure $p$ and does only expansion work, then $\mathrm{d} w=-p \mathrm{~d} V$, which cancels the other $p \mathrm{~d} V$ term, leaving

$$
\mathrm{d} H=\mathrm{d} q+V \mathrm{~d} p
$$

At constant pressure, $\mathrm{d} p=0$, so\\
$\mathrm{d} H=\mathrm{d} q$ (at constant pressure, no additional work)\\
The constraint of constant pressure is denoted by a $p$, so this equation can be written

\[
\begin{array}{ll}
\mathrm{d} H=\mathrm{d} q_{p} & \begin{array}{l}
\text { Heat transferred at constant } \\
\text { pressure [infinitesimal change] }
\end{array} \tag{2B.2a}
\end{array}
\]

This equation states that, provided there is no additional (non-expansion) work done, the change in enthalpy is equal to the energy supplied as heat at constant pressure.

\subsection*{Step 4 Evaluate $\Delta H$ by integration}
For a measurable change between states i and $f$ along a path at constant pressure, the preceding expression is integrated as follows

$$
\overbrace{\int_{i}^{f}}^{H_{f}-H_{i}}=\overbrace{\int_{i}^{f} \mathrm{~d} q}^{q_{p}}
$$

Note that the integral over $\mathrm{d} q$ is not written as $\Delta q$ because $q$, unlike $H$, is not a state function and $q_{\mathrm{f}}-q_{\mathrm{i}}$ is meaningless. The final result is\\
\includegraphics[max width=\textwidth, center]{2025_02_27_d4b95d9718f0b9e2e6b9g-079}

\subsection*{Brief illustration 2B. 1}
Water is heated to boiling under a pressure of 1.0 atm . When an electric current of 0.50 A from a 12 V supply is passed for 300 s through a resistance in thermal contact with the water, it is found that 0.798 g of water is vaporized. The enthalpy change is

$$
\begin{aligned}
\Delta H & =q_{p}=I t \Delta \phi=(0.50 \mathrm{~A}) \times(300 \mathrm{~s}) \times(12 \mathrm{~V}) \\
& =0.50 \times 300 \mathrm{~J} \times 12
\end{aligned}
$$

where $1 \mathrm{AVs}=1 \mathrm{~J}$. Because 0.798 g of water is $(0.798 \mathrm{~g})$ / $\left(18.02 \mathrm{~g} \mathrm{~mol}^{-1}\right)=(0.798 / 18.02) \mathrm{mol} \mathrm{H}_{2} \mathrm{O}$, the enthalpy of vaporization per mole of $\mathrm{H}_{2} \mathrm{O}$ is

$$
\Delta H_{\mathrm{m}}=\frac{0.50 \times 12 \times 300 \mathrm{~J}}{(0.798 / 18.02) \mathrm{mol}}=+41 \mathrm{~kJ} \mathrm{~mol}^{-1}
$$

\subsection*{(b) Calorimetry}
An enthalpy change can be measured calorimetrically by monitoring the temperature change that accompanies a physical or chemical change at constant pressure. A calorimeter for\\
\includegraphics[max width=\textwidth, center]{2025_02_27_d4b95d9718f0b9e2e6b9g-079(1)}

Figure 2B. 2 A constant-pressure flame calorimeter consists of this component immersed in a stirred water bath. Combustion occurs as a known amount of reactant is passed through to fuel the flame, and the rise of temperature is monitored.\\
studying processes at constant pressure is called an isobaric calorimeter. A simple example is a thermally insulated vessel open to the atmosphere: the energy released as heat in the reaction is monitored by measuring the change in temperature of the contents. For a combustion reaction an adiabatic flame calorimeter may be used to measure $\Delta T$ when a given amount of substance burns in a supply of oxygen (Fig. 2B.2). The most sophisticated way to measure enthalpy changes, however, is to use a differential scanning calorimeter (DSC), as explained in Topic 2C. Changes in enthalpy and internal energy may also be measured by non-calorimetric methods (Topic 6C).

One route to $\Delta H$ is to measure the internal energy change by using a bomb calorimeter (Topic 2 A ), and then to convert $\Delta U$ to $\Delta H$. Because solids and liquids have small molar volumes, for them $p V_{\mathrm{m}}$ is so small that the molar enthalpy and molar internal energy are almost identical ( $\left.H_{\mathrm{m}}=U_{\mathrm{m}}+p V_{\mathrm{m}} \approx U_{\mathrm{m}}\right)$. Consequently, if a process involves only solids or liquids, the values of $\Delta H$ and $\Delta U$ are almost identical. Physically, such processes are accompanied by a very small change in volume; the system does negligible work on the surroundings when the process occurs, so the energy supplied as heat stays entirely within the system.

\subsection*{Example 2B.1 Relating $\Delta H$ and $\Delta U$}
The change in molar internal energy when $\mathrm{CaCO}_{3}(\mathrm{~s})$ as calcite converts to its polymorph aragonite, is $+0.21 \mathrm{~kJ} \mathrm{~mol}^{-1}$. Calculate the difference between the molar enthalpy and internal energy changes when the pressure is 1.0 bar. The mass densities of the polymorphs are $2.71 \mathrm{~g} \mathrm{~cm}^{-3}$ (calcite) and $2.93 \mathrm{~g} \mathrm{~cm}^{-3}$ (aragonite).

Collect your thoughts The starting point for the calculation is the relation between the enthalpy of a substance and its internal energy (eqn 2B.1). You need to express the difference between the two quantities in terms of the pressure and the difference of their molar volumes. The latter can be calculated\\
from their molar masses, $M$, and their mass densities, $\rho$, by using $\rho=M / V_{\mathrm{m}}$.

The solution The change in enthalpy when the transition occurs is

$$
\begin{aligned}
\Delta H_{\mathrm{m}} & =H_{\mathrm{m}}(\text { aragonite })-H_{\mathrm{m}}(\text { calcite }) \\
& =\left\{U_{\mathrm{m}}(\mathrm{a})+p V_{\mathrm{m}}(\mathrm{a})\right\}-\left\{U_{\mathrm{m}}(\mathrm{c})+p V_{\mathrm{m}}(\mathrm{c})\right\} \\
& =\Delta U_{\mathrm{m}}+p\left\{V_{\mathrm{m}}(\mathrm{a})-V_{\mathrm{m}}(\mathrm{c})\right\}
\end{aligned}
$$

where a denotes aragonite and calcite. It follows by substituting $V_{\mathrm{m}}=M / \rho$ that

$$
\Delta H_{\mathrm{m}}-\Delta U_{\mathrm{m}}=p M\left(\frac{1}{\rho(\mathrm{a})}-\frac{1}{\rho(\mathrm{c})}\right)
$$

Substitution of the data, using $M=100.09 \mathrm{~g} \mathrm{~mol}^{-1}$, gives

$$
\begin{aligned}
\Delta H_{\mathrm{m}}-\Delta U_{\mathrm{m}}= & \left(1.0 \times 10^{5} \mathrm{~Pa}\right) \times\left(100.09 \mathrm{~g} \mathrm{~mol}^{-1}\right) \\
& \times\left(\frac{1}{2.93 \mathrm{~g} \mathrm{~cm}^{-3}}-\frac{1}{2.71 \mathrm{gcm}^{-3}}\right) \\
= & -2.8 \times 10^{5} \mathrm{Pacm}^{3} \mathrm{~mol}^{-1}=-0.28 \mathrm{Pam}^{3} \mathrm{~mol}^{-1}
\end{aligned}
$$

Hence (because $1 \mathrm{~Pa} \mathrm{~m}=1 \mathrm{~J}$ ), $\Delta H_{\mathrm{m}}-\Delta U_{\mathrm{m}}=-0.28 \mathrm{~J} \mathrm{~mol}^{-1}$, which is only 0.1 per cent of the value of $\Delta U_{\mathrm{m}}$.

Comment. It is usually justifiable to ignore the difference between the molar enthalpy and internal energy of condensed phases except at very high pressures when $p \Delta V_{\mathrm{m}}$ is no longer negligible.

Self-test 2B. 1 Calculate the difference between $\Delta H$ and $\Delta U$ when $1.0 \mathrm{~mol} \mathrm{Sn}\left(\mathrm{s}\right.$, grey) of density $5.75 \mathrm{~g} \mathrm{~cm}^{-3}$ changes to Sn (s, white) of density $7.31 \mathrm{~g} \mathrm{~cm}^{-3}$ at 10.0 bar.\\
\includegraphics[max width=\textwidth, center]{2025_02_27_d4b95d9718f0b9e2e6b9g-080}

In contrast to processes involving condensed phases, the values of the changes in internal energy and enthalpy might differ significantly for processes involving gases. The enthalpy of a perfect gas is related to its internal energy by using $p V=$ $n R T$ in the definition of $H$ :


\begin{equation*}
H=U+p V=U+n R T \tag{2B.3}
\end{equation*}


This relation implies that the change of enthalpy in a reaction that produces or consumes gas under isothermal conditions is

\[
\Delta H=\Delta U+\Delta n_{\mathrm{g}} R T \quad \begin{align*}
& \text { Relation between } \Delta H \text { and } \Delta U  \tag{2B.4}\\
& \text { [isothermal process, perfect gas] }
\end{align*}
\]

where $\Delta n_{\mathrm{g}}$ is the change in the amount of gas molecules in the reaction. For molar quantities, replace $\Delta n_{\mathrm{g}}$ by $\Delta \nu_{\mathrm{g}}$.

\subsection*{Brief illustration 2B. 2}
In the reaction $2 \mathrm{H}_{2}(\mathrm{~g})+\mathrm{O}_{2}(\mathrm{~g}) \rightarrow 2 \mathrm{H}_{2} \mathrm{O}(\mathrm{l}), 3 \mathrm{~mol}$ of gas-phase molecules are replaced by 2 mol of liquid-phase molecules,\\
so $\Delta n_{\mathrm{g}}=-3 \mathrm{~mol}$ and $\Delta \nu_{\mathrm{g}}=-3$. Therefore, at 298 K , when $R T=$ $2.5 \mathrm{~kJ} \mathrm{~mol}^{-1}$, the enthalpy and internal energy changes taking place in the system are related by

$$
\Delta H_{\mathrm{m}}-\Delta U_{\mathrm{m}}=(-3) \times R T \approx-7.5 \mathrm{~kJ} \mathrm{~mol}^{-1}
$$

Note that the difference is expressed in kilojoules, not joules as in Example 2B.1. The enthalpy change is smaller than the change in internal energy because, although energy escapes from the system as heat when the reaction occurs, the system contracts as the liquid is formed, so energy is restored to it as work from the surroundings.

\subsection*{2B.2 The variation of enthalpy with temperature}
The enthalpy of a substance increases as its temperature is raised. The reason is the same as for the internal energy: molecules are excited to states of higher energy so their total energy increases. The relation between the increase in enthalpy and the increase in temperature depends on the conditions (e.g. whether the pressure or the volume is constant).

\subsection*{(a) Heat capacity at constant pressure}
The most frequently encountered condition in chemistry is constant pressure. The slope of the tangent to a plot of enthalpy against temperature at constant pressure is called the heat capacity at constant pressure (or isobaric heat capacity), $C_{p}$, at a given temperature (Fig. 2B.3). More formally:

\[
C_{p}=\left(\frac{\partial H}{\partial T}\right)_{p} \quad \begin{align*}
& \text { Heat capacity at constant pressure }  \tag{2B.5}\\
& {[\text { definition }]}
\end{align*}
\]

\begin{center}
\includegraphics[max width=\textwidth]{2025_02_27_d4b95d9718f0b9e2e6b9g-080(1)}
\end{center}

Figure 2B. 3 The constant-pressure heat capacity at a particular temperature is the slope of the tangent to a curve of the enthalpy of a system plotted against temperature (at constant pressure). For gases, at a given temperature the slope of enthalpy versus temperature is steeper than that of internal energy versus temperature, and $C_{p, m}$ is larger than $C_{V, m}$.

The heat capacity at constant pressure is the analogue of the heat capacity at constant volume (Topic 2A) and is an extensive property. The molar heat capacity at constant pressure, $C_{p, m}$, is the heat capacity per mole of substance; it is an intensive property.

The heat capacity at constant pressure relates the change in enthalpy to a change in temperature. For infinitesimal changes of temperature, eqn 2B. 5 implies that


\begin{equation*}
\mathrm{d} H=C_{p} \mathrm{~d} T \text { (at constant pressure) } \tag{2B.6a}
\end{equation*}


If the heat capacity is constant over the range of temperatures of interest, then for a measurable increase in temperature

$$
\Delta H=\int_{T_{1}}^{T_{2}} C_{p} \mathrm{~d} T=C_{p} \int_{T_{1}}^{T_{2}} \mathrm{~d} T=C_{p} \overbrace{\left(T_{2}-T_{1}\right)}^{\Delta T}
$$

which can be summarized as


\begin{equation*}
\Delta H=C_{p} \Delta T(\text { at constant pressure }) \tag{2B.6b}
\end{equation*}


Because a change in enthalpy can be equated to the energy supplied as heat at constant pressure, the practical form of this equation is


\begin{equation*}
q_{p}=C_{p} \Delta T \tag{2B.7}
\end{equation*}


This expression shows how to measure the constant-pressure heat capacity of a sample: a measured quantity of energy is supplied as heat under conditions of constant pressure (as in a sample exposed to the atmosphere and free to expand), and the temperature rise is monitored.

The variation of heat capacity with temperature can sometimes be ignored if the temperature range is small; this is an excellent approximation for a monatomic perfect gas (for instance, one of the noble gases at low pressure). However, when it is necessary to take the variation into account for other substances, a convenient approximate empirical expression is


\begin{equation*}
C_{p, \mathrm{~m}}=a+b T+\frac{c}{T^{2}} \tag{2B.8}
\end{equation*}


The empirical parameters $a, b$, and $c$ are independent of temperature (Table 2B.1) and are found by fitting this expression to experimental data.

Table 2B. 1 Temperature variation of molar heat capacities, $C_{p, \mathrm{~m}} /\left(\mathrm{JK}^{-1} \mathrm{~mol}^{-1}\right)=a+b T+c / T^{2 *}$

\begin{center}
\begin{tabular}{llll}
\hline
 & $\boldsymbol{a}$ & $\boldsymbol{b} /\left(10^{-3} \mathrm{~K}^{-1}\right)$ & $\boldsymbol{c} /\left(10^{5} \mathrm{~K}^{2}\right)$ \\
\hline
$\mathrm{C}(\mathrm{s}$, graphite $)$ & 16.86 & 4.77 & -8.54 \\
$\mathrm{CO}_{2}(\mathrm{~g})$ & 44.22 & 8.79 & -8.62 \\
$\mathrm{H}_{2} \mathrm{O}(\mathrm{l})$ & 75.29 & 0 & 0 \\
$\mathrm{~N}_{2}(\mathrm{~g})$ & 28.58 & 3.77 & -0.50 \\
\end{tabular}
\end{center}

\footnotetext{\begin{itemize}
  \item More values are given in the Resource section
\end{itemize}
}Example 2B. 2\\
Evaluating an increase in enthalpy with temperature

What is the change in molar enthalpy of $\mathrm{N}_{2}$ when it is heated from $25^{\circ} \mathrm{C}$ to $100^{\circ} \mathrm{C}$ ? Use the heat capacity information in Table 2B.1.

Collect your thoughts The heat capacity of $\mathrm{N}_{2}$ changes with temperature significantly in this range, so you cannot use eqn 2B.6b (which assumes that the heat capacity of the substance is constant). Therefore, use eqn 2B.6a, substitute eqn 2B. 8 for the temperature dependence of the heat capacity, and integrate the resulting expression from $25^{\circ} \mathrm{C}(298 \mathrm{~K})$ to $100^{\circ} \mathrm{C}(373 \mathrm{~K})$.

The solution For convenience, denote the two temperatures $T_{1}$ ( 298 K ) and $T_{2}(373 \mathrm{~K})$. The required relation is

$$
\int_{H_{\mathrm{m}}\left(T_{1}\right)}^{H_{\mathrm{m}}\left(T_{2}\right)} \mathrm{d} H_{\mathrm{m}}=\int_{T_{1}}^{T_{2}}\left(a+b T+\frac{c}{T^{2}}\right) \mathrm{d} T
$$

By using Integral A. 1 in the Resource section for each term, it follows that

$$
H_{\mathrm{m}}\left(T_{2}\right)-H_{\mathrm{m}}\left(T_{1}\right)=a\left(T_{2}-T_{1}\right)+\frac{1}{2} b\left(T_{2}^{2}-T_{1}^{2}\right)-c\left(\frac{1}{T_{2}}-\frac{1}{T_{1}}\right)
$$

Substitution of the numerical data results in

$$
H_{\mathrm{m}}(373 \mathrm{~K})=H_{\mathrm{m}}(298 \mathrm{~K})+2.20 \mathrm{~kJ} \mathrm{~mol}^{-1}
$$

Comment. If a constant heat capacity of $29.14 \mathrm{~J} \mathrm{~K}^{-1} \mathrm{~mol}^{-1}$ (the value given by eqn 2 B .8 for $T=298 \mathrm{~K}$ ) had been assumed, then the difference between the two enthalpies would have been calculated as $2.19 \mathrm{~kJ} \mathrm{~mol}^{-1}$, only slightly different from the more accurate value.

Self-test 2B. 2 At very low temperatures the heat capacity of a solid is proportional to $T^{3}$, and $C_{p, \mathrm{~m}}=a T^{3}$. What is the change in enthalpy of such a substance when it is heated from 0 to a temperature $T$ (with $T$ close to 0 )?

$$
{ }_{\forall} L^{p} \frac{\downarrow}{T}={ }^{\mathrm{w}} H \nabla:{\iota \partial \mu s u_{V}}
$$

\subsection*{(b) The relation between heat capacities}
Most systems expand when heated at constant pressure. Such systems do work on the surroundings and therefore some of the energy supplied to them as heat escapes back to the surroundings as work. As a result, the temperature of the system rises less than when the heating occurs at constant volume. A smaller increase in temperature implies a larger heat capacity, so in most cases the heat capacity at constant pressure of a system is larger than its heat capacity at constant volume. As shown in Topic 2D, there is a simple relation between the two heat capacities of a perfect gas:

$$
\begin{array}{ll}
C_{p}-C_{V}=n R & \begin{array}{l}
\text { Relation between heat capacities } \\
\text { [perfect gas] }
\end{array}
\end{array}
$$

(2B.9)

It follows that the molar heat capacity of a perfect gas is about $8 \mathrm{~J} \mathrm{~K}^{-1} \mathrm{~mol}^{-1}$ larger at constant pressure than at constant volume. Because the molar constant-volume heat capacity of a monatomic gas is about $\frac{3}{2} R=12 \mathrm{~J} \mathrm{~K}^{-1} \mathrm{~mol}^{-1}$ (Topic 2A),\\
the difference is highly significant and must be taken into account. The two heat capacities are typically very similar for condensed phases, and for them the difference can normally be ignored.

\subsection*{Checklist of concepts}
\begin{enumerate}
  \item Energy transferred as heat at constant pressure is equal to the change in enthalpy of a system.2. Enthalpy changes can be measured in a constant-pressure calorimeter.
  \item The heat capacity at constant pressure is equal to the slope of enthalpy with temperature.
\end{enumerate}

\subsection*{Checklist of equations}
\begin{center}
\begin{tabular}{llll}
\hline
Property & Equation & Comment & Equation number \\
\hline
Enthalpy & $H=U+p V$ & Definition & 2B.1 \\
Heat transfer at constant pressure & $\mathrm{d} H=\mathrm{d} q_{p}$, & No additional work &  \\
 & $\Delta H=q_{p}$ &  & 2B.2 \\
\begin{tabular}{ll}
Relation between $\Delta H$ and $\Delta U$ at a volumes of the participating condensed & phases are negligible \\
\end{tabular} & 2B.4 &  &  \\
temperature $T$ & $\Delta H=\Delta U+\Delta n_{\mathrm{g}} R T$ & Definition & 2B.5 \\
Heat capacity at constant pressure & $C_{p}=(\partial H / \partial T)_{p}$ & $C_{p}-C_{V}=n R$ & Perfect gas \\
Relation between heat capacities &  &  & 2B.9 \\
\hline
\end{tabular}
\end{center}

\section*{TOPIC 2C Thermochemistry}
\subsection*{Why do you need to know this material?}
Thermochemistry is one of the principal applications of thermodynamics in chemistry. Thermochemical data provide a way of assessing the heat output of chemical reactions, including those involved with the combustion of fuels and the consumption of foods. The data are also used widely in other chemical applications of thermodynamics.

\subsection*{What is the key idea?}
Reaction enthalpies can be combined to provide data on other reactions of interest.

\subsection*{> What do you need to know already?}
You need to be aware of the definition of enthalpy and its status as a state function (Topic 2B). The material on temperature dependence of reaction enthalpies makes use of information about heat capacities (Topic 2B).

The study of the energy transferred as heat during the course of chemical reactions is called thermochemistry. Thermochemistry is a branch of thermodynamics because a reaction vessel and its contents form a system, and chemical reactions result in the exchange of energy between the system and the surroundings. Thus calorimetry can be used to measure the energy supplied or discarded as heat by a reaction, with $q$ identified with a change in internal energy if the reaction occurs at constant volume (Topic 2A) or with a change in enthalpy if the reaction occurs at constant pressure (Topic 2B). Conversely, if $\Delta U$ or $\Delta H$ for a reaction is known, it is possible to predict the heat the reaction can produce.

As pointed out in Topic 2A, a process that releases energy as heat is classified as exothermic, and one that absorbs energy as heat is classified as endothermic. Because the release of heat into the surroundings at constant pressure signifies a decrease in the enthalpy of a system, it follows that an exothermic process is one for which $\Delta H<0$; such a process is exenthalpic. Conversely, because the absorption of heat from the surroundings results in an increase in enthalpy, an endothermic process has $\Delta H>0$; such a process is endenthalpic:\\
exothermic (exenthalpic) process: $\Delta H<0$\\
endothermic (endenthalpic) process: $\Delta H>0$

\subsection*{2C.1 Standard enthalpy changes}
Changes in enthalpy are normally reported for processes taking place under a set of standard conditions. The standard enthalpy change, $\Delta H^{\ominus}$, is the change in enthalpy for a process in which the initial and final substances are in their standard states:

The standard state of a substance at a specified temperature is its pure form at 1 bar.

For example, the standard state of liquid ethanol at 298 K is pure liquid ethanol at 298 K and 1 bar; the standard state of solid iron at 500 K is pure iron at 500 K and 1 bar. The definition of standard state is more sophisticated for solutions (Topic 5E). The standard enthalpy change for a reaction or a physical process is the difference in enthalpy between the products in their standard states and the reactants in their standard states, all at the same specified temperature.

An example of a standard enthalpy change is the standard enthalpy of vaporization, $\Delta_{\text {vap }} H^{\ominus}$, which is the enthalpy change per mole of molecules when a pure liquid at 1 bar vaporizes to a gas at 1 bar , as in

$$
\mathrm{H}_{2} \mathrm{O}(\mathrm{l}) \rightarrow \mathrm{H}_{2} \mathrm{O}(\mathrm{~g}) \quad \Delta_{\text {vap }} H^{\ominus}(373 \mathrm{~K})=+40.66 \mathrm{~kJ} \mathrm{~mol}^{-1}
$$

As implied by the examples, standard enthalpies may be reported for any temperature. However, the conventional temperature for reporting thermodynamic data is 298.15 K . Unless otherwise mentioned or indicated by attaching the temperature to $\Delta H^{\ominus}$, all thermodynamic data in this text are for this conventional temperature.

A note on good practice The attachment of the name of the transition to the symbol $\Delta$, as in $\Delta_{\text {vap }} H$, is the current convention. However, the older convention, $\Delta H_{\text {vap }}$, is still widely used. The current convention is more logical because the subscript identifies the type of change, not the physical observable related to the change.

\subsection*{(a) Enthalpies of physical change}
The standard molar enthalpy change that accompanies a change of physical state is called the standard enthalpy of transition and is denoted $\Delta_{\text {trs }} H^{\ominus}$ (Table 2C.1). The standard enthalpy of vaporization, $\Delta_{\text {vap }} H^{\ominus}$, is one example. Another is

Table 2C. 1 Standard enthalpies of fusion and vaporization at the transition temperature*

\begin{center}
\begin{tabular}{lcccc}
\hline
 & $\boldsymbol{T}_{\mathbf{f}} / \mathbf{K}$ & Fusion & $\boldsymbol{T}_{\mathbf{b}} / \mathbf{K}$ & Vaporization \\
\hline
Ar & 83.81 & 1.188 & 87.29 & 6.506 \\
$\mathrm{C}_{6} \mathrm{H}_{6}$ & 278.61 & 10.59 & 353.2 & 30.8 \\
$\mathrm{H}_{2} \mathrm{O}$ & 273.15 & 6.008 & 373.15 & \begin{tabular}{c}
$40.656(44.016$ \\
at 298 K$)$ \\
\end{tabular} \\
He & 3.5 & 0.021 & 4.22 & 0.084 \\
\hline
\end{tabular}
\end{center}

\begin{itemize}
  \item More values are given in the Resource section.\\
the standard enthalpy of fusion, $\Delta_{\text {fus }} H^{\ominus}$, the standard molar enthalpy change accompanying the conversion of a solid to a liquid, as in
\end{itemize}

$$
\mathrm{H}_{2} \mathrm{O}(\mathrm{~s}) \rightarrow \mathrm{H}_{2} \mathrm{O}(\mathrm{l}) \quad \Delta_{\text {fus }} H^{\ominus}(273 \mathrm{~K})=+6.01 \mathrm{~kJ} \mathrm{~mol}^{-1}
$$

As in this case, it is sometimes convenient to know the standard molar enthalpy change at the transition temperature as well as at the conventional temperature of 298 K . The different types of enthalpy changes encountered in thermochemistry are summarized in Table 2C.2.

Because enthalpy is a state function, a change in enthalpy is independent of the path between the two states. This feature is of great importance in thermochemistry, because it implies that the same value of $\Delta H^{\ominus}$ will be obtained however the change is brought about between specified initial and final states. For example, the conversion of a solid to a vapour can be pictured either as occurring by sublimation (the direct conversion from solid to vapour),

$$
\mathrm{H}_{2} \mathrm{O}(\mathrm{~s}) \rightarrow \mathrm{H}_{2} \mathrm{O}(\mathrm{~g}) \quad \Delta_{\text {sub }} H^{\ominus}
$$

Table 2C. 2 Enthalpies of reaction and transition

\begin{center}
\begin{tabular}{lll}
\hline
Transition & Process & Symbol $^{*}$ \\
\hline
Transition & Phase $\alpha \rightarrow$ phase $\beta$ & $\Delta_{\text {trs }} H$ \\
Fusion & $\mathrm{s} \rightarrow 1$ & $\Delta_{\text {fus }} H$ \\
Vaporization & $\mathrm{l} \rightarrow \mathrm{g}$ & $\Delta_{\text {vap }} H$ \\
Sublimation & $\mathrm{s} \rightarrow \mathrm{g}$ & $\Delta_{\text {sub }} H$ \\
Mixing & Pure $\rightarrow$ mixture & $\Delta_{\text {mix }} H$ \\
Solution & Solute $\rightarrow$ solution & $\Delta_{\text {sol }} H$ \\
Hydration & $\mathrm{X}^{ \pm}(\mathrm{g}) \rightarrow \mathrm{X}^{ \pm}(\mathrm{aq})$ & $\Delta_{\text {hyd }} H$ \\
Atomization & Species $(\mathrm{s}, \mathrm{l}, \mathrm{g}) \rightarrow$ atoms $(\mathrm{g})$ & $\Delta_{\mathrm{at}} H$ \\
Ionization & $\mathrm{X}(\mathrm{g}) \rightarrow \mathrm{X}^{+}(\mathrm{g})+\mathrm{e}^{-}(\mathrm{g})$ & $\Delta_{\text {ion }} H$ \\
Electron gain & $\mathrm{X}(\mathrm{g})+\mathrm{e}^{-}(\mathrm{g}) \rightarrow \mathrm{X}^{-}(\mathrm{g})$ & $\Delta_{\mathrm{eg}} H$ \\
Reaction & Reactants $\rightarrow$ products & $\Delta_{\mathrm{r}} H$ \\
Combustion & Compound $(\mathrm{s}, 1, \mathrm{~g})+\mathrm{O}_{2}(\mathrm{~g}) \rightarrow \mathrm{CO}_{2}(\mathrm{~g})+$ & $\Delta_{\mathrm{c}} H$ \\
 & $\quad \mathrm{H}, \mathrm{O}(\mathrm{l}, \mathrm{g})$ &  \\
Formation & Elements $\rightarrow$ compound & $\Delta_{\mathrm{f}} H$ \\
Activation & Reactants $\rightarrow$ activated complex & $\Delta^{\ddagger} H$ \\
\hline
\end{tabular}
\end{center}

*IUPAC recommendations. In common usage, the process subscript is often attached to $\Delta H$, as in $\Delta H_{\text {trs }}$ and $\Delta H_{\mathrm{f}}$. All are molar quantities.\\
or as occurring in two steps, first fusion (melting) and then vaporization of the resulting liquid:

$$
\begin{array}{ll} 
& \mathrm{H}_{2} \mathrm{O}(\mathrm{~s}) \rightarrow \mathrm{H}_{2} \mathrm{O}(\mathrm{l}) \\
& \Delta_{\text {fus }} H^{\ominus} \\
\mathrm{H}_{2} \mathrm{O}(\mathrm{l}) \rightarrow \mathrm{H}_{2} \mathrm{O}(\mathrm{~g}) & \Delta_{\text {vap }} H^{\ominus} \\
\text { Overall: } & \mathrm{H}_{2} \mathrm{O}(\mathrm{~s}) \rightarrow \mathrm{H}_{2} \mathrm{O}(\mathrm{~g}) \\
\Delta_{\text {fus }} H^{\ominus}+\Delta_{\text {vap }} H^{\ominus}
\end{array}
$$

Because the overall result of the indirect path is the same as that of the direct path, the overall enthalpy change is the same in each case (1), and (for processes occurring at the same temperature)

It follows that, because all enthalpies of fusion are positive, the enthalpy of sublimation of a substance is greater than its enthalpy of vaporization (at a given temperature).

Another consequence of $H$ being a state function is that the standard enthalpy change of a forward process is the negative of its reverse (2):


\begin{equation*}
\Delta H^{\ominus}(\mathrm{A} \rightarrow \mathrm{~B})=-\Delta H^{\ominus}(\mathrm{A} \leftarrow \mathrm{~B}) \tag{2C.2}
\end{equation*}


For instance, because the enthalpy of vaporization of water is $+44 \mathrm{~kJ} \mathrm{~mol}^{-1}$ at 298 K , the enthalpy of condensation of water vapour at that temperature is $-44 \mathrm{~kJ} \mathrm{~mol}^{-1}$.

\subsection*{(b) Enthalpies of chemical change}
There are two ways of reporting the change in enthalpy that accompanies a chemical reaction. One is to write the thermochemical equation, a combination of a chemical equation and the corresponding change in standard enthalpy:

$$
\mathrm{CH}_{4}(\mathrm{~g})+2 \mathrm{O}_{2}(\mathrm{~g}) \rightarrow \mathrm{CO}_{2}(\mathrm{~g})+2 \mathrm{H}_{2} \mathrm{O}(\mathrm{~g}) \Delta H^{\ominus}=-890 \mathrm{~kJ}
$$

$\Delta H^{\ominus}$ is the change in enthalpy when reactants in their standard states change to products in their standard states:

Pure, separate reactants in their standard states $\rightarrow$ pure, separate products in their standard states

Except in the case of ionic reactions in solution, the enthalpy changes accompanying mixing and separation are insignificant in comparison with the contribution from the reaction itself. For the combustion of methane, the standard value refers to the reaction in which $1 \mathrm{~mol} \mathrm{CH}_{4}$ in the form of pure methane gas at 1 bar reacts completely with $2 \mathrm{~mol} \mathrm{O}_{2}$ in the form of pure oxygen gas to produce $1 \mathrm{~mol} \mathrm{CO}_{2}$ as pure carbon dioxide at 1 bar and $2 \mathrm{~mol} \mathrm{H}_{2} \mathrm{O}$ as pure liquid water at 1 bar ; the numerical value quoted is for the reaction at 298.15 K .

Alternatively, the chemical equation is written and the standard reaction enthalpy, $\Delta_{\mathrm{r}} H^{\ominus}$ (or 'standard enthalpy of reaction') reported. Thus, for the combustion of methane at 298 K, write

$$
\mathrm{CH}_{4}(\mathrm{~g})+2 \mathrm{O}_{2}(\mathrm{~g}) \rightarrow \mathrm{CO}_{2}(\mathrm{~g})+2 \mathrm{H}_{2} \mathrm{O}(\mathrm{l}) \Delta_{\mathrm{r}} H^{\ominus}=-890 \mathrm{~kJ} \mathrm{~mol}^{-1}
$$

For a reaction of the form $2 \mathrm{~A}+\mathrm{B} \rightarrow 3 \mathrm{C}+\mathrm{D}$ the standard reaction enthalpy would be

$$
\Delta_{\mathrm{r}} H^{\ominus}=\left\{3 H_{\mathrm{m}}^{\ominus}(\mathrm{C})+H_{\mathrm{m}}^{\ominus}(\mathrm{D})\right\}-\left\{2 H_{\mathrm{m}}^{\ominus}(\mathrm{A})+H_{\mathrm{m}}^{\ominus}(\mathrm{B})\right\}
$$

where $H_{\mathrm{m}}^{\ominus}(\mathrm{J})$ is the standard molar enthalpy of species J at the temperature of interest. Note how the 'per mole' of $\Delta_{r} H^{\ominus}$ comes directly from the fact that molar enthalpies appear in this expression. The 'per mole' is interpreted by noting the stoichiometric coefficients in the chemical equation. In this case, 'per mole' in $\Delta_{\mathrm{r}} H^{\ominus}$ means 'per $2 \mathrm{~mol} \mathrm{A'}, \mathrm{'per} \mathrm{~mol} \mathrm{B'}, \mathrm{'per} 3 \mathrm{~mol} \mathrm{C'}$ ', or 'per mol D'. In general,

\[
\Delta_{\mathrm{r}} H^{\ominus}=\sum_{\text {Products }} v H_{\mathrm{m}}^{\ominus}-\sum_{\text {Reactants }} v H_{\mathrm{m}}^{\ominus} \quad \begin{align*}
& \begin{array}{l}
\text { Standard reaction } \\
\text { enthalpy } \\
\text { [definition] }
\end{array} \tag{2C.3}
\end{align*}
\]

where in each case the molar enthalpies of the species are multiplied by their (dimensionless and positive) stoichiometric coefficients, $v$. This formal definition is of little practical value, however, because the absolute values of the standard molar enthalpies are unknown; this problem is overcome by following the techniques of Section 2C.2a.

Some standard reaction enthalpies have special names and significance. For instance, the standard enthalpy of combustion, $\Delta_{c} H^{\ominus}$, is the standard reaction enthalpy for the complete oxidation of an organic compound to $\mathrm{CO}_{2}$ gas and liquid $\mathrm{H}_{2} \mathrm{O}$ if the compound contains $\mathrm{C}, \mathrm{H}$, and O , and to $\mathrm{N}_{2}$ gas if N is also present.

\subsection*{Brief illustration 2C. 1}
The combustion of glucose is

$$
\begin{gathered}
\mathrm{C}_{6} \mathrm{H}_{12} \mathrm{O}_{6}(\mathrm{~s})+6 \mathrm{O}_{2}(\mathrm{~g}) \rightarrow 6 \mathrm{CO}_{2}(\mathrm{~g})+6 \mathrm{H}_{2} \mathrm{O}(\mathrm{l}) \\
\Delta_{\mathrm{c}} H^{\ominus}=-2808 \mathrm{~kJ} \mathrm{~mol}^{-1}
\end{gathered}
$$

The value quoted shows that 2808 kJ of heat is released when $1 \mathrm{~mol} \mathrm{C}_{6} \mathrm{H}_{12} \mathrm{O}_{6}$ burns under standard conditions (at 298 K ). More values are given in Table 2C.3.

Table 2C. 3 Standard enthalpies of formation and combustion of organic compounds at $298 \mathrm{~K}^{*}$

\begin{center}
\begin{tabular}{lcl}
\hline
 & $\Delta_{\mathrm{f}} \boldsymbol{H}^{\ominus} /\left(\mathbf{k j ~ m o l}^{-1}\right)$ & $\Delta_{c} \boldsymbol{H}^{\ominus} /(\mathbf{k J ~ m o l}$ \\
 &  &  \\
\hline
Benzene, $\mathrm{C}_{6} \mathrm{H}_{6}(\mathrm{l})$ & +49.0 & -3268 \\
Ethane, $\mathrm{C}_{2} \mathrm{H}_{6}(\mathrm{~g})$ & -84.7 & -1560 \\
Glucose, $\mathrm{C}_{6} \mathrm{H}_{12} \mathrm{O}_{6}(\mathrm{~s})$ & -1274 & -2808 \\
Methane, $\mathrm{CH}_{4}(\mathrm{~g})$ & -74.8 & -890 \\
Methanol, $\mathrm{CH}_{3} \mathrm{OH}(\mathrm{l})$ & -238.7 & -721 \\
\hline
\end{tabular}
\end{center}

\begin{itemize}
  \item More values are given in the Resource section.
\end{itemize}

\subsection*{(c) Hess's law}
Standard reaction enthalpies can be combined to obtain the value for another reaction. This application of the First Law is called Hess's law:

The standard reaction enthalpy is the sum of the values for the individual reactions into which the overall reaction may be divided.

The individual steps need not be realizable in practice: they may be 'hypothetical' reactions, the only requirement being that their chemical equations should balance. The thermodynamic basis of the law is the path-independence of the value of $\Delta_{\mathrm{r}} H^{\ominus}$. The importance of Hess's law is that information about a reaction of interest, which may be difficult to determine directly, can be assembled from information on other reactions.

\subsection*{Example 2C. 1 Using Hess's law}
The standard reaction enthalpy for the hydrogenation of propene,

$$
\mathrm{CH}_{2}=\mathrm{CHCH}_{3}(\mathrm{~g})+\mathrm{H}_{2}(\mathrm{~g}) \rightarrow \mathrm{CH}_{3} \mathrm{CH}_{2} \mathrm{CH}_{3}(\mathrm{~g})
$$

is $-124 \mathrm{~kJ} \mathrm{~mol}^{-1}$. The standard reaction enthalpy for the combustion of propane,

$$
\mathrm{CH}_{3} \mathrm{CH}_{2} \mathrm{CH}_{3}(\mathrm{~g})+5 \mathrm{O}_{2}(\mathrm{~g}) \rightarrow 3 \mathrm{CO}_{2}(\mathrm{~g})+4 \mathrm{H}_{2} \mathrm{O}(\mathrm{l})
$$

is $-2220 \mathrm{~kJ} \mathrm{~mol}^{-1}$. The standard reaction enthalpy for the formation of water,

$$
\mathrm{H}_{2}(\mathrm{~g})+\frac{1}{2} \mathrm{O}_{2}(\mathrm{~g}) \rightarrow \mathrm{H}_{2} \mathrm{O}(\mathrm{l})
$$

is $-286 \mathrm{~kJ} \mathrm{~mol}^{-1}$. Calculate the standard enthalpy of combustion of propene.

Collect your thoughts The skill you need to develop is the ability to assemble a given thermochemical equation from others. Add or subtract the reactions given, together with any others needed, so as to reproduce the reaction required. Then add or subtract the reaction enthalpies in the same way.

The solution The combustion reaction is

$$
\mathrm{C}_{3} \mathrm{H}_{6}(\mathrm{~g})+\frac{9}{2} \mathrm{O}_{2}(\mathrm{~g}) \rightarrow 3 \mathrm{CO}_{2}(\mathrm{~g})+3 \mathrm{H}_{2} \mathrm{O}(\mathrm{l})
$$

This reaction can be recreated from the following sum:

\begin{center}
\begin{tabular}{lc}
\hline
 & $\Delta_{\mathrm{r}} \mathrm{H}^{\ominus} /\left(\mathrm{kJ} \mathrm{mol}^{-1}\right)$ \\
\hline
$\mathrm{C}_{3} \mathrm{H}_{6}(\mathrm{~g})+\mathrm{H}_{2}(\mathrm{~g}) \rightarrow \mathrm{C}_{3} \mathrm{H}_{8}(\mathrm{~g})$ & -124 \\
$\mathrm{C}_{3} \mathrm{H}_{8}(\mathrm{~g})+5 \mathrm{O}_{2}(\mathrm{~g}) \rightarrow 3 \mathrm{CO}_{2}(\mathrm{~g})+4 \mathrm{H}_{2} \mathrm{O}(\mathrm{l})$ & -2220 \\
$\mathrm{H}_{2} \mathrm{O}(\mathrm{l}) \rightarrow \mathrm{H}_{2}(\mathrm{~g})+\frac{1}{2} \mathrm{O}_{2}(\mathrm{~g})$ & +286 \\
\hline
$\mathrm{C}_{3} \mathrm{H}_{6}(\mathrm{~g})+\frac{9}{2} \mathrm{O}_{2}(\mathrm{~g}) \rightarrow 3 \mathrm{CO}_{2}(\mathrm{~g})+3 \mathrm{H}_{2} \mathrm{O}(\mathrm{l})$ & -2058 \\
\hline
\end{tabular}
\end{center}

Self-test 2C. 1 Calculate the standard enthalpy of hydrogenation of liquid benzene from its standard enthalpy of combustion $\left(-3268 \mathrm{~kJ} \mathrm{~mol}^{-1}\right)$ and the standard enthalpy of combustion of liquid cyclohexane $\left(-3920 \mathrm{~kJ} \mathrm{~mol}^{-1}\right)$.

\subsection*{2C. 2 Standard enthalpies of formation}
The standard enthalpy of formation, $\Delta_{\mathrm{f}} H^{\ominus}$, of a substance is the standard reaction enthalpy for the formation of the compound from its elements in their reference states:

The reference state of an element is its most stable state at the specified temperature and 1 bar.

For example, at 298 K the reference state of nitrogen is a gas of $\mathrm{N}_{2}$ molecules, that of mercury is liquid mercury, that of carbon is graphite, and that of tin is the white (metallic) form. There is one exception to this general prescription of reference states: the reference state of phosphorus is taken to be white phosphorus despite this allotrope not being the most stable form but simply the most reproducible form of the element. Standard enthalpies of formation are expressed as enthalpies per mole of molecules or (for ionic substances) formula units of the compound. The standard enthalpy of formation of liquid benzene at 298 K , for example, refers to the reaction

$$
6 \mathrm{C}(\mathrm{~s}, \text { graphite })+3 \mathrm{H}_{2}(\mathrm{~g}) \rightarrow \mathrm{C}_{6} \mathrm{H}_{6}(\mathrm{l})
$$

and is $+49.0 \mathrm{~kJ} \mathrm{~mol}^{-1}$. The standard enthalpies of formation of elements in their reference states are zero at all temperatures because they are the enthalpies of such 'null' reactions as $\mathrm{N}_{2}(\mathrm{~g}) \rightarrow \mathrm{N}_{2}(\mathrm{~g})$. Some enthalpies of formation are listed in Tables 2C. 4 and 2C. 5 and a much longer list will be found in the Resource section.

The standard enthalpy of formation of ions in solution poses a special problem because it is not possible to prepare a solution of either cations or anions alone. This problem is overcome by defining one ion, conventionally the hydrogen\\
ion, to have zero standard enthalpy of formation at all temperatures:

\[
\Delta_{\mathrm{f}} H^{\ominus}\left(\mathrm{H}^{+}, \mathrm{aq}\right)=0 \quad \begin{align*}
& \text { Ions in solution }  \tag{2C.4}\\
& \text { [convention] }
\end{align*}
\]

\subsection*{Brief illustration 2C. 2}
If the enthalpy of formation of $\operatorname{HBr}(\mathrm{aq})$ is found to be $-122 \mathrm{~kJ} \mathrm{~mol}^{-1}$, then the whole of that value is ascribed to the formation of $\mathrm{Br}^{-}(\mathrm{aq})$, and $\Delta_{\mathrm{f}} H^{\ominus}\left(\mathrm{Br}^{-}, \mathrm{aq}\right)=-122 \mathrm{~kJ} \mathrm{~mol}^{-1}$. That value may then be combined with, for instance, the enthalpy of formation of $\operatorname{AgBr}(\mathrm{aq})$ to determine the value of $\Delta_{\mathrm{f}} H^{\ominus}\left(\mathrm{Ag}^{+}, \mathrm{aq}\right)$, and so on. In essence, this definition adjusts the actual values of the enthalpies of formation of ions by a fixed value, which is chosen so that the standard value for one of them, $\mathrm{H}^{+}(\mathrm{aq})$, is zero.

Conceptually, a reaction can be regarded as proceeding by decomposing the reactants into their elements in their reference states and then forming those elements into the products. The value of $\Delta_{\mathrm{r}} H^{\ominus}$ for the overall reaction is the sum of these 'unforming' and forming enthalpies. Because 'unforming' is the reverse of forming, the enthalpy of an unforming step is

Table 2C. 4 Standard enthalpies of formation of inorganic compounds at $298 \mathrm{~K}^{*}$

\begin{center}
\begin{tabular}{lc}
\hline
 & $\Delta_{\mathrm{f}} \mathrm{H}^{\ominus} /\left(\mathbf{k} \mathbf{~ m o l}^{-1}\right)$ \\
\hline
$\mathrm{H}_{2} \mathrm{O}(\mathrm{l})$ & -285.83 \\
$\mathrm{H}_{2} \mathrm{O}(\mathrm{g})$ & -241.82 \\
$\mathrm{NH}_{3}(\mathrm{~g})$ & -46.11 \\
$\mathrm{~N}_{2} \mathrm{H}_{4}(\mathrm{l})$ & +50.63 \\
$\mathrm{NO}_{2}(\mathrm{~g})$ & +33.18 \\
$\mathrm{~N}_{2} \mathrm{O}_{4}(\mathrm{~g})$ & +9.16 \\
$\mathrm{NaCl}(\mathrm{s})$ & -411.15 \\
$\mathrm{KCl}(\mathrm{s})$ & -436.75 \\
\hline
\end{tabular}
\end{center}

\begin{itemize}
  \item More values are given in the Resource section.
\end{itemize}

Table 2C. 5 Standard enthalpies of formation of organic compounds at $298 \mathrm{~K}^{*}$

\begin{center}
\begin{tabular}{lc}
\hline
 & $\Delta_{\mathrm{f}} \boldsymbol{H}^{\ominus} /\left(\mathbf{k} \mathbf{~ m o l}^{-1}\right)$ \\
\hline
$\mathrm{CH}_{4}(\mathrm{~g})$ & -74.81 \\
$\mathrm{C}_{6} \mathrm{H}_{6}(\mathrm{l})$ & +49.0 \\
$\mathrm{C}_{6} \mathrm{H}_{12}(\mathrm{l})$ & -156 \\
$\mathrm{CH}_{3} \mathrm{OH}(\mathrm{l})$ & -238.66 \\
$\mathrm{CH}_{3} \mathrm{CH}_{2} \mathrm{OH}(\mathrm{l})$ & -277.69 \\
\hline
\end{tabular}
\end{center}

\begin{itemize}
  \item More values are given in the Resource section.\\
the negative of the enthalpy of formation (3). Hence, in the enthalpies of formation of substances, there is enough information to calculate the enthalpy of any reaction by using
\end{itemize}

$$
\begin{aligned}
\Delta_{\mathrm{r}} H^{\ominus}= & \sum_{\text {Products }} v \Delta_{\mathrm{f}} H^{\ominus} \\
& -\sum_{\text {Reactants }} v \Delta_{\mathrm{f}} H^{\ominus}
\end{aligned}
$$

Standard reaction enthalpy [practical implementation]\\
(2C.5a)\\
where in each case the enthalpies of formation of the species that occur are multiplied by their stoichiometric coefficients. This procedure is the practical implementation of the formal definition in eqn 2C.3. A more sophisticated way of expressing the same result is to introduce the stoichiometric numbers $v_{\mathrm{J}}$ (as distinct from the stoichiometric coefficients) which are positive for products and negative for reactants. Then


\begin{equation*}
\Delta_{\mathrm{r}} H^{\ominus}=\sum_{\mathrm{J}} v_{\mathrm{J}} \Delta_{\mathrm{f}} H^{\ominus}(\mathrm{J}) \tag{2C.5b}
\end{equation*}


Stoichiometric numbers, which have a sign, are denoted $v_{\mathrm{J}}$ or $v(\mathrm{~J})$. Stoichiometric coefficients, which are all positive, are denoted simply $v$ (with no subscript).

\subsection*{Brief illustration 2C. 3}
According to eqn 2C.5a, the standard enthalpy of the reaction $2 \mathrm{HN}_{3}(\mathrm{l})+2 \mathrm{NO}(\mathrm{g}) \rightarrow \mathrm{H}_{2} \mathrm{O}_{2}(\mathrm{l})+4 \mathrm{~N}_{2}(\mathrm{~g})$ is calculated as follows:

$$
\begin{aligned}
\Delta_{\mathrm{r}} H^{\ominus}= & \left\{\Delta_{\mathrm{f}} H^{\ominus}\left(\mathrm{H}_{2} \mathrm{O}_{2}, \mathrm{l}\right)+4 \Delta_{\mathrm{f}} H^{\ominus}\left(\mathrm{N}_{2}, \mathrm{~g}\right)\right\} \\
& -\left\{2 \Delta_{\mathrm{f}} H^{\ominus}\left(\mathrm{HN}_{3}, \mathrm{l}\right)+2 \Delta_{\mathrm{f}} H^{\ominus}(\mathrm{NO}, \mathrm{~g})\right\} \\
= & \{-187.78+4(0)\} \mathrm{kJ} \mathrm{~mol}^{-1} \\
& -\{2(264.0)+2(90.25)\} \mathrm{kJ} \mathrm{~mol}^{-1} \\
= & -896.3 \mathrm{~kJ} \mathrm{~mol}^{-1}
\end{aligned}
$$

To use eqn 2C.5b, identify $v\left(\mathrm{HN}_{3}\right)=-2, v(\mathrm{NO})=-2, v\left(\mathrm{H}_{2} \mathrm{O}_{2}\right)$ $=+1$, and $v\left(\mathrm{~N}_{2}\right)=+4$, and then write

$$
\begin{aligned}
\Delta_{\mathrm{r}} H^{\ominus}= & \Delta_{\mathrm{f}} H^{\ominus}\left(\mathrm{H}_{2} \mathrm{O}_{2}, \mathrm{l}\right)+4 \Delta_{\mathrm{f}} H^{\ominus}\left(\mathrm{N}_{2}, \mathrm{~g}\right)-2 \Delta_{\mathrm{f}} H^{\ominus}\left(\mathrm{HN}_{3}, \mathrm{l}\right) \\
& -2 \Delta_{\mathrm{f}} H^{\ominus}(\mathrm{NO}, \mathrm{~g})
\end{aligned}
$$

which gives the same result.

\subsection*{2c.3 The temperature dependence of reaction enthalpies}
Many standard reaction enthalpies have been measured at different temperatures. However, in the absence of this information, standard reaction enthalpies at different temperatures can be calculated from heat capacities and the reaction enthalpy at some other temperature (Fig. 2C.1). In many cases heat capacity data are more accurate than reaction enthalpies. Therefore, providing the information is available, the procedure about to be described is more accurate than the direct measurement of a reaction enthalpy at an elevated temperature.

It follows from eqn 2B.6a $\left(\mathrm{d} H=C_{p} \mathrm{~d} T\right)$ that, when a substance is heated from $T_{1}$ to $T_{2}$, its enthalpy changes from $H\left(T_{1}\right)$ to


\begin{equation*}
H\left(T_{2}\right)=H\left(T_{1}\right)+\int_{T_{1}}^{T_{2}} C_{p} \mathrm{~d} T \tag{2C.6}
\end{equation*}


(It has been assumed that no phase transition takes place in the temperature range of interest.) Because this equation applies to each substance in the reaction, the standard reaction enthalpy changes from $\Delta_{\mathrm{r}} H^{\ominus}\left(T_{1}\right)$ to


\begin{equation*}
\Delta_{\mathrm{r}} H^{\ominus}\left(T_{2}\right)=\Delta_{\mathrm{r}} H^{\ominus}\left(T_{1}\right)+\int_{T_{1}}^{T_{2}} \Delta_{\mathrm{r}} C_{p}^{\ominus} \mathrm{d} T \quad \text { Kirchhoff's law } \tag{2C.7a}
\end{equation*}


where $\Delta_{\mathrm{r}} C_{\mathrm{p}}^{\ominus}$ is the difference of the molar heat capacities of products and reactants under standard conditions weighted by the stoichiometric coefficients that appear in the chemical equation:


\begin{equation*}
\Delta_{\mathrm{r}} C_{p}^{\ominus}=\sum_{\text {Products }} v C_{p, \mathrm{~m}}^{\ominus}-\sum_{\text {Reactants }} v C_{p, \mathrm{~m}}^{\ominus} \tag{2C.7b}
\end{equation*}


\begin{center}
\includegraphics[max width=\textwidth]{2025_02_27_d4b95d9718f0b9e2e6b9g-087}
\end{center}

Figure 2C.1 When the temperature is increased, the enthalpy of the products and the reactants both increase, but may do so to different extents. In each case, the change in enthalpy depends on the heat capacities of the substances. The change in reaction enthalpy reflects the difference in the changes of the enthalpies of the products and reactants.\\
or, in the notation of eqn 2C.5b,


\begin{equation*}
\Delta_{\mathrm{r}} C_{p}^{\ominus}=\sum_{\mathrm{J}} v_{\mathrm{J}} C_{p, \mathrm{~m}}^{\ominus}(\mathrm{J}) \tag{2C.7c}
\end{equation*}


Equation 2C.7a is known as Kirchhoff's law. It is normally a good approximation to assume that $\Delta_{\mathrm{r}} C_{p}^{\ominus}$ is independent of the temperature, at least over reasonably limited ranges. Although the individual heat capacities might vary, their difference varies less significantly. In some cases the temperature dependence of heat capacities is taken into account by using eqn 2C.7a. If $\Delta_{\mathrm{r}} C_{p}^{\ominus}$ is largely independent of temperature in the range $T_{1}$ to $T_{2}$, the integral in eqn 2C.7a evaluates to $\left(T_{2}-T_{1}\right) \Delta_{\mathrm{r}} C_{p}^{\ominus}$ and that equation becomes

\[
\Delta_{\mathrm{r}} H^{\ominus}\left(T_{2}\right)=\Delta_{\mathrm{r}} H^{\ominus}\left(T_{1}\right)+\Delta_{\mathrm{r}} C_{p}^{\ominus}\left(T_{2}-T_{1}\right) \quad \begin{align*}
& \text { Integrated }  \tag{2C.7d}\\
& \text { form of } \\
& \text { Kirchhoff's law }
\end{align*}
\]

\subsection*{Example 2C. 2 Using Kirchhoff's law}
The standard enthalpy of formation of $\mathrm{H}_{2} \mathrm{O}(\mathrm{g})$ at 298 K is $-241.82 \mathrm{~kJ} \mathrm{~mol}^{-1}$. Estimate its value at $100^{\circ} \mathrm{C}$ given the following values of the molar heat capacities at constant pressure: $\mathrm{H}_{2} \mathrm{O}(\mathrm{g}): 33.58 \mathrm{~J} \mathrm{~K}^{-1} \mathrm{~mol}^{-1} ; \mathrm{H}_{2}(\mathrm{~g}): 28.84 \mathrm{~J} \mathrm{~K}^{-1} \mathrm{~mol}^{-1} ; \mathrm{O}_{2}(\mathrm{~g})$ : $29.37 \mathrm{~J} \mathrm{~K}^{-1} \mathrm{~mol}^{-1}$. Assume that the heat capacities are independent of temperature.

Collect your thoughts When $\Delta_{\mathrm{r}} \mathrm{C}_{p}^{\ominus}$ is independent of temperature in the range $T_{1}$ to $T_{2}$, you can use the integrated form of the Kirchhoff equation, eqn 2 C .7 d . To proceed, write the chemical equation, identify the stoichiometric coefficients, and calculate $\Delta_{\mathrm{r}} C_{p}^{\ominus}$ from the data.\\
The solution The reaction is $\mathrm{H}_{2}(\mathrm{~g})+\frac{1}{2} \mathrm{O}_{2}(\mathrm{~g}) \rightarrow \mathrm{H}_{2} \mathrm{O}(\mathrm{g})$, so

$$
\begin{aligned}
\Delta_{\mathrm{r}} C_{p}^{\ominus} & =C_{p, \mathrm{~m}}^{\ominus}\left(\mathrm{H}_{2} \mathrm{O}, \mathrm{~g}\right)-\left\{C_{p, \mathrm{~m}}^{\ominus}\left(\mathrm{H}_{2}, \mathrm{~g}\right)+\frac{1}{2} C_{p, \mathrm{~m}}^{\ominus}\left(\mathrm{O}_{2}, \mathrm{~g}\right)\right\} \\
& =-9.94 \mathrm{~J} \mathrm{~K}^{-1} \mathrm{~mol}^{-1}
\end{aligned}
$$

It then follows that

$$
\begin{aligned}
\Delta_{\mathrm{r}} H^{\ominus}(373 \mathrm{~K})= & -241.82 \mathrm{~kJ} \mathrm{~mol}^{-1}+(75 \mathrm{~K}) \\
& \times\left(-9.94 \mathrm{~J} \mathrm{~K}^{-1} \mathrm{~mol}^{-1}\right)=-242.6 \mathrm{~kJ} \mathrm{~mol}^{-1}
\end{aligned}
$$

Self-test 2C.2 Estimate the standard enthalpy of formation of cyclohexane, $\mathrm{C}_{6} \mathrm{H}_{12}(\mathrm{l})$, at 400 K from the data in Table 2C. 5 and heat capacity data given in the Resource section.

\subsection*{2c. 4 Experimental techniques}
The classic tool of thermochemistry is the calorimeter (Topics 2 A and 2B). However, technological advances have been made that allow measurements to be made on samples with mass as little as a few milligrams.

\subsection*{(a) Differential scanning calorimetry}
A differential scanning calorimeter (DSC) measures the energy transferred as heat to or from a sample at constant pressure during a physical or chemical change. The term 'differential' refers to the fact that measurements on a sample are compared to those on a reference material that does not undergo a physical or chemical change during the analysis. The term 'scanning' refers to the fact that the temperatures of the sample and reference material are increased, or scanned, during the analysis.

A DSC consists of two small compartments that are heated electrically at a constant rate. The temperature, $T$, at time $t$ during a linear scan is $T=T_{0}+\alpha t$, where $T_{0}$ is the initial temperature and $\alpha$ is the scan rate. A computer controls the electrical power supply that maintains the same temperature in the sample and reference compartments throughout the analysis (Fig. 2C.2).

If no physical or chemical change occurs in the sample at temperature $T$, the heat transferred to the sample is written as $q_{p}=C_{p} \Delta T$, where $\Delta T=T-T_{0}$ and $C_{p}$ is assumed to be independent of temperature. Because $T=T_{0}+\alpha t$, it follows that $\Delta T$ $=\alpha t$. If a chemical or physical process takes place, the energy required to be transferred as heat to attain the same change in temperature of the sample as the control is $q_{p}+q_{p, \mathrm{ex}}$.

The quantity $q_{p, \text { ex }}$ is interpreted in terms of an apparent change in the heat capacity at constant pressure, from $C_{p}$ to $C_{p}+C_{p, \text { ex }}$ of the sample during the temperature scan:


\begin{equation*}
C_{p, \mathrm{ex}}=\frac{q_{p, \mathrm{ex}}}{\Delta T}=\frac{q_{p, \mathrm{ex}}}{\alpha t}=\frac{P_{\mathrm{ex}}}{\alpha} \tag{2C.8}
\end{equation*}


where $P_{\text {ex }}=q_{p, \text { ex }} / t$ is the excess electrical power necessary to equalize the temperature of the sample and reference compartments. A DSC trace, also called a thermogram, consists of\\
\includegraphics[max width=\textwidth, center]{2025_02_27_d4b95d9718f0b9e2e6b9g-088}

Figure 2C. 2 A differential scanning calorimeter. The sample and a reference material are heated in separate but identical metal heat sinks. The output is the difference in power needed to maintain the heat sinks at equal temperatures as the temperature rises.\\
\includegraphics[max width=\textwidth, center]{2025_02_27_d4b95d9718f0b9e2e6b9g-089(2)}

Figure 2C. 3 A thermogram for the protein ubiquitin at $\mathrm{pH}=$ 2.45. The protein retains its native structure up to about $45^{\circ} \mathrm{C}$ and then undergoes an endothermic conformational change. (Adapted from B. Chowdhry and S. LeHarne, J. Chem. Educ. 74, 236 (1997).)\\
a plot of $C_{p, \text { ex }}$ against $T$ (Fig. 2C.3). The enthalpy change associated with the process is


\begin{equation*}
\Delta H=\int_{T_{1}}^{T_{2}} C_{p, \mathrm{ex}} \mathrm{~d} T \tag{2C.9}
\end{equation*}


where $T_{1}$ and $T_{2}$ are, respectively, the temperatures at which the process begins and ends. This relation shows that the enthalpy change is equal to the area under the plot of $C_{p, \mathrm{ex}}$ against $T$.

\subsection*{(b) Isothermal titration calorimetry}
Isothermal titration calorimetry (ITC) is also a 'differential' technique in which the thermal behaviour of a sample is compared with that of a reference. The apparatus is shown in Fig. 2C.4. One of the thermally conducting vessels, which have a volume of a few cubic centimetres, contains the reference (water for instance) and a heater rated at a few milliwatts. The second vessel contains one of the reagents, such as a solution of a macromolecule with binding sites; it also contains a heater. At the start of the experiment, both heaters are activated, and then precisely determined amounts (of volume of about a cubic millimetre) of the second reagent are added to the reaction cell. The power required to maintain the same temperature differential with the reference cell is monitored.\\
\includegraphics[max width=\textwidth, center]{2025_02_27_d4b95d9718f0b9e2e6b9g-089}

Figure 2C.4 A schematic diagram of the apparatus used for isothermal titration calorimetry.\\
\includegraphics[max width=\textwidth, center]{2025_02_27_d4b95d9718f0b9e2e6b9g-089(1)}

Figure 2C. 5 (a) The record of the power applied as each injection is made, and (b) the sum of successive enthalpy changes in the course of the titration.

If the reaction is exothermic, less power is needed; if it is endothermic, then more power must be supplied.

A typical result is shown in Fig. 2C.5, which shows the power needed to maintain the temperature differential: from the power and the length of time, $\Delta t$, for which it is supplied, the heat supplied, $q_{i}$, for the injection $i$ can be calculated from $q_{i}=P_{i} \Delta t$. If the volume of solution is $V$ and the molar concentration of unreacted reagent A is $c_{i}$ at the time of the $i$ th injection, then the change in its concentration at that injection is $\Delta c_{i}$ and the heat generated (or absorbed) by the reaction is $V \Delta_{\mathrm{r}} H \Delta c_{i}=q_{i}$. The sum of all such quantities, given that the sum of $\Delta c_{i}$ is the known initial concentration of the reactant, can then be interpreted as the value of $\Delta_{\mathrm{r}} H$ for the reaction.

\subsection*{Checklist of concepts}
$\square$ 1. The standard enthalpy of transition is equal to the energy transferred as heat at constant pressure in the transition under standard conditions.\\
$\square$ 2. The standard state of a substance at a specified temperature is its pure form at 1 bar.3. A thermochemical equation is a chemical equation and its associated change in enthalpy.4. Hess's law states that the standard reaction enthalpy is the sum of the values for the individual reactions into which the overall reaction may be divided.5. Standard enthalpies of formation are defined in terms of the reference states of elements.6. The reference state of an element is its most stable state at the specified temperature and 1 bar.7. The standard reaction enthalpy is expressed as the difference of the standard enthalpies of formation of products and reactants. expressed by Kirchhoff's law.

\subsection*{Checklist of equations}
\begin{center}
\begin{tabular}{|c|c|c|c|}
\hline
Property & Equation & Comment & Equation number \\
\hline
The standard reaction enthalpy & \( \Delta_{\mathrm{r}} H^{\ominus}=\sum_{\text {Products }} v \Delta_{\mathrm{f}} H^{\ominus}-\sum_{\text {Reacants }} v \Delta_{\mathrm{f}} H^{\ominus} \) & \begin{tabular}{l}
$v$ : stoichiometric coefficients; \\
$v_{\mathrm{J}}$ : (signed) stoichiometric numbers \\
\end{tabular} & 2C. 5 \\
\hline
 & \( \Delta_{\mathrm{r}} H^{\ominus}=\sum_{\mathrm{J}} v_{\mathrm{J}} \Delta_{\mathrm{f}} H^{\ominus}(\mathrm{J}) \) &  &  \\
\hline
\multirow[t]{3}{*}{Kirchhoff's law} & \( \Delta_{\mathrm{r}} H^{\ominus}\left(T_{2}\right)=\Delta_{\mathrm{r}} H^{\ominus}\left(T_{1}\right)+\int_{T_{1}}^{T_{2}} \Delta_{\mathrm{r}} C_{p}^{\ominus} \mathrm{d} T \) &  & 2C.7a \\
\hline
 & $\Delta_{\mathrm{r}} C_{p}^{\ominus}=\sum_{\mathrm{J}} v_{\mathrm{J}} C_{p, \mathrm{~m}}^{\ominus}(\mathrm{J})$ &  & 2C.7c \\
\hline
 & $\Delta_{\mathrm{r}} H^{\ominus}\left(T_{2}\right)=\Delta_{\mathrm{r}} H^{\ominus}\left(T_{1}\right)+\left(T_{2}-T_{1}\right) \Delta_{\mathrm{r}} C_{p}^{\ominus}$ & If $\Delta_{\mathrm{r}} C_{p}^{\ominus}$ independent of temperature & 2C.7d \\
\hline
\end{tabular}
\end{center}

\section*{TOPIC 2D State functions and exact differentials}
\subsection*{Why do you need to know this material?}
Thermodynamics has the power to provide relations between a variety of properties. This Topic introduces its key procedure, the manipulation of equations involving state functions.

\subsection*{What is the key idea?}
The fact that internal energy and enthalpy are state functions leads to relations between thermodynamic properties.

\subsection*{> What do you need to know already?}
You need to be aware that the internal energy and enthalpy are state functions (Topics 2 B and 2 C ) and be familiar with the concept of heat capacity. You need to be able to make use of several simple relations involving partial derivatives (The chemist's toolkit 9 in Topic 2A).

A state function is a property that depends only on the current state of a system and is independent of its history. The internal energy and enthalpy are two examples. Physical quantities with values that do depend on the path between two states are called path functions. Examples of path functions are the work and the heating that are done when preparing a state. It is not appropriate to speak of a system in a particular state as possessing work or heat. In each case, the energy transferred as work or heat relates to the path being taken between states, not the current state itself.

A part of the richness of thermodynamics is that it uses the mathematical properties of state functions to draw far-reaching conclusions about the relations between physical properties and thereby establish connections that may be completely unexpected. The practical importance of this ability is the possibility of combining measurements of different properties to obtain the value of a desired property.

\subsection*{2D.1 Exact and inexact differentials}
Consider a system undergoing the changes depicted in Fig. 2D.1. The initial state of the system is $i$ and in this state the internal energy is $U_{\mathrm{i}}$. Work is done by the system as it expands\\
\includegraphics[max width=\textwidth, center]{2025_02_27_d4b95d9718f0b9e2e6b9g-091}

Figure 2D. 1 As the volume and temperature of a system are changed, the internal energy changes. An adiabatic and a nonadiabatic path are shown as Path 1 and Path 2, respectively: they correspond to different values of $q$ and $w$ but to the same value of $\Delta U$.\\
adiabatically to a state $f$. In this state the system has an internal energy $U_{\mathrm{f}}$ and the work done on the system as it changes along Path 1 from ito $f$ is $w$. Notice the use of language: $U$ is a property of the state; $w$ is a property of the path. Now consider another process, Path 2, in which the initial and final states are the same as those in Path 1 but in which the expansion is not adiabatic. The internal energy of both the initial and the final states are the same as before (because $U$ is a state function). However, in the second path an energy $q^{\prime}$ enters the system as heat and the work $w^{\prime}$ is not the same as $w$. The work and the heat are path functions.

If a system is taken along a path (e.g. by heating it), $U$ changes from $U_{\mathrm{i}}$ to $U_{\mathrm{f}}$, and the overall change is the sum (integral) of all the infinitesimal changes along the path:


\begin{equation*}
\Delta U=\int_{\mathrm{i}}^{\mathrm{f}} \mathrm{~d} U \tag{2D.1}
\end{equation*}


The value of $\Delta U$ depends on the initial and final states of the system but is independent of the path between them. This path-independence of the integral is expressed by saying that $\mathrm{d} U$ is an 'exact differential'. In general, an exact differential is an infinitesimal quantity that, when integrated, gives a result that is independent of the path between the initial and final states.

When a system is heated, the total energy transferred as heat is the sum of all individual contributions at each point of the path:


\begin{equation*}
q=\int_{\mathrm{i}, \mathrm{path}}^{\mathrm{f}} \mathrm{~d} q \tag{2D.2}
\end{equation*}


Notice the differences between this equation and eqn 2D.1. First, the result of integration is $q$ and not $\Delta q$, because $q$ is not a state function and the energy supplied as heat cannot be expressed as $q_{\mathrm{f}}-q_{\mathrm{i}}$. Secondly, the path of integration must be specified because $q$ depends on the path selected (e.g. an adiabatic path has $q=0$, whereas a non-adiabatic path between the same two states would have $q \neq 0$ ). This path dependence is expressed by saying that $\mathrm{d} q$ is an 'inexact differential'. In general, an inexact differential is an infinitesimal quantity that, when integrated, gives a result that depends on the path between the initial and final states. Often $\mathrm{d} q$ is written $đ q$ to emphasize that it is inexact and requires the specification of a path.

The work done on a system to change it from one state to another depends on the path taken between the two specified states. For example, in general the work is different if the change takes place adiabatically and non-adiabatically. It follows that $\mathrm{d} w$ is an inexact differential. It is often written $\AA w$.

\subsection*{Example 2D. 1 Calculating work, heat, and change in internal energy}
Consider a perfect gas inside a cylinder fitted with a piston. Let the initial state be $T, V_{\mathrm{i}}$ and the final state be $T, V_{\mathrm{f}}$. The change of state can be brought about in many ways, of which the two simplest are the following:

\begin{itemize}
  \item Path 1 , in which there is free expansion against zero external pressure;
  \item Path 2, in which there is reversible, isothermal expansion.
\end{itemize}

Calculate $w, q$, and $\Delta U$ for each process.\\
Collect your thoughts To find a starting point for a calculation in thermodynamics, it is often a good idea to go back to first principles and to look for a way of expressing the quantity to be calculated in terms of other quantities that are easier to calculate. It is argued in Topic 2B that the internal energy of a perfect gas depends only on the temperature and is independent of the volume those molecules occupy, so for any isothermal change, $\Delta U=0$. Also, $\Delta U=q+w$ in general. To solve the problem you need to combine the two expressions, selecting the appropriate expression for the work done from the discussion in Topic 2A.

The solution Because $\Delta U=0$ for both paths and $\Delta U=q+w$, in each case $q=-w$. The work of free expansion is zero (eqn 2 A .7 of Topic $2 \mathrm{~A}, w=0$ ); so in Path $1, w=0$ and therefore $q=$ 0 too. For Path 2, the work is given by eqn 2A. 9 of Topic 2A ( $w$ $=-n R T \ln \left(V_{\mathrm{f}} / V_{\mathrm{i}}\right)$ ) and consequently $q=n R T \ln \left(V_{\mathrm{f}} / V_{\mathrm{i}}\right)$.

Self-test 2D. 1 Calculate the values of $q, w$, and $\Delta U$ for an irreversible isothermal expansion of a perfect gas against a constant non-zero external pressure.

$$
0=\Omega \nabla \nabla^{\prime} \Lambda \nabla^{x} d-=\mu ' \Lambda \nabla^{x \infty} d=b: \iota \partial м s u_{V}
$$

\subsection*{2D. 2 Changes in internal energy}
Consider a closed system of constant composition (the only type of system considered in the rest of this Topic). The internal energy $U$ can be regarded as a function of $V, T$, and $p$, but, because there is an equation of state that relates these quantities (Topic 1A), choosing the values of two of the variables fixes the value of the third. Therefore, it is possible to write $U$ in terms of just two independent variables: $V$ and $T, p$ and $T$, or $p$ and $V$. Expressing $U$ as a function of volume and temperature turns out to result in the simplest expressions.

\subsection*{(a) General considerations}
Because the internal energy is a function of the volume and the temperature, when these two quantities change, the internal energy changes by

\[
\mathrm{d} U=\left(\frac{\partial U}{\partial V}\right)_{T} \mathrm{~d} V+\left(\frac{\partial U}{\partial T}\right)_{V} \mathrm{~d} T \quad \begin{align*}
& \text { General expression }  \tag{2D.3}\\
& \text { for a change in } U \\
& \text { with } T \text { and } V
\end{align*}
\]

The interpretation of this equation is that, in a closed system of constant composition, any infinitesimal change in the internal energy is proportional to the infinitesimal changes of volume and temperature, the coefficients of proportionality being the two partial derivatives (Fig. 2D.2).\\
\includegraphics[max width=\textwidth, center]{2025_02_27_d4b95d9718f0b9e2e6b9g-092}

Figure 2D. 2 An overall change in $U$, which is denoted $\mathrm{d} U$, arises when both $V$ and $T$ are allowed to change. If second-order infinitesimals are ignored, the overall change is the sum of changes for each variable separately.\\
\includegraphics[max width=\textwidth, center]{2025_02_27_d4b95d9718f0b9e2e6b9g-093(1)}

Figure 2D. 3 The internal pressure, $\pi_{T}$, is the slope of $U$ with respect to $V$ with the temperature $T$ held constant.

In many cases partial derivatives have a straightforward physical interpretation, and thermodynamics gets shapeless and difficult only when that interpretation is not kept in sight. The term $(\partial U / \partial T)_{V}$ occurs in Topic 2A, as the constant-volume heat capacity, $C_{V}$. The other coefficient, $(\partial U / \partial V)_{T}$, denoted $\pi_{T}$, plays a major role in thermodynamics because it is a measure of the variation of the internal energy of a substance as its volume is changed at constant temperature (Fig. 2D.3). Because $\pi_{T}$ has the same dimensions as pressure but arises from the interactions between the molecules within the sample, it is called the internal pressure:

\[
\pi_{T}=\left(\frac{\partial U}{\partial V}\right)_{T} \quad \begin{align*}
& \text { Internal pressure }  \tag{2D.4}\\
& \text { [definition] }
\end{align*}
\]

In terms of the notation $C_{V}$ and $\pi_{T}$, eqn 2D. 3 can now be written


\begin{equation*}
\mathrm{d} U=\pi_{T} \mathrm{~d} V+C_{V} \mathrm{~d} T \tag{2D.5}
\end{equation*}


It is shown in Topic 3D that the statement $\pi_{T}=0$ (i.e. the internal energy is independent of the volume occupied by the sample) can be taken to be the definition of a perfect gas, because it implies the equation of state $p V \propto T$. In molecular terms, when there are no interactions between the molecules, the internal energy is independent of their separation and hence independent of the volume of the sample and $\pi_{T}=0$. If the gas is described by the van der Waals equation with $a$, the parameter corresponding to attractive interactions, dominant, then an increase in volume increases the average separation of the molecules and therefore raises the internal energy. In this case, it is expected that $\pi_{T}>0$ (Fig. 2D.4). This expectation is confirmed in Topic 3D, where it is shown that $\pi_{T}=n a / V^{2}$.

James Joule thought that he could measure $\pi_{T}$ by observing the change in temperature of a gas when it is allowed to expand into a vacuum. He used two metal vessels immersed in a water bath (Fig. 2D.5). One was filled with air at about\\
\includegraphics[max width=\textwidth, center]{2025_02_27_d4b95d9718f0b9e2e6b9g-093}

Figure 2D. 4 For a perfect gas, the internal energy is independent of the volume (at constant temperature). If attractions are dominant in a real gas, the internal energy increases with volume because the molecules become farther apart on average. If repulsions are dominant, the internal energy decreases as the gas expands.

22 atm and the other was evacuated. He then tried to measure the change in temperature of the water of the bath when a stopcock was opened and the air expanded into a vacuum. He observed no change in temperature.

The thermodynamic implications of the experiment are as follows. No work was done in the expansion into a vacuum, so $w=0$. No energy entered or left the system (the gas) as heat because the temperature of the bath did not change, so $q=0$. Consequently, within the accuracy of the experiment, $\Delta U=0$. Joule concluded that $U$ does not change when a gas expands isothermally and therefore that $\pi_{T}=0$. His experiment, however, was crude. The heat capacity of the apparatus was so large that the temperature change, which would in fact occur for a real gas, is simply too small to measure. Joule had extracted an essential limiting property of a gas, a property of a perfect gas, without detecting the small deviations characteristic of real gases.\\
\includegraphics[max width=\textwidth, center]{2025_02_27_d4b95d9718f0b9e2e6b9g-093(2)}

Figure 2D. 5 A schematic diagram of the apparatus used by Joule in an attempt to measure the change in internal energy when a gas expands isothermally. The heat absorbed by the gas is proportional to the change in temperature of the bath.

\subsection*{(b) Changes in internal energy at constant pressure}
Partial derivatives have many useful properties and some are reviewed in The chemist's toolkit 9 of Topic 2A. Skilful use of them can often turn some unfamiliar quantity into a quantity that can be recognized, interpreted, or measured.

As an example, to find how the internal energy varies with temperature when the pressure rather than the volume of the system is kept constant, begin by dividing both sides of eqn 2 D .5 by dT. Then impose the condition of constant pressure on the resulting differentials, so that $\mathrm{d} U / \mathrm{d} T$ on the left becomes $(\partial U / \partial T)_{p}$. At this stage the equation becomes

$$
\left(\frac{\partial U}{\partial T}\right)_{p}=\pi_{T}\left(\frac{\partial V}{\partial T}\right)_{p}+C_{V}
$$

As already emphasized, it is usually sensible in thermodynamics to inspect the output of a manipulation to see if it contains any recognizable physical quantity. The partial derivative on the right in this expression is the slope of the plot of volume against temperature (at constant pressure). This property is normally tabulated as the expansion coefficient, $\alpha$, of a substance, which is defined as


\begin{equation*}
\alpha=\frac{1}{V}\left(\frac{\partial V}{\partial T}\right)_{p} \tag{2D.6}
\end{equation*}


Expansion coefficient [definition]\\
and physically is the fractional change in volume that accompanies a rise in temperature. A large value of $\alpha$ means that the volume of the sample responds strongly to changes in temperature. Table 2D. 1 lists some experimental values of $\alpha$. For future reference, it also lists the isothermal compressibility, $\kappa_{T}$ (kappa), which is defined as

\[
\kappa_{T}=-\frac{1}{V}\left(\frac{\partial V}{\partial p}\right)_{T} \quad \begin{align*}
& \text { Isothermal compressibility }  \tag{2D.7}\\
& \text { [definition] }
\end{align*}
\]

The isothermal compressibility is a measure of the fractional change in volume when the pressure is increased; the nega-

Table 2D. 1 Expansion coefficients ( $\alpha$ ) and isothermal compressibilities ( $\kappa_{T}$ ) at $298 \mathrm{~K}^{*}$

\begin{center}
\begin{tabular}{lll}
\hline
 & $\alpha /\left(10^{-4} \mathrm{~K}^{-1}\right)$ & $\kappa_{T} /\left(1 \mathbf{1 0}^{-6} \mathbf{b a r}^{-1}\right)$ \\
\hline
Liquids: &  &  \\
Benzene & 12.4 & 90.9 \\
Water & 2.1 & 49.0 \\
Solids: &  &  \\
Diamond & 0.030 & 0.185 \\
Lead & 0.861 & 2.18 \\
\hline
\end{tabular}
\end{center}

\begin{itemize}
  \item More values are given in the Resource section.\\
tive sign in the definition ensures that the compressibility is a positive quantity, because an increase of pressure, implying a positive $\mathrm{d} p$, brings about a reduction of volume, a negative $\mathrm{d} V$.
\end{itemize}

Example 2D. 2 Calculating the expansion coefficient of a gas

Derive an expression for the expansion coefficient of a perfect gas.\\
Collect your thoughts The expansion coefficient is defined in eqn 2D.6. To use this expression, you need to substitute the expression for $V$ in terms of $T$ obtained from the equation of state for the gas. As implied by the subscript in eqn 2D.6, the pressure, $p$, is treated as a constant.

The solution Because $p V=n R T$, write

$$
\alpha=\frac{1}{V}\left(\frac{\partial V}{\partial T}\right)_{p}=\frac{1}{V}\left(\frac{\partial(n R T / p)}{\partial T}\right)_{p}=\frac{1}{V} \times \frac{n R}{p}=\frac{n R}{p V}=\frac{n R}{n R T}=\frac{1}{T}
$$

The physical interpretation of this result is that the higher the temperature, the less responsive is the volume of a perfect gas to a change in temperature.\\
Self-test 2D. 2 Derive an expression for the isothermal compressibility of a perfect gas.

Introduction of the definition of $\alpha$ into the equation for $(\partial U / \partial T)_{p}$ gives


\begin{equation*}
\left(\frac{\partial U}{\partial T}\right)_{p}=\alpha \pi_{T} V+C_{V} \tag{2D.8}
\end{equation*}


This equation is entirely general (provided the system is closed and its composition is constant). It expresses the dependence of the internal energy on the temperature at constant pressure in terms of $C_{\mathrm{v}}$, which can be measured in one experiment, in terms of $\alpha$, which can be measured in another, and in terms of the internal pressure $\pi_{T}$. For a perfect gas, $\pi_{T}=0$, so then


\begin{equation*}
\left(\frac{\partial U}{\partial T}\right)_{p}=C_{V} \tag{2D.9}
\end{equation*}


That is, although the constant-volume heat capacity of a perfect gas is defined as the slope of a plot of internal energy against temperature at constant volume, for a perfect gas $C_{V}$ is also the slope of a plot of internal energy against temperature at constant pressure.

Equation 2D. 9 provides an easy way to derive the relation between $C_{p}$ and $C_{V}$ for a perfect gas (they differ, as explained in Topic 2B, because some of the energy supplied as heat\\
escapes back into the surroundings as work of expansion when the volume is not constant). First, write

$$
C_{p}-C_{V}=\overbrace{\left(\frac{\partial H}{\partial T}\right)_{p}}^{\begin{array}{c}
\text { Definition } \\
\text { of } C_{p}
\end{array}}-\overbrace{\left(\frac{\partial U}{\partial T}\right)_{p}}^{\text {eqn2D. } 9}
$$

Then introduce $H=U+p V=U+n R T$ into the first term and obtain


\begin{equation*}
C_{p}-C_{V}=\left(\frac{\partial(U+n R T)}{\partial T}\right)_{p}-\left(\frac{\partial U}{\partial T}\right)_{p}=n R \tag{2D.10}
\end{equation*}


The general result for any substance (the proof makes use of the Second Law, which is introduced in Focus 3) is


\begin{equation*}
C_{p}-C_{V}=\frac{\alpha^{2} T V}{\kappa_{T}} \tag{2D.11}
\end{equation*}


This relation reduces to eqn 2D. 10 for a perfect gas when $\alpha=$ $1 / T$ and $\kappa_{T}=1 / p$. Because expansion coefficients $\alpha$ of liquids and solids are small, it is tempting to deduce from eqn 2D. 11 that for them $C_{p} \approx C_{\mathrm{v}}$. But this is not always so, because the compressibility $\kappa_{T}$ might also be small, so $\alpha^{2} / \kappa_{T}$ might be large. That is, although only a little work need be done to push back the atmosphere, a great deal of work may have to be done to pull atoms apart from one another as the solid expands.

\subsection*{Brief illustration 2D. 1}
The expansion coefficient and isothermal compressibility of water at $25^{\circ} \mathrm{C}$ are given in Table 2D. 1 as $2.1 \times 10^{-4} \mathrm{~K}^{-1}$ and $49.0 \times 10^{-6} \mathrm{bar}^{-1}\left(4.90 \times 10^{-10} \mathrm{~Pa}^{-1}\right)$, respectively. The molar volume of water at that temperature, $V_{\mathrm{m}}=M / \rho$ (where $\rho$ is the mass density), is $18.1 \mathrm{~cm}^{3} \mathrm{~mol}^{-1}\left(1.81 \times 10^{-5} \mathrm{~m}^{3} \mathrm{~mol}^{-1}\right)$. Therefore, from eqn 2D.11, the difference in molar heat capacities (which is given by using $V_{\mathrm{m}}$ in place of $V$ ) is

$$
\begin{aligned}
C_{p, \mathrm{~m}}-C_{V, \mathrm{~m}} & =\frac{\left(2.1 \times 10^{-4} \mathrm{~K}^{-1}\right)^{2} \times(298 \mathrm{~K}) \times\left(1.81 \times 10^{-5} \mathrm{~m}^{3} \mathrm{~mol}^{-1}\right)}{4.90 \times 10^{-10} \mathrm{~Pa}^{-1}} \\
& =0.485 \mathrm{Pam}^{3} \mathrm{~K}^{-1} \mathrm{~mol}^{-1}=0.485 \mathrm{JK}^{-1} \mathrm{~mol}^{-1}
\end{aligned}
$$

For water, $C_{p, \mathrm{~m}}=75.3 \mathrm{~J} \mathrm{~K}^{-1} \mathrm{~mol}^{-1}$, so $C_{V, \mathrm{~m}}=74.8 \mathrm{~J} \mathrm{~K}^{-1} \mathrm{~mol}^{-1}$. In some cases, the two heat capacities differ by as much as 30 per cent.

\subsection*{2D.3 Changes in enthalpy}
A similar set of operations can be carried out on the enthalpy, $H=U+p V$. The quantities $U, p$, and $V$ are all state functions; therefore $H$ is also a state function and $\mathrm{d} H$ is an exact differential. It turns out that $H$ is a useful thermodynamic function when the pressure can be controlled: a sign of that is the relation $\Delta H=q_{p}$ (eqn 2B.2b). Therefore, $H$ can be regarded as\\
a function of $p$ and $T$, and the argument in Section 2D. 2 for the variation of $U$ can be adapted to find an expression for the variation of $H$ with temperature at constant volume.

How is that done? 2D. 1\\
Deriving an expression for the variation of enthalpy with pressure and temperature

Consider a closed system of constant composition. Because $H$ is a function of $p$ and $T$, when these two quantities change by an infinitesimal amount, the enthalpy changes by

$$
\mathrm{d} H=\left(\frac{\partial H}{\partial p}\right)_{T} \mathrm{~d} p+\left(\frac{\partial H}{\partial T}\right)_{p} \mathrm{~d} T
$$

The second partial derivative is $C_{p}$. The task at hand is to express $(\partial H / \partial p)_{T}$ in terms of recognizable quantities. If the enthalpy is constant, then $\mathrm{d} H=0$ and

$$
\left(\frac{\partial H}{\partial p}\right)_{T} \mathrm{~d} p=-C_{p} \mathrm{~d} T \text { at constant } H
$$

Division of both sides by $\mathrm{d} p$ then gives

$$
\left(\frac{\partial H}{\partial p}\right)_{T}=-C_{p}\left(\frac{\partial T}{\partial p}\right)_{H}=-C_{p} \mu
$$

where the Joule-Thomson coefficient, $\mu(\mathrm{mu})$, is defined as

\[
\mu=\left(\frac{\partial T}{\partial p}\right)_{H} \quad \begin{align*}
& \text { Joule-Thomson coefficient }  \tag{2D.12}\\
& \text { [definition] }
\end{align*}
\]

It follows that

$$
-\mathrm{d} H=-\mu C_{p} \mathrm{~d} p+C_{p} \mathrm{~d} T \longmapsto \begin{aligned}
& \text { The variation of enthalpy with } \\
& \text { temperature and pressure }
\end{aligned}
$$

\subsection*{Brief illustration 2D. 2}
The Joule-Thomson coefficient for nitrogen at 298 K and 1 atm (Table 2D.2) is $+0.27 \mathrm{Kbar}^{-1}$. (Note that $\mu$ is an intensive property.) It follows that the change in temperature the gas undergoes when its pressure changes by -10 bar under isenthalpic conditions is

$$
\Delta T \approx \mu \Delta p=+\left(0.27 \mathrm{~K} \mathrm{bar}^{-1}\right) \times(-10 \text { bar })=-2.7 \mathrm{~K}
$$

Table 2D. 2 Inversion temperatures ( $T_{\mathrm{T}}$ ), normal freezing $\left(T_{f}\right)$ and boiling ( $T_{b}$ ) points, and Joule-Thomson coefficients ( $\mu$ ) at 1 atm and $298 \mathrm{~K}^{*}$

\begin{center}
\begin{tabular}{lrrcl}
\hline
 & $T_{\mathrm{I}} / \mathrm{K}$ & \multicolumn{1}{c}{$T_{\mathrm{f}} / \mathrm{K}$} & \multicolumn{1}{c}{$T_{\mathrm{b}} / \mathrm{K}$} & $\mu /\left(\mathrm{K} \mathrm{atm}^{-1}\right)$ \\
\hline
Ar & 723 & 83.8 & 87.3 &  \\
$\mathrm{CO}_{2}$ & 1500 & 194.7 & +1.10 & +1.11 at 300 K \\
He & 40 & 4.2 & 4.22 & -0.062 \\
$\mathrm{~N}_{2}$ & 621 & 63.3 & 77.4 & +0.27 \\
\hline
\end{tabular}
\end{center}

\footnotetext{\begin{itemize}
  \item More values are given in the Resource section.
\end{itemize}
}\subsection*{2D.4 The Joule-Thomson effect}
The analysis of the Joule-Thomson coefficient is central to the technological problems associated with the liquefaction of gases. To determine the coefficient, it is necessary to measure the ratio of the temperature change to the change of pressure, $\Delta T / \Delta p$, in a process at constant enthalpy. The cunning required to impose the constraint of constant enthalpy, so that the expansion is isenthalpic, was supplied by James Joule and William Thomson (later Lord Kelvin). They let a gas expand through a porous barrier from one constant pressure to another and monitored the difference of temperature that arose from the expansion (Fig. 2D.6). The change of temperature that they observed as a result of isenthalpic expansion is called the Joule-Thomson effect.

The 'Linde refrigerator' makes use of the Joule-Thomson effect to liquefy gases (Fig. 2D.7). The gas at high pressure is allowed to expand through a throttle; it cools and is circulated past the incoming gas. That gas is cooled, and its subsequent expansion cools it still further. There comes a stage when the circulating gas becomes so cold that it condenses to a liquid.

\subsection*{(a) The observation of the Joule-Thomson effect}
The apparatus Joule and Thomson used was insulated so that the process was adiabatic. By considering the work done at each stage it is possible to show that the expansion is isenthalpic.\\
\includegraphics[max width=\textwidth, center]{2025_02_27_d4b95d9718f0b9e2e6b9g-096(2)}

Figure 2D. 6 The apparatus used for measuring the JouleThomson effect. The gas expands through the porous barrier, which acts as a throttle, and the whole apparatus is thermally insulated. As explained in the text, this arrangement corresponds to an isenthalpic expansion (expansion at constant enthalpy). Whether the expansion results in a heating or a cooling of the gas depends on the conditions.\\
\includegraphics[max width=\textwidth, center]{2025_02_27_d4b95d9718f0b9e2e6b9g-096}

Figure 2D. 7 The principle of the Linde refrigerator is shown in this diagram. The gas is recirculated, and so long as it is beneath its inversion temperature it cools on expansion through the throttle. The cooled gas cools the high-pressure gas, which cools still further as it expands. Eventually liquefied gas drips from the throttle.

\subsection*{How is that done? 2D. 2 \\
 Establishing that the expansion is isenthalpic}
Because all changes to the gas occur adiabatically, $q=0$ and, consequently, $\Delta U=w$.

\subsection*{Step 1 Calculate the total work}
Consider the work done as the gas passes through the barrier by focusing on the passage of a fixed amount of gas from the high pressure side, where the pressure is $p_{\mathrm{i}}$, the temperature $T_{\mathrm{i}}$, and the gas occupies a volume $V_{\mathrm{i}}$ (Fig. 2D.8). The gas\\
\includegraphics[max width=\textwidth, center]{2025_02_27_d4b95d9718f0b9e2e6b9g-096(1)}

Figure 2D. 8 The thermodynamic basis of Joule-Thomson expansion. The pistons represent the upstream and downstream gases, which maintain constant pressures either side of the throttle. The transition from the top diagram to the bottom diagram, which represents the passage of a given amount of gas through the throttle, occurs without change of enthalpy.\\
emerges on the low pressure side, where the same amount of gas has a pressure $p_{\mathrm{f}}$, a temperature $T_{\mathrm{f}}$, and occupies a volume $V_{\mathrm{f}}$. The gas on the left is compressed isothermally by the upstream gas acting as a piston. The relevant pressure is $p_{\mathrm{i}}$ and the volume changes from $V_{\mathrm{i}}$ to 0 ; therefore, the work done on the gas is

$$
w_{1}=-p_{\mathrm{i}}\left(0-V_{\mathrm{i}}\right)=p_{\mathrm{i}} V_{\mathrm{i}}
$$

The gas expands isothermally on the right of the barrier (but possibly at a different constant temperature) against the pressure $p_{\mathrm{f}}$ provided by the downstream gas acting as a piston to be driven out. The volume changes from 0 to $V_{f}$, so the work done on the gas in this stage is

$$
w_{2}=-p_{\mathrm{f}}\left(V_{\mathrm{f}}-0\right)=-p_{\mathrm{f}} V_{\mathrm{f}}
$$

The total work done on the gas is the sum of these two quantities, or

$$
w=w_{1}+w_{2}=p_{\mathrm{i}} V_{\mathrm{i}}-p_{\mathrm{f}} V_{\mathrm{f}}
$$

\subsection*{Step 2 Calculate the change in internal energy}
It follows that the change of internal energy of the gas as it moves adiabatically from one side of the barrier to the other is

$$
U_{\mathrm{f}}-U_{\mathrm{i}}=w=p_{\mathrm{i}} V_{\mathrm{i}}-p_{\mathrm{f}} V_{\mathrm{f}}
$$

\subsection*{Step 3 Calculate the initial and final enthalpies}
Reorganization of the preceding expression, and noting that $H=U+p V$, gives

$$
U_{\mathrm{f}}+p_{\mathrm{f}} V_{\mathrm{f}}=U_{\mathrm{i}}+p_{\mathrm{i}} V_{\mathrm{i}} \text { or } H_{\mathrm{f}}=H_{\mathrm{i}}
$$

Therefore, the expansion occurs without change of enthalpy.

For a perfect gas, $\mu=0$; hence, the temperature of a perfect gas is unchanged by Joule-Thomson expansion. This characteristic points clearly to the involvement of intermolecular forces in determining the size of the effect.

Real gases have non-zero Joule-Thomson coefficients. Depending on the identity of the gas, the pressure, the relative magnitudes of the attractive and repulsive intermolecular forces, and the temperature, the sign of the coefficient may be either positive or negative (Fig. 2D.9). A positive sign implies that $\mathrm{d} T$ is negative when $\mathrm{d} p$ is negative, in which case the gas cools on expansion. However, the Joule-Thomson coefficient of a real gas does not necessarily approach zero as the pressure is reduced even though the equation of state of the gas approaches that of a perfect gas. The coefficient behaves like the properties discussed in Topic 1C in the sense that it depends on derivatives and not on $p, V$, and $T$ themselves.

Gases that show a heating effect $(\mu<0)$ at one temperature show a cooling effect $(\mu>0)$ when the temperature is below their upper inversion temperature, $T_{\mathrm{I}}$ (Table 2D.2, Fig. 2D.10). As indicated in Fig. 2D.10, a gas typically has two inversion temperatures.\\
\includegraphics[max width=\textwidth, center]{2025_02_27_d4b95d9718f0b9e2e6b9g-097}

Figure 2D. 9 The sign of the Joule-Thomson coefficient, $\mu$, depends on the conditions. Inside the boundary, the blue area, it is positive and outside it is negative. The temperature corresponding to the boundary at a given pressure is the 'inversion temperature' of the gas at that pressure. Reduction of pressure under adiabatic conditions moves the system along one of the isenthalps, or curves of constant enthalpy (the blue lines). The inversion temperature curve runs through the points of the isenthalps where their slope changes from negative to positive.\\
\includegraphics[max width=\textwidth, center]{2025_02_27_d4b95d9718f0b9e2e6b9g-097(1)}

Figure 2D. 10 The inversion temperatures for three real gases, nitrogen, hydrogen, and helium.

\subsection*{(b) The molecular interpretation of the Joule-Thomson effect}
The kinetic model of gases (Topic 1B) and the equipartition theorem (The chemist's toolkit 7 of Topic 2 A ) jointly imply that the mean kinetic energy of molecules in a gas is proportional to the temperature. It follows that reducing the average speed of the molecules is equivalent to cooling the gas. If the speed of the molecules can be reduced to the point that neighbours can capture each other by their intermolecular attractions, then the cooled gas will condense to a liquid.

Slowing gas molecules makes use of an effect similar to that seen when a ball is thrown up into the air: as it rises it slows\\
in response to the gravitational attraction of the Earth and its kinetic energy is converted into potential energy. As seen in Topic 1C, molecules in a real gas attract each other (the attraction is not gravitational, but the effect is the same). It follows that, if the molecules move apart from each other, like a ball rising from a planet, then they should slow. It is very easy to move molecules apart from each other by simply allowing the gas to expand, which increases the average separation of the molecules. To cool a gas, therefore, expansion must occur without allowing any energy to enter from outside as heat. As the gas expands, the molecules move apart to fill the available volume, struggling as they do so against the attraction of\\
their neighbours. Because some kinetic energy must be converted into potential energy to reach greater separations, the molecules travel more slowly as their separation increases, and the temperature drops. The cooling effect, which corresponds to $\mu>0$, is observed in real gases under conditions when attractive interactions are dominant $(Z<1$, where $Z$ is the compression factor defined in eqn 1C.1, $\left.Z=V_{\mathrm{m}} / V_{\mathrm{m}}^{\circ}\right)$, because the molecules have to climb apart against the attractive force in order for them to travel more slowly. For molecules under conditions when repulsions are dominant ( $Z>1$ ), the Joule-Thomson effect results in the gas becoming warmer, or $\mu<0$.

\subsection*{Checklist of concepts}
\begin{enumerate}
  \item The quantity $\mathrm{d} U$ is an exact differential, $\mathrm{d} w$ and $\mathrm{d} q$ are not.2. The change in internal energy may be expressed in terms of changes in temperature and volume.
  \item The internal pressure is the variation of internal energy with volume at constant temperature.4. Joule's experiment showed that the internal pressure of a perfect gas is zero.
  \item The change in internal energy with pressure and temperature is expressed in terms of the internal pressure and the heat capacity and leads to a general expression for the relation between heat capacities.
  \item The Joule-Thomson effect is the change in temperature of a gas when it undergoes isenthalpic expansion.
\end{enumerate}

\subsection*{Checklist of equations}
\begin{center}
\begin{tabular}{llll}
\hline
Property & Equation & Comment & Equation number \\
\hline
Change in $U(V, T)$ & $\mathrm{d} U=(\partial U / \partial V)_{T} \mathrm{~d} V+(\partial U / \partial T)_{V} \mathrm{~d} T$ & Constant composition & 2D.3 \\
Internal pressure & $\pi_{T}=(\partial U / \partial V)_{T}$ & Definition; for a perfect gas, $\pi_{T}=0$ & 2 D .4 \\
Change in $U(V, T)$ & $\mathrm{d} U=\pi_{T} \mathrm{~d} V+C_{V} \mathrm{~d} T$ & Constant composition & 2 D .5 \\
Expansion coefficient & $\alpha=(1 / V)(\partial V / \partial T)_{p}$ & Definition & 2 D .6 \\
Isothermal compressibility & $\kappa_{T}=-(1 / V)(\partial V / \partial p)_{T}$ & Definition & 2D.7 \\
Relation between heat capacities & $C_{p}-C_{V}=n R$ & Perfect gas & 2 D .10 \\
 & $C_{p}-C_{V}=\alpha^{2} T V / \kappa_{T}$ & For a perfect gas, $\mu=0$ & 2 D .11 \\
Joule-Thomson coefficient & $\mu=(\partial T / \partial p)_{H}$ & Constant composition & 2 D .12 \\
Change in $H(p, T)$ & $\mathrm{d} H=-\mu C_{p} \mathrm{~d} p+C_{p} \mathrm{~d} T$ &  & 2 D .13 \\
\hline
\end{tabular}
\end{center}

\section*{TOPIC 2E Adiabatic changes}
\subsection*{- Why do you need to know this material?}
Adiabatic processes complement isothermal processes, and are used in the discussion of the Second Law of thermodynamics.

\subsection*{- What is the key idea?}
The temperature of a perfect gas falls when it does work in an adiabatic expansion.

\subsection*{What do you need to know already?}
This Topic makes use of the discussion of the properties of gases (Topic 1A), particularly the perfect gas law. It also uses the definition of heat capacity at constant volume (Topic 2A) and constant pressure (Topic 2B) and the relation between them (Topic 2D).

The temperature falls when a gas expands adiabatically (in a thermally insulated container). Work is done, but as no heat enters the system, the internal energy falls, and therefore the temperature of the working gas also falls. In molecular terms, the kinetic energy of the molecules falls as work is done, so their average speed decreases, and hence the temperature falls too.

\subsection*{2E. 1 The change in temperature}
The change in internal energy of a perfect gas when the temperature is changed from $T_{\mathrm{i}}$ to $T_{\mathrm{f}}$ and the volume is changed from $V_{\mathrm{i}}$ to $V_{\mathrm{f}}$ can be expressed as the sum of two steps (Fig. 2E.1). In the first step, only the volume changes and the temperature is held constant at its initial value. However, because the internal energy of a perfect gas is independent of the volume it occupies (Topic 2A), the overall change in internal energy arises solely from the second step, the change in temperature at constant volume. Provided the heat capacity is independent of temperature, the change in the internal energy is

$$
\Delta U=\left(T_{\mathrm{f}}-T_{\mathrm{i}}\right) C_{V}=C_{V} \Delta T
$$

Because the expansion is adiabatic, $q=0$; then because $\Delta U=$ $q+w$, it follows that $\Delta U=w_{\mathrm{ad}}$. The subscript 'ad' denotes an\\
\includegraphics[max width=\textwidth, center]{2025_02_27_d4b95d9718f0b9e2e6b9g-099}

Figure 2E. 1 To achieve a change of state from one temperature and volume to another temperature and volume, treat the overall change as composed of two steps. In the first step, the system expands at constant temperature; there is no change in internal energy if the system consists of a perfect gas. In the second step, the temperature of the system is reduced at constant volume. The overall change in internal energy is the sum of the changes for the two steps.\\
adiabatic process. Therefore, by equating the two expressions for $\Delta U$,


\begin{equation*}
w_{\mathrm{ad}}=C_{V} \Delta T \tag{2E.1}
\end{equation*}


Work of adiabatic change [perfect gas]

That is, the work done during an adiabatic expansion of a perfect gas is proportional to the temperature difference between the initial and final states. That is exactly what is expected on molecular grounds, because the mean kinetic energy is proportional to $T$, so a change in internal energy arising from temperature alone is also expected to be proportional to $\Delta T$. From these considerations it is possible to calculate the temperature change of a perfect gas that undergoes reversible adiabatic expansion (reversible expansion in a thermally insulated container).

\subsection*{How is that done? 2E. 1 Deriving an expression for the}
 temperature change in a reversible adiabatic expansionConsider a stage in a reversible adiabatic expansion of a perfect gas when the pressure inside and out is $p$. When considering reversible processes, it is usually appropriate to consider infinitesimal changes in the conditions, because pressures and temperatures typically change during the process. Then follow these steps.

Step 1 Write an expression relating temperature and volume changes\\
The work done when the gas expands reversibly by $\mathrm{d} V$ is $\mathrm{d} w=-p \mathrm{~d} V$. This expression applies to any reversible change, including an adiabatic change, so specifically $\mathrm{d} w_{\mathrm{ad}}=-p \mathrm{~d} V$. Therefore, because $\mathrm{d} q=0$ for an adiabatic change, $\mathrm{d} U=\mathrm{d} w_{\mathrm{ad}}$ (the infinitesimal version of $\Delta U=w_{\mathrm{ad}}$ ).\\
For a perfect gas, $\mathrm{d} U=C_{V} \mathrm{~d} T$ (the infinitesimal version of $\left.\Delta U=C_{V} \Delta T\right)$. Equating these expressions for $\mathrm{d} U$ gives

$$
C_{V} \mathrm{~d} T=-p \mathrm{~d} V
$$

Because the gas is perfect, $p$ can be replaced by $n R T / V$ to give $C_{V} \mathrm{~d} T=-(n R T / V) \mathrm{d} V$ and therefore

$$
\frac{C_{V} \mathrm{~d} T}{T}=-\frac{n R \mathrm{~d} V}{V}
$$

\subsection*{Step 2 Integrate the expression to find the overall change}
To integrate this expression, ensure that the limits of integration match on each side of the equation. Note that $T$ is equal to $T_{\mathrm{i}}$ when $V$ is equal to $V_{\mathrm{i}}$, and is equal to $T_{\mathrm{f}}$ when $V$ is equal to $V_{\mathrm{f}}$ at the end of the expansion. Therefore,

$$
C_{V} \int_{T_{i}}^{T_{i}} \frac{\mathrm{~d} T}{T}=-n R \int_{V_{\mathrm{i}}}^{V_{\mathrm{i}}} \frac{\mathrm{~d} V}{V}
$$

where $C_{V}$ is taken to be independent of temperature. Use Integral A. 2 in each case, and obtain

$$
C_{V} \ln \frac{T_{\mathrm{f}}}{T_{\mathrm{i}}}=-n R \ln \frac{V_{\mathrm{f}}}{V_{\mathrm{i}}}
$$

Step 3 Simplify the expression\\
Because $\ln (x / y)=-\ln (y / x)$, the preceding expression rearranges to

$$
\frac{C_{V}}{n R} \ln \frac{T_{\mathrm{f}}}{T_{\mathrm{i}}}=\ln \frac{V_{\mathrm{i}}}{V_{\mathrm{f}}}
$$

Next, note that $C_{V} / n R=C_{V, \mathrm{~m}} / R=c$ and use $\ln x^{a}=a \ln x$ to obtain

$$
\ln \left(\frac{T_{\mathrm{f}}}{T_{\mathrm{i}}}\right)^{c}=\ln \frac{V_{\mathrm{i}}}{V_{\mathrm{f}}}
$$

This relation implies that $\left(T_{\mathrm{f}} / T_{\mathrm{i}}\right)^{c}=\left(V_{\mathrm{i}} / V_{\mathrm{f}}\right)$ and, upon rearrangement,

$$
-\left|T_{\mathrm{f}}=T_{\mathrm{i}}\left(\frac{V_{\mathrm{i}}}{V_{\mathrm{f}}}\right)^{1 / c} c=C_{V, \mathrm{~m}} / R\right| \begin{aligned}
& \text { Temperature change } \\
& \begin{array}{l}
\text { reversible adiabatic } \\
\text { expansion, perfect gas] }
\end{array}
\end{aligned}
$$

By raising each side of this expression to the power $c$ and reorganizing it slightly, an equivalent expression is

\[
V_{\mathrm{i}} T_{\mathrm{i}}^{c}=V_{\mathrm{f}} T_{\mathrm{f}}^{c} \quad c=C_{V, \mathrm{~m}} / R \quad \begin{align*}
& \text { Temperature change } \\
& \text { [reversible adiabatic }  \tag{2E.2b}\\
& \text { expansion, perfect gas] }
\end{align*}
\]

This result is often summarized in the form $V T^{c}=$ constant.

\subsection*{Brief illustration 2E. 1}
Consider the adiabatic, reversible expansion of 0.020 mol Ar , initially at $25^{\circ} \mathrm{C}$, from $0.50 \mathrm{dm}^{3}$ to $1.00 \mathrm{dm}^{3}$. The molar heat capacity of argon at constant volume is $12.47 \mathrm{~J} \mathrm{~K}^{-1} \mathrm{~mol}^{-1}$, so $c=1.501$. Therefore, from eqn 2E. 2 a ,

$$
T_{\mathrm{f}}=(298 \mathrm{~K}) \times\left(\frac{0.50 \mathrm{dm}^{3}}{1.00 \mathrm{dm}^{3}}\right)^{1 / 1.501}=188 \mathrm{~K}
$$

It follows that $\Delta T=-110 \mathrm{~K}$, and therefore, from eqn 2E.1, that

$$
w_{\mathrm{ad}}=\left\{(0.020 \mathrm{~mol}) \times\left(12.47 \mathrm{~J} \mathrm{~K}^{-1} \mathrm{~mol}^{-1}\right)\right\} \times(-110 \mathrm{~K})=-27 \mathrm{~J}
$$

Note that temperature change is independent of the amount of gas but the work is not.

\subsection*{2E. 2 The change in pressure}
Equation 2E.2a may be used to calculate the pressure of a perfect gas that undergoes reversible adiabatic expansion.

\subsection*{How is that done? 2E. 2 Deriving the relation between}
 pressure and volume for a reversible adiabatic expansionThe initial and final states of a perfect gas satisfy the perfect gas law regardless of how the change of state takes place, so $p V=n R T$ can be used to write

$$
\frac{p_{\mathrm{i}} V_{\mathrm{i}}}{p_{\mathrm{f}} V_{\mathrm{f}}}=\frac{T_{\mathrm{i}}}{T_{\mathrm{f}}}
$$

However, $T_{\mathrm{i}} / T_{\mathrm{f}}=\left(V_{\mathrm{f}} / V_{\mathrm{i}}\right)^{1 / c}($ eqn 2E.2a). Therefore,

$$
\frac{p_{\mathrm{i}} V_{\mathrm{i}}}{p_{\mathrm{f}} V_{\mathrm{f}}}=\left(\frac{V_{\mathrm{f}}}{V_{\mathrm{i}}}\right)^{1 / c}, \text { so } \frac{p_{\mathrm{i}}}{p_{\mathrm{f}}}\left(\frac{V_{\mathrm{i}}}{V_{\mathrm{f}}}\right)^{\frac{1}{c}+1}=1
$$

For a perfect gas $C_{p, \mathrm{~m}}-C_{V, \mathrm{~m}}=R$ (Topic 2B). It follows that

$$
\frac{1}{c}+1=\frac{1+c}{c}=\frac{R+C_{V, \mathrm{~m}}}{C_{V, \mathrm{~m}}}=\frac{C_{p, \mathrm{~m}}}{C_{V, \mathrm{~m}}}=\gamma
$$

and therefore that

$$
\frac{p_{\mathrm{i}}}{p_{\mathrm{f}}}\left(\frac{V_{\mathrm{i}}}{V_{\mathrm{f}}}\right)^{\gamma}=1
$$

which rearranges to\\
\includegraphics[max width=\textwidth, center]{2025_02_27_d4b95d9718f0b9e2e6b9g-100}

This result is commonly summarized in the form $p V^{\gamma}=$ constant.\\
\includegraphics[max width=\textwidth, center]{2025_02_27_d4b95d9718f0b9e2e6b9g-101}

Figure 2E. 2 An adiabat depicts the variation of pressure with volume when a gas expands adiabatically and, in this case, reversibly. Note that the pressure declines more steeply for an adiabat than it does for an isotherm because in an adiabatic change the temperature falls.

For a monatomic perfect gas, $C_{V, \mathrm{~m}}=\frac{3}{2} R$ (Topic 2 A ), and $C_{p, \mathrm{~m}}=\frac{5}{2} R\left(\right.$ from $C_{p, \mathrm{~m}}-C_{V, \mathrm{~m}}=R$ ), so $\gamma=\frac{5}{3}$. For a gas of nonlinear polyatomic molecules (which can rotate as well as translate;\\
vibrations make little contribution at normal temperatures), $C_{V, \mathrm{~m}}=3 R$ and $C_{p, \mathrm{~m}}=4 R$, so $\gamma=\frac{4}{3}$. The curves of pressure versus volume for adiabatic change are known as adiabats, and one for a reversible path is illustrated in Fig. 2E.2. Because $\gamma>1$, an adiabat falls more steeply $\left(p \propto 1 / V^{\gamma}\right)$ than the corresponding isotherm $(p \propto 1 / V)$. The physical reason for the difference is that, in an isothermal expansion, energy flows into the system as heat and maintains the temperature; as a result, the pressure does not fall as much as in an adiabatic expansion.

\subsection*{Brief illustration 2E. 2}
When a sample of argon (for which $\gamma=\frac{5}{3}$ ) at 100 kPa expands reversibly and adiabatically to twice its initial volume the final pressure will be

$$
p_{\mathrm{f}}=\left(\frac{V_{\mathrm{i}}}{V_{\mathrm{f}}}\right)^{\gamma} p_{\mathrm{i}}=\left(\frac{1}{2}\right)^{5 / 3} \times(100 \mathrm{kPa})=31 \mathrm{kPa}
$$

For an isothermal expansion in which the volume doubles the final pressure would be 50 kPa .

\subsection*{Checklist of concepts}
\begin{enumerate}
  \item The temperature of a gas falls when it undergoes an adiabatic expansion in which work is done.2. An adiabat is a curve showing how pressure varies with volume in an adiabatic process.
\end{enumerate}

\subsection*{Checklist of equations}
\begin{center}
\begin{tabular}{llll}
\hline
Property & Equation & Comment & Equation number \\
\hline
Work of adiabatic expansion & $w_{\mathrm{ad}}=C_{V} \Delta T$ & Perfect gas & 2E.1 \\
Final temperature & $T_{\mathrm{f}}=T_{\mathrm{i}}\left(V_{\mathrm{i}} / V_{\mathrm{f}}\right)^{1 / c}$ & \begin{tabular}{l}
Perfect gas, reversible adiabatic \\
expansion \\
\end{tabular} & 2 E .2 a \\
 & $c=C_{V, \mathrm{~m}} / R$ &  & 2 E .2 b \\
Adiabats & $V_{\mathrm{i}} T_{\mathrm{i}}^{c}=V_{\mathrm{f}} T_{\mathrm{f}}^{c}$ &  & 2E.3 \\
 & $p_{\mathrm{f}} V_{\mathrm{f}}^{\gamma}=p_{\mathrm{i}} V_{\mathrm{i}}^{\gamma}$ &  &  \\
\hline
\end{tabular}
\end{center}

\chapter{FOCUS 3 The Second and Third Laws}
Some things happen naturally, some things don't. Some aspect of the world determines the spontaneous direction of change, the direction of change that does not require work to bring it about. An important point, though, is that throughout this text 'spontaneous' must be interpreted as a natural tendency which might or might not be realized in practice. Thermodynamics is silent on the rate at which a spontaneous change in fact occurs, and some spontaneous processes (such as the conversion of diamond to graphite) may be so slow that the tendency is never realized in practice whereas others (such as the expansion of a gas into a vacuum) are almost instantaneous.

\subsection*{3A Entropy}
The direction of change is related to the distribution of energy and matter, and spontaneous changes are always accompanied by a dispersal of energy or matter. To quantify this concept we introduce the property called 'entropy', which is central to the formulation of the 'Second Law of thermodynamics'. That law governs all spontaneous change.\\
3A. 1 The Second Law; 3A. 2 The definition of entropy;\\
3A. 3 The entropy as a state function

\subsection*{3B Entropy changes accompanying specific processes}
This Topic shows how to use the definition of entropy change to calculate its value for a number of common physical processes, such as the expansion of a gas, a phase transition, and heating a substance.\\
3B. 1 Expansion; 3B. 2 Phase transitions; 3B. 3 Heating; 3B. 4 Composite processes

\subsection*{3C The measurement of entropy}
To make the Second Law quantitative, it is necessary to measure the entropy of a substance. The measurement of heat capacities, and the energy transferred as heat during physical\\
processes, makes it possible to determine the entropies of substances. The discussion in this Topic also leads to the 'Third Law of thermodynamics', which relates to the properties of matter at very low temperatures and is used to set up an absolute measure of the entropy of a substance.\\
3C. 1 The calorimetric measurement of entropy; 3C. 2 The Third Law

\subsection*{3D Concentrating on the system}
One problem with dealing with the entropy is that it requires separate calculations of the changes taking place in the system and the surroundings. Providing certain restrictions on the system can be accepted, that problem can be overcome by introducing the 'Gibbs energy'. Indeed, most thermodynamic calculations in chemistry focus on the change in Gibbs energy rather than the entropy change itself.\\
3D. 1 The Helmholtz and Gibbs energies; 3D. 2 Standard molar Gibbs energies

\subsection*{3E Combining the First and Second Laws}
In this Topic the First and Second Laws are combined, which leads to a very powerful way of applying thermodynamics to the properties of matter.\\
3E. 1 Properties of the internal energy; 3E. 2 Properties of the Gibbs energy

\subsection*{Web resources What are the applications of this material?}
The Second Law is at the heart of the operation of engines of all types, including devices resembling engines that are used to cool objects. See Impact 4 on the website of this book for an application to the technology of refrigeration. Entropy considerations are also important in modern electronic materials for they permit a quantitative discussion of the concentration of impurities. See Impact 5 for a note about how measurement of the entropy at low temperatures gives insight into the purity of materials used as superconductors.

\section*{TOPIC 3A Entropy}
\subsection*{> Why do you need to know this material?}
Entropy is the concept on which almost all applications of thermodynamics in chemistry are based: it explains why some physical transformations and chemical reactions are spontaneous and others are not.

\subsection*{- What is the key idea?}
The change in entropy of a system can be calculated from the heat transferred to it reversibly; a spontaneous process in an isolated system is accompanied by an increase in entropy.

\subsection*{> What do you need to know already?}
You need to be familiar with the First-Law concepts of work, heat, and internal energy (Topic 2A). The Topic draws on the expression for work of expansion of a perfect gas (Topic 2A) and on the changes in volume and temperature that accompany the reversible adiabatic expansion of a perfect gas (Topic 2E).

What determines the direction of spontaneous change? It is not a tendency to achieve a lower energy, because the First Law asserts that the total energy of the universe does not change in any process. It turns out that the direction is determined by the manner in which energy and matter are distributed. This concept is made precise by the Second Law of thermodynamics and made quantitative by introducing the property known as 'entropy'.

\subsection*{3A. 1 The Second Law}
The role of the distribution of energy and matter can be appreciated by thinking about a ball bouncing on a floor. The ball does not rise as high after each bounce because some of the energy associated with its motion spreads out-is dispersedinto the thermal motion of the particles in the ball and the floor. The direction of spontaneous change is towards a state in which the ball is at rest with all its energy dispersed into disorderly thermal motion of the particles in the surroundings (Fig. 3A.1).\\
\includegraphics[max width=\textwidth, center]{2025_02_27_d4b95d9718f0b9e2e6b9g-110}

Figure 3A. 1 The direction of spontaneous change for a ball bouncing on a floor. On each bounce some of its energy is degraded into the thermal motion of the atoms of the floor, and that energy then disperses. The reverse process, a ball rising from the floor as a result of acquiring energy from the thermal motion of the atoms in the floor, has never been observed to take place.

A ball resting on a warm floor has never been observed to start bouncing as a result of energy transferred to the ball from the floor. For bouncing to begin, something rather special would need to happen. In the first place, some of the thermal motion of the atoms in the floor (the surroundings) would have to accumulate in a single, small object, the ball (the system). This accumulation requires a spontaneous localization of energy from the myriad vibrations of the atoms of the floor into the much smaller number of atoms that constitute the ball (Fig. 3A.2). Furthermore, whereas the thermal motion is random, for the ball to move upwards its atoms must all move in the same direction. The localization of random, disorderly motion as directed, orderly motion is so unlikely that it can be dismissed as virtually impossible. ${ }^{1}$

The signpost of spontaneous change has been identified: look for the direction of change that leads to the dispersal of energy. This principle accounts for the direction of change of the bouncing ball, because its energy is spread out as thermal motion of the atoms of the floor. The reverse process is not spontaneous because it is highly improbable that energy will become localized, leading to uniform motion of the ball's atoms.

\footnotetext{${ }^{1}$ Orderly motion, but on a much smaller scale and continued only very briefly, is observed as Brownian motion, the jittering motion of small particles suspended in a liquid or gas.
}
\includegraphics[max width=\textwidth, center]{2025_02_27_d4b95d9718f0b9e2e6b9g-111(1)}

Figure 3A. 2 (a) A ball resting on a warm surface; the atoms are undergoing thermal motion (vibration, in this instance), as indicated by the arrows. (b) For the ball to fly upwards, some of the random vibrational motion would have to change into coordinated, directed motion. Such a conversion is highly improbable.

Matter also has a tendency to disperse. A gas does not contract spontaneously because to do so the random motion of its molecules would have to take them all into the same region of the container. The opposite change, spontaneous expansion, is a natural consequence of matter becoming more dispersed as the gas molecules are free to occupy a larger volume.

The Second Law of thermodynamics expresses these conclusions more precisely and without referring to the behaviour of the molecules that are responsible for the properties of bulk matter. One statement was formulated by Kelvin:

\begin{displayquote}
No process is possible in which the sole result is the absorption of heat from a reservoir and its complete conversion into work.
\end{displayquote}

Statements like this are commonly explored by thinking about an idealized device called a heat engine (Fig. 3A.3(a)). A heat engine consists of two reservoirs, one hot (the 'hot source') and one cold (the 'cold sink'), connected in such a way that some of the energy flowing as heat between the two reservoirs can be converted into work. The Kelvin statement implies that it is not possible to construct a heat engine in which all the heat drawn from the hot source is completely converted into work (Fig. 3A.3(b)): all working heat engines must have a cold sink. The Kelvin statement is a generalization of the everyday observation that a ball at rest on a surface has never been observed to leap spontaneously upwards. An upward leap of the ball would be equivalent to the spontaneous conversion of heat from the surface into the work of raising the ball.

Another statement of the Second Law is due to Rudolf Clausius (Fig. 3A.4):

\footnotetext{Heat does not flow spontaneously from a cool body to a hotter body.
}
\includegraphics[max width=\textwidth, center]{2025_02_27_d4b95d9718f0b9e2e6b9g-111}

Figure 3A. 3 (a) A heat engine is a device in which energy is extracted from a hot reservoir (the hot source) as heat and then some of that energy is converted into work and the rest discarded into a cold reservoir (the cold sink) as heat. (b) The Kelvin statement of the Second Law denies the possibility of the process illustrated here, in which heat is changed completely into work, there being no other change.

To achieve the transfer of heat to a hotter body, it is necessary to do work on the system, as in a refrigerator. Although they appear somewhat different, it can be shown that the Clausius statement is logically equivalent to the Kelvin statement. One way to do so is to show that the two observations can be summarized by a single statement.

First, the system and its surroundings are regarded as a single (and possibly huge) isolated system sometimes referred to as 'the universe'. Energy can be transferred within this isolated system between the actual system and its surroundings, but none can enter or leave it. Then the Second Law is expressed in terms of a new state function, the entropy, $S$ :\\
\includegraphics[max width=\textwidth, center]{2025_02_27_d4b95d9718f0b9e2e6b9g-111(2)}

Figure 3A. 4 According to the Clausius statement of the Second Law, the process shown here, in which energy as heat migrates from a cool source to a hot sink, does not take place spontaneously. The process is not in conflict with the First Law because energy is conserved.

The entropy of an isolated system increases in the course of a spontaneous change: $\Delta S_{\text {tot }}>0$\\
where $S_{\text {tot }}$ is the total entropy of the overall isolated system. That is, if $S$ is the entropy of the system of interest, and $S_{\text {sur }}$ the entropy of the surroundings, then $S_{\text {tot }}=S+S_{\text {sur }}$. It is vitally important when considering applications of the Second Law to remember that it is a statement about the total entropy of the overall isolated system (the 'universe'), not just about the entropy of the system of interest. The following section defines entropy and interprets it as a measure of the dispersal of energy and matter, and relates it to the empirical observations discussed so far.

In summary, the First Law uses the internal energy to identify permissible changes; the Second Law uses the entropy to identify which of these permissible changes are spontaneous.

\subsection*{3A. 2 The definition of entropy}
To make progress, and to turn the Second Law into a quantitatively useful expression, the entropy change accompanying various processes needs to be defined and calculated. There are two approaches, one classical and one molecular. They turn out to be equivalent, but each one enriches the other.

\subsection*{(a) The thermodynamic definition of entropy}
The thermodynamic definition of entropy concentrates on the change in entropy, $\mathrm{d} S$, that occurs as a result of a physical or chemical change (in general, as a result of a 'process'). The definition is motivated by the idea that a change in the extent to which energy is dispersed in a disorderly way depends on how much energy is transferred as heat, not as work. As explained in Topic 2A, heat stimulates random motion of atoms whereas work stimulates their uniform motion and so does not change the extent of their disorder.\\
The thermodynamic definition of entropy is based on the expression

$$
\mathrm{d} S=\frac{\mathrm{d} q_{\mathrm{rev}}}{T}
$$

Entropy change [definition]\\
(3A.1a)\\
where $q_{\text {rev }}$ is the energy transferred as heat reversibly to the system at the absolute temperature $T$. For a measurable change between two states i and f,


\begin{equation*}
\Delta S=\int_{\mathrm{i}}^{\mathrm{f}} \frac{\mathrm{~d} q_{\mathrm{rev}}}{T} \tag{3A.1b}
\end{equation*}


That is, to calculate the difference in entropy between any two states of a system, find a reversible path between them, and integrate the energy supplied as heat at each stage of the path divided by the temperature at which that heat is transferred.

According to the definition of an entropy change given in eqn 3A.1a, when the energy transferred as heat is expressed in joules and the temperature is in kelvins, the units of entropy are joules per kelvin $\left(\mathrm{J} \mathrm{K}^{-1}\right)$. Entropy is an extensive property. Molar entropy, the entropy divided by the amount of substance, $S_{\mathrm{m}}=S / n$, is expressed in joules per kelvin per mole ( $\mathrm{JK}^{-1} \mathrm{~mol}^{-1}$ ); molar entropy is an intensive property.

\subsection*{Example 3A. 1 Calculating the entropy change for the isothermal expansion of a perfect gas}
Calculate the entropy change of a sample of perfect gas when it expands isothermally from a volume $V_{\mathrm{i}}$ to a volume $V_{\mathrm{f}}$.

Collect your thoughts The definition of entropy change in eqn 3A.1b instructs you to find the energy supplied as heat for a reversible path between the stated initial and final states regardless of the actual manner in which the process takes place. The process is isothermal, so $T$ can be treated as a constant and taken outside the integral in eqn 3A.1b. Moreover, because the internal energy of a perfect gas is independent of its volume (Topic 2 A ), $\Delta U=0$ for the expansion. Then, because $\Delta U=q+w$, it follows that $q=-w$, and therefore that $q_{\mathrm{rev}}=-w_{\mathrm{rev}}$. The work of reversible isothermal expansion is calculated in Topic 2A. Finally, calculate the change in molar entropy from $\Delta S_{\mathrm{m}}=\Delta S / n$.

The solution The temperature is constant, so eqn 3A.1b becomes

$$
\Delta S=\frac{1}{T} \int_{\mathrm{i}}^{\mathrm{f}} \mathrm{~d} q_{\mathrm{rev}}=\frac{q_{\mathrm{rev}}}{T}
$$

From Topic 2A the reversible work in an isothermal expansion is $w_{\text {rev }}=-n R T \ln \left(V_{\mathrm{f}} / V_{\mathrm{i}}\right)$, hence $q_{\text {rev }}=n R T \ln \left(V_{\mathrm{f}} / V_{\mathrm{i}}\right)$. It follows, after dividing $q_{\text {rev }}$ by $T$, that

$$
\Delta S=n R \ln \frac{V_{\mathrm{f}}}{V_{\mathrm{i}}} \text { and } \Delta S_{\mathrm{m}}=R \ln \frac{V_{\mathrm{f}}}{V_{\mathrm{i}}}
$$

Self-test 3A. 1 Calculate the change in entropy when the pressure of a fixed amount of perfect gas is changed isothermally from $p_{\mathrm{i}}$ to $p_{\mathrm{f}}$. What is the origin of this change?\\
\includegraphics[max width=\textwidth, center]{2025_02_27_d4b95d9718f0b9e2e6b9g-112}\\
\includegraphics[max width=\textwidth, center]{2025_02_27_d4b95d9718f0b9e2e6b9g-112(1)}

To see how the definition in eqn 3A.1a is used to formulate an expression for the change in entropy of the surroundings, $\Delta S_{\text {sur }}$, consider an infinitesimal transfer of heat $\mathrm{d} q_{\text {sur }}$ from the system to the surroundings. The surroundings consist of a reservoir of constant volume, so the energy supplied to them by heating can be identified with the change in the internal energy of the surroundings, $\mathrm{d} U_{\text {sur }}{ }^{2}$ The internal energy is a state function, and $\mathrm{d} U_{\text {sur }}$ is an exact differential. These properties

\footnotetext{${ }^{2}$ Alternatively, the surroundings can be regarded as being at constant pressure, in which case $\mathrm{d} q_{\text {sur }}=\mathrm{d} H_{\text {sur }}$.
}
imply that $\mathrm{d} U_{\text {sur }}$ is independent of how the change is brought about and in particular it is independent of whether the process is reversible or irreversible. The same remarks therefore apply to $\mathrm{d} q_{\text {sur }}$, to which $\mathrm{d} U_{\text {sur }}$ is equal. Therefore, the definition in eqn \href{http://3A.la}{3A.la} can be adapted simply by deleting the constraint 'reversible' and writing\\
\$\$

\begin{equation*}
\mathrm{d} S_{\text {sur }}=\frac{\mathrm{d} q_{\text {sur }}}{T_{\text {sur }}} \tag{3A.2a}
\end{equation*}

\$\$

Entropy change of the surroundings\\
Furthermore, because the temperature of the surroundings is constant whatever the change, for a measurable change


\begin{equation*}
\Delta S_{\mathrm{sur}}=\frac{q_{\mathrm{sur}}}{T_{\mathrm{sur}}} \tag{3A.2b}
\end{equation*}


That is, regardless of how the change is brought about in the system, reversibly or irreversibly, the change of entropy of the surroundings is calculated simply by dividing the heat transferred by the temperature at which the transfer takes place.

Equation 3A.2b makes it very simple to calculate the changes in entropy of the surroundings that accompany any process. For instance, for any adiabatic change, $q_{\text {sur }}=0$, so


\begin{equation*}
\Delta S_{\text {sur }}=0 \tag{3A.3}
\end{equation*}


Adiabatic change\\
This expression is true however the change takes place, reversibly or irreversibly, provided no local hot spots are formed in the surroundings. That is, it is true (as always assumed) provided the surroundings remain in internal equilibrium. If hot spots do form, then the localized energy may subsequently disperse spontaneously and hence generate more entropy.

\subsection*{Brief illustration 3A. 1}
To calculate the entropy change in the surroundings when $1.00 \mathrm{~mol} \mathrm{H}_{2} \mathrm{O}(\mathrm{l})$ is formed from its elements under standard conditions at 298 K , use $\Delta_{\mathrm{f}} H^{\ominus}=-286 \mathrm{~kJ} \mathrm{~mol}^{-1}$ from Table 2C.4. The energy released as heat from the system is supplied to the surroundings, so $q_{\text {sur }}=+286 \mathrm{~kJ}$. Therefore,

$$
\Delta S_{\text {sur }}=\frac{2.86 \times 10^{5} \mathrm{~J}}{298 \mathrm{~K}}=+960 \mathrm{JK}^{-1}
$$

This strongly exothermic reaction results in an increase in the entropy of the surroundings as energy is released as heat into them.

You are now in a position to see how the definition of entropy is consistent with Kelvin's and Clausius's statements of the Second Law and unifies them. In Fig. 3A.3(b) the entropy of the hot source is reduced as energy leaves it as heat. The transfer of energy as work does not result in the production of entropy, so the overall result is that the entropy of the (overall isolated) system decreases. The Second Law asserts that such\\
a process is not spontaneous, so the arrangement shown in Fig. 3A.3(b) does not produce work. In the Clausius version, the entropy of the cold source in Fig 3A. 4 decreases when energy leaves it as heat, but when that heat enters the hot sink the rise in entropy is not as great (because the temperature is higher). Overall there is a decrease in entropy and so the transfer of heat from a cold source to a hot sink is not spontaneous.

\subsection*{(b) The statistical definition of entropy}
The molecular interpretation of the Second Law and the 'statistical' definition of entropy start from the idea, introduced in the Prologue, that atoms and molecules are distributed over the energy states available to them in accord with the Boltzmann distribution. Then it is possible to predict that as the temperature is increased the molecules populate higher energy states. Boltzmann proposed that there is a link between the spread of molecules over the available energy states and the entropy, which he expressed as ${ }^{3}$


\begin{equation*}
S=k \ln \mathscr{W} \quad \text { Boltzmann formula for the entropy } \tag{3A.4}
\end{equation*}


where $k$ is Boltzmann's constant $\left(k=1.381 \times 10^{-23} \mathrm{~J} \mathrm{~K}^{-1}\right)$ and $\mathcal{W}$ is the number of microstates, the number of ways in which the molecules of a system can be distributed over the energy states for a specified total energy. When the properties of a system are measured, the outcome is an average taken over the many microstates the system can occupy under the prevailing conditions. The concept of the number of microstates makes quantitative the ill-defined qualitative concepts of 'disorder' and 'the dispersal of matter and energy' used to introduce the concept of entropy: a more disorderly distribution of matter and a greater dispersal of energy corresponds to a greater number of microstates associated with the same total energy. This point is discussed in much greater detail in Topic 13E.

Equation 3A. 4 is known as the Boltzmann formula and the entropy calculated from it is called the statistical entropy. If all the molecules are in one energy state there is only one way of achieving this distribution, so $\mathcal{W}=1$ and, because $\ln 1=0$, it follows that $S=0$. As the molecules spread out over the available energy states, $\mathcal{W}$ increases and therefore so too does the entropy. The value of $\mathcal{W}$ also increases if the separation of energy states decreases, because more states become accessible. An example is a gas confined to a container, because its translational energy levels get closer together as the container expands (Fig. 3A.5; this is a conclusion from quantum theory which is verified in Topic 7D). The value of $\mathcal{W}$, and hence the entropy, is expected to increase as the gas expands, which is in accord with the conclusion drawn from the thermodynamic definition of entropy (Example 3A.1).

\footnotetext{${ }^{3}$ He actually wrote $\mathrm{S}=\mathrm{k} \log \mathrm{W}$, and it is carved on his tombstone in Vienna.
}
\includegraphics[max width=\textwidth, center]{2025_02_27_d4b95d9718f0b9e2e6b9g-114(2)}

Figure 3A. 5 When a container expands from (b) to (a), the translational energy levels of gas molecules in it come closer together and, for the same temperature, more become accessible to the molecules. As a result the number of ways of achieving the same energy (the value of $\mathcal{W}$ ) increases, and so therefore does the entropy.

The molecular interpretation of entropy helps to explain why, in the thermodynamic definition given by eqn 3A.1, the entropy change depends inversely on the temperature. In a system at high temperature the molecules are spread out over a large number of energy states. Increasing the energy of the system by the transfer of heat makes more states accessible, but given that very many states are already occupied the proportionate change in $\mathcal{W}$ is small (Fig. 3A.6). In contrast, for a system at a low temperature fewer states are occupied, and so the transfer of the same energy results in a proportionately larger increase in the number of accessible states, and hence a larger increase in $\mathcal{W}$. This argument suggests that the change in entropy for a given transfer of energy as heat should be greater at low temperatures than at high, as in eqn 3A.1a.

There are several final points. One is that the Boltzmann definition of entropy makes it possible to calculate the absolute\\
\includegraphics[max width=\textwidth, center]{2025_02_27_d4b95d9718f0b9e2e6b9g-114}

Figure 3A. 6 The supply of energy as heat to the system results in the molecules moving to higher energy states, so increasing the number of microstates and hence the entropy. The increase in the entropy is smaller for (a) a system at a high temperature than (b) one at a low temperature because initially the number of occupied states is greater.\\
value of the entropy of a system, whereas the thermodynamic definition leads only to values for a change in entropy. This point is developed in Focus 13 where it is shown how to relate values of $S$ to the structural properties of atoms and molecules. The second point is that the Boltzmann formula cannot readily be applied to the surroundings, which are typically far too complex for $\mathcal{W}$ to be a meaningful quantity.

\subsection*{3A.3 The entropy as a state function}
Entropy is a state function. To prove this assertion, it is necessary to show that the integral of $\mathrm{d} S$ between any two states is independent of the path between them. To do so, it is sufficient to prove that the integral of eqn 3A.1a round an arbitrary cycle is zero, for that guarantees that the entropy is the same at the initial and final states of the system regardless of the path taken between them (Fig. 3A.7). That is, it is necessary to show that


\begin{equation*}
\oint \mathrm{d} S=\oint \frac{\mathrm{d} q_{\mathrm{rev}}}{T}=0 \tag{3A.5}
\end{equation*}


where the symbol $\oint$ denotes integration around a closed path. There are three steps in the argument:

\begin{enumerate}
  \item First, to show that eqn 3 A .5 is true for a special cycle (a 'Carnot cycle') involving a perfect gas.
  \item Then to show that the result is true whatever the working substance.
  \item Finally, to show that the result is true for any cycle.
\end{enumerate}

\subsection*{(a) The Carnot cycle}
A Carnot cycle, which is named after the French engineer Sadi Carnot, consists of four reversible stages in which a gas (the working substance) is either expanded or compressed in\\
\includegraphics[max width=\textwidth, center]{2025_02_27_d4b95d9718f0b9e2e6b9g-114(1)}

Figure 3A. 7 In a thermodynamic cycle, the overall change in a state function (from the initial state to the final state and then back to the initial state again) is zero.\\
\includegraphics[max width=\textwidth, center]{2025_02_27_d4b95d9718f0b9e2e6b9g-115}

Figure 3A. 8 The four stages which make up the Carnot cycle. In stage 1 the gas (the working substance) is in thermal contact with the hot reservoir, and in stage 3 contact is with the cold reservoir; both stages are isothermal. Stages 2 and 4 are adiabatic, with the gas isolated from both reservoirs.\\
various ways; in two of the stages energy as heat is transferred to or from a hot source or a cold sink (Fig. 3A.8).

Figure 3A. 9 shows how the pressure and volume change in each stage:

\begin{enumerate}
  \item The gas is placed in thermal contact with the hot source (which is at temperature $T_{\mathrm{h}}$ ) and undergoes reversible isothermal expansion from A to B ; the entropy change is $q_{\mathrm{h}} / T_{\mathrm{h}}$, where $q_{\mathrm{h}}$ is the energy supplied to the system as heat from the hot source.
  \item Contact with the hot source is broken and the gas then undergoes reversible adiabatic expansion from B to C. No energy leaves the system as heat, so the change in entropy is zero. The expansion is carried on until the temperature of the gas falls from $T_{\mathrm{h}}$ to $T_{\mathrm{c}}$, the temperature of the cold sink.
  \item The gas is placed in contact with the cold sink and then undergoes a reversible isothermal compression from C to D at $T_{\mathrm{c}}$. Energy is released as heat to the cold sink; the change in entropy of the system is $q_{c} / T_{c}$; in this expression $q_{c}$ is negative.
  \item Finally, contact with the cold sink is broken and the gas then undergoes reversible adiabatic compression from D to A such that the final temperature is $T_{\mathrm{h}}$. No energy enters the system as heat, so the change in entropy is zero.
\end{enumerate}

The total change in entropy around the cycle is the sum of the changes in each of these four steps:

$$
\oint \mathrm{d} S=\frac{q_{\mathrm{h}}}{T_{\mathrm{h}}}+\frac{q_{\mathrm{c}}}{T_{\mathrm{c}}}
$$

The next task is to show that the sum of the two terms on the right of this expression is zero for a perfect gas and so confirming, for that substance at least, that entropy is a state function.\\
\includegraphics[max width=\textwidth, center]{2025_02_27_d4b95d9718f0b9e2e6b9g-115(1)}

Figure 3A. 9 The basic structure of a Carnot cycle. Stage 1 is the isothermal reversible expansion at the temperature $T_{h}$. Stage 2 is a reversible adiabatic expansion in which the temperature falls from $T_{\mathrm{h}}$ to $T_{\mathrm{c}}$. Stage 3 is an isothermal reversible compression at $T_{\mathrm{c}}$. Stage 4 is an adiabatic reversible compression, which restores the system to its initial state.

\subsection*{How is that done? 3A. 1 Showing that the entropy is a}
 state function for a perfect gasFirst, you need to note that a reversible adiabatic expansion (stage 2 in Fig. 3A.9) takes the system from $T_{\mathrm{h}}$ to $T_{c}$. You can then use the properties of such an expansion, specifically $V T^{c}$ $=$ constant (Topic 2E), to relate the two volumes at the start and end of the expansion. You also need to note that energy as heat is transferred by reversible isothermal processes (stages 1 and 3) and, as derived in Example 3A.1, for a perfect gas

$$
\overbrace{q_{\mathrm{h}}=n R T_{\mathrm{h}} \ln \frac{V_{\mathrm{B}}}{V_{\mathrm{A}}}}^{\text {stage }} \overbrace{q_{\mathrm{c}}=n R T_{\mathrm{c}} \ln \frac{V_{\mathrm{D}}}{V_{\mathrm{C}}}}^{\text {stage }}
$$

Step 1 Relate the volumes in the adiabatic expansions\\
For a reversible adiabatic process the temperature and volume are related by $V T^{c}=$ constant (Topic 2E). Therefore\\
for the path D to $\mathrm{A}($ stage 4$): V_{\mathrm{A}} T_{\mathrm{h}}^{c}=V_{\mathrm{D}} T_{\mathrm{c}}^{c}$\\
for the path B to C (stage 2): $V_{\mathrm{C}} T_{\mathrm{c}}^{c}=V_{\mathrm{B}} T_{\mathrm{h}}^{c}$\\
Multiplication of the first of these expressions by the second gives

$$
V_{\mathrm{A}} V_{\mathrm{C}} T_{\mathrm{h}}^{c} T_{\mathrm{c}}^{c}=V_{\mathrm{D}} V_{\mathrm{B}} T_{\mathrm{h}}^{c} T_{\mathrm{c}}^{c}
$$

which, on cancellation of the temperatures, simplifies to

$$
\frac{V_{\mathrm{D}}}{V_{\mathrm{C}}}=\frac{V_{\mathrm{A}}}{V_{\mathrm{B}}}
$$

Step 2 Establish the relation between the two heat transfers\\
You can now use this relation to write an expression for energy discarded as heat to the cold sink in terms of $V_{\mathrm{A}}$ and $V_{\mathrm{B}}$

$$
q_{\mathrm{c}}=n R T_{\mathrm{c}} \ln \frac{V_{\mathrm{D}}}{V_{\mathrm{C}}}=n R T_{\mathrm{c}} \ln \frac{V_{\mathrm{A}}}{V_{\mathrm{B}}}=-n R T_{\mathrm{c}} \ln \frac{V_{\mathrm{B}}}{V_{\mathrm{A}}}
$$

\subsection*{It follows that}
$$
\frac{q_{\mathrm{h}}}{q_{\mathrm{c}}}=\frac{n R T_{\mathrm{h}} \ln \left(V_{\mathrm{B}} / V_{\mathrm{A}}\right)}{-n R T_{\mathrm{c}} \ln \left(V_{\mathrm{B}} / V_{\mathrm{A}}\right)}=-\frac{T_{\mathrm{h}}}{T_{\mathrm{c}}}
$$

Note that $q_{\mathrm{h}}$ is negative (heat is withdrawn from the hot source) and $q_{c}$ is positive (heat is deposited in the cold sink), so their ratio is negative. This expression can be rearranged into


\begin{equation*}
\frac{q_{\mathrm{h}}}{T_{\mathrm{h}}}+\frac{q_{\mathrm{c}}}{T_{\mathrm{c}}}=0 \tag{3A.6}
\end{equation*}


Because the total change in entropy around the cycle is $q_{\mathrm{h}} / T_{\mathrm{h}}+q_{\mathrm{c}} / T_{\mathrm{c}}$, it follows immediately from eqn 3A. 6 that, for a perfect gas, this entropy change is zero.

\subsection*{Brief illustration 3A. 2}
The Carnot cycle can be regarded as a representation of the changes taking place in a heat engine in which part of the energy extracted as heat from the hot reservoir is converted into work. Consider an engine running in accord with the Carnot cycle, and in which 100 J of energy is withdrawn from the hot source $\left(q_{\mathrm{h}}=-100 \mathrm{~J}\right)$ at 500 K . Some of this energy is used to do work and the remainder is deposited in the cold sink at 300 K . According to eqn 3A.6, the heat deposited is

$$
q_{\mathrm{c}}=-q_{\mathrm{h}} \times \frac{T_{\mathrm{c}}}{T_{\mathrm{h}}}=-(-100 \mathrm{~J}) \times \frac{300 \mathrm{~K}}{500 \mathrm{~K}}=+60 \mathrm{~J}
$$

This value implies that 40 J was used to do work.

It is now necessary to show that eqn 3A. 5 applies to any material, not just a perfect gas. To do so, it is helpful to introduce the efficiency, $\eta$ (eta), of a heat engine:

\[
\eta=\frac{\text { work performed }}{\text { heat absorbed from hot source }}=\frac{|w|}{\left|q_{\mathrm{h}}\right|} \quad \begin{align*}
& \text { Efficiency }  \tag{3A.7}\\
& {[\text { definition }]}
\end{align*}
\]

Modulus signs (|...|) have been used to avoid complications with signs: all efficiencies are positive numbers. The definition implies that the greater the work output for a given supply of heat from the hot source, the greater is the efficiency of the engine. The definition can be expressed in terms of the heat transactions alone, because (as shown in Fig. 3A.10) the energy supplied as work by the engine is the difference between the energy supplied as heat by the hot source and that returned to the cold sink:


\begin{equation*}
\eta=\frac{\left|q_{\mathrm{h}}\right|-\left|q_{\mathrm{c}}\right|}{\left|q_{\mathrm{h}}\right|}=1-\frac{\left|q_{\mathrm{c}}\right|}{\left|q_{\mathrm{h}}\right|} \tag{3A.8}
\end{equation*}


It then follows from eqn 3A.6, written as $\left|q_{\mathrm{c}}\right| /\left|q_{\mathrm{h}}\right|=T_{\mathrm{c}} / T_{\mathrm{h}}$ that


\begin{equation*}
\eta=1-\frac{T_{\mathrm{c}}}{T_{\mathrm{h}}} \tag{3A.9}
\end{equation*}


Carnot efficiency\\
\includegraphics[max width=\textwidth, center]{2025_02_27_d4b95d9718f0b9e2e6b9g-116}

Figure 3A. 10 In a heat engine, an energy $q_{\mathrm{h}}$ (for example, $\left|q_{\mathrm{h}}\right|=$ 20 kJ ) is extracted as heat from the hot source and $q_{c}$ is discarded into the cold sink (for example, $\left|q_{\mathrm{c}}\right|=15 \mathrm{~kJ}$ ). The work done by the engine is equal to $\left|q_{\mathrm{h}}\right|-\left|q_{\mathrm{c}}\right|$ (e.g. $20 \mathrm{~kJ}-15 \mathrm{~kJ}=5 \mathrm{~kJ}$ ).

\subsection*{Brief illustration 3A. 3}
A certain power station operates with superheated steam at $300^{\circ} \mathrm{C}\left(T_{\mathrm{h}}=573 \mathrm{~K}\right)$ and discharges the waste heat into the environment at $20^{\circ} \mathrm{C}\left(T_{\mathrm{c}}=293 \mathrm{~K}\right)$. The theoretical efficiency is therefore

$$
\eta=1-\frac{293 \mathrm{~K}}{573 \mathrm{~K}}=0.489
$$

or 48.9 per cent. In practice, there are other losses due to mechanical friction and the fact that the turbines do not operate reversibly.

Now this conclusion can be generalized. The Second Law of thermodynamics implies that all reversible engines have the same efficiency regardless of their construction. To see the truth of this statement, suppose two reversible engines are coupled together and run between the same hot source and cold sink (Fig. 3A.11). The working substances and details of construction of the two engines are entirely arbitrary. Initially, suppose that engine $A$ is more efficient than engine $B$, and that a setting of the controls has been chosen that causes engine $B$ to acquire energy as heat $q_{c}$ from the cold sink and to release a certain quantity of energy as heat into the hot source. However, because engine $A$ is more efficient than engine $B$, not all the work that A produces is needed for this process and the difference can be used to do work. The net result is that the cold reservoir is unchanged, work has been done, and the hot reservoir has lost a certain amount of energy. This outcome is contrary to the Kelvin statement of the Second Law, because some heat has been converted directly into work. Because the conclusion is contrary to experience, the initial assumption that engines A and B can have different efficiencies must be false. It follows\\
\includegraphics[max width=\textwidth, center]{2025_02_27_d4b95d9718f0b9e2e6b9g-117}

Figure 3A. 11 (a) The demonstration of the equivalence of the efficiencies of all reversible engines working between the same thermal reservoirs is based on the flow of energy represented in this diagram. (b) The net effect of the processes is the conversion of heat into work without there being a need for a cold sink. This is contrary to the Kelvin statement of the Second Law.\\
that the relation between the heat transfers and the temperatures must also be independent of the working material, and therefore that eqn 3A. 9 is true for any substance involved in a Carnot cycle.

For the final step of the argument note that any reversible cycle can be approximated as a collection of Carnot cycles. This approximation is illustrated in Fig. 3A.12, which shows three Carnot cycles A, B, and C fitted together in such a way that their perimeter approximates the cycle indicated by the\\
\includegraphics[max width=\textwidth, center]{2025_02_27_d4b95d9718f0b9e2e6b9g-117(1)}

Figure 3A. 12 The path indicated by the purple line can be approximated by traversing the overall perimeter of the area created by the three Carnot cycles A, B, and C; for each individual cycle the overall entropy change is zero. The entropy changes along the adiabatic segments (such as $a_{1}-a_{4}$ and $c_{2}-c_{3}$ ) are zero, so it follows that the entropy changes along the isothermal segments of any one cycle (such as $a_{1}-a_{2}$ and $a_{3}-a_{4}$ ) cancel. The entropy change resulting from traversing the overall perimeter of the three cycles is therefore zero.\\
purple line. The entropy change around each individual cycle is zero (as already demonstrated), so the sum of entropy changes for all the cycles is zero. However, in the sum, the entropy change along any individual path is cancelled by the entropy change along the path it shares with the neighbouring cycle (because neighbouring paths are traversed in opposite directions). Therefore, all the entropy changes cancel except for those along the perimeter of the overall cycle and therefore the sum $q_{\text {rev }} / T$ around the perimeter is zero.

The path shown by the purple line can be approximated more closely by using more Carnot cycles, each of which is much smaller, and in the limit that they are infinitesimally small their perimeter matches the purple path exactly. Equation 3A. 5 (that the integral of $\mathrm{d} q_{\mathrm{rev}} / T$ round a general cycle is zero) then follows immediately. This result implies that $\mathrm{d} S$ is an exact differential and therefore that $S$ is a state function.

\subsection*{(b) The thermodynamic temperature}
Suppose an engine works reversibly between a hot source at a temperature $T_{\mathrm{h}}$ and a cold sink at a temperature $T$, then it follows from eqn 3A. 9 that


\begin{equation*}
T=(1-\eta) T_{\mathrm{h}} \tag{3A.10}
\end{equation*}


This expression enabled Kelvin to define the thermodynamic temperature scale in terms of the efficiency of a heat engine: construct an engine in which the hot source is at a known temperature and the cold sink is the object of interest. The temperature of the latter can then be inferred from the measured efficiency of the engine. The Kelvin scale (which is a special case of the thermodynamic temperature scale) is currently defined by using water at its triple point as the notional hot source and defining that temperature as 273.16 K exactly. ${ }^{4}$

\subsection*{(c) The Clausius inequality}
To show that the definition of entropy is consistent with the Second Law, note that more work is done when a change is reversible than when it is irreversible. That is, $\left|\mathrm{d} w_{\text {rev }}\right| \geq|\mathrm{d} w|$. Because $\mathrm{d} w$ and $\mathrm{d} w_{\text {rev }}$ are negative when energy leaves the system as work, this expression is the same as $-\mathrm{d} w_{\text {rev }} \geq-\mathrm{d} w$, and hence $\mathrm{d} w-\mathrm{d} w_{\text {rev }} \geq 0$. The internal energy is a state function, so its change is the same for irreversible and reversible paths between the same two states, and therefore

$$
\mathrm{d} U=\mathrm{d} q+\mathrm{d} w=\mathrm{d} q_{\mathrm{rev}}+\mathrm{d} w_{\mathrm{rev}}
$$

and hence $\mathrm{d} q_{\text {rev }}-\mathrm{d} q=\mathrm{d} w-\mathrm{d} w_{\text {rev }}$. Then, because $\mathrm{d} w-\mathrm{d} w_{\text {rev }} \geq 0$, it follows that $\mathrm{d} q_{\text {rev }}-\mathrm{d} q \geq 0$ and therefore $\mathrm{d} q_{\text {rev }} \geq \mathrm{d} q$. Division\\
${ }^{4}$ The international community has agreed to replace this definition by another that is independent of the specification of a particular substance, but the new definition has not yet (in 2018) been implemented.\\
by $T$ then results in $\mathrm{d} q_{\text {rev }} / T \geq \mathrm{d} q / T$. From the thermodynamic definition of the entropy $\left(\mathrm{d} S=\mathrm{d} q_{\mathrm{rev}} / T\right)$ it then follows that


\begin{equation*}
\mathrm{d} S \geq \frac{\mathrm{d} q}{T} \tag{3A.11}
\end{equation*}


Clausius inequality\\
This expression is the Clausius inequality. It proves to be of great importance for the discussion of the spontaneity of chemical reactions (Topic 3D).

Suppose a system is isolated from its surroundings, so that $\mathrm{d} q=0$. The Clausius inequality implies that


\begin{equation*}
\mathrm{d} S \geq 0 \tag{3A.12}
\end{equation*}


That is, in an isolated system the entropy cannot decrease when a spontaneous change occurs. This statement captures the content of the Second Law.

The Clausius inequality also implies that spontaneous processes are also necessarily irreversible processes. To confirm this conclusion, the inequality is introduced into the expression for the total entropy change that accompanies a process:

$$
\mathrm{d} S_{\text {tot }}=\overbrace{\mathrm{d} S}^{\geq \mathrm{d} q / T}+\overbrace{\mathrm{d} S_{\text {sur }}}^{-\mathrm{d} q / T} \geq 0
$$

where the inequality corresponds to an irreversible process and the equality to a reversible process. That is, a spontaneous process ( $\mathrm{d} S_{\text {tot }}>0$ ) is an irreversible process. A reversible process, for which $\mathrm{d} S_{\text {tot }}=0$, is spontaneous in neither direction: it is at equilibrium.\\
Apart from its fundamental importance in linking the definition of entropy to the Second Law, the Clausius inequality can also be used to show that a familiar process, the cooling of an object to the temperature of its surroundings, is indeed spontaneous. Consider the transfer of energy as heat from one system-the hot source-at a temperature $T_{\mathrm{h}}$ to another system—the cold sink-at a temperature $T_{c}$ (Fig. 3A.13). When\\
\includegraphics[max width=\textwidth, center]{2025_02_27_d4b95d9718f0b9e2e6b9g-118}

Figure 3A. 13 When energy leaves a hot source as heat, the entropy of the source decreases. When the same quantity of energy enters a cooler sink, the increase in entropy is greater. Hence, overall there is an increase in entropy and the process is spontaneous. Relative changes in entropy are indicated by the sizes of the arrows.\\
$|\mathrm{d} q|$ leaves the hot source (so $\mathrm{d} q_{\mathrm{h}}<0$ ), the Clausius inequality implies that $\mathrm{d} S \geq \mathrm{d} q_{\mathrm{h}} / T_{\mathrm{h}}$. When $|\mathrm{d} q|$ enters the cold sink the Clausius inequality implies that $\mathrm{d} S \geq \mathrm{d} q_{\mathrm{c}} / T_{\mathrm{c}}$ (with $\mathrm{d} q_{\mathrm{c}}>0$ ). Overall, therefore,

$$
\mathrm{d} S \geq \frac{\mathrm{d} q_{\mathrm{h}}}{T_{\mathrm{h}}}+\frac{\mathrm{d} q_{\mathrm{c}}}{T_{\mathrm{c}}}
$$

However, $\mathrm{d} q_{\mathrm{h}}=-\mathrm{d} q_{\mathrm{c}}$, so

$$
\mathrm{d} S \geq-\frac{\mathrm{d} q_{\mathrm{c}}}{T_{\mathrm{h}}}+\frac{\mathrm{d} q_{\mathrm{c}}}{T_{\mathrm{c}}}=\left(\frac{1}{T_{\mathrm{c}}}-\frac{1}{T_{\mathrm{h}}}\right) \mathrm{d} q_{\mathrm{c}}
$$

which is positive (because $\mathrm{d} q_{\mathrm{c}}>0$ and $T_{\mathrm{h}} \geq T_{\mathrm{c}}$ ). Hence, cooling (the transfer of heat from hot to cold) is spontaneous, in accord with experience.

\subsection*{Checklist of concepts}
\begin{enumerate}
  \item The entropy is a signpost of spontaneous change: the entropy of the universe increases in a spontaneous process.\\
$\square$ 2. A change in entropy is defined in terms of reversible heat transactions.
  \item The Boltzmann formula defines entropy in terms of the number of ways that the molecules can be arranged amongst the energy states, subject to the arrangements having the same overall energy.
  \item The Carnot cycle is used to prove that entropy is a state function.
  \item The efficiency of a heat engine is the basis of the definition of the thermodynamic temperature scale and one realization of such a scale, the Kelvin scale.
  \item The Clausius inequality is used to show that the entropy of an isolated system increases in a spontaneous change and therefore that the Clausius definition is consistent with the Second Law.\\
$\square$ 7. Spontaneous processes are irreversible processes; processes accompanied by no change in entropy are at equilibrium.
\end{enumerate}

\subsection*{Checklist of equations}
\begin{center}
\begin{tabular}{llll}
\hline
Property & Equation & Comment & Equation number \\
\hline
Thermodynamic entropy & $\mathrm{d} S=\mathrm{d} q_{\mathrm{rev}} / T$ & Definition & 3A.1a \\
Entropy change of surroundings & $\Delta S_{\text {sur }}=q_{\mathrm{sur}} / T_{\mathrm{sur}}$ &  & 3A.2b \\
Boltzmann formula & $S=k \ln \mathcal{W}$ & Definition & Reversible processes \\
Carnot efficiency & $\eta=1-T_{\mathrm{c}} / T_{\mathrm{h}}$ &  & 3 A .9 \\
Thermodynamic temperature & $T=(1-\eta) T_{\mathrm{h}}$ &  & 3A.10 \\
Clausius inequality & $\mathrm{d} S \geq \mathrm{d} q / T$ &  & 3A.11 \\
\hline
\end{tabular}
\end{center}

\section*{TOPIC 3B Entropy changes accompanying specific processes }
\subsection*{- Why do you need to know this material?}
The changes in entropy accompanying a variety of basic physical processes occur throughout the application of the Second Law to chemistry.

\subsection*{> What is the key idea?}
The change in entropy accompanying a process is calculated by identifying a reversible path between the initial and final states.

\subsection*{- What do you need to know already?}
You need to be familiar with the thermodynamic definition of entropy (Topic 3A), the First-Law concepts of work, heat, and internal energy (Topic 2A), and heat capacity (Topic 2B). The Topic makes use of the expressions for the work and heat transactions during the reversible, isothermal expansion of a perfect gas (Topic 2A).

The thermodynamic definition of entropy change given in eqn 3A.1,


\begin{equation*}
\mathrm{d} S=\frac{\mathrm{d} q_{\mathrm{rev}}}{T} \quad \Delta S=\int_{\mathrm{i}}^{\mathrm{f}} \frac{\mathrm{~d} q_{\mathrm{rev}}}{T} \tag{3B.1a}
\end{equation*}


Entropy change [definition]\\
where $q_{\text {rev }}$ is the energy supplied reversibly as heat to the system at a temperature $T$, is the basis of all calculations relating to entropy in thermodynamics. When applied to the surroundings, this definition implies eqn 3A.2b, which is repeated here as

\[
\Delta S_{\text {sur }}=\frac{q_{\text {sur }}}{T_{\text {sur }}} \quad \begin{array}{ll}
\text { Entropy change }  \tag{3B.1b}\\
\text { of surroundings }
\end{array}
\]

where $q_{\text {sur }}$ is the energy supplied as heat to the surroundings and $T_{\text {sur }}$ is their temperature; note that the entropy change of the surroundings is the same whether or not the process is reversible or irreversible for the system. The total change in entropy of an (overall) isolated system (the 'universe') is

$$
\Delta S_{\text {tot }}=\Delta S+\Delta S_{\text {sur }} \quad \text { Total entropy change }
$$

The entropy changes accompanying some physical changes are of particular importance and are treated here. As explained\\
in Topic 3A, a spontaneous process is also irreversible (in the thermodynamic sense) and a process for which $\Delta S_{\text {tot }}=0$ is at equilibrium.

\subsection*{3B. 1 Expansion}
In Topic 3A (specifically Example 3A.1) it is established that the change in entropy of a perfect gas when it expands isothermally from $V_{\mathrm{i}}$ to $V_{\mathrm{f}}$ is

\[
\Delta S=n R \ln \frac{V_{\mathrm{f}}}{V_{\mathrm{i}}} \quad \begin{align*}
& \text { Entropy change for the }  \tag{3B.2}\\
& \text { isothermal expansion of } \\
& \text { a perfect gas }
\end{align*}
\]

Because $S$ is a state function, the value of $\Delta S$ of the system is independent of the path between the initial and final states, so this expression applies whether the change of state occurs reversibly or irreversibly. The logarithmic dependence of entropy on volume is illustrated in Fig. 3B.1.

The total change in entropy, however, does depend on how the expansion takes place. For any process the energy lost as heat from the system is acquired by the surroundings, so $\mathrm{d} q_{\text {sur }}=-\mathrm{d} q$. For the reversible isothermal expansion of a perfect gas $q_{\text {rev }}=n R T \ln \left(V_{\mathrm{f}} / V_{\mathrm{i}}\right)$, so $q_{\text {sur }}=-n R T \ln \left(V_{\mathrm{f}} / V_{\mathrm{i}}\right)$, and consequently


\begin{equation*}
\Delta S_{\mathrm{sur}}=-\frac{q_{\mathrm{rev}}}{T}=-n R \ln \frac{V_{\mathrm{f}}}{V_{\mathrm{i}}} \tag{3B.3a}
\end{equation*}


\begin{center}
\includegraphics[max width=\textwidth]{2025_02_27_d4b95d9718f0b9e2e6b9g-120}
\end{center}

Figure 3B. 1 The logarithmic increase in entropy of a perfect gas as it expands isothermally.

This change is the negative of the change in the system, so $\Delta S_{\text {tot }}=0$, as expected for a reversible process. If, on the other hand, the isothermal expansion occurs freely (if the expansion is into a vacuum) no work is done ( $w=0$ ). Because the expansion is isothermal, $\Delta U=0$, and it follows from the First Law, $\Delta U=q+w$, that $q=0$. As a result, $q_{\text {sur }}=0$ and hence $\Delta S_{\text {sur }}=0$. For this expansion doing no work the total entropy change is therefore given by eqn 3B. 1 itself:


\begin{equation*}
\Delta S_{\mathrm{tot}}=n R \ln \frac{V_{\mathrm{f}}}{V_{\mathrm{i}}} \tag{3B.3b}
\end{equation*}


In this case, $\Delta S_{\text {tot }}>0$, as expected for an irreversible process.

\subsection*{Brief illustration 3B. 1}
When the volume of any perfect gas is doubled at constant temperature, $V_{\mathrm{f}} / V_{\mathrm{i}}=2$, and hence the change in molar entropy of the system is

$$
\Delta S_{\mathrm{m}}=\left(8.3145 \mathrm{~J} \mathrm{~K}^{-1} \mathrm{~mol}^{-1}\right) \times \ln 2=+5.76 \mathrm{~J} \mathrm{~K}^{-1} \mathrm{~mol}^{-1}
$$

If the change is carried out reversibly, the change in entropy of the surroundings is $-5.76 \mathrm{~J} \mathrm{~K}^{-1} \mathrm{~mol}^{-1}$ (the 'per mole' meaning per mole of gas molecules in the sample). The total change in entropy is 0 . If the expansion is free, the change in molar entropy of the gas is still $+5.76 \mathrm{~J} \mathrm{~K}^{-1} \mathrm{~mol}^{-1}$, but that of the surroundings is 0 , and the total change is $+5.76 \mathrm{~J} \mathrm{~K}^{-1} \mathrm{~mol}^{-1}$.

\subsection*{3B.2 Phase transitions}
When a substance freezes or boils the degree of dispersal of matter and the associated energy changes reflect the order with which the molecules pack together and the extent to which the energy is localized. Therefore, a transition is expected to be accompanied by a change in entropy. For example, when a substance vaporizes, a compact condensed phase changes into a widely dispersed gas, and the entropy of the substance can be expected to increase considerably. The entropy of a solid also increases when it melts to a liquid.

Consider a system and its surroundings at the normal transition temperature, $T_{\text {trs }}$, the temperature at which two phases are in equilibrium at 1 atm . This temperature is $0^{\circ} \mathrm{C}(273 \mathrm{~K})$ for ice in equilibrium with liquid water at 1 atm , and $100^{\circ} \mathrm{C}$ ( 373 K ) for water in equilibrium with its vapour at 1 atm . At the transition temperature, any transfer of energy as heat between the system and its surroundings is reversible because the two phases in the system are in equilibrium. Because at constant pressure $q=\Delta_{\text {trs }} H$, the change in molar entropy of the system is ${ }^{1}$

\footnotetext{${ }^{1}$ According to Topic 2C, $\Delta_{\text {trs }} H$ is an enthalpy change per mole of substance, so $\Delta_{\text {trs }} S$ is also a molar quantity.
}
\$\$

\begin{equation*}
\Delta_{\mathrm{trs}} S=\frac{\Delta_{\mathrm{trs}} H}{T_{\mathrm{trs}}} \tag{3B.4}
\end{equation*}

\$\$

$$
\begin{aligned}
& \text { Entropy of phase transition } \\
& \text { [at } T_{\text {trs }} \text { ] }
\end{aligned}
$$

If the phase transition is exothermic ( $\Delta_{\text {trs }} H<0$, as in freezing or condensing), then the entropy change of the system is negative. This decrease in entropy is consistent with the increased order of a solid compared with a liquid, and with the increased order of a liquid compared with a gas. The change in entropy of the surroundings, however, is positive because energy is released as heat into them. At the transition temperature the total change in entropy is zero because the two phases are in equilibrium. If the transition is endothermic ( $\Delta_{\text {trs }} H>0$, as in melting and vaporization), then the entropy change of the system is positive, which is consistent with dispersal of matter in the system. The entropy of the surroundings decreases by the same amount, and overall the total change in entropy is zero.

Table 3B. 1 lists some experimental entropies of phase transitions. Table 3B. 2 lists in more detail the standard entropies of vaporization of several liquids at their normal boiling points. An interesting feature of the data is that a wide range of liquids give approximately the same standard entropy of vaporization (about $85 \mathrm{~J} \mathrm{~K}^{-1} \mathrm{~mol}^{-1}$ ): this empirical observation is called Trouton's rule. The explanation of Trouton's rule is that a similar change in volume occurs when any liquid evaporates and becomes a gas. Hence, all liquids can be expected to have similar standard entropies of vaporization.

Liquids that show significant deviations from Trouton's rule do so on account of strong molecular interactions that result

Table 3B. 1 Standard entropies of phase transitions, $\Delta_{\text {trs }} S^{\ominus} /\left(\mathrm{JK}^{-1}\right.$ $\left.\mathrm{mol}^{-1}\right)$, at the corresponding normal transition temperatures*

\begin{center}
\begin{tabular}{llc}
\hline
 & Fusion (at $T_{\mathrm{f}}$ ) & Vaporization (at $T_{\mathrm{b}}$ ) \\
\hline
Argon, Ar & $14.17($ at 83.8 K$)$ & 74.53 (at 87.3 K ) \\
Benzene, $\mathrm{C}_{6} \mathrm{H}_{6}$ & $38.00($ at 279 K$)$ & 87.19 (at 353 K ) \\
Water, $\mathrm{H}_{2} \mathrm{O}$ & 22.00 (at 273.15 K ) & 109.0 (at 373.15 K ) \\
Helium, He & 4.8 (at 8 K and 30 bar) & 19.9 (at 4.22 K ) \\
\hline
\end{tabular}
\end{center}

\begin{itemize}
  \item More values are given in the Resource section.
\end{itemize}

Table 3B. 2 The standard enthalpies and entropies of vaporization of liquids at their boiling temperatures*

\begin{center}
\begin{tabular}{llrl}
\hline
 & $\Delta_{\text {vap }} H^{\ominus} /\left(\mathbf{k J ~ m o l}^{-1}\right)$ & $\theta_{\mathbf{b}} /{ }^{\circ} \mathrm{C}$ & \begin{tabular}{l}
$\Delta_{\text {vap }} S^{\ominus} /$ \\
$\left(\mathrm{JK}^{-1} \mathrm{~mol}^{-1}\right)$ \\
\end{tabular} \\
\hline
Benzene & 30.8 & 80.1 & 87.2 \\
Carbon tetrachloride & 30 & 76.7 & 85.8 \\
Cyclohexane & 30.1 & 80.7 & 85.1 \\
Hydrogen sulfide & 18.7 & -60.4 & 87.9 \\
Methane & 8.18 & -161.5 & 73.2 \\
Water & 40.7 & 100.0 & 109.1 \\
\hline
\end{tabular}
\end{center}

\begin{itemize}
  \item More values are given in the Resource section.\\
in a partial ordering of their molecules. As a result, there is a greater change in disorder when the liquid turns into a vapour than for when a fully disordered liquid vaporizes. An example is water, where the large entropy of vaporization reflects the presence of structure arising from hydrogen bonding in the liquid. Hydrogen bonds tend to organize the molecules in the liquid so that they are less random than, for example, the molecules in liquid hydrogen sulfide (in which there is no hydrogen bonding). Methane has an unusually low entropy of vaporization. A part of the reason is that the entropy of the gas itself is slightly low ( $186 \mathrm{~J} \mathrm{~K}^{-1} \mathrm{~mol}^{-1}$ at 298 K ; the entropy of $\mathrm{N}_{2}$ under the same conditions is $192 \mathrm{~J} \mathrm{~K}^{-1} \mathrm{~mol}^{-1}$ ). As explained in Topic 13B, fewer translational and rotational states are accessible at room temperature for molecules with low mass and moments of inertia (like $\mathrm{CH}_{4}$ ) than for molecules with relatively high mass and moments of inertia (like $\mathrm{N}_{2}$ ), so their molar entropy is slightly lower.
\end{itemize}

\subsection*{Brief illustration 3B. 2}
There is no hydrogen bonding in liquid bromine and $\mathrm{Br}_{2}$ is a heavy molecule which is unlikely to display unusual behaviour in the gas phase, so it is safe to use Trouton's rule. To predict the standard molar enthalpy of vaporization of bromine given that it boils at $59.2^{\circ} \mathrm{C}$, use Trouton's rule in the form

$$
\Delta_{\text {vap }} H^{\ominus}=T_{\mathrm{b}} \times\left(85 \mathrm{~J} \mathrm{~K}^{-1} \mathrm{~mol}^{-1}\right)
$$

Substitution of the data then gives

$$
\begin{aligned}
\Delta_{\text {vap }} H^{\ominus} & =(332.4 \mathrm{~K}) \times\left(85 \mathrm{~J} \mathrm{~K}^{-1} \mathrm{~mol}^{-1}\right) \\
& =+2.8 \times 10^{4} \mathrm{~J} \mathrm{~mol}^{-1}=+28 \mathrm{~kJ} \mathrm{~mol}^{-1}
\end{aligned}
$$

The experimental value is $+29.45 \mathrm{~kJ} \mathrm{~mol}^{-1}$.

\subsection*{зв. 3 Heating}
The thermodynamic definition of entropy change in eqn 3B.1a is used to calculate the entropy of a system at a temperature $T_{\mathrm{f}}$ from a knowledge of its entropy at another temperature $T_{\mathrm{i}}$ and the heat supplied to change its temperature from one value to the other:


\begin{equation*}
S\left(T_{\mathrm{f}}\right)=S\left(T_{\mathrm{i}}\right)+\int_{T_{\mathrm{i}}}^{T_{\mathrm{f}}} \frac{\mathrm{~d} q_{\mathrm{rev}}}{T} \tag{3B.5}
\end{equation*}


The most common version of this expression is for a system subjected to constant pressure (such as from the atmosphere) during the heating, so then $\mathrm{d} q_{\mathrm{rev}}=\mathrm{d} H$. From the definition of constant-pressure heat capacity (eqn 2B.5, $C_{p}=(\partial H / \partial T)_{p}$ ) it follows that $\mathrm{d} H=C_{p} \mathrm{~d} T$, and hence $\mathrm{d} q_{\mathrm{rev}}=C_{p} \mathrm{~d} T$. Substitution into eqn 3B. 5 gives

\[
S\left(T_{\mathrm{f}}\right)=S\left(T_{\mathrm{i}}\right)+\int_{T_{\mathrm{i}}}^{T_{\mathrm{i}}} \frac{C_{p} \mathrm{~d} T}{T} \quad \begin{align*}
& \text { Entropy variation with }  \tag{3B.6}\\
& \text { temperature } \\
& \text { [constant } p \text { ] }
\end{align*}
\]

\begin{center}
\includegraphics[max width=\textwidth]{2025_02_27_d4b95d9718f0b9e2e6b9g-122}
\end{center}

Figure 3B. 2 The logarithmic increase in entropy of a substance as it is heated at either constant volume or constant pressure. Different curves are labelled with the corresponding value of $C_{m} / R$, taken to be constant over the temperature range. For constant volume conditions $C_{m}=C_{V, m}$ and at constant pressure $C_{m}=C_{p, m}$.

The same expression applies at constant volume, but with $C_{p}$ replaced by $C_{V}$. When $C_{p}$ is independent of temperature over the temperature range of interest, it can be taken outside the integral to give


\begin{equation*}
S\left(T_{\mathrm{f}}\right)=S\left(T_{\mathrm{i}}\right)+C_{p} \int_{T_{\mathrm{i}}}^{T_{\mathrm{f}}} \frac{\mathrm{~d} T}{T}=S\left(T_{\mathrm{i}}\right)+C_{p} \ln \frac{T_{\mathrm{f}}}{T_{\mathrm{i}}} \tag{3B.7}
\end{equation*}


with a similar expression for heating at constant volume. The logarithmic dependence of entropy on temperature is illustrated in Fig. 3B.2.

\subsection*{Brief illustration 3B. 3}
The molar constant-volume heat capacity of water at 298 K is $75.3 \mathrm{~J} \mathrm{~K}^{-1} \mathrm{~mol}^{-1}$. The change in molar entropy when it is heated from $20^{\circ} \mathrm{C}(293 \mathrm{~K})$ to $50^{\circ} \mathrm{C}(323 \mathrm{~K})$, supposing the heat capacity to be constant in that range, is therefore

$$
\begin{aligned}
\Delta S_{\mathrm{m}} & =S_{\mathrm{m}}(323 \mathrm{~K})-S_{\mathrm{m}}(293 \mathrm{~K})=\left(75.3 \mathrm{JK}^{-1} \mathrm{~mol}^{-1}\right) \times \ln \frac{323 \mathrm{~K}}{293 \mathrm{~K}} \\
& =+7.34 \mathrm{JK}^{-1} \mathrm{~mol}^{-1}
\end{aligned}
$$

\subsection*{3B. 4 Composite processes}
In many processes, more than one parameter changes. For instance, it might be the case that both the volume and the temperature of a gas are different in the initial and final states. Because $S$ is a state function, the change in its value can be calculated by considering any reversible path between the initial and final states. For example, it might be convenient to split the path into two steps: an isothermal expansion to the final\\
volume, followed by heating at constant volume to the final temperature. Then the total entropy change when both variables change is the sum of the two contributions.

Example 3B. 1 Calculating the entropy change for a composite process

Calculate the entropy change when argon at $25^{\circ} \mathrm{C}$ and 1.00 bar in a container of volume $0.500 \mathrm{dm}^{3}$ is allowed to expand to $1.000 \mathrm{dm}^{3}$ and is simultaneously heated to $100^{\circ} \mathrm{C}$. (Take the molar heat capacity at constant volume to be $\frac{3}{2} R$.)\\
Collect your thoughts As remarked in the text, you can break the overall process down into two steps: isothermal expansion to the final volume, followed by heating at constant volume to the final temperature. The entropy change in the first step is given by eqn 3B. 2 and that of the second step, provided $C_{V}$ is independent of temperature, by eqn 3B. 7 (with $C_{V}$ in place of $C_{p}$ ). In each case you need to know $n$, the amount of gas molecules, which can be calculated from the perfect gas equation and the data for the initial state by using $n=p_{\mathrm{i}} V_{\mathrm{i}} / R T_{\mathrm{i}}$.

The solution The amount of gas molecules is

$$
n=\frac{\left(1.00 \times 10^{5} \mathrm{~Pa}\right) \times\left(0.500 \times 10^{-3} \mathrm{~m}^{3}\right)}{\left(8.3145 \mathrm{JK}^{-1} \mathrm{~mol}^{-1}\right) \times 298 \mathrm{~K}}=0.0201 \ldots \mathrm{~mol}
$$

From eqn 3B. 2 the entropy change in the isothermal expansion from $V_{\mathrm{i}}$ to $V_{\mathrm{f}}$ is

$$
\begin{aligned}
\Delta S(\text { Step } 1) & =n R \ln \frac{V_{\mathrm{f}}}{V_{\mathrm{i}}} \\
& =0.0201 \ldots \mathrm{~mol} \times\left(8.3145 \mathrm{JK}^{-1} \mathrm{~mol}^{-1}\right) \ln \frac{1.000 \mathrm{dm}^{3}}{0.500 \mathrm{dm}^{3}} \\
& =+0.116 \ldots \mathrm{JK}^{-1}
\end{aligned}
$$

From eqn 3B.6, the entropy change in the second step, heating from $T_{\mathrm{i}}$ to $T_{\mathrm{f}}$ at constant volume, is

$$
\begin{aligned}
\Delta S(\text { Step } 2) & =n C_{V, m} \ln \frac{T_{\mathrm{f}}}{T_{\mathrm{i}}}=\frac{3}{2} n R \ln \frac{T_{\mathrm{f}}}{T_{\mathrm{i}}} \\
& =\frac{3}{2} \times(0.0201 \ldots \mathrm{~mol}) \times\left(8.3145 \mathrm{JK}^{-1} \mathrm{~mol}^{-1}\right) \ln \frac{373 \mathrm{~K}}{298 \mathrm{~K}} \\
& =+0.0564 \ldots \mathrm{JK}^{-1}
\end{aligned}
$$

The overall entropy change of the system, the sum of these two changes, is

$$
\Delta S=0.116 \ldots \mathrm{~J} \mathrm{~K}^{-1}+0.0564 \ldots \mathrm{~J} \mathrm{~K}^{-1}=+0.173 \mathrm{~J} \mathrm{~K}^{-1}
$$

Self-test 3B.1 Calculate the entropy change when the same initial sample is compressed to $0.0500 \mathrm{dm}^{3}$ and cooled to $-25^{\circ} \mathrm{C}$.\\
\includegraphics[max width=\textwidth, center]{2025_02_27_d4b95d9718f0b9e2e6b9g-123}

\subsection*{Checklist of concepts}
$\square$ 1. The entropy of a perfect gas increases when it expands isothermally.2. The change in entropy of a substance accompanying a change of state at its transition temperature is calculated from its enthalpy of transition.\\
3. The increase in entropy when a substance is heated is calculated from its heat capacity.

\subsection*{Checklist of equations}
\begin{center}
\begin{tabular}{|c|c|c|c|}
\hline
Property & Equation & Comment & Equation number \\
\hline
Entropy of isothermal expansion & $\Delta S=n R \ln \left(V_{\mathrm{f}} / V_{\mathrm{i}}\right)$ & Perfect gas & 3B. 2 \\
\hline
Entropy of transition & $\Delta_{\text {trs }} S=\Delta_{\text {trs }} H / T_{\text {trs }}$ & At the transition temperature & 3B. 4 \\
\hline
Variation of entropy with temperature & $S\left(T_{\mathrm{f}}\right)=S\left(T_{\mathrm{i}}\right)+C \ln \left(T_{\mathrm{f}} / T_{\mathrm{i}}\right)$ & The heat capacity, $C$, is independent of temperature and no phase transitions occur; $C=C_{p}$ for constant pressure and $C_{V}$ for constant volume. & 3B. 7 \\
\hline
\end{tabular}
\end{center}

\section*{TOPIC 3C The measurement of entropy}
\subsection*{- Why do you need to know this material?}
For entropy to be a quantitatively useful concept it is important to be able to measure it: the calorimetric procedure is described here. The Third Law of thermodynamics is used to report the measured values.

\subsection*{> What is the key idea?}
The entropy of a perfectly crystalline solid is zero at $T=0$.

\subsection*{> What do you need to know already?}
You need to be familiar with the expression for the temperature dependence of entropy and how entropies of phase changes are calculated (Topic 3B). The discussion of residual entropy draws on the Boltzmann formula for the entropy (Topic 3A).

The entropy of a substance can be determined in two ways. One, which is the subject of this Topic, is to make calorimetric measurements of the heat required to raise the temperature of a sample from $T=0$ to the temperature of interest. There are then two equations to use. One is the dependence of entropy on temperature, which is eqn 3B. 7 reproduced here as

$$
S\left(T_{2}\right)=S\left(T_{1}\right)+\int_{T_{1}}^{T_{2}} \frac{C_{p}(T)}{T} \mathrm{~d} T \quad \text { Entropy and temperature (3C.1a) }
$$

The second is the contribution of a phase change to the entropy, which according to eqn 3B. 4 is


\begin{equation*}
\Delta S\left(T_{\mathrm{trs}}\right)=\frac{\Delta_{\mathrm{trs}} H\left(T_{\text {trs }}\right)}{T_{\text {trs }}} \quad \text { Entropy of phase transition } \tag{3C.1b}
\end{equation*}


where $\Delta_{\text {trs }} H\left(T_{\text {trs }}\right)$ is the enthalpy of transition at the transition temperature $T_{\text {trs }}$. The other method, which is described in Topic 13E, is to use calculated parameters or spectroscopic data to calculate the entropy by using Boltzmann's statistical definition.

\subsection*{3C. 1 The calorimetric measurement of entropy}
According to eqn 3C.1a, the entropy of a system at a temperature $T$ is related to its entropy at $T=0$ by measuring its heat\\
capacity $C_{p}$ at different temperatures and evaluating the integral. The entropy of transition for each phase transition that occurs between $T=0$ and the temperature of interest must then be included in the overall sum. For example, if a substance melts at $T_{\mathrm{f}}$ and boils at $T_{\mathrm{b}}$, then its molar entropy at a particular temperature $T$ above its boiling temperature is given by


\begin{align*}
& S_{\mathrm{m}}(T)=S_{\mathrm{m}}(0)+\overbrace{\int_{0}^{T_{\mathrm{f}} C_{p, \mathrm{~m}}\left(\mathrm{~s}, T^{\prime}\right)} T^{\prime} \mathrm{d} T^{\prime}}^{\begin{array}{c}
\text { Heat solid } \\
\text { to its } \\
\text { melting point }
\end{array}}+\overbrace{\frac{\Delta_{\text {fus }} H}{T_{\mathrm{f}}}}^{\text {Entropy of }} \text { fusion } \\
& \text { Heat liquid } \\
& \text { to its Entropy of } \\
& \text { boiling point vaporization } \\
& +\overbrace{\int_{T_{\mathrm{f}}}^{T_{\mathrm{b}}} \frac{C_{p, \mathrm{~m}}\left(1, T^{\prime}\right)}{T^{\prime}} \mathrm{d} T^{\prime}}+\frac{\overbrace{\Delta_{\mathrm{vap}} H}^{T_{\mathrm{b}}}}{}  \tag{3C.2}\\
& \text { Heat vapour } \\
& \text { to the } \\
& +\overbrace{\int_{T_{\mathrm{b}}}^{T} \frac{C_{p, \mathrm{~m}}\left(\mathrm{~g}, T^{\prime}\right)}{T^{\prime}} \mathrm{d} T^{\prime}}^{\text {final temperature }}
\end{align*}


The variable of integration has been changed to $T^{\prime}$ to avoid confusion with the temperature of interest, $T$. All the properties required, except $S_{\mathrm{m}}(0)$, can be measured calorimetrically, and the integrals can be evaluated either graphically or, as is now more usual, by fitting a polynomial to the data and integrating the polynomial analytically. The former procedure is illustrated in Fig. 3C.1: the area under the curve of $C_{p, \mathrm{~m}}(T) / T$ against $T$ is the integral required. Provided all measurements are made at 1 bar on a pure material, the final value is the standard entropy, $S^{\ominus}(T)$; division by the amount of substance, $n$, gives the standard molar entropy, $S_{\mathrm{m}}^{\ominus}(T)=S^{\ominus}(T) / n$. Because $\mathrm{d} T / T=\mathrm{d} \ln T$, an alternative procedure is to evaluate the area under a plot of $C_{p, \mathrm{~m}}(T)$ against $\ln T$.

\subsection*{Brief illustration 3C. 1}
The standard molar entropy of nitrogen gas at $25^{\circ} \mathrm{C}$ has been calculated from the following data:

\begin{center}
\begin{tabular}{lc}
\hline
 & Contribution to $S_{\mathrm{m}}^{\ominus} /\left(\mathrm{J} \mathrm{K}^{-1} \mathrm{~mol}^{-1}\right)$ \\
\hline
Debye extrapolation & 1.92 \\
Integration, from 10 K to 35.61 K & 25.25 \\
Phase transition at 35.61 K & 6.43 \\
Integration, from 35.61 K to 63.14 K & 23.38 \\
Fusion at 63.14 K & 11.42 \\
Integration, from 63.14 K to 77.32 K & 11.41 \\
Vaporization at 77.32 K & 72.13 \\
Integration, from 77.32 K to 298.15 K & 39.20 \\
Correction for gas imperfection & 0.92 \\
Total & 192.06 \\
\hline
\end{tabular}
\end{center}

Therefore, $S_{\mathrm{m}}^{\ominus}(298.15 \mathrm{~K})=S_{\mathrm{m}}(0)+192.1 \mathrm{~J} \mathrm{~K}^{-1} \mathrm{~mol}^{-1}$. The Debye extrapolation is explained in the next paragraph.

One problem with the determination of entropy is the difficulty of measuring heat capacities near $T=0$. There are good theoretical grounds for assuming that the heat capacity of a non-metallic solid is proportional to $T^{3}$ when $T$ is low (see Topic 7A), and this dependence is the basis of the Debye extrapolation (or the Debye $T^{3}$ law). In this method, $C_{p}$ is measured down to as low a temperature as possible and a curve of the form $a T^{3}$ is fitted to the data. The fit determines the value of $a$, and the expression $C_{p, \mathrm{~m}}(T)=a T^{3}$ is then assumed to be valid down to $T=0$.\\
\includegraphics[max width=\textwidth, center]{2025_02_27_d4b95d9718f0b9e2e6b9g-125}

Figure 3C. 1 The variation of $C_{p} / T$ with the temperature for a sample is used to evaluate the entropy, which is equal to the area beneath the upper curve up to the corresponding temperature, plus the entropy of each phase transition encountered between $T=0$ and the temperature of interest. For instance, the entropy denoted by the yellow dot on the lower curve is given by the dark shaded area in the upper graph.

\subsection*{Example 3C. 1 Calculating the entropy at low \\
 temperatures}
The molar constant-pressure heat capacity of a certain nonmetallic solid at 4.2 K is $0.43 \mathrm{~J} \mathrm{~K}^{-1} \mathrm{~mol}^{-1}$. What is its molar entropy at that temperature?

Collect your thoughts Because the temperature is so low, you can assume that the heat capacity varies with temperature according to $C_{p, \mathrm{~m}}(T)=a T^{3}$, in which case you can use eqn 3C.1a to calculate the entropy at a temperature $T$ in terms of the entropy at $T=0$ and the constant $a$. When the integration is carried out, it turns out that the result can be expressed in terms of the heat capacity at the temperature $T$, so the data can be used directly to calculate the entropy.

The solution The integration required is

$$
\begin{aligned}
S_{\mathrm{m}}(T) & =S_{\mathrm{m}}(0)+\int_{0}^{T} \frac{a T^{\prime 3}}{T^{\prime}} \mathrm{d} T^{\prime}=S_{\mathrm{m}}(0)+a \overbrace{\int_{0}^{T} T^{\prime 2} \mathrm{~d} T^{\prime}}^{\text {Integral A. } 1} \\
& =S_{\mathrm{m}}(0)+\frac{1}{3} a T^{3}=S_{\mathrm{m}}(0)+\frac{1}{3} C_{p, \mathrm{~m}}(T)
\end{aligned}
$$

from which it follows that

$$
S_{\mathrm{m}}(4.2 \mathrm{~K})=S_{\mathrm{m}}(0)+0.14 \mathrm{~J} \mathrm{~K}^{-1} \mathrm{~mol}^{-1}
$$

Self-test 3C. 1 For metals, there is also a contribution to the heat capacity from the electrons which is linearly proportional to $T$ when the temperature is low; that is, $C_{p, \mathrm{~m}}(T)=b T$. Evaluate its contribution to the entropy at low temperatures.

\subsection*{3c. 2 The Third Law}
At $T=0$, all energy of thermal motion has been quenched, and in a perfect crystal all the atoms or ions are in a regular, uniform array. The localization of matter and the absence of thermal motion suggest that such materials also have zero entropy. This conclusion is consistent with the molecular interpretation of entropy (Topic 3A) because there is only one way of arranging the molecules when they are all in the ground state, which is the case at $T=0$. Thus, at $T=0, \mathcal{W}=1$ and from $S=$ $k \ln \mathcal{W}$ it follows that $S=0$.

\subsection*{(a) The Nernst heat theorem}
The Nernst heat theorem summarizes a series of experimental observations that turn out to be consistent with the view that the entropy of a regular array of molecules is zero at $T=0$ :

The entropy change accompanying any physical or chemical transformation approaches zero as the temperature approaches zero: $\Delta S \rightarrow 0$ as $T \rightarrow 0$ provided all the substances involved are perfectly ordered.

\subsection*{Brief illustration 3C. 2}
The entropy of the transition between orthorhombic sulfur, $\alpha$, and monoclinic sulfur, $\beta$, can be calculated from the transition enthalpy ( $402 \mathrm{~J} \mathrm{~mol}^{-1}$ ) at the transition temperature ( 369 K ):

$$
\begin{aligned}
\Delta_{\text {trs }} S(369 \mathrm{~K}) & =S_{\mathrm{m}}(\beta, 369 \mathrm{~K})-S_{\mathrm{m}}(\alpha, 369 \mathrm{~K}) \\
& =\frac{\Delta_{\text {trs }} H}{T_{\text {trs }}}=\frac{402 \mathrm{Jmol}^{-1}}{369 \mathrm{~K}}=1.09 \mathrm{JK}^{-1} \mathrm{~mol}^{-1}
\end{aligned}
$$

The entropies of the $\alpha$ and $\beta$ allotropes can also be determined by measuring their heat capacities from $T=0$ up to $T=369 \mathrm{~K}$. It is found that $S_{\mathrm{m}}(\alpha, 369 \mathrm{~K})=S_{\mathrm{m}}(\alpha, 0)+37 \mathrm{~J} \mathrm{~K}^{-1} \mathrm{~mol}^{-1}$ and $S_{\mathrm{m}}(\beta, 369 \mathrm{~K})=S_{\mathrm{m}}(\beta, 0)+38 \mathrm{~J} \mathrm{~K}^{-1} \mathrm{~mol}^{-1}$. These two values imply that at the transition temperature

$$
\begin{aligned}
\Delta_{\mathrm{trs}} S(369 \mathrm{~K})= & \left\{S_{\mathrm{m}}(\beta, 0)+38 \mathrm{~J} \mathrm{~K}^{-1} \mathrm{~mol}^{-1}\right\}- \\
& \left\{S_{\mathrm{m}}(\alpha, 0)+37 \mathrm{~J} \mathrm{~K}^{-1} \mathrm{~mol}^{-1}\right\} \\
= & S_{\mathrm{m}}(\beta, 0)-S_{\mathrm{m}}(\alpha, 0)+1 \mathrm{~J} \mathrm{~K}^{-1} \mathrm{~mol}^{-1}
\end{aligned}
$$

On comparing this value with the one above, it follows that $S_{\mathrm{m}}(\beta, 0)-S_{\mathrm{m}}(\alpha, 0) \approx 0$, in accord with the theorem.

It follows from the Nernst theorem that, if the value zero is ascribed to the entropies of elements in their perfect crystalline form at $T=0$, then all perfect crystalline compounds also have zero entropy at $T=0$ (because the change in entropy that accompanies the formation of the compounds, like the entropy of all transformations at that temperature, is zero). This conclusion is summarized by the Third Law of thermodynamics:

The entropy of all perfect crystalline\\
Third Law of substances is zero at $T=0$. thermodynamics

As far as thermodynamics is concerned, choosing this common value as zero is a matter of convenience. As noted above, the molecular interpretation of entropy justifies the value $S=0$ at $T=0$ because at this temperature $\mathcal{W}=1$.

In certain cases $\mathcal{W}>1$ at $T=0$ and therefore $S(0)>0$. This is the case if there is no energy advantage in adopting a particular orientation even at absolute zero. For instance, for a diatomic molecule $A B$ there may be almost no energy difference between the arrangements ... $\mathrm{AB} A B \mathrm{AB} \ldots$ and ... $\mathrm{BA} A B$ BA... in a solid, so $\mathcal{W}>1$ even at $T=0$. If $S(0)>0$ the substance is said to have a residual entropy. Ice has a residual entropy of $3.4 \mathrm{~J} \mathrm{~K}^{-1} \mathrm{~mol}^{-1}$. It stems from the arrangement of the hydrogen bonds between neighbouring water molecules: a given O\\
atom has two short $\mathrm{O}-\mathrm{H}$ bonds and two long $\mathrm{O} \cdots \mathrm{H}$ bonds to its neighbours, but there is a degree of randomness in which two bonds are short and which two are long.

\subsection*{(b) Third-Law entropies}
Entropies reported on the basis that $S(0)=0$ are called ThirdLaw entropies (and commonly just 'entropies'). When the substance is in its standard state at the temperature $T$, the standard (Third-Law) entropy is denoted $S^{\ominus}(T)$. A list of values at 298 K is given in Table 3C.1.

The standard reaction entropy, $\Delta_{\mathrm{r}} S^{\ominus}$, is defined, like the standard reaction enthalpy in Topic 2C, as the difference between the molar entropies of the pure, separated products and the pure, separated reactants, all substances being in their standard states at the specified temperature:

\[
\Delta_{\mathrm{r}} S^{\ominus}=\sum_{\text {Products }} v S_{\mathrm{m}}^{\ominus}-\sum_{\text {Reactants }} v S_{\mathrm{m}}^{\ominus} \quad \begin{align*}
& \begin{array}{l}
\text { Standard reaction } \\
\text { entropy } \\
\text { [definition] }
\end{array} \tag{3C.3a}
\end{align*}
\]

In this expression, each term is weighted by the appropriate stoichiometric coefficient. A more sophisticated approach is to adopt the notation introduced in Topic 2C and to write


\begin{equation*}
\Delta_{\mathrm{r}} S^{\ominus}=\sum_{\mathrm{J}} v_{\mathrm{J}} S_{\mathrm{m}}^{\ominus}(\mathrm{J}) \tag{3C.3b}
\end{equation*}


where the $v_{\mathrm{J}}$ are signed (+ for products, - for reactants) stoichiometric numbers. Standard reaction entropies are likely to be positive if there is a net formation of gas in a reaction, and are likely to be negative if there is a net consumption of gas.

Table 3C. 1 Standard Third-Law entropies at $298 \mathrm{~K}^{*}$

\begin{center}
\begin{tabular}{lc}
\hline
 & $S_{\mathrm{m}}^{\ominus} /\left(\mathrm{J} \mathrm{K}^{-1} \mathrm{~mol}^{-1}\right)$ \\
\hline
Solids &  \\
Graphite, $\mathrm{C}(\mathrm{s})$ & 5.7 \\
Diamond, $\mathrm{C}(\mathrm{s})$ & 2.4 \\
Sucrose, $\mathrm{C}_{12} \mathrm{H}_{22} \mathrm{O}_{11}(\mathrm{~s})$ & 360.2 \\
Iodine, $\mathrm{I}_{2}(\mathrm{~s})$ & 116.1 \\
Liquids &  \\
Benzene, $\mathrm{C}_{6} \mathrm{H}_{6}(\mathrm{l})$ & 173.3 \\
Water, $\mathrm{H}_{2} \mathrm{O}(\mathrm{l})$ & 69.9 \\
Mercury, $\mathrm{Hg}^{2}(\mathrm{l})$ & 76.0 \\
Gases &  \\
Methane, $\mathrm{CH}_{4}(\mathrm{~g})$ & 186.3 \\
Carbon dioxide, $\mathrm{CO}_{2}(\mathrm{~g})$ & 213.7 \\
Hydrogen, $\mathrm{H}_{2}(\mathrm{~g})$ & 130.7 \\
Helium, $\mathrm{He}^{(\mathrm{g})}$ & 126.2 \\
Ammonia, $\mathrm{NH}_{3}(\mathrm{~g})$ & 192.4 \\
\hline
\end{tabular}
\end{center}

\footnotetext{\begin{itemize}
  \item More values are given in the Resource section.
\end{itemize}
}\subsection*{Brief illustration 3C. 3}
To calculate the standard reaction entropy of $\mathrm{H}_{2}(\mathrm{~g})+\frac{1}{2} \mathrm{O}_{2}(\mathrm{~g})$ $\rightarrow \mathrm{H}_{2} \mathrm{O}(\mathrm{l})$ at 298 K , use the data in Table 2C. 4 of the Resource section to write

$$
\begin{aligned}
\Delta_{\mathrm{r}} \mathrm{~S}^{\ominus} & =S_{\mathrm{m}}^{\ominus}\left(\mathrm{H}_{2} \mathrm{O}, \mathrm{l}\right)-\left\{S_{\mathrm{m}}^{\ominus}\left(\mathrm{H}_{2}, \mathrm{~g}\right)+\frac{1}{2} S_{\mathrm{m}}^{\ominus}\left(\mathrm{O}_{2}, \mathrm{~g}\right)\right\} \\
& =69.9 \mathrm{~J} \mathrm{~K}^{-1} \mathrm{~mol}^{-1}-\left\{130.7+\frac{1}{2}(205.1)\right\} \mathrm{JK}^{-1} \mathrm{~mol}^{-1} \\
& =-163.4 \mathrm{JK} \mathrm{~K}^{-1} \mathrm{~mol}^{-1}
\end{aligned}
$$

The negative value is consistent with the conversion of two gases to a compact liquid.

A note on good practice Do not make the mistake of setting the standard molar entropies of elements equal to zero: they have non-zero values (provided $T>0$ ).

Just as in the discussion of enthalpies in Topic 2C, where it is acknowledged that solutions of cations cannot be prepared in the absence of anions, the standard molar entropies of ions in solution are reported on a scale in which by convention the standard entropy of the $\mathrm{H}^{+}$ions in water is taken as zero at all temperatures:

$$
\begin{array}{ll}
S^{\ominus}\left(\mathrm{H}^{+}, \mathrm{aq}\right)=0 & \text { lons in solution } \\
\text { [convention] }
\end{array}
$$

(3C.4)\\
Table 2C. 4 in the Resource section lists some values of standard entropies of ions in solution using this convention. ${ }^{1}$ Because the entropies of ions in water are values relative to the hydrogen ion in water, they may be either positive or negative. A positive entropy means that an ion has a higher molar entropy than $\mathrm{H}^{+}$in water and a negative entropy means that the ion has a lower molar entropy than $\mathrm{H}^{+}$in water. Ion entropies vary as expected on the basis that they are related to the degree to which the ions order the water molecules around them in the solution. Small, highly charged ions induce local structure in the surrounding water, and the disorder of the solution is decreased more than in the case of large, singly charged ions. The absolute, Third-Law standard molar entropy of the proton in water can be estimated by proposing a model of the structure it induces, and there is some agreement on the value $-21 \mathrm{~J} \mathrm{~K}^{-1} \mathrm{~mol}^{-1}$. The negative value indicates that the proton induces order in the solvent.

\subsection*{Brief illustration 3C. 4}
The standard molar entropy of $\mathrm{Cl}^{-}(\mathrm{aq})$ is $+57 \mathrm{~J} \mathrm{~K}^{-1} \mathrm{~mol}^{-1}$ and that of $\mathrm{Mg}^{2+}(\mathrm{aq})$ is $-128 \mathrm{~J} \mathrm{~K}^{-1} \mathrm{~mol}^{-1}$. That is, the molar entropy of $\mathrm{Cl}^{-}(\mathrm{aq})$ is $57 \mathrm{~J} \mathrm{~K}^{-1} \mathrm{~mol}^{-1}$ higher than that of the proton in water (presumably because it induces less local structure

\footnotetext{${ }^{1}$ In terms of the language introduced in Topic 5A, the entropies of ions in solution are actually partial molar entropies, for their values include the consequences of their presence on the organization of the solvent molecules around them.
}
in the surrounding water), whereas that of $\mathrm{Mg}^{2+}(\mathrm{aq})$ is $128 \mathrm{~J} \mathrm{~K}^{-1} \mathrm{~mol}^{-1}$ lower (presumably because its higher charge induces more local structure in the surrounding water).

\subsection*{(c) The temperature dependence of reaction entropy}
The temperature dependence of entropy is given by eqn \href{http://3C.la}{3C.la}, which for the molar entropy becomes

$$
S_{\mathrm{m}}\left(T_{2}\right)=S_{\mathrm{m}}\left(T_{1}\right)+\int_{T_{1}}^{T_{2}} \frac{C_{p, \mathrm{~m}}(T)}{T} \mathrm{~d} T
$$

This equation applies to each substance in the reaction, so from eqn 3C. 3 the temperature dependence of the standard reaction entropy, $\Delta_{\mathrm{r}} S^{\ominus}$, is


\begin{equation*}
\Delta_{\mathrm{r}} S^{\ominus}\left(T_{2}\right)=\Delta_{\mathrm{r}} S^{\ominus}\left(T_{1}\right)+\int_{T_{1}}^{T_{2}} \frac{\Delta_{\mathrm{r}} C_{p}^{\ominus}}{T} \mathrm{~d} T \tag{3C.5a}
\end{equation*}


where $\Delta_{\mathrm{r}} C_{p}^{\ominus}$ is the difference of the molar heat capacities of products and reactants under standard conditions weighted by the stoichiometric numbers that appear in the chemical equation:


\begin{equation*}
\Delta_{\mathrm{r}} C_{p}^{\ominus}=\sum_{\mathrm{r}} v_{\mathrm{J}} C_{p, \mathrm{~m}}^{\ominus}(\mathrm{J}) \tag{3C.5b}
\end{equation*}


Equation 3C.5a is analogous to Kirchhoff's law for the temperature dependence of $\Delta_{\mathrm{r}} H^{\ominus}$ (eqn 2C.7a in Topic 2C). If $\Delta_{\mathrm{r}} C_{p}^{\ominus}$ is independent of temperature in the range $T_{1}$ to $T_{2}$, the integral in eqn 3C.5a evaluates to $\Delta_{\mathrm{r}} C_{p}^{\ominus} \ln \left(T_{2} / T_{1}\right)$ and


\begin{equation*}
\Delta_{\mathrm{r}} S^{\ominus}\left(T_{2}\right)=\Delta_{\mathrm{r}} S^{\ominus}\left(T_{1}\right)+\Delta_{\mathrm{r}} C_{p}^{\ominus} \ln \frac{T_{2}}{T_{1}} \tag{3C.5c}
\end{equation*}


\subsection*{Brief illustration 3C. 5}
The standard reaction entropy for $\mathrm{H}_{2}(\mathrm{~g})+\frac{1}{2} \mathrm{O}_{2}(\mathrm{~g}) \rightarrow \mathrm{H}_{2} \mathrm{O}(\mathrm{g})$ at 298 K is $-44.42 \mathrm{~J} \mathrm{~K}^{-1} \mathrm{~mol}^{-1}$, and the molar heat capacities at constant pressure of the molecules are $\mathrm{H}_{2} \mathrm{O}(\mathrm{g})$ : $33.58 \mathrm{~J} \mathrm{~K}^{-1} \mathrm{~mol}^{-1} ; \mathrm{H}_{2}(\mathrm{~g}): 28.84 \mathrm{~J} \mathrm{~K}^{-1} \mathrm{~mol}^{-1} ; \mathrm{O}_{2}(\mathrm{~g}): 29.37 \mathrm{~J} \mathrm{~K}^{-1} \mathrm{~mol}^{-1}$. It follows that

$$
\begin{aligned}
\Delta_{\mathrm{r}} C_{p}^{\ominus} & =C_{p, \mathrm{~m}}^{\ominus}\left(\mathrm{H}_{2} \mathrm{O}, \mathrm{~g}\right)-C_{p, \mathrm{~m}}^{\ominus}\left(\mathrm{H}_{2}, \mathrm{~g}\right)-\frac{1}{2} C_{p, \mathrm{~m}}^{\ominus}\left(\mathrm{O}_{2}, \mathrm{~g}\right) \\
& =-9.94 \mathrm{~J} \mathrm{~K}^{-1} \mathrm{~mol}^{-1}
\end{aligned}
$$

This value of $\Delta_{\mathrm{r}} C_{p}^{\ominus}$ is used in eqn 3C.5c to find $\Delta_{\mathrm{r}} S^{\ominus}$ at another temperature, for example at 373 K

$$
\begin{aligned}
\Delta_{\mathrm{r}} S^{\ominus}(373 \mathrm{~K}) & =-44.42 \mathrm{~J} \mathrm{~K}^{-1} \mathrm{~mol}^{-1}+\left(-9.94 \mathrm{~J} \mathrm{~K}^{-1} \mathrm{~mol}^{-1}\right) \times \ln \frac{373 \mathrm{~K}}{298 \mathrm{~K}} \\
& =-46.65 \mathrm{~J} \mathrm{~K}^{-1} \mathrm{~mol}^{-1}
\end{aligned}
$$

\subsection*{Checklist of concepts}
$\square$ 1. Entropies are determined calorimetrically by measuring the heat capacity of a substance from low temperatures up to the temperature of interest and taking into account any phase transitions in that range.2. The Debye extrapolation (or the Debye $T^{3}$-law) is used to estimate heat capacities of non-metallic solids close to $T=0$.3. The Nernst heat theorem states that the entropy change accompanying any physical or chemical transformation approaches zero as the temperature approaches zero: $\Delta S \rightarrow 0$ as $T \rightarrow 0$ provided all the substances involved are perfectly ordered.\\
4. The Third Law of thermodynamics states that the entropy of all perfect crystalline substances is zero at $T=0$.\\
5. The residual entropy of a solid is the entropy arising from disorder that persists at $T=0$.\\
6. Third-law entropies are entropies based on $S(0)=0$.\\
7. The standard entropies of ions in solution are based on setting $S^{\ominus}\left(\mathrm{H}^{+}, \mathrm{aq}\right)=0$ at all temperatures.8. The standard reaction entropy, $\Delta_{\mathrm{t}} S^{\ominus}$, is the difference between the molar entropies of the pure, separated products and the pure, separated reactants, all substances being in their standard states.

\subsection*{Checklist of equations}
\begin{center}
\begin{tabular}{|c|c|c|c|}
\hline
Property & Equation & Comment & Equation number \\
\hline
Standard molar entropy from calorimetry & See eqn 3C. 2 & Sum of contributions from $T=0$ to temperature of interest & 3C. 2 \\
\hline
Standard reaction entropy & \( \Delta_{\mathrm{r}} S^{\ominus}=\sum_{\text {Products }} v S_{\mathrm{m}}^{\ominus}-\sum_{\text {Reactants }} v S_{\mathrm{m}}^{\ominus} \) & $v$ : (positive) stoichiometric coefficients; & 3C. 3 \\
\hline
\multirow{3}{*}{Temperature dependence of the standard reaction entropy} & \( \Delta_{\mathrm{r}} S^{\ominus}=\sum_{\mathrm{J}} v_{\mathrm{J}} S_{\mathrm{m}}^{\ominus}(\mathrm{J}) \) & $V_{\mathrm{T}}$ : (signed) stoichiometric numbers &  \\
\hline
 & \( \Delta_{\mathrm{r}} S^{\ominus}\left(T_{2}\right)=\Delta_{\mathrm{r}} S^{\ominus}\left(T_{1}\right)+\int_{T_{1}}^{T_{2}}\left(\Delta_{\mathrm{r}} C_{p}^{\ominus} / T\right) \mathrm{d} T \) &  & 3C.5a \\
\hline
 & $\Delta_{\mathrm{r}} S^{\ominus}\left(T_{2}\right)=\Delta_{\mathrm{r}} S^{\ominus}\left(T_{1}\right)+\Delta_{\mathrm{r}} C_{p}^{\ominus} \ln \left(T_{2} / T_{1}\right)$ & $\Delta_{\mathrm{r}} C_{p}^{\ominus}$ independent of temperature & 3C.5c \\
\hline
\end{tabular}
\end{center}

\section*{TOPIC 3D Concentrating on the system}
\textbackslash begin\{abstract\}\\
> Why do you need to know this material?\\
Most processes of interest in chemistry occur at constant temperature and pressure. Under these conditions, thermodynamic processes are discussed in terms of the Gibbs energy, which is introduced in this Topic. The Gibbs energy is the foundation of the discussion of phase equilibria, chemical equilibrium, and bioenergetics.\\
\textbackslash end\{abstract\}

\subsection*{> What is the key idea?}
The Gibbs energy is a signpost of spontaneous change at constant temperature and pressure, and is equal to the maximum non-expansion work that a system can do.

\subsection*{> What do you need to know already?}
This Topic develops the Clausius inequality (Topic 3A) and draws on information about standard states and reaction enthalpy introduced in Topic 2C. The derivation of the Born equation makes use of the Coulomb potential energy between two electric charges (The chemist's toolkit 6 in Topic 2A).

Entropy is the basic concept for discussing the direction of natural change, but to use it the changes in both the system and its surroundings must be analysed. In Topic 3A it is shown that it is always very simple to calculate the entropy change in the surroundings (from $\Delta S_{\text {sur }}=q_{\text {sur }} / T_{\text {sur }}$ ) and this Topic shows that it is possible to devise a simple method for taking this contribution into account automatically. This approach focuses attention on the system and simplifies discussions. Moreover, it is the foundation of all the applications of chemical thermodynamics that follow.

\subsection*{3D. 1 The Helmholtz and Gibbs energies}
Consider a system in thermal equilibrium with its surroundings at a temperature $T$. When a change in the system occurs and there is a transfer of energy as heat between the system and the surroundings, the Clausius inequality (eqn 3A.11, $\mathrm{d} S \geq \mathrm{d} q / T$ ) reads


\begin{equation*}
\mathrm{d} S-\frac{\mathrm{d} q}{T} \geq 0 \tag{3D.1}
\end{equation*}


This inequality can be developed in two ways according to the conditions (of constant volume or constant pressure) under which the process occurs.\\
(a) Criteria of spontaneity

First, consider heating at constant volume. Under these conditions and in the absence of additional (non-expansion) work $\mathrm{d} q_{V}=\mathrm{d} U$; consequently

$$
\mathrm{d} S-\frac{\mathrm{d} U}{T} \geq 0
$$

The importance of the inequality in this form is that it expresses the criterion for spontaneous change solely in terms of the state functions of the system. The inequality is easily rearranged into


\begin{equation*}
T \mathrm{~d} S \geq \mathrm{d} U \quad \text { (constant } V, \text { no additional work) } \tag{3D.2}
\end{equation*}


If the internal energy is constant, meaning that $\mathrm{d} U=0$, then it follows that $T \mathrm{~d} S \geq 0$, but as $T>0$, this relation can be written $\mathrm{d} S_{U, V} \geq 0$, where the subscripts indicate the constant conditions. This expression is a criterion for spontaneous change in terms of properties relating to the system. It states that in a system at constant volume and constant internal energy (such as an isolated system), the entropy increases in a spontaneous change. That statement is essentially the content of the Second Law.

When energy is transferred as heat at constant pressure and there is no work other than expansion work, $\mathrm{d} q_{p}=\mathrm{d} H$. Then eqn 3D. 1 becomes


\begin{equation*}
T \mathrm{~d} S \geq \mathrm{d} H \quad \text { (constant } p, \text { no additional work) } \tag{3D.3}
\end{equation*}


If the enthalpy is constant as well as the pressure, this relation becomes $T \mathrm{~d} S \geq 0$ and therefore $\mathrm{d} S \geq 0$, which may be written $\mathrm{d} S_{H, p} \geq 0$. That is, in a spontaneous process the entropy of the system at constant pressure must increase if its enthalpy remains constant (under these circumstances there can then be no change in entropy of the surroundings).

The criteria of spontaneity at constant volume and pressure can be expressed more simply by introducing two more thermodynamic quantities. One is the Helmholtz energy, $A$, which is defined as

$$
A=U-T S
$$ [definition]

The other is the Gibbs energy, $G$ :

\[
\begin{array}{ll}
G=H-T S & \begin{array}{l}
\text { Gibbs energy } \\
\text { [definition] }
\end{array} \tag{3D.4b}
\end{array}
\]

All the symbols in these two definitions refer to the system.\\
When the state of the system changes at constant temperature, the two properties change as follows:\\
(a) $\mathrm{d} A=\mathrm{d} U-T \mathrm{~d} S$\\
(b) $\mathrm{d} G=\mathrm{d} H-T \mathrm{~d} S$

At constant volume, $T \mathrm{~d} S \geq \mathrm{d} U$ (eqn 3D.2) which, by using (a), implies $\mathrm{d} A \leq 0$. At constant pressure, $T \mathrm{~d} S \geq \mathrm{d} H$ (eqn 3D.3) which, by using (b), implies $\mathrm{d} G \leq 0$. Using the subscript notation to indicate which variables are held constant, the criteria of spontaneous change in terms of $\mathrm{d} A$ and $\mathrm{d} G$ are

$$
\begin{array}{lll}
\text { (a) } \mathrm{d} A_{T, V} \leq 0 & \text { (b) } \mathrm{d} G_{T, p} \leq 0 & \begin{array}{l}
\text { Criteria of spontaneous } \\
\text { change }
\end{array}
\end{array}
$$

(3D.6)\\
These criteria, especially the second, are central to chemical thermodynamics. For instance, in an endothermic reaction $H$ increases, $\mathrm{d} H>0$, but if such a reaction is to be spontaneous at constant temperature and pressure, $G$ must decrease. Because $\mathrm{d} G=\mathrm{d} H-T \mathrm{~d} S$, it is possible for $\mathrm{d} G$ to be negative provided that the entropy of the system increases so much that $T \mathrm{~d} S$ outweighs $\mathrm{d} H$. Endothermic reactions are therefore driven by the increase of entropy of the system, which overcomes the reduction of entropy brought about in the surroundings by the inflow of heat into the system in an endothermic process $\left(\mathrm{d} S_{\text {sur }}=-\mathrm{d} H / T\right.$ at constant pressure). Exothermic reactions are commonly spontaneous because $\mathrm{d} H<0$ and then $\mathrm{d} G<0$ provided $T \mathrm{~d} S$ is not so negative that it outweighs the decrease in enthalpy.

\subsection*{(b) Some remarks on the Helmholtz energy}
At constant temperature and volume, a change is spontaneous if it corresponds to a decrease in the Helmholtz energy: $\mathrm{d} A_{T, V}$ $\leq 0$. Such systems move spontaneously towards states of lower $A$ if a path is available. The criterion of equilibrium, when neither the forward nor reverse process has a tendency to occur, is $\mathrm{d} A_{T, V}=0$.

The expressions $\mathrm{d} A=\mathrm{d} U-T \mathrm{~d} S$ and $\mathrm{d} A_{T, V} \leq 0$ are sometimes interpreted as follows. A negative value of $\mathrm{d} A$ is favoured by a negative value of $\mathrm{d} U$ and a positive value of $T \mathrm{~d} S$. This observation suggests that the tendency of a system to move to lower $A$ is due to its tendency to move towards states of lower internal energy and higher entropy. However, this interpretation is false because the tendency to lower $A$ is solely a tendency towards states of greater overall entropy. Systems change spontaneously if in doing so the total entropy of the system and its surroundings increases, not because they tend to lower internal energy. The form of $\mathrm{d} A$ may give the impression that systems favour lower energy, but that is misleading: $\mathrm{d} S$ is the entropy\\
change of the system, $-\mathrm{d} U / T$ is the entropy change of the surroundings (when the volume of the system is constant), and their total tends to a maximum.

\subsection*{(c) Maximum work}
As well as being the signpost of spontaneous change, a short argument can be used to show that the change in the Helmholtz energy is equal to the maximum work obtainable from a system at constant temperature.

\subsection*{How is that done? 3D. 1 Relating the change in the}
 Helmholtz energy to the maximum workTo demonstrate that maximum work can be expressed in terms of the change in Helmholtz energy, you need to combine the Clausius inequality $\mathrm{d} S \geq \mathrm{d} q / T$ in the form $T \mathrm{~d} S \geq \mathrm{d} q$ with the First Law, $\mathrm{d} U=\mathrm{d} q+\mathrm{d} w$, and obtain

$$
\mathrm{d} U \leq T \mathrm{~d} S+\mathrm{d} w
$$

The term $\mathrm{d} U$ is smaller than the sum of the two terms on the right because $\mathrm{d} q$ has been replaced by $T \mathrm{~d} S$, which in general is larger than $\mathrm{d} q$. This expression rearranges to

$$
\mathrm{d} w \geq \mathrm{d} U-T \mathrm{~d} S
$$

It follows that the most negative value of $\mathrm{d} w$ is obtained when the equality applies, which is for a reversible process. Thus a reversible process gives the maximum amount of energy as work, and this maximum work is given by

$$
\mathrm{d} w_{\max }=\mathrm{d} U-T \mathrm{~d} S
$$

Because at constant temperature $\mathrm{d} A=\mathrm{d} U-T \mathrm{~d} S$ (eqn 3D.5), it follows that

$$
\longrightarrow \mathrm{d} w_{\max }=\mathrm{d} A \longmapsto \begin{aligned}
& \text { Maximum work } \\
& \text { [constant } T \text { ] }
\end{aligned}
$$

In recognition of this relation, $A$ is sometimes called the 'maximum work function', or the 'work function'.'

When a measurable isothermal change takes place in the system, eqn 3D. 7 becomes $w_{\max }=\Delta A$ with $\Delta A=\Delta U-T \Delta S$. These relations show that, depending on the sign of $T \Delta S$, not all the change in internal energy may be available for doing work. If the change occurs with a decrease in entropy (of the system), in which case $T \Delta S<0$, then $\Delta U-T \Delta S$ is not as negative as $\Delta U$ itself, and consequently the maximum work is less than $\Delta U$. For the change to be spontaneous, some of the energy must escape as heat in order to generate enough entropy in the surroundings to overcome the reduction in entropy in the system (Fig. 3D.1). In this case, Nature is demanding a tax

\footnotetext{${ }^{1}$ Arbeit is the German word for work; hence the symbol A.
}
\includegraphics[max width=\textwidth, center]{2025_02_27_d4b95d9718f0b9e2e6b9g-131(2)}

Figure 3D. 1 In a system not isolated from its surroundings, the work done may be different from the change in internal energy. In the process depicted here, the entropy of the system decreases, so for the process to be spontaneous the entropy of the surroundings must increase, so energy must pass from the system to the surroundings as heat. Therefore, less work than $\Delta U$ can be obtained.\\
on the internal energy as it is converted into work. This interpretation is the origin of the alternative name 'Helmholtz free energy' for $A$, because $\Delta A$ is that part of the change in internal energy free to do work.

Further insight into the relation between the work that a system can do and the Helmholtz energy is to recall that work is energy transferred to the surroundings as the uniform motion of atoms. The expression $A=U-T S$ can be interpreted as showing that $A$ is the total internal energy of the system, $U$, less a contribution that is stored as energy of thermal motion (the quantity TS). Because energy stored in random thermal motion cannot be used to achieve uniform motion in the surroundings, only the part of $U$ that is not stored in that way, the quantity $U-T S$, is available for conversion into work.

If the change occurs with an increase of entropy of the system (in which case $T \Delta S>0$ ), $\Delta U-T \Delta S$ is more negative than $\Delta U$. In this case, the maximum work that can be obtained from the system is greater than $\Delta U$. The explanation of this apparent paradox is that the system is not isolated and energy\\
\includegraphics[max width=\textwidth, center]{2025_02_27_d4b95d9718f0b9e2e6b9g-131}

Figure 3D. 2 In this process, the entropy of the system increases; hence some reduction in the entropy of the surroundings can be tolerated. That is, some of their energy may be lost as heat to the system. This energy can be returned to them as work, and hence the work done can exceed $\Delta U$.\\
may flow in as heat as work is done. Because the entropy of the system increases, a reduction of the entropy of the surroundings can be afforded yet still have, overall, a spontaneous process. Therefore, some energy (no more than the value of $T \Delta S$ ) may leave the surroundings as heat and contribute to the work the change is generating (Fig. 3D.2). Nature is now providing a tax refund.

Example 3D. 1 Calculating the maximum available work\\
When $1.000 \mathrm{~mol} \mathrm{C}_{6} \mathrm{H}_{12} \mathrm{O}_{6}$ (glucose) is oxidized completely to carbon dioxide and water at $25^{\circ} \mathrm{C}$ according to the equation $\mathrm{C}_{6} \mathrm{H}_{12} \mathrm{O}_{6}(\mathrm{~s})+6 \mathrm{O}_{2}(\mathrm{~g}) \rightarrow 6 \mathrm{CO}_{2}(\mathrm{~g})+6 \mathrm{H}_{2} \mathrm{O}(\mathrm{l})$, calorimetric measurements give $\Delta_{\mathrm{r}} U=-2808 \mathrm{~kJ} \mathrm{~mol}^{-1}$ and $\Delta_{\mathrm{r}} S=+182.4 \mathrm{~J} \mathrm{~K}^{-1} \mathrm{~mol}^{-1}$ at $25^{\circ} \mathrm{C}$ and 1 bar. How much of this change in internal energy can be extracted as (a) heat at constant pressure, (b) work?

Collect your thoughts You know that the heat released at constant pressure is equal to the value of $\Delta H$, so you need to relate $\Delta_{\mathrm{r}} H$ to the given value of $\Delta_{\mathrm{r}} U$. To do so, suppose that all the gases involved are perfect, and use eqn 2B. $4\left(\Delta H=\Delta U+\Delta n_{\mathrm{g}} R T\right)$ in the form $\Delta_{\mathrm{r}} H=\Delta_{\mathrm{r}} U+\Delta_{\mathrm{g}} R T$. For the maximum work available from the process use $w_{\max }=\Delta A$ in the form $w_{\max }=\Delta_{\mathrm{r}} A$.

The solution (a) Because $\Delta v_{\mathrm{g}}=0, \Delta_{\mathrm{r}} H=\Delta_{\mathrm{r}} U=-2808 \mathrm{~kJ} \mathrm{~mol}^{-1}$. Therefore, at constant pressure, the energy available as heat is $2808 \mathrm{~kJ} \mathrm{~mol}^{-1}$. (b) Because $T=298 \mathrm{~K}$, the value of $\Delta_{\mathrm{r}} A$ is

$$
\Delta_{\mathrm{r}} A=\Delta_{\mathrm{r}} U-T \Delta_{\mathrm{r}} S=-2862 \mathrm{~kJ} \mathrm{~mol}^{-1}
$$

Therefore, the complete oxidation of $1.000 \mathrm{~mol} \mathrm{C}_{6} \mathrm{H}_{12} \mathrm{O}_{6}$ at constant temperature can be used to produce up to 2862 kJ of work.

Comment. The maximum work available is greater than the change in internal energy on account of the positive entropy of reaction (which is partly due to there being a significant increase in the number of molecules as the reaction proceeds). The system can therefore draw in energy from the surroundings (so reducing their entropy) and make it available for doing work.

Self-test 3D.1 Repeat the calculation for the combustion of $1.000 \mathrm{~mol} \mathrm{CH}_{4}(\mathrm{~g})$ under the same conditions, using data from Table 2C. 3 and that $\Delta_{\mathrm{r}} S$ for the reaction is $-243 \mathrm{~J} \mathrm{~K}^{-1} \mathrm{~mol}^{-1}$ at 298 K .\\
\includegraphics[max width=\textwidth, center]{2025_02_27_d4b95d9718f0b9e2e6b9g-131(1)}

\subsection*{(d) Some remarks on the Gibbs energy}
The Gibbs energy (the 'free energy') is more common in chemistry than the Helmholtz energy because, at least in laboratory chemistry, changes occurring at constant pressure are more common than at constant volume. The criterion $\mathrm{d} G_{T, p} \leq 0$ carries over into chemistry as the observation that, at constant temperature and pressure, chemical reactions are spontaneous\\
in the direction of decreasing Gibbs energy. Therefore, to decide whether a reaction is spontaneous, the pressure and temperature being constant, it is necessary to assess the change in the Gibbs energy. If $G$ decreases as the reaction proceeds, then the reaction has a spontaneous tendency to convert the reactants into products. If $G$ increases, the reverse reaction is spontaneous. The criterion for equilibrium, when neither the forward nor reverse process is spontaneous, under conditions of constant temperature and pressure, is $\mathrm{d} G_{T, p}=0$.\\
The existence of spontaneous endothermic reactions provides an illustration of the role of $G$. In such reactions, $H$ increases, the system rises spontaneously to states of higher enthalpy, and $\mathrm{d} H>0$. Because the reaction is spontaneous, $\mathrm{d} G$ $<0$ despite $\mathrm{d} H>0$; it follows that the entropy of the system increases so much that $T \mathrm{~d} S$ outweighs $\mathrm{d} H$ in $\mathrm{d} G=\mathrm{d} H-T \mathrm{~d} S$. Endothermic reactions are therefore driven by the increase of entropy of the system, and this entropy change overcomes the reduction of entropy brought about in the surroundings by the inflow of heat into the system $\left(\mathrm{d} S_{\text {sur }}=-\mathrm{d} H / T\right.$ at constant pressure). Exothermic reactions are commonly spontaneous because $\mathrm{d} H<0$ and then $\mathrm{d} G<0$ provided $T \mathrm{~d} S$ is not so negative that it outweighs the decrease in enthalpy.

\subsection*{(e) Maximum non-expansion work}
The analogue of the maximum work interpretation of $\Delta A$, and the origin of the name 'free energy', can be found for $\Delta G$. By an argument like that relating the Helmholtz energy to maximum work, it can be shown that, at constant temperature and pressure, the change in Gibbs energy is equal to the maximum additional (non-expansion) work.

How is that done? 3D. 2 Relating the change in Gibbs energy to maximum non-expansion work

Because $H=U+p V$ and $\mathrm{d} U=\mathrm{d} q+\mathrm{d} w$, the change in enthalpy for a general change in conditions is

$$
\mathrm{d} H=\mathrm{d} q+\mathrm{d} w+\mathrm{d}(p V)
$$

The corresponding change in Gibbs energy ( $G=H-T S$ ) is

$$
\mathrm{d} G=\mathrm{d} H-T \mathrm{~d} S-S \mathrm{~d} T=\mathrm{d} q+\mathrm{d} w+\mathrm{d}(p V)-T \mathrm{~d} S-S \mathrm{~d} T
$$

Step 1 Confine the discussion to constant temperature When the change is isothermal $\mathrm{d} T=0$; then

$$
\mathrm{d} G=\mathrm{d} q+\mathrm{d} w+\mathrm{d}(p V)-T \mathrm{~d} S
$$

\subsection*{Step 2 Confine the change to a reversible process}
When the change is reversible, $\mathrm{d} w=\mathrm{d} w_{\mathrm{rev}}$ and $\mathrm{d} q=\mathrm{d} q_{\mathrm{rev}}=T \mathrm{~d} S$, so for a reversible, isothermal process

$$
\mathrm{d} G=\stackrel{\mathrm{d} q_{\text {rev }}}{T \mathrm{~d} S}+\mathrm{d} w_{\text {rev }}+\mathrm{d}(p V)-T \mathrm{~d} S=\mathrm{d} w_{\text {rev }}+\mathrm{d}(p V)
$$

\subsection*{Step 3 Divide the work into different types}
The work consists of expansion work, which for a reversible change is given by $-p \mathrm{~d} V$, and possibly some other kind of work (for instance, the electrical work of pushing electrons through a circuit or of raising a column of liquid); this additional work is denoted $\mathrm{d} w_{\text {add }}$. Therefore, with $\mathrm{d}(p V)=p \mathrm{~d} V+V \mathrm{~d} p$,

$$
\mathrm{d} G=\overbrace{\left(-p \mathrm{~d} V+\mathrm{d} w_{\text {add,rev }}\right)}^{\mathrm{d} w_{\text {rev }}}+\overbrace{p \mathrm{~d} V+V \mathrm{~d} p}^{\mathrm{d}(p V)}=\mathrm{d} w_{\text {add,rev }}+V \mathrm{~d} p
$$

\subsection*{Step 4 Confine the process to constant pressure}
If the change occurs at constant pressure (as well as constant temperature), $\mathrm{d} p=0$ and hence $\mathrm{d} G=\mathrm{d} w_{\text {add,rev }}$. Therefore, at constant temperature and pressure, $\mathrm{d} w_{\text {add,rev }}=\mathrm{d} G$. However, because the process is reversible, the work done must now have its maximum value, so it follows that\\
\includegraphics[max width=\textwidth, center]{2025_02_27_d4b95d9718f0b9e2e6b9g-132}

For a measurable change, the corresponding expression is $w_{\text {add,max }}=\Delta G$. This is particularly useful for assessing the maximum electrical work that can be produced by fuel cells and electrochemical cells (Topic 6C).

\subsection*{3D.2 Standard molar Gibbs energies}
Standard entropies and enthalpies of reaction (which are introduced in Topics 2C and 3C) can be combined to obtain the standard Gibbs energy of reaction (or 'standard reaction Gibbs energy'), $\Delta_{\mathrm{r}} G^{\ominus}$ :

\[
\Delta_{\mathrm{r}} G^{\ominus}=\Delta_{\mathrm{r}} H^{\ominus}-T \Delta_{\mathrm{r}} S^{\ominus} \quad \begin{align*}
& \text { Standard Gibbs energy of reaction }  \tag{3D.9}\\
& \text { [definition] }
\end{align*}
\]

The standard Gibbs energy of reaction is the difference in standard molar Gibbs energies of the products and reactants in their standard states for the reaction as written and at the temperature specified.

Calorimetry (for $\Delta H$ directly, and for $S$ from heat capacities) is only one of the ways of determining Gibbs energies. They may also be obtained from equilibrium constants (Topic 6A) and electrochemical measurements (Topic 6D), and for gases they may be calculated using data from spectroscopic observations (Topic 13E).

Example 3D. 2 Calculating the maximum non-expansion work of a reaction

How much energy is available for sustaining muscular and nervous activity from the oxidation of 1.00 mol of glucose molecules under standard conditions at $37^{\circ} \mathrm{C}$ (blood temperature)? The standard entropy of reaction is $+182.4 \mathrm{~J} \mathrm{~K}^{-1} \mathrm{~mol}^{-1}$.

Collect your thoughts The non-expansion work available from the reaction at constant temperature and pressure is equal to the change in standard Gibbs energy for the reaction, $\Delta_{\mathrm{r}} G^{\ominus}$. To calculate this quantity, you can (at least approximately) ignore the temperature dependence of the reaction enthalpy, and obtain $\Delta_{\mathrm{r}} H^{\ominus}$ from Table 2C. 4 (where the data are for $25^{\circ} \mathrm{C}$, not $37^{\circ} \mathrm{C}$ ), and substitute the data into $\Delta_{\mathrm{r}} G^{\ominus}=\Delta_{\mathrm{r}} H^{\ominus}-$ $T \Delta_{\mathrm{r}} S^{\ominus}$.

The solution Because the standard reaction enthalpy is $-2808 \mathrm{~kJ} \mathrm{~mol}^{-1}$, it follows that the standard reaction Gibbs energy is\\
$\Delta_{\mathrm{r}} G^{\ominus}=-2808 \mathrm{~kJ} \mathrm{~mol}^{-1}-(310 \mathrm{~K}) \times\left(182.4 \mathrm{JK}^{-1} \mathrm{~mol}^{-1}\right)=-2865 \mathrm{~kJ} \mathrm{~mol}^{-1}$\\
Therefore, $w_{\text {add,max }}=-2865 \mathrm{~kJ}$ for the oxidation of 1 mol glucose molecules, and the reaction can be used to do up to 2865 kJ of non-expansion work.

Comment. To place this result in perspective, consider that a person of mass 70 kg needs to do 2.1 kJ of work to climb vertically through 3.0 m ; therefore, at least 0.13 g of glucose is needed to complete the task (and in practice significantly more).

Self-test 3D. 2 How much non-expansion work can be obtained from the combustion of $1.00 \mathrm{~mol} \mathrm{CH}_{4}(\mathrm{~g})$ under standard conditions at 298 K ? Use $\Delta_{\mathrm{r}} \mathrm{S}^{\ominus}=-243 \mathrm{~J} \mathrm{~K}^{-1} \mathrm{~mol}^{-1}$.\\
fy8I8:ءวмsuV

\subsection*{(a) Gibbs energies of formation}
As in the case of standard reaction enthalpies (Topic 2C), it is convenient to define the standard Gibbs energies of formation, $\Delta_{\mathrm{f}} G^{\ominus}$, the standard reaction Gibbs energy for the formation of a compound from its elements in their reference states, as specified in Topic 2C. Standard Gibbs energies of formation of the elements in their reference states are zero, because their formation is a 'null' reaction. A selection of values for compounds is given in Table 3D.1. The standard Gibbs energy of a reaction is then found by taking the appropriate combination:

Table 3D. 1 Standard Gibbs energies of formation at 298 K $^{*}$

\begin{center}
\begin{tabular}{lc}
\hline
 & $\Delta_{\mathrm{f}} G^{\ominus} /(\mathbf{k ~ J ~ m o l}$ \\
 &  \\
\hline
Diamond, $\mathrm{C}(\mathrm{s})$ & +2.9 \\
Benzene, $\mathrm{C}_{6} \mathrm{H}_{6}(\mathrm{l})$ & +124.3 \\
Methane, $\mathrm{CH}_{4}(\mathrm{~g})$ & -50.7 \\
Carbon dioxide, $\mathrm{CO}_{2}(\mathrm{~g})$ & -394.4 \\
Water, $\mathrm{H}_{2} \mathrm{O}(\mathrm{l})$ & -237.1 \\
Ammonia, $\mathrm{NH}_{3}(\mathrm{~g})$ & -16.5 \\
Sodium chloride, $\mathrm{NaCl}(\mathrm{s})$ & -384.1 \\
\hline
\end{tabular}
\end{center}

\begin{itemize}
  \item More values are given in the Resource section.
\end{itemize}

$$
\Delta_{\mathrm{r}} G^{\ominus}=\sum_{\text {Products }} v \Delta_{\mathrm{f}} G^{\ominus}-\sum_{\text {Reactants }} v \Delta_{\mathrm{f}} G^{\ominus} \quad \begin{aligned}
& \text { Standard Gibbs } \\
& \text { energy of reaction } \\
& \text { [practical } \\
& \text { implementation] }
\end{aligned}
$$

(3D.10a)

In the notation introduced in Topic 2C,


\begin{equation*}
\Delta_{\mathrm{r}} G^{\ominus}=\sum_{\mathrm{J}} v_{\mathrm{J}} \Delta_{\mathrm{f}} G^{\ominus}(\mathrm{J}) \tag{3D.10b}
\end{equation*}


where the $v_{\mathrm{J}}$ are the (signed) stoichiometric numbers in the chemical equation.

\subsection*{Brief illustration 3D. 1}
To calculate the standard Gibbs energy of the reaction $\mathrm{CO}(\mathrm{g})$ $+\frac{1}{2} \mathrm{O}_{2}(\mathrm{~g}) \rightarrow \mathrm{CO}_{2}(\mathrm{~g})$ at $25^{\circ} \mathrm{C}$, write

$$
\begin{aligned}
\Delta_{\mathrm{r}} G^{\ominus} & =\Delta_{\mathrm{f}} G^{\ominus}\left(\mathrm{CO}_{2}, \mathrm{~g}\right)-\left\{\Delta_{\mathrm{f}} G^{\ominus}(\mathrm{CO}, \mathrm{~g})+\frac{1}{2} \Delta_{\mathrm{f}} G^{\ominus}\left(\mathrm{O}_{2}, \mathrm{~g}\right)\right\} \\
& =-394.4 \mathrm{~kJ} \mathrm{~mol}^{-1}-\left\{(-137.2)+\frac{1}{2}(0)\right\} \mathrm{kJ} \mathrm{~mol}^{-1} \\
& =-257.2 \mathrm{~kJ} \mathrm{~mol}^{-1}
\end{aligned}
$$

As explained in Topic 2C the standard enthalpy of formation of $\mathrm{H}^{+}$in water is by convention taken to be zero; in Topic 3 C , the absolute entropy of $\mathrm{H}^{+}(\mathrm{aq})$ is also by convention set equal to zero (at all temperatures in both cases). These conventions are needed because it is not possible to prepare cations without their accompanying anions. For the same reason, the standard Gibbs energy of formation of $\mathrm{H}^{+}(\mathrm{aq})$ is set equal to zero at all temperatures:

\[
\begin{array}{ll}
\Delta_{\mathrm{f}} G^{\ominus}\left(\mathrm{H}^{+}, \mathrm{aq}\right)=0 & \begin{array}{l}
\text { lons in solution } \\
\text { [convention] }
\end{array} \tag{3D.11}
\end{array}
\]

This definition effectively adjusts the actual values of the Gibbs energies of formation of ions by a fixed amount, which is chosen so that the standard value for one of them, $\mathrm{H}^{+}(\mathrm{aq})$, has the value zero.

\subsection*{Brief illustration 3D. 2}
For the reaction\\
$\frac{1}{2} \mathrm{H}_{2}(\mathrm{~g})+\frac{1}{2} \mathrm{Cl}_{2}(\mathrm{~g}) \rightarrow \mathrm{H}^{+}(\mathrm{aq})+\mathrm{Cl}^{-}(\mathrm{aq}) \quad \Delta_{\mathrm{r}} G^{\ominus}=-131.23 \mathrm{~kJ} \mathrm{~mol}^{-1}$\\
the value of $\Delta_{\mathrm{r}} G^{\ominus}$ can be written in terms of standard Gibbs energies of formation as

$$
\Delta_{\mathrm{r}} G^{\ominus}=\Delta_{\mathrm{f}} G^{\ominus}\left(\mathrm{H}^{+}, \mathrm{aq}\right)+\Delta_{\mathrm{f}} G^{\ominus}\left(\mathrm{Cl}^{-}, \mathrm{aq}\right)
$$

where the $\Delta_{f} G^{\ominus}$ of the elements on the left of the chemical equation are zero. Because by convention $\Delta_{f} G^{\ominus}\left(\mathrm{H}^{+}, \mathrm{aq}\right)$ $=0$, it follows that $\Delta_{\mathrm{r}} G^{\ominus}=\Delta_{\mathrm{f}} G^{\ominus}\left(\mathrm{Cl}^{-}, \mathrm{aq}\right)$ and therefore that $\Delta_{\mathrm{f}} G^{\ominus}\left(\mathrm{Cl}^{-}, \mathrm{aq}\right)=-131.23 \mathrm{~kJ} \mathrm{~mol}^{-1}$.\\
\includegraphics[max width=\textwidth, center]{2025_02_27_d4b95d9718f0b9e2e6b9g-134}

Figure 3D.3 A thermodynamic cycle for discussion of the Gibbs energies of hydration and formation of chloride ions in aqueous solution. The changes in Gibbs energies around the cycle sum to zero because $G$ is a state function.

The factors responsible for the Gibbs energy of formation of an ion in solution can be identified by analysing its formation in terms of a thermodynamic cycle. As an illustration, consider the standard Gibbs energy of formation of $\mathrm{Cl}^{-}$in water. The formation reaction $\frac{1}{2} \mathrm{H}_{2}(\mathrm{~g})+\frac{1}{2} \mathrm{Cl}_{2}(\mathrm{~g}) \rightarrow \mathrm{H}^{+}(\mathrm{aq})+\mathrm{Cl}^{-}(\mathrm{aq})$ is treated as the outcome of the sequence of steps shown in Fig. 3D. 3 (with values taken from the Resource section). The sum of the Gibbs energies for all the steps around a closed cycle is zero, so

$$
\Delta_{\mathrm{f}} G^{\ominus}\left(\mathrm{Cl}^{-}, \mathrm{aq}\right)=1287 \mathrm{~kJ} \mathrm{~mol}^{-1}+\Delta_{\text {solv }} G^{\ominus}\left(\mathrm{H}^{+}\right)+\Delta_{\text {solv }} G^{\ominus}\left(\mathrm{Cl}^{-}\right)
$$

The standard Gibbs energies of formation of the gas-phase ions are unknown and have been replaced by energies and electron affinities and the assumption that any differences from the Gibbs energies arising from conversion to enthalpy and the inclusion of entropies to obtain Gibbs energies in the formation of $\mathrm{H}^{+}$are cancelled by the corresponding terms in the electron gain of Cl . The conclusions from the cycles are therefore only approximate. An important point to note is that the value of $\Delta_{\mathrm{f}} G^{\ominus}$ of $\mathrm{Cl}^{-}$is not determined by the properties of Cl alone but includes contributions from the dissociation, ionization, and hydration of hydrogen.

\subsection*{(b) The Born equation}
Gibbs energies of solvation of individual ions may be estimated on the basis of a model in which solvation is expressed as an electrostatic property.

How is that done? 3D. 3\\
Developing an electrostatic model for solvation

The model treats the interaction between the ion and the solvent using elementary electrostatics: the ion is regarded as a charged sphere and the solvent is treated as a continuous medium (a continuous dielectric). The key step is to use the result from Section 3D.1(e) to identify the Gibbs energy of solvation with the work of transferring an ion from a vacuum into the solvent. That work is calculated by taking the difference of the work of charging an ion when it is in the solution and the work of charging the same ion when it is in a vacuum.

The derivation uses concepts developed in The chemist's toolkit 6 in Topic 2A, where it is seen that the Coulomb potential energy of two point electric charges $Q_{1}$ and $Q_{2}$ separated by a distance $r$ in a medium with permittivity $\varepsilon$ is

$$
V(r)=\frac{Q_{1} Q_{2}}{4 \pi \varepsilon r}
$$

The energy of this interaction may also be expressed in terms of the Coulomb potential $\phi$ that the point charge $Q_{2}$ experiences at a distance $r$ from the point charge $Q_{1}$. Then $V(r)=$ $Q_{2} \phi(r)$, with

$$
\phi(r)=\frac{Q_{1}}{4 \pi \varepsilon r}
$$

With the distance $r$ in metres and the charge $Q_{1}$ in coulombs (C), the potential is obtained in $\mathrm{JC}^{-1}$. By definition, $1 \mathrm{JC}^{-1}=$ 1 V (volt), so $\phi$ can also be expressed in volts.

Step 1 Obtain an expression for charging a spherical ion to its final value in a medium\\
The Coulomb potential, $\phi$, at the surface of a sphere (representing the ion) of radius $r_{\mathrm{i}}$ and charge $Q$ is the same as the potential due to a point charge at its centre, so

$$
\phi\left(r_{\mathrm{i}}\right)=\frac{Q}{4 \pi \varepsilon r_{\mathrm{i}}}
$$

The work of bringing up a charge $\mathrm{d} Q$ to the sphere is $\phi\left(r_{\mathrm{i}}\right) \mathrm{d} Q$. If the charge number of the ion is $z_{\mathrm{i}}$, the total work of charging the sphere from 0 to $z_{\mathrm{i}} e$ is

$$
w=\int_{0}^{z_{i} e} \phi\left(r_{\mathrm{i}}\right) \mathrm{d} Q=\frac{1}{4 \pi \varepsilon r_{\mathrm{i}}} \int_{0}^{z_{i} e} Q \mathrm{~d} Q=\frac{z_{\mathrm{i}}^{2} e^{2}}{8 \pi \varepsilon r_{\mathrm{i}}}
$$

This electrical work of charging, when multiplied by Avogadro's constant, $N_{\mathrm{A}}$, is the molar Gibbs energy for charging the ions.

\subsection*{Step 2 Apply the result to solution and a vacuum}
The work of charging an ion in a vacuum is obtained by setting $\varepsilon=\varepsilon_{0}$, the vacuum permittivity. The corresponding value for charging the ion in a medium is obtained by setting $\varepsilon=$ $\varepsilon_{\mathrm{r}} \varepsilon_{0}$, where $\varepsilon_{\mathrm{r}}$ is the relative permittivity of the medium.

$$
w(\text { vacuum })=\frac{z_{\mathrm{i}}^{2} e^{2}}{8 \pi \varepsilon_{0} r_{\mathrm{i}}} \quad w(\text { medium })=\frac{z_{\mathrm{i}}^{2} e^{2}}{8 \pi \varepsilon_{\mathrm{r}} \varepsilon_{0} r_{\mathrm{i}}}
$$

Step 3 Identify the Gibbs energy of solvation as the work needed to move the ion from a vacuum into the medium\\
It follows that the change in molar Gibbs energy that accompanies the transfer of ions from a vacuum to a solvent is the difference of these two expressions for the work of charging:

$$
\Delta_{\text {solv }} G^{\ominus}=\frac{z_{\mathrm{i}}^{2} e^{2} N_{\mathrm{A}}}{8 \pi \varepsilon r_{\mathrm{i}}}-\frac{z_{\mathrm{i}}^{2} e^{2} N_{\mathrm{A}}}{8 \pi \varepsilon_{0} r_{\mathrm{i}}}=\frac{z_{\mathrm{i}}^{2} e^{2} N_{\mathrm{A}}}{8 \pi \varepsilon_{\mathrm{r}} \varepsilon_{0} r_{\mathrm{i}}}-\frac{z_{\mathrm{i}}^{2} e^{2} N_{\mathrm{A}}}{8 \pi \varepsilon_{0} r_{\mathrm{i}}}
$$

A minor rearrangement of the right-hand side gives the Born equation:\\
$-\Delta_{\text {solv }} G^{\ominus}=-\frac{z_{\mathrm{i}}^{2} e^{2} N_{\mathrm{A}}}{8 \pi \varepsilon_{0} r_{\mathrm{i}}}\left(1-\frac{1}{\varepsilon_{\mathrm{r}}}\right)$\\
(3D.12a)

Note that $\Delta_{\text {solv }} G^{\ominus}<0$, and that $\Delta_{\text {solv }} G^{\ominus}$ is strongly negative for small, highly charged ions in media of high relative\\
permittivity. For water, for which $\varepsilon_{\mathrm{r}}=78.54$ at $25^{\circ} \mathrm{C}$, the Born equation becomes


\begin{equation*}
\Delta_{\text {solv }} G^{\ominus}=-\frac{z_{\mathrm{i}}^{2}}{r_{\mathrm{i}} / \mathrm{pm}} \times 6.86 \times 10^{4} \mathrm{~kJ} \mathrm{~mol}^{-1} \tag{3D.12b}
\end{equation*}


\subsection*{Brief illustration 3D. 3}
To estimate the difference in the values of $\Delta_{\mathrm{f}} G^{\ominus}$ for $\mathrm{Cl}^{-}$and $\mathrm{I}^{-}$in water at $25^{\circ} \mathrm{C}$, given their radii as 181 pm and 220 pm , respectively, write

$$
\begin{aligned}
\Delta_{\text {solv }} G^{\ominus}\left(\mathrm{Cl}^{-}\right)-\Delta_{\text {solv }} G^{\ominus}\left(\mathrm{I}^{-}\right) & =-\left(\frac{1}{181}-\frac{1}{220}\right) \times 6.86 \times 10^{4} \mathrm{~kJ} \mathrm{~mol}^{-1} \\
& =-67 \mathrm{~kJ} \mathrm{~mol}^{-1}
\end{aligned}
$$

\subsection*{Checklist of concepts}
$\square$ 1. The Clausius inequality implies a number of criteria for spontaneous change under a variety of conditions which may be expressed in terms of the properties of the system alone; they are summarized by introducing the Helmholtz and Gibbs energies.\\
$\square$ 2. A spontaneous process at constant temperature and volume is accompanied by a decrease in the Helmholtz energy.\\
$\square$ 3. The change in the Helmholtz energy is equal to the maximum work obtainable from a system at constant temperature.\\
$\square$ 4. A spontaneous process at constant temperature and pressure is accompanied by a decrease in the Gibbs energy.\\
5. The change in the Gibbs energy is equal to the maximum non-expansion work obtainable from a system at constant temperature and pressure.6. Standard Gibbs energies of formation are used to calculate the standard Gibbs energies of reactions.\\
$\square$ 7. The standard Gibbs energies of formation of ions may be estimated from a thermodynamic cycle and the Born equation.

\subsection*{Checklist of equations}
\begin{center}
\begin{tabular}{|c|c|c|c|}
\hline
Property & Equation & Comment & Equation number \\
\hline
\multirow[t]{2}{*}{Criteria of spontaneity} & $\mathrm{d} S_{U, V} \geq 0$ & Subscripts show which variables are held constant, here and below &  \\
\hline
 & $\mathrm{d} S_{H, p} \geq 0$ &  &  \\
\hline
Helmholtz energy & $A=U-T S$ & Definition & 3D.4a \\
\hline
Gibbs energy & $G=H-T S$ & Definition & 3D.4b \\
\hline
Criteria of spontaneous change & $\begin{array}{ll}\text { (a) } \mathrm{d} A_{T, V} \leq 0 & \text { (b) } \mathrm{d} G_{T, p} \leq 0\end{array}$ & Equality refers to equilibrium & 3D. 6 \\
\hline
Maximum work & $\mathrm{d} w_{\text {max }}=\mathrm{d} A, w_{\text {max }}=\Delta A$ & Constant temperature & 3D. 7 \\
\hline
Maximum non-expansion work & $\mathrm{d} w_{\text {add,max }}=\mathrm{d} G, w_{\text {add,max }}=\Delta G$ & Constant temperature and pressure & 3D. 8 \\
\hline
\multirow[t]{2}{*}{Standard Gibbs energy of reaction} & $\Delta_{\mathrm{r}} G^{\ominus}=\Delta_{\mathrm{r}} H^{\ominus}-T \Delta_{\mathrm{r}} S^{\ominus}$ & Definition & 3D. 9 \\
\hline
 & \( \Delta_{\mathrm{r}} G^{\ominus}=\sum_{\mathrm{J}} v_{\mathrm{I}} \Delta_{\mathrm{f}} G^{\ominus}(\mathrm{J}) \) & Practical implementation & 3D.10b \\
\hline
Ions in solution & $\Delta_{\mathrm{f}} G^{\ominus}\left(\mathrm{H}^{+}, \mathrm{aq}\right)=0$ & Convention & 3D. 11 \\
\hline
Born equation & $\Delta_{\text {solv }} G^{\ominus}=-\left(z_{\mathrm{i}}^{2} e^{2} N_{\mathrm{A}} / 8 \pi \varepsilon_{0} r_{\mathrm{i}}\right)\left(1-1 / \varepsilon_{\mathrm{r}}\right)$ & Solvent treated as a continuum and the ion as a sphere & 3D.12a \\
\hline
\end{tabular}
\end{center}

\section*{TOPIC 3E Combining the First and Second Laws }
\subsection*{> Why do you need to know this material?}
The First and Second Laws of thermodynamics are both relevant to the behaviour of bulk matter, and the whole force of thermodynamics can be brought to bear on a problem by setting up a formulation that combines them.

\subsection*{> What is the key idea?}
The fact that infinitesimal changes in thermodynamic functions are exact differentials leads to relations between a variety of properties.

\subsection*{- What do you need to know already?}
You need to be aware of the definitions of the state functions $U$ (Topic 2A), $H$ (Topic 2B), $S$ (Topic 3A), and $A$ and $G$ (Topic 3D). The mathematical derivations in this Topic draw frequently on the properties of partial derivatives, which are described in The chemist's toolkit 9 in Topic 2A.

The First Law of thermodynamics may be written $\mathrm{d} U=\mathrm{d} q+$ $\mathrm{d} w$. For a reversible change in a closed system of constant composition, and in the absence of any additional (non-expansion) work, $\mathrm{d} w_{\mathrm{rev}}=-p \mathrm{~d} V$ and (from the definition of entropy) $\mathrm{d} q_{\mathrm{rev}}=$ $T \mathrm{~d} S$, where $p$ is the pressure of the system and $T$ its temperature. Therefore, for a reversible change in a closed system,


\begin{equation*}
\mathrm{d} U=T \mathrm{~d} S-p \mathrm{~d} V \quad \text { The fundamental equation } \tag{3E.1}
\end{equation*}


However, because $d U$ is an exact differential, its value is independent of path. Therefore, the same value of $\mathrm{d} U$ is obtained whether the change is brought about irreversibly or reversibly. Consequently, this equation applies to any change-reversible or irreversible-of a closed system that does no additional (nonexpansion) work. This combination of the First and Second Laws is called the fundamental equation.

The fact that the fundamental equation applies to both reversible and irreversible changes may be puzzling at first sight. The reason is that only in the case of a reversible change may $T \mathrm{~d} S$ be identified with $\mathrm{d} q$ and $-p \mathrm{~d} V$ with $\mathrm{d} w$. When the change is irreversible, $T \mathrm{~d} S>\mathrm{d} q$ (the Clausius inequality) and $-p \mathrm{~d} V>$ $\mathrm{d} w$. The sum of $\mathrm{d} w$ and $\mathrm{d} q$ remains equal to the sum of $T \mathrm{~d} S$ and $-p \mathrm{~d} V$, provided the composition is constant.

\subsection*{3E.1 Properties of the internal energy}
Equation 3E. 1 shows that the internal energy of a closed system changes in a simple way when either $S$ or $V$ is changed ( $\mathrm{d} U \propto \mathrm{~d} S$ and $\mathrm{d} U \propto \mathrm{~d} V$ ). These simple proportionalities suggest that $U$ is best regarded as a function of $S$ and $V$. It could be regarded as a function of other variables, such as $S$ and $p$ or $T$ and $V$, because they are all interrelated; but the simplicity of the fundamental equation suggests that $U(S, V)$ is the best choice.

The mathematical consequence of $U$ being a function of $S$ and $V$ is that an infinitesimal change $\mathrm{d} U$ can be expressed in terms of changes $\mathrm{d} S$ and $\mathrm{d} V$ by


\begin{equation*}
\mathrm{d} U=\left(\frac{\partial U}{\partial S}\right)_{V} \mathrm{~d} S+\left(\frac{\partial U}{\partial V}\right)_{S} \mathrm{~d} V \tag{3E.2}
\end{equation*}


The two partial derivatives (see The chemist's toolkit 9 in Topic 2A) are the slopes of the plots of $U$ against $S$ at constant $V$, and $U$ against $V$ at constant $S$. When this expression is compared term-by-term to the thermodynamic relation, eqn 3E.1, it follows that for systems of constant composition,


\begin{equation*}
\left(\frac{\partial U}{\partial S}\right)_{V}=T \quad\left(\frac{\partial U}{\partial V}\right)_{S}=-p \tag{3E.3}
\end{equation*}


The first of these two equations is a purely thermodynamic definition of temperature as the ratio of the changes in the internal energy (a First-Law concept) and entropy (a Second-Law concept) of a constant-volume, closed, constant-composition system. Relations between the properties of a system are starting to emerge.

\subsection*{(a) The Maxwell relations}
An infinitesimal change in a function $f(x, y)$ can be written $\mathrm{d} f=g \mathrm{~d} x+h \mathrm{~d} y$ where $g$ and $h$ may be functions of $x$ and $y$. The mathematical criterion for $\mathrm{d} f$ being an exact differential (in the sense that its integral is independent of path) is that


\begin{equation*}
\left(\frac{\partial g}{\partial y}\right)_{x}=\left(\frac{\partial h}{\partial x}\right)_{y} \tag{3E.4}
\end{equation*}


This criterion is derived in The chemist's toolkit 10. Because the fundamental equation, eqn 3E.1, is an expression for

\subsection*{The chemist's toolkit 10 Exact differentials}
Suppose that $\mathrm{d} f$ can be expressed in the following way:

$$
\mathrm{d} f=g(x, y) \mathrm{d} x+h(x, y) \mathrm{d} y
$$

Is $\mathrm{d} f$ is an exact differential? If it is exact, then it can be expressed in the form

$$
\mathrm{d} f=\left(\frac{\partial f}{\partial x}\right)_{y} \mathrm{~d} x+\left(\frac{\partial f}{\partial y}\right)_{x} \mathrm{~d} y
$$

Comparing these two expressions gives

$$
\left(\frac{\partial f}{\partial x}\right)_{y}=g(x, y) \quad\left(\frac{\partial f}{\partial y}\right)_{x}=h(x, y)
$$

It is a property of partial derivatives that successive derivatives may be taken in any order:

$$
\left(\frac{\partial}{\partial y}\left(\frac{\partial f}{\partial x}\right)_{y}\right)_{x}=\left(\frac{\partial}{\partial x}\left(\frac{\partial f}{\partial y}\right)_{x}\right)_{y}
$$

an exact differential, the functions multiplying $\mathrm{d} S$ and $\mathrm{d} V$ (namely $T$ and $-p$ ) must pass this test. Therefore, it must be the case that


\begin{equation*}
\left(\frac{\partial T}{\partial V}\right)_{S}=-\left(\frac{\partial p}{\partial S}\right)_{V} \tag{3E.5}
\end{equation*}


A Maxwell relation\\
A relation has been generated between quantities which, at first sight, would not seem to be related.

Equation 3E. 5 is an example of a Maxwell relation. However, apart from being unexpected, it does not look particularly interesting. Nevertheless, it does suggest that there might be other similar relations that are more useful. Indeed, the fact that $H, G$, and $A$ are all state functions can be used to derive three more Maxwell relations. The argument to obtain them runs in the same way in each case: because $H, G$, and $A$ are state functions, the expressions for $\mathrm{d} H, \mathrm{~d} G$, and $\mathrm{d} A$ satisfy relations like eqn 3E.4. All four relations are listed in Table 3E.1.

Table 3E. 1 The Maxwell relations

\begin{center}
\begin{tabular}{lll}
\hline
State function & Exact differential & Maxwell relation \\
\hline
$U$ & $\mathrm{~d} U=T \mathrm{~d} S-p \mathrm{~d} V$ & $\left(\frac{\partial T}{\partial V}\right)_{S}=-\left(\frac{\partial p}{\partial S}\right)_{V}$ \\
$H$ & $\mathrm{~d} H=T \mathrm{~d} S+V \mathrm{~d} p$ & $\left(\frac{\partial T}{\partial p}\right)_{S}=\left(\frac{\partial V}{\partial S}\right)_{p}$ \\
$A$ & $\mathrm{~d} A=-p \mathrm{~d} V-S \mathrm{~d} T$ & $\left(\frac{\partial p}{\partial T}\right)_{V}=\left(\frac{\partial S}{\partial V}\right)_{T}$ \\
$G$ & $\mathrm{~d} G=V \mathrm{~d} p-S \mathrm{~d} T$ & $\left(\frac{\partial V}{\partial T}\right)_{p}=-\left(\frac{\partial S}{\partial p}\right)_{T}$ \\
\end{tabular}
\end{center}

Taking the partial derivative with respect to $x$ of the first equation, and with respect to $y$ of the second gives

$$
\left(\frac{\partial}{\partial y}\left(\frac{\partial f}{\partial x}\right)_{y}\right)_{x}=\left(\frac{\partial g(x, y)}{\partial y}\right)_{x}\left(\frac{\partial}{\partial x}\left(\frac{\partial f}{\partial y}\right)_{x}\right)_{y}=\left(\frac{\partial h(x, y)}{\partial x}\right)_{y}
$$

By the property of partial derivatives these two successive derivatives of $f$ with respect to $x$ and $y$ must be the same, hence

$$
\left(\frac{\partial g(x, y)}{\partial y}\right)_{x}=\left(\frac{\partial h(x, y)}{\partial y}\right)_{y}
$$

If this equality is satisfied, then $\mathrm{d} f=g(x, y) \mathrm{d} x+h(x, y) \mathrm{d} y$ is an exact differential. Conversely, if it is known from other arguments that $\mathrm{d} f$ is exact, then this relation between the partial derivatives follows.

\subsection*{Example 3E. 1}
\subsection*{Using the Maxwell relations}
Use the Maxwell relations in Table 3E. 1 to show that the entropy of a perfect gas is linearly dependent on $\ln V$, that is, $S=a+b \ln V$.

Collect your thoughts The natural place to start, given that you are invited to use the Maxwell relations, is to consider the relation for $(\partial S / \partial V)_{T}$, as that differential coefficient shows how the entropy varies with volume at constant temperature. Be alert for an opportunity to use the perfect gas equation of state.

The solution From Table 3E.1,

$$
\left(\frac{\partial S}{\partial V}\right)_{T}=\left(\frac{\partial p}{\partial T}\right)_{V}
$$

Now use the perfect gas equation of state, $p V=n R T$, to write $p=n R T / V$ :

$$
\left(\frac{\partial p}{\partial T}\right)_{V}=\left(\frac{\partial(n R T / V)}{\partial T}\right)_{V}=\frac{n R}{V}
$$

At this point, write

$$
\left(\frac{\partial S}{\partial V}\right)_{T}=\frac{n R}{V}
$$

and therefore, at constant temperature,

$$
\int \mathrm{d} S=n R \int \frac{\mathrm{~d} V}{V}=n R \ln V+\text { constant }
$$

The integral on the left is $S+$ constant, which completes the demonstration.

Self-test 3E. 1 How does the entropy depend on the volume of a van der Waals gas? Suggest a reason.\\
\includegraphics[max width=\textwidth, center]{2025_02_27_d4b95d9718f0b9e2e6b9g-137}

\subsection*{(b) The variation of internal energy with volume}
The internal pressure, $\pi_{T}$ (introduced in Topic 2D), is defined as $\pi_{T}=(\partial U / \partial V)_{T}$ and represents how the internal energy changes as the volume of a system is changed isothermally; it plays a central role in the manipulation of the First Law. By using a Maxwell relation, $\pi_{T}$ can be expressed as a function of pressure and temperature.

\subsection*{How is that done? 3E. 1 Deriving a thermodynamic}
 equation of stateTo construct the partial differential $(\partial U / \partial V)_{T}$ you need to start from eqn 3E.2, divide both sides by $\mathrm{d} V$, and impose the constraint of constant temperature:

$$
\overbrace{\left(\frac{\partial U}{\partial V}\right)_{T}}^{\pi_{T}}=\overbrace{\left(\frac{\partial U}{\partial S}\right)_{V}}^{T}\left(\frac{\partial S}{\partial V}\right)_{T}+\overbrace{\left(\frac{\partial U}{\partial V}\right)_{S}}^{-p}
$$

Next, introduce the two relations in eqn 3E. 3 (as indicated by the annotations) and the definition of $\pi_{T}$ to obtain

$$
\pi_{T}=T\left(\frac{\partial S}{\partial V}\right)_{T}-p
$$

The third Maxwell relation in Table 3E. 1 turns $(\partial S / \partial V)_{T}$ into $(\partial p / \partial T)_{V}$, to give\\
$\left.\rightarrow \pi_{T}=T\left(\frac{\partial p}{\partial T}\right)_{V}-p \right\rvert\, \begin{array}{r}\text { A thermodynamic equation of state }\end{array}$\\
Equation 3E.6a is called a thermodynamic equation of state because, when written in the form


\begin{equation*}
p=T\left(\frac{\partial p}{\partial T}\right)_{V}-\pi_{T} \tag{3E.6b}
\end{equation*}


it is an expression for pressure in terms of a variety of thermodynamic properties of the system.

\subsection*{Example 3E. 2}
\subsection*{Deriving a thermodynamic relation}
Show thermodynamically that $\pi_{T}=0$ for a perfect gas, and compute its value for a van der Waals gas.

Collect your thoughts Proving a result 'thermodynamically' means basing it entirely on general thermodynamic relations and equations of state, without drawing on molecular arguments (such as the existence of intermolecular forces). You know that for a perfect gas, $p=n R T / V$, so this relation should be used in eqn 3E.6. Similarly, the van der Waals equation is given in Table 1C.4, and for the second part of the question it should be used in eqn 3E.6.

The solution For a perfect gas write

$$
\left(\frac{\partial p}{\partial T}\right)_{V}=\left(\frac{\partial n R T / V}{\partial T}\right)_{V}=\frac{n R}{V}
$$

Then, eqn 3E. 6 becomes

$$
\pi_{T}=T\left(\frac{\partial p}{\partial T}\right)_{V}-p=\overbrace{\left(\frac{n R T}{V}\right)}^{p}-p=0
$$

The equation of state of a van der Waals gas is

$$
p=\frac{n R T}{V-n b}-a \frac{n^{2}}{V^{2}}
$$

Because $a$ and $b$ are independent of temperature,

$$
\left(\frac{\partial p}{\partial T}\right)_{V}=\left(\frac{\partial n R T /(V-n b)}{\partial T}\right)_{V}=\frac{n R}{V-n b}
$$

Therefore, from eqn 3E.6,

$$
\pi_{T}=\frac{n R T}{V-n b}-p=\frac{n R T}{V-n b}-\overbrace{\left(\frac{n R T}{V-n b}-a \frac{n^{2}}{V^{2}}\right)}^{p}=a \frac{n^{2}}{V^{2}}
$$

Comment. This result for $\pi_{T}$ implies that the internal energy of a van der Waals gas increases when it expands isothermally, that is, $(\partial U / \partial V)_{T}>0$, and that the increase is related to the parameter $a$, which models the attractive interactions between the particles. A larger molar volume, corresponding to a greater average separation between molecules, implies weaker mean intermolecular attractions, so the total energy is greater.\\
Self-test $3 E .2$ Calculate $\pi_{T}$ for a gas that obeys the virial equation of state (Table 1 C .4 ), retaining only the term in $B$.

\subsection*{3E.2 Properties of the Gibbs energy}
The same arguments that were used for $U$ can also be used for the Gibbs energy, $G=H-T S$. They lead to expressions showing how $G$ varies with pressure and temperature and which are important for discussing phase transitions and chemical reactions.

\subsection*{(a) General considerations}
When the system undergoes a change of state, $G$ may change because $H, T$, and $S$ all change:\\
$\mathrm{d} G=\mathrm{d} H-\mathrm{d}(T S)=\mathrm{d} H-T \mathrm{~d} S-S \mathrm{~d} T$

Because $H=U+p V$,

$$
\mathrm{d} H=\mathrm{d} U+\mathrm{d}(p V)=\mathrm{d} U+p \mathrm{~d} V+V \mathrm{~d} p
$$

and therefore

$$
\mathrm{d} G=\mathrm{d} U+p \mathrm{~d} V+V \mathrm{~d} p-T \mathrm{~d} S-S \mathrm{~d} T
$$

For a closed system doing no non-expansion work, $\mathrm{d} U$ can be replaced by the fundamental equation $\mathrm{d} U=T \mathrm{~d} S-p \mathrm{~d} V$ to give

$$
\mathrm{d} G=T \mathrm{~d} S-p \mathrm{~d} V+p \mathrm{~d} V+V \mathrm{~d} p-T \mathrm{~d} S-S \mathrm{~d} T
$$

Four terms now cancel on the right, and so for a closed system in the absence of non-expansion work and at constant composition

\[
\begin{array}{ll}
\mathrm{d} G=V \mathrm{~d} p-S \mathrm{~d} T & \begin{array}{l}
\text { The fundamental equation of } \\
\text { chemical thermodynamics }
\end{array} \tag{3E.7}
\end{array}
\]

This expression, which shows that a change in $G$ is proportional to a change in $p$ or $T$, suggests that $G$ may be best regarded as a function of $p$ and $T$. It may be regarded as the fundamental equation of chemical thermodynamics as it is so central to the application of thermodynamics to chemistry. It also suggests that $G$ is an important quantity in chemistry because the pressure and temperature are usually the variables that can be controlled. In other words, $G$ carries around the combined consequences of the First and Second Laws in a way that makes it particularly suitable for chemical applications.

The same argument that led to eqn 3E.3, when applied to the exact differential $\mathrm{d} G=V \mathrm{~d} p-S \mathrm{~d} T$, now gives

\[
\left(\frac{\partial G}{\partial T}\right)_{p}=-S\left(\frac{\partial G}{\partial p}\right)_{T}=V \quad \begin{align*}
& \text { The variation of }  \tag{3E.8}\\
& G \text { with } T \text { and } p
\end{align*}
\]

These relations show how the Gibbs energy varies with temperature and pressure (Fig. 3E.1).\\
\includegraphics[max width=\textwidth, center]{2025_02_27_d4b95d9718f0b9e2e6b9g-139(1)}

Figure 3E. 1 The variation of the Gibbs energy of a system with (a) temperature at constant pressure and (b) pressure at constant temperature. The slope of the former is equal to the negative of the entropy of the system and that of the latter is equal to the volume.\\
\includegraphics[max width=\textwidth, center]{2025_02_27_d4b95d9718f0b9e2e6b9g-139}

Figure 3E. 2 The variation of the Gibbs energy with the temperature is determined by the entropy. Because the entropy of the gaseous phase of a substance is greater than that of the liquid phase, and the entropy of the solid phase is smallest, the Gibbs energy changes most steeply for the gas phase, followed by the liquid phase, and then the solid phase of the substance.

The first implies that:

\begin{itemize}
  \item Because $S>0$ for all substances, $G$ always decreases when the temperature is raised (at constant pressure and composition).
  \item Because $(\partial G / \partial T)_{p}$ becomes more negative as $S$ increases, $G$ decreases most sharply with increasing temperature when the entropy of the system is large.
\end{itemize}

Therefore, the Gibbs energy of the gaseous phase of a substance, which has a high molar entropy, is more sensitive to temperature than its liquid and solid phases (Fig. 3E.2).

Similarly, the second relation implies that:

\begin{itemize}
  \item Because $V>0$ for all substances, $G$ always increases when the pressure of the system is increased (at constant temperature and composition).
  \item Because $(\partial G / \partial p)_{T}$ increases with $V, G$ is more sensitive to pressure when the volume of the system is large.\\
Because the molar volume of the gaseous phase of a substance is greater than that of its condensed phases, the molar Gibbs energy of a gas is more sensitive to pressure than its liquid and solid phases (Fig. 3E.3).
\end{itemize}

\subsection*{Brief illustration 3E. 1}
The mass density of liquid water is $0.9970 \mathrm{~g} \mathrm{~cm}^{-3}$ at 298 K . It follows that when the pressure is increased by 0.1 bar (at constant temperature), the molar Gibbs energy changes by

$$
\begin{aligned}
\Delta G_{\mathrm{m}} & \approx\left(\frac{\partial G_{\mathrm{m}}}{\partial p}\right)_{T} \Delta p=V_{\mathrm{m}} \Delta p=\overbrace{\frac{18.0 \mathrm{~g} \mathrm{~mol}^{-1}}{0.9970 \times 10^{6} \mathrm{gm}^{-3}}}^{V_{\mathrm{m}}} \times\left(0.1 \times 10^{5} \mathrm{Nm}^{-2}\right) \\
& =+0.18 \mathrm{Jmol}^{-1}
\end{aligned}
$$

\begin{center}
\includegraphics[max width=\textwidth]{2025_02_27_d4b95d9718f0b9e2e6b9g-140}
\end{center}

Figure 3E. 3 The variation of the Gibbs energy with the pressure is determined by the volume of the sample. Because the volume of the gaseous phase of a substance is greater than that of the same amount of liquid phase, and the volume of the solid phase is smallest (for most substances), the Gibbs energy changes most steeply for the gas phase, followed by the liquid phase, and then the solid phase of the substance. Because the molar volumes of the solid and liquid phases of a substance are similar, their molar Gibbs energies vary by similar amounts as the pressure is changed.

\subsection*{(b) The variation of the Gibbs energy with temperature}
Because the equilibrium composition of a system depends on the Gibbs energy, in order to discuss the response of the composition to temperature it is necessary to know how $G$ varies with temperature.

The first relation in eqn 3E.8, $(\partial G / \partial T)_{p}=-S$, is the starting point for this discussion. Although it expresses the variation of $G$ in terms of the entropy, it can be expressed in terms of the enthalpy by using the definition of $G$ to write $S=(H-G) / T$. Then


\begin{equation*}
\left(\frac{\partial G}{\partial T}\right)_{p}=\frac{G-H}{T} \tag{3E.9}
\end{equation*}


In Topic 6A it is shown that the equilibrium constant of a reaction is related to $G / T$ rather than to $G$ itself. With this application in mind, eqn 3E. 9 can be developed to show how $G / T$ varies with temperature.

\subsection*{How is that done? 3E. 2 Deriving an expression for the temperature variation of $G / T$}
First, note that

$$
\begin{aligned}
\left(\frac{\partial G / T}{\partial T}\right)_{p} & =\frac{1}{T}\left(\frac{\partial G}{\partial T}\right)_{p}+G \underbrace{\frac{\mathrm{~d}(f g) / d \mathrm{~d} x=f(\mathrm{~d} / \mathrm{d} / \mathrm{d})+g(\mathrm{~d} / \mathrm{d} / \mathrm{d} \mathrm{~d})}{\mathrm{d} T}}=\frac{1}{T}\left(\frac{\partial G}{\partial T}\right)_{p}-\frac{G}{T^{2}} \\
& =\frac{1}{T}\left\{\left(\frac{\partial G}{\partial T}\right)_{p}-\frac{G}{T}\right\}
\end{aligned}
$$

Now replace the term $(\partial G / \partial T)_{p}$ on the right by eqn 3E. 9

$$
\left(\frac{\partial G / T}{\partial T}\right)_{p}=\frac{1}{T}\left\{\left(\frac{\partial G}{\partial T}\right)_{p}-\frac{G}{T}\right\}=\frac{1}{T}\left\{\frac{G-H}{T}-\frac{G}{T}\right\}=\frac{1}{T}\left\{\frac{-H}{T}\right\}
$$

from which follows the Gibbs-Helmholtz equation

$$
\left.-\left(\frac{\partial G / T}{\partial T}\right)_{p}=-\frac{H}{T^{2}} \right\rvert\, \begin{aligned}
& \text { (3E.10) } \\
&
\end{aligned}
$$

The Gibbs-Helmholtz equation is most useful when it is applied to changes, including changes of physical state, and chemical reactions at constant pressure. Then, because $\Delta G=$ $G_{f}-G_{i}$ for the change of Gibbs energy between the final and initial states, and because the equation applies to both $G_{f}$ and $G_{i}$,


\begin{equation*}
\left(\frac{\partial \Delta G / T}{\partial T}\right)_{p}=-\frac{\Delta H}{T^{2}} \tag{3E.11}
\end{equation*}


This equation shows that if the change in enthalpy of a system that is undergoing some kind of transformation (such as vaporization or reaction) is known, then how the corresponding change in Gibbs energy varies with temperature is also known. This turns out to be a crucial piece of information in chemistry.

\subsection*{(c) The variation of the Gibbs energy with pressure}
To find the Gibbs energy at one pressure in terms of its value at another pressure, the temperature being constant, set $\mathrm{d} T=0$ in eqn 3E.7, which gives $\mathrm{d} G=V \mathrm{~d} p$, and integrate:


\begin{equation*}
G\left(p_{\mathrm{f}}\right)=G\left(p_{\mathrm{i}}\right)+\int_{p_{\mathrm{i}}}^{p_{\mathrm{f}}} V \mathrm{~d} p \tag{3E.12a}
\end{equation*}


For molar quantities,


\begin{equation*}
G_{\mathrm{m}}\left(p_{\mathrm{f}}\right)=G_{\mathrm{m}}\left(p_{\mathrm{i}}\right)+\int_{p_{\mathrm{i}}}^{p_{\mathrm{f}}} V_{\mathrm{m}} \mathrm{~d} p \tag{3E.12b}
\end{equation*}


This expression is applicable to any phase of matter, but it is necessary to know how the molar volume, $V_{\mathrm{m}}$, depends on the pressure before the integral can be evaluated.

The molar volume of a condensed phase changes only slightly as the pressure changes, so in this case $V_{\mathrm{m}}$ can be treated as constant and taken outside the integral:

$$
G_{\mathrm{m}}\left(p_{\mathrm{f}}\right)=G_{\mathrm{m}}\left(p_{\mathrm{i}}\right)+V_{\mathrm{m}} \int_{p_{\mathrm{i}}}^{p_{\mathrm{f}}} \mathrm{~d} p
$$

That is,

$$
G_{\mathrm{m}}\left(p_{\mathrm{f}}\right)=G_{\mathrm{m}}\left(p_{\mathrm{i}}\right)+\left(p_{\mathrm{f}}-p_{\mathrm{i}}\right) V_{\mathrm{m}} \quad \begin{aligned}
& \text { Molar Gibbs energy } \\
& \text { [incompressible } \\
& \text { substance] }
\end{aligned}
$$

\begin{center}
\includegraphics[max width=\textwidth]{2025_02_27_d4b95d9718f0b9e2e6b9g-141(1)}
\end{center}

Figure 3E. 4 At constant temperature, the difference in Gibbs energy of a solid or liquid between two pressures is equal to the rectangular area shown. The variation of volume with pressure has been assumed to be negligible.

The origin of the term $\left(p_{\mathrm{f}}-p_{\mathrm{i}}\right) V_{\mathrm{m}}$ is illustrated graphically in Fig. 3E.4. Under normal laboratory conditions $\left(p_{\mathrm{f}}-p_{\mathrm{i}}\right) V_{\mathrm{m}}$ is very small and may be neglected. Hence, the Gibbs energies of solids and liquids are largely independent of pressure. However, in geophysical problems, because pressures in the Earth's interior are huge, their effect on the Gibbs energy cannot be ignored. If the pressures are so great that there are substantial volume changes over the range of integration, then the complete expression, eqn 3E.12, must be used.

\subsection*{Example 3E. 3 \\
 Evaluating the pressure dependence of a Gibbs energy of transition}
Suppose that for a certain phase transition of a solid $\Delta_{\text {trs }} V=$ $+1.0 \mathrm{~cm}^{3} \mathrm{~mol}^{-1}$ independent of pressure. By how much does that Gibbs energy of transition change when the pressure is increased from $1.0 \operatorname{bar}\left(1.0 \times 10^{5} \mathrm{~Pa}\right)$ to $3.0 \mathrm{Mbar}\left(3.0 \times 10^{11} \mathrm{~Pa}\right)$ ?

Collect your thoughts You need to start with eqn 3E.12b to obtain expressions for the Gibbs energy of each of the phases 1 and 2 of the solid

$$
\begin{aligned}
& G_{\mathrm{m}, 1}\left(p_{\mathrm{f}}\right)=G_{\mathrm{m}, 1}\left(p_{\mathrm{i}}\right)+\int_{p_{\mathrm{i}}}^{p_{\mathrm{f}}} V_{\mathrm{m}, 1} \mathrm{~d} p \\
& G_{\mathrm{m}, 2}\left(p_{\mathrm{f}}\right)=G_{\mathrm{m}, 2}\left(p_{\mathrm{i}}\right)+\int_{p_{\mathrm{i}}}^{p_{\mathrm{f}}} V_{\mathrm{m}, 2} \mathrm{~d} p
\end{aligned}
$$

Then, to obtain $\Delta_{\mathrm{trs}} G=G_{\mathrm{m}, 2}-G_{\mathrm{m}, 1}$ subtract the second expression from the first, noting that $V_{\mathrm{m}, 2}-V_{\mathrm{m}, 1}=\Delta_{\text {trs }} V$ :

$$
\Delta_{\mathrm{trS}} G_{\mathrm{m}}\left(p_{\mathrm{f}}\right)=\Delta_{\mathrm{trs}} G_{\mathrm{m}}\left(p_{\mathrm{i}}\right)+\int_{p_{\mathrm{i}}}^{p_{\mathrm{f}}} \Delta_{\mathrm{trs}} V \mathrm{~d} p
$$

Use the data to complete the calculation.

The solution Because $\Delta_{\text {trs }} V_{\mathrm{m}}$ is independent of pressure,

$$
\Delta_{\mathrm{trs}} G_{\mathrm{m}}\left(p_{\mathrm{f}}\right)=\Delta_{\mathrm{trs}} G_{\mathrm{m}}\left(p_{\mathrm{i}}\right)+\overbrace{\Delta_{\mathrm{trs}} V}^{\text {constant }} \int_{p_{\mathrm{i}}}^{p_{\mathrm{f}}} \mathrm{~d} p=\Delta_{\mathrm{trs}} G_{\mathrm{m}}\left(p_{\mathrm{i}}\right)+\Delta_{\mathrm{trs}} V\left(p_{\mathrm{f}}-p_{\mathrm{i}}\right)
$$

Inserting the data and using $1 \mathrm{Pam}^{3}=1 \mathrm{~J}$ gives

$$
\begin{aligned}
\Delta_{\mathrm{trs}} G(3 \mathrm{Mbar})= & \Delta_{\mathrm{tr}} G(1 \mathrm{bar})+\left(1.0 \times 10^{-6} \mathrm{~m}^{3} \mathrm{~mol}^{-1}\right) \\
& \times\left(3.0 \times 10^{11} \mathrm{~Pa}-1.0 \times 10^{5} \mathrm{~Pa}\right) \\
= & \Delta_{\mathrm{trs}} G(1 \mathrm{bar})+3.0 \times 10^{2} \mathrm{~kJ} \mathrm{~mol}^{-1}
\end{aligned}
$$

Self-test $3 E .3$ Calculate the change in $G_{m}$ for ice at $-10^{\circ} \mathrm{C}$, with density $917 \mathrm{~kg} \mathrm{~m}^{-3}$, when the pressure is increased from 1.0 bar to 2.0 bar.\\
\includegraphics[max width=\textwidth, center]{2025_02_27_d4b95d9718f0b9e2e6b9g-141}

The molar volumes of gases are large, so the Gibbs energy of a gas depends strongly on the pressure. Furthermore, because the volume also varies markedly with the pressure, the volume cannot be treated as a constant in the integral in eqn 3E.12b (Fig. 3E.5).

For a perfect gas, substitute $V_{\mathrm{m}}=R T / p$ into the integral, note that $T$ is constant, and find


\begin{equation*}
G_{\mathrm{m}}\left(p_{\mathrm{f}}\right)=G_{\mathrm{m}}\left(p_{\mathrm{i}}\right)+R T \overbrace{\int_{p_{\mathrm{i}}}^{p_{\mathrm{i}}} \frac{1}{p} \mathrm{~d} p}^{\text {Integral A.2 }}=G_{\mathrm{m}}\left(p_{\mathrm{i}}\right)+R T \ln \frac{p_{\mathrm{f}}}{p_{\mathrm{i}}} \tag{3E.14}
\end{equation*}


This expression shows that when the pressure is increased tenfold at room temperature, the molar Gibbs energy increases by $R T \ln 10 \approx 6 \mathrm{~kJ} \mathrm{~mol}^{-1}$. It also follows from this equation that if $p_{\mathrm{i}}=p^{\ominus}$ (the standard pressure of 1 bar ), then the molar Gibbs energy of a perfect gas at a pressure $p\left(\right.$ set $\left.p_{\mathrm{f}}=p\right)$ is related to its standard value by

\[
G_{\mathrm{m}}(p)=G_{\mathrm{m}}^{\ominus}+R T \ln \frac{p}{p^{\ominus}} \quad \begin{align*}
& \text { Molar Gibbs energy }  \tag{3E.15}\\
& {[\text { perfect gas, constant } T]}
\end{align*}
\]

\begin{center}
\includegraphics[max width=\textwidth]{2025_02_27_d4b95d9718f0b9e2e6b9g-141(2)}
\end{center}

Figure 3E. 5 At constant temperature, the change in Gibbs energy for a perfect gas between two pressures is equal to the area shown below the perfect-gas isotherm.

\subsection*{Brief illustration 3E. 2}
When the pressure is increased isothermally on water vapour (treated as a perfect gas) from 1.0 bar to 2.0 bar at 298 K , then according to eqn 3 E .15

$$
\begin{aligned}
G_{\mathrm{m}}(2.0 \mathrm{bar})= & G_{\mathrm{m}}^{\ominus}(1.0 \mathrm{bar})+\left(8.3145 \mathrm{JK}^{-1} \mathrm{~mol}^{-1}\right) \\
& \times(298 \mathrm{~K}) \times \ln \left(\frac{2.0 \mathrm{bar}}{1.0 \mathrm{bar}}\right) \\
= & G_{\mathrm{m}}^{\ominus}(1.0 \mathrm{bar})+1.7 \mathrm{kJmol}^{-1}
\end{aligned}
$$

Note that whereas the change in molar Gibbs energy for a condensed phase is a few joules per mole, for a gas the change is of the order of kilojoules per mole.

The logarithmic dependence of the molar Gibbs energy on the pressure predicted by eqn 3E. 15 is illustrated in Fig. 3E.6. This very important expression applies to perfect gases (which is usually a good approximation).\\
\includegraphics[max width=\textwidth, center]{2025_02_27_d4b95d9718f0b9e2e6b9g-142}

Figure 3E. 6 At constant temperature, the molar Gibbs energy of a perfect gas varies as $\ln p$, and the standard state is reached at $p^{\ominus}$. Note that, as $p \rightarrow 0$, the molar Gibbs energy becomes negatively infinite.

\subsection*{Checklist of concepts}
$\square$ 1. The fundamental equation, a combination of the First and Second Laws, is an expression for the change in internal energy that accompanies changes in the volume and entropy of a system.\\
$\square$ 2. Relations between thermodynamic properties are generated by combining thermodynamic and mathematical expressions for changes in their values.\\
$\square$ 3. The Maxwell relations are a series of relations between partial derivatives of thermodynamic properties based on criteria for changes in the properties being exact differentials.\\
4. The Maxwell relations are used to derive the thermodynamic equation of state and to determine how the internal energy of a substance varies with volume.\\
5. The variation of the Gibbs energy of a system suggests that it is best regarded as a function of pressure and temperature.\\
$\square$ 6. The Gibbs energy of a substance decreases with temperature and increases with pressure.\\
$\square$ 7. The variation of Gibbs energy with temperature is related to the enthalpy by the Gibbs-Helmholtz equation.\\
8. The Gibbs energies of solids and liquids are almost independent of pressure; those of gases vary linearly with the logarithm of the pressure.

\subsection*{Checklist of equations}
\begin{center}
\begin{tabular}{llll}
\hline
Property & Equation & Comment & Equation number \\
\hline
Fundamental equation & $\mathrm{d} U=T \mathrm{~d} S-p \mathrm{~d} V$ & No additional work & 3E.1 \\
Fundamental equation of chemical thermodynamics & $\mathrm{d} G=V \mathrm{~d} p-S \mathrm{~d} T$ & No additional work & 3E.7 \\
Variation of $G$ & $(\partial G / \partial p)_{T}=V$ and $(\partial G / \partial T)_{p}=-S$ & Composition constant & 3E.8 \\
Gibbs-Helmholtz equation & $(\partial(G / T) / \partial T)_{p}=-H / T^{2}$ & Composition constant & 3E.10 \\
Pressure dependence of $G_{\mathrm{m}}$ & $G_{\mathrm{m}}\left(p_{\mathrm{f}}\right)=G_{\mathrm{m}}\left(p_{\mathrm{i}}\right)+V_{\mathrm{m}}\left(p_{\mathrm{f}}-p_{\mathrm{i}}\right)$ & Incompressible substance & 3E.13 \\
 & $G_{\mathrm{m}}\left(p_{\mathrm{f}}\right)=G_{\mathrm{m}}\left(p_{\mathrm{i}}\right)+R T \ln \left(p_{\mathrm{f}} / p_{\mathrm{i}}\right)$ & Perfect gas, isothermal & 3E.14 \\
 & $G_{\mathrm{m}}(p)=G_{\mathrm{m}}^{\ominus}+R T \ln \left(p / p^{\ominus}\right)$ & Perfect gas, isothermal & 3E.15 \\
\end{tabular}
\end{center}

\chapter{FOCUS 4 Physical transformations of pure substances}
Vaporization, melting (fusion), and the conversion of graphite to diamond are all examples of changes of phase without change of chemical composition. The discussion of the phase transitions of pure substances is among the simplest applications of thermodynamics to chemistry, and is guided by the principle that, at constant temperature and pressure, the tendency of systems is to minimize their Gibbs energy.

\subsection*{4A Phase diagrams of pure substances}
One type of phase diagram is a map of the pressures and temperatures at which each phase of a substance is the most stable. The thermodynamic criterion for phase stability leads to a very general result, the 'phase rule', which summarizes the constraints on the equilibria between phases. In preparation for later chapters, this rule is expressed in a general way that can be applied to systems of more than one component. This Topic also introduces the 'chemical potential', a property that is at the centre of discussions of mixtures and chemical reactions. The Topic then describes the interpretation of the phase diagrams of a representative selection of substances.

\footnotetext{4A. 1 The stabilities of phases; 4A. 2 Phase boundaries; 4A. 3 Three
} representative phase diagrams

\subsection*{4B Thermodynamic aspects of phase transitions}
This Topic considers the factors that determine the positions and shapes of the phase boundaries. The expressions derived show how the vapour pressure of a substance varies with temperature and how the melting point varies with pressure.\\
4B. 1 The dependence of stability on the conditions; 4B. 2 The location of phase boundaries

\subsection*{Web resource What is an application of this material?}
The properties of carbon dioxide in its supercritical fluid phase can form the basis for novel and useful chemical separation methods, and have considerable promise for the synthetic procedures adopted in 'green' chemistry. Its properties and applications are discussed in Impact 6 on the website of this book.

\section*{TOPIC 4A Phase diagrams of pure substances }
\subsection*{Why do you need to know this material?}
Phase diagrams summarize the behaviour of substances under different conditions, and identify which phase or phases are the most stable at a particular temperature and pressure. Such diagrams are important tools for understanding the behaviour of both pure substances and mixtures.

\subsection*{What is the key idea?}
A pure substance tends to adopt the phase with the lowest chemical potential.

\subsection*{What do you need to know already?}
This Topic builds on the fact that the Gibbs energy is a signpost of spontaneous change under conditions of constant temperature and pressure (Topic 3D).

One of the most succinct ways of presenting the physical changes of state that a substance can undergo is in terms of its 'phase diagram'. This material is also the basis of the discussion of mixtures in Focus 5.

\subsection*{4A. 1 The stabilities of phases}
Thermodynamics provides a powerful framework for describing and understanding the stabilities and transformations of phases, but the terminology must be used carefully. In particular, it is necessary to understand the terms 'phase', 'component', and 'degree of freedom'.

\subsection*{(a) The number of phases}
A phase is a form of matter that is uniform throughout in chemical composition and physical state. Thus, there are the solid, liquid, and gas phases of a substance, as well as various solid phases, such as the white and black allotropes of phosphorus, or the aragonite and calcite polymorphs of calcium carbonate.

A note on good practice An allotrope is a particular molecular form of an element (such as $\mathrm{O}_{2}$ and $\mathrm{O}_{3}$ ) and may be solid, liquid, or gas. A polymorph is one of a number of solid phases of an element or compound.

The number of phases in a system is denoted $P$. A gas, or a gaseous mixture, is a single phase $(P=1)$, a crystal of a substance is a single phase, and two fully mixed liquids form a single phase.

\subsection*{Brief illustration 4A. 1}
A solution of sodium chloride in water is a single phase $(P=1)$. Ice is a single phase even though it might be chipped into small fragments. A slurry of ice and water is a two-phase system $(P=2)$ even though it is difficult to map the physical boundaries between the phases. A system in which calcium carbonate undergoes the thermal decomposition $\mathrm{CaCO}_{3}(\mathrm{~s}) \rightarrow$ $\mathrm{CaO}(\mathrm{s})+\mathrm{CO}_{2}(\mathrm{~g})$ consists of two solid phases (one consisting of calcium carbonate and the other of calcium oxide) and one gaseous phase (consisting of carbon dioxide), so $P=3$.

Two metals form a two-phase system $(P=2)$ if they are immiscible, but a single-phase system $(P=1)$, an alloy, if they are miscible (and actually mixed). A solution of solid B in solid A-a homogeneous mixture of the two miscible substancesis uniform on a molecular scale. In a solution, atoms of $A$ are surrounded by atoms of A and B, and any sample cut from the sample, even microscopically small, is representative of the composition of the whole. It is therefore a single phase.

A dispersion is uniform on a macroscopic scale but not on a microscopic scale, because it consists of grains or droplets of one substance in a matrix of the other (Fig. 4A.1). A small sample could come entirely from one of the minute grains of pure A and would not be representative of the whole. A dispersion therefore consists of two phases.

\subsection*{(b) Phase transitions}
A phase transition, the spontaneous conversion of one phase into another phase, occurs at a characteristic transition temperature, $T_{\text {trs }}$, for a given pressure. At the transition temperature\\
\includegraphics[max width=\textwidth, center]{2025_02_27_d4b95d9718f0b9e2e6b9g-153(1)}\\
(a)\\
\includegraphics[max width=\textwidth, center]{2025_02_27_d4b95d9718f0b9e2e6b9g-153(2)}\\
(b)

Figure 4A. 1 The difference between (a) a single-phase solution, in which the composition is uniform on a molecular scale, and (b) a dispersion, in which microscopic regions of one component are embedded in a matrix of a second component.\\
the two phases are in equilibrium and the Gibbs energy of the system is a minimum at the prevailing pressure.

\subsection*{Brief illustration 4A. 2}
At 1 atm , ice is the stable phase of water below $0^{\circ} \mathrm{C}$, but above $0^{\circ} \mathrm{C}$ liquid water is more stable. This difference indicates that below $0^{\circ} \mathrm{C}$ the Gibbs energy decreases as liquid water changes into ice, but that above $0^{\circ} \mathrm{C}$ the Gibbs energy decreases as ice changes into liquid water. The numerical values of the Gibbs energies are considered in the next Brief illustration.

The detection of a phase transition is not always straightforward as there may be nothing to see, especially if the two phases are both solids. Thermal analysis, which takes advantage of the heat that is evolved or absorbed during a transition, can be used. Thus, if the phase transition is exothermic and the temperature of a sample is monitored as it cools, the presence of the transition can be recognized by a pause in the otherwise steady fall of the temperature (Fig. 4A.2). Similarly, if a sample is heated steadily and the transition is endothermic, there will\\
\includegraphics[max width=\textwidth, center]{2025_02_27_d4b95d9718f0b9e2e6b9g-153(3)}

Figure 4A. 2 A cooling curve at constant pressure. The flat section corresponds to the pause in the fall of temperature while an exothermic transition (freezing) occurs. This pause enables $T_{f}$ to be located even if the transition cannot be observed visually.\\
be a pause in the temperature rise at the transition temperature. Differential scanning calorimetry (Topic 2C) is also used to detect phase transitions, and X-ray diffraction (Topic 15B) is useful for detecting phase transitions in a solid, because the two phases will have different structures.

As always, it is important to distinguish between the thermodynamic description of a process and the rate at which the process occurs. A phase transition that is predicted by thermodynamics to be spontaneous might occur too slowly to be significant in practice. For instance, at normal temperatures and pressures the molar Gibbs energy of graphite is lower than that of diamond, so there is a thermodynamic tendency for diamond to change into graphite. However, for this transition to take place, the C atoms must change their locations, which, except at high temperatures, is an immeasurably slow process in a solid. The discussion of the rate of attainment of equilibrium is a kinetic problem and is outside the range of thermodynamics. In gases and liquids the mobilities of the molecules allow phase transitions to occur rapidly, but in solids thermodynamic instability may be frozen in. Thermodynamically unstable phases that persist because the transition is kinetically hindered are called metastable phases. Diamond is a metastable but persistent phase of carbon under normal conditions.

\subsection*{(c) Thermodynamic criteria of phase stability}
All the following considerations are based on the Gibbs energy of a substance, and in particular on its molar Gibbs energy, $G_{\mathrm{m}}$. In fact, this quantity plays such an important role in this Focus and elsewhere in the text that it is given a special name and symbol, the chemical potential, $\mu(\mathrm{mu})$. For a system that consists of a single substance, the 'molar Gibbs energy' and the 'chemical potential' are exactly the same: $\mu=G_{\mathrm{m}}$. In Topic 5A the chemical potential is given a broader significance and a more general definition. The name 'chemical potential' is also instructive: as the concept is developed it will become clear that $\mu$ is a measure of the potential that a substance has for undergoing change. In this Focus, and in Focus 5, it reflects the potential of a substance to undergo physical change. In Focus $6, \mu$ is the potential of a substance to undergo chemical change.

The discussion in this Topic is based on the following consequence of the Second Law (Fig. 4A.3):

At equilibrium, the chemical potential of a substance is the same in and throughout every phase present in the system.\\
\includegraphics[max width=\textwidth, center]{2025_02_27_d4b95d9718f0b9e2e6b9g-153}

To see the validity of this remark, consider a system in which the chemical potential of a substance is $\mu_{1}$ at one location and $\mu_{2}$ at another location. The locations may be in the same or in different phases. When an infinitesimal amount $\mathrm{d} n$ of the substance is transferred from one location to the other, the Gibbs energy of the system changes by $-\mu_{1} \mathrm{~d} n$ (i.e. $\mathrm{d} G=$ $-G_{\mathrm{m}, 1} \mathrm{~d} n$ ) when material is removed from location 1. It changes\\
\includegraphics[max width=\textwidth, center]{2025_02_27_d4b95d9718f0b9e2e6b9g-154(1)}

Figure 4A. 3 When two or more phases are in equilibrium, the chemical potential of a substance (and, in a mixture, a component) is the same in each phase, and is the same at all points in each phase.\\
by $+\mu_{2} \mathrm{~d} n$ (i.e. $\mathrm{d} G=G_{\mathrm{m}, 2} \mathrm{~d} n$ ) when that material is added to location 2. The overall change is therefore $\mathrm{d} G=\left(\mu_{2}-\mu_{1}\right) \mathrm{d} n$. If the chemical potential at location 1 is higher than that at location 2 , the transfer is accompanied by a decrease in $G$, and so has a spontaneous tendency to occur. Only if $\mu_{1}=\mu_{2}$ is there no change in $G$, and only then is the system at equilibrium.

\subsection*{Brief illustration 4A. 3}
The standard molar Gibbs energy of formation of water vapour at $298 \mathrm{~K}\left(25^{\circ} \mathrm{C}\right)$ is $-229 \mathrm{~kJ} \mathrm{~mol}^{-1}$, and that of liquid water at the same temperature is $-237 \mathrm{~kJ} \mathrm{~mol}^{-1}$. It follows that there is a decrease in Gibbs energy when water vapour condenses to the liquid at 298 K , so condensation is spontaneous at that temperature (and 1 bar ).

\subsection*{4A. 2 Phase boundaries}
The phase diagram of a pure substance shows the regions of pressure and temperature at which its various phases are thermodynamically stable (Fig. 4A.4). In fact, any two intensive variables may be used (such as temperature and magnetic field; in Topic 5A mole fraction is another variable), but this Topic focuses on pressure and temperature. The lines separating the regions, which are called phase boundaries (or coexistence curves), show the values of $p$ and $T$ at which two phases coexist in equilibrium and their chemical potentials are equal. A single phase is represented by an area on a phase diagram.

\subsection*{(a) Characteristic properties related to phase transitions}
Consider a liquid sample of a pure substance in a closed vessel. The pressure of a vapour in equilibrium with the liquid\\
\includegraphics[max width=\textwidth, center]{2025_02_27_d4b95d9718f0b9e2e6b9g-154(2)}

Figure 4A. 4 The general regions of pressure and temperature where solid, liquid, or gas is stable (that is, has minimum molar Gibbs energy) are shown on this phase diagram. For example, the solid phase is the most stable phase at low temperatures and high pressures.\\
is its vapour pressure (the property introduced in Topic 1 C ; Fig. 4A.5). Therefore, the liquid-vapour phase boundary in a phase diagram shows how the vapour pressure of the liquid varies with temperature. Similarly, the solid-vapour phase boundary shows the temperature variation of the sublimation vapour pressure, the vapour pressure of the solid phase. The vapour pressure of a substance increases with temperature because at higher temperatures more molecules have sufficient energy to escape from their neighbours.

When a liquid is in an open vessel and subject to an external pressure, it is possible for the liquid to vaporize from its surface. However, only when the temperature is such that the vapour pressure is equal to the external pressure will it be possible for vaporization to occur throughout the bulk of the liquid and for the vapour to expand freely into the surroundings. This condition of free vaporization throughout the liquid is called boiling. The temperature at which the vapour pres-\\
\includegraphics[max width=\textwidth, center]{2025_02_27_d4b95d9718f0b9e2e6b9g-154}

Figure 4A. 5 The vapour pressure of a liquid or solid is the pressure exerted by the vapour in equilibrium with the condensed phase.\\
sure of a liquid is equal to the external pressure is called the boiling temperature at that pressure. For the special case of an external pressure of 1 atm , the boiling temperature is called the normal boiling point, $T_{\mathrm{b}}$. With the replacement of 1 atm by 1 bar as standard pressure, there is some advantage in using the standard boiling point instead: this is the temperature at which the vapour pressure reaches 1 bar. Because 1 bar is slightly less than 1 atm ( $1.00 \mathrm{bar}=0.987 \mathrm{~atm}$ ), the standard boiling point of a liquid is slightly lower than its normal boiling point. For example, the normal boiling point of water is $100.0^{\circ} \mathrm{C}$, but its standard boiling point is $99.6^{\circ} \mathrm{C}$.

Boiling does not occur when a liquid is heated in a rigid, closed vessel. Instead, the vapour pressure, and hence the density of the vapour, rises as the temperature is raised (Fig. 4A.6). At the same time, the density of the liquid decreases slightly as a result of its expansion. There comes a stage when the density of the vapour is equal to that of the remaining liquid and the surface between the two phases disappears. The temperature at which the surface disappears is the critical temperature, $T_{c}$, of the substance. The vapour pressure at the critical temperature is called the critical pressure, $p_{c}$. At and above the critical temperature, a single uniform phase called a supercritical fluid fills the container and an interface no longer exists. That is, above the critical temperature, the liquid phase of the substance does not exist.

The temperature at which, under a specified pressure, the liquid and solid phases of a substance coexist in equilibrium is called the melting temperature. Because a substance melts at exactly the same temperature as it freezes, the melting temperature of a substance is the same as its freezing temperature. The freezing temperature when the pressure is 1 atm is called\\
\includegraphics[max width=\textwidth, center]{2025_02_27_d4b95d9718f0b9e2e6b9g-155}

Figure 4A. 6 (a) A liquid in equilibrium with its vapour. (b) When a liquid is heated in a sealed container, the density of the vapour phase increases and the density of the liquid decreases slightly. There comes a stage, (c), at which the two densities are equal and the interface between the fluids disappears. This disappearance occurs at the critical temperature.\\
the normal freezing point, $T_{\mathrm{f}}$, and its freezing point when the pressure is 1 bar is called the standard freezing point. The normal and standard freezing points are negligibly different for most purposes. The normal freezing point is also called the normal melting point.

There is a set of conditions under which three different phases of a substance (typically solid, liquid, and vapour) all simultaneously coexist in equilibrium. These conditions are represented by the triple point, a point at which the three phase boundaries meet. The temperature at the triple point is denoted $T_{3}$. The triple point of a pure substance cannot be changed: it occurs at a single definite pressure and temperature characteristic of the substance.

As can be seen from Fig. 4A.4, the triple point marks the lowest pressure at which a liquid phase of a substance can exist. If (as is common) the slope of the solid-liquid phase boundary is as shown in the diagram, then the triple point also marks the lowest temperature at which the liquid can exist.

\subsection*{Brief illustration 4A. 4}
The triple point of water lies at 273.16 K and 611 Pa ( 6.11 mbar , 4.58 Torr), and the three phases of water (ice, liquid water, and water vapour) coexist in equilibrium at no other combination of pressure and temperature. This invariance of the triple point was the basis of its use in the now superseded definition of the Kelvin scale of temperature (Topic 3A).

\subsection*{(b) The phase rule}
In one of the most elegant arguments in the whole of chemical thermodynamics, J.W. Gibbs deduced the phase rule, which gives the number of parameters that can be varied independently (at least to a small extent) while the number of phases in equilibrium is preserved. The phase rule is a general relation between the variance, $F$, the number of components, $C$, and the number of phases at equilibrium, $P$, for a system of any composition. Each of these quantities has a precisely defined meaning:

\begin{itemize}
  \item The variance (or number of degrees of freedom), $F$, of a system is the number of intensive variables that can be changed independently without disturbing the number of phases in equilibrium.
  \item A constituent of a system is any chemical species that is present.
  \item A component is a chemically independent constituent of a system.
  \item The number of components, $C$, in a system is the minimum number of types of independent species (ions or molecules) necessary to define the composition of all the phases present in the system.
\end{itemize}

\subsection*{Brief illustration 4A. 5}
A mixture of ethanol and water has two constituents. A solution of sodium chloride has three constituents: water, $\mathrm{Na}^{+}$ions, and $\mathrm{Cl}^{-}$ions, but only two components because the numbers of $\mathrm{Na}^{+}$and $\mathrm{Cl}^{-}$ions are constrained to be equal by the requirement of charge neutrality.

The relation between these quantities, which is called the phase rule, is established by considering the conditions for equilibrium to exist between the phases in terms of the chemical potentials of all the constituents.

\subsection*{How is that done? 4A. 1 Deducing the phase rule}
The argument that leads to the phase rule is most easily appreciated by first thinking about the simpler case when only one component is present and then generalizing the result to an arbitrary number of components.

Step 1 Consider the case where only one component is present When only one phase is present $(P=1)$, both $p$ and $T$ can be varied independently, so $F=2$. Now consider the case where two phases $\alpha$ and $\beta$ are in equilibrium ( $P=2$ ). If the phases are in equilibrium at a given pressure and temperature, their chemical potentials must be equal:

$$
\mu(\alpha ; p, T)=\mu(\beta ; p, T)
$$

This equation relates $p$ and $T$ : when the pressure changes, the changes in the chemical potentials are different in general, so in order to keep them equal, the temperature must change too. To keep the two phases in equilibrium only one variable can be changed arbitrarily, so $F=1$.

If three phases of a one-component system are in mutual equilibrium, the chemical potentials of all three phases ( $\alpha, \beta$, and $\gamma$ ) must be equal:

$$
\mu(\alpha ; p, T)=\mu(\beta ; p, T)=\mu(\gamma ; p, T)
$$

This relation is actually two equations $\mu(\alpha ; p, T)=\mu(\beta ; p, T)$ and $\mu(\beta ; p, T)=\mu(\gamma ; p, T)$, in which there are two variables: pressure and temperature. With two equations for two unknowns, there is a single solution (just as the pair of algebraic equations $x+y=x y$ and $3 x-y=x y$ have the single, fixed solutions $x=2$ and $y=2$ ). There is therefore only one single, unchangeable value of the pressure and temperature as a solution. The conclusion is that there is no freedom to choose these variables, so $F=0$.

Four phases cannot be in mutual equilibrium in a onecomponent system because the three equalities

$$
\begin{aligned}
\mu(\alpha ; p, T)= & \mu(\beta ; p, T), \mu(\beta ; p, T)=\mu(\gamma ; p, T), \\
& \text { and } \mu(\gamma ; p, T)=\mu(\delta ; p, T)
\end{aligned}
$$

are three equations with only two unknowns ( $p$ and $T$ ), which are not consistent because no values of $p$ and $T$ satisfy all three\\
equations (just as the three equations $x+y=x y, 3 x-y=x y$, and $4 x-y=2 x y^{2}$ have no solution).

In summary, for a one-component system $(C=1)$ it has been shown that: $F=2$ when $P=1 ; F=1$ when $P=2$; and $F$ $=0$ when $P=3$. The general result is that for $C=1, F=3-P$.

Step 2 Consider the general case of any number of components, $C$\\
Begin by counting the total number of intensive variables. The pressure, $p$, and temperature, $T$, count as 2 . The composition of a phase is specified by giving the mole fractions of the $C$ components, but as the sum of the mole fractions must be 1 , only $C-1$ mole fractions are independent. Because there are $P$ phases, the total number of composition variables is $P(C-1)$. At this stage, the total number of intensive variables is $P(C-1)+2$.

At equilibrium, the chemical potential of a component $J$ is the same in every phase:

$$
\mu_{\mathrm{J}}(\alpha ; p, T)=\mu_{\mathrm{J}}(\beta ; p, T)=\cdots \text { for } P \text { phases }
$$

There are $P-1$ equations of this kind to be satisfied for each component J. As there are $C$ components, the total number of equations is $C(P-1)$. Each equation reduces the freedom to vary one of the $P(C-1)+2$ intensive variables. It follows that the total number of degrees of freedom is

$$
F=P(C-1)+2-C(P-1)
$$

The right-hand side simplifies to give the phase rule in the form derived by Gibbs:\\
\includegraphics[max width=\textwidth, center]{2025_02_27_d4b95d9718f0b9e2e6b9g-156}

The implications of the phase rule for a one-component system, when

\[
\begin{array}{ll}
F=3-P & \begin{array}{l}
\text { The phase rule } \\
{[C=1]}
\end{array} \tag{4A.2}
\end{array}
\]

are summarized in Fig. 4A.7. When only one phase is present in a one-component system, $F=2$ and both $p$ and $T$ can be varied independently (at least over a small range) without changing the number of phases. The system is said to be bivariant, meaning having two degrees of freedom. In other words, a single phase is represented by an area on a phase diagram.

When two phases are in equilibrium $F=1$, which implies that pressure is not freely variable if the temperature is set; indeed, at a given temperature, a liquid has a characteristic vapour pressure. It follows that the equilibrium of two phases is represented by a line in the phase diagram. Instead of selecting the temperature, the pressure could be selected, but having done so the two phases would be in equilibrium only at a single definite temperature. Therefore, freezing (or any other phase transition) occurs at a definite temperature at a given pressure.

When three phases are in equilibrium, $F=0$ and the system is invariant, meaning that it has no degrees of freedom. This\\
\includegraphics[max width=\textwidth, center]{2025_02_27_d4b95d9718f0b9e2e6b9g-157(1)}

Figure 4A. 7 The typical regions of a one-component phase diagram. The lines represent conditions under which the two adjoining phases are in equilibrium. A point represents the unique set of conditions under which three phases coexist in equilibrium. Four phases cannot mutually coexist in equilibrium when only one component is present.\\
special condition can be established only at a definite temperature and pressure that is characteristic of the substance and cannot be changed. The equilibrium of three phases is therefore represented by a point, the triple point, on a phase diagram. Four phases cannot be in equilibrium in a one-component system because $F$ cannot be negative.

\subsection*{4A. 3 Three representative phase diagrams}
Carbon dioxide, water, and helium illustrate the significance of the various features of a phase diagram.

\subsection*{(a) Carbon dioxide}
Figure 4A. 8 shows the phase diagram for carbon dioxide. The features to notice include the positive slope (up from left to right) of the solid-liquid phase boundary; the direction of this line is characteristic of most substances. This slope indicates that the melting temperature of solid carbon dioxide rises as the pressure is increased. Notice also that, as the triple point lies above 1 atm , the liquid cannot exist at normal atmospheric pressures whatever the temperature. As a result, the solid sublimes when left in the open (hence the name 'dry ice'). To obtain the liquid, it is necessary to exert a pressure of at least 5.11 atm . Cylinders of carbon dioxide generally contain the liquid or compressed gas; at $25^{\circ} \mathrm{C}$ that implies a vapour pressure of 67 atm if both gas and liquid are present in equilibrium. When the gas is released through a tap (which acts as a throttle) the gas cools by the Joule-Thomson effect, so when it emerges into a region where the pressure is only 1 atm , it condenses into a finely divided snow-like solid. That carbon\\
\includegraphics[max width=\textwidth, center]{2025_02_27_d4b95d9718f0b9e2e6b9g-157}

Figure 4A. 8 The experimental phase diagram for carbon dioxide; note the break in the vertical scale. As the triple point lies at pressures well above atmospheric, liquid carbon dioxide does not exist under normal conditions; a pressure of at least 5.11 atm must be applied for liquid to be formed. The path ABCD is discussed in Brief illustration 4A.6.\\
dioxide gas cannot be liquefied except by applying high pressure reflects the weakness of the intermolecular forces between the nonpolar carbon dioxide molecules (Topic 14B).

\subsection*{Brief illustration 4A. 6}
Consider the path ABCD in Fig. 4A.8. At A the carbon dioxide is a gas. When the temperature and pressure are adjusted to $B$, the vapour condenses directly to a solid. Increasing the pressure and temperature to C results in the formation of the liquid phase, which evaporates to the vapour when the conditions are changed to D.

\subsection*{(b) Water}
Figure 4A. 9 shows the phase diagram for water. The liquidvapour boundary in the phase diagram summarizes how the vapour pressure of liquid water varies with temperature. It also summarizes how the boiling temperature varies with pressure: simply read off the temperature at which the vapour pressure is equal to the prevailing atmospheric pressure. The solid (ice I)-liquid boundary shows how the melting temperature varies with the pressure. Its very steep slope indicates that enormous pressures are needed to bring about significant changes. Notice that the line has a negative slope (down from left to right) up to 2 kbar , which means that the melting temperature falls as the pressure is raised.

The reason for this almost unique behaviour can be traced to the decrease in volume that occurs on melting: it is more favourable for the solid to transform into the liquid as the pressure is raised. The decrease in volume is a result of the very open structure of ice: as shown in Fig. 4A.10, the water molecules are held apart, as well as together, by the hydrogen bonds\\
\includegraphics[max width=\textwidth, center]{2025_02_27_d4b95d9718f0b9e2e6b9g-158(1)}

Figure 4A. 9 The phase diagram for water showing the different solid phases, which are indicated with Roman numerals I, II, ...; solid phase I (ice I) is ordinary ice. The path ABCD is discussed in Brief illustration 4A.7.\\
between them, but the hydrogen-bonded structure partially collapses on melting and the liquid is denser than the solid. Other consequences of its extensive hydrogen bonding are the anomalously high boiling point of water for a molecule of its molar mass and its high critical temperature and pressure.

The diagram shows that water has one liquid phase but many different solid phases other than ordinary ice ('ice I'). Some of these phases melt at high temperatures. Ice VII, for instance, melts at $100^{\circ} \mathrm{C}$ but exists only above 25 kbar . Two further phases, Ice XIII and XIV, were identified in 2006 at $-160^{\circ} \mathrm{C}$ but have not yet been allocated regions in the phase diagram. Note that five more triple points occur in the diagram other than the one where vapour, liquid, and ice I coexist. Each one occurs at a definite pressure and temperature that cannot be changed. The solid phases of ice differ in the arrangement of the water molecules: under the influence of very high pressures, hydrogen bonds buckle and the $\mathrm{H}_{2} \mathrm{O}$ molecules adopt different arrangements. These polymorphs of ice may contribute to the advance of glaciers, for ice at the bottom of glaciers experiences very high pressures where it rests on jagged rocks.\\
\includegraphics[max width=\textwidth, center]{2025_02_27_d4b95d9718f0b9e2e6b9g-158}

Figure 4A. 10 A fragment of the structure of ice I . Each O atom is linked by two covalent bonds to H atoms and by two hydrogen bonds to a neighbouring O atom, in a tetrahedral array.

\subsection*{Brief illustration 4A. 7}
Consider the path ABCD in Fig. 4A.9. Water is present at A as ice V. Increasing the pressure to $B$ at the same temperature results in the formation of ice VIII. Heating to C leads to the formation of ice VII, and reduction in pressure to D results in the solid melting to liquid.

\subsection*{(c) Helium}
The two isotopes of helium, ${ }^{3} \mathrm{He}$ and ${ }^{4} \mathrm{He}$, behave differently at low temperatures because ${ }^{4} \mathrm{He}$ is a boson whereas ${ }^{3} \mathrm{He}$ is a fermion, and are treated differently by the Pauli principle (Topic 8B). Figure 4A. 11 shows the phase diagram of he-lium-4. Helium behaves unusually at low temperatures because the mass of its atoms is so low and there are only very weak interactions between neighbours. At 1 atm , the solid and gas phases of helium are never in equilibrium however low the temperature: the atoms are so light that they vibrate with a large-amplitude motion even at very low temperatures and the solid simply shakes itself apart. Solid helium can be obtained, but only by holding the atoms together by applying pressure.

Pure helium-4 has two liquid phases. The phase marked He-I in the diagram behaves like a normal liquid; the other phase, He-II, is a superfluid. It is so called because it flows without viscosity. ${ }^{1}$ The liquid-liquid phase boundary is called the $\lambda$-line (lambda line) for reasons related to the shape of a\\
\includegraphics[max width=\textwidth, center]{2025_02_27_d4b95d9718f0b9e2e6b9g-158(2)}

Figure 4A. 11 The phase diagram for helium $\left({ }^{4} \mathrm{He}\right)$. The $\lambda$-line marks the conditions under which the two liquid phases are in equilibrium; He -II is the superfluid phase. Note that a pressure of over 20 bar must be exerted before solid helium can be obtained. The labels hcp and bcc denote different solid phases in which the atoms pack together differently: hcp denotes hexagonal closed packing and bce denotes body-centred cubic (Topic 15A). The path ABCD is discussed in Briefillustration 4A.8.

\footnotetext{${ }^{1}$ Water might also have a superfluid liquid phase.
}
\includegraphics[max width=\textwidth, center]{2025_02_27_d4b95d9718f0b9e2e6b9g-159}

Figure 4A. 12 The heat capacity of superfluid He -II increases with temperature and rises steeply as the transition temperature to He -l is approached. The appearance of the plot has led the transition to be described as a $\lambda$-transition and the line on the phase diagram to be called a $\lambda$-line.\\
plot of the heat capacity of helium-4 against temperature at the transition temperature (Fig. 4A.12).

Helium-3 also has a superfluid phase. Helium-3 is unusual in that melting is exothermic $\left(\Delta_{\text {fus }} H<0\right)$ and therefore (from $\Delta_{\text {fus }} S=\Delta_{\text {fus }} H / T_{\mathrm{f}}$ ) at the melting point the entropy of the liquid is lower than that of the solid.

\subsection*{Brief illustration 4A. 8}
Consider the path ABCD in Fig. 4A.11. At A, helium is present as a vapour. On cooling to B it condenses to helium-I, and further cooling to C results in the formation of heliumII. Adjustment of the pressure and temperature to $D$ results in a system in which three phases, helium-I, helium-II, and vapour are in mutual equilibrium.

\subsection*{Checklist of concepts}
1. A phase is a form of matter that is uniform throughout in chemical composition and physical state.2. A phase transition is the spontaneous conversion of one phase into another.3. The thermodynamic analysis of phases is based on the fact that at equilibrium, the chemical potential of a substance is the same throughout a sample.4. A phase diagram indicates the values of the pressure and temperature at which a particular phase is most stable, or is in equilibrium with other phases.$\square$ 5. The phase rule relates the number of variables that may be changed while the phases of a system remain in mutual equilibrium.

\subsection*{Checklist of equations}
\begin{center}
\begin{tabular}{llll}
\hline
Property & Equation & Comment & Equation number \\
\hline
Chemical potential & $\mu=G_{\mathrm{m}}$ & For a single substance &  \\
Phase rule & $F=C-P+2$ & $F$ is the variance, $C$ the number of &  \\
components, and $P$ the number of phases &  &  &  \\
\end{tabular}
\end{center}

\section*{TOPIC 4B Thermodynamic aspects of phase transitions }
\subsection*{Why do you need to know this material?}
Thermodynamic arguments explain the appearance of phase diagrams and can be used to make predictions about the effect of pressure on phase transitions. They provide insight into the properties that account for the behaviour of matter under different conditions.

\subsection*{What is the key idea?}
The effect of temperature and pressure on the chemical potential of a substance in each phase depends on its molar entropy and molar volume, respectively.

\subsection*{What do you need to know already?}
You need to be aware that phases are in equilibrium when their chemical potentials are equal (Topic 4A) and that the variation of the molar Gibbs energy of a substance depends on its molar volume and entropy (Topic 3E). The Topic makes use of expressions for the entropy of transition (Topic 3B) and of the perfect gas law (Topic 1A).

As explained in Topic 4A, the thermodynamic criterion for phase equilibrium is the equality of the chemical potentials of each substance in each phase. For a one-component system, the chemical potential is the same as the molar Gibbs energy ( $\mu=G_{\mathrm{m}}$ ). In Topic 3E it is explained how the Gibbs energy varies with temperature and pressure:\\
$\mathrm{d} G=-S \mathrm{~d} T$ at constant pressure;\\
$\mathrm{d} G=V \mathrm{~d} p$ at constant temperature

These expressions also apply to the molar Gibbs energy, and therefore to the chemical potential. By using the notation of partial derivatives (The chemist's toolkit 9 in Topic 2A) they can be expressed as

$$
\left(\frac{\partial \mu}{\partial T}\right)_{p}=-S_{\mathrm{m}}
$$

Variation of chemical potential with $T$\\
(4B.1a)

\[
\left(\frac{\partial \mu}{\partial p}\right)_{T}=V_{\mathrm{m}} \quad \begin{align*}
& \text { Variation of chemical }  \tag{4B.1b}\\
& \text { potential with } p \\
& {[\text { constant } T \text { ] }}
\end{align*}
\]

By combining the equality of chemical potentials of a substance in each phase with these expressions for the variation of $\mu$ with temperature and pressure it is possible to deduce how phase equilibria respond to changes in the conditions.

\subsection*{4B. 1 The dependence of stability on the conditions}
At sufficiently low temperatures the solid phase of a substance commonly has the lowest chemical potential and is therefore the most stable phase. However, the chemical potentials of different phases depend on temperature to different extents (because the molar entropy of each phase is different), and above a certain temperature the chemical potential of another phase (perhaps another solid phase, a liquid, or a gas) might turn out to be lower. Then a transition to the second phase becomes spontaneous and occurs if it is kinetically feasible.

\subsection*{(a) The temperature dependence of phase stability}
Because $S_{\mathrm{m}}>0$ for all substances above $T=0$, eqn 4B.1a shows that the chemical potential of a pure substance decreases as the temperature is raised. That is, a plot of chemical potential against temperature slopes down from left to right. It also implies that because $S_{\mathrm{m}}(\mathrm{g})>S_{\mathrm{m}}(\mathrm{l})$, the slope is steeper for gases than for liquids. Because it is almost always the case that $S_{\mathrm{m}}(\mathrm{l})$ $>S_{\mathrm{m}}(\mathrm{s})$, the slope is also steeper for a liquid than the corresponding solid. These features are illustrated in Fig. 4B.1. The steeper slope of $\mu(\mathrm{l})$ compared with that of $\mu(\mathrm{s})$ results in $\mu(\mathrm{l})$ falling below $\mu(\mathrm{s})$ when the temperature is high enough; then the liquid becomes the stable phase, and melting is spontaneous. The chemical potential of the gas phase plunges steeply downwards as the temperature is raised (because the molar entropy of the vapour is so high), and there comes a temperature at which it lies below that of the liquid. Then the gas is the stable phase and vaporization is spontaneous.\\
\includegraphics[max width=\textwidth, center]{2025_02_27_d4b95d9718f0b9e2e6b9g-161}

Figure 4B. 1 The schematic temperature dependence of the chemical potential of the solid, liquid, and gas phases of a substance (in practice, the lines are curved). The phase with the lowest chemical potential at a specified temperature is the most stable one at that temperature. The transition temperatures, the freezing (melting) and boiling temperatures ( $T_{\mathrm{f}}$ and $T_{\mathrm{b}}$, respectively), are the temperatures at which the chemical potentials of the two phases are equal.

\subsection*{Brief illustration 4B. 1}
The standard molar entropy of liquid water at $100^{\circ} \mathrm{C}$ is $86.8 \mathrm{~J} \mathrm{~K}^{-1} \mathrm{~mol}^{-1}$ and that of water vapour at the same temperature is $195.98 \mathrm{~J} \mathrm{~K}^{-1} \mathrm{~mol}^{-1}$. It follows that when the temperature is raised by 1.0 K the changes in chemical potential are

$$
\begin{aligned}
& \Delta \mu(\mathrm{l}) \approx-S_{\mathrm{m}}(\mathrm{l}) \Delta T=-87 \mathrm{~J} \mathrm{~mol}^{-1} \\
& \Delta \mu(\mathrm{~g}) \approx-S_{\mathrm{m}}(\mathrm{~g}) \Delta T=-196 \mathrm{~J} \mathrm{~mol}^{-1}
\end{aligned}
$$

At $100^{\circ} \mathrm{C}$ the two phases are in equilibrium with equal chemical potentials. At $101^{\circ} \mathrm{C}$ the chemical potential of both vapour and liquid are lower than at $100^{\circ} \mathrm{C}$, but the chemical potential of the vapour has decreased by a greater amount. It follows that the vapour is the stable phase at the higher temperature, so vaporization will be spontaneous.

\subsection*{(b) The response of melting to applied pressure}
Equation 4B.1b shows that because $V_{\mathrm{m}}>0$, an increase in pressure raises the chemical potential of any pure substance. In most cases, $V_{\mathrm{m}}(\mathrm{l})>V_{\mathrm{m}}(\mathrm{s})$, so an increase in pressure increases the chemical potential of the liquid phase of a substance more than that of its solid phase. As shown in Fig. 4B.2(a), the effect of pressure in such a case is to raise the freezing temperature slightly. For water, however, $V_{\mathrm{m}}(\mathrm{l})<V_{\mathrm{m}}(\mathrm{s})$, and an increase in pressure increases the chemical potential of the solid more than that of the liquid. In this case, the freezing temperature is lowered slightly (Fig. 4B.2(b)).\\
\includegraphics[max width=\textwidth, center]{2025_02_27_d4b95d9718f0b9e2e6b9g-161(1)}

Figure 4B. 2 The pressure dependence of the chemical potential of a substance depends on the molar volume of the phase. The lines show schematically the effect of increasing pressure on the chemical potential of the solid and liquid phases (in practice, the lines are curved), and the corresponding effects on the freezing temperatures. (a) In this case the molar volume of the solid is smaller than that of the liquid and $\mu(\mathrm{s})$ increases less than $\mu(\mathrm{I})$. As a result, the freezing temperature rises. (b) Here the molar volume is greater for the solid than the liquid (as for water), $\mu(\mathrm{s})$ increases more strongly than $\mu(\mathrm{I})$, and the freezing temperature is lowered.

\subsection*{Example 4B. 1 Assessing the effect of pressure on the chemical potential}
Calculate the effect on the chemical potentials of ice and water of increasing the pressure from 1.00 bar to 2.00 bar at $0^{\circ} \mathrm{C}$. The mass density of ice is $0.917 \mathrm{~g} \mathrm{~cm}^{-3}$ and that of liquid water is $0.999 \mathrm{~g} \mathrm{~cm}^{-3}$ under these conditions.

Collect your thoughts From $\mathrm{d} \mu=V_{\mathrm{m}} \mathrm{d} p$, you can infer that the change in chemical potential of an incompressible substance when the pressure is changed by $\Delta p$ is $\Delta \mu=V_{\mathrm{m}} \Delta p$. Therefore, you need to know the molar volumes of the two phases of water. These values are obtained from the mass density, $\rho$, and the molar mass, $M$, by using $V_{\mathrm{m}}=M / \rho$. Then $\Delta \mu=M \Delta p / \rho$. To keep the units straight, you will need to express the mass densities in kilograms per cubic metre $\left(\mathrm{kg} \mathrm{m}^{-3}\right)$ and the molar mass in kilograms per mole $\left(\mathrm{kg} \mathrm{mol}^{-1}\right)$, and use $1 \mathrm{Pam}^{3}=1 \mathrm{~J}$.

The solution The molar mass of water is $18.02 \mathrm{~g} \mathrm{~mol}^{-1}$ (i.e. $1.802 \times 10^{-2} \mathrm{~kg} \mathrm{~mol}^{-1}$ ); therefore, when the pressure is increased by $1.00 \mathrm{bar}\left(1.00 \times 10^{5} \mathrm{~Pa}\right)$

$$
\begin{aligned}
& \Delta \mu(\text { ice })=\frac{\left(1.802 \times 10^{-2} \mathrm{~kg} \mathrm{~mol}^{-1}\right) \times\left(1.00 \times 10^{5} \mathrm{~Pa}\right)}{917 \mathrm{~kg} \mathrm{~m}^{-3}}=+1.97 \mathrm{~J} \mathrm{~mol}^{-1} \\
& \begin{aligned}
\Delta \mu(\text { water }) & =\frac{\left(1.802 \times 10^{-2} \mathrm{~kg} \mathrm{~mol}^{-1}\right) \times\left(1.00 \times 10^{5} \mathrm{~Pa}\right)}{999 \mathrm{~kg} \mathrm{~m}^{-3}} \\
& =+1.80 \mathrm{~J} \mathrm{~mol}^{-1}
\end{aligned}
\end{aligned}
$$

Comment. The chemical potential of ice rises by more than that of water, so if they are initially in equilibrium at 1 bar, then there is a tendency for the ice to melt at 2 bar.

Self-test 4B. 1 Calculate the effect of an increase in pressure of 1.00 bar on the liquid and solid phases of carbon dioxide (molar mass $44.0 \mathrm{~g} \mathrm{~mol}^{-1}$ ) in equilibrium with mass densities $2.35 \mathrm{~g} \mathrm{~cm}^{-3}$ and $2.50 \mathrm{~g} \mathrm{~cm}^{-3}$, respectively.

\begin{itemize}
  \item யuof ot spuət p!̣os
\end{itemize}

\subsection*{(c) The vapour pressure of a liquid subjected to pressure}
Pressure can be exerted on the condensed phase mechanically or by subjecting it to the applied pressure of an inert gas (Fig. 4B.3). In the latter case, the partial vapour pressure is the partial pressure of the vapour in equilibrium with the condensed phase. When pressure is applied to a condensed phase, its vapour pressure rises: in effect, molecules are squeezed out of the phase and escape as a gas. The effect can be explored thermodynamically and a relation established between the applied pressure $P$ and the vapour pressure $p$.

\subsection*{How is that done? 4B. 1 Deriving an expression for the vapour pressure of a pressurized liquid}
At equilibrium the chemical potentials of the liquid and its vapour are equal: $\mu(\mathrm{l})=\mu(\mathrm{g})$. It follows that, for any change that preserves equilibrium, the resulting change in $\mu(\mathrm{l})$ must be equal to the change in $\mu(\mathrm{g})$; therefore, $\mathrm{d} \mu(\mathrm{g})=\mathrm{d} \mu(\mathrm{l})$.

Step 1 Express changes in the chemical potentials that arise from changes in pressure\\
When the pressure $P$ on the liquid is increased by $\mathrm{d} P$, the chemical potential of the liquid changes by $\mathrm{d} \mu(\mathrm{l})=V_{\mathrm{m}}(1) \mathrm{d} P$.\\
\includegraphics[max width=\textwidth, center]{2025_02_27_d4b95d9718f0b9e2e6b9g-162(1)}

Figure 4B. 3 Pressure may be applied to a condensed phase either (a) by compressing it or (b) by subjecting it to an inert pressurizing gas. When pressure is applied, the vapour pressure of the condensed phase increases.

The chemical potential of the vapour changes by $\mathrm{d} \mu(\mathrm{g})=$ $V_{\mathrm{m}}(\mathrm{g}) \mathrm{d} p$, where $\mathrm{d} p$ is the change in the vapour pressure. If the vapour is treated as a perfect gas, the molar volume can be replaced by $V_{\mathrm{m}}(\mathrm{g})=R T / p$, to give $\mathrm{d} \mu(\mathrm{g})=(R T / p) \mathrm{d} p$.

Step 2 Equate the changes in chemical potentials of the vapour and the liquid\\
Equate $\mathrm{d} \mu(\mathrm{l})=V_{\mathrm{m}}(\mathrm{l}) \mathrm{d} P$ and $\mathrm{d} \mu(\mathrm{g})=(R T / p) \mathrm{d} p$ :

$$
\frac{R T \mathrm{~d} p}{p}=V_{\mathrm{m}}(\mathrm{l}) \mathrm{d} P
$$

Be careful to distinguish between $P$, the total pressure, and $p$, the partial vapour pressure.

Step 3 Set up the integration of this expression by identifying the appropriate limits\\
When there is no additional pressure acting on the liquid, $P$ (the pressure experienced by the liquid) is equal to the normal vapour pressure $p^{*}$, so when $P=p^{*}, p=p^{*}$ too. When there is an additional pressure $\Delta P$ on the liquid, so $P=p+\Delta P$, the vapour pressure is $p$ (the value required). Provided the effect of pressure on the vapour pressure is small (as will turn out to be the case) a good approximation is to replace the $p$ in $p+\Delta P$ by $p^{*}$ itself, and to set the upper limit of the integral to $p^{\star}+\Delta P$. The integrations required are therefore as follows:

$$
R T \int_{p^{*}}^{p} \frac{\mathrm{~d} p^{\prime}}{p^{\prime}}=\int_{p^{*}}^{p^{*}+\Delta P} V_{\mathrm{m}}(\mathrm{l}) \mathrm{d} P
$$

(In the first integral, the variable of integration has been changed from $p$ to $p^{\prime}$ to avoid confusion with the $p$ at the upper limit.)

\subsection*{Step 4 Carry out the integrations}
Divide both sides by $R T$ and assume that the molar volume of the liquid is the same throughout the small range of pressures involved:

$$
\overbrace{\int_{p^{*}}^{p} \frac{\mathrm{~d} p^{\prime}}{p^{\prime}}}^{\text {Integral A. } 2}=\frac{1}{R T} \int_{p^{*}}^{p^{*}+\Delta P} V_{\mathrm{m}}(1) \mathrm{d} P=\frac{V_{\mathrm{m}}(\mathrm{l})}{R T} \overbrace{\int_{p^{*}}^{p^{*}+\Delta P} \mathrm{~d} P}^{\text {Integral A. } 1}
$$

Both integrations are straightforward, and lead to

$$
\ln \frac{p}{p^{*}}=\frac{V_{\mathrm{m}}(1)}{R T} \Delta P
$$

which (by using $\mathrm{e}^{\ln x}=x$ ) rearranges to\\
\includegraphics[max width=\textwidth, center]{2025_02_27_d4b95d9718f0b9e2e6b9g-162}

One complication that has been ignored is that, if the condensed phase is a liquid, then the pressurizing gas might dissolve and change its properties. Another complication is that the gas-phase molecules might attract molecules out of the liquid by the process of gas solvation, the attachment of molecules to gas-phase species.

\subsection*{Brief illustration 4B. 2}
For water, which has mass density $0.997 \mathrm{~g} \mathrm{~cm}^{-3}$ at $25^{\circ} \mathrm{C}$ and therefore molar volume $18.1 \mathrm{~cm}^{3} \mathrm{~mol}^{-1}$, when the applied pressure is increased by 10 bar (i.e. $\Delta P=1.0 \times 10^{6} \mathrm{~Pa}$ )

$$
\begin{aligned}
\ln \frac{p}{p^{*}} & =\frac{V_{\mathrm{m}}(1) \Delta P}{R T}=\frac{\left(1.81 \times 10^{-5} \mathrm{~m}^{3} \mathrm{~mol}^{-1}\right) \times\left(1.0 \times 10^{6} \mathrm{~Pa}\right)}{\left(8.3145 \mathrm{JK}^{-1} \mathrm{~mol}^{-1}\right) \times(298 \mathrm{~K})} \\
& =0.0073 \ldots
\end{aligned}
$$

where $1 \mathrm{~J}=1 \mathrm{Pam}^{3}$. It follows that $p=1.0073 p^{*}$, an increase of only 0.73 per cent.

\subsection*{4B.2 The location of phase boundaries}
The precise locations of the phase boundaries-the pressures and temperatures at which two phases can coexist-can be found by making use once again of the fact that, when two phases are in equilibrium, their chemical potentials must be equal. Therefore, when the phases $\alpha$ and $\beta$ are in equilibrium,


\begin{equation*}
\mu(\alpha ; p, T)=\mu(\beta ; p, T) \tag{4B.3}
\end{equation*}


Solution of this equation for $p$ in terms of $T$ gives an equation for the phase boundary (the coexistence curve).

\subsection*{(a) The slopes of the phase boundaries}
Imagine that at some particular pressure and temperature the two phases are in equilibrium: their chemical potentials are then equal. Now $p$ and $T$ are changed infinitesimally, but in such a way that the phases remain in equilibrium: after these changes, the chemical potentials of the two phases change but remain equal (Fig. 4B.4). It follows that the change in the\\
\includegraphics[max width=\textwidth, center]{2025_02_27_d4b95d9718f0b9e2e6b9g-163}

Figure 4B. 4 When pressure is applied to a system in which two phases are in equilibrium (at a), the equilibrium is disturbed. It can be restored by changing the temperature, so moving the state of the system to $b$. It follows that there is a relation between $\mathrm{d} p$ and $\mathrm{d} T$ that ensures that the system remains in equilibrium as either variable is changed.\\
chemical potential of phase $\alpha$ must be the same as the change in chemical potential of phase $\beta$, so $\mathrm{d} \mu(\alpha)=\mathrm{d} \mu(\beta)$.

Equation 3E. $7(\mathrm{~d} G=V \mathrm{~d} p-S \mathrm{~d} T)$ gives the variation of $G$ with $p$ and $T$, so with $\mu=G_{\mathrm{m}}$, it follows that $\mathrm{d} \mu=V_{\mathrm{m}} \mathrm{d} p-S_{\mathrm{m}} \mathrm{d} T$ for each phase. Therefore the relation $\mathrm{d} \mu(\alpha)=\mathrm{d} \mu(\beta)$ can be written

$$
V_{\mathrm{m}}(\alpha) \mathrm{d} p-S_{\mathrm{m}}(\alpha) \mathrm{d} T=V_{\mathrm{m}}(\beta) \mathrm{d} p-S_{\mathrm{m}}(\beta) \mathrm{d} T
$$

where $S_{\mathrm{m}}(\alpha)$ and $S_{\mathrm{m}}(\beta)$ are the molar entropies of the two phases, and $V_{\mathrm{m}}(\alpha)$ and $V_{\mathrm{m}}(\beta)$ are their molar volumes. Hence

$$
\left\{S_{\mathrm{m}}(\beta)-S_{\mathrm{m}}(\alpha)\right\} \mathrm{d} T=\left\{V_{\mathrm{m}}(\beta)-V_{\mathrm{m}}(\alpha)\right\} \mathrm{d} p
$$

The change in (molar) entropy accompanying the phase transition, $\Delta_{\text {trs }} S$, is the difference in the molar entropies $\Delta_{\text {trs }} S=$ $S_{\mathrm{m}}(\beta)-S_{\mathrm{m}}(\alpha)$, and likewise for the change in (molar) volume, $\Delta_{\text {trs }} V=V_{\mathrm{m}}(\beta)-V_{\mathrm{m}}(\alpha)$. Therefore,

$$
\Delta_{\mathrm{trs}} S \mathrm{~d} T=\Delta_{\mathrm{trs}} V \mathrm{~d} p
$$

This relation turns into the Clapeyron equation:

$$
\frac{\mathrm{d} p}{\mathrm{~d} T}=\frac{\Delta_{\mathrm{trs}} S}{\Delta_{\mathrm{trs}} V}
$$

Clapeyron equation\\
(4B.4a)

The Clapeyron equation is an exact expression for the slope of the tangent to the phase boundary at any point and applies to any phase equilibrium of any pure substance. It implies that thermodynamic data can be used to predict the appearance of phase diagrams and to understand their form. A more practical application is to the prediction of the response of freezing and boiling points to the application of pressure, when it can be used in the form obtained by inverting both sides:


\begin{equation*}
\frac{\mathrm{d} T}{\mathrm{~d} p}=\frac{\Delta_{\mathrm{trs}} V}{\Delta_{\mathrm{trs}} S} \tag{4B.4b}
\end{equation*}


\subsection*{Brief illustration 4B.3}
For water at $0^{\circ} \mathrm{C}$, the standard volume of transition of ice to liquid is $-1.6 \mathrm{~cm}^{3} \mathrm{~mol}^{-1}$, and the corresponding standard entropy of transition is $+22 \mathrm{~J} \mathrm{~K}^{-1} \mathrm{~mol}^{-1}$. The slope of the solidliquid phase boundary at that temperature is therefore

$$
\begin{aligned}
\frac{\mathrm{d} T}{\mathrm{~d} p} & =\frac{-1.6 \times 10^{-6} \mathrm{~m}^{3} \mathrm{~mol}^{-1}}{22 \mathrm{JK}^{-1} \mathrm{~mol}^{-1}}=-7.3 \times 10^{-8} \frac{\mathrm{~K}}{\mathrm{Jm}^{-3}} \\
& =-7.3 \times 10^{-8} \mathrm{KPa}^{-1}
\end{aligned}
$$

which corresponds to $-7.3 \mathrm{mK} \mathrm{bar}^{-1}$. An increase of 100 bar therefore results in a lowering of the freezing point of water by 0.73 K .

\subsection*{(b) The solid-liquid boundary}
Melting (fusion) is accompanied by a molar enthalpy change $\Delta_{\text {fus }} H$, and if it occurs at a temperature $T$ the molar entropy of melting is $\Delta_{\text {fus }} H / T$ (Topic 3B); all points on the phase boundary correspond to equilibrium, so $T$ is in fact a transition temperature, $T_{\mathrm{trs}}$. The Clapeyron equation for this phase transition then becomes


\begin{equation*}
\frac{\mathrm{d} p}{\mathrm{~d} T}=\frac{\Delta_{\text {fus }} H}{T \Delta_{\text {fus }} V} \quad \text { Slope of solid-liquid boundary } \tag{4B.5}
\end{equation*}


where $\Delta_{\text {fus }} V$ is the change in molar volume that accompanies melting. The enthalpy of melting is positive (the only exception is helium-3); the change in molar volume is usually positive and always small. Consequently, the slope $\mathrm{d} p / \mathrm{d} T$ is steep and usually positive (Fig. 4B.5).

The equation for the phase boundary is found by integrating $\mathrm{d} p / \mathrm{d} T$ and assuming that $\Delta_{\text {fus }} H$ and $\Delta_{\text {fus }} V$ change so little with temperature and pressure that they can be treated as constant. If the melting temperature is $T^{*}$ when the pressure is $p^{*}$, and $T$ when the pressure is $p$, the integration required is

$$
\int_{p^{*}}^{p} \mathrm{~d} p=\frac{\Delta_{\text {fus }} H}{\Delta_{\text {fus }} V} \overbrace{\int_{T^{*}}}^{\text {Integral A. } 2}
$$

Therefore, the approximate equation of the solid-liquid boundary is


\begin{equation*}
p=p^{*}+\frac{\Delta_{\text {fus }} H}{\Delta_{\text {fus }}} \ln \frac{T}{T^{*}} \tag{4B.6}
\end{equation*}


This equation was originally obtained by yet another Thomson-James, the brother of William, Lord Kelvin.

When $T$ is close to $T^{*}$, the logarithm can be approximated by using the expansion $\ln (1+x)=x-\frac{1}{2} x^{2}+\cdots$ (see The\\
\includegraphics[max width=\textwidth, center]{2025_02_27_d4b95d9718f0b9e2e6b9g-164}

Figure 4B. 5 A typical solid-liquid phase boundary slopes steeply upwards. This slope implies that, as the pressure is raised, the melting temperature rises. Most substances behave in this way, water being the notable exception.\\
chemist's toolkit 12 in Topic 5B) and neglecting all but the leading term:

$$
\ln \frac{T}{T^{\star}}=\ln \left(1+\frac{T-T^{*}}{T^{*}}\right) \approx \frac{T-T^{*}}{T^{*}}
$$

Therefore


\begin{equation*}
p \approx p^{\star}+\frac{\Delta_{\text {fus }} H}{T^{\star} \Delta_{\text {fus }} V}\left(T-T^{\star}\right) \tag{4B.7}
\end{equation*}


This expression is the equation of a steep straight line when $p$ is plotted against $T$ (as in Fig. 4B.5).

\subsection*{Brief illustration 4B.4}
The enthalpy of fusion of ice at $0^{\circ} \mathrm{C}(273 \mathrm{~K})$ and 1 bar is $6.008 \mathrm{~kJ} \mathrm{~mol}^{-1}$ and the volume of fusion is $-1.6 \mathrm{~cm}^{3} \mathrm{~mol}^{-1}$. It follows that the solid-liquid phase boundary is given by the equation

$$
\begin{aligned}
p & \approx 1.0 \times 10^{5} \mathrm{~Pa}+\frac{6.008 \times 10^{3} \mathrm{Jmol}^{-1}}{(273 \mathrm{~K}) \times\left(-1.6 \times 10^{-6} \mathrm{~m}^{3} \mathrm{~mol}^{-1}\right)}\left(T-T^{\star}\right) \\
& \approx 1.0 \times 10^{5} \mathrm{~Pa}-1.4 \times 10^{7} \mathrm{PaK}^{-1}\left(T-T^{\star}\right)
\end{aligned}
$$

That is,

$$
p / \text { bar }=1-140\left(T-T^{*}\right) / \mathrm{K}
$$

with $T^{*}=273 \mathrm{~K}$. This expression is plotted in Fig. 4B.6.\\
\includegraphics[max width=\textwidth, center]{2025_02_27_d4b95d9718f0b9e2e6b9g-164(1)}

Figure 4B.6 The solid-liquid phase boundary (the melting point curve) for water as calculated in Brief illustration 4B.4. For comparison, the boundary for benzene is included.

\subsection*{(c) The liquid-vapour boundary}
The entropy of vaporization at a temperature $T$ is equal to $\Delta_{\text {vap }} H / T$ (as before, all points on the phase boundary correspond to equilibrium, so $T$ is a transition temperature, $T_{\text {trs }}$ ), so\\
the Clapeyron equation for the liquid-vapour boundary can therefore be written


\begin{equation*}
\frac{\mathrm{d} p}{\mathrm{~d} T}=\frac{\Delta_{\text {vap }} H}{T \Delta_{\text {vap }} V} \tag{4B.8}
\end{equation*}


Slope of liquid-vapour boundary

The enthalpy of vaporization is positive and $\Delta_{\text {vap }} V$ is large and positive, so $\mathrm{d} p / \mathrm{d} T$ is positive, but much smaller than for the solid-liquid boundary. Consequently $\mathrm{d} T / \mathrm{d} p$ is large, and the boiling temperature is more responsive to pressure than the freezing temperature.

\subsection*{Example 4B. 2}
Estimating the effect of pressure on the boiling temperature

Estimate the typical size of the effect of increasing pressure on the boiling point of a liquid.

Collect your thoughts To use eqn 4B. 8 you need to estimate the right-hand side. At the boiling point, the term $\Delta_{\text {vap }} H / T$ is Trouton's constant (Topic 3B). Because the molar volume of a gas is so much greater than the molar volume of a liquid, you can write $\Delta_{\text {vap }} V=V_{\mathrm{m}}(\mathrm{g})-V_{\mathrm{m}}(1) \approx V_{\mathrm{m}}(\mathrm{g})$ and take for $V_{\mathrm{m}}(\mathrm{g})$ the molar volume of a perfect gas (at low pressures, at least). You will need to use $1 \mathrm{~J}=1 \mathrm{Pam}^{3}$.

The solution Trouton's constant has the value $85 \mathrm{~J} \mathrm{~K}^{-1} \mathrm{~mol}^{-1}$. The molar volume of a perfect gas is about $25 \mathrm{dm}^{3} \mathrm{~mol}^{-1}$ at 1 atm and near but above room temperature. Therefore,

$$
\frac{\mathrm{d} p}{\mathrm{~d} T} \approx \frac{85 \mathrm{JK}^{-1} \mathrm{~mol}^{-1}}{2.5 \times 10^{-2} \mathrm{~m}^{3} \mathrm{~mol}^{-1}}=3.4 \times 10^{3} \mathrm{PaK}^{-1}
$$

This value corresponds to $0.034 \mathrm{~atm} \mathrm{~K}^{-1}$ and hence to $\mathrm{d} T / \mathrm{d} p=$ $29 \mathrm{Katm}^{-1}$. Therefore, a change of pressure of +0.1 atm can be expected to change a boiling temperature by about +3 K .\\
Self-test 4B. 2 Estimate $\mathrm{d} T / \mathrm{d} p$ for water at its normal boiling point using the information in Table 3B. 2 and $V_{\mathrm{m}}(\mathrm{g})=R T / p$.

Because the molar volume of a gas is so much greater than the molar volume of a liquid, $\Delta_{\text {vap }} V \approx V_{\mathrm{m}}(\mathrm{g})$ (as in Example 4B.2). Moreover, if the gas behaves perfectly, $V_{\mathrm{m}}(\mathrm{g})=R T / p$. These two approximations turn the exact Clapeyron equation into

$$
\frac{\mathrm{d} p}{\mathrm{~d} T}=\frac{\Delta_{\mathrm{vap}} H}{T(R T / p)}=\frac{p \Delta_{\mathrm{vap}} H}{R T^{2}}
$$

By using $\mathrm{d} x / x=\mathrm{d} \ln x$, this expression can be rearranged into the Clausius-Clapeyron equation for the variation of vapour pressure with temperature:

$$
\frac{\mathrm{d} \ln p}{\mathrm{~d} T}=\frac{\Delta_{\text {vap }} H}{R T^{2}}
$$

Clausius-Clapeyron equation\\
\includegraphics[max width=\textwidth, center]{2025_02_27_d4b95d9718f0b9e2e6b9g-165}

Figure 4B.7 A typical liquid-vapour phase boundary. The boundary can be interpreted as a plot of the vapour pressure against the temperature. This phase boundary terminates at the critical point (not shown).

Like the Clapeyron equation (which is exact), the ClausiusClapeyron equation (which is an approximation) is important for understanding the appearance of phase diagrams, particularly the location and shape of the liquid-vapour and solid-vapour phase boundaries. It can be used to predict how the vapour pressure varies with temperature and how the boiling temperature varies with pressure. For instance, if it is also assumed that the enthalpy of vaporization is independent of temperature, eqn 4B. 9 can be integrated as follows:

$$
\overbrace{\int_{\ln p^{*}}^{\ln p} \mathrm{~d} \ln p}^{\begin{array}{l}
\text { Integral A.1, } \\
\text { with } x=\ln p
\end{array}}=\frac{\Delta_{\text {vap }} H}{R} \overbrace{\int_{T^{*}}^{T} \frac{\mathrm{~d} T}{T^{2}}}^{\text {Integral A. } 1}
$$

hence

$$
\ln \frac{p}{p^{*}}=-\frac{\Delta_{\text {vap }} H}{R}\left(\frac{1}{T}-\frac{1}{T^{*}}\right)
$$

where $p^{*}$ is the vapour pressure when the temperature is $T^{*}$, and $p$ the vapour pressure when the temperature is $T$. It follows that


\begin{equation*}
p=p^{\star} \mathrm{e}^{-\chi} \quad \chi=\frac{\Delta_{\text {vap }} H}{R}\left(\frac{1}{T}-\frac{1}{T^{\star}}\right) \tag{4B.10}
\end{equation*}


Equation 4B. 10 is plotted as the liquid-vapour boundary in Fig. 4B.7. The line does not extend beyond the critical temperature, $T_{c}$, because above this temperature the liquid does not exist.

\subsection*{Brief illustration 4B. 5}
Equation 4B. 10 can be used to estimate the vapour pressure of a liquid at any temperature from knowledge of its normal boiling point, the temperature at which the vapour pressure is $1.00 \mathrm{~atm}(101 \mathrm{kPa})$. The normal boiling point of benzene is $80^{\circ} \mathrm{C}(353 \mathrm{~K})$ and (from Table 3B.2) $\Delta_{\text {vap }} H^{\ominus}=30.8 \mathrm{~kJ} \mathrm{~mol}^{-1}$.

Therefore, to calculate the vapour pressure at $20^{\circ} \mathrm{C}(293 \mathrm{~K})$, write

$$
\begin{aligned}
\chi & =\frac{3.08 \times 10^{4} \mathrm{Jmol}^{-1}}{8.3145 \mathrm{JK}^{-1} \mathrm{~mol}^{-1}}\left(\frac{1}{293 \mathrm{~K}}-\frac{1}{353 \mathrm{~K}}\right) \\
& =2.14 \ldots
\end{aligned}
$$

and substitute this value into eqn 4 B .10 with $p^{*}=101 \mathrm{kPa}$. The result is 12 kPa . The experimental value is 10 kPa .

A note on good practice Because exponential functions are so sensitive, it is good practice to carry out numerical calculations like this without evaluating the intermediate steps and using rounded values.

\subsection*{(d) The solid-vapour boundary}
The only difference between the solid-vapour and the liquidvapour boundary is the replacement of the enthalpy of vaporization by the enthalpy of sublimation, $\Delta_{\text {sub }} H$. Because the enthalpy of sublimation is greater than the enthalpy of vaporization (recall that $\Delta_{\text {sub }} H=\Delta_{\text {fus }} H+\Delta_{\text {vap }} H$ ), at similar temperatures the equation predicts a steeper slope for the sublimation curve than for the vaporization curve. These two boundaries meet at the triple point (Fig. 4B.8).

\subsection*{Brief illustration 4B. 6}
The enthalpy of fusion of ice at the triple point of water ( $6.1 \mathrm{mbar}, 273 \mathrm{~K}$ ) is negligibly different from its standard enthalpy of fusion at its freezing point, which is $6.008 \mathrm{~kJ} \mathrm{~mol}^{-1}$. The enthalpy of vaporization at that temperature is $45.0 \mathrm{~kJ} \mathrm{~mol}^{-1}$\\
(once again, ignoring differences due to the pressure not being 1 bar). The enthalpy of sublimation is therefore $51.0 \mathrm{~kJ} \mathrm{~mol}^{-1}$. Therefore, the equations for the slopes of (a) the liquid-vapour and (b) the solid-vapour phase boundaries at the triple point are\\
(a) $\frac{\mathrm{d} \ln p}{\mathrm{~d} T}=\frac{45.0 \times 10^{3} \mathrm{Jmol}^{-1}}{\left(8.3145 \mathrm{~J} \mathrm{~K}^{-1} \mathrm{~mol}^{-1}\right) \times(273 \mathrm{~K})^{2}}=0.0726 \mathrm{~K}^{-1}$\\
(b) $\frac{\mathrm{d} \ln p}{\mathrm{~d} T}=\frac{51.0 \times 10^{3} \mathrm{Jmol}^{-1}}{\left(8.3145 \mathrm{~J} \mathrm{~K}^{-1} \mathrm{~mol}^{-1}\right) \times(273 \mathrm{~K})^{2}}=0.0823 \mathrm{~K}^{-1}$

The slope of $\ln p$ plotted against $T$ is greater for the solidvapour boundary than for the liquid-vapour boundary at the triple point.\\
\includegraphics[max width=\textwidth, center]{2025_02_27_d4b95d9718f0b9e2e6b9g-166}

Figure 4B. 8 At temperatures close to the triple point the solidvapour boundary is steeper than the liquid-vapour boundary because the enthalpy of sublimation is greater than the enthalpy of vaporization.

\subsection*{Checklist of concepts}
\begin{enumerate}
  \item The chemical potential of a substance decreases with increasing temperature in proportion to its molar entropy.
  \item The chemical potential of a substance increases with increasing pressure in proportion to its molar volume.
  \item The vapour pressure of a condensed phase increases when pressure is applied.
  \item The Clapeyron equation is an exact expression for the slope of a phase boundary.
  \item The Clausius-Clapeyron equation is an approximate expression for the boundary between a condensed phase and its vapour.
\end{enumerate}

\subsection*{Checklist of equations}
\begin{center}
\begin{tabular}{llll}
\hline
Property & Equation & Comment & Equation number \\
\hline
Variation of $\mu$ with temperature & $(\partial \mu / \partial T)_{p}=-S_{\mathrm{m}}$ & $\mu=G_{\mathrm{m}}$ & 4 B.1a \\
Variation of $\mu$ with pressure & $(\partial \mu / \partial)_{T}=V_{\mathrm{m}}$ &  & $4 P=P-p^{*}$ \\
Vapour pressure in the presence of applied pressure & $p=p^{*} \mathrm{e}^{V_{\mathrm{m}}(1) \Delta P / R T}$ &  &  \\
Clapeyron equation & $\mathrm{d} p / \mathrm{d} T=\Delta_{\mathrm{trs}} S / \Delta_{\mathrm{ts}} V$ &  & 4 B.1b \\
Clausius-Clapeyron equation & $\mathrm{d} \ln p / \mathrm{d} T=\Delta_{\text {vap }} H / R T^{2}$ & \begin{tabular}{c}
Assumes $V_{\mathrm{m}}(\mathrm{g}) \gg V_{\mathrm{m}}(\mathrm{l})$ or $V_{\mathrm{m}}(\mathrm{s})$, \\
and vapour is a perfect gas \\
\end{tabular} & 4 4B.9 \\
\hline
\end{tabular}
\end{center}

\chapter{FOCUS 5 Simple mixtures}
Mixtures are an essential part of chemistry, either in their own right or as starting materials for chemical reactions. This group of Topics deals with the rich physical properties of mixtures and shows how to express them in terms of thermodynamic quantities.

\subsection*{5A The thermodynamic description of mixtures}
The first Topic in this Focus develops the concept of chemical potential as an example of a partial molar quantity and explores how to use the chemical potential of a substance to describe the physical properties of mixtures. The underlying principle to keep in mind is that at equilibrium the chemical potential of a species is the same in every phase. By making use of the experimental observations known as Raoult's and Henry's laws, it is possible to express the chemical potential of a substance in terms of its mole fraction in a mixture.\\
5A. 1 Partial molar quantities; 5A. 2 The thermodynamics of mixing;\\
5A. 3 The chemical potentials of liquids

\subsection*{5B The properties of solutions}
In this Topic, the concept of chemical potential is applied to the discussion of the effect of a solute on certain thermodynamic properties of a solution. These properties include the lowering of vapour pressure of the solvent, the elevation of its boiling point, the depression of its freezing point, and the origin of osmotic pressure. It is possible to construct a model of a certain class of real solutions called 'regular solutions', which have properties that diverge from those of ideal solutions.\\
5B. 1 Liquid mixtures; 5B. 2 Colligative properties

\subsection*{5C Phase diagrams of binary systems: liquids}
One widely employed device used to summarize the equilibrium properties of mixtures is the phase diagram. The Topic describes phase diagrams of systems of liquids with gradually increasing complexity. In each case the phase diagram for the system summarizes empirical observations on the conditions under which the liquid and vapour phases of the system are stable.\\
5C. 1 Vapour pressure diagrams; 5C. 2 Temperature-composition diagrams; 5C. 3 Distillation; 5C. 4 Liquid-liquid phase diagrams

\subsection*{5D Phase diagrams of binary systems: solids}
In this Topic it is seen how the phase diagrams of solid mixtures summarize experimental results on the conditions under which the liquid and solid phases of the system are stable.\\
5D. 1 Eutectics; 5D. 2 Reacting systems; 5D. 3 Incongruent melting

\subsection*{5E Phase diagrams of ternary systems}
Many modern materials (and ancient ones too) have more than two components. This Topic shows how phase diagrams are extended to the description of systems of three components and how to interpret triangular phase diagrams.\\
5E. 1 Triangular phase diagrams; 5E.2 Ternary systems

\subsection*{5F Activities}
The extension of the concept of chemical potential to real solutions involves introducing an effective concentration\\
called an 'activity'. In certain cases, the activity may be interpreted in terms of intermolecular interactions. An important example is an electrolyte solution. Such solutions often deviate considerably from ideal behaviour on account of the strong, long-range interactions between ions. This Topic shows how a model can be used to estimate the deviations from ideal behaviour when the solution is very dilute, and how to extend the resulting expressions to more concentrated solutions.

\footnotetext{5F. 1 The solvent activity; 5F. 2 The solute activity; 5F. 3 The activities of
} regular solutions; 5F. 4 The activities of ions

\subsection*{Web resources What is an application of this material?}
Two applications of this material are discussed, one from biology and the other from materials science, from among the huge number that could be chosen for this centrally important field. Impact 7 shows how the phenomenon of osmosis contributes to the ability of biological cells to maintain their shapes. In Impact 8, phase diagrams of the technologically important liquid crystals are discussed.

\section*{TOPIC 5A The thermodynamic description of mixtures }
\subsection*{Why do you need to know this material?}
Chemistry deals with a wide variety of mixtures, including mixtures of substances that can react together. Therefore, it is important to generalize the concepts introduced in Focus 4 to deal with substances that are mingled together.

\subsection*{> What is the key idea?}
The chemical potential of a substance in a mixture is a logarithmic function of its concentration.

\subsection*{> What do you need to know already?}
This Topic extends the concept of chemical potential to substances in mixtures by building on the concept introduced in the context of pure substances (Topic 4A). It makes use of the relation between the temperature dependence of the Gibbs energy and entropy (Topic 3E), and the concept of partial pressure (Topic 1A). Throughout this and related Topics various measures of concentration of a solute in a solution are used: they are summarized in The chemist's toolkit 11.

The consideration of mixtures of substances that do not react together is a first step towards dealing with chemical reactions (which are treated in Topic 6A). At this stage the discussion centres on binary mixtures, which are mixtures of two components, $A$ and $B$. In Topic 1A it is shown how the partial pressure, which is the contribution of one component to the total pressure, is used to discuss the properties of mixtures of gases. For a more general description of the thermodynamics of mixtures other analogous 'partial' properties need to be introduced.

\subsection*{5A. 1 Partial molar quantities}
The easiest partial molar property to visualize is the 'partial molar volume', the contribution that a component of a mixture makes to the total volume of a sample.

\subsection*{(a) Partial molar volume}
Imagine a huge volume of pure water at $25^{\circ} \mathrm{C}$. When a further $1 \mathrm{~mol} \mathrm{H}_{2} \mathrm{O}$ is added, the volume increases by $18 \mathrm{~cm}^{3}$ and it follows that the molar volume of pure water is $18 \mathrm{~cm}^{3} \mathrm{~mol}^{-1}$. However, upon adding $1 \mathrm{~mol} \mathrm{H}_{2} \mathrm{O}$ to a huge volume of pure ethanol, the volume is found to increase by only $14 \mathrm{~cm}^{3}$. The reason for the different increase in volume is that the volume occupied by a given number of water molecules depends on the identity of the molecules that surround them. In the latter case there is so much ethanol present that each $\mathrm{H}_{2} \mathrm{O}$ molecule is surrounded by ethanol molecules. The network of hydrogen bonds that normally hold $\mathrm{H}_{2} \mathrm{O}$ molecules at certain distances from each other in pure water does not form; as a result the $\mathrm{H}_{2} \mathrm{O}$ molecules are packed more tightly and so increase the volume by only $14 \mathrm{~cm}^{3}$. The quantity $14 \mathrm{~cm}^{3} \mathrm{~mol}^{-1}$ is the 'partial molar volume' of water in pure ethanol. In general, the partial molar volume of a substance A in a mixture is the change in volume per mole of A added to a large volume of the mixture.

The partial molar volumes of the components of a mixture vary with composition because the environment of each type of molecule changes as the composition changes from pure A to pure B . This changing molecular environment, and the consequential modification of the forces acting between molecules, results in the variation of the thermodynamic properties of a mixture as its composition is changed. The partial molar volumes of water and ethanol across the full composition range at $25^{\circ} \mathrm{C}$ are shown in Fig. 5A.1.

The partial molar volume, $V_{\mathrm{J}}$, of a substance J at some general composition is defined formally as follows:

\[
V_{\mathrm{J}}=\left(\frac{\partial V}{\partial n_{\mathrm{J}}}\right)_{p, T, n^{\prime}} \quad \begin{align*}
& \text { Partial molar volume }  \tag{5A.1}\\
& \text { [definition] }
\end{align*}
\]

where the subscript $n^{\prime}$ signifies that the amounts of all other substances present are constant. The partial molar volume is the slope of the plot of the total volume as the amount of J is changed, the pressure, temperature, and amount of the other components being constant (Fig. 5A.2). Its value depends on the composition, as seen for water and ethanol.\\
\includegraphics[max width=\textwidth, center]{2025_02_27_d4b95d9718f0b9e2e6b9g-176(1)}

Figure 5A. 1 The partial molar volumes of water and ethanol at $25^{\circ} \mathrm{C}$. Note the different scales (water on the left, ethanol on the right).\\
\includegraphics[max width=\textwidth, center]{2025_02_27_d4b95d9718f0b9e2e6b9g-176}

Figure 5A. 2 The partial molar volume of a substance is the slope of the variation of the total volume of the sample plotted against the amount of that substance. In general, partial molar quantities vary with the composition, as shown by the different slopes at $a$ and $b$. Note that the partial molar volume at $b$ is negative: the overall volume of the sample decreases as $A$ is added.

A note on good practice The IUPAC recommendation is to denote a partial molar quantity by $\bar{X}$, but only when there is the possibility of confusion with the quantity $X$. For instance, to avoid confusion, the partial molar volume of NaCl in water could be written $\bar{V}(\mathrm{NaCl}, \mathrm{aq})$ to distinguish it from the total volume of the solution, $V$.

The definition in eqn 5A. 1 implies that when the composition of a binary mixture is changed by the addition of $\mathrm{d} n_{\mathrm{A}}$ of A and $\mathrm{d} n_{\mathrm{B}}$ of B , then the total volume of the mixture changes by

$$
\begin{aligned}
\mathrm{d} V & =\left(\frac{\partial V}{\partial n_{\mathrm{A}}}\right)_{p, T, n_{\mathrm{B}}} \mathrm{~d} n_{\mathrm{A}}+\left(\frac{\partial V}{\partial n_{\mathrm{B}}}\right)_{p, T, n_{\mathrm{A}}} \mathrm{~d} n_{\mathrm{B}} \\
& =V_{\mathrm{A}} \mathrm{~d} n_{\mathrm{A}}+V_{\mathrm{B}} \mathrm{~d} n_{\mathrm{B}}
\end{aligned}
$$

This equation can be integrated with respect to $n_{\mathrm{A}}$ and $n_{\mathrm{B}}$ provided that the amounts of A and B are both increased in such a way as to keep their ratio constant. This linkage ensures that the partial molar volumes $V_{\mathrm{A}}$ and $V_{\mathrm{B}}$ are constant and so can be taken outside the integrals:


\begin{align*}
V & =\int_{0}^{n_{\mathrm{A}}} V_{\mathrm{A}} \mathrm{~d} n_{\mathrm{A}}+\int_{0}^{n_{\mathrm{B}}} V_{\mathrm{B}} \mathrm{~d} n_{\mathrm{B}}=V_{\mathrm{A}} \int_{0}^{n_{\mathrm{A}}} \mathrm{~d} n_{\mathrm{A}}+V_{\mathrm{B}} \int_{0}^{n_{\mathrm{B}}} \mathrm{~d} n_{\mathrm{B}}  \tag{5A.3}\\
& =V_{\mathrm{A}} n_{\mathrm{A}}+V_{\mathrm{B}} n_{\mathrm{B}}
\end{align*}


Although the two integrations are linked (in order to preserve constant relative composition), because $V$ is a state function the final result in eqn 5 A .3 is valid however the solution is in fact prepared.

Partial molar volumes can be measured in several ways. One method is to measure the dependence of the volume on the composition and to fit the observed volume to a function of the amount of the substance. Once the function has been found, its slope can be determined at any composition of interest by differentiation.

\subsection*{Example 5A. 1 Determining a partial molar volume}
A polynomial fit to measurements of the total volume of a water/ethanol mixture at $25^{\circ} \mathrm{C}$ that contains 1.000 kg of water is

$$
v=1002.93+54.6664 z-0.36394 z^{2}+0.028256 z^{3}
$$

where $v=V / \mathrm{cm}^{3}, z=n_{\mathrm{E}} / \mathrm{mol}$, and $n_{\mathrm{E}}$ is the amount of $\mathrm{CH}_{3} \mathrm{CH}_{2} \mathrm{OH}$ present. Determine the partial molar volume of ethanol.

Collect your thoughts Apply the definition in eqn 5A.1, taking care to convert the derivative with respect to $n$ to a derivative with respect to $z$ and keeping the units intact.

The solution The partial molar volume of ethanol, $V_{\mathrm{E}}$, is

$$
\begin{aligned}
V_{\mathrm{E}} & =\left(\frac{\partial V}{\partial n_{\mathrm{E}}}\right)_{p, T, n_{\mathrm{W}}}=\left(\frac{\partial\left(V / \mathrm{cm}^{3}\right)}{\partial\left(n_{\mathrm{E}} / \mathrm{mol}\right)}\right)_{p, T, n_{\mathrm{W}}} \frac{\mathrm{~cm}^{3}}{\mathrm{~mol}} \\
& =\left(\frac{\partial v}{\partial z}\right)_{p, T, n_{\mathrm{W}}} \mathrm{~cm}^{3} \mathrm{~mol}^{-1}
\end{aligned}
$$

Then, because

$$
\frac{\mathrm{d} v}{\mathrm{~d} z}=54.6664-2(0.36394) z+3(0.028256) z^{2}
$$

it follows that

$$
V_{\mathrm{E}} /\left(\mathrm{cm}^{3} \mathrm{~mol}^{-1}\right)=54.6664-0.72788 z+0.084768 z^{2}
$$

Figure 5A. 3 shows a graph of this function.\\
\includegraphics[max width=\textwidth, center]{2025_02_27_d4b95d9718f0b9e2e6b9g-177(3)}

Figure 5A. 3 The partial molar volume of ethanol as expressed by the polynomial in Example 5A.1.

Self-test 5 A .1 At $25^{\circ} \mathrm{C}$, the mass density of a 50 per cent by mass ethanol/water solution is $0.914 \mathrm{~g} \mathrm{~cm}^{-3}$. Given that the partial molar volume of water in the solution is $17.4 \mathrm{~cm}^{3} \mathrm{~mol}^{-1}$, what is the partial molar volume of the ethanol?\\
\includegraphics[max width=\textwidth, center]{2025_02_27_d4b95d9718f0b9e2e6b9g-177(1)}\\
\includegraphics[max width=\textwidth, center]{2025_02_27_d4b95d9718f0b9e2e6b9g-177}

Molar volumes are always positive, but partial molar quantities need not be. For example, the limiting partial molar volume of $\mathrm{MgSO}_{4}$ in water (its partial molar volume in the limit of zero concentration) is $-1.4 \mathrm{~cm}^{3} \mathrm{~mol}^{-1}$, which means that the addition of $1 \mathrm{~mol}_{\mathrm{MgSO}}^{4}$ to a large volume of water results in a decrease in volume of $1.4 \mathrm{~cm}^{3}$. The mixture contracts because the salt breaks up the open structure of water as the $\mathrm{Mg}^{2+}$ and $\mathrm{SO}_{4}^{2-}$ ions become hydrated, so the structure collapses slightly.

\subsection*{(b) Partial molar Gibbs energies}
The concept of a partial molar quantity can be broadened to any extensive state function. For a substance in a mixture, the chemical potential is defined as the partial molar Gibbs energy:


\begin{equation*}
\mu_{\mathrm{J}}=\left(\frac{\partial G}{\partial n_{\mathrm{J}}}\right)_{p, T, n^{\prime}} \tag{5А.4}
\end{equation*}


Chemical potential [definition]\\
where $n^{\prime}$ is used to denote that the amounts of all other components of the mixture are held constant. That is, the chemical potential is the slope of a plot of Gibbs energy against the amount of the component J , with the pressure, temperature, and the amounts of the other substances held constant (Fig. 5A.4). For a pure substance $G=n_{\mathrm{J}} G_{\mathrm{J}, \mathrm{m}}$, and from eqn 5A. 4 it follows that $\mu_{\mathrm{J}}=G_{\mathrm{J}, \mathrm{m}}$ : in this case, the chemical potential is simply the molar Gibbs energy of the substance, as is used in Topic 4A.

By the same argument that led to eqn 5A.3, it follows that the total Gibbs energy of a binary mixture is


\begin{equation*}
G=n_{\mathrm{A}} \mu_{\mathrm{A}}+n_{\mathrm{B}} \mu_{\mathrm{B}} \tag{5A.5}
\end{equation*}


\begin{center}
\includegraphics[max width=\textwidth]{2025_02_27_d4b95d9718f0b9e2e6b9g-177(2)}
\end{center}

Figure 5A. 4 The chemical potential of a substance is the slope of the total Gibbs energy of a mixture with respect to the amount of substance of interest. In general, the chemical potential varies with composition, as shown for the two values at $a$ and $b$. In this case, both chemical potentials are positive.\\
where $\mu_{\mathrm{A}}$ and $\mu_{\mathrm{B}}$ are the chemical potentials at the composition of the mixture. That is, the chemical potential of a substance, multiplied by the amount of that substance present in the mixture, is its contribution to the total Gibbs energy of the mixture. Because the chemical potentials depend on composition (and the pressure and temperature), the Gibbs energy of a mixture may change when these variables change, and for a system of components A, B, $\ldots$, eqn 3E. $7(\mathrm{~d} G=V \mathrm{~d} p-S \mathrm{~d} T)$ for a general change in $G$ becomes


\begin{align*}
& \mathrm{d} G=V \mathrm{~d} p-S \mathrm{~d} T+\mu_{\mathrm{A}} \mathrm{~d} n_{\mathrm{A}}+\mu_{\mathrm{B}} \mathrm{~d} n_{\mathrm{B}}+\cdots \\
& \text { Fundamental equation of chemical thermodynamics } \tag{5A.6}
\end{align*}


This expression is the fundamental equation of chemical thermodynamics. Its implications and consequences are explored and developed in this and the next Focus.

At constant pressure and temperature, eqn 5A. 6 simplifies to


\begin{equation*}
\mathrm{d} G=\mu_{\mathrm{A}} \mathrm{~d} n_{\mathrm{A}}+\mu_{\mathrm{B}} \mathrm{~d} n_{\mathrm{B}}+\cdots \tag{5A.7}
\end{equation*}


As established in Topic 3E, under the same conditions $\mathrm{d} G=$ $\mathrm{d} w_{\text {add,max }}$. Therefore, at constant temperature and pressure,


\begin{equation*}
\mathrm{d} w_{\mathrm{add}, \max }=\mu_{\mathrm{A}} \mathrm{~d} n_{\mathrm{A}}+\mu_{\mathrm{B}} \mathrm{~d} n_{\mathrm{B}}+\cdots \tag{5A.8}
\end{equation*}


That is, additional (non-expansion) work can arise from the changing composition of a system. For instance, in an electrochemical cell the chemical reaction is arranged to take place in two distinct sites (at the two electrodes) and the electrical work the cell performs can be traced to its changing composition as products are formed from reactants.

\subsection*{(c) The wider significance of the chemical potential}
The chemical potential does more than show how $G$ varies with composition. Because $G=U+p V-T S$, and therefore $U=-p V+T S+G$, the general form of an infinitesimal change in $U$ for a system of variable composition is

$$
\begin{aligned}
\mathrm{d} U= & -p \mathrm{~d} V-V \mathrm{~d} p+S \mathrm{~d} T+T \mathrm{~d} S+\mathrm{d} G \\
= & -p \mathrm{~d} V-V \mathrm{~d} p+S \mathrm{~d} T+T \mathrm{~d} S \\
& +\left(V \mathrm{~d} p-S \mathrm{~d} T+\mu_{\mathrm{A}} \mathrm{~d} n_{\mathrm{A}}+\mu_{\mathrm{B}} \mathrm{~d} n_{\mathrm{B}}+\cdots\right) \\
= & -p \mathrm{~d} V+T \mathrm{~d} S+\mu_{\mathrm{A}} \mathrm{~d} n_{\mathrm{A}}+\mu_{\mathrm{B}} \mathrm{~d} n_{\mathrm{B}}+\cdots
\end{aligned}
$$

This expression is the generalization of eqn 3E. 1 (that $\mathrm{d} U=$ $T \mathrm{~d} S-p \mathrm{~d} V)$ to systems in which the composition may change. It follows that at constant volume and entropy,


\begin{equation*}
\mathrm{d} U=\mu_{\mathrm{A}} \mathrm{~d} n_{\mathrm{A}}+\mu_{\mathrm{B}} \mathrm{~d} n_{\mathrm{B}}+\cdots \tag{5A.9}
\end{equation*}


and hence that


\begin{equation*}
\mu_{\mathrm{J}}=\left(\frac{\partial U}{\partial n_{\mathrm{J}}}\right)_{S, V, n^{\prime}} \tag{5A.10}
\end{equation*}


Therefore, not only does the chemical potential show how $G$ changes when the composition changes, it also shows how the internal energy changes too (but under a different set of conditions). In the same way it is possible to deduce that\\
(a) $\mu_{\mathrm{J}}=\left(\frac{\partial H}{\partial n_{\mathrm{J}}}\right)_{S, p, n^{\prime}}$\\
(b) $\mu_{\mathrm{J}}=\left(\frac{\partial A}{\partial n_{\mathrm{J}}}\right)_{T, V, n^{\prime}}$

Thus, $\mu_{\mathrm{J}}$ shows how all the extensive thermodynamic properties $U, H, A$, and $G$ depend on the composition. This is why the chemical potential is so central to chemistry.

\subsection*{(d) The Gibbs-Duhem equation}
Because the total Gibbs energy of a binary mixture is given by eqn 5A. $5\left(G=n_{\mathrm{A}} \mu_{\mathrm{A}}+n_{\mathrm{B}} \mu_{\mathrm{B}}\right)$, and the chemical potentials depend on the composition, when the compositions are changed infinitesimally the Gibbs energy of a binary system is expected to change by

$$
\mathrm{d} G=\mu_{\mathrm{A}} \mathrm{~d} n_{\mathrm{A}}+\mu_{\mathrm{B}} \mathrm{~d} n_{\mathrm{B}}+n_{\mathrm{A}} \mathrm{~d} \mu_{\mathrm{A}}+n_{\mathrm{B}} \mathrm{~d} \mu_{\mathrm{B}}
$$

However, at constant pressure and temperature the change in Gibbs energy is given by eqn 5A.7. Because $G$ is a state function, these two expressions for $\mathrm{d} G$ must be equal, which implies that at constant temperature and pressure


\begin{equation*}
n_{\mathrm{A}} \mathrm{~d} \mu_{\mathrm{A}}+n_{\mathrm{B}} \mathrm{~d} \mu_{\mathrm{B}}=0 \tag{5A.12a}
\end{equation*}


This equation is a special case of the Gibbs-Duhem equation:


\begin{equation*}
\sum_{\mathrm{J}} n_{\mathrm{J}} \mathrm{~d} \mu_{\mathrm{J}}=0 \tag{5A.12b}
\end{equation*}


Gibbs-Duhem equation\\
The significance of the Gibbs-Duhem equation is that the chemical potential of one component of a mixture cannot change independently of the chemical potentials of the other components. In a binary mixture, if one chemical potential increases, then the other must decrease, with the two changes related by eqn 5A.12a and therefore


\begin{equation*}
\mathrm{d} \mu_{\mathrm{B}}=-\frac{n_{\mathrm{A}}}{n_{\mathrm{B}}} \mathrm{~d} \mu_{\mathrm{A}} \tag{5A.13}
\end{equation*}


\subsection*{Brief illustration 5A. 1}
If the composition of a mixture is such that $n_{\mathrm{A}}=2 n_{\mathrm{B}}$, and a small change in composition results in $\mu_{\mathrm{A}}$ changing by $\Delta \mu_{\mathrm{A}}=$ $+1 \mathrm{~J} \mathrm{~mol}^{-1}, \mu_{\mathrm{B}}$ will change by

$$
\Delta \mu_{\mathrm{B}}=-2 \times\left(1 \mathrm{~J} \mathrm{~mol}^{-1}\right)=-2 \mathrm{~J} \mathrm{~mol}^{-1}
$$

The same line of reasoning applies to all partial molar quantities. For instance, changes in the partial molar volumes of the species in a mixture are related by


\begin{equation*}
\sum_{\mathrm{J}} n_{\mathrm{J}} \mathrm{~d} V_{\mathrm{J}}=0 \tag{5A.14a}
\end{equation*}


For a binary mixture,


\begin{equation*}
\mathrm{d} V_{\mathrm{B}}=-\frac{n_{\mathrm{A}}}{n_{\mathrm{B}}} \mathrm{~d} V_{\mathrm{A}} \tag{5A.14b}
\end{equation*}


As seen in Fig. 5A.1, where the partial molar volume of water increases, the partial molar volume of ethanol decreases. Moreover, as eqn 5A.14b implies, and as seen from Fig. 5A.1, a small change in the partial molar volume of A corresponds to a large change in the partial molar volume of B if $n_{\mathrm{A}} / n_{\mathrm{B}}$ is large, but the opposite is true when this ratio is small. In practice, the Gibbs-Duhem equation is used to determine the partial molar volume of one component of a binary mixture from measurements of the partial molar volume of the second component.

\subsection*{Example 5A. 2}
Using the Gibbs-Duhem equation\\
The experimental values of the partial molar volume of $\mathrm{K}_{2} \mathrm{SO}_{4}(\mathrm{aq})$ at 298 K are found to fit the expression

$$
v_{\mathrm{B}}=32.280+18.216 z^{1 / 2}
$$

where $v_{\mathrm{B}}=V_{\mathrm{K}_{2} \mathrm{SO}_{4}} /\left(\mathrm{cm}^{3} \mathrm{~mol}^{-1}\right)$ and $z$ is the numerical value of the molality of $\mathrm{K}_{2} \mathrm{SO}_{4}\left(z=b / b^{\ominus}\right.$; see The chemist's toolkit 11$)$. Use the Gibbs-Duhem equation to derive an equation for the molar volume of water in the solution. The molar volume of pure water at 298 K is $18.079 \mathrm{~cm}^{3} \mathrm{~mol}^{-1}$.

Collect your thoughts Let A denote $\mathrm{H}_{2} \mathrm{O}$, the solvent, and B denote $\mathrm{K}_{2} \mathrm{SO}_{4}$, the solute. Because the Gibbs-Duhem equation for the partial molar volumes of two components implies that $\mathrm{d} v_{\mathrm{A}}=-\left(n_{\mathrm{B}} / n_{\mathrm{A}}\right) \mathrm{d} v_{\mathrm{B}}, v_{\mathrm{A}}$ can be found by integration:

$$
v_{\mathrm{A}}=v_{\mathrm{A}}^{\star}-\int_{0}^{\nu_{\mathrm{B}}} \frac{n_{\mathrm{B}}}{n_{\mathrm{A}}} \mathrm{~d} v_{\mathrm{B}}
$$

where $v_{\mathrm{A}}^{*}=V_{\mathrm{A}}^{*} /\left(\mathrm{cm}^{3} \mathrm{~mol}^{-1}\right)$ is the numerical value of the molar volume of pure A . The first step is to change the variable of integration from $v_{\mathrm{B}}$ to $z=b / b^{\ominus}$; then integrate the right-hand side between $z=0$ (pure A) and the molality of interest.

The solution It follows from the information in the question that, with $\mathrm{B}=\mathrm{K}_{2} \mathrm{SO}_{4}, \mathrm{~d} v_{\mathrm{B}} / \mathrm{d} z=9.108 z^{-1 / 2}$. Therefore, the integration required is

$$
v_{\mathrm{A}}=v_{\mathrm{A}}^{*}-9.108 \int_{0}^{b / 6^{\ominus}} \frac{n_{\mathrm{B}}}{n_{\mathrm{A}}} z^{-1 / 2} \mathrm{~d} z
$$

The amount of $\mathrm{A}\left(\mathrm{H}_{2} \mathrm{O}\right)$ is $n_{\mathrm{A}}=(1 \mathrm{~kg}) / M_{\mathrm{A}}$, where $M_{\mathrm{A}}$ is the molar mass of water, and $n_{\mathrm{B}} /(1 \mathrm{~kg})$, which then occurs in the ratio $n_{\mathrm{B}} / n_{\mathrm{A}}$, will be recognized as the molality $b$ of B :

$$
\frac{n_{\mathrm{B}}}{n_{\mathrm{A}}}=\frac{n_{\mathrm{B}}}{(1 \mathrm{~kg}) / M_{\mathrm{A}}}=\frac{n_{\mathrm{B}} M_{\mathrm{A}}}{1 \mathrm{~kg}}=(1 \mathrm{~kg}) / M_{\mathrm{A}}=b M_{\mathrm{A}}=z b^{\ominus} M_{\mathrm{A}}
$$

Hence

$$
\begin{aligned}
v_{\mathrm{A}} & =v_{\mathrm{A}}^{*}-9.108 M_{\mathrm{A}} b^{\ominus} \int_{0}^{b / b^{\ominus}} z^{1 / 2} \mathrm{~d} z \\
& =v_{\mathrm{A}}^{*}-\frac{2}{3}\left(9.108 M_{\mathrm{A}} b^{\ominus}\right)\left(b / b^{\ominus}\right)^{3 / 2}
\end{aligned}
$$

It then follows, by substituting the data (including $M_{\mathrm{A}}=1.802 \times$ $10^{-2} \mathrm{~kg} \mathrm{~mol}^{-1}$, the molar mass of water), that

$$
V_{\mathrm{A}} /\left(\mathrm{cm}^{3} \mathrm{~mol}^{-1}\right)=18.079-0.1094\left(b / b^{\ominus}\right)^{3 / 2}
$$

The partial molar volumes are plotted in Fig. 5A.5.\\
\includegraphics[max width=\textwidth, center]{2025_02_27_d4b95d9718f0b9e2e6b9g-179(1)}

Figure 5A. 5 The partial molar volumes of the components of an aqueous solution of potassium sulfate.

Self-test 5A.2 Repeat the calculation for a salt B for which $V_{\mathrm{B}} /\left(\mathrm{cm}^{3} \mathrm{~mol}^{-1}\right)=6.218+5.146 z-7.147 z^{2}$ with $z=b / b^{\ominus}$.\\
\includegraphics[max width=\textwidth, center]{2025_02_27_d4b95d9718f0b9e2e6b9g-179}

\subsection*{5A. 2 The thermodynamics of mixing}
The dependence of the Gibbs energy of a mixture on its composition is given by eqn 5 A .5 , and, as established in Topic 3E, at constant temperature and pressure systems tend towards lower Gibbs energy. This is the link needed in order to apply thermodynamics to the discussion of spontaneous changes of composition, as in the mixing of two substances. One simple example of a spontaneous mixing process is that of two gases introduced into the same container. The mixing is spontaneous, so it must correspond to a decrease in $G$.

\subsection*{(a) The Gibbs energy of mixing of perfect gases}
Let the amounts of two perfect gases in the two containers before mixing be $n_{\mathrm{A}}$ and $n_{\mathrm{B}}$; both are at a temperature $T$ and a pressure $p$ (Fig. 5A.6). At this stage, the chemical potentials of the two gases have their 'pure' values, which are obtained by applying the definition $\mu=G_{\mathrm{m}}$ to eqn 3E. $15\left(G_{\mathrm{m}}(p)=G_{\mathrm{m}}^{\ominus}+\right.$ $\left.R T \ln \left(p / p^{\circ}\right)\right)$ :

\[
\mu=\mu^{\ominus}+R T \ln \frac{p}{p^{\ominus}} \quad \begin{align*}
& \text { Variation of chemical }  \tag{5A.15a}\\
& \text { potential with pressure } \\
& \text { [perfect gas] }
\end{align*}
\]

where $\mu^{\ominus}$ is the standard chemical potential, the chemical potential of the pure gas at 1 bar.

The notation is simplified by replacing $p / p^{\ominus}$ by $p$ itself, for eqn 5 A .15 a then becomes


\begin{equation*}
\mu=\mu^{\ominus}+R T \ln p \tag{5A.15b}
\end{equation*}


\begin{center}
\includegraphics[max width=\textwidth]{2025_02_27_d4b95d9718f0b9e2e6b9g-179(2)}
\end{center}

Figure 5A. 6 The arrangement for calculating the thermodynamic functions of mixing of two perfect gases.

\subsection*{The chemist's toolkit 11}
\subsection*{Measures of concentration}
Let A be the solvent and B the solute. The molar concentration (informally: 'molarity'), $c_{\mathrm{B}}$ or $[\mathrm{B}]$, is the amount of solute molecules (in moles) divided by the volume, $V$, of the solution:

$$
c_{\mathrm{B}}=\frac{n_{\mathrm{B}}}{V}
$$

It is commonly reported in moles per cubic decimetre $\left(\mathrm{mol} \mathrm{dm}^{-3}\right)$ or, equivalently, in moles per litre ( $\mathrm{mol} \mathrm{L}^{-1}$ ). It is convenient to define its 'standard' value as $c^{\ominus}=1 \mathrm{~mol} \mathrm{dm}^{-3}$.

The molality, $b_{\mathrm{B}}$, of a solute is the amount of solute species (in moles) in a solution divided by the total mass of the solvent (in kilograms), $m_{\mathrm{A}}$ :

$$
b_{\mathrm{B}}=\frac{n_{\mathrm{B}}}{m_{\mathrm{A}}}
$$

Both the molality and mole fraction are independent of temperature; in contrast, the molar concentration is not. It is convenient to define the 'standard' value of the molality as $b^{\ominus}=1 \mathrm{~mol} \mathrm{~kg}^{-1}$.

\subsection*{1. The relation between molality and mole fraction}
Consider a solution with one solute and having a total amount $n$ of molecules. If the mole fraction of the solute is $x_{\mathrm{B}}$, the amount of solute molecules is $n_{\mathrm{B}}=x_{\mathrm{B}} n$. The mole fraction of solvent molecules is $x_{\mathrm{A}}=1-x_{\mathrm{B}}$, so the amount of solvent molecules is $n_{\mathrm{A}}=x_{\mathrm{A}} n=\left(1-x_{\mathrm{B}}\right) n$. The mass of solvent, of molar mass $M_{\mathrm{A}}$, present is $m_{\mathrm{A}}=n_{\mathrm{A}} M_{\mathrm{A}}=\left(1-x_{\mathrm{B}}\right) n M_{\mathrm{A}}$. The molality of the solute is therefore

$$
b_{\mathrm{B}}=\frac{n_{\mathrm{B}}}{m_{\mathrm{A}}}=\frac{x_{\mathrm{B}} n}{\left(1-x_{\mathrm{B}}\right) n M_{\mathrm{A}}}=\frac{x_{\mathrm{B}}}{\left(1-x_{\mathrm{B}}\right) M_{\mathrm{A}}}
$$

The inverse of this relation, the mole fraction in terms of the molality, is

$$
x_{\mathrm{B}}=\frac{b_{\mathrm{B}} M_{\mathrm{A}}}{1+b_{\mathrm{B}} M_{\mathrm{A}}}
$$

\subsection*{2. The relation between molality and molar concentration}
The total mass of a volume $V$ of solution (not solvent) of mass density $\rho$ is $m=\rho V$. The amount of solute molecules in this volume is $n_{\mathrm{B}}=c_{\mathrm{B}} V$. The mass of solute present is $m_{\mathrm{B}}=n_{\mathrm{B}} M_{\mathrm{B}}=$ $c_{\mathrm{B}} V M_{\mathrm{B}}$. The mass of solvent present is therefore $m_{\mathrm{A}}=m-m_{\mathrm{B}}$ $=\rho V-c_{\mathrm{B}} V M_{\mathrm{B}}=\left(\rho-c_{\mathrm{B}} M_{\mathrm{B}}\right) V$. The molality is therefore

$$
b_{\mathrm{B}}=\frac{n_{\mathrm{B}}}{m_{\mathrm{A}}}=\frac{c_{\mathrm{B}} V}{\left(\rho-c_{\mathrm{B}} M_{\mathrm{B}}\right) V}=\frac{c_{\mathrm{B}}}{\rho-c_{\mathrm{B}} M_{\mathrm{B}}}
$$

The inverse of this relation, the molar concentration in terms of the molality, is

$$
c_{\mathrm{B}}=\frac{b_{\mathrm{B}} \rho}{1+b_{\mathrm{B}} M_{\mathrm{B}}}
$$

\subsection*{3. The relation between molar concentration and mole fraction}
By inserting the expression for $b_{\mathrm{B}}$ in terms of $x_{\mathrm{B}}$ into the expression for $c_{\mathrm{B}}$, the molar concentration of B in terms of its mole fraction is

$$
c_{\mathrm{B}}=\frac{x_{\mathrm{B}} \rho}{x_{\mathrm{A}} M_{\mathrm{A}}+x_{\mathrm{B}} M_{\mathrm{B}}}
$$

with $x_{\mathrm{A}}=1-x_{\mathrm{B}}$. For a dilute solution in the sense that $x_{\mathrm{B}} M_{\mathrm{B}} \ll x_{\mathrm{A}} M_{\mathrm{A}}$,

$$
c_{\mathrm{B}} \approx\left(\frac{\rho}{x_{\mathrm{A}} M_{\mathrm{A}}}\right) x_{\mathrm{B}}
$$

If, moreover, $x_{\mathrm{B}} \ll 1$, so $x_{\mathrm{A}} \approx 1$, then

$$
c_{\mathrm{B}} \approx\left(\frac{\rho}{M_{\mathrm{A}}}\right) x_{\mathrm{B}}
$$


\begin{equation*}
\Delta_{\text {mix }} G=n_{\mathrm{A}} R T \ln \frac{p_{\mathrm{A}}}{p}+n_{\mathrm{B}} R T \ln \frac{p_{\mathrm{B}}}{p} \tag{5A.16c}
\end{equation*}


At this point $n_{\mathrm{J}}$ can be replaced by $x_{\mathrm{J}} n$, where $n$ is the total amount of A and B , and the relation between partial pressure and mole fraction (Topic $1 \mathrm{~A}, p_{\mathrm{J}}=x_{\mathrm{J}} p$ ) can be used to write $p_{\mathrm{J}} / p=x_{\mathrm{J}}$ for each component. The result is

\[
\Delta_{\text {mix }} G=n R T\left(x_{\mathrm{A}} \ln x_{\mathrm{A}}+x_{\mathrm{B}} \ln x_{\mathrm{B}}\right) \quad \begin{align*}
& \text { Gibbs energy }  \tag{5A.17}\\
& \text { of mixing } \\
& \text { [perfect gas] }
\end{align*}
\]

Because mole fractions are never greater than 1, the logarithms in this equation are negative, and $\Delta_{\text {mix }} G<0$ (Fig. 5A.7). The conclusion that $\Delta_{\text {mix }} G$ is negative for all compositions confirms that perfect gases mix spontaneously in all proportions.\\
\includegraphics[max width=\textwidth, center]{2025_02_27_d4b95d9718f0b9e2e6b9g-181(2)}

Figure 5A. 7 The Gibbs energy of mixing of two perfect gases at constant temperature and pressure, and (as discussed later) of two liquids that form an ideal solution. The Gibbs energy of mixing is negative for all compositions, so perfect gases mix spontaneously in all proportions.

\subsection*{Example 5A. 3}
\subsection*{Calculating a Gibbs energy of mixing}
A container is divided into two equal compartments (Fig. 5A.8). One contains $3.0 \mathrm{~mol} \mathrm{H}_{2}(\mathrm{~g})$ at $25^{\circ} \mathrm{C}$; the other contains $1.0 \mathrm{~mol}_{2}(\mathrm{~g})$ at $25^{\circ} \mathrm{C}$. Calculate the Gibbs energy of mixing when the partition is removed. Assume that the gases are perfect.\\
\includegraphics[max width=\textwidth, center]{2025_02_27_d4b95d9718f0b9e2e6b9g-181(1)}

Figure 5A. 8 The initial and final states considered in the calculation of the Gibbs energy of mixing of gases at different initial pressures.

Collect your thoughts Equation 5A. 17 cannot be used directly because the two gases are initially at different pressures, so proceed by calculating the initial Gibbs energy from the chemical potentials. To do so, calculate the pressure of each gas: write the pressure of nitrogen as $p$, then the pressure of hydrogen as a multiple of $p$ can be found from the gas laws. Next, calculate the Gibbs energy for the system when the partition is removed. The volume occupied by each gas doubles, so its final partial pressure is half its initial pressure.

The solution Given that the pressure of nitrogen is $p$, the pressure of hydrogen is $3 p$. Therefore, the initial Gibbs energy is

$$
\begin{aligned}
G_{\mathrm{i}}= & (3.0 \mathrm{~mol})\left\{\mu^{\ominus}\left(\mathrm{H}_{2}\right)+R T \ln 3 p\right\} \\
& +(1.0 \mathrm{~mol})\left\{\mu^{\ominus}\left(\mathrm{N}_{2}\right)+R T \ln p\right\}
\end{aligned}
$$

When the partition is removed and each gas occupies twice the original volume, the final total pressure is $2 p$. The partial pressure of nitrogen falls to $\frac{1}{2} p$ and that of hydrogen falls to $\frac{3}{2} p$. Therefore, the Gibbs energy changes to

$$
\begin{aligned}
G_{\mathrm{f}}= & (3.0 \mathrm{~mol})\left\{\mu^{\ominus}\left(\mathrm{H}_{2}\right)+R T \ln \frac{3}{2} p\right\} \\
& +(1.0 \mathrm{~mol})\left\{\mu^{\ominus}\left(\mathrm{N}_{2}\right)+R T \ln \frac{1}{2} p\right\}
\end{aligned}
$$

The Gibbs energy of mixing is the difference of these two quantities:

$$
\begin{aligned}
\Delta_{\text {mix }} G & =(3.0 \mathrm{~mol}) R T \ln \frac{\frac{3}{2} p}{3 p}+(1.0 \mathrm{~mol}) R T \ln \frac{\frac{1}{2} p}{p} \\
& =-(3.0 \mathrm{~mol}) R T \ln 2-(1.0 \mathrm{~mol}) R T \ln 2 \\
& =-(4.0 \mathrm{~mol}) R T \ln 2=-6.9 \mathrm{~kJ}
\end{aligned}
$$

Comment. In this example, the value of $\Delta_{\text {mix }} G$ is the sum of two contributions: the mixing itself, and the changes in pressure of the two gases to their final total pressure, $2 p$. Do not be misled into interpreting this negative change in Gibbs energy as a sign of spontaneity: in this case, the pressure changes, and $\Delta G<0$ is a signpost of spontaneous change only at constant temperature and pressure. When $3.0 \mathrm{~mol} \mathrm{H}_{2}$ mixes with $1.0 \mathrm{~mol} \mathrm{~N}_{2}$ at the same pressure, with the volumes of the vessels adjusted accordingly, the change of Gibbs energy is -5.6 kJ . Because this value is for a change at constant pressure and temperature, the fact that it is negative does imply spontaneity.

Self-test 5 A. 3 Suppose that 2.0 mol H a 2.0 atm and $25^{\circ} \mathrm{C}$ and $4.0 \mathrm{~mol} \mathrm{~N} \mathrm{~N}_{2}$ at 3.0 atm and $25^{\circ} \mathrm{C}$ are mixed by removing the partition between them. Calculate $\Delta_{\text {mix }} G$.\\
\includegraphics[max width=\textwidth, center]{2025_02_27_d4b95d9718f0b9e2e6b9g-181}

\subsection*{(b) Other thermodynamic mixing functions}
In Topic 3E it is shown that $(\partial G / \partial T)_{p}=-S$. It follows immediately from eqn 5A. 17 that, for a mixture of perfect gases initially at the same pressure, the entropy of mixing, $\Delta_{\text {mix }} S$, is


\begin{align*}
\Delta_{\operatorname{mix}} S=-\left(\frac{\partial \Delta_{\operatorname{mix}} G}{\partial T}\right)_{p}= & n R\left(x_{\mathrm{A}} \ln x_{\mathrm{A}}+x_{\mathrm{B}} \ln x_{\mathrm{B}}\right) \\
& \begin{array}{l}
\text { Entropy of mixing } \\
\\
\\
\text { [perfect gases, constant } T \text { and } p]
\end{array}
\end{align*}


Because $\ln x<0$, it follows that $\Delta_{\text {mix }} S>0$ for all compositions (Fig. 5A.9).\\
\includegraphics[max width=\textwidth, center]{2025_02_27_d4b95d9718f0b9e2e6b9g-182(1)}

Figure 5A. 9 The entropy of mixing of two perfect gases at constant temperature and pressure, and (as discussed later) of two liquids that form an ideal solution. The entropy increases for all compositions, and because there is no transfer of heat to the surroundings when perfect gases mix, the entropy of the surroundings is unchanged. Hence, the graph also shows the total entropy of the system plus the surroundings; because the total entropy of mixing is positive at all compositions, perfect gases mix spontaneously in all proportions.

\subsection*{Brief illustration 5A. 2}
For equal amounts of perfect gas molecules that are mixed at the same pressure, set $x_{\mathrm{A}}=x_{\mathrm{B}}=\frac{1}{2}$ and obtain

$$
\Delta_{\operatorname{mix}} S=-n R\left\{\frac{1}{2} \ln \frac{1}{2}+\frac{1}{2} \ln \frac{1}{2}\right\}=n R \ln 2
$$

with $n$ the total amount of gas molecules. For 1 mol of each species, so $n=2 \mathrm{~mol}$,

$$
\Delta_{\operatorname{mix}} S=(2 \mathrm{~mol}) \times R \ln 2=+11.5 \mathrm{~J} \mathrm{~K}^{-1}
$$

An increase in entropy is expected when one gas disperses into the other and the disorder increases.

Under conditions of constant pressure and temperature, the enthalpy of mixing, $\Delta_{\text {mix }} H$, the enthalpy change accompanying mixing, of two perfect gases can be calculated from $\Delta G=$ $\Delta H-T \Delta S$. It follows from eqns 5A. 17 and 5A. 18 that

\[
\begin{array}{ll}
\Delta_{\text {mix }} H=0 & \begin{array}{l}
\text { Enthalpy of mixing } \\
\text { [perfect gases, constant } T \text { and } p]
\end{array} \tag{5A.19}
\end{array}
\]

The enthalpy of mixing is zero, as expected for a system in which there are no interactions between the molecules forming the gaseous mixture. It follows that, because the entropy of the surroundings is unchanged, the whole of the driving force for mixing comes from the increase in entropy of the system.

\subsection*{5A.3 The chemical potentials of liquids}
To discuss the equilibrium properties of liquid mixtures it is necessary to know how the Gibbs energy of a liquid varies with composition. The calculation of this dependence uses the fact that, as established in Topic 4A, at equilibrium the chemical potential of a substance present as a vapour must be equal to its chemical potential in the liquid.

\subsection*{(a) Ideal solutions}
Quantities relating to pure substances are denoted by a superscript ${ }^{*}$, so the chemical potential of pure A is written $\mu_{\mathrm{A}}^{*}$ and as $\mu_{\mathrm{A}}^{\star}(\mathrm{l})$ when it is necessary to emphasize that A is a liquid. Because the vapour pressure of the pure liquid is $p_{A}^{*}$ it follows from eqn 5A.15b that the chemical potential of A in the vapour (treated as a perfect gas) is $\mu_{\mathrm{A}}^{\ominus}+R T \ln p_{\mathrm{A}}$ (with $p_{\mathrm{A}}$ to be interpreted as $\left.p_{\mathrm{A}} / p^{\ominus}\right)$. These two chemical potentials are equal at equilibrium (Fig. 5A.10), so


\begin{equation*}
\overbrace{\mu_{\mathrm{A}}^{*}(1)}^{\text {liquid }}=\overbrace{\mu_{\mathrm{A}}^{\ominus}(\mathrm{g})+R T \ln p_{\mathrm{A}}^{*}}^{\text {vapour }} \tag{5A.20a}
\end{equation*}


If another substance, a solute, is also present in the liquid, the chemical potential of A in the liquid is changed to $\mu_{\mathrm{A}}$ and its vapour pressure is changed to $p_{\mathrm{A}}$. The vapour and solvent are still in equilibrium, so


\begin{equation*}
\mu_{\mathrm{A}}(1)=\mu_{\mathrm{A}}^{\ominus}(\mathrm{g})+R T \ln p_{\mathrm{A}} \tag{5A.20b}
\end{equation*}


\begin{center}
\includegraphics[max width=\textwidth]{2025_02_27_d4b95d9718f0b9e2e6b9g-182}
\end{center}

Figure 5A. 10 At equilibrium, the chemical potential of the gaseous form of a substance $A$ is equal to the chemical potential of its condensed phase. The equality is preserved if a solute is also present. Because the chemical potential of $A$ in the vapour depends on its partial vapour pressure, it follows that the chemical potential of liquid A can be related to its partial vapour pressure.

The next step is the combination of these two equations to eliminate the standard chemical potential of the gas, $\mu_{\mathrm{A}}^{\ominus}(\mathrm{g})$. To do so, write eqn 5A.20a as $\mu_{\mathrm{A}}^{\ominus}(\mathrm{g})=\mu_{\mathrm{A}}^{*}(1)-R T \ln p_{\mathrm{A}}^{*}$ and substitute this expression into eqn 5A. 20 b to obtain


\begin{equation*}
\mu_{\mathrm{A}}(1)=\overbrace{\mu_{\mathrm{A}}^{*}(1)-R T \ln p_{\mathrm{A}}^{\star}}^{\left.\mu_{\mathrm{A}}^{\star} \mathrm{g}\right)}+R T \ln p_{\mathrm{A}}^{\star}=\mu_{\mathrm{A}}^{\star}(1)+R T \ln \frac{p_{\mathrm{A}}}{p_{\mathrm{A}}^{*}} \tag{5A.21}
\end{equation*}


The final step draws on additional experimental information about the relation between the ratio of vapour pressures and the composition of the liquid. In a series of experiments on mixtures of closely related liquids (such as benzene and methylbenzene), François Raoult found that the ratio of the partial vapour pressure of each component to its vapour pressure when present as the pure liquid, $p_{\mathrm{A}} / p_{\mathrm{A}}^{\star}$, is approximately equal to the mole fraction of $A$ in the liquid mixture. That is, he established what is now called Raoult's law:

\[
p_{\mathrm{A}}=x_{\mathrm{A}} p_{\mathrm{A}}^{*} \quad l \begin{align*}
& \text { Raoult's law }  \tag{5A.22}\\
& \text { [ideal solution] }
\end{align*}
\]

This law is illustrated in Fig. 5A.11. Some mixtures obey Raoult's law very well, especially when the components are structurally similar (Fig. 5A.12). Mixtures that obey the law throughout the composition range from pure $A$ to pure $B$ are called ideal solutions.

\subsection*{Brief illustration 5A. 3}
The vapour pressure of pure benzene at $20^{\circ} \mathrm{C}$ is 75 Torr and that of pure methylbenzene is 25 Torr at the same temperature. In an equimolar mixture $x_{\text {benzene }}=x_{\text {methylbenzene }}=\frac{1}{2}$ so the partial vapour pressure of each one in the mixture is

$$
\begin{aligned}
& p_{\text {benzene }}=\frac{1}{2} \times 80 \text { Torr }=40 \text { Torr } \\
& p_{\text {methylbenzene }}=\frac{1}{2} \times 25 \text { Torr }=12.5 \text { Torr }
\end{aligned}
$$

The total vapour pressure of the mixture is 48 Torr. Given the two partial vapour pressures, it follows from the definition of partial pressure (Topic 1A) that the mole fractions in the vapour are

$$
x_{\text {vap,benzene }}=(40 \text { Torr }) /(48 \text { Torr })=0.83
$$

and

$$
x_{\text {vap,methylbenzene }}=(12.5 \text { Torr }) /(48 \text { Torr })=0.26
$$

The vapour is richer in the more volatile component (benzene).\\
\includegraphics[max width=\textwidth, center]{2025_02_27_d4b95d9718f0b9e2e6b9g-183}

Figure 5A. 11 The partial vapour pressures of the two components of an ideal binary mixture are proportional to the mole fractions of the components, in accord with Raoult's law. The total pressure is also linear in the mole fraction of either component.\\
\includegraphics[max width=\textwidth, center]{2025_02_27_d4b95d9718f0b9e2e6b9g-183(1)}

Figure 5A. 12 Two similar liquids, in this case benzene and methylbenzene (toluene), behave almost ideally, and the variation of their vapour pressures with composition resembles that for an ideal solution.

For an ideal solution, it follows from eqns 5A. 21 and 5A. 22 that

\[
\mu_{\mathrm{A}}(1)=\mu_{\mathrm{A}}^{*}(1)+R T \ln x_{\mathrm{A}} \quad \begin{align*}
& \text { Chemical potential }  \tag{5A.23}\\
& \text { [ideal solution] }
\end{align*}
\]

This important equation can be used as the definition of an ideal solution (so that it implies Raoult's law rather than stemming from it). It is in fact a better definition than eqn 5A. 22 because it does not assume that the vapour is a perfect gas.

The molecular origin of Raoult's law is the effect of the solute on the entropy of the solution. In the pure solvent, the molecules have a certain disorder and a corresponding entropy; the vapour pressure then represents the tendency of the system and its surroundings to reach a higher entropy. When a solute is present, the solution has a greater disorder than the pure solvent because a molecule chosen at random might or might not be a solvent molecule. Because the entropy of the solution is higher than that of the pure solvent, the solution\\
\includegraphics[max width=\textwidth, center]{2025_02_27_d4b95d9718f0b9e2e6b9g-184(2)}

Figure 5A. 13 Strong deviations from ideality are shown by dissimilar liquids (in this case carbon disulfide and acetone (propanone)). The dotted lines show the values expected from Raoult's law.\\
has a lower tendency to acquire an even higher entropy by the solvent vaporizing. In other words, the vapour pressure of the solvent in the solution is lower than that of the pure solvent.

Some solutions depart significantly from Raoult's law (Fig. 5A.13). Nevertheless, even in these cases the law is obeyed increasingly closely for the component in excess (the solvent) as it approaches purity. The law is another example of a limiting law (in this case, achieving reliability as $x_{\mathrm{A}} \rightarrow 1$ ) and is a good approximation for the properties of the solvent if the solution is dilute.

\subsection*{(b) Ideal-dilute solutions}
In ideal solutions the solute, as well as the solvent, obeys Raoult's law. However, William Henry found experimentally that, for real solutions at low concentrations, although the vapour pressure of the solute is proportional to its mole fraction, the constant of proportionality is not the vapour pressure of the pure substance (Fig. 5A.14). Henry's law is:

\[
p_{\mathrm{B}}=x_{\mathrm{B}} K_{\mathrm{B}} \quad l \begin{align*}
& \text { Henry's law }  \tag{5A.24}\\
& \text { [ideal-dilute solution] }
\end{align*}
\]

In this expression $x_{\mathrm{B}}$ is the mole fraction of the solute and $K_{\mathrm{B}}$ is an empirical constant (with the dimensions of pressure) chosen so that the plot of the vapour pressure of $B$ against its mole fraction is tangent to the experimental curve at $x_{\mathrm{B}}=0$. Henry's law is therefore also a limiting law, achieving reliability as $x_{\mathrm{B}} \rightarrow 0$.

Mixtures for which the solute B obeys Henry's law and the solvent A obeys Raoult's law are called ideal-dilute solutions. The difference in behaviour of the solute and solvent at low concentrations (as expressed by Henry's and Raoult's laws, respectively) arises from the fact that in a dilute solution the solvent molecules are in an environment very much like the one they have in the pure liquid (Fig. 5A.15). In contrast, the solute molecules are surrounded by solvent molecules,\\
\includegraphics[max width=\textwidth, center]{2025_02_27_d4b95d9718f0b9e2e6b9g-184(1)}

Figure 5A. 14 When a component (the solvent) is nearly pure, it has a vapour pressure that is proportional to the mole fraction with a slope $p_{\mathrm{B}}^{\neq}$(Raoult's law). When it is the minor component (the solute) its vapour pressure is still proportional to the mole fraction, but the constant of proportionality is now $K_{B}$ (Henry's law).\\
\includegraphics[max width=\textwidth, center]{2025_02_27_d4b95d9718f0b9e2e6b9g-184}

Figure 5A. 15 In a dilute solution, the solvent molecules (the blue spheres) are in an environment that differs only slightly from that of the pure solvent. The solute particles (the red spheres), however, are in an environment totally unlike that of the pure solute.\\
which is entirely different from their environment when it is in its pure form. Thus, the solvent behaves like a slightly modified pure liquid, but the solute behaves entirely differently from its pure state unless the solvent and solute molecules happen to be very similar. In the latter case, the solute also obeys Raoult's law.

\subsection*{Example 5A. 4 \\
 Investigating the validity of Raoult's and}
 Henry's lawsThe vapour pressures of each component in a mixture of propanone (acetone, A) and trichloromethane (chloroform, C) were measured at $35^{\circ} \mathrm{C}$ with the following results:

\begin{center}
\begin{tabular}{lcccccc}
$x_{\mathrm{C}}$ & 0 & 0.20 & 0.40 & 0.60 & 0.80 & 1 \\
$p_{\mathrm{C}} / \mathrm{kPa}$ & 0 & 4.7 & 11 & 18.9 & 26.7 & 36.4 \\
$p_{\mathrm{A}} / \mathrm{kPa}$ & 46.3 & 33.3 & 23.3 & 12.3 & 4.9 & 0 \\
\end{tabular}
\end{center}

Confirm that the mixture conforms to Raoult's law for the component in large excess and to Henry's law for the minor component. Find the Henry's law constants.

Collect your thoughts Both Raoult's and Henry's laws are statements about the form of the graph of partial vapour pressure against mole fraction. Therefore, plot the partial vapour pressures against mole fraction. Raoult's law is tested by comparing the data with the straight line $p_{\mathrm{J}}=x_{\mathrm{J}} p_{\mathrm{J}}^{*}$ for each component in the region in which it is in excess (and acting as the solvent). Henry's law is tested by finding a straight line $p_{\mathrm{J}}=x_{\mathrm{J}} K_{\mathrm{J}}$ that is tangent to each partial vapour pressure curve at low $x$, where the component can be treated as the solute.

The solution The data are plotted in Fig. 5A. 16 together with the Raoult's law lines. Henry's law requires $K_{\mathrm{A}}=24.5 \mathrm{kPa}$ for acetone and $K_{\mathrm{C}}=23.5 \mathrm{kPa}$ for chloroform.\\
\includegraphics[max width=\textwidth, center]{2025_02_27_d4b95d9718f0b9e2e6b9g-185}

Figure 5A. 16 The experimental partial vapour pressures of a mixture of chloroform (trichloromethane) and acetone (propanone) based on the data in Example 5A.4. The values of $K$ are obtained by extrapolating the dilute solution vapour pressures, as explained in the Example.

Comment. Notice how the system deviates from both Raoult's and Henry's laws even for quite small departures from $x=1$ and $x=0$, respectively. These deviations are discussed in Topic 5E.

Self-test 5A.4 The vapour pressure of chloromethane at various mole fractions in a mixture at $25^{\circ} \mathrm{C}$ was found to be as follows:

\begin{center}
\begin{tabular}{lcccc}
$x$ & 0.005 & 0.009 & 0.019 & 0.024 \\
$p / \mathrm{kPa}$ & 27.3 & 48.4 & 101 & 126 \\
\end{tabular}
\end{center}

Estimate the Henry's law constant for chloromethane.

For practical applications, Henry's law is expressed in terms of the molality, $b$, of the solute, $p_{\mathrm{B}}=b_{\mathrm{B}} K_{\mathrm{B}}$. Some Henry's law data for this convention are listed in Table 5A.1. As well as providing a link between the mole fraction of the solute and its partial pressure, the data in the table may also be used to calculate gas solubilities. Knowledge of Henry's law constants for gases in blood and fats is important for the discussion of respiration, especially when the partial pressure of oxygen is abnormal, as in diving and mountaineering, and for the discussion of the action of gaseous anaesthetics.

Table 5A. 1 Henry's law constants for gases in water at $298 \mathrm{~K}^{*}$

\begin{center}
\begin{tabular}{ll}
\hline
 & $K /\left(\mathrm{kPa} \mathrm{kg} \mathrm{mol}^{-1}\right)$ \\
\hline
$\mathrm{CO}_{2}$ & $3.01 \times 10^{3}$ \\
$\mathrm{H}_{2}$ & $1.28 \times 10^{5}$ \\
$\mathrm{~N}_{2}$ & $1.56 \times 10^{5}$ \\
$\mathrm{O}_{2}$ & $7.92 \times 10^{4}$ \\
\hline
\end{tabular}
\end{center}

${ }^{*}$ More values are given in the Resource section.

\subsection*{Brief illustration 5A. 4}
To estimate the molar solubility of oxygen in water at $25^{\circ} \mathrm{C}$ and a partial pressure of 21 kPa , its partial pressure in the atmosphere at sea level, write

$$
b_{\mathrm{O}_{2}}=\frac{p_{\mathrm{O}_{2}}}{K_{\mathrm{O}_{2}}}=\frac{21 \mathrm{kPa}}{7.9 \times 10^{4} \mathrm{kPa} \mathrm{~kg} \mathrm{~mol}^{-1}}=2.9 \times 10^{-4} \mathrm{~mol} \mathrm{~kg}^{-1}
$$

The molality of the saturated solution is therefore 0.29 mmol $\mathrm{kg}^{-1}$. To convert this quantity to a molar concentration, assume that the mass density of this dilute solution is essentially that of pure water at $25^{\circ} \mathrm{C}$, or $\rho=0.997 \mathrm{~kg} \mathrm{dm}^{-3}$. It follows that the molar concentration of oxygen is

$$
\begin{aligned}
{\left[\mathrm{O}_{2}\right] } & =b_{\mathrm{O}_{2}} \rho=\left(2.9 \times 10^{-4} \mathrm{molkg}^{-1}\right) \times\left(0.997 \mathrm{~kg} \mathrm{dm}^{-3}\right) \\
& =0.29 \mathrm{mmoldm}^{-3}
\end{aligned}
$$

\subsection*{Checklist of concepts}
\begin{enumerate}
  \item The partial molar volume of a substance is the contribution to the volume that a substance makes when it is part of a mixture.\\
$\square$ 2. The chemical potential is the partial molar Gibbs energy and is the contribution to the total Gibbs energy that a substance makes when it is part of a mixture.
  \item The chemical potential also expresses how, under a variety of different conditions, the thermodynamic functions vary with composition.4. The Gibbs-Duhem equation shows how the changes in chemical potentials (and, by extension, of other partial molar quantities) of the components of a mixture are related.5. The Gibbs energy of mixing is negative for perfect gases at the same pressure and temperature.6. The entropy of mixing of perfect gases initially at the same pressure is positive and the enthalpy of mixing is zero.
  \item Raoult's law provides a relation between the vapour pressure of a substance and its mole fraction in a mixture.8. An ideal solution is a solution that obeys Raoult's law over its entire range of compositions; for real solutions it is a limiting law valid as the mole fraction of the species approaches 1.9. Henry's law provides a relation between the vapour pressure of a solute and its mole fraction in a mixture; it is the basis of the definition of an ideal-dilute solution.10. An ideal-dilute solution is a solution that obeys Henry's law at low concentrations of the solute, and for which the solvent obeys Raoult's law.
\end{enumerate}

\subsection*{Checklist of equations}
\begin{center}
\begin{tabular}{|c|c|c|c|}
\hline
Property & Equation & Comment & Equation number \\
\hline
Partial molar volume & $V_{\mathrm{J}}=\left(\partial V / \partial n_{\mathrm{J}}\right)_{p, T, n^{\prime}}$ & Definition & 5A. 1 \\
\hline
Chemical potential & $\mu_{\mathrm{J}}=\left(\partial G / \partial n_{\mathrm{J}}\right)_{p, T, n^{\prime}}$ & Definition & 5A. 4 \\
\hline
Total Gibbs energy & $G=n_{A} \mu_{\mathrm{A}}+n_{\mathrm{B}} \mu_{\mathrm{B}}$ & Binary mixture & 5A. 5 \\
\hline
Fundamental equation of chemical thermodynamics & $\mathrm{d} G=V \mathrm{~d} p-S \mathrm{~d} T+\mu_{\mathrm{A}} \mathrm{d} n_{\mathrm{A}}+\mu_{\mathrm{B}} \mathrm{d} n_{\mathrm{B}}+\cdots$ &  & 5A. 6 \\
\hline
Gibbs-Duhem equation & $\sum_{j} n_{j} \mathrm{~d} \mu_{\mathrm{J}}=0$ &  & 5A.12b \\
\hline
Chemical potential of a gas & $\mu=\mu^{\ominus}+R T \ln \left(p / p^{\ominus}\right)$ & Perfect gas & 5A.15a \\
\hline
Gibbs energy of mixing & $\Delta_{\text {mix }} G=n R T\left(x_{\mathrm{A}} \ln x_{\mathrm{A}}+x_{\mathrm{B}} \ln x_{\mathrm{B}}\right)$ & Perfect gases and ideal solutions & 5A. 17 \\
\hline
Entropy of mixing & $\Delta_{\text {mix }} S=-n R\left(x_{\mathrm{A}} \ln x_{\mathrm{A}}+x_{\mathrm{B}} \ln x_{\mathrm{B}}\right)$ & Perfect gases and ideal solutions & 5A. 18 \\
\hline
Enthalpy of mixing & $\Delta_{\text {mix }} H=0$ & Perfect gases and ideal solutions & 5A. 19 \\
\hline
Raoult's law & $p_{\mathrm{A}}=x_{\mathrm{A}} p_{\mathrm{A}}^{*}$ & True for ideal solutions; limiting law as $x_{\mathrm{A}} \rightarrow 1$ & 5A. 22 \\
\hline
Chemical potential of component & $\mu_{\mathrm{A}}(1)=\mu_{\mathrm{A}}^{*}(1)+R T \ln x_{\mathrm{A}}$ & Ideal solution & 5A. 23 \\
\hline
Henry's law & $p_{\mathrm{B}}=x_{\mathrm{B}} K_{\mathrm{B}}$ & True for ideal-dilute solutions; limiting law as $x_{\mathrm{B}} \rightarrow 0$ & 5A. 24 \\
\hline
\end{tabular}
\end{center}

\section*{TOPIC 5B The properties of solutions}
\subsection*{Why do you need to know this material?}
Mixtures and solutions play a central role in chemistry, and so it is important to understand how their compositions affect their thermodynamic properties, such as their boiling and freezing points. One very important physical property of a solution is its osmotic pressure, which is used, for example, to determine the molar masses of macromolecules.

\subsection*{What is the key idea?}
The chemical potential of a substance in a mixture is the same in every phase in which it occurs.

\subsection*{What do you need to know already?}
This Topic is based on the expression derived from Raoult's law (Topic 5A) in which chemical potential is related to mole fraction. The derivations make use of the GibbsHelmholtz equation (Topic 3E) and the effect of pressure on chemical potential (Topic 5A). Some of the derivations are the same as those used in the discussion of the mixing of perfect gases (Topic 5A).

Thermodynamics can provide insight into the properties of liquid mixtures, and a few simple ideas can unify the whole field of study.

\subsection*{5B.1 Liquid mixtures}
The development here is based on the relation derived in Topic 5A between the chemical potential of a component (which here is called J , with $\mathrm{J}=\mathrm{A}$ or B in a binary mixture) in an ideal mixture or solution, $\mu_{\mathrm{j}}$, its value when pure, $\mu_{\mathrm{J}}^{*}$, and its mole fraction in the mixture, $x_{\mathrm{J}}$ :

\[
\mu_{\mathrm{J}}=\mu_{\mathrm{J}}^{*}+R T \ln x_{\mathrm{J}} \quad l l l \begin{align*}
& \text { Chemical potential }  \tag{5B.1}\\
& \text { [ideal solution] }
\end{align*}
\]

\subsection*{(a) Ideal solutions}
The Gibbs energy of mixing of two liquids to form an ideal solution is calculated in exactly the same way as for two\\
gases (Topic 5A). The total Gibbs energy before the liquids are mixed is


\begin{equation*}
G_{\mathrm{i}}=n_{\mathrm{A}} \mu_{\mathrm{A}}^{\star}+n_{\mathrm{B}} \mu_{\mathrm{B}}^{\star} \tag{5B.2a}
\end{equation*}


where the ${ }^{*}$ denotes the pure liquid. When they are mixed, the individual chemical potentials are given by eqn 5B. 1 and the total Gibbs energy is


\begin{equation*}
G_{\mathrm{f}}=n_{\mathrm{A}}\left(\mu_{\mathrm{A}}^{\star}+R T \ln x_{\mathrm{A}}\right)+n_{\mathrm{B}}\left(\mu_{\mathrm{B}}^{\star}+R T \ln x_{\mathrm{B}}\right) \tag{5B.2b}
\end{equation*}


Consequently, the Gibbs energy of mixing, the difference of these two quantities, is

$$
\Delta_{\mathrm{mix}} G=n R T\left(x_{\mathrm{A}} \ln x_{\mathrm{A}}+x_{\mathrm{B}} \ln x_{\mathrm{B}}\right) \quad \begin{align*}
& \text { Gibbs energy of mixing } \\
& \text { [ideal solution] }
\end{align*}
$$

where $n=n_{\mathrm{A}}+n_{\mathrm{B}}$. As for gases, it follows that the ideal entropy of mixing of two liquids is

\[
\Delta_{\mathrm{mix}} S=-n R\left(x_{\mathrm{A}} \ln x_{\mathrm{A}}+x_{\mathrm{B}} \ln x_{\mathrm{B}}\right) \quad \begin{align*}
& \text { Entropy of mixing }  \tag{5B.4}\\
& \text { [ideal solution] }
\end{align*}
\]

Then from $\Delta_{\text {mix }} G=\Delta_{\text {mix }} H-T \Delta_{\text {mix }} S$ it follows that the ideal enthalpy of mixing is zero, $\Delta_{\text {mix }} H=0$. The ideal volume of mixing, the change in volume on mixing, is also zero. To see why, consider that, because $(\partial G / \partial p)_{T}=V$ (eqn 3E.8), $\Delta_{\text {mix }} V=$ $\left(\partial \Delta_{\text {mix }} G / \partial p\right)_{T}$. But $\Delta_{\text {mix }} G$ in eqn 5 B. 3 is independent of pressure, so the derivative with respect to pressure is zero, and therefore $\Delta_{\text {mix }} V=0$.

Equations 5B. 3 and 5B. 4 are the same as those for the mixing of two perfect gases and all the conclusions drawn there are valid here: because the enthalpy of mixing is zero there is no change in the entropy of the surroundings so the driving force for mixing is the increasing entropy of the system as the molecules mingle. It should be noted, however, that solution ideality means something different from gas perfection. In a perfect gas there are no interactions between the molecules. In ideal solutions there are interactions, but the average energy of $\mathrm{A}-\mathrm{B}$ interactions in the mixture is the same as the average energy of $\mathrm{A}-\mathrm{A}$ and $\mathrm{B}-\mathrm{B}$ interactions in the pure liquids. The variation of the Gibbs energy and entropy of mixing with composition is the same as that for gases (Figs. 5A. 7 and 5A.9); both graphs are repeated here (as Figs. 5B. 1 and 5B.2).\\
\includegraphics[max width=\textwidth, center]{2025_02_27_d4b95d9718f0b9e2e6b9g-188(2)}

Figure 5B. 1 The Gibbs energy of mixing of two liquids that form an ideal solution.\\
\includegraphics[max width=\textwidth, center]{2025_02_27_d4b95d9718f0b9e2e6b9g-188}

Figure 5B. 2 The entropy of mixing of two liquids that form an ideal solution.

A note on good practice It is on the basis of this distinction that the term 'perfect gas' is preferable to the more common 'ideal gas'. In an ideal solution there are interactions, but they are effectively the same between the various species. In a perfect gas, not only are the interactions the same, but they are also zero. Few people, however, trouble to make this valuable distinction.

\subsection*{Brief illustration 5B. 1}
Consider a mixture of benzene and methylbenzene, which form an approximately ideal solution, and suppose 1.0 mol $\mathrm{C}_{6} \mathrm{H}_{6}(\mathrm{l})$ is mixed with $2.0 \mathrm{~mol} \mathrm{C}_{6} \mathrm{H}_{5} \mathrm{CH}_{3}(\mathrm{l})$. For the mixture, $x_{\text {benzene }}=0.33$ and $x_{\text {methybenzene }}=0.67$. The Gibbs energy and entropy of mixing at $25^{\circ} \mathrm{C}$, when $R T=2.48 \mathrm{~kJ} \mathrm{~mol}^{-1}$, are

$$
\begin{aligned}
\Delta_{\text {mix }} G / n & =\left(2.48 \mathrm{~kJ} \mathrm{~mol}^{-1}\right) \times(0.33 \ln 0.33+0.67 \ln 0.67) \\
& =-1.6 \mathrm{~kJ} \mathrm{~mol}^{-1} \\
\Delta_{\text {mix }} S / n & =-\left(8.3145 \mathrm{~J} \mathrm{~K}^{-1} \mathrm{~mol}^{-1}\right) \times(0.33 \ln 0.33+0.67 \ln 0.67) \\
& =+5.3 \mathrm{~J} \mathrm{~K}^{-1} \mathrm{~mol}^{-1}
\end{aligned}
$$

The enthalpy of mixing is zero (presuming that the solution is ideal).

Real solutions are composed of molecules for which the $\mathrm{A}-\mathrm{A}, \mathrm{A}-\mathrm{B}$, and $\mathrm{B}-\mathrm{B}$ interactions are all different. Not only may there be enthalpy and volume changes when such liquids mix, but there may also be an additional contribution to the entropy arising from the way in which the molecules of one type might cluster together instead of mingling freely with the others. If the enthalpy change is large and positive, or if the entropy change is negative (because of a reorganization of the molecules that results in an orderly mixture), the Gibbs energy of mixing might be positive. In that case, separation is spontaneous and the liquids are immiscible. Alternatively, the liquids might be partially miscible, which means that they are miscible only over a certain range of compositions.

\subsection*{(b) Excess functions and regular solutions}
The thermodynamic properties of real solutions are expressed in terms of the excess functions, $X^{\mathrm{E}}$, the difference between the observed thermodynamic function of mixing and the function for an ideal solution:

\[
\begin{array}{ll}
X^{\mathrm{E}}=\Delta_{\mathrm{mix}} X-\Delta_{\mathrm{mix}} X^{\text {ideal }} & \text { Excess function }  \tag{5B.5}\\
\text { [definition] }
\end{array}
\]

The excess entropy, $S^{\mathrm{E}}$, for example, is calculated by using the value of $\Delta_{\text {mix }}$ Sedeal $^{\text {given by eqn 5B.4. The excess enthalpy and }}$ volume are both equal to the observed enthalpy and volume of mixing, because the ideal values are zero in each case.

Figure 5B. 3 shows two examples of the composition dependence of excess functions. Figure 5B.3(a) shows data for\\
\includegraphics[max width=\textwidth, center]{2025_02_27_d4b95d9718f0b9e2e6b9g-188(1)}

Figure 5 B .3 Experimental excess functions at $25^{\circ} \mathrm{C}$. (a) $H^{\mathrm{E}}$ for benzene/cyclohexane; this graph shows that the mixing is endothermic (because $\Delta_{\text {mix }} H=0$ for an ideal solution). (b) The excess volume, $V^{\mathrm{E}}$, for tetrachloroethene/cyclopentane; this graph shows that there is a contraction at low tetrachloroethene mole fractions, but an expansion at high mole fractions (because $\Delta_{\text {mix }} V=0$ for an ideal mixture).\\
a benzene/cyclohexane mixture: the positive values of $H^{\mathrm{E}}$, which implies that $\Delta_{\text {mix }} H>0$, indicate that the A-B interactions in the mixture are less attractive than the $\mathrm{A}-\mathrm{A}$ and $\mathrm{B}-\mathrm{B}$ interactions in the pure liquids. The symmetrical shape of the curve reflects the similar strengths of the $\mathrm{A}-\mathrm{A}$ and $\mathrm{B}-\mathrm{B}$ interactions. Figure 5B.3(b) shows the composition dependence of the excess volume, $V^{\mathrm{E}}$, of a mixture of tetrachloroethene and cyclopentane. At high mole fractions of cyclopentane, the solution contracts as tetrachloroethene is added because the ring structure of cyclopentane results in inefficient packing of the molecules, but as tetrachloroethene is added, the molecules in the mixture pack together more tightly. Similarly, at high mole fractions of tetrachloroethene, the solution expands as cyclopentane is added because tetrachloroethene molecules are nearly flat and pack efficiently in the pure liquid, but become disrupted as the bulky ring cyclopentane is added.

Deviations of the excess enthalpy from zero indicate the extent to which the solutions are non-ideal. In this connection a useful model system is the regular solution, a solution for which $H^{\mathrm{E}} \neq 0$ but $S^{\mathrm{E}}=0$. A regular solution can be thought of as one in which the two kinds of molecules are distributed randomly (as in an ideal solution) but have different energies of interaction with each other. To express this concept more quantitatively, suppose that the excess enthalpy depends on composition as


\begin{equation*}
H^{\mathrm{E}}=n \xi R T x_{\mathrm{A}} x_{\mathrm{B}} \tag{5B.6}
\end{equation*}


where $\xi$ (xi) is a dimensionless parameter that is a measure of the energy of $\mathrm{A}-\mathrm{B}$ interactions relative to that of the $\mathrm{A}-\mathrm{A}$ and $\mathrm{B}-\mathrm{B}$ interactions. (For $H^{\mathrm{E}}$ expressed as a molar quantity, discard the $n$.) The function given by eqn 5 B .6 is plotted in Fig. 5B.4; it resembles the experimental curve in Fig. 5B.3a. If $\xi<0$, then mixing is exothermic and the $\mathrm{A}-\mathrm{B}$ interactions are more favourable than the $\mathrm{A}-\mathrm{A}$ and $\mathrm{B}-\mathrm{B}$ interactions. If $\xi>0$,\\
\includegraphics[max width=\textwidth, center]{2025_02_27_d4b95d9718f0b9e2e6b9g-189(1)}

Figure 5B. 4 The excess enthalpy according to a model in which it is proportional to $\xi x_{A} X_{B}$, for different values of the parameter $\xi$.\\
\includegraphics[max width=\textwidth, center]{2025_02_27_d4b95d9718f0b9e2e6b9g-189}

Figure 5B. 5 The Gibbs energy of mixing for different values of the parameter $\xi$.\\
then the mixing is endothermic. Because the entropy of mixing has its ideal value for a regular solution, the Gibbs energy of mixing is


\begin{align*}
\Delta_{\text {mix }} G & =\overbrace{n \xi R T x_{\mathrm{A}} x_{\mathrm{B}}}^{\Delta_{\text {mix }} H}-T \overbrace{\left[-n R\left(x_{\mathrm{A}} \ln x_{\mathrm{A}}+x_{\mathrm{B}} \ln x_{\mathrm{B}}\right)\right.}^{\Delta_{\text {mix }} \mathrm{S}}]  \tag{5B.7}\\
& =n R T\left(x_{\mathrm{A}} \ln x_{\mathrm{A}}+x_{\mathrm{B}} \ln x_{\mathrm{B}}+\xi x_{\mathrm{A}} x_{\mathrm{B}}\right)
\end{align*}


Figure 5B. 5 shows how $\Delta_{\text {mix }} G$ varies with composition for different values of $\xi$. The important feature is that for $\xi>2$ the graph shows two minima separated by a maximum. The implication of this observation is that, provided $\xi>2$, the system will separate spontaneously into two phases with compositions corresponding to the two minima, because such a separation corresponds to a reduction in Gibbs energy. This point is developed in Topic 5C.

\subsection*{Example 5B. 1 Identifying the parameter for a \\
 regular solution}
Identify the value of the parameter $\xi$ that would be appropriate to model a mixture of benzene and cyclohexane at $25^{\circ} \mathrm{C}$, and estimate the Gibbs energy of mixing for an equimolar mixture.\\
Collect your thoughts Refer to Fig. 5B.3a and identify the value of the maximum in the curve; then relate it to eqn 5B. 6 written as a molar quantity $\left(H^{\mathrm{E}}=\xi R T x_{\mathrm{A}} x_{\mathrm{B}}\right)$. For the second part, assume that the solution is regular and that the Gibbs energy of mixing is given by eqn 5B.7.

The solution In the experimental data the maximum occurs close to $x_{\mathrm{A}}=x_{\mathrm{B}}=\frac{1}{2}$ and its value is close to $701 \mathrm{Jmol}^{-1}$. It follows that

$$
\begin{aligned}
\xi & =\frac{H^{\mathrm{E}}}{R T x_{\mathrm{A}} x_{\mathrm{B}}}=\frac{701 \mathrm{~J} \mathrm{~mol}^{-1}}{\left(8.3145 \mathrm{JK}^{-1} \mathrm{~mol}^{-1}\right) \times(298 \mathrm{~K}) \times \frac{1}{2} \times \frac{1}{2}} \\
& =1.13
\end{aligned}
$$

The total Gibbs energy of mixing to achieve the stated composition (provided the solution is regular) is therefore

$$
\begin{aligned}
\Delta_{\text {mix }} G / n & =\frac{1}{2} R T \ln \frac{1}{2}+\frac{1}{2} R T \ln \frac{1}{2}+701 \mathrm{~J} \mathrm{~mol}^{-1} \\
& =-R T \ln 2+701 \mathrm{~J} \mathrm{~mol}^{-1} \\
& =-1.72 \mathrm{~kJ} \mathrm{~mol}^{-1}+0.701 \mathrm{~kJ} \mathrm{~mol}^{-1}=-1.02 \mathrm{~kJ} \mathrm{~mol}^{-1}
\end{aligned}
$$

Self-test 5B. 1 The graph in Fig. 5B.3a suggests the following values:

\begin{center}
\begin{tabular}{lccccccccc}
$x$ & 0.1 & 0.2 & 0.3 & 0.4 & 0.5 & 0.6 & 0.7 & 0.8 & 0.9 \\
$H^{\mathrm{E}} /\left(\mathrm{J} \mathrm{mol}^{-1}\right)$ & 150 & 350 & 550 & 680 & 700 & 690 & 600 & 500 & 280 \\
\end{tabular}
\end{center}

Use a curve-fitting procedure to fit these data to an expression of the form in eqn 5B. 6 written as $H^{\mathrm{E}} / n=A x(1-x)$.\\
\includegraphics[max width=\textwidth, center]{2025_02_27_d4b95d9718f0b9e2e6b9g-190}

\subsection*{5B.2 Colligative properties}
A colligative property is a physical property that depends on the relative number of solute particles present but not their chemical identity ('colligative' denotes 'depending on the collection'). They include the lowering of vapour pressure, the elevation of boiling point, the depression of freezing point, and the osmotic pressure arising from the presence of a solute. In dilute solutions these properties depend only on the number of solute particles present, not their identity.\\
In this development, the solvent is denoted by A and the solute by B. There are two assumptions. First, the solute is not volatile, so it does not contribute to the vapour. Second, the solute does not dissolve in the solid solvent: that is, the pure solid solvent separates when the solution is frozen. The latter assumption is quite drastic, although it is true of many mixtures; it can be avoided at the expense of more algebra, but that introduces no new principles.

\subsection*{(a) The common features of colligative properties}
All the colligative properties stem from the reduction of the chemical potential of the liquid solvent as a result of the presence of solute. For an ideal solution (one that obeys Raoult's law, Topic $5 \mathrm{~A} ; p_{\mathrm{A}}=x_{\mathrm{A}} p_{\mathrm{A}}^{*}$ ), the reduction is from $\mu_{\mathrm{A}}^{*}$ for the pure solvent to $\mu_{\mathrm{A}}=\mu_{\mathrm{A}}^{\star}+R T \ln x_{\mathrm{A}}$ when a solute is present $\left(\ln x_{\mathrm{A}}\right.$ is negative because $\left.x_{\mathrm{A}}<1\right)$. There is no direct influence of the solute on the chemical potential of the solvent vapour and the solid solvent because the solute appears in neither the vapour nor the solid. As can be seen from Fig. 5B.6, the reduction in chemical potential of the solvent implies that the liquid-vapour equilibrium occurs at a higher\\
\includegraphics[max width=\textwidth, center]{2025_02_27_d4b95d9718f0b9e2e6b9g-190(1)}

Figure 5B. 6 The chemical potential of the liquid solvent in a solution is lower than that of the pure liquid. As a result, the temperature at which the chemical potential of the solvent is equal to that of the solid solvent (the freezing point) is lowered, and the temperature at which it is equal to the vapour (the boiling point) is raised. The lowering of the liquid's chemical potential has a greater effect on the freezing point than on the boiling point because of the angles at which the lines intersect.\\
temperature (the boiling point is raised) and the solid-liquid equilibrium occurs at a lower temperature (the freezing point is lowered).\\
The molecular origin of the lowering of the chemical potential is not the energy of interaction of the solute and solvent particles, because the lowering occurs even in an ideal solution (for which the enthalpy of mixing is zero). If it is not an enthalpy effect, it must be an entropy effect. ${ }^{1}$ When a solute is present, there is an additional contribution to the entropy of the solvent which results is a weaker tendency to form the vapour (Fig. 5B.7). This weakening of the tendency to form a vapour lowers the vapour pressure and hence raises the boiling point. Similarly, the enhanced molecular randomness of the solution opposes the tendency to freeze. Consequently, a lower temperature must be reached before equilibrium between solid and solution is achieved. Hence, the freezing point is lowered.

The strategy for the quantitative discussion of the elevation of boiling point and the depression of freezing point is to look for the temperature at which, at 1 atm , one phase (the pure solvent vapour or the pure solid solvent) has the same chemical potential as the solvent in the solution. This is the new equilibrium temperature for the phase transition at 1 atm , and hence corresponds to the new boiling point or the new freezing point of the solvent.

\footnotetext{${ }^{1}$ More precisely, if it is not an enthalpy effect (that is, an effect arising from changes in the entropy of the surroundings due to the transfer of energy as heat into or from them), then it must be an effect arising from the entropy of the system.
}
\includegraphics[max width=\textwidth, center]{2025_02_27_d4b95d9718f0b9e2e6b9g-191(1)}

Figure 5B.7 The vapour pressure of a pure liquid represents a balance between the increase in disorder arising from vaporization and the decrease in disorder of the surroundings. (a) Here the structure of the liquid is represented highly schematically by the grid of squares. (b) When solute (the dark green squares) is present, the disorder of the condensed phase is higher than that of the pure liquid, and there is a decreased tendency to acquire the disorder characteristic of the vapour.

\subsection*{(b) The elevation of boiling point}
The equilibrium of interest when considering boiling is between the solvent vapour and the solvent in solution at 1 atm (Fig. 5B.8). The equilibrium is established at a temperature for which


\begin{equation*}
\mu_{\mathrm{A}}^{\star}(\mathrm{g})=\mu_{\mathrm{A}}^{*}(\mathrm{l})+R T \ln x_{\mathrm{A}} \tag{5B.8}
\end{equation*}


where $\mu_{\mathrm{A}}^{*}(\mathrm{~g})$ is the chemical potential of the pure vapour; the pressure of 1 atm is the same throughout, and will not be written explicitly. It can be shown that a consequence of this relation is that the normal boiling point of the solvent is raised and that in a dilute solution the increase is proportional to the mole fraction of solute.\\
\includegraphics[max width=\textwidth, center]{2025_02_27_d4b95d9718f0b9e2e6b9g-191}

Figure 5B. 8 The equilibrium involved in the calculation of the elevation of boiling point is between $A$ present as pure vapour and $A$ in the mixture, $A$ being the solvent and $B$ a non-volatile solute.

How is that done? 5B. 1\\
Deriving an expression for the elevation of the boiling point

The starting point for the calculation is the equality of the chemical potentials of the solvent in the liquid and vapour phases, eqn 5B.8. The strategy then involves examining how the temperature must be changed to maintain that equality when solute is added. You need to follow these steps.

Step 1 Relate $\ln x_{\mathrm{A}}$ to the Gibbs energy of vaporization\\
Equation 5B. 8 can be rearranged into

$$
\ln x_{\mathrm{A}}=\frac{\mu_{\mathrm{A}}^{*}(\mathrm{~g})-\mu_{\mathrm{A}}^{*}(\mathrm{l})}{R T}=\frac{\Delta_{\text {vap }} G}{R T}
$$

where $\Delta_{\text {vap }} G$ is the (molar) Gibbs energy of vaporization of the pure solvent (A).\\
Step 2 Write an expression for the variation of $\ln x_{\mathrm{A}}$ with temperature\\
Differentiating both sides of the expression from Step 1 with respect to temperature and using the Gibbs-Helmholtz equation (Topic 3E, $\left.(\partial(G / T) / \partial T)_{p}=-H / T^{2}\right)$ to rewrite the term on the right gives

$$
\frac{\mathrm{d} \ln x_{\mathrm{A}}}{\mathrm{~d} T}=\frac{1}{R} \frac{\mathrm{~d}\left(\Delta_{\mathrm{vap}} G / T\right)}{\mathrm{d} T}=-\frac{\Delta_{\mathrm{vap}} H}{R T^{2}}
$$

The change in temperature $\mathrm{d} T$ needed to maintain equilibrium when solute is added and the change in $\ln x_{\mathrm{A}}$ by $\mathrm{d} \ln x_{\mathrm{A}}$ are therefore related by

$$
\mathrm{d} \ln x_{\mathrm{A}}=-\frac{\Delta_{\mathrm{vap}} H}{R T^{2}} \mathrm{~d} T
$$

Step 3 Find the relation between the measurable changes in $\ln x_{\mathrm{A}}$ and $T$ by integration\\
To integrate the preceding expression, integrate from $x_{\mathrm{A}}=1$, corresponding to $\ln x_{\mathrm{A}}=0$ (and when $T=T^{*}$, the boiling point of pure A) to $x_{\mathrm{A}}$ (when the boiling point is $T$ ). As usual, to avoid confusing the variables of integration with the final value they reach, replace $\ln x_{\mathrm{A}}$ by $\ln x_{\mathrm{A}}^{\prime}$ and $T$ by $T^{\prime}$ :

$$
\int_{0}^{\ln x_{A}} \mathrm{~d} \ln x_{\mathrm{A}}^{\prime}=-\frac{1}{R} \int_{T^{*}}^{T} \frac{\Delta_{\text {vap }} H}{T^{\prime 2}} \mathrm{~d} T^{\prime}
$$

The left-hand side integrates to $\ln x_{\mathrm{A}}$, which is equal to $\ln \left(1-x_{\mathrm{B}}\right)$. The right-hand side can be integrated if the enthalpy of vaporization is assumed to be constant over the small range of temperatures involved, so can be taken outside the integral:

$$
\ln \left(1-x_{\mathrm{B}}\right)=-\frac{\Delta_{\text {vap }} H}{R} \overbrace{\int_{T^{*}}^{T} \frac{1}{T^{\prime 2}} \mathrm{~d} T^{\prime}}^{\text {Integral A.1 }}
$$

\subsection*{Therefore}
$$
\ln \left(1-x_{\mathrm{B}}\right)=\frac{\Delta_{\mathrm{vap}} H}{R}\left(\frac{1}{T}-\frac{1}{T^{*}}\right)
$$

\subsection*{Step 4 Approximate the expression for dilute solutions}
Suppose that the amount of solute present is so small that $x_{\mathrm{B}} \ll 1$; the approximation $\ln (1-x) \approx-x$ (The chemist's toolkit 12) can then be used. It follows that

$$
x_{\mathrm{B}}=\frac{\Delta_{\text {vap }} H}{R}\left(\frac{1}{T^{*}}-\frac{1}{T}\right)
$$

Finally, because the increase in the boiling point is small, $T \approx T^{*}$, it also follows that

$$
\frac{1}{T^{\star}}-\frac{1}{T}=\frac{T-T^{\star}}{T T^{\star}} \approx \frac{T-T^{\star}}{T^{* 2}}=\frac{\Delta T_{\mathrm{b}}}{T^{*^{2}}}
$$

with $\Delta T_{\mathrm{b}}=T-T^{*}$. The previous equation then becomes


\begin{equation*}
x_{\mathrm{B}}=\frac{\Delta_{\mathrm{vap}} H}{R} \times \frac{\Delta T_{\mathrm{b}}}{T^{\star^{2}}} \tag{5B.9a}
\end{equation*}


which confirms that the elevation of boiling point and the mole fraction of solute are proportional to each other.

\subsection*{Step 5 Rearrange the expression}
The calculation has shown that the presence of a solute at a mole fraction $x_{\mathrm{B}}$ causes an increase in normal boiling point\\
from $T^{*}$ to $T^{*}+\Delta T$, and after minor rearrangement of eqn 5B.9a the relation is\\
$\Delta T_{\mathrm{b}}=K x_{\mathrm{B}} \quad K=\frac{R T^{\not{ }^{2}}}{\Delta_{\text {vap }} H} \quad \begin{aligned} & \text { (5B.9b) } \\ & \\ & \text { Elevation of boiling point } \\ & \text { lideal solution] }\end{aligned}$\\
Because eqn 5B.9b makes no reference to the identity of the solute, only to its mole fraction, it follows that the elevation of boiling point is a colligative property. The value of $\Delta T$ does depend on the properties of the solvent, and the biggest changes occur for solvents with high boiling points. By Trouton's rule (Topic 3B), $\Delta_{\text {vap }} H / T^{*}$ is a constant; therefore eqn 5 B. 9 b has the form $\Delta T \propto T^{*}$ and is independent of $\Delta_{\text {vap }} H$ itself. If $x_{\mathrm{B}} \ll 1$ it follows that the mole fraction of B is proportional to its molality, $b$ (see The chemist's toolkit 11 in Topic 5A). Equation 5B.9b can therefore be written as

$$
\Delta T_{\mathrm{b}}=K_{\mathrm{b}} b
$$

Boiling point elevation [empirical relation]\\
(5B.9c)\\
where $K_{\mathrm{b}}$ is the empirical boiling-point constant of the solvent (Table 5B.1).

Table 5B. 1 Freezing-point $\left(K_{\mathrm{f}}\right)$ and boiling-point $\left(K_{\mathrm{b}}\right)$ constants*

\begin{center}
\begin{tabular}{lcc}
\hline
 & $K_{\mathrm{f}} /\left(\mathrm{K} \mathrm{kg} \mathrm{mol}^{-1}\right)$ & $K_{\mathrm{b}} /\left(\mathrm{K} \mathrm{kg} \mathrm{mol}^{-1}\right)$ \\
\hline
Benzene & 5.12 & 2.53 \\
Camphor & 40 &  \\
Phenol & 7.27 & 3.04 \\
Water & 1.86 & 0.51 \\
\hline
\end{tabular}
\end{center}

\begin{itemize}
  \item More values are given in the Resource section.
\end{itemize}

\subsection*{The chemist's toolkit 12 Series expansions}
A function $f(x)$ can be expressed in terms of its value in the vicinity of $x=a$ by using the Taylor series

$$
\begin{aligned}
f(x) & =f(a)+\left(\frac{\mathrm{d} f}{\mathrm{~d} x}\right)_{a}(x-a)+\frac{1}{2!}\left(\frac{\mathrm{d}^{2} f}{\mathrm{~d} x^{2}}\right)_{a}(x-a)^{2}+\cdots \\
& =\sum_{n=0}^{\infty} \frac{1}{n!}\left(\frac{\mathrm{d}^{n} f}{\mathrm{~d} x^{n}}\right)_{a}(x-a)^{n}
\end{aligned}
$$

Taylor series\\
where the notation $(\ldots)_{a}$ means that the derivative is evaluated at $x=a$ and $n!$ denotes a factorial defined as

$$
n!=n(n-1)(n-2) \ldots 1, \quad 0!\equiv 1
$$

Factorial\\
The Maclaurin series for a function is a special case of the Taylor series in which $a=0$. The following Maclaurin series are used at various stages in the text:

$$
(1+x)^{-1}=1-x+x^{2}-\cdots=\sum_{n=0}^{\infty}(-1)^{n} x^{n}
$$

$$
\begin{aligned}
& \mathrm{e}^{x}=1+x+\frac{1}{2} x^{2}+\cdots=\sum_{n=0}^{\infty} \frac{x^{n}}{n!} \\
& \ln (1+x)=x-\frac{1}{2} x^{2}+\frac{1}{3} x^{3}-\cdots=\sum_{n=1}^{\infty}(-1)^{n+1} \frac{x^{n}}{n}
\end{aligned}
$$

Series expansions are used to simplify calculations, because when $|x| \ll 1$ it is possible, to a good approximation, to terminate the series after one or two terms. Thus, provided $|x| \ll 1$,

$$
\begin{aligned}
& (1+x)^{-1} \approx 1-x \\
& \mathrm{e}^{x} \approx 1+x \\
& \ln (1+x) \approx x
\end{aligned}
$$

A series is said to converge if the sum approaches a finite, definite value as $n$ approaches infinity. If it does not, the series is said to diverge. Thus, the series expansion of $(1+x)^{-1}$ converges for $|x|<1$ and diverges for $|x| \geq 1$. Tests for convergence are explained in mathematical texts.

\subsection*{Brief illustration 5B. 2}
The boiling-point constant of water is $0.51 \mathrm{Kkg} \mathrm{mol}^{-1}$, so a solute present at a molality of $0.10 \mathrm{~mol} \mathrm{~kg}^{-1}$ would result in an elevation of boiling point of only 0.051 K . The boiling-point constant of benzene is significantly larger, at $2.53 \mathrm{Kkg} \mathrm{mol}^{-1}$, so the elevation would be 0.25 K .

\subsection*{(c) The depression of freezing point}
The equilibrium now of interest is between pure solid solvent A and the solution with solute present at a mole fraction $x_{B}$ (Fig. 5B.9). At the freezing point, the chemical potentials of A in the two phases are equal:


\begin{equation*}
\mu_{\mathrm{A}}^{\star}(\mathrm{s})=\mu_{\mathrm{A}}^{\not}(\mathrm{l})+R T \ln x_{\mathrm{A}} \tag{5B.10}
\end{equation*}


where $\mu_{\mathrm{A}}^{*}(\mathrm{~s})$ is the chemical potential of pure solid A . The only difference between this calculation and the last is the appearance of the chemical potential of the solid in place of that of the vapour. Therefore the result can be written directly from eqn 5 B. 9 b :


\begin{equation*}
\Delta T_{\mathrm{f}}=K^{\prime} x_{\mathrm{B}} \quad K^{\prime}=\frac{R T^{\not{ }^{2}}}{\Delta_{\text {fus }} H} \quad \text { Freezing point depression } \tag{5B.11}
\end{equation*}


where $T^{*}$ is the freezing point of the pure liquid, $\Delta T_{\mathrm{f}}$ is the freezing point depression, $T^{*}-T$, and $\Delta_{\text {fus }} H$ is the enthalpy of fusion of the solvent. Larger depressions are observed in solvents with low enthalpies of fusion and high melting points. When the solution is dilute, the mole fraction is proportional to the molality of the solute, $b$, and it is common to write the last equation as

\[
\Delta T_{\mathrm{f}}=K_{\mathrm{f}} b \quad \begin{align*}
& \text { Freezing point depression }  \tag{5B.12}\\
& \text { [empirical relation] }
\end{align*}
\]

where $K_{\mathrm{f}}$ is the empirical freezing-point constant (Table 5B.1).\\
\includegraphics[max width=\textwidth, center]{2025_02_27_d4b95d9718f0b9e2e6b9g-193}

Figure 5B.9 The equilibrium involved in the calculation of the lowering of freezing point is between $A$ present as pure solid and $A$ in the mixture, $A$ being the solvent and $B$ a solute that is insoluble in solid A .

\subsection*{Brief illustration 5B.3}
The freezing-point constant of water is $1.86 \mathrm{Kkg} \mathrm{mol}^{-1}$, so a solute present at a molality of 0.10 mol kg -1 would result in a depression of freezing point of only 0.19 K . The freezing-point constant of camphor is significantly larger, at $40 \mathrm{Kkg} \mathrm{mol}^{-1}$, so the depression would be 4.0 K .

\subsection*{(d) Solubility}
Although solubility is not a colligative property (because solubility varies with the identity of the solute), it may be estimated in a similar way. When a solid solute is left in contact with a solvent, it dissolves until the solution is saturated. Saturation is a state of equilibrium, with the undissolved solute in equilibrium with the dissolved solute. Therefore, in a saturated solution the chemical potential of the pure solid solute, $\mu_{\mathrm{B}}^{*}(\mathrm{~s})$, and the chemical potential of B in solution, $\mu_{\mathrm{B}}$, are equal (Fig. 5B.10). Because the latter is related to the mole fraction in the solution by $\mu_{\mathrm{B}}=\mu_{\mathrm{B}}^{\star}(\mathrm{l})+R T \ln x_{\mathrm{B}}$, it follows that


\begin{equation*}
\mu_{\mathrm{B}}^{\nsucc}(\mathrm{s})=\mu_{\mathrm{B}}^{\ngtr}(\mathrm{l})+R T \ln x_{\mathrm{B}} \tag{5B.13}
\end{equation*}


This expression is the same as the starting equation of the last section, except that the quantities refer to the solute $B$, not the solvent $A$. It can be used in a similar way to derive the relation between the solubility and the temperature.

\subsection*{How is that done? 5B. 2 Deriving a relation between the}
 solubility and the temperatureIn the present case, the goal is to find the mole fraction of $B$ in solution at equilibrium when the temperature is $T$. Therefore, start by rearranging eqn 5B. 13 into

$$
\ln x_{\mathrm{B}}=\frac{\mu_{\mathrm{B}}^{\star}(\mathrm{s})-\mu_{\mathrm{B}}^{\star}(\mathrm{l})}{R T}=-\frac{\Delta_{\mathrm{fus}} G}{R T}
$$

As in the derivation of eqn 5B.9, differentiate both side of this equation with respect to $T$ to relate the change in composition\\
\includegraphics[max width=\textwidth, center]{2025_02_27_d4b95d9718f0b9e2e6b9g-193(1)}

Figure 5B. 10 The equilibrium involved in the calculation of the solubility is between pure solid $B$ and $B$ in the mixture.\\
to the change in temperature, and use the Gibbs-Helmholtz equation. Then integrate the resulting expression from the melting temperature of B (when $x_{\mathrm{B}}=1$ and $\ln x_{\mathrm{B}}=0$ ) to the temperature of interest (when $x_{\mathrm{B}}$ has a value between 0 and 1 ):

$$
\int_{0}^{\ln x_{\mathrm{B}}} \mathrm{~d} \ln x_{\mathrm{B}}^{\prime}=\frac{1}{R} \int_{T_{\mathrm{f}}}^{T} \frac{\Delta_{\mathrm{fus}} H}{T^{\prime 2}} \mathrm{~d} T^{\prime}
$$

where $\Delta_{\text {fus }} H$ is the enthalpy of fusion of the solute and $T_{\mathrm{f}}$ is its melting point.

In the final step, suppose that the enthalpy of fusion of B is constant over the range of temperatures of interest, and take it outside the integral. The result of the calculation is then\\
$\left.-\ln x_{\mathrm{B}}=\frac{\Delta_{\text {fus }} H}{R}\left(\frac{1}{T_{\mathrm{f}}}-\frac{1}{T}\right) \right\rvert\,$ (5B.14)\\
This equation is plotted in Fig. 5B.11. It shows that the solubility of B decreases as the temperature is lowered from its melting point. The illustration also shows that solutes with high melting points and large enthalpies of melting have low solubilities at normal temperatures. However, the detailed content of eqn 5B. 14 should not be treated too seriously because it is based on highly questionable approximations, such as the ideality of the solution. One aspect of its approximate character is that it fails to predict that solutes will have different solubilities in different solvents, for no solvent properties appear in the expression.\\
\includegraphics[max width=\textwidth, center]{2025_02_27_d4b95d9718f0b9e2e6b9g-194(2)}

Figure 5B. 11 The variation of solubility, the mole fraction of solute in a saturated solution, with temperature; $T^{*}$ is the freezing temperature of the solute. Individual curves are labelled with the value of $\Delta_{\mathrm{fus}} H / R T^{*}$.

\subsection*{Brief illustration 5B. 4}
The ideal solubility of naphthalene in benzene is calculated from eqn 5B. 14 by noting that the enthalpy of fusion of naphthalene is $18.80 \mathrm{~kJ} \mathrm{~mol}^{-1}$ and its melting point is 354 K . Then, at $20^{\circ} \mathrm{C}$,

$$
\ln x_{\text {naphthalene }}=\frac{1.880 \times 10^{4} \mathrm{Jmol}^{-1}}{8.3145 \mathrm{JK}^{-1} \mathrm{~mol}^{-1}}\left(\frac{1}{354 \mathrm{~K}}-\frac{1}{293 \mathrm{~K}}\right)=-1.32 \ldots
$$

and therefore $x_{\text {naphthalene }}=0.26$. This mole fraction corresponds to a molality of $4.5 \mathrm{~mol} \mathrm{~kg}^{-1}$ ( 580 g of naphthalene in 1 kg of benzene).

\subsection*{(e) Osmosis}
The phenomenon of osmosis (from the Greek word for 'push') is the spontaneous passage of a pure solvent into a solution separated from it by a semipermeable membrane, a membrane permeable to the solvent but not to the solute (Fig. 5B.12). The osmotic pressure, $\Pi$ (uppercase pi), is the pressure that must be applied to the solution to stop the influx of solvent. Important examples of osmosis include transport of fluids through cell membranes, dialysis, and osmometry, the determination of molar mass by the measurement of osmotic pressure. Osmometry is widely used to determine the molar masses of macromolecules.

In the simple arrangement shown in Fig. 5B.13, the opposing pressure arises from the column of solution that the osmosis itself produces. Equilibrium is reached when the pressure\\
\includegraphics[max width=\textwidth, center]{2025_02_27_d4b95d9718f0b9e2e6b9g-194(1)}

Figure 5B. 12 The equilibrium involved in the calculation of osmotic pressure, $\Pi$, is between pure solvent $A$ at a pressure $p$ on one side of the semipermeable membrane and $A$ as a component of the mixture on the other side of the membrane, where the pressure is $p+\Pi$.\\
\includegraphics[max width=\textwidth, center]{2025_02_27_d4b95d9718f0b9e2e6b9g-194}

Figure 5B. 13 In a simple version of the osmotic pressure experiment, $A$ is at equilibrium on each side of the membrane when enough has passed into the solution to cause a hydrostatic pressure difference.\\
due to that column matches the osmotic pressure. The complicating feature of this arrangement is that the entry of solvent into the solution results in its dilution, and so it is more difficult to treat than the arrangement in Fig. 5B.12, in which there is no flow and the concentrations remain unchanged.

The thermodynamic treatment of osmosis depends on noting that, at equilibrium, the chemical potential of the solvent must be the same on each side of the membrane. The chemical potential of the solvent is lowered by the solute, but is restored to its 'pure' value by the application of pressure. The challenge in this instance is to show that, provided the solution is dilute, the extra pressure to be exerted is proportional to the molar concentration of the solute in the solution.

\subsection*{How is that done? 5B. 3}
Deriving a relation between the osmotic pressure and the molar concentration of solute

On the pure solvent side the chemical potential of the solvent, which is at a pressure $p$, is $\mu_{\mathrm{A}}^{*}(p)$. On the solution side, the chemical potential is lowered by the presence of the solute, which reduces the mole fraction of the solvent from 1 to $x_{\mathrm{A}}$. However, the chemical potential of A is raised on account of the greater pressure, $p+\Pi$, that the solution experiences. Now follow these steps, and be prepared to make a number of approximations by supposing that the solution is dilute ( $x_{\mathrm{B}} \ll 1$ ).

Step 1 Write an expression for the chemical potential of the solvent in the solution\\
At equilibrium the chemical potential of A is the same in both compartments:

$$
\mu_{\mathrm{A}}^{*}(p)=\mu_{\mathrm{A}}\left(x_{\mathrm{A}}, p+\Pi\right)
$$

The presence of solute is taken into account in the normal way by using eqn 5B.1:

$$
\mu_{\mathrm{A}}\left(x_{\mathrm{A}}, p+\Pi\right)=\mu_{\mathrm{A}}^{*}(p+\Pi)+R T \ln x_{\mathrm{A}}
$$

By combining these two expressions it follows that

$$
\mu_{\mathrm{A}}^{\star}(p)=\mu_{\mathrm{A}}^{\star}(p+\Pi)+R T \ln x_{\mathrm{A}}
$$

and therefore

$$
\mu_{\mathrm{A}}^{\star}(p+\Pi)=\mu_{\mathrm{A}}^{\star}(p)-R T \ln x_{\mathrm{A}}
$$

Step 2 Evaluate the effect of pressure on the chemical potential of the solvent\\
The effect of pressure is taken into account by using eqn 3E.12b,

$$
G_{\mathrm{m}}\left(p_{\mathrm{f}}\right)=G_{\mathrm{m}}\left(p_{\mathrm{i}}\right)+\int_{p_{\mathrm{i}}}^{p_{\mathrm{f}}} V_{\mathrm{m}} \mathrm{~d} p
$$

written as

$$
\mu_{\mathrm{A}}^{*}(p+\Pi)=\mu_{\mathrm{A}}^{*}(p)+\int_{p}^{p+\Pi} V_{\mathrm{m}} \mathrm{~d} p
$$

where $V_{\mathrm{m}}$ is the molar volume of the pure solvent A. On substituting $\mu_{\mathrm{A}}^{*}(p+\Pi)=\mu_{\mathrm{A}}^{*}(p)-R T \ln x_{\mathrm{A}}$ into this expression and cancelling the $\mu_{\mathrm{A}}^{*}(p)$, it follows that


\begin{equation*}
-R T \ln x_{\mathrm{A}}=\int_{p}^{p+\Pi} V_{\mathrm{m}} \mathrm{~d} p \tag{5B.15}
\end{equation*}


\subsection*{Step 3 Evaluate the integral}
Suppose that the pressure range in the integration is so small that the molar volume of the solvent is a constant. Then the right-hand side of eqn 5B. 15 simplifies to

$$
\int_{p}^{p+\Pi} V_{\mathrm{m}} \mathrm{~d} p=V_{\mathrm{m}} \int_{p}^{p+\Pi} \mathrm{d} p=V_{\mathrm{m}} \Pi
$$

which implies that

$$
-R T \ln x_{\mathrm{A}}=V_{\mathrm{m}} \Pi
$$

On the left-hand side of this expression, $\ln x_{\mathrm{A}}$ may be replaced by $\ln \left(1-x_{\mathrm{B}}\right)$, and if it is assumed that the solution is dilute $\ln \left(1-x_{\mathrm{B}}\right) \approx-x_{\mathrm{B}}$ (The chemist's toolkit 12), then

$$
-R T \ln x_{\mathrm{A}}=-R T \ln \left(1-x_{\mathrm{B}}\right) \approx R T x_{B}
$$

The equation then becomes

$$
R T x_{\mathrm{B}}=\Pi V_{\mathrm{m}}
$$

Step 4 Simplify the expression for the osmotic pressure for dilute solutions\\
When the solution is dilute, $x_{\mathrm{B}} \approx n_{\mathrm{B}} / n_{\mathrm{A}}$, and therefore $R T n_{\mathrm{B}} \approx n_{\mathrm{A}} \Pi V_{\mathrm{m}}$. Moreover, $n_{\mathrm{A}} V_{\mathrm{m}}=V$, the total volume of the solvent, so $R T n_{\mathrm{B}} \approx \Pi V$. At this stage $n_{\mathrm{B}} / V$ can be recognized as the molar concentration $[\mathrm{B}]$ of the solute B . It follows that for dilute solutions the osmotic pressure is given by

$$
-\Pi=[\mathrm{B}] R T \longmapsto \quad \text { (5B.16) }
$$

This relation, which is called the van't Hoff equation, is valid only for ideal solutions. However, one of the most common applications of osmometry is to the measurement of molar masses of macromolecules, such as proteins and synthetic polymers. As these huge molecules dissolve to produce solutions that are far from ideal, it is assumed that the van 't Hoff equation is only the first term of a virial-like expansion, much like the extension of the perfect gas equation to real gases (in Topic 1C) to take into account molecular interactions:


\begin{equation*}
\Pi=[\mathrm{J}] R T\{1+B[\mathrm{~J}]+\cdots\} \quad \text { Osmotic virial expansion } \tag{5B.17}
\end{equation*}


(The solute is denoted as J to avoid too many different Bs in this expression.) The additional terms take the non-ideality into account; the empirical constant $B$ is called the osmotic virial coefficient. When it is possible to ignore corrections beyond the term depending on $B$, the osmotic pressure is written as


\begin{equation*}
\Pi=[\mathrm{J}] R T\{1+B[\mathrm{~J}]\} \quad \text { or } \quad \Pi /[\mathrm{J}]=R T+B R T[\mathrm{~J}] \tag{5B.18}
\end{equation*}


It follows that the osmotic virial coefficient may be calculated from the slope, $B R T$, of a plot of $\Pi /[\mathrm{J}]$ against [J], as shown in Fig. 5B.14a.\\
\includegraphics[max width=\textwidth, center]{2025_02_27_d4b95d9718f0b9e2e6b9g-196}

Figure 5B. 14 The plot and extrapolation made to analyse the results of an osmometry experiment using (a) the molar concentration and (b) the mass concentration.

\subsection*{Example 5B. 2 Using osmometry to determine the molar mass of a macromolecule}
The osmotic pressures of solutions of a polymer, denoted J, in water at 298 K are given below. Determine the molar mass of the polymer.

\begin{center}
\begin{tabular}{lccccc}
$c_{\text {mass } / J} /\left(\mathrm{g} \mathrm{dm}^{-3}\right)$ & 1.00 & 2.00 & 4.00 & 7.00 & 9.00 \\
$\Pi / \mathrm{Pa}$ & 27 & 70 & 197 & 500 & 785 \\
\end{tabular}
\end{center}

Collect your thoughts This example is an application of eqn 5B.18, but as the data are in terms of the mass concentration, that equation must first be converted. To do so, note that the molar concentration [J] and the mass concentration $c_{\text {mass, }}$, are related by $[\mathrm{J}]=c_{\text {mass }} / M$, where $M$ is the molar mass of J . Then identify the appropriate plot and the quantity (it will turn out to be the intercept on the vertical axis at $c_{\text {mass }, \mathrm{J}}=0$ ) that gives you the value of $M$.

The solution To express eqn 5B. 18 in terms of the mass concentration, substitute [J] $=c_{\text {mass, }} / M$ and obtain

$$
\frac{\overbrace{\Pi M}^{c_{\text {mass }, \mathrm{J}}}}{\Pi /[\mathrm{J}]}=R T+\overbrace{\frac{B R T c_{\text {mass }, \mathrm{J}}}{B R T[\mathrm{~J}]}}^{B}
$$

Division through by $M$ gives

$$
\overbrace{\frac{\Pi}{c_{\text {mass }, \mathrm{J}}}}^{y}=\overbrace{\frac{R T}{M}}^{\text {intercept }}+\overbrace{\left(\frac{B R T}{M^{2}}\right)}^{+ \text {slope }} \overbrace{c_{\text {mass, } \mathrm{J}}}^{x}+\cdots
$$

It follows that, by plotting $\Pi / c_{\text {mass }, \text {, }}$ against $c_{\text {mass }, \text {, }}$ the results should fall on a straight line with intercept $R T / M$ on the vertical axis at $c_{\text {mass, },}=0$. The following values of $\Pi / c_{\text {mass, }}$ can be calculated from the data:

\begin{center}
\begin{tabular}{lccccc}
$c_{\text {mass }, J} /\left(\mathrm{g} \mathrm{dm}^{-3}\right)$ & 1.00 & 2.00 & 4.00 & 7.00 & 9.00 \\
$(\Pi / \mathrm{Pa}) /\left(c_{\text {mass., }} / \mathrm{g} \mathrm{dm}^{-3}\right)$ & 27 & 35 & 49.2 & 71.4 & 87.2 \\
\end{tabular}
\end{center}

The intercept with the vertical axis at $c_{\text {mass }, \mathrm{J}}=0$ (which is best found by using linear regression and mathematical software) is at

$$
\frac{\Pi / \mathrm{Pa}}{c_{\text {mass }, ~} /\left(\mathrm{g} \mathrm{dm}^{-3}\right)}=19.8
$$

which rearranges into

$$
\Pi / c_{\text {mas } \mathrm{J}, \mathrm{~J}}=19.8 \mathrm{Pag}^{-1} \mathrm{dm}^{3}
$$

Therefore, because this intercept is equal to $R T / M$,

$$
M=\frac{R T}{19.8 \mathrm{~Pa} \mathrm{~g}^{-1} \mathrm{dm}^{3}}=\frac{R T}{1.98 \times 10^{-2} \mathrm{~Pa} \mathrm{~g}^{-1} \mathrm{~m}^{3}}
$$

It follows that

$$
M=\frac{\left(8.3145 \mathrm{~J} \mathrm{~K}^{-1} \mathrm{~mol}^{-1}\right) \times(298 \mathrm{~K})}{1.98 \times 10^{-2} \mathrm{~Pa} \mathrm{~g}^{-1} \mathrm{~m}^{3}}=1.25 \times 10^{5} \mathrm{~g} \mathrm{~mol}^{-1}
$$

The molar mass of the polymer is therefore $125 \mathrm{~kg} \mathrm{~mol}^{-1}$.\\
Comment. Note that once $M$ is known, the coefficient $B$ can be determined from the slope of the graph, which is equal to $B R T / M^{2}$, as shown in Fig. 5B.14b.

Self-test 5B. 2 The osmotic pressures of solutions of poly(vinyl chloride), PVC, in dioxane at $25^{\circ} \mathrm{C}$ were as follows:

\begin{center}
\begin{tabular}{lccccc}
$c_{\text {mass },} /\left(\mathrm{g} \mathrm{dm}^{-3}\right)$ & 0.50 & 1.00 & 1.50 & 2.00 & 2.50 \\
$\Pi / \mathrm{Pa}$ & 33.6 & 35.2 & 36.8 & 38.4 & 40.0 \\
\end{tabular}
\end{center}

Determine the molar mass of the polymer.\\
\includegraphics[max width=\textwidth, center]{2025_02_27_d4b95d9718f0b9e2e6b9g-196(1)}

\subsection*{Checklist of concepts}
\begin{enumerate}
  \item The Gibbs energy of mixing of two liquids to form an ideal solution is calculated in the same way as for two perfect gases.
  \item The enthalpy of mixing for an ideal solution is zero and the Gibbs energy is due entirely to the entropy of mixing.3. A regular solution is one in which the entropy of mixing is the same as for an ideal solution but the enthalpy of mixing is non-zero.4. A colligative property depends only on the number of solute particles present, not their identity.5. All the colligative properties stem from the reduction of the chemical potential of the liquid solvent as a result of the presence of solute.6. The elevation of boiling point is proportional to the molality of the solute.
  \item The depression of freezing point is also proportional to the molality of the solute.
  \item The osmotic pressure is the pressure that when applied to a solution prevents the influx of solvent through a semipermeable membrane.
  \item The relation of the osmotic pressure to the molar concentration of the solute is given by the van 't Hoff equation and is a sensitive way of determining molar mass.
\end{enumerate}

\subsection*{Checklist of equations}
\begin{center}
\begin{tabular}{llll}
\hline
Property & Equation & Comment & Equation number \\
\hline
Gibbs energy of mixing & $\Delta_{\operatorname{mix}} G=n R T\left(x_{\mathrm{A}} \ln x_{\mathrm{A}}+x_{\mathrm{B}} \ln x_{\mathrm{B}}\right)$ & Ideal solutions & 5B.3 \\
Entropy of mixing & $\Delta_{\operatorname{mix}} S=-n R\left(x_{\mathrm{A}} \ln x_{\mathrm{A}}+x_{\mathrm{B}} \ln x_{\mathrm{B}}\right)$ & Ideal solutions & 5B.4 \\
Enthalpy of mixing & $\Delta_{\operatorname{mix}} H=0$ & Ideal solutions &  \\
Excess function & $X^{\mathrm{E}}=\Delta_{\text {mix }} X-\Delta_{\text {mix }} X^{\text {ideal }}$ & Definition & 5B.5 \\
Regular solution & $H^{\mathrm{E}}=n \xi R T x_{\mathrm{A}} x_{\mathrm{B}}$ & Model; $S^{\mathrm{E}}=0$ & 5B.6 \\
Elevation of boiling point & $\Delta T_{\mathrm{b}}=K_{\mathrm{b}} b$ & Empirical, non-volatile solute & 5B.9c \\
Depression of freezing point & $\Delta T_{\mathrm{f}}=K_{\mathrm{f}} b$ & Empirical, solute insoluble in solid solvent & 5B.12 \\
Ideal solubility & $\ln x_{\mathrm{B}}=\left(\Delta_{\mathrm{fus}} H / R\right)\left(1 / T_{\mathrm{f}}-1 / T\right)$ & Ideal solution & 5B.14 \\
van 't Hoff equation & $\Pi=[\mathrm{B}] R T$ & Valid as [B] $\rightarrow 0$ & 5B.16 \\
Osmotic virial expansion & $\Pi=[J] R T\{1+B[J]+\cdots\}$ & Empirical & 5B.17 \\
\hline
\end{tabular}
\end{center}

\section*{TOPIC 5C Phase diagrams of binary systems: liquids }
\subsection*{Why do you need to know this material?}
The separation of complex mixtures is a common task in the chemical industry. The information needed to formulate efficient separation methods is contained in phase diagrams, so it is important to be able to interpret them.

\subsection*{What is the key idea?}
The phase diagram of a liquid mixture can be understood in terms of the variation with temperature and pressure of the composition of the liquid and vapour in mutual equilibrium.

\subsection*{What do you need to know already?}
It would be helpful to review the interpretation of onecomponent phase diagrams and the phase rule (Topic 4A). This Topic also draws on Raoult's law (Topic 5A) and the concept of partial pressure (Topic 1A).

One-component phase diagrams are described in Topic 4A. The phase equilibria of binary systems are more complex because composition is an additional variable. However, they provide very useful summaries of phase equilibria for both ideal and empirically established real systems. This Topic focuses on binary mixtures of liquids. The phase diagrams of liquid-solid mixtures are discussed in Topic 5D.

\subsection*{5c. 1 Vapour pressure diagrams}
The partial vapour pressures of the components of an ideal solution of two volatile liquids are related to the composition of the liquid mixture by Raoult's law (Topic 5A):


\begin{equation*}
p_{\mathrm{A}}=x_{\mathrm{A}} p_{\mathrm{A}}^{\star} \quad p_{\mathrm{B}}=x_{\mathrm{B}} p_{\mathrm{B}}^{\star} \tag{5C.1}
\end{equation*}


where $p_{\mathrm{J}}^{*}$, with $\mathrm{J}=\mathrm{A}, \mathrm{B}$, is the vapour pressure of pure J and $x_{\mathrm{J}}$ is the mole fraction of $J$ in the liquid. The total vapour pressure $p$ of the mixture is therefore\\
\includegraphics[max width=\textwidth, center]{2025_02_27_d4b95d9718f0b9e2e6b9g-198(1)}

Figure 5C. 1 The variation of the total vapour pressure of a binary mixture with the mole fraction of A in the liquid when Raoult's law is obeyed.\\
\includegraphics[max width=\textwidth, center]{2025_02_27_d4b95d9718f0b9e2e6b9g-198}\\
(5C.2)\\
This expression shows that the total vapour pressure (at some fixed temperature) changes linearly with the composition from $p_{\mathrm{B}}^{\star}$ to $p_{\mathrm{A}}^{*}$ as $x_{\mathrm{A}}$ changes from 0 to 1 (Fig. 5C.1).

The compositions of the liquid and vapour that are in mutual equilibrium are not necessarily the same. Common sense suggests that the vapour should be richer in the more volatile component. This expectation can be confirmed as follows. If the mole fractions of the components in the vapour are $y_{J}$ with $\mathrm{J}=\mathrm{A}$ and B , then their partial pressures are $p_{\mathrm{J}}=y_{\mathrm{J}} p$, with $p$ the total pressure. Therefore


\begin{equation*}
y_{\mathrm{A}}=\frac{p_{\mathrm{A}}}{p} \quad y_{\mathrm{B}}=\frac{p_{\mathrm{B}}}{p} \tag{5С.3}
\end{equation*}


Provided the mixture is ideal, the partial pressures and the total pressure may be expressed in terms of the mole fractions in the liquid by using eqn 5C. 1 for $p_{\mathrm{J}}$ and eqn 5C. 2 for the total vapour pressure $p$. The result of combining these relations is

$$
y_{\mathrm{A}}=\frac{x_{\mathrm{A}} p_{\mathrm{A}}^{*}}{p_{\mathrm{B}}^{*}+\left(p_{\mathrm{A}}^{*}-p_{\mathrm{B}}^{*}\right) x_{\mathrm{A}}} \quad y_{\mathrm{B}}=1-y_{\mathrm{A}}
$$

\begin{center}
\includegraphics[max width=\textwidth]{2025_02_27_d4b95d9718f0b9e2e6b9g-199}
\end{center}

Figure 5C. 2 The mole fraction of $A$ in the vapour of a binary ideal solution expressed in terms of its mole fraction in the liquid, calculated using eqn 5C. 4 for various values of $p_{\mathrm{A}}^{*} / p_{\mathrm{B}}^{*}$. For A more volatile than $\mathrm{B}\left(p_{\mathrm{A}}^{*} / p_{\mathrm{B}}^{\star}>1\right)$, the vapour is richer in A compared with the liquid..

Figure 5C. 2 shows the composition of the vapour plotted against the composition of the liquid for various values of $p_{\mathrm{A}}^{*} / p_{\mathrm{B}}^{*} \geq 1$. Provided that $p_{\mathrm{A}}^{*} / p_{\mathrm{B}}^{*}>1$, then $y_{\mathrm{A}}>x_{\mathrm{A}}$ : the vapour is richer than the liquid in the more volatile component. Note that if B is not volatile, so $p_{\mathrm{B}}^{*}=0$ at the temperature of interest, then it makes no contribution to the vapour $\left(y_{B}=0\right)$.

\subsection*{Brief illustration 5C. 1}
The vapour pressures of pure benzene and methylbenzene at $20^{\circ} \mathrm{C}$ are 75 Torr and 21 Torr, respectively. The composition of the vapour in equilibrium with an equimolar liquid mixture $\left(x_{\text {benzene }}=x_{\text {methylbenzene }}=\frac{1}{2}\right)$ is

$$
\begin{aligned}
& y_{\text {benzene }}=\frac{\frac{1}{2} \times(75 \text { Torr })}{21 \text { Torr }+(75-21 \text { Torr }) \times \frac{1}{2}}=0.78 \\
& y_{\text {methylbenzene }}=1-0.78=0.22
\end{aligned}
$$

The partial vapour pressure of each component is

$$
p_{\text {benzene }}=\frac{1}{2} \times(75 \text { Torr })=37.5 \text { Torr }
$$

$$
p_{\text {methylbenzene }}=\frac{1}{2} \times(21 \text { Torr })=10.5 \text { Torr }
$$

and the total vapour pressure is the sum of these two values, 48 Torr.

Equations 5C. 2 and 5C. 4 can be combined to express the total vapour pressure in terms of the composition of the vapour.

How is that done? 5C. 1 Deriving an expression for the total vapour pressure of a binary mixture in terms of the composition of the vapour

Equation 5C. 4 can be rearranged as follows to express $x_{\mathrm{A}}$ in terms of $y_{A}$. First, multiply both sides by $p_{\mathrm{B}}^{*}+\left(p_{\mathrm{A}}^{*}-p_{\mathrm{B}}^{*}\right) x_{\mathrm{A}}$ to obtain

$$
p_{\mathrm{B}}^{*} y_{\mathrm{A}}+\left(p_{\mathrm{A}}^{*}-p_{\mathrm{B}}^{*}\right) x_{\mathrm{A}} y_{\mathrm{A}}=x_{\mathrm{A}} p_{\mathrm{A}}^{*}
$$

Then collect terms in $x_{\mathrm{A}}$ :

$$
p_{\mathrm{B}}^{*} y_{\mathrm{A}}=\left\{p_{\mathrm{A}}^{\star}+\left(p_{\mathrm{B}}^{*}-p_{\mathrm{A}}^{*}\right) y_{\mathrm{A}}\right\} x_{\mathrm{A}}
$$

which rearranges to

$$
x_{\mathrm{A}}=\frac{p_{\mathrm{B}}^{*} y_{\mathrm{A}}}{p_{\mathrm{A}}^{*}+\left(p_{\mathrm{B}}^{*}-p_{\mathrm{A}}^{*}\right) y_{\mathrm{A}}}
$$

From eqn 5C. 2 and the expression for $x_{A}$,

$$
p=p_{\mathrm{B}}^{\star}+\left(p_{\mathrm{A}}^{\star}-p_{\mathrm{B}}^{\star}\right) x_{\mathrm{A}}=p_{\mathrm{B}}^{\star}+\frac{\left(p_{\mathrm{A}}^{\star}-p_{\mathrm{B}}^{\star}\right) p_{\mathrm{B}}^{\star} y_{\mathrm{A}}}{p_{\mathrm{A}}^{\star}+\left(p_{\mathrm{B}}^{\star}-p_{\mathrm{A}}^{\star}\right) y_{\mathrm{A}}}
$$

Finally, after some algebra,

$$
p=\frac{p_{\mathrm{A}}^{\star} p_{\mathrm{B}}^{\star}+\left(p_{\mathrm{B}}^{\star}-p_{\mathrm{A}}^{*}\right) p_{\mathrm{B}}^{*} y_{\mathrm{A}}+\left(p_{\mathrm{A}}^{\star}-p_{\mathrm{B}}^{*}\right) p_{\mathrm{B}}^{\star} y_{\mathrm{A}}}{p_{\mathrm{A}}^{\star}+\left(p_{\mathrm{B}}^{\star}-p_{\mathrm{A}}^{*}\right) y_{\mathrm{A}}}
$$

which simplifies to

$-p=\frac{p_{\mathrm{A}}^{\star} p_{\mathrm{B}}^{\star}}{p_{\mathrm{A}}^{\star}+\left(p_{\mathrm{B}}^{\star}-p_{\mathrm{A}}^{\star}\right) y_{\mathrm{A}}} \quad$\begin{tabular}{|l}
(5C.5) \\
\end{tabular}

This expression is plotted in Fig. 5C.3.\\
\includegraphics[max width=\textwidth, center]{2025_02_27_d4b95d9718f0b9e2e6b9g-199(1)}

Figure 5C. 3 The dependence of the vapour pressure of the same system as in Fig. 5C.2, but expressed in terms of the mole fraction of $A$ in the vapour by using eqn 5C.5. Individual curves are labelled with the value of $p_{A}^{*} / p_{\mathrm{B}}^{\star}$.

\subsection*{5c.2 Temperature-composition diagrams}
A temperature-composition diagram is a phase diagram in which the boundaries show the composition of the phases that are in equilibrium at various temperatures (and a given pressure, typically 1 atm ). An example is shown in Fig. 5C.4. Note that the liquid phase lies in the lower part of the diagram. Temperature-composition diagrams are central to the discussion of distillation. In the following discussion, it will be best to keep in mind a system consisting of a liquid and its vapour confined inside a cylinder fitted with a movable piston that exerts a constant pressure, which in most cases is 1 atm. In this arrangement, the liquid and its vapour are in equilibrium at the normal boiling point of the mixture.

\subsection*{(a) The construction of the diagrams}
Although in principle a temperature-composition diagram could be constructed from vapour-pressure diagrams by examining the temperature dependence of the vapour pressures of the components and identifying the temperature at which the total vapour pressure becomes equal to 1 atm (or whatever ambient pressure is of interest), they are normally constructed from empirical data on the composition of the phases in equilibrium at each temperature.

Provided the ambient pressure is 1 atm , the points representing liquid/vapour equilibrium for each of the pure liquid components are their normal boiling points. The line labelled 'Liquid' displays the boiling temperature (the temperature at which the total vapour pressure is 1 atm ) of the mixture across the range of compositions. The line labelled 'Vapour' is the composition of the vapour in equilibrium with the liquid at\\
\includegraphics[max width=\textwidth, center]{2025_02_27_d4b95d9718f0b9e2e6b9g-200}

Figure 5C. 4 The temperature-composition diagram corresponding to an ideal mixture with the component $A$ more volatile than component B. As described in Section 5C.2, successive boilings and condensations of a liquid originally of composition $a_{1}$ lead to a condensate that is pure A .\\
each temperature. As remarked in the preceding discussion, for ideal solutions the vapour is richer in the more volatile component, so the curve is necessarily displaced towards the pure component that has the higher vapour pressure and therefore the lower boiling temperature.

Example 5C. 1 Constructing a temperature-composition diagram

The following temperature/composition data were obtained for a mixture of octane ( O ) and methylbenzene ( M ) at 1.00 atm , where $x_{\mathrm{M}}$ is the mole fraction of M in the liquid and $y_{\mathrm{M}}$ the mole fraction in the vapour at equilibrium.

\begin{center}
\begin{tabular}{lccccccccc}
$\theta /{ }^{\circ} \mathrm{C}$ & 110.9 & 112.0 & 114.0 & 115.8 & 117.3 & 119.0 & 121.1 & 123.0 \\
$x_{\mathrm{M}}$ & 0.908 & 0.795 & 0.615 & 0.527 & 0.408 & 0.300 & 0.203 & 0.097 \\
$y_{\mathrm{M}}$ & 0.923 & 0.836 & 0.698 & 0.624 & 0.527 & 0.410 & 0.297 & 0.164 \\
\end{tabular}
\end{center}

The boiling points are $110.6^{\circ} \mathrm{C}$ and $125.6^{\circ} \mathrm{C}$ for M and O , respectively. Plot the temperature/composition diagram for the mixture.

Collect your thoughts Plot the composition of each phase (on the horizontal axis) against the temperature (on the vertical axis). The two boiling points give two further points corresponding to $x_{\mathrm{M}}=1$ and $x_{\mathrm{M}}=0$, respectively. Use a spreadsheet or mathematical software to draw the phase boundaries.

The solution The points are plotted in Fig. 5C.5. The two sets of points are fitted to the polynomials $a+b z+c z^{2}+d z^{3}$ with $z=x_{\mathrm{M}}$ for the liquid line and $z=y_{\mathrm{M}}$ for the vapour line.

For the liquid line: $\theta /{ }^{\circ} \mathrm{C}=125.422-22.9494 x_{\mathrm{M}}+6.64602 x_{\mathrm{M}}^{2}$

$$
+1.32623 x_{M}^{3}
$$

For the vapour line: $\theta /{ }^{\circ} \mathrm{C}=125.485-11.9387 y_{\mathrm{M}}-12.5626 y_{\mathrm{M}}^{2}$

$$
+9.36542 y_{M}^{3}
$$

\begin{center}
\includegraphics[max width=\textwidth]{2025_02_27_d4b95d9718f0b9e2e6b9g-200(1)}
\end{center}

Figure 5C. 5 The plot of data and the fitted curves for a mixture of octane $(\mathrm{O})$ and methylbenzene $(\mathrm{M})$ in Example 5C.1.

Self-test 5C.1 Repeat the analysis for the following data on hexane and heptane:

\begin{center}
\begin{tabular}{|c|c|c|c|c|c|c|}
\hline
$\theta /{ }^{\circ} \mathrm{C}$ & 65 & 66 & 70 & 77 & 85 & 100 \\
\hline
$x_{\text {hexane }}$ & 0 & 0.20 & 0.40 & 0.60 & 0.80 & 1 \\
\hline
$y_{\text {hexane }}$ & 0 & 0.02 & 0.08 & 0.20 & 0.48 & 1 \\
\hline
\end{tabular}
\end{center}

\begin{center}
\includegraphics[max width=\textwidth]{2025_02_27_d4b95d9718f0b9e2e6b9g-201(1)}
\end{center}

Figure 5C. 6 The plot of data and the fitted curves for a mixture of hexane $(\mathrm{Hx})$ and heptane in Self-test 5C.1.

\subsection*{(b) The interpretation of the diagrams}
The horizontal axis of the diagram denotes the value of the mole fraction $x_{\mathrm{A}}$ when interpreting the 'Liquid' line and the mole fraction $y_{\mathrm{A}}$ when interpreting the 'Vapour' line, as illustrated in Example 5C.1. That is, a vertical line at $x_{\mathrm{A}}$ intersects the 'Liquid' line at the boiling point of the mixture as it was prepared. The horizontal line at that temperature, which is called a tie line, intersects the 'Vapour' line at a composition that represents the mole fraction $y_{\mathrm{A}}$ of A in the vapour phase in equilibrium with the boiling liquid. When appropriate, the horizontal axis will be labelled $z_{\mathrm{A}}$ and interpreted as $x_{\mathrm{A}}$ or $y_{\mathrm{A}}$ according to which line, 'Liquid' or 'Vapour' respectively, is of interest.

A point in the diagram below the 'Liquid' line at a given temperature corresponds to the mixture being at a temperature below its boiling point. If the ambient pressure is 1 atm , which is greater than the vapour pressure at that temperature, the entire sample is liquid and $x_{\mathrm{A}}$ is its composition. Similarly, if a point is above the 'Vapour' line at a given temperature, then that temperature is above the boiling point of the mixture, its vapour pressure is greater than 1 atm , and the entire sample is a vapour with a composition that is the same as that of the original mixture (because it has become entirely vapour). If the temperature is such that the point lies on the 'Liquid' curve, then the liquid and its vapour are in equilibrium and the composition of the vapour is represented by noting where the tie line meets the 'Vapour' curve. Note that the phase boundary (the 'coexistence curve') representing the frontier between the\\
\includegraphics[max width=\textwidth, center]{2025_02_27_d4b95d9718f0b9e2e6b9g-201}

Figure 5C. 7 The points of the temperature-composition diagram discussed in the text. The vertical line through $a$ is an isopleth, a line of constant composition of the entire system.\\
regions where either the liquid or the vapour is the more stable phase is the 'Liquid' line: the 'Vapour' line simply provides additional information.

Points that lie between the two lines do provide additional information if the horizontal axis denotes the overall composition of the mixture in equilibrium at a given temperature rather than the liquid or vapour composition separately. Thus, consider what happens when a mixture in which the mole fraction of A is $z_{\mathrm{A}}$ is heated. The overall composition does not change regardless of how much liquid vaporizes, so the system moves up the vertical line at $a$ in Fig. 5C.7. Such a vertical line is called an isopleth (from the Greek words for 'equal abundance').

At $a_{1}$ the liquid boils and initially is in equilibrium with its vapour of composition $a_{1}^{\prime}$, as given by the tie line. This vapour is richer in the more volatile component (B), so the liquid is depleted in $B$. Being richer in $A$, the boiling point of the remaining liquid moves to $a_{2}$ and the composition of the vapour in equilibrium with that liquid changes to $a_{2}^{\prime}$. Further heating migrates the composition of the liquid further towards pure A, the boiling point rises and the composition of the vapour changes accordingly to $a_{3}^{\prime}$. At $a_{4}^{\prime}$ the composition of the vapour is the same as the overall composition of the mixture, which implies that all the liquid has vaporized. Above that temperature, only vapour is present and has the initial overall composition.

It is also possible to predict the abundances of liquid and vapour at any stage of heating when the temperature and overall composition correspond to a point between the 'Liquid' and 'Vapour' lines, where a liquid of one composition is in equilibrium with a vapour of another composition.

\subsection*{How is that done? 5C. 2 Establishing the lever rule}
If the amount of A molecules in the vapour is $n_{A, V}$ and the amount in the liquid phase is $n_{\mathrm{A}, \mathrm{L}}$, the total amount of A molecules is $n_{\mathrm{A}}=n_{\mathrm{A}, \mathrm{L}}+n_{\mathrm{A}, \mathrm{V}}$ and likewise for B molecules.

The overall mole fraction of A is $z_{\mathrm{A}}=\left(n_{\mathrm{A}, \mathrm{L}}+n_{\mathrm{A}, \mathrm{V}}\right) /\left(n_{\mathrm{A}}+n_{\mathrm{B}}\right)$. The total amount of molecules in the liquid (both A and B ) is $n_{\mathrm{L}}=n_{\mathrm{A}, \mathrm{L}}+n_{\mathrm{B}, \mathrm{L}}$, and the total amount of molecules in the vapour is likewise $n_{\mathrm{V}}=n_{\mathrm{A}, \mathrm{V}}+n_{\mathrm{B}, \mathrm{V}}$. These relations can be written in terms of the mole fractions in the vapour $\left(y_{A}\right)$ and liquid $\left(x_{\mathrm{A}}\right)$ phases. Thus, the amount of A in the liquid phase is $n_{\mathrm{L}} x_{\mathrm{A}}$. Similarly, the amount of A in the vapour phase is $n_{\mathrm{V}} y_{\mathrm{A}}$. The total amount of A is therefore

$$
n_{\mathrm{A}}=n_{\mathrm{L}} x_{\mathrm{A}}+n_{\mathrm{V}} y_{\mathrm{A}}
$$

The total amount of A molecules is also

$$
n_{\mathrm{A}}=n z_{\mathrm{A}}=n_{\mathrm{L}} z_{\mathrm{A}}+n_{\mathrm{V}} z_{\mathrm{A}}
$$

By equating these two expressions it follows that $n_{\mathrm{L}} x_{\mathrm{A}}+n_{\mathrm{V}} y_{\mathrm{A}}=$ $n_{\mathrm{L}} z_{\mathrm{A}}+n_{\mathrm{V}} z_{\mathrm{A}}$, and therefore

$$
n_{\mathrm{L}} \overbrace{\left.z_{\mathrm{A}}-x_{\mathrm{A}}\right)}^{\mathrm{I}_{\mathrm{L}}}=n_{\mathrm{V}} \overbrace{\left(y_{\mathrm{A}}-z_{\mathrm{A}}\right)}^{I_{\mathrm{V}}}
$$

As shown in Fig. 5C.8, with $z_{\mathrm{A}}-x_{\mathrm{A}}$ defined as the 'length' $l_{\mathrm{L}}$, and $y_{\mathrm{A}}-z_{\mathrm{A}}$ defined as the 'length' $l_{\mathrm{V}}$, this relation can be expressed as the lever rule:

$$
-\mid n_{\mathrm{L}} l_{\mathrm{L}}=n_{\mathrm{v}} l_{\mathrm{V}} \downharpoonright \quad \text { (5C.6) }
$$

The lever rule applies to any phase diagram, not only to liquid-vapour equilibria.\\
\includegraphics[max width=\textwidth, center]{2025_02_27_d4b95d9718f0b9e2e6b9g-202}

Figure 5C. 8 The lever rule. The distances $I_{V}$ and $I_{L}$ are used to find the proportions of the amounts of the vapour and liquid present at equilibrium. The lever rule is so called because a similar rule relates the masses at two ends of a lever to their distances from a pivot (in that case $m_{V} l_{V}=m_{L} l_{L}$ for balance).

\subsection*{Brief illustration 5C. 2}
In the case illustrated in Fig. 5C.7, because $l_{\mathrm{V}} \approx 2 l_{\mathrm{L}}$ at the tie line at $a_{3}$, the amount of molecules in the liquid phase is about twice the amount of molecules in the vapour phase. At $a_{1}$ in Fig. 5C.7, the ratio $l_{\mathrm{V}} / l_{\mathrm{L}}$ is almost infinite for this tie line, so\\
$n_{\mathrm{L}} / n_{\mathrm{V}}$ is also almost infinite, and there is only a trace of vapour present. When the temperature is raised to $a_{2}$, the value of $l_{\mathrm{V}} / l_{\mathrm{L}}$ is about 6.9 , so $n_{\mathrm{L}} / n_{\mathrm{V}} \approx 0.15$ and the amount of molecules present in the liquid is about 0.15 times the amount in the vapour. When the temperature has increased to $a_{4}$ and $l_{\mathrm{V}} / l_{\mathrm{L}} \approx 0$ there is only a trace of liquid present.

\subsection*{5c. 3 Distillation}
Consider what happens when a liquid of composition $a_{1}$ in Fig. 5C. 4 is heated. It boils when the temperature reaches $T_{2}$. Then the liquid has composition $a_{2}$ (the same as $a_{1}$ ) and the vapour (which is present only as a trace) has composition $a_{2}^{\prime}$. The vapour is richer in the more volatile component A (the component with the lower boiling point). The composition of the vapour at the boiling point follows from the location of $a_{2}$, and from the location of the tie line joining $a_{2}$ and $a_{2}^{\prime}$ it is possible to read off the boiling temperature $\left(T_{2}\right)$ of the original liquid mixture.

\subsection*{(a) Simple and fractional distillation}
In a simple distillation, the vapour is withdrawn and condensed. This technique is used to separate a volatile liquid from a non-volatile solute or solid. In fractional distillation, the boiling and condensation cycle is repeated successively. This technique is used to separate volatile liquids.

Consider what happens if the vapour at $a_{2}^{\prime}$ in Fig. 5C. 4 is condensed, and then this condensate (of composition $a_{3}$ ) is reheated. The phase diagram shows that this mixture boils at $T_{3}$ and yields a vapour of composition $a_{3}^{\prime}$, which is even richer in the more volatile component. That vapour is drawn off, and the first drop condenses to a liquid of composition $a_{4}$. The cycle can then be repeated until in due course almost pure $A$ is obtained in the vapour and pure $B$ remains in the liquid.

The efficiency of a fractionating column is expressed in terms of the number of theoretical plates, the number of effective vaporization and condensation steps that are required to achieve a condensate of given composition from a given distillate.

\subsection*{Brief illustration 5C. 3}
To achieve the degree of separation shown in Fig. 5C.9a, the fractionating column must correspond to three theoretical plates. To achieve the same separation for the system shown in Fig. 5C.9b, in which the components have more similar normal boiling points, the fractionating column must be designed to correspond to four theoretical plates.\\
\includegraphics[max width=\textwidth, center]{2025_02_27_d4b95d9718f0b9e2e6b9g-203}

Figure 5C. 9 The number of theoretical plates is the number of steps needed to bring about a specified degree of separation of two components in a mixture. The two systems shown correspond to (a) 3, (b) 4 theoretical plates.

\subsection*{(b) Azeotropes}
Although many liquids have temperature-composition phase diagrams resembling the ideal version shown in Fig. 5C.4, in a number of important cases there are marked deviations. A maximum in the phase diagram (Fig. 5C.10) may occur when the favourable interactions between $A$ and $B$ molecules reduce the vapour pressure of the mixture below the ideal value and so raise its boiling temperature: in effect, the A-B interactions stabilize the liquid. In such cases the excess Gibbs energy, $G^{\mathrm{E}}$ (Topic 5B), is negative (more favourable to mixing than ideal). Phase diagrams showing a minimum (Fig. 5C.11) indicate that the mixture is destabilized relative to the ideal solution, the A-B interactions then being unfavourable; in this case, the boiling temperature is lowered. For such mixtures $G^{\mathrm{E}}$ is positive (less favourable to mixing than ideal), and there may be contributions from both enthalpy and entropy effects.\\
\includegraphics[max width=\textwidth, center]{2025_02_27_d4b95d9718f0b9e2e6b9g-203(2)}

Figure 5C. 10 A high-boiling azeotrope. When the liquid of composition $a$ is distilled, the composition of the remaining liquid changes towards $b$ but no further.\\
\includegraphics[max width=\textwidth, center]{2025_02_27_d4b95d9718f0b9e2e6b9g-203(1)}

Figure 5C. 11 A low-boiling azeotrope. When the mixture at $a$ is fractionally distilled, the vapour in equilibrium in the fractionating column moves towards $b$ and then remains unchanged.

Deviations from ideality are not always so strong as to lead to a maximum or minimum in the phase diagram, but when they do there are important consequences for distillation. Consider a liquid of composition $a$ on the right of the maximum in Fig. 5C.10. The vapour (at $a_{2}^{\prime}$ ) of the boiling mixture (at $a_{2}$ ) is richer in A. If that vapour is removed (and condensed elsewhere), then the remaining liquid will move to a composition that is richer in $B$, such as that represented by $a_{3}$, and the vapour in equilibrium with this mixture will have composition $a_{3}^{\prime}$. As that vapour is removed, the composition of the boiling liquid shifts to a point such as $a_{4}$, and the composition of the vapour shifts to $a_{4}^{\prime}$. Hence, as evaporation proceeds, the composition of the remaining liquid shifts towards B as A is drawn off. The boiling point of the liquid rises, and the vapour becomes richer in B . When so much A has been evaporated that the liquid has reached the composition $b$, the vapour has the same composition as the liquid. Evaporation then occurs without change of composition. The mixture is said to form an azeotrope. ${ }^{1}$ When the azeotropic composition has been reached, distillation cannot separate the two liquids because the condensate has the same composition as the azeotropic liquid.

The system shown in Fig. 5C. 11 is also azeotropic, but shows its azeotropy in a different way. Suppose we start with a mixture of composition $a_{1}$, and follow the changes in the composition of the vapour that rises through a fractionating column (essentially a vertical glass tube packed with glass rings to give a large surface area). The mixture boils at $a_{2}$ to give a vapour of composition $a_{2}^{\prime}$. This vapour condenses in the column to a liquid of the same composition (now marked $a_{3}$ ). That liquid reaches equilibrium with its vapour at $a_{3}^{\prime}$, which condenses higher up the tube to give a liquid of the same composition, which we now call $a_{4}$. The fractionation therefore shifts the

\footnotetext{${ }^{1}$ The name comes from the Greek words for 'boiling without changing'.
}
vapour towards the azeotropic composition at $b$, but not beyond, and the azeotropic vapour emerges from the top of the column.

\subsection*{Brief illustration 5C. 4}
Examples of the behaviour of the type shown in Fig. 5C. 10 include (a) trichloromethane/propanone and (b) nitric acid/ water mixtures. Hydrochloric acid/water is azeotropic at 80 per cent by mass of water and boils unchanged at $108.6^{\circ} \mathrm{C}$. Examples of the behaviour of the type shown in Fig. 5C. 11 include (c) dioxane/water and (d) ethanol/water mixtures. Ethanol/water boils unchanged when the water content is 4 per cent by mass and the temperature is $78^{\circ} \mathrm{C}$.

\subsection*{(c) Immiscible liquids}
Consider the distillation of two immiscible liquids, such as octane and water. At equilibrium, there is a tiny amount of A dissolved in $B$, and similarly a tiny amount of $B$ dissolved in A: both liquids are saturated with the other component (Fig. 5C.12(a)). As a result, the total vapour pressure of the mixture is close to $p=p_{\mathrm{A}}^{*}+p_{\mathrm{B}}^{\star}$. If the temperature is raised to the value at which this total vapour pressure is equal to the atmospheric pressure, boiling commences and the dissolved substances are purged from their solution. However, this boiling results in a vigorous agitation of the mixture, so each component is kept saturated in the other component, and the purging continues as the very dilute solutions are replenished. This intimate contact is essential: two immiscible liquids heated in a container like that shown in Fig. 5C.12(b) would not boil at the same temperature. The presence of the saturated solutions means that the 'mixture' boils at a lower temperature than either component would alone because boiling begins when the total vapour pressure reaches 1 atm, not when either vapour pressure reaches 1 atm . This distinction is the basis of steam distillation, which enables some heat-sensitive,\\
\includegraphics[max width=\textwidth, center]{2025_02_27_d4b95d9718f0b9e2e6b9g-204}

Figure 5C. 12 The distillation of (a) two immiscible liquids is quite different from (b) the joint distillation of the separated components, because in the former, boiling occurs when the sum of the partial pressures equals the external pressure.\\
water-insoluble organic compounds to be distilled at a lower temperature than their normal boiling point. The only snag is that the composition of the condensate is in proportion to the vapour pressures of the components, so oils of low volatility distil in low abundance.

\subsection*{5c. 4 Liquid-liquid phase diagrams}
Consider temperature-composition diagrams for systems that consist of pairs of partially miscible liquids, which are liquids that do not mix in all proportions at all temperatures. An example is hexane and nitrobenzene. The same principles of interpretation apply as to liquid-vapour diagrams.

\subsection*{(a) Phase separation}
Suppose a small amount of a liquid B is added to a sample of another liquid A at a temperature $T^{\prime}$. Liquid B dissolves completely, and the binary system remains a single phase. As more $B$ is added, a stage comes at which no more dissolves. The sample now consists of two phases in equilibrium with each other, the most abundant one consisting of $A$ saturated with $B$, the minor one a trace of $B$ saturated with $A$. In the temperaturecomposition diagram drawn in Fig. 5C.13, the composition of the former is represented by the point $a^{\prime}$ and that of the latter by the point $a^{\prime \prime}$. The relative abundances of the two phases are given by the lever rule. When more $B$ is added the composition $a$ moves to the right on the diagram, A dissolves in the added $B$ slightly, and the compositions of the two phases in equilibrium remain $a^{\prime}$ and $a^{\prime \prime}$. As yet more B is added, composition $a$ moves further to the right and eventually crosses the phase\\
\includegraphics[max width=\textwidth, center]{2025_02_27_d4b95d9718f0b9e2e6b9g-204(1)}

Figure 5C. 13 The temperature-composition diagram for a mixture of $A$ and $B$. The region below the curve corresponds to the compositions and temperatures at which the liquids are partially miscible. The upper critical temperature, $T_{u c}$, is the temperature above which the two liquids are miscible in all proportions.\\
boundary into the one-phase region. So much $B$ is now present that it can dissolve all the $A$ and the system reverts to a single phase. The addition of more B now simply dilutes the solution, and from then on a single phase remains.

The composition of the two phases at equilibrium varies with the temperature. For the system shown in Fig. 5C.13, raising the temperature increases the miscibility of A and B . The two-phase region therefore becomes narrower because each phase in equilibrium is richer in its minor component: the A-rich phase is richer in B and the B-rich phase is richer in A . The entire phase diagram can be constructed by repeating the observations at different temperatures and drawing the envelope of the two-phase region.

\subsection*{Example 5C. 2}
Interpreting a liquid-liquid phase diagram\\
The phase diagram for the system nitrobenzene/hexane at 1 atm is shown in Fig. 5C.14. A mixture of 50 g of hexane $\left(0.59 \mathrm{~mol} \mathrm{C}_{6} \mathrm{H}_{14}\right)$ and 50 g of nitrobenzene $\left(0.41 \mathrm{~mol} \mathrm{C}_{6} \mathrm{H}_{5} \mathrm{NO}_{2}\right)$ was prepared at 290 K . What are the compositions of the phases, and in what proportions do they occur? To what temperature must the sample be heated in order to obtain a single phase?\\
\includegraphics[max width=\textwidth, center]{2025_02_27_d4b95d9718f0b9e2e6b9g-205}

Figure 5C. 14 The temperature-composition diagram for hexane and nitrobenzene at 1 atm , with the points and lengths discussed in the text.

Collect your thoughts The compositions of phases in equilibrium are given by the points where the tie line at the relevant temperature intersects the phase boundary. Their proportions are given by the lever rule (eqn 5C.6). The temperature at which the components are completely miscible is found by following the isopleth upwards and noting the temperature at which it enters the one-phase region of the phase diagram.

The solution Denote hexane by H and nitrobenzene by N , then refer to Fig. 5C.14. The mole fraction of N in the mixture is $0.41 /(0.41+0.59)=0.41$. The point $x_{\mathrm{N}}=0.41, T=290 \mathrm{~K}$ occurs in the two-phase region of the phase diagram. The horizontal tie line cuts the phase boundary at $x_{\mathrm{N}}=0.35$ and $x_{\mathrm{N}}=0.83$, so those are the compositions of the two phases.

According to the lever rule, the ratio of amounts of each phase, which are now denoted $\alpha$ and $\beta$, is equal to the ratio of the distances $l_{\alpha}$ and $l_{\beta}$ :

$$
\frac{n_{\alpha}}{n_{\beta}}=\frac{l_{\beta}}{l_{\alpha}}=\frac{0.83-0.41}{0.41-0.35}=\frac{0.42}{0.06}=7
$$

That is, there is about 7 times more hexane-rich phase than nitrobenzene-rich phase. Heating the sample to 292 K takes it into the single-phase region. Because the phase diagram has been constructed experimentally, these conclusions are not based on any assumptions about ideality. They would be modified if the system were subjected to a different pressure.

Self-test 5C. 2 Repeat the problem for 50 g of hexane and 100 g of nitrobenzene at 273 K .\\
\includegraphics[max width=\textwidth, center]{2025_02_27_d4b95d9718f0b9e2e6b9g-205(1)}

\subsection*{(b) Critical solution temperatures}
The upper critical solution temperature, $T_{\text {uc }}$ (or upper consolute temperature), is the highest temperature at which phase separation occurs. Above the upper critical temperature the two components are fully miscible. This temperature exists because the greater thermal motion overcomes any potential energy advantage in molecules of one type being close together. An example is the nitrobenzene/hexane system shown in Fig. 5C. 14.

The thermodynamic interpretation of the upper critical solution temperature focuses on the Gibbs energy of mixing and its variation with temperature. The simple model of a real solution (specifically, of a regular solution) discussed in Topic 5B results in a Gibbs energy of mixing that behaves as shown in Fig. 5C.15.\\
\includegraphics[max width=\textwidth, center]{2025_02_27_d4b95d9718f0b9e2e6b9g-205(2)}

Figure 5C. 15 The temperature variation of the Gibbs energy of mixing of a system composed of two components that are partially miscible at low temperatures. When two minima are present in one of these curves, the system separates into two phases with compositions corresponding to the position of the local minima. This illustration is a duplicate of Fig. 5B.5.

Provided the parameter $\xi$ introduced in eqn 5B.6 $\left(H^{\mathrm{E}}=\right.$ $n \xi R T x_{\mathrm{A}} x_{\mathrm{B}}$ ) is greater than 2 , the Gibbs energy of mixing has a double minimum. As a result, for $\xi>2$ phase separation is expected to occur. The compositions corresponding to the minima are obtained by looking for the conditions at which $\partial \Delta_{\text {mix }} G / \partial x_{\mathrm{A}}=0$. A simple manipulation of eqn 5B. 7 $\left(\Delta_{\text {mix }} G=n R T\left(x_{\mathrm{A}} \ln x_{\mathrm{A}}+x_{\mathrm{B}} \ln x_{\mathrm{B}}+\xi x_{\mathrm{A}} x_{\mathrm{B}}\right)\right.$, with $\left.x_{\mathrm{B}}=1-x_{\mathrm{A}}\right)$ shows that

$$
\begin{aligned}
& \left(\frac{\partial \Delta_{\text {mix }} G}{\partial x_{\mathrm{A}}}\right)_{T, p} \\
& =n R T\left(\frac{\partial\left\{x_{\mathrm{A}} \ln x_{\mathrm{A}}+\left(1-x_{\mathrm{A}}\right) \ln \left(1-x_{\mathrm{A}}\right)+\xi x_{\mathrm{A}}\left(1-x_{\mathrm{A}}\right)\right\}}{\partial x_{\mathrm{A}}}\right)_{T, p} \\
& =n R T\left\{\ln x_{\mathrm{A}}+1-\ln \left(1-x_{\mathrm{A}}\right)-1+\xi\left(1-2 x_{\mathrm{A}}\right)\right\} \\
& =n R T\left\{\ln \frac{x_{\mathrm{A}}}{1-x_{\mathrm{A}}}+\xi\left(1-2 x_{\mathrm{A}}\right)\right\}
\end{aligned}
$$

The Gibbs-energy minima therefore occur where


\begin{equation*}
\ln \frac{x_{\mathrm{A}}}{1-x_{\mathrm{A}}}=-\xi\left(1-2 x_{\mathrm{A}}\right) \tag{5C.7}
\end{equation*}


This equation is an example of a 'transcendental equation', an equation that does not have a solution that can be expressed in a closed form. The solutions (the values of $x_{\mathrm{A}}$ that satisfy the equation) can be found numerically by using mathematical software or by plotting the terms on the left and right against $x_{\mathrm{A}}$ for a choice of values of $\xi$ and identifying the values of $x_{\mathrm{A}}$ where the plots intersect, which is where the two expressions are equal (Fig. 5C.16). The solutions found in this way are plotted in Fig. 5C.17. As $\xi$ decreases, the two minima move together and merge when $\xi=2$.\\
\includegraphics[max width=\textwidth, center]{2025_02_27_d4b95d9718f0b9e2e6b9g-206}

Figure 5C. 16 The graphical procedure for solving eqn 5C.7. When $\xi<2$, the only intersection occurs at $x=0$. When $\xi \geq 2$, there are two solutions (those for $\xi=3$ are marked).\\
\includegraphics[max width=\textwidth, center]{2025_02_27_d4b95d9718f0b9e2e6b9g-206(1)}

Figure 5C. 17 The location of the phase boundary as computed on the basis of the $\xi$-parameter model introduced in Topic 5B.

\subsection*{Brief illustration 5C. 5}
In the system composed of benzene and cyclohexane treated in Example 5B. 1 it is established that $\xi=1.13$, so a two-phase system is not expected. That is, the two components are completely miscible at the temperature of the experiment. The single solution of the equation

$$
\ln \frac{x_{\mathrm{A}}}{1-x_{\mathrm{A}}}+1.13\left(1-2 x_{\mathrm{A}}\right)=0
$$

is $x_{\mathrm{A}}=\frac{1}{2}$, corresponding to a single minimum of the Gibbs energy of mixing, and there is no phase separation.

Some systems show a lower critical solution temperature, $T_{\text {lc }}$ (or lower consolute temperature), below which they mix in all proportions and above which they form two phases. An example is water and triethylamine (Fig. 5C.18). In this\\
\includegraphics[max width=\textwidth, center]{2025_02_27_d4b95d9718f0b9e2e6b9g-206(2)}

Figure 5C. 18 The temperature-composition diagram for water and triethylamine. This system shows a lower critical solution temperature at 292 K . The labels indicate the interpretation of the boundaries.\\
\includegraphics[max width=\textwidth, center]{2025_02_27_d4b95d9718f0b9e2e6b9g-207(2)}

Figure 5C. 19 The temperature-composition diagram for water and nicotine, which has both upper and lower critical temperatures. Note the high temperatures for the liquid (especially the water): the diagram corresponds to a sample under pressure.\\
case, at low temperatures the two components are more miscible because they form a weak complex; at higher temperatures the complexes break up and the two components are less miscible.

Some systems have both upper and lower critical solution temperatures. They occur because, after the weak complexes have been disrupted, leading to partial miscibility, the thermal motion at higher temperatures homogenizes the mixture again, just as in the case of ordinary partially miscible liquids. The most famous example is nicotine and water, which are partially miscible between $61^{\circ} \mathrm{C}$ and $210^{\circ} \mathrm{C}$ (Fig. 5C.19).

\subsection*{(c) The distillation of partially miscible liquids}
Consider a pair of liquids that are partially miscible and form a low-boiling azeotrope. This combination is quite common because both properties reflect the tendency of the two kinds of molecule to avoid each other. There are two possibilities: one in which the liquids become fully miscible before they boil; the other in which boiling occurs before mixing is complete.

Figure 5C. 20 shows the phase diagram for two components that become fully miscible before they boil. Distillation of a mixture of composition $a_{1}$ leads to a vapour of composition $b_{1}$, which condenses to the completely miscible single-phase solution at $b_{2}$. Phase separation occurs only when this distillate is cooled to a point in the two-phase liquid region, such as $b_{3}$. This description applies only to the first drop of distillate. If distillation continues, the composition of the remaining liquid changes. In the end, when the whole sample has evaporated and condensed, the composition is back to $a_{1}$.

Figure 5C. 21 shows the second possibility, in which there is no upper critical solution temperature. The distillate obtained\\
\includegraphics[max width=\textwidth, center]{2025_02_27_d4b95d9718f0b9e2e6b9g-207(1)}

Figure 5C. 20 The temperature-composition diagram for a binary system in which the upper critical solution temperature is less than the boiling point at all compositions. The mixture forms a low-boiling azeotrope.\\
\includegraphics[max width=\textwidth, center]{2025_02_27_d4b95d9718f0b9e2e6b9g-207}

Figure 5C. 21 The temperature-composition diagram for a binary system in which boiling occurs before the two liquids are fully miscible.\\
from a liquid initially of composition $a_{1}$ has composition $b_{3}$ and is a two-phase mixture. One phase has composition $b_{3}^{\prime}$ and the other has composition $b_{3}^{\prime \prime}$.

The behaviour of a system of composition represented by the isopleth $e$ in Fig. 5C. 21 is interesting. A system at $e_{1}$ forms two phases, which persist (but with changing proportions) up to the boiling point at $e_{2}$. The vapour of this mixture has the same composition as the liquid (the liquid is an azeotrope). Similarly, condensing a vapour of composition $e_{3}$ gives a two-phase liquid of the same overall composition. At a fixed temperature, the mixture vaporizes and condenses like a single substance.

\subsection*{Example 5C. 3}
 Interpreting a phase diagramState the changes that occur when a mixture of composition $x_{\mathrm{B}}=0.95\left(a_{1}\right)$ in Fig. 5C. 22 is boiled and the vapour condensed.\\
\includegraphics[max width=\textwidth, center]{2025_02_27_d4b95d9718f0b9e2e6b9g-208}

Figure 5C. 22 The points of the phase diagram in Fig. 5C. 20 that are discussed in Example 5C.3.

Collect your thoughts The area in which the point lies gives the number of phases; the compositions of the phases are given by the points at the intersections of the horizontal tie line with the phase boundaries; the relative abundances are given by the lever rule.

The solution The initial point is in the one-phase region. When heated it boils at $350 \mathrm{~K}\left(a_{2}\right)$ giving a vapour of composition $x_{\mathrm{B}}=0.66\left(b_{1}\right)$. The liquid gets richer in B , and the last drop (of pure B) evaporates at 390 K . The boiling range of the liquid is therefore $350-390 \mathrm{~K}$. If the initial vapour is drawn off, it has a composition $x_{\mathrm{B}}=0.66$. Cooling the distillate corresponds to moving down the $x_{\mathrm{B}}=0.66$ isopleth. At 330 K , for instance, the liquid phase has composition $x_{\mathrm{B}}=0.87$, the vapour $x_{\mathrm{B}}=0.49$; their relative proportions are 1:6. At 320 K the sample consists of three phases: the vapour and two liquids. One liquid phase has composition $x_{\mathrm{B}}=0.30$; the other has composition $x_{\mathrm{B}}=$ 0.80 in the ratio $0.52: 1$. Further cooling moves the system into the two-phase region, and at 298 K the compositions are 0.20 and 0.90 in the ratio $0.82: 1$. As further distillate boils over, the overall composition of the distillate becomes richer in B. When the last drop has been condensed the phase composition is the same as at the beginning.

Self-test 5C. 3 Repeat the discussion, beginning at the point $x_{\mathrm{B}}=0.4, T=298 \mathrm{~K}$.

\subsection*{Checklist of concepts}
\begin{enumerate}
  \item Raoult's law is used to calculate the total vapour pressure of a binary system of two volatile liquids.2. A temperature-composition diagram is a phase diagram in which the boundaries show the composition of the phases that are in equilibrium at various temperatures.3. The composition of the vapour and the liquid phase in equilibrium are located at each end of a tie line.4. The lever rule is used to deduce the relative abundances of each phase in equilibrium.
  \item Separation of a liquid mixture by fractional distillation involves repeated cycles of boiling and condensation.6. An azeotrope is a liquid mixture that evaporates without change of composition.
  \item Phase separation of partially miscible liquids may occur when the temperature is below the upper critical solution temperature or above the lower critical solution temperature; the process may be discussed in terms of the model of a regular solution.
\end{enumerate}

\subsection*{Checklist of equations}
\begin{center}
\begin{tabular}{llll}
\hline
Property & Equation & Comment & Equation number \\
\hline
Composition of vapour & $y_{\mathrm{A}}=x_{\mathrm{A}} p_{\mathrm{A}}^{*} /\left\{p_{\mathrm{B}}^{*}+\left(p_{\mathrm{A}}^{*}-p_{\mathrm{B}}^{*}\right) x_{\mathrm{A}}\right\}$ & Ideal solution &  \\
 & $y_{\mathrm{B}}=1-y_{\mathrm{A}}$ & 5C.4 &  \\
Total vapour pressure & $p=p_{\mathrm{A}}^{*} p_{\mathrm{B}}^{*} /\left\{p_{\mathrm{A}}^{*}+\left(p_{\mathrm{B}}^{*}-p_{\mathrm{A}}^{*}\right) y_{\mathrm{A}}\right\}$ & Ideal solution & 5C.5 \\
Lever rule & $n_{\mathrm{L}} l_{\mathrm{L}}=n_{\mathrm{V}} l_{\mathrm{V}}$ (liquid and vapour phase at equilibrium) & In general, $n_{\alpha} l_{\alpha}=n_{\beta} l_{\beta}$ for phases $\alpha$ and $\beta$ & 5C.6 \\
\hline
\end{tabular}
\end{center}

\section*{TOPIC 5D Phase diagrams of binary}
\subsection*{systems: solids}
\subsection*{Why do you need to know this material?}
Phase diagrams of solid mixtures are used widely in materials science, metallurgy, geology, and the chemical industry to summarize the composition of the various phases of mixtures, and it is important to be able to interpret them.

\subsection*{What is the key idea?}
A phase diagram is a map showing the conditions under which each phase of a system is the most stable.

\subsection*{- What do you need to know already?}
It would be helpful to review the interpretation of liquidliquid phase diagrams and the significance of the lever rule (Topic 5C).

This Topic considers systems where solid and liquid phases might both be present at temperatures below the boiling point.

\subsection*{5D. 1 Eutectics}
Consider the two-component liquid of composition $a_{1}$ in Fig. 5D.1. The changes that occur as the system is cooled may be expressed as follows:

\begin{itemize}
  \item $a_{1} \rightarrow a_{2}$. The system enters the two-phase region labelled 'Liquid + B'. Pure solid B begins to come out of solution and the remaining liquid becomes richer in A.
  \item $a_{2} \rightarrow a_{3}$. More of the solid B forms and the relative amounts of the solid and liquid (which are in equilibrium) are given by the lever rule (Topic 5C). At this stage there are roughly equal amounts of each. The liquid phase is richer in A than before (its composition is given by $b_{3}$ ) because some B has been deposited.
  \item $a_{3} \rightarrow a_{4}$. At the end of this step, there is less liquid than at $a_{3}$, and its composition is given by $e_{2}$. This liquid now freezes to give a two-phase system of pure B and pure A.\\
\includegraphics[max width=\textwidth, center]{2025_02_27_d4b95d9718f0b9e2e6b9g-209}
\end{itemize}

Figure 5D. 1 The temperature-composition phase diagram for two almost immiscible solids and their completely miscible liquids. Note the similarity to Fig. 5C.21. The isopleth through $e_{2}$ corresponds to the eutectic composition, the mixture with lowest melting point.

The isopleth (constant-composition line) at $e_{2}$ in Fig. 5D. 1 corresponds to the eutectic composition, the mixture with the lowest melting point. ${ }^{1}$ A liquid with the eutectic composition freezes at a single temperature, without previously depositing solid A or B. A solid with the eutectic composition melts, without change of composition, at the lowest temperature of any mixture. Solutions of composition to the right of $e_{2}$ deposit $B$ as they cool, and solutions to the left deposit A: only the eutectic mixture (apart from pure A or pure B) solidifies at a single definite temperature without gradually unloading one or other of the components from the liquid.

One eutectic that was technologically important until replaced by modern materials is a formulation of solder in which the mass composition is about 67 per cent tin and 33 per cent lead and melts at $183^{\circ} \mathrm{C}$. The eutectic formed by 23 per cent NaCl and 77 per cent $\mathrm{H}_{2} \mathrm{O}$ by mass melts at $-21.1^{\circ} \mathrm{C}$. When salt is added to ice under isothermal conditions (for example, when spread on an icy road) the mixture melts if the temperature is above $-21.1^{\circ} \mathrm{C}$ (and the eutectic composition has been achieved). When salt is added to ice under adiabatic conditions (for example, when added to ice in a vacuum

\footnotetext{${ }^{1}$ The name comes from the Greek words for 'easily melted'.
}
\includegraphics[max width=\textwidth, center]{2025_02_27_d4b95d9718f0b9e2e6b9g-210(2)}

Figure 5D. 2 The cooling curves for the system shown in Fig. 5D.1. For isopleth $a$, the rate of cooling slows at $a_{2}$ because solid B deposits from solution. There is a complete halt between $a_{3}$ and $a_{4}$ while the eutectic solidifies. This halt is longest for the eutectic isopleth, $e$. The eutectic halt shortens again for compositions beyond $e$ (richer in A). Cooling curves are used to construct the phase diagram.\\
flask) the ice melts, but in doing so it absorbs heat from the rest of the mixture. The temperature of the system falls and, if enough salt is added, cooling continues down to the eutectic temperature. Eutectic formation occurs in the great majority of binary alloy systems, and is of great importance for the microstructure of solid materials. Although a eutectic solid is a two-phase system, it crystallizes out in a nearly homogeneous mixture of microcrystals. The two microcrystalline phases can be distinguished by microscopy and structural techniques such as X-ray diffraction (Topic 15B).

Thermal analysis is a very useful practical way of detecting eutectics. How it is used can be understood by considering the rate of cooling down the isopleth through $a_{1}$ in Fig. 5D.1. The liquid cools steadily until it reaches $a_{2}$, when B begins to be deposited (Fig. 5D.2). Cooling is now slower because the solidification of B is exothermic and retards the cooling. When the remaining liquid reaches the eutectic composition, the temperature remains constant until the whole sample has solidified: this region of constant temperature is the eutectic halt. If the liquid has the eutectic composition $e$ initially, the liquid cools steadily down to the freezing temperature of the eutectic, when there is a long eutectic halt as the entire sample solidifies (like the freezing of a pure liquid).

\subsection*{Brief illustration 5D. 1}
Figure 5D. 3 shows the phase diagram for the binary system silver/tin. The regions have been labelled to show which each one represents. When a liquid of composition $a$ is cooled, solid silver with dissolved tin begins to precipitate at $a_{1}$ and the sample solidifies completely at $a_{2}$.\\
\includegraphics[max width=\textwidth, center]{2025_02_27_d4b95d9718f0b9e2e6b9g-210}

Figure 5D. 3 The phase diagram for silver/tin discussed in Brief illustration 5D.1.

Monitoring the cooling curves at different overall compositions gives a clear indication of the structure of the phase diagram. The solid-liquid boundary is given by the points at which the rate of cooling changes. The longest eutectic halt gives the location of the eutectic composition and its melting temperature.

\subsection*{5D. 2 Reacting systems}
Many binary mixtures react to produce compounds, and technologically important examples of this behaviour include the Group 13/15 (III/V) semiconductors, such as the gallium arsenide system, which forms the compound GaAs. Although three constituents are present, there are only two components because GaAs is formed from the reaction $\mathrm{Ga}+\mathrm{As} \rightarrow \mathrm{GaAs}$. To illustrate some of the principles involved, consider a system that forms a compound C that also forms eutectic mixtures with the species A and B (Fig. 5D.4).\\
\includegraphics[max width=\textwidth, center]{2025_02_27_d4b95d9718f0b9e2e6b9g-210(1)}

Figure 5D. 4 The phase diagram for a system in which $A$ and $B$ react to form a compound $C=A B$. This resembles two versions of Fig. 5D. 1 in each half of the diagram. The constituent $C$ is a true compound, not just an equimolar mixture.

A system prepared by mixing an excess of B with A consists of C and unreacted B. This is a binary B,C system, which in this illustration is supposed to form a eutectic. The principal change from the eutectic phase diagram in Fig. 5D. 1 is that the whole of the phase diagram is squeezed into the range of compositions lying between equal amounts of $A$ and B ( $x_{\mathrm{B}}=0.5$, marked C in Fig. 5D.4) and pure B. The interpretation of the information in the diagram is obtained in the same way as for Fig. 5D.1. The solid deposited on cooling along the isopleth $a$ is the compound C. At temperatures below $a_{4}$ there are two solid phases, one consisting of C and the other of $B$. The pure compound $C$ melts congruently, that is, the composition of the liquid it forms is the same as that of the solid compound.

\subsection*{5D. 3 Incongruent melting}
In some cases the compound C is not stable as a liquid. An example is the alloy $\mathrm{Na}_{2} \mathrm{~K}$, which survives only as a solid (Fig. 5D.5). Consider what happens as a liquid at $a_{1}$ is cooled:

\begin{itemize}
  \item $a_{2} \rightarrow a_{3}$. A solid solution rich in Na is deposited, and the remaining liquid is richer in K .
  \item Below $a_{3}$. The sample is now entirely solid and consists of a solid solution rich in Na and solid $\mathrm{Na}_{2} \mathrm{~K}$.
\end{itemize}

Now consider the isopleth through $b_{1}$ :

\begin{itemize}
  \item $b_{1} \rightarrow b_{2}$. No obvious change occurs until the phase boundary is reached at $b_{2}$ when a solid solution rich in Na begins to deposit.
  \item $b_{2} \rightarrow b_{3}$. A solid solution rich in Na deposits, but at $b_{3}$ a reaction occurs to form $\mathrm{Na}_{2} \mathrm{~K}$ : this compound is formed by the K atoms diffusing into the solid Na .\\
\includegraphics[max width=\textwidth, center]{2025_02_27_d4b95d9718f0b9e2e6b9g-211}
\end{itemize}

Figure 5D. 5 The phase diagram for an actual system (sodium and potassium) like that shown in Fig. 5D.4, but with two differences. One is that the compound is $\mathrm{Na}_{2} \mathrm{~K}$, corresponding to $\mathrm{A}_{2} \mathrm{~B}$ and not $A B$ as in that illustration. The second is that the compound exists only as the solid, not as the liquid. The transformation of the compound at its melting point is an example of incongruent melting.

\begin{itemize}
  \item At $b_{3}$, three phases are in mutual equilibrium: the liquid, the compound $\mathrm{Na}_{2} \mathrm{~K}$, and a solid solution rich in Na . The horizontal line representing this three-phase equilibrium is called a peritectic line. At this stage the liquid $\mathrm{Na} / \mathrm{K}$ mixture is in equilibrium with a little solid $\mathrm{Na}_{2} \mathrm{~K}$, but there is still no liquid compound.
  \item $b_{3} \rightarrow b_{4}$. As cooling continues, the amount of solid compound increases until at $b_{4}$ the liquid reaches its eutectic composition. It then solidifies to give a two-phase solid consisting of a solid solution rich in K and solid $\mathrm{Na}_{2} \mathrm{~K}$.
\end{itemize}

If the solid is reheated, the sequence of events is reversed. No liquid $\mathrm{Na}_{2} \mathrm{~K}$ forms at any stage because it is too unstable to exist as a liquid. This behaviour is an example of incongruent melting, in which a compound melts into its components and does not itself form a liquid phase.

\subsection*{Checklist of concepts}
\begin{enumerate}
  \item At the eutectic composition the liquid phase solidifies without change of composition.
  \item A peritectic line in a phase diagram represents an equilibrium between three phases.
\end{enumerate}

\section*{TOPIC 5E Phase diagrams of ternary}
\subsection*{systems}
\subsection*{Why do you need to know this material?}
Ternary phase diagrams have become important in materials science as more complex materials are investigated, such as the ceramics found to have superconducting properties.

\subsection*{What is the key idea?}
A phase diagram is a map showing the conditions under which each phase of a system is the most stable.

\subsection*{What do you need to know already?}
It would be helpful to review the interpretation of twocomponent phase diagrams (Topics 5C and 5D) and the phase rule (Topic 4A). The interpretation of the phase diagrams presented here uses the lever rule (Topic 5C).

Consider the phases of a ternary system, a system with three components so $C=3$. In terms of the phase rule (Topic 4A), $F=5-P$. If the system is restricted to constant temperature and pressure, two degrees of freedom are discarded and $F^{\prime \prime}=3-P$. If two phases are present $(P=2)$, then $F^{\prime \prime}=1$ and the system has one degree of freedom: changing the amount of one component results in changes in the amounts of the other two components. This condition is represented by an area in the phase diagram. If three phases are present $(P=3)$, then $F^{\prime \prime}=0$, and the system is represented by a single point on the phase diagram.

Lines in ternary phase diagrams represent conditions under which two phases may coexist. Two phases are in equilibrium when they are connected by tie lines, as in binary phase diagrams.

\subsection*{5E. 1 Triangular phase diagrams}
The mole fractions of the three components of a ternary system satisfy $x_{\mathrm{A}}+x_{\mathrm{B}}+x_{\mathrm{C}}=1$. A phase diagram drawn as an equilateral triangle ensures that this property is satisfied automatically because the sum of the distances to a point inside an equilateral triangle of side 1 and measured parallel to the edges is equal to 1 (Fig. 5E.1).

Figure 5E. 1 shows how this approach works in practice. The edge AB corresponds to $x_{\mathrm{C}}=0$, and likewise for the other two edges. Hence, each of the three edges corresponds to one of\\
\includegraphics[max width=\textwidth, center]{2025_02_27_d4b95d9718f0b9e2e6b9g-212}

Figure 5E. 1 The triangular coordinates used for the discussion of three-component systems. Each edge corresponds to a binary system. All points along the dotted line a correspond to mole fractions of $C$ and $B$ in the same ratio.\\
the three binary systems (A,B), (B,C), and (C,A). An interior point corresponds to a system in which all three components are present. The point P , for instance, represents $x_{\mathrm{A}}=0.50, x_{\mathrm{B}}=$ $0.10, x_{\mathrm{C}}=0.40$.

Any point on a straight line joining the A apex to a point on the opposite edge (the dotted line $a$ in Fig. 5E.1) represents a composition that is progressively richer in A the closer the point is to the A apex, but for which the concentration ratio $\mathrm{B}: \mathrm{C}$ remains constant. Therefore, to represent the changing composition of a system as A is added, draw a line from the A apex to the point on BC representing the initial binary system. Any ternary system formed by adding A then lies at some point on this line.

\subsection*{Brief illustration 5E. 1}
The following points are represented on Fig. 5E.2.

\begin{center}
\begin{tabular}{llll}
\hline
Point & $x_{\mathrm{A}}$ & $x_{\mathrm{B}}$ & $\boldsymbol{x}_{\mathrm{C}}$ \\
\hline
$a$ & 0.20 & 0.80 & 0 \\
$b$ & 0.42 & 0.26 & 0.32 \\
$c$ & 0.80 & 0.10 & 0.10 \\
$d$ & 0.10 & 0.20 & 0.70 \\
$e$ & 0.20 & 0.40 & 0.40 \\
$f$ & 0.30 & 0.60 & 0.10 \\
\hline
\end{tabular}
\end{center}

Note that the points $d, e$, and $f$ have $x_{\mathrm{A}} / x_{\mathrm{B}}=0.50$ and lie on a straight line.\\
\includegraphics[max width=\textwidth, center]{2025_02_27_d4b95d9718f0b9e2e6b9g-213(1)}

Figure 5E. 2 The points referred to in Brief illustration 5E.1.\\
\includegraphics[max width=\textwidth, center]{2025_02_27_d4b95d9718f0b9e2e6b9g-213(2)}

Figure 5E. 3 When temperature is included as a variable, the phase diagram becomes a triangular prism. Horizontal sections through the prism correspond to the triangular phase diagrams being discussed and illustrated in Fig. 5E.1.

A single triangle represents the equilibria when one of the discarded degrees of freedom (the temperature, for instance) has a certain value. Different temperatures give rise to different equilibrium behaviour and therefore different triangular phase diagrams. Each one may therefore be regarded as a horizontal slice through a three-dimensional triangular prism, such as that shown in Fig. 5E.3.

\subsection*{5E. 2 Ternary systems}
Ternary phase diagrams are widely used in metallurgy and materials science. Although they can become quite complex, they can be interpreted in much the same way as binary diagrams.

\subsection*{(a) Partially miscible liquids}
The phase diagram for a ternary system in which W (in due course: water) and E (in due course: ethanoic acid (acetic acid)) are fully miscible, E and T (in due course: trichloromethane (chloroform)) are fully miscible, but W and T are only par-\\
tially miscible is shown in Fig. 5E.4. This illustration is for the system water/ethanoic acid/trichloromethane at room temperature, which behaves in this way:

\begin{itemize}
  \item The two fully miscible pairs, (E,W) and (E,T), form single-phase regions.
  \item The (W,T) system (along the base of the triangle) has a two-phase region.
\end{itemize}

The base of the triangle corresponds to one of the horizontal lines in a two-component phase diagram. The tie lines in the two-phase regions are constructed experimentally by determining the compositions of the two phases that are in equilibrium, marking them on the diagram, and then joining them with a straight line.

A single-phase system is formed when enough $E$ is added to the binary ( $\mathrm{W}, \mathrm{T}$ ) mixture. This effect is illustrated by following the line $a$ in Fig. 5E.4:

\begin{itemize}
  \item $a_{1}$. The system consists of two phases and the relative amounts of the two phases can be read off by using the lever rule.
  \item $a_{1} \rightarrow a_{2}$. The addition of E takes the system along the line joining $a_{1}$ to the E apex. At $a_{2}$ the solution still has two phases, but there is slightly more W in the largely T phase (represented by the point $a_{2}^{\prime \prime}$ ) and more T in the largely W phase ( $a_{2}^{\prime}$ ) because the presence of E helps both to dissolve. The phase diagram shows that there is more E in the W -rich phase than in the T-rich phase ( $a_{2}^{\prime}$ is closer than $a_{2}^{\prime \prime}$ to the E apex).
  \item $a_{2} \rightarrow a_{3}$. At $a_{3}$ two phases are present, but the T-rich layer is present only as a trace (lever rule).
  \item $a_{3} \rightarrow a_{4}$. Further addition of E takes the system towards and beyond $a_{4}$, and only a single phase is present.\\
\includegraphics[max width=\textwidth, center]{2025_02_27_d4b95d9718f0b9e2e6b9g-213}
\end{itemize}

Figure 5E. 4 The phase diagram, at fixed temperature and pressure, of the three-component system ethanoic acid (E), trichloromethane (T), and water (W). Only some of the tie lines have been drawn in the two-phase region. All points along the line $a$ correspond to trichloromethane and water present in the same ratio.

\subsection*{Brief illustration 5E. 2}
Consider a mixture of water (W in Fig. 5E.4) and trichloromethane (T) with $x_{\mathrm{W}}=0.40$ and $x_{\mathrm{T}}=0.60$, and ethanoic acid (E) is added to it. The relative proportions of W and T remain constant, so the point representing the overall composition moves along the straight line $b$ from $x_{\mathrm{T}}=0.60$ on the base to the ethanoic acid apex. The initial composition is in a twophase region: one phase has the composition $\left(x_{\mathrm{W}}, x_{\mathrm{T}}, x_{\mathrm{E}}\right)=$ $(0.05,0.95,0)$ and the other has composition $\left(x_{\mathrm{W}}, x_{\mathrm{T}}, x_{\mathrm{E}}\right)=$ ( $0.88,0.12,0$ ). When sufficient ethanoic acid has been added to raise its mole fraction to 0.18 the system consists of two phases of composition ( $0.07,0.82,0.11$ ) and ( $0.57,0.20,0.23$ ) in the ratio 1:3.

The point marked P in Fig. 5E. 4 is called the plait point: at this point the compositions of the two phases in equilibrium become identical. It is yet another example of a critical point. For convenience, the general interpretation of a triangular phase diagram is summarized in Fig. 5E.5.\\
\includegraphics[max width=\textwidth, center]{2025_02_27_d4b95d9718f0b9e2e6b9g-214}

Figure 5E. 5 The interpretation of a triangular phase diagram. The region inside the curved line consists of two phases, and the compositions of the two phases in equilibrium are given by the points at the ends of the tie lines (the tie lines are determined experimentally).

\subsection*{(b) Ternary solids}
The triangular phase diagram in Fig. 5E. 6 is typical of that for a solid alloy with varying compositions of three metals, $\mathrm{A}, \mathrm{B}$, and C.

\subsection*{Brief illustration 5E. 3}
Figure 5E. 6 is a simplified version of the phase diagram for a stainless steel consisting of iron, chromium, and nickel. The axes denote the mass percentage compositions instead of the mole fractions, but as the three percentages add up to 100 per cent, the interpretation of points in the triangle is essentially the same as for mole fractions. The point $a$ corresponds to the composition 74 per cent Fe , 18 per cent Cr , and 8 per cent Ni . It corresponds to the most common form of stainless steel, '18-8 stainless steel'. The composition corresponding to point $b$ lies in the two-phase region, one phase consisting of Cr and the other of the alloy $\gamma$-FeNi.\\
\includegraphics[max width=\textwidth, center]{2025_02_27_d4b95d9718f0b9e2e6b9g-214(1)}

Figure 5E. 6 A simplified triangular phase diagram of the ternary system represented by a stainless steel composed of iron, chromium, and nickel.

\subsection*{Checklist of concepts}
\begin{enumerate}
  \item A phase diagram drawn as an equilateral triangle ensures that the property $x_{\mathrm{A}}+x_{\mathrm{B}}+x_{\mathrm{C}}=1$ is satisfied automatically.
  \item At the plait point, the compositions of the two phases in equilibrium are identical.
\end{enumerate}

\section*{TOPIC 5F Activities}
\subsection*{Why do you need to know this material?}
The concept of an ideal solution is a good starting point for the discussion of mixtures, but to understand real solutions it is important to be able to describe deviations from ideal behaviour and to express them in terms of molecular interactions.

\subsection*{What is the key idea?}
The activity of a species, its effective concentration, helps to preserve the form of the expressions derived on the basis of ideal behaviour but extends their reach to real mixtures.

\subsection*{What do you need to know already?}
This Topic is based on the expression for chemical potential of a species derived from Raoult's and Henry's laws (Topic 5A). It also uses the formulation of a model of a regular solution introduced in Topic 5B.

This Topic shows how to adjust the expressions developed in Topics 5A and 5B to take into account deviations from ideal behaviour. As in other Topics collected in this Focus, the solvent is denoted by A , the solute by B , and a general component by J.

\subsection*{5F. 1 The solvent activity}
The general form of the chemical potential of a real or ideal solvent is given by a straightforward modification of eqn 5A. 21 $\left(\mu_{\mathrm{A}}=\mu_{\mathrm{A}}^{*}+R T \ln \left(p_{\mathrm{A}} / p_{\mathrm{A}}^{*}\right)\right)$, where $p_{\mathrm{A}}^{*}$ is the vapour pressure of pure A and $p_{\mathrm{A}}$ is the vapour pressure of A when it is a component of a solution. The solvent in an ideal solution obeys Raoult's law (Topic $5 \mathrm{~A}, p_{\mathrm{A}}=x_{\mathrm{A}} p_{\mathrm{A}}^{*}$ ) at all concentrations and the chemical potential is expressed as eqn $5 \mathrm{~A} .23\left(\mu_{\mathrm{A}}=\mu_{\mathrm{A}}^{\star}+R T \ln x_{\mathrm{A}}\right)$. The form of this relation can be preserved when the solution does not obey Raoult's law by writing

\[
\mu_{\mathrm{A}}=\mu_{\mathrm{A}}^{*}+R T \ln a_{\mathrm{A}} \quad \begin{align*}
& \text { Activity of solvent }  \tag{5F.1}\\
& \text { [definition] }
\end{align*}
\]

The quantity $a_{\mathrm{A}}$ is the activity of A , a kind of 'effective' mole fraction.

Because eqn 5F. 1 is true for both real and ideal solutions, comparing it with $\mu_{\mathrm{A}}=\mu_{\mathrm{A}}^{*}+R T \ln \left(p_{\mathrm{A}} / p_{\mathrm{A}}^{*}\right)$ gives

\[
\begin{array}{ll}
a_{\mathrm{A}}=\frac{p_{\mathrm{A}}}{p_{\mathrm{A}}^{\star}} & \begin{array}{l}
\text { Activity of solvent } \\
{[\text { measurement }]}
\end{array} \tag{5F.2}
\end{array}
\]

There is nothing mysterious about the activity of a solvent: it can be determined experimentally simply by measuring the vapour pressure and then using this relation.

\subsection*{Brief illustration 5F. 1}
The vapour pressure of $0.500 \mathrm{~mol} \mathrm{dm}^{-3} \mathrm{KNO}_{3}(\mathrm{aq})$ at $100^{\circ} \mathrm{C}$ is 99.95 kPa , and the vapour pressure of pure water at this temperature is $1.00 \mathrm{~atm}(101 \mathrm{kPa})$. It follows that the activity of water in this solution at this temperature is

$$
a_{\mathrm{A}}=\frac{99.95 \mathrm{kPa}}{101 \mathrm{kPa}}=0.990
$$

Because all solvents obey Raoult's law more closely as the concentration of solute approaches zero, the activity of the solvent approaches the mole fraction as $x_{\mathrm{A}} \rightarrow 1$ :


\begin{equation*}
a_{\mathrm{A}} \rightarrow x_{\mathrm{A}} \text { as } x_{\mathrm{A}} \rightarrow 1 \tag{5F.3}
\end{equation*}


A convenient way of expressing this convergence is to introduce the activity coefficient, $\gamma$ (gamma), by the definition

\[
a_{\mathrm{A}}=\gamma_{\mathrm{A}} x_{\mathrm{A}} \quad \gamma_{\mathrm{A}} \rightarrow 1 \text { as } x_{\mathrm{A}} \rightarrow 1 \quad \begin{align*}
& \text { Activity coefficient }  \tag{5F.4}\\
& \text { of solvent } \\
& \text { [definition] }
\end{align*}
\]

at all temperatures and pressures. The chemical potential of the solvent is then

\[
\mu_{\mathrm{A}}=\mu_{\mathrm{A}}^{\star}+R T \ln x_{\mathrm{A}}+R T \ln \gamma_{\mathrm{A}} \quad \begin{align*}
& \text { Chemical potential }  \tag{5F.5}\\
& \text { of solvent }
\end{align*}
\]

The standard state of the solvent is established when $x_{A}=1$ (the pure solvent) and the pressure is 1 bar.

\subsection*{5F.2 The solute activity}
The problem with defining activity coefficients and standard states for solutes is that they approach ideal-dilute (Henry's law) behaviour as $x_{\mathrm{B}} \rightarrow 0$, not as $x_{\mathrm{B}} \rightarrow 1$ (corresponding to pure solute).

\subsection*{(a) Ideal-dilute solutions}
A solute B that satisfies Henry's law (Topic 5A) has a vapour pressure given by $p_{\mathrm{B}}=K_{\mathrm{B}} x_{\mathrm{B}}$, where $K_{\mathrm{B}}$ is an empirical constant. In this case, the chemical potential of $B$ is\\
$\mu_{\mathrm{B}}=\mu_{\mathrm{B}}^{\star}+R T \ln \frac{p_{\mathrm{B}}}{p_{\mathrm{B}}^{\star}} \stackrel{\overbrace{p_{\mathrm{B}}=K_{\mathrm{B}} x_{\mathrm{B}}}^{=}}{\mu_{\mathrm{B}}^{\star}+R T \ln \frac{K_{\mathrm{B}} x_{\mathrm{B}}}{p_{\mathrm{B}}^{\star}}=\overbrace{\mu_{\mathrm{B}}^{\star}+R T \ln \frac{K_{\mathrm{B}}}{p_{\mathrm{B}}^{*}}+R T \ln x_{\mathrm{B}}}^{\mu_{\mathrm{B}}^{\ominus}}{ }^{\rho_{0}}}$\\
(5F.6)\\
Both $K_{\mathrm{B}}$ and $p_{\mathrm{B}}^{*}$ are characteristics of the solute, so the two blue terms may be combined to give a new standard chemical potential, $\mu_{\mathrm{B}}^{\ominus}$


\begin{equation*}
\mu_{\mathrm{B}}^{\ominus}=\mu_{\mathrm{B}}^{\star}+R T \ln \frac{K_{\mathrm{B}}}{p_{\mathrm{B}}^{\star}} \tag{5F.7}
\end{equation*}


It then follows that the chemical potential of a solute in an ideal-dilute solution is related to its mole fraction by


\begin{equation*}
\mu_{\mathrm{B}}=\mu_{\mathrm{B}}^{\ominus}+R T \ln x_{\mathrm{B}} \tag{5F.8}
\end{equation*}


If the solution is ideal, $K_{\mathrm{B}}=p_{\mathrm{B}}^{*}$ (Raoult's law) and eqn 5F. 7 reduces to $\mu_{\mathrm{B}}^{\ominus}=\mu_{\mathrm{B}}^{\star}$, as expected.

\subsection*{Brief illustration 5F. 2}
In Example 5A. 4 it is established that in a mixture of propanone and trichloromethane at $298 \mathrm{~K} K_{\text {propanone }}=24.5 \mathrm{kPa}$, whereas $p_{\text {propanone }}^{*}=46.3 \mathrm{kPa}$. It follows from eqn 5F. 7 that

$$
\begin{aligned}
\mu_{\text {propanone }}^{\ominus} & =\mu_{\text {propanone }}^{\star}+R T \ln \frac{24.5 \mathrm{kPa}}{46.3 \mathrm{kPa}} \\
& =\mu_{\text {propanone }}^{\star}+\left(8.3145 \mathrm{~J} \mathrm{~K}^{-1} \mathrm{~mol}^{-1}\right) \times(298 \mathrm{~K}) \times \ln \frac{24.5}{46.3} \\
& =\mu_{\text {propanone }}^{\star}-1.58 \mathrm{~kJ} \mathrm{~mol}^{-1}
\end{aligned}
$$

and the standard value differs from the value for the pure liquid by $-1.58 \mathrm{~kJ} \mathrm{~mol}^{-1}$.

\subsection*{(b) Real solutes}
Real solutions deviate from ideal-dilute, Henry's law behaviour. For the solute, the introduction of $a_{\mathrm{B}}$ in place of $x_{\mathrm{B}}$ in eqn 5F. 8 gives

$$
\mu_{\mathrm{B}}=\mu_{\mathrm{B}}^{\ominus}+R T \ln a_{\mathrm{B}} \quad \begin{aligned}
& \text { Chemical potential of solute } \\
& \text { [definition] }
\end{aligned}
$$

(5F.9)

The standard state remains unchanged in this last stage, and all the deviations from ideality are captured in the activity $a_{\mathrm{B}}$.

It remains true that $\mu_{\mathrm{B}}=\mu_{\mathrm{B}}^{\star}+R T \ln \left(p_{\mathrm{B}} / p_{\mathrm{B}}^{\star}\right)$, but now, from eqn 5 F .7 written as $\mu_{\mathrm{B}}^{\star}=\mu_{\mathrm{B}}^{\ominus}-R T \ln \left(K_{\mathrm{B}} / p_{\mathrm{B}}^{*}\right)$ it follows that

$$
\mu_{\mathrm{B}}=\overbrace{\mu_{\mathrm{B}}^{\odot}-R T \ln \frac{K_{\mathrm{B}}}{p_{\mathrm{B}}^{\star}}}^{\mu_{\mathrm{B}}^{\star}}+R T \ln \frac{p_{\mathrm{B}}}{p_{\mathrm{B}}^{\star}}=\mu_{\mathrm{B}}^{\star}+R T \ln \frac{p_{\mathrm{B}}}{K_{\mathrm{B}}}
$$

Comparison of this expression with eqn 5F. 9 identifies the activity $a_{\mathrm{B}}$ as

$$
\begin{array}{ll}
a_{\mathrm{B}}=\frac{p_{\mathrm{B}}}{K_{\mathrm{B}}} & \begin{array}{l}
\text { Activity of solute } \\
\text { [measurement] }
\end{array}
\end{array}
$$

(5F.10)

As for the solvent, it is sensible to introduce an activity coefficient through

\[
\begin{array}{ll}
a_{\mathrm{B}}=\gamma_{\mathrm{B}} x_{\mathrm{B}} & \begin{array}{l}
\text { Activity coefficient of solute } \\
\text { [definition] }
\end{array} \tag{5F.11}
\end{array}
\]

Now all the deviations from ideality are captured in the activity coefficient $\gamma_{\mathrm{B}}$. Because the solute obeys Henry's law ( $p_{\mathrm{B}}=K_{\mathrm{B}} x_{\mathrm{B}}$ ) as its concentration goes to zero. It follows that


\begin{equation*}
a_{\mathrm{B}} \rightarrow x_{\mathrm{B}} \text { and } \gamma_{\mathrm{B}} \rightarrow 1 \text { as } x_{\mathrm{B}} \rightarrow 0 \tag{5F.12}
\end{equation*}


at all temperatures and pressures. Deviations of the solute from ideality disappear as its concentration approaches zero.

\subsection*{Example 5F. 1 Measuring activity}
Use the following information to calculate the activity and activity coefficient of trichloromethane (chloroform, C) in propanone (acetone, A) at $25^{\circ} \mathrm{C}$, treating it first as a solvent and then as a solute.

\begin{center}
\begin{tabular}{lcccccc}
$x_{\mathrm{C}}$ & 0 & 0.20 & 0.40 & 0.60 & 0.80 & 1 \\
$p_{\mathrm{C}} / \mathrm{kPa}$ & 0 & 4.7 & 11 & 18.9 & 26.7 & 36.4 \\
$p_{\mathrm{A}} / \mathrm{kPa}$ & 46.3 & 33.3 & 23.3 & 12.3 & 4.9 & 0 \\
\end{tabular}
\end{center}

Collect your thoughts For the activity of chloroform as a solvent (the Raoult's law activity), write $a_{\mathrm{C}}=p_{\mathrm{C}} / p_{\mathrm{C}}^{\star}$ and $\gamma_{\mathrm{C}}=$ $a_{\mathrm{C}} / x_{\mathrm{C}}$. For its activity as a solute (the Henry's law activity), write $a_{\mathrm{C}}=p_{\mathrm{C}} / K_{\mathrm{C}}$ and $\gamma_{\mathrm{C}}=a_{\mathrm{C}} / x_{\mathrm{C}}$ with the new activity.

The solution Because $p_{\mathrm{C}}^{\star}=36.4 \mathrm{kPa}$ and $K_{\mathrm{C}}=23.5 \mathrm{kPa}$ (from Example 5A.4), construct the following tables. For instance, at $x_{\mathrm{C}}=0.20$, in the Raoult's law case $a_{\mathrm{C}}=(4.7 \mathrm{kPa}) /(36.4 \mathrm{kPa})=$ 0.13 and $\gamma_{\mathrm{C}}=0.13 / 0.20=0.65$; likewise, in the Henry's law case, $a_{\mathrm{C}}=(4.7 \mathrm{kPa}) /(23.5 \mathrm{kPa})=0.20$ and $\gamma_{\mathrm{C}}=0.20 / 0.20=1.0$.

From Raoult's law (chloroform regarded as the solvent):

\begin{center}
\begin{tabular}{lllllll}
$x_{\mathrm{C}}$ & 0 & 0.20 & 0.40 & 0.60 & 0.80 & 1 \\
$a_{\mathrm{C}}$ & 0 & 0.13 & 0.30 & 0.52 & 0.73 & 1.00 \\
$\gamma_{\mathrm{C}}$ &  & 0.65 & 0.75 & 0.87 & 0.92 & 1.00 \\
\end{tabular}
\end{center}

From Henry's law (chloroform regarded as the solute):

\begin{center}
\begin{tabular}{lllllll}
$x_{\mathrm{C}}$ & 0 & 0.20 & 0.40 & 0.60 & 0.80 & 1 \\
$a_{\mathrm{C}}$ & 0 & 0.20 & 0.47 & 0.80 & 1.14 & 1.55 \\
$\gamma_{\mathrm{C}}$ & 1 & 1.00 & 1.17 & 1.34 & 1.42 & 1.55 \\
\end{tabular}
\end{center}

These values are plotted in Fig. 5F.1. Notice that $\gamma_{\mathrm{C}} \rightarrow 1$ as $x_{\mathrm{C}} \rightarrow 1$ in the Raoult's law case, but that $\gamma_{\mathrm{C}} \rightarrow 1$ as $x_{\mathrm{C}} \rightarrow 0$ in the Henry's law case.\\
\includegraphics[max width=\textwidth, center]{2025_02_27_d4b95d9718f0b9e2e6b9g-217(1)}

Figure 5F. 1 The variation of activity and activity coefficient for a trichloromethane/propanone (chloroform/acetone) mixture with composition according to (a) Raoult's law, (b) Henry's law.

Self-test 5F. 1 Calculate the activities and activity coefficients for acetone according to the two conventions (use $p_{\mathrm{A}}^{*}=$ 46.3 kPa and $K_{\mathrm{A}}=24.5 \mathrm{kPa}$ )\\
\includegraphics[max width=\textwidth, center]{2025_02_27_d4b95d9718f0b9e2e6b9g-217}

\subsection*{(c) Activities in terms of molalities}
The selection of a standard state is entirely arbitrary and can be chosen in a way that suits the description of the composition of the system best. Because compositions are often expressed as molalities, $b$, in place of mole fractions (see The chemist's toolkit 11 in Topic 5A) it is then convenient to write


\begin{equation*}
\mu_{\mathrm{B}}=\mu_{\mathrm{B}}^{\ominus}+R T \ln \frac{b_{\mathrm{B}}}{b^{\ominus}} \tag{5F.13}
\end{equation*}


where $\mu_{\mathrm{B}}^{\ominus}$ has a different value from the standard value introduced earlier. According to this definition, the chemical potential of the solute has its standard value $\mu_{\mathrm{B}}^{\ominus}$ when the molality of B is $b^{\ominus}$ (i.e. at $1 \mathrm{~mol} \mathrm{~kg}^{-1}$ ). Note that as $b_{\mathrm{B}} \rightarrow 0, \mu_{\mathrm{B}} \rightarrow-\infty$; that is, as the solution becomes diluted, so the solute becomes increasingly thermodynamically stable. The practical consequence of this result is that it is very difficult to remove the last traces of a solute from a solution.

As before, deviations from ideality are incorporated by introducing a dimensionless activity $a_{\mathrm{B}}$ and a dimensionless activity coefficient $\gamma_{\mathrm{B}}$, and writing


\begin{equation*}
a_{\mathrm{B}}=\gamma_{\mathrm{B}} \frac{b_{\mathrm{B}}}{b^{\ominus}} \text {, where } \gamma_{\mathrm{B}} \rightarrow 1 \text { as } b_{\mathrm{B}} \rightarrow 0 \tag{5F.14}
\end{equation*}


at all temperatures and pressures. The standard state remains unchanged in this last stage and, as before, all the deviations from ideality are captured in the activity coefficient $\gamma_{B}$. The final expression for the chemical potential of a real solute at any molality is then


\begin{equation*}
\mu_{\mathrm{B}}=\mu_{\mathrm{B}}^{\ominus}+R T \ln a_{\mathrm{B}} \tag{5F.15}
\end{equation*}


\subsection*{5F. 3 The activities of regular solutions}
The concept of regular solutions (Topic 5B) gives further insight into the origin of deviations from Raoult's law and its relation to activity coefficients. The starting point is the model expression for the excess enthalpy (eqn 5B.6, $H^{\mathrm{E}}=$ $n \xi R T x_{\mathrm{A}} x_{\mathrm{B}}$ ) and its implication for the Gibbs energy of mixing for a regular solution (eqn 5B.7, $\Delta_{\text {mix }} G=n R T\left\{x_{\mathrm{A}} \ln x_{\mathrm{A}}+\right.$ $\left.x_{\mathrm{B}} \ln x_{\mathrm{B}}+\xi x_{\mathrm{A}} x_{\mathrm{B}}\right\}$ ). On the basis of this model it is possible to develop expressions for the activity coefficients in terms of the parameter $\xi$.

\subsection*{How is that done? 5 F. 1}
Developing expressions for the activity coefficients of a regular solution

The Gibbs energy of mixing to form an ideal solution is given in eqn. 5B.3:

$$
\Delta_{\text {mix }} G=n R T\left\{x_{\mathrm{A}} \ln x_{\mathrm{A}}+x_{\mathrm{B}} \ln x_{\mathrm{B}}\right\}
$$

The corresponding expression for a non-ideal solution is

$$
\Delta_{\text {mix }} G=n R T\left\{x_{\mathrm{A}} \ln a_{\mathrm{A}}+x_{\mathrm{B}} \ln a_{\mathrm{B}}\right\}
$$

This relation follows in the same way as for an ideal mixture but with activities in place of mole fractions. However, in Topic 5B. 7 it is established (in eqn 5B.7) that for a regular solution

$$
\Delta_{\operatorname{mix}} G=n R T\left\{x_{\mathrm{A}} \ln x_{\mathrm{A}}+x_{\mathrm{B}} \ln x_{\mathrm{B}}+\xi x_{\mathrm{A}} x_{\mathrm{B}}\right\}
$$

The last two equations can be made consistent as follows. First replace each activity by $\gamma_{j} x_{\mathrm{J}}$ :

$$
\begin{aligned}
\Delta_{\text {mix }} G & =n R T\left\{x_{\mathrm{A}} \ln x_{\mathrm{A}} \gamma_{\mathrm{A}}+x_{\mathrm{B}} \ln x_{\mathrm{B}} \gamma_{\mathrm{B}}\right\} \\
& =n R T\left\{x_{\mathrm{A}} \ln x_{\mathrm{A}}+x_{\mathrm{B}} \ln x_{\mathrm{B}}+x_{\mathrm{A}} \ln \gamma_{\mathrm{A}}+x_{\mathrm{B}} \ln \gamma_{\mathrm{B}}\right\}
\end{aligned}
$$

For consistency, the sum of the two terms in blue must be equal to $\xi_{x_{\mathrm{A}} x_{\mathrm{B}}}$, which can be achieved by writing $\ln \gamma_{\mathrm{A}}=\xi x_{\mathrm{B}}^{2}$ and $\ln \gamma_{B}=\xi x_{A}^{2}$, because then

$$
x_{\mathrm{A}} \overbrace{\ln \gamma_{\mathrm{A}}}^{\xi x_{\mathrm{B}}^{2}}+x_{\mathrm{B}} \overbrace{\ln \gamma_{\mathrm{B}}}^{\xi x_{\mathrm{A}}^{2}}=\xi x_{\mathrm{A}} x_{\mathrm{B}}^{2}+\xi x_{\mathrm{B}} x_{\mathrm{A}}^{2}=\xi \overbrace{\left(x_{\mathrm{A}}+x_{\mathrm{B}}\right)}^{1} x_{\mathrm{A}} x_{\mathrm{B}}=\xi x_{\mathrm{A}} x_{\mathrm{B}}
$$

It follows that the activity coefficients of a regular solution are given by what are known as the Margules equations:


\begin{equation*}
-\ln \gamma_{\mathrm{A}}=\xi x_{\mathrm{B}}^{2} \quad \ln \gamma_{\mathrm{B}}=\xi x_{\mathrm{A}}^{2} \tag{5F.16}
\end{equation*}


Margules equations

Note that the activity coefficients behave correctly for dilute solutions: $\gamma_{\mathrm{A}} \rightarrow 1$ as $x_{\mathrm{B}} \rightarrow 0$ and $\gamma_{\mathrm{B}} \rightarrow 1$ as $x_{\mathrm{A}} \rightarrow 0$. Also note that $A$ and $B$ are treated here as equal components of a mixture, not as solvent and solute.

At this point the Margules equations can be used to write the activity of $A$ as\\
\includegraphics[max width=\textwidth, center]{2025_02_27_d4b95d9718f0b9e2e6b9g-218(2)}\\
with a similar expression for $a_{\mathrm{B}}$. The activity of A , though, is just the ratio of the vapour pressure of $A$ in the solution to the vapour pressure of pure A (eqn 5F.2, $a_{\mathrm{A}}=p_{\mathrm{A}} / p_{\mathrm{A}}^{*}$ ), so


\begin{equation*}
p_{\mathrm{A}}=p_{\mathrm{A}}^{*} x_{\mathrm{A}} \mathrm{e}^{\xi\left(1-x_{\mathrm{A}}\right)^{2}} \tag{5F.18}
\end{equation*}


This function is plotted in Fig. 5F.2, and interpreted as follows:

\begin{itemize}
  \item When $\xi=0$, corresponding to an ideal solution, $p_{\mathrm{A}}=p_{\mathrm{A}}^{*} x_{\mathrm{A}}$, in accord with Raoult's law.
  \item Positive values of $\xi$ (endothermic mixing, unfavourable solute-solvent interactions) give vapour pressures higher than for an ideal solution.
  \item Negative values of $\xi$ (exothermic mixing, favourable solute-solvent interactions) give a vapour pressure lower than for an ideal solution.
\end{itemize}

All the plots of eqn 5F. 18 approach linearity and coincide with the Raoult's law line as $x_{\mathrm{A}} \rightarrow 1$ and the exponential function in eqn 5F. 18 approaches 1. When $x_{\mathrm{A}} \ll 1$, eqn 5F. 18 approaches


\begin{equation*}
p_{\mathrm{A}}=p_{\mathrm{A}}^{\star} x_{\mathrm{A}} \mathrm{e}^{\xi} \tag{5F.19}
\end{equation*}


\begin{center}
\includegraphics[max width=\textwidth]{2025_02_27_d4b95d9718f0b9e2e6b9g-218}
\end{center}

Figure 5F. 2 The vapour pressure of a mixture based on a model in which the excess enthalpy is $n \xi R T x_{\mathrm{A}} x_{\mathrm{B}}$; the lines are labelled with the value of $\xi$. An ideal solution corresponds to $\xi=0$ and gives a straight line, in accord with Raoult's law. Positive values of $\xi$ give vapour pressures higher than ideal. Negative values of $\xi$ give a lower vapour pressure.

This expression has the form of Henry's law once $K$ is identified with $\mathrm{e}^{\xi} p_{\mathrm{A}}^{\star}$, which is different for each solute-solvent system.

\subsection*{Brief illustration 5F. 3}
In Example 5B.1 of Topic 5B it is established that $\xi=1.13$ for a mixture of benzene and cyclohexane at $25^{\circ} \mathrm{C}$. Because $\xi>0$ the vapour pressure of the mixture is expected to be greater than its ideal value. The total vapour pressure of the mixture is therefore

$$
p=p_{\text {benzene }}^{*} x_{\text {benzene }} \mathrm{e}^{1.13\left(1-x_{\text {bexpene }}\right)^{2}}+p_{\text {cyydohexane }}^{*} x_{\text {cyclohexane }} \mathrm{e}^{1.13\left(1-x_{\text {gcdobexane }}\right)^{2}}
$$

This expression is plotted in Fig. 5F.3, using $p_{\text {benzene }}^{\star}=10.0 \mathrm{kPa}$ and $p_{\text {cyclohexane }}^{*}=10.4 \mathrm{kPa}$.\\
\includegraphics[max width=\textwidth, center]{2025_02_27_d4b95d9718f0b9e2e6b9g-218(1)}

Figure 5F. 3 The computed vapour pressure curves for a mixture of benzene and cyclohexane at $25^{\circ} \mathrm{C}$ as derived in Brief illustration 5F.3.

\subsection*{5F.4 The activities of ions}
Interactions between ions are so strong that the approximation of replacing activities by molalities is valid only in very dilute solutions (less than $1 \mathrm{mmolkg}^{-1}$ in total ion concentration), and in precise work activities themselves must be used.

If the chemical potential of the cation $\mathrm{M}^{+}$is denoted $\mu_{+}$and that of the anion $\mathrm{X}^{-}$is denoted $\mu_{-}$, the molar Gibbs energy of the ions in the electrically neutral solution is the sum of these partial molar quantities. The molar Gibbs energy of an ideal solution of such ions is


\begin{equation*}
G_{\mathrm{m}}^{\text {ideal }}=\mu_{+}^{\text {ideal }}+\mu_{-}^{\text {ideal }} \tag{5F.20}
\end{equation*}


with $\mu_{\mathrm{J}}^{\text {ideal }}=\mu_{\mathrm{J}}^{\ominus}+R T \ln x_{\mathrm{J}}$. However, for a real solution of $\mathrm{M}^{+}$and $\mathrm{X}^{-}$of the same molality it is necessary to write $\mu_{\mathrm{I}}=$ $\mu_{\mathrm{J}}^{\ominus}+R T \ln a_{\mathrm{J}}$ with $a_{\mathrm{J}}=\gamma_{\mathrm{J}} x_{\mathrm{J}}$, which implies that $\mu_{\mathrm{J}}=\mu_{\mathrm{J}}^{\text {ideal }}+$ $R T \ln \gamma_{\mathrm{j}}$. It then follows that


\begin{align*}
G_{\mathrm{m}} & =\mu_{+}+\mu_{-}=\mu_{+}^{\text {ideal }}+\mu_{-}^{\text {ideal }}+R T \ln \gamma_{+}+R T \ln \gamma_{-} \\
& =G_{\mathrm{m}}^{\text {ideal }}+R T \ln \gamma_{+} \gamma_{-} \tag{5F.21}
\end{align*}


All the deviations from ideality are contained in the last term.

\subsection*{(a) Mean activity coefficients}
There is no experimental way of separating the product $\gamma_{+} \gamma_{-}$into contributions from the cations and the anions. The best that can be done experimentally is to assign responsibility for the non-ideality equally to both kinds of ion. Therefore, the 'mean activity coefficient' is introduced as the geometric mean of the individual coefficients, where the geometric mean of $x^{p}$ and $y^{q}$ is $\left(x^{p} y^{q}\right)^{1 /(p+q)}$. For a 1,1-electrolyte $p=1, q=1$ and the required geometric mean is


\begin{equation*}
\gamma_{ \pm}=\left(\gamma_{+} \gamma_{-}\right)^{1 / 2} \tag{5F.22}
\end{equation*}


The individual chemical potentials of the ions are then written


\begin{equation*}
\mu_{+}=\mu_{+}^{\text {ideal }}+R T \ln \gamma_{ \pm} \quad \mu_{-}=\mu_{-}^{\text {ideal }}+R T \ln \gamma_{ \pm} \tag{5F.23}
\end{equation*}


The sum of these two chemical potentials is the same as before, eqn 5F.21, but now the non-ideality is shared equally.

To generalize this approach to the case of a compound $M_{p} X_{q}$ that dissolves to give a solution of $p$ cations and $q$ anions from each formula unit, the molar Gibbs energy of the ions is written as the sum of their partial molar Gibbs energies (i.e. their chemical potentials):


\begin{equation*}
G_{\mathrm{m}}=p \mu_{+}+q \mu_{-}=G_{\mathrm{m}}^{\mathrm{ideal}}+p R T \ln \gamma_{+}+q R T \ln \gamma_{-} \tag{5F.24}
\end{equation*}


The mean activity coefficient can now be defined in a more general way as

\[
\gamma_{ \pm}=\left(\gamma_{+}^{p} \gamma_{-}^{q}\right)^{1 / s} \quad s=p+q \quad \begin{align*}
& \text { Mean activity coefficient }  \tag{5F.25}\\
& \text { [definition] }
\end{align*}
\]

and the chemical potential of each ion written as


\begin{equation*}
\mu_{i}=\mu_{i}^{\text {ideal }}+R T \ln \gamma_{ \pm} \tag{5F.26}
\end{equation*}


\subsection*{(b) The Debye-Hückel limiting law}
The long range and strength of the Coulombic interaction between ions means that it is likely to be primarily responsible for the departures from ideality in ionic solutions and to dominate all the other contributions to non-ideality. This domination is the basis of the Debye-Hückel theory of ionic solutions, which was devised by Peter Debye and Erich Hückel in 1923. The following is a qualitative account of the theory and its principal conclusions. For a quantitative treatment, see A deeper look 1 on the website for this text.

Oppositely charged ions attract one another. As a result, anions are more likely to be found near cations in solution, and vice versa (Fig. 5F.4). Overall, the solution is electrically neutral, but near any given ion there is an excess of counter ions (ions of opposite charge). Averaged over time, counter ions are more likely to be found near any given ion. This time-averaged, spherical haze around the central ion, in which counter ions outnumber ions of the same charge as the central ion, has a net charge equal in magnitude but opposite in sign to that on the central ion, and is called its ionic atmosphere. The energy, and therefore the chemical potential, of any given central ion are lowered as a result of its electrostatic interaction with its ionic atmosphere. This lowering of energy appears as the difference\\
\includegraphics[max width=\textwidth, center]{2025_02_27_d4b95d9718f0b9e2e6b9g-219}

Figure 5F.4 The model underlying the Debye-Hückel theory is of a tendency for anions to be found around cations, and of cations to be found around anions (one such local clustering region is shown by the shaded sphere). The ions are in ceaseless motion, and the diagram represents a snapshot of their motion. The solutions to which the theory applies are far less concentrated than shown here.\\
between the molar Gibbs energy $G_{m}$ and the ideal value of the molar Gibbs energy $G_{\mathrm{m}}^{\text {ideal }}$ of the solute, and hence can be identified with the term $R T \ln \gamma_{ \pm}$in eqn 5F.21. The stabilization of ions by their interaction with their ionic atmospheres is part of the explanation why chemists commonly use dilute solutions, in which the stabilization is minimized, to achieve precipitation of ions from electrolyte solutions.

The model leads to the result that at very low concentrations the activity coefficient can be calculated from the DebyeHückel limiting law


\begin{equation*}
\log \gamma_{ \pm}=-A\left|z_{+} z_{-}\right| I^{1 / 2} \tag{5F.27}
\end{equation*}


Debye-Hückel limiting law\\
where $A=0.509$ for an aqueous solution at $25^{\circ} \mathrm{C}$ and $I$ is the dimensionless ionic strength of the solution:

\[
I=\frac{1}{2} \sum_{i} z_{i}^{2}\left(b_{i} / b^{\ominus}\right) \quad \begin{align*}
& \text { lonic strength }  \tag{5F.28}\\
& \text { [definition] }
\end{align*}
\]

In this expression $z_{i}$ is the charge number of an ion $i$ (positive for cations and negative for anions) and $b_{i}$ is its molality. The ionic strength occurs widely, and often as its square root (as in eqn 5F.27) wherever ionic solutions are discussed. The sum extends over all the ions present in the solution. For solutions consisting of two types of ion at molalities $b_{+}$and $b_{-}$,


\begin{equation*}
I=\frac{1}{2}\left(b_{+} z_{+}^{2}+b_{-} z_{-}^{2}\right) / b^{\ominus} \tag{5F.29}
\end{equation*}


The ionic strength emphasizes the charges of the ions because the charge numbers occur as their squares. Table 5 F. 1 summarizes the relation of ionic strength and molality in an easily usable form.

Table 5F.1 Ionic strength and molality, $I=k b / b^{\ominus}$

\begin{center}
\begin{tabular}{lrrrl}
\hline
$\boldsymbol{k}$ & $\mathrm{X}^{-}$ & $\mathrm{X}^{2-}$ & $\mathrm{X}^{3-}$ & $\mathrm{X}^{4-}$ \\
\hline
$\mathrm{M}^{+}$ & 1 & 3 & 6 & 10 \\
$\mathrm{M}^{2+}$ & 3 & 4 & 15 & 12 \\
$\mathrm{M}^{3+}$ & 6 & 15 & 9 & 42 \\
$\mathrm{M}^{4+}$ & 10 & 12 & 42 & 16 \\
\end{tabular}
\end{center}

For example, the ionic strength of an $\mathrm{M}_{2} \mathrm{X}_{3}$ solution of molality $b$, which is understood to give $\mathrm{M}^{3+}$ and $\mathrm{X}^{2-}$ ions in solution, is $15 b / b^{\ominus}$.

\subsection*{Brief illustration 5F. 4}
The mean activity coefficient of $5.0 \mathrm{mmol} \mathrm{kg}{ }^{-1} \mathrm{KCl}(\mathrm{aq})$ at $25^{\circ} \mathrm{C}$ is calculated by writing $I=\frac{1}{2}\left(b_{+}+b_{-}\right) / b^{\ominus}=b / b^{\ominus}$, where $b$ is the molality of the solution (and $b_{+}=b_{-}=b$ ). Then, from eqn 5F.27,

$$
\log \gamma_{ \pm}=-0.509 \times\left(5.0 \times 10^{-3}\right)^{1 / 2}=-0.03 \ldots
$$

Hence, $\gamma_{ \pm}=0.92$. The experimental value is 0.927 .

Table 5F. 2 Mean activity coefficients in water at $298 \mathrm{~K}^{*}$

\begin{center}
\begin{tabular}{lll}
\hline
$\boldsymbol{b} / \boldsymbol{b}^{\ominus}$ & KCl & $\mathrm{CaCl}_{2}$ \\
\hline
0.001 & 0.966 & 0.888 \\
0.01 & 0.902 & 0.732 \\
0.1 & 0.770 & 0.524 \\
1.0 & 0.607 & 0.725 \\
\hline
\end{tabular}
\end{center}

\begin{itemize}
  \item More values are given in the Resource section.
\end{itemize}

The name 'limiting law' is applied to eqn 5F. 27 because ionic solutions of moderate molalities may have activity coefficients that differ from the values given by this expression, but all solutions are expected to conform as $b \rightarrow 0$. Table 5F. 2 lists some experimental values of activity coefficients for salts of various valence types. Figure 5F. 5 shows some of these values plotted against $I^{1 / 2}$, and compares them with the theoretical straight lines calculated from eqn 5F.27. The agreement at very low molalities (less than about $1 \mathrm{mmol} \mathrm{kg}^{-1}$, depending on charge type) is impressive and convincing evidence in support of the model. Nevertheless, the departures from the theoretical curves above these molalities are large, and show that the approximations are valid only at very low concentrations.

\subsection*{(c) Extensions of the limiting law}
When the ionic strength of the solution is too high for the limiting law to be valid, the activity coefficient may be estimated\\
\includegraphics[max width=\textwidth, center]{2025_02_27_d4b95d9718f0b9e2e6b9g-220}

Figure 5F. 5 An experimental test of the Debye-Hückel limiting law. Although there are marked deviations for moderate ionic strengths, the limiting slopes (shown as dotted lines) as $I \rightarrow 0$ are in good agreement with the theory, so the law can be used for extrapolating data to very low molalities. The numbers in parentheses are the charge numbers of the ions.\\
\includegraphics[max width=\textwidth, center]{2025_02_27_d4b95d9718f0b9e2e6b9g-221}

Figure 5F. 6 The Davies equation gives agreement with experiment over a wider range of molalities than the limiting law (shown as a dotted line), but it fails at higher molalities. The data are for a 1,1-electrolyte.\\
from the extended Debye-Hückel law (sometimes called the Truesdell-Jones equation):

$$
\log \gamma_{ \pm}=-\frac{A\left|z_{+} z_{-}\right| I^{1 / 2}}{1+B I^{1 / 2}}
$$

Extended DebyeHückel law\\
(5F.30a)\\
where $B$ is a dimensionless constant. A more flexible extension is the Davies equation proposed by C.W. Davies in 1938:

$$
\log \gamma_{ \pm}=-\frac{A\left|z_{+} z_{-}\right| I^{1 / 2}}{1+B I^{1 / 2}}+C I \quad \quad \text { Davies equation }
$$

(5F.30b)\\
where $C$ is another dimensionless constant. Although $B$ can be interpreted as a measure of the closest approach of the ions, it (like $C$ ) is best regarded as an adjustable empirical parameter. A graph drawn on the basis of the Davies equation is shown in Fig. 5F.6. It is clear that eqn 5F.30b accounts for some activity coefficients over a moderate range of dilute solutions (up to about $0.1 \mathrm{~mol} \mathrm{~kg}^{-1}$ ); nevertheless it remains very poor near $1 \mathrm{molkg}^{-1}$.

Current theories of activity coefficients for ionic solutes take an indirect route. They set up a theory for the dependence of the activity coefficient of the solvent on the concentration of the solute, and then use the Gibbs-Duhem equation (eqn 5A.12a, $n_{\mathrm{A}} \mathrm{d} \mu_{\mathrm{A}}+n_{\mathrm{B}} \mathrm{d} \mu_{\mathrm{B}}=0$ ) to estimate the activity coefficient of the solute. The results are reasonably reliable for solutions with molalities greater than about $0.1 \mathrm{~mol} \mathrm{~kg}^{-1}$ and are valuable for the discussion of mixed salt solutions, such as sea water.

Table 5F. 3 Activities and standard states: a summary*

\begin{center}
\begin{tabular}{lllll}
\hline
Component & Basis & Standard state & Activity & Limits \\
\hline
Solid or liquid &  & Pure, 1 bar & $a=1$ &  \\
Solvent & Raoult & Pure solvent, 1 bar & $a=p / p^{*}, a=\gamma x$ & $\gamma \rightarrow 1$ as $x \rightarrow 1$ (pure solvent) \\
Solute & Henry & (1) A hypothetical state of the pure solute & $a=p / K, a=\gamma x$ & $\gamma \rightarrow 1$ as $x \rightarrow 0$ \\
 &  & (2) A hypothetical state of the solute at molality $b^{\ominus}$ & $a=\gamma b / b^{\ominus}$ & $\gamma \rightarrow 1$ as $b \rightarrow 0$ \\
Gas & Fugacity $^{\dagger}$ & Pure, a hypothetical state of 1 bar and behaving as a perfect gas & $f=\gamma p$ & $\gamma \rightarrow 1$ as $p \rightarrow 0$ \\
\hline
\end{tabular}
\end{center}

\begin{itemize}
  \item In each case, $\mu=\mu^{\ominus}+R T \ln a$.\\
${ }^{\dagger}$ Fugacity is discussed in A deeper look 2 on the website for this text.
\end{itemize}

\subsection*{Checklist of concepts}
1. The activity is an effective concentration that preserves the form of the expression for the chemical potential. See Table 5F.3.2. The chemical potential of a solute in an ideal-dilute solution is defined on the basis of Henry's law.3. The activity of a solute takes into account departures from Henry's law behaviour.4. The Margules equations relate the activities of the components of a model regular solution to its composition. They lead to expressions for the vapour pressures of the components of a regular solution,$\square$ 5. Mean activity coefficients apportion deviations from ideality equally to the cations and anions in an ionic solution.6. An ionic atmosphere is the time average accumulation of counter ions that exists around an ion in solution.7. The Debye-Hückel theory ascribes deviations from ideality to the Coulombic interaction of an ion with the ionic atmosphere around it.\\
$\square$ 8. The Debye-Hückel limiting law is extended by including two further empirical constants.

\subsection*{Checklist of equations}
\begin{center}
\begin{tabular}{llll}
\hline
Property & Equation & Comment & Equation number \\
\hline
Chemical potential of solvent & $\mu_{\mathrm{A}}=\mu_{\mathrm{A}}^{*}+R T \ln a_{\mathrm{A}}$ & Definition & 5 F .1 \\
Activity of solvent & $a_{\mathrm{A}}=p_{\mathrm{A}} / p_{\mathrm{A}}^{*}$ & $a_{\mathrm{A}} \rightarrow x_{\mathrm{A}}$ as $x_{\mathrm{A}} \rightarrow 1$ & 5 F .2 \\
Activity coefficient of solvent & $a_{\mathrm{A}}=\gamma_{\mathrm{A}} x_{\mathrm{A}}$ & $\gamma_{\mathrm{A}} \rightarrow 1$ as $x_{\mathrm{A}} \rightarrow 1$ & 5 F .4 \\
Chemical potential of solute & $\mu_{\mathrm{B}}=\mu_{\mathrm{B}}^{\ominus}+R T \ln a_{\mathrm{B}}$ & Definition & 5 F .9 \\
Activity of solute & $a_{\mathrm{B}}=p_{\mathrm{B}} / K_{\mathrm{B}}$ & $a_{\mathrm{B}} \rightarrow x_{\mathrm{B}}$ as $x_{\mathrm{B}} \rightarrow 0$ & 5 F .10 \\
Activity coefficient of solute & $a_{\mathrm{B}}=\gamma_{\mathrm{B}} x_{\mathrm{B}}$ & $\gamma_{\mathrm{B}} \rightarrow 1$ as $x_{\mathrm{B}} \rightarrow 0$ & 5 F .11 \\
Margules equations & $\ln \gamma_{\mathrm{A}}=\xi x_{\mathrm{B}}^{2}, \ln \gamma_{\mathrm{B}}=\xi x_{\mathrm{A}}^{2}$ & Regular solution & 5 F .16 \\
Vapour pressure & $p_{\mathrm{A}}=p_{\mathrm{A}}^{*} x_{\mathrm{A}} \mathrm{e}^{\xi\left(1-x_{\mathrm{A}}\right)^{2}}$ & Regular solution & 5 F .18 \\
Mean activity coefficient & $\gamma_{ \pm}=\left(\gamma_{+}^{p} \gamma_{-}^{q}\right)^{1 / s,} s=p+q$ & Definition & 5 F .25 \\
Debye-Hückel limiting law & $\log \gamma_{ \pm}=-A\left|z_{+} z_{-}\right| I^{1 / 2}$ & Valid as $I \rightarrow 0$ & 5 F .27 \\
Ionic strength & $I=\frac{1}{2} \sum_{i} z_{i}^{2}\left(b_{i} / b^{\ominus}\right)$ & Definition & 5 F .28 \\
Davies equation & $\log \gamma_{ \pm}=-A\left|z_{+} z_{-}\right| I^{1 / 2} /\left(1+B I^{1 / 2}\right)+C I$ & $A, B, C$ empirical constants & 5 F .30 b \\
\hline
\end{tabular}
\end{center}

\chapter{FOCUS 6 Chemical equilibrium}
Chemical reactions tend to move towards a dynamic equilibrium in which both reactants and products are present but have no further tendency to undergo net change. In some cases, the concentration of products in the equilibrium mixture is so much greater than that of the unchanged reactants that for all practical purposes the reaction is 'complete'. However, in many important cases the equilibrium mixture has significant concentrations of both reactants and products.

\subsection*{6A The equilibrium constant}
This Topic develops the concept of chemical potential and shows how it is used to account for the equilibrium composition of chemical reactions. The equilibrium composition corresponds to a minimum in the Gibbs energy plotted against the extent of reaction. By locating this minimum it is possible to establish the relation between the equilibrium constant and the standard Gibbs energy of reaction.\\
6A. 1 The Gibbs energy minimum; 6A. 2 The description of equilibrium

\subsection*{6B The response of equilibria to the conditions}
The thermodynamic formulation of equilibrium establishes the quantitative effects of changes in the conditions. One very important aspect of equilibrium is the control that can be exercised by varying the conditions, such as the pressure or temperature. 6B. 1 The response to pressure; 6B. 2 The response to temperature

\subsection*{6C Electrochemical cells}
Because many reactions involve the transfer of electrons, they can be studied (and utilized) by allowing them to take place in\\
a cell equipped with electrodes, with the spontaneous reaction forcing electrons through an external circuit. The electric potential of the cell is related to the reaction Gibbs energy, so its measurement provides an electrical procedure for the determination of thermodynamic quantities.\\
6C. 1 Half-reactions and electrodes; 6C. 2 Varieties of cells; 6C. 3 The cell potential; 6C. 4 The determination of thermodynamic functions

\subsection*{6D Electrode potentials}
Electrochemistry is in part a major application of thermodynamic concepts to chemical equilibria as well as being of great technological importance. As elsewhere in thermodynamics, electrochemical data can be reported in a compact form and applied to problems of chemical significance, especially to the prediction of the spontaneous direction of reactions and the calculation of equilibrium constants.\\
6D. 1 Standard potentials; 6D. 2 Applications of standard potentials

\subsection*{Web resources What is an application of this material?}
The thermodynamic description of spontaneous reactions has numerous practical and theoretical applications. One is to the discussion of biochemical processes, where one reaction drives another (Impact 9). Ultimately that is why we have to eat, for the reaction that takes place when one substance is oxidized can drive non-spontaneous reactions, such as protein synthesis, forward. Another makes use of the great sensitivity of electrochemical processes to the concentration of electroactive materials, and leads to the design of electrodes used in chemical analysis (Impact 10).

\section*{TOPIC 6A The equilibrium constant}
\subsection*{Why do you need to know this material?}
Equilibrium constants lie at the heart of chemistry and are a key point of contact between thermodynamics and laboratory chemistry. To understand the behaviour of reactions you need to see how equilibrium constants arise and understand how thermodynamic properties account for their values.

\subsection*{What is the key idea?}
At constant temperature and pressure, the composition of a reaction mixture tends to change until the Gibbs energy is a minimum.

\subsection*{What do you need to know already?}
Underlying the whole discussion is the expression of the direction of spontaneous change in terms of the Gibbs energy of a system (Topic 3D). This material draws on the concept of chemical potential and its dependence on the concentration or pressure of the substance (Topic 5A). You need to know how to express the total Gibbs energy of a mixture in terms of the chemical potentials of its components (Topic 5A).

As explained in Topic 3D, the direction of spontaneous change at constant temperature and pressure is towards lower values of the Gibbs energy, $G$. The idea is entirely general, and in this Topic it is applied to the discussion of chemical reactions. At constant temperature and pressure, a mixture of reactants has a tendency to undergo reaction until the Gibbs energy of the mixture has reached a minimum: that condition corresponds to a state of chemical equilibrium. The equilibrium is dynamic in the sense that the forward and reverse reactions continue, but at matching rates. As always in the application of thermodynamics, spontaneity is a tendency: there might be kinetic reasons why that tendency is not realized.

\subsection*{6A.1 The Gibbs energy minimum}
The equilibrium composition of a reaction mixture is located by calculating the Gibbs energy of the reaction mixture and then identifying the composition that corresponds to minimum $G$.

\subsection*{(a) The reaction Gibbs energy}
Consider the equilibrium $\mathrm{A} \rightleftharpoons \mathrm{B}$. Even though this reaction looks trivial, there are many examples of it, such as the isomerization of pentane to 2-methylbutane and the conversion of L -alanine to D -alanine.

If an infinitesimal amount $d \xi$ of $A$ turns into $B$, the change in the amount of A present is $\mathrm{d} n_{\mathrm{A}}=-\mathrm{d} \xi$ and the change in the amount of B present is $\mathrm{d} n_{\mathrm{B}}=+\mathrm{d} \xi$. The quantity $\xi$ (xi) is called the extent of reaction; it has the dimensions of amount of substance and is reported in moles. When the extent of reaction changes by a measurable amount $\Delta \xi$, the amount of A present changes from $n_{\mathrm{A}, 0}$ to $n_{\mathrm{A}, 0}-\Delta \xi$ and the amount of B changes from $n_{\mathrm{B}, 0}$ to $n_{\mathrm{B}, 0}+\Delta \xi$. In general, the amount of a component J changes by $v_{\mathrm{J}} \Delta \xi$, where $v_{\mathrm{J}}$ is the stoichiometric number of the species J (positive for products, negative for reactants). For example, if initially 2.0 mol A is present and after a period of time $\Delta \xi=+1.5 \mathrm{~mol}$, then the amount of A remaining is 0.5 mol . The amount of B formed is 1.5 mol .

The reaction Gibbs energy, $\Delta_{\mathrm{r}} G$, is defined as the slope of the graph of the Gibbs energy plotted against the extent of reaction:

\[
\Delta_{\mathrm{r}} G=\left(\frac{\partial G}{\partial \xi}\right)_{p, T} \quad \begin{align*}
& \text { Reaction Gibbs energy }  \tag{6A.1}\\
& {[\text { definition] }}
\end{align*}
\]

Although $\Delta$ normally signifies a difference in values, here it signifies a derivative, the slope of $G$ with respect to $\xi$. However, to see that there is a close relationship with the normal usage, suppose the reaction advances by $\mathrm{d} \xi$. The corresponding change in Gibbs energy is

$$
\mathrm{d} G=\mu_{\mathrm{A}} \mathrm{~d} n_{\mathrm{A}}+\mu_{\mathrm{B}} \mathrm{~d} n_{\mathrm{B}}=-\mu_{\mathrm{A}} \mathrm{~d} \xi+\mu_{\mathrm{B}} \mathrm{~d} \xi=\left(\mu_{\mathrm{B}}-\mu_{\mathrm{A}}\right) \mathrm{d} \xi
$$

This equation can be reorganized into

$$
\left(\frac{\partial G}{\partial \xi}\right)_{p, T}=\mu_{\mathrm{B}}-\mu_{\mathrm{A}}
$$

That is,


\begin{equation*}
\Delta_{\mathrm{r}} G=\mu_{\mathrm{B}}-\mu_{\mathrm{A}} \tag{6A.2}
\end{equation*}


and $\Delta_{r} G$ can also be interpreted as the difference between the chemical potentials (the partial molar Gibbs energies) of the reactants and products at the current composition of the reaction mixture.

Because chemical potentials vary with composition, the slope of the plot of Gibbs energy against extent of reaction, and\\
\includegraphics[max width=\textwidth, center]{2025_02_27_d4b95d9718f0b9e2e6b9g-237(1)}

Fig. 6A. 1 As the reaction advances, represented by the extent of reaction $\xi$ increasing, the slope of a plot of total Gibbs energy of the reaction mixture against $\xi$ changes. Equilibrium corresponds to the minimum in the Gibbs energy, which is where the slope is zero.\\
therefore the reaction Gibbs energy, changes as the reaction proceeds. The spontaneous direction of reaction lies in the direction of decreasing $G$ (i.e. down the slope of $G$ plotted against $\xi$ ). Thus, the reaction $\mathrm{A} \rightarrow \mathrm{B}$ is spontaneous when $\mu_{\mathrm{A}}>$ $\mu_{\mathrm{B}}$, whereas the reverse reaction is spontaneous when $\mu_{\mathrm{B}}>\mu_{\mathrm{A}}$. The slope is zero, and the reaction is at equilibrium and spontaneous in neither direction, when

$$
\Delta_{\mathrm{r}} G=0 \quad \text { Condition of equilibrium } \quad(6 \mathrm{~A} .3)
$$

This condition occurs when $\mu_{\mathrm{B}}=\mu_{\mathrm{A}}$ (Fig. 6A.1). It follows that if the composition of the reaction mixture that ensures $\mu_{\mathrm{B}}=\mu_{\mathrm{A}}$ can be found, then that will be the composition of the reaction mixture at equilibrium. Note that the chemical potential is now fulfilling the role its name suggests: it represents the potential for chemical change, and equilibrium is attained when these potentials are in balance.

\subsection*{(b) Exergonic and endergonic reactions}
The spontaneity of a reaction at constant temperature and pressure can be expressed in terms of the reaction Gibbs energy:

\begin{displayquote}
If $\Delta_{\mathrm{r}} G<0$, the forward reaction is spontaneous.\\
If $\Delta_{\mathrm{r}} G>0$, the reverse reaction is spontaneous.\\
If $\Delta_{\mathrm{r}} G=0$, the reaction is at equilibrium.
\end{displayquote}

A reaction for which $\Delta_{r} G<0$ is called exergonic (from the Greek words for 'work producing'). The name signifies that, because the process is spontaneous, it can be used to drive another process, such as another reaction, or used to do non-expansion work. A simple mechanical analogy is a pair of weights joined by a string (Fig. 6A.2): the lighter of the pair of weights will be pulled up as the heavier weight falls\\
\includegraphics[max width=\textwidth, center]{2025_02_27_d4b95d9718f0b9e2e6b9g-237}

Fig. 6A. 2 If two weights are coupled as shown here, then the heavier weight will move the lighter weight in its nonspontaneous direction: overall, the process is still spontaneous. The weights are the analogues of two chemical reactions: a reaction with a large negative $\Delta G$ can force another reaction with a smaller $\Delta G$ to run in its non-spontaneous direction.\\
down. Although the lighter weight has a natural tendency to move down, its coupling to the heavier weight results in it being raised. In biological cells, the oxidation of carbohydrates acts as the heavy weight that drives other reactions forward and results in the formation of proteins from amino acids, muscle contraction, and brain activity. A reaction for which $\Delta_{\mathrm{r}} G>0$ is called endergonic (signifying 'work consuming'); such a reaction can be made to occur only by doing work on it.

\subsection*{Brief illustration 6A. 1}
The reaction Gibbs energy of a certain reaction is $-200 \mathrm{~kJ} \mathrm{~mol}^{-1}$, so the reaction is exergonic, and in a suitable device (a fuel cell, for instance) operating at constant temperature and pressure, it could produce 200 kJ of electrical work for each mole of reaction events. The reverse reaction, for which $\Delta_{\mathrm{r}} G=$ $+200 \mathrm{~kJ} \mathrm{~mol}^{-1}$ is endergonic and at least 200 kJ of work must be done to achieve it, perhaps through electrolysis.

\subsection*{6A. 2 The description of equilibrium}
With the background established, it is now possible to apply thermodynamics to the description of chemical equilibrium.

\subsection*{(a) Perfect gas equilibria}
When A and B are perfect gases, eqn 5A.15a $\left(\mu=\mu^{\ominus}+R T \ln \left(p / p^{\ominus}\right)\right)$ can be used to write


\begin{align*}
\Delta_{\mathrm{r}} G & =\mu_{\mathrm{B}}-\mu_{\mathrm{A}}=\left(\mu_{\mathrm{B}}^{\ominus}+R T \ln \frac{p_{\mathrm{B}}}{p^{\ominus}}\right)-\left(\mu_{\mathrm{A}}^{\ominus}+R T \ln \frac{p_{\mathrm{A}}}{p^{\ominus}}\right) \\
& =\Delta_{\mathrm{r}} G^{\ominus}+R T \ln \frac{p_{\mathrm{B}}}{p_{\mathrm{A}}} \tag{6A.4}
\end{align*}


If the ratio of partial pressures is denoted by $Q$, then it follows that


\begin{equation*}
\Delta_{\mathrm{r}} G=\Delta_{\mathrm{r}} G^{\ominus}+R T \ln Q \quad Q=\frac{p_{\mathrm{B}}}{p_{\mathrm{A}}} \tag{6A.5}
\end{equation*}


The ratio $Q$ is an example of a 'reaction quotient', a quantity to be defined more formally shortly. It ranges from 0 when $p_{\mathrm{B}}=0$ (corresponding to pure A) to infinity when $p_{\mathrm{A}}=0$ (corresponding to pure B ). The standard reaction Gibbs energy, $\Delta_{\mathrm{r}} G^{\ominus}$ (Topic 3D), is the difference in the standard molar Gibbs energies of the reactants and products, so


\begin{equation*}
\Delta_{\mathrm{r}} G^{\ominus}=G_{\mathrm{m}}^{\ominus}(\mathrm{B})-G_{\mathrm{m}}^{\ominus}(\mathrm{A})=\mu_{\mathrm{B}}^{\ominus}-\mu_{\mathrm{A}}^{\ominus} \tag{6A.6}
\end{equation*}


Note that in the definition of $\Delta_{\mathrm{r}} G^{\ominus}$, the $\Delta_{\mathrm{r}}$ has its normal meaning as the difference 'products - reactants'. As seen in Topic 3D, the difference in standard molar Gibbs energies of the products and reactants is equal to the difference in their standard Gibbs energies of formation, so in practice $\Delta_{r} G^{\ominus}$ is calculated from


\begin{equation*}
\Delta_{\mathrm{r}} G^{\ominus}=\Delta_{\mathrm{f}} G^{\ominus}(\mathrm{B})-\Delta_{\mathrm{f}} G^{\ominus}(\mathrm{A}) \tag{6A.7}
\end{equation*}


At equilibrium $\Delta_{r} G=0$. The ratio of partial pressures, the reaction quotient $Q$, at equilibrium has a certain value $K$, and eqn 6 A .5 becomes

$$
0=\Delta_{\mathrm{r}} G^{\ominus}+R T \ln K
$$

which rearranges to


\begin{equation*}
R T \ln K=-\Delta_{\mathrm{r}} G^{\ominus} \quad K=\left(\frac{p_{\mathrm{B}}}{p_{\mathrm{A}}}\right)_{\text {equilibrium }} \tag{6A.8}
\end{equation*}


This relation is a special case of one of the most important equations in chemical thermodynamics: it is the link between tables of thermodynamic data, such as those in the Resource section and the chemically important 'equilibrium constant', $K$ (again, a quantity that will be defined formally shortly).

\subsection*{Brief illustration 6A. 2}
The standard Gibbs energy for the isomerization of pentane to 2-methylbutane at 298 K , the reaction $\mathrm{CH}_{3}\left(\mathrm{CH}_{2}\right)_{3} \mathrm{CH}_{3}(\mathrm{~g}) \rightarrow$ $\left(\mathrm{CH}_{3}\right)_{2} \mathrm{CHCH}_{2} \mathrm{CH}_{3}(\mathrm{~g})$, is close to $-6.7 \mathrm{~kJ} \mathrm{~mol}^{-1}$ (this is an estimate based on enthalpies of formation; its actual value is not listed). Therefore, the equilibrium constant for the reaction is

$$
K=\mathrm{e}^{-\left(-6.7 \times 10^{3} \mathrm{Jmol}^{-1}\right) /\left(8.3145 \mathrm{JK}^{-1} \mathrm{~mol}^{-1}\right) \times(298 \mathrm{~K})}=\mathrm{e}^{2.7 \ldots}=15
$$

In molecular terms, the minimum in the Gibbs energy, which corresponds to $\Delta_{\mathrm{r}} G=0$, stems from the Gibbs energy of mixing of the two gases. To see the role of mixing, consider the reaction $\mathrm{A} \rightarrow \mathrm{B}$. If only the enthalpy were important, then $H$, and therefore $G$, would change linearly from its value for pure\\
\includegraphics[max width=\textwidth, center]{2025_02_27_d4b95d9718f0b9e2e6b9g-238}

Fig. 6A. 3 If the mixing of reactants and products is ignored, the Gibbs energy changes linearly from its initial value (pure reactants) to its final value (pure products) and the slope of the line is $\Delta_{t} G^{\ominus}$. However, as products are produced, there is a further contribution to the Gibbs energy arising from their mixing (lowest curve). The sum of the two contributions has a minimum, which corresponds to the equilibrium composition of the system.\\
reactants to its value for pure products. The slope of this straight line is a constant and equal to $\Delta_{\mathrm{r}} G^{\ominus}$ at all stages of the reaction and there is no intermediate minimum in the graph (Fig. 6A.3). However, when the entropy is taken into account, there is an additional contribution to the Gibbs energy that is given by eqn 5 A. $17\left(\Delta_{\text {mix }} G=n R T\left(x_{\mathrm{A}} \ln x_{\mathrm{A}}+x_{\mathrm{B}} \ln x_{\mathrm{B}}\right)\right)$. This expression makes a U-shaped contribution to the total Gibbs energy. As can be seen from Fig. 6A.3, when it is included there is an intermediate minimum in the total Gibbs energy, and its position corresponds to the equilibrium composition of the reaction mixture.

It follows from eqn 6A. 8 that, when $\Delta_{\mathrm{r}} G^{\ominus}>0, K<1$. Therefore, at equilibrium the partial pressure of A exceeds that of B , which means that the reactant A is favoured in the equilibrium. When $\Delta_{\mathrm{r}} G^{\ominus}<0, K>1$, so at equilibrium the partial pressure of $B$ exceeds that of $A$. Now the product $B$ is favoured in the equilibrium.

A note on good practice A common remark is that 'a reaction is spontaneous if $\Delta_{\mathrm{r}} G^{\ominus}<0$ '. However, whether or not a reaction is spontaneous at a particular composition depends on the value of $\Delta_{\mathrm{r}} G$ at that composition, not $\Delta_{\mathrm{r}} G^{\ominus}$. The forward reaction is spontaneous ( $\Delta_{\mathrm{r}} G<0$ ) when $Q<K$ and the reverse reaction is spontaneous when $Q>K$. It is far better to interpret the sign of $\Delta_{\mathrm{r}} G^{\ominus}$ as indicating whether $K$ is greater or smaller than 1 .

\subsection*{(b) The general case of a reaction}
To extend the argument that led to eqn 6A. 8 to a general reaction, first note that a chemical reaction may be expressed symbolically in terms of stoichiometric numbers as

$$
0=\sum_{\mathrm{J}} v_{\mathrm{J}} \mathrm{~J}
$$

where $J$ denotes the substances and the $v_{\mathrm{J}}$ are the corresponding stoichiometric numbers in the chemical equation, which are positive for products and negative for reactants. In the reaction $2 \mathrm{~A}+\mathrm{B} \rightarrow 3 \mathrm{C}+\mathrm{D}$, for instance, $v_{\mathrm{A}}=-2, v_{\mathrm{B}}=-1, v_{\mathrm{C}}=+3$, and $v_{\mathrm{D}}=+1$.

With these points in mind, it is possible to write an expression for the reaction Gibbs energy, $\Delta_{\mathrm{r}} G$, at any stage during the reaction.

How is that done? 6A. 1 Deriving an expression for the dependence of the reaction Gibbs energy on the reaction quotient

Consider a reaction with stoichiometric numbers $v_{\mathrm{J}}$. When the reaction advances by $\mathrm{d} \xi$, the amounts of reactants and products change by $\mathrm{d} n_{\mathrm{J}}=v_{\mathrm{J}} \mathrm{d} \xi$. The resulting infinitesimal change in the Gibbs energy at constant temperature and pressure is

$$
\mathrm{d} G=\sum_{\mathrm{J}} \mu_{\mathrm{J}} \mathrm{~d} n_{\mathrm{J}}=\sum_{\mathrm{J}} \mu_{\mathrm{J}} v_{\mathrm{J}} \mathrm{~d} \xi=\left(\sum_{\mathrm{J}} \mu_{\mathrm{J}} v_{\mathrm{J}}\right) \mathrm{d} \xi
$$

It follows that

$$
\Delta_{\mathrm{r}} G=\left(\frac{\partial G}{\partial \xi}\right)_{p, T}=\sum_{\mathrm{J}} v_{\mathrm{J}} \mu_{\mathrm{J}}
$$

Step 1 Write the chemical potential in terms of the activity\\
To make progress, note that the chemical potential of a species J is related to its activity by eqn $5 \mathrm{~F} .9\left(\mu_{\mathrm{J}}=\mu_{\mathrm{J}}^{\ominus}+R T \ln a_{\mathrm{J}}\right)$. When this relation is substituted into the expression for $\Delta_{\mathrm{r}} G$ the result is

$$
\begin{aligned}
\Delta_{\mathrm{r}} G & =\overbrace{\sum_{\mathrm{J}} v_{\mathrm{J}} \mu_{\mathrm{J}}^{\ominus}}^{\Delta_{\mathrm{J}} G^{\ominus}}+R T \sum_{\mathrm{J}} v_{\mathrm{J}} \ln a_{\mathrm{J}} \\
& =\Delta_{\mathrm{r}} G^{\ominus}+R T \sum_{\mathrm{J}} v_{\mathrm{J}} \ln a_{\mathrm{J}}=\Delta_{\mathrm{r}} G^{\ominus}+R T \sum_{\mathrm{J}} \ln a_{\mathrm{J}}^{v_{\mathrm{J}}}
\end{aligned}
$$

Because $\ln x+\ln y+\cdots=\ln x y \ldots$, it follows that

$$
\sum_{i} \ln x_{i}=\ln \left(\prod_{i} x_{i}\right)
$$

The symbol $\Pi$ denotes the product of what follows it (just as $\Sigma$ denotes the sum). The expression for the Gibbs energy change then simplifies to

$$
\Delta_{\mathrm{r}} G=\Delta_{\mathrm{r}} G^{\ominus}+R T \ln \prod_{J} a_{J}^{v_{J}}
$$

Step 2 Introduce the reaction quotient\\
Now define the reaction quotient as

$$
Q=\prod_{J} a_{J}^{v_{j}}
$$

Reaction quotient [definition]

Because reactants have negative stoichiometric numbers, they automatically appear as the denominator when the product is written out explicitly and this expression has the form

\[
Q=\frac{\text { activities of products }}{\text { activities of reactants }} \quad \begin{align*}
& \text { Reaction quotient }  \tag{6A.11}\\
& \text { [general form] }
\end{align*}
\]

with the activity of each species raised to the power given by its stoichiometric coefficient.

It follows that the expression for the reaction Gibbs energy simplifies to\\
$-\Delta_{\mathrm{r}} G=\Delta_{\mathrm{r}} G^{\ominus}+R T \ln Q$\\
Reaction Gibbs energy at an arbitrary stage

\subsection*{Brief illustration 6A. 3}
Consider the reaction $2 \mathrm{~A}+3 \mathrm{~B} \rightarrow \mathrm{C}+2 \mathrm{D}$, in which case $v_{\mathrm{A}}=$ $-2, v_{\mathrm{B}}=-3, v_{\mathrm{C}}=+1$, and $v_{\mathrm{D}}=+2$. The reaction quotient is then

$$
Q=a_{\mathrm{A}}^{-2} a_{\mathrm{B}}^{-3} a_{\mathrm{C}} a_{\mathrm{D}}^{2}=\frac{a_{\mathrm{C}} a_{\mathrm{D}}^{2}}{a_{\mathrm{A}}^{2} a_{\mathrm{B}}^{3}}
$$

As in Topic 3D, the standard reaction Gibbs energy is calculated from


\begin{equation*}
\Delta_{\mathrm{r}} G^{\ominus}=\sum_{\text {Products }} v \Delta_{\mathrm{f}} G^{\ominus}-\sum_{\text {Reactants }} v \Delta_{\mathrm{f}} G^{\ominus} \tag{6A.13a}
\end{equation*}


Reaction Gibbs energy [practical implementation]\\
where the $v$ are the (positive) stoichiometric coefficients. More formally,

\[
\Delta_{\mathrm{r}} G^{\ominus}=\sum_{\mathrm{J}} v_{\mathrm{J}} \Delta_{\mathrm{f}} G^{\ominus}(\mathrm{J}) \quad \begin{align*}
& \text { Reaction Gibbs energy }  \tag{6A.13b}\\
& \text { [formal expression] }
\end{align*}
\]

where the $v_{\mathrm{J}}$ are the (signed) stoichiometric numbers.\\
At equilibrium, the slope of $G$ is zero: $\Delta_{\mathrm{r}} G=0$. The activities then have their equilibrium values and

\[
K=\left(\prod_{J} a_{J}^{v_{J}}\right)_{\text {equilibrium }} \quad \begin{align*}
& \text { Equilibrium constant }  \tag{6A.14}\\
& \text { [definition] }
\end{align*}
\]

This expression has the same form as $Q$ but is evaluated using equilibrium activities. From now on, the 'equilibrium' subscript will not be written explicitly, but it will be clear from the context that $Q$ is defined in terms of the activities at an arbitrary stage of the reaction and $K$ is the value of $Q$ at equilibrium. An equilibrium constant $K$ expressed in terms of activities is called a thermodynamic equilibrium constant. Note that, because activities are dimensionless, the thermodynamic equilibrium constant is also dimensionless. In elementary applications, the activities that occur in eqn 6A. 14 are often replaced as follows:

\begin{center}
\begin{tabular}{llll}
\hline
State & Measure & \begin{tabular}{l}
Approximation \\
for $a_{\mathrm{J}}$ \\
\end{tabular} & Definition \\
\hline
Solute & molality & $b_{\mathrm{J}} / b_{\mathrm{J}}^{\ominus}$ & $b^{\ominus}=1 \mathrm{molkg}^{-1}$ \\
 & molar concentration & $[\mathrm{J}] / c^{\ominus}$ & $c^{\ominus}=1 \mathrm{~mol} \mathrm{dm}^{-3}$ \\
Gas phase & partial pressure & $p_{\mathrm{J}} / p^{\ominus}$ & $p^{\ominus}=1 \mathrm{bar}$ \\
Pure solid, liquid &  & 1 (exact) &  \\
\hline
\end{tabular}
\end{center}

Note that the activity is 1 for pure solids and liquids, so such substances make no contribution to $Q$ even though they might appear in the chemical equation. When the approximations are made, the resulting expressions for $Q$ and $K$ are only approximations. The approximation is particularly severe for electrolyte solutions, for in them activity coefficients differ from 1 even in very dilute solutions (Topic 5F).

\subsection*{Brief illustration 6A. 4}
The equilibrium constant for the heterogeneous equilibrium $\mathrm{CaCO}_{3}(\mathrm{~s}) \rightleftharpoons \mathrm{CaO}(\mathrm{s})+\mathrm{CO}_{2}(\mathrm{~g})$ is

$$
K=a_{\mathrm{CaCO}_{3}(\mathrm{~s})}^{-1} a_{\mathrm{CaO}_{(\mathrm{s}}} a_{\mathrm{CO}_{2}(\mathrm{~g})}=\frac{\overbrace{\mathrm{CaO}_{\mathrm{Ca}(\mathrm{~s})}}^{1} a_{\mathrm{CO}_{2}(\mathrm{~g})}}{\underbrace{a_{\mathrm{CaCO}_{3}(\mathrm{~s})}}_{1}}=a_{\mathrm{CO}_{2}(\mathrm{~g})}
$$

Provided the carbon dioxide can be treated as a perfect gas, go on to write

$$
K=p_{\mathrm{CO}_{2}} / p^{\ominus}
$$

and conclude that in this case the equilibrium constant is the numerical value of the equilibrium pressure of $\mathrm{CO}_{2}$ above the solid sample.

At equilibrium $\Delta_{\mathrm{r}} G=0$ in eqn 6A. 12 and $Q$ is replaced by $K$. The result is


\begin{equation*}
\Delta_{\mathrm{r}} G^{\ominus}=-R T \ln K \quad \text { Thermodynamic equilibrium constant } \tag{6A.15}
\end{equation*}


This is an exact and highly important thermodynamic relation, for it allows the calculation of the equilibrium constant of any reaction from tables of thermodynamic data, and hence the prediction of the equilibrium composition of the reaction mixture.

\subsection*{Example 6A. 1 Calculating an equilibrium constant}
Calculate the equilibrium constant for the ammonia synthesis reaction, $\mathrm{N}_{2}(\mathrm{~g})+3 \mathrm{H}_{2}(\mathrm{~g}) \rightleftharpoons 2 \mathrm{NH}_{3}(\mathrm{~g})$, at 298 K , and show how $K$ is related to the partial pressures of the species at equilibrium when the overall pressure is low enough for the gases to be treated as perfect.

Collect your thoughts Calculate the standard reaction Gibbs energy from eqn 6A. 13 and use its value in eqn 6A. 15 to evaluate the equilibrium constant. The expression for the equilibrium constant is obtained from eqn 6A.14, and because the gases are taken to be perfect, replace each activity by the ratio $p_{\mathrm{I}} / p^{\ominus}$, where $p_{\mathrm{J}}$ is the partial pressure of species J.

The solution The standard Gibbs energy of the reaction is

$$
\begin{aligned}
\Delta_{\mathrm{r}} G^{\ominus} & =2 \Delta_{\mathrm{f}} G^{\ominus}\left(\mathrm{NH}_{3}, \mathrm{~g}\right)-\left\{\Delta_{\mathrm{f}} G^{\ominus}\left(\mathrm{N}_{2}, \mathrm{~g}\right)+3 \Delta_{\mathrm{f}} G^{\ominus}\left(\mathrm{H}_{2}, \mathrm{~g}\right)\right\} \\
& =2 \Delta_{\mathrm{f}} G^{\ominus}\left(\mathrm{NH}_{3}, \mathrm{~g}\right)=2 \times\left(-16.45 \mathrm{~kJ} \mathrm{~mol}^{-1}\right)
\end{aligned}
$$

Then,\\
$\ln K=-\frac{2 \times\left(-1.645 \times 10^{4} \mathrm{Jmol}^{-1}\right)}{\left(8.3145 \mathrm{JK}^{-1} \mathrm{~mol}^{-1}\right) \times(298 \mathrm{~K})}=\frac{2 \times 1.645 \times 10^{4}}{8.3145 \times 298}=13.2 \ldots$\\
Hence, $K=5.8 \times 10^{5}$. This result is thermodynamically exact. The thermodynamic equilibrium constant for the reaction is

$$
K=\frac{a_{\mathrm{NH}_{3}}^{2}}{a_{\mathrm{N}_{2}} a_{\mathrm{H}_{2}}^{3}}
$$

and has the value just calculated. At low overall pressures, the activities can be replaced by the ratios $p_{\mathrm{J}} / p^{\ominus}$ and an approximate form of the equilibrium constant is

$$
K=\frac{\left(p_{\mathrm{NH}_{3}} / p^{\ominus}\right)^{2}}{\left(p_{\mathrm{N}_{2}} / p^{\ominus}\right)\left(p_{\mathrm{H}_{2}} / p^{\ominus}\right)^{3}}=\frac{p_{\mathrm{NH}_{3}}^{2} p^{\ominus 2}}{p_{\mathrm{N}_{2}} p_{\mathrm{H}_{2}}^{3}}
$$

Self-test 6 A. 1 Evaluate the equilibrium constant for $\mathrm{N}_{2} \mathrm{O}_{4}(\mathrm{~g})$ $\rightleftharpoons 2 \mathrm{NO}_{2}(\mathrm{~g})$ at 298 K .\\
$\varsigma I^{\circ} 0=y: \iota \partial м s u_{V}$

\subsection*{Example 6A. 2 Estimating the degree of dissociation}
 at equilibriumThe degree of dissociation (or extent of dissociation, $\alpha$ ) is defined as the fraction of reactant that has decomposed; if the initial amount of reactant is $n$ and the amount at equilibrium is $n_{\mathrm{eq}}$, then $\alpha=\left(n-n_{\text {eq }}\right) / n$. The standard reaction Gibbs energy for the decomposition $\mathrm{H}_{2} \mathrm{O}(\mathrm{g}) \rightarrow \mathrm{H}_{2}(\mathrm{~g})+\frac{1}{2} \mathrm{O}_{2}(\mathrm{~g})$ is $+118.08 \mathrm{~kJ} \mathrm{~mol}^{-1}$ at 2300 K . What is the degree of dissociation of $\mathrm{H}_{2} \mathrm{O}$ at 2300 K when the reaction is allowed to come to equilibrium at a total pressure of 1.00 bar ?

Collect your thoughts The equilibrium constant is obtained from the standard Gibbs energy of reaction by using eqn 6A.15, so your task is to relate the degree of dissociation, $\alpha$, to $K$ and then to find its numerical value. Proceed by expressing the equilibrium compositions in terms of $\alpha$, and solve for $\alpha$ in terms of $K$. Because the standard reaction Gibbs energy is large and positive, you can anticipate that $K$ will be small, and hence that $\alpha \ll 1$, which opens the way to making approximations to obtain its numerical value.

The solution The equilibrium constant is obtained from eqn 6A. 15 in the form

$$
\begin{aligned}
\ln K & =-\frac{\Delta_{\mathrm{r}} G^{\ominus}}{R T}=-\frac{1.1808 \times 10^{5} \mathrm{Jmol}^{-1}}{\left(8.3145 \mathrm{JK}^{-1} \mathrm{~mol}^{-1}\right) \times(2300 \mathrm{~K})} \\
& =-\frac{1.1808 \times 10^{5}}{8.3145 \times 2300}=-6.17 \ldots
\end{aligned}
$$

It follows that $K=2.08 \times 10^{-3}$. The equilibrium composition is expressed in terms of $\alpha$ by drawing up the following table:

\begin{center}
\begin{tabular}{lllll}
\hline
 & $\mathrm{H}_{2} \mathrm{O} \rightarrow$ & $\mathrm{H}_{2}$ & + & $\frac{1}{2} \mathrm{O}_{2}$ \\
\hline
\begin{tabular}{l}
Initial amount \\
\end{tabular} & $n$ & 0 & 0 &  \\
\begin{tabular}{l}
Change to reach \\
equilibrium \\
\end{tabular} & $-\alpha n$ & $+\alpha n$ & $+\frac{1}{2} \alpha n$ &  \\
\begin{tabular}{l}
Amount at \\
equilibrium \\
\end{tabular} & $(1-\alpha) n$ & $\alpha n$ & $\frac{1}{2} \alpha n$ & Total: $\left(1+\frac{1}{2} \alpha\right) n$ \\
\begin{tabular}{llll}
Mole fraction, $x_{\mathrm{J}}$ & $\frac{1-\alpha}{1+\frac{1}{2} \alpha}$ & $\frac{\alpha}{1+\frac{1}{2} \alpha}$ & $\frac{\frac{1}{2} \alpha}{1+\frac{1}{2} \alpha}$ \\
 &  &  &  \\
Partial pressure, $p_{\mathrm{J}}$ & $\frac{(1-\alpha) p}{1+\frac{1}{2} \alpha}$ & $\frac{\alpha p}{1+\frac{1}{2} \alpha}$ & $\frac{\frac{1}{2} \alpha p}{1+\frac{1}{2} \alpha}$ \\
\hline
\end{tabular} &  &  &  &  \\
\hline
\end{tabular}
\end{center}

where, for the entries in the last row, $p_{\mathrm{J}}=x_{\mathrm{p}} p$ (eqn 1 A .6 ) has been used. The equilibrium constant is therefore

$$
K=\frac{p_{\mathrm{H}_{2}} p_{\mathrm{O}_{2}}^{1 / 2}}{p_{\mathrm{H}_{2} \mathrm{O}}}=\frac{\alpha^{3 / 2} p^{1 / 2}}{(1-\alpha)(2+\alpha)^{1 / 2}}
$$

In this expression, $p$ has been used in place of $p / p^{\ominus}$, to simplify its appearance. Now make the approximation that $\alpha \ll 1$, so $1-\alpha \approx 1$ and $2+\alpha \approx 2$, and hence obtain

$$
K \approx \frac{\alpha^{3 / 2} p^{1 / 2}}{2^{1 / 2}}
$$

Under the stated conditions, $p=1.00 \mathrm{bar}$ (that is, $p / p^{\ominus}=1.00$ ), so $\alpha \approx\left(2^{1 / 2} K\right)^{2 / 3}=0.0205$. That is, about 2 per cent of the water has decomposed.

A note on good practice Always check that the approximation is consistent with the final answer. In this case, $\alpha \ll 1$ in accord with the original assumption.

Self-test 6A.2 For the same reaction, the standard Gibbs energy of reaction at 2000 K is $+135.2 \mathrm{~kJ} \mathrm{~mol}^{-1}$. Suppose that steam at 200 kPa is passed through a furnace tube at that temperature. Calculate the mole fraction of $\mathrm{O}_{2}$ present in the output gas stream.\\
\includegraphics[max width=\textwidth, center]{2025_02_27_d4b95d9718f0b9e2e6b9g-241}

\subsection*{(c) The relation between equilibrium constants}
Equilibrium constants in terms of activities are exact, but it is often necessary to relate them to concentrations. Formally, it is necessary to know the activity coefficients $\gamma_{\mathrm{J}}$ (Topic 5F), and then to use $a_{\mathrm{J}}=\gamma_{\mathrm{J}} x_{\mathrm{J}}, a_{\mathrm{J}}=\gamma_{\mathrm{J}} b_{\mathrm{J}} / b^{\ominus}$, or $a_{\mathrm{J}}=\gamma_{\mathrm{J}}[\mathrm{J}] / c^{\ominus}$, where $x_{\mathrm{J}}$ is\\
a mole fraction, $b_{\mathrm{J}}$ is a molality, and [J] is a molar concentration. For example, if the composition is expressed in terms of molality for an equilibrium of the form $\mathrm{A}+\mathrm{B} \rightleftharpoons \mathrm{C}+\mathrm{D}$, where all four species are solutes, then


\begin{equation*}
K=\frac{a_{\mathrm{C}} a_{\mathrm{D}}}{a_{\mathrm{A}} a_{\mathrm{B}}}=\frac{\gamma_{\mathrm{C}} \gamma_{\mathrm{D}}}{\gamma_{\mathrm{A}} \gamma_{\mathrm{B}}} \times \frac{b_{\mathrm{C}} b_{\mathrm{D}}}{b_{\mathrm{A}} b_{\mathrm{B}}}=K_{\gamma} K_{b} \tag{6A.16}
\end{equation*}


The activity coefficients must be evaluated at the equilibrium composition of the mixture (for instance, by using one of the Debye-Hückel expressions, Topic 5F), which may involve a complicated calculation, because the activity coefficients are known only if the equilibrium composition is already known. In elementary applications, and to begin the iterative calculation of the concentrations in a real example, the assumption is often made that the activity coefficients are all so close to unity that $K_{\gamma}=1$. Given these difficulties, it is common in elementary chemistry to assume that $K \approx K_{b}$, which allows equilibria to be discussed in terms of the molalities (or molar concentrations) themselves.

A special case arises when the equilibrium constant of a gas-phase reaction is to be expressed in terms of molar concentrations instead of the partial pressures that appear in the thermodynamic equilibrium constant. Provided the gases are perfect, the $p_{\mathrm{J}}$ that appear in $K$ can be replaced by [J]RT, and

$$
\begin{aligned}
K & =\prod_{J} a_{J}^{v_{J}}=\prod_{J}\left(\frac{p_{J}}{p^{\ominus}}\right)^{v_{J}}=\prod_{J}[J]^{v_{J}}\left(\frac{R T}{p^{\ominus}}\right)^{v_{J}} \\
& =\prod_{J}[J]^{v_{j}} \times \prod_{J}\left(\frac{R T}{p^{\ominus}}\right)^{v_{J}}
\end{aligned}
$$

(Products can always be factorized in this way: $a b c d e f$ is the same as $a b c \times d e f$.) The (dimensionless) equilibrium constant $K_{c}$ is defined as


\begin{equation*}
K_{c}=\prod_{\mathrm{J}}\left(\frac{[\mathrm{~J}]}{c^{\ominus}}\right)^{v_{1}} \tag{6A.17}
\end{equation*}


$K_{c}$ for gas-phase reactions [definition]

It follows that


\begin{equation*}
K=K_{c} \times \prod_{J}\left(\frac{c^{\ominus} R T}{p^{\ominus}}\right)^{v_{1}} \tag{6A.18a}
\end{equation*}


With $\Delta v=\sum_{\mathrm{J}} v_{\mathrm{J}}$, which is easier to think of as $v$ (products) - $v$ (reactants), the relation between $K$ and $K_{c}$ for a gas-phase reaction is

\[
K=K_{c} \times\left(\frac{c^{\ominus} R T}{p^{\ominus}}\right)^{\Delta v} \quad \begin{align*}
& \text { Relation between } K \text { and }  \tag{6A.18b}\\
& K_{c} \text { for gas-phase reactions }
\end{align*}
\]

For numerical calculations, note that $p^{\ominus} / c^{\ominus} R$ evaluates to 12.03 K .

\subsection*{Brief illustration 6A. 5}
For the reaction $\mathrm{N}_{2}(\mathrm{~g})+3 \mathrm{H}_{2}(\mathrm{~g}) \rightarrow 2 \mathrm{NH}_{3}(\mathrm{~g}), \Delta v=2-3-1=-2$, so

$$
K=K_{c} \times\left(\frac{T}{12.03 \mathrm{~K}}\right)^{-2}=K_{c} \times\left(\frac{12.03 \mathrm{~K}}{T}\right)^{2}
$$

At 298.15 K the relation is

$$
K=K_{c} \times\left(\frac{12.03 \mathrm{~K}}{298.15 \mathrm{~K}}\right)^{2}=\frac{K_{c}}{614.2}
$$

so $K_{c}=614.2 K$. Note that both $K$ and $K_{c}$ are dimensionless.

\subsection*{(d) Molecular interpretation of the equilibrium constant}
Deeper insight into the origin and significance of the equilibrium constant can be obtained by considering the Boltzmann distribution of molecules over the available states of a system composed of reactants and products (see the Prologue to this text). When atoms can exchange partners, as in a reaction, the species present include atoms bonded together as molecules of both reactants and products. These molecules have their characteristic sets of energy levels, but the Boltzmann distribution does not distinguish between their identities, only their energies. The available atoms distribute themselves over both sets of energy levels in accord with the Boltzmann distribution (Fig. 6A.4). At a given temperature, there will be a specific distribution of populations, and hence a specific composition of the reaction mixture.

It can be appreciated from the illustration that, if the reactants and products both have similar arrays of molecular energy levels, then the dominant species in a reaction mixture at equilibrium is the species with the lower set of energy levels (Fig. 6A.4(a)). However, the fact that the Gibbs energy occurs in the expression for the equilibrium constant is a signal that entropy plays a role as well as energy. Its role can be appreciated by referring to Fig. 6A.4. Figure 6A.4(b) shows that, although the $B$ energy levels lie higher than the A energy levels, in this instance they are much more closely spaced. As a result, their total population may be considerable and B could even dominate in the reaction mixture at equilibrium. Closely spaced energy levels correlate with a high entropy (Topic 13E), so in this case entropy effects dominate adverse energy effects. This competition is mirrored in eqn 6A.15, as can be\\
\includegraphics[max width=\textwidth, center]{2025_02_27_d4b95d9718f0b9e2e6b9g-242}

Fig. 6A. 4 The Boltzmann distribution of populations over the energy levels of two species $A$ and $B$. The reaction $A \rightarrow B$ is endothermic in this example. In (a) the two species have similar densities of energy levels: the bulk of the population is associated with the species $A$, so that species is dominant at equilibrium. In (b) the density of energy levels in $B$ is much greater than that in $A$, and as a result, even though the reaction $A \rightarrow B$ is endothermic, the population associated with $B$ is greater than that associated with $A$, so $B$ is dominant at equilibrium.\\
seen most clearly by using $\Delta_{\mathrm{r}} G^{\ominus}=\Delta_{\mathrm{r}} H^{\ominus}-T \Delta_{\mathrm{r}} S^{\ominus}$ and writing it in the form


\begin{equation*}
K=\mathrm{e}^{-\Delta_{\mathrm{r}} H^{\ominus} / R T} \mathrm{e}^{\Delta_{r} S^{\ominus} / R} \tag{6A.19}
\end{equation*}


Note that a positive reaction enthalpy results in a lowering of the equilibrium constant (that is, an endothermic reaction can be expected to have an equilibrium composition that favours the reactants). However, if there is positive reaction entropy, then the equilibrium composition may favour products, despite the endothermic character of the reaction.

\subsection*{Brief illustration 6A. 6}
From data provided in the Resource section it is found that for the reaction $\mathrm{N}_{2}(\mathrm{~g})+3 \mathrm{H}_{2}(\mathrm{~g}) \rightleftharpoons 2 \mathrm{NH}_{3}(\mathrm{~g})$ at $298 \mathrm{~K}, \Delta_{\mathrm{r}} G^{\ominus}$ $=-32.9 \mathrm{~kJ} \mathrm{~mol}^{-1}, \Delta_{\mathrm{r}} H^{\ominus}=-92.2 \mathrm{~kJ} \mathrm{~mol}^{-1}$, and $\Delta_{\mathrm{r}} S^{\ominus}=-198.8 \mathrm{~J} \mathrm{~K}^{-1}$ $\mathrm{mol}^{-1}$. The contributions to $K$ are therefore

$$
\begin{aligned}
K= & \left.\mathrm{e}^{-\left(-9.22 \times 10^{4} \mathrm{Jmol}\right.} \mathrm{mol}^{-1}\right)\left(8.3145 \mathrm{JK}^{-1} \mathrm{~mol}^{-1}\right) \times(298 \mathrm{~K}) \\
& \times \mathrm{e}^{\left(-198.8 \mathrm{~J} \mathrm{~K}^{-1} \mathrm{~mol}^{-1}\right) /\left(8.3145 \mathrm{~J} \mathrm{~K}^{-1} \mathrm{~mol}^{-1}\right)} \\
= & \mathrm{e}^{37.2 \ldots} \times \mathrm{e}^{-23.9 \ldots}
\end{aligned}
$$

Note that the exothermic character of the reaction encourages the formation of products (it results in a large increase in entropy of the surroundings) but the decrease in entropy of the system as H atoms are pinned to N atoms opposes their formation.

\subsection*{Checklist of concepts}
$\square$ 1. The reaction Gibbs energy $\Delta_{\mathrm{r}} G$ is the slope of the plot of Gibbs energy against extent of reaction.2. Reactions that have $\Delta_{\mathrm{r}} G<0$ are classified as exergonic, and those with $\Delta_{\mathrm{r}} G>0$ are classified as endergonic.3. The reaction quotient is a combination of activities used to express the current value of the reaction Gibbs energy.4. The equilibrium constant is the value of the reaction quotient at equilibrium.

\subsection*{Checklist of equations}
\begin{center}
\begin{tabular}{|c|c|c|c|}
\hline
Property & Equation & Comment & Equation number \\
\hline
Reaction Gibbs energy & $\Delta_{\mathrm{r}} G=(\partial G / \partial \xi)_{p, T}$ & Definition & 6A. 1 \\
\hline
Reaction Gibbs energy & $\Delta_{\mathrm{r}} G=\Delta_{\mathrm{r}} G^{\ominus}+R T \ln Q, Q=\prod_{\mathrm{J}} a_{\mathrm{J}}^{v_{\mathrm{J}}}$ & Evaluated at arbitrary stage of reaction & 6A. 12 \\
\hline
Standard reaction Gibbs energy & \( \begin{aligned} \Delta_{\mathrm{r}} G^{\ominus} & =\sum_{\text {Products }} v \Delta_{\mathrm{f}} G^{\ominus}-\sum_{\text {Reactants }} v \Delta_{\mathrm{f}} G^{\ominus} \\ & =\sum_{\mathrm{J}} v_{\mathrm{J}} \Delta_{\mathrm{f}} G^{\ominus}(\mathrm{J}) \end{aligned} \) & $v$ are positive; $v_{\mathrm{J}}$ are signed & 6A. 13 \\
\hline
Equilibrium constant & \( K=\left(\prod_{J} a_{\mathrm{J}}^{v_{\mathrm{J}}}\right)_{\text {equilibrium }} \) & Definition & 6A. 14 \\
\hline
Thermodynamic equilibrium constant & $\Delta_{\mathrm{r}} G^{\ominus}=-R T \ln K$ &  & 6A. 15 \\
\hline
Relation between $K$ and $K_{c}$ & $K=K_{c}\left(c^{\ominus} R T / p^{\ominus}\right)^{\Delta v}$ & Gas-phase reactions; perfect gases & 6A.18b \\
\hline
\end{tabular}
\end{center}

\section*{TOPIC 6B The response of equilibria to the conditions }
\subsection*{Why do you need to know this material?}
Chemists, and chemical engineers designing a chemical plant, need to know how the position of equilibrium will respond to changes in the conditions, such as a change in pressure or temperature. The variation with temperature also provides a way to determine the standard enthalpy and entropy of a reaction.

What is the key idea?\\
A system at equilibrium, when subjected to a disturbance, tends to respond in a way that minimizes the effect of the disturbance.

\subsection*{What do you need to know already?}
This Topic builds on the relation between the equilibrium constant and the standard Gibbs energy of reaction (Topic 6 A). To express the temperature dependence of $K$ it draws on the Gibbs-Helmholtz equation (Topic 3E).

The equilibrium constant for a reaction is not affected by the presence of a catalyst. As explained in detail in Topics 17F and 19C, catalysts increase the rate at which equilibrium is attained but do not affect its position. However, it is important to note that in industry reactions rarely reach equilibrium, partly on account of the rates at which reactants mix and products are extracted. The equilibrium constant is also independent of pressure, but as will be seen, that does not necessarily mean that the composition at equilibrium is independent of pressure. The equilibrium constant does depend on the temperature in a manner that can be predicted from the standard reaction enthalpy.

\subsection*{6B.1 The response to pressure}
The equilibrium constant depends on the value of $\Delta_{\mathrm{r}} G^{\ominus}$, which is defined at a single, standard pressure. The value of $\Delta_{\mathrm{r}} G^{\ominus}$, and hence of $K$, is therefore independent of the pressure at which\\
the equilibrium is actually established. In other words, at a given temperature, $K$ is a constant.

The effect of pressure depends on how the pressure is applied. The pressure within a reaction vessel can be increased by injecting an inert gas into it. However, so long as the gases are perfect, this addition of gas leaves all the partial pressures of the reacting gases unchanged: the partial pressure of a perfect gas is the pressure it would exert if it were alone in the container, so the presence of another gas has no effect on its value. It follows that pressurization by the addition of an inert gas has no effect on the equilibrium composition of the system (provided the gases are perfect).

Alternatively, the pressure of the system may be increased by confining the gases to a smaller volume (that is, by compression). Now the individual partial pressures are changed but their ratio (raised to the various powers that appear in the equilibrium constant) remains the same. Consider, for instance, the perfect gas equilibrium $\mathrm{A}(\mathrm{g}) \rightleftharpoons 2 \mathrm{~B}(\mathrm{~g})$, for which the equilibrium constant is

$$
K=\frac{p_{\mathrm{B}}^{2}}{p_{\mathrm{A}} p^{\ominus}}
$$

The right-hand side of this expression remains constant when the mixture is compressed only if an increase in $p_{\mathrm{A}}$ cancels an increase in the square of $p_{\mathrm{B}}$. This relatively steep increase of $p_{\mathrm{A}}$ compared to $p_{\mathrm{B}}$ will occur if the equilibrium composition shifts in favour of $A$ at the expense of $B$. Then the number of A molecules will increase as the volume of the container is decreased and the partial pressure of A will rise more rapidly than can be ascribed to a simple change in volume alone (Fig. 6B.1).\\
The increase in the number of A molecules and the corresponding decrease in the number of $B$ molecules in the equilibrium $\mathrm{A}(\mathrm{g}) \rightleftharpoons 2 \mathrm{~B}(\mathrm{~g})$ is a special case of a principle proposed by the French chemist Henri Le Chatelier, which states that:

A system at equilibrium, when subjected to a disturbance, tends to respond in a way that minimizes the effect of the disturbance.

The principle implies that, if a system at equilibrium is compressed, then the reaction will tend to adjust so as to mini-\\
\includegraphics[max width=\textwidth, center]{2025_02_27_d4b95d9718f0b9e2e6b9g-245}

Figure 6B. 1 When a reaction at equilibrium is compressed (from $a$ to $b$ ), the reaction responds by reducing the number of molecules in the gas phase (in this case by producing the dimers represented by the linked spheres).\\
mize the increase in pressure. This it can do by reducing the number of particles in the gas phase, which implies a shift $\mathrm{A}(\mathrm{g}) \leftarrow 2 \mathrm{~B}(\mathrm{~g})$.

To treat the effect of compression quantitatively, suppose that there is an amount $n$ of A present initially (and no B). At equilibrium the amount of A is $(1-\alpha) n$ and the amount of B is $2 \alpha n$, where $\alpha$ is the degree of dissociation of A into 2 B . It follows that the mole fractions present at equilibrium are

$$
x_{\mathrm{A}}=\frac{n_{\mathrm{A}}}{n_{\mathrm{tot}}}=\frac{(1-\alpha) n}{(1-\alpha) n+2 \alpha n}=\frac{1-\alpha}{1+\alpha} \quad x_{\mathrm{B}}=\frac{2 \alpha}{1+\alpha}
$$

The equilibrium constant for the reaction is

$$
K=\frac{p_{\mathrm{B}}^{2}}{p_{\mathrm{A}} p^{\ominus}}=\frac{x_{\mathrm{B}}^{2} p^{2}}{x_{\mathrm{A}} p p^{\ominus}}=\frac{4 \alpha^{2}\left(p / p^{\ominus}\right)}{1-\alpha^{2}}
$$

where $p$ is the total pressure. This expression rearranges to


\begin{equation*}
\alpha=\left(\frac{1}{1+4 p / K p^{\ominus}}\right)^{1 / 2} \tag{6B.1}
\end{equation*}


This formula shows that, even though $K$ is independent of pressure, the amounts of A and B do depend on pressure (Fig. 6B.2). It also shows that as $p$ is increased, $\alpha$ decreases, in accord with Le Chatelier's principle.

\subsection*{Brief illustration 6B. 1}
To predict the effect of an increase in pressure on the composition of the ammonia synthesis at equilibrium, $\mathrm{N}_{2}(\mathrm{~g})+$ $3 \mathrm{H}_{2}(\mathrm{~g}) \rightleftharpoons 2 \mathrm{NH}_{3}(\mathrm{~g})$, note that the number of gas molecules decreases (from 4 to 2). Le Chatelier's principle predicts that an increase in pressure favours the product. The equilibrium constant is

$$
K=\frac{p_{\mathrm{NH}_{3}}^{2} p^{\ominus 2}}{p_{\mathrm{N}_{2}} p_{\mathrm{H}_{2}}^{3}}=\frac{x_{\mathrm{NH}_{3}}^{2} p^{2} p^{\ominus 2}}{x_{\mathrm{N}_{2}} x_{\mathrm{H}_{2}}^{3} p^{4}}=\frac{x_{\mathrm{NH}_{3}}^{2} p^{\ominus 2}}{x_{\mathrm{N}_{2}} x_{\mathrm{H}_{2}}^{3} p^{2}}=K_{x} \times \frac{p^{\ominus 2}}{p^{2}}
$$

where $K_{x}$ is the part of the equilibrium constant expression that contains the equilibrium mole fractions of reactants and products (note that, unlike $K$ itself, $K_{x}$ is not an equilibrium constant). Therefore, doubling the pressure must increase $K_{x}$ by a factor of 4 to preserve the value of $K$.\\
\includegraphics[max width=\textwidth, center]{2025_02_27_d4b95d9718f0b9e2e6b9g-245(1)}

Figure 6B.2 The pressure dependence of the degree of dissociation, $\alpha$, at equilibrium for an $\mathrm{A}(\mathrm{g}) \rightleftharpoons 2 \mathrm{~B}(\mathrm{~g})$ reaction for different values of the equilibrium constant $K$ (the line labels). The value $\alpha=0$ corresponds to pure $\mathrm{A} ; \alpha=1$ corresponds to pure B .

\subsection*{6B.2 The response to temperature}
Le Chatelier's principle predicts that a system at equilibrium tends to shift in the endothermic direction if the temperature is raised, for then energy is absorbed as heat and the rise in temperature is opposed. Conversely, an equilibrium can be expected to tend to shift in the exothermic direction if the temperature is lowered, for then energy is released and the reduction in temperature is opposed. These conclusions can be summarized as follows:

Exothermic reactions: increased temperature favours the reactants.

Endothermic reactions: increased temperature favours the products.

\subsection*{(a) The van't Hoff equation}
The response to temperature can be explored quantitatively by deriving an expression for the slope of a plot of the equilibrium constant (specifically, of $\ln K$ ) as a function of temperature.

How is that done? 6B. 1 Deriving an expression for the variation of $\ln K$ with temperature

The starting point for this derivation is eqn 6A. 15 ( $\Delta_{\mathrm{t}} G^{\ominus}=-R T \ln K$ ), in the form

$$
\ln K=-\frac{\Delta_{G} G^{\ominus}}{R T}
$$

Now follow these steps.\\
Step 1 Differentiate the expression for $\ln K$\\
Differentiation of $\ln K$ with respect to temperature gives

$$
\frac{\mathrm{d} \ln K}{\mathrm{~d} T}=-\frac{1}{R} \frac{\mathrm{~d}\left(\Delta_{\mathrm{r}} G^{\ominus} / T\right)}{\mathrm{d} T}
$$

which can be rearranged into

$$
\frac{\mathrm{d}\left(\Delta_{T} G^{\ominus} / T\right)}{\mathrm{d} T}=-R \frac{\mathrm{~d} \ln K}{\mathrm{~d} T}
$$

The differentials are complete (i.e. they are not partial derivatives) because $K$ and $\Delta_{t} G^{\ominus}$ depend only on temperature, not on pressure.\\
Step 2 Use the Gibbs-Helmholtz equation\\
To develop the preceding equation, use the Gibbs-Helmholtz equation (eqn $3 \mathrm{E} .10, \mathrm{~d}(G / T) / \mathrm{d} T=-H / T^{2}$ ) in the form

$$
\frac{\mathrm{d}\left(\Delta_{\mathrm{r}} G^{\ominus} / T\right)}{\mathrm{d} T}=-\frac{\Delta_{\mathrm{r}} H^{\ominus}}{T^{2}}
$$

where $\Delta_{r} H^{\ominus}$ is the standard reaction enthalpy at the temperature $T$. Combining this equation with the expression from Step 1 gives

$$
R \frac{\mathrm{~d} \ln K}{\mathrm{~d} T}=\frac{\Delta_{\mathrm{r}} H^{\ominus}}{T^{2}}
$$

which rearranges into

$$
-\left|\frac{\mathrm{d} \ln K}{\mathrm{~d} T}=\frac{\Delta_{r} H^{\ominus}}{R T^{2}}\right| \text { van't Hoff equation }^{\text {(6B.2) }}
$$

Equation 6 B. 2 is known as the van 't Hoff equation. For a reaction that is exothermic under standard conditions ( $\Delta_{\mathrm{r}} H^{\ominus}<0$ ), it implies that $\mathrm{d} \ln K / \mathrm{d} T<0$ (and therefore that $\mathrm{d} K / \mathrm{d} T<0$ ). A negative slope means that $\ln K$, and therefore $K$ itself, decreases as the temperature rises. Therefore, in line with Le Chatelier's principle, in the case of an exothermic reaction the equilibrium shifts away from products. The opposite occurs in the case of endothermic reactions.\\
Insight into the thermodynamic basis of this behaviour comes from the expression $\Delta_{\mathrm{r}} G^{\ominus}=\Delta_{\mathrm{r}} H^{\ominus}-T \Delta_{\mathrm{r}} S^{\ominus}$ written in the form $-\Delta_{\mathrm{t}} G^{\ominus} / T=-\Delta_{\mathrm{r}} H^{\ominus} / T+\Delta_{\mathrm{t}} S^{\ominus}$. When the reaction is exothermic, $-\Delta_{\mathrm{r}} H^{\circ} / T$ corresponds to a positive change of entropy of the surroundings and favours the formation of products. When the temperature is raised, $-\Delta_{\mathrm{r}} H^{\ominus} / T$ decreases and the\\
\includegraphics[max width=\textwidth, center]{2025_02_27_d4b95d9718f0b9e2e6b9g-246}

Figure 6B.3 The effect of temperature on a chemical equilibrium can be interpreted in terms of the change in the Boltzmann distribution with temperature and the effect of that change in the population of the species. (a) In an endothermic reaction, the population of B increases at the expense of A as the temperature is raised. (b) In an exothermic reaction, the opposite happens.\\
increasing entropy of the surroundings has a less important role. As a result, the equilibrium lies less to the right. When the reaction is endothermic, the contribution of the unfavourable change of entropy of the surroundings is reduced if the temperature is raised (because then $\Delta_{\mathrm{t}} H^{\circ} / T$ is smaller), and the reaction then shifts towards products.\\
These remarks have a molecular basis that stems from the Boltzmann distribution of molecules over the available energy levels (see the Prologue to this text). The typical arrangement of energy levels for an endothermic reaction is shown in Fig. 6B.3a. When the temperature is increased, the Boltzmann distribution adjusts and the populations change as shown. The change corresponds to an increased population of the higher energy states at the expense of the population of the lower energy states. The states that arise from the B molecules become more populated at the expense of the A molecules. Therefore, the total population of $B$ states increases, and $B$ becomes more abundant in the equilibrium mixture. Conversely, if the reaction is exothermic (Fig. 6B.3b), then an increase in temperature increases the population of the A states (which start at higher energy) at the expense of the B states, so the reactants become more abundant.

\subsection*{Example 6B.1 Measuring a standard reaction enthalpy}
The data below show the temperature variation of the equilibrium constant of the reaction $\mathrm{Ag}_{2} \mathrm{CO}_{3}(\mathrm{~s}) \rightleftharpoons \mathrm{Ag}_{2} \mathrm{O}(\mathrm{s})+$ $\mathrm{CO}_{2}(\mathrm{~g})$. Calculate the standard reaction enthalpy of the decomposition.

\begin{center}
\begin{tabular}{lllll}
\hline
$\mathbf{T} / \mathbf{K}$ & $\mathbf{3 5 0}$ & $\mathbf{4 0 0}$ & $\mathbf{4 5 0}$ & $\mathbf{5 0 0}$ \\
\hline
$K$ & $3.98 \times 10^{-4}$ & $1.41 \times 10^{-2}$ & $1.86 \times 10^{-1}$ & 1.48 \\
\hline
\end{tabular}
\end{center}

Collect your thoughts You need to adapt the van 't Hoff equation into a form that corresponds to a straight line. So note that $\mathrm{d}(1 / T) / \mathrm{d} T=-1 / T^{2}$, which implies that $\mathrm{d} T=-T^{2} \mathrm{~d}(1 / T)$. Then, after cancelling the $T^{2}$, eqn 6B. 2 becomes

$$
-\frac{\mathrm{d} \ln K}{\mathrm{~d}(1 / T)}=\frac{\Delta_{\mathrm{r}} H^{\ominus}}{R}
$$

Therefore, provided the standard reaction enthalpy can be assumed to be independent of temperature, a plot of $-\ln K$ against $1 / T$ should be a straight line of slope $\Delta_{\mathrm{r}} H^{\ominus} / R$. The actual dimensionless plot is of $-\ln K$ against $1 /(T / K)$, so equate $\Delta_{\mathrm{r}} H^{\ominus} / R$ to slope $\times \mathrm{K}$.

The solution Draw up the following table:

\begin{center}
\begin{tabular}{llllc}
\hline
$\boldsymbol{T} / \mathbf{K}$ & $\mathbf{3 5 0}$ & $\mathbf{4 0 0}$ & $\mathbf{4 5 0}$ & $\mathbf{5 0 0}$ \\
\hline
$\left(10^{3} \mathrm{~K}\right) / T$ & 2.86 & 2.50 & 2.22 & 2.00 \\
$-\ln K$ & 7.83 & 4.26 & 1.68 & -0.392 \\
\hline
\end{tabular}
\end{center}

These points are plotted in Fig. 6B.4. The slope of the graph is $+9.6 \times 10^{3}$, and it follows from slope $\times \mathrm{K}=\Delta_{\mathrm{r}} H^{\ominus} / R$ that\\
\includegraphics[max width=\textwidth, center]{2025_02_27_d4b95d9718f0b9e2e6b9g-247}

Figure 6B. 4 When $-\ln K$ is plotted against $1 / T$, a straight line is expected with slope equal to $\Delta_{\mathrm{r}} H^{\ominus} / R$ if the standard reaction enthalpy does not vary appreciably with temperature. This is a non-calorimetric method for the measurement of standard reaction enthalpies. The data plotted are from Example 6B.1.

Self-test 6B.1 The equilibrium constant of the reaction $2 \mathrm{SO}_{2}(\mathrm{~g})+\mathrm{O}_{2}(\mathrm{~g}) \rightleftharpoons 2 \mathrm{SO}_{3}(\mathrm{~g})$ is $4.0 \times 10^{24}$ at $300 \mathrm{~K}, 2.5 \times 10^{10}$ at 500 K , and $3.0 \times 10^{4}$ at 700 K . Estimate the standard reaction enthalpy at 500 K .

The temperature dependence of the equilibrium constant provides a non-calorimetric method of determining $\Delta_{\mathrm{r}} H^{\ominus}$. A drawback is that the standard reaction enthalpy is actually temperature-dependent, so the plot is not expected to be perfectly linear. However, the temperature dependence is weak in many cases, so the plot is reasonably straight. In practice, the method is not very accurate, but it is often the only one available.

\subsection*{(b) The value of $K$ at different temperatures}
To find the value of the equilibrium constant at a temperature $T_{2}$ in terms of its value $K_{1}$ at another temperature $T_{1}$, integrate eqn 6B. 2 between these two temperatures:


\begin{equation*}
\ln K_{2}-\ln K_{1}=\frac{1}{R} \int_{T_{1}}^{T_{2}} \frac{\Delta_{\mathrm{r}} H^{\ominus}}{T^{2}} \mathrm{~d} T \tag{6B.3}
\end{equation*}


If $\Delta_{\mathrm{r}} H^{\ominus}$ is supposed to vary only slightly with temperature over the temperature range of interest, it may be taken outside the integral. It follows that

$$
\ln K_{2}-\ln K_{1}=\frac{\Delta_{\mathrm{r}} H^{\ominus}}{R} \overbrace{\int_{T_{1}}^{T_{2}} \frac{1}{T^{2}} \mathrm{~d} T}^{\begin{array}{l}
\text { Integral A.1, } \\
n=-2
\end{array}}
$$

and therefore that

\[
\ln K_{2}-\ln K_{1}=-\frac{\Delta_{\mathrm{r}} H^{\ominus}}{R}\left(\frac{1}{T_{2}}-\frac{1}{T_{1}}\right) \quad \begin{align*}
& \text { Temperature }  \tag{6B.4}\\
& \text { dependence of } K
\end{align*}
\]

\subsection*{Brief illustration 6B. 2}
To estimate the equilibrium constant for the synthesis of ammonia at 500 K from its value at $298 \mathrm{~K}\left(6.1 \times 10^{5}\right.$ for the reaction written as $\left.\mathrm{N}_{2}(\mathrm{~g})+3 \mathrm{H}_{2}(\mathrm{~g}) \rightleftharpoons 2 \mathrm{NH}_{3}(\mathrm{~g})\right)$, use the standard reaction enthalpy, which can be obtained from Table 2C. 4 in the Resource section by using $\Delta_{\mathrm{r}} H^{\ominus}=2 \Delta_{\mathrm{f}} H^{\ominus}\left(\mathrm{NH}_{3}, \mathrm{~g}\right)$, and assume that its value is constant over the range of temperatures. Then, with $\Delta_{\mathrm{r}} H^{\ominus}=-92.2 \mathrm{~kJ} \mathrm{~mol}^{-1}$, from eqn 6 B. 4 it follows that

$$
\begin{aligned}
\ln K_{2} & =\ln \left(6.1 \times 10^{5}\right)-\left(\frac{-9.22 \times 10^{4} \mathrm{~J} \mathrm{~mol}^{-1}}{8.3145 \mathrm{JK}^{-1} \mathrm{~mol}^{-1}}\right) \times\left(\frac{1}{500 \mathrm{~K}}-\frac{1}{298 \mathrm{~K}}\right) \\
& =-1.7 \ldots
\end{aligned}
$$

That is, $K_{2}=0.18$, a lower value than at 298 K , as expected for this exothermic reaction.

\subsection*{Checklist of concepts}
1. The thermodynamic equilibrium constant is independent of the presence of a catalyst and independent of pressure.2. The response of composition to changes in the conditions is summarized by Le Chatelier's principle.\begin{enumerate}
  \setcounter{enumi}{2}
  \item The dependence of the equilibrium constant on the temperature is expressed by the van 't Hoff equation and can be explained in terms of the distribution of molecules over the available states.
\end{enumerate}

\subsection*{Checklist of equations}
\begin{center}
\begin{tabular}{lll}
\hline
Property & Equation & Comment \\
\hline
van 't Hoff equation & $\mathrm{d} \ln K / \mathrm{d} T=\Delta_{\mathrm{r}} H^{\ominus} / R T^{2}$ &  \\
 & $\mathrm{~d} \ln K / \mathrm{d}(1 / T)=-\Delta_{\mathrm{r}} H^{\ominus} / R$ & Alternative version \\
Temperature dependence of equilibrium constant & $\ln K_{2}-\ln K_{1}=-\left(\Delta_{\mathrm{r}} H^{\ominus} / R\right)\left(1 / T_{2}-1 / T_{1}\right)$ & $\Delta_{\mathrm{r}} H^{\ominus}$ assumed constant \\
\hline
\end{tabular}
\end{center}

\section*{TOPIC 6C Electrochemical cells }
\subsection*{Why do you need to know this material?}
One very special case of the material treated in Topic 6B, with enormous fundamental, technological, and economic significance, concerns reactions that take place in electrochemical cells. Moreover, the ability to make very precise measurements of potential differences ('voltages') means that electrochemical methods can be used to determine thermodynamic properties of reactions that may be inaccessible by other methods.

\subsection*{What is the key idea?}
The electrical work that a reaction can perform at constant pressure and temperature is equal to the reaction Gibbs energy.

\subsection*{> What do you need to know already?}
This Topic develops the relation between the Gibbs energy and non-expansion work (Topic 3D). You need to be aware of how to calculate the work of moving a charge through an electrical potential difference (Topic 2A). The equations make use of the definition of the reaction quotient $Q$ and the equilibrium constant $K$ (Topic 6A).

An electrochemical cell consists of two electrodes, or metallic conductors, in contact with an electrolyte, an ionic conductor (which may be a solution, a liquid, or a solid). An electrode and its electrolyte comprise an electrode compartment; the two electrodes may share the same compartment. The various kinds of electrode are summarized in Table 6C.1. Any 'inert

Table 6C. 1 Varieties of electrode

\begin{center}
\begin{tabular}{llll}
\hline
\begin{tabular}{l}
Electrode \\
type \\
\end{tabular} & Designation & \begin{tabular}{l}
Redox \\
couple \\
\end{tabular} & Half-reaction \\
\hline
\begin{tabular}{c}
Metal $/$ \\
metal \\
ion \\
\end{tabular} & $\mathrm{M}(\mathrm{s}) \mid \mathrm{M}^{+}(\mathrm{aq})$ & $\mathrm{M}^{+} / \mathrm{M}$ & $\mathrm{M}^{+}(\mathrm{aq})+\mathrm{e}^{-} \rightarrow \mathrm{M}(\mathrm{s})$ \\
Gas & $\mathrm{Pt}(\mathrm{s})\left|\mathrm{X}_{2}(\mathrm{~g})\right| \mathrm{X}^{+}(\mathrm{aq})$ & $\mathrm{X}^{+} / \mathrm{X}_{2}$ & $\mathrm{X}^{+}(\mathrm{aq})+\mathrm{e}^{-} \rightarrow \frac{1}{2} \mathrm{X}_{2}(\mathrm{~g})$ \\
 & $\mathrm{Pt}(\mathrm{s})\left|\mathrm{X}_{2}(\mathrm{~g})\right| \mathrm{X}^{-}(\mathrm{aq})$ & $\mathrm{X}_{2} / \mathrm{X}^{-}$ & $\frac{1}{2} \mathrm{X}_{2}(\mathrm{~g})+\mathrm{e}^{-} \rightarrow \mathrm{X}^{-}(\mathrm{aq})$ \\
\begin{tabular}{l}
Metal/ \\
insoluble \\
salt \\
\end{tabular} & $\mathrm{M}(\mathrm{s})|\mathrm{MX}(\mathrm{s})| \mathrm{X}^{-}(\mathrm{aq})$ & $\mathrm{MX} / \mathrm{M}, \mathrm{X}^{-}$ & $\mathrm{MX}(\mathrm{s})+\mathrm{e}^{-} \rightarrow \mathrm{M}(\mathrm{s})$ \\
\begin{tabular}{l}
Redox \\
\end{tabular} & $\mathrm{Pt}(\mathrm{s}) \mid \mathrm{M}^{+}(\mathrm{aq}), \mathrm{M}^{2+}(\mathrm{aq})$ & $\mathrm{M}^{2+} / \mathrm{M}^{+}$ & $\mathrm{M}^{2+}(\mathrm{aq})+\mathrm{e}^{-} \rightarrow \mathrm{M}^{+}(\mathrm{aq})$ \\
\hline
\end{tabular}
\end{center}

metal' shown as part of the specification is present to act as a source or sink of electrons, but takes no other part in the reaction other than perhaps acting as a catalyst for it. If the electrolytes are different, the two compartments may be joined by a salt bridge, which is a tube containing a concentrated electrolyte solution (for instance, potassium chloride in agar jelly) that completes the electrical circuit and enables the cell to function. A galvanic cell is an electrochemical cell that produces electricity as a result of the spontaneous reaction occurring inside it. An electrolytic cell is an electrochemical cell in which a nonspontaneous reaction is driven by an external source of current.

\subsection*{6C. 1 Half-reactions and electrodes}
It will be familiar from introductory chemistry courses that oxidation is the removal of electrons from a species, reduction is the addition of electrons to a species, and a redox reaction is a reaction in which there is a transfer of electrons from one species to another. The electron transfer may be accompanied by other events, such as atom or ion transfer, but the net effect is electron transfer and hence a change in oxidation number of an element. The reducing agent (or reductant) is the electron donor; the oxidizing agent (or oxidant) is the electron acceptor. It should also be familiar that any redox reaction may be expressed as the difference of two reduction halfreactions, which are conceptual reactions showing the gain of electrons. Even reactions that are not redox reactions may often be expressed as the difference of two reduction halfreactions. The reduced and oxidized species in a half-reaction form a redox couple. A couple is denoted $\mathrm{Ox} /$ Red and the corresponding reduction half-reaction is written


\begin{equation*}
\mathrm{Ox}+v \mathrm{e}^{-} \rightarrow \mathrm{Red} \tag{6C.1}
\end{equation*}


\subsection*{Brief illustration 6C. 1}
The dissolution of silver chloride in water $\mathrm{AgCl}(\mathrm{s}) \rightarrow \mathrm{Ag}^{+}(\mathrm{aq})+$ $\mathrm{Cl}^{-}(\mathrm{aq})$, which is not a redox reaction, can be expressed as the difference of the following two reduction half-reactions:

$$
\begin{aligned}
& \mathrm{AgCl}(\mathrm{~s})+\mathrm{e}^{-} \rightarrow \mathrm{Ag}(\mathrm{~s})+\mathrm{Cl}^{-}(\mathrm{aq}) \\
& \mathrm{Ag}^{+}(\mathrm{aq})+\mathrm{e}^{-} \rightarrow \mathrm{Ag}(\mathrm{~s})
\end{aligned}
$$

The redox couples are $\mathrm{AgCl} / \mathrm{Ag}, \mathrm{Cl}^{-}$and $\mathrm{Ag}^{+} / \mathrm{Ag}$, respectively.

It is often useful to express the composition of an electrode compartment in terms of the reaction quotient, $Q$, for the halfreaction. This quotient is defined like the reaction quotient for the overall reaction (Topic $6 \mathrm{~A}, Q=\prod a_{\mathrm{J}}^{v_{\mathrm{J}}}$ ), but the electrons are ignored because they are stateless.

\subsection*{Brief illustration 6C. 2}
The reaction quotient for the reduction of $\mathrm{O}_{2}$ to $\mathrm{H}_{2} \mathrm{O}$ in acid solution, $\mathrm{O}_{2}(\mathrm{~g})+4 \mathrm{H}^{+}(\mathrm{aq})+4 \mathrm{e}^{-} \rightarrow 2 \mathrm{H}_{2} \mathrm{O}(\mathrm{l})$, is

$$
Q=\frac{a_{\mathrm{H}_{2} \mathrm{O}}^{2}}{a_{\mathrm{H}^{+}}^{4} a_{\mathrm{O}_{2}}} \approx \frac{p^{\ominus}}{a_{\mathrm{H}^{+}}^{4} p_{\mathrm{O}_{2}}}
$$

The approximations used in the second step are that the activity of water is 1 (because the solution is dilute) and the oxygen behaves as a perfect gas, so $a_{\mathrm{O}_{2}} \approx p_{\mathrm{O}_{2}} / p^{\ominus}$.

The reduction and oxidation processes responsible for the overall reaction in a cell are separated in space: oxidation takes place at one electrode and reduction takes place at the other. As the reaction proceeds, the electrons released in the oxidation $\operatorname{Red}_{1} \rightarrow \mathrm{Ox}_{1}+v \mathrm{e}^{-}$at one electrode travel through the external circuit and re-enter the cell through the other electrode. There they bring about reduction $\mathrm{Ox}_{2}+v \mathrm{e}^{-} \rightarrow \operatorname{Red}_{2}$. The electrode at which oxidation occurs is called the anode; the electrode at which reduction occurs is called the cathode. In a galvanic cell, the cathode has a higher potential than the anode: the species undergoing reduction, $\mathrm{Ox}_{2}$, withdraws electrons from its electrode (the cathode, Fig. 6C.1), so leaving a relative positive charge on it (corresponding to a high potential). At the anode, oxidation results in the transfer of electrons to the electrode, so giving it a relative negative charge (corresponding to a low potential).\\
\includegraphics[max width=\textwidth, center]{2025_02_27_d4b95d9718f0b9e2e6b9g-250(1)}

Figure 6C. 1 When a spontaneous reaction takes place in a galvanic cell, electrons are deposited in one electrode (the site of oxidation, the anode) and collected from another (the site of reduction, the cathode), and so there is a net flow of current which can be used to do work. Note that the + sign of the cathode can be interpreted as indicating the electrode at which electrons enter the cell, and the - sign of the anode is where the electrons leave the cell.\\
\includegraphics[max width=\textwidth, center]{2025_02_27_d4b95d9718f0b9e2e6b9g-250}

Figure 6C. 2 One version of the Daniell cell. The copper electrode is the cathode and the zinc electrode is the anode. Electrons leave the cell from the zinc electrode and enter it again through the copper electrode.

\subsection*{6C. 2 Varieties of cells}
The simplest type of cell has a single electrolyte common to both electrodes (as in Fig. 6C.1). In some cases it is necessary to immerse the electrodes in different electrolytes, as in the 'Daniell cell' in which the redox couple at one electrode is $\mathrm{Cu}^{2+} / \mathrm{Cu}$ and at the other is $\mathrm{Zn}^{2+} / \mathrm{Zn}$ (Fig. 6C.2). In an electrolyte concentration cell, the electrode compartments are identical except for the concentrations of the electrolytes. In an electrode concentration cell the electrodes themselves have different concentrations, either because they are gas electrodes operating at different pressures or because they are amalgams (solutions in mercury) or analogous materials with different concentrations.

\subsection*{(a) Liquid junction potentials}
In a cell with two different electrolyte solutions in contact, as in the Daniell cell, there is an additional source of potential difference across the interface of the two electrolytes. This contribution is called the liquid junction potential, $E_{1 \mathrm{j}}$. Another example of a junction potential is that at the interface between different concentrations of hydrochloric acid. At the junction, the mobile $\mathrm{H}^{+}$ions diffuse into the more dilute solution. The bulkier $\mathrm{Cl}^{-}$ions follow, but initially do so more slowly, which results in a potential difference at the junction. The potential then settles down to a value such that, after that brief initial period, the ions diffuse at the same rates. Electrolyte concentration cells always have a liquid junction; electrode concentration cells do not.\\
The contribution of the liquid junction to the potential difference can be reduced (to about $1-2 \mathrm{mV}$ ) by joining the electrolyte compartments through a salt bridge (Fig. 6C.3). The reason for the success of the salt bridge is that, provided the ions dissolved in the jelly have similar mobilities, then the liquid junction potentials at either end are largely independent of the concentrations of the two dilute solutions, and so nearly cancel.\\
\includegraphics[max width=\textwidth, center]{2025_02_27_d4b95d9718f0b9e2e6b9g-251}

Figure 6C. 3 The salt bridge, essentially an inverted U-tube full of concentrated salt solution in a jelly, has two opposing liquid junction potentials that almost cancel.

\subsection*{(b) Notation}
The following notation is used for electrochemical cells:\\
| An interface between components or phases\\
$\vdots \quad$ A liquid junction\\
|| An interface for which it is assumed that the junction potential has been eliminated

\subsection*{Brief illustration 6C. 3}
A cell in which two electrodes share the same electrolyte is

$$
\operatorname{Pt}(\mathrm{s})\left|\mathrm{H}_{2}(\mathrm{~g})\right| \mathrm{HCl}(\mathrm{aq})|\operatorname{AgCl}(\mathrm{s})| \operatorname{Ag}(\mathrm{s})
$$

The cell in Fig. 6C. 2 is denoted

$$
\mathrm{Zn}(\mathrm{~s})\left|\mathrm{ZnSO}_{4}(\mathrm{aq}): \mathrm{CuSO}_{4}(\mathrm{aq})\right| \mathrm{Cu}(\mathrm{~s})
$$

The cell in Fig. 6C. 3 is denoted

$$
\mathrm{Zn}(\mathrm{~s})\left|\mathrm{ZnSO}_{4}(\mathrm{aq}) \| \mathrm{CuSO}_{4}(\mathrm{aq})\right| \mathrm{Cu}(\mathrm{~s})
$$

An example of an electrolyte concentration cell in which the liquid junction potential is assumed to be eliminated is

$$
\mathrm{Pt}(\mathrm{~s})\left|\mathrm{H}_{2}(\mathrm{~g})\right| \mathrm{HCl}\left(\mathrm{aq}, b_{1}\right)\left|\| \mathrm{HCl}\left(\mathrm{aq}, b_{2}\right)\right| \mathrm{H}_{2}(\mathrm{~g}) \mid \mathrm{Pt}(\mathrm{~s})
$$

\subsection*{6c.3 The cell potential}
The current produced by a galvanic cell arises from the spontaneous chemical reaction taking place inside it. The cell reaction is the reaction in the cell written on the assumption that the right-hand electrode is the cathode, and hence the assumption that the spontaneous reaction is one in which reduction is taking place in the right-hand compartment. If the right-hand electrode is in fact the cathode, then the cell reaction is spontaneous as written. If the left-hand electrode turns\\
out to be the cathode, then the reverse of the corresponding cell reaction is spontaneous.

To write the cell reaction corresponding to a cell diagram, first write the right-hand half-reaction as a reduction. Then subtract from it the left-hand reduction half-reaction (because, by implication, that electrode is the site of oxidation). If necessary, adjust the number of electrons in the two halfreactions to be the same.

\subsection*{Brief illustration 6C. 4}
For the cell $\mathrm{Zn}(\mathrm{s})\left|\mathrm{ZnSO}_{4}(\mathrm{aq}) \| \mathrm{CuSO}_{4}(\mathrm{aq})\right| \mathrm{Cu}(\mathrm{s})$ the two electrodes and their reduction half-reactions are

Right-hand electrode: $\mathrm{Cu}^{2+}(\mathrm{aq})+2 \mathrm{e}^{-} \rightarrow \mathrm{Cu}(\mathrm{s})$\\
Left-hand electrode: $\mathrm{Zn}^{2+}(\mathrm{aq})+2 \mathrm{e}^{-} \rightarrow \mathrm{Zn}(\mathrm{s})$\\
The same number of electrons is involved in each half-reaction. The overall cell reaction is the difference Right - Left:

$$
\mathrm{Cu}^{2+}(\mathrm{aq})+2 \mathrm{e}^{-}-\mathrm{Zn}^{2+}(\mathrm{aq})-2 \mathrm{e}^{-} \rightarrow \mathrm{Cu}(\mathrm{~s})-\mathrm{Zn}(\mathrm{~s})
$$

which, after cancellation of the $2 \mathrm{e}^{-}$, rearranges to

$$
\mathrm{Cu}^{2+}(\mathrm{aq})+\mathrm{Zn}(\mathrm{~s}) \rightarrow \mathrm{Cu}(\mathrm{~s})+\mathrm{Zn}^{2+}(\mathrm{aq})
$$

\subsection*{(a) The Nernst equation}
A cell in which the overall cell reaction has not reached chemical equilibrium can do electrical work as the reaction drives electrons through an external circuit. The work that a given transfer of electrons can accomplish depends on the potential difference between the two electrodes. When the potential difference is large, a given number of electrons travelling between the electrodes can do a lot of electrical work. When the potential difference is small, the same number of electrons can do only a little work. A cell in which the overall reaction is at equilibrium can do no work, and then the potential difference is zero.

According to the discussion in Topic 3D, the maximum non-expansion work a system can do at constant temperature and pressure is given by eqn 3D. $8\left(w_{\text {add,max }}=\Delta G\right)$. In electrochemistry, the additional (non-expansion) work is identified with electrical work, $w_{\mathrm{e}}$ : the system is the cell, and $\Delta G$ is the Gibbs energy of the cell reaction, $\Delta_{\mathrm{r}} G$. Because maximum work is produced when a change occurs reversibly, it follows that, to draw thermodynamic conclusions from measurements of the work that a cell can do, it is necessary to ensure that the cell is operating reversibly. Moreover, it is established in Topic 6A that the reaction Gibbs energy is actually a property relating, through the term $R T \ln Q$, to a specified composition of the reaction mixture. Therefore, the cell must be operating reversibly at a specific, constant composition. Both these conditions are achieved by measuring the potential difference generated by the cell when it is balanced by an exactly opposing source of\\
potential difference so that the cell reaction occurs reversibly, the composition is constant, and no current flows: in effect, the cell reaction is poised for change, but not actually changing. The resulting potential difference is called the cell potential, $E_{\text {cell }}$, of the cell.

A note on good practice The cell potential was formerly, and is still widely, called the electromotive force (emf) of the cell. IUPAC prefers the term 'cell potential' because a potential difference is not a force.

As this introduction has indicated, there is a close relation between the cell potential and the reaction Gibbs energy. It can be established by considering the electrical work that a cell can do.

\subsection*{How is that done? 6C. 1 Establishing the relation between}
 the cell potential and the reaction Gibbs energyConsider the change in $G$ when the cell reaction advances by an infinitesimal amount $\mathrm{d} \xi$ at some composition. From Topic 6 A , specifically the equation $\Delta_{\mathrm{r}} G=(\partial G / \partial \xi)_{T, p}$, it follows that (at constant temperature and pressure)

$$
\mathrm{d} G=\Delta_{\mathrm{r}} \mathrm{Gd} \xi
$$

The maximum non-expansion (electrical) work, $w_{\mathrm{e}}$, that the reaction can do as it advances by $\mathrm{d} \xi$ at constant temperature and pressure is therefore

$$
\mathrm{d} w_{\mathrm{e}}=\Delta_{\mathrm{r}} \mathrm{Gd} \xi
$$

This work is infinitesimal, and the composition of the system is virtually constant when it occurs.

Suppose that the reaction advances by $\mathrm{d} \xi$, then $v \mathrm{~d} \xi$ electrons must travel from the anode to the cathode, where $v$ is the stoichiometric coefficient of the electrons in the half-reactions into which the cell reaction can be divided. The total charge transported between the electrodes when this change occurs is $-v e N_{A} \mathrm{~d} \xi$ (because $v \mathrm{~d} \xi$ is the amount of electrons in moles and the charge per mole of electrons is $-e N_{\mathrm{A}}$ ). Hence, the total charge transported is $-v F d \xi$ because $e N_{\mathrm{A}}=F$, Faraday's constant. The work done when an infinitesimal charge $-v F d \xi$ travels from the anode to the cathode is equal to the product of the charge and the potential difference, $E_{\text {cell }}$ (see Table 2A.1, the entry $\mathrm{d} w=\phi \mathrm{d} Q$ ):

$$
\mathrm{d} w_{\mathrm{e}}=-v F E_{\text {cell }} \mathrm{d} \xi
$$

When this relation is equated to the one above ( $\mathrm{d} w_{\mathrm{e}}=\Delta_{\mathrm{r}} G \mathrm{~d} \xi$ ), the advancement $\mathrm{d} \xi$ cancels, and the resulting expression is

$$
--v F E_{\text {cell }}=\Delta_{\mathrm{r}} G \longmapsto \text { (6C.2) }
$$

This equation is the key connection between electrical measurements on the one hand and thermodynamic properties on the other. It is the basis of all that follows.\\
\includegraphics[max width=\textwidth, center]{2025_02_27_d4b95d9718f0b9e2e6b9g-252}

Figure 6C.4 A spontaneous reaction occurs in the direction of decreasing Gibbs energy. When expressed in terms of a cell potential, the spontaneous direction of change can be expressed in terms of the cell potential, $E_{\text {cell }}$. The cell reaction is spontaneous as written when $E_{\text {cell }}>0$. The reverse reaction is spontaneous when $E_{\text {cell }}<0$. When the cell reaction is at equilibrium, the cell potential is zero.

It follows from eqn 6C. 2 that, by knowing the reaction Gibbs energy at a specified composition, the cell potential is known at that composition. Note that a negative reaction Gibbs energy, signifying a spontaneous cell reaction, corresponds to a positive cell potential, one in which a voltmeter connected to the cell shows that the right-hand electrode (as in the specification of the cell, not necessarily how the cell is arranged on the bench) is the positive electrode. Another way of looking at the content of eqn 6C. 2 is that it shows that the driving power of a cell (that is, its potential difference) is proportional to the slope of the Gibbs energy with respect to the extent of reaction (the significance of $\Delta_{\mathrm{r}} G$ ). It is plausible that a reaction that is far from equilibrium (when the slope is steep) has a strong tendency to drive electrons through an external circuit (Fig. 6C.4). When the slope is close to zero (when the cell reaction is close to equilibrium), the cell potential is small.

\subsection*{Brief illustration 6C. 5}
Equation 6C. 2 provides an electrical method for measuring a reaction Gibbs energy at any composition of the reaction mixture: simply measure the cell potential and convert it to $\Delta_{\mathrm{r}} G$. Conversely, if the value of $\Delta_{\mathrm{r}} G$ is known at a particular composition, then it is possible to predict the cell potential. For example, if $\Delta_{\mathrm{r}} G=-1.0 \times 10^{2} \mathrm{~kJ} \mathrm{~mol}^{-1}$ and $v=1$, then (using $1 \mathrm{~J}=1 \mathrm{CV}$ ):

$$
E_{\text {cell }}=-\frac{\Delta_{\mathrm{r}} G}{v F}=-\frac{\left(-1.0 \times 10^{5} \mathrm{Jmol}^{-1}\right)}{1 \times\left(9.6485 \times 10^{4} \mathrm{Cmol}^{-1}\right)}=1.0 \mathrm{~V}
$$

The reaction Gibbs energy is related to the composition of the reaction mixture by eqn $6 \mathrm{~A} .12\left(\Delta_{\mathrm{r}} G=\Delta_{\mathrm{r}} G^{\ominus}+R T \ln Q\right)$. It\\
follows, on division of both sides by $-v F$ and recognizing that $\Delta_{\mathrm{r}} G /(-v F)=E_{\text {cell }}$, that

$$
E_{\text {cell }}=-\frac{\Delta_{\mathrm{r}} G^{\ominus}}{v F}-\frac{R T}{v F} \ln Q
$$

The first term on the right is written


\begin{equation*}
E_{\text {cell }}^{\ominus}=-\frac{\Delta_{\mathrm{r}} G^{\ominus}}{v F} \tag{6C.3}
\end{equation*}


Standard cell potential\\[0pt]
[definition]\\
and called the standard cell potential. That is, the standard cell potential is the standard reaction Gibbs energy expressed as a potential difference (in volts). It follows that

$$
E_{\text {cell }}=E_{\text {cell }}^{\ominus}-\frac{R T}{v F} \ln Q
$$

Nernst equation (6C.4)\\
This equation for the cell potential in terms of the composition is called the Nernst equation; the dependence that it predicts is summarized in Fig. 6C.5.

Through eqn 6C.4, the standard cell potential can be interpreted as the cell potential when all the reactants and products in the cell reaction are in their standard states, for then all activities are 1 , so $Q=1$ and $\ln Q=0$. However, the fact that the standard cell potential is merely a disguised form of the standard reaction Gibbs energy (eqn 6C.3) should always be kept in mind and underlies all its applications.

\subsection*{Brief illustration 6C. 6}
Because $R T / F=25.7 \mathrm{mV}$ at $25^{\circ} \mathrm{C}$, a practical form of the Nernst equation at this temperature is

$$
E_{\text {cell }}=E_{\text {cell }}^{\ominus}-\frac{25.7 \mathrm{mV}}{v} \ln Q
$$

It then follows that, for a reaction in which $v=1$, if $Q$ is increased by a factor of 10 , then the cell potential decreases by 59.2 mV .\\
\includegraphics[max width=\textwidth, center]{2025_02_27_d4b95d9718f0b9e2e6b9g-253}

Figure 6C. 5 The variation of cell potential with the value of the reaction quotient for the cell reaction for different values of $v$ (the number of electrons transferred). At $298 \mathrm{~K}, R T / F=25.69 \mathrm{mV}$, so the vertical scale refers to multiples of this value.

An important feature of a standard cell potential is that it is unchanged if the chemical equation for the cell reaction is multiplied by a numerical factor. A numerical factor increases the value of the standard Gibbs energy for the reaction. However, it also increases the number of electrons transferred by the same factor, and by eqn 6C. 3 the value of $E_{\text {cell }}^{\ominus}$ remains unchanged. A practical consequence is that a cell potential is independent of the physical size of the cell. In other words, the cell potential is an intensive property.

\subsection*{(b) Cells at equilibrium}
A special case of the Nernst equation has great importance in electrochemistry and provides a link to the discussion of equilibrium in Topic 6A. Suppose the reaction has reached equilibrium; then $Q=K$, where $K$ is the equilibrium constant of the cell reaction. However, a chemical reaction at equilibrium cannot do work, and hence it generates zero potential difference between the electrodes of a galvanic cell. Therefore, setting $E_{\text {cell }}=0$ and $Q=K$ in the Nernst equation gives

\[
E_{\text {cell }}^{\ominus}=\frac{R T}{v F} \ln K \quad \begin{align*}
& \text { Equilibrium constant and }  \tag{6C.5}\\
& \text { standard cell potential }
\end{align*}
\]

This very important equation (which could also have been obtained more directly by substituting eqn $6 \mathrm{~A} .15, \Delta_{\mathrm{r}} G^{\ominus}=$ $-R T \ln K$, into eqn 6 C .3 ) can be used to predict equilibrium constants from measured standard cell potentials.

\subsection*{Brief illustration 6C. 7}
Because the standard potential of the Daniell cell is +1.10 V , the equilibrium constant for the cell reaction $\mathrm{Cu}^{2+}(\mathrm{aq})+$ $\mathrm{Zn}(\mathrm{s}) \rightarrow \mathrm{Cu}(\mathrm{s})+\mathrm{Zn}^{2+}(\mathrm{aq})$, for which $v=2$, is $K=1.5 \times 10^{37}$ at 298 K . That is, the displacement of copper by zinc goes virtually to completion. Note that a cell potential of about 1 V is easily measurable but corresponds to an equilibrium constant that would be impossible to measure by direct chemical analysis.

\subsection*{6c. 4 The determination of thermodynamic functions}
The standard potential of a cell is related to the standard reaction Gibbs energy through eqn 6C. 3 (written as $-v F E_{\text {cell }}^{\ominus}=$ $\Delta_{\mathrm{r}} G^{\ominus}$ ). Therefore, this important thermodynamic quantity can be obtained by measuring $E_{\text {cell }}^{\ominus}$. Its value can then be used to calculate the Gibbs energy of formation of ions by using the convention explained in Topic 3D, that $\Delta_{\mathrm{f}} G^{\ominus}\left(\mathrm{H}^{+}, \mathrm{aq}\right)=0$.

\subsection*{Brief illustration 6C. 8}
The reaction taking place in the cell

$$
\operatorname{Pt}(\mathrm{s})\left|\mathrm{H}_{2}(\mathrm{~g})\right| \mathrm{H}^{+}(\mathrm{aq}) \| \mathrm{Ag}^{+}(\mathrm{aq}) \mid \mathrm{Ag}(\mathrm{~s}) \quad E_{\text {cell }}^{\ominus}=+0.7996 \mathrm{~V}
$$

is

$$
\mathrm{Ag}^{+}(\mathrm{aq})+\frac{1}{2} \mathrm{H}_{2}(\mathrm{~g}) \rightarrow \mathrm{H}^{+}(\mathrm{aq})+\mathrm{Ag}(\mathrm{~s}) \quad \Delta_{\mathrm{r}} G^{\ominus}=-\Delta_{\mathrm{f}} G^{\ominus}\left(\mathrm{Ag}^{+}, \mathrm{aq}\right)
$$

Therefore, with $v=1$,

$$
\begin{aligned}
\Delta_{\mathrm{f}} G^{\ominus}\left(\mathrm{Ag}^{+}, \mathrm{aq}\right) & =-\left(-F E^{\ominus}\right) \\
& =\left(9.6485 \times 10^{4} \mathrm{C} \mathrm{~mol}^{-1}\right) \times(0.7996 \mathrm{~V}) \\
& =+77.15 \mathrm{~kJ} \mathrm{~mol}^{-1}
\end{aligned}
$$

which is in close agreement with the value in Table 2C. 4 of the Resource section.

The temperature coefficient of the standard cell potential, $\mathrm{d} E_{\text {cell }}^{\ominus} / \mathrm{d} T$, gives the standard entropy of the cell reaction. This conclusion follows from the thermodynamic relation $(\partial G / \partial T)_{p}$ $=-S$ derived in Topic 3E and eqn 6C.3, which combine to give

$$
\frac{\mathrm{d} E_{\text {cell }}^{\ominus}}{\mathrm{d} T}=\frac{\Delta_{\mathrm{r}} S^{\ominus}}{v F}
$$

Temperature coefficient\\
of standard cell potential\\
(6C.6)

The derivative is complete (not partial) because $E_{\text {cell }}^{\ominus}$, like $\Delta_{\mathrm{r}} G^{\ominus}$, is independent of the pressure. This is an electrochemical technique for obtaining standard reaction entropies and through them the entropies of ions in solution.

Finally, the combination of the results obtained so far leads to an expression for the standard reaction enthalpy:


\begin{equation*}
\Delta_{\mathrm{r}} H^{\ominus}=\Delta_{\mathrm{r}} G^{\ominus}+T \Delta_{\mathrm{r}} S^{\ominus}=-v F\left(E_{\text {cell }}^{\ominus}-T \frac{\mathrm{~d} E_{\text {cell }}^{\ominus}}{\mathrm{d} T}\right) \tag{6C.7}
\end{equation*}


This expression provides a non-calorimetric method for measuring $\Delta_{\mathrm{r}} H^{\ominus}$ and, through the convention $\Delta_{\mathrm{f}} H^{\ominus}\left(\mathrm{H}^{+}, \mathrm{aq}\right)=0$, the standard enthalpies of formation of ions in solution (Topic 2C).

\subsection*{Example 6C. 1 Using the temperature coefficient of the}
 standard cell potentialThe standard potential of the cell $\mathrm{Pt}(\mathrm{s})\left|\mathrm{H}_{2}(\mathrm{~g})\right| \mathrm{HBr}(\mathrm{aq}) \mid \mathrm{AgBr}(\mathrm{s})$ $\mid \mathrm{Ag}(\mathrm{s})$ was measured over a range of temperatures, and the data were found to fit the following polynomial:

$$
E_{\text {cell }}^{\ominus} / \mathrm{V}=0.07131-4.99 \times 10^{-4}(T / \mathrm{K}-298)-3.45 \times 10^{-6}(T / \mathrm{K}-298)^{2}
$$

The cell reaction is $\operatorname{AgBr}(\mathrm{s})+\frac{1}{2} \mathrm{H}_{2}(\mathrm{~g}) \rightarrow \mathrm{Ag}(\mathrm{s})+\mathrm{HBr}(\mathrm{aq})$, and has $v=1$. Evaluate the standard reaction Gibbs energy, enthalpy, and entropy at 298 K .\\
Collect your thoughts The standard Gibbs energy of reaction is obtained by using eqn 6C. 3 after evaluating $E_{\text {cell }}^{\ominus}$ at 298 K and by using $1 \mathrm{VC}=1 \mathrm{~J}$. The standard reaction entropy is obtained by using eqn 6C.6, which involves differentiating the polynomial with respect to $T$ and then setting $T=298 \mathrm{~K}$. The standard reaction enthalpy is obtained by combining the values of the standard Gibbs energy and entropy.

The solution At $T=298 \mathrm{~K}, E_{\text {cell }}^{\ominus}=0.07131 \mathrm{~V}$, so

$$
\begin{aligned}
\Delta_{\mathrm{r}} G^{\ominus} & =-v F E_{\text {cell }}^{\ominus}=-(1) \times\left(9.6485 \times 10^{4} \mathrm{C} \mathrm{~mol}^{-1}\right) \times(0.07131 \mathrm{~V}) \\
& =-6.880 \times 10^{3} \mathrm{CV} \mathrm{~mol}^{-1}=-6.880 \mathrm{~kJ} \mathrm{~mol}^{-1}
\end{aligned}
$$

The temperature coefficient of the standard cell potential is

$$
\frac{\mathrm{d} E_{\text {cell }}^{\ominus}}{\mathrm{d} T}=-4.99 \times 10^{-4} \mathrm{VK}^{-1}-2\left(3.45 \times 10^{-6}\right)(T / \mathrm{K}-298) \mathrm{VK}^{-1}
$$

At $T=298 \mathrm{~K}$ this expression evaluates to

$$
\frac{\mathrm{d} E_{\text {cell }}^{\ominus}}{\mathrm{d} T}=-4.99 \times 10^{-4} \mathrm{VK}^{-1}
$$

So, from eqn 6C. 6 the standard reaction entropy is

$$
\begin{aligned}
\Delta_{\mathrm{r}} S^{\ominus} & =\nu F \frac{\mathrm{~d} E_{\text {cell }}^{\ominus}}{\mathrm{d} T}=(1) \times\left(9.6485 \times 10^{4} \mathrm{Cmol}^{-1}\right) \times\left(-4.99 \times 10^{-4} \mathrm{VK}^{-1}\right) \\
& =-48.1 \mathrm{~J} \mathrm{~K}^{-1} \mathrm{~mol}^{-1}
\end{aligned}
$$

The negative value stems in part from the elimination of gas in the cell reaction. It then follows that

$$
\begin{aligned}
\Delta_{\mathrm{r}} H^{\ominus}= & \Delta_{\mathrm{r}} G^{\ominus}+T \Delta_{\mathrm{r}} S^{\ominus}=-6.880 \mathrm{~kJ} \mathrm{~mol}^{-1} \\
& +(298 \mathrm{~K}) \times\left(-0.0481 \mathrm{kJK}^{-1} \mathrm{~mol}^{-1}\right) \\
= & -21.2 \mathrm{~kJ} \mathrm{~mol}^{-1}
\end{aligned}
$$

Comment. One difficulty with this procedure lies in the accurate measurement of small temperature coefficients of cell potential. Nevertheless, it is another example of the striking ability of thermodynamics to relate the apparently unrelated, in this case to relate electrical measurements to thermal properties.

Self-test 6C. 1 Predict the standard potential of the Harned cell, $\mathrm{Pt}(\mathrm{s})\left|\mathrm{H}_{2}(\mathrm{~g})\right| \mathrm{HCl}(\mathrm{aq})|\mathrm{AgCl}(\mathrm{s})| \operatorname{Ag}(\mathrm{s})$, at 303 K from tables of thermodynamic data.\\
\includegraphics[max width=\textwidth, center]{2025_02_27_d4b95d9718f0b9e2e6b9g-254}

\subsection*{Checklist of concepts}
$\square$ 1. A cell reaction is expressed as the difference of two reduction half-reactions; each one defines a redox couple.2. Galvanic cells can have different electrodes or electrodes that differ in either the electrolyte or electrode concentration.3. A liquid junction potential arises at the junction of two electrolyte solutions.4. The cell potential is the potential difference measured under reversible conditions. The cell potential is posi-\\
tive if a voltmeter shows that the right-hand electrode (in the specification of the cell) is the positive electrode.5. The Nernst equation relates the cell potential to the composition of the reaction mixture.6. The standard cell potential may be used to calculate the standard Gibbs energy of the cell reaction and hence its equilibrium constant.7. The temperature coefficient of the standard cell potential is used to measure the standard entropy and standard enthalpy of the cell reaction.

\subsection*{Checklist of equations}
\begin{center}
\begin{tabular}{llll}
\hline
Property & Equation & Comment & Equation number \\
\hline
Cell potential and reaction Gibbs energy & $-v F E_{\text {cell }}=\Delta_{\mathrm{r}} G$ & Constant temperature and pressure & 6 C .2 \\
Standard cell potential & $E_{\text {cell }}^{\ominus}=-\Delta_{\mathrm{r}} G^{\ominus} / v F$ & Definition & 6 C .3 \\
Nernst equation & $E_{\text {cell }}=E_{\text {cell }}^{\ominus}-(R T / v F) \ln Q$ & 6 C .4 &  \\
Equilibrium constant of cell reaction & $E_{\text {cell }}^{\ominus}=(R T / v F) \ln K$ & 6 C .5 &  \\
Temperature coefficient of cell potential & $\mathrm{d} E_{\text {cell }}^{\ominus} / \mathrm{d} T=\Delta_{\mathrm{r}} S^{\ominus} / v F$ & 6 C .6 &  \\
\hline
\end{tabular}
\end{center}

\section*{TOPIC 6D Electrode potentials}
Why do you need to know this material?\\
A very powerful, compact, and widely used way to report standard cell potentials is to ascribe a potential to each electrode. Electrode potentials are used in chemistry to assess the oxidizing and reducing power of redox couples and to infer thermodynamic properties, including equilibrium constants.

\subsection*{What is the key idea?}
Each electrode of a cell can be supposed to make a characteristic contribution to the cell potential; redox couples with low electrode potentials tend to reduce those with higher potentials.

What do you need to know already?\\
This Topic develops the concepts in Topic 6C, so you need to understand the concept of cell potential and standard cell potential; it also makes use of the Nernst equation. The measurement of standard potentials makes use of the Debye-Hückel limiting law (Topic 5F).

As explained in Topic 6C, a galvanic cell is a combination of two electrodes. Each electrode can be considered to make a characteristic contribution to the overall cell potential. Although it is not possible to measure the contribution of a single electrode, the potential of one of the electrodes can be defined as zero, so values can be assigned to others on that basis.

\subsection*{6D. 1 Standard potentials}
The specially selected electrode is the standard hydrogen electrode (SHE):


\begin{equation*}
\operatorname{Pt}(\mathrm{s})\left|\mathrm{H}_{2}(\mathrm{~g})\right| \mathrm{H}^{+}(\mathrm{aq}) \quad E^{\ominus}=0 \text { at all temperatures } \tag{6D.1}
\end{equation*}


Standard potentials\\[0pt]
[convention]\\
To achieve standard conditions, the activity of the hydrogen ions must be 1 (i.e. $\mathrm{pH}=0$ ) and the pressure of the hydrogen gas must be 1 bar. ${ }^{1}$ The standard potential, $E^{\ominus}(\mathrm{X})$, of another redox couple X is then equal to the cell potential in which it\\
forms the right-hand electrode and the standard hydrogen electrode is the left-hand electrode:

\[
\operatorname{Pt}(\mathrm{s})\left|\mathrm{H}_{2}(\mathrm{~g})\right| \mathrm{H}^{+}(\mathrm{aq}) \| \mathrm{X} \quad E^{\ominus}(\mathrm{X})=E_{\text {cell }}^{\ominus} \quad \begin{align*}
& \begin{array}{l}
\text { Standard } \\
\text { potentials } \\
\text { [convention] }
\end{array} \tag{6D.2}
\end{align*}
\]

The standard potential of a cell of the form $L \| R$, where $L$ is the left-hand electrode of the cell as written (not as arranged on the bench) and R is the right-hand electrode, is then given by the difference of the two standard (electrode) potentials:

$$
\mathrm{L} \| \mathrm{R} \quad E_{\text {cell }}^{\ominus}=E^{\ominus}(\mathrm{R})-E^{\ominus}(\mathrm{L}) \quad \text { Standard cell potential }
$$

(6D.3)

A list of standard potentials at 298 K is given in Table 6D.1, and longer lists in numerical and alphabetical order are in the Resource section.

Table 6D. 1 Standard potentials at $298 \mathrm{~K}^{*}$

\begin{center}
\begin{tabular}{lc}
\hline
Couple & $\boldsymbol{E}^{\ominus} / \mathbf{V}$ \\
\hline
$\mathrm{Ce}^{4+}(\mathrm{aq})+\mathrm{e}^{-} \rightarrow \mathrm{Ce}^{3+}(\mathrm{aq})$ & +1.61 \\
$\mathrm{Cu}^{2+}(\mathrm{aq})+2 \mathrm{e}^{-} \rightarrow \mathrm{Cu}(\mathrm{s})$ & +0.34 \\
$\mathrm{AgCl}(\mathrm{s})+\mathrm{e}^{-} \rightarrow \mathrm{Ag}(\mathrm{s})+\mathrm{Cl}^{-}(\mathrm{aq})$ & +0.22 \\
$\mathrm{H}^{+}(\mathrm{aq})+\mathrm{e}^{-} \rightarrow \frac{1}{2} \mathrm{H}_{2}(\mathrm{~g})$ & 0 \\
$\mathrm{Zn}^{2+}(\mathrm{aq})+2 \mathrm{e}^{-} \rightarrow \mathrm{Zn}(\mathrm{s})$ & -0.76 \\
$\mathrm{Na}^{+}(\mathrm{aq})+\mathrm{e}^{-} \rightarrow \mathrm{Na}(\mathrm{s})$ & -2.71 \\
\hline
\end{tabular}
\end{center}

\begin{itemize}
  \item More values are given in the Resource section.
\end{itemize}

\subsection*{Brief illustration 6D. 1}
The cell $\mathrm{Ag}(\mathrm{s})|\mathrm{AgCl}(\mathrm{s})| \mathrm{HCl}(\mathrm{aq})\left|\mathrm{O}_{2}(\mathrm{~g})\right| \mathrm{Pt}(\mathrm{s})$ can be regarded as formed from the following two electrodes, with their standard potentials taken from the Resource section:

\begin{center}
\begin{tabular}{lll}
\hline
Electrode & Half-reaction & \begin{tabular}{l}
Standard \\
potential \\
\end{tabular} \\
\hline
$\mathrm{R}: \mathrm{Pt}(\mathrm{s})\left|\mathrm{O}_{2}(\mathrm{~g})\right| \mathrm{H}^{+}(\mathrm{aq})$ & $\mathrm{O}_{2}(\mathrm{~g})+4 \mathrm{H}^{+}(\mathrm{aq})+4 \mathrm{e}^{-} \rightarrow 2 \mathrm{H}_{2} \mathrm{O}(\mathrm{l})$ & +1.23 V \\
$\mathrm{~L}: \mathrm{Ag}(\mathrm{s})|\operatorname{AgCl}(\mathrm{s})| \mathrm{Cl}^{-}(\mathrm{aq})$ & $\mathrm{AgCl}(\mathrm{s})+\mathrm{e}^{-} \rightarrow \mathrm{Ag}(\mathrm{s})+\mathrm{Cl}^{-}(\mathrm{aq})$ & +0.22 V \\
 &  & $E_{\text {cell }}^{\ominus}=$ \\
 & +1.01 V &  \\
\hline
\end{tabular}
\end{center}

\footnotetext{${ }^{1}$ Strictly speaking the fugacity, which is the equivalent of activity for a gas (see A deeper look 2 on the website for this text), should be 1 . This complication is ignored here, which is equivalent to assuming perfect gas behaviour.
}\subsection*{(a) The measurement procedure}
The procedure for measuring a standard potential can be illustrated by considering a specific case, the silver/silver chloride electrode. The measurement is made on the 'Harned cell':

$$
\begin{aligned}
& \mathrm{Pt}(\mathrm{~s})\left|\mathrm{H}_{2}(\mathrm{~g})\right| \mathrm{HCl}(\mathrm{aq}, b)|\mathrm{AgCl}(\mathrm{~s})| \mathrm{Ag}(\mathrm{~s}) \\
& \frac{1}{2} \mathrm{H}_{2}(\mathrm{~g})+\mathrm{AgCl}(\mathrm{~s}) \rightarrow \mathrm{HCl}(\mathrm{aq})+\mathrm{Ag}(\mathrm{~s}) \\
& E_{\text {cell }}^{\ominus}=E^{\ominus}\left(\mathrm{AgCl} / \mathrm{Ag}, \mathrm{Cl}^{-}\right)-E^{\ominus}(\mathrm{SHE})=E^{\ominus}\left(\mathrm{AgCl} / \mathrm{Ag}, \mathrm{Cl}^{-}\right), v=1
\end{aligned}
$$

for which the Nernst equation is

$$
E_{\text {cell }}=E^{\ominus}\left(\mathrm{AgCl} / \mathrm{Ag}, \mathrm{Cl}^{-}\right)-\frac{R T}{F} \ln \frac{a_{\mathrm{H}^{+}} a_{\mathrm{Cl}^{-}}}{a_{\mathrm{H}_{2}}^{1 / 2}}
$$

If the hydrogen gas is at the standard pressure of 1 bar, then $a_{\mathrm{H}_{2}}=1$. For simplicity, writing the standard potential of the $\mathrm{AgCl} / \mathrm{Ag}, \mathrm{Cl}^{-}$electrode as $E^{\ominus}$, turns this equation into

$$
E_{\text {cell }}=E^{\ominus}-\frac{R T}{F} \ln a_{\mathrm{H}^{+}} a_{\mathrm{Cl}^{-}}
$$

The activities in this expression can be written in terms of the molality $b$ of $\mathrm{HCl}(\mathrm{aq})$ through $a_{\mathrm{H}^{+}}=\gamma_{ \pm} b / b^{\ominus}$ and $a_{\mathrm{Cl}^{-}}=\gamma_{ \pm} b / b^{\ominus}$, as established in Topic 5F:

$$
\begin{aligned}
E_{\text {cell }} & =E^{\ominus}-\frac{R T}{F} \ln \frac{\gamma_{ \pm}^{2} b^{2}}{b^{\ominus^{2}}} \\
& =E^{\ominus}-\frac{2 R T}{F} \ln \frac{b}{b^{\ominus}}-\frac{2 R T}{F} \ln \gamma_{ \pm}
\end{aligned}
$$

and therefore

$$
E_{\text {cell }}+\frac{2 R T}{F} \ln \frac{b}{b^{\ominus}}=E^{\ominus}-\frac{2 R T}{F} \ln \gamma_{ \pm}
$$

From the Debye-Hückel limiting law for a 1,1-electrolyte (eqn 5F.27, $\log \gamma_{ \pm}=-A\left|z_{+} z_{-}\right| I^{1 / 2}$, it follows that as $b \rightarrow 0$

$$
\log \gamma_{ \pm}=-A\left|z_{+} z_{-}\right| I^{1 / 2}=-A\left(b / b^{\ominus}\right)^{1 / 2}
$$

Therefore, because $\ln x=\ln 10 \log x$,

$$
\ln \gamma_{ \pm}=\ln 10 \log \gamma_{ \pm}=-(A \ln 10)\left(b / b^{\ominus}\right)^{1 / 2}
$$

The equation for $E_{\text {cell }}$ then becomes

$$
E_{\text {cell }}+\frac{2 R T}{F} \ln \frac{b}{b^{\ominus}}=E^{\ominus}+\frac{2 A R T \ln 10}{F}\left(\frac{b}{b^{\ominus}}\right)^{1 / 2} \quad \text { as } b \rightarrow 0
$$

With the term in blue denoted $C$, this equation becomes.


\begin{equation*}
\overbrace{E_{\text {cell }}+\frac{2 R T}{F} \ln \frac{b}{b^{\ominus}}}^{y}=\overbrace{E^{\ominus}}^{\text {intercept }}+\overbrace{C \times\left(\frac{b}{b^{\ominus}}\right)^{1 / 2}}^{\text {slope } x x} \tag{6D.4}
\end{equation*}


where $C$ is a constant. To use this equation, which has the form $y=$ intercept + slope $\times x$ with $x=\left(b / b^{\ominus}\right)^{1 / 2}$, the expression on the left is evaluated at a range of molalities, plotted against $\left(b / b^{\ominus}\right)^{1 / 2}$, and extrapolated to $b=0$. The intercept at $b^{1 / 2}=0$ is the value\\
of $E^{\ominus}$ for the silver/silver-chloride electrode. In precise work, the $\left(b / b^{\ominus}\right)^{1 / 2}$ term is brought to the left, and a higher-order correction term from extended versions of the Debye-Hückel law (Topic 5 F ) is used on the right.

\subsection*{Example 6D. 1 Evaluating a standard potential}
The potential of the Harned cell at $25^{\circ} \mathrm{C}$ has the following values:

\begin{center}
\begin{tabular}{llllc}
$b /\left(10^{-3} b^{\ominus}\right)$ & 3.215 & 5.619 & 9.138 & 25.63 \\
$E_{\text {cell }} / \mathrm{V}$ & 0.52053 & 0.49257 & 0.46860 & 0.41824 \\
\end{tabular}
\end{center}

Determine the standard potential of the silver/silver chloride electrode.

Collect your thoughts As explained in the text, evaluate $y=E_{\text {cell }}$ $+(2 R T / F) \ln \left(b / b^{\ominus}\right)$ and plot it against $\left(b / b^{\ominus}\right)^{1 / 2}$; then extrapolate to $b=0$.

The solution To determine the standard potential of the cell, draw up the following table, using $2 R T / F=0.05139 \mathrm{~V}$ :

\begin{center}
\begin{tabular}{llllc}
$b /\left(10^{-3} b^{\ominus}\right)$ & 3.215 & 5.619 & 9.138 & 25.63 \\
$\left\{b /\left(10^{-3} b^{\ominus}\right)\right\}^{1 / 2}$ & 1.793 & 2.370 & 3.023 & 5.063 \\
$E_{\text {cell }} / \mathrm{V}$ & 0.52053 & 0.49257 & 0.46860 & 0.41824 \\
$y / \mathrm{V}$ & 0.2256 & 0.2263 & 0.2273 & 0.2299 \\
\end{tabular}
\end{center}

The data are plotted in Fig. 6D.1; as can be seen, they extrapolate to $E^{\ominus}=+0.2232 \mathrm{~V}$ (the value obtained, to preserve the precision of the data, by linear regression).\\
\includegraphics[max width=\textwidth, center]{2025_02_27_d4b95d9718f0b9e2e6b9g-257}

Figure 6D. 1 The plot and the extrapolation used for the experimental measurement of a standard cell potential. The intercept at $b^{1 / 2}=0$ is $E_{\text {cell }}^{\ominus}$.

Self-test 6D.1 The following data are for the cell $\operatorname{Pt}(\mathrm{s})\left|\mathrm{H}_{2}(\mathrm{~g})\right|$ $\operatorname{HBr}(\mathrm{aq}, b)|\operatorname{AgBr}(\mathrm{s})| \operatorname{Ag}(\mathrm{s})$ at $25^{\circ} \mathrm{C}$ and with the hydrogen gas at 1 bar. Determine the standard cell potential.

\begin{center}
\begin{tabular}{lllc}
$b /\left(10^{-4} b^{\ominus}\right)$ & 4.042 & 8.444 & 37.19 \\
$E_{\text {cell }} / \mathrm{V}$ & 0.46942 & 0.43636 & 0.36173 \\
\end{tabular}
\end{center}

\subsection*{(b) Combining measured values}
The standard potentials in Table 6D. 1 may be combined to give values for couples that are not listed there. However, to do so, it is necessary to take into account the fact that different couples might correspond to the transfer of different numbers of electrons. The procedure is illustrated in the following Example.

\subsection*{Example 6D. 2}
\subsection*{Evaluating a standard potential from}
\subsection*{two others}
Given that the standard potentials of the $\mathrm{Cu}^{2+} / \mathrm{Cu}$ and $\mathrm{Cu}^{+} / \mathrm{Cu}$ couples are +0.340 V and +0.522 V , respectively, evaluate $E^{\ominus}\left(\mathrm{Cu}^{2+}, \mathrm{Cu}^{+}\right)$.

Collect your thoughts First, note that reaction Gibbs energies may be added (as in a Hess's law analysis of reaction enthalpies). Therefore, you should convert the $E^{\ominus}$ values to $\Delta_{\mathrm{r}} G^{\ominus}$ values by using eqn $6 \mathrm{C} .3\left(-v F E^{\ominus}=\Delta_{\mathrm{r}} G^{\ominus}\right)$, add them appropriately, and then convert the overall $\Delta_{\mathrm{r}} G^{\ominus}$ to the required $E^{\ominus}$ by using eqn 6 C. 3 again. This roundabout procedure is necessary because, as seen below, although the factor $F$ cancels (and should be kept in place until it cancels), the factor $v$ in general does not cancel.

The solution The electrode half-reactions are as follows:\\
(a) $\mathrm{Cu}^{2+}(\mathrm{aq})+2 \mathrm{e}^{-} \rightarrow \mathrm{Cu}(\mathrm{s})$

$$
E^{\ominus}(\mathrm{a})=+0.340 \mathrm{~V} \text {, so } \Delta_{\mathrm{r}} G^{\ominus}(\mathrm{a})=-2(0.340 \mathrm{~V}) F
$$

(b) $\mathrm{Cu}^{+}(\mathrm{aq})+\mathrm{e}^{-} \rightarrow \mathrm{Cu}(\mathrm{s})$

$$
E^{\ominus}(\mathrm{b})=+0.522 \mathrm{~V} \text {, so } \Delta_{\mathrm{r}} G^{\ominus}(\mathrm{b})=-(0.522 \mathrm{~V}) F
$$

The required reaction is\\
(c) $\mathrm{Cu}^{2+}(\mathrm{aq})+\mathrm{e}^{-} \rightarrow \mathrm{Cu}^{+}(\mathrm{aq}) \quad E^{\ominus}(\mathrm{c})=-\Delta_{\mathrm{r}} G^{\ominus}(\mathrm{c}) / F$

Because $(c)=(a)-(b)$, the standard Gibbs energy of reaction (c) is

$$
\begin{aligned}
\Delta_{\mathrm{r}} G^{\ominus}(\mathrm{c}) & =\Delta_{\mathrm{r}} G^{\ominus}(\mathrm{a})-\Delta_{\mathrm{r}} G^{\ominus}(\mathrm{b})=-(0.680 \mathrm{~V}) F-(-0.522 \mathrm{~V}) F \\
& =(-0.158 \mathrm{~V}) F
\end{aligned}
$$

Therefore, $E^{\ominus}(\mathrm{c})=-\Delta_{\mathrm{r}} G^{\ominus}(\mathrm{c}) / F=+0.158 \mathrm{~V}$.\\
Self-test 6D.2 Evaluate $E^{\ominus}\left(\mathrm{Fe}^{3+}, \mathrm{Fe}^{2+}\right)$ from $E^{\ominus}\left(\mathrm{Fe}^{3+}, \mathrm{Fe}\right)$ and $E^{\ominus}\left(\mathrm{Fe}^{2+}, \mathrm{Fe}\right)$.\\
$\Lambda 9 L^{\circ} 0+: \iota \partial s u_{V}$

The generalization of the calculation in the Example is

\[
v_{\mathrm{c}} E^{\ominus}(\mathrm{c})=v_{\mathrm{a}} E^{\ominus}(\mathrm{a})-v_{\mathrm{b}} E^{\ominus}(\mathrm{b}) \quad \begin{align*}
& \text { Combination of }  \tag{6D.5}\\
& \text { standard potentials }
\end{align*}
\]

with the $v_{\mathrm{r}}$ the stoichiometric coefficients of the electrons in each half-reaction.

\subsection*{6D.2 Applications of standard potentials}
Cell potentials are a convenient source of data on equilibrium constants and the Gibbs energies, enthalpies, and entropies of reactions. In practice the standard values of these quantities are the ones normally determined.

\subsection*{(a) The electrochemical series}
For two redox couples, $\mathrm{Ox}_{\mathrm{L}} / \operatorname{Red}_{\mathrm{L}}$ and $\mathrm{Ox}_{\mathrm{R}} / \operatorname{Red}_{\mathrm{R}}$, and the cell\\
$\mathrm{L}\left\|\mathrm{R}=\mathrm{Ox}_{\mathrm{L}} / \operatorname{Red}_{\mathrm{L}}\right\| \mathrm{Ox}_{\mathrm{R}} / \operatorname{Red}_{\mathrm{R}}$\\
$\mathrm{Ox}_{\mathrm{R}}+v \mathrm{e}^{-} \rightarrow \operatorname{Red}_{\mathrm{R}} \quad \mathrm{Ox}_{\mathrm{L}}+v \mathrm{e}^{-} \rightarrow \operatorname{Red}_{\mathrm{L}} \quad$ Cell convention (6D.6a)\\
$E_{\text {cell }}^{\ominus}=E^{\ominus}(\mathrm{R})-E^{\ominus}(\mathrm{L})$\\
the cell reaction


\begin{equation*}
\mathrm{R}-\mathrm{L}: \operatorname{Red}_{\mathrm{L}}+\mathrm{Ox}_{\mathrm{R}} \rightarrow \mathrm{Ox}_{\mathrm{L}}+\operatorname{Red}_{\mathrm{R}} \tag{6D.6b}
\end{equation*}


has $K>1$ if $E_{\text {cell }}^{\ominus}>0$, and therefore if $E^{\ominus}(\mathrm{L})<E^{\ominus}(\mathrm{R})$. Because in the cell reaction $\operatorname{Red}_{\mathrm{L}}$ reduces $\mathrm{Ox}_{\mathrm{R}}$, it follows that\\
$\operatorname{Red}_{\mathrm{L}}$ has a thermodynamic tendency (in the sense $K>1$ ) to reduce $\mathrm{Ox}_{\mathrm{R}}$ if $E^{\ominus}(\mathrm{L})<E^{\ominus}(\mathrm{R})$.

More briefly: low reduces high.\\
Table 6D. 2 shows a part of the electrochemical series, the metallic elements (and hydrogen) arranged in the order of their reducing power as measured by their standard potentials in aqueous solution. A metal low in the series (with a lower standard potential) can reduce the ions of metals with higher standard potentials. This conclusion is qualitative. The quantitative value of $K$ is obtained by doing the calculations described previously and reviewed below.

\subsection*{Brief illustration 6D. 2}
Zinc lies above magnesium in the electrochemical series, so zinc cannot reduce magnesium ions in aqueous solution. Zinc can reduce hydrogen ions, because hydrogen lies higher in the series. However, even for reactions that are thermodynamically favourable, there may be kinetic factors that result in very slow rates of reaction.

\subsection*{(b) The determination of activity coefficients}
Once the standard potential of an electrode in a cell is known, it can be used to determine mean activity coefficients by measuring the cell potential with the ions at the concentration of

Table 6D. 2 The electrochemical series*

\begin{verbatim}
Least strongly reducing
Gold (Au 3+}/\textrm{Au}
Platinum ( }\mp@subsup{\textrm{Pt}}{}{2+}/\textrm{Pt}\mathrm{ )
Silver ( }\mp@subsup{\textrm{Ag}}{}{+}/\textrm{Ag}
Mercury (Hg}\mp@subsup{}{}{2+}/\textrm{Hg}
Copper (Cu 2+}/\textrm{Cu}
Hydrogen (H+
Tin(Sn
Nickel (Ni' }\mp@subsup{}{}{2+}/\textrm{Ni}
Iron (Fe
Zinc( }\mp@subsup{\textrm{Zn}}{}{2+}/\textrm{Zn}
Chromium ( }\mp@subsup{\textrm{Cr}}{}{3+}/\textrm{Cr}\mathrm{ )
Aluminium (Al }\mp@subsup{}{}{3+}/\textrm{Al}
Magnesium (Mg 2+}/\textrm{Mg}
Sodium ( }\mp@subsup{\textrm{Na}}{}{+}/\textrm{Na}
Calcium (Ca }\mp@subsup{}{}{2+}/\textrm{Ca}
Potassium ( }\mp@subsup{\textrm{K}}{}{+}/\textrm{K}
Most strongly reducing
\end{verbatim}

\begin{itemize}
  \item The complete series can be inferred from Table 6D. 1 in the Resource section.\\
interest. For example, in the Harned cell analysed in Section 6D.1, the mean activity coefficient of the ions in hydrochloric acid of molality $b$ is obtained from the relation
\end{itemize}

$$
E_{\text {cell }}+\frac{2 R T}{F} \ln \frac{b}{b^{\ominus}}=E^{\ominus}-\frac{2 R T}{F} \ln \gamma_{ \pm}
$$

which can be rearranged into


\begin{equation*}
\ln \gamma_{ \pm}=\frac{E^{\ominus}-E_{\text {cell }}}{2 R T / F}-\ln \frac{b}{b^{\ominus}} \tag{6D.7}
\end{equation*}


\subsection*{Checklist of concepts}
\begin{enumerate}
  \item The standard potential of a couple is the potential of a cell in which the couple forms the right-hand electrode and the left-hand electrode is a standard hydrogen electrode, all species being present at unit activity.2. The electrochemical series lists the metallic elements in the order of their reducing power as measured by their standard potentials in aqueous solution: low reduces high.
\end{enumerate}

\subsection*{Brief illustration 6D. 3}
The data in Example 6D. 1 include the fact that $E_{\text {cell }}=0.46860 \mathrm{~V}$ when $b=9.138 \times 10^{-3} b^{\ominus}$. Because $2 R T / F=0.05139 \mathrm{~V}$, and in the Example it is established that $E_{\text {cell }}^{\ominus}=0.2232 \mathrm{~V}$, the mean activity coefficient at this molality is

$$
\ln \gamma_{ \pm}=\frac{0.2232 \mathrm{~V}-0.46860 \mathrm{~V}}{0.05139 \mathrm{~V}}-\ln \left(9.138 \times 10^{-3}\right)=-0.0799 \ldots
$$

Therefore, $\gamma_{ \pm}=0.9232$.

\subsection*{(c) The determination of equilibrium constants}
The principal use for standard potentials is to calculate the standard potential of a cell formed from any two electrodes and then to use that value to evaluate the equilibrium constant of the cell reaction. To do so, construct $E_{\text {cell }}^{\ominus}=E^{\ominus}(\mathrm{R})-E^{\ominus}(\mathrm{L})$ and then use eqn 6 C .5 of Topic $6 \mathrm{C}\left(E_{\text {cell }}^{\ominus}=(R T / \nu F) \ln K\right.$, arranged into $\left.\ln K=v F E_{\text {cell }}^{\ominus} / R T\right)$.

\subsection*{Brief illustration 6D. 4}
A disproportionation reaction is a reaction in which a species is both oxidized and reduced. To study the disproportionation $2 \mathrm{Cu}^{+}(\mathrm{aq}) \rightarrow \mathrm{Cu}(\mathrm{s})+\mathrm{Cu}^{2+}(\mathrm{aq})$ at 298 K , combine the following electrodes:\\
$\begin{array}{lll}\mathrm{R}: \mathrm{Cu}(\mathrm{s}) \mid \mathrm{Cu}^{+}(\mathrm{aq}) & \mathrm{Cu}^{+}(\mathrm{aq})+\mathrm{e}^{-} \rightarrow \mathrm{Cu}(\mathrm{s}) & E^{\ominus}(\mathrm{R})=+0.52 \mathrm{~V} \\ \mathrm{~L}: \mathrm{Pt}(\mathrm{s}) \mid \mathrm{Cu}^{2+}(\mathrm{aq}), \mathrm{Cu}^{+}(\mathrm{aq}) & \mathrm{Cu}^{2+}(\mathrm{aq})+\mathrm{e}^{-} \rightarrow \mathrm{Cu}^{+}(\mathrm{aq}) & E^{\ominus}(\mathrm{L})=+0.16 \mathrm{~V}\end{array}$\\
The cell reaction is therefore $2 \mathrm{Cu}^{+}(\mathrm{aq}) \rightarrow \mathrm{Cu}(\mathrm{s})+\mathrm{Cu}^{2+}(\mathrm{aq})$, and the standard cell potential is

$$
E_{\text {cell }}^{\ominus}=0.52 \mathrm{~V}-0.16 \mathrm{~V}=+0.36 \mathrm{~V}
$$

Now calculate the equilibrium constant of the cell reaction. Because $v=1$, from eqn 6C. 5 with $R T / F=0.025693 \mathrm{~V}$,\\
$\ln K=\frac{0.36 \mathrm{~V}}{0.025693 \mathrm{~V}}=14.0 \ldots$\\
Hence, $K=1.2 \times 10^{6}$.\\
3. The difference of the cell potential from its standard value is used to measure the activity coefficient of ions in solution.4. Standard potentials are used to calculate the standard cell potential and then to calculate the equilibrium constant of the cell reaction.

\subsection*{Checklist of equations}
\begin{center}
\begin{tabular}{llll}
\hline
Property & Equation & Comment & Equation number \\
\hline
Standard cell potential from standard potentials & $E_{\text {cell }}^{\ominus}=E^{\ominus}(\mathrm{R})-E^{\ominus}(\mathrm{L})$ & Cell: $\mathrm{L} \| \mathrm{R}$ & 6 D .3 \\
Combined standard potentials & $v_{\mathrm{c}} E^{\ominus}(\mathrm{c})=v_{\mathrm{a}} E^{\ominus}(\mathrm{a})-v_{\mathrm{b}} E^{\ominus}(\mathrm{b})$ &  & 6 D .5 \\
\hline
\end{tabular}
\end{center}

\chapter{FOCUS 7 Quantum theory}
It was once thought that the motion of atoms and subatomic particles could be expressed using 'classical mechanics', the laws of motion introduced in the seventeenth century by Isaac Newton, for these laws were very successful at explaining the motion of everyday objects and planets. However, a proper description of electrons, atoms, and molecules requires a different kind of mechanics, 'quantum mechanics', which is introduced in this Focus and applied widely throughout the text.

\subsection*{7A The origins of quantum mechanics}
Experimental evidence accumulated towards the end of the nineteenth century showed that classical mechanics failed when it was applied to particles as small as electrons. More specifically, careful measurements led to the conclusion that particles may not have an arbitrary energy and that the classical concepts of a particle and wave blend together. This Topic shows how these observations set the stage for the development of the concepts and equations of quantum mechanics in the early twentieth century.\\
7A. 1 Energy quantization; 7A. 2 Wave-particle duality

\subsection*{7B Wavefunctions}
In quantum mechanics, all the properties of a system are expressed in terms of a wavefunction which is obtained by solving the equation proposed by Erwin Schrödinger. This Topic focuses on the interpretation of the wavefunction, and specifically what it reveals about the location of a particle.\\
7B. 1 The Schrödinger equation; 7B. 2 The Born interpretation

\subsection*{7C Operators and observables}
A central feature of quantum theory is its representation of observables by 'operators', which act on the wavefunction and extract the information it contains. This Topic shows how op-\\
erators are constructed and used. One consequence of their use is the 'uncertainty principle', one of the most profound departures of quantum mechanics from classical mechanics.\\
7C. 1 Operators; 7C. 2 Superpositions and expectation values; 7C. 3 The uncertainty principle; 7C. 4 The postulates of quantum mechanics

\subsection*{7D Translational motion}
Translational motion, motion through space, is one of the fundamental types of motion treated by quantum mechanics. According to quantum theory, a particle constrained to move in a finite region of space is described by only certain wavefunctions and can possess only certain energies. That is, quantization emerges as a natural consequence of solving the Schrödinger equation and the conditions imposed on it. The solutions also expose a number of non-classical features of particles, especially their ability to tunnel into and through regions where classical physics would forbid them to be found.\\
7D. 1 Free motion in one dimension; 7D. 2 Confined motion in one dimension; 7D. 3 Confined motion in two and more dimensions; 7D. 4 Tunnelling

\subsection*{7E Vibrational motion}
This Topic introduces the 'harmonic oscillator', a simple but very important model for the description of vibrations. It shows that the energies of an oscillator are quantized and that an oscillator may be found at displacements that are forbidden by classical physics.\\
7E. 1 The harmonic oscillator; 7E. 2 Properties of the harmonic oscillator

\subsection*{7F Rotational motion}
The constraints on the wavefunctions of a body rotating in two and three dimensions result in the quantization of its energy.

In addition, because the energy is related to the angular momentum, it follows that angular momentum is also restricted to certain values. The quantization of angular momentum is a very important aspect of the quantum theory of electrons in atoms and of rotating molecules.\\
7F. 1 Rotation in two dimensions; 7F. 2 Rotation in three dimensions

\subsection*{Web resources What is an application of this material?}
Impact 11 highlights an application of quantum mechanics which still requires much research before it becomes a useful technology. It is based on the expectation that a 'quantum\\
computer' can carry out calculations on many states of a system simultaneously, leading to a new generation of very fast computers. 'Nanoscience' is the study of atomic and molecular assemblies with dimensions ranging from 1 nm to about 100 nm , and 'nanotechnology' is concerned with the incorporation of such assemblies into devices. Impact 12 explores quantum mechanical effects that show how the properties of a nanometre-sized assembly depend on its size.

\section*{TOPIC 7A The origins of quantum mechanics }
\subsection*{Why do you need to know this material?}
Quantum theory is central to almost every explanation in chemistry. It is used to understand atomic and molecular structure, chemical bonds, and most of the properties of matter.

\subsection*{What is the key idea?}
Experimental evidence led to the conclusion that energy can be transferred only in discrete amounts, and that the classical concepts of a 'particle' and a 'wave' blend together.

\subsection*{- What do you need to know already?}
You should be familiar with the basic principles of classical mechanics, especially momentum, force, and energy set out in The chemist's toolkits 3 (in Topic 1B) and 6 (in Topic 2A). The discussion of heat capacities of solids makes light use of material in Topic 2A.

The classical mechanics developed by Newton in the seventeenth century is an extraordinarily successful theory for describing the motion of everyday objects and planets. However,\\
late in the nineteenth century scientists started to make observations that could not be explained by classical mechanics. They were forced to revise their entire conception of the nature of matter and replace classical mechanics by a theory that became known as quantum mechanics.

\subsection*{7A. 1 Energy quantization}
Three experiments carried out near the end of the nineteenth century drove scientists to the view that energy can be transferred only in discrete amounts.

\subsection*{(a) Black-body radiation}
The key features of electromagnetic radiation according to classical physics are described in The chemist's toolkit 13. It is observed that all objects emit electromagnetic radiation over a range of frequencies with an intensity that depends on the temperature of the object. A familiar example is a heated metal bar that first glows red and then becomes 'white hot' upon further heating. As the temperature is raised, the colour shifts from red towards blue and results in the white glow.

\subsection*{The chemist's toolkit 13}
\subsection*{Electromagnetic radiation}
Electromagnetic radiation consists of oscillating electric and magnetic disturbances that propagate as waves. The two components of an electromagnetic wave are mutually perpendicular and are also perpendicular to the direction of propagation (Sketch 1). Electromagnetic waves travel through a vacuum at a constant speed called the speed of light, $c$, which has the defined value of exactly $2.99792458 \times 10^{8} \mathrm{~m} \mathrm{~s}^{-1}$.\\
\includegraphics[max width=\textwidth, center]{2025_02_27_d4b95d9718f0b9e2e6b9g-269(1)}

Sketch 1

A wave is characterized by its wavelength, $\lambda$ (lambda), the distance between consecutive peaks of the wave (Sketch 2). The classification of electromagnetic radiation according to its wavelength is shown in Sketch 3. Light, which is electromagnetic radiation that is visible to the human eye, has a wavelength in the range 420 nm (violet light) to 700 nm (red light). The properties of a wave may also be expressed in terms of its frequency, $v(\mathrm{nu})$, the number of oscillations in a time interval divided by the duration of the interval. Frequency is reported in hertz, Hz , with $1 \mathrm{~Hz}=1 \mathrm{~s}^{-1}$ (i.e. 1 cycle per second). Light spans the frequency range from 710 THz (violet light) to 430 THz (red light).\\
\includegraphics[max width=\textwidth, center]{2025_02_27_d4b95d9718f0b9e2e6b9g-269}

Sketch 2\\
\includegraphics[max width=\textwidth, center]{2025_02_27_d4b95d9718f0b9e2e6b9g-270(1)}

Sketch 3

The wavelength and frequency of an electromagnetic wave are related by:

$$
c=\lambda v
$$

The relation between wavelength and frequency in a vacuum

It is also common to describe a wave in terms of its wavenumber, $\tilde{v}$ (nu tilde), which is defined as

$$
\tilde{v}=\frac{1}{\lambda} \text { or equivalently } \tilde{v}=\frac{v}{c}
$$

Wavenumber [definition]

Thus, wavenumber is the reciprocal of the wavelength and can be interpreted as the number of wavelengths in a given distance. In spectroscopy, for historical reasons, wavenumber is usually reported in units of reciprocal centimetres $\left(\mathrm{cm}^{-1}\right)$. Visible light therefore corresponds to electromagnetic radiation with a wavenumber of $14000 \mathrm{~cm}^{-1}$ (red light) to $24000 \mathrm{~cm}^{-1}$ (violet light).\\
Electromagnetic radiation that consists of a single frequency (and therefore single wavelength) is monochromatic, because it corresponds to a single colour. White light consists of electromagnetic waves with a continuous, but not uniform, spread of frequencies throughout the visible region of the spectrum.\\
A characteristic property of waves is that they interfere with one another, which means that they result in a greater amplitude where their displacements add and a smaller amplitude

The radiation emitted by hot objects is discussed in terms of a black body, a body that emits and absorbs electromagnetic radiation without favouring any wavelengths. A good approximation to a black body is a small hole in an empty container (Fig. 7A.1). Figure 7A. 2 shows how the intensity of the radiation from a black body varies with wavelength at several temperatures. At each temperature $T$ there is a wavelength, $\lambda_{\text {max }}$, at which the intensity of the radiation is a maximum, with $T$ and $\lambda_{\text {max }}$ related by the empirical Wien's law:

$$
\lambda_{\max } T=2.9 \times 10^{-3} \mathrm{mK}
$$

Wien's law\\
(7A.1)\\
where their displacements subtract (Sketch 4). The former is called 'constructive interference' and the latter 'destructive interference'. The regions of constructive and destructive interference show up as regions of enhanced and diminished intensity. The phenomenon of diffraction is the interference caused by an object in the path of waves and occurs when the dimensions of the object are comparable to the wavelength of the radiation. Light waves, with wavelengths of the order of 500 nm , are diffracted by narrow slits.\\
\includegraphics[max width=\textwidth, center]{2025_02_27_d4b95d9718f0b9e2e6b9g-270}

Sketch 4

The intensity of the emitted radiation at any temperature declines sharply at short wavelengths (high frequencies). The intensity is effectively a window on to the energy present inside the container, in the sense that the greater the intensity at a given wavelength, the greater is the energy inside the container due to radiation at that wavelength.

The energy density, $\mathcal{E}(T)$, is the total energy inside the container divided by its volume. The energy spectral density, $\rho(\lambda, T)$, is defined so that $\rho(\lambda, T) \mathrm{d} \lambda$ is the energy density at temperature $T$ due to the presence of electromagnetic radiation with wavelengths between $\lambda$ and $\lambda+\mathrm{d} \lambda$. A high energy\\
\includegraphics[max width=\textwidth, center]{2025_02_27_d4b95d9718f0b9e2e6b9g-271(1)}

Figure 7A. 1 Black-body radiation can be detected by allowing it to leave an otherwise closed container through a pinhole. The radiation is reflected many times within the container and comes to thermal equilibrium with the wall. Radiation leaking out through the pinhole is characteristic of the radiation inside the container.\\
\includegraphics[max width=\textwidth, center]{2025_02_27_d4b95d9718f0b9e2e6b9g-271(2)}

Figure 7A. 2 The energy spectral density of radiation from a black body at several temperatures. Note that as the temperature increases, the maximum in the energy spectral density moves to shorter wavelengths and increases in intensity overall.\\
spectral density at the wavelength $\lambda$ and temperature $T$ simply means that there is a lot of energy associated with wavelengths lying between $\lambda$ and $\lambda+\mathrm{d} \lambda$ at that temperature. The energy density is obtained by summing (integrating) the energy spectral density over all wavelengths:


\begin{equation*}
\mathcal{E}(T)=\int_{0}^{\infty} \rho(\lambda, T) \mathrm{d} \lambda \tag{7A.2}
\end{equation*}


The units of $\mathcal{E}(T)$ are joules per metre cubed $\left(\mathrm{Jm}^{-3}\right)$, so the units of $\rho(\lambda, T)$ are $\mathrm{Jm}^{-4}$. Empirically, the energy density is found to vary as $T^{4}$, an observation expressed by the Stefan-Boltzmann law:

$$
\mathcal{E}(T)=\text { constant } \times T^{4}
$$

Stefan-Boltzmann law\\
(7A.3)\\
with the constant equal to $7.567 \times 10^{-16} \mathrm{Jm}^{-3} \mathrm{~K}^{-4}$.\\
The container in Fig. 7A. 1 emits radiation that can be thought of as oscillations of the electromagnetic field stimulated by the oscillations of electrical charges in the material of the wall. According to classical physics, every oscillator is excited to some extent, and according to the equipartition principle (The chemist's toolkit 7 in Topic 2A) every oscillator,\\
\includegraphics[max width=\textwidth, center]{2025_02_27_d4b95d9718f0b9e2e6b9g-271}

Figure 7A. 3 Comparison of the experimental energy spectral density with the prediction of the Rayleigh-Jeans law (eqn 7A.4). The latter predicts an infinite energy spectral density at short wavelengths and infinite overall energy density.\\
regardless of its frequency, has an average energy of $k T$. On this basis, the physicist Lord Rayleigh, with minor help from James Jeans, deduced what is now known as the Rayleigh-Jeans law:

$$
\rho(\lambda, T)=\frac{8 \pi k T}{\lambda^{4}}
$$

Rayleigh-Jeans law\\
(7A.4)\\
where $k$ is Boltzmann's constant ( $k=1.381 \times 10^{-23} \mathrm{~J} \mathrm{~K}^{-1}$ ).\\
The Rayleigh-Jeans law is not supported by the experimental measurements. As is shown in Fig. 7A.3, although there is agreement at long wavelengths, it predicts that the energy spectral density (and hence the intensity of the radiation emitted) increases without going through a maximum as the wavelength decreases. That is, the Rayleigh-Jeans law is inconsistent with Wien's law. Equation 7A. 4 also implies that the radiation is intense at very short wavelengths and becomes infinitely intense as the wavelength tends to zero. The concentration of radiation at short wavelengths is called the ultraviolet catastrophe, and is an unavoidable consequence of classical physics.

In 1900, Max Planck found that the experimentally observed intensity distribution of black-body radiation could be explained by proposing that the energy of each oscillator is limited to discrete values. In particular, Planck assumed that for an electromagnetic oscillator of frequency $v$, the permitted energies are integer multiples of $h v$ :


\begin{equation*}
E=n h v \quad n=0,1,2, \ldots \tag{7A.5}
\end{equation*}


In this expression $h$ is a fundamental constant now known as Planck's constant. The limitation of energies to discrete values is called energy quantization. On this basis Planck was able to derive an expression for the energy spectral density which is now called the Planck distribution:

$$
\rho(\lambda, T)=\frac{8 \pi h c}{\lambda^{5}\left(\mathrm{e}^{h c / \lambda k T}-1\right)}
$$

Planck distribution\\
(7A.6a)

This expression is plotted in Fig. 7A. 4 and fits the experimental data very well at all wavelengths. The value of $h$, which is an\\
\includegraphics[max width=\textwidth, center]{2025_02_27_d4b95d9718f0b9e2e6b9g-272}

Figure 7A. 4 The Planck distribution (eqn 7A.6a) accounts for the experimentally determined energy distribution of black-body radiation. It coincides with the Rayleigh-Jeans distribution at long wavelengths.\\
undetermined parameter in the theory, can be found by varying its value until the best fit is obtained between the eqn 7A.6a and experimental measurements. The currently accepted value is $h=6.626 \times 10^{-34} \mathrm{~J} \mathrm{~s}$.

For short wavelengths, $h c / \lambda k T \gg 1$, and because $\mathrm{e}^{h c / \lambda k T} \rightarrow \infty$ faster than $\lambda^{5} \rightarrow 0$ it follows that $\rho \rightarrow 0$ as $\lambda \rightarrow 0$. Hence, the energy spectral density approaches zero at short wavelengths, and so the Planck distribution avoids the ultraviolet catastrophe. For long wavelengths in the sense $h c / \lambda k T \ll 1$, the denominator in the Planck distribution can be replaced by (see The chemist's toolkit 12 in Topic 5B)

$$
\mathrm{e}^{h c \mid \lambda k T}-1=\left(1+\frac{h c}{\lambda k T}+\cdots\right)-1 \approx \frac{h c}{\lambda k T}
$$

When this approximation is substituted into eqn 7A.6a, the Planck distribution reduces to the Rayleigh-Jeans law, eqn 7A.4. The wavelength at the maximum can be found by differentiation, and is given by $\lambda_{\max } T=$ constant, in accord with Wien's law; the value of the constant found in this way, $h c / 5 k$, agrees with the experimentally determined value. Finally, the total energy density is


\begin{equation*}
\mathcal{E}(T)=\int_{0}^{\infty} \frac{8 \pi h c}{\lambda^{5}\left(\mathrm{e}^{h c / \lambda k T}-1\right)} \mathrm{d} \lambda=a T^{4} \text { with } a=\frac{8 \pi^{5} k^{4}}{15(h c)^{3}} \tag{7A.7}
\end{equation*}


which is finite and agrees with the Stefan-Boltzmann law (eqn 7A.3), including predicting the value of its constant correctly.

\subsection*{Brief illustration 7A. 1}
Consider eqn 7A.6a with $\lambda_{1}=450 \mathrm{~nm}$ (blue light) and $\lambda_{2}=$ 700 nm (red light), and $T=298 \mathrm{~K}$. It follows that

$$
\begin{aligned}
& \frac{h c}{\lambda_{1} k T}=\frac{\left(6.626 \times 10^{-34} \mathrm{Js}\right) \times\left(2.998 \times 10^{8} \mathrm{~ms}^{-1}\right)}{\left(450 \times 10^{-9} \mathrm{~m}\right) \times\left(1.381 \times 10^{-23} \mathrm{JK}^{-1}\right) \times(298 \mathrm{~K})}=107.2 \ldots \\
& \frac{h c}{\lambda_{2} k T}=\frac{\left(6.626 \times 10^{-34} \mathrm{Js}\right) \times\left(2.998 \times 10^{8} \mathrm{~ms}^{-1}\right)}{\left(700 \times 10^{-9} \mathrm{~m}\right) \times\left(1.381 \times 10^{-23} \mathrm{JK}^{-1}\right) \times(298 \mathrm{~K})}=68.9 \ldots
\end{aligned}
$$

and

$$
\begin{aligned}
\frac{\rho(450 \mathrm{~nm}, 298 \mathrm{~K})}{\rho(700 \mathrm{~nm}, 298 \mathrm{~K})} & =\left(\frac{700 \times 10^{-9} \mathrm{~m}}{450 \times 10^{-9} \mathrm{~m}}\right)^{5} \times \frac{\mathrm{e}^{68.9 \cdots}-1}{\mathrm{e}^{107.2 \cdots}-1} \\
& =9.11 \times\left(2.30 \times 10^{-17}\right)=2.10 \times 10^{-16}
\end{aligned}
$$

At room temperature, the proportion of shorter wavelength radiation is insignificant.

There is a single reason why Planck's approach is successful but Rayleigh's is not. Instead of allowing each oscillator to have the same average energy, regardless of its frequency, Planck used the Boltzmann distribution (see the Prologue to this text) to argue that higher frequency oscillators, which generate shorter wavelength radiation, are less likely to be excited than lower frequency oscillators. Indeed, for very high frequencies the minimum excitation energy of $h v$ is too large for the oscillator to be excited at all. This elimination of the contribution from very high frequency oscillators avoids the ultraviolet catastrophe.

It is sometimes convenient to express the Planck distribution in terms of the frequency. Then $\rho(v, T) \mathrm{d} v$ is the energy density at temperature $T$ due to the presence of electromagnetic radiation with frequencies between $v$ and $v+\mathrm{d} v$, and

\[
\rho(v, T)=\frac{8 \pi h \nu^{3}}{c^{3}\left(\mathrm{e}^{h \nu / k T}-1\right)} \quad \begin{align*}
& \text { Planck distribution in }  \tag{7A.6b}\\
& \text { terms of frequency }
\end{align*}
\]

\subsection*{(b) Heat capacity}
When energy is supplied as heat to a substance its temperature rises; the heat capacity (Topic 2 A ) is the constant of proportionality between the energy supplied and the temperature rise $\left(C=\mathrm{d} q / \mathrm{d} T\right.$ and, at constant volume, $\left.C_{V, \mathrm{~m}}=\left(\partial U_{\mathrm{m}} / \partial T\right)_{V}\right)$. Experimental measurements made during the nineteenth century had shown that at room temperature the molar heat capacities of many monatomic solids are about $3 R$, where $R$ is the gas constant. ${ }^{1}$ However, when measurements were made at much lower temperatures it was found that the heat capacity decreased, tending to zero as the temperature approached zero.

Classical physics was unable to explain this temperature dependence. The classical picture of a solid is of atoms oscillating about fixed positions, with the expectation that each oscillating atom will have the same average energy $k T$. This model predicts that a solid consisting of $N$ atoms, each free to oscillate in three dimensions, will have energy $U=3 N k T$ and hence heat capacity $C_{V}=(\partial U / \partial T)_{V}=3 N k$. The molar heat capacity is therefore predicted to be $3 N_{\mathrm{A}} k$ which, recognizing that $N_{\mathrm{A}} k=R$, is equal to $3 R$ at all temperatures. In 1905, Einstein suggested applying Planck's hypothesis and supposing that each oscillating atom

\footnotetext{${ }^{1}$ The gas constant occurs in the context of solids because it is actually the more fundamental Boltzmann's constant in disguise: $R=N_{\mathrm{A}} k$.
}
could have an energy $n h v$, where $n$ is an integer and $v$ is the frequency of the oscillation. Einstein went on to show by using the Boltzmann distribution that each oscillator is unlikely to be excited to high energies and at low temperatures few oscillators can be excited at all. As a consequence, because the oscillators cannot be excited, the heat capacity falls to zero. The quantitative result that Einstein obtained (as shown in Topic 13E) is
$$
C_{V, \mathrm{~m}}(T)=3 R f_{\mathrm{E}}(T), \quad f_{\mathrm{E}}(T)=\left(\frac{\theta_{\mathrm{E}}}{T}\right)^{2}\left(\frac{\mathrm{e}^{\theta_{\mathrm{E}} / 2 T}}{\mathrm{e}_{\mathrm{E} / T}-1}\right)^{2}
$$

Einstein formula\\
(7A.8a)\\
In this expression $\theta_{\mathrm{E}}$ is the Einstein temperature, $\theta_{\mathrm{E}}=h v / k$.\\
At high temperatures (in the sense $T \gg \theta_{\mathrm{E}}$ ) the exponentials in $f_{\mathrm{E}}$ can be expanded as $\mathrm{e}^{x}=1+x+\cdots$ and higher terms ignored (The chemist's toolkit 12 in Topic 5B). The result is

$$
f_{\mathrm{E}}(T)=\left(\frac{\theta_{\mathrm{E}}}{T}\right)^{2}\left\{\frac{1+\theta_{\mathrm{E}} / 2 T+\cdots}{\left(1+\theta_{\mathrm{E}} / T+\cdots\right)-1}\right\}^{2} \approx\left(\frac{\theta_{\mathrm{E}}}{T}\right)^{2}\left\{\frac{1}{\theta_{\mathrm{E}} / T}\right\}^{2} \approx 1
$$

(7A.8b)\\
and the classical result $\left(C_{V, \mathrm{~m}}=3 R\right)$ is obtained. At low temperatures (in the sense $T \ll \theta_{\mathrm{E}}$ ), $\mathrm{e}^{\theta_{\mathrm{E}} / T} \gg 1$ and


\begin{equation*}
f_{\mathrm{E}}(T) \approx\left(\frac{\theta_{\mathrm{E}}}{T}\right)^{2}\left(\frac{\mathrm{e}^{\theta_{\mathrm{E}} / 2 T}}{\mathrm{e}^{\theta_{\mathrm{E}} / T}}\right)^{2}=\left(\frac{\theta_{\mathrm{E}}}{T}\right)^{2} \mathrm{e}^{-\theta_{\mathrm{E}} / T} \tag{7A.8c}
\end{equation*}


The strongly decaying exponential function goes to zero more rapidly than $1 / T^{2}$ goes to infinity; so $f_{\mathrm{E}} \rightarrow 0$ as $T \rightarrow 0$, and the heat capacity approaches zero, as found experimentally. The physical reason for this success is that as the temperature is lowered, less energy is available to excite the atomic oscillations. At high temperatures many oscillators are excited into high energy states leading to classical behaviour.

Figure 7A. 5 shows the temperature dependence of the heat capacity predicted by the Einstein formula and some experi-\\
\includegraphics[max width=\textwidth, center]{2025_02_27_d4b95d9718f0b9e2e6b9g-273}

Figure 7A. 5 Experimental low-temperature molar heat capacities (open circles) and the temperature dependence predicted on the basis of Einstein's theory (solid line). His equation (eqn 7A.8) accounts for the dependence fairly well, but is everywhere too low.\\
mental data; the value of the Einstein temperature is adjusted to obtain the best fit to the data. The general shape of the curve is satisfactory, but the numerical agreement is in fact quite poor. This discrepancy arises from Einstein's assumption that all the atoms oscillate with the same frequency. A more sophisticated treatment, due to Peter Debye, allows the oscillators to have a range of frequencies from zero up to a maximum. This approach results in much better agreement with the experimental data and there can be little doubt that mechanical motion as well as electromagnetic radiation is quantized.

\subsection*{(c) Atomic and molecular spectra}
The most compelling and direct evidence for the quantization of energy comes from spectroscopy, the detection and analysis of the electromagnetic radiation absorbed, emitted, or scattered by a substance. The record of the variation of the intensity of this radiation with frequency $(v)$, wavelength $(\lambda)$, or wavenumber ( $\tilde{v}=v / c$, see The chemist's toolkit 13) is called its spectrum (from the Latin word for appearance).

An atomic emission spectrum is shown in Fig. 7A.6, and a molecular absorption spectrum is shown in Fig. 7A.7. The obvious feature of both is that radiation is emitted or absorbed at a series of discrete frequencies. This observation can be understood if the energy of the atoms or molecules is also confined to discrete values, because then the energies that a molecule can discard or acquire are also confined to discrete values (Fig. 7A.8). If the energy of an atom or molecule decreases by $\Delta E$, and this energy is carried away as radiation, the frequency of the radiation $v$ and the change in energy are related by the Bohr frequency condition:

$$
\Delta E=h v
$$

Bohr frequency condition\\
(7A.9)\\
A molecule is said to undergo a spectroscopic transition, a change of state, and as a result an emission 'line', a sharply defined peak, appears in the spectrum at frequency $v$.\\
\includegraphics[max width=\textwidth, center]{2025_02_27_d4b95d9718f0b9e2e6b9g-273(1)}

Figure 7A. 6 A region of the spectrum of radiation emitted by excited iron atoms consists of radiation at a series of discrete wavelengths (or frequencies).\\
\includegraphics[max width=\textwidth, center]{2025_02_27_d4b95d9718f0b9e2e6b9g-274(1)}

Figure 7A. 7 A molecule can change its state by absorbing radiation at definite frequencies. This spectrum is due to the electronic, vibrational, and rotational excitation of sulfur dioxide $\left(\mathrm{SO}_{2}\right)$ molecules. The observation of discrete spectral lines suggests that molecules can possess only discrete energies, not an arbitrary energy.\\
\includegraphics[max width=\textwidth, center]{2025_02_27_d4b95d9718f0b9e2e6b9g-274}

Figure 7A. 8 Spectroscopic transitions, such as those shown in Fig. 7A.6, can be accounted for by supposing that an atom (or molecule) emits electromagnetic radiation as it changes from a discrete level of high energy to a discrete level of lower energy. High-frequency radiation is emitted when the energy change is large. Transitions like those shown in Fig. 7A. 7 can be explained by supposing that a molecule (or atom) absorbs radiation as it changes from a low-energy level to a higher-energy level.

\subsection*{Brief illustration 7A. 2}
Atomic sodium produces a yellow glow (as in some street lamps) resulting from the emission of radiation of 590 nm . The spectroscopic transition responsible for the emission involves electronic energy levels that have a separation given by eqn 7A.9:

$$
\begin{aligned}
\Delta E & =h v=\frac{h c}{\lambda}=\frac{\left(6.626 \times 10^{-34} \mathrm{Js}\right) \times\left(2.998 \times 10^{8} \mathrm{~ms}^{-1}\right)}{590 \times 10^{-9} \mathrm{~m}} \\
& =3.37 \times 10^{-19} \mathrm{~J}
\end{aligned}
$$

This energy difference can be expressed in a variety of ways. For instance, multiplication by Avogadro's constant results in\\
an energy separation per mole of atoms, of $203 \mathrm{~kJ} \mathrm{~mol}^{-1}$, comparable to the energy of a weak chemical bond.

\subsection*{7A. 2 Wave-particle duality}
The experiments about to be described show that electromagnetic radiation-which classical physics treats as wave-likeactually also displays the characteristics of particles. Another experiment shows that electrons-which classical physics treats as particles-also display the characteristics of waves. This waveparticle duality, the blending together of the characteristics of waves and particles, lies at the heart of quantum mechanics.

\subsection*{(a) The particle character of electromagnetic radiation}
The Planck treatment of black-body radiation introduced the idea that an oscillator of frequency $v$ can have only the energies $0, h v, 2 h v, \ldots$. This quantization leads to the suggestion (and at this stage it is only a suggestion) that the resulting electromagnetic radiation of that frequency can be thought of as consisting of $0,1,2, \ldots$ particles, each particle having an energy $h v$. These particles of electromagnetic radiation are now called photons. Thus, if an oscillator of frequency $v$ is excited to its first excited state, then one photon of that frequency is present, if it is excited to its second excited state, then two photons are present, and so on. The observation of discrete spectra from atoms and molecules can be pictured as the atom or molecule generating a photon of energy $h \nu$ when it discards an energy of magnitude $\Delta E$, with $\Delta E=h v$.

\subsection*{Example 7A. 1 Calculating the number of photons}
Calculate the number of photons emitted by a 100 W yellow lamp in 1.0 s . Take the wavelength of yellow light as 560 nm , and assume 100 per cent efficiency.

Collect your thoughts Each photon has an energy $h v$, so the total number $N$ of photons needed to produce an energy $E$ is $N=E / h v$. To use this equation, you need to know the frequency of the radiation (from $\nu=c / \lambda$ ) and the total energy emitted by the lamp. The latter is given by the product of the power ( $P$, in watts) and the time interval, $\Delta t$, for which the lamp is turned on: $E=P \Delta t$ (see The chemist's toolkit 8 in Topic 2A).

The solution The number of photons is

$$
N=\frac{E}{h \nu}=\frac{P \Delta t}{h(c / \lambda)}=\frac{\lambda P \Delta t}{h c}
$$

Substitution of the data gives

$$
N=\frac{\left(5.60 \times 10^{-7} \mathrm{~m}\right) \times\left(100 \mathrm{~J} \mathrm{~s}^{-1}\right) \times(1.0 \mathrm{~s})}{\left(6.626 \times 10^{-34} \mathrm{Js}\right) \times\left(2.998 \times 10^{8} \mathrm{~m} \mathrm{~s}^{-1}\right)}=2.8 \times 10^{20}
$$

A note on good practice To avoid rounding and other numerical errors, it is best to carry out algebraic calculations first, and to substitute numerical values into a single, final formula. Moreover, an analytical result may be used for other data without having to repeat the entire calculation.

Self-test 7A. 1 How many photons does a monochromatic (single frequency) infrared rangefinder of power 1 mW and wavelength 1000 nm emit in 0.1 s ?\\
\includegraphics[max width=\textwidth, center]{2025_02_27_d4b95d9718f0b9e2e6b9g-275}

So far, the existence of photons is only a suggestion. Experimental evidence for their existence comes from the measurement of the energies of electrons produced in the photoelectric effect, the ejection of electrons from metals when they are exposed to ultraviolet radiation. The experimental characteristics of the photoelectric effect are as follows:

\begin{itemize}
  \item No electrons are ejected, regardless of the intensity of the radiation, unless its frequency exceeds a threshold value characteristic of the metal.
  \item The kinetic energy of the ejected electrons increases linearly with the frequency of the incident radiation but is independent of the intensity of the radiation.
  \item Even at low radiation intensities, electrons are ejected immediately if the frequency is above the threshold value.
\end{itemize}

Figure 7A. 9 illustrates the first and second characteristics.\\
These observations strongly suggest that in the photoelectric effect a particle-like projectile collides with the metal and, if the kinetic energy of the projectile is high enough, an electron is ejected. If the projectile is a photon of energy $h v$ ( $v$ is the frequency of the radiation), the kinetic energy of the\\
\includegraphics[max width=\textwidth, center]{2025_02_27_d4b95d9718f0b9e2e6b9g-275(2)}

Figure 7A. 9 In the photoelectric effect, it is found that no electrons are ejected when the incident radiation has a frequency below a certain value that is characteristic of the metal. Above that value, the kinetic energy of the photoelectrons varies linearly with the frequency of the incident radiation.\\
electron is $E_{\mathrm{k}}$, and the energy needed to remove an electron from the metal, which is called its work function, is $\Phi$ (uppercase phi), then as illustrated in Fig. 7A.10, the conservation of energy implies that


\begin{equation*}
h v=E_{\mathrm{k}}+\Phi \quad \text { or } E_{\mathrm{k}}=h v-\Phi \quad \text { Photoelectric effect } \tag{7A.10}
\end{equation*}


This model explains the three experimental observations:

\begin{itemize}
  \item Photoejection cannot occur if $h v<\Phi$ because the photon brings insufficient energy.
  \item The kinetic energy of an ejected electron increases linearly with the frequency of the photon.
  \item When a photon collides with an electron, it gives up all its energy, so electrons should appear as soon as the collisions begin, provided the photons have sufficient energy.
\end{itemize}

A practical application of eqn 7A. 10 is that it provides a technique for the determination of Planck's constant, because the slopes of the lines in Fig. 7A. 9 are all equal to $h$.

The energies of photoelectrons, the work function, and other quantities are often expressed in the alternative energy unit the electronvolt $(\mathrm{eV}): 1 \mathrm{eV}$ is defined as the kinetic energy acquired when an electron (of charge $-e$ ) is accelerated from rest through a potential difference $\Delta \phi=1 \mathrm{~V}$. That kinetic energy is $e \Delta \phi$, so

$$
E_{\mathrm{k}}=e \Delta \phi=\left(1.602 \times 10^{-19} \mathrm{C}\right) \times 1 \mathrm{~V}=1.602 \times 10^{-19} \mathrm{CV}=1 \mathrm{eV}
$$

Because $1 \mathrm{CV}=1 \mathrm{~J}$, it follows that the relation between electronvolts and joules is

$$
1 \mathrm{eV}=1.602 \times 10^{-19} \mathrm{~J}
$$

\begin{center}
\includegraphics[max width=\textwidth]{2025_02_27_d4b95d9718f0b9e2e6b9g-275(1)}
\end{center}

Figure 7A. 10 The photoelectric effect can be explained if it is supposed that the incident radiation is composed of photons that have energy proportional to the frequency of the radiation. (a) The energy of the photon is insufficient to drive an electron out of the metal. (b) The energy of the photon is more than enough to eject an electron, and the excess energy is carried away as the kinetic energy of the photoelectron (the ejected electron).

\subsection*{Example 7A. 2 Calculating the longest wavelength capable of photoejection}
A photon of radiation of wavelength 305 nm ejects an electron with a kinetic energy of 1.77 eV from a metal. Calculate the longest wavelength of radiation capable of ejecting an electron from the metal.

Collect your thoughts You can use eqn 7A.10, rearranged into $\Phi=h v-E_{\mathrm{k}}$, to compute the work function because you know the frequency of the photon from $v=c / \lambda$. The threshold for photoejection is the lowest frequency at which electron ejection occurs without there being any excess energy; that is, the kinetic energy of the ejected electron is zero. Setting $E_{\mathrm{k}}=0$ in $E_{\mathrm{k}}=h v-\Phi$ gives the minimum photon frequency as $\nu_{\text {min }}=$ $\Phi / h$. Use this value of the frequency to calculate the corresponding wavelength, $\lambda_{\text {max }}$.

The solution The minimum frequency for photoejection is

$$
v_{\min }=\frac{\Phi}{h}=\frac{h v-E_{\mathrm{k}}}{h}=\frac{c}{\lambda}-\frac{E_{\mathrm{k}}}{h}
$$

The longest wavelength that can cause photoejection is therefore

$$
\lambda_{\max }=\frac{c}{v_{\min }}=\frac{c}{c / \lambda-E_{\mathrm{k}} / h}=\frac{1}{1 / \lambda-E_{\mathrm{k}} / h c}
$$

Now substitute the data. The kinetic energy of the electron is

$$
E_{\mathrm{k}}=1.77 \mathrm{eV} \times\left(1.602 \times 10^{-19} \mathrm{JeV}^{-1}\right)=2.83 \ldots \times 10^{-19} \mathrm{~J}
$$

SO

$$
\frac{E_{\mathrm{k}}}{h c}=\frac{2.83 \ldots \times 10^{-19} \mathrm{~J}}{\left(6.626 \times 10^{-34} \mathrm{Js}\right) \times\left(2.998 \times 10^{8} \mathrm{~ms}^{-1}\right)}=1.42 \ldots \times 10^{6} \mathrm{~m}^{-1}
$$

Therefore, with $1 / \lambda=1 / 305 \mathrm{~nm}=3.27 \ldots \times 10^{6} \mathrm{~m}^{-1}$,

$$
\lambda_{\max }=\frac{1}{\left(3.27 \ldots \times 10^{6} \mathrm{~m}^{-1}\right)-\left(1.42 \ldots \times 10^{6} \mathrm{~m}^{-1}\right)}=5.40 \times 10^{-7} \mathrm{~m}
$$

or 540 nm .\\
Self-test 7A. 2 When ultraviolet radiation of wavelength 165 nm strikes a certain metal surface, electrons are ejected with a speed of $1.24 \mathrm{Mm} \mathrm{s}^{-1}$. Calculate the speed of electrons ejected by radiation of wavelength 265 nm .

$$
{ }_{\mathrm{I}-} \mathrm{S} \text { ury sé }: \iota \partial \mathrm{Ms} u_{V}
$$

\subsection*{(b) The wave character of particles}
Although contrary to the long-established wave theory of radiation, the view that radiation consists of particles had been held before, but discarded. No significant scientist, however, had taken the view that matter is wave-like. Nevertheless,\\
\includegraphics[max width=\textwidth, center]{2025_02_27_d4b95d9718f0b9e2e6b9g-276}

Figure 7A. 11 The Davisson-Germer experiment. The scattering of an electron beam from a nickel crystal shows a variation in intensity characteristic of a diffraction experiment in which waves interfere constructively and destructively in different directions.\\
experiments carried out in 1925 forced people to consider that possibility. The crucial experiment was performed by Clinton Davisson and Lester Germer, who observed the diffraction of electrons by a crystal (Fig. 7A.11). As remarked in The chemist's toolkit 13, diffraction is the interference caused by an object in the path of waves. Davisson and Germer's success was a lucky accident, because a chance rise of temperature caused their polycrystalline sample to anneal, and the ordered planes of atoms then acted as a diffraction grating. The DavissonGermer experiment, which has since been repeated with other particles (including $\alpha$ particles, molecular hydrogen, and neutrons), shows clearly that particles have wave-like properties. At almost the same time, G.P. Thomson showed that a beam of electrons was diffracted when passed through a thin gold foil.

Some progress towards accounting for wave-particle duality had already been made by Louis de Broglie who, in 1924, suggested that any particle, not only photons, travelling with a linear momentum $p=m v$ (with $m$ the mass and $v$ the speed of the particle) should have in some sense a wavelength given by what is now called the de Broglie relation:


\begin{equation*}
\lambda=\frac{h}{p} \tag{7A.11}
\end{equation*}


de Broglie relation\\
That is, a particle with a high linear momentum has a short wavelength. Macroscopic bodies have such high momenta even when they are moving slowly (because their mass is so great), that their wavelengths are undetectably small, and the wave-like properties cannot be observed. This undetectability is why classical mechanics can be used to explain the behaviour of macroscopic bodies. It is necessary to invoke quantum mechanics only for microscopic bodies, such as atoms and molecules, in which masses are small.

\subsection*{Example 7A. 3 Estimating the de Broglie wavelength}
Estimate the wavelength of electrons that have been accelerated from rest through a potential difference of 40 kV .

Collect your thoughts To use the de Broglie relation, you need to know the linear momentum, $p$, of the electrons. To calculate the linear momentum, note that the energy acquired by an electron accelerated through a potential difference $\Delta \phi$ is $e \Delta \phi$, where $e$ is the magnitude of its charge. At the end of the period of acceleration, all the acquired energy is in the form of kinetic energy, $E_{\mathrm{k}}=\frac{1}{2} m_{\mathrm{e}} v^{2}=p^{2} / 2 m_{\mathrm{e}}$. You can therefore calculate $p$ by setting $p^{2} / 2 m_{\mathrm{e}}$ equal to $e \Delta \phi$. For the manipulation of units use $1 \mathrm{VC}=1 \mathrm{~J}$ and $1 \mathrm{~J}=1 \mathrm{~kg} \mathrm{~m}^{2} \mathrm{~s}^{-2}$.\\
The solution The expression $p^{2} / 2 m_{\mathrm{e}}=e \Delta \phi$ implies that $p=$ $\left(2 m_{\mathrm{e}} e \Delta \phi\right)^{1 / 2}$ then, from the de Broglie relation $\lambda=h / p$,

$$
\lambda=\frac{h}{\left(2 m_{e} e \Delta \phi\right)^{1 / 2}}
$$

Substitution of the data and the fundamental constants gives

$$
\begin{aligned}
\lambda & =\frac{6.626 \times 10^{-34} \mathrm{Js}}{\left\{2 \times\left(9.109 \times 10^{-31} \mathrm{~kg}\right) \times\left(1.602 \times 10^{-19} \mathrm{C}\right) \times\left(4.0 \times 10^{4} \mathrm{~V}\right)\right\}^{1 / 2}} \\
& =6.1 \times 10^{-12} \mathrm{~m}
\end{aligned}
$$

or 6.1 pm .\\
Comment. Electrons accelerated in this way are used in the technique of electron diffraction for imaging biological systems and for the determination of the structures of solid surfaces (Topic 19A).

Self-test 7A. 3 Calculate the wavelength of (a) a neutron with a translational kinetic energy equal to $k T$ at 300 K , (b) a tennis ball of mass 57 g travelling at $80 \mathrm{~km} \mathrm{~h}^{-1}$.

\subsection*{Checklist of concepts}
$\square$ 1. A black body is an object capable of emitting and absorbing all wavelengths of radiation without favouring any wavelength.2. An electromagnetic field of a given frequency can take up energy only in discrete amounts.3. Atomic and molecular spectra show that atoms and molecules can take up energy only in discrete amounts.\\
4. The photoelectric effect establishes the view that electromagnetic radiation, regarded in classical physics as wave-like, consists of particles (photons).5. The diffraction of electrons establishes the view that electrons, regarded in classical physics as particles, are wavelike with a wavelength given by the de Broglie relation.6. Wave-particle duality is the recognition that the concepts of particle and wave blend together.

\subsection*{Checklist of equations}
\begin{center}
\begin{tabular}{|c|c|c|c|}
\hline
Property & Equation & Comment & Equation number \\
\hline
Wien's law & $\lambda_{\text {max }} T=2.9 \times 10^{-3} \mathrm{mK}$ &  & 7A. 1 \\
\hline
Stefan-Boltzmann law & $\mathcal{E}(T)=$ constant $\times T^{4}$ &  & 7A. 3 \\
\hline
Planck distribution & $\rho(\lambda, T)=8 \pi h c /\left\{\lambda^{5}\left(\mathrm{e}^{h c / \lambda k T}-1\right)\right\}$ & Black-body radiation & 7A. 6 \\
\hline
 & \( \rho(v, T)=8 \pi h v^{3} /\left\{c^{3}\left(\mathrm{e}^{h \nu / k T}-1\right)\right\} \) &  &  \\
\hline
Einstein formula for heat capacity of a solid & \( \begin{aligned} & C_{V, \mathrm{~m}}(T)=3 R f_{\mathrm{E}}(T) \\ & f_{\mathrm{E}}(T)=\left(\theta_{\mathrm{E}} / T\right)^{2}\left\{\mathrm{e}^{\theta_{\mathrm{E}} / 2 T} /\left(\mathrm{e}^{\theta_{\mathrm{E}} / T}-1\right)\right\}^{2} \end{aligned} \) & Einstein temperature: \( \theta_{\mathrm{E}}=h v / k \) & 7A. 8 \\
\hline
Bohr frequency condition & $\Delta E=h \nu$ &  & 7A. 9 \\
\hline
Photoelectric effect & $E_{\mathrm{k}}=h v-\Phi$ & $\Phi$ is the work function & 7A. 10 \\
\hline
de Broglie relation & $\lambda=h / p$ & $\lambda$ is the wavelength of a particle of linear momentum $p$ & 7A. 11 \\
\hline
\end{tabular}
\end{center}

\section*{TOPIC 7B Wavefunctions}
\subsection*{Why do you need to know this material?}
Wavefunctions provide the essential foundation for understanding the properties of electrons in atoms and molecules, and are central to explanations in chemistry.

\subsection*{What is the key idea?}
All the dynamical properties of a system are contained in its wavefunction, which is obtained by solving the Schrödinger equation.

\subsection*{What do you need to know already?}
You need to be aware of the shortcomings of classical physics that drove the development of quantum theory (Topic 7A).

In classical mechanics an object travels along a definite path or trajectory. In quantum mechanics a particle in a particular state is described by a wavefunction, $\psi(\mathrm{psi})$, which is spread out in space, rather than being localized. The wavefunction contains all the dynamical information about the object in that state, such as its position and momentum.

\subsection*{7B. 1 The Schrödinger equation}
In 1926 Erwin Schrödinger proposed an equation for finding the wavefunctions of any system. The time-independent Schrödinger equation for a particle of mass $m$ moving in one dimension with energy $E$ in a system that does not change with time (for instance, its volume remains constant) is

\[
-\frac{\hbar^{2}}{2 m} \frac{\mathrm{~d}^{2} \psi}{\mathrm{~d} x^{2}}+V(x) \psi=E \psi \quad \begin{align*}
& \text { Time-independent }  \tag{7B.1}\\
& \text { Schrödinger equation }
\end{align*}
\]

The constant $\hbar=h / 2 \pi$ (which is read $h$-cross or $h$-bar) is a convenient modification of Planck's constant used widely in quantum mechanics; $V(x)$ is the potential energy of the particle at $x$. Because the total energy $E$ is the sum of potential and kinetic energies, the first term on the left must be related (in a manner explored later) to the kinetic energy of the particle. The Schrödinger equation can be regarded as a fundamental postulate of quantum mechanics, but its plausibility can be\\
demonstrated by showing that, for the case of a free particle, it is consistent with the de Broglie relation (Topic 7A).

\subsection*{How is that done? 7B. 1 Showing that the Schrödinger equation is consistent with the de Broglie relation}
The potential energy of a freely moving particle is zero everywhere, $V(x)=0$, so the Schrödinger equation (eqn 7B.1) becomes

$$
\frac{\mathrm{d}^{2} \psi}{\mathrm{~d} x^{2}}=-\frac{2 m E}{\hbar^{2}} \psi
$$

Step 1 Find a solution of the Schrödinger equation for a free particle\\
A solution of this equation is $\psi=\cos k x$, as you can confirm by noting that

$$
\frac{\mathrm{d}^{2} \psi}{\mathrm{~d} x^{2}}=\frac{\mathrm{d}^{2} \cos k x}{\mathrm{~d} x^{2}}=-k^{2} \cos k x=-k^{2} \psi
$$

It follows that $-k^{2}=-2 m E / \hbar^{2}$ and hence

$$
k=\left(\frac{2 m E}{\hbar^{2}}\right)^{1 / 2}
$$

The energy, which is only kinetic in this instance, is related to the linear momentum of the particle by $E=p^{2} / 2 m$ (The chemist's toolkit 6 in Topic 2A), so it follows that

$$
k=\left(\frac{2 m\left(p^{2} / 2 m\right)}{\hbar^{2}}\right)^{1 / 2}=\frac{p}{\hbar}
$$

The linear momentum is therefore related to $k$ by $p=k \hbar$.\\
Step 2 Interpret the wavefunction in terms of a wavelength\\
Now recognize that a wave (more specifically, a 'harmonic wave') can be described mathematically by a sine or cosine function. It follows that $\cos k x$ can be regarded as a wave that goes through a complete cycle as $k x$ increases by $2 \pi$. The wavelength is therefore given by $k \lambda=2 \pi$, so $k=2 \pi / \lambda$. Therefore, the linear momentum is related to the wavelength of the wavefunction by

$$
p=k \hbar=\frac{2 \pi}{\lambda} \times \frac{h}{2 \pi}=\frac{h}{\lambda}
$$

which is the de Broglie relation. The Schrödinger equation therefore has solutions consistent with the de Broglie relation.

\subsection*{7B. 2 The Born interpretation}
One piece of dynamical information contained in the wavefunction is the location of the particle. Max Born used an analogy with the wave theory of radiation, in which the square of the amplitude of an electromagnetic wave in a region is interpreted as its intensity and therefore (in quantum terms) as a measure of the probability of finding a photon present in the region. The Born interpretation of the wavefunction is:

If the wavefunction of a particle has the value $\psi$ at $x$, then the probability of finding the particle between $x$ and $x+\mathrm{d} x$ is proportional to $|\psi|^{2} \mathrm{~d} x$ (Fig. 7B.1).\\
The quantity $|\psi|^{2}=\psi^{\star} \psi$ allows for the possibility that $\psi$ is complex (see The chemist's toolkit 14). If the wavefunction is real (such as $\cos k x$ ), then $|\psi|^{2}=\psi^{2}$.

Because $|\psi|^{2} \mathrm{~d} x$ is a (dimensionless) probability, $|\psi|^{2}$ is the probability density, with the dimensions of $1 /$ length (for a onedimensional system). The wavefunction $\psi$ itself is called the probability amplitude. For a particle free to move in three dimensions (for example, an electron near a nucleus in an atom), the wavefunction depends on the coordinates $x, y$, and $z$ and is denoted $\psi(r)$. In this case the Born interpretation is (Fig. 7B.2):

If the wavefunction of a particle has the value $\psi$ at $r$, then the probability of finding the particle in an infinitesimal volume $\mathrm{d} \tau=\mathrm{d} x \mathrm{~d} y \mathrm{~d} z$ at that position is proportional to $|\psi|^{2} \mathrm{~d} \tau$.

\subsection*{The chemist's toolkit 14 Complex numbers}
Complex numbers have the general form

$$
z=x+\mathrm{i} y
$$

where $\mathrm{i}=\sqrt{-1}$. The real number $x$ is the 'real part of $z$ ', denoted $\operatorname{Re}(z)$; likewise, the real number $y$ is 'the imaginary part of $z$ ', $\operatorname{Im}(z)$. The complex conjugate of $z$, denoted $z^{\star}$, is formed by replacing i by -i :

$$
z^{*}=x-\mathrm{i} y
$$

The product of $z^{*}$ and $z$ is denoted $|z|^{2}$ and is called the square modulus of $z$. From the definition of $z$ and $z^{*}$ and $\mathrm{i}^{2}=$ -1 it follows that

$$
|z|^{2}=z^{\star} z=(x+\mathrm{i} y)(x-\mathrm{i} y)=x^{2}+y^{2}
$$

The square modulus is a real, non-negative number. The absolute value or modulus is denoted $|z|$ and is given by:

$$
|z|=\left(z^{\star} z\right)^{1 / 2}=\left(x^{2}+y^{2}\right)^{1 / 2}
$$

For further information about complex numbers, see The chemist's toolkit 16 in Topic 7C.\\
\includegraphics[max width=\textwidth, center]{2025_02_27_d4b95d9718f0b9e2e6b9g-279(1)}

Figure 7B. 1 The wavefunction $\psi$ is a probability amplitude in the sense that its square modulus ( $\psi^{*} \psi$ or $|\psi|^{2}$ ) is a probability density. The probability of finding a particle in the region between $x$ and $x+\mathrm{d} x$ is proportional to $|\psi|^{2} \mathrm{~d} x$. Here, the probability density is represented by the density of shading in the superimposed band.

In this case, $|\psi|^{2}$ has the dimensions of $1 /$ length $^{3}$ and the wavefunction itself has dimensions of $1 /$ length ${ }^{3 / 2}$ (and units such as $\mathrm{m}^{-3 / 2}$ ).

The Born interpretation does away with any worry about the significance of a negative (and, in general, complex) value of $\psi$ because $|\psi|^{2}$ is always real and nowhere negative. There is no direct significance in the negative (or complex) value of a wavefunction: only the square modulus is directly physically significant, and both negative and positive regions of a wavefunction may correspond to a high probability of finding a particle in a region (Fig. 7B.3). However, the presence of positive and negative regions of a wavefunction is of great indirect significance, because it gives rise to the possibility of constructive and destructive interference between different wavefunctions.

A wavefunction may be zero at one or more points, and at these locations the probability density is also zero. It is important to distinguish a point at which the wavefunction is zero (for instance, far from the nucleus of a hydrogen atom) from the point at which it passes through zero. The latter is called a node. A location where the wavefunction approaches zero without actually passing through zero is not a node. Thus, the\\
\includegraphics[max width=\textwidth, center]{2025_02_27_d4b95d9718f0b9e2e6b9g-279}

Figure 7B. 2 The Born interpretation of the wavefunction in threedimensional space implies that the probability of finding the particle in the volume element $\mathrm{d} \tau=\mathrm{d} x \mathrm{~d} y \mathrm{~d} z$ at some position $r$ is proportional to the product of $d \tau$ and the value of $|\psi|^{2}$ at that position.\\
\includegraphics[max width=\textwidth, center]{2025_02_27_d4b95d9718f0b9e2e6b9g-280(1)}

Figure 7B. 3 The sign of a wavefunction has no direct physical significance: the positive and negative regions of this wavefunction both correspond to the same probability distribution (as given by the square modulus of $\psi$ and depicted by the density of the shading).\\
wavefunction $\cos k x$ has nodes wherever $k x$ is an odd integral multiple of $\frac{1}{2} \pi$ (where the wave passes through zero), but the wavefunction $\mathrm{e}^{-k x}$ has no nodes, despite becoming zero as $x \rightarrow \infty$.

\subsection*{Example 7B. 1}
\subsection*{Interpreting a wavefunction}
The wavefunction of an electron in the lowest energy state of a hydrogen atom is proportional to $\mathrm{e}^{-r / a_{0}}$, where $a_{0}$ is a constant and $r$ the distance from the nucleus. Calculate the relative probabilities of finding the electron inside a region of volume $\delta V=1.0 \mathrm{pm}^{3}$, which is small even on the scale of the atom, located at (a) the nucleus, (b) a distance $a_{0}$ from the nucleus.

Collect your thoughts The region of interest is so small on the scale of the atom that you can ignore the variation of $\psi$ within it and write the probability, $P$, as proportional to the probability density ( $\psi^{2}$; note that $\psi$ is real) evaluated at the point of interest multiplied by the volume of interest, $\delta V$. That is, $P$ $\propto \psi^{2} \delta V$, with $\psi^{2} \propto \mathrm{e}^{-2 r / a_{0}}$.

The solution In each case $\delta V=1.0 \mathrm{pm}^{3}$. (a) At the nucleus, $r=0$, so

$$
P \propto \mathrm{e}^{0} \times\left(1.0 \mathrm{pm}^{3}\right)=1 \times\left(1.0 \mathrm{pm}^{3}\right)=1.0 \mathrm{pm}^{3}
$$

(b) At a distance $r=a_{0}$ in an arbitrary direction,

$$
P \propto \mathrm{e}^{-2} \times\left(1.0 \mathrm{pm}^{3}\right)=0.14 \ldots \times\left(1.0 \mathrm{pm}^{3}\right)=0.14 \mathrm{pm}^{3}
$$

Therefore, the ratio of probabilities is $1.0 / 0.14=7.1$.\\
Comment. Note that it is more probable (by a factor of 7) that the electron will be found at the nucleus than in a volume element of the same size located at a distance $a_{0}$ from the nucleus. The negatively charged electron is attracted to the positively charged nucleus, and is likely to be found close to it.

Self-test 7B. 1 The wavefunction for the electron in its lowest energy state in the ion $\mathrm{He}^{+}$is proportional to $\mathrm{e}^{-2 r / a_{0}}$. Repeat the calculation for this ion and comment on the result.\\
\includegraphics[max width=\textwidth, center]{2025_02_27_d4b95d9718f0b9e2e6b9g-280}

\subsection*{(a) Normalization}
A mathematical feature of the Schrödinger equation is that if $\psi$ is a solution, then so is $N \psi$, where $N$ is any constant. This feature is confirmed by noting that because $\psi$ occurs in every term in eqn 7B.1, it can be replaced by $N \psi$ and the constant factor $N$ cancelled to recover the original equation. This freedom to multiply the wavefunction by a constant factor means that it is always possible to find a normalization constant, $N$, such that rather than the probability density being proportional to $|\psi|^{2}$ it becomes equal to $|\psi|^{2}$.

A normalization constant is found by noting that, for a normalized wavefunction $N \psi$, the probability that a particle is in the region $\mathrm{d} x$ is equal to $\left(N \psi^{*}\right)(N \psi) \mathrm{d} x$ ( $N$ is taken to be real). Furthermore, the sum over all space of these individual probabilities must be 1 (the probability of the particle being somewhere is 1). Expressed mathematically, the latter requirement is


\begin{equation*}
N^{2} \int_{-\infty}^{\infty} \psi^{\star} \psi \mathrm{d} x=1 \tag{7B.2}
\end{equation*}


and therefore


\begin{equation*}
N=\frac{1}{\left(\int_{-\infty}^{\infty} \psi^{\star} \psi \mathrm{d} x\right)^{1 / 2}} \tag{7B.3}
\end{equation*}


Provided this integral has a finite value (that is, the wavefunction is 'square integrable'), the normalization factor can be found and the wavefunction 'normalized' (and specifically 'normalized to 1 '). From now on, unless stated otherwise, all wavefunctions are assumed to have been normalized to 1 , in which case in one dimension


\begin{equation*}
\int_{-\infty}^{\infty} \psi^{\star} \psi \mathrm{d} x=1 \tag{7B.4a}
\end{equation*}


and in three dimensions


\begin{equation*}
\int_{-\infty}^{\infty} \int_{-\infty}^{\infty} \int_{-\infty}^{\infty} \psi^{\star} \psi \mathrm{d} x \mathrm{~d} y \mathrm{~d} z=1 \tag{7B.4b}
\end{equation*}


In quantum mechanics it is common to write all such integrals in a short-hand form as


\begin{equation*}
\int \psi^{\star} \psi \mathrm{d} \tau=1 \tag{7B.4c}
\end{equation*}


where $\mathrm{d} \tau$ is the appropriate volume element and the integration is understood as being over all space.

\subsection*{Example 7B. 2 Normalizing a wavefunction}
Carbon nanotubes are thin hollow cylinders of carbon with diameters between 1 nm and 2 nm , and lengths of several micrometres. According to one simple model, the lowestenergy electrons of the nanotube are described by the wavefunction $\sin (\pi x / L)$, where $L$ is the length of the nanotube. Find the normalized wavefunction.

Collect your thoughts Because the wavefunction is onedimensional, you need to find the factor $N$ that guarantees\\
that the integral in eqn 7B.4a is equal to 1 . The wavefunction is real, so $\psi^{*}=\psi$. Relevant integrals are found in the Resource section.

The solution Write the wavefunction as $\psi=N \sin (\pi x / L)$, where $N$ is the normalization factor. The limits of integration are $x=0$ to $x=L$ because the wavefunction spans the length of the tube. It follows that

$$
\int \psi^{\star} \psi \mathrm{d} \tau=N^{2} \overbrace{\int_{0}^{L} \sin ^{2} \frac{\pi x}{L} \mathrm{~d} x}^{\text {Integral T. } 2}
$$

For the wavefunction to be normalized, this integral must be equal to 1 ; that is, $\frac{1}{2} N^{2} L=1$, and hence

$$
N=\left(\frac{2}{L}\right)^{1 / 2}
$$

The normalized wavefunction is therefore

$$
\psi=\left(\frac{2}{L}\right)^{1 / 2} \sin \frac{\pi x}{L}
$$

Comment. Because $L$ is a length, the dimensions of $\psi$ are $1 /$ length $^{1 / 2}$, and therefore those of $\psi^{2}$ are $1 /$ length, as is appropriate for a probability density in one dimension.

Self-test 7B. 2 The wavefunction for the next higher energy level for the electrons in the same tube is $\sin (2 \pi x / L)$. Normalize this wavefunction.

$$
z_{z / I}(T / Z)=N: \iota_{2} s u_{V}
$$

To calculate the probability of finding the system in a finite region of space the probability density is summed (integrated) over the region of interest. Thus, for a one-dimensional system, the probability $P$ of finding the particle between $x_{1}$ and $x_{2}$ is given by


\begin{equation*}
P=\int_{x_{1}}^{x_{2}}|\psi(x)|^{2} \mathrm{~d} x \tag{7B.5}
\end{equation*}


[one-dimensional region]

\subsection*{Example 7B. 3 \\
 Determining a probability}
As seen in Example 7B.2, the lowest-energy electrons of a carbon nanotube of length $L$ can be described by the normalized wavefunction $(2 / L)^{1 / 2} \sin (\pi x / L)$. What is the probability of finding the electron between $x=L / 4$ and $x=L / 2$ ?

Collect your thoughts Use eqn 7B. 5 and the normalized wavefunction to write an expression for the probability of finding the electron in the region of interest. Relevant integrals are given in the Resource section.

The solution From eqn 7B. 5 the probability is

$$
P=\frac{2}{L} \overbrace{\int_{L / 4}^{L / 2} \sin ^{2}(\pi x / L) \mathrm{d} x}^{\text {Integral } \mathrm{T} .2}
$$

It follows that

$$
P=\left.\frac{2}{L}\left(\frac{x}{2}-\frac{\sin (2 \pi x / L)}{4 \pi / L}\right)\right|_{L / 4} ^{L / 2}=\frac{2}{L}\left(\frac{L}{4}-\frac{L}{8}-0+\frac{L}{4 \pi}\right)=0.409
$$

Comment. There is a chance of about 41 per cent that the electron will be found in the region.

Self-test 7B. 3 As remarked in Self-test 7B.2, the normalized wavefunction of the next higher energy level of the electron in this model of the nanotube is $(2 / L)^{1 / 2} \sin (2 \pi x / L)$. What is the probability of finding the electron between $x=L / 4$ and $x=L / 2$ ?\\
\includegraphics[max width=\textwidth, center]{2025_02_27_d4b95d9718f0b9e2e6b9g-281}

\subsection*{(b) Constraints on the wavefunction}
The Born interpretation puts severe restrictions on the acceptability of wavefunctions. The first constraint is that $\psi$ must not be infinite over a finite region, because if it were, the Born interpretation would fail. This requirement rules out many possible solutions of the Schrödinger equation, because many mathematically acceptable solutions rise to infinity and are therefore physically unacceptable. The Born interpretation also rules out solutions of the Schrödinger equation that give rise to more than one value of $|\psi|^{2}$ at a single point because it would be absurd to have more than one value of the probability density for the particle at a point. This restriction is expressed by saying that the wavefunction must be single-valued; that is, it must have only one value at each point of space.

The Schrödinger equation itself also implies some mathematical restrictions on the type of functions that can occur. Because it is a second-order differential equation (in the sense that it depends on the second derivative of the wavefunction), $\mathrm{d}^{2} \psi / \mathrm{d} x^{2}$ must be well-defined if the equation is to be applicable everywhere. The second derivative is defined only if the first derivative is continuous: this means that (except as specified below) there can be no kinks in the function. In turn, the first derivative is defined only if the function is continuous: no sharp steps are permitted.

Overall, therefore, the constraints on the wavefunction, which are summarized in Fig. 7B.4, are that it

\begin{itemize}
  \item must not be infinite over a finite region;
  \item must be single-valued;
  \item must be continuous;
  \item must have a continuous first derivative (slope).
\end{itemize}

The last of these constraints does not apply if the potential energy has abrupt, infinitely high steps (as in the particle-in-abox model treated in Topic 7D).\\
\includegraphics[max width=\textwidth, center]{2025_02_27_d4b95d9718f0b9e2e6b9g-282}

Figure 7B. 4 The wavefunction must satisfy stringent conditions for it to be acceptable: (a) unacceptable because it is infinite over a finite region; (b) unacceptable because it is not single-valued; (c) unacceptable because it is not continuous; (d) unacceptable because its slope is discontinuous.

\subsection*{(c) Quantization}
The constraints just noted are so severe that acceptable solutions of the Schrödinger equation do not in general exist for arbitrary values of the energy $E$. In other words, a particle may possess only certain energies, for otherwise its wavefunction would be physically unacceptable. That is,

As a consequence of the restrictions on its wavefunction, the energy of a particle is quantized.

These acceptable energies are found by solving the Schrödinger equation for motion of various kinds, and selecting the solutions that conform to the restrictions listed above.

\subsection*{Checklist of concepts}
1. A wavefunction is a mathematical function that contains all the dynamical information about a system.2. The Schrödinger equation is a second-order differential equation used to calculate the wavefunction of a system.3. According to the Born interpretation, the probability density at a point is proportional to the square of the wavefunction at that point.4. A node is a point where a wavefunction passes through zero.\begin{enumerate}
  \setcounter{enumi}{4}
  \item A wavefunction is normalized if the integral over all space of its square modulus is equal to 1 .6. A wavefunction must be single-valued, continuous, not infinite over a finite region of space, and (except in special cases) have a continuous slope.7. The quantization of energy stems from the constraints that an acceptable wavefunction must satisfy.
\end{enumerate}

\subsection*{Checklist of equations}
\begin{center}
\begin{tabular}{|c|c|c|c|}
\hline
Property & Equation & Comment & Equation number \\
\hline
The time-independent Schrödinger equation & $-\left(\hbar^{2} / 2 m\right)\left(\mathrm{d}^{2} \psi / \mathrm{d} x^{2}\right)+V(x) \psi=E \psi$ & One-dimensional system* & 7B. 1 \\
\hline
Normalization & $\int \psi^{*} \psi \mathrm{~d} \tau=1$ & Integration over all space & 7B.4c \\
\hline
Probability of a particle being between $x_{1}$ and $x_{2}$ & \( P=\int_{x_{1}}^{x_{2}}|\psi(x)|^{2} \mathrm{~d} x \) & One-dimensional region & 7B. 5 \\
\hline
\end{tabular}
\end{center}

\footnotetext{\begin{itemize}
  \item Higher dimensions are treated in Topics 7D, 7F, and 8A.
\end{itemize}
}

\section*{TOPIC 7C Operators and observables}
\subsection*{Why do you need to know this material?}
To interpret the wavefunction fully it is necessary to be able to extract dynamical information from it. The predictions of quantum mechanics are often very different from those of classical mechanics, and those differences are essential for understanding the structures and properties of atoms and molecules.

\subsection*{What is the key idea?}
The dynamical information in the wavefunction is extracted by calculating the expectation values of hermitian operators.

\begin{displayquote}
What do you need to know already?\\
You need to know that the state of a system is fully described by a wavefunction (Topic 7B), and that the probability density is proportional to the square modulus of the wavefunction.
\end{displayquote}

A wavefunction contains all the information it is possible to obtain about the dynamical properties of a particle (for example, its location and momentum). The Born interpretation (Topic 7B) provides information about location, but the wavefunction contains other information, which is extracted by using the methods described in this Topic.

\subsection*{7.. 1 Operators}
The Schrödinger equation can be written in the succinct form


\begin{equation*}
\hat{H} \psi=E \psi \tag{7С.1a}
\end{equation*}


Operator form of Schrödinger equation

Comparison of this expression with the one-dimensional Schrödinger equation

$$
-\frac{\hbar^{2}}{2 m} \frac{\mathrm{~d}^{2} \psi}{\mathrm{~d} x^{2}}+V(x) \psi=E \psi
$$

shows that in one dimension

$$
\hat{H}=-\frac{\hbar^{2}}{2 m} \frac{\mathrm{~d}^{2}}{\mathrm{~d} x^{2}}+V(x)
$$

Hamiltonian operator

The quantity $\hat{H}$ (commonly read h-hat) is an operator, an expression that carries out a mathematical operation on a function. In this case, the operation is to take the second derivative of $\psi$, and (after multiplication by $-\hbar^{2} / 2 m$ ) to add the result to the outcome of multiplying $\psi$ by $V(x)$.

The operator $\hat{H}$ plays a special role in quantum mechanics, and is called the hamiltonian operator after the nineteenth century mathematician William Hamilton, who developed a form of classical mechanics which, it subsequently turned out, is well suited to the formulation of quantum mechanics. The hamiltonian operator (and commonly simply 'the hamiltonian') is the operator corresponding to the total energy of the system, the sum of the kinetic and potential energies. In eqn 7C.1b the second term on the right is the potential energy, so the first term (the one involving the second derivative) must be the operator for kinetic energy.

In general, an operator acts on a function to produce a new function, as in

$$
\text { (operator)(function) }=(\text { new function })
$$

In some cases the new function is the same as the original function, perhaps multiplied by a constant. Combinations of operators and functions that have this property are of great importance in quantum mechanics.

\subsection*{Brief illustration 7C. 1}
For example, when the operator $\mathrm{d} / \mathrm{d} x$, which means 'take the derivative of the following function with respect to $x$ ', acts on the function $\sin a x$, it generates the new function $a \cos a x$. However, when $\mathrm{d} / \mathrm{d} x$ operates on $\mathrm{e}^{-a x}$ it generates $-a \mathrm{e}^{-a x}$, which is the original function multiplied by the constant $-a$.

\subsection*{(a) Eigenvalue equations}
The Schrödinger equation written as in eqn 7C.1a is an eigenvalue equation, an equation of the form\\
$($ operator $)($ function $)=($ constant factor $) \times($ same function $)$

In an eigenvalue equation, the action of the operator on the function generates the same function, multiplied by a constant. If a general operator is denoted $\hat{\Omega}$ (where $\Omega$ is uppercase omega) and the constant factor by $\omega$ (lowercase omega), then an eigenvalue equation has the form\\
$\hat{\Omega} \psi=\omega \psi$\\
Eigenvalue equation\\
(7C.2b)

If this relation holds, the function $\psi$ is said to be an eigenfunction of the operator $\hat{\Omega}$, and $\omega$ is the eigenvalue associated with that eigenfunction. With this terminology, eqn 7C.2a can be written\\
$($ operator $)($ eigenfunction $)=($ eigenvalue $) \times($ eigenfunction $)$\\
(7C.2c)\\
Equation 7C.1a is therefore an eigenvalue equation in which $\psi$ is an eigenfunction of the hamiltonian and $E$ is the associated eigenvalue. It follows that 'solving the Schrödinger equation' can be expressed as 'finding the eigenfunctions and eigenvalues of the hamiltonian operator for the system'.\\
Just as the hamiltonian is the operator corresponding to the total energy, there are operators that represent other observables, the measurable properties of the system, such as linear momentum or electric dipole moment. For each such operator $\hat{\Omega}$ there is an eigenvalue equation of the form $\hat{\Omega} \psi=\omega \psi$, with the following significance:

\begin{displayquote}
If the wavefunction is an eigenfunction of the operator $\hat{\Omega}$ corresponding to the observable $\Omega$, then the outcome of a measurement of the property $\Omega$ will be the eigenvalue corresponding to that eigenfunction.
\end{displayquote}

Quantum mechanics is formulated by constructing the operator corresponding to the observable of interest and then predicting the outcome of a measurement by examining the eigenvalues of the operator.

\subsection*{(b) The construction of operators}
A basic postulate of quantum mechanics specifies how to set up the operator corresponding to a given observable.

Observables are represented by operators built from the following position and linear momentum operators:


\begin{equation*}
\hat{x}=x \times \quad \hat{p}_{x}=\frac{\hbar}{i} \frac{\mathrm{~d}}{\mathrm{~d} x} \quad \text { Specification of operators } \tag{7C.3}
\end{equation*}


That is, the operator for location along the $x$-axis is multiplication (of the wavefunction) by $x$, and the operator for linear momentum parallel to the $x$-axis is $\hbar / \mathrm{i}$ times the derivative (of the wavefunction) with respect to $x$.

The definitions in eqn 7C. 3 are used to construct operators for other spatial observables. For example, suppose the potential energy has the form $V(x)=\frac{1}{2} k_{\mathrm{f}} x^{2}$, where $k_{\mathrm{f}}$ is a constant (this potential energy describes the vibrations of atoms in molecules). Because the operator for $x$ is multiplication by $x$, by extension the operator for $x^{2}$ is multiplication by $x$ and then by $x$ again, or multiplication by $x^{2}$. The operator corresponding to $\frac{1}{2} k_{\mathrm{f}} x^{2}$ is therefore


\begin{equation*}
\hat{V}(x)=\frac{1}{2} k_{\mathrm{f}} x^{2} \times \tag{7С.4}
\end{equation*}


In practice, the multiplication sign is omitted and multiplication is understood. To construct the operator for kinetic en-\\
ergy, the classical relation between kinetic energy and linear momentum, $E_{\mathrm{k}}=p_{x}^{2} / 2 m$ is used. Then, by using the operator for $p_{x}$ from eqn 7C.3:


\begin{equation*}
\hat{E}_{\mathrm{k}}=\frac{1}{2 m} \overbrace{\left(\frac{\hbar}{\mathrm{i}} \frac{\mathrm{~d}}{\mathrm{~d} x}\right)}^{\hat{p}_{x}} \overbrace{\left(\frac{\hbar}{\mathrm{i}} \frac{\mathrm{~d}}{\mathrm{~d} x}\right)}^{\hat{p}_{x}}=-\frac{\hbar^{2}}{2 m} \frac{\mathrm{~d}^{2}}{\mathrm{~d} x^{2}} \tag{7С.5}
\end{equation*}


It follows that the operator for the total energy, the hamiltonian operator, is


\begin{equation*}
\hat{H}=\hat{E}_{\mathrm{k}}+\hat{V}=-\frac{\hbar^{2}}{2 m} \frac{\mathrm{~d}^{2}}{\mathrm{~d} x^{2}}+\hat{V}(x) \quad \text { Hamiltonian operator } \tag{7C.6}
\end{equation*}


where $\hat{V}(x)$ is the operator corresponding to whatever form the potential energy takes, exactly as in eqn 7C.1b.

\subsection*{Example 7C. 1 Determining the value of an observable}
What is the linear momentum of a free particle described by the wavefunctions (a) $\psi(x)=\mathrm{e}^{\mathrm{i} k x}$ and (b) $\psi(x)=\mathrm{e}^{-\mathrm{i} k x}$ ?

Collect your thoughts You need to operate on $\psi$ with the operator corresponding to linear momentum (eqn 7C.3), and inspect the result. If the outcome is the original wavefunction multiplied by a constant (that is, if the application of the operator results in an eigenvalue equation), then you can identify the constant with the value of the observable.

The solution (a) For $\psi(x)=\mathrm{e}^{\mathrm{i} k x}$,

$$
\hat{p}_{x} \psi=\frac{\hbar}{\mathrm{i}} \frac{\mathrm{~d} \psi}{\mathrm{~d} x}=\frac{\hbar}{\mathrm{i}} \frac{\mathrm{de}}{\mathrm{~d} k x}=\frac{\hbar}{\mathrm{i}} \times \mathrm{i} k \mathrm{e}^{\mathrm{i} k x}=\overbrace{+k \hbar}^{\text {Eigenvalue }}
$$

This is an eigenvalue equation, with eigenvalue $+k \hbar$. It follows that a measurement of the momentum will give the value $p_{x}=$ $+k \hbar$.\\
(b) For $\psi(x)=\mathrm{e}^{-\mathrm{i} k x}$,

$$
\hat{p}_{x} \psi=\frac{\hbar}{\mathrm{i}} \frac{\mathrm{~d} \psi}{\mathrm{~d} x}=\frac{\hbar}{\mathrm{i}} \frac{\mathrm{de}}{\mathrm{-} \mathrm{i} k x} \frac{\hbar}{\mathrm{~d} x} \times(-\mathrm{i} k) \mathrm{e}^{-\mathrm{i} k x}=\overbrace{-k \hbar}^{\text {Eigenvalue }} \psi
$$

Now the eigenvalue is $-k \hbar$, so $p_{x}=-k \hbar$. In case (a) the momentum is positive, meaning that the particle is travelling in the positive $x$-direction, whereas in (b) the particle is moving in the opposite direction.

Comment. A general feature of quantum mechanics is that taking the complex conjugate of a wavefunction reverses the direction of travel. An implication is that if the wavefunction is real (such as $\cos k x$ ), then taking the complex conjugate leaves the wavefunction unchanged: there is no net direction of travel.

Self-test 7C. 1 What is the kinetic energy of a particle described by the wavefunction $\cos k x$ ?\\
\includegraphics[max width=\textwidth, center]{2025_02_27_d4b95d9718f0b9e2e6b9g-285(2)}

Figure 7C. 1 The average kinetic energy of a particle can be inferred from the average curvature of the wavefunction. This figure shows two wavefunctions: the sharply curved function corresponds to a higher kinetic energy than the less sharply curved function.

The expression for the kinetic energy operator (eqn 7C.5) reveals an important point about the Schrödinger equation. In mathematics, the second derivative of a function is a measure of its curvature: a large second derivative indicates a sharply curved function (Fig. 7C.1). It follows that a sharply curved wavefunction is associated with a high kinetic energy, and one with a low curvature is associated with a low kinetic energy.

The curvature of a wavefunction in general varies from place to place (Fig. 7C.2): wherever a wavefunction is sharply curved, its contribution to the total kinetic energy is large; wherever the wavefunction is not sharply curved, its contribution to the overall kinetic energy is low. The observed kinetic energy of the particle is an average of all the contributions of the kinetic energy from each region. Hence, a particle can be expected to have a high kinetic energy if the average curvature of its wavefunction is high. Locally there can be both positive and negative contributions to the kinetic energy (because\\
\includegraphics[max width=\textwidth, center]{2025_02_27_d4b95d9718f0b9e2e6b9g-285}

Figure 7C. 2 The observed kinetic energy of a particle is an average of contributions from the entire space covered by the wavefunction. Sharply curved regions contribute a high kinetic energy to the average; less sharply curved regions contribute only a small kinetic energy.\\
\includegraphics[max width=\textwidth, center]{2025_02_27_d4b95d9718f0b9e2e6b9g-285(1)}

Figure 7C. 3 The wavefunction of a particle with a potential energy $V$ that decreases towards the right. As the total energy is constant, the kinetic energy $E_{k}$ increases to the right, which results in a faster oscillation and hence greater curvature of the wavefunction.\\
the curvature can be either positive, $\cup$, or negative, $\cap$ ) locally, but the average is always positive.

The association of high curvature with high kinetic energy is a valuable guide to the interpretation of wavefunctions and the prediction of their shapes. For example, suppose the wavefunction of a particle with a given total energy and a potential energy that decreases with increasing $x$ is required. Because the difference $E-V=E_{\mathrm{k}}$ increases from left to right, the wavefunction must become more sharply curved by oscillating more rapidly as $x$ increases (Fig. 7C.3). It is therefore likely that the wavefunction will look like the function sketched in the illustration, and more detailed calculation confirms this to be so.

\subsection*{(c) Hermitian operators}
All the quantum mechanical operators that correspond to observables have a very special mathematical property: they are 'hermitian'. A hermitian operator is one for which the following relation is true:

\[
\int \psi_{i}^{\star} \hat{\Omega} \psi_{j} \mathrm{~d} \tau=\left\{\int \psi_{j}^{\star} \hat{\Omega} \psi_{i} \mathrm{~d} \tau\right\}^{\star} \quad \begin{align*}
& \text { Hermiticity }  \tag{7С.7}\\
& \text { [definition] }
\end{align*}
\]

As stated in Topic 7B, in quantum mechanics $\int \ldots \mathrm{d} \tau$ implies integration over the full range of all relevant spatial variables.

It is easy to confirm that the position operator $(x \times)$ is hermitian because in this case the order of the factors in the integrand can be changed:

$$
\int \psi_{i}^{\star} x \psi_{j} \mathrm{~d} \tau=\int \psi_{j} x \psi_{i}^{\star} \mathrm{d} \tau=\left\{\int \psi_{j}^{\star} x \psi_{i} \mathrm{~d} \tau\right\}^{\star}
$$

The final step uses $\left(\psi^{*}\right)^{*}=\psi$. The demonstration that the linear momentum operator is hermitian is more involved because the order of functions being differentiated cannot be changed.

\subsection*{How is that done? 7C. 1 Showing that the linear}
 momentum operator is hermitianThe task is to show that

$$
\int \psi_{i}^{\star} \hat{p}_{x} \psi_{j} \mathrm{~d} \tau=\left\{\int \psi_{j}^{\star} \hat{p}_{x} \psi_{i} \mathrm{~d} \tau\right\}^{*}
$$

with $\hat{p}_{x}$ given in eqn 7C.3. To do so, use 'integration by parts' (see The chemist's toolkit 15) which, when applied to the present case, gives

$$
\begin{aligned}
\int \psi_{i}^{*} \hat{p}_{x} \psi_{j} \mathrm{~d} \tau & =\frac{\hbar}{\mathrm{i}} \int_{-\infty}^{\infty} \overbrace{\psi_{i}^{*}}^{f} \frac{\mathrm{~d} \psi_{j}}{\mathrm{~d} x} \mathrm{~d} / \mathrm{d} x \\
& =\left.\overbrace{\overbrace{\mathrm{i}}}^{0} \overbrace{\psi_{i}^{*}}^{\psi_{j}}\right|_{-\infty} ^{\infty}-\frac{\hbar}{\mathrm{i}} \int_{-\infty}^{\infty} \overbrace{\psi_{j}}^{g} \overbrace{\frac{\mathrm{~d} \psi_{i}^{*}}{\mathrm{~d} x}}^{\mathrm{d} / \mathrm{d} x} \mathrm{~d} x
\end{aligned}
$$

The blue term is zero because all wavefunctions are either zero at $x= \pm \infty$ (see Topic 7B) or the product $\psi_{i}^{*} \psi_{j}$ converges to the same value at $x=+\infty$ and $x=-\infty$. As a result

$$
\begin{aligned}
\int \psi_{i}^{*} \hat{p}_{x} \psi_{j} \mathrm{~d} \tau & =-\frac{\hbar}{\mathrm{i}} \int_{-\infty}^{\infty} \psi_{j} \frac{\mathrm{~d} \psi_{i}^{*}}{\mathrm{~d} x} \mathrm{~d} x=\left\{\frac{\hbar}{\mathrm{i}} \int_{-\infty}^{\infty} \psi_{j}^{\star} \frac{\mathrm{d} \psi_{i}}{\mathrm{~d} x} \mathrm{~d} x\right\}^{*} \\
& =\left\{\int \psi_{j}^{*} \hat{p}_{x} \psi_{i} \mathrm{~d} \tau\right\}^{*}
\end{aligned}
$$

as was to be proved. The final line uses $\left(\psi^{*}\right)^{\star}=\psi$ and $\mathrm{i}^{*}=-\mathrm{i}$.

Hermitian operators are enormously important in quantum mechanics because their eigenvalues are real: that is, $\omega^{*}=\omega$. Any measurement must yield a real value because a position, momentum, or an energy cannot be complex or imaginary. Because the outcome of a measurement of an observable is one of the eigenvalues of the corresponding operator, those eigenvalues must be real. It therefore follows that an operator that represents an observable must be hermitian. The\\
proof that their eigenfunctions are real makes use of the definition of hermiticity in eqn 7C.7.

How is that done? 7C. 2 Showing that the eigenvalues of hermitian operators are real

Begin by setting $\psi_{i}$ and $\psi_{j}$ to be the same, writing them both as $\psi$. Then eqn 7C. 7 becomes

$$
\int \psi^{\star} \hat{\Omega} \psi \mathrm{d} \tau=\left\{\int \psi^{\star} \hat{\Omega} \psi \mathrm{d} \tau\right\}^{\star}
$$

Next suppose that $\psi$ is an eigenfunction of $\hat{\Omega}$ with eigenvalue $\omega$. That is, $\hat{\Omega} \psi=\omega \psi$. Now use this relation in both integrals on the left- and right-hand sides:\\
$\int \psi^{*} \omega \psi \mathrm{~d} \tau=\left\{\int \psi^{\star} \omega \psi \mathrm{d} \tau\right\}^{*}$\\
The eigenvalue is a constant that can be taken outside the integrals:

$$
\omega \int \psi^{\star} \psi \mathrm{d} \tau=\left\{\omega \int \psi^{\star} \psi \mathrm{d} \tau\right\}^{\star}=\omega^{\star} \int \psi \psi^{\star} \mathrm{d} \tau
$$

Finally, the (blue) integrals cancel, leaving $\omega=\omega^{*}$. It follows that $\omega$ is real.

\subsection*{(d) Orthogonality}
To say that two different functions $\psi_{i}$ and $\psi_{j}$ are orthogonal means that the integral (over all space) of $\psi_{i}^{*} \psi_{j}$ is zero:

\[
\int \psi_{i}^{\star} \psi_{j} \mathrm{~d} \tau=0 \text { for } i \neq j \quad \begin{align*}
& \text { Orthogonality }  \tag{7C.8}\\
& \text { [definition] }
\end{align*}
\]

Functions that are both normalized and mutually orthogonal are called orthonormal. Hermitian operators have the important property that

\subsection*{Eigenfunctions that correspond to different eigenvalues of a hermitian operator are orthogonal.}
The proof of this property also follows from the definition of hermiticity (eqn 7C.7).

\subsection*{The chemist's toolkit 15 Integration by parts}
Many integrals in quantum mechanics have the form $\int f(x) h(x) \mathrm{d} x$ where $f(x)$ and $h(x)$ are two different functions. Such integrals can often be evaluated by regarding $h(x)$ as the derivative of another function, $g(x)$, such that $h(x)=\mathrm{d} g(x) / \mathrm{d} x$. For instance, if $h(x)=x$, then $g(x)=\frac{1}{2} x^{2}$. The integral is then found using integration by parts:

$$
\int f \frac{\mathrm{~d} g}{\mathrm{~d} x} \mathrm{~d} x=f g-\int g \frac{\mathrm{~d} f}{\mathrm{~d} x} \mathrm{~d} x
$$

The procedure is successful only if the integral on the right turns out to be one that can be evaluated more easily than the one on the left. The procedure is often summarized by expressing this relation as

$$
\int f \mathrm{~d} g=f g-\int g \mathrm{~d} f
$$

As an example, consider integration of $x \mathrm{e}^{-a x}$. In this case, $f(x)=$ $x$, so $\mathrm{d} f(x) / \mathrm{d} x=1$ and $\mathrm{d} g(x) / \mathrm{d} x=\mathrm{e}^{-a x}$, so $g(x)=-(1 / a) \mathrm{e}^{-a x}$. Then

$$
\begin{aligned}
& \int \overbrace{x \mathrm{e}^{-a x}}^{f \mathrm{~d} g / \mathrm{d} x}=\overbrace{x}^{f} \overbrace{\frac{-\mathrm{e}^{-a x}}{a}}^{g}-\int \overbrace{\frac{-\mathrm{e}^{-a x}}{a}}^{g} \mathrm{~d} / \mathrm{d} \mathrm{~d} x \\
& =-\frac{x \mathrm{e}^{-a x}}{a}+\frac{1}{a} \int \mathrm{e}^{-a x} \mathrm{~d} x=-\frac{x \mathrm{e}^{-a x}}{a}-\frac{\mathrm{e}^{-a x}}{a^{2}}+\text { constant }
\end{aligned}
$$

How is that done? 7C. 3 Showing that the eigenfunctions of hermitian operators are orthogonal\\
Start by supposing that $\psi_{j}$ is an eigenfunction of $\hat{\Omega}$ with eigenvalue $\omega_{j}\left(\right.$ i.e. $\left.\Omega \psi_{j}=\omega_{j} \psi_{j}\right)$ and that $\psi_{i}$ is an eigenfunction with a different eigenvalue $\omega_{i}$ (i.e. $\hat{\Omega} \psi_{i}=\omega_{i} \psi_{i}$, with $\omega_{i} \neq \omega_{j}$ ). Then eqn 7C. 7 becomes

$$
\int \psi_{i}^{\star} \omega_{j} \psi_{j} \mathrm{~d} \tau=\left\{\int \psi_{j}^{\star} \omega_{i} \psi_{i} \mathrm{~d} \tau\right\}^{\star}
$$

The eigenvalues are constants and can be taken outside the integrals; moreover, they are real (being the eigenvalues of hermitian operators), so $\omega_{i}^{*}=\omega_{i}$. Then

$$
\omega_{j} \int \psi_{i}^{\star} \psi_{j} \mathrm{~d} \tau=\omega_{i}\left\{\int \psi_{j}^{\star} \psi_{i} \mathrm{~d} \tau\right\}^{*}
$$

Next, note that $\left\{\int \psi_{j}^{*} \psi_{i} \mathrm{~d} \tau\right\}^{*}=\int \psi_{j} \psi_{i}^{*} \mathrm{~d} \tau$, so\\
$\omega_{j} \int \psi_{i}^{*} \psi_{j} \mathrm{~d} \tau=\omega_{i} \int \psi_{j} \psi_{i}^{*} \mathrm{~d} \tau$, hence $\left(\omega_{j}-\omega_{i}\right) \int \psi_{i}^{*} \psi_{j} \mathrm{~d} \tau=0$\\
The two eigenvalues are different, so $\omega_{j}-\omega_{i} \neq 0$; therefore it must be the case that $\int \psi_{i}^{\star} \psi_{j} \mathrm{~d} \tau=0$. That is, the two eigenfunctions are orthogonal, as was to be proved.

The hamiltonian operator is hermitian (it corresponds to an observable, the energy, but its hermiticity can be proved specifically). Therefore, if two of its eigenfunctions correspond to different energies, the two functions must be orthogonal. The property of orthogonality is of great importance in quantum mechanics because it eliminates a large number of integrals from calculations. Orthogonality plays a central role in the theory of chemical bonding (Focus 9) and spectroscopy (Focus 11).

\subsection*{Example 7C. 2 Verifying orthogonality}
Two possible wavefunctions for a particle constrained to move along the $x$ axis between $x=0$ and $x=L$ are $\psi_{1}=\sin (\pi x / L)$ and $\psi_{2}=\sin (2 \pi x / L)$. Outside this region the wavefunctions are zero. The wavefunctions correspond to different energies. Verify that the two wavefunctions are mutually orthogonal.

Collect your thoughts To verify the orthogonality of two functions, you need to integrate $\psi_{2}^{*} \psi_{1}=\sin (2 \pi x / L) \sin (\pi x / L)$ over all space, and show that the result is zero. In principle the integral is taken from $x=-\infty$ to $x=+\infty$, but the wavefunctions are zero outside the range $x=0$ to $L$ so you need integrate only over this range. Relevant integrals are given in the Resource section.\\
The solution To evaluate the integral, use Integral T. 5 from the Resource section with $a=2 \pi / L$ and $b=\pi / L$ :\\
$\int_{0}^{L} \sin (2 \pi x / L) \sin (\pi x / L) \mathrm{d} x=\left.\frac{\sin (\pi x / L)}{2(\pi / L)}\right|_{0} ^{L}-\left.\frac{\sin (3 \pi x / L)}{2(3 \pi / L)}\right|_{0} ^{L}=0$\\
The sine functions have been evaluated by using $\sin n \pi=0$ for $n=0, \pm 1, \pm 2, \ldots$. The two functions are therefore mutually orthogonal.

Self-test 7C. 2 The next higher energy level has $\psi_{3}=\sin (3 \pi x / L)$. Confirm that the functions $\psi_{1}=\sin (\pi x / L)$ and $\psi_{3}=\sin (3 \pi x / L)$ are mutually orthogonal.

\subsection*{7C. 2 Superpositions and expectation values}
The hamiltonian for a free particle moving in one dimension is

$$
\hat{H}=-\frac{\hbar^{2}}{2 m} \frac{\mathrm{~d}^{2}}{\mathrm{~d} x^{2}}
$$

The particle is 'free' in the sense that there is no potential to constrain it, hence $V(x)=0$. It is easily confirmed that $\psi(x)=\cos k x$ is an eigenfunction of this operator

$$
\hat{H} \psi(x)=-\frac{\hbar^{2}}{2 m} \frac{\mathrm{~d}^{2}}{\mathrm{~d} x^{2}} \cos k x=\frac{k^{2} \hbar^{2}}{2 m} \cos k x
$$

The energy associated with this wavefunction, $k^{2} \hbar^{2} / 2 m$, is therefore well defined, as it is the eigenvalue of an eigenvalue equation. However, the same is not necessarily true of other observables. For instance, $\cos k x$ is not an eigenfunction of the linear momentum operator:


\begin{equation*}
\hat{p}_{x} \psi(x)=\frac{\hbar}{\mathrm{i}} \frac{\mathrm{~d} \psi}{\mathrm{~d} x}=\frac{\hbar}{\mathrm{i}} \frac{\mathrm{~d} \cos k x}{\mathrm{~d} x}=-\frac{k \hbar}{\mathrm{i}} \sin k x \tag{7C.9}
\end{equation*}


This expression is not an eigenvalue equation, because the function on the right $(\sin k x)$ is different from that on the left $(\cos k x)$.

When the wavefunction of a particle is not an eigenfunction of an operator, the corresponding observable does not have a definite value. However, in the current example the momentum is not completely indefinite because the cosine wavefunction can be written as a linear combination, or sum, ${ }^{1}$ of $\mathrm{e}^{\mathrm{i} k x}$ and $\mathrm{e}^{-\mathrm{i} k x}$ : $\cos k x=\frac{1}{2}\left(\mathrm{e}^{\mathrm{i} k x}+\mathrm{e}^{-\mathrm{i} k x}\right)$ (see The chemist's toolkit 16). As shown in Example 7C.1, these two exponential functions are eigenfunctions of $\hat{p}_{x}$ with eigenvalues $+k \hbar$ and $-k \hbar$, respectively. They therefore each correspond to a state of definite but different momentum. The wavefunction $\cos k x$ is said to be a superposition of the two individual wavefunctions $\mathrm{e}^{\mathrm{i} k x}$ and $\mathrm{e}^{-\mathrm{i} k x}$, and is written

$$
\psi=\underbrace{\mathrm{e}^{+\mathrm{i} k x}}_{\begin{array}{c}
\text { Particle with linear } \\
\text { momentum }+k \hbar
\end{array}}+\underbrace{\mathrm{e}^{-\mathrm{i} k x}}_{\begin{array}{c}
\text { Particle with linear } \\
\text { momentum }-k \hbar
\end{array}}
$$

The interpretation of this superposition is that if many repeated measurements of the momentum are made, then half the measurements would give the value $p_{x}=+k \hbar$, and half would give the value $p_{x}=-k \hbar$. The two values $\pm k \hbar$ occur equally often since $\mathrm{e}^{\mathrm{i} k x}$ and $\mathrm{e}^{-\mathrm{i} k x}$ contribute equally to the superposition. All that can be inferred from the wavefunction $\cos k x$ about the linear momentum is that the particle it describes is equally

\footnotetext{${ }^{1}$ A linear combination is more general than a sum, for it includes weighted sums of the form $a x+b y+\cdots$ where $a, b, \ldots$ are constants. A sum is a linear combination with $a=b=\cdots=1$.
}\subsection*{The chemist's toolkit 16 \\
 Euler's formula}
A complex number $z=x+\mathrm{i} y$ can be represented as a point in a plane, the complex plane, with $\operatorname{Re}(z)$ along the $x$-axis and $\operatorname{Im}(z)$ along the $y$-axis (Sketch 1). The position of the point can also be specified in terms of a distance $r$ and an angle $\phi$ (the polar coordinates). Then $x=r \cos \phi$ and $y=r \sin \phi$, so it follows that

$$
z=r(\cos \phi+\mathrm{i} \sin \phi)
$$

The angle $\phi$, called the argument of $z$, is the angle that $r$ makes with the $x$-axis. Because $y / x=\tan \phi$, it follows that

$$
r=\left(x^{2}+y^{2}\right)^{1 / 2}=|z| \quad \phi=\arctan \frac{y}{x}
$$

\begin{center}
\includegraphics[max width=\textwidth]{2025_02_27_d4b95d9718f0b9e2e6b9g-288(2)}
\end{center}

Sketch 1

One of the most useful relations involving complex numbers is Euler's formula:

$$
\mathrm{e}^{\mathrm{i} \phi}=\cos \phi+\mathrm{i} \sin \phi
$$

likely to be found travelling in the positive and negative $x$ directions, with the same magnitude, $k \hbar$, of the momentum.

A similar interpretation applies to any wavefunction written as a linear combination of eigenfunctions of an operator. In general, a wavefunction can be written as the following linear combination

$$
\psi=c_{1} \psi_{1}+c_{2} \psi_{2}+\cdots=\sum_{k} c_{k} \boldsymbol{\psi}_{k} \quad \begin{aligned}
& \text { Linear combination } \\
& \text { of eigenfunctions }
\end{aligned}
$$

(7C.10)\\
where the $c_{k}$ are numerical (possibly complex) coefficients and the $\psi_{k}$ are different eigenfunctions of the operator $\hat{\Omega}$ corresponding to the observable of interest. The functions $\psi_{k}$ are said to form a complete set in the sense that any arbitrary function can be expressed as a linear combination of them. Then, according to quantum mechanics:

\begin{itemize}
  \item A single measurement of the observable corresponding to the operator $\hat{\Omega}$ will give one of the eigenvalues corresponding to the $\psi_{k}$ that contribute to the superposition.
  \item The probability of measuring a specific eigenvalue in a series of measurements is proportional to the square modulus $\left(\left|c_{k}\right|^{2}\right)$ of the corresponding coefficient in the linear combination.\\
from which it follows that $z=r(\cos \phi+\mathrm{i} \sin \phi)$ can be written
\end{itemize}

$$
z=r \mathrm{e}^{\mathrm{i} \phi}
$$

Two more useful relations arise by noting that $\mathrm{e}^{-\mathrm{i} \phi}=\cos (-\phi)+$ $i \sin (-\phi)=\cos \phi-i \sin \phi$; it then follows that

$$
\cos \phi=\frac{1}{2}\left(\mathrm{e}^{\mathrm{i} \phi}+\mathrm{e}^{-\mathrm{i} \phi}\right) \quad \sin \phi=-\frac{1}{2} \mathrm{i}\left(\mathrm{e}^{\mathrm{i} \phi}-\mathrm{e}^{-\mathrm{i} \phi}\right)
$$

The polar form of a complex number is commonly used to perform arithmetical operations. For instance, the product of two complex numbers in polar form is

$$
z_{1} z_{2}=\left(r_{1} \mathrm{e}^{\mathbf{i} \phi_{1}}\right)\left(r_{2} \mathrm{e}^{\mathrm{i} \phi_{2}}\right)=r_{1} r_{2} \mathrm{e}^{\mathrm{i}\left(\phi_{1}+\phi_{2}\right)}
$$

This construction is illustrated in Sketch 2.\\
\includegraphics[max width=\textwidth, center]{2025_02_27_d4b95d9718f0b9e2e6b9g-288(1)}

Sketch 2

The average value of a large number of measurements of an observable $\Omega$ is called the expectation value of the operator $\hat{\Omega}$, and is written $\langle\Omega\rangle$. For a normalized wavefunction $\psi$, the expectation value of $\hat{\Omega}$ is calculated by evaluating the integral

$$
\langle\Omega\rangle=\int \psi^{\star} \hat{\Omega} \psi \mathrm{d} \tau
$$

Expectation value [normalized wavefunction, (7C.11) definition]

This definition can be justified by considering two cases, one where the wavefunction is an eigenfunction of the operator $\hat{\Omega}$ and another where the wavefunction is a superposition of that operator's eigenfunctions.

\subsection*{How is that done? 7C. 4 Justifying the expression for the expectation value of an operator}
If the wavefunction $\psi$ is an eigenfunction of $\hat{\Omega}$ with eigenvalue $\omega$ (so $\hat{\Omega} \psi=\omega \psi$ ),\\
\includegraphics[max width=\textwidth, center]{2025_02_27_d4b95d9718f0b9e2e6b9g-288}

The interpretation of this expression is that, because the wavefunction is an eigenfunction of $\hat{\Omega}$, each observation of the property $\Omega$ results in the same value $\omega$; the average value of all the observations is therefore $\omega$.

Now suppose the (normalized) wavefunction is the linear combination of two eigenfunctions of the operator $\hat{\Omega}$, each of which is individually normalized to 1 . Then

$$
\begin{aligned}
& \langle\Omega\rangle=\int\left(c_{1} \psi_{1}+c_{2} \psi_{2}\right)^{*} \hat{\Omega}\left(c_{1} \psi_{1}+c_{2} \psi_{2}\right) \mathrm{d} \tau \\
& =\int\left(c_{1} \psi_{1}+c_{2} \psi_{2}\right)^{*}(\overbrace{c_{1} \hat{\Omega} \psi_{1}}^{\omega_{1} \psi_{1}}+c_{2} \overbrace{\hat{\Omega} \psi_{2}}^{\omega_{2} \psi_{2}}) \mathrm{d} \tau \\
& =\int\left(c_{1} \psi_{1}+c_{2} \psi_{2}\right)^{*}\left(c_{1} \omega_{1} \psi_{1}+c_{2} \omega_{2} \psi_{2}\right) \mathrm{d} \tau \\
& =c_{1}^{*} c_{1} \omega_{1} \overbrace{\int \psi_{1}^{*} \psi_{1} \mathrm{~d} \tau}^{1}+c_{2}^{*} c_{2} \omega_{2} \overbrace{\int \psi_{2}^{*} \psi_{2} \mathrm{~d} \tau}^{1} \\
& +c_{1}^{*} c_{2} \omega_{2} \overbrace{\int \psi_{1}^{*} \psi_{2} \mathrm{~d} \tau}^{0}+c_{2}^{*} c_{1} \omega_{1} \overbrace{\int \psi_{2}^{*} \psi_{1} \mathrm{~d} \tau}^{0}
\end{aligned}
$$

The first two integrals on the right are both equal to 1 because the wavefunctions $\psi_{1}$ and $\psi_{2}$ are individually normalized. Because $\psi_{1}$ and $\psi_{2}$ correspond to different eigenvalues of a hermitian operator, they are orthogonal, so the third and fourth integrals on the right are zero. Therefore

$$
\langle\Omega\rangle=\left|c_{1}\right|^{2} \omega_{1}+\left|c_{2}\right|^{2} \omega_{2}
$$

The interpretation of this expression is that in a series of measurements each individual measurement yields either $\omega_{1}$ or $\omega_{2}$, but that the probability of $\omega_{1}$ occurring is $\left|c_{1}\right|^{2}$, and likewise the probability of $\omega_{2}$ occurring is $\left|c_{2}\right|^{2}$. The average is the sum of the two eigenvalues, but with each weighted according to the probability that it will occur in a measurement:

$$
\begin{aligned}
\text { average }= & \left(\text { probability of } \omega_{1} \text { occurring }\right) \times \omega_{1} \\
& +\left(\text { probability of } \omega_{2} \text { occurring }\right) \times \omega_{2}
\end{aligned}
$$

The expectation value therefore predicts the result of taking a series of measurements, each of which gives an eigenvalue, and then taking the weighted average of these values. This justifies the form of eqn 7C.11.

\subsection*{Example 7C. 3}
Calculating an expectation value\\
Calculate the average value of the position of an electron in the lowest energy state of a one-dimensional box of length $L$, with the (normalized) wavefunction $\psi=(2 / L)^{1 / 2} \sin (\pi x / L)$ inside the box and zero outside it.

Collect your thoughts The average value of the position is the expectation value of the operator corresponding to position, which is multiplication by $x$. To evaluate $\langle x\rangle$, you need to evaluate the integral in eqn 7C. 11 with $\hat{\Omega}=\hat{x}=x \times$

The solution The expectation value of position is

$$
\langle x\rangle=\int_{0}^{L} \psi^{*} \hat{x} \psi \mathrm{~d} x \quad \text { with } \psi=\left(\frac{2}{L}\right)^{1 / 2} \sin \frac{\pi x}{L} \quad \text { and } \hat{x}=x \times
$$

The integral is restricted to the region $x=0$ to $x=L$ because outside this region the wavefunction is zero. Use Integral T. 11 from the Resources section to obtain

$$
\langle x\rangle=\frac{2}{L} \overbrace{\int_{0}^{L} x \sin ^{2} \frac{\pi x}{L} \mathrm{~d} x}^{\text {Integral T. } 11}=\frac{2}{L} \frac{L^{2}}{4}=\frac{1}{2} L
$$

Comment. This result means that if a very large number of measurements of the position of the electron are made, then the mean value will be at the centre of the box. However, each different observation will give a different and unpredictable individual result somewhere in the range $0 \leq x \leq L$ because the wavefunction is not an eigenfunction of the operator corresponding to $x$.

Self-test 7C. 3 Evaluate the mean square position, $\left\langle x^{2}\right\rangle$, of the electron; you will need Integral T. 12 from the Resource section.

$$
\left.{ }_{\tau} T \angle I Z \cdot 0={ }_{\tau} \nu^{2} \frac{\tau}{T}-\frac{\varepsilon}{T}\right\}_{\tau} T: \iota \partial M s u_{V}
$$

The mean kinetic energy of a particle in one dimension is the expectation value of the operator given in eqn 7C.5. Therefore,


\begin{equation*}
\left\langle E_{\mathrm{k}}\right\rangle=\int_{-\infty}^{\infty} \psi^{\star} \hat{E}_{\mathrm{k}} \psi \mathrm{~d} x=-\frac{\hbar^{2}}{2 m} \int_{-\infty}^{\infty} \psi^{\star} \frac{\mathrm{d}^{2} \psi}{\mathrm{~d} x^{2}} \mathrm{~d} x \tag{7С.12}
\end{equation*}


This conclusion confirms the previous assertion that the kinetic energy is a kind of average over the curvature of the wavefunction: a large contribution to the observed value comes from regions where the wavefunction is sharply curved ( $\operatorname{so~}^{2} \psi / \mathrm{d} x^{2}$ is large) and the wavefunction itself is large (so that $\psi^{*}$ is large there too).

\subsection*{7C. 3 The uncertainty principle}
The wavefunction $\psi=\mathrm{e}^{\mathrm{ikx}}$ is an eigenfunction of $\hat{p}_{x}$ with eigenvalue $+k \hbar$ : in this case the wavefunction describes a particle with a definite state of linear momentum. Where, though, is the particle? The probability density is proportional to $\psi^{*} \psi$, so if the particle is described by the wavefunction $\mathrm{e}^{\mathrm{i} k x}$ the probability density is proportional to $\left(\mathrm{e}^{\mathrm{i} k}\right) * \mathrm{e}^{\mathrm{i} k x}=\mathrm{e}^{-\mathrm{i} k x} \mathrm{e}^{\mathrm{i} k x}=\mathrm{e}^{-\mathrm{i} k x+\mathrm{i} k x}=\mathrm{e}^{0}=1$. In other words, the probability density is the same for all values of $x$ : the location of the particle is completely unpredictable. In summary, if the momentum of the particle is known precisely, it is not possible to predict its location.

This conclusion is an example of the consequences of the Heisenberg uncertainty principle, one of the most celebrated results of quantum mechanics:

It is impossible to specify simultaneously, with arbitrary precision, both the linear momentum and the position of a particle.\\
\includegraphics[max width=\textwidth, center]{2025_02_27_d4b95d9718f0b9e2e6b9g-289}\\
\includegraphics[max width=\textwidth, center]{2025_02_27_d4b95d9718f0b9e2e6b9g-290}

Figure 7C. 4 The wavefunction of a particle at a well-defined location is a sharply spiked function that has zero amplitude everywhere except at the position of the particle.

Note that the uncertainty principle also implies that if the position is known precisely, then the momentum cannot be predicted. The argument runs as follows.

Suppose the particle is known to be at a definite location, then its wavefunction must be large there and zero everywhere else (Fig. 7C.4). Such a wavefunction can be created by superimposing a large number of harmonic (sine and cosine) functions, or, equivalently, a number of $\mathrm{e}^{\mathrm{i} k x}$ functions (because $\mathrm{e}^{\mathrm{i} k x}=$ $\cos k x+\mathrm{i} \sin k x)$. In other words, a sharply localized wavefunction, called a wavepacket, can be created by forming a linear combination of wavefunctions that correspond to many different linear momenta.

The superposition of a few harmonic functions gives a wavefunction that spreads over a range of locations (Fig. 7C.5). However, as the number of wavefunctions in the superposition increases, the wavepacket becomes sharper on account of the more complete interference between the positive and nega-\\
\includegraphics[max width=\textwidth, center]{2025_02_27_d4b95d9718f0b9e2e6b9g-290(1)}

Figure 7C. 5 The wavefunction of a particle with an illdefined location can be regarded as a superposition of several wavefunctions of definite wavelength that interfere constructively in one place but destructively elsewhere. As more waves are used in the superposition (as given by the numbers attached to the curves), the location becomes more precise at the expense of uncertainty in the momentum of the particle. An infinite number of waves are needed in the superposition to construct the wavefunction of the perfectly localized particle.\\
tive regions of the individual waves. When an infinite number of components are used, the wavepacket is a sharp, infinitely narrow spike, which corresponds to perfect localization of the particle. Now the particle is perfectly localized but all information about its momentum has been lost. A measurement of the momentum will give a result corresponding to any one of the infinite number of waves in the superposition, and which one it will give is unpredictable. Hence, if the location of the particle is known precisely (implying that its wavefunction is a superposition of an infinite number of momentum eigenfunctions), then its momentum is completely unpredictable.

The quantitative version of the uncertainty principle is

\[
\Delta p_{q} \Delta q \geq \frac{1}{2} \hbar \quad \begin{align*}
& \text { Heisenberg }  \tag{7C.13a}\\
& \text { uncertainty principle }
\end{align*}
\]

In this expression $\Delta p_{q}$ is the 'uncertainty' in the linear momentum parallel to the axis $q$, and $\Delta q$ is the uncertainty in position along that axis. These 'uncertainties' are given by the root-mean-square deviations of the observables from their mean values:


\begin{equation*}
\Delta p_{q}=\left\{\left\langle p_{q}^{2}\right\rangle-\left\langle p_{q}\right\rangle^{2}\right\}^{1 / 2} \quad \Delta q=\left\{\left\langle q^{2}\right\rangle-\langle q\rangle^{2}\right\}^{1 / 2} \tag{7C.13b}
\end{equation*}


If there is complete certainty about the position of the particle $(\Delta q=0)$, then the only way that eqn 7C.13a can be satisfied is for $\Delta p_{q}=\infty$, which implies complete uncertainty about the momentum. Conversely, if the momentum parallel to an axis is known exactly ( $\Delta p_{q}=0$ ), then the position along that axis must be completely uncertain $(\Delta q=\infty)$.

The $p$ and $q$ that appear in eqn 7C.13a refer to the same direction in space. Therefore, whereas simultaneous specification of the position on the $x$-axis and momentum parallel to the $x$-axis are restricted by the uncertainty relation, simultaneous location of position on $x$ and motion parallel to $y$ or $z$ are not restricted.

\subsection*{Example 7C. 4}
\subsection*{Using the uncertainty principle}
Suppose the speed of a projectile of mass 1.0 g is known to within $1 \mu \mathrm{~m} \mathrm{~s}^{-1}$. What is the minimum uncertainty in its position?\\
Collect your thoughts You can estimate $\Delta p$ from $m \Delta v$, where $\Delta v$ is the uncertainty in the speed; then use eqn 7C.13a to estimate the minimum uncertainty in position, $\Delta q$, by using it in the form $\Delta p \Delta q=\frac{1}{2} \hbar$ rearranged into $\Delta q=\frac{1}{2} \hbar / \Delta p$. You will need to use $1 \mathrm{~J}=1 \mathrm{~kg} \mathrm{~m}^{2} \mathrm{~s}^{-2}$.

The solution The minimum uncertainty in position is

$$
\begin{aligned}
\Delta q & =\frac{\hbar}{2 m \Delta v} \\
& =\frac{1.055 \times 10^{-34} \mathrm{Js}}{2 \times\left(1.0 \times 10^{-3} \mathrm{~kg}\right) \times\left(1 \times 10^{-6} \mathrm{~ms}^{-1}\right)}=5 \times 10^{-26} \mathrm{~m}
\end{aligned}
$$

Comment. This uncertainty is completely negligible for all practical purposes. However, if the mass is that of an electron,\\
then the same uncertainty in speed implies an uncertainty in position far larger than the diameter of an atom (the analogous calculation gives $\Delta q=60 \mathrm{~m}$ ).\\
Self-test 7C. 4 Estimate the minimum uncertainty in the speed of an electron in a one-dimensional region of length $2 a_{0}$, the approximate diameter of a hydrogen atom, where $a_{0}$ is the Bohr radius, 52.9 pm .\\
\includegraphics[max width=\textwidth, center]{2025_02_27_d4b95d9718f0b9e2e6b9g-291}

The Heisenberg uncertainty principle is more general than even eqn 7C.13a suggests. It applies to any pair of observables, called complementary observables, for which the corresponding operators $\hat{\Omega}_{1}$ and $\hat{\Omega}_{2}$ have the property

\[
\hat{\Omega}_{1} \hat{\Omega}_{2} \psi \neq \hat{\Omega}_{2} \hat{\Omega}_{1} \psi \quad \begin{align*}
& \text { Complementarity }  \tag{7C.14}\\
& \text { of observables }
\end{align*}
\]

The term on the left implies that $\hat{\Omega}_{2}$ acts first, then $\hat{\Omega}_{1}$ acts on the result, and the term on the right implies that the operations are performed in the opposite order. When the effect of two operators applied in succession depends on their order (as this equation implies), they do not commute. The different outcomes of the effect of applying $\hat{\Omega}_{1}$ and $\hat{\Omega}_{2}$ in a different order are expressed by introducing the commutator of the two operators, which is defined as

\[
\left[\hat{\Omega}_{1}, \hat{\Omega}_{2}\right]=\hat{\Omega}_{1} \hat{\Omega}_{2}-\hat{\Omega}_{2} \hat{\Omega}_{1} \quad \begin{align*}
& \text { Commutator }  \tag{7С.15}\\
& \text { [definition] }
\end{align*}
\]

By using the definitions of the operators for position and momentum, an explicit value of this commutator can be found.

How is that done? 7C. 5 Evaluating the commutator of position and momentum

You need to consider the effect of $\hat{x} \hat{p}_{x}$ (i.e. the effect of $\hat{p}_{x}$ followed by the effect on the outcome of multiplication by $x$ ) on an arbitrary wavefunction $\psi$, which need not be an eigenfunction of either operator.

$$
\hat{x} \hat{p}_{x} \psi=x \times \frac{\hbar}{\mathrm{i}} \frac{\mathrm{~d} \psi}{\mathrm{~d} x}
$$

Then you need to consider the effect of $\hat{p}_{x} \hat{x}$ on the same function (that is, the effect of multiplication by $x$ followed by the effect of $\hat{p}_{x}$ on the outcome):

$$
\hat{p}_{x} \hat{x} \psi=\frac{\hbar}{\mathrm{i}} \frac{\mathrm{~d}(x \psi)}{\mathrm{d} x}=\frac{\hbar}{\mathrm{i}}\left(\psi+x \frac{\mathrm{~d} \psi}{\mathrm{~d} x}\right)
$$

The second expression is different from the first, so $\hat{p}_{x} \hat{x} \psi \neq \hat{x} \hat{p}_{x} \psi$ and therefore the two operators do not commute. You can infer the value of the commutator from the difference of the two expressions:

$$
\left[\hat{x}, \hat{p}_{x}\right] \psi=\hat{x} \hat{p}_{x} \psi-\hat{p}_{x} \hat{x} \psi=-\frac{\hbar}{\mathrm{i}} \psi=\mathrm{i} \hbar \psi, \quad \text { so }\left[\hat{x}, \hat{p}_{x}\right] \psi=\mathrm{i} \hbar \psi
$$

This relation is true for any wavefunction $\psi$, so the commutator is\\
\includegraphics[max width=\textwidth, center]{2025_02_27_d4b95d9718f0b9e2e6b9g-291(1)}

The commutator in eqn 7C. 16 is of such central significance in quantum mechanics that it is taken as a fundamental distinction between classical mechanics and quantum mechanics. In fact, this commutator may be taken as a postulate of quantum mechanics and used to justify the choice of the operators for position and linear momentum in eqn 7C.3.

Classical mechanics supposed, falsely as is now known, that the position and momentum of a particle could be specified simultaneously with arbitrary precision. However, quantum mechanics shows that position and momentum are complementary, and that a choice must be made: position can be specified, but at the expense of momentum, or momentum can be specified, but at the expense of position.

\subsection*{7c. 4 The postulates of quantum mechanics}
The principles of quantum theory can be summarized as a series of postulates, which will form the basis for chemical applications of quantum mechanics throughout the text.

The wavefunction: All dynamical information is contained in the wavefunction $\psi$ for the system, which is a mathematical function found by solving the appropriate Schrödinger equation for the system.

The Born interpretation: If the wavefunction of a particle has the value $\psi$ at some position $r$, then the probability of finding the particle in an infinitesimal volume $\mathrm{d} \tau=\mathrm{d} x \mathrm{~d} y \mathrm{~d} z$ at that position is proportional to $|\psi|^{2} \mathrm{~d} \tau$.

Acceptable wavefunctions: An acceptable wavefunction must be single-valued, continuous, not infinite over a finite region of space, and (except in special cases) have a continuous slope.

Observables: Observables, $\Omega$, are represented by hermitian operators, $\hat{\Omega}$, built from the position and momentum operators specified in eqn 7C.3.

Observations and expectation values: A single measurement of the observable represented by the operator $\hat{\Omega}$ gives one of the eigenvalues of $\hat{\Omega}$. If the wavefunction is not an eigenfunction of $\hat{\Omega}$, the average of many measurements is given by the expectation value, $\langle\Omega\rangle$, defined in eqn 7C.11.

The Heisenberg uncertainty principle: It is impossible to specify simultaneously, with arbitrary precision, both the linear momentum and the position of a particle and, more generally, any pair of observables represented by operators that do not commute.

\subsection*{Checklist of concepts}
$\square$ 1. The Schrödinger equation is an eigenvalue equation.2. An operator carries out a mathematical operation on a function.3. The hamiltonian operator is the operator corresponding to the total energy of the system, the sum of the kinetic and potential energies.4. The wavefunction corresponding to a specific energy is an eigenfunction of the hamiltonian operator.5. Two different functions are orthogonal if the integral (over all space) of their product is zero.6. Hermitian operators have real eigenvalues and orthogonal eigenfunctions.7. Observables are represented by hermitian operators.8. Sets of functions that are normalized and mutually orthogonal are called orthonormal.9. When the system is not described by a single eigenfunction of an operator, it may be expressed as a superposition of such eigenfunctions.10. The mean value of a series of observations is given by the expectation value of the corresponding operator.11. The uncertainty principle restricts the precision with which complementary observables may be specified and measured simultaneously.\\
$\square$ 12. Complementary observables are observables for which the corresponding operators do not commute.

\subsection*{Checklist of equations}
\begin{center}
\begin{tabular}{llll}
\hline
Property & Equation & Comment & \begin{tabular}{c}
Equation \\
number \\
\end{tabular} \\
\hline
Eigenvalue equation & $\hat{\Omega} \psi=\omega \psi$ & $\psi$ eigenfunction; $\omega$ eigenvalue & 7C.2b \\
Hermiticity & $\int \psi_{i}^{*} \hat{\Omega}_{j} \mathrm{~d} \tau=\left\{\int \psi_{j}^{*} \hat{\Omega} \psi_{i} \mathrm{~d} \tau\right\}^{*}$ & \begin{tabular}{l}
Hermitian operators have real eigenvalues and orthogonal \\
eigenfunctions \\
\end{tabular} & 7C.7 \\
Orthogonality & $\int \psi_{i}^{*} \psi_{j} \mathrm{~d} \tau=0$ for $i \neq j$ & Integration over all space & Definition; assumes $\psi$ normalized \\
\end{tabular}
\end{center}

\section*{TOPIC 7D Translational motion}
\subsection*{Why do you need to know this material?}
The application of quantum theory to translational motion reveals the origin of quantization and non-classical features, such as tunnelling and zero-point energy. This material is important for the discussion of atoms and molecules that are free to move within a restricted volume, such as a gas in a container.

\subsection*{What is the key idea?}
The translational energy levels of a particle confined to a finite region of space are quantized, and under certain conditions particles can pass into and through classically forbidden regions.

\subsection*{What do you need to know already?}
You should know that the wavefunction is the solution of the Schrödinger equation (Topic 7B), and be familiar, in one instance, with the techniques of deriving dynamical properties from the wavefunction by using the operators corresponding to the observables (Topic 7C).

Translation, motion through space, is one of the basic types of motion. Quantum mechanics, however, shows that translation can have a number of non-classical features, such as its confinement to discrete energies and passage into and through classically forbidden regions.

\subsection*{7D. 1 Free motion in one dimension}
A free particle is unconstrained by any potential, which may be taken to be zero everywhere. In one dimension $V(x)=0$ everywhere, so the Schrödinger equation becomes (Topic 7B)

$$
-\frac{\hbar^{2}}{2 m} \frac{\mathrm{~d}^{2} \psi(x)}{\mathrm{d} x^{2}}=E \psi(x)
$$

Free motion in one dimension\\
(7D.1)\\
The most straightforward way to solve this simple secondorder differential equation is to take the known general form of solutions of equations of this kind, and then show that it does indeed satisfy eqn 7D.1.

How is that done? 7D. 1\\
Finding the solutions to the\\
Schrödinger equation for a free particle in one dimension\\
The general solution of a second-order differential equation of the kind shown in eqn 7D. 1 is

$$
\psi_{k}(x)=A \mathrm{e}^{\mathrm{i} k x}+B \mathrm{e}^{-\mathrm{i} k x}
$$

where $k, A$, and $B$ are constants. You can verify that $\psi_{k}(x)$ is a solution of eqn 7D. 1 by substituting it into the left-hand side of the equation, evaluating the derivatives, and then confirming that you have generated the right-hand side. Because $\mathrm{de}^{ \pm a x} / \mathrm{d} x= \pm a \mathrm{e}^{ \pm a x}$, the left-hand side becomes

$$
\begin{aligned}
-\frac{\hbar^{2}}{2 m} \frac{\mathrm{~d}^{2}}{\mathrm{~d} x^{2}} \overbrace{\left(A \mathrm{e}^{\mathrm{i} k x}+B \mathrm{e}^{-\mathrm{i} k x}\right.}^{\psi_{k}(x)} & =-\frac{\hbar^{2}}{2 m}\left\{A(\mathrm{i} k)^{2} \mathrm{e}^{\mathrm{i} k x}+B(-\mathrm{i} k)^{2} \mathrm{e}^{-\mathrm{i} k x}\right\} \\
& =\frac{\overbrace{k^{2} \hbar^{2}}^{2 m}}{E_{k}}(\overbrace{A \mathrm{e}^{\mathrm{i} k x}+B \mathrm{e}^{-\mathrm{i} k x}}^{\psi_{k}(x)})
\end{aligned}
$$

The left-hand side is therefore equal to a constant $\times \psi_{k}(x)$, which is the same as the term on the right-hand side of eqn 7D. 1 provided the constant, the term in blue, is identified with $E$. The value of the energy depends on the value of $k$, so henceforth it will be written $E_{k}$. The wavefunctions and energies of a free particle are therefore\\
\includegraphics[max width=\textwidth, center]{2025_02_27_d4b95d9718f0b9e2e6b9g-293}

The wavefunctions in eqn 7D. 2 are continuous, have continuous slope everywhere, are single-valued, and do not go to infinity: they are therefore acceptable wavefunctions for all values of $k$. Because $k$ can take any value, the energy can take any non-negative value, including zero. As a result, the translational energy of a free particle is not quantized.

In Topic 7C it is explained that in general a wavefunction can be written as a superposition (a linear combination) of the eigenfunctions of an operator. The wavefunctions of eqn 7D. 2 can be recognized as superpositions of the two functions $\mathrm{e}^{ \pm i k x}$ which are eigenfunctions of the linear momentum operator with eigenvalues $\pm k \hbar$ (Topic 7C). These eigenfunctions correspond to states with definite linear momentum:

$$
\psi_{k}(x)=\underbrace{A \mathrm{e}^{+\mathrm{i} k x}}_{\begin{array}{c}
\text { Particle with linear } \\
\text { momentum }+k \hbar
\end{array}}+\underbrace{B \mathrm{e}^{-\mathrm{i} k x}}_{\begin{array}{c}
\text { Particle with linear } \\
\text { momentum }-k \hbar
\end{array}}
$$

According to the interpretation given in Topic 7C, if a system is described by the wavefunction $\psi_{k}(x)$, then repeated measurements of the momentum will give $+k \hbar$ (that is, the particle travelling in the positive $x$-direction) with a probability proportional to $A^{2}$, and $-k \hbar$ (that is, the particle travelling in the negative $x$-direction) with a probability proportional to $B^{2}$. Only if $A$ or $B$ is zero does the particle have a definite momentum of $-k \hbar$ or $+k \hbar$, respectively.

\subsection*{Brief illustration 7D. 1}
Suppose an electron emerges from an accelerator moving towards positive $x$ with kinetic energy $1.0 \mathrm{eV}(1 \mathrm{eV}=1.602 \times$ $10^{-19} \mathrm{~J}$ ). The wavefunction for such a particle is given by eqn 7 D .2 with $B=0$ because the momentum is definitely in the positive $x$-direction. The value of $k$ is found by rearranging the expression for the energy in eqn 7D. 2 into

$$
\begin{aligned}
k & =\left(\frac{2 m_{\mathrm{e}} E_{k}}{\hbar^{2}}\right)^{1 / 2}=\left(\frac{2 \times\left(9.109 \times 10^{-31} \mathrm{~kg}\right) \times\left(1.6 \times 10^{-19} \mathrm{~J}\right)}{\left(1.055 \times 10^{-34} \mathrm{Js}\right)^{2}}\right)^{1 / 2} \\
& =5.1 \times 10^{9} \mathrm{~m}^{-1}
\end{aligned}
$$

or $5.1 \mathrm{~nm}^{-1}$ (with $1 \mathrm{~nm}=10^{-9} \mathrm{~m}$ ). Therefore, the wavefunction is $\psi(x)=A \mathrm{e}^{5.1 i x / n \mathrm{~m}}$.

So far, the motion of the particle has been confined to the $x$-axis. In general, the linear momentum is a vector (see The chemist's toolkit 17) directed along the line of travel of the particle. Then $\boldsymbol{p}=\boldsymbol{k} \hbar$ and the magnitude of the vector is $p=k \hbar$ and its component on each axis is $p_{q}=k_{q} \hbar$, with the wavefunction for each component proportional to $\mathrm{e}^{\mathrm{i} k_{q} q}$ with $q=x, y$, or $z$ and overall equal to $\mathrm{e}^{\mathrm{i}\left(k_{x} x+k_{y} y+k_{z} z\right)}$.

\subsection*{The chemist's toolkit 17 Vectors}
A vector is a quantity with both magnitude and direction. The vector $v$ shown in Sketch 1 has components on the $x, y$, and $z$ axes with values $v_{x}, v_{y}$, and $v_{z}$, respectively, which may be positive or negative. For example, if $v_{x}=-1.0$, the $x$-component of the vector $v$ has a magnitude of 1.0 and points in the $-x$ direction. The magnitude of a vector is denoted $v$ or $|v|$ and is given by

$$
v=\left(v_{x}^{2}+v_{y}^{2}+v_{z}^{2}\right)^{1 / 2}
$$

Thus, a vector with components $v_{x}=-1.0, v_{y}=+2.5$, and $v_{z}=+1.1$ has magnitude 2.9 and would be represented by an arrow of length 2.9 units and the appropriate orientation (as in the inset in the Sketch). Velocity and momentum are vectors; the magnitude of a velocity vector is called the speed. Force, too, is a vector. Electric and magnetic fields are two more examples of vectors.

\subsection*{7D. 2 Confined motion in one dimension}
Consider a particle in a box in which a particle of mass $m$ is confined to a region of one-dimensional space between two impenetrable walls. The potential energy is zero inside the box but rises abruptly to infinity at the walls located at $x=0$ and $x=L$ (Fig. 7D.1). When the particle is between the walls, the Schrödinger equation is the same as for a free particle (eqn 7D.1), so the general solutions given in eqn 7D. 2 are also the same. However, it will prove convenient to rewrite the wavefunction in terms of sines and cosines by using $\mathrm{e}^{ \pm i k x}=\cos k x \pm$ i $\sin k x$ (The chemist's toolkit 16 in Topic 7C)

$$
\begin{aligned}
\psi_{k}(x) & =A \mathrm{e}^{\mathrm{i} k x}+B \mathrm{e}^{-\mathrm{i} k x} \\
& =A(\cos k x+\mathrm{i} \sin k x)+B(\cos k x-\mathrm{i} \sin k x) \\
& =(A+B) \cos k x+\mathrm{i}(A-B) \sin k x
\end{aligned}
$$

\begin{center}
\includegraphics[max width=\textwidth]{2025_02_27_d4b95d9718f0b9e2e6b9g-294}
\end{center}

Figure 7D. 1 The potential energy for a particle in a onedimensional box. The potential is zero between $x=0$ and $x=L$, and then rises to infinity outside this region, resulting in impenetrable walls which confine the particle.\\
\includegraphics[max width=\textwidth, center]{2025_02_27_d4b95d9718f0b9e2e6b9g-294(1)}

Sketch 1

The operations involving vectors (addition, multiplication, etc.) needed for this text are described in The chemist's toolkit 22 in Topic 8C.

\footnotetext{${ }^{1}$ In terms of scalar products, this overall wavefunction would be written $\mathrm{e}^{\mathrm{i} k r}$.
}The constants $\mathrm{i}(A-\mathrm{B})$ and $A+B$ can be denoted $C$ and $D$, respectively, in which case


\begin{equation*}
\psi_{k}(x)=C \sin k x+D \cos k x \tag{7D.3}
\end{equation*}


Outside the box the wavefunctions must be zero as the particle will not be found in a region where its potential energy would be infinite:


\begin{equation*}
\text { For } x<0 \text { and } x>L, \psi_{k}(x)=0 \tag{7D.4}
\end{equation*}


\subsection*{(a) The acceptable solutions}
One of the requirements placed on a wavefunction is that it must be continuous. It follows that since the wavefunction is zero when $x<0$ (where the potential energy is infinite) the wavefunction must be zero at $x=0$ which is the point where the potential energy rises to infinity. Likewise, the wavefunction is zero where $x>L$ and so must be zero at $x=L$ where the potential energy also rises to infinity. These two restrictions are the boundary conditions, or constraints on the function:

$$
\psi_{k}(0)=0 \text { and } \psi_{k}(L)=0
$$

Boundary conditions\\
(7D.5)\\
Now it is necessary to show that the requirement that the wavefunction must satisfy these boundary conditions implies that only certain wavefunctions are acceptable, and that as a consequence only certain energies are allowed.

\subsection*{How is that done? 7D. 2 Showing that the boundary}
 conditions lead to quantized levelsYou need to start from the general solution and then explore the consequences of imposing the boundary conditions.

Step 1 Apply the boundary conditions\\
At $x=0, \psi_{k}(0)=C \sin 0+D \cos 0=D$ (because $\sin 0=0$ and $\cos 0=1)$. One boundary condition is $\psi_{k}(0)=0$, so it follows that $D=0$.

At $x=L, \psi_{k}(L)=C \sin k L$. The boundary condition $\psi_{k}(L)=0$ therefore requires that $\sin k L=0$, which in turn requires that $k L=n \pi$ with $n=1,2, \ldots$. Although $n=0$ also satisfies the boundary condition it is ruled out because the wavefunction would be $C \sin 0=0$ for all values of $x$, and the particle would be found nowhere. Negative integral values of $n$ also satisfy the boundary condition, but simply result in a change of sign of the wavefunction (because $\sin (-\theta)=-\sin \theta$ ). It therefore follows that the wavefunctions that satisfy the two boundary conditions are $\psi_{k}(x)=C \sin (n \pi x / L)$ with $n=1,2, \ldots$ and $k=n \pi / L$.

\subsection*{Step 2 Normalize the wavefunctions}
To normalize the wavefunction, write it as $N \sin (n \pi x / L)$ and require that the integral of the square of the wavefunction over all space is equal to 1 . The wavefunction is zero outside the range $0 \leq x \leq L$, so the integration needs to be carried out only inside this range:

$$
\int_{0}^{L} \psi^{2} \mathrm{~d} x=N^{2} \overbrace{\int_{0}^{L} \sin ^{2} \frac{n \pi x}{L} \mathrm{~d} x}^{\text {Integral } \mathrm{T} .2}=N^{2} \times \frac{L}{2}=1 \text {, so } N=\left(\frac{2}{L}\right)^{1 / 2}
$$

\subsection*{Step 3 Identify the allowed energies}
According to eqn 7D.2, $E_{k}=k^{2} \hbar^{2} / 2 m$, but because $k$ is limited to the values $k=n \pi / L$ with $n=1,2, \ldots$ the energies are restricted to the values

$$
E_{k}=\frac{k^{2} \hbar^{2}}{2 m}=\frac{(n \pi / L)^{2}(h / 2 \pi)^{2}}{2 m}=\frac{n^{2} h^{2}}{8 m L^{2}}
$$

At this stage it is sensible to replace the label $k$ by the label $n$, and to label the wavefunctions and energies as $\psi_{n}(x)$ and $E_{n}$. The allowed normalized wavefunctions and energies are therefore

$$
\left.-\psi_{n}(x)=\left(\frac{2}{L}\right)^{1 / 2} \sin \left(\frac{n \pi x}{L}\right) \quad E_{n}=\frac{n^{2} h^{2}}{8 m L^{2}} \quad n=1,2, \ldots\right) \left\lvert\, \begin{array}{|l}
\text { Particle } \\
\text { in a one- } \\
\text { dimensional } \\
\text { box }
\end{array}\right.
$$

The fact that $n$ is restricted to positive integer values implies that the energy of the particle in a one-dimensional box is quantized. This quantization arises from the boundary conditions that $\psi$ must satisfy. This is a general conclusion: the need to satisfy boundary conditions implies that only certain wavefunctions are acceptable, and hence restricts the eigenvalues to discrete values.

The integer $n$ that has been used to label the wavefunctions and energies is an example of a 'quantum number'. In general, a quantum number is an integer (in some cases, Topic 8B, a half-integer) that labels the state of the system. For a particle in a one-dimensional box there are an infinite number of acceptable solutions, and the quantum number $n$ specifies the one of interest (Fig. 7D.2). ${ }^{2}$ As well as acting as a label, a\\
\includegraphics[max width=\textwidth, center]{2025_02_27_d4b95d9718f0b9e2e6b9g-295}

Figure 7D. 2 The energy levels for a particle in a box. Note that the energy levels increase as $n^{2}$, and that their separation increases as the quantum number increases. Classically, the particle is allowed to have any value of the energy in the continuum shown as a tinted area.

\footnotetext{${ }^{2}$ You might object that the wavefunctions have a discontinuous slope at the edges of the box, and so do not qualify as acceptable according to the criteria in Topic 7B. This is a rare instance where the requirement does not apply, because the potential energy suddenly jumps to an infinite value.
}
quantum number can often be used to calculate the value of a property, such as the energy corresponding to the state, as in eqn 7D.6b.

\subsection*{(b) The properties of the wavefunctions}
Figure 7D. 3 shows some of the wavefunctions of a particle in a one-dimensional box. The points to note are as follows.

\begin{itemize}
  \item The wavefunctions are all sine functions with the same maximum amplitude but different wavelengths; the wavelength gets shorter as $n$ increases.
  \item Shortening the wavelength results in a sharper average curvature of the wavefunction and therefore an increase in the kinetic energy of the particle (recall that, as $V=0$ inside the box, the energy is entirely kinetic).
  \item The number of nodes (the points where the wavefunction passes through zero) also increases as $n$ increases; the wavefunction $\psi_{n}$ has $n-1$ nodes.
\end{itemize}

The probability density for a particle in a one-dimensional box is


\begin{equation*}
\psi_{n}^{2}(x)=\frac{2}{L} \sin ^{2}\left(\frac{n \pi x}{L}\right) \tag{7D.7}
\end{equation*}


and varies with position. The non-uniformity in the probability density is pronounced when $n$ is small (Fig. 7D.4). The maxima in the probability density give the locations at which the particle has the greatest probability of being found.\\
\includegraphics[max width=\textwidth, center]{2025_02_27_d4b95d9718f0b9e2e6b9g-296(2)}

Figure 7D. 3 The first five normalized wavefunctions of a particle in a box. As the energy increases the wavelength decreases, and successive functions possess one more half wave. The wavefunctions are zero outside the box.\\
\includegraphics[max width=\textwidth, center]{2025_02_27_d4b95d9718f0b9e2e6b9g-296}

Figure 7D. 4 (a) The first two wavefunctions for a particle in a box, (b) the corresponding probability densities, and (c) a representation of the probability density in terms of the darkness of shading.

\subsection*{Brief illustration 7D. 2}
As explained in Topic 7B, the total probability of finding the particle in a specified region is the integral of $\psi(x)^{2} \mathrm{~d} x$ over that region. Therefore, the probability of finding the particle with $n=1$ in a region between $x=0$ and $x=L / 2$ is

$$
\begin{aligned}
P=\int_{0}^{L / 2} \psi_{1}^{2} \mathrm{~d} x & =\frac{2}{L} \overbrace{\int_{0}^{L / 2} \sin ^{2}\left(\frac{\pi x}{L}\right) \mathrm{d} x}^{\text {Integral } \mathrm{T} .2}=\frac{2}{L}\left[\frac{x}{2}-\frac{1}{2 \pi / L} \sin \left(\frac{2 \pi x}{L}\right)\right]_{0}^{L / 2} \\
& =\frac{2}{L}(\frac{L}{4}-\frac{1}{2 \pi / L} \overbrace{\sin \pi}^{0})=\frac{1}{2}
\end{aligned}
$$

The result should not be a surprise, because the probability density is symmetrical around $x=L / 2$. The probability of finding the particle between $x=0$ and $x=L / 2$ must therefore be half of the probability of finding the particle between $x=0$ and $x=L$, which is 1 .\\
\includegraphics[max width=\textwidth, center]{2025_02_27_d4b95d9718f0b9e2e6b9g-296(1)}

Figure 7D. 5 The probability density $\psi^{2}(x)$ for large quantum number (here $n=50$, blue, compared with $n=1$, red). Notice that for high $n$ the probability density is nearly uniform, provided the fine detail of the increasingly rapid oscillations is ignored.

The probability density $\psi_{n}^{2}(x)$ becomes more uniform as $n$ increases provided the fine detail of the increasingly rapid oscillations is ignored (Fig. 7D.5). The probability density at high quantum numbers reflects the classical result that a particle bouncing between the walls spends equal times at all points. This conclusion is an example of the correspondence principle, which states that as high quantum numbers are reached, the classical result emerges from quantum mechanics.

\subsection*{(c) The properties of the energy}
The linear momentum of a particle in a box is not well defined because the wavefunction $\sin k x$ is not an eigenfunction of the linear momentum operator. However, because $\sin k x=\left(\mathrm{e}^{\mathrm{i} k x}-\mathrm{e}^{-\mathrm{i} k x}\right) / 2 \mathrm{i}$,


\begin{equation*}
\psi_{n}(x)=\left(\frac{2}{L}\right)^{1 / 2} \sin \left(\frac{n \pi x}{L}\right)=\frac{1}{2 \mathrm{i}}\left(\frac{2}{L}\right)^{1 / 2}\left(\mathrm{e}^{\mathrm{i} n \pi x / L}-\mathrm{e}^{-\mathrm{i} n \pi x / L}\right) \tag{7D.8}
\end{equation*}


It follows that, if repeated measurements are made of the linear momentum, half will give the value $+n \pi \hbar / L$ and half will give the value $-n \pi \hbar / L$. This conclusion is the quantum mechanical version of the classical picture in which the particle bounces back and forth in the box, spending equal times travelling to the left and to the right.

Because $n$ cannot be zero, the lowest energy that the particle may possess is not zero (as allowed by classical mechanics, corresponding to a stationary particle) but

$$
E_{1}=\frac{h^{2}}{8 m L^{2}}
$$

Zero-point energy\\
(7D.9)\\
This lowest, irremovable energy is called the zero-point energy. The physical origin of the zero-point energy can be explained in two ways:

\begin{itemize}
  \item The Heisenberg uncertainty principle states that $\Delta p_{x} \Delta x$ $\geq \frac{1}{2} \hbar$. For a particle confined to a box, $\Delta x$ has a finite value, therefore $\Delta p_{x}$ cannot be zero, as that would violate the uncertainty principle. Therefore the kinetic energy cannot be zero.
  \item If the wavefunction is to be zero at the walls, but smooth, continuous, and not zero everywhere, then it must be curved, and curvature in a wavefunction implies the possession of kinetic energy.
\end{itemize}

\subsection*{Brief illustration 7D. 3}
The lowest energy of an electron in a region of length 100 nm is given by eqn 7D. 6 with $n=1$ :

$$
E_{1}=\frac{(1)^{2} \times\left(6.626 \times 10^{-34} \mathrm{Js}\right)^{2}}{8 \times\left(9.109 \times 10^{-31} \mathrm{~kg}\right) \times\left(100 \times 10^{-9} \mathrm{~m}\right)^{2}}=6.02 \times 10^{-24} \mathrm{~J}
$$

where $1 \mathrm{~J}=1 \mathrm{~kg} \mathrm{~m}^{2} \mathrm{~s}^{-2}$ has been used. The energy $E_{1}$ can be expressed as $6.02 \mathrm{yJ}\left(1 \mathrm{yJ}=10^{-24} \mathrm{~J}\right)$.

The separation between adjacent energy levels with quantum numbers $n$ and $n+1$ is


\begin{equation*}
E_{n+1}-E_{n}=\frac{(n+1)^{2} h^{2}}{8 m L^{2}}-\frac{n^{2} h^{2}}{8 m L^{2}}=(2 n+1) \frac{h^{2}}{8 m L^{2}} \tag{7D.10}
\end{equation*}


This separation decreases as the length of the container increases, and is very small when the container has macroscopic dimensions. The separation of adjacent levels becomes zero when the walls are infinitely far apart. Atoms and molecules free to move in normal laboratory-sized vessels may therefore be treated as though their translational energy is not quantized.

\subsection*{Example 7D. 1}
\subsection*{Estimating an absorption wavelength}
$\beta$-Carotene (1) is a linear polyene in which 10 single and 11 double bonds alternate along a chain of 22 carbon atoms. If each CC bond length is taken to be 140 pm , the length of the molecular box in $\beta$-carotene is $L=2.94 \mathrm{~nm}$. Estimate the wavelength of the light absorbed by this molecule when it undergoes a transition from its ground state to the next higher excited state.\\
\includegraphics{smile-86e3ce9c92a3bb828d138210b7899ebc546e03f8}\\
$1 \beta$-Carotene

Collect your thoughts For reasons that will be familiar from introductory chemistry, each $\pi$-bonded $C$ atom contributes one p electron to the $\pi$-orbitals and two electrons occupy each state. Use eqn 7D. 10 to calculate the energy separation between the highest occupied and the lowest unoccupied levels, and convert that energy to a wavelength by using the Bohr frequency condition (eqn 7A.9, $\Delta E=h v$ ).

The solution There are 22 C atoms in the conjugated chain; each contributes one $p$ electron to the levels, so each level up to $n=11$ is occupied by two electrons. The separation in energy between the ground state and the state in which one electron is promoted from $n=11$ to $n=12$ is

$$
\begin{aligned}
\Delta E & =E_{12}-E_{11} \\
& =(2 \times 11+1) \frac{\left(6.626 \times 10^{-34} \mathrm{Js}\right)^{2}}{8 \times\left(9.109 \times 10^{-31} \mathrm{~kg}\right) \times\left(2.94 \times 10^{-9} \mathrm{~m}\right)^{2}} \\
& =1.60 \ldots \times 10^{-19} \mathrm{~J}
\end{aligned}
$$

or 0.160 aJ . It follows from the Bohr frequency condition ( $\Delta E=h v$ ) that the frequency of radiation required to cause this transition is

$$
v=\frac{\Delta E}{h}=\frac{1.60 \ldots \times 10^{-19} \mathrm{~J}}{6.626 \times 10^{-34} \mathrm{~J} \mathrm{~s}}=2.42 \times 10^{14} \mathrm{~s}^{-1}
$$

or $242 \mathrm{THz}\left(1 \mathrm{THz}=10^{12} \mathrm{~Hz}\right)$, corresponding to a wavelength $\lambda=1240 \mathrm{~nm}$. The experimental value is $603 \mathrm{THz}(\lambda=497 \mathrm{~nm})$, corresponding to radiation in the visible range of the electromagnetic spectrum.

Comment. The model is too crude to expect quantitative agreement, but the calculation at least predicts a wavelength in the right general range.

Self-test 7D. 1 Estimate a typical nuclear excitation energy in electronvolts ( $1 \mathrm{eV}=1.602 \times 10^{-19} \mathrm{~J} ; 1 \mathrm{GeV}=10^{9} \mathrm{eV}$ ) by calculating the first excitation energy of a proton confined to a one-dimensional box with a length equal to the diameter of a nucleus (approximately $1 \times 10^{-15} \mathrm{~m}$, or 1 fm ).\\
\includegraphics[max width=\textwidth, center]{2025_02_27_d4b95d9718f0b9e2e6b9g-298}

\subsection*{7D. 3 Confined motion in two and more dimensions}
Now consider a rectangular two-dimensional region, between 0 and $L_{1}$ along $x$, and between 0 and $L_{2}$ along $y$. Inside this region the potential energy is zero, but at the edges it rises to infinity (Fig. 7D.6). As in the one-dimensional case, the wavefunction can be expected to be zero at the edges of this region (at $x=0$ and $L_{1}$, and at $y=0$ and $L_{2}$ ), and to be zero outside the region. Inside the region the particle has contributions to its kinetic energy from its motion along both the $x$ and $y$ directions, and so the Schrödinger equation has two kinetic energy terms, one for each axis. For a particle of mass $m$ the equation is


\begin{equation*}
-\frac{\hbar^{2}}{2 m}\left(\frac{\partial^{2} \psi}{\partial x^{2}}+\frac{\partial^{2} \psi}{\partial y^{2}}\right)=E \psi \tag{7D.11}
\end{equation*}


Equation 7D. 11 is a partial differential equation, and the resulting wavefunctions are functions of both $x$ and $y$, denoted $\psi(x, y)$.\\
\includegraphics[max width=\textwidth, center]{2025_02_27_d4b95d9718f0b9e2e6b9g-298(1)}

Figure 7D. 6 A two-dimensional rectangular well. The potential goes to infinity at $x=0$ and $x=L_{1}$, and $y=0$ and $y=L_{2}$, but in between these values the potential is zero. The particle is confined to this rectangle by impenetrable walls.

\subsection*{(a) Energy levels and wavefunctions}
The procedure for finding the allowed wavefunctions and energies involves starting with the two-dimensional Schrödinger equation, and then applying the 'separation of variables' technique to turn it into two separate one-dimensional equations.

\subsection*{How is that done? 7D. 3 \\
 Constructing the wavefunctions}
 for a particle in a two-dimensional boxThe 'separation of variables' technique, which is explained and used here, is used in several cases in quantum mechanics.

\subsection*{Step 1 Apply the separation of variables technique}
First, recognize the presence of two operators, each of which acts on functions of only $x$ or $y$ :

$$
\hat{H}_{x}=-\frac{\hbar^{2}}{2 m} \frac{\partial^{2}}{\partial x^{2}} \quad \hat{H}_{y}=-\frac{\hbar^{2}}{2 m} \frac{\partial^{2}}{\partial y^{2}}
$$

Equation 7D.11, which is

$$
\overbrace{-\frac{\hbar^{2}}{2 m} \frac{\partial^{2}}{\partial x^{2}}}^{A_{x}} \overbrace{\psi-\frac{\hbar^{2}}{2 m} \frac{\partial^{2}}{\partial y^{2}}}^{A_{y}} \psi=E \psi
$$

then becomes

$$
\hat{H}_{x} \psi+\hat{H}_{y} \psi=E \psi
$$

Now suppose that the wavefunction $\psi$ can be expressed as the product of two functions, $\psi(x, y)=X(x) Y(y)$, one depending only on $x$ and the other depending only on $y$. This assumption is the central step of the procedure, and does not work for all partial differential equations: that it works here must be demonstrated. With this substitution the preceding equation becomes

$$
\hat{H}_{x} X(x) Y(y)+\hat{H}_{y} X(x) Y(y)=E X(x) Y(y)
$$

Then, because $H_{x}$ operates on (takes the second derivatives with respect to $x$ of) $X(x)$, and likewise for $H_{y}$ and $Y(y)$, this equation is the same as

$$
Y(y) \hat{H}_{x} X(x)+X(x) \hat{H}_{y} Y(y)=E X(x) Y(y)
$$

Division by both sides by $X(x) Y(y)$ then gives

$$
\overbrace{\frac{1}{X(x)} \hat{H}_{x} X(x)}^{\text {Depends only on } x}+\overbrace{\frac{1}{Y(y)} \hat{H}_{y} Y(y)}^{\text {Depends only on } y}=\overbrace{E}^{\text {A constant }}
$$

If $x$ is varied, only the first term can change; but the other two terms do not change, so the first term must be a constant for the equality to remain true. The same is true of the second term when $y$ is varied. Therefore, denoting these constants as $E_{X}$ and $E_{Y}$,

$$
\begin{aligned}
& \frac{1}{X(x)} \hat{H}_{x} X(x)=E_{X} \text {, so } \hat{H}_{x} X(x)=E_{X} X(x) \\
& \frac{1}{Y(y)} \hat{H}_{y} Y(y)=E_{Y} \text {, so } \hat{H}_{y} Y(y)=E_{Y} Y(y)
\end{aligned}
$$

with $E_{X}+E_{Y}=E$. The procedure has successfully separated the partial differential equation into two ordinary differential equations, one in $x$ and the other in $y$.

\subsection*{Step 2 Recognize the two ordinary differential equations}
Each of the two equations is identical to the Schrödinger equation for a particle in a one-dimensional box, one for the coordinate $x$ and the other for the coordinate $y$. The boundary conditions are also essentially the same (that the wavefunction must be zero at the walls). Consequently, the two solutions are

$$
\begin{array}{ll}
X_{n_{1}}(x)=\left(\frac{2}{L_{1}}\right)^{1 / 2} \sin \left(\frac{n_{1} \pi x}{L_{1}}\right) & E_{X, n_{1}}=\frac{n_{1}^{2} h^{2}}{8 m L_{1}^{2}} \\
Y_{n_{2}}(y)=\left(\frac{2}{L_{2}}\right)^{1 / 2} \sin \left(\frac{n_{2} \pi y}{L_{2}}\right) & E_{Y, n_{2}}=\frac{n_{2}^{2} h^{2}}{8 m L_{2}^{2}}
\end{array}
$$

with each of $n_{1}$ and $n_{2}$ taking the values $1,2, \ldots$ independently.\\
Step 3 Assemble the complete wavefunction\\
Inside the box, which is when $0 \leq x \leq L_{1}$ and $0 \leq x \leq L_{2}$, the wavefunction is the product $X_{n_{1}}(x) Y_{n_{2}}(y)$, and is given by eqn 7D.12a below. Outside the box, the wavefunction is zero. The energies are the sum $E_{X, n_{1}}+E_{Y, n_{2}}$. The two quantum numbers take the values $n_{1}=1,2, \ldots$ and $n_{2}=1,2, \ldots$ independently. Overall, therefore,

$$
\left\lvert\, \begin{array}{l|l}
\psi_{n_{1}, n_{2}}(x, y)=\frac{2}{\left(L_{1} L_{2}\right)^{1 / 2}} \sin \left(\frac{n_{1} \pi x}{L_{1}}\right) \sin \left(\frac{n_{2} \pi y}{L_{2}}\right) & \begin{array}{l}
\text { (7D.12a) } \\
\text { Wavefunctions } \\
\text { [two dimensions] }
\end{array} \\
E_{n_{1}, n_{2}}=\left(\frac{n_{1}^{2}}{L_{1}^{2}}+\frac{n_{2}^{2}}{L_{2}^{2}}\right) \frac{h^{2}}{8 m} & \begin{array}{l}
\text { Energy levels } \\
\text { [two dimensions] }
\end{array}
\end{array}\right.
$$

Some of the wavefunctions are plotted as contours in Fig. 7D.7. They are the two-dimensional versions of the wavefunctions shown in Fig. 7D.3. Whereas in one dimension the wavefunctions resemble states of a vibrating string with ends fixed, in two dimensions the wavefunctions correspond to vibrations of a rectangular plate with fixed edges.

\subsection*{Brief illustration 7D. 4}
Consider an electron confined to a square cavity of side $L$ (that is $L_{1}=L_{2}=L$ ), and in the state with quantum numbers $n_{1}=1$, $n_{2}=2$. Because the probability density is

$$
\psi_{1,2}^{2}(x, y)=\frac{4}{L^{2}} \sin ^{2}\left(\frac{\pi x}{L}\right) \sin ^{2}\left(\frac{2 \pi y}{L}\right)
$$

the most probable locations correspond to $\sin ^{2}(\pi x / L)=1$ and $\sin ^{2}(2 \pi y / L)=1$, or $(x, y)=(L / 2, L / 4)$ and $(L / 2,3 L / 4)$. The least probable locations (the nodes, where the wavefunction passes through zero) correspond to zeroes in the probability density within the box, which occur along the line $y=L / 2$.\\
\includegraphics[max width=\textwidth, center]{2025_02_27_d4b95d9718f0b9e2e6b9g-299}

Figure 7D. 7 The wavefunctions for a particle confined to a rectangular surface depicted as contours of equal amplitude. (a) $n_{1}=1, n_{2}=1$, the state of lowest energy; (b) $n_{1}=1, n_{2}=2$; (c) $n_{1}=2$, $n_{2}=1 ;$ (d) $n_{1}=2, n_{2}=2$.

A three-dimensional box can be treated in the same way: the wavefunctions are products of three terms and the energy is a sum of three terms. As before, each term is analogous to that for the one-dimensional case. Overall, therefore,

$$
\begin{array}{r}
\psi_{n_{1}, n_{2}, n_{3}}(x, y, z) \\
=\left(\frac{8}{L_{1} L_{2} L_{3}}\right)^{1 / 2} \sin \left(\frac{n_{1} \pi x}{L_{1}}\right) \sin \left(\frac{n_{2} \pi y}{L_{2}}\right) \sin \left(\frac{n_{3} \pi z}{L_{3}}\right) \\
\begin{array}{l}
\text { Wavefunctions } \\
\text { [three dimensions] }
\end{array}
\end{array}
$$

for $0 \leq x \leq L_{1}, 0 \leq y \leq L_{2}, 0 \leq z \leq L_{3}$

\[
E_{n_{1}, n_{2}, n_{3}}=\left(\frac{n_{1}^{2}}{L_{1}^{2}}+\frac{n_{2}^{2}}{L_{2}^{2}}+\frac{n_{3}^{2}}{L_{3}^{2}}\right) \frac{h^{2}}{8 m} \quad \begin{align*}
& \text { Energy levels }  \tag{7D.13b}\\
& \text { [three dimensions] }
\end{align*}
\]

The quantum numbers $n_{1}, n_{2}$, and $n_{3}$ are all positive integers 1 , $2, \ldots$ that can be chosen independently. The system has a zeropoint energy, the value of $E_{1,1,1}$.

\subsection*{(b) Degeneracy}
A special feature of the solutions arises when a two-dimensional box is not merely rectangular but square, with $L_{1}=L_{2}=$ $L$. Then the wavefunctions and their energies are

$$
\begin{gathered}
\psi_{n_{1}, n_{2}}(x, y)=\frac{2}{L} \sin \left(\frac{n_{1} \pi x}{L}\right) \sin \left(\frac{n_{2} \pi y}{L}\right) \\
\quad \text { for } 0 \leq x \leq L, 0 \leq y \leq L
\end{gathered}
$$

(7D.14a)

$$
\begin{array}{ll}
\psi_{n_{1}, n_{2}}(x, y)=0 \quad \text { outside box } & \\
E_{n_{1}, n_{2}}=\left(n_{1}^{2}+n_{2}^{2}\right) \frac{h^{2}}{8 m L^{2}} & \text { Energy levels } \\
\text { [square] }
\end{array}
$$

(7D.14b)\\
Consider the cases $n_{1}=1, n_{2}=2$ and $n_{1}=2, n_{2}=1$ :

$$
\begin{array}{ll}
\psi_{1,2}=\frac{2}{L} \sin \left(\frac{\pi x}{L}\right) \sin \left(\frac{2 \pi y}{L}\right) & E_{1,2}=\left(1^{2}+2^{2}\right) \frac{h^{2}}{8 m L^{2}}=\frac{5 h^{2}}{8 m L^{2}} \\
\psi_{2,1}=\frac{2}{L} \sin \left(\frac{2 \pi x}{L}\right) \sin \left(\frac{\pi y}{L}\right) & E_{2,1}=\left(2^{2}+1^{2}\right) \frac{h^{2}}{8 m L^{2}}=\frac{5 h^{2}}{8 m L^{2}}
\end{array}
$$

Although the wavefunctions are different, they correspond to the same energy. The technical term for different wavefunctions corresponding to the same energy is degeneracy, and in this case energy level $5 h^{2} / 8 m L^{2}$ is 'doubly degenerate'. In general, if $N$ wavefunctions correspond to the same energy, then that level is ' $N$-fold degenerate'.

The occurrence of degeneracy is related to the symmetry of the system. Figure 7D. 8 shows contour diagrams of the two degenerate functions $\psi_{1,2}$ and $\psi_{2,1}$. Because the box is square, one wavefunction can be converted into the other simply by rotating the plane by $90^{\circ}$. Interconversion by rotation through $90^{\circ}$ is not possible when the plane is not square, and $\psi_{1,2}$ and $\psi_{2,1}$ are then not degenerate. Similar arguments account for the degeneracy of the energy levels of a particle in a cubic box. Other examples of degeneracy occur in quantum mechanical systems (for instance, in the hydrogen atom, Topic 8A), and all of them can be traced to the symmetry properties of the system.

\subsection*{Brief illustration 7D. 5}
The energy of a particle in a two-dimensional square box of side $L$ in the energy level with $n_{1}=1, n_{2}=7$ is

$$
E_{1,7}=\left(1^{2}+7^{2}\right) \frac{h^{2}}{8 m L^{2}}=\frac{50 h^{2}}{8 m L^{2}}
$$

The level with $n_{1}=7$ and $n_{2}=1$ has the same energy. Thus, at first sight the energy level $50 h^{2} / 8 m L^{2}$ is doubly degenerate. However, in certain systems there may be levels that are not apparently related by symmetry but have the same energy and are said to be 'accidentally' degenerate. Such is the case here, for the level with $n_{1}=5$ and $n_{2}=5$ also has energy $50 h^{2} / 8 m L^{2}$. The level is therefore actually three-fold degenerate. Accidental degeneracy is also encountered in the hydrogen atom (Topic 8A) and can always be traced to a 'hidden' symmetry, one that is not immediately obvious.

\subsection*{7D. 4 Tunnelling}
A new quantum-mechanical feature appears when the potential energy does not rise abruptly to infinity at the walls (Fig. 7D.9). Consider the case in which there are two regions where\\
\includegraphics[max width=\textwidth, center]{2025_02_27_d4b95d9718f0b9e2e6b9g-300}

Figure 7D. 8 Two of the wavefunctions for a particle confined to a geometrically square well: (a) $n_{1}=2, n_{2}=1$; (b) $n_{1}=1, n_{2}=2$. The two functions correspond to the same energy and are said to be degenerate. Note that one wavefunction can be converted into the other by rotation of the box by $90^{\circ}$ : degeneracy is always a consequence of symmetry.\\
the potential energy is zero separated by a barrier where it rises to a finite value, $V_{0}$. Suppose the energy of the particle is less than $V_{0}$. A particle arriving from the left of the barrier has an oscillating wavefunction but inside the barrier the wavefunction decays rather than oscillates. Provided the barrier is not too wide the wavefunction emerges to the right, but with reduced amplitude; it then continues to oscillate once it is back in a region where it has zero potential energy. As a result of this behaviour the particle has a non-zero probability of passing through the barrier, which is forbidden classically because a particle cannot have a potential energy that exceeds its total energy. The ability of a particle to penetrate into, and possibly pass through, a classically forbidden region is called tunnelling.

The Schrödinger equation can be used to calculate the probability of tunnelling of a particle of mass $m$ incident from the left on a rectangular potential energy barrier of width $W$. On the left of the barrier $(x<0)$ the wavefunctions are those of a particle with $V=0$, so from eqn 7D.2,

$$
\begin{aligned}
\psi=A \mathrm{e}^{\mathrm{i} k x}+B \mathrm{e}^{-\mathrm{i} k x} \quad k \hbar= & (2 m E)^{1 / 2} \\
& \text { Wavefunction to left of barrier }
\end{aligned}
$$

(7D.15)\\
\includegraphics[max width=\textwidth, center]{2025_02_27_d4b95d9718f0b9e2e6b9g-300(1)}

Figure 7D. 9 The wavefunction for a particle encountering a potential barrier. Provided that the barrier is neither too wide nor too tall, the wavefunction will be non-zero as it exits to the right.

The Schrödinger equation for the region representing the barrier $(0 \leq x \leq W)$, where the potential energy has the constant value $V_{0}$, is


\begin{equation*}
-\frac{\hbar^{2}}{2 m} \frac{\mathrm{~d}^{2} \psi(x)}{\mathrm{d} x^{2}}+V_{0} \psi(x)=E \psi(x) \tag{7D.16}
\end{equation*}


Provided $E<V_{0}$ the general solutions of eqn 7D. 16 are


\begin{align*}
\psi=C \mathrm{e}^{\kappa x}+D \mathrm{e}^{-\kappa x} \quad \kappa \hbar= & \left\{2 m\left(V_{0}-E\right)\right\}^{1 / 2} \\
& \text { Wavefunction inside barrier } \tag{7D.17}
\end{align*}


as can be verified by substituting this solution into the lefthand side of eqn 7D.16. The important feature to note is that the two exponentials in eqn 7D. 17 are now real functions, as distinct from the complex, oscillating functions for the region where $V=0$. To the right of the barrier $(x>W)$, where $V=0$ again, the wavefunctions are

\[
\psi=A^{\prime} \mathrm{e}^{\mathrm{i} k x} \quad k \hbar=(2 m E)^{1 / 2} \quad \begin{align*}
& \text { Wavefunction to }  \tag{7D.18}\\
& \text { right of barrier }
\end{align*}
\]

Note that to the right of the barrier, the particle can be moving only to the right and therefore only the term $\mathrm{e}^{\mathrm{i} k x}$ contributes as it corresponds to a particle with positive linear momentum (moving to the right).

The complete wavefunction for a particle incident from the left consists of (Fig. 7D.10):

\begin{itemize}
  \item an incident wave ( $A \mathrm{e}^{\mathrm{i} k x}$ corresponds to positive linear momentum);
  \item a wave reflected from the barrier ( $B \mathrm{e}^{-\mathrm{i} k x}$ corresponds to negative linear momentum, motion to the left);
  \item the exponentially changing amplitude inside the barrier (eqn 7D.17);
  \item an oscillating wave (eqn 7D.18) representing the propagation of the particle to the right after tunnelling through the barrier successfully.\\
\includegraphics[max width=\textwidth, center]{2025_02_27_d4b95d9718f0b9e2e6b9g-301(1)}
\end{itemize}

Figure 7D. 10 When a particle is incident on a barrier from the left, the wavefunction consists of a wave representing linear momentum to the right, a reflected component representing momentum to the left, a varying but not oscillating component inside the barrier, and a (weak) wave representing motion to the right on the far side of the barrier.

The probability that a particle is travelling towards positive $x$ (to the right) on the left of the barrier $(x<0)$ is proportional to $|A|^{2}$, and the probability that it is travelling to the right after passing through the barrier $(x>W)$ is proportional to $\left|A^{\prime}\right|^{2}$. The ratio of these two probabilities, $\left|A^{\prime}\right|^{2} /|A|^{2}$, which expresses the probability of the particle tunnelling through the barrier, is called the transmission probability, $T$.

The values of the coefficients $A, B, C$, and $D$ are found by applying the usual criteria of acceptability to the wavefunction. Because an acceptable wavefunction must be continuous at the edges of the barrier (at $x=0$ and $x=W$ )


\begin{equation*}
\text { at } x=0: A+B=C+D \quad \text { at } x=W: C \mathrm{e}^{\kappa W}+D \mathrm{e}^{-\kappa W}=A^{\prime} \mathrm{e}^{\mathrm{i} k W} \tag{7D.19a}
\end{equation*}


Their slopes (their first derivatives) must also be continuous at these positions (Fig. 7D.11):


\begin{align*}
& \text { at } x=0: \mathrm{i} k A-\mathrm{i} k B=\kappa C-\kappa D \\
& \text { at } x=W: \kappa C \mathrm{e}^{\kappa W}-\kappa D \mathrm{e}^{-\kappa W}=\mathrm{i} k A^{\prime} \mathrm{e}^{\mathrm{i} k W} \tag{7D.19b}
\end{align*}


After straightforward but lengthy algebraic manipulations of these four equations 7D. 19 (see Problem P7D.12), the transmission probability turns out to be

$$
T=\left\{1+\frac{\left(\mathrm{e}^{\kappa W}-\mathrm{e}^{-\kappa W}\right)^{2}}{16 \varepsilon(1-\varepsilon)}\right\}^{-1} \quad \begin{aligned}
& \text { Transmission probability } \\
& \text { [rectangular barrier] }
\end{aligned}
$$

(7D.20a)\\
where $\varepsilon=E / V_{0}$. This function is plotted in Fig. 7D.12. The transmission probability for $E>V_{0}$ is shown there\\
\includegraphics[max width=\textwidth, center]{2025_02_27_d4b95d9718f0b9e2e6b9g-301}

Figure 7D. 11 The wavefunction and its slope must be continuous at the edges of the barrier. The conditions for continuity enable the wavefunctions at the junctions of the three zones to be connected and hence relations between the coefficients that appear in the solutions of the Schrödinger equation to be obtained.\\
\includegraphics[max width=\textwidth, center]{2025_02_27_d4b95d9718f0b9e2e6b9g-302}

Figure 7D. 12 The transmission probabilities $T$ for passage through a rectangular potential barrier. The horizontal axis is the energy of the incident particle expressed as a multiple of the barrier height. The curves are labelled with the value of $W\left(2 m V_{0}\right)^{1 / 2} / \hbar$. (a) $E<V_{0}$; (b) $E>V_{0}$.\\
too. The transmission probability has the following properties:

\begin{itemize}
  \item $T \approx 0$ for $E \ll V_{0}$ : there is negligible tunnelling when the energy of the particle is much lower than the height of the barrier;
  \item $T$ increases as $E$ approaches $V_{0}$ : the probability of tunnelling increases as the energy of the particle rises to match the height of the barrier;
  \item $T$ approaches 1 for $E>V_{0}$, but the fact that it does not immediately reach 1 means that there is a probability of the particle being reflected by the barrier even though according to classical mechanics it can pass over it;
  \item $T \approx 1$ for $E \gg V_{0}$, as expected classically: the barrier is invisible to the particle when its energy is much higher than the barrier.
\end{itemize}

For high, wide barriers (in the sense that $\kappa W \gg 1$ ), eqn 7D.20a simplifies to

$$
T \approx 16 \varepsilon(1-\varepsilon) \mathrm{e}^{-2 \kappa W}
$$

Rectangular potential barrier; $\kappa W \gg 1$\\
(7D.20b)

The transmission probability decreases exponentially with the thickness of the barrier and with $m^{1 / 2}$ (because $\kappa \propto m^{1 / 2}$ ). It follows that particles of low mass are more able to tunnel through barriers than heavy ones (Fig. 7D.13). Tunnelling is very important for electrons and muons ( $m_{\mu} \approx 207 m_{\mathrm{e}}$ ), and moderately important for protons ( $m_{\mathrm{p}} \approx 1840 m_{\mathrm{e}}$ ); for heavier particles it is less important.

A number of effects in chemistry depend on the ability of the proton to tunnel more readily than the deuteron. The very rapid equilibration of proton transfer reactions is also a mani-\\
\includegraphics[max width=\textwidth, center]{2025_02_27_d4b95d9718f0b9e2e6b9g-302(1)}

Figure 7D. 13 The wavefunction of a heavy particle decays more rapidly inside a barrier than that of a light particle. Consequently, a light particle has a greater probability of tunnelling through the barrier.\\
festation of the ability of protons to tunnel through barriers and transfer quickly from an acid to a base. Tunnelling of protons between acidic and basic groups is also an important feature of the mechanism of some enzyme-catalysed reactions.

\subsection*{Brief illustration 7D. 6}
Suppose that a proton of an acidic hydrogen atom is confined to an acid that can be represented by a barrier of height 2.000 eV and length 100 pm . The probability that a proton with energy 1.995 eV (corresponding to 0.3195 aJ ) can escape from the acid is computed using eqn 7D.20a, with $\varepsilon=E / V_{0}=$ $1.995 \mathrm{eV} / 2.000 \mathrm{eV}=0.9975$ and $V_{0}-E=0.005 \mathrm{eV}$ (corresponding to $8.0 \times 10^{-22} \mathrm{~J}$ ). The quantity $\kappa$ is given by eqn 7D.17:

$$
\begin{aligned}
\kappa & =\frac{\left\{2 \times\left(1.67 \times 10^{-27} \mathrm{~kg}\right) \times\left(8.0 \times 10^{-22} \mathrm{~J}\right)\right\}^{1 / 2}}{1.055 \times 10^{-34} \mathrm{Js}} \\
& =1.54 \cdots \times 10^{10} \mathrm{~m}^{-1}
\end{aligned}
$$

It follows that

$$
\kappa W=\left(1.54 \ldots \times 10^{10} \mathrm{~m}^{-1}\right) \times\left(100 \times 10^{-12} \mathrm{~m}\right)=1.54 \ldots
$$

Equation 7D.20a then yields

$$
\begin{aligned}
T & =\left\{1+\frac{\left(\mathrm{e}^{1.54 \ldots}-\mathrm{e}^{-1.54 \ldots}\right)^{2}}{16 \times 0.9975 \times(1-0.9975)}\right\}^{-1} \\
& =1.97 \times 10^{-3}
\end{aligned}
$$

A problem related to tunnelling is that of a particle in a square-well potential of finite depth (Fig. 7D.14). Inside the well the potential energy is zero and the wavefunctions oscillate as they do for a particle in an infinitely deep box. At the edges, the potential energy rises to a finite value $V_{0}$. If $E<V_{0}$ the wavefunction decays as it penetrates into the walls, just as it does when it enters a barrier. The wavefunctions are found by ensuring, as in the discussion of the potential barrier, that they and their slopes are continuous at the edges of the potential. The two lowest energy solutions are shown in Fig. 7D.15.\\
\includegraphics[max width=\textwidth, center]{2025_02_27_d4b95d9718f0b9e2e6b9g-303}

Figure 7D. 14 A potential well with a finite depth.

For a potential well of finite depth, there are a finite number of wavefunctions with energy less than $V_{0}$ : they are referred to as bound states, in the sense that the particle is mainly confined to the well. Detailed consideration of the Schrödinger equation for the problem shows that the number of bound states is equal to $N$, with


\begin{equation*}
N-1<\frac{\left(8 m V_{0} L\right)^{1 / 2}}{h}<N \tag{7D.21}
\end{equation*}


\begin{center}
\includegraphics[max width=\textwidth]{2025_02_27_d4b95d9718f0b9e2e6b9g-303(1)}
\end{center}

Figure 7D. 15 Wavefunctions of the lowest two bound levels for a particle in the potential well shown in Fig. 7D.14.\\
where $V_{0}$ is the depth of the well and $L$ is its width. This relation shows that the deeper and wider the well, the greater the number of bound states. As the depth becomes infinite, so the number of bound states also becomes infinite, as for the particle in a box treated earlier in this Topic.

\subsection*{Checklist of concepts}
$\square$ 1. The translational energy of a free particle is not quantized.2. The need to satisfy boundary conditions implies that only certain wavefunctions are acceptable and restricts observables, specifically the energy, to discrete values.3. A quantum number is an integer (in certain cases, a half-integer) that labels the state of the system.4. A particle in a box possesses a zero-point energy, an irremovable minimum energy.5. The correspondence principle states that the quantum mechanical result with high quantum numbers should agree with the predictions of classical mechanics.6. The wavefunction for a particle in a two- or threedimensional box is the product of wavefunctions for the particle in a one-dimensional box.\\
$\square$ 7. The energy of a particle in a two- or three-dimensional box is the sum of the energies for the particle in two or three one-dimensional boxes.8. Energy levels are $N$-fold degenerate if $N$ wavefunctions correspond to the same energy.\\
9. The occurrence of degeneracy is a consequence of the symmetry of the system.10. Tunnelling is penetration into or through a classically forbidden region.11. The probability of tunnelling decreases with an increase in the height and width of the potential barrier.12. Light particles are more able to tunnel through barriers than heavy ones.

\subsection*{Checklist of equations}
\begin{center}
\begin{tabular}{|c|c|c|c|}
\hline
Property & Equation & Comment & Equation number \\
\hline
Free-particle wavefunctions and energies & $\psi_{k}=A \mathrm{e}^{\mathrm{i} k x}+B \mathrm{e}^{-\mathrm{i} k x} \quad E_{k}=k^{2} \hbar^{2} / 2 m$ & All values of $k$ allowed & 7D. 2 \\
\hline
\multicolumn{4}{|l|}{Particle in a box} \\
\hline
\multicolumn{4}{|l|}{One dimension:} \\
\hline
Wavefunctions & \( \begin{aligned} & \psi_{n}(x)=(2 / L)^{1 / 2} \sin (n \pi x / L), \quad 0 \leq x \leq L \\ & \psi_{n}(x)=0, \quad x<0 \text { and } x>L \end{aligned} \) & $n=1,2, \ldots$ & 7D. 6 \\
\hline
Energies & $E_{n}=n^{2} h^{2} / 8 m L^{2}$ &  &  \\
\hline
\multicolumn{4}{|l|}{Two dimensions:} \\
\hline
Wavefunctions & \( \begin{aligned} & \psi_{n_{1}, n_{2}}(x, y)=\psi_{n_{1}}(x) \psi_{n_{2}}(y) \\ & \psi_{n_{1}}(x)=\left(2 / L_{1}\right)^{1 / 2} \sin \left(n_{1} \pi x / L_{1}\right), \quad 0 \leq x \leq L_{1} \\ & \psi_{n_{2}}(y)=\left(2 / L_{2}\right)^{1 / 2} \sin \left(n_{2} \pi y / L_{2}\right), \quad 0 \leq y \leq L_{2} \end{aligned} \) & $n_{1}, n_{2}=1,2, \ldots$ & 7D.12a \\
\hline
Energies & $E_{n_{1}, n_{2}}=\left(n_{1}^{2} / L_{1}^{2}+n_{2}^{2} / L_{2}^{2}\right) h^{2} / 8 m$ &  & 7D.12b \\
\hline
\multicolumn{4}{|l|}{Three dimensions:} \\
\hline
Wavefunctions & $\psi_{n_{1}, n_{2}, n_{3}}(x, y, z)=\psi_{n_{1}}(x) \psi_{n_{2}}(y) \psi_{n_{3}}(z)$ & $n_{1}, n_{2}, n_{3}=1,2, \ldots$ & 7D.13a \\
\hline
Energies & $E_{n_{1}, n_{2}, n_{3}}=\left(n_{1}^{2} / L_{1}^{2}+n_{2}^{2} / L_{2}^{2}+n_{3}^{2} / L_{3}^{2}\right) h^{2} / 8 m$ &  & 7D.13b \\
\hline
\multirow[t]{2}{*}{Transmission probability} & $T=\left\{1+\left(\mathrm{e}^{\kappa W}-\mathrm{e}^{-\kappa W}\right)^{2} / 16 \varepsilon(1-\varepsilon)\right\}^{-1}$ & Rectangular potential barrier & 7D.20a \\
\hline
 & $T=16 \varepsilon(1-\varepsilon) \mathrm{e}^{-2 \kappa W}$ & High, wide rectangular barrier & 7D.20b \\
\hline
\end{tabular}
\end{center}

\section*{TOPIC 7E Vibrational motion}
\subsection*{Why do you need to know this material?}
Molecular vibration plays a role in the interpretation of thermodynamic properties, such as heat capacities (Topics 2 A and 13 E ), and of the rates of chemical reactions (Topic 18C). The measurement and interpretation of the vibrational frequencies of molecules is the basis of infrared spectroscopy (Topics 11C and 11D).

\subsection*{What is the key idea?}
The energy of vibrational motion is quantized.

\subsection*{What do you need to know already?}
You should know how to formulate the Schrödinger equation for a given potential energy. You should also be familiar with the concepts of tunnelling (Topic 7D) and the expectation value of an observable (Topic 7C).

Atoms in molecules and solids vibrate around their equilibrium positions as bonds stretch, compress, and bend. The simplest model for this kind of motion is the 'harmonic oscillator', which is considered in detail in this Topic.

\subsection*{The chemist's toolkit 18}
\subsection*{The classical harmonic oscillator}
A harmonic oscillator consists of a particle of mass $m$ that experiences a 'Hooke's law' restoring force, one that is proportional to the displacement of the particle from equilibrium. For a one-dimensional system,

$$
F_{x}=-k_{\mathrm{f}} x
$$

From Newton's second law of motion $\left(F=m a=m\left(\mathrm{~d}^{2} x / \mathrm{d} t^{2}\right)\right.$; see The chemist's toolkit 3 in Topic 1B),

$$
m \frac{\mathrm{~d}^{2} x}{\mathrm{~d} t^{2}}=-k_{\mathrm{f}} x
$$

If $x=0$ at $t=0$, a solution (as may be verified by substitution) is

$$
x(t)=A \sin 2 \pi v t \quad v=\frac{1}{2 \pi}\left(\frac{k_{\mathrm{f}}}{m}\right)^{1 / 2}
$$

This solution shows that the position of the particle oscillates harmonically (i.e. as a sine function) with frequency $v$ (units: Hz ). The angular frequency of the oscillator is $\omega=2 \pi \nu$ (units: radians per second). It follows that the angular frequency of a classical harmonic oscillator is $\omega=\left(k_{f} / m\right)^{1 / 2}$.

\subsection*{7E. 1 The harmonic oscillator}
In classical mechanics a harmonic oscillator is a particle of mass $m$ that experiences a restoring force proportional to its displacement, $x$, from the equilibrium position. As is shown in The chemist's toolkit 18, the particle oscillates about the equilibrium position at a characteristic frequency, $v$. The potential energy of the particle is


\begin{equation*}
V(x)=\frac{1}{2} k_{\mathrm{f}} x^{2} \quad \text { Parabolic potential energy } \tag{7E.1}
\end{equation*}


where $k_{\mathrm{f}}$ is the force constant, which characterizes the strength of the restoring force (Fig. 7E.1) and is expressed in newtons per metre $\left(\mathrm{Nm}^{-1}\right)$. This form of potential energy is called a 'harmonic potential energy' or a 'parabolic potential energy'. The Schrödinger equation for the oscillator is therefore

\[
-\frac{\hbar^{2}}{2 m} \frac{\mathrm{~d}^{2} \psi(x)}{\mathrm{d} x^{2}}+\frac{1}{2} k_{\mathrm{f}} x^{2} \psi(x)=E \psi(x) \quad \begin{align*}
& \text { Schrödinger }  \tag{7E.2}\\
& \text { equation }
\end{align*}
\]

The potential energy becomes infinite at $x= \pm \infty$, and so the wavefunction is zero at these limits. However, as the potential energy rises smoothly rather than abruptly to infinity, as it does for a particle in a box, the wavefunction decreases

The potential energy $V$ is related to force by $F=-\mathrm{d} V / \mathrm{d} x$ (The chemist's toolkit 6 in Topic 2A), so the potential energy corresponding to a Hooke's law restoring force is

$$
V(x)=\frac{1}{2} k_{\mathrm{f}} x^{2}
$$

As the particle moves away from the equilibrium position its potential energy increases and so its kinetic energy, and hence its speed, decreases. At some point all the energy is potential and the particle comes to rest at a turning point. The particle then accelerates back towards and through the equilibrium position. The greatest probability of finding the particle is where it is moving most slowly, which is close to the turning points.\\
The turning point, $x_{\text {tp }}$, of a classical oscillator occurs when its potential energy $\frac{1}{2} k_{\mathrm{f}} x^{2}$ is equal to its total energy, so

$$
x_{\mathrm{tp}}= \pm\left(\frac{2 E}{k_{\mathrm{f}}}\right)^{1 / 2}
$$

The turning point increases with the total energy: in classical terms, the amplitude of the swing of a pendulum or the displacement of a mass on a spring increases.\\
\includegraphics[max width=\textwidth, center]{2025_02_27_d4b95d9718f0b9e2e6b9g-306(1)}

Figure 7E. 1 The potential energy for a harmonic oscillator is the parabolic function $V_{\text {HO }}(x)=\frac{1}{2} k_{Y} x^{2}$, where $x$ is the displacement from equilibrium. The larger the force constant $k_{\mathrm{f}}$ the steeper the curve and narrower the curve becomes.\\
smoothly towards zero rather than becoming zero abruptly. The boundary conditions $\psi( \pm \infty)=0$ imply that only some solutions of the Schrödinger equation are acceptable, and therefore that the energy of the oscillator is quantized.

\subsection*{(a) The energy levels}
Equation 7E. 2 is a standard form of differential equation and its solutions are well known to mathematicians. ${ }^{1}$ The energies permitted by the boundary conditions are

\[
\begin{array}{cl}
E_{v}=\left(v+\frac{1}{2}\right) \hbar \omega & \omega=\left(k_{\mathrm{f}} / m\right)^{1 / 2} \quad \text { Energy levels }  \tag{7E.3}\\
& v=0,1,2, \ldots
\end{array}
\]

where $v$ is the vibrational quantum number. Note that the energies depend on $\omega$, which has the same dependence on the mass and the force constant as the angular frequency of a classical oscillator (see The chemist's toolkit 18) and is high when the force constant is large and the mass small. The separation of adjacent levels is


\begin{equation*}
E_{\nu+1}-E_{v}=\hbar \omega \tag{7E.4}
\end{equation*}


for all $v$. The energy levels therefore form a uniform ladder with spacing $\hbar \omega$ (Fig. 7E.2). The energy separation $\hbar \omega$ is negligibly small for macroscopic objects (with large mass) but significant for objects with mass similar to that of an atom.

The energy of the lowest level, with $v=0$, is not zero:


\begin{equation*}
E_{0}=\frac{1}{2} \hbar \omega \quad \text { Zero-point energy } \tag{7E.5}
\end{equation*}


\footnotetext{${ }^{1}$ For details, see our Molecular quantum mechanics, Oxford University Press, Oxford (2011).
}
\includegraphics[max width=\textwidth, center]{2025_02_27_d4b95d9718f0b9e2e6b9g-306}

Figure 7E. 2 The energy levels of a harmonic oscillator are evenly spaced with separation $\hbar \omega$, where $\omega=\left(k_{f} / m\right)^{1 / 2}$. Even in its lowest energy state, an oscillator has an energy greater than zero.

The physical reason for the existence of this zero-point energy is the same as for the particle in a box (Topic 7D). The particle is confined, so its position is not completely uncertain. It follows that its momentum, and hence its kinetic energy, cannot be zero. A classical interpretation of the zero-point energy is that the quantum oscillator is never completely at rest and therefore has kinetic energy; moreover, because its motion samples the potential energy away from the equilibrium position, it also has non-zero potential energy.

The model of a particle oscillating in a parabolic potential is used to describe the vibrational motion of a diatomic molecule A-B (and, with elaboration, Topic 11D, polyatomic molecules). In this case both atoms move as the bond between them is stretched and compressed and the mass $m$ is replaced by the effective mass, $\mu$, given by

$$
\mu=\frac{m_{\mathrm{A}} m_{\mathrm{B}}}{m_{\mathrm{A}}+m_{\mathrm{B}}} \quad \begin{aligned}
& \text { Effective mass } \\
& \text { [diatomic molecule] }
\end{aligned}
$$

(7E.6)\\
When A is much heavier than $\mathrm{B}, m_{\mathrm{B}}$ can be neglected in the denominator and the effective mass is $\mu \approx m_{\mathrm{B}}$, the mass of the lighter atom. In this case, only the light atom moves and the heavy atom acts as a stationary anchor.

\subsection*{Brief illustration 7E. 1}
The effective mass of ${ }^{1} \mathrm{H}^{35} \mathrm{Cl}$ is

$$
\mu=\frac{m_{\mathrm{H}} m_{\mathrm{Cl}}}{m_{\mathrm{H}}+m_{\mathrm{Cl}}}=\frac{\left(1.0078 m_{\mathrm{u}}\right) \times\left(34.9688 m_{\mathrm{u}}\right)}{\left(1.0078 m_{\mathrm{u}}\right)+\left(34.9688 m_{\mathrm{u}}\right)}=0.9796 m_{\mathrm{u}}
$$

which is close to the mass of the hydrogen atom. The force constant of the bond is $k_{\mathrm{f}}=516.3 \mathrm{Nm}^{-1}$. It follows from eqn 7 E .3 and $1 \mathrm{~N}=1 \mathrm{~kg} \mathrm{~s} \mathrm{~s}^{-2}$, with $\mu$ in place of $m$, that

$$
\omega=\left(\frac{k_{\mathrm{f}}}{\mu}\right)^{1 / 2}=\left(\frac{516.3 \mathrm{Nm}^{-1}}{0.9796 \times\left(1.66054 \times 10^{-27} \mathrm{~kg}\right)}\right)^{1 / 2}=5.634 \times 10^{14} \mathrm{~s}^{-1}
$$

or (after division by $2 \pi$ ) 89.67 THz . Therefore, the separation of adjacent levels is (eqn 7E.4)

$$
\begin{aligned}
E_{v+1}-E_{v} & =\left(1.05457 \times 10^{-34} \mathrm{~J} \mathrm{~s}\right) \times\left(5.634 \times 10^{14} \mathrm{~s}^{-1}\right) \\
& =5.941 \times 10^{-20} \mathrm{~J}
\end{aligned}
$$

or 59.41 zJ , about 0.37 eV . This energy separation corresponds to $36 \mathrm{~kJ} \mathrm{~mol}^{-1}$, which is chemically significant. The zero-point energy (eqn 7E.5) of this molecular oscillator is 29.71 zJ , which corresponds to 0.19 eV , or $18 \mathrm{~kJ} \mathrm{~mol}^{-1}$.

\subsection*{(b) The wavefunctions}
The acceptable solutions of eqn 7E.2, all have the form

$$
\begin{aligned}
& \psi(x)=N \times(\text { polynomial in } x) \times(\text { bell-shaped Gaussian } \\
& \text { function })
\end{aligned}
$$

where $N$ is a normalization constant. A Gaussian function is a bell-shaped function of the form $\mathrm{e}^{-x^{2}}$ (Fig. 7E.3). The precise form of the wavefunctions is

$$
\begin{aligned}
& \psi_{v}(x)=N_{v} H_{v}(y) \mathrm{e}^{-y^{2} / 2} \\
& y=\frac{x}{\alpha} \quad \alpha=\left(\frac{\hbar^{2}}{m k_{\mathrm{f}}}\right)^{1 / 4}
\end{aligned}
$$

Wavefunctions (7E.7)

The factor $H_{v}(y)$ is a Hermite polynomial; the form of these polynomials and some of their properties are listed in Table 7E.1. Note that the first few Hermite polynomials are rather simple: for instance, $H_{0}(y)=1$ and $H_{1}(y)=2 y$. Hermite polynomials, which are members of a class of functions called 'orthogonal polynomials', have a wide range of important properties which allow a number of quantum mechanical calculations to be done with relative ease.

The wavefunction for the ground state, which has $v=0$, is

$$
\psi_{0}(x)=N_{0} \mathrm{e}^{-y^{2} / 2}=N_{0} \mathrm{e}^{-x^{2} / 2 \alpha^{2}} \quad \begin{aligned}
& \text { Ground-state } \\
& \text { wavefunction }
\end{aligned}
$$

(7E.8a)\\
\includegraphics[max width=\textwidth, center]{2025_02_27_d4b95d9718f0b9e2e6b9g-307(1)}

Figure 7E. 3 The graph of the Gaussian function, $f(x)=\mathrm{e}^{-x^{2}}$.

Table 7E. 1 The Hermite polynomials

\begin{center}
\begin{tabular}{ll}
\hline
$\boldsymbol{v}$ & $\boldsymbol{H}_{v}(\boldsymbol{y})$ \\
\hline
0 & 1 \\
1 & $2 y$ \\
2 & $4 y^{2}-2$ \\
3 & $8 y^{3}-12 y$ \\
4 & $16 y^{4}-48 y^{2}+12$ \\
5 & $32 y^{5}-160 y^{3}+120 y$ \\
6 & $64 y^{6}-480 y^{4}+720 y^{2}-120$ \\
\hline
\end{tabular}
\end{center}

The Hermite polynomials are solutions of the differential equation

$$
H_{v}^{\prime \prime}-2 y H_{v}^{\prime}+2 v H_{v}=0
$$

where primes denote differentiation. They satisfy the recursion relation

$$
H_{v+1}-2 y H_{v}+2 \nu H_{v-1}=0
$$

An important integral is

$$
\int_{-\infty}^{\infty} H_{v^{\prime}} H_{v} \mathrm{e}^{-y^{2}} \mathrm{~d} y=\left\{\begin{array}{lr}
0 & \text { if } v^{\prime} \neq v \\
\pi^{1 / 2} 2^{v} v!\text { if } v^{\prime}=v
\end{array}\right.
$$

and the corresponding probability density is

\[
\psi_{0}^{2}(x)=N_{0}^{2} \mathrm{e}^{-y^{2}}=N_{0}^{2} \mathrm{e}^{-x^{2} / \alpha^{2}} \quad \begin{array}{ll}
\text { Ground-state }  \tag{7E.8b}\\
\text { probability density }
\end{array}
\]

The wavefunction and the probability density are shown in Fig. 7E.4. The probability density has its maximum value at $x=0$, the equilibrium position, but is spread about this position. The curvature is consistent with the kinetic energy being non-zero and the spread is consistent with the potential energy also being non-zero, so resulting in a zero-point energy.

The wavefunction for the first excited state, $v=1$, is

$$
\psi_{1}(x)=N_{1} 2 y \mathrm{e}^{-y^{2} / 2}=N_{1}\left(\frac{2}{\alpha}\right) x \mathrm{e}^{-x^{2} / 2 \alpha^{2}} \quad \begin{aligned}
& \text { First excited-state } \\
& \text { wavefunction }
\end{aligned}
$$

(7E.9)\\
\includegraphics[max width=\textwidth, center]{2025_02_27_d4b95d9718f0b9e2e6b9g-307}

Figure 7E. 4 The normalized wavefunction and probability density (shown also by shading) for the lowest energy state of a harmonic oscillator.\\
\includegraphics[max width=\textwidth, center]{2025_02_27_d4b95d9718f0b9e2e6b9g-308}

Figure 7E. 5 The normalized wavefunction and probability density (shown also by shading) for the first excited state of a harmonic oscillator.

This function has a node at zero displacement ( $x=0$ ), and the probability density has maxima at $x= \pm \alpha$ (Fig. 7E.5).

\subsection*{Example 7E. 1 Confirming that a wavefunction is a solution of the Schrödinger equation}
Confirm that the ground-state wavefunction (eqn 7E.8a) is a solution of the Schrödinger equation (eqn 7E.2).

Collect your thoughts You need to substitute the wavefunction given in eqn 7E.8a into eqn 7E. 2 and see that the lefthand side generates the right-hand side of the equation; use the definition of $\alpha$ in eqn 7E.7. Confirm that the factor that multiplies the wavefunction on the right-hand side agrees with eqn 7E.5.

The solution First, evaluate the second derivative of the ground-state wavefunction by differentiating it twice in succession:

$$
\begin{aligned}
& \frac{\mathrm{d}}{\mathrm{~d} x} N_{0} \mathrm{e}^{-x^{2} / 2 \alpha^{2}}=-N_{0}\left(\frac{x}{\alpha^{2}}\right) \mathrm{e}^{-x^{2} / 2 \alpha^{2}} \\
& \begin{aligned}
\frac{\mathrm{d}^{2}}{\mathrm{~d} x^{2}} N_{0} \mathrm{e}^{-x^{2} / 2 \alpha^{2}} & =\frac{\mathrm{d}}{\mathrm{~d} x}\{\overbrace{-N_{0}\left(\frac{x}{\alpha^{2}}\right)}^{f} \overbrace{\mathrm{e}^{-x^{2} / 2 \alpha^{2}}}^{g}\} \\
& \overbrace{\mathrm{d}(f g) / \mathrm{d} x=f \mathrm{~d} g / \mathrm{d} x+g \mathrm{~d} f / \mathrm{d} x} \\
& =-\frac{N_{0}}{\alpha^{2}} \mathrm{e}^{-x^{2} / 2 \alpha^{2}}+N_{0}\left(\frac{x}{\alpha^{2}}\right)^{2} \mathrm{e}^{-x^{2} / 2 \alpha^{2}} \\
& =-\left(1 / \alpha^{2}\right) \psi_{0}+\left(x^{2} / \alpha^{4}\right) \psi_{0}
\end{aligned}
\end{aligned}
$$

Now substitute this expression and $\alpha^{2}=\left(\hbar^{2} / m k_{\mathrm{f}}\right)^{1 / 2}$ into the left-hand side of eqn 7E.2, which then becomes

$$
\overbrace{\frac{\hbar^{2}}{2 m}\left(\frac{m k_{\mathrm{f}}}{\hbar^{2}}\right)^{1 / 2}}^{(\hbar / 2)\left(k_{\mathrm{f}} / m\right)^{1 / 2}} \psi_{0}-\overbrace{\frac{\hbar^{2}}{2 m}\left(\frac{m k_{\mathrm{f}}}{\hbar^{2}}\right)}^{k_{\mathrm{f}} / 2} x^{2} \psi_{0}+\frac{1}{2} k_{\mathrm{f}} x^{2} \psi_{0}=E \psi_{0}
$$

and therefore (keeping track of the blue terms)

$$
\frac{\hbar}{2}\left(\frac{k_{\mathrm{f}}}{m}\right)^{1 / 2} \psi_{0}-\frac{1}{2} k_{\mathrm{f}} x^{2} \psi_{0}+\frac{1}{2} k_{\mathrm{f}} x^{2} \psi_{0}=E \psi_{0}
$$

The blue terms cancel, leaving

$$
\frac{\hbar}{2}\left(\frac{k_{\mathrm{f}}}{m}\right)^{1 / 2} \psi_{0}=E \psi_{0}
$$

It follows that $\psi_{0}$ is a solution to the Schrödinger equation for the harmonic oscillator with energy $E=\frac{1}{2} \hbar\left(k_{\mathrm{f}} / m\right)^{1 / 2}$, in accord with eqn 7E. 5 for the zero-point energy.

Self-test 7E. 1 Confirm that the wavefunction in eqn 7E. 9 is a solution of eqn 7E. 2 and evaluate its energy.

The shapes of several of the wavefunctions are shown in Fig. 7E. 6 and the corresponding probability densities are shown in Fig. 7E.7. These probability densities show that, as the quantum number increases, the positions of highest probability migrate towards the classical turning points (see The chemist's toolkit 18). This behaviour is another example of the correspondence principle (Topic 7D) in which at high quantum numbers the classical behaviour emerges from the quantum behaviour.\\
\includegraphics[max width=\textwidth, center]{2025_02_27_d4b95d9718f0b9e2e6b9g-308(1)}

Figure 7E. 6 The normalized wavefunctions for the first seven states of a harmonic oscillator. Note that the number of nodes is equal to $v$. The wavefunctions with even $v$ are symmetric about $y=0$, and those with odd $v$ are anti-symmetric. The wavefunctions are shown superimposed on the potential energy function, and the horizontal axis for each wavefunction is set at the corresponding energy.\\
\includegraphics[max width=\textwidth, center]{2025_02_27_d4b95d9718f0b9e2e6b9g-309(1)}

Figure 7E. 7 The probability densities for the states of a harmonic oscillator with $v=0,5,10,15$, and 20 . Note how the regions of highest probability density move towards the turning points of the classical motion as $v$ increases.

The wavefunctions have the following features:

\begin{itemize}
  \item The Gaussian function decays quickly to zero as the displacement in either direction increases, so all the wavefunctions approach zero at large displacements: the particle is unlikely to be found at large displacements.
  \item The wavefunction oscillates between the classical turning points but decays without oscillating outside them.
  \item The exponent $y^{2}$ is proportional to $x^{2} \times\left(m k_{\mathrm{f}}\right)^{1 / 2}$, so the wavefunctions decay more rapidly for large masses and strong restoring forces (stiff springs).
  \item As $v$ increases, the Hermite polynomials become larger at large displacements (as $x^{\nu}$ ), so the wavefunctions grow large before the Gaussian function damps them down to zero: as a result, the wavefunctions spread over a wider range as $v$ increases (Fig. 7E.6).
\end{itemize}

\subsection*{Example 7E. 2 Normalizing a harmonic oscillator wavefunction}
Find the normalization constant for the harmonic oscillator wavefunctions.

Collect your thoughts A wavefunction is normalized (to 1) by evaluating the integral of $|\psi|^{2}$ over all space and then finding the normalization factor from eqn 7B. $3\left(N=1 /\left(\int \psi^{*} \psi \mathrm{~d} \tau\right)^{1 / 2}\right)$. The normalized wavefunction is then equal to $N \psi$. In this one-dimensional problem, the volume element is $\mathrm{d} x$ and the integration is from $-\infty$ to $+\infty$. The wavefunctions are expressed in terms of the dimensionless variable $y=x / \alpha$, so begin by expressing the integral in terms of $y$ by using $\mathrm{d} x=$ $\alpha \mathrm{d} y$. The integrals required are given in Table 7E.1.

The solution The unnormalized wavefunction is

$$
\psi_{v}(x)=H_{v}(y) \mathrm{e}^{-y^{2} / 2}
$$

It follows from the integrals given in Table 7E. 1 that

$$
\int_{-\infty}^{\infty} \psi_{v}^{*} \psi_{v} \mathrm{~d} x=\alpha \int_{-\infty}^{\infty} \psi_{v}^{*} \psi_{v} \mathrm{~d} y=\alpha \int_{-\infty}^{\infty} H_{v}^{2}(y) \mathrm{e}^{-y^{2}} \mathrm{~d} y=\alpha \pi^{1 / 2} 2^{v} v!
$$

where $v!=v(v-1)(v-2) \ldots 1$ and $0!\equiv 1$. Therefore,


\begin{equation*}
N_{v}=\left(\frac{1}{\alpha \pi^{1 / 2} 2^{v} v!}\right)^{1 / 2} \quad \text { Normalization constant } \tag{7E.10}
\end{equation*}


Note that $N_{v}$ is different for each value of $v$.\\
Self-test 7E. 2 Confirm, by explicit evaluation of the integral, that $\psi_{0}$ and $\psi_{1}$ are orthogonal.\\
\includegraphics[max width=\textwidth, center]{2025_02_27_d4b95d9718f0b9e2e6b9g-309}

\subsection*{7E. 2 Properties of the harmonic oscillator}
The average value of a property is calculated by evaluating the expectation value of the corresponding operator (eqn 7C.11, $\langle\Omega\rangle=\int \psi^{\star} \hat{\Omega} \psi \mathrm{d} \tau$ for a normalized wavefunction). For a harmonic oscillator,


\begin{equation*}
\langle\Omega\rangle_{v}=\int_{-\infty}^{\infty} \psi_{v}^{*} \hat{\Omega} \psi_{v} \mathrm{~d} x \tag{7E.11}
\end{equation*}


When the explicit wavefunctions are substituted, the integrals might look fearsome, but the Hermite polynomials have many features that simplify the calculation.

\subsection*{(a) Mean values}
Equation 7E. 11 can be used to calculate the mean displacement, $\langle x\rangle$, and the mean square displacement, $\left\langle x^{2}\right\rangle$, for a harmonic oscillator in a state with quantum number $v$.

\subsection*{How is that done? 7E. 1 Finding the mean values of $x$ and $x^{2}$ for the harmonic oscillator}
The evaluation of the integrals needed to compute $\langle x\rangle$ and $\left\langle x^{2}\right\rangle$ is simplified by recognizing the symmetry of the problem and using the special properties of the Hermite polynomials.

Step 1 Use a symmetry argument to find the mean displacement The mean displacement $\langle x\rangle$ is expected to be zero because the probability density of the oscillator is symmetrical about zero; that is, there is equal probability of positive and negative displacements.

Step 2 Confirm the result by examining the necessary integral

More formally, the mean value of $x$, which is expectation value of $x$, is

$$
\begin{aligned}
\langle x\rangle_{v} & =\int_{-\infty}^{\infty} \psi_{v}^{*} x \psi_{v} \mathrm{~d} x=N_{v}^{2} \int_{-\infty}^{\infty}\left(H_{v} \mathrm{e}^{-y^{2} / 2}\right) x\left(H_{v} \mathrm{e}^{-y^{2} / 2}\right) \mathrm{d} x \\
& \underbrace{x=\alpha y \mathrm{~d} x=\alpha \mathrm{d} y}_{\text {An odd function }} \\
& =\alpha^{2} N_{v}^{2} \int_{-\infty}^{\infty} \overbrace{y\left(H_{v} \mathrm{e}^{-y^{2} / 2}\right)^{2}} \mathrm{~d} y
\end{aligned}
$$

The integrand is an odd function because when $y \rightarrow-y$ it changes sign (the squared term does not change sign, but the term $y$ does). The integral of an odd function over a symmetrical range is necessarily zero, so\\
$\langle x\rangle_{v}=0$ for all $v \quad$ Mean displacement (7E.12a)

\subsection*{Step 3 Find the mean square displacement}
The mean square displacement, the expectation value of $x^{2}$, is

$$
\begin{aligned}
\left\langle x^{2}\right\rangle_{v} & =N_{v}^{2} \int_{-\infty}^{\infty}\left(H_{v} \mathrm{e}^{-y^{2} / 2}\right) x^{2}\left(H_{v} \mathrm{e}^{-y^{2} / 2}\right) \mathrm{d} x \\
& -x=\alpha y \mathrm{~d} x=\alpha \mathrm{d} y \\
& =\alpha^{3} N_{v}^{2} \int_{-\infty}^{\infty}\left(H_{v} \mathrm{e}^{-y^{2} / 2}\right) y^{2}\left(H_{v} \mathrm{e}^{-y^{2} / 2}\right) \mathrm{d} y
\end{aligned}
$$

You can develop the factor $y^{2} H_{v}$ by using the recursion relation given in Table 7E. 1 rearranged into $y H_{v}=v H_{v-1}+\frac{1}{2} H_{v+1}$. After multiplying this expression by $y$ it becomes

$$
y^{2} H_{v}=v y H_{v-1}+\frac{1}{2} y H_{v+1}
$$

Now use the recursion relation (with $v$ replaced by $v-1$ or $v+1)$ again for both $y H_{v-1}$ and $y H_{v+1}$ :

$$
\begin{aligned}
& y H_{v-1}=(v-1) H_{v-2}+\frac{1}{2} H_{v} \\
& y H_{v+1}=(v+1) H_{v}+\frac{1}{2} H_{v+2}
\end{aligned}
$$

It follows that

$$
\begin{aligned}
& y^{2} H_{v}=v y H_{v-1}+\frac{1}{2} y H_{v+1}=v\left\{(v-1) H_{v-2}+\frac{1}{2} H_{v}\right\} \\
& +\frac{1}{2}\left\{(v+1) H_{v}+\frac{1}{2} H_{v+2}\right\} \\
& \quad=v(v-1) H_{v-2}+\left(v+\frac{1}{2}\right) H_{v}+\frac{1}{4} H_{v+2}
\end{aligned}
$$

Substitution of this result into the integral gives

$$
\left\langle x^{2}\right\rangle_{v}=\alpha^{3} N_{v}^{2} \int_{-\infty}^{\infty}\left(H_{v} \mathrm{e}^{-y^{2} / 2}\right) \overbrace{\left\{v(v-1) H_{v-2}+\left(v+\frac{1}{2}\right) H_{v}+\frac{1}{4} H_{v+2}\right\}}^{y^{2} H_{v}} \mathrm{e}^{-y^{2} / 2} \mathrm{~d} y
$$

$$
=\alpha^{3} N_{v}^{2} v(v-1) \overbrace{\int_{-\infty}^{\infty} H_{v} H_{v-2} \mathrm{e}^{-y^{2}} \mathrm{~d} y}^{0}
$$

$$
\begin{aligned}
& +\alpha^{3} N_{v}^{2}\left(v+\frac{1}{2}\right) \overbrace{\int_{-\infty}^{\infty} H_{v} H_{v} \mathrm{e}^{-y^{2}} \mathrm{~d} y+\frac{1}{4} \alpha^{3} N_{v}^{2} \overbrace{\int_{-\infty}^{\infty} H_{v} H_{v+2} \mathrm{e}^{-y^{2}} \mathrm{~d} y}^{\pi^{1 / 2} 2^{v} v!}}^{0} \\
& =\alpha^{3} N_{v}^{2}\left(v+\frac{1}{2}\right) \pi^{1 / 2} 2^{v} v!
\end{aligned}
$$

Each of the three integrals is evaluated by making use of the information in Table 7E.1. Therefore, after noting the expression for $N_{v}$ in eqn 7E.10,

$$
\left\langle x^{2}\right\rangle_{v}=\frac{\alpha^{3}\left(v+\frac{1}{2}\right) \pi^{1 / 2} 2^{v} v!}{\alpha \pi^{1 / 2} 2^{v} v!}=\left(v+\frac{1}{2}\right) \alpha^{2}
$$

Finally, with, $\alpha^{2}=\left(\hbar^{2} / m k_{f}\right)^{1 / 2}$\\
\includegraphics[max width=\textwidth, center]{2025_02_27_d4b95d9718f0b9e2e6b9g-310}\\
The result for $\langle x\rangle_{v}$ shows that the oscillator is equally likely to be found on either side of $x=0$ (like a classical oscillator). The result for $\left\langle x^{2}\right\rangle_{v}$ shows that the mean square displacement increases with $v$. This increase is apparent from the probability densities in Fig. 7E.7, and corresponds to the amplitude of a classical harmonic oscillator becoming greater as its energy increases.

The mean potential energy of an oscillator, which is the expectation value of $V=\frac{1}{2} k_{\mathrm{f}} x^{2}$, can now be calculated:

$$
\langle V\rangle_{v}=\frac{1}{2}\left\langle k_{\mathrm{f}} x^{2}\right\rangle_{v}=\frac{1}{2} k_{\mathrm{f}}\left\langle x^{2}\right\rangle_{v}=\frac{1}{2}\left(v+\frac{1}{2}\right) \hbar\left(\frac{k_{\mathrm{f}}}{m}\right)^{1 / 2}
$$

or

$$
\langle V\rangle_{v}=\frac{1}{2}\left(v+\frac{1}{2}\right) \hbar \omega \quad \text { Mean potential energy }
$$

(7E.13a)\\
Because the total energy in the state with quantum number $v$ is $\left(v+\frac{1}{2}\right) \hbar \omega$, it follows that

$$
\langle V\rangle_{v}=\frac{1}{2} E_{v}
$$

Mean potential energy\\
(7E.13b)\\
The total energy is the sum of the potential and kinetic energies, $E_{v}=\langle V\rangle_{v}+\left\langle E_{\mathrm{k}}\right\rangle_{v}$, so it follows that the mean kinetic energy of the oscillator is

$$
\left\langle E_{\mathrm{k}}\right\rangle_{v}=E_{v}-\langle V\rangle_{v}=E_{v}-\frac{1}{2} E_{v}=\frac{1}{2} E_{v} \quad \begin{aligned}
& \text { Mean kinetic } \\
& \text { energy }
\end{aligned}
$$

(7E.13c)\\
The result that the mean potential and kinetic energies of a harmonic oscillator are equal (and therefore that both are equal to half the total energy) is a special case of the virial theorem:

If the potential energy of a particle has the form $V=a x^{b}$, then its mean potential and kinetic energies are related by

$$
2\left\langle E_{\mathrm{k}}\right\rangle=b\langle V\rangle
$$

Virial theorem\\
(7E.14)\\
For a harmonic oscillator $b=2$, so $\left\langle E_{\mathrm{k}}\right\rangle_{v}=\langle V\rangle_{v}$. The virial theorem is a short cut to the establishment of a number of useful results, and it is used elsewhere (e.g. in Topic 8A).

\subsection*{(b) Tunnelling}
A quantum oscillator may be found at displacements with $V>E$, which are forbidden by classical physics because they correspond to negative kinetic energy. That is, a harmonic\\
oscillator can tunnel into classically forbidden displacements. As shown in Example 7E.3, for the lowest energy state of the harmonic oscillator, there is about an 8 per cent chance of finding the oscillator at classically forbidden displacements in either direction. These tunnelling probabilities are independent of the force constant and mass of the oscillator.

Example 7E. 3 Calculating the tunnelling probability for the harmonic oscillator

Calculate the probability that the ground-state harmonic oscillator will be found in a classically forbidden region.

Collect your thoughts Find the expression for the classical turning point, $x_{\mathrm{t} \text {, }}$, where the kinetic energy goes to zero, by equating the potential energy to the total energy of the harmonic oscillator. You can then calculate the probability of finding the oscillator at a displacement beyond $x_{\text {tp }}$ by integrating $\psi^{2} \mathrm{~d} x$ between $x_{\mathrm{tp}}$ and infinity

$$
P=\int_{x_{\mathrm{p}}}^{\infty} \psi_{v}^{2} \mathrm{~d} x
$$

By symmetry, the probability of the particle being found in the classically forbidden region from $-x_{\mathrm{tp}}$ to $-\infty$ is the same.

The solution According to classical mechanics, the turning point, $x_{\mathrm{tp}}$, of an oscillator occurs when its potential energy $\frac{1}{2} k_{\mathrm{f}} x^{2}$ is equal to its total energy. When that energy is one of the allowed values $E_{v}$, the turning point is at

$$
E_{v}=\frac{1}{2} k_{\mathrm{f}} x_{\mathrm{tp}}^{2} \quad \text { and therefore at } \quad x_{\mathrm{tp}}= \pm\left(\frac{2 E_{v}}{k_{\mathrm{f}}}\right)^{1 / 2}
$$

The variable of integration in the integral $P$ is best expressed in terms of $y=x / \alpha$ with $\alpha=\left(\hbar^{2} / m k_{\mathrm{f}}\right)^{1 / 4}$. With these substitutions, and also using $E_{v}=\left(v+\frac{1}{2}\right) \hbar \omega$, the turning points are given by

$$
y_{\mathrm{tp}}=\frac{x_{\mathrm{tp}}}{\alpha}=\left\{\frac{2\left(v+\frac{1}{2}\right) \hbar \omega}{\alpha^{2} k_{\mathrm{f}}}\right\}^{1 / 2} \stackrel{\nabla}{=}=(2 v+1)^{1 / 2}
$$

For the state of lowest energy $(v=0), y_{\mathrm{tp}}=1$ and the probability of being beyond that point is

$$
\begin{aligned}
& \qquad P=\int_{x_{\mathrm{tp}}}^{\infty} \psi_{0}^{2} \mathrm{~d} x=\alpha \int_{1}^{\infty} \psi_{0}^{2} \mathrm{~d} y=\alpha N_{0}^{2} \int_{1}^{\infty} \mathrm{e}^{-y^{2}} \mathrm{~d} y \\
& \text { with } \\
& \qquad N_{0}=\left(\frac{1}{\alpha \pi^{1 / 2} 2^{0} 0!}\right)^{1 / 2}=\left(\frac{1}{\alpha \pi^{1 / 2}}\right)^{1 / 2}
\end{aligned}
$$

Therefore

$$
P=\frac{1}{\pi^{1 / 2}} \int_{1}^{\infty} \mathrm{e}^{-y^{2}} \mathrm{~d} y
$$

The integral must be evaluated numerically (by using mathematical software), and is equal to $0.139 \ldots$ It follows that $P=$ 0.079 .

Comment. In 7.9 per cent of a large number of observations of an oscillator in the state with quantum number $v=0$, the particle will be found beyond the (positive) classical turning point. It will be found with the same probability at negative forbidden displacements. The total probability of finding the oscillator in a classically forbidden region is about 16 per cent.\\
Self-test 7E. 3 Calculate the probability that a harmonic oscillator in the state with quantum number $v=1$ will be found at a classically forbidden extension. You will need to use mathematical software to evaluate the integral.

$$
9 \subseteq 0^{\circ} 0={ }_{d}: \iota \partial м s u_{V}
$$

The probability of finding the oscillator in classically forbidden regions decreases quickly with increasing $v$, and vanishes entirely as $v$ approaches infinity, as is expected from the correspondence principle. Macroscopic oscillators (such as pendulums) are in states with very high quantum numbers, so the tunnelling probability is wholly negligible and classical mechanics is reliable. Molecules, however, are normally in their vibrational ground states, and for them the probability is very significant and classical mechanics is misleading.

\subsection*{Checklist of concepts}
\begin{enumerate}
  \item The energy levels of a quantum mechanical harmonic oscillator are evenly spaced.2. The wavefunctions of a quantum mechanical harmonic oscillator are products of a Hermite polynomial and a Gaussian (bell-shaped) function.
  \item A quantum mechanical harmonic oscillator has zeropoint energy, an irremovable minimum energy.\\
$\square$ 4. The probability of finding a quantum mechanical harmonic oscillator at classically forbidden displacements is significant for the ground vibrational state $(v=0)$ but decreases quickly with increasing $v$.
\end{enumerate}

\subsection*{Checklist of equations}
\begin{center}
\begin{tabular}{|c|c|c|c|}
\hline
Property & Equation & Comment & Equation number \\
\hline
Energy levels & $E_{v}=\left(\nu+\frac{1}{2}\right) \hbar \omega \quad \omega=\left(k_{\mathrm{f}} / m\right)^{1 / 2}$ & $v=0,1,2, \ldots$ & 7E. 3 \\
\hline
Zero-point energy & $E_{0}=\frac{1}{2} \hbar \omega$ &  & 7E. 5 \\
\hline
Wavefunctions & \( \begin{aligned} & \psi_{v}(x)=N_{v} H_{v}(y) \mathrm{e}^{-y^{2} / 2} \\ & y=x / \alpha \quad \alpha=\left(\hbar^{2} / m k_{\mathrm{f}}\right)^{1 / 4} \end{aligned} \) & $v=0,1,2, \ldots$ & 7E. 7 \\
\hline
Normalization constant & $N_{v}=\left(1 / \alpha \pi^{1 / 2} 2^{v} v!\right)^{1 / 2}$ &  & 7E. 10 \\
\hline
Mean displacement & $\langle x\rangle_{v}=0$ &  & 7E.12a \\
\hline
Mean square displacement & $\left\langle x^{2}\right\rangle_{\nu}=\left(v+\frac{1}{2}\right) \hbar /\left(m k_{\mathrm{f}}\right)^{1 / 2}$ &  & 7E.12b \\
\hline
Virial theorem & $2\left\langle E_{\mathrm{k}}\right\rangle=b\langle V\rangle$ & $V=a x^{b}$ & 7E. 14 \\
\hline
\end{tabular}
\end{center}

\section*{TOPIC 7F Rotational motion}
\subsection*{Why do you need to know this material?}
Angular momentum is central to the description of the electronic structure of atoms and molecules and the interpretation of molecular spectra.

\subsection*{What is the main idea?}
The energy, angular momentum, and orientation of the angular momentum of a rotating body are quantized.

\subsection*{What do you need to know already?}
You should be aware of the postulates of quantum mechanics and the role of boundary conditions (Topics 7C and 7D). Background information on the description of rotation and the coordinate systems used to describe it are given in three Toolkits.

Rotational motion is encountered in many aspects of chemistry, including the electronic structures of atoms, because electrons orbit (in a quantum mechanical sense) around nuclei and spin on their axis. Molecules also rotate; transitions between their rotational states affect the appearance of spectra and their detection gives valuable information about the structures of molecules.

\subsection*{7F. 1 Rotation in two dimensions}
Consider a particle of mass $m$ constrained to move in a circular path (a 'ring') of radius $r$ in the $x y$-plane with constant potential energy, which may be taken to be zero (Fig. 7F.1); the energy is entirely kinetic. The Schrödinger equation is


\begin{equation*}
-\frac{\hbar^{2}}{2 m}\left(\frac{\partial^{2}}{\partial x^{2}}+\frac{\partial^{2}}{\partial y^{2}}\right) \psi(x, y)=E \psi(x, y) \tag{7F.1}
\end{equation*}


with the particle confined to a path of constant radius $r$. The equation is best expressed in cylindrical coordinates $r$ and $\phi$ with $z=0$ (The chemist's toolkit 19) because they reflect the symmetry of the system. In cylindrical coordinates


\begin{equation*}
\frac{\partial^{2}}{\partial x^{2}}+\frac{\partial^{2}}{\partial y^{2}}=\frac{\partial^{2}}{\partial r^{2}}+\frac{1}{r} \frac{\partial}{\partial r}+\frac{1}{r^{2}} \frac{\partial^{2}}{\partial \phi^{2}} \tag{7F.2}
\end{equation*}


\begin{center}
\includegraphics[max width=\textwidth]{2025_02_27_d4b95d9718f0b9e2e6b9g-313}
\end{center}

Figure 7F. 1 A particle on a ring is free to move in the $x y$-plane around a circular path of radius $r$.

However, because the radius of the path is fixed, the (blue) derivatives with respect to $r$ can be discarded. Only the last term in eqn 7F. 2 then survives and the Schrödinger equation becomes


\begin{equation*}
-\frac{\hbar^{2}}{2 m r^{2}} \frac{\mathrm{~d}^{2} \psi(\phi)}{\mathrm{d} \phi^{2}}=E \psi(\phi) \tag{7F.3a}
\end{equation*}


The partial derivative has been replaced by a complete derivative because $\phi$ is now the only variable. The term $m r^{2}$ is the moment of inertia, $I=m r^{2}$ (The chemist's toolkit 20), and so the Schrödinger equation becomes

\[
-\frac{\hbar^{2}}{2 I} \frac{\mathrm{~d}^{2} \psi(\phi)}{\mathrm{d} \phi^{2}}=E \psi(\phi) \quad \begin{align*}
& \text { Schrödinger equation }  \tag{7F.3b}\\
& \text { [particle on a ring] }
\end{align*}
\]

\subsection*{The chemist's toolkit 19}
\subsection*{Cylindrical coordinates}
For systems with cylindrical symmetry it is best to work in cylindrical coordinates $r, \phi$, and $z$ (Sketch 1), with

$$
x=r \cos \phi \quad y=r \sin \phi
$$

and where

$$
0 \leq r \leq \infty \quad 0 \leq \phi \leq 2 \pi \quad-\infty \leq z \leq+\infty
$$

The volume element is

$$
\mathrm{d} \tau=r \mathrm{~d} r \mathrm{~d} \phi \mathrm{~d} z
$$

For motion in a plane, $z=0$ and the volume element is

$$
\mathrm{d} \tau=r \mathrm{~d} r \mathrm{~d} \phi
$$

\begin{center}
\includegraphics[max width=\textwidth]{2025_02_27_d4b95d9718f0b9e2e6b9g-313(1)}
\end{center}

Sketch 1

\subsection*{The chemist's toolkit 20}
\subsection*{Angular momentum}
Angular velocity, $\omega$ (omega), is the rate of change of angular position; it is reported in radians per second ( $\mathrm{rad} \mathrm{s}^{-1}$ ). There are $2 \pi$ radians in a circle, so 1 cycle per second is the same as $2 \pi$ radians per second. For convenience, the 'rad' is often dropped, and the units of angular velocity are denoted $\mathrm{s}^{-1}$.\\
Expressions for other angular properties follow by analogy with the corresponding equations for linear motion (The chemist's toolkit 3 in Topic 1B). Thus, the magnitude, J, of the angular momentum, $J$, is defined, by analogy with the magnitude of the linear momentum $(p=m v)$ :

$$
J=I \omega
$$

The quantity $I$ is the moment of inertia of the object. It represents the resistance of the object to a change in the state of rotation in the same way that mass represents the resistance of the object to a change in the state of translation. In the case of a rotating molecule the moment of inertia is defined as

$$
I=\sum_{i} m_{i} r_{i}^{2}
$$

where $m_{i}$ is the mass of atom $i$ and $r_{i}$ is its perpendicular distance from the axis of rotation (Sketch 1). For a point particle of mass $m$ moving in a circle of radius $r$, the moment of inertia about the axis of rotation is

$$
I=m r^{2}
$$

The SI units of moment of inertia are therefore kilogram metre ${ }^{2}$ $\left(\mathrm{kg} \mathrm{m}^{2}\right)$, and those of angular momentum are kilogram metre ${ }^{2}$ per second ( $\mathrm{kg} \mathrm{m}^{2} \mathrm{~s}^{-1}$ ).\\
\includegraphics[max width=\textwidth, center]{2025_02_27_d4b95d9718f0b9e2e6b9g-314}

Sketch 1\\
The angular momentum is a vector, a quantity with both magnitude and direction (The chemist's toolkit 17 in Topic 7D). For rotation in three dimensions, the angular momentum has three components: $J_{x}, J_{y}$, and $J_{z}$. For a particle travelling on a circular path of radius $r$ about the $z$-axis, and therefore confined to the $x y$-plane, the angular momentum vector points in the $z$-direction only (Sketch 2), and its only component is

$$
J_{z}= \pm p r
$$

where $p$ is the magnitude of the linear momentum in the $x y$ plane at any instant. When $J_{z}>0$, the particle travels in a clockwise direction as viewed from below; when $J_{z}<0$, the motion is anticlockwise. A particle that is travelling at high speed in a circle has a higher angular momentum than a particle of the same mass travelling more slowly. An object with a high angular momentum (like a flywheel) requires a strong braking force (more precisely, a strong 'torque') to bring it to a standstill.\\
\includegraphics[max width=\textwidth, center]{2025_02_27_d4b95d9718f0b9e2e6b9g-314(1)}

Sketch 2\\
The components of the angular momentum vector $J$ when it lies in a general orientation are

$$
J_{x}=y p_{z}-z p_{y} \quad J_{y}=z p_{x}-x p_{z} \quad J_{z}=x p_{y}-y p_{x}
$$

where $p_{x}$ is the component of the linear momentum in the $x$-direction at any instant, and likewise $p_{y}$ and $p_{z}$ in the other directions. The square of the magnitude of the angular momentum vector is given by

$$
J^{2}=J_{x}^{2}+J_{y}^{2}+J_{z}^{2}
$$

By analogy with the expression for linear motion $\left(E_{\mathrm{k}}=\right.$ $\frac{1}{2} m v^{2}=p^{2} / 2 m$ ), the kinetic energy of a rotating object is

$$
E_{\mathrm{k}}=\frac{1}{2} I \omega^{2}=\frac{J^{2}}{2 I}
$$

For a given moment of inertia, high angular momentum corresponds to high kinetic energy. As may be verified, the units of rotational energy are joules (J).\\
The analogous roles of $m$ and $I$, of $v$ and $\omega$, and of $p$ and $J$ in the translational and rotational cases respectively provide a ready way of constructing and recalling equations. These analogies are summarized below:

\begin{center}
\begin{tabular}{llll}
\hline
Translation &  & Rotation &  \\
\hline
Property & Significance & Property & Significance \\
\hline
Mass, $m$ & \begin{tabular}{l}
Resistance to \\
the effect of a \\
force \\
\end{tabular} & \begin{tabular}{l}
Moment of \\
inertia, $I$ \\
\end{tabular} & \begin{tabular}{l}
Resistance to the \\
effect of a twisting \\
force (torque) \\
\end{tabular} \\
Speed, $v$ & \begin{tabular}{l}
Rate of change \\
of position \\
\end{tabular} & \begin{tabular}{l}
Angular velocity, \\
$\omega$ \\
\end{tabular} & \begin{tabular}{l}
Rate of change of \\
angle \\
\end{tabular} \\
\begin{tabular}{l}
Magnitude \\
of linear \\
momentum, $p$ \\
\end{tabular} & $p=m v$ & \begin{tabular}{l}
Magnitude \\
of angular \\
momentum, $J$ \\
\end{tabular} & $J=I \omega$ \\
\begin{tabular}{l}
Translational \\
kinetic energy, \\
$E_{\mathrm{k}}$ \\
\end{tabular} & $E_{\mathrm{k}}=\frac{1}{2} m v^{2}=$ & \begin{tabular}{l}
Rotational \\
$p^{2} / 2 m$ \\
\end{tabular} & $E_{\mathrm{k}}=\frac{1}{2} I \omega^{2}=J^{2} / 2 I$ \\
\hline
\end{tabular}
\end{center}

\subsection*{(a) The solutions of the Schrödinger equation}
The most straightforward way of finding the solutions of eqn 7 F .3 b is to take the known general solution to a second-order differential equation of this kind and show that it does indeed satisfy the equation. Then find the allowed solutions and energies by imposing the relevant boundary conditions.

How is that done? 7F. 1 Finding the solutions of the Schrödinger equation for a particle on a ring

A solution of eqn 7F.3b is

$$
\psi(\phi)=\mathrm{e}^{\mathrm{i} m_{l} \phi}
$$

where, as yet, $m_{l}$ is an arbitrary dimensionless number (the notation is explained later). This is not the most general solution, which would be $\psi(\phi)=A e^{\mathrm{im} / \phi}+B \mathrm{e}^{-\mathrm{i} m, \phi}$, but is sufficiently general for the present purpose.

Step 1 Verify that the function satisfies the equation To verify that $\psi(\phi)$ is a solution note that

$$
\frac{\mathrm{d}^{2}}{\mathrm{~d} \phi^{2}} \mathrm{e}^{\mathrm{i} m_{l} \phi}=\frac{\mathrm{d}}{\mathrm{~d} \phi}\left(\mathrm{i} m_{l}\right) \mathrm{e}^{\mathrm{i} m_{l} \phi}=\left(\mathrm{i} m_{l}\right)^{2} \mathrm{e}^{\mathrm{i} m_{l} \phi}=-m_{l}^{2} \mathrm{e}^{\mathrm{i} m_{l} \phi}=-m_{l}^{2} \psi
$$

Then

$$
-\frac{\hbar^{2}}{2 I} \frac{\mathrm{~d}^{2} \psi}{\mathrm{~d} \phi^{2}}=-\frac{\hbar^{2}}{2 I}\left(-m_{l}^{2} \psi\right)=\frac{m_{l}^{2} \hbar^{2}}{2 I} \psi
$$

which has the form constant $\times \psi$, so the proposed wavefunction is indeed a solution and the corresponding energy is $m_{l}^{2} \hbar^{2} / 2 I$.

Step 2 Impose the appropriate boundary conditions\\
The requirement that a wavefunction must be single-valued implies the existence of a cyclic boundary condition, the requirement that the wavefunction must be the same after a complete revolution: $\psi(\phi+2 \pi)=\psi(\phi)$ (Fig. 7F.2). In this case

$$
\begin{aligned}
\psi(\phi+2 \pi) & =\mathrm{e}^{\mathrm{i} m_{l}(\phi+2 \pi)}=\mathrm{e}^{\mathrm{i} m_{l} \phi} \mathrm{e}^{2 \pi \mathrm{i} m_{l}} \\
& =\psi(\phi) \mathrm{e}^{2 \pi m_{l}}=\psi(\phi)\left(\mathrm{e}^{\mathrm{i} \pi}\right)^{2 m_{l}}
\end{aligned}
$$

As $\mathrm{e}^{\mathrm{i} \pi}=-1$, this relation is equivalent to

$$
\psi(\phi+2 \pi)=(-1)^{2 m} \psi(\phi)
$$

The cyclic boundary condition $\psi(\phi+2 \pi)=\psi(\phi)$ requires $(-1)^{2 m_{l}}=1$; this requirement is satisfied for any positive or negative integer value of $m_{l}$, including 0 .\\
Step 3 Normalize the wavefunction\\
A one-dimensional wavefunction is normalized (to 1) by finding the normalization constant $N$ given by eqn 7B. 3\\
\includegraphics[max width=\textwidth, center]{2025_02_27_d4b95d9718f0b9e2e6b9g-315}

Figure 7F. 2 Two possible solutions of the Schrödinger equation for a particle on a ring. The circumference has been opened out into a straight line; the points at $\phi=0$ and $2 \pi$ are identical. The solution in (a), $e^{i \phi}=\cos \phi+i \sin \phi$, is acceptable because after a complete revolution the wavefunction has the same value. The solution in $(b), \mathrm{e}^{0.9 i \phi}=\cos (0.9 \phi)+\mathrm{i} \sin (0.9 \phi)$ is unacceptable because its value, both for the real and imaginary parts, is not the same at $\phi=0$ and $2 \pi$.\\
( $N=\left(\int_{-\infty}^{\infty} \psi^{\star} \psi \mathrm{d} x\right)^{-1 / 2}$ ). In this case, the wavefunction depends only on the angle $\phi$ and the range of integration is from $\phi=0$ to $2 \pi$, so the normalization constant is

$$
N=\frac{1}{\left(\int_{0}^{2 \pi} \psi^{*} \psi \mathrm{~d} \phi\right)^{1 / 2}}=\frac{1}{(\int_{0}^{2 \pi} \underbrace{\mathrm{e}^{-\mathrm{i} m m_{1} \phi} \mathrm{e}^{\mathrm{i} m \phi \phi}}_{1} \mathrm{~d} \phi)^{1 / 2}}=\frac{1}{(2 \pi)^{1 / 2}}
$$

The normalized wavefunctions and corresponding energies are labelled with the integer $m_{l}$, which is playing the role of a quantum number, and are therefore

$$
-\begin{aligned}
& \psi_{m_{l}}(\phi)=\frac{\mathrm{e}^{\mathrm{i} m_{l} \phi}}{(2 \pi)^{1 / 2}} \\
& E_{m_{l}}=\frac{m_{l}^{2} \hbar^{2}}{2 I} \quad m_{l}=0, \pm 1, \pm 2, \ldots
\end{aligned}
$$

Wavefunctions and energy levels of a particle on a ring

Apart from the level with $m_{l}=0$, each of the energy levels is doubly degenerate because the dependence of the energy on $m_{l}^{2}$ means that two values of $m_{l}$ (such as +1 and -1 ) correspond to the same energy.

A note on good practice Note that, when quoting the value of $m_{l}$, it is good practice always to give the sign, even if $m_{l}$ is positive. Thus, write $m_{l}=+1$, not $m_{l}=1$.

\subsection*{(b) Quantization of angular momentum}
Classically, a particle moving around a circular path possesses 'angular momentum' analogous to the linear momentum possessed by a particle moving in a straight line (The chemist's toolkit 20). Although in general angular momentum is represented by the vector $J$, when considering orbital angular momentum, the angular momentum of a particle around a fixed point in space, it is denoted $\boldsymbol{l}$. It can be shown that angular momentum also occurs in quantum mechanical systems, including a particle on a ring, but its magnitude is confined to discrete values.

\subsection*{How is that done? 7F. 2}
Showing that angular momentum is quantized

As explained in Topic 7C, the outcome of a measurement of a property is one of the eigenfunctions of the corresponding operator. The first step is therefore to identify the operator corresponding to angular momentum, and then to identify its eigenvalues.

\subsection*{Step 1 Construct the operator for angular momentum}
Because the particle is confined to the $x y$-plane, its angular momentum is directed along the $z$-axis, so only this component need be considered. According to The chemist's toolkit 20 , the $z$-component of the orbital angular momentum is

$$
l_{z}=x p_{y}-y p_{x}
$$

where $x$ and $y$ specify the position and $p_{x}$ and $p_{y}$ are the components of the linear momentum of the particle. The corresponding operator is formed by replacing $x, y, p_{x}$, and $p_{y}$ by their corresponding operators (Topic $7 \mathrm{C} ; \hat{q}=q \times$ and $\hat{p}_{q}=(\hbar / \mathrm{i}) \partial / \partial q$, with $q=x$ and $\left.y\right)$, which gives

\[
\hat{l}_{z}=\frac{\hbar}{\mathrm{i}}\left(x \frac{\partial}{\partial y}-y \frac{\partial}{\partial x}\right) \quad \begin{align*}
& \text { Operator for the z-component }  \tag{7F.5a}\\
& \text { of the angular momentum }
\end{align*}
\]

In cylindrical coordinates (see The chemist's toolkit 19) this operator becomes


\begin{equation*}
\hat{l}_{z}=\frac{\hbar}{\mathrm{i}} \frac{\mathrm{~d}}{\mathrm{~d} \phi} \tag{7F.5b}
\end{equation*}


Step 2 Verify that the wavefunctions are eigenfunctions of this operator\\
To decide whether the wavefunctions in eqn 7F. 4 are eigenfunctions of $\hat{l}_{z}$, allow it to act on the wavefunction:

$$
\hat{l}_{z} \psi_{m_{l}}=\frac{\hbar}{\mathrm{i}} \frac{\mathrm{~d}}{\mathrm{~d} \phi} \mathrm{e}^{\mathrm{i} m_{l, \phi}}=\frac{\hbar}{\mathrm{i}} \mathrm{i} m_{l} \overbrace{\mathrm{e}^{\mathrm{i} m_{l / \phi}}}^{\psi_{m_{l}}}=m_{l} \hbar \psi_{m_{l}}
$$

The wavefunction is an eigenfunction of the angular momentum, with the eigenvalue $m_{l} \hbar$. In summary,\\
$-\hat{l}_{z} \psi_{m_{l}}(\phi)=m_{l} \hbar \psi_{m_{l}}(\phi) \quad m_{l}=0, \pm 1, \pm 2,\left.\ldots\right|_{\text {Eigenfunctions of } \hat{l}_{z}}$

Because $m_{l}$ is confined to discrete values, the $z$-component of angular momentum is quantized. When $m_{l}$ is positive, the $z$-component of angular momentum is positive (clockwise rotation when seen from below); when $m_{l}$ is negative, the $z$-component of angular momentum is negative (anticlockwise when seen from below).

The important features of the results so far are:

\begin{itemize}
  \item The energies are quantized because $m_{l}$ is confined to integer values.
  \item The occurrence of $m_{l}$ as its square means that the energy of rotation is independent of the sense of rotation (the sign of $m_{l}$ ), as expected physically.
  \item Apart from the state with $m_{l}=0$, all the energy levels are doubly degenerate; rotation can be clockwise or anticlockwise with the sane energy.
  \item There is no zero-point energy: the particle can be stationary.
  \item As $m_{l}$ increases the wavefunctions oscillate with shorter wavelengths and so have greater curvature, corresponding to increasing kinetic energy (Fig. 7F.3).
  \item As pointed out in Topic 7D, a wavefunction that is complex represents a direction of motion, and taking its complex conjugate reverses the direction. The wavefunctions with $m_{l}>0$ and $m_{l}<0$ are each other's complex conjugate, and so they correspond to motion in opposite directions.
\end{itemize}

The probability density predicted by the wavefunctions of eqn 7F. 4 is uniform around the ring:

$$
\begin{aligned}
\psi_{m_{l}}^{\star} \psi_{m_{l}} & =\left(\frac{\mathrm{e}^{\mathrm{i} m m_{\phi} \phi}}{(2 \pi)^{1 / 2}}\right)^{*}\left(\frac{\mathrm{e}^{\mathrm{i} m m_{\phi}}}{(2 \pi)^{1 / 2}}\right) \\
& =\left(\frac{\mathrm{e}^{-\mathrm{i} m_{l} \phi}}{(2 \pi)^{1 / 2}}\right)\left(\frac{\mathrm{e}^{\mathrm{i} m m_{\phi}}}{(2 \pi)^{1 / 2}}\right)=\frac{1}{2 \pi}
\end{aligned}
$$

\begin{center}
\includegraphics[max width=\textwidth]{2025_02_27_d4b95d9718f0b9e2e6b9g-316}
\end{center}

Figure 7F. 3 The real parts of the wavefunctions of a particle on a ring. As the energy increases, so does the number of nodes and the curvature.

Angular momentum and angular position are a pair of complementary observables (in the sense defined in Topic 7C), and the inability to specify them simultaneously with arbitrary precision is another example of the uncertainty principle. In this case the $z$-component of angular momentum is known exactly (as $\left.m_{l} \hbar\right)$ but the location of the particle on the ring is completely unknown, which is reflected by the uniform probability density.

\subsection*{Example 7F. 1 Using the particle on a ring model}
The particle-on-a-ring is a crude but illustrative model of cyclic, conjugated molecular systems. Treat the $\pi$ electrons in benzene as particles freely moving over a circular ring of six carbon atoms and calculate the minimum energy required for the excitation of a $\pi$ electron. The carbon-carbon bond length in benzene is 140 pm .

Collect your thoughts Because each carbon atom contributes one $\pi$ electron, there are six electrons to accommodate. Each state is occupied by two electrons, so only the $m_{l}=0,+1$, and -1 states are occupied (with the last two being degenerate). The minimum energy required for excitation corresponds to a transition of an electron from the $m_{l}=+1$ (or -1 ) state to the $m_{l}=+2$ (or -2 ) state. Use eqn 7F.4, and the mass of the electron, to calculate the energies of the states. A hexagon can be inscribed inside a circle with a radius equal to the side of the hexagon, so take $r=140 \mathrm{pm}$.

The solution From eqn 7F.4, the energy separation between the states with $m_{l}=+1$ and $m_{l}=+2$ is

$$
\begin{aligned}
\Delta E & =E_{+2}-E_{+1}=(4-1) \times \frac{\left(1.055 \times 10^{-34} \mathrm{Js}\right)^{2}}{2 \times\left(9.109 \times 10^{-31} \mathrm{~kg}\right) \times\left(1.40 \times 10^{-10} \mathrm{~m}\right)^{2}} \\
& =9.35 \times 10^{-19} \mathrm{~J}
\end{aligned}
$$

Therefore the minimum energy required to excite an electron is 0.935 aJ or $563 \mathrm{~kJ} \mathrm{~mol}^{-1}$. This energy separation corresponds to an absorption frequency of $1410 \mathrm{THz}\left(1 \mathrm{THz}=10^{12} \mathrm{~Hz}\right)$ and a wavelength of 213 nm ; the experimental value for a transition of this kind is 260 nm . Such a crude model cannot be expected to give quantitative agreement, but the value is at least of the right order of magnitude.

Self-test 7F. 1 Use the particle-on-a-ring model to calculate the minimum energy required for the excitation of a $\pi$ electron in coronene, $\mathrm{C}_{24} \mathrm{H}_{12}(\mathbf{1})$. Assume that the radius of the ring is three times the carbon-carbon bond length in benzene and that the electrons are confined to the periphery of the molecule.\\
\includegraphics[max width=\textwidth, center]{2025_02_27_d4b95d9718f0b9e2e6b9g-317(2)}

\subsection*{1 Coronene}
 (model ring in red)\includegraphics[max width=\textwidth, center]{2025_02_27_d4b95d9718f0b9e2e6b9g-317(1)}\\
\includegraphics[max width=\textwidth, center]{2025_02_27_d4b95d9718f0b9e2e6b9g-317}

\subsection*{7F. 2 Rotation in three dimensions}
Now consider a particle of mass $m$ that is free to move anywhere on the surface of a sphere of radius $r$.

\subsection*{(a) The wavefunctions and energy levels}
The potential energy of a particle on the surface of a sphere is the same everywhere and may be taken to be zero. The Schrödinger equation is therefore


\begin{equation*}
-\frac{\hbar^{2}}{2 m} \nabla^{2} \psi=E \psi \tag{7F.7a}
\end{equation*}


where the sum of the three second derivatives, denoted $\nabla^{2}$ and read 'del squared', is called the 'laplacian':


\begin{equation*}
\nabla^{2}=\frac{\partial^{2}}{\partial x^{2}}+\frac{\partial^{2}}{\partial y^{2}}+\frac{\partial^{2}}{\partial z^{2}} \quad \text { Laplacian } \tag{7F.7b}
\end{equation*}


To take advantage of the symmetry of the problem it is appropriate to change to spherical polar coordinates (The chemist's toolkit 21) when the laplacian becomes

$$
\nabla^{2}=\frac{1}{r} \frac{\partial^{2}}{\partial r^{2}} r+\frac{1}{r^{2}} \Lambda^{2}
$$

where the derivatives with respect to the colatitude $\theta$ and the azimuth $\phi$ are collected in $\Lambda^{2}$, which is called the 'legendrian' and is given by

$$
\Lambda^{2}=\frac{1}{\sin ^{2} \theta} \frac{\partial^{2}}{\partial \phi^{2}}+\frac{1}{\sin \theta} \frac{\partial}{\partial \theta} \sin \theta \frac{\partial}{\partial \theta}
$$

In the present case, $r$ is fixed, so the derivatives with respect to $r$ in the laplacian can be ignored and only the term $\Lambda^{2} / r^{2}$ survives. The Schrödinger equation then becomes

$$
-\frac{\hbar^{2}}{2 m} \frac{1}{r^{2}} \Lambda^{2} \psi(\theta, \phi)=E \psi(\theta, \phi)
$$

The term $m r^{2}$ in the denominator can be recognized as the moment of inertia, $I$, of the particle, so the Schrödinger equation is

\[
-\frac{\hbar^{2}}{2 I} \Lambda^{2} \psi(\theta, \phi)=E \psi(\theta, \phi) \quad \begin{align*}
& \text { Schrödinger equation }  \tag{7F.8}\\
& \text { [particle on a sphere] }
\end{align*}
\]

There are two cyclic boundary conditions to fulfil. The first is the same as for the two-dimensional case, where the wavefunction must join up on completing a circuit around the equator, as specified by the angle $\phi$. The second is a similar requirement that the wavefunction must join up on encircling over the poles, as specified by the angle $\theta$. These two conditions are illustrated in Fig. 7F.4. Once again, it can be shown that the need to satisfy them leads to the conclusion that the energy and the angular momentum are quantized.

\subsection*{The chemist's toolkit 21 \\
 Spherical polar coordinates}
The mathematics of systems with spherical symmetry (such as atoms) is often greatly simplified by using spherical polar coordinates (Sketch 1): $r$, the distance from the origin (the radius), $\theta$, the colatitude, and $\phi$, the azimuth. The ranges of these coordinates are (with angles in radians, Sketch 2): $0 \leq r \leq+\infty, 0 \leq \theta$ $\leq \pi, 0 \leq \phi \leq 2 \pi$.\\
\includegraphics[max width=\textwidth, center]{2025_02_27_d4b95d9718f0b9e2e6b9g-318}

Sketch 1\\
\includegraphics[max width=\textwidth, center]{2025_02_27_d4b95d9718f0b9e2e6b9g-318(1)}

Figure 7F. 4 The wavefunction of a particle on the surface of a sphere must satisfy two cyclic boundary conditions. This requirement leads to two quantum numbers for its state of angular momentum.

\subsection*{How is that done? 7F. 3 Finding the solutions of the}
Schrödinger equation for a particle on a sphere\\
The functions known as spherical harmonics, $Y_{l, m_{l}}(\theta, \phi)$ (Table 7F.1), are well known to mathematicians and are the solutions of the equation ${ }^{1}$


\begin{align*}
& \Lambda^{2} Y_{l, m_{l}}(\theta, \phi)=-l(l+1) Y_{l, m_{l}}(\theta, \phi) \\
& l=0,1,2, \ldots m_{l}=0, \pm 1, \ldots, \pm l \tag{7F.9}
\end{align*}


${ }^{1}$ See the first section of $A$ deeper look 3 on the website for this text for details of how the separation of variables procedure is used to find the form of the spherical harmonics.\\
\includegraphics[max width=\textwidth, center]{2025_02_27_d4b95d9718f0b9e2e6b9g-318(2)}

Sketch 2\\
Cartesian and polar coordinates are related by

$$
x=r \sin \theta \cos \phi \quad y=r \sin \theta \sin \phi \quad z=r \cos \theta
$$

The volume element in Cartesian coordinates is $\mathrm{d} \tau=\mathrm{d} x \mathrm{~d} y \mathrm{~d} z$, and in spherical polar coordinates it becomes

$$
\mathrm{d} \tau=r^{2} \sin \theta \mathrm{~d} r \mathrm{~d} \theta \mathrm{~d} \phi
$$

An integral of a function $f(r, \theta, \phi)$ over all space in polar coordinates therefore has the form

$$
\int f \mathrm{~d} \tau=\int_{r=0}^{\infty} \int_{\theta=0}^{\pi} \int_{\phi=0}^{2 \pi} f(r, \theta, \phi) r^{2} \sin \theta \mathrm{~d} r \mathrm{~d} \theta \mathrm{~d} \phi
$$

\subsection*{Table 7F. 1 The spherical harmonics}
\begin{center}
\begin{tabular}{|c|c|c|}
\hline
$l$ & $m_{l}$ & $\boldsymbol{Y}_{l, m_{l}}(\theta, \phi)$ \\
\hline
0 & 0 & $\left(\frac{1}{4 \pi}\right)^{1 / 2}$ \\
\hline
1 & 0 & $\left(\frac{3}{4 \pi}\right)^{1 / 2} \cos \theta$ \\
\hline
\multirow[b]{2}{*}{2} & $\pm 1$ & $\mp\left(\frac{3}{8 \pi}\right)^{1 / 2} \sin \theta \mathrm{e}^{ \pm \text {i }}$ \\
\hline
 & 0 & $\left(\frac{5}{16 \pi}\right)^{1 / 2}\left(3 \cos ^{2} \theta-1\right)$ \\
\hline
\multirow{6}{*}{3} & $\pm 1$ & $\mp\left(\frac{15}{8 \pi}\right)^{1 / 2} \cos \theta \sin \theta \mathrm{e}^{ \pm \mathrm{i} \phi}$ \\
\hline
 & $\pm 2$ & $\left(\frac{15}{32 \pi}\right)^{1 / 2} \sin ^{2} \theta \mathrm{e}^{ \pm 2 i \phi}$ \\
\hline
 & 0 & $\left(\frac{7}{16 \pi}\right)^{1 / 2}\left(5 \cos ^{3} \theta-3 \cos \theta\right)$ \\
\hline
 & $\pm 1$ & $\mp\left(\frac{21}{64 \pi}\right)^{1 / 2}\left(5 \cos ^{2} \theta-1\right) \sin \theta \mathrm{e}^{ \pm i \phi}$ \\
\hline
 & $\pm 2$ & $\left(\frac{105}{32 \pi}\right)^{1 / 2} \sin ^{2} \theta \cos \theta \mathrm{e}^{ \pm 2 i \phi}$ \\
\hline
 & $\pm 3$ & $\mp\left(\frac{35}{64 \pi}\right)^{1 / 2} \sin ^{3} \theta \mathrm{e}^{ \pm 3 i \phi}$ \\
\hline
\end{tabular}
\end{center}

These functions satisfy the two cyclic boundary conditions and are normalized (to 1).

Step 1 Show that the spherical harmonics solve the Schrödinger equation\\
It follows from eqn 7F. 8 that

$$
-\frac{\hbar^{2}}{2 I} \overbrace{\Lambda^{2} Y_{l, m_{l}}(\theta, \phi)}^{-l(l+1) Y_{l, m_{l}}}=\overbrace{l(l+1) \frac{\hbar^{2}}{2 I}}^{E} Y_{l, m_{l}}(\theta, \phi)
$$

The spherical harmonics are therefore solutions of the Schrödinger equation with energies $E=l(l+1) \hbar^{2} / 2 I$. Note that the energies depend only on $l$ and not on $m_{l}$.

Step 2 Show that the wavefunctions are also eigenfunctions of the $z$-component of angular momentum\\
The operator for the $z$-component of angular momentum is $\hat{l}_{z}=(\hbar / \mathrm{i}) \partial / \partial \phi$. From Table 7F. 1 note that each spherical harmonic is of the form $Y_{l, m_{l}}(\theta, \phi)=\mathrm{e}^{\mathrm{i} m \phi} f(\theta)$. It then follows that

$$
\begin{aligned}
& \hat{l}_{z} Y_{l, m_{l}}(\theta, \phi)=\hat{l}_{z} \overbrace{\mathrm{e}^{\mathrm{i} m_{l} \phi}}^{Y_{l, m_{l}}(\theta, \phi)} \\
&=m_{l} \hbar \times Y_{l, m_{l}}(\theta) \\
&(\theta) \frac{\hbar}{i} \frac{\partial}{\partial \phi} \mathrm{e}^{\mathrm{i} m, \phi} f(\theta)=m_{l} \hbar \times \mathrm{e}^{\mathrm{i} m, \phi} f(\theta) \\
&
\end{aligned}
$$

Therefore, the $Y_{l, m_{l}}(\theta, \phi)$ are eigenfunctions of $\hat{l}_{z}$ with eigenvalues $m_{l} \hbar$.

In summary, the $Y_{l, m_{l}}(\theta, \phi)$ are solutions to the Schrödinger equation for a particle on a sphere, with the corresponding energies given by

$$
-E_{l, m_{l}}=l(l+1) \frac{\hbar^{2}}{2 I} \quad l=0,1,2 \ldots m_{l}=0, \pm 1, \ldots \pm l \left\lvert\, \begin{aligned}
& \text { Energy levels } \\
& \begin{array}{c}
\text { [particle on a } \\
\text { sphere] }
\end{array}
\end{aligned}\right.
$$

The integers $l$ and $m_{l}$ are now identified as quantum numbers: $l$ is the orbital angular momentum quantum number and $m_{l}$ is the magnetic quantum number. The energy is specified by $l$ alone, but for each value of $l$ there are $2 l+1$ values of $m_{l}$, so each energy level is $(2 l+1)$-fold degenerate. Each wavefunction is also an eigenfunction of $\hat{l}_{z}$ and therefore corresponds to a definite value, $m_{l} \hbar$, of the $z$-component of the angular momentum.

Figure 7F. 5 shows a representation of the spherical harmonics for $l=0-4$ and $m_{l}=0$. The use of different colours for different signs of the wavefunction emphasizes the location of the angular nodes (the positions at which the wavefunction passes through zero). Note that:

\begin{itemize}
  \item There are no angular nodes around the $z$-axis for functions with $m_{l}=0$. The spherical harmonic with $l=0, m_{l}=0$ has no nodes: it has a constant value at all positions of the surface and corresponds to a stationary particle.
  \item The number of angular nodes for states with $m_{l}=0$ is equal to $l$. As the number of nodes increases, the wavefunctions become more buckled, and with this increasing curvature the kinetic energy of the particle increases.\\
\includegraphics[max width=\textwidth, center]{2025_02_27_d4b95d9718f0b9e2e6b9g-319}
\end{itemize}

Figure 7F. 5 A representation of the wavefunctions of a particle on the surface of a sphere that emphasizes the location of angular nodes: blue and grey shading correspond to different signs of the wavefunction. Note that the number of nodes increases as the value of $/$ increases. All these wavefunctions correspond to $m_{l}=0$; a path round the vertical $z$-axis of the sphere does not cut through any nodes.

\subsection*{According to eqn 7F.10,}
\begin{itemize}
  \item Because $l$ is confined to non-negative integral values, the energy is quantized.
  \item The energies are independent of the value of $m_{l}$, because the energy is independent of the direction of the rotational motion.
  \item There are $2 l+1$ different wavefunctions (one for each value of $m_{l}$ ) that correspond to the same energy, so it follows that a level with quantum number $l$ is $(2 l+1)$-fold degenerate.
  \item There is no zero-point energy: $E_{0,0}=0$.
\end{itemize}

\subsection*{Example 7F. 2 Using the rotational energy levels}
The particle on a sphere is a good starting point for developing a model for the rotation of a diatomic molecule. Treat the rotation of ${ }^{1} \mathrm{H}^{127} \mathrm{I}$ as a hydrogen atom rotating around a stationary I atom (this is a good first approximation as the I atom is so heavy it hardly moves). The bond length is 160 pm . Evaluate the energies and degeneracies of the lowest four rotational energy levels of ${ }^{1} \mathrm{H}^{127}$ I. What is the frequency of the transition between the lowest two rotational levels?

Collect your thoughts The moment of inertia is $I=m_{1_{\mathrm{H}}} R^{2}$, with $R=160 \mathrm{pm}$; the rotational energies are given in eqn 7F.10. When describing the rotational energy levels of a molecule it is usual to denote the angular momentum quantum number by $J$ rather than $l$; as a result the degeneracy is $2 J+1$ (the analogue of $2 l+1$ ). A transition between two rotational levels can be brought about by the emission\\
or absorption of a photon with a frequency given by the Bohr frequency condition (Topic 7A, $h v=\Delta E$ ).

The solution The moment of inertia is

$$
I=\overbrace{\left(1.675 \times 10^{-27} \mathrm{~kg}\right)}^{m_{1_{H}}} \times \overbrace{\left(1.60 \times 10^{-12} \mathrm{~m}\right)^{2}}^{R^{2}}=4.29 \times 10^{-47} \mathrm{~kg} \mathrm{~m}^{2}
$$

It follows that

$$
\frac{\hbar^{2}}{2 I}=\frac{\left(1.055 \times 10^{-34} \mathrm{Js}\right)^{2}}{2 \times\left(4.29 \times 10^{-47} \mathrm{~kg} \mathrm{~m}^{2}\right)}=1.30 \times 10^{-22} \mathrm{~J}
$$

or 0.130 zJ . Draw up the following table, where the molar energies are obtained by multiplying the individual energies by Avogadro's constant:

\begin{center}
\begin{tabular}{llcl}
\hline
$\boldsymbol{J}$ & $\boldsymbol{E} / \mathbf{z J}$ & $\boldsymbol{E} /\left(\mathbf{J ~ m o l}^{-1}\right)$ & Degeneracy \\
\hline
0 & 0 & 0 & 1 \\
1 & 0.260 & 156 & 3 \\
2 & 0.780 & 470 & 5 \\
3 & 1.56 & 939 & 7 \\
\hline
\end{tabular}
\end{center}

The energy separation between the two lowest rotational energy levels ( $J=0$ and 1 ) is $2.60 \times 10^{-22}$ J, which corresponds to a photon of frequency

$$
v=\frac{\Delta E}{h}=\frac{2.60 \times 10^{-22} \mathrm{~J}}{6.626 \times 10^{-34} \mathrm{Js}}=3.92 \times 10^{11} \stackrel{\mathrm{~s}^{-1}}{\mathrm{~Hz}}=392 \mathrm{GHz}
$$

Comment. Radiation of this frequency belongs to the microwave region of the electromagnetic spectrum, so microwave spectroscopy is used to study molecular rotations (Topic 11B). Because the transition frequencies depend on the moment of inertia and frequencies can be measured with great precision, microwave spectroscopy is a very precise technique for the determination of bond lengths.

Self-test 7F. 2 What is the frequency of the transition between the lowest two rotational levels in ${ }^{2} \mathrm{H}^{127} \mathrm{I}$ ? (Assume that the bond length is the same as for ${ }^{1} \mathrm{H}^{127} \mathrm{I}$ and that the iodine atom is stationary.)\\
\includegraphics[max width=\textwidth, center]{2025_02_27_d4b95d9718f0b9e2e6b9g-320}

\subsection*{(b) Angular momentum}
According to classical mechanics (The chemist's toolkit 20) the kinetic energy of a particle circulating on a ring is $E_{\mathrm{k}}=J^{2} / 2 I$, where $J$ is the magnitude of the angular momentum. By comparing this relation with eqn 7F.10, it follows that the square of the magnitude of the angular momentum is $l(l+1) \hbar^{2}$, so the magnitude of the angular momentum is


\begin{align*}
& \text { Magnitude }=\{l(l+1)\}^{1 / 2} \hbar \\
& \quad l=0,1,2 \ldots \tag{7F.11}
\end{align*}


Magnitude of angular momentum

The spherical harmonics are also eigenfunctions of $\hat{l}_{z}$ with eigenvalues

\[
\begin{array}{cl}
z \text {-Component }=m_{l} \hbar & \text { z-Component of angular } \\
m_{l}=0, \pm 1, \ldots \pm l & \text { momentum } \tag{7F.12}
\end{array}
\]

So, both the magnitude and the $z$-component of angular momentum are quantized.

\subsection*{Brief illustration 7F. 1}
The lowest four rotational energy levels of any object rotating in three dimensions correspond to $l=0,1,2,3$. The following table can be constructed by using eqns 7F.11 and 7F.12.

\begin{center}
\begin{tabular}{llll}
\hline
$\boldsymbol{l}$ & \begin{tabular}{l}
Magnitude of angular \\
momentum $/ \hbar$ \\
\end{tabular} & Degeneracy & \begin{tabular}{l}
$z$-Component \\
of angular \\
momentum $/ \hbar$ \\
\end{tabular} \\
\hline
0 & 0 & 1 & 0 \\
1 & $2^{1 / 2}$ & 3 & $0, \pm 1$ \\
2 & $6^{1 / 2}$ & 5 & $0, \pm 1, \pm 2$ \\
3 & $12^{1 / 2}$ & 7 & $0, \pm 1, \pm 2, \pm 3$ \\
\hline
\end{tabular}
\end{center}

\subsection*{(c) The vector model}
The result that $m_{l}$ is confined to the values $0, \pm 1, \ldots \pm l$ for a given value of $l$ means that the component of angular momentum about the $z$-axis-the contribution to the total angular momentum of rotation around that axis-may take only $2 l+1$ values. If the angular momentum is represented by a vector of length $\{l(l+1)\}^{1 / 2}$, it follows that this vector must be oriented so that its projection on the $z$-axis is $m_{l}$ and that it can have only $2 l+1$ orientations rather than the continuous range of orientations of a rotating classical body (Fig. 7F.6). The remarkable implication is that

\subsection*{The orientation of a rotating body is quantized.}
The quantum mechanical result that a rotating body may not take up an arbitrary orientation with respect to some specified axis (e.g. an axis defined by the direction of an externally applied electric or magnetic field) is called space quantization.

The preceding discussion has referred to the $z$-component of angular momentum and there has been no reference to the $x$ - and $y$-components. The reason for this omission is found by examining the operators for the three components, each one being given by a term like that in eqn 7F.5a: ${ }^{2}$

\footnotetext{${ }^{2}$ Each one is in fact a component of the vector product of $\boldsymbol{r}$ and $\boldsymbol{p}, \boldsymbol{l}=\boldsymbol{r} \times \boldsymbol{p}$, and the replacement of $\boldsymbol{r}$ and $\boldsymbol{p}$ by their operator equivalents.
}
\includegraphics[max width=\textwidth, center]{2025_02_27_d4b95d9718f0b9e2e6b9g-321(1)}

Figure 7F. 6 The permitted orientations of angular momentum when $I=2$. This representation is too specific because the azimuthal orientation of the vector (its angle around $z$ ) is indeterminate.


\begin{align*}
& \hat{l}_{x}=\frac{\hbar}{\mathrm{i}}\left(y \frac{\partial}{\partial z}-z \frac{\partial}{\partial y}\right) \\
& \hat{l}_{y}=\frac{\hbar}{\mathrm{i}}\left(z \frac{\partial}{\partial x}-x \frac{\partial}{\partial z}\right) \quad \text { Angular momentum operators }  \tag{7F.13}\\
& \hat{l}_{z}=\frac{\hbar}{\mathrm{i}}\left(x \frac{\partial}{\partial y}-y \frac{\partial}{\partial x}\right)
\end{align*}


Each of these expressions can be derived in the same way as eqn 7F.5a by converting the classical expressions for the components of the angular momentum into their quantum mechanical equivalents. The commutation relations among the three operators (Problem P7F.9), are

\[
\left[\hat{l}_{x}, \hat{l}_{y}\right]=\mathrm{i} \hbar \hat{l}_{z} \quad\left[\hat{l}_{y}, \hat{l}_{z}\right]=\mathrm{i} \hbar \hat{l}_{x} \quad\left[\hat{l}_{z}, \hat{l}_{x}\right]=\mathrm{i} \hbar \hat{l}_{y} \begin{align*}
& \text { Angular } \\
& \begin{array}{l}
\text { momentum } \\
\text { commatation }
\end{array}  \tag{7F.14}\\
& \text { relations }
\end{align*}
\]

Because the three operators do not commute, they represent complementary observables (Topic 7C). Therefore, the more precisely any one component is known, the greater the uncertainty in the other two. It is possible to have precise knowledge of only one of the components of the angular momentum, so if $l_{z}$ is specified exactly (as in the preceding discussion), neither $l_{x}$ nor $l_{y}$ can be specified.

The operator for the square of the magnitude of the angular momentum is

$$
\hat{l}^{2}=\hat{l}_{x}^{2}+\hat{l}_{y}^{2}+\hat{l}_{z}^{2} \quad \begin{aligned}
& \text { Operator for the square of the } \\
& \text { magnitude of angular momentum }
\end{aligned}
$$

(7F.15)\\
\includegraphics[max width=\textwidth, center]{2025_02_27_d4b95d9718f0b9e2e6b9g-321}

Figure 7F. 7 (a) A summary of Fig. 7F.6. However, because the azimuthal angle of the vector around the $z$-axis is indeterminate, a better representation is as in (b), where each vector lies at an unspecified azimuthal angle on its cone.

This operator commutes with all three components (Problem P7F.11):

\[
\left[\hat{l}^{2}, \hat{l}_{q}\right]=0 \quad q=x, y \text {, and } z \quad \begin{align*}
& \text { Commutators of angular }  \tag{7F.16}\\
& \text { momentum operators }
\end{align*}
\]

It follows that both the square magnitude and one component, commonly the $z$-component, of the angular momentum can be specified precisely. The illustration in Fig. 7F.6, which is summarized in Fig. 7F.7(a), therefore gives a false impression of the state of the system, because it suggests definite values for the $x$ - and $y$-components too. A better picture must reflect the impossibility of specifying $l_{x}$ and $l_{y}$ if $l_{z}$ is known.

The vector model of angular momentum uses pictures like that in Fig. 7F.7(b). The cones are drawn with side $\{l(l+1)\}^{1 / 2}$ units, and represent the magnitude of the angular momentum. Each cone has a definite projection (of $m_{l}$ units) on to the $z$-axis, representing the precisely known value of $l_{z}$. The projections of the vector on to the $x$ - and $y$-axes, which give the values of $l_{x}$ and $l_{y}$, are indefinite: the vector representing angular momentum can be thought of as lying with its tip on any point on the mouth of the cone. At this stage it should not be thought of as sweeping round the cone; that aspect of the model will be added when the picture is allowed to convey more information (Topics 8B and 8C).

\subsection*{Brief illustration 7F. 2}
If the wavefunction of a rotating molecule is given by the spherical harmonic $Y_{3,2}$ then the angular momentum can be represented by a cone

\begin{itemize}
  \item with a side of length $12^{1 / 2}$ (representing the magnitude of $12^{1 / 2} \hbar$ ); and
  \item with a projection of +2 on the $z$-axis (representing the $z$-component of $+2 \hbar$ ).
\end{itemize}

\subsection*{Checklist of concepts}
$\square$ 1. The energy and angular momentum for a particle rotating in two- or three-dimensions are quantized; quantization results from the requirement that the wavefunction satisfies cyclic boundary conditions.2. All energy levels of a particle rotating in two dimensions are doubly-degenerate except for the lowest level ( $m_{l}=0$ ).3. There is no zero-point energy for a rotating particle.4. It is impossible to specify simultaneously the angular momentum and location of a particle with arbitrary precision.5. For a particle rotating in three dimensions, the cyclic boundary conditions imply that the magnitude and $z$-component of the angular momentum are quantized.6. Space quantization refers to the quantum mechanical result that a rotating body may not take up an arbitrary orientation with respect to some specified axis.7. The three components of the angular momentum are mutually complementary observables.8. Because the operators that represent the components of angular momentum do not commute, only the magnitude of the angular momentum and one of its components can be specified simultaneously with arbitrary precision.9. In the vector model of angular momentum, the angular momentum is represented by a cone with a side of length $\{l(l+1)\}^{1 / 2}$ and a projection of $m_{l}$ on the $z$-axis. The vector can be thought of as lying with its tip on an indeterminate point on the mouth of the cone.

\subsection*{Checklist of equations}
\textbackslash begin\{tabular\}\{|c|c|c|c|\}\\
\textbackslash hline Property \& Equation \& Comment \& Equation number \\
\\
\textbackslash hline Wavefunction of particle on a ring \& $\psi_{m_{l}}(\phi)=\mathrm{e}^{\mathrm{i} m m_{l} /} /(2 \pi)^{1 / 2}$ \& $m_{l}=0, \pm 1, \pm 2, \ldots$ \& 7F. 4 \\


\textbackslash hline Energy of particle on a ring \& $E_{m_{l}}=m_{l}^{2} \hbar^{2} / 2 I$ \& $$
\begin{aligned}
& m_{l}=0, \pm 1, \pm 2, \ldots \\
& I=m r^{2}
\end{aligned}
$$ \& 7F. 4 \\


\textbackslash hline $z$-Component of angular momentum of particle on a ring \& $m_{l} \hbar$ \& $m_{l}=0, \pm 1, \pm 2, \ldots$ \& 7F. 6 \\
\\
\textbackslash hline Wavefunction of particle on a sphere \& $\psi(\theta, \phi)=Y_{l, m_{l}}(\theta, \phi)$ \& $Y$ is a spherical harmonic (Table 7F.1) \& \\
\\
\textbackslash hline Energy of particle on a sphere \& $E_{l, m_{l}}=l(l+1) \hbar^{2} / 2 I$ \& $l=0,1,2, \ldots$ \& 7F. 10 \\
\\
\textbackslash hline Magnitude of angular momentum \& $\{l(l+1)\}^{1 / 2} \hbar$ \& $l=0,1,2, \ldots$ \& 7F. 11 \\
\\
\textbackslash hline $z$-Component of angular momentum \& $m_{l} \hbar$ \& $m_{l}=0, \pm 1, \pm 2, \ldots \pm l$ \& 7F. 12 \\


\textbackslash hline Angular momentum commutation relations \& $$
\begin{aligned}
& {\left[\hat{l}_{x}, \hat{l}_{y}\right]=\mathrm{i} \hbar \hat{l}_{z}} \\
& {\left[\hat{l_{y}}, \hat{l}_{z}\right]=\mathrm{i} \hbar \hat{l}_{x}} \\
& {\left[\hat{l}_{z}, \hat{l}_{x}\right]=\mathrm{i} \hbar \hat{l}_{y}} \\
& {\left[\hat{l}^{2}, \hat{l}_{q}\right]=0, q=x, y, \text { and } z}
\end{aligned}
$$ \& \& $7 F .14$

$7 F .16$ \\
\\
\textbackslash hline\\
\textbackslash end\{tabular\}

\chapter{FOCUS 8 Atomic structure and spectra}
This Focus discusses the use of quantum mechanics to describe and investigate the 'electronic structure' of atoms, the arrangement of electrons around their nuclei. The concepts are of central importance for understanding the properties of atoms and molecules, and hence have extensive chemical applications.

\subsection*{8A Hydrogenic atoms}
This Topic uses the principles of quantum mechanics introduced in Focus 7 to describe the electronic structure of a 'hydrogenic atom', a one-electron atom or ion of general atomic number $Z$. Hydrogenic atoms are important because their Schrödinger equations can be solved exactly and they provide a set of concepts that are used to describe the structures of many-electron atoms and molecules. Solving the Schrödinger equation for an electron in an atom involves the separation of the wavefunction into angular and radial parts and the resulting wavefunctions are the hugely important 'atomic orbitals' of hydrogenic atoms.\\
8A. 1 The structure of hydrogenic atoms; 8A. 2 Atomic orbitals and their energies

\subsection*{8B Many-electron atoms}
A 'many-electron atom' is an atom or ion with more than one electron. Examples include all neutral atoms other than H ; so even He , with only two electrons, is a many-electron atom.

This Topic uses hydrogenic atomic orbitals to describe the structures of many-electron atoms. Then, in conjunction with the concept of 'spin' and the 'Pauli exclusion principle', it describes the origin of the periodicity of atomic properties and the structure of the periodic table.\\
8B. 1 The orbital approximation; 8B. 2 The Pauli exclusion principle;\\
8B. 3 The building-up principle; 8 B. 4 Self-consistent field orbitals

\subsection*{8C Atomic spectra}
The spectra of many-electron atoms are more complicated than that of hydrogen. Similar principles apply, but Coulombic and magnetic interactions between the electrons give rise to a variety of energy differences, which are summarized by constructing 'term symbols'. These symbols act as labels that display the total orbital and spin angular momentum of a many-electron atom and are used to express the selection rules that govern their spectroscopic transitions.\\
8C. 1 The spectra of hydrogenic atoms; 8C. 2 The spectra of manyelectron atoms

\subsection*{Web resource What is an application of this material?}
Impact 13 focuses on the use of atomic spectroscopy to examine stars. By analysing their spectra it is possible to determine the composition of their outer layers and the surrounding gases and to determine features of their physical state.

\section*{TOPIC 8A Hydrogenic atoms}
\subsection*{Why do you need to know this material?}
An understanding of the structure of hydrogenic atoms is central to the description of all other atoms, the periodic table, and bonding. All accounts of the structures of molecules are based on the language and concepts introduced here.

\subsection*{What is the key idea?}
Atomic orbitals are one-electron wavefunctions for atoms and are labelled by three quantum numbers that specify the energy and angular momentum of the electron.

\subsection*{What do you need to know already?}
You need to be aware of the concept of a wavefunction (Topic 7B) and its interpretation. You also need to know how to set up a Schrödinger equation and how boundary conditions result in only certain solutions being acceptable (Topic 7D).

When an electric discharge is passed through gaseous hydrogen, the $\mathrm{H}_{2}$ molecules are dissociated and the energetically excited H atoms that are produced emit electromagnetic radiation at a number of discrete frequencies (and therefore discrete wavenumbers), producing a spectrum of a series of 'lines' (Fig. 8A.1).\\
\includegraphics[max width=\textwidth, center]{2025_02_27_d4b95d9718f0b9e2e6b9g-336}

Figure 8A. 1 The spectrum of atomic hydrogen. Both the observed spectrum and its resolution into overlapping series are shown. Note that the Balmer series lies in the visible region.

The Swedish spectroscopist Johannes Rydberg noted (in 1890) that the wavenumbers of all the lines are given by the expression


\begin{equation*}
\tilde{v}=\tilde{R}_{\mathrm{H}}\left(\frac{1}{n_{1}^{2}}-\frac{1}{n_{2}^{2}}\right) \quad \text { Spectral lines of a hydrogen atom } \tag{8A.1}
\end{equation*}


with $n_{1}=1$ (the Lyman series), 2 (the Balmer series), and 3 (the Paschen series), and that in each case $n_{2}=n_{1}+1, n_{1}+2, \ldots$. The constant $\tilde{R}_{\mathrm{H}}$ is now called the Rydberg constant for the hydrogen atom and is found empirically to have the value $109677 \mathrm{~cm}^{-1}$.

\subsection*{8A.1 The structure of hydrogenic atoms}
Consider a hydrogenic atom, an atom or ion of arbitrary atomic number $Z$ but having a single electron. Hydrogen itself is an example (with $Z=1$ ). The Coulomb potential energy of an electron in a hydrogenic atom of atomic number $Z$ and therefore nuclear charge $Z e$ is


\begin{equation*}
V(r)=-\frac{Z e^{2}}{4 \pi \varepsilon_{0} r} \tag{8A.2}
\end{equation*}


where $r$ is the distance of the electron from the nucleus and $\varepsilon_{0}$ is the vacuum permittivity. The hamiltonian for the entire atom, which consists of an electron and a nucleus of mass $m_{\mathrm{N}}$, is therefore

\[
\begin{array}{rlrl}
\hat{H} & =\hat{E}_{\mathrm{k}, \text { electron }}+\hat{E}_{\mathrm{k}, \text { nucleus }}+\hat{V}(r) & & \\
& =-\frac{\hbar^{2}}{2 m_{\mathrm{e}}} \nabla_{\mathrm{e}}^{2}-\frac{\hbar^{2}}{2 m_{\mathrm{N}}} \nabla_{\mathrm{N}}^{2}-\frac{Z e^{2}}{4 \pi \varepsilon_{0} r} & & \text { hymiltonian for a }  \tag{8A.3}\\
\text { hydrogenic atom }
\end{array}
\]

The subscripts e and N on $\nabla^{2}$ indicate differentiation with respect to the electron or nuclear coordinates.

\subsection*{(a) The separation of variables}
Physical intuition suggests that the full Schrödinger equation ought to separate into two equations, one for the motion of the atom as a whole through space and the other for the motion of the electron relative to the nucleus. The Schrödinger\\
equation for the internal motion of the electron relative to the nucleus is ${ }^{1}$

$$
\begin{aligned}
& -\frac{\hbar^{2}}{2 \mu} \nabla^{2} \psi-\frac{Z e^{2}}{4 \pi \varepsilon_{0} r} \psi=E \psi \\
& \frac{1}{\mu}=\frac{1}{m_{\mathrm{e}}}+\frac{1}{m_{\mathrm{N}}}
\end{aligned}
$$

Schrödinger\\
equation for a\\
hydrogenic atom\\
where differentiation is now with respect to the coordinates of the electron relative to the nucleus. The quantity $\mu$ is called the reduced mass. The reduced mass is very similar to the electron mass because $m_{\mathrm{N}}$, the mass of the nucleus, is much larger than the mass of an electron, so $1 / \mu \approx 1 / m_{\mathrm{e}}$ and therefore $\mu \approx m_{\mathrm{e}}$. In all except the most precise work, the reduced mass can be replaced by $m_{\mathrm{e}}$.

Because the potential energy is centrosymmetric (independent of angle), the equation for the wavefunction is expected to be separable into radial and angular components, as in


\begin{equation*}
\psi(r, \theta, \phi)=R(r) Y(\theta, \phi) \tag{8A.5}
\end{equation*}


with $R(r)$ the radial wavefunction and $Y(\theta, \phi)$ the angular wavefunction. The equation does separate, and the two contributions to the wavefunction are solutions of two equations:


\begin{align*}
& \Lambda^{2} Y=-l(l+1) Y  \tag{8А.6a}\\
& -\frac{\hbar^{2}}{2 \mu}\left(\frac{\mathrm{~d}^{2} R}{\mathrm{~d} r^{2}}+\frac{2}{r} \frac{\mathrm{~d} R}{\mathrm{~d} r}\right)+V_{\mathrm{eff}} R=E R \tag{8A.6b}
\end{align*}


where


\begin{equation*}
V_{\text {eff }}(r)=-\frac{Z e^{2}}{4 \pi \varepsilon_{0} r}+\frac{l(l+1) \hbar^{2}}{2 \mu r^{2}} \tag{8A.6c}
\end{equation*}


Equation 8A.6a is the same as the Schrödinger equation for a particle free to move at constant radius around a central point, and is considered in Topic 7F. The allowed solutions are the spherical harmonics (Table 7F.1), and are specified by the quantum numbers $l$ and $m_{l}$. Equation 8A.6b is called the radial wave equation. The radial wave equation describes the motion of a particle of mass $\mu$ in a one-dimensional region $0 \leq$ $r<\infty$ where the potential energy is $V_{\text {eff }}(r)$.

\subsection*{(b) The radial solutions}
Some features of the shapes of the radial wavefunctions can be anticipated by examining the form of $V_{\text {eff }}(r)$. The first term in eqn 8 A .6 c is the Coulomb potential energy of the electron in the field of the nucleus. The second term stems from what

\footnotetext{${ }^{1}$ See the first section of $A$ deeper look 3 on the website for this text for full details of this separation procedure and then the second section for the calculations that lead to eqn 8A. 6 .
}
\includegraphics[max width=\textwidth, center]{2025_02_27_d4b95d9718f0b9e2e6b9g-337}

Figure 8A. 2 The effective potential energy of an electron in the hydrogen atom. When the electron has zero orbital angular momentum, the effective potential energy is the Coulombic potential energy. When the electron has non-zero orbital angular momentum, the centrifugal effect gives rise to a positive contribution which is very large close to the nucleus. The $I=0$ and $I \neq 0$ wavefunctions are therefore very different near the nucleus.\\
in classical physics would be called the centrifugal force arising from the angular momentum of the electron around the nucleus. When $l=0$, the electron has no angular momentum, and the effective potential energy is purely Coulombic and the force exerted on the electron is attractive at all radii (Fig. 8A.2). When $l \neq 0$, the centrifugal term gives a positive contribution to the effective potential energy, corresponding to a repulsive force at all radii. When the electron is close to the nucleus ( $r \approx 0$ ), the latter contribution to the potential energy, which is proportional to $1 / r^{2}$, dominates the Coulombic contribution, which is proportional to $1 / r$, and the net result is an effective repulsion of the electron from the nucleus. The two effective potential energies, the one for $l=0$ and the one for $l \neq 0$, are therefore qualitatively very different close to the nucleus. However, they are similar at large distances because the centrifugal contribution tends to zero more rapidly (as $1 / r^{2}$ ) than the Coulombic contribution (as $1 / r$ ). Therefore, the solutions with $l=0$ and $l \neq 0$ are expected to be quite different near the nucleus but similar far away from it.

Two features of the radial wavefunction are important:

\begin{itemize}
  \item Close to the nucleus the radial wavefunction is proportional to $r^{l}$, and the higher the orbital angular momentum, the less likely it is that the electron will be found there (Fig. 8A.3).
  \item Far from the nucleus all radial wavefunctions approach zero exponentially.
\end{itemize}

The detailed solution of the radial equation for the full range of radii shows how the form $r^{l}$ close to the nucleus blends\\
\includegraphics[max width=\textwidth, center]{2025_02_27_d4b95d9718f0b9e2e6b9g-338}

Figure 8A. 3 Close to the nucleus, orbitals with $I=1$ are proportional to $r$, orbitals with $/=2$ are proportional to $r^{2}$, and orbitals with $I=3$ are proportional to $r^{3}$. Electrons are progressively excluded from the neighbourhood of the nucleus as / increases. An orbital with / = 0 has a finite, non-zero value at the nucleus.\\
into the exponentially decaying form at great distances. It turns out that the two regions are bridged by a polynomial in $r$ and that

$$
R(r)=\overbrace{r^{l}}^{\begin{array}{c}
\text { Dominant } \\
\text { close to the } \\
\text { nucleus }
\end{array}} \times \overbrace{(\text { polynomial in } r)}^{\begin{array}{c}
\text { Bridges the two } \\
\text { ends of the function }
\end{array}} \times \overbrace{(\text { decaying exponential in } r)}^{\begin{array}{c}
\text { Dominant far } \\
\text { from the nucleus }
\end{array}}
$$

(8A.7)

The radial wavefunction therefore has the form

$$
R(r)=r^{l} L(r) \mathrm{e}^{-r}
$$

with various constants and where $L(r)$ is the bridging polynomial. Close to the nucleus $(r \approx 0)$ the polynomial is a constant and $\mathrm{e}^{-r} \approx 1$, so $R(r) \propto r^{l}$; far from the nucleus the dominant term in the polynomial is proportional to $r^{n-l-1}$, where $n$ is an integer, so regardless of the value of $l$, all the wavefunctions of a given value of $n$ are proportional to $r^{n-1} \mathrm{e}^{-r}$ and decay exponentially to zero in the same way (exponential functions $\mathrm{e}^{-x}$ always dominate simple powers, $x^{n}$ ).

The detailed solution also shows that, for the wavefunction to be acceptable, the value of $n$ that appears in the polynomial can take only positive integral values, and specifically $n=1,2, \ldots$. This number also determines the allowed energies through the expression:

$$
E_{n}=-\frac{\mu e^{4}}{32 \pi^{2} \varepsilon_{0}^{2} \hbar^{2}} \times \frac{Z^{2}}{n^{2}}
$$

Bound-state energies\\
(8A.8)

So far, only the general form of the radial wavefunctions has been given. It is now time to show how they depend on various fundamental constants and the atomic number of the atom. They are most simply written in terms of the dimensionless quantity $\rho$ (rho), where


\begin{equation*}
\rho=\frac{2 Z r}{n a} \quad a=\frac{m_{\mathrm{e}}}{\mu} a_{0} \quad a_{0}=\frac{4 \pi \varepsilon_{0} \hbar^{2}}{m_{\mathrm{e}} e^{2}} \tag{8A.9}
\end{equation*}


The Bohr radius, $a_{0}$, has the value 52.9 pm ; it is so called because the same quantity appeared in Bohr's early model of the hydrogen atom as the radius of the electron orbit of lowest energy. In practice, because $m_{\mathrm{e}} \ll m_{\mathrm{N}}$ (so $m_{\mathrm{e}} / \mu \approx 1$ ) there is so little difference between $a$ and $a_{0}$ that it is safe to use $a_{0}$ in the definition of $\rho$ for all atoms (even for ${ }^{1} \mathrm{H}, a$ $=1.0005 a_{0}$ ). In terms of these quantities and with the various quantum numbers displayed, the radial wavefunctions for an electron with quantum numbers $n$ and $l$ are the (real) functions


\begin{equation*}
R_{n, l}(r)=N_{n, l} \rho^{l} L_{n, l}(\rho) \mathrm{e}^{-\rho / 2} \quad \text { Radial wavefunctions } \tag{8A.10}
\end{equation*}


where $L_{n, l}(\rho)$ is an associated Laguerre polynomial. These polynomials have quite simple forms, such as $1, \rho$, and $2-\rho$ (they can be picked out in Table 8A.1). The factor $N_{n, l}$ ensures that the radial wavefunction is normalized to 1 in the sense that


\begin{equation*}
\int_{0}^{\infty} R_{n, l}(r)^{2} r^{2} \mathrm{~d} r=1 \tag{8A.11}
\end{equation*}


Table 8A. 1 Hydrogenic radial wavefunctions

\begin{center}
\begin{tabular}{lll}
\hline
$n$ & $l$ & $R_{n, l}(r)$ \\
1 & 0 & $2\left(\frac{Z}{a}\right)^{3 / 2} \mathrm{e}^{-\rho / 2}$ \\
2 & 0 & $\frac{1}{8^{1 / 2}}\left(\frac{Z}{a}\right)^{3 / 2}(2-\rho) \mathrm{e}^{-\rho / 2}$ \\
2 & 1 & $\frac{1}{24^{1 / 2}}\left(\frac{Z}{a}\right)^{3 / 2} \rho \mathrm{e}^{-\rho / 2}$ \\
3 & 0 & $\frac{1}{243^{1 / 2}}\left(\frac{Z}{a}\right)^{3 / 2}\left(6-6 \rho+\rho^{2}\right) \mathrm{e}^{-\rho / 2}$ \\
3 & 1 & $\frac{1}{486^{1 / 2}}\left(\frac{Z}{a}\right)^{3 / 2}(4-\rho) \rho \mathrm{e}^{-\rho / 2}$ \\
3 & 2 & $\frac{1}{2430^{1 / 2}}\left(\frac{Z}{a}\right)^{3 / 2} \rho^{2} \mathrm{e}^{-\rho / 2}$ \\
\hline
\end{tabular}
\end{center}

\footnotetext{$\rho=(2 Z / n a) r$ with $a=4 \pi \varepsilon_{0} \hbar^{2} / \mu e^{2}$. For an infinitely heavy nucleus (or one that may be assumed to be), $\mu=m_{\mathrm{e}}$ and $a=a_{0}$, the Bohr radius.
}
(The $r^{2}$ comes from the volume element in spherical coordinates; see The chemist's toolkit 21 in Topic 7F.) Specifically, the components of eqn 8 A .10 can be interpreted as follows:

\begin{itemize}
  \item The exponential factor ensures that the wavefunction approaches zero far from the nucleus.
  \item The factor $\rho^{l}$ ensures that (provided $l>0$ ) the wavefunction vanishes at the nucleus. The zero at $r=0$ is not a radial node because the radial wavefunction does not pass through zero at that point (because $r$ cannot be negative).
  \item The associated Laguerre polynomial is a function that in general oscillates from positive to negative values and accounts for the presence of radial nodes.
\end{itemize}

Expressions for some radial wavefunctions are given in Table 8A. 1 and illustrated in Fig. 8A.4. Finally, with the form of\\
the radial wavefunction established, the total wavefunction, eqn 8 A .5 , in full dress becomes


\begin{equation*}
\psi_{n, l, m_{l}}(r, \theta, \phi)=R_{n, l}(r) Y_{l, m_{l}}(\theta, \phi) \tag{8A.12}
\end{equation*}


\subsection*{Brief illustration 8A. 1}
To calculate the probability density at the nucleus for an electron with $n=1, l=0$, and $m_{l}=0$, evaluate $\psi$ at $r=0$ :

$$
\psi_{1,0,0}(0, \theta, \phi)=R_{1,0}(0) Y_{0,0}(\theta, \phi)=2\left(\frac{Z}{a_{0}}\right)^{3 / 2}\left(\frac{1}{4 \pi}\right)^{1 / 2}
$$

The probability density is therefore

$$
\psi_{1,0,0}^{2}(0, \theta, \phi)=\frac{Z^{3}}{\pi a_{0}^{3}}
$$

which evaluates to $2.15 \times 10^{-6} \mathrm{pm}^{-3}$ when $Z=1$.\\
\includegraphics[max width=\textwidth, center]{2025_02_27_d4b95d9718f0b9e2e6b9g-339}

Figure 8A. 4 The radial wavefunctions of the first few states of hydrogenic atoms of atomic number $Z$. Note that the orbitals with $/=0$ have a non-zero and finite value at the nucleus. The horizontal scales are different in each case: as the principal quantum number increases, so too does the size of the orbital.

\subsection*{8A.2 Atomic orbitals and their energies}
An atomic orbital is a one-electron wavefunction for an electron in an atom, and for hydrogenic atoms has the form specified in eqn 8A.12. Each hydrogenic atomic orbital is defined by three quantum numbers, designated $n, l$, and $m_{l}$. An electron described by one of the wavefunctions in eqn 8 A .12 is said to 'occupy' that orbital. For example, an electron described by the wavefunction $\psi_{1,0,0}$ is said to 'occupy' the orbital with $n=1$, $l=0$, and $m_{l}=0$.

\subsection*{(a) The specification of orbitals}
Each of the three quantum numbers specifies a different attribute of the orbital:

\begin{itemize}
  \item The principal quantum number, $n$, specifies the energy of the orbital (through eqn 8A.8); it takes the values $n=$ $1,2,3, \ldots$.
  \item The orbital angular momentum quantum number, $l$, specifies the magnitude of the angular momentum of the electron as $\{l(l+1)\}^{1 / 2} \hbar$, with $l=0,1,2, \ldots, n-1$.
  \item The magnetic quantum number, $m_{l}$, specifies the $z$-component of the angular momentum as $m_{l} \hbar$, with $m_{l}=0, \pm 1$, $\pm 2, \ldots, \pm l$.
\end{itemize}

Note how the value of the principal quantum number controls the maximum value of $l$, and how the value of $l$ controls the range of values of $m_{l}$.

\subsection*{(b) The energy levels}
The energy levels predicted by eqn 8A. 8 are depicted in Fig. 8A.5. The energies, and also the separation of neighbouring levels, are proportional to $Z^{2}$, so the levels are four times as wide apart (and the ground state four times lower in energy) in $\mathrm{He}^{+}(Z=2)$ than in $\mathrm{H}(Z=1)$. All the energies given by eqn 8 A .8 are negative. They refer to the bound states of the atom, in which the energy of the atom is lower than that of the infinitely separated, stationary electron and nucleus (which corresponds to the zero of energy). There are also solutions of the Schrödinger equation with positive energies. These solutions correspond to unbound states of the electron, the states to which an electron is raised when it is ejected from the atom by a high-energy collision or photon. The energies of the unbound electron are not quantized and form the continuum states of the atom.

Equation 8A.8, which can be written as


\begin{equation*}
E_{n}=-\frac{h c Z^{2} \tilde{R}_{\mathrm{N}}}{n^{2}} \quad \tilde{R}_{\mathrm{N}}=\frac{\mu e^{4}}{8 \varepsilon_{0}^{2} c h^{3}} \quad \text { Bound-state energies } \tag{8A.13}
\end{equation*}


\begin{center}
\includegraphics[max width=\textwidth]{2025_02_27_d4b95d9718f0b9e2e6b9g-340}
\end{center}

Figure 8A. 5 The energy levels of a hydrogen atom. The values are relative to an infinitely separated, stationary electron and a proton.\\
is consistent with the spectroscopic result summarized by eqn 8A.1, with the Rydberg constant for the atom identified as


\begin{equation*}
\tilde{R}_{\mathrm{N}}=\frac{\mu}{m_{\mathrm{e}}} \times \tilde{R}_{\infty} \quad \tilde{R}_{\infty}=\frac{m_{\mathrm{e}} e^{4}}{8 \varepsilon_{0}^{2} h^{3} c} \quad \text { Rydberg constant } \tag{8A.14}
\end{equation*}


where $\mu$ is the reduced mass of the atom and $\tilde{R}_{\infty}$ is the Rydberg constant; the constant $\tilde{R}_{\mathrm{N}}$ is the value that constant takes for a specified atom N (not nitrogen!), such as hydrogen, when N is replaced by H and $\mu$ takes the appropriate value. Insertion of the values of the fundamental constants into the expression for $\tilde{R}_{\mathrm{H}}$ gives almost exact agreement with the experimental value for hydrogen. The only discrepancies arise from the neglect of relativistic corrections (in simple terms, the increase of mass with speed), which the non-relativistic Schrödinger equation ignores.

\subsection*{Brief illustration 8A. 2}
The value of $\tilde{R}_{\infty}$ is given inside the front cover and is $109737 \mathrm{~cm}^{-1}$. The reduced mass of a hydrogen atom with $m_{\mathrm{p}}=$ $1.67262 \times 10^{-27} \mathrm{~kg}$ and $m_{\mathrm{e}}=9.10938 \times 10^{-31} \mathrm{~kg}$ is

$$
\begin{aligned}
\mu & =\frac{m_{\mathrm{e}} m_{\mathrm{p}}}{m_{\mathrm{e}}+m_{\mathrm{p}}}=\frac{\left(9.10938 \times 10^{-31} \mathrm{~kg}\right) \times\left(1.67262 \times 10^{-27} \mathrm{~kg}\right)}{\left(9.10938 \times 10^{-31} \mathrm{~kg}\right)+\left(1.67262 \times 10^{-27} \mathrm{~kg}\right)} \\
& =9.10442 \times 10^{-31} \mathrm{~kg}
\end{aligned}
$$

It then follows that

$$
\tilde{R}_{\mathrm{H}}=\frac{9.10442 \times 10^{-31} \mathrm{~kg}}{9.10938 \times 10^{-31} \mathrm{~kg}} \times 109737 \mathrm{~cm}^{-1}=109677 \mathrm{~cm}^{-1}
$$

and that the ground state of the electron $(n=1)$ lies at

$$
\begin{aligned}
E_{1}=-h c \tilde{R}_{\mathrm{H}}= & -\left(6.62608 \times 10^{-34} \mathrm{Js}\right) \times\left(2.997945 \times 10^{10} \mathrm{~cm} \mathrm{~s}^{-1}\right) \\
& \times\left(109677 \mathrm{~cm}^{-1}\right)=-2.17870 \times 10^{-18} \mathrm{~J}
\end{aligned}
$$

or 2.17870 aJ. This energy corresponds to -13.598 eV .

\subsection*{(c) Ionization energies}
The ionization energy, $I$, of an element is the minimum energy required to remove an electron from the ground state, the state of lowest energy, of one of its atoms in the gas phase. Because the ground state of hydrogen is the state with $n=1$, with energy $E_{1}=-h c \tilde{R}_{\mathrm{H}}$ and the atom is ionized when the electron has been excited to the level corresponding to $n=\infty$ (see Fig. 8A.5), the energy that must be supplied is


\begin{equation*}
I=h c \tilde{R}_{\mathrm{H}} \tag{8A.15}
\end{equation*}


The value of $I$ is $2.179 \mathrm{aJ}\left(1 \mathrm{aJ}=10^{-18} \mathrm{~J}\right)$, which corresponds to 13.60 eV .

A note on good practice Ionization energies are sometimes referred to as ionization potentials. That is incorrect, but not uncommon. If the term is used at all, it should denote the electrical potential difference through which an electron must be moved for the change in its potential energy to be equal to the ionization energy, and reported in volts: the ionization energy of hydrogen is 13.60 eV ; its ionization potential is 13.60 V .

\subsection*{Example 8A. 1 Measuring an ionization energy spectroscopically}
The emission spectrum of atomic hydrogen shows lines at $82259,97492,102824,105292,106632$, and $107440 \mathrm{~cm}^{-1}$, which correspond to transitions to the same lower state from successive upper states with $n=2,3, \ldots$. Determine the ionization energy of the lower state.

Collect your thoughts The spectroscopic determination of ionization energies depends on the identification of the 'series limit', the wavenumber at which the series terminates and becomes a continuum. If the upper state lies at an energy $-h c \tilde{R}_{\mathrm{H}} / n^{2}$, then the wavenumber of the photon emitted when the atom makes a transition to the lower state, with energy $E_{\text {lower }}$, is

$$
\tilde{v}=-\frac{\tilde{R}_{\mathrm{H}}}{n^{2}}-\frac{E_{\text {lower }}}{h c}=-\frac{\tilde{R}_{\mathrm{H}}+\frac{I}{l=-E_{\text {lower }}}}{n^{2}}+\frac{1}{h c}
$$

A plot of the wavenumbers against $1 / n^{2}$ should give a straight line of slope $-\tilde{R}_{\mathrm{H}}$ and intercept $I / h c$. Use software to calculate a least-squares fit of the data in order to obtain a result that reflects the precision of the data.

The solution The wavenumbers are plotted against $1 / n^{2}$ in Fig. 8A.6. From the (least-squares) intercept, it follows that $I / h c=109679 \mathrm{~cm}^{-1}$, so the ionization energy is

$$
\begin{aligned}
I & =h c \times\left(109679 \mathrm{~cm}^{-1}\right) \\
& =\left(6.62608 \times 10^{-34} \mathrm{Js}\right) \times\left(2.997945 \times 10^{10} \mathrm{~cm} \mathrm{~s}^{-1}\right) \times\left(109679 \mathrm{~cm}^{-1}\right) \\
& =2.1787 \times 10^{-18} \mathrm{~J}
\end{aligned}
$$

\begin{center}
\includegraphics[max width=\textwidth]{2025_02_27_d4b95d9718f0b9e2e6b9g-341}
\end{center}

Figure 8A. 6 The plot of the data in Example 8A. 1 used to determine the ionization energy of an atom (in this case, of H ).\\
or 2.1787 aJ , corresponding to $1312.1 \mathrm{~kJ} \mathrm{~mol}^{-1}$ (the negative of the value of $E$ calculated in Brief illustration 8A.2).

Self-test 8 A. 1 The emission spectrum of atomic deuterium shows lines at $15238,20571,23039$, and $24380 \mathrm{~cm}^{-1}$, which correspond to transitions from successive upper states with $n=3,4, \ldots$ to the same lower state. Determine (a) the ionization energy of the lower state, (b) the ionization energy of the ground state, (c) the mass of the deuteron (by expressing the Rydberg constant in terms of the reduced mass of the electron and the deuteron, and solving for the mass of the deuteron).

\subsection*{(d) Shells and subshells}
All the orbitals of a given value of $n$ are said to form a single shell of the atom. In a hydrogenic atom (and only in a hydrogenic atom), all orbitals of given $n$, and therefore belonging to the same shell, have the same energy. It is common to refer to successive shells by letters:

$$
n=\begin{array}{llll}
1 & 2 & 3 & 4 \ldots \\
\mathrm{~K} & \mathrm{~L} & \mathrm{M} & \mathrm{~N} \ldots
\end{array}
$$

Specification of shells

Thus, all the orbitals of the shell with $n=2$ form the L shell of the atom, and so on.

The orbitals with the same value of $n$ but different values of $l$ are said to form a subshell of a given shell. These subshells are also generally referred to by letters:

$$
l=\begin{array}{llllllll}
0 & 1 & 2 & 3 & 4 & 5 & 6 \ldots & \\
\mathrm{~s} & \mathrm{p} & \mathrm{~d} & \mathrm{f} & \mathrm{~g} & \mathrm{~h} & \mathrm{i} \ldots & \text { Specification of subshells }
\end{array}
$$

All orbitals of the same subshell have the same energy in all kinds of atoms, not only hydrogenic atoms. After $l=3$ the letters run alphabetically ( j is not used because in some lan-\\
\includegraphics[max width=\textwidth, center]{2025_02_27_d4b95d9718f0b9e2e6b9g-342(2)}

Figure 8A. 7 The energy levels of a hydrogenic atom showing the subshells and (in square brackets) the numbers of orbitals in each subshell. All orbitals of a given shell have the same energy.\\
\includegraphics[max width=\textwidth, center]{2025_02_27_d4b95d9718f0b9e2e6b9g-342(1)}

Figure 8A. 8 The organization of orbitals (white squares) into subshells (characterized by $l$ ) and shells (characterized by $n$ ).\\
guages i and j are not distinguished). Figure 8 A .7 is a version of Fig. 8A. 5 which shows the subshells explicitly. Because $l$ can range from 0 to $n-1$, giving $n$ values in all, it follows that there are $n$ subshells of a shell with principal quantum number $n$. The organization of orbitals in the shells is summarized in Fig. 8A.8. The number of orbitals in a shell of principal quantum number $n$ is $n^{2}$, so in a hydrogenic atom each energy level is $n^{2}$-fold degenerate.

\subsection*{Brief illustration 8A. 3}
When $n=1$ there is only one subshell, that with $l=0$, and that subshell contains only one orbital, with $m_{l}=0$ (the only value of $m_{l}$ permitted). When $n=2$, there are four orbitals, one in the s subshell with $l=0$ and $m_{l}=0$, and three in the $l=1$ subshell with $m_{l}=+1,0,-1$. When $n=3$ there are nine orbitals (one with $l=0$, three with $l=1$, and five with $l=2$ ).

\subsection*{(e) $\mathbf{s}$ Orbitals}
The orbital occupied in the ground state is the one with $n=1$ (and therefore with $l=0$ and $m_{l}=0$, the only possible values of these quantum numbers when $n=1$ ). From Table 8A. 1 and with $Y_{0,0}=(1 / 4 \pi)^{1 / 2}$ (Table 7F.1) it follows that (for $Z=1$ ):


\begin{equation*}
\psi=\frac{1}{\left(\pi a_{0}^{3}\right)^{1 / 2}} \mathrm{e}^{-r / a_{0}} \tag{8A.16}
\end{equation*}


This wavefunction is independent of angle and has the same value at all points of constant radius; that is, the 1 s orbital (the $s$ orbital with $n=1$, and in general $n s$ ) is 'spherically symmetrical'. The wavefunction decays exponentially from a maximum value of $1 /\left(\pi a_{0}^{3}\right)^{1 / 2}$ at the nucleus (at $r=0$ ). It follows that the probability density of the electron is greatest at the nucleus itself.

The general form of the ground-state wavefunction can be understood by considering the contributions of the potential and kinetic energies to the total energy of the atom. The closer the electron is to the nucleus on average, the lower (more negative) its average potential energy. This dependence suggests that the lowest potential energy should be obtained with a sharply peaked wavefunction that has a large amplitude at the nucleus and is zero everywhere else (Fig. 8A.9). However, this shape implies a high kinetic energy, because such a wavefunction has a very high average curvature. The electron would have very low kinetic energy if its wavefunction had only a very low average curvature. However, such a wavefunction spreads to great distances from the nucleus and the average potential energy of the electron is correspondingly high. The actual ground-state wavefunction is a compromise between these two extremes: the wavefunction spreads away\\
\includegraphics[max width=\textwidth, center]{2025_02_27_d4b95d9718f0b9e2e6b9g-342}

Figure 8A. 9 The balance of kinetic and potential energies that accounts for the structure of the ground state of hydrogenic atoms. (a) The sharply curved but localized orbital has high mean kinetic energy, but low mean potential energy; (b) the mean kinetic energy is low, but the potential energy is not very favourable; (c) the compromise of moderate kinetic energy and moderately favourable potential energy.\\
\includegraphics[max width=\textwidth, center]{2025_02_27_d4b95d9718f0b9e2e6b9g-343(1)}

Figure 8A. 10 Representations of cross-sections through the (a) 1 s and (b) 2 s hydrogenic atomic orbitals in terms of their electron probability densities (as represented by the density of shading).\\
\includegraphics[max width=\textwidth, center]{2025_02_27_d4b95d9718f0b9e2e6b9g-343}

Figure 8A. 11 The boundary surface of a 1 s orbital, within which there is a 90 per cent probability of finding the electron. All s orbitals have spherical boundary surfaces.\\
from the nucleus (so the expectation value of the potential energy is not as low as in the first example, but nor is it very high) and has a reasonably low average curvature (so the expectation of the kinetic energy is not very low, but nor is it as high as in the first example).

One way of depicting the probability density of the electron is to represent $|\psi|^{2}$ by the density of shading (Fig. 8A.10). A simpler procedure is to show only the boundary surface, the surface that mirrors the shape of the orbital and captures a high proportion (typically about 90 per cent) of the electron probability. For the 1 s orbital, the boundary surface is a sphere centred on the nucleus (Fig. 8A.11).

\subsection*{Example 8A. 2}
\subsection*{Calculating the mean radius of an orbital}
Calculate the mean radius of a hydrogenic 1 s orbital.\\
Collect your thoughts The mean radius is the expectation value

$$
\langle r\rangle=\int \psi^{\star} r \psi \mathrm{~d} \tau=\int r|\psi|^{2} \mathrm{~d} \tau
$$

You need to evaluate the integral by using the wavefunctions given in Table 8A. 1 and $\mathrm{d} \tau=r^{2} \mathrm{~d} r \sin \theta \mathrm{~d} \theta \mathrm{~d} \phi$ (The chemist's toolkit 21 in Topic 7F). The angular parts of the wavefunction (Table 7F.1) are normalized in the sense that

$$
\int_{\theta=0}^{\pi} \int_{\phi=0}^{2 \pi}\left|Y_{l, m_{l}}\right|^{2} \sin \theta \mathrm{~d} \theta \mathrm{~d} \phi=1
$$

The relevant integral over $r$ is given in the Resource section.\\
The solution With the wavefunction written in the form $\psi=$ $R Y$, the integration (with the integral over the angular variables, which is equal to 1 , in blue) is

$$
\langle r\rangle=\int_{0}^{\infty} \int_{0}^{\pi} \int_{0}^{2 \pi} r R_{n, l}^{2}\left|Y_{l, m_{l}}\right|^{2} r^{2} \mathrm{~d} r \sin \theta \mathrm{~d} \theta \mathrm{~d} \phi=\int_{0}^{\infty} r^{3} R_{n, l}^{2} \mathrm{~d} r
$$

For a 1s orbital

$$
R_{1,0}=2\left(\frac{Z}{a_{0}}\right)^{3 / 2} \mathrm{e}^{-Z r / a_{0}}
$$

Hence

$$
\langle r\rangle=\frac{4 Z^{3}}{a_{0}^{3}} \overbrace{0}^{\text {Integral E. } 3} r^{3} \mathrm{e}^{-2 Z r / a_{0}} \mathrm{~d} r=\frac{4 Z^{3}}{a_{0}^{3}} \times \frac{3!}{\left(2 Z / a_{0}\right)^{4}}=\frac{3 a_{0}}{2 Z}
$$

Self-test 8A.2 Evaluate the mean radius of a 3 s orbital by integration.

$$
Z Z /{ }^{0} v_{L Z}: \iota \partial м s u_{V}
$$

All $s$ orbitals are spherically symmetric, but differ in the number of radial nodes. For example, the $1 \mathrm{~s}, 2 \mathrm{~s}$, and 3 s orbitals have 0,1 , and 2 radial nodes, respectively. In general, an $n s$ orbital has $n-1$ radial nodes. As $n$ increases, the radius of the spherical boundary surface that captures a given fraction of the probability also increases.

\subsection*{Brief illustration 8A. 4}
The radial nodes of a $2 s$ orbital lie at the locations where the associated Laguerre polynomial factor (Table 8A.1) is equal to zero. In this case the factor is simply $2-\rho$ so there is a node at $\rho=2$. For a 2 s orbital, $\rho=Z r / a_{0}$, so the radial node occurs at $r$ $=2 a_{0} / Z$ (see Fig. 8A.4).

\subsection*{(f) Radial distribution functions}
The wavefunction yields, through the value of $|\psi|^{2}$, the probability of finding an electron in any region. As explained in Topic $7 \mathrm{~B},|\psi|^{2}$ is a probability density (dimensions: $1 /$ volume) and can be interpreted as a (dimensionless) probability when multiplied by the (infinitesimal) volume of interest. Imagine a probe with a fixed volume $\mathrm{d} \tau$ and sensitive to electrons that\\
\includegraphics[max width=\textwidth, center]{2025_02_27_d4b95d9718f0b9e2e6b9g-344}

Figure 8A. 12 A constant-volume electron-sensitive detector (the small cube) gives its greatest reading at the nucleus, and a smaller reading elsewhere. The same reading is obtained anywhere on a circle of given radius at any orientation: the s orbital is spherically symmetrical.\\
can move around near the nucleus of a hydrogenic atom. Because the probability density in the ground state of the atom is proportional to $\mathrm{e}^{-2 Z r / a_{0}}$, the reading from the detector decreases exponentially as the probe is moved out along any radius but is constant if the probe is moved on a circle of constant radius (Fig. 8A.12).

Now consider the total probability of finding the electron anywhere between the two walls of a spherical shell of thickness $\mathrm{d} r$ at a radius $r$. The sensitive volume of the probe is now the volume of the shell (Fig. 8A.13), which is $4 \pi r^{2} \mathrm{~d} r$ (the product of its surface area, $4 \pi r^{2}$, and its thickness, $\mathrm{d} r$ ). Note that the volume probed increases with distance from the nucleus and is zero at the nucleus itself, when $r=0$. The probability that the\\
\includegraphics[max width=\textwidth, center]{2025_02_27_d4b95d9718f0b9e2e6b9g-344(1)}

Figure 8A. 13 The radial distribution function $P(r)$ is the probability density that the electron will be found anywhere in a shell of radius $r$; the probability itself is $P(r) \mathrm{d} r$, where $\mathrm{d} r$ is the thickness of the shell. For a 1 s electron in hydrogen, $P(r)$ is a maximum when $r$ is equal to the Bohr radius $a_{0}$. The value of $P(r) \mathrm{d} r$ is equivalent to the reading that a detector shaped like a spherical shell of thickness $\mathrm{d} r$ would give as its radius is varied.\\
electron will be found between the inner and outer surfaces of this shell is the probability density at the radius $r$ multiplied by the volume of the probe, or $|\psi(r)|^{2} \times 4 \pi r^{2} \mathrm{~d} r$. This expression has the form $P(r) \mathrm{d} r$, where

\[
P(r)=4 \pi r^{2}|\psi(r)|^{2} \quad \begin{align*}
& \text { Radial distribution function }  \tag{8А.17a}\\
& {[\text { s orbitals only] }}
\end{align*}
\]

The function $P(r)$ is called the radial distribution function (in this case, for an $s$ orbital). It is also possible to devise a more general expression which applies to orbitals that are not spherically symmetrical.

\subsection*{How is that done? 8A. 1 Deriving the general form of the radial distribution function}
The probability of finding an electron in a volume element $\mathrm{d} \tau$ when its wavefunction is $\psi=R Y$ is $|R Y|^{2} \mathrm{~d} \tau$ with $\mathrm{d} \tau=$ $r^{2} \mathrm{~d} r \sin \theta \mathrm{~d} \theta \mathrm{~d} \phi$. The total probability of finding the electron at any angle in a shell of radius $r$ and thickness $\mathrm{d} r$ is the integral of this probability over the entire surface, and is written $P(r) \mathrm{d} r$; so

$$
P(r) \mathrm{d} r=\int_{0}^{\pi} \int_{0}^{2 \pi} R(r)^{2}\left|Y_{l, m_{l}}\right|^{2} r^{2} \mathrm{~d} r \sin \theta \mathrm{~d} \theta \mathrm{~d} \phi
$$

Because the spherical harmonics are normalized to 1 (the blue integration, as in Example 8A.2, gives 1), the final result is

$$
-\left\lvert\, P(r)=r^{2} R(r)^{2} \longmapsto \begin{aligned}
& \text { Radial distribution function } \\
& \text { [general form] }
\end{aligned}\right.
$$

The radial distribution function is a probability density in the sense that, when it is multiplied by $\mathrm{d} r$, it gives the probability of finding the electron anywhere between the two walls of a spherical shell of thickness $\mathrm{d} r$ at the radius $r$. For a 1 s orbital,


\begin{equation*}
P(r)=\frac{4 Z^{3}}{a_{0}^{3}} r^{2} \mathrm{e}^{-2 Z r a_{0}} \tag{8A.18}
\end{equation*}


This expression can be interpreted as follows:

\begin{itemize}
  \item Because $r^{2}=0$ at the nucleus, $P(0)=0$. The volume of the shell is zero when $r=0$ so the probability of finding the electron in the shell is zero.
  \item As $r \rightarrow \infty, P(r) \rightarrow 0$ on account of the exponential term. The wavefunction has fallen to zero at great distances from the nucleus and there is little probability of finding the electron even in a large shell.
  \item The increase in $r^{2}$ and the decrease in the exponential factor means that $P$ passes through a maximum at an intermediate radius (see Fig. 8A.13); it marks the most probable radius at which the electron will be found regardless of direction.
\end{itemize}

\subsection*{Example 8A. 3 Calculating the most probable radius}
Calculate the most probable radius, $r_{\mathrm{mp}}$, at which an electron will be found when it occupies a 1 s orbital of a hydrogenic atom of atomic number $Z$, and tabulate the values for the oneelectron species from H to $\mathrm{Ne}^{9+}$.

Collect your thoughts You need to find the radius at which the radial distribution function of the hydrogenic 1s orbital has a maximum value by solving $\mathrm{d} P / \mathrm{d} r=0$. If there are several maxima, you should choose the one corresponding to the greatest amplitude.

The solution The radial distribution function is given in eqn 8A.18. It follows that

$$
\frac{\mathrm{d} P}{\mathrm{~d} r}=\frac{4 Z^{3}}{a_{0}^{3}}\left(2 r-\frac{2 Z r^{2}}{a_{0}}\right) \mathrm{e}^{-2 Z r / a_{0}}=\frac{8 r Z^{3}}{a_{0}^{3}}\left(1-\frac{Z r}{a_{0}}\right) \mathrm{e}^{-2 Z r / a_{0}}
$$

This function is zero other than at $r=0$ where the term in parentheses is zero, which is at

$$
r_{\mathrm{mp}}=\frac{a_{0}}{Z}
$$

Then, with $a_{0}=52.9 \mathrm{pm}$, the most probable radii are

\begin{center}
\begin{tabular}{lllllllllll}
\hline
 & H & $\mathrm{He}^{+}$ & $\mathrm{Li}^{2+}$ & $\mathrm{Be}^{3+}$ & $\mathrm{B}^{4+}$ & $\mathrm{C}^{5+}$ & $\mathbf{N}^{6+}$ & $\mathbf{O}^{7+}$ & $\mathrm{F}^{8+}$ & $\mathrm{Ne}^{9+}$ \\
\hline
$r_{\mathrm{mp}} / \mathrm{pm}$ & 52.9 & 26.5 & 17.6 & 13.2 & 10.6 & 8.82 & 7.56 & 6.61 & 5.88 & 5.29 \\
\hline
\end{tabular}
\end{center}

Comment. Notice how the 1 s orbital is drawn towards the nucleus as the nuclear charge increases. At uranium the most probable radius is only 0.58 pm , almost 100 times closer than for hydrogen. (On a scale where $r_{\mathrm{mp}}=10 \mathrm{~cm}$ for $\mathrm{H}, r_{\mathrm{mp}}=1 \mathrm{~mm}$ for U.) However, extending this result to very heavy atoms neglects important relativistic effects that complicate the calculation.

Self-test 8A. 3 Find the most probable distance of a 2 s electron from the nucleus in a hydrogenic atom.\\
\includegraphics[max width=\textwidth, center]{2025_02_27_d4b95d9718f0b9e2e6b9g-345}\\
\includegraphics[max width=\textwidth, center]{2025_02_27_d4b95d9718f0b9e2e6b9g-345(1)}

\subsection*{(g) p Orbitals}
All three 2 p orbitals have $l=1$, and therefore the same magnitude of angular momentum; they are distinguished by different values of $m_{l}$, the quantum number that specifies the component of angular momentum around a chosen axis (conventionally taken to be the $z$-axis). The orbital with $m_{l}=0$, for instance, has zero angular momentum around the $z$-axis. Its angular variation is given by the spherical harmonic $Y_{1,0}$, which is proportional to $\cos \theta$ (see Table 7F.1). Therefore, the probability density, which is proportional to $\cos ^{2} \theta$, has its maximum value on either side of the nucleus along the $z$-axis\\
(at $\theta=0$ and $180^{\circ}$, where $\cos ^{2} \theta=1$ ). Specifically, the wavefunction of a 2 p orbital with $m_{l}=0$ is


\begin{align*}
\psi_{2,1,0} & =R_{2,1}(r) Y_{1,0}(\theta, \phi)=\frac{1}{4(2 \pi)^{1 / 2}}\left(\frac{Z}{a_{0}}\right)^{5 / 2} r \cos \theta \mathrm{e}^{-Z r / 2 a_{0}}  \tag{8A.19a}\\
& =r \cos \theta f(r)
\end{align*}


where $f(r)$ is a function only of $r$. Because in spherical polar coordinates $z=r \cos \theta$ (The chemist's toolkit 21 in Topic 7F), this wavefunction may also be written


\begin{equation*}
\psi_{2,1,0}=z f(r) \tag{8A.19b}
\end{equation*}


All p orbitals with $m_{l}=0$ and any value of $n$ have wavefunctions of this form, but $f(r)$ depends on the value of $n$. This way of writing the orbital is the origin of the name ' $p_{z}$ orbital': its boundary surface is shown in Fig. 8A.14. The wavefunction is zero everywhere in the $x y$-plane, where $z=0$, so the $x y$-plane is a nodal plane of the orbital: the wavefunction changes sign on going from one side of the plane to the other.

The wavefunctions of 2 p orbitals with $m_{l}= \pm 1$ have the following form:


\begin{align*}
\psi_{2,1, \pm 1}=R_{2,1}(r) Y_{1, \pm 1}(\theta, \phi) & =\mp \frac{1}{8 \pi^{1 / 2}}\left(\frac{Z}{a_{0}}\right)^{5 / 2} r \sin \theta \mathrm{e}^{ \pm i \phi} \mathrm{e}^{-Z r / 2 a_{0}} \\
& =\mp \frac{1}{2^{1 / 2}} r \sin \theta \mathrm{e}^{ \pm i \phi} f(r) \tag{8A.20}
\end{align*}


In Topic 7D it is explained that a particle described by a complex wavefunction has net motion. In the present case, the\\
\includegraphics[max width=\textwidth, center]{2025_02_27_d4b95d9718f0b9e2e6b9g-345(2)}

Figure 8A. 14 The boundary surfaces of $2 p$ orbitals. A nodal plane passes through the nucleus and separates the two lobes of each orbital. The dark and light lobes denote regions of opposite sign of the wavefunction. The angles of the spherical polar coordinate system are also shown. All p orbitals have boundary surfaces like those shown here.\\
functions correspond to non-zero angular momentum about the $z$-axis: $\mathrm{e}^{+\mathrm{it}}$ corresponds to clockwise rotation when viewed from below, and $\mathrm{e}^{-\mathrm{i} \phi}$ corresponds to anticlockwise rotation (from the same viewpoint). They have zero amplitude where $\theta=0$ and $180^{\circ}$ (along the $z$-axis) and maximum amplitude at $90^{\circ}$, which is in the $x y$-plane. To draw the functions it is usual to represent them by forming the linear combinations


\begin{align*}
& \psi_{2 \mathrm{p}_{x}}=\frac{1}{2^{1 / 2}}\left(\psi_{2,1,+1}-\psi_{2,1,-1}\right) \stackrel{\downarrow}{=} r \sin \theta \cos \phi f(r)=x f(r) \\
& \psi_{2 \mathrm{p}_{x}}=\frac{\mathrm{i}}{2^{1 / 2}+\mathrm{e}^{-\mathrm{i} \phi}=2 \cos \phi} \\
& \underbrace{}_{2,1,+1}+\psi_{2,1,-1})=r \sin \theta \sin \phi f(r)=y f(r) \tag{8A.21}
\end{align*}


These linear combinations correspond to zero orbital angular momentum around the $z$-axis, as they are superpositions of states with equal and opposite values of $m_{l}$. The $\mathrm{p}_{x}$ orbital has the same shape as a $p_{z}$ orbital, but it is directed along the $x$-axis (see Fig. 8A.14); the $p_{y}$ orbital is similarly directed along the $y$-axis. The wavefunction of any p orbital of a given shell can be written as a product of $x, y$, or $z$ and the same function $f$ (which depends on the value of $n$ ).

\subsection*{(h) d Orbitals}
When $n=3, l$ can be 0,1 , or 2 . As a result, this shell consists of one 3 s orbital, three 3 p orbitals, and five 3d orbitals. Each value of the quantum number $m_{l}=0, \pm 1, \pm 2$ corresponds to a different value of the component of angular momentum about the $z$-axis. As for the p orbitals, d orbitals with opposite values of $m_{l}$ (and hence opposite senses of motion around the $z$-axis) may be combined in pairs to give real wavefunctions, and the boundary surfaces of the resulting shapes are shown\\
\includegraphics[max width=\textwidth, center]{2025_02_27_d4b95d9718f0b9e2e6b9g-346}

Figure 8A. 15 The boundary surfaces of 3d orbitals. The purple and yellow areas denote regions of opposite sign of the wavefunction. All d orbitals have boundary surfaces like those shown here.\\
in Fig. 8A.15. The real linear combinations have the following forms, with the function $f(r)$ depending on the value of $n$ :


\begin{align*}
& \psi_{\mathrm{d}_{x y}}=x y f(r) \quad \psi_{\mathrm{d}_{y z}}=y z f(r) \quad \psi_{\mathrm{d}_{z x}}=z x f(r) \\
& \psi_{\mathrm{d}_{x^{2}-y^{2}}}=\frac{1}{2}\left(x^{2}-y^{2}\right) f(r) \quad \psi_{\mathrm{d}_{z^{2}}}=\frac{3^{1 / 2}}{2}\left(3 z^{2}-r^{2}\right) f(r) \tag{8А.22}
\end{align*}


These linear combinations give rise to the notation $\mathrm{d}_{x y}, \mathrm{~d}_{y z}$, etc. for the d-orbitals. With the exception of the $\mathrm{d}_{z^{2}}$ orbital, each combination has two angular nodes which divide the orbital into four lobes. For the $\mathrm{d}_{z^{2}}$ orbital, the two angular nodes combine to give a conical surface that separates the main lobes from a smaller toroidal component encircling the nucleus.

\subsection*{Checklist of concepts}
1. The Schrödinger equation for a hydrogenic atom separates into angular and radial equations.2. Close to the nucleus the radial wavefunction is proportional to $r^{l}$; far from the nucleus all hydrogenic wavefunctions approach zero exponentially.3. An atomic orbital is a one-electron wavefunction for an electron in an atom.4. An atomic orbital is specified by the values of the quantum numbers $n, l$, and $m_{l}$.5. The energies of the bound states of hydrogenic atoms are proportional to $-Z^{2} / n^{2}$.6. The ionization energy of an element is the minimum energy required to remove an electron from the ground state of one of its atoms.7. Orbitals of a given value of $n$ form a shell of an atom, and within that shell orbitals of the same value of $l$ form subshells.8. Orbitals of the same shell all have the same energy in hydrogenic atoms; orbitals of the same subshell of a shell are degenerate in all types of atoms.9. s Orbitals are spherically symmetrical and have nonzero probability density at the nucleus.10. A radial distribution function is the probability density for the distribution of the electron as a function of distance from the nucleus.11. There are three $\mathbf{p}$ orbitals in a given subshell; each one has one angular node.12. There are five d orbitals in a given subshell; each one has two angular nodes.\subsection*{Checklist of equations}
\begin{center}
\begin{tabular}{|c|c|c|c|}
\hline
Property & Equation & Comment & Equation number \\
\hline
Wavenumbers of the spectral lines of a hydrogen atom & $\tilde{v}=\tilde{R}_{\mathrm{H}}\left(1 / n_{1}^{2}-1 / n_{2}^{2}\right)$ & $\tilde{R}_{\mathrm{H}}$ is the Rydberg constant for hydrogen (expressed as a wavenumber) & 8A. 1 \\
\hline
Bohr radius & $a_{0}=4 \pi \varepsilon_{0} \hbar^{2} / m_{e} e^{2}$ & $a_{0}=52.9 \mathrm{pm}$ & 8A. 9 \\
\hline
Wavefunctions of hydrogenic atoms & $\psi_{n, l, m_{l}}(r, \theta, \phi)=R_{n, l}(r) Y_{l, m_{l}}(\theta, \phi)$ & $Y_{l, m_{l}}$ are spherical harmonics & 8A. 12 \\
\hline
Energies of hydrogenic atoms & \( \begin{aligned} & E_{n}=-h c Z^{2} \tilde{R}_{\mathrm{N}} / n^{2}, \\ & \tilde{R}_{\mathrm{N}}=\mu e^{4} / 8 \varepsilon_{0}^{2} c h^{3} \end{aligned} \) & $\tilde{R}_{\mathrm{N}} \approx \tilde{R}_{\infty}$, the Rydberg constant; \( \mu=m_{\mathrm{e}} m_{\mathrm{N}} /\left(m_{\mathrm{e}}+m_{\mathrm{N}}\right) \) & 8A. 13 \\
\hline
Radial distribution function & $P(r)=r^{2} R(r)^{2}$ & $P(r)=4 \pi r^{2} \psi^{2}$ for s orbitals & 8A.17b \\
\hline
\end{tabular}
\end{center}

\section*{TOPIC 8B Many-electron atoms}
\subsection*{Why do you need to know this material?}
Many-electron atoms are the building blocks of all compounds, and to understand their properties, including their ability to participate in chemical bonding, it is essential to understand their electronic structure. Moreover, a knowledge of that structure explains the structure of the periodic table and all that it summarizes.

\subsection*{What is the key idea?}
Electrons occupy the orbitals that result in the lowest energy of the atom, subject to the requirements of the Pauli exclusion principle.

\subsection*{- What do you need to know already?}
This Topic builds on the account of the structure of hydrogenic atoms (Topic 8A), especially their shell structure.

A many-electron atom (or polyelectron atom) is an atom with more than one electron. The Schrödinger equation for a manyelectron atom is complicated because all the electrons interact with one another. One very important consequence of these interactions is that orbitals of the same value of $n$ but different values of $l$ are no longer degenerate. Moreover, even for a helium atom, with just two electrons, it is not possible to find analytical expressions for the orbitals and energies, so it is necessary to use various approximations.

\subsection*{8B. 1 The orbital approximation}
The wavefunction of a many-electron atom is a very complicated function of the coordinates of all the electrons, written as $\Psi\left(\boldsymbol{r}_{1}, \boldsymbol{r}_{2}, \ldots\right)$, where $\boldsymbol{r}_{i}$ is the vector from the nucleus to electron $i$ (uppercase psi, $\Psi$, is commonly used to denote a manyelectron wavefunction). The orbital approximation states that a reasonable first approximation to this exact wavefunction is obtained by thinking of each electron as occupying its 'own' orbital, and writing


\begin{equation*}
\Psi\left(\boldsymbol{r}_{1}, \boldsymbol{r}_{2}, \ldots\right)=\psi\left(\boldsymbol{r}_{1}\right) \psi\left(\boldsymbol{r}_{2}\right) \ldots \quad \text { Orbital approximation } \tag{8B.1}
\end{equation*}


The individual orbitals can be assumed to resemble the hydrogenic orbitals based on nuclei with charges modified by the presence of all the other electrons in the atom. This assumption can be justified if, to a first approximation, electronelectron interactions are ignored.

\subsection*{How is that done? 8B. 1}
Justifying the orbital\\
approximation\\
Consider a system in which the hamiltonian for the energy is the sum of two contributions, one for electron 1 and the other for electron 2: $\hat{H}=\hat{H}_{1}+\hat{H}_{2}$. In an actual two-electron atom (such as a helium atom), there is an additional term (proportional to $1 / r_{12}$, where $r_{12}$ is the distance between the two electrons) corresponding to their interaction:

$$
\hat{H}=\overbrace{-\frac{\hbar^{2}}{2 m_{\mathrm{e}}} \nabla_{1}^{2}-\frac{2 e^{2}}{4 \pi \varepsilon_{0} r_{1}}}^{\hat{H}_{1}} \overbrace{-\frac{\hbar^{2}}{2 m_{\mathrm{e}}} \nabla_{2}^{2}-\frac{2 e^{2}}{4 \pi \varepsilon_{0} r_{2}}}^{\hat{H}_{2}}+\frac{e^{2}}{4 \pi \varepsilon_{0} r_{12}}
$$

In the orbital approximation the final term is ignored. Then the task is to show that if $\psi\left(r_{1}\right)$ is an eigenfunction of $\hat{H}_{1}$ with energy $E_{1}$, and $\psi\left(r_{2}\right)$ is an eigenfunction of $\hat{H}_{2}$ with energy $E_{2}$, then the product $\Psi\left(\boldsymbol{r}_{1}, \boldsymbol{r}_{2}\right)=\psi\left(\boldsymbol{r}_{1}\right) \psi\left(\boldsymbol{r}_{2}\right)$ is an eigenfunction of the combined hamiltonian $\hat{H}$. To do so write

$$
\begin{aligned}
\hat{H} \Psi\left(\boldsymbol{r}_{1}, \boldsymbol{r}_{2}\right) & =\left(\hat{H}_{1}+\hat{H}_{2}\right) \psi\left(\boldsymbol{r}_{1}\right) \psi\left(\boldsymbol{r}_{2}\right) \\
& =\overbrace{\hat{H}_{1} \psi\left(\boldsymbol{r}_{1}\right) \psi\left(\boldsymbol{r}_{2}\right)}^{\psi\left(\boldsymbol{r}_{2}\right) \hat{H}_{1} \psi\left(\boldsymbol{r}_{1}\right)}+\overbrace{\hat{H}_{2} \psi\left(\boldsymbol{r}_{1}\right) \psi\left(\boldsymbol{r}_{2}\right)}^{\psi\left(\boldsymbol{r}_{1}\right) \hat{H}_{2} \psi\left(\boldsymbol{r}_{2}\right)} \\
& =\psi\left(\boldsymbol{r}_{2}\right) \overbrace{\hat{H}_{1} \psi\left(\boldsymbol{r}_{1}\right)}^{E_{1} \psi\left(\boldsymbol{r}_{1}\right)}+\psi\left(\boldsymbol{r}_{1}\right) \overbrace{\hat{H}_{2} \psi\left(\boldsymbol{r}_{2}\right)}^{E_{2} \psi\left(\boldsymbol{r}_{2}\right)} \\
& =\psi\left(\boldsymbol{r}_{2}\right) E_{1} \psi\left(\boldsymbol{r}_{1}\right)+\psi\left(\boldsymbol{r}_{1}\right) E_{2} \psi\left(\boldsymbol{r}_{2}\right)=\left(E_{1}+E_{2}\right) \psi\left(\boldsymbol{r}_{1}\right) \psi\left(\boldsymbol{r}_{2}\right) \\
& =E \Psi\left(\boldsymbol{r}_{1}, \boldsymbol{r}_{2}\right)
\end{aligned}
$$

where $E=E_{1}+E_{2}$, which is the desired result. Note how each hamiltonian operates on only its 'own' wavefunction. If the electrons interact (as they do in fact), then the term in $1 / r_{12}$ must be included, and the proof fails. Therefore, this description is only approximate, but it is a useful model for discussing the chemical properties of atoms and is the starting point for more sophisticated descriptions of atomic structure.

The orbital approximation can be used to express the electronic structure of an atom by reporting its configuration, a statement of its occupied orbitals (usually, but not necessarily, in its ground state). Thus, as the ground state of a hydrogenic atom consists of the single electron in a 1 s orbital, its configuration is reported as $1 \mathrm{~s}^{1}$ (read 'one-ess-one').

A He atom has two electrons. The first electron occupies a 1s hydrogenic orbital, but because $Z=2$ that orbital is more compact than in H itself. The second electron joins the first in the 1 s orbital, so the electron configuration of the ground state of He is $1 \mathrm{~s}^{2}$.

\subsection*{Brief illustration 8B.1}
According to the orbital approximation, each electron in He occupies a hydrogenic $1 s$ orbital of the kind given in Topic 8A. Anticipating (see below) that the electrons experience an effective nuclear charge $Z_{\text {eff }} e$ rather than the actual charge on the nucleus with $Z=2$ (specifically, as seen later, a charge $1.69 e$ rather than $2 e$ ), then the two-electron wavefunction of the atom is

$$
\begin{aligned}
\Psi\left(\boldsymbol{r}_{1}, \boldsymbol{r}_{2}\right) & =\overbrace{\frac{Z_{\text {eff }}^{3 / 2}}{\left(\pi a_{0}^{3}\right)^{1 / 2}} \mathrm{e}^{-Z_{\text {eff }} / a_{0}}}^{\psi_{1 s}\left(\boldsymbol{r}_{1}\right)} \times \overbrace{\frac{Z_{\text {eff }}^{3 / 2}}{\left(\pi a_{0}^{3}\right)^{1 / 2}} \mathrm{e}^{-Z_{\text {eff }} / a_{0}}}^{\psi_{1 s}\left(\boldsymbol{r}_{2}\right)} \\
& =\frac{Z_{\text {eff }}^{3}}{\pi a_{0}^{3}} \mathrm{e}^{-Z_{\text {eff }}\left(r_{1}+r_{2}\right) / a_{0}}
\end{aligned}
$$

There is nothing particularly mysterious about a two-electron wavefunction: in this case it is a simple exponential function of the distances of the two electrons from the nucleus.

\subsection*{8B.2 The Pauli exclusion principle}
It is tempting to suppose that the electronic configurations of the atoms of successive elements with atomic numbers $Z=$ $3,4, \ldots$, and therefore with $Z$ electrons, are simply $1 s^{Z}$. That, however, is not the case. The reason lies in two aspects of nature: that electrons possess 'spin' and that they must obey the very fundamental 'Pauli principle'.

\subsection*{(a) Spin}
The quantum mechanical property of electron spin, the possession of an intrinsic angular momentum, was identified by an experiment performed by Otto Stern and Walther Gerlach in 1921, who shot a beam of silver atoms through an inhomogeneous magnetic field (Fig. 8B.1). The idea behind the experiment was that each atom possesses a certain electronic angular momentum and (because moving charges generate a magnetic field) as a result behaves like a small bar magnet aligned with\\
\includegraphics[max width=\textwidth, center]{2025_02_27_d4b95d9718f0b9e2e6b9g-349}

Figure 8B. 1 (a) The experimental arrangement for the SternGerlach experiment: the magnet provides an inhomogeneous field. (b) The classically expected result. (c) The observed outcome using silver atoms.\\
the direction of the angular momentum vector. As the atoms pass through the inhomogeneous magnetic field they are deflected, with the deflection depending on the relative orientation of the applied magnetic field and the atomic magnet.

The classical expectation is that the electronic angular momentum, and hence the resulting magnet, can be oriented in any direction. Each atom would be deflected into a direction that depends on the orientation and the beam should spread out into a broad band as it emerges from the magnetic field. In contrast, the expectation from quantum mechanics is that the angular momentum, and hence the atomic magnet, has only discrete orientations (Topic 7F). Each of these orientations results in the atoms being deflected in a specific direction, so the beam should split into a number of sharp bands, each corresponding to a different orientation of the angular momentum of the electrons in the atom.

In their first experiment, Stern and Gerlach appeared to confirm the classical prediction. However, the experiment is difficult because collisions between the atoms in the beam blur the bands. When they repeated the experiment with a beam of very low intensity (so that collisions were less frequent), they observed discrete bands, and so confirmed the quantum prediction. However, Stern and Gerlach observed two bands of Ag atoms in their experiment. This observation seems to conflict with one of the predictions of quantum mechanics, because an angular momentum $l$ gives rise to $2 l+1$ orientations, which is equal to 2 only if $l=\frac{1}{2}$, contrary to the requirement that $l$ is an integer. The conflict was resolved by the suggestion that the angular momentum they were observing was not due to orbital angular momentum (the motion of an electron around the atomic nucleus) but arose instead from the rotation of the electron about its own axis, its 'spin'.

The spin of an electron does not have to satisfy the same boundary conditions as those for a particle circulating through space around a central point, so the quantum number for spin angular momentum is subject to different restrictions. The spin\\
\includegraphics[max width=\textwidth, center]{2025_02_27_d4b95d9718f0b9e2e6b9g-350(1)}

Figure 8B. 2 The vector representation of the spin of an electron. The length of the side of the cone is $3^{1 / 2} / 2$ units and the projections on to the $z$-axis are $\pm \frac{1}{2}$ units.\\
quantum number $s$ is used in place of the orbital angular momentum quantum number $l$ (Topic 7 F ; like $l, s$ is a non-negative number) and $m_{s}$, the spin magnetic quantum number, is used in place of $m_{l}$ for the projection on the $z$-axis. The magnitude of the spin angular momentum is $\{s(s+1)\}^{1 / 2} \hbar$ and the component $m_{s} \hbar$ is restricted to the $2 s+1$ values $m_{s}=s, s-1, \ldots,-s$. To account for Stern and Gerlach's observation, $s=\frac{1}{2}$ and $m_{s}= \pm \frac{1}{2}$.

A note on good practice You will sometimes see the quantum number $s$ used in place of $m_{s}$, and written $s= \pm \frac{1}{2}$. That is wrong: like $l, s$ is never negative and denotes the magnitude of the spin angular momentum. For the $z$-component, use $m_{s}$.

The detailed analysis of the spin of a particle is sophisticated and shows that the property should not be taken to be an actual spinning motion. It is better to regard 'spin' as an intrinsic property like mass and charge: every electron has exactly the same value and the magnitude of the spin angular momentum of an electron cannot be changed. However, the picture of an actual spinning motion can be very useful when used with care. In the vector model of angular momentum (Topic 7F), the spin may lie in two different orientations (Fig. 8B.2). One orientation corresponds to $m_{s}=+\frac{1}{2}$ (this state is often denoted $\alpha$ or $\uparrow$ ); the other orientation corresponds to $m_{s}=-\frac{1}{2}$ (this state is denoted $\beta$ or $\downarrow$ ).

Other elementary particles have characteristic spin. For example, protons and neutrons are spin- $\frac{1}{2}$ particles (i.e. $s=\frac{1}{2}$ ). Because the masses of a proton and a neutron are so much greater than the mass of an electron, yet they all have the same spin angular momentum, the classical picture would be of these two particles spinning much more slowly than an electron. Some mesons, another variety of fundamental particle, are spin-1 particles (i.e. $s=1$ ), as are some atomic nuclei, but for our purposes the most important spin-1 particle is the photon. The importance of photon spin in spectroscopy is explained in Topic 11A; nuclear spin is the basis of nuclear magnetic resonance (Topic 12A).

\subsection*{Brief illustration 8B.2}
The magnitude of the spin angular momentum, like any angular momentum, is $\{s(s+1)\}^{1 / 2} \hbar$. For any spin- $\frac{1}{2}$ particle, not only electrons, this angular momentum is $\left(\frac{3}{4}\right)^{1 / 2} \hbar=0.866 \hbar$, or $9.13 \times 10^{-35} \mathrm{~J} \mathrm{~s}$. The component on the $z$-axis is $m_{s} \hbar$, which for a spin $-\frac{1}{2}$ particle is $\pm \frac{1}{2} \hbar$, or $\pm 5.27 \times 10^{-35} \mathrm{~J} \mathrm{~s}$.

Particles with half-integral spin are called fermions and those with integral spin (including 0) are called bosons. Thus, electrons and protons are fermions; photons are bosons. It is a very deep feature of nature that all the elementary particles that constitute matter are fermions whereas the elementary particles that transmit the forces that bind fermions together are all bosons. Photons, for example, transmit the electromagnetic force that binds together electrically charged particles. Matter, therefore, is an assembly of fermions held together by forces conveyed by bosons.

\subsection*{(b) The Pauli principle}
With the concept of spin established, it is possible to resume discussion of the electronic structures of atoms. Lithium, with $Z=3$, has three electrons. The first two occupy a 1 s orbital drawn even more closely than in He around the more highly charged nucleus. The third electron, however, does not join the first two in the 1s orbital because that configuration is forbidden by the Pauli exclusion principle:

No more than two electrons may occupy any given orbital, and if two do occupy one orbital, then their spins must be paired.\\
\includegraphics[max width=\textwidth, center]{2025_02_27_d4b95d9718f0b9e2e6b9g-350}

Electrons with paired spins, denoted $\uparrow \downarrow$, have zero net spin angular momentum because the spin of one electron is cancelled by the spin of the other. Specifically, one electron has $m_{s}=+\frac{1}{2}$ the other has $m_{s}=-\frac{1}{2}$ and in the vector model they are orientated on their respective cones so that the resultant spin is zero (Fig. 8B.3). The exclusion principle is the key to the structure of complex atoms, to chemical periodicity, and to molecular structure. It was proposed by Wolfgang Pauli in 1924 when he was trying to account for the absence of some lines in the spectrum of helium. Later he was able to derive a very general form of the principle from theoretical considerations.

The Pauli exclusion principle is a special case of a general statement called the Pauli principle:\\
\includegraphics[max width=\textwidth, center]{2025_02_27_d4b95d9718f0b9e2e6b9g-350(2)}

Figure 8B. 3 Electrons with paired spins have zero resultant spin angular momentum. They can be represented by two vectors that lie at an indeterminate position on the cones shown here, but wherever one lies on its cone, the other points in the opposite direction; their resultant is zero.

When the labels of any two identical fermions are exchanged, the total wavefunction changes sign; when the labels of any two identical bosons are exchanged, the sign of the total wavefunction remains the same.

By 'total wavefunction' is meant the entire wavefunction, including the spin of the particles.

To see that the Pauli principle implies the Pauli exclusion principle, consider the wavefunction for two electrons, $\Psi(1,2)$. The Pauli principle implies that it is a fact of nature (which has its roots in the theory of relativity) that the wavefunction must change sign if the labels 1 and 2 are interchanged wherever they occur in the function:


\begin{equation*}
\Psi(2,1)=-\Psi(1,2) \tag{8B.2}
\end{equation*}


Suppose the two electrons in a two-electron atom occupy the same orbital $\psi$, then in the orbital approximation the overall spatial wavefunction is $\psi\left(\boldsymbol{r}_{1}\right) \psi\left(\boldsymbol{r}_{2}\right)$, which for simplicity will be denoted $\psi(1) \psi(2)$. To apply the Pauli principle, it is necessary to consider the total wavefunction, the wavefunction including spin. There are several possibilities for two electrons: both $\alpha$, denoted $\alpha(1) \alpha(2)$, both $\beta$, denoted $\beta(1) \beta(2)$, and one $\alpha$ and the other $\beta$, denoted either $\alpha(1) \beta(2)$ or $\alpha(2) \beta(1)$. Because it is not possible to know which electron is $\alpha$ and which is $\beta$, in the last case it is appropriate to express the spin states as the (normalized) linear combinations ${ }^{1}$


\begin{align*}
& \sigma_{+}(1,2)=\left(\frac{1}{2^{1 / 2}}\right)\{\alpha(1) \beta(2)+\beta(1) \alpha(2)\} \\
& \sigma_{-}(1,2)=\left(\frac{1}{2^{1 / 2}}\right)\{\alpha(1) \beta(2)-\beta(1) \alpha(2)\} \tag{8B.3}
\end{align*}


These combinations allow one spin to be $\alpha$ and the other $\beta$ with equal probability; the former corresponds to parallel spins (the individual spins do not cancel) and the latter to paired spins (the individual spins cancel). The total wavefunction of the system is therefore the product of the orbital part and one of the four spin states:

\[
\begin{array}{ll}
\psi(1) \psi(2) \alpha(1) \alpha(2) & \psi(1) \psi(2) \beta(1) \beta(2) \\
\psi(1) \psi(2) \sigma_{+}(1,2) & \psi(1) \psi(2) \sigma_{-}(1,2) \tag{8B.4}
\end{array}
\]

The Pauli principle says that for a wavefunction to be acceptable (for electrons), it must change sign when the electrons are exchanged. In each case, exchanging the labels 1 and 2 converts $\psi(1) \psi(2)$ into $\psi(2) \psi(1)$, which is the same, because the order of multiplying the functions does not change the value of the product. The same is true of $\alpha(1) \alpha(2)$ and $\beta(1) \beta(2)$. Therefore, $\psi(1) \psi(2) \alpha(1) \alpha(2)$ and $\psi(1) \psi(2) \beta(1) \beta(2)$ are not

\footnotetext{${ }^{1}$ A stronger justification for taking these linear combinations is that they correspond to eigenfunctions of the total spin operators $S^{2}$ and $S_{z}$, with $M_{\mathrm{S}}=$ 0 and, respectively, $S=1$ and 0 .
}
allowed, because they do not change sign. When the labels are exchanged the combination $\sigma_{+}(1,2)$ becomes
$$
\sigma_{+}(2,1)=\left(\frac{1}{2^{1 / 2}}\right)\{\alpha(2) \beta(1)+\beta(2) \alpha(1)\}=\sigma_{+}(1,2)
$$
because the central term is simply the original function written in a different order. The product $\psi(1) \psi(2) \sigma_{+}(1,2)$ is therefore also disallowed. Finally, consider $\sigma_{-}(1,2)$ :\\
\$\$

\[
\begin{aligned}
\sigma_{-}(2,1) & =\left(\frac{1}{2^{1 / 2}}\right)\{\alpha(2) \beta(1)-\beta(2) \alpha(1)\} \\
& =-\left(\frac{1}{2^{1 / 2}}\right)\{\alpha(1) \beta(2)-\beta(1) \alpha(2)\}=-\sigma_{-}(1,2)
\end{aligned}
\]

\$\$

The combination $\psi(1) \psi(2) \sigma_{-}(1,2)$ therefore does change sign (it is 'antisymmetric') and is acceptable.

In summary, only one of the four possible states is allowed by the Pauli principle: the one that survives has paired $\alpha$ and $\beta$ spins. This is the content of the Pauli exclusion principle. The exclusion principle (but not the more general Pauli principle) is irrelevant when the orbitals occupied by the electrons are different, and both electrons may then have, but need not have, the same spin state. In each case the overall wavefunction must still be antisymmetric and must satisfy the Pauli principle itself.

Now returning to lithium, $\mathrm{Li}(Z=3)$, the third electron cannot enter the 1 s orbital because that orbital is already full: the K shell (the shell with $n=1$, Topic 8 A ) is complete and the two electrons form a closed shell, a shell in which all the orbitals are fully occupied. Because a similar closed shell is characteristic of the He atom, it is commonly denoted [He]. The third electron cannot enter the $K$ shell and must occupy the next available orbital, which is one with $n=2$ and hence belonging to the L shell (which consists of the four orbitals with $n=2$ ). It is now necessary to decide whether the next available orbital is the 2 s orbital or a 2 p orbital, and therefore whether the lowest energy configuration of the atom is $[\mathrm{He}] 2 \mathrm{~s}^{1}$ or $[\mathrm{He}] 2 \mathrm{p}^{1}$.

\subsection*{8B.3 The building-up principle}
Unlike in hydrogenic atoms, the 2 s and 2 p orbitals (and, in general, the subshells of a given shell) do not have the same energy in many-electron atoms.

\subsection*{(a) Penetration and shielding}
An electron in a many-electron atom experiences a Coulombic repulsion from all the other electrons present. If the electron is at a distance $r$ from the nucleus, it experiences an average repulsion that can be represented by a point negative charge located at the nucleus and equal in magnitude to the total charge\\
of all the other electrons within a sphere of radius $r$ (Fig. 8B.4). This property is a conclusion of classical electrostatics, where the effect of a spherical distribution of charge can be represented by a point charge of the same magnitude located at its centre. The effect of this point negative charge is to reduce the full charge of the nucleus from $Z e$ to $Z_{\text {eff }} e$, the effective nuclear charge. In everyday parlance, $Z_{\text {eff }}$ itself is commonly referred to as the 'effective nuclear charge'. The electron is said to experience a shielded nuclear charge, and the difference between $Z$ and $Z_{\text {eff }}$ is called the shielding constant, $\sigma$ :


\begin{equation*}
Z_{\text {eff }}=Z-\sigma \quad \text { Nuclear shielding } \tag{8B.5}
\end{equation*}


The electrons do not actually 'block' the full Coulombic attraction of the nucleus: the shielding constant is simply a way of expressing the net outcome of the nuclear attraction and the electronic repulsions in terms of a single equivalent charge at the centre of the atom.

The shielding constant is different for $s$ and $p$ electrons because they have different radial distribution functions and therefore respond to the other electrons in the atom to different extents (Fig. 8B.5). An selectron has a greater penetration through inner shells than a $p$ electron, in the sense that an $s$ electron is more likely to be found close to the nucleus than a p electron of the same shell. Because only electrons inside the sphere defined by the location of the electron of interest contribute to shielding, an electron experiences less shielding than a p electron. Consequently, as a result of the combined effects of penetration and shielding, an $s$ electron is more tightly bound than a p electron of the same shell. Similarly, a d electron penetrates less than a p electron of the same shell (recall that a d orbital is proportional to $r^{2}$ close to the nucleus, whereas a p orbital is proportional to $r$, so the amplitude of a d orbital is smaller there than that of a p orbital), and therefore experiences more shielding.

Shielding constants for different types of electrons in atoms have been calculated from wavefunctions obtained by nu-\\
\includegraphics[max width=\textwidth, center]{2025_02_27_d4b95d9718f0b9e2e6b9g-352}

Figure 8B. 4 An electron at a distance $r$ from the nucleus experiences a Coulombic repulsion from all the electrons within a sphere of radius $r$. This repulsion is equivalent to that from a point negative charge located on the nucleus. The negative charge reduces the effective nuclear charge of the nucleus from $Z e$ to $Z_{\text {eff }} e$.\\
\includegraphics[max width=\textwidth, center]{2025_02_27_d4b95d9718f0b9e2e6b9g-352(1)}

Figure 8B.5 An electron in an s orbital (here a 3 s orbital) is more likely to be found close to the nucleus than an electron in a $p$ orbital of the same shell (note the closeness of the innermost peak of the 3 s orbital to the nucleus at $r=0$ ). Hence an s electron experiences less shielding and is more tightly bound than a p electron of the same shell.

Table 8B. 1 Effective nuclear charge*

\begin{center}
\begin{tabular}{llll}
\hline
Element & $\boldsymbol{Z}$ & Orbital & $\boldsymbol{Z}_{\text {eff }}$ \\
\hline
He & 2 & 1 s & 1.6875 \\
C & 6 & 1 s & 5.6727 \\
 &  & 2 s & 3.2166 \\
 & 2 p & 3.1358 &  \\
\hline
\end{tabular}
\end{center}

\begin{itemize}
  \item More values are given in the Resource section.\\
merical solution of the Schrödinger equation (Table 8B.1). In general, valence-shell s electrons do experience higher effective nuclear charges than $p$ electrons, although there are some discrepancies.
\end{itemize}

\subsection*{Brief illustration 8B.3}
The effective nuclear charge for $1 \mathrm{~s}, 2 \mathrm{~s}$, and 2 p electrons in a carbon atom are $5.6727,3.2166$, and 3.1358 , respectively. The radial distribution functions for these orbitals (Topic 8A) are generated by forming $P(r)=r^{2} R(r)^{2}$, where $R(r)$ is the radial wavefunction, which are given in Table 8A.1. The three radial distribution functions are plotted in Fig. 8B.6. As can be seen (especially in the magnified view close to the nucleus), the $s$ orbital has greater penetration than the $p$ orbital. The average radii of the 2 s and 2 p orbitals are 99 pm and 84 pm , respective$l y$, which shows that the average distance of a $2 s$ electron from the nucleus is greater than that of a 2 p orbital. To account for the lower energy of the 2 s orbital, the extent of penetration is more important than the average distance from the nucleus.

The consequence of penetration and shielding is that the energies of subshells of a shell in a many-electron atom (those\\
\includegraphics[max width=\textwidth, center]{2025_02_27_d4b95d9718f0b9e2e6b9g-353}

Figure 8B. 6 The radial distribution functions for electrons in a carbon atom, as calculated in Brief illustration 8B.3.\\
with the same values of $n$ but different values of $l$ ) in general lie in the order $\mathrm{s}<\mathrm{p}<\mathrm{d}<\mathrm{f}$. The individual orbitals of a given subshell (those with the same value of $l$ but different values of $m_{l}$ ) remain degenerate because they all have the same radial characteristics and so experience the same effective nuclear charge.

To complete the Li story, consider that, because the shell with $n=2$ consists of two subshells, with the 2 s subshell lower in energy than the 2 p subshell, the third electron occupies the $2 s$ orbital (the only orbital in that subshell). This occupation results in the ground-state configuration $1 s^{2} 2 s^{1}$, with the central nucleus surrounded by a complete helium-like shell of two 1 s electrons, and around that a more diffuse 2 s electron. The electrons in the outermost shell of an atom in its ground state are called the valence electrons because they are largely responsible for the chemical bonds that the atom forms (and 'valence', as explained in Focus 9, refers to the ability of an atom to form bonds). Thus, the valence electron in Li is a 2 s electron and its other two electrons belong to its core.

\subsection*{(b) Hund's rules}
The extension of the argument used to account for the structures of $\mathrm{H}, \mathrm{He}$, and Li is called the building-up principle, or the Aufbau principle, from the German word for "building up", and should be familiar from introductory courses. In brief, imagine the bare nucleus of atomic number $Z$, and then feed into the orbitals $Z$ electrons in succession. The order of occupation, following the shells and their subshells arranged in order of increasing energy, is

1s 2s 2p 3s 3p 4s 3d 4p 5s 4d 5p 6s\\
Each orbital may accommodate up to two electrons.

\subsection*{Brief illustration 8B.4}
Consider the carbon atom, for which $Z=6$ and there are six electrons to accommodate. Two electrons enter and fill the 1 s\\
orbital, two enter and fill the 2 s orbital, leaving two electrons to occupy the orbitals of the 2 p subshell. Hence the groundstate configuration of C is $1 \mathrm{~s}^{2} 2 \mathrm{~s}^{2} 2 \mathrm{p}^{2}$, or more succinctly $[\mathrm{He}] 2 \mathrm{~s}^{2} 2 \mathrm{p}^{2}$, with [He] the helium-like $1 \mathrm{~s}^{2}$ core.

It is possible to be more precise about the configuration of a carbon atom than in Brief illustration 8B.4. The last two electrons are expected to occupy different 2 p orbitals because they are then farther apart on average and repel each other less than if they were in the same orbital. Thus, one electron can be thought of as occupying the $2 \mathrm{p}_{x}$ orbital and the other the $2 \mathrm{p}_{y}$ orbital (the $x, y, z$ designation is arbitrary, and it would be equally valid to use the complex forms of these orbitals), and the lowest energy configuration of the atom is [He] $2 \mathrm{~s}^{2} 2 \mathrm{p}_{x}^{1} 2 \mathrm{p}_{y}^{1}$. The same rule applies whenever degenerate orbitals of a subshell are available for occupation. Thus, another rule of the building-up principle is:

\subsection*{Electrons occupy different orbitals of a given subshell before doubly occupying any one of them.}
For instance, nitrogen $(Z=7)$ has the ground-state configuration $[\mathrm{He}] 2 \mathrm{~s}^{2} 2 \mathrm{p}_{x}^{1} 2 \mathrm{p}_{y}^{1} 2 \mathrm{p}_{z}^{1}$, and only at oxygen $(Z=8)$ is a 2 p orbital doubly occupied, giving [He] $2 \mathrm{~s}^{2} 2 \mathrm{p}_{x}^{2} 2 \mathrm{p}_{y}^{1} 2 \mathrm{p}_{z}^{1}$.

When electrons occupy orbitals singly it is necessary to invoke Hund's maximum multiplicity rule:

An atom in its ground state adopts a configuration with the greatest number of unpaired electrons.

The explanation of Hund's rule is subtle, but it reflects the quantum mechanical property of spin correlation. In essence, the effect of spin correlation is to allow the atom to shrink slightly when the spins are parallel, so the electron-nucleus interaction is improved. As a consequence, in the ground state of the carbon atom, the two 2 p electrons have parallel spins, all three 2 p electrons in the N atoms have parallel spins, and the two 2 p electrons in different orbitals in the O atom have parallel spins (the two in the $2 \mathrm{p}_{x}$ orbital are necessarily paired). The effect can be explained by considering the Pauli principles and showing that electrons with parallel spins behave as if they have a tendency to stay apart, and hence repel each other less.

\subsection*{How is that done? 8B.2 Exploring the origins of spin correlation}
Suppose electron 1 is in orbital $a$ and described by a wavefunction $\psi_{a}\left(\boldsymbol{r}_{1}\right)$, and electron 2 is in orbital $b$ with wavefunction $\psi_{b}\left(\boldsymbol{r}_{2}\right)$. Then, in the orbital approximation, the joint spatial wavefunction of the electrons is the product $\Psi=\psi_{a}\left(\boldsymbol{r}_{1}\right) \psi_{b}\left(\boldsymbol{r}_{2}\right)$. However, this wavefunction is not acceptable, because it suggests that it is possible to know which electron is in\\
which orbital. According to quantum mechanics, the correct description is either of the two following wavefunctions:

$$
\Psi_{ \pm}=\left(\frac{1}{2^{1 / 2}}\right)\left\{\psi_{a}\left(\boldsymbol{r}_{1}\right) \psi_{b}\left(\boldsymbol{r}_{2}\right) \pm \psi_{b}\left(\boldsymbol{r}_{1}\right) \psi_{a}\left(\boldsymbol{r}_{2}\right)\right\}
$$

According to the Pauli principle, because $\Psi_{+}$is symmetrical under particle interchange, it must be multiplied by an antisymmetric spin state (the one denoted $\sigma_{-}$). That combination corresponds to a spin-paired state. Conversely, $\Psi_{-}$is antisymmetric, so it must be multiplied by one of the three symmetric spin states. These three symmetric states correspond to electrons with parallel spins (see Topic 8C for an explanation of this point).\\
Now consider the behaviour of the two wavefunctions $\Psi_{ \pm}$ when one electron approaches another, and $\boldsymbol{r}_{1}=\boldsymbol{r}_{2}$. As a result, $\Psi$ \_ vanishes, which means that there is zero probability of finding the two electrons at the same point in space when they have parallel spins. In contrast, the wavefunction $\Psi_{+}$does not vanish when the two electrons are at the same point in space. Because the two electrons have different relative spatial distributions depending on whether their spins are parallel or not, it follows that their Coulombic interaction is different, and hence that the two states described by these wavefunctions have different energies, with the spin-parallel state lower in energy than the spin-paired state.

Neon, with $Z=10$, has the configuration $[\mathrm{He}] 2 \mathrm{~s}^{2} 2 \mathrm{p}^{6}$, which completes the L shell. This closed-shell configuration is denoted [ Ne ], and acts as a core for subsequent elements. The next electron must enter the 3 s orbital and begin a new shell, so an Na atom, with $Z=11$, has the configuration $[\mathrm{Ne}] 3 \mathrm{~s}^{1}$. Like lithium with the configuration $[\mathrm{He}] 2 \mathrm{~s}^{1}$, sodium has a single $s$ electron outside a complete core. This analysis hints at the origin of chemical periodicity. The L shell is completed by eight electrons, so the element with $Z=3(\mathrm{Li})$ should have similar properties to the element with $Z=11(\mathrm{Na})$. Likewise, $\mathrm{Be}(Z=$ 4) should be similar to $Z=12(\mathrm{Mg})$, and so on, up to the noble gases $\mathrm{He}(Z=2), \mathrm{Ne}(Z=10)$, and $\mathrm{Ar}(Z=18)$.

At potassium $(Z=19)$ the next orbital in line for occupation is 4 s : this orbital is brought below 3d by the effects of penetration and shielding, and the ground state configuration is $[\mathrm{Ar}] 4 \mathrm{~s}^{1}$. Calcium $(Z=20)$ is likewise $[\mathrm{Ar}] 4 \mathrm{~s}^{2}$. At this stage the five 3d orbitals are in line for occupation, but there are complications arising from the energy changes arising from the interaction of the electrons in the valence shell, and penetration arguments alone are no longer reliable.

Calculations of the type discussed in Section 8B. 4 show that for the atoms from scandium to zinc the energies of the 3d orbitals are always lower than the energy of the 4 s orbital, in spite of the greater penetration of a 4 s electron. However, spectroscopic results show that Sc has the configuration $[\mathrm{Ar}] 3 \mathrm{~d}^{1} 4 \mathrm{~s}^{2}$, not $[\mathrm{Ar}] 3 \mathrm{~d}^{3}$ or $[\mathrm{Ar}] 3 \mathrm{~d}^{2} 4 \mathrm{~s}^{1}$. To understand this observation, consider the nature of electron-electron repulsions in 3d and 4 s orbitals. Because the average distance of a 3d electron from\\
the nucleus is less than that of a 4 s electron, two 3d electrons are so close together that they repel each other more strongly than two 4 s electrons do and $3 \mathrm{~d}^{2}$ and $3 \mathrm{~d}^{3}$ configurations are disfavoured. As a result, Sc has the configuration [Ar] $3 \mathrm{~d}^{1} 4 \mathrm{~s}^{2}$ rather than the two alternatives, for then the strong electronelectron repulsions in the 3d orbitals are minimized. The total energy of the atom is lower despite the cost of allowing electrons to populate the high energy 4 s orbital (Fig. 8B.7). The effect just described is generally true for scandium to zinc, so their electron configurations are of the form $[\mathrm{Ar}] 3 \mathrm{~d}^{n} 4 \mathrm{~s}^{2}$, where $n=1$ for scandium and $n=10$ for zinc. Two notable exceptions, which are observed experimentally, are Cr , with electron configuration $[\mathrm{Ar}] 3 \mathrm{~d}^{5} 4 \mathrm{~s}^{1}$, and Cu , with electron configuration $[\mathrm{Ar}] 3 \mathrm{~d}^{10} 4 \mathrm{~s}^{1}$. At gallium, these complications disappear and the building-up principle is used in the same way as in preceding periods. Now the 4 s and 4 p subshells constitute the valence shell, and the period terminates with krypton. Because 18 electrons have intervened since argon, this row is the first 'long period' of the periodic table.

At this stage it becomes apparent that sequential occupation of the orbitals in successive shells results in periodic similarities in the electronic configurations. This periodicity of structure accounts for the formulation of the periodic table (see inside the back cover). The vertical columns of the periodic table are called groups and (in the modern convention) numbered from 1 to 18 . Successive rows of the periodic table are called periods, the number of the period being equal to the principal quantum number of the valence shell.

The periodic table is divided into $s, p, d$, and $f$ blocks, according to the subshell that is last to be occupied in the formulation of the electronic configuration of the atom. The members of the d block (specifically the members of Groups $3-11$ in the d block) are also known as the transition metals; those of the f block (which is not divided into numbered groups) are sometimes called the inner transition metals. The upper row of the f block (Period 6) consists of the lanthanoids\\
\includegraphics[max width=\textwidth, center]{2025_02_27_d4b95d9718f0b9e2e6b9g-354}

Figure 8B.7 Strong electron-electron repulsions in the 3d orbitals are minimized in the ground state of Sc if the atom has the configuration $[\mathrm{Ar}] 3 \mathrm{~d}^{1} 4 \mathrm{~s}^{2}$ (shown on the left) instead of [Ar] $3 d^{2} 4 s^{1}$ (shown on the right). The total energy of the atom is lower when it has the [Ar] $3 d^{1} 4 s^{2}$ configuration despite the cost of populating the high energy 4 s orbital.\\
(still commonly the 'lanthanides') and the lower row (Period 7) consists of the actinoids (still commonly the 'actinides').

The configurations of cations of elements in the $s, p$, and $d$ blocks of the periodic table are derived by removing electrons from the ground-state configuration of the neutral atom in a specific order. First, remove valence p electrons, then valence $s$ electrons, and then as many d electrons as are necessary to achieve the specified charge. The configurations of anions of the p-block elements are derived by continuing the buildingup procedure and adding electrons to the neutral atom until the configuration of the next noble gas has been reached.

\subsection*{Brief illustration 8B.5}
Because the configuration of vanadium is $[\mathrm{Ar}] 3 \mathrm{~d}^{3} 4 \mathrm{~s}^{2}$, the $\mathrm{V}^{2+}$ cation has the configuration $[\mathrm{Ar}] 3 \mathrm{~d}^{3}$. It is reasonable to remove the more energetic $4 s$ electrons in order to form the cation, but it is not obvious why the $[\mathrm{Ar}] 3 \mathrm{~d}^{3}$ configuration is preferred in $\mathrm{V}^{2+}$ over the $[\mathrm{Ar}] 3 \mathrm{~d}^{1} 4 \mathrm{~s}^{2}$ configuration, which is found in the isoelectronic Sc atom. Calculations show that the energy difference between $[\mathrm{Ar}] 3 \mathrm{~d}^{3}$ and $[\mathrm{Ar}] 3 \mathrm{~d}^{1} 4 \mathrm{~s}^{2}$ depends on $Z_{\text {eff }}$. As $Z_{\text {eff }}$ increases, transfer of a 4 s electron to a 3d orbital becomes more favourable because the electron-electron repulsions are compensated by attractive interactions between the nucleus and the electrons in the spatially compact 3d orbital. Indeed, calculations reveal that, for a sufficiently large $Z_{\text {eff }}$, Ar$] 3 \mathrm{~d}^{3}$ is lower in energy than $[\mathrm{Ar}] 3 \mathrm{~d}^{1} 4 \mathrm{~s}^{2}$. This conclusion explains why $\mathrm{V}^{2+}$ has a $[\mathrm{Ar}] 3 \mathrm{~d}^{3}$ configuration and also accounts for the observed [Ar] $4 \mathrm{~s}^{0} 3 \mathrm{~d}^{n}$ configurations of the $\mathrm{M}^{2+}$ cations of Sc through Zn.

\subsection*{(c) Atomic and ionic radii}
The atomic radius of an element is half the distance between the centres of neighbouring atoms in a solid (such as Cu ) or, for non-metals, in a homonuclear molecule (such as $\mathrm{H}_{2}$ or $\mathrm{S}_{8}$ ). As seen in Table 8B. 2 and Fig. 8B.8, atomic radii tend to decrease from left to right across a period of the periodic table, and increase down each group. The decrease across a period can be traced to the increase in nuclear charge, which draws the electrons in closer to the nucleus. The increase in nuclear charge is partly cancelled by the increase in the number of electrons, but because electrons are spread over a region of space, one electron does not fully shield one nuclear charge, so the increase in nuclear charge dominates. The increase in atomic radius down a group (despite the increase in nuclear charge) is explained by the fact that the valence shells of successive periods correspond to higher principal quantum numbers. That is, successive periods correspond to the start and then completion of successive (and more distant) shells of the atom that surround each other like the successive layers of an onion. The need to occupy a more distant shell leads to a larger atom despite the increased nuclear charge.

Table 8B. 2 Atomic radii of main-group elements, r/pm*

\begin{center}
\begin{tabular}{lllllll}
\hline
Li & Be & B & C & N & O & F \\
157 & 112 & 88 & 77 & 74 & 66 & 64 \\
Na & Mg & Al & Si & P & S & Cl \\
191 & 160 & 143 & 118 & 110 & 104 & 99 \\
K & Ca & Ga & Ge & As & Se & Br \\
235 & 197 & 153 & 122 & 121 & 117 & 114 \\
Rb & Sr & In & Sn & Sb & Te & I \\
250 & 215 & 167 & 158 & 141 & 137 & 133 \\
Cs & Ba & Tl & Pb & Bi & Po &  \\
272 & 224 & 171 & 175 & 182 & 167 &  \\
\hline
\end{tabular}
\end{center}

\begin{itemize}
  \item More values are given in the Resource section.
\end{itemize}

A modification of the increase down a group is encountered in Period 6, for the radii of the atoms in the d block and in the following atoms of the $p$ block are not as large as would be expected by simple extrapolation down the group. The reason can be traced to the fact that in Period 6 the forbitals are in the process of being occupied. An $f$ electron is a very inefficient shielder of nuclear charge (for reasons connected with its radial extension), and as the atomic number increases from La to Lu , there is a considerable contraction in radius. By the time the d block resumes (at hafnium, Hf), the poorly shielded but considerably increased nuclear charge has drawn in the surrounding electrons, and the atoms are compact. They are so compact, that the metals in this region of the periodic table (iridium to lead) are very dense. The reduction in radius below that expected by extrapolation from preceding periods is called the lanthanide contraction.

The ionic radius of an element is its share of the distance between neighbouring ions in an ionic solid. That is, the distance between the centres of a neighbouring cation and anion is the sum of the two ionic radii. The size of the 'share' leads\\
\includegraphics[max width=\textwidth, center]{2025_02_27_d4b95d9718f0b9e2e6b9g-355}

Figure 8B. 8 The variation of atomic radius through the periodic table. Note the contraction of radius following the lanthanoids in Period 6 (following Lu, lutetium).

Table 8B. 3 Ionic radii, r/pm*

\begin{center}
\begin{tabular}{llllll}
\hline
$\mathrm{Li}^{+}(4)$ & $\mathrm{Be}^{2+}(4)$ & $\mathrm{B}^{3+}(4)$ & $\mathrm{N}^{3-}$ & $\mathrm{O}^{2-}(6)$ & $\mathrm{F}^{-}(6)$ \\
59 & 27 & 12 & 171 & 140 & 133 \\
$\mathrm{Na}^{+}(6)$ & $\mathrm{Mg}^{2+}(6)$ & $\mathrm{Al}^{3+}(6)$ & $\mathrm{P}^{3-}$ & $\mathrm{S}^{2-}(6)$ & $\mathrm{Cl}^{-}(6)$ \\
102 & 72 & 53 & 212 & 184 & 181 \\
$\mathrm{~K}^{+}(6)$ & $\mathrm{Ca}^{2+}(6)$ & $\mathrm{Ga}^{3+}(6)$ & $\mathrm{As}^{3-}(6)$ & $\mathrm{Se}^{2-}(6)$ & $\mathrm{Br}^{-}(6)$ \\
138 & 100 & 62 & 222 & 198 & 196 \\
$\mathrm{Rb}^{+}(6)$ & $\mathrm{Sr}^{2+}(6)$ & $\mathrm{In}^{3+}(6)$ &  & $\mathrm{Te}^{2-}(6)$ & $\mathrm{I}^{-}(6)$ \\
149 & 116 & 79 &  & 221 & 220 \\
$\mathrm{Cs}^{+}(6)$ & $\mathrm{Ba}^{2+}(6)$ & $\mathrm{Tl}^{3+}(6)$ &  &  &  \\
167 & 136 & 88 &  &  &  \\
\end{tabular}
\end{center}

*Numbers in parentheses are the coordination numbers of the ions, the numbers of species (for example, counterions, solvent molecules) around the ions. Values for ions without a coordination number stated are estimates. More values are given in the Resource section.\\
to some ambiguity in the definition. One common definition sets the ionic radius of $\mathrm{O}^{2-}$ equal to 140 pm , but there are other scales, and care must be taken not to mix them. Ionic radii also vary with the number of counterions (ions of opposite charge) around a given ion; unless otherwise stated, the values in this text have been corrected to correspond to an environment of six counterions.

When an atom loses one or more valence electrons to form a cation, the remaining atomic core is smaller than the parent atom. Therefore, a cation is invariably smaller than its parent atom. For example, the atomic radius of Na , with the configuration $[\mathrm{Ne}] 3 \mathrm{~s}^{1}$, is 191 pm , but the ionic radius of $\mathrm{Na}^{+}$, with the configuration [Ne], is only 102 pm (Table 8B.3). Like atomic radii, cation radii increase down each group because electrons are occupying shells with higher principal quantum numbers.\\
An anion is larger than its parent atom because the electrons added to the valence shell repel one another. Without a compensating increase in the nuclear charge, which would draw the electrons closer to the nucleus and each other, the ion expands. The variation in anion radii shows the same trend as that for atoms and cations, with the smallest anions at the upper right of the periodic table, close to fluorine (Table 8B.3).

\subsection*{Brief illustration 8B.6}
The $\mathrm{Ca}^{2+}, \mathrm{K}^{+}$, and $\mathrm{Cl}^{-}$ions have the configuration [Ar]. However, their radii differ because they have different nuclear charges. The $\mathrm{Ca}^{2+}$ ion has the largest nuclear charge, so it has the strongest attraction for the electrons and the smallest radius. The $\mathrm{Cl}^{-}$ion has the lowest nuclear charge of the three ions and, as a result, the largest radius.

\subsection*{(d) Ionization energies and electron affinities}
The minimum energy necessary to remove an electron from a many-electron atom in the gas phase is the first ionization energy, $I_{1}$, of the element. The second ionization energy, $I_{2}$, is the minimum energy needed to remove a second electron (from the singly charged cation). The variation of the first ionization energy through the periodic table is shown in Fig. 8B. 9 and some numerical values are given in Table 8B.4.

The electron affinity, $E_{\text {ea }}$, is the energy released when an electron attaches to a gas-phase atom (Table 8B.5). In a common, logical (given its name), but not universal convention (which is adopted here), the electron affinity is positive if energy is released when the electron attaches to the atom. That is, $E_{\text {ea }}>0$ implies that electron attachment is exothermic.

As will be familiar from introductory chemistry, ionization energies and electron affinities show periodicities. The former is more regular and concentrated on here. Lithium has a low first ionization energy because its outermost electron is well shielded from the nucleus by the core ( $Z_{\text {eff }}=1.3$, compared with $Z=3)$. The ionization energy of $\operatorname{Be}(Z=4)$ is greater but that of $B$ is lower because in the latter the outermost electron occupies a 2 p orbital and is less strongly bound than if it had been a $2 s$ electron. The ionization energy increases from $B$ to N on account of the increasing nuclear charge. However, the ionization energy of $O$ is less than would be expected by simple extrapolation. The explanation is that at oxygen a 2 p orbital must become doubly occupied, and the electron-electron repulsions are increased above what would be expected by simple extrapolation along the row. In addition, the loss of a 2 p electron results in a configuration with a half-filled subshell (like that of N ), which is an arrangement of low energy, so the energy of $\mathrm{O}^{+}+\mathrm{e}^{-}$is lower than might be expected, and the ionization energy is correspondingly low too. (The kink is less pronounced in the next row, between phosphorus and sulfur\\
\includegraphics[max width=\textwidth, center]{2025_02_27_d4b95d9718f0b9e2e6b9g-356}

Figure 8B. 9 The first ionization energies of the elements plotted against atomic number.

Table 8B.4 First and second ionization energies*

\begin{center}
\begin{tabular}{lll}
\hline
Element & $I_{1} /\left(\mathbf{k J ~ m o l}^{-1}\right)$ & $\boldsymbol{I}_{2} /\left(\mathbf{k J ~ m o l}^{-1}\right)$ \\
\hline
H & 1312 &  \\
He & 2372 & 5251 \\
Mg & 738 & 1451 \\
Na & 496 & 4562 \\
\hline
\end{tabular}
\end{center}

\begin{itemize}
  \item More values are given in the Resource section.
\end{itemize}

Table 8B.5 Electron affinities, $E_{\mathrm{a}} /\left(\mathrm{kJ} \mathrm{mol}^{-1}\right)^{*}$

\begin{center}
\begin{tabular}{lrll}
\hline
Cl & 349 &  &  \\
F & 322 &  &  \\
H & 73 &  &  \\
O & 141 & $\mathrm{O}^{-}$ & -844 \\
\hline
\end{tabular}
\end{center}

\begin{itemize}
  \item More values are given in the Resource section.\\
because their orbitals are more diffuse.) The values for $\mathrm{O}, \mathrm{F}$, and Ne fall roughly on the same line, the increase of their ionization energies reflecting the increasing attraction of the more highly charged nuclei for the outermost electrons.
\end{itemize}

The outermost electron in sodium $(Z=11)$ is 3 s . It is far from the nucleus, and the latter's charge is shielded by the compact, complete neon-like core, with the result that $Z_{\text {eff }} \approx 2.5$. As a result, the ionization energy of Na is substantially lower than that of $\mathrm{Ne}\left(Z=10, Z_{\text {eff }} \approx 5.8\right)$. The periodic cycle starts again along this row, and the variation of the ionization energy can be traced to similar reasons.

Electron affinities are greatest close to fluorine, for the incoming electron enters a vacancy in a compact valence shell and can interact strongly with the nucleus. The attachment of an electron to an anion (as in the formation of $\mathrm{O}^{2-}$ from $\mathrm{O}^{-}$) is invariably endothermic, so $E_{\text {ea }}$ is negative. The incoming electron is repelled by the charge already present. Electron affinities are also small, and may be negative, when an electron enters an orbital that is far from the nucleus (as in the heavier alkali metal atoms) or is forced by the Pauli principle to occupy a new shell (as in the noble gas atoms).

\subsection*{88.4 Self-consistent field orbitals}
The preceding treatment of the electronic configuration of many-electron species is only approximate because of the complications introduced by electron-electron interactions. However, computational techniques are available that give reliable approximate solutions for the wavefunctions and energies. The techniques were originally introduced by D.R. Hartree (before computers were available) and then modified by V. Fock to take into account the Pauli principle correctly. In broad outline, the Hartree-Fock self-consistent field (HFSCF) procedure is as follows.

Start with an idea of the structure of the atom as suggested by the building-up principle. In the Ne atom, for instance, the principle suggests the configuration $1 s^{2} 2 s^{2} 2 p^{6}$ with the orbitals approximated by hydrogenic atomic orbitals with the appropriate effective nuclear charges. Now consider one of the 2 p electrons. A Schrödinger equation can be written for this electron by ascribing to it a potential energy due to the nuclear attraction and the average repulsion from the other electrons. Although the equation is for the 2 p orbital, that repulsion, and therefore the equation, depends on the wavefunctions of all the other occupied orbitals in the atom. To solve the equation, guess an approximate form of the wavefunctions of all the other orbitals and then solve the Schrödinger equation for the $2 p$ orbital. The procedure is then repeated for the 1 s and 2 s orbitals. This sequence of calculations gives the form of the $2 p, 2 s$, and $1 s$ orbitals, and in general they will differ from the set used to start the calculation. These improved orbitals can be used in another cycle of calculation, and a second improved set of orbitals and a better energy are obtained. The recycling continues until the orbitals and energies obtained are insignificantly different from those used at the start of the current cycle. The solutions are then self-consistent and accepted as solutions of the problem.

The outcomes of HF-SCF calculations are radial distribution functions that show the grouping of electron density into shells, as the building-up principle suggests. These calculations therefore support the qualitative discussions that are used to explain chemical periodicity. They also extend that discussion considerably by providing detailed wavefunctions and precise energies.

\subsection*{Checklist of concepts}
\begin{enumerate}
  \item In the orbital approximation, each electron is regarded as being described by its own wavefunction; the overall wavefunction of a many-electron atom is the product of the orbital wavefunctions.2. The configuration of an atom is the statement of its occupied orbitals.3. The Pauli exclusion principle, a special case of the Pauli principle, limits to two the number of electrons that can occupy a given orbital.4. In many-electron atoms, s orbitals lie at a lower energy than p orbitals of the same shell due to the combined effects of penetration and shielding.5. The building-up principle is a procedure for predicting the ground state electron configuration of an atom.6. Electrons occupy different orbitals of a given subshell before doubly occupying any one of them.7. An atom in its ground state adopts a configuration with the greatest number of unpaired electrons.8. The atomic radius of an element is half the distance between the centres of neighbouring atoms in a solid or in a homonuclear molecule.9. The ionic radius of an element is its share of the distance between neighbouring ions in an ionic solid.10. The first ionization energy is the minimum energy necessary to remove an electron from a many-electron atom in the gas phase.\\
$\square$ 11. The second ionization energy is the minimum energy needed to remove an electron from a singly charged cation.\\
$\square$ 12. The electron affinity is the energy released when an electron attaches to a gas-phase atom.13. The atomic radius, ionization energy, and electron affinity vary periodically through the periodic table.\\
$\square$ 14. The Schrödinger equation for many-electron atoms is solved numerically and iteratively until the solutions are self-consistent.
\end{enumerate}

\subsection*{Checklist of equations}
\begin{center}
\begin{tabular}{llll}
\hline
Property & Equation & Comment & Equation number \\
\hline
Orbital approximation & $\Psi\left(\boldsymbol{r}_{1}, \boldsymbol{r}_{2}, \ldots\right)=\psi\left(\boldsymbol{r}_{1}\right) \psi\left(\boldsymbol{r}_{2}\right) \ldots$ &  & 8 B .1 \\
Effective nuclear charge & $Z_{\text {eff }}=Z-\sigma$ & The charge is this number times $e$ & 8 B .5 \\
\hline
\end{tabular}
\end{center}

\section*{TOPIC 8C Atomic spectra}
\subsection*{> Why do you need to know this material?}
A knowledge of the energies of electrons in atoms is essential for understanding many chemical properties and chemical bonding.

\subsection*{What is the key idea?}
The frequency and wavenumber of radiation emitted or absorbed when atoms undergo electronic transitions provide detailed information about their electronic energy states.

\subsection*{> What do you need to know already?}
This Topic draws on knowledge of the energy levels of hydrogenic atoms (Topic 8A) and the configurations of many-electron atoms (Topic 8B). In places, it uses the properties of angular momentum (Topic 7F).

The general idea behind atomic spectroscopy is straightforward: lines in the spectrum (in either emission or absorption) occur when the electron distribution in an atom undergoes a transition, a change of state, in which its energy changes by $\Delta E$. This transition leads to the emission or is accompanied by absorption of a photon of frequency $v=|\Delta E| / h$ and wavenumber $\tilde{v}=|\Delta E| / h c$. In spectroscopy, transitions are said to take place between two terms. Broadly speaking, a term is simply another name for the energy level of an atom, but as this Topic progresses its full significance will become clear.

\subsection*{8C. 1 The spectra of hydrogenic atoms}
Not all transitions between the possible terms are observed. Spectroscopic transitions are allowed, if they can occur, or forbidden, if they cannot occur. A selection rule is a statement about which transitions are allowed.

The origin of selection rules can be identified by considering transitions in hydrogenic atoms. A photon has an intrinsic spin angular momentum corresponding to $s=1$ (Topic 8 B ). Because total angular momentum is conserved in a transition, the angular momentum of the electron must change to compensate for the angular momentum carried away by the photon. Thus, an electron in a d orbital $(l=2)$ cannot make a transition into an $s$ orbital $(l=0)$ because the photon can-\\
not carry away enough angular momentum. Similarly, an $s$ electron cannot make a transition to another s orbital, because there would then be no change in the angular momentum of the electron to make up for the angular momentum carried away by the photon. A more formal treatment of selection rules requires mathematical manipulation of the wavefunctions for the initial and final states of the atom.

\subsection*{How is that done? 8C. 1 Identifying selection rules}
The underlying classical idea behind a spectroscopic transition is that, for an atom or molecule to be able to interact with the electromagnetic field and absorb or create a photon of frequency $v$, it must possess, at least transiently, a dipole oscillating at that frequency. The consequences of this idea are explored in the following steps.

Step 1 Write an expression for the transition dipole moment\\
The transient dipole is expressed quantum mechanically as the transition dipole moment, $\mu_{\mathrm{f}}$, between the initial and final states i and f , where ${ }^{1}$


\begin{equation*}
\mu_{\mathrm{fi}}=\int \psi_{\mathrm{f}}^{\star} \hat{\mu} \psi_{\mathrm{i}} \mathrm{~d} \tau \tag{8C.1}
\end{equation*}


and $\hat{\mu}$ is the electric dipole moment operator. For a oneelectron atom $\hat{\mu}$ is multiplication by -er. Because $\boldsymbol{r}$ is a vector with components $x, y$, and $z, \hat{\mu}$ is also a vector, with components $\mu_{x}=-e x, \mu_{y}=-e y$, and $\mu_{z}=-e z$. If the transition dipole moment is zero, then the transition is forbidden; the transition is allowed if the transition moment is non-zero.

Step 2 Formulate the integrand in terms of spherical harmonics\\
To evaluate a transition dipole moment, consider each component in turn. For example, for the $z$-component,

$$
\mu_{z, \mathrm{fi}}=-e \int \psi_{\mathrm{f}}^{\star} z \psi_{\mathrm{i}} \mathrm{~d} \tau
$$

In spherical polar coordinates (see The chemist's toolkit 21 in Topic 7F) $z=r \cos \theta$. Then, according to Table 7F.1, $z=$ $(4 \pi / 3)^{1 / 2} r Y_{1,0}$. The wavefunctions for the initial and final states are atomic orbitals of the form $R_{n, l}(r) Y_{l, m_{l}}(\theta, \phi)$ (Topic 8A). With these substitutions the integral becomes\\
$\int \psi_{\mathrm{f}}^{\star} z \psi_{\mathrm{i}} \mathrm{d} \tau=$

$$
\int_{0}^{\infty} \int_{0}^{2 \pi} \int_{0}^{\pi} \overbrace{R_{n_{\mathrm{f}}, l_{\mathrm{f}}} Y_{l_{\mathrm{f}}, m_{l, f}}^{*}}^{\psi_{f}^{*}} \overbrace{\left(\frac{4 \pi}{3}\right)^{1 / 2} r Y_{1,0}} \overbrace{R_{n_{\mathrm{i}}, l_{\mathrm{i}}} Y_{l_{\mathrm{l}, ~}, m_{l, i}}}^{\psi_{\mathrm{i}}} \overbrace{r^{2} \mathrm{~d} r \sin \theta \mathrm{~d} \theta \mathrm{~d} \phi}^{\mathrm{Z}}
$$

\footnotetext{${ }^{1}$ See our Physical chemistry: Quanta, matter, and change (2014) for a detailed development of the form of eqn 8C.1.
}This multiple integral is the product of three factors, an integral over $r$ and two integrals (in blue) over the angles, so the factors on the right can be grouped as follows:

$$
\begin{aligned}
& \int \psi_{\mathrm{f}}^{*} z \psi_{\mathrm{i}} \mathrm{~d} \tau= \\
& \qquad\left(\frac{4 \pi}{3}\right)^{1 / 2} \int_{0}^{\infty} R_{n_{i}, l_{\mathrm{f}}} r^{3} R_{n_{i}, l_{l}} \mathrm{~d} r \int_{0}^{2 \pi} \int_{0}^{\pi} Y_{l_{t}, m_{l, Y}}^{*} Y_{1,0} Y_{l_{i, ~}, m_{l, l}} \sin \theta \mathrm{~d} \theta \mathrm{~d} \phi
\end{aligned}
$$

Step 3 Evaluate the angular integral\\
It follows from the properties of the spherical harmonics that the integral

$$
I=\int_{0}^{\pi} \int_{0}^{2 \pi} Y_{l_{\mathrm{s}}, m_{l, t}}^{*} Y_{l, m} Y_{l_{i}, m_{l, t}} \sin \theta \mathrm{~d} \theta \mathrm{~d} \phi
$$

is zero unless $l_{\mathrm{f}}=l_{\mathrm{i}} \pm l$ and $m_{l, \mathrm{f}}=m_{l, \mathrm{i}}+m$. Because in the present case $l=1$ and $m=0$, the angular integral, and hence the $z$-component of the transition dipole moment, is zero unless $\Delta l= \pm 1$ and $\Delta m_{l}=0$, which is a part of the set of selection rules. The same procedure, but considering the $x$ - and $y$-components, results in the complete set of rules:

\[
\begin{array}{lll}
\Delta l= \pm 1 & \Delta m_{l}=0, \pm 1 & \begin{array}{l}
\text { Selection rules for } \\
\text { hydrogenic atoms }
\end{array} \tag{8C.2}
\end{array}
\]

The principal quantum number $n$ can change by any amount consistent with the value of $\Delta l$ for the transition, because it does not relate directly to the angular momentum.

\subsection*{Brief illustration 8C. 1}
To identify the orbitals to which a 4 d electron may make radiative transitions, first identify the value of $l$ and then apply the selection rule for this quantum number. Because $l=$ 2 , the final orbital must have $l=1$ or 3 . Thus, an electron may make a transition from a 4 d orbital to any $n \mathrm{p}$ orbital (subject to $\Delta m_{l}=0, \pm 1$ ) and to any $n \mathrm{f}$ orbital (subject to the same rule). However, it cannot undergo a transition to any other orbital, such as an $n$ s or an $n$ d orbital.

The selection rules and the atomic energy levels jointly account for the structure of a Grotrian diagram (Fig. 8C.1), which summarizes the energies of the states and the transitions between them. In some versions, the thicknesses of the transition lines in the diagram denote their relative intensities in the spectrum.

\subsection*{8c. 2 The spectra of many-electron atoms}
The spectra of atoms rapidly become very complicated as the number of electrons increases, in part because their energy levels, their terms, are not given solely by the energies of the orbitals but depend on the interactions between the electrons.\\
\includegraphics[max width=\textwidth, center]{2025_02_27_d4b95d9718f0b9e2e6b9g-360}

Figure 8C. 1 A Grotrian diagram that summarizes the appearance and analysis of the spectrum of atomic hydrogen. The wavenumbers of some transitions (in $\mathrm{cm}^{-1}$ ) are indicated. The colours of the lines are for reference only: they are not the colours of the transitions.

\subsection*{(a) Singlet and triplet terms}
Consider the energy levels of a He atom, with its two electrons. The ground-state configuration is $1 s^{2}$, and an excited configuration is one in which an electron has been promoted into a different orbital to give, for instance, the configuration $1 s^{1} 2 s^{1}$. The two electrons need not be paired because they occupy different orbitals. According to Hund's maximum multiplicity rule (Topic 8B), the state of the atom with the spins parallel lies lower in energy than the state in which they are paired. Both states are permissible, correspond to different terms, and can contribute to the spectrum of the atom.

Parallel and antiparallel (paired) spins differ in their total spin angular momentum. In the paired case, the two spin momenta cancel, and there is zero net spin (as depicted in Fig. 8C.2(a)). Its state is the one denoted $\sigma_{-}$in the discussion of the Pauli principle (Topic 8B):


\begin{equation*}
\sigma_{-}(1,2)=\left(\frac{1}{2^{1 / 2}}\right)\{\alpha(1) \beta(2)-\beta(1) \alpha(2)\} \tag{8С.3a}
\end{equation*}


The angular momenta of two parallel spins add to give a nonzero total spin. As illustrated in Fig. 8C.2(b), there are three ways of achieving non-zero total spin. The three spin states are the symmetric combinations introduced in Topic 8B:


\begin{align*}
& \alpha(1) \alpha(2) \\
& \sigma_{+}(1,2)=\left(\frac{1}{2^{1 / 2}}\right)\{\alpha(1) \beta(2)+\beta(1) \alpha(2)\} \tag{8C.3b}
\end{align*}


$$
\beta(1) \beta(2)
$$

The state of the He atom in which the two electrons are paired and their spins are described by eqn 8C.3a gives rise to a singlet term. The alternative arrangement, in which the spins are parallel and are described by any of the three expressions in eqn 8C.3b, gives rise to a triplet term. The fact that the parallel\\
\includegraphics[max width=\textwidth, center]{2025_02_27_d4b95d9718f0b9e2e6b9g-361(1)}

Figure 8C. 2 (a) Electrons with paired spins have zero resultant spin angular momentum $(S=0)$. They can be represented by two vectors that lie at an indeterminate position on the cones shown here, but wherever one lies on its cone, the other points in the opposite direction; their resultant is zero. (b) When two electrons have parallel spins, they have a nonzero total spin angular momentum $(S=1)$. There are three ways of achieving this resultant, which are shown by these vector representations. The red vectors show the total spin angular momentum. Note that, whereas two paired spins are precisely antiparallel, two 'parallel' spins are not strictly parallel. The notation $S, M_{S}$ is explained later.\\
arrangement of spins in the triplet term of the $1 s^{1} 2 s^{1}$ configuration of the He atom lies lower in energy than the antiparallel arrangement, the singlet term, can now be expressed by saying that the triplet term of the $1 s^{1} 2 s^{1}$ configuration of He lies lower in energy than the singlet term. This is a general conclusion and applies to other atoms (and molecules):

For states arising from the same configuration, the triplet term generally lies lower than the singlet term.

The origin of the energy difference lies in the effect of spin correlation on the Coulombic interactions between electrons, as in the case of Hund's maximum multiplicity rule for groundstate configurations (Topic 8B): electrons with parallel spins tend to avoid each other. Because the Coulombic interaction between electrons in an atom is strong, the difference in energies between singlet and triplet terms of the same configuration can be large. The singlet and triplet terms of the configuration $1 \mathrm{~s}^{1} 2 \mathrm{~s}^{1}$ of He , for instance, differ by $6421 \mathrm{~cm}^{-1}$ (corresponding to 0.80 eV ).

The spectrum of atomic helium is more complicated than that of atomic hydrogen, but there are two simplifying features. One is that the only excited configurations to consider are of the form $1 \mathrm{~s}^{1} n l^{1}$; that is, only one electron is excited. Excitation of two electrons requires an energy greater than the ionization energy of the atom, so the $\mathrm{He}^{+}$ion is formed instead of the doubly excited atom. Second, and as seen later in this Topic, no radiative transitions take place between singlet and triplet terms because the relative orientation of the two electron spins\\
\includegraphics[max width=\textwidth, center]{2025_02_27_d4b95d9718f0b9e2e6b9g-361(2)}

Figure 8C. 3 Some of the transitions responsible for the spectrum of atomic helium. The labels give the wavelengths (in nanometres) of the transitions.\\
cannot change during a transition. Thus, there is a spectrum arising from transitions between singlet terms (including the ground state) and between triplet terms, but not between the two. Spectroscopically, helium behaves like two distinct species. The Grotrian diagram for helium in Fig. 8C. 3 shows the two sets of transitions.

\subsection*{(b) Spin-orbit coupling}
An electron has a magnetic moment that arises from its spin. Similarly, an electron with orbital angular momentum (that is, an electron in an orbital with $l>0$ ) is in effect a circulating current, and possesses a magnetic moment that arises from its orbital momentum. The interaction of the spin magnetic moment with the magnetic field arising from the orbital angular momentum is called spin-orbit coupling. The strength of the coupling, and its effect on the energy levels of the atom, depend on the relative orientations of the spin and orbital magnetic moments, and therefore on the relative orientations of the two angular momenta (Fig. 8C.4).\\
\includegraphics[max width=\textwidth, center]{2025_02_27_d4b95d9718f0b9e2e6b9g-361}

Figure 8C. 4 Spin-orbit coupling is a magnetic interaction between spin and orbital magnetic moments; the black arrows show the direction of the angular momentum and the green arrows show the direction of the associated magnetic moments When the angular momenta are parallel, as in (a), the magnetic moments are aligned unfavourably; when they are opposed, as in (b), the interaction is favourable. This magnetic coupling is the cause of the splitting of a term into levels.\\
\includegraphics[max width=\textwidth, center]{2025_02_27_d4b95d9718f0b9e2e6b9g-362}

Figure 8C. 5 The coupling of the spin and orbital angular momenta of a d electron $(l=2)$ gives two possible values of $j$ depending on the relative orientations of the spin and orbital angular momenta of the electron.

One way of expressing the dependence of the spin-orbit interaction on the relative orientation of the spin and orbital momenta is to say that it depends on the total angular momentum of the electron, the vector sum of its spin and orbital momenta. Thus, when the spin and orbital angular momenta are nearly parallel, the total angular momentum is high; when the two angular momenta are opposed, the total angular momentum is low.

The total angular momentum of an electron is described by the quantum numbers $j$ and $m_{j}$, with $j=l+\frac{1}{2}$ (when the orbital and spin angular momenta are in the same direction) or $j=l-\frac{1}{2}$ (when they are opposed; both cases are illustrated in Fig. 8C.5). The different values of $j$ that can arise for a given value of $l$ label the levels of a term. For $l=0$, the only permitted value is $j=\frac{1}{2}$ (the total angular momentum is the same as the spin angular momentum because there is no other source of\\
angular momentum in the atom). When $l=1, j$ may be either $\frac{3}{2}$ (the spin and orbital angular momenta are in the same sense) or $\frac{1}{2}$ (the spin and angular momenta are in opposite senses).

\subsection*{Brief illustration 8C. 2}
To identify the levels that may arise from the configurations (a) $d^{1}$ and (b) $s^{1}$, identify the value of $l$ and then the possible values of $j$. (a) For a d electron, $l=2$ and there are two levels in the configuration, one with $j=2+\frac{1}{2}=\frac{5}{2}$ and the other with $j=$ $2-\frac{1}{2}=\frac{3}{2}$. (b) For an s electron $l=0$, so only one level is possible, and $j=\frac{1}{2}$.

With a little work, it is possible to incorporate the effect of spin-orbit coupling on the energies of the levels.

\subsection*{How is that done? 8C. 2 \\
 Deriving an expression for the energy of spin-orbit interaction}
Classically, the energy of a magnetic moment $\mu$ in a magnetic field $\mathcal{B}$ is equal to their scalar product $-\mu \cdot \mathcal{B}$. Follow these steps to arrive at an expression for the spin-orbit interaction energy. The procedures for manipulating vectors are described in The chemist's toolkit 22.

\subsection*{Step 1 Write an expression for the energy of interaction}
If the magnetic field arises from the orbital angular momentum of the electron, it is proportional to $l$; if the magnetic moment $\mu$ is that of the electron spin, then it is proportional to $s$. It follows that the energy of interaction is proportional to the scalar product s.l:

\subsection*{The chemist's toolkit 22}
\subsection*{The manipulation of vectors}
In three dimensions, the vectors $\boldsymbol{u}$ (with components $u_{x}, u_{y}$, and $u_{z}$ ) and $\boldsymbol{v}$ (with components $v_{x}, v_{y}$, and $v_{z}$ ) have the general form:

$$
\boldsymbol{u}=u_{x} \boldsymbol{i}+u_{y} \boldsymbol{j}+u_{z} \boldsymbol{k} \quad \boldsymbol{v}=v_{x} \boldsymbol{i}+v_{y} \boldsymbol{j}+v_{z} \boldsymbol{k}
$$

where $\boldsymbol{i}, \boldsymbol{j}$, and $\boldsymbol{k}$ are unit vectors, vectors of magnitude 1 , pointing along the positive directions on the $x, y$, and $z$ axes. The operations of addition, subtraction, and multiplication are as follows:

\begin{enumerate}
  \item Addition:\\
$\boldsymbol{v}+\boldsymbol{u}=\left(v_{x}+u_{x}\right) \boldsymbol{i}+\left(v_{y}+u_{y}\right) \boldsymbol{j}+\left(v_{z}+u_{z}\right) \boldsymbol{k}$
  \item Subtraction:
\end{enumerate}

$$
\boldsymbol{v}-\boldsymbol{u}=\left(v_{x}-u_{x}\right) \boldsymbol{i}+\left(v_{y}-u_{y}\right) \boldsymbol{j}+\left(v_{z}-u_{z}\right) \boldsymbol{k}
$$

\begin{enumerate}
  \setcounter{enumi}{2}
  \item Multiplication:\\
(a) The scalar product, or dot product, of the two vectors $\boldsymbol{u}$ and $\boldsymbol{v}$ is
\end{enumerate}

$$
\boldsymbol{u} \cdot \boldsymbol{v}=u_{x} v_{x}+u_{y} v_{y}+u_{z} v_{z}
$$

nitude of the vector.

$$
\boldsymbol{u} \cdot \boldsymbol{u}=u_{x}^{2}+u_{y}^{2}+u_{z}^{2}=u^{2}
$$

(b) The vector product, or cross product, of two vectors is

$$
\begin{aligned}
\boldsymbol{u} \times \boldsymbol{v} & =\left|\begin{array}{ccc}
\boldsymbol{i} & \boldsymbol{j} & \boldsymbol{k} \\
u_{x} & u_{y} & u_{z} \\
v_{x} & v_{y} & v_{z}
\end{array}\right| \\
& =\left(u_{y} v_{z}-u_{z} v_{y}\right) \boldsymbol{i}-\left(u_{x} v_{z}-u_{z} v_{x}\right) \boldsymbol{j}+\left(u_{x} v_{y}-u_{y} v_{x}\right) \boldsymbol{k}
\end{aligned}
$$

(Determinants are discussed in The chemist's toolkit 23 in Topic 9D.) If the two vectors lie in the plane defined by the unit vectors $\boldsymbol{i}$ and $\boldsymbol{j}$, their vector product lies parallel to the unit vector $\boldsymbol{k}$.

\subsection*{Energy of interaction $=-\mu \cdot \mathcal{B} \propto \boldsymbol{s} \cdot \boldsymbol{l}$}
Step 2 Express the scalar product in terms of the magnitudes of the vectors\\
Note that the total angular momentum is the vector sum of the spin and orbital momenta: $\boldsymbol{j}=\boldsymbol{l}+\boldsymbol{s}$. The magnitude of the vector $\boldsymbol{j}$ is calculated by evaluating

$$
\overbrace{\boldsymbol{j} \cdot \boldsymbol{j}}^{j^{2}}=(\boldsymbol{l}+\boldsymbol{s}) \cdot(\boldsymbol{l}+\boldsymbol{s})=\overbrace{\boldsymbol{l} \cdot \boldsymbol{l}}^{\mathcal{l}^{2}}+\overbrace{\boldsymbol{s} \cdot \boldsymbol{s}}^{s^{2}}+2 \boldsymbol{s} \cdot \boldsymbol{l}
$$

so

$$
j^{2}=l^{2}+s^{2}+2 s \cdot l
$$

That is,

$$
\boldsymbol{s} \cdot \boldsymbol{l}=\frac{1}{2}\left\{j^{2}-l^{2}-s^{2}\right\}
$$

This equation is a classical result.\\
Step 3 Replace the classical magnitudes by their quantum mechanical versions

To derive the quantum mechanical version of this expression, replace all the quantities on the right with their quantummechanical values, which are of the form $j(j+1) \hbar^{2}$, etc (Topic 7F):

$$
\boldsymbol{s} \cdot \boldsymbol{l}=\frac{1}{2}\{j(j+1)-l(l+1)-s(s+1)\} \hbar^{2}
$$

Then, by inserting this expression into the formula for the energy of interaction $(E \propto \boldsymbol{s} \cdot \boldsymbol{l})$ and writing the constant of proportionality as $h c \tilde{A} / \hbar^{2}$, obtain an expression for the energy in terms of the quantum numbers and the spin-orbit coupling constant, $\tilde{A}$ (a wavenumber):

$$
\left.-\left\lvert\, E_{l, s, j}=\frac{1}{2} h c \tilde{A}\{j(j+1)-l(l+1)-s(s+1)\} \nmid \begin{array}{c}
\text { Spin-orbit } \\
\text { interaction energy }
\end{array}\right.\right)
$$

\subsection*{Brief illustration 8C. 3}
The unpaired electron in the ground state of an alkali metal atom has $l=0$, so $j=\frac{1}{2}$. Because the orbital angular momentum is zero in this state, the spin-orbit coupling energy is zero (as is confirmed by setting $j=s$ and $l=0$ in eqn 8C.4). When the electron is excited to an orbital with $l=1$, it has orbital angular momentum and can give rise to a magnetic field that interacts with its spin. In this configuration the electron can have $j=\frac{3}{2}$ or $j=\frac{1}{2}$, and the energies of these levels are

$$
\begin{aligned}
& E_{1,1 / 2,3 / 2}=\frac{1}{2} h c \tilde{A}\left\{\frac{3}{2} \times \frac{5}{2}-1 \times 2-\frac{1}{2} \times \frac{3}{2}\right\}=\frac{1}{2} h c \tilde{A} \\
& E_{1,1 / 2,1 / 2}=\frac{1}{2} h c \tilde{A}\left\{\frac{1}{2} \times \frac{3}{2}-1 \times 2-\frac{1}{2} \times \frac{3}{2}\right\}=-h c \tilde{A}
\end{aligned}
$$

\begin{center}
\includegraphics[max width=\textwidth]{2025_02_27_d4b95d9718f0b9e2e6b9g-363}
\end{center}

Figure 8C. 6 The levels of a $2 p^{1}$ configuration arising from spin-orbit coupling. Note that the low-j level lies below the high- $j$ level in energy. The number of states in a level with quantum number $j$ is $2 j+1$.

The corresponding energies are shown in Fig. 8C.6. Note that the barycentre (the 'centre of gravity') of the levels is unchanged, because there are four states of energy $\frac{1}{2} h c \tilde{A}$ and two of energy $-h c \tilde{A}$

The strength of the spin-orbit coupling depends on the nuclear charge. To understand why this is so, imagine riding on the orbiting electron and seeing a charged nucleus apparently orbiting around you (like the Sun rising and setting). As a result, you find yourself at the centre of a ring of current. The greater the nuclear charge, the greater is this current, and therefore the stronger is the magnetic field you detect. Because the spin magnetic moment of the electron interacts with this orbital magnetic field, it follows that the greater the nuclear charge, the stronger is the spin-orbit interaction. It turns out that the coupling increases sharply with atomic number (as $Z^{4}$ ) because not only is the current greater but the electron is drawn closer to the nucleus. Whereas the coupling is only weak in H (giving rise to shifts of energy levels of no more than about $0.4 \mathrm{~cm}^{-1}$ ), in heavy atoms like Pb it is very strong (giving shifts of the order of thousands of reciprocal centimetres).

Two spectral lines are observed when the $p$ electron of an electronically excited alkali metal atom undergoes a transition into a lower s orbital. One line is due to a transition starting in a $j=\frac{3}{2}$ level of the upper term and the other line is due to a transition starting in the $j=\frac{1}{2}$ level of the same term. The two lines are jointly an example of the fine structure of a spectrum, the structure due to spin-orbit coupling. Fine structure can be seen in the emission spectrum from sodium vapour excited by an electric discharge (for example, in one kind of street lighting). The yellow line at 589 nm (close to $17000 \mathrm{~cm}^{-1}$ ) is actually a doublet composed of one line at $589.76 \mathrm{~nm}\left(16956.2 \mathrm{~cm}^{-1}\right)$ and another at $589.16 \mathrm{~nm}\left(16973.4 \mathrm{~cm}^{-1}\right)$; the components of this doublet are the ' $D$ lines' of the spectrum (Fig. 8C.7). Therefore, in Na, the spin-orbit coupling affects the energies by about $17 \mathrm{~cm}^{-1}$.\\
\includegraphics[max width=\textwidth, center]{2025_02_27_d4b95d9718f0b9e2e6b9g-364(1)}

Figure 8C. 7 The energy-level diagram for the formation of the sodium $D$ lines. The splitting of the spectral lines (by $17 \mathrm{~cm}^{-1}$ ) reflects the splitting of the levels of the ${ }^{2} P$ term.

\subsection*{Example 8C. 1 Analysing a spectrum for the spin-orbit coupling constant}
The origin of the D lines in the spectrum of atomic sodium is shown in Fig. 8C.7. Calculate the spin-orbit coupling constant for the upper configuration of the Na atom.

Collect your thoughts It follows from Fig. 8C. 7 that the splitting of the lines is equal to the energy separation of the $j=\frac{3}{2}$ and $\frac{1}{2}$ levels of the excited configuration. You need to express this separation in terms of $\tilde{A}$ by using eqn 8C.4.

The solution The two levels are split by

$$
\Delta \tilde{v}=\left(E_{1, \frac{1}{2}, \frac{3}{2}}-E_{1, \frac{1}{2}, \frac{1}{2}}\right) / h c=\frac{1}{2} \tilde{A}\left\{\frac{3}{2}\left(\frac{3}{2}+1\right)-\frac{1}{2}\left(\frac{1}{2}+1\right)\right\}=\frac{3}{2} \tilde{A}
$$

The experimental value of $\Delta \tilde{v}$ is $17.2 \mathrm{~cm}^{-1}$; therefore

$$
\tilde{A}=\frac{2}{3} \times\left(17.2 \mathrm{~cm}^{-1}\right)=11.5 \mathrm{~cm}^{-1}
$$

Comment. The same calculation repeated for the atoms of other alkali metals gives Li: $0.23 \mathrm{~cm}^{-1}$, K: $38.5 \mathrm{~cm}^{-1}, \mathrm{Rb}: 158 \mathrm{~cm}^{-1}$, Cs: $370 \mathrm{~cm}^{-1}$. Note the increase of $\tilde{A}$ with atomic number (but more slowly than $Z^{4}$ for these many-electron atoms).

Self-test 8C. 1 The configuration... $4 \mathrm{p}^{6} 5 \mathrm{~d}^{1}$ of rubidium has two levels at $25700.56 \mathrm{~cm}^{-1}$ and $25703.52 \mathrm{~cm}^{-1}$ above the ground state. What is the spin-orbit coupling constant in this excited state?\\
\includegraphics[max width=\textwidth, center]{2025_02_27_d4b95d9718f0b9e2e6b9g-364}

\subsection*{(c) Term symbols}
The discussion so far has used expressions such as 'the $j=\frac{3}{2}$ level of a doublet term with $l=l^{\prime}$. A term symbol, which is a\\
symbol looking like ${ }^{2} \mathrm{P}_{3 / 2}$ or ${ }^{3} \mathrm{D}_{2}$, conveys this information, specifically the total spin, total orbital angular momentum, and total overall angular momentum, very succinctly.

A term symbol gives three pieces of information:

\begin{itemize}
  \item The letter ( P or D in the examples) indicates the total orbital angular momentum quantum number, $L$.
  \item The left superscript in the term symbol (the 2 in ${ }^{2} \mathrm{P}_{3 / 2}$ ) gives the multiplicity of the term.
  \item The right subscript on the term symbol (the $\frac{3}{2}$ in ${ }^{2} \mathrm{P}_{3 / 2}$ ) is the value of the total angular momentum quantum number, $J$, and labels the level of the term.
\end{itemize}

The meaning of these statements can be discussed in the light of the contributions to the energies summarized in Fig. 8C.8.

When several electrons are present, it is necessary to judge how their individual orbital angular momenta add together to augment or oppose each other. The total orbital angular momentum quantum number, $L$, gives the magnitude of the angular momentum through $\{L(L+1)\}^{1 / 2} \hbar$. It has $2 L+1$ orientations distinguished by the quantum number $M_{L}$, which can take the values $0, \pm 1, \ldots, \pm L$. Similar remarks apply to the total spin quantum number, $S$, and the quantum number $M_{S}$, and the total angular momentum quantum number, $J$, and the quantum number $M_{J}$.

The value of $L$ (a non-negative integer) is obtained by coupling the individual orbital angular momenta by using the Clebsch-Gordan series:


\begin{equation*}
L=l_{1}+l_{2}, l_{1}+l_{2}-1, \ldots,\left|l_{1}-l_{2}\right| \quad \text { Clebsch-Gordan series } \tag{8C.5}
\end{equation*}


The modulus signs are attached to $l_{1}-l_{2}$ to ensure that $L$ is nonnegative. The maximum value, $L=l_{1}+l_{2}$, is obtained when the two orbital angular momenta are in the same direction; the lowest value, $\left|l_{1}-l_{2}\right|$, is obtained when they are in opposite\\
\includegraphics[max width=\textwidth, center]{2025_02_27_d4b95d9718f0b9e2e6b9g-364(2)}

Figure 8 C .8 A summary of the types of interaction that are responsible for the various kinds of splitting of energy levels in atoms. For light (low $Z$ ) atoms, magnetic interactions are small, but in heavy (high Z) atoms they may dominate the electrostatic (charge-charge) interactions.\\
\includegraphics[max width=\textwidth, center]{2025_02_27_d4b95d9718f0b9e2e6b9g-365(2)}

Figure 8C.9 The total orbital angular momenta of a p electron and a d electron correspond to $L=3,2$, and 1 and reflect the different relative orientations of the two momenta.\\
directions. The intermediate values represent possible intermediate relative orientations of the two momenta (Fig. 8C.9). For two p electrons (for which $l_{1}=l_{2}=1$ ), $L=2,1,0$. The code for converting the value of $L$ into a letter is the same as for the $s, p, d, f, \ldots$ designation of orbitals, but uses uppercase Roman letters ${ }^{2}$ :

\begin{center}
\begin{tabular}{llllllll}
$L:$ & 0 & 1 & 2 & 3 & 4 & 5 & $6 \ldots$ \\
 & S & P & D & F & G & H & $\mathrm{I} \ldots$ \\
\end{tabular}
\end{center}

Thus, a $p^{2}$ configuration has $L=2,1,0$ and gives rise to $\mathrm{D}, \mathrm{P}$, and $S$ terms. The terms differ in energy on account of the different spatial distribution of the electrons and the consequent differences in repulsion between them.

A note on good practice Throughout this discussion of atomic spectroscopy, distinguish italic $S$, the total spin quantum number, from Roman S, the term label.

A closed shell has zero orbital angular momentum because all the individual orbital angular momenta sum to zero. Therefore, when working out term symbols, only the electrons of the unfilled shell need to be considered. In the case of a single electron outside a closed shell, the value of $L$ is the same as the value of $l$; so the configuration $[\mathrm{Ne}] 3 s^{1}$ has only an $S$ term.

Example 8C. 2 Deriving the total orbital angular momentum of a configuration\\
Find the terms that can arise from the configurations (a) $d^{2}$, (b) $\mathrm{p}^{3}$.

Collect your thoughts Use the Clebsch-Gordan series and begin by finding the minimum value of $L$ (so that you know where the series terminates). When there are more than two electrons to couple together, you need to use two series in

\footnotetext{${ }^{2}$ The convention of using lowercase letters to label orbitals and uppercase letters to label overall states applies throughout spectroscopy, not just to atoms.
}
succession: first to couple two electrons, and then to couple the third to each combined state, and so on.

The solution (a) The minimum value is $\left|l_{1}-l_{2}\right|=|2-2|=0$. Therefore,

$$
L=2+2,2+2-1, \ldots, 0=4,3,2,1,0
$$

corresponding to $\mathrm{G}, \mathrm{F}, \mathrm{D}, \mathrm{P}$, and S terms, respectively. (b) Coupling two p electrons gives a minimum value of $|1-1|=0$. Therefore,

$$
L^{\prime}=1+1,1+1-1, \ldots, 0=2,1,0
$$

Now couple $l_{3}=1$ with $L^{\prime}=2$, to give $L=3,2,1$; with $L^{\prime}=1$, to give $L=2,1,0$; and with $L^{\prime}=0$, to give $L=1$. The overall result is

$$
L=3,2,2,1,1,1,0
$$

giving one F , two D , three P , and one S term.\\
Self-test 8C. 2 Repeat the question for the configurations (a) $f^{1} d^{1}$ and (b) $d^{3}$.\\
\includegraphics[max width=\textwidth, center]{2025_02_27_d4b95d9718f0b9e2e6b9g-365}

When there are several electrons to be taken into account, their total spin angular momentum quantum number, $S$ (a non-negative integer or half-integer), must be assessed. Once again the Clebsch-Gordan series is used, but now in the form


\begin{equation*}
S=s_{1}+s_{2}, s_{1}+s_{2}-1, \ldots,\left|s_{1}-s_{2}\right| \tag{8C.6}
\end{equation*}


to decide on the value of $S$, noting that each electron has $s=\frac{1}{2}$. For two electrons the possible values of $S$ are 1 and 0 (Fig. 8C.10). If there are three electrons, the total spin angular momentum is obtained by coupling the third spin to each of the values of $S$ for the first two spins, which results in $S=\frac{3}{2}$ and $\frac{1}{2}$.

The multiplicity of a term is the value of $2 S+1$. When $S=0$ (as for a closed shell, like $1 s^{2}$ ) the electrons are all paired and there is no net spin: this arrangement gives a singlet term, ${ }^{1}$.\\
\includegraphics[max width=\textwidth, center]{2025_02_27_d4b95d9718f0b9e2e6b9g-365(1)}

Figure 8C. 10 For two electrons (each of which has $s=\frac{1}{2}$ ), only two total spin states are permitted $(S=0,1)$. (a) The state with $S=0$ can have only one value of $M_{S}\left(M_{S}=0\right)$ and gives rise to a singlet term; (b) the state with $S=1$ can have any of three values of $M_{s}(+1,0,-1)$ and gives rise to a triplet term. The vector representations of the $S=0$ and 1 states are shown in Fig. 8C.2.

A lone electron has $S=s=\frac{1}{2}$, so a configuration such as $[\mathrm{Ne}] 3 \mathrm{~s}^{1}$ can give rise to a doublet term, ${ }^{2}$ S. Likewise, the configuration [ Ne ] $3 \mathrm{p}{ }^{1}$ is a doublet, ${ }^{2} \mathrm{P}$. When there are two unpaired (parallel spin) electrons $S=1$, so $2 S+1=3$, giving a triplet term, such as ${ }^{3} \mathrm{D}$. The relative energies of singlets and triplets are discussed earlier in the Topic, where it is seen that their energies differ on account of spin correlation.

As already explained, the quantum number $j$ gives the relative orientation of the spin and orbital angular momenta of a single electron. The total angular momentum quantum number, $J$ (a non-negative integer or half-integer), does the same for several electrons. If there is a single electron outside a closed shell, $J=j$, with $j$ either $l+\frac{1}{2}$ or $\left|l-\frac{1}{2}\right|$. The $[\mathrm{Ne}] 3 \mathrm{~s}^{1}$ configuration has $j=\frac{1}{2}$ (because $l=0$ and $s=\frac{1}{2}$ ), so the ${ }^{2}$ S term has a single level, denoted ${ }^{2} S_{1 / 2}$. The $[\mathrm{Ne}] 3 \mathrm{p}^{1}$ configuration has $l=1$; therefore $j=\frac{3}{2}$ and $\frac{1}{2}$; the ${ }^{2} \mathrm{P}$ term therefore has two levels, ${ }^{2} \mathrm{P}_{3 / 2}$ and ${ }^{2} \mathrm{P}_{1 / 2}$. These levels lie at different energies on account of the spin-orbit interaction.

If there are several electrons outside a closed shell it is necessary to consider the coupling of all the spins and all the orbital angular momenta. This complicated problem can be simplified when the spin-orbit coupling is weak (for atoms of low atomic number), by using the Russell-Saunders coupling scheme. This scheme is based on the view that, if spin-orbit coupling is weak, then it is effective only when all the orbital momenta are operating cooperatively. That is, all the orbital angular momenta of the electrons couple to give a total $L$, and all the spins are similarly coupled to give a total S. Only at this stage do the two kinds of momenta couple through the spinorbit interaction to give a total $J$. The permitted values of $J$ are given by the Clebsch-Gordan series


\begin{equation*}
J=L+S, L+S-1, \ldots,|L-S| \tag{8С.7}
\end{equation*}


For example, in the case of the ${ }^{3} \mathrm{D}$ term of the configuration [Ne] $2 p^{1} 3 p^{1}$, the permitted values of $J$ are $3,2,1$ (because ${ }^{3} \mathrm{D}$ has $L=2$ and $S=1$ ), so the term has three levels, ${ }^{3} \mathrm{D}_{3},{ }^{3} \mathrm{D}_{2}$, and ${ }^{3} \mathrm{D}_{1}$.

When $L \geq S$, the multiplicity is equal to the number of levels. For example, a ${ }^{2} \mathrm{P}$ term $\left(L=1>S=\frac{1}{2}\right)$ has the two levels ${ }^{2} \mathrm{P}_{3 / 2}$ and ${ }^{2} \mathrm{P}_{1 / 2}$, and ${ }^{3} \mathrm{D}(L=2>S=1)$ has the three levels ${ }^{3} \mathrm{D}_{3},{ }^{3} \mathrm{D}_{2}$, and ${ }^{3} \mathrm{D}_{1}$. However, this is not the case when $L<S$ : the term ${ }^{2} \mathrm{~S}(L=0$ $<S=\frac{1}{2}$ ), for example, has only the one level ${ }^{2} S_{1 / 2}$.

\subsection*{Example 8C. 3}
Deriving term symbols\\
Write the term symbols arising from the ground-state configurations of (a) Na and (b) F, and (c) the excited configuration $1 s^{2} 2 s^{2} 2 p^{1} 3 p^{1}$ of C.

Collect your thoughts Begin by writing the configurations, but ignore inner closed shells. Then couple the orbital momenta to find $L$ and the spins to find $S$. Next, couple $L$ and $S$ to find $J$. Finally, express the term as ${ }^{2 S+1}\{L\}$, where $\{L\}$ is the appropriate\\
letter. For F, for which the valence configuration is $2 p^{5}$, treat the single gap in the closed-shell $2 \mathrm{p}^{6}$ configuration as a single spin- $\frac{1}{2}$ particle.

The solution (a) For Na , the configuration is $[\mathrm{Ne}] 3 \mathrm{~s}^{1}$, and consider only the single 3 s electron. Because $L=l=0$ and $S=s=\frac{1}{2}$, the only possible value is $J=\frac{1}{2}$. Hence the term symbol is ${ }^{2} \mathrm{~S}_{1 / 2}$. (b) For F, the configuration is $[\mathrm{He}] 2 s^{2} 2 p^{5}$, which can be treated as $[\mathrm{Ne}] 2 \mathrm{p}^{-1}$ (where the notation $2 \mathrm{p}^{-1}$ signifies the absence of a 2 p electron). Hence $L=l=1$, and $S=s=\frac{1}{2}$. Two values of $J$ are possible: $J=\frac{3}{2}, \frac{1}{2}$. Hence, the term symbols for the two levels are ${ }^{2} \mathrm{P}_{3 / 2}$ and ${ }^{2} \mathrm{P}_{1 / 2}$. (c) This is a two-electron problem, and $l_{1}=$ $l_{2}=1, s_{1}=s_{2}=\frac{1}{2}$. It follows that $L=2,1,0$ and $S=1,0$. The terms are therefore ${ }^{3} \mathrm{D}$ and ${ }^{1} \mathrm{D},{ }^{3} \mathrm{P}$ and ${ }^{1} \mathrm{P}$, and ${ }^{3} \mathrm{~S}$ and ${ }^{1} \mathrm{~S}$. For ${ }^{3} \mathrm{D}, L=2$ and $S=1$; hence $J=3,2,1$ and the levels are ${ }^{3} \mathrm{D}_{3},{ }^{3} \mathrm{D}_{2}$, and ${ }^{3} \mathrm{D}_{1}$. For ${ }^{1} \mathrm{D}, L=2$ and $S=0$, so the single level is ${ }^{1} \mathrm{D}_{2}$. The triplet of levels of ${ }^{3} \mathrm{P}$ is ${ }^{3} \mathrm{P}_{2},{ }^{3} \mathrm{P}_{1}$, and ${ }^{3} \mathrm{P}_{0}$, and the singlet is ${ }^{1} \mathrm{P}_{1}$. For the ${ }^{3} \mathrm{~S}$ term there is only one level, ${ }^{3} \mathrm{~S}_{1}$ (because $J=1$ only), and the singlet term is ${ }^{1} S_{0}$.

Comment. Fewer terms arise from a configuration like $\ldots 2 \mathrm{p}^{2}$ or $\ldots 3 \mathrm{p}^{2}$ than from a configuration like. $.2 \mathrm{p}^{1} 3 \mathrm{p}^{1}$ because the Pauli exclusion principle forbids parallel arrangements of spins when two electrons occupy the same orbital. The analysis of the terms arising in such cases requires more detail than given here.

Self-test 8C. 3 Identify the terms arising from the configurations (a) $2 s^{1} 2 p^{1}$, (b) $2 p^{1} 3 \mathrm{~d}^{1}$.

Russell-Saunders coupling fails when the spin-orbit coupling is large (in heavy atoms, those with high $Z$ ). In that case, the individual spin and orbital momenta of the electrons are coupled into individual $j$ values; then these momenta are combined into a grand total, J, given by a Clebsch-Gordan series. This scheme is called $\mathbf{j} \mathbf{j}$-coupling. For example, in a $p^{2}$ configuration, the individual values of $j$ are $\frac{3}{2}$ and $\frac{1}{2}$ for each electron. If the spin and the orbital angular momentum of each electron are coupled together strongly, it is best to consider each electron as a particle with angular momentum $j=\frac{3}{2}$ or $\frac{1}{2}$. These individual total momenta then couple as follows:

\begin{center}
\begin{tabular}{lll}
\hline
$\boldsymbol{j}_{1}$ & $\boldsymbol{j}_{2}$ & $\boldsymbol{J}$ \\
\hline
$\frac{3}{2}$ & $\frac{3}{2}$ & $3,2,1,0$ \\
$\frac{3}{2}$ & $\frac{1}{2}$ & 2,1 \\
$\frac{1}{2}$ & $\frac{3}{2}$ & 2,1 \\
$\frac{1}{2}$ & $\frac{1}{2}$ & 1,0 \\
\hline
\end{tabular}
\end{center}

For heavy atoms, in which $j j$-coupling is appropriate, it is best to discuss their energies by using these quantum numbers.

Although jj-coupling should be used for assessing the energies of heavy atoms, the term symbols derived from RussellSaunders coupling can still be used as labels. To see why this\\
\includegraphics[max width=\textwidth, center]{2025_02_27_d4b95d9718f0b9e2e6b9g-367}

Figure 8C. 11 The correlation diagram for some of the states of a two-electron system. All atoms lie between the two extremes, but the heavier the atom, the closer it lies to the pure $j j$-coupling case.\\
procedure is valid, it is useful to examine how the energies of the atomic states change as the spin-orbit coupling increases in strength. Such a correlation diagram is shown in Fig. 8C.11. It shows that there is a correspondence between the low spinorbit coupling (Russell-Saunders coupling) and high spinorbit coupling (jj-coupling) schemes, so the labels derived by using the Russell-Saunders scheme can be used to label the states of the jj-coupling scheme.

\subsection*{(d) Hund's rules}
As already remarked, the terms arising from a given configuration differ in energy because they represent different relative orientations of the angular momenta of the electrons and therefore different spatial distributions. The terms arising from the ground-state configuration of an atom (and less reliably from other configurations) can be put into the order of increasing energy by using Hund's rules, which summarize the preceding discussion:

\begin{enumerate}
  \item For a given configuration, the term of greatest multiplicity lies lowest in energy.
\end{enumerate}

As discussed in Topic 8B, this rule is a consequence of spin correlation, the quantum-mechanical tendency of electrons with parallel spins to stay apart from one another.\\
2. For a given multiplicity, the term with the highest value of $L$ lies lowest in energy.

This rule can be explained classically by noting that two electrons have a high orbital angular momentum if they circulate in the same direction, in which case they can stay apart. If they circulate in opposite directions, they meet. Thus, a D term is expected to lie lower in energy than an $S$ term of the same multiplicity.\\
3. For atoms with less than half-filled shells, the level with the lowest value of $J$ lies lowest in energy; for more than half-filled shells, the highest value of $J$ lies lowest.

This rule arises from considerations of spin-orbit coupling. Thus, for a state of low $J$, the orbital and spin angular momenta lie in opposite directions, and so too do the corresponding magnetic moments. In classical terms the magnetic moments are then antiparallel, with the N pole of one close to the S pole of the other, which is a low-energy arrangement.

\subsection*{(e) Selection rules}
Any state of the atom, and any spectral transition, can be specified by using term symbols. For example, the transitions giving rise to the yellow sodium doublet (which are shown in Fig. 8C.7) are

$$
3 \mathrm{p}^{12} \mathrm{P}_{3 / 2} \rightarrow 3 \mathrm{~s}^{12} \mathrm{~S}_{1 / 2} \quad 3 \mathrm{p}^{12} \mathrm{P}_{1 / 2} \rightarrow 3 \mathrm{~s}^{12} \mathrm{~S}_{1 / 2}
$$

By convention, the upper term precedes the lower. The corresponding absorptions are therefore denoted ${ }^{2} \mathrm{P}_{3 / 2} \leftarrow{ }^{2} \mathrm{~S}_{1 / 2}$ and ${ }^{2} \mathrm{P}_{1 / 2} \leftarrow{ }^{2} \mathrm{~S}_{1 / 2}$. (The configurations have been omitted.)

As seen in Section 8C.1, selection rules arise from the conservation of angular momentum during a transition and from the fact that a photon has a spin of 1 . They can therefore be expressed in terms of the term symbols, because the latter carry information about angular momentum. A detailed analysis leads to the following rules:


\begin{align*}
& \Delta S=0 \\
& \Delta L=0, \pm 1, \Delta l= \pm 1 \\
& \Delta J=0, \pm 1 \text { but } J=0 \longleftrightarrow J=0 \quad \text { Selection rules for atoms } \tag{8С.8}
\end{align*}


where the symbol $\longleftrightarrow$ denotes a forbidden transition. The rule about $\Delta S$ (no change of overall spin) stems from the fact that electromagnetic radiation does not affect the spin directly. The rules about $\Delta L$ and $\Delta l$ express the fact that the orbital angular momentum of an individual electron must change (so $\Delta l= \pm 1$ ), but whether or not this results in an overall change of orbital momentum depends on the coupling.

The selection rules given above apply when RussellSaunders coupling is valid (in light atoms, those of low Z). If labelling the terms of heavy atoms with symbols like ${ }^{3} \mathrm{D}$, then the selection rules progressively fail as the atomic number increases because the quantum numbers $S$ and $L$ become ill defined as $j j$-coupling becomes more appropriate. As explained above, Russell-Saunders term symbols are only a convenient way of labelling the terms of heavy atoms: they do not bear any direct relation to the actual angular momenta of the electrons in a heavy atom. For this reason, transitions between singlet and triplet states (for which $\Delta S= \pm 1$ ), while forbidden in light atoms, are allowed in heavy atoms.

\subsection*{Checklist of concepts}
$\square$ 1. Two electrons with paired spins in a configuration give rise to a singlet term; if their spins are parallel, they give rise to a triplet term.2. The orbital and spin angular momenta interact magnetically.3. Spin-orbit coupling results in the levels of a term having different energies.4. Fine structure in a spectrum is due to transitions to different levels of a term.5. A term symbol specifies the angular momentum states of an atom.6. Angular momenta are combined into a resultant by using the Clebsch-Gordan series.\\
7. The multiplicity of a term is the value of $2 S+1$.8. The total angular momentum in light atoms is obtained on the basis of Russell-Saunders coupling; in heavy atoms, $\mathbf{j} \mathbf{j}$-coupling is used.9. The term with the maximum multiplicity lies lowest in energy.\\
$\square$ 10. For a given multiplicity, the term with the highest value of $L$ lies lowest in energy.11. For atoms with less than half-filled shells, the level with the lowest value of $J$ lies lowest in energy; for more than half-filled shells, the highest value of $J$ lies lowest.\\
$\square$ 12. Selection rules for light atoms include the fact that changes of total spin do not occur.

\subsection*{Checklist of equations}
\begin{center}
\begin{tabular}{|c|c|c|c|}
\hline
Property & Equation & Comment & Equation number \\
\hline
Spin-orbit interaction energy & $E_{l, s, j}=\frac{1}{2} h c \tilde{A}\{j(j+1)-l(l+1)-s(s+1)\}$ &  & 8C. 4 \\
\hline
Clebsch-Gordan series & $J=j_{1}+j_{2}, j_{1}+j_{2}-1, \ldots,\left|j_{1}-j_{2}\right|$ & $J, j$ denote any kind of angular momentum & 8C. 5 \\
\hline
Selection rules & \( \begin{aligned} & \Delta S=0 \\ & \Delta L=0, \pm 1, \Delta l= \pm 1, \\ & \Delta J=0, \pm 1, \text { but } J=0 \longleftrightarrow J=0 \end{aligned} \) & Light atoms & 8C. 8 \\
\hline
\end{tabular}
\end{center}

\chapter{FOCUS 9 Molecular structure}
The concepts developed in Focus 8, particularly those of orbitals, can be extended to a description of the electronic structures of molecules. There are two principal quantum mechanical theories of molecular electronic structure: 'valence-bond theory' is centred on the concept of the shared electron pair; 'molecular orbital theory' treats electrons as being distributed over all the nuclei in a molecule.

\subsection*{Prologue The Born-Oppenheimer approximation}
The starting point for the theories discussed here and the interpretation of spectroscopic results (Focus 11) is the 'BornOppenheimer approximation', which separates the relative motions of nuclei and electrons in a molecule.

\subsection*{9A Valence-bond theory}
The key concept of this Topic is the wavefunction for a shared electron pair, which is then used to account for the structures of a wide variety of molecules. The theory introduces the concepts of $\sigma$ and $\pi$ bonds, promotion, and hybridization, which are used widely in chemistry.\\
9A. 1 Diatomic molecules; 9A. 2 Resonance; 9A. 3 Polyatomic molecules

\subsection*{9B Molecular orbital theory: the hydrogen molecule-ion}
In molecular orbital theory the concept of an atomic orbital is extended to that of a 'molecular orbital', which is a wave-\\
function that spreads over all the atoms in a molecule. This Topic focuses on the hydrogen molecule-ion, setting the scene for the application of the theory to more complicated molecules. 9B. 1 Linear combinations of atomic orbitals; 9B. 2 Orbital notation

\subsection*{9C Molecular orbital theory: homonuclear diatomic molecules}
The principles established for the hydrogen molecule-ion are extended to other homonuclear diatomic molecules and ions. The principal differences are that all the valence-shell atomic orbitals must be included and that they give rise to a more varied collection of molecular orbitals. The building-up principle for atoms is extended to the occupation of molecular orbitals and used to predict the electronic configurations of molecules and ions.\\
9C. 1 Electron configurations; 9C. 2 Photoelectron spectroscopy

\subsection*{9D Molecular orbital theory: heteronuclear diatomic molecules}
The molecular orbital theory of heteronuclear diatomic molecules introduces the possibility that the atomic orbitals on the two atoms contribute unequally to the molecular orbital. As a result, the molecule is polar. The polarity can be expressed in terms of the concept of electronegativity. This Topic shows how quantum mechanics is used to calculate the form of a molecular orbital arising from the overlap of different atomic orbitals and its energy.\\
9D. 1 Polar bonds and electronegativity; 9D. 2 The variation principle

\subsection*{9E Molecular orbital theory: polyatomic molecules}
Most molecules are polyatomic, so it is important to be able to account for their electronic structure. An early approach to the electronic structure of planar conjugated polyenes is the 'Hückel method', which uses severe approximations but sets the scene for more sophisticated procedures. The latter have given rise to the huge and vibrant field of computational theoretical chemistry in which elaborate computations are used to predict molecular properties. This Topic describes briefly how those calculations are formulated and displayed.

\subsection*{Web resources What is an application of this material?}
The concepts introduced in this chapter pervade the whole of chemistry and are encountered throughout the text. Two biochemical aspects are discussed here. In Impact 14 simple concepts are used to account for the reactivity of small molecules that occur in organisms. Impact 15 provides a glimpse of the contribution of computational chemistry to the explanation of the thermodynamic and spectroscopic properties of several biologically significant molecules.

\subsection*{PROLOGUE The Born-Oppenheimer approximation}
All theories of molecular structure make the same simplification at the outset. Whereas the Schrödinger equation for a hydrogen atom can be solved exactly, an exact solution is not possible for any molecule because even the simplest molecule consists of three particles (two nuclei and one electron). Therefore, it is common to adopt the Born-Oppenheimer approximation in which it is supposed that the nuclei, being so much heavier than an electron, move relatively slowly and may be treated as stationary while the electrons move in their field. That is, the nuclei are assumed to be fixed at arbitrary locations, and the Schrödinger equation is then solved for the wavefunction of the electrons alone.

To use the Born-Oppenheimer approximation for a diatomic molecule, the nuclear separation is set at a chosen value, the Schrödinger equation for the electrons is then solved and the energy calculated. Then a different separation is selected, the calculation repeated, and so on for other values of the separation. In this way the variation of the energy of the molecule with bond length is explored, and a molecular potential energy curve is obtained (see the illustration). It is called a potential energy curve because the kinetic energy of the stationary nuclei is zero. Once the curve has been calculated or determined experimentally (by using the spectroscopic techniques described in Focus 11), it is possible to identify the\\
\includegraphics[max width=\textwidth, center]{2025_02_27_d4b95d9718f0b9e2e6b9g-375}

A molecular potential energy curve. The equilibrium bond length corresponds to the energy minimum.\\
equilibrium bond length, $R_{\mathrm{e}}$, the internuclear separation at the minimum of the curve, and the bond dissociation energy, $h c \tilde{D}_{0}$, which is closely related to the depth, $h c \tilde{D}_{e}$, of the minimum below the energy of the infinitely widely separated and stationary atoms. When more than one molecular parameter is changed in a polyatomic molecule, such as its various bond lengths and angles, a potential energy surface is obtained. The overall equilibrium shape of the molecule corresponds to the global minimum of the surface.

\section*{TOPIC 9A Valence-bond theory}
\subsection*{Why do you need to know this material?}
The language introduced by valence-bond theory is used throughout chemistry, especially in the description of the properties and reactions of organic compounds.

\subsection*{What is the key idea?}
A bond forms when an electron in an atomic orbital on one atom pairs its spin with that of an electron in an atomic orbital on another atom.

\subsection*{What do you need to know already?}
You need to know about atomic orbitals (Topic 8A) and the concepts of normalization and orthogonality (Topic 7C). This Topic also makes use of the Pauli principle (Topic 8B).

Valence-bond theory (VB theory) begins by considering the chemical bond in molecular hydrogen, $\mathrm{H}_{2}$. The basic concepts are then applied to all diatomic and polyatomic molecules and ions.

\subsection*{9A. 1 Diatomic molecules}
The spatial wavefunction for an electron on each of two widely separated H atoms is


\begin{equation*}
\Psi(1,2)=\psi_{\mathrm{Hls}_{\mathrm{A}}}\left(\boldsymbol{r}_{1}\right) \psi_{\mathrm{Hls}_{\mathrm{B}}}\left(\boldsymbol{r}_{2}\right) \tag{9A.1}
\end{equation*}


if electron 1 is in the H1s atomic orbital on atom A and electron 2 is in the H1s atomic orbital on atom B. For simplicity, this wavefunction will be written $\Psi(1,2)=\psi_{A}(1) \psi_{B}(2)$. When the atoms are close together, it is not possible to know whether it is electron 1 or electron 2 that is on A. An equally valid description is therefore $\Psi(1,2)=\psi_{\mathrm{A}}(2) \psi_{\mathrm{B}}(1)$, in which electron 2 is on A and electron 1 is on B . When two outcomes are equally probable in quantum mechanics, the true state of the system is described as a superposition of the wavefunctions for each possibility (Topic 7C), so a better description of the molecule than either wavefunction alone is one of the (unnormalized) linear combinations $\Psi(1,2)=$ $\psi_{A}(1) \psi_{B}(2) \pm \psi_{A}(2) \psi_{B}(1)$. The combination with lower energy turns out to be the one with a + sign, so the valence-bond wavefunction of the electrons in an $\mathrm{H}_{2}$ molecule is

\[
\Psi(1,2)=\psi_{\mathrm{A}}(1) \psi_{\mathrm{B}}(2)+\psi_{\mathrm{A}}(2) \psi_{\mathrm{B}}(1) \quad \begin{align*}
& \text { A valence-bond }  \tag{9A.2}\\
& \text { wavefunction }
\end{align*}
\]

The reason why this linear combination has a lower energy than either the separate atoms or the linear combination with a negative sign can be traced to the constructive interference between the wave patterns represented by the terms $\psi_{A}(1) \psi_{B}(2)$ and $\psi_{\mathrm{A}}(2) \psi_{\mathrm{B}}(1)$, and the resulting enhancement of the probability density of the electrons in the internuclear region (Fig. 9A.1).

\subsection*{Brief illustration 9A. 1}
The wavefunction in eqn 9A. 2 might look abstract, but in fact it can be expressed in terms of simple exponential functions. Thus, if the wavefunction for an H1s orbital $(Z=1)$ given in Topic 8 A is used, then, with the distances $r$ measured from their respective nuclei,

$$
\begin{aligned}
& \Psi(1,2)=\overbrace{\frac{1}{\left(\pi a_{0}^{3}\right)^{1 / 2}} \mathrm{e}^{-r_{\mathrm{A}} / a_{0}}}^{\psi_{\mathrm{A}}(1)} \times \overbrace{\frac{1}{\left(\pi a_{0}^{3}\right)^{1 / 2}} \mathrm{e}^{-r_{r_{2}} / a_{0}}}^{\psi_{B}(2)} \\
& +\overbrace{\frac{1}{\left(\pi a_{0}^{3}\right)^{1 / 2}} \mathrm{e}^{-r_{R_{2} / 2} / a_{0}}}^{\psi_{A}(2)} \cdot \overbrace{\frac{1}{\left(\pi a_{0}^{3}\right)^{1 / 2}} \mathrm{e}^{-r_{B_{1} / a_{0}}}}^{\psi_{B}(1)} \\
& =\frac{1}{\pi a_{0}^{3}}\left\{\mathrm{e}^{-\left(r_{\mathrm{r}_{11}}+r_{\mathrm{r}_{2}}\right) / a_{0}}+\mathrm{e}^{-\left(r_{\mathrm{r}_{2}}+r_{\mathrm{B}}\right) / a_{0}}\right\}
\end{aligned}
$$

where $r_{\mathrm{Al}}$ is the distance of electron 1 from nucleus A , etc.\\
\includegraphics[max width=\textwidth, center]{2025_02_27_d4b95d9718f0b9e2e6b9g-376}

Figure 9A. 1 It is very difficult to represent valence-bond wavefunctions because they refer to two electrons simultaneously. However, this illustration is an attempt. The atomic orbital for electron 1 is represented by the purple shading, and that of electron 2 is represented by the green shading. The left illustration represents $\psi_{A}(1) \psi_{B}(2)$ and the right illustration represents the contribution $\psi_{A}(2) \psi_{B}(1)$. When the two contributions are superimposed, there is interference between the purple contributions and between the green contributions, resulting in an enhanced (two-electron) density in the internuclear region.

The electron distribution described by the wavefunction in eqn 9A. 2 is called a $\sigma$ bond. A $\sigma$ bond has cylindrical symmetry around the internuclear axis, and is so called because, when viewed along the internuclear axis, it resembles a pair of electrons in an s orbital (and $\sigma$ is the Greek equivalent of s).

A chemist's picture of a covalent bond is one in which the spins of two electrons pair as the atomic orbitals overlap. It can be shown that the origin of the role of spin is that the wavefunction in eqn 9A. 2 can be formed only by two spin-paired electrons.

How is that done? 9A. 1 Establishing the origin of electron pairs in VB theory

The Pauli principle requires the overall wavefunction of two electrons, the wavefunction including spin, to change sign when the labels of the electrons are interchanged (Topic 8B). The overall VB wavefunction for two electrons is

$$
\Psi(1,2)=\left\{\psi_{A}(1) \psi_{B}(2)+\psi_{A}(2) \psi_{B}(1)\right\} \sigma(1,2)
$$

where $\sigma$ represents the spin component of the wavefunction. When the labels 1 and 2 are interchanged, this wavefunction becomes

$$
\begin{aligned}
\Psi(2,1) & =\left\{\psi_{A}(2) \psi_{\mathrm{B}}(1)+\psi_{\mathrm{A}}(1) \psi_{\mathrm{B}}(2)\right\} \sigma(2,1) \\
& =\left\{\psi_{A}(1) \psi_{\mathrm{B}}(2)+\psi_{A}(2) \psi_{\mathrm{B}}(1)\right\} \sigma(2,1)
\end{aligned}
$$

The Pauli principle requires that $\Psi(2,1)=-\Psi(1,2)$, which is satisfied only if $\sigma(2,1)=-\sigma(1,2)$. The combination of two spins that has this property is

$$
\sigma_{-}(1,2)=\frac{1}{2^{1 / 2}}\{\alpha(1) \beta(2)-\beta(1) \alpha(2)\}
$$

which corresponds to paired electron spins (Topic 8B). Therefore, the state of lower energy (and hence the formation of a chemical bond) is achieved if the electron spins are paired. Spin pairing is not an end in itself: it is a means of achieving a wavefunction, and the probability distribution it implies, that corresponds to a low energy.

The VB description of $\mathrm{H}_{2}$ can be applied to other homonuclear diatomic molecules. The starting point for the discussion of $\mathrm{N}_{2}$, for instance, is the valence electron configuration of each atom, which is $2 \mathrm{~s}^{2} 2 \mathrm{p}_{x}^{1} 2 \mathrm{p}_{y}^{1} 2 \mathrm{p}_{z}^{1}$. It is conventional to take the $z$-axis to be the internuclear axis in a linear molecule, so each atom is imagined as having a $2 \mathrm{p}_{z}$ orbital pointing towards a $2 \mathrm{p}_{z}$ orbital on the other atom (Fig. 9A.2), with the $2 \mathrm{p}_{x}$ and $2 \mathrm{p}_{y}$ orbitals perpendicular to the axis. A $\sigma$ bond is then formed by spin pairing between the two electrons in the two $2 \mathrm{p}_{z}$ orbitals. Its spatial wavefunction is given by eqn 9A.2, but now $\psi_{\mathrm{A}}$ and $\psi_{\mathrm{B}}$ stand for the two $2 \mathrm{p}_{z}$ orbitals.

The remaining N 2 p orbitals ( $2 \mathrm{p}_{x}$ and $2 \mathrm{p}_{y}$ ) cannot merge to give $\sigma$ bonds as they do not have cylindrical symmetry around the internuclear axis. Instead, they merge to form two ' $\pi$ bonds'. A $\pi$ bond arises from the spin pairing of electrons\\
\includegraphics[max width=\textwidth, center]{2025_02_27_d4b95d9718f0b9e2e6b9g-377}

Figure 9A. 2 The orbital overlap and spin pairing between electrons in two collinear $p$ orbitals that results in the formation of a $\sigma$ bond.\\
\includegraphics[max width=\textwidth, center]{2025_02_27_d4b95d9718f0b9e2e6b9g-377(1)}

Figure 9A. 3 A $\pi$ bond results from orbital overlap and spin pairing between electrons in $p$ orbitals with their axes perpendicular to the internuclear axis. The bond has two lobes of electron density separated by a nodal plane.\\
in two p orbitals that approach side-by-side (Fig. 9A.3). It is so called because, viewed along the internuclear axis, a $\pi$ bond resembles a pair of electrons in a p orbital (and $\pi$ is the Greek equivalent of $p$ ). ${ }^{1}$

There are two $\pi$ bonds in $\mathrm{N}_{2}$, one formed by spin pairing in two neighbouring $2 \mathrm{p}_{x}$ orbitals and the other by spin pairing in two neighbouring $2 \mathrm{p}_{y}$ orbitals. The overall bonding pattern in $\mathrm{N}_{2}$ is therefore a $\sigma$ bond plus two $\pi$ bonds (Fig. 9A.4), which is consistent with the Lewis structure : $\mathrm{N} \equiv \mathrm{N}$ : for dinitrogen.\\
\includegraphics[max width=\textwidth, center]{2025_02_27_d4b95d9718f0b9e2e6b9g-377(2)}

Figure 9A. 4 The structure of bonds in a nitrogen molecule, with one $\sigma$ bond and two $\pi$ bonds. The overall electron density has cylindrical symmetry around the internuclear axis.

\footnotetext{${ }^{1} \pi$ bonds can also be formed from d orbitals in the appropriate orientation.
}\subsection*{9A. 2 Resonance}
Another term introduced into chemistry by VB theory is resonance, the superposition of the wavefunctions representing different electron distributions in the same nuclear framework. To understand what this means, consider the VB description of a purely covalently bonded HCl molecule, which could be written as $\Psi_{\mathrm{H}-\mathrm{Cl}}=\psi_{\mathrm{A}}(1) \psi_{\mathrm{B}}(2)+\psi_{\mathrm{A}}(2) \psi_{\mathrm{B}}(1)$, with $\psi_{\mathrm{A}}$ now a H1s orbital and $\psi_{\mathrm{B}}$ a Cl3p orbital. This description allows electron 1 to be on the H atom when electron 2 is on the Cl atom, and vice versa, but it does not allow for the possibility that both electrons are on the Cl atom $\left(\Psi_{\mathrm{H}^{+} \mathrm{Cl}^{-}}=\psi_{\mathrm{B}}(1) \psi_{\mathrm{B}}(2)\right.$, representing the ionic form $\mathrm{H}^{+} \mathrm{Cl}^{-}$) or both are on the H atom $\left(\Psi_{\mathrm{H}^{-} \mathrm{Cl}^{+}}=\psi_{\mathrm{A}}(1) \psi_{\mathrm{A}}(2)\right.$, representing the much less likely ionic form $\left.\mathrm{H}^{-} \mathrm{Cl}^{+}\right)$. A better description of the wavefunction for the molecule is as a superposition of the covalent and ionic descriptions, written as (with a slightly simplified notation, and ignoring the less likely $\mathrm{H}^{-} \mathrm{Cl}^{+}$possibility) $\Psi_{\mathrm{HCl}}=\Psi_{\mathrm{H}-\mathrm{Cl}}+\lambda \Psi_{\mathrm{H}^{+} \mathrm{Cl}^{-}}$with $\lambda$ (lambda) some numerical coefficient. In general,


\begin{equation*}
\Psi=\Psi_{\text {covalent }}+\lambda \Psi_{\text {ionic }} \tag{9A.3}
\end{equation*}


where $\Psi_{\text {covalent }}$ is the two-electron wavefunction for the purely covalent form of the bond and $\Psi_{\text {ionic }}$ is the twoelectron wavefunction for the ionic form of the bond. In this case, where one structure is pure covalent and the other pure ionic, it is called ionic-covalent resonance. The interpretation of the (un-normalized) wavefunction, which is called a resonance hybrid, is that if the molecule is inspected, then the probability that it would be found with an ionic structure is proportional to $\lambda^{2}$. If $\lambda^{2} \ll 1$, the covalent description is dominant. If $\lambda^{2} \gg 1$, the ionic description is dominant. Resonance is not a flickering between the contributing states: it is a blending of their characteristics. It is only a mathematical device for achieving a closer approximation to the true wavefunction of the molecule than that represented by any single contributing electronic structure alone.

A systematic way of calculating the value of $\lambda$ is provided by the variation principle:

If an arbitrary wavefunction is used to calculate the energy, then the value calculated is never less than the true energy.\\
(This principle is derived and used in Topic 9C.) The arbitrary wavefunction is called the trial wavefunction. The principle implies that if the energy, the expectation value of the hamiltonian, is calculated for various trial wavefunctions with different values of the parameter $\lambda$, then the best value of $\lambda$ is the one that results in the lowest energy. The ionic contribution to the resonance is then proportional to $\lambda^{2}$.

\subsection*{Brief illustration 9A. 2}
Consider a bond described by eqn 9A.3. If the lowest energy is reached when $\lambda=0.1$, then the best description of the bond in the molecule is a resonance structure described by the wavefunction $\Psi=\Psi_{\text {covalent }}+0.1 \Psi_{\text {ionic. }}$. This wavefunction implies that the probabilities of finding the molecule in its covalent and ionic forms are in the ratio 100:1 (because $0.1^{2}=0.01$ ).

\subsection*{9A. 3 Polyatomic molecules}
Each $\sigma$ bond in a polyatomic molecule is formed by the spin pairing of electrons in atomic orbitals with cylindrical symmetry around the relevant internuclear axis. Likewise, $\pi$ bonds are formed by pairing electrons that occupy atomic orbitals of the appropriate symmetry.

\subsection*{Brief illustration 9A. 3}
The VB description of $\mathrm{H}_{2} \mathrm{O}$ is as follows. The valence-electron configuration of an O atom is $2 \mathrm{~s}^{2} 2 \mathrm{p}_{x}^{2} 2 \mathrm{p}_{y}^{1} 2 \mathrm{p}_{z}^{1}$. The two unpaired electrons in the O2p orbitals can each pair with an electron in an H1s orbital, and each combination results in the formation of a $\sigma$ bond (each bond has cylindrical symmetry about the respective $\mathrm{O}-\mathrm{H}$ internuclear axis). Because the $2 \mathrm{p}_{y}$ and $2 \mathrm{p}_{z}$ orbitals lie at $90^{\circ}$ to each other, the two $\sigma$ bonds also lie at $90^{\circ}$ to each other (Fig. 9A.5). Therefore, $\mathrm{H}_{2} \mathrm{O}$ is predicted to be an angular molecule, which it is. However, the theory predicts a bond angle of $90^{\circ}$, whereas the actual bond angle is $104.5^{\circ}$.\\
\includegraphics[max width=\textwidth, center]{2025_02_27_d4b95d9718f0b9e2e6b9g-378}

Figure 9A. 5 In a primitive view of the structure of an $\mathrm{H}_{2} \mathrm{O}$ molecule, each bond is formed by the overlap and spin pairing of an H 1 s electron and an O 2 p electron.

Resonance plays an important role in the VB description of polyatomic molecules. One of the most famous examples of resonance is in the VB description of benzene, where the wavefunction of the molecule is written as a superposition of the many-electron wavefunctions of the two covalent Kekule structures:


\begin{equation*}
\Psi=\Psi(\triangleq)+\Psi(\Omega) \tag{9A.4}
\end{equation*}


The two contributing structures have identical energies, so they contribute equally to the superposition. The effect of resonance (which is represented by a double-headed arrow (1)), in this case, is to distribute double-bond character around the ring and to make the lengths\\
\includegraphics[max width=\textwidth, center]{2025_02_27_d4b95d9718f0b9e2e6b9g-379(1)}

1 and strengths of all the carbon-carbon bonds identical. The wavefunction is improved by allowing resonance because it allows the electrons to adjust into a distribution of lower energy. This lowering is called the resonance stabilization of the molecule and, in the context of VB theory, is largely responsible for the unusual stability of aromatic rings. Resonance always lowers the energy, and the lowering is greatest when the contributing structures have similar energies. The wavefunction of benzene is improved still further, and the calculated energy of the molecule is lowered still further, if ionic-covalent resonance is also considered, by allowing a small admixture of ionic structures, such as (2).\\
\includegraphics[max width=\textwidth, center]{2025_02_27_d4b95d9718f0b9e2e6b9g-379}

\subsection*{(a) Promotion}
A deficiency of this initial formulation of VB theory is its inability to account for the common tetravalence of carbon (its ability to form four bonds). The ground-state configuration of carbon is $2 \mathrm{~s}^{2} 2 \mathrm{p}_{x}^{1} 2 \mathrm{p}_{y}^{1}$, which suggests that a carbon atom should be capable of forming only two bonds, not four.

This deficiency is overcome by allowing for promotion, the excitation of an electron to an orbital of higher energy. In carbon, for example, the promotion of a $2 s$ electron to a $2 p$ orbital can be thought of as leading to the configuration $2 \mathrm{~s}^{1} 2 \mathrm{p}_{x}^{1} 2 \mathrm{p}_{y}^{1} 2 \mathrm{p}_{z}^{1}$, with four unpaired electrons in separate orbitals. These electrons may pair with four electrons in orbitals provided by four other atoms (such as four H1s orbitals if the molecule is $\mathrm{CH}_{4}$ ), and hence form four $\sigma$ bonds. Although energy is required to promote the electron, it is more than recovered by the promoted atom's ability to form four bonds in place of the two bonds of the unpromoted atom.

Promotion, and the formation of four bonds, is a characteristic feature of carbon because the promotion energy is quite small: the promoted electron leaves a doubly occupied 2 s orbital and enters a vacant 2 p orbital, hence significantly relieving the electron-electron repulsion it experiences in the ground state. However, it is important to remember that promotion is not a 'real' process in which an atom somehow becomes excited and then forms bonds: it is a notional contribution to the overall energy change that occurs when bonds form.

\subsection*{Brief illustration 9A. 4}
Sulfur can form six bonds (an 'expanded octet'), as in the molecule $\mathrm{SF}_{6}$. Because the ground-state electron configuration of sulfur is $[\mathrm{Ne}] 3 \mathrm{~s}^{2} 3 \mathrm{p}^{4}$, this bonding pattern requires the promotion of a 3 s electron and a 3 p electron to two\\
different 3d orbitals, which are nearby in energy, to produce the notional configuration $[\mathrm{Ne}] 3 \mathrm{~s}^{1} 3 \mathrm{p}^{3} 3 \mathrm{~d}^{2}$ with all six of the valence electrons in different orbitals and capable of bond formation with six electrons provided by six F atoms.

\subsection*{(b) Hybridization}
The description of the bonding in $\mathrm{CH}_{4}$ (and other alkanes) is still incomplete because it implies the presence of three $\sigma$ bonds of one type (formed from H1s and C2p orbitals) and a fourth $\sigma$ bond of a distinctly different character (formed from H1s and $\mathrm{C} 2 \mathrm{~s})$. This problem is overcome by realizing that the electron density distribution in the promoted atom is equivalent to the electron density in which each electron occupies a hybrid orbital formed by interference between the C2s and C2p orbitals of the same atom. The origin of the hybridization can be appreciated by thinking of the four atomic orbitals centred on a nucleus as waves that interfere destructively and constructively in different regions, and give rise to four new shapes.

The specific linear combinations that give rise to four equivalent hybrid orbitals can be constructed by considering their tetrahedral arrangement.

\subsection*{How is that done? 9A. 2}
Constructing tetrahedral hybrid orbitals

Each tetrahedral bond can be regarded as directed to one corner of a unit cube (3). Suppose that each hybrid can be written in the form $h=a s+b_{x} \mathrm{p}_{x}+b_{y} \mathrm{p}_{y}+b_{z} \mathrm{p}_{z}$. The hybrid $h_{1}$ that points to the corner with coordinates ( $1,1,1$ ) must have equal contributions from all three p orbitals, so the three $b$ coefficients can be set equal to each other and $h_{1}=a s+b\left(\mathrm{p}_{x}+\mathrm{p}_{y}+\mathrm{p}_{z}\right)$. The other three hybrids have the same composition (they are equivalent, apart from their direction in space), but are orthogonal to $h_{1}$. This orthogonality is achieved by choosing different signs for the p orbitals but the same overall composition. For instance, choosing $h_{2}=a s+b\left(-\mathrm{p}_{x}-\mathrm{p}_{y}+\mathrm{p}_{z}\right)$, the orthogonality condition is\\
$\int h_{1} h_{2} \mathrm{~d} \tau=\int\left\{a \mathrm{~s}+b\left(\mathrm{p}_{x}+\mathrm{p}_{y}+\mathrm{p}_{z}\right)\right\}\left\{a \mathrm{~s}+b\left(-\mathrm{p}_{x}-\mathrm{p}_{y}+\mathrm{p}_{z}\right)\right\} \mathrm{d} \tau$

$$
\begin{aligned}
& =a^{2} \overbrace{\int \mathrm{~s}^{2} \mathrm{~d} \tau}^{1}-b^{2} \overbrace{\int \mathrm{p}_{x}^{2} \mathrm{~d} \tau}^{1}-\cdots-a b \overbrace{\int \mathrm{sp}_{x} \mathrm{~d} \tau}^{0}-\cdots-b^{2} \overbrace{\int \mathrm{p}_{x} \mathrm{p}_{y} \mathrm{~d} \tau}^{0}+\cdots \\
& =a^{2}-b^{2}-b^{2}+b^{2}=a^{2}-b^{2}=0
\end{aligned}
$$

\begin{center}
\includegraphics[max width=\textwidth]{2025_02_27_d4b95d9718f0b9e2e6b9g-379(2)}
\end{center}

The values of the integrals come from the fact that the atomic orbitals are normalized and mutually orthogonal (Topic 7C). It follows that a solution is $a=b$ (the alternative solution, $a=$ $-b$, simply corresponds to choosing different absolute phases for the p orbitals) and that the two hybrid orbitals are $h_{1}=\mathrm{s}+$ $\mathrm{p}_{x}+\mathrm{p}_{y}+\mathrm{p}_{z}$ and $h_{2}=\mathrm{s}-\mathrm{p}_{x}-\mathrm{p}_{y}+\mathrm{p}_{z}$. A similar argument but with $h_{3}=a s+b\left(-\mathrm{p}_{x}+\mathrm{p}_{y}-\mathrm{p}_{z}\right)$ or $h_{4}=a s+b\left(\mathrm{p}_{x}-\mathrm{p}_{y}-\mathrm{p}_{z}\right)$ leads to two other hybrids. In sum,

$$
-\left|\begin{array}{ll}
h_{1}=\mathrm{s}+\mathrm{p}_{x}+\mathrm{p}_{y}+\mathrm{p}_{z} & h_{2}=\mathrm{s}-\mathrm{p}_{x}-\mathrm{p}_{y}+\mathrm{p}_{z} \\
h_{3}=\mathrm{s}-\mathrm{p}_{x}+\mathrm{p}_{y}-\mathrm{p}_{z} & h_{4}=\mathrm{s}+\mathrm{p}_{x}-\mathrm{p}_{y}-\mathrm{p}_{z}
\end{array}\right|>\begin{aligned}
& \text { sp hybrid } \\
& \text { orbitals }
\end{aligned}
$$

As a result of the interference between the component orbitals, each hybrid orbital consists of a large lobe pointing in the direction of one corner of a regular tetrahedron (Fig. 9A.6). The angle between the axes of the hybrid orbitals is the tetrahedral angle, $\arccos \left(-\frac{1}{3}\right)=109.47^{\circ}$. Because each hybrid is built from one s orbital and three p orbitals, it is called an sp ${ }^{\mathbf{3}}$ hybrid orbital.

It is now straightforward to see how the VB description of the $\mathrm{CH}_{4}$ molecule leads to a tetrahedral molecule containing four equivalent $\mathrm{C}-\mathrm{H}$ bonds. Each hybrid orbital of the promoted C atom contains a single unpaired electron; an H1s electron can pair with each one, giving rise to a $\sigma$ bond pointing to a corner of a tetrahedron. For example, the (un-normalized) two-electron wavefunction for the bond formed by the hybrid orbital $h_{1}$ and the $1 \mathrm{~s}_{\mathrm{A}}$ orbital is


\begin{equation*}
\Psi(1,2)=h_{1}(1) \psi_{\mathrm{H} 1 \mathrm{~s}}(2)+h_{1}(2) \psi_{\mathrm{H} 1 \mathrm{~s}}(1) \tag{9A.6}
\end{equation*}


As for $\mathrm{H}_{2}$, to achieve this wavefunction, the two electrons it describes must be paired. Because each $\mathrm{sp}^{3}$ hybrid orbital has the same composition, all four $\sigma$ bonds are identical apart from their orientation in space (Fig. 9A.7).

A hybrid orbital has enhanced amplitude in the internuclear region, which arises from the constructive interference between the s orbital and the positive lobes of the p orbitals. As a result, the bond strength is greater than for a bond formed\\
\includegraphics[max width=\textwidth, center]{2025_02_27_d4b95d9718f0b9e2e6b9g-380}

Figure 9A. $6 \mathrm{An} \mathrm{sp}^{3}$ hybrid orbital formed from the superposition of $s$ and $p$ orbitals on the same atom. There are four such hybrids: each one points towards the corner of a regular tetrahedron. The overall electron density remains spherically symmetrical.\\
\includegraphics[max width=\textwidth, center]{2025_02_27_d4b95d9718f0b9e2e6b9g-380(1)}

Figure 9A. 7 Each $\mathrm{sp}^{3}$ hybrid orbital forms a $\sigma$ bond by overlap with an H1s orbital located at the corner of the tetrahedron. This model is consistent with the equivalence of the four bonds in $\mathrm{CH}_{4}$.\\
from an $s$ or $p$ orbital alone. This increased bond strength is another factor that helps to repay the promotion energy.

The hybridization of $N$ atomic orbitals always results in the formation of $N$ hybrid orbitals, which may either form bonds or may contain lone pairs of electrons, pairs of electrons that do not participate directly in bond formation (but may influence the shape of the molecule).

\subsection*{Brief illustration 9A. 5}
To accommodate the observed bond angle of $104.5^{\circ}$ in $\mathrm{H}_{2} \mathrm{O}$ in VB theory it is necessary to suppose that the oxygen 2 s and three 2 p orbitals hybridize. As a first approximation, suppose they hybridize to form four equivalent $\mathrm{sp}^{3}$ orbitals. Four electrons pair and occupy two of the hybrids, and so become lone pairs. The remaining two pair with the two electrons on the H atoms, and so form two $\mathrm{O}-\mathrm{H}$ bonds at $109.5^{\circ}$. The actual hybridization will be slightly different to account for the observed bond angle not being exactly the tetrahedral angle.

Hybridization is also used to describe the structure of an ethene molecule, $\mathrm{H}_{2} \mathrm{C}=\mathrm{CH}_{2}$, and the torsional rigidity of double bonds. An ethene molecule is planar, with HCH and HCC bond angles close to $120^{\circ}$. To reproduce the $\sigma$ bonding structure, each $C$ atom is regarded as being promoted to a $2 s^{1} 2 p^{3}$ configuration. However, instead of using all four orbitals to form hybrids, $\mathbf{s p}^{2}$ hybrid orbitals are formed:

$$
\begin{aligned}
& h_{1}=s+2^{1 / 2} \mathrm{p}_{y} \\
& h_{2}=s+\left(\frac{3}{2}\right)^{1 / 2} \mathrm{p}_{x}-\left(\frac{1}{2}\right)^{1 / 2} \mathrm{p}_{y} \\
& h_{3}=\mathrm{s}-\left(\frac{3}{2}\right)^{1 / 2} \mathrm{p}_{x}-\left(\frac{1}{2}\right)^{1 / 2} \mathrm{p}_{y}
\end{aligned}
$$

These hybrids lie in a plane and point towards the corners of an equilateral triangle at $120^{\circ}$ to each other (Fig. 9A.8). The\\
\includegraphics[max width=\textwidth, center]{2025_02_27_d4b95d9718f0b9e2e6b9g-381(2)}

Figure 9A. 8 (a) An s orbital and two p orbitals can be hybridized to form three equivalent orbitals that point towards the corners of an equilateral triangle. (b) The remaining, unhybridized $p$ orbital is perpendicular to the plane.\\
third 2 p orbital $\left(2 \mathrm{p}_{z}\right)$ is not included in the hybridization; it lies along an axis perpendicular to the plane formed by the hybrids. The different signs of the coefficients, as well as ensuring that the hybrids are mutually orthogonal, also ensure that constructive interference takes place in different regions of space, so giving the patterns in the illustration. The $\mathrm{sp}^{2}$ hybridized C atoms each form three $\sigma$ bonds by spin pairing with either a hybrid orbital on the other C atom or with H 1 s orbitals. The $\sigma$ framework therefore consists of $\mathrm{C}-\mathrm{H}$ and $\mathrm{C}-\mathrm{C}$ $\sigma$ bonds at $120^{\circ}$ to each other. When the two $\mathrm{CH}_{2}$ groups lie in the same plane, each electron in the two unhybridized $p$ orbitals can pair and form a $\pi$ bond (Fig. 9A.9). The formation of this $\pi$ bond locks the framework into the planar arrangement, because any rotation of one $\mathrm{CH}_{2}$ group relative to the other leads to a weakening of the $\pi$ bond (and consequently an increase in energy of the molecule).

A similar description applies to ethyne, $\mathrm{HC} \equiv \mathrm{CH}$, a linear molecule. Now the C atoms are sp hybridized, and the $\sigma$ bonds are formed using hybrid atomic orbitals of the form

\[
h_{1}=\mathrm{s}+\mathrm{p}_{z} \quad h_{2}=\mathrm{s}-\mathrm{p}_{z} \quad \begin{align*}
& \text { sp hybrid }  \tag{9A.8}\\
& \text { orbitals }
\end{align*}
\]

These two hybrids lie along the internuclear axis (conventionally the $z$-axis in a linear molecule). The electrons in them pair either with an electron in the corresponding hybrid orbital on\\
\includegraphics[max width=\textwidth, center]{2025_02_27_d4b95d9718f0b9e2e6b9g-381(1)}

Figure 9A. 9 A representation of the structure of a double bond in ethene; only the $\pi$ bond is shown explicitly.\\
\includegraphics[max width=\textwidth, center]{2025_02_27_d4b95d9718f0b9e2e6b9g-381}

Figure 9A.10 A representation of the structure of a triple bond in ethyne; only the $\pi$ bonds are shown explicitly. The overall electron density has cylindrical symmetry around the axis of the molecule.\\
the other C atom or with an electron in one of the H1s orbitals. Electrons in the two remaining $p$ orbitals on each atom, which are perpendicular to the molecular axis, pair to form two perpendicular $\pi$ bonds (Fig. 9A.10).

Other hybridization schemes, particularly those involving d orbitals, are often invoked in VB descriptions of molecular structure to be consistent with other molecular geometries (Table 9A.1).

\subsection*{Brief illustration 9A. 6}
Consider an octahedral molecule, such as $\mathrm{SF}_{6}$. The promotion of sulfur's electrons as in Brief illustration 9A.4, followed by $s p^{3} \mathrm{~d}^{2}$ hybridization results in six equivalent hybrid orbitals pointing towards the corners of a regular octahedron.

Table 9A. 1 Some hybridization schemes

\begin{center}
\begin{tabular}{lll}
\hline
\begin{tabular}{l}
Coordination \\
number \\
\end{tabular} & Arrangement & Composition \\
\hline
2 & Linear & $\mathrm{sp}, \mathrm{pd}, \mathrm{sd}$ \\
 & Angular & sd \\
3 & Trigonal planar & $\mathrm{sp}^{2}, \mathrm{p}^{2} \mathrm{~d}$ \\
 & Unsymmetrical planar & spd \\
 & Trigonal pyramidal & $\mathrm{pd}^{2}$ \\
4 & Tetrahedral & $\mathrm{sp}^{3}, \mathrm{sd}^{3}$ \\
 & Irregular tetrahedral & $\mathrm{spd}^{2}, \mathrm{p}^{3} \mathrm{~d}, \mathrm{dp}^{3}$ \\
 & Square planar & $\mathrm{p}^{2} \mathrm{~d}^{2}, \mathrm{sp}^{2} \mathrm{~d}$ \\
5 & Trigonal bipyramidal & $\mathrm{sp}^{3} \mathrm{~d}, \mathrm{spd}^{3}$ \\
 & Tetragonal pyramidal & $\mathrm{sp}^{2} \mathrm{~d}^{2}, \mathrm{sd}^{4}, \mathrm{pd}^{4}, \mathrm{p}^{3} \mathrm{~d}^{2}$ \\
 & Pentagonal planar & $\mathrm{p}^{2} \mathrm{~d}^{3}$ \\
 & Octahedral & $\mathrm{sp}^{3} \mathrm{~d}^{2}$ \\
6 & Trigonal prismatic & $\mathrm{spd}^{4}, \mathrm{pd}^{5}$ \\
 & Trigonal antiprismatic & $\mathrm{p}^{3} \mathrm{~d}^{3}$ \\
\hline
\end{tabular}
\end{center}

\subsection*{Checklist of concepts}
1. A bond forms when an electron in an atomic orbital on one atom pairs its spin with that of an electron in an atomic orbital on another atom.2. A $\sigma$ bond has cylindrical symmetry around the internuclear axis.3. Resonance is the superposition of structures with different electron distributions but the same nuclear arrangement.4. A $\pi$ bond has symmetry like that of a $p$ orbital perpendicular to the internuclear axis.5. Promotion is the notional excitation of an electron to an empty orbital to enable the formation of additional bonds.6. Hybridization is the blending together of atomic orbitals on the same atom to achieve the appropriate directional properties and enhanced overlap.\subsection*{Checklist of equations}
\begin{center}
\begin{tabular}{llll}
\hline
Property & Equation & Comment & Equation number \\
\hline
Valence-bond wavefunction & $\Psi=\psi_{\mathrm{A}}(1) \psi_{\mathrm{B}}(2)+\psi_{\mathrm{A}}(2) \psi_{\mathrm{B}}(1)$ & Spins must be paired ${ }^{*}$ & 9 A .2 \\
Resonance & $\Psi=\Psi_{\text {covalent }}+\lambda \Psi_{\text {ionic }}$ & Ionic-covalent resonance & 9 A .3 \\
Hybridization & $h=a \mathrm{~s}+b \mathrm{p}+\cdots$ & All atomic orbitals on the same atom; specific forms in the text & $9 \mathrm{~A} .5,9 \mathrm{~A} .7$, and 9A.8 \\
\hline
\end{tabular}
\end{center}

\begin{itemize}
  \item The spin contribution is $\sigma_{-}(1,2)=\frac{1}{2^{1 / 2}}\{\alpha(1) \beta(2)-\beta(1) \alpha(2)\}$
\end{itemize}

\section*{TOPIC 9B Molecular orbital theory: the hydrogen molecule-ion }
\subsection*{- Why do you need to know this material?}
Molecular orbital theory is the basis of almost all descriptions of chemical bonding, in both individual molecules and solids.

\subsection*{- What is the key idea?}
Molecular orbitals are wavefunctions that spread over all the atoms in a molecule and are commonly represented as linear combinations of atomic orbitals.

\subsection*{$>$ What do you need to know already?}
You need to be familiar with the shapes of atomic orbitals (Topic 8A) and how an energy is calculated from a wavefunction (Topic 7C). The entire discussion is within the framework of the Born-Oppenheimer approximation (see the Prologue for this Focus).

In molecular orbital theory (MO theory), electrons do not belong to particular bonds but spread throughout the entire molecule. This theory has been more fully developed than valence-bond theory (Topic 9A) and provides the language that is widely used in modern discussions of bonding. To introduce it, the strategy of Topic 8 A is followed, where the one-electron H atom is taken as the fundamental species for discussing atomic structure and then developed into a description of many-electron atoms. This Topic uses the simplest molecular species of all, the hydrogen molecule-ion, $\mathrm{H}_{2}^{+}$, to introduce the essential features of the theory, which are then used in subsequent Topics to describe the structures of more complex systems.

\subsection*{98. 1 Linear combinations of atomic orbitals}
The hamiltonian for the single electron in $\mathrm{H}_{2}^{+}$is


\begin{equation*}
\hat{H}=-\frac{\hbar^{2}}{2 m_{\mathrm{e}}} \nabla_{1}^{2}+V \quad V=-\frac{e^{2}}{4 \pi \varepsilon_{0}}\left(\frac{1}{r_{\mathrm{A} 1}}+\frac{1}{r_{\mathrm{B} 1}}-\frac{1}{R}\right) \tag{9B.1}
\end{equation*}


where $r_{\mathrm{A} 1}$ and $r_{\mathrm{B} 1}$ are the distances of the electron from the two nuclei $A$ and $B$ (1) and $R$ is the distance between the two nuclei. In the expression for $V$, the first two terms in parentheses are\\
\includegraphics[max width=\textwidth]{2025_02_27_d4b95d9718f0b9e2e6b9g-383} the attractive contribution from the interaction between the electron and the nuclei; the remaining term is the repulsive interaction between the nuclei. The collection of fundamental constant $e^{2} / 4 \pi \varepsilon_{0}$ occurs widely throughout this chapter, and is denoted $j_{0}$.

The one-electron wavefunctions obtained by solving the Schrödinger equation $\hat{H} \psi=E \psi$ are called molecular orbitals. A molecular orbital $\psi$ gives, through the value of $|\psi|^{2}$, the distribution of the electron in the molecule. A molecular orbital is like an atomic orbital, but spreads throughout the molecule.

\subsection*{(a) The construction of linear combinations}
The Schrödinger equation can be solved analytically for $\mathrm{H}_{2}^{+}$ (within the Born-Oppenheimer approximation), but the wavefunctions are very complicated functions; moreover, the solution cannot be extended to polyatomic systems. The simpler procedure adopted here, while more approximate, can be extended readily to other molecules.

If an electron can be found in an atomic orbital $\psi_{\mathrm{A}}$ belonging to atom A and also in an atomic orbital $\psi_{\mathrm{B}}$ belonging to atom $B$, then the overall wavefunction is a superposition of the two atomic orbitals:

\[
\begin{array}{ll}
\psi_{ \pm}=N_{ \pm}\left(\psi_{\mathrm{A}} \pm \psi_{\mathrm{B}}\right) & \begin{array}{l}
\text { Linear combination } \\
\text { of atomic orbitals }
\end{array} \tag{9В.2}
\end{array}
\]

where, for $\mathrm{H}_{2}^{+}, \psi_{\mathrm{A}}$ and $\psi_{\mathrm{B}}$ are 1 s atomic orbitals on atom A and B , respectively, and $N_{ \pm}$is a normalization factor. The technical term for the type of superposition in eqn 9B. 2 is a linear combination of atomic orbitals (LCAO). An approximate molecular orbital formed from a linear combination of atomic orbitals is called an LCAO-MO. A molecular orbital that has cylindrical symmetry around the internuclear axis, such as the one being discussed, is called a $\sigma$ orbital because it resembles an s orbital when viewed along the axis and, more precisely, because it has zero orbital angular momentum around the internuclear axis.

\subsection*{Example 9B. 1 Normalizing a molecular orbital}
Normalize (to 1 ) the molecular orbital $\psi_{+}$in eqn 9B.2.\\
Collect your thoughts You need to find the factor $N_{+}$such that $\int \psi^{\star} \psi \mathrm{d} \tau=1$, where the integration is over the whole of space. To proceed, you should substitute the LCAO into this integral and make use of the fact that the atomic orbitals are individually normalized.

The solution Substitution of the wavefunction gives

$$
\int \psi^{\star} \psi \mathrm{d} \tau=N_{+}^{2}\{\overbrace{\int \psi_{\mathrm{A}}^{2} \mathrm{~d} \tau}^{1}+\overbrace{\int \psi_{\mathrm{B}}^{2} \mathrm{~d} \tau}^{1}+2 \overbrace{\int \psi_{\mathrm{A}} \psi_{\mathrm{B}} \mathrm{~d} \tau}^{S}\}=2(1+S) N_{+}^{2}
$$

where $S=\int \psi_{\mathrm{A}} \psi_{\mathrm{B}} \mathrm{d} \tau$ and has a value that depends on the nuclear separation (this 'overlap integral' will play a significant role later). For the integral to be equal to 1 ,

$$
N_{+}=\frac{1}{\{2(1+S)\}^{1 / 2}}
$$

For $\mathrm{H}_{2}^{+}$at its equilibrium bond length $S \approx 0.59$, so $N_{+}=0.56$.\\
Self-test 9B. 1 Normalize the orbital $\psi_{-}$in eqn 9B. 2 and evaluate $N_{-}$for $S=0.59$.

Figure 9B. 1 shows the contours of constant amplitude for the molecular orbital $\psi_{+}$in eqn 9B.2. Plots like these are readily obtained using commercially available software. The calculation is quite straightforward, because all that it is necessary to do is to feed in the mathematical forms of the two atomic orbitals and then let the software do the rest.\\
\includegraphics[max width=\textwidth, center]{2025_02_27_d4b95d9718f0b9e2e6b9g-384}

Figure 9B. 1 (a) The amplitude of the bonding molecular orbital in a hydrogen molecule-ion in a plane containing the two nuclei and (b) a contour representation of the amplitude.

\subsection*{Brief illustration 9B. 1}
The surfaces of constant amplitude shown in Fig. 9B. 2 have been calculated using the two H1s orbitals

$$
\psi_{\mathrm{A}}=\frac{1}{\left(\pi a_{0}^{3}\right)^{1 / 2}} \mathrm{e}^{-r_{\mathrm{A}} / a_{0}} \quad \psi_{\mathrm{B}}=\frac{1}{\left(\pi a_{0}^{3}\right)^{1 / 2}} \mathrm{e}^{-r_{\mathrm{B}} / a_{0}}
$$

and noting that $r_{\mathrm{A} 1}$ and $r_{\mathrm{B} 1}$ are not independent (1). When expressed in Cartesian coordinates based on atom A (2), these radii are given by $r_{\mathrm{A} 1}=\left\{x^{2}+y^{2}+z^{2}\right\}^{1 / 2}$ and $r_{\mathrm{B} 1}=\left\{x^{2}+y^{2}+\right.$ $\left.(z-R)^{2}\right\}^{1 / 2}$, where $R$ is the bond length. A repeat of the analysis for $\psi_{-}$gives the results shown in Fig. 9B.3.\\
\includegraphics[max width=\textwidth, center]{2025_02_27_d4b95d9718f0b9e2e6b9g-384(2)}

Figure 9B. 2 Surfaces of constant amplitude of the wavefunction $\psi_{+}$of the hydrogen molecule-ion.\\
\includegraphics[max width=\textwidth, center]{2025_02_27_d4b95d9718f0b9e2e6b9g-384(1)}

Figure 9B. 3 Surfaces of constant amplitude of the wavefunction $\psi_{\text {- }}$ of the hydrogen molecule-ion.

\subsection*{(b) Bonding orbitals}
According to the Born interpretation, the probability density of the electron at each point in $\mathrm{H}_{2}^{+}$is proportional to the square modulus of its wavefunction at that point. The probability density corresponding to the (real) wavefunction $\psi_{+}$in eqn 9B. 2 is

\[
\psi_{+}^{2} \propto \psi_{\mathrm{A}}^{2}+\psi_{\mathrm{B}}^{2}+2 \psi_{\mathrm{A}} \psi_{\mathrm{B}} \quad \begin{align*}
& \text { Bonding probability }  \tag{9B.3}\\
& \text { density }
\end{align*}
\]

This probability density is plotted in Fig. 9B.4. An important feature becomes apparent in the internuclear region, where both atomic orbitals have similar amplitudes. According to eqn 9B.3, the total probability density is proportional to the sum of:

\begin{itemize}
  \item $\psi_{A}^{2}$, the probability density if the electron were confined to atom A;
  \item $\psi_{B}^{2}$, the probability density if the electron were confined to atom B;
  \item $2 \psi_{A} \psi_{B}$, an extra contribution to the density from both atomic orbitals.
\end{itemize}

The last contribution, the overlap density, is crucial, because it represents an enhancement of the probability of finding the electron in the internuclear region. The enhancement can be traced to the constructive interference of the two atomic orbitals: each has a positive amplitude in the internuclear region, so the total amplitude is greater there than if the electron were confined to a single atom. This observation is summarized as

\subsection*{Bonds form as a result of the build-up of electron density where atomic orbitals overlap and interfere constructively.}
The conventional explanation of this observation is based on the notion that accumulation of electron density between the nuclei puts the electron in a position where it interacts strongly with both nuclei. Hence, the energy of the molecule is lower than that of the separate atoms, where each electron can interact strongly with only one nucleus. This conventional explanation, however, has been called into question, because shifting an electron away from a nucleus into the internuclear region raises its potential energy. The modern (and still controversial) explanation does not emerge from the simple LCAO treatment given here. It seems that, at the same time as the electron shifts into the internuclear region, the atomic orbitals shrink. This orbital shrinkage improves the electron-nucleus attraction more than it is decreased by the migration to the internuclear region, so there is a net lowering of potential energy. The kinetic energy of the electron is also modified because the curvature of the wavefunction is changed, but the change\\
\includegraphics[max width=\textwidth, center]{2025_02_27_d4b95d9718f0b9e2e6b9g-385}

Figure 9B. 4 The electron density calculated by forming the square of the wavefunction used to construct Fig. 9B.2. Note the accumulation of electron density in the internuclear region.\\
in kinetic energy is dominated by the change in potential energy. Throughout the following discussion the strength of chemical bonds is ascribed to the accumulation of electron density in the internuclear region. In molecules more complicated than $\mathrm{H}_{2}^{+}$the true source of energy lowering may be this accumulation of electron density or some indirect but related effect.

The $\sigma$ orbital just described is an example of a bonding orbital, an orbital which, if occupied, helps to bind two atoms together. An electron that occupies a $\sigma$ orbital is called a $\sigma$ electron, and if that is the only electron present in the molecule (as in the ground state of $\mathrm{H}_{2}^{+}$), then the configuration of the molecule is $\sigma^{1}$.

The energy $E_{\sigma}$ of the $\sigma$ orbital is: ${ }^{1}$

\[
E_{\sigma}=E_{\mathrm{H} 1 \mathrm{~s}}+\frac{j_{0}}{R}-\frac{j+k}{1+S} \quad \begin{align*}
& \text { Energy of }  \tag{9B.4}\\
& \text { bonding orbital }
\end{align*}
\]

where $E_{\mathrm{H} 1 \mathrm{~s}}$ is the energy of a H1s orbital, $j_{0} / R$ is the potential energy of repulsion between the two nuclei (remember that $j_{0}$ is shorthand for $e^{2} / 4 \pi \varepsilon_{0}$ ), and


\begin{align*}
& S=\int \psi_{\mathrm{A}} \psi_{\mathrm{B}} \mathrm{~d} \tau=\left\{1+\frac{R}{a_{0}}+\frac{1}{3}\left(\frac{R}{a_{0}}\right)^{2}\right\} \mathrm{e}^{-R / a_{0}}  \tag{9B.5a}\\
& j=j_{0} \int \frac{\psi_{\mathrm{A}}^{2}}{r_{\mathrm{B}}} \mathrm{~d} \tau=\frac{j_{0}}{R}\left\{1-\left(1+\frac{R}{a_{0}}\right) \mathrm{e}^{-2 R / a_{0}}\right\}  \tag{9B.5b}\\
& k=j_{0} \int \frac{\psi_{\mathrm{A}} \psi_{\mathrm{B}}}{r_{\mathrm{B}}} \mathrm{~d} \tau=\frac{j_{0}}{a_{0}}\left(1+\frac{R}{a_{0}}\right) \mathrm{e}^{-R / a_{0}} \tag{9B.5c}
\end{align*}


Note that


\begin{equation*}
\frac{j_{0}}{a_{0}}=\frac{e^{2}}{4 \pi \varepsilon_{0} a_{0}}=\frac{e^{2}}{4 \pi \varepsilon_{0}} \times \frac{\pi m_{e} e^{2}}{\varepsilon_{0} h^{2}}=\frac{m_{e} e^{4}}{4 \varepsilon_{0}^{2} h^{2}}=2 h c \tilde{R}_{\infty} \tag{9B.5d}
\end{equation*}


\footnotetext{${ }^{1}$ For a derivation of eqn 9B.4, see $A$ deeper look 4 on the website for this text.
}
\includegraphics[max width=\textwidth, center]{2025_02_27_d4b95d9718f0b9e2e6b9g-386(2)}

Figure 9B. 5 The dependence of the integrals (a) $S$, (b) $j$ and $k$ on the internuclear distance, each calculated for $\mathrm{H}_{2}^{+}$.

The numerical value of $2 h c \tilde{R}_{\infty}$ (when expressed in electronvolts) is 27.21 eV . The integrals are plotted in Fig. 9B.5, and are interpreted as follows:

\begin{itemize}
  \item All three integrals are positive and decline towards zero at large internuclear separations ( $S$ and $k$ on account of the exponential term, $j$ on account of the factor $1 / R$ ). The integral $S$ is discussed in more detail in Topic 9C.
  \item The integral $j$ is a measure of the interaction between a nucleus and electron density centred on the other nucleus.
  \item The integral $k$ is a measure of the interaction between a nucleus and the excess electron density in the internuclear region arising from overlap.
\end{itemize}

\subsection*{Brief illustration 9B. 2}
It turns out (see below) that the minimum value of $E_{\sigma}$ occurs at $R=2.49 a_{0}$. At this separation

$$
\begin{aligned}
& S=\left\{1+2.49+\frac{2.49^{2}}{3}\right\} \mathrm{e}^{-2.49}=0.46 \\
& j=\frac{j_{0} / a_{0}}{2.49}\left\{1-3.49 \mathrm{e}^{-4.98}\right\}=0.39 j_{0} / a_{0} \\
& k=\frac{j_{0}}{a_{0}}(1+2.49) \mathrm{e}^{-2.49}=0.29 j_{0} / a_{0}
\end{aligned}
$$

Therefore, with $j_{0} / a_{0}=27.21 \mathrm{eV}, j=10.7 \mathrm{eV}$, and $k=7.9 \mathrm{eV}$. The energy separation between the bonding MO and the H1s atomic orbital (being cautious with rounding) is $E_{\sigma}-E_{\mathrm{H} 1 \mathrm{~s}}=$ -1.76 eV .

Figure 9B. 6 shows a plot of $E_{\sigma}$ against $R$ relative to the energy of the separated atoms. The energy of the $\sigma$ orbital decreases\\
\includegraphics[max width=\textwidth, center]{2025_02_27_d4b95d9718f0b9e2e6b9g-386(1)}

Figure 9B. 6 The calculated molecular potential energy curves for a hydrogen molecule-ion showing the variation of the energies of the bonding and antibonding orbitals as the internuclear distance is changed. The energy $E_{\sigma}$ is that of the $\sigma$ orbital and $E_{\sigma^{*}}$ is that of $\sigma^{*}$.\\
as the internuclear separation is decreased from large values because electron density accumulates in the internuclear region as the constructive interference between the atomic orbitals increases (Fig. 9B.7). However, at small separations there is too little space between the nuclei for significant accumulation of electron density there. In addition, the nucleus-nucleus repulsion (which is proportional to $1 / R$ ) becomes large. As a result, the energy of the molecular orbital rises at short distances, resulting in a minimum in the potential energy curve of depth $h c \tilde{D}_{\mathrm{e}}$. Calculations on $\mathrm{H}_{2}^{+}$give $R_{\mathrm{e}}=2.49 a_{0}=132 \mathrm{pm}$ and $h c \tilde{D}_{\mathrm{e}}$ $=1.76 \mathrm{eV}\left(171 \mathrm{~kJ} \mathrm{~mol}^{-1}\right)$; the experimental values are 106 pm and 2.6 eV , so this simple LCAO-MO description of the molecule, while inaccurate, is not absurdly wrong.

\subsection*{(c) Antibonding orbitals}
The linear combination $\psi_{-}$in eqn 9B. 2 has higher energy than $\psi_{+}$, and for now it is labelled $\sigma^{*}$ because it is also a $\sigma$ orbital. This orbital has a nodal plane perpendicular to the\\
\includegraphics[max width=\textwidth, center]{2025_02_27_d4b95d9718f0b9e2e6b9g-386}

Figure 9B.7 A representation of the constructive interference that occurs when two H1s orbitals overlap and form a bonding $\sigma$ orbital.\\
\includegraphics[max width=\textwidth, center]{2025_02_27_d4b95d9718f0b9e2e6b9g-387}

Figure 9B. 8 A representation of the destructive interference that occurs when two H1s orbitals overlap and form an antibonding $\sigma$ orbital.\\
\includegraphics[max width=\textwidth, center]{2025_02_27_d4b95d9718f0b9e2e6b9g-387(2)}

Figure 9B. 9 (a) The amplitude of the antibonding molecular orbital in a hydrogen molecule-ion in a plane containing the two nuclei and (b) a contour representation of the amplitude. Note the internuclear nodal plane.\\
internuclear axis and passing through the mid-point of the bond where $\psi_{\mathrm{A}}$ and $\psi_{\mathrm{B}}$ cancel exactly (Figs. 9B. 8 and 9B.9).

The probability density is

\[
\psi_{-}^{2} \propto \psi_{A}^{2}+\psi_{B}^{2}-2 \psi_{A} \psi_{\mathrm{B}} \quad \begin{align*}
& \text { Antibonding }  \tag{9B.6}\\
& \text { probability density }
\end{align*}
\]

There is a reduction in probability density between the nuclei due to the term $-2 \psi_{A} \psi_{B}$ (Fig. 9B.10); in physical terms, there is\\
\includegraphics[max width=\textwidth, center]{2025_02_27_d4b95d9718f0b9e2e6b9g-387(1)}

Figure 9B. 10 The electron density calculated by forming the square of the wavefunction used to construct Fig. 9B.9. Note the reduction of electron density in the internuclear region.\\
\includegraphics[max width=\textwidth, center]{2025_02_27_d4b95d9718f0b9e2e6b9g-387(3)}

Figure 9B.11 A partial explanation of the origin of bonding and antibonding effects. (a) In a bonding orbital, the nuclei are attracted to the accumulation of electron density in the internuclear region. (b) In an antibonding orbital, the nuclei are attracted to an accumulation of electron density outside the internuclear region.\\
destructive interference where the two atomic orbitals overlap. The $\sigma^{*}$ orbital is an example of an antibonding orbital, an orbital that, if occupied, contributes to a reduction in the cohesion between two atoms and helps to raise the energy of the molecule relative to the separated atoms.

The energy $E_{\sigma^{*}}$ of the $\sigma^{\star}$ antibonding orbital is ${ }^{2}$


\begin{equation*}
E_{\sigma^{*}}=E_{\mathrm{H} 1 \mathrm{~s}}+\frac{j_{0}}{R}-\frac{j-k}{1-S} \tag{9B.7}
\end{equation*}


where the integrals $S, j$, and $k$ are the same as in eqn 9B.5. The variation of $E_{\sigma^{*}}$ with $R$ is shown in Fig. 9B.6, which shows the destabilizing effect of an antibonding electron. The effect is partly due to the fact that an antibonding electron is excluded from the internuclear region and hence is distributed largely outside the bonding region. In effect, whereas a bonding electron pulls two nuclei together, an antibonding electron pulls the nuclei apart (Fig. 9B.11). The illustration also shows another feature drawn on later: $\left|E_{\sigma^{*}}-E_{\mathrm{H} 1 \mathrm{~s}}\right|>\left|E_{\sigma}-E_{\mathrm{H} 1 \mathrm{~s}}\right|$, which indicates that the antibonding orbital is more antibonding than the bonding orbital is bonding. This important conclusion stems in part from the presence of the nucleus-nucleus repulsion $\left(j_{0} / R\right)$ : this contribution raises the energy of both molecular orbitals.

\subsection*{Brief illustration 9B. 3}
At the minimum of the bonding orbital energy $R=2.49 a_{0}$, and, from Brief illustration 9B.2, $S=0.46, j=10.7 \mathrm{eV}$, and $k=$ 7.9 eV . It follows that at that separation, the energy of the antibonding orbital relative to that of a hydrogen atom 1 s orbital is

$$
\left(E_{\sigma^{*}}-E_{\mathrm{H} 1 \mathrm{~s}}\right) / \mathrm{eV}=\frac{27.2}{2.49}-\frac{10.7-7.9}{1-0.46}=5.7
$$

That is, the antibonding orbital lies $(5.7+1.76) \mathrm{eV}=7.5 \mathrm{eV}$ above the bonding orbital at this internuclear separation.

\footnotetext{${ }^{2}$ This result is obtained by applying the strategy in A deeper look 4 on the text's website.
}\subsection*{98.2 Orbital notation}
For homonuclear diatomic molecules (molecules consisting of two atoms of the same element, such as $\mathrm{N}_{2}$ ), it proves helpful to label a molecular orbital according to its inversion symmetry, the behaviour of the wavefunction when it is inverted through the centre (more formally, the centre of inversion, Topic 10A) of the molecule. Thus, any point on the bonding $\sigma$ orbital that is projected through the centre of the molecule and out an equal distance on the other side leads to an identical value (and sign) of the wavefunction (Fig. 9B.12). This so-called gerade symmetry (from the German word for 'even') is denoted by a subscript $g$, as in $\sigma_{g}$. The same procedure applied to the antibonding $\sigma^{*}$ orbital results in the same amplitude but opposite sign of the wavefunction. This ungerade symmetry ('odd symmetry') is denoted by a subscript $u$, as in $\sigma_{u}$.\\
\includegraphics[max width=\textwidth, center]{2025_02_27_d4b95d9718f0b9e2e6b9g-388}

Figure 9B. 12 The parity of an orbital is even (g) if its wavefunction is unchanged under inversion through the centre of symmetry of the molecule, but odd ( $u$ ) if the wavefunction changes sign. Heteronuclear diatomic molecules do not have a centre of inversion, so for them the g , u classification is irrelevant.

The inversion symmetry classification is not applicable to heteronuclear diatomic molecules (diatomic molecules formed by atoms from two different elements, such as CO) because these molecules do not have a centre of inversion.

\subsection*{Checklist of concepts}
\begin{enumerate}
  \item A molecular orbital is constructed from a linear combination of atomic orbitals.
  \item A bonding orbital arises from the constructive overlap of neighbouring atomic orbitals.3. An antibonding orbital arises from the destructive overlap of neighbouring atomic orbitals.
  \item $\sigma$ Orbitals have cylindrical symmetry and zero orbital angular momentum around the internuclear axis.
  \item A molecular orbital in a homonuclear diatomic molecule is labelled 'gerade' ( g ) or 'ungerade' ( u ) according to its behaviour under inversion symmetry.
\end{enumerate}

\subsection*{Checklist of equations}
\begin{center}
\begin{tabular}{llll}
\hline
Property & Equation & Comment & Equation number \\
\hline
Linear combination of atomic orbitals & $\psi_{ \pm}=N_{ \pm}\left(\psi_{\mathrm{A}} \pm \psi_{\mathrm{B}}\right)$ & Homonuclear diatomic molecule & 9 B .2 \\
\begin{tabular}{ll}
Energies of $\sigma$ orbitals formed from two 1s &  \\
atomic orbitals & $E_{\sigma}=E_{\mathrm{H} 1 \mathrm{~s}}+j_{0} / R-(j+k) /(1+S)$ \\
 & $E_{\sigma^{*}}=E_{\mathrm{H1s}}+j_{0} / R-(j-k) /(1-S)$ \\
Molecular integrals & $S=\int \psi_{\mathrm{A}} \psi_{\mathrm{B}} \mathrm{d} \tau$, \\
9 B .4 &  \\
 & $j=j_{0} \int\left(\psi_{\mathrm{A}}^{2} / r_{\mathrm{B}}\right) \mathrm{d} \tau$ \\
 & $k=j_{0} \int\left(\psi_{\mathrm{A}} \psi_{\mathrm{B}} / r_{\mathrm{B}}\right) \mathrm{d} \tau$ \\
\end{tabular} & 9 B .7 &  &  \\
\hline
\end{tabular}
\end{center}

\section*{TOPIC 9C Molecular orbital theory: homonuclear diatomic molecules }
\textbackslash begin\{abstract\}\\
> Why do you need to know this material?\\
Almost all chemically significant molecules have more than one electron, so you need to see how to construct their electron configurations. This Topic shows how to use molecular orbital theory when more than one electron is present in a molecule.\\
\textbackslash end\{abstract\}

\subsection*{What is the key idea?}
Each molecular orbital can accommodate up to two electrons, and the ground state of the molecule is the configuration of lowest energy.

\subsection*{> What do you need to know already?}
You need to be familiar with the discussion of the bonding and antibonding linear combinations of atomic orbitals in Topic 9B and the building-up principle for atoms (Topic 8B).

Just as hydrogenic atomic orbitals and the building-up principle can be used as a basis for the discussion and prediction of the ground electronic configurations of many-electron atoms, the molecular orbitals for the one-electron hydrogen mol-ecule-ion introduced in Topic 9B and a version of the build-ing-up principle introduced in Topic 8B can be developed to account for the configurations of many-electron diatomic molecules and ions.

\subsection*{9C. 1 Electron configurations}
The starting point of the molecular orbital theory (MO theory) of bonding in diatomic molecules (and ions) is the construction of molecular orbitals as linear combinations of the available atomic orbitals. Once the molecular orbitals have been formed, a building-up principle, like that for atoms, can\\
be used to establish their ground-state electron configurations (Topic 8B):

\begin{itemize}
  \item The electrons supplied by the atoms are accommodated in the molecular orbitals so as to achieve the lowest overall energy subject to the constraint of the Pauli exclusion principle that no more than two electrons may occupy a single orbital (and then their spins must be paired).
  \item If several degenerate molecular orbitals are available, electrons are added singly to each individual orbital before any one orbital is completed (because that minimizes electron-electron repulsions).
  \item According to Hund's maximum multiplicity rule (Topic 8B), if two electrons do occupy different degenerate orbitals, then a lower energy is obtained if their spins are parallel.
\end{itemize}

\subsection*{(a) $\sigma$ Orbitals and $\pi$ orbitals}
Consider $\mathrm{H}_{2}$, the simplest many-electron diatomic molecule. Each H atom contributes a 1s orbital (as in $\mathrm{H}_{2}^{+}$), which combine to form bonding $\sigma$ and antibonding $\sigma^{\star}$ orbitals, as explained in Topic 9B. At the equilibrium nuclear separation these orbitals have the energies shown in Fig. 9C.1, which is called a\\
\includegraphics[max width=\textwidth, center]{2025_02_27_d4b95d9718f0b9e2e6b9g-389}

Figure 9C. 1 A molecular orbital energy level diagram for orbitals constructed from the overlap of H 1 s orbitals. The energies of the atomic orbitals are indicated by the lines at the outer edges of the diagram, and the energies of the molecular orbitals are shown in the middle. The ground electronic configuration of $\mathrm{H}_{2}$ is obtained by accommodating the two electrons in the lowest available orbital (the bonding orbital).\\
\includegraphics[max width=\textwidth, center]{2025_02_27_d4b95d9718f0b9e2e6b9g-390}

Figure 9C. 2 The ground-state electronic configuration of the hypothetical four-electron molecule $\mathrm{He}_{2}$ (at an arbitrary internuclear separation) has two bonding electrons and two antibonding electrons. It has a higher energy than the separated atoms, and so is unstable.\\
molecular orbital energy level diagram. Note that from two atomic orbitals two molecular orbitals are built. In general, from $N$ atomic orbitals $N$ molecular orbitals can be built.

There are two electrons to accommodate, and both can enter the $\sigma$ orbital by pairing their spins, as required by the Pauli principle (just as for atoms, Topic 8B). The ground-state configuration is therefore $\sigma^{2}$ and the bond consists of an electron pair in a bonding $\sigma$ orbital. This approach shows that an electron pair, which was the focus of Lewis's account of chemical bonding, represents the maximum number of electrons that can enter a bonding molecular orbital.

A straightforward extension of this argument explains why helium does not form diatomic molecules. Each He atom contributes a 1 s orbital, so $\sigma$ and $\sigma^{*}$ molecular orbitals can be constructed. Although these orbitals differ in detail from those in $\mathrm{H}_{2}$, their general shapes are the same and the same qualitative energy level diagram can be used in the discussion. There are four electrons to accommodate. Two can enter the $\sigma$ orbital, but then it is full, and the next two must enter the $\sigma^{*}$ orbital (Fig. 9C.2). The ground electronic configuration of $\mathrm{He}_{2}$ is therefore $\sigma^{2} \sigma^{\star 2}$. Because $\sigma^{\star}$ lies higher in energy above the separate atoms more than $\sigma$ lies below them, an $\mathrm{He}_{2}$ molecule has a higher energy than the separated atoms, so it is unstable relative to them and dihelium does not form.

The concepts introduced so far also apply to homonuclear diatomics in general. In the elementary treatment used here, only the orbitals of the valence shell are used to form molecular orbitals so, for molecules formed with atoms from Period 2 elements, only the 2 s and 2 p atomic orbitals are considered.\\
A general principle of MO theory is that

\subsection*{All orbitals of the appropriate symmetry contribute to a}
 molecular orbital.Thus, $\sigma$ orbitals are built by forming linear combinations of all atomic orbitals that have cylindrical symmetry about the internuclear axis. These orbitals include the 2 s orbitals on each atom and the $2 \mathrm{p}_{z}$ orbitals on the two atoms (Fig. 9C.3; the $z$ axis on each atom lies along the internuclear axis and points\\
\includegraphics[max width=\textwidth, center]{2025_02_27_d4b95d9718f0b9e2e6b9g-390(1)}

Figure 9C. 3 According to molecular orbital theory, $\sigma$ orbitals are built from all orbitals that have the appropriate symmetry. In homonuclear diatomic molecules of Period 2, that means that two 2 s and two $2 p_{z}$ orbitals should be used. From these four orbitals, four molecular orbitals can be built.\\
towards the neighbouring atom). The general form of the $\sigma$ orbitals that may be formed is therefore


\begin{equation*}
\psi=c_{\mathrm{A} 2 \mathrm{~s}} \psi_{\mathrm{A} 2 \mathrm{~s}}+c_{\mathrm{B} 2 \mathrm{~s}} \psi_{\mathrm{B} 2 \mathrm{~s}}+c_{\mathrm{A} 2 \mathrm{p}_{z}} \psi_{\mathrm{A} 2 \mathrm{p}_{z}}+c_{\mathrm{B} 2 \mathrm{p}_{z}} \psi_{\mathrm{B} 2 \mathrm{p}_{z}} \tag{9С.1}
\end{equation*}


From these four atomic orbitals four molecular orbitals of $\sigma$ symmetry can be formed by an appropriate choice of the coefficients $c$.

Because the 2 s and 2 p orbitals on each atom have such different energies, they may be treated separately (this approximation is removed later). That is, the four $\sigma$ molecular orbitals fall approximately into two sets, one consisting of two molecular orbitals formed from the 2 s orbitals


\begin{equation*}
\psi=c_{\mathrm{A} 2 \mathrm{~s}} \psi_{\mathrm{A} 2 \mathrm{~s}}+c_{\mathrm{B} 2 \mathrm{~s}} \psi_{\mathrm{B} 2 \mathrm{~s}} \tag{9С.2a}
\end{equation*}


and another consisting of two orbitals formed from the $2 \mathrm{p}_{z}$ orbitals


\begin{equation*}
\psi=c_{\mathrm{A} 2 \mathrm{p}_{z}} \psi_{{\mathrm{A} 2 \mathrm{p}_{z}}}+c_{\mathrm{B} 2 \mathrm{p}_{z}} \psi_{\mathrm{B} 2 \mathrm{p}_{z}} \tag{9C.2b}
\end{equation*}


In a homonuclear diatomic molecule the energies of the 2 s orbitals on atoms A and B are the same. Their coefficients are therefore equal (apart from a possible difference in sign). The same is true of the $2 \mathrm{p}_{z}$ orbitals on each atom. Therefore, the two\\
\includegraphics[max width=\textwidth]{2025_02_27_d4b95d9718f0b9e2e6b9g-390(2)} combination being bonding and the - combination antibonding in each case.

At this stage it is useful to adopt a more formal system for denoting molecular orbitals. First, the orbitals are labelled with $g$ and $u$ to indicate their inversion symmetry, as explained in Topic 9B. Then each set of orbitals of the same inversion symmetry is numbered separately. Therefore, the $\sigma$ orbital formed from the 2 s orbitals is labelled $1 \sigma_{\mathrm{g}}$ and the $\sigma^{*}$ orbital formed from the same atomic orbitals is denoted $1 \sigma_{u}$.\\
The two $2 \mathrm{p}_{z}$ orbitals directed along the internuclear axis also overlap strongly. They may interfere either constructively or destructively, and give a bonding or antibonding $\sigma$ orbital that lie higher in energy than the $1 \sigma_{g}$ and $1 \sigma_{u}$ orbitals because it has been supposed that the 2 p atomic orbitals lie significantly higher in energy than the 2 s orbitals (Fig. 9C.4). These two $\sigma$ orbitals are labelled $2 \sigma_{\mathrm{g}}$ and $2 \sigma_{\mathrm{u}}$, respectively. Note how the numbering follows the order of increasing energy and orbitals of different symmetry are labelled separately.\\
\includegraphics[max width=\textwidth, center]{2025_02_27_d4b95d9718f0b9e2e6b9g-391}

Figure 9C. 4 A representation of the form of the bonding and antibonding $\sigma$ orbitals built from the overlap of $p$ orbitals. These illustrations are schematic.

Now consider the $2 \mathrm{p}_{x}$ and $2 \mathrm{p}_{y}$ orbitals of each atom. These orbitals are perpendicular to the internuclear axis and overlap broadside-on when the atoms are close together. This overlap may be constructive or destructive and results in a bonding or an antibonding $\pi$ orbital (Fig. 9C.5). The notation $\pi$ is the analogue of p in atoms: when viewed along the axis of the molecule, a $\pi$ orbital looks like a p orbital and has one unit of orbital angular momentum around the internuclear axis. The two neighbouring $2 \mathrm{p}_{x}$ orbitals overlap to give a bonding and antibonding $\pi_{x}$ orbital, and the two $2 \mathrm{p}_{y}$ orbitals overlap to give two $\pi_{y}$ orbitals. The $\pi_{x}$ and $\pi_{y}$ bonding orbitals are degenerate; so too are their antibonding partners. As seen in Fig. 9C.5, a bonding $\pi$ orbital has odd parity ( $u$ ) and the antibonding $\pi$ orbital has even parity (g). The lower two doubly degenerate orbitals are therefore labelled $1 \pi_{\mathrm{u}}$ and their higher energy antibonding partners are labelled $1 \pi_{\mathrm{g}}$.

\subsection*{(b) The overlap integral}
As in the discussion of the hydrogen molecule-ion, the lowering of energy that results from constructive interference between neighbouring atomic orbitals (and the raising of energy that results from destructive interference) correlates with the extent of overlap of the orbitals. As explained in Topic 9B, the extent to which two atomic orbitals overlap is measured by the overlap integral, $S$ :

$$
S=\int \psi_{\mathrm{A}}^{*} \psi_{\mathrm{B}} \mathrm{~d} \tau
$$

Overlap integral [definition]\\
(9С.3)\\
\includegraphics[max width=\textwidth, center]{2025_02_27_d4b95d9718f0b9e2e6b9g-391(1)}

Figure 9C. 5 The parity of $\pi$ bonding and antibonding molecular orbitals.\\
\includegraphics[max width=\textwidth, center]{2025_02_27_d4b95d9718f0b9e2e6b9g-391(2)}

Figure 9C. 6 (a) When two orbitals are on atoms that are far apart, the wavefunctions are small where they overlap, so $S$ is small.\\
(b) When the atoms are closer, both orbitals have significant amplitudes where they overlap, and $S$ may approach 1 . Note that $S$ will decrease again as the two atoms approach more closely than shown here, because the region of negative amplitude of the $p$ orbital starts to overlap the positive amplitude of the $s$ orbital. When the centres of the atoms coincide, $S=0$.

If the atomic orbital $\psi_{\mathrm{A}}$ on A is small wherever the orbital $\psi_{\mathrm{B}}$ on B is large, or vice versa, then the product of their amplitudes is everywhere small and the integral-the sum of these products-is small (Fig. 9C.6). If $\psi_{\mathrm{A}}$ and $\psi_{\mathrm{B}}$ are both large in some region of space, then $S$ may approach 1 . If the two normalized atomic orbitals are identical (for instance, 1 s orbitals on the same nucleus), then $S=1$. In some cases, simple formulas can be given for overlap integrals (Table 9C.1) and illustrated in Fig. 9C.7.

Now consider the arrangement in which an s orbital spreads into the same region of space as a $p_{x}$ orbital of a different atom (Fig. 9C.8). The integral over the region where the product of the wavefunctions is positive exactly cancels the integral over the region where the product is negative, so overall $S=0$ exactly. Therefore, there is no net overlap between the $s$ and $\mathrm{p}_{x}$ orbitals in this arrangement.

The extent of overlap as measured by the overlap integral is suggestive of the contribution that different kinds of orbital overlap makes to bond formation, but the value of the integral must be treated with caution. Thus, the overlap integral for broadside overlap of $2 \mathrm{p}_{x}$ or $2 \mathrm{p}_{y}$ orbitals is typically greater than that for the overlap of $2 \mathrm{p}_{z}$ orbitals, suggesting weaker $\sigma$ than $\pi$ bonding.

Table 9C. 1 Overlap integrals between hydrogenic orbitals

\begin{center}
\begin{tabular}{ll}
\hline
Orbitals & Overlap integral \\
$1 \mathrm{~s}, 1 \mathrm{~s}$ & $S=\left\{1+\frac{Z R}{a_{0}}+\frac{1}{3}\left(\frac{Z R}{a_{0}}\right)^{2}\right\} \mathrm{e}^{-Z R / a_{0}}$ \\
$2 \mathrm{~s}, 2 \mathrm{~s}$ & $S=\left\{1+\frac{Z R}{2 a_{0}}+\frac{1}{12}\left(\frac{Z R}{a_{0}}\right)^{2}+\frac{1}{240}\left(\frac{Z R}{a_{0}}\right)^{4}\right\} \mathrm{e}^{-Z \mathrm{R} / 2 a_{0}}$ \\
$2 \mathrm{p}_{x}, 2 \mathrm{p}_{x}(\pi)$ & $S=\left\{1+\frac{Z R}{2 a_{0}}+\frac{1}{10}\left(\frac{Z R}{a_{0}}\right)^{2}+\frac{1}{120}\left(\frac{Z R}{a_{0}}\right)^{3}\right\} \mathrm{e}^{-Z R / 2 a_{0}}$ \\
$2 \mathrm{p}_{z}, 2 \mathrm{p}_{z}(\sigma)$ & $S=-\left\{1+\frac{Z R}{2 a_{0}}+\frac{1}{20}\left(\frac{Z R}{a_{0}}\right)^{2}-\frac{1}{60}\left(\frac{Z R}{a_{0}}\right)^{3}-\frac{1}{240}\left(\frac{Z R}{a_{0}}\right)^{4}\right\} \mathrm{e}^{-Z R / 2 a_{0}}$ \\
\end{tabular}
\end{center}

\begin{center}
\includegraphics[max width=\textwidth]{2025_02_27_d4b95d9718f0b9e2e6b9g-392(3)}
\end{center}

Figure 9C. 7 The variation of the overlap integral, $S$, between two hydrogenic orbitals with the internuclear separation. A negative value of $S$ corresponds to separations at which the contribution to the overlap of the positive region of one $2 p$ orbital with the negative lobe of the other $2 p$ orbital outweighs that from the regions where both have the same sign.

However, the constructive overlap in the region between the nuclei and on the axis is greater in $\sigma$ interactions, and its effect on bonding is more important than the overall extent of overlap. As a result, the separation of $1 \pi_{\mathrm{u}}$ and $1 \pi_{\mathrm{g}}$ orbitals is likely to be smaller than the separation of $2 \sigma_{g}$ and $2 \sigma_{u}$ orbitals in the same molecule. The relative energies of these orbitals is therefore likely to be as shown in Fig. 9C.9, and electrons occupying $\pi$ orbitals are likely to be less effective at bonding than those occupying the $\sigma$ orbitals derived from the same p orbitals.\\
\includegraphics[max width=\textwidth, center]{2025_02_27_d4b95d9718f0b9e2e6b9g-392(1)}

Figure 9C. 8 A p orbital in the orientation shown here has zero net overlap ( $S=0$ ) with the s orbital at all internuclear separations.\\
\includegraphics[max width=\textwidth, center]{2025_02_27_d4b95d9718f0b9e2e6b9g-392}

Figure 9C.9 As explained in the text, the separation of $1 \pi_{u}$ and $1 \pi_{\mathrm{g}}$ orbitals is likely to be smaller than the separation of $2 \sigma_{\mathrm{g}}$ and $2 \sigma_{u}$ orbitals in the same molecule, leading to the relative energies shown here.

\subsection*{(c) Period 2 diatomic molecules}
To construct the molecular orbital energy level diagram for Period 2 homonuclear diatomic molecules, eight molecular orbitals are formed from the eight valence shell orbitals (four from each atom). The ordering suggested by the extent of overlap is shown in Fig. 9C.10. However, remember that this scheme assumes that the 2 s and $2 \mathrm{p}_{z}$ orbitals contribute to different sets of molecular orbitals. In fact all four atomic orbitals have the same symmetry around the internuclear axis and contribute jointly to the four $\sigma$ orbitals. Hence, there is no guarantee that this order of energies will be found, and detailed calculation shows that the order varies along Period 2 (Fig. 9C.11). The order shown in Fig. 9C. 12 is appropriate as far as $\mathrm{N}_{2}$, and Fig. 9 C .10 is appropriate for $\mathrm{O}_{2}$ and $\mathrm{F}_{2}$. The relative order is controlled by the energy separation of the 2 s and\\
\includegraphics[max width=\textwidth, center]{2025_02_27_d4b95d9718f0b9e2e6b9g-392(2)}

Figure 9C. 10 The molecular orbital energy level diagram for homonuclear diatomic molecules. The lines in the middle are an indication of the energies of the molecular orbitals that can be formed by overlap of atomic orbitals. Energy increases upwards. As remarked in the text, this diagram is appropriate for $\mathrm{O}_{2}$ (the configuration shown) and $\mathrm{F}_{2}$.\\
\includegraphics[max width=\textwidth, center]{2025_02_27_d4b95d9718f0b9e2e6b9g-393(1)}

Figure 9C. 11 The variation of the orbital energies of Period 2 homonuclear diatomics.\\
\includegraphics[max width=\textwidth, center]{2025_02_27_d4b95d9718f0b9e2e6b9g-393}

Figure 9C. 12 An alternative molecular orbital energy level diagram for homonuclear diatomic molecules. Energy increases upwards. As remarked in the text, this diagram is appropriate for Period 2 homonuclear diatomics up to and including $\mathrm{N}_{2}$ (the configuration shown).

2 p orbitals in the atoms, which increases across the period. The change in the order of the $1 \pi_{\mathrm{u}}$ and $2 \sigma_{\mathrm{g}}$ orbitals occurs at about $\mathrm{N}_{2}$.

With the molecular orbital energy level diagram established, the probable ground-state configurations of the molecules are deduced by adding the appropriate number of electrons to the orbitals and following the building-up rules. Anionic species (such as the peroxide ion, $\mathrm{O}_{2}^{2-}$ ) need more electrons than the parent neutral molecules; cationic species (such as $\mathrm{O}_{2}^{+}$) need fewer.

Consider $\mathrm{N}_{2}$, which has 10 valence electrons. Two electrons pair, occupy, and fill the $1 \sigma_{\mathrm{g}}$ orbital; the next two occupy and fill the $1 \sigma_{u}$ orbital. Six electrons remain. There are two $1 \pi_{u}$ orbitals, so four electrons can be accommodated in them. The last two enter the $2 \sigma_{\mathrm{g}}$ orbital. Therefore, the ground-state configuration of $\mathrm{N}_{2}$ is $1 \sigma_{\mathrm{g}}^{2} 1 \sigma_{\mathrm{u}}^{2} 1 \pi_{\mathrm{u}}^{4} 2 \sigma_{\mathrm{g}}^{2}$. It is sometimes helpful to include an asterisk to denote an antibonding orbital, in which case this configuration would be denoted $1 \sigma_{\mathrm{g}}^{2} 1 \sigma_{\mathrm{u}}^{\not 22} 1 \pi_{\mathrm{u}}^{4} 2 \sigma_{\mathrm{g}}^{2}$.

A measure of the net bonding in a diatomic molecule is its bond order, $b$ :

$$
\begin{array}{ll}
b=\frac{1}{2}\left(N-N^{\star}\right) & \text { Bond order } \\
\text { [definition] }
\end{array}
$$

(9С.4)\\
where $N$ is the number of electrons in bonding orbitals and $N^{*}$ is the number of electrons in antibonding orbitals.

\subsection*{Brief illustration 9C. 1}
Each electron pair in a bonding orbital increases the bond order by 1 and each pair in an antibonding orbital decreases $b$ by 1 . For $\mathrm{H}_{2}, b=1$, corresponding to a single bond, $\mathrm{H}-\mathrm{H}$, between the two atoms. In $\mathrm{He}_{2}, b=0$, and there is no bond. In $\mathrm{N}_{2}, b=\frac{1}{2}(8-2)=3$. This bond order accords with the Lewis structure of the molecule (: $\mathrm{N} \equiv \mathrm{N}$ :).

The ground-state electron configuration of $\mathrm{O}_{2}$, with 12 valence electrons, is based on Fig. 9C.10, and is $1 \sigma_{\mathrm{g}}^{2} 1 \sigma_{\mathrm{u}}^{2} 2 \sigma_{\mathrm{g}}^{2} 1 \pi_{\mathrm{u}}^{4} 1 \pi_{\mathrm{g}}^{2}$ (or $1 \sigma_{\mathrm{g}}^{2} 1 \sigma_{\mathrm{u}}^{\not 2} 2 \sigma_{\mathrm{g}}^{2} 1 \pi_{\mathrm{u}}^{4} 1 \pi_{\mathrm{g}}^{\not{ }^{2}}$ ). The bond order is $b=\frac{1}{2}(8-4)=2$. According to the building-up principle, however, the two $1 \pi_{\mathrm{g}}$ electrons occupy two different orbitals: one will enter $1 \pi_{\mathrm{g}, x}$ and the other will enter $1 \pi_{g, y}$. Because the electrons are in different orbitals, they will have parallel spins. Therefore, an $\mathrm{O}_{2}$ molecule is predicted to have a net spin angular momentum with $S=1$ and, in the language introduced in Topic 8C, to be in a triplet state. As electron spin is the source of a magnetic moment, oxygen is also predicted to be paramagnetic, a substance that tends to be drawn into a magnetic field (see Topic 15C). This prediction, which VB theory does not make, is confirmed by experiment.

An $\mathrm{F}_{2}$ molecule has two more electrons than an $\mathrm{O}_{2}$ molecule. Its configuration is therefore $1 \sigma_{\mathrm{g}}^{2} 1 \sigma_{\mathrm{u}}^{\star^{2}} 2 \sigma_{\mathrm{g}}^{2} 1 \pi_{\mathrm{u}}^{4} 1 \pi_{\mathrm{g}}^{\star^{4}}$ and $b=1$, so $F_{2}$ is a singly-bonded molecule, in agreement with its Lewis structure. The hypothetical molecule dineon, $\mathrm{Ne}_{2}$, has two more electrons than $F_{2}$ : its configuration is $1 \sigma_{\mathrm{g}}^{2} 1 \sigma_{\mathrm{u}}^{\not 22} 2 \sigma_{\mathrm{g}}^{2} 1 \pi_{\mathrm{u}}^{4} 1 \pi_{\mathrm{g}}^{\star^{4}} 2 \sigma_{\mathrm{u}}^{\not{ }^{2}}$ and $b=0$. The zero bond order is consistent with the fact that neon occurs as a monatomic gas.

The bond order is a useful parameter for discussing the characteristics of bonds, because it correlates with bond length and bond strength. For bonds between atoms of a given pair of elements:

\begin{itemize}
  \item The greater the bond order, the shorter the bond.
  \item The greater the bond order, the greater the bond strength.
\end{itemize}

Table 9C. 2 lists some typical bond lengths in diatomic and polyatomic molecules. The strength of a bond is measured by its bond dissociation energy, $h c \tilde{D}_{0}$, the energy required to separate the atoms to infinity or by the well depth, $h c \tilde{D}_{\mathrm{e}}$, with $h c \tilde{D}_{0}=h c \tilde{D}_{\mathrm{e}}-$ $\frac{1}{2} \hbar \omega$. Table 9C. 3 lists some experimental values of $h c \tilde{D}_{0}$.

Table 9C. 2 Bond lengths*

\begin{center}
\begin{tabular}{lll}
\hline
Bond & Order & $R_{\mathrm{e}} / \mathrm{pm}$ \\
\hline
HH & 1 & 74.14 \\
NN & 3 & 109.76 \\
HCl & 1 & 127.45 \\
CH & 1 & 114 \\
CC & 1 & 154 \\
 & 2 & 134 \\
 & 3 & 120 \\
\end{tabular}
\end{center}

\begin{itemize}
  \item More values will be found in the Resource section. Numbers in italics are mean values for polyatomic molecules.
\end{itemize}

Table 9C. 3 Bond dissociation energies*

\begin{center}
\begin{tabular}{lll}
\hline
Bond & Order & $N_{\mathrm{A}} h c \tilde{D}_{0} /(\mathrm{kJ} \mathrm{mol}$ \\
 &  &  \\
\hline
$1 \mathbf{1})$ &  &  \\
\hline
HH & 1 & 432.1 \\
NN & 3 & 941.7 \\
HCl & 1 & 427.7 \\
CH & 1 & 435 \\
CC & 1 & 368 \\
 & 2 & 720 \\
 & 3 & 962 \\
\end{tabular}
\end{center}

\begin{itemize}
  \item More values will be found in the Resource section. Numbers in italics are mean values for polyatomic molecules.
\end{itemize}

\subsection*{Brief illustration 9C. 2}
From Fig. 9C.12, the electron configurations and bond orders of $\mathrm{N}_{2}$ and $\mathrm{N}_{2}{ }^{+}$are

$$
\begin{array}{lll}
\mathrm{N}_{2} & 1 \sigma_{\mathrm{g}}^{2} 1 \sigma_{\mathrm{u}}^{\not 2} 1 \pi_{\mathrm{u}}^{4} 2 \sigma_{\mathrm{g}}^{2} & b=3 \\
\mathrm{~N}_{2}^{+} & 1 \sigma_{\mathrm{g}}^{2} 1 \sigma_{\mathrm{u}}^{\not 2} 1 \pi_{\mathrm{u}}^{4} 2 \sigma_{\mathrm{g}}^{1} & b=2 \frac{1}{2}
\end{array}
$$

Because the cation has the smaller bond order, you should expect it to have the smaller dissociation energy. The experimental dissociation energies are $945 \mathrm{~kJ} \mathrm{~mol}^{-1}$ for $\mathrm{N}_{2}$ and $842 \mathrm{~kJ} \mathrm{~mol}^{-1}$ for $\mathrm{N}_{2}^{+}$.

\subsection*{9c.2 Photoelectron spectroscopy}
So far, molecular orbitals have been regarded as purely theoretical constructs, but is there experimental evidence for their existence? Photoelectron spectroscopy (PES) measures the ionization energies of molecules when electrons are ejected from different orbitals by absorption of a photon of known energy, and uses the information to infer the energies of molecular orbitals.

Because energy is conserved when a photon ionizes a sample, the sum of the ionization energy, $I$, of the sample and the\\
\includegraphics[max width=\textwidth, center]{2025_02_27_d4b95d9718f0b9e2e6b9g-394}

Figure 9C. 13 An incoming photon carries an energy $h v$; an energy $l_{i}$ is needed to remove an electron from an orbital $i$, and the difference appears as the kinetic energy of the electron.\\
kinetic energy of the photoelectron, the ejected electron, must be equal to the energy of the incident photon $h v$ (Fig. 9C.13):


\begin{equation*}
h \nu=\frac{1}{2} m_{\mathrm{e}} \nu^{2}+I \tag{9C.5}
\end{equation*}


This equation can be refined in two ways. First, photoelectrons may originate from one of a number of different orbitals, and each one has a different ionization energy. Hence, a series of photoelectrons with different kinetic energies will be obtained, each one satisfying $h v=\frac{1}{2} m_{\mathrm{e}} v^{2}+I_{i}$, where $I_{i}$ is the ionization energy for ejection of an electron from an orbital $i$. Therefore, by measuring the kinetic energies of the photoelectrons, and knowing the frequency $v$, these ionization energies can be determined. Photoelectron spectra are interpreted in terms of an approximation called Koopmans' theorem, which states that the ionization energy $I_{i}$ is equal to the orbital energy of the ejected electron (formally: $I_{i}=-\varepsilon_{i}$ ). That is, the ionization energy can be identified with the energy of the orbital from which it is ejected. The theorem is only an approximation because it ignores the fact that the remaining electrons adjust their distributions when ionization occurs.

The ionization energies of molecules are several electronvolts even for valence electrons, so it is essential to work in at least the ultraviolet region of the spectrum and with wavelengths of less than about 200 nm . Much work has been done with radiation generated by a discharge through helium: the $\mathrm{He}(\mathrm{I})$ line $\left(1 \mathrm{~s}^{1} 2 \mathrm{p}^{1} \rightarrow 1 \mathrm{~s}^{2}\right)$ lies at 58.43 nm , corresponding to a photon energy of 21.22 eV . Its use gives rise to the technique of ultraviolet photoelectron spectroscopy (UPS). When core electrons are being studied, photons of even higher energy are needed to expel them: X-rays are used, and the technique is denoted XPS.

The kinetic energies of the photoelectrons are measured using an electrostatic deflector that produces different deflections in the paths of the photoelectrons as they pass between charged plates (Fig. 9C.14). As the field strength between the\\
\includegraphics[max width=\textwidth, center]{2025_02_27_d4b95d9718f0b9e2e6b9g-395}

Figure 9C. 14 A photoelectron spectrometer consists of a source of ionizing radiation (such as a helium discharge lamp for UPS and an X-ray source for XPS), an electrostatic analyser, and an electron detector. The deflection of the electron path caused by the analyser depends on the speed of the electrons.\\
plates is increased, electrons of different speeds, and therefore kinetic energies, reach the detector. The electron flux can be recorded and plotted against kinetic energy to obtain the photoelectron spectrum (Fig. 9C.15).

\subsection*{Brief illustration 9C. 3}
The photoelectrons of highest kinetic energy ejected from $\mathrm{N}_{2}$ in a spectrometer using $\mathrm{He}(\mathrm{I})$ radiation have kinetic energies of 5.63 eV . Because photons of helium(I) radiation have energy 21.22 eV it follows that $21.22 \mathrm{eV}=5.63 \mathrm{eV}+I_{i}$, so $I_{i}=$ 15.59 eV . This ionization energy is the energy needed to\\
\includegraphics[max width=\textwidth, center]{2025_02_27_d4b95d9718f0b9e2e6b9g-395(1)}

Figure 9C. 15 The photoelectron spectrum of $\mathrm{N}_{2}$ recorded using $\mathrm{He}(\mathrm{I})$ radiation.\\
remove an electron from the occupied molecular orbital with the highest energy of the $\mathrm{N}_{2}$ molecule, the $2 \sigma_{\mathrm{g}}$ bonding orbital. Photoelectrons are also detected at 4.53 eV , corresponding to an ionization energy of 16.7 eV . The likely origin of these electrons is the $1 \pi_{u}$ orbital.

Photoejection commonly results in cations that are excited vibrationally. Because different energies are needed to excite different vibrational states of the ion, the photoelectrons appear with different kinetic energies. The result is vibrational fine structure, a progression of lines with a spacing in energy that corresponds to the vibrational frequency of the molecular ion. This fine structure occurs between 16.7 eV and 18 eV in the photoelectron spectrum of $\mathrm{N}_{2}$ shown in Fig. 9C.15.

\subsection*{Checklist of concepts}
1. Molecular orbitals are constructed as linear combinations of all valence orbitals of the appropriate symmetry.2. As a first approximation, $\sigma$ orbitals are constructed separately from valence $s$ and $p$ orbitals.3. $\pi$ Orbitals are constructed from the sideways overlap of p orbitals.4. An overlap integral is a measure of the extent of orbital overlap.5. According to the building-up principle, electrons occupy the available molecular orbitals so as to achievethe lowest total energy subject to the Pauli exclusion principle.6. If electrons occupy different orbitals, the lowest energy is obtained if their spins are parallel.\\
7. The greater the bond order of a molecule or ion between the same two atoms, the shorter and stronger is the bond.8. Photoelectron spectroscopy is a technique for determining the energies of electrons in molecular orbitals.

\subsection*{Checklist of equations}
\begin{center}
\begin{tabular}{|c|c|c|c|}
\hline
Property & Equation & Comment & Equation number \\
\hline
Overlap integral & $S=\int \psi_{\mathrm{A}}^{*} \psi_{\mathrm{B}} \mathrm{d} \tau$ & Integration over all space & 9C. 3 \\
\hline
Bond order & $b=\frac{1}{2}\left(N-N^{\star}\right)$ & $N$ and $N^{*}$ are the numbers of electrons in bonding and antibonding orbitals, respectively & 9C. 4 \\
\hline
Photoelectron spectroscopy & $h \nu=\frac{1}{2} m_{e} \nu^{2}+I$ & Interpret $I$ as $I_{i}$, the ionization energy from orbital $i$ & 9C. 5 \\
\hline
\end{tabular}
\end{center}

\section*{TOPIC 9D Molecular orbital theory: heteronuclear diatomic molecules }
\subsection*{Why do you need to know this material?}
Most diatomic molecules are heteronuclear, so you need to appreciate the differences in their electronic structure from homonuclear species, and how to treat those differences quantitatively.

\subsection*{- What is the key idea?}
The bonding molecular orbital of a heteronuclear diatomic molecule is composed mostly of the atomic orbital of the more electronegative atom; the opposite is true of the antibonding orbital.

\subsection*{> What do you need to know already?}
You need to know about the molecular orbitals of homonuclear diatomic molecules (Topic 9C) and the concepts of normalization and orthogonality (Topic 7C). This Topic makes light use of determinants (The chemist's toolkit 23) and the rules of differentiation (The chemist's toolkit 5 in Topic 1C).

The electrons in a covalent bond in a heteronuclear diatomic species are not distributed equally over the atoms because it is energetically favourable for the electron pair to be found closer to one atom than to the other. This imbalance results in a polar bond, a bond in which the bonding electron density is shared unequally between the bonded atoms. The bond in HF, for instance, is polar, with the bonding electron density greater near the F atom than the H atom. The accumulation of bonding electron density near the F atom results in that atom having a net negative charge, which is called a partial negative charge and denoted $\delta$-. There is a matching partial positive charge, $\delta+$, on the H atom (Fig. 9D.1).

\subsection*{9D. 1 Polar bonds and electronegativity}
The description of polar bonds is a straightforward extension of the molecular orbital theory of homonuclear diatomic molecules (Topic 9C). The principal difference is that the atomic\\
\includegraphics[max width=\textwidth, center]{2025_02_27_d4b95d9718f0b9e2e6b9g-397}

Figure 9D. 1 The electron density of the molecule HF, computed using one of the methods described in Topic 9E. Different colours show the variation in the electrostatic potential and hence the net charge, with blue representing the region with largest partial positive charge, and red the region with largest partial negative charge.\\
orbitals on the two atoms have different energies and spatial extents.

A polar bond consists of two electrons in a bonding molecular orbital of the form


\begin{equation*}
\psi=c_{\mathrm{A}} \psi_{\mathrm{A}}+c_{\mathrm{B}} \psi_{\mathrm{B}} \quad \text { Wavefunction of a polar bond } \tag{9D.1}
\end{equation*}


with unequal coefficients. It follows that the contribution of the atomic orbital $\psi_{A}$ to the bond, in the sense that on inspection of the location of the electron the probability that it will be found on atom A , is $\left|c_{\mathrm{A}}\right|^{2}$, and that of $\psi_{\mathrm{B}}$ is $\left|c_{\mathrm{B}}\right|^{2}$. A nonpolar bond has $\left|c_{A}\right|^{2}=\left|c_{B}\right|^{2}$, and a pure ionic bond has one coefficient equal to zero (so the species $\mathrm{A}^{+} \mathrm{B}^{-}$would have $c_{\mathrm{A}}=0$ and $c_{\mathrm{B}}=1$ ). The atomic orbital with the lower energy makes the larger contribution to the bonding molecular orbital. The opposite is true of the antibonding orbital, for which the dominant component comes from the atomic orbital with higher energy.

The distribution of partial charges in bonds is commonly discussed in terms of the electronegativity, $\chi$ (chi), of the elements involved. The electronegativity is a parameter introduced by Linus Pauling as a measure of the power of an atom in a bond to attract electrons to itself. Pauling used valencebond arguments to suggest that an appropriate numerical scale of electronegativities could be defined in terms of bond dissociation energies, $h c \tilde{D}_{0}$, and proposed that the difference in electronegativities could be expressed as

$$
\begin{aligned}
&\left|\chi_{\mathrm{A}}-\chi_{\mathrm{B}}\right|=\left\{h c \tilde{D}_{0}(\mathrm{AB})-\frac{1}{2}\left[h c \tilde{D}_{0}(\mathrm{AA})+h c \tilde{D}_{0}(\mathrm{BB})\right]\right\}^{1 / 2} / \mathrm{eV} \\
& \text { Pauling electronegativity } \\
& \text { [definition] }
\end{aligned}
$$

Table 9D. 1 Pauling electronegativities*

\begin{center}
\begin{tabular}{ll}
\hline
Element & $\chi_{\mathrm{P}}$ \\
\hline
H & 2.2 \\
C & 2.6 \\
N & 3.0 \\
O & 3.4 \\
F & 4.0 \\
Cl & 3.2 \\
Cs & 0.79 \\
\hline
\end{tabular}
\end{center}

\begin{itemize}
  \item More values are given in the Resource section.\\
where $h c \tilde{D}_{0}(\mathrm{XY})$ is the dissociation energy of an $\mathrm{X}-\mathrm{Y}$ bond. This expression gives differences of electronegativities; to establish an absolute scale Pauling chose individual values that gave the best match to the values obtained from eqn 9D.2. Electronegativities based on this definition are called Pauling electronegativities (Table 9D.1). The most electronegative elements are those close to F (excluding the noble gases). It is found that the greater the difference in electronegativities, the greater the polar character of the bond. The difference for HF, for instance, is 1.8 ; a C-H bond, which is commonly regarded as almost nonpolar, has an electronegativity difference of 0.4.
\end{itemize}

\subsection*{Brief illustration 9D. 1}
The bond dissociation energies of $\mathrm{H}_{2}, \mathrm{Cl}_{2}$, and HCl are 4.52 eV , 2.51 eV , and 4.47 eV , respectively. From eqn 9D.2,

$$
\left|\chi_{\text {Pauling }}(\mathrm{H})-\chi_{\text {Pauling }}(\mathrm{Cl})\right|=\left\{4.47-\frac{1}{2}(4.52+2.51)\right\}^{1 / 2}=0.98 \approx 1.0
$$

The spectroscopist Robert Mulliken proposed an alternative definition of electronegativity. He argued that an element is likely to be highly electronegative if it has a high ionization energy (so it will not release electrons readily) and a high electron affinity (so it is energetically favourable to acquire electrons). The Mulliken electronegativity scale is therefore based on the definition

\[
\chi=\frac{1}{2}\left(I+E_{\mathrm{ea}}\right) / \mathrm{eV} \quad \begin{align*}
& \text { Mulliken electronegativity }  \tag{9D.3}\\
& \text { [definition] }
\end{align*}
\]

where $I$ is the ionization energy of the element and $E_{\text {ea }}$ is its electron affinity. The greater the value of the Mulliken electronegativity the greater is the contribution of that atom to the electron distribution in the bond. There is one word of caution: the values of $I$ and $E_{\text {ea }}$ in eqn 9D. 3 are strictly those for a special 'valence state' of the atom, not a true spectroscopic state, but that complication is ignored here. The Mulliken and Pauling scales are approximately in line with each other. A reasonably reliable conversion between the two is


\begin{equation*}
\chi_{\text {Pauling }}=1.35 \chi_{\text {Mulliken }}^{1 / 2}-1.37 \tag{9D.4}
\end{equation*}


\subsection*{9D. 2 The variation principle}
The systematic way of discussing bond polarity and finding the coefficients in the linear combinations used to build molecular orbitals is provided by the variation principle:

If an arbitrary wavefunction is used to calculate the energy, the value calculated is never less than the true energy.

It can be justified by setting up an arbitrary 'trial function' and showing that the corresponding energy is not less than the true energy (it might be the same).

\subsection*{How is that done? 9D. 1 Justifying the variation principle}
Any arbitrary function can be expressed as a linear combination of the eigenfunctions $\psi_{n}$ of the exact hamiltonian for a molecule. In the present case, consider a normalized trial wavefunction written as a linear combination $\psi_{\text {trial }}=\sum_{n} c_{n} \psi_{n}$ and suppose that the $\psi_{n}$ are themselves normalized and mutually orthogonal.

Step 1 Write an expression for the difference between the calculated and true energy

The energy associated with the normalized trial function is the expectation value

$$
E=\int \psi_{\text {trial }}^{\star} \hat{H} \psi_{\text {trial }} \mathrm{d} \tau
$$

The lowest energy of the system is $E_{0}$, the eigenvalue of $\hat{H}$ corresponding to $\psi_{0}$. Consider the following difference:

$$
\begin{aligned}
E-E_{0} & =\int \psi_{\text {trial }}^{*} \hat{H} \psi_{\text {trial }} \mathrm{d} \tau-E_{0} \overbrace{\int \psi_{\text {trial }}^{*} \psi_{\text {trial }} \mathrm{d} \tau}^{1} \\
& =\int \psi_{\text {trial }}^{*} \hat{H} \psi_{\text {trial }} \mathrm{d} \tau-\int \psi_{\text {trial }}^{\star} E_{0} \psi_{\text {trial }} \mathrm{d} \tau \\
& =\int \psi_{\text {trial }}^{\star}\left(\hat{H}-E_{0}\right) \psi_{\text {trial }} \mathrm{d} \tau \\
& =\int\left(\sum_{n}^{\left.c_{n}^{\star} \psi_{n}^{*}\right)\left(\hat{H}-E_{0}\right)\left(\sum_{n^{\prime}} c_{n^{\prime}} \psi_{n^{\prime}}\right) \mathrm{d} \tau}\right. \\
& =\sum_{n, n^{\prime}} c_{n}^{\star} c_{n^{\prime}} \int \psi_{n}^{\star}\left(\hat{H}-E_{0}\right) \psi_{n^{\prime}} \mathrm{d} \tau
\end{aligned}
$$

Step 2 Simplify the expression\\
Because $\int \psi_{n}^{\star} \hat{H} \psi_{n^{\prime}} \mathrm{d} \tau=E_{n^{\prime}} \int \psi_{n}^{\star} \psi_{n^{\prime}} \mathrm{d} \tau$ and $\int \psi_{n}^{\star} E_{0} \psi_{n^{\prime}} \mathrm{d} \tau=$ $E_{0} \int \psi_{n}^{\star} \psi_{n^{\prime}} \mathrm{d} \tau$, write

$$
\int \psi_{n}^{\star}\left(\hat{H}-E_{0}\right) \psi_{n^{\prime}} \mathrm{d} \tau=\left(E_{n^{\prime}}-E_{0}\right) \int \psi_{n}^{\star} \psi_{n^{\prime}} \mathrm{d} \tau
$$

It follows that

$$
E-E_{0}=\sum_{n, n^{\prime}} c_{n}^{\star} c_{n^{\prime}}\left(E_{n^{\prime}}-E_{0}\right) \overbrace{\int \psi_{n}^{*} \psi_{n^{\prime}} \mathrm{d} \tau}^{\substack{1 \text { if } n^{\prime}=n \\ 0 \text { otherwise }}}
$$

\subsection*{Step 3 Analyse the final expression}
The eigenfunctions $\psi_{n}$ are orthogonal, so only terms with $n^{\prime}=n$ contribute to this sum. Because each eigenfunction is normalized, each surviving integral is 1 . Consequently

$$
E-E_{0}=\sum_{n} \overbrace{c_{n}^{*} c_{n}}^{\geq 0} \overbrace{\left(E_{n}-E_{0}\right)}^{\geq 0} \geq 0
$$

The quantity $c_{n}^{*} c_{n}$ is necessarily real and greater than or equal to zero, and because $E_{0}$ is the lowest energy, $E_{n}-E_{0}$ is also greater than or equal to zero. It follows that the product of the two terms on the right is greater than or equal to zero. Therefore, $E \geq E_{0}$, as asserted.

The variation principle is the basis of all modern molecular structure calculations. The principle implies that, if the coefficients in the trial wavefunction are varied until the lowest energy is achieved (by evaluating the expectation value of the hamiltonian for the wavefunction in each case), then those coefficients will be the best for that particular form of trial function. A lower energy might be obtained with a more complicated wavefunction, for example, by taking a linear combination of several atomic orbitals on each atom. However, for a molecular orbital constructed from a given basis set, a given set of atomic orbits, the variation principle gives the optimum molecular orbital of that kind.

\subsection*{(a) The procedure}
The practical application of the variation principle can be illustrated by applying it to the trial wavefunction in eqn 9D.1, where the coefficients define the trial function.

\subsection*{How is that done? 9D. 2}
Applying the variation principle to a heteronuclear diatomic molecule

The trial wavefunction in eqn 9D. 1 is real but not normalized because at this stage the coefficients can take arbitrary values. Because it is real, write $\psi^{*}=\psi$. To normalize it, multiply it by $N=1 /\left(\int \psi^{\star} \psi \mathrm{d} \tau\right)^{1 / 2}$. So from now on use $\psi /\left(\int \psi^{2} \mathrm{~d} \tau\right)^{1 / 2}$ as the trial function. Then follow these steps.

\subsection*{Step 1 Write an expression for the energy}
The expectation value of the hamiltonian, the energy, using the normalized real trial function, is


\begin{equation*}
E=\frac{\int \psi \hat{H} \psi \mathrm{~d} \tau}{\int \psi^{2} \mathrm{~d} \tau} \tag{9D.5}
\end{equation*}


The denominator is

$$
\int \psi^{2} \mathrm{~d} \tau=\int\left(c_{\mathrm{A}} \psi_{\mathrm{A}}+c_{\mathrm{B}} \psi_{\mathrm{B}}\right)^{2} \mathrm{~d} \tau
$$

$$
\begin{aligned}
& =c_{\mathrm{A}}^{2} \overbrace{\int \psi_{\mathrm{A}}^{2} \mathrm{~d} \tau}^{1}+c_{\mathrm{B}}^{2} \overbrace{\int \psi_{\mathrm{B}}^{2} \mathrm{~d} \tau}^{1}+2 c_{\mathrm{A}} c_{\mathrm{B}} \overbrace{\psi_{\mathrm{A}} \psi_{\mathrm{B}} \mathrm{~d} \tau}^{\mathrm{d}} \\
& =c_{\mathrm{A}}^{2}+c_{\mathrm{B}}^{2}+2 c_{\mathrm{A}} c_{\mathrm{B}} S
\end{aligned}
$$

because the individual atomic orbitals are normalized to 1 and the third integral is the overlap integral $S$ (eqn 9C.3, $\left.S=\int \psi_{\mathrm{A}} \psi_{\mathrm{B}} \mathrm{d} \tau\right)$. The numerator is

$$
\begin{aligned}
\int \psi \hat{H} \psi \mathrm{~d} \tau & =\int\left(c_{\mathrm{A}} \psi_{\mathrm{A}}+c_{\mathrm{B}} \psi_{\mathrm{B}}\right) \hat{H}\left(c_{\mathrm{A}} \psi_{\mathrm{A}}+c_{\mathrm{B}} \psi_{\mathrm{B}}\right) \mathrm{d} \tau \\
& =c_{\mathrm{A}}^{2} \overbrace{\int \psi_{\mathrm{A}} \hat{H} \psi_{\mathrm{A}} \mathrm{~d} \tau}^{\alpha_{\mathrm{A}}}+c_{\mathrm{B}}^{2} \overbrace{\int \psi_{\mathrm{B}} \hat{H} \psi_{\mathrm{B}} \mathrm{~d} \tau}^{\alpha_{\mathrm{B}}}+c_{\mathrm{A}} c_{c_{\mathrm{B}}} \overbrace{\psi_{\mathrm{A}} \hat{H} \psi_{\mathrm{B}} \mathrm{~d} \tau}^{\beta} \\
& +c_{\mathrm{A}} c_{\mathrm{B}} \overbrace{\psi_{\mathrm{B}} \hat{H} \psi_{\mathrm{A}} \mathrm{~d} \tau}^{\beta}
\end{aligned}
$$

The significance of the quantities $\alpha_{\mathrm{A}}, \alpha_{\mathrm{B}}$, and $\beta$ (which are all energies) is discussed shortly. Because the hamiltonian is hermitian (Topic 7C), the third and fourth integrals are equal. Therefore

$$
\int \psi \hat{H} \psi \mathrm{~d} \tau=c_{\mathrm{A}}^{2} \alpha_{\mathrm{A}}+c_{\mathrm{B}}^{2} \alpha_{\mathrm{B}}+2 c_{\mathrm{A}} c_{\mathrm{B}} \beta
$$

At this point the complete expression for $E$ is

$$
E=\frac{c_{\mathrm{A}}^{2} \alpha_{\mathrm{A}}+c_{\mathrm{B}}^{2} \alpha_{\mathrm{B}}+2 c_{\mathrm{A}} c_{\mathrm{B}} \beta}{c_{\mathrm{A}}^{2}+c_{\mathrm{B}}^{2}+2 c_{\mathrm{A}} c_{\mathrm{B}} S}
$$

Step 2 Minimize the energy\\
Now search for values of the coefficients in the trial function that minimize the value of $E$. This is a standard problem in calculus, and is solved by finding the coefficients for which

$$
\frac{\partial E}{\partial c_{\mathrm{A}}}=0 \quad \frac{\partial E}{\partial c_{\mathrm{B}}}=0
$$

After some straightforward application of the rules of differentiation (The chemist's toolkit 5 in Topic 1C), the result is

$$
\begin{aligned}
& \frac{\partial E}{\partial c_{\mathrm{A}}}=\frac{2\left\{\left(\alpha_{\mathrm{A}}-E\right) c_{\mathrm{A}}+(\beta-S E) c_{\mathrm{B}}\right\}}{c_{\mathrm{A}}^{2}+c_{\mathrm{B}}^{2}+2 c_{\mathrm{A}} c_{\mathrm{B}} S} \\
& \frac{\partial E}{\partial c_{\mathrm{B}}}=\frac{2\left\{\left(\alpha_{\mathrm{B}}-E\right) c_{\mathrm{B}}+(\beta-S E) c_{\mathrm{A}}\right\}}{c_{\mathrm{A}}^{2}+c_{\mathrm{B}}^{2}+2 c_{\mathrm{A}} c_{\mathrm{B}} S}
\end{aligned}
$$

For the derivatives to be equal to 0 , the numerators, and specifically the terms in blue, of these expressions must vanish, leading to the secular equations: ${ }^{1}$

$-\left\lvert\,$\begin{tabular}{l|l}
$\left(\alpha_{\mathrm{A}}-E\right) c_{\mathrm{A}}+(\beta-S E) c_{\mathrm{B}}=0$ \\
$\left(\alpha_{\mathrm{B}}-E\right) c_{\mathrm{B}}+(\beta-S E) c_{\mathrm{A}}=0$ \\
\end{tabular}$\quad\right.$ Secular equations $\quad$ (9D.6a)

\footnotetext{${ }^{1}$ The name 'secular' is derived from the Latin word for age or generation. The term comes from astronomy, where the same equations appear in connection with slowly accumulating modifications of planetary orbits.
}The quantities $\alpha_{A}, \alpha_{B}, \beta$, and $S$ in the secular equations are

\[
\begin{array}{lrl}
\alpha_{\mathrm{A}}=\int \psi_{\mathrm{A}} \hat{H} \psi_{\mathrm{A}} \mathrm{~d} \tau \quad \alpha_{\mathrm{B}}=\int \psi_{\mathrm{B}} \hat{H} \psi_{\mathrm{B}} \mathrm{~d} \tau & \begin{array}{l}
\text { Coulomb } \\
\text { integrals }
\end{array} \\
\beta=\int \psi_{\mathrm{A}} \hat{H} \psi_{\mathrm{B}} \mathrm{~d} \tau=\int \psi_{\mathrm{B}} \hat{H} \psi_{\mathrm{A}} \mathrm{~d} \tau & \begin{array}{l}
\text { Resonance } \\
\text { integral }
\end{array} & \text { (9D.7b) }  \tag{9D.7b}\\
S=\int \psi_{\mathrm{A}} \psi_{\mathrm{B}} \mathrm{~d} \tau & \begin{array}{l}
\text { Overlap } \\
\text { integral }
\end{array} & \text { (9D.7c) }
\end{array}
\]

The parameter $\alpha$ is called a Coulomb integral. It is negative and can be interpreted as the energy of the electron when it occupies $\psi_{\mathrm{A}}\left(\right.$ for $\alpha_{\mathrm{A}}$ ) or $\psi_{\mathrm{B}}$ (for $\alpha_{\mathrm{B}}$ ). In a homonuclear diatomic molecule, $\alpha_{\mathrm{A}}=\alpha_{\mathrm{B}}$. The parameter $\beta$ is called a resonance integral (for classical reasons). It vanishes when the orbitals do not overlap, and at equilibrium bond lengths it is normally negative. The overlap integral $S$ is introduced and discussed in Topic 9C.

In order to solve the secular equations for the coefficients it is necessary to know the energy $E$ and then use its value in eqn 9D.6. As for any set of simultaneous equations, the secular equations have a solution if the secular determinant, the determinant of the coefficients (The chemist's toolkit 23), is zero. That is, if


\begin{align*}
\left.\begin{array}{cc}
\alpha_{\mathrm{A}}-E & \beta-S E \\
\beta-S E & \alpha_{\mathrm{B}}-E
\end{array} \right\rvert\,= & \left(\alpha_{\mathrm{A}}-E\right)\left(\alpha_{\mathrm{B}}-E\right)-(\beta-S E)^{2} \\
= & \left(1-S^{2}\right) E^{2}+\left\{2 \beta S-\left(\alpha_{\mathrm{A}}+\alpha_{\mathrm{B}}\right)\right\} E \\
& +\left(\alpha_{\mathrm{A}} \alpha_{\mathrm{B}}-\beta^{2}\right)=0 \tag{9D.8}
\end{align*}


This is a quadratic equation for $E$. A quadratic equation of the form $a x^{2}+b x+c=0$ has the solutions

$$
x=\frac{-b \pm\left(b^{2}-4 a c\right)^{1 / 2}}{2 a}
$$

In the present case, $a=1-S^{2}, b=2 \beta S-\left(\alpha_{\mathrm{A}}+\alpha_{\mathrm{B}}\right)$, and $c=\alpha_{\mathrm{A}} \alpha_{\mathrm{B}}-$ $\beta^{2}$, so the solutions (the energies) are

$$
E_{ \pm}=\frac{\alpha_{\mathrm{A}}+\alpha_{\mathrm{B}}-2 \beta S \pm\left\{\left(2 \beta S-\left(\alpha_{\mathrm{A}}+\alpha_{\mathrm{B}}\right)\right)^{2}-4\left(1-S^{2}\right)\left(\alpha_{\mathrm{A}} \alpha_{\mathrm{B}}-\beta^{2}\right)\right\}^{1 / 2}}{2\left(1-S^{2}\right)}
$$

(9D.9a)\\
which, according to the variation principle, are the closest approximations to the true energy for a trial function of the form given in eqn 9D.1. They are the energies of the bonding and antibonding molecular orbitals formed from the two atomic orbitals.

Equation 9D.9a can be simplified. For a homonuclear diatomic, $\alpha_{\mathrm{A}}=\alpha_{\mathrm{B}}=\alpha$ and then

$$
\begin{aligned}
E_{ \pm} & =\frac{2 \alpha-2 \beta S \pm\{\overbrace{(2 \beta S-2 \alpha)^{2}-4\left(1-S^{2}\right)\left(\alpha^{2}-\beta^{2}\right)}^{(2 \beta-2 \alpha S)^{2}}\}^{1 / 2}}{2 \underbrace{\left(1-S^{2}\right)}_{(1+S)(1-S)}} \\
& =\frac{\alpha-\beta S \pm(\beta-\alpha S)}{(1+S)(1-S)}=\frac{(\alpha \pm \beta)(1 \mp S)}{(1+S)(1-S)}
\end{aligned}
$$

That is,

$$
E_{+}=\frac{\alpha+\beta}{1+S} \quad E_{-}=\frac{\alpha-\beta}{1-S} \quad \text { Homonuclear diatomics }
$$

(9D.9b)\\
For $\beta<0, E_{+}$is the lower energy solution.\\
For heteronuclear diatomic molecules, making the approximation that $S=0$ (simply to obtain a more transparent expression) gives

$$
E_{ \pm}=\frac{1}{2}\left(\alpha_{\mathrm{A}}+\alpha_{\mathrm{B}}\right) \pm \frac{1}{2}\left(\alpha_{\mathrm{A}}-\alpha_{\mathrm{B}}\right)\left\{1+\left(\frac{2 \beta}{\alpha_{\mathrm{A}}-\alpha_{\mathrm{B}}}\right)^{2}\right\}^{1 / 2}
$$

Zero overlap approximation\\
(9D.9c)\\
The values of the Coulomb integrals $\alpha_{\mathrm{A}}$ and $\alpha_{\mathrm{B}}$ may be estimated as follows. The extreme case of an atom X in a molecule is $\mathrm{X}^{+}$if it has lost control of the electron it supplied, X if it is sharing the electron pair equally with its bonded partner, and $\mathrm{X}^{-}$if it has gained control of both electrons in the bond. If $\mathrm{X}^{+}$is taken as defining the energy 0 , then X lies at $-I(\mathrm{X})$ and $\mathrm{X}^{-}$lies at $-\left\{I(\mathrm{X})+E_{\text {ea }}(\mathrm{X})\right\}$, where $I$ is the ionization energy and $E_{\text {ea }}$ the electron affinity (Fig. 9D.2). The actual energy of the electron in the molecule lies at an intermediate value, and in the absence of further information, it is reasonable to estimate it as half-way down to the lowest of these values, namely at $-\frac{1}{2}\{I(\mathrm{X})$ $\left.+E_{\text {ea }}(\mathrm{X})\right\}$. This quantity should be recognized (apart from its sign) as the Mulliken definition of electronegativity.

\subsection*{The chemist's toolkit 23 Determinants}
A $2 \times 2$ determinant is the entity

$$
\left.\begin{array}{ll}
a & b \\
c & d
\end{array} \right\rvert\,=a d-b c
$$

$2 \times 2$ Determinant

A $3 \times 3$ determinant is evaluated by expanding it as a sum of $2 \times 2$ determinants:

$$
\begin{aligned}
\left|\begin{array}{lll}
a & b & c \\
d & e & f \\
g & h & i
\end{array}\right| & =a\left|\begin{array}{cc}
e & f \\
h & i
\end{array}\right|-b\left|\begin{array}{cc}
d & f \\
g & i
\end{array}\right|+c\left|\begin{array}{ll}
d & e \\
g & h
\end{array}\right| \\
& =a(e i-f h)-b(d i-f g)+c(d h-e g) \quad 3 \times 3 \text { Determinant }
\end{aligned}
$$

Note the sign change in alternate columns ( $b$ occurs with a negative sign in the expansion). An important property of a determinant is that if any two rows or any two columns are interchanged, then the determinant changes sign:\\
Exchange columns: $\left|\begin{array}{ll}b & a \\ d & c\end{array}\right|=b c-a d=-(a d-b c)=-\left|\begin{array}{ll}a & b \\ c & d\end{array}\right|$\\
Exchange rows: $\left|\begin{array}{ll}c & d \\ a & b\end{array}\right|=c b-d a=-(a d-b c)=-\left|\begin{array}{ll}a & b \\ c & d\end{array}\right|$\\
An implication is that if any two columns or rows are identical, then the determinant is zero.\\
\includegraphics[max width=\textwidth, center]{2025_02_27_d4b95d9718f0b9e2e6b9g-401(1)}

Figure 9D. 2 The procedure for estimating the Coulomb integral in terms of the ionization energy and electron affinity.

\subsection*{Brief illustration 9D. 2}
Consider HF. The general form of the molecular orbital is $\psi=$ $c_{\mathrm{H}} \psi_{\mathrm{H}}+c_{\mathrm{F}} \psi_{\mathrm{F},}$, where $\psi_{\mathrm{H}}$ is an H1s orbital and $\psi_{\mathrm{F}}$ is an F2 $p_{z}$ orbital (with $z$ along the internuclear axis, the convention for linear molecules). The relevant data are as follows:

\begin{center}
\begin{tabular}{llll}
\hline
 & $I / \mathrm{eV}$ & $E_{\text {ea }} / \mathrm{eV}$ & $-\frac{1}{2}\left\{I+E_{\text {ea }}\right\} / \mathrm{eV}$ \\
\hline
H & 13.6 & 0.75 & -7.2 \\
F & 17.4 & 3.34 & -10.4 \\
\hline
\end{tabular}
\end{center}

Therefore set $\alpha_{\mathrm{A}}=\alpha_{\mathrm{H}}=-7.2 \mathrm{eV}$ and $\alpha_{\mathrm{B}}=\alpha_{\mathrm{F}}=-10.4 \mathrm{eV}$. Taking $\beta=-1.0 \mathrm{eV}$ as a typical value and setting $S=0$ for simplicity, substitution into eqn 9D.9c gives

$$
\begin{aligned}
E_{ \pm} / \mathrm{eV} & =\frac{1}{2}(-7.2-10.4) \pm \frac{1}{2}(-7.2+10.4)\left\{1+\left(\frac{-2.0}{-7.2+10.4}\right)^{2}\right\}^{1 / 2} \\
& =-8.8 \pm 1.9=-10.7 \text { and }-6.9
\end{aligned}
$$

These values, representing a bonding orbital at -10.7 eV and an antibonding orbital at -6.9 eV , are shown in Fig. 9D.3.

\subsection*{(b) The features of the solutions}
An important feature of eqn 9D.9c is that as the energy difference $\left|\alpha_{\mathrm{A}}-\alpha_{\mathrm{B}}\right|$ between the interacting atomic orbitals increases, the bonding and antibonding effects decrease (Fig. 9D.4). When $\left|\alpha_{\mathrm{B}}-\alpha_{\mathrm{A}}\right| \gg 2|\beta|$ it is possible to use the approximation $(1+x)^{1 / 2} \approx 1+\frac{1}{2} x$ (see The chemist's toolkit 12 in Topic 5B) to obtain


\begin{equation*}
E_{+} \approx \alpha_{\mathrm{A}}+\frac{\beta^{2}}{\alpha_{\mathrm{A}}-\alpha_{\mathrm{B}}} \quad E_{-} \approx \alpha_{\mathrm{B}}-\frac{\beta^{2}}{\alpha_{\mathrm{A}}-\alpha_{\mathrm{B}}} \tag{9D.10}
\end{equation*}


As these expressions show, and as can be seen from the graph, when the energy difference $\left|\alpha_{\mathrm{A}}-\alpha_{\mathrm{B}}\right|$ is very large, the energies of the resulting molecular orbitals differ only slightly from those of the atomic orbitals, which implies in turn that the bonding and antibonding effects are small. That is:\\
\includegraphics[max width=\textwidth, center]{2025_02_27_d4b95d9718f0b9e2e6b9g-401}

Figure 9D. 3 The estimated energies of the Coulomb integrals $\alpha$ in HF and the molecular orbitals they form.

The strongest bonding and antibonding effects are obtained when the two contributing orbitals have similar energies.

The large difference in energy between core and valence orbitals is the justification for neglecting the contribution of core orbitals to molecular orbitals constructed from valence atomic orbitals. Although the core orbitals of one atom have a similar energy to the core orbitals of the other atom, so might be expected to combine strongly, core-core interaction is largely negligible because the core orbitals are so contracted that the interaction between them, as measured by the value of $|\beta|$, is negligible. It is also a justification for treating the $s$ and $\mathrm{p}_{z}$ contributions to $\sigma$ orbital formation separately, an approximation used in Topic 9C in the discussion of homonuclear diatomic molecules.

The values of the coefficients in the linear combination in eqn 9D. 1 are obtained by solving the secular equations after substituting the two energies obtained from the secular determinant. The lower energy, $E_{+}$, gives the coefficients for the bonding molecular orbital, the upper energy, $E_{-}$, the coefficients for the antibonding molecular orbital. The secular\\
\includegraphics[max width=\textwidth, center]{2025_02_27_d4b95d9718f0b9e2e6b9g-401(2)}

Figure 9D.4 The variation of the energies of the molecular orbitals as the energy difference of the contributing atomic orbitals is changed. The plots are for $\beta=-1 \mathrm{eV}$; the blue lines are for the energies in the absence of mixing (i.e. $\beta=0$ ).\\
equations give expressions for the ratio of the coefficients. Thus, the first of the two secular equations in eqn 9D.6a, $\left(\alpha_{\mathrm{A}}-E\right) c_{\mathrm{A}}+(\beta-E S) c_{\mathrm{B}}=0$, gives


\begin{equation*}
c_{\mathrm{B}}=-\left(\frac{\alpha_{\mathrm{A}}-E}{\beta-E S}\right) c_{\mathrm{A}} \tag{9D.11}
\end{equation*}


The wavefunction should also be normalized. It has already been shown that $\int \psi^{2} \mathrm{~d} \tau=c_{A}^{2}+c_{B}^{2}+2 c_{A} c_{B} S$, so normalization requires that


\begin{equation*}
c_{\mathrm{A}}^{2}+c_{\mathrm{B}}^{2}+2 c_{\mathrm{A}} c_{\mathrm{B}} S=1 \tag{9D.12}
\end{equation*}


When eqn 9D. 11 is substituted into this expression, the result is


\begin{equation*}
c_{\mathrm{A}}=\frac{1}{\left\{1+\left(\frac{\alpha_{\mathrm{A}}-E}{\beta-E S}\right)^{2}-2 S\left(\frac{\alpha_{\mathrm{A}}-E}{\beta-E S}\right)\right\}^{1 / 2}} \tag{9D.13}
\end{equation*}


which, together with eqn 9D.11, gives explicit expressions for the coefficients once the appropriate values of $E=E_{ \pm}$given in eqn 9D. 9 a are substituted.

As before, this expression becomes more transparent in two cases. First, for a homonuclear diatomic, with $\alpha_{\mathrm{A}}=\alpha_{\mathrm{B}}=\alpha$ and $E_{ \pm}$given in eqn 9D.9b, the results are

\[
\left.\begin{array}{lll}
E_{+}=\frac{\alpha+\beta}{1+S} & c_{\mathrm{A}}=\frac{1}{\{2(1+S)\}^{1 / 2}} & c_{\mathrm{B}}=c_{\mathrm{A}}
\end{array} \begin{array}{l}
\text { Homonuclear } \\
\text { diatomics }
\end{array}\right] \begin{array}{lll}
E_{-}=\frac{\alpha-\beta}{1-S} & c_{\mathrm{A}}=\frac{1}{\{2(1-S)\}^{1 / 2}} & c_{\mathrm{B}}=-c_{\mathrm{A}} \tag{9D.14b}
\end{array}
\]

For a heteronuclear diatomic with $S=0$,

$$
c_{\mathrm{A}}=\frac{1}{\left\{1+\left(\frac{\alpha_{\mathrm{A}}-E}{\beta}\right)^{2}\right\}^{1 / 2}}
$$

with the appropriate values of $E=E_{ \pm}$taken from eqn 9D.9c. The coefficient $c_{\mathrm{B}}$ is then calculated from eqn 9D.11.

\subsection*{Brief illustration 9D. 3}
Consider HF again. In the previous Briefillustration, with $\alpha_{\mathrm{H}}=$ $-7.2 \mathrm{eV}, \alpha_{\mathrm{F}}=-10.4 \mathrm{eV}, \beta=-1.0 \mathrm{eV}$, and $S=0$, the two orbital energies were found to be $E_{+}=-10.7 \mathrm{eV}$ and $E_{-}=-6.9 \mathrm{eV}$. When these values are substituted into eqn 9D. 15 the following coefficients are found:

$$
\begin{array}{ll}
E_{+}=-10.7 \mathrm{eV} & \psi_{+}=0.28 \psi_{\mathrm{H}}+0.96 \psi_{\mathrm{F}} \\
E_{-}=-6.9 \mathrm{eV} & \psi_{-}=0.96 \psi_{\mathrm{H}}-0.28 \psi_{\mathrm{F}}
\end{array}
$$

Notice how the lower energy orbital (the one with energy -10.7 eV ), and belonging to the atom with the greater electronegativity, has a composition that is more F2p orbital than H1s, and that the opposite is true of the higher energy, antibonding orbital.

\subsection*{Checklist of concepts}
\begin{enumerate}
  \item A polar bond can be regarded as arising from a molecular orbital that is concentrated more on one atom than its partner.2. The electronegativity of an element is a measure of the power of an atom to attract electrons to itself in a bond.
  \item The electron pair in a bonding orbital is more likely to be found on the more electronegative atom; the opposite is true for electrons in an antibonding orbital.
  \item The variation principle provides a criterion for optimizing a trial wavefunction.\\
$\square$ 4. A basis set is the set of atomic orbitals from which the molecular orbitals are constructed.
  \item The bonding and antibonding effects are strongest when contributing atomic orbitals have similar energies.
\end{enumerate}

\subsection*{Checklist of equations}
\begin{center}
\begin{tabular}{llll}
\hline
Property & Equation & Comment & Equation number \\
\hline
Molecular orbital & $\psi=c_{\mathrm{A}} \psi_{\mathrm{A}}+c_{\mathrm{B}} \psi_{\mathrm{B}}$ & 9 D .1 &  \\
Pauling electronegativity & $\left|\chi_{\mathrm{A}}-\chi_{\mathrm{B}}\right|=\left\{h c \tilde{D}_{0}(\mathrm{AB})-\frac{1}{2}\left[h c \tilde{D}_{0}(\mathrm{AA})+h c \tilde{D}_{0}(\mathrm{BB})\right]\right\}^{1 / 2} / \mathrm{eV}$ &  & 9 D .2 \\
Mulliken electronegativity & $\chi=\frac{1}{2}\left(I+E_{\text {ea }}\right) / \mathrm{eV}$ & 9 D .3 &  \\
Coulomb integral & $\alpha_{\mathrm{A}}=\int \psi_{\mathrm{A}} \hat{H} \psi_{\mathrm{A}} \mathrm{d} \tau$ & Definition & 9 D .7 a \\
Resonance integral & $\beta=\int \psi_{\mathrm{A}} \hat{H} \psi_{\mathrm{B}} \mathrm{d} \tau$ & Definition & 9 D .7 b \\
Variation principle & $E=\int \psi_{\text {trial }}^{\star} \hat{H} \psi_{\text {trial }} \mathrm{d} \tau / \int \psi_{\text {trial }}^{\star} \psi_{\text {trial }} \mathrm{d} \tau ; \partial E / \partial c=0$ &  &  \\
\hline
\end{tabular}
\end{center}

\section*{TOPIC 9E Molecular orbital theory: polyatomic molecules }
\begin{abstract}
-Why do you need to know this material? Most molecules of interest in chemistry are polyatomic, so it is important to be able to discuss their electronic structure. Although computational procedures are now widely available, to understand them it is helpful to see how they emerged from the more primitive approach described here.
\end{abstract}

\subsection*{- What is the key idea?}
Molecular orbitals can be expressed as linear combinations of all the atomic orbitals of the appropriate symmetry.

\subsection*{> What do you need to know already?}
This Topic extends the approach used for heteronuclear diatomic molecules in Topic 9D, particularly the concepts of secular equations and secular determinants. The principal mathematical technique used is matrix algebra (The chemist's toolkits 24 and 25). You should become familiar with the use of mathematical software to manipulate matrices numerically.

The molecular orbitals of polyatomic molecules are built in the same way as in diatomic molecules (Topic 9D), the only difference being that more atomic orbitals are used to construct them. As for diatomic molecules, polyatomic molecular orbitals spread over the entire molecule. A molecular orbital has the general form

$$
\psi=\sum_{i} c_{i} \psi_{i}
$$

General form of LCAO-MO\\
(9E.1)\\
where $\psi_{i}$ is an atomic orbital and the sum extends over all the valence orbitals of all the atoms in the molecule. The coefficients are found by setting up the secular equations, just as for diatomic molecules, then solving them for the energies (Topic 9D). That step involves formulating the secular determinant and finding the values of the energy that ensure the determinant is equal to 0 . Finally these energies are used in the secular equations to find the coefficients of the atomic orbitals for each molecular orbital.

The principal difference between diatomic and polyatomic molecules lies in the greater range of shapes that are possible: a diatomic molecule is necessarily linear, but a triatomic\\
molecule, for instance, may be either linear or angular (bent) with a characteristic bond angle. The shape of a polyatomic molecule-the specification of its bond lengths and its bond angles-can be predicted by calculating the total energy of the molecule for a variety of nuclear positions, and then identifying the conformation that corresponds to the lowest energy. Such calculations are best done using software, which handles the minimization problem automatically and generates the molecular orbital coefficients. However, a more primitive approach gives useful insights into the electronic structure of polyatomic molecules and its interpretation.

Symmetry considerations play a central role in the construction of molecular orbitals of polyatomic molecules, for only atomic orbitals of matching symmetry have non-zero overlap and contribute to a molecular orbital. To discuss these symmetry requirements fully requires the machinery developed in Focus 10, especially Topic 10C. There is one type of symmetry, however, that is intuitive: the planarity of conjugated hydrocarbons. That symmetry provides a distinction between the $\sigma$ and $\pi$ orbitals of the molecule, and in elementary approaches such molecules are commonly discussed in terms of the characteristics of their $\pi$ orbitals, with the $\sigma$ bonds providing a rigid framework that determines the general shape of the molecule.

\subsection*{9E. 1 The Hückel approximation}
The $\pi$ molecular orbital energy level diagrams of conjugated molecules can be constructed by using a set of approximations suggested by Erich Hückel in 1931. All the C atoms are treated identically, so all the Coulomb integrals $\alpha$ (Topic 9D) for the atomic orbitals that contribute to the $\pi$ orbitals are set equal. For example, in ethene, which is used here to introduce the method, the $\sigma$ bonds are regarded as fixed, and the calculation leads to the energies of $\pi$ bonding and antibonding molecular orbitals.

\subsection*{(a) An introduction to the method}
The $\pi$ orbitals are expressed as linear combinations of the C2p orbitals that lie perpendicular to the molecular plane. In ethene, for instance,


\begin{equation*}
\psi=c_{\mathrm{A}} \psi_{\mathrm{A}}+c_{\mathrm{B}} \psi_{\mathrm{B}} \tag{9E.2}
\end{equation*}


where $\psi_{\mathrm{A}}$ is a C2p orbital on atom A , and likewise for $\psi_{\mathrm{B}}$. Next, the optimum coefficients and energies are found by the variation principle as explained in Topic 9D. That is, the appropriate secular determinant is set up, equated to 0 , and the equation solved for the energies. For ethene, with $\alpha_{\mathrm{A}}=\alpha_{\mathrm{B}}=\alpha$, the secular determinant is

\[
\left|\begin{array}{lr}
\alpha-E & \beta-E S  \tag{9E.3}\\
\beta-E S & \alpha-E
\end{array}\right|=0
\]

In a modern computation all the resonance integrals and overlap integrals would be included, but an indication of the molecular orbital energy level diagram can be obtained more readily by making the following additional Hückel approximations:

\begin{itemize}
  \item All overlap integrals are set equal to zero.
  \item All resonance integrals between non-neighbours are set equal to zero.
  \item All remaining resonance integrals are set equal (to $\beta$ ).
\end{itemize}

These approximations are obviously very severe, but they give at least a general picture of the molecular orbital energy levels. The approximations result in the following structure of the secular determinant:

\begin{itemize}
  \item All diagonal elements: $\alpha-E$
  \item Off-diagonal elements between neighbouring atoms: $\beta$
  \item All other elements: 0
\end{itemize}

These approximations convert eqn 9E. 3 into

\[
\begin{array}{cc}
\alpha-E & \beta  \tag{9E.4}\\
\beta & \alpha-E
\end{array}=(\alpha-E)^{2}-\beta^{2}=(\alpha-E+\beta)(\alpha-E-\beta)=0
\]

where the determinant has been expanded as explained in The chemist's toolkit 23 in Topic 9D. The roots of the equation are $E=\alpha \pm \beta$. The $+\operatorname{sign}$ corresponds to the bonding combination ( $\beta$ is negative) and the - sign corresponds to the antibonding combination (Fig. 9E.1).

The building-up principle results in the configuration $1 \pi^{2}$, because each carbon atom supplies one electron to the $\pi$ system and both electrons can occupy the bonding orbital. The\\
\includegraphics[max width=\textwidth, center]{2025_02_27_d4b95d9718f0b9e2e6b9g-404}

Figure 9E. 1 The Hückel molecular orbital energy levels of ethene. Two electrons occupy the lower $\pi$ orbital.\\
highest occupied molecular orbital in ethene, its HOMO, is the $1 \pi$ orbital; the lowest unoccupied molecular orbital, its LUMO, is the $2 \pi$ orbital (or, as it is sometimes denoted, the $2 \pi^{*}$ orbital). These two orbitals jointly form the frontier orbitals of the molecule. The frontier orbitals are important because they are largely responsible for many of the chemical and spectroscopic properties of this and analogous molecules.

\subsection*{Brief illustration 9E. 1}
Within the Hückel framework the energy needed to excite a $\pi^{*} \leftarrow \pi$ transition is equal to the separation of the $1 \pi$ and $2 \pi$ orbitals, which is $2|\beta|$. This transition is known to occur at close to $40000 \mathrm{~cm}^{-1}$, corresponding to 5.0 eV . It follows that a plausible value of $\beta$ is about $-2.5 \mathrm{eV}\left(-240 \mathrm{~kJ} \mathrm{~mol}^{-1}\right)$.

\subsection*{(b) The matrix formulation of the method}
To make the Hückel theory readily applicable to bigger molecules, it helps to reformulate it in terms of matrices (see The chemist's toolkit 24). The starting point is the pair of secular equations developed for a heteronuclear diatomic molecule in Topic 9D:

$$
\begin{aligned}
& \left(\alpha_{\mathrm{A}}-E\right) c_{\mathrm{A}}+(\beta-E S) c_{\mathrm{B}}=0 \\
& (\beta-E S) c_{\mathrm{A}}+\left(\alpha_{\mathrm{B}}-E\right) c_{\mathrm{B}}=0
\end{aligned}
$$

To prepare to generalize this expression write $\alpha_{\mathrm{I}}=H_{\mathrm{JJ}}$ (with J $=\mathrm{A}$ or B ), $\beta=H_{\mathrm{AB}}$, and label the overlap integrals with their respective atoms, so $S$ becomes $S_{\mathrm{AB}}$. More symmetry can be introduced into the equations (which makes it simpler to generalize them) by replacing the $E$ in $\alpha_{\mathrm{J}}-E$ by $E S_{\mathrm{JJ}}$, with $S_{\mathrm{JJ}}=1$. At this point, the two equations are

$$
\begin{aligned}
& \left(H_{\mathrm{AA}}-E S_{\mathrm{AA}}\right) c_{\mathrm{A}}+\left(H_{\mathrm{AB}}-E S_{\mathrm{AB}}\right) c_{\mathrm{B}}=0 \\
& \left(H_{\mathrm{BA}}-E S_{\mathrm{BA}}\right) c_{\mathrm{A}}+\left(H_{\mathrm{BB}}-E S_{\mathrm{BB}}\right) c_{\mathrm{B}}=0
\end{aligned}
$$

There is one further notational change. The coefficients $c_{\mathrm{J}}$ depend on the value of $E$, so it is necessary to distinguish the two sets corresponding to the two energies, denoted $E_{n}$ with $n=1$ and 2 . The coefficients are written as $c_{n, \mathrm{~J}}$, with $n=1$ (the coefficients $c_{1, \mathrm{~A}}$ and $c_{1, \mathrm{~B}}$ for energy $E_{1}$ ) or 2 (the coefficients $c_{2, \mathrm{~A}}$ and $c_{2, \mathrm{~B}}$ for energy $E_{2}$ ). With this notational change, the two equations become


\begin{align*}
& \left(H_{\mathrm{AA}}-E_{n} S_{\mathrm{AA}}\right) c_{n, \mathrm{~A}}+\left(H_{\mathrm{AB}}-E_{n} S_{\mathrm{AB}}\right) c_{n, \mathrm{~B}}=0  \tag{9‥5a}\\
& \left(H_{\mathrm{BA}}-E_{n} S_{\mathrm{BA}}\right) c_{n, \mathrm{~A}}+\left(H_{\mathrm{BB}}-E_{n} S_{\mathrm{BB}}\right) c_{n, \mathrm{~B}}=0 \tag{9E.5b}
\end{align*}


with $n=1$ and 2, giving four equations in all. Each pair of equations can be written in matrix form as

\[
\left(\begin{array}{ll}
H_{\mathrm{AA}}-E_{n} S_{\mathrm{AA}} & H_{\mathrm{AB}}-E_{n} S_{\mathrm{AB}}  \tag{9E.5c}\\
H_{\mathrm{BA}}-E_{n} S_{\mathrm{BA}} & H_{\mathrm{BB}}-E_{n} S_{\mathrm{BB}}
\end{array}\right)\binom{c_{n, \mathrm{~A}}}{c_{n, \mathrm{~B}}}=0
\]

\subsection*{The chemist's toolkit 24}
\subsection*{Matrices}
A matrix is an array of numbers arranged in a certain number of rows and a certain number of columns; the numbers of rows and columns may be different. The rows and columns are numbered $1,2, \ldots$ so that the number at each position in the matrix, called the matrix element, has a unique row and column index. The element of a matrix $M$ at row $r$ and column c is denoted $M_{r c}$. For instance, a $3 \times 3$ matrix is

$$
\boldsymbol{M}=\left(\begin{array}{lll}
M_{11} & M_{12} & M_{13} \\
M_{21} & M_{22} & M_{23} \\
M_{31} & M_{32} & M_{33}
\end{array}\right)
$$

The trace of a matrix, $\operatorname{Tr} M$, is the sum of the diagonal elements. In this case

$$
\operatorname{Tr} M=M_{11}+M_{22}+M_{33}
$$

A unit matrix has diagonal elements equal to 1 and all other elements zero. A $3 \times 3$ unit matrix is therefore

$$
\mathbf{1}=\left(\begin{array}{lll}
1 & 0 & 0 \\
0 & 1 & 0 \\
0 & 0 & 1
\end{array}\right)
$$

Matrices are added by adding the corresponding matrix elements. Thus, to add the matrices $A$ and $B$ to give the sum $S=A+B$, each element of $S$ is given by

$$
S_{r c}=A_{r c}+B_{r c}
$$

Only matrices of the same dimensions can be added together.\\
Matrices are multiplied to obtain the product $P=A B$; each element of $P$ is given by

$$
P_{r c}=\sum_{n} A_{r n} B_{n c}
$$

as can be verified by multiplying out the matrices to give the two expressions in eqns 9E.5a and 9E.5b. Now introduce the hamiltonian matrix $H$ and the overlap matrix $S$, and write the coefficients corresponding to the energy $E_{n}$ as a column vector $c_{n}$ :

\[
H=\left(\begin{array}{ll}
H_{\mathrm{AA}} & H_{\mathrm{AB}}  \tag{9E.6}\\
H_{\mathrm{BA}} & H_{\mathrm{BB}}
\end{array}\right) \quad \boldsymbol{S}=\left(\begin{array}{ll}
S_{\mathrm{AA}} & S_{\mathrm{AB}} \\
S_{\mathrm{BA}} & S_{\mathrm{BB}}
\end{array}\right) \quad \boldsymbol{c}_{n}=\binom{c_{n, \mathrm{~A}}}{c_{n, \mathrm{~B}}}
\]

Then

$$
H-E_{n} S=\left(\begin{array}{ll}
H_{\mathrm{AA}}-E_{n} S_{\mathrm{AA}} & H_{\mathrm{AB}}-E_{n} S_{\mathrm{AB}} \\
H_{\mathrm{BA}}-E_{n} S_{\mathrm{BA}} & H_{\mathrm{BB}}-E_{n} S_{\mathrm{BB}}
\end{array}\right)
$$

and eqn 9E.5c may be written more succinctly as


\begin{equation*}
\left(H-E_{n} S\right) c_{n}=0 \quad \text { or } \quad H c_{n}=S c_{n} E_{n} \tag{9E.7}
\end{equation*}


Matrices can be multiplied only if the number of columns in $A$ is equal to the number of rows in $\boldsymbol{B}$. Square matrices (those with the same number of rows and columns) can therefore be multiplied only if both matrices have the same dimension (that is, both are $n \times n$ ). The products $A B$ and $B A$ are not necessarily the same, so matrix multiplication is in general 'non-commutative'.\\
An $n \times 1$ matrix (with $n$ elements in one column) is called a column vector. It may be multiplied by a square $n \times n$ matrix to generate a new column vector, as in

$$
\left(\begin{array}{c}
P_{1} \\
P_{2} \\
P_{3}
\end{array}\right)=\left(\begin{array}{lll}
A_{11} & A_{12} & A_{13} \\
A_{21} & A_{22} & A_{23} \\
A_{31} & A_{32} & A_{33}
\end{array}\right) \times\left(\begin{array}{c}
B_{1} \\
B_{2} \\
B_{3}
\end{array}\right)
$$

The elements of the two column vectors need only one index to indicate their row. Each element of $P$ is given by

$$
P_{r}=\sum_{n} A_{r n} B_{n}
$$

A $1 \times n$ matrix (a single row with $n$ elements) is called a row vector. It may be multiplied by a square $n \times n$ matrix to generate a new row vector, as in

$$
\left(P_{1} P_{2} P_{3}\right)=\left(B_{1} B_{2} B_{3}\right) \times\left(\begin{array}{ccc}
A_{11} & A_{12} & A_{13} \\
A_{21} & A_{22} & A_{23} \\
A_{31} & A_{32} & A_{33}
\end{array}\right)
$$

In general the elements of $P$ are

$$
P_{c}=\sum_{n} B_{n} A_{n c}
$$

Note that a column vector is multiplied 'from the left' by the square matrix and a row vector is multiplied 'from the right'. The inverse of a matrix $A$, denoted $A^{-1}$, has the property that $A A^{-1}=$ $A^{-1} A=1$, where 1 is a unit matrix with the same dimensions as $A$.

These two sets of equations (with $n=1$ and 2 ) can be combined into a single matrix equation of the form


\begin{equation*}
H c=S c E \tag{9E.8}
\end{equation*}


by introducing the matrices

\[
\boldsymbol{c}=\left(\begin{array}{ll}
c_{1} & c_{2}
\end{array}\right)=\left(\begin{array}{ll}
c_{1, \mathrm{~A}} & c_{2, \mathrm{~A}}  \tag{9E.9}\\
c_{1, \mathrm{~B}} & c_{2, \mathrm{~B}}
\end{array}\right) \quad \boldsymbol{E}=\left(\begin{array}{cc}
E_{1} & 0 \\
0 & E_{2}
\end{array}\right)
\]

\subsection*{How is that done? 9E. 1 Justifying the matrix formulation}
Substitution of the matrices defined in eqn 9E. 9 into eqn 9E. 8 gives

$$
\overbrace{\left(\begin{array}{cc}
H_{\mathrm{AA}} & H_{\mathrm{AB}} \\
H_{\mathrm{BA}} & H_{\mathrm{BB}}
\end{array}\right)}^{\mathrm{H}} \overbrace{\left(\begin{array}{cc}
c_{1, \mathrm{~A}} & c_{2, \mathrm{~A}} \\
c_{1, \mathrm{~B}} & c_{2, \mathrm{~B}}
\end{array}\right)}^{c}=\overbrace{\left(\begin{array}{cc}
S_{\mathrm{AA}} & S_{\mathrm{AB}} \\
S_{\mathrm{BA}} & S_{\mathrm{BB}}
\end{array}\right)}^{s} \overbrace{\left(\begin{array}{cc}
c_{1, \mathrm{~A}} & c_{2, \mathrm{~A}} \\
c_{1, \mathrm{~B}} & c_{2, \mathrm{~B}}
\end{array}\right)}^{c} \overbrace{\left(\begin{array}{cc}
E_{1} & 0 \\
0 & E_{2}
\end{array}\right)}^{E}
$$

The product on the left is\\
$\left(\begin{array}{ll}H_{\mathrm{AA}} & H_{\mathrm{AB}} \\ H_{\mathrm{BA}} & H_{\mathrm{BB}}\end{array}\right)\left(\begin{array}{cc}c_{1, \mathrm{~A}} & c_{2, \mathrm{~A}} \\ c_{1, \mathrm{~B}} & c_{2, \mathrm{~B}}\end{array}\right)=\left(\begin{array}{cc}H_{\mathrm{AA}} c_{1, \mathrm{~A}}+H_{\mathrm{AB}} c_{1, \mathrm{~B}} & H_{\mathrm{AA}} c_{2, \mathrm{~A}}+H_{\mathrm{AB}} c_{2, \mathrm{~B}} \\ H_{\mathrm{BA}} c_{1, \mathrm{~A}}+H_{\mathrm{BB}} c_{1, \mathrm{~B}} & H_{\mathrm{BA}} c_{2, \mathrm{~A}}+H_{\mathrm{BB}} c_{2, \mathrm{~B}}\end{array}\right)$\\
The product on the right is

$$
\begin{gathered}
\left(\begin{array}{cc}
S_{\mathrm{AA}} & S_{\mathrm{AB}} \\
S_{\mathrm{BA}} & S_{\mathrm{BB}}
\end{array}\right)\left(\begin{array}{cc}
c_{1, \mathrm{~A}} & c_{2, \mathrm{~A}} \\
c_{1, \mathrm{~B}} & c_{2, \mathrm{~B}}
\end{array}\right)\left(\begin{array}{cc}
E_{1} & 0 \\
0 & E_{2}
\end{array}\right)=\left(\begin{array}{ll}
S_{\mathrm{AA}} & S_{\mathrm{AB}} \\
S_{\mathrm{BA}} & S_{\mathrm{BB}}
\end{array}\right)\left(\begin{array}{ll}
c_{1, \mathrm{~A}} E_{1} & c_{2, \mathrm{~A}} E_{2} \\
c_{1, \mathrm{~B}} E_{1} & c_{2, \mathrm{~B}} E_{2}
\end{array}\right) \\
=\left(\begin{array}{ll}
E_{1} S_{\mathrm{AA}} c_{1, \mathrm{~A}}+E_{1} S_{\mathrm{AB}} c_{1, \mathrm{~B}} & E_{2} S_{\mathrm{AA}} c_{2, \mathrm{~A}}+E_{2} S_{\mathrm{AB}} c_{2, \mathrm{~B}} \\
E_{1} S_{\mathrm{BA}} c_{1, \mathrm{~A}}+E_{1} S_{\mathrm{BB}} c_{1, \mathrm{~B}} & E_{2} S_{\mathrm{BA}} c_{2, \mathrm{~A}}+E_{2} S_{\mathrm{BB}} c_{2, \mathrm{~B}}
\end{array}\right)
\end{gathered}
$$

Comparison of matching terms (like those in blue) recreates the four secular equations (two for each value of $n$ ) given in eqns 9E.5a and 9E.5b.

In the Hückel approximation, $H_{\mathrm{AA}}=H_{\mathrm{BB}}=\alpha, H_{\mathrm{AB}}=H_{\mathrm{BA}}=\beta$, and overlap is neglected by setting $S=1$, the unit matrix (with 1 on the diagonal and 0 elsewhere). Then the first two matrices in eqn 9 E .6 become

$$
H=\left(\begin{array}{ll}
\alpha & \beta \\
\beta & \alpha
\end{array}\right) \quad S=\left(\begin{array}{ll}
1 & 0 \\
0 & 1
\end{array}\right)
$$

and because $S$ is now a unit matrix, multiplication by which has no effect, eqn 9E. 8 becomes

$$
H c=c E
$$

At this point, multiplication from the left by the inverse matrix $c^{-1}$ gives, after using $c^{-1} c=1$,


\begin{equation*}
c^{-1} H c=E \tag{9E.10}
\end{equation*}


The matrix $E$ is diagonal, with diagonal elements $E_{n}$, so an interpretation of this equation is that the energies are calculated by finding a transformation of $H$, its conversion to $c^{-1} \mathrm{Hc}$, that makes it diagonal. This procedure is called matrix diagonalization. The columns of the matrix $c$ that brings about this diagonalization are the coefficients of the orbitals used as the basis set, and give the composition of the molecular orbitals.

\subsection*{Example 9E. 1 Finding molecular orbitals by matrix}
 diagonalizationSet up and solve the matrix equations within the Hückel approximation for the $\pi$ orbitals of butadiene (1).\\
\includegraphics{smile-a872ecbaeacdcafcd5678f1add4567f1dcc236d1}

1 Butadiene\\
Collect your thoughts The matrices are four-dimensional for this four-atom system. You need to construct the matrix $H$ by using the Hückel approximation and the parameters $\alpha$ and $\beta$. Once you have the hamiltonian matrix, you need to find the matrix $c$ that diagonalizes it: for this step, use mathematical\\
software. Full details are given in The chemist's toolkit 25, but note that if $H=\alpha 1+M$, where $M$ is a non-diagonal matrix, then because $\alpha c^{-1} 1 c=\alpha c^{-1} c=\alpha 1$, whatever matrix $c$ diagonalizes $M$ leaves $\alpha \mathbf{1}$ unchanged, so to achieve the overall diagonalization of $H$ you need to diagonalize only $M$.

The solution With C atoms labelled A, B, C, and D, the hamiltonian matrix $H$ is

$$
H=\left(\begin{array}{cccc}
\overbrace{H_{\mathrm{AA}}}^{\alpha} & \overbrace{H_{\mathrm{AB}}}^{\beta} & \overbrace{H_{\mathrm{AC}}}^{0} & \overbrace{H_{\mathrm{AD}}}^{0} \\
H_{\mathrm{BA}} & H_{\mathrm{BB}} & H_{\mathrm{BC}} & H_{\mathrm{BD}} \\
H_{\mathrm{CA}} & H_{\mathrm{CB}} & H_{\mathrm{CC}} & H_{\mathrm{CD}} \\
H_{\mathrm{DA}} & H_{\mathrm{DB}} & H_{\mathrm{DC}} & H_{\mathrm{DD}}
\end{array}\right) \xrightarrow{\overbrace{\text { Hückel }}^{\text {approximation }}}\left(\begin{array}{cccc}
\alpha & \beta & 0 & 0 \\
\beta & \alpha & \beta & 0 \\
0 & \beta & \alpha & \beta \\
0 & 0 & \beta & \alpha
\end{array}\right)
$$

which is written as

$$
H=\alpha \mathbf{1}+\beta \overbrace{\left(\begin{array}{cccc}
0 & 1 & 0 & 0 \\
1 & 0 & 1 & 0 \\
0 & 1 & 0 & 1 \\
0 & 0 & 1 & 0
\end{array}\right)}^{M}
$$

The diagonalized form of the matrix $M$ (using software) is

$$
\left(\begin{array}{cccc}
+1.62 & 0 & 0 & 0 \\
0 & +0.62 & 0 & 0 \\
0 & 0 & -0.62 & 0 \\
0 & 0 & 0 & -1.62
\end{array}\right)
$$

so the diagonalized hamiltonian matrix is

$$
E=\left(\begin{array}{cccc}
\alpha+1.62 \beta & 0 & 0 & 0 \\
0 & \alpha+0.62 \beta & 0 & 0 \\
0 & 0 & \alpha-0.62 \beta & 0 \\
0 & 0 & 0 & \alpha-1.62 \beta
\end{array}\right)
$$

The matrix that achieves the diagonalization is

$$
\boldsymbol{c}=\left(\begin{array}{rrrr}
0.372 & 0.602 & 0.602 & 0.372 \\
0.602 & 0.372 & -0.372 & -0.602 \\
0.602 & -0.372 & -0.372 & 0.602 \\
0.372 & -0.602 & 0.602 & -0.372
\end{array}\right)
$$

with each column giving the coefficients of the atomic orbitals for the corresponding molecular orbital. It follows that the energies and molecular orbitals are

$$
\begin{array}{ll}
E_{1}=\alpha+1.62 \beta & \psi_{1}=0.372 \psi_{\mathrm{A}}+0.602 \psi_{\mathrm{B}}+0.602 \psi_{\mathrm{C}}+0.372 \psi_{\mathrm{D}} \\
E_{2}=\alpha+0.62 \beta & \psi_{2}=0.602 \psi_{\mathrm{A}}+0.372 \psi_{\mathrm{B}}-0.372 \psi_{\mathrm{C}}-0.602 \psi_{\mathrm{D}} \\
E_{3}=\alpha-0.62 \beta & \psi_{3}=0.602 \psi_{\mathrm{A}}-0.372 \psi_{\mathrm{B}}-0.372 \psi_{\mathrm{C}}+0.602 \psi_{\mathrm{D}} \\
E_{4}=\alpha-1.62 \beta & \psi_{4}=0.372 \psi_{\mathrm{A}}-0.602 \psi_{\mathrm{B}}+0.602 \psi_{\mathrm{C}}-0.372 \psi_{\mathrm{D}}
\end{array}
$$

\subsection*{The chemist's toolkit 25 Matrix methods for solving eigenvalue equations}
In matrix form, an eigenvalue equation is

$$
M x=\lambda x
$$

Eigenvalue equation\\
(1a)\\
where $M$ is a square matrix with $n$ rows and $n$ columns, $\lambda$ is a constant, the eigenvalue, and $x$ is the eigenvector, an $n \times 1$ (column) matrix that satisfies the conditions of the eigenvalue equation and has the form:

$$
\boldsymbol{x}=\left(\begin{array}{c}
x_{1} \\
x_{2} \\
\vdots \\
x_{n}
\end{array}\right)
$$

In general, there are $n$ eigenvalues $\lambda^{(\mathrm{i})}, i=1,2, \ldots, n$, and $n$ corresponding eigenvectors $\boldsymbol{x}^{(\mathrm{i})}$. Equation la can be rewritten as


\begin{equation*}
(M-\lambda 1) x=0 \tag{1b}
\end{equation*}


where 1 is an $n \times n$ unit matrix, and where the property $1 x=x$ has been used. This equation has a solution only if the determinant $|M-\lambda 1|$ of the matrix $M-\lambda 1$ is zero. It follows that the $n$ eigenvalues may be found from the solution of the secular equation:


\begin{equation*}
|M-\lambda 1|=0 \tag{2}
\end{equation*}


The $n$ eigenvalues found by solving the secular equations are used to find the corresponding eigenvectors. To do so, begin by considering an $n \times n$ matrix $X$ the columns of which are formed from the eigenvectors corresponding to all the eigenvalues. Thus, if the eigenvalues are $\lambda_{1}, \lambda_{2}, \ldots$, and the corresponding eigenvectors are

\[
\boldsymbol{x}^{(1)}=\left(\begin{array}{c}
x_{1}^{(1)}  \tag{3a}\\
x_{2}^{(1)} \\
\vdots \\
x_{n}^{(1)}
\end{array}\right) \quad \boldsymbol{x}^{(2)}=\left(\begin{array}{c}
x_{1}^{(2)} \\
x_{2}^{(2)} \\
\vdots \\
x_{n}^{(2)}
\end{array}\right) \quad \cdots \boldsymbol{x}^{(n)}=\left(\begin{array}{c}
x_{1}^{(n)} \\
x_{2}^{(n)} \\
\vdots \\
x_{n}^{(n)}
\end{array}\right)
\]

where the C2p atomic orbitals are denoted by $\psi_{A}, \ldots, \psi_{\mathrm{D}}$. The molecular orbitals are mutually orthogonal and, with overlap neglected, normalized.

Comment. Note that $\psi_{1}, \ldots, \psi_{4}$ correspond to the $1 \pi, \ldots, 4 \pi$ molecular orbitals of butadiene.

Self-test 9E. 1 Repeat the exercise for the allyl radical, $\cdot \mathrm{CH}_{2}-\mathrm{CH}=\mathrm{CH}_{2}$; assume that each carbon is $\mathrm{sp}^{2}$ hybridized, and take as a basis one out-of-plane 2 p orbital on each atom.\\
\includegraphics[max width=\textwidth, center]{2025_02_27_d4b95d9718f0b9e2e6b9g-407(1)}\\
\includegraphics[max width=\textwidth, center]{2025_02_27_d4b95d9718f0b9e2e6b9g-407}

\subsection*{9E. 2 Applications}
Although the Hückel method is very primitive, it can be used to account for some of the properties of conjugated polyenes.\\
then the matrix $X$ is

\[
\boldsymbol{X}=\left(\boldsymbol{x}^{(1)} \boldsymbol{x}^{(2)} \cdots \boldsymbol{x}^{(n)}\right)=\left(\begin{array}{cccc}
x_{1}^{(1)} & x_{1}^{(2)} & \cdots & x_{1}^{(n)}  \tag{3b}\\
x_{2}^{(1)} & x_{2}^{(2)} & \cdots & x_{2}^{(n)} \\
\vdots & \vdots & & \vdots \\
x_{n}^{(1)} & x_{n}^{(2)} & \cdots & x_{n}^{(n)}
\end{array}\right)
\]

Similarly, form an $n \times n$ matrix $\Lambda$ with the eigenvalues $\lambda$ along the diagonal and zeroes elsewhere:

\[
\Lambda=\left(\begin{array}{cccc}
\lambda_{1} & 0 & \cdots & 0  \tag{4}\\
0 & \lambda_{2} & \cdots & 0 \\
\vdots & \vdots & & \vdots \\
0 & 0 & \cdots & \lambda_{n}
\end{array}\right)
\]

Now all the eigenvalue equations $M x^{(i)}=\lambda_{i} x^{(i)}$ may be combined into the single matrix equation


\begin{equation*}
M X=X \Lambda \tag{5}
\end{equation*}


Finally, form $X^{-1}$ from $X$ and multiply eqn 5 by it from the left:


\begin{equation*}
X^{-1} M X=X^{-1} X \Lambda=\Lambda \tag{6}
\end{equation*}


A structure of the form $X^{-1} M X$ is called a similarity transformation. In this case the similarity transformation $X^{-1} M X$ makes $M$ diagonal (because $\Lambda$ is diagonal). It follows that if the matrix $X$ that causes $X^{-1} M X$ to be diagonal is known, then the problem is solved: the diagonal matrix so produced has the eigenvalues as its only nonzero elements, and the matrix $X$ used to bring about the transformation has the corresponding eigenvectors as its columns. In practice, the eigenvalues and eigenvectors are obtained by using mathematical software.

\subsection*{(a) $\pi$-Electron binding energy}
As seen in Example 9E.1, the energies of the four LCAO-MOs for butadiene are


\begin{equation*}
E=\alpha \pm 1.62 \beta, \alpha \pm 0.62 \beta \tag{9E.11}
\end{equation*}


These orbitals and their energies are drawn in Fig. 9E.2. Note that:

\begin{itemize}
  \item The greater the number of internuclear nodes, the higher the energy of the orbital.
  \item There are four electrons to accommodate, so the groundstate configuration is $1 \pi^{2} 2 \pi^{2}$.
  \item The frontier orbitals of butadiene are the $2 \pi$ orbital (the HOMO, which is largely bonding) and the $3 \pi$ orbital (the LUMO, which is largely antibonding).\\
'Largely bonding' means that an orbital has both bonding and antibonding interactions between various neighbours, but the\\
\includegraphics[max width=\textwidth, center]{2025_02_27_d4b95d9718f0b9e2e6b9g-408(1)}
\end{itemize}

Figure 9E. 2 The Hückel molecular orbital energy levels of butadiene and the top view of the corresponding $\pi$ orbitals. The four $p$ electrons (one supplied by each C) occupy the two lower $\pi$ orbitals. Note that all the orbitals are delocalized.\\
bonding effects dominate. 'Largely antibonding' indicates that the antibonding effects dominate.

An important point emerges by calculating the total $\pi$-electron binding energy, $E_{\pi}$, the sum of the energies of each $\pi$ electron, and comparing it with the value for ethene. In ethene the $\pi$-electron binding energy is

$$
E_{\pi}=2(\alpha+\beta)=2 \alpha+2 \beta
$$

In butadiene it is

$$
E_{\pi}=2(\alpha+1.62 \beta)+2(\alpha+0.62 \beta)=4 \alpha+4.48 \beta
$$

Therefore, the energy of the butadiene molecule lies lower by $0.48 \beta$ (about $115 \mathrm{~kJ} \mathrm{~mol}^{-1}$ ) than the sum of two individual $\pi$ bonds (recall that $\beta$ is negative). This extra stabilization of a conjugated system compared with a set of localized $\pi$ bonds is called the delocalization energy of the molecule.

A closely related quantity is the $\pi$-bond formation energy, $E_{\mathrm{bf}}$, the energy released when a $\pi$ bond is formed. Because the contribution of $\alpha$ is the same in the molecule as in the atoms, the $\pi$-bond formation energy can be calculated from the $\pi$-electron binding energy by writing

\[
E_{\mathrm{bf}}=E_{\pi}-N_{\mathrm{C}} \alpha \quad \begin{align*}
& \pi \text {-Bond formation energy }  \tag{9Е.12}\\
& {[\text { definition }]}
\end{align*}
\]

where $N_{\mathrm{C}}$ is the number of carbon atoms in the molecule. The $\pi$-bond formation energy in butadiene, for instance, is $4.48 \beta$.

\subsection*{Example 9E. 2 Estimating the delocalization energy}
Use the Hückel approximation to find the energies of the $\pi$ orbitals of cyclobutadiene, and estimate the delocalization energy.

Collect your thoughts Set up the hamiltonian matrix using the same basis as for butadiene, but note that atoms A and D are also now neighbours. Then diagonalize the matrix to find the energies. For the delocalization energy, subtract from the total $\pi$-bond energy the energy of two $\pi$-bonds.

The solution The hamiltonian matrix is

$$
\begin{aligned}
H & =\left(\begin{array}{cccc}
\alpha & \beta & 0 & \beta \\
\beta & \alpha & \beta & 0 \\
0 & \beta & \alpha & \beta \\
\beta & 0 & \beta & \alpha
\end{array}\right) \\
& =\alpha \mathbf{1}+\beta\left(\begin{array}{cccc}
0 & 1 & 0 & 1 \\
1 & 0 & 1 & 0 \\
0 & 1 & 0 & 1 \\
1 & 0 & 1 & 0
\end{array}\right) \xrightarrow{\text { Diagonalize }}\left(\begin{array}{cccc}
2 & 0 & 0 & 0 \\
0 & 0 & 0 & 0 \\
0 & 0 & 0 & 0 \\
0 & 0 & 0 & -2
\end{array}\right)
\end{aligned}
$$

Diagonalization gives the energies of the orbitals as

$$
E=\alpha+2 \beta, \alpha, \alpha, \alpha-2 \beta
$$

Four electrons must be accommodated. Two occupy the lowest orbital (of energy $\alpha+2 \beta$ ), and two occupy the doubly degenerate orbitals (of energy $\alpha$ ). The total energy is therefore $4 \alpha+4 \beta$. Two isolated $\pi$ bonds would have an energy $4 \alpha+4 \beta ;$ therefore, in this case, the delocalization energy is zero.

Self-test 9E. 2 Repeat the calculation for benzene (use software!).\\
\includegraphics[max width=\textwidth, center]{2025_02_27_d4b95d9718f0b9e2e6b9g-408(2)}

\subsection*{(b) Aromatic stability}
The most notable example of delocalization conferring extra stability is benzene and the aromatic molecules based on its structure. In elementary accounts, the structure of benzene, and other aromatic compounds, is often expressed in a mixture of valence-bond and molecular orbital terms, with typically va-lence-bond language (Topic 9A) used for its $\sigma$ framework and molecular orbital language used to describe its $\pi$ electrons.

First, the valence-bond component. The six C atoms are regarded as $\mathrm{sp}^{2}$ hybridized, with a single unhydridized perpendicular 2p orbital. One H atom is bonded by ( $\mathrm{Csp}^{2}, \mathrm{H} 1 \mathrm{~s}$ ) overlap to each C carbon, and the remaining hybrids overlap to give a regular hexagon of atoms (Fig. 9E.3). The internal angle of a\\
\includegraphics[max width=\textwidth, center]{2025_02_27_d4b95d9718f0b9e2e6b9g-408}

Figure 9E. 3 The $\sigma$ framework of benzene is formed by the overlap of $\mathrm{Csp}^{2}$ hybrids, which fit without strain into a hexagonal arrangement.\\
regular hexagon is $120^{\circ}$, so $\mathrm{sp}^{2}$ hybridization is ideally suited for forming $\sigma$ bonds. The hexagonal shape of benzene permits strain-free $\sigma$ bonding.

Now consider the molecular orbital component of the description. The six C2p orbitals overlap to give six $\pi$ orbitals that spread all round the ring. Their energies are calculated within the Hückel approximation by diagonalizing the hamiltonian matrix

$$
H=\left(\begin{array}{llllll}
\alpha & \beta & 0 & 0 & 0 & \beta \\
\beta & \alpha & \beta & 0 & 0 & 0 \\
0 & \beta & \alpha & \beta & 0 & 0 \\
0 & 0 & \beta & \alpha & \beta & 0 \\
0 & 0 & 0 & \beta & \alpha & \beta \\
\beta & 0 & 0 & 0 & \beta & \alpha
\end{array}\right)
$$

$=\alpha \mathbf{1}+\beta\left(\begin{array}{cccccc}0 & 1 & 0 & 0 & 0 & 1 \\ 1 & 0 & 1 & 0 & 0 & 0 \\ 0 & 1 & 0 & 1 & 0 & 0 \\ 0 & 0 & 1 & 0 & 1 & 0 \\ 0 & 0 & 0 & 1 & 0 & 1 \\ 1 & 0 & 0 & 0 & 1 & 0\end{array}\right) \xrightarrow{\text { Diagonalize }}\left(\begin{array}{cccccc}2 & 0 & 0 & 0 & 0 & 0 \\ 0 & 1 & 0 & 0 & 0 & 0 \\ 0 & 0 & 1 & 0 & 0 & 0 \\ 0 & 0 & 0 & -1 & 0 & 0 \\ 0 & 0 & 0 & 0 & -1 & 0 \\ 0 & 0 & 0 & 0 & 0 & -2\end{array}\right)$\\
The MO energies, the diagonal elements of this matrix, are


\begin{equation*}
E=\alpha \pm 2 \beta, \alpha \pm \beta, \alpha \pm \beta \tag{9E.13}
\end{equation*}


as shown in Fig. 9E.4. The orbitals there have been given symmetry labels that are explained in Topic 10B. Note that the lowest energy orbital is bonding between all neighbouring atoms, the highest energy orbital is antibonding between each\\
\includegraphics[max width=\textwidth, center]{2025_02_27_d4b95d9718f0b9e2e6b9g-409}

Figure 9E. 4 The Hückel orbitals of benzene and the corresponding energy levels. The orbital labels are explained in Topic 10B. The bonding and antibonding character of the delocalized orbitals reflects the numbers of nodes between the atoms. In the ground state, only the bonding orbitals are occupied.\\
pair of neighbours, and the intermediate orbitals are a mixture of bonding, nonbonding, and antibonding character between adjacent atoms.

Now apply the building-up principle to the $\pi$ system. There are six electrons to accommodate (one from each C atom), so the three lowest orbitals ( $\mathrm{a}_{2 \mathrm{u}}$ and the doubly-degenerate pair $\mathrm{e}_{1 \mathrm{~g}}$ ) are fully occupied, giving the ground-state configuration $\mathrm{a}_{2 \mathrm{u}}^{2} \mathrm{e}_{\mathrm{lg}}^{4}$. A significant point is that the only molecular orbitals occupied are those with net bonding character (the analogy with the strongly bonded $\mathrm{N}_{2}$ molecule, Topic 9B, should be noted).

The $\pi$-electron binding energy of benzene is

$$
E_{\pi}=2(\alpha+2 \beta)+4(\alpha+\beta)=6 \alpha+8 \beta
$$

If delocalization is ignored and the molecule is thought of as having three isolated $\pi$ bonds, it would be ascribed a $\pi$-electron energy of only $3(2 \alpha+2 \beta)=6 \alpha+6 \beta$. The delocalization energy is therefore $2 \beta \approx-480 \mathrm{~kJ} \mathrm{~mol}^{-1}$, which is considerably more than for butadiene. The $\pi$-bond formation energy in benzene is $8 \beta$.

This discussion suggests that aromatic stability can be traced to two main contributions. First, the shape of the regular hexagon is ideal for the formation of strong $\sigma$ bonds: the $\sigma$ framework is relaxed and without strain. Second, the $\pi$ orbitals are such as to be able to accommodate all the electrons in bonding orbitals, and the delocalization energy is large.

\subsection*{Brief illustration 9E. 2}
The energies of the four molecular orbitals of cyclobutadiene are $E=\alpha \pm 2 \beta, \alpha, \alpha$ (see Example 9E.2). There are four $\pi$ electrons to accommodate in $\mathrm{C}_{4} \mathrm{H}_{4}$, so the total $\pi$-electron binding energy is $2(\alpha+2 \beta)+2 \alpha=4(\alpha+\beta)$. The energy of two localized $\pi$-bonds is $4(\alpha+\beta)$. Therefore, the delocalization energy is zero, so the molecule is not aromatic. There are only two $\pi$ electrons to accommodate in $\mathrm{C}_{4} \mathrm{H}_{4}^{2+}$, so the total $\pi$-electron binding energy is $2(\alpha+2 \beta)=2 \alpha+4 \beta$. The energy of a single localized $\pi$-bond is $2(\alpha+\beta)$, so the delocalization energy is $2 \beta$ and the molecule-ion is aromatic.

\subsection*{9е. 3 Computational chemistry}
The severe assumptions of the Hückel method are now easy to avoid by using a variety of software packages that can be used not only to calculate the shapes and energies of molecular orbitals but also to predict with reasonable accuracy the structure and reactivity of molecules. The full treatment of molecular electronic structure has received an enormous amount of attention by chemists and has become a keystone of modern chemical research. However, the calculations are very complex, and all this section seeks to do is to provide a brief\\
introduction. ${ }^{1}$ In every case, the procedures focus on the calculation or estimation of integrals like $H_{\mathrm{JJ}}$ and $H_{\mathrm{IJ}}$ rather than setting them equal to the constants $\alpha$ or $\beta$, or ignoring them entirely.

In all cases the Schrödinger equation is solved iteratively and self-consistently, just as for the self-consistent field (SCF) approach to atoms (Topic 8B). First, the molecular orbitals for the electrons present in the molecule are formulated as LCAOs. One molecular orbital is then selected and all the others are used to set up an expression for the potential energy of an electron in the chosen orbital. The resulting Schrödinger equation is then solved numerically to obtain a better version of the chosen molecular orbital and its energy. The procedure is repeated for all the molecular orbitals and used to calculate the total energy of the molecule. The process is repeated until the computed orbitals and energy are constant to within some tolerance.

\subsection*{(a) Semi-empirical and ab initio methods}
In a semi-empirical method, many of the integrals are estimated by appealing to spectroscopic data or physical properties such as ionization energies, and using a series of rules to set certain integrals equal to zero. A primitive form of this procedure is used in Brief illustration 9D. 1 of Topic 9D where the integral $\alpha$ is identified with a combination of the ionization energy and electron affinity of an atom. In an $\boldsymbol{a b}$ initio method an attempt is made to calculate all the integrals, including overlap integrals. Both procedures employ a great deal of computational effort. The integrals that are required involve atomic orbitals that in general may be centred on different nuclei. It can be appreciated that, if there are several dozen atomic orbitals used to build the molecular orbitals, then there will be tens of thousands of integrals of this form to evaluate (the number of integrals increases as the fourth power of the number of atomic orbitals in the basis, so even for a 10 -atom molecule there are $10^{4}$ integrals to evaluate). Some kind of approximation scheme is necessary.

One severe semi-empirical approximation used in the early days of computational chemistry was called complete neglect of differential overlap (CNDO), in which all molecular integrals of the form

$$
j_{0} \iint \psi_{\mathrm{A}}\left(\boldsymbol{r}_{1}\right) \psi_{\mathrm{B}}\left(\boldsymbol{r}_{1}\right) \frac{1}{r_{12}} \psi_{\mathrm{C}}\left(\boldsymbol{r}_{2}\right) \psi_{\mathrm{D}}\left(\boldsymbol{r}_{2}\right) \mathrm{d} \tau_{1} \mathrm{~d} \tau_{2}
$$

are set to zero unless $\psi_{\mathrm{A}}$ and $\psi_{\mathrm{B}}$ are the same orbitals centred on the same nucleus, and likewise for $\psi_{\mathrm{C}}$ and $\psi_{\mathrm{D}}$. The surviving integrals are then adjusted until the energy levels are in good agreement with experiment or the computed enthalpy of formation of the compound is in agreement with experiment.

\footnotetext{${ }^{1}$ A more complete account with detailed examples will be found in our companion volume, Physical chemistry: Quanta, matter, and change (2014).
}
\includegraphics[max width=\textwidth, center]{2025_02_27_d4b95d9718f0b9e2e6b9g-410}

Figure 9E. 5 The product of two Gaussian functions on different centres is itself a Gaussian function located at a point between the two contributing Gaussians. The scale of the product has been increased relative to that of its two components.

More recent semi-empirical methods make less severe decisions about which integrals are to be ignored, but they are all descendants of the early CNDO technique.

Commercial packages are also available for $a b$ initio calculations. Here the problem is to evaluate as efficiently as possible thousands of integrals that arise from the Coulombic interaction between two electrons like that displayed above, with the possibility that each of the atomic orbitals is centred on a different atom, a so-called 'four-centre integral'. This task is greatly facilitated by expressing the atomic orbitals used in the LCAOs as linear combinations of Gaussian orbitals. A Gaussian type orbital (GTO) is a function of the form $\mathrm{e}^{-r^{2}}$. The advantage of GTOs over the correct orbitals (which for hydrogenic systems are proportional to exponential functions of the form $\mathrm{e}^{-r}$ ) is that the product of two Gaussian functions is itself a Gaussian function that lies between the centres of the two contributing functions (Fig. 9E.5). In this way, the four-centre integrals become two-centre integrals of the form

$$
j_{0} \iint X\left(\boldsymbol{r}_{1}\right) \frac{1}{r_{12}} Y\left(\boldsymbol{r}_{2}\right) \mathrm{d} \tau_{1} \mathrm{~d} \tau_{2}
$$

where $X$ is the Gaussian corresponding to the product $\psi_{\mathrm{A}} \psi_{\mathrm{B}}$, and $Y$ is the corresponding Gaussian from $\psi_{\mathrm{C}} \psi_{\mathrm{D}}$. Integrals of this form are much easier and faster to evaluate numerically than the original four-centre integrals. Although more GTOs have to be used to simulate the atomic orbitals, there is an overall increase in speed of computation.

\subsection*{Brief illustration 9E. 3}
Consider a one-dimensional 'homonuclear' system, with GTOs of the form $\mathrm{e}^{-a x^{2}}$ located at 0 and $R$. Then one of the integrals that would have to be evaluated would include the term

$$
\psi_{\mathrm{A}}\left(\boldsymbol{r}_{1}\right) \psi_{\mathrm{B}}\left(\boldsymbol{r}_{1}\right)=\mathrm{e}^{-a x^{2}} \mathrm{e}^{-a(x-R)^{2}}=\mathrm{e}^{-2 a x^{2}+2 a x k-a R^{2}}
$$

Next note that $-2 a\left(x-\frac{1}{2} R\right)^{2}=-2 a x^{2}+2 a x R-\frac{1}{2} a R^{2}$, so

$$
\psi_{\mathrm{A}}\left(r_{1}\right) \psi_{\mathrm{B}}\left(r_{1}\right)=\mathrm{e}^{-2 a(x-R / 2)^{2}-a R^{2} / 2}=\mathrm{e}^{-2 a(x-R / 2)^{2}} \mathrm{e}^{-a R^{2} / 2}
$$

which is proportional to a single Gaussian (the term in blue) centred on the mid-point of the internuclear separation, at $x=\frac{1}{2} R$.

\subsection*{(b) Density functional theory}
A technique that has gained considerable ground in recent years to become one of the most widely used techniques for the calculation of molecular structure is density functional theory (DFT). Its advantages include less demanding computational effort, less computer time, and-in some cases (particularly d-metal complexes)-better agreement with experimental values than is obtained from other procedures.

The central focus of DFT is the electron density, $\rho$, rather than the wavefunction, $\psi$. The 'functional' part of the name comes from the fact that the energy of the molecule is a function of the electron density, written $E[\rho]$, and the electron density is itself a function of position, $\rho(r)$ : in mathematics a function of a function is called a 'functional'. The occupied orbitals are used to construct the electron density from


\begin{equation*}
\rho(r)=\sum_{m, o c c u p i e d}\left|\psi_{m}(r)\right|^{2} \quad \text { Electron probability density } \tag{9E.14}
\end{equation*}


and are calculated from modified versions of the Schrödinger equation known as the Kohn-Sham equations.

The Kohn-Sham equations are solved iteratively and selfconsistently. First, the electron density is guessed. For this step it is common to use a superposition of atomic electron densities. Next, the Kohn-Sham equations are solved to obtain an initial set of orbitals. This set of orbitals is used to obtain a better approximation to the electron density and the process is repeated until the density and the computed energy are constant to within some tolerance.

\subsection*{(c) Graphical representations}
One of the most significant developments in computational chemistry has been the introduction of graphical representations of molecular orbitals and electron densities. The raw output of a molecular structure calculation is a list of the coefficients of the atomic orbitals in each molecular orbital and the energies of these orbitals. The graphical representation of a molecular orbital uses stylized shapes to represent the basis set, and then scales their size to indicate the coefficient in the linear combination. Different signs of the wavefunctions are represented by different colours.

Once the coefficients are known, it is possible to construct a representation of the electron density in the molecule by\\
\includegraphics[max width=\textwidth, center]{2025_02_27_d4b95d9718f0b9e2e6b9g-411}

Figure 9E. 6 Various representations of an isodensity surface of ethanol: (a) solid surface, (b) transparent surface, and (c) mesh surface.\\
noting which orbitals are occupied and then forming the squares of those orbitals. The total electron density at any point is then the sum of the squares of the wavefunctions evaluated at that point. The outcome is commonly represented by an isodensity surface, a surface of constant total electron density (Fig. 9E.6). As shown in the illustration, there are several styles of representing an isodensity surface, as a solid form, as a transparent form with a ball-and-stick representation of the molecule within, or as a mesh. A related representation is a solvent-accessible surface in which the shape represents the shape of the molecule by imagining a sphere representing a solvent molecule rolling across the surface and plotting the locations of the centre of that sphere.

One of the most important aspects of a molecule other than its geometrical shape is the distribution of charge over its surface, which is commonly depicted as an electrostatic potential surface (an 'elpot surface'). The potential energy, $E_{\mathrm{p}}$, of an imaginary positive charge $Q$ at a point is calculated by taking into account its interaction with the nuclei and the electron density throughout the molecule. Then, because $E_{\mathrm{p}}=Q \phi$, where $\phi$ is the electric potential, the potential energy can be interpreted as a potential and depicted as an appropriate colour (Fig. 9E.7). Electron-rich regions usually have negative potentials and electron-poor regions usually have positive potentials.

Representations such as those illustrated here are of critical importance in a number of fields. For instance, they may be used to identify an electron-poor region of a molecule that is susceptible to association with or chemical attack by an electron-rich region of another molecule. Such considerations are important for assessing the pharmacological activity of potential drugs.\\
\includegraphics[max width=\textwidth, center]{2025_02_27_d4b95d9718f0b9e2e6b9g-411(1)}

Figure 9E. 7 An elpot diagram of ethanol; the molecule has the same orientation as in Fig. 9E.6. Red denotes regions of negative electrostatic potential and blue regions of positive potential (as in ${ }^{\delta-} \mathrm{O}-\mathrm{H}^{\delta+}$ ).

\subsection*{Checklist of concepts}
1. The Hückel method neglects overlap and interactions between orbitals on atoms that are not neighbours.2. The highest occupied molecular orbital (HOMO) and the lowest unoccupied molecular orbital (LUMO) are the frontier orbitals of a molecule.3. The Hückel method may be expressed in a compact manner by introducing matrices.4. The $\pi$-bond formation energy is the energy released when a $\pi$ bond is formed.5. The $\pi$-electron binding energy is the sum of the energies of each $\pi$ electron.6. The delocalization energy is the difference between the $\pi$-electron binding energy and the energy of the same molecule with localized $\pi$ bonds.7. The stability of benzene arises from the geometry of the ring and the high delocalization energy.8. Semi-empirical calculations approximate integrals by estimating them by using empirical data; ab initio methods evaluate all integrals numerically.\begin{enumerate}
  \setcounter{enumi}{8}
  \item Density functional theories develop equations based on the electron density rather than the wavefunction itself.\\
$\square$ 10. Graphical techniques are used to plot a variety of surfaces based on electronic structure calculations.
\end{enumerate}

\subsection*{Checklist of equations}
\begin{center}
\begin{tabular}{llll}
\hline
Property & Equation & Comment & Equation number \\
LCAO-MO & $\psi=\sum_{i} c_{i} \psi_{i}$ & $\psi_{i}$ are atomic orbitals & 9 E .1 \\
Hückel equations & $H c=S c E$ & \begin{tabular}{l}
Hückel approximations: $H_{\mathrm{AB}}=0$ except \\
between neighbours; $S=1$. \\
\end{tabular} & 9 E .8 \\
Diagonalization & $c^{-1} H c=E$ & Definition &  \\
$\pi$-Electron binding energy & $E_{\pi}=$ sum of energies of $\pi$ electrons & Definition; $N_{\mathrm{C}}$ is the number of carbon atoms & 9 E .12 \\
$\pi$-Bond formation energy & $E_{\mathrm{bf}}=E_{\pi}-N_{\mathrm{C}} \alpha$ &  &  \\
$\pi$-Delocalization energy & $E_{\text {deloc }}=E_{\pi}-N_{\mathrm{C}}(\alpha+\beta)$ &  &  \\
\hline
\end{tabular}
\end{center}

\chapter{FOCUS 10 Molecular symmetry}
In this Focus the concept of 'shape' is sharpened into a precise definition of 'symmetry'. As a result, symmetry and its consequences can be discussed systematically, thereby providing a very powerful tool for the prediction and analysis of molecular structure and properties.

\subsection*{10A Shape and symmetry}
This Topic shows how to classify any molecule according to its symmetry. Two immediate applications of this classification are the identification of whether or not a molecule can have an electric dipole moment (and so be polar) and whether or not it can be chiral (and so be optically active).\\
10A. 1 Symmetry operations and symmetry elements; 10A. 2 The symmetry classification of molecules; 10A. 3 Some immediate consequences of symmetry

\subsection*{10B Group theory}
The systematic treatment of symmetry is an application of 'group theory'. This theory represents the outcome of symme-\\
try operations (such as rotations and reflections) by matrices. This step is important, for once symmetry operations are expressed numerically they can be manipulated quantitatively. This Topic introduces 'character tables' which are exceptionally important in the application of group theory to chemical problems.\\
10B. 1 The elements of group theory; 10B. 2 Matrix representations; 10B. 3 Character tables

\subsection*{10C Applications of symmetry}
Group theory provides simple criteria for deciding whether certain integrals necessarily vanish. One application is to decide whether the overlap integral between two atomic orbitals is necessarily zero and therefore to decide which atomic orbitals can contribute to molecular orbitals. Symmetry is also used to identify linear combinations of atomic orbitals that match the symmetry of the nuclear framework. By considering the symmetry properties of integrals, it is also possible to derive the selection rules that govern spectroscopic transitions.\\
10C. 1 Vanishing integrals; 10C. 2 Applications to molecular orbital theory; 10C. 3 Selection rules

\section*{TOPIC 10A Shape and symmetry}
\subsection*{- Why do you need to know this material?}
Symmetry arguments can be used to make immediate assessments of the properties of molecules; the initial step is to identify the symmetry a molecule possesses and then to classify it accordingly.

\subsection*{- What is the key idea?}
Molecules can be classified into groups according to their symmetry elements.

\subsection*{> What do you need to know already?}
This Topic does not draw on others directly, but it will be useful to be aware of the shapes of a variety of simple molecules and ions encountered in introductory chemistry courses.

Some objects are 'more symmetrical' than others. A sphere is more symmetrical than a cube because it looks the same after it has been rotated through any angle about any axis passing through the centre. A cube looks the same only if it is rotated through certain angles about specific axes, such as $90^{\circ}, 180^{\circ}$, or $270^{\circ}$ about an axis passing through the centres of any of its opposite faces (Fig. 10A.1), or by $120^{\circ}$ or $240^{\circ}$ about an axis passing through any of its opposite corners. Similarly, an $\mathrm{NH}_{3}$ molecule is 'more symmetrical' than an $\mathrm{H}_{2} \mathrm{O}$ molecule because $\mathrm{NH}_{3}$ looks the same after rotations of $120^{\circ}$ or $240^{\circ}$ about\\
\includegraphics[max width=\textwidth, center]{2025_02_27_d4b95d9718f0b9e2e6b9g-420}

Figure 10A. 1 Some of the symmetry elements of a cube. The twofold, threefold, and fourfold axes are labelled with the conventional symbols.\\
\includegraphics[max width=\textwidth, center]{2025_02_27_d4b95d9718f0b9e2e6b9g-420(1)}

Figure 10A. 2 (a) $\mathrm{An}_{\mathrm{NH}}^{3}$ molecule has a threefold ( $C_{3}$ ) axis and (b) an $\mathrm{H}_{2} \mathrm{O}$ molecule has a twofold $\left(C_{2}\right)$ axis. Both have other symmetry elements too.\\
the axis shown in Fig. 10A.2, whereas $\mathrm{H}_{2} \mathrm{O}$ looks the same only after a rotation of $180^{\circ}$.

This Topic puts these intuitive notions on a more formal foundation. It will be seen that molecules can be grouped together according to their symmetry, with the tetrahedral species $\mathrm{CH}_{4}$ and $\mathrm{SO}_{4}^{2-}$ in one group and the pyramidal species $\mathrm{NH}_{3}$ and $\mathrm{SO}_{3}^{2-}$ in another. It turns out that molecules in the same group share certain physical properties, so powerful predictions can be made about whole series of molecules once the group to which they belong has been identified.

\subsection*{10A. 1 Symmetry operations and symmetry elements}
An action that leaves an object looking the same after it has been carried out is called a symmetry operation. Typical symmetry operations include rotations, reflections, and inversions. There is a corresponding symmetry element for each symmetry operation, which is the point, line, or plane with respect to which the symmetry operation is performed. For instance, a rotation (a symmetry operation) is carried out around an axis (the corresponding symmetry element). Molecules can be classified by identifying all their symmetry elements, and then grouping together molecules that possess the same set of symmetry elements. This procedure, for example, puts the trigonal planar species $\mathrm{BF}_{3}$ and $\mathrm{CO}_{3}^{2-}$ into one group and the species $\mathrm{H}_{2} \mathrm{O}$ (bent) and $\mathrm{ClF}_{3}$ (T-shaped) into another group.

An $\boldsymbol{n}$-fold rotation (the operation) about an $\boldsymbol{n}$-fold axis of symmetry, $C_{n}$ (the corresponding element) is a rotation through $360^{\circ} / n$. An $\mathrm{H}_{2} \mathrm{O}$ molecule has one twofold axis, $C_{2}$. An $\mathrm{NH}_{3}$ molecule has one threefold axis, $C_{3}$, with which is associated two symmetry operations, one being a rotation by $120^{\circ}$ in a clockwise sense, and the other a rotation by $120^{\circ}$ in an anticlockwise sense. There is only one twofold rotation associated with a $C_{2}$ axis because clockwise and anticlockwise $180^{\circ}$ rotations are identical. A pentagon has a $C_{5}$ axis, with two rotations (one clockwise, the other anticlockwise) through $72^{\circ}$ associated with it. It also has an operation denoted $C_{5}^{2}$, corresponding to two successive $C_{5}$ rotations; there are two such operations, one through $144^{\circ}$ in a clockwise sense and the other through $144^{\circ}$ in an anticlockwise sense. A cube has three $C_{4}$ axes, four $C_{3}$ axes, and six $C_{2}$ axes. However, even this high symmetry is exceeded by that of a sphere, which possesses an infinite number of symmetry axes (along any axis passing through the centre) of all possible integral values of $n$.

If a molecule possesses several rotation axes, then the one with the highest value of $n$ is called the principal axis. The principal axis of a benzene molecule is the sixfold axis perpendicular to the hexagonal ring (1). If a molecule has more than one rotation axis with this highest value of $n$, and it is wished to designate one of them as the principal axis, then it is common to choose the axis that passes through the greatest number of atoms or, in the case of a planar molecule (such as naphthalene, 2 , which has three $C_{2}$ axes competing for the title), to choose the axis perpendicular to the plane.\\
\includegraphics[max width=\textwidth, center]{2025_02_27_d4b95d9718f0b9e2e6b9g-421(3)}

1 Benzene, $\mathrm{C}_{6} \mathrm{H}_{6}$\\
\includegraphics[max width=\textwidth, center]{2025_02_27_d4b95d9718f0b9e2e6b9g-421}

2 Naphthalene, $\mathrm{C}_{10} \mathrm{H}_{8}$

A reflection is the operation corresponding to a mirror plane, $\sigma$ (the element). If the plane contains the principal axis, it is called 'vertical' and denoted $\sigma_{\mathrm{r}}$. An $\mathrm{H}_{2} \mathrm{O}$ molecule has two vertical mirror planes (Fig. 10A.3) and an $\mathrm{NH}_{3}$ molecule has three. A vertical mirror plane that bisects the angle between two $C_{2}$ axes is called a 'dihedral plane' and is denoted\\
\includegraphics[max width=\textwidth, center]{2025_02_27_d4b95d9718f0b9e2e6b9g-421(2)}

Figure 10A. $3 \mathrm{An} \mathrm{H}_{2} \mathrm{O}$ molecule has two mirror planes. They are both vertical (i.e. contain the principal axis), so are denoted $\sigma_{v}$ and $\sigma_{\mathrm{v}}^{\prime}$.\\
\includegraphics[max width=\textwidth, center]{2025_02_27_d4b95d9718f0b9e2e6b9g-421(4)}

Figure 10A.4 Dihedral mirror planes $\left(\sigma_{d}\right)$ bisect the $C_{2}$ axes perpendicular to the principal axis.\\
$\sigma_{d}$ (Fig. 10A.4). When the mirror plane is perpendicular to the principal axis it is called 'horizontal' and denoted $\sigma_{h}$. The benzene molecule has such a horizontal mirror plane, perpendicular to the $C_{6}$ (principal) axis.

In an inversion (the operation) through a centre of symmetry, $i$ (the element), each point in a molecule is imagined as being moved in a straight line to the centre of the molecule and then out the same distance on the other side; that is, the point $(x, y, z)$ is taken into the point $(-x,-y,-z)$. Neither an $\mathrm{H}_{2} \mathrm{O}$ molecule nor an $\mathrm{NH}_{3}$ molecule has a centre of inversion, but a sphere and a cube do have one. A benzene molecule has a centre of inversion, as does a regular octahedron (Fig. 10A.5); a regular tetrahedron and a $\mathrm{CH}_{4}$ molecule do not.\\
\includegraphics[max width=\textwidth, center]{2025_02_27_d4b95d9718f0b9e2e6b9g-421(1)}

Figure 10A. 5 A regular octahedron has a centre of inversion (i).\\
\includegraphics[max width=\textwidth, center]{2025_02_27_d4b95d9718f0b9e2e6b9g-422}

Figure 10A.6 (a) $\mathrm{A} \mathrm{CH}_{4}$ molecule has a fourfold improper rotation axis $\left(S_{4}\right)$ : the molecule is indistinguishable after a $90^{\circ}$ rotation followed by a reflection across the horizontal plane, but neither operation alone is a symmetry operation. (b) The staggered form of ethane has an $S_{6}$ axis composed of a $60^{\circ}$ rotation followed by a reflection.

An $n$-fold improper rotation (the operation) about an $\boldsymbol{n}$-fold axis of improper rotation or an $\boldsymbol{n}$-fold improper rotation axis, $S_{n}$ (the symmetry element) is composed of two successive transformations. The first is a rotation through $360 \%$, and the second is a reflection through a plane perpendicular to the axis of that rotation. Neither transformation alone needs to be a symmetry operation. $\mathrm{A} \mathrm{CH}_{4}$ molecule has three $S_{4}$ axes, and the staggered conformation of ethane has an $S_{6}$ axis (Fig. 10A.6).\\
The identity, $E$, consists of doing nothing; the corresponding symmetry element is the entire object. Because every molecule is indistinguishable from itself if nothing is done to it, every object possesses at least the identity element. One reason for including the identity is that some molecules have only this symmetry element (3).

\subsection*{Brief illustration 10A. 1}
To identify the symmetry elements of a naphthalene molecule (2), note that:

\begin{itemize}
  \item Like all molecules, it has the identity element, $E$.
  \item There are three twofold axes of rotation, $C_{2}$ : one perpendicular to the plane of the molecule, and two others lying in the plane.
  \item With the $C_{2}$ axis perpendicular to the plane of the molecule chosen as the principal axis, there is a $\sigma_{\mathrm{h}}$ plane perpendicular to the principal axis, and two $\sigma_{\mathrm{v}}$ planes which contain the principal axis.
  \item There is also a centre of inversion, $i$, through the midpoint of the molecule, which is mid-way along the $\mathrm{C}-\mathrm{C}$ bond at the ring junction.
\end{itemize}

\subsection*{10A. 2 The symmetry classification of molecules}
Objects are classified into groups according to the symmetry elements they possess. Point groups arise when objects are classified according to symmetry elements that correspond to operations leaving at least one common point unchanged. The five kinds of symmetry element identified so far are of this kind. When crystals are considered (Topic 15A), symmetries arising from translation through space also need to be taken into account, and the classification according to these elements gives rise to the more extensive space groups.

All molecules with the same set of symmetry elements belong to the same point group, and the name of the group is determined by this set of symmetry elements. There are two systems of notation (Table 10A.1). The Schoenflies system (in which a name looks like $C_{4 v}$ ) is more common for the discussion of individual molecules, and the Hermann-Mauguin system, or International system (in which a name looks like 4 mm ), is used almost exclusively in the discussion of crystal symmetry. The identification of the point group to which a molecule belongs (in the Schoenflies system) is simplified by referring to the flow diagram in Fig. 10A. 7 and to the shapes shown in Fig. 10A.8.

Table 10A. 1 The notations for point groups*

\begin{center}
\begin{tabular}{llllllllll}
\hline
$C_{\mathrm{i}}$ & $\overline{1}$ &  &  &  &  &  &  &  &  \\
$C_{\mathrm{s}}$ & $m$ &  &  &  &  &  &  &  &  \\
$C_{1}$ & 1 & $C_{2}$ & 2 & $C_{3}$ & 3 & $C_{4}$ & 4 & $C_{6}$ & 6 \\
 &  & $C_{2 \mathrm{v}}$ & $2 m m$ & $C_{3 \mathrm{v}}$ & $3 m$ & $C_{4 \mathrm{v}}$ & $4 m m$ & $C_{6 \mathrm{v}}$ & 6 mm \\
 &  & $C_{2 \mathrm{~h}}$ & $2 / m$ & $C_{3 \mathrm{~h}}$ & $\overline{6}$ & $C_{4 \mathrm{~h}}$ & $4 / m$ & $C_{6 \mathrm{~h}}$ & $6 / m$ \\
 &  & $D_{2}$ & 222 & $D_{3}$ & 32 & $D_{4}$ & 422 & $D_{6}$ & 622 \\
 &  & $D_{2 \mathrm{~h}}$ & $m m m$ & $D_{3 \mathrm{~h}}$ & $\overline{6} 2 m$ & $D_{4 \mathrm{~h}}$ & $4 / m m m$ & $D_{6 \mathrm{~h}}$ & $6 / m m m$ \\
 &  & $D_{2 \mathrm{~d}}$ & $\overline{4} 2 m$ & $D_{3 \mathrm{~d}}$ & $\overline{3} m$ & $S_{4}$ & $\overline{4}$ & $S_{6}$ & $\overline{3}$ \\
$T$ & 23 & $T_{\mathrm{d}}$ & $\overline{4} 3 m$ & $T_{\mathrm{h}}$ & $m 3$ &  &  &  &  \\
$O$ & 432 & $O_{\mathrm{h}}$ & $m 3 m$ &  &  &  &  &  &  \\
\hline
\end{tabular}
\end{center}

\footnotetext{\begin{itemize}
  \item Schoenflies notation in black, Hermann-Mauguin (International system) in blue. In the Hermann-Mauguin system, a number $n$ denotes the presence of an $n$-fold axis and $m$ denotes a mirror plane. A slash (/) indicates that the mirror plane is perpendicular to the symmetry axis. It is important to distinguish symmetry elements of the same type but of different classes, as in $4 / \mathrm{mmm}$, in which there are three classes of mirror plane. A bar over a number indicates that the element is combined with an inversion. The only groups listed here are the so-called 'crystallographic point groups'.
\end{itemize}
}
\includegraphics[max width=\textwidth, center]{2025_02_27_d4b95d9718f0b9e2e6b9g-423(3)}

Figure 10A. 7 A flow diagram for determining the point group of a molecule. Start at the top and answer the question posed in each diamond ( $Y=$ yes, $N=n o$ ). The blue line refers to the path taken in Brief illustration 10A.2.

\subsection*{Brief illustration 10A. 2}
To identify the point group to which a ruthenocene molecule (4) belongs, first identify the symmetry elements present and the use the flow diagram in Fig. 10A.7. Note that:

\begin{itemize}
  \item The molecule has a fivefold axis, and five twofold axes which pass through the Ru and are perpendicular to the $C_{5}$ axis.
  \item There is a mirror plane, $\sigma_{h}$, perpendicular to the $C_{5}$ axis and passing through the Ru.
  \item There are five $\sigma_{\mathrm{v}}$ planes containing the principal axis: each passes through one carbon in a ring and the midpoint of the C-C bond on the opposite side. Each one of these planes contains one of the twofold axes.
\end{itemize}

The path to trace in Fig. 10A. 7 is shown by a blue line; it ends at $D_{n \mathrm{~h}}$, and because the molecule has a fivefold axis, it belongs to the point group $D_{5 h}$.\\
\includegraphics[max width=\textwidth, center]{2025_02_27_d4b95d9718f0b9e2e6b9g-423}

4 Ruthenocene, $\mathrm{Ru}(\mathrm{Cp})_{2}$\\
If the rings are staggered, as they are in an excited state of ferrocene (5), the $\sigma_{\mathrm{h}}$ plane is absent. The other mirror planes are still present, but now they bisect the angles between twofold axes and so are described as $\sigma_{\mathrm{d}}$. Tracing the appropriate path in Fig. 10A. 7 gives the point group as $D_{5 d}$.\\
\includegraphics[max width=\textwidth, center]{2025_02_27_d4b95d9718f0b9e2e6b9g-423(1)}

5 Ferrocene, $\mathrm{Fe}(\mathrm{Cp})_{2}$ (excited state)\\
\includegraphics[max width=\textwidth, center]{2025_02_27_d4b95d9718f0b9e2e6b9g-423(2)}

Figure 10A. 8 A summary of the shapes corresponding to different point groups. The group to which a molecule belongs can often be identified from this diagram without going through the formal procedure in Fig. 10A.7.

\subsection*{(a) The groups $C_{1}, C_{i}$, and $C_{s}$}
\begin{center}
\begin{tabular}{ll}
\hline
Name & Elements \\
\hline
$C_{1}$ & $E$ \\
$C_{i}$ & $E, i$ \\
$C_{s}$ & $E, \sigma$ \\
\hline
\end{tabular}
\end{center}

A molecule belongs to the group $C_{1}$ if it has no element other than the identity. It belongs to $C_{\mathrm{i}}$ if it has the identity and the inversion alone, and to $C_{\mathrm{s}}$ if it has the identity and a mirror plane alone.

\subsection*{Brief illustration 10A. 3}
\begin{itemize}
  \item The CBrClFI molecule (3) has only the identity element, and so belongs to the group $C_{1}$.
  \item Meso-tartaric acid (6) has the identity and inversion elements, and so belongs to the group $C_{\mathrm{i}}$.\\
\includegraphics[max width=\textwidth, center]{2025_02_27_d4b95d9718f0b9e2e6b9g-424(1)}
\end{itemize}

6 Meso-tartaric acid,\\
$\mathrm{HOOCCH}(\mathrm{OH}) \mathrm{CH}(\mathrm{OH}) \mathrm{COOH}$

\begin{itemize}
  \item Quinoline (7) has the elements $(E, \sigma)$, and so belongs to the group $C_{s}$.\\
\includegraphics[max width=\textwidth, center]{2025_02_27_d4b95d9718f0b9e2e6b9g-424(3)}
\end{itemize}

\subsection*{(b) The groups $C_{n^{\prime}} C_{n v}$, and $C_{n h}$}
A molecule belongs to the group $C_{n}$ if it possesses an $n$-fold axis. Note that symbol $C_{n}$ is now playing a triple role: as the label of a symmetry element, a symmetry operation, and a group. If in addition to the identity and a $C_{n}$ axis a molecule has $n$ vertical mirror planes $\sigma_{v}$, then it belongs to the group $C_{n v}$. Molecules that in addition to the identity and an $n$-fold principal axis also have a horizontal mirror plane $\sigma_{\mathrm{h}}$ belong to the groups $C_{n h}$. The presence of certain symmetry elements may be implied by the presence of others: thus, in $C_{2 \mathrm{~h}}$ the elements $C_{2}$ and $\sigma_{\mathrm{h}}$ jointly imply the presence

\begin{center}
\begin{tabular}{ll}
\hline
Name & Elements \\
\hline
$C_{n}$ & $E, C_{n}$ \\
$C_{n v}$ & $E, C_{n} n \sigma_{v}$ \\
$C_{n \mathrm{~h}}$ & $E, C_{n}, \sigma_{\mathrm{h}}$ \\
\hline
\end{tabular}
\end{center} of a centre of inversion (Fig. 10A.9). Note also that the tables specify the elements, not the operations: for instance, there are two operations associated with a single $C_{3}$ axis (rotations by $+120^{\circ}$ and $-120^{\circ}$ ).

\begin{center}
\includegraphics[max width=\textwidth]{2025_02_27_d4b95d9718f0b9e2e6b9g-424(4)}
\end{center}

Figure 10A. 9 The presence of a twofold axis and a horizontal mirror plane jointly imply the presence of a centre of inversion in the molecule.

\subsection*{Brief illustration 10A. 4}
\begin{itemize}
  \item In the $\mathrm{H}_{2} \mathrm{O}_{2}$ molecule (8) the two $\mathrm{O}-\mathrm{H}$ bonds make an angle of about $115^{\circ}$ to one another when viewed down the $\mathrm{O}-\mathrm{O}$ bond direction. The molecule has the symmetry elements $E$ and $C_{2}$, and so belongs to the group $C_{2}$.\\
\includegraphics[max width=\textwidth, center]{2025_02_27_d4b95d9718f0b9e2e6b9g-424(2)}
\end{itemize}

$$
8 \text { Hydrogen peroxide, } \mathrm{H}_{2} \mathrm{O}_{2}
$$

\begin{itemize}
  \item An $\mathrm{H}_{2} \mathrm{O}$ molecule has the symmetry elements $E, C_{2}$, and $2 \sigma_{v}$, so it belongs to the group $C_{2 v}$.
  \item An $\mathrm{NH}_{3}$ molecule has the elements $E, C_{3}$, and $3 \sigma_{v}$, so it belongs to the group $C_{3 v}$.
  \item A heteronuclear diatomic molecule such as HCl belongs to the group $C_{\infty v}$ because rotations around the internuclear axis by any angle and reflections in any of the infinite number of planes that contain this axis are symmetry operations. Other members of the group $C_{\infty}$ include the linear OCS molecule and a cone.
  \item The molecule trans- $\mathrm{CHCl}=\mathrm{CHCl}$ (9) has the elements $E$, $C_{2}$, and $\sigma_{\mathrm{h}}$, so belongs to the group $C_{2 \mathrm{~h}}$.\\
\includegraphics[max width=\textwidth, center]{2025_02_27_d4b95d9718f0b9e2e6b9g-424}
  \item The molecule $\mathrm{B}(\mathrm{OH})_{3}$, in the planar conformation shown in (10), has a $C_{3}$ axis and a $\sigma_{\mathrm{h}}$ plane, and so belongs to the point group $C_{3 h}$.
\end{itemize}

\subsection*{(c) The groups $D_{n^{\prime}} D_{n h^{\prime}}$ and $D_{n d}$}
Figure 10A. 7 shows that a molecule that has an $n$-fold principal axis and $n$ twofold axes perpendicular to $C_{n}$ belongs to the group $D_{n}$. A molecule belongs to $D_{n \mathrm{~h}}$ if it also possesses a horizontal mirror plane. The linear molecules OCO and HCCH,

\begin{center}
\begin{tabular}{ll}
\hline
Name & Elements \\
\hline
$D_{n}$ & $E, C_{n}, n C_{2}^{\prime}$ \\
$D_{n \mathrm{~h}}$ & $E, C_{n}, n C_{2}^{\prime}, \sigma_{\mathrm{h}}$ \\
$D_{n \mathrm{~d}}$ & $E, C_{n}, n C_{2}^{\prime}, n \sigma_{\mathrm{d}}$ \\
\hline
\end{tabular}
\end{center} and a uniform cylinder, all belong to the group $D_{\text {oh }}$. A molecule belongs to the group $D_{n \mathrm{~d}}$ if in addition to the elements of $D_{n}$ it possesses $n$ dihedral mirror planes $\sigma_{\mathrm{d}}$.

\subsection*{Brief illustration 10A. 5}
\begin{itemize}
  \item The planar trigonal $\mathrm{BF}_{3}$ molecule (11) has the elements $E$, $C_{3}, 3 C_{2}$ (with one $C_{2}$ axis along each $\mathrm{B}-\mathrm{F}$ bond), and $\sigma_{\mathrm{h}}$, so belongs to $D_{3 h}$.\\
\includegraphics[max width=\textwidth, center]{2025_02_27_d4b95d9718f0b9e2e6b9g-425(1)}
  \item The $\mathrm{C}_{6} \mathrm{H}_{6}$ molecule has the elements $E, C_{6}, 3 C_{2}, 3 C_{2}^{\prime}$, and $\sigma_{\mathrm{h}}$ together with some others that these elements imply, so it belongs to $D_{66}$. Three of the $C_{2}$ axes bisect C-C bonds on opposite sides of the hexagonal ring formed by the carbon atoms, and the other three pass through vertices on opposite sides of the ring. The prime on $3 C_{2}^{\prime}$ indicates that these axes are different from the other three $C_{2}$ axes.
  \item All homonuclear diatomic molecules, such as $\mathrm{N}_{2}$, belong to the group $D_{\text {oh }}$ because all rotations around the internuclear axis are symmetry operations, as are end-overend rotations by $180^{\circ}$.
  \item $\mathrm{PCl}_{5}(12)$ is another example of a $D_{3 \mathrm{~h}}$ species.\\
\includegraphics[max width=\textwidth, center]{2025_02_27_d4b95d9718f0b9e2e6b9g-425(2)}
\end{itemize}

12 Phosphorus pentachloride, $\mathrm{PCl}_{5}\left(D_{3 \mathrm{~h}}\right)$

\begin{itemize}
  \item Propadiene (an allene, 13), in which the two $\mathrm{CH}_{2}$ groups lie in perpendicular planes, belongs to the point group $D_{2 d}$.\\
\includegraphics[max width=\textwidth, center]{2025_02_27_d4b95d9718f0b9e2e6b9g-425(3)}
\end{itemize}

\subsection*{(d) The groups $S_{n}$}
Molecules that have not been classified into one of the groups mentioned so far, but which possess one $S_{n}$ axis, belong to the groups $S_{n}$. Note that the group $S_{2}$ is the same as $C_{\mathrm{i}}$, so such a molecule will already have been classified as $C_{\mathrm{i}}$. Tetraphenylmethane (14) belongs to the point group $S_{4}$; molecules belonging to $S_{n}$ with $n>4$ are rare.

\begin{center}
\begin{tabular}{ll}
\hline
Name & Elements \\
\hline
$S_{n}$ & $E, S_{n}$ and not previously \\
classified &  \\
\end{tabular}
\end{center}

\begin{center}
\includegraphics[max width=\textwidth]{2025_02_27_d4b95d9718f0b9e2e6b9g-425}
\end{center}

14 Tetraphenylmethane, $\left.\mathrm{C}_{\left(\mathrm{C}_{6}\right.} \mathrm{H}_{5}\right)_{4}\left(\mathrm{~S}_{4}\right)$

\subsection*{(e) The cubic groups}
A number of very important molecules possess more than one principal axis. Most belong to the cubic groups, and in particular to the tetrahedral groups $T, T_{\mathrm{d}}$, and $T_{\mathrm{h}}$ (Fig. 10A.10a) or to the octahedral groups $O$ and $O_{\mathrm{h}}$ (Fig. 10A.10b). A few icosahedral ( 20 -faced) molecules belonging to the icosahedral group, $I$ (Fig. 10A.10c), are also known. The groups $T_{\mathrm{d}}$ and $O_{\mathrm{h}}$ are the groups of the regular tetrahedron and the regular octahedron, respectively. If the object possesses the rotational symmetry of

\begin{center}
\begin{tabular}{ll}
\hline
Name & Elements \\
\hline
$T$ & $E, 4 C_{3}, 3 C_{2}$ \\
$T_{\mathrm{d}}$ & $E, 3 C_{2}, 4 C_{3}, 3 S_{4}, 6 \sigma_{\mathrm{d}}$ \\
$T_{\mathrm{h}}$ & $E, 3 C_{2}, 4 C_{3}, i, 4 S_{6}, 3 \sigma_{\mathrm{h}}$ \\
$O$ & $E, 3 C_{4}, 4 C_{3}, 6 C_{2}$ \\
$O_{\mathrm{h}}$ & $E, 3 S_{4}, 3 C_{4}, 6 C_{2}, 4 S_{6}, 4 C_{3}, 3 \sigma_{\mathrm{h}}, 6 \sigma_{\mathrm{d}}, i$ \\
$I$ & $E, 6 C_{5}, 10 C_{3}, 15 C_{2}$ \\
$I_{\mathrm{h}}$ & $E, 6 S_{10}, 10 S_{6}, 6 C_{5}, 10 C_{3}, 15 C_{2}, 15 \sigma, i$ \\
\hline
\end{tabular}
\end{center}

\begin{center}
\includegraphics[max width=\textwidth]{2025_02_27_d4b95d9718f0b9e2e6b9g-426(2)}
\end{center}

Figure 10A. 10 (a) Tetrahedral, (b) octahedral, and (c) icosahedral shapes drawn to show their relation to a cube: they belong to the cubic groups $T_{\mathrm{d}}, O_{\mathrm{h}}$, and $I_{\mathrm{h}}$, respectively.\\
\includegraphics[max width=\textwidth, center]{2025_02_27_d4b95d9718f0b9e2e6b9g-426(3)}

Figure 10A. 11 Shapes corresponding to the point groups (a) $T$ and (b) $O$. The presence of the decorated slabs reduces the symmetry of the object from $T_{d}$ and $O_{h}$, respectively.\\
the tetrahedron or the octahedron, but none of their planes of reflection, then it belongs to the simpler groups $T$ or $O$ (Fig. 10A.11). The group $T_{\mathrm{h}}$ is based on $T$ but also contains a centre of inversion (Fig. 10A.12).

\subsection*{Brief illustration 10A. 6}
\begin{itemize}
  \item The molecules $\mathrm{CH}_{4}$ and $\mathrm{SF}_{6}$ belong, respectively, to the groups $T_{\mathrm{d}}$ and $O_{\mathrm{h}}$.
  \item Molecules belonging to the icosahedral group $I$ include some of the boranes and buckminsterfullerene, $\mathrm{C}_{60}$ (15).\\
\includegraphics[max width=\textwidth, center]{2025_02_27_d4b95d9718f0b9e2e6b9g-426}
\end{itemize}

15 Buckminsterfullerene, $\mathrm{C}_{60}$ (I)

\begin{itemize}
  \item The objects shown in Fig. 10A. 11 belong to the groups $T$ and $O$, respectively.\\
\includegraphics[max width=\textwidth, center]{2025_02_27_d4b95d9718f0b9e2e6b9g-426(1)}
\end{itemize}

Figure 10A. 12 The shape of an object belonging to the group $T_{h}$.

\subsection*{(f) The full rotation group}
\begin{center}
\begin{tabular}{ll}
\hline
Name & Elements \\
\hline
$R_{3}$ & $E, \infty C_{2}, \infty C_{3}, \ldots$ \\
\hline
\end{tabular}
\end{center}

The full rotation group, $R_{3}$ (the 3 refers to rotation in three dimensions), consists of an infinite number of rotation axes with all possible values of $n$. A sphere and an atom belong to $R_{3}$, but no molecule does. Exploring the consequences of $R_{3}$ is a very important way of applying symmetry arguments to atoms, and is an alternative approach to the theory of orbital angular momentum.

\subsection*{10A. 3 Some immediate consequences of symmetry}
Some statements about the properties of a molecule can be made as soon as its point group has been identified.

\subsection*{(a) Polarity}
A polar molecule is one with a permanent electric dipole moment $\left(\mathrm{HCl}, \mathrm{O}_{3}\right.$, and $\mathrm{NH}_{3}$ are examples). A dipole moment is a property of the molecule, so it follows that the dipole moment (which is represented by a vector) must be unaffected by any symmetry operation of the molecule because, by definition, such an operation leaves the molecule apparently unchanged. If a molecule possesses a $C_{n}$ axis $(n>1)$ then it is not possible for there to be a dipole moment perpendicular to this axis because such a dipole moment would change its orientation on rotation about the axis. It is however possible for there to be a dipole parallel to the axis, because it would not be affected by the rotation. For example, in $\mathrm{H}_{2} \mathrm{O}$ the dipole lies in the plane of the molecule, pointing along the bisector of the HOH bond, which is the direction of the $C_{2}$ axis. Similarly, if a molecule possesses a mirror plane there can be no dipole moment perpendicular to this plane, because reflection in the plane would reverse its direction. A molecule that possesses a centre of symmetry cannot have a dipole moment in any direction because the inversion operation would reverse it.

These considerations lead to the conclusion that\\
Only molecules belonging to the groups $C_{n}, C_{n v}$, and $C_{s}$ may have a permanent electric dipole moment.

For $C_{n}$ and $C_{n v}$, the dipole moment must lie along the principal axis.

\subsection*{Brief illustration 10A. 7}
\begin{itemize}
  \item Ozone, $\mathrm{O}_{3}$, which has an angular structure, belongs to the group $C_{2 v}$ and is polar.
  \item Carbon dioxide, $\mathrm{CO}_{2}$, which is linear and belongs to the group $D_{\text {oh }}$, is not polar.
  \item Tetraphenylmethane (14) belongs to the point group $S_{4}$ and so is not polar.
\end{itemize}

\subsection*{(b) Chirality}
A chiral molecule (from the Greek word for 'hand') is a molecule that cannot be superimposed on its mirror image. An achiral molecule is a molecule that can be superimposed on its mirror image. Chiral molecules are optically active in the sense that they rotate the plane of polarized light. A chiral molecule and its mirror-image partner constitute an enantiomeric pair (from the Greek word for 'both') of isomers and rotate the plane of polarization by equal amounts but opposite directions.

\begin{displayquote}
A molecule may be chiral, and therefore optically active, only if it does not possess an axis of improper rotation, $S_{n}$.
\end{displayquote}

An $S_{n}$ improper rotation axis may be present under a different name, and be implied by other symmetry elements that are present. For example, molecules belonging to the groups $C_{n \mathrm{~h}}$ possess an $S_{n}$ axis implicitly because they possess both $C_{n}$ and $\sigma_{h}$, which are the two components of an improper rotation axis. A centre of inversion, $i$, is in fact the same as $S_{2}$ because the two corresponding operations achieve exactly the same result (Fig. 10A.13). Furthermore, a mirror plane is the same as $S_{1}$ (rotation through $360^{\circ}$ followed by reflection). Thus molecules possessing a mirror plane or a centre of inversion effectively possess an axis of improper rotation and so, by the above rule, are achiral.\\
\includegraphics[max width=\textwidth, center]{2025_02_27_d4b95d9718f0b9e2e6b9g-427}

Figure 10A. 13 The operations $i$ and $S_{2}$ are equivalent in the sense that they achieve exactly the same outcome when applied to a point in the object.

\subsection*{Brief illustration 10A. 8}
\begin{itemize}
  \item The amino acid alanine (16) does not possess a centre of inversion nor does it have any mirror planes: it is therefore chiral.\\
\includegraphics[max width=\textwidth, center]{2025_02_27_d4b95d9718f0b9e2e6b9g-427(1)}
\end{itemize}

16 L-Alanine, $\mathrm{NH}_{2} \mathrm{CH}\left(\mathrm{CH}_{3}\right) \mathrm{COOH}$

\begin{itemize}
  \item In contrast, glycine (17) has a mirror plane and so is achiral.\\
\includegraphics[max width=\textwidth, center]{2025_02_27_d4b95d9718f0b9e2e6b9g-427(2)}
  \item Tetraphenylmethane (14) belongs to the point group $S_{4} ;$ it does not possess a centre of inversion or any mirror planes, but it is still achiral since it possesses an axis of improper rotation $\left(S_{4}\right)$.
\end{itemize}

\subsection*{Checklist of concepts}
\begin{enumerate}
  \item A symmetry operation is an action that leaves an object looking the same after it has been carried out.\\
$\square$ 2. A symmetry element is a point, line, or plane with respect to which a symmetry operation is performed.\\
$\square$ 3. The notation for point groups commonly used for molecules and solids is summarized in Table 10A.1.\\
$\square$ 4. To be polar, a molecule must belong to $C_{n}, C_{n v}$, or $C_{s}$ (and have no higher symmetry).\\
$\square$ 5. A molecule will be chiral only if it does not possess an axis of improper rotation, $S_{n}$.
\end{enumerate}

\subsection*{Checklist of operations and elements}
\begin{center}
\begin{tabular}{lll}
\hline
Symmetry operation & Symbol & Symmetry element \\
\hline
$n$-Fold rotation & $C_{n}$ & $n$-Fold axis of rotation \\
Reflection & $\sigma$ & Mirror plane \\
Inversion & $i$ & Centre of symmetry \\
$n$-Fold improper rotation & $S_{n}$ & $n$-Fold improper axis of rotation \\
Identity & $E$ & Entire object \\
\hline
\end{tabular}
\end{center}

\section*{TOPIC 10B Group theory}
\subsection*{- Why do you need to know this material?}
Group theory expresses qualitative ideas about symmetry mathematically and can be applied systematically to a wide variety of problems. The theory is also the origin of the labelling of atomic and molecular orbitals that is used throughout chemistry.

\subsection*{- What is the key idea?}
Symmetry operations may be represented by the effect of matrices on a basis.

\subsection*{> What do you need to know already?}
You need to know about the types of symmetry operation and element introduced in Topic 10A. This discussion uses matrix algebra, and especially matrix multiplication, as set out in The chemist's toolkit 24 in Topic 9E.

The systematic discussion of symmetry is called group theory. Much of group theory is a summary of common sense about the symmetries of objects. However, because group theory is systematic, its rules can be applied in a straightforward, mechanical way. In most cases the theory gives a simple, direct method for arriving at useful conclusions with the minimum of calculation, and this is the aspect that is stressed here.

\subsection*{10B. 1 The elements of group theory}
A group in mathematics is a collection of transformations that satisfy four criteria. If the transformations are written as $R, R^{\prime}$, ... (which might be the reflections, rotations, and so on introduced in Topic 10A), then they form a group if:

\begin{enumerate}
  \item One of the transformations is the identity (i.e. 'do nothing').
  \item For every transformation $R$, the inverse transformation $R^{-1}$ is included in the collection so that the combination $R R^{-1}$ (the transformation $R^{-1}$ followed by $R$ ) is equivalent to the identity.
  \item The combination $R R^{\prime}$ (the transformation $R^{\prime}$ followed by $R$ ) is equivalent to a single member of the collection of transformations.
  \item The combination $R\left(R^{\prime} R^{\prime \prime}\right)$, the transformation ( $R^{\prime} R^{\prime \prime}$ ) followed by $R$, is equivalent to $\left(R R^{\prime}\right) R^{\prime \prime}$, the transformation $R^{\prime \prime}$ followed by ( $R R^{\prime}$ ).
\end{enumerate}

\subsection*{Example 10B. 1 \\
 Showing that the symmetry operations of}
 a molecule form a groupThe point group $C_{2 v}$ consists of the elements $\left\{E, C_{2}, \sigma_{v}, \sigma_{v}^{\prime}\right\}$ and correspond to the operations $\left\{E, C_{2}, \sigma_{v}, \sigma_{v}^{\prime}\right\}$. Show that this set of operations is a group in the mathematical sense.

Collect your thoughts You need to show that combinations of the operations match the criteria set out above. The operations are specified in Topic 10A, and illustrated in Figs. 10A. 2 and 10A. 3 for $\mathrm{H}_{2} \mathrm{O}$, which belongs to this group.

\subsection*{The solution}
\begin{itemize}
  \item Criterion 1 is fulfilled because the collection of symmetry operations includes the identity $E$.
  \item Criterion 2 is fulfilled because in each case the inverse of an operation is the operation itself. Thus, two successive twofold rotations is equivalent to the identity: $C_{2} C_{2}=E$ and likewise for the two reflections and the identity itself.
  \item Criterion 3 is fulfilled, because in each case one operation followed by another is the same as one of the four symmetry operations. For instance, a twofold rotation $C_{2}$ followed by the reflection $\sigma_{v}$ is the same as the single reflection $\sigma_{\mathrm{v}}^{\prime}$ (Fig. 10B.1); thus, $\sigma_{\mathrm{v}} C_{2}=\sigma_{\mathrm{v}}^{\prime}$. A 'group multiplication table' can be constructed in a similar way for all possible products of symmetry operations $R R^{\prime}$; as required, each product is equivalent to another symmetry operation.
\end{itemize}

\begin{center}
\begin{tabular}{lllll}
\hline
$R \downarrow R^{\prime} \rightarrow$ & $E$ & $C_{2}$ & $\sigma_{\mathrm{v}}$ & $\sigma_{\mathrm{v}}^{\prime}$ \\
\hline
$E$ & $E$ & $C_{2}$ & $\sigma_{\mathrm{v}}$ & $\sigma_{\mathrm{v}}^{\prime}$ \\
$C_{2}$ & $C_{2}$ & $E$ & $\sigma_{\mathrm{v}}^{\prime}$ & $\sigma_{\mathrm{v}}$ \\
$\sigma_{\mathrm{v}}$ & $\sigma_{\mathrm{v}}$ & $\sigma_{\mathrm{v}}^{\prime}$ & $E$ & $C_{2}$ \\
$\sigma_{\mathrm{v}}^{\prime}$ & $\sigma_{\mathrm{v}}^{\prime}$ & $\sigma_{\mathrm{v}}$ & $C_{2}$ & $E$ \\
\hline
\end{tabular}
\end{center}

\begin{itemize}
  \item Criterion 4 is fulfilled, as it is immaterial how the operations are grouped together. Thus $\left(\sigma_{\mathrm{v}} \sigma_{\mathrm{v}}^{\prime}\right) C_{2}=C_{2} C_{2}=E$ and $\sigma_{\mathrm{v}}\left(\sigma_{\mathrm{v}}^{\prime} \mathrm{C}_{2}\right)=\sigma_{\mathrm{v}} \sigma_{\mathrm{v}}=E$, and likewise for all other combinations.
\end{itemize}

Self-test 10B. 1 Confirm that $C_{2 \mathrm{~h}}$, which has the elements $\left\{E, C_{2}, i, \sigma_{\mathrm{h}}\right\}$ and hence the corresponding operations $\left\{E, C_{2}, i, \sigma_{\mathrm{h}}\right\}$, is a group (construct the group multiplication table).\\
\includegraphics[max width=\textwidth, center]{2025_02_27_d4b95d9718f0b9e2e6b9g-429}\\
\includegraphics[max width=\textwidth, center]{2025_02_27_d4b95d9718f0b9e2e6b9g-430}

Figure 10B. 1 A twofold rotation $C_{2}$ followed by the reflection $\sigma_{v}$ gives the same result as the reflection $\sigma_{v}^{\prime}$.

One potentially confusing point needs to be clarified at the outset. The entities that make up a group are its 'elements'. For applications in chemistry, these elements are almost always symmetry operations. However, as explained in Topic 10A, 'symmetry operations' are distinct from 'symmetry elements', the latter being the points, axes, and planes with respect to which the operations are carried out. A third use of the word 'element' is to denote the number lying in a particular location in a matrix. Be very careful to distinguish element (of a group), symmetry element, and matrix element.

Symmetry operations fall into the same class if they are of the same type (for example, rotations) and can be transformed into one another by a symmetry operation of the group. The two threefold rotations in $C_{3 \mathrm{v}}$ belong to the same class because one can be converted into the other by a reflection (Fig. 10B.2); the three reflections all belong to the same class because each can be rotated into another by a threefold rotation. The formal definition of a class is that two operations $R$ and $R^{\prime}$ belong to the same class if there is a member $S$ of the group such that

$$
R^{\prime}=S^{-1} R S
$$

Membership of a class\\
(10B.1)\\
where $S^{-1}$ is the inverse of $S$.\\
\includegraphics[max width=\textwidth, center]{2025_02_27_d4b95d9718f0b9e2e6b9g-430(1)}

Figure 10B. 2 Symmetry operations in the same class are related to one another by the symmetry operations of the group. Thus, the three mirror planes shown here are related by threefold rotations, and the two rotations shown here are related by reflection in $\sigma_{\mathrm{v}}$.\\
\includegraphics[max width=\textwidth, center]{2025_02_27_d4b95d9718f0b9e2e6b9g-430(2)}

Figure 10B.3 (a) The sequence of operations $\sigma_{v}^{-1} C_{3}^{+} \sigma_{v}$ when applied to the point 1 takes it through the sequence $1 \rightarrow 2 \rightarrow$ $3 \rightarrow 4$ ( $\sigma_{v}^{-1}$ has the same effect as $\sigma_{v}$ ). The single operation $C_{3}^{-}$ (dotted curve) takes point $1 \rightarrow 4$, so $C_{3}^{+}$and $C_{3}^{-}$are in the same class. (b) The sequence of operations $\left(C_{3}^{+}\right)^{-1} \sigma_{\mathrm{v}} C_{3}^{+}$takes point $1 \rightarrow 4$, $\left(\left(C_{3}^{+}\right)^{-1}\right.$ has the same effect as $\left.C_{3}^{-}\right)$but the same transformation can be achieved with the single operation $\sigma_{v}^{\prime}$ (dotted line); so $\sigma_{v}$ and $\sigma_{v}^{\prime}$ are in the same class.

Figure 10B.3(a) shows how eqn 10B. 1 can be used to confirm that $C_{3}^{+}$and $C_{3}^{-}$belong to the same class in the group $C_{3 v}$ by considering how an arbitrary point, 1 , behaves under the various operations. The transformation of interest is $\sigma_{\mathrm{v}}^{-1} C_{3}^{+} \sigma_{\mathrm{v}}$. Start at 1: the operation $\sigma_{v}$ moves the point to 2 , and then $C_{3}^{+}$ moves the point to 3 . The inverse of a reflection is itself, $\sigma_{\mathrm{v}}^{-1}=$ $\sigma_{v}$, so the effect of $\sigma_{v}^{-1}$ is to move the point to 4 . From the diagram it can be seen that point 4 can be reached by applying $C_{3}^{-}$ to point 1 , thus demonstrating that $\sigma_{\mathrm{v}}^{-1} C_{3}^{+} \sigma_{\mathrm{v}}=C_{3}^{-}$, and hence that $C_{3}^{+}$and $C_{3}^{-}$do indeed belong in the same class.

\subsection*{Brief illustration 10B. 1}
To show that $\sigma_{\mathrm{v}}$ and $\sigma_{\mathrm{v}}^{\prime}$ are in the same class in the group $C_{3 \mathrm{v}}$, consider the transformation $\left(C_{3}^{+}\right)^{-1} \sigma_{\mathrm{v}} C_{3}^{+}$. Because $C_{3}^{-}$is the inverse of $C_{3}^{+}$, this transformation is the same as $C_{3}^{-} \sigma_{\mathrm{v}} C_{3}^{+}$; the effect of this sequence of operations on an arbitrary point 1 is shown in Fig. 10B.3(b). The final position, 4, can also be reached from 1 by applying the operation $\sigma_{v}^{\prime}$, thus showing that $C_{3}^{-} \sigma_{\mathrm{v}} C_{3}^{+}=\sigma_{\mathrm{v}}^{\prime}$ and hence that $\sigma_{\mathrm{v}}$ and $\sigma_{\mathrm{v}}^{\prime}$ are in the same class.

\subsection*{10B.2 Matrix representations}
Group theory takes on great power when the notional ideas presented so far are expressed in terms of collections of numbers in the form of matrices. For basic information about how to handle matrices, see The chemist's toolkit 24 in Topic 9E.

\subsection*{(a) Representatives of operations}
Consider the set of five p orbitals shown on the $\mathrm{C}_{2 \mathrm{v}} \mathrm{SO}_{2}$ molecule in Fig. 10B. 4 and how they are affected by the reflection operation $\sigma_{\mathrm{v}}$. The corresponding symmetry element is the\\
mirror plane perpendicular to the plane of the molecule and passing through the $S$ atom. The effect of this reflection is to leave $\mathrm{p}_{x}$ and $\mathrm{p}_{z}$ unaffected, to change the sign of $\mathrm{p}_{y}$, and to exchange $\mathrm{p}_{\mathrm{A}}$ and $\mathrm{p}_{\mathrm{B}}$. Its effect can be written $\left(\mathrm{p}_{x}-\mathrm{p}_{y} \mathrm{p}_{z} \mathrm{p}_{\mathrm{B}} \mathrm{p}_{\mathrm{A}}\right) \leftarrow$ $\left(\mathrm{p}_{x} \mathrm{p}_{y} \mathrm{p}_{z} \mathrm{p}_{\mathrm{A}} \mathrm{p}_{\mathrm{B}}\right)$. This transformation can be expressed by using matrix multiplication:


\begin{align*}
\left(\mathrm{p}_{x}-\mathrm{p}_{y} \mathrm{p}_{z} \mathrm{p}_{\mathrm{B}} \mathrm{p}_{\mathrm{A}}\right) & =\left(\mathrm{p}_{x} \mathrm{p}_{y} \mathrm{p}_{z} \mathrm{p}_{\mathrm{A}} \mathrm{p}_{\mathrm{B}}\right) \overbrace{\left(\begin{array}{ccccc}
1 & 0 & 0 & 0 & 0 \\
0 & -1 & 0 & 0 & 0 \\
0 & 0 & 1 & 0 & 0 \\
0 & 0 & 0 & 0 & 1 \\
0 & 0 & 0 & 1 & 0
\end{array}\right)}^{D\left(\sigma_{v}\right)} \\
& =\left(\mathrm{p}_{x} \mathrm{p}_{y} \mathrm{p}_{z} \mathrm{p}_{\mathrm{A}} \mathrm{p}_{\mathrm{B}}\right) \boldsymbol{D}\left(\sigma_{\mathrm{v}}\right) \tag{10B.2a}
\end{align*}


The matrix $\boldsymbol{D}\left(\sigma_{\mathrm{v}}\right)$ is called a representative of the operation $\sigma_{\mathrm{v}}$. Representatives take different forms according to the basis, the set of orbitals that has been adopted. In this case, the basis is the row vector $\left(\mathrm{p}_{x} \mathrm{p}_{y} \mathrm{p}_{z} \mathrm{p}_{\mathrm{A}} \mathrm{p}_{\mathrm{B}}\right)$. Note that the matrix $\boldsymbol{D}$ appears to the right of the basis functions on which it acts.

The same technique can be used to find matrices that reproduce the other symmetry operations. For instance, $C_{2}$ has the effect $\left(-\mathrm{p}_{x}-\mathrm{p}_{y} \mathrm{p}_{z}-\mathrm{p}_{\mathrm{B}}-\mathrm{p}_{\mathrm{A}}\right) \leftarrow\left(\mathrm{p}_{x} \mathrm{p}_{y} \mathrm{p}_{z} \mathrm{p}_{\mathrm{A}} \mathrm{p}_{\mathrm{B}}\right)$, and its representative is

\[
\boldsymbol{D}\left(C_{2}\right)=\left(\begin{array}{ccccc}
-1 & 0 & 0 & 0 & 0  \tag{10B.2b}\\
0 & -1 & 0 & 0 & 0 \\
0 & 0 & 1 & 0 & 0 \\
0 & 0 & 0 & 0 & -1 \\
0 & 0 & 0 & -1 & 0
\end{array}\right)
\]

The effect of $\sigma_{\mathrm{v}}^{\prime}$ (reflection in the plane of the molecule) is $\left(-\mathrm{p}_{x} \mathrm{p}_{y} \mathrm{p}_{z}-\mathrm{p}_{\mathrm{A}}-\mathrm{p}_{\mathrm{B}}\right) \leftarrow\left(\mathrm{p}_{x} \mathrm{p}_{y} \mathrm{p}_{z} \mathrm{p}_{\mathrm{A}} \mathrm{p}_{\mathrm{B}}\right)$; the oxygen orbitals remain in the same places, but change sign. The representative of this operation is

\[
\boldsymbol{D}\left(\sigma_{2}^{\prime}\right)=\left(\begin{array}{ccccc}
-1 & 0 & 0 & 0 & 0  \tag{10B.2c}\\
0 & 1 & 0 & 0 & 0 \\
0 & 0 & 1 & 0 & 0 \\
0 & 0 & 0 & -1 & 0 \\
0 & 0 & 0 & 0 & -1
\end{array}\right)
\]

The identity operation has no effect on the basis, so its representative is the $5 \times 5$ unit matrix:

\[
\boldsymbol{D}(E)=\left(\begin{array}{lllll}
1 & 0 & 0 & 0 & 0  \tag{10B.2d}\\
0 & 1 & 0 & 0 & 0 \\
0 & 0 & 1 & 0 & 0 \\
0 & 0 & 0 & 1 & 0 \\
0 & 0 & 0 & 0 & 1
\end{array}\right)
\]

\begin{center}
\includegraphics[max width=\textwidth]{2025_02_27_d4b95d9718f0b9e2e6b9g-431}
\end{center}

Figure 10B. 4 The five p orbitals (three on the sulfur and one on each oxygen) that are used to illustrate the construction of a matrix representation in a $\mathrm{C}_{2 \mathrm{v}}$ molecule $\left(\mathrm{SO}_{2}\right)$.

\subsection*{(b) The representation of a group}
The set of matrices that represents all the operations of the group is called a matrix representation, $\Gamma$ (uppercase gamma), of the group in the basis that has been chosen. In the current example, there are five members of the basis and the representation is five-dimensional in the sense that the matrices are all $5 \times 5$ arrays. The matrices of a representation multiply together in the same way as the operations they represent. Thus, if for any two operations $R$ and $R^{\prime}, R R^{\prime}=R^{\prime \prime}$, then $\boldsymbol{D}(R) \boldsymbol{D}\left(R^{\prime}\right)=$ $D\left(R^{\prime \prime}\right)$ for a given basis.

\subsection*{Brief illustration 10B. 2}
In the group $C_{2 v}$, a twofold rotation followed by a reflection in a mirror plane is equivalent to a reflection in the second mirror plane: specifically, $\sigma_{\mathrm{v}}^{\prime} \mathrm{C}_{2}=\sigma_{\mathrm{v}}$. Multiplying out the representatives specified in eqn 10B. 2 gives

$$
\begin{aligned}
\boldsymbol{D}\left(\sigma_{v}^{\prime}\right) \boldsymbol{D}\left(C_{2}\right) & =\left(\begin{array}{ccccc}
-1 & 0 & 0 & 0 & 0 \\
0 & 1 & 0 & 0 & 0 \\
0 & 0 & 1 & 0 & 0 \\
0 & 0 & 0 & -1 & 0 \\
0 & 0 & 0 & 0 & -1
\end{array}\right)\left(\begin{array}{ccccc}
-1 & 0 & 0 & 0 & 0 \\
0 & -1 & 0 & 0 & 0 \\
0 & 0 & 1 & 0 & 0 \\
0 & 0 & 0 & 0 & -1 \\
0 & 0 & 0 & -1 & 0
\end{array}\right) \\
& =\left(\begin{array}{ccccc}
1 & 0 & 0 & 0 & 0 \\
0 & -1 & 0 & 0 & 0 \\
0 & 0 & 1 & 0 & 0 \\
0 & 0 & 0 & 0 & 1 \\
0 & 0 & 0 & 1 & 0
\end{array}\right)=\boldsymbol{D}\left(\sigma_{v}\right)
\end{aligned}
$$

As expected, this multiplication reproduces the same result as the group multiplication table. The same is true for multiplication of any two representatives, so the four matrices form a representation of the group.

The discovery of a matrix representation of the group means that a link has been established between symbolic manipulations of operations and algebraic manipulations of numbers. This link is the basis of the power of group theory in chemistry.

\subsection*{(c) Irreducible representations}
Inspection of the representatives found above shows that they are all of block-diagonal form:

\[
\boldsymbol{D}=\left(\begin{array}{ccccc}
\boldsymbol{\square} & 0 & 0 & 0 & 0  \tag{10B.3}\\
0 & \square & 0 & 0 & 0 \\
0 & 0 & \boldsymbol{\square} & 0 & 0 \\
0 & 0 & 0 & \boldsymbol{\square} & \boldsymbol{\square} \\
0 & 0 & 0 & \boldsymbol{\square} & \boldsymbol{\square}
\end{array}\right) \quad \text { Block-diagonal form }
\]

The block-diagonal form of the representatives implies that the symmetry operations of $C_{2 \mathrm{v}}$ never mix $\mathrm{p}_{x}, \mathrm{p}_{y}$, and $\mathrm{p}_{z}$ together, nor do they mix these three orbitals with $p_{A}$ and $p_{B}$, but $p_{A}$ and $p_{B}$ are mixed together by the operations of the group. Consequently, the basis can be cut into four parts: three for the individual $p$ orbitals on $S$ and the fourth for the two oxygen orbitals $\left(p_{A}, p_{B}\right)$. The representations in these three onedimensional bases are

$$
\begin{array}{llll}
\text { For } \mathrm{p}_{x}: \boldsymbol{D}(E)=1 & \boldsymbol{D}\left(C_{2}\right)=-1 & \boldsymbol{D}\left(\sigma_{\mathrm{v}}\right)=1 & \boldsymbol{D}\left(\sigma_{\mathrm{v}}^{\prime}\right)=-1 \\
\text { For } \mathrm{p}_{y}: \boldsymbol{D}(E)=1 & \boldsymbol{D}\left(C_{2}\right)=-1 & \boldsymbol{D}\left(\sigma_{\mathrm{v}}\right)=-1 & \boldsymbol{D}\left(\sigma_{\mathrm{v}}^{\prime}\right)=1 \\
\text { For } \mathrm{p}_{z}: \boldsymbol{D}(E)=1 & \boldsymbol{D}\left(C_{2}\right)=1 & \boldsymbol{D}\left(\sigma_{\mathrm{v}}\right)=1 & \boldsymbol{D}\left(\sigma_{\mathrm{v}}^{\prime}\right)=1
\end{array}
$$

These representations will be called $\Gamma^{(1)}, \Gamma^{(2)}$, and $\Gamma^{(3)}$, respectively. The remaining two functions ( $p_{A} p_{B}$ ) are a basis for a two-dimensional representation denoted $\Gamma^{\prime}$ :

$$
\begin{array}{ll}
\boldsymbol{D}(E)=\left(\begin{array}{ll}
1 & 0 \\
0 & 1
\end{array}\right) & \boldsymbol{D}\left(C_{2}\right)=\left(\begin{array}{cc}
0 & -1 \\
-1 & 0
\end{array}\right) \\
\boldsymbol{D}\left(\sigma_{\mathrm{v}}\right)=\left(\begin{array}{ll}
0 & 1 \\
1 & 0
\end{array}\right) & \boldsymbol{D}\left(\sigma_{\mathrm{v}}^{\prime}\right)=\left(\begin{array}{cc}
-1 & 0 \\
0 & -1
\end{array}\right)
\end{array}
$$

The original five-dimensional representation has been reduced to the 'direct sum' of three one-dimensional representations 'spanned' by each of the $p$ orbitals on $S$, and a two-dimensional representation spanned by $\left(p_{A} p_{B}\right)$. The reduction is represented symbolically by writing ${ }^{1}$

$$
\Gamma=\Gamma^{(1)}+\Gamma^{(2)}+\Gamma^{(3)}+\Gamma^{\prime}
$$

Direct sum (10B.4)

\footnotetext{${ }^{1}$ The symbol $\oplus$ is sometimes used to denote a direct sum to distinguish it from an ordinary sum, in which case eqn 10B. 4 would be written $\Gamma=\Gamma^{(1)} \oplus$ $\Gamma^{(2)} \oplus \Gamma^{(3)} \oplus \Gamma^{\prime}$.
}The representations $\Gamma^{(1)}, \Gamma^{(2)}$, and $\Gamma^{(3)}$ cannot be reduced any further, and each one is called an irreducible representation of the group (an 'irrep'). That the two-dimensional representation $\Gamma^{\prime}$ is reducible (for this basis, in this group) is demonstrated by switching attention to the linear combinations $p_{1}=$ $p_{A}+p_{B}$ and $p_{2}=p_{A}-p_{B}$ in (Fig. 10B.5). The effect of the operation $\sigma_{\mathrm{v}}$ is to exchange $\mathrm{p}_{\mathrm{A}}$ and $\mathrm{p}_{\mathrm{B}}:\left(\mathrm{p}_{\mathrm{B}} \mathrm{p}_{\mathrm{A}}\right) \leftarrow\left(\mathrm{p}_{\mathrm{A}} \mathrm{p}_{\mathrm{B}}\right)$. Therefore $\left(p_{B}+p_{A}\right) \leftarrow\left(p_{A}+p_{B}\right)$, corresponding to $\left(p_{1}\right) \leftarrow\left(p_{1}\right)$. Similarly, $\left(p_{B}-p_{A}\right) \leftarrow\left(p_{A}-p_{B}\right)$, corresponding to $\left(-p_{2}\right) \leftarrow\left(p_{2}\right)$. It follows from these results, and similar ones for the other operations, that the representation in the basis $\left(p_{1} p_{2}\right)$ is

$$
\begin{array}{ll}
\boldsymbol{D}(E)=\left(\begin{array}{ll}
1 & 0 \\
0 & 1
\end{array}\right) & \boldsymbol{D}\left(C_{2}\right)=\left(\begin{array}{cc}
-1 & 0 \\
0 & 1
\end{array}\right) \\
\boldsymbol{D}\left(\sigma_{v}\right)=\left(\begin{array}{ll}
1 & 0 \\
0 & -1
\end{array}\right) & \boldsymbol{D}\left(\sigma_{v}^{\prime}\right)=\left(\begin{array}{cc}
-1 & 0 \\
0 & -1
\end{array}\right)
\end{array}
$$

The new representatives are all in block-diagonal form, in this case in the form $\left(\begin{array}{cc}\boldsymbol{\square} & 0 \\ 0 & \boldsymbol{\square}\end{array}\right)$, and the two combinations are not mixed with each other by any operation of the group. The representation $\Gamma^{\prime}$ has therefore been reduced to the sum of two one-dimensional representations. Thus, $\mathrm{p}_{1}$ spans the onedimensional representation

$$
\boldsymbol{D}(E)=1 \quad \boldsymbol{D}\left(C_{2}\right)=-1 \quad \boldsymbol{D}\left(\sigma_{v}\right)=1 \quad \boldsymbol{D}\left(\sigma_{v}^{\prime}\right)=-1
$$

which is the same as the representation $\Gamma^{(1)}$ spanned by $\mathrm{p}_{x}$. The combination $\mathrm{p}_{2}$ spans

$$
\boldsymbol{D}(E)=1 \quad \boldsymbol{D}\left(C_{2}\right)=1 \quad \boldsymbol{D}\left(\sigma_{\mathrm{v}}\right)=-1 \quad \boldsymbol{D}\left(\sigma_{\mathrm{v}}^{\prime}\right)=-1
$$

which is a new one-dimensional representation and denoted $\Gamma^{(4)}$. At this stage the original representation has been reduced into five one-dimensional representations as follows:

$$
\Gamma=2 \Gamma^{(1)}+\Gamma^{(2)}+\Gamma^{(3)}+\Gamma^{(4)}
$$

\begin{center}
\includegraphics[max width=\textwidth]{2025_02_27_d4b95d9718f0b9e2e6b9g-432}
\end{center}

Figure 10B. 5 Two symmetry-adapted linear combinations of the oxygen basis orbitals shown in Fig. 10B.4: on the left $p_{1}=p_{A}+$ $p_{B}$, on the right $p_{2}=p_{A}-p_{B}$. The two combinations each span a one-dimensional irreducible representation, and their symmetry species are different.

\subsection*{(d) Characters}
The character, $\chi$ (chi), of an operation in a particular matrix representation is the sum of the diagonal elements of the representative of that operation. Thus, in the original basis $\left(p_{x} p_{y}\right.$ $\mathrm{p}_{z} \mathrm{p}_{\mathrm{A}} \mathrm{p}_{\mathrm{B}}$ ) the characters of the representatives are

\begin{center}
\begin{tabular}{lcc}
\hline
$\boldsymbol{R}$ & $\boldsymbol{E}$ & \multicolumn{3}{c}{$C_{2}$} \\
\hline
$\boldsymbol{D}(R)$ &  &  &  &  \\
 & $\left(\begin{array}{ccccc}1 & 0 & 0 & 0 & 0 \\ 0 & 1 & 0 & 0 & 0 \\ 0 & 0 & 1 & 0 & 0 \\ 0 & 0 & 0 & 1 & 0 \\ 0 & 0 & 0 & 0 & 1\end{array}\right)$ & $\left(\begin{array}{ccccc}-1 & 0 & 0 & 0 & 0 \\ 0 & -1 & 0 & 0 & 0 \\ 0 & 0 & 1 & 0 & 0 \\ 0 & 0 & 0 & 0 & -1 \\ 0 & 0 & 0 & -1 & 0\end{array}\right)$ &  &  \\
$\chi(R)$ &  &  &  &  \\
\hline
\end{tabular}
\end{center}

\begin{center}
\begin{tabular}{lc}
\hline
$\boldsymbol{R}$ & $\sigma_{\mathrm{v}}$ \\
\hline
$\boldsymbol{D}(R)$ &  \\
 &  \\
$\chi(R)$ & $\left(\begin{array}{ccccc}1 & 0 & 0 & 0 & 0 \\ 0 & -1 & 0 & 0 & 0 \\ 0 & 0 & 1 & 0 & 0 \\ 0 & 0 & 0 & 0 & 1 \\ 0 & 0 & 0 & 1 & 0\end{array}\right)$ \\
\end{tabular}
\end{center}\$\textbackslash quad\textbackslash left(\textbackslash begin\{array\}\{ccccc\}-1 \& 0 \& 0 \& 0 \& 0 \\


0 \& 1 \& 0 \& 0 \& 0 \\
\\
0 \& 0 \& 1 \& 0 \& 0 \\
\\
0 \& 0 \& 0 \& -1 \& 0 \\
\\
0 \& 0 \& 0 \& 0 \& -1\textbackslash end\{array\}\textbackslash right)\$

The characters of one-dimensional representatives are just the representatives themselves. For each operation, the sum of the characters of the reduced representations is the same as the character of the original representation (allowing for the appearance of $\Gamma^{(1)}$ twice in the reduction $\Gamma=2 \Gamma^{(1)}+\Gamma^{(2)}+\Gamma^{(3)}$ $+\Gamma^{(4)}$ ):

\begin{center}
\begin{tabular}{llrrr}
\hline
$R$ & $E$ & $C_{2}$ & $\sigma_{\mathrm{v}}$ & $\sigma_{\mathrm{v}}^{\prime}$ \\
\hline
$\chi(R)$ for $\Gamma^{(1)}$ & 1 & -1 & 1 & -1 \\
$\chi(R)$ for $\Gamma^{(1)}$ & 1 & -1 & 1 & -1 \\
$\chi(R)$ for $\Gamma^{(2)}$ & 1 & -1 & -1 & 1 \\
$\chi(R)$ for $\Gamma^{(3)}$ & 1 & 1 & 1 & 1 \\
$\chi(R)$ for $\Gamma^{(4)}$ & 1 & 1 & -1 & -1 \\
Sum for $\Gamma:$ & 5 & -1 & 1 & -1 \\
\hline
\end{tabular}
\end{center}

At this point, four irreducible representations of the group $C_{2 v}$ have been found. Are these the only irreducible representations of the group $C_{2 \mathrm{v}}$ ? There are in fact no more irreducible representations in this group, a fact that can be deduced from a surprising theorem of group theory, which states that

Number of irreducible representations $=$ number of classes\\
Number of irreducible representations\\
(10B.5)

In $C_{2 \mathrm{v}}$ there are four classes of operations (the four columns in the table), so there must be four irreducible representations. The ones already found are the only ones for this group.

Another powerful result from group theory, which applies to all groups other than the pure rotation groups $C_{n}$ with $n>2$, relates the sum of the squares of the dimensions, $d_{i}$, of all the irreducible representations $\Gamma^{(i)}$ to the order of the group, which is the total number of symmetry operations, $h$ :


\begin{equation*}
\sum_{\substack{\text { irreducible } \\ \text { representations, } i}} d_{i}^{2}=h \quad \text { Dimensionality and order } \tag{10B.6}
\end{equation*}


The four irreducible representations of $C_{2 v}$ are all onedimensional, so

$$
\sum_{\substack{\text { irreducible } \\ \text { representations, } i}} d_{i}^{2}=1^{2}+1^{2}+1^{2}+1^{2}=4
$$

and there are indeed four symmetry operations of the group.

\subsection*{Brief illustration 10B. 3}
The group $C_{3 v}$ has three classes of operations $\left\{E, 2 C_{3}, 3 \sigma_{v}\right\}$, so there are three irreducible representations. The order of the group is $1+2+3=6$, so if it is already known that two of the irreducible representations are one-dimensional, the remaining irreducible representation must be two-dimensional by using eqn 10B.6: $1^{2}+1^{2}+d_{3}^{2}=6$, hence $d_{3}=2$.

\subsection*{10B. 3 Character tables}
Tables showing all the characters of the operations of a group are called character tables and from now on they move to the centre of the discussion. The columns of a character table are labelled with the symmetry operations of the group. Although the notation $\Gamma^{(i)}$ is used to label general irreducible representations, in chemical applications and for displaying character tables it is more common to distinguish different irreducible representations by the use of the labels $\mathrm{A}, \mathrm{B}, \mathrm{E}$, and T to denote the symmetry species of each representation:

A: one-dimensional representation, character +1 under the principal rotation\\
B: one-dimensional representation, character -1 under the principal rotation

E: two-dimensional irreducible representation\\
T : three-dimensional irreducible representation\\
Subscripts are used to distinguish the irreducible representations if there is more than one of the same type: $A_{1}$ is reserved for the representation with character 1 for all operations

Table 10B. 1 The $C_{2 v}$ character table*

\begin{center}
\begin{tabular}{lccccll}
\hline
$\boldsymbol{C}_{2 \mathrm{v}}, \mathbf{2 m m}$ & $\boldsymbol{E}$ & $\boldsymbol{C}_{2}$ & $\sigma_{\mathrm{v}}(x z)$ & $\sigma_{\mathrm{v}}^{\prime}(y z)$ & $\boldsymbol{h}=4$ &  \\
\hline
$\mathrm{~A}_{1}$ & 1 & 1 & 1 & 1 & $z$ & $z^{2}, y^{2}, x^{2}$ \\
$\mathrm{~A}_{2}$ & 1 & 1 & -1 & -1 &  & $x y$ \\
$\mathrm{~B}_{1}$ & 1 & -1 & 1 & -1 & $x$ & $z x$ \\
$\mathrm{~B}_{2}$ & 1 & -1 & -1 & 1 & $y$ & $y z$ \\
\hline
\end{tabular}
\end{center}

\begin{itemize}
  \item More character tables are given in the Resource section.\\
(called the totally symmetric irreducible representation); $\mathrm{A}_{2}$ has 1 for the principal rotation but -1 for reflections. There appears to be no systematic way of attaching subscripts to B symmetry species, so care must be used when referring to character tables from different sources.
\end{itemize}

Table 10B. 1 shows the character table for the group $C_{2 v}$, with its four symmetry species (irreducible representations) and its four columns of symmetry operations. Table 10B. 2 shows the table for the group $C_{3 \mathrm{v}^{\circ}}$. The columns are headed $E, 2 C_{3}$, and $3 \sigma_{\mathrm{v}}$ : the numbers multiplying each operation are the number of members of each class. As inferred in Brief illustration 10B.3, there are three symmetry species, with one of them two-dimensional (E).

Character tables, and some of the data contained in them, are constructed on the assumption that the axis system is arranged in a particular way and is specified in the character table when there is ambiguity. There is ambiguity in $C_{2 \mathrm{v}}$ (and certain other groups), and so a more detailed specification of the symmetry operations is then necessary. The principal axis (a unique $C_{n}$ axis with the greatest value of $n$ ), is taken to be the $z$-direction. If the molecule is planar the molecule is taken to lie in the $y z$-plane (referring to Fig. 10B.6). Then $\sigma_{\mathrm{v}}^{\prime}$ is a reflection in the $y z$-plane and henceforth will be denoted $\sigma_{\mathrm{v}}^{\prime}(y z)$, and $\sigma_{\mathrm{v}}$ is a reflection in the $x z$-plane, and henceforth is denoted $\sigma_{\mathrm{v}}(x z)$.

The irreducible representations are mutually orthogonal in the sense that if the set of characters is regarded as forming a row vector, the dot product (or scalar product) of the vectors corresponding to different irreducible representations is zero: the vectors are mutually perpendicular. ${ }^{2}$ The vectors are

Table 10B. 2 The $C_{3 v}$ character table*

\begin{center}
\begin{tabular}{llcrll}
\hline
$\boldsymbol{C}_{3 \mathbf{v}}, \mathbf{3 m}$ & $\boldsymbol{E}$ & $2 \boldsymbol{C}_{3}$ & $3 \sigma_{\mathrm{v}}$ & $\boldsymbol{h}=\mathbf{6}$ &  \\
\hline
$\mathrm{A}_{1}$ & 1 & 1 & 1 & $z$ & $z^{2}, x^{2}+y^{2}$ \\
$\mathrm{~A}_{2}$ & 1 & 1 & -1 &  &  \\
E & 2 & -1 & 0 & $(x, y)$ & $\left(x y, x^{2}-y^{2}\right),(y z, z x)$ \\
\hline
\end{tabular}
\end{center}

\begin{itemize}
  \item More character tables are given in the Resource section.
\end{itemize}

\footnotetext{${ }^{2}$ This result is a consequence of the 'great orthogonality theorem' of group theory; see our Molecular quantum mechanics (2011). In this Topic, the characters are taken to be real.
}
\includegraphics[max width=\textwidth, center]{2025_02_27_d4b95d9718f0b9e2e6b9g-434}

Figure 10B. 6 A $p_{x}$ orbital on the central atom of a $C_{2 v}$ molecule and the symmetry elements of the group.\\
also each normalized to 1 , in the sense that the dot product of a vector with itself is equal to 1 . Vectors that are both orthogonal and normalized (to 1) are said to be 'orthonormal'. Formally, this orthonormality is expressed as

$$
\frac{1}{h} \sum_{C} N(C) \chi^{\Gamma^{(i)}}(C) \chi^{\Gamma^{(j)}}(C)=\left\{\begin{array}{l}
0 \text { for } i \neq j \\
1 \text { for } i=j
\end{array}\right.
$$

Orthonormality of irreducible representations\\
(10B.7)\\
where the sum is over the classes of the group, $N(C)$ is the number of operations in class $C$, and $h$ is the number of operations in the group (its order).

\subsection*{Brief illustration 10B.4}
In the point group $C_{3 \mathrm{v}}$ with elements $\left\{E, 2 C_{3}, 3 \sigma_{\mathrm{v}}\right\}$ and $h=6$, for the two irreducible representations with labels $\mathrm{A}_{2}$ (with characters $\{1,1,-1\}$ ) and E (with characters $\{2,-1,0\}$ ), eqn 10B. 7 is

$$
\frac{1}{6}\{1 \times 1 \times 2+2 \times 1 \times(-1)+3 \times(-1) \times 0\}=0
$$

If the two irreducible representations are both E , the sum in eqn 10B. 7 is

$$
\frac{1}{6}\{1 \times 2 \times 2+2 \times(-1) \times(-1)+3 \times 0 \times 0\}=1
$$

The sum is also 1 if both irreducible representations are $A_{2}$ :

$$
\frac{1}{6}\{1 \times 1 \times 1+2 \times 1 \times 1+3 \times(-1) \times(-1)\}=1
$$

\subsection*{(a) The symmetry species of atomic orbitals}
The characters in the rows of one-dimensional irreducible representations (the rows labelled A or B ) and in the columns headed by symmetry operations indicate the behaviour of an orbital under the corresponding operations: a 1 indicates that an orbital is unchanged, and a-1 indicates that it changes sign. It follows that the symmetry label of the orbital can be identi-\\
fied by comparing the changes that occur to an orbital under each operation, and then comparing the resulting 1 or -1 with the entries in a row of the character table for the relevant point group. By convention, the orbitals are labelled with the lower case equivalent of the symmetry species label (so an orbital of symmetry species $A_{1}$ is called an $a_{1}$ orbital).

\subsection*{Brief illustration 10B. 5}
Consider an $\mathrm{H}_{2} \mathrm{O}$ molecule, point group $C_{2 v}$, shown in Fig. 10B.6. The effect of $C_{2}$ on the oxygen $2 p_{x}$ orbital is to cause it to change sign, so the character is $-1 ; \sigma_{\mathrm{v}}^{\prime}(y z)$ has the same effect and so has character -1 . In contrast, $\sigma_{\mathrm{v}}(x z)$ leaves the orbital unaffected and so has character 1 , and of course the same is true of the identity operation. The characters of the operations $\left\{E, C_{2}, \sigma_{v}, \sigma_{v}^{\prime}\right\}$ are therefore $\{1,-1,1,-1\}$. Reference to the $C_{2 v}$ character table (Table 10B.1) shows that $\{1,-1,1,-1\}$ are the characters for the symmetry species $B_{1}$; the orbital is therefore labelled $b_{1}$. A similar procedure gives the characters for oxygen $2 p_{y}$ as $\{1,-1,-1,1\}$, which corresponds to $B_{2}$ : the orbital is labelled $\mathrm{b}_{2}$, therefore. Both the oxygen $2 \mathrm{p}_{z}$ and 2 s are $\mathrm{a}_{1}$.

The characters in a row of the table for irreducible representations of dimensionality greater than 1 (typically, but not only, the E and T symmetry species) are the sums of the characters for the behaviour of the individual orbitals in the basis. Thus, if one member of a pair remains unchanged under a symmetry operation but the other changes sign (Fig. 10B.7), then the entry is reported as $\chi=1-1=0$.

The behaviour of $\mathrm{s}, \mathrm{p}$, and d orbitals on a central atom under the symmetry operations of the molecule is so important that the symmetry species of these orbitals are generally indicated in a character table. To make these assignments, identify the symmetry species of $x, y$, and $z$, which appear on the right hand side of the character table. Thus, the position of $z$ in Table 10B. 2 shows that $\mathrm{p}_{z}$ (which is proportional to $z f(r)$ ),\\
has symmetry species $\mathrm{A}_{1}$ in $C_{3 \mathrm{v}}$, whereas $\mathrm{p}_{x}$ and $\mathrm{p}_{y}$ (which are proportional to $x f(r)$ and $y f(r)$, respectively) are jointly of E symmetry. In technical terms, it is said that $\mathrm{p}_{x}$ and $\mathrm{p}_{y}$ jointly span an irreducible representation of symmetry species E. An s orbital on the central atom always spans the totally symmetric irreducible representation of a group as it is unchanged under all symmetry operations; in $C_{3 \mathrm{v}}$ it has symmetry species $\mathrm{A}_{1}$.

The five d orbitals of a shell are represented by $x y$ for $\mathrm{d}_{x y}$ etc. and are also listed on the right of the character table. It can be seen at a glance that in $C_{3 \mathrm{v}} \mathrm{d}_{x y}$ and $\mathrm{d}_{x^{2}-y^{2}}$ on a central atom jointly span E.

\subsection*{(b) The symmetry species of linear combinations of orbitals}
The same technique may be applied to identify the symmetry species of linear combinations of orbitals, such as the combination $\psi_{1}=\mathrm{s}_{\mathrm{A}}+\mathrm{s}_{\mathrm{B}}+\mathrm{s}_{\mathrm{C}}$ of the three H1s orbitals in the $C_{3 \mathrm{v}}$ molecule $\mathrm{NH}_{3}$ (Fig. 10B.8). This combination remains unchanged under a $C_{3}$ rotation and under any of the three vertical reflections of the group, so its characters are

$$
\chi(E)=1 \quad \chi\left(C_{3}\right)=1 \quad \chi\left(\sigma_{v}\right)=1
$$

Comparison with the $C_{3 v}$ character table shows that $\psi_{1}$ is of symmetry species $A_{1}$, and therefore has the label $a_{1}$.\\
\includegraphics[max width=\textwidth, center]{2025_02_27_d4b95d9718f0b9e2e6b9g-435}

Figure 10B. 8 The three H 1 s orbitals used to construct symmetryadapted linear combinations in a $\mathrm{C}_{3 \mathrm{v}}$ molecule such as $\mathrm{NH}_{3}$.

\subsection*{Example 10B. 2 \\
 Identifying the symmetry species \\
 of orbitals}
Identify the symmetry species of the orbital $\psi=\psi_{\mathrm{A}}-\psi_{\mathrm{B}}$ in a $C_{2 \mathrm{v}} \mathrm{NO}_{2}$ molecule, where $\psi_{\mathrm{A}}$ is an $\mathrm{O}_{2} \mathrm{p}_{x}$ orbital on one O atom (and perpendicular to the molecular plane) and $\psi_{B}$ that on the other O atom.

Collect your thoughts The negative sign in $\psi$ indicates that the sign of $\psi_{B}$ is opposite to that of $\psi_{A}$. You need to consider how the combination changes under each operation of the group, and then write the character as $1,-1$, or 0 as specified above.

Then compare the resulting characters with each row in the character table for the point group, and hence identify the symmetry species.

The solution The combination is shown in Fig. 10B.9. Under $C_{2}, \psi$ changes into itself, implying a character of 1 . Under the reflection $\sigma_{\mathrm{v}}(x z)$ both atomic orbitals change sign, so $\psi \rightarrow-\psi$, implying a character of -1 . Under $\sigma_{\mathrm{v}}^{\prime}(y z) \psi \rightarrow-\psi$, so the character for this operation is also -1 . The characters are therefore

$$
\chi(E)=1 \quad \chi\left(C_{2}\right)=1 \quad \chi\left(\sigma_{\mathrm{v}}(x z)\right)=-1 \quad \chi\left(\sigma_{\mathrm{v}}^{\prime}(y z)\right)=-1
$$

These values match the characters of the $\mathrm{A}_{2}$ symmetry species, so $\psi$ is labelled $\mathrm{a}_{2}$.\\
\includegraphics[max width=\textwidth, center]{2025_02_27_d4b95d9718f0b9e2e6b9g-436(1)}

Figure 10B.9 One symmetry-adapted linear combination of $\mathrm{O}_{2} \mathrm{p}_{x}$ orbitals in the $\mathrm{C}_{2 \mathrm{v}} \mathrm{NO}_{2}$ molecule.

Self-test 10B.2 Consider $\mathrm{PtCl}_{4}^{2-}$, in which the Cl ligands form a square planar array and the ion belongs to the point group $D_{4 \mathrm{~h}}$ (1). Identify the symmetry species of the combination $\psi_{A}-\psi_{\mathrm{B}}+$ $\psi_{\mathrm{C}}-\psi_{\mathrm{D}}$. Note that in this group the $C_{2}$ axes coincide with the $x$ and $y$ axes, and $\sigma_{v}$ planes coincide with the $x z$ - and $y z$ planes; choose the $x$ - and $y$-axes to pass through the corners of the square.\\
\includegraphics[max width=\textwidth, center]{2025_02_27_d4b95d9718f0b9e2e6b9g-436(2)}\\
\includegraphics[max width=\textwidth, center]{2025_02_27_d4b95d9718f0b9e2e6b9g-436(3)}

\subsection*{(c) Character tables and degeneracy}
In Topic 7D it is pointed out that degeneracy, which is when different wavefunctions have the same energy, is always related to symmetry, and that an energy level is degenerate if the wavefunctions corresponding to that energy can be transformed into each other by a symmetry operation (such as rotating a square well through $90^{\circ}$ ). Clearly, group theory should have a role in the identification of degeneracy.

A geometrically square well belongs to the group $C_{4}$ (Fig. 10B. 10 and Table 10B.3), with the $C_{4}$ rotations (through $90^{\circ}$ ) converting $x$ into $y$ and vice versa. ${ }^{3}$ As explained in Topic 7D,

\footnotetext{${ }^{3}$ More complicated groups could be used, such as $C_{4 \mathrm{v}}$ or $D_{4 \mathrm{~h}}$, but $C_{4}$ captures the symmetry sufficiently.
}
the two wavefunctions $\psi_{1,2}=(2 / L) \sin (\pi x / L) \sin (2 \pi y / L)$ and $\psi_{2,1}=(2 / L) \sin (2 \pi x / L) \sin (\pi y / L)$ both correspond to the energy $5 h^{2} / 8 m L^{2}$, so that level is doubly degenerate. Under the operations of the group, these two functions transform as follows:\\
\$\$

\[
\begin{array}{ll}
E:\left(\psi_{1,2} \psi_{2,1}\right) \rightarrow\left(\psi_{1,2} \psi_{2,1}\right) & C_{4}^{+}:\left(\psi_{1,2} \psi_{2,1}\right) \rightarrow\left(\psi_{2,1}-\psi_{1,2}\right) \\
C_{4}^{-}:\left(\psi_{1,2} \psi_{2,1}\right) \rightarrow\left(-\psi_{2,1} \psi_{1,2}\right) & C_{2}:\left(\psi_{1,2} \psi_{2,1}\right) \rightarrow\left(-\psi_{1,2}-\psi_{2,1}\right)
\end{array}
\]

\$\$

The corresponding matrix representatives are

$$
\begin{array}{ll}
\boldsymbol{D}(E)=\left(\begin{array}{ll}
1 & 0 \\
0 & 1
\end{array}\right) & \boldsymbol{D}\left(C_{4}^{+}\right)=\left(\begin{array}{cc}
0 & 1 \\
-1 & 0
\end{array}\right) \\
\boldsymbol{D}\left(C_{4}^{-}\right)=\left(\begin{array}{cc}
0 & -1 \\
1 & 0
\end{array}\right) & \boldsymbol{D}\left(C_{2}\right)=\left(\begin{array}{cc}
-1 & 0 \\
0 & -1
\end{array}\right)
\end{array}
$$

and their characters are

$$
\chi(E)=2 \quad \chi\left(C_{4}^{+}\right)=0 \quad \chi\left(C_{4}^{-}\right)=0 \quad \chi\left(C_{2}\right)=-2
$$

A glance at the character table in Table 10B. 3 (noting that the rotations $C_{4}^{+}$and $C_{4}^{-}$belong to the same class and appear in the column labelled $2 C_{4}$ ) shows that the basis spans the irreducible representation of symmetry species $E$. The same is true of all the doubly-degenerate energy levels, and there are no triplydegenerate (or higher) energy levels in the system. Notice too that in the group $C_{4}$ there are no irreducible representations of dimension 3 or higher. These two observations illustrate the general principle that:

\begin{displayquote}
The highest dimensionality of irreducible representation in a group is the maximum degree of degeneracy in the group.
\end{displayquote}

Thus, if there is an E irreducible representation in a group, 2 is the highest degree of degeneracy; if there is a T irreducible representation in a group, then 3 is the highest degree of\\
\includegraphics[max width=\textwidth, center]{2025_02_27_d4b95d9718f0b9e2e6b9g-436}

Figure 10B.10 A geometrically square well can be treated as belonging to the group $C_{4}$ (with elements $\left\{E_{,} 2 C_{4}, C_{2}\right\}$ ).

Table 10B. 3 The $C_{4}$ character table ${ }^{*}$

\begin{center}
\begin{tabular}{llcrll}
\hline
$\boldsymbol{C}_{4}, \mathbf{4}$ & $\boldsymbol{E}$ & $\mathbf{2 C}$ & $\boldsymbol{C}_{\mathbf{2}}$ & $\boldsymbol{h}=\mathbf{4}$ &  \\
\hline
A & 1 & 1 & 1 & $z$ & $z^{2}, x^{2}+y^{2}$ \\
B & 1 & -1 & 1 &  & $x y, x^{2}-y^{2}$ \\
E & 2 & 0 & -2 & $(x, y)$ & $(y z, z x)$ \\
\hline
\end{tabular}
\end{center}

\begin{itemize}
  \item More character tables are given in the Resource section.\\
degeneracy. Some groups have irreducible representations of higher dimension, and therefore allow higher degrees of degeneracy. Furthermore, because the character of the identity operation is always equal to the dimensionality of the representation, the maximum degeneracy can be identified by noting the maximum value of $\chi(E)$ in the relevant character table.
\end{itemize}

\subsection*{Brief illustration 10B.6}
\begin{itemize}
  \item A trigonal planar molecule such as $\mathrm{BF}_{3}$ cannot have triply degenerate orbitals because its point group is $D_{3 \mathrm{~h}}$ and the character table for this group (in the Resource section) does not have a T symmetry species.
  \item A methane molecule belongs to the tetrahedral point group $T_{\mathrm{d}}$ and because that group has irreducible representations of T symmetry, it can have triply-degenerate orbitals. The same is true of tetrahedral $P_{4}$, which, with just four atoms, is the simplest kind of molecule with triply-degenerate orbitals.
  \item A buckminsterfullerene molecule, $\mathrm{C}_{60}$, belongs to the icosahedral point group $\left(I_{\mathrm{h}}\right)$ and its character table (in the Resource section) shows that the maximum dimensionality of its irreducible representations is 5 , so it can have fivefold degenerate orbitals.
\end{itemize}

\subsection*{Checklist of concepts}
$\square$ 1. A group is a collection of transformations that satisfy the four criteria set out at the start of the Topic.\\
2. The order of a group is the number of its symmetry operations.\\
$\square$ 3. A matrix representative is a matrix that represents the effect of an operation on a basis.4. The character is the sum of the diagonal elements of a matrix representative of an operation.\\
$\square$ 5. A matrix representation is the collection of matrix representatives for the operations in the group.\\
$\square$ 6. A character table consists of entries showing the characters of all the irreducible representations of a group.\\
$\square$ 7. A symmetry species is a label for an irreducible representation of a group.\\
$\square$ 8. The highest dimensionality of irreducible representation in a group is the maximum degree of degeneracy in the group.

Checklist of equations

\begin{center}
\begin{tabular}{|c|c|c|c|}
\hline
Property & Equation & Comment & Equation number \\
\hline
Class membership & $R^{\prime}=S^{-1} R S$ & All elements members of the group & 10B. 1 \\
\hline
Number of irreducible representations & Number of irreducible representations $=$ number of classes &  & 10B. 5 \\
\hline
Dimensionality and order & \( \sum_{\text {irreps } i} d_{i}^{2}=h \) & For groups other than pure rotation groups with $n>2$ & 10B.6 \\
\hline
Orthonormality of irreducible representations & \( \begin{array}{r} \frac{1}{h} \sum_{C} N(C) \chi^{\Gamma^{(i)}}(C) \chi^{\Gamma^{(j)}}(C) \\ =\left\{\begin{array}{l} 0 \text { for } i \neq j \\ 1 \text { for } i=j \end{array}\right. \end{array} \) & Sum over classes & 10B. 7 \\
\hline
\end{tabular}
\end{center}

\section*{TOPIC 10C Applications of symmetry}
\subsection*{- Why do you need to know this material?}
Group theory is a key tool for constructing molecular orbitals and formulating spectroscopic selection rules.

\subsection*{> What is the key idea?}
An integral can be non-zero only if the integrand is invariant under the symmetry transformations of a molecule.

\subsection*{> What do you need to know already?}
This Topic develops the material in Topic 10A, where the classification of molecules on the basis of their symmetry elements is introduced, and draws heavily on the properties of characters and character tables described in Topic 10B. You need to be aware that many quan-tum-mechanical properties, including transition dipole moments (Topic 8C), depend on integrals involving products of wavefunctions (Topic 7C).

Group theory shows its power when brought to bear on a variety of problems in chemistry, among them the construction of molecular orbitals and the formulation of spectroscopic selection rules.

\subsection*{10C. 1 Vanishing integrals}
Any integral, $I$, of a function $f(x)$ over a symmetric range around $x=0$ is zero if the function is antisymmetric in the sense that $f(-x)=-f(x)$. In two dimensions the integral (over a symmetrical range) of the integrand $f(x, y)$ has contributions from regions that are related by symmetry operations of the area of integration (Fig. 10C.1). If $f(x, y)$ changes sign under one of these operations, the contribution of the first region is cancelled by that from the symmetry-related region and the integral is zero. The integral may be non-zero only if the integrand is invariant (or at least can be expressed as a sum of terms at least one of which is invariant) under each symmetry operation of the group that reflects the shape of the area (and in general, the volume) of the range of integration. In group-theoretical terms:

\footnotetext{An integral over a region of space can be non-zero only if the integrand (or a contribution to it) spans the totally symmetric irreducible representation of the point group of the region.
}
\includegraphics[max width=\textwidth, center]{2025_02_27_d4b95d9718f0b9e2e6b9g-438(1)}

Figure 10C. 1 (a) Only if the integrand is unchanged under each symmetry operation of the group (here $C_{4}$ ) can its integral over the region indicated be non-zero. (b) If the integrand changes sign under any operation, its integral is necessarily zero.

The totally symmetric irreducible representation has all characters equal to 1 , and is typically the symmetry species denoted $A_{1}$.

\subsection*{Brief illustration 10C. 1}
To decide whether the integral of the function $f=x y$ may be non-zero when evaluated over a region the shape of an equilateral triangle centred on the origin (Fig. 10C.2), recognize that the triangle belongs to the group $C_{3}$. Reference to the character table of the group shows that $x y$ is a member of a basis that spans the irreducible representation E. Therefore, its integral must be zero, because the integrand has no component that spans $\mathrm{A}_{1}$.\\
\includegraphics[max width=\textwidth, center]{2025_02_27_d4b95d9718f0b9e2e6b9g-438}

Figure 10C. 2 The integral of the function $f=x y$ over the yellow region (of symmetry $C_{3}$ ) is zero. In this case, the result is obvious by inspection, but group theory can be used to establish similar results in less obvious cases. The insert shows the shape of the function in three dimensions.

\subsection*{(a) Integrals of the product of functions}
Suppose the integral of interest is of a product of two functions, $f_{1}$ and $f_{2}$, taken over all space and over all relevant variables (represented, as is usual in quantum mechanics, by integration over $\mathrm{d} \tau$ ):


\begin{equation*}
I=\int f_{1} f_{2} \mathrm{~d} \tau \tag{10C.1}
\end{equation*}


For example, $f_{1}$ and $f_{2}$ might be atomic orbitals on different atoms, in which case $I$ would be their overlap integral. The implication of such an integral being zero is that a molecular orbital does not result from the overlap of these two orbitals. It follows from the general point made above that the integral may be non-zero only if the integrand itself, the product $f_{1} f_{2}$, is unchanged by any symmetry operation of the molecular point group and so spans the totally symmetric irreducible representation (typically the symmetry species with the label $\mathrm{A}_{1}$ ). To decide whether the product $f_{1} f_{2}$ does indeed span $\mathrm{A}_{1}$, it is necessary to form the direct product of the symmetry species spanned by $f_{1}$ and $f_{2}$ separately. The procedure is as follows:

\begin{itemize}
  \item Write down a table with columns headed by the symmetry operations, $R$, of the group.
  \item In the first row write down the characters of the symmetry species spanned by $f_{1}$; in the second row write down the characters of the symmetry species spanned by $f_{2}$.
  \item Multiply the numbers in the two rows together, column by column. The resulting set of numbers are the characters of the representation spanned by $f_{1} f_{2}$.
\end{itemize}

\subsection*{Brief illustration 10C. 2}
Suppose that in the point group $C_{2 v} f_{1}$ has the symmetry species $\mathrm{A}_{2}$, and $f_{2}$ has the symmetry species $\mathrm{B}_{1}$. From the character table the characters for these species are $1,1,-1,-1$ and $1,-1,1,-1$, respectively. The direct product of these two species is found by setting up the following table

\begin{center}
\begin{tabular}{llrlc}
\hline
 & $\boldsymbol{E}$ & $\boldsymbol{C}_{2}$ & $\sigma_{\mathrm{v}}(x z)$ & $\sigma_{\mathrm{v}}^{\prime}(y z)$ \\
\hline
$\mathrm{A}_{2}$ & 1 & 1 & -1 & -1 \\
$\mathrm{~B}_{1}$ & 1 & -1 & 1 & -1 \\
product & 1 & -1 & -1 & 1 \\
\hline
\end{tabular}
\end{center}

Now recognize that the characters in the final row are those of the symmetry species $B_{2}$. It follows that the symmetry species of the product $f_{1} f_{2}$ is $\mathrm{B}_{2}$. Because the direct product does not contain $\mathrm{A}_{1}$, the integral of $f_{1} f_{2}$ over all space must be zero.

Direct products have some simplifying features.

\begin{itemize}
  \item The direct product of the totally symmetric irreducible representation with any other representation is the latter irreducible representation itself: $\mathrm{A}_{1} \times \Gamma^{(i)}=\Gamma^{(i)}$.
\end{itemize}

All the characters of $\mathrm{A}_{1}$ are 1 , so multiplication by them leaves the characters of $\Gamma^{(i)}$ unchanged. It follows that if one of the functions in eqn 10C. 1 transforms as $\mathrm{A}_{1}$, then the integral will vanish if the other function is not $A_{1}$.

\begin{itemize}
  \item The direct product of two irreducible representations is $\mathrm{A}_{1}$ only if the two irreducible representations are identical: $\Gamma^{(i)} \times \Gamma^{(j)}$ contains $\mathrm{A}_{1}$ only if $i=j$.
\end{itemize}

For one-dimensional irreducible representations the characters are either 1 or -1 , and the character 1 is obtained only if the characters of $\Gamma^{(i)}$ and $\Gamma^{(j)}$ are the same (both 1 or both -1 ). For example, in $C_{2 v}, A_{1} \times A_{1}, A_{2} \times A_{2}, B_{1} \times B_{1}$, and $B_{2} \times B_{2}$, but no other combination, all give $A_{1}$. That the requirement also holds for higher-dimensional representations requires more work and is demonstrated at the end of Section 10C.1b.

It follows that if $f_{1}$ and $f_{2}$ transform as different symmetry species, then the product cannot transform as the totally symmetric irreducible representation and so the integral of $f_{1} f_{2}$ is necessarily zero. If, on the other hand, both functions transform as the same symmetry species, then the product transforms as the totally symmetric irreducible representation (and possibly has contributions from other symmetry species too) and the integral is not necessarily zero.

An important point is that group theory is specific about when an integral must be zero, but integrals that it allows to be non-zero may be zero for reasons unrelated to symmetry. For example, the $\mathrm{N}-\mathrm{H}$ distance in ammonia may be so long that the $\left(s_{N}, s_{1}+s_{2}+s_{3}\right)$ overlap integral, in which $f_{1}$ is a 2 s orbital on N and $f_{2}$ is a combination of 1 s orbitals on the three H atoms with the same symmetry, is zero simply because the orbitals are so far apart.

Integrals of the form


\begin{equation*}
I=\int f_{1} f_{2} f_{3} \mathrm{~d} \tau \tag{10C.2}
\end{equation*}


are also common in quantum mechanics, and it is important to know when they are necessarily zero. For example, they appear in the calculation of transition dipole moments (Topic 8 C ). As for integrals over two functions, for $I$ to be non-zero, the product $f_{1} f_{2} f_{3}$ must span the totally symmetric irreducible representation or contain a component that spans that representation. To test whether this is so, the characters of all three irreducible representations are multiplied together in the same way as in the rules set out above.

\subsection*{Example 10C. 1 Deciding if an integral must be zero}
Does the integral $\int\left(\mathrm{d}_{z^{2}}\right) x\left(\mathrm{~d}_{x y}\right) \mathrm{d} \tau$ vanish in a $C_{2 \mathrm{v}}$ molecule?\\
Collect your thoughts Use the $C_{2 v}$ character table to find the characters of the irreducible representations spanned by $3 z^{2}-r^{2}$ (the form of the $\mathrm{d}_{z^{2}}$ orbital), $x$, and $x y$. Then set up a table to work out the triple direct product and identify whether the symmetry species it spans includes $A_{1}$.

The solution The $C_{2 v}$ character table shows that the function $x y$, and hence the orbital $\mathrm{d}_{x y}$, transforms as $\mathrm{A}_{2}$, that $z^{2}$ transforms as $\mathrm{A}_{1}$, and that $x$ transforms as $\mathrm{B}_{1}$. The table is therefore

\begin{center}
\begin{tabular}{llrlll}
\hline
 & $\boldsymbol{E}$ & $\boldsymbol{C}_{2}$ & $\sigma_{v}(\boldsymbol{x} \boldsymbol{z})$ & $\sigma_{v}^{\prime}(\boldsymbol{y z})$ &  \\
\hline
$\mathrm{A}_{2}$ & 1 & 1 & -1 & -1 & $f_{3}=\mathrm{d}_{x y}$ \\
$\mathrm{~B}_{1}$ & 1 & -1 & 1 & -1 & $f_{2}=x$ \\
$\mathrm{~A}_{1}$ & 1 & 1 & 1 & 1 & $f_{1}=\mathrm{d}_{z^{2}}$ \\
 & 1 & -1 & -1 & 1 & product \\
\hline
\end{tabular}
\end{center}

The characters in the bottom row are those of $\mathrm{B}_{2}$, not of $\mathrm{A}_{1}$. Therefore, the integral is necessarily zero.

Comment. A quicker solution involves noting that $\mathrm{A}_{1}\left(\right.$ for $\left.f_{1}\right)$ has no effect on the outcome of the triple direct product (by the first feature mentioned above), and therefore, by the second feature, the symmetry species of the two functions $f_{2}$ and $f_{3}$ must be the same for their direct product to be $\mathrm{A}_{1}$; but in this example they are not the same.

Self-test 10C. 1 Does the integral $\int\left(\mathrm{d}_{x z}\right) x\left(\mathrm{p}_{z}\right) \mathrm{d} \tau$ necessarily vanish in a $C_{2 \mathrm{v}}$ molecule?\\
$\mathrm{o}_{\mathrm{N}}: \iota \partial м u_{V}$

\subsection*{(b) Decomposition of a representation}
In some cases, it turns out that the direct product is a sum of symmetry species, not just a single species. For instance, in $C_{3 v}$ the characters of the direct product $\mathrm{E} \times \mathrm{E}$ are $\{4,1,0\}$, which can be decomposed as $\mathrm{A}_{1}, \mathrm{~A}_{2}$, and E :

\begin{center}
\begin{tabular}{lccc}
\hline
 & $\boldsymbol{E}$ & $2 \boldsymbol{C}_{3}$ & $3 \sigma_{\mathrm{v}}$ \\
\hline
$\mathrm{A}_{1}$ & 1 & 1 & 1 \\
$\mathrm{~A}_{2}$ & 1 & 1 & -1 \\
E & 2 & -1 & 0 \\
sum & 4 & 1 & 0 \\
\hline
\end{tabular}
\end{center}

This decomposition is written symbolically $\mathrm{E} \times \mathrm{E}=\mathrm{A}_{1}+\mathrm{A}_{2}$ + E. ${ }^{1}$

In simple cases the decomposition can be done by inspection. Group theory, however, provides a systematic way of using the characters of the representation to find the symmetry species of the irreducible representations of which it is composed. The formal recipe for finding the number of times, $n(\Gamma)$, that irreducible representation $\Gamma$ occurs is based on a

\footnotetext{${ }^{1}$ As mentioned in Topic 10B, a direct sum is sometimes denoted $\oplus$. The analogous symbol for a direct product is $\otimes$. The symbolic expression is then written $\mathrm{E} \otimes \mathrm{E}=\mathrm{A}_{1} \oplus \mathrm{~A}_{2} \oplus \mathrm{E}$.
}
general expression derived from a very deep result of group theory: ${ }^{2}$\\
\$\$

n(\textbackslash Gamma)=\textbackslash frac\{1\}\{h\} \textbackslash sum\_\{C\} N(C) \textbackslash chi\^{}\{(\textbackslash Gamma)\}(C) \textbackslash chi(C) \textbackslash quad \[
\begin{aligned}
& \text { Decomposition of } \\
& \text { a representation }
\end{aligned}
\] \textbackslash quad \text{(10C.3a)}

\$\$

Here $\Gamma$ is the symmetry species of the irreducible representation of interest, $h$ is the order of the group, $\chi^{(\mathrm{T})}(C)$ is the character of the members of class $C$ of operations for that irreducible representation, and $\chi(C)$ is the corresponding character of the representation being decomposed. Note that the sum is over the classes of operations. In the character table the number of operations in each class, $N(C)$, is indicated in the header of the columns. All the characters of the totally symmetric irreducible representation of symmetry species $\mathrm{A}_{1}$ are 1 , so setting $\Gamma=$ $\mathrm{A}_{1}$ and $\chi^{\left(A_{1}\right)}(C)=1$ for all $C$ in eqn 10C.3a gives

$$
n\left(\mathrm{~A}_{1}\right)=\frac{1}{h} \sum_{C} N(C) \chi(C) \quad \text { Occurrence of } \mathrm{A}_{1}
$$

(10C.3b)

\subsection*{Brief illustration 10C. 3}
In the character table for $C_{3 v}$, the columns are headed $E, 2 C_{3}$, and $3 \sigma_{\mathrm{v}}$ indicating that the numbers in each class are 1,2 , and 3 , respectively and $h=1+2+3=6$. To decide whether $\mathrm{A}_{1}$ occurs in the representation with characters $\{4,1,0\}$ in $C_{3 \mathrm{v}}$ form

$$
\begin{aligned}
n\left(\mathrm{~A}_{1}\right) & =\frac{1}{6}\left\{1 \times \chi(E)+2 \times \chi\left(C_{3}\right)+3 \times \chi\left(\sigma_{v}\right)\right\} \\
& =\frac{1}{6}\{1 \times 4+2 \times 1+3 \times 0\}=1
\end{aligned}
$$

$\mathrm{A}_{1}$ therefore occurs once in the decomposition.

It is asserted in Section 10C.1a that the direct product of two irreducible representations is $A_{1}$ only if the two irreducible representations are identical. That this is so can now be shown with the aid of eqn 10C.3b.

\subsection*{How is that done? 10C. 1 Confirming the criterion for a}
 direct product to contain the totally symmetric irreducible representationStart by considering the characters of the direct product between irreducible representations $\Gamma^{(i)}$ and $\Gamma^{(j)}$. The character of a class of operations in a direct product is the product of the characters of the two contributing representations: $\chi(\mathrm{C})$ $=\chi^{\Gamma^{(i)}}(C) \chi^{\Gamma^{(j)}}(C)$, where $\chi^{\Gamma^{(i)}}(C)$ is the character for the operation in class $C$ of irreducible representation $\Gamma^{(i)}$ and likewise

\footnotetext{${ }^{2}$ This result arises from the 'great orthogonality theorem': see our Molecular quantum mechanics (2011). In this Topic, the characters are taken to be real.
}
for $\chi^{\Gamma^{(j)}}(C)$. The number of times that the totally symmetric irreducible representation $\left(\mathrm{A}_{1}\right)$ occurs in this direct-product representation is given by eqn 10 C .3 b as
$$
n\left(\mathrm{~A}_{1}\right)=\frac{1}{h} \sum_{C} N(C) \chi^{\Gamma^{(i)}}(C) \chi^{\Gamma^{(j)}}(C)
$$

Irreducible representations are orthonormal in the sense that (eqn 10B.7)

$$
\frac{1}{h} \sum_{C} N(C) \chi^{\Gamma^{(i)}}(C) \chi^{\Gamma^{(j)}}(C)=\left\{\begin{array}{c}
0 \text { if } i \neq j \\
1 \text { if } i=j
\end{array}\right.
$$

It follows that

$$
n\left(\mathrm{~A}_{1}\right)=\left\{\begin{array}{l}
0 \text { if } i \neq j \\
1 \text { if } i=j
\end{array}\right.
$$

In other words, the direct product of two irreducible representations has a component that spans $\mathrm{A}_{1}$ only if the two irreducible representations belong to the same symmetry species. This result is independent of the dimensionality of the irreducible representations.

\subsection*{10c. 2 Applications to molecular orbital theory}
The rules outlined so far can be used to decide which atomic orbitals may have non-zero overlap in a molecule. Group theory also provides procedures for constructing linear combinations of atomic orbitals of a specified symmetry.

\subsection*{(a) Orbital overlap}
The overlap integral, $S$, between orbitals $\psi_{1}$ and $\psi_{2}$ is

$$
S=\int \psi_{2}^{\star} \psi_{1} \mathrm{~d} \tau \quad \text { Overlap integral }
$$

It follows from the discussion of eqn 10C. 1 that this integral can be non-zero only if the two orbitals span the same symmetry species. In other words,

\begin{displayquote}
Only orbitals of the same symmetry species may have non-zero overlap ( $S \neq 0$ ) and hence go on to form bonding and antibonding combinations.
\end{displayquote}

The selection of atomic orbitals with non-zero overlap is the central and initial step in the construction of molecular orbitals as LCAOs.

Example 10C. 2\\
Identifying which orbitals can contribute to bonding

The four H1s orbitals of methane span $A_{1}+T_{2}$. With which of the C2s and C2p atomic orbitals can they overlap? What additional overlap would be possible if $d$ orbitals on the $C$ atom were also considered?

Collect your thoughts Refer to the $T_{\mathrm{d}}$ character table (in the Resource section) and look for $\mathrm{s}, \mathrm{p}$, and d orbitals spanning $\mathrm{A}_{1}$ or $\mathrm{T}_{2}$. Recall that the symmetry species can be identified by looking for the appropriate Cartesian functions listed on the right of the table.

The solution A C2s orbital spans $\mathrm{A}_{1}$ in the group $T_{\mathrm{d}}$, so it may have non-zero overlap with the $\mathrm{A}_{1}$ combination of H1s orbitals. From the table $(x, y, z)$ jointly span $\mathrm{T}_{2}$, so the three C2p orbitals together transform in the same way; they may have non-zero overlap with the $\mathrm{T}_{2}$ combination of H1s orbitals. The combinations ( $x y, y z, x z$ ) span $\mathrm{T}_{2}$, therefore the $\mathrm{d}_{x y}, \mathrm{~d}_{y z}$, and $\mathrm{d}_{x z}$ orbitals do the same and so they may overlap with the $\mathrm{T}_{2}$ combination of H1s orbitals. The other two d orbitals span E and so they cannot overlap with the $\mathrm{A}_{1}$ or $\mathrm{T}_{2} \mathrm{H} 1$ s orbitals and remain nonbonding. It follows that in methane there are $a_{1}$ orbitals arising from (C2s,H1s)-overlap and $\mathrm{t}_{2}$ orbitals arising from (C2p,H1s)-overlap. The C3d orbitals might contribute to the latter. The lowest energy configuration is probably $\mathrm{a}_{1}^{2} \mathrm{t}_{2}^{6}$, with all bonding orbitals occupied.

Self-test 10C. 2 Consider the octahedral $\mathrm{SF}_{6}$ molecule, with the bonding arising from overlap of $s$ orbitals and a $2 p$ orbital on each fluorine directed towards the central sulfur atom. The latter span $A_{1 g}+E_{g}+T_{\text {lu }}$. Which sulfur orbitals have non-zero overlap with these F orbitals? Suggest what the ground-state configuration is likely to be.\\
\includegraphics[max width=\textwidth, center]{2025_02_27_d4b95d9718f0b9e2e6b9g-441}

\subsection*{(b) Symmetry-adapted linear combinations}
Topic 10B introduces the idea of generating a combination of atomic orbitals designed to transform as a particular symmetry species. Such a combination is an example of a symmetryadapted linear combination (SALC), which is a combination of orbitals constructed from equivalent atoms and having a specified symmetry. SALCs are very useful in constructing molecular orbitals because a given SALC has non-zero overlap only with other orbitals of the same symmetry.

The technique for building SALCs is derived by using the full power of group theory and involves the use of a projection operator, $P^{(\Gamma)}$, an operator that takes one of the basis orbitals and generates from it-projects from it-an SALC of the symmetry species $\Gamma$ :

$$
P^{(\Gamma)}=\frac{1}{h} \sum_{R} \chi^{(\Gamma)}(R) R \quad \psi^{(\Gamma)}=P^{(\Gamma)} \psi_{i} \quad \begin{aligned}
& \text { Projection } \\
& \text { operator }
\end{aligned}
$$

(10C.5)

Here $\psi_{i}$ is one of the basis orbitals and $\psi^{(\Gamma)}$ is a SALC (there might be more than one) that transforms as the symmetry species $\Gamma$; the sum is over the operations (not the classes) of the group of order $h$. To implement this rule, do the following:

\begin{itemize}
  \item Construct a table with the columns headed by each symmetry operation $R$ of the group; include a column for each operation, not just for each class.
  \item Select a basis function and work out the effect that each operation has on it. Enter the resulting function beneath each operation.
  \item On the next row enter the characters of the symmetry species of interest, $\chi^{(\mathrm{T})}(R)$.
  \item Multiply the entries in the previous two rows, operation by operation.
  \item Sum the result, and divide it by the order of the group, $h$.
\end{itemize}

\subsection*{Brief illustration 10C. 4}
To construct the $\mathrm{B}_{1}$ SALC from the two $\mathrm{O} 2 \mathrm{p}_{x}$ orbitals in $\mathrm{NO}_{2}$, point group $C_{2 v}$ (Fig. 10C.3), draw up the following table:

\begin{center}
\begin{tabular}{lllll}
\hline
 & $\boldsymbol{E}$ & $\boldsymbol{C}_{2}$ & $\sigma_{\mathrm{v}}(\boldsymbol{x} \boldsymbol{z})$ & $\sigma_{\mathrm{v}}^{\prime}(\boldsymbol{y} \boldsymbol{z})$ \\
\hline
Effect on $\mathrm{p}_{\mathrm{A}}$ & $\mathrm{p}_{\mathrm{A}}$ & $-\mathrm{p}_{\mathrm{B}}$ & $\mathrm{p}_{\mathrm{B}}$ & $-\mathrm{p}_{\mathrm{A}}$ \\
Characters for $\mathrm{B}_{1}$ & 1 & -1 & 1 & -1 \\
Product of rows 1 and 2 & $\mathrm{p}_{\mathrm{A}}$ & $\mathrm{p}_{\mathrm{B}}$ & $\mathrm{p}_{\mathrm{B}}$ & $\mathrm{p}_{\mathrm{A}}$ \\
\hline
\end{tabular}
\end{center}

The sum of the final row, divided by the order of the group $(h=4)$, gives $\psi^{\left(\mathrm{B}_{1}\right)}=\frac{1}{2}\left(\mathrm{p}_{\mathrm{A}}+\mathrm{p}_{\mathrm{B}}\right)$.\\
\includegraphics[max width=\textwidth, center]{2025_02_27_d4b95d9718f0b9e2e6b9g-442(1)}

Figure 10C. 3 The two $02 p_{x}$ atomic orbitals in $\mathrm{NO}_{2}$ (point group $C_{2 v}$ ) can be used as a basis for forming SALCs.

If an attempt is made to generate a SALC with symmetry that is not spanned by the basis functions, the result is zero. For example, if in the Brief illustration an attempt is made to project an $\mathrm{A}_{1}$ symmetry orbital, all the characters in the second row of the table will be 1 , so when the product of rows 1 and 2 is formed the result is $p_{A}-p_{B}+p_{B}-p_{A}=0$.

A difficulty is encountered when aiming to generate an SALC of symmetry species of dimension higher than 1 , because then the rules generate sums of SALCs. Consider, for\\
instance, the generation of SALCs from the three H1s atomic orbitals in $\mathrm{NH}_{3}$, point group $C_{3 v^{*}}$. The molecule and the orbitals are shown in Fig. 10C.4. The table below shows the effect of applying the projection operator to $\mathrm{s}_{\mathrm{A}}, \mathrm{s}_{\mathrm{B}}$, and $\mathrm{s}_{\mathrm{C}}$ in turn to give the SALC of symmetry species E.

\begin{center}
\begin{tabular}{|c|c|c|c|c|c|c|c|}
\hline
\multicolumn{2}{|l|}{Row} & E & $\mathrm{C}_{3}^{+}$ & $\mathrm{C}_{3}^{-}$ & $\sigma_{\mathrm{v}}$ & $\sigma_{\mathrm{v}}^{\prime}$ & $\sigma_{\mathrm{v}}^{\prime \prime}$ \\
\hline
1 & effect on $\mathrm{s}_{\text {A }}$ & $\mathrm{S}_{\text {A }}$ & $S_{\text {B }}$ & $\mathrm{s}_{\mathrm{C}}$ & $\mathrm{S}_{\text {A }}$ & $\mathrm{s}_{\mathrm{C}}$ & $\mathrm{s}_{\text {B }}$ \\
\hline
2 & characters for E & 2 & -1 & -1 & 0 & 0 & 0 \\
\hline
3 & product of rows 1 and 2 & $2 \mathrm{~s}_{\text {A }}$ & $-S_{\text {B }}$ & $-\mathrm{s}_{\mathrm{C}}$ &  &  &  \\
\hline
4 & effect on $s_{B}$ & $\mathrm{S}_{\text {B }}$ & $s_{C}$ & $\mathrm{s}_{\text {A }}$ & $\mathrm{s}_{\mathrm{C}}$ & $\mathrm{S}_{\text {B }}$ & $\mathrm{s}_{\text {A }}$ \\
\hline
5 & product of rows 4 and 2 & $2 \mathrm{~s}_{\text {B }}$ & $-s_{C}$ & $-S_{\text {A }}$ &  &  &  \\
\hline
6 & effect on $\mathrm{s}_{\mathrm{C}}$ & $\mathrm{s}_{\mathrm{C}}$ & $\mathrm{s}_{\text {A }}$ & $\mathrm{s}_{\text {B }}$ & $\mathrm{S}_{\text {B }}$ & $\mathrm{s}_{\text {A }}$ & $\mathrm{s}_{\mathrm{C}}$ \\
\hline
7 & product of rows 6 and 2 & $2 s_{\text {C }}$ & $-\mathrm{s}_{\text {A }}$ & $-S_{B}$ &  &  &  \\
\hline
\end{tabular}
\end{center}

Application of the projection operator to a different basis function gives a different SALC in each case (rows 3, 5, and 7).

$$
\frac{1}{6}\left(2 s_{A}-s_{B}-s_{C}\right) \quad \frac{1}{6}\left(2 s_{B}-s_{C}-s_{A}\right) \quad \frac{1}{6}\left(2 s_{C}-s_{A}-s_{B}\right)
$$

However, any one of these SALCs can be expressed as a sum of the other two (the three are not 'linearly independent'). The difference of the first and second gives $\frac{1}{2}\left(s_{A}-s_{B}\right)$. This combination and the third, $\frac{1}{6}\left(2 s_{C}-s_{A}-s_{B}\right)$, are the two (now linearly independent) SALCs that are used in the construction of the molecular orbitals.

According to the discussion in Topic 9E concerning the construction of molecular orbitals of polyatomic molecules, only orbitals with the same symmetry can overlap to give a molecular orbital. In the language introduced here, this means that only SALCs of the same symmetry species have non-zero overlap and contribute to a molecular orbital. In $\mathrm{NH}_{3}$, for instance, the molecular orbitals will have the form

$$
\begin{aligned}
& \psi\left(\mathrm{a}_{1}\right)=c_{\mathrm{a} 1} \mathrm{~s}_{\mathrm{N}}+c_{\mathrm{a} 2}\left(\mathrm{~s}_{\mathrm{A}}+\mathrm{s}_{\mathrm{B}}+\mathrm{s}_{\mathrm{C}}\right) \\
& \psi\left(\mathrm{e}_{x}\right)=c_{\mathrm{e} 1} \mathrm{p}_{\mathrm{N} x}+c_{\mathrm{e} 2}\left(2 \mathrm{~s}_{\mathrm{C}}-\mathrm{s}_{\mathrm{A}}-\mathrm{s}_{\mathrm{B}}\right) \\
& \psi\left(\mathrm{e}_{y}\right)=c_{\mathrm{e} 1} \mathrm{p}_{\mathrm{N} y}+c_{\mathrm{e} 2}\left(\mathrm{~s}_{\mathrm{A}}-\mathrm{s}_{\mathrm{B}}\right)
\end{aligned}
$$

\begin{center}
\includegraphics[max width=\textwidth]{2025_02_27_d4b95d9718f0b9e2e6b9g-442}
\end{center}

Figure 10C. 4 The three H 1 s atomic orbitals in $\mathrm{NH}_{3}$ (point group $C_{3 v}$ ) can be used as a basis for forming SALCs.

Group theory is silent on the values of the coefficients: they have to be determined by one of the methods outlined in Topic 9E.

\subsection*{10C. 3 Selection rules}
The intensity of a spectral line arising from a molecular transition between some initial state with wavefunction $\psi_{\mathrm{i}}$ and a final state with wavefunction $\psi_{f}$ depends on the (electric) transition dipole moment, $\mu_{\mathrm{fi}}$ (Topic 8C). The $q$-component, where $q$ is $x, y$, or $z$, of this vector, is defined through

$$
\mu_{q, \mathrm{fi}}=-e \int \psi_{\mathrm{f}}^{\star} q \psi_{\mathrm{i}} \mathrm{~d} \tau \quad \begin{aligned}
& \text { Transition dipole moment } \\
& \text { [definition] }
\end{aligned}
$$

(10C.6)\\
where $-e$ is the charge of the electron. The transition moment has the form of the integral over $f_{1} f_{2} f_{3}$ (eqn 10C.2). Therefore, once the symmetry species of the wavefunctions and the operator are known, group theory can be used to formulate the selection rules for the transitions.

\subsection*{Example 10C. 3 Deducing a selection rule}
Is $\mathrm{p}_{y} \rightarrow \mathrm{p}_{z}$ an allowed electric dipole transition in a molecule with $C_{2 v}$ symmetry?

Collect your thoughts You need to decide whether the product $\mathrm{p}_{z} q \mathrm{p}_{y}$, with $q=x, y$, or $z$, spans $\mathrm{A}_{1}$. The symmetry species for $\mathrm{p}_{y}, \mathrm{p}_{z}$, and $q$ can be read off from the right-hand side of the character table.

The solution The $\mathrm{p}_{y}$ orbital spans $\mathrm{B}_{2}$ and $\mathrm{p}_{z}$ spans $\mathrm{A}_{1}$, so the required direct product is $\mathrm{A}_{1} \times \Gamma^{(q)} \times \mathrm{B}_{2}$, where $\Gamma^{(q)}$ is the symmetry species of $x, y$, or $z$. It does not matter in which order the direct products are calculated, so noting that $A_{1} \times B_{2}=B_{2}$ implies that $A_{1} \times \Gamma^{(q)} \times B_{2}=\Gamma^{(q)} \times B_{2}$. This direct product can be equal to $\mathrm{A}_{1}$ only if $\Gamma^{(q)}$ is $\mathrm{B}_{2}$, which is the symmetry species of $y$. Therefore, provided $q=y$ the integral may be non-zero and the transition allowed.

Comment. The analysis implies that the electromagnetic radiation involved in the transition has a component of its electric vector in the $y$-direction.

Self-test 10C. 3 Are (a) $\mathrm{p}_{x} \rightarrow \mathrm{p}_{y}$ and (b) $\mathrm{p}_{x} \rightarrow \mathrm{p}_{z}$ allowed electric dipole transitions in a molecule with $C_{2 v}$ symmetry?

$$
x=b \text { ЧІ!М ‘səૂ (q) ؛oN (e) :ıวмs }
$$

\subsection*{Checklist of concepts}
\begin{enumerate}
  \item Character tables can be used to decide whether an integral is necessarily zero.2. For an integral to be non-zero, the integrand must include a component that is a basis for the totally symmetric irreducible representation $\left(A_{1}\right)$.
  \item Only orbitals of the same symmetry species may have non-zero overlap.\\
$\square$ 4. A symmetry-adapted linear combination (SALC) is a linear combination of atomic orbitals constructed from equivalent atoms and having a specified symmetry.
\end{enumerate}

\subsection*{Checklist of equations}
\begin{center}
\begin{tabular}{llll}
\hline
Property & Equation & Comment & Equation number \\
\hline
Decomposition of a representation & $n(\Gamma)=\frac{1}{h} \sum_{C} N(C) \chi^{(\Gamma)}(C) \chi(C)$ & Real characters &  \\
Presence of $\mathrm{A}_{1}$ in a decomposition & $n\left(\mathrm{~A}_{1}\right)=\frac{1}{h} \sum_{C} N(C) \chi(C)$ & Real characters ${ }^{*}$ & 10 C .3 a \\
Overlap integral & $S=\int \psi_{2}^{\star} \psi_{1} \mathrm{~d} \tau$ & Definition & 10 C .3 b \\
Projection operator & $P^{(\Gamma)}=\frac{1}{h} \sum_{R} \chi^{(\Gamma)}(R) R$ & To generate $\psi^{(\Gamma)}=P^{(\Gamma)} \psi_{i}$ & 10 C .4 \\
Transition dipole moment & $\mu_{q, \mathrm{fi}}=-e \int \psi_{\mathrm{f}}^{*} q \psi_{\mathrm{i}} \mathrm{d} \tau$ & $q$-Component, $q=x, y$ or $z$ & 10 C .5 \\
\hline
\end{tabular}
\end{center}

\footnotetext{\begin{itemize}
  \item In general, characters may have complex values; throughout this text only real values are encountered.
\end{itemize}
}

\chapter{FOCUS 11 Molecular spectroscopy}
The origin of spectral lines in molecular spectroscopy is the absorption, emission, or scattering of a photon accompanied by a change in the energy of a molecule. The difference from atomic spectroscopy (Topic 8 C ) is that the energy of a molecule can change not only as a result of electronic transitions but also because it can undergo changes of rotational and vibrational state. Molecular spectra are therefore more complex than atomic spectra. However, they contain information relating to more properties, and their analysis leads to values of bond strengths, lengths, and angles. They also provide a way of determining a variety of molecular properties, such as dissociation energies and dipole moments.

\subsection*{11A General features of molecular spectroscopy}
This Topic begins with a discussion of the theory of absorption and emission of radiation, leading to the factors that determine the intensities and widths of spectral lines. The features of the instrumentation used to monitor the absorption, emission, and scattering of radiation spanning a wide range of frequencies are also described.\\
11A. 2 The absorption and emission of radiation; 11A. 2 Spectral linewidths; 11A. 3 Experimental techniques

\subsection*{11B Rotational spectroscopy}
This Topic shows how expressions for the values of the rotational energy levels of diatomic and polyatomic molecules are derived. The most direct procedure, which is used here, is to identify the expressions for the energy and angular momentum obtained in classical physics, and then to transform these\\
expressions into their quantum mechanical counterparts. The Topic then focuses on the interpretation of pure rotational and rotational Raman spectra, in which only the rotational state of a molecule changes. The observation that not all molecules can occupy all rotational states is shown to arise from symmetry constraints resulting from the presence of nuclear spin.\\
11B. 1 Rotational energy levels; 11B. 2 Microwave spectroscopy; 11B. 3 Rotational Raman spectroscopy; 11B. 4 Nuclear statistics and rotational states

\subsection*{11C Vibrational spectroscopy of diatomic molecules}
The harmonic oscillator (Topic 7E) is a good starting point for modelling the vibrations of diatomic molecules, but it is shown that the description of real molecules requires deviations from harmonic behaviour to be taken into account. The vibrational spectra of gaseous samples show features due to the rotational transitions that accompany the excitation of vibrations.\\
11C. 1 Vibrational motion; 11C. 2 Infrared spectroscopy;\\
11C. 3 Anharmonicity; 11C. 4 Vibration-rotation spectra;\\
11C. 5 Vibrational Raman spectra

\subsection*{11D Vibrational spectroscopy of polyatomic molecules}
The vibrational spectra of polyatomic molecules may be discussed as though they consisted of a set of independent harmonic oscillators. Their spectra can then be understood in much the same way as those of diatomic molecules.\\
11D. 1 Normal modes; 11D. 2 Infrared absorption spectra;\\
11D. 3 Vibrational Raman spectra

\subsection*{11E Symmetry analysis of vibrational}
 spectraThe atomic displacements involved in the vibrations of polyatomic molecules can be classified according to the symmetry possessed by the molecule. This classification makes it possible to decide which vibrations can be studied spectroscopically.\\
11E. 1 Classification of normal modes according to symmetry;\\
11E. 2 Symmetry of vibrational wavefunctions

\subsection*{11F Electronic spectra}
This Topic introduces the key idea that electronic transitions occur within a stationary nuclear framework. The electronic spectra of diatomic molecules are considered first, and it is seen that in the gas phase it is possible to observe simultaneous vibrational and rotational transitions that accompany the electronic transition. The general features of the electronic spectra of polyatomic molecules are also described. 11F. 1 Diatomic molecules; 11F. 2 Polyatomic molecules

\subsection*{11G Decay of excited states}
This Topic begins with an account of spontaneous emission by molecules, including the phenomena of 'fluorescence' and 'phosphorescence'. It is also seen how non-radiative decay of excited states can result in the transfer of energy as heat to the surroundings or result in molecular dissociation. The stimulated radiative decay of excited states is the key process responsible for the action of lasers.\\
11G. 1 Fluorescence and phosphorescence; 11G. 2 Dissociation and predissociation; 11G. 3 Lasers

\subsection*{Web resources What is an application of this material?}
Molecular spectroscopy is also useful to astrophysicists and environmental scientists. Impact 16 discusses how the identities of molecules found in interstellar space can be inferred from their rotational and vibrational spectra. Impact 17 focuses back on Earth and shows how the vibrational properties of its atmospheric constituents can affect its climate.

\section*{TOPIC 11A General features of molecular spectroscopy }
\begin{abstract}
Why do you need to know this material? To interpret data from the wide range of varieties of molecular spectroscopy you need to understand the experimental and theoretical features shared by them all.
\end{abstract}

\subsection*{What is the key idea?}
A transition from a low energy state to one of higher energy can be stimulated by absorption of radiation; a transition from a higher to a lower state, resulting in emission of a photon, may be either spontaneous or stimulated by radiation.

\subsection*{What do you need to know already?}
You need to be familiar with the fact that molecular energy is quantized (Topics 7E and 7F) and be aware of the concept of selection rules (Topic 8C).

In emission spectroscopy the electromagnetic radiation that arises from molecules undergoing a transition from a higher energy state to a lower energy state is detected and its frequency analysed. In absorption spectroscopy, the net absorption of radiation passing through a sample is monitored over a range of frequencies. It is necessary to specify the net absorption because not only can radiation be absorbed but it can also stimulate the emission of radiation, so the net absorption is detected. In Raman spectroscopy, the frequencies of radiation scattered by molecules is analysed to determine the changes in molecular states that accompany the scattering process. Throughout this discussion it is important to be able to express the characteristics of radiation variously as a frequency, $v$, a wavenumber, $\tilde{v}=v / c$, or a wavelength, $\lambda=c / v$, as set out in The chemist's toolkit 13 in Topic 7A.

In each case, the emission, absorption, or scattering of radiation can be interpreted in terms of individual photons. When a molecule undergoes a transition between states of energy $E_{1}$, the energy of the lower state, and $E_{\mathrm{u}}$, the energy of the upper state, the energy, $h v$, of the photon emitted or absorbed is given by the Bohr frequency condition (eqn 7A. 9 of

Topic 7A, $h v=E_{\mathrm{u}}-E_{1}$ ), where $v$ is the frequency of the radiation emitted or absorbed. Emission and absorption spectroscopy give the same information about electronic, vibrational, or rotational energy level separations, but practical considerations generally determine which technique is employed.

In Raman spectroscopy the sample is exposed to monochromatic (single frequency) radiation and therefore photons of the same energy. When the photons encounter the molecules, most are scattered elastically (without change in their energy): this process is called Rayleigh scattering. About 1 in $10^{7}$ of the photons are scattered inelastically (with different energy). In Stokes scattering the photons lose energy to the molecules and the emerging radiation has a lower frequency. In anti-Stokes scattering, a photon gains energy from a molecule and the emerging radiation has a higher frequency (Fig. 11A.1). By analysing the frequencies of the scattered radiation it is possible to gather information about the energy levels of the molecules. Raman spectroscopy is used to study molecular vibrations and rotations.\\
\includegraphics[max width=\textwidth, center]{2025_02_27_d4b95d9718f0b9e2e6b9g-451}

Figure 11A. 1 In Raman spectroscopy, incident photons are scattered from a molecule. Most photons are scattered elastically and so have the same energy as the incident photons. Some photons lose energy to the molecule and so emerge as Stokes radiation; others gain energy and so emerge as anti-Stokes radiation. The scattering can be regarded as taking place by an excitation of the molecule from its initial state to a series of excited states (represented by the shaded band), and the subsequent return to a final state. Any net energy change is either supplied from, or carried away by, the photon.

\subsection*{11A. 1 The absorption and emission of radiation}
The separation of rotational energy levels (in small molecules, $\Delta E \approx 0.01 \mathrm{zJ}$, corresponding to about $0.01 \mathrm{~kJ} \mathrm{~mol}^{-1}$ ) is smaller than that of vibrational energy levels ( $\Delta E \approx 10 \mathrm{zJ}$, corresponding to $10 \mathrm{~kJ} \mathrm{~mol}^{-1}$ ), which itself is smaller than that of electronic energy levels ( $\Delta E \approx 0.1-1 \mathrm{a}$, corresponding to about $10^{2}-10^{3} \mathrm{~kJ} \mathrm{~mol}^{-1}$ ). From the Bohr frequency condition in the form $v=\Delta E / h$, the corresponding frequencies of the photons involved in these different kinds of transitions is about $10^{10} \mathrm{~Hz}$ for rotation, $10^{13} \mathrm{~Hz}$ for vibration, and in the range $10^{14}-10^{15} \mathrm{~Hz}$ for electronic transitions. It follows that rotational, vibrational, and electronic transitions result from the absorption or emission of microwave, infrared, and ultraviolet/visible radiation, respectively.

\subsection*{(a) Stimulated and spontaneous radiative processes}
Albert Einstein identified three processes by which radiation could be either generated or absorbed by matter as a result of transitions between states. In stimulated absorption a transition from a lower energy state $l$ to a higher energy state $u$ is driven by an electromagnetic field oscillating at the frequency $v$ corresponding to the energy separation of the two states: $h v=E_{\mathrm{u}}-E_{\mathrm{l}}$. The rate of such transitions is proportional to the intensity of the incident radiation at the transition frequency: the more intense the radiation, the greater the number of photons impinging on the molecules and the greater the probability that a photon will be absorbed. The rate is also proportional to the number of molecules in the lower state, $N_{1}$, because the greater the population of that state the more likely it is that a photon will encounter a molecule in that state. The rate of stimulated absorption, $W_{u \leftarrow-1}$, can therefore be written


\begin{equation*}
W_{\mathrm{u} \leftarrow 1}=B_{\mathrm{u}, 1} N_{\mathrm{l}} \rho(v) \quad \text { Rate of stimulated absorption } \tag{11A.1a}
\end{equation*}


In this expression, $\rho(v)$ is the energy spectral density, such that $\rho(v) \mathrm{d} v$ is the energy density of radiation in the frequency range from $v$ to $v+\mathrm{d} v$. The constant $B_{\mathrm{u}, \mathrm{l}}$ is the Einstein coefficient of stimulated absorption.

Einstein also supposed that the radiation could induce the molecule in an upper state to undergo a transition to a lower state and thereby generate a photon of frequency $v$. This process is called stimulated emission, and its rate depends on the number of molecules in the upper level $N_{u}$ and the intensity of the radiation at the transition frequency. As in eqn \href{http://11A.la}{11A.la} he wrote


\begin{equation*}
W_{\mathrm{u} \rightarrow 1}^{\prime}=B_{1, \mathrm{u}} N_{\mathrm{u}} \rho(v) \tag{11A.1b}
\end{equation*}


Rate of stimulated emission

In this expression $B_{1, \mathrm{u}}$ is the Einstein coefficient of stimulated emission.

Einstein went on to suppose that a molecule could lose energy by spontaneous emission in which the molecule makes a transition to a lower state without it being driven by the presence of radiation. The rate of spontaneous emission is written

$$
W_{\mathrm{u} \rightarrow 1}^{\prime \prime}=A_{\mathrm{l}, \mathrm{u}} N_{\mathrm{u}} \quad \text { Rate of spontaneous emission } \quad \text { (11A.1c) }
$$

where $A_{1, \mathrm{u}}$ is the Einstein coefficient of spontaneous emission. When both stimulated and spontaneous emission are taken into account, the total rate of emission is


\begin{equation*}
W_{\mathrm{u} \rightarrow 1}=B_{1, \mathrm{u}} N_{\mathrm{u}} \rho(v)+A_{1, \mathrm{u}} N_{\mathrm{u}} \quad \text { Total rate of emission } \tag{11A.1d}
\end{equation*}


When the molecules and radiation are in equilibrium, the rates given in eqns 11A.1a and 11A.1d must be equal, and the populations then have their equilibrium values $N_{1}^{\mathrm{eq}}$ and $N_{\mathrm{u}}^{\mathrm{eq}}$. Therefore


\begin{equation*}
B_{\mathrm{u}, 1} N_{\mathrm{l}}^{\mathrm{eq}} \rho(v)=B_{1, \mathrm{u}} N_{\mathrm{u}}^{\mathrm{eq}} \rho(v)+A_{1, \mathrm{u}} N_{\mathrm{u}}^{\mathrm{eq}} \tag{11A.2a}
\end{equation*}


and so


\begin{equation*}
\rho(v)=\frac{A_{\mathrm{l}, \mathrm{u}} / B_{\mathrm{u}, \mathrm{l}}}{N_{\mathrm{l}}^{\mathrm{eq}} / N_{\mathrm{u}}^{\mathrm{eq}}-B_{1, \mathrm{l}} / B_{\mathrm{u}, \mathrm{l}}} \tag{11A.2b}
\end{equation*}


However, the ratio of the equilibrium populations must be in accord with the Boltzmann distribution (as specified in the Prologue of this text and Topic 13A):


\begin{equation*}
\frac{N_{\mathrm{u}}^{\mathrm{eq}}}{N_{\mathrm{l}}^{\mathrm{eq}}}=\mathrm{e}^{-\left(E_{\mathrm{u}}-E_{\mathrm{l}}\right) / k T}=\mathrm{e}^{-h \nu / k T} \tag{11A.3}
\end{equation*}


and therefore, at equilibrium,


\begin{equation*}
\rho(v)=\frac{A_{\mathrm{l}, \mathrm{u}} / B_{\mathrm{u}, \mathrm{l}}}{\mathrm{e}^{h \mathrm{v} / k T}-B_{\mathrm{l}, \mathrm{u}} / B_{\mathrm{u}, \mathrm{l}}} \tag{11A.4}
\end{equation*}


Moreover, at equilibrium the radiation density is given by the Planck distribution of radiation in equilibrium with a black body (eqn 7A.6b, Topic 7A):


\begin{equation*}
\rho(v)=\frac{8 \pi h v^{3} / c^{3}}{\mathrm{e}^{h \nu / k T}-1} \tag{11A.5}
\end{equation*}


Planck distribution

It then follows that


\begin{equation*}
A_{\mathrm{l}, \mathrm{u}}=\left(\frac{8 \pi h v^{3}}{c^{3}}\right) B_{\mathrm{l}, \mathrm{u}} \quad \text { and } \quad B_{\mathrm{l}, \mathrm{u}}=B_{\mathrm{u}, \mathrm{l}} \tag{11A.6a}
\end{equation*}


Although these relations have been derived on the assumption that the molecules and radiation are in equilibrium, they are properties of the molecules themselves and are independent of the spectral distribution of the radiation (that is, whether it is black-body or not) and can be used in eqn 11A. 1 for any energy densities.

The ratio of the rate of spontaneous to stimulated emission can be found by combining eqns 11A.1b, 11A.1c, and 11A.6a to give


\begin{equation*}
\frac{W_{\mathrm{u} \rightarrow 1}^{\prime \prime}}{W_{\mathrm{u} \rightarrow 1}^{\prime}}=\frac{A_{\mathrm{l}, \mathrm{u}}}{B_{\mathrm{l}, \mathrm{u}} \rho(v)}=\frac{8 \pi h v^{3}}{c^{3} \rho(v)} \tag{11A.6b}
\end{equation*}


This relation shows that, for a given spectral density, the relative importance of spontaneous emission increases as the cube of the transition frequency and therefore that spontaneous emission is most likely to be of importance at high frequencies. Conversely, spontaneous emission can be ignored at low frequencies, in which case the intensities of such transitions can be discussed in terms of stimulated emission and absorption alone.

\subsection*{Brief illustration 11A. 1}
On going from infrared to visible radiation, the frequency increases by a factor of about 100 , so the ratio of the rates of spontaneous to stimulated emission increases by a factor of $10^{6}$ for the same spectral density. This strong increase accounts for the observation that whereas electronic transitions are often monitored by emission spectroscopy, vibrational spectroscopy is an absorption technique and spontaneous (but not stimulated) emission is negligible.

\subsection*{(b) Selection rules and transition moments}
A 'selection rule' is a statement about whether a transition is forbidden or allowed (Topic 8C). The underlying idea is that, for the molecule to be able to interact with the electromagnetic field and absorb or create a photon of frequency $v$, it must possess, at least transiently, an electric dipole oscillating at that frequency. This transient dipole is expressed quantum mechanically in terms of the transition dipole moment, $\mu_{\mathrm{f}}$, between the initial and final states with wavefunctions $\psi_{\mathrm{i}}$ and $\psi_{f}$ :

\[
\boldsymbol{\mu}_{\mathrm{fi}}=\int \psi_{\mathrm{f}}^{\star} \hat{\mu} \psi_{\mathrm{i}} \mathrm{~d} \tau \quad \begin{align*}
& \text { Transition dipole moment }  \tag{11A.7}\\
& \text { [definition] }
\end{align*}
\]

where $\hat{\mu}$ is the electric dipole moment operator. The magnitude of the transition dipole moment can be regarded as a measure of the charge redistribution that accompanies a transition and a transition is active (and generates or absorbs a photon) only if the accompanying charge redistribution is dipolar (Fig. 11A.2). It follows that, to identify the selection rules, the conditions for which $\mu_{\mathrm{fi}} \neq 0$ must be established.

A gross selection rule specifies the general features that a molecule must have if it is to have a spectrum of a given kind. For instance, in Topic 11B it is shown that a molecule gives a rotational spectrum only if it has a permanent electric dipole\\
\includegraphics[max width=\textwidth, center]{2025_02_27_d4b95d9718f0b9e2e6b9g-453}

Figure 11A. 2 (a) When a 1 s electron becomes a 2 s electron, there is a spherical migration of charge. There is no dipole moment associated with this migration of charge, so this transition is electric-dipole forbidden. (b) In contrast, when a 1 s electron becomes a 2 p electron, there is a dipole associated with the charge migration; this transition is allowed.\\
moment. This rule, and others like it for other types of transition, is explained in the relevant Topic. A detailed study of the transition moment leads to the specific selection rules which express the allowed transitions in terms of the changes in quantum numbers or various symmetry features of the molecule.

\subsection*{(c) The Beer-Lambert law}
It is found empirically that when electromagnetic radiation passes through a sample of length $L$ and molar concentration [J] of the absorbing species J, the incident and transmitted intensities, $I_{0}$ and $I$, are related by the Beer-Lambert law:


\begin{equation*}
I=I_{0} 10^{-\varepsilon[]] L} \quad \text { Beer-Lambert law } \tag{11A.8}
\end{equation*}


The quantity $\varepsilon$ (epsilon) is called the molar absorption coefficient (formerly, and still widely, the 'extinction coefficient'); it depends on the frequency (or wavenumber and wavelength) of the incident radiation and is greatest where the absorption is most intense. The dimensions of $\varepsilon$ are $1 /($ concentration $\times$ length), and it is normally convenient to express it in cubic decimetres per mole per centimetre $\left(\mathrm{dm}^{3} \mathrm{~mol}^{-1} \mathrm{~cm}^{-1}\right)$; in SI base units it is expressed in metre squared per mole $\left(\mathrm{m}^{2} \mathrm{~mol}^{-1}\right)$. The latter units imply that $\varepsilon$ may be regarded as a (molar) crosssection for absorption, and that the greater the cross-sectional area of the molecule for absorption, the greater is its ability to block the passage of the incident radiation at a given frequency. The Beer-Lambert law is an empirical result. However, its form can be derived on the basis of a simple model.

How is that done? 11A. 1 Justifying the Beer-Lambert law\\
You need to imagine the sample as consisting of a stack of infinitesimal slices, like sliced bread (Fig. 11A.3). The thickness of each layer is $\mathrm{d} x$.\\
\includegraphics[max width=\textwidth, center]{2025_02_27_d4b95d9718f0b9e2e6b9g-454}

Figure 11A. 3 To establish the Beer-Lambert law, the sample is supposed to be divided into a large number of thin slices. The reduction in intensity caused by one slice is proportional to the intensity incident on it (after passing through the preceding slices), the thickness of the slice, and the concentration of the absorbing species.

Step 1 Calculate the change in intensity due to passage through one slice

The change in intensity, $\mathrm{d} I$, that occurs when electromagnetic radiation passes through one particular slice is proportional to the thickness of the slice, the molar concentration of the absorber J, and (because the absorption is stimulated) the intensity of the incident radiation at that slice of the sample, so $\mathrm{d} I \propto[J] I \mathrm{~d} x$. The intensity is reduced by absorption, which means that $\mathrm{d} I$ is negative and can therefore be written

$$
\mathrm{d} I=-\kappa[J] I \mathrm{~d} x
$$

where $\kappa$ (kappa) is the proportionality coefficient. Division of both sides by I gives

$$
\frac{\mathrm{d} I}{I}=-\kappa[J] \mathrm{d} x
$$

This expression applies to each successive slice.\\
Step 2 Evaluate the total change in intensity due to passage through successive slices\\
To obtain the intensity that emerges from a sample of thickness $L$ when the intensity incident on one face of the sample is $I_{0}$, you need the sum of all the successive changes. Assume that the molar concentration of the absorbing species is uniform and may be treated as a constant. Because a sum over infinitesimally small increments is an integral, it follows that:\\
\includegraphics[max width=\textwidth, center]{2025_02_27_d4b95d9718f0b9e2e6b9g-454(1)}

Therefore

$$
\ln \frac{I}{I_{0}}=-\kappa[J] L
$$

Now express the natural logarithm as a common logarithm (to base 10 ) by using $\ln x=(\ln 10) \log x$, and a new constant $\varepsilon$ defined as $\varepsilon=\kappa / \ln 10$ to give

$$
\log \frac{I}{I_{0}}=-\varepsilon[J] L
$$

Raising each side as a power of 10 gives the Beer-Lambert law (eqn 11A.8).

The spectral characteristics of a sample are commonly reported as the transmittance, $T$, of the sample at a given frequency:


\begin{equation*}
T=\frac{I}{I_{0}} \tag{11A.9a}
\end{equation*}


Transmittance [definition]\\
or its absorbance, $A$ :

\[
A=\log \frac{I_{0}}{I} \quad \begin{align*}
& \text { Absorbance }  \tag{11A.9b}\\
& \text { [definition] }
\end{align*}
\]

The two quantities are related by $A=-\log T$ (note the common logarithm) and the Beer-Lambert law becomes


\begin{equation*}
A=\varepsilon[J] L \tag{11A.9c}
\end{equation*}


The product $\varepsilon[J] L$ was known formerly as the optical density of the sample.

\subsection*{Example 11A. 1 Determining a molar absorption coefficient}
Radiation of wavelength 280 nm passed through 1.0 mm of a solution that contained an aqueous solution of the amino acid tryptophan at a molar concentration of $0.50 \mathrm{mmoldm}^{-3}$. The intensity is reduced to 54 per cent of its initial value (so $T=0.54$ ). Calculate the absorbance and the molar absorption coefficient of tryptophan at 280 nm . What would be the transmittance through a cell of thickness 2.0 mm ?

Collect your thoughts From $A=-\log T=\varepsilon[J] L$, it follows that $\varepsilon=A /[J] L$. For the transmittance through the thicker cell, you need to calculate the absorbance by using $A=-\log T=\varepsilon[J] L$ and the computed value of $\varepsilon$; the transmittance is $T=10^{-A}$.

The solution The absorbance is $A=-\log 0.54=0.27$, and so the molar absorption coefficient is

$$
\varepsilon=\frac{-\log 0.54}{\left(5.0 \times 10^{-4} \mathrm{~mol} \mathrm{dm}^{-3}\right) \times(1.0 \mathrm{~mm})}=5.4 \times 10^{2} \mathrm{dm}^{3} \mathrm{~mol}^{-1} \mathrm{~mm}^{-1}
$$

These units are convenient for the rest of the calculation (but the outcome could be reported as $5.4 \times 10^{3} \mathrm{dm}^{3} \mathrm{~mol}^{-1} \mathrm{~cm}^{-1}$ if\\
desired or even as $\left.5.4 \times 10^{2} \mathrm{~m}^{2} \mathrm{~mol}^{-1}\right)$. The absorbance of a sample of length 2.0 mm is

$$
\begin{aligned}
A= & \left(5.4 \times 10^{2} \mathrm{dm}^{3} \mathrm{~mol}^{-1} \mathrm{~mm}^{-1}\right) \times\left(5.0 \times 10^{-4} \mathrm{~mol} \mathrm{dm}^{-3}\right) \\
& \times(2.0 \mathrm{~mm})=0.54
\end{aligned}
$$

The transmittance is now $T=10^{-A}=10^{-0.54}=0.29$.\\
Self-test 11A. 1 The transmittance of an aqueous solution containing the amino acid tyrosine at a molar concentration of $0.10 \mathrm{mmoldm}^{-3}$ was measured as 0.14 at 240 nm in a cell of length 5.0 mm . Calculate the absorbance of the solution and the molar absorption coefficient of tyrosine at that wavelength. What would be the transmittance through a cell of length 1.0 mm ?\\
\includegraphics[max width=\textwidth, center]{2025_02_27_d4b95d9718f0b9e2e6b9g-455}

The maximum value of the molar absorption coefficient, $\varepsilon_{\max }$, is an indication of the intensity of a transition. However, because absorption bands generally spread over a range of wavenumbers, quoting the absorption coefficient at a single wavenumber might not give a true indication of the intensity of a transition. The integrated absorption coefficient, $\mathcal{A}$, is the sum of the absorption coefficients over the entire band (Fig. 11A.4), and corresponds to the area under the plot of the molar absorption coefficient against wavenumber:

\[
\mathcal{A}=\int_{\text {band }} \varepsilon(\tilde{v}) \mathrm{d} \tilde{v} \quad \begin{align*}
& \text { Integrated absorption coefficient }  \tag{11A.10}\\
& \text { [definition] }
\end{align*}
\]

For bands of similar widths, the integrated absorption coefficients are proportional to the heights of the bands. Equation 11A. 10 also applies to the individual lines that contribute to a band: a spectroscopic line is not a geometrically thin line, but has a width.\\
\includegraphics[max width=\textwidth, center]{2025_02_27_d4b95d9718f0b9e2e6b9g-455(1)}

Figure 11A. 4 The integrated absorption coefficient of a transition is the area under a plot of the molar absorption coefficient against the wavenumber of the incident radiation.

\subsection*{11A.2 Spectral linewidths}
A number of effects contribute to the widths of spectroscopic lines. The design of the spectrometer itself affects the linewidth, and there are other contributions that arise from physical processes in the sample. Some of the latter can be minimized by altering the conditions, while others are intrinsic to the molecules and cannot be altered.

\subsection*{(a) Doppler broadening}
One important broadening process in gaseous samples is the Doppler effect, in which radiation is shifted in frequency when its source is moving towards or away from the observer. When a molecule emitting electromagnetic radiation of frequency $v$ moves with a speed $s$ relative to an observer, the observer detects radiation of frequency

$$
\begin{array}{r}
v_{\text {receding }}=\left(\frac{1-s / c}{1+s / c}\right)^{1 / 2} v_{0} \quad v_{\text {approaching }}=\left(\frac{1+s / c}{1-s / c}\right)^{1 / 2} v_{0} \\
\\
\text { Doppler shifts }
\end{array}
$$

(11A.11a)\\
where $c$ is the speed of light. For nonrelativistic speeds $(s \ll c)$, these expressions simplify to


\begin{equation*}
v_{\text {receding }} \approx(1-s / c) v_{0} \quad v_{\text {approaching }} \approx(1+s / c) v_{0} \tag{11A.11b}
\end{equation*}


Atoms and molecules reach high speeds in all directions in a gas, and a stationary observer detects the corresponding Doppler-shifted range of frequencies. Some molecules approach the observer, some move away; some move quickly, others slowly. The detected spectral 'line' is the absorption or emission profile arising from all the resulting Doppler shifts. The challenge is to relate the observed linewidth to the spread of speeds in the gas, and in turn to see how that spread depends on the temperature.

\subsection*{How is that done? 11A. 2 Deriving an expression for Doppler broadening}
You need to relate the spread of Doppler shifts to the distribution of molecular kinetic energy as expressed by the Boltzmann distribution.

Step 1 Establish the relation between the observed frequency and the molecular speed\\
It follows from the Boltzmann distribution (see the Prologue to the text) that the probability that an atom or molecule of mass $m$ and speed $s$ in a gas phase sample at a temperature $T$\\
has kinetic energy $E_{\mathrm{k}}=\frac{1}{2} m s^{2}$ is proportional to $\mathrm{e}^{-m s^{2} / 2 k T}$. When $s \ll c$, the Doppler shifts for receding and approaching molecules are given by the expressions in eqn 11A.11b. It follows that the shift between the observed frequency and the true frequency is $v_{\text {obs }}-v_{0} \approx \pm v_{0} s / c$. This expression can be rearranged to give

$$
s= \pm c\left(v_{\mathrm{obs}}-v_{0}\right) / v_{0}
$$

Step 2 Evaluate the distribution of frequencies arising from a distribution of speeds\\
The intensity $I$ of a transition at $v_{\text {obs }}$ is proportional to the probability of there being a molecule that emits or absorbs at $v_{\text {obs }}$. Such a molecule would have a speed given by the above expression, so it follows from the Boltzmann distribution that

$$
I\left(v_{\text {obs }}\right) \propto \mathrm{e}^{-m s^{2} / 2 k T}=\mathrm{e}^{-m c^{2}\left(v_{\text {obs }}-v_{0}\right)^{2} / 2 \nu_{0}^{2} k T}
$$

which has the form of a Gaussian function. The width at halfheight, $\delta v_{\text {obs }}$, can be inferred directly from the general form of such a function (as specified in The chemist's toolkit 26):

$$
\left.-\delta v_{\mathrm{obs}}=\frac{2 v_{0}}{c}\left(\frac{2 k T \ln 2}{m}\right)^{1 / 2} \right\rvert\, \begin{array}{r}
(11 \mathrm{~A} .12 \mathrm{a}) \\
\end{array}
$$

Doppler broadening increases with temperature (Fig. 11A.5) because the molecules then acquire a wider range of speeds. Conversely, reducing the temperature results in narrower lines. Note that the Doppler linewidth is proportional to the frequency, so Doppler broadening becomes more important as higher frequencies are observed.\\
\includegraphics[max width=\textwidth, center]{2025_02_27_d4b95d9718f0b9e2e6b9g-456}

Figure 11A. 5 The Gaussian shape of a Doppler-broadened spectral line reflects the Boltzmann distribution of translational kinetic energies in the sample at the temperature of the experiment. The line broadens as the temperature is increased.

The chemist's toolkit 26\\
Exponential and Gaussian functions

An exponential function is a function of the form

$$
f(x)=a e^{-b x}
$$

Exponential function\\
This function has the value $a$ at $x=0$ and decays toward zero as $x \rightarrow \infty$. This decay is faster when $b$ is large than when it is small. The function rises rapidly to infinity as $x \rightarrow-\infty$. See Sketch 1 .\\
\includegraphics[max width=\textwidth, center]{2025_02_27_d4b95d9718f0b9e2e6b9g-456(1)}

The general form of a Gaussian function is

$$
f(x)=a \mathrm{e}^{-(x-b)^{2} / 2 \sigma^{2}}
$$

Gaussian function\\
The graph of this function is a symmetrical bell-shaped curve centred on $x=b$; the function has is maximum values of $a$ at its centre. The width of the function, measured at half its height, is $\delta x=2 \sigma(2 \ln 2)^{1 / 2}$; the greater $\sigma$, the greater is the width at halfheight. Sketch 1 also shows a Gaussian function with $b=0$.

\subsection*{Brief illustration 11A. 2}
For a molecule such as CO at $T=300 \mathrm{~K}$, and noting that $1 \mathrm{~J}=$ $1 \mathrm{~kg} \mathrm{~m}^{2} \mathrm{~s}^{-2}$,

$$
\begin{aligned}
\frac{\delta v_{\text {obs }}}{v_{0}}= & \frac{2}{c}\left(\frac{2 k T \ln 2}{m_{\mathrm{CO}}}\right)^{1 / 2} \\
= & \frac{2}{2.998 \times 10^{8} \mathrm{~m} \mathrm{~s}^{-1}} \\
& \times\left(\frac{2 \times\left(1.380 \times 10^{-23} \mathrm{~J} \mathrm{~K}^{-1}\right) \times(300 \mathrm{~K}) \times \ln 2}{4.651 \times 10^{-26} \mathrm{~kg}}\right)^{1 / 2} \\
= & 2.34 \times 10^{-6}
\end{aligned}
$$

For a transition wavenumber of $2150 \mathrm{~cm}^{-1}$ from the infrared spectrum of CO , corresponding to a frequency of 64.4 THz $\left(1 \mathrm{THz}=10^{12} \mathrm{~Hz}\right)$, the linewidth is 151 MHz or $5.0 \times 10^{-3} \mathrm{~cm}^{-1}$.

The linewidth due to Doppler broadening can also be expressed in terms of wavelength as

$$
\delta \lambda_{\mathrm{obs}}=\frac{2 \lambda_{0}}{c}\left(\frac{2 k T \ln 2}{m}\right)^{1 / 2} \quad \text { Doppler broadening }
$$

(11A.12b)

\subsection*{(b) Lifetime broadening}
At any instant a molecule exists in a specific state but it does not remain in that state indefinitely. For example, the molecule might collide with another and in the process change its state. Alternatively, the molecule may fall to a lower level and emit a photon by the process of spontaneous emission. How long the state persists depends on the rates of these processes and is characterized by a lifetime, $\tau$. When the Schrödinger equation is analysed for a state that persists for a time $\tau$, it turns out that the energy is uncertain to an extent $\delta E \approx \hbar / \tau$. Therefore, spectroscopic transitions that involve this state have a linewidth of the order of $\delta E / h=1 / 2 \pi \tau$; the shorter the lifetime, the broader is the line. This process is called lifetime broadening. When the lifetime is limited by the process of spontaneous emission rather than external causes such as collisions, the resulting linewidth is called the natural linewidth.

\subsection*{Brief illustration 11A. 3}
Excited electronic states of molecules often have short lifetimes due to the high rate of spontaneous emission. A typical lifetime might be 10 ns , which would lead to a natural linewidth of $1 /\left(2 \pi \times 10 \times 10^{-9} \mathrm{~s}\right)=16 \mathrm{MHz}$ or $5.3 \times 10^{-4} \mathrm{~cm}^{-1}$. As can be inferred from the previous Brief illustration, the Doppler linewidth is typically much greater than the natural linewidth.

Collisions between molecules are generally efficient at changing their rotational or vibrational energies, so a good estimate of the resulting collisional lifetime, $\tau_{\text {col }}$, is to equate it to $1 / z$, where $z$ is the collision frequency (Topic 1B). If it is assumed that each collision results in a change of rotational or vibrational state, the lifetime of a state can be taken as $\tau_{\text {col }}$, and hence the resulting broadening is $\delta E / h=1 / 2 \pi \tau_{\text {col }}=z / 2 \pi$; this contribution to the linewidth is often referred to as collisional line broadening. The collision frequency for two molecules with masses $m_{\mathrm{A}}$ and $m_{\mathrm{B}}$ is given by eqn 1B.12b as

$$
z=\frac{\sigma v_{\mathrm{rel}} p}{k T} \text { with } v_{\mathrm{rel}}=\left(\frac{8 k T}{\pi \mu}\right)^{1 / 2} \text { and } \mu=\frac{m_{\mathrm{A}} m_{\mathrm{B}}}{m_{\mathrm{A}}+m_{\mathrm{B}}}
$$

Note that the collision frequency, and hence the linewidth, is proportional to the pressure, which is why the broadening due

\footnotetext{${ }^{1}$ See our Molecular quantum mechanics (2011) for a discussion of the origin of this relation.
}
to collisions is sometimes referred to as pressure broadening. This contribution to the linewidth can be minimized by lowering the pressure as much as possible, although doing so decreases the intensity of the absorption because there are fewer molecules to absorb the radiation. In contrast to Doppler broadening, pressure broadening is independent of the transition frequency.

\subsection*{Brief illustration 11A. 4}
The linewidth due to pressure broadening in methane gas at 1 bar and 298 K can be estimated by using the expressions just quoted; the collision cross-section $\sigma$ is $0.46 \mathrm{~nm}^{2}$. Taking $m_{\mathrm{A}}$ and $m_{\mathrm{B}}$ both to be the mass of a methane molecule, $v_{\text {rel }}$ is $888 \mathrm{~m} \mathrm{~s}^{-1}$. Hence

$$
\begin{aligned}
z & =\frac{\sigma v_{\text {rel }} p}{k T} \\
& =\frac{\left(0.46 \times 10^{-18} \mathrm{~m}^{2}\right) \times\left(888 \mathrm{~m} \mathrm{~s}^{-1}\right) \times\left(1 \times 10^{5} \mathrm{~N} \mathrm{~m}^{-2}\right)}{\left(1.381 \times 10^{-23} \mathrm{~J} \mathrm{~K}^{-1}\right) \times(300 \mathrm{~K})} \\
& =9.9 \times 10^{9} \mathrm{~s}^{-1}
\end{aligned}
$$

The linewidth is therefore $z / 2 \pi=1.6 \mathrm{GHz}$ or $0.053 \mathrm{~cm}^{-1}$. In Brief illustration 11A. 2 the Doppler linewidth for a transition in the infrared is estimated as 150 MHz , which is much less than the pressure broadening estimated for the present set of conditions. As the frequency is raised, the Doppler broadening increases in proportion, but the pressure broadening remains unchanged, so Doppler broadening might become dominant.

\subsection*{11A.3 Experimental techniques}
Common to all spectroscopic techniques is a spectrometer, an instrument used to detect the characteristics of radiation scattered, emitted, or absorbed by atoms and molecules. Figure 11A. 6 shows the general layout of an absorption spectrometer. Radiation from an appropriate source is directed\\
\includegraphics[max width=\textwidth, center]{2025_02_27_d4b95d9718f0b9e2e6b9g-457}

Figure 11A. 6 The layout of a typical absorption spectrometer. Radiation from the source passes through the sample; the radiation is then dispersed according to frequency and the intensity at each frequency is measured with a detector.\\
towards a sample and the radiation transmitted strikes a device that separates it into different frequencies or wavelengths. The intensity of radiation at each frequency is then analysed by a suitable detector. It is usual to record one spectrum with the sample in place, and one with the sample removed (the 'background spectrum'): the difference between these two spectra eliminates any absorption not due to the sample itself.

\subsection*{(a) Sources of radiation}
Sources of radiation are either monochromatic, those spanning a very narrow range of frequencies around a central value, or polychromatic, those spanning a wide range of frequencies. In the microwave region frequency synthesizers and various solid state devices can be used to generate monochromatic radiation that can be tuned over a wide range of frequencies. Certain kinds of lasers and light-emitting diodes are often used to provide monochromatic radiation from the infrared to the ultraviolet region. Polychromatic black-body radiation from hot materials (Topic 7A) can be used over the same range. Examples include mercury arcs inside a quartz envelope (usable in the range $35-200 \mathrm{~cm}^{-1}$ ), Nernst filaments and globars (200-4000 $\mathrm{cm}^{-1}$ ), and quartz-tungsten-halogen lamps (320-2500 nm).

A gas discharge lamp is a common source of ultraviolet and visible radiation. In a xenon discharge lamp, an electrical discharge excites xenon atoms to excited states, which then emit ultraviolet radiation. In a deuterium lamp, excited $\mathrm{D}_{2}$ molecules dissociate into electronically excited D atoms that emit intense radiation in the range $200-400 \mathrm{~nm}$.

For certain applications, radiation is generated in a synchrotron storage ring, which consists of an electron beam travelling in a circular path with circumferences of up to several hundred metres. As electrons travelling in a circle are constantly accelerated by the forces that constrain them to their path, they generate radiation (Fig. 11A.7). This 'synchrotron\\
\includegraphics[max width=\textwidth, center]{2025_02_27_d4b95d9718f0b9e2e6b9g-458(1)}

Figure 11A. 7 A simple synchrotron storage ring. The electrons injected into the ring from the linear accelerator and booster synchrotron are accelerated to high speed in the main ring. An electron in a curved path is subject to constant acceleration, and an accelerated charge radiates electromagnetic energy.\\
radiation' spans a wide range of frequencies, including infrared radiation and X-rays. Except in the microwave region, synchrotron radiation is much more intense than can be obtained by most conventional sources.

\subsection*{(b) Spectral analysis}
A common device for the analysis of the frequencies, wavenumbers, or wavelengths in a beam of radiation is a diffraction grating, which consists of a glass or ceramic plate into which fine grooves have been cut and covered with a reflective aluminium coating. For work in the visible region of the spectrum, the grooves are cut about 1000 nm apart (a spacing comparable to the wavelength of visible light). The grating causes interference between waves reflected from its surface, and constructive interference occurs at specific angles that depend on the wavelength of the radiation being used. Thus, each wavelength of light is diffracted into a specific direction (Fig. 11A.8). In a monochromator, a narrow exit slit allows only a narrow range of wavelengths to reach the detector. Turning the grating on an axis perpendicular to the incident and diffracted beams allows different wavelengths to be analysed; in this way, the absorption spectrum is built up one narrow wavelength range at a time. In a polychromator, there is no slit and a broad range of wavelengths can be analysed simultaneously by array detectors, such as those discussed below.

Currently, almost all spectrometers operating in the infrared and near-infrared use 'Fourier transform' techniques for spectral detection and analysis. (Fourier transforms are discussed, but in more detail than needed here, in The chemist's toolkit 28 in Topic 12C.) The heart of a Fourier transform spectrometer is a Michelson interferometer, a device for analysing the wavelengths present in a composite signal. The Michelson interferometer works by splitting the beam from the sample\\
\includegraphics[max width=\textwidth, center]{2025_02_27_d4b95d9718f0b9e2e6b9g-458}

Figure 11A. 8 A polychromatic beam is dispersed by a diffraction grating into three component wavelengths $\lambda_{1}, \lambda_{2}$, and $\lambda_{3}$. In the configuration shown, only radiation with $\lambda_{2}$ passes through a narrow slit and reaches the detector. Rotation of the diffraction grating (as shown by the arrows on the dotted circle) allows $\lambda_{1}$ or $\lambda_{3}$ to reach the detector.\\
\includegraphics[max width=\textwidth, center]{2025_02_27_d4b95d9718f0b9e2e6b9g-459}

Figure 11A. 9 A Michelson interferometer. The beam-splitting element divides the incident beam into two beams with a path difference that depends on the location of the movable mirror $M_{1}$. The compensator ensures that both beams pass through the same thickness of material. The beams have been kept separate and coloured to distinguish them: they do not indicate different wavelengths.\\
into two and arranging for them to take different routes through the instrument before eventually recombining at the detector (Fig. 11A.9). One beam is reflected from mirror $M_{1}$ and one from mirror $M_{2}$; by moving $M_{1}$ it is therefore possible to introduce a difference in the length of the path traversed by the two beams.

Consider first the simplest case in which a beam of monochromatic light of wavelength $\lambda$ is passed into the interferometer. If the path length difference $p$ is 0 , the two beams interfere constructively; the same is true if $p$ is an integer number of wavelengths: $\lambda, 2 \lambda, 3 \lambda, \ldots$. If $p$ is one half of a wavelength, the two beams interfere destructively and cancel; the same is true if $p$ is an odd multiple of half-wavelengths: $\lambda / 2,3 \lambda / 2,5 \lambda / 2, \ldots$. Therefore, as the mirror $\mathrm{M}_{1}$ is moved the detected signal goes through a series of peaks and troughs depending on whether the two beams interfere constructively or destructively, and the net signal varies as $1+\cos (2 \pi p / \lambda)$, or $1+\cos (2 \pi p \tilde{v})$ (Fig. 11A.10).\\
\includegraphics[max width=\textwidth, center]{2025_02_27_d4b95d9718f0b9e2e6b9g-459(1)}

Figure 11A. 10 An interferogram produced as the path length difference $p$ is changed in the interferometer shown in Fig. 11A.9. Only a single wavelength component is present in the signal, so the graph is a plot of $1+\cos (2 \pi \widetilde{v} p)$.

In a spectroscopic observation a mixture of radiation of different wavelengths and intensities is passed into the spectrometer. Each component gives rise an interference pattern proportional to $1+\cos (2 \pi p \tilde{v})$, and the signal recorded by the detector is their sum. Thus, if the intensity of the radiation entering the spectrometer consists of a mixture of wavenumber $\tilde{v}_{i}$ with intensities $I\left(\tilde{v}_{i}\right)$, the signal measured at the detector is given by the sum


\begin{equation*}
\tilde{I}(p)=\sum_{i} I\left(\tilde{v}_{i}\right)\left\{1+\cos \left(2 \pi \tilde{v}_{i} p\right)\right\} \tag{11A.13}
\end{equation*}


A plot of $\tilde{I}(p)$ against $p$, which is detected by the system and recorded, is called an interferogram. The problem is to find $I(\tilde{v})$, the variation of the intensity with wavenumber, which is the spectrum, from $\tilde{I}(p)$. This conversion can be carried out by using a standard mathematical technique, called the Fourier transform, which involves evaluating the integral

$$
I(\tilde{v})=4 \int_{0}^{\infty}\left\{\tilde{I}(p)-\frac{1}{2} \tilde{I}(0)\right\} \cos (2 \pi \tilde{v} p) \mathrm{d} p
$$

In practice, the measured values of $\tilde{I}(p)$ are digitized, stored in a computer attached to the spectrometer, and then the Fourier transform is computed numerically.

\subsection*{Example 11A. 2}
Relating a spectrum to an interferogram\\
Suppose the light entering the interferometer consists of three components with the following characteristics:

\begin{center}
\begin{tabular}{llll}
$\tilde{v}_{i} / \mathrm{cm}^{-1}$ & 150 & 250 & 450 \\
$I\left(\tilde{v}_{i}\right)$ & 1 & 3 & 6 \\
\end{tabular}
\end{center}

where the intensities are relative to the first value listed. Plot the interferogram associated with this signal. Then calculate and plot the Fourier transform of the interferogram.\\
Collect your thoughts For a signal consisting of just these three component beams, you can use eqn 11A. 13 directly. Although in this case (where $\tilde{I}(p)$ is simply the sum of trigonometric functions) the Fourier transform $I(\tilde{v})$ can be carried out exactly, in general it is best done numerically by using mathematical software.

The solution From the data, the interferogram is

$$
\begin{aligned}
\tilde{I}(p) & =\left(1+\cos 2 \pi \tilde{v}_{1} p\right)+3 \times\left(1+\cos 2 \pi \tilde{v}_{2} p\right)+6 \times\left(1+\cos 2 \pi \tilde{v}_{3} p\right) \\
& =10+\cos 2 \pi \tilde{v}_{1} p+3 \cos 2 \pi \tilde{v}_{2} p+6 \cos 2 \pi \tilde{v}_{3} p
\end{aligned}
$$

This function is plotted in Fig. 11A.11. The result of evaluating the Fourier transform numerically is shown in Fig. 11A.12.\\
\includegraphics[max width=\textwidth, center]{2025_02_27_d4b95d9718f0b9e2e6b9g-460(2)}

Figure 11A. 11 The interferogram calculated from data in Example 11A.2.\\
\includegraphics[max width=\textwidth, center]{2025_02_27_d4b95d9718f0b9e2e6b9g-460(3)}

Figure 11A. 12 The Fourier transform of the interferogram shown in Fig. 11A.11. The oscillations arise from the way that the signal in Fig. 11A. 11 is sampled. As the sampling is extended to greater path-length differences, the oscillations disappear and the peaks become sharper.

Self-test 11A.2 Explore the effect of varying the wavenumbers of the three components of the radiation on the shape of the interferogram by changing the value of $\tilde{v}_{3}$ to $550 \mathrm{~cm}^{-1}$.\\
\includegraphics[max width=\textwidth, center]{2025_02_27_d4b95d9718f0b9e2e6b9g-460}\\
\includegraphics[max width=\textwidth, center]{2025_02_27_d4b95d9718f0b9e2e6b9g-460(1)}

Path length difference, $p$\\
Figure 11A. 13 The interferogram calculated from the data in Self-test 11A. 2 superimposed on the interferogram obtained in the Example itself (in pale blue).

\subsection*{(c) Detectors}
A detector is a device that converts radiation into an electric signal for processing and display. Detectors may consist of a single radiation sensing element or of several small elements arranged in one- or two-dimensional arrays.

A microwave detector is typically a crystal diode consisting of a tungsten tip in contact with a semiconductor. The most common detectors found in commercial infrared spectrometers are sensitive in the mid-infrared region. In a photovoltaic device the potential difference changes upon exposure to infrared radiation. In a pyroelectric device the capacitance is sensitive to temperature and hence to the presence of infrared radiation.

A common detector for work in the ultraviolet and visible ranges is a photomultiplier tube (PMT), in which the photoelectric effect (Topic 7A) is used to generate an electrical signal proportional to the intensity of light that strikes the detector. A common, but less sensitive, alternative to the PMT is a photodiode, a solid-state device that conducts electricity when struck by photons because light-induced electron transfer reactions in the detector material create mobile charge carriers (negatively charged electrons and positively charged 'holes').

A charge-coupled device (CCD) is a two-dimensional array of several million small photodiode detectors. With a CCD, a wide range of wavelengths that emerge from a polychromator are detected simultaneously, thus eliminating the need to measure the radiation intensity one narrow wavelength range at a time.

\subsection*{(d) Examples of spectrometers}
With an appropriate choice of spectrometer, absorption spectroscopy can be used to probe electronic, vibrational, and rotational transitions in molecules. It is often necessary to modify the general design of Fig. 11A. 6 in order to detect weak signals. For example, to detect rotational transitions with a microwave spectrometer it is useful to modulate the transmitted intensity by varying the energy levels with an oscillating electric field. In this Stark modulation, an electric field of about $10^{5} \mathrm{~V} \mathrm{~m}^{-1}\left(1 \mathrm{kV} \mathrm{cm}^{-1}\right)$ and a frequency of between 10 and 100 kHz is applied to the sample.

In a typical Raman spectroscopy experiment, a monochromatic incident laser beam is scattered from the front face of the sample and monitored (Fig. 11A.14). Lasers are used as the source of the incident radiation because the scattered beam is then more intense. The monochromatic character of laser radiation makes possible the observation of Stokes and antiStokes lines with frequencies that differ only slightly from that of the incident radiation. Such high resolution is particularly useful for observing rotational transitions by Raman spectroscopy.\\
\includegraphics[max width=\textwidth, center]{2025_02_27_d4b95d9718f0b9e2e6b9g-461}

Figure 11A. 14 A common arrangement adopted in Raman spectroscopy. A laser beam passes through a lens and then through a small hole in a mirror with a curved reflecting surface. The focused beam strikes the sample and scattered light is both deflected and focused by the mirror. The spectrum is analysed by a monochromator or an interferometer.

\subsection*{Checklist of concepts}
\begin{enumerate}
  \item In emission spectroscopy the electromagnetic radiation that arises from molecules undergoing a transition from a higher energy state to a lower energy state is detected.2. In absorption spectroscopy, the net absorption of radiation passing through a sample is monitored.3. In Raman spectroscopy, changes in molecular state are explored by examining the energies (frequencies) of the photons scattered by molecules.4. Photons which are scattered elastically give rise to Raleigh scattering.5. In Stokes scattering a photon gives up some of its energy to a molecule; in anti-Stokes scattering the photon gains energy from the molecule.6. Stimulated absorption is a process in which a transition from a low energy state to one of higher energy is driven by an electromagnetic field oscillating at the transition frequency; its rate is determined in part by the Einstein coefficient of stimulated absorption.
  \item Stimulated emission is a process in which a transition from a high energy state to one of lower energy is driven by an electromagnetic field oscillating at the transition frequency; its rate is determined in part by the Einstein coefficient of stimulated emission.8. Spontaneous emission is the transition from a high energy state to a lower energy state at a rate independent of any radiation also present. The relative importance of spontaneous emission increases as the cube of the transition frequency.9. A gross selection rule specifies the general features a molecule must have if it is to have a spectrum of a given kind; a specific selection rule expresses the allowed transitions in terms of the changes in quantum numbers.10. Collisional line broadening arises from the shortened lifetime due to collisions. The broadening is proportional to the pressure, and is often termed pressure broadening.
\end{enumerate}

\subsection*{Checklist of equations}
\begin{center}
\begin{tabular}{|c|c|c|c|}
\hline
Property & Equation & Comment & Equation number \\
\hline
Ratio of Einstein coefficients of spontaneous and stimulated emission & $A_{\mathrm{l}, \mathrm{u}}=\left(8 \pi h v^{3} / c^{3}\right) B_{\mathrm{l}, \mathrm{u}}$ & $B_{\mathrm{u}, 1}=B_{\mathrm{l}, \mathrm{u}}$ & 11A.6a \\
\hline
Transition dipole moment & $\mu_{\mathrm{fi}}=\int \psi_{\mathrm{f}}^{\star} \hat{\mu} \psi_{\mathrm{i}} \mathrm{d} \tau$ & Electric dipole transitions & 11A. 7 \\
\hline
Beer-Lambert law & $I=I_{0} 10^{-\varepsilon[]] L}$ & Uniform sample & 11A. 8 \\
\hline
Absorbance and transmittance & $A=\log \left(I_{0} / I\right)=-\log T$ & Definition & 11A.9b \\
\hline
Integrated absorption coefficient & $\mathcal{A}=\int_{\text {band }} \varepsilon(\tilde{v}) \mathrm{d} \tilde{v}$ & Definition & 11A.10 \\
\hline
Doppler broadening & $\delta v_{\text {obs }}=\left(2 v_{0} / c\right)(2 k T \ln 2 / m)^{1 / 2}$ &  & 11A.12a \\
\hline
Lifetime broadening & $\delta E / h=1 / 2 \pi \tau$ & $\tau$ is the lifetime of the state &  \\
\hline
\end{tabular}
\end{center}

\section*{TOPIC 11B Rotational spectroscopy}
\subsection*{Why do you need to know this material?}
Rotational spectroscopy provides very precise values of bond lengths, bond angles, and dipole moments of molecules in the gas phase.

\subsection*{What is the key idea?}
The spacing of the lines in rotational spectra is used to determine the rotational constants of molecules and, through them, values of their bond lengths and angles.

\subsection*{What do you need to know already?}
You need to be familiar with the classical description of rotational motion (The chemist's toolkit 20 in Topic 7F), the quantization of angular momentum (Topic 7F), the general principles of molecular spectroscopy (Topic 11A), and the Pauli principle (Topic 8B).

Pure rotational spectra, in which only the rotational state of a molecule changes, can be observed only in the gas phase. In spite of this limitation, rotational spectroscopy can provide a wealth of information about molecules, including precise bond lengths, angles, and dipole moments.

\subsection*{11B.1 Rotational energy levels}
The classical expression for the energy of a body rotating about an axis $q$ (The chemist's toolkit 20 in Topic 7F) is


\begin{equation*}
E_{q}=\frac{1}{2} I_{q} \omega_{q}^{2} \tag{11B.1}
\end{equation*}


where $\omega_{q}$ is the angular velocity about the axis $q=x, y, z$ and $I_{q}$ is the corresponding moment of inertia. The moment of inertia, $I$, of a molecule about an axis passing through the centre of mass is defined as (Fig. 11B.1)


\begin{equation*}
I=\sum_{i} m_{i} x_{i}^{2} \tag{11B.2}
\end{equation*}


Moment of inertia [definition]\\
where $m_{i}$ is the mass of the atom $i$ treated as a point and $x_{i}$ is its perpendicular distance from the axis of rotation. In general, the rotational properties of any molecule can be expressed in\\
\includegraphics[max width=\textwidth, center]{2025_02_27_d4b95d9718f0b9e2e6b9g-462}

Figure 11B. 1 The definition of moment of inertia. In this molecule there are three atoms with mass $m_{A}$ attached to the $B$ atom and three atoms with mass $m_{D}$ attached to the $C$ atom. The moment of inertia about the axis passing through the $B$ and $C$ atoms depends on the perpendicular distance $x_{\mathrm{A}}$ from this axis to the A atoms, and the perpendicular distance $x_{D}$ to the D atoms.\\
terms of its three principal moments of inertia $I_{q}$ about three mutually perpendicular axes, $q=x, y, z$. For linear molecules, the moment of inertia around the internuclear axis is zero (because $x_{i}=0$ for all the atoms) and the two remaining moments of inertia, which are equal, are denoted simply $I$. Explicit expressions for the moments of inertia of some symmetrical molecules are given in Table 11B.1. The principal moments of inertia are also commonly recorded as $I_{a}, I_{b}$, and $I_{c}$, with $I_{c} \geq I_{b} \geq I_{a}$.

A note on good practice The mass to use in the calculation of the moment of inertia is the actual atomic mass, not the element's molar mass; don't forget to convert from relative masses to actual masses by using the atomic mass constant $m_{\mathrm{u}}$.

The energy of a body free to rotate about three axes is


\begin{equation*}
E=\frac{1}{2} I_{x} \omega_{x}^{2}+\frac{1}{2} I_{y} \omega_{y}^{2}+\frac{1}{2} I_{z} \omega_{z}^{2} \tag{11B.3}
\end{equation*}


Because the classical angular momentum about the axis $q$ is $J_{q}=I_{q} \omega_{q}$, it follows that

$$
E=\frac{J_{x}^{2}}{2 I_{x}}+\frac{J_{y}^{2}}{2 I_{y}}+\frac{J_{z}^{2}}{2 I_{z}} \quad \begin{aligned}
& \text { Rotational energy: } \\
& \text { classical expression }
\end{aligned}
$$

(11B.4)

Table 11B. 1 Moments of inertia*\\
\includegraphics[max width=\textwidth, center]{2025_02_27_d4b95d9718f0b9e2e6b9g-463(4)}\\
3. Symmetric rotors\\
\includegraphics[max width=\textwidth, center]{2025_02_27_d4b95d9718f0b9e2e6b9g-463(3)}\\
$I_{\|}=2 m_{\mathrm{A}} f_{1}(\theta) R^{2}$\\
$I_{\perp}=m_{\mathrm{A}} f_{1}(\theta) R^{2}$\\
$+\frac{m_{\mathrm{A}} m_{\mathrm{B}}}{m} f_{2}(\theta) R^{2}$\\
\includegraphics[max width=\textwidth, center]{2025_02_27_d4b95d9718f0b9e2e6b9g-463(1)}\\
4. Spherical rotors\\
\includegraphics[max width=\textwidth, center]{2025_02_27_d4b95d9718f0b9e2e6b9g-463(2)}\\
\includegraphics[max width=\textwidth, center]{2025_02_27_d4b95d9718f0b9e2e6b9g-463}

$$
I=4 m_{A} R^{2}
$$

\begin{itemize}
  \item $f_{1}(\theta)=1-\cos \theta, \quad f_{2}(\theta)=1+2 \cos \theta$; in each case, $m$ is the total mass of the molecule.
\end{itemize}

\subsection*{Example 11B. 1}
Calculating the moment of inertia of a molecule

Calculate the moment of inertia of an $\mathrm{H}_{2} \mathrm{O}$ molecule around the axis defined by the bisector of the HOH angle (1). The HOH bond angle is $104.5^{\circ}$ and the OH bond length is 95.7 pm . Use $m\left({ }^{1} \mathrm{H}\right)=1.0078 m_{\mathrm{u}}$.\\
\includegraphics[max width=\textwidth, center]{2025_02_27_d4b95d9718f0b9e2e6b9g-463(5)}

Collect your thoughts You can compute the moment of inertia using eqn 11B.2. In this equation the $x_{i}$ are the perpendicular distances from each atom to the axis of rotation, and you will be able to calculate these distances by using trigonometry and the bond angle and bond length.

The solution From eqn 11B.2,

$$
I=\sum_{i} m_{i} x_{i}^{2}=m_{\mathrm{H}} x_{\mathrm{H}}^{2}+0+m_{\mathrm{H}} x_{\mathrm{H}}^{2}=2 m_{\mathrm{H}} x_{\mathrm{H}}^{2}
$$

If the bond angle of the molecule is $\phi$ and the OH bond length is $R$, trigonometry gives $x_{\mathrm{H}}=R \sin \frac{1}{2} \phi$. It follows that

$$
I=2 m_{\mathrm{H}} R^{2} \sin ^{2} \frac{1}{2} \phi
$$

Substitution of the data gives

$$
\begin{aligned}
I= & 2 \times\left(1.0078 \times 1.6605 \times 10^{-27} \mathrm{~kg}\right) \times\left(9.57 \times 10^{-11} \mathrm{~m}\right)^{2} \\
& \times \sin ^{2}\left(\frac{1}{2} \times 104.5^{\circ}\right) \\
= & 1.92 \times 10^{-47} \mathrm{~kg} \mathrm{~m}^{2}
\end{aligned}
$$

Note that the mass of the O atom makes no contribution to the moment of inertia because the axis passes through this atom and so it does not move.

Self-test 11B. 1 Calculate the moment of inertia of a $\mathrm{CH}^{35} \mathrm{Cl}_{3}$ molecule around a rotational axis that contains the $\mathrm{C}-\mathrm{H}$ bond. The $\mathrm{C}-\mathrm{Cl}$ bond length is 177 pm and the HCCl angle is $107^{\circ} ; m\left({ }^{35} \mathrm{Cl}\right)=34.97 m_{\mathrm{u}}$.

\subsection*{(a) Spherical rotors}
Spherical rotors have all three moments of inertia equal, as in $\mathrm{CH}_{4}$ and $\mathrm{SF}_{6}$. If these moments of inertia have the value $I$, the classical expression for the energy is


\begin{equation*}
E=\frac{J_{x}^{2}+J_{y}^{2}+J_{z}^{2}}{2 I}=\frac{J^{2}}{2 I} \tag{11B.5}
\end{equation*}


where $\mathcal{J}^{2}$ is the square of the magnitude of the angular momentum. The corresponding quantum expression is generated by making the replacement

$$
J^{2} \rightarrow J(J+1) \hbar^{2} \quad J=0,1,2, \ldots
$$

where $J$ is the angular momentum quantum number. Therefore, the energy of a spherical rotor is confined to the values

\[
E_{J}=J(J+1) \frac{\hbar^{2}}{2 I} \quad J=0,1,2, \ldots \quad \begin{align*}
& \text { Energy levels of }  \tag{11B.6}\\
& \text { a spherical rotor }
\end{align*}
\]

The resulting ladder of energy levels is illustrated in Fig. 11B.2. The energy is normally expressed in terms of the rotational constant, $\tilde{B}$ (a wavenumber), of the molecule, where

\[
h c \tilde{B}=\frac{\hbar^{2}}{2 I} \quad \text { so } \quad \tilde{B}=\frac{\hbar}{4 \pi c I} \quad \begin{align*}
& \text { Rotational constant }  \tag{11B.7}\\
& \text { [definition] }
\end{align*}
\]

The expression for the energy is then

\[
E_{J}=h c \tilde{B} J(J+1) \quad J=0,1,2, \ldots \quad \begin{align*}
& \text { Energy levels of }  \tag{11B.8}\\
& \text { a spherical rotor }
\end{align*}
\]

It is also common to express the rotational constant as a frequency and to denote it $B$. Then $B=\hbar / 4 \pi I$ and the energy is $E_{J}=h B J(J+1)$. The two quantities are related by $B=c \tilde{B}$.

The energy of a rotational state is normally reported as the rotational term, $\tilde{F}(J)$, a wavenumber, by division of both sides of eqn 11B. 6 by $h c$ :

\[
\tilde{F}(J)=\tilde{B} J(J+1) \quad \begin{array}{ll}
\text { Rotational terms }  \tag{11B.9}\\
\text { of spherical rotor }
\end{array}
\]

To express the rotational term as a frequency, use $F=c \tilde{F}$. The separation of adjacent terms is


\begin{equation*}
\tilde{F}(J+1)-\tilde{F}(J)=\tilde{B}(J+1)(J+2)-\tilde{B} J(J+1)=2 \tilde{B}(J+1) \tag{11B.10}
\end{equation*}


Because the rotational constant is inversely proportional to $I$, large molecules have closely spaced rotational energy levels.

\subsection*{Brief illustration 11B. 1}
Consider ${ }^{12} \mathrm{C}^{35} \mathrm{Cl}_{4}$ : from Table 11 B .1 and given the $\mathrm{C}-\mathrm{Cl}$ bond length ( $R_{\mathrm{C}-\mathrm{Cl}}=177 \mathrm{pm}$ ) and the mass of the ${ }^{35} \mathrm{Cl}$ nuclide $\left(m\left({ }^{35} \mathrm{Cl}\right)=34.97 m_{\mathrm{u}}\right)$, find

$$
\begin{aligned}
I & =\frac{8}{3} m\left({ }^{35} \mathrm{Cl}\right) R_{\mathrm{C}-\mathrm{Cl}}^{2}=\frac{8}{3} \times \overbrace{\left(5.807 \times 10^{-26} \mathrm{~kg}\right)}^{34.97 m_{\mathrm{u}}} \times\left(1.77 \times 10^{-10} \mathrm{~m}\right)^{2} \\
& =4.85 \ldots \times 10^{-45} \mathrm{~kg} \mathrm{~m}^{2}
\end{aligned}
$$

and, from eqn 11B. 7 ,

$$
\begin{aligned}
\tilde{B} & =\frac{\mathrm{kg} \mathrm{~m}^{2} \mathrm{~s}^{-2}}{4 \pi \times\left(2.998 \times 10^{8} \mathrm{~m} \mathrm{~s}^{-1}\right) \times\left(4.85 \ldots \times 10^{-45} \mathrm{~kg} \mathrm{~m}^{2}\right)} \\
& =5.77 \mathrm{~m}^{-1}=0.0577 \mathrm{~cm}^{-1}
\end{aligned}
$$

It follows from eqn 11B. 10 that the separation between the $J=0$ and $J=1$ terms is $\tilde{F}(1)-\tilde{F}(0)=2 \tilde{B}=0.1154 \mathrm{~cm}^{-1}$, corresponding to 3.46 GHz .

\subsection*{(b) Symmetric rotors}
Symmetric rotors have two equal moments of inertia and a third that is non-zero. In group theoretical terms (Topic 10 A ), such rotors have an $n$-fold axis of rotation, with $n>2$.

The unique axis of a symmetric rotor (such as $\mathrm{CH}_{3} \mathrm{Cl}$, $\mathrm{NH}_{3}$, and $\mathrm{C}_{6} \mathrm{H}_{6}$ ) is its principal axis (or figure axis). If the moment of inertia about the principal axis is larger than the other two, the rotor is classified as oblate (like a pancake, and $\mathrm{C}_{6} \mathrm{H}_{6}$ ). If the moment of inertia around the principal axis is smaller than the other two, the rotor is classified as prolate (like a cigar, and $\mathrm{CH}_{3} \mathrm{Cl}$ ). The two equal moments of inertia ( $I_{x}$ and $I_{y}$ ) are denoted $I_{\perp}$ and $I_{z}$ is denoted $I_{\|}$; then eqn 11B. 4 becomes


\begin{equation*}
E=\frac{J_{x}^{2}+J_{y}^{2}}{2 I_{\perp}}+\frac{J_{z}^{2}}{2 I_{\|}} \tag{11B.11}
\end{equation*}


This expression can be written in terms of $\mathcal{J}^{2}=J_{x}^{2}+J_{y}^{2}+J_{z}^{2}$ :


\begin{equation*}
E=\frac{\mathcal{J}^{2}-J_{z}^{2}}{2 I_{\perp}}+\frac{J_{z}^{2}}{2 I_{\|}}=\frac{\mathcal{J}^{2}}{2 I_{\perp}}+\left(\frac{1}{2 I_{\|}}-\frac{1}{2 I_{\perp}}\right) J_{z}^{2} \tag{11B.12}
\end{equation*}


The quantum expression is generated by replacing $\mathcal{I}^{2}$ by $J(J+1) \hbar^{2}$. The quantum theory of angular momentum (Topic 7 F ) also restricts the component of angular momentum about any axis to the values $K \hbar$, with $K=0, \pm 1, \ldots, \pm J$. (The quantum number $K$ is used to signify the component on the principal axis, as distinct from the quantum number $M_{\mathrm{J}}$ which is used to signify the component on an externally defined axis.) Then, after making the replacements $\mathcal{J}^{2} \rightarrow J(J+1) \hbar^{2}$ and $J_{z}^{2} \rightarrow K^{2} \hbar^{2}$ the rotational terms are

\[
\begin{array}{ll}
\tilde{F}(J, K)=\tilde{B} J(J+1)+(\tilde{A}-\tilde{B}) K^{2} & \text { Rotational terms of }  \tag{11B.13a}\\
J=0,1,2, \ldots \quad K=0, \pm 1, \ldots, \pm J & \text { a symmetric rotor }
\end{array}
\]

with


\begin{equation*}
\tilde{A}=\frac{\hbar}{4 \pi c I_{\|}} \quad \tilde{B}=\frac{\hbar}{4 \pi c I_{\perp}} \tag{11B.13b}
\end{equation*}


Equation 11B.13a matches what is expected for the dependence of the energy levels on the two distinct moments of inertia of the molecule:

\begin{itemize}
  \item When $K=0$, there is no component of angular momentum about the principal axis, and the energy levels depend only on $I_{\perp}$ (Fig. 11B.3a).
  \item When $K= \pm J$, almost all the angular momentum arises from rotation around the principal axis, and the energy levels are determined largely by $I_{\|}$ (Fig. 11B.3b).
  \item The sign of $K$ does not affect the energy because opposite values of $K$ correspond to opposite senses of rotation, and the energy does not depend on the sense of rotation.\\
\includegraphics[max width=\textwidth, center]{2025_02_27_d4b95d9718f0b9e2e6b9g-465(2)}
\end{itemize}

Figure 11B. 3 The significance of the quantum number $K$. (a) When $K=0$ the molecule has no angular momentum about its principal axis: it is undergoing end-over-end rotation. (b) When $|K|$ is close to its maximum value, $J$, most of the molecular rotation is around the principal axis.

Example 11B. 2 Calculating the rotational energy levels of a symmetric rotor

A ${ }^{14} \mathrm{~N}^{1} \mathrm{H}_{3}$ molecule is a symmetric rotor with bond length 101.2 pm and HNH bond angle $106.7^{\circ}$. Calculate its rotational terms.

Collect your thoughts Begin by calculating the moments of inertia by using the expressions given in Table 11B.1. Then use eqn 11B.13a to find the rotational terms. The rotational constants are found using eqn 11B.13b.

The solution Substitution of $m_{\mathrm{A}}=1.0078 m_{\mathrm{u}}, m_{\mathrm{B}}=14.0031 m_{\mathrm{u}}$, $R=101.2 \mathrm{pm}$, and $\theta=106.7^{\circ}$ into the second of the symmetric rotor expressions in Table 11B. 1 gives $I_{\|}=4.4128 \times 10^{-47} \mathrm{~kg} \mathrm{~m}^{2}$ and $I_{\stackrel{\perp}{ }}=2.8059 \times 10^{-47} \mathrm{~kg} \mathrm{~m}^{2}$. The expressions in eqn. 11B.13b give $\tilde{\tilde{A}}=6.344 \mathrm{~cm}^{-1}$ and $\tilde{B}=9.977 \mathrm{~cm}^{-1}$. It follows from eqn 11B.13a that

$$
\tilde{F}(J, K) / \mathrm{cm}^{-1}=9.977 J(J+1)-3.633 K^{2}
$$

Multiplication by $c$ converts $\tilde{F}(J, K)$ to a frequency, denoted $F(J, K)$ :

$$
F(J, K) / \mathrm{GHz}=299.1 J(J+1)-108.9 K^{2}
$$

For $J=1$, the energy needed for the molecule to rotate mainly about its principal axis $(K= \pm J)$ is equivalent to $16.32 \mathrm{~cm}^{-1}$ ( 489.3 GHz ), but end-over-end rotation $(K=0)$ corresponds to $19.95 \mathrm{~cm}^{-1}(598.1 \mathrm{GHz})$.\\
Self-test 11B. $2 \quad \mathrm{~A}^{12} \mathrm{C}^{1} \mathrm{H}_{3}^{35} \mathrm{Cl}$ molecule has a $\mathrm{C}-\mathrm{Cl}$ bond length of $178 \mathrm{pm}, \mathrm{a}-\mathrm{H}$ bond length of 111 pm , and an HCH angle of $110.5^{\circ}$. Identify whether the molecule is oblate or prolate, and calculate its rotational energy terms.\\
\includegraphics[max width=\textwidth, center]{2025_02_27_d4b95d9718f0b9e2e6b9g-465}\\
\includegraphics[max width=\textwidth, center]{2025_02_27_d4b95d9718f0b9e2e6b9g-465(1)}\\
\includegraphics[max width=\textwidth, center]{2025_02_27_d4b95d9718f0b9e2e6b9g-465(3)}\\
\includegraphics[max width=\textwidth, center]{2025_02_27_d4b95d9718f0b9e2e6b9g-466}

Figure 11B. 4 The significance of the quantum number $M_{J}$. (a) When $M$, is close to its maximum value, $J$, most of the molecular rotation is around the laboratory axis (taken as the $z$-axis). (b) An intermediate value of $M_{J}$. (c) When $M_{J}=0$ the molecule has no angular momentum about the $z$-axis. All three diagrams correspond to a state with $K=0$; there are corresponding diagrams for different values of $K$, in which the angular momentum makes a different angle to the principal axis of the molecule.

The energy of a symmetric rotor depends on $J$ and $K$. Because the states with $K$ and $-K$ have the same energy, each level, except those with $K=0$, is doubly degenerate. In addition, the angular momentum of the molecule has a component on an external, laboratory-fixed axis. This component is quantized, and its permitted values are $M_{J} \hbar$, with $M_{J}=0, \pm 1, \ldots, \pm J$, giving $2 J+1$ values in all (Fig. 11B.4). The quantum number $M_{J}$ does not appear in the expression for the energy, but it is necessary for a complete specification of the state of the rotor. Consequently, all $2 J+1$ orientations of the rotating molecule have the same energy. It follows that a symmetric rotor level is $2(2 J+1)$-fold degenerate for $K \neq 0$ and $(2 J+1)$-fold degenerate for $K=0$.

A spherical rotor can be regarded as a version of a symmetric rotor in which $I_{\perp}=I_{\|}$and therefore $\tilde{A}=\tilde{B}$. The quantum number $K$ still takes any one of $2 J+1$ values, but the energy is independent of which value it takes. Therefore, as well as having a $(2 J+1)$-fold degeneracy arising from its orientation in space, the rotor also has a $(2 J+1)$-fold degeneracy arising from its orientation with respect to an arbitrary axis in the molecule. The overall degeneracy of a symmetric rotor energy level with quantum number $J$ is therefore $(2 J+1)^{2}$. This degeneracy increases very rapidly: when $J=10$, for instance, there are 441 states of the same energy.

\subsection*{(c) Linear rotors}
For a linear rotor (such as $\mathrm{CO}_{2}, \mathrm{HCl}$, and $\mathrm{C}_{2} \mathrm{H}_{2}$ ), in which the atoms are regarded as mass points, the rotation occurs only about an axis perpendicular to the internuclear axis and there is no rotation around that axis. Therefore the component of angular momentum around the internuclear axis of a linear\\
rotor is identically zero, and $K \equiv 0$ in eqn 11B.13a. The rotational terms of a linear molecule are therefore

\[
\tilde{F}(J)=\tilde{B} J(J+1) \quad J=0,1,2, \ldots \quad \begin{align*}
& \text { Rotational terms }  \tag{11B.14}\\
& \text { of linear rotor }
\end{align*}
\]

This expression is the same as eqn 11B. 9 but arrived at it in a significantly different way: here $K \equiv 0$, but for a spherical rotor $\tilde{A}=\tilde{B}$ and $K$ has a range of values. The angular momentum of a linear rotor has $2 J+1$ components on an external axis, so its degeneracy is just $2 J+1$ rather than the $(2 J+1)^{2}$-fold degeneracy of a spherical rotor.

\subsection*{Brief illustration 11B. 2}
Equation 11B. 10 for the energy separation of adjacent levels of a spherical rotor also applies to linear rotors, so $\tilde{F}(3)-\tilde{F}(2)=$ $6 \tilde{B}$. Spectroscopic measurements on ${ }^{1} \mathrm{H}^{35} \mathrm{Cl}$ gives $\tilde{F}(3)-\tilde{F}(2)=$ $63.56 \mathrm{~cm}^{-1}$, so it follows that $6 \tilde{B}=63.56 \mathrm{~cm}^{-1}, \tilde{B}=10.59 \mathrm{~cm}^{-1}$, and therefore

$$
\begin{aligned}
I & =\frac{\hbar}{4 \pi c \tilde{B}}=\frac{1.05457 \times 10^{-34} \mathrm{~J} \mathrm{~s}}{4 \pi \times\left(2.998 \times 10^{10} \mathrm{~cm} \mathrm{~s}^{-1}\right) \times\left(10.59 \mathrm{~cm}^{-1}\right)} \\
& =2.643 \times 10^{-47} \mathrm{~kg} \mathrm{~m}^{2}
\end{aligned}
$$

\subsection*{(d) Centrifugal distortion}
In the discussion so far molecules have been treated as rigid rotors. However, the atoms of rotating molecules are subject to centrifugal forces which tend to distort the molecular geometry and change its moments of inertia (Fig. 11B.5). The effect of centrifugal distortion on a diatomic molecule is to stretch the bond and hence to increase the moment of inertia. As a result, the rotational constant is reduced and the energy levels are slightly closer together than the rigid-rotor expressions predict. The effect is usually taken into account by including in the energy expression a negative term that becomes more important as $J$ increases:

\[
\tilde{F}(J)=\tilde{B} J(J+1)-\tilde{D}_{J} J^{2}(J+1)^{2} \quad \begin{align*}
& \text { Rotational terms affected }  \tag{11B.15}\\
& \text { by centrifugal distortion }
\end{align*}
\]

The parameter $\tilde{D}_{J}$ is the centrifugal distortion constant. The centrifugal distortion constant of a diatomic molecule is related to the vibrational wavenumber of the bond, $\tilde{v}$ (which, as seen in Topic 11C, is a measure of its stiffness), through the approximate relation (see Problem P11C.16)


\begin{equation*}
\tilde{D}_{J}=\frac{4 \tilde{B}^{3}}{\tilde{v}^{2}} \tag{11B.16}
\end{equation*}


Centrifugal distortion constant

As expected, a bond that is easily stretched, and therefore has a low vibrational wavenumber, has a high centrifugal distortion constant.\\
\includegraphics[max width=\textwidth, center]{2025_02_27_d4b95d9718f0b9e2e6b9g-467(1)}

Figure 11B. 5 The effect of rotation on a molecule. The centrifugal force arising from rotation distorts the molecule, opening out bond angles and stretching bonds slightly. The effect is to increase the moment of inertia of the molecule and hence to decrease its rotational constant.

\subsection*{Brief illustration 11B. 3}
For ${ }^{12} \mathrm{C}^{16} \mathrm{O}, \tilde{B}=1.931 \mathrm{~cm}^{-1}$ and $\tilde{v}=2170 \mathrm{~cm}^{-1}$. It follows that

$$
\tilde{D}_{I}=\frac{4 \times\left(1.931 \mathrm{~cm}^{-1}\right)^{3}}{\left(2170 \mathrm{~cm}^{-1}\right)^{2}}=6.116 \times 10^{-6} \mathrm{~cm}^{-1}
$$

Because $\tilde{D}_{J} \ll \tilde{B}$, centrifugal distortion has a very small effect on the energy levels until $J$ is large. For $J=20, \tilde{D}_{J} J^{2}(J+1)^{2}=$ $1.08 \mathrm{~cm}^{-1}$ (corresponding to 32 GHz ).

\subsection*{118.2 Microwave spectroscopy}
Typical values of the rotational constant $\tilde{B}$ for small molecules are in the region of $0.1-10 \mathrm{~cm}^{-1}$; two examples are $0.356 \mathrm{~cm}^{-1}$ for $\mathrm{NF}_{3}$ and $10.59 \mathrm{~cm}^{-1}$ for HCl . It follows that rotational transitions can be studied with microwave spectroscopy, a technique that monitors the absorption of radiation in the microwave region of the spectrum.

\subsection*{(a) Selection rules}
As usual in spectroscopy, the selection rules can be established by considering the relevant transition dipole moment. The details of the calculation are shown in A deeper look 5 on the website of this text. The conclusion is that the gross selection rule for the observation of a pure rotational transition is that a molecule must have a permanent electric dipole moment. The classical basis of this rule is that a polar molecule appears to possess a fluctuating dipole when rotating, but a nonpolar molecule does not (Fig. 11B.6). The permanent dipole can be regarded as a handle with which the molecule stirs the electromagnetic field into oscillation (and vice versa for absorption).\\
\includegraphics[max width=\textwidth, center]{2025_02_27_d4b95d9718f0b9e2e6b9g-467}

Figure 11B. 6 To a stationary observer, a rotating polar molecule looks like an oscillating dipole which will generate an oscillating electromagnetic wave (or, in the case of absorption, interact with such a wave). This picture is the classical origin of the gross selection rule for rotational transitions.

\subsection*{Brief illustration 11B. 4}
Homonuclear diatomic molecules and nonpolar polyatomic molecules such as $\mathrm{CO}_{2}, \mathrm{CH}_{2}=\mathrm{CH}_{2}$, and $\mathrm{C}_{6} \mathrm{H}_{6}$ do not give rise to microwave spectra. On the other hand, OCS and $\mathrm{H}_{2} \mathrm{O}$ are polar and have microwave spectra. Spherical rotors cannot have electric dipole moments unless they become distorted by rotation, so they are rotationally inactive except in special cases. An example of a spherical rotor that does become sufficiently distorted for it to acquire a dipole moment is $\mathrm{SiH}_{4}$, which has a dipole moment of about $8.3 \mu \mathrm{D}$ by virtue of its rotation when $J \approx 10$ (for comparison, HCl has a permanent dipole moment of 1.1 D ; molecular dipole moments and their units are discussed in Topic 14A).

The analysis also shows that, for a linear molecule, the transition moment vanishes unless the following conditions are fulfilled:

$$
\begin{array}{lll}
\Delta J= \pm 1 & \Delta M_{J}=0, \pm 1 & \begin{array}{l}
\text { Rotational selection } \\
\text { rules: linear rotors }
\end{array}
\end{array}
$$

The transition $\Delta J=+1$ corresponds to absorption and the transition $\Delta J=-1$ corresponds to emission.

\begin{itemize}
  \item The allowed change in $J$ arises from the conservation of angular momentum when a photon, a spin-1 particle, is emitted or absorbed (Fig. 11B.7).
  \item The allowed change in $M_{J}$ also arises from the conservation of angular momentum when a photon is emitted or absorbed in a specific direction.
\end{itemize}

When the transition moment is evaluated for all possible relative orientations of the molecule to the line of flight of the photon, it is found that the total $J+1 \leftrightarrow J$ transition intensity is proportional to


\begin{equation*}
\left.\mu_{J+1, J}\right|^{2}=\left(\frac{J+1}{2 J+1}\right) \mu_{0}^{2} \tag{11B.18}
\end{equation*}


\begin{center}
\includegraphics[max width=\textwidth]{2025_02_27_d4b95d9718f0b9e2e6b9g-468(1)}
\end{center}

Figure 11B. 7 When a photon is absorbed by a molecule, the angular momentum of the combined system is conserved. If the molecule is rotating in the same sense as the spin of the incoming photon, then $J$ increases by 1.\\
where $\mu_{0}$ is the permanent electric dipole moment of the molecule. The intensity is proportional to the square of $\mu_{0}$, so strongly polar molecules give rise to much more intense rotational lines than less polar molecules.

Rotation of a symmetric rotor about its principal (figure) axis does not lead to any change in the orientation of the dipole; there is no fluctuating dipole to interact with the radiation, and therefore no change in $K$ is possible. For symmetric rotors the selection rules are therefore:

\[
\Delta J= \pm 1 \quad \Delta M_{J}=0, \pm 1 \quad \Delta K=0 \quad \begin{align*}
& \text { Rotational selection }  \tag{11B.19}\\
& \text { rules: symmetric rotors }
\end{align*}
\]

The degeneracy associated with the quantum number $M_{J}$ (the orientation of the rotation in space) is partly removed when an electric field is applied to a polar molecule (Fig. 11B.8). The splitting of states by an electric field is called the Stark effect. The energy shift depends on the permanent electric dipole moment, $\mu_{0}$, so the observation of the Stark effect in a rotational spectrum can be used to measure the magnitudes of electric dipole moments.\\
\includegraphics[max width=\textwidth, center]{2025_02_27_d4b95d9718f0b9e2e6b9g-468(2)}

Figure 11B. 8 The effect of an electric field on the energy level with $J=7$ of a polar linear rotor. All levels are doubly degenerate except that with $M_{J}=0$.

\subsection*{(b) The appearance of microwave spectra}
When the selection rules are applied to the expressions for the energy levels of a linear rigid rotor (eqn 11B.14), it follows that the wavenumbers of the allowed $J+1 \leftarrow J$ absorptions are

$$
\tilde{\tilde{V}(J+1 \leftarrow J)=\tilde{F}(J+1)-\tilde{F}(J)=2 \tilde{B}(J+1) \quad J=0,1,2, \ldots}
$$

(11B.20a)

When centrifugal distortion is taken into account, the corresponding expression obtained from eqn 11B. 15 is


\begin{equation*}
\tilde{v}(J+1 \leftarrow J)=2 \tilde{B}(J+1)-4 \tilde{D}_{J}(J+1)^{3} \tag{11B.20b}
\end{equation*}


However, because the second term is typically very small compared with the first, the appearance of the spectrum closely resembles that predicted from eqn 11B.20a.

\subsection*{Example 11B. 3 \\
 Predicting the appearance of a rotational spectrum}
Predict the form of the rotational spectrum of ${ }^{14} \mathrm{NH}_{3}$, which is an oblate symmetric rotor with $\tilde{B}=9.977 \mathrm{~cm}^{-1}$.

Collect your thoughts The rotational terms are given by eqn 11B.13a. Because $\Delta J= \pm 1$ and $\Delta K=0$, the expression for the wavenumbers of the rotational transitions is identical to eqn 11B.20a and depends only on $\tilde{B}$.

The solution The following table can be drawn up for the $J+1 \leftarrow J$ transitions.

\begin{center}
\begin{tabular}{lccccc}
$J$ & 0 & 1 & 2 & 3 & $\ldots$ \\
$\tilde{v} / \mathrm{cm}^{-1}$ & 19.95 & 39.91 & 59.86 & 79.82 & $\ldots$ \\
$v / \mathrm{GHz}$ & 598.2 & 1196 & 1795 & 2393 & $\ldots$ \\
\end{tabular}
\end{center}

The line spacing is $19.95 \mathrm{~cm}^{-1}(598.1 \mathrm{GHz})$.\\
Self-test 11B.3 Repeat the problem for $\mathrm{CH}_{3}^{35} \mathrm{Cl}$, a prolate symmetric rotor for which $\tilde{B}=0.444 \mathrm{~cm}^{-1}$.\\
\includegraphics[max width=\textwidth, center]{2025_02_27_d4b95d9718f0b9e2e6b9g-468}

The form of the spectrum predicted by eqn 11B.20a is shown in Fig. 11B.9. The most significant feature is that it consists of a series of lines with wavenumbers $2 \tilde{B}, 4 \tilde{B}, 6 \tilde{B}, \ldots$ and of separation $2 \tilde{B}$. The measurement of the line spacing therefore gives $\tilde{B}$, and hence the moment of inertia $I_{\perp}$ perpendicular to the principal axis of the molecule. Because the masses of the atoms are known, it is a simple matter to deduce the bond length of a diatomic molecule. However, in the case of a polyatomic molecule such as OCS or $\mathrm{NH}_{3}$, a knowledge of one moment of inertia is insufficient data from which to infer, for example, the two bond lengths in OCS, or the bond length and bond angle in $\mathrm{NH}_{3}$.\\
\includegraphics[max width=\textwidth, center]{2025_02_27_d4b95d9718f0b9e2e6b9g-469}

Figure 11B. 9 The rotational energy levels of a linear rotor, the transitions allowed by the selection rule $\Delta J=+1$, and a typical pure rotational absorption spectrum (displayed here in terms of the radiation transmitted through the sample). The intensities reflect the populations of the initial level in each case and the strengths of the transition dipole moments.

This difficulty can be overcome by measuring the spectra of isotopologues, isotopically substituted molecules. The spectrum from each isotopologue gives a separate moment of inertia and, if it is assumed that the bond lengths and angles are unaffected by isotopic substitution, the extra data make it possible to extract values of the bond lengths and angles. A good example of this procedure is the study of OCS; the actual calculation is worked through in Problem P11B.7. The assumption that bond lengths are unchanged in isotopologues is only an approximation, but it is a good one in most cases. Nuclear spin (Topic 12A), which differs from one isotope to another, also affects the appearance of high-resolution rotational spectra because spin is a source of angular momentum and can couple with the rotation of the molecule itself and hence affect the rotational energy levels.

The intensities of spectral lines increase with increasing $J$ and pass through a maximum before tailing off as $J$ becomes large. The most important reason for this behaviour is the existence of a maximum in the population of rotational levels. The Boltzmann distribution (see the Prologue to this text and Topic 13A) implies that the population of a state decreases exponentially as its energy increases. However, the population of a level is also proportional to its degeneracy, and in the case of rotational levels this degeneracy increases with $J$. These two opposite trends result in the population of the energy levels (as distinct from the individual states) passing through a maximum. Specifically, the population $N_{J}$ of a rotational energy level $J$ is given by the Boltzmann expression

$$
N_{J} \propto N g_{J} \mathrm{e}^{-E_{J} / k T}
$$

where $N$ is the total number of molecules in the sample and $g_{J}$ is the degeneracy of the level $J$. The value of $J$ corresponding\\
to a maximum of this expression is found by treating $J$ as a continuous variable, differentiating with respect to $J$, and then setting the result equal to zero. The result for a linear rotor (see Problem P11B.11) is

\[
J_{\max } \approx\left(\frac{k T}{2 h c \tilde{B}}\right)^{1 / 2}-\frac{1}{2} \quad \begin{align*}
& \text { Rotational level with largest }  \tag{11B.21}\\
& \text { population: linear rotor }
\end{align*}
\]

For a typical molecule (e.g. OCS, with $\tilde{B}=0.2 \mathrm{~cm}^{-1}$ ) $k T / 2 h c \tilde{B} \approx 500$ at room temperature, so $J_{\max } \approx 22$. However, it must be recalled that the transition dipole moment depends on the value of $J$ (eqn 11B.18) and, because the radiation can also cause stimulated emission (Topic 11A), the intensity also depends on the population difference between the two states involved in the transition. Hence the value of $J$ corresponding to the most intense line is not quite the same as the value of $J$ for the most highly populated level.

\subsection*{11B.3 Rotational Raman spectroscopy}
Raman scattering (Topic 11A) can also arise as a result of rotational transitions. The gross selection rule for rotational Raman transitions is that the molecule must be anisotropically polarizable. To understand this criterion it is necessary to know that the distortion of a molecule in an electric field is determined by its polarizability, $\alpha$ (Topic 14A). More precisely, if the strength of the field is $\mathcal{E}$, then the molecule acquires an induced dipole moment of magnitude


\begin{equation*}
\mu=\alpha E \tag{11B.22}
\end{equation*}


in addition to any permanent dipole moment it might have. An atom is isotropically polarizable: that is, the same distortion is induced whatever the direction of the applied field. The polarizability of a spherical rotor is also isotropic. However, non-spherical rotors have polarizabilities that do depend on the direction of the field relative to the molecule, so these molecules are anisotropically polarizable (Fig. 11B.10). The electron distribution in $\mathrm{H}_{2}$, for example, is more distorted when the field is applied parallel to the bond than when it is applied perpendicular to it, and so $\alpha_{\|}>\alpha_{\perp}$.

All linear molecules, including both heteronuclear and homonuclear diatomics, have anisotropic polarizabilities and so are rotationally Raman active. This activity is one reason for the importance of rotational Raman spectroscopy, because the technique can be used to study many of the molecules that are inaccessible to microwave spectroscopy. Spherical rotors such as $\mathrm{CH}_{4}$ and $\mathrm{SF}_{6}$, however, are rotationally Raman inactive as well as microwave inactive. This inactivity does not mean that such molecules are never found in rotationally excited states. Molecular collisions do not have to obey such\\
\includegraphics[max width=\textwidth, center]{2025_02_27_d4b95d9718f0b9e2e6b9g-470(1)}

Figure 11B. 10 An electric field $\mathcal{E}$ applied to a molecule results in its distortion, and the distorted molecule acquires a contribution to its dipole moment (even if it is nonpolar initially). The polarizability may be different when the field is applied (a) parallel or (b) perpendicular to the molecular axis (or, in general, in different directions relative to the molecule); if that is so, then the molecule has an anisotropic polarizability.\\
restrictive selection rules, and hence collisions between molecules can result in the population of any rotational state.

As usual, to establish the selection rules, it is necessary to consider the transition dipole moment. The full calculation can be found in A deeper look 5 on the website of this text and leads to the conclusion that the specific rotational Raman selection rules are

$$
\begin{array}{lll}
\text { Linear rotors: } & \Delta J=0, \pm 2 & \\
\text { Symmetric rotors: } & \Delta J=0, \pm 1, \pm 2 & \begin{array}{l}
\text { Rotational } \\
\text { Raman selection } \\
\text { rules }
\end{array} \\
& \Delta K=0 &
\end{array}
$$

The $\Delta J=0$ transitions do not lead to a shift in frequency of the scattered photon and therefore contribute to the unshifted radiation (the Rayleigh radiation, Topic 11A). A classical argument can be used to give physical insight into the quantum mechanical calculation.

\subsection*{How is that done? 11B. 1 Justifying the rotational Raman selection rules}
The incident electric field, $\mathcal{E}$, of a wave of electromagnetic radiation of frequency $\omega_{i}$ induces a molecular dipole moment given by

$$
\mu_{\mathrm{ind}}=\alpha \mathcal{E}(t)=\alpha \mathcal{E} \cos \omega_{\mathrm{i}} t
$$

If the molecule is rotating at an angular frequency $\omega_{\mathrm{R}}$, it appears to an external observer that the polarizability is also time dependent (if it is anisotropic). This dependence can be written

$$
\alpha=\alpha_{0}+\Delta \alpha \cos 2 \omega_{\mathrm{R}} t
$$

where $\Delta \alpha=\alpha_{\|}-\alpha_{\perp}$ and $\alpha$ ranges from $\alpha_{0}+\Delta \alpha$ to $\alpha_{0}-\Delta \alpha$ as the molecule rotates. The $2 \omega_{\mathrm{R}}$ appears because the polar-\\
izability returns to its initial value twice each revolution (Fig. 11B.11). Combining these expressions gives

$$
\begin{aligned}
& \mu_{\mathrm{ind}}=\left(\alpha_{0}+\Delta \alpha \cos 2 \omega_{\mathrm{R}} t\right) \times\left(\mathcal{E} \cos \omega_{\mathrm{i}} t\right) \\
&=\alpha_{0} E \cos \omega_{\mathrm{i}} t+\mathbb{E} \Delta \alpha \cos \omega_{\mathrm{i}} t \cos 2 \omega_{\mathrm{R}} t \\
&=\begin{array}{c}
\cos x \cos y \\
=\frac{1}{2}\{\cos (x+y)+\cos (x-y)\} \\
\end{array} \\
&=\alpha_{0} E \cos \omega_{\mathrm{i}} t+\frac{1}{2} \mathcal{E} \Delta \alpha\left\{\cos \left(\omega_{\mathrm{i}}+2 \omega_{\mathrm{R}}\right) t+\cos \left(\omega_{\mathrm{i}}-2 \omega_{\mathrm{R}}\right) t\right\}
\end{aligned}
$$

This calculation shows that the induced dipole has a component oscillating at the incident frequency (which results in Rayleigh scattering), and that it also has components at $\omega_{i} \pm 2 \omega_{R}$, which give rise to the shifted Raman lines. These lines appear only if $\Delta \alpha \neq 0$; hence the polarizability must be anisotropic for there to be Raman lines. This is the gross selection rule for rotational Raman spectroscopy.\\
\includegraphics[max width=\textwidth, center]{2025_02_27_d4b95d9718f0b9e2e6b9g-470}

Figure 11B. 11 The distortion induced in a molecule by an applied electric field returns the polarizability to its initial value after a rotation of only $180^{\circ}$ (i.e. twice a revolution). This is the origin of the $\Delta J= \pm 2$ selection rule in rotational Raman spectroscopy.

The distortion induced in the molecule by the incident electric field returns to its initial value after a rotation of $180^{\circ}$ (i.e. twice a revolution). This is the classical origin of the specific selection rule $\Delta J= \pm 2$.

To predict the form of the Raman spectrum of a linear rotor the selection rule $\Delta J= \pm 2$ is applied to the rotational energy levels (Fig. 11B.12). For Stokes lines, $\Delta J=+2$ and the scattered radiation is at a lower wavenumber than the incident radiation at $\tilde{v}_{\mathrm{i}}$, the shift being the difference $\tilde{F}(J+2)-\tilde{F}(J)$

\[
\begin{array}{rlrl}
\tilde{v}(J+2 \leftarrow J) & =\tilde{v}_{\mathrm{i}}-\{\tilde{F}(J+2)-\tilde{F}(J)\} & & \text { Wavenumbers } \\
& =\tilde{v}_{\mathrm{i}}-2 \tilde{B}(2 J+3) & & \text { of Stokes lines: }  \tag{11B.24a}\\
\text { linear rotor }
\end{array}
\]

\begin{center}
\includegraphics[max width=\textwidth]{2025_02_27_d4b95d9718f0b9e2e6b9g-471(2)}
\end{center}

Figure 11B. 12 The rotational energy levels of a linear rotor and the transitions allowed by the $\Delta J= \pm 2$ Raman selection rules. The form of a typical rotational Raman spectrum is also shown. In practice the Rayleigh line is much stronger than depicted in the figure.

For anti-Stokes lines $\Delta J=-2$ and the scattered radiation is at a higher wavenumber, the shift being the difference $\tilde{F}(J)-\tilde{F}(J-2)$ :

\[
\begin{array}{rlrl}
\tilde{v}(J-2 \leftarrow J) & =\tilde{v}_{\mathrm{i}}+\{\tilde{F}(J)-\tilde{F}(J-2)\} & & \text { Wavenumbers } \\
& =\tilde{v}_{\mathrm{i}}+2 \tilde{B}(2 J-1) & & \text { of anti-Stokes }  \tag{11B.24b}\\
\text { lines: linear rotor }
\end{array}
\]

The Stokes lines appear to low frequency of the incident radiation and at displacements $6 \tilde{B}, 10 \tilde{B}, 14 \tilde{B}, \ldots$ from $\tilde{v}_{\mathrm{i}}$ for $J=$ $0,1,2, \ldots$. The anti-Stokes lines occur at displacements of $6 \tilde{B}$, $10 \tilde{B}, 14 \tilde{B}, \ldots$ (for $J=2,3,4, \ldots ; J=2$ is the lowest state that can contribute under the selection rule $\Delta J=-2$ ) to high frequency of the incident radiation. The separation of adjacent lines in both the Stokes and the anti-Stokes regions is $4 \tilde{B}$, so from the spacing $I_{\perp}$ can be determined and then used to find the bond length exactly as in the case of microwave spectroscopy.

\subsection*{Example 11B. 4}
Predicting the form of a Raman spectrum\\
Predict the form of the rotational Raman spectrum of ${ }^{14} \mathrm{~N}_{2}$, for which $\tilde{B}=1.99 \mathrm{~cm}^{-1}$, when it is exposed to 336.732 nm laser radiation.

Collect your thoughts The molecule is rotationally Raman active because end-over-end rotation modulates its polarizability as viewed by a stationary observer. The wavenumbers of the Stokes and anti-Stokes lines are given by eqn 11B.24.

The solution The incident radiation with wavelength 336.732 nm corresponds to a wavenumber of $\tilde{v}_{\mathrm{i}}=29697.2 \mathrm{~cm}^{-1}$; eqns 11B.24a and 11B. 24 b give the following line positions:

\begin{center}
\begin{tabular}{lcccc}
$J$ & 0 & 1 & 2 & 3 \\
Stokes lines &  &  &  &  \\
$\tilde{v} / \mathrm{cm}^{-1}$ & 29685.3 & 29677.3 & 29669.3 & 29661.4 \\
$\lambda / \mathrm{nm}$ & 336.867 & 336.958 & 337.048 & 337.139 \\
\end{tabular}
\end{center}

J\\
0\\
1\\
2\\
3\\
Anti-Stokes lines\\
$\tilde{v} / \mathrm{cm}^{-1}$\\
29709.1\\
336.597\\
29717.1\\
$\lambda / \mathrm{nm}$\\
336.507

There will be a strong central line at 336.732 nm accompanied on either side by lines of increasing and then decreasing intensity (as a result of transition moment and population effects). The spread of the entire spectrum is very small, so the incident light must be highly monochromatic.\\
Self-test 11B.4 Repeat the calculation for the rotational Raman spectrum of ${ }^{35} \mathrm{Cl}_{2}\left(\tilde{B}=0.9752 \mathrm{~cm}^{-1}\right)$.\\
\includegraphics[max width=\textwidth, center]{2025_02_27_d4b95d9718f0b9e2e6b9g-471}\\
\includegraphics[max width=\textwidth, center]{2025_02_27_d4b95d9718f0b9e2e6b9g-471(1)}

\subsection*{11B. 4 Nuclear statistics and rotational states}
If eqn 11B. 24 is used to analyse the rotational Raman spectrum of $\mathrm{C}^{16} \mathrm{O}_{2}$, the rotational constant derived from the spacing of the lines is inconsistent with other measurements of $\mathrm{C}-\mathrm{O}$ bond lengths. The results are consistent if it is supposed that the molecule can exist in states with only even values of $J$, so the observed Stokes lines are $2 \leftarrow 0,4 \leftarrow 2, \ldots$; the lines $3 \leftarrow 1,5 \leftarrow 3, \ldots$ are missing.

The explanation of the missing lines lies in the Pauli principle (Topic 8 B ) and the fact that ${ }^{16} \mathrm{O}$ nuclei are spin- 0 bosons: just as the Pauli principle excludes certain electronic states, so too does it exclude certain molecular rotational states. The Pauli principle states that, when two identical bosons are exchanged, the overall wavefunction must remain unchanged. When a $\mathrm{C}^{16} \mathrm{O}_{2}$ molecule rotates through $180^{\circ}$, two identical ${ }^{16} \mathrm{O}$ nuclei are interchanged, so the overall wavefunction of the molecule must remain unchanged. However, inspection of the form of the rotational wavefunctions (which have the same angular dependence as the $s, p$, etc. orbitals of atoms) shows that they change sign by $(-1)^{J}$ under such a rotation (Fig. 11B.13). Therefore, only even values of $J$ are permissible for $\mathrm{C}^{16} \mathrm{O}_{2}$, and hence the Raman spectrum shows only alternate lines.\\
\includegraphics[max width=\textwidth, center]{2025_02_27_d4b95d9718f0b9e2e6b9g-471(3)}

Figure 11B. 13 The symmetries of rotational wavefunctions (shown here, for simplicity as a two-dimensional rotor) under a rotation through $180^{\circ}$ depend on the value of $J$. Wavefunctions with $J$ even do not change sign; those with J odd do change sign.\\
\includegraphics[max width=\textwidth, center]{2025_02_27_d4b95d9718f0b9e2e6b9g-472}

Figure 11B. 14 The rotational Raman spectrum of a homonuclear diatomic molecule with two identical spin- $\frac{1}{2}$ nuclei shows an alternation in intensity as a result of nuclear statistics. In practice the Rayleigh line is much stronger than depicted in the figure.

The selective existence of rotational states that stems from the Pauli principle is termed nuclear statistics. Nuclear statistics must be taken into account whenever a rotation interchanges equivalent nuclei. However, the consequences are not always as simple as for $\mathrm{C}^{16} \mathrm{O}_{2}$ because there are complicating features when the nuclei have non-zero spin: it is found that there are several different relative nuclear spin orientations consistent with even values of $J$ and a different number of spin orientations consistent with odd values of $J$. For ${ }^{1} \mathrm{H}_{2}$ and ${ }^{19} \mathrm{~F}_{2}$, which have two identical spin- $\frac{1}{2}$ nuclei, by using the Pauli principle it can be shown that there are three times as many ways of achieving a state with odd $J$ than with even $J$, and there is a corresponding 3:1 alternation in intensity in their rotational Raman spectra (Fig. 11B.14).

\subsection*{How is that done? 11B. 2 \\
 Identifying the effect of nuclear statistics}
Because ${ }^{1} \mathrm{H}$ nuclei have $I=\frac{1}{2}$, like electrons, they are fermions and the Pauli principle requires the overall wavefunction to change sign under particle interchange. However, the rotation of a ${ }^{1} \mathrm{H}_{2}$ molecule through $180^{\circ}$ has a more complicated effect than simply relabelling the nuclei (Fig. 11B.15).

There are four nuclear spin wavefunctions: three correspond to a total nuclear spin $I_{\text {total }}=1$ (parallel spins, $\uparrow \uparrow$ ); and one with $I_{\text {total }}=0$ (paired spins, $\uparrow \downarrow$ ). The three wavefunctions with $I_{\text {total }}=1$ are $\alpha(\mathrm{A}) \alpha(\mathrm{B}), \alpha(\mathrm{A}) \beta(\mathrm{B})+\alpha(\mathrm{B}) \beta(\mathrm{A})$, and $\beta(\mathrm{A}) \beta(\mathrm{B})$ with $M_{I}=+1,0$, and -1 , respectively. Rotation of the molecule through $180^{\circ}$ interchanges the labels A and B , but overall these three wavefunctions are unchanged. Therefore, to achieve an overall change of sign, the rotational wavefunction must change sign, and so only odd values of $J$ are allowed.\\
\includegraphics[max width=\textwidth, center]{2025_02_27_d4b95d9718f0b9e2e6b9g-472(1)}

Figure 11B. 15 The interchange of two identical fermion nuclei results in the change in sign of the overall wavefunction. The relabelling can be thought of as occurring in two steps: the first is a rotation of the molecule; the second is the interchange of unlike spins (represented by the different colours of the nuclei). The wavefunction changes sign in the second step if the nuclei have antiparallel spins.

The fourth wavefunction, with $I_{\text {total }}=0$ and $M_{I}=0$, is $\alpha(A) \beta(B)-\alpha(B) \beta(A)$. When the labels $A$ and $B$ are interchanged the nuclear spin wavefunction changes sign: $\alpha(A) \beta(B)-$ $\alpha(\mathrm{B}) \beta(\mathrm{A}) \rightarrow \alpha(\mathrm{B}) \beta(\mathrm{A})-\alpha(\mathrm{A}) \beta(\mathrm{B}) \equiv-\{\alpha(\mathrm{A}) \beta(\mathrm{B})-\alpha(\mathrm{B}) \beta(\mathrm{A})\}$. Therefore, in this case for the overall wavefunction to change sign requires that the rotational wavefunction not change sign. Hence, only even values of $J$ are allowed.

The analysis leads to the conclusion that there are three nuclear spin wavefunctions that can be combined with odd values of $J$, and one wavefunction that can be combined with even values of $J$. In accord with the prediction of eqn 11B.25, the ratio of the number of ways of achieving odd $J$ to even $J$ is 3:1. In general, for a homonuclear diatomic molecule with nuclei of spin $I$, the numbers of ways of achieving states of odd and even $J$ are in the ratio

$$
\frac{\text { Number of ways of achieving odd } J}{\text { Number of ways of achieving even } J}
$$

\[
=\left\{\begin{array}{l}
(I+1) / I \text { for half-integral spin nuclei } \\
I /(I+1) \text { for integral spin nuclei }  \tag{11B.25}\\
\text { Nuclear statistics: homonuclear diatomics }
\end{array}\right.
\]

For ${ }^{1} \mathrm{H}_{2}, I=\frac{1}{2}$ and the ratio is $3: 1$. For ${ }^{14} \mathrm{~N}_{2}$, with $I=1$ the ratio is $1: 2$. Additional complications arise when the electronic state of the molecule is not totally symmetric (as for $\mathrm{O}_{2}$, Topic 11F).

Nuclear statistics have consequences outside spectroscopy. Different relative nuclear spin orientations change into one another only very slowly, so a ${ }^{1} \mathrm{H}_{2}$ molecule with parallel nuclear spins remains distinct from one with paired nuclear spins for long periods. The form with parallel nuclear spins is called ortho-hydrogen and the form with paired nuclear spins is called para-hydrogen. Because ortho-hydrogen cannot exist in a state with $J=0$, it continues to rotate at very low temperatures and has an effective rotational zero-point energy.

\subsection*{Checklist of concepts}
1. A rigid rotor is a body that does not distort under the stress of rotation.2. Rigid rotors are classified as spherical, symmetric, linear, or asymmetric by noting the number of equal principal moments of inertia (or their symmetry).3. Symmetric rotors are classified as prolate or oblate.4. Centrifugal distortion arises from forces that change the geometry of a molecule.5. The gross selection rule for a molecule to give a pure rotational spectrum is that it must be polar.6. The specific selection rules for microwave spectroscopy are $\Delta J= \pm 1, \Delta M_{J}=0, \pm 1$; for symmetric rotors the additional rule $\Delta K=0$ also applies.7. A molecule must be anisotropically polarizable for it to give rise to rotational Raman scattering.8. The specific selection rules for rotational Raman spectroscopy are: (i) linear rotors, $\Delta J=0, \pm 2$; (ii) symmetric rotors, $\Delta J=0, \pm 1, \pm 2 ; \Delta K=0$.9. The appearance of rotational spectra is affected by nuclear statistics, the selective occupation of rotational states that stems from the Pauli principle.\subsection*{Checklist of equations}
\begin{center}
\begin{tabular}{|c|c|c|c|}
\hline
Property & Equation & Comment & Equation number \\
\hline
Moment of inertia & \( I=\sum_{i} m_{i} x_{i}^{2} \) & $x_{i}$ is the perpendicular distance of atom $i$ from the axis of rotation & 11B. 2 \\
\hline
Rotational terms of a spherical or linear rotor & $\tilde{F}(J)=\tilde{B} J(J+1)$ & \( \begin{aligned} & J=0,1,2, \ldots \\ & \tilde{B}=\hbar / 4 \pi c I \end{aligned} \) & 11B.9, 11B. 14 \\
\hline
Rotational terms of a symmetric rotor & $\tilde{F}(J, K)=\tilde{B} J(J+1)+(\tilde{A}-\tilde{B}) K^{2}$ & \( \begin{aligned} & J=0,1,2, \ldots \\ & K=0, \pm 1, \ldots, \pm J \\ & \tilde{A}=\hbar / 4 \pi c I_{\|} \\ & \tilde{B}=\hbar / 4 \pi c I_{\perp} \end{aligned} \) & \begin{tabular}{l}
11B.13a \\
11B.13b \\
\end{tabular} \\
\hline
Centrifugal distortion & $\tilde{F}(J)=\tilde{B} J(J+1)-\tilde{D}_{J} J^{2}(J+1)^{2}$ & Spherical or linear rotor & 11B. 15 \\
\hline
Centrifugal distortion constant & $\tilde{D}_{J}=4 \tilde{B}^{3} / \tilde{v}^{2}$ &  & 11B. 16 \\
\hline
Wavenumbers of rotational transitions & $\tilde{v}(J+1 \leftarrow J)=2 \tilde{B}(J+1)$ & $J=0,1,2, \ldots$; linear rigid rotors & 11B.20a \\
\hline
Rotational state with largest population & $J_{\text {max }} \approx(k T / 2 h c \tilde{B})^{1 / 2}-\frac{1}{2}$ & Linear rotors & 11B. 21 \\
\hline
Wavenumbers of (i) Stokes and (ii) anti-Stokes lines in the rotational Raman spectrum of linear rotors & \begin{tabular}{l}
(i) $\tilde{v}(J+2 \leftarrow J)=\tilde{v}_{\mathrm{i}}-2 \tilde{B}(2 J+3)$ \\
(ii) $\tilde{v}(J-2 \leftarrow J)=\tilde{v}_{\mathrm{i}}+2 \tilde{B}(2 J-1)$ \\
\end{tabular} & \begin{tabular}{l}
(i) $J=0,1,2, \ldots$ \\
(ii) $J=2,3,4, \ldots$ \\
\end{tabular} & \( \begin{aligned} & \text { 11B.24a } \\ & \text { 11B.24b } \end{aligned} \) \\
\hline
\end{tabular}
\end{center}

\section*{TOPIC 11C Vibrational spectroscopy of diatomic molecules }
\subsection*{Why do you need to know this material?}
The observation of vibrational transition frequencies is used to determine the strengths and rigidities of bonds. Measurements in the gas phase can also be used to measure the bond lengths of diatomic molecules.

\subsection*{> What is the key idea?}
The vibrational spectrum of a diatomic molecule can be interpreted by using the harmonic oscillator model, with modifications that account for bond dissociation and the coupling of rotational and vibrational motion.

\subsection*{What do you need to know already?}
You need to be familiar with the harmonic oscillator (Topic 7E) and rigid rotor (Topic 11B) models of molecular motion and the general principles of spectroscopy (Topic 11A).

One internal mode of motion of a diatomic molecule is its vibration, in which the internuclear separation increases and decreases periodically. This motion, and the transitions between the allowed quantum states, can be treated initially as an example of harmonic motion like that described in Topic 7E.

\subsection*{11C. 1 Vibrational motion}
Figure 11C. 1 shows a typical potential energy curve of a diatomic molecule (it is essentially a reproduction of Fig. 7E. 1 of Topic 7E). The potential energy $V(x)$, where $x=R-R_{e}$ (the displacement from equilibrium), can be expanded around its minimum by using a Taylor series (see The chemist's toolkit 12 in Topic 5B):


\begin{equation*}
V(x)=V(0)+\left(\frac{\mathrm{d} V}{\mathrm{~d} x}\right)_{0} x+\frac{1}{2}\left(\frac{\mathrm{~d}^{2} V}{\mathrm{~d} x^{2}}\right)_{0} x^{2}+\cdots \tag{11C.1a}
\end{equation*}


The notation (...) $)_{0}$ means that the derivative is evaluated at $x=0$. The term $V(0)$ can be set arbitrarily to zero, and the first\\
\includegraphics[max width=\textwidth, center]{2025_02_27_d4b95d9718f0b9e2e6b9g-474}

Figure 11C. 1 A molecular potential energy curve can be approximated by a parabola near the bottom of the well. The parabolic potential energy results in harmonic oscillations. At high excitation energies the parabolic approximation is poor (the true potential energy is less confining), and is totally wrong near the dissociation limit.\\
derivative of $V$ is zero at the minimum. Therefore, the first surviving term is proportional to the square of the displacement. For small displacements all the higher terms can be ignored so the potential energy can be written


\begin{equation*}
V(x) \approx \frac{1}{2}\left(\frac{\mathrm{~d}^{2} V}{\mathrm{~d} x^{2}}\right)_{0} x^{2} \tag{11C.1b}
\end{equation*}


Therefore, the first approximation to a molecular potential energy curve is a parabolic potential of the form


\begin{equation*}
V(x)=\frac{1}{2} k_{\mathrm{f}} x^{2} \quad x=R-R_{\mathrm{e}} \quad \text { Parabolic potential energy } \tag{11C.2a}
\end{equation*}


where $k_{\mathrm{f}}$ is the force constant of the bond, a measure of its stiffness:

$$
k_{\mathrm{f}}=\left(\frac{\mathrm{d}^{2} V}{\mathrm{~d} x^{2}}\right)_{0}
$$

$$
\begin{aligned}
& \text { Force constant } \\
& \text { [definition] }
\end{aligned}
$$

(11c.2b)

If the potential energy curve is sharply curved close to its minimum, then $k_{\mathrm{f}}$ will be large and the bond stiff. Conversely, if the potential energy curve is wide and shallow, then $k_{\mathrm{f}}$ will be small and the bond easily stretched or compressed (Fig. 11C.2).\\
\includegraphics[max width=\textwidth, center]{2025_02_27_d4b95d9718f0b9e2e6b9g-475}

Figure 11C. 2 The force constant is a measure of the curvature of the potential energy close to the equilibrium extension of the bond. A strongly confining well (one with steep sides, a stiff bond) corresponds to high values of $k_{\mathrm{f}}$.

The Schrödinger equation for the relative motion of two atoms of masses $m_{1}$ and $m_{2}$ with a parabolic potential energy is


\begin{equation*}
-\frac{\hbar^{2}}{2 m_{\mathrm{eff}}} \frac{\mathrm{~d}^{2} \psi}{\mathrm{~d} x^{2}}+\frac{1}{2} k_{\mathrm{f}} x^{2} \psi=E \psi \tag{11C.3a}
\end{equation*}


where $m_{\text {eff }}$ is the effective mass:

$$
m_{\mathrm{eff}}=\frac{m_{1} m_{2}}{m_{1}+m_{2}}
$$

Effective mass\\[0pt]
[definition]\\
(11C.3b)

These equations are derived by using the separation of variables procedure (see A deeper look 3 on the website for this text) to separate the relative motion of the atoms from the motion of the molecule as a whole.

A note on good practice Distinguish effective mass from reduced mass. The former is a measure of the mass that is moved during a vibration. The latter is the quantity that emerges from the separation of relative internal and overall translational motion. For a diatomic molecule the two are the same, but that is not true in general for vibrations of polyatomic molecules. Many, however, do not make this distinction and refer to both quantities as the 'reduced mass'.

Apart from the appearance of the effective mass, the Schrödinger equation in eqn 11C.3a is the same as eqn 7E. 2 for a particle of mass $m$ undergoing harmonic motion. Therefore, the results from that Topic can be used to write the permitted vibrational energy levels:

$$
E_{v}=\left(v+\frac{1}{2}\right) \hbar \omega \quad \omega=\left(\frac{k_{\mathrm{f}}}{m_{\text {eff }}}\right)^{1 / 2} \quad v=0,1,2, \ldots
$$

Vibrational energy levels [diatomic molecule]\\
(11C.4a)\\
The vibrational terms of a molecule, the energies of its vibrational states expressed as wavenumbers, are denoted\\
$\tilde{G}(v)$, with $E_{v}=h c \tilde{G}(v)$. Therefore, with $\omega=2 \pi v$ and $\tilde{v}=v / c=$ $\omega / 2 \pi c$ :

$$
\tilde{G}(v)=\left(v+\frac{1}{2}\right) \tilde{v} \quad \tilde{v}=\frac{1}{2 \pi c}\left(\frac{k_{\mathrm{f}}}{m_{\mathrm{eff}}}\right)^{1 / 2}
$$

Vibrational terms [diatomic]\\
(11C.4b)\\
The vibrational wavefunctions are the same as those discussed in Topic 7E for a harmonic oscillator.

The vibrational terms depend on the effective mass of the molecule, not directly on its total mass. This dependence is physically reasonable, because if atom 1 is very much heavier than atom 2, then the effective mass is close to $m_{2}$, the mass of the lighter atom. The vibration would then be that of the light atom relative to an essentially stationary heavy atom. For a homonuclear diatomic molecule $m_{1}=m_{2}$, and the effective mass is half the total mass: $m_{\text {eff }}=\frac{1}{2} m$.

\subsection*{Brief illustration 11C. 1}
The force constant of the bond in HCl is $516 \mathrm{Nm}^{-1}$, a reasonably typical value for a single bond. The effective mass of ${ }^{1} \mathrm{H}^{35} \mathrm{Cl}$ is $1.63 \times 10^{-27} \mathrm{~kg}$ (note that this mass is very close to the mass of the hydrogen atom, $1.67 \times 10^{-27} \mathrm{~kg}$, implying that the H atom is essentially vibrating against a stationary Cl atom). The vibrational frequency is therefore

$$
\omega=\left(\frac{516 \stackrel{\sim}{\mathrm{~N} \mathrm{~m} \mathrm{~m}}^{-1}}{1.63 \times 10^{-27} \mathrm{~kg}}\right)^{1 / 2}=5.63 \times 10^{14} \mathrm{~s}^{-1}
$$

and the corresponding wavenumber is

$$
\tilde{v}=\frac{\omega}{2 \pi c}=\frac{5.63 \times 10^{14} \mathrm{~s}^{-1}}{2 \pi \times\left(2.998 \times 10^{10} \mathrm{~cm} \mathrm{~s}^{-1}\right)}=2.99 \times 10^{3} \mathrm{~cm}^{-1}
$$

\subsection*{11c. 2 Infrared spectroscopy}
The gross and specific selection rules for vibrational transitions are established, as usual, by considering the properties of the electric transition dipole moment. The detailed calculation is shown in A deeper look 6 on the website of this text. The conclusion is that

\begin{displayquote}
The gross selection rule for a change in vibrational state brought about by absorption or emission of radiation is that the electric dipole moment of the molecule must change when the atoms are displaced relative to one another.
\end{displayquote}

Such vibrations are said to be infrared active. The classical basis of this rule is that an oscillating electric dipole generates\\
an electromagnetic wave, and an oscillating electric field of such a wave generates an oscillating electric dipole.

Note that the molecule need not have a permanent dipole moment: the rule requires only a change in dipole moment. Some vibrations do not affect the dipole moment of the molecule (for instance, the stretching motion of a homonuclear diatomic molecule), so they neither absorb nor generate radiation: such vibrations are said to be infrared inactive. Weak infrared transitions can be observed from homonuclear diatomic molecules trapped within various nanomaterials. For instance, when incorporated into solid $\mathrm{C}_{60}, \mathrm{H}_{2}$ molecules interact through van der Waals forces with the surrounding $\mathrm{C}_{60}$ molecules and acquire dipole moments, with the result that they have observable infrared spectra.

The calculation also shows that the specific selection rule is

$$
\Delta v= \pm 1 \quad \text { Specific selection rule [harmonic oscillator] (11C.5) }
$$

Transitions for which $\Delta v=+1$ correspond to absorption and those with $\Delta v=-1$ correspond to emission. It follows that the wavenumbers of allowed vibrational transitions, which are denoted $\Delta \tilde{G}_{v+\frac{1}{2}}$ for the transition $v+1 \leftarrow v$, are

$$
\Delta \tilde{G}_{v+\frac{1}{2}}=\tilde{G}(v+1)-\tilde{G}(v)=\tilde{v} \quad \text { Harmonic oscillator }
$$

(11C.6)

The wavenumbers of vibrational transitions correspond to radiation in the infrared region of the electromagnetic spectrum, so vibrational transitions absorb and generate infrared radiation.

At room temperature $k T / h c \approx 200 \mathrm{~cm}^{-1}$, and because most vibrational wavenumbers are significantly greater than $200 \mathrm{~cm}^{-1}$ it follows from the Boltzmann distribution that almost all the molecules are in their vibrational ground states. Hence, the dominant spectral transition will be the fundamental transition, $1 \leftarrow 0$. As a result, the spectrum is expected to consist of a single absorption line. If the molecules are formed in a vibrationally excited state, such as when vibrationally excited HF molecules are formed in the reaction $\mathrm{H}_{2}+\mathrm{F}_{2} \rightarrow 2 \mathrm{HF}^{*}$, where the star indicates a vibrationally 'hot' molecule, the transitions $5 \rightarrow 4,4 \rightarrow 3, \ldots$ may also appear (in emission). In the harmonic approximation, all these lines lie at the same frequency, and the spectrum is also a single line. However, the breakdown of the harmonic approximation causes the transitions to lie at slightly different frequencies, so several lines are observed.

\subsection*{11C. 3 Anharmonicity}
The vibrational terms in eqn 11C.4b are only approximate because they are based on a parabolic approximation to the actual potential energy curve. A parabola cannot be correct at all extensions because it does not allow a bond to dissociate.

At high vibrational excitations the separation of the atoms (more precisely, the spread of the vibrational wavefunction) allows the molecule to explore regions of the potential energy curve where the parabolic approximation is poor and additional terms in the Taylor expansion of $V$ (eqn 11C.1a) must be retained. The motion then becomes anharmonic, in the sense that the restoring force is no longer proportional to the displacement. Because the actual curve is less confining than a parabola, it can be anticipated that the energy levels become more closely spaced at high excitations.

\subsection*{(a) The convergence of energy levels}
One approach to the calculation of the energy levels in the presence of anharmonicity is to use a function that resembles the true potential energy more closely. The Morse potential energy is

\[
V(x)=h c \tilde{D}_{\mathrm{e}}\left\{1-\mathrm{e}^{-a x}\right\}^{2} \quad a=\left(\frac{m_{\mathrm{eff}} \omega^{2}}{2 h c \tilde{D}_{\mathrm{e}}}\right)^{1 / 2} \quad \begin{align*}
& \text { Morse }  \tag{11C.7}\\
& \text { potential } \\
& \text { energy }
\end{align*}
\]

At $x=0, V(0)=0$; at large displacements $V(x)$ approaches $h c \tilde{D}_{e}$ (Fig. 11C.3). Near the well minimum the variation of $V$ with displacement resembles a parabola (as can be checked by expanding the exponential and retaining the first two terms). The Schrödinger equation can be solved for the Morse potential energy and the permitted levels are

\[
\begin{array}{rlrl}
\tilde{G}(v) & =\left(v+\frac{1}{2}\right) \tilde{v}-\left(v+\frac{1}{2}\right)^{2} x_{\mathrm{e}} \tilde{v} & & \\
v & =0,1,2, \ldots, v_{\max } & \begin{array}{l}
\text { Vibrational terms } \\
\text { [Morse potential energy] }
\end{array}  \tag{11C.8}\\
x_{\mathrm{e}} & =\frac{a^{2} \hbar}{2 m_{\mathrm{eff}} \omega}=\frac{\tilde{v}}{4 \tilde{D}_{\mathrm{e}}} & &
\end{array}
\]

The positive dimensionless parameter $x_{\mathrm{e}}$ is called the anharmonicity constant. The number of vibrational levels\\
\includegraphics[max width=\textwidth, center]{2025_02_27_d4b95d9718f0b9e2e6b9g-476}

Figure 11C. 3 The dissociation energy of a molecule, $h c \tilde{D}_{0}$, differs from the depth of the potential well, $h c \tilde{D}_{e}$, on account of the zeropoint energy of the vibration of the bond.\\
\includegraphics[max width=\textwidth, center]{2025_02_27_d4b95d9718f0b9e2e6b9g-477(1)}

Figure 11C. 4 The Morse potential energy curve reproduces the general shape of a molecular potential energy curve. The corresponding Schrödinger equation can be solved, and the values of the energies obtained. The number of bound levels is finite.\\
of a Morse oscillator is finite, as shown in Fig. 11C. 4 (see also Problem P11C.8). The second term in the expression for $\tilde{G}$ subtracts from the first with increasing effect as $v$ increases, and hence gives rise to the convergence of the levels at high quantum numbers. In addition to the depth of the well, $h c \tilde{D}_{\mathrm{e}}$, the dissociation energy $h c \tilde{D}_{0}$ is the energy difference between the lowest vibrational state $(v=0)$ and the infinitely separated atoms. As can be seen in Fig. 11C.3, the two quantities are related by $\tilde{D}_{\mathrm{e}}=\tilde{D}_{0}+\tilde{G}(0)$.

Although the Morse oscillator is quite useful theoretically, in practice the more general expression


\begin{equation*}
\tilde{G}(v)=\left(v+\frac{1}{2}\right) \tilde{v}-\left(v+\frac{1}{2}\right)^{2} x_{e} \tilde{v}+\left(v+\frac{1}{2}\right)^{3} y_{\mathrm{e}} \tilde{v}+\cdots \tag{11C.9a}
\end{equation*}


where $x_{\mathrm{e}}, y_{\mathrm{e}}, \ldots$ are empirical dimensionless constants characteristic of the molecule, is used to fit the experimental data and to determine the dissociation energy of the molecule. In this case the wavenumbers of transitions with $\Delta v=+1$ are


\begin{equation*}
\Delta \tilde{G}_{v+\frac{1}{2}}=\tilde{G}(v+1)-\tilde{G}(v)=\tilde{v}-2(v+1) x_{e} \tilde{v}+\cdots \tag{11C.9b}
\end{equation*}


Equation 11C.9b shows that, because $x_{\mathrm{e}}>0$, the transitions move to lower wavenumbers as $v$ increases.

In addition to the strong fundamental transition $1 \leftarrow 0$, a set of weaker absorption lines are also seen and correspond to the transitions $2 \leftarrow 0,3 \leftarrow 0, \ldots$. These transitions are forbidden for a harmonic oscillator, but become weakly allowed as a result of anharmonicity. The transition $2 \leftarrow 0$ is known as the first overtone, $3 \leftarrow 0$ is the second overtone, and so on. The wavenumber of the first overtone is given by


\begin{equation*}
\tilde{G}(v+2)-\tilde{G}(v)=2 \tilde{v}-2(2 v+3) x_{\mathrm{e}} \tilde{v}+\cdots \tag{11C.10}
\end{equation*}


The reason for the appearance of overtones is that the selection rule is derived from the properties of harmonic oscillator wavefunctions, which are only approximately valid when anharmonicity is present. Therefore, the selection rule is also only an approximation. For an anharmonic oscillator, all values of $\Delta v$ are allowed, but transitions with $\Delta v>1$ are allowed only weakly if the anharmonicity is slight. Typically, the first overtone is only about one-tenth as intense as the fundamental.

\subsection*{Example 11C. 1 Estimating an anharmonicity constant}
Estimate the anharmonicity constant $x_{\mathrm{e}}$ for ${ }^{35} \mathrm{Cl}^{19} \mathrm{~F}$ given that the wavenumbers of the fundamental and first overtones are found to be 773.8 and $1535.3 \mathrm{~cm}^{-1}$, respectively.

Collect your thoughts You can find an expression for the wavenumber of the fundamental transition $1 \leftarrow 0$, by using eqn 11C.9b with $v=0$, and for the wavenumber of the first overtone $2 \leftarrow 0$ by using eqn 11 C .10 with $v=0$. You then need to solve the two equations to give values for $\tilde{v}$ and $x_{e} \tilde{v}$, and hence find $x_{\mathrm{e}}$ itself.

The solution From eqn 11C.9b the expression for the wavenumber of the fundamental is $\tilde{v}-2 x_{\mathrm{e}} \tilde{v}$, and from eqn 11C. 10 the expression for the first overtone is $2 \tilde{v}-6 x_{e} \tilde{v}$. From the data it follows that $773.8 \mathrm{~cm}^{-1}=\tilde{v}-2 x_{\mathrm{e}} \tilde{v}$ and $1535.3 \mathrm{~cm}^{-1}=$ $2 \tilde{v}-6 x_{e} \tilde{v}$. The terms in $\tilde{v}$ are eliminated by noting that $\left(\tilde{v}-2 x_{\mathrm{e}} \tilde{v}\right)-\frac{1}{2}\left(2 \tilde{v}-6 x_{\mathrm{e}} \tilde{v}\right)=x_{\mathrm{e}} \tilde{v}$ to give $x_{\mathrm{e}} \tilde{v}=773.8 \mathrm{~cm}^{-1}-$ $\frac{1}{2} \times 1535.3 \mathrm{~cm}^{-1}=6.15 \mathrm{~cm}^{-1}$. This value for $x_{\mathrm{e}} \tilde{v}$ can then be substituted into $773.8 \mathrm{~cm}^{-1}=\tilde{v}-2 x_{\mathrm{e}} \tilde{v}$ to give $\tilde{v}=773.8 \mathrm{~cm}^{-1}+2 x_{\mathrm{e}} \tilde{v}=$ $773.8 \mathrm{~cm}^{-1}+2 \times 6.15 \mathrm{~cm}^{-1}=786.1 \mathrm{~cm}^{-1}$. It follows that

$$
x_{\mathrm{e}}=\frac{x_{\mathrm{e}} \tilde{v}}{\tilde{v}}=\frac{6.15 \mathrm{~cm}^{-1}}{786.1 \mathrm{~cm}^{-1}}=7.82 \times 10^{-3}
$$

Self-test 11C. 1 Predict the wavenumber of the second overtone for this molecule.\\
\includegraphics[max width=\textwidth, center]{2025_02_27_d4b95d9718f0b9e2e6b9g-477}

\subsection*{(b) The Birge-Sponer plot}
When several vibrational transitions are detectable, a graphical technique called a Birge-Sponer plot can be used to determine the dissociation energy of the bond. The basis of the Birge-Sponer plot is that the sum of the successive intervals $\Delta \tilde{G}_{v+\frac{1}{2}}$ (eqn 11C.9b) from $v=0$ to the dissociation limit is the dissociation wavenumber $\tilde{D}_{0}$ :


\begin{equation*}
\tilde{D}_{0}=\Delta \tilde{G}_{1 / 2}+\Delta \tilde{G}_{3 / 2}+\cdots=\sum_{v} \Delta \tilde{G}_{v+\frac{1}{2}} \tag{11C.11}
\end{equation*}


just as the height of a ladder is the sum of the separations of its rungs (Fig. 11C.5). The construction in Fig. 11C. 6 shows that the area under the plot of $\Delta \tilde{G}_{v+\frac{1}{2}}$ against $v+\frac{1}{2}$ is equal to the\\
\includegraphics[max width=\textwidth, center]{2025_02_27_d4b95d9718f0b9e2e6b9g-478(1)}

Figure 11C. 5 The dissociation wavenumber is the sum of the separations $\Delta \tilde{G}_{v+\frac{1}{2}}$ of the vibrational terms up to the dissociation limit, just as the height of a ladder is the sum of the separations of its rungs.\\
\includegraphics[max width=\textwidth, center]{2025_02_27_d4b95d9718f0b9e2e6b9g-478(2)}

Figure 11C. 6 The area under a plot of $\Delta \tilde{G}_{v+\frac{1}{2}}$ against vibrational quantum number is equal to the dissociation wavenumber of the molecule. The assumption that the differences approach zero linearly is the basis of the Birge-Sponer extrapolation.\\
sum, and therefore to $\tilde{D}_{0}$. The successive terms decrease linearly when only the $x_{\mathrm{e}}$ anharmonicity constant is taken into account and the inaccessible part of the spectrum can be estimated by linear extrapolation. Most actual plots differ from the linear plot as shown in Fig. 11C.6, so the value of $\tilde{D}_{0}$ obtained in this way is usually an overestimate of the true value.

\subsection*{Example 11C. 2}
\subsection*{Using a Birge-Sponer plot}
The observed vibrational intervals of $\mathrm{H}_{2}^{+}$lie at the following values for $1 \leftarrow 0,2 \leftarrow 1, \ldots$, respectively (in $\mathrm{cm}^{-1}$ ): 2191, 2064, 1941, 1821, 1705, 1591, 1479, 1368, 1257, 1145, 1033, 918, 800, 677, 548, 411. Determine the dissociation energy of the molecule.

Collect your thoughts Plot the separations against $v+\frac{1}{2}$, extrapolate linearly to the point cutting the horizontal axis, and then measure the area under the curve.

The solution The points are plotted in Fig. 11C.7, and a linear extrapolation is shown as a green line. The area under the curve (use the formula for the area of a triangle or count the\\
squares) is 216 . Each square corresponds to $100 \mathrm{~cm}^{-1}$ (refer to the scale of the vertical axis); hence the dissociation energy, expressed as a wavenumber, is $21600 \mathrm{~cm}^{-1}$ (corresponding to $258 \mathrm{~kJ} \mathrm{~mol}^{-1}$ ).\\
\includegraphics[max width=\textwidth, center]{2025_02_27_d4b95d9718f0b9e2e6b9g-478(4)}

Figure 11C. 7 The Birge-Sponer plot used in Example 11C.2. The area is obtained simply by counting the squares beneath the line or using the formula for the area of a triangle (area $=$ $\frac{1}{2} \times$ base $\times$ height).

Self-test 11C. 2 The vibrational levels of HgH converge rapidly, and successive intervals are 1203.7 (which corresponds to the $1 \leftarrow 0$ transition), $965.6,632.4$, and $172 \mathrm{~cm}^{-1}$. Estimate the molar dissociation energy.\\
\includegraphics[max width=\textwidth, center]{2025_02_27_d4b95d9718f0b9e2e6b9g-478}

\subsection*{11c.4 Vibration-rotation spectra}
Each line of the high resolution vibrational spectrum of a gasphase heteronuclear diatomic molecule is found to consist of a large number of closely spaced components (Fig. 11C.8). Hence, molecular spectra are often called band spectra. The separation between the components is less than $10 \mathrm{~cm}^{-1}$, which\\
\includegraphics[max width=\textwidth, center]{2025_02_27_d4b95d9718f0b9e2e6b9g-478(3)}

Figure 11C. 8 A high-resolution vibration-rotation spectrum of HCl . The lines appear in pairs because $\mathrm{H}^{35} \mathrm{Cl}$ and $\mathrm{H}^{37} \mathrm{Cl}$ both contribute (their abundance ratio is 3:1). There is no Q branch (see below), because $\Delta J=0$ is forbidden for this molecule.\\
suggests that the structure is due to rotational transitions accompanying the vibrational transition. A rotational change should be expected because classically a vibrational transition can be thought of as leading to a sudden increase or decrease in the instantaneous bond length. Just as ice-skaters rotate more rapidly when they bring their arms in, and more slowly when they throw them out, so the molecular rotation is either accelerated or retarded by a vibrational transition.

\subsection*{(a) Spectral branches}
A detailed analysis of the quantum mechanics of simultaneous vibrational and rotational changes shows that the rotational quantum number $J$ changes by $\pm 1$ during the vibrational transition of a diatomic molecule. If the molecule also possesses angular momentum about its axis, as in the case of the electronic orbital angular momentum of the molecule NO with its configuration $\ldots \pi^{1}$, then the selection rules also allow $\Delta J=0$.

The appearance of the vibration-rotation spectrum of a diatomic molecule can be discussed by using the combined vibration-rotation terms, $\tilde{S}$ :


\begin{equation*}
\tilde{S}(v, J)=\tilde{G}(v)+\tilde{F}(J) \tag{11C.12a}
\end{equation*}


If anharmonicity and centrifugal distortion are ignored, $\tilde{G}(v)$ can be replaced by the expression in eqn 11C.4b, and $\tilde{F}(J)$ can be replaced by the expression in eqn 11B. $9(\tilde{F}(J)=\tilde{B} J(J+1))$ to give

$$
\tilde{S}(v, J)=\left(v+\frac{1}{2}\right) \tilde{v}+\tilde{B} J(J+1)
$$

(11C.12b)

In a more detailed treatment, $\tilde{B}$ is allowed to depend on the vibrational state and written $\tilde{B}_{v}$.

In the vibrational transition $v+1 \leftarrow v$, J changes by $\pm 1$ and in some cases by 0 (when $\Delta J=0$ is allowed). The absorptions then fall into three groups called branches of the spectrum. The P branch consists of all transitions with $\Delta J=-1$ :

$$
\begin{aligned}
\tilde{v}_{\mathrm{P}}(J) & =\tilde{S}(v+1, J-1)-\tilde{S}(v, J)=\tilde{v}-2 \tilde{B} J \\
J & =1,2,3, \ldots
\end{aligned}
$$

(11C.13a)

This branch consists of lines extending to the low wavenumber side of $\tilde{v}$ at $\tilde{v}-2 \tilde{B}, \tilde{v}-4 \tilde{B}, \ldots$ with an intensity distribution reflecting both the populations of the rotational levels and the magnitude of the $J-1 \leftarrow J$ transition moment (Fig. 11C.9). The Q branch consists of all transitions with $\Delta J=0$, and its wavenumbers are the same for all values of $J$ :

$$
\tilde{v}_{\mathrm{Q}}(J)=\tilde{S}(v+1, J)-\tilde{S}(v, J)=\tilde{v} \quad \text { Q branch transitions } \quad \text { (11C.13b) }
$$

This branch, when it is allowed (as in NO), appears at the vibrational transition wavenumber $\tilde{v}$. In Fig. 11C. 8 there is a gap at the expected location of the Q branch because it is\\
\includegraphics[max width=\textwidth, center]{2025_02_27_d4b95d9718f0b9e2e6b9g-479}

Figure 11C. 9 The formation of $P, Q$, and $R$ branches in a vibrationrotation spectrum. The intensities reflect the populations of the initial rotational levels and magnitudes of the transition moments.\\
forbidden in HCl because it has zero electronic angular momentum around its internuclear axis. The $\mathbf{R}$ branch consists of lines with $\Delta J=+1$ :


\begin{align*}
\tilde{V}_{\mathrm{R}}(J) & =\tilde{S}(v+1, J+1)-\tilde{S}(v, J)=\tilde{v}+2 \tilde{B}(J+1) \\
J & =0,1,2, \ldots \quad \text { R branch transitions } \tag{11C.13c}
\end{align*}


This branch consists of lines extending to the high-wavenumber side of $\tilde{v}$ at $\tilde{v}+2 \tilde{B}, \tilde{v}+4 \tilde{B}, \ldots$.

The separation between the lines in the P and R branches of a vibrational transition gives the value of $\tilde{B}$. Therefore, the bond length can be deduced in the same way as from microwave spectra (Topic 11B). However, the latter technique gives more precise bond lengths because microwave frequencies can be measured with greater precision than infrared frequencies.

\subsection*{Brief illustration 11C. 2}
The infrared absorption spectrum of ${ }^{1} \mathrm{H}^{81} \mathrm{Br}$ contains a band arising from $v=0$. It follows from eqn 11C.13c and the data in Table 11C. 1 that the wavenumber of the line in the R branch originating from the rotational state with $J=2$ is

$$
\begin{aligned}
\tilde{V}_{\mathrm{R}}(2) & =\tilde{v}+6 \tilde{B}=2648.98 \mathrm{~cm}^{-1}+6 \times\left(8.465 \mathrm{~cm}^{-1}\right) \\
& =2699.77 \mathrm{~cm}^{-1}
\end{aligned}
$$

Table 11C. 1 Properties of diatomic molecules*

\begin{center}
\begin{tabular}{lcccll}
\hline
 & $\tilde{V} / \mathrm{cm}^{-1}$ & $R_{\mathrm{e}} / \mathrm{pm}$ & $\tilde{B} / \mathrm{cm}^{-1}$ & $\boldsymbol{k}_{\mathrm{f}} /\left(\mathrm{N} \mathrm{m}^{-1}\right)$ & $\tilde{D}_{0} /\left(10^{4} \mathrm{~cm}^{-1}\right)$ \\
\hline
${ }^{1} \mathrm{H}_{2}$ & 4400 & 74 & 60.86 & 575 & 3.61 \\
${ }^{1} \mathrm{H}^{35} \mathrm{Cl}$ & 2991 & 127 & 10.59 & 516 & 3.58 \\
${ }^{1} \mathrm{H}^{127} \mathrm{I}$ & 2308 & 161 & 6.51 & 314 & 2.46 \\
${ }^{35} \mathrm{Cl}_{2}$ & 560 & 199 & 0.244 & 323 & 2.00 \\
\hline
\end{tabular}
\end{center}

\begin{itemize}
  \item More values are given in the Resource section.
\end{itemize}

\subsection*{(b) Combination differences}
A more detailed analysis of the rotational fine structure shows that the rotational constant decreases as the vibrational quantum number $v$ increases. The origin of this effect is that the average value of $1 / R^{2}$ decreases because the asymmetry of the potential well results in the average bond length increasing with vibrational energy. A harmonic oscillator also shows this effect because although the average value of $R$ is unchanged with increasing $v$, the average value of $1 / R^{2}$ does change (see Problem P11C.13). Typically, $\tilde{B}_{1}$ is $1-2$ per cent smaller than $\tilde{B}_{0}$.\\
The result of $\tilde{B}_{1}$ being smaller than $\tilde{B}_{0}$ is that the Q branch (if it is present) consists of a series of closely spaced lines; the lines of the R branch converge slightly as $J$ increases, and those of the P branch diverge. It follows from eqn 11C.12b with $\tilde{B}_{v}$ in place of $\tilde{B}$


\begin{align*}
& \tilde{v}_{\mathrm{P}}(J)=\tilde{v}-\left(\tilde{B}_{1}+\tilde{B}_{0}\right) J+\left(\tilde{B}_{1}-\tilde{B}_{0}\right) J^{2} \\
& \tilde{v}_{\mathrm{Q}}(J)=\tilde{v}+\left(\tilde{B}_{1}-\tilde{B}_{0}\right) J(J+1)  \tag{11C.14}\\
& \tilde{v}_{\mathrm{R}}(J)=\tilde{v}+\left(\tilde{B}_{1}+\tilde{B}_{0}\right)(J+1)+\left(\tilde{B}_{1}-\tilde{B}_{0}\right)(J+1)^{2}
\end{align*}


To determine the two rotational constants individually, the method of combination differences is used, which involves setting up expressions for the difference in the wavenumbers of transitions to a common state. The resulting expression then depends solely on properties of the other states.

As can be seen from Fig. 11C.10, the transitions $\tilde{v}_{\mathrm{R}}(J-1)$ and $\tilde{v}_{\mathrm{p}}(J+1)$ have a common upper state, and hence the difference between these transitions can be anticipated to depend on $\tilde{B}_{0}$. From the diagram it can be seen that $\tilde{v}_{\mathrm{R}}(J-1)-\tilde{v}_{\mathrm{P}}(J+1)=\tilde{S}(0, J+1)-\tilde{S}(0, J-1)$. The right-hand side is evaluated by using the expression for $\tilde{S}(\nu, J)$ in eqn 11C.12b (with $\tilde{B}_{0}$ in place of $\tilde{B}$ ) to give


\begin{equation*}
\tilde{v}_{\mathrm{R}}(J-1)-\tilde{v}_{\mathrm{P}}(J+1)=4 \tilde{B}_{0}\left(J+\frac{1}{2}\right) \tag{11C.15a}
\end{equation*}


Therefore, a plot of the combination difference against $J+\frac{1}{2}$ should be a straight line of slope $4 \tilde{B}_{0}$ and intercept (with the vertical axis) zero; the value of $\tilde{B}_{0}$ can therefore be determined\\
\includegraphics[max width=\textwidth, center]{2025_02_27_d4b95d9718f0b9e2e6b9g-480}

Figure 11C. 10 The method of combination differences makes use of the fact that certain pairs of transitions share a common level.\\
from the slope. The presence of centrifugal distortion results in the intercept deviating from zero, but has little effect on the quality of the straight line.\\
The two lines $\tilde{v}_{\mathrm{R}}(J)$ and $\tilde{v}_{\mathrm{P}}(J)$ have a common lower state, and hence their combination difference depends on $\tilde{B}_{1}$. As before, from Fig. 11C. 10 it can be seen that $\tilde{v}_{\mathrm{R}}(J)-\tilde{v}_{\mathrm{P}}(J)=$ $\tilde{S}(1, J+1)-\tilde{S}(1, J-1)$ which is


\begin{equation*}
\tilde{v}_{\mathrm{R}}(J)-\tilde{v}_{\mathrm{p}}(J)=4 \tilde{B}_{1}\left(J+\frac{1}{2}\right) \tag{11C.15b}
\end{equation*}


\subsection*{Brief illustration 11C. 3}
The rotational constants of $\tilde{B}_{0}$ and $\tilde{B}_{1}$ can be estimated from a calculation involving only a few transitions. For ${ }^{1} \mathrm{H}^{35} \mathrm{Cl}$, $\tilde{V}_{\mathrm{R}}(0)-\tilde{v}_{\mathrm{P}}(2)=62.6 \mathrm{~cm}^{-1}$, and it follows from eqn 11C.15a, with $J=1$, that $\tilde{B}_{0}=62.6 / 4\left(1+\frac{1}{2}\right) \mathrm{cm}^{-1}=10.4 \mathrm{~cm}^{-1}$. Similarly, $\tilde{V}_{\mathrm{R}}(1)-\tilde{v}_{\mathrm{P}}(1)=60.8 \mathrm{~cm}^{-1}$, and it follows from eqn 11C.15b, again with $J=1$ that $\tilde{B}_{1}=60.8 / 4\left(1+\frac{1}{2}\right) \mathrm{cm}^{-1}=10.1 \mathrm{~cm}^{-1}$. If more lines are used to make combination difference plots, the values $\tilde{B}_{0}=10.440 \mathrm{~cm}^{-1}$ and $\tilde{B}_{1}=10.136 \mathrm{~cm}^{-1}$ are found. The two rotational constants differ by about 3 per cent of $\tilde{B}_{0}$.

\subsection*{11c. 5 Vibrational Raman spectra}
The gross and specific selection rules for vibrational Raman transitions are established, as usual, by considering the appropriate transition dipole moment. The details are set out in $A$ deeper look 6 on the website of this text. The conclusion is that the gross selection rule for vibrational Raman transitions is that the polarizability must change as the molecule vibrates. The polarizability plays a role in vibrational Raman spectroscopy because the molecule must be squeezed and stretched by the incident radiation in order for vibrational excitation to occur during the inelastic photon-molecule collision. Both homonuclear and heteronuclear diatomic molecules swell and contract during a vibration, the control of the nuclei over the electrons varies, and hence the molecular polarizability changes. Both types of diatomic molecule are therefore vibrationally Raman active. The analysis also shows that the specific selection rule for vibrational Raman transitions in the harmonic approximation is $\Delta v= \pm 1$, just as for infrared transitions.

The lines to high frequency of the incident radiation, in the language introduced in Topic 11A, the 'anti-Stokes lines', are those for which $\Delta v=-1$. The lines to low frequency, the 'Stokes lines', correspond to $\Delta v=+1$. The intensities of the anti-Stokes and Stokes lines are governed largely by the Boltzmann populations of the vibrational states involved in the transition. It follows that anti-Stokes lines are usually weak because the populations of the excited vibrational states are very small.\\
\includegraphics[max width=\textwidth, center]{2025_02_27_d4b95d9718f0b9e2e6b9g-481(1)}

Figure 11C. 11 The formation of $O, Q$, and $S$ branches in a vibration-rotation Raman spectrum of a diatomic molecule (its Stokes lines). Note that the frequency scale runs in the opposite direction to that in Fig. 11C.9, because the higher energy transitions (on the right) extract more energy from the incident beam and leave it at lower frequency.

In gas-phase spectra, the Stokes and anti-Stokes lines have a branch structure arising from the simultaneous rotational transitions that accompany the vibrational excitation (Fig. 11C.11). The selection rules are $\Delta J=0, \pm 2$ (as in pure rotational Raman spectroscopy), and give rise to the O branch $(\Delta J=-2)$, the $\mathbf{Q}$ branch $(\Delta J=0)$, and the S branch $(\Delta J=+2)$ :

$$
\begin{array}{lll}
\tilde{v}_{\mathrm{O}}(J)=\tilde{v}_{\mathrm{i}}-\tilde{v}+4 \tilde{B}\left(J-\frac{1}{2}\right) & J=2,3,4, \ldots & \begin{array}{l}
\text { O branch } \\
\text { transitions }
\end{array} \\
\tilde{v}_{\mathrm{Q}}(J)=\tilde{v}_{\mathrm{i}}-\tilde{v} & \begin{array}{l}
\text { Q branch } \\
\text { transitions }
\end{array} \\
\tilde{v}_{\mathrm{S}}(J)=\tilde{v}_{\mathrm{i}}-\tilde{v}+4 \tilde{B}\left(J-\frac{3}{2}\right) & J=0,1,2, \ldots & \begin{array}{l}
\text { S branch } \\
\text { transitions }
\end{array}
\end{array}
$$

\begin{center}
\includegraphics[max width=\textwidth]{2025_02_27_d4b95d9718f0b9e2e6b9g-481}
\end{center}

Figure 11C. 12 The structure of a vibrational line in the vibrational Raman spectrum (the Stokes lines) of carbon monoxide, showing the $\mathrm{O}, \mathrm{Q}$, and S branches. The horizontal axis represents the wavenumber difference between the incident and scattered radiation. For these Stokes lines the wavenumber of the scattered radiation (as distinct from the difference, which represents the energy deposited in the molecule) increases to the left, as in Fig. 11C. 11.\\
where $\tilde{v}_{\mathrm{i}}$ is the wavenumber of the incident radiation. Note that, unlike in infrared spectroscopy, a Q branch is obtained for all linear molecules. The spectrum of CO, for instance, is shown in Fig. 11C.12: rather than being a single line, as implied by eqn 11C.16, the Q branch appears as a broad feature. Its breadth arises from the presence of several overlapping lines arising from the difference in rotational constants of the upper and lower vibrational states.

The information available from vibrational Raman spectra adds to that from infrared spectroscopy because homonuclear diatomics can also be studied. The spectra can be interpreted in terms of the force constants, dissociation energies, and bond lengths, and some of the information obtained is included in Table 11C.1.

\subsection*{Checklist of concepts}
\begin{enumerate}
  \item The vibrational energy levels of a diatomic molecule modelled as a harmonic oscillator depend on the force constant $k_{\mathrm{f}}$ (a measure of the stiffness of the bond) and the effective mass of the vibration.2. The gross selection rule for infrared spectra is that the electric dipole moment of the molecule must depend on the bond length.3. The specific selection rule for infrared spectra (within the harmonic approximation) is $\Delta v= \pm 1$.4. The Morse potential energy can be used to model anharmonic vibration.5. The strongest infrared transitions are the fundamental transitions ( $v=1 \leftarrow v=0$ ).
  \item Anharmonicity gives rise to weaker overtone transitions $(v=2 \leftarrow v=0, v=3 \leftarrow v=0, \ldots)$.7. A Birge-Sponer plot may be used to determine the dissociation energy of a diatomic molecule.8. In the gas phase, vibrational transitions have a $\mathbf{P}, \mathbf{R}$ branch structure due to simultaneous rotational transitions; some molecules also have a Q branch.9. For a vibration to be Raman active, the polarizability must change as the molecule vibrates.10. The specific selection rule for vibrational Raman spectra (within the harmonic approximation) is $\Delta v= \pm 1$.$\square$ 11. In gas-phase spectra, the Stokes and anti-Stokes lines in a Raman spectrum have an $\mathrm{O}, \mathrm{Q}, \mathrm{S}$ branch structure.
\end{enumerate}

\subsection*{Checklist of equations}
\begin{center}
\begin{tabular}{|c|c|c|c|}
\hline
Property & Equation & Comment & Equation number \\
\hline
Vibrational terms & \( \begin{aligned} & \tilde{G}(v)=\left(v+\frac{1}{2}\right) \tilde{v}, \tilde{v}=(1 / 2 \pi c)\left(k_{\mathrm{f}} / m_{\text {eff }}\right)^{1 / 2} \\ & \quad m_{\text {eff }}=m_{1} m_{2} /\left(m_{1}+m_{2}\right) \end{aligned} \) & Diatomic molecules; harmonic approximation & 11C.4b \\
\hline
Infrared spectra (vibrational) & $\Delta \tilde{G}_{\nu+\frac{1}{2}}=\tilde{v}$ & Diatomic molecules; harmonic approximation & 11C. 6 \\
\hline
Morse potential energy & \( \begin{aligned} & V(x)=h c \tilde{D}_{e}\left\{1-\mathrm{e}^{-a x}\right\}^{2} \\ & \quad a=\left(m_{\mathrm{eff}} \omega^{2} / 2 h c \tilde{D}_{\mathrm{e}}\right)^{1 / 2} \end{aligned} \) &  & 11C. 7 \\
\hline
Vibrational terms (diatomic molecules) & $\tilde{G}(v)=\left(v+\frac{1}{2}\right) \tilde{v}-\left(v+\frac{1}{2}\right)^{2} x_{\mathrm{e}} \tilde{v}, x_{\mathrm{e}}=\tilde{v} / 4 \tilde{D}_{\mathrm{e}}$ & Morse potential energy & 11C. 8 \\
\hline
\multirow[t]{2}{*}{Infrared spectra (vibrational)} & $\Delta \tilde{G}_{\nu+\frac{1}{2}}=\tilde{v}-2(\nu+1) x_{e} \tilde{v}+\cdots$ & Anharmonic oscillator & 11C.9b \\
\hline
 & $\tilde{G}(v+2)-\tilde{G}(v)=2 \tilde{v}-2(2 v+3) x_{\mathrm{e}} \tilde{v}+\cdots$ & First overtone & 11C. 10 \\
\hline
Dissociation wavenumber & $\tilde{D}_{0}=\Delta \tilde{G}_{1 / 2}+\Delta \tilde{G}_{3 / 2}+\cdots=\sum_{\nu} \Delta \tilde{G}_{\nu+\frac{1}{2}}$ & Birge-Sponer plot & 11C. 11 \\
\hline
Vibration-rotation terms (diatomic molecules) & $\tilde{S}(\nu, J)=\left(\nu+\frac{1}{2}\right) \tilde{v}+\tilde{B} J(J+1)$ & Rotation coupled to vibration & 11C.12b \\
\hline
\multirow[t]{4}{*}{Infrared spectra (vibrationrotation)} & \( \begin{aligned} & \tilde{v}_{\mathrm{P}}(J)=\tilde{S}(v+1, J-1)-\tilde{S}(v, J)=\tilde{v}-2 \tilde{B} J \\ & \quad J=1,2,3, \ldots \end{aligned} \) & P branch $(\Delta J=-1)$ & 11C.13a \\
\hline
 & \( \tilde{v}_{\mathrm{Q}}(J)=\tilde{S}(v+1, J)-\tilde{S}(v, J)=\tilde{v} \) & Q branch $(\Delta J=0)$ & 11C.13b \\
\hline
 & \( \begin{aligned} & \tilde{V}_{\mathrm{R}}(J)=\tilde{S}(v+1, J+1)-\tilde{S}(\nu, J)=\tilde{v}+2 \tilde{B}(J+1) \\ & \quad J=0,1,2, \ldots \end{aligned} \) & R branch $(\Delta J=+1)$ & 11C.13c \\
\hline
 & \( \begin{aligned} & \tilde{v}_{\mathrm{R}}(J-1)-\tilde{v}_{\mathrm{P}}(J+1)=4 \tilde{B}_{0}\left(J+\frac{1}{2}\right) \\ & \tilde{v}_{\mathrm{R}}(J)-\tilde{v}_{\mathrm{P}}(J)=4 \tilde{B}_{1}\left(J+\frac{1}{2}\right) \end{aligned} \) & Combination differences & 11C. 15 \\
\hline
\multirow[t]{2}{*}{Raman spectra (vibrationrotation)} & \( \begin{aligned} & \tilde{v}_{\mathrm{O}}(J)=\tilde{v}_{\mathrm{i}}-\tilde{v}+4 \tilde{B}\left(J-\frac{1}{2}\right) \quad J=2,3,4, \ldots \\ & \tilde{v}_{\mathrm{Q}}(J)=\tilde{v}_{\mathrm{i}}-\tilde{v} \end{aligned} \) & \begin{tabular}{l}
O branch $(\Delta J=-2)$ \\
Q branch $(\Delta J=0)$ \\
\end{tabular} & 11C. 16 \\
\hline
 & $\tilde{v}_{\mathrm{S}}(J)=\tilde{v}_{\mathrm{i}}-\tilde{v}-4 \tilde{B}\left(J+\frac{3}{2}\right) \quad J=0,1,2, \ldots$ & S branch ( $\Delta J=+2$ ) &  \\
\hline
\end{tabular}
\end{center}

\section*{TOPIC 11D Vibrational spectroscopy of polyatomic molecules }
\subsection*{Why do you need to know this material?}
The analysis of vibrational spectra is a widely used analytical technique that provides information about the identity and shapes of polyatomic molecules in the gas and condensed phases.

\subsection*{What is the key idea?}
The vibrational spectrum of a polyatomic molecule can be interpreted in terms of its normal modes.

\subsection*{What do you need to know already?}
You need to be familiar with the harmonic oscillator (Topic 7 E ), the general principles of spectroscopy (Topic 11A), and the selection rules for vibrational infrared and Raman spectroscopy (Topic 11C).

There is only one mode of vibration for a diatomic molecule: the periodic stretching and compression of the bond. In polyatomic molecules there are many bond lengths and angles that can change, and as a result the vibrational motion of the molecule is very complex. Some order can be brought to this complexity by introducing the concept of 'normal modes'.

\subsection*{11D. 1 Normal modes}
The first step in the analysis of the vibrations of a polyatomic molecule is to calculate the total number of vibrational modes.

\subsection*{How is that done? 11D. 1 \\
 Counting the number of \\
 vibrational modes}
The total number of coordinates needed to specify the locations of $N$ atoms is $3 N$. Each atom may change its location by varying each of its three coordinates ( $x, y$, and $z$ ), so the total number of displacements available is 3 N . These displacements can be grouped together in a physically sensible way. For example, three coordinates are needed to specify the location of the centre of mass of the molecule, so three of the 3 N\\
displacements correspond to the translational motion of the molecule as a whole. The remaining $3 N-3$ displacements are 'internal' modes of the molecule.

Two angles are needed to specify the orientation of a linear molecule in space: in effect, only the latitude and longitude of the direction in which the molecular axis is pointing need be specified (Fig. 11D.1a). However, three angles are needed for a nonlinear molecule because the orientation of the molecule around the direction defined by the latitude and longitude also needs to be specified (Fig. 11D.1b). Therefore, for linear molecules two of the $3 N-3$ internal displacements are rotational, whereas for nonlinear molecules three of the displacements are rotational. That leaves $3 N-5$ (linear) or $3 N-6$ (nonlinear) non-rotational internal displacements of the atoms: these are the vibrational modes. It follows that the number of modes of vibration is:

Linear molecule: $\quad N_{\text {vib }}=3 N-5$\\
(11D.1)\\
Nonlinear molecule: $N_{v i b}=3 N-6$\\
Numbers of vibrational modes\\
\includegraphics[max width=\textwidth, center]{2025_02_27_d4b95d9718f0b9e2e6b9g-483}\\
(b)

Figure 11D. 1 (a) The orientation of a linear molecule requires the specification of two angles. (b) The orientation of a nonlinear molecule requires the specification of three angles.

\subsection*{Brief illustration 11D. 1}
Water, $\mathrm{H}_{2} \mathrm{O}$, is a nonlinear triatomic molecule with $N=3$, and so has $3 N-6=3$ modes of vibration; $\mathrm{CO}_{2}$ is a linear triatomic molecule, and has $3 N-5=4$ modes of vibration. Methylbenzene has 15 atoms and 39 modes of vibration.

The greatest simplification of the description of vibrational motion is obtained by analysing it in terms of 'normal modes'. A normal mode is a vibration of the molecule in which the centre of mass remains fixed, the orientation is unchanged, and the atoms move synchronously. When a normal mode is excited, the energy remains in that mode and does not migrate into other normal modes of the molecule.

A normal mode analysis is possible only if it is assumed that the potential energy is parabolic (as in a harmonic oscillator, Topic 11C). In reality, the potential energy is not parabolic, the vibrations are anharmonic (Topic 11C), and the normal modes are not completely independent. Nevertheless, a normal mode analysis remains a good starting point for the description of the vibrations of polyatomic molecules.

Figure 11D. 2 shows the four normal modes of $\mathrm{CO}_{2}$. Mode $v_{1}$ is the symmetric stretch in which the two oxygen atoms move in and out synchronously but the carbon atom remains stationary. Mode $v_{2}$, the antisymmetric stretch, in which the two oxygen atoms always move in the same direction and opposite to that of the carbon. Finally, there are two bending modes $v_{3}$ in which the oxygen atoms move perpendicular to the internuclear axis in one direction and the carbon atom moves in the opposite direction: this bending motion can take place in either of two perpendicular planes. In all these modes, the position of the centre of mass and orientation of the molecule are unchanged by the vibration.

In the harmonic approximation, each normal mode, $q$, behaves like an independent harmonic oscillator and has an energy characterized by the quantum number $v_{q}$. Expressed as a wavenumber, these terms are

$$
\tilde{G}_{q}(v)=\left(v_{q}+\frac{1}{2}\right) \tilde{v}_{q} \quad v_{q}=0,1,2, \ldots \quad \tilde{v}_{q}=\frac{1}{2 \pi c}\left(\frac{k_{\mathrm{f}, q}}{m_{q}}\right)^{1 / 2}
$$

Vibrational terms of normal modes [harmonic]\\
(11D.2)\\
where $\tilde{v}_{q}$ is the wavenumber of mode $q$; this quantity depends on the force constant $k_{\mathrm{f}, q}$ for the mode and on the effective mass $m_{q}$ of the mode: stiff bonds and low effective masses correspond to high wavenumbers and hence to high frequency vibrations. The effective mass of the mode is a measure of the mass that moves in the vibration and in general is a combination of the\\
\includegraphics[max width=\textwidth, center]{2025_02_27_d4b95d9718f0b9e2e6b9g-484}

Figure 11D. 2 The four normal modes of $\mathrm{CO}_{2}$. The two bending motions ( $v_{3}$ ) have the same vibrational frequency.\\
\includegraphics[max width=\textwidth, center]{2025_02_27_d4b95d9718f0b9e2e6b9g-484(1)}

Figure 11D. 3 The three normal modes of $\mathrm{H}_{2} \mathrm{O}$. The mode $v_{2}$ is predominantly bending, and occurs at lower wavenumber than the other two.\\
masses of the atoms. For example, in the symmetric stretch of $\mathrm{CO}_{2}$, the carbon atom is stationary, and the effective mass depends on the masses of only the oxygen atoms. In the antisymmetric stretch and in the bends all three atoms move, so the masses of all three atoms contribute (but to different extents) to the effective mass of each mode.\\
The three normal modes of $\mathrm{H}_{2} \mathrm{O}$ are shown in Fig. 11D.3: note that the predominantly bending mode $\left(v_{2}\right)$ has a lower frequency (and wavenumber) than the others, which are predominantly stretching modes. It is generally the case that the frequencies of bending motions are lower than those of stretching modes. Only in special cases (such as the $\mathrm{CO}_{2}$ molecule) are the normal modes purely stretches or purely bends. In general, a normal mode is a composite motion involving simultaneous stretching of bonds and changes to bond angles. In a given normal mode, heavy atoms generally move less than light atoms.

The vibrational state of a polyatomic molecule is specified by the vibrational quantum number $v_{q}$ for each of the normal modes. For example, for $\mathrm{H}_{2} \mathrm{O}$ with three normal modes, the vibrational state is designated $\left(v_{1}, v_{2}, v_{3}\right)$. The vibrational ground state of an $\mathrm{H}_{2} \mathrm{O}$ molecule is therefore $(0,0,0)$; the state $(0,1,0)$ implies that modes 1 and 3 are in their ground states, and mode 2 is in the first excited state.

\subsection*{11D. 2 Infrared absorption spectra}
The gross selection rule for infrared activity is a straightforward generalization of the rule for diatomic molecules (Topic 11C):

\begin{displayquote}
The motion corresponding to a normal mode must be accompanied by a change of electric dipole moment.
\end{displayquote}

Simple inspection of atomic motions is sometimes all that is needed in order to assess whether a normal mode is infrared active. For example, the symmetric stretch of $\mathrm{CO}_{2}$ leaves the dipole moment unchanged (at zero, see Fig. 11D.2), so this mode is infrared inactive. The antisymmetric stretch, however, changes the dipole moment because the molecule becomes unsymmetrical as it vibrates, so this mode is infrared active. Because the dipole moment change is parallel to the principal axis, the transitions arising from this mode are classified as parallel bands in the spectrum. Both bending modes\\
are infrared active: they are accompanied by a changing dipole perpendicular to the principal axis, so transitions involving them lead to a perpendicular band in the spectrum.

\subsection*{Example 11D. 1 Using the gross selection rule for infrared}
 spectroscopyState which of the following molecules are infrared active: $\mathrm{N}_{2} \mathrm{O}, \mathrm{OCS}, \mathrm{H}_{2} \mathrm{O}, \mathrm{CH}_{2}=\mathrm{CH}_{2}$.

Collect your thoughts Molecules that are infrared active have a normal mode (or modes) in which there is a change in dipole moment during the course of the motion. Therefore, to work out if a molecule is infrared active you need to decide whether there is any distortion of the molecule that results in a change in its electric dipole moment (including changes from zero).

The solution The linear molecules $\mathrm{N}_{2} \mathrm{O}$ and OCS both have permanent electric dipole moments that change as a result of stretching any of the bonds; in addition, bending perpendicular to the internuclear axis results in a dipole in that direction: both molecules are therefore infrared active. An $\mathrm{H}_{2} \mathrm{O}$ molecule also has a permanent dipole moment which changes either by stretching the bonds or by altering the bond angle: the molecule is infrared active. $\mathrm{A} \mathrm{CH}_{2}=\mathrm{CH}_{2}$ molecule does not have a permanent dipole moment (it possesses a centre of symmetry) but there are vibrations in which the symmetry is reduced and a dipole moment forms, for example, by stretching the two $\mathrm{C}-\mathrm{H}$ bonds on one carbon atom and simultaneously compressing the two $\mathrm{C}-\mathrm{H}$ bonds on the other carbon atom.

Comment. Topic 11E describes a systematic procedure based on group theory for deciding whether a vibrational mode is infrared active.

Self-test 11D. 1 Identify an infrared inactive normal mode of $\mathrm{CH}_{2}=\mathrm{CH}_{2}$.\\
\includegraphics[max width=\textwidth, center]{2025_02_27_d4b95d9718f0b9e2e6b9g-485}\\
\includegraphics[max width=\textwidth, center]{2025_02_27_d4b95d9718f0b9e2e6b9g-485(1)}

The specific selection rule in the harmonic approximation is $\Delta v_{q}= \pm 1$. In this approximation the quantum number of only one active mode can change in the interaction of a molecule with a photon. A fundamental transition is a transition from the ground state of the molecule to the next higher energy level of the specified mode. For example, in $\mathrm{H}_{2} \mathrm{O}$ there are three such fundamentals corresponding to the excitation of each of the three normal modes: $(1,0,0) \leftarrow(0,0,0),(0,1,0) \leftarrow$ $(0,0,0)$, and $(0,0,1) \leftarrow(0,0,0)$.

Anharmonicity also allows transitions in which more than one quantum of excitation takes place: such transitions are referred to as overtones. A transition such as $(0,0,2) \leftarrow$ $(0,0,0)$ in $\mathrm{H}_{2} \mathrm{O}$ is described as a first overtone, and a transition such as $(0,0,3) \leftarrow(0,0,0)$ is a second overtone of the mode $v_{3}$. Combination bands (or combination lines) corresponding to the simultaneous excitation of more than one normal mode\\
in the transition, as in $(1,1,0) \leftarrow(0,0,0)$, are also possible in the presence of anharmonicity.

As for diatomic molecules (Topic 11C), transitions between vibrational levels can be accompanied by simultaneous changes in rotational state, so giving rise to band spectra rather than the single absorption line of a pure vibrational transition. The spectra of linear polyatomic molecules show branches similar to those of diatomic molecules. For nonlinear molecules, the rotational fine structure is considerably more complex and difficult to analyse: even in moderately complex molecules the presence of several normal modes gives rise to several fundamental transitions, many overtones, and many combination lines, each with associated rotational fine structure, and results in infrared spectra of considerable complexity.

These complications are eliminated (or at least concealed) by recording infrared spectra of samples in the condensed phase (liquids, solutions, or solids). Molecules in liquids do not rotate freely but collide with each other after only a small change of orientation. As a result, the lifetimes of rotational states in liquids are very short, which results in a broadening of the associated energies (Topic 11A). Collisions occur at a rate of about $10^{13} \mathrm{~s}^{-1}$ and, even allowing for only a 10 per cent success rate in changing the molecule into another rotational state, a lifetime broadening of more than $1 \mathrm{~cm}^{-1}$ can easily result. The rotational structure of the vibrational spectrum is blurred by this effect, so the infrared spectra of molecules in condensed phases usually consist of bands without any resolved branch structure.

Infrared spectroscopy is commonly used in routine chemical analysis, most usually on samples in solution, prepared as a fine dispersion (a 'mull'), or compressed as solids into a very thin layer. The resulting spectra show many absorption bands, even for moderately complex molecules. There is no chance of analysing such complex spectra in terms of the normal modes. However, they are of great utility for it turns out that certain groups within a molecule (such as a carbonyl group or an $-\mathrm{NH}_{2}$ group) give rise to absorption bands in a particular range of wavenumbers. The spectra of a very large number of molecules have been recorded and these data have been used to draw up charts and tables of the expected range of the wavenumbers of absorptions from different groups. Comparison of the features in the spectrum of an unknown molecule or the product of a chemical reaction with entries in these data tables is a common first step towards identifying the molecule.

\subsection*{11D. 3 Vibrational Raman spectra}
As for diatomic molecules, the normal modes of vibration of molecules are Raman active if they are accompanied by a changing polarizability. A closer analysis of infrared and

Raman activity of normal modes based on considerations of symmetry leads to the exclusion rule:

If the molecule has a centre of symmetry, then no mode can be both infrared and Raman active. Exclusion rule\\
(A mode may be inactive in both.) Because it is often possible to judge intuitively if a mode changes the molecular dipole moment, this rule can be used to identify modes that are not Raman active.

\subsection*{Brief illustration 11D. 2}
The symmetric stretch of $\mathrm{CO}_{2}$ alternately swells and contracts the molecule: this motion changes the size and hence the polarizability of the molecule, so the mode is Raman active. Because $\mathrm{CO}_{2}$ has a centre of symmetry the exclusion rule applies, so the stretching mode cannot be infrared active. The antisymmetric stretch and the two bends are infrared active, and so cannot be Raman active.

The assignment of Raman lines to particular vibrational modes is aided by noting the state of polarization of the scattered light. The depolarization ratio, $\rho$, of a line is the ratio of the intensities of the scattered light with polarization perpendicular and parallel, $I_{\perp}$ and $I_{\|}$, to the plane of the incident radiation (Fig. 11D.4):

$$
\rho=\frac{I_{\perp}}{I_{\|}}
$$

Depolarization ratio [definition]\\
(11D.3)\\
\includegraphics[max width=\textwidth, center]{2025_02_27_d4b95d9718f0b9e2e6b9g-486}

Figure 11D. 4 The definition of the planes used for the specification of the depolarization ratio, $\rho$, in Raman scattering.

To measure $\rho$, the intensity of a Raman line is measured with a polarizing filter (a 'half-wave plate') first parallel and then perpendicular to the polarization of the incident beam. If the emergent light is not polarized, then both intensities are the same and $\rho$ is close to 1 ; if the light retains its initial polarization, then $I_{\perp}=0$, so $\rho=0$. A line is classified as depolarized if it has $\rho$ close to or greater than 0.75 and as polarized if $\rho<0.75$. Only totally symmetrical vibrations give rise to polarized lines in which the incident polarization is largely preserved. Vibrations that are not totally symmetrical give rise to depolarized lines because the incident radiation can give rise to radiation in the perpendicular direction too.

\subsection*{Checklist of concepts}
$\square$ 1. A normal mode is a synchronous displacement of the atoms in which the centre of mass and orientation of the molecule remains fixed. In the harmonic approximation, normal modes are mutually independent.2. A normal mode is infrared active if it is accompanied by a change of electric dipole moment; the specific selection rule is $\Delta v_{q}= \pm 1$.\\
3. A normal mode is Raman active if it is accompanied by a change in polarizability; the specific selection rule is $\Delta v_{q}= \pm 1$.\\
4. The exclusion rule states that, if the molecule has a centre of symmetry, then no mode can be both infrared and Raman active.\\
5. Polarized lines preserve the polarization of the incident radiation in the Raman spectrum and arise from totally symmetrical vibrations.

\subsection*{Checklist of equations}
\begin{center}
\begin{tabular}{|c|c|c|c|}
\hline
Property & Equation & Comment & Equation number \\
\hline
Number of normal modes & \( \begin{aligned} & N_{\text {vib }}=3 N-5 \text { (linear) } \\ & N_{\text {vib }}=3 N-6 \text { (nonlinear) } \end{aligned} \) & Independent if harmonic; $N$ is the number of atoms & 11D. 1 \\
\hline
Vibrational terms of normal modes & \( \begin{aligned} & \tilde{G}_{q}\left(v_{q}\right)=\left(v_{q}+\frac{1}{2}\right) \tilde{v}_{q}, \\ & \tilde{v}_{q}=(1 / 2 \pi c)\left(k_{\mathrm{f}, q} / m_{q}\right)^{1 / 2} \end{aligned} \) & Harmonic approximation & 11D. 2 \\
\hline
Depolarization ratio & $\rho=I_{\perp} / I_{\|}$ & \begin{tabular}{l}
Depolarized lines: $\rho$ close to or greater than 0.75 \\
Polarized lines: $\rho<0.75$ \\
\end{tabular} & 11D. 3 \\
\hline
\end{tabular}
\end{center}

\section*{TOPIC 11E Symmetry analysis of vibrational spectra }
\subsection*{Why do you need to know this material?}
The analysis of vibrational spectra is aided by understanding the relationship between the symmetry of the molecule, its normal modes, and the selection rules that govern the transitions.

\subsection*{What is the key idea?}
The vibrational modes of a molecule can be classified according to the symmetry of the molecule.

\subsection*{What do you need to know already?}
You need to be familiar with the vibrational spectra of polyatomic molecules (Topic 11D) and the treatment of symmetry in Focus 10.

The classification of the normal modes of vibration of a polyatomic molecule according to their symmetry makes it possible to predict in a very straightforward way which are infrared or Raman active.

\subsection*{11E. 1 Classification of normal modes according to symmetry}
Each normal mode can be classified as belonging to one of the symmetry species of the irreducible representations of the molecular point group. The classification proceeds as follows:

\begin{enumerate}
  \item The basis functions are the three displacement vectors ( $x$, $y$, and $z$ ) on each atom: there are $3 N$ such basis functions for a molecule with $N$ atoms.
  \item The character, $\chi(C)$, for each class, $C$, of operations in the group is found by considering the effect of one operation of the class and counting 1 for each basis function that is unchanged by the operation, -1 if the basis function changes sign, and 0 if it changes into some other displacement.
  \item The resulting representation is decomposed into its component irreducible representations, denoted $\Gamma$ with\\
characters $\chi^{(\Gamma)}(C)$, by using the relevant character table in conjunction with eqn 10C.3a:
\end{enumerate}

$$
n(\Gamma)=\frac{1}{h} \sum_{C} N(C) \chi^{(\Gamma)}(C) \chi(C)
$$

where $h$ is the order of the group and $N(C)$ the number of operations in the class $C$.\\
4. The symmetry species corresponding to $x, y$, and $z$ (corresponding to translations) and those corresponding to the rotations about $x, y$, and $z$ (denoted $R_{x}, R_{y}$, and $R_{z}$ ) are removed. Their symmetry species are listed in the character tables.\\
5. The remaining symmetry species correspond to the normal modes.

\subsection*{Example 11E. 1 Identifying the symmetry species of the normal modes of $\mathrm{H}_{2} \mathrm{O}$}
Identify the symmetry species of the normal modes of $\mathrm{H}_{2} \mathrm{O}$, which belongs to the point group $C_{2 v}$.

Collect your thoughts You need to identify the axes in the molecule and then refer to the character table (in the Resource section) for the symmetry operations and their characters. You need consider only one symmetry operation of each class, because all members of the same class have the same character. (In $C_{2 v}$, there is only one member of each class anyway.) Then follow the five steps outlined in the text. Note from the character table the symmetry species of the translations and rotations, which are given in the right-hand column.

The solution The molecule lies in the $y z$-plane, with the $z$-axis bisecting the HOH bond angle. The three displacement vectors of each atom are shown in the Fig. 11E.1. From the character table for $C_{2 v}$, the symmetry operations are $E, C_{2}, \sigma_{v}(x z)$, and $\sigma_{\mathrm{v}}^{\prime}(y z)$. None of the nine displacement vectors is affected by the operation $E$, so $\chi(E)=9$. The $C_{2}$ operation moves all the displacement vectors on the H atoms to other positions, so these count 0 ; the $x$ and $y$ displacement vectors on the O atom change sign, giving a count of -1 each, whereas the $z$ displacement vector is unaffected, giving a count of +1 . Hence $\chi\left(C_{2}\right)=$ $-1-1+1=-1$. The operation $\sigma_{\mathrm{v}}(x z)$ moves all the displace-\\
ment vectors on the H atoms to other positions, changes the sign of the $y$ displacement vector on the O atom, and leaves the $x$ and $z$ vectors unaffected: hence $\chi\left(\sigma_{v}\right)=-1+1+1=1$. The operation $\sigma_{\mathrm{v}}^{\prime}(y z)$ changes the sign of the $x$ displacement vectors on both H atoms and leaves the sign of the $y$ and $z$ displacement vectors unaffected. For the O atom, the $x$ vector changes sign and the $y$ and $z$ displacement vectors are unaffected: hence $\chi\left(\sigma_{v}^{\prime}\right)=-1-1+1+1+1+1-1+1+1=3$. The characters of the representation are therefore $9,-1,1,3$. Decomposition by using eqn 10C.3a shows that the reducible representation spans the symmetry species $3 \mathrm{~A}_{1}+\mathrm{A}_{2}+2 \mathrm{~B}_{1}+$ $3 \mathrm{~B}_{2}$. The translations have symmetry species $\mathrm{B}_{1}, \mathrm{~B}_{2}$, and $\mathrm{A}_{1}$ and the rotations are $B_{1}, B_{2}$, and $A_{2}$; their removal leaves $2 A_{1}+B_{2}$ as the symmetry species of the normal modes. As expected, there are three such modes.\\
\includegraphics[max width=\textwidth, center]{2025_02_27_d4b95d9718f0b9e2e6b9g-488}

Fig. 11E. 1 The atomic displacements of $\mathrm{H}_{2} \mathrm{O}$ and the symmetry elements used to calculate the characters.

Comment. In Fig. 11D. 3 (Topic 11D), $v_{1}$ and $v_{2}$ have symmetry $\mathrm{A}_{1}, v_{3}$ has symmetry $\mathrm{B}_{2}$. This assignment is evident from the fact that the combination of displacements for both $v_{1}$ and $v_{2}$ are unchanged by any of the operations of the group, so the characters are all 1 as required for $\mathrm{A}_{1}$. In contrast, for $v_{3}$ the displacements shown change sign under $C_{2}$ and $\sigma_{v}$ giving the characters $1,-1,-1,1$, which correspond to $B_{2}$.

Self-test 11E. 1 Identify the symmetry species of the normal modes of methanal, $\mathrm{H}_{2} \mathrm{C}=\mathrm{O}$, point group $\mathrm{C}_{2 \mathrm{v}}$ (orientate the molecule in the same way as $\mathrm{H}_{2} \mathrm{O}$, with the $\mathrm{CH}_{2}$ group in the $y z$-plane).\\
${ }^{\mathrm{r}} \mathrm{g} z+{ }^{\mathrm{I}} \mathrm{g}+{ }^{\mathrm{I}} \forall \varepsilon$ : ıวмs $\boldsymbol{u}_{V}$

All the normal modes of $\mathrm{H}_{2} \mathrm{O}$ are either A or B and therefore non-degenerate. There are no two- or higher-dimensional irreducible representations in $C_{2 v}$ molecules, so vibrational degeneracy never arises. Degeneracy can arise in molecules with higher symmetry, as illustrated in the following example.

Example 11E. 2 Identifying the symmetry species of the normal modes of $\mathrm{BF}_{3}$

Identify the symmetry species of the normal modes of vibration of $\mathrm{BF}_{3}$, which is trigonal planar and belongs to the point group $D_{3 h}$.

Collect your thoughts The overall procedure is the same as in Example 11E.1. However, because the molecule is $D_{3 \mathrm{~h}}$, which has two-dimensional irreducible representations ( $\mathrm{E}^{\prime}$ and $E^{\prime \prime}$ ), you need to be alert for the possibility that there are doubly degenerate pairs of normal modes. You can treat the displacement vectors on the $B$ atom separately from those on the F atoms because no symmetry operation interconverts these two sets: this separation simplifies the calculations. Because the molecule is nonlinear with 4 atoms, there are 6 normal modes.

The solution The $C_{3}$ axis is the principal axis and defines the $z$-direction; the molecule lies in the $x y$-plane. The three $C_{2}$ axes pass along the $\mathrm{B}-\mathrm{F}$ bonds, and the three $\sigma_{\mathrm{v}}$ planes contain the $\mathrm{B}-\mathrm{F}$ bonds and are perpendicular to the plane of the molecule. The $\sigma_{\mathrm{h}}$ plane lies in the plane of the molecule, and the $S_{3}$ axis is coincident with the $C_{3}$ axis.\\
First, consider the displacement vectors on the $B$ atom. Because this atom lies on the principal axis, the $z$ displacement vector must transform as the function $z$, which from the character table has the symmetry species $\mathrm{A}_{2}^{\prime \prime}$. Similarly, the $x$ and $y$ displacement vectors together transform as $\mathrm{E}^{\prime}$.

Next consider the nine displacement vectors on the F atoms. The identity operation has no effect, so $\chi(E)=9$. The $C_{3}$ operation moves all these vectors, so $\chi\left(C_{3}\right)=0$. A $C_{2}$ operation about a particular B-F bond has no effect on the displacement vector that points along the bond, but the other two vectors change sign; the displacement vectors on the other $F$ atoms are moved, hence $\chi\left(C_{2}\right)=1-1-1=-1$. Under $\sigma_{\mathrm{h}}$ the $z$ displacement vector on each F changes sign, but the $x$ and $y$ vectors do not. The character is therefore $\chi\left(\sigma_{\mathrm{h}}\right)=3 \times(-1+$ $1+1)=3$. The character for $S_{3}$ is the same as for $C_{3}, \chi\left(S_{3}\right)=0$. A $\sigma_{\mathrm{v}}$ reflection in a plane containing a particular B-F bond has no effect on the displacement vector that points along the bond, nor on the $z$ displacement vector; however, the other vector changes sign. The displacement vectors on the other atoms are moved, hence $\chi\left(\sigma_{v}\right)=1+1-1=1$. The characters of the reducible representation are therefore $9,0,-1,3,0,1$; this set can be decomposed into the symmetry species $\mathrm{A}_{1}^{\prime}+\mathrm{A}_{2}^{\prime}+2 \mathrm{E}^{\prime}+\mathrm{A}_{2}^{\prime \prime}+\mathrm{E}^{\prime \prime}$ for the displacement vectors on the F atoms. The displacement vectors on the B atom transform as $A_{2}^{\prime \prime}+E^{\prime}$, so the complete set of symmetry species is $A_{1}^{\prime}+A_{2}^{\prime}+$ $3 \mathrm{E}^{\prime}+2 \mathrm{~A}_{2}^{\prime \prime}+\mathrm{E}^{\prime \prime}$.\\
The character table shows that $z$ transforms as $\mathrm{A}_{2}^{\prime \prime}$, and $x$ and $y$ together span $\mathrm{E}^{\prime}$. The rotation about $z, R_{z}$, transforms as $\mathrm{A}_{2}^{\prime}$ and rotations about $x$ and $y$ together $\left(R_{x}, R_{y}\right)$ transform as $\mathrm{E}^{\prime \prime}$. Removing these symmetry species from the complete set leaves $\mathrm{A}_{1}^{\prime}+2 \mathrm{E}^{\prime}+\mathrm{A}_{2}^{\prime \prime}$ as the symmetry species of the vibrational modes. Figure 11E. 2 shows these normal modes.\\
Comment. Because the $\mathrm{E}^{\prime}$ symmetry species is two-dimensional, the corresponding normal mode is doubly degenerate. The above analysis shows that there are two $\mathrm{E}^{\prime}$ symmetry species present, which correspond to two different doublydegenerate normal modes. The total number of normal modes represented by $\mathrm{A}_{1}^{\prime}+2 \mathrm{E}^{\prime}+\mathrm{A}_{2}^{\prime \prime}$ is therefore $1+2 \times 2+1=6$.\\
\includegraphics[max width=\textwidth, center]{2025_02_27_d4b95d9718f0b9e2e6b9g-489(1)}

Fig. 11E. 2 The normal modes of vibration of $\mathrm{BF}_{3}$.

Self-test 11E. 2 Identify the symmetry species of the normal modes of ammonia $\mathrm{NH}_{3}$, point group $C_{3 v}$.\\
\includegraphics[max width=\textwidth, center]{2025_02_27_d4b95d9718f0b9e2e6b9g-489}

\subsection*{11E. 2 Symmetry of vibrational wavefunctions}
For a one-dimensional harmonic oscillator the ground-state wavefunction (with $v=0$ ) is proportional to $\mathrm{e}^{-x^{2} / 2 \alpha^{2}}$, where $x$ is the displacement from the equilibrium position and $\alpha$ is a constant (Topic 7E). For the first excited state, with $v=1$, the wavefunction is proportional to $x \mathrm{e}^{-x^{2} / 2 \alpha^{2}}$. The same wavefunctions apply to a normal mode $q$ of a more complex molecule provided that $x$ is replaced by the normal coordinate, $Q_{q}$, which is the combination of displacements that corresponds to the normal mode. For example, in the case of the symmetric stretch of $\mathrm{CO}_{2}$ the normal coordinate is $z_{\mathrm{O}, 1}-z_{\mathrm{O}, 2}$, where $z_{\mathrm{O}, i}$ is the $z$-displacement of oxygen atom $i$.

The effect of any symmetry operation on the normal coordinate of a non-degenerate normal mode is either to leave it unchanged or at most to change its sign. In other words, all the characters are either 1 or -1 . The ground-state wavefunction is a function of the square of the normal coordinate, so regardless of whether $Q_{q} \rightarrow+Q_{q}$, or $Q_{q} \rightarrow-Q_{q}$ as a result of any symmetry operation, the effect on $Q_{q}^{2}$ is to leave it unaffected. Therefore, all the characters for $Q_{q}^{2}$ are 1 , so the ground-state wavefunction transforms as the totally symmetric irreducible representation (typically $\mathrm{A}_{1}$ ).

The first excited state wavefunction is a product of a part that depends on $Q_{q}^{2}$ (the exponential term) and a factor proportional to $Q_{q}$. As has already been seen, $Q_{q}^{2}$ transforms as the totally symmetric irreducible representation and $Q_{q}$ has the same symmetry species as the normal mode. The direct product of the totally symmetric irreducible representation with any symmetry species leaves the latter unaffected, hence it follows that the symmetry of the first excited state wavefunction is the same as that of the normal mode.

\subsection*{(a) Infrared activity of normal modes}
Once the symmetry of a particular normal mode is known it is a simple matter to determine from the appropriate character table whether or not the fundamental transition of that mode is allowed and therefore is infrared active.

How is that done? 11E. 1 Determining the infrared activity of a normal mode

You need to note that the fundamental transition of a particular normal mode is the transition from the ground state, $v_{q}=0$, to the first excited state, $v_{q}=1$. You already know that the state with $v_{q}=0$ transforms as the totally symmetric irreducible representation, and the state with $v_{q}=1$ has the same symmetry as the corresponding normal mode.

Step 1 Formulate the integral used to identify the selection rule Whether or not the transition between $v_{q}=0$ and $v_{q}=1$ is allowed is assessed by evaluating the transition dipole between $\psi_{0}$ and $\psi_{1}, \mu_{10}=\int \psi_{1}^{*} \hat{\mu} \psi_{0} \mathrm{~d} \tau$ (Topic 11A); the dipole moment operator transforms as $x, y$, or $z$ (Topic 10C). As is shown in Topic 10C, this integral may be non-zero only if the product $\psi_{1}^{*} \hat{\mu} \psi_{0}$ spans the totally symmetric irreducible representation.\\
Step 2 Identify the symmetry species spanned by the integrand You can find the symmetry species of $\psi_{1}^{\star} \hat{\mu} \psi_{0}$ by taking the direct product of the symmetry species spanned by each of $\psi_{1}, \hat{\mu}$, and $\psi_{0}$ separately (Topic 10C). Because $\psi_{0}$ transforms as the totally symmetric irreducible representation it has no effect of the symmetry of $\psi_{1}^{\star} \hat{\mu} \psi_{0}$, and so you need consider only the symmetry of the product $\psi_{1}^{\star} \hat{\mu}$. As is shown in Topic 10 C , the only way for this product to span the totally symmetric irreducible representation is for $\psi_{1}^{*}$ and $\hat{\mu}$ to span the same symmetry species. In other words, the integral can be non-zero only if $\psi_{1}^{*}$, and hence the normal mode, has the same symmetry species as $x, y$, or $z$.

The result of the analysis can be summarized by the following rule:

A mode is infrared active only if its symmetry species is the same as the symmetry species of any of $x, y$, or $z$.

Symmetry test for IR activity

\subsection*{Brief illustration 11E. 1}
The normal modes of $\mathrm{BF}_{3}$ (point group $D_{3 \mathrm{~h}}$ ) have symmetry species $\mathrm{A}_{1}^{\prime}+2 \mathrm{E}^{\prime}+\mathrm{A}_{2}^{\prime \prime}$ (Example 11E.2). From the character table it can be seen that $z$ transforms as $\mathrm{A}_{2}^{\prime \prime}$ and $(x, y)$ jointly transform as $\mathrm{E}^{\prime}$. The $\mathrm{A}_{2}^{\prime \prime}$ normal mode and the two doublydegenerate $\mathrm{E}^{\prime}$ normal modes are therefore infrared active. The $\mathrm{A}_{1}^{\prime}$ mode is not.

\subsection*{(b) Raman activity of normal modes}
Symmetry arguments also provide a systematic way of deciding whether or not the fundamental of a normal mode gives rise to Raman scattering; that is, whether or not the mode is Raman active. The argument is similar to that for assessing infrared activity except that it is based on the symmetry of the polarizability operator rather than the dipole moment operator. That operator transforms in the same way as the quadratic forms ( $x^{2}, x y$, and so on, which are listed in the character tables), and leads to the following rule:

A mode is Raman active only if its symmetry species is the same as the symmetry species of a quadratic form.

Symmetry test for Raman activity

\subsection*{Brief illustration 11E. 2}
The normal modes of $\mathrm{BF}_{3}$ (point group $D_{3 \mathrm{~h}}$ ) have symmetry species $\mathrm{A}_{1}^{\prime}+2 \mathrm{E}^{\prime}+\mathrm{A}_{2}^{\prime \prime}$ (Example 11E.2). From the character table it can be seen that $x^{2}, y^{2}$, and $z^{2}$ all transform as $\mathrm{A}_{1}^{\prime}$, and that ( $x^{2}-y^{2}, 2 x y$ ) jointly transform as $\mathrm{E}^{\prime}$. The $\mathrm{A}_{1}^{\prime}$ normal mode and the two doubly-degenerate $\mathrm{E}^{\prime}$ normal modes are therefore Raman active. The $A_{2}^{\prime \prime}$ mode is not active because no quadratic form has this symmetry species. The $\mathrm{E}^{\prime}$ modes are both infrared and Raman active: the exclusion rule does not apply because $\mathrm{BF}_{3}$ does not have a centre of symmetry. The $\mathrm{A}_{1}^{\prime}$ normal mode is the highly-symmetrical breathing mode in which all the B-F bonds stretch together. The corresponding Raman\\
line is expected to be polarized. The $\mathrm{E}^{\prime}$ modes are expected to give depolarized lines.

\subsection*{(c) The symmetry basis of the exclusion rule}
The exclusion rule, Section 11D.3, can be derived by using a symmetry argument. If a molecule has a centre of symmetry, then all the symmetry species of its displacements are either g or u according to their behaviour under inversion. If the character for this operation is positive, indicating that the displacement or displacements are unchanged by the operation, the label is g , whereas if the character is negative, indicating that the sign of the displacement or displacements are changed, the label is u .

The functions $x, y$, and $z$ (which occur in the transition dipole moment) all change sign under inversion, so they must correspond to symmetry species with a label $u$. In contrast, the quadratic forms (which govern the Raman activity) are all unchanged by inversion and so have the label g. For example, the effect of the inversion on $x z$ is to transform it into $(-x)(-z)=x z$.

Any normal mode in a molecule with a centre of symmetry corresponds to a symmetry species that is either g or u . If the normal mode has the same symmetry species as $x, y$, or $z$, it is infrared active; such a mode must be $u$. If the normal mode has the same symmetry species as a quadratic form, it is Raman active; such a mode must be g. Because a normal mode cannot be both $g$ and $u$, no mode can be both infrared and Raman active.

\subsection*{Checklist of concepts}
\begin{enumerate}
  \item A normal mode is infrared active if its symmetry species is the same as the symmetry species of $x, y$, or $z$.
  \item A normal mode is Raman active if its symmetry species is the same as the symmetry species of a quadratic form.
\end{enumerate}

\section*{TOPIC 11F Electronic spectra}
\subsection*{Why do you need to know this material?}
The study of spectroscopic transitions between different electronic states of molecules gives access to data on the electronic structure of molecules, and hence insight into bonding, as well as vibrational frequencies and bond lengths.

\subsection*{What is the key idea?}
Electronic transitions occur within a stationary nuclear framework.

\subsection*{What do you need to know already?}
You need to be familiar with the general features of spectroscopy (Topic 11A), the quantum mechanical origins of selection rules (Topics $8 \mathrm{C}, 11 \mathrm{~B}$, and 11C), and vibrationrotation spectra (Topic 11C). It would be helpful to be aware of atomic term symbols (Topic 8C).

Electronic spectra arise from transitions between the electronic energy levels of molecules. These transitions may also be accompanied by simultaneous changes in vibrational energy; for small molecules in the gas phase the resulting spectral features can be resolved (Fig. 11F.1a), but in a liquid or solid the individual lines usually merge together and result in a broad, almost featureless band (Fig. 11F.1b).

The energies needed to change the electron distributions of molecules are of the order of several electronvolts $(1 \mathrm{eV}$ is equivalent to about $8000 \mathrm{~cm}^{-1}$ or $100 \mathrm{~kJ} \mathrm{~mol}^{-1}$ ). Consequently, the photons emitted or absorbed when such changes occur lie in the visible and ultraviolet regions of the spectrum (Table 11F.1).

\subsection*{11F. 1 Diatomic molecules}
Topic 8C explains how the states of atoms are described by using term symbols. The electronic states of diatomic molecules are also specified by using term symbols, the key difference being that the full spherical symmetry of atoms is replaced by the cylindrical symmetry defined by the axis of the molecule.\\
\includegraphics[max width=\textwidth, center]{2025_02_27_d4b95d9718f0b9e2e6b9g-491}

Figure 11F. 1 Electronic absorption spectra recorded in the visible region. (a) The spectrum of $I_{2}$ in the gas phase shows resolved vibrational structure. (b) The spectrum of chlorophyll recorded in solution, shows only broad bands with no resolved structure. (Absorbance, $A$, is defined in Topic 11A.)

Table 11F. 1 Colour, frequency, and energy of light*

\begin{center}
\begin{tabular}{lrcc}
\hline
Colour & $\lambda / \mathrm{nm}$ & $v /\left(10^{14} \mathrm{~Hz}\right)$ & $E /\left(\mathbf{k J ~ m o l}^{-1}\right)$ \\
\hline
Infrared & $>1000$ & $<3.0$ & $<120$ \\
Red & 700 & 4.3 & 170 \\
Yellow & 580 & 5.2 & 210 \\
Blue & 470 & 6.4 & 250 \\
Ultraviolet & $<400$ & $>7.5$ & $>300$ \\
\hline
\end{tabular}
\end{center}

${ }^{*}$ More values are given in the Resource section.

\subsection*{(a) Term symbols}
In a diatomic molecule only the component of the total orbital angular momentum around the internuclear axis can be specified: the quantum number for this component is $\Lambda$ (uppercase lambda). Its value is found by adding together the component of the orbital angular momentum along the internuclear axis, $\lambda_{i}$, for each electron present:


\begin{equation*}
\Lambda=\lambda_{1}+\lambda_{2}+\cdots \tag{11F.1}
\end{equation*}


For an electron in a $\sigma$ molecular orbital (which is cylindrically symmetric), $\lambda=0$; for an electron in one of the degenerate pair of $\pi$ orbitals $\lambda= \pm 1$. In the molecular term symbol the value of $\Lambda$ is represented by an uppercase Greek letter in the following way

\begin{center}
\begin{tabular}{lllll}
$|\Lambda|$ & 0 & 1 & 2 & $\cdots$ \\
 & $\Sigma$ & $\Pi$ & $\Delta$ & $\cdots$ \\
\end{tabular}
\end{center}

These labels are the analogues of $\mathrm{S}, \mathrm{P}, \mathrm{D}, \ldots$ used for atomic states with $L=0,1,2, \ldots$. The total spin, $S$, of a linear molecule is specified in the same way as for an atom. As in an atomic term symbol, the value of $2 S+1$ is shown as a left superscript and denotes the multiplicity of the term.

Configurations such $\sigma^{2}$ and $\pi^{4}$ with all electrons paired have $S=0$. Such configurations do not contribute to the total orbital angular momentum, either because both electrons have $\lambda=0$ or because there are equal numbers of electrons with $\lambda=+1$ and -1 . The term symbol for the ground state of $\mathrm{H}_{2}$, configuration $1 \sigma_{\mathrm{g}}^{2}$, is therefore ${ }^{1} \Sigma$, and the same is true of the ground state of $\mathrm{N}_{2}$, configuration $1 \sigma_{\mathrm{g}}^{2} 1 \sigma_{\mathrm{u}}^{2} 1 \pi_{\mathrm{u}}^{4} 2 \sigma_{\mathrm{g}}^{2}$.

\subsection*{Brief illustration 11F.1}
The ground-state configuration of $\mathrm{H}_{2}^{+}$is $1 \sigma_{\mathrm{g}}^{1}$. The single $\sigma$ electron has $\lambda=0$, so $\Lambda=0$; for a single electron $S=\frac{1}{2}$, so $2 S+1$ $=2$. The term symbol is therefore ${ }^{2} \Sigma$ (read as 'doublet sigma').\\
The ground-state configuration of NO is ... $1 \pi^{1}$, where ... indicates completed orbitals that make no contribution to either $S$ or $\Lambda$. The single electron can occupy either of the degenerate $\pi$ orbitals, so $\Lambda=+1$ or -1 ; for a single electron, $S=\frac{1}{2}$. The term symbol is therefore ${ }^{2} \Pi$ (read as 'doublet pi').

The ground-state configuration of $\mathrm{O}_{2}$ is $\ldots 1 \pi_{\mathrm{g}}^{2}$. If the two electrons occupy the same $\pi$ orbital, $\Lambda=(+1)+(+1)=+2$ (or $\Lambda=(-1)+(-1)=-2)$; in this arrangement the electron spins must be paired, so $S=0$. The resulting term is ${ }^{1} \Delta$; a ${ }^{3} \Delta$ term is not possible as it would require two electrons with parallel spins to occupy one of the $\pi$ orbitals. If the electrons occupy different $\pi$ orbitals, $\Lambda=(+1)+(-1)=0$; in this arrangement the spins can be paired or parallel, so $S=0$ or $S=1$. Two further terms therefore arise: ${ }^{1} \Sigma$ and ${ }^{3} \Sigma$; the latter turns out to be the lowest in energy of all the three terms.

As explained in Topic 9B, homonuclear diatomic molecules (but not heteronuclear diatomic molecules) possess a centre of symmetry and their orbitals are labelled g or u according to their parity (the behaviour under inversion through a centre of symmetry). Orbitals that are unchanged upon inversion are $g$ and orbitals that change sign are $u$. Parity labels also apply to centrosymmetric polyatomic linear molecules, such as $\mathrm{CO}_{2}$ and $\mathrm{HC} \equiv \mathrm{CH}$. The overall parity of a configuration of a manyelectron homonuclear diatomic molecule is found by noting the parity of each occupied orbital and using


\begin{equation*}
\mathrm{g} \times \mathrm{g}=\mathrm{g} \quad \mathrm{u} \times \mathrm{u}=\mathrm{g} \quad \mathrm{u} \times \mathrm{g}=\mathrm{u} \tag{11F.2}
\end{equation*}


for each electron. (These rules are generated by interpreting $g$ as +1 and $u$ as -1 .) The resulting parity label $g$ or $u$ is added as a right-subscript to the term symbol. Any molecule in which the occupied orbitals are full (in the sense of being occupied by a pair of electrons) must have overall parity $g$ because there is an even number of electrons. The term symbol for such a homonuclear diatomic molecule is therefore ${ }^{1} \Sigma_{\mathrm{g}}$.

\subsection*{Brief illustration 11F. 2}
Dinitrogen, $\mathrm{N}_{2}$, has the configuration $1 \sigma_{\mathrm{g}}^{2} 1 \sigma_{\mathrm{u}}^{2} 1 \pi_{\mathrm{u}}^{4} 2 \sigma_{\mathrm{g}}^{2}$ in which all the occupied orbitals are full; the same is true of $\mathrm{H}_{2}$ and $\mathrm{F}_{2}$ : all three therefore have the term symbol ${ }^{1} \Sigma_{\mathrm{g}}$.\\
The configuration of $\mathrm{He}_{2}^{+}$is $1 \sigma_{\mathrm{g}}^{2} 1 \sigma_{\mathrm{u}}^{1}$. There is one electron outside the doubly occupied bonding orbital, and the parity of that orbital is $u$. Because $S=\frac{1}{2}$ and $\Lambda=0$, its term symbol is ${ }^{2} \Sigma_{u}$.\\
The ground-state configuration of $\mathrm{O}_{2}$ is $\ldots 1 \pi_{\mathrm{g}}^{2}$. Although the $\pi_{\mathrm{g}}$ orbitals may be both singly occupied, both electrons are in orbitals with $g$ parity, so the overall parity is $g \times g=g$. The three terms arising from this configuration are therefore ${ }^{1} \Sigma_{g}$, ${ }^{3} \Sigma_{\mathrm{g}}$, and ${ }^{1} \Delta_{\mathrm{g}}$ (see Brief illustration 11F.1).

Diatomic molecules (and all linear molecules) possess a mirror plane containing the internuclear axis. All $\sigma$ orbitals (both bonding and antibonding) are symmetric with respect to reflection in this plane. The overall symmetry of a configuration is found by assigning +1 to an electron in a symmetric orbital and -1 to an electron in an orbital that changes sign under reflection and then multiplying the numbers for all the electrons. For example, for the ground state of $\mathrm{H}_{2}$, in which both electrons are in $\sigma$ orbitals, the overall symmetry is $(+1) \times(+1)=+1$. A + sign is added as a right-superscript to the term symbol: ${ }^{1} \Sigma_{g}^{+}$. Any configuration consisting of electrons solely in $\sigma$ orbitals necessarily has + overall reflection symmetry; for example, the ground state of $\mathrm{He}_{2}^{+}$(Brief illustration 11F.2) is ${ }^{2} \Sigma_{\mathrm{u}}^{+}$.

The behaviour of the degenerate pair of $\pi$ molecular orbitals under reflection is more complex: as is shown in Fig. 11F. 2 one of the orbitals changes sign but the other does not. The consequences of this observation can be explored by considering the mathematical form of the $\pi$ orbitals and how they depend on the angle $\phi$ shown in Fig. 11F.3. The orbital $\pi_{x}$ is proportional to $\cos \phi$ and therefore has a nodal plane at $\phi=$ $\pi / 2$ (the $y z$-plane) with positive and negative lobes on either side of this plane; it is unchanged by reflection in the $x z$-plane. The orbital $\pi_{y}$ is proportional to $\sin \phi$, so the $x z$-plane at $\phi=0$ is a nodal plane; the orbital changes sign on reflection in the $x z$ plane. These two wavefunctions are degenerate, so any linear combination of them is also an acceptable wavefunction. For the present discussion the combinations $\pi_{+}=\cos \phi+i \sin \phi=e^{i \phi}$ and $\pi_{-}=\cos \phi-\mathrm{i} \sin \phi=\mathrm{e}^{-\mathrm{i} \phi}$ are convenient for they correspond\\
\includegraphics[max width=\textwidth, center]{2025_02_27_d4b95d9718f0b9e2e6b9g-493(2)}

Figure 11F. 2 A molecular orbital can be classified as symmetric $(+)$ or antisymmetric ( - ) according to whether it changes sign under reflection in a plane containing the internuclear axis.\\
\includegraphics[max width=\textwidth, center]{2025_02_27_d4b95d9718f0b9e2e6b9g-493(1)}

Figure 11F. 3 In a linear molecule, the molecular orbital depends on the azimuthal angle $\phi$. Reflection in the mirror plane is equivalent to reversing the sign of $\phi$.\\
to $\lambda=+1$, and $\lambda=-1$, respectively. As shown in Fig. 11F.3, reflection in the $x z$-plane leaves $\cos \phi$ unchanged but changes the sign of $\sin \phi$. It follows that this reflection interconverts $\pi_{+}$and $\pi_{-}$.

Now consider $\mathrm{O}_{2}$, which has the configuration $\ldots 1 \pi_{\mathrm{g}}^{2}$. The triplet state $(S=1)$, in which the two electrons have parallel spins and necessarily occupy different orbitals, is the state of lowest energy. The triplet spin wavefunction is symmetric with respect to the interchange of the two electrons (it consists of spin states $\alpha(1) \alpha(2)$, and so on), so it follows from the Pauli principle (Topic 8B) that the spatial part of the wavefunction must be antisymmetric with respect to interchange. Such a wavefunction, in which one electron occupies the $\pi_{+}$orbital and the other occupies $\pi_{-}$, is $\Psi_{-}(1,2)=\pi_{+}\left(\boldsymbol{r}_{1}\right) \pi_{-}\left(\boldsymbol{r}_{2}\right)-\pi_{+}\left(\boldsymbol{r}_{2}\right) \pi_{-}\left(\boldsymbol{r}_{1}\right)$. Reflection in the mirror plane gives $\pi_{-}\left(\boldsymbol{r}_{1}\right) \pi_{+}\left(\boldsymbol{r}_{2}\right)-\pi_{-}\left(\boldsymbol{r}_{2}\right) \pi_{+}\left(\boldsymbol{r}_{1}\right)=$ $-\Psi_{+}(1,2)$. That is, the spatial wavefunction of the triplet state is antisymmetric with respect to reflection in the mirror plane, and so a right-superscript - is attached to the term symbol, to give ${ }^{3} \Sigma_{g}^{-}$.

\subsection*{Brief illustration 11F. 3}
An alternative, higher energy configuration of $\mathrm{O}_{2}$ is with the outermost two electrons in separate $\pi$ orbitals but with their spins paired (a ${ }^{1} \Sigma_{\mathrm{g}}$ term). The spin state, which is proportional to $\alpha(1) \beta(2)-\alpha(2) \beta(1)$, is antisymmetric with respect to the interchange of the electrons. The spatial function therefore must be symmetric. A suitable wavefunction is $\Psi_{+}(1,2)=$ $\pi_{+}\left(\boldsymbol{r}_{1}\right) \pi_{-}\left(\boldsymbol{r}_{2}\right)+\pi_{+}\left(\boldsymbol{r}_{2}\right) \pi_{-}\left(\boldsymbol{r}_{1}\right)$. Reflection changes this function to $\pi_{-}\left(\boldsymbol{r}_{1}\right) \pi_{+}\left(\boldsymbol{r}_{2}\right)+\pi_{-}\left(\boldsymbol{r}_{2}\right) \pi_{+}\left(\boldsymbol{r}_{1}\right)$ which is $+\Psi_{+}(1,2)$. The state is symmetric with respect to this reflection, and a superscript + is added to the term symbol, to give ${ }^{1} \Sigma_{g}^{+}$.\\
\includegraphics[max width=\textwidth, center]{2025_02_27_d4b95d9718f0b9e2e6b9g-493}

Figure 11F. 4 The coupling of spin and orbital angular momenta in a linear molecule: only their components along the internuclear axis ( $\Sigma$ and $\Lambda$ ) are well defined.

As for atoms, it is sometimes necessary to specify the total electronic angular momentum, the sum of the orbital and spin contributions, and hence the different 'levels' of a term. In a linear molecule, only the total electronic angular momentum about the internuclear axis is well defined, and is specified by the quantum number $\Omega$ (uppercase omega). For light molecules, where the spin-orbit coupling is weak, $\Omega$ is obtained by adding together the components of orbital angular momentum around the axis (the value of $\Lambda$ ) and the component of the electron spin on that axis (Fig. 11F.4). The latter is denoted $\Sigma$, where $\Sigma=S, S-1, S-2, \ldots,-S$. $^{1}$ Then


\begin{equation*}
\Omega=\Lambda+\Sigma \tag{11F.3}
\end{equation*}


The value of $|\Omega|$ is then be attached to the term symbol as a right subscript (just like $J$ is used in atoms) to denote the different levels. These levels differ in energy, as in atoms, as a result of spin-orbit coupling.

\subsection*{Brief illustration 11F. 4}
The ground-state configuration of NO is $\ldots 1 \pi^{1}$, so it is a ${ }^{2} \Pi$ term with $\Lambda= \pm 1$ and $S=\frac{1}{2}$; from the latter it follows that $\Sigma=$ $\pm \frac{1}{2}$. There are two levels of the term, one with $\Omega= \pm \frac{1}{2}$ and the other with $\pm \frac{3}{2}$, denoted ${ }^{2} \Pi_{1 / 2}$ and ${ }^{2} \Pi_{3 / 2}$, respectively. Each level is doubly degenerate (corresponding to the opposite signs of $\Omega$ ). It turns out that, in $\mathrm{NO},{ }^{2} \Pi_{1 / 2}$ lies slightly lower in energy than ${ }^{2} \Pi_{3 / 2}$.

\subsection*{(b) Selection rules}
A number of selection rules govern which transitions can be observed in the electronic spectrum of a molecule. The selection rules concerned with changes in angular momentum in a linear molecule are

$$
\begin{aligned}
\Delta \Lambda= & 0, \pm 1 \quad \Delta S=0 \quad \Delta \Sigma=0 \quad \Delta \Omega=0, \pm 1 \\
& \text { Selection rules for electronic spectra of linear molecules }
\end{aligned}
$$

(11F.4)

\footnotetext{${ }^{1}$ It is important to distinguish between the upright term symbol $\Sigma$ and the sloping quantum number $\Sigma$.
}As in atoms (Topic 8C), the origins of these rules are conservation of angular momentum during a transition and the fact that a photon has a spin of 1 .\\
Two selection rules can be deduced on the basis of symmetry.

\subsection*{How is that done? 11F. 1}
Establishing symmetry-based\\
selection rules\\
As usual when establishing selection rules, you need to consider properties of the electric-dipole transition moment introduced in Topic $8 \mathrm{C}, \mu_{\mathrm{fi}}=\int \psi_{\mathrm{f}}^{\star} \hat{\mu} \psi_{\mathrm{i}} \mathrm{d} \tau$, and to note that it vanishes unless the integrand is invariant under all symmetry operations of the molecule (Topic 10C).\\
The $z$-component (the component parallel to the axis of the molecules) of the electric dipole moment operator is responsible for $\Sigma \leftrightarrow \Sigma$ transitions (the other components of $\boldsymbol{\mu}$ perpendicular to the axis have $\Pi$ symmetry and cannot make a contribution to this transition). The $z$-component of $\boldsymbol{\mu}$ has $(+)$ symmetry with respect to reflection in a plane containing the internuclear axis. Therefore, for a $(+) \leftrightarrow(-)$ transition, the overall symmetry of the integrand is $(+) \times(+) \times(-)=(-)$, so the integral must be zero and hence $\Sigma^{+} \leftrightarrow \Sigma^{-}$transitions are not allowed. The integrands for $\Sigma^{+} \leftrightarrow \Sigma^{+}$and $\Sigma^{-} \leftrightarrow \Sigma^{-}$transitions transform as $(+) \times(+) \times(+)=(+)$ and $(-) \times(+) \times(-)=(+)$, respectively. The integrals are therefore not necessarily zero and so both transitions are allowed.\\
The three components of the dipole moment operator transform like $x, y$, and $z$, and in a centrosymmetric molecule are all $u$. Therefore, for a $g \rightarrow g$ transition, the overall parity of the integrand is $\mathrm{g} \times \mathrm{u} \times \mathrm{g}=\mathrm{u}$, so the integral must be zero. Likewise, for $a u \rightarrow u$ transition, the overall parity is $u \times u \times u=u$, so the integral is again zero. Hence, transitions without a change of parity are forbidden. For a $\mathrm{g} \leftrightarrow \mathrm{u}$ transition the integrand transforms as $\mathrm{g} \times \mathrm{u} \times \mathrm{u}=\mathrm{g}$, so the transition is allowed.

The first part of this analysis can be summarized as follows:\\
For $\Sigma$ terms, only $\Sigma^{+} \leftrightarrow \Sigma^{+}$and $\Sigma^{-} \leftrightarrow \Sigma^{-}$are allowed.\\
The second part is in fact the Laporte selection rule for centrosymmetric molecules (those with a centre of inversion, not only linear molecules) which states that the only transitions allowed are accompanied by a change of parity. That is,

For centrosymmetric molecules, only $\mathrm{u} \rightarrow \mathrm{g}$ and $\mathrm{g} \rightarrow \mathrm{u}$ transitions are allowed.

Laporte selection rule\\
A forbidden $\mathrm{g} \rightarrow \mathrm{g}$ transition can become allowed if the centre of symmetry is eliminated by an asymmetrical vibration, such as the one shown in Fig. 11F.5. When the centre of symmetry is lost, $\mathrm{g} \rightarrow \mathrm{g}$ and $\mathrm{u} \rightarrow \mathrm{u}$ transitions are no longer parity-forbidden and become weakly allowed. A transition that derives its intensity from an asymmetrical vibration of a molecule is called a vibronic transition.\\
\includegraphics[max width=\textwidth, center]{2025_02_27_d4b95d9718f0b9e2e6b9g-494}

Figure 11F.5 A d-d transition is parity-forbidden because it corresponds to a g-g transition. However, a vibration of the molecule can destroy the inversion symmetry of the molecule and the $\mathrm{g}, \mathrm{u}$ classification no longer applies. The removal of the centre of symmetry gives rise to a vibronically allowed transition.

\subsection*{Brief illustration 11F. 5}
Three possible transitions in the electronic spectrum of $\mathrm{O}_{2}$, ${ }^{3} \Sigma_{\mathrm{g}}^{-} \leftarrow{ }^{3} \Sigma_{\mathrm{u}}^{-},{ }^{3} \Sigma_{\mathrm{g}}^{-} \leftarrow{ }^{1} \Delta_{\mathrm{g}},{ }^{3} \Sigma_{\mathrm{g}}^{-} \leftarrow{ }^{3} \Sigma_{\mathrm{u}}^{+}$, can be considered in the light of the selection rules in eqn 11F. 4 to see which are allowed. A table can be drawn up, in which forbidden values are shown in red.

\begin{center}
\begin{tabular}{lrrlll}
\hline
 & $\Delta S$ & $\Delta \Lambda$ & $\Sigma^{ \pm} \leftarrow \Sigma^{ \pm}$ & \begin{tabular}{l}
Change \\
of parity \\
\end{tabular} &  \\
\hline
${ }^{3} \Sigma_{\mathrm{g}}^{-} \leftarrow{ }^{3} \Sigma_{\mathrm{u}}^{-}$ & 0 & 0 & $\Sigma^{-} \leftarrow \Sigma^{-}$ & $\mathrm{g} \leftarrow \mathrm{u}$ & Allowed \\
${ }^{3} \Sigma_{\mathrm{g}}^{-} \leftarrow{ }^{1} \Delta_{\mathrm{g}}$ & +1 & -2 & Not applicable & $\mathrm{g} \leftarrow \mathrm{g}$ & Forbidden \\
${ }^{3} \Sigma_{\mathrm{g}}^{-} \leftarrow{ }^{3} \Sigma_{\mathrm{u}}^{+}$ & 0 & 0 & $\Sigma^{-} \leftarrow \Sigma^{+}$ & $\mathrm{g} \leftarrow \mathrm{u}$ & Forbidden \\
\hline
\end{tabular}
\end{center}

\subsection*{(c) Vibrational fine structure}
An electronic transition may be accompanied by a simultaneous change in the vibrational state of a molecule, giving rise to vibrational fine structure in the spectrum. In the case of absorption spectra, the transitions are from the ground electronic state, and typically it is only the ground vibrational level, $v^{\prime \prime}=0$, of this state that is occupied significantly. In some cases, the transition from $v^{\prime \prime}=0$ in the electronic ground state to $v^{\prime}=0$ in the upper electronic state is found to be the strongest, with a sharp decline in intensity as $v^{\prime}$ increases. In other cases, transitions with significant intensity to a range of $v^{\prime}$ levels are seen (as in Fig. 11F.1a).

The Franck-Condon principle accounts for the vibrational fine structure in the electronic spectra of molecules:

Because nuclei are so much more massive than electrons, an electronic transition takes place very much faster than the nuclei can respond.

Franck-Condon principle\\
\includegraphics[max width=\textwidth, center]{2025_02_27_d4b95d9718f0b9e2e6b9g-495}

Figure 11F. 6 According to the Franck-Condon principle, the most intense vibronic transition is from the ground vibrational state to the vibrational state lying vertically above it. As a result of the vertical transition, the nuclei suddenly experience a new force field, to which they respond through their vibrational motion. The equilibrium separation of the nuclei in the initial electronic state therefore becomes a turning point in the final electronic state. Transitions to other vibrational levels also occur, but with lower intensity.

The physical basis of this principle is as follows. As a result of the electronic transition, electron density is built up rapidly in new regions of the molecule and removed from others. In classical terms, the initially stationary nuclei suddenly experience a new force field, to which they respond by beginning to vibrate and (in classical terms) swing backwards and forwards from their original separation which was maintained during the rapid electronic excitation. The stationary equilibrium separation of the nuclei in the initial electronic state therefore becomes a stationary turning point in the final electronic state. The transition can be thought of as taking place up the vertical line in Fig. 11F.6. This interpretation is the origin of the expression vertical transition, which denotes an electronic transition that occurs without change of nuclear geometry and, in classical terms, with the nuclei remaining stationary.

Now consider the two potential energy curves shown in Fig. 11F.7a in which the equilibrium bond lengths are the same and initially the molecule is not vibrating. The vertical transition takes place from the minimum of the lower curve, the nuclei remain at the same separation, and ends at the minimum of the upper curve.

Next consider the case shown in Fig. 11F.7b, in which the equilibrium bond length in the upper state is greater than that in the ground electronic state, and the molecule is initially not vibrating. Preservation of the nuclear separation during the transition takes the molecule up the vertical line. The nuclei are not moving initially, and do not start to move during the transition, so the transition terminates at the turning point of the upper electronic state where the nuclei are still stationary.

The quantum mechanical version of the Franck-Condon principle refines this picture. Instead of saying that the nuclei stay at the same locations and are stationary during the transi-\\
\includegraphics[max width=\textwidth, center]{2025_02_27_d4b95d9718f0b9e2e6b9g-495(2)}

Figure 11F. 7 (a) If the equilibrium bond lengths of the ground and excited electronic states are the same, a vertical transition leaves the vibrational state of the molecule unexcited. (b) If the equilibrium bond length is greater in the upper electronic state, the vertical transition ends at a compressed state of the bond and results in vibrational excitation.\\
tion, it replaces that statement by the assertion that the nuclei retain their initial dynamical state. In quantum mechanics, the dynamical state is expressed by the wavefunction, so an equivalent statement is that the vibrational wavefunction does not change during the electronic transition. Initially the molecule is in the lowest vibrational state of its ground electronic state with a bell-shaped wavefunction centred on the equilibrium bond length (Fig. 11F.8). To find the nuclear state to which\\
\includegraphics[max width=\textwidth, center]{2025_02_27_d4b95d9718f0b9e2e6b9g-495(1)}

Figure 11F. 8 (a) If the equilibrium bond lengths of the ground and excited electronic states are the same, the wavefunctions for $v^{\prime \prime}=0$ and $v^{\prime}=0$ are similar and the most probable transition leaves the molecule vibrationally unexcited. (b) If the equilibrium bond length of the upper state is greater than that of the ground state, the wavefunction that most resembles the ground state vibrational wavefunction is that of an excited state. Other transitions also occur with lower intensity. In the inserts, the black curve is the initial vibrational wavefunction and the blue curves are those of the upper electronic state.\\
the transition takes place, it is necessary to look for the vibrational wavefunction of the upper electronic state that most closely resembles this initial wavefunction, for that corresponds to the nuclear dynamical state that is least changed in the transition. That final wavefunction is the one with a large peak close to the position of the initial bell-shaped function. As explained in Topic 7E, provided the vibrational quantum number is not zero, the biggest peaks of vibrational wavefunctions occur close to the edges of the confining potential, so the transition can be expected to occur to those vibrational states, in accord with the classical description. However, several vibrational states have their major peaks in similar positions, so transitions can be expected to occur to a range of vibrational states, to give rise to a vibrational progression, a series of transitions to different vibrational states of the upper electronic state. In a typical progression, the vertical transition is the most intense.

The quantitative version of the Franck-Condon principle involves considering how the transition dipole moment for a given electronic transition varies with the vibrational levels in the two electronic states.

\subsection*{How is that done? 11F. 2 Expressing the Franck-Condon principle quantitatively}
Once again, you need to consider the properties of the electricdipole transition moment. First, note that the electric dipole moment operator is a sum over all nuclei and electrons in the molecule:

$$
\hat{\boldsymbol{\mu}}=-e \sum_{i} r_{i}+e \sum_{N} Z_{N} \boldsymbol{R}_{N}
$$

where the origin of the vectors is the centre of charge of the molecule, $i$ labels the electrons, and $N$ labels the nuclei. Within the Born-Oppenheimer approximation (the separation of electronic and vibrational motion, the Prologue to Focus 9), the overall state of the molecule consists of an electronic contribution, labelled $\varepsilon$, and a vibrational contribution, labelled $v$. Therefore, the transition dipole moment factorizes as follows:

$$
\begin{aligned}
\mu_{\mathrm{fi}}= & \int \psi_{\varepsilon, \mathrm{f}}^{*} \psi_{v, \mathrm{f}}^{*}\left\{-e \sum_{i} \boldsymbol{r}_{i}+e \sum_{N} Z_{N} \boldsymbol{R}_{N}\right\} \psi_{\varepsilon, \mathrm{i}} \psi_{v, \mathrm{i}} \mathrm{~d} \tau \\
= & -e \sum_{i} \int \psi_{\varepsilon, \mathrm{f}}^{*} \boldsymbol{r}_{i} \psi_{\varepsilon, \mathrm{i}} \mathrm{~d} \tau_{\mathrm{e}} \int \psi_{v, \mathrm{f}}^{*} \psi_{v, \mathrm{i}} \mathrm{~d} \tau_{\mathrm{N}} \\
& +e \sum_{N} Z_{N} \overbrace{\int \psi_{\varepsilon, \mathrm{f}}^{*} \psi_{\varepsilon, \mathrm{i}} \mathrm{~d} \tau_{\mathrm{e}} \int}^{0} \psi_{v, \mathrm{f}}^{*} \boldsymbol{R}_{N} \psi_{v, \mathrm{i}} \mathrm{~d} \tau_{\mathrm{N}}
\end{aligned}
$$

where $\mathrm{d} \tau_{\mathrm{e}}$ indicates integration over the electronic coordinates, and $\mathrm{d} \tau_{\mathrm{N}}$ integration over the nuclear coordinates. Because the two different electronic states are orthogonal\\
(they are eigenstates of the same hamiltonian but correspond to different eigenvalues) the integral in blue is zero, which leaves

$$
\mu_{\mathrm{fi}}=-e \sum_{i} \overbrace{\int \psi_{\varepsilon, \mathrm{f}}^{*} r_{i} \psi_{\varepsilon, \mathrm{i}} \mathrm{~d} \tau_{\mathrm{e}} \int \overbrace{\int \psi_{v, \mathrm{f}} \psi_{v, \mathrm{i}}}^{\mu_{\mathrm{e}}} \mathrm{~d} \tau_{\mathrm{N}}}^{S\left(v_{\mathrm{v}}, v_{\mathrm{i}}\right)}=\mu_{\varepsilon, \mathrm{f}} S\left(v_{\mathrm{f}}, v_{\mathrm{i}}\right)
$$

The quantity $\boldsymbol{\mu}_{\varepsilon, f i}$ is the electric-dipole transition moment arising from the change in the electronic wavefunction: this term describes the interaction of the electrons with the electromagnetic field. The factor $S\left(v_{\mathrm{f}}, v_{\mathrm{i}}\right)$, is the overlap integral between the vibrational level with quantum number $v_{\mathrm{i}}$ in the initial electronic state of the molecule, and the vibrational level with quantum number $v_{\mathrm{f}}$ in the final electronic state of the molecule.

The transition intensity is proportional to the square of the magnitude of the transition dipole moment, so is proportional to the square of $S\left(v_{\mathrm{f}}, v_{\mathrm{i}}\right)$, and specifically


\begin{equation*}
\left.S\left(v_{\mathrm{f}}, v_{\mathrm{i}}\right)\right|^{2}=\left(\int \psi_{v, \mathrm{f}}^{\star} \psi_{\nu, \mathrm{i}} \mathrm{~d} \tau_{\mathrm{N}}\right)^{2} \quad \text { Franck-Condon factor } \tag{11F.5}
\end{equation*}


The integral on the right-hand side of eqn 11F. 5 is the overlap between the two vibrational wavefunctions: the greater this overlap (physically, the greater the resemblance of the vibrational wavefunctions), the greater is the intensity of the transition.

\subsection*{Example 11F. 1}
\subsection*{Calculating a Franck-Condon factor}
Consider the transition from one electronic state to another, their equilibrium bond lengths being $R_{\mathrm{e}}$ and $R_{e}^{\prime}$, and their force constants equal. Calculate the Franck-Condon factor for the $v^{\prime \prime}=0$ to $v^{\prime}=0$ transition (the $0-0$ transition) and show that the transition is most intense when the bond lengths are equal.\\
Collect your thoughts You need to calculate $S(0,0)$, the overlap integral of the two ground-state vibrational wavefunctions, and then take its square. The difference between harmonic and anharmonic vibrational wavefunctions is negligible for $v=0$, so it is safe for you to use the harmonic oscillator wavefunctions.

The solution The (real) wavefunctions are (Topic 7E)

$$
\psi_{0}=\left(\frac{1}{\alpha \pi^{1 / 2}}\right)^{1 / 2} \mathrm{e}^{-x^{2} / 2 \alpha^{2}} \quad \psi_{0}^{\prime}=\left(\frac{1}{\alpha \pi^{1 / 2}}\right)^{1 / 2} \mathrm{e}^{-x^{2} / 2 \alpha^{2}}
$$

where $x=R-R_{\mathrm{e}}$ and $x^{\prime}=R-R_{\mathrm{e}}^{\prime}$, with $\alpha=\left(\hbar^{2} / \mu k_{\mathrm{f}}\right)^{1 / 4}$. The overlap integral is

$$
S(0,0)=\int_{-\infty}^{\infty} \psi_{0}^{\prime} \psi_{0} \mathrm{~d} R=\frac{1}{\alpha \pi^{1 / 2}} \int_{-\infty}^{\infty} \mathrm{e}^{-\left(x^{2}+x^{2}\right) / 2 \alpha^{2}} \mathrm{~d} x
$$

Now recognize that

$$
\begin{aligned}
x^{2}+x^{\prime 2} & =\left(R-R_{\mathrm{e}}\right)^{2}+\left(R-R_{\mathrm{e}}^{\prime}\right)^{2} \\
& =2 R^{2}+R_{\mathrm{e}}^{2}+R_{\mathrm{e}}^{\prime 2}-2 R\left(R_{\mathrm{e}}+R_{\mathrm{e}}^{\prime}\right) \\
& =2\left\{R-\frac{1}{2}\left(R_{\mathrm{e}}+R_{\mathrm{e}}^{\prime}\right)\right\}^{2}+\frac{1}{2}\left(R_{\mathrm{e}}-R_{\mathrm{e}}^{\prime}\right)^{2}
\end{aligned}
$$

and write $\alpha z=R-\frac{1}{2}\left(R_{\mathrm{e}}+R_{\mathrm{e}}^{\prime}\right)$, so

$$
\frac{x^{2}+x^{\prime 2}}{2 \alpha^{2}}=z^{2}+\frac{\left(R_{e}-R_{e}^{\prime}\right)^{2}}{4 \alpha^{2}}
$$

and $\mathrm{d} R=\alpha \mathrm{d} z$. Then

$$
S(0,0)=\frac{1}{\pi^{1 / 2}} \mathrm{e}^{-\left(R_{e}-R_{e}^{\prime}\right)^{2} / 4 \alpha^{\alpha^{2}}} \overbrace{\int_{-\infty}^{\infty} \mathrm{e}^{-z^{2}} \mathrm{~d} z=\mathrm{e}^{-\left(R_{e}-R_{e}^{\prime}\right)^{2} / 4 \alpha^{2}}}^{\text {Integral G. } 1 \pi^{1 / 2}}
$$

and the Franck-Condon factor is

$$
S(0,0)^{2}=\mathrm{e}^{-\left(R_{e}-R_{e}^{\prime}\right)^{2} / 2 \alpha^{2}}
$$

This factor is equal to 1 when $R_{\mathrm{e}}^{\prime}=R_{\mathrm{e}}$ and decreases as the equilibrium bond lengths diverge from each other (Fig. 11F.9).\\
\includegraphics[max width=\textwidth, center]{2025_02_27_d4b95d9718f0b9e2e6b9g-497(1)}

Figure 11F. 9 The Franck-Condon factor for the arrangement discussed in Example 11F.1.

For ${ }^{79} \mathrm{Br}_{2}, R_{\mathrm{e}}=228 \mathrm{pm}$ and there is an upper state with $R_{\mathrm{e}}^{\prime}=$ 266 pm . Taking the vibrational wavenumber as $250 \mathrm{~cm}^{-1}$ gives $\alpha^{2}=3.42 \times 10^{-23} \mathrm{~m}^{2}$ and hence $S(0,0)^{2}=6.7 \times 10^{-10}$, so the intensity of the $0-0$ transition is only $6.7 \times 10^{-10}$ what it would have been if the potential curves had been directly above each other.

Self-test 11F. 1 Suppose the normalized vibrational wavefunctions can be approximated by rectangular functions of width $W$ and $W^{\prime}$, centred on the equilibrium bond lengths (Fig. 11F.10). Find the corresponding Franck-Condon factors when the centres are coincident and $W^{\prime}<W$.\\
\includegraphics[max width=\textwidth, center]{2025_02_27_d4b95d9718f0b9e2e6b9g-497}

Figure 11F. 10 The model wavefunctions used in Self-test 11F.1.

$$
M /, M={ }_{\tau} S: \swarrow \partial M s u_{V}
$$

\subsection*{(d) Rotational fine structure}
Electronic transitions may be accompanied by simultaneous changes in both vibrational and rotational energy. Therefore, when the lines dues to vibrational fine structure are inspected at higher resolution they are found to have rotational fine structure and to consist of $\mathrm{P}, \mathrm{Q}$, and R branches of the type discussed in Topic 11C. Because electronic excitation can result in much larger changes in bond length than vibrational excitation causes alone, the rotational branches have a more complex structure than in vibration-rotation spectra.

The rotational constants of the electronic ground and excited states are denoted $\tilde{B}$ and $\tilde{B}^{\prime}$, respectively. The rotational terms of the initial and final states are


\begin{equation*}
\tilde{F}(J)=\tilde{B} J(J+1) \quad \tilde{F}\left(J^{\prime}\right)=\tilde{B}^{\prime} J^{\prime}\left(J^{\prime}+1\right) \tag{11F.6}
\end{equation*}


When a transition occurs with $\Delta J=-1$ the wavenumber of the vibrational contribution to the electronic transition is shifted from $\tilde{v}$ to

$$
\tilde{v}+\tilde{B}^{\prime}(J-1) J-\tilde{B} J(J+1)=\tilde{v}-\left(\tilde{B}^{\prime}+\tilde{B}\right) J+\left(\tilde{B}^{\prime}-\tilde{B}\right) J^{2}
$$

This transition is a contribution to the P branch (just as in Topic 11C). There are corresponding transitions for the Q and R branches with wavenumbers that may be calculated in a similar way. All three branches are:


\begin{align*}
& \mathrm{P} \text { branch }(\Delta J=-1): \tilde{v}_{\mathrm{P}}(J)=\tilde{v}-\left(\tilde{B}^{\prime}+\tilde{B}\right) J+\left(\tilde{B}^{\prime}-\tilde{B}\right) J^{2}  \tag{11F.7a}\\
& \mathrm{Q} \text { branch }(\Delta J=0): \tilde{v}_{\mathrm{Q}}(J)=\tilde{v}+\left(\tilde{B}^{\prime}-\tilde{B}\right) J(J+1) \tag{11F.7b}
\end{align*}



\begin{equation*}
\mathrm{R} \text { branch }(\Delta J=+1): \tilde{v}_{\mathrm{R}}(J)=\tilde{v}+\left(\tilde{B}^{\prime}+\tilde{B}\right)(J+1)+\left(\tilde{B}^{\prime}-\tilde{B}\right)(J+1)^{2} \tag{11F.7c}
\end{equation*}


\subsection*{Example 11F. 2 Estimating rotational constants from \\
 electronic spectra}
The following rotational transitions were observed in the 0-0 band of the ${ }^{1} \Sigma^{+} \leftarrow{ }^{1} \Sigma^{+}$electronic transition of ${ }^{63} \mathrm{Cu}^{2} \mathrm{H}: \tilde{v}_{\mathrm{R}}(3)=$ $23347.69 \mathrm{~cm}^{-1}, \tilde{v}_{\mathrm{P}}(3)=23298.85 \mathrm{~cm}^{-1}$, and $\tilde{v}_{\mathrm{P}}(5)=23275.77 \mathrm{~cm}^{-1}$. Estimate the values of $\tilde{B}^{\prime}$ and $\tilde{B}$.

Collect your thoughts You need to be aware that Topic 11C introduces a procedure, the method of 'combination differences', for analysing transitions that have a common state. According to that method, form the differences $\tilde{v}_{\mathrm{R}}(J)-\tilde{v}_{\mathrm{P}}(J)$ and $\tilde{v}_{\mathrm{R}}(J-1)-\tilde{v}_{\mathrm{P}}(J+1)$ from eqns 11F.7a and 11F.7c, then use the resulting expressions to calculate the rotational constants $\tilde{B}^{\prime}$ and $\tilde{B}$ from the data provided.

The solution From eqns 11F.7a and 11F.7c it follows that

$$
\begin{aligned}
\tilde{V}_{\mathrm{R}}(J)-\tilde{V}_{\mathrm{P}}(J)= & \left(\tilde{B}^{\prime}+\tilde{B}\right)(J+1)+\left(\tilde{B}^{\prime}-\tilde{B}\right)(J+1)^{2} \\
& -\left\{-\left(\tilde{B}^{\prime}+\tilde{B}\right) J+\left(\tilde{B}^{\prime}-\tilde{B}\right) J^{2}\right\}=4 \tilde{B}^{\prime}\left(J+\frac{1}{2}\right) \\
\tilde{V}_{\mathrm{R}}(J-1)-\tilde{v}_{\mathrm{P}}(J+1)= & \left(\tilde{B}^{\prime}+\tilde{B}\right) J+\left(\tilde{B}^{\prime}-\tilde{B}\right) J^{2} \\
& -\left\{-\left(\tilde{B}^{\prime}+\tilde{B}\right)(J+1)+\left(\tilde{B}^{\prime}-\tilde{B}\right)(J+1)^{2}\right\} \\
= & 4 \tilde{B}\left(J+\frac{1}{2}\right)
\end{aligned}
$$

(These equations are analogous to eqn 11C.14.) The data provided can be used in the following way:

$$
\text { For } J=3: \quad \tilde{v}_{\mathrm{R}}(3)-\tilde{v}_{\mathrm{P}}(3)=\overbrace{48.84 \mathrm{~cm}^{-1}=14 \tilde{B}^{\prime}}^{23347.69-23298.85}
$$

23347.69-23275.77

For $J=4: \quad \tilde{V}_{\mathrm{R}}(3)-\tilde{V}_{\mathrm{p}}(5)=\overbrace{71.92} \mathrm{~cm}^{-1}=18 \tilde{B}$

Hence $\tilde{B}^{\prime}=3.489 \mathrm{~cm}^{-1}$ and $\tilde{B}=3.996 \mathrm{~cm}^{-1}$.\\
Self-test 11F. 2 The following rotational transitions were observed in the ${ }^{1} \Sigma^{+} \leftarrow{ }^{1} \Sigma^{+}$electronic transition of RhN: $\tilde{V}_{\mathrm{R}}(5)=$ $22387.06 \mathrm{~cm}^{-1}, \tilde{v}_{\mathrm{P}}(5)=22376.87 \mathrm{~cm}^{-1}$, and $\tilde{v}_{\mathrm{P}}(7)=22373.95 \mathrm{~cm}^{-1}$. Estimate the values of $\tilde{B}^{\prime}$ and $\tilde{B}$.\\
\includegraphics[max width=\textwidth, center]{2025_02_27_d4b95d9718f0b9e2e6b9g-498}

Suppose that the bond length in the electronically excited state is greater than that in the ground state; it follows that $\tilde{B}^{\prime}<\tilde{B}$ and hence $\tilde{B}^{\prime}-\tilde{B}<0$. In this case the lines of the R branch converge with increasing $J$ and at sufficiently high values of $J$ the negative term in $(J+1)^{2}$ in eqn 11F.7c will dominate the positive term in $(J+1)$ and the lines will start to appear at successively decreasing wavenumbers. That is, the R branch has a band head (Fig. 11F.11a).\\
\includegraphics[max width=\textwidth, center]{2025_02_27_d4b95d9718f0b9e2e6b9g-498(1)}

Frequency, $v \longrightarrow$\\
(a) $\widetilde{B}^{\prime}<\widetilde{B}$\\
\includegraphics[max width=\textwidth, center]{2025_02_27_d4b95d9718f0b9e2e6b9g-498(2)}\\
(b) $\widetilde{B}>\widetilde{B}$

Figure 11F. 11 (a) The formation of a head in the R branch when $\tilde{B}^{\prime}<\tilde{B} ;$ (b) the formation of a head in the P branch when $\tilde{B}^{\prime}>\tilde{B}$. The red curve shows the wavenumbers of the lines in the two branches as they spread away from the centre of the band.

The value of $J$ at which the band head appears can be found by finding the maximum wavenumber of the lines in the R branch, in other words an integral value of $J$ close to where $\mathrm{d} \tilde{v}_{\mathrm{R}}(J) / \mathrm{d} J=0$. This maximum occurs at $J_{\max } \approx\left(\tilde{B}-3 \tilde{B}^{\prime}\right) / 2\left(\tilde{B}^{\prime}-\tilde{B}\right)$. When the bond is shorter in the excited state than in the ground state, $\tilde{B}^{\prime}>\tilde{B}$ and $\tilde{B}^{\prime}-\tilde{B}>0$. In this case, the lines of the P branch begin to converge and form a band head when $J_{\max } \approx\left(\tilde{B}^{\prime}+\tilde{B}\right) / 2\left(\tilde{B}^{\prime}-\tilde{B}\right)$, as shown in Fig. 11F.11b.

\subsection*{Brief illustration 11F. 6}
For the transition described in Example 11F.2, $\tilde{B}^{\prime}<\tilde{B}$, so a band head is expected in the R branch. The approximate value of $J$ at which this occurs is given by

$$
J_{\max } \approx \frac{\tilde{B}-3 \tilde{B}^{\prime}}{2\left(\tilde{B}^{\prime}-\tilde{B}\right)}=\frac{3.996 \mathrm{~cm}^{-1}-3 \times 3.489 \mathrm{~cm}^{-1}}{2\left(3.489 \mathrm{~cm}^{-1}-3.996 \mathrm{~cm}^{-1}\right)}=6.38
$$

The closest integer is $J=6$.

\subsection*{11F. 2 Polyatomic molecules}
The absorption of a photon can often be traced to the excitation of electrons that belong to a small group of atoms in a polyatomic molecule. For example, when a carbonyl group ( $\mathrm{C}=\mathrm{O}$ ) is present, an absorption at about 290 nm is normally observed, although its precise location depends on the nature of the rest of the molecule. Groups with characteristic optical absorption bands are called chromophores (from the Greek for 'colour bringer'), and their presence often accounts for the colours of substances (Table 11F.2).

Table 11F.2 Absorption characteristics of some groups and molecules*

\begin{center}
\begin{tabular}{llll}
\hline
Group & $\tilde{v} / \mathrm{cm}^{-1}$ & $\lambda_{\max } / \mathrm{nm}$ & $\varepsilon_{\max } /\left(\mathrm{dm}^{3} \mathrm{~mol}^{-1} \mathrm{~cm}^{-1}\right)$ \\
\hline
$\mathrm{C}=\mathrm{C}\left(\pi^{*} \leftarrow \pi\right)$ & 61000 & 163 & 15000 \\
$\mathrm{C}=\mathrm{O}\left(\pi^{*} \leftarrow \mathrm{n}\right)$ & $35000-37000$ & $270-290$ & $10-20$ \\
$\mathrm{H}_{2} \mathrm{O}$ & 60000 & 167 & 7000 \\
\end{tabular}
\end{center}

${ }^{*}$ More values are given in the Resource section. $\varepsilon_{\text {max }}$ is the molar absorption coefficient (see Topic 11A). The wavenumbers and wavelengths are the values for maximum absorption.

\subsection*{(a) d-Metal complexes}
In a free atom, all five $d$ orbitals of a given shell are degenerate. In a d-metal complex, where the immediate environment of the atom is no longer spherical, the d orbitals are not all degenerate, and electrons can absorb energy by making transitions between them.

To see the origin of this splitting in an octahedral complex such as $\left[\mathrm{Ti}\left(\mathrm{OH}_{2}\right)_{6}\right]^{3+}(1)$, the six ligands can be regarded as point negative charges that repel the delectrons of the central ion (Fig. 11F.12). As a result, the orbitals fall into two groups, with $\mathrm{d}_{x^{2}-y^{2}}$ and $\mathrm{d}_{z^{2}}$ pointing directly towards the ligand positions, and $\mathrm{d}_{x y}, \mathrm{~d}_{y z}$, and $\mathrm{d}_{z x}$ pointing between them. An electron occupying an orbital of the former group has a less favourable potential energy than when it occupies any of the three orbitals of the other group, and so the d orbitals split into the two sets shown in (2): a triply degenerate set comprising the $\mathrm{d}_{x y}$, $\mathrm{d}_{y z}$, and $\mathrm{d}_{z x}$ orbitals and labelled $\mathrm{t}_{2 \mathrm{~g}}$, and a doubly degenerate set comprising the $\mathrm{d}_{x^{2}-y^{2}}$ and $\mathrm{d}_{z^{2}}$ orbitals and labelled $\mathrm{e}_{\mathrm{g}}$ (these symmetry labels are discussed in Topic 10B). The $\mathrm{t}_{2 \mathrm{~g}}$ orbitals lie below the $e_{g}$ orbitals in energy; the difference in energy $\Delta_{o}$ is called the ligand-field splitting parameter (the o denotes octahedral symmetry). The ligand field splitting is typically about 10 per cent of the overall energy of interaction between the ligands and the central metal atom, which is largely responsible for the existence of the complex. The d orbitals also divide into two sets in a tetrahedral complex, but in this case the two e orbitals lie below the three $t_{2}$ orbitals (no $g$ or $u$ label is given because a tetrahedral complex has no centre of inversion); the separation of these groups of orbitals is written $\Delta_{t}$.\\
\includegraphics[max width=\textwidth, center]{2025_02_27_d4b95d9718f0b9e2e6b9g-499}\\
$1\left[\mathrm{Ti}\left(\mathrm{OH}_{2}\right)_{6}\right]^{3+}$\\
\includegraphics[max width=\textwidth, center]{2025_02_27_d4b95d9718f0b9e2e6b9g-499(2)}

2\\
\includegraphics[max width=\textwidth, center]{2025_02_27_d4b95d9718f0b9e2e6b9g-499(3)}

Figure 11F. 12 The classification of $d$ orbitals in an octahedral environment. The open circles represent the positions of the six (point-charge) ligands.

The values of $\Delta_{\mathrm{o}}$ and $\Delta_{\mathrm{t}}$ are such that transitions between the two sets of orbitals typically occur in the visible region of the spectrum. The transitions are responsible for many of the colours that are so characteristic of d-metal complexes.

\subsection*{Brief illustration 11F. 7}
The spectrum of $\left[\mathrm{Ti}\left(\mathrm{OH}_{2}\right)_{6}\right]^{3+}$ near $24000 \mathrm{~cm}^{-1}(500 \mathrm{~nm})$ is shown in Fig. 11F.13, and can be ascribed to the promotion of its single $d$ electron from a $t_{2 g}$ orbital to an $e_{g}$ orbital. The wavenumber of the absorption maximum suggests that $\Delta_{\mathrm{o}} \approx 24000 \mathrm{~cm}^{-1}$ for this complex, which corresponds to about 3.0 eV .\\
\includegraphics[max width=\textwidth, center]{2025_02_27_d4b95d9718f0b9e2e6b9g-499(1)}

Figure 11F. 13 The electronic absorption spectrum of $\left[\mathrm{Ti}\left(\mathrm{OH}_{2}\right)_{6}\right]^{3+}$ in aqueous solution.

According to the Laporte rule (Section 11F.1b), d-d transitions are parity-forbidden in octahedral complexes because they are $\mathrm{g} \rightarrow \mathrm{g}$ transitions (more specifically, $\mathrm{e}_{\mathrm{g}} \leftarrow \mathrm{t}_{2 \mathrm{~g}}$ transitions).

However, d-d transitions become weakly allowed as vibronic transitions as a result of coupling to asymmetrical vibrations such as that shown in Fig. 11F.5.

A d-metal complex may also absorb radiation as a result of the transfer of an electron from the ligands into the d orbitals of the central atom, or vice versa. In such charge-transfer transitions the electron moves through a considerable distance, which means that the transition dipole moment may be large and the absorption correspondingly intense. In the permanganate ion, $\mathrm{MnO}_{4}^{-}$, the charge redistribution that accompanies the migration of an electron from the O atoms to the central Mn atom results in a strong transition in the range $475-575 \mathrm{~nm}$ that accounts for the intense purple colour of the ion. Such an electronic migration from the ligands to the metal corresponds to a ligand-to-metal charge-transfer transition (LMCT). The reverse migration, a metal-to-ligand chargetransfer transition (MLCT), can also occur. An example is the migration of a d electron onto the antibonding $\pi$ orbitals of an aromatic ligand. The resulting excited state may have a very long lifetime if the electron is extensively delocalized over several aromatic rings.

In common with other transitions, the intensities of chargetransfer transitions are proportional to the square of the transition dipole moment. The transition moment can be thought of as a measure of the distance moved by the electron as it migrates from metal to ligand or vice versa, with a large distance of migration corresponding to a large transition dipole moment and therefore a high intensity of absorption. However, when calculating the transition dipole moment the integrand is proportional to the product of the initial and final wavefunctions; this product is zero unless the two wavefunctions have non-zero values in the same region of space. Therefore, although large distances of migration favour high intensities, the diminished overlap of the initial and final wavefunctions for large separations of metal and ligands favours low intensities (see Problem P11F.9).

\subsection*{(b) $\pi^{\star} \leftarrow \pi$ and $\pi^{\star} \leftarrow \mathbf{n}$ transitions}
Absorption by a $\mathrm{C}=\mathrm{C}$ double bond results in the excitation of a $\pi$ electron into an antibonding $\pi^{*}$ orbital (Fig. 11F.14). The chromophore activity is therefore due to a $\pi^{\star} \leftarrow \pi$ transition (a ' $\pi$ to $\pi$-star transition'). Its energy is about 6.9 eV for an unconjugated double bond, which corresponds to an absorption at 180 nm (in the ultraviolet). When the double bond is part of a conjugated chain, the energies of the molecular orbitals lie closer together and the $\pi^{\star} \leftarrow \pi$ transition moves to a longer wavelength; it may even lie in the visible region if the conjugated system is long enough.\\
\includegraphics[max width=\textwidth, center]{2025_02_27_d4b95d9718f0b9e2e6b9g-500}

Figure 11F. 14 A C=C double bond acts as a chromophore. In the $\pi^{*} \leftarrow \pi$ transition illustrated here, an electron is promoted from a $\pi$ orbital to the corresponding antibonding orbital.\\
\includegraphics[max width=\textwidth, center]{2025_02_27_d4b95d9718f0b9e2e6b9g-500(1)}

Figure 11F. 15 A carbonyl group ( $\mathrm{C}=\mathrm{O}$ ) acts as a chromophore partly on account of the excitation of a nonbonding O lonepair electron to an antibonding CO $\pi^{*}$ orbital; this transition is denoted $\pi^{*} \leftarrow \mathrm{n}$.

One of the transitions responsible for absorption in carbonyl compounds can be traced to the lone pairs of electrons on the O atom. The Lewis concept of a 'lone pair' of electrons is represented in molecular orbital theory by a pair of electrons in an orbital confined largely to one atom and not appreciably involved in bond formation. One of these electrons may be excited into an empty $\pi^{\star}$ orbital of the carbonyl group (Fig. 11F.15), which gives rise to an $\pi^{*} \leftarrow \mathbf{n}$ transition (an ' $n$ to $\pi$-star transition'). Typical absorption energies are about $4.3 \mathrm{eV}(290 \mathrm{~nm})$. Because $\pi^{\star} \leftarrow \mathrm{n}$ transitions in carbonyls are symmetry forbidden, the absorptions are weak. By contrast, the $\pi^{*} \leftarrow \pi$ transition in a carbonyl, which corresponds to excitation of a $\pi$ electron of the $\mathrm{C}=\mathrm{O}$ double bond, is allowed by symmetry and results in relatively strong absorption.

\subsection*{Brief illustration 11F. 8}
The compound $\mathrm{CH}_{3} \mathrm{CH}=\mathrm{CHCHO}$ has a strong absorption in the ultraviolet at $46950 \mathrm{~cm}^{-1}(213 \mathrm{~nm})$ and a weak absorption at $30000 \mathrm{~cm}^{-1}(330 \mathrm{~nm})$. The former is a $\pi^{\star} \leftarrow \pi$ transition associated with the delocalized $\pi$ system $\mathrm{C}=\mathrm{C}-\mathrm{C}=\mathrm{O}$. Delocalization extends the range of the $\mathrm{C}=\mathrm{O} \pi^{\star} \leftarrow \pi$ transition to lower wavenumbers (longer wavelengths). The latter is an $\pi^{*} \leftarrow \mathrm{n}$ transition associated with the carbonyl chromophore.

\subsection*{Checklist of concepts}
$\square$ 1. The term symbols of linear molecules give the components of various kinds of angular momentum around the internuclear axis along with relevant symmetry labels.2. The Laporte selection rule states that, for centrosymmetric molecules, only $\mathrm{u} \rightarrow \mathrm{g}$ and $\mathrm{g} \rightarrow \mathrm{u}$ transitions are allowed.3. The Franck-Condon principle asserts that electronic transitions occur within an unchanging nuclear framework.4. Vibrational fine structure is the structure in a spectrum that arises from changes in vibrational energy accompanying an electronic transition.5. Rotational fine structure is the structure in a spectrum that arises from changes in rotational energy accompanying an electronic transition.6. In gas phase samples, rotational fine structure can be resolved and under some circumstances band heads are formed.7. In d-metal complexes, the presence of ligands removes the degeneracy of $d$ orbitals and vibrationally allowed d-d transitions can occur between them.8. Charge-transfer transitions typically involve the migration of electrons between the ligands and the central metal atom.\\
9. A chromophore is a group with a characteristic optical absorption band.

\subsection*{Checklist of equations}
\begin{center}
\begin{tabular}{llll}
\hline
Property & Equation & Comment & Equation number \\
\hline
Selection rules (angular momentum) & $\Delta \Lambda=0, \pm 1 ; \Delta S=0 ; \Delta \Sigma=0 ; \Delta \Omega=0, \pm 1$ & Linear molecules & 11 F .4 \\
Franck-Condon factor & $\left|S\left(v_{\mathrm{f}}, v_{\mathrm{i}}\right)\right|^{2}=\left(\int \psi_{v, \mathrm{f}}^{*} \psi_{v, \mathrm{~s}} \mathrm{~d} \tau_{\mathrm{N}}\right)^{2}$ &  & 11 F .5 \\
\begin{tabular}{l}
Rotational structure of electronic spectra \\
(diatomic molecules) \\
\end{tabular} & $\tilde{v}_{\mathrm{p}}(J)=\tilde{v}-\left(\tilde{B}^{\prime}+\tilde{B}\right) J+\left(\tilde{B}^{\prime}-\tilde{B}\right) J^{2}$ & P branch $(\Delta J=-1)$ & 11 F .7 a \\
 & $\tilde{v}_{\mathrm{Q}}(J)=\tilde{v}+\left(\tilde{B}^{\prime}-\tilde{B}\right) J(J+1)$ & Q branch $(\Delta J=0)$ & 11 F .7 b \\
 & $\tilde{v}_{\mathrm{R}}(J)=\tilde{v}+\left(\tilde{B}^{\prime}+\tilde{B}\right)(J+1)+\left(\tilde{B}^{\prime}-\tilde{B}\right)(J+1)^{2}$ & R branch $(\Delta J=+1)$ & 11 F .7 c \\
\hline
\end{tabular}
\end{center}

\section*{TOPIC 11G Decay of excited states}
\subsection*{Why do you need to know this material?}
Much information about the electronic structure of a molecule can be obtained by observing the radiative decay of excited electronic states back to the ground state. Such decay is also used in lasers, which are of exceptional technological importance.

\subsection*{What is the key idea?}
Molecules in excited electronic states discard their excess energy by emission of electromagnetic radiation, transfer as heat to the surroundings, or fragmentation.

\subsection*{- What do you need to know already?}
You need to be familiar with electronic transitions in molecules (Topic 11F), the difference between spontaneous and stimulated emission of radiation (Topic 11A), and the general features of spectroscopy (Topic 11A). You need to be aware of the difference between singlet and triplet states (Topic 8C) and of the Franck-Condon principle (Topic 11F).

Radiative decay is a process in which a molecule discards its excitation energy as a photon (Topic 11A); depending on the nature of the excited state, this process is classified as either fluorescence or phosphorescence. A more common fate of an electronically excited molecule is non-radiative decay, in which the excess energy is transferred into the vibration, rotation, and translation of the surrounding molecules. This thermal degradation converts the excitation energy into thermal motion of the environment (i.e. to 'heat'). An excited molecule may also dissociate or take part in a chemical reaction (Topic 17G). Stimulated emission from excited states is the key process that can lead to laser action.

\subsection*{11G. 1 Fluorescence and phosphorescence}
In fluorescence, spontaneous emission of radiation occurs while the sample is being irradiated and ceases as soon as the exciting radiation is extinguished (Fig. 11G.1). In phosphorescence, the spontaneous emission may persist for long periods (even hours, but more commonly seconds or fractions\\
\includegraphics[max width=\textwidth, center]{2025_02_27_d4b95d9718f0b9e2e6b9g-502}

Figure 11G. 1 The empirical (observation-based) distinction between fluorescence and phosphorescence is that the former is extinguished very quickly after the exciting radiation is removed, whereas the latter continues with relatively slowly diminishing intensity.\\
of seconds). The difference suggests that fluorescence is a fast conversion of absorbed radiation into re-emitted energy, and that phosphorescence involves the storage of energy in a reservoir from which it slowly leaks.

Figure 11G. 2 shows the sequence of steps involved in fluorescence from molecules in solution. The initial stimulated absorption takes the molecule to an excited electronic state; if the absorption spectrum were monitored it would look like the\\
\includegraphics[max width=\textwidth, center]{2025_02_27_d4b95d9718f0b9e2e6b9g-502(1)}

Figure 11G. 2 The sequence of steps leading to fluorescence by molecules in solution. After the initial absorption, the upper vibrational states undergo radiationless decay by giving up energy to the surrounding molecules. A radiative transition then occurs from the vibrational ground state of the upper electronic state.\\
\includegraphics[max width=\textwidth, center]{2025_02_27_d4b95d9718f0b9e2e6b9g-503}

Figure 11G.3 An absorption spectrum (blue) shows vibrational structure characteristic of the upper electronic state. A fluorescence spectrum (purple) shows structure characteristic of the lower state; it is also displaced to lower frequencies (but the $0-0$ transitions are coincident) and is often a mirror image of the absorption spectrum.\\
one shown in Fig. 11G.3a. The excited molecule is subjected to collisions with the surrounding molecules, and as it gives up energy to them non-radiatively it steps down (typically within picoseconds) the ladder of vibrational levels to the lowest vibrational level of the excited electronic state. The surrounding molecules, however, might now be unable to accept the larger energy difference needed to lower the molecule to the ground electronic state. The excited electronic state might therefore survive long enough to undergo spontaneous emission and emit the remaining excess energy as radiation. The downward electronic transition is vertical, in accord with the Franck-Condon principle (Topic 11F), and the fluorescence spectrum has vibrational structure characteristic of the lower electronic state (Fig. 11G.3b).

Provided they can be seen, the $0-0$ absorption and fluorescence transitions (where the numbers are the values of $v_{\mathrm{f}}$ and $v_{\mathrm{i}}$, the vibrational quantum numbers for the final and initial states) can be expected to be coincident. The absorption spectrum arises from $0 \leftarrow 0,1 \leftarrow 0,2 \leftarrow 0, \ldots$ transitions which occur at progressively higher wavenumbers (shorter wavelengths) and with intensities governed by the FranckCondon principle. The fluorescence spectrum arises from $0 \rightarrow 0,0 \rightarrow 1, \ldots$ downward transitions which occur with decreasing wavenumbers (longer wavelengths). The $0-0 \mathrm{ab}-$ sorption and fluorescence peaks are not always exactly coincident, however, because the solvent may interact differently with the solute in the ground and excited electronic states (for instance, the hydrogen bonding pattern might differ). Because the solvent molecules do not have time to rearrange during the transition, the absorption occurs in an environment characteristic of the solvated ground state; however, the fluorescence occurs in an environment characteristic of the solvated excited state (Fig. 11G.4).\\
\includegraphics[max width=\textwidth, center]{2025_02_27_d4b95d9718f0b9e2e6b9g-503(1)}

Figure 11G. 4 The solvent can shift the fluorescence spectrum relative to the absorption spectrum. On the left absorption occurs with the solvent (depicted by the ellipses) in the arrangement characteristic of the ground electronic state of the molecule (the central blob). However, before fluorescence occurs, the solvent molecules relax into a new arrangement, and that arrangement is preserved during the subsequent radiative transition, the fluorescence.

Fluorescence occurs at lower frequencies than the incident radiation because the emissive transition occurs after some vibrational energy has been discarded into the surroundings. The vivid oranges and greens of fluorescent dyes are an everyday manifestation of this effect: they absorb in the ultraviolet and blue, and fluoresce in the visible. The mechanism also suggests that the intensity of the fluorescence ought to depend on the ability of the solvent molecules to accept the electronic and vibrational quanta. It is indeed found that a solvent composed of molecules with widely spaced vibrational levels (such as water) can in some cases accept the large quantum of electronic energy and so extinguish, or 'quench', the fluorescence. The rate at which fluorescence is quenched by other molecules also gives valuable kinetic information (Topic 17G).

Figure 11G. 5 shows the sequence of events leading to phosphorescence for a molecule with a singlet ground state (denoted $\mathrm{S}_{0}$ ). The first steps are the same as in fluorescence, but the presence of a triplet excited state $\left(\mathrm{T}_{1}\right)$ at an energy close to that of the singlet excited state $\left(S_{1}\right)$ plays a decisive role. The singlet and triplet excited states share a common geometry at the point where their potential energy curves intersect. Hence, if there is a mechanism for unpairing two electron spins (and achieving the conversion of $\uparrow \downarrow$ to $\uparrow \uparrow$ ), the molecule may undergo intersystem crossing, a non-radiative transition between states of different multiplicity, and become a triplet state. As in the discussion of atomic spectra (Topic 8C), singlet-triplet transitions may occur in the presence of spinorbit coupling. Intersystem crossing is expected to be important when a molecule contains a moderately heavy atom (such as sulfur), because then the spin-orbit coupling is large.

Once an excited molecule has crossed into a triplet state, it continues to discard energy into the surroundings. However, it is now stepping down the triplet's vibrational ladder and ends\\
\includegraphics[max width=\textwidth, center]{2025_02_27_d4b95d9718f0b9e2e6b9g-504(2)}

Figure 11G. 5 The sequence of steps leading to phosphorescence. The important step is the intersystem crossing (ISC), the switch from a singlet state $\left(S_{1}\right)$ to a triplet state $\left(T_{1}\right)$ brought about by spin-orbit coupling. The triplet state acts as a slowly radiating reservoir because the return to the ground state is spin-forbidden.\\
in its lowest vibrational level. The triplet state is lower in energy than the corresponding singlet state (Hund's rule, Topic 8 B ). The solvent cannot absorb the final, large quantum of electronic excitation energy, and the molecule cannot radiate its energy because return to the ground state is spin-forbidden. The radiative transition, however, is not totally forbidden because the spin-orbit coupling that was responsible for the intersystem crossing also weakens the selection rule. The molecules are therefore able to emit weakly, and the emission may continue long after the original excited state was formed.

This mechanism accounts for the observation that the excitation energy seems to get trapped in a slowly leaking reservoir. It also suggests (as is confirmed experimentally) that phosphorescence should be most intense from solid samples: energy transfer is then less efficient and intersystem crossing has time to occur as the singlet excited state steps slowly past the intersection point. The mechanism also suggests that the phosphorescence efficiency should depend on the presence of a moderately heavy atom (with strong spin-orbit coupling), which is in fact the case.\\
The various types of non-radiative and radiative transitions that can occur in molecules are often represented on a schematic Jablonski diagram of the type shown in Fig. 11G.6.

\subsection*{Brief illustration 11G. 1}
Fluorescence efficiency decreases, and the phosphorescence efficiency increases, in the series of compounds: naphthalene, 1-chloronaphthalene, 1-bromonaphthalene, 1-iodonaphthalene. The replacement of an H atom by successively heavier atoms enhances intersystem crossing from $S_{1}$ into $T_{1}$, thereby decreasing the efficiency of fluorescence. The rate of the radiative transition from $\mathrm{T}_{1}$ to $\mathrm{S}_{0}$ is also enhanced by the presence of heavier atoms, thereby increasing the efficiency of phosphorescence.\\
\includegraphics[max width=\textwidth, center]{2025_02_27_d4b95d9718f0b9e2e6b9g-504}

Figure 11G.6 A Jablonski diagram (here, for naphthalene) is a simplified portrayal of the relative positions of the electronic energy levels of a molecule. Vibrational levels of states of a given electronic state lie above each other, but the relative horizontal locations of the columns bear no relation to the nuclear separations in the states. The ground vibrational states of each electronic state are correctly located vertically but the other vibrational states are shown only schematically. (IC: internal conversion; ISC: intersystem crossing.)

\subsection*{11G.2 Dissociation and predissociation}
A chemically important fate for an electronically excited molecule is dissociation, the breaking of bonds (Fig. 11G.7). The onset of dissociation can be detected in an absorption spectrum by noting that the vibrational fine structure of a band terminates at a certain frequency. Absorption occurs in a continuous band above this dissociation limit because the final state is an unquantized translational motion of the fragments. Locating the dissociation limit is a valuable way of determining the bond dissociation energy.

In some cases, the vibrational structure disappears but resumes at higher frequencies of the incident radiation. This\\
\includegraphics[max width=\textwidth, center]{2025_02_27_d4b95d9718f0b9e2e6b9g-504(1)}

Figure 11G. 7 When absorption occurs to unbound states of the upper electronic state, the molecule dissociates and the absorption spectrum is a continuum. Below the dissociation limit the electronic spectrum shows a normal vibrational structure.\\
\includegraphics[max width=\textwidth, center]{2025_02_27_d4b95d9718f0b9e2e6b9g-505}

Figure 11G.8 When a dissociative state crosses a bound state, as in the upper part of the illustration, molecules excited to levels near the crossing may dissociate. This predissociation is detected in the spectrum as a loss of vibrational structure that resumes at higher frequencies.\\
effect provides evidence of predissociation, which can be interpreted in terms of the molecular potential energy curves shown in Fig. 11G.8. When a molecule is excited to a high vibrational level of the upper electronic state, its electrons may undergo a redistribution that results in it undergoing an internal conversion, a radiationless conversion to another electronic state of the same multiplicity. An internal conversion occurs most readily at the point of intersection of the two molecular potential energy curves, because there the nuclear geometries of the two electronic states are the same. The state into which the molecule converts may be dissociative, so the states near the intersection have a finite lifetime and hence their energies are imprecisely defined (as a result of lifetime broadening, Topic 11 A ). As a result, the absorption spectrum is blurred. When the incoming photon brings enough energy to excite the molecule to a vibrational level high above the intersection, the internal conversion does not occur (the nuclei are unlikely to have the same geometry). Consequently, the levels resume their well-defined, vibrational character with correspondingly well-defined energies, and the line structure resumes on the high-frequency side of the blurred region.

\subsection*{Brief illustration 11G. 2}
The $\mathrm{O}_{2}$ molecule absorbs ultraviolet radiation in a transition from its ${ }^{3} \Sigma_{\mathrm{g}}^{-}$ground electronic state to a ${ }^{3} \Sigma_{\mathrm{u}}^{-}$excited state that is energetically close to a dissociative ${ }^{3} \Pi_{u}$ state. In this case, the effect of predissociation is more subtle than the abrupt loss of vibrational-rotational structure in the spectrum; instead, the vibrational structure simply broadens rather than being lost completely. As before, the broadening is explained by the short lifetimes of the excited vibrational states near the intersection of the curves describing the bound and dissociative excited electronic states.

\subsection*{11G. 3 Lasers}
An excited state can be driven to discard its excess energy by using radiation to induce stimulated emission. The word laser is an acronym formed from light amplification by stimulated emission of radiation. In stimulated emission (Topic 11A), an excited state is stimulated to emit a photon by radiation of the same frequency: the more photons that are present, the greater is the probability of emission.

Laser radiation has a number of striking characteristics (Table 11G.1). Each of them (sometimes in combination with the others) opens up interesting opportunities in physical chemistry. Raman spectroscopy (Topics 11B-D) has flourished on account of the high intensity of monochromatic radiation available from lasers, and the ultrashort pulses that lasers can generate make possible the study of light-initiated reactions on timescales of femtoseconds and even attoseconds.

One requirement of laser action is the existence of a metastable excited state, an excited state with a long enough lifetime for it to undergo stimulated emission. For stimulated emission to dominate absorption, it is necessary for there to be a population inversion in which the population of the excited state is greater than that of the lower state. Figure 11G. 9 illustrates one way to achieve population inversion indirectly through an intermediate state I. Thus, the molecule or atom is excited to I, which then gives up some of its energy nonradiatively (for example, by passing energy on to vibrations of the surroundings) and changes into a lower state $B$. The laser

Table 11G. 1 Characteristics of laser radiation and their chemical applications

\begin{center}
\begin{tabular}{lll}
\hline
Characteristic & Advantage & Application \\
\hline
High power & Multiphoton process & \begin{tabular}{l}
Spectroscopy \\
Low detector noise \\
Improved sensitivity \\
\end{tabular} \\
 & \begin{tabular}{c}
High scattering \\
intensity \\
\end{tabular} & \begin{tabular}{c}
Raman spectroscopy \\
(Topics 11B-11D) \\
\end{tabular} \\
Monochromatic & \begin{tabular}{l}
High resolution \\
State selection \\
\end{tabular} & \begin{tabular}{l}
Spectroscopy \\
Photochemical studies \\
(Topic 17G) \\
\end{tabular} \\
 &  & \begin{tabular}{c}
Reaction dynamics \\
(Topic 18D) \\
\end{tabular} \\
 &  & \begin{tabular}{c}
Improved sensitivity \\
\end{tabular} \\
Collimated beam & \begin{tabular}{c}
Long path lengths \\
Forward-scattering \\
observable \\
\end{tabular} & \begin{tabular}{c}
Raman spectroscopy \\
(Topics 11B-11D) \\
\end{tabular} \\
 & \begin{tabular}{c}
Precise timing of \\
excitation \\
\end{tabular} & \begin{tabular}{c}
Fast reactions \\
(Topics 17G, 18C) \\
\end{tabular} \\
 &  & Relaxation (Topic 17C) \\
\end{tabular}
\end{center}

\begin{center}
\includegraphics[max width=\textwidth]{2025_02_27_d4b95d9718f0b9e2e6b9g-506(1)}
\end{center}

Figure 11G. 9 The transitions involved in a four-level laser. Because the laser transition terminates in an excited state (A), the population inversion between $A$ and $B$ is much easier to achieve than when the lower state of the laser transition is the ground state, X .\\
transition is the return of $B$ to a lower state $A$. Because four levels are involved overall, this arrangement leads to a four-level laser. One advantage of this arrangement is that the population inversion of the A and B levels is easier to achieve than one involving the heavily populated ground state. The transition from X to I is caused by irradiation with intense light (either continuously or as a flash) in the process called pumping. In some cases pumping is achieved with an electric discharge through xenon or with the radiation from another laser.

\subsection*{Brief illustration 11G. 3}
The neodymium laser is an example of a four-level solidstate laser (Fig. 11G.10). In one form it consists of $\mathrm{Nd}^{3+}$ ions at low concentration in yttrium aluminium garnet (YAG, specifically $\mathrm{Y}_{3} \mathrm{Al}_{5} \mathrm{O}_{12}$ ), and is then known as an Nd:YAG laser. A neodymium laser operates at a number of wavelengths in the infrared. The most common wavelength of operation is 1064 nm , which corresponds to the electronic transition from the ${ }^{4} \mathrm{~F}$ to the ${ }^{4} \mathrm{I}$ state of the $\mathrm{Nd}^{3+}$ ion.\\
\includegraphics[max width=\textwidth, center]{2025_02_27_d4b95d9718f0b9e2e6b9g-506}

Figure 11G. 10 The transitions involved in a neodymium laser.

Many of the most important laser systems are solid-state devices; they are discussed in Topic 15G.

\subsection*{Checklist of concepts}
\begin{enumerate}
  \item Fluorescence is radiative decay between states of the same multiplicity; it ceases soon after the exciting radiation is removed.2. Phosphorescence is radiative decay between states of different multiplicity; it persists after the exciting radiation is removed.3. Intersystem crossing is the non-radiative conversion to an electronic state of different multiplicity.4. A Jablonski diagram is a schematic diagram showing the types of non-radiative and radiative transitions that can occur in molecules.5. An additional fate of an electronically excited species is dissociation.6. Internal conversion is a non-radiative conversion to an electronic state of the same multiplicity.\\
$\square$ 7. Predissociation is the observation of the effects of dissociation before the dissociation limit is reached.\\
$\square$ 8. Laser action is the stimulated emission of coherent radiation between states related by a population inversion.9. A metastable excited state is an excited state with a long enough lifetime for it to undergo stimulated emission.10. A population inversion is a condition in which the population of an upper state is greater than that of a relevant lower state.11. Pumping, the stimulation of an absorption with an external source of intense radiation, is a process by which a population inversion is created.
\end{enumerate}

\chapter{FOCUS 12 Magnetic resonance}
The techniques of 'magnetic resonance' observe transitions between spin states of nuclei and electrons in molecules. 'Nuclear magnetic resonance' (NMR) spectroscopy observes nuclear spin transitions and is one of the most widely used spectroscopic techniques for the exploration of the structures and dynamics of molecules ranging from simple organic species to biopolymers. 'Electron paramagnetic resonance' (EPR) spectroscopy is a similar technique that probes electron spin transitions in species with unpaired electrons.

\subsection*{12A General principles}
This Topic gives an account of the principles that govern the energies and spectroscopic transitions between spin states of nuclei and electrons in molecules when a magnetic field is present. It describes simple experimental arrangements for the detection of these transitions.\\
12A. 1 Nuclear magnetic resonance; 12A. 2 Electron paramagnetic resonance

\subsection*{12B Features of NMR spectra}
This Topic contains a discussion of conventional NMR spectroscopy, showing how the properties of a magnetic nucleus are affected by its electronic environment and the presence of magnetic nuclei in its vicinity. These concepts explain how molecular structure governs the appearance of NMR spectra both in solution and in the solid state.\\
12B. 1 The chemical shift; 12B. 2 The origin of shielding constants; 12B.3 The fine structure; 12B. 4 Exchange processes; 12B. 5 Solid-state NMR

\subsection*{12C Pulse techniques in NMR}
The modern implementation of NMR spectroscopy employs pulses of radiofrequency radiation followed by analysis of the\\
resulting signal. This approach opens up many possibilities for the development of more sophisticated experiments. The Topic includes a discussion of spin relaxation in NMR and how it can be exploited, through the 'nuclear Overhauser effect', for structural studies.\\
12C. 1 The magnetization vector; 12C. 2 Spin relaxation; 12C. 3 Spin decoupling; 12C. 4 The nuclear Overhauser effect

\subsection*{12D Electron paramagnetic resonance}
The detailed form of an EPR spectrum reflects the molecular environment of the unpaired electron and the nuclei with which it interacts. From an analysis of the spectrum it is possible to infer how the electron spin density is distributed.\\
12D. 1 The $g$-value; 12D. 2 Hyperfine structure

\subsection*{Web resources What is an application of this material?}
Magnetic resonance is ubiquitous in chemistry, as it is an enormously powerful analytical and structural technique, especially in organic chemistry and biochemistry. One of the most striking applications of nuclear magnetic resonance is in medicine. 'Magnetic resonance imaging' (MRI) is a portrayal of the distribution of protons in a solid object (Impact 18), and this technique has proved to be particularly useful for diagnosing disease. Impact 19 highlights an application of electron paramagnetic resonance in materials science and biochemistry: the use of a 'spin probe', a radical that interacts with biopolymers or nanostructures, and has an EPR spectrum that is sensitive to the local structure and dynamics of its environment.

\section*{TOPIC 12A General principles}
\subsection*{Why do you need to know this material?}
Nuclear magnetic resonance spectroscopy and electron paramagnetic resonance spectroscopy are used widely in chemistry to identify molecules and to determine their structures. To understand these techniques, you need to understand how a magnetic field affects the energies of the spin states of nuclei and electrons.

\subsection*{What is the key idea?}
The spin states of nuclei and electrons have different energies in a magnetic field, and resonant transitions can take place between them when electromagnetic radiation of the appropriate frequency is applied.

\subsection*{What do you need to know already?}
You need to be familiar with the quantum mechanical concept of spin (Topic 8B), the Boltzmann distribution (see the Prologue to this text and Topic 13A), and the general features of spectroscopy (Topic 11A).

Electrons and many nuclei have the property called 'spin', an intrinsic angular momentum. This spin gives rise to a magnetic moment and results in them behaving like small bar magnets. The energies of these magnetic moments depend on their orientation with respect to an applied magnetic field.\\
Spectroscopic techniques that measure transitions between nuclear and electron spin energy levels rely on the phenomenon of resonance, the strong coupling of oscillators of the same frequency. In fact, all spectroscopy is a form of resonant coupling between the electromagnetic field and the molecules, but in magnetic resonance, at least in its original form, the energy levels are adjusted to match the electromagnetic field rather than vice versa.

\subsection*{12A. 1 Nuclear magnetic resonance}
The nuclear spin quantum number, $I$, is a fixed characteristic property of a nucleus ${ }^{1}$ and, depending on the nuclide, it is either an integer (including zero) or a half-integer (Table

\footnotetext{${ }^{1}$ Excited nuclear states, which are states in which the nucleons are arranged differently from the ground state, can have different spin from the ground state. Only ground states are considered here.
}12A.1). The angular momentum associated with nuclear spin has the same properties as other kinds of angular momentum (Topic 7F):

\begin{itemize}
  \item The magnitude of the angular momentum is $\{I(I+$ 1) ${ }^{1 / 2} \hbar$.
  \item The component of the angular momentum on a specified axis ('the $z$-axis') is $m_{I} \hbar$ where $m_{I}=I, I-1, \ldots,-I$.
  \item The orientation of the angular momentum, and hence of the magnetic moment, is determined by the value of $m_{r}$.
\end{itemize}

According to the second property, the angular momentum, and hence the magnetic moment, of the nucleus may lie in $2 I+1$ different orientations relative to an axis. A ${ }^{1} \mathrm{H}$ nucleus has $I=\frac{1}{2}$ so its magnetic moment may adopt either of two orientations ( $m_{I}=+\frac{1}{2},-\frac{1}{2}$ ). The $m_{I}=+\frac{1}{2}$ state is commonly denoted $\alpha$ and the $m_{I}=-\frac{1}{2}$ state is commonly denoted $\beta$. $\mathrm{A}^{14} \mathrm{~N}$ nucleus has $I=1$ so there are three orientations ( $m_{I}=+1,0,-1$ ). Examples of nuclei with $I=0$, and hence no magnetic moment, are ${ }^{12} \mathrm{C}$ and ${ }^{16} \mathrm{O}$.

\subsection*{(a) The energies of nuclei in magnetic fields}
The energy of a magnetic moment $\boldsymbol{\mu}$ in a magnetic field $\mathscr{B}$ is equal to their scalar product (see The chemist's toolkit 22 in Topic 8C):


\begin{equation*}
E=-\mu \cdot \mathcal{B} \tag{12A.1}
\end{equation*}


More formally, $\mathscr{B}$ is the 'magnetic induction' and is measured in tesla, $\mathrm{T} ; 1 \mathrm{~T}=1 \mathrm{~kg} \mathrm{~s}^{-2} \mathrm{~A}^{-1}$. The (non-SI) unit gauss, G , is also occasionally used: $1 \mathrm{~T}=10^{4} \mathrm{G}$. The corresponding expression for the hamiltonian is


\begin{equation*}
\hat{H}=-\hat{\mu} \cdot \mathscr{B} \tag{12A.2}
\end{equation*}


Table 12A. 1 Nuclear constitution and the nuclear spin quantum number*

\begin{center}
\begin{tabular}{lll}
\hline
Number of protons & Number of neutrons & $\boldsymbol{I}$ \\
\hline
Even & even & 0 \\
Odd & odd & integer $(1,2,3, \ldots)$ \\
Even & odd & half-integer $\left(\frac{1}{2}, \frac{3}{2}, \frac{5}{2}, \ldots\right)$ \\
Odd & even & half-integer $\left(\frac{1}{2}, \frac{3}{2}, \frac{5}{2}, \ldots\right)$ \\
\hline
\end{tabular}
\end{center}

\footnotetext{\begin{itemize}
  \item For nuclear ground states.
\end{itemize}
}The magnetic moment operator of a nucleus is proportional to its spin angular momentum operator and is written


\begin{equation*}
\hat{\mu}=\gamma_{\mathrm{N}} \hat{I} \tag{12A.3a}
\end{equation*}


The constant of proportionality, $\gamma_{\mathrm{N}}$, is the nuclear magnetogyric ratio (also called the 'gyromagnetic ratio'); its value depends on the identity of the nucleus and is determined empirically (Table 12A.2). If the magnetic field defines the $z$-direction and has magnitude $\mathscr{B}_{0}$, then eqn 12A. 2 becomes


\begin{equation*}
\hat{H}=-\hat{\mu}_{z} \mathcal{B}_{0}=-\gamma_{\mathrm{N}} \mathcal{B}_{0} \hat{I}_{z} \tag{12A.3b}
\end{equation*}


The eigenvalues of the operator $\hat{I}_{z}$ for the $z$-component of the spin angular momentum are $m_{I} \hbar$. The eigenvalues of the hamiltonian in eqn 12A.3b, the allowed energy levels of the nucleus in a magnetic field, are therefore

$$
E_{m_{I}}=-\gamma_{\mathrm{N}} \hbar \mathscr{B}_{0} m_{I} \quad \begin{array}{ll}
\text { Energies of a nuclear spin } \\
\text { in a magnetic field }
\end{array}
$$

(12A.4a)\\
It is common to rewrite this expression in terms of the nuclear magneton, $\mu_{\mathrm{N}}$,

$$
\mu_{\mathrm{N}}=\frac{e \hbar}{2 m_{\mathrm{p}}} \quad \begin{aligned}
& \text { Nuclear magneton } \\
& \text { [definition] }
\end{aligned}
$$

(12A.4b)\\
(where $m_{\mathrm{p}}$ is the mass of the proton) and an experimentally determined dimensionless constant called the nuclear $g$-factor, $g_{I}$,

$$
g_{I}=\frac{\gamma_{\mathrm{N}} \hbar}{\mu_{\mathrm{N}}}
$$

Nuclear $g$-factor [definition]\\
(12A.4c)\\
Equation 12A.4a then becomes

$$
E_{m_{I}}=-g_{I} \mu_{\mathrm{N}} \mathcal{B}_{0} m_{I} \quad \begin{aligned}
& \text { Energies of a nuclear spin } \\
& \text { in a magnetic field }
\end{aligned}
$$

(12A.4d)\\
The value of the nuclear magneton is $\mu_{\mathrm{N}}=5.051 \times 10^{-27} \mathrm{~J} \mathrm{~T}^{-1}$. Typical values of nuclear $g$-factors range between -6 and +6 (Table 12A.2). Positive values of $g_{I}$ and $\gamma_{\mathrm{N}}$ denote a magnetic moment that lies in the same direction as the spin angular momentum; negative values indicate that the magnetic moment and spin lie in opposite directions.

When $\gamma_{\mathrm{N}}>0$, as is the case for the most commonly observed nuclei ${ }^{1} \mathrm{H}$ and ${ }^{13} \mathrm{C}$, in a magnetic field the energies of states with $m_{I}>0$ lie below states with $m_{I}<0$. For a spin- $\frac{1}{2}$ nucleus, a\\
\includegraphics[max width=\textwidth, center]{2025_02_27_d4b95d9718f0b9e2e6b9g-521}

Figure 12A. 1 The nuclear spin energy levels of a spin- $\frac{1}{2}$ nucleus with positive magnetogyric ratio (e.g. ${ }^{1} \mathrm{H}$ or ${ }^{13} \mathrm{C}$ ) in a magnetic field. Resonant absorption of radiation occurs when the energy separation of the levels matches the energy of the photons.\\
nucleus for which $I=\frac{1}{2}$, the $\alpha$ state lies lower in energy than the $\beta$ state, and the separation between them is


\begin{equation*}
\Delta E=E_{-1 / 2}-E_{+1 / 2}=\frac{1}{2} \gamma_{\mathrm{N}} \hbar \mathcal{B}_{0}-\left(-\frac{1}{2} \gamma_{\mathrm{N}} \hbar \mathcal{B}_{0}\right)=\gamma_{\mathrm{N}} \hbar \mathcal{B}_{0} \tag{12A.5}
\end{equation*}


The corresponding frequency of electromagnetic radiation for a transition between these states is given by the Bohr frequency condition $\Delta E=h v$ (Fig. 12A.1). Therefore,


\begin{equation*}
h v=\gamma_{\mathrm{N}} \hbar \mathcal{B}_{0} \quad \text { or } \quad v=\frac{\gamma_{\mathrm{N}} \mathscr{B}_{0}}{2 \pi} \text { Resonance condition in NMR } \tag{12A.6}
\end{equation*}


This relation is called the resonance condition, and $v$ is called the NMR frequency for that nucleus. Although eqn 12A. 6 has been derived for a spin $-\frac{1}{2}$ nucleus, the same expression applies for any nucleus with non-zero spin.

It is sometimes useful to compare the quantum mechanical treatment with the classical picture in which magnetic nuclei are pictured as tiny bar magnets. A bar magnet in a magnetic field undergoes the motion called precession as it twists round the direction of the field and sweeps out the surface of a cone (Fig. 12A.2). The rate of precession $v_{\mathrm{L}}$ is called the Larmor precession frequency:

$$
v_{\mathrm{L}}=\frac{\gamma_{\mathrm{N}} \mathcal{B}_{0}}{2 \pi} \quad \begin{aligned}
& \text { Larmor frequency of a nucleus } \\
& \text { [definition] }
\end{aligned}
$$

(12A.7)

Table 12A. 2 Nuclear spin properties*

\begin{center}
\begin{tabular}{lccccc}
\hline
Nucleus & Natural abundance $/ \%$ & Spin, $I$ & $g$-factor, $g_{I}$ & Magnetogyric ratio, $\gamma_{\mathrm{N}} /\left(10^{7} \mathrm{~T}^{-1} \mathrm{~s}^{-1}\right)$ & NMR frequency at $1 \mathrm{~T}, v / \mathrm{MHz}$ \\
\hline
${ }^{1} \mathrm{H}$ & 99.98 & $\frac{1}{2}$ & 5.586 & 26.75 & 42.576 \\
${ }^{2} \mathrm{H}$ & 0.02 & 1 & 0.857 & 4.11 & 6.536 \\
${ }^{13} \mathrm{C}$ & 1.11 & $\frac{1}{2}$ & 1.405 & 6.73 & 10.708 \\
${ }^{11} \mathrm{~B}$ & 80.4 & $\frac{3}{2}$ & 1.792 & 8.58 & 13.663 \\
${ }^{14} \mathrm{~N}$ & 99.64 & 1 & 0.404 & 1.93 & 3.078 \\
\hline
\end{tabular}
\end{center}

\footnotetext{\begin{itemize}
  \item More values are given in the Resource section.
\end{itemize}
}
\includegraphics[max width=\textwidth, center]{2025_02_27_d4b95d9718f0b9e2e6b9g-522(1)}

Figure 12A. 2 The classical view of magnetic nuclei pictures them as behaving as tiny bar magnets. In an externally applied magnetic field the resulting magnetic moment, here represented as a vector, precesses round the direction of the field.

The Larmor precession frequency is the same as the resonance frequency given by eqn 12A.6. In other words, the frequency of radiation that causes resonant transitions between the $\alpha$ and $\beta$ states is the same as the Larmor precession frequency. The achievement of resonance absorption can therefore be pictured as changing the applied magnetic field until the bar magnet representing the nuclear magnetic moment precesses at the same frequency as the magnetic component of the electromagnetic field to which it is exposed.

\subsection*{Brief illustration 12A. 1}
The NMR frequency for ${ }^{1} \mathrm{H}$ nuclei $\left(I=\frac{1}{2}\right)$ in a 12.0 T magnetic field can be found using eqn 12A.6, with the relevant value of $\gamma_{\mathrm{N}}$ taken from Table 12A.1:

$$
\nu=\frac{\overbrace{\left(2.6752 \times 10^{8} \mathrm{~T}^{-1} \mathrm{~s}^{-1}\right)}^{\gamma_{N}} \times \overbrace{(12.0 \mathrm{~T})}^{\mathcal{B}_{0}}}{2 \pi}=5.11 \times 10^{8} \mathrm{~s}^{-1}=511 \mathrm{MHz}
$$

This radiation lies in the radiofrequency region of the electromagnetic spectrum, close to frequencies used for radio communication.

\subsection*{(b) The NMR spectrometer}
The key component of an NMR spectrometer (Fig. 12A.3) is the magnet into which the sample is placed. Most modern spectrometers use superconducting magnets capable of producing fields of 12 T or more. Such magnets have the advantages that the field they produce is stable over time and no electrical power is needed to maintain the field. With the currently available magnets, all NMR frequencies fall in the radiofrequency range (see the previous Brief illustration). Therefore, a radiofrequency transmitter and receiver are needed to excite and detect the transitions taking place between nuclear spin states. The details of how the transitions are excited and detected are discussed in Topic 12C.

The sample being studied is most commonly in the form of a solution contained in a glass tube placed within the magnet. It is also possible to study solid samples by using more specialized techniques. Although the superconducting magnet itself has to be held close to the temperature of liquid helium (4K),\\
\includegraphics[max width=\textwidth, center]{2025_02_27_d4b95d9718f0b9e2e6b9g-522}

Figure 12A. 3 The layout of a typical NMR spectrometer. The sample is held within the probe, which is placed at the centre of the magnetic field.\\
the magnet is designed so as to have a room-temperature clear space into which the sample can be placed.

The intensity of an NMR transition depends on a number of factors which can be identified by considering the populations of the two spin states.

\subsection*{How is that done? 12A. 1 Identifying the contributions to the absorption intensity}
The rate of absorption of electromagnetic radiation is proportional to the population of the lower energy state ( $N_{\alpha}$ in the case of a spin $-\frac{1}{2}$ nucleus) and the rate of stimulated emission is proportional to the population of the upper state $\left(N_{\beta}\right)$. At the low frequencies typical of magnetic resonance, spontaneous emission can be neglected as it is very slow. Therefore, the net rate of absorption is proportional to the difference in populations:

$$
\text { Rate of absorption } \propto N_{\alpha}-N_{\beta}
$$

Step 1 Write an expression for the intensity of absorption in terms of the population difference

The intensity of absorption, the rate at which energy is absorbed, is proportional to the product of the rate of absorption (the rate at which photons are absorbed) and the energy of each photon. The latter is proportional to the frequency $v$ of the incident radiation (through $E=h v$ ). At resonance, this frequency is proportional to the applied magnetic field (through $\left.v=\gamma_{\mathrm{N}} \mathcal{B}_{0} / 2 \pi\right)$, so it follows that


\begin{align*}
& \text { Intensity of absorption } \propto \text { rate of absorption } \\
& \times \text { energy of photon }  \tag{12A.8a}\\
& \propto\left(N_{\alpha}-N_{\beta}\right) \times h \gamma_{\mathrm{N}} \mathcal{B}_{0} / 2 \pi
\end{align*}


\subsection*{Step 2 Write an expression for the ratio of populations}
Now use the Boltzmann distribution (see the Prologue to this text and Topic 13A) to write an expression for the ratio of populations:

$$
\frac{N_{\beta}}{N_{\alpha}}=\mathrm{e}^{-\overbrace{\gamma_{\mathrm{N}} \hbar \mathcal{B}_{0} / k T}^{\Delta E}} \stackrel{\underbrace{\mathrm{e}^{-x}=1-x+\cdots}}{\approx 1-\frac{\gamma_{\mathrm{N}} \hbar \mathcal{B}_{0}}{k T}}
$$

The expansion of the exponential term (see The chemist's toolkit 12 in Topic 5 B ) is appropriate for $\Delta E=\gamma_{\mathrm{N}} \hbar \mathcal{B}_{0} \ll k T$, a condition usually met for nuclear spins.

Step 3 Use that ratio to write an expression for the population difference\\
Consider the ratio of the population difference to the total number of spins, $N:\left(N_{\alpha}-N_{\beta}\right) / N$ and use the expression for the ratio of populations to write this ratio as

$$
\begin{aligned}
\frac{N_{\alpha}-N_{\beta}}{N_{\alpha}-N_{\beta}} & =\frac{N_{\alpha}\left(1-N_{\beta} / N_{\alpha}\right)}{N_{\alpha}\left(1+N_{\beta} / N_{\alpha}\right)}=\frac{1-\overbrace{N} / N_{\alpha}}{1+\underbrace{N_{\beta} / N_{\alpha}}_{1-\gamma_{N} \hbar \mathcal{B}_{0} / k T}} \\
& \approx \frac{1-\gamma_{N} \hbar \mathcal{B}_{0} / k T}{1+\underbrace{\left(1-\gamma_{N} \hbar \mathcal{B}_{0} / k T\right)}_{\approx 1}}=\frac{\gamma_{\mathrm{N}} \hbar \mathcal{B}_{0} / k T}{2}
\end{aligned}
$$

Therefore

$$
N_{\alpha}-N_{\beta} \approx \frac{N \gamma_{\mathrm{N}} \hbar \mathscr{B}_{0}}{2 k T} \quad \begin{aligned}
& \text { Population difference } \\
& \text { [spin-2 }-\frac{1}{2} \text { nuclei] }
\end{aligned}
$$

(12A.8b)

The intensity is now obtained by substituting this expression into eqn 12A.8a which gives, after discarding constants that do not refer to the spins,\\
$\Longrightarrow$ Intensity $\propto \frac{N \gamma_{\mathrm{N}}^{2} \mathcal{B}_{0}^{2}}{T}$\\
$\frac{\text { (12A. } 8 \mathrm{c})}{\text { intensity }}$

Because the intensity is proportional to $\mathscr{B}_{0}^{2}$ it follows that the signal can be enhanced significantly by increasing the strength of the applied magnetic field. The use of high magnetic fields also simplifies the appearance of spectra (a point explained in Topic 12B) and so allows them to be interpreted more readily. The intensity is also proportional to $\gamma_{\mathrm{N}}^{2}$, so, all other things being equal, nuclei with large magnetogyric ratios ( ${ }^{1} \mathrm{H}$, for instance) give more intense signals than those with small magnetogyric ratios ( ${ }^{13} \mathrm{C}$, for instance).

\subsection*{Brief illustration 12A. 2}
For ${ }^{1} \mathrm{H}$ nuclei $\gamma_{\mathrm{N}}=2.675 \times 10^{8} \mathrm{~T}^{-1} \mathrm{~s}^{-1}$. Therefore, for 1000000 protons in a field of 10 T at $20^{\circ} \mathrm{C}$,

$$
\begin{aligned}
N_{\alpha}-N_{\beta} & \approx \frac{\overbrace{1000000}^{N} \times \overbrace{\left(2.675 \times 10^{8} \mathrm{~T}^{-1} \mathrm{~s}^{-1}\right)}^{\gamma_{N}} \times \overbrace{\left(1.055 \times 10^{-34} \mathrm{Js}\right)}^{\hbar} \times \overbrace{(10 \mathrm{~T})}^{\mathcal{B}_{0}}}{2 \times \underbrace{\left(1.381 \times 10^{-23} \mathrm{~J} \mathrm{~K}^{-1}\right)}_{k} \times \underbrace{(293 \mathrm{~K})}_{T}} \\
& \approx 35
\end{aligned}
$$

Even in such a strong field there is only a tiny imbalance of population of about 35 in a million. This small population difference means that special techniques had to be developed before NMR became a viable technique.

\subsection*{12A.2 Electron paramagnetic resonance}
The observation of resonant transitions between the energy levels of an electron in a magnetic field is the basis of electron paramagnetic resonance (EPR; or electron spin resonance, ESR). This kind of spectroscopy has several features in common with NMR.

\subsection*{(a) The energies of electrons in magnetic fields}
The magnetic moment of an electron is proportional to its spin angular momentum. Its magnetic moment operator and the hamiltonian for its interaction with a magnetic field are


\begin{equation*}
\hat{\boldsymbol{\mu}}=\gamma_{\mathrm{e}} \hat{\boldsymbol{s}} \quad \text { and } \quad \hat{H}=-\gamma_{\mathrm{e}} \mathcal{B} \cdot \hat{\boldsymbol{s}} \tag{12A.9a}
\end{equation*}


where $\hat{\boldsymbol{s}}$ is the spin angular momentum operator and $\gamma_{\mathrm{e}}$ is the magnetogyric ratio of the electron:

$$
\gamma_{\mathrm{e}}=-\frac{g_{\mathrm{e}} e}{2 m_{\mathrm{e}}}
$$

Magnetogyric ratio of electron\\
(12A.9b)\\
with $g_{\mathrm{e}}=2.002319 \ldots$ as the $\boldsymbol{g}$-value of the free electron. (Note that the current convention is to include the $g$-value in the definition of the magnetogyric ratio.) Dirac's relativistic theory, his modification of the Schrödinger equation to make it consistent with Einstein's special relativity, gives $g_{e}=2$; the additional $0.002319 \ldots$ arises from interactions of the electron with the electromagnetic fluctuations of the vacuum that surrounds it. The negative sign of $\gamma_{e}$ (arising from the sign of the charge on the electron) shows that the magnetic moment is opposite in direction to the vector representing its angular momentum.

The hamiltonian when the magnetic field lies in the $z$-direction and has magnitude $\mathcal{B}_{0}$ is


\begin{equation*}
\hat{H}=-\gamma_{\mathrm{e}} \mathcal{B}_{0} \hat{s}_{z} \tag{12A.10}
\end{equation*}


where $\hat{s}_{z}$ is the operator for the $z$-component of the spin angular momentum. It follows that the energies of an electron spin in a magnetic field are

\[
E_{m_{s}}=-\gamma_{\mathrm{e}} \hbar \mathcal{B}_{0} m_{s} \quad \begin{align*}
& \text { Energies of an electron spin in }  \tag{12A.11a}\\
& \text { a magnetic field }
\end{align*}
\]

with $m_{s}= \pm \frac{1}{2}$. It is common to write this expression in terms of the Bohr magneton, $\mu_{B}$, defined as

$$
\mu_{\mathrm{B}}=\frac{e \hbar}{2 m_{\mathrm{e}}}
$$

Bohr magneton (12A.11b)\\
\includegraphics[max width=\textwidth, center]{2025_02_27_d4b95d9718f0b9e2e6b9g-524}

Figure 12A. 4 Electron spin levels in a magnetic field. Note that the $\beta$ state is lower in energy than the $\alpha$ state (because the magnetogyric ratio of an electron is negative). Resonant absorption occurs when the frequency of the incident radiation matches the frequency corresponding to the energy separation.

Its value is $9.274 \times 10^{-24} \mathrm{JT}^{-1}$. This positive quantity is often regarded as the fundamental quantum of magnetic moment. Note that the Bohr magneton is about 2000 times bigger than the nuclear magneton, so electron magnetic moments are that much bigger than nuclear magnetic moments. By using eqn 12A.9b the term $\gamma_{e} \hbar$ in eqn 12A.11a can be expressed as $-g_{e} e \hbar / 2 m_{e}$, which in turn can be written $-g_{e} \mu_{\mathrm{B}}$ by introducing the definition of the Bohr magneton from eqn 12A.11b. It then follows that

\[
E_{m_{\mathrm{s}}}=g_{\mathrm{e}} \mu_{\mathrm{B}} \mathscr{B}_{0} m_{\mathrm{s}} \quad \begin{align*}
& \text { Energies of an electron spin in }  \tag{12A.11c}\\
& \text { a magnetic field }
\end{align*}
\]

and the energy separation between the $m_{\mathrm{s}}=+\frac{1}{2}(\alpha)$ and $m_{\mathrm{s}}=-\frac{1}{2}$ $(\beta)$ states is


\begin{equation*}
\Delta E=E_{+1 / 2}-E_{-1 / 2}=\frac{1}{2} g_{\mathrm{e}} \mu_{\mathrm{B}} \mathcal{B}_{0}-\left(-\frac{1}{2} g_{\mathrm{e}} \mu_{\mathrm{B}} \mathcal{B}_{0}\right)=g_{\mathrm{e}} \mu_{\mathrm{B}} \mathcal{B}_{0} \tag{12A.12a}
\end{equation*}


with $\beta$ the lower state. This energy separation comes into resonance with electromagnetic radiation of frequency $v$ when (Fig. 12A.4)

$$
h v=g_{\mathrm{e}} \mu_{\mathrm{B}} \mathscr{B}_{0} \quad \text { Resonance condition for EPR }
$$

(12A.12b)

\subsection*{Brief illustration 12A. 3}
A typical commercial EPR spectrometer uses a magnetic field of about 0.33 T . The EPR resonance frequency is

$$
\begin{aligned}
& v=\frac{\overbrace{(2.0023)}^{g_{e}} \times \overbrace{\left(9.274 \times 10^{-24} \mathrm{~J} \mathrm{~T}^{-1}\right)}^{\mu_{\mathrm{B}}} \times \overbrace{(0.33 \mathrm{~T})}^{\mathcal{B}_{0}}}{\underbrace{}_{h}} \\
& =9.2 \times 10^{9} \mathrm{~s}^{-1}=9.2 \mathrm{GHz}
\end{aligned}
$$

This frequency corresponds to a wavelength of 3.2 cm , which is in the microwave region, and specifically in the ' X band' of frequencies.

\subsection*{(b) The EPR spectrometer}
Most commercial EPR spectrometers use magnetic field strengths that result in EPR frequencies in the microwave region (see the preceding Brief illustration). The layout of a typical EPR spectrometer is shown in Fig. 12A.5. It consists of a fixed-frequency microwave source (typically a Gunn oscillator, based on a solid-state device), a cavity into which the sample (held in a glass or quartz tube) is inserted, a microwave detector, and an electromagnet with a variable magnetic field. The sample in an EPR observation must have unpaired electrons, so is either a radical or a d-metal complex. For technical reasons related to the detection procedure, the spectrum shows the first derivative of the absorption line (Fig. 12A.6).\\
\includegraphics[max width=\textwidth, center]{2025_02_27_d4b95d9718f0b9e2e6b9g-524(1)}

Figure 12A. 5 The layout of a typical EPR spectrometer. A typical magnetic field is 0.3 T , which requires $9 \mathrm{GHz}(3 \mathrm{~cm})$ microwaves for resonance.\\
\includegraphics[max width=\textwidth, center]{2025_02_27_d4b95d9718f0b9e2e6b9g-524(2)}

Figure 12A. 6 When phase-sensitive detection is used, the signal is the first derivative of the absorption intensity. Note that the peak of the absorption corresponds to the point where the derivative passes through zero.

As in NMR, the intensities of spectral lines in EPR depend on the difference in populations between the ground and excited states. For an electron, the $\beta$ state lies below the $\alpha$ state in energy and, by a similar argument that led to eqn 12A. 8 b for nuclei,

\[
N_{\beta}-N_{\alpha} \approx \frac{N g_{\mathrm{e}} \mu_{\mathrm{B}} \mathcal{B}_{0}}{2 k T} \quad \begin{align*}
& \text { Population difference }  \tag{12A.13}\\
& \text { [electrons] }
\end{align*}
\]

where $N$ is the total number of electron spins.

\subsection*{Brief illustration 12A. 4}
When 1000 electron spins experience a magnetic field of 1.0 T at $20^{\circ} \mathrm{C}(293 \mathrm{~K})$, the population difference is

$$
\begin{aligned}
N_{\beta}-N_{\alpha} & \approx \frac{\overbrace{1000 \times}^{N}}{\frac{\overbrace{2.0023}^{g_{\mathrm{e}}} \times \overbrace{\left(9.274 \times 10^{-24} \mathrm{~J} \mathrm{~T}^{-1}\right)}^{\mu_{\mathrm{B}}} \times \overbrace{(1.0 \mathrm{~T})}^{\mathcal{B}_{0}}}{2 \times \underbrace{\left(1.381 \times 10^{-23} \mathrm{JK}^{-1}\right)}_{k} \times \underbrace{(293 \mathrm{~K})}_{T}}} \\
& \approx 2.3
\end{aligned}
$$

There is an imbalance of populations of only about two electrons in a thousand. However, the imbalance is much larger for electron spins than for nuclear spins (Brief illustration 12A.2) because the energy separation between the spin states of electrons is larger than that for nuclear spins even at the lower magnetic field strengths normally employed for EPR.

\subsection*{Checklist of concepts}
1. The nuclear spin quantum number, $I$, of a nucleus is either a non-negative integer or half-integer; $I$ can be zero.2. In the presence of a magnetic field a nucleus has $2 I+1$ energy levels characterized by different values of $m_{I}$.3. Nuclear magnetic resonance (NMR) is the observation of the absorption of radiofrequency electromagnetic radiation by nuclei in a magnetic field.4. In NMR the absorption intensity increases with the strength of the applied magnetic field (as $\mathscr{B}_{0}^{2}$ ) and is also proportional to the square of the magnetogyric ratio of the nucleus.5. In the presence of a magnetic field, an electron has two energy levels corresponding to the $\alpha$ and $\beta$ spin states.6. Electron paramagnetic resonance (EPR) is the observation of the resonant absorption of microwave electromagnetic radiation by unpaired electrons in a magnetic field.\subsection*{Checklist of equations}
\begin{center}
\begin{tabular}{|c|c|c|c|}
\hline
Property & Equation & Comment & Equation number \\
\hline
\multirow[t]{2}{*}{Energies of a nuclear spin in a magnetic field} & $E_{m_{I}}=-\gamma_{\mathrm{N}} \hbar \mathcal{B}_{0} m_{I}$ &  & 12A.4a \\
\hline
 & $=-g_{I} \mu_{\mathrm{N}} \mathcal{B}_{0} m_{I}$ &  & 12A.4d \\
\hline
Nuclear magneton & $\mu_{\mathrm{N}}=e \hbar / 2 m_{\mathrm{p}}$ & $\mu_{\mathrm{N}}=5.051 \times 10^{-27} \mathrm{JT}^{-1}$ & 12A.4b \\
\hline
Resonance condition (spin- $\frac{1}{2}$ nuclei) & $h \nu=\gamma_{N} \hbar \mathcal{B}_{0}$ &  & 12A. 6 \\
\hline
Larmor frequency & $v_{\mathrm{L}}=\gamma_{\mathrm{N}} \mathcal{B}_{0} / 2 \pi$ &  & 12A. 7 \\
\hline
Magnetogyric ratio (electron) & $\gamma_{\mathrm{e}}=-g_{\mathrm{e}} e / 2 m_{\mathrm{e}}$ & $g_{\text {e }}=2.002319 \ldots$ & 12A.9b \\
\hline
\multirow[t]{2}{*}{Energies of an electron spin in a magnetic field} & $E_{m_{s}}=-\gamma_{\mathrm{e}} \hbar \mathcal{B}_{0} m_{s}$ &  & 12A.11a \\
\hline
 & $=g_{e} \mu_{\mathrm{B}} \mathcal{B}_{0} m_{s}$ &  & 12A.11c \\
\hline
Bohr magneton & $\mu_{\mathrm{B}}=e \hbar / 2 m_{\text {e }}$ & $\mu_{\mathrm{B}}=9.274 \times 10^{-24} \mathrm{JT}^{-1}$ & 12A.11b \\
\hline
Resonance condition (electrons) & $h \nu=g_{\mathrm{e}} \mu_{\mathrm{B}} \mathcal{B}_{0}$ &  & 12A.12b \\
\hline
\end{tabular}
\end{center}

\section*{TOPIC 12B Features of NMR spectra}
\subsection*{Why do you need to know this material?}
To analyse NMR spectra and extract the wealth of information they contain you need to understand how the appearance of a spectrum correlates with molecular structure.

\subsection*{What is the key idea?}
The resonance frequency of a magnetic nucleus is affected by its electronic environment and the presence of magnetic nuclei in the vicinity.

\subsection*{What do you need to know already?}
You need to be familiar with the general principles of magnetic resonance (Topic 12A).

Nuclear magnetic moments interact with the local magnetic field, the field at the location of the nucleus in question. The local field may differ from the applied field due to the effects of the electrons surrounding the nucleus and the presence of other magnetic nuclei in the molecule. The overall effect is that the NMR frequency of a given nucleus is sensitive to its molecular environment.

\subsection*{12B. 1 The chemical shift}
The applied magnetic field can be thought of as causing a circulation of electrons through the molecule. This circulation is analogous to an electric current and so gives rise to a magnetic field. The local magnetic field, $\mathscr{B}_{\text {loc }}$, the total field experienced by the nucleus, is the sum of the applied field $\mathscr{B}_{0}$ and the additional field $\delta \mathscr{B}$ due to the circulation of the electrons


\begin{equation*}
\mathscr{B}_{\mathrm{loc}}=\mathscr{B}_{0}+\delta \mathscr{B} \tag{12B.1}
\end{equation*}


The additional field is proportional to the applied field, and it is conventional to write

$$
\delta \mathscr{B}=-\sigma \mathcal{B}_{0}
$$

Shielding constant [definition]\\
(12B.2)\\
where the dimensionless quantity $\sigma$ is called the shielding constant of the nucleus. The ability of the applied field to in-\\
duce an electronic current in the molecule, and hence affect the strength of the resulting local magnetic field, depends on the details of the electronic structure near the magnetic nucleus of interest, so nuclei in different chemical groups have different shielding constants. As a result, the Larmor frequency $v_{\mathrm{L}}$ of the nucleus (and therefore its resonance frequency) changes from $\gamma_{\mathrm{N}} \mathcal{B}_{0} / 2 \pi$ to


\begin{equation*}
v_{\mathrm{L}}=\frac{\gamma_{\mathrm{N}} \mathcal{B}_{\mathrm{loc}}}{2 \pi}=\frac{\gamma_{\mathrm{N}}\left(\mathscr{B}_{0}+\delta \mathscr{B}\right)}{2 \pi}=\frac{\gamma_{\mathrm{N}} \mathscr{B}_{0}}{2 \pi}(1-\sigma) \tag{12B.3}
\end{equation*}


The Larmor frequency is different for nuclei in different environments, even if those nuclei are of the same element.

The chemical shift of a nucleus is the difference between its resonance frequency and that of a reference standard. The standard for ${ }^{1} \mathrm{H}$ and ${ }^{13} \mathrm{C}$ is the resonance in tetramethylsilane, $\mathrm{Si}\left(\mathrm{CH}_{3}\right)_{4}$, commonly referred to as TMS. The frequency separation between the resonance from a particular nucleus and that from the standard increases with the strength of the applied magnetic field because the Larmor frequency (eqn 12B.3) is proportional to the applied field.

Chemical shifts are reported on the $\delta$ scale, which is defined as

$$
\delta=\frac{v-v^{\circ}}{v^{\circ}} \times 10^{6} \quad \begin{aligned}
& \delta \text { scale } \\
& \text { [definition] }
\end{aligned}
$$

(12B.4a)\\
where $v^{\circ}$ is the resonance frequency (Larmor frequency) of the standard. Because $v^{\circ}$ is very close to the operating frequency of the spectrometer $v_{\text {spect }}$, which is typically chosen to be in the middle of the range of Larmor frequencies exhibited by the nucleus being studied, the $v^{\circ}$ in the denominator of eqn 12B.4a can safely be replaced by $v_{\text {spect }}$ to give


\begin{equation*}
\delta=\frac{v-v^{\circ}}{v_{\text {spect }}} \times 10^{6} \tag{12B.4b}
\end{equation*}


The advantage of the $\delta$ scale is that shifts reported on it are independent of the applied field (because both numerator and denominator are proportional to the applied field). The resonance frequencies themselves, however, do depend on the applied field through


\begin{equation*}
v=v^{\circ}+\left(\frac{v_{\text {spect }}}{10^{6}}\right) \delta \tag{12B.5}
\end{equation*}


\subsection*{Brief illustration 12B. 1}
In an NMR spectrometer operating at 500.13000 MHz the resonance from TMS is found to be at a frequency of 500.12750 MHz . The chemical shift of a resonance at a frequency of 500.12825 MHz is

$$
\begin{aligned}
\delta=\frac{v-v^{\circ}}{v_{\text {spect }}} \times 10^{6} & =\frac{(500.12825 \mathrm{MHz})-(500.12750 \mathrm{MHz})}{500.13000 \mathrm{MHz}} \times 10^{6} \\
& =1.5
\end{aligned}
$$

On the same spectrometer the frequency separation of two resonances with chemical shifts $\delta_{1}=1.25$ and $\delta_{2}=5.75$ is found using eqn 12B.5b

$$
\begin{aligned}
v_{2}-v_{1} & =\left(\frac{v_{\text {spect }}}{10^{6}}\right)\left(\delta_{2}-\delta_{1}\right)=\left(\frac{500.13000 \times 10^{6} \mathrm{~Hz}}{10^{6}}\right)(5.75-1.25) \\
& =2250 \mathrm{~Hz}
\end{aligned}
$$

A note on good practice In much of the literature, chemical shifts are reported in parts per million, ppm , in recognition of the factor of $10^{6}$ in the definition; this is unnecessary. If you see ' $\delta=10$ ppm ', interpret it, and use it in eqn 12B.5, as $\delta=10$.

The relation between $\delta$ and $\sigma$ is obtained by substituting eqn 12B. 3 into eqn 12B.4a:


\begin{align*}
& \delta=\frac{(1-\sigma) \mathscr{B}_{0}-\left(1-\sigma^{\circ}\right) \mathscr{B}_{0}}{\left(1-\sigma^{\circ}\right) \mathscr{B}_{0}} \times 10^{6} \\
& =\frac{\sigma^{\circ}-\sigma}{1-\sigma^{\circ}} \times 10^{6} \approx\left(\sigma^{\circ}-\sigma\right) \times 10^{6} \quad \underbrace{\left|\sigma^{\circ}\right| \ll 1} \tag{12B.6}
\end{align*}


where $\sigma^{\circ}$ is the shielding constant of the reference standard. A decrease in $\sigma$ (reduction in shielding) therefore leads to an increase in $\delta$. Therefore, nuclei with large chemical shifts are said to be strongly deshielded. Some typical chemical shifts are given in Fig. 12B.1. As can be seen from the illustration, the nuclei of different elements have very different ranges of\\
\includegraphics[max width=\textwidth, center]{2025_02_27_d4b95d9718f0b9e2e6b9g-527(1)}

Figure 12B. 1 The range of typical chemical shifts for (a) ${ }^{1} \mathrm{H}$ resonances and (b) ${ }^{13} \mathrm{C}$ resonances.\\
chemical shifts. The ranges exhibit the variety of electronic environments of the nuclei in molecules: the higher the atomic number of the element, the greater the number of electrons around the nucleus and hence the greater the range of the extent of shielding. By convention, NMR spectra are plotted with $\delta$ decreasing from left to right.

\subsection*{Example 12B. 1 Interpreting a ${ }^{1} \mathrm{H}$ NMR spectrum}
Figure 12B. 2 shows the ${ }^{1} \mathrm{H}$ (proton) NMR spectrum of 1-methoxy-2-propanone, $\mathrm{CH}_{3} \mathrm{OCH}_{2} \mathrm{COCH}_{3}$. Account for the observed chemical shifts.

Collect your thoughts You need to consider the effect of any electron-withdrawing atom: it deshields strongly the protons to which it is bound and has a diminishing effect on more distant protons. To identify which of the B and C resonances correspond to H atoms 2 and 3 you can take either of two approaches. One is to look at the large compilations of chemical shift data available. The second approach is to make use of the 'integral' of a line, the area under the resonance peak, which is proportional to the number of nuclei giving rise to the peak. These integrals are commonly shown by step-like curves superimposed on the spectrum, as is the case in Fig. 12B.2: the integral is proportional to the height of the step.

The solution The H atoms labelled 2 and 3 are all attached to a C atom that is attached to the strongly electron-withdrawing O atom, whereas the H atoms labelled 1 are further away from any O atoms. You can expect the deshielding of H atoms 2 and 3 to be greater than that of H atoms 1 , and so the chemical shifts of 2 and 3 will be larger than that of 1 . Resonance $A$ at $\delta=2.2$ can therefore confidently be assigned to the H atoms at position 1. It is evident from the spectrum that the integral of peak $B$ is greater than that of $C$ (in fact they are in the ratio 3:2), immediately identifying peak $B$ as corresponding to the H atoms at position 3, and peak C to those at position 2.

Self-test 12B. 1 The NMR spectrum of ethanal (acetaldehyde) has lines at $\delta=2.20$ and $\delta=9.80$. Which feature can be assigned to the CHO proton?\\
\includegraphics[max width=\textwidth, center]{2025_02_27_d4b95d9718f0b9e2e6b9g-527}

Figure 12B. 2 The ${ }^{1} \mathrm{H}$ (proton) NMR spectrum of 1-methoxy-2-propanone. The step-like curve indicates the integral of the peak (the area under the peak) with the height of the step being proportional to the integral.

\subsection*{12B.2 The origin of shielding constants}
The calculation of shielding constants (and hence chemical shifts) is difficult because it requires detailed knowledge of the distribution of electron density in the ground and excited states and the electronic excitation energies of the molecule. Nevertheless, it is helpful to understand the different contributions to chemical shifts so that patterns and trends can be identified.

A useful approach is to assume that the observed shielding constant is the sum of three contributions:


\begin{equation*}
\sigma=\sigma(\text { local })+\sigma(\text { neighbour })+\sigma(\text { solvent }) \tag{12B.7}
\end{equation*}


The local contribution, $\sigma$ (local), is essentially the contribution of the electrons of the atom that contains the nucleus in question. The neighbouring group contribution, $\sigma$ (neighbour), is the contribution from the groups of atoms that form the rest of the molecule. The solvent contribution, $\sigma$ (solvent), is the contribution from the solvent molecules.

\subsection*{(a) The local contribution}
It is convenient to regard the local contribution to the shielding constant as the sum of a diamagnetic contribution, $\sigma_{d}$, and a paramagnetic contribution, $\sigma_{\mathrm{p}}$ :

\[
\sigma(\text { local })=\sigma_{\mathrm{d}}+\sigma_{\mathrm{p}} \quad \begin{array}{ll}
\text { Local contribution to }  \tag{12B.8}\\
\text { the shielding constant }
\end{array}
\]

The diamagnetic contribution arises from additional fields that oppose the applied magnetic field and hence shield the nucleus; $\sigma_{\mathrm{d}}$ is therefore positive. The paramagnetic contribution arises from additional fields that reinforce the applied field and hence lead to deshielding; $\sigma_{\mathrm{p}}$ is therefore negative.

The diamagnetic contribution arises from the ability of the applied field to generate a circulation of charge in the groundstate electron distribution. The circulation generates a magnetic field which opposes the applied field and hence shields the nucleus. The magnitude of $\sigma_{\mathrm{d}}$ depends on the electron density close to the nucleus and for atoms it can be calculated from the Lamb formula: ${ }^{1}$

$$
\sigma_{\mathrm{d}}=\frac{e^{2} \mu_{0}}{12 \pi m_{\mathrm{e}}}\langle 1 / r\rangle
$$

Lamb formula\\
(12B.9)\\
where $\mu_{0}$ is the vacuum permeability, $r$ is the electron-nucleus distance, and the angle brackets $\langle\ldots\rangle$ indicate an expectation value.

\footnotetext{${ }^{1}$ For a derivation, see our Molecular quantum mechanics (2011).
}\subsection*{Example 12B. 2 Using the Lamb formula}
Calculate the shielding constant for the nucleus in a free H atom.

Collect your thoughts To calculate $\sigma_{\mathrm{d}}$ from the Lamb formula, you need to calculate the expectation value of $1 / r$ for a hydrogen 1 s orbital. The radial part of the wavefunction can be found from Table 8A.1 and the angular part from Table 7F.1.

The solution The normalized wavefunction for a hydrogen 1s orbital is, in spherical polar coordinates (see The chemist's toolkit 21 in Topic 7F),

$$
\psi=\left(\frac{1}{\pi a_{0}^{3}}\right)^{1 / 2} \mathrm{e}^{-r / a_{0}}
$$

In this coordinate system the volume element is $\mathrm{d} \tau=$ $\sin \theta r^{2} \mathrm{~d} r \mathrm{~d} \theta \mathrm{~d} \phi$, so the expectation value of $1 / r$ is

$$
\begin{aligned}
\langle 1 / r\rangle & =\int \frac{\psi^{\star} \psi}{r} \mathrm{~d} \tau=\frac{1}{\pi a_{0}^{3}} \overbrace{\int_{0}^{2 \pi} \mathrm{~d} \phi \int_{0}^{\pi} \sin \theta \mathrm{d} \theta \int_{0}^{\infty} r \mathrm{e}^{-2 r / a_{0}} \mathrm{~d} r}^{2 \pi} \\
& =\frac{1}{\pi a_{0}^{3}} \times 4 \pi \times \frac{a_{0}^{2}}{4}=\frac{1}{a_{0}}
\end{aligned}
$$

Therefore,

$$
\begin{aligned}
\sigma_{\mathrm{d}} & =\frac{e^{2} \mu_{0}}{12 \pi m_{\mathrm{e}} a_{0}}=\frac{\left(1.602 \times 10^{-19} \mathrm{C}\right)^{2} \times\left(4 \pi \times 10^{-7} \mathrm{Js}^{2} \mathrm{C}^{-2} \mathrm{~m}^{-1}\right)}{12 \pi \times\left(9.109 \times 10^{-31} \mathrm{~kg}\right) \times\left(5.292 \times 10^{-11} \mathrm{~m}\right)} \\
& =\frac{\left(1.602 \times 10^{-19}\right)^{2} \times\left(4 \pi \times 10^{-7}\right)}{12 \pi \times\left(9.109 \times 10^{-31}\right) \times\left(5.292 \times 10^{-11}\right)} \times \frac{\mathrm{C}^{2} \mathrm{Js}^{2} \mathrm{C}^{-2} \mathrm{~m}^{-1}}{\mathrm{kgm}} \\
& =1.775 \times 10^{-5}
\end{aligned}
$$

where $1 \mathrm{~J}=1 \mathrm{~kg} \mathrm{~m}^{2} \mathrm{~s}^{-2}$ has been used.\\
Self-test 12B.2 Derive an expression for $\sigma_{\mathrm{d}}$ for a hydrogenic atom with nuclear charge $Z$.

$$
{ }^{0} v^{\partial} u \Psi Z \mathrm{I} /{ }^{0} n_{\tau}^{\prime} \partial Z={ }^{\mathrm{P}} \rho: \not \partial m s u V
$$

The diamagnetic contribution is the only contribution in closed-shell free atoms and when the electron distribution is spherically symmetric. In a molecule the core electrons near to a particular nucleus are likely to have spherical symmetry, even if the valence electron distribution is highly distorted. Therefore, core electrons contribute only to the diamagnetic part of the shielding. The diamagnetic contribution is broadly proportional to the electron density of the atom containing the nucleus of interest. It follows that the shielding is decreased if the electron density on the atom is reduced by the influence of an electronegative atom nearby. That reduction in shielding as the electronegativity of a neighbouring atom increases translates into an increase in the chemical shift $\delta$ (Fig. 12B.3).\\
\includegraphics[max width=\textwidth, center]{2025_02_27_d4b95d9718f0b9e2e6b9g-529(1)}

Figure 12B. 3 The variation of chemical shielding with electronegativity. The shifts for the methyl protons follow the simple expectation that increasing the electronegativity of the halogen will increase the chemical shift. However, to emphasize that chemical shifts are subtle phenomena, notice that the trend for the methylene protons is opposite to that expected. For these protons another contribution (the magnetic anisotropy of $\mathrm{C}-\mathrm{H}$ and $\mathrm{C}-\mathrm{X}$ bonds) is dominant.

The local paramagnetic contribution, $\sigma_{p}$, arises from the ability of the applied field to force electrons to circulate through the molecule by making use of orbitals that are unoccupied in the ground state. It is absent in free atoms and also in linear molecules (such as ethyne, $\mathrm{HC} \equiv \mathrm{CH}$ ) when the applied field lies along the symmetry axis; in this arrangement the electrons can circulate freely and the applied field is unable to force them into other orbitals. Large paramagnetic contributions can be expected for light atoms (because the valence electrons, and hence the induced currents, are close to the nucleus) and in molecules with low-lying excited states (because an applied field can then induce significant currents). In fact, the paramagnetic contribution is the dominant local contribution for atoms other than hydrogen.

\subsection*{(b) Neighbouring group contributions}
The neighbouring group contribution arises from the currents induced in nearby groups of atoms. Consider the influence of the neighbouring group X on the hydrogen atom in a molecule such as $\mathrm{H}-\mathrm{X}$. The applied field generates currents in the electron distribution of X and gives rise to an induced magnetic moment (an induced magnetic dipole) proportional to the applied field; the constant of proportionality is the magnetic susceptibility, $\chi$ (chi), of the group X: $\mu_{\text {induced }}=\chi \mathcal{B}_{0}$. The susceptibility is negative for a diamagnetic group because the induced moment is opposite to the direction of the applied field.

The induced moment gives rise to a magnetic field which is experienced by neighbouring nuclei. As is explained in The chemist's toolkit 27, a nucleus at distance $R$ and angle $\theta$ (defined in 1 ) from the induced moment experiences a local field that has the form

\subsection*{The chemist's toolkit 27}
\subsection*{Dipolar magnetic fields}
Standard electromagnetic theory gives the magnetic field at a point $r$ from a point magnetic dipole $\mu$ as

$$
\mathscr{B}=\frac{\mu_{0}}{4 \pi r^{3}}\left(\mu-\frac{3(\boldsymbol{\mu} \cdot \boldsymbol{r}) \boldsymbol{r}}{r^{2}}\right)
$$

where $\mu_{0}$ is the vacuum permeability (a fundamental constant with the defined value $4 \pi \times 10^{-7} \mathrm{~T}^{2} \mathrm{~J}^{-1} \mathrm{~m}^{3}$ ). The component of magnetic field in the $z$-direction is

$$
\mathscr{B}_{z}=\frac{\mu_{0}}{4 \pi r^{3}}\left(\mu_{z}-\frac{3(\boldsymbol{\mu} \cdot \boldsymbol{r}) z}{r^{2}}\right)
$$

with $z=r \cos \theta$, the $z$-component of the distance vector $r$. If the magnetic dipole is also parallel to the $z$-direction, it follows that

$$
\mathscr{B}_{z}=\frac{\mu_{0}}{4 \pi r^{3}}\left(\begin{array}{l}
\mu_{z} \\
\mu-\frac{3(\mu r \cos \theta)}{r^{2}} \overbrace{(r \cos \theta)}^{z} \\
r^{2} \cdot r
\end{array}\right)=\frac{\mu \mu_{0}}{4 \pi r^{3}}\left(1-3 \cos ^{2} \theta\right)
$$

\begin{center}
\includegraphics[max width=\textwidth]{2025_02_27_d4b95d9718f0b9e2e6b9g-529}
\end{center}

$$
\mathscr{B}_{\text {local }} \propto \frac{\mu_{\text {induced }}}{r^{3}}\left(1-3 \cos ^{2} \theta\right)
$$

Local dipolar field\\
(12B.10a)\\
This local field is parallel to the applied field, and the angle $\theta$ is measured from the direction of the applied field. Note that the strength of the field is inversely proportional to the cube of the distance $r$ between H and X . If the magnetic susceptibility is independent of the orientation of the molecule (that is, it is 'isotropic'), the local field averages to zero because, when averaged over a sphere, $1-3 \cos ^{2} \theta$ is zero (see Problem 12B.8). However, if the magnetic susceptibility varies with the orientation of the molecule with respect to the magnetic field, the local field may average to a non-zero value. For instance, suppose that the neighbouring group has axial symmetry (as might be the case for a triple bond): when the applied field is parallel to the symmetry axis the susceptibility is $\chi_{\|}$, and when it is perpendicular the susceptibility is $\chi_{\perp}$. After averaging over all orientations of the molecule the contribution to the shielding constant of a nucleus at a distance $R$ has the following form

$$
\sigma(\text { neighbour }) \propto\left(\chi_{\|}-\chi_{\perp}\right)\left(\frac{1-3 \cos ^{2} \Theta}{R^{3}}\right)
$$

\begin{center}
\includegraphics[max width=\textwidth]{2025_02_27_d4b95d9718f0b9e2e6b9g-530(1)}
\end{center}

Figure 12B.4 A depiction of the field arising from a point magnetic dipole. The three shades of colour represent the strength of field declining with distance (as $1 / R^{3}$ ), and each surface shows the angle dependence of the $z$-component of the field for each distance.\\
where $\Theta$ (uppercase theta) is the angle between the symmetry axis and the vector to the nucleus (2). Equation 12B.10b shows that the neighbouring group contribution may be positive or negative according to the relative magnitudes of the two magnetic susceptibilities and the direction given by $\Theta$. If $54.7^{\circ}<\Theta$ $<125.3^{\circ}$, then $1-3 \cos ^{2} \Theta$ is positive, but it is negative otherwise (Figs. 12B. 4 and 12B.5).\\
\includegraphics[max width=\textwidth, center]{2025_02_27_d4b95d9718f0b9e2e6b9g-530}

A special case of a neighbouring group effect is found in aromatic compounds. The strong anisotropy of the magnetic susceptibility of the benzene ring is ascribed to the ability of the field to induce a ring current, a circulation of electrons around the ring, when the field is applied perpendicular to the molecular plane. Protons in the plane are deshielded (Fig. 12B.6), but\\
\includegraphics[max width=\textwidth, center]{2025_02_27_d4b95d9718f0b9e2e6b9g-530(4)}

Figure 12B.5 The variation of the function $1-3 \cos ^{2} \Theta$ with the angle $\Theta$.\\
\includegraphics[max width=\textwidth, center]{2025_02_27_d4b95d9718f0b9e2e6b9g-530(2)}

Figure 12B.6 The shielding and deshielding effects of the ring current induced in the benzene ring by the applied field. Protons attached to the ring are deshielded but a proton attached to a substituent that projects above the ring is shielded.\\
any that happen to lie above or below the plane (as members of substituents of the ring) are shielded.

\subsection*{(c) The solvent contribution}
A solvent can influence the local magnetic field experienced by a nucleus in a variety of ways. Some of these effects arise from specific interactions between the solute and the solvent (such as hydrogen bond formation and other forms of Lewis acid-base complex formation). The anisotropy of the magnetic susceptibility of the solvent molecules, especially if they are aromatic, can also be the source of a local magnetic field. Moreover, if there are steric interactions that result in a loose but specific interaction between a solute molecule and a solvent molecule, then protons in the solute molecule may experience shielding or deshielding effects according to their location relative to the solvent molecule. An aromatic solvent such as benzene can give rise to local currents that shield or deshield a proton in a solute molecule. The arrangement shown in Fig. 12B. 7 leads to shielding of a proton on the solute molecule.\\
\includegraphics[max width=\textwidth, center]{2025_02_27_d4b95d9718f0b9e2e6b9g-530(3)}

Figure 12B.7 An aromatic solvent (benzene here) can give rise to local currents that shield or deshield a proton in a solute molecule. In this relative orientation of the solvent and solute, the proton on the solute molecule is shielded.

\subsection*{12B.3 The fine structure}
Figure 12B. 8 shows the ${ }^{1} \mathrm{H}$ (proton) NMR spectrum of chloroethane. In this molecule there are two different types of ${ }^{1} \mathrm{H}$, the methylene $\left(\mathrm{CH}_{2}\right)$ and the methyl $\left(\mathrm{CH}_{3}\right)$ protons, each with a characteristic chemical shift. In addition, the spectrum shows fine structure, the splitting of a resonance line into several components. The groups of lines are called multiplets.

This fine structure arises from scalar coupling in which the resonance frequency of one nucleus is affected by the spin state of another nucleus. Qualitatively, the effect of scalar coupling arises when the local magnetic field at one nucleus depends on the relative orientation of the other spin. If the spin is in one state ( $\alpha$, for instance) the local field is increased, whereas when the spin is in the other state ( $\beta$ in this case), the local field is decreased. Therefore, rather than there being one line in the spectrum, there are two because there are two possible values for the local field, corresponding to the second nucleus being either $\alpha$ or $\beta$.

The scalar coupling interaction is represented by a term $\left(h J / \hbar^{2}\right) \hat{I}_{1} \cdot \hat{I}_{2}$ in the hamiltonian, where $\hat{I}_{N}$, with $N=1$ or 2 , is the operator for the nuclear spin angular momentum of nucleus $N$. That the coupling term is a scalar product simply expresses the fact that the energy of interaction depends on the relative orientation of the spins of the two nuclei. The strength of the interaction is given by the value of the scalar coupling constant, $J$. The presence of $\hbar^{2}$ in $\left(h J / \hbar^{2}\right) \hat{I}_{1} \cdot \hat{\boldsymbol{I}}_{2}$ cancels the $\hbar^{2}$ arising from the eigenvalues of the two angular momenta, leaving the energy as $h J$, so $J$ is a frequency (measured in hertz, Hz ). The coupling constant can be positive or negative, and it is independent of the field strength.

If the Larmor frequencies of the two coupled nuclei are significantly different, they precess at very different rates and the $x$ - and $y$-components of their magnetic moments are never in step. Only the $z$-components remain in alignment whatever the precession rates, and so the only surviving term in the scalar product is $\left(h J / \hbar^{2}\right) \hat{I}_{1 z} \hat{I}_{2 z}$. The eigenvalues of each $\hat{I}_{N z}$ are $m_{I_{N}} \hbar$,\\
\includegraphics[max width=\textwidth, center]{2025_02_27_d4b95d9718f0b9e2e6b9g-531}

Figure 12B. 8 The ${ }^{1} \mathrm{H}$ (proton) NMR spectrum of chloroethane. The coloured letters denote the protons giving rise to each multiplet; the multiplets arise due to scalar coupling between the protons.\\
so it follows that the eigenvalues (the energies) of the coupling term are


\begin{equation*}
E_{m_{I_{1} m_{12}}}=h J m_{I_{1}} m_{I_{2}} \quad \text { Spin-spin coupling energy } \tag{12B.11}
\end{equation*}


\subsection*{(a) The appearance of the spectrum}
In NMR, letters far apart in the alphabet (typically A and X) are used to indicate nuclei with very different chemical shifts in the sense that the difference in chemical shift corresponds to a frequency that is large compared to $J$; letters close together (such as A and B) are used for nuclei with similar chemical shifts.

Consider first an AX system, a molecule that contains two spin- $\frac{1}{2}$ nuclei A and X with very different chemical shifts, so eqn 12B. 11 can be used for the spin-spin coupling energy. Nucleus A has two spin states with $m_{\mathrm{A}}= \pm \frac{1}{2}$ corresponding to the states denoted $\alpha_{\mathrm{A}}$ and $\beta_{\mathrm{A}}$. The X nucleus also has two spin states with $m_{\mathrm{x}}= \pm \frac{1}{2}\left(\alpha_{\mathrm{x}}\right.$ and $\left.\beta_{\mathrm{x}}\right)$. In the AX system there are therefore four spin states: $\alpha_{A} \alpha_{X}, \alpha_{A} \beta_{X}, \beta_{A} \alpha_{X}$, and $\beta_{A} \beta_{X}$. The energies of these states, neglecting any scalar coupling, are therefore


\begin{align*}
E_{m_{A} m_{x}} & =-\gamma_{\mathrm{N}} \hbar\left(1-\sigma_{\mathrm{A}}\right) \mathcal{B}_{0} m_{\mathrm{A}}-\gamma_{\mathrm{N}} \hbar\left(1-\sigma_{\mathrm{X}}\right) \mathcal{B}_{0} m_{\mathrm{X}} \\
& =-h v_{\mathrm{A}} m_{\mathrm{A}}-h v_{\mathrm{x}} m_{\mathrm{x}} \tag{12B.12a}
\end{align*}


where $v_{\mathrm{A}}$ and $v_{\mathrm{X}}$ are the Larmor frequencies of A and X (eqn 12B.3). This expression gives the four levels illustrated on the left of Fig. 12B.9. When spin-spin coupling is included (by using eqn 12B.11) the energy levels are


\begin{equation*}
E_{m_{\mathrm{A}} m_{x}}=-h v_{\mathrm{A}} m_{\mathrm{A}}-h v_{\mathrm{x}} m_{\mathrm{x}}+h J m_{\mathrm{A}} m_{\mathrm{x}} \tag{12B.12b}
\end{equation*}


The resulting energy level diagram (for $J>0$ ) is shown on the right of Fig. 12B.9. The $\alpha_{A} \alpha_{X}$ and $\beta_{A} \beta_{\mathrm{X}}$ states are both raised by $\frac{1}{4} h J$ and the $\alpha_{A} \beta_{x}$ and $\beta_{A} \alpha_{x}$ states are both lowered by $\frac{1}{4} h J$. For $J>0$, the effect of the coupling term is to lower the energy of the $\alpha_{A} \beta_{X}$ and $\beta_{A} \alpha_{X}$ states, and raise the energy of the other two states. The opposite is the case for $J<0$.

In a transition, only one nucleus changes its orientation, so the selection rule is that either $m_{\mathrm{A}}$ or $m_{\mathrm{x}}$ can change by $\pm 1$, but not both. There are two transitions in which the spin state of A changes while that of $X$ remains fixed: $\beta_{A} \alpha_{X} \leftarrow \alpha_{A} \alpha_{X}$ and $\beta_{A} \beta_{x} \leftarrow \alpha_{A} \beta_{x}$. They are shown in Fig. 12B. 9 and in a slightly different form in Fig. 12B.10. The energies of the transitions are


\begin{equation*}
\Delta E=h v_{\mathrm{A}} \pm \frac{1}{2} h J \tag{12B.13a}
\end{equation*}


The spectrum due to A transitions therefore consists of a doublet of separation $J$ centred on the Larmor frequency of A (Fig. 12B.11). Similar remarks apply to the transitions in which the spin state of X changes while that of A remains fixed. These are also shown in Figs 12B. 9 and 12B.10, and the transition energies are


\begin{equation*}
\Delta E=h v_{\mathrm{x}} \pm \frac{1}{2} h J \tag{12B.13b}
\end{equation*}


\begin{center}
\includegraphics[max width=\textwidth]{2025_02_27_d4b95d9718f0b9e2e6b9g-532(4)}
\end{center}

Figure 12B. 9 The energy levels of an $A X$ spin system. The four levels on the left are those in the case of no spin-spin coupling. The four levels on the right show how a positive spin-spin coupling constant affects the energies. The red arrows show the allowed transitions in which A goes from the $\alpha$ to the $\beta$ spin state, while the spin state of $X$ remains unchanged; the blue arrows show the corresponding transitions of the $X$ nucleus. The effect of the coupling on the energy levels has been exaggerated greatly for clarity; in practice, the change in energy caused by spin-spin coupling is much smaller than that caused by the applied field.

It follows that there is a doublet with the same separation $J$, but now centred on the Larmor frequency of X (as shown in Fig. 12B.11). Overall, the spectrum of an AX spin system consists of two doublets.\\
If there is another X nucleus in the molecule with the same chemical shift as the first X (giving an $\mathrm{AX}_{2}$ spin system), the X resonance is split into a doublet by A, just as for AX (Fig. 12B.12). The resonance of $A$ is split into a doublet by one $X$, and each line of the doublet is split again by the same amount by the second X (Fig. 12B.13). This splitting results in three\\
\includegraphics[max width=\textwidth, center]{2025_02_27_d4b95d9718f0b9e2e6b9g-532(1)}

Figure 12B.10 An alternative depiction of the energy levels and transitions shown in Fig. 12B.9. Once again, the effect of spin-spin coupling has been exaggerated.\\
\includegraphics[max width=\textwidth, center]{2025_02_27_d4b95d9718f0b9e2e6b9g-532(3)}

Figure 12B. 11 The effect of spin-spin coupling on an AX spectrum. Each resonance is split into two lines, a doublet, separated by J. There is a doublet centred on the Larmor frequency (chemical shift) of A, and one centred on the Larmor frequency of $B$.\\
\includegraphics[max width=\textwidth, center]{2025_02_27_d4b95d9718f0b9e2e6b9g-532}

Figure 12B. 12 The $X$ resonance of an $A X_{2}$ spin system is also a doublet, because the two equivalent X nuclei behave like a single nucleus; however, the overall absorption is twice as intense as that of an AX spin system.\\
\includegraphics[max width=\textwidth, center]{2025_02_27_d4b95d9718f0b9e2e6b9g-532(2)}

Figure 12B. 13 The origin of the 1:2:1 triplet in the A resonance of an $\mathrm{AX}_{2}$ spin system. The resonance of $A$ is split into two by coupling with one $X$ nucleus (as shown in the inset), and then each of those two lines is split into two by coupling to the second $X$ nucleus. Because each $X$ nucleus causes the same splitting, the two central transitions are coincident and give rise to an absorption line of double the intensity of the outer lines.\\
\includegraphics[max width=\textwidth, center]{2025_02_27_d4b95d9718f0b9e2e6b9g-533(3)}

Figure 12B. 14 The origin of the 1:3:3:1 quartet in the $A$ resonance of an $\mathrm{AX}_{3}$ species. The third X nucleus splits each of the lines shown in Fig. 12B. 13 for an $\mathrm{AX}_{2}$ species into a doublet, and the intensity distribution reflects the number of transitions that have the same energy.\\
lines in the intensity ratio 1:2:1 (because the central frequency can be obtained in two ways).

Three equivalent X nuclei (an $\mathrm{AX}_{3}$ spin system) split the resonance of A into four lines of intensity ratio 1:3:3:1 (Fig. 12B.14). The X resonance remains a doublet as a result of the splitting caused by A. In general, $N$ equivalent spin- $\frac{1}{2}$ nuclei split the resonance of a nearby spin or group of equivalent spins into $N+1$ lines with an intensity distribution given by Pascal's triangle (3). Successive rows of this triangle are formed by adding together the two adjacent numbers in the line above.\\
\includegraphics[max width=\textwidth, center]{2025_02_27_d4b95d9718f0b9e2e6b9g-533(4)}

\subsection*{Example 12B. 3}
Accounting for the fine structure in a spectrum

Account for the fine structure in the ${ }^{1} \mathrm{H}$ (proton) NMR spectrum of chloroethane shown in Fig. 12B.8.\\
Collect your thoughts You need to consider how each group of equivalent protons (for instance, the three methyl protons) splits the resonances of the other groups of protons. There is no splitting within groups of equivalent protons. You can identify the pattern of intensities within a multiplet by referring to Pascal's triangle.

The solution The three protons of the $\mathrm{CH}_{3}$ group split the resonance of the $\mathrm{CH}_{2}$ protons into a 1:3:3:1 quartet with a splitting $J$. Likewise, the two protons of the $\mathrm{CH}_{2}$ group split the resonance of the $\mathrm{CH}_{3}$ protons into a 1:2:1 triplet with the same splitting $J$.

Self-test 12B. 3 What fine-structure can be expected in the ${ }^{1} \mathrm{H}$ spectrum, and in the ${ }^{15} \mathrm{~N}$ spectrum, of ${ }^{15} \mathrm{NH}_{4}^{+}$? Nitrogen- 15 is a spin- $\frac{1}{2}$ nucleus.\\
\includegraphics[max width=\textwidth, center]{2025_02_27_d4b95d9718f0b9e2e6b9g-533(2)}\\
\includegraphics[max width=\textwidth, center]{2025_02_27_d4b95d9718f0b9e2e6b9g-533}

\subsection*{(b) The magnitudes of coupling constants}
The scalar coupling constant of two nuclei separated by $N$ bonds is denoted ${ }^{N} J$, with subscripts to indicate the types of nuclei involved. Thus, ${ }^{1} J_{\mathrm{CH}}$ is the coupling constant for a proton joined directly to a ${ }^{13} \mathrm{C}$ atom, and ${ }^{2} J_{\mathrm{CH}}$ is the coupling constant when the same two nuclei are separated by two bonds (as in $\left.{ }^{13} \mathrm{C}-\mathrm{C}-\mathrm{H}\right)$. A typical value of ${ }^{1} J_{\mathrm{CH}}$ is in the range $120-250 \mathrm{~Hz}$; ${ }^{2} J_{\mathrm{CH}}$ is between 10 and 20 Hz . Both ${ }^{3} J$ and ${ }^{4} J$ can give detectable effects in a spectrum, but couplings over larger numbers of bonds can generally be ignored.

As remarked in the discussion following eqn 12B.12b, the sign of $J_{\mathrm{XY}}$ determines whether a particular energy level is raised or lowered as a result of the coupling interaction. If $J>0$, the levels with antiparallel spins are lowered in energy, whereas if $J<0$ the levels with parallel spins are lowered. Experimentally, it is found that ${ }^{1} J_{\mathrm{CH}}$ is\\
\includegraphics[max width=\textwidth, center]{2025_02_27_d4b95d9718f0b9e2e6b9g-533(1)}

4 invariably positive, ${ }^{2} J_{\mathrm{HH}}$ is often negative, and ${ }^{3} J_{\mathrm{HH}}$ is often positive. An additional point is that $J$ varies with the dihedral angle between the bonds (Fig. 12B.15). Thus, a ${ }^{3} J_{\mathrm{HH}}$ coupling constant is often found to depend on the dihedral angle $\phi$ (4) according to the Karplus equation:


\begin{equation*}
{ }^{3} J_{\mathrm{HH}}=A+B \cos \phi+C \cos 2 \phi \quad \text { Karplus equation } \tag{12B.14}
\end{equation*}


with $A, B$, and $C$ empirical constants with values close to $+7 \mathrm{~Hz},-1 \mathrm{~Hz}$, and +5 Hz , respectively, for an HCCH fragment. It follows that the measurement of ${ }^{3} J_{\mathrm{HH}}$ in a series of related compounds can be used to determine their conformations. The coupling constant ${ }^{1} J_{\mathrm{CH}}$ also depends on the hybridization of the C atom, as the following values indicate:

$$
\begin{array}{llll} 
& \mathrm{sp} & \mathrm{sp}^{2} & \mathrm{sp}^{3} \\
{ }^{1} J_{\mathrm{CH}} / \mathrm{Hz} & 250 & 160 & 125
\end{array}
$$

\begin{center}
\includegraphics[max width=\textwidth]{2025_02_27_d4b95d9718f0b9e2e6b9g-534}
\end{center}

Figure 12B. 15 The variation of the spin-spin coupling constant with dihedral angle predicted by the Karplus equation for an HCCH group and an HNCH group.

\subsection*{Brief illustration 12B.2}
The investigation of $\mathrm{H}-\mathrm{N}-\mathrm{C}-\mathrm{H}$ couplings in polypeptides can help reveal their conformation. For ${ }^{3} J_{\mathrm{HH}}$ coupling in such a group, $A=+5.1 \mathrm{~Hz}, B=-1.4 \mathrm{~Hz}$, and $C=+3.2 \mathrm{~Hz}$. For a helical polymer, $\phi$ is close to $120^{\circ}$, which gives ${ }^{3} J_{\mathrm{HH}} \approx 4 \mathrm{~Hz}$. For the sheet-like conformation, $\phi$ is close to $180^{\circ}$, which gives ${ }^{3} J_{\mathrm{HH}} \approx$ 10 Hz . Experimental measurements of the value of ${ }^{3} J_{\mathrm{HH}}$ should therefore make it possible to distinguish between the two possible structures.

\subsection*{(c) The origin of spin-spin coupling}
Some insight into the origin of coupling, if not its precise magnitude-or always reliably its sign-can be obtained by considering the magnetic interactions within molecules. A nucleus with the $z$-component of its spin angular momentum specified by the quantum number $m_{I}$ gives rise to a magnetic field with $z$-component $\mathcal{B}_{\text {nuc }}$ at a distance $R$, where, to a good approximation,


\begin{equation*}
\mathcal{B}_{\mathrm{nuc}}=\frac{\gamma_{\mathrm{N}} \hbar \mu_{0}}{4 \pi R^{3}}\left(1-3 \cos ^{2} \theta\right) m_{I} \tag{12B.15}
\end{equation*}


The angle $\theta$ is defined in (1); this expression is a version of eqn 12B.10a. However, in solution molecules tumble rapidly so it is necessary to average $\mathcal{B}_{\text {nuc }}$ over all values of $\theta$. As has already been noted, the average of $1-3 \cos ^{2} \theta$ is zero, therefore the direct dipolar interaction between spins cannot account for the fine structure seen in the spectra of molecules in solution.

\subsection*{Brief illustration 12B. 3}
There can be a direct dipolar interaction between nuclei in solids, where the molecules do not rotate. The $z$-component of the magnetic field arising from a ${ }^{1} \mathrm{H}$ nucleus with $m_{I}=+\frac{1}{2}$, at $R=0.30 \mathrm{~nm}$, and at an angle $\theta=0$ is

$$
\begin{aligned}
\mathcal{B}_{\text {nuc }} & =\frac{\overbrace{\left(2.821 \times 10^{-26} \mathrm{JT}^{-1}\right)}^{\gamma_{1} \hbar} \times \overbrace{\left(4 \pi \times 10^{-7} \mathrm{~T}^{2} \mathrm{~J}^{-1} \mathrm{~m}^{3}\right)}^{\mu_{0}} \times \overbrace{(-1)}^{\left(1-3 \cos ^{2} \theta\right) m_{/}}}{4 \pi \times \underbrace{\left.3.0 \times 10^{-10} \mathrm{~m}\right)^{3}}_{R^{3}}} \\
& =1.0 \times 10^{-4} \mathrm{~T}=0.10 \mathrm{mT}
\end{aligned}
$$

Spin-spin coupling in molecules in solution can be explained in terms of the polarization mechanism, in which the interaction is transmitted through the bonds. The simplest case to consider is that of $J_{\mathrm{XY}}$, where X and Y are spin- $\frac{1}{2}$ nuclei joined by an electron-pair bond. The coupling mechanism depends on the fact that the energy depends on the relative orientation of the bonding electrons and the nuclear spins. This electron-nucleus coupling is magnetic in origin, and may be either a dipolar interaction or a Fermi contact interaction. A pictorial description of the latter is as follows. First, regard the magnetic moment of the nucleus as arising from the circulation of a current in a tiny loop with a radius similar to that of the nucleus (Fig. 12B.16). Far from the nucleus the field generated by this loop is indistinguishable from the field generated by a point magnetic dipole. Close to the loop, however, the field differs from that of a point dipole. The magnetic interaction between this non-dipolar field and the electron's magnetic moment is the contact interaction. The contact in-teraction-essentially the failure of the point-dipole approxi-mation-depends on the very close approach of an electron to the nucleus and hence can occur only if the electron occupies an s orbital (which is the reason why ${ }^{1} J_{\mathrm{CH}}$ depends on the hybridization ratio).\\
\includegraphics[max width=\textwidth, center]{2025_02_27_d4b95d9718f0b9e2e6b9g-534(1)}

Figure 12B. 16 The origin of the Fermi contact interaction. From far away, the magnetic field pattern arising from a ring of current (representing the rotating charge of the nucleus, the pale grey sphere) is that of a point dipole. However, if an electron can sample the field close to the region indicated by the sphere, the field distribution differs significantly from that of a point dipole. For example, if the electron can penetrate the sphere, then the spherical average of the field it experiences is not zero.

Suppose that it is energetically favourable for an electron spin and a nuclear spin to be antiparallel (as is the case for a proton and an electron in a hydrogen atom). If the X nucleus is $\alpha$, a $\beta$ electron of the bonding pair will tend to be found nearby, because that is an energetically favourable arrangement (Fig. 12B.17). The second electron in the bond, which must have $\alpha$ spin if the other is $\beta$ (by the Pauli principle; Topic $8 B$ ), will be found mainly at the far end of the bond because electrons tend to stay apart to reduce their mutual repulsion. Because it is energetically favourable for the spin of Y to be antiparallel to an electron spin, a Y nucleus with $\beta$ spin has a lower energy than when it has $\alpha$ spin. The opposite is true when $X$ is $\beta$, for now the $\alpha$ spin of $Y$ has the lower energy. In other words, the antiparallel arrangement of nuclear spins lies lower in energy than the parallel arrangement as a result of their magnetic coupling with the bond electrons. That is, ${ }^{1} J_{\mathrm{CH}}$ is positive.

To account for the value of ${ }^{2} J_{\mathrm{XY}}$, as for ${ }^{2} J_{\mathrm{HH}}$ in $\mathrm{H}-\mathrm{C}-\mathrm{H}$, a mechanism is needed that can transmit the spin alignments through the central C atom (which may be ${ }^{12} \mathrm{C}$, with no nuclear spin of its own). In this case (Fig. 12B.18), an X nucleus with $\alpha$ spin polarizes the electrons in its bond, and the $\alpha$ electron is likely to be found closer to the C nucleus. The more favourable arrangement of two electrons on the same atom is with their spins parallel (Hund's rule, Topic 8B), so the more favourable arrangement is for the $\alpha$ electron of the neighbouring bond to be close to the $C$ nucleus. Consequently, the $\beta$ electron of that bond is more likely to be found close to the Y nucleus, and therefore that nucleus will have a lower energy if it is $\alpha$. Hence, according to this mechanism, the lower energy will be obtained if the Y spin is parallel to that of X. That is, ${ }^{2} J_{\mathrm{HH}}$ is negative.

The coupling of nuclear spin to electron spin by the Fermi contact interaction is most important for proton spins, but it is not necessarily the most important mechanism for other nuclei. These nuclei may also interact by a dipolar mechanism with the electron magnetic moments and with their orbital motion, and there is no simple way of specifying whether $J$ will be positive or negative. The dipolar interaction does not\\
\includegraphics[max width=\textwidth, center]{2025_02_27_d4b95d9718f0b9e2e6b9g-535}

Figure 12B. 17 The polarization mechanism for spin-spin coupling ( ${ }^{\prime} J_{\text {CH }}$ ). The two arrangements have slightly different energies. In this case, $J$ is positive, corresponding to a lower energy when the nuclear spins are antiparallel.\\
\includegraphics[max width=\textwidth, center]{2025_02_27_d4b95d9718f0b9e2e6b9g-535(1)}

Figure 12B. 18 The polarization mechanism for ${ }^{2} J_{\mathrm{HH}} \mathrm{spin}-\mathrm{spin}$ coupling. The spin information is transmitted from one bond to the next by a version of the mechanism that accounts for the lower energy of electrons with parallel spins in different atomic orbitals (Hund's rule of maximum multiplicity). In this case, $J<$ 0 , corresponding to a lower energy when the nuclear spins are parallel.\\
average to zero as the molecule tumbles if it accounts for the interaction of both nuclei with their surrounding electrons because then $1-3 \cos ^{2} \theta$ appears as its square and therefore with a non-negative value at all orientations, and its average value is no longer zero.

\subsection*{(d) Equivalent nuclei}
A group of identical nuclei are chemically equivalent if they are related by a symmetry operation of the molecule and have the same chemical shifts. Chemically equivalent nuclei are from atoms that would be regarded as 'equivalent' according to ordinary chemical criteria. Nuclei are magnetically equivalent if, as well as being chemically equivalent, they also have identical spin-spin interactions with any other magnetic nuclei in the molecule.

\subsection*{Brief illustration 12B.4}
The difference between chemical and magnetic equivalence is illustrated by $\mathrm{CH}_{2} \mathrm{~F}_{2}$ and $\mathrm{H}_{2} \mathrm{C}=\mathrm{CF}_{2}$ (recall that ${ }^{19} \mathrm{~F}$ is a spin- $\frac{1}{2}$ nucleus). In each of these molecules the ${ }^{1} \mathrm{H}$ nuclei (protons) are chemically equivalent because they are related by symmetry. The protons in $\mathrm{CH}_{2} \mathrm{~F}_{2}$ are magnetically equivalent, but those in $\mathrm{CH}_{2}=\mathrm{CF}_{2}$ are not. One proton in the latter has a cis spin-coupling interaction with a given $F$ nucleus whereas the other proton has a trans interaction with the same nucleus. In contrast, in $\mathrm{CH}_{2} \mathrm{~F}_{2}$ each proton has the same coupling to both fluorine nuclei since the bonding pathway between them is the same.

Strictly speaking, in a molecule such as $\mathrm{CH}_{3} \mathrm{CH}_{2} \mathrm{Cl}$ the three $\mathrm{CH}_{3}$ protons are magnetically inequivalent because each may have a different coupling to the $\mathrm{CH}_{2}$ protons on account of the\\
different dihedral angles between the protons. However, the three $\mathrm{CH}_{3}$ protons are in practice made magnetically equivalent by the rapid rotation of the $\mathrm{CH}_{3}$ group, which averages out any differences. The spectra of molecules with chemically equivalent but magnetically inequivalent sets of spins can become very complicated (e.g. the proton and ${ }^{19} \mathrm{~F}$ spectra of $\mathrm{H}_{2} \mathrm{C}=\mathrm{CF}_{2}$ each consist of 12 lines); we shall not consider such spectra further.

An important feature of chemically equivalent magnetic nuclei is that, although they do couple together, the coupling has no effect on the appearance of the spectrum. How this comes about can be illustrated by considering the case of an $A_{2}$ spin system. The first step is to establish the energy levels.

\subsection*{How is that done? 12B. 1}
Deriving the energy levels of an $\mathrm{A}_{2}$ system

Consider an $\mathrm{A}_{2}$ system of two spin- $\frac{1}{2}$ nuclei, and first consider the energy levels in the absence of spin-spin coupling. When considering spin-spin coupling, be prepared to use the complete expression for the energy (the one proportional to $\boldsymbol{I}_{1} \cdot \boldsymbol{I}_{2}$ ), because the Larmor frequencies are the same and the approximate form $\left(I_{12} I_{2 z}\right)$ cannot be used as it is applicable only when the Larmor frequencies are very different.

Step 1 Identify the states and their energies in the absence of spin-spin coupling\\
There are four energy levels; they can be classified according to their total spin angular momentum $I_{\text {tot }}$ (the analogue of $S$ for several electrons) and its projection on to the $z$-axis, given by the quantum number $M_{I}$. There are three states with $I_{\text {tot }}=$ 1 , and one further state with $I_{\text {tot }}=0$ :

\begin{center}
\begin{tabular}{lll}
Spins parallel, $I_{\text {tot }}=1:$ & $M_{I}=+1$ & $\alpha \alpha$ \\
 & $M_{I}=0$ & $\left(1 / 2^{1 / 2}\right)\{\alpha \beta+\beta \alpha\}$ \\
Spins paired, $I_{\text {tot }}=0:$ & $M_{I}=-1$ & $\beta \beta$ \\
 & $M_{I}=0$ & $\left(1 / 2^{1 / 2}\right)\{\alpha \beta-\beta \alpha\}$ \\
\end{tabular}
\end{center}

The effect of a magnetic field on these four states is shown on the left-hand side of Fig. 12B.19: the two states with $M_{I}=0$ are unaffected by the field as they are composed of equal proportions of $\alpha$ and $\beta$ spins, and both spins have the same Larmor frequency.

Step 2 Allow for spin-spin interaction\\
The scalar product in the expression $E=\left(h J / \hbar^{2}\right) \boldsymbol{I}_{1} \cdot \boldsymbol{I}_{2}$ can be expressed in terms of the total nuclear spin $\boldsymbol{I}_{\text {tot }}=\boldsymbol{I}_{1}+\boldsymbol{I}_{2}$ by noting that

$$
I_{\mathrm{tot}}^{2}=\left(\boldsymbol{I}_{1}+\boldsymbol{I}_{2}\right) \cdot\left(\boldsymbol{I}_{1}+\boldsymbol{I}_{2}\right)=I_{1}^{2}+I_{2}^{2}+2 \boldsymbol{I}_{1} \cdot \boldsymbol{I}_{2}
$$

Rearranging this expression to

$$
\boldsymbol{I}_{1} \cdot \boldsymbol{I}_{2}=\frac{1}{2}\left\{I_{\mathrm{tot}}^{2}-I_{1}^{2}-I_{2}^{2}\right\}
$$

\begin{center}
\includegraphics[max width=\textwidth]{2025_02_27_d4b95d9718f0b9e2e6b9g-536}
\end{center}

Figure 12B. 19 The energy levels of an $\mathrm{A}_{2}$ spin system in the absence of spin-spin coupling are shown on the left. When spinspin coupling is taken into account, the energy levels on the right are obtained. Note that the effect of spin-spin coupling is to raise the three states with total nuclear spin $I_{\text {total }}=1$ (the triplet) by the same amount ( $J$ is positive); in contrast, the one state with $I_{\text {total }}=$ 0 (the singlet) is lowered in energy. The only allowed transitions, indicated by red arrows, are those for which $\Delta \mathrm{I}_{\text {total }}=0$ and $\Delta M_{l}=$ $\pm 1$. These two transitions occur at the same resonance frequency as they would have in the absence of spin-spin coupling.\\
and replacing the square magnitudes by their quantum mechanical values gives:

$$
\boldsymbol{I}_{1} \cdot \boldsymbol{I}_{2}=\frac{1}{2}\left\{I_{\text {tot }}\left(I_{\text {tot }}+1\right)-I_{1}\left(I_{1}+1\right)-I_{2}\left(I_{2}+1\right)\right\} \hbar^{2}
$$

Then, because $I_{1}=I_{2}=\frac{1}{2}$, it follows that

$$
E=\frac{1}{2} h J\left\{I_{\mathrm{tot}}\left(I_{\mathrm{tot}}+1\right)-\frac{3}{2}\right\}
$$

For parallel spins, $I_{\text {tot }}=1$ and $E=+\frac{1}{4} h J$; for antiparallel spins $I_{\text {tot }}=0$ and $E=-\frac{3}{4} h J$, as shown on the right-hand side of Fig. 12B.19.

The calculation shows that the three states with $I_{\text {total }}=1$ all move in energy in the same direction and by the same amount. The single state with $I_{\text {total }}=0$ moves three times as much in the opposite direction. In the resonance transition, the relative orientation of the nuclei cannot change, so there are no transitions between states of different $I_{\text {total }}$. The selection rule $\Delta M_{I}= \pm 1$ also applies, and arises from the conservation of angular momentum and the unit spin of the photon. As shown in Fig. 12B.19, there are only two allowed transitions and because they have the same energy spacing they appear at the same frequency in the spectrum. Hence, the spin-spin coupling interaction does not affect the appearance of the spectrum of an $\mathrm{A}_{2}$ molecule.

\subsection*{(e) Strongly coupled nuclei}
The multiplets seen in NMR spectra due to the presence of spin-spin coupling are relatively simple to analyse provided the difference in chemical shifts between any two coupled spins is\\
much greater than the value of the spin-spin coupling constant between them. This limit is often described as weak coupling, and the resulting spectra are described as first-order spectra.

When the difference in chemical shifts is comparable to the value of the spin-spin coupling constant, the multiplets take on a more complex form. Such spin systems are said to be strongly coupled, and the spectra are described as secondorder. In such spectra the lines shift from where they are expected in the weak coupling case, their intensities change, and in some cases additional lines appear. Strongly coupled spectra are more difficult to analyse in the sense that the relation between the frequencies of the lines in the spectrum and the values of chemical shifts and coupling constants is not as straightforward as in the weakly coupled case.

Figure 12B. 20 shows NMR spectra for two coupled spins as a function of the difference in chemical shift between the two spins. In Fig. 12B.20a (an AX species) this difference is large enough for the weak coupling limit to apply and two doublets are detected, with all lines having the same intensity. As the shift difference decreases the inner two lines gain intensity at the expense of the outer lines, and in the limit that the shift difference is zero (an $\mathrm{A}_{2}$ species), the outer lines disappear and the inner lines converge.

If the two nuclei belong to different elements (e.g. ${ }^{1} \mathrm{H}$ and ${ }^{13} \mathrm{C}$ ), or different isotopes of the same element (e.g. ${ }^{1} \mathrm{H}$ and ${ }^{2} \mathrm{H}$ ), the fact that they have widely different Larmor frequencies means that the spin system will always be weakly coupled, and hence described as AX. If the two nuclei are of the same element the spin system is described as homonuclear, whereas if they are of different elements the system is described as heteronuclear.\\
\includegraphics[max width=\textwidth, center]{2025_02_27_d4b95d9718f0b9e2e6b9g-537(1)}

Figure 12B.20 The NMR spectra of (a) an AX system and (d) a 'nearly $\mathrm{A}_{2}$ ' system are simple 'first-order' spectra (for an actual $\mathrm{A}_{2}$ system, $\Delta \delta=0$ ). At intermediate relative values of the chemical shift difference and the spin-spin coupling (b and c), more complex 'strongly coupled' spectra are obtained. Note how the inner two lines of the $A X$ spectrum move together, grow in intensity, and form the single central line of the $A_{2}$ spectrum.

\subsection*{12B.4 Exchange processes}
The appearance of an NMR spectrum is changed if magnetic nuclei can jump rapidly between different environments. For example, consider the molecule $\mathrm{N}, \mathrm{N}$-dimethylmethanamide, $\mathrm{HCON}\left(\mathrm{CH}_{3}\right)_{2}$, in which the $\mathrm{O}-\mathrm{C}-\mathrm{N}$ fragment is planar, and there is restricted rotation about the $\mathrm{C}-\mathrm{N}$ bond. The lowest energy conformation is shown in Fig. 12B.21. In this conformation the two methyl groups are not equivalent because one is cis and the other is trans to the carbonyl group. The two groups therefore have different environments and hence different chemical shifts.

Rotation by $180^{\circ}$ about the $\mathrm{C}-\mathrm{N}$ bond gives the same conformation, but it exchanges the $\mathrm{CH}_{3}$ groups between the two environments. When the jumping rate of this process is low, the spectrum shows a distinct line for each $\mathrm{CH}_{3}$ environment. When the rate is fast, the spectrum shows a single line at the mean of the two chemical shifts. At intermediate rates, the lines start to broaden and eventually coalesce into a single broad line. Coalescence of the two lines occurs when


\begin{equation*}
\tau=\frac{2^{1 / 2}}{\pi \delta v} \quad \text { Condition for coalescence of two NMR lines } \tag{12B.16}
\end{equation*}


where $\tau$ is the lifetime of an environment and $\delta v$ is the difference between the Larmor frequencies of the two environments.\\
\includegraphics[max width=\textwidth, center]{2025_02_27_d4b95d9718f0b9e2e6b9g-537(2)}

Figure 12B. 21 In this molecule the two methyl groups are in different environments and so will have different chemical shifts. Rotation about the $\mathrm{C}-\mathrm{N}$ bond interchanges the two groups, so that a particular methyl group is swapped between environments.

\subsection*{Brief illustration 12B. 5}
The NO group in $\mathrm{N}, \mathrm{N}$-dimethylnitrosamine, $\left(\mathrm{CH}_{3}\right)_{2} \mathrm{~N}-\mathrm{NO}$ (5), rotates about the $\mathrm{N}-\mathrm{N}$ bond and, as a result, the magnetic environments of the two $\mathrm{CH}_{3}$ groups are interchanged. The two $\mathrm{CH}_{3}$ resonances are separated by 390 Hz in a 600 MHz spectrometer. According to eqn 12B.16,

$$
\tau=\frac{2^{1 / 2}}{\pi \times\left(390 \mathrm{~s}^{-1}\right)}=1.2 \mathrm{~ms}
$$

\begin{center}
\includegraphics[max width=\textwidth]{2025_02_27_d4b95d9718f0b9e2e6b9g-537}
\end{center}

\footnotetext{$5 \mathrm{~N}, \mathrm{~N}$-DimethyInitrosamine
}Such a lifetime corresponds to a (first-order) rate constant of $1 / \tau=870 \mathrm{~s}^{-1}$. It follows that the signal will collapse to a single line when the rate constant for interconversion exceeds this value.

A similar explanation accounts for the loss of fine structure for protons that can exchange with the solvent. For example, the resonance from the OH group in the spectrum of ethanol appears as a single line (Fig. 12B.22). In this molecule the hydroxyl protons are able to exchange with the protons in water which, unless special precautions are taken, is inevitably present as an impurity in the (organic) solvent. When this chemical exchange, an exchange of atoms, occurs, a molecule ROH with an $\alpha$-spin proton (written as $\mathrm{ROH}_{\alpha}$ ) rapidly converts to $\mathrm{ROH}_{\beta}$ and then perhaps to $\mathrm{ROH}_{\alpha}$ again because the protons provided by the water molecules in successive exchanges have random spin orientations.

If the rate constant for the exchange process is fast compared to the value of the coupling constant $J$, in the sense $1 / \tau$ $\gg J$, the two lines merge and no splitting is seen. Because the values of coupling constants are typically just a few hertz, even rather slow exchange leads to the loss of the splitting. In the case of OH groups, only by rigorously excluding water from the solvent can the exchange rate be made slow enough that splittings due to coupling to OH protons are observed.

\subsection*{12B.5 Solid-state NMR}
In contrast to the narrow lines seen in the NMR spectra of samples in solution, the spectra from solid samples give broad lines, often to the extent that chemical shifts are not resolved. Nevertheless, there are good reasons for seeking to overcome these difficulties. They include the possibility that a compound is unstable in solution or that it is insoluble. Moreover, many species, such as polymers (both synthetic and naturally\\
\includegraphics[max width=\textwidth, center]{2025_02_27_d4b95d9718f0b9e2e6b9g-538}

Figure 12B.22 The ${ }^{1} \mathrm{H}$ (proton) NMR spectrum of ethanol. The coloured letters denote the protons giving rise to each multiplet. Due to chemical exchange between the OH proton and water molecules present in the solvent, no splittings due to coupling to the OH proton are seen.\\
occurring), are intrinsically interesting as solids and might not be open to study by X-ray diffraction: in these cases, solidstate NMR provides a useful alternative way of probing both structure and dynamics.\\
There are three principal contributions to the linewidths of solids. One is the direct magnetic dipolar interaction between nuclear spins. As pointed out in the discussion of spin-spin coupling, a nuclear magnetic moment gives rise to a local magnetic field which points in different directions at different locations around the nucleus. If the only component of interest is parallel to the direction of the applied magnetic field (because only this component has a significant effect), then provided certain subtle effects arising from transformation from the static to the rotating frame are neglected, the classical expression in The chemist's toolkit 27 can be used to write the magnitude of the local magnetic field as


\begin{equation*}
\mathcal{B}_{\mathrm{loc}}=\frac{\gamma_{\mathrm{N}} \hbar \mu_{0} m_{I}}{4 \pi R^{3}}\left(1-3 \cos ^{2} \theta\right) \tag{12B.17a}
\end{equation*}


Unlike in solution, in a solid this field is not averaged to zero by the molecular motion. Many nuclei may contribute to the total local field experienced by a nucleus of interest, and different nuclei in a sample may experience a wide range of fields. Typical dipole fields are of the order of 1 mT , which corresponds to splittings and linewidths of the order of 10 kHz for ${ }^{1} \mathrm{H}$. When the angle $\theta$ can vary only between 0 and $\theta_{\max }$, the average value of $1-3 \cos ^{2} \theta$ can be shown to be $-\left(\cos ^{2} \theta_{\text {max }}+\right.$ $\cos \theta_{\max }$ ). This result, in conjunction with eqn 12B.17a, gives the average local field as


\begin{equation*}
{\mathcal{B}_{\text {loc,av }}}=\frac{\gamma_{\mathrm{N}} \hbar \mu_{0} m_{I}}{4 \pi R^{3}}\left(\cos ^{2} \theta_{\max }+\cos \theta_{\max }\right) \tag{12B.17b}
\end{equation*}


\subsection*{Brief illustration 12B.6}
When $\theta_{\text {max }}=30^{\circ}$ and $R=160 \mathrm{pm}$, the local field generated by a proton is

$$
\begin{aligned}
\mathcal{B}_{\text {loc,av }} & =\frac{\overbrace{\left(3.546 \ldots \times 10^{-32} \mathrm{Tm}^{3}\right)}^{\gamma_{N} \hbar \mu_{0}} \times \overbrace{\left(\frac{1}{2}\right)}^{m_{1}} \times \overbrace{(1.616)}^{\cos ^{2} \theta_{\max }}+\cos \theta_{\max }}{4 \pi \times \underbrace{\left.1.60 \times 10^{-10} \mathrm{~m}\right)^{3}}_{R^{3}}} \\
& =5.57 \times 10^{-4} \mathrm{~T}=0.557 \mathrm{mT}
\end{aligned}
$$

A second source of linewidth is the anisotropy of the chemical shift. Chemical shifts arise from the ability of the applied field to generate electron currents in molecules. In general, this ability depends on the orientation of the molecule relative to the applied field. In solution, when the molecule is tumbling rapidly, only the average value of the chemical shift is relevant. However, the anisotropy is not averaged to zero for stationary\\
molecules in a solid, and molecules in different orientations have resonances at different frequencies. The chemical shift anisotropy also varies with the angle $\theta$ between the applied field and the principal axis of the molecule as $1-3 \cos ^{2} \theta$.

The third contribution is the electric quadrupole interaction. Nuclei with $I>\frac{1}{2}$ have an 'electric quadrupole moment', a measure of the extent to which the distribution of charge over the nucleus is not uniform (for instance, the positive charge may be concentrated around the equator or at the poles). An electric quadrupole interacts with an electric field gradient, such as may arise from a non-spherical distribution of charge around the nucleus. This interaction also varies as $1-3 \cos ^{2} \theta$.

Fortunately, there are techniques available for reducing the linewidths of solid samples. One technique, magic-angle spinning (MAS), takes note of the $1-3 \cos ^{2} \theta$ dependence of the dipole-dipole interaction, the chemical shift anisotropy, and the electric quadrupole interaction. The 'magic angle' is the angle at which $1-3 \cos ^{2} \theta=0$, and corresponds to $54.74^{\circ}$. In the technique, the sample is spun at high speed around an axis at the magic angle to the applied field (Fig. 12B.23). All the dipolar interactions and the anisotropies average to the value they would have at the magic angle, but at that angle they are zero. In principle, MAS therefore removes completely the linebroadening due to dipole-dipole interactions and chemical\\
\includegraphics[max width=\textwidth, center]{2025_02_27_d4b95d9718f0b9e2e6b9g-539}

Figure 12B.23 In magic-angle spinning, the sample spins on an axis at $54.74^{\circ}$ (i.e. $\arccos 1 / 3^{1 / 2}$ ) to the applied magnetic field. Rapid motion at this angle averages dipole-dipole interactions and chemical shift anisotropies to zero.\\
shift anisotropy. The difficulty with MAS is that the spinning frequency must not be less than the width of the spectrum, which is of the order of kilohertz. However, gas-driven sample spinners that can be rotated at up to 50 kHz are now routinely available.

\subsection*{Checklist of concepts}
$\square$ 1. The chemical shift of a nucleus is the difference between its resonance frequency and that of a reference standard.2. The shielding constant is the sum of a local contribution, a neighbouring group contribution, and a solvent contribution.3. The local contribution is the sum of a diamagnetic contribution and a paramagnetic contribution.4. The neighbouring group contribution arises from the currents induced in nearby groups of atoms.5. The solvent contribution can arise from specific molecular interactions between the solute and the solvent.6. Fine structure is the splitting of resonances into individual lines by spin-spin coupling; these splittings give rise to multiplets.7. Spin-spin coupling is expressed in terms of the spinspin coupling constant $J$; coupling leads to the splitting of lines in the spectrum.8. The coupling constant decreases as the number of bonds separating two nuclei increases.\\
9. Spin-spin coupling can be explained in terms of the polarization mechanism and the Fermi contact interaction.10. If the shift difference between two nuclei is large compared to the coupling constant between the nuclei the spin system is said to be weakly coupled; if the shift difference is small compared to the coupling, the spin system is strongly coupled.11. Chemically equivalent nuclei have the same chemical shifts; the same is true of magnetically equivalent nuclei, but in addition the coupling constant to any other nucleus is the same for each of the equivalent nuclei.12. Coalescence of two NMR lines occurs when nuclei are exchanged rapidly between environments by either conformational or chemical process.\\
13. Magic-angle spinning (MAS) is a technique in which the NMR linewidths in a solid sample are reduced by spinning at an angle of $54.74^{\circ}$ to the applied magnetic field.

\subsection*{Checklist of equations}
\begin{center}
\begin{tabular}{|c|c|c|c|}
\hline
Property & Equation & Comment & Equation number \\
\hline
$\delta$-Scale of chemical shifts & $\delta=\left\{\left(v-v^{\circ}\right) / v^{\circ}\right\} \times 10^{6}$ & Definition & 12B.4a \\
\hline
Relation between chemical shift and shielding constant & $\delta \approx\left(\sigma^{\circ}-\sigma\right) \times 10^{6}$ &  & 12B. 6 \\
\hline
Local contribution to the shielding constant & $\sigma($ local $)=\sigma_{\text {d }}+\sigma_{\mathrm{p}}$ &  & 12B. 8 \\
\hline
Lamb formula & $\sigma_{\mathrm{d}}=\left(e^{2} \mu_{0} / 12 \pi m_{\mathrm{e}}\right)\langle 1 / r\rangle$ & Applies to atoms & 12B. 9 \\
\hline
Neighbouring group contribution to the shielding constant & $\sigma($ neighbour $) \propto\left(\chi_{\|}-\chi_{\perp}\right)\left\{\left(1-3 \cos ^{2} \Theta\right) / R^{3}\right\}$ &  & 12B.10b \\
\hline
Karplus equation & ${ }^{3} J_{\mathrm{HH}}=A+B \cos \phi+C \cos 2 \phi$ & $A, B$, and $C$ are empirical constants & 12B. 14 \\
\hline
Condition for coalescence of two NMR lines & $\tau=2^{1 / 2} / \pi \delta v$ & $\tau$ is the lifetime of the exchange process & 12B. 16 \\
\hline
\end{tabular}
\end{center}

\section*{TOPIC 12C Pulse techniques in NMR}
\subsection*{Why do you need to know this material?}
To appreciate the power and scope of modern nuclear magnetic resonance techniques you need to understand how radiofrequency pulses can be used to obtain spectra.

\subsection*{What is the key idea?}
Sequences of pulses of radiofrequency radiation manipulate nuclear spins, leading to efficient acquisition of NMR spectra and the measurement of relaxation times.

\subsection*{What do you need to know already?}
You need to be familiar with the general principles of magnetic resonance (Topics 12A and 12B), and the vector model of angular momentum (Topic 7F). The development makes use of the concept of precession at the Larmor frequency (Topic 12A).

In modern forms of NMR spectroscopy the nuclear spins are first excited by a short, intense burst of radiofrequency radiation (a 'pulse'), applied at or close to the Larmor frequency. As a result of the excitation caused by the pulse, the spins emit radiation as they return to equilibrium. This timedependent signal is recorded and its 'Fourier transform' computed (as will be described) to give the spectrum. The technique is known as Fourier-transform NMR (FT-NMR). An analogy for the difference between conventional spectroscopy and pulsed NMR is the detection of the frequencies at which a bell vibrates. The 'conventional' option is to connect an audio oscillator to a loudspeaker and direct the sound towards the bell. The frequency of the sound source is then scanned until the bell starts to ring in resonance. The 'pulse' analogy is to strike the bell with a hammer and then Fourier transform the signal to identify the resonance frequencies of the bell.

One advantage of FT-NMR over conventional NMR is that it improves the sensitivity. However, the real power of the technique comes from the possibility of manipulating the nuclear spins by applying a sequence of several pulses. In this way it is possible to record spectra in which particular features are emphasized, or from which other properties of the molecule can be determined.

\subsection*{12C. 1 The magnetization vector}
To understand the pulse procedure, consider a sample composed of many identical spin- $\frac{1}{2}$ nuclei. According to the vector model of angular momentum (Topic 7F), a nuclear spin can be represented by a vector of length $\{I(I+1)\}^{1 / 2}$ with a component of length $m_{I}$ along the $z$-axis. As the three components of the angular momentum are complementary variables, the $x$ - and $y$-components cannot be specified if the $z$-component is known, so the vector lies anywhere on a cone around the $z$-axis. For $I=\frac{1}{2}$, the length of the vector is $3^{1 / 2} / 2$ and when $m_{I}=+\frac{1}{2}$ it makes an angle of $\arccos \left\{\frac{1}{2} /\left(3^{1 / 2} / 2\right)\right\}=54.7^{\circ}$ to the $z$-axis (Fig. 12C.1); when $m_{I}=-\frac{1}{2}$ the cone makes the same angle to the $-z$-axis.

In the absence of a magnetic field, the sample consists of equal numbers of $\alpha$ and $\beta$ nuclear spins with their vectors lying at random, stationary positions on their cones. The magnetization, $M$, of the sample, its net nuclear magnetic moment, is zero (Fig. 12C.2a). There are two changes when a magnetic field of magnitude $\mathscr{B}_{0}$ is applied along the $z$-direction:

\begin{itemize}
  \item The energies of the two spin states change, the $\alpha$ spins moving to a lower energy and the $\beta$ spins to a higher energy (provided $\gamma_{\mathrm{N}}>0$ ).
\end{itemize}

In the vector model, the two vectors are pictured as precessing at the Larmor frequency (Topic $12 \mathrm{~A}, v_{\mathrm{L}}=\gamma_{\mathrm{N}} \mathcal{B}_{0} / 2 \pi$ ). At 10 T , the Larmor frequency for ${ }^{1} \mathrm{H}$ nuclei (commonly referred to as 'protons') is 427 MHz . As the strength of the field is increased, the Larmor frequency increases and the precession becomes faster.

\begin{itemize}
  \item The populations of the two spin states (the numbers of $\alpha$ and $\beta$ spins) at thermal equilibrium change, with slightly more $\alpha$ spins than $\beta$ spins (Topic 12A).\\
\includegraphics[max width=\textwidth, center]{2025_02_27_d4b95d9718f0b9e2e6b9g-541}
\end{itemize}

Figure 12C. 1 The vector model of angular momentum for a single spin $-\frac{1}{2}$ nucleus with $m_{l}=+\frac{1}{2}$. The position of the vector on the cone is indeterminate.\\
\includegraphics[max width=\textwidth, center]{2025_02_27_d4b95d9718f0b9e2e6b9g-542(1)}

Figure 12C. 3 (a) In a pulsed NMR experiment the net magnetization is rotated away from the $z$-axis by applying radiofrequency radiation with its magnetic component $\mathscr{B}_{1}$ rotating in the $x y$-plane at the Larmor frequency. (b) When viewed in a frame also rotating about $z$ at the Larmor frequency, $\mathscr{B}_{1}$ appears to be stationary. The magnetization rotates around the $\mathscr{B}_{1}$ field, thus moving the magnetization away from the $z$-axis and generating transverse components.\\
at the bottom of the $\alpha^{\prime}$ and $\beta^{\prime}$ cones, but precess similarly. At thermal equilibrium there are fewer $\beta$ spins than $\alpha$ spins, so the net effect is a magnetization vector initially along the $z$ axis that rotates around the $\mathscr{B}_{1}$ direction at the $\mathscr{B}_{1}$ Larmor frequency and into the $x y$-plane.

When the radiofrequency field is applied in a pulse of duration $\Delta \tau$ the magnetization rotates through an angle (in radians) of $\phi=\Delta \tau \times\left(\gamma_{\mathrm{N}} \mathcal{B}_{1} / 2 \pi\right) \times 2 \pi$; this angle is known as the flip angle of the pulse. Therefore, to achieve a flip angle $\phi$, the duration of the pulse must be $\Delta \tau=\phi / \gamma_{\mathrm{N}} \mathcal{B}_{1}$. A $90^{\circ}$ pulse (with $90^{\circ}$\\
\includegraphics[max width=\textwidth, center]{2025_02_27_d4b95d9718f0b9e2e6b9g-542}

Figure 12C.4 When attention switches to the rotating frame, the vectors representing the spins are in states referring to the axis defined by $\mathcal{B}_{1}$, and precess on their cones. A uniform distribution in the $\mathscr{B}_{0}$ frame (blue) is actually a superposition of $\alpha^{\prime}$ and $\beta^{\prime}$ states that seem to be bunched together on their cones (pink). As the latter precess on the pink cones, the magnetization vector rotates into the $x y$-plane. Vectors representing $\beta$ spins in the original frame behave similarly.\\
(a)\\
\includegraphics[max width=\textwidth, center]{2025_02_27_d4b95d9718f0b9e2e6b9g-543}\\
\includegraphics[max width=\textwidth, center]{2025_02_27_d4b95d9718f0b9e2e6b9g-543(2)}\\
(b)

Figure 12C. 5 (a) If the radiofrequency field is applied for the appropriate time, the magnetization vector is rotated into the $x y$-plane; this is termed a $90^{\circ}$ pulse. (b) Once the magnetization is in the transverse plane it rotates about the $\mathscr{B}_{0}$ field at the Larmor frequency (when observed in a static frame). The magnetization vector periodically rotates past a small coil, inducing in it an oscillating current, which is the detected signal.\\
corresponding to $\phi=\pi / 2$ radians) of duration $\Delta \tau_{90}=\pi / 2 \gamma_{\mathrm{N}} \mathcal{B}_{1}$ ) rotates the magnetization from the $z$-axis into the $x y$-plane (Fig. 12C.5a).

\subsection*{Brief illustration 12C. 1}
The duration of a radiofrequency pulse depends on the strength of the $\mathscr{B}_{1}$ field. If a $90^{\circ}$ pulse requires $10 \mu \mathrm{~s}$, then for protons

$$
\mathscr{B}_{1}=\frac{\pi}{2 \times \underbrace{\left(2.675 \times 10^{8} \mathrm{~T}^{-1} \mathrm{~s}^{-1}\right)}_{\gamma_{N}} \times \underbrace{\left(1.0 \times 10^{-5} \mathrm{~s}\right)}_{\Delta \tau}}=5.9 \times 10^{-4} \mathrm{~T}
$$

or 0.59 mT .

Immediately after a $90^{\circ}$ pulse the magnetization lies in the $x y$-plane. Next, imagine stepping out of the rotating frame. The magnetization vector is now rotating in the $x y$ plane at the $\mathscr{B}_{0}$ Larmor frequency (Fig. 12C.5b). In an NMR spectrometer a small coil is wrapped around the sample and perpendicular to the $\mathscr{B}_{0}$ field in such a way that the precessing magnetization vector periodically 'cuts' the coil, thereby inducing in it a small current oscillating at the $\mathcal{B}_{0}$ Larmor frequency. This oscillating current is detected by a radiofrequency receiver.

As time passes the magnetization returns to an equilibrium state in which it has no transverse components; as it does so the oscillating signal induced in the coil decays to zero. This decay is exponential with a time constant denoted $T_{2}$. The overall form of the signal is therefore an oscillating-decaying free-induction decay (FID) like that shown in Fig. 12C. 6 and of the form

$$
S(t)=S_{0} \cos \left(2 \pi v_{\mathrm{L}} t\right) \mathrm{e}^{-t / T_{2}} \quad \text { Free-induction decay }
$$

(12C.1)\\
\includegraphics[max width=\textwidth, center]{2025_02_27_d4b95d9718f0b9e2e6b9g-543(1)}

Figure 12C.6 A freeinduction decay from of a sample of spins with a single resonance frequency.

So far it has been assumed that the radiofrequency radiation is exactly at the $\mathscr{B}_{0}$ Larmor frequency. However, virtually the same effect is obtained if the separation of the radiofrequency from the Larmor frequency is small compared to the inverse of the duration of the $90^{\circ}$ pulse. In practice, a spectrum with several peaks, each with a slightly different Larmor frequency, can be excited by selecting a radiofrequency somewhere in the centre of the spectrum and then making sure that the $90^{\circ}$ pulse is sufficiently short (which entails using an intense radiofrequency field so that $\mathscr{B}_{1}$ is still large enough to achieve rotation through $90^{\circ}$ ).

\subsection*{(b) Time- and frequency-domain signals}
Each line in an NMR spectrum can be thought of as arising from its own magnetization vector. Once that vector has been rotated into the $x y$-plane it precesses at the frequency of the corresponding line. Each vector therefore contributes a decay-ing-oscillating term to the observed signal and the FID is the sum of many such contributions. If there is only one line it is possible to determine its frequency simply by inspecting the FID, but that is rarely possible when the signal is composite. In such cases, Fourier transformation (The chemist's toolkit 28), as mentioned in the introduction, is used to analyse the signal.

The input to the Fourier transform is the oscillating-decaying 'time-domain' function $S(t)$. The output is the absorption spectrum, the 'frequency-domain' function, $I(v)$, which is obtained by computing the integral


\begin{equation*}
I(v)=\int_{0}^{\infty} S(t) \cos (2 \pi v t) \mathrm{d} t \tag{12C.2}
\end{equation*}


where $I(v)$ is the intensity at the frequency $v$. The complete fre-quency-domain function, which is the spectrum, is built up by evaluating this integral over a range of frequencies.

\subsection*{The chemist's toolkit 28}
\subsection*{The Fourier transform}
A Fourier transform expresses any waveform as a superposition of harmonic (sine and cosine) waves. If the waveform is the real function $S(t)$, then the contribution $I(v)$ of the oscillating function $\cos (2 \pi v t)$ is given by the 'cosine transform'


\begin{equation*}
I(v)=\int_{0}^{\infty} S(t) \cos (2 \pi v t) \mathrm{d} t \tag{1}
\end{equation*}


There is an analogous transform appropriate for complex functions: see the additional information for this Toolkit available on the website. If the signal varies slowly, then the greatest contribution comes from low-frequency waves; rapidly changing features in the signal are reproduced by high-frequency contributions. If the signal is a simple exponential decay of the form $S(t)=S_{0} \mathrm{e}^{-t / \tau}$, the contribution of the wave of frequency $v$ is


\begin{equation*}
I(v)=S_{0} \int_{0}^{\infty} \mathrm{e}^{-t / \tau} \cos (2 \pi v t) \mathrm{d} t=\frac{S_{0} \tau}{1+(2 \pi v \tau)^{2}} \tag{2}
\end{equation*}


Sketch 1 shows a fast and slow decay and the corresponding frequency contributions: note that a slow decay has predominantly low-frequency contributions and a fast decay has contributions at higher frequencies.\\
If an experimental procedure results in the function $I(v)$ itself, then the corresponding signal can be reconstructed by forming the inverse Fourier transform:


\begin{equation*}
S(t)=\frac{2}{\pi} \int_{0}^{\infty} I(v) \cos (2 \pi v t) \mathrm{d} v \tag{3}
\end{equation*}


Fourier transforms are applicable to spatial functions too. Their interpretation is similar but it is more appropriate to think in terms of the wavelengths of the contributing waves. Thus, if the function varies only slowly with distance, then its Fourier transform has mainly long-wavelength contributions. If the features vary quickly with distance (as in the electron density in a crystal), then short-wavelength contributions feature.\\
\includegraphics[max width=\textwidth, center]{2025_02_27_d4b95d9718f0b9e2e6b9g-544(1)}

\subsection*{Sketch 1}
If the time-domain contains a single oscillating-decaying term, as in eqn 12C.1, the Fourier transform (see the extended Chemist's toolkit 28 on the website) is


\begin{equation*}
I(v)=\frac{S_{0} T_{2}}{1+\left(v_{\mathrm{L}}-v\right)^{2}\left(2 \pi T_{2}\right)^{2}} \tag{12C.3}
\end{equation*}


The graph of this expression has a so-called 'Lorentzian' shape, a symmetrical peak centred at $v=v_{\mathrm{L}}$ and height $S_{0} T_{2}$; its width at half the peak height is $1 / \pi T_{2}$ (Fig. 12C.7). If the FID consists of a sum of decaying-oscillating functions, Fourier transformation gives a spectrum consisting of a series of peaks at the various frequencies.

In practice, the FID is sampled digitally and the integral of eqn 12C. 2 is evaluated numerically by a computer in the NMR spectrometer. Figure 12C. 8 shows the time- and frequency-domain functions for three different FIDs of increasing complexity.\\
\includegraphics[max width=\textwidth, center]{2025_02_27_d4b95d9718f0b9e2e6b9g-544}

Figure 12C. 7 A Lorentzian absorption line. The width at halfheight depends of the time constant $T_{2}$ which characterizes the decay of the time-domain signal.\\
\includegraphics[max width=\textwidth, center]{2025_02_27_d4b95d9718f0b9e2e6b9g-545}

Figure 12C. 8 Free induction decays (the time domain) and the corresponding spectra (the frequency domain) obtained by Fourier transformation. (a) An uncoupled A resonance, (b) the A resonance of an $A X$ system, (c) the $A$ and $X$ resonances of an $A_{2} X_{3}$ system.

\subsection*{12C.2 Spin relaxation}
Relaxation is the process by which the magnetization returns to its equilibrium value, at which point it is entirely along the $z$-axis, with no $x$ - and no $y$-component (no transverse components). In terms of the behaviour of individual spins, the approach to equilibrium involves transitions between the two spin states in order to establish the thermal equilibrium populations of $\alpha$ and $\beta$. The attainment of equilibrium also requires the magnetic moments of individual nuclei becoming distributed at random angles on their two cones.

As already explained, after a $90^{\circ}$ pulse the magnetization vector lies in the $x y$-plane. This orientation implies that, from the viewpoint of the laboratory quantization axis, there are now equal numbers of $\alpha$ and $\beta$ spins because otherwise there would be a component of the magnetization in the $z$-direction. At thermal equilibrium, however, there is a Boltzmann distribution of spins, with more $\alpha$ spins than $\beta$ spins (provided $\gamma_{\mathrm{N}}>0$ ) and a non-zero $z$-component of magnetization. The return of the $z$-component of magnetization to its equilibrium value is termed longitudinal relaxation. It is usually assumed that this process follows an exponential recovery curve, with a time constant called the longitudinal relaxation time, $T_{1}$. Because longitudinal relaxation involves the transfer of energy between the spins and the surroundings (the 'lattice'), the time constant $T_{1}$ is also called the spin-lattice relaxation time.

If the $z$-component of magnetization at time $t$ is $M_{z}(t)$, then the recovery to the equilibrium magnetization $M_{0}$ takes the form

$$
M_{z}(t)-M_{0} \propto \mathrm{e}^{-t / T_{1}} \quad \begin{aligned}
& \text { Longitudinal relaxation time } \\
& {[\text { definition] }}
\end{aligned}
$$

(12C.4)

Immediately after a $90^{\circ}$ pulse the fact that the magnetization vector lies in the $x y$-plane implies that the $\alpha$ and $\beta$ spins are aligned in a particular way so as to give a net (and rotating) component of magnetization in the $x y$-plane. At thermal equilibrium, however, the spins are at random angles on their cones and there is no transverse component of magnetization. The return of the transverse magnetization to its equilibrium value of zero is termed transverse relaxation. It is usually assumed that this process is an exponential decay with a time constant called the transverse relaxation time, $T_{2}$ (or spinspin relaxation time). If the transverse magnetization at time $t$ is $M_{x y}(t)$, then the decay takes the form

$$
M_{x y}(t) \propto \mathrm{e}^{-t / T_{2}} \quad \begin{aligned}
& \text { Transverse relaxation time } \\
& \text { [definition] }
\end{aligned}
$$

(12C.5)

This $T_{2}$ is the same as that describing the decay of the free induction signal which is proportional to $M_{x y}(t)$.

\subsection*{(a) The mechanism of relaxation}
The return of the $z$-component of magnetization to its equilibrium value involves transitions between $\alpha$ and $\beta$ spin states so as to achieve the populations required by the Boltzmann distribution. These transitions are caused by local magnetic fields that fluctuate at a frequency close to the resonance frequency of the $\beta \leftrightarrow \alpha$ transition. The local fields can have a variety of origins, but commonly arise from nearby magnetic nuclei or unpaired electrons. These fields fluctuate due to the tumbling motion of molecules in a fluid sample. If molecular tumbling is too slow or too fast compared with the resonance frequency, it gives rise to a fluctuating magnetic field with a frequency that is either too low or too high to stimulate a transition between $\beta$ and $\alpha$, so $T_{1}$ is long. Only if the molecule tumbles at about the resonance frequency is the fluctuating magnetic field able to induce spin flips effectively, and then $T_{1}$ is short.

The rate of molecular tumbling increases with temperature and with reducing viscosity of the solvent, so a dependence of the relaxation time like that shown in Fig. 12C. 9 can be expected. The quantitative treatment of relaxation times depends on setting up models of molecular motion and using, for instance, the diffusion equation (Topic 16C) adapted for rotational motion.

Transverse relaxation is the result of the individual magnetic moments of the spins losing their relative alignment as they spread out on their cones. One contribution to this randomization is any process that involves a transition between the\\
\includegraphics[max width=\textwidth, center]{2025_02_27_d4b95d9718f0b9e2e6b9g-546(1)}

Figure 12C.9 The variation of the two relaxation times with the rate at which the molecules tumble in solution. The horizontal axis can be interpreted as representing temperature or viscosity. Note that the two relaxation times coincide when the motion is fast.\\
two spin states. That is, any process that causes longitudinal relaxation also contributes to transverse relaxation. Another contribution is the variation in the local magnetic fields experienced by the nuclei. When the fluctuations in these fields are slow, each molecule lingers in its local magnetic environment, and because the precessional rates depend on the strength of the field, the spin orientations randomize quickly around their cones. In other words, slow molecular motion corresponds to short $T_{2}$. If the molecules move rapidly from one magnetic environment to another, the effects of differences in local magnetic field average to zero: individual spins do not precess at very different rates, remain bunched for longer, and transverse relaxation does not take place as quickly. This fast motion corresponds to long $T_{2}$ (as shown in Fig. 12C.9). Calculations show that, when the motion is fast, transverse and longitudinal relaxation have similar time constants.

\subsection*{Brief illustration 12C. 2}
For a small molecule dissolved in a non-viscous solvent the value of $T_{2}$ for ${ }^{1} \mathrm{H}$ can be as long as several seconds. The width (measured at half the peak height) of the corresponding line in the spectrum is $1 / \pi T_{2}$. For $T_{2}=3.0 \mathrm{~s}$ the width is

$$
\Delta v_{1 / 2}=\frac{1}{\pi T_{2}}=\frac{1}{\pi \times(3.0 \mathrm{~s})}=0.11 \mathrm{~Hz}
$$

In contrast, for a larger molecule such as a protein dissolved in water $T_{2}$ is much shorter, and a value of 30 ms is not unusual. The corresponding linewidth is 11 Hz .

So far, it has been assumed that the applied magnetic field is homogeneous (uniform) in the sense of having the same value across the sample so that the differences in Larmor frequencies arise solely from interactions within the sample. In practice, due to the limitations in the design of the magnet, the\\
field is not perfectly uniform and is different at different locations in the sample. The inhomogeneity results in a inhomogeneous broadening of the resonance. It is common to express the extent of inhomogeneous broadening in terms of an effective transverse relaxation time, $T_{2}^{*}$, by using a relation like eqn 12C.5, but writing

\[
T_{2}^{*}=\frac{1}{\pi \Delta v_{1 / 2}} \quad \begin{align*}
& \text { Effective transverse relaxation time }  \tag{12C.6}\\
& {[\text { definition] }}
\end{align*}
\]

where $\Delta \nu_{1 / 2}$ is the observed width at half-height of the line (which is assumed to be Lorentzian). In practice an inhomogeneously broadened line is unlikely to be Lorentzian, so the assumption that the decay is exponential and characterized by a time constant $T_{2}^{*}$ is an approximation. The observed linewidth has a contribution from both the inhomogeneous broadening and the transverse relaxation. The latter is usually referred to as homogeneous broadening. Which contributions dominate depends on the molecular system being studied and the quality of the magnet.

\subsection*{(b) The measurement of $T_{1}$ and $T_{2}$}
The longitudinal relaxation time $T_{1}$ can be measured by the inversion recovery technique. The first step is to apply a $180^{\circ}$ pulse to the sample by applying the $\mathcal{B}_{1}$ field for twice as long as for a $90^{\circ}$ pulse. As a result of the pulse, the magnetization vector is rotated into the $-z$-direction (Fig. 12C.10a). The effect of the pulse is to invert the population of the two levels and to result in more $\beta$ spins than $\alpha$ spins.\\
Immediately after the $180^{\circ}$ pulse no signal can be detected because the magnetization has no transverse component. The $\beta$ spins immediately begin to relax back into $\alpha$ spins, and the magnetization vector first shrinks towards zero and then increases in the opposite direction until it reaches its thermal equilibrium value. Before that has happened, after an interval\\
\includegraphics[max width=\textwidth, center]{2025_02_27_d4b95d9718f0b9e2e6b9g-546}

Figure 12C. 10 (a) The result of applying a $180^{\circ}$ pulse to the magnetization in the rotating frame, and the effect of a subsequent $90^{\circ}$ pulse. (b) The amplitude of the frequency-domain spectrum varies with the interval between the two pulses because there has been time for longitudinal relaxation to occur.\\
$\tau$, a $90^{\circ}$ pulse is applied. That pulse rotates the remaining $z$ component of magnetization into the $x y$-plane, where it generates an FID signal. The frequency-domain spectrum is then obtained by Fourier transformation in the usual way.

The intensity of the resulting spectrum depends on the magnitude of the magnetization vector that has been rotated into the $x y$-plane. That magnitude changes exponentially with a time constant $T_{1}$ as the interval $\tau$ is increased, so the intensity of the spectrum also changes exponentially with increasing $\tau$. The longitudinal relaxation time can therefore be measured by fitting an exponential curve to the series of spectra obtained with different values of $\tau$.

The measurement of $T_{2}$ (as distinct from $T_{2}^{*}$ ) depends on being able to eliminate the effects of inhomogeneous broadening. The cunning required is at the root of some of the most important advances made in NMR since its introduction.

A spin echo is the magnetic analogue of an audible echo. The sequence of events is shown in Fig. 12C.11. The overall magnetization can be regarded as made up of a number of different magnetizations, each of which arises from a spin packet of nuclei with very similar precession frequencies. The spread in these frequencies arises from the inhomogeneity of $\mathscr{B}_{0}$ (which is responsible for inhomogeneous broadening), so different parts of the sample experience different fields. The precession frequencies also differ if there is more than one chemical shift present.

First, a $90^{\circ}$ pulse is applied to the sample. The subsequent events are best followed in a rotating frame in which $\mathscr{B}_{1}$ is stationary along the $x$-axis and causes the magnetization to rotate on to the $y$-axis of the $x y$-plane. Immediately after the\\
\includegraphics[max width=\textwidth, center]{2025_02_27_d4b95d9718f0b9e2e6b9g-547}

Figure 12C. 11 The action of the spin echo pulse sequence $90^{\circ}-\tau-180^{\circ}-\tau$, viewed in a rotating frame at the Larmor frequency. Note that the $90^{\circ}$ pulse is applied about the $x$-axis, but the $180^{\circ}$ pulse is applied about the $y$-axis. 'Slow' and 'Fast' refer to the speed of the spin packet relative to the rotating frame frequency.\\
pulse, the spin packets begin to fan out because they have different Larmor frequencies. In Fig. 12C. 11 the magnetization vectors of two representative packets are shown and are described as 'fast' and 'slow', indicating their frequency relative to the rotating frame frequency (the nominal Larmor frequency). Because the rotating frame is at the Larmor frequency, the 'fast' and 'slow' vectors rotate in opposite senses when viewed in this frame.

First, suppose that there is no transverse relaxation but that the field is inhomogeneous. After an evolution period $\tau$, a $180^{\circ}$ pulse is applied along the $y$-axis of the rotating frame. The pulse rotates the magnetization vectors around that axis into mirror-image positions with respect to the $y z$-plane. Once there, the packets continue to move in the same direction as they did before, and so migrate back towards the $y$-axis. After an interval $\tau$ all the packets are again aligned along the axis. The resultant signal grows in magnitude, reaching a maximum, the 'spin echo', at the end of the second period $\tau$. The fanning out caused by the field inhomogeneity is said to have been 'refocused'.

The important feature of the technique is that the size of the echo is independent of any local fields that remain constant during the two $\tau$ intervals. If a spin packet is 'fast' because it happens to be composed of spins in a region of the sample that experience a higher than average field, then it remains fast throughout both intervals, and so the angle through which it rotates is the same in the two intervals. Hence, the size of the echo is independent of inhomogeneities in the magnetic field, because these remain constant.

Now consider the consequences of transverse relaxation. This relaxation arises from fields that vary on a molecular scale, and there is no guarantee that an individual 'fast' spin will remain 'fast' in the refocusing phase: the spins within the packets therefore spread with a time constant $T_{2}$. Consequently, the effects of the relaxation are not refocused, and the size of the echo decays with the time constant $T_{2}$. The intensity of the signal after a spin echo is measured for a series of values of the delay $\tau$, and the resulting data analysed to determine $T_{2}$.

\subsection*{12C. 3 Spin decoupling}
Carbon-13 has a natural abundance of only 1.1 per cent and is therefore described as a dilute-spin species. The probability that any one molecule contains more than one ${ }^{13} \mathrm{C}$ nucleus is rather low, and so the possibility of observing the effects of ${ }^{13} \mathrm{C}-{ }^{13} \mathrm{C}$ spin-spin coupling can be ignored. In contrast the abundance of the isotope ${ }^{1} \mathrm{H}$ is very close to 100 per cent and is therefore described as an abundant-spin species. All the hydrogen nuclei in a molecule can be assumed to be ${ }^{1} \mathrm{H}$ and the effects of coupling between them is observed.

Dilute-spin species are observed in NMR spectroscopy and show coupling to abundant-spin species present in the molecule. Generally speaking, ${ }^{13} \mathrm{C}$-NMR spectra are very complex on account of the spin couplings of each ${ }^{13} \mathrm{C}$ nucleus with the numerous ${ }^{1} \mathrm{H}$ nuclei in the molecule. However, a dramatic simplification of the spectrum can be obtained by the use of proton decoupling. In this technique radiofrequency radiation is applied at (or close to) the ${ }^{1} \mathrm{H}$ Larmor frequency while the ${ }^{13} \mathrm{C}$ FID is being observed. This stimulation of the ${ }^{1} \mathrm{H}$ nuclei causes their spins state to change rapidly, so averaging the ${ }^{1} \mathrm{H}-{ }^{13} \mathrm{C}$ couplings to zero. As a result, each ${ }^{13} \mathrm{C}$ nucleus gives a single line rather than a complex multiplet. Not only is the spectrum simplified, but the sensitivity is also improved as all the intensity is concentrated into a single line.

\subsection*{12c. 4 The nuclear Overhauser effect}
A common source of the local magnetic fields that are responsible for relaxation is the dipole-dipole interaction between two magnetic nuclei (see The chemist's toolkit 27 in Topic 12B). In this interaction the magnetic dipole of the first spin generates a magnetic field that interacts with the magnetic dipole of the second spin. The strength of the interaction is proportional to $1 / R^{3}$, where $R$ is the distance between the two spins, and is also proportional to the product of the magnetogyric ratios of the spins. As a result, the interaction is characterized as being short-range and significant for nuclei with high magnetogyric ratios. In typical organic and biological molecules, which have many ${ }^{1} \mathrm{H}$ nuclei, the local fields due to the dipole-dipole interaction are likely to be the dominant source of relaxation.

The nuclear Overhauser effect (NOE) makes use of the relaxation caused by the dipole-dipole interaction of nuclear spins. In this effect, irradiation of one spin leads to a change in the intensity of the resonance from a second spin provided the two spins are involved in mutual dipole-dipole relaxation. Because the dipole-dipole interaction has only a short range, the observation of an NOE is indicative of the closeness of the two nuclei involved and can be interpreted in terms of the structure of the molecule.

To understand the effect, consider the populations of the four levels of a homonuclear AX spin system shown in Fig. 12C.12. At thermal equilibrium, the population of the $\alpha_{\mathrm{A}} \alpha_{\mathrm{x}}$ level is the greatest, and that of the $\beta_{\mathrm{A}} \beta_{\mathrm{x}}$ level is the least; the other two levels have the same energy and an intermediate population. For the purposes of this discussion it is sufficient to consider the deviations of the populations from the average value for all four levels: the $\alpha_{\mathrm{A}} \alpha_{\mathrm{X}}$ level has a greater population than the average, the $\beta_{\mathrm{A}} \beta_{\mathrm{x}}$ level has a smaller population, and the other two levels have a population equal to the average. The deviations from the average are $\Delta N$ for $\alpha_{A} \alpha_{\mathrm{x}},-\Delta N$ for\\
\includegraphics[max width=\textwidth, center]{2025_02_27_d4b95d9718f0b9e2e6b9g-548}

Figure 12C. 12 The energy levels of an $A X$ system and an indication of their relative populations. Each green square above the line represents an excess population above the average, and each white square below the line represents a lower population than the average. The symbols in red show the deviations in population from their average value.\\
$\beta_{\mathrm{A}} \beta_{\mathrm{x}}$, and 0 for the other two levels. These population differences are represented in Fig. 12C.12. The intensity of a transition reflects the difference in the population of the two energy levels involved, and for all four transitions this population difference (lower - upper) is $\Delta N$, implying that all four transitions have the same intensity.

The NOE experiment involves irradiating the two X spin transitions ( $\alpha_{\mathrm{A}} \alpha_{\mathrm{x}} \leftrightarrow \alpha_{\mathrm{A}} \beta_{\mathrm{x}}$ and $\beta_{\mathrm{A}} \alpha_{\mathrm{X}} \leftrightarrow \beta_{\mathrm{A}} \beta_{\mathrm{X}}$ ) with a radiofrequency field, but making sure that the field is sufficiently weak that the two A spin transitions are not affected. When applied for a long time this field saturates the X spin transitions in the sense that the populations of the two energy levels are equalized. In the case of the $\alpha_{\mathrm{A}} \alpha_{\mathrm{x}} \leftrightarrow \alpha_{\mathrm{A}} \beta_{\mathrm{x}}$ transition the populations of the two levels become $\frac{1}{2} \Delta N$, and for the other X spin transition $-\frac{1}{2} \Delta N$, as shown in Fig. 12C.13a. These changes in population do not affect the population differences across the A spin transitions ( $\alpha_{\mathrm{A}} \alpha_{\mathrm{x}} \leftrightarrow \beta_{\mathrm{A}} \alpha_{\mathrm{x}}$ and $\alpha_{\mathrm{A}} \beta_{\mathrm{x}} \leftrightarrow \beta_{\mathrm{A}} \beta_{\mathrm{x}}$ ) which remain at $\Delta N$, therefore the intensity of the A spin transitions is unaffected.\\
Now consider the effect of spin relaxation. One source of relaxation is dipole-dipole interaction between the two spins. The hamiltonian for this interaction is proportional to the spin operators of the two nuclei and contains terms that can flip both spins simultaneously and convert $\alpha_{\mathrm{A}} \alpha_{\mathrm{x}}$ into $\beta_{\mathrm{A}} \beta_{\mathrm{x}}$. This double-flipping process tends to restore the populations of these two levels to their equilibrium values of $\Delta N$ and $-\Delta N$, respectively, as shown in Fig. 12C.13b. A consequence is that the population difference across each of the A spin transitions is now $\frac{3}{2} \Delta N$, which is greater than it is at equilibrium. In summary, the combination of saturating the X spin transitions and dipole-dipole relaxation leads to an enhancement of the intensity of the A spin transitions.

The hamiltonian for the dipole-dipole interaction also contains combinations of spin operators that flip the spins in op-\\
\includegraphics[max width=\textwidth, center]{2025_02_27_d4b95d9718f0b9e2e6b9g-549}

Figure 12C. 13 (a) When an $X$ transition is saturated, the populations of the two states are equalized, leading to the populations shown (using the same symbols as in Fig. 12C.12). (b) Dipole-dipole relaxation can cause transitions between the $\alpha \alpha$ and $\beta \beta$ states, such that they return to their original populations: the result is that the population difference across the A transitions is increased. (c) Dipole-dipole relaxation can also cause the populations of the $\alpha \beta$ and $\beta \alpha$ states to return to their equilibrium values: this decreases the population difference across the A transitions.\\
posite directions, so taking the system from $\alpha_{\mathrm{A}} \beta_{\mathrm{x}}$ to $\beta_{\mathrm{A}} \alpha_{\mathrm{x}}$. This process drives the populations of these states to their equilibrium values, as shown in Fig. 12C.13c. As before, the population differences across the A spins transitions are affected, but now they are reduced to $+\frac{1}{2} \Delta N$, meaning that the A spin transitions are less intense than they would have been in the absence of dipole-dipole relaxation.

It should be clear that there are two opposing effects: the relaxation induced transitions between $\alpha_{\mathrm{A}} \alpha_{\mathrm{x}}$ and $\beta_{\mathrm{A}} \beta_{\mathrm{x}}$ which enhance the A spin transitions, and the transitions between $\alpha_{\mathrm{A}} \beta_{\mathrm{x}}$ and $\beta_{\mathrm{A}} \alpha_{\mathrm{x}}$ which reduce the intensity of these transitions. Which effect dominates depends on the relative rates of these two relaxation pathways. As in the discussion of relaxation times in Section 12C.2, the efficiency of the $\beta_{\mathrm{A}} \beta_{\mathrm{x}} \leftrightarrow$ $\alpha_{\mathrm{A}} \alpha_{\mathrm{x}}$ relaxation is high if the dipole field oscillates close to the transition frequency, which in this case is about $2 v_{\mathrm{L}}$; likewise, the efficiency of the $\alpha_{\mathrm{A}} \beta_{\mathrm{x}} \leftrightarrow \beta_{\mathrm{A}} \alpha_{\mathrm{x}}$ relaxation is high if the dipole field is stationary (in this case there is no frequency difference between the initial and final states). Small molecules tumble rapidly and have substantial motion at $2 v_{\mathrm{L}}$. As a result, the $\beta_{\mathrm{A}} \beta_{\mathrm{X}} \leftrightarrow \alpha_{\mathrm{A}} \alpha_{\mathrm{x}}$ pathway dominates and results in an increase in the intensity of the A spin transitions. This increase is called a 'positive NOE enhancement'. On the other hand, large molecules tumble more slowly so there is less motion at $2 \nu_{\mathrm{L}}$. In this case, the $\alpha_{\mathrm{A}} \beta_{\mathrm{x}} \leftrightarrow \beta_{\mathrm{A}} \alpha_{\mathrm{x}}$ pathway dominates and results in a decrease in the intensity of the A spin transitions. This decrease is called a 'negative NOE enhancement'.

The NOE enhancement is usually reported in terms of the parameter $\eta$ (eta), where

$$
\eta=\frac{I_{\mathrm{A}}-I_{\mathrm{A}}^{\circ}}{I_{\mathrm{A}}^{\circ}}
$$

NOE enhancement parameter (12C.7)

Here $I_{\mathrm{A}}^{\circ}$ is the intensity of the signals due to nucleus A before saturation, and $I_{\mathrm{A}}$ is the intensity after the X spins have been\\
saturated for long enough for the NOE to build up (typically several multiples of $T_{1}$ ). For a homonuclear system, and if the only source of relaxation is due to the dipole-dipole interaction, $\eta$ lies between -1 (a negative enhancement) for slow tumbling molecules and $+\frac{1}{2}$ (a positive enhancement) for fast tumbling molecules. In practice, other sources of relaxation are invariably present so these limiting values are rarely achieved.

The utility of the NOE in NMR spectroscopy comes about because only dipole-dipole relaxation can give rise to the effect: only this type of relaxation causes both spins to flip simultaneously. Thus, if nucleus X is saturated and a change in the intensity of the transitions from nucleus A is observed, then there must be a dipole-dipole interaction between the two nuclei. As that interaction is proportional to $1 / R^{3}$, where $R$ is the distance between the two spins, and the relaxation it causes is proportional to the square of the interaction, the NOE is proportional to $1 / R^{6}$. Therefore, for there to be significant dipole-dipole relaxation between two spins they must be close (for ${ }^{1} \mathrm{H}$ nuclei, not more than 0.5 nm apart), so the observation of an NOE is used as qualitative indication of the proximity of nuclei. In principle it is possible to make a quantitative estimate of the distance from the size of the NOE, but to do so requires the effects of other kinds of relaxation to be taken into account.

The value of the NOE enhancement $\eta$ also depends on the values of the magnetogyric ratios of A and X , because these properties affect both the populations of the levels and the relaxation rates. For a heteronuclear spin system the maximal enhancement is


\begin{equation*}
\eta=\frac{\gamma_{\mathrm{x}}}{2 \gamma_{\mathrm{A}}} \tag{12C.8}
\end{equation*}


where $\gamma_{\mathrm{A}}$ and $\gamma_{\mathrm{x}}$ are the magnetogyric ratios of nuclei A and X , respectively.

\subsection*{Brief illustration 12C. 3}
From eqn 12C. 8 and the data in Table 12A.2, the maximum NOE enhancement parameter for a ${ }^{13} \mathrm{C}-{ }^{1} \mathrm{H}$ pair is

$$
\eta=\frac{\overbrace{2.675 \times 10^{8} \mathrm{~T}^{-1} \mathrm{~s}^{-1}}^{\gamma_{1 \mathrm{H}}}}{2 \times \underbrace{\left(6.73 \times 10^{7} \mathrm{~T}^{-1} \mathrm{~s}^{-1}\right)}_{\gamma_{1 \mathrm{~B}_{\mathrm{C}}}}}=1.99
$$

\subsection*{Checklist of concepts}
1. The free-induction decay (FID) is the time-domain signal resulting from the precession of transverse magnetization.2. Fourier transformation of the FID (the time domain) gives the NMR spectrum (the frequency domain).3. Longitudinal (or spin-lattice) relaxation is the process by which the $z$-component of the magnetization returns to its equilibrium value.4. Transverse (or spin-spin) relaxation is the process by which the $x$ - and $y$-components of the magnetization return to their equilibrium values of zero.It is common to take advantage of this enhancement to improve the sensitivity of ${ }^{13} \mathrm{C}$ NMR spectra. Prior to recording the spectrum, the ${ }^{1} \mathrm{H}$ nuclei are irradiated so that they become saturated, leading to a build-up of the NOE enhancement on the ${ }^{13} \mathrm{C}$ nuclei.\\
5. The longitudinal relaxation time $T_{1}$ can be measured by the inversion recovery technique.6. The transverse relaxation time $T_{2}$ can be measured by observing spin echoes.7. In proton decoupling of ${ }^{13} \mathrm{C}$-NMR spectra, the protons are continuously irradiated; the effect is to collapse the splittings due to the ${ }^{13} \mathrm{C}-{ }^{1} \mathrm{H}$ couplings.8. The nuclear Overhauser effect is the modification of the intensity of one resonance by the saturation of another: it occurs only if the two spins are involved in mutual dipole-dipole relaxation.

\subsection*{Checklist of equations}
\begin{center}
\begin{tabular}{llll}
\hline
Property & Equation & Comment & Equation number \\
\hline
Free-induction decay & $S(t)=S_{0} \cos \left(2 \pi v_{\mathrm{L}} t\right) \mathrm{e}^{-t / T_{2}}$ & $T_{2}$ is the transverse relaxation time & 12 C .1 \\
Width at half-height of an NMR line & $\Delta v_{1 / 2}=1 / \pi T_{2}$ & Assumed Lorentzian & $T_{1}$ is the longitudinal relaxation time \\
Longitudinal relaxation & $M_{z}(t)-M_{0} \propto \mathrm{e}^{-t / T_{1}}$ &  & 12 C .4 \\
Transverse relaxation & $M_{x y}(t) \propto \mathrm{e}^{-t / T_{2}}$ & Definition & 12 C .5 \\
NOE enhancement parameter & $\eta=\left(I_{\mathrm{A}}-I_{\mathrm{A}}^{\circ}\right) / I_{\mathrm{A}}^{\circ}$ &  & 12 C .7 \\
\hline
\end{tabular}
\end{center}

\section*{TOPIC 12D Electron paramagnetic resonance }
\subsection*{Why do you need to know this material?}
Many materials and biological systems contain species bearing unpaired electrons and some chemical reactions generate intermediates with unpaired electrons. Electron paramagnetic resonance is a key spectroscopic tool for studying them.

\subsection*{What is the key idea?}
The details of an EPR spectrum give information on the distribution of the density of the unpaired electron.

\subsection*{> What do you need to know already?}
You need to be familiar with the concepts of electron spin (Topic 8 B ) and the general principles of magnetic resonance (Topic 12A). The discussion refers to spin-orbit coupling in atoms (Topic 8C) and the Fermi contact interaction in molecules (Topic 12B).

Electron paramagnetic resonance (EPR), which is also known as electron spin resonance (ESR), is used to study species that contain unpaired electrons. Both solids and liquids can be studied, but the study of gas-phase samples is complicated by the free rotation of the molecules.

\subsection*{12D. 1 The $g$-value}
According to the discussion in Topic 12 A , the resonance frequency for a transition between the $m_{s}=-\frac{1}{2}$ and the $m_{s}=+\frac{1}{2}$ levels of a free electron is

$$
\begin{array}{ll}
h \nu=g_{\mathrm{e}} \mu_{\mathrm{B}} \mathcal{B}_{0} & \begin{array}{l}
\text { Resonance condition } \\
\text { [free electron] }
\end{array}
\end{array}
$$

(12D.1)\\
where $g_{\mathrm{e}} \approx 2.0023$. If the electron is in a radical the field it experiences differs from the applied field due to the presence of local magnetic fields arising from electronic currents induced in the molecular framework. This difference is taken into account by replacing $g_{\mathrm{e}}$ by $g$ and expressing the resonance condition as

$$
h \nu=g \mu_{\mathrm{B}} \mathcal{B}_{0}
$$

EPR resonance condition\\
(12D.2)\\
where $g$ is the $g$-value of the radical.

Electron paramagnetic resonance spectra are usually recorded by keeping the frequency of the microwave radiation fixed and then varying the magnetic field so as to bring the electron into resonance with the microwave frequency. The positions of the peaks, and the horizontal scale on spectra, are therefore specified in terms of the magnetic field.

\subsection*{Brief illustration 12D. 1}
The centre of the EPR spectrum of the methyl radical occurs at 329.40 mT in a spectrometer operating at 9.2330 GHz (radiation belonging to the X band of the microwave region). Its $g$-value is therefore

$$
g=\overbrace{\frac{\left(6.62608 \times 10^{-34} \mathrm{Js}\right)}{h} \times \overbrace{\left(9.2330 \times 10^{9} \mathrm{~s}^{-1}\right)}^{(\underbrace{\left(9.2740 \times 10^{-24} \mathrm{JT}^{-1}\right)}_{\mu_{\mathrm{B}}}} \times \underbrace{(0.32940 \mathrm{~T})}_{\mathcal{B}_{0}}}^{v}=2.0027
$$

The $g$-value is related to the ease with which the applied field can generate currents through the molecular framework and the strength of the magnetic field these currents generate. Therefore, the $g$-value gives some information about electronic structure and plays a similar role in EPR to that played by shielding constants in NMR. A $g$-value smaller than $g_{\mathrm{e}}$ implies that in the molecule the electron experiences a magnetic field smaller than the applied field, whereas a value greater than $g_{\text {e }}$ implies that the magnetic field is greater. Both outcomes are possible, depending on the details of the electronic excited states.

Two factors are responsible for the difference of the $g$-value from $g_{\mathrm{e}}$. Electrons migrate through the molecular framework by making use of excited states (Fig. 12D.1). This circulation gives rise to a local magnetic field that can add to or subtract from the applied field. The extent to which these currents are induced is inversely proportional to the separation of energy levels, $\Delta E$, in the radical or complex. Secondly, the strength of the field experienced by the electron spin as a result of these orbital currents is proportional to the spin-orbit coupling constant, $\xi$ (Topic 8C). It follows that the difference of the $g$-value from $g_{\mathrm{e}}$ is proportional to $\xi / \Delta E$. This proportionality is widely\\
\includegraphics[max width=\textwidth, center]{2025_02_27_d4b95d9718f0b9e2e6b9g-552}

Figure 12D. 1 An applied magnetic field can induce circulation of electrons that makes use of excited state orbitals (shown with a white line).\\
observed. Many organic radicals, for which $\Delta E$ is large and $\xi$ (for carbon) is small, have $g$-values close to 2.0027, not far removed from $g_{e}$ itself. Inorganic radicals, which commonly are built from heavier atoms and therefore have larger spinorbit coupling constants, have $g$-values typically in the range 1.9-2.1. The $g$-values of paramagnetic d-metal complexes often differ considerably from $g_{e}$, varying from 0 to 6 , because in them $\Delta E$ is small on account of the small splitting of d-orbitals brought about by interactions with ligands (Topic 11F).

The $g$-value is anisotropic: that is, its magnitude depends on the orientation of the radical with respect to the applied field. The anisotropy arises from the fact that the extent to which an applied field induces currents in the molecule, and therefore the magnitude of the local field, depends on the relative orientation of the molecules and the field. In solution, when the molecule is tumbling rapidly, only the average value of the $g$-value is observed. Therefore, the anisotropy of the $g$-value is observed only for radicals trapped in solids and crystalline dmetal complexes.

\subsection*{12D. 2 Hyperfine structure}
The most important feature of an EPR spectrum is its hyperfine structure, the splitting of individual resonance lines into components. In general in spectroscopy, the term 'hyperfine structure' means the structure of a spectrum that can be traced to interactions of the electrons with nuclei other than as a result of the point electric charge of the nucleus. The source of the hyperfine structure in EPR is the magnetic interaction between the electron spin and the magnetic dipole moments of the nuclei present in the radical that give rise to local magnetic fields.

\subsection*{(a) The effects of nuclear spin}
Consider the effect on the EPR spectrum of a single ${ }^{1} \mathrm{H}$ nucleus located somewhere in a radical. The proton spin is a source of\\
magnetic field and, depending on the orientation of the nuclear spin, the field it generates either adds to or subtracts from the applied field. The total local field is therefore


\begin{equation*}
\mathscr{B}_{\mathrm{loc}}=\mathscr{B}_{0}+a m_{I} \quad m_{I}= \pm \frac{1}{2} \tag{12D.3}
\end{equation*}


where $a$ is the hyperfine coupling constant (or hyperfine splitting constant); from eqn 12D. 3 it follows that $a$ has the same units as the magnetic field, for example tesla. Half the radicals in a sample have $m_{I}=+\frac{1}{2}$, so half resonate when the applied field satisfies the condition


\begin{equation*}
h \nu=g \mu_{\mathrm{B}}\left(\mathscr{B}_{0}+\frac{1}{2} a\right), \quad \text { or } \mathscr{B}_{0}=\frac{h \nu}{g \mu_{\mathrm{B}}}-\frac{1}{2} a \tag{12D.4a}
\end{equation*}


The other half (which have $m_{I}=-\frac{1}{2}$ ) resonate when


\begin{equation*}
h \nu=g \mu_{\mathrm{B}}\left(\mathcal{B}_{0}-\frac{1}{2} a\right), \quad \text { or } \mathcal{B}_{0}=\frac{h \nu}{g \mu_{\mathrm{B}}}+\frac{1}{2} a \tag{12D.4b}
\end{equation*}


Therefore, instead of a single line, the spectrum shows two lines of half the original intensity separated by $a$ and centred on the field determined by $g$ (Fig. 12D.2).

If the radical contains an ${ }^{14} \mathrm{~N}$ atom ( $I=1$ ), its EPR spectrum consists of three lines of equal intensity, because the ${ }^{14} \mathrm{~N}$ nucleus has three possible spin orientations, and each spin orientation is possessed by one-third of all the radicals in the sample. In general, a spin-I nucleus splits the spectrum into $2 I+1$ hyperfine lines of equal intensity.\\
\includegraphics[max width=\textwidth, center]{2025_02_27_d4b95d9718f0b9e2e6b9g-552(1)}

Figure 12D. 2 The hyperfine interaction between an electron and a spin $-\frac{1}{2}$ nucleus results in four energy levels in place of the original two; $\alpha_{N}$ and $\beta_{N}$ indicate the spin states of the nucleus. As a result, the spectrum consists of two lines (of equal intensity) instead of one. The intensity distribution can be summarized by a simple stick diagram. The diagonal lines show the energies of the states as the applied field is increased, and resonance occurs when the separation of states matches the fixed energy of the microwave photon.\\
\includegraphics[max width=\textwidth, center]{2025_02_27_d4b95d9718f0b9e2e6b9g-553}

Figure 12D. 3 The EPR spectrum of the benzene radical anion, $\mathrm{C}_{6} \mathrm{H}_{6}^{-}$, in solution; $a$ is the hyperfine coupling constant. The centre of the spectrum is determined by the $g$-value of the radical.

When there are several magnetic nuclei present in the radical, each one contributes to the hyperfine structure. In the case

1\\
11\\
121\\
$1 \begin{array}{lll}1 & 3 & 1\end{array}$\\
14641\\
1 of equivalent protons (for example, the two $\mathrm{CH}_{2}$ protons in the radical $\mathrm{CH}_{3} \mathrm{CH}_{2}$ ) some of the hyperfine lines are coincident. If the radical contains $N$ equivalent protons, then there are $N+1$ hyperfine lines with an intensity distribution given by Pascal's triangle (1). The spectrum of the benzene radical anion in Fig. 12D.3, which has seven lines with intensity ratio $1: 6: 15: 20: 15: 6: 1$, is consistent with a radical containing six equivalent protons. More generally, if the radical contains $N$ equivalent nuclei with spin quantum number $I$, then there are $2 N I+1$ hyperfine lines.

\subsection*{Example 12D. 1}
Predicting the hyperfine structure of an\\
EPR spectrum\\
A radical contains one ${ }^{14} \mathrm{~N}$ nucleus ( $I=1$ ) with hyperfine constant 1.61 mT and two equivalent protons ( $I=\frac{1}{2}$ ) with hyperfine constant 0.35 mT . Predict the form of the EPR spectrum.

Collect your thoughts You will need to consider the hyperfine structure that arises from each type of nucleus or group of equivalent nuclei in succession. First, split a line with one nucleus, then split each of the lines again by a second nucleus (or group of nuclei), and so on. It is best to start with the nucleus with the largest hyperfine splitting; however, any choice could be made, and the order in which nuclei are considered does not affect the conclusion.

The solution The ${ }^{14} \mathrm{~N}$ nucleus gives three hyperfine lines of equal intensity separated by 1.61 mT . Each line is split into a doublet of spacing 0.35 mT by the first proton, and each line of these doublets is split into a doublet with the same 0.35 mT splitting by the second proton (Fig. 12D.4). Two of the lines in the centre coincide, so splitting by the two protons gives a 1:2:1 triplet of internal splitting 0.35 mT . Overall the spectrum consists of three identical 1:2:1 triplets.

Self-test 12D. 1 Predict the form of the EPR spectrum of a radical containing three equivalent ${ }^{14} \mathrm{~N}$ nuclei.\\
\includegraphics[max width=\textwidth, center]{2025_02_27_d4b95d9718f0b9e2e6b9g-553(1)}

Figure 12D. 4 The analysis of the hyperfine structure of a radical containing one ${ }^{14} \mathrm{~N}$ nucleus ( $I=1$ ) and two equivalent protons.\\
\includegraphics[max width=\textwidth, center]{2025_02_27_d4b95d9718f0b9e2e6b9g-553(2)}

Figure 12D. 5 The analysis of the hyperfine structure of a radical containing three equivalent ${ }^{14} \mathrm{~N}$ nuclei.

\subsection*{(b) The McConnell equation}
The hyperfine structure of an EPR spectrum is a kind of fingerprint that helps to identify the radicals present in a sample. Moreover, because the magnitude of the splitting depends on the distribution of the unpaired electron in the vicinity of the magnetic nuclei, the spectrum can be used to map the molecular orbital occupied by the unpaired electron.

The hyperfine splitting observed in $\mathrm{C}_{6} \mathrm{H}_{6}^{-}$is 0.375 mT . If it is assumed that the unpaired electron is in an orbital with equal probability at each $C$ atom, this hyperfine splitting can be attributed to the interaction between a proton and onesixth of the unpaired electron spin density. If all the electron density were located on the neighbouring C atom, a hyperfine coupling of $6 \times 0.375 \mathrm{mT}=2.25 \mathrm{mT}$ would be expected. If in another aromatic radical the hyperfine coupling constant is found to be $a$, then the spin density, $\rho$, the probability that an unpaired electron is on the neighbouring C atom, can be calculated from the McConnell equation:

$$
a=Q \rho
$$

McConnell equation\\
(12D.5)\\
with $Q=2.25 \mathrm{mT}$. In this equation, $\rho$ is the spin density on a C atom and $a$ is the hyperfine splitting observed for the H atom to which it is attached.

\subsection*{Brief illustration 12D. 2}
The hyperfine structure of the EPR spectrum of the naphthalene radical anion $\mathrm{C}_{10} \mathrm{H}_{8}^{-}$(2) can be interpreted as arising from\\
\includegraphics[max width=\textwidth, center]{2025_02_27_d4b95d9718f0b9e2e6b9g-554}

2 two groups of four equivalent protons. Those at the $1,4,5$, and 8 positions in the ring have $a=0.490 \mathrm{mT}$ and those in the $2,3,6$, and 7 positions have $a=0.183 \mathrm{mT}$. The spin densities obtained by using the McConnell equation are, respectively,

$$
\rho=\overbrace{\overbrace{0}^{\frac{\overbrace{0} .490 \mathrm{mT}}{2.25 \mathrm{mT}}}}^{a}=0.218 \quad \text { and } \quad \rho=\frac{0.183 \mathrm{mT}}{2.25 \mathrm{mT}}=0.0813
$$

\subsection*{(c) The origin of the hyperfine interaction}
An electron in a p orbital centred on a nucleus does not approach the nucleus very closely, so the electron experiences a magnetic field that appears to arise from a point magnetic dipole. The resulting interaction is called the dipole-dipole interaction. The contribution of a magnetic nucleus to the local field experienced by the unpaired electron is given by an expression like that in eqn 12B. 15 (a dependence proportional to $\left.\left(1-3 \cos ^{2} \theta\right) / r^{3}\right)$. A characteristic of this type of interaction is that it is anisotropic and averages to zero when the radical is free to tumble. Therefore, hyperfine structure due to the di-pole-dipole interaction is observed only for radicals trapped in solids.

There is a second contribution to the hyperfine splitting. An s electron is spherically distributed around a nucleus and so has zero average dipole-dipole interaction with the nucleus even in a solid sample. However, because an $s$ electron has a non-zero probability of being at the nucleus itself, it is incorrect to treat the interaction as one between two point dipoles. As explained in Topic 12B, an s electron has a 'Fermi contact interaction' with the nucleus, a magnetic interaction that occurs when the point dipole approximation fails. The contact interaction is isotropic (that is, independent of the orientation of the radical), and consequently is shown even by rapidly tumbling molecules in fluids (provided the spin density has some s character).

The dipole-dipole interactions of p electrons and the Fermi contact interaction of $s$ electrons can be quite large. For example, a $2 p$ electron in a nitrogen atom experiences an average

Table 12D. 1 Hyperfine coupling constants for atoms, $a / \mathrm{mT}^{*}$

\begin{center}
\begin{tabular}{lcc}
\hline
Nuclide & Isotropic coupling & Anisotropic coupling \\
\hline
${ }^{1} \mathrm{H}$ & $50.8(1 \mathrm{~s})$ &  \\
${ }^{2} \mathrm{H}$ & $7.8(1 \mathrm{~s})$ &  \\
${ }^{14} \mathrm{~N}$ & $55.2(2 \mathrm{~s})$ & $4.8(2 \mathrm{p})$ \\
${ }^{19} \mathrm{~F}$ & $1720(2 \mathrm{~s})$ & $108.4(2 \mathrm{p})$ \\
\hline
\end{tabular}
\end{center}

${ }^{*}$ More values are given in the Resource section.\\
field of about 3.4 mT from the ${ }^{14} \mathrm{~N}$ nucleus. A 1s electron in a hydrogen atom experiences a field of about 50 mT as a result of its Fermi contact interaction with the central proton. More values are listed in Table 12D.1. The magnitudes of the contact interactions in radicals can be interpreted in terms of the $s$ orbital character of the molecular orbital occupied by the unpaired electron, and the dipole-dipole interaction can be interpreted in terms of the p character. The analysis of hyperfine structure therefore gives information about the composition of the orbital, and especially the hybridization of the atomic orbitals.

\subsection*{Brief illustration 12D. 3}
From Table 12D.1, the hyperfine interaction between a 2 s electron and the nucleus of a nitrogen atom is 55.2 mT . The EPR spectrum of $\mathrm{NO}_{2}$ shows an isotropic hyperfine interaction of 5.7 mT . The s character of the molecular orbital occupied by the unpaired electron is the ratio $5.7 / 55.2=0.10$. For a continuation of this story, see Problem P12D.7.

Neither interaction appears to account for the hyperfine structure of the $\mathrm{C}_{6} \mathrm{H}_{6}^{-}$anion and other aromatic radical anions. The sample is fluid, and as the radicals are tumbling the hyperfine structure cannot be due to the dipole-dipole interaction. Moreover, the protons lie in the nodal plane of the $\pi$ orbital occupied by the unpaired electron, so the structure cannot be due to a Fermi contact interaction. The explanation lies in a polarization mechanism similar to the one responsible for spin-spin coupling in NMR. The magnetic interaction between a proton and the electrons favours one of the electrons\\
\includegraphics[max width=\textwidth, center]{2025_02_27_d4b95d9718f0b9e2e6b9g-554(1)}

Figure 12D. 6 The polarization mechanism for the hyperfine interaction in $\pi$-electron radicals. The arrangement in (a) is lower in energy than that in (b), so there is an effective coupling between the unpaired electron and the proton.\\
being found with a greater probability nearby (Fig. 12D.6). The electron with opposite spin is therefore more likely to be close to the C atom at the other end of the bond. The unpaired electron on the C atom has a lower energy if it is parallel to that\\
electron (Hund's rule favours parallel electrons on atoms), so the unpaired electron can detect the spin of the proton indirectly. Calculation using this model leads to a hyperfine interaction in agreement with the observed value of 2.25 mT .

\subsection*{Checklist of concepts}
1. The EPR resonance condition is expressed in terms of the $\boldsymbol{g}$-value of the radical.2. The value of $g$ depends on the ability of the applied field to induce local electron currents in the radical and the magnetic field experienced by the electron as a result of these currents.3. The hyperfine structure of an EPR spectrum is the splitting of individual resonance lines into components by the magnetic interaction between the electron and nuclei with spin.4. If a radical contains $N$ equivalent nuclei with spin quantum number $I$, then there are $2 N I+1$ hyperfine lines.5. Hyperfine structure arises from dipole-dipole interactions, Fermi contact interactions, and the polarization mechanism.6. The spin density on an atom is the probability that an unpaired electron is on that atom.\subsection*{Checklist of equations}
\begin{center}
\begin{tabular}{llll}
\hline
Property & Equation & Comment & Equation number \\
\hline
EPR resonance condition & $h \nu=g \mu_{\mathrm{B}} \mathcal{B}_{0}$ & No hyperfine interaction & 12 D .2 \\
 & $h \nu=g \mu_{\mathrm{B}}\left(\mathcal{B}_{0} \pm \frac{1}{2} a\right)$ & Hyperfine interaction between an electron and a proton & 12 D .4 \\
McConnell equation & $a=Q \rho$ & $Q=2.25 \mathrm{mT}$ & 12 D .5 \\
\hline
\end{tabular}
\end{center}

\chapter{FOCUS 13 Statistical thermodynamics}
Statistical thermodynamics provides the link between the microscopic properties of matter and its bulk properties. It provides a means of calculating thermodynamic properties from structural and spectroscopic data and gives insight into the molecular origins of chemical properties.

\subsection*{13A The Boltzmann distribution}
The 'Boltzmann distribution', which is used to predict the populations of states in systems at thermal equilibrium, is among the most important equations in chemistry for it summarizes the populations of states. It also provides insight into the nature of 'temperature'.\\
13A. 1 Configurations and weights; 13A. 2 The relative populations of states

\subsection*{13B Molecular partition functions}
The Boltzmann distribution introduces the central mathematical concept of a 'partition function'. The Topic shows how to interpret the partition function and how to calculate it in a number of simple cases.\\
13B. 1 The significance of the partition function; 13B. 2 Contributions to the partition function

\subsection*{13C Molecular energies}
A partition function is the thermodynamic version of a wavefunction, and contains all the thermodynamic information about a system. In this Topic partition functions are used to calculate the mean values of the energy of the basic modes of motion of a collection of independent molecules.\\
13C. 1 The basic equations; 13B. 2 Contributions of the fundamental modes of motion

\subsection*{13D The canonical ensemble}
Molecules do interact with one another, and statistical thermodynamics would be incomplete without being able to take\\
these interactions into account. This Topic shows how that is done in principle by introducing the 'canonical ensemble', and hints at how this concept can be used.\\
13D. 1 The concept of ensemble; 13D. 2 The mean energy of a system;\\
13D. 3 Independent molecules revisited; 13D. 4 The variation of the energy with volume

\subsection*{13E The internal energy and the entropy}
This Topic shows how molecular partition functions are used to calculate (and give insight into) the two basic thermodynamic functions, the internal energy and the entropy. The latter is based on another central equation introduced by Boltzmann, his definition of 'statistical entropy'.\\
13E. 1 The internal energy; 13E. 2 The entropy

\subsection*{13F Derived functions}
With expressions relating internal energy and entropy to partition functions, it is possible to develop expressions for the derived thermodynamic functions, such as the Helmholtz and Gibbs energies. Then, with the Gibbs energy available, the final step is taken into the calculations of chemically significant expressions by showing how equilibrium constants can be calculated from structural and spectroscopic data.\\
13F. 1 The derivations; 13F. 2 Equilibrium constants

\subsection*{Web resource What is an application of this material?}
There are numerous applications of statistical arguments in biochemistry. One of the most directly related to partition functions is explored in Impact 20: the helix-coil equilibrium in a polypeptide and the role of cooperative behaviour.

\section*{TOPIC 13A The Boltzmann distribution}
\subsection*{Why do you need to know this material?}
The Boltzmann distribution is the key to understanding a great deal of chemistry. All thermodynamic properties can be interpreted in its terms, as can the temperature dependence of equilibrium constants and the rates of chemical reactions. There is, perhaps, no more important unifying concept in chemistry.

\subsection*{What is the key idea?}
The most probable distribution of molecules over the available energy levels subject to certain restraints depends on a single parameter, the temperature.

\subsection*{What do you need to know already?}
You need to be aware that molecules can exist only in certain discrete energy levels (Topic 7A) and that in some cases more than one state has the same energy.

The problem addressed in this Topic is the calculation of the populations of states for any type of molecule in any mode of motion at any temperature. The only restriction is that the molecules should be independent, in the sense that the total energy of the system is a sum of their individual energies. In a real system a contribution to the total energy may arise from interactions between molecules, but that possibility is discounted at this stage. The development is based on the principle of equal a priori probabilities, the assumption that all possibilities for the distribution of energy are equally probable. 'A priori' means loosely in this context 'as far as one knows'. There is no reason to presume otherwise than that for a collection of molecules at thermal equilibrium, a vibrational state of a certain energy, for instance, is as likely to be populated as a rotational state of the same energy.

One very important conclusion that will emerge from the following analysis is that the overwhelmingly most probable populations of the available states depend on a single parameter, the 'temperature'. That is, the work done here provides a molecular justification for the concept of temperature and some insight into this crucially important quantity.

\subsection*{13A. 1 Configurations and weights}
Any individual molecule may exist in states with energies $\varepsilon_{0}$, $\varepsilon_{1}, \ldots$. For reasons that will become clear, the lowest available state is always taken as the zero of energy (that is, $\varepsilon_{0} \equiv 0$ ), and all other energies are measured relative to that state. To obtain the actual energy of the system it is necessary to add a constant to the energy calculated on this basis. For example, when considering the vibrational contribution to the energy, the total zero-point energy of any oscillators in the system must be added.

\subsection*{(a) Instantaneous configurations}
At any instant there will be $N_{0}$ molecules in the state 0 with en$\operatorname{ergy} \varepsilon_{0}, N_{1}$ in the state 1 with $\varepsilon_{1}$, and so on, with $N_{0}+N_{1}+\cdots=$ $N$, the total number of molecules in the system. The specification of the set of populations $N_{0}, N_{1}, \ldots$ in the form $\left\{N_{0}, N_{1}, \ldots\right\}$ is a statement of the instantaneous configuration of the system. The instantaneous configuration fluctuates with time because the populations change, perhaps as a result of collisions.

Initially suppose that all the states have exactly the same energy. The energies of all the configurations are then identical, so there is no restriction on how many of the $N$ molecules are in each state. Now picture a large number of different instantaneous configurations. One, for example, might be $\{N, 0,0, \ldots\}$, corresponding to every molecule being in state 0 . Another might be $\{N-2,2,0,0, \ldots\}$, in which two molecules are in state 1 . The latter configuration is intrinsically more likely to be found than the former because it can be achieved in more ways: $\{N, 0,0, \ldots\}$ can be achieved in only one way, but $\{N-2,2,0, \ldots\}$ can be achieved in $\frac{1}{2} N(N-1)$ different ways. Thus, one candidate for migration to state 1 can be selected in $N$ ways. There are $N-1$ candidates for the second choice, so the total number of choices is $N(N-1)$. However, the choice (Jack, Jill) cannot be distinguished from the choice (Jill, Jack) because they lead to the same configuration. Therefore, only half the choices lead to distinguishable configurations, and the total number of distinguishable choices is $\frac{1}{2} N(N-1)$. If, as a result of collisions, the system were to fluctuate between the configurations $\{N, 0,0, \ldots\}$ and $\{N-2,2,0, \ldots\}$, it would almost always be found in the second, more likely configuration, especially if $N$ were large. In other words, a system free to switch between the two configurations would show properties characteristic almost exclusively of the second configuration.\\
\includegraphics[max width=\textwidth, center]{2025_02_27_d4b95d9718f0b9e2e6b9g-565}

Figure 13A. 1 Eighteen molecules (the vertical bars) shown here are distributed into four receptacles (distinguished by the three vertical lines) such that there are 3 molecules in the first, 6 in the second, and so on. There are 18! different ways in which this distribution can be achieved. However, there are 3! equivalent ways of putting three molecules in the first receptacle, and likewise 6! equivalent ways of putting six molecules into the second receptacle, and so on. Hence the number of distinguishable arrangements is $18!/ 3!6!5!4!$, or about 515 million.

The next step is to develop an expression for the number of ways that a general configuration $\left\{N_{0}, N_{1}, \ldots\right\}$ can be achieved. This number is called the weight of the configuration and denoted $\mathcal{W}$.

\subsection*{How is that done? 13A. 1 Evaluating the weight of a configuration}
Consider the number of ways of distributing $N$ balls into bins. The first ball can be selected in $N$ different ways, the next ball in $N-1$ different ways from the balls remaining, and so on. Therefore, there are $N(N-1) \ldots 1=N$ ! ways of selecting the balls for distribution over the bins. However, if there are $N_{0}$ balls in the bin labelled $\varepsilon_{0}$, there would be $N_{0}$ ! different ways in which the same balls could have been chosen (Fig. 13A.1). Similarly, there are $N_{1}$ ! ways in which the $N_{1}$ balls in the bin labelled $\varepsilon_{1}$ can be chosen, and so on. Therefore, the total number of distinguishable ways of distributing the balls so that there are $N_{0}$ in bin $\varepsilon_{0}, N_{1}$ in bin $\varepsilon_{1}$, etc. regardless of the order in which the balls were chosen is

$$
\left.-\mathcal{W}=\frac{N!}{N_{0}!N_{1}!N_{2}!\cdots} \right\rvert\, \quad \text { Weight of a configuration }
$$

\subsection*{Brief illustration 13A. 1}
To calculate the number of ways of distributing 20 identical objects with the arrangement $1,0,3,5,10,1$, note that the configuration is $\{1,0,3,5,10,1\}$ with $N=20$. Remember that $0!\equiv 1$, Therefore the weight is

$$
\mathcal{W}=\frac{20!}{1!0!3!5!10!1!}=9.31 \times 10^{8}
$$

It will turn out to be more convenient to deal with the natural logarithm of the weight, $\ln \mathcal{W}$, rather than with the weight itself:

$$
\begin{aligned}
\ln \mathscr{W} & =\ln \frac{N!}{N_{0}!N_{1}!N_{2}!\cdots}=\ln N!-\ln N_{0}!N_{1}!N_{2}!\cdots \\
& =\ln x y=\ln x+\ln y \\
& =\ln N!-\ln N_{0}!-\ln N_{1}!-\ln N_{2}!-\cdots=\ln N!-\sum_{i} \ln N_{i}!
\end{aligned}
$$

One reason for introducing $\ln \mathcal{W}$ is that it is easier to make approximations. In particular, the factorials can be simplified by using Stirling's approximation ${ }^{1}$

\[
\begin{array}{ll}
\ln x!\approx x \ln x-x & \begin{array}{l}
\text { Stirling's approximation } \\
{[x \gg 1]}
\end{array} \tag{13A.2}
\end{array}
\]

Then the approximate expression for the weight is


\begin{align*}
\ln \mathcal{W} & =\{N \ln N-N\}-\sum_{i}\left\{N_{i} \ln N_{i}-N_{i}\right\} \\
& =N \ln N-N-\sum_{i} N_{i} \ln N_{i}+N \quad\left[\text { because } \sum_{i} N_{i}=N\right] \\
& =N \ln N-\sum_{i} N_{i} \ln N_{i} \tag{13A.3}
\end{align*}


\subsection*{(b) The most probable distribution}
The configuration $\{N-2,2,0, \ldots\}$ has much greater weight than $\{N, 0,0, \ldots\}$, and it should be easy to believe that there may be other configurations that have a much greater weight than both. In fact, for large $N$ there is a configuration with so great a weight that it overwhelms all the rest in importance to such an extent that the system will almost always be found in it. The properties of the system will therefore be characteristic of that particular dominating configuration. This dominating configuration can be found by looking for the values of $N_{i}$ that lead to a maximum value of $\mathcal{W}$. Because $\mathcal{W}$ is a function of all the $N_{i}$, this search is done by varying the $N_{i}$ and looking for the values that correspond to $\mathrm{d} \mathcal{W}=0$ (just as in the search for the maximum of any function), or equivalently a maximum value of $\ln \mathcal{W}$. Because $\ln \mathcal{W}$ depends on all the $N_{i}$, when a configuration changes and the $N_{i}$ change to $N_{i}+\mathrm{d} N_{i}$, the function $\ln \mathcal{W}$ changes to $\ln \mathcal{W}+\mathrm{d} \ln \mathcal{W}$, where


\begin{equation*}
\mathrm{d} \ln \mathcal{W}=\sum_{i}\left(\frac{\partial \ln \mathcal{W}}{\partial N_{i}}\right) \mathrm{d} N_{i}=0 \tag{13A.4}
\end{equation*}


when $\mathcal{W}$ is a maximum

The derivative $\partial \ln \mathcal{W} / \partial N_{i}$ expresses how $\ln \mathcal{W}$ changes when $N_{i}$ changes: if $N_{i}$ changes by $\mathrm{d} N_{i}$, then $\ln \mathcal{W}$ changes by $\left(\partial \ln \mathcal{W} / \partial N_{i}\right) \times \mathrm{d} N_{i}$. The total change in $\ln \mathcal{W}$ is then the sum of

\footnotetext{${ }^{1}$ A more precise form of this approximation is $\ln x!\approx \ln (2 \pi)^{1 / 2}+\left(x+\frac{1}{2}\right) \ln x-x$.
}
all these changes. However, there are two difficulties with this procedure.

Until now it has been assumed that all the states have the same energy. That restriction must now be removed and only configurations that correspond to the specified, constant, total energy of the system must be retained. This requirement rules out many configurations; $\{N, 0,0, \ldots\}$ and $\{N-2,2,0, \ldots\}$, for instance, have different energies (unless $\varepsilon_{0}$ and $\varepsilon_{1}$ happen to have the same energy), so both cannot occur in the same isolated system. It follows that the configuration with the greatest weight must also satisfy the condition

$$
\sum_{i} N_{i} \varepsilon_{i}=E \quad \begin{aligned}
& \text { Energy constraint } \\
& \text { [constant total energy] }
\end{aligned}
$$

(13A.5a)\\
where $E$ is the total energy of the system. Therefore, when the $N_{i}$ change by $\mathrm{d} N_{i}$, the total energy must not change, so

$$
\sum_{i} \varepsilon_{i} \mathrm{~d} N_{i}=0
$$

Energy constraint\\
(13A.5b)

The second constraint is that, because the total number of molecules present is also fixed (at $N$ ), not all the populations can be varied independently. Thus, increasing the population of one state by 1 demands that the population of another state must be reduced by 1 . Therefore, the search for the maximum value of $\mathcal{W}$ is also subject to the condition

\[
\sum_{i} N_{i}=N \quad \begin{align*}
& \text { Number constraint }  \tag{13A.6a}\\
& \text { [constant total number of molecules] }
\end{align*}
\]

It follows that when the $N_{i}$ change by $\mathrm{d} N_{i}$, this sum too must not change, so

\[
\sum_{i} \mathrm{~d} N_{i}=0 \quad \begin{align*}
& \text { Number }  \tag{13A.6b}\\
& \text { constraint }
\end{align*}
\]

The challenge now is to see how to solve eqn 13A. 4 subject to these two constraints.

\subsection*{How is that done? 13A. 2 \\
 Imposing constraints}
The way to take constraints into account was devised by the mathematician Joseph-Louis Lagrange, and is called the 'method of undetermined multipliers':

\begin{itemize}
  \item Multiply each constraint by a constant and then add it to the main variation equation.
  \item Now treat the variables as though they are all independent.
  \item Evaluate the constants at a later stage of the calculation.
\end{itemize}

\subsection*{Step 1 Introduce the constants}
There are two constraints, so introduce two constants $\alpha$ and $-\beta$ (the negative sign will be helpful later), leading to

$$
\begin{aligned}
& \overbrace{\sum_{i}\left(\frac{\partial \ln \mathcal{W}}{\partial N_{i}}\right) \mathrm{d} N_{i}}^{\begin{array}{c}
\text { Original } \\
\text { expression }
\end{array}}+\alpha \overbrace{\sum_{i} \mathrm{~d} N_{i}}-\beta \overbrace{\sum_{i} \varepsilon_{i} \mathrm{~d} N_{i}}^{\begin{array}{c}
\text { Number } \\
\text { constraint }
\end{array}} \\
& =\sum_{i}\left\{\left(\frac{\partial \ln \mathcal{W}}{\partial N_{i}}\right)+\alpha-\beta \varepsilon_{i}\right\} \mathrm{d} N_{i} \\
& =0 \text { at a maximum }
\end{aligned}
$$

Step 2 Treat the variables as independent\\
The $\mathrm{d} N_{i}$ are now treated as independent. Hence, the only way of satisfying $\mathrm{d} \ln \mathcal{W}=0$ is to require that for each $i$,

$$
\left.-\left(\frac{\partial \ln \mathcal{W}}{\partial N_{i}}\right)+\alpha-\beta \varepsilon_{i}=0 \right\rvert\, \begin{gathered}
\text { (13A.7) } \\
\end{gathered}
$$

The next step in this lengthy derivation is to insert the expression for $\ln \mathcal{W}$ (eqn 13A.3) into this equation. That involves evaluating the differentiation of $\ln \mathcal{W}$ with respect to $N_{i}$.

\subsection*{How is that done? 13A. 3 \\
 Evaluating the derivative of $\ln \mathscr{W}$}
In preparation for this calculation it is useful to change eqn 13A. 3 from

$$
\ln \mathcal{W}=N \ln N-\sum_{i} N_{i} \ln N_{i}
$$

to

$$
\ln \mathcal{W}=N \ln N-\sum_{j} N_{j} \ln N_{j}
$$

by using $j$ instead of $i$ as the 'name' of the states. In this way the $i$ in the differentiation variable $\left(N_{i}\right)$ will not be confused with the $i$ in the summation. Differentiation of this expression gives

$$
\frac{\partial \ln \mathcal{W}}{\partial N_{i}}=\overbrace{\frac{\partial(N \ln N)}{\partial N_{i}}}^{\text {First term }}-\overbrace{\sum_{j} \frac{\partial\left(N_{j} \ln N_{j}\right)}{\partial N_{i}}}^{\text {Second term }}
$$

\subsection*{Step 1 Evaluate the first term in the expression}
The first term on the right is obtained (by using the product rule) as follows:

$$
\frac{\partial(N \ln N)}{\partial N_{i}}=\left(\frac{\partial N}{\partial N_{i}}\right) \ln N+N\left(\frac{\partial \ln N}{\partial N_{i}}\right)
$$

Now note that

$$
\frac{\partial N}{\partial N_{i}}=\frac{\partial}{\partial N_{i}}\left(N_{1}+N_{2}+\cdots\right)=1
$$

because $N_{i}$ will match one and only one term in the sum whatever the value of $i$. Next, note that

$$
\frac{\partial \ln N}{\partial N_{i}}=\frac{\underbrace{d \ln y / d x=(1 / y) d y / d x}}{\frac{1}{N} \frac{\overbrace{N}}{\partial N_{i}}=\frac{1}{N}}
$$

Therefore

$$
\frac{\partial(N \ln N)}{\partial N_{i}}=\ln N+1
$$

Step 2 Evaluate the second term\\
For the derivative of the second term, first note that

$$
\begin{aligned}
\sum_{j} \frac{\partial\left(N_{j} \ln N_{j}\right)}{\partial N_{i}} & =\sum_{j}\left\{\left(\frac{\partial N_{j}}{\partial N_{i}}\right) \ln N_{j}+N_{j}\left(\frac{\partial \ln N_{j}}{\partial N_{i}}\right)\right\} \\
& \underbrace{\mathrm{d} \ln y / \mathrm{d} x=(1 / y) \mathrm{d} y / \mathrm{d} x=\mathrm{dd} / \mathrm{d} x+g \mathrm{~d} / \mathrm{d} x} \\
& =\sum_{j}\left\{\left(\frac{\partial N_{j}}{\partial N_{i}}\right) \ln N_{j}+\left(\frac{\partial N_{j}}{\partial N_{i}}\right)\right\} \\
& =\sum_{j}\left\{\ln N_{j}+1\right\}\left(\frac{\partial N_{j}}{\partial N_{i}}\right)
\end{aligned}
$$

All the $N_{j}$ are independent, so the only term that survives in the differentiation $\partial N_{j} / \partial N_{i}$ is the one with $j=i$, and then $\partial N_{i} / \partial N_{i}=1$. It follows that

$$
\sum_{j} \frac{\partial\left(N_{j} \ln N_{j}\right)}{\partial N_{i}}=\ln N_{i}+1
$$

Step 3 Bring the two terms together\\
Bringing the first and second terms together gives

$$
\frac{\partial \ln \mathcal{W}}{\partial N_{i}}=\ln N+1-\left(\ln N_{i}+1\right)
$$

That is,

$$
\frac{\partial \ln \mathcal{W}}{\partial N_{i}}=-\ln \frac{N_{i}}{N}
$$

It now follows from eqn 13A. 7 that

$$
-\ln \frac{N_{i}}{N}+\alpha-\beta \varepsilon_{i}=0
$$

and therefore that


\begin{equation*}
\frac{N_{i}}{N}=\mathrm{e}^{\alpha-\beta \varepsilon_{i}} \tag{13A.8}
\end{equation*}


which is very close to being the Boltzmann distribution.

\subsection*{(b) The values of the constants}
At this stage note that

$$
N=\sum_{i} N_{i}=\sum_{i} N \mathrm{e}^{\alpha-\beta \varepsilon_{i}}=N \mathrm{e}^{\alpha} \sum_{i} \mathrm{e}^{-\beta \varepsilon_{i}}
$$

Because the $N$ cancels on each side of this equality, it follows that


\begin{equation*}
\mathrm{e}^{\alpha}=\frac{1}{\sum_{i} \mathrm{e}^{-\beta \varepsilon_{i}}} \tag{13A.9}
\end{equation*}


and therefore


\begin{equation*}
\frac{N_{i}}{N}=\mathrm{e}^{\alpha-\beta \varepsilon_{i}}=\mathrm{e}^{\alpha} \mathrm{e}^{-\beta \varepsilon_{i}}=\frac{\mathrm{e}^{-\beta \varepsilon_{i}}}{\sum_{i} \mathrm{e}^{-\beta \varepsilon_{i}}} \quad \text { Boltzmann distribution } \tag{13A.10a}
\end{equation*}


which is the Boltzmann distribution. This distribution is commonly written


\begin{equation*}
\frac{N_{i}}{N}=\frac{\mathrm{e}^{-\beta \varepsilon_{i}}}{q} \tag{13A.10b}
\end{equation*}


Boltzmann distribution\\
where $g$ is called the partition function:

\[
q=\sum_{i} \mathrm{e}^{-\beta \varepsilon_{i}} \quad \begin{align*}
& \text { Partition function }  \tag{13A.11}\\
& \text { [definition] }
\end{align*}
\]

At this stage the partition function is no more than a convenient abbreviation for the sum; but Topic 13B shows that it is central to the statistical interpretation of thermodynamic properties.

Equation 13A. 10 is the justification of the remark that a single parameter, here denoted $\beta$, governs the most probable populations of the states of the system, which strongly suggests that it is related to the temperature. The formal deduction of the value of $\beta$ depends on using the Boltzmann distribution to deduce the perfect gas equation of state (that is done in Topic 13F, and specifically in Example 13F.1), which confirms this relation and shows that


\begin{equation*}
\beta=\frac{1}{k T} \tag{13A.12}
\end{equation*}


where $T$ is the thermodynamic temperature and $k$ is Boltzmann's constant. In other words:

The temperature is the unique parameter that governs the most probable populations of states of a system at thermal equilibrium.

\subsection*{Brief illustration 13A. 2}
Suppose that two conformations of a molecule differ in energy by $5.0 \mathrm{~kJ} \mathrm{~mol}^{-1}$ (corresponding to 8.3 zJ for a single molecule; $1 \mathrm{zJ}=10^{-21} \mathrm{~J}$ ), so conformation A lies at energy 0 and conformation B lies at $\varepsilon=8.3 \mathrm{zJ}$. At $20^{\circ} \mathrm{C}(293 \mathrm{~K})$ the denominator in eqn 13A.10a is

$$
\left.\sum_{i} \mathrm{e}^{-\beta \varepsilon_{i}}=1+\mathrm{e}^{-\varepsilon / k T}=1+\mathrm{e}^{-\left(8.3 \times 10^{-21} \mathrm{~J}\right) /\left(1.381 \times 10^{-23} \mathrm{~J}\right.} \mathrm{K}^{-1}\right) \times(293 \mathrm{~K})=1.12 \ldots
$$

The proportion of molecules in conformation B at this temperature is therefore

$$
\frac{N_{\mathrm{B}}}{N}=\frac{\mathrm{e}^{-\left(8.3 \times 10^{-21} \mathrm{~J}\right) /\left(1.381 \times 10^{-23} \mathrm{~J} \mathrm{~K}^{-1}\right) \times(293 \mathrm{~K})}}{1.12 \ldots}=0.11
$$

or 11 per cent of the molecules.

\subsection*{13A. 2 The relative population of states}
When considering only the relative populations of states by using eqn 13A.10b, the partition function need not be evaluated, because it cancels when the ratio is taken:

\[
\frac{N_{i}}{N_{j}}=\frac{\mathrm{e}^{-\beta \varepsilon_{i}}}{\mathrm{e}^{-\beta \varepsilon_{j}}}=\mathrm{e}^{-\beta\left(\varepsilon_{i}-\varepsilon_{j}\right)} \quad \begin{align*}
& \text { Boltzmann population ratio }  \tag{13A.13a}\\
& {[\text { thermal equilibrium }]}
\end{align*}
\]

Note that for a given energy separation the ratio of populations $N_{1} / N_{0}$ decreases as $\beta$ increases (and the temperature decreases). At $T=0(\beta=\infty)$ all the population is in the ground state and the ratio is zero. Equation 13A.13a is enormously important for understanding a wide range of chemical phenomena and is the form in which the Boltzmann distribution is commonly employed (for instance, in the discussion of the intensities of spectral transitions, Topic 11A). It implies that the relative population of two states decreases exponentially with their difference in energy.\\
A very important point to note is that the Boltzmann distribution gives the relative populations of states, not energy levels. Several states might have the same energy, and each state has a population given by eqn 13A.13a. When calculating the relative populations of energy levels rather than states, it is necessary to take into account this degeneracy. Thus, if the level of energy $\varepsilon_{i}$ is $g_{i}$-fold degenerate (in the sense that there are $g_{i}$ states with that energy), and the level of energy $\varepsilon_{j}$ is $g_{j}-$ fold degenerate, then the relative total populations of the levels are given by


\begin{equation*}
\frac{N_{i}}{N_{j}}=\frac{g_{i} \mathrm{e}^{-\beta \varepsilon_{i}}}{g_{j} \mathrm{e}^{-\beta \varepsilon_{j}}}=\frac{g_{i}}{g_{j}} \mathrm{e}^{-\beta\left(\varepsilon_{i}-\varepsilon_{j}\right)} \tag{13A.13b}
\end{equation*}


Boltzmann population ratio [thermal equilibrium,\\[0pt]
degeneracies]

\subsection*{Example 13A. 1 Calculating the relative populations of rotational states}
Calculate the relative populations of the $J=1$ and $J=0$ rotational levels of HCl at $25^{\circ} \mathrm{C}$; for $\mathrm{HCl}, \tilde{B}=10.591 \mathrm{~cm}^{-1}$.

Collect your thoughts Although the ground state is nondegenerate, you need to note that the level with $J=1$ is triply degenerate ( $M_{J}=0, \pm 1$ ); the energy of the state with quantum number $J$ is $\varepsilon_{J}=h c \tilde{B} J(J+1)$ (Topic 11B). A useful relation is $k T / h c=207.22 \mathrm{~cm}^{-1}$ at 298.15 K .

The solution The energy separation of states with $J=1$ and $J=0$ is $\varepsilon_{1}-\varepsilon_{0}=2 h c \tilde{B}$. The ratio of the population of a level with $J=1$ and any one of its three states $M_{I}$ to the population of the single state with $J=0$ is therefore

$$
\frac{N_{J, M_{J}}}{N_{0}}=\mathrm{e}^{-2 h c \tilde{B} \beta}
$$

The relative populations of the levels, taking into account the three-fold degeneracy of the upper level, is

$$
\frac{N_{J}}{N_{0}}=3 \mathrm{e}^{-2 h c \tilde{\beta} \beta}
$$

Insertion of $h c \tilde{B} \beta=h c \tilde{B} / k T=\left(10.591 \mathrm{~cm}^{-1}\right) /\left(207.22 \mathrm{~cm}^{-1}\right)=$ $0.0511 \ldots$ then gives

$$
\frac{N_{J}}{N_{0}}=3 \mathrm{e}^{-2 \times 0.0511 \ldots}=2.708
$$

Comment. Because the $J=1$ level is triply degenerate, it has a higher population than the level with $J=0$, despite being of higher energy. As the example illustrates, it is very important to take note of whether you are asked for the relative populations of individual states or of a (possibly degenerate) energy level.

Self-test 13A. 1 What is the ratio of the populations of the levels with $J=2$ and $J=1$ of HCl at the same temperature?


\subsection*{Checklist of concepts}
\begin{enumerate}
  \item The principle of equal a priori probabilities assumes that all possibilities for the distribution of energy are equally probable in the sense that the distribution is blind to the type of motion involved.2. The instantaneous configuration of a system of $N$ molecules is the specification of the set of populations $N_{0}, N_{1}, \ldots$ of the energy levels $\varepsilon_{0}, \varepsilon_{1}, \ldots$.
  \item The Boltzmann distribution gives the numbers of molecules in each state of a system at any temperature.\\
$\square$ 4. The relative populations of energy levels, as opposed to states, must take into account the degeneracies of the energy levels.
\end{enumerate}

\subsection*{Checklist of equations}
\begin{center}
\begin{tabular}{llll}
\hline
Property & Equation & Comment & Equation number \\
\hline
Boltzmann distribution & $N_{i} / N=\mathrm{e}^{-\beta \varepsilon_{i}} / \mathcal{G}$ & $\beta=1 / k T$ & 13 A .10 b \\
Partition function & $\mathcal{q}=\sum_{i} \mathrm{e}^{-\beta \varepsilon_{i}}$ & See Topic 13B & 13 A .11 \\
Boltzmann population ratio & $N_{i} / N_{j}=\left(g_{i} / g_{j}\right) \mathrm{e}^{-\beta\left(\varepsilon_{i}-\varepsilon_{j}\right)}$ & $g_{i}, g_{j}$ are degeneracies & 13A.13b \\
\hline
\end{tabular}
\end{center}

\section*{TOPIC 13B Molecular partition functions}
Why do you need to know this material?\\
Through the partition function, statistical thermodynamics provides the link between thermodynamic data and molecular properties that have been calculated or derived from spectroscopy. Therefore, this material is an essential foundation for understanding physical and chemical properties of bulk matter in terms of the properties of the constituent molecules.

\subsection*{What is the key idea?}
The partition function is calculated by drawing on calculated or spectroscopically derived structural information about molecules.

What do you need to know already?\\
You need to know that the Boltzmann distribution expresses the most probable distribution of molecules over the available energy levels (Topic 13A). The concept of the partition function is introduced in that Topic, and is developed here. You need to be aware of the expressions for the rotational and vibrational levels of molecules (Topics 11B-11D) and the energy levels of a particle in a box (Topic 7D).

The partition function $q=\sum \mathrm{e}^{-\beta \varepsilon_{i}}$ is introduced in Topic 13A simply as a symbol to denote the sum over states that occurs in the denominator of the Boltzmann distribution (eqn $13 \mathrm{~A} .10 \mathrm{~b}, p_{i}=\mathrm{e}^{-\beta \varepsilon_{i}} / \mathcal{G}$, with $\left.p_{i}=N_{i} / N\right)$. But it is far more important than that might suggest. For instance, it contains all the information needed to calculate the bulk properties of a system of independent particles. In this respect $g$ plays a role for bulk matter very similar to that played by the wavefunction in quantum mechanics for individual molecules: $q$ is a kind of thermal wavefunction.

\subsection*{13B.1 The significance of the partition function}
The molecular partition function is

$$
\mathcal{G}=\sum_{\text {states } i} \mathrm{e}^{-\beta \varepsilon_{i}}
$$

Molecular partition function [definition]\\
(13B.1a)\\
where $\beta=1 / k T$. As emphasized in Topic 13A, the sum is over states, not energy levels. If $g_{i}$ states have the same energy $\varepsilon_{i}$ (so the level is $g_{i}$-fold degenerate), then

\[
\mathcal{G}=\sum_{\text {levels } i} g_{i} e^{-\beta \varepsilon_{i}} \quad \begin{align*}
& \text { Molecular partition function }  \tag{13B.1b}\\
& \text { [alternative definition] }
\end{align*}
\]

where the sum is now over energy levels (sets of states with the same energy), not individual states. Also as emphasized in Topic 13A, the lowest available state is taken as the zero of energy, so $\varepsilon_{0} \equiv 0$.

\subsection*{Brief illustration 13B. 1}
Suppose a molecule is confined to the following non-degenerate energy levels: $0, \varepsilon, 2 \varepsilon, \ldots$ (Fig. 13B.1). Then the molecular partition function is

$$
q=1+\mathrm{e}^{-\beta \varepsilon}+\mathrm{e}^{-2 \beta \varepsilon}+\cdots=1+\mathrm{e}^{-\beta \varepsilon}+\left(\mathrm{e}^{-\beta \varepsilon}\right)^{2}+\cdots
$$

\begin{center}
\includegraphics[max width=\textwidth]{2025_02_27_d4b95d9718f0b9e2e6b9g-570}
\end{center}

Figure 13B. 1 The equally spaced infinite array of energy levels used in the calculation of the partition function. A harmonic oscillator has the same array of levels.

Provided $|x|<1$, the sum to infinity of the geometrical series $1+x+x^{2}+\cdots$ is $1 /(1-x)$. In this case $x=\mathrm{e}^{-\beta \varepsilon}<1$, so the series evaluates to

$$
\mathcal{G}=\frac{1}{1-\mathrm{e}^{-\beta \varepsilon}}
$$

This function is plotted in Fig. 13B. 2 (with $\beta=1 / k T$ ).\\
\includegraphics[max width=\textwidth, center]{2025_02_27_d4b95d9718f0b9e2e6b9g-571}

Figure 13B.2 The partition function for the system shown in Fig. 13B. 1 (a harmonic oscillator) as a function of temperature.

The result in the Brief illustration is an important expression for the partition function for a uniform ladder of states of spacing $\varepsilon$ :

\[
\mathcal{q}=\frac{1}{1-\mathrm{e}^{-\beta \varepsilon}} \quad \begin{align*}
& \text { Partition function }  \tag{13B.2a}\\
& \text { [uniform ladder] }
\end{align*}
\]

This expression can be used to interpret the physical significance of a partition function. To do so, first note that the Boltzmann distribution for this arrangement of energy levels gives the fraction, $p_{i}=N_{i} / N$, of molecules in the state with energy $\varepsilon_{i}$ as

\[
p_{i}=\frac{\mathrm{e}^{-\beta \varepsilon_{i}}}{q}=\left(1-\mathrm{e}^{-\beta \varepsilon}\right) \mathrm{e}^{-\beta \varepsilon_{i}} \quad \begin{align*}
& \text { Fractional population }  \tag{13B.2b}\\
& {[\text { uniform ladder }]}
\end{align*}
\]

Figure 13B. 3 shows how $p_{i}$ varies with temperature. At very low temperatures (high $\beta$ ), where $q$ is close to 1 , only the lowest state is significantly populated. As the temperature is raised, the population breaks out of the lowest state, and the upper states become progressively more highly populated. At the same time, the partition function rises from 1 , so its value gives an indication of the range of states populated at any given temperature. The name 'partition function' reflects the sense in which $\mathcal{G}$ measures how the total number of molecules is distributed-partitioned-over the available states.

The corresponding expressions for a system in which there are just two states, with energies $\varepsilon_{0}=0$ and $\varepsilon_{1}=\varepsilon$ (a'two-level system'), are\\
\includegraphics[max width=\textwidth, center]{2025_02_27_d4b95d9718f0b9e2e6b9g-571(1)}

Figure 13B. 3 The populations of the energy levels of the system shown in Fig. 13B. 1 at different temperatures, and the corresponding values of the partition function calculated from eqn 13B.2a. Note that $\beta=1 / k T$.

$$
\begin{array}{ll}
\mathcal{q}=1+\mathrm{e}^{-\beta \varepsilon} & \begin{array}{l}
\text { Partition function } \\
\text { [two-level system] }
\end{array} \\
p_{i}=\frac{\mathrm{e}^{-\beta \varepsilon_{i}}}{\mathcal{q}}=\frac{\mathrm{e}^{-\beta \varepsilon_{i}}}{1+\mathrm{e}^{-\beta \varepsilon}} & \begin{array}{l}
\text { Fractional population } \\
{[\text { two-level system, } i=0,1]}
\end{array}
\end{array}
$$

The fractional populations of the two states are therefore


\begin{equation*}
p_{0}=\frac{1}{1+\mathrm{e}^{-\beta \varepsilon}} \quad p_{1}=\frac{\mathrm{e}^{-\beta \varepsilon}}{1+\mathrm{e}^{-\beta \varepsilon}} \tag{13B.4}
\end{equation*}


Figure 13B. 4 shows the variation of the partition function with temperature and Fig. 13B. 5 shows how the fractional populations change. Notice how at $T=0$ the fractional populations are $p_{0}=1$ and $p_{1}=0$, and the partition function is $q=1$ (one state occupied). However, the fractional populations tend towards equality $\left(p_{0}=\frac{1}{2}, p_{1}=\frac{1}{2}\right)$ and $q=2$ (two states occupied) as $T \rightarrow \infty(\beta \rightarrow 0)$.\\
\includegraphics[max width=\textwidth, center]{2025_02_27_d4b95d9718f0b9e2e6b9g-571(2)}

Figure 13B. 4 The dependence of the partition function of a twolevel system on temperature. The two graphs differ in the scale of the temperature axis to show the approach to 1 as $T \rightarrow 0$ and the slow approach to 2 as $T \rightarrow \infty$.\\
\includegraphics[max width=\textwidth, center]{2025_02_27_d4b95d9718f0b9e2e6b9g-572}

Figure 13B. 5 The variation with temperature of the fractional populations of the two states of a two-level system (eqn 13B.4). Note that as the temperature approaches infinity, the populations of the two states become equal (and the fractional populations both approach 0.5 ).

A note on good practice A common error is to suppose that when $T=\infty$ all the molecules in the system will be found in the upper energy state. However, as seen from eqn 13B.4, as $T \rightarrow \infty$ the populations of states become equal. The same conclusion is also true of multi-level systems: as $T \rightarrow \infty$, all states become equally populated.

Now consider the general case of a system with an infinite number of energy levels as $T$ approaches zero and therefore as the parameter $\beta=1 / k T$ approaches infinity. Then every term except one in the sum defining $q$ in eqn 13B.1a is zero because each one has the form $\mathrm{e}^{-x}$ with $x \rightarrow \infty$. The exception is the term with $\varepsilon_{0} \equiv 0$ (or the $g_{0}$ states at zero energy if this level is $g_{0}$-fold degenerate), because then $\varepsilon_{0} / k T=0$ whatever the temperature, including zero. As there is only one surviving term when $T=0$, and its value is $g_{0}$, it follows that

$$
\lim _{T \rightarrow 0} \mathfrak{q}=g_{0}
$$

That is, at $T=0$, the partition function is equal to the degeneracy of the ground state (commonly, but not necessarily, 1 ).

When $T$ is so high that for each term in the sum $\beta \varepsilon_{i}=$ $\varepsilon_{i} / k T \approx 0$, each term in the sum now contributes 1 because $\mathrm{e}^{-x}=1$ when $x=0$. It follows that the sum is equal to the number of molecular states, which in general is infinite:

$$
\lim _{T \rightarrow \infty} \mathscr{G}=\infty
$$

In some idealized cases, the molecule may have only a finite number of states; then the upper limit of $q$ is equal to the number of states, as for the two-level system.

In summary,\\
The molecular partition function gives an indication of the number of states that are thermally accessible to a molecule at the temperature of the system.

\subsection*{138.2 Contributions to the partition function}
The energy of an isolated molecule is the sum of contributions from its different modes of motion:


\begin{equation*}
\varepsilon_{i}=\varepsilon_{i}^{\mathrm{T}}+\varepsilon_{i}^{\mathrm{R}}+\varepsilon_{i}^{\mathrm{V}}+\varepsilon_{i}^{\mathrm{E}} \tag{13B.5}
\end{equation*}


where T denotes translation, R rotation, V vibration, and E the electronic contribution. The possibility that the molecules are interacting with each other is ignored throughout this Topic, because it adds considerable complexity: that is, the molecules are treated as 'independent'. The electronic contribution is not actually a 'mode of motion', but it is convenient to include it here. The separation of terms in eqn 13B. 5 is only approximate (except for translation) because the modes are not completely independent, but in most cases it is satisfactory.

Given that the energy is a sum of independent contributions, the partition function factorizes into a product of contributions:

$$
\begin{aligned}
\mathcal{q} & =\sum_{i} \mathrm{e}^{-\beta \varepsilon_{i}}=\sum_{i(\text { all states })} \mathrm{e}^{-\beta \varepsilon_{i}^{\mathrm{T}}-\beta \varepsilon_{i}^{\mathrm{R}}-\beta \varepsilon_{i}^{\mathrm{V}}-\beta \varepsilon_{i}^{\mathrm{E}}} \\
= & \sum_{i(\text { translational })} \sum_{i(\text { rotational })} \sum_{i(\text { vibrational })} \sum_{i(\text { electronic })} \mathrm{e}^{-\beta \varepsilon_{i}^{\mathrm{T}}-\beta \varepsilon_{i}^{\mathrm{R}}-\beta \varepsilon_{i}^{\mathrm{V}}-\beta \varepsilon_{i}^{\mathrm{E}}} \\
= & \left(\sum_{i(\text { translational })} \mathrm{e}^{-\beta \varepsilon_{i}^{\mathrm{T}}}\right)\left(\sum_{i(\text { rotational })} \mathrm{e}^{-\beta \varepsilon_{i}^{\mathrm{R}}}\right)\left(\sum_{i(\text { vibrational })} \mathrm{e}^{-\beta \varepsilon_{i}^{\mathrm{V}}}\right) \\
& \times\left(\sum_{i(\text { electronic })} \mathrm{e}^{-\beta \varepsilon_{i}^{\mathrm{E}}}\right)
\end{aligned}
$$

That is,


\begin{equation*}
\mathfrak{q}=\mathfrak{q}^{\mathrm{T}} \mathfrak{q}^{\mathrm{R}} \mathfrak{q}^{\mathrm{V}} \mathfrak{q}^{\mathrm{E}} \quad \text { Factorization of the partition function } \tag{13B.6}
\end{equation*}


This factorization means that each contribution can be investigated separately. In general, exact analytical expressions for partition functions cannot be obtained. However, approximate expressions can often be found and prove to be very important for understanding chemical phenomena; they are derived in the following sections and collected at the end of this Topic.

\subsection*{(a) The translational contribution}
The translational partition function for a particle of mass $m$ free to move in a one-dimensional container of length $X$ can be evaluated by making use of the fact that the separation of energy levels is very small and that large numbers of states are accessible at normal temperatures.

\subsection*{How is that done? 13 B. 1 Deriving an expression for the}
 translational partition functionThe starting point of the derivation is eqn7D. $6\left(E_{n}=n^{2} h^{2} / 8 m L^{2}\right)$ for the energy levels of a particle in a one-dimensional box. For a molecule of mass $m$ in a container of length $X$ it follows that:

$$
E_{n}=\frac{n^{2} h^{2}}{8 m X^{2}}
$$

Step 1 Write an expression for the sum in eqn 13B.1a\\
The lowest level $(n=1)$ has energy $h^{2} / 8 m X^{2}$, so the energies relative to that level are

$$
\varepsilon_{n}=\left(n^{2}-1\right) \varepsilon \quad \varepsilon=h^{2} / 8 m X^{2}
$$

The sum to evaluate is therefore

$$
q_{x}^{\mathrm{T}}=\sum_{n=1}^{\infty} \mathrm{e}^{-\left(n^{2}-1\right) \beta \varepsilon}
$$

Step 2 Convert the sum to an integral\\
The translational energy levels are very close together in a container the size of a typical laboratory vessel. Therefore, the sum can be approximated by an integral:\\
\includegraphics[max width=\textwidth, center]{2025_02_27_d4b95d9718f0b9e2e6b9g-573}

The extension of the lower limit to $n=0$ and the replacement of $n^{2}-1$ by $n^{2}$ introduces negligible error but turns the integral into a standard form.

Step 3 Evaluate the integral\\
Make the substitution $x^{2}=n^{2} \beta \varepsilon$, implying $\mathrm{d} n=\mathrm{d} x /(\beta \varepsilon)^{1 / 2}$, and therefore that

$$
\mathcal{q}_{X}^{\mathrm{T}}=\left(\frac{1}{\beta \varepsilon}\right)^{1 / 2} \overbrace{\int_{0}^{\infty} \mathrm{e}^{-x^{2}} \mathrm{~d} x}^{\pi^{1 / 2 / 2}}=\left(\frac{1}{\beta \varepsilon}\right)^{1 / 2} \frac{\pi^{1 / 2}}{2}=\overbrace{\left(\frac{2 \pi m}{h^{2} \beta}\right)^{2 / 2} X}=
$$

With $\beta=1 / k T$ this relation has the form

$$
\dashv q_{X}^{\mathrm{T}}=\frac{X}{\Lambda} \quad \Lambda=\frac{h}{(2 \pi m k T)^{1 / 2}} \xlongequal[\text { Translational partition function }]{ }
$$

The quantity $\Lambda$ (uppercase lambda) has the dimensions of length and is called the thermal wavelength (sometimes the 'thermal de Broglie wavelength') of the molecule. The thermal wavelength decreases with increasing mass and temperature. This expression shows that:

\begin{itemize}
  \item The partition function for translational motion increases with the length of the box and the mass of the particle, because in each case the separation of the energy levels becomes smaller and more levels become thermally accessible.
  \item For a given mass and length of the box, the partition function also increases with increasing temperature (decreasing $\beta$ ), because more states become accessible.
\end{itemize}

The total energy of a molecule free to move in three dimensions is the sum of its translational energies in all three directions:


\begin{equation*}
\varepsilon_{n_{1} n_{2} n_{3}}=\varepsilon_{n_{1}}^{(X)}+\varepsilon_{n_{2}}^{(Y)}+\varepsilon_{n_{3}}^{(Z)} \tag{13B.8}
\end{equation*}


where $n_{1}, n_{2}$, and $n_{3}$ are the quantum numbers for motion in the $x$-, $y$-, and $z$-directions, respectively. Therefore, because $\mathrm{e}^{a+b+c}=\mathrm{e}^{a} \mathrm{e}^{b} \mathrm{e}^{c}$, the partition function factorizes as follows:

$$
\begin{aligned}
\mathcal{q}^{\mathrm{T}} & =\sum_{\text {all } n} \mathrm{e}^{-\beta \varepsilon_{n_{1}}^{(X)}-\beta \varepsilon_{n_{2}}^{(\gamma)}-\beta \varepsilon_{n_{3}}^{(Z)}}=\sum_{\text {all } n} \mathrm{e}^{-\beta \varepsilon_{n_{1}}^{(X)}} \mathrm{e}^{-\beta \varepsilon_{n_{2}}^{(X)}} \mathrm{e}^{-\beta \varepsilon_{n_{3}}^{(Z)}} \\
& =\left(\sum_{n_{1}} \mathrm{e}^{-\beta \varepsilon_{n_{1}}^{(X)}}\right)\left(\sum_{n_{2}} \mathrm{e}^{-\beta \varepsilon \varepsilon_{n_{2}}^{(Y)}}\right)\left(\sum_{n_{3}} \mathrm{e}^{-\beta \varepsilon_{n_{3}}^{(Z)}}\right)
\end{aligned}
$$

That is,


\begin{equation*}
q^{\mathrm{T}}=q_{X}^{\mathrm{T}} \mathscr{q}_{Y}^{\mathrm{T}} q_{Z}^{\mathrm{T}} \tag{13B.9}
\end{equation*}


Equation 13B. 7 gives the partition function for translational motion in the $x$-direction. The only change for the other two directions is to replace the length $X$ by the lengths $Y$ or $Z$. Hence the partition function for motion in three dimensions is


\begin{equation*}
q^{\mathrm{T}}=\left(\frac{2 \pi m}{h^{2} \beta}\right)^{3 / 2} X Y Z=\frac{(2 \pi m k T)^{3 / 2}}{h^{3}} X Y Z \tag{13B.10a}
\end{equation*}


The product of lengths $X Y Z$ is the volume, $V$, of the container, so

\[
q^{\mathrm{T}}=\frac{V}{\Lambda^{3}} \quad \begin{align*}
& \text { Translational partition function }  \tag{13B.10b}\\
& \text { [three-dimensional] }
\end{align*}
\]

with $\Lambda$ defined in eqn 13B.7. As in the one-dimensional case, the partition function increases with the mass of the particle (as $\mathrm{m}^{3 / 2}$ ) and the volume of the container (as $V$ ); for a given mass and volume, the partition function increases with temperature (as $T^{3 / 2}$ ). Moreover, $q^{\mathrm{T}} \rightarrow \infty$ as $T \rightarrow \infty$ because there is no limit to the number of states that become accessible as the temperature is raised. Even at room temperature, $q^{T} \approx 2 \times 10^{28}$ for an $\mathrm{O}_{2}$ molecule in a vessel of volume $100 \mathrm{~cm}^{3}$.

\subsection*{Brief illustration 13B. 2}
To calculate the translational partition function of an $\mathrm{H}_{2}$ molecule confined to a $100 \mathrm{~cm}^{3}$ vessel at $25^{\circ} \mathrm{C}$, use $m=2.016 m_{\mathrm{u}}$. Then, from $\Lambda=h /(2 \pi m k T)^{1 / 2}$,

$$
\begin{aligned}
\Lambda & =\frac{6.626 \times 10^{-34} \overbrace{\mathrm{~J} \mathrm{~s}}^{\mathrm{kg} \mathrm{~m}^{2} \mathrm{~s}^{-2}}}{\{2 \pi \times\left(2.016 \times 1.6605 \times 10^{-27} \mathrm{~kg}\right) \times(1.381 \times 10^{-23} \underbrace{-1}_{\mathrm{kg} \mathrm{~m}^{2} \mathrm{~J}^{-2}}) \times(298 \mathrm{~K})\}^{1 / 2}} \\
& =7.12 \ldots \times 10^{-11} \mathrm{~m}
\end{aligned}
$$

Therefore,

$$
q^{\mathrm{T}}=\frac{1.00 \times 10^{-4} \mathrm{~m}^{3}}{\left(7.12 \ldots \times 10^{-11} \mathrm{~m}\right)^{3}}=2.77 \times 10^{26}
$$

About $10^{26}$ quantum states are thermally accessible, even at room temperature and for this light molecule. Many states are occupied if the thermal wavelength (which in this case is 71.2 pm ) is small compared with the linear dimensions of the container.

Equation 13B.10b can be interpreted in terms of the average separation, $d$, of the particles in the container. Because $q$ is the total number of accessible states, the average number of translational states per molecule is $q^{\mathrm{T}} / N$. For this quantity to be large, the condition $V / N \Lambda^{3} \gg 1$ must be met. However, $V / N$ is the volume occupied by a single particle, and therefore the average separation of the particles is $d=(V / N)^{1 / 3}$. The condition for there being many states available per molecule is therefore $d^{3} / \Lambda^{3} \gg 1$, and therefore $d \gg \Lambda$. That is, for eqn 13B.10b to be valid, the average separation of the particles must be much greater than their thermal wavelength. For $1 \mathrm{~mol} \mathrm{H}_{2}$ molecules at 1 bar and 298 K , the average separation is 3 nm , which is significantly larger than their thermal wavelength ( 71.2 pm ).

The validity of eqn 13B.10b can be expressed in a different way by noting that the approximations that led to it are valid if many states are occupied, which requires $V / \Lambda^{3}$ to be large. That will be so if $\Lambda$ is small compared with the linear dimensions of the container. For $\mathrm{H}_{2}$ at $298 \mathrm{~K}, \Lambda=71 \mathrm{pm}$, which is far smaller than any conventional container is likely to be (but comparable to pores in zeolites or cavities in clathrates). For $\mathrm{O}_{2}$, a heavier molecule, $\Lambda=18 \mathrm{pm}$.

\subsection*{(b) The rotational contribution}
The energy levels of a linear rotor are $\varepsilon_{J}=h c \tilde{B} J(J+1)$, with $J=0$, $1,2, \ldots$ (Topic 11B). The state of lowest energy has zero energy, so no adjustment need be made to the energies given by this expression. Each level consists of $2 J+1$ degenerate states. Therefore, the partition function of a non-symmetrical (AB) linear rotor is


\begin{equation*}
q^{\mathrm{R}}=\sum_{J} \overbrace{(2 J+1)}^{g_{j}} \mathrm{e}^{-\beta \overbrace{h \bar{c} B(J+1)}^{\varepsilon_{j}}} \tag{13B.11}
\end{equation*}


The direct method of calculating $g^{R}$ is to substitute the experimental values of the rotational energy levels into this expression and to sum the series numerically. (The case of symmetrical $\mathrm{A}_{2}$ molecules is dealt with later.)

\subsection*{Example 13B.1 Evaluating the rotational partition}
 function explicitlyEvaluate the rotational partition function of ${ }^{1} \mathrm{H}^{35} \mathrm{Cl}$ at $25^{\circ} \mathrm{C}$, given that $\tilde{B}=10.591 \mathrm{~cm}^{-1}$.

Collect your thoughts You need to evaluate eqn 13B. 11 term by term, using $k T / h c=207.224 \mathrm{~cm}^{-1}$ at 298.15 K . The sum is readily evaluated by using mathematical software.

The solution To show how successive terms contribute, draw up the following table by using $h c \tilde{B} / k T=0.05111$ (Fig. 13B.6):

\begin{center}
\begin{tabular}{llllllll}
$J$ & 0 & 1 & 2 & 3 & 4 & $\ldots$ & 10 \\
$(2 J+1) \mathrm{e}^{-0.05111 J(J+1)}$ & 1 & 2.71 & 3.68 & 3.79 & 3.24 & $\ldots$ & 0.08 \\
\end{tabular}
\end{center}

The sum required by eqn 13B.11 (the sum of the numbers in the second row of the table) is 19.9 , hence $q^{\mathrm{R}}=19.9$ at this temperature. Taking $J$ up to 50 gives $q^{R}=19.903$.\\
\includegraphics[max width=\textwidth, center]{2025_02_27_d4b95d9718f0b9e2e6b9g-574}

Figure 13B. 6 The contributions to the rotational partition function of an HCl molecule at $25^{\circ} \mathrm{C}$. The vertical axis is the value of $(2 J+1) \mathrm{e}^{-\beta h \bar{\beta} /(J+1)}$. Successive terms (which are proportional to the populations of the levels) pass through a maximum because the population of individual states decreases exponentially, but the degeneracy of the levels increases with J.

Comment. Notice that about ten $J$-levels are significantly populated but the number of populated states is larger on account of the $(2 J+1)$-fold degeneracy of each level.

Self-test 13B. 1 Evaluate the rotational partition function for ${ }^{1} \mathrm{H}^{35} \mathrm{Cl}$ at $0^{\circ} \mathrm{C}$.

9で8I :»วмsuV

At room temperature, $k T / h c \approx 200 \mathrm{~cm}^{-1}$. The rotational constants of many molecules are close to $1 \mathrm{~cm}^{-1}$ (Table 11C.1) and often smaller (though the very light $\mathrm{H}_{2}$ molecule, for which $\tilde{B}=60.9 \mathrm{~cm}^{-1}$, is one important exception). It follows that many rotational levels are populated at normal temperatures. When that is the case, explicit expressions for the rotational partition function can be derived.

Consider a linear rotor. If many rotational states are occupied and $k T$ is much larger than the separation between neighbouring states, the sum that defines the partition function can be approximated by an integral:

$$
q^{\mathrm{R}}=\int_{0}^{\infty}(2 J+1) \mathrm{e}^{-\beta h c \tilde{B}(J+1)} \mathrm{d} J
$$

This integral can be evaluated without much effort by making the substitution $x=\beta h c \tilde{B} J(J+1)$, so that $\mathrm{d} x / \mathrm{d} J=\beta h c \tilde{B}(2 J+1)$ and therefore $(2 J+1) \mathrm{d} J=\mathrm{d} x / \beta h c \tilde{B}$. Then

$$
q^{\mathrm{R}}=\frac{1}{\beta h c \tilde{B}} \overbrace{\int_{0}^{\infty} \mathrm{e}^{-x} \mathrm{~d} x}^{\text {Integral E.1: }}=\frac{1}{\beta h c \tilde{B}}
$$

which (because $\beta=1 / k T$ ) is

\[
q^{\mathrm{R}}=\frac{k T}{h c \tilde{B}} \quad \begin{align*}
& \text { Rotational partition function }  \tag{13B.12a}\\
& {[\text { unsymmetrical linear molecule] }}
\end{align*}
\]

A similar but more elaborate approach can be used for nonlinear molecules.

\subsection*{How is that done? 13B. 2 Deriving an expression for the}
 rotational partition function of a nonlinear moleculeConsider a symmetric rotor (Topic 11B) for which the energy levels are

$$
E_{J, K, M_{J}}=h c \tilde{B} J(J+1)+h c(\tilde{A}-\tilde{B}) K^{2}
$$

with $J=0,1,2, \ldots, K=J, J-1, \ldots,-J$, and $M_{J}=J, J-1, \ldots,-J$. Instead of considering these ranges, the same values can be covered by allowing $K$ to range from $-\infty$ to $\infty$, with $J$ confined to $|K|,|K|+1, \ldots, \infty$ for each value of $K$ (Fig. 13B.7).

\subsection*{Step 1 Write an expression for the sum over energy states}
Because the energy is independent of $M_{p}$, and there are $2 J+1$ values of $M_{J}$ for each value of $J$, each value of $J$ is $(2 J+1)$-fold degenerate. It follows that the partition function

$$
\mathfrak{q}=\sum_{J=0}^{\infty} \sum_{K=-J}^{J} \sum_{M_{J}=-J}^{J} \mathrm{e}^{-\beta E_{J, K, M_{J}}}
$$

\begin{center}
\includegraphics[max width=\textwidth]{2025_02_27_d4b95d9718f0b9e2e6b9g-575}
\end{center}

Figure 13B. 7 The calculation of the rotational partition function includes a contribution, indicated by the circles, for all possible combinations of $J$ and $K$. The sum is formed either (a) by allowing $J$ to take the values $0,1,2, \ldots$ and then for each $J$ allowing $K$ to range from $J$ to $-J$, or (b) by allowing $K$ to range from $-\infty$ to $\infty$, and for each value of $K$ allowing $J$ to take the values $|K|,|K|+1, \ldots, \infty$. The tinted areas show in (a) the sum from $K=-2$ to +2 for $J=2$, and in (b) the sum with $J=2,3 \ldots$ for $K=-2$.\\
can be written equivalently as

$$
\begin{aligned}
\mathscr{G} & =\sum_{J=0}^{\infty} \sum_{K=-J}^{J}(2 J+1) \mathrm{e}^{-\beta E_{J, K, M_{J}}}=\sum_{K=-\infty}^{\infty} \sum_{J=|K|}^{\infty}(2 J+1) \mathrm{e}^{-\beta E_{J, K, M_{J}}} \\
& =\sum_{K=-\infty}^{\infty} \mathrm{e}^{-h c \beta(\tilde{A}-\tilde{B}) K^{2}} \sum_{J=K \mid}^{\infty}(2 J+1) \mathrm{e}^{-h c \tilde{\beta} /(J+1)}
\end{aligned}
$$

Step 2 Convert the sums to integrals\\
As for linear molecules, assume that the temperature is so high that numerous states are occupied, in which case the sums may be approximated by integrals. Then

$$
\mathcal{q}=\int_{-\infty}^{\infty} \mathrm{e}^{-h c \beta(\tilde{A}-\tilde{B}) K^{2}} \int_{|K|}^{\infty}(2 J+1) \mathrm{e}^{-h c \beta \tilde{B} /(J+1)} \mathrm{d} J \mathrm{~d} K
$$

Step 3 Evaluate the integrals\\
You should recognize the integral over $J$ as the integral of the derivative of a function, which is the function itself, so

$$
\int_{|K|}^{\infty}(2 J+1) \mathrm{e}^{-h c \tilde{B}(J+1)} \mathrm{d} J=\frac{1}{h c \beta \tilde{B}} \mathrm{e}^{-h c \tilde{B} \tilde{K} \mid(|| |+1)} \approx \frac{1}{h c \beta \tilde{B}} \mathrm{e}^{-h c \tilde{\beta} \tilde{B} K^{2}}
$$

Use this result in the integral over $J$ in Step 2:

$$
\begin{aligned}
\mathscr{G} & =\frac{1}{h c \beta \tilde{B}} \int_{-\infty}^{\infty} \mathrm{e}^{-h c \beta(\tilde{A}-\tilde{B}) K^{2}} \mathrm{e}^{-h c c \tilde{B} K^{2}} \mathrm{~d} K=\frac{1}{h c \beta \tilde{B}} \overbrace{-\infty}^{\text {Integral G. } 1} \mathrm{e}^{-h c \beta \tilde{\beta} K^{2}} \mathrm{~d} K \\
& =\frac{1}{h c \beta \tilde{B}}\left(\frac{\pi}{h c \beta \tilde{A}}\right)^{1 / 2}
\end{aligned}
$$

now rearrange this expression into

$$
q=\frac{1}{(h c \beta)^{3 / 2}}\left(\frac{\pi}{\tilde{A} \tilde{B}^{2}}\right)^{1 / 2}=\left(\frac{k T}{h c}\right)^{3 / 2}\left(\frac{\pi}{\tilde{A} \tilde{B}^{2}}\right)^{1 / 2}
$$

For an asymmetric rotor with its three moments of inertia, one of the $\tilde{B}$ is replaced by $\tilde{C}$, to give

$$
\left.-q^{\mathrm{R}}=\left(\frac{k T}{h c}\right)^{3 / 2}\left(\frac{\pi}{\tilde{A} \tilde{B} \tilde{C}}\right)^{1 / 2} \right\rvert\, \begin{aligned}
& \text { Rotational partition function } \\
& \text { [nonlinear molecule] }
\end{aligned}
$$

\subsection*{Brief illustration 13B.3}
For ${ }^{1} \mathrm{H}^{35} \mathrm{Cl}$ at 298.15 K , use $k T / h c=207.224 \mathrm{~cm}^{-1}$ and $\tilde{B}=10.591$ $\mathrm{cm}^{-1}$. Then

$$
q^{\mathrm{R}}=\frac{k T}{h c \tilde{B}}=\frac{207.224 \mathrm{~cm}^{-1}}{10.591 \mathrm{~cm}^{-1}}=19.57
$$

The value is in good agreement with the exact value (19.903) and obtained with much less effort.

A useful way of expressing the temperature above which eqns 13B.12a and 13B.12b are valid is to introduce the characteristic rotational temperature, $\theta^{\mathrm{R}}=h c \tilde{B} / k$. Then 'high temperature' means $T \gg \theta^{\mathrm{R}}$ and under these conditions the rotational partition function of a linear molecule is simply $T / \theta^{\mathrm{R}}$. Some typical values of $\theta^{\mathrm{R}}$ are given in Table 13B.1. The value for ${ }^{1} \mathrm{H}_{2}(87.6 \mathrm{~K})$ is abnormally high, so the approximation must be used carefully for this molecule. However, before using eqn 13B.12a for symmetrical molecules, such as $\mathrm{H}_{2}$, read on (to eqn 13B.13a).

The general conclusion at this stage is that\\
Molecules with large moments of inertia (and hence small rotational constants and low characteristic rotational temperatures) have large rotational partition functions.

A large value of $q^{R}$ reflects the closeness in energy (compared with $k T$ ) of the rotational levels in large, heavy molecules, and the large number of rotational states that are accessible at normal temperatures.

It is important not to include too many rotational states in the sum that defines the partition function. For a homonuclear diatomic molecule or a symmetrical linear molecule

Table 13B. 1 Rotational temperatures of diatomic molecules*

\begin{center}
\begin{tabular}{lc}
\hline
 & $\theta^{\mathrm{R}} / \mathrm{K}$ \\
\hline
${ }^{1} \mathrm{H}_{2}$ & 87.6 \\
${ }^{1} \mathrm{H}^{35} \mathrm{Cl}$ & 15.2 \\
${ }^{14} \mathrm{~N}_{2}$ & 2.88 \\
${ }^{35} \mathrm{Cl}_{2}$ & 0.351 \\
\hline
\end{tabular}
\end{center}

\footnotetext{\begin{itemize}
  \item More values are given in the Resource section, Table 11C.1.
\end{itemize}
}
(such as $\mathrm{CO}_{2}$ or $\mathrm{HC} \equiv \mathrm{CH}$ ), a rotation through $180^{\circ}$ results in an indistinguishable state of the molecule. Hence, the number of thermally accessible states is only half the number that can be occupied by a heteronuclear diatomic molecule, where rotation through $180^{\circ}$ does result in a distinguishable state. Therefore, for a symmetrical linear molecule,\\
\$\$

q\^{}\{\textbackslash mathrm\{R\}\}=\textbackslash frac\{k T\}\{2 h c \textbackslash tilde\{B\}\}=\textbackslash frac\{T\}\{2 \textbackslash theta\^{}\{\textbackslash mathrm\{R\}\}\} \textbackslash quad \begin{align*}
& \text { Rotational partition function }  \tag{13B.13a}\\
& {[\text { symmetrical linear rotor }]}
\end{align*}

\$\$

The equations for symmetrical and non-symmetrical molecules can be combined into a single expression by introducing the symmetry number, $\sigma$, which is the number of indistinguishable orientations of the molecule. Then

\[
q^{\mathrm{R}}=\frac{T}{\sigma \theta^{\mathrm{R}}} \quad \begin{align*}
& \text { Rotational partition function }  \tag{13B.13b}\\
& {[\text { linear rotor }]}
\end{align*}
\]

For a heteronuclear diatomic molecule $\sigma=1$; for a homonuclear diatomic molecule or a symmetrical linear molecule, $\sigma=2$. The formal justification of this rule depends on assessing the role of the Pauli principle.

\subsection*{How is that done? 13B. 3 \\
 Identifying the origin of the}
 symmetry numberThe Pauli principle forbids the occupation of certain states. It is shown in Topic 11B, for example, that ${ }^{1} \mathrm{H}_{2}$ may occupy rotational states with even $J$ only if its nuclear spins are paired (para-hydrogen), and odd $J$ states only if its nuclear spins are parallel (ortho-hydrogen). In ortho- $\mathrm{H}_{2}$ there are three nuclear spin states for each value of $J$ (because there are three 'parallel' spin states of the two nuclei); in para- $\mathrm{H}_{2}$ there is just one nuclear spin state for each value of $J$.

To set up the rotational partition function and take into account the Pauli principle, note that 'ordinary' molecular hydrogen is a mixture of one part para- $\mathrm{H}_{2}$ (with only its even- $J$ rotational states occupied) and three parts ortho $-\mathrm{H}_{2}$ (with only its odd- $J$ rotational states occupied). Therefore, the average partition function for each molecule is

$$
q^{\mathrm{R}}=\frac{1}{4} \sum_{\text {even } J}(2 J+1) \mathrm{e}^{-\beta h \tilde{B} J(J+1)}+\frac{3}{4} \sum_{\text {odd } J}(2 J+1) \mathrm{e}^{-\beta h \tilde{\tilde{B}}(J+1)}
$$

The odd $-J$ states are three times more heavily weighted than the even $-J$ states (Fig. 13B.8). The illustration shows that approximately the same answer would be obtained for the partition function (the sum of all the populations) if each $J$ term contributed half its normal value to the sum. That is, the last equation can be approximated as

$$
q^{\mathrm{R}}=\frac{1}{2} \sum_{J}(2 J+1) \mathrm{e}^{-\beta h c \tilde{B} J(J+1)}
$$

and this approximation is very good when many terms contribute (at high temperatures, $T \gg \theta^{\mathrm{R}}$ ). At such high temperatures the sum can be approximated by the integral that led to\\
\includegraphics[max width=\textwidth, center]{2025_02_27_d4b95d9718f0b9e2e6b9g-577(7)}

Figure 13B. 8 The values of the individual terms $(2 J+1) \mathrm{e}^{-\beta h c \bar{\beta}(J+1)}$ contributing to the mean partition function of a 3:1 mixture of ortho- and para $-\mathrm{H}_{2}$. The partition function is the sum of all these terms. At high temperatures, the sum is approximately equal to the sum of the terms over all values of $J$, each with a weight of $\frac{1}{2}$. This sum is indicated by the curve.\\
eqn 13B.12a. Therefore, on account of the factor of $\frac{1}{2}$ in this expression, the rotational partition function for 'ordinary' molecular hydrogen at high temperatures is one-half this value, as in eqn 13B.13a.

The same type of argument may be used for linear symmetrical molecules in which identical bosons are interchanged by rotation (such as $\mathrm{CO}_{2}$ ). As pointed out in Topic 11B, if the nuclear spin of the bosons is 0 , then only even $-J$ states are admissible. Because only half the rotational states are occupied, the rotational partition function is only half the value of the sum obtained by allowing all values of $J$ to contribute (Fig. 13B.9).

The same care must be exercised for other types of symmetrical molecules, and for a nonlinear molecule eqn 13B.12b is replaced by

\[
q^{\mathrm{R}}=\frac{1}{\sigma}\left(\frac{k T}{h c}\right)^{3 / 2}\left(\frac{\pi}{\tilde{A} \tilde{B} \tilde{C}}\right)^{1 / 2} \quad \begin{align*}
& \text { Rotational partition }  \tag{13B.14}\\
& \text { function } \\
& \text { [nonlinear rotor] }
\end{align*}
\]

\begin{center}
\includegraphics[max width=\textwidth]{2025_02_27_d4b95d9718f0b9e2e6b9g-577(1)}
\end{center}

Figure 13B.9 The values of the individual terms contributing to the rotational partition function of $\mathrm{CO}_{2}$. Only states with even $J$ values are allowed. The full line shows the smoothed, averaged contributions of the levels.

Table 13B. 2 Symmetry numbers of molecules*

\begin{center}
\begin{tabular}{lc}
\hline
 & $\sigma$ \\
\hline
${ }^{1} \mathrm{H}_{2}$ & 2 \\
${ }^{1} \mathrm{H}^{2} \mathrm{H}$ & 1 \\
$\mathrm{NH}_{3}$ & 3 \\
$\mathrm{C}_{6} \mathrm{H}_{6}$ & 12 \\
\hline
\end{tabular}
\end{center}

\begin{itemize}
  \item More values are given in the Resource section, Table 11C.1.
\end{itemize}

Some typical values of the symmetry numbers are given in Table 13B.2. To see how group theory is used to identify the value of the symmetry number, see Integrated activity I13.1; the following Brief illustration outlines the approach.

\subsection*{Brief illustration 13B. 4}
The value $\sigma\left(\mathrm{H}_{2} \mathrm{O}\right)=2$ reflects the fact that a $180^{\circ}$ rotation about the bisector of the $\mathrm{H}-\mathrm{O}-\mathrm{H}$ angle interchanges two indistinguishable atoms. In $\mathrm{NH}_{3}$, there are three indistinguishable orientations around the axis shown in (1). For $\mathrm{CH}_{4}$, any of three $120^{\circ}$ rotations about any of its four C-H bonds leaves the molecule in an indistinguishable state (2), so the symmetry number is $3 \times 4=12$.\\
\includegraphics[max width=\textwidth, center]{2025_02_27_d4b95d9718f0b9e2e6b9g-577(2)}\\
\includegraphics[max width=\textwidth, center]{2025_02_27_d4b95d9718f0b9e2e6b9g-577}

For benzene, any of six orientations around the axis perpendicular to the plane of the molecule leaves it apparently unchanged (Fig. 13B.10), as does a rotation of $180^{\circ}$ around any of six axes in the plane of the molecule (three of which pass through C atoms diametrically opposite across the ring and the remaining three pass through the mid-points of $\mathrm{C}-\mathrm{C}$ bonds on opposite sides of the ring).\\
\includegraphics[max width=\textwidth, center]{2025_02_27_d4b95d9718f0b9e2e6b9g-577(6)}\\
\includegraphics[max width=\textwidth, center]{2025_02_27_d4b95d9718f0b9e2e6b9g-577(3)}\\
\includegraphics[max width=\textwidth, center]{2025_02_27_d4b95d9718f0b9e2e6b9g-577(4)}\\
\includegraphics[max width=\textwidth, center]{2025_02_27_d4b95d9718f0b9e2e6b9g-577(5)}

Figure 13B. 10 The 12 equivalent orientations of a benzene molecule that can be reached by pure rotations and give rise to a symmetry number of 12 . The six pale colours are the underside of the hexagon after that face has been rotated into view.

\subsection*{(c) The vibrational contribution}
The vibrational partition function of a molecule is calculated by substituting the measured vibrational energy levels into the definition of $q^{v}$, and summing them numerically. However, provided it is permissible to assume that the vibrations are harmonic, there is a much simpler way. In that case, the vibrational energy levels form a uniform ladder of separation $h c \tilde{v}$ (Topics 7E and 11C), which is exactly the problem treated in Brief illustration 13B. 1 and lead to eqn 13B.2a. Therefore that result can be used by setting $\varepsilon=h c \tilde{v}$, giving

\[
\mathcal{q}^{\mathrm{V}}=\frac{1}{1-\mathrm{e}^{-\beta h c \tilde{v}}} \quad \begin{align*}
& \text { Vibrational partition function }  \tag{13B.15}\\
& {[\text { harmonic approximation] }}
\end{align*}
\]

This function is plotted in Fig. 13B. 11 (which is essentially the same as Fig. 13B.1). Similarly, the population of each state is given by eqn 13B.2b.

\subsection*{Brief illustration 13B. 5}
To calculate the partition function of $\mathrm{I}_{2}$ molecules at 298.15 K note from Table 11C. 1 that their vibrational wavenumber is $214.6 \mathrm{~cm}^{-1}$. Then, because at $298.15 \mathrm{~K}, k T / h c=207.224 \mathrm{~cm}^{-1}$,

$$
\beta h c \tilde{v}=\frac{h c \tilde{v}}{k T}=\frac{214.6 \mathrm{~cm}^{-1}}{207.244 \mathrm{~cm}^{-1}}=1.035 \ldots
$$

Then it follows from eqn 13B. 15 that

$$
q^{\mathrm{V}}=\frac{1}{1-\mathrm{e}^{-1.035 \ldots \ldots}}=1.55
$$

From this value it can be inferred that only the ground and first excited states are significantly populated.\\
\includegraphics[max width=\textwidth, center]{2025_02_27_d4b95d9718f0b9e2e6b9g-578}

Figure 13B. 11 The vibrational partition function of a molecule in the harmonic approximation. Note that the partition function is proportional to the temperature when the temperature is high ( $T \gg \theta^{V}$ ).

In a polyatomic molecule each normal mode (Topic 11D) has its own partition function (provided the anharmonicities are so small that the modes are independent). The overall vibrational partition function is the product of the individual partition functions, and so $q^{\mathrm{V}}=q^{\mathrm{V}}(1) q^{\mathrm{V}}(2) \ldots$, where $q^{\mathrm{V}}(K)$ is the partition function for normal mode $K$ and is calculated by direct summation of the observed spectroscopic levels.

\subsection*{Example 13B. 2 \\
 Calculating a vibrational partition function}
The wavenumbers of the three normal modes of $\mathrm{H}_{2} \mathrm{O}$ are $3656.7 \mathrm{~cm}^{-1}, 1594.8 \mathrm{~cm}^{-1}$, and $3755.8 \mathrm{~cm}^{-1}$. Evaluate the vibrational partition function at 1500 K .

Collect your thoughts You need to use eqn 13B. 15 for each mode, and then form the product of the three contributions. At $1500 \mathrm{~K}, k T / h c=1042.6 \mathrm{~cm}^{-1}$.

The solution Draw up the following table displaying the contributions of each mode:

\begin{center}
\begin{tabular}{llrl}
\hline
Mode & \multicolumn{1}{l}{} & \multicolumn{1}{l}{} & \multicolumn{1}{l}{2} \\
\multicolumn{1}{l}{} &  &  &  \\
\hline
$\tilde{v} / \mathrm{cm}^{-1}$ & 3656.7 & 1594.8 & \multicolumn{1}{l}{3755.8} \\
$h c \tilde{v} / k T$ & 3.507 & 1.530 & 3.602 \\
$q^{V}$ & 1.031 & 1.276 & 1.028 \\
\hline
\end{tabular}
\end{center}

The overall vibrational partition function is therefore

$$
q^{v}=1.031 \times 1.276 \times 1.028=1.352
$$

The three normal modes of $\mathrm{H}_{2} \mathrm{O}$ are at such high wavenumbers that even at 1500 K most of the molecules are in their vibrational ground state.

Comment. There may be so many normal modes in a large molecule that their overall contribution may be significant even though each mode is not appreciably excited. For example, a nonlinear molecule containing 10 atoms has $3 N-6=24$ normal modes (Topic 11D). If a value of about 1.1 is assumed for the vibrational partition function of one normal mode, the overall vibrational partition function is about $\mathcal{q}^{\mathrm{V}} \approx(1.1)^{24}=9.8$, which indicates significant overall vibrational excitation relative to a smaller molecule, such as $\mathrm{H}_{2} \mathrm{O}$.

Self-test 13B. 2 Repeat the calculation for $\mathrm{CO}_{2}$ at the same temperature. The vibrational wavenumbers are $1388 \mathrm{~cm}^{-1}$, $667.4 \mathrm{~cm}^{-1}$, and $2349 \mathrm{~cm}^{-1}$, the second being the doublydegenerate bending mode which contributes twice to the overall vibrational partition function.

6L'9: :дммиин

In many molecules the vibrational wavenumbers are so great that $\beta h c \tilde{v}>1$. For example, the lowest vibrational wavenumber of $\mathrm{CH}_{4}$ is $1306 \mathrm{~cm}^{-1}$, so $\beta h c \tilde{v}=6.3$ at room temperature. Most C-H stretches normally lie in the range $2850-2960 \mathrm{~cm}^{-1}$,

Table 13B.3 Vibrational temperatures of diatomic molecules*

\begin{center}
\begin{tabular}{lr}
\hline
 & $\theta^{\mathrm{V}} / \mathrm{K}$ \\
\hline
${ }^{1} \mathrm{H}_{2}$ & 6332 \\
${ }^{1} \mathrm{H}^{35} \mathrm{Cl}$ & 4304 \\
${ }^{14} \mathrm{~N}_{2}$ & 3393 \\
${ }^{35} \mathrm{Cl}_{2}$ & 805 \\
\hline
\end{tabular}
\end{center}

\begin{itemize}
  \item More values are given in the Resource section, Table 11C. 1\\
so for them $\beta h c \tilde{v} \approx 14$. In these cases, $\mathrm{e}^{-\beta h c \tilde{v}}$ in the denominator of $q^{V}$ is very close to zero (e.g. $e^{-6.3}=0.002$ ), and the vibrational partition function for a single mode is very close to 1 $\left(\mathcal{q}^{\mathrm{V}}=1.002\right.$ when $\left.\beta h c \tilde{v}=6.3\right)$, implying that only the lowest level is significantly occupied.
\end{itemize}

Now consider the case of modes with such low vibrational frequencies that $\beta h c \tilde{v} \ll 1$. When this condition is satisfied, the partition function may be approximated by expanding the exponential ( $\mathrm{e}^{x}=1+x+\cdots$ ):

$$
q^{\mathrm{v}}=\frac{1}{1-\mathrm{e}^{-\beta h c \tilde{v}}}=\frac{1}{1-(1-\beta h c \tilde{v}+\cdots)}
$$

That is, for low-frequency modes at high temperatures,

\[
q^{\mathrm{V}} \approx \frac{k T}{h c \tilde{v}} \quad \begin{align*}
& \text { Vibrational partition function }  \tag{13B.16}\\
& \text { [high-temperature } \\
& \text { approximation] }
\end{align*}
\]

The temperatures for which eqn $13 B .16$ is valid can be expressed in terms of the characteristic vibrational temperature, $\theta^{\mathrm{V}}=h c \tilde{v} / k$ (Table 13B.3). The value for $\mathrm{H}_{2}(6332 \mathrm{~K})$ is abnormally high because the atoms are so light and the vibrational frequency is correspondingly high. In terms of the vibrational temperature, 'high temperature' means $T \gg \theta^{\vee}$, and when this condition is satisfied eqn 13B. 16 is valid and can be written $\mathscr{q}^{\mathrm{V}}=T / \theta^{\vee}$ (the analogue of the rotational expression).

\subsection*{(d) The electronic contribution}
Electronic energy separations from the ground state are usually very large, so for most cases $q^{E}=1$ because only the ground state is occupied. An important exception arises in the case of atoms and molecules having electronically degenerate ground states, in which case $q^{\mathrm{E}}=g^{\mathrm{E}}$, where $g^{\mathrm{E}}$ is the degeneracy of the electronic ground state. Alkali metal atoms, for example, have doubly degenerate ground states (corresponding to the two orientations of their electron spin), so $\mathfrak{q}^{\mathrm{E}}=2$.

\subsection*{Brief illustration 13B. 6}
Some atoms and molecules have degenerate ground states and low-lying electronically excited degenerate states. An example is NO, which has a configuration of the form ... $\pi^{1}$ (Topic 11F).

The energy of the two degenerate states in which the orbital and spin momenta are parallel (giving the ${ }^{2} \Pi_{3 / 2}$ term, Fig. 13B.12) is slightly greater than that of the two degenerate states in which they are antiparallel (giving the ${ }^{2} \Pi_{1 / 2}$ term). The separation, which arises from spin-orbit coupling, is only $121 \mathrm{~cm}^{-1}$.\\
\includegraphics[max width=\textwidth, center]{2025_02_27_d4b95d9718f0b9e2e6b9g-579(1)}

Figure 13B. 12 The doubly-degenerate ground electronic level of NO (with the spin and orbital angular momentum around the axis in opposite directions) and the doubly-degenerate first excited level (with the spin and orbital momenta parallel). The upper level is thermally accessible at room temperature.

If the energies of the two levels are denoted as $E_{1 / 2}=0$ and $E_{3 / 2}=\varepsilon$, then the partition function is

$$
\mathcal{q}^{\mathrm{E}}=\sum_{\text {levels } i} g_{i} \mathrm{e}^{-\beta \varepsilon_{i}}=2+2 \mathrm{e}^{-\beta \varepsilon}
$$

This function is plotted in Fig. 13B.13. At $T=0, q^{\mathrm{E}}=2$, because only the doubly degenerate ground state is accessible. At high temperatures, $q^{\mathrm{E}}$ approaches 4 because all four states are accessible. At $25^{\circ} \mathrm{C}, 9^{\mathrm{E}}=3.1$.\\
\includegraphics[max width=\textwidth, center]{2025_02_27_d4b95d9718f0b9e2e6b9g-579}

Figure 13B. 13 The variation with temperature of the electronic partition function of an NO molecule. Note that the curve resembles that for a two-level system (Fig.13B.4), but rises from 2 (the degeneracy of the lower level) and approaches 4 (the total number of states) at high temperatures.

\subsection*{Checklist of concepts}
$\square$ 1. The molecular partition function is an indication of the number of thermally accessible states at the temperature of interest.2. If the energy of a molecule is given by the sum of contributions from different modes, then the molecular partition function is a product of the partition functions for each of the modes.3. The symmetry number takes into account the number of indistinguishable orientations of a symmetrical molecule.\\
4. The vibrational partition function of a molecule is found by evaluating the contribution from each normal mode treated as a harmonic oscillator.\\
5. Because electronic energy separations from the ground state are usually very large, in most cases the electronic partition function is equal to the degeneracy of the electronic ground state.

\subsection*{Checklist of equations}
\begin{center}
\begin{tabular}{|c|c|c|c|}
\hline
Property & Equation & Comment & Equation number \\
\hline
\multirow[t]{2}{*}{Molecular partition function} & \( \mathcal{G}=\sum_{\text {states } i} \mathrm{e}^{-\beta \varepsilon_{i}} \) & Definition, independent molecules & 13B.1a \\
\hline
 & \( G=\sum_{\text {levels } i} g_{i} \mathrm{e}^{-\beta \varepsilon_{i}} \) & Definition, independent molecules & 13B.1b \\
\hline
Uniform ladder & $\mathcal{q}=1 /\left(1-\mathrm{e}^{-\beta \varepsilon}\right)$ &  & 13B.2a \\
\hline
Two-level system & $q=1+\mathrm{e}^{-\beta \varepsilon}$ &  & 13B.3a \\
\hline
Thermal wavelength & $\Lambda=h /(2 \pi m k T)^{1 / 2}$ &  & 13B. 7 \\
\hline
Translation & $q^{\text {T }}=V / \Lambda^{3}$ &  & 13B.10b \\
\hline
\multirow[t]{2}{*}{Rotation} & $q^{\mathrm{R}}=k T / \sigma h c \tilde{B}$ & $T \gg \theta^{\mathrm{R}}$, linear rotor & 13B.13a \\
\hline
 & $q^{\mathrm{R}}=(1 / \sigma)(k T / h c)^{3 / 2}(\pi / \tilde{A} \tilde{B} \tilde{C})^{1 / 2}$ & $T \gg \theta^{\mathrm{R}}$, nonlinear rotor, $\theta^{\mathrm{R}}=h c \tilde{X} / k, \tilde{X}=\tilde{A}, \tilde{B}$, or $\tilde{C}$ & 13B. 14 \\
\hline
Vibration & $q^{\mathrm{V}}=1 /\left(1-\mathrm{e}^{-\beta h c \bar{v}}\right)$ & Harmonic approximation, $\theta^{\mathrm{v}}=h c \tilde{v} / k$ & 13B. 15 \\
\hline
\end{tabular}
\end{center}

\section*{TOPIC 13C Molecular energies}
\subsection*{Why do you need to know this material?}
For statistical thermodynamics to be useful, you need to know how to extract thermodynamic information from a partition function.

\subsection*{What is the key idea?}
The average energy of a molecule in a collection of independent molecules can be calculated from the molecular partition function alone.

\subsection*{What do you need to know already?}
You need to know how to calculate a molecular partition function from calculated or spectroscopic data (Topic 13B) and its significance as a measure of the number of accessible states. The Topic also draws on expressions for the rotational and vibrational energies of molecules (Topics 11B-11D).

A partition function in statistical thermodynamics is like a wavefunction in quantum mechanics. A wavefunction contains all the dynamical information about a system; a partition function contains all the thermodynamic information about a system. As in quantum mechanics, it is important to know how to extract that information. One of the simplest thermodynamic properties is the mean energy, and the equations are simplest for a system composed of non-interacting molecules.

\subsection*{13C. 1 The basic equations}
Consider a collection of $N$ molecules that do not interact with one another. Any member of the collection can exist in a state $i$ of energy $\varepsilon_{i}$ measured from the lowest energy state of the molecule. The mean energy of a molecule, $\langle\varepsilon\rangle$, relative to its energy in its ground state, is the total energy of the collection, $E$, divided by the total number of molecules:


\begin{equation*}
\langle\varepsilon\rangle=\frac{E}{N}=\frac{1}{N} \sum_{i} N_{i} \varepsilon_{i} \tag{13C.1}
\end{equation*}


where $N_{i}$ is the population of state $i$. In Topic 13A it is shown that the overwhelmingly most probable population of a state in a collection at a temperature $T$ is given by the Boltzmann distribution, eqn 13A.10b $\left(N_{i} / N=(1 / \mathcal{q}) \mathrm{e}^{-\beta \varepsilon_{i}}\right)$, so


\begin{equation*}
\langle\varepsilon\rangle=\frac{1}{q} \sum_{i} \varepsilon_{i} \mathrm{e}^{-\beta \varepsilon_{i}} \tag{13C.2}
\end{equation*}


with $\beta=1 / k T$. This expression can be manipulated into a form involving only $\mathfrak{q}$. First note that

$$
\varepsilon_{i} \mathrm{e}^{-\beta \varepsilon_{i}}=-\frac{\mathrm{d}}{\mathrm{~d} \beta} \mathrm{e}^{-\beta \varepsilon_{i}}
$$

It follows that


\begin{equation*}
\langle\varepsilon\rangle=-\frac{1}{\mathfrak{G}} \sum_{i} \frac{\mathrm{~d}}{\mathrm{~d} \beta} \mathrm{e}^{-\beta \varepsilon_{i}}=-\frac{1}{\mathfrak{g}} \frac{\mathrm{~d}}{\mathrm{~d} \beta} \overbrace{\sum_{i} \mathrm{e}^{-\beta \varepsilon_{i}}}^{\mathcal{G}}=-\frac{1}{\mathfrak{g}} \frac{\mathrm{~d} \mathfrak{G}}{\mathrm{~d} \beta} \tag{13C.3}
\end{equation*}


Several points need to be made in relation to eqn 13C.3. Because $\varepsilon_{0} \equiv 0$, (all energies are measured from the lowest available level), $\langle\varepsilon\rangle$ is the value of the energy relative to the actual ground-state energy. If the lowest energy of the molecule is in fact $\varepsilon_{\mathrm{gs}}$ rather than 0 , then the true mean energy is $\varepsilon_{\mathrm{gs}}+$ $\langle\varepsilon\rangle$. For instance, for a harmonic oscillator, $\varepsilon_{\mathrm{gs}}$ is the zero-point energy, $\frac{1}{2} h c \tilde{v}$. Secondly, because the partition function might depend on variables other than the temperature (e.g. the volume), the derivative with respect to $\beta$ in eqn 13C. 3 is actually a partial derivative with these other variables held constant. The complete expression relating the molecular partition function to the mean energy of a molecule is therefore

$$
\langle\varepsilon\rangle=\varepsilon_{\mathrm{gs}}-\frac{1}{\mathscr{q}}\left(\frac{\partial \mathscr{G}}{\partial \beta}\right)_{V} \quad \text { Mean molecular energy }
$$

(13C.4a)

An equivalent form is obtained by noting that $\mathrm{d} x / x=\mathrm{d} \ln x$ :

$$
\langle\varepsilon\rangle=\varepsilon_{\mathrm{gs}}-\left(\frac{\partial \ln _{\mathscr{q}}}{\partial \beta}\right)_{V} \quad \text { Mean molecular energy }
$$

(13C.4b)

These two equations confirm that only the partition function (as a function of temperature) is needed in order to calculate the mean energy.

\subsection*{Brief illustration 13C. 1}
If a molecule has only two available states, one at 0 and the other at an energy $\varepsilon$, its partition function is

$$
\mathfrak{q}=1+\mathrm{e}^{-\beta \varepsilon}
$$

Therefore, the mean energy of a collection of these molecules at a temperature $T$ is

$$
\langle\varepsilon\rangle=-\frac{1}{1+\mathrm{e}^{-\beta \varepsilon}} \frac{\mathrm{d}\left(1+\mathrm{e}^{-\beta \varepsilon}\right)}{\mathrm{d} \beta}=\frac{\varepsilon \mathrm{e}^{-\beta \varepsilon}}{1+\mathrm{e}^{-\beta \varepsilon}}=\frac{\varepsilon}{\mathrm{e}^{\beta \varepsilon}+1}
$$

This function is plotted in Fig. 13C.1. Notice how the mean energy is zero at $T=0$, when only the lower state (at the zero of energy) is occupied, and rises to $\frac{1}{2} \varepsilon$ as $T \rightarrow \infty$, when the two states become equally populated.\\
\includegraphics[max width=\textwidth, center]{2025_02_27_d4b95d9718f0b9e2e6b9g-582}

Figure 13C. 1 The variation with temperature of the mean energy of a two-level system. The graph on the left shows the slow rise away from zero energy at low temperatures; the slope of the graph at $T=0$ is 0 . The graph on the right shows the slow approach to 0.5 as $T \rightarrow \infty$ as both states become equally populated.

\subsection*{13C. 2 Contributions of the fundamental modes of motion}
The remainder of this Topic explains how to write expressions for the contributions to the energy of three fundamental types of motion, namely translation (T), rotation (R), and vibration (V). It also shows how to incorporate the contribution of the electronic states of molecules (E) and electron spin (S).

\subsection*{(a) The translational contribution}
For a one-dimensional container of length $X$, the partition function is $\mathscr{G}_{X}^{\mathrm{T}}=X / \Lambda$ with $\Lambda=h /(2 \pi m / \beta)^{1 / 2}$ (Topic 13B, with 'constant volume $V$ ' replaced by 'constant length $X$ '). The\\
partition function can be written as constant $\times \beta^{-1 / 2}$, a form convenient for the calculation of the average energy using eqn 13C.4b:

$$
\begin{aligned}
\left\langle\varepsilon_{X}^{\mathrm{T}}\right\rangle & =-\left(\frac{\partial \ln \mathscr{G}}{\partial \beta}\right)_{X}=-\left(\frac{\partial \ln \left(\text { constant } \times \beta^{-1 / 2}\right)}{\partial \beta}\right)_{X} \\
& =-\left(\frac{\partial\left(\ln (\text { constant })+\ln \beta^{-1 / 2}\right)}{\partial \beta}\right)_{X} \\
& =-\left(\frac{\partial\left(-\frac{1}{2} \ln \beta\right)}{\partial \beta}\right)_{X}=\frac{1}{2 \beta}
\end{aligned}
$$

That is,

\[
\left\langle\varepsilon_{X}^{\mathrm{T}}\right\rangle=\frac{1}{2} k T \quad \begin{align*}
& \text { Mean translational energy }  \tag{13C.5a}\\
& \text { [one dimension] }
\end{align*}
\]

For a molecule free to move in three dimensions, the analogous calculation leads to

$$
\left\langle\varepsilon^{\mathrm{T}}\right\rangle=\frac{3}{2} k T \quad \begin{aligned}
& \text { Mean translational energy } \\
& \text { [three dimensions] }
\end{aligned}
$$

(13C.5b)

\subsection*{(b) The rotational contribution}
The mean rotational energy of a linear molecule is obtained from the rotational partition function (eqn 13B.11):

$$
q^{\mathrm{R}}=\sum_{J}(2 J+1) \mathrm{e}^{-\beta h \tilde{\tilde{B}}(J+1)}
$$

When the temperature is not high (in the sense that it is not true that $\left.T \gg \theta^{\mathrm{R}}=h c \tilde{B} / k\right)$ the series must be summed term by term, which for a heteronuclear diatomic molecule or other non-symmetrical linear molecule gives

$$
q^{\mathrm{R}}=1+3 \mathrm{e}^{-2 \beta h c \tilde{B}}+5 \mathrm{e}^{-6 \beta h c \tilde{B}}+\cdots
$$

Hence, because

$$
\frac{\mathrm{d} \mathcal{q}^{\mathrm{R}}}{\mathrm{~d} \beta}=-h c \tilde{B}\left(6 \mathrm{e}^{-2 \beta h c \tilde{B}}+30 \mathrm{e}^{-6 \beta h c \tilde{B}}+\cdots\right)
$$

( $q^{R}$ is independent of $V$, so the partial derivative has been replaced by a complete derivative) it follows that

$$
\begin{array}{r}
\left\langle\varepsilon^{\mathrm{R}}\right\rangle=-\frac{1}{q^{\mathrm{R}}} \frac{\mathrm{~d} q^{\mathrm{R}}}{\mathrm{~d} \beta}=\frac{h c \tilde{B}\left(6 \mathrm{e}^{-2 \beta h c \tilde{B}}+30 \mathrm{e}^{-6 \beta h c \tilde{B}}+\cdots\right)}{1+3 \mathrm{e}^{-2 \beta h c \tilde{B}}+5 \mathrm{e}^{-6 \beta h c \tilde{B}}+\cdots} \\
\text { Mean rotational energy } \\
\text { [unsymmetrical linear molecule] }
\end{array}
$$

This ungainly function is plotted in Fig. 13C.2. At high temperatures $\left(T \gg \theta^{\mathrm{R}}\right), q^{\mathrm{R}}$ is given by eqn $13 \mathrm{~B} .13 \mathrm{~b}\left(q^{\mathrm{R}}=T / \sigma \theta^{\mathrm{R}}\right)$ in the form $q^{\mathrm{R}}=1 / \sigma \beta h c \tilde{B}$, where $\sigma=1$ for a heteronuclear diatomic molecule. It then follows that\\
\includegraphics[max width=\textwidth, center]{2025_02_27_d4b95d9718f0b9e2e6b9g-583}

Figure 13C. 2 The variation with temperature of the mean rotational energy of an unsymmetrical linear rotor. At high temperatures ( $T \gg \theta^{R}$ ), the energy is proportional to the temperature, in accord with the equipartition theorem.

$$
\left\langle\varepsilon^{\mathrm{R}}\right\rangle=-\frac{1}{q^{\mathrm{R}}} \frac{\mathrm{~d} q^{\mathrm{R}}}{\mathrm{~d} \beta}=-\sigma \beta h c \tilde{B} \frac{\mathrm{~d}}{\mathrm{~d} \beta}\left(\frac{1}{\sigma \beta h c \tilde{B}}\right)=-\beta \frac{\mathrm{d}^{-1 / \beta^{2}}}{\mathrm{~d} \beta} \frac{1}{\beta}
$$

and therefore that

$$
\left\langle\varepsilon^{\mathrm{R}}\right\rangle=\frac{1}{\beta}=k T
$$

Mean rotational energy [linear molecule, $T \gg \theta^{\text {R }}$ ]\\
(13C.6b)

The high-temperature result, which is valid when many rotational states are occupied, is also in agreement with the equipartition theorem (The chemist's toolkit 7 in Topic 2A), because the classical expression for the energy of a linear rotor is $E_{\mathrm{k}}=\frac{1}{2} I_{\perp} \omega_{\mathrm{a}}^{2}+\frac{1}{2} I_{\perp} \omega_{\mathrm{b}}^{2}$ and therefore has two quadratic contributions. (There is no rotation around the line of atoms.) It follows from the equipartition theorem that the mean rotational energy is $2 \times \frac{1}{2} k T=k T$.

\subsection*{Brief illustration 13C. 2}
To estimate the mean energy of a nonlinear molecule in the high-temperature limit recognize that its rotational kinetic energy (the only contribution to its rotational energy) is $E_{\mathrm{k}}=\frac{1}{2} I_{a} \omega_{a}^{2}+\frac{1}{2} I_{b} \omega_{b}^{2}+\frac{1}{2} I_{c} \omega_{c}^{2}$. As there are three quadratic contributions, its mean rotational energy is $\frac{3}{2} k T$. The molar contribution is $\frac{3}{2} R T$. At $25^{\circ} \mathrm{C}$, this contribution is $3.7 \mathrm{~kJ} \mathrm{~mol}^{-1}$, the same as the translational contribution, giving a total of $7.4 \mathrm{~kJ} \mathrm{~mol}^{-1}$. A monatomic gas has no rotational contribution.

\subsection*{(c) The vibrational contribution}
The vibrational partition function in the harmonic approximation is given in eqn 13B.15 $\left(\mathcal{q}^{\mathrm{v}}=1 /\left(1-\mathrm{e}^{-\beta h c \tilde{v}}\right)\right)$. Because $\mathcal{q}^{\mathrm{V}}$ is independent of the volume, it follows that


\begin{equation*}
\frac{\mathrm{d} \mathcal{q}^{\mathrm{v}}}{\mathrm{~d} \beta}=\frac{\mathrm{d}}{\mathrm{~d} \beta}\left(\frac{1}{1-\mathrm{e}^{-\beta h c \tilde{v}}}\right)=-\frac{h c \tilde{v} \mathrm{e}^{-\beta h c \tilde{v}}}{\left(1-\mathrm{e}^{-\beta h c \tilde{v}}\right)^{2}} \tag{13C.7}
\end{equation*}


and hence

$$
\begin{aligned}
\left\langle\varepsilon^{\mathrm{V}}\right\rangle & =-\frac{1}{q^{\mathrm{V}}} \frac{\mathrm{~d} \mathcal{q}^{\mathrm{V}}}{\mathrm{~d} \beta}=\left(1-\mathrm{e}^{-\beta h c \tilde{v}}\right) \frac{h c \tilde{v} \mathrm{e}^{-\beta h c \tilde{v}}}{\left(1-\mathrm{e}^{-\beta h c \tilde{v}}\right)^{2}} \\
& =\frac{h c \tilde{v} \mathrm{e}^{-\beta h c \tilde{v}}}{1-\mathrm{e}^{-\beta h c \tilde{v}}}
\end{aligned}
$$

The final result, after multiplying the numerator and denominator by $e^{\beta h c \tilde{v}}$ is

\[
\left\langle\varepsilon^{\mathrm{v}}\right\rangle=\frac{h c \tilde{v}}{\mathrm{e}^{\beta h c \tilde{v}}-1} \quad \begin{align*}
& \text { Mean vibrational energy }  \tag{13C.8}\\
& \text { [harmonic approximation] }
\end{align*}
\]

The zero-point energy, $\frac{1}{2} h c \tilde{v}$, can be added to the right-hand side if the mean energy is to be measured from 0 rather than the lowest attainable level (the zero-point level). The variation of the mean energy with temperature is illustrated in Fig. 13C.3. At high temperatures, when $T \gg \theta^{V}=h c \tilde{v} / k$ or $\beta h c \tilde{v} \ll 1$, the exponential function can be expanded ( $\mathrm{e}^{x}=1+x+\cdots$ ) and all but the leading terms discarded. This approximation leads to

\[
\left\langle\varepsilon^{\mathrm{V}}\right\rangle=\frac{h c \tilde{v}}{(1+\beta h c \tilde{v}+\cdots)-1} \approx \frac{1}{\beta}=k T \quad \begin{align*}
& \text { Mean vibrational }  \tag{13C.9}\\
& \text { energy } \\
& {\left[T \gg \theta^{v}\right]}
\end{align*}
\]

This result is in agreement with the value predicted by the classical equipartition theorem, because the energy of a onedimensional oscillator is $E=\frac{1}{2} m v_{x}^{2}+\frac{1}{2} k_{\mathrm{f}} x^{2}$ and the mean energy of each quadratic term is $\frac{1}{2} k T$. Bear in mind, however, that the condition $T \gg \theta^{\mathrm{V}}$ is rarely satisfied.\\
\includegraphics[max width=\textwidth, center]{2025_02_27_d4b95d9718f0b9e2e6b9g-583(1)}

Figure 13C. 3 The variation with temperature of the mean vibrational energy of a molecule in the harmonic approximation. At high temperatures ( $T \gg \theta^{v}$ ), the energy is proportional to the temperature, in accord with the equipartition theorem.

\subsection*{Brief illustration 13C. 3}
To calculate the mean vibrational energy of $I_{2}$ molecules at 298.15 K note from Table 11C. 1 that their vibrational wavenumber is $214.6 \mathrm{~cm}^{-1}$. At $298.15 \mathrm{~K}, k T / h c=207.224 \mathrm{~cm}^{-1}$,

$$
\beta h c \tilde{v}=\frac{h c \tilde{v}}{k T}=\frac{214.6 \mathrm{~cm}^{-1}}{207.244 \mathrm{~cm}^{-1}}=1.035 \ldots
$$

Because $\beta h c \tilde{v}=\theta^{v} / T$ is not large compared with 1 (implying that $T$ is not high compared with $\theta^{\vee}$ ), equipartition cannot be used. It then follows from eqn 13C. 8 that

$$
\left\langle\varepsilon^{\mathrm{v}}\right\rangle / h c=\frac{214.6 \mathrm{~cm}^{-1}}{\mathrm{e}^{1.035 \ldots}-1}=118.1 \mathrm{~cm}^{-1}
$$

The addition of the zero-point energy (corresponding to $\frac{1}{2} \times$ $214.6 \mathrm{~cm}^{-1}$ ) increases this value to $225.4 \mathrm{~cm}^{-1}$.

When there are several normal modes that can be treated as harmonic, the overall vibrational partition function is the product of each individual partition function, and the total mean vibrational energy is the sum of the mean energy of each mode.

\subsection*{(d) The electronic contribution}
In most cases of interest, the electronic states of atoms and molecules are so widely separated that only the electronic ground state is occupied. Because all energies are measured from the ground state of each mode, it follows that

$$
\left\langle\varepsilon^{\mathrm{E}}\right\rangle=0
$$

Mean electronic energy\\
(13C.10)

In certain cases, there are thermally accessible states at the temperature of interest. In that case, the partition function and hence the mean electronic energy are best calculated by direct summation over the available states. Care must be taken to take any degeneracies into account, as illustrated in the following example.

\subsection*{Example 13C. 1 \\
 Calculating the electronic contribution to the energy}
A certain atom has a doubly-degenerate electronic ground state and a fourfold degenerate excited state at $\varepsilon / h c=\tilde{v}=$ $600 \mathrm{~cm}^{-1}$ above the ground state. What is its mean electronic energy at $25^{\circ} \mathrm{C}$, expressed as a wavenumber?

Collect your thoughts You need to write the expression for the partition function at a general temperature $T$ (in terms of $\beta$ ) and then derive the mean energy by using eqn 13C.3. Doing so involves differentiating the partition function with respect\\
to $\beta$. Finally, substitute the data. Use $\varepsilon=h c \tilde{v},\left\langle\varepsilon^{\mathrm{E}}\right\rangle=h c\left\langle\tilde{v}^{\mathrm{E}}\right\rangle$, and (from inside the front cover), $k T / h c=207.224 \mathrm{~cm}^{-1}$ at $25^{\circ} \mathrm{C}$.

The solution The partition function is $q^{\mathrm{E}}=2+4 \mathrm{e}^{-\beta \varepsilon}$. The mean energy is therefore

$$
\begin{aligned}
\left\langle\varepsilon^{\mathrm{E}}\right\rangle & =-\frac{1}{\mathcal{G}^{\mathrm{E}}} \frac{\mathrm{~d} \mathcal{q}^{\mathrm{E}}}{\mathrm{~d} \beta}=-\frac{1}{2+4 \mathrm{e}^{-\beta \varepsilon}} \frac{\mathrm{d}}{\mathrm{~d} \beta} \overbrace{\left(2+4 \mathrm{e}^{-\beta \varepsilon}\right)}^{-4 \varepsilon e^{-\beta \varepsilon}} \\
& =\frac{4 \varepsilon \mathrm{e}^{-\beta \varepsilon}}{2+4 \mathrm{e}^{-\beta \varepsilon}}=\frac{\varepsilon}{\frac{1}{2} \mathrm{e}^{\beta \varepsilon}+1}
\end{aligned}
$$

Expressed as a wavenumber the mean energy is $\left\langle\varepsilon^{\mathrm{E}}\right\rangle / h c=\left\langle\tilde{v}^{\mathrm{E}}\right\rangle$

$$
\left\langle\tilde{v}^{\mathrm{E}}\right\rangle=\frac{\tilde{v}}{\frac{1}{2} e^{h=h c \tilde{v} / k T}+1}
$$

From the data,

$$
\left\langle\tilde{v}^{\mathrm{E}}\right\rangle=\frac{600 \mathrm{~cm}^{-1}}{\frac{1}{2} \mathrm{e}^{\left(600 \mathrm{~cm}^{-1}\right) /\left(207.224 \mathrm{~cm}^{-1}\right)}+1}=59.7 \mathrm{~cm}^{-1}
$$

Self-test 13C. 1 Repeat the problem for an atom that has a threefold degenerate ground state and a sevenfold degenerate excited state $400 \mathrm{~cm}^{-1}$ above.

\subsection*{(e) The spin contribution}
An electron spin in a magnetic field $\mathscr{B}$ has two possible energy states that depend on its orientation as expressed by the magnetic quantum number $m_{s}$, and which are given by

$$
E_{m_{s}}=g_{e} \mu_{\mathrm{B}} B m_{s} \quad \text { Electron spin energies }
$$

where $\mu_{\mathrm{B}}$ is the Bohr magneton (see inside front cover) and $g_{\mathrm{e}}=2.0023$. These energies are discussed in more detail in Topic 12A. The lower state has $m_{s}=-\frac{1}{2}$, so the two energy levels available to the electron lie (according to the convention that $\left.\varepsilon_{0} \equiv 0\right)$ at $\varepsilon_{-1 / 2}=0$ and at $\varepsilon_{+1 / 2}=g_{e} \mu_{\mathrm{B}} \mathcal{B}$. The spin partition function is therefore


\begin{equation*}
q^{\mathrm{S}}=\sum_{m_{s}} \mathrm{e}^{-\beta \varepsilon_{m_{s}}}=1+\mathrm{e}^{-\beta g_{c} \mu_{\mathrm{B}} \mathcal{B}} \quad \text { Spin partition function } \tag{13C.12}
\end{equation*}


The mean energy of the spin is therefore

$$
\begin{aligned}
&\left\langle\varepsilon^{\mathrm{s}}\right\rangle=-\frac{1}{q^{\mathrm{s}}} \frac{\mathrm{~d} \mathscr{q}^{\mathrm{s}}}{\mathrm{~d} \beta}=-\frac{1}{1+\mathrm{e}^{-\beta_{g} \mu_{\mathrm{B}} \mathcal{B}}} \frac{\mathrm{~d}}{\mathrm{~d} \beta}\left(1+\mathrm{e}^{-\beta g_{c} \mu_{\mathrm{B}} \mathcal{B}}\right) \\
&-g_{\mathrm{e}} \mu_{\mathrm{B}} \mathfrak{B e} e^{-\beta g_{e} \mu_{\mathrm{B}} \mathcal{B}} \\
&=\frac{g_{\mathrm{e}} \mu_{\mathrm{B}} B \mathrm{~B}^{-\beta_{g_{c}} \mu_{\mathrm{B}} \mathcal{B}}}{1+\mathrm{e}^{-\beta g_{c} \mu_{\mathrm{B}} \mathcal{B}}}
\end{aligned}
$$

The final expression, after multiplying the numerator and denominator by $\mathrm{e}^{\beta_{g_{e}} \mu_{B} \mathcal{B}^{\mathcal{B}}}$, is

$$
\left\langle\varepsilon^{\mathrm{S}}\right\rangle=\frac{g_{\mathrm{e}} \mu_{\mathrm{B}} \mathcal{B}}{\mathrm{e}^{\beta g_{\mathrm{e}} \mu_{\mathrm{B}} \mathcal{B}}+1}
$$

Mean spin energy (13C.13)

This function is essentially the same as that plotted in Fig. 13C.1.

\subsection*{Brief illustration 13C. 4}
Suppose a collection of radicals is exposed to a magnetic field of 2.5 T (T denotes tesla) at $25^{\circ} \mathrm{C}$. With $\mu_{\mathrm{B}}=9.274 \times 10^{-24} \mathrm{~J} \mathrm{~T}^{-1}$,

$$
\begin{aligned}
& g_{e} \mu_{\mathrm{B}} \mathcal{B}=2.0023 \times\left(9.274 \times 10^{-24} \mathrm{JT}^{-1}\right) \times 2.5 \mathrm{~T}=4.6 \ldots \times 10^{-23} \mathrm{~J} \\
& \beta g_{\mathrm{e}} \mu_{\mathrm{B}} \mathcal{B}=\frac{2.0023 \times\left(9.274 \times 10^{-24} \mathrm{JT}^{-1}\right) \times(2.5 \mathrm{~T})}{\left(1.381 \times 10^{-23} \mathrm{JK}^{-1}\right) \times(298 \mathrm{~K})}=0.011 \ldots
\end{aligned}
$$

The mean energy is therefore

$$
\left\langle\varepsilon^{\mathrm{s}}\right\rangle=\frac{4.6 \ldots \times 10^{-23} \mathrm{~J}}{\mathrm{e}^{0.011 \ldots}+1}=2.3 \times 10^{-23} \mathrm{~J}
$$

This energy is equivalent to $14 \mathrm{~J} \mathrm{~mol}^{-1}$ (note joules, not kilojoules).

\subsection*{Checklist of concepts}
1. The mean molecular energy can be calculated from the molecular partition function.2. Individual contributions to the mean molecular energy from each mode of motion are calculated from the relevant partition functions,3. In the high temperature limit, the results obtained for the mean molecular energy are in accord with the equipartition principle.\subsection*{Checklist of equations}
\begin{center}
\begin{tabular}{llll}
\hline
Property & Equation & Comment & Equation number \\
\hline
Mean energy & $\langle\varepsilon\rangle=\varepsilon_{\mathrm{gs}}-(1 / \mathcal{q})(\partial \mathcal{q} / \partial \beta)_{V}$ & $\beta=1 / k T$ & 13 C .4 a \\
Translation & $\langle\varepsilon\rangle=\varepsilon_{\mathrm{gs}}-(\partial \ln \mathcal{q} / \partial \beta)_{V}$ & Alternative version & 13C.4b \\
Rotation & $\left\langle\varepsilon^{\mathrm{T}}\right\rangle=\frac{d}{2} k T$ & In $d$ dimensions, $d=1,2,3$ & 13 C .5 \\
Vibration & $\left\langle\varepsilon^{\mathrm{R}}\right\rangle=k T$ & Linear molecule, $T \gg \theta^{\mathrm{R}}$ & 13C.6b \\
 & $\left\langle\varepsilon^{\mathrm{V}}\right\rangle=h c \tilde{v} /\left(\mathrm{e}^{\beta h c \tilde{v}}-1\right)$ & Harmonic approximation & 13 C .8 \\
Spin & $\left\langle\varepsilon^{\mathrm{V}}\right\rangle=k T$ & $T \gg \theta^{\mathrm{V}}$ & Electron in a magnetic field \\
\hline
\end{tabular}
\end{center}

\section*{TOPIC 13D The canonical ensemble}
\subsection*{Why do you need to know this material?}
Whereas Topics 13B and 13C deal with independent molecules, in practice molecules do interact. Therefore, this material is essential for constructing models of real gases, liquids, and solids and of any system in which intermolecular interactions cannot be neglected.

\subsection*{What is the key idea?}
A system composed of interacting molecules is described in terms of a canonical partition function, from which its thermodynamic properties may be deduced.

\subsection*{What do you need to know already?}
The calculations here, which are not carried through in detail, are essentially the same as in Topic 13A. Calculations of mean energies are also essentially the same as in Topic 13C and are not repeated in detail.

The crucial concept needed in the treatment of systems of interacting particles, as in real gases and liquids, is the 'ensemble'. Like so many scientific terms, the term has basically its normal meaning of 'collection', but in statistical thermodynamics it has been sharpened and refined into a precise significance.

\subsection*{13D. 1 The concept of ensemble}
To set up an ensemble, take a closed system of specified volume, composition, and temperature, and think of it as replicated $\tilde{N}$ times (Fig. 13D.1). All the identical closed systems are regarded as being in thermal contact with one another, so they can exchange energy. The total energy of the ensemble is $\tilde{E}$. Because the members of the ensemble are all in thermal equilibrium with each other, they have the same temperature, $T$. The volume of each member of the ensemble is the same, so the energy levels available to the molecules are the same in each system, and each member contains the same number of molecules, so there is a fixed number of molecules to distribute within each system. This imaginary collection of replications of the actual system with a common temperature is called the canonical ensemble. ${ }^{1}$

\footnotetext{${ }^{1}$ The word 'canon' means 'according to a rule'.
}
\includegraphics[max width=\textwidth, center]{2025_02_27_d4b95d9718f0b9e2e6b9g-586}

Figure 13D. 1 A representation of the canonical ensemble, in this case for $\tilde{N}=20$. The individual replications of the actual system all have the same composition and volume. They are all in mutual thermal contact, and so all have the same temperature. Energy may be transferred between them as heat, and so they do not all have the same energy. The total energy of all 20 replications is a constant because the ensemble is isolated overall.

There are two other important types of ensembles. In the microcanonical ensemble the condition of constant temperature is replaced by the requirement that all the systems should have exactly the same energy: each system is individually isolated. In the grand canonical ensemble the volume and temperature of each system is the same, but they are open, which means that matter can be imagined as able to pass between them; the composition of each one may fluctuate, but now the property known as the chemical potential ( $\mu$, Topic 5A) is the same in each system. In summary:

\begin{center}
\begin{tabular}{ll}
\hline
Ensemble & Common properties \\
\hline
Microcanonical & $V, E, N$ \\
Canonical & $V, T, N$ \\
Grand canonical & $V, T, \mu$ \\
\hline
\end{tabular}
\end{center}

The microcanonical ensemble is the basis of the discussion in Topic 13A; the grand canonical ensemble will not be considered explicitly.

The important point about an ensemble is that it is a collection of imaginary replications of the system, so the number of members can be as large as desired; when appropriate, $\tilde{N}$ can be taken as infinite. The number of members of the ensemble in a state with energy $E_{i}$ is denoted $\tilde{N}_{i}$, and it is possible to speak of the configuration of the ensemble (by analogy with the configuration of the system used in Topic 13A) and\\
its weight, $\tilde{\mathcal{W}}$, the number of ways of achieving the configuration $\left\{\tilde{N}_{0}, \tilde{N}_{1}, \ldots\right\}$. Note that $\tilde{N}$ is unrelated to $N$, the number of molecules in the actual system; $\tilde{N}$ is the number of imaginary replications of that system. In summary:\\
$N$ is the number of molecules in the system, the same in each member of the ensemble.\\
$T$ is the common temperature of the members.\\
$E_{i}$ is the total energy of one member of the ensemble (the one labelled $i$ ).\\
$\tilde{N}_{i}$ is the number of replicas that have the energy $E_{i}$. $\tilde{E}$ is the total energy of the entire ensemble.\\
$\tilde{N}$ is the total number of replicas (the number of members of the ensemble).\\
$\tilde{\mathcal{W}}$ is the weight of the configuration $\left\{\tilde{N}_{0}, \tilde{N}_{1}, \ldots\right\}$.

\subsection*{(a) Dominating configurations}
Just as in Topic 13A, some of the configurations of the canonical ensemble are very much more probable than others. For instance, it is very unlikely that the whole of the total energy of the ensemble will accumulate in one system to give the configuration $\{\tilde{N}, 0,0, \ldots\}$. By analogy with the discussion in Topic 13 A , there is a dominating configuration, and thermodynamic properties can be calculated by taking the average over the ensemble using that single, most probable, configuration. In the thermodynamic limit of $\tilde{N} \rightarrow \infty$, this dominating configuration is overwhelmingly the most probable, and dominates its properties.

The quantitative discussion follows the argument in Topic 13A (leading to the Boltzmann distribution) with the modification that $N$ and $N_{i}$ are replaced by $\tilde{N}$ and $\tilde{N}_{i}$. The weight $\tilde{\mathcal{W}}$ of a configuration $\left\{\tilde{N}_{0}, \tilde{N}_{1}, \ldots\right\}$ is

$$
\tilde{\mathcal{W}}=\frac{\tilde{N}!}{\tilde{N}_{0}!\tilde{N}_{1}!\ldots}
$$

Weight (13D.1)

The configuration of greatest weight, subject to the constraints that the total energy of the ensemble is constant at $\tilde{E}$ and that the total number of members is fixed at $\tilde{N}$, is given by the canonical distribution:


\begin{equation*}
\frac{\tilde{N}_{i}}{\tilde{N}}=\frac{\mathrm{e}^{-\beta E_{i}}}{Q} \quad Q=\sum_{i} \mathrm{e}^{-\beta E_{i}} \quad \text { Canonical distribution } \tag{13D.2}
\end{equation*}


where the sum is over all members of the ensemble. The quantity $Q$, which is a function of the temperature, is called the canonical partition function, and $\beta=1 / k T$. Like the molecular partition function, the canonical partition function contains all the thermodynamic information about a system but allows for the possibility of interactions between the constituent molecules.

\subsection*{(b) Fluctuations from the most probable distribution}
The canonical distribution in eqn 13D. 2 is only apparently an exponentially decreasing function of the energy of the system. Just as the Boltzmann distribution gives the occupation of a single state of a molecule, the canonical distribution function gives the probability of occurrence of members in a single state $i$ of energy $E_{i}$. There may in fact be numerous states with almost identical energies. For example, in a gas the identities of the molecules moving slowly or quickly can change without necessarily affecting the total energy. The energy density of states, the number of states in an energy range divided by the width of the range (Fig. 13D.2), is a very sharply increasing function of energy. It follows that the probability of a member of an ensemble having a specified energy (as distinct from being in a specified state) is given by eqn 13D.2, a sharply decreasing function, multiplied by a sharply increasing function (Fig. 13D.3). Therefore, the overall distribution is a sharply peaked function. That is, most members of the ensemble have an energy very close to the mean value.\\
\includegraphics[max width=\textwidth, center]{2025_02_27_d4b95d9718f0b9e2e6b9g-587}

Figure 13D. 2 The energy density of states is the number of states in an energy range divided by the width of the range.\\
\includegraphics[max width=\textwidth, center]{2025_02_27_d4b95d9718f0b9e2e6b9g-587(1)}

Figure 13D. 3 To construct the form of the distribution of members of the canonical ensemble in terms of their energies, the probability that any one is in a state of given energy (eqn 13D.2) is multiplied by the number of states corresponding to that energy (a steeply rising function). The product is a sharply peaked function at the mean energy (here considerably magnified), which shows that almost all the members of the ensemble have that energy.

\subsection*{Brief illustration 13D. 1}
A function that increases rapidly is $x^{n}$, with $n$ large. A function that decreases rapidly is $\mathrm{e}^{-n x}$, once again, with $n$ large. The product of these two functions, normalized so that the maxima for different values of $n$ all coincide,

$$
f(x)=\mathrm{e}^{n} x^{n} \mathrm{e}^{-n x}
$$

is plotted for three values of $n$ in Fig. 13D.4. Note that the width of the product does indeed decrease as $n$ increases.\\
\includegraphics[max width=\textwidth, center]{2025_02_27_d4b95d9718f0b9e2e6b9g-588}

Figure 13D. 4 The product of the two functions discussed in Brief illustration 13D.1, for three different values of $n$.

\subsection*{13D.2 The mean energy of a system}
Just as the molecular partition function can be used to calculate the mean value of a molecular property, so the canonical partition function can be used to calculate the mean energy of an entire system composed of molecules which might or might not be interacting with one another. Thus, $Q$ is more general than $g$ because it does not assume that the molecules are independent. Therefore $Q$ can be used to discuss the properties of condensed phases and real gases where molecular interactions are important.

Because the total energy of the ensemble is $\tilde{E}$, and there are $\tilde{N}$ members, the mean energy of a member is $\langle E\rangle=\tilde{E} / \tilde{N}$. Because the fraction, $\tilde{p}_{i}$, of members of the ensemble in a state $i$ with energy $E_{i}$ is given by the analogue of eqn 13A.10b ( $p_{i}=\mathrm{e}^{-\beta \varepsilon_{i}} / g$ with $\left.p_{i}=N_{\mathrm{i}} / N\right)$ as


\begin{equation*}
\tilde{p}_{i}=\frac{\mathrm{e}^{-\beta E_{i}}}{Q} \tag{13D.3}
\end{equation*}


it follows that


\begin{equation*}
\langle E\rangle=\sum_{i} \tilde{p}_{i} E_{i}=\frac{1}{Q} \sum_{i} E_{i} \mathrm{e}^{-\beta E_{i}} \tag{13D.4}
\end{equation*}


By the same argument that led to eqn 13C.4a $(\langle\varepsilon\rangle=-(1 / q)$ $(\partial q / \partial \beta)_{V}$, when $\left.\varepsilon_{\mathrm{gs}} \equiv 0\right)$,

\[
\langle E\rangle=-\frac{1}{Q}\left(\frac{\partial Q}{\partial \beta}\right)_{V}=-\left(\frac{\partial \ln Q}{\partial \beta}\right)_{V} \quad \begin{align*}
& \text { Mean energy }  \tag{13D.5}\\
& \text { of a system }
\end{align*}
\]

As in the case of the mean molecular energy, the ground-state energy of the entire system must be added to this expression if it is not zero.

\subsection*{13D. 3 Independent molecules revisited}
When the molecules are in fact independent of each other, $Q$ can be shown to be related to the molecular partition function $q$.

\subsection*{How is that done? 13D. 1 Establishing the relation}
between $Q$ and $q$\\
There are two cases you need to consider because it turns out to be important to consider initially a system in which the particles are distinguishable (such as when they are at fixed locations in a solid) and then one in which they are indistinguishable (such as when they are in a gas and able to exchange places).

Step 1 Consider a system of independent, distinguishable molecules

The total energy of a collection of $N$ independent molecules is the sum of the energies of the molecules. Therefore, the total energy of a state $i$ of the system is written as

$$
E_{i}=\varepsilon_{i}(1)+\varepsilon_{i}(2)+\cdots+\varepsilon_{i}(N)
$$

In this expression, $\varepsilon_{i}(1)$ is the energy of molecule 1 when the system is in the state $i, \varepsilon_{i}(2)$ the energy of molecule 2 when the system is in the same state $i$, and so on. The canonical partition function is then

$$
Q=\sum_{i} \mathrm{e}^{-\beta \varepsilon_{i}(1)-\beta \varepsilon_{i}(2)-\cdots-\beta \varepsilon_{i}(N)}
$$

Provided the molecules are distinguishable (in a sense described below), the sum over the states of the system can be reproduced by letting each molecule enter all its own individual states. Therefore, instead of summing over the states $i$ of the system, sum over all the individual states $j$ of molecule 1 , all the states $j$ of molecule 2 , and so on. This rewriting of the original expression leads to

$$
Q=\overbrace{\left(\sum_{j} \mathrm{e}^{-\beta \varepsilon_{j}}\right)}^{q} \overbrace{\left(\sum_{j} \mathrm{e}^{-\beta \varepsilon_{j}}\right)}^{q} \overbrace{\left(\sum_{j} \mathrm{e}^{-\beta \varepsilon_{j}}\right)}^{q}
$$

and therefore to $Q=q^{N}$.

Step 2 Consider a system of independent, indistinguishable molecules

If all the molecules are identical and free to move through space, it is not possible to distinguish them and the relation $Q=q^{N}$ is not valid. Suppose that molecule 1 is in some state $a$, molecule 2 is in $b$, and molecule 3 is in $c$, then one member of the ensemble has an energy $E=\varepsilon_{a}+\varepsilon_{b}+\varepsilon_{c}$. This member, however, is indistinguishable from one formed by putting molecule 1 in state $b$, molecule 2 in state $c$, and molecule 3 in state $a$, or some other permutation. There are six such permutations in all, and $N$ ! in general. In the case of indistinguishable molecules, it follows that too many states have been counted in going from the sum over system states to the sum over molecular states, so writing $Q=q^{N}$ overestimates the value of $Q$. The detailed argument is quite involved, but at all except very low temperatures it turns out that the correction factor is $1 / N!$, so $Q=q^{N} / N!$.

Step 3 Summarize the results\\
It follows that:

For distinguishable independent molecules:


\begin{equation*}
Q=q^{N} \tag{13D.6a}
\end{equation*}


For indistinguishable independent molecules:\\
(13D.6b) $Q=q^{N} / N$ !

For molecules to be indistinguishable, they must be of the same kind: an Ar atom is never indistinguishable from a Ne atom. Their identity, however, is not the only criterion. Each identical molecule in a crystal lattice, for instance, can be 'named' with a set of coordinates. Identical molecules in a lattice can therefore be treated as distinguishable because their sites are distinguishable, and eqn 13D.6a can be used. On the other hand, identical molecules in a gas are free to move to different locations, and there is no way of keeping track of the identity of a given molecule; therefore eqn 13D.6b must be used.

\subsection*{Brief illustration 13D. 2}
For a gas of $N$ indistinguishable molecules, $Q=q^{N} / N$ !. The energy of the system is therefore

$$
\begin{aligned}
\langle E\rangle & =-\left(\frac{\partial \ln \left(\mathcal{q}^{N} / N!\right)}{\partial \beta}\right)_{V}=-\left(\frac{\partial\left(\ln \mathcal{q}^{N}-\ln N!\right)}{\partial \beta}\right)_{V}=-\left(\frac{\partial(N \ln \mathcal{q})}{\partial \beta}\right)_{V} \\
& =-N\left(\frac{\partial \ln \mathcal{q}}{\partial \beta}\right)_{V}=N\langle\varepsilon\rangle
\end{aligned}
$$

That is, the mean energy of the gas is $N$ times the mean energy of a single molecule.

\subsection*{13D.4 The variation of the energy with volume}
When there are interactions between molecules, the energy of a collection depends on the average distance between them, and therefore on the volume that a fixed number occupy. This dependence on volume is particularly important for the discussion of real gases (Topic 1C).

To discuss the dependence of the energy on the volume at constant temperature it is necessary to evaluate $(\partial\langle E\rangle / \partial V)_{T}$. (In Topics 2D and 3E, this quantity is identified as the 'internal pressure' of a gas and denoted $\pi_{T}$.) To proceed, substitute eqn 13D. 5 and obtain


\begin{equation*}
\left(\frac{\partial\langle E\rangle}{\partial V}\right)_{T}=-\left(\frac{\partial}{\partial V}\left(\frac{\partial \ln Q}{\partial \beta}\right)_{V}\right)_{T} \tag{13D.7}
\end{equation*}


If the gas were perfect, $Q=q^{N} / N$ !, with $q$ the partition function for the translational and any internal (such as rotational) modes. If the gas is monatomic, only the translational mode is present and $\mathfrak{q}=V / \Lambda^{3}$ with $\Lambda=h /(2 \pi m k T)^{1 / 2}$ and the canonical partition function would be $V^{N} / \Lambda^{3 N} N$ !. The presence of interactions is taken into account by replacing $V^{N} / N$ ! by a factor called the configuration integral, $Z$, which depends on the intermolecular potentials (don't confuse this $Z$ with the compression factor $Z$ in Topic 1C), and writing


\begin{equation*}
Q=\frac{Z}{\Lambda^{3 N}} \tag{13D.8}
\end{equation*}


It then follows that


\begin{align*}
&\left(\frac{\partial\langle E\rangle}{\partial V}\right)_{T}=-\left(\frac{\partial}{\partial V}\left(\frac{\partial \ln \left(Z / \Lambda^{3 N}\right)}{\partial \beta}\right)_{V}\right)_{T} \\
&=-\left(\frac{\partial}{\partial V}\left(\frac{\partial \ln Z}{\partial \beta}\right)_{V}\right)_{T}-\left(\frac{\partial}{\partial V}\left(\frac{\partial \ln \left(1 / \Lambda^{3 N}\right)}{\partial \beta}\right)_{V}\right)_{T} \\
&=-(\underbrace{\partial^{2} f / \partial x \partial y=\partial^{2} f / \partial y \partial x \text { for the second (blue) term }} \\
&=-\left(\frac{\partial}{\partial V}\left(\frac{\partial \ln Z}{\partial \beta}\right)_{V}\right)_{T}-(\frac{\partial}{\partial \beta} \overbrace{\left(\frac{\partial \ln \left(1 / \Lambda^{3 N}\right)}{\partial V}\right)_{T}}^{\partial \ln y / \partial x=(1 / y \partial y / \partial x)}  \tag{13D.9}\\
&=-\left(\frac{\partial}{\partial V}\left(\frac{\partial \ln Z}{\partial \beta}\right)_{V}\right)_{T} \\
&=-\left(\frac{\partial}{\partial V} \frac{1}{Z}\left(\frac{\partial Z}{\partial \beta}\right)_{V}\right)_{T}
\end{align*}


In the third line, $\Lambda$ is independent of volume, so its derivative with respect to volume is zero.

For a real gas of atoms (for which the intermolecular interactions are isotropic), $Z$ is related to the total potential energy $E_{\mathrm{p}}$ of interaction of all the particles, which depends on all their relative locations, by


\begin{equation*}
Z=\frac{1}{N!} \int \mathrm{e}^{-\beta E_{\mathrm{p}}} \mathrm{~d} \tau_{1} \mathrm{~d} \tau_{2} \ldots \mathrm{~d} \tau_{N} \quad \text { Configuration integral } \tag{13D.10}
\end{equation*}


where $\mathrm{d} \tau_{i}$ is the volume element for atom $i$ and the integration is over all the variables. The physical origin of this term is that the probability of occurrence of each arrangement of molecules possible in the sample is given by a Boltzmann distribution in which the exponent is given by the potential energy corresponding to that arrangement.

\subsection*{Brief illustration 13D. 3}
Equation 13D. 10 is very difficult to manipulate in practice, even for quite simple intermolecular potentials, except for a perfect gas for which $E_{\mathrm{P}}=0$. In that case, the exponential function becomes 1 and

$$
Z=\frac{1}{N!} \int \mathrm{d} \tau_{1} \mathrm{~d} \tau_{2} \ldots \mathrm{~d} \tau_{N}=\frac{1}{N!}\left(\int \mathrm{d} \tau\right)^{N}=\frac{V^{N}}{N!}
$$

just as it should be for a perfect gas.

If the potential energy has the form of a central hard sphere surrounded by a shallow attractive well (Fig. 13D.5), then detailed calculation, which is too involved to reproduce here (see A deeper look 7 on the website of this text), leads to


\begin{equation*}
\left(\frac{\partial\langle E\rangle}{\partial V}\right)_{T}=\frac{a n^{2}}{V^{2}} \quad \text { Attractive potential } \tag{13D.11}
\end{equation*}


where $n$ is the amount of molecules present in the volume $V$ and $a$ is a constant that is proportional to the area under the attractive part of the potential. In Example 3E. 2 of Topic 3E exactly the same expression (in the form $\pi_{T}=a n^{2} / V^{2}$ ) was derived from the van der Waals equation of state. The conclusion at this point is that if there are attractive interactions between molecules in a gas, then its energy increases as it expands isothermally (because $(\partial\langle E\rangle / \partial V)_{T}>0$, and the slope of $\langle E\rangle$ with respect to $V$ is positive). The energy rises because, at greater average separations, the molecules spend less time in regions where they interact favourably.\\
\includegraphics[max width=\textwidth, center]{2025_02_27_d4b95d9718f0b9e2e6b9g-590}

Figure 13D. 5 The intermolecular potential energy of molecules in a real gas can be modelled by a central hard sphere that determines the van der Waals parameter $b$ surrounded by a shallow attractive well that determines the parameter $a$. As mentioned in the text, calculations of the canonical partition function based on this are consistent with the van der Waals equation of state (Topic 1C).

\subsection*{Checklist of concepts}
$\square$ 1. The canonical ensemble is a collection of imaginary replications of the actual system with a common temperature and number of particles.2. The canonical distribution gives the most probable number of members of the ensemble with a specified total energy.\\
3. The mean energy of the members of the ensemble can be calculated from the canonical partition function.

\subsection*{Checklist of equations}
\begin{center}
\begin{tabular}{llll}
\hline
Property & Equation & Comment & Equation number \\
\hline
Canonical partition function & $Q=\sum_{i} \mathrm{e}^{-\beta E_{i}}$ & Definition & 13 D .2 \\
Canonical distribution & $N_{i} / \tilde{N}=\mathrm{e}^{-\beta E_{i}} / Q$ &  & 13 D .2 \\
Mean energy & $\langle E\rangle=-(1 / Q)(\partial Q / \partial \beta)_{V}=-(\partial \ln Q / \partial \beta)_{V}$ &  & 13 D .5 \\
Canonical partition function & $Q=Z / \Lambda^{3 N}$ & 13 D .8 &  \\
Configuration integral & $Z=(1 / N!) \int \mathrm{e}^{-\beta E_{\mathrm{F}}} \mathrm{d} \tau_{1} \mathrm{~d} \tau_{2} \ldots \mathrm{~d} \tau_{N}$ & Isotropic interaction & 13 D .10 \\
Variation of mean energy with volume & $(\partial\langle E\rangle / \partial V)_{T}=a n^{2} / V^{2}$ & Potential energy as specified in Fig. 13D.5 & 13 D .11 \\
\hline
\end{tabular}
\end{center}

\section*{TOPIC 13E The internal energy and the entropy }
\subsection*{Why do you need to know this material?}
The importance of this discussion is the insight that a molecular interpretation provides into thermodynamic properties.

\subsection*{What is the key idea?}
The partition function contains all the thermodynamic information about a system and thus provides a bridge between spectroscopy and thermodynamics.

What do you need to know already?\\
You need to know how to calculate a molecular partition function from structural data (Topic 13B); you should also be familiar with the concepts of internal energy (Topic 2A) and entropy (Topic 3A). This Topic makes use of the calculations of mean molecular energies in Topic 13C.

Any thermodynamic function can be obtained once the partition function is known. The two fundamental properties of thermodynamics are the internal energy, $U$, and the entropy, $S$. Once these two properties have been calculated, it is possible to turn to the derived functions, such as the Gibbs energy, $G$ (Topic 13F) and all the chemically interesting properties that stem from them.

\subsection*{13E. 1 The internal energy}
The first example of the importance of the molecular partition function, $q$, is the derivation of an expression for the internal energy.

\subsection*{(a) The calculation of internal energy}
It is established in Topic 13C that the mean energy of a molecule in a system composed of independent molecules is related to the molecular partition function by


\begin{equation*}
\langle\varepsilon\rangle=-\frac{1}{q}\left(\frac{\partial \mathscr{G}}{\partial \beta}\right)_{V} \tag{13E.1}
\end{equation*}


with $\beta=1 / k T$. The total energy of a system composed of $N$ molecules is therefore $N\langle\varepsilon\rangle$. This total energy is the energy above the ground state (the calculation of $q$ is based on the convention $\varepsilon_{0} \equiv 0$ ) and so the internal energy is $U(T)=U(0)+N\langle\varepsilon\rangle$, where $U(0)$ is the internal energy when only the ground state is occupied, which is the case at $T=0$. It follows that the internal energy is related to the molecular partition function by

$$
U(T)=U(0)+N\langle\varepsilon\rangle=U(0)-\frac{N}{q}\left(\frac{\partial \mathfrak{q}}{\partial \beta}\right)_{V}
$$

Internal energy\\
(13E.2a)

In many cases, the expression for $\langle\varepsilon\rangle$ established for each mode of motion in Topic 13C can be used and it is not necessary to go back to $q$ itself except for some formal manipulations. For instance, it is established in Topic 13C (eqn 13C.8) that the mean energy of a harmonic oscillator is $\left\langle\varepsilon^{\mathrm{v}}\right\rangle=h c \tilde{v} /\left(\mathrm{e}^{\beta h c \tilde{v}}-1\right)$. It follows that the molar internal energy of a system composed of oscillators is

$$
U_{\mathrm{m}}^{\mathrm{V}}(T)=U_{\mathrm{m}}^{\mathrm{V}}(0)+\frac{N_{\mathrm{A}} h c \tilde{v}}{\mathrm{e}^{\beta h c \tilde{v}}-1}
$$

\subsection*{Brief illustration 13E. 1}
The vibrational wavenumber of an $\mathrm{I}_{2}$ molecule is $214.6 \mathrm{~cm}^{-1}$. At $298.15 \mathrm{~K}, k T / h c=207.224 \mathrm{~cm}^{-1}$ and $h c \tilde{v}=4.26 \ldots \mathrm{zJ}$. With these values

$$
\beta h c \tilde{v}=\frac{h c \tilde{v}}{k T}=\frac{214.6 \mathrm{~cm}^{-1}}{207.224 \mathrm{~cm}^{-1}}=1.035 \ldots
$$

It follows that the vibrational contribution to the molar internal energy of $\mathrm{I}_{2}$ is

$$
\begin{aligned}
U_{\mathrm{m}}^{\mathrm{V}}(T) & =U_{\mathrm{m}}^{\mathrm{V}}(0)+\frac{\left(6.022 \times 10^{23} \mathrm{~mol}^{-1}\right) \times\left(4.26 \ldots \times 10^{-21} \mathrm{~J}\right)}{\mathrm{e}^{1.035 \ldots}-1} \\
& =U_{\mathrm{m}}^{\mathrm{V}}(0)+1.41 \mathrm{~kJ} \mathrm{~mol}^{-1}
\end{aligned}
$$

An alternative form of eqn 13E. 2 a is

$$
U(T)=U(0)-N\left(\frac{\partial \ln _{\mathcal{G}}}{\partial \beta}\right)_{V} \quad \begin{aligned}
& \text { Internal energy } \\
& \text { [independent molecules] }
\end{aligned}
$$

(13E.2b)

A very similar expression is used for a system of interacting molecules. In that case the canonical partition function (Topic 13D), Q, is used to write

$$
U(T)=U(0)-\left(\frac{\partial \ln Q}{\partial \beta}\right)_{V} \quad \begin{aligned}
& \text { Internal energy } \\
& \text { [interacting molecules] }
\end{aligned}
$$

(13E.2c)

\subsection*{(b) Heat capacity}
The constant-volume heat capacity (Topic 2 A ) is defined as $C_{V}=(\partial U / \partial T)_{V}$. Then, because the mean vibrational energy of a harmonic oscillator (eqn 13C.8, quoted above as $\left.\left\langle\varepsilon^{\mathrm{V}}\right\rangle=h c \tilde{v} /\left(\mathrm{e}^{\beta h c \tilde{v}}-1\right)\right)$ can be written in terms of the vibrational temperature $\theta^{\mathrm{V}}=h c \tilde{v} / k$ as

$$
\left\langle\varepsilon^{\mathrm{V}}\right\rangle=\frac{k \theta^{\mathrm{v}}}{\mathrm{e}^{\theta^{\mathrm{V}} / T}-1}
$$

it follows that the vibrational contribution to the molar con-stant-volume heat capacity is

$$
\begin{array}{rr}
U_{\mathrm{m}}(T)=U_{\mathrm{m}}(0)+N_{\mathrm{A}}\left\langle\varepsilon^{\mathrm{V}}\right\rangle & \begin{array}{l}
\mathrm{d}(1 / f) / \mathrm{d} x=-\left(1 / f^{2}\right) \mathrm{d} / \mathrm{d} x \\
\text { used twice }
\end{array} \\
C_{V, \mathrm{~m}}^{\mathrm{V}}=\frac{\mathrm{d} N_{\mathrm{A}}\left\langle\varepsilon^{\mathrm{V}}\right\rangle}{\mathrm{d} T}=R \theta^{\mathrm{V}} \frac{\mathrm{~d}}{\mathrm{~d} T} \frac{1}{\mathrm{e}^{\theta^{V} / T}-1}=R\left(\frac{\theta^{\mathrm{V}}}{T}\right)^{2} \frac{\mathrm{e}^{\theta^{\mathrm{V}} / T}}{\left(\mathrm{e}^{\theta^{V} / T}-1\right)^{2}}
\end{array}
$$

By noting that $\mathrm{e}^{\theta^{V} / T}=\left(\mathrm{e}^{\theta^{\mathrm{V}} / 2 T}\right)^{2}$, this expression can be rearranged into

$$
C_{V, \mathrm{~m}}^{\mathrm{V}}=R f(T) \quad f(T)=\left(\frac{\theta^{\mathrm{V}}}{T}\right)^{2}\left(\frac{\mathrm{e}^{-\theta^{\mathrm{V}} / 2 T}}{1-\mathrm{e}^{-\theta^{\mathrm{V}} / T}}\right)^{2}
$$


\begin{equation*}
\text { Vibrational contribution to } C_{V, m} \tag{13E.3}
\end{equation*}


The graph in Fig. 13E. 1 shows how the vibrational heat capacity depends on temperature. Note that even when the temperature\\
\includegraphics[max width=\textwidth, center]{2025_02_27_d4b95d9718f0b9e2e6b9g-592}

Figure 13E. 1 The temperature dependence of the vibrational heat capacity of a molecule in the harmonic approximation calculated by using eqn 13E.3. Note that the heat capacity is within 10 per cent of its classical value for temperatures greater than $\theta^{\vee}$.\\
is only slightly above $\theta^{\mathrm{V}}$ the heat capacity is close to its equipartition value. Equation 13E. 3 is essentially the same as the Einstein formula for the heat capacity of a solid (eqn 7A.8a) with $\theta^{\vee}$ the Einstein temperature, $\theta_{\mathrm{E}}$. The only difference is that in a solid the vibrations take place in three dimensions.

It is sometimes more convenient to convert the derivative with respect to $T$ into a derivative with respect to $\beta=1 / k T$ by using


\begin{equation*}
\frac{\mathrm{d}}{\mathrm{~d} T}=\frac{\mathrm{d} \beta}{\mathrm{~d} T} \frac{\mathrm{~d}}{\mathrm{~d} \beta}=-\frac{1}{k T^{2}} \frac{\mathrm{~d}}{\mathrm{~d} \beta}=-k \beta^{2} \frac{\mathrm{~d}}{\mathrm{~d} \beta} \tag{13E.4}
\end{equation*}


It follows that

$$
C_{V}=-k \beta^{2}\left(\frac{\partial U}{\partial \beta}\right)_{V}=-N k \beta^{2}\left(\frac{\partial\langle\varepsilon\rangle}{\partial \beta}\right)_{V}=N k \beta^{2}\left(\frac{\partial^{2} \ln q}{\partial \beta^{2}}\right)_{V}^{\text {eqn 13E.1 }}
$$

There is a much simpler route to finding $C_{V}$ when the equipartition principle can be applied, which is the case when $T$ $\gg \theta^{\mathrm{M}}$, where $\theta^{\mathrm{M}}$ is the characteristic temperature of the mode $\mathrm{M}\left(\theta^{\mathrm{v}}=h c \tilde{v} / k\right.$ for vibration, $\theta^{\mathrm{R}}=h c \tilde{B} / k$ for rotation). The heat capacity can then be estimated simply by counting the numbers of modes that are active, noting that the average molar energy of each mode is a multiple of $\frac{1}{2} R T$, differentiating that energy with respect to $T$, and getting that same multiple of $\frac{1}{2} R$. In gases, all three translational modes are always active and contribute $\frac{3}{2} R$ to the molar heat capacity. If the number of active rotational modes is denoted by $v^{\mathrm{R}}$ (so for most molecules at normal temperatures $v^{\mathrm{R} *}=2$ for linear molecules, and 3 for nonlinear molecules), then the rotational contribution is $\frac{1}{2} v^{R *} R$. If the temperature is high enough for $v^{v \star}$ vibrational modes to be active, then the vibrational contribution to the molar heat capacity is $v^{v *} R$. In most cases, $v^{v *} \approx 0$. It follows that the total molar heat capacity of a gas is approximately

$$
C_{V, \mathrm{~m}}=\frac{1}{2}\left(3+v^{\mathrm{R*}}+2 v^{\mathrm{V} *}\right) R \quad \begin{aligned}
& \text { Total heat capacity } \\
& {\left[T \gg \theta^{M}\right]}
\end{aligned}
$$

(13E.6)

\subsection*{Brief illustration 13E. 2}
The characteristic temperatures (in round numbers) of the vibrations of $\mathrm{H}_{2} \mathrm{O}$ are $5300 \mathrm{~K}, 2300 \mathrm{~K}$, and 5400 K ; the vibrations are therefore not excited at 373 K . The three rotational modes of $\mathrm{H}_{2} \mathrm{O}$ have characteristic temperatures $40 \mathrm{~K}, 21 \mathrm{~K}$, and 13 K , so they are fully excited, like the three translational modes. The translational contribution is $\frac{3}{2} R=12.5 \mathrm{~J} \mathrm{~K}^{-1} \mathrm{~mol}^{-1}$. Fully excited rotations contribute a further $12.5 \mathrm{~J} \mathrm{~K}^{-1} \mathrm{~mol}^{-1}$. Therefore, a value close to $25 \mathrm{~J} \mathrm{~K}^{-1} \mathrm{~mol}^{-1}$ is predicted. The experimental value is $26.1 \mathrm{~J} \mathrm{~K}^{-1} \mathrm{~mol}^{-1}$. The discrepancy is probably due to deviations from perfect gas behaviour.

\subsection*{13E.2 The entropy}
One of the most celebrated equations in statistical thermodynamics, the 'Boltzmann formula' for the entropy, is obtained by establishing the relation between the entropy and the weight of the most probable configuration.

\subsection*{How is that done? 13E. 1 \\
 Deriving the Boltzmann formula}
 for the entropyThe starting point for this derivation is the expression for the internal energy, $U(T)=U(0)+N\langle\varepsilon\rangle$, which, with $\langle\varepsilon\rangle=(1 / N) \sum_{i} N_{i} \varepsilon_{i}$ can be written

$$
U(T)=U(0)+\sum_{i} N_{i} \varepsilon_{i}
$$

The strategy is to use classical thermodynamics to establish a relation between $\mathrm{d} S$ and $\mathrm{d} U$, and then to use this relation to express $\mathrm{d} U$ in terms of the weight of the most probable configuration.

\subsection*{Step 1 Write an expression for $\mathrm{d} U(T)$}
A change in $U(T)$ may arise from either a modification of the energy levels of a system (when $\varepsilon_{i}$ changes to $\varepsilon_{i}+\mathrm{d} \varepsilon_{i}$ ) or from a modification of the populations (when $N_{i}$ changes to $N_{i}+\mathrm{d} N_{i}$ ). The most general change is therefore

$$
\mathrm{d} U(T)=\mathrm{d} U(0)+\sum_{i} N_{i} \mathrm{~d} \varepsilon_{i}+\sum_{i} \varepsilon_{i} \mathrm{~d} N_{i}
$$

Step 2 Write an expression for $\mathrm{d} S$\\
Because neither $U(0)$ nor the energy levels change when a system is heated at constant volume (Fig. 13E.2), in the absence of all changes other than heating, only the third (blue) term on the right survives. Moreover, from eqn 3E.1 $(\mathrm{d} U=T \mathrm{~d} S-p \mathrm{~d} V)$, $\mathrm{d} U=T \mathrm{~d} S$ under the same conditions. Therefore,

$$
\mathrm{d} S=\frac{\mathrm{d} U}{T}=\frac{1}{T} \sum_{i} \varepsilon_{i} \mathrm{~d} N_{i}=k \beta \sum_{i} \varepsilon_{i} \mathrm{~d} N_{i}
$$

Step 3 Write an expression for $\mathrm{d} S$ in terms of $\mathcal{W}$\\
Changes in the most probable configuration (the only one to consider) are given by eqn 13A.7 $\left(\partial(\ln \mathcal{W}) / \partial N_{i}+\alpha-\beta \varepsilon_{i}=0\right)$. It follows that $\beta \varepsilon_{i}=\partial(\ln \mathcal{W}) / \partial N_{i}+\alpha$ and, because the system contains a fixed number of molecules,

$$
\mathrm{d} S=k \sum_{i} \overbrace{\beta \varepsilon_{i}} \mathrm{~d} N_{i}=k \overbrace{\sum_{i}\left(\frac{\partial \ln \mathcal{W}}{\partial N_{i}}\right) \mathrm{d} N_{i}}^{\mathrm{d} \ln W}+k \alpha \overbrace{\sum_{i} \mathrm{~d} N_{i}}^{0}=k(\mathrm{~d} \ln \mathcal{W})
$$

This relation strongly suggests that\\
$-S=k \ln \mathcal{W}$\\
(13E.7)\\
\includegraphics[max width=\textwidth, center]{2025_02_27_d4b95d9718f0b9e2e6b9g-593}

Figure 13E. 2 (a) When a system is heated, the energy levels are unchanged but their populations are changed. (b) When work is done on a system, the energy levels themselves are changed. The levels in this case are the one-dimensional particle-in-a-box energy levels of Topic 7D: they depend on the size of the container and move apart as its length is decreased.

This very important expression is the Boltzmann formula for the entropy. In it, $\mathcal{W}$ is the weight of the most probable configuration of the system (as discussed in Topic 13A).

\subsection*{(a) Entropy and the partition function}
The statistical entropy, the entropy calculated from the Boltzmann formula, behaves in exactly the same way as the thermodynamic entropy. Thus, as the temperature is lowered, the value of $\mathcal{W}$, and hence of $S$, decreases because fewer configurations are consistent with the total energy. In the limit $T \rightarrow 0, \mathcal{W}=1$, so $\ln \mathcal{W}=0$, because only one configuration (every molecule in the lowest level) is compatible with $U(T)=$ $U(0)$. It follows that $S \rightarrow 0$ as $T \rightarrow 0$, which is compatible with the Third Law of thermodynamics, that the entropies of all perfect crystals approach the same value as $T \rightarrow 0$ (Topic 3C).

The challenge now is to establish a relation between the Boltzmann formula and the partition function.

How is that done? 13E. 2\\
Relating the statistical entropy to the partition function

It is necessary to separate the calculation into two parts, one for distinguishable independent molecules and the other for indistinguishable independent molecules. Interactions can be taken into account, in principle at least, by a simple generalization.

Step 1 Derive the relation for distinguishable independent molecules\\
For a system composed of $N$ distinguishable molecules, $\ln \mathscr{W}$ is given by eqn 13A.3 ( $\ln \mathcal{W}=N \ln N-\sum_{i} N_{i} \ln N_{i}$ ); the substitution $N=\sum_{i} N_{i}$ then gives

$$
\begin{aligned}
\ln \mathcal{W} & =\stackrel{\overbrace{N}^{N}}{\sum N_{i}} \ln N-\sum_{i} N_{i} \ln N_{i} \\
& =\sum_{i} N_{i} \ln N-\sum_{i} N_{i} \ln N_{i} \\
& =\sum_{i} N_{i}\left(\ln N-\ln N_{i}\right)=-\sum_{i} N_{i} \ln \frac{N_{i}}{N}
\end{aligned}
$$

Equation 13E. $7(S=k \ln \mathcal{W})$ then becomes

$$
S=-k \sum_{i} N_{i} \ln \frac{N_{i}}{N}
$$

The value of $N_{i} / N$ for the most probable distribution is given by the Boltzmann distribution, $N_{i} / N=\mathrm{e}^{-\beta \varepsilon_{i}} / \mathcal{G}$, and so

$$
\ln \frac{N_{i}}{N}=\ln \mathrm{e}^{-\beta \varepsilon_{i}}-\ln \mathscr{G}=-\beta \varepsilon_{i}-\ln \mathscr{G}
$$

Therefore,

$$
S=k \beta \overbrace{i}^{N\{\varepsilon\rangle} N_{i} \varepsilon_{i}+k \sum_{i} N_{i} \ln \mathscr{G}=N k \beta \varepsilon+N k \ln \mathscr{q}
$$

Finally, because $N\langle\varepsilon\rangle=U(T)-U(0)$ and $\beta=1 / k T$, it follows that\\
\includegraphics[max width=\textwidth, center]{2025_02_27_d4b95d9718f0b9e2e6b9g-594(1)}

Step 2 Derive the relation for indistinguishable molecules\\
To treat a system composed of $N$ indistinguishable molecules, the weight $\mathcal{W}$ is reduced by a factor of $N$ ! because the $N$ ! permutations of the molecules among the states result in the same state of the system. Therefore, $\ln \mathcal{W}$, and therefore the entropy, is reduced by $k \ln N$ ! from the 'distinguishable' value. Because $N$ is so large, Stirling's approximation $(\ln N!=$ $N \ln N-N$ ) can be used to convert eqn 13E.8a into

$$
\begin{aligned}
S(T) & =\frac{U(T)-U(0)}{T}+N k \ln \mathfrak{q}-k \overbrace{(N \ln N-N)}^{\ln N!} \\
& =\frac{U(T)-U(0)}{T}+N k \ln \frac{\mathfrak{q}}{N}+k N
\end{aligned}
$$

The $k N$ can be combined with the logarithm by writing it as $k N \ln$ e, to give\\
$-\left|S=\frac{U(T)-U(0)}{T}+N k \ln \frac{g \mathrm{e}}{N}\right|$\\[0pt]
(13E.8b) The entropy [independent, indistinguishable molecules]

\subsection*{Step 3 Generalize to interacting molecules}
For completeness, the corresponding expression for interacting molecules, based on the canonical partition function in place of the molecular partition function, is\\
\includegraphics[max width=\textwidth, center]{2025_02_27_d4b95d9718f0b9e2e6b9g-594}

Equation 13E.8a expresses the entropy of a collection of independent molecules in terms of the internal energy and the molecular partition function. However, because the energy of a molecule is a sum of contributions, such as translational (T), rotational ( R ), vibrational ( V ), and electronic ( E ), the partition function factorizes into a product of contributions. As a result, the entropy is also the sum of the individual contributions. In a gas, the molecules are free to change places, so they are indistinguishable; therefore, for the translational contribution to the entropy, use eqn 13E.8b. For the other modes (specifically R, V, and E ), which do not involve the exchange of molecules and therefore do not require the weight to be reduced by $N$ !, use eqn 13E.8a.

\subsection*{Brief illustration 13E. 3}
For a system with two states, with energies 0 and $\varepsilon$, it is shown in Topics 13B and 13C that the partition function and mean energy are $q=1+\mathrm{e}^{-\beta \varepsilon}$ and $\langle\varepsilon\rangle=\varepsilon /\left(\mathrm{e}^{\beta \varepsilon}+1\right)$. The contribution to the molar entropy, with $1 / T=k \beta$, is therefore

$$
S_{\mathrm{m}}=R\left\{\frac{\beta \varepsilon}{1+\mathrm{e}^{\beta \varepsilon}}+\ln \left(1+\mathrm{e}^{-\beta \varepsilon}\right)\right\}
$$

This awkward function is plotted in Fig. 13E.3. It should be noted that as $T \rightarrow \infty$ (corresponding to $\beta \rightarrow 0$ ), the molar entropy approaches $R \ln 2$.\\
\includegraphics[max width=\textwidth, center]{2025_02_27_d4b95d9718f0b9e2e6b9g-594(2)}

Figure 13E. 3 The temperature variation of the molar entropy of a collection of two-level systems expressed as a multiple of $R=N_{A} k$. As $T \rightarrow \infty$ the two states become equally populated and $S_{m}$ approaches $R \ln 2$.

\subsection*{(b) The translational contribution}
The expressions derived for the entropy are in line with what is expected for entropy as a measure of the spread of the populations of molecules over the available states. This interpretation can be illustrated by deriving an expression for the molar entropy of a monatomic perfect gas.

\subsection*{How is that done? 13 E .3 \\
 Deriving the expression for the}
 entropy of a monatomic perfect gasYou need to start with eqn 13E.8b for a collection of independent, indistinguishable atoms and write $N=n N_{A}$, where $N_{\mathrm{A}}$ is Avogadro's constant and $n$ is their amount (in moles). The only mode of motion for a gas of atoms is translation and $U(T)-U(0)=\frac{3}{2} n R T$. The partition function is $q=V / \Lambda^{3}$ (eqn 13B.10b), where $\Lambda$ is the thermal wavelength. Therefore,

$$
\begin{aligned}
S & =\frac{\overbrace{U(T)-U(0)}^{2}}{T}+n N_{\mathrm{A}} k \ln \frac{g \mathrm{e}}{n N_{\mathrm{A}}}=\frac{3}{2} n R+\overbrace{n N_{\mathrm{A}}}^{n R} \ln \frac{V \mathrm{e}}{n N_{\mathrm{A}} \Lambda^{3}} \\
& =n R\left\{\frac{3}{2}+\ln \frac{V_{\mathrm{m}} \mathrm{e}}{N_{\mathrm{A}} \Lambda^{3}}\right\}=n R \ln \frac{V_{\mathrm{m}} \mathrm{e}^{5 / 2}}{N_{\mathrm{A}} \Lambda^{3}}
\end{aligned}
$$

where $V_{\mathrm{m}}=V / n$ is the molar volume of the gas and $\frac{3}{2}$ has been replaced by $\ln \mathrm{e}^{3 / 2}$. Division of both sides by $n$ then results in the Sackur-Tetrode equation:\\
$-\left|S_{\mathrm{m}}=R \ln \left(\frac{V_{\mathrm{m}} \mathrm{e}^{5 / 2}}{N_{\mathrm{A}} \Lambda^{3}}\right)\right| \begin{array}{r}\text { Sackur-Tetrode equation } \\ \text { [monatomic perfect gas] }\end{array}$\\
where $\Lambda$ is the thermal wavelength $\left(\Lambda=h /(2 \pi m k T)^{1 / 2}\right)$. To calculate the standard molar entropy, note that $V_{\mathrm{m}}=R T / p$, and set $p=p^{\ominus}$ :


\begin{equation*}
S_{\mathrm{m}}^{\ominus}=R \ln \left(\frac{R T e^{5 / 2}}{p^{\ominus} N_{\mathrm{A}} \Lambda^{3}}\right)=R \ln \left(\frac{k T e^{5 / 2}}{p^{\ominus} \Lambda^{3}}\right) \tag{13E.9b}
\end{equation*}


These expressions are based on the high-temperature approximation of the partition functions, which assumes that many levels are occupied; therefore, they do not apply when $T$ is equal to or very close to zero.

\subsection*{Brief illustration 13E. 4}
The mass of an Ar atom is $m=39.95 m_{u}$. At $25^{\circ} \mathrm{C}$, the thermal wavelength of Ar is 16.0 pm and $k T=4.12 \times 10^{-21} \mathrm{~J}$. Therefore, the molar entropy of argon at this temperature is

$$
\begin{aligned}
S_{\mathrm{m}}^{\ominus} & =R \ln \left\{\frac{\left(4.12 \times 10^{-21} \mathrm{~J}\right) \times \mathrm{e}^{5 / 2}}{\left(10^{5} \mathrm{Nm}^{-2}\right) \times\left(1.60 \times 10^{-11} \mathrm{~m}\right)^{3}}\right\} \\
& =18.6 R=155 \mathrm{JK}^{-1} \mathrm{~mol}^{-1}
\end{aligned}
$$

On the basis that there are fewer accessible translational states for a lighter atom than for a heavy atom under the same conditions (see below), it can be anticipated that the standard molar entropy of Ne is likely to be smaller than for Ar; its actual value is 17.60 R at 298 K .

The physical interpretation of these equations is as follows:

\begin{itemize}
  \item Because the molecular mass appears in the numerator (because it appears in the denominator of $\Lambda$ ), the molar entropy of a perfect gas of heavy molecules is greater than that of a perfect gas of light molecules under the same conditions. This feature can be understood in terms of the energy levels of a particle in a box being closer together for heavy particles than for light particles, so more states are thermally accessible.
  \item Because the molar volume appears in the numerator, the molar entropy increases with the molar volume of the gas. The reason is similar: large containers have more closely spaced energy levels than small containers, so once again more states are thermally accessible.
  \item Because the temperature appears in the numerator (because, like $m$, it appears in the denominator of $\Lambda$ ), the molar entropy increases with increasing temperature. The reason for this behaviour is that more energy levels become accessible as the temperature is raised.
\end{itemize}

The Sackur-Tetrode equation written in the form

$$
S=n R \ln \frac{V \mathrm{e}^{5 / 2}}{n N_{\mathrm{A}} \Lambda^{3}}=n R \ln a V, a=\frac{\mathrm{e}^{5 / 2}}{n N_{\mathrm{A}} \Lambda^{3}}
$$

implies that when a monatomic perfect gas expands isothermally from $V_{\mathrm{i}}$ to $V_{\mathrm{f}}$, its entropy changes by

\[
\begin{array}{rlrl}
\Delta S & =n R \ln a V_{\mathrm{f}}-n R \ln a V_{\mathrm{i}} & & \\
& =n R \ln \frac{V_{\mathrm{f}}}{V_{\mathrm{i}}} & \begin{array}{l}
\text { Change of entropy on } \\
\text { expansion } \\
\text { [perfect gas, isothermal] }
\end{array} \tag{13E.10}
\end{array}
\]

This expression is the same as that obtained starting from the thermodynamic definition of entropy (Topic 3B).

\subsection*{(c) The rotational contribution}
The rotational contribution to the molar entropy, $S_{\mathrm{m}}^{\mathrm{R}}$, can be calculated once the molecular partition function is known.

For a linear molecule, the high-temperature limit of $q^{R}$ is $k T / \sigma h c \tilde{B}$ (eqn 13B.13b, $q^{\mathrm{R}}=T / \sigma \theta^{\mathrm{R}}$ with $\theta^{\mathrm{R}}=h c \tilde{B} / k$ ) and the equipartition theorem gives the rotational contribution to the molar internal energy as RT; therefore, from eqn 13E.8a:

$$
S_{\mathrm{m}}^{\mathrm{R}}=\frac{\overbrace{U_{\mathrm{m}}(T)-U_{\mathrm{m}}(0)}^{T}}{T}+R \ln \overbrace{\mathcal{q}^{\mathrm{R}}}^{\mathrm{kT} / \sigma h c \tilde{B}}
$$

and the contribution at high temperatures is

\[
S_{\mathrm{m}}^{\mathrm{R}}=R\left\{1+\ln \frac{k T}{\sigma h c \tilde{B}}\right\} \quad \begin{align*}
& \text { Rotational contribution }  \tag{13E.11a}\\
& {[\text { linear molecule, high }} \\
& \text { temperature } \left.\left(T \gg \theta^{\mathrm{R}}\right)\right]
\end{align*}
\]

In terms of the rotational temperature,

\[
S_{\mathrm{m}}^{\mathrm{R}}=R\left\{1+\ln \frac{T}{\sigma \theta^{\mathrm{R}}}\right\} \quad \begin{align*}
& \text { Rotational contribution }  \tag{13E.11b}\\
& {[\text { Linear molecule, high }} \\
& \text { temperature } \left.\left(T \gg \theta^{\mathrm{R}}\right)\right]
\end{align*}
\]

This function is plotted in Fig. 13E.4. It is seen that:

\begin{itemize}
  \item The rotational contribution to the entropy increases with temperature because more rotational states become accessible.
  \item The rotational contribution is large when $\tilde{B}$ is small, because then the rotational energy levels are close together.
\end{itemize}

It follows that large, heavy molecules have a large rotational contribution to their entropy. As shown in the following Brief illustration, the rotational contribution to the molar entropy of ${ }^{35} \mathrm{Cl}_{2}$ is $58.6 \mathrm{~J} \mathrm{~K}^{-1} \mathrm{~mol}^{-1}$ whereas that for $\mathrm{H}_{2}$ is only $12.7 \mathrm{~J} \mathrm{~K}^{-1} \mathrm{~mol}^{-1}$. That is, it is appropriate to regard $\mathrm{Cl}_{2}$ as a more rotationally disordered gas than $\mathrm{H}_{2}$, in the sense that at a given temperature $\mathrm{Cl}_{2}$ occupies a greater number of rotational states than $\mathrm{H}_{2}$ does.\\
\includegraphics[max width=\textwidth, center]{2025_02_27_d4b95d9718f0b9e2e6b9g-596}

Figure 13E. 4 The variation of the rotational contribution to the molar entropy of a linear molecule ( $\sigma=1$ ) using the hightemperature approximation and the exact expression (the latter evaluated up to $J=20$ ).

\subsection*{Brief illustration 13E. 5}
The rotational contribution for ${ }^{35} \mathrm{Cl}_{2}$ at $25^{\circ} \mathrm{C}$, for instance, is calculated by noting that $\sigma=2$ for this homonuclear diatomic molecule and taking $\tilde{B}=0.2441 \mathrm{~cm}^{-1}$ (corresponding to $24.41 \mathrm{~m}^{-1}$ ). The rotational temperature of the molecule is\\
$\theta^{\mathrm{R}}=\frac{\left(6.626 \times 10^{-34} \mathrm{Js}\right) \times\left(2.998 \times 10^{8} \mathrm{~m} \mathrm{~s}^{-1}\right) \times\left(24.41 \mathrm{~m}^{-1}\right)}{1.381 \times 10^{23} \mathrm{JK}^{-1}}=0.351 \mathrm{~K}$\\
Therefore,

$$
S_{\mathrm{m}}^{\mathrm{R}}=R\left\{1+\ln \frac{298 \mathrm{~K}}{2 \times(0.351 \mathrm{~K})}\right\}=7.05 R=58.6 \mathrm{JK}^{-1} \mathrm{~mol}^{-1}
$$

Equation 13E. 11 is valid at high temperatures $\left(T \gg \theta^{\mathrm{R}}\right)$. To track the rotational contribution down to low temperatures it is necessary to use the full form of the rotational partition function (Topic 13B; see Problem P13E.12). The resulting graph has the form shown in Fig. 13E.4, and it is seen that the approximate curve matches the exact curve very well for $T / \theta^{\mathrm{R}}$ greater than about 1 .

\subsection*{(d) The vibrational contribution}
The vibrational contribution to the molar entropy, $S_{\mathrm{m}}^{\mathrm{V}}$, is obtained by combining the expression for the molecular partition function (eqn 13B.15, $q^{v}=1 /\left(1-\mathrm{e}^{-\beta h c \tilde{v}}\right)=1 /\left(1-\mathrm{e}^{-\beta \varepsilon}\right)$ for $\varepsilon=h c \tilde{v}$ ) with the expression for the mean energy (eqn 13C.8, $\left.\left\langle\varepsilon^{\mathrm{v}}\right\rangle=\varepsilon /\left(\mathrm{e}^{\beta \varepsilon}-1\right)\right)$, to obtain

$$
\begin{aligned}
S_{\mathrm{m}}^{\mathrm{V}} & =\frac{\overbrace{U_{\mathrm{m}}(T)-U_{\mathrm{m}}(0)}^{N_{\mathrm{A}}\langle\varepsilon\rangle}}{\underbrace{T}_{1 / k \beta}}+R \ln \mathscr{q}^{\mathrm{V}}=\frac{\overbrace{N_{\mathrm{A}} k} \beta \varepsilon}{\mathrm{e}^{\beta \varepsilon}-1}+R \ln \frac{1}{1-\mathrm{e}^{-\beta \varepsilon}} \\
& =R\left\{\frac{\beta \varepsilon}{\mathrm{e}^{\beta \varepsilon}-1}-\ln \left(1-\mathrm{e}^{-\beta \varepsilon}\right)\right\}
\end{aligned}
$$

That is,

\[
S_{\mathrm{m}}^{\mathrm{V}}=R\left\{\frac{\beta h c \tilde{v}}{\mathrm{e}^{\beta h \tilde{v}}-1}-\ln \left(1-\mathrm{e}^{-\beta h c \tilde{v}}\right)\right\} \quad \begin{align*}
& \text { Vibrational }  \tag{13E.12a}\\
& \text { contribution to } \\
& \text { the molar entropy }
\end{align*}
\]

Once again it is convenient to express this formula in terms of a characteristic temperature, in this case the vibrational temperature $\theta^{\mathrm{V}}=h c \tilde{v} / k$ :

\[
S_{\mathrm{m}}^{\mathrm{V}}=R\left\{\frac{\theta^{\mathrm{v}} / T}{\mathrm{e}^{\theta^{\mathrm{v}} / T}-1}-\ln \left(1-\mathrm{e}^{\theta^{\mathrm{v}} / T}\right)\right\} \quad \begin{align*}
& \text { Vibrational }  \tag{13E.12b}\\
& \text { contribution to } \\
& \text { the molar entropy }
\end{align*}
\]

This function is plotted in Fig. 13E.5. As usual, it is helpful to interpret it, with the graph in mind:\\
\includegraphics[max width=\textwidth, center]{2025_02_27_d4b95d9718f0b9e2e6b9g-597(1)}

Figure 13E. 5 The temperature variation of the molar entropy of a collection of harmonic oscillators expressed as a multiple of $R=N_{\mathrm{A}} k$. The molar entropy approaches zero as $T \rightarrow 0$, and increases without limit as $T \rightarrow \infty$.

\begin{itemize}
  \item Both terms multiplying $R$ become zero as $T \rightarrow 0$, so the entropy is zero at $T=0$.
  \item The molar entropy rises as the temperature is increased as more vibrational states become accessible.
  \item The molar entropy is higher at a given temperature for molecules with heavy atoms or low force constant than one with light atoms or high force constant. The vibrational energy levels are closer together in the former case than in the latter, so more are thermally accessible.
\end{itemize}

\subsection*{Brief illustration 13E. 6}
The vibrational wavenumber of $\mathrm{I}_{2}$ is $214.5 \mathrm{~cm}^{-1}$, corresponding to $2.145 \times 10^{4} \mathrm{~m}^{-1}$, so its vibrational temperature is 309 K . At $25^{\circ} \mathrm{C}$

$$
\begin{aligned}
S_{\mathrm{m}}^{\mathrm{V}} & =R\left\{\frac{(309 \mathrm{~K}) /(298 \mathrm{~K})}{\mathrm{e}^{(309 \mathrm{~K}) /(298 \mathrm{~K})}-1}-\ln \left(1-\mathrm{e}^{-(309 \mathrm{~K}) /(298 \mathrm{~K})}\right)\right\}=1.01 R \\
& =8.38 \mathrm{JK}^{-1} \mathrm{~mol}^{-1}
\end{aligned}
$$

\subsection*{(e) Residual entropies}
Entropies may be calculated from spectroscopic data; they may also be measured experimentally (Topic 3C). In many cases there is good agreement, but in some the experimental entropy is less than the calculated value. One possibility is that the experimental determination failed to take a phase transition into account and a contribution of the form $\Delta_{\text {trs }} H / T_{\text {trs }}$ was incorrectly omitted from the sum. Another possibility is that some disorder is present in the solid even at $T=0$. The entropy at $T=0$ is then greater than zero and is called the residual entropy.

The origin and magnitude of the residual entropy can be explained by considering a crystal composed of AB molecules,\\
where A and B are similar atoms (such as CO, with its very small electric dipole moment). There may be so little energy difference between... AB AB AB AB ..., ... AB BA BA AB..., and other arrangements that the molecules adopt the orientations AB and BA at random in the solid. The entropy arising from residual disorder can be calculated readily by using the Boltzmann formula, $S=k \ln \mathcal{W}$. To do so, suppose that two orientations are equally probable, and that the sample consists of $N$ molecules. Because the same energy can be achieved in $2^{N}$ different ways (because each molecule can take either of two orientations), the total number of ways of achieving the same energy is $\mathcal{W}=2^{N}$. It follows that


\begin{equation*}
S=k \ln 2^{N}=N k \ln 2=n R \ln 2 \tag{13E.13a}
\end{equation*}


and that a residual molar entropy of $R \ln 2=5.8 \mathrm{~J} \mathrm{~K}^{-1} \mathrm{~mol}^{-1}$ is expected for solids composed of molecules that can adopt either of two orientations at $T=0$. If $s$ orientations are possible, the residual molar entropy is

$$
S_{\mathrm{m}}(0)=R \ln s \quad \text { Residual entropy }
$$

(13E.13b)

For CO , the measured residual entropy is $5 \mathrm{~J} \mathrm{~K}^{-1} \mathrm{~mol}^{-1}$, which is close to $R \ln 2$, the value expected for a random structure of the form ...CO CO OC CO OC OC....

Similar arguments apply to more complicated cases. Consider a sample of ice with $\mathrm{NH}_{2} \mathrm{O}$ molecules. Each O atom is surrounded tetrahedrally by four H atoms, two of which are attached by short $\sigma$ bonds, the other two being attached by long hydrogen bonds (Fig. 13E.6). It follows that each of the 2 NH atoms can be in one of two positions (either close to or far from an O atom as shown in Fig. 13E.7), resulting in $2^{2 N}$ possible arrangements. However, not all these arrangements are acceptable. Indeed, of the $2^{4}=16$ ways of arranging four H atoms around one O atom, only 6 have two short and two long OH distances and hence are acceptable (Fig. 13E.7). Therefore, the number of permitted arrangements is $\mathcal{W}=2^{2 N}\left(\frac{6}{16}\right)^{N}=\left(\frac{3}{2}\right)^{N}$.\\
\includegraphics[max width=\textwidth, center]{2025_02_27_d4b95d9718f0b9e2e6b9g-597}

Figure 13E. 6 The possible locations of H atoms around a central O atom in an ice crystal are shown by the white spheres. Only one of the locations on each bond may be occupied by an atom, and two H atoms must be close to the O atom and two H atoms must be distant from it.\\
\includegraphics[max width=\textwidth, center]{2025_02_27_d4b95d9718f0b9e2e6b9g-598}

Figure 13E. 7 The six possible arrangements of H atoms in the locations identified in Fig.13E.6. Occupied locations are denoted by grey spheres and unoccupied locations by white spheres.

It then follows that the residual entropy is $S(0) \approx k \ln \left(\frac{3}{2}\right)^{N}=$ $k N \ln \frac{3}{2}$, and its molar value is $S_{\mathrm{m}}(0) \approx R \ln \frac{3}{2}=3.4 \mathrm{~J} \mathrm{~K}^{-1} \mathrm{~mol}^{-1}$, which is in good agreement with the experimental value of $3.4 \mathrm{~J} \mathrm{~K}^{-1} \mathrm{~mol}^{-1}$. The model, however, is not exact because it ignores the possibility that next-nearest neighbours and those beyond can influence the local arrangement of bonds.

\subsection*{Checklist of concepts}
1. The internal energy is proportional to the derivative of the logarithm of the partition function with respect to temperature.2. The total heat capacity of a molecular substance is the sum of the contribution of each mode.3. The statistical entropy is defined by the Boltzmann formula and expressed in terms of the molecular partition function.4. The residual entropy is a non-zero entropy at $T=0$ arising from molecular disorder.\subsection*{Checklist of equations}
\begin{center}
\begin{tabular}{|c|c|c|c|}
\hline
Property & Equation & Comment & Equation number \\
\hline
Internal energy & $U(T)=U(0)-(N / \mathcal{q})(\partial \mathcal{q} / \partial \beta)_{V}=U(0)-N(\partial \ln q / \partial \beta)_{V}$ & Independent molecules & 13E.2a \\
\hline
\multirow[t]{2}{*}{Heat capacity} & $C_{V}=N k \beta^{2}\left(\partial^{2} \ln q / \partial \beta^{2}\right)_{V}$ & Independent molecules & 13E. 5 \\
\hline
 & \( C_{V, \mathrm{~m}}=\frac{1}{2}\left(3+v^{\mathrm{R}^{*}}+2 v^{V^{*}}\right) R \) & $T \gg \theta^{\text {M }}$ & 13E. 6 \\
\hline
Boltzmann formula for the entropy & $S=k \ln \mathcal{W}$ & Definition & 13 E .7 \\
\hline
\multirow[t]{2}{*}{Entropy} & $S=\{U(T)-U(0)\} / T+N k \ln \mathcal{G}$ & Distinguishable molecules & 13E.8a \\
\hline
 & $S=\{U(T)-U(0)\} / T+N k \ln (\mathrm{ge} / N)$ & Indistinguishable molecules & 13E.8b \\
\hline
Sackur-Tetrode equation & $S_{\mathrm{m}}=R \ln \left(V_{\mathrm{m}} \mathrm{e}^{5 / 2} / N_{\mathrm{A}} \Lambda^{3}\right)$ & Molar entropy of a monatomic perfect gas & 13E.9a \\
\hline
Residual molar entropy & $S_{\mathrm{m}}(0)=R \ln s$ & $s$ is the number of equivalent orientations & 13E.13b \\
\hline
\end{tabular}
\end{center}

\section*{TOPIC 13F Derived functions}
\subsection*{Why do you need to know this material?}
The power of chemical thermodynamics stems from its deployment of a variety of derived functions, particularly the enthalpy and Gibbs energy. It is therefore important to relate these functions to structural features through partition functions.

\subsection*{What is the key idea?}
The partition function provides a link between spectroscopic and structural data and the derived functions of thermodynamics, particularly the equilibrium constant.

\subsection*{What do you need to know already?}
This Topic develops the discussion of internal energy and entropy (Topic 13E). You need to know the relations between those properties and the enthalpy (Topic 2B) and the Helmholtz and Gibbs energies (Topic 3D). The final section makes use of the relation between the standard reaction Gibbs energy and the equilibrium constant (Topic 6A). Although the equations are introduced in terms of the canonical partition function (Topic 13D), all the applications are in terms of the molecular partition function (Topic 13B).

Classical thermodynamics makes extensive use of various derived functions. Thus, in thermochemistry the focus is on the enthalpy and, provided the pressure and temperature are constant, in discussions of spontaneity the focus is on the Gibbs energy. All these functions are derived from the internal energy and the entropy, which in terms of the canonical partition function, $Q$, are given by

$$
\begin{aligned}
U(T) & =U(0)-\left(\frac{\partial \ln Q}{\partial \beta}\right)_{V} \text { with } \beta=\frac{1}{k T} \text { Internal energy (13F.1a) } \\
S(T) & =\frac{U(T)-U(0)}{T}+k \ln Q \\
& =-k \beta\left(\frac{\partial \ln Q}{\partial \beta}\right)_{V}+k \ln Q
\end{aligned}
$$

There is no need to worry about this appearance of the canonical partition function (Topic 13D): the only applications in this Topic make use of the molecular partition function ( $q$, Topic 13B). Equation 13F. 1 can be regarded simply as a succinct way of expressing the relations for collections of independent molecules by writing $Q=q^{N}$ for distinguishable molecules and $Q=q^{N} / N$ ! for indistinguishable molecules (as in a gas). The advantage of using $Q$ is that the equations to be derived are simple and, if necessary, can be used to calculate thermodynamic properties when there are interactions between the molecules.

\subsection*{13F. 1 The derivations}
The Helmholtz energy, $A$, is defined as $A=U-T S$. This relation implies that at $T=0, A(0)=U(0)$, so substitution of the expressions for $U(T)$ and $S(T)$ in eqn 13E. 1 gives

$$
\begin{aligned}
A(T) & =A(0)-\left(\frac{\partial \ln Q}{\partial \beta}\right)_{V}-T\left\{-k \beta\left(\frac{\partial \ln Q}{\partial \beta}\right)_{V}+k \ln Q\right\} \\
& =A(0)-k T \ln Q
\end{aligned}
$$

That is,

$$
A(T)=A(0)-k T \ln Q \quad \text { Helmholtz energy }
$$

(13F.2)

An infinitesimal change in conditions changes the Helmholtz energy by $\mathrm{d} A=-p \mathrm{~d} V-S \mathrm{~d} T$ (this is the analogue of the expression for $\mathrm{d} G$ derived in Topic 3E (eqn 3E.7, $\mathrm{d} G=V \mathrm{~d} p-$ $\mathrm{Sd} T)$ ). It follows that on imposing constant temperature ( $\mathrm{d} T=0$ ), the pressure and the Helmholtz energy are related by $p=-(\partial A / \partial V)_{T}$. It then follows from eqn 13F. 2 that

$$
p=k T\left(\frac{\partial \ln Q}{\partial V}\right)_{T}
$$

Pressure (13F.3)

This relation is entirely general, and may be used for any type of substance, including perfect gases, real gases, and liquids. Because $Q$ is in general a function of the volume, temperature, and amount of substance, eqn 13F. 3 is an equation of state of the kind discussed in Topics 1A and 3E.

\subsection*{Example 13F.1 Deriving an equation of state}
Derive an expression for the pressure of a gas of independent particles.

Collect your thoughts You should suspect that the pressure is that given by the perfect gas law, $p=n R T / V$. To proceed systematically, substitute the explicit formula for $Q$ for a gas of independent, indistinguishable molecules. Only the translational partition function depends on the volume, so there is no need to include the partition functions for the internal modes of the molecules.

The solution For a gas of independent molecules, $Q=q^{N} / N$ ! with $\mathcal{q}=V / \Lambda^{3}$ :

$$
\begin{aligned}
& p=k T\left(\frac{\partial \ln Q}{\partial V}\right)_{T}=k T\left(\frac{\partial \ln \left(\mathcal{G}^{N} / N!\right)}{\partial V}\right)_{T} \\
&=k T\left(\frac{\partial\left(\ln \mathscr{G}^{N}-\ln N!\right)}{\partial V}\right)_{T}=k T\left(\frac{\partial(N \ln \mathcal{G}-\ln N!)}{\partial V}\right)_{T} \\
&=N k T\left(\frac{\partial \ln \mathscr{G}}{\partial V}\right)_{T}-k T \overbrace{\left(\frac{\partial \ln N!}{\partial V}\right)_{T}}^{0} \\
&=\frac{\mathrm{d} \ln f / \mathrm{d} x=(1 / f) \mathrm{d} / \mathrm{d} x}{\mathscr{q}}\left(\frac{\partial \mathfrak{q}}{\partial V}\right)_{T} \\
&=\frac{N k T}{V / \Lambda^{3}}\left(\frac{\partial\left(V / \Lambda^{3}\right)}{\partial V}\right)_{T}=\frac{N k T}{V}=\frac{n N_{\mathrm{A}} k T}{V}=\frac{n R T}{V} \\
& N_{A} k=R \\
& V
\end{aligned}
$$

The calculation shows that the equation of state of a gas of independent particles is indeed the perfect gas law, $p V=n R T$.\\
Comment. If $\beta$ is left undefined in eqn 13F.1, then the same calculation carried through results in $p=N / \beta V$. For this result to be the perfect gas law, it follows that $\beta=1 / k T$, as anticipated and used in Topics 13A-13E. This is the formal proof of that relation. A similar approach can be used to derive an equation of state resembling the van der Waals equation: see $A$ deeper look 7 on the website of this text.

Self-test 13F. 1 Derive the equation of state of a gas for which $Q=q^{N} f / N!$, with $q=V / \Lambda^{3}$, where $f$ depends on the volume.\\
\includegraphics[max width=\textwidth, center]{2025_02_27_d4b95d9718f0b9e2e6b9g-600}

At this stage the expressions for $U$ and $p$ and the definition $H=U+p V$, with $H(0)=U(0)$, can be used to obtain an expression for the enthalpy, $H$, of any substance:

$$
H(T)=H(0)-\left(\frac{\partial \ln Q}{\partial \beta}\right)_{V}+k T V\left(\frac{\partial \ln Q}{\partial V}\right)_{T} \quad \text { Enthalpy }
$$

(13F.4)

The fact that eqn 13F. 4 is rather cumbersome is a sign that the enthalpy is not a fundamental property: as shown in Topic 2B, it is more of an accounting convenience. For a gas of independent structureless particles $U(T)-U(0)=\frac{3}{2} n R T$ and $p V=n R T$. Therefore, for such a gas, it follows directly from $H=U+p V$ that


\begin{equation*}
H(T)-H(0)=\frac{5}{2} n R T \tag{13F.5}
\end{equation*}


One of the most important thermodynamic functions for chemistry is the Gibbs energy, $G=H-T S$. This definition implies that $G=U+p V-T S$ and therefore that $G=A+p V$. Note that $G(0)=A(0)$, both being equal to $U(0)$. The Gibbs energy can now be written in terms of the partition function by combining the expressions for $A$ and $p$ :


\begin{equation*}
G(T)=G(0)-k T \ln Q+k T V\left(\frac{\partial \ln Q}{\partial V}\right)_{T} \quad \text { Gibbs energy } \tag{13F.6}
\end{equation*}


This expression takes a simple form for a gas of independent molecules because $p V$ in the expression $G=A+p V$ can be replaced by $n R T$ :


\begin{equation*}
G(T)=G(0)-k T \ln Q+n R T \tag{13F.7}
\end{equation*}


Furthermore, because $Q=q^{N} / N$ ! for the indistinguishable particles in a gas and therefore $\ln Q=N \ln g-\ln N!$, it follows by using Stirling's approximation $(\ln N!=N \ln N-N)$ that


\begin{align*}
G(T) & =G(0)-\overbrace{N k}^{n N_{A} k=n R} T \ln \mathscr{G}+k T \overbrace{\ln N!}^{N \ln N-N}+n R T \\
& =G(0)-n R T \ln \mathscr{G}+k T(N \ln N-N)+n R T \\
& =G(0)-n R T \ln \mathscr{G}+\overbrace{N k T}^{n R T} \ln N \\
& =G(0)-n R T \ln \frac{\mathscr{G}}{N} \tag{13F.8}
\end{align*}


Note that a statistical interpretation of the Gibbs energy is now possible: because $q$ is the number of thermally accessible states and $N$ is the number of molecules, the difference $G(T)-G(0)$ is proportional to the logarithm of the average number of thermally accessible states available to a molecule. As this average number increases, the Gibbs energy falls further and further below $G(0)$. The thermodynamic tendency to lower Gibbs energy can now be seen to be a tendency to maximize the number of thermally accessible states.

It turns out to be convenient to define the molar partition function, $\mathfrak{q}_{\mathrm{m}}=\mathfrak{q} / n$ (with units $\mathrm{mol}^{-1}$ and $n=N / N_{\mathrm{A}}$ ), for then

$$
G(T)=G(0)-n R T \ln \frac{\mathcal{G}_{\mathrm{m}}}{N_{\mathrm{A}}} \quad \begin{aligned}
& \text { Gibbs energy } \\
& \text { [indistinguishable, } \\
& \text { independent molecules] }
\end{aligned}
$$

(13F.9a)

To use this expression, $G(0)=U(0)$ is identified with the energy of the system when all the molecules are in their ground state, $E_{0}$. To calculate the standard Gibbs energy, the partition function has its standard value, $\mathscr{q}_{\mathrm{m}}^{\ominus}$, which is evaluated by setting the molar volume in the translational contribution equal to the standard molar volume, so $\mathscr{q}_{\mathrm{m}}^{\ominus}=\left(V_{\mathrm{m}}^{\ominus} / \Lambda^{3}\right) \mathcal{q}^{\mathrm{R}} \mathcal{q}^{\mathrm{V}}$ with $V_{\mathrm{m}}^{\ominus}=R T / p^{\ominus}$. The standard molar Gibbs energy is then obtained with these substitutions and after dividing through by $n$ :

$$
G_{\mathrm{m}}^{\ominus}(T)=G_{\mathrm{m}}^{\ominus}(0)-R T \ln \frac{\mathcal{q}_{\mathrm{m}}^{\ominus}}{N_{\mathrm{A}}} \quad \begin{aligned}
& \text { Standard molar Gibbs } \\
& \text { energy } \\
& \text { [indistinguishable, } \\
& \text { independent molecules] }]
\end{aligned}
$$

(13F.9b)\\
where $G_{\mathrm{m}}^{\ominus}(0)=E_{0, \mathrm{~m}}$, the molar ground-state energy of the system.

\subsection*{Example 13F.2 Calculating a standard Gibbs energy of formation from partition functions}
Calculate the standard Gibbs energy of formation of $\mathrm{H}_{2} \mathrm{O}(\mathrm{g})$ at $25^{\circ} \mathrm{C}$.

Collect your thoughts Write the chemical equation for the formation reaction, and then the expression for the standard Gibbs energy of formation in terms of the Gibbs energy of each molecule; then express those Gibbs energies in terms of the molecular partition functions. Ignore molecular vibration as it is unlikely to be excited at $25^{\circ} \mathrm{C}$. The rotational constant for $\mathrm{H}_{2}$ is $60.864 \mathrm{~cm}^{-1}$, that for $\mathrm{O}_{2}$ is $1.4457 \mathrm{~cm}^{-1}$, and those for $\mathrm{H}_{2} \mathrm{O}$ are 27.877, 14.512, and $9.285 \mathrm{~cm}^{-1}$. Before using the approximate form of the rotational partition functions, you need to judge whether the temperature is high enough: if it is not, use the full expression. For the values of $E_{0, \mathrm{~m}}$, use bond enthalpies from Table 9C.3; for precise calculations, bond energies (at $T=0$ ) should be used instead. You will need the following expressions from Topic 13B:

$$
\Lambda(\mathrm{J})=\frac{h}{\{2 \pi m(\mathrm{~J}) k T\}^{1 / 2}} \overbrace{q^{\mathrm{R}}=\frac{k T}{2 h c \tilde{B}}}^{\begin{array}{c}
\text { Linear molecule. } \\
\sigma=2
\end{array}} \text { and } \overbrace{q^{\mathrm{R}}=\frac{1}{2}\left(\frac{k T}{h c}\right)^{1 / 2}\left(\frac{\pi}{\tilde{A} \tilde{B} \tilde{C})^{1 / 2}}\right.}^{\begin{array}{c}
\text { Nonlinear molecule, } \\
\sigma=2
\end{array}}
$$

The solution The chemical reaction is $\mathrm{H}_{2}(\mathrm{~g})+\frac{1}{2} \mathrm{O}_{2}(\mathrm{~g}) \rightarrow \mathrm{H}_{2} \mathrm{O}(\mathrm{g})$. Therefore,

$$
\Delta_{\mathrm{f}} G^{\ominus}=G_{\mathrm{m}}^{\ominus}\left(\mathrm{H}_{2} \mathrm{O}, \mathrm{~g}\right)-G_{\mathrm{m}}^{\ominus}\left(\mathrm{H}_{2}, \mathrm{~g}\right)-\frac{1}{2} G_{\mathrm{m}}^{\ominus}\left(\mathrm{O}_{2}, \mathrm{~g}\right)
$$

Now write the standard molar Gibbs energies in terms of the standard molar partition functions of each species J:

$$
G_{\mathrm{m}}^{\ominus}(\mathrm{J})=E_{0, \mathrm{~m}}(\mathrm{~J})-R T \ln \frac{\mathfrak{q}_{\mathrm{m}}^{\ominus}(\mathrm{J})}{N_{\mathrm{A}}}, \mathscr{q}_{\mathrm{m}}^{\ominus}(\mathrm{J})=\mathscr{q}_{\mathrm{m}}^{\mathrm{T} \ominus}(\mathrm{~J}) \mathcal{q}^{\mathrm{R}}(\mathrm{~J})=\frac{V_{\mathrm{m}}^{\ominus}}{\Lambda(\mathrm{J})^{3}} \mathcal{q}^{\mathrm{R}}(\mathrm{~J})
$$

Therefore

$$
\begin{aligned}
& \Delta_{\mathrm{f}} G^{\ominus}=\left\{E_{0, \mathrm{~m}}\left(\mathrm{H}_{2} \mathrm{O}\right)-R T \ln \frac{\mathcal{q}_{\mathrm{m}}^{\ominus}\left(\mathrm{H}_{2} \mathrm{O}\right)}{N_{\mathrm{A}}}\right\}-\left\{E_{0, \mathrm{~m}}\left(\mathrm{H}_{2}\right)-R T \ln \frac{\mathcal{q}_{\mathrm{m}}^{\ominus}\left(\mathrm{H}_{2}\right)}{N_{\mathrm{A}}}\right\} \\
&-\frac{1}{2}\left\{E_{0, \mathrm{~m}}\left(\mathrm{O}_{2}\right)-R T \ln \frac{\mathcal{G}_{\mathrm{m}}^{\ominus}\left(\mathrm{O}_{2}\right)}{N_{\mathrm{A}}}\right\} \\
&=\Delta E_{0, \mathrm{~m}}-R T \ln \frac{q_{\mathrm{m}}^{\circ}\left(\mathrm{H}_{2} \mathrm{O}\right) / N_{\mathrm{A}}}{[\underbrace{\left\{V_{\mathrm{m}}^{\ominus} / N_{\mathrm{A}} \Lambda\left(\mathrm{H}_{2}\right)^{3}\right\} \mathcal{q}^{\mathrm{R}}\left(\mathrm{H}_{2}\right)}_{\mathscr{q}_{\mathrm{m}}^{\circ}\left(\mathrm{H}_{2}\right) / N_{\mathrm{A}}}]}] \underbrace{\left\{V_{\mathrm{m}}^{\ominus} / N_{\mathrm{A}} \Lambda\left(\mathrm{O}_{2}\right)^{3}\right\} q^{\mathrm{R}}\left(\mathrm{O}_{2}\right)}_{\mathscr{q}_{\mathrm{m}}^{\circ}\left(\mathrm{O}_{2}\right) / N_{\mathrm{A}}}]^{1 / 2} \\
&= \Delta E_{0, \mathrm{~m}}-R T \ln \frac{N_{\mathrm{A}}^{1 / 2}\left\{\Lambda\left(\mathrm{H}_{2}\right) \Lambda\left(\mathrm{O}_{2}\right)^{1 / 2} / \Lambda\left(\mathrm{H}_{2} \mathrm{O}\right)\right\}^{3}}{V_{\mathrm{m}}^{\ominus 1 / 2}\left\{\mathfrak{q}^{\mathrm{R}}\left(\mathrm{H}_{2}\right) \mathcal{q}^{\mathrm{R}}\left(\mathrm{O}_{2}\right)^{1 / 2} / \mathcal{q}^{\mathrm{R}}\left(\mathrm{H}_{2} \mathrm{O}\right)\right\}}
\end{aligned}
$$

Now substitute the data. The molar energy difference (with bond dissociation energies in kilojoules per mole indicated) is

$$
\Delta E_{0, \mathrm{~m}}=\overbrace{E_{0, \mathrm{~m}}\left(\mathrm{H}_{2} \mathrm{O}\right)}^{-492-428=-920}-\overbrace{E_{0, \mathrm{~m}}\left(\mathrm{H}_{2}\right)}^{-436}-\frac{1}{2} \overbrace{E_{0, \mathrm{~m}}\left(\mathrm{O}_{2}\right)}^{-497}=-236 \mathrm{k} \mathrm{~mol}^{-1}
$$

Next, establish whether the temperature is high enough for the approximate expressions for the rotational partition functions (Topic 13B) to be reliable:

\begin{center}
\begin{tabular}{lllcll}
\hline
 & $\mathrm{H}_{2} \mathrm{O}$ & $\mathrm{H}_{2} \mathrm{O}$ & $\mathrm{H}_{2} \mathrm{O}$ & $\mathrm{H}_{2}$ & $\mathrm{O}_{2}$ \\
\hline
$\tilde{X} / \mathrm{cm}^{-1}, X=A, B$, or $C$ & 27.877 & 14.512 & 9.285 & 60.864 & 1.4457 \\
$\theta^{\mathrm{R}} / \mathrm{K}$ & 40.1 & 20.9 & 13.4 & 87.5 & 2.1 \\
\hline
\end{tabular}
\end{center}

Only $\mathrm{H}_{2}$ is marginal at 298 K , so for it, use the full calculation of the rotational partition function; for the others, use the approximate forms quoted above:

$$
\begin{array}{lll}
\Lambda\left(\mathrm{H}_{2}\right)=71.21 \mathrm{pm} & \Lambda\left(\mathrm{O}_{2}\right)=17.87 \mathrm{pm} & \Lambda\left(\mathrm{H}_{2} \mathrm{O}\right)=23.82 \mathrm{pm} \\
q^{\mathrm{R}}\left(\mathrm{H}_{2}\right)=1.88 & q^{\mathrm{R}}\left(\mathrm{O}_{2}\right)=71.60 & q^{\mathrm{R}}\left(\mathrm{H}_{2} \mathrm{O}\right)=42.13
\end{array}
$$

It then follows that

$$
\Delta_{\mathrm{f}} G^{\ominus}=-236 \mathrm{~kJ} \mathrm{~mol}^{-1}-\overbrace{R T}^{2.48 \mathrm{~kJ} \mathrm{~mol}^{-1}} \ln 0.0291=-227 \mathrm{~kJ} \mathrm{~mol}^{-1}
$$

The value quoted in Table 2C. 1 of the Resource section is $-228.57 \mathrm{~kJ} \mathrm{~mol}^{-1}$.

Self-test 15E. 2 Estimate the standard Gibbs energy of formation of $\mathrm{NH}_{3}(\mathrm{~g})$ at $25^{\circ} \mathrm{C}$. The rotational constants of $\mathrm{NH}_{3}$ are $\tilde{B}=10.001 \mathrm{~cm}^{-1}$ and $\tilde{A}=6.449 \mathrm{~cm}^{-1}$.

\subsection*{13F.2 Equilibrium constants}
The following discussion is focused on gas-phase reactions in which the equilibrium constant is defined in terms of the partial pressures of the reactants and products.

\subsection*{(a) The relation between $K$ and the partition function}
The Gibbs energy of a gas of independent molecules is given by eqn 13F. 9 in terms of the molar partition function, $q_{\mathrm{m}}=$ $q / n$. The equilibrium constant $K$ of a reaction is related to the standard reaction Gibbs energy. The task is to combine these two relations and so obtain an expression for the equilibrium constant in terms of the molecular partition functions of the reactants and products.

\subsection*{How is that done? 13F. 1 Relating the equilibrium constant}
 to partition functionsYou need to use the expressions for the standard molar Gibbs energies, $G^{\ominus} / n$, of each species to find an expression for the standard reaction Gibbs energy. Then find the equilibrium constant $K$ by using eqn 6 A. $15\left(\Delta_{\mathrm{r}} G^{\ominus}=-R T \ln K\right)$.

Step 1 Write an expression for $\Delta_{\mathrm{r}} G^{\ominus}$\\
From eqn 13F.9b, the standard molar reaction Gibbs energy for the reaction $a \mathrm{~A}+b \mathrm{~B} \rightarrow c \mathrm{C}+d \mathrm{D}$ is

$$
\begin{aligned}
\Delta_{\mathrm{r}} G^{\ominus}= & c G_{\mathrm{m}}^{\ominus}(\mathrm{C})+d G_{\mathrm{m}}^{\ominus}(\mathrm{D})-\left\{a G_{\mathrm{m}}^{\ominus}(\mathrm{A})+b G_{\mathrm{m}}^{\ominus}(\mathrm{B})\right\} \\
= & c G_{\mathrm{m}}^{\ominus}(\mathrm{C}, 0)+d G_{\mathrm{m}}^{\ominus}(\mathrm{D}, 0)-\left\{a G_{\mathrm{m}}^{\ominus}(\mathrm{A}, 0)+b G_{\mathrm{m}}^{\ominus}(\mathrm{B}, 0)\right\} \\
& -R T\left\{c \ln \frac{q_{\mathrm{G}, \mathrm{~m}}^{\ominus}}{N_{\mathrm{A}}}+d \ln \frac{\mathcal{q}_{\mathrm{D}, \mathrm{~m}}^{\ominus}}{N_{\mathrm{A}}}-a \ln \frac{\mathcal{q}_{\mathrm{A}, \mathrm{~m}}^{\ominus}}{N_{\mathrm{A}}}-b \ln \frac{\mathcal{q}_{\mathrm{B}, \mathrm{~m}}^{\ominus}}{N_{\mathrm{A}}}\right\}
\end{aligned}
$$

Because $G_{\mathrm{m}}(\mathrm{J}, 0)=E_{0, \mathrm{~m}}(\mathrm{~J})$, the molar ground-state energy of the species J, the first (blue) term on the right is

$$
c E_{0, \mathrm{~m}}(\mathrm{C}, 0)+d E_{0, \mathrm{~m}}(\mathrm{D}, 0)-\left\{a E_{0, \mathrm{~m}}(\mathrm{~A}, 0)+b E_{0, \mathrm{~m}}(\mathrm{~B}, 0)\right\}=\Delta_{\mathrm{r}} E_{0}
$$

Then, by using $a \ln x=\ln x^{a}$ and $\ln x+\ln y=\ln x y$,

$$
\begin{aligned}
\Delta_{\mathrm{r}} G^{\ominus} & =\Delta_{\mathrm{r}} E_{0}-R T \ln \frac{\left(\mathfrak{q}_{\mathrm{C}, \mathrm{~m}}^{\ominus} / N_{\mathrm{A}}\right)^{c}\left(\mathfrak{q}_{\mathrm{D}, \mathrm{~m}}^{\ominus} / N_{\mathrm{A}}\right)^{d}}{\left(\mathfrak{q}_{\mathrm{A}, \mathrm{~m}}^{\ominus} / N_{\mathrm{A}}\right)^{a}\left(\mathfrak{q}_{\mathrm{B}, \mathrm{~m}}^{\ominus} / N_{\mathrm{A}}\right)^{b}} \\
& =-R T\left\{-\frac{\Delta_{\mathrm{r}} E_{0}}{R T}+\ln \frac{\left(\mathfrak{q}_{\mathrm{C}, \mathrm{~m}}^{\ominus} / N_{\mathrm{A}}\right)^{c}\left(\mathfrak{q}_{\mathrm{D}, \mathrm{~m}}^{\ominus} / N_{\mathrm{A}}\right)^{d}}{\left(\mathfrak{q}_{\mathrm{A}, \mathrm{~m}}^{\ominus} / N_{\mathrm{A}}\right)^{a}\left(\mathfrak{q}_{\mathrm{B}, \mathrm{~m}}^{\mathrm{A}} / N_{\mathrm{A}}\right)^{b}}\right\}
\end{aligned}
$$

Step 2 Write an expression for $K$\\
At this stage pick out an expression for $K$ by comparing this equation with $\Delta_{r} G^{\ominus}=-R T \ln K$, which gives

$$
\ln K=-\frac{\Delta_{\mathrm{r}} E_{0}}{R T}+\ln \frac{\left(\mathfrak{q}_{\mathrm{C}, \mathrm{~m}}^{\ominus} / N_{\mathrm{A}}\right)^{c}\left(\mathfrak{q}_{\mathrm{D}, \mathrm{~m}}^{\ominus} / N_{\mathrm{A}}\right)^{d}}{\left(\mathfrak{q}_{\mathrm{A}, \mathrm{~m}}^{\ominus} / N_{\mathrm{A}}\right)^{a}\left(\mathfrak{q}_{\mathrm{B}, \mathrm{~m}}^{\ominus} / N_{\mathrm{A}}\right)^{b}}
$$

Finally, by forming the exponential of both sides, this equation becomes

$$
K=\frac{\left(\mathfrak{q}_{\mathrm{C}, \mathrm{~m}}^{\ominus} / N_{\mathrm{A}}\right)^{c}\left(\mathscr{q}_{\mathrm{D}, \mathrm{~m}}^{\ominus} / N_{\mathrm{A}}\right)^{d}}{\left(\mathfrak{q}_{\mathrm{A}, \mathrm{~m}}^{\ominus} / N_{\mathrm{A}}\right)^{a}\left(\mathfrak{q}_{\mathrm{B}, \mathrm{~m}}^{\ominus} / N_{\mathrm{A}}\right)^{b}} \mathrm{e}^{-\Delta_{\mathrm{t}} E_{0} / R T} \quad \begin{aligned}
& \text { The } \\
& \begin{array}{l}
\text { equilibrium } \\
\text { constant }
\end{array}
\end{aligned}
$$

(13F.10a)\\
where $\Delta_{\mathrm{r}} E_{0}$ is the difference in molar energies of the ground states of the products and reactants and is calculated from the bond dissociation energies of the species (Fig. 13F.1). In terms of the (signed) stoichiometric numbers introduced in Topic 2C, eqn 13F.10a can be written\\
\includegraphics[max width=\textwidth, center]{2025_02_27_d4b95d9718f0b9e2e6b9g-602}\\
\includegraphics[max width=\textwidth, center]{2025_02_27_d4b95d9718f0b9e2e6b9g-602(1)}

Figure 13F. 1 The definition of $\Delta_{r} E_{0}$ for the calculation of equilibrium constants. The reactants are imagined as dissociating into atoms and then forming the products from the atoms.

\subsection*{(b) A dissociation equilibrium}
Equation 13F.10a can be used to write an expression for the equilibrium constant for the dissociation of a diatomic molecule $\mathrm{X}_{2}$ :

$$
\mathrm{X}_{2}(\mathrm{~g}) \rightleftharpoons 2 \mathrm{X}(\mathrm{~g}) \quad K=\frac{p_{\mathrm{x}}^{2}}{p_{\mathrm{x}_{2}} p^{\ominus}}
$$

According to eqn 13F.10a (with $a=1, b=0, c=2$, and $d=0$ ):


\begin{equation*}
K=\frac{\left(\mathscr{q}_{\mathrm{x}, \mathrm{~m}}^{\ominus} / N_{\mathrm{A}}\right)^{2}}{\mathscr{q}_{\mathrm{x}_{2}, \mathrm{~m}}^{\ominus} / N_{\mathrm{A}}} \mathrm{e}^{-\Delta_{\mathrm{r}} E_{0} / R T}=\frac{\left(\mathscr{q}_{\mathrm{x}, \mathrm{~m}}^{\ominus}\right)^{2}}{\mathfrak{q}_{\mathrm{x}_{2}, \mathrm{~m}}^{\ominus} N_{\mathrm{A}}} \mathrm{e}^{-\Delta_{\mathrm{t}} E_{0} / R T} \tag{13F.11a}
\end{equation*}


with


\begin{equation*}
\Delta_{\mathrm{r}} E_{0}=2 E_{0, \mathrm{~m}}(\mathrm{X}, 0)-E_{0, \mathrm{~m}}\left(\mathrm{X}_{2}, 0\right)=N_{\mathrm{A}} h c \tilde{D}_{0}(\mathrm{X}-\mathrm{X}) \tag{13F.11b}
\end{equation*}


where $N_{\mathrm{A}} h c \tilde{D}_{0}(\mathrm{X}-\mathrm{X})$ is the (molar) dissociation energy of the $\mathrm{X}-\mathrm{X}$ bond. The standard molar partition functions of the atoms X are

$$
\mathscr{q}_{\mathrm{x}, \mathrm{~m}}^{\ominus}=\overbrace{g_{\mathrm{x}}}^{q^{\mathrm{E}}} \times \frac{\overbrace{V_{\mathrm{m}}^{\ominus}}^{\Lambda_{\mathrm{x}}^{3}}=\frac{g_{\mathrm{x}} R T}{p^{\ominus} \Lambda_{\mathrm{x}}^{3}}=R T / p^{\circ}}{v^{\circ}}
$$

where $g_{\mathrm{x}}$ is the degeneracy of the electronic ground state of X . The diatomic molecule $\mathrm{X}_{2}$ also has rotational and vibrational degrees of freedom, so its standard molar partition function is

$$
q_{\mathrm{x}_{2}, \mathrm{~m}}^{\ominus}=g_{\mathrm{x}_{2}} \frac{V_{\mathrm{m}}^{\ominus}}{\Lambda_{\mathrm{x}_{2}}^{3}} q_{\mathrm{x}_{2}}^{\mathrm{R}} q_{\mathrm{x}_{2}}^{\mathrm{V}}=\frac{g_{\mathrm{x}_{2}} R T q_{\mathrm{x}_{2}}^{\mathrm{R}} q_{\mathrm{x}_{2}}^{\mathrm{V}}}{p^{\ominus} \Lambda_{\mathrm{x}_{2}}^{3}}
$$

where $g_{\mathrm{x}_{2}}$ is the degeneracy of the electronic ground state of $\mathrm{X}_{2}$. It follows that


\begin{align*}
& K=\frac{\left(g_{\mathrm{X}} R T / p^{\ominus} \Lambda_{\mathrm{X}}^{3}\right)^{2}}{g_{\mathrm{X}_{2}} N_{\mathrm{A}} R T \mathscr{q}_{\mathrm{x}_{2}}^{\mathrm{R}} \mathscr{q}_{\mathrm{x}_{2}}^{\mathrm{v}} / p^{\ominus} \Lambda_{\mathrm{X}_{2}}^{3}} \mathrm{e}^{-N_{\mathrm{A}} h e \tilde{0}_{0}(\mathrm{X}-\mathrm{X}) / R T} \\
& \begin{array}{l}
\quad \begin{aligned}
R=N_{A} k \\
g_{X}^{2} k T \Lambda_{\mathrm{X}_{2}}^{3} \\
g_{\mathrm{X}_{2}} \rho^{\circ} q_{\mathrm{X}_{2}}^{\mathrm{R}} q_{\mathrm{X}_{2}} \\
\Lambda_{\mathrm{X}}^{6}
\end{aligned} \mathrm{e}^{-h \tilde{D}_{0}(\mathrm{X}-\mathrm{X}) k T}
\end{array} \tag{13F.12}
\end{align*}


All the quantities in this expression can be calculated from spectroscopic data.

\subsection*{Example 13F.3 Evaluating an equilibrium constant}
Evaluate the equilibrium constant for the dissociation $\mathrm{Na}_{2}(\mathrm{~g})$ $\rightarrow 2 \mathrm{Na}(\mathrm{g})$ at 1000 K . The data for $\mathrm{Na}_{2}$ are: $\tilde{B}=0.1547 \mathrm{~cm}^{-1}$, $\tilde{v}=159.2 \mathrm{~cm}^{-1}$, and $N_{\mathrm{A}} h c \tilde{D}_{0}=70.4 \mathrm{~kJ} \mathrm{~mol}^{-1}$.

Collect your thoughts Recall that Na has a doublet ground state (term symbol ${ }^{2} \mathrm{~S}_{1 / 2}$ ). You need to use eqn 13 F .12 and the expressions for the partition functions assembled in Topic 13B. For such a heavy diatomic molecule it is safe to use the approximate expression for its rotational partition function (but check that assumption for consistency). Remember that for a homonuclear diatomic molecule, $\sigma=2$.

The solution The partition functions and other quantities required are as follows:

$$
\begin{array}{ll}
\Lambda\left(\mathrm{Na}_{2}\right)=8.14 \mathrm{pm} & \Lambda(\mathrm{Na})=11.5 \mathrm{pm} \\
q^{\mathrm{R}}\left(\mathrm{Na}_{2}\right)=2246 & q^{\mathrm{V}}\left(\mathrm{Na}_{2}\right)=4.885 \\
g(\mathrm{Na})=2 & g\left(\mathrm{Na}_{2}\right)=1 \\
h c \widetilde{D} / k T=8.47 \ldots &
\end{array}
$$

There are many rotational states occupied, so the use of the approximate formula for $q^{R}\left(\mathrm{Na}_{2}\right)$ is valid. Then, from eqn 13F.12,

$$
\begin{aligned}
& K=\frac{2^{2} \times(1.381 \times 10^{-19} \overbrace{\mathrm{~J} \mathrm{~K}}}{\mathrm{kg} \mathrm{~m}^{2} \mathrm{~s}^{-2}}) \times(1000 \mathrm{~K}) \times\left(8.14 \times 10^{-12} \mathrm{~m}\right)^{3} \\
&(10^{5} \underbrace{\mathrm{~Pa}}_{\mathrm{kg} \mathrm{~m}^{-1} \mathrm{~s}^{-2}}) \times 2246 \times 4.885 \times\left(1.15 \times 10^{-11} \mathrm{~m}\right)^{6}
\end{aligned} \mathrm{e}^{-8.47 \ldots}
$$

Self- test $13 F .3$ Evaluate $K$ at 1500 K . Is the answer consistent with the dissociation being endothermic?\\
səK 〔̌s :ıамsuн

\subsection*{(c) Contributions to the equilibrium constant}
To appreciate the physical basis of equilibrium constants, consider a simple $\mathrm{R} \rightleftharpoons \mathrm{P}$ gas-phase equilibrium ( R for reactants, $P$ for products).

Figure 13F. 2 shows two sets of energy levels; one set of states belongs to R , and the other belongs to P . The populations of the states are given by the Boltzmann distribution, and are independent of whether any given state happens to belong to R or to P. A single Boltzmann distribution spreads, without distinction, over the two sets of states. If the spacings of R and P are similar (as in Fig. 13F.2), and the ground state of P lies above that of $R$, the diagram indicates that R will dominate in the equilibrium mixture. However, if P has a high density of states (a large number of states in a given energy range, as in Fig. 13F.3), then, even though its ground-state energy lies above that of R , the species P might still dominate at equilibrium.

It is quite easy to show that the ratio of numbers of R and P molecules at equilibrium is given by a Boltzmann-like expression.\\
\includegraphics[max width=\textwidth, center]{2025_02_27_d4b95d9718f0b9e2e6b9g-603(1)}

Figure 13F. 2 The array of R (eactant) and P (roduct) energy levels. At equilibrium all are accessible (to differing extents, depending on the temperature), and the equilibrium composition of the system reflects the overall Boltzmann distribution of populations. As $\Delta_{r} E_{0}$ increases, $R$ becomes dominant.\\
\includegraphics[max width=\textwidth, center]{2025_02_27_d4b95d9718f0b9e2e6b9g-603}

Figure 13F. 3 It is important to take into account the densities of states of the molecules. Even though P might lie above R in energy (that is, $\Delta_{r} E_{0}$ is positive), P might have so many states that its total population dominates in the mixture. In classical thermodynamic terms, when considering equilibria, entropies must be taken into account as well as enthalpies.

\subsection*{How is that done? 13F. 2}
Relating the equilibrium constant to state populations

You need to start the derivation by noting that the population in a state $i$ of the composite ( $\mathrm{R}, \mathrm{P}$ ) system is $N_{i}=\mathrm{Ne}^{-\beta \varepsilon} / \mathrm{g}$, where $N$ is the total number of molecules.\\
Step 1 Write expressions for the numbers of $R$ and $P$ molecules The total number of R molecules is the sum of the populations of the ( $\mathrm{R}, \mathrm{P}$ ) system taken over the states belonging to $R$; these states are labelled $r$ with energies $\varepsilon_{r}$. The total number of P molecules is the sum over the states belonging to P ; these states are labelled p with energies $\varepsilon_{\mathrm{p}}^{\prime}$ (the prime is explained in a moment):

$$
N_{\mathrm{R}}=\sum_{\mathrm{r}} N_{\mathrm{r}}=\frac{N}{q} \sum_{\mathrm{r}} \mathrm{e}^{-\beta \varepsilon_{\mathrm{r}}} \quad N_{\mathrm{P}}=\sum_{\mathrm{p}} N_{\mathrm{p}}=\frac{N}{\mathrm{q}} \sum_{\mathrm{p}} \mathrm{e}^{-\beta \varepsilon_{\mathrm{p}}}
$$

The sum over the states of R is its partition function, $\mathscr{G}_{\mathrm{R}}$, so $N_{\mathrm{R}}=N_{G_{\mathrm{R}}} / \mathrm{g}$. The sum over the states of P is also a partition function, but the energies are measured from the ground state of the combined system, which is the ground state of R. However, because $\varepsilon_{\mathrm{p}}^{\prime}=\varepsilon_{\mathrm{p}}+\Delta \varepsilon_{0}$, where $\Delta \varepsilon_{0}$ is the separation of zero-point energies (as in Fig. 13F.3),

$$
\begin{aligned}
N_{\mathrm{P}} & =\frac{N}{q} \sum_{\mathrm{p}} \mathrm{e}^{-\beta\left(\varepsilon_{\mathrm{p}}+\Delta \varepsilon_{0}\right)}=\frac{N}{G}\left(\sum_{\mathrm{p}} \mathrm{e}^{-\beta \varepsilon_{\mathrm{p}}}\right) \mathrm{e}^{-\beta \Delta \varepsilon_{0}}=\frac{N \mathcal{G}_{\mathrm{P}}}{q} \mathrm{e}^{-\beta \Delta \varepsilon_{0}} \\
& =\frac{N \mathscr{G}_{\mathrm{P}}}{q} \mathrm{e}^{-\Delta_{\mathrm{r}} E_{0} / R T}
\end{aligned}
$$

The switch from $\Delta \varepsilon_{0} / k T$ to $\Delta_{\mathrm{r}} E_{0} / R T$ in the last step is the conversion of molecular energies to molar energies.

Step 2 Write an expression for the equilibrium constant\\
The ratio of populations is

$$
\frac{N_{\mathrm{P}}}{N_{\mathrm{R}}}=\frac{\mathscr{q}_{\mathrm{P}}}{G_{\mathrm{R}}} \mathrm{e}^{-\Delta_{\mathrm{r}} E_{0} / R T}
$$

The equilibrium constant of the $\mathrm{R} \rightleftharpoons \mathrm{P}$ reaction is proportional to the ratio of the numbers of the two types of molecule. Therefore,\\
\includegraphics[max width=\textwidth, center]{2025_02_27_d4b95d9718f0b9e2e6b9g-604}

For an $\mathrm{R} \rightleftharpoons \mathrm{P}$ equilibrium, the $V_{\mathrm{m}}^{\ominus}$ factors in the partition functions cancel, so the appearance of $q$ in place of $q^{\ominus}$ has no effect. In the case of a more general reaction, the conversion from $q$ to $q^{\ominus}$ comes about at the stage of converting the pressures that occur in $K$ to numbers of molecules.\\
\includegraphics[max width=\textwidth, center]{2025_02_27_d4b95d9718f0b9e2e6b9g-604(1)}

Figure 13F. 4 The model used in the text for exploring the effects of energy separations and densities of states on equilibria. The products $P$ can dominate provided $\Delta_{r} E_{0}$ is not too large and $P$ has an appreciable density of states.

The implications of eqn 13 F .13 can be seen most clearly by exaggerating the molecular features that contribute to it. Suppose that R has only a single accessible level, which implies that $g_{\mathrm{R}}=1$. Also suppose that P has a large number of evenly, closely spaced levels, with spacing $\varepsilon$ (Fig. 13F.4). The partition function for such an array is calculated in Brief illustration 13B.1, and is $q=1 /\left(1-\mathrm{e}^{-\beta \varepsilon}\right)$. Provided the levels are close (in the sense $\varepsilon \ll k T$ ), $q \approx k T / \varepsilon$ in the high-temperature limit (Topic 13B). In this model system, the equilibrium constant is


\begin{equation*}
K=\frac{k T}{\varepsilon} \mathrm{e}^{-\Delta_{\tau} E_{0} / R T} \tag{13F.14}
\end{equation*}


When $\Delta_{\mathrm{r}} E_{0}$ is very large, the exponential term dominates and $K \ll 1$, which implies that very little P is present at equilibrium. When $\Delta_{\mathrm{r}} E_{0}$ is small but still positive, $K$ can exceed 1 because the factor $k T / \varepsilon$ may be large enough to overcome the small size of the exponential term. The size of $K$ then reflects the predominance of P at equilibrium on account of its high density of states. At low temperatures $K \ll 1$ and the system consists entirely of R. At high temperatures the exponential function approaches 1 and the factor $k T / \varepsilon$ is large. Now P is dominant. In this endothermic reaction (endothermic because $P$ lies above $R$ ) a rise in temperature favours $P$, because its states become accessible. This is the behaviour described, from a macroscopic perspective, in Topic 6B.

The model also shows why the Gibbs energy, $G$, and not just the enthalpy, determines the position of equilibrium. It shows that the density of states (and hence the entropy) of each species as well as their relative energies controls the distribution of populations and hence the value of the equilibrium constant.

\subsection*{Checklist of concepts}
1. The thermodynamic functions $A, p, H$, and $G$ can be calculated from the canonical partition function.2. The equilibrium constant can be written in terms of the partition functions of the reactants and products.3. The equilibrium constant for dissociation of a diatomic molecule in the gas phase may be calculated from spectroscopic data.\begin{enumerate}
  \setcounter{enumi}{3}
  \item The physical basis of chemical equilibrium can be understood in terms of a competition between energy separations and densities of states.
\end{enumerate}

\subsection*{Checklist of equations}
\begin{center}
\begin{tabular}{lll}
\hline
Property & Equation & Comment \\
\hline
Helmholtz energy & $A(T)=A(0)-k T \ln Q$ &  \\
Pressure & $p=k T(\partial \ln Q / \partial V)_{T}$ &  \\
Enthalpy & $H(T)=H(0)-(\partial \ln Q / \partial \beta)_{V}+k T V(\partial \ln Q / \partial V)_{T}$ & 13 F .2 \\
Gibbs energy & $G(T)=G(0)-k T \ln Q+k T V(\partial \ln Q / \partial V)_{T}$ & 13 F .3 \\
 & $G(T)=G(0)-n R T \ln \left(\mathscr{q}_{\mathrm{m}} / N_{\mathrm{A}}\right)$ & Perfect gas \\
Equilibrium constant & $K=\left\{\prod_{\mathrm{J}}\left(\mathcal{q}_{J, \mathrm{~m}}^{\ominus} / N_{\mathrm{A}}\right)^{v_{\mathrm{J}}}\right\} \mathrm{e}^{-\Delta_{\mathrm{r}} E_{0} / R T}$ & Perfect gas \\
\hline
\end{tabular}
\end{center}

\chapter{FOCUS 14 Molecular interactions}
The electric properties of molecules give rise to molecular interactions. In turn, those interactions govern the formation of condensed phases and the structures and functions of macromolecules and molecular assemblies.

Macromolecules are built from covalently linked components. They are everywhere, inside us and outside us. Some are natural: they include polysaccharides (such as cellulose), polypeptides (such as protein enzymes), and polynucleotides (such as deoxyribonucleic acid, DNA). Others are synthetic (such as nylon and polystyrene). Molecules both large and small may also gather together in a process called 'self-assembly' and give rise to aggregates that to some extent behave like macromolecules.

\subsection*{14A The electric properties of molecules}
Important electric properties of molecules include 'electric dipole moments' and 'polarizabilities'. These properties reflect the degree to which the nuclei of atoms exert control over the electrons in a molecule, either by causing electrons to accumulate in particular regions, or by permitting them to respond more or less strongly to the effects of external electric fields.\\
14A. 1 Electric dipole moments; 14A. 2 Polarizabilities;\\
14A. 3 Polarization

\subsection*{14B Interactions between molecules}
This Topic describes the basic theory of several important molecular interactions, with a special focus on 'van der Waals interactions' between closed-shell molecules and 'hydrogen bonding'. Many liquids and solids are bound together by one or more of the cohesive interactions explored in this Topic. These interactions are also important for the structural organization of macromolecules.

\subsection*{14C Liquids}
This Topic begins with the basic theory of molecular interactions in liquids, and then turns to a description of the properties of liquid surfaces. It is seen that important effects, such as 'surface tension', 'capillary action', the formation of 'surface films', and condensation, can be explained by thermodynamics arguments.\\
14C. 1 Molecular interactions in liquids; 14C. 2 The liquid-vapour interface; 14C. 3 Surface films; 14C. 4 Condensation

\subsection*{14D Macromolecules}
Macromolecules adopt shapes that are governed by molecular interactions. This Topic considers a range of shapes, but focuses on the structureless 'random coil' and partially structured coils. It also explores the connection between structure and the mechanical and thermal properties of macromolecules.\\
14D. 1 Average molar masses; 14D. 2 The different levels of structure; 14D. 3 Random coils; 14D. 4 Mechanical properties; 14D. 5 Thermal properties

\subsection*{14E Self-assembly}
Atoms, small molecules, and macromolecules can form large aggregates, sometimes by processes involving self-assembly. This Topic explores 'colloids', 'micelles', and biological membranes, which are assemblies with some of the typical properties of molecules but also with their own characteristic features. It also introduces an important type of molecular interaction, the 'hydrophobic interaction', which is driven by changes in entropy of the solvent.\\
14E. 1 Colloids; 14E. 2 Micelles and biological membranes

\subsection*{Web resources What is an application of this material?}
Molecular interactions play important roles in biochemistry and biomedicine. Natural macromolecules differ in certain respects from synthetic macromolecules, particularly in their\\
composition and the resulting structure. The different levels of structure in proteins and nucleic acids are explored in Impact 21. In Impact 22 attention shifts to the binding of a drug, a small molecule or protein, to a specific receptor site of a target molecule, such as a larger protein or nucleic acid. The chemical result of the formation of this assembly is the inhibition of the progress of disease.

\section*{TOPIC 14A The electric properties of molecules }
\subsection*{Why do you need to know this material?}
Because the molecular interactions responsible for the formation of condensed phases and large molecular assemblies arise from the electric properties of molecules, you need to know how the electronic structures of molecules account for these interactions.

\subsection*{What is the key idea?}
The nuclei of atoms exert control over the electrons in a molecule, and can cause electrons to accumulate in various regions, or permit them to respond to external fields.

\subsection*{What do you need to know already?}
You need to be familiar with the Coulomb law (The chemist's toolkit 6 of Topic 2A), molecular geometry, and molecular orbital theory, especially the existence of the energy gap between a HOMO and LUMO (Topic 9E). One calculation draws on the Boltzmann distribution (Topic 13A).

The electric properties of molecules are responsible for many of the properties of bulk matter. The small imbalances of charge distributions in molecules and the ability of electron distributions to be distorted allow molecules to interact with one another and to respond to externally applied fields.

\subsection*{14A.1 Electric dipole moments}
An electric dipole consists of two electric charges $+Q$ and $-Q$ with a vector separation $R$. A point electric dipole is an electric dipole in which $R$ is very small compared with its distance to the observer. The electric dipole moment is a vector $\mu$ (1) that points from the negative charge to the positive charge and has a magnitude given by\\
$\mu=Q R$\\[0pt]
Magnitude of the electric dipole moment [definition]\\
\includegraphics[max width=\textwidth, center]{2025_02_27_d4b95d9718f0b9e2e6b9g-617}

1 Electric dipole

Although the SI unit of dipole moment is coulomb metre ( C m ), it is still commonly reported in the non-SI unit debye, D, named after Peter Debye, a pioneer in the study of dipole moments of molecules:


\begin{equation*}
1 \mathrm{D}=3.33564 \times 10^{-30} \mathrm{Cm} \tag{14A.2}
\end{equation*}


The magnitude of the dipole moment formed by a pair of charges $+e$ and $-e$ separated by 100 pm is $1.6 \times 10^{-29} \mathrm{Cm}$, which corresponds to 4.8 D . The magnitudes of the dipole moments of small molecules are typically about 1 D .

A polar molecule is a molecule with a permanent electric dipole moment. A permanent dipole moment stems from the partial charges on the atoms in the molecule, which arise from differences in electronegativity or, in more sophisticated treatments, variations in electron density through the molecule (Topic 9E). Nonpolar molecules acquire an induced dipole moment in an electric field on account of the distortion the field causes in their electronic distributions and nuclear positions. However, this induced moment is only temporary, and disappears as soon as the perturbing field is removed. Polar molecules also have their permanent dipole moments temporarily modified by an applied field.

All heteronuclear diatomic molecules are polar and typical values of $\mu$ are 1.08 D for HCl and 0.42 D for HI (Table 14A.1). Molecular symmetry is of the greatest importance in deciding whether a polyatomic molecule is polar or not. Indeed, molecular symmetry is more important than the question of whether or not the atoms in the molecule belong to the same element. For this reason, and as seen in the following Brief illustration, homonuclear polyatomic molecules may be polar if they have low symmetry and the atoms are in inequivalent positions.

Table 14A. 1 Dipole moments and polarizability volumes*

\begin{center}
\begin{tabular}{lll}
\hline
 & $\mu / \mathrm{D}$ & $\alpha^{\prime} /\left(10^{-30} \mathrm{~m}^{3}\right)$ \\
\hline
$\mathrm{CCl}_{4}$ & 0 & 10.5 \\
$\mathrm{H}_{2}$ & 0 & 0.819 \\
$\mathrm{H}_{2} \mathrm{O}$ & 1.85 & 1.48 \\
HCl & 1.08 & 2.63 \\
HI & 0.42 & 5.45 \\
\hline
\end{tabular}
\end{center}

\footnotetext{\begin{itemize}
  \item More values are given in the Resource section
\end{itemize}
}\subsection*{Brief illustration 14A. 1}
The angular molecule ozone (2) is homonuclear. However, it is polar because the central O atom is different from the outer two (it is bonded to two atoms, which are bonded only to one). Moreover, the dipole moments associated with each bond make an angle to each other and do not cancel. The heteronuclear linear triatomic molecule $\mathrm{CO}_{2}$ is nonpolar because,\\
\includegraphics[max width=\textwidth, center]{2025_02_27_d4b95d9718f0b9e2e6b9g-618(3)}

2 Ozone, $\mathrm{O}_{3}$\\
\includegraphics[max width=\textwidth, center]{2025_02_27_d4b95d9718f0b9e2e6b9g-618}

3 Carbon dioxide, $\mathrm{CO}_{2}$ although there are partial charges on all three atoms, the dipole moment associated with the OC bond points in the opposite direction to the dipole moment associated with the CO bond, and the two cancel (3).

The dipole moment of a polyatomic molecule can be resolved into contributions from various groups of atoms in the molecule and their relative locations (Fig. 14A.1). Thus, 1,4 -dichlorobenzene is nonpolar by symmetry on account of the cancellation of two equal but opposing $\mathrm{C}-\mathrm{Cl}$ moments (exactly as in carbon dioxide). 1,2-Dichlorobenzene, however, has a dipole moment which is approximately the resultant of two chlorobenzene dipole moments arranged at $60^{\circ}$ to each other. This technique of 'vector addition' can be applied with fair success to other series of related molecules. The magnitude of the resultant moment $\mu_{\text {res }}$ of $\mu_{1}$ and $\mu_{2}$ that make an angle $\Theta$ to each other (4) is approximately (see The chemist's toolkit 22 of Topic 8C)


\begin{equation*}
\mu_{\mathrm{res}} \approx\left(\mu_{1}^{2}+\mu_{2}^{2}+2 \mu_{1} \mu_{2} \cos \Theta\right)^{1 / 2} \tag{14A.3a}
\end{equation*}


\begin{center}
\includegraphics[max width=\textwidth]{2025_02_27_d4b95d9718f0b9e2e6b9g-618(2)}
\end{center}

Figure 14A. 1 The resultant dipole moments (red in (c) and (d)) of the dichlorobenzene isomers, (b) to (d), can be obtained approximately by vectorial addition of two chlorobenzene dipole moments (shown in (a), with $\mu_{\text {obs }}=1.57 \mathrm{D}$ ).\\
\includegraphics[max width=\textwidth, center]{2025_02_27_d4b95d9718f0b9e2e6b9g-618(1)}

4 Addition of dipole moments

When the two contributing dipole moments have the same magnitude (as in the dichlorobenzenes), this equation simplifies to

$$
\mu_{\mathrm{res}} \approx\left\{2 \mu_{1}^{2}(1+\cos \Theta)\right\}^{1 / 2}=2 \mu_{1} \cos \frac{1}{2} \Theta
$$

\subsection*{Brief illustration 14A. 2}
Consider ortho (1,2-) and meta (1,3-) disubstituted benzenes, for which $\Theta_{\text {ortho }}=60^{\circ}$ and $\Theta_{\text {meta }}=120^{\circ}$. It follows from eqn 14 A .3 b that the ratio of the magnitudes of the electric dipole moments is:

$$
\frac{\mu_{\text {res }, \text { ortho }}}{\mu_{\text {res }, \text { meta }}}=\frac{\cos \frac{1}{2} \Theta_{\text {ortho }}}{\cos \frac{1}{2} \Theta_{\text {meta }}}=\frac{\cos \frac{1}{2}\left(60^{\circ}\right)}{\cos \frac{1}{2}\left(120^{\circ}\right)}=\frac{\frac{3^{1 / 2}}{2}}{\frac{1}{2}}=3^{1 / 2} \approx 1.7
$$

A more reliable approach to the calculation of dipole moments is to take into account the locations and magnitudes of the partial charges on all the atoms. These partial charges are included in the output of many molecular structure software packages. To calculate the $x$-component of the dipole moment, for instance, it is necessary to know the partial charge on each atom and the atom's $x$-coordinate relative to a point in the molecule and form the sum


\begin{equation*}
\mu_{x}=\sum_{\mathrm{J}} Q_{\mathrm{J}} x_{\mathrm{J}} \tag{14A.4a}
\end{equation*}


Here $Q_{\mathrm{J}}$ is the partial charge of atom $\mathrm{J}, x_{\mathrm{J}}$ is the $x$-coordinate of atom J , and the sum is over all the atoms in the molecule. Analogous expressions are used for the $y$ - and $z$-components. For an electrically neutral molecule, the origin of the coordinates is arbitrary, so it is best chosen to simplify the calculations. In common with all vectors, the magnitude of $\mu$ is related to the three components $\mu_{x}, \mu_{y}$, and $\mu_{z}$ by


\begin{equation*}
\mu=\left(\mu_{x}^{2}+\mu_{y}^{2}+\mu_{z}^{2}\right)^{1 / 2} \tag{14A.4b}
\end{equation*}


\subsection*{Example 14A. 1 Calculating a molecular dipole moment}
Estimate the magnitude and orientation of the electric dipole moment of the planar amide group shown in (5) by using the partial charges (as multiples of $e$ ) and the locations of\\
the atoms shown as $(x, y, z)$ coordinates, with distances in picometres.\\
\includegraphics[max width=\textwidth, center]{2025_02_27_d4b95d9718f0b9e2e6b9g-619(4)}

Collect your thoughts You need to use eqn 14A.4a to calculate each of the components of the dipole moment and then eqn 14 A .4 b to assemble the three components into the magnitude of the dipole moment. Note that the partial charges are multiples of the fundamental charge, $e=1.602 \times 10^{-19} \mathrm{C}$. This group is a fragment of a neutral molecule, so the choice of origin of the coordinates is not arbitrary. Here the origin is taken to be coincident with the carbon atom.

The solution The expression for $\mu_{x}$ is

$$
\begin{aligned}
\mu_{x}= & (-0.36 e) \times(132 \mathrm{pm})+(0.45 e) \times(0 \mathrm{pm})+(0.18 e) \times(182 \mathrm{pm}) \\
& +(-0.38 e) \times(-62.0 \mathrm{pm}) \\
= & 8.8 e \mathrm{pm} \\
= & 8.8 \times\left(1.602 \times 10^{-19} \mathrm{C}\right) \times\left(10^{-12} \mathrm{~m}\right)=1.4 \times 10^{-30} \mathrm{Cm}
\end{aligned}
$$

corresponding to $\mu_{x}=+0.42 \mathrm{D}$. The expression for $\mu_{y}$ is:

$$
\begin{aligned}
\mu_{y}= & (-0.36 e) \times(0 \mathrm{pm})+(0.45 e) \times(0 \mathrm{pm})+(0.18 e) \times(-87 \mathrm{pm}) \\
& +(-0.38 e) \times(107 \mathrm{pm}) \\
= & -56 e \mathrm{pm} \\
= & -9.02 \times 10^{-30} \mathrm{Cm}
\end{aligned}
$$

It follows that $\mu_{y}=-2.7 \mathrm{D}$. The amide group is planar, so $\mu_{z}=$ 0 and

$$
\mu=\left\{(0.42 \mathrm{D})^{2}+(-2.7 \mathrm{D})^{2}\right\}^{1 / 2}=2.7 \mathrm{D}
$$

The orientation of the dipole moment is found by arranging an arrow of length 2.7 units of length to have $x, y$, and $z$ components of $0.42,-2.7$, and 0 units; the orientation is superimposed on (5).

Self-test 14A. 1 Estimate the magnitude of the electric dipole moment of methanal (formaldehyde) by using the information in (6).\\
\includegraphics[max width=\textwidth, center]{2025_02_27_d4b95d9718f0b9e2e6b9g-619}\\
\includegraphics[max width=\textwidth, center]{2025_02_27_d4b95d9718f0b9e2e6b9g-619(1)}

Molecules may have higher multipoles, or arrays of point charges (Fig. 14A.2). Specifically, an $\boldsymbol{n}$-pole is an array of point charges with an $n$-pole moment but no lower moment. Thus, a monopole ( $n=1$ ) is a point charge, and the monopole moment is what is normally called the overall charge. A dipole ( $n=2$ ), as already seen, is an array of charges that has no monopole moment (no net charge). A quadrupole $(n=3)$ consists of\\
\includegraphics[max width=\textwidth, center]{2025_02_27_d4b95d9718f0b9e2e6b9g-619(3)}

Figure 14A. 2 Typical charge arrays corresponding to electric multipoles. The field arising from an arbitrary charge distribution can be expressed as the superposition of the fields arising from a superposition of multipoles.\\
an array of point charges that has neither net charge nor dipole moment (as for a $\mathrm{CO}_{2}$ molecule, 3). An octupole ( $n=4$ ) consists of an array of point charges that sum to zero and which has neither a dipole moment nor a quadrupole moment (as for a $\mathrm{CH}_{4}$ molecule, 7).\\
\includegraphics[max width=\textwidth, center]{2025_02_27_d4b95d9718f0b9e2e6b9g-619(2)}

7 Methane, $\mathrm{CH}_{4}$

\subsection*{14A.2 Polarizabilities}
The failure of nuclear charges to control the surrounding electrons totally means that those electrons can respond to external fields. Therefore, an applied electric field can distort a molecule as well as align its permanent electric dipole moment. The magnitude $\mu^{*}$ of the induced dipole moment, $\mu^{*}$, is proportional to the electric field strength, $\mathcal{E}$, so

$$
\mu^{*}=\alpha \mathcal{E} \quad \begin{aligned}
& \text { Polarizability } \\
& \text { [definition] }
\end{aligned}
$$

(14A.5a)\\
The constant of proportionality $\alpha$ is the polarizability of the molecule. The greater the polarizability, the larger is the induced dipole moment for a given applied field. In a formal treatment, vector quantities are used to allow for the possibility that the induced dipole moment might not lie parallel to the applied field, in which case the scalar $\alpha$ is replaced by $\alpha$, a $3 \times 3$ matrix.

When the applied field is very strong (as in tightly focused laser beams), the magnitude of the induced dipole moment is not strictly linear in the strength of the field, and

$$
\mu^{*}=\alpha E+\frac{1}{2} \beta E^{2}+\cdots \quad \begin{aligned}
& \text { Hyperpolarizability } \\
& \text { [definition] }
\end{aligned}
$$

(14A.5b)\\
The coefficient $\beta$ is the (first) hyperpolarizability of the molecule.\\
Polarizability has the units (coulomb metre) ${ }^{2}$ per joule $\left(\mathrm{C}^{2} \mathrm{~m}^{2} \mathrm{~J}^{-1}\right)$. That collection of units is awkward, so $\alpha$ is often expressed as a polarizability volume, $\alpha^{\prime}$, by using the relation

$$
\begin{array}{ll}
\alpha^{\prime}=\frac{\alpha}{4 \pi \varepsilon_{0}} & \begin{array}{l}
\text { Polarizability volume } \\
\text { [definition] }
\end{array}
\end{array}
$$

(14A.6)\\
where $\varepsilon_{0}$ is the vacuum permittivity (The chemist's toolkit 6 of Topic 2 A ). Because the units of $4 \pi \varepsilon_{0}$ are coulomb-squared per joule per metre ( $\mathrm{C}^{2} \mathrm{~J}^{-1} \mathrm{~m}^{-1}$ ), it follows that $\alpha^{\prime}$ has the dimensions of volume (hence its name). Polarizability volumes are similar in magnitude to actual molecular volumes (of the order of $10^{-30} \mathrm{~m}^{3}, 10^{-3} \mathrm{~nm}^{3}, 1 \AA^{3}$ ).

\subsection*{Brief illustration 14A. 3}
The polarizability volume of $\mathrm{H}_{2} \mathrm{O}$ is $1.48 \times 10^{-30} \mathrm{~m}^{3}$. It follows from eqns 14A.5a and 14A. 6 that $\mu^{*}=4 \pi \varepsilon_{0} \alpha^{\prime} E$ and the magnitude of the dipole moment of the molecule (in addition to the permanent dipole moment) induced by an applied electric field of strength $1.0 \times 10^{5} \mathrm{~V} \mathrm{~m}^{-1}$ is

$$
\begin{aligned}
\mu^{\star}= & 4 \pi \times\left(8.854 \times 10^{-12} \mathrm{~J}^{-1} \mathrm{C}^{2} \mathrm{~m}^{-1}\right) \times\left(1.48 \times 10^{-30} \mathrm{~m}^{3}\right) \\
& \times\left(1.0 \times 10^{5} \mathrm{Vm}^{-1}\right) \\
& \\
& 1 \mathrm{~V}=1 \mathrm{JC}^{-1} \\
= & 1.65 \times 10^{-35} \mathrm{Cm}=4.9 \times 10^{-6} \mathrm{D}=4.9 \mu \mathrm{D}
\end{aligned}
$$

The polarizability volumes of some molecules are given in Table 14A.1. It is possible to establish a correlation between these values and the electronic structure of atoms and molecules.

\subsection*{How is that done? 14A. 1 Correlating polarizability and molecular structure}
The argument starts from the quantum mechanical expression for the molecular polarizability in the $z$-direction: ${ }^{1}$

$$
\alpha=2 \sum_{n \neq 0} \frac{\left|\mu_{z, 0 n}\right|^{2}}{E_{n}^{(0)}-E_{0}^{(0)}}
$$

where $\mu_{z, 0 n}=\int \psi_{n}^{\star} \hat{\mu}_{z} \psi_{0} \mathrm{~d} \tau$ is the $z$-component of the transition electric dipole moment, a measure of the extent to which electric charge is shifted when an electron migrates from the ground state to create an excited state. The sum is over the excited states, with energies $E_{n}^{(0)}$.

\subsection*{Step 1 Introduce approximations}
Now approximate the excitation energies by a mean value $\Delta E$ (an indication of the HOMO-LUMO separation). Also suppose that the most important transition dipole moment is approximately equal to the charge of an electron multiplied by the molecular radius $R$. Then

$$
\alpha \approx \frac{2 e^{2} R^{2}}{\Delta E}
$$

\footnotetext{${ }^{1}$ For a derivation of this equation see our Physical chemistry: Quanta, matter, and change (2014).
}This expression shows that $\alpha$ increases with the size of the molecule and with the ease with which it can be excited electronically. Therefore, the polarizability increases as the HOMO-LUMO separation decreases.

Step 2 Express the excitation energy in terms of the atomic radius\\
If the excitation energy is approximated by the energy needed to remove an electron to infinity from a distance $R$ from a single positive charge, then $\Delta E \approx e^{2} / 4 \pi \varepsilon_{0} R$. When you insert this value into the expression for $\alpha$ you find that

$$
\alpha \approx 2 e^{2} R^{2} \times \frac{4 \pi \varepsilon_{0} R}{e^{2}}=2 \times 4 \pi \varepsilon_{0} R^{3}
$$

To obtain the polarizability volume, divide $\alpha$ by $4 \pi \varepsilon_{0}$, and ignore the factor of 2 in this approximation. The result is $\alpha^{\prime} \approx R^{3}$, which is of the same order of magnitude as the molecular volume.

As just shown, polarizability volumes correlate with the HOMO-LUMO separations in atoms and molecules. The electron distribution can be distorted readily if the LUMO lies close to the HOMO in energy, so the polarizability is then large. If the LUMO lies high above the HOMO, an applied field cannot perturb the electron distribution significantly, and the polarizability is low. Molecules with small HOMO-LUMO gaps are typically large, and have numerous electrons.

For most molecules, the polarizability is 'anisotropic', which means that its value depends on the orientation of the molecule relative to the applied field. The polarizability volume of benzene when the field is applied perpendicular to the ring is $0.0067 \mathrm{~nm}^{3}$ and it is $0.0123 \mathrm{~nm}^{3}$ when the field is applied in the plane of the ring. The anisotropy of the polarizability determines whether a molecule is rotationally Raman active (Topic 11B).

\subsection*{14A. 3 Polarization}
The polarization, $P$, of a bulk sample is the electric dipole moment density, which is the mean electric dipole moment of the molecules, $\langle\mu\rangle$, multiplied by the number density, $\mathcal{N}=N / V$ :

$$
\begin{array}{ll}
P=\langle\mu\rangle \mathcal{N} & \begin{array}{l}
\text { Polarization } \\
\text { [definition] }
\end{array}
\end{array}
$$

(14A.7)\\
A dielectric is a polarizable, non-conducting medium.

\subsection*{(a) The frequency dependence of the polarization}
The polarization of a fluid dielectric is zero in the absence of an applied field because the molecules adopt ceaselessly changing random orientations due to thermal motion, so $\langle\mu\rangle=0$. In the presence of a weak electric field, the energy depends on the orientation of the dipole with respect to the field, with\\
lower energy orientations being more populated; as a result, the mean dipole moment is no longer zero. The Boltzmann distribution can be used to find an expression for $\langle\mu\rangle$.

\subsection*{How is that done? 14A. 2}
Deriving an expression for the mean dipole moment

In spherical polar coordinates, the area of an infinitesimal patch of surface (at constant radius) is $\sin \theta \mathrm{d} \theta \mathrm{d} \phi$. The number of molecules with their electric dipole moments pointing into this patch is proportional to this area multiplied by the Boltzmann factor $\mathrm{e}^{-E(\theta) / k T}$. If the electric field is in the $z$-direction, the energy is independent of the azimuthal angle $\phi$. The probability $\mathrm{d} p$ that a dipole moment has an orientation in the range $\theta$ to $\theta+\mathrm{d} \theta$ and at any azimuthal angle $\phi$ around the direction of the applied field is therefore

$$
\mathrm{d} p=\frac{\mathrm{e}^{-E(\theta) / k T} \sin \theta \mathrm{~d} \theta}{\int_{0}^{\pi} \mathrm{e}^{-E(\theta) / k T} \sin \theta \mathrm{~d} \theta}
$$

where $0 \leq \theta \leq \pi$. If the applied electric field $E$ is in the $z$-direction, then the dipole moment is also aligned in the $z$-direction, and its mean value is

$$
\left\langle\mu_{z}\right\rangle=\int \mu_{z} \mathrm{~d} p
$$

Step 1 Write an expression for the energy of a dipole in an electric field\\
The energy $E(\theta)$ of a dipole depends on the angle $\theta$ it makes with the electric field $\mathcal{E}$ as

$$
E(\theta)=-\mu E \cos \theta
$$

Step 2 Set up an expression for the average of the $z$-component of the dipole moment\\
The average value of the component of the dipole moment parallel to the applied electric field is

$$
\begin{aligned}
&\left\langle\mu_{z}\right\rangle=\int \overbrace{\mu \cos \theta}^{\mu_{z}} \mathrm{~d} p=\mu \frac{\int_{0}^{\pi} \cos \theta \mathrm{e}^{-E(\theta) / k T} \sin \theta \mathrm{~d} \theta}{\int_{0}^{\pi} \mathrm{e}^{-E(\theta) / k T} \sin \theta \mathrm{~d} \theta} \\
&=\mu \frac{\int_{0}^{E(\theta)=-\mu E \cos \theta}}{\int_{0}^{\pi} \mathrm{e}^{\mu E \cos \theta / k T} \cos \theta \sin \theta \mathrm{~d} \theta} \\
&=\mu \cos \theta / k T \\
& \sin \theta \mathrm{~d} \theta
\end{aligned}
$$

Step 3 Evaluate the integrals\\
To simplify the appearance of this expression, write $x=\mu E / k T$ and obtain

$$
\left\langle\mu_{z}\right\rangle=\frac{\mu \int_{0}^{\pi} \mathrm{e}^{x \cos \theta} \cos \theta \sin \theta \mathrm{~d} \theta}{\int_{0}^{\pi} \mathrm{e}^{x \cos \theta} \sin \theta \mathrm{~d} \theta}
$$

Then write $y=\cos \theta$ and $\mathrm{d} y=-\sin \theta \mathrm{d} \theta$, and change the limits of integration to $y=-1$ (at $\theta=\pi$ ) and $y=1$ (at $\theta=0$ ):

$$
\begin{aligned}
\left\langle\mu_{z}\right\rangle & =\frac{\mu \overbrace{\int_{-1}^{1} y \mathrm{e}^{x y} \mathrm{~d} y}^{\text {Integral E.4 }}}{\underbrace{\int_{-1}^{1} \mathrm{e}^{x y} \mathrm{~d} y}_{\text {Integral E.3 }}}=\mu \frac{\left\{\frac{\mathrm{e}^{x}+\mathrm{e}^{-x}}{x}-\frac{\mathrm{e}^{x}-\mathrm{e}^{-x}}{x^{2}}\right\}}{\left\{\frac{\mathrm{e}^{x}-\mathrm{e}^{-x}}{x}\right\}} \\
& =\mu\left\{\frac{\mathrm{e}^{x}+\mathrm{e}^{-x}}{\mathrm{e}^{x}-\mathrm{e}^{-x}}-\frac{1}{x}\right\}
\end{aligned}
$$

That is, noting again that $x=\mu E / k T$,

$$
\left.-\left\langle\mu_{z}\right\rangle=\mu L(x) \quad L(x)=\frac{\mathrm{e}^{x}+\mathrm{e}^{-x}}{\mathrm{e}^{x}-\mathrm{e}^{-x}}-\frac{1}{x} \right\rvert\, \quad \text { (14A.8a) } \begin{aligned}
& \text { Mean electric } \\
& \text { dipole moment }
\end{aligned}
$$

The function $L(x)$ is called the Langevin function (Fig. 14A.3). Under most circumstances, $x$ is very small. For example, if $\mu=1 \mathrm{D}$ and $T=300 \mathrm{~K}$, then $x$ exceeds 0.01 only if the field strength exceeds $100 \mathrm{kV} \mathrm{cm}^{-1}$, and most measurements are done at much lower strengths. The exponentials in the Langevin function can be expanded when the field is so weak that $x \ll 1$, and the largest term that survives is $L(x)=\frac{1}{3} x$. Therefore, the average molecular dipole moment is

$$
\left\langle\mu_{z}\right\rangle=\frac{\mu^{2} E}{3 k T} \quad \begin{aligned}
& \text { Mean value of the dipole moment } \\
& {[\text { weak electric field }]}
\end{aligned}
$$

(14A.8b)

As the electric field strength is increased to very high values, the orientations of molecular dipole moments fluctuate less about the field direction and the mean dipole moment approaches its maximum value of $\left\langle\mu_{z}\right\rangle=\mu$.\\
\includegraphics[max width=\textwidth, center]{2025_02_27_d4b95d9718f0b9e2e6b9g-621}

Figure 14A. 3 The Langevin function (in purple) used in the calculation of the mean electric dipole moment. When $x$ is small the weak-field approximation (in blue) is appropriate.

When the applied field changes direction slowly, the orientation of the permanent dipole moment has time to changethe whole molecule rotates into a new orientation-and follows the field. However, when the electric field changes direction rapidly, a molecule cannot change orientation fast enough to follow and the permanent dipole moment then makes no contribution to the polarization of the sample. The orientation polarization, the polarization arising from the permanent dipole moments, is lost at such high frequencies. Because a molecule takes about 1 ps to turn through about 1 radian in a fluid, the loss of the contribution of orientation polarization to the total polarization occurs when measurements are made at frequencies greater than about $10^{11} \mathrm{~Hz}$ (in the microwave region).

The next contribution to the polarization to be lost as the frequency increases is the distortion polarization, the polarization that arises from the distortion of the positions of the nuclei by the applied field. The molecule is bent and stretched by the applied field, and the molecular dipole moment changes accordingly. The time taken for a molecule to bend is approximately the inverse of the molecular vibrational frequency, so the distortion polarization disappears when the frequency of the radiation is increased through the infrared.

Polarization disappears in stages: each successive stage occurs as the frequency of oscillation of the electric field rises above the frequency of a particular mode of vibration. At even higher frequencies, in the visible region, only the electrons are mobile enough to respond to the rapidly changing direction of the applied field. The polarization that remains is now due entirely to the distortion of the electron distribution, and the surviving contribution to the molecular polarizability is called the electronic polarizability. This behaviour can be explored by noting that the quantum mechanical expression for the polarizability of a molecule in the presence of an electric field that is oscillating at a frequency $\omega$ in the $z$-direction is ${ }^{2}$

$$
\alpha(\omega)=\frac{2}{\hbar} \sum_{n \neq 0} \frac{\omega_{n 0}\left|\mu_{z, 0 n}\right|^{2}}{\omega_{n 0}^{2}-\omega^{2}}
$$

where the excitation frequency is defined by $\hbar \omega_{n 0}=E_{n}^{(0)}-E_{0}^{(0)}$. Two conclusions can be made:

\begin{itemize}
  \item As $\omega \rightarrow 0$, the equation reduces to the expression for the static polarizability
  \item As $\omega$ becomes very high (and much higher than any excitation frequency of the molecule so that the $\omega_{n 0}^{2}$ in the denominator can be ignored), the polarizability becomes
\end{itemize}

$$
\alpha(\omega)=-\frac{2}{\hbar \omega^{2}} \sum_{n} \omega_{z, n 0}\left|\mu_{z, 0 n}\right|^{2} \rightarrow 0 \text { as } \omega \rightarrow \infty
$$

\footnotetext{${ }^{2}$ For a derivation of this equation, see our Physical chemistry: Quanta, matter, and change (2014).
}The conclusion applies to each type of excitation, vibrational as well as electronic, and accounts for the successive decreases in polarizability as the frequency is increased.

\subsection*{(b) Molar polarization}
When two charges $Q_{1}$ and $Q_{2}$ are separated by a distance $r$ in a medium, the Coulomb potential energy of their interaction is


\begin{equation*}
V=\frac{Q_{1} Q_{2}}{4 \pi \varepsilon r} \tag{14A.9}
\end{equation*}


where $\varepsilon$ is the permittivity of the medium, which is reported by introducing the relative permittivity and writing $\varepsilon=\varepsilon_{\mathrm{r}} \varepsilon_{0}$ (The chemist's toolkit 6 in Topic 2A). The relative permittivity of a substance is measured by comparing the capacitance of a capacitor with and without the sample present ( $C$ and $C_{0}$, respectively) and using $\varepsilon_{\mathrm{r}}=C / C_{0}$. The relative permittivity can have a very significant effect on the strength of the interactions between ions in solution. For instance, water has a relative permittivity of 78 at $25^{\circ} \mathrm{C}$, so the interionic Coulombic interaction energy is reduced by nearly two orders of magnitude from its vacuum value.

The relative permittivity of a substance is large if its molecules are polar or highly polarizable. The quantitative relation between the relative permittivity and the electric properties of the molecules is obtained by considering the polarization of a medium, and is expressed by the Debye equation:


\begin{equation*}
\frac{\varepsilon_{\mathrm{r}}-1}{\varepsilon_{\mathrm{r}}+2}=\frac{\rho P_{\mathrm{m}}}{M} \tag{14A.10}
\end{equation*}


Debye equation\\
where $\rho$ is the mass density of the sample, $M$ is the molar mass of the molecules, and $P_{\mathrm{m}}$ is the molar polarization, which is defined as

$$
P_{\mathrm{m}}=\frac{N_{\mathrm{A}}}{3 \varepsilon_{0}}\left(\alpha+\frac{\mu^{2}}{3 k T}\right)
$$

Molar polarization [definition]\\
(14A.11)\\
where $\alpha$ is the polarizability. The term $\mu^{2} / 3 k T$ stems from the thermal averaging of the electric dipole moment in the presence of the applied field (eqn 14A.8b). The corresponding expression without the contribution from the permanent dipole moment is called the Clausius-Mossotti equation:


\begin{equation*}
\frac{\varepsilon_{\mathrm{r}}-1}{\varepsilon_{\mathrm{r}}+2}=\frac{\rho N_{\mathrm{A}} \alpha}{3 M \varepsilon_{0}} \tag{14A.12}
\end{equation*}


Clausius-Mossotti equation

The Clausius-Mossotti equation is used when there is no contribution from permanent electric dipole moments to the polarization, either because the molecules are nonpolar or because the frequency of the applied field is so high that the molecules cannot orientate quickly enough to follow the change in direction of the field.

\subsection*{Example 14A. 2}
Determining dipole moment and polarizability

The relative permittivity of camphor (8) was measured at a series of temperatures with the results given below. Determine the dipole moment and the polarizability volume of the molecule.\\
\includegraphics[max width=\textwidth, center]{2025_02_27_d4b95d9718f0b9e2e6b9g-623}

8 Camphor

\begin{center}
\begin{tabular}{lll}
\hline
$\theta /{ }^{\circ} \mathrm{C}$ & $\rho /\left(\mathrm{g} \mathrm{cm}^{-3}\right)$ & $\varepsilon_{\mathrm{r}}$ \\
\hline
0 & 0.99 & 12.5 \\
20 & 0.99 & 11.4 \\
40 & 0.99 & 10.8 \\
60 & 0.99 & 10.0 \\
80 & 0.99 & 9.50 \\
100 & 0.99 & 8.90 \\
120 & 0.97 & 8.10 \\
140 & 0.96 & 7.60 \\
160 & 0.95 & 7.11 \\
200 & 0.91 & 6.21 \\
\hline
\end{tabular}
\end{center}

Collect your thoughts The relative permittivity depends on the molar polarization (eqn 14A.10), which in turn depends on the temperature, polarizability, and the magnitude of the permanent dipole moment (eqn 14A.11). These relations suggest that you should

\begin{itemize}
  \item Calculate $\left(\varepsilon_{\mathrm{r}}-1\right) /\left(\varepsilon_{\mathrm{r}}+2\right)$ at each temperature, and then multiply by $M / \rho$ to form $P_{\mathrm{m}}$ from eqn 14A.10.
  \item Plot $P_{\mathrm{m}}$ against 1/T.\\
because eqn 14A. 11 rearranges to
\end{itemize}

$$
P_{\mathrm{m}}=\frac{\overbrace{\frac{N_{\mathrm{A}} \alpha}{3 \varepsilon_{0}}}^{\text {intercept }}}{\text { 虹 }}+\overbrace{N_{\mathrm{A}} \mu^{2}}^{9 \varepsilon_{0} k} \times \frac{1}{T}
$$

The slope of the graph is $N_{\mathrm{A}} \mu^{2} / 9 \varepsilon_{0} k$ and its intercept at $1 / T=$ 0 is $N_{\mathrm{A}} \alpha / 3 \varepsilon_{0}$.\\
The solution Use the data to draw up the following table, with $M=152.23 \mathrm{~g} \mathrm{~mol}^{-1}$ for camphor.

\begin{center}
\begin{tabular}{lllll}
\hline
$\theta /{ }^{\circ} \mathrm{C}$ & $\left(10^{3} \mathrm{~K}\right) / \boldsymbol{T}$ & $\varepsilon_{\mathrm{r}}$ & $\left(\varepsilon_{\mathrm{r}}-1\right) /\left(\varepsilon_{\mathrm{r}}+2\right)$ & $\boldsymbol{P}_{\mathrm{m}} /\left(\mathrm{cm}^{3} \mathrm{~mol}^{-1}\right)$ \\
\hline
0 & 3.66 & 12.5 & 0.793 & 122 \\
20 & 3.41 & 11.4 & 0.776 & 119 \\
40 & 3.19 & 10.8 & 0.766 & 118 \\
60 & 3.00 & 10.0 & 0.750 & 115 \\
80 & 2.83 & 9.50 & 0.739 & 114 \\
100 & 2.68 & 8.90 & 0.725 & 111 \\
120 & 2.54 & 8.10 & 0.703 & 110 \\
140 & 2.42 & 7.60 & 0.688 & 109 \\
160 & 2.31 & 7.11 & 0.671 & 108 \\
200 & 2.11 & 6.21 & 0.635 & 106 \\
\hline
\end{tabular}
\end{center}

\begin{center}
\includegraphics[max width=\textwidth]{2025_02_27_d4b95d9718f0b9e2e6b9g-623(1)}
\end{center}

Figure 14A. 4 The plot of $P_{\mathrm{m}} /\left(\mathrm{cm}^{3} \mathrm{~mol}^{-1}\right)$ against $\left(10^{3} \mathrm{~K}\right) / T$ used in Example 14A. 2 for the determination of the polarizability and the magnitude of the dipole moment of camphor.

The points are plotted in Fig. 14A.4. The intercept on the vertical axis lies at $P_{\mathrm{m}} /\left(\mathrm{cm}^{3} \mathrm{~mol}^{-1}\right)=83.5$, so $N_{\mathrm{A}} \alpha / 3 \varepsilon_{0}=$ $83.5 \mathrm{~cm}^{3} \mathrm{~mol}^{-1}=8.35 \times 10^{-5} \mathrm{~m}^{3} \mathrm{~mol}^{-1}$. It then follows that

$$
\begin{aligned}
\alpha & =\frac{3 \times \overbrace{\left(8.854 \times 10^{-12} \mathrm{~J}^{-1} \mathrm{C}^{2} \mathrm{~m}^{-1}\right)}^{\varepsilon_{0}}}{\underbrace{6.02 \times 10^{23} \mathrm{~mol}^{-1}}_{N_{A}}} \times \overbrace{8.35 \times 10^{-5} \mathrm{~m}^{3} \mathrm{~mol}^{-1}}^{\text {intercept }} \\
& =3.68 \times 10^{-39} \mathrm{C}^{2} \mathrm{~m}^{2} \mathrm{~J}^{-1}
\end{aligned}
$$

From eqn 14A.6, it follows that $\alpha=3.31 \times 10^{-29} \mathrm{~m}^{3}$. The slope is 10.55 , so $N_{A} \mu^{2} / 9 \varepsilon_{0} k=1.055 \times 10^{4} \mathrm{~cm}^{3} \mathrm{~mol}^{-1} \mathrm{~K}=1.055 \times 10^{-2}$ $\mathrm{m}^{3} \mathrm{~mol}^{-1} \mathrm{~K}$, so from the expression for $P_{\mathrm{m}}$ it follows that

$$
\begin{aligned}
\mu & =(\frac{9 \times \overbrace{\left(8.854 \times 10^{-12} \mathrm{~J}^{-1} \mathrm{C}^{2} \mathrm{~m}^{-1}\right)}^{\varepsilon_{0}} \times \overbrace{\left(1.381 \times 10^{-23} \mathrm{JK}^{-1}\right)}^{6.022 \times 10^{23} \mathrm{~mol}^{-1}}}{k})_{N_{\mathrm{A}}}^{1 / 2} \\
& \times \overbrace{\left(1.055 \times 10^{-2} \mathrm{~m}^{3} \mathrm{~mol}^{-1} \mathrm{~K}\right)^{1 / 2}}^{\text {slope }} \\
& =4.39 \times 10^{-30} \mathrm{C} \mathrm{~m}=1.32 \mathrm{D}
\end{aligned}
$$

Because the Debye equation describes molecules that are free to rotate, the data show that camphor, which does not melt until $175^{\circ} \mathrm{C}$, is rotating even in the solid. It is an approximately spherical molecule.

Self-test 14A. 2 The relative permittivity of chlorobenzene is 5.71 at $20^{\circ} \mathrm{C}$ and 5.62 at $25^{\circ} \mathrm{C}$. Assuming a constant density $\left(1.11 \mathrm{~g} \mathrm{~cm}^{-3}\right)$, estimate its polarizability volume and the magnitude of its dipole moment.

$$
\mathrm{CZ} \cdot \mathrm{I}{ }_{\varepsilon} \mathrm{w}_{6 \tau-} 0 \mathrm{I} \times \bigoplus^{\prime} \mathrm{I}: \iota \partial s u \forall
$$

The refractive index, $n_{r}$, of the medium is the ratio of the speed of light in a vacuum, $c$, to its speed $c^{\prime}$ in the medium: $n_{\mathrm{r}}=c / c^{\prime}$. According to Maxwell's theory of electromagnetic radiation, the refractive index at a specified (visible or ultraviolet) wavelength is related to the relative permittivity at that frequency by

$$
n_{\mathrm{r}}=\varepsilon_{\mathrm{r}}^{1 / 2}
$$

Relation between refractive index and relative permittivity

A beam of light changes direction ('bends') when it passes from a region of one refractive index to a region with a different refractive index. Therefore, the molar polarization, $P_{\mathrm{m}}$, and the molecular polarizability, $\alpha$, can be measured at frequencies typical of visible light (about $10^{15}-10^{16} \mathrm{~Hz}$ ) by measuring the refractive index of the sample and using the ClausiusMossotti equation.

\subsection*{Checklist of concepts}
$\square$ 1. An electric dipole consists of two electric charges $+Q$ and $-Q$ separated by a vector $\boldsymbol{R}$.2. The electric dipole moment $\mu$ is a vector that points from the negative charge to the positive charge of a dipole; its magnitude is $\mu=Q R$.3. A polar molecule is a molecule with a permanent electric dipole moment.4. Molecules may have higher electric multipoles: an $n$-pole is an array of point charges with an n-pole moment but no lower moment.5. The polarizability is a measure of the ability of an electric field to induce a dipole moment in a molecule.6. Polarizabilities (and polarizability volumes) correlate with the HOMO-LUMO separations in molecules.7. The polarization of a medium is the electric dipole moment density.8. Orientation polarization is the polarization arising from the permanent dipole moments.9. Distortion polarization is the polarization arising from the distortion of the positions of the nuclei by the applied field.10. Electronic polarizability is the polarizability due to the distortion of the electron distribution.

\subsection*{Checklist of equations}
\begin{center}
\begin{tabular}{llll}
\hline
Property & Equation & Comment & Equation number \\
Magnitude of the electric dipole moment & $\mu=Q R$ & Definition & 14 A .1 \\
Magnitude of the resultant of two dipole moments & $\mu_{\mathrm{res}} \approx\left(\mu_{1}^{2}+\mu_{2}^{2}+2 \mu_{1} \mu_{2} \cos \Theta\right)^{1 / 2}$ &  & 14 A .3 a \\
Magnitude of the induced dipole moment & $\mu^{*}=\alpha E$ & Linear approximation; $\alpha$ is the polarizability & 14 A .5 a \\
 & $\mu^{*}=\alpha \mathcal{E}+\frac{1}{2} \beta \mathcal{E}^{2}$ & \begin{tabular}{l}
Quadratic approximation; $\beta$ is the \\
hyperpolarizability \\
\end{tabular} &  \\
Polarizability volume & $\alpha^{\prime}=\alpha / 4 \pi \varepsilon_{0}$ & Definition & Definition \\
Polarization & $P=\langle\mu\rangle \mathcal{N}$ &  & 14 A .5 b \\
Debye equation & $\left(\varepsilon_{\mathrm{r}}-1\right) /\left(\varepsilon_{\mathrm{r}}+2\right)=\rho P_{\mathrm{m}} / M$ & 14 A .7 &  \\
Molar polarization & $P_{\mathrm{m}}=\left(N_{\mathrm{A}} / 3 \varepsilon_{0}\right)\left(\alpha+\mu^{2} / 3 k T\right)$ & 14 A .10 &  \\
Clausius-Mossotti equation & $\left(\varepsilon_{\mathrm{r}}-1\right) /\left(\varepsilon_{\mathrm{r}}+2\right)=\rho N_{\mathrm{A}} \alpha / 3 M \varepsilon_{0}$ &  & 14 A .11 \\
\hline
\end{tabular}
\end{center}

\section*{TOPIC 14B Interactions between molecules }
\subsection*{Why do you need to know this material?}
Many types of molecular interactions are responsible for the formation of condensed phases and large molecular assemblies.

\begin{itemize}
  \item What is the key idea?
\end{itemize}

Attractive interactions result in cohesion but repulsive interactions prevent the complete collapse of matter to nuclear densities.

What do you need to know already?\\
You need to be familiar with elementary aspects of electrostatics, specifically the Coulomb interaction (The chemist's toolkit 6 of Topic 2A), and the relationships between the structure and electric properties of a molecule, specifically its dipole moment and polarizability (Topic 14A).

A van der Waals interaction is an attractive (energy lowering) interaction between closed-shell molecules that depends on the separation of the molecules as the inverse sixth power $\left(V \propto 1 / r^{6}\right)$. This precise criterion is often relaxed to include all nonbonding interactions. They occur in various guises and can all be traced to the interaction of partial charges. The underlying interaction throughout this discussion is the Coulomb potential energy, $V=Q_{1} Q_{2} / 4 \pi \varepsilon r$, discussed in The chemist's toolkit 6 of Topic 2A.

\subsection*{14B. 1 The interactions of dipoles}
A point dipole is a dipole in which the separation between the charges $l$ is much smaller than the distance $r$ from which the dipole is being observed $(l \ll r)$.

\subsection*{(a) Charge-dipole interactions}
The Coulomb potential energy of one charge near another can be adapted to find the potential energy of a point charge and a dipole, and extended to the interaction between two dipoles.

How is that done? 14B. 1 Deriving the expression for the interaction between a point charge and a point dipole

You need to consider the interaction between the two charges $\pm Q_{1}$ of a point dipole, with a dipole moment of magnitude $\mu_{1}=$ $Q_{1} l$, and the point charge $Q_{2}$ as shown in (1). Suppose that the arrangement is in a vacuum, so use $\varepsilon=\varepsilon_{0}$.\\
\includegraphics[max width=\textwidth, center]{2025_02_27_d4b95d9718f0b9e2e6b9g-625}

1\\
Step 1 Write an expression for the potential energy due to interaction with both charges\\
The sum of the potential energies due to repulsion between like charges and attraction between opposite charges is

$$
V=\frac{1}{4 \pi \varepsilon_{0}}\left(-\frac{Q_{1} Q_{2}}{r-\frac{1}{2} l}+\frac{Q_{1} Q_{2}}{r+\frac{1}{2} l}\right)=\frac{Q_{1} Q_{2}}{4 \pi \varepsilon_{0} r}\left(-\frac{1}{1-x}+\frac{1}{1+x}\right)
$$

Step 2 Treat the dipole as a point dipole\\
Because $l \ll r$ for a point dipole, this expression can be simplified by expanding the terms in $x$ by using

$$
\frac{1}{1+x}=1-x+x^{2}-\cdots \quad \frac{1}{1-x}=1+x+x^{2}+\cdots
$$

and retaining only the first two terms:

$$
V=\frac{Q_{1} Q_{2}}{4 \pi \varepsilon_{0} r}\{-(1+x+\cdots)+(1-x+\cdots)\} \approx-\frac{2 x Q_{1} Q_{2}}{4 \pi \varepsilon_{0} r}=-\frac{Q_{1} Q_{2} l}{4 \pi \varepsilon_{0} r^{2}}
$$

With $\mu_{1}=Q_{1} l$ this expression becomes

$$
-\left\lvert\, V=-\frac{\mu_{1} Q_{2}}{4 \pi \varepsilon_{0} r^{2}}\right.
$$

$\qquad$\\
$\qquad$ Point dipole-point charge interaction [as in (1)]

With $\mu$ in coulomb metres, $Q_{2}$ in coulombs, and $r$ in metres, $V$ is obtained in joules. In the orientation shown in (1), $V$ is negative, representing a net attraction. The expression should be multiplied by $\cos \Theta$ when the point charge lies at an angle $\Theta$ to the axis of the dipole and $\varepsilon_{0}$ replaced by $\varepsilon$ if the medium is not a vacuum.

The potential energy approaches zero (the value at infinite separation of the charge and the dipole) more rapidly (as $1 / r^{2}$ )\\
\includegraphics[max width=\textwidth, center]{2025_02_27_d4b95d9718f0b9e2e6b9g-626}

Figure 14B. 1 There are two contributions to the diminishing field of an electric dipole with distance (here seen from the side). The potentials of the charges decrease (shown here by a fading intensity) and the two charges appear to merge, so their combined effect approaches zero more rapidly than by the distance effect alone.\\
than that between two point charges (which varies as $1 / r$ ) because, from the viewpoint of the point charge, the partial charges of the dipole seem to merge and cancel as the distance $r$ increases (Fig. 14B.1).

\subsection*{Brief illustration 14B. 1}
Consider a $\mathrm{Li}^{+}$ion and a water molecule ( $\mu=1.85 \mathrm{D}$ ) separated by 1.0 nm in a vacuum, with the point charge on the ion and the dipole of the molecule arranged as in (1). The energy of interaction is given by eqn 14B. 1 as

$$
\begin{aligned}
V & =-\frac{\overbrace{\left(1.602 \times 10^{-19} \mathrm{C}\right)}^{Q_{\mathrm{L}}+} \times \overbrace{\left(1.85 \times 3.336 \times 10^{-30} \mathrm{Cm}\right)}^{\mu_{\mathrm{H}_{2} \mathrm{O}}}}{4 \pi \times \underbrace{\left(8.854 \times 10^{-12} \mathrm{~J}^{-1} \mathrm{C}^{-1} \mathrm{~m}^{-1}\right)}_{\varepsilon_{0}} \times \underbrace{\left(1.0 \times 10^{-9} \mathrm{~m}\right)^{2}}_{r}} \\
& =-8.9 \times 10^{-21} \mathrm{~J}
\end{aligned}
$$

This energy corresponds to $-5.4 \mathrm{~kJ} \mathrm{~mol}^{-1}$.

\subsection*{(b) Dipole-dipole interactions}
The preceding discussion can be extended to the interaction of two dipoles arranged as in (2).\\
\includegraphics[max width=\textwidth, center]{2025_02_27_d4b95d9718f0b9e2e6b9g-626(1)}

How is that done? 14B. 2 Deriving the expression for the interaction energy of two point dipoles

To calculate the potential energy of interaction of two point dipoles separated by $r$ in a vacuum in the arrangement shown in (2) proceed in exactly the same way as before. In this case the total interaction energy is the sum of four pairwise terms. Two are attractions between opposite charges, which contribute negative terms to the potential energy, and two are repulsions between like charges, which contribute positive terms.

Step 1 Write an expression for the potential energy due to interaction of the charges on the two dipoles\\
The sum of the four contributions is

$$
\begin{aligned}
V & =\frac{1}{4 \pi \varepsilon_{0}}\left(-\frac{Q_{1} Q_{2}}{r+l}+\frac{Q_{1} Q_{2}}{r}+\frac{Q_{1} Q_{2}}{r}-\frac{Q_{1} Q_{2}}{r-l}\right) \\
& x=1 / r \\
& =-\frac{Q_{1} Q_{2}}{4 \pi \varepsilon_{0} r}\left(\frac{1}{1+x}-2+\frac{1}{1-x}\right)
\end{aligned}
$$

Step 2 Treat the dipoles as point dipoles\\
As before, provided $l \ll r$ the two terms in $x$ may be expanded, leading to

$$
V=-\frac{Q_{1} Q_{2}}{4 \pi \varepsilon_{0} r}(\overbrace{1-x+x^{2}+\cdots}^{1 /(1+x)}-2+\overbrace{\left.1+x+x^{2}+\cdots\right)}^{1 /(1-x)}
$$

The terms in blue sum to zero, so the only surviving term is $2 x^{2}$. It follows that

$$
V=-\frac{2 x^{2} Q_{1} Q_{2}}{4 \pi \varepsilon_{0} r}=-\frac{\underbrace{x=1 / r}}{2 l^{2} Q_{1} Q_{2}} 4 \pi \varepsilon_{0} r^{3}
$$

Because $\mu_{1}=Q_{1} l$ and $\mu_{2}=Q_{2} l$, it follows that the potential energy of interaction in the alignment shown in (2) is given by

$$
-V=-\frac{\mu_{1} \mu_{2}}{2 \pi \varepsilon_{0} r^{3}} \left\lvert\, \begin{aligned}
& \text { Point dipole-point dipole interaction } \\
& {[\text { as in (2)] }}
\end{aligned}\right.
$$

This interaction energy approaches zero more rapidly (as $1 / r^{3}$ ) than for the previous case: now both interacting entities appear neutral to each other at large separations.

Equation 14B. 2 applies only to the arrangement in (2). More generally, as in the arrangement in (3), the potential energy of interaction between two polar molecules separated by a vector $r$ is

$$
\left.-V=\frac{1}{4 \pi \varepsilon_{0} r^{3}}\left\{\mu_{1} \cdot \mu_{2}-\frac{3\left(\mu_{1} \cdot \boldsymbol{r}\right)\left(\boldsymbol{r} \cdot \boldsymbol{\mu}_{2}\right)}{r^{2}}\right\} \right\rvert\, \begin{aligned}
& \begin{array}{l}
\text { Point dipole-point } \\
\text { dipole interaction } \\
\text { [as in (3)] }
\end{array}
\end{aligned}
$$

\includegraphics[max width=\textwidth, center]{2025_02_27_d4b95d9718f0b9e2e6b9g-627(1)}\\
(For the origin of this expression, see $A$ deeper look 8 on the website of this book.) When the two dipoles are parallel and arranged as in (4), the potential energy is simply

$$
V=\frac{\mu_{1} \mu_{2} f(\Theta)}{4 \pi \varepsilon_{0} r^{3}} \quad f(\Theta)=1-3 \cos ^{2} \Theta
$$

Point dipole-point dipole interaction\\[0pt]
[as in (4)]

\subsection*{Brief illustration 14B. 2}
Equation 14B.3b can be used to calculate the potential energy of the dipolar interaction between two amide groups. Supposing that the groups are separated by 3.0 nm with $\Theta=180^{\circ}$ (so $\cos \Theta=$ -1 and $1-3 \cos ^{2} \Theta=-2$ ). Take $\mu_{1}=\mu_{2}=2.7 \mathrm{D}$, corresponding to $9.0 \times 10^{-30} \mathrm{Cm}$, and find

$$
\begin{aligned}
V & =\frac{\overbrace{\left(9.0 \times 10^{-30} \mathrm{Cm}\right)^{2}}^{\mu_{1} \mu_{2}} \times \overbrace{(-2)}^{1-3 \cos ^{2} \theta}}{4 \pi \times(\underbrace{8.854 \times 10^{-12} \mathrm{~J}^{-1} \mathrm{C}^{2} \mathrm{~m}^{-1}}_{\varepsilon_{0}}) \times \underbrace{\left.3.0 \times 10^{-9} \mathrm{~m}\right)^{3}}_{r^{3}}} \\
& =\frac{\left(9.0 \times 10^{-30}\right)^{2} \times(-2)}{4 \pi \times\left(8.854 \times 10^{-12}\right) \times\left(3.0 \times 10^{-9}\right)^{3}} \frac{\mathrm{C}^{-1} \mathrm{~J}^{2}}{\mathrm{C}^{2} \mathrm{~m}^{-1} \mathrm{~m}^{3}} \\
& =-5.4 \times 10^{-23} \mathrm{~J}
\end{aligned}
$$

This value corresponds to $-33 \mathrm{~J} \mathrm{~mol}^{-1}$.

Equation 14B.3b applies to polar molecules in a fixed, parallel, orientation in a solid. In a fluid of freely rotating molecules, the interaction between dipoles averages to zero because like partial charges of two freely rotating molecules are close together as much as the two opposite partial charges, and the repulsion of the former is cancelled by the attraction of the latter. For instance, if the dipoles are arranged as in (5) with the second free to rotate,

$$
V=\frac{1}{4 \pi \varepsilon_{0} r^{3}}\left\{\mu_{1} \mu_{2} \cos \theta-3 \mu_{1} \mu_{2} \cos \theta\right\}=-\frac{\mu_{1} \mu_{2}}{2 \pi \varepsilon_{0} r^{3}} \cos \theta
$$

\begin{center}
\includegraphics[max width=\textwidth]{2025_02_27_d4b95d9718f0b9e2e6b9g-627}
\end{center}

The average value of $\cos \theta$ is zero, because

$$
\int_{0}^{2 \pi} \int_{0}^{\pi} \cos \theta \overbrace{\sin \theta \mathrm{d} \theta \mathrm{~d} \phi}^{\begin{array}{c}
\text { spherical polar } \\
\text { coordinates }
\end{array}}=\overbrace{\int_{0}^{2 \pi} \mathrm{~d} \phi}^{\begin{array}{c}
\text { Surface area } \\
\text { element in }
\end{array}} \times \overbrace{-1}^{2 \pi} x \mathrm{~d} x ; x ; \mathrm{d} \cos \theta=-\sin \theta \mathrm{d} \theta
$$

Therefore, the average energy of interaction is also zero. This conclusion is applicable at any relative location of the two dipoles.

The average interaction energy of two freely rotating dipoles is zero. However, because their mutual potential energy depends on their relative orientation, the molecules do not in fact rotate completely freely, even in a gas. The lower energy orientations are marginally favoured, so there is a non-zero average interaction between polar molecules. The detailed calculation of the interaction energy between two polar molecules is quite complicated, but the form of the final answer can be constructed quite simply.

\subsection*{How is that done? 14B. 3}
Deriving the expression for the energy of interaction between rotating polar molecules

You can use the simplified model of the interaction for the arrangement shown in (5), with the second dipole free to rotate, but not sampling all orientations equally.

Step 1 Write an expression for the average interaction energy The average interaction energy of two polar molecules rotating at a fixed separation $r$ is given by

$$
\langle V\rangle=-\frac{\mu_{1} \mu_{2}}{2 \pi \varepsilon_{0} r^{3}}\langle\cos \theta\rangle
$$

where $\langle\cos \theta\rangle$ now includes a weighting factor in the averaging that recognizes that not all orientations are equally probable.

Step 2 Write an expression for the probability of finding a particular orientation\\
The probability that dipole 2 lies in a patch of orientation $\sin \theta \mathrm{d} \theta \mathrm{d} \phi$ of the surface of a sphere (Fig. 14B.2) is

$$
\mathrm{d} p=\frac{\mathrm{e}^{-V(\theta) / k T} \sin \theta \mathrm{~d} \theta \mathrm{~d} \phi}{\int_{0}^{2 \pi} \int_{0}^{\pi} \mathrm{e}^{-V(\theta) / k T} \sin \theta \mathrm{~d} \theta \mathrm{~d} \phi} \quad V(\theta)=-\frac{\mu_{1} \mu_{2}}{2 \pi \varepsilon_{0} r^{3}} \cos \theta
$$

\begin{center}
\includegraphics[max width=\textwidth]{2025_02_27_d4b95d9718f0b9e2e6b9g-628}
\end{center}

Figure 14B.2 The surface of a sphere showing the area element $\sin \theta \mathrm{d} \theta \mathrm{d} \phi$.\\
where $\mathrm{e}^{-V(\theta) / k T}$ is the Boltzmann factor. The (weighted) average value of $\cos \theta$ is

$$
\langle\cos \theta\rangle=\int \cos \theta \mathrm{d} p=\frac{\int_{0}^{2 \pi} \int_{0}^{\pi} \mathrm{e}^{-V(\theta) / k T} \cos \theta \sin \theta \mathrm{~d} \theta \mathrm{~d} \phi}{\int_{0}^{2 \pi} \int_{0}^{\pi} \mathrm{e}^{-V(\theta) / k T} \sin \theta \mathrm{~d} \theta \mathrm{~d} \phi}
$$

Step 3 Evaluate the integrals\\
With $a=\mu_{1} \mu_{2} / 2 \pi k T \varepsilon_{0} r^{3}$ the denominator (after integration over $\phi$, which gives a factor of $2 \pi$, and writing $x=\cos \theta$ ) gives

$$
2 \pi \int_{0}^{\pi} \mathrm{e}^{-a \cos \theta} \overbrace{\sin \theta \mathrm{~d} \theta}^{-\mathrm{d} \cos \theta}=2 \pi \int_{-1}^{1} \mathrm{e}^{-a x} \mathrm{~d} x=-\frac{2 \pi}{a}\left(\mathrm{e}^{a}-\mathrm{e}^{-a}\right)
$$

Likewise, the numerator is

$$
\begin{aligned}
2 \pi \int_{0}^{\pi} \mathrm{e}^{-a \cos \theta} \cos \theta \sin \theta \mathrm{~d} \theta & =2 \pi \int_{-1}^{1} x \mathrm{e}^{-a x} \mathrm{~d} x \\
& =\frac{2 \pi}{a^{2}}\left(\mathrm{e}^{a}-\mathrm{e}^{-a}\right)-\frac{2 \pi}{a}\left(\mathrm{e}^{a}+\mathrm{e}^{-a}\right)
\end{aligned}
$$

It follows that

$$
\langle\cos \theta\rangle=-\frac{1}{a}+\frac{\mathrm{e}^{a}+\mathrm{e}^{-a}}{\mathrm{e}^{a}-\mathrm{e}^{-a}}=L(a)
$$

where $L(a)$ is the Langevin function introduced in Topic 14A. Therefore,

$$
\langle V\rangle=-\frac{\mu_{1} \mu_{2}}{2 \pi \varepsilon_{0} r^{3}} L(a)
$$

As in Topic 14A, for $a \ll 1, L(a)=a / 3$, so

$$
\langle V\rangle=-\frac{a \mu_{1} \mu_{2}}{6 \pi \varepsilon_{0} r^{3}}=-\frac{\mu_{1}^{2} \mu_{2}^{2}}{12 \pi^{2} \varepsilon_{0}^{2} k T r^{6}}
$$

In a more realistic calculation, where the second dipole is allowed to roam around the first (at the same distance), a further factor of $\frac{1}{2}$ is introduced, and the final outcome is the Keesom interaction:\\
$\left.-\langle V\rangle=-\frac{C}{r^{6}} \quad C=\frac{2 \mu_{1}^{2} \mu_{2}^{2}}{3\left(4 \pi \varepsilon_{0}\right)^{2} k T} \right\rvert\, \begin{array}{r}\text { (14B.4) } \\ \end{array}$

The important features of eqn 14B. 4 are:

\begin{itemize}
  \item The negative sign shows that the average interaction is attractive.
  \item The dependence of the average interaction energy on the inverse sixth power of the separation identifies it as a van der Waals interaction.
  \item The inverse dependence on the temperature reflects how the greater thermal motion overcomes the mutual orientating effects of the dipoles at higher temperatures.
  \item The inverse sixth power arises from the inverse third power of the interaction potential energy weighted by the energy in the Boltzmann term, which is also proportional to the inverse third power of the separation.
\end{itemize}

\subsection*{Brief illustration 14B. 3}
Suppose a water molecule ( $\mu_{1}=1.85 \mathrm{D}$ ) can rotate 1.0 nm from an amide group ( $\mu_{2}=2.7 \mathrm{D}$ ). The average energy of their interaction at $25^{\circ} \mathrm{C}(298 \mathrm{~K})$ is

$$
\begin{aligned}
V & =-\frac{2 \times(\overbrace{\left(1.85 \times 3.336 \times 10^{-30} \mathrm{Cm}\right.})^{2} \times(\overbrace{\left.2.7 \times 3.336 \times 10^{-30} \mathrm{Cm}\right)^{2}}^{\mu_{1}}}{3 \times \underbrace{\left.1.710 \times 10^{-43} \mathrm{~J}^{-1} \mathrm{C}^{4} \mathrm{~m}^{-2} \mathrm{~K}^{-1}\right)}_{\left(4 \pi \varepsilon_{0}\right)^{2} k} \times \underbrace{(\underbrace{298 \mathrm{~K})}_{2} \times(\underbrace{1.0 \times 10^{-9} \mathrm{~m}}_{r})^{\epsilon}}_{T}} \\
& =-4.0 \times 10^{-23} \mathrm{~J}
\end{aligned}
$$

This interaction energy corresponds (after multiplication by Avogadro's constant) to $-24 \mathrm{~J} \mathrm{~mol}^{-1}$, and it is much smaller than the energies involved in the making and breaking of chemical bonds.

Table 14B. 1 summarizes the various expressions for the interaction of charges and dipoles. It is quite easy to extend the formulas given there to obtain expressions for the energy of interaction of higher multipoles (electric multipoles are described in Topic 14A). The feature to remember is that the interaction energy approaches zero more rapidly the higher the order of the multipole. For the interaction of a stationary $n$ pole with a stationary $m$-pole, the potential energy varies with distance as

$$
V \propto \frac{1}{r^{n+m-1}}
$$

\footnotetext{Energy of interaction between
} stationary multipoles\\
(14B.5)

Table 14B. 1 Interaction potential energies\\
\$\textbackslash left.\[
\begin{array}{llcl}\hline \begin{array}{l}\text { Interaction } \\
\text { type }\end{array} & \begin{array}{l}\text { Distance } \\
\text { dependence of } \\
\text { potential energy }\end{array} & \begin{array}{l}\text { Typical } \\
\text { energy/ } \\
\left(\mathbf{k J ~ m o l}^{-1}\right)\end{array} & \text { Comment } \\
\hline \text { Ion-ion } & 1 / r & 250 & \begin{array}{l}\text { Only between ions } \\
\text { Hydrogen bond }\end{array}
\] \\
\\
\text{Occurs in X-H...Y,} \\
\\
\text{where X, Y = N,} \\


\text{O, or F}\textbackslash end\{array\}\textbackslash right]\$\begin{tabular}{l}
Ion-dipole \\
Dipole-dipole \\
$1 / r^{2}$ \\
\end{tabular}

The reason for the even steeper decrease with distance is the same as before: the array of charges appears to blend together into neutrality more rapidly with distance the higher the number of individual charges that contribute to the multipole. Note that a given molecule may have a charge distribution that corresponds to a combination of several different multipoles, and in such cases the energy of interaction is the sum of terms given by eqn 14B.5.

\subsection*{(c) Dipole-induced dipole interactions}
A polar molecule can induce a dipole in a neighbouring polarizable molecule (Fig. 14B.3). The induced dipole interacts with the permanent dipole of the first molecule, and the two are attracted together. The average interaction energy when the separation of the centres of the molecules is $r$ is

$$
V=-\frac{C}{r^{6}} \quad C=\frac{\mu_{1}^{2} \alpha_{2}^{\prime}}{4 \pi \varepsilon_{0}}
$$

Potential energy of a polar molecule and a polarizable molecule\\
(14B.6)\\
where $\alpha_{2}^{\prime}$ is the polarizability volume (Topic 14A) of molecule 2 and $\mu_{1}$ is the magnitude of the permanent dipole moment of molecule 1 . Note that the $C$ in this expression is different from the $C$ in eqn 14B. 4 and other expressions below: the use of the same symbol in $C / r^{6}$ emphasizes the similarity of form of each expression.\\
\includegraphics[max width=\textwidth, center]{2025_02_27_d4b95d9718f0b9e2e6b9g-629}

Figure 14B. 3 (a) A polar molecule (yellow arrow) can induce a dipole (grey arrow) in a nonpolar molecule, and (b) the orientation of the latter follows that of the former, so the interaction does not average to zero.

The dipole-induced dipole interaction energy is independent of the temperature because thermal motion has no effect on the averaging process. Moreover, like the dipole-dipole interaction, the potential energy depends on $1 / r^{6}$ : this distance dependence stems from the $1 / r^{3}$ dependence of the distorting electric field of molecule 1 (and hence the magnitude of the dipole induced in molecule 2) and the $1 / r^{3}$ dependence of the potential energy of interaction between the permanent and induced dipoles.

\subsection*{Brief illustration 14B.4}
For a molecule with $\mu=1.0 \mathrm{D}\left(3.3 \times 10^{-30} \mathrm{Cm}\right.$, such as HCl$)$ separated by 0.30 nm from a molecule of polarizability volume $\alpha^{\prime}=10 \times 10^{-30} \mathrm{~m}^{3}$ (such as benzene, Table 14A.1), the average interaction energy is

$$
\begin{aligned}
V & =-\frac{\left(3.3 \times 10^{-30} \mathrm{Cm}\right)^{2} \times\left(10 \times 10^{-30} \mathrm{~m}^{3}\right)}{4 \pi \times\left(8.854 \times 10^{-12} \mathrm{~J}^{-1} \mathrm{C}^{2} \mathrm{~m}^{-1}\right) \times\left(3.0 \times 10^{-10} \mathrm{~m}\right)^{6}} \\
& =-1.4 \times 10^{-21} \mathrm{~J}
\end{aligned}
$$

which, upon multiplication by Avogadro's constant, corresponds to $-0.83 \mathrm{~kJ} \mathrm{~mol}^{-1}$.

\subsection*{(d) Induced dipole-induced dipole interactions}
Nonpolar molecules (including closed-shell atoms, such as Ar) attract one another even though neither has a permanent dipole moment. The abundant evidence for the existence of interactions between nonpolar molecules is their ability to exist as condensed phases, such as liquid hydrogen or argon and the fact that benzene is a liquid at normal temperatures.

The interaction between nonpolar molecules arises from the transient dipoles which all molecules possess as a result of fluctuations in the instantaneous positions of electrons. To appreciate the origin of the interaction, suppose that the electrons in one molecule flicker into an arrangement that gives the molecule an instantaneous dipole moment $\boldsymbol{\mu}_{1}^{*}$. This dipole generates an electric field which polarizes the other molecule, and induces in that molecule an instantaneous dipole moment $\mu_{2}$. The two dipoles attract each other and the potential energy of the pair is lowered. Although the first molecule will go on to change the size and direction of its instantaneous dipole, the electron distribution of the second molecule will follow; that is, the two dipoles are correlated in direction (Fig. 14B.4). Because of this correlation, the attraction between the two instantaneous dipoles does not average to zero, and gives rise to an induced dipole-induced dipole interaction. This interaction is called either the dispersion\\
\includegraphics[max width=\textwidth, center]{2025_02_27_d4b95d9718f0b9e2e6b9g-630(2)}\\
\includegraphics[max width=\textwidth, center]{2025_02_27_d4b95d9718f0b9e2e6b9g-630(1)}

Figure 14B. 4 (a) In the dispersion interaction, an instantaneous dipole on one molecule induces a dipole on another molecule, and the two dipoles then interact to lower the energy. (b) The two instantaneous dipoles are correlated, and although they occur in different orientations at different instants, the interaction does not average to zero.\\
interaction or the London interaction (for Fritz London, who first described it).

The strength of the dispersion interaction depends on the polarizability of the first molecule because the instantaneous dipole moment of magnitude $\mu_{1}^{*}$ depends on the looseness of the control that the nuclear charge exercises over the outer electrons. The strength of the interaction also depends on the polarizability of the second molecule, for that polarizability determines how readily a dipole can be induced by the electric field of another molecule. The actual calculation of the dispersion interaction is quite involved, but a reasonable approximation to the interaction energy is given by the London formula:

$$
V=-\frac{C}{r^{6}} \quad C=\frac{3}{2} \alpha_{1}^{\prime} \alpha_{2}^{\prime} \frac{I_{1} I_{2}}{I_{1}+I_{2}}
$$

London formula (14B.7)\\
where $I_{1}$ and $I_{2}$ are the ionization energies of the two molecules. This interaction energy is also proportional to the inverse sixth power of the separation of the molecules, which identifies it as a third contribution to the van der Waals interaction. The dispersion interaction generally dominates all the interactions between molecules other than hydrogen bonds.

\subsection*{Brief illustration 14B. 5}
For two $\mathrm{CH}_{4}$ molecules separated by 0.30 nm , use eqn 14B. 7 with $\alpha^{\prime}=2.6 \times 10^{-30} \mathrm{~m}^{3}$ and $I \approx 700 \mathrm{~kJ} \mathrm{~mol}^{-1}$ and obtain

$$
\begin{aligned}
V & =-\frac{\frac{3}{2} \times\left(2.6 \times 10^{-30} \mathrm{~m}^{3}\right)^{2}}{\left(0.30 \times 10^{-9} \mathrm{~m}\right)^{6}} \times \frac{\left(7.00 \times 10^{5} \mathrm{Jmol}^{-1}\right)^{2}}{2 \times\left(7.00 \times 10^{5} \mathrm{Jmol}^{-1}\right)} \\
& =-4.9 \mathrm{~kJ} \mathrm{~mol}^{-1}
\end{aligned}
$$

A very approximate check on this figure is the enthalpy of vaporization of methane, which is $8.2 \mathrm{~kJ} \mathrm{~mol}^{-1}$. However, this comparison is questionable, partly because the total energy of interaction between molecules in a liquid is not due only to pairwise interactions and partly because the long-distance assumption breaks down.

\subsection*{14B.2 Hydrogen bonding}
The interactions described so far are universal in the sense that they are possessed by all molecules independent of their specific identity. However, there is a type of interaction possessed by molecules that have a particular constitution. A hydrogen bond is an attractive interaction between two species that arises from a link of the form $\mathrm{A}-\mathrm{H} \cdots \mathrm{B}$, where A and B are highly electronegative elements and $B$ possesses a lone pair of electrons. Hydrogen bonding is conventionally regarded as being limited to $\mathrm{N}, \mathrm{O}$, and F but if B is an anionic species (such as $\mathrm{Cl}^{-}$) it may also participate in hydrogen bonding. There is no strict cut-off for an ability to participate in hydrogen bonding, but N, O, and F participate most effectively.\\
The formation of a hydrogen bond can be regarded either as the approach between a partial positive charge of H and a partial negative charge of B or as a particular example of delocalized molecular orbital formation in which $A, H$, and $B$ each supply one atomic orbital from which three molecular orbitals are constructed (Fig. 14B.5). Experimental evidence and theoretical arguments have been presented in favour of both views and the matter has not yet been resolved.

In the molecular orbital model, the $\mathrm{A}-\mathrm{H}$ bond is regarded as formed from the overlap of an orbital on $A, \psi_{A}$, and a hydrogen 1 s orbital, $\psi_{\mathrm{H}}$, and the orbital on $\mathrm{B}, \psi_{\mathrm{B}}$, which is occupied by a lone pair. When the two molecules are close together, build three molecular orbitals are built from the three basis orbitals and writing: $\psi=c_{1} \psi_{\mathrm{A}}+c_{2} \psi_{\mathrm{H}}+c_{3} \psi_{\mathrm{B}}$. One of the molecular orbitals is bonding, one almost nonbonding, and the third antibonding (Topic 9E). These three orbitals need to accommodate four electrons (two from the original A-H bond and two from the lone pair of B), so two enter the bonding orbital and two enter the nonbonding orbital. Because the\\
\includegraphics[max width=\textwidth, center]{2025_02_27_d4b95d9718f0b9e2e6b9g-630}

Figure 14B. 5 The molecular orbital interpretation of the formation of an $A-H \cdots B$ hydrogen bond. From the three $A, H$, and $B$ orbitals, three molecular orbitals can be formed (their relative contributions are represented by the sizes of the spheres). Only the two lower energy orbitals are occupied, and there may therefore be a net lowering of energy compared with the separate AH and B species.\\
antibonding orbital remains empty, the net effect-depending on the precise energy of the almost nonbonding orbitalmay be a lowering of energy.

In practice, the strength of the bond is found to be about $20 \mathrm{~kJ} \mathrm{~mol}^{-1}$ (there are two hydrogen bonds per molecule in liquid water, and its standard enthalpy of vaporization, from Table 2C.1, is $44 \mathrm{~kJ} \mathrm{~mol}^{-1}$ ). Because the bonding depends on orbital overlap, it is a contact-like interaction that is turned on when AH touches B and is zero as soon as the contact is broken. If hydrogen bonding is present, it dominates the other intermolecular interactions. The properties of liquid and solid water, for example, are dominated by the hydrogen bonding between $\mathrm{H}_{2} \mathrm{O}$ molecules. The structural evidence for hydrogen bonding comes from noting that the internuclear distance between formally non-bonded atoms is less than expected on the basis of their van der Waals radii, the radii based on the closest approach of non-bonded atoms, which suggests that a dominating attractive interaction is present. For example, the $\mathrm{O}-\mathrm{O}$ distance in $\mathrm{O}-\mathrm{H} \cdots \mathrm{O}$ is expected to be 280 pm on the basis of van der Waals radii, but is found to be 270 pm in typical compounds. Moreover, the $\mathrm{H} \cdots \mathrm{O}$ distance is expected to be 260 pm but is found to be only 170 pm .

Hydrogen bonds may be either symmetric or non-symmetric. In a symmetric hydrogen bond, the H atom lies midway between the two other atoms. This arrangement is rare, but occurs in $\mathrm{F}-\mathrm{H} \cdots \mathrm{F}^{-}$, where both bond lengths are 120 pm . More common is the non-symmetric arrangement, where the $\mathrm{A}-\mathrm{H}$ bond is shorter than the H$\cdots B$ bond. Simple electrostatic arguments, treating A-H $\cdots \mathrm{B}$ as an array of point charges (partial negative charges on $A$ and $B$, partial positive on $H$ ), suggest that the lowest energy is achieved when the bond is linear, because then the two partial negative charges are farthest apart. The experimental evidence from structural studies supports a linear or near-linear arrangement.

\subsection*{Brief illustration 14B. 6}
A common hydrogen bond is that formed between $\mathrm{O}-\mathrm{H}$ groups and O atoms, as in liquid water and ice. In Problem P14B.8, you are invited to use the electrostatic model to calculate the dependence of\\
\includegraphics[max width=\textwidth, center]{2025_02_27_d4b95d9718f0b9e2e6b9g-631}

6 the potential energy of interaction on the OOH angle, denoted $\Theta$ in (6), and the results are plotted in Fig. 14B.6. The strength of bonding is greatest at $\Theta=0$ when the OHO atoms lie in a straight line; the molar potential energy is then $-19 \mathrm{~kJ} \mathrm{~mol}^{-1}$. Note that the interaction energy is negative (and the interaction is attractive) only between $-12^{\circ}$ and $+12^{\circ}$, so the atoms adopt a nearly linear arrangement.\\
\includegraphics[max width=\textwidth, center]{2025_02_27_d4b95d9718f0b9e2e6b9g-631(1)}

Figure 14B.6 The variation of the energy of interaction (according to the electrostatic model) of a hydrogen bond as the angle between the $\mathrm{O}-\mathrm{H}$ and $: \mathrm{O}$ groups is changed.

\subsection*{14B.3 The total interaction}
Consider molecules that are unable to participate in the formation of a hydrogen bond. The total favourable (energy lowering) interaction energy between rotating molecules is then the sum of the dipole-dipole, dipole-induced dipole, and dispersion interactions. Only the dispersion interaction contributes if both molecules are nonpolar. In a fluid phase, all three contributions to the potential energy vary as the inverse sixth power of the separation of the molecules, so for all of them and their sum


\begin{equation*}
V=-\frac{C_{6}}{r^{6}} \tag{14B.8}
\end{equation*}


where $C_{6}$ is a coefficient that depends on the identity of the molecules.

Although attractive interactions between molecules are often expressed as in eqn 14B.8, remember that this equation has only limited validity. First, only dipolar interactions of various kinds are taken into account, for they have the longest range and are dominant if the average separation of the molecules is large. However, in a complete treatment, quadrupolar and higher-order multipole interactions should also be considered, particularly if the molecules do not have permanent dipole moments. Secondly, the expressions have been derived by assuming that the molecules can rotate reasonably freely. That is not the case in most solids, and in rigid media the dipole-dipole interaction is proportional to $1 / r^{3}$ (as in eqn 14B.3b) because the Boltzmann averaging procedure is irrelevant when the molecules are trapped into a fixed orientation.

A different kind of limitation is that eqn 14B. 8 relates to the interactions of pairs of molecules. There is no reason to suppose that the energy of interaction of three (or more) molecules is the sum of the pairwise interaction energies alone.

The total dispersion energy of three closed-shell atoms, for instance, is given approximately by the Axilrod-Teller formula:

\[
V=-\frac{C_{6}}{r_{\mathrm{AB}}^{6}}-\frac{C_{6}}{r_{\mathrm{BC}}^{6}}-\frac{C_{6}}{r_{\mathrm{CA}}^{6}}+\frac{C^{\prime}}{\left(r_{\mathrm{AB}} r_{\mathrm{BC}} r_{\mathrm{CA}}\right)^{3}} \quad \begin{align*}
& \text { Axilrod-Teller }  \tag{14B.9a}\\
& \text { formula }
\end{align*}
\]

where

$$
C^{\prime}=a\left(3 \cos \theta_{\mathrm{A}} \cos \theta_{\mathrm{B}} \cos \theta_{\mathrm{C}}+1\right)
$$

(14B.9b)\\
The parameter $a$ is approximately equal to $\frac{3}{4} \alpha^{\prime} C_{6}$; the angles $\theta$ are the internal angles of the triangle formed by the three atoms (7). The term in $C^{\prime}$ (which represents the non-additivity of the pairwise interactions) is negative for a lin-\\
\includegraphics[max width=\textwidth]{2025_02_27_d4b95d9718f0b9e2e6b9g-632} ear arrangement of atoms (so that arrangement is stabilized) and positive for an equilateral triangular cluster (so that arrangement is destabilized). The three-body term contributes about 10 per cent of the total interaction energy in liquid argon.

When molecules are squeezed together, the nuclear and electronic repulsions begin to dominate the attractive forces. The repulsions increase steeply with decreasing separation in a way that can be deduced only by very extensive, complicated molecular structure calculations of the kind described in Topic 9E (Fig. 14B.7).

In many cases, however, progress can be made by using a greatly simplified representation of the potential energy, where the details are ignored and the general features expressed by a few adjustable parameters. One such approximation is the hard-sphere potential energy, in which it is assumed that the potential energy rises abruptly to infinity as soon as the particles come within a separation $d$ :

\[
V= \begin{cases}\infty \text { for } r \leq d & \text { Hard-sphere }  \tag{14B.10}\\ 0 \text { for } r>d & \text { potential energy }\end{cases}
\]

\begin{center}
\includegraphics[max width=\textwidth]{2025_02_27_d4b95d9718f0b9e2e6b9g-632(1)}
\end{center}

Figure 14B. 7 The general form of an intermolecular potential energy curve (the graph of the potential energy of two closed shell species as the distance between them is changed). The attractive (negative) contribution has a long range, but the repulsive (positive) interaction increases more sharply once the molecules come into contact.\\
\includegraphics[max width=\textwidth, center]{2025_02_27_d4b95d9718f0b9e2e6b9g-632(2)}

Fig 14B. 8 The Lennard-Jones potential energy.

This very simple expression for the potential energy is surprisingly useful for assessing a number of properties. Another widely used approximation is the Mie potential energy:

\[
V=\frac{C_{n}}{r^{n}}-\frac{C_{m}}{r^{m}} \quad \begin{align*}
& \text { Mie potential }  \tag{14B.11}\\
& \text { energy }
\end{align*}
\]

with $n>m$. The first term represents repulsions and the second term attractions. The Lennard-Jones potential energy is a special case of the Mie potential energy with $n=12$ and $m=6$ (Fig. 14B.8); it is often written in the form

\[
V=4 \varepsilon\left\{\left(\frac{r_{0}}{r}\right)^{12}-\left(\frac{r_{0}}{r}\right)^{6}\right\} \quad \begin{align*}
& \text { Lennard-Jones }  \tag{14B.12}\\
& \text { potential energy }
\end{align*}
\]

The two parameters are $\varepsilon$, the depth of the well, and $r_{0}$, the separation other than infinity at which $V=0$ (Table 14B.2).

Although the Lennard-Jones potential energy has been used in many calculations, there is plenty of evidence to show that $1 / r^{12}$ is a very poor representation of the repulsive potential energy, and that an exponential form, $\mathrm{e}^{-r / r_{0}}$, is greatly superior. An exponential function is more faithful to the exponential decay of atomic wavefunctions at large distances, and hence to the overlap that is responsible for repulsion. The potential energy with an exponential repulsive term and a $1 / r^{6}$ attractive term is known as an exp-6 potential energy. These expressions for the

Table 14B. 2 Lennard-Jones-( 12,6 ) potential energy parameters*

\begin{center}
\begin{tabular}{lcl}
\hline
 & $(\varepsilon / k) / \mathrm{K}$ & $r_{0} / \mathbf{p m}$ \\
\hline
Ar & 111.84 & 362.3 \\
$\mathrm{BF}_{2}$ & 104.29 & 357.1 \\
$\mathrm{C}_{6} \mathrm{H}_{6}$ & 377.46 & 617.4 \\
$\mathrm{Cl}_{2}$ & 296.27 & 448.5 \\
$\mathrm{~N}_{2}$ & 91.85 & 391.9 \\
$\mathrm{O}_{2}$ & 113.27 & 365.4 \\
Xe & 213.96 & 426 \\
\hline
\end{tabular}
\end{center}

\footnotetext{\begin{itemize}
  \item More values are given in the Resource section.
\end{itemize}
}
potential energy can be used to calculate the virial coefficients of gases, as explained in Topics 1C and 13D, and through them various properties of real gases. They are also used to model the structures of condensed fluids.

With the advent of atomic force microscopy (AFM), in which the force between a molecular-sized probe and a surface is monitored (Topic 19A), it has become possible to measure directly the forces acting between molecules. The force, $F$, is the negative slope of the potential energy, so for the LennardJones potential energy between individual molecules


\begin{equation*}
F=-\frac{\mathrm{d} V}{\mathrm{~d} r}=\frac{24 \varepsilon}{r_{0}}\left\{2\left(\frac{r_{0}}{r}\right)^{13}-\left(\frac{r_{0}}{r}\right)^{7}\right\} \tag{14B.13}
\end{equation*}


Example 14B. 1 Calculating an intermolecular force from the Lennard-Jones potential energy

Use the expression for the Lennard-Jones potential energy to estimate the greatest net attractive force between two $\mathrm{N}_{2}$ molecules.

Collect your thoughts The force is greatest when $\mathrm{d} F / \mathrm{d} r=0$. Therefore you need to differentiate eqn 14B. 13 with respect to $r$, set the resulting expression to zero, and then solve for $r$. Finally, use the value of $r$ in eqn 14B. 13 to calculate the corresponding value of $F$.

The solution Because $\mathrm{d} x^{n} / \mathrm{d} x=n x^{n-1}$,

$$
\frac{\mathrm{d} F}{\mathrm{~d} r}=\frac{24 \varepsilon}{r_{0}}\left\{2\left(-\frac{13 r_{0}^{13}}{r^{14}}\right)-\left(-\frac{7 r_{0}^{7}}{r^{8}}\right)\right\}=24 \varepsilon r_{0}^{6}\left\{\frac{7}{r^{8}}-\frac{26 r_{0}^{6}}{r^{14}}\right\}
$$

It follows that $\mathrm{d} F / \mathrm{d} r=0$ when

$$
\frac{7}{r^{8}}-\frac{26 r_{0}^{6}}{r^{14}}=0 \text { or } 7 r^{6}-26 r_{0}^{6}=0
$$

That is, when

$$
r=\left(\frac{26}{7}\right)^{1 / 6} r_{0}=1.244 r_{0}
$$

At this separation the force is

$$
F=\frac{24 \varepsilon}{r_{0}}\left\{2\left(\frac{r_{0}}{1.244 r_{0}}\right)^{13}-\left(\frac{r_{0}}{1.244 r_{0}}\right)^{7}\right\}=-\frac{2.396 \varepsilon}{r_{0}}
$$

From Table 14B.2, $\varepsilon=1.268 \times 10^{-21} \mathrm{~J}$ and $r_{0}=3.919 \times 10^{-10} \mathrm{~m}$. It follows that

$$
F=-\frac{2.396 \times\left(1.268 \times 10^{-21} \mathrm{~J}\right)}{3.919 \times 10^{-10} \mathrm{~m}} \stackrel{1 \mathrm{~N}=1 \mathrm{Jm}^{-1}}{=}-7.752 \times 10^{-12} \mathrm{~N}
$$

That is, the magnitude of the force is about 8 pN .\\
Self-test 14B. 1 At what separation $r$ does a Lennard-Jones potential energy have its minimum value?\\
${ }^{0} \iota_{9 / 1} \tau=\boldsymbol{\imath}: \iota \partial$ мsu ${ }_{V}$

\subsection*{Checklist of concepts}
$\square$ 1. A van der Waals interaction is an attractive interaction between closed-shell molecules; the corresponding potential energy is inversely proportional to the sixth power of their separation.2. The following molecular interactions are important: charge-dipole, dipole-dipole, dipole-induced dipole, dispersion (London), and hydrogen bonding.\\
$\square 3$ The van der Waals radius of an atom is based on the closest approach of non-bonded atoms.4. A hydrogen bond is an interaction of the form $\mathrm{A}-\mathrm{H} \cdots \mathrm{B}$, where A and B are typically $\mathrm{N}, \mathrm{O}$, or F .5. The Lennard-Jones potential energy is a model of the total intermolecular potential energy, including repulsion.

\subsection*{Checklist of equations}
\begin{center}
\begin{tabular}{lll}
\hline
Property & Equation & Comment \\
\hline
Energy of interaction between a point dipole and a point charge & $V=-\mu_{1} Q_{2} / 4 \pi \varepsilon_{0} r^{2}$ & Linear arrangement \\
Energy of interaction between two fixed dipoles & $V=\mu_{1} \mu_{2} f(\Theta) / 4 \pi \varepsilon_{0} r^{3}$, & Parallel dipoles \\
 & $f(\Theta)=1-3 \cos ^{2} \Theta$ & 14 B .1 \\
Energy of interaction between two rotating dipoles & $V=-2 \mu_{1}^{2} \mu_{2}^{2} / 3\left(4 \pi \varepsilon_{0}\right)^{2} k T r^{6}$ &  \\
Energy of interaction between a polar molecule and a polarizable molecule & $V=-\mu_{1}^{2} \alpha_{2}^{\prime} / 4 \pi \varepsilon_{0} r^{6}$ & 14 B .4 \\
London formula & $V=-\frac{3}{2} \alpha_{1}^{\prime} \alpha_{2}^{\prime}\left(I_{1} I_{2} /\left(I_{1}+I_{2}\right)\right) / r^{6}$ & 14 B .6 \\
Lennard-Jones potential energy & $V=4 \varepsilon\left\{\left(r_{0} / r\right)^{12}-\left(r_{0} / r\right)^{6}\right\}$ & 14 B .7 \\
\hline
\end{tabular}
\end{center}

\section*{TOPIC 14C Liquids}
\subsection*{Why do you need to know this material?}
Many substances are liquids under normal conditions and many chemical reactions take place in liquids, so it is important to be able to describe and understand the structure of the liquid phase and its interface with its vapour.

\subsection*{What is the key idea?}
The properties of liquids reflect the short-range order of their molecules in the bulk and the behaviour of their molecules at the mobile surface.

\subsection*{What do you need to know already?}
You need to be aware of the ways in which molecules interact with each other (Topic 14B), and to be familiar with the Helmholtz and Gibbs energies (Topic 3D) and their significance. The Topic makes light use of the Boltzmann distribution (the Prologue and Topic 13A).

Molecules attract each other when they are less than a few diameters apart, but as soon as they come into contact they repel each other. The attraction is responsible for the formation of condensed phases, including liquids, and the repulsion is responsible for the fact that liquids (and solids) have a definite bulk. In a liquid the kinetic energies of the molecules are comparable to their potential energies and, as a result, although the molecules of a liquid are not free to escape completely from the bulk at low temperatures, the structure is very mobile. The cohesive forces responsible for the formation of liquids also result in the interface between the liquid and another phase having an effect on the thermodynamic and physical properties of the liquid.

\subsection*{14C. 1 Molecular interactions in liquids}
The starting point for the discussion of solids is the wellordered structure of a perfect crystal (Topic 15A). The starting point for the discussion of gases is the completely disordered distribution of the molecules of a perfect gas (Topic 1A). Liquids lie between these two extremes. The structural and\\
thermodynamic properties of liquids depend on the nature of intermolecular interactions and, just as for a real gas, a model based on these interactions can be used to develop an equation of state.

\subsection*{(a) The radial distribution function}
The average relative locations of the particles of a liquid are expressed in terms of the radial distribution function (or pair distribution function), $g(r)$. This function is defined so that $4 \pi \mathcal{N} g(r) r^{2} \mathrm{~d} r$ is the number of molecules in a shell of thickness $\mathrm{d} r$ at radius $r$ from a given molecule; $\mathcal{N}=N / V$ is the overall number density. For a uniform material, $g(r)=1$. When defined in this way the radial distribution function is also the ratio of the number of molecules in a shell to the number expected in the same shell for a uniform material. If, at a particular distance, $g(r)>1$ there are more molecules than in a uniform material, whereas if $g(r)<1$, there are fewer.

In a perfect crystal, $g(r)$ is a periodic array of sharp spikes, representing the certainty (in the absence of defects and thermal motion) that molecules (or ions) lie at definite locations. This regularity continues out to the edges of the crystal, so crystals are said to have long-range order. When the crystal melts, the long-range order is lost and the probability of finding a second molecule at long distances from the first is independent of the distance. However, close to the first molecule the nearest neighbours might still adopt approximately their original relative positions and, even if thermal motion drives them away, incoming molecules adopt the vacated positions. It is therefore still possible to detect a shell of nearest neighbours at a distance $r_{1}$, and perhaps beyond them a shell of next-nearest neighbours at $r_{2}$. The existence of this short-range order means that the radial distribution function in a liquid can be expected to oscillate at short distances, with a peak at $r_{1}$, a smaller peak at $r_{2}$, and perhaps some more structure beyond that.

The experimentally determined radial distribution function of the oxygen atoms in liquid water is shown in Fig. 14C.1. Close analysis of a more elaborate form of the distribution function shows that any given $\mathrm{H}_{2} \mathrm{O}$ molecule is surrounded by other molecules at the corners of a tetrahedron. The form of $g(r)$ at $100^{\circ} \mathrm{C}$ shows that the intermolecular interactions (in this case, principally hydrogen bonds) are strong enough to affect the local structure right up to the boiling point. Raman spectra indicate that in liquid water most molecules participate in either three or four hydrogen bonds. Infrared spectra show that about 90 per cent of hydrogen bonds are intact at the\\
\includegraphics[max width=\textwidth, center]{2025_02_27_d4b95d9718f0b9e2e6b9g-635(1)}

Figure 14C. 1 The radial distribution function of the oxygen atoms in liquid water at three temperatures. Note the expansion as the temperature is raised. (Based on A.H. Narten, M.D. Danford, and H.A. Levy, Discuss. Faraday Soc. 43, 97 (1967).)\\
melting point of ice, falling to about 20 per cent at the boiling point.

The formal expression for the radial distribution function for molecules 1 and 2 in a fluid consisting of $N$ particles is

$$
g\left(r_{12}\right)=\frac{1}{(N-2)!\mathcal{N}^{2} Z} \int \mathrm{e}^{-\beta V_{N}} \mathrm{~d} \tau_{3} \mathrm{~d} \tau_{4} \ldots \mathrm{~d} \tau_{N}
$$

Radial distribution function\\
(14C.1a)\\
where $\mathrm{d} \tau_{i}$ is the volume element for molecule $i, \beta=1 / k T, V_{N}$ is the $N$-particle potential energy, and $Z$ is the 'configuration integral' (this quantity is introduced in Topic 13D):

$$
Z=\frac{1}{N!} \int \mathrm{e}^{-\beta V_{N}} \mathrm{~d} \tau_{1} \mathrm{~d} \tau_{2} \ldots \mathrm{~d} \tau_{N} \quad \begin{aligned}
& \text { Configuration integral } \\
& \text { [definition] }
\end{aligned}
$$

(14C.1b)\\
Equation 14C.1a is nothing more than the Boltzmann distribution for the relative locations of two molecules in a field provided by all the molecules in the system. Thus, if there are no interactions between molecules (so $V_{N}=0$ ), $Z=V^{N} / N$ ! and

$$
\begin{aligned}
g\left(r_{12}\right) & =\frac{N!}{(N-2)!\mathcal{N}^{2} V^{N}} \overbrace{\int \mathrm{~d} \tau_{3} \mathrm{~d} \tau_{4} \ldots \mathrm{~d} \tau_{N}}^{V^{N-2}} \\
& =\frac{N(N-1)}{\mathcal{N}^{2} V^{2}}=\frac{N(N-1)}{N^{2}}=1
\end{aligned}
$$

In the absence of interactions the fluid is expected to be uniform, which is consistent with this value of $g\left(r_{12}\right)$.

\subsection*{(b) The calculation of $g(r)$}
The integrals in eqn 14C. 1 are very difficult to evaluate so various numerical procedures are used to calculate the radial distribution function. Such calculations involve specifying the form\\
\includegraphics[max width=\textwidth, center]{2025_02_27_d4b95d9718f0b9e2e6b9g-635}

Figure 14C. 2 In a two-dimensional simulation of a liquid that uses periodic boundary conditions, when one particle leaves the cell (on the left here) its mirror image enters through the opposite face (on the right).\\
of the intermolecular potential energy, for example by specifying it as a pairwise Lennard-Jones interaction (Topic 14B).

Numerical methods typically approach the calculation of the distribution function by considering a box containing about $10^{3}$ particles (the number increases as computers grow more powerful), and then mimicking the rest of the liquid by surrounding the box with replications of the original box (Fig. 14C.2). Whenever a particle leaves the box through one of its faces, its image arrives through the opposite face. When calculating the interactions of a molecule in a box, it interacts with all the molecules in the box and all the periodic replications of those molecules (and of itself) in the other boxes.

In the Monte Carlo method, the particles in the box are first distributed at random and then moved through small random distances. The change in total potential energy of the $N$ particles in the box, $\Delta V_{N}$, is calculated, and if the new arrangement has lower potential energy than the original one, it is accepted. If the potential energy increases, implying that $\Delta V_{N}$ is positive, the new arrangement is accepted only if the Boltzmann-type factor $\mathrm{e}^{-\Delta V_{\mathrm{N}} / k T}$ is greater than a random number (chosen to lie somewhere in the range 0 to 1 ). If this criterion is not met, a new arrangement is generated from the original one and tested in the same way.

The result of applying this selection process is that moving to an arrangement in which the energy is lower is always allowed but moving to one of higher energy, although possible, become increasing unlikely the higher its energy. As a result, the simulation explores a wide range of arrangements, but tends to include fewer with high energies. For each accepted arrangement the number of pairs of molecules with a separation $r$ is counted, and the result is then averaged over the whole collection of accepted arrangements; by repeating this for different values of $r$, the form of $g(r)$ is built up.

In the molecular dynamics approach, the history of an initial arrangement is followed by calculating the trajectories of all the molecules under the influence of the intermolecular\\
\includegraphics[max width=\textwidth, center]{2025_02_27_d4b95d9718f0b9e2e6b9g-636}

Figure 14C. 3 The radial distribution function for a simulation of a liquid using impenetrable hard spheres (such as ball bearings).\\
potentials and the forces they exert. The calculation gives a series of snapshots of the liquid, and $g(r)$ can be calculated as before. The temperature of the system is inferred by computing the mean translational kinetic energy of the molecules and using the equipartition result (The chemist's toolkit 7, in Topic 2A) that $\left\langle\frac{1}{2} m v_{q}^{2}\right\rangle=\frac{1}{2} k T$, for each coordinate $q$.

Such numerical calculations for a fluid of hard spheres without attractive interactions (a collection of ball-bearings in a container) give a radial distribution function that oscillates for small separations of the molecules (Fig. 14C.3). It appears that one of the factors influencing, and sometimes dominating, the structure of a liquid is the geometrical problem of stacking together reasonably hard spheres. Indeed, the radial distribution function of a liquid composed of hard spheres shows more pronounced oscillations at a given temperature than that of any other model liquid. The attractive part of the potential modifies this basic structure, and one of the reasons behind the difficulty of describing actual liquids theoretically is the similar importance of both the attractive and repulsive (hard core) components of the potential energy.

\subsection*{(c) The thermodynamic properties of liquids}
Once $g(r)$ is known, it can be used to calculate the thermodynamic properties of liquids, but for dense fluids the calculations are very complicated. The calculations are simpler for fluids that are so dilute that they are little more than real gases. In such a case, if it is assumed that the interaction of each pair of molecules is given by an isotropic pairwise potential energy function, $V_{2}(r)$, the result is a contribution to the internal energy given by


\begin{align*}
& U_{\text {interaction }}(T)=\frac{2 \pi N^{2}}{V} \int_{0}^{\infty} V_{2}(r) g(r) r^{2} \mathrm{~d} r \\
& \quad \text { Contribution of pairwise interactions to the internal energy } \tag{14C.2a}
\end{align*}


which it is more revealing to write as


\begin{equation*}
U_{\text {interaction }}(T)=\frac{1}{2} N \int_{0}^{\infty} V_{2}(r)\{4 \pi \mathcal{N} g(r)\} r^{2} \mathrm{~d} r \tag{14C.2b}
\end{equation*}


This formulation shows that the internal energy is given by $4 \pi \mathcal{N} g(r) r^{2} \mathrm{~d} r$, which is the number of molecules in a shell of radius $r$ and thickness $\mathrm{d} r$, multiplied by $V_{2}(r)$, and then integrated over $r$, which gives the total energy of interaction of one molecule with all the others. Multiplication by $N$ then gives the total energy of interaction of all the molecules; the factor $\frac{1}{2}$ is needed to avoid counting each interaction twice. Likewise, the equation of state of the dilute fluid including contributions from the pairwise interactions is


\begin{equation*}
\frac{p V}{n R T}=1-\frac{2 \pi \mathcal{N}}{3 k T} \int_{0}^{\infty} v_{2} g(r) r^{2} \mathrm{~d} r \quad v_{2}(r)=r \frac{\mathrm{~d} V_{2}(r)}{\mathrm{d} r} \tag{14C.3a}
\end{equation*}


The quantity $v_{2}(r)$ is called the virial (hence the term 'virial equation of state'). To understand the physical significance of this expression, it can be rewritten as

$$
p=\frac{n R T}{V}-\frac{2 \pi}{3} \mathcal{N}^{2} \int_{0}^{\infty} v_{2}(r) g(r) r^{2} \mathrm{~d} r \quad \begin{aligned}
& \text { Pressure in } \\
& \text { terms of } g(r)
\end{aligned}
$$

(14C.3b)\\
and interpreted as follows:

\begin{itemize}
  \item The first term on the right is the kinetic pressure, the contribution to the pressure from the impact of the molecules in free flight, as in a perfect gas.
  \item The second term, as explained below, is essentially the internal pressure, $\pi_{T}=(\partial U / \partial V)_{T}$ (Topic 2D), representing the contribution of the intermolecular forces to the pressure.
\end{itemize}

To see the connection to the internal pressure, the term $-\mathrm{d} V_{2} / \mathrm{d} r$ (in $V_{2}$ ) should be recognized as the force required to move two molecules apart, and therefore $-r\left(\mathrm{~d} V_{2} / \mathrm{d} r\right)$ is the work required to separate the molecules through a distance $r$. The second term is therefore the average of this work over the range of pairwise separations in the fluid, with the contributions to the average weighted by the probability of finding two molecules at separations between $r$ and $r+\mathrm{d} r$, which is $4 \pi g(r) r^{2} \mathrm{~d} r$. That is, the integral, when multiplied by the square of the number density, is the change in internal energy of the system as it expands, $(\partial U / \partial V)_{T}$, and is therefore equal to the internal pressure. A deeper look 9 on the website of this text explains how this interpretation can be used to infer the virial equation of state of a real gas and interpret the van der Waals parameters.

The pressure given by eqn 14C.3b has nothing to do with the hydrostatic pressure, the pressure experienced at the foot of a column of incompressible liquid. The mass of a column of liquid of mass density $\rho$, height $h$, and cross-sectional area $A$ is $\rho h A$. In a gravitational field the downward force is $\rho h A g_{\text {acc }}$, where $g_{\text {acc }}$ is the acceleration of free fall (it is normally denoted simply $g$, but in this Topic it is necessary to distinguish it from the radial distribution function). The hydrostatic pressure, the force divided by the area $A$ on which it is exerted, is therefore

$$
p=\rho g_{\mathrm{acc}} h
$$

Hydrostatic pressure [incompressible fluid]

The molecular origin of this pressure is as follows. Bear in mind that the fluid is incompressible, so if a molecule moves, space must be made available for it. If a molecule moves up through a molecular diameter from a certain point at a location a distance $h$ down in the liquid, the entire column of height $h$ above it must move up by a molecular diameter. The force required is $\rho g_{\text {acc }} h a$, where $a$ is the cross-sectional area of that molecular column: this force corresponds to a pressure $\rho g_{\text {acc }} h$. If a molecule moves to one side through a molecular diameter, the incompressible column above the new position must move up to make room for it. Once again, the force required corresponds to a pressure $\rho g_{\text {acc }} h$. Even if a molecule moves down, a molecule must move to one side to make room for it, and therefore a whole column must move up to allow that movement, and the force once again corresponds to a pressure $\rho g_{\text {acc }} h$. This interpretation shows why the hydrostatic pressure is isotropic even though gravity operates downwards. It also shows why the interior of a solid column does not have an analogous hydrostatic pressure: the molecules cannot move past each other, so there are no consequent forces involved.

\subsection*{14c. 2 The liquid-vapour interface}
The distinctive feature of the interface between the liquid and its vapour is that it is mobile and molecules there experience attractive forces that no longer pull equally in all directions.

\subsection*{(a) Surface tension}
Liquids tend to adopt shapes that minimize their surface area, because this maximizes the number of molecules that are in the bulk and hence are surrounded by and interact with neighbours. Droplets of liquids therefore tend to be spherical, because a sphere is the shape with the smallest surface-to-volume ratio. However, the presence of other forces distort this ideal shape: the combined effect of gravity and adhesion to a surface results in small droplets on a surface becoming flattened, and gravity results in larger volumes of liquid adopting the shape of the lower part of its container.

Surface effects can be analysed by using the properties of the Helmholtz and Gibbs energies, $A$ and $G$ (Topic 3D). As shown there, under appropriate conditions, $\mathrm{d} A$ and $\mathrm{d} G$ are equal to the work done on the system, including the work done when the surface area changes. The work needed to change the surface area, $\sigma$, of a sample by an infinitesimal amount $\mathrm{d} \sigma$ is proportional to $\mathrm{d} \sigma$, and is written

\[
\begin{array}{ll}
\mathrm{d} w=\gamma \mathrm{d} \sigma & \begin{array}{l}
\text { Surface tension } \\
\text { [definition] }
\end{array} \tag{14C.5}
\end{array}
\]

The constant of proportionality, $\gamma$, is called the surface tension; its dimensions are energy/area and so in SI its units are joules per metre squared $\left(\mathrm{J} \mathrm{m}^{-2}\right)$. However, for reasons that will

Table 14C. 1 Surface tensions of liquids at $293 \mathrm{~K}^{*}$

\begin{center}
\begin{tabular}{lc}
\hline
 & $\gamma /\left(\mathrm{mN} \mathrm{m}^{-1}\right)$ \\
\hline
Benzene & 28.88 \\
Mercury & 472 \\
Methanol & 22.6 \\
Water & 72.75 \\
\hline
\end{tabular}
\end{center}

\begin{itemize}
  \item More values are given in the Resource section. Note that $1 \mathrm{mN} \mathrm{m}^{-1}=1 \mathrm{mJm}^{-2}$.\\
become clear, values of $\gamma$ are usually reported in newtons per metre (because $1 \mathrm{~J}=1 \mathrm{Nm}$, it follows that $1 \mathrm{Jm}^{-2}=1 \mathrm{Nm}^{-1}$ ); Table 14C. 1 gives some typical values. The maximum work of surface formation at constant volume and temperature can be identified with the change in the Helmholtz energy:
\end{itemize}


\begin{equation*}
\mathrm{d} A=\gamma \mathrm{d} \sigma \tag{14C.6}
\end{equation*}


The Helmholtz energy decreases $(\mathrm{d} A<0)$ if the surface area decreases $(\mathrm{d} \sigma<0)$. A process in which $A$ decreases is spontaneous, so it follows that surfaces have a natural tendency to contract, which is the thermodynamic explanation of the observations made at the start of this section.

\subsection*{Example 14C. 1 Using the surface tension}
Consider the arrangement shown in Fig. 14C. 4 in which a wire frame of width $l$ is raised out of a liquid to a height $h$, thereby generating a rectangular film of liquid within the frame. Calculate the work needed to draw a frame of width 5.0 cm out of water at $20^{\circ} \mathrm{C}$ (when $\gamma=72.75 \mathrm{~mJ} \mathrm{~m}^{-2}$ ) through 2.0 cm ; disregard gravitational potential energy.

Collect your thoughts If you assume that the surface tension does not vary with the area, eqn 14C. 5 becomes $w=\gamma \Delta \sigma$ for an increase in surface area $\Delta \sigma$. The increase in surface area is from zero to the area of the rectangle, but you need to recognize that two surfaces are created, one on each side of the frame. Once you have the appropriate expression, insert the data.\\
\includegraphics[max width=\textwidth, center]{2025_02_27_d4b95d9718f0b9e2e6b9g-637}

Figure 14C. 4 The model used for calculating the work of forming a liquid film when a wire frame of width / is raised from a liquid to a height $h$.

The solution The area of the rectangle is $l h$, and hence the increase in surface area of the film is $2 l h$; the work done is therefore $2 \gamma \mathrm{lh}$. A frame of width 5.0 cm pulled out of water at $20^{\circ} \mathrm{C}$ through 2.0 cm requires

$$
w=2 \times\left(72.75 \mathrm{~mJ} \mathrm{~m}^{-2}\right) \times\left(5 \times 10^{-2} \mathrm{~m}\right) \times\left(2 \times 10^{-2} \mathrm{~m}\right)=0.15 \mathrm{~mJ}
$$

Comment. The expression $2 \gamma l h$ can be thought of as $2 \gamma l \times h$, which is force $\times$ distance. The term $2 \gamma l$ can be identified as the opposing force on the top of the frame, which has length $l$. This interpretation is why $\gamma$ is called a tension and why its units are often chosen to be newtons per metre (so $\gamma l$ is a force in newtons).

Self-test 14C. 1 Derive an expression for the work of creating a spherical cavity of radius $r$ in a liquid of surface tension $\gamma$; evaluate this work for a cavity of radius 1.0 cm in water at $20^{\circ} \mathrm{C}$.

\subsection*{(b) Curved surfaces}
The minimization of the surface area of a liquid commonly results in the formation of a curved surface. A bubble is a region in which vapour (and possibly air too) is trapped by a thin film; a cavity is a vapour-filled hole in a liquid. What are widely called 'bubbles' in liquids are therefore strictly cavities. True bubbles have two surfaces (one on each side of the film); cavities have only one. The treatments of both are similar, but a factor of 2 is required for bubbles to take into account the doubled surface area. A droplet is a small volume of liquid surrounded by its vapour (and possibly also air).\\
The tendency for a cavity to minimize its surface area results in the pressure inside the cavity (the concave side of the surface) being greater than that outside (the convex side of the surface). The challenge is to find the relation between these two pressures.

\subsection*{How is that done? 14C. 1}
Relating the pressures inside and outside a cavity

A cavity is at equilibrium when the tendency for its surface area to decrease is balanced by the rise in internal pressure that would result. Equilibrium is achieved when the inward and outward forces on the surface are equal.

Step 1 Evaluate the force due to the external pressure\\
The force on the surface of a spherical cavity of radius $r$ due to the external pressure $p_{\text {out }}$ is given by area $\times$ pressure $=4 \pi r^{2} p_{\text {out }}$.

\subsection*{Step 2 Evaluate the force due to surface tension}
The change in surface area when the radius changes from $r$ to $r+\mathrm{d} r$ is\\
$\mathrm{d} \sigma=4 \pi(r+\mathrm{d} r)^{2}-4 \pi r^{2}=8 \pi r \mathrm{~d} r$\\
(The second-order infinitesimal, ( $\mathrm{d} r)^{2}$, has been ignored.) The work done when the surface is stretched by this amount is therefore

$$
\mathrm{d} w=\gamma \mathrm{d} \sigma=8 \pi \gamma r \mathrm{~d} r
$$

As force $\times$ distance is work, the force opposing an increase in the radius by $\mathrm{d} r$ is

$$
F=8 \pi \gamma r
$$

The total inward force is therefore $4 \pi r^{2} p_{\text {out }}+8 \pi \gamma r$.\\
Step 3 Balance the inward and outward forces\\
If the pressure inside the cavity is $p_{\mathrm{in}}$, the outward force on the surface is $4 \pi r^{2} p_{\mathrm{in}}$. At equilibrium, the outward and inward forces are balanced:

$$
4 \pi r^{2} p_{\text {in }}=4 \pi r^{2} p_{\text {out }}+8 \pi \gamma r
$$

Division of both sides by $4 \pi r^{2}$ results in the Laplace equation:

$$
-\left|p_{\text {in }}=p_{\text {out }}+\frac{2 \gamma}{r}\right| \quad \text { (14C.7) }
$$

The Laplace equation shows that the difference in pressure decreases to zero as the radius of curvature becomes infinite (when the surface is flat, Fig. 14C.5). Small cavities have small radii of curvature, so the pressure difference across their surface is quite large.

\subsection*{Brief illustration 14C. 1}
The pressure difference across the surface of a spherical droplet of water of radius 200 nm at $20^{\circ} \mathrm{C}$ can be calculated using the Laplace equation:

$$
\begin{aligned}
p_{\text {in }}-p_{\text {out }} & =\frac{2 \times \overbrace{\left(72.75 \times 10^{-3} \mathrm{Nm}^{-1}\right)}^{\gamma_{\text {water }} \text { at } 20^{\circ} \mathrm{C}}}{\underbrace{2.00 \times 10^{-7} \mathrm{~m}}_{r}} \\
& =7.28 \times 10^{5} \mathrm{Nm}^{-2}=728 \mathrm{kPa}
\end{aligned}
$$

\subsection*{(c) Capillary action}
When a narrow-bore tube (a 'capillary tube'; the name comes from the Latin word for 'hair') is dipped in a liquid there is a tendency for the liquid to rise up the tube. This tendency is called capillary action and can be understood as a consequence of surface tension.

Water has a tendency to adhere to the surface of glass, so when a glass capillary tube is first immersed in water a film of\\
\includegraphics[max width=\textwidth, center]{2025_02_27_d4b95d9718f0b9e2e6b9g-639}

Figure 14C. 5 The dependence of the pressure inside a curved surface on the radius of the surface, for increasing values of the surface tension.\\
solvent spreads along the surface: the further the film spreads, the lower is the energy due to the interaction. As the film spreads up the inside walls of the capillary tube the surface of the liquid becomes curved and, as has just been discussed, this curvature results in a pressure difference across the surface. The pressure just above the meniscus, the concave side, is greater than that just below, the convex side.

If it is assumed that the surface is hemispherical and that the tube has radius $r$, the pressure difference is given by the Laplace equation as $2 \gamma / r$. The pressure just above the surface is the atmospheric pressure, $p$, so the pressure just below the surface is $p-2 \gamma / r$. The atmospheric pressure pushing on the liquid outside the tube causes the liquid inside the tube to rise until hydrostatic equilibrium is achieved, which is when there are equal pressures at equal depths (Fig. 14C.6). When the liquid has risen to a height $h$ the column exerts a hydrostatic pressure\\
\includegraphics[max width=\textwidth, center]{2025_02_27_d4b95d9718f0b9e2e6b9g-639(1)}

Figure 14C. 6 When a capillary tube is first stood in a liquid, the liquid climbs up the walls, so curving the surface. The pressure just under the meniscus is less than that arising from the atmosphere by $2 \gamma / r$. The pressure is equal at equal heights throughout the liquid provided the hydrostatic pressure (which is equal to $\rho g_{\text {acc }} h$ ) cancels the pressure difference arising from the curvature.\\
\includegraphics[max width=\textwidth, center]{2025_02_27_d4b95d9718f0b9e2e6b9g-639(2)}

Figure 14C. 7 The variation of the surface tension of water with temperature.\\
given by eqn $14 \mathrm{C} .4, \rho g_{\text {acc }} h$, where $\rho$ is the mass density of the liquid and $g_{\text {acc }}$ is the acceleration of free fall. The total pressure at the base of the column is therefore $(p-2 \gamma / r)+\rho g_{\text {acc }} h$. At the same level outside the tube the pressure is $p$, but at equilibrium these two pressures must be the same: $p=p-2 \gamma / r+\rho g_{\text {acc }} h$. It therefore follows that $2 \gamma / r=\rho g_{\text {acc }} h$ which rearranges to


\begin{equation*}
h=\frac{2 \gamma}{\rho g_{\mathrm{acc}} r} \tag{14C.8}
\end{equation*}


This simple expression provides a reasonably accurate way of measuring the surface tension of liquids. Surface tension decreases with increasing temperature (Fig. 14C.7), which can be understood as arising from the increase in thermal motion moving molecules more rapidly between the surface and the bulk.

\subsection*{Brief illustration 14C. 2}
If water at $25^{\circ} \mathrm{C}$ (with mass density $997.1 \mathrm{~kg} \mathrm{~m}^{-3}$ ) rises through 7.36 cm in a capillary of radius 0.20 mm , its surface tension at that temperature is

$$
\begin{aligned}
& \gamma=\frac{1}{2} \rho g_{\text {acc }} h r \\
&=\frac{1}{2} \times\left(997.1 \mathrm{~kg} \mathrm{~m}^{-3}\right) \times\left(9.81 \mathrm{~m} \mathrm{~s}^{-2}\right) \times\left(7.36 \times 10^{-2} \mathrm{~m}\right) \times\left(2.0 \times 10^{-4} \mathrm{~m}\right) \\
& \\
& \mathrm{kg} \mathrm{~ms}^{-2}=\mathrm{N} \\
&=72 \mathrm{mNm}^{-1}
\end{aligned}
$$

When the adhesive forces between the liquid and the material of the capillary wall are weaker than the cohesive forces within the liquid (as for mercury in glass), it is energetically favourable for the liquid in the tube to retract from the walls. This retraction curves the surface with the concave, high pressure side downwards. To equalize the pressure at the same\\
\includegraphics[max width=\textwidth, center]{2025_02_27_d4b95d9718f0b9e2e6b9g-640}

Figure 14C. 8 The balance of forces that results in a contact angle, $\theta_{c}$.\\
depth throughout the liquid the surface must fall to compensate for the increased pressure arising from its curvature. This compensation results in a capillary depression.

The angle between the surface and the wall, where the two meet, is called the contact angle, $\theta_{c}$ (Fig. 14C.8). In many cases this angle is found to be non-zero, and eqn 14C. 8 must then be modified by multiplying the right-hand side by $\cos \theta_{\mathrm{c}}$. The origin of the contact angle can be traced to the balance of forces at the line of contact between the liquid and the solid (Fig. 14C.8).

If the solid-gas, solid-liquid, and liquid-gas surface tensions are denoted $\gamma_{\mathrm{sg}}, \gamma_{\mathrm{sl}}$, and $\gamma_{\mathrm{lg}}$, respectively, then the vertical forces are in balance if


\begin{equation*}
\gamma_{\mathrm{sg}}=\gamma_{\mathrm{sl}}+\gamma_{\mathrm{lg}} \cos \theta_{\mathrm{c}} \tag{14C.9a}
\end{equation*}


and therefore


\begin{equation*}
\cos \theta_{\mathrm{c}}=\frac{\gamma_{\mathrm{sg}}-\gamma_{\mathrm{sl}}}{\gamma_{\mathrm{lg}}} \tag{14C.9b}
\end{equation*}


The 'superficial work' of adhesion of the liquid to the solid, $w_{\mathrm{ad}}$, is the work of adhesion divided by the area of contact. As the liquid-solid interface expands it does so at the expense of the solid-gas and liquid-gas interfaces, and the net work involved is


\begin{equation*}
w_{\mathrm{ad}}=\gamma_{\mathrm{sg}}+\gamma_{\mathrm{lg}}-\gamma_{\mathrm{sl}} \tag{14C.10a}
\end{equation*}


and therefore, from the previous equation,

$$
\cos \theta_{\mathrm{c}}=\frac{w_{\mathrm{ad}}}{\gamma_{\mathrm{lg}}}-1
$$

Contact angle (14C.10b)

It is now seen that:

\begin{itemize}
  \item When the contact angle is between 0 and $90^{\circ}$ (implying that $0<\cos \theta_{c}<1$ ) the liquid 'wets' the surface, meaning that it spreads over the surface. From eqn 14C.10b wetting occurs when $1<w_{\text {ad }} / \gamma_{\text {gg }}<2$ (Fig. 14C.9).
  \item When the contact angle is between $90^{\circ}$ and $180^{\circ}$ (implying that $-1<\cos \theta_{c}<0$ ) the liquid does not wet the surface; this condition corresponds to $0<w_{\mathrm{ad}} / \gamma_{\mathrm{lg}}<1$.\\
\includegraphics[max width=\textwidth, center]{2025_02_27_d4b95d9718f0b9e2e6b9g-640(1)}
\end{itemize}

Figure 14C. 9 The variation of contact angle as the ratio $w_{\text {ad }} / \gamma_{l g}$ changes.

For mercury in contact with glass, $\theta_{c}=140^{\circ}$, which corresponds to $w_{\text {ad }} / \gamma_{\mathrm{lg}}=0.23$, indicating a relatively low work of adhesion of the mercury to glass on account of the strong cohesive forces within the liquid.

\subsection*{14C. 3 Surface films}
The compositions of surface layers have been investigated by the simple but technically elegant procedure of slicing thin layers off the surfaces of solutions and analysing their compositions. The physical properties of surface films have also been investigated. Surface films one molecule thick are called monolayers, and when such a monolayer has been transferred to a solid support, it is called a Langmuir-Blodgett film, after Irving Langmuir and Katherine Blodgett, who developed experimental techniques for studying them.

\subsection*{(a) Surface pressure}
The principal apparatus used for the study of surface monolayers is a surface-film balance (or Langmuir-Blodgett trough, Fig. 14C.10). This device consists of a shallow trough and a barrier that can be moved along the surface of the liquid in the trough, thus compressing any monolayer on the surface. The surface pressure, $\pi$, the difference between the surface tension of the pure solvent and the solution $\left(\pi=\gamma^{\star}-\gamma\right)$, is measured by using a torsion wire attached to a strip of mica that rests on the surface and pressing against one edge of the monolayer. The parts of the apparatus that are in touch with liquids are coated in polytetrafluoroethene to eliminate effects arising from the liquid-solid interface. In an actual experiment, a small amount (about 0.01 mg ) of the substance under investigation is dissolved in a volatile solvent and then poured on to the surface of the water; the compression barrier is then moved across the surface and the surface pressure exerted on the mica bar is monitored.\\
\includegraphics[max width=\textwidth, center]{2025_02_27_d4b95d9718f0b9e2e6b9g-641(1)}

Figure 14C.10 A schematic diagram of the apparatus used to measure the surface pressure and other characteristics of a surface film. The surfactant is spread on the surface of the liquid in the trough, and then compressed horizontally by moving the compression barrier towards the mica float. The latter is connected to a torsion wire, so the net force on the float can be monitored.

When the surface coverage is low it is found that the surface pressure is inversely proportional to the total area of the surface. This behaviour is the analogue in two dimensions of a perfect gas (where $p \propto 1 / V$ ), and can be interpreted as arising when the average separation between the molecules at the surface is so large that the interactions between them are not important. As the area is further decreased, the surface pressure eventually starts to increase more rapidly, as is illustrated in Fig. 14C.11. This behaviour can be thought of as arising from the formation of a monolayer in which the molecules are in relatively close contact; like a liquid, such a layer is almost incompressible. The area corresponding to a complete closepacked monolayer is found by extrapolating the steepest part of the isotherm.

As can be seen from Fig. 14C.11, even though stearic acid (1) and isostearic acid (2) are chemically very similar (they differ only in the location of a methyl group at the end of a long hy-\\
\includegraphics[max width=\textwidth, center]{2025_02_27_d4b95d9718f0b9e2e6b9g-641}

Figure 14C. 11 The variation of surface pressure with the average area occupied by a surfactant molecule. The collapse pressures are indicated by the horizontal arrows.\\
drocarbon chain), they occupy significantly different areas in the monolayer. Neither, though, occupies as much area as the tri- $p$-cresyl phosphate molecule (3), which is like a wide bush rather than a lanky tree.\\
\includegraphics{smile-8e3e1102e5d056cc3ce239cce0a2616a8c8aeeac}\\
\includegraphics{smile-e6509c1f28b1aaec22cda190256863e4226758d9}\\
\includegraphics{smile-b940278be855eeb56c30972f2677ad03dc4554af}

3 Tri-p-cresyl phosphate

The second feature to note from Fig. 14C. 11 is that the tri- $p$ cresyl phosphate isotherm is much less steep than the stearic acid isotherms. This difference indicates that the tri- $p$-cresyl phosphate film is more compressible than the stearic acid films, which is consistent with their different molecular structures.

A third feature of the isotherms is the collapse pressure, the highest surface pressure at which a monolayer can be maintained. When the monolayer is compressed beyond the collapse pressure, the monolayer buckles and collapses into a film several molecules thick. As can be seen from the isotherms in Fig. 14C.11, stearic acid has a high collapse pressure, but that of tri- $p$-cresyl phosphate is significantly lower, indicating a much weaker film.

\subsection*{(b) The thermodynamics of surface layers}
A surfactant is a species that accumulates at the interface between two phases, such as that between hydrophilic and hydrophobic phases, and modifies the surface tension. The relation between the concentration of surfactant at the surface and the change in surface tension it brings about can be established by considering a model in which two phases $\alpha$ and $\beta$ come into contact, therefore creating an interface. Within each bulk phase the composition is constant, but the\\
composition in the interfacial region may be different due to the accumulation of the surfactant.

The total Gibbs energy can be thought of as having a contribution from the two phases, $G(\alpha)$ and $G(\beta)$, together with a contribution $G(\sigma)$, the surface Gibbs energy, from the interfacial region

$$
G=G(\alpha)+G(\beta)+G(\sigma) \quad \begin{aligned}
& \text { Surface Gibbs energy } \\
& \text { [definition] }
\end{aligned}
$$

(14C.11)

In a similar way, the total amount of substance J, $n_{\mathrm{J}}$, can be thought of as being divided between the amounts in phases $\alpha$ and $\beta, n_{J}(\alpha)$ and $n_{J}(\beta)$, and the amount at the interfacial region, $n_{\mathrm{J}}(\sigma): n_{\mathrm{J}}=n_{\mathrm{J}}(\alpha)+n_{\mathrm{J}}(\beta)+n_{\mathrm{J}}(\sigma)$. The amount at the interface can be expressed in terms of the surface excess, $\Gamma_{\mathrm{J}}$ :

\[
\Gamma_{\mathrm{J}}=\frac{n_{\mathrm{J}}(\sigma)}{\sigma} \quad \begin{align*}
& \text { Surface excess }  \tag{14C.12}\\
& \text { [definition] }
\end{align*}
\]

where $\sigma$ is the area of the surface. It is possible to relate the surface tension to the surface excess and therefore to the concentration of surfactant at the interface.

\subsection*{How is that done? 14C. 2}
Relating the surface tension to\\
the surfactant concentration\\
The change in $G$ brought about by changes in $T, p$, and $n_{\mathrm{J}}$ is given by eqn 5A.6:

$$
\mathrm{d} G=-S \mathrm{~d} T+V \mathrm{~d} p+\sum_{\mathrm{J}} \mu_{\mathrm{J}} \mathrm{~d} n_{\mathrm{J}}
$$

where $\mu_{\mathrm{J}}$ is the chemical potential of substance J . At constant pressure, $V \mathrm{~d} p=0$ and at constant temperature $S \mathrm{~d} T=0$, so the first two terms on the right vanish.

Step 1 Write an expression for the change in Gibbs energy of the interface\\
To apply what remains of this relation to the interface, it is necessary to introduce an additional term $\gamma \mathrm{d} \sigma$ (eqn 14C.5) arising from the work done expanding the interface. The expression for the change in Gibbs energy of the interface, $\mathrm{d} G(\sigma)$, becomes

$$
\mathrm{d} G(\sigma)=\gamma \mathrm{d} \sigma+\sum_{\mathrm{J}} \mu_{\mathrm{J}} \mathrm{~d} n_{\mathrm{J}}(\sigma)
$$

At equilibrium the chemical potential of each component is the same in each phase and is written $\mu_{\text {J }}$.

\subsection*{Step 2 Integrate the infinitesimal change}
Following the same argument as in the discussion of partial molar quantities (Topic 5A), this equation can be integrated at constant temperature, surface tension, and composition to give

$$
G(\sigma)=\gamma \sigma+\sum_{\mathrm{J}} \mu_{\mathrm{J}} n_{\mathrm{J}}(\sigma)
$$

Step 3 Identify the total change in Gibbs energy implied by this expression\\
An infinitesimal change in each of the quantities on the right of this expression gives the following total change in $G(\sigma)$ :

$$
\mathrm{d} G(\sigma)=\gamma \mathrm{d} \sigma+\sigma \mathrm{d} \gamma+\sum_{\mathrm{J}} \mu_{\mathrm{J}} \mathrm{~d} n_{\mathrm{J}}(\sigma)+\sum_{\mathrm{J}} \eta_{\mathrm{J}}(\sigma) \mathrm{d} \mu_{\mathrm{J}}
$$

\subsection*{Step 4 Use the fact that $G$ is a state function}
Because Gibbs energy is a state function, the two expressions for $\mathrm{d} G(\sigma)$ must be the same. Their comparison implies that

$$
\sigma \mathrm{d} \gamma+\sum_{\mathrm{J}} \eta_{\mathrm{J}}(\sigma) \mathrm{d} \mu_{\mathrm{J}}=0
$$

Division by $\sigma$ and introduction of the definition of the surface excess $\left(\Gamma_{\mathrm{J}}=n_{\mathrm{J}}(\sigma) / \sigma\right)$ gives the Gibbs isotherm, which relates the change in surface tension to the changes in the chemical potentials of the substances present in the interface

$$
\mathrm{d} \gamma=-\sum_{\mathrm{J}} \Gamma_{\mathrm{J}} \mathrm{~d} \mu_{\mathrm{J}} \quad \text { Gibbs isotherm }
$$

Step 5 Relate the change in chemical potentials to the composition If just one species, the surfactant S , accumulates at the surface, the Gibbs isotherm becomes

$$
\mathrm{d} \gamma=-\Gamma_{\mathrm{s}} \mathrm{~d} \mu_{\mathrm{s}}
$$

The chemical potential for species $J$ in a dilute solution can be written as $\mu_{\mathrm{J}}=\mu_{\mathrm{J}}^{\ominus}+R T \ln \left(c_{\mathrm{J}} / c^{\ominus}\right)$, where $c_{\mathrm{J}}$ is the molar concentration and $c^{\ominus}$ is its standard value. It follows that at constant temperature $\mathrm{d} \mu_{\mathrm{S}}=R T \mathrm{~d} \ln \left(c / c^{\ominus}\right)$. This expression for $\mathrm{d} \mu_{\mathrm{S}}$ is used in eqn 14C. 13 to give

$$
\mathrm{d} \gamma=-R T \Gamma_{\mathrm{S}} \mathrm{~d} \ln \left(c / c^{\ominus}\right)
$$

which can be rearranged to\\
$-\left(\frac{\partial \gamma}{\partial \ln \left(c / c^{\ominus}\right)}\right)_{T}=-R T \Gamma_{S} \left\lvert\, \begin{aligned} & \text { (14C.14) } \\ & \\ & \begin{array}{l}\text { Dependence of the surface tension } \\ \text { on surfactant concentration }\end{array}\end{aligned}\right.$

When the surfactant accumulates at the interface, its surface excess is positive and eqn 14C. 14 implies that $\left(\partial \gamma / \partial \ln \left(c / c^{\ominus}\right)\right)_{T}<0$. That is, accumulation of the surfactant causes the surface tension to decrease. If the variation of $\gamma$ with concentration is measured, eqn 14C. 14 can be used to determine the surface excess, and this value can be used to infer the area occupied by each surfactant molecule on the surface, as illustrated in the following example.

\subsection*{Example 14C. 2 Determining the surface excess and the}
 surface concentration of surfactant moleculesMeasurements of the surface tension of an aqueous solution of 1-aminobutanoic acid as a function of concentration give $\mathrm{d} \gamma / \mathrm{d} \ln \left(c / c^{\ominus}\right)=-40 \mu \mathrm{Nm}^{-1}$ at $20^{\circ} \mathrm{C}$. Calculate the surface excess of 1 -aminobutanoic acid and the number of molecules per square metre.

Collect your thoughts Equation 14C. 14 relates the measured value of $\mathrm{d} \gamma / \mathrm{d} \ln \left(c / c^{\ominus}\right)$ directly to the surface excess. Multiplication of the surface excess by Avogadro's constant gives the number of molecules per square metre.

The solution From eqn 14C. 14 it follows that

$$
\begin{aligned}
\Gamma_{\mathrm{S}} & =-\frac{1}{R T}\left(\frac{\partial \gamma}{\partial \ln \left(c / c^{\ominus}\right)}\right)_{T} \\
& =-\frac{1}{\left(8.3145 \mathrm{JK}^{-1} \mathrm{~mol}^{-1}\right) \times(293 \mathrm{~K})} \times\left(-4.0 \times 10^{-5} \mathrm{Nm}^{-1}\right) \\
& =1.6 \times 10^{-8} \mathrm{molm}^{-2}
\end{aligned}
$$

The number of molecules per square metre is $N_{\mathrm{A}} \Gamma_{\mathrm{S}}$ :\\
$N_{\mathrm{A}} \Gamma_{\mathrm{S}}=\left(6.022 \times 10^{23} \mathrm{~mol}^{-1}\right) \times\left(1.6 \times 10^{-8} \mathrm{molm}^{-2}\right)=9.6 \times 10^{15} \mathrm{~m}^{-2}$\\
Self-test 14C. 2 Use the result obtained to calculate the area occupied by each molecule of 1 -aminobutanoic acid at the surface.\\
\includegraphics[max width=\textwidth, center]{2025_02_27_d4b95d9718f0b9e2e6b9g-643}

\subsection*{14C. 4 Condensation}
The concepts from this Topic together with some from Topic 4B can be used to explain aspects of the condensation of a gas to a liquid. In Topic 4B it is shown that the vapour pressure of a liquid, $p$, is increased when additional pressure $\Delta P$ is applied to the liquid: according to eqn $4 \mathrm{~B} .2, p=p^{*} \mathrm{e}^{V_{\mathrm{m}}(1) \Delta P / R T}$, where $p^{*}$ is the vapour pressure when no additional pressure is applied and $V_{\mathrm{m}}(\mathrm{l})$ is the molar volume of the liquid. Because of its curved surface, a droplet experiences an additional pressure, given by the Laplace equation (eqn 14C.7) as $2 \gamma / r$, where $r$ is the radius of the surface. When this value is used for $\Delta P$ in eqn 4B. 2 the result is the Kelvin equation for the vapour pressure of a liquid when it is dispersed as spherical droplets:

$$
p=p^{\star} \mathrm{e}^{2 \gamma V_{\mathrm{m}}(1) / r R T}
$$

Kelvin equation\\
(14C.15)\\
For a cavity, the pressure of the liquid outside is less than the pressure inside, so the sign of the exponent in eqn 14C. 15 is changed to obtain an expression for the vapour pressure in a cavity. For droplets of water of radius $1 \mu \mathrm{~m}$ and 1 nm the ratios $p / p^{*}$ at $25^{\circ} \mathrm{C}$ are about 1.001 and 3 , respectively. The second figure, although quite large, is unreliable because at that radius\\
the droplet is less than about 10 molecules in diameter and the basis of the calculation is suspect. The first figure shows that the effect is usually small; nevertheless it may have important consequences.

For instance, consider the formation of a cloud. Warm, moist air rises into the cooler regions higher in the atmosphere. At some altitude the temperature drops to the point that the vapour becomes thermodynamically unstable with respect to the liquid and it is then expected that the vapour will condense into a cloud of liquid droplets. The initial step can be imagined as a swarm of water molecules congregating into a microscopic droplet. Because the initial droplet is so small it has an enhanced vapour pressure; therefore, instead of growing it evaporates. This effect stabilizes the vapour because an initial tendency to condense is overcome by a heightened tendency to evaporate. The vapour phase is then said to be supersaturated. It is thermodynamically unstable with respect to the liquid but not unstable with respect to the small droplets that need to form before the bulk liquid phase can appear, so the formation of the latter by a simple, direct mechanism is hindered.

Two processes are responsible for overcoming this tendency of droplets to evaporate and thus for allowing clouds to form. The first is that a sufficiently large number of molecules might congregate into a droplet so big that the enhanced evaporative effect is unimportant. The chance of one of these spontaneous nucleation centres forming is low, and in rain formation it is not a dominant mechanism. The more important process depends on the presence of minute dust particles or other kinds of foreign matter. These nucleate the condensation (that is, provide centres at which it can occur) by providing surfaces to which the water molecules can attach.

Liquids may be superheated above their boiling temperatures and supercooled below their freezing temperatures. In each case the thermodynamically stable phase is not achieved on account of the kinetic stabilization that occurs in the absence of nucleation centres. For example, superheating occurs because the vapour pressure inside a cavity is artificially low, so any cavity that does form tends to collapse. This instability is encountered when an unstirred beaker of water is heated, for its temperature may be raised above its boiling point. Violent bumping often ensues as spontaneous nucleation leads to bubbles big enough to survive. To ensure smooth boiling at the true boiling temperature, nucleation centres, such as small pieces of sharp-edged glass or bubbles (cavities) of air, should be introduced.

\subsection*{Checklist of concepts}
1. The radial distribution function, $g(r)$, is defined such that $4 \pi \mathcal{N} g(r) r^{2} \mathrm{~d} r$ is the number of molecules in a shell of thickness $\mathrm{d} r$ at radius $r$ from a given molecule; $\mathcal{N}$ is the overall number density.2. The radial distribution function may be calculated numerically by using Monte Carlo and molecular dynamics techniques.3. Liquids tend to adopt shapes that minimize their surface area.4. Capillary action is the tendency of liquids to rise up (and in some cases, drop down) narrow tubes.5. The surface pressure is the difference between the surface tension of a pure solvent and a solution.6. The collapse pressure is the highest surface pressure that a surface film can sustain.7. A surfactant is a species that accumulates at the interface between phases and modifies the surface tension and surface pressure.8. Nucleation provides surfaces to which molecules can attach and thereby induce condensation.\subsection*{Checklist of equations}
\begin{center}
\begin{tabular}{|c|c|c|c|}
\hline
Property & Equation & Comment & Equation number \\
\hline
Hydrostatic pressure & $p=\rho g_{\text {acc }} h$ & Incompressible fluid & 14C. 4 \\
\hline
Laplace equation & $p_{\text {in }}=p_{\text {out }}+2 \gamma / r$ & $\gamma$ is the surface tension & 14C. 7 \\
\hline
Contact angle & $\cos \theta_{\mathrm{c}}=w_{\mathrm{ad}} / \gamma_{\mathrm{lg}}-1$ &  & 14C.10b \\
\hline
Surface Gibbs energy & $G=G(\alpha)+G(\beta)+G(\sigma)$ & Definition & 14C. 11 \\
\hline
Surface excess & $\Gamma_{\mathrm{J}}=n_{\mathrm{J}}(\sigma) / \sigma$ & Definition & 14C. 12 \\
\hline
Gibbs isotherm & \( \mathrm{d} \gamma=-\sum_{\mathrm{J}} \Gamma_{\mathrm{J}} \mathrm{~d} \mu_{\mathrm{J}} \) &  & 14C. 13 \\
\hline
Dependence of the surface tension on surfactant concentration & $\left(\partial \gamma / \partial \ln \left(c / c^{\ominus}\right)\right)_{T}=-R T \Gamma_{\mathrm{S}}$ &  & 14C. 14 \\
\hline
Kelvin equation & $p=p^{*} \mathrm{e}^{2 \gamma \mathrm{~V}_{\mathrm{m}}(1) / r R T}$ &  & 14C. 15 \\
\hline
\end{tabular}
\end{center}

\section*{TOPIC 14D Macromolecules}
\subsection*{Why do you need to know this material?}
Macromolecules give rise to special problems that include the investigation and description of their molar masses and shapes. You need to know how to describe the structural features of macromolecules in order to understand their physical and chemical properties.

\subsection*{What is the key idea?}
The structure of a macromolecule takes on different meanings at the different levels at which the arrangement of the chain or network of its building blocks is considered.

\subsection*{- What do you need to know already?}
The discussion of the shapes of macromolecules depends on an understanding of the nonbonding interactions between molecules (Topic 14B). You also need to be familiar with the statistical interpretation of entropy (Topic 13E) and the concept of internal energy (Topic 2A). Some of the calculations draw on statistical arguments like those used in the discussion of the Boltzmann distribution (Topic 13A).

Macromolecules are very large molecules assembled from smaller molecules biosynthetically in organisms, by chemists in the laboratory, or in an industrial reactor. Naturally occurring macromolecules include polysaccharides such as cellulose, polypeptides such as protein enzymes, and polynucleotides such as deoxyribonucleic acid (DNA). This Topic deals principally with synthetic macromolecules. They include polymers, such as nylon and polystyrene, that are manufactured by stringing together, and in some cases crosslinking, smaller units known as monomers (Fig. 14D.1).\\
\includegraphics[max width=\textwidth, center]{2025_02_27_d4b95d9718f0b9e2e6b9g-645}

Figure 14D. 1 Three varieties of polymer: (a) a simple linear polymer, (b) a cross-linked polymer, and (c) one variety of copolymer.

\subsection*{14D. 1 Average molar masses}
A monodisperse polymer has a single, definite molar mass. A synthetic polymer, however, is polydisperse, in the sense that a sample is a mixture of molecules with various chain lengths and molar masses. The various techniques that are used to measure molar masses result in different types of mean values of polydisperse systems.

The number-average molar mass, $\bar{M}_{\mathrm{n}}$, is obtained by weighting each molar mass by the number of molecules of that mass present in the sample:

\[
\bar{M}_{\mathrm{n}}=\frac{1}{N_{\text {total }}} \sum_{i} N_{i} M_{i} \quad \begin{align*}
& \text { Number-average molar mass }  \tag{14D.1a}\\
& \text { [definition] }
\end{align*}
\]

where $N_{i}$ is the number of molecules with molar mass $M_{i}$ and $N_{\text {total }}$ is the total number of molecules. This type of average is typically obtained by mass spectroscopic determinations of molar mass. The weight-average molar mass is the average calculated by weighting the molar masses of the molecules by the mass present in the sample:

\[
\bar{M}_{\mathrm{w}}=\frac{1}{m_{\text {total }}} \sum_{i} m_{i} M_{i} \quad \begin{align*}
& \text { Weight-average molar mass }  \tag{14D.1b}\\
& \text { [definition] }
\end{align*}
\]

In this expression, $m_{i}$ is the total mass of molecules of molar mass $M_{i}$ and $m_{\text {total }}$ is the total mass of the sample. This type of average is typically obtained by measurements that make use of the ability of molecules to scatter light and by measurements that make use of the distribution of particles in solutions rotated at high speed in an ultracentrifuge.

\subsection*{Example 14D. 1}
Calculating number and mass averages\\
Evaluate the number-average and the weight-average molar masses of a sample of poly(vinyl chloride) from the following data:

\begin{center}
\begin{tabular}{llrrrrr}
\hline
Interval & 1 & \multicolumn{1}{c}{2} & \multicolumn{1}{c}{3} & \multicolumn{1}{c}{4} & \multicolumn{1}{c}{5} & \multicolumn{1}{c}{6} \\
$\boldsymbol{M}_{\boldsymbol{i}} /\left(\mathrm{kg} \mathrm{mol}^{-1}\right)$ & 7.5 & 12.5 & 17.5 & 22.5 & 27.5 & 32.5 \\
$\boldsymbol{m}_{\boldsymbol{i}} / \mathbf{g}$ & 9.6 & 8.7 & 8.9 & 5.6 & 3.1 & 1.7 \\
\hline
\end{tabular}
\end{center}

Collect your thoughts The relevant equations are eqns 14D.1a and 14D.1b. Note that because $N_{i}=n_{i} N_{A}$, you can express the number average in terms of amounts (in moles):

$$
\bar{M}_{\mathrm{n}}=\frac{1}{n_{\text {total }} N_{\mathrm{A}}} \sum_{i} n_{i} N_{\mathrm{A}} M_{i}=\frac{1}{n_{\text {total }}} \sum_{i} n_{i} M_{i}
$$

where $n_{i}$ is the amount (in moles) of molecules with molar mass $M_{i}$ and $n_{\text {total }}$ is the total amount of molecules. Calculate the amounts in each interval by dividing the mass of the sample in each interval by the average molar mass for that interval; $n_{i}=m_{i} / M_{i}$. Then calculate the two averages by weighting the molar mass $M_{i}$ within each interval by the amount $\left(n_{i}\right)$ and mass ( $m_{i}$ ), respectively, of the molecules in each interval.

The solution The amounts in each interval are as follows:

\begin{center}
\begin{tabular}{llccccc}
\hline
Interval & 1 & 2 & 3 & 4 & 5 & 6 \\
$\boldsymbol{M}_{\boldsymbol{i}} /\left(\mathrm{kg} \mathrm{mol}^{-1}\right)$ & 7.5 & 12.5 & 17.5 & 22.5 & 27.5 & 32.5 \\
$\boldsymbol{n}_{\boldsymbol{i}} / \mathrm{mmol}$ & 1.3 & 0.70 & 0.51 & 0.25 & 0.11 & 0.052 \\
\hline
\end{tabular}
\end{center}

The total amount is $n_{\text {total }}=2.92 \mathrm{mmol}$ and the number-average molar mass is

$$
\begin{aligned}
\bar{M}_{\mathrm{n}} /\left(\mathrm{kg} \mathrm{~mol}^{-1}\right)= & \frac{1}{2.92}(1.3 \times 7.5+0.70 \times 12.5+0.51 \times 17.5 \\
& +0.25 \times 22.5+0.11 \times 27.5+0.052 \times 32.5)=13
\end{aligned}
$$

The weight-average molar mass is calculated directly from the data after noting that the total mass of the sample is 37.6 g :

$$
\begin{aligned}
\bar{M}_{\mathrm{w}} /\left(\mathrm{kg} \mathrm{~mol}^{-1}\right)= & \frac{1}{37.6}(9.6 \times 7.5+8.7 \times 12.5+8.9 \times 17.5 \\
& +5.6 \times 22.5+3.1 \times 27.5+1.7 \times 32.5)=16
\end{aligned}
$$

Comment. Note the different values of the two averages. In this instance, $\bar{M}_{\mathrm{w}} / \bar{M}_{\mathrm{n}}=1.2$.

Self-test 14D. 1 The Z-average molar mass, which is obtained in certain sedimentation experiments, is defined as $\bar{M}_{\mathrm{Z}}=\sum_{i} N_{i} M_{i}^{3} / \sum_{i} N_{i} M_{i}^{2}$. Evaluate its value for the sample in this example.\\
\includegraphics[max width=\textwidth, center]{2025_02_27_d4b95d9718f0b9e2e6b9g-646}

The ratio $\bar{M}_{\mathrm{w}} / \bar{M}_{\mathrm{n}}$ is called the (molar-mass) dispersity, $Đ$ (previously the 'polydispersity index', PDI), read 'd-stroke' and defined as

$$
Ð=\frac{\bar{M}_{\mathrm{w}}}{\bar{M}_{\mathrm{n}}}
$$

Dispersity [definition]\\
(14D.2)

The term 'monodisperse' is conventionally applied to synthetic polymers for which the dispersity is less than 1.1; commercial polyethene samples might be much more heterogeneous, with a dispersity close to 30 . One feature of a narrow molar-mass distribution for synthetic polymers is often a higher degree of long-range order in the solid and therefore higher density and melting point. The spread of values is controlled by the choice of catalyst and reaction conditions.

A note on good practice The masses of macromolecules are often reported in daltons $(\mathrm{Da})$, where $1 \mathrm{Da}=m_{\mathrm{u}}$ (with $m_{\mathrm{u}}=1.661$ $\times 10^{-27} \mathrm{~kg}$ ). Note that daltons are used to report molecular mass\\
not molar mass. So the mass (not the molar mass) of a certain macromolecule may be reported as 100 kDa (i.e. its mass is $100 \times$ $10^{3} \times m_{\mathrm{u}}$ ), and its molar mass as $100 \mathrm{~kg} \mathrm{~mol}^{-1}$. But it should not be said (even though it is common practice) that its molar mass is 100 kDa .

\subsection*{14D. 2 The different levels of structure}
The concept of the 'structure' of a macromolecule takes on different meanings at the different levels of the arrangement of the chain or network of monomers. The primary structure of a macromolecule is the sequence of small molecular residues making up the polymer. The residues may form either a chain, as in polyethene, or a more complex network in which cross-links connect different chains, as in cross-linked polyacrylamide. In a synthetic polymer, virtually all the residues are identical and it is sufficient to name the monomer used in the synthesis. Thus, the repeating unit of polyethene and its derivatives is $-\mathrm{CHXCH}_{2}-$, and the primary structure of the chain is specified by denoting it as $-\left(\mathrm{CHXCH}_{2}\right)_{n}-$. The concept of primary structure ceases to be trivial in the case of synthetic copolymers and biological macromolecules, because in general these substances are chains formed from different molecules. For example, proteins are polypeptides formed from different amino acids (about twenty occur naturally) strung together by the peptide link, -CONH- (1). The degradation of a biological macromolecule is a disruption of its primary structure, when the\\
\includegraphics{smile-10ff6c3e5518e5e801bb2e628e5082870d3f2e54}

1 Peptide link chain breaks into shorter components.

The term conformation refers to the spatial arrangement of the different parts of a chain, and one conformation can be changed into another by rotating one part of a chain around a bond. The conformation of a macromolecule is relevant at three levels of structure. The secondary structure of a macromolecule is the (often local) spatial arrangement of a chain. The secondary structure of a molecule of polyethene in some solvents may be a random coil (see below). In the absence of a solvent, polyethene forms crystals consisting of stacked sheets with a hairpin-like bend about every 100 monomer units, presumably because for that number of monomers the intermolecular (in this case intramolecular) potential energy is sufficient to overcome thermal disordering. The secondary structure of a protein is a highly organized arrangement determined largely by hydrogen bonds, and taking the form of random coils, helices (Fig. 14D.2a), or sheets in various segments of the molecule.

The tertiary structure is the overall three-dimensional structure of a macromolecule. For instance, the hypothetical protein shown in Fig. 14D.2b has helical regions connected\\
\includegraphics[max width=\textwidth, center]{2025_02_27_d4b95d9718f0b9e2e6b9g-647(3)}

Figure 14D. 2 (a) A polymer may adopt a highly organized helical conformation, an example of a secondary structure. The helix is represented as a cylinder. (b) Several helical segments connected by short random coils pack together, an example of tertiary structure.\\
\includegraphics[max width=\textwidth, center]{2025_02_27_d4b95d9718f0b9e2e6b9g-647(2)}

Figure 14D. 3 Several subunits with specific tertiary structures pack together, an example of quaternary structure.\\
by short random-coil sections. The helices interact to form a compact tertiary structure.

The quaternary structure of a macromolecule is the manner in which large molecules are formed by the aggregation of others. Figure 14D. 3 shows how four molecular subunits, each with a specific tertiary structure, aggregate. Quaternary structure can be very important in biology. For example, the oxygentransport protein haemoglobin consists of four myoglobin-like subunits that work cooperatively to take up and release $\mathrm{O}_{2}$.

\subsection*{14D.3 Random coils}
The most likely conformation of a chain of identical units not capable of forming hydrogen bonds or any other type of specific bond is a random coil. Polyethene is a simple example. The simplest model of a random coil is a 'freely jointed chain', in which any bond is free to make any angle with respect to the preceding one (Fig. 14D.4). It is also assumed that the mono-\\
\includegraphics[max width=\textwidth, center]{2025_02_27_d4b95d9718f0b9e2e6b9g-647}

Figure 14D. 4 A freely jointed chain is like a three-dimensional random walk, each step being in an arbitrary direction but of the same length.\\
\includegraphics[max width=\textwidth, center]{2025_02_27_d4b95d9718f0b9e2e6b9g-647(1)}

Figure 14D. 5 A better description is obtained by fixing the bond angle (for example, at the tetrahedral angle) and allowing free rotation about a bond direction.\\
mer units occupy zero volume, so different parts of the chain can occupy the same region of space. The model is obviously an oversimplification because a bond is actually constrained to a cone of angles around a direction defined by its neighbour (Fig. 14D.5) and real chains are self-avoiding in the sense that distant parts of the same chain cannot fold back and occupy the same space. In a hypothetical one-dimensional freely jointed chain all the monomer units lie in a straight line, and the angle between neighbours is either $0^{\circ}$ or $180^{\circ}$. The units of a three-dimensional freely jointed chain are not restricted to lie in a line or a plane.

\subsection*{(a) Measures of size}
The size of a freely jointed chain is related to the probability that its ends are a certain distance apart. That probability can be calculated by considering a one-dimensional random coil.

\subsection*{How is that done? 14D. 1 Calculating the probability distribution in a one-dimensional random coil}
Your goal is to calculate the probability, $P$, that the ends of a long one-dimensional freely jointed chain composed of $N$\\
units of length $l$ (and therefore of total length $N l$ ) are a distance $n l$ apart.\\
Step 1 Write expressions for the numbers of bonds pointing to the left or right\\
The conformation of a one-dimensional freely jointed chain can be described by stating the number of bonds pointing to the right $\left(N_{\mathrm{R}}\right)$ and the number pointing to the left $\left(N_{\mathrm{L}}\right)$. The distance between the ends of the chain is $\left(N_{\mathrm{R}}-N_{\mathrm{L}}\right) l$; it follows that $n=N_{\mathrm{R}}-N_{\mathrm{L}}$. The total number of units is $N=N_{\mathrm{R}}+N_{\mathrm{L}}$, therefore, $N_{\mathrm{R}}=\frac{1}{2}(N+n)$ and $N_{\mathrm{L}}=\frac{1}{2}(N-n)$.

Step 2 Write an expression for the probability that a polymer has a specified end-to-end separation\\
The probability, $P$, that the end-to-end separation of a randomly selected polymer is $n l$ is

$$
P=\frac{\text { number of conformations with end-to-end distance } n l}{\text { total number of possible conformations }}
$$

Each of the $N$ monomer units of the polymer may in principle lie to the left or the right, so the total number of possible conformations is $2^{N}$. The total number of ways, $W$, of forming a chain of $N$ units with the end-to-end distance $n l$ is the number of ways of having $N_{\mathrm{R}}$ right-pointing units, the rest being left-pointing units. Therefore, to calculate $W$, count the number of ways of achieving $N_{\mathrm{R}}$ right-pointing units given a total of $N$ units. This is the same problem as selecting $N_{\mathrm{R}}$ objects from a collection of $N$ objects (see Topic 13A), and is

$$
W=\frac{N!}{N_{\mathrm{R}}!\left(N-N_{\mathrm{R}}\right)!}=\frac{N!}{N_{\mathrm{R}}!N_{\mathrm{L}}!}=\frac{N!}{\left\{\frac{1}{2}(N+n)\right\}!\left\{\frac{1}{2}(N-n)\right\}!}
$$

It follows that

$$
P=\frac{W}{2^{N}}=\frac{N!}{\left\{\frac{1}{2}(N+n)\right\}!\left\{\frac{1}{2}(N-n)\right\}!2^{N}}
$$

\subsection*{Step 3 Consider the case of compact chains}
When the chain is compact in the sense that $n \ll N$, it is more convenient to evaluate $\ln P$ : the factorials are then large and it is possible to use Stirling's approximation (Topic 13A). Although the approximation used there is $\ln x!=x \ln x-x$ (with $x=N$ ), here it is appropriate to use the more precise form

$$
\ln x!\approx \ln (2 \pi)^{1 / 2}+\left(x+\frac{1}{2}\right) \ln x-x
$$

The result, after quite a lot of algebra, is

$$
\ln P \approx \ln \left(\frac{2}{\pi N}\right)^{1 / 2}-\frac{1}{2}(N+n+1) \ln (1+v)-\frac{1}{2}(N-n+1) \ln (1-v)
$$

where $v=n / N$. For a compact coil ( $v \ll 1$ ), use the approximation $\ln (1 \pm v) \approx \pm v-\frac{1}{2} v^{2}$ and obtain

$$
\ln P \approx \ln \left(\frac{2}{\pi N}\right)^{1 / 2}-\frac{1}{2} N v^{2}
$$

\begin{center}
\includegraphics[max width=\textwidth]{2025_02_27_d4b95d9718f0b9e2e6b9g-648}
\end{center}

Figure 14D. 6 The probability distribution for the separation of the ends of a one-dimensional random coil. The separation of the ends is $n l$, where $l$ is the length of each monomer unit.\\
which rearranges into\\
\includegraphics[max width=\textwidth, center]{2025_02_27_d4b95d9718f0b9e2e6b9g-648(1)}

This function is plotted in Fig. 14D.6.

\subsection*{Brief illustration 14D. 1}
Suppose that $N=1000$ and $l=150 \mathrm{pm}$, then the probability that the ends of a one-dimensional random coil are $n l=$ 3.00 nm apart is given by eqn 14D. 3 by setting $n=(3.00 \times$ $\left.10^{3} \mathrm{pm}\right) /(150 \mathrm{pm})=20.0$ :

$$
P=\left(\frac{2}{\pi \times 1000}\right)^{1 / 2} \mathrm{e}^{-20.0^{2} /(2 \times 1000)}=0.0207
$$

meaning that there is a 1 in 48 chance of being found there.

Equation 14D. 3 can be adapted to calculate the probability that the ends of a long three-dimensional freely jointed chain lie in the range $r$ to $r+\mathrm{d} r$. The probability is written as $f(r) \mathrm{d} r$, where

$$
\begin{aligned}
f(r) & =4 \pi\left(\frac{a}{\pi^{1 / 2}}\right)^{3} r^{2} \mathrm{e}^{-a^{2} r^{2}} \\
a & =\left(\frac{3}{2 N l^{2}}\right)^{1 / 2}
\end{aligned}
$$

Probability distribution [3D random coil]\\
(14D.4)

Here and elsewhere the fact that the chain cannot be longer than $N l$ is ignored. Although eqn 14D. 4 gives a non-zero probability for $r>N l$, the values are so small that the errors in pretending that $r$ can range up to infinity are negligible. For a narrow range of distances $\delta r$, the probability density can be treated as a constant and the probability calculated from $f(r) \delta r$. An alternative\\
interpretation of this expression is to regard each molecule in a sample as ceaselessly writhing from one conformation to another; then $f(r) \mathrm{d} r$ is the probability that at any instant the chain will be found with the separation of its ends between $r$ and $r+\mathrm{d} r$.

\subsection*{Brief illustration 14D. 2}
Consider the chain described in Brief illustration 14D.1, with $N=1000$ and $l=150 \mathrm{pm}$ but now in three dimensions. Then

$$
a=\left(\frac{3}{2 \times 1000 \times(150 \mathrm{pm})^{2}}\right)^{1 / 2}=2.58 \ldots \times 10^{-4} \mathrm{pm}^{-1}
$$

Then the probability density at $r=3.00 \mathrm{~nm}$ is given by eqn 14D. 4 as

$$
\begin{aligned}
f(3.00 \mathrm{~nm})= & 4 \pi \times\left(\frac{2.58 \ldots \times 10^{-4} \mathrm{pm}^{-1}}{\pi^{1 / 2}}\right)^{3} \\
& \times\left(3.00 \times 10^{3} \mathrm{pm}\right)^{2} \times \mathrm{e}^{-\left(2.58 . \ldots 10^{-4} \mathrm{pm}^{-1}\right)^{2}\left(3.0 \times 10^{3} \mathrm{pm}\right)^{2}} \\
= & 1.92 \times 10^{-4} \mathrm{pm}^{-1}
\end{aligned}
$$

The probability that the ends will be found in a narrow range of width $\delta r=10.0 \mathrm{pm}$ at 3.00 nm (regardless of direction) is therefore

$$
f(3.00 \mathrm{~nm}) \delta r=\left(1.92 \times 10^{-4} \mathrm{pm}^{-1}\right) \times(10.0 \mathrm{pm})=1.92 \times 10^{-3}
$$

or about 1 in 520 .

There are several measures of the geometrical size of a random coil. The contour length, $R_{c}$, is the length of the polymer (not only a random coil) measured along its backbone from atom to atom. For a polymer of $N$ monomer units each of length $l$, the contour length is

$$
R_{c}=N l
$$

Contour length\\
(14D.5)

The root-mean-square separation, $R_{\mathrm{rms}}$, is the square root of the mean value of the square of the separation of the ends of the coil. Thus, if the vector joining the ends of the coil is $\boldsymbol{R}$, and each monomer is represented by the vector $\boldsymbol{r}_{\boldsymbol{i}}$, then $\boldsymbol{R}=\sum_{i=1}^{N} \boldsymbol{r}_{i}$ and

$$
\left\langle R^{2}\right\rangle=\langle\boldsymbol{R} \cdot \boldsymbol{R}\rangle=\sum_{i, j=1}^{N}\left\langle\boldsymbol{r}_{i} \cdot \boldsymbol{r}_{j}\right\rangle=\sum_{i=1}^{N} \overbrace{\left\langle r_{i}^{2}\right\rangle}^{I^{2}}+\sum_{i \neq 1}^{N}\left\langle\boldsymbol{r}_{i} \cdot \boldsymbol{r}_{j}\right\rangle
$$

When $N$ is large (as assumed throughout), the term in blue is zero because the individual vectors lie in random directions. The remaining term is $N l^{2}$. It follows that for a random coil of any dimensionality

$$
R_{\mathrm{rms}}=N^{1 / 2} l
$$

Root-mean-square separation [random coil]\\
(14D.6)\\
\includegraphics[max width=\textwidth, center]{2025_02_27_d4b95d9718f0b9e2e6b9g-649}

Figure 14D. 7 The variation of the root-mean-square separation of the ends of a three-dimensional random coil, $R_{\text {rms }}$, with the number of monomers.

As the number of monomer units increases, the root-meansquare separation of the ends of the polymer increases as $N^{1 / 2}$ (Fig. 14D.7), and consequently the volume of a three-dimensional coil increases as $N^{3 / 2}$. The result must be multiplied by a factor when the chain is not freely jointed (see below).

Another convenient measure of size is the radius of gyration, $R_{\mathrm{g}}$, which is the radius of a hollow sphere that has the same mass and moment of inertia (and therefore rotational characteristics) as the actual molecule. Once again, a one-dimensional random coil can be used to illustrate the procedure for the calculation of $R_{g}$.

\subsection*{How is that done? 14D. 2 Deriving an expression for the}
 radius of gyrationYou need to set up an expression for the moment of inertia of the random one-dimensional coil of $N$ monomer units each of mass $m$ and then equate it to $m_{\text {total }} R_{\mathrm{g}}^{2}$, where $m_{\text {total }}$ is the total mass of the polymer molecule, $m_{\text {total }}=\mathrm{Nm}$.

Step 1 Set up the expression for the moment of inertia\\
For a one-dimensional random coil with $N$ identical monomers each of mass $m$, the moment of inertia around the centre of the chain (which is at the origin of the vector representing the first monomer, because the vectors point in equal numbers to left and right) is

$$
I=\sum_{i=1}^{N} m_{i} d_{i}^{2}=m \sum_{i=1}^{N} d_{i}^{2}
$$

where $d_{i}$ is the distance of mass $m_{i}$ from the origin. This distance is the length of the vector $\boldsymbol{d}_{i}$, the sum of $i$ steps from the origin, $\boldsymbol{d}_{i}=\sum_{j=1}^{i} \boldsymbol{r}_{j}$.

Step 2 Evaluate the average distance of a monomer from the origin\\
As in the calculation that led to eqn 14D.6, write

$$
\left\langle d_{i}^{2}\right\rangle=\left\langle\boldsymbol{d}_{i} \cdot \boldsymbol{d}_{i}\right\rangle=\sum_{j, k=1}^{i}\left\langle\boldsymbol{r}_{j} \cdot \boldsymbol{r}_{k}\right\rangle=\overbrace{\sum_{j=1}^{i} \overbrace{\left\langle r_{j}^{2}\right\rangle}^{\rho^{2}}}^{\left.i\right|^{2}}+\sum_{j \neq k}^{i}\left\langle\boldsymbol{r}_{j} \cdot \boldsymbol{r}_{k}\right\rangle
$$

Again, the blue term is zero for a random coil, so $\left\langle d_{i}^{2}\right\rangle=i l^{2}$ and the average moment of inertia of the coil (recognizing that there is a monomer on both sides of the origin at a given distance) is

$$
\langle I\rangle=2 m \sum_{i=1}^{N}\left\langle d_{i}^{2}\right\rangle=2 m l^{2} \sum_{i=1}^{N} i=N^{2} m l^{2}=N m_{\text {total }} l^{2}
$$

\subsection*{Step 3 Identify the radius of gyration}
Finally, set this moment of inertia equal to $m_{\text {total }} R_{\mathrm{g}}^{2}$, which implies that $R_{\mathrm{g}}^{2}=\mathrm{Nl}^{2}$ and therefore that\\
\includegraphics[max width=\textwidth, center]{2025_02_27_d4b95d9718f0b9e2e6b9g-650}

A similar calculation for a three-dimensional random coil gives

$$
R_{\mathrm{g}}=\left(\frac{N}{6}\right)^{1 / 2} l
$$

Radius of gyration [3D random coil]\\
(14D.7b)

The radius of gyration is smaller in this case because the extra dimensions enable the coil to be more compact.\\
The radius of gyration may also be calculated for other geometries. For example, a solid uniform sphere of radius $R$ has $R_{\mathrm{g}}=\left(\frac{3}{5}\right)^{1 / 2} R$, and a long thin uniform rod of length $L$ has $R_{\mathrm{g}}$ $=L / 12^{1 / 2}$ for rotation about an axis perpendicular to the long axis. A solid sphere with the same radius and mass as a random coil has a greater radius of gyration as it is entirely dense throughout.

\subsection*{Brief illustration 14D. 3}
Consider a polymer that writhes as if it were a three-dimensional random coil. However, suppose that small segments of the macromolecule resist bending, so it is more appropriate to visualize it as a freely jointed chain with $N$ and $l$ as the number and length, respectively, of these rigid units. With the length $l=45 \mathrm{~nm}$ and $N=200$ (and using $10^{3} \mathrm{~nm}=1 \mu \mathrm{~m}$ ),

$$
\begin{aligned}
& \text { From eqn 14D.5: } R_{\mathrm{c}}=200 \times 45 \mathrm{~nm}=9.0 \mu \mathrm{~m} \\
& \text { From eqn 14D.6: } R_{\mathrm{rms}}=(200)^{1 / 2} \times 45 \mathrm{~nm}=0.64 \mu \mathrm{~m} \\
& \text { From eqn 14D.7b: } R_{\mathrm{g}}=\left(\frac{200}{6}\right)^{1 / 2} \times 45 \mathrm{~nm}=0.26 \mu \mathrm{~m}
\end{aligned}
$$

The random coil model ignores the role of the solvent: a poor solvent tends to cause the coil to tighten so that solutesolvent contacts are minimized; a good solvent does the opposite. Therefore, calculations based on this model are better regarded as lower bounds to the dimensions for a polymer in a good solvent and as an upper bound for a polymer in a poor solvent. The model is most reliable for a polymer in a bulk solid sample, where the coil is likely to have its natural dimensions.

\subsection*{(b) Constrained chains}
The freely jointed chain model is improved by removing the freedom of bond angles to take any value. For long chains, it is convenient to take groups of neighbouring bonds and consider the direction of their resultant. Although each successive individual bond is constrained to a single cone of angle $\theta$ relative to its neighbour, the resultant of several bonds lies in a random direction. By concentrating on such groups rather than individuals, it turns out that for long chains the expressions for the root-mean-square separation and the radius of gyration given above should be multiplied by


\begin{equation*}
F=\left(\frac{1-\cos \theta}{1+\cos \theta}\right)^{1 / 2} \tag{14D.8}
\end{equation*}


For a tetrahedral arrangement of bonds, for which $\cos \theta=-\frac{1}{3}$ (i.e. $\theta=109.5^{\circ}$ ), $F=2^{1 / 2}$. Therefore:

$$
R_{\mathrm{rms}}=(2 N)^{1 / 2} l \quad R_{\mathrm{g}}=\left(\frac{N}{3}\right)^{1 / 2} l \begin{aligned}
& \text { Dimensions of } \\
& \text { a constrained } \\
& \text { tetrahedral chain }
\end{aligned}
$$

(14D.9)

The model of a randomly coiled molecule is still an approximation, even after the bond angles have been restricted, because it does not take into account the impossibility of two or more atoms occupying the same place. Such self-avoidance tends to swell the coil, so (in the absence of solvent effects) it is better to regard $R_{\mathrm{rms}}$ and $R_{\mathrm{g}}$ as lower bounds to the actual values.

\subsection*{(c) Partly rigid coils}
An important measure of the flexibility of a chain is the persistence length, $l_{\mathrm{p}}$, a measure of the length over which the direction of the first monomer-monomer direction is sustained. If the chain is a rigid rod, then the persistence length is the same as the contour length. For a freely-jointed random coil, the persistence length is just the length of one monomer unit. Therefore, the persistence length can be regarded as a measure of the stiffness of the chain.

The mean square distance between the ends of a chain that has a persistence length greater than the monomer length can be expected to be greater than for a random coil because the partial rigidity of the coil does not let it roll up so tightly. A detailed calculation shows that


\begin{equation*}
R_{\mathrm{rms}}=N^{1 / 2} l F \text { where } F=\left(\frac{2 l_{\mathrm{p}}}{l}-1\right)^{1 / 2} \tag{14D.10}
\end{equation*}


For a random coil, $l_{\mathrm{p}}=l$, so $R_{\mathrm{rms}}=N^{1 / 2} l$, as already found. For $l_{\mathrm{p}}>l, F>1$, so the coil has swollen, as anticipated.

\subsection*{Example 14D. 2}
Calculating the root-mean-square\\
separation of a partly rigid coil\\
By what percentage does the root-mean-square separation of the ends of a polymer chain with $N=1000$ increase or decrease when the persistence length changes from $l$ (the length of one monomer unit) to 2.5 per cent of the contour length?

Collect your thoughts The contour length is $R_{\mathrm{c}}=N l$. When $l_{\mathrm{p}}=l$, the chain is a random coil and $R_{\mathrm{rms}, \text { random coil }}=N^{1 / 2} l$, so eqn 14D. 10 can be expressed as $R_{\mathrm{rms}}=F R_{\mathrm{rms}, \text { random coil }}$. The fractional change in root-mean-square separation is therefore

$$
\frac{R_{\mathrm{rms}}-R_{\mathrm{rmss}, \text { random coil }}}{R_{\mathrm{rms}, \text { andom coil }}}=\frac{R_{\mathrm{rms}}}{R_{\mathrm{rms}, \text { random coil }}}-1=\left(\frac{2 l_{\mathrm{p}}}{l}-1\right)^{1 / 2}-1
$$

In the final step you should express this fractional change as a percentage.

The solution Because $l_{\mathrm{p}}=0.025 R_{\mathrm{c}}=0.025 \mathrm{Nl}$, the fractional change is

$$
\begin{aligned}
\frac{R_{\mathrm{rms}}-R_{\mathrm{rms}, \text { random coil }}}{R_{\mathrm{rms}, \text { andom coil }}} & =\left(\frac{2 \times 0.025 \mathrm{Nl}}{l}-1\right)^{1 / 2}-1 \\
& =(0.050 \mathrm{~N}-1)^{1 / 2}-1
\end{aligned}
$$

With $N=1000$, the fractional change is 6.00 , so the root-mean-square separation increases by 600 per cent.

Self-test 14D. 2 Calculate the fractional change in the volume of the same three-dimensional coil.

\subsection*{14D.4 Mechanical properties}
Insight into the consequences of stretching and contracting a polymer can be obtained on the basis of the freely jointed chain as a model.

\subsection*{(a) Conformational entropy}
A random coil is the least structured conformation of a polymer chain and therefore corresponds to the state of greatest entropy. Any stretching of the coil reduces disorder and reduces the entropy. Conversely, the formation of a random coil from a more extended form is spontaneous (provided enthalpy contributions do not interfere). The same model can be used to deduce an expression for the change in conformational entropy, the statistical entropy arising from the arrangement of bonds, when a one-dimensional chain is stretched or compressed.

\subsection*{How is that done? 14D. 3 Deriving an expression for the}
 conformational entropy of a freely jointed chainConsider a freely jointed one-dimensional chain containing $N$ units of length $l$ that is stretched or compressed through a distance $x$. You then need to use the Boltzmann formula (eqn 13E.7, $S=k \ln W$ ) to calculate the conformational entropy of the chain, which involves assessing the value of $W$, the number of ways of achieving a particular conformation.

\subsection*{Step 1 Calculate $W$}
To achieve an extension, the number of steps to the right $\left(N_{R}\right)$ must be greater than the number to the left $\left(N_{\mathrm{L}}\right)$, so with $N_{\mathrm{L}}+$ $N_{\mathrm{R}}=N$ write $N_{\mathrm{R}}-N_{\mathrm{L}}=\lambda N$, with $\lambda$ between -1 (all to the left) and 1 (all to the right). Then $N_{\mathrm{R}}=\frac{1}{2}(1+\lambda) N$ and $N_{\mathrm{L}}=\frac{1}{2}(1-\lambda) N$ and the distance stretched is $x=\lambda N l$, or $\lambda R_{\mathrm{c}}$. The number of ways of taking these numbers of steps (as in the earlier discussion of the random coil) is

$$
W=\frac{N!}{\mathrm{N}_{\mathrm{R}}!\mathrm{N}_{\mathrm{L}}!}=\frac{N!}{\left\{\frac{1}{2}(1+\lambda) N\right\}!\left\{\frac{1}{2}(1+\lambda) N\right\}!}
$$

Step 2 Write an expression for $S$\\
It follows from the expression for $W$ and the Boltzmann formula that

$$
S / k=\ln N!-\ln \left\{\frac{1}{2}(1+\lambda) N\right\}!-\ln \left\{\frac{1}{2}(1-\lambda) N\right\}!
$$

Because the factorials are large (except for large extensions), use Stirling's approximation in the form $\ln x!\approx$ $\left(x+\frac{1}{2}\right) \ln x-x+\frac{1}{2} \ln (2 \pi)$ to obtain

$$
\begin{aligned}
S / k= & -\ln (2 \pi)^{1 / 2}+(N+1) \ln 2 \\
& +\left(N+\frac{1}{2}\right) \ln N-\frac{1}{2} \ln \left\{N^{2 N+2}(1+\lambda)^{N(1+\lambda)+1}(1+\lambda)^{N(1-\lambda)+1}\right\}
\end{aligned}
$$

Step 3 Write an expression for the change in entropy When the coil is not extended, and adopts its most random conformation ( $\lambda=0$ ), the entropy is

$$
S / k=-\ln (2 \pi)^{1 / 2}+(N+1) \ln 2-\frac{1}{2} \ln N
$$

\begin{center}
\includegraphics[max width=\textwidth]{2025_02_27_d4b95d9718f0b9e2e6b9g-652}
\end{center}

Figure 14D. 8 The change in entropy of a perfect elastomer as its extension changes; $\lambda= \pm 1$ corresponds to complete extension in either direction; $\lambda=0$, the conformation of highest entropy, corresponds to a random coil.

The change in entropy when the chain is stretched or compressed by the distance $\lambda R_{c}$ is therefore the difference between this quantity and that from Step 2. The resulting expression, after some algebraic manipulation and using $N \gg 1$, is

$$
\left.\left.\begin{array}{rl}
-\Delta S=-\frac{1}{2} k N \ln \left\{(1+\lambda)^{1+\lambda}(1-\lambda)^{1-\lambda}\right\} & \text { with } \lambda=x / R_{c}
\end{array} \right\rvert\, \text { (14D.11) }\right) ~\left[\begin{array}{l}
\text { Conformational entropy change } \\
\text { [1D random coil] }
\end{array}\right.
$$

This function is plotted in Fig. 14D.8, and it is seen that minimum extension corresponds to maximum entropy.

\subsection*{Brief illustration 14D. 4}
Suppose that $N=1000$ and $l=150 \mathrm{pm}$, so $R_{c}=150 \mathrm{~nm}$. The change in entropy when the (one-dimensional) random coil is stretched through 1.5 nm (corresponding to $\lambda=1 / 100$ ) is

$$
\begin{aligned}
\Delta S & =-\frac{1}{2} k \times(1000) \times \ln \left\{\left(1+\frac{1}{100}\right)^{1+(1 / 100)}\left(1-\frac{1}{100}\right)^{1-(1 / 100)}\right\} \\
& =-0.050 k
\end{aligned}
$$

Because $R=N_{\mathrm{A}} k$, the change in molar entropy is $\Delta S_{\mathrm{m}}=-0.050 R$, or $-0.42 \mathrm{~J} \mathrm{~K}^{-1} \mathrm{~mol}^{-1}$.

\subsection*{(b) Elastomers}
An elastomer is a flexible polymer that can expand or contract easily upon application of an external force. Elastomers are polymers with numerous crosslinks that pull them back into their original shape when a stress is removed. The weak directional constraints on silicon-oxygen bonds are responsible for the high elasticity of silicones. Even a freely jointed chain be-\\
haves as an elastomer for small extensions. It is a model of a 'perfect elastomer', a polymer in which the internal energy is independent of the extension, and can be used to deduce the restoring force associated with stretching or compression of the chain.

\subsection*{How is that done? 14D. 4 Deriving an expression for the restoring force of a perfect elastomer}
Your goal is to find an expression for the restoring force, $F$, of an elastomer, modelled as a one-dimensional random coil composed of $N$ units each of length $l$, when the chain is stretched or compressed by a distance $x=v l$.

Step 1 Use thermodynamics to relate the restoring force to the entropy\\
The work done on an elastomer when it is extended reversibly through a distance $\mathrm{d} x$ is $F \mathrm{~d} x$, The change in internal energy, from $\mathrm{d} U=\mathrm{d} w_{\text {rev }}+\mathrm{d} q_{\text {rev }}$ with $\mathrm{d} q_{\text {rev }}=T \mathrm{~d} S$ is therefore $\mathrm{d} U=F \mathrm{~d} x+$ $T \mathrm{~d} S$. It follows that for an isothermal extension

$$
\left(\frac{\partial U}{\partial x}\right)_{T}=F+T\left(\frac{\partial S}{\partial x}\right)_{T}
$$

In a perfect elastomer, as in a perfect gas, the internal energy is independent of the dimensions (at constant temperature), so $(\partial U / \partial x)_{T}=0$. The restoring force is therefore

$$
F=-T\left(\frac{\partial S}{\partial x}\right)_{T}
$$

Step 2 Evaluate the force from the change in conformational entropy\\
The conformational entropy (eqn 14D.11) is expressed in terms of the parameter $\lambda$ used to express the extension $x$ as $x=\lambda N l$ (or $x=\lambda R_{c}$ ). Therefore, replace the derivative with respect to $x$ by the derivative with respect to $\lambda$ by noting that $\mathrm{d} x=N l \mathrm{~d} \lambda$. Then

$$
F=-\frac{T}{N l}\left(\frac{\partial S}{\partial \lambda}\right)_{T}=-\frac{T}{N l}\left(\frac{\partial \Delta S}{\partial \lambda}\right)_{T}
$$

The replacement of $S$ by the change $\Delta S$ is valid because the initial value of the entropy is independent of the extension being applied. Now use eqn 14D. 11 to obtain

$$
\begin{aligned}
F & =\frac{T}{N l} \times \frac{N k}{2} \times \frac{\mathrm{d}}{\mathrm{~d} \lambda} \ln \left\{(1+\lambda)^{1+\lambda}(1-\lambda)^{1-\lambda}\right\} \\
& =\frac{k T}{2 l} \frac{\mathrm{~d}}{\mathrm{~d} \lambda}\{(1+\lambda) \ln (1+\lambda)+(1-\lambda) \ln (1-\lambda)\} \\
& =\frac{k T}{2 l}\{\ln (1+\lambda)-\ln (1-\lambda)\}
\end{aligned}
$$

That is,


\begin{equation*}
\left.-F=\frac{k T}{2 l} \ln \left(\frac{1+\lambda}{1-\lambda}\right) \quad \lambda=\frac{x}{N l} \right\rvert\, \tag{14D.12a}
\end{equation*}


[1D random coil]\\
\includegraphics[max width=\textwidth, center]{2025_02_27_d4b95d9718f0b9e2e6b9g-653}

Figure 14D. 9 The restoring force, $F$, of a one-dimensional perfect elastomer. For small extensions, $F$ is proportional to the extension, corresponding to Hooke's law.

For small displacements $(\lambda \ll 1$, corresponding to $x \ll N l$ and therefore $x \ll R_{c}$ ) the logarithms can be expanded by using $\ln (1+\lambda) \approx \lambda$ and $\ln (1-\lambda) \approx-\lambda$, to give

$$
F \approx \frac{\lambda k T}{l}=\frac{k T}{N l^{2}} x \quad \begin{aligned}
& \text { Restoring force } \\
& \text { [1D random coil] }
\end{aligned}
$$

(14D.12b)

That is, for small displacements the sample obeys Hooke's law (Fig. 14D.9): the restoring force is proportional to the displacement and the force constant $k_{\mathrm{f}}$ (the constant of proportionality between the force and the displacement) is


\begin{equation*}
k_{\mathrm{f}}=\frac{k T}{N l^{2}} \tag{14D.12c}
\end{equation*}


\subsection*{Brief illustration 14D. 5}
Consider a polymer chain with $N=5000$ and $l=0.15 \mathrm{~nm}$. If the ends of the chain are moved apart by $x=1.5 \mathrm{~nm}$, then $\lambda=(1.5 \mathrm{~nm}) /(5000 \times 0.15 \mathrm{~nm})=2.0 \times 10^{-3}$. Because $\lambda \ll 1$, the restoring force at 293 K is given by eqn 14D.12b as

$$
F=\frac{\left(1.381 \times 10^{-23} \stackrel{\sim}{\mathrm{Jm}} \mathrm{~J}^{-1}\right) \times(293 \mathrm{~K})}{5000 \times\left(1.5 \times 10^{-10} \mathrm{~m}\right)^{2}} \times 1.5 \times 10^{-9} \mathrm{~m}=5.4 \times 10^{-14} \mathrm{~N}
$$

or 54 fN .\\
\includegraphics[max width=\textwidth, center]{2025_02_27_d4b95d9718f0b9e2e6b9g-653(1)}

Figure 14D. 10 The variation of specific volume with temperature of a synthetic polymer. The glass transition temperature, $T_{g}$, is at the point of intersection of extrapolations of the two linear parts of the curve.

\subsection*{14D. 5 Thermal properties}
The crystallinity of synthetic polymers can be destroyed by thermal motion at sufficiently high temperatures. This change in crystallinity may be thought of as a kind of intramolecular melting from a crystalline solid to a more fluid random coil. Polymer melting also occurs at a specific melting temperature, $T_{\mathrm{m}}$, which increases with the strength and number of intermolecular interactions in the material. Thus, polyethene, which has chains that interact only weakly in the solid, has $T_{\mathrm{m}}=414 \mathrm{~K}$ and nylon- 66 fibres, in which there are strong hydrogen bonds between chains, have $T_{\mathrm{m}}=530 \mathrm{~K}$. High melting temperatures are desirable in most practical applications involving fibres and plastics.

All synthetic polymers undergo a transition from a state of high to low chain mobility at the glass transition temperature, $T_{\mathrm{g}}$. To visualize the glass transition, consider what happens to an elastomer as its temperature is lowered. There is sufficient energy available at normal temperatures for limited bond rotation to occur and the flexible chains writhe. At lower temperatures, the amplitudes of the writhing motion decrease until a specific temperature, $T_{\mathrm{g}}$, is reached at which motion is frozen completely and the sample forms a glass. Glass transition temperatures well below 300 K are desirable in elastomers that are to be used at normal temperatures. Both the glass transition temperature and the melting temperature of a polymer may be measured by calorimetric methods. Because the motion of the segments of a polymer chain increase at the glass transition temperature, $T_{\mathrm{g}}$ may also be determined from a plot of the specific volume of a polymer (the reciprocal of its mass density) against temperature (Fig. 14D.10).

\subsection*{Checklist of concepts}
$\square$ 1. Macromolecules are very large molecules assembled from smaller molecules.2. Synthetic polymers are manufactured by stringing together and in some cases cross-linking smaller units known as monomers.3. Macromolecules can be monodisperse, with a single molar mass, or polydisperse, with a spread of molar mass.4. The conformation of a macromolecule is the spatial arrangement of the different parts of a chain.5. The primary structure of a macromolecule is the sequence of small molecular residues making up the polymer.5. The secondary structure is the spatial arrangement of a chain of residues.6. The tertiary structure is the overall three-dimensional structure of a macromolecule.7. The quaternary structure is the manner in which large molecules are formed by the aggregation of others.8. In a freely jointed chain any bond in a polymer is free to make any angle with respect to the preceding one.10. The least structured conformation of a macromolecule is a random coil, which can be modelled as a freely jointed chain.11. An elastomer is a flexible polymer that can expand or contract easily upon application of an external force.\\
$\square$ 12. The disruption of long-range order in a polymer occurs at a melting temperature.\\
$\square$ 13. Synthetic polymers undergo a transition from a state of high to low chain mobility at the glass transition temperature.

\subsection*{Checklist of equations}
\begin{center}
\begin{tabular}{|c|c|c|c|}
\hline
Property & Equation & Comment & Equation number \\
\hline
Number-average molar mass & $\bar{M}_{\mathrm{n}}=\frac{1}{N_{\text {total }}} \sum_{i} N_{i} M_{i}$ & Definition & 14D.1a \\
\hline
Weight-average molar mass & $\bar{M}_{\text {w }}=\frac{1}{m_{\text {total }}} \sum_{i} m_{i} M_{i}$ & Definition & 14D.1b \\
\hline
Dispersity & $甲=\bar{M}_{\mathrm{w}} / \bar{M}_{\mathrm{n}}$ & Definition & 14D. 2 \\
\hline
\multirow[t]{2}{*}{Probability distribution} & $P=(2 / \pi N)^{1 / 2} \mathrm{e}^{-n^{2} / 2 N}$ & One-dimensional random coil & 14D. 3 \\
\hline
 & \( \begin{aligned} & f(r)=4 \pi\left(a / \pi^{1 / 2}\right)^{3} r^{2} \mathrm{e}^{-a^{2} r^{2}} \\ & a=\left(3 / 2 N l^{2}\right)^{1 / 2} \end{aligned} \) & Three-dimensional random coil & 14D. 4 \\
\hline
Contour length of a random coil & $R_{\text {c }}=N l$ &  & 14D. 5 \\
\hline
Root-mean-square separation of a random coil & $R_{\text {rms }}=N^{1 / 2} l$ & Unconstrained chain & 14D. 6 \\
\hline
\multirow[t]{2}{*}{Radius of gyration of a random coil} & $R_{\mathrm{g}}=N^{1 / 2} l$ & Unconstrained one-dimensional chain & 14D.7a \\
\hline
 & $R_{\mathrm{g}}=(N / 6)^{1 / 2} l$ & Unconstrained three-dimensional chain & 14D.7b \\
\hline
Root-mean-square separation of a random coil & $R_{\text {rms }}=(2 N)^{1 / 2} l$ & Constrained tetrahedral chain & 14D. 9 \\
\hline
Change in conformational entropy on extending a random coil & $\Delta S=-\frac{1}{2} k N \ln \left\{(1+\lambda)^{1+\lambda}(1-\lambda)^{1-\lambda}\right\}$ &  & 14D. 11 \\
\hline
\multirow[t]{2}{*}{Restoring force of a one-dimensional random coil} & $F=(k T / 2 l) \ln \{(1+\lambda) /(1-\lambda)\}$ &  & 14D.12a \\
\hline
 & $F \approx\left(k T / N l^{2}\right) x$ & $x \ll R_{\text {c }}$ & 14D.12b \\
\hline
\end{tabular}
\end{center}

\subsection*{Temperature dependence of polymer melt viscosity: the WLF equation}

The Arrhenius expression $\eta = \eta_0 \mathrm{e}^{E_\mathrm{a}/RT}$ (eqn 16B.2) describes the viscosity of simple liquids but fails for polymer melts near the glass transition temperature $T_\mathrm{g}$, where free volume becomes the controlling factor. Above $T_\mathrm{g}$ the temperature dependence is well described by the Williams--Landel--Ferry (WLF) equation:

\begin{equation*}
\log a_T = \log \frac{\eta(T)}{\eta(T_\mathrm{ref})} = \frac{-C_1 (T - T_\mathrm{ref})}{C_2 + (T - T_\mathrm{ref})} \tag{14D.13}
\end{equation*}

where $a_T$ is the shift factor. When the reference temperature is chosen as $T_\mathrm{ref} = T_\mathrm{g}$, the universal (empirical) constants are $C_1 = 17.44$ and $C_2 = 51.6\ \mathrm{K}$. These values apply to a wide range of amorphous polymers above $T_\mathrm{g}$. At temperatures well above $T_\mathrm{g}$ the WLF equation converges toward Arrhenius behaviour. The WLF equation is used to construct time--temperature superposition master curves and to estimate relaxation times, creep compliances, and diffusion coefficients of polymer melts at temperatures not directly measured.

\subsection*{Equilibrium swelling and crosslink density: the Flory--Rehner equation}

A crosslinked polymer network swells in a compatible solvent until an equilibrium is reached between the osmotic driving force for mixing and the elastic restoring force of the network. The crosslink density $\nu_e$ (moles of network strands per unit volume) can be determined from the equilibrium volume fraction of polymer $\phi_2$ using the Flory--Rehner equation:

\begin{equation*}
\nu_e = \frac{-\left[\ln(1-\phi_2) + \phi_2 + \chi \phi_2^2\right]}{V_1\!\left(\phi_2^{1/3} - \phi_2/2\right)} \tag{14D.14}
\end{equation*}

where $V_1$ is the molar volume of the solvent, and $\chi$ is the Flory--Huggins polymer--solvent interaction parameter (dimensionless; typically $0.3$--$0.5$ for good solvents, $0.5$ for the theta condition). The volume fraction of polymer at swelling equilibrium is $\phi_2 = V_{2,\mathrm{dry}}/V_{2,\mathrm{swollen}}$, obtained from the dry and swollen masses together with the polymer and solvent densities. The mean molar mass between crosslinks is $M_c = \rho_2 / \nu_e$, where $\rho_2$ is the dry polymer density.

\subsection*{Copolymer composition: the Mayo--Lewis equation}

In free-radical copolymerization of two monomers $\mathrm{M}_1$ and $\mathrm{M}_2$, the instantaneous mole fraction $F_1$ of $\mathrm{M}_1$ units incorporated into the copolymer is related to the feed composition by the Mayo--Lewis (terminal) equation:

\begin{equation*}
F_1 = \frac{r_1 f_1^2 + f_1 f_2}{r_1 f_1^2 + 2 f_1 f_2 + r_2 f_2^2} \tag{14D.15}
\end{equation*}

where $f_1 = [\mathrm{M}_1]/([\mathrm{M}_1]+[\mathrm{M}_2])$ and $f_2 = 1 - f_1$ are the mole fractions of the two monomers in the feed, and $r_1 = k_{11}/k_{12}$, $r_2 = k_{22}/k_{21}$ are the reactivity ratios ($k_{ij}$ is the rate constant for addition of monomer $\mathrm{M}_j$ to a chain ending in $\mathrm{M}_i$). When $r_1 r_2 = 1$ (ideal copolymerization) the equation simplifies to $F_1 = r_1 f_1 / (r_1 f_1 + f_2)$. The feed ratio $f_1/f_2$ required to achieve a desired copolymer composition $F_1$ can be solved from eqn 14D.15.

\section*{TOPIC 14E Self-assembly}
\subsection*{Why do you need to know this material?}
Aggregates of small and large molecules form the basis of many established and emerging technologies. To see why this is the case, you need to understand their structures and properties.

\subsection*{What is the key idea?}
Colloids and micelles form spontaneously by self-assembly of molecules or macromolecules and are held together by molecular interactions.

\subsection*{What do you need to know already?}
You need to be familiar with molecular interactions (Topic 14B) and interactions between ions (Topic 5E).

Self-assembly is the spontaneous formation of complex structures of molecules or macromolecules that are held together by molecular interactions, such as Coulombic, dispersion, hydrogen bonding, and hydrophobic interactions. Examples of self-assembly include the formation of liquid crystals, and of protein quaternary structures from two or more polypeptide chains (Topic 14C).

\subsection*{14E. 1 Colloids}
A colloid, or disperse phase, is a dispersion of small particles of one material in another that does not settle out under gravity. In this context, 'small' means that one dimension at least is smaller than about 500 nm (about the wavelength of visible light). Many colloids are suspensions of nanoparticles (particles of diameter up to about 100 nm ). In general, colloidal particles are aggregates of numerous atoms or molecules, but are commonly but not universally too small to be seen with an ordinary optical microscope.

\subsection*{(a) Classification and preparation}
The name given to the colloid depends on the two phases involved:

\begin{itemize}
  \item A sol is a dispersion of a solid in a liquid (such as clusters of gold atoms in water) or of a solid in a solid (such as\\
ruby glass, which is a gold-in-glass sol, and achieves its colour by light scattering).
  \item An aerosol is a dispersion of a liquid in a gas (like fog and many sprays) or a solid in a gas (such as smoke): the particles are often large enough to be seen with a microscope.
  \item An emulsion is a dispersion of a liquid in a liquid (such as milk).
  \item A foam is a dispersion of a gas in a liquid.
\end{itemize}

A further classification of colloids is as lyophilic, or solvent attracting, and lyophobic, solvent repelling. If the solvent is water, the terms hydrophilic and hydrophobic, respectively, are used instead. Lyophobic colloids include the metal sols. Lyophilic colloids generally have some chemical similarity to the solvent, such as -OH groups able to form hydrogen bonds. A gel is a semi-rigid mass of a lyophilic sol.

The preparation of aerosols can be as simple as sneezing (which produces an imperfect aerosol). Laboratory and commercial methods make use of several techniques. Material (e.g. quartz) may be ground in the presence of the dispersion medium. Passing a heavy electric current through a cell may lead to the sputtering (crumbling) of an electrode into colloidal particles. Arcing between electrodes immersed in the support medium also produces a colloid. Chemical precipitation sometimes results in a colloid. A precipitate (e.g. silver iodide) already formed may be dispersed by the addition of a 'peptizing agent' (e.g. potassium iodide). Clays may be peptized by alkalis, the $\mathrm{OH}^{-}$ion being the active agent.

Emulsions are normally prepared by shaking the two components together vigorously, although some kind of emulsifying agent usually has to be added to stabilize the product. This emulsifying agent may be a soap (the salt of a long-chain carboxylic acid) or other surfactant (surface active) species, or a lyophilic sol that forms a protective film around the dispersed phase. In milk, which is an emulsion of fats in water, the emulsifying agent is casein, a protein containing phosphate groups. It is clear from the formation of cream on the surface of milk that casein is not completely successful in stabilizing milk: the dispersed fats coalesce into oily droplets which float to the surface. This coagulation may be prevented by ensuring that the emulsion is dispersed very finely initially: intense agitation with ultrasonics brings this dispersion about, the product being 'homogenized' milk.

One way to form an aerosol is to tear apart a spray of liquid with a jet of gas. The dispersal is aided if a charge is applied to the liquid, for then electrostatic repulsions help to blast it\\
apart into droplets. This procedure may also be used to produce emulsions, for the charged liquid phase may be directed into another liquid.

Colloids are often purified by dialysis, the process of squeezing the solution though a membrane. The aim is to remove much (but not all, for reasons explained later) of the ionic material that may have accompanied their formation. A membrane (for example, cellulose) is selected that is permeable to solvent and ions, but not to the colloid particles. Dialysis is very slow, and is normally accelerated by applying an electric field and making use of the charges carried by many colloidal particles; the technique is then called electrodialysis.

\subsection*{(b) Structure and stability}
Colloids are thermodynamically unstable with respect to the bulk. This instability can be expressed thermodynamically by noting that because the change in Helmholtz energy, $\mathrm{d} A$, when the surface area of the sample changes by $\mathrm{d} \sigma$ at constant temperature and pressure is $\mathrm{d} A=\gamma \mathrm{d} \sigma$, where $\gamma$ is the interfacial surface tension (Topic 14C), it follows that $\mathrm{d} A<0$ if $\mathrm{d} \sigma<0$. That is, the contraction of the surface $(\mathrm{d} \sigma<0)$ is spontaneous ( $\mathrm{d} A<0$ ). The survival of colloids must therefore be a consequence of the kinetics of collapse: colloids are thermodynamically unstable but kinetically non-labile.

At first sight, even the kinetic argument seems to fail: colloidal particles attract each other over large distances, so there is a long-range force that tends to condense them into a single blob. The reasoning behind this remark is as follows. The energy of attraction between two individual atoms $i$ and $j$ separated by a distance $R_{i j}$, one in each colloidal particle, varies with their separation as $1 / R_{i j}^{6}$ (Topic 14B). The sum of all these pairwise interactions, however, decreases only as approximately $1 / R^{2}$ (the precise variation depending on the shape of the particles and their closeness), where $R$ is the separation of the centres of the particles. The change in the power from 6 to 2 stems from the fact that at short distances only a few molecules interact but at large distances many individual molecules are at about the same distance from one another, and contribute equally to the sum (Fig. 14E.1), so the total interaction does not fall off as fast as the single molecule-molecule interaction.\\
Several factors oppose the long-range dispersion attraction. For example, there may be a protective film at the surface of the colloid particles that stabilizes the interface and cannot be penetrated when two particles touch. Thus, the surface atoms of a platinum sol in water react chemically and are turned into $-\mathrm{Pt}(\mathrm{OH})_{3} \mathrm{H}_{3}$; this layer encases the particle like a shell. A fat can be emulsified by a soap because the long hydrocarbon tails penetrate the oil droplet but the carboxylate head groups (or other hydrophilic groups in synthetic detergents) surround the surface, form hydrogen bonds with water, and give rise to a shell of negative charge that repels a possible approach from another similarly charged particle.\\
\includegraphics[max width=\textwidth, center]{2025_02_27_d4b95d9718f0b9e2e6b9g-656}

Figure 14E. 1 Although the attraction between individual molecules is proportional to $1 / R^{6}$, more molecules are within range at large separations (pale region) than at small separation (dark region), so the total interaction energy declines more slowly and is proportional to a lower power of $1 / R$.

\subsection*{(c) The electrical double layer}
A major source of kinetic non-lability of colloids is the existence of an electric charge on the surfaces of the particles. Ions of opposite charge tend to cluster near each other, and form an ionic atmosphere around the particles, just as for individual ions (Topic 5F).

There are two regions of charge. First, there is a fairly immobile layer of ions that adhere tightly to the surface of the colloidal particle, and which may include water molecules (if that is the support medium). The radius of the sphere that captures this rigid layer is called the radius of shear and is the major factor determining the mobility of the particles. The electric potential at the radius of shear relative to its value in the distant, bulk medium is called the electrokinetic potential, $\zeta$ (or the zeta potential). Second, the charged unit attracts an oppositely charged atmosphere of mobile ions. The inner shell of charge and the outer ionic atmosphere constitute the electrical double layer.\\
The theory of the stability of lyophobic dispersions was developed by B. Derjaguin and L. Landau and independently by E. Verwey and J.T.G. Overbeek, and is known as the DLVO theory. It assumes that there is a balance between the repulsive interaction between the charges of the electrical double layers on neighbouring particles and the attractive interactions arising from van der Waals interactions between the molecules in the colloidal particles. The potential energy arising from the repulsion of double layers on particles of radius $a$ has the form


\begin{equation*}
V_{\text {repulsion }}=+\frac{A a^{2} \zeta^{2}}{R} \mathrm{e}^{-s / r_{\mathrm{D}}} \tag{14E.1}
\end{equation*}


where $A$ is a constant, $\zeta$ is the zeta potential, $R$ is the separation of centres, $s$ is the separation of the surfaces of the two particles ( $s=R-2 a$ for spherical particles of radius $a$ ), and $r_{D}$ is the thickness of the double layer. This expression is valid for\\
small particles with a thick double layer ( $r_{\mathrm{D}} \gg a$ ). When the double layer is thin $\left(r_{\mathrm{D}} \ll a\right)$, the expression is replaced by


\begin{equation*}
V_{\text {repulsion }}=+\frac{1}{2} B a^{2} \zeta^{2} \ln \left(1+\mathrm{e}^{-s / / r_{\mathrm{D}}}\right) \tag{14E.2}
\end{equation*}


where $B$ is another constant. In each case, the thickness of the double layer can be estimated from an expression like that derived for the thickness of the ionic atmosphere in the DebyeHückel theory (Topic 5F and A deeper look 1 on the website for this text) in which there is a competition between the assembling influences of the attraction between opposite charges and the disruptive effect of thermal motion:

\[
r_{\mathrm{D}}=\left(\frac{\varepsilon R T}{2 \rho F^{2} I b^{\ominus}}\right)^{1 / 2} \quad \begin{align*}
& \text { Thickness of the }  \tag{14E.3}\\
& \text { electrical double layer }
\end{align*}
\]

where $I$ is the ionic strength of the solution (eqn 5 F .28 , $I=\frac{1}{2} \sum_{i} z_{i}^{2} b_{i} / b^{\ominus}$ with $b^{\ominus}=1 \mathrm{molkg}^{-1}$ ) and $\rho$ its mass density. As usual, $F$ is Faraday's constant and $\varepsilon$ is the permittivity, $\varepsilon=\varepsilon_{\mathrm{r}} \varepsilon_{0}$. The potential energy arising from the attractive interaction has the form


\begin{equation*}
V_{\text {attraction }}=-\frac{C}{s} \tag{14E.4}
\end{equation*}


where $C$ is yet another constant. The variation of the total potential energy with separation is shown in Fig. 14E.2.

At high ionic strengths, the ionic atmosphere is dense and the potential shows a secondary minimum at large separations. Aggregation of the particles arising from the stabilizing effect of this secondary minimum is called flocculation. The flocculated material can often be redispersed by agitation because the well is so shallow. Coagulation, the irreversible aggregation of distinct particles into large particles, occurs when the separation of the particles is so small that they enter the primary minimum of the potential energy curve and van der Waals forces are dominant.\\
\includegraphics[max width=\textwidth, center]{2025_02_27_d4b95d9718f0b9e2e6b9g-657}

Figure 14E. 2 The variation of the potential energy of interaction with separation of the centres of the two particles and with the ratio of the particle size $a$ to the thickness of the electrical double layer, $r_{D}$. The regions labelled coagulation and flocculation show the dips in the potential energy curves where these processes occur.

The ionic strength is increased by the addition of ions, particularly those of high charge type, so such ions act as flocculating agents. This increase is the basis of the empirical Schulze-Hardy rule, that hydrophobic colloids are flocculated most efficiently by ions of opposite charge type and high charge number. The $\mathrm{Al}^{3+}$ ions in alum are very effective, and are used to induce the congealing of blood. When river water containing colloidal clay flows into the sea, the salt water induces flocculation and coagulation, and is a major cause of silting in estuaries.

Metal oxide sols tend to be positively charged whereas sulfur and the noble metals tend to be negatively charged. Naturally occurring macromolecules also acquire a charge when dispersed in water, and an important feature of proteins and other natural macromolecules is that their overall charge depends on the pH of the medium. For instance, in acidic environments protons attach to basic groups, and the net charge of the macromolecule is positive; in basic media the net charge is negative as a result of proton loss. At the isoelectric point the pH is such that there is no net charge on the macromolecule.

\subsection*{Example 14E. 1 Determining the isoelectric point of a protein}
The velocity with which the protein bovine serum albumin (BSA) moves through water under the influence of an electric field was monitored at several values of pH , and the data are listed below. What is the isoelectric point of the protein?

\begin{center}
\begin{tabular}{lrrrrrr}
pH & 4.20 & 4.56 & 5.20 & 5.65 & 6.30 & 7.00 \\
Velocity/( $\left.\mathrm{m} \mathrm{s} \mathrm{s}^{-1}\right)$ & 0.50 & 0.18 & -0.25 & -0.65 & -0.90 & -1.25 \\
\end{tabular}
\end{center}

Collect your thoughts Plot velocity against pH , then use interpolation to find the pH at which the velocity is zero, which is the pH at which the molecule has zero net charge.

The solution The data are plotted in Fig.14E.3. The velocity passes through zero at $\mathrm{pH}=4.8$; hence $\mathrm{pH}=4.8$ is the isoelectric point.\\
\includegraphics[max width=\textwidth, center]{2025_02_27_d4b95d9718f0b9e2e6b9g-657(1)}

Figure 14E. 3 The plot of the velocity of a moving macromolecule against pH allows the isoelectric point to be detected as the pH at which the velocity is zero. The data are from Example 14E.1.

Self-test 14E. 1 The following data were obtained for another protein:

\begin{center}
\begin{tabular}{lccccc}
pH & 3.5 & 4.5 & 5.0 & 5.5 & 6.0 \\
Velocity $/\left(\mu \mathrm{m} \mathrm{s}^{-1}\right)$ & 0.10 & -0.10 & -0.20 & -0.30 & -0.40 \\
\end{tabular}
\end{center}

Estimate the pH of the isoelectric point.

The primary role of the electrical double layer is to confer kinetic non-lability. Colliding colloidal particles break through the double layer and coalesce only if the collision is sufficiently energetic to disrupt the layers of ions and solvating molecules, or if thermal motion has stirred away the surface accumulation of charge. This disruption may occur at high temperatures, which is one reason why sols precipitate when they are heated.

\subsection*{14E. 2 Micelles and biological membranes}
In aqueous solutions surfactant molecules or ions can cluster together as micelles, which are colloid-sized clusters of molecules, for their hydrophobic tails tend to congregate, and their hydrophilic head groups provide protection (Fig. 14E.4).

\subsection*{(a) The hydrophobic interaction}
Consider a long-chained alcohol, such as pentan-1-ol $\left(\mathrm{CH}_{3} \mathrm{CH}_{2} \mathrm{CH}_{2} \mathrm{CH}_{2} \mathrm{CH}_{2} \mathrm{OH}\right)$. The hydrocarbon chain is hydrophobic and the -OH group is hydrophilic. A species with both\\
\includegraphics[max width=\textwidth, center]{2025_02_27_d4b95d9718f0b9e2e6b9g-658(2)}

Figure 14E. 4 A schematic version of a spherical micelle. The hydrophilic groups are represented by spheres and the hydrophobic hydrocarbon chains are represented by the stalks; these stalks are mobile.\\
\includegraphics[max width=\textwidth, center]{2025_02_27_d4b95d9718f0b9e2e6b9g-658(1)}

Figure 14E. 5 When a hydrocarbon molecule is surrounded by water, the $\mathrm{H}_{2} \mathrm{O}$ molecules form a cage. As a result of this acquisition of structure, the entropy of the water decreases, so the dispersal of the hydrocarbon into the water is accompanied by a local decrease in entropy. However, the aggregation of these individual caged hydrocarbon molecules into a micelle releases many of the caging water molecules back into the bulk and results in an increase in entropy.\\
hydrophobic and hydrophilic regions is called amphipathic. ${ }^{1}$ Amphipathic substances do dissolve slightly in water, and an understanding of the process gives insight into the formation of micelles and biological structures in general.\\
To understand the dissolution process in more detail, imagine a hypothetical initial state in which the alcohol is present in water as individual molecules. Each hydrophobic chain is surrounded by a cage of water molecules (Fig. 14E.5). This order reduces the entropy of the water below its 'pure' value. Now consider the final state, in which the hydrophobic chains have clustered together. Although the clustering contributes to the lowering of the entropy of the system, fewer (but larger) cages are required, and more water molecules are free to move. The net effect of the formation of clusters of hydrophobic chains is therefore a decrease in the organization of water molecules and therefore a net increase in entropy of the system. This increase in entropy of the solvent (water) means that the association of hydrophobic groups in an aqueous environment is spontaneous (provided there are no overwhelming enthalpy effects). This spontaneous clustering of hydrophobic groups in the presence of water gives the appearance of it being the outcome of an actual intermolecular force and is called the hydrophobic interaction.

Some insight into the processes involved can be obtained from studies of the thermodynamics of dissolving (as distinct from micelle formation). The entropy of dissolution of largely hydrophobic molecules in water is positive $\left(\Delta_{\text {diss }} S^{\ominus}>0\right)$ as the molecules disperse and the structure of the water changes to accommodate them. The process is commonly endothermic ( $\Delta_{\text {diss }} H^{\ominus}>0$ ), but the Gibbs energy of dissolution $\left(\Delta_{\text {diss }} G^{\ominus}\right)$ is typically negative, as is illustrated by the following data (at 298 K ):

\footnotetext{${ }^{1}$ The amphi- part of the name is from the Greek word for 'both', and the -pathic part is from the same root (meaning 'feeling') as sympathetic.
}\begin{center}
\begin{tabular}{|c|c|c|c|}
\hline
 & \( \begin{aligned} & \Delta_{\text {diss }} G^{\ominus} / \\ & \left(\mathrm{kJ} \mathrm{~mol}^{-1}\right) \end{aligned} \) & \( \begin{aligned} & \Delta_{\text {diss }} H^{\ominus} / \\ & \left(\mathrm{kJ} \mathrm{~mol}^{-1}\right) \end{aligned} \) & \( \begin{aligned} & \Delta_{\text {diss }} S^{\ominus} / \\ & \left(\mathrm{JK}^{-1} \mathrm{~mol}^{-1}\right) \end{aligned} \) \\
\hline
$\mathrm{CH}_{3} \mathrm{CH}_{2} \mathrm{CH}_{2} \mathrm{CH}_{2} \mathrm{OH}$ & -10 & +8 & +61 \\
\hline
$\mathrm{CH}_{3} \mathrm{CH}_{2} \mathrm{CH}_{2} \mathrm{CH}_{2} \mathrm{CH}_{2} \mathrm{OH}$ & -13 & +8 & +70 \\
\hline
\end{tabular}
\end{center}

In other words, the tendency to dissolve (at least to a small extent) is entropy-driven, with contributions from the dispersion of the solute molecules and the restructuring of the water. Once dissolved, further reorganization of the water occurs to drive the formation of micelles. The experimental values are consistent with a general rule that each additional $-\mathrm{CH}_{2}-$ group contributes a further $-3 \mathrm{~kJ} \mathrm{~mol}^{-1}$ to the standard Gibbs energy of dissolution.

A further aspect of this discussion is that it is possible to establish a scale of hydrophobicities. The hydrophobicity of a small molecular group R is reported by defining the hydrophobicity constant, $\pi$, as

$$
\pi=\log \frac{s(\mathrm{RX})}{s(\mathrm{HX})} \quad \begin{aligned}
& \text { Hydrophobicity constant } \\
& \text { [definition] }
\end{aligned}
$$

(14E.5)\\
where $s(\mathrm{RX})$ is the ratio of the molar solubility of the hydrophobic compound RX in the largely hydrocarbon solvent oc-tan-1-ol to that in water, and $s(\mathrm{HX})$ is the analogous ratio for the compound HX. A positive value of $\pi$ indicates that RX is more hydrophobic than RH.

It is found that the $\pi$ values of most compounds do not depend on the identity of X (which might be $\mathrm{OH}, \mathrm{NH}_{2}$, and so on). However, measurements suggest that $\pi$ increases by the same amount each time a $\mathrm{CH}_{2}$ group is added:

\begin{center}
\begin{tabular}{llllll}
\hline
-R & $-\mathrm{CH}_{3}$ & $-\mathrm{CH}_{2} \mathrm{CH}_{3}$ & $-\left(\mathrm{CH}_{2}\right)_{2} \mathrm{CH}_{3}$ & $-\left(\mathrm{CH}_{2}\right)_{3} \mathrm{CH}_{3}$ & $-\left(\mathrm{CH}_{2}\right)_{4} \mathrm{CH}_{3}$ \\
$\pi$ & 0.5 & 1.0 & 1.5 & 2.0 & 2.5 \\
\hline
\end{tabular}
\end{center}

It follows that acyclic saturated hydrocarbons become more hydrophobic as the carbon chain length increases. This trend can be rationalized by noting that $\Delta_{\text {diss }} G^{\ominus}$ becomes more negative as the number of carbon atoms in the chain increases, with the data on butan-1-ol and pentan-1-ol (see above) suggesting that the principal effect is due to the entropy.

\subsection*{(b) Micelle formation}
Micelles form only above the critical micelle concentration (CMC) and above the Krafft temperature. The CMC is detected by noting a pronounced change in physical properties of the solution, particularly the molar conductivity (Fig. 14E.6). There is no abrupt change in some properties at the CMC; rather, there is a transition region corresponding to a range of concentrations around the CMC where physical properties vary smoothly but nonlinearly with the concentra-\\
\includegraphics[max width=\textwidth, center]{2025_02_27_d4b95d9718f0b9e2e6b9g-659}

Figure 14E. 6 The typical variation of some physical properties of an aqueous solution of sodium dodecyl sulfate close to the critical micelle concentration (CMC).\\
tion. The hydrocarbon interior of a micelle is like a droplet of oil. Nuclear magnetic resonance shows that the hydrocarbon tails are mobile, but slightly more restricted than in the bulk. Micelles are important in industry and biology on account of their solubilizing function: matter can be transported by water after it has been dissolved in their hydrocarbon interiors. For this reason, micellar systems are used as detergents, for organic synthesis, froth flotation for the treatment of ores, and petroleum recovery.

The self-assembly of a micelle has the characteristics of a cooperative process in which the addition of a surfactant molecule to a cluster that is forming becomes more probable the larger the size of the aggregate, so after a slow start there is a cascade of formation of micelles. If it is supposed that the dominant micelle $\mathrm{M}_{N}$ consists of $N$ monomers M , then the dominant equilibrium to consider is


\begin{equation*}
N \mathrm{M} \rightleftharpoons \mathrm{M}_{N} \quad K=\frac{\left[\mathrm{M}_{N}\right] / c^{\ominus}}{\left([\mathrm{M}] / c^{\ominus}\right)^{N}} \tag{14E.6a}
\end{equation*}


where it has been assumed, probably dangerously on account of the large sizes of monomers, that the solution is ideal and that activities can be replaced by molar concentrations. The total concentration of surfactant, $[\mathrm{M}]_{\text {total }}$, is $[\mathrm{M}]+N\left[\mathrm{M}_{N}\right]$ because each micelle consists of $N$ monomer molecules. Therefore (omitting the $c^{\ominus}$ for clarity),


\begin{equation*}
K=\frac{\left[\mathrm{M}_{N}\right]}{\left([\mathrm{M}]_{\text {total }}-N\left[\mathrm{M}_{N}\right]\right)^{N}} \tag{14E.6b}
\end{equation*}


\subsection*{Brief illustration 14E. 1}
Equation 14E.6b can be solved numerically for the variation of the fraction of molecules present as micelles with the number of molecules present in a micelle and some results for $K=1$ are shown in Fig. 14E.7. For large $N$, there is a reasonably sharp transition in the fractions of surfactant molecules that are present in micelles, which corresponds to the existence of a CMC.\\
\includegraphics[max width=\textwidth, center]{2025_02_27_d4b95d9718f0b9e2e6b9g-660(1)}

Figure 14E. 7 The dependence of the fraction of surfactant molecules present as micelles on the number of molecules in the micelle for $K=1$.

Non-ionic surfactant molecules may cluster together in clumps of 1000 or more, but ionic species tend to be disrupted by the electrostatic repulsions between head groups and are normally limited to groups of less than about 100 . However, the disruptive effect depends more on the effective size of the head group than the charge. For example, ionic surfactants such as sodium dodecyl sulfate (SDS) and cetyltrimethylammonium bromide (CTAB) form rods at moderate concentrations whereas sugar surfactants form small, approximately spherical micelles. The micelle population commonly spans a wide range of particle sizes (i.e. it is polydisperse), and the shapes of the individual micelles vary with shape of the constituent surfactant molecules, surfactant concentration, and temperature. A useful predictor of the shape of the micelle is the surfactant parameter, $N_{s}$, defined as

$$
N_{\mathrm{s}}=\frac{V}{A l} \quad \begin{aligned}
& \text { Surfactant parameter } \\
& {[\text { definition] }}
\end{aligned}
$$

(14E.7)\\
where $V$ is the volume of the hydrophobic surfactant tail, $A$ is the area of the hydrophilic surfactant head group, and $l$ is the maximum length of the surfactant tail. Table 14E. 1 summarizes the dependence of aggregate structure on the surfactant parameter.

In aqueous solutions spherical micelles form, as shown in Fig. 14E.4, with the polar head groups of the surfactant mol-

Table 14E. 1 Micelle shape and the surfactant parameter

\begin{center}
\begin{tabular}{ll}
\hline
$\boldsymbol{N}_{\mathrm{s}}$ & Micelle shape \\
\hline
$<0.33$ & Spherical \\
$0.33-0.50$ & Cylindrical rods \\
$0.50-1.00$ & Vesicles \\
1.00 & Planar bilayers \\
$>1.00$ & Reverse micelles and other shapes \\
\hline
\end{tabular}
\end{center}

\begin{center}
\includegraphics[max width=\textwidth]{2025_02_27_d4b95d9718f0b9e2e6b9g-660}
\end{center}

Figure 14E. 8 The cross-sectional structure of a spherical liposome.\\
ecules on the micellar surface and interacting favourably with solvent and ions in solution. Hydrophobic interactions stabilize the aggregation of the hydrophobic surfactant tails in the micellar core. Under certain experimental conditions, a liposome may form, with an inward pointing inner surface of molecules surrounded by an outward pointing outer layer (Fig. 14E.8). Liposomes may be used to carry nonpolar drug molecules in blood.

Increasing the ionic strength of the aqueous solution reduces repulsions between surface head groups, and cylindrical micelles can form. These cylinders may stack together in reasonably close-packed (hexagonal) arrays, forming lyotropic mesomorphs and, more colloquially, 'liquid crystalline phases'.

Reverse micelles form in nonpolar solvents, with small polar surfactant head groups in a micellar core and more voluminous hydrophobic surfactant tails extending into the organic bulk phase. These spherical aggregates can solubilize water in organic solvents by creating a pool of trapped water molecules in the micellar core. As aggregates arrange at high surfactant concentrations to yield long-range positional order, many other types of structures are possible including cubic and hexagonal shapes.\\
As already noted, micelle formation is driven by hydrophobic interactions. The enthalpy of formation reflects contributions of interactions between micelle chains within the micelles and between the polar head groups and the surrounding medium. Consequently, enthalpies of micelle formation display no readily discernible pattern and may be positive (endothermic) or negative (exothermic). Many non-ionic micelles form endothermically, with $\Delta H$ of the order of 10 kJ per mole of surfactant molecules. That such micelles do form above the CMC indicates that the entropy change accompanying their formation must then be positive, and measurements suggest a value of about $+140 \mathrm{~J} \mathrm{~K}^{-1} \mathrm{~mol}^{-1}$ at room temperature.

\subsection*{(c) Bilayers, vesicles, and membranes}
Some micelles at concentrations well above the CMC form extended parallel sheets two molecules thick, called planar bilayers. The individual molecules lie perpendicular to the sheets, with hydrophilic groups on the outside in aqueous so-\\
lution and on the inside in nonpolar media. When segments of planar bilayers fold back on themselves, unilamellar vesicles may form where the spherical hydrophobic bilayer shell separates an inner aqueous compartment from the external aqueous environment.

Bilayers show a close resemblance to biological membranes, and are often a useful model on which to base investigations of biological structures. However, actual membranes are highly sophisticated structures. The basic structural element of a membrane is a phospholipid, such as phosphatidyl choline (1), which contains long hydrocarbon chains (typically in the range $\mathrm{C}_{14}-\mathrm{C}_{24}$ ) and a variety of polar groups, such as $-\mathrm{CH}_{2} \mathrm{CH}_{2} \mathrm{~N}\left(\mathrm{CH}_{3}\right)_{3}^{+}$. The hydrophobic chains stack together to form an extensive layer about 5 nm across. The lipid molecules form layers instead of micelles because the hydrocarbon chains are too bulky to allow packing into nearly spherical clusters.\\
\includegraphics[max width=\textwidth, center]{2025_02_27_d4b95d9718f0b9e2e6b9g-661}

The bilayer is a highly mobile structure. Not only are the hydrocarbon chains ceaselessly twisting and turning in the region between the polar groups, but the phospholipid molecules migrate over the surface. It is better to think of the membrane as a viscous fluid rather than a permanent structure, with a viscosity about 100 times that of water. Typically, a phospholipid molecule in a membrane migrates through about $1 \mu \mathrm{~m}$ in about 1 min .

All lipid bilayers undergo a transition from a state of high to low chain mobility at a temperature that depends on the structure of the lipid. To visualize the transition, consider what happens to a membrane as its temperature is lowered (Fig. 14E.9). There is sufficient energy available at normal temperatures for limited bond rotation to occur and the flexible chains writhe. However, the membrane is still highly organized in the sense that the bilayer structure does not come apart and the system is best described as a liquid crystal. At lower temperatures, the amplitudes of the writhing motion decrease until a specific temperature is reached at which motion is largely frozen. The membrane then exists as a gel. Biological membranes exist as liquid crystals at physiological temperatures.

Phase transitions in membranes are often observed as 'melting' from gel to liquid crystal by calorimetric methods. The\\
\includegraphics[max width=\textwidth, center]{2025_02_27_d4b95d9718f0b9e2e6b9g-661(1)}

Figure 14E. 9 A depiction of the variation with temperature of the flexibility of hydrocarbon chains in a lipid bilayer. (a) At physiological temperature, the bilayer exists as a liquid crystal, in which some order exists but the chains writhe. (b) At a specific temperature, the chains are largely frozen and the bilayer exists as a gel.\\
data show relations between the structure of the lipid and the melting temperature. Interspersed among the phospholipids of biological membranes are sterols, such as cholesterol (2), which is largely hydrophobic but does contain a hydrophilic -OH group. Sterols, which are present in different proportions in different types of cells, prevent the hydrophobic chains of lipids from 'freezing' into a gel and, by disrupting the packing of the chains, spread the melting point of the membrane over a range of temperatures.\\
\includegraphics[max width=\textwidth, center]{2025_02_27_d4b95d9718f0b9e2e6b9g-661(2)}

\subsection*{Brief illustration 14E. 2}
To predict trends in melting temperatures you need to assess the strengths of the interactions between molecules. Longer chains can be expected to be held together more strongly by hydrophobic interactions than shorter chains, so you should expect the melting temperature to increase with the length of the hydrophobic chain of the lipid. On the other hand, any structural elements that prevent alignment of the hydrophobic chains in the gel phase lead to low melting temperatures. Indeed, lipids containing unsaturated chains, those containing some $\mathrm{C}=\mathrm{C}$ bonds, form membranes with lower melting temperatures than those formed from lipids with fully saturated chains, those consisting of $\mathrm{C}-\mathrm{C}$ bonds only.

\subsection*{Checklist of concepts}
1. A disperse system is a dispersion of small particles of one material in another.2. Colloids are classified as lyophilic and lyophobic.3. A surfactant is a species that accumulates at the interface of two phases or substances.4. Many colloidal particles are thermodynamically unstable but kinetically non-labile.5. The radius of shear is the radius of the sphere that captures the rigid layer of charge attached to a colloid particle.6. The electrokinetic potential is the electric potential at the radius of shear relative to its value in the distant, bulk medium.7. The inner shell of charge and the outer atmosphere jointly constitute the electrical double layer.8. Flocculation is the reversible aggregation of colloidal particles.9. Coagulation is the irreversible aggregation of colloidal particles.$\square$ 10. The Schulze-Hardy rule states that hydrophobic colloids are flocculated most efficiently by ions of opposite charge type and high charge number.11. An amphipathic species has both hydrophobic and hydrophilic regions.\\
$\square$ 12. The hydrophobic interaction results in the clustering of nonpolar solutes in water.\\
$\square$ 13. A micelle is a colloid-sized cluster of molecules that forms at and above the critical micelle concentration and above the Krafft temperature.\\
$\square$ 14. Unilamellar vesicles are micelles that exist as extended parallel sheets.

\subsection*{Checklist of equations}
\begin{center}
\begin{tabular}{llll}
\hline
Property & Equation & Comment & Equation number \\
\hline
Thickness of the electrical double layer & $r_{\mathrm{D}}=\left(\varepsilon R T / 2 \rho F^{2} I b^{\ominus}\right)^{1 / 2}$ & Debye-Hückel theory & 14 E .3 \\
Hydrophobicity constant & $\pi=\log \{s(\mathrm{RX}) / s(\mathrm{HX})\}$ & Definition & 14 E .5 \\
Surfactant parameter & $N_{\mathrm{s}}=V / A l$ & Definition & 14 E .7 \\
\hline
\end{tabular}
\end{center}

\chapter{FOCUS 15 Solids}
This Focus explores the structures and physical properties of solids. The solid state includes most of the materials that make modern technology possible. It includes the wide varieties of steel that are used in architecture and engineering, the semiconductors and metallic conductors that are used in information technology and power distribution, the ceramics that increasingly are replacing metals, and the synthetic and natural polymers discussed in Focus 14 that are used in the textile industry and in the fabrication of many of the common objects of the modern world.

\subsection*{15A Crystal structure}
The characteristic feature of a crystal is the regular arrangement of its constituents. This Topic describes how that regularity is described in terms of the symmetry of the arrangement and then explains how the arrangements are described quantitatively.\\
15A. 1 Periodic crystal lattices; 15A. 2 The identification of lattice planes

\subsection*{15B Diffraction techniques}
Diffraction techniques enable the structures of solids to be determined in great detail. This Topic considers the basic principles of 'X-ray diffraction' and describes how the diffraction pattern can be interpreted in terms of the distribution of electron density. The diffraction of electrons and neutrons adds to the information that can be obtained about the structures of individual molecules and atoms in solids.\\
15B. 1 X-ray crystallography; 15B. 2 Neutron and electron diffraction

\subsection*{15C Bonding in solids}
The constituents of solids are held together by a variety of interactions, which impart characteristic properties. This Topic explores the interactions and prepares the ground for a discussion of the resulting properties.\\
15C. 1 Metals; 15C. 2 Ionic solids; 15C. 3 Covalent and molecular solids

\subsection*{15D The mechanical properties of solids}
The characteristic mechanical properties of a solid include various aspects of its rigidity. These properties are reported in terms of several parameters that can be related to the structure of the solid.

\subsection*{15E The electrical properties of solids}
One very important property of a solid is its ability to transport an electric current. This Topic explores how solids are classified according to their electrical conductivity and how the different behaviours seen can be rationalized using the 'band theory' of electronic structure. It goes on to show how the introduction of low concentrations of impurities can have a profound effect on the properties of semiconductors, and how this effect is exploited in making the semiconductor devices which are ubiquitous in modern electronics.\\
15E. 1 Metallic conductors; 15E. 2 Insulators and semiconductors;\\
15E. 3 Superconductors

\subsection*{15F The magnetic properties of solids}
The magnetic properties of solids are reported in terms of their 'susceptibility'. If the magnetic centres are independent, this property can be traced to individual electron spins. If the centres interact, properties such as ferromagnetism emerge.\\
15F. 1 Magnetic susceptibility; 15F. 2 Permanent and induced magnetic moments; 15F. 3 Magnetic properties of superconductors

\subsection*{15G The optical properties of solids}
Spectroscopy is a key tool for exploring the electronic structure of solids. As well as exploring the electronic band structure of a solid material, spectroscopic observations give insight into phenomena that arise from the interactions present in such materials. When subject to very intense radiation some\\
solids respond nonlinearly, leading to useful phenomena such as frequency doubling.\\
15G. 1 Excitons; 15G. 2 Metals and semiconductors; 15G. 3 Nonlinear optical phenomena

\subsection*{Web resources What is an application of this material?}
The deployment of X-ray diffraction techniques for the determination of the location of all the atoms in biological macromolecules has revolutionized the study of biochemistry and molecular biology. Impact 23 demonstrates the power of the techniques by exploring one of the most seminal X-ray images of all: the characteristic pattern obtained from strands of DNA and used in the construction of the double-helix model of DNA. Nanometre-sized assemblies that conduct electricity are currently of great technological interest, and Impact 24 describes their synthesis.

\section*{TOPIC 15A Crystal structure}
\subsection*{Why do you need to know this material?}
Crystalline solids are important in many technologies, and to be able to account for their mechanical, electrical, optical, and magnetic properties you need to understand their microscopic structures.

\subsection*{What is the key idea?}
The regular arrangement of the atoms in periodic crystals can be described in terms of unit cells.

What do you need to know already?\\
Light use is made of some of the language used to describe symmetry (Topic 10A).

The internal structure of a crystal is a regular array of its atoms, ions, or molecules. The features of this regular array, such as the details of the stacking pattern and its characteristic dimensions, are a crucial aspect of the link between the structure and properties of the solid.

\subsection*{15A. 1 Periodic crystal lattices}
A periodic crystal is built up from regularly repeating 'structural motifs', which may be atoms, molecules, or groups of atoms, molecules, or ions. A space lattice is the pattern formed by points representing the locations of these motifs (Fig. 15A.1). A space lattice is, in effect, an abstract scaffolding for the crystal structure. More formally, a space lattice is a three-dimensional, infinite array of points, each of which is surrounded in an identical way by its neighbours, and which defines the basic structure of the crystal. In some cases there may be a structural motif centred on each lattice point, but that is not necessary. The crystal structure itself is obtained by associating with each lattice point an identical structural motif. The solids known as quasicrystals are 'aperiodic', in the sense that the space lattice, though still filling space, does not have translational symmetry. This Topic deals only with periodic crystals.

A unit cell is an imaginary parallelepiped (parallel-sided figure) from which the entire space lattice can be constructed\\
\includegraphics[max width=\textwidth, center]{2025_02_27_d4b95d9718f0b9e2e6b9g-673(1)}

Figure 15A. 1 Each lattice point specifies the location of a structural motif (e.g. a molecule or a group of molecules). The space lattice is the entire array of lattice points; the crystal structure is the collection of structural motifs arranged according to the lattice.\\
by purely translational displacements (Fig. 15A.2). The cell is commonly formed by joining neighbouring lattice points by straight lines, and such unit cells are described as primitive (Fig. 15A.3). If each of the four points of a two-dimensional unit cell in Fig. 15A. 2 is regarded as shared with its four neighbours, then the cell has only one lattice point overall. The same definition applies in three dimensions, where each of the eight points of a primitive unit cell is shared by eight neighbours, giving one lattice point overall. It is often more convenient to draw larger non-primitive unit cells that also have lattice\\
\includegraphics[max width=\textwidth, center]{2025_02_27_d4b95d9718f0b9e2e6b9g-673}

Figure 15A. 2 A unit cell is a parallel-sided (but not necessarily rectangular) figure from which the entire space lattice can be constructed by using only translations (not reflections, rotations, or inversions). In the two-dimensional case shown here, each lattice point is shared by four neighbouring cells.\\
\includegraphics[max width=\textwidth, center]{2025_02_27_d4b95d9718f0b9e2e6b9g-674(4)}

Figure 15A.3 A primitive unit cell, an example of which is shown by the shaded volume, has lattice points only at its vertices. If each of the eight points is regarded as shared with its eight neighbours, the unit cell has only one lattice point overall.\\
points at their centres or on pairs of opposite faces. An infinite number of different unit cells can describe the same lattice, but the one with sides that have the shortest lengths and that are most nearly perpendicular to one another is normally chosen. The lengths of the sides of a unit cell are denoted $a$, $b$, and $c$, and the angles between them are denoted $\alpha, \beta$, and $\gamma$ (Fig. 15A.4).

Unit cells are classified into seven crystal systems by noting the rotational symmetry elements they possess:

\begin{itemize}
  \item A cubic unit cell has four threefold axes pointing to the corners of a tetrahedron, and passing through the centre of the cube (Fig. 15A.5).
  \item A monoclinic unit cell has one twofold axis (Fig. 15A.6).
  \item A triclinic unit cell has no rotational symmetry, and typically all three sides and angles are different (Fig. 15A.7).
\end{itemize}

Table 15A. 1 lists the essential symmetries, the elements that must be present for the unit cell to belong to a particular crystal system.\\
\includegraphics[max width=\textwidth, center]{2025_02_27_d4b95d9718f0b9e2e6b9g-674}

Figure 15A. 4 The notation for the sides and angles of a unit cell. Note that the angle $\alpha$ lies in the plane ( $b, c$ ).\\
\includegraphics[max width=\textwidth, center]{2025_02_27_d4b95d9718f0b9e2e6b9g-674(3)}

Figure 15A. 5 A unit cell belonging to the cubic system has four threefold axes, denoted $C_{3}$, arranged tetrahedrally. The insert shows the threefold symmetry.\\
\includegraphics[max width=\textwidth, center]{2025_02_27_d4b95d9718f0b9e2e6b9g-674(2)}

Figure 15A. 6 A unit cell belonging to the monoclinic system has a twofold axis (denoted $C_{2}$ and shown in more detail in the insert).\\
\includegraphics[max width=\textwidth, center]{2025_02_27_d4b95d9718f0b9e2e6b9g-674(1)}

Figure 15A. 7 A triclinic unit cell has no axes of rotational symmetry.

Table 15A. 1 The seven crystal systems*

\begin{center}
\begin{tabular}{ll}
\hline
System & Essential symmetries \\
\hline
Triclinic & None \\
Monoclinic & One $C_{2}$ axis \\
Orthorhombic & Three perpendicular $C_{2}$ axes \\
Rhombohedral & One $C_{3}$ axis \\
Tetragonal & One $C_{4}$ axis \\
Hexagonal & One $C_{6}$ axis \\
Cubic & Four $C_{3}$ axes in a tetrahedral arrangement \\
\hline
\end{tabular}
\end{center}

${ }^{*} C_{n}$ denotes an $n$-fold rotation, in which identical structures are obtained after rotation by $360^{\circ} / n$.

In three dimensions there are only 14 distinct space lattices. The unit cells of these Bravais lattices are illustrated in Fig. 15A.8. It is conventional to portray the lattices by primitive unit cells in some cases and by non-primitive unit cells in others. The following notation is used:

\begin{itemize}
  \item A primitive unit cell ( P ) has lattice points only at the corners.
  \item A body-centred unit cell (I) also has a lattice point at its centre.
  \item A face-centred unit cell (F) has lattice points at its corners and also at the centres of its six faces.
  \item A side-centred unit cell (A, B, or C) has lattice points at its corners and at the centres of two opposite faces.
\end{itemize}

For simple structures, it is often convenient to choose an atom belonging to the structural motif, or the centre of a molecule, as the location of a lattice point or the vertex of a unit cell, but that is not a necessary requirement. Symmetry-related lattice points within the unit cell of a Bravais lattice have identical surroundings.\\
\includegraphics[max width=\textwidth, center]{2025_02_27_d4b95d9718f0b9e2e6b9g-675(2)}

Orthorhombic P Orthorhombic C Orthorhombic I Orthorhombic F\\
\includegraphics[max width=\textwidth, center]{2025_02_27_d4b95d9718f0b9e2e6b9g-675}

Figure 15A. 8 The 14 Bravais lattices. The points are lattice points, and are not necessarily occupied by atoms. P denotes a primitive unit cell ( $R$ is used for a trigonal lattice), I a body-centred unit cell, $F$ a face-centred unit cell, and $C$ (or A or B) a cell with lattice points on two opposite faces. Trigonal lattices may belong to the rhombohedral or hexagonal systems (Table 15A.1).

\subsection*{Brief illustration 15A. 1}
The two-dimensional lattice shown in Fig. 15A. 9 consists of a rectangular array of lattice points, with one additional point at the centre of each rectangle; a (non-primitive) unit cell is indicated. This cell has twofold axes of symmetry passing through the mid-points of opposite sides of the rectangle. Rotations about these axes interchange the lattice points at the corners of the rectangle, but the lattice point at the centre is not affected. It therefore follows that the lattice points at the corners are equivalent, but the lattice point at the centre is distinct.\\
\includegraphics[max width=\textwidth, center]{2025_02_27_d4b95d9718f0b9e2e6b9g-675(1)}

Figure 15A. 9 The two-dimensional lattice used in Brief illustration 15A.1; a (non-primitive) unit cell is indicated by the shaded area. The lattice points at the corners of the unit cell are related by rotations about the twofold axes of symmetry shown; the lattice point at the centre is unaffected by these operations.

\subsection*{15A. 2 The identification of lattice planes}
The interpretation of the diffraction techniques that are used to measure the size of unit cells and the arrangement of molecules within them makes use of the orientation and separation of planes that pass through the crystal (Topic 15B). Two-dimensional lattices are easier to visualize than threedimensional lattices, so in this discussion the concepts involved in identifying lattice planes are introduced for two dimensions initially, and then the results are extended by analogy to three dimensions. Note that lattice planes do not necessarily pass through lattice points.

\subsection*{(a) The Miller indices}
Consider a two-dimensional rectangular lattice formed from a unit cell of sides $a$ and $b$ (Fig. 15A.10). Each panel in the illustration shows a set of evenly spaced planes that can be identified by considering, for each set, the plane lying closest to the origin (but not passing through it) and then quoting the distances at which this plane intersects the $a$ and $b$ axes. These distances are: (a) $(1 a, 1 b)$, (b) $\left(\frac{1}{2} a, \frac{1}{3} b\right)$, (c) $(-1 a, 1 b)$, and (d) $(\infty a, 1 b)$, with $\infty$ indicating that the plane is parallel to an axis and intersects it (notionally) at infinity. If it is agreed to quote distances along the axes as multiples of the\\
\includegraphics[max width=\textwidth, center]{2025_02_27_d4b95d9718f0b9e2e6b9g-676(2)}

Figure 15A. 10 Some of the sets of equally-spaced planes that can be drawn in a rectangular space lattice; the origin is indicated by the purple lattice point. The Miller indices $\{h k /\}$ of each set of planes are: (a) $\{110\}$, (b) $\{230\}$, (c) $\{\overline{1} 10\}$, and (d) $\{010\}$.\\
corresponding dimensions of the unit cell, then these intersections can be expressed more simply as $(1,1),\left(\frac{1}{2}, \frac{1}{3}\right),(-1,1)$, and $(\infty, 1)$, respectively. If the lattice in Fig. 15A. 10 is the top view of a three-dimensional orthorhombic lattice, all four sets of planes intersect the $c$ axis at infinity. Therefore, in the three-dimensional case the labels are $(1,1, \infty),\left(\frac{1}{2}, \frac{1}{3}, \infty\right)$, $(-1,1, \infty)$, and $(\infty, 1, \infty)$.\\
The inconvenience of fractions and infinity in these labels can be avoided by specifying a plane by using its Miller indices, $(h k l)$, where $h, k$, and $l$ are the reciprocals of the intersection distances along the $a, b$, and $c$ axes, respectively. For example, the plane $\left(\frac{1}{2}, \frac{1}{3}, \infty\right)$ has Miller indices (230). As will be seen, this notation brings with it additional advantages. The Miller notation has the following features:

\begin{itemize}
  \item Negative indices are written with a bar over the number, as in ( $\overline{1} 10$ ).
  \item The notation ( $h k l$ ) denotes an individual plane. A set of parallel planes with identical spacing, is denoted $\{h k l\}$.
\end{itemize}

For example,

\begin{center}
\begin{tabular}{lllll}
\hline
Intersect axes at & $(a, b, \infty c)$ & $\left(\frac{1}{2} a, \frac{1}{3} b, \infty c\right)$ & $(-a, b, \infty c)$ & $(\infty a, b, \infty c)$ \\
Remove cell dimensions & $(1,1, \infty)$ & $\left(\frac{1}{2}, \frac{1}{3}, \infty\right)$ & $(-1,1, \infty)$ & $(\infty, 1, \infty)$ \\
Take reciprocals & $(1,1,0)$ & $(2,3,0)$ & $(-1,1,0)$ & $(0,1,0)$ \\
Express as indices & $(110)$ & $(230)$ & $(\overline{1} 10)$ & $(010)$ \\
Sets of parallel planes & $\{110\}$ & $\{230\}$ & $\{\overline{1} 10\}$ & $\{010\}$ \\
\hline
\end{tabular}
\end{center}

A helpful feature to remember is that the smaller the absolute value of $h$ in $\{h k l\}$, the more nearly parallel the set of planes is to the $a$ axis (the $\{h 00\}$ planes are an exception). The same is true of $k$ and the $b$ axis, and $l$ and the $c$ axis. When $h=0$, the planes intersect the $a$ axis at infinity, so the $\{0 \mathrm{kl}\}$ planes are parallel to the $a$ axis. Similarly, the $\{h 0 l\}$ planes are parallel to the $b$ axis and the $\{h k 0\}$ planes are parallel to the $c$ axis.\\
\includegraphics[max width=\textwidth, center]{2025_02_27_d4b95d9718f0b9e2e6b9g-676(1)}

Figure 15A. 11 Some representative planes in three dimensions and their Miller indices; the origin is indicated by the purple lattice point. Note that the index 0 indicates that a plane is parallel to the corresponding axis, and that the indexing may also be used for unit cells with non-orthogonal axes.

Figure 15A. 11 shows a three-dimensional representation of a selection of planes, including one in a lattice with nonorthogonal axes.

\subsection*{(b) The separation of neighbouring planes}
The Miller indices are very useful for expressing the separation of neighbouring planes.

\subsection*{How is that done? 15A. 1}
Deriving an expression for the separation of planes

Consider the $\{h k 0\}$ planes of a square lattice built from a unit cell with sides of length $a$ (Fig. 15A.12). The separation between the lattice planes is equal to the perpendicular distance from the ( $h k 0$ ) plane to the origin. Expressions for the sine and cosine of the angle $\phi$ are found by considering the sides of the two right-angle triangles shown in the figure

$$
\sin \phi=\frac{d_{h k 0}}{(a / h)}=\frac{h d_{h k 0}}{a} \quad \cos \phi=\frac{d_{h k 0}}{(a / k)}=\frac{k d_{h k 0}}{a}
$$

The length of the hypotenuse of the lower triangle is $a / h$ because a Miller index $h$ indicates that the plane intersects the $a$\\
\includegraphics[max width=\textwidth, center]{2025_02_27_d4b95d9718f0b9e2e6b9g-676}

Figure 15A. 12 The construction used to find the spacing of the ( $h k 0$ ) plane in a square unit cell.\\
axis at a distance $a / h$ from the origin. Likewise, the hypotenuse of the upper triangle is $a / k$. Then, because $\sin ^{2} \phi+\cos ^{2} \phi=1$, it follows that

$$
\left(\frac{h d_{h k 0}}{a}\right)^{2}+\left(\frac{k d_{h k 0}}{a}\right)^{2}=1
$$

which can be rearranged by dividing both sides by $d_{h k 0}^{2}$ into

$$
\frac{1}{d_{h k 0}^{2}}=\frac{h^{2}+k^{2}}{a^{2}} \text { or } d_{h k 0}=\frac{a}{\left(h^{2}+k^{2}\right)^{1 / 2}}
$$

By extension to three dimensions, the separation of the $\{h k l\}$ planes, $d_{h k l}$, of a cubic lattice is given by\\
$-\left|\begin{array}{l}\frac{1}{d_{h k l}^{2}}=\frac{h^{2}+k^{2}+l^{2}}{a^{2}} \\ d_{h k l}=\frac{a}{\left(h^{2}+k^{2}+l^{2}\right)^{1 / 2}}\end{array}\right| \begin{aligned} & \text { Separation of planes } \\ & \text { [cubic lattice] }\end{aligned}$\\
The corresponding expression for a general orthorhombic lattice (one in which the axes are mutually perpendicular, but not equal in length) is the generalization of this expression:

\[
\frac{1}{d_{h k l}^{2}}=\frac{h^{2}}{a^{2}}+\frac{k^{2}}{b^{2}}+\frac{l^{2}}{c^{2}} \quad \begin{align*}
& \text { Separation of planes }  \tag{15A.1b}\\
& \text { [orthorhombic lattice] }
\end{align*}
\]

\subsection*{Example 15A. 1 Using the Miller indices}
Calculate the separation of (a) the $\{123\}$ planes and (b) the $\{246\}$ planes of an orthorhombic unit cell with $a=0.82 \mathrm{~nm}$, $b=0.94 \mathrm{~nm}$, and $c=0.75 \mathrm{~nm}$.

Collect your thoughts For the first part, all you need to do is substitute the given values into eqn 15A.1b. You could do the same for part (b), but note that the Miller indices for the second set of planes are just twice those of the first part. By referring to eqn 15A.1b you can see that multiplying the values of $h, k$, and $l$ by $n$ gives the following expression for the spacing of the $\{n h n k n l\}$ planes

$$
\frac{1}{d_{n h, n k, n l}^{2}}=\frac{(n h)^{2}}{a^{2}}+\frac{(n k)^{2}}{b^{2}}+\frac{(n l)^{2}}{c^{2}}=n^{2} \overbrace{\left(\frac{h^{2}}{a^{2}}+\frac{k^{2}}{b^{2}}+\frac{l^{2}}{c^{2}}\right)}^{=1 / d_{h k l}^{2}}=\frac{n^{2}}{d_{h k l}^{2}}
$$

which implies that

$$
d_{n h, n k, n l}=\frac{d_{h k l}}{n}
$$

The solution Substituting the indices into eqn 15A.1b gives

$$
\frac{1}{d_{123}^{2}}=\frac{1^{2}}{(0.82 \mathrm{~nm})^{2}}+\frac{2^{2}}{(0.94 \mathrm{~nm})^{2}}+\frac{3^{2}}{(0.75 \mathrm{~nm})^{2}}=22.0 \ldots \mathrm{~nm}^{-2}
$$

Hence, $d_{123}=0.21 \mathrm{~nm}$. It then follows immediately that $d_{246}$ is one-half this value, or 0.11 nm .

Self-test 15A. 1 Calculate the separation of (a) the $\{133\}$ planes and (b) the $\{399\}$ planes in the same lattice.\\
\includegraphics[max width=\textwidth, center]{2025_02_27_d4b95d9718f0b9e2e6b9g-677}

\subsection*{Checklist of concepts}
1. A periodic crystal is built up from regularly repeating structural motifs.\begin{enumerate}
  \setcounter{enumi}{1}
  \item A space lattice is the pattern formed by points (lattice points) representing the locations of structural motifs (atoms, molecules, or groups of atoms, molecules, or ions).3. A unit cell is an imaginary parallel-sided figure from which the entire space lattice can be constructed by purely translational displacements.4. A primitive unit cell has lattice points only at its vertices and only one lattice point overall; non-primitive unit cells also have lattice points at their centres or on pairs of opposite faces.
  \item Unit cells are classified into seven crystal systems according to their rotational symmetries: unit cells are classified as cubic, monoclinic, or triclinic according to the essential symmetries they possess.6. The Bravais lattices are the 14 distinct space lattices in three dimensions (Fig. 15A.8).7. The unit cells of the Bravais lattices are classed as primitive ( P ), body-centred ( I ), face-centred ( F ), and side-centred ( $A, B$, or $C$ ).8. A lattice plane is specified by a set of Miller indices $(h k l)$; sets of planes are denoted $\{h k l\}$.
\end{enumerate}

\subsection*{Checklist of equations}
\begin{center}
\begin{tabular}{llll}
\hline
Property & Equation & Comment & Equation number \\
\hline
Separation of planes in a cubic lattice & $1 / d_{h k l}^{2}=\left(h^{2}+k^{2}+l^{2}\right) / a^{2}$ & $h, k$, and $l$ are Miller indices & 15A.1a \\
Separation of planes in an orthorhombic lattice & $1 / d_{h k l}^{2}=h^{2} / a^{2}+k^{2} / b^{2}+l^{2} / c^{2}$ &  & 15 A .1 b \\
\hline
\end{tabular}
\end{center}

\section*{TOPIC 15B Diffraction techniques}
\subsection*{Why do you need to know this material?}
To account for the properties of solids it is necessary to understand their detailed structures and how they are determined by a variety of diffraction techniques.

\subsection*{What is the key idea?}
The regular arrangement of the atoms in periodic crystals can be determined by techniques based on diffraction.

\subsection*{What do you need to know already?}
You need to be familiar with the description of crystal structures and the use of Miller indices to identify lattice planes (Topic 15A). You also need to be familiar with the wave description of electromagnetic radiation (The chemist's toolkit 13 in Topic 7A), and the basic properties of the Fourier transform (The chemist's toolkit 28 in Topic 12C). Light use is made of the de Broglie relation (Topic 7 A ) and the equipartition theorem (The chemist's toolkit 7 in Topic 2A).

Diffraction techniques can be used to determine the details of the arrangement of ions, atoms, and molecules in a crystalline solid to high precision. Such techniques are now so well developed that both the collection of the diffraction data and its interpretation in terms of a structure are automated to a high degree.

\subsection*{15B.1 X-ray crystallography}
As explained in The chemist's toolkit 13 (in Topic 7A), a characteristic property of waves is that when they are present in the same region of space they interfere with one another. A greater displacement is obtained where peaks or troughs of the waves coincide and a smaller displacement where peaks coincide with troughs. Diffraction is the interference caused by an object in the path of waves; it occurs when the dimensions of the diffracting object are comparable to the wavelength of the radiation.

\subsection*{(a) X-ray diffraction}
X-rays diffract when passed through a crystal because their wavelengths are comparable to the separation of lattice planes.

Consequently, X-ray diffraction is a very powerful technique for structural studies of solid materials. The actual process of going from the observed diffraction pattern to a structure is rather involved, but such is the degree of integration of computers into the experimental apparatus that the technique is almost fully automated, even for large molecules and complex solids. The analysis is aided by molecular modelling techniques, which can guide the investigation towards a plausible structure.\\
X-rays are electromagnetic radiation with wavelengths of the order of $10^{-10} \mathrm{~m}$. They are typically generated by bombarding a metal with high-energy electrons (Fig. 15B.1). The electrons decelerate as they plunge into the metal and generate radiation with a continuous range of wavelengths called Bremsstrahlung (Bremse is German for deceleration, Strahlung for ray). Superimposed on the continuum are a few high-intensity, sharp peaks (Fig. 15B.2). These peaks arise from collisions of the incoming electrons with the electrons in the inner shells of the atoms. The collision expels an electron from an inner shell, and an electron of higher energy drops into the vacancy, emitting the excess energy as an X-ray photon (Fig. 15B.3). If the electron falls into a K shell (a shell with $n=1$ ), the X-rays are classified as 'K-radiation', and similarly for transitions into the $\mathrm{L}(n=2)$ and $\mathrm{M}(n=3)$ shells. Strong, distinct lines are labelled $K_{\alpha}, K_{\beta}$, and so on. Synchrotrons (Topic 11A) generate high-intensity X-ray radiation which is increasingly used in diffraction experiments on account of the resulting greater intensity in the diffraction pattern, and hence higher sensitivity.

An early method of observing diffraction consisted of passing a beam containing X-rays with a range of wavelengths into\\
\includegraphics[max width=\textwidth, center]{2025_02_27_d4b95d9718f0b9e2e6b9g-678}

Figure 15B. 1 X-rays are generated by directing an electron beam on to a cooled metal target. Beryllium is transparent to X-rays (on account of the small number of electrons in each atom) and is used for the windows.\\
\includegraphics[max width=\textwidth, center]{2025_02_27_d4b95d9718f0b9e2e6b9g-679(2)}

Figure 15B.2 The X-ray emission from a metal consists of a broad, featureless Bremsstrahlung background, with sharp peaks superimposed on it. The label K indicates that the radiation comes from a transition in which an electron falls into a vacancy in the K shell of the atom.\\
\includegraphics[max width=\textwidth, center]{2025_02_27_d4b95d9718f0b9e2e6b9g-679(1)}

Figure 15B. 3 The processes that contribute to the generation of X -rays. An incoming electron collides with an electron (in the K shell), and ejects it. Another electron (from the L shell in this illustration) falls into the vacancy and emits its excess energy as an X-ray photon.\\
a single crystal, and recording the diffraction pattern photographically. The idea behind this approach is that a crystal might not be suitably orientated to cause diffraction for a single wavelength but, whatever its orientation, diffraction would be achieved for at least one of the wavelengths present in the beam. There is currently a resurgence of interest in this approach because synchrotron radiation spans a range of X-ray wavelengths.

A more common technique uses monochromatic radiation and a powdered sample, which consists of many tiny crystallites, oriented at random. At least some of the crystallites will be appropriately orientated to give diffraction. In modern 'powder diffractometers' the intensities of the reflections are monitored electronically as the detector is rotated around the sample in a plane containing the incident ray. Powder diffraction techniques are used to identify the composition of a sample of a solid substance by comparison of the positions of the diffraction lines and their intensities with previously recorded\\
\includegraphics[max width=\textwidth, center]{2025_02_27_d4b95d9718f0b9e2e6b9g-679(3)}

Figure 15B.4 X-ray powder diffraction patterns of two polymorphs of $\mathrm{CaCO}_{3}$, calcite and aragonite. The patterns are distinctive and can be used to identify the polymorph present in an unknown sample.\\
patterns of known structures (Fig. 15B.4); large databases of such information are available. This approach can be used to determine the composition of mixed phases, and hence to construct a phase diagram. The technique is also used for the initial determination of the dimensions and symmetries of unit cells.

The method developed by the Braggs (William and his son Lawrence) is the foundation of almost all modern work in X-ray crystallography. They used a single crystal and a monochromatic beam of X-rays, and rotated the crystal until a reflection was detected. There are many different sets of planes in a crystal, so there are many angles at which a reflection occurs. The complete set of data consists of the list of angles at which reflections are observed and their intensities.

Single-crystal diffraction patterns are measured by using a 'four-circle diffractometer' (Fig. 15B.5). Once the dimensions and symmetry of the unit cell of the crystal in question have been identified, the angular setting of the detector on the four circles is adjusted so that the precise position and intensity of\\
\includegraphics[max width=\textwidth, center]{2025_02_27_d4b95d9718f0b9e2e6b9g-679}

Figure 15B.5 A four-circle diffractometer. The settings of the orientations ( $\phi, \chi, \theta$, and $\Omega$ ) of the components are controlled by computer; each ( $h k l$ ) reflection is monitored in turn, and their intensities are recorded.\\
each peak in the diffraction pattern can be measured. Modern instruments use area detectors and image plates, which sample whole regions of diffraction patterns simultaneously, rather than peak by peak, thus increasing the speed with which data is collected.

\subsection*{(b) Bragg's law}
An early approach to the analysis of diffraction patterns produced by crystals was to regard a lattice plane as a semi-transparent mirror and to model a crystal as a stack of reflecting lattice planes of separation $d$ (Fig. 15B.6). The model makes it easy to calculate the angle the crystal must make to the incoming beam of X-rays for constructive interference to occur. It has also given rise to the name reflection to denote an intense beam arising from constructive interference.

Consider the reflection of two parallel rays of the same wavelength and phase by two adjacent planes of a lattice, as shown in Fig. 15B.6. One ray strikes point $D$ on the upper plane but the other ray must travel an additional distance $A B$ before striking the plane immediately below. The reflected rays also differ in path length by the distance BC . As is evident from the inset in Fig. 15B.6, both the lengths $A B$ and $B C$ are $d \sin \theta$; the total path length difference of the two rays is then

$$
\mathrm{AB}+\mathrm{BC}=2 d \sin \theta
$$

where $2 \theta$ is the glancing angle ( $2 \theta$ rather than $\theta$, because the beam is deflected through $2 \theta$ from its initial direction). For many glancing angles the path-length difference is not an integer number of wavelengths, and the waves interfere largely destructively. However, when the path-length difference is an integer number of wavelengths $(\mathrm{AB}+\mathrm{BC}=n \lambda)$, the reflected waves are in phase and interfere constructively. It follows that a reflection should be observed when $\theta$ satisfies Bragg's law:

$$
n \lambda=2 d \sin \theta \quad \text { Bragg's law }
$$

(15B.1a)\\
\includegraphics[max width=\textwidth, center]{2025_02_27_d4b95d9718f0b9e2e6b9g-680}

Figure 15B. 6 The conventional derivation of Bragg's law treats each lattice plane as a plane reflecting the incident radiation. The path lengths for reflection from adjacent planes differ by AB + $B C$, which depends on the angle $\theta$. Constructive interference (a 'reflection') occurs when $A B+B C$ is equal to an integer number of wavelengths.

Reflections with $n=2,3, \ldots$ are called second order, third order, and so on; they correspond to path-length differences of 2 , $3, \ldots$ wavelengths. In modern work it is normal to absorb the $n$ into $d$, to write Bragg's law as

\[
\begin{array}{ll}
\lambda=2 d \sin \theta & \begin{array}{l}
\text { Bragg's law } \\
\text { [alternative form] }
\end{array} \tag{15B.1b}
\end{array}
\]

and to regard the $n$ th-order reflection as arising from the $\{n h n k n l\}$ planes. As discussed in Example 15A. 1 in Topic 15A, the spacing of the $\{n h n k n l\}$ planes is $d_{h k l} / n$, where $d_{h k l}$ is the spacing of the $\{h k l\}$ planes. The primary use of Bragg's law is in the determination of the spacing between the layers in the lattice because $d$ may readily be calculated from a measured value of the angle $\theta$.

\subsection*{Brief illustration 15B. 1}
A first-order reflection from the $\{111\}$ planes of a cubic crystal was observed at a glancing angle, $2 \theta$, of $22.4^{\circ}$ when X-rays of wavelength 154 pm were used. According to eqn 15B.1b, the $\{111\}$ planes responsible for the diffraction have separation $d_{111}=\lambda /(2 \sin \theta)$, therefore

$$
d_{111}=\frac{\lambda}{2 \sin \theta}=\frac{154 \mathrm{pm}}{2 \sin 11.2^{\circ}}=396 \mathrm{pm}
$$

The separation of the $\{111\}$ planes of a cubic lattice of side $a$ is given by eqn 15A.1a in Topic 15 A as $d_{111}=a / 3^{1 / 2}$. It therefore follows that $a=3^{1 / 2} d_{111}=687 \mathrm{pm}$.

Some types of unit cell give characteristic patterns of lines. In a cubic lattice with dimension $a$, the spacing of the $\{h k l\}$ planes, $d_{h k l}$, is given by eqn 15A.1a in Topic 15A ( $d_{h k l}=$ $\left.a /\left(h^{2}+k^{2}+l^{2}\right)^{1 / 2}\right)$; the angles at which the $\{h k l\}$ planes give first-order reflections are given by

$$
\sin \theta=\frac{\lambda}{2 d_{h k l}}=\frac{d_{h k l}=a /\left(h^{2}+k^{2}+l\right)^{1 / 2}}{\left(h^{2}+k^{2}+l^{2}\right)^{1 / 2} \frac{\lambda}{2 a}}
$$

Not all integral values of $h^{2}+k^{2}+l^{2}$ are obtained when integer values of the indices are substituted:

\begin{center}
\begin{tabular}{llllll}
\hline
$\{h k l\}$ & $\{100\}$ & $\{110\}$ & $\{111\}$ & $\{200\}$ & $\{210\}$ \\
$h^{2}+k^{2}+l^{2}$ & 1 & 2 & 3 & 4 & 5 \\
$\{h k l\}$ & $\{211\}$ & $\{220\}$ & $\{300\}$ & $\{221\}$ & $\{310\}$ \\
$h^{2}+k^{2}+l^{2}$ & 6 & 8 & 9 & 9 & 10 \\
\hline
\end{tabular}
\end{center}

Notice that $h^{2}+k^{2}+l^{2}=7$ (and $15, \ldots$ ) does not appear. As a result, in the pattern of diffraction lines there is a larger gap between the $\{211\}$ and $\{220\}$ reflections than between nearby lines, and likewise between $\{321\}$ (for which $h^{2}+k^{2}+l^{2}=14$ ) and $\{400\}$ (for which $h^{2}+k^{2}+l^{2}=16$ ). The absence of these lines leads to a characteristic pattern which helps in identifying the type of unit cell.

\subsection*{(c) Scattering factors}
X-ray scattering is caused by the oscillations an incoming electromagnetic wave generates in the electrons of atoms. Heavy, electron-rich atoms give rise to stronger scattering than light atoms. This dependence on the number of electrons is expressed in terms of the scattering factor, $f$, of the element. If the scattering factor is large, then the atoms scatter X-rays strongly. It turns out that, in a spherically symmetrical atom, the scattering factor for the atom is related to the electron density, $\rho(r)$, and the glancing angle, $2 \theta$, by

\[
f(\theta)=4 \pi \int_{0}^{\infty} \rho(r) \frac{\sin k r}{k r} r^{2} \mathrm{~d} r \quad k=\frac{4 \pi}{\lambda} \sin \theta \quad \begin{align*}
& \text { Scattering }  \tag{15B.2}\\
& \text { factor }
\end{align*}
\]

The scattering factor is greatest in the forward direction $(\theta=0$, Fig. 15B.7), and it can be shown that in that direction it is equal to the total number of electrons in the atom, $N_{e}$.\\
\includegraphics[max width=\textwidth, center]{2025_02_27_d4b95d9718f0b9e2e6b9g-681}

Figure 15B.7 The variation of the scattering factor of atoms and ions with atomic number and angle. The scattering factor in the forward direction (at $\theta=0$, and hence at $(\sin \theta) / \lambda=0$ ) is equal to the number of electrons present in the species.

\subsection*{How is that done? 15B. 1 Evaluating the scattering factor}
 in the forward directionThe first step is to note that because $\sin k r$ cannot exceed 1 , the maximum value of $(\sin k r) / k r$ occurs as $k \rightarrow 0$, which corresponds to $\sin \theta \rightarrow 0$ and therefore $\theta \rightarrow 0$. Therefore, the scattering factor has its maximum value in the forward direction.

To evaluate the scattering factor in the forward direction you need to take the limit $k \rightarrow 0$, and therefore $k r \rightarrow 0$ in eqn 15B.2. Therefore, use $\sin x=x-\frac{1}{6} x^{3}+\cdots$ and write

$$
\lim _{k r \rightarrow 0} \frac{\sin k r}{k r}=\lim _{k r \rightarrow 0} \frac{k r-\frac{1}{6}(k r)^{3}+\cdots}{k r}=\lim _{k r \rightarrow 0}\left\{1-\frac{1}{6}(k r)^{2}+\cdots\right\}=1
$$

In this limit eqn 15B. 2 simplifies to

$$
\begin{aligned}
f(0) & =\lim _{k r \rightarrow 0} 4 \pi \int_{0}^{\infty} \rho(r) \overbrace{\frac{\sin k r}{k r}}^{\rightarrow 1} r^{2} \mathrm{~d} r \\
& =4 \pi \int_{0}^{\infty} \rho(r) r^{2} \mathrm{~d} r=\int_{0}^{\infty} \rho(r) \overbrace{4 \pi r^{2} \mathrm{~d} r}^{\text {volume element }}
\end{aligned}
$$

The factor $4 \pi r^{2} \mathrm{~d} r$ is the volume of a spherical shell of radius $r$ and thickness $\mathrm{d} r$. The total number of electrons in this shell is the electron density at $r$ multiplied by the volume of the shell, $4 \pi r^{2} \rho(r) \mathrm{d} r$, and this number summed over shells of all radii is the total number of electrons in the atom. Hence, in the forward direction, $f=N_{\mathrm{e}}$. For example, the scattering factors of $\mathrm{Na}^{+}, \mathrm{K}^{+}$, and $\mathrm{Cl}^{-}$are 8,18 , and 18 , respectively.

\subsection*{(d) The electron density}
The structure factor, $F_{h k l}$, is the net amplitude of a given $\{h k l\}$ reflection that takes into account the positions and types of all the atoms in the unit cell. It can be expressed in terms of the locations and scattering factors of the atoms.

How is that done? 15B. 2 Relating the structure factor to the location of the atoms and their scattering factors

Suppose that a unit cell contains several atoms with scattering factors $f_{j}$ and coordinates $\left(x_{j} a, y_{j} b, z_{j} c\right)$, where $x_{j}$ is the coordinate of the atom $j$ in the $a$ direction, expressed as a fraction of the length $a$, and likewise for the other coordinates.

Step 1 Consider the (h00) reflection with $h=1$\\
The reflection shown in Fig. 15B. 8 corresponds to two waves from adjacent A planes; the phase difference of the waves is $2 \pi$. If there is a $B$ atom at a fraction $x$ of the distance between the two A planes, then it gives rise to a wave with a phase difference $2 \pi x$ relative to an A reflection. To see this conclusion, note that, if $x=0$, there is no phase difference; if $x=\frac{1}{2}$ the phase difference is $\pi$; if $x=1$, the B atom lies where the lower A atom is and the phase difference is $2 \pi$.

Step 2 Consider the (h00) reflection with $h=2$\\
In this case, there is a $2 \times 2 \pi$ difference between the waves from the two A layers, and if B were to lie at $x=\frac{1}{2}$ it would give rise to a wave that differed in phase by $2 \pi$ from the wave from the lower A layer. Thus, for a general fractional position $x$, the phase difference for a (200) reflection is $2 \times 2 \pi x$.

\subsection*{Step 3 Generalize these conclusions}
For a general ( $h 00$ ) reflection, the phase difference is $h \times 2 \pi x$. For three dimensions, this result generalizes to $\phi_{h k l}=2 \pi(h x+$ $k y+l z)$.\\
\includegraphics[max width=\textwidth, center]{2025_02_27_d4b95d9718f0b9e2e6b9g-682(2)}

Figure 15B. 8 Diffraction from a crystal containing two kinds of atoms. (a) For a (100) reflection from the A planes, there is a phase difference of $2 \pi$ between waves reflected by neighbouring planes. (b) For a (200) reflection, the phase difference is $4 \pi$. The reflection from a B plane at a fractional distance $x a$ from an $A$ plane has a phase that is $x$ times these phase differences.

\subsection*{Step 4 Formulate the total amplitude of the scattered waves}
 If the amplitude of the waves scattered from A is $f_{\mathrm{A}}$ at the detector, that of the waves scattered from B with phase difference $\phi_{h k l}$ is $f_{\mathrm{B}}{ }^{\mathrm{i} \phi_{h k l}}$. The total amplitude at the detector is therefore$$
F_{h k l}=f_{\mathrm{A}}+f_{\mathrm{B}} \mathrm{e}^{\mathrm{i} \mathrm{p}_{h k l}}
$$

When there are several atoms present, each with scattering factor $f_{j}$ and phase $\phi_{h k l}(j)=2 \pi\left(h x_{j}+k y_{j}+l z_{j}\right)$, the total amplitude of the $(h k l)$ reflection, the structure factor, is\\
\includegraphics[max width=\textwidth, center]{2025_02_27_d4b95d9718f0b9e2e6b9g-682}

\subsection*{Example 15B. 1}
Calculating a structure factor\\
Calculate the structure factor for the unit cell of NaCl depicted in Fig. 15B.9.\\
\includegraphics[max width=\textwidth, center]{2025_02_27_d4b95d9718f0b9e2e6b9g-682(1)}

Figure 15B. 9 The location of the atoms for the structure factor calculation in Example 15B.1. The red spheres are $\mathrm{Na}^{+}$, the green spheres are $\mathrm{Cl}^{-}$.

Collect your thoughts You need to evaluate the sum in eqn 15B.3. The sum is over all atoms in the unit cell, so you need to know the location of each atom, expressed as a fraction of the unit cell parameters. Several of the atomic coordinates are already marked in the figure. The atoms at the corners of the cube are shared between eight adjacent unit cells, so in the calculation of the structure factor these atoms is given a weight of $\frac{1}{8}$ and their scattering factors taken to be $\frac{1}{8} f$. The atoms on the faces are shared between two cells and so have a weight of $\frac{1}{2}$; those on the edges are shared between four cells, and so have a weight of $\frac{1}{4}$. Write $f^{+}$for the $\mathrm{Na}^{+}$scattering factor and $f^{-}$for the $\mathrm{Cl}^{-}$scattering factor; for simplicity, ignore the fact that scattering occurs in non-forward directions and suppose that all the $\mathrm{Na}^{+}$have the same scattering factor, and likewise for the $\mathrm{Cl}^{-}$ions. The best way to proceed is to draw up a table showing the weights, positions, and phases. The phase factors $\mathrm{e}^{\mathrm{i} \phi_{h l}}$ are evaluated by noting that $h, k$, and $l$ are integers. A useful identity is that $\mathrm{e}^{\mathrm{i} n \pi}$ is +1 for even $n$, and -1 for odd $n$, more succinctly expressed as $\mathrm{e}^{\mathrm{i} n \pi}=(-1)^{n}$.

The solution The table for the $\mathrm{Na}^{+}$ions is

\begin{center}
\begin{tabular}{|c|c|c|c|c|c|}
\hline
Atom & Weight & $x$ & $y$ & $z$ & $\phi_{h k l}$ \\
\hline
1 & $\frac{1}{8}$ & 0 & 0 & 0 & 0 \\
\hline
2 & $\frac{1}{8}$ & 1 & 0 & 0 & $2 \pi h$ \\
\hline
3 & $\frac{1}{8}$ & 0 & 1 & 0 & $2 \pi k$ \\
\hline
4 & $\frac{1}{8}$ & 1 & 1 & 0 & $2 \pi(h+k)$ \\
\hline
5 & $\frac{1}{8}$ & 0 & 0 & 1 & $2 \pi l$ \\
\hline
6 & $\frac{1}{8}$ & 1 & 0 & 1 & $2 \pi(h+l)$ \\
\hline
7 & $\frac{1}{8}$ & 0 & 1 & 1 & $2 \pi(k+l)$ \\
\hline
8 & $\frac{1}{8}$ & 1 & 1 & 1 & $2 \pi(h+k+l)$ \\
\hline
9 & $\frac{1}{2}$ & $\frac{1}{2}$ & $\frac{1}{2}$ & 0 & $2 \pi\left(\frac{1}{2} h+\frac{1}{2} k\right)$ \\
\hline
10 & $\frac{1}{2}$ & $\frac{1}{2}$ & 0 & $\frac{1}{2}$ & $2 \pi\left(\frac{1}{2} h+\frac{1}{2} l\right)$ \\
\hline
11 & $\frac{1}{2}$ & 0 & $\frac{1}{2}$ & $\frac{1}{2}$ & $2 \pi\left(\frac{1}{2} k+\frac{1}{2} l\right)$ \\
\hline
12 & $\frac{1}{2}$ & 1 & $\frac{1}{2}$ & $\frac{1}{2}$ & $2 \pi\left(h+\frac{1}{2} k+\frac{1}{2} l\right)$ \\
\hline
13 & $\frac{1}{2}$ & $\frac{1}{2}$ & 1 & $\frac{1}{2}$ & $2 \pi\left(\frac{1}{2} h+k+\frac{1}{2} l\right)$ \\
\hline
14 & $\frac{1}{2}$ & $\frac{1}{2}$ & $\frac{1}{2}$ & 1 & $2 \pi\left(\frac{1}{2} h+\frac{1}{2} k+l\right)$ \\
\hline
\end{tabular}
\end{center}

The phase factors for the first eight Na atoms in the table are all +1 , and as they each have a weight of $\frac{1}{8}$, the total contribution to the structure factor is $f^{+}$. The remaining six atoms all have weight $\frac{1}{2}$, and their contribution to the structure factor is

$$
\begin{gathered}
\frac{1}{2}\left[\mathrm{e}^{\mathrm{i} 2 \pi(h / 2+k / 2)}+\mathrm{e}^{\mathrm{i} 2 \pi(h / 2+l / 2)}+\mathrm{e}^{\mathrm{i} 2 \pi(k / 2+l / 2)}+\mathrm{e}^{\mathrm{i} 2 \pi(h+k / 2+l / 2)}\right. \\
\left.\quad+\mathrm{e}^{\mathrm{i} 2 \pi(h / 2+k+l / 2)}+\mathrm{e}^{\mathrm{i} 2 \pi(h / 2+k / 2+l)}\right]
\end{gathered}
$$

The terms $\mathrm{e}^{\mathrm{i} 2 \pi h}, \mathrm{e}^{\mathrm{i} 2 \pi k}$, and $\mathrm{e}^{\mathrm{i} 2 \pi l}$ are all +1 , so the last three terms can be simplified

$$
\begin{aligned}
& \frac{1}{2}\left[\mathrm{e}^{\mathrm{i} 2 \pi(h / 2+k / 2)}+\mathrm{e}^{\mathrm{i} 2 \pi(h / 2+l / 2)}+\mathrm{e}^{\mathrm{i} 2 \pi(k / 2+l / 2)}+\mathrm{e}^{\mathrm{i} 2 \pi(k / 2+l / 2)}\right. \\
& \left.\quad+\mathrm{e}^{\mathrm{i} 2 \pi(h / 2+l / 2)}+\mathrm{e}^{\mathrm{i} 2 \pi(h / 2+k / 2)}\right]
\end{aligned}
$$

A further simplification is to use $\mathrm{e}^{\mathrm{i} n \pi}=(-1)^{n}$ :

$$
\begin{aligned}
& \frac{1}{2}\left[(-1)^{h+k}+(-1)^{h+l}+(-1)^{k+l}+(-1)^{k+l}+(-1)^{h+l}+(-1)^{h+k}\right] \\
& \quad=(-1)^{h+k}+(-1)^{h+l}+(-1)^{k+l}
\end{aligned}
$$

The overall contribution of the $\mathrm{Na}^{+}$ions to the structure factor is therefore

$$
f^{+}\left[1+(-1)^{h+k}+(-1)^{h+l}+(-1)^{k+l}\right]
$$

The table for the $\mathrm{Cl}^{-}$ions is

\begin{center}
\begin{tabular}{lcccll}
\hline
Atom & Weight & $x$ & $y$ & $z$ & $\phi_{h k l}$ \\
\hline
1 & $\frac{1}{4}$ & $\frac{1}{2}$ & 0 & 0 & $\pi h$ \\
2 & $\frac{1}{4}$ & 0 & $\frac{1}{2}$ & 0 & $\pi k$ \\
3 & $\frac{1}{4}$ & 1 & $\frac{1}{2}$ & 0 & $2 \pi\left(h+\frac{1}{2} k\right)$ \\
4 & $\frac{1}{4}$ & $\frac{1}{2}$ & 1 & 0 & $2 \pi\left(\frac{1}{2} h+k\right)$ \\
5 & $\frac{1}{4}$ & 0 & 0 & $\frac{1}{2}$ & $\pi l$ \\
6 & $\frac{1}{4}$ & 1 & 0 & $\frac{1}{2}$ & $2 \pi\left(h+\frac{1}{2} l\right)$ \\
7 & $\frac{1}{4}$ & 0 & 1 & $\frac{1}{2}$ & $2 \pi\left(k+\frac{1}{2} l\right)$ \\
8 & $\frac{1}{4}$ & 1 & 1 & $\frac{1}{2}$ & $2 \pi\left(h+k+\frac{1}{2} l\right)$ \\
9 & 1 & $\frac{1}{2}$ & $\frac{1}{2}$ & $\frac{1}{2}$ & $2 \pi\left(\frac{1}{2} h+\frac{1}{2} k+\frac{1}{2} l\right)$ \\
10 & $\frac{1}{4}$ & $\frac{1}{2}$ & 0 & 1 & $2 \pi\left(\frac{1}{2} h+l\right)$ \\
11 & $\frac{1}{4}$ & 0 & $\frac{1}{2}$ & 1 & $2 \pi\left(\frac{1}{2} k+l\right)$ \\
12 & $\frac{1}{4}$ & 1 & $\frac{1}{2}$ & 1 & $2 \pi\left(h+\frac{1}{2} k+l\right)$ \\
13 & $\frac{1}{4}$ & $\frac{1}{2}$ & 1 & 1 & $2 \pi\left(\frac{1}{2} h+k+l\right)$ \\
\hline
\end{tabular}
\end{center}

A similar procedure to that used for the $\mathrm{Na}^{+}$ions gives the following contribution of the $\mathrm{Cl}^{-}$ions to the structure factor

$$
f^{-}\left[(-1)^{h+k+l}+(-1)^{h}+(-1)^{k}+(-1)^{l}\right]
$$

The structure factor is therefore

$$
\begin{aligned}
F_{h k l}=f^{+}[ & \left.+(-1)^{h+k}+(-1)^{h+l}+(-1)^{k+l}\right] \\
& +f^{-}\left[(-1)^{h+k+l}+(-1)^{h}+(-1)^{k}+(-1)^{l}\right]
\end{aligned}
$$

Now note that:

\begin{itemize}
  \item if $h, k$, and $l$ are all even, $F_{h k l}=f^{+}\{1+1+1+1\}+f^{-}\{1+$ $1+1+1\}=4\left(f^{+}+f^{-}\right)$
  \item if $h, k$, and $l$ are all odd, $F_{h k l}=4\left(f^{+}-f^{-}\right)$
  \item if one index is odd and two are even, or vice versa, $F_{h k l}=0$
\end{itemize}

The $h k l$ all-odd reflections are therefore less intense than the $h k l$ all even, and some of the reflections are absent.\\
Comment. If $f^{+}=f^{-}$, which is the case for identical atoms, the $h k l$ all-odd reflections have zero intensity; such a structure would be a cubic $P$ lattice with lattice parameter $a / 2$.

Self-test 15B. 1 Which reflections cannot be observed for a cubic I lattice?

The intensity of a reflection is proportional to the square modulus of the amplitude of the wave, which is in turn proportional to the structure factor, $F_{h k l}$. If the structure factor is $f_{\mathrm{A}}+f_{\mathrm{B}} \mathrm{e}^{\mathrm{i} \mathrm{i}_{h k}}$, the intensity, $I_{h k l}$, is

$$
\begin{aligned}
I_{h k l} \propto F_{h k l}^{*} F_{h k l} & =\left(f_{\mathrm{A}}+f_{\mathrm{B}} \mathrm{e}^{-\mathrm{i} \phi_{h l l}}\right)\left(f_{\mathrm{A}}+f_{\mathrm{B}} \mathrm{e}^{\mathrm{i} \phi_{h l l}}\right) \\
& =f_{\mathrm{A}}^{2}+f_{\mathrm{B}}^{2}+f_{\mathrm{A}} f_{\mathrm{B}}\left(\mathrm{e}^{\mathrm{i} \phi_{h l l}}+\mathrm{e}^{-\mathrm{i} \phi_{h l l}}\right) \\
& =\mathrm{e}^{\mathrm{ix}}+\mathrm{e}^{-\mathrm{ix}}=2 \cos x \\
& =f_{\mathrm{A}}^{2}+f_{\mathrm{B}}^{2}+2 f_{\mathrm{A}} f_{\mathrm{B}} \cos \phi_{h k l}
\end{aligned}
$$

The cosine term either adds to or subtracts from $f_{\mathrm{A}}^{2}+f_{\mathrm{B}}^{2}$ depending on the value of $\phi_{h k}$, which in turn depends on $h, k$, and $l$ and $x, y$, and $z$. Hence, there is a variation in the intensities of the reflections with different $h k l$. The A and B reflections interfere destructively when the phase difference is $\pi$, and in this case the total intensity is zero if the atoms have the same scattering power. For example, if the unit cell is cubic I with a B atom at $x=y=z=\frac{1}{2}$, then the $\mathrm{A}, \mathrm{B}$ phase difference is $(h+k+l) \pi$. Therefore, all reflections for odd values of $h+k+l$ vanish if A and B are identical atoms because the waves from $A$ and $B$ are displaced in phase by $\pi$.

For a cubic P lattice diffraction is possible for all $\{h k l\}$, therefore the diffraction pattern for a cubic I lattice can be constructed from that for cubic P by striking out all reflections with odd values of $h+k+l$. Similarly, for a cubic F lattice the missing lines are ones with two out of $h, k$, and $l$ odd, and the remaining one even, or two even and one odd. Recognition of these systematic absences in a powder spectrum can be used to assign the lattice type (Fig. 15B.10).

The intensity of the $\{h k l\}$ reflection is proportional to $\left|F_{h k l}\right|^{2}$, so in principle the structure factors can be determined experimentally by taking the square root of the corresponding intensities (but see below). Then, once all the structure factors $F_{h k l}$ are known the electron density distribution, $\rho(r)$, in the unit cell can be calculated by using

$$
\rho(\boldsymbol{r})=\frac{1}{V} \sum_{h k l} F_{h k l} \mathrm{e}^{-2 \pi \mathrm{i}(h x+k y+l z)} \quad \text { Fourier synthesis }
$$

(15B.4)\\
where $V$ is the volume of the unit cell. Equation 15B. 4 is called a Fourier synthesis of the electron density. Fourier transforms occur throughout chemistry in a variety of guises, and are described in more detail in The chemist's toolkit 28 in Topic 12C.\\
\includegraphics[max width=\textwidth, center]{2025_02_27_d4b95d9718f0b9e2e6b9g-684(2)}

Cubic P\\
\includegraphics[max width=\textwidth, center]{2025_02_27_d4b95d9718f0b9e2e6b9g-684(1)}

Cubic F\\
Figure 15B. 10 The powder diffraction patterns and the systematic absences of three versions of a cubic cell as a function of angle: cubic F (fcc; $h, k, l$ all even or all odd are present), cubic I (bcc; $h+k+l=$ odd are absent), cubic P. Comparison of the observed pattern with patterns like these enables the unit cell to be identified. The locations of the lines give the cell dimensions.

The essence of the procedure in this case is to express the varying electron density in a unit cell as a superposition of sine and cosine waves.

\subsection*{Example 15B. 2}
Calculating an electron density by Fourier synthesis

Consider the $\{h 00\}$ planes of a crystal extending indefinitely in the $x$ direction. In an X-ray analysis the structure factors were found as follows:

\begin{center}
\begin{tabular}{lrrrrrrrrrr}
$h$ & 0 & 1 & 2 & 3 & 4 & 5 & 6 & 7 & 8 & 9 \\
$F_{h}$ & 16 & -10 & 2 & -1 & 7 & -10 & 8 & -3 & 2 & -3 \\
$h$ & 10 & 11 & 12 & 13 & 14 & 15 &  &  &  &  \\
$F_{h}$ & 6 & -5 & 3 & -2 & 2 & -3 &  &  &  &  \\
\end{tabular}
\end{center}

It was also found that $F_{-h}=F_{h}$. Construct a plot of the electron density projected on to the xaxis of the unit cell.

Collect your thoughts You need to substitute these values into eqn 15B.4, but because the problem is one-dimensional, the sum is over only the index $h$ and only the terms $\mathrm{e}^{-2 \pi i h x}$ need be considered.

The solution Because $F_{-h}=F_{h}$, the sum, rather than running from $h=-\infty$ to $+\infty$, can be written as running from 1 to $+\infty$ :

$$
\begin{aligned}
V \rho(x) & =\sum_{h=-\infty}^{\infty} F_{h} \mathrm{e}^{-2 \pi i h x}=F_{0}+\sum_{h=1}^{\infty}(F_{h} \mathrm{e}^{-2 \pi i h x}+\overbrace{F_{-h}}^{=F_{h}} \mathrm{e}^{2 \pi i h x}) \\
& =F_{0}+\sum_{h=1}^{\infty} F_{h}\left(\mathrm{e}^{-2 \pi i h x}+\mathrm{e}^{2 \pi i h x}\right) \\
& =\underbrace{}_{\mathrm{e}_{0}^{\mathrm{ix} x}+\mathrm{e}^{-i x}=2 \cos x}+2 \sum_{h=1}^{\infty} F_{h} \cos 2 \pi h x
\end{aligned}
$$

\begin{center}
\includegraphics[max width=\textwidth]{2025_02_27_d4b95d9718f0b9e2e6b9g-684(3)}
\end{center}

Figure 15B. 11 The plot of the electron density calculated in Example 15B. 2 (green) and Self-test 15B. 2 (purple).

Only 15 values of $F_{h}$ are given, so $\rho(x)$ will be approximate: the result (computed using mathematical software) is plotted in Fig. 15B. 11 (green line). There are three clear maxima in this function, which can be identified as the positions of three atoms.

Comment. The more terms that are included (meaning the more reflections that are measured), the more accurate is the density plot. Terms corresponding to high values of $h$ (which correspond to short-wavelength cosine terms in the sum) account for the finer details of the electron density; low values of $h$ account for the broad features.

Self-test 15B.2 Use mathematical software to experiment with the result of altering the structure factors in the table: consider the effect of both changing the signs and amplitudes. For example, use the same values of $F_{h}$ as above, but change all the signs for $h \geq 6$.\\
\includegraphics[max width=\textwidth, center]{2025_02_27_d4b95d9718f0b9e2e6b9g-684}

\subsection*{(e) The determination of structure}
The structure factors used in eqn 15B. 4 to compute the electron density are in general complex quantities that can be written $\left|F_{h k \mid}\right| e^{\mathrm{i} \alpha}$, where $\left|F_{h k l}\right|$ is the amplitude and $\alpha$ the phase ('phase' in the sense used to express a complex number by a diagram in two dimensions; see The chemist's toolkit 16 in Topic 7C). However, the observed intensity $I_{h k l}$ is proportional to the square modulus of the structure factor, $\left|F_{h k l}\right|^{2}$, so from the experiment no information is available about the phase, which may lie anywhere from 0 to $2 \pi$. This ambiguity is called the phase problem; its consequences are illustrated by comparing the two plots in Fig. 15B. 11 in which the phases of the structure factors have been changed but the amplitudes kept the same. Some way must be found to assign phases to the structure factors, because unless these are known $\rho$ cannot be evaluated using eqn 15B.4. The phase problem is less severe for centrosymmetric unit cells, for then the structure factors are real. It still remains a problem though, to decide whether $F_{h k l}$ is positive or negative.

The phase problem can be overcome to some extent by a variety of methods. One procedure that is widely used for inorganic materials with a reasonably small number of atoms in a unit cell, and for organic molecules with a small number of heavy atoms, is the Patterson synthesis. Instead of the structure factors $F_{h k l}$, the values of $\left|F_{h k l}\right|^{2}$, which can be obtained without ambiguity from the intensities, are used in an expression that resembles eqn 15B.4:


\begin{equation*}
P(\boldsymbol{r})=\frac{1}{V} \sum_{h k l}\left|F_{h k l}\right|^{2} \mathrm{e}^{-2 \pi \mathrm{i}(h x+k y+l z)} \quad \text { Patterson synthesis } \tag{15B.5}
\end{equation*}


where the $r$ are the vector separations between the atoms in the unit cell; that is, the distances and directions between atoms. Whereas the electron density function $\rho(r)$ is the probability density of the positions of atoms, the function $P(r)$ is a map of the probability density of the separations between atoms; $P(r)$ is often called the Patterson map. In such a map, a peak at a position specified by a vector $r$ from the origin arises from pairs of atoms that are separated by $r$. Thus, if atom A is at the coordinates $\left(x_{\mathrm{A}}, y_{\mathrm{A}}, z_{\mathrm{A}}\right)$ and atom B is at $\left(x_{\mathrm{B}}, y_{\mathrm{B}}, z_{\mathrm{B}}\right)$, then there will be a peak at $\left(x_{\mathrm{A}}-x_{\mathrm{B}}, y_{\mathrm{A}}-y_{\mathrm{B}}, z_{\mathrm{A}}-z_{\mathrm{B}}\right)$ in the Patterson map. There will also be a peak at the negative of these coordinates, because there is a separation vector from $B$ to $A$ as well as a separation vector from A to B. The height of the peak in the map is proportional to the product of the atomic numbers of the two atoms, $Z_{\mathrm{A}} Z_{\mathrm{B}}$. The Patterson map also shows a strong feature at its origin arising from the separation between each atom and itself, which is necessarily zero.

\subsection*{Brief illustration 15B. 2}
For the electron density shown in Fig. 15B.12a, the corresponding Patterson map is shown in Fig. 15B.12b. The location of each peak, relative to the origin, corresponds to the\\
\includegraphics[max width=\textwidth, center]{2025_02_27_d4b95d9718f0b9e2e6b9g-685}

Figure 15B. 12 The Patterson map corresponding to the electron density in (a) is shown in (b). The distance and orientation of each spot from the origin gives the orientation and separation of one atom-atom separation in (a); in addition, there is a large spot at the origin. Some of the typical distances and their contribution to (b) are shown as $R_{1}$, etc.\\
separation and relative orientation of a pair of atoms in the cell. Note that the Patterson map is centrosymmetric and has a strong feature at the origin.

Heavy atoms dominate the scattering because their scattering factors are large, of the order of their atomic numbers, and their locations may be deduced quite readily. The sign of $F_{h k l}$ can now be calculated from the known locations of the heavy atoms in the unit cell, and to a high probability the phase calculated for them will be the same as the phase for the entire unit cell. To see why this is so, consider a centrosymmetric cell for which each term in the structure factor is either positive or negative. The structure factor has the form


\begin{equation*}
F=( \pm) f_{\text {heavy }}+( \pm) f_{\text {light }}+( \pm) f_{\text {light }}+\cdots \tag{15B.6}
\end{equation*}


where $f_{\text {heavy }}$ is the scattering factor of the heavy atom and $f_{\text {light }}$ the scattering factors of the light atoms. The $f_{\text {light }}$ are all much smaller than $f_{\text {heavy }}$, and their phases are more or less random if the atoms are distributed throughout the unit cell. Therefore, the net effect of the $f_{\text {light }}$ is to change $F$ only slightly from $f_{\text {heary }}$, and with reasonable confidence $F$ will have the same sign as that calculated from the location of the heavy atoms. This phase can then be combined with the observed $|F|$ (from the reflection intensity) to perform a Fourier synthesis of the full electron density in the unit cell, and hence to locate the light atoms as well as the heavy atoms.

Modern structural analyses make extensive use of direct methods. Direct methods are based on the possibility of treating the atoms in a unit cell as being virtually randomly distributed (from the radiation's point of view), and then to use statistical techniques to compute the probabilities that the phases have a particular value. It is possible to deduce relations between some structure factors and sums (and sums of squares) of others, which have the effect of constraining the phases to particular values (with high probability, so long as the structure factors are large). For example, the Sayre probability relation has the form


\begin{align*}
& \text { sign of } \left.F_{h+h^{\prime}, k+k^{\prime}, l \mu^{\prime}} \text { is probably equal to (sign of } F_{h k l}\right) \\
& \times\left(\operatorname{sign} \text { of } F_{h^{\prime} k \prime}\right) \tag{15B.7}
\end{align*}


For example, if $F_{122}$ and $F_{232}$ are both large and negative, then it is highly likely that $F_{354}$, provided it is large, will be positive.

In the final stages of the determination of a crystal structure, the parameters describing the structure (atom positions, for instance) are adjusted systematically to give the best fit between the observed intensities and those calculated from the model of the structure deduced from the diffraction pattern. This process is called structure refinement. Not only does the procedure give accurate positions for all the atoms in the unit cell, but it also gives an estimate of the errors in those positions and in the bond lengths and angles derived from them. The procedure also provides information on the vibrational amplitudes of the atoms.

\subsection*{15B.2 Neutron and electron diffraction}
Neutrons and electrons can also give rise to diffraction on account of their wave nature; their wavelength is given by the de Broglie relation (Topic 7A, $\lambda=h / p$ ). Neutrons generated in a nuclear reactor and then slowed to thermal velocities have wavelengths similar to those of X-rays and may also be used for diffraction studies. For instance, a neutron generated in a reactor and slowed to thermal velocities by repeated collisions with a moderator (such as graphite) until it is travelling at about $4 \mathrm{~km} \mathrm{~s}^{-1}$ has a wavelength of about 100 pm . In practice, a range of wavelengths occurs in a neutron beam, but a monochromatic beam can be selected by diffraction from a crystal, such as a single crystal of germanium.

\subsection*{Example 15B.3 Calculating the typical wavelength of thermal neutrons}
Calculate the typical wavelength of neutrons after reaching thermal equilibrium with their surroundings at 373 K. For simplicity, assume that the particles are travelling in one dimension.

Collect your thoughts In order to use the de Broglie relation, you need to know the momentum of the neutrons, and therefore their velocity. The velocity can be computed from the kinetic energy, which you can assume has its equipartition value for translation in one dimension, $E_{\mathrm{k}}=\frac{1}{2} k T$ (see The chemist's toolkit 7 in Topic 2A). The mass of a neutron is given inside the front cover.

The solution From the equipartition principle, the mean translational kinetic energy of a neutron at a temperature $T$ travelling in the $x$-direction is $E_{\mathrm{k}}=\frac{1}{2} k T$. The kinetic energy is also equal to $p^{2} / 2 m$, where $p$ is the momentum of the neutron and $m$ is its mass. Hence, $p=(m k T)^{1 / 2}$. It follows from the de Broglie relation $\lambda=h / p$ that the wavelength of the neutron is

$$
\lambda=\frac{h}{(m k T)^{1 / 2}}
$$

Therefore, at 373 K ,

$$
\begin{aligned}
\lambda & =\frac{6.626 \times 10^{-34} \mathrm{Js}}{\left\{\left(1.675 \times 10^{-27} \mathrm{~kg}\right) \times\left(1.381 \times 10^{-23} \mathrm{JK}^{-1}\right) \times(373 \mathrm{~K})\right\}^{1 / 2}} \\
& =\frac{11 \mathrm{~kg} \mathrm{~m}^{2} \mathrm{~s}^{-2}}{\left(1.675 \times 10^{-27} \times 1.381 \times 10^{-23} \times 373\right)^{1 / 2}} \frac{\mathrm{~kg} \mathrm{~m}^{2} \mathrm{~s}^{-1}}{\left(\mathrm{~kg}^{2} \mathrm{~m}^{2} \mathrm{~s}^{-2}\right)^{1 / 2}} \\
& =2.26 \times 10^{-10} \mathrm{~m}=226 \mathrm{pm}
\end{aligned}
$$

Self-test 15B.3 Calculate the temperature needed for the average wavelength of the neutrons to be 100 pm .

Neutron diffraction differs from X-ray diffraction in two main respects. First, the scattering of neutrons is a nuclear phenomenon. Neutrons pass through the extranuclear electrons of atoms and interact with the nuclei through the 'strong force' that is responsible for binding nucleons together. As a result, the intensity with which neutrons are scattered is independent of the number of electrons and neighbouring elements in the periodic table might scatter neutrons with markedly different intensities. Neutron diffraction can be used to distinguish atoms of elements such as Ni and Co that are present in the same compound and to study order-disorder phase transitions in FeCo . A second difference is that neutrons possess a magnetic moment due to their spin. This magnetic moment can couple to the magnetic fields of atoms or ions in a crystal (if the ions have unpaired electrons) and modify the diffraction pattern. One consequence is that neutron diffraction is well suited to the investigation of magnetically ordered lattices in which neighbouring atoms may be of the same element but have different orientations of their electronic spin (Fig. 15B.13).

Electrons accelerated through a potential difference of 40 kV have wavelengths of about 6 pm , and so are also suitable for diffraction studies of molecules. Consider the scattering of electrons from pairs of atoms with centres separated by a distance $R_{i j}$ and orientated at a definite angle $\theta$ to an incident beam of electrons. When the molecule consists of a number of atoms, the scattering intensity can be calculated by summing over the contribution from all pairs. The total intensity $I(\theta)$ is given by the Wierl equation:

$$
I(\theta)=\sum_{i, j} f_{i} f_{j} \frac{\sin s R_{i j}}{s R_{i j}} \quad s=\frac{4 \pi}{\lambda} \sin \frac{1}{2} \theta \quad \text { Wierl equation }
$$

(15B.8)\\
where $\lambda$ is the wavelength of the electrons in the beam, and $f$ is the electron scattering factor, a measure of the electron scattering power of the atom. The main application of electron diffraction techniques is to the study of surfaces (Topic 19A), and you are invited to explore the Wierl equation in Problem P15B.8.\\
\includegraphics[max width=\textwidth, center]{2025_02_27_d4b95d9718f0b9e2e6b9g-686}

Figure 15B. 13 If the spins of atoms at lattice points are orderly, as in this material, where the spins of one set of atoms are aligned antiparallel to those of the other set, neutron diffraction detects two interpenetrating simple cubic lattices on account of the magnetic interaction of the neutron with the atoms, but X-ray diffraction would see only a single bcc lattice.

\subsection*{Checklist of concepts}
$\square$ 1. A reflection refers to an intense beam emerging in a particular direction and arising from constructive interference.2. The glancing angle, $2 \theta$, is the angle through which a beam is deflected.3. Bragg's law relates the angle of a diffracted beam to the spacing of a given set of lattice planes.4. The scattering factor is a measure of the ability of an atom to scatter electromagnetic radiation.5. The structure factor is the overall amplitude of a wave diffracted by the $\{h k l\}$ planes and atoms distributed through the unit cell.6. The electron density and the diffraction pattern are related by a Fourier transform.7. Fourier synthesis is the construction of the electron density distribution from structure factors.\\
$\square$ 8. The phase problem arises because it is possible to measure only the intensity of the reflections and not their phases; as a result Fourier synthesis cannot be used in a straightforward way to determine the electron density.9. A Patterson map is a map of the interatomic vectors.10. Direct methods use statistical techniques to determine the likely phases of the reflections.\\
$\square$ 11. Structure refinement is the adjustment of structural parameters to give the best fit between the observed intensities and those calculated from the model of the structure deduced from the diffraction pattern.12. The Wierl equation relates the intensity of electron scattering to the distances between pairs of atoms in the sample.

\subsection*{Checklist of equations}
\begin{center}
\begin{tabular}{|c|c|c|c|}
\hline
Property & Equation & Comment & Equation number \\
\hline
Bragg's law & $\lambda=2 d \sin \theta$ & $d$ is the lattice spacing, $2 \theta$ the glancing angle & 15B.1b \\
\hline
Scattering factor & $f=4 \pi \int_{0}^{\infty}[\{\rho(r) \sin k r\} / k r] r^{2} \mathrm{~d} r, k=(4 \pi / \lambda) \sin \theta$ & Spherically symmetrical atom & 15B. 2 \\
\hline
Structure factor & $F_{h k l}=\sum f_{j} \mathrm{e}^{\mathrm{i} \phi_{h k l}(j)}, \quad \phi_{h k l}(j)=2 \pi\left(h x_{j}+k y_{j}+l z_{j}\right)$ & Definition & 15B. 3 \\
\hline
Fourier synthesis & \( \rho(\boldsymbol{r})=(1 / V) \sum_{h k l} F_{h k l} \mathrm{e}^{-2 \pi \mathrm{i}(h x+k y+l z)} \) & $V$ is the volume of the unit cell & 15B. 4 \\
\hline
Patterson synthesis & \( P(\boldsymbol{r})=(1 / V) \sum_{h k l}\left|F_{h k l}\right|^{2} \mathrm{e}^{-2 \pi \mathrm{i}(h x+k y+l z)} \) &  & 15B. 5 \\
\hline
Wierl equation & \( I(\theta)=\sum_{i, j} f_{i} f_{j}\left(\sin s R_{i j} / s R_{i j}\right), s=(4 \pi / \lambda) \sin \frac{1}{2} \theta \) &  & 15B. 8 \\
\hline
\end{tabular}
\end{center}

\section*{TOPIC 15C Bonding in solids}
\subsection*{Why do you need to know this material?}
To understand the properties and structures of solid materials you need to know about the type of bonding that holds together the atoms, ions, and molecules.

\subsection*{What is the key idea?}
Four characteristic types of bonding result in metals, ionic solids, covalent solids, and molecular solids.

\subsection*{What do you need to know already?}
You need to be familiar with molecular interactions (Topic 14B) and the general features of crystal structures (Topic 15A). For the discussion of metallic bonding you should be aware of the principles of Hückel molecular orbital theory (Topic 9E). The discussion of ionic bonding makes use of the concept of enthalpy (Topic 2B).

Solids may be classified into four broad types, namely metals, ionic solids, covalent (or network) solids, and molecular solids. Each is characterized by the nature of the bonding between the constituents.

\subsection*{15C. 1 Metals}
In a metal the electrons are delocalized over arrays of identical cations and bind the whole together into a rigid but ductile and malleable structure. The crystalline forms of metallic elements can be discussed in terms of a model in which their atoms are treated as identical hard spheres. Most metallic elements crystallize in one of three simple forms, two of which can be explained in terms of the hard spheres packing together in the closest possible arrangement.

\subsection*{(a) Close packing}
Figure 15C. 1 shows a close-packed layer of identical spheres, one with maximum utilization of space. A close-packed threedimensional structure is obtained by stacking such layers on top of one another. However, this stacking can be done in different ways and results in close-packed polytypes, which are\\
\includegraphics[max width=\textwidth, center]{2025_02_27_d4b95d9718f0b9e2e6b9g-688(1)}

Figure 15C. 1 A layer of close-packed spheres used to build a three-dimensional close-packed structure.\\
structures that are identical in two dimensions (the closepacked layers) but differ in the third dimension.

In all polytypes, the spheres of the second close-packed layer lie in the depressions of the first layer (Fig. 15C.2). The third layer may be added in either of two ways. In one, the spheres are placed directly above the first layer to give an ABA pattern of layers (Fig. 15C.3a). Alternatively, the spheres may be placed over the depressions in the first layer that are not occupied by the second layer (these depressions are visible in Fig. 15C.2), so giving an ABC pattern (Fig. 15C.3b). Two polytypes are formed if the two stacking patterns are repeated in the vertical direction:

\begin{itemize}
  \item Hexagonally close-packed (hcp): the ABA pattern is repeated, to give the sequence of layers ABABAB ....
  \item Cubic close-packed (ccp): the ABC pattern is repeated, to give the sequence $A B C A B C$....\\
The origins of these names can be seen by referring to Fig. 15C.4. The ccp structure gives rise to a face-centred unit\\
\includegraphics[max width=\textwidth, center]{2025_02_27_d4b95d9718f0b9e2e6b9g-688}
\end{itemize}

Figure 15C. 2 To achieve the greatest packing density, the second layer of close-packed spheres must sit in the depressions of the first layer. The two layers are the $A B$ component of the close-packed structure.\\
\includegraphics[max width=\textwidth, center]{2025_02_27_d4b95d9718f0b9e2e6b9g-689}

Figure 15C. 3 (a) The third layer of close-packed spheres might occupy the depressions lying directly above the spheres in the first layer, resulting in an ABA structure, which corresponds to hexagonal close-packing. (b) Alternatively, the third layer might lie in the depressions that are not above the spheres in the first layer, resulting in an ABC structure, which corresponds to cubic close-packing.\\
\includegraphics[max width=\textwidth, center]{2025_02_27_d4b95d9718f0b9e2e6b9g-689(1)}

Figure 15C. 4 Fragments of the structures shown in Fig. 15C. 3 revealing the (a) hexagonal (b) cubic symmetry. The colours of the spheres are the same as for the layers in Fig. 15C.3.\\
cell, so may also be denoted cubic F (or fcc, for face-centred cubic). It is also possible to have random sequences of layers; however, the hcp and ccp polytypes are the most important. Table 15C. 1 lists the structures adopted by a selection of elements.

The compactness of close-packed structures is indicated by their coordination number, the number of spheres immediately surrounding any selected sphere, which is 12 for both the ccp and hcp structures. Another measure of their

Table 15C. 1 The crystal structures of some elements*

\begin{center}
\begin{tabular}{ll}
\hline
Structure & Element \\
\hline
hcp $^{\ddagger}$ & $\mathrm{Be}, \mathrm{Cd}, \mathrm{Co}, \mathrm{He}, \mathrm{Mg}, \mathrm{Sc}, \mathrm{Ti}, \mathrm{Zn}$ \\
$\mathrm{fcc}^{\ddagger}(\mathrm{ccp}, \mathrm{cubic} \mathrm{F})$ & $\mathrm{Ag}, \mathrm{Al}, \mathrm{Ar}, \mathrm{Au}, \mathrm{Ca}, \mathrm{Cu}, \mathrm{Kr}, \mathrm{Ne}, \mathrm{Ni}, \mathrm{Pd}, \mathrm{Pb}, \mathrm{Pt}, \mathrm{Rh}$, \\
$\mathrm{Rn}, \mathrm{Sr}, \mathrm{Xe}$ &  \\
\end{tabular}
\end{center}, \begin{tabular}{ll}
$\mathrm{Ba}, \mathrm{Cs}, \mathrm{Cr}, \mathrm{Fe}, \mathrm{K}, \mathrm{Li}, \mathrm{Mn}, \mathrm{Mo}, \mathrm{Rb}, \mathrm{Na}, \mathrm{Ta}, \mathrm{W}, \mathrm{V}$ &  \\
cubic (cubic I) & Po \\
\hline
\end{tabular}

\footnotetext{\begin{itemize}
  \item The notation used to describe unit cells is introduced in Topic 15A. The structures are for the elements at 298 K and 1 bar.\\
${ }^{\ddagger}$ Close-packed structures.
\end{itemize}
}
compactness is the packing fraction, the fraction of space occupied by the spheres, which is 0.740 (see Example 15C.1). That is, in a close-packed solid of identical hard spheres, only 26.0 per cent of the volume is empty space. The fact that many metals are close-packed accounts for their high mass densities.

Example 15C. 1

\subsection*{Calculating a packing fraction}
Calculate the packing fraction of a ccp structure formed from hard spheres.\\
Collect your thoughts You need to calculate the volume of the unit cell and the volume of the spheres that are wholly or partly contained within the cell, and then take the ratio of the two volumes. The key step is to establish a relation between the radius of the spheres, $R$, and the cell dimension, $a$. Figure 15C. 5 shows that the sphere in the middle of a face just touches the two spheres at opposite corners of the face. The length of the face diagonal is therefore $4 R$. To calculate the volume of the spheres in the cell, you need to account for the fraction of each sphere that is contained within the cell. Spheres at the corners contribute $\frac{1}{8}$ of their volume to the cell; those on the faces contribute $\frac{1}{2}$ to the cell.

The solution As can be seen from Fig. 15C.5, the length of the face diagonal is $4 R$. From Pythagoras' theorem it follows that $a^{2}+a^{2}=(4 R)^{2}$, so $2 a^{2}=16 R^{2}$ and therefore $a=8^{1 / 2} R$. The volume of the unit cell is $a^{3}$, which is therefore $8^{3 / 2} R^{3}$. There are eight spheres at the corners, each contributing $\frac{1}{8}$ of their volume to the cell (for a net contribution of 1 sphere), and six spheres on the faces, each contributing $\frac{1}{2}$ (for a net contribution of 3 spheres). The total volume occupied by the spheres is equivalent to 4 complete spheres. Because the volume of each sphere is $\frac{4}{3} \pi R^{3}$, the total occupied volume is $\frac{16}{3} \pi R^{3}$. The fraction of space occupied is therefore

$$
\frac{\left(\frac{16}{3}\right) \pi R^{3}}{8^{3 / 2} R^{3}}=\frac{\pi}{3 \sqrt{2}}=0.740
$$

\begin{center}
\includegraphics[max width=\textwidth]{2025_02_27_d4b95d9718f0b9e2e6b9g-689(2)}
\end{center}

Figure 15C. 5 In a ccp unit cell a sphere located on the face of the cube just touches the spheres at opposite corners of the face.

Because an hcp structure has the same coordination number, its packing fraction is the same.\\
Self-test 15C. 1 Calculate the packing fraction of the cubic I (body-centred cubic, bcc) structure in which one sphere is at the centre of a cube formed by eight others. The spheres touch along the body diagonal of the cube.\\
\includegraphics[max width=\textwidth, center]{2025_02_27_d4b95d9718f0b9e2e6b9g-690}

As shown in Table 15C.1, a number of common metals adopt structures that are less than close-packed. The departure from close packing suggests that factors, such as specific covalent bonding between neighbouring atoms, are beginning to influence the structure and impose a particular geometrical arrangement. One such arrangement results in a cubic I (bcc, for body-centred cubic) structure, with one sphere at the centre of a cube formed by eight others. The coordination number of a bcc structure is only 8 , but there are six more atoms not much further away than the eight nearest neighbours. The packing fraction of 0.680 (Self-test 15C.1) is not much smaller than the value for a close-packed structure (0.740), and shows that about two-thirds of the available space is occupied by the atoms.

\subsection*{(b) Electronic structure of metals}
The central aspect of solids that determines their electrical properties (Topic 15E) is the distribution of their electrons. There are two models of this distribution. In one, the nearly free-electron approximation, the valence electrons are assumed to be trapped in a box with a periodic potential, with low energy corresponding to the locations of cations. In the tight-binding approximation, the valence electrons are assumed to occupy molecular orbitals delocalized throughout the solid. The latter model is more in accord with the discussion of electrical properties of solids discussed in Topic 15E, and is described here.

As a starting point, consider a one-dimensional solid, which consists of a single, infinitely long line of atoms. Suppose that each atom has one s orbital available for forming molecular orbitals. The LCAO-MOs of the solid are constructed by adding $N$ atoms in succession to a line, and then inferring the electronic structure by using the building-up principle. One atom contributes one s orbital at a certain energy (Fig. 15C.6). When a second atom is brought up it overlaps the first and forms a bonding and an antibonding orbital. The third atom overlaps its nearest neighbour (and only slightly the next-nearest), and from these three atomic orbitals, three molecular orbitals are formed: one is fully bonding, one fully antibonding, and the intermediate orbital is nonbonding between neighbours. The fourth atom leads to the formation of a fourth molecular orbital. At this stage, it can be seen that the effect of bringing up successive atoms is to spread the range of energies covered\\
\includegraphics[max width=\textwidth, center]{2025_02_27_d4b95d9718f0b9e2e6b9g-690(1)}

Figure 15C. 6 The formation of a band of $N$ molecular orbitals by successive addition of $N$ atomic orbitals in a line. When $N$ becomes infinite the band covers a finite range of energy and the orbitals within it are very closely spaced, but still discrete.\\
by the molecular orbitals and to fill in the range of energies with more and more orbitals (one more for each atom). When $N$ atoms have been added to the line, there are $N$ molecular orbitals covering a finite range of energies: this set of orbitals is said to form a band.

The energies of the molecular orbitals that form the band can be found by using the Hückel approximations described in Topic 9E. They are found by solving the Hückel secular determinant

$$
\begin{array}{ccccc}
\alpha-E & \beta & 0 & \ldots & 0 \\
\beta & \alpha-E & \beta & \ldots & 0 \\
0 & \beta & \alpha-E & \ldots & 0 \\
\vdots & \vdots & \vdots & \ldots & \vdots \\
0 & 0 & 0 & \ldots & \alpha-E
\end{array}=0
$$

where $\alpha$ is the Coulomb integral and $\beta$ is the ( $\mathrm{s}, \mathrm{s}$ ) resonance integral. The general expression for the solutions of this 'tridiagonal determinant' gives the energies $E_{k}$ of the molecular orbitals:

\[
E_{k}=\alpha+2 \beta \cos \frac{k \pi}{N+1} \quad k=1,2, \ldots, N \quad \begin{align*}
& \text { Energy levels }  \tag{15C.1}\\
& {[\text { linear array of }} \\
& \text { s orbitals] }
\end{align*}
\]

It is not hard to show that this expression implies that when $N$ is infinitely large, the separation between neighbouring levels, $E_{k+1}-E_{k}$, is infinitely small, but the band still has finite width overall, with $E_{N}-E_{1} \rightarrow-4 \beta$.

\subsection*{How is that done? 15C. 1 Evaluating the separation}
 between neighbouring levels and width of a bandThis calculation requires looking at a specific limiting case of eqn 15C.1.

Step 1 Write an expression for the energy difference between two neighbouring levels\\
From eqn 15C. 1 it follows that the energy separation between neighbouring energy levels $k$ and $k+1$ is

$$
\begin{aligned}
E_{k+1}-E_{k} & =\left(\alpha+2 \beta \cos \frac{(k+1) \pi}{N+1}\right)-\left(\alpha+2 \beta \cos \frac{k \pi}{N+1}\right) \\
& =2 \beta\left(\cos \frac{(k+1) \pi}{N+1}-\cos \frac{k \pi}{N+1}\right)
\end{aligned}
$$

Step 2 Find the limit as $N \rightarrow \infty$\\
By using the trigonometric identity $\cos (A+B)=\cos A \cos B$ $-\sin A \sin B$, followed by $\cos 0=1$ and $\sin 0=0$, the first (blue) term in parentheses is

$$
\cos \frac{(k+1) \pi}{N+1}=\cos \frac{k \pi}{N+1} \overbrace{\cos \frac{\pi}{N+1}}^{\rightarrow 1 \text { as } N \rightarrow \infty}-\sin \frac{k \pi}{N+1} \overbrace{\sin \frac{\pi}{N+1}}^{\rightarrow 0 \text { as } N \rightarrow \infty}
$$

Therefore, as $N \rightarrow \infty$,

$$
E_{k+1}-E_{k} \rightarrow 2 \beta\left(\cos \frac{k \pi}{N+1}-\cos \frac{k \pi}{N+1}\right)=0
$$

It follows that when $N$ is infinitely large, the difference between neighbouring energy levels is infinitely small.

Step 3 Write an expression for the width of the band for $N \rightarrow \infty$ The width of the band is simply $E_{N}-E_{1}$. Each of the energies can be approximated as follows for the case that $N \rightarrow \infty$.

$$
E_{1}=\alpha+2 \beta \cos \frac{\pi}{N+1}
$$

As $N \rightarrow \infty$, the term $\pi /(N+1)$ tends to zero so the cosine tends to 1 ; therefore in this limit

$$
E_{1}=\alpha+2 \beta
$$

When $k$ has its maximum value of $N$,

$$
E_{N}=\alpha+2 \beta \cos \frac{N \pi}{N+1}
$$

As $N \rightarrow \infty$, the 1 in the denominator can be ignored, so the cosine term becomes $\cos \pi=-1$. Therefore, in this limit $E_{N}=$ $\alpha-2 \beta$, and so the band width is $E_{N}-E_{1} \rightarrow-4 \beta$. Recall that $\beta$ is negative, so that the band width, $-4 \beta$, is positive.

The band can be thought of as consisting of $N$ different molecular orbitals, the lowest-energy orbital $(k=1)$ being fully bonding, and the highest-energy orbital $(k=N)$ being fully antibonding between adjacent atoms (Fig. 15C.7). The molecular orbitals of intermediate energy have $k-1$ nodes distributed along the chain of atoms. Similar bands form in three-dimensional solids.\\
\includegraphics[max width=\textwidth, center]{2025_02_27_d4b95d9718f0b9e2e6b9g-691}

Figure 15C. 7 The overlap of $s$ orbitals gives rise to an $s$ band and the overlap of $p$ orbitals gives rise to a $p$ band. In this case, the $s$ and $p$ orbitals of the atoms are so widely spaced in energy that there is a band gap; it is also possible that the separation will be less, resulting in the bands overlapping.

\subsection*{Brief illustration 15C. 1}
To illustrate the dependence of $E_{k+1}-E_{k}$ on $N$, note that

$$
\begin{aligned}
& \text { for } N=3: \quad E_{2}-E_{1}=2 \beta\left(\cos \frac{2 \pi}{4}-\cos \frac{\pi}{4}\right) \approx-1.414 \beta \\
& \text { for } N=30: \quad E_{2}-E_{1}=2 \beta\left(\cos \frac{2 \pi}{31}-\cos \frac{\pi}{31}\right) \approx-0.0307 \beta \\
& \text { for } N=300: \quad E_{2}-E_{1}=2 \beta\left(\cos \frac{2 \pi}{301}-\cos \frac{\pi}{301}\right) \approx-0.000327 \beta
\end{aligned}
$$

The energy difference decreases with increasing $N$, as expected.

The band formed from overlap of $s$ orbitals is called the $s$ band. If the atoms have $p$ orbitals available, the same procedure leads to a $\mathbf{p}$ band (as shown in the upper half of Fig. 15C.7). If the atomic $p$ orbitals lie higher in energy than the $s$ orbitals, then the $p$ band lies higher than the $s$ band, and there may be a band gap, a range of energies to which no orbital corresponds. However, it is also possible for the bands to touch, with the highest orbital of the $s$ band coincident with the lowest level of the $p$ band, or even overlap (as is the case for the 3 s and 3 p bands in magnesium).

Now consider the electronic structure of a solid formed from $N$ atoms each able to contribute one electron (for example, the alkali metals). The $N$ atomic orbitals give rise to a band consisting of $N$ molecular orbitals. Each of these orbitals can accommodate two spin-paired electrons, so at $T=0$ only the lowest $\frac{1}{2} N$ molecular orbitals are occupied (Fig. 15C.8). The HOMO is called the Fermi level. Only the small number of electrons close to the Fermi level can undergo thermal excitation, so only these electrons contribute to the heat capacity of the metal. It is for this reason that Dulong and Petit's law\\
\includegraphics[max width=\textwidth, center]{2025_02_27_d4b95d9718f0b9e2e6b9g-692(1)}

Figure 15C. 8 When $N$ electrons occupy a band of $N$ orbitals, only half of the orbitals are occupied (at $T=0$ ) because two electrons will occupy each orbital. The highest occupied level is known as the Fermi level.\\
for heat capacities (Topic 7A) gives reasonable agreement with experiment at normal temperatures by considering only the atoms in a sample, not the atoms plus the 'free' electrons. The presence of an incompletely filled band is responsible for electrical conductivity, as explained in Topic 15E.

\subsection*{15C. 2 Ionic solids}
An ionic solid consists of cations and anions held together by electrostatic interactions. Two key questions arise in considering such solids: the relative locations adopted by the ions and the energetics of the resulting structure.

\subsection*{(a) Structure}
When crystals of compounds of monatomic ions (such as NaCl and MgO ) are modelled by stacks of hard spheres it is necessary to allow for the possibility that the ions have different radii (typically with the cations smaller than the anions) and different electrical charges. The coordination number of an ion is the number of nearest neighbours of opposite charge; the structure itself is characterized as having $\left(N_{+}, N_{\perp}\right)$ coordination, where $N_{+}$is the coordination number of the cation and $N_{-}$that of the anion.

Even if, by chance, the ions have the same size, the requirement that the unit cells are electrically neutral make it impossible to achieve 12 -coordinate close-packed ionic structures. As a result, ionic solids are generally less dense than metals. The best packing that can be achieved is the $(8,8)$-coordinate caesium chloride structure in which each cation is surrounded by eight anions and each anion is surrounded by eight cations (Fig. 15C.9). In this structure, an ion of one charge occupies the centre of a cubic unit cell with eight counter ions at its corners. The cell is electrically neutral because the ions at the corners of the cell are shared between eight cells and so contribute one eighth of their charge to each. The structure shown in Fig. 15C. 9 is adopted by CsCl itself and also by CaS .\\
\includegraphics[max width=\textwidth, center]{2025_02_27_d4b95d9718f0b9e2e6b9g-692(2)}

Figure 15C. 9 The caesium chloride structure consists of two interpenetrating simple cubic arrays of ions, one of cations and the other of anions, so that each cube of ions of one kind has a counter-ion at its centre.

When the radii of the ions differ more than they do in CsCl , even eight-coordinate packing cannot be achieved. One common structure adopted is the $(6,6)$-coordinate rock salt structure typified by rock salt itself, NaCl (Fig. 15C.10). In this structure, each cation is surrounded by six anions and each anion is surrounded by six cations. The rock salt structure can be pictured as consisting of two interpenetrating slightly expanded cubic $F$ (fcc) arrays, one composed of cations and the other of anions. This structure is adopted by NaCl itself and also by several other MX compounds, including $\mathrm{KBr}, \mathrm{AgCl}$, MgO , and ScN .

The switch from the caesium chloride structure to the rock salt structure is related to the value of the radius ratio, $\gamma$ :

$$
\gamma=\frac{r_{\text {smaller }}}{r_{\text {larger }}}
$$

Radius ratio [definition]\\
(15C.2)\\
where $r_{\text {smaller }}$ is the radius of the smaller ions in the crystal and $r_{\text {larger }}$ that of the larger ions. The radius-ratio rule, which is derived by considering the geometrical problem of packing the maximum number of hard spheres of one radius around a hard sphere of a different radius, can be summarized as follows:

\begin{center}
\begin{tabular}{ll}
\hline
Radius ratio & Structural type \\
\hline
$\gamma<2^{1 / 2}-1=0.414$ & sphalerite (Fig. 15C.11) \\
$0.414<\gamma<3^{1 / 2}-1=0.732$ & rock salt (Fig. 15C.10) \\
$\gamma>0.732$ & caesium chloride (Fig. 15C.9) \\
\hline
\end{tabular}
\end{center}

\begin{center}
\includegraphics[max width=\textwidth]{2025_02_27_d4b95d9718f0b9e2e6b9g-692}
\end{center}

Figure 15C. 10 The rock salt ( NaCl ) structure consists of two mutually interpenetrating slightly expanded face-centred cubic arrays of ions. The assembly shown here is a unit cell.\\
\includegraphics[max width=\textwidth, center]{2025_02_27_d4b95d9718f0b9e2e6b9g-693}

Figure 15C. 11 The structure of the sphalerite form of ZnS showing the location of the Zn atoms in half the tetrahedral holes formed by the fcc array of $S$ atoms.

The deviation of a structure from that expected on the basis of this rule is often taken to be an indication of a shift from ionic towards covalent bonding. A major source of unreliability, though, is the arbitrariness of ionic radii (as explained in a moment) and their variation with coordination number.

Experimental measurements give the distances between the centres of two ions, and a decision has to be made about how to apportion this difference between the two ions. One approach is simply to assign a value to the radius of one ion and then use this value to infer the radii of other ions. A scale based on the value 140 pm for the radius of the $\mathrm{O}^{2-}$ ion is widely used (Table 15C.2). Other scales are also available (such as one based on $\mathrm{F}^{-}$for discussing halides), and it is essential not to mix values from different scales. Because ionic radii are so arbitrary, predictions based on them must be viewed cautiously.

\subsection*{Brief illustration 15C. 2}
From the values of ionic radii in the Resource section, the radius ratio for MgO is

$$
\gamma=\overbrace{\underbrace{\frac{\overbrace{72 \mathrm{pm}}^{140 \mathrm{pm}}}{\text { radius of } \mathrm{Mg}^{2+}}}_{\text {radius of } \mathrm{O}^{2-}}=0.51 . . .2{ }^{2+}}
$$

which is consistent with the observed rock salt structure of MgO crystals.

Table 15C. 2 Ionic radii, r/pm*

\begin{center}
\begin{tabular}{ll}
\hline
$\mathrm{Na}^{+}$ & $102\left(6^{\ddagger}\right), 116(8)$ \\
$\mathrm{K}^{+}$ & $138(6), 151(8)$ \\
$\mathrm{F}^{-}$ & $128(2), 131(4)$ \\
$\mathrm{Cl}^{-}$ & 181 (close packing) \\
\hline
\end{tabular}
\end{center}

*This scale is based on a value 140 pm for the radius of the $\mathrm{O}^{2-}$ ion. More values are given in the Resource section.\\
$\ddagger$ Coordination number.

\subsection*{(b) Energetics}
The lattice energy of a solid is the change in potential energy when the ions go from being packed together in a solid to being widely separated as a gas. All lattice energies are positive; a high lattice energy indicates that the ions interact strongly with one another to give a tightly bonded solid. The lattice enthalpy, $\Delta H_{\mathrm{L}}$, is the change in standard molar enthalpy for the process $\mathrm{MX}(\mathrm{s}) \rightarrow \mathrm{M}^{+}(\mathrm{g})+\mathrm{X}^{-}(\mathrm{g})$ and its equivalent for other charge types and stoichiometries. At $T=0$ the lattice enthalpy is equal to the lattice energy; at normal temperatures they differ by only a few kilojoules per mole, an amount so small compared to typical lattice energies that the difference is normally neglected.

Each ion in a solid experiences favourable (energy lowering) electrostatic interactions from all the other oppositely charged ions and unfavourable (energy raising) interactions from all the other like-charged ions. The total Coulomb potential energy is the sum of all the electrostatic contributions. Each cation is surrounded by anions, and so there is a large negative contribution to the potential energy from the interaction of opposite charges. Beyond those nearest neighbours, there are cations that contribute a positive term to the total potential energy of the central cation. There is also a negative contribution from the anions beyond those cations, a positive contribution from the cations beyond them, and so on, to the edge of the solid. These favourable and unfavourable interactions become progressively weaker as the distance from the central ion increases, but the net outcome of all these contributions is dominated by the interaction between nearest neighbours and is therefore a lowering of the potential energy.

First, consider a simple one-dimensional model of an ionic solid that consists of a long line of uniformly spaced alternating cations and anions; the distance between neighbouring centres is $d$, the sum of the ionic radii (Fig. 15C.12). If the charge numbers of the ions have the same absolute value $(+1$ and -1 , or +2 and -2 , for instance), then $z_{1}=+z, z_{2}=-z$, and $z_{1} z_{2}=-z^{2}$. The potential energy of the central ion is calculated by summing all the terms, with negative terms representing favourable interactions with oppositely charged ions and positive terms representing unfavourable interactions with likecharged ions. For the interaction with ions extending in a line to the right of the central ion, the contribution of the Coulomb interaction to the lattice energy is

$$
\begin{aligned}
E_{\mathrm{p}} & =\frac{1}{4 \pi \varepsilon_{0}} \times\left(-\frac{z^{2} e^{2}}{d}+\frac{z^{2} e^{2}}{2 d}-\frac{z^{2} e^{2}}{3 d}+\frac{z^{2} e^{2}}{4 d}-\cdots\right) \\
& =\frac{z^{2} e^{2}}{4 \pi \varepsilon_{0} d} \times \overbrace{\left(-1+\frac{1}{2}-\frac{1}{3}+\frac{1}{4}-\cdots\right)}^{-\ln 2} \\
& =-\frac{z^{2} e^{2}}{4 \pi \varepsilon_{0} d} \times \ln 2
\end{aligned}
$$

\begin{center}
\includegraphics[max width=\textwidth]{2025_02_27_d4b95d9718f0b9e2e6b9g-694}
\end{center}

Figure 15C. 12 A line of alternating cations and anions used in the calculation of the Madelung constant in one dimension.

To complete the calculation $E_{\mathrm{p}}$ is multiplied by 2 to obtain the total energy arising from interactions on both sides of the ion, and then by Avogadro's constant, $N_{\mathrm{A}}$, to obtain an expression for the Coulomb contribution to the (molar) lattice energy. The outcome is

$$
E_{\mathrm{P}}=-2 \ln 2 \times \frac{z^{2} N_{\mathrm{A}} e^{2}}{4 \pi \varepsilon_{0} d}
$$

with $d=r_{\text {cation }}+r_{\text {anion }}$. This energy is negative, corresponding to a net favourable interaction. The calculation can be extended to three-dimensional arrays of ions with different charge numbers $z_{\mathrm{A}}$ and $z_{\mathrm{B}}$ :


\begin{equation*}
E_{\mathrm{p}}=-A \times \frac{\left|z_{\mathrm{A}} z_{\mathrm{B}}\right| N_{\mathrm{A}} e^{2}}{4 \pi \varepsilon_{0} d} \tag{15C.3}
\end{equation*}


The factor $A$ is a positive numerical constant called the Madelung constant; its value depends on how the ions are arranged in the crystal. For a rock salt structure, $A=1.748$; Table 15C. 3 lists Madelung constants for other common structures.

The Coulomb interaction is not the only contribution to the lattice energy. When atomic orbitals overlap to form bonding and antibonding molecular orbitals and both kinds of orbitals are full, there is an increase in energy because the antibonding orbital is raised in energy more than the bonding orbital is lowered (Topic 9D). This positive contribution to the potential energy depends on the overlap of the atomic orbitals, and, because orbitals decay exponentially with distance, at large distances from the nucleus it is often modelled by writing


\begin{equation*}
E_{\mathrm{p}}^{*}=N_{\mathrm{A}} C^{\prime} \mathrm{e}^{-d / d^{*}} \tag{15C.4}
\end{equation*}


where $d$ is the distance between the atoms, and $C^{\prime}$ and $d^{*}$ are constants. It turns out that the value of $C^{\prime}$ is not needed (it cancels in expressions that make use of this formula; see below); $d^{*}$ is commonly taken to be 34.5 pm .\\
The total potential energy is the sum of $E_{\mathrm{p}}$ and $E_{\mathrm{p}}^{*}$, and passes through a minimum when $\mathrm{d}\left(E_{\mathrm{p}}+E_{\mathrm{p}}^{\star}\right) / \mathrm{d} d=0$ (Fig. 15C.13).

Table 15C. 3 Madelung constants

\begin{center}
\begin{tabular}{ll}
\hline
Structural type & $A$ \\
\hline
Caesium chloride & 1.763 \\
Fluorite & 2.519 \\
Rock salt & 1.748 \\
Rutile & 2.408 \\
Sphalerite (zinc blende) & 1.638 \\
Wurtzite & 1.641 \\
\hline
\end{tabular}
\end{center}

\begin{center}
\includegraphics[max width=\textwidth]{2025_02_27_d4b95d9718f0b9e2e6b9g-694(1)}
\end{center}

Figure 15C. 13 The contributions to the total potential energy of an ionic crystal.

A short calculation leads to the Born-Mayer equation for the minimum total potential energy (see Problem P15C.9):


\begin{equation*}
E_{\mathrm{p}, \text { min }}=-\frac{N_{\mathrm{A}}\left|z_{\mathrm{A}} z_{\mathrm{B}}\right| e^{2}}{4 \pi \varepsilon_{0} d}\left(1-\frac{d^{*}}{d}\right) A \quad \text { Born-Mayer equation } \tag{15C.5}
\end{equation*}


Provided zero-point contributions to the energy are ignored, the negative of this potential energy can be identified with the lattice energy. The important features of this equation are:

\begin{itemize}
  \item Because $E_{\mathrm{p}, \min } \propto\left|z_{\mathrm{A}} z_{\mathrm{B}}\right|$, the potential energy decreases (becomes more negative) with increasing charge number of the ions.
  \item Because the electrostatic (and dominant) contribution to $E_{\mathrm{p}, \text { min }}$ is proportional to $1 / d$, the potential energy decreases (becomes more negative) with decreasing ionic radius.
\end{itemize}

The second conclusion follows from the fact that the smaller the ionic radii, the smaller is the value of $d$. High lattice energies are expected when the ions are highly charged (so $\left|z_{\mathrm{A}} z_{\mathrm{B}}\right|$ is large) and small (so $d$ is small).

\subsection*{Brief illustration 15C. 3}
To estimate $E_{\mathrm{p}, \text { min }}$ for MgO , which has a rock salt structure ( $A=1.748$ ), use the following values: $d=r\left(\mathrm{Mg}^{2+}\right)+r\left(\mathrm{O}^{2-}\right)=72+$ $140 \mathrm{pm}=212 \mathrm{pm}$. Note that

$$
\begin{aligned}
\frac{N_{\mathrm{A}} e^{2}}{4 \pi \varepsilon_{0}} & =\frac{\left(6.02214 \times 10^{23} \mathrm{~mol}^{-1}\right) \times\left(1.602176 \times 10^{-19} \mathrm{C}\right)^{2}}{4 \pi \times\left(8.85419 \times 10^{-12} \mathrm{~J}^{-1} \mathrm{C}^{2} \mathrm{~m}^{-1}\right)} \\
& =1.38935 \times 10^{-4} \mathrm{Jm} \mathrm{~mol}^{-1}
\end{aligned}
$$

Then

$$
\begin{aligned}
E_{\mathrm{p}, \text { min }}= & -\underbrace{\frac{z_{\mathrm{Mg}^{2+} z_{\mathrm{o}^{2-}} \mid}^{4}}{2.12 \times 10^{-10} \mathrm{~m}} \times \overbrace{\left(1.38935 \times 10^{-4} \mathrm{Jm} \mathrm{~mol}^{-1}\right)}^{N_{\mathrm{A}} \mathrm{e}^{2} / 4 \pi \varepsilon_{0}}}_{d} \\
& \times \overbrace{\left(1-\frac{34.5 \mathrm{pm}}{212 \mathrm{pm}}\right)}^{1-d^{*} / d} \times \overbrace{1.748}^{A} \\
= & -3.84 \times 10^{3} \mathrm{kJmol}^{-1}
\end{aligned}
$$

It is not possible to measure the lattice enthalpy directly, but values can be obtained by combining experimental values of other enthalpy changes by using a Born-Haber cycle. Such a cycle is a closed path of transformations starting and ending at the same point, one step of which is the formation of the solid compound from a gas of widely separated ions.

\subsection*{Example 15C. 2 Using the Born-Haber cycle}
Calculate the lattice enthalpy of KCl .\\
Collect your thoughts You need to construct a suitable BornHaber cycle, such as the one shown in Fig. 15C.14. For the cycle to be useful, experimental values for the enthalpy changes for all the steps need to be available (for example, from tabulated data), apart, of course, from the step involving the formation of the lattice from the ions. For the cycle in Fig. 15C. 14 the enthalpy changes are (for convenience, starting at the elements):

\begin{center}
\begin{tabular}{|c|c|c|}
\hline
\multicolumn{3}{|c|}{$\Delta H /\left(\mathrm{kJ} \mathrm{mol}^{-1}\right)$} \\
\hline
1. Sublimation of $\mathrm{K}(\mathrm{s})$ & +89 & [dissociation enthalpy of $\mathrm{K}(\mathrm{s})$ ] \\
\hline
2. Dissociation of $\frac{1}{2} \mathrm{Cl}_{2}(\mathrm{~g})$ & +122 & [ $\frac{1}{2} \times$ dissociation enthalpy of $\left.\mathrm{Cl}_{2}(\mathrm{~g})\right]$ \\
\hline
3. Ionization of $\mathrm{K}(\mathrm{g})$ & +418 & [ionization enthalpy of $\mathrm{K}(\mathrm{g})$ ] \\
\hline
4. Electron attachment to $\mathrm{Cl}(\mathrm{g})$ & -349 & [electron-gain enthalpy of $\mathrm{Cl}(\mathrm{g})]$ \\
\hline
5. Formation of solid from gaseous ions & \( \begin{gathered} -\Delta H_{\mathrm{L}}^{\mathrm{L} /(\mathrm{kJ}} \\ \left.\mathrm{mol}^{-1}\right) \end{gathered} \) & [value to be determined] \\
\hline
6. Decomposition of compound to its elements in their reference states & +437 & [negative of enthalpy of formation of $\mathrm{KCl}(\mathrm{s})$ ] \\
\hline
\end{tabular}
\end{center}

Because this is a closed cycle, the sum of these enthalpy changes is equal to zero, and the lattice enthalpy can be inferred from the resulting equation.

The solution The sum of contributions around the cycle is

$$
89+122+418-349-\Delta H_{\mathrm{L}} /\left(\mathrm{kJ} \mathrm{~mol}^{-1}\right)+437=0
$$

It follows that $\Delta H_{\mathrm{L}}=+717 \mathrm{~kJ} \mathrm{~mol}^{-1}$.\\
\includegraphics[max width=\textwidth, center]{2025_02_27_d4b95d9718f0b9e2e6b9g-695(1)}

Figure 15C. 14 The Born-Haber cycle for KCl at 298 K . Enthalpy changes are in kilojoules per mole.

Self-test 15C. 2 Calculate the lattice enthalpy of CaO from the following data:

\begin{center}
\begin{tabular}{lc}
\hline
 & $\Delta H /\left(\mathrm{kJ} \mathrm{mol}^{-1}\right)$ \\
\hline
Sublimation of $\mathrm{Ca}(\mathrm{s})$ & +178 \\
Ionization of $\mathrm{Ca}(\mathrm{g})$ to $\mathrm{Ca}^{2+}(\mathrm{g})$ & +1735 \\
Dissociation of $\frac{1}{2} \mathrm{O}_{2}(\mathrm{~g})$ & +249 \\
Electron attachment to $\mathrm{O}(\mathrm{g})$ & -141 \\
Electron attachment to $\mathrm{O}^{-}(\mathrm{g})$ & +844 \\
Formation of $\mathrm{CaO}(\mathrm{s})$ from $\mathrm{Ca}(\mathrm{s})$ and $\frac{1}{2} \mathrm{O}_{2}(\mathrm{~g})$ & -635 \\
\hline
\end{tabular}
\end{center}

\begin{center}
\includegraphics[max width=\textwidth]{2025_02_27_d4b95d9718f0b9e2e6b9g-695}
\end{center}

Some lattice enthalpies obtained by the Born-Haber cycle are listed in Table 15C.4. As can be seen from the data, the trends in values are in general accord with the predictions of the Born-Mayer equation. Agreement is typically taken to imply that the ionic model of bonding is valid for the substance; disagreement implies that there is a covalent contribution to the bonding. It is important, though, to be cautious, because numerical agreement might be coincidental and, as noted above, the values for ionic radii are subject to significant uncertainty.

\subsection*{15C.3 Covalent and molecular solids}
X-ray diffraction studies of solids reveal a huge amount of information, including interatomic distances, bond angles, stereochemistry, and vibrational parameters. This section can do no more than hint at the diversity of types of solids found when molecules pack together or atoms link together in extended networks.

In covalent solids (or covalent network solids), covalent bonds in a definite spatial orientation link the atoms together into a network that extends through the crystal; effectively the crystal is a giant molecule. The demands of directional bonding, which have only a small effect on the structures of many metals, now override the geometrical problem of packing spheres together, resulting in a wide variety of often quite elaborate structures.

Table 15C. 4 Lattice enthalpies at $298 \mathrm{~K}, \Delta H_{\mathrm{L}} /\left(\mathrm{kJ} \mathrm{mol}^{-1}\right)^{*}$

\begin{center}
\begin{tabular}{lr}
\hline
NaF & 926 \\
NaBr & 751 \\
MgO & 3850 \\
MgS & 3406 \\
\hline
\end{tabular}
\end{center}

\footnotetext{\begin{itemize}
  \item More values are given in the Resource section
\end{itemize}
}\subsection*{Brief illustration 15C. 4}
Diamond and graphite are two allotropes of carbon. In diamond each $\mathrm{sp}^{3}$-hybridized carbon is bonded tetrahedrally to its four neighbours (Fig. 15C.15). The network of strong C-C bonds is repeated throughout the crystal and, as a result, diamond is very hard (in fact, the hardest known substance). In graphite, $\sigma$ bonds between $\mathrm{sp}^{2}$-hybridized carbon atoms form hexagonal rings which, when repeated throughout a plane, give rise to 'graphene' sheets (Fig. 15C.16). Because the sheets can slide against each other when impurities are present, impure graphite is used widely as a lubricant.\\
\includegraphics[max width=\textwidth, center]{2025_02_27_d4b95d9718f0b9e2e6b9g-696(2)}

Figure 15C. 15 A fragment of the structure of diamond. Each $C$ atom is tetrahedrally bonded to four neighbours. This framework-like structure results in a rigid crystal.\\
\includegraphics[max width=\textwidth, center]{2025_02_27_d4b95d9718f0b9e2e6b9g-696(1)}\\
(a)\\
\includegraphics[max width=\textwidth, center]{2025_02_27_d4b95d9718f0b9e2e6b9g-696}\\
(b)

Figure 15C. 16 Graphite consists of flat planes of hexagons of carbon atoms lying above one another. (a) The arrangement of carbon atoms in a 'graphene' sheet; (b) the relative arrangement of neighbouring sheets. The planes can slide over one another easily when impurities are present.

Molecular solids, which are the subject of the overwhelming majority of modern structural determinations, are held together by van der Waals interactions between the individual molecular components (Topic 14B). The observed crystal structure is nature's solution to the problem of condensing objects of various shapes into an aggregate of minimum energy (actually, for $T>0$, of minimum Gibbs energy). The prediction of the structure is difficult, but software specifically designed to explore interaction energies can now make reasonably reliable predictions. The problem is made more complicated by the role of hydrogen bonds, which in some cases dominate the crystal structure, as in ice (Fig. 15C.17), but in others (for example, in solid phenol) distort a structure that is determined largely by the van der Waals interactions.\\
\includegraphics[max width=\textwidth, center]{2025_02_27_d4b95d9718f0b9e2e6b9g-696(3)}

Figure 15C. 17 A fragment of the crystal structure of ice (ice-I). Each O atom is at the centre of a tetrahedron of four O atoms at a distance of 276 pm . The central O atom is attached by two short $\mathrm{O}-\mathrm{H}$ bonds to two H atoms and by two long hydrogen bonds to the H atoms of two of the neighbouring molecules. Both alternative H atoms locations are shown for each $\mathrm{O}-\mathrm{O}$ separation. Overall, the structure consists of planes of hexagonal puckered rings of $\mathrm{H}_{2} \mathrm{O}$ molecules (like the chair form of cyclohexane).

\subsection*{Checklist of concepts}
1. A close-packed layer is a layer of spheres arranged so there is maximum utilization of space.2. A hexagonally close-packed structure is one where the sequence of close-packed layers is ABABAB....3. A cubic close-packed structure is one where the sequence of close-packed layers is $\mathrm{ABCABC} . .$.\begin{enumerate}
  \setcounter{enumi}{3}
  \item The coordination number is the number of spheres immediately surrounding any selected sphere.
  \item In the nearly free-electron approximation the valence electrons are assumed to be trapped in a box with a periodic potential energy, with low energy corresponding to the locations of cations.6. In the tight-binding approximation the valence electrons are assumed to occupy molecular orbitals delocalized throughout the solid.7. In metals atomic orbitals overlap to form a band, which is a set of molecular orbitals that are closely spaced and cover a finite range of energy; electrons occupy the orbitals within the band.8. A band gap is a range of energies to which no orbital corresponds.9. The Fermi level is the highest occupied molecular orbital at $T=0$.10. The coordination number of an ionic lattice is denoted $\left(N_{+}, N_{-}\right)$, where $N_{+}$is the number of nearest neighbour\\
anions around a cation and $N_{-}$the number of nearest neighbour cations around an anion.11. The lattice energy of a solid is the change in potential energy when the ions go from being packed together in a solid to being widely separated as a gas.12. A Born-Haber cycle is a closed path of transformations starting and ending at the same point, one step of which is the formation of the solid compound from a gas of widely separated ions.13. A molecular solid is a solid consisting of discrete molecules held together by van der Waals interactions, and possibly hydrogen bonds.
\end{enumerate}

\subsection*{Checklist of equations}
\begin{center}
\begin{tabular}{llll}
\hline
Property & Equation & Comment & Equation number \\
\hline
Energy levels of a linear array of orbitals & $E_{k}=\alpha+2 \beta \cos (k \pi /(N+1)), k=1,2, \ldots, N$ & Hückel approximation & $15 C .1$ \\
Band width & $E_{N}-E_{1} \rightarrow-4 \beta$ as $N \rightarrow \infty$ & Hückel approximation &  \\
Radius ratio & $\gamma=r_{\text {smaller }} / r_{\text {larger }}$ & For criteria, see Section 15C.2 & 15 C .2 \\
Born-Mayer equation & $E_{\mathrm{p}, \text { min }}=-\left\{N_{\mathrm{A}}\left|z_{\mathrm{A}} z_{\mathrm{B}}\right| e^{2} / 4 \pi \varepsilon_{0} d\right\}\left(1-d^{*} / d\right) A$ &  & 15 C .5 \\
\hline
\end{tabular}
\end{center}

\section*{TOPIC 15D The mechanical properties of solids }
\subsection*{Why do you need to know this material?}
An understanding of the mechanical properties of solid materials is crucial to the development of modern materials.

\subsection*{What is the key idea?}
The mechanical properties of solids are expressed in terms of various 'moduli' that are related to the intermolecular potential energy of the constituents.

What do you need to know already?\\
You need to be familiar with the Lennard-Jones potential energy (Topic 14B).

The fundamental concepts needed for the discussion of the mechanical properties of solids are stress and strain. The stress on an object is the applied force divided by the area to which the force is applied. For instance, if a mass $m$ hangs from a wire of radius $r$, and therefore cross-section $\pi r^{2}$, the mass exerts a gravitational force $m g$ and the uniaxial stress (along the length of the wire) would be reported as $m g / \pi r^{2}$. The strain is the resulting fractional distortion of the object. The study of the relation between stress and strain is called rheology from the Greek word for 'flow'.\\
Stress may be applied in a number of different ways (Fig. 15D.1):

\begin{itemize}
  \item Uniaxial stress is a simple compression or extension in one direction.\\
\includegraphics[max width=\textwidth, center]{2025_02_27_d4b95d9718f0b9e2e6b9g-698}
\end{itemize}

Figure 15D. 1 Types of stress applied to a body. (a) Uniaxial stress, (b) shear stress, (c) hydrostatic pressure.

\begin{itemize}
  \item Hydrostatic stress is a stress applied simultaneously in all directions, as in a body immersed in a fluid.
  \item Pure shear is a stress that tends to push opposite faces of the sample in opposite directions.
\end{itemize}

A sample subjected to a low stress typically undergoes elastic deformation in the sense that it recovers its original shape when the stress is removed. For low stresses, the strain is linearly proportional to the stress, and the stress-strain relation is a Hooke's law of force (Fig. 15D.2). The response becomes nonlinear at high stresses but may remain elastic. Above a certain threshold, the strain becomes plastic in the sense that recovery does not occur when the stress is removed. Plastic deformation occurs when bond breaking takes place and, in pure metals, typically takes place through the agency of dislocations. Brittle materials, such as ionic solids, exhibit sudden fracture as the stress focused by cracks causes them to spread catastrophically.

The response of a solid to an applied stress is commonly summarized by a number of coefficients of proportionality known as moduli:


\begin{align*}
& \text { Young's modulus: } E=\frac{\text { uniaxial stress }}{\text { uniaxial strain }}  \tag{15D.1a}\\
& \text { Bulk modulus: } K=\frac{\text { pressure }}{\text { fractional change in }}  \tag{15D.1b}\\
& \text { Shear modulus: } G=\frac{\text { shear stress }}{\text { shear strain }}
\end{align*}


\begin{center}
\includegraphics[max width=\textwidth]{2025_02_27_d4b95d9718f0b9e2e6b9g-698(1)}
\end{center}

Figure 15D. 2 At small strains, a body obeys Hooke's law (stress proportional to strain) and is elastic (recovers its shape when the stress is removed). At high strains, the body is no longer elastic, may become plastic, and finally yield.\\
\includegraphics[max width=\textwidth, center]{2025_02_27_d4b95d9718f0b9e2e6b9g-699}

Figure 15D. 3 (a) Uniaxial stress and the resulting uniaxial and transverse strain; Poisson's ratio indicates the extent to which a body changes shape when subjected to a uniaxial stress. (b) Shear stress and the resulting strain.\\
'Uniaxial strain' refers to stretching and compression of the material in one direction, as shown in Fig. 15D.3a, and 'shear strain' refers to the distortion arising from a shear stress, as depicted in Fig. 15D.3b. The fractional change in volume is $\delta V / V$, where $\delta V$ is the change in volume of a sample of volume $V$; similarly, the uniaxial strain and the shear strain are (dimensionless) fractional changes in dimensions. The bulk modulus is the inverse of the isothermal compressibility, $\kappa_{T}$, discussed in Topic 2D (eqn 2D.7, $\left.\kappa_{T}=-(\partial V / \partial p)_{T} / V\right)$.

A third ratio, called Poisson's ratio, indicates how the sample changes its shape:


\begin{equation*}
v_{\mathrm{P}}=\frac{\text { transverse strain }}{\text { normal strain }} \tag{15D.2}
\end{equation*}


Poisson's ratio [definition]

Transverse and normal strains are illustrated in Fig. 15D.3a: they are the mutually perpendicular uniaxial distortions arising from the 'normal' uniaxial stress. The three moduli introduced in eqn 15D. 1 are interrelated in the following way (see Problem P15D.1):

\[
G=\frac{E}{2\left(1+v_{\mathrm{p}}\right)} \quad K=\frac{E}{3\left(1-2 v_{\mathrm{P}}\right)} \quad \begin{align*}
& \text { Relations between }  \tag{15D.3}\\
& \text { moduli }
\end{align*}
\]

\subsection*{Brief illustration 15D. 1}
The uniaxial stress when a mass of $m=10.0 \mathrm{~kg}$ is suspended from an iron wire of radius $r=0.050 \mathrm{~mm}$ is

$$
\begin{aligned}
\text { uniaxial stress } & =\frac{m g}{\pi r^{2}}=\frac{(10.0 \mathrm{~kg}) \times\left(9.81 \mathrm{~m} \mathrm{~s}^{-2}\right)}{\pi\left(5.0 \times 10^{-5} \mathrm{~m}\right)^{2}} \\
& =1.24 \ldots \times 10^{10} \mathrm{~kg} \mathrm{~m}^{-1} \mathrm{~s}^{-2}
\end{aligned}
$$

The Young's modulus of iron at room temperature is 215 GPa . Therefore

$$
\text { uniaxial strain }=\frac{1.24 \ldots \times 10^{10} \mathrm{~kg} \mathrm{~m}^{-1} \mathrm{~s}^{-2}}{2.15 \times 10^{11} \underbrace{\mathrm{~kg} \mathrm{~m}^{-1} \mathrm{~s}^{-2}}_{\mathrm{Pa}}}=0.0581
$$

which corresponds to elongation of the wire by 5.81 per cent.

If neighbouring molecules interact by a Lennard-Jones potential energy (Topic 14B), then the bulk modulus and the compressibility of the solid are related to the Lennard-Jones parameter $\varepsilon$ (the depth of the potential well) by


\begin{equation*}
K=\frac{8 N_{\mathrm{A}} \varepsilon}{V_{\mathrm{m}}} \quad \kappa_{T}=\frac{V_{\mathrm{m}}}{8 N_{\mathrm{A}} \varepsilon} \tag{15D.4}
\end{equation*}


For the derivation of these relations, see $A$ deeper look 10 on the website of this text. The bulk modulus is large and the compressibility low (the solid stiff) if the potential well is deep and the solid is dense (its molar volume small).

The differing rheological characteristics of metals can be traced to the presence of slip planes, which are planes of atoms that, when under stress, may slip or slide relative to one another. The slip planes of a ccp structure are the close-packed planes, and careful inspection of a unit cell shows that there are eight sets of slip planes in different directions. As a result, metals with ccp structures, like copper, are malleable, meaning they can easily be bent, flattened, or hammered into shape. In contrast, a hexagonal close-packed structure has only one set of slip planes so that metals with hexagonal close packing, such as zinc or cadmium, tend to be more brittle.

\subsection*{Checklist of concepts}
\begin{enumerate}
  \item Uniaxial stress is a simple compression or extension applied to a sample in one direction.
  \item Hydrostatic stress is a stress applied simultaneously in all directions, as in a body immersed in a fluid.3. Pure shear is a stress that tends to push opposite faces of the sample in opposite directions.4. A sample subjected to a small stress typically undergoes elastic deformation; as the stress increases the sample becomes plastic.
  \item The response of a solid to an applied stress is summarized by the Young's modulus, the bulk modulus, the shear modulus, and Poisson's ratio.
  \item The differing rheological characteristics of metals can be traced to the presence of slip planes.
\end{enumerate}

\subsection*{Checklist of equations}
\begin{center}
\begin{tabular}{llll}
\hline
Property & Equation & Comment & Equation number \\
\hline
Young's modulus & $E=$ uniaxial stress/uniaxial strain & Definition & 15D.1a \\
Bulk modulus & $K=$ pressure/fractional change in volume & Definition & $15 D .1 \mathrm{~b}$ \\
Shear modulus & $G=$ shear stress/shear strain & Definition & $15 D .1 \mathrm{c}$ \\
Poisson's ratio & $v_{\mathrm{P}}=$ transverse strain/normal strain & Definition & 15 D .2 \\
\hline
\end{tabular}
\end{center}

\section*{TOPIC 15E The electrical properties of solids }
\subsection*{Why do you need to know this material?}
The electrical properties of solids underlie numerous technological applications on which the infrastructure of the modern world depends.

\subsection*{What is the key idea?}
Electrons in solids occupy bands that determine the electrical conductivities of various types of solid.

\subsection*{What do you need to know already?}
You need to be familiar with the formation of bands in solids (Topic 15C).

The electrical conductivity of common materials arises from the motion of electrons, but some ionic solids display ionic conductivity in which whole ions migrate through the lattice. Three types of solid are distinguished by the temperature dependence of their electrical conductivity (Fig. 15E.1):

\begin{itemize}
  \item A metallic conductor is a substance with an electrical conductivity that decreases as the temperature is raised.\\
\includegraphics[max width=\textwidth, center]{2025_02_27_d4b95d9718f0b9e2e6b9g-701(1)}
\end{itemize}

Figure 15E. 1 The variation of the electrical conductivity of a substance with temperature is the basis of its classification as a metallic conductor, a semiconductor, or a superconductor. Conductivity is expressed in siemens per metre ( $\mathrm{S} \mathrm{m}^{-1}$ or, as here, $\mathrm{Scm}^{-1}$ ), where $1 \mathrm{~S}=1 \Omega^{-1}$ (the resistance is expressed in ohms, $\Omega$ ); note the log scale.

\begin{itemize}
  \item A semiconductor is a substance with an electrical conductivity that increases as the temperature is raised.
  \item A superconductor is a solid that, below a critical temperature, conducts electricity without resistance.
\end{itemize}

A semiconductor generally has a lower conductivity than that typical of metallic conductors, but the magnitude of the conductivity is not the criterion for distinguishing between them. It is conventional to classify semiconductors with very low electrical conductivities, such as most synthetic polymers, as insulators. The term is one of convenience rather than one of fundamental significance.

\subsection*{15E. 1 Metallic conductors}
To understand the origins of the electric conductivity in conductors and semiconductors, it is necessary to explore the consequences of the formation of bands (Topic 15C). The starting point is Fig. 15C.8, which is repeated here for convenience as Fig. 15E.2. It shows the electronic structure of a solid formed from a line of $N$ atoms, each of which contributes one electron (such as the alkali metals). At $T=0$, only the lowest $\frac{1}{2} N$ molecular orbitals are occupied, up to the Fermi level. The levels are very closely spaced, so there are unoccupied molecular orbitals just above the Fermi level. A solid with a partially filled band is expected to be a metallic conductor, an observation which can be understood in the following way.

The key point is that each molecular orbital in a band can be regarded as the superposition of two waves travelling in\\
\includegraphics[max width=\textwidth, center]{2025_02_27_d4b95d9718f0b9e2e6b9g-701}

Figure 15E. 2 (A reproduction of Fig. 15C.8.) When $N$ electrons occupy a band of $N$ orbitals at $T=0$, it is only half full. The highest occupied level is the Fermi level.\\
\includegraphics[max width=\textwidth, center]{2025_02_27_d4b95d9718f0b9e2e6b9g-702}

Figure 15E. 3 Each part of this illustration separates out the waves travelling in opposite directions. (a) The two sets of waves have the same energy, are equally occupied, and there is no net motion. (b) When a potential difference is applied (positive on the right) the two sets no longer have the same energy. There are now more electrons moving to the right than to the left and therefore there is a net current. (c) If the band is full, the two sets remain equally populated and there is no net flow even when a potential difference is applied.\\
opposite directions (in the same sense that $\cos x \propto \mathrm{e}^{\mathrm{i} x}+\mathrm{e}^{-\mathrm{i} x}$ ). Figure 15E.3a is an adaptation of Fig. 15E. 2 which separates the two contributions. In the absence of any applied field, electrons occupy both components equally and there is no net motion through the solid. This absence of net motion is true whether the band is full or incomplete. When a potential difference is applied the energies of the components differ, as it is energetically favourable for electrons to travel towards regions of positive potential. Now the two components are no longer equally occupied (Fig. 15E.3b) and provided the band is not full there are more electrons travelling in one direction than the other and electrical conduction occurs. If the band is full, however, the populations of the two components remain equal (Fig. 15E.3c) and there is no net motion in either direction. The material does not conduct: it is an insulator.

The conductivity of a metallic conductor decreases as the temperature is raised. This decrease is due to collisions between the moving electrons and the atoms. The greater the temperature, the more vigorous is the thermal motion of the atoms, so collisions between the moving electrons and an atom are more likely. That is, the electrons are scattered out of their paths through the solid, and are less efficient at transporting charge.

\subsection*{15E. 2 Insulators and semiconductors}
Now consider a solid which has an arrangement of bands as shown in Fig. 15E.4. At $T=0$ the lower band is full, and the Fermi level lies at the top of the band. A second empty band lies at a higher energy, separated from the top of the lower band by an energy $E_{\mathrm{g}}$, known as the band gap. At $T=0$ this material is an insulator because there are no partially filled bands. If the temperature is high enough, though, electrons\\
\includegraphics[max width=\textwidth, center]{2025_02_27_d4b95d9718f0b9e2e6b9g-702(1)}

Figure 15E. 4 (a) Typical band structure for a semiconductor: at $T=0$ the valence band is full and the conduction band is empty.\\
(b) At higher temperatures electrons populate the levels of the conduction band leading to electrical conductivity which increases with temperature.\\
are excited out of the lower band into the upper. There are now incomplete bands and conduction can occur.

The question that now arises is how the populations of the two bands, and therefore the conductivity of the semiconducting material, depend on the temperature. The discussion starts by introducing the density of states, $\rho(E)$, defined such that the number of states between $E$ and $E+\mathrm{d} E$ is $\rho(E) \mathrm{d} E$. Note that the 'state' of an electron includes its spin, so each spatial orbital counts as two states. To obtain the number of electrons $\mathrm{d} N(E)$ that occupy states between $E$ and $E+\mathrm{d} E, \rho(E) \mathrm{d} E$ is multiplied by the probability $f(E)$ of occupation of the state with energy $E$. That is,

$$
\mathrm{d} N(E)=\overbrace{\rho(E) \mathrm{d} E}^{\begin{array}{c}
\text { Number of } \\
\text { states } \\
\text { between }
\end{array}} \times \begin{gathered}
\begin{array}{c}
\text { Probability of } \\
\text { occupation of } \\
\text { a state with } \\
\text { energy } E
\end{array}
\end{gathered}
$$

The function $f(E)$ is the Fermi-Dirac distribution, a version of the Boltzmann distribution that takes into account the Pauli exclusion principle, that each orbital can be occupied by no more than two electrons:


\begin{equation*}
f(E)=\frac{1}{\mathrm{e}^{(E-\mu) / k T}+1} \quad \text { Fermi-Dirac distribution } \tag{15E.2a}
\end{equation*}


In this expression $\mu$ is a temperature-dependent parameter known as the 'chemical potential' (it has a subtle relation to the familiar chemical potential of thermodynamics); provided $T>0, \mu$ is the energy of the state for which $f=\frac{1}{2}$. At $T=0$, only states up to a certain energy, known as the Fermi energy, $E_{\mathrm{F}}$, are occupied (Fig. 15E.2). Provided the temperature is not so high that many electrons are excited to states above the Fermi\\
energy, the chemical potential can be identified with $E_{\mathrm{F}}$, in which case the Fermi-Dirac distribution becomes

$$
f(E)=\frac{1}{\mathrm{e}^{\left(E-E_{\mathrm{F}}\right) / k T}+1}
$$

Fermi-Dirac distribution\\
(15E.2b)

This expression implies that $f\left(E_{\mathrm{F}}\right)=\frac{1}{2}$. For energies well above $E_{\mathrm{F}}$, the exponential term is so large that the 1 in the denominator can be neglected, and then

$$
f(E) \approx \mathrm{e}^{-\left(E-E_{\mathrm{F}}\right) / k T}
$$

Fermi-Dirac distribution [approximate form for $E>E_{F}$ ]\\
(15E.2c)

The function now resembles a Boltzmann distribution, decaying exponentially with increasing energy; the higher the temperature, the longer is the exponential tail.

There is a distinction between the Fermi energy and the Fermi level:

\begin{itemize}
  \item The Fermi level is the uppermost occupied level at $T=0$.
  \item The Fermi energy is the energy level at which $f(E)=\frac{1}{2}$ at any temperature.
\end{itemize}

The Fermi energy coincides with the Fermi level as $T \rightarrow 0$.\\
Figure 15E. 5 shows the form of $f(E)$ at different temperatures. At $T=0$ the probability distribution is a step function, equal to 1 for $E<E_{\mathrm{F}}$, and 0 at higher energies, as in Fig. 15E.2. At higher temperatures the probability of occupation of levels above $E_{\mathrm{F}}$ increases at the expense of those below $E_{\mathrm{F}}$, with the greatest changes occurring in the energies close to $E_{\mathrm{F}}$. As the temperature is raised, electrons are promoted from the lower band to the upper. This promotion is represented by the tail of the Fermi-Dirac distribution extending across the band gap and is significant only when $k T$ is comparable to or greater\\
\includegraphics[max width=\textwidth, center]{2025_02_27_d4b95d9718f0b9e2e6b9g-703}

Figure 15E. 5 The Fermi-Dirac distribution, which gives the probability of occupation of a state with energy $E$ and at a temperature $T$. At higher energies the probability decays exponentially towards zero. The curves are labelled with the value of $E_{\mathrm{F}} / k T$. The tinted region shows the occupation of levels at $T=0$.\\
than the band gap. The material, an insulator at $T=0$, is now a conductor, because both bands are partially filled. As the temperature is increased, the conductivity increases as more electrons are promoted across the band gap, so the material is a semiconductor.

The lower band, which is full at $T=0$, is called the valence band and the upper band, which is empty at $T=0$ and to which electrons are thermally excited, is called the conduction band. When electrons leave the valence band they can be thought of as creating positively charged 'holes' in that band (i.e. the absence of an electron), and the electrical conductivity arises from the movement of these holes and the promoted electrons.

Figure 15E. 4 depicts an intrinsic semiconductor, in which semiconduction is a property of the band structure of the pure material. Examples of intrinsic semiconductors include silicon and germanium. A compound semiconductor is an intrinsic semiconductor formed from a combination of different elements, such as $\mathrm{GaN}, \mathrm{CdS}$, and many d-metal oxides.

An extrinsic semiconductor is one in which charge carriers (electrons or holes) are present as a result of the replacement of some atoms (to the extent of about 1 in $10^{9}$ ) by dopant atoms, the atoms of another element. If, for example, pure silicon (a Group 14 element) is doped with atoms of indium (a Group 13 element) an electron can be transferred from a Si atom to a neighbouring In atom, thereby creating a hole in the valence band and increasing the conductivity. The resulting semiconductor is described as p-type, the p indicating that the positive holes are responsible for conduction. Figure 15E.6a shows the band structure of such a semiconductor. The dopant atoms result in a set of empty levels, called acceptor levels, which lie just above the top of the valence band. Electrons from the valance band are transferred into these levels, so generating holes in the band.\\
\includegraphics[max width=\textwidth, center]{2025_02_27_d4b95d9718f0b9e2e6b9g-703(1)}

Figure 15E. 6 (a) A dopant with fewer electrons than its host contributes levels that accept electrons from the valence band. The resulting holes in the band give rise to electrical conductivity; the doped semiconductor is classified as p-type. (b) A dopant with more electrons than its host contributes occupied levels that can supply electrons to the conduction band, thus giving rise to electrical conductivity; the substance is classed as an n-type semiconductor.

If the dopant atoms are from a Group 15 element (e.g. phosphorus), an electron can be transferred from a P atom into the otherwise empty conduction band, thereby increasing the conductivity. This type of doping results in an n-type semiconductor, where n refers to the negative charge of the carriers. The band structure is shown in Fig. 15E.6b. The dopant atoms create a set of occupied levels, called donor levels, just below the bottom of the conduction band, and electrons from these levels are transferred into the conduction band. In practical cases the level of doping is such that the charge carriers created by the dopant atoms are greatly in excess of those arising from thermal excitation across the band gap: electrical conductivity is therefore dominated by the type and extent of the doping.

Doped semiconductors are of great technological importance because they are the materials from which the active components of electronic circuits are made. The simplest example of an electronic device constructed from doped semiconductors is the ' $\mathrm{p}-\mathrm{n}$ diode' which consists of a p -type semiconductor in contact with an n-type semiconductor, thereby creating a $\mathbf{p}-\mathbf{n}$ junction. A p-n junction conducts electricity only in one direction. To understand this property consider first the arrangement shown in Fig. 15E.7a where the p-type semiconductor is attached to the negative electrode and the n-type is attached to the positive electrode; this arrangement is known as 'reverse bias'. The positively charged holes in the p-type semiconductor are attracted to the negative electrode, and the negatively charged electrons in the n-type semiconductor are attracted to the positive electrode. As a consequence, charge does not flow across the junction so the device does not conduct. Now consider what happens when the charges on the electrodes are reversed, as shown in Fig. 15E.7b, an arrangement known as 'forward bias'. Electrons in the n-type semiconductor move towards the positive electrode, and holes move in the opposite direction: as a result charge flows across the junction. The $\mathrm{p}-\mathrm{n}$ junction therefore conducts only under forward bias.

As electrons and holes move across a $\mathrm{p}-\mathrm{n}$ junction under forward bias, they recombine and release energy. However, as long as the forward bias persists, the flow of charge from the electrodes to the junction replenishes them with electrons and holes. In some solids, the energy of electron-hole recombination is released as heat and the device becomes warm. The\\
\includegraphics[max width=\textwidth, center]{2025_02_27_d4b95d9718f0b9e2e6b9g-704}

Figure 15E. 7 A p-n junction under (a) reverse bias, (b) forward bias.\\
reason lies in the fact that the return of the electron to a hole involves a change in the electron's linear momentum, which the atoms of the lattice must absorb, and therefore electronhole recombination stimulates lattice vibrations. This is the case for silicon semiconductors, and is one reason why computers need efficient cooling systems.

Another electronic device, a 'transistor', consists of a p-type semiconductor sandwiched between two n-type semiconductors, and as such has two $\mathrm{p}-\mathrm{n}$ junctions. Under the correct conditions it is possible to control the current flowing between the two n-type semiconductors by varying the current flowing into the p-type semiconductor. Most significantly, the change in the current between the n-type semiconductors can be larger than the change in the current in the p-type semiconductor; in other words the device can act as an amplifier. It is the exploitation of this property that has led to the development of modern solid-state electronics.

\subsection*{15E. 3 Superconductors}
The resistance to flow of electrical current of a normal metallic conductor decreases smoothly with decreasing temperature but never vanishes. However, a superconductor conducts electricity without resistance once the temperature is below the critical temperature, $T_{c}$. Following the discovery in 1911 that mercury is a superconductor below 4.2 K , the normal boiling point of liquid helium, physicists and chemists made slow but steady progress in the discovery of superconductors with higher values of $T_{c}$. Metals, such as tungsten, mercury, and lead, have $T_{c}$ values below about 10 K . Intermetallic compounds, such as $\mathrm{Nb}_{3} \mathrm{X}(\mathrm{X}=\mathrm{Sn}, \mathrm{Al}$, or Ge$)$, and alloys, such as $\mathrm{Nb} / \mathrm{Ti}$ and $\mathrm{Nb} / \mathrm{Zr}$, have intermediate $T_{\mathrm{c}}$ values ranging between 10 K and 23 K . In 1986, high-temperature superconductors (HTSCs) were discovered. Several ceramics, inorganic powders that have been fused and hardened by heating to a high temperature, containing oxocuprate motifs, $\mathrm{Cu}_{m} \mathrm{O}_{n}$, are now known with $T_{\mathrm{c}}$ values well above 77 K , the boiling point of the inexpensive refrigerant liquid nitrogen. For example, $\mathrm{HgBa}_{2} \mathrm{Ca}_{2} \mathrm{Cu}_{3} \mathrm{O}_{8}$ has $T_{\mathrm{c}}=153 \mathrm{~K}$.

The elements that exhibit superconductivity cluster in certain parts of the periodic table. The metals iron, cobalt, nickel, copper, silver, and gold do not display superconductivity; nor do the alkali metals. One of the most widely studied oxocuprate superconductors $\mathrm{YBa}_{2} \mathrm{Cu}_{3} \mathrm{O}_{7}$ (informally known as ' 123 ' on account of the proportions of the metal atoms in the compound) has the structure shown in Fig. 15E.8. The square-pyramidal $\mathrm{CuO}_{5}$ units arranged as two-dimensional layers and the square planar $\mathrm{CuO}_{4}$ units arranged in sheets are common structural features of oxocuprate HTSCs.

The mechanism of superconduction is well-understood for low-temperature materials, and is based on the properties of\\
\includegraphics[max width=\textwidth, center]{2025_02_27_d4b95d9718f0b9e2e6b9g-705(1)}

Figure 15E. 8 Structure of the $\mathrm{YBa}_{2} \mathrm{Cu}_{3} \mathrm{O}_{7}$ superconductor. (a) Metal atom positions. (b) The polyhedra show the positions of oxygen atoms and indicate that the Cu ions are either in squareplanar or square-pyramidal coordination environments.\\
a Cooper pair, a pair of electrons that exists on account of the indirect electron-electron interactions mediated by the nuclei of the atoms in the lattice. Thus, if one electron is in a particular region of a solid, the nuclei there move towards it to give a distorted local structure (Fig. 15E.9). Because that local distortion is rich in positive charge, it is favourable for a second electron to join the first. Hence, there is a virtual attraction between the two electrons and they move together as a pair. The local distortion is disrupted by thermal motion of the ions in the solid, so the virtual attraction occurs only at very low temperatures. A Cooper pair undergoes less scattering than an individual electron as it travels through the solid because the distortion caused by one electron can attract back\\
\includegraphics[max width=\textwidth, center]{2025_02_27_d4b95d9718f0b9e2e6b9g-705}

Figure 15E. 9 The formation of a Cooper pair. One electron distorts the crystal lattice and the second electron has a lower energy if it goes to that region. These electron-lattice interactions effectively bind the two electrons into a pair.\\
the other electron should it be scattered out of its path in a collision. Because the Cooper pair is stable against scattering, it can carry charge freely through the solid, and hence give rise to superconduction.

The Cooper pairs responsible for low-temperature superconductivity are likely to be important in HTSCs, but the mechanism for pairing is hotly debated. There is evidence implicating the arrangement of $\mathrm{CuO}_{5}$ layers and $\mathrm{CuO}_{4}$ sheets in the mechanism. It is believed that movement of electrons along the linked $\mathrm{CuO}_{4}$ units accounts for superconductivity, whereas the linked $\mathrm{CuO}_{5}$ units act as 'charge reservoirs' that maintain an appropriate number of electrons in the superconducting layers.

\subsection*{Checklist of concepts}
$\square$ 1. Electronic conductors are classified as metallic conductors or semiconductors according to the temperature dependence of their conductivities; an insulator is a semiconductor with very low conductivity.2. Superconductors conduct electricity without resistance below a critical temperature $T_{c}$.\\
3. The Fermi-Dirac distribution gives the probability that a state with a particular energy is occupied by an electron.4. The Fermi energy is the energy of the level for which the probability of occupation is $\frac{1}{2}$.5. In a semiconductor at $T=0$ there is a full valence band and, at higher energy, an empty conduction band.6. In an intrinsic semiconductor electrical conductivity is due to electrons thermally promoted from the valence band to the conduction band.7. In an extrinsic semiconductor electrical conductivity is due to electrons or holes generated by the inclusion of dopant atoms.8. Semiconductors are classified as p-type or n-type according to whether conduction is due to holes in the valence band or electrons in the conduction band.

\subsection*{Checklist of equations}
\begin{center}
\begin{tabular}{llll}
\hline
Property & Equation & Comment & Equation number \\
\hline
Fermi-Dirac distribution & $f(E)=1 /\left\{\mathrm{e}^{\left(E-E_{\mathrm{F}}\right) / k T}+1\right\}$ & $E_{\mathrm{F}}$ is the Fermi energy & 15 E .2 b \\
\hline
\end{tabular}
\end{center}

\section*{TOPIC 15F The magnetic properties of solids }
\subsection*{Why do you need to know this material?}
The magnetic properties of solids give an indication of the electronic structures of individual molecules and many modern information storage devices make use of the additional properties that arise when the spins on different centres interact.

\subsection*{What is the key idea?}
The principal magnetic properties of solids arise from the spins of unpaired electrons and their interactions.

\subsection*{What do you need to know already?}
You need to be aware of the properties of electron angular momentum (Topic 8B) and the relation of magnetic moments to angular momenta (Topic 8C).

The magnetic properties of metallic solids and semiconductors depend strongly on the band structures of the material. In this section, attention is confined largely to the much simpler magnetic properties of collections of individual molecules or ions, such as d-metal complexes. Much of the discussion therefore applies to liquid and gas-phase samples, as well as to solids.

\subsection*{15F. 1 Magnetic susceptibility}
Some molecules possess permanent magnetic dipole moments. In the absence of an external magnetic field, the orientation of the dipole is random and the material has no net magnetic moment. That changes when a magnetic field is applied and certain orientations are favoured. The magnetization, $\mathcal{M}$, the net dipole-moment density, is the resulting average molecular magnetic dipole moment multiplied by the number density of molecules in the sample. The magnetization induced by a magnetic field of strength $\mathcal{H}$ is proportional to $\mathcal{H}$, and is written

$$
\mathscr{M}=\chi \mathcal{H}
$$

Volume magnetic susceptibility [definition]\\
where $\chi$ is the dimensionless volume magnetic susceptibility. A closely related quantity is the molar magnetic susceptibility, $\chi_{\mathrm{m}}$ :

$$
\chi_{\mathrm{m}}=\chi V_{\mathrm{m}} \quad \begin{aligned}
& \text { Molar magnetic susceptibility } \\
& \text { [definition] }
\end{aligned}
$$

(15F.2)\\
where $V_{\mathrm{m}}$ is the molar volume of the substance.\\
The magnetization can be thought of as contributing to the density of lines of force in the material (Fig. 15F.1). Materials for which $\chi>0$ are called paramagnetic; they tend to move into a magnetic field and the density of lines of force within them is greater than in a vacuum. Those for which $\chi<0$ are called diamagnetic and tend to move out of a magnetic field; the density of lines of force within them is lower than in a vacuum. A paramagnetic material consists of ions or molecules with unpaired electrons, such as radicals and many d-metal complexes; a diamagnetic substance (which is far more common) is one with no unpaired electrons.

The magnetic susceptibility is traditionally measured with a 'Gouy balance'. This instrument consists of a sensitive balance from which the sample, contained in a narrow tube, hangs between the poles of a magnet. If the sample is paramagnetic, it is drawn into the field and its apparent weight is greater when the field is turned on. A diamagnetic sample tends to be expelled from the field and appears to weigh less when the field is turned on. The balance is normally calibrated against a sample of known susceptibility. The modern version of the determination makes use of a 'superconducting quantum interference device’ (a SQUID, Fig. 15F.2). A SQUID makes use\\
\includegraphics[max width=\textwidth, center]{2025_02_27_d4b95d9718f0b9e2e6b9g-706}

Figure 15F. 1 (a) In a vacuum, the strength of a magnetic field can be represented by the density of lines of force; (b) in a diamagnetic material, the density is reduced; (c) in a paramagnetic material, the density is increased.\\
\includegraphics[max width=\textwidth, center]{2025_02_27_d4b95d9718f0b9e2e6b9g-707}

Figure 15F.2 The arrangement used to measure magnetic susceptibility with a SQUID. The sample is moved upwards in small increments and the potential difference across the SQUID is measured.

Table 15F. 1 Magnetic susceptibilities at $298 \mathrm{~K}^{*}$

\begin{center}
\begin{tabular}{lll}
\hline
 & $\chi / 10^{-6}$ & $\chi_{\mathrm{m}} /\left(10^{-10} \mathrm{~m}^{3} \mathrm{~mol}^{-1}\right)$ \\
\hline
$\mathrm{H}_{2} \mathrm{O}(\mathrm{l})$ & -9.02 & -1.63 \\
$\mathrm{NaCl}(\mathrm{s})$ & -16 & -3.8 \\
$\mathrm{Cu}(\mathrm{s})$ & -9.7 & -0.69 \\
$\mathrm{CuSO}_{4} \cdot 5 \mathrm{H}_{2} \mathrm{O}(\mathrm{s})$ & +167 & +183 \\
\hline
\end{tabular}
\end{center}

\begin{itemize}
  \item More values are given in the Resource section.\\
of the quantization of magnetic flux and the property of current loops in superconductors that, as part of the circuit, include a weakly conducting link through which electrons must tunnel. The current that flows in the loop in a magnetic field depends on the value of the magnetic flux, and a SQUID can be exploited as a very sensitive magnetometer. Table 15F. 1 lists some experimental values of the magnetic susceptibility.
\end{itemize}

\subsection*{15F. 2 Permanent and induced magnetic moments}
The permanent magnetic moment of a molecule arises from any unpaired electron spins in the molecule. The magnitude, $m$, of the magnetic moment of an electron is proportional to the magnitude of the spin angular momentum, $\{s(s+1)\}^{1 / 2} \hbar$ :

\[
m=g_{\mathrm{e}}\{s(s+1)\}^{1 / 2} \mu_{\mathrm{B}} \quad \mu_{\mathrm{B}}=\frac{e \hbar}{2 m_{\mathrm{e}}} \quad \begin{align*}
& \text { Magnetic moment }  \tag{15F.3}\\
& {[\text { magnitude }]}
\end{align*}
\]

where $g_{\mathrm{e}}=2.0023$ and $\mu_{\mathrm{B}}$, the Bohr magneton, has the value $9.274 \times 10^{-24} \mathrm{JT}^{-1}$. If there are several electron spins in each molecule, they combine to give a total spin $S$, and then $s(s+1)$ is replaced by $S(S+1)$.

The magnetization, and consequently the magnetic susceptibility, depends on the temperature because the orientations\\
of the electron spins fluctuate. Some orientations have lower energy than others, and the magnetization depends on the randomizing influence of thermal motion. Thermal averaging of the permanent magnetic moments in the presence of an applied magnetic field results in a contribution to the magnetic susceptibility that is proportional to $m^{2} / 3 k T$. It follows that the spin contribution to the molar magnetic susceptibility is

$$
\chi_{\mathrm{m}}=\frac{N_{\mathrm{A}} g_{\mathrm{e}}^{2} \mu_{0} \mu_{\mathrm{B}}^{2} S(S+1)}{3 k T} \quad \begin{aligned}
& \text { Molar magnetic susceptibility } \\
& \text { [spin contribution] }
\end{aligned}
$$

(15F.4a)\\
where $\mu_{0}$ is the vacuum permeability. This susceptibility is positive, so the spin magnetic moments contribute to the paramagnetic susceptibilities of materials. Equation 15F.4a is commonly written as the Curie law:


\begin{equation*}
\chi_{\mathrm{m}}=\frac{C}{T} \quad C=\frac{N_{\mathrm{A}} g_{\mathrm{e}}^{2} \mu_{0} \mu_{\mathrm{B}}^{2} S(S+1)}{3 k} \quad \text { Curie law } \tag{15F.4b}
\end{equation*}


The spin contribution to the susceptibility decreases with increasing temperature because the thermal motion randomizes the spin orientations. In practice, a contribution to the paramagnetism also arises from the orbital angular momenta of electrons: here only the spin contribution has been considered.

\subsection*{Brief illustration 15F. 1}
Consider a complex salt with three unpaired electrons per complex cation and molar volume $61.7 \mathrm{~cm}^{3} \mathrm{~mol}^{-1}$; its molar magnetic susceptibility can be calculated using eqn 15F.4b. First, note that

$$
\frac{N_{\mathrm{A}} g_{\mathrm{e}}^{2} \mu_{0} \mu_{\mathrm{B}}^{2}}{3 k}=6.3001 \times 10^{-6} \mathrm{~m}^{3} \mathrm{~K} \mathrm{~mol}^{-1}
$$

Then with $S=\frac{3}{2}$ eqn 15 F .4 b gives

$$
\chi_{\mathrm{m}}=\left(6.3001 \times 10^{-6} \mathrm{~m}^{3} \mathrm{Kmol}^{-1}\right) \times \frac{\frac{3}{3}\left(\frac{3}{2}+1\right)}{298 \mathrm{~K}}=7.92 \ldots \times 10^{-8} \mathrm{~m}^{3} \mathrm{~mol}^{-1}
$$

and from eqn 15F. 2

$$
\chi=\frac{\chi_{\mathrm{m}}}{V_{\mathrm{m}}}=\frac{7.92 \ldots \times 10^{-8} \mathrm{~m}^{3} \mathrm{~mol}^{-1}}{6.17 \times 10^{-5} \mathrm{~m}^{3} \mathrm{~mol}^{-1}}=1.28 \times 10^{-3}
$$

At low temperatures, some paramagnetic solids make a phase transition to a state in which large domains of spins align with parallel orientations. This cooperative alignment gives rise to a very strong magnetization and is called ferromagnetism (Fig. 15F.3). In other cases, exchange interactions lead to alternating spin orientations: the spins are locked into a low-magnetization arrangement to give an antiferromagnetic

\footnotetext{${ }^{1}$ See our Physical chemistry: Quanta, matter, and change (2014) for the derivation of this contribution.
}
\includegraphics[max width=\textwidth, center]{2025_02_27_d4b95d9718f0b9e2e6b9g-708}

Figure15F. 3 (a) In a paramagnetic material, the electron spins are aligned at random in the absence of an applied magnetic field. (b) In a ferromagnetic material, the electron spins are locked into a parallel alignment over large domains. (c) In an antiferromagnetic material, the electron spins are locked into an antiparallel arrangement. The latter two arrangements survive even in the absence of an applied field.\\
phase. The ferromagnetic phase has a non-zero magnetization in the absence of an applied field, but the antiferromagnetic phase has zero magnetization because the spin magnetic moments cancel. The ferromagnetic transition occurs at the Curie temperature, and the antiferromagnetic transition occurs at the Néel temperature. Which type of cooperative behaviour occurs depends on the details of the band structure of the solid.

Magnetic moments can also be induced in molecules. To see how this effect arises, it is necessary to note that the circulation of electronic currents induced by an applied field gives rise to a magnetic field which usually opposes the applied field, thus making the substance diamagnetic. In these cases, the induced electron currents occur within the molecular orbitals that are occupied in its ground state. There are a few cases in which molecules are paramagnetic despite having no unpaired electrons. In these materials the induced electron currents flow in the opposite direction because they can make use of unoccupied orbitals that lie close to the HOMO in energy (a similar effect is the paramagnetic contribution to the chemical shift, Topic 12B). This orbital paramagnetism is distinguished from spin paramagnetism by the fact that it is temperature independent and is called temperature-independent paramagnetism (TIP).

These remarks can be summarized as follows. All molecules have a diamagnetic component to their susceptibility, but this contribution is dominated by spin paramagnetism if the molecules have unpaired electrons. In a few cases (where there are low-lying excited states) TIP is strong enough to make the molecules paramagnetic even though all their electrons are paired.

\subsection*{15F. 3 Magnetic properties of superconductors}
Superconductors have unique magnetic properties. Superconductors classed as Type I show abrupt loss of superconductivity when an applied magnetic field exceeds a critical value $\mathcal{H}_{c}$ characteristic of the material. An empirical relation between the value of $\mathcal{H}_{\mathfrak{c}}$, the temperature $T$, and the critical temperature $T_{c}$ is


\begin{equation*}
\mathcal{H}_{\mathrm{c}}(T)=\mathcal{H}_{\mathrm{c}}(0)\left(1-\frac{T^{2}}{T_{\mathrm{c}}^{2}}\right) \quad \text { Dependence of } \mathcal{H}_{\mathrm{c}} \text { on } T \tag{15F.5}
\end{equation*}


provided $T \leq T_{c}$. Note that the critical field falls as $T$ rises from 0 towards $T_{c}$. Therefore, to maintain superconductivity in the presence of a magnetic field, it is best to keep $T$ well below $T_{c}$ and to select a material with a high $\mathcal{H}_{c}(0)$.

\subsection*{Brief illustration 15F. 2}
Lead has $T_{\mathrm{c}}=7.19 \mathrm{~K}$ and $\mathcal{H}_{\mathrm{c}}(0)=63.9 \mathrm{kA} \mathrm{m}^{-1}$. At $T=6.0 \mathrm{~K}$ the magnetic field that would quench its superconductivity would be

$$
\mathcal{H}_{\mathrm{c}}(6.0 \mathrm{~K})=\left(63.9 \mathrm{kA} \mathrm{~m}^{-1}\right) \overbrace{\left(1-\frac{(6.0 \mathrm{~K})^{2}}{(7.19 \mathrm{~K})^{2}}\right)}^{0.303}=19 \mathrm{kA} \mathrm{~m}^{-1}
$$

The lead remains superconducting at 6.0 K for this and weaker applied field strengths. If the temperature is lowered to 5.0 K the corresponding calculation gives $\mathcal{H}_{c}(5.0 \mathrm{~K})=33 \mathrm{kAm}^{-1}$, and superconductivity survives at higher field strengths. In each case, superconductivity would survive to higher field strengths for a material with a higher $\mathcal{H}_{c}(0)$.

Type I superconductors are also completely diamagnetic below $\mathcal{H}_{c}$, meaning that the magnetic field does not penetrate into the material. This complete exclusion of a magnetic field from a material is known as the Meissner effect, which can be demonstrated by the levitation of a superconductor above a magnet. Type II superconductors, which include the HTSCs, show a gradual loss of diamagnetism with increasing magnetic field.

\subsection*{Checklist of concepts}
\begin{enumerate}
  \item The magnetization of a material is the average molecular magnetic dipole moment multiplied by the number density of the molecules.\\
$\square$ 2. The magnetic susceptibility expresses the relation between the magnetization and the applied magnetic field strength.3. Diamagnetic materials tend to move out of a magnetic field and have negative magnetic susceptibilities.4. Paramagnetic materials tend to move into a magnetic field and have positive magnetic susceptibilities.5. The Curie law describes the temperature dependence of the magnetic susceptibility.6. Ferromagnetism is the cooperative alignment of electron spins in a material and gives rise to strong permanent magnetization.
  \item Antiferromagnetism results from alternating spin orientations in a material and leads to weak magnetization.
  \item Temperature-independent paramagnetism arises from induced electron currents.
  \item The Meissner effect is the exclusion of a magnetic field from a Type I superconductor.
\end{enumerate}

\subsection*{Checklist of equations}
\begin{center}
\begin{tabular}{llll}
\hline
Property & Equation & Comment & Equation number \\
\hline
Magnetization & $\mathcal{M}=\chi \mathcal{H}$ & Definition & 15 F .1 \\
Molar magnetic susceptibility & $\chi_{\mathrm{m}}=\chi V_{\mathrm{m}}$ & Definition & 15 F .2 \\
Magnetic moment & $m=g_{\mathrm{e}}\{s(s+1)\}^{1 / 2} \mu_{\mathrm{B}}$ & $\mu_{\mathrm{B}}=e \hbar / 2 m_{\mathrm{e}}$ & 15 F .3 \\
Curie law & $\chi_{\mathrm{m}}=C / T, C=N_{\mathrm{A}} g_{\mathrm{e}}^{2} \mu_{0} \mu_{\mathrm{B}}^{2} S(S+1) / 3 k$ & Paramagnetism & 15 F .4 b \\
Dependence of $\mathcal{H}_{\mathrm{c}}$ on $T_{\mathrm{c}}$ & $\mathcal{H}_{\mathrm{c}}(T)=\mathcal{H}_{\mathrm{c}}(0)\left(1-T^{2} / T_{\mathrm{c}}^{2}\right)$ & Empirical & 15 F .5 \\
\hline
\end{tabular}
\end{center}

\section*{TOPIC 15G The optical properties of solids }
\subsection*{Why do you need to know this material?}
The optical properties of solids are of ever increasing importance in modern technology, not only for the generation of light but for the propagation and manipulation of information.

\subsection*{What is the key idea?}
The optical properties of molecules in solids differ from those of isolated molecules as a result of the interaction of their transition dipoles.

What do you need to know already?\\
You need to be familiar with the concept of transition dipole (Topics 8C and 11A) and of the band theory of solids (Topic 15C).

Topic 11A explains the factors that determine the energy and intensity of light absorbed by isolated atoms and molecules in the gas phase and in solution. However, significant differences arise when the molecules are neighbours in a solid.

\subsection*{15G. 1 Excitons}
Consider an electronic excitation of a molecule (or an ion) in a crystal. If the excitation corresponds to the removal of an electron from one orbital of a molecule and its elevation to an orbital of higher energy, then the excited state of the molecule can be envisaged as the coexistence of an electron and a hole. This electron-hole pair, which behaves as a particlelike exciton, migrates from molecule to molecule in the crystal (Fig. 15G.1). A migrating excitation of this kind is called a Frenkel exciton, and is commonly found in molecular solids. The electron and hole can also be on different molecules, but in each other's vicinity. A migrating excitation of this kind, which is now spread over several molecules (more usually ions), is called a Wannier exciton. Exciton formation causes spectral lines to shift, split, and change intensity.\\
\includegraphics[max width=\textwidth, center]{2025_02_27_d4b95d9718f0b9e2e6b9g-710}

Figure 15G. 1 The electron-hole pair shown on the left can migrate through a solid lattice as the excitation hops from molecule to molecule. The mobile excitation is called an exciton.

The migration of a Frenkel exciton (the only type considered here) implies that there is an interaction between the molecules that constitute the crystal: if this were not the case the excitation on one molecule could not move to another. This interaction affects the energy levels of the system. The strength of the interaction also governs the rate at which an exciton moves through the crystal: a strong interaction results in fast migration and a vanishingly small interaction leaves the exciton localized on its original molecule. The specific mechanism of interaction that leads to exciton migration is the interaction between the transition dipoles of the excitation (Topic 11A). Thus, an electric dipole transition in a molecule is accompanied by a shift of charge, and this transient dipole exerts a force on an adjacent molecule. The latter responds by shifting its charge. This process continues and the excitation migrates through the crystal.

The energy shift arising from the interaction between transition dipoles can be explained as follows. The potential energy of interaction between two parallel electric dipole moments $\mu_{1}$ and $\mu_{2}$ separated by a distance $r$ is $V=\mu_{1} \mu_{2}\left(1-3 \cos ^{2} \theta\right) / 4 \pi \varepsilon_{0} r^{3}$, where the angle $\theta$ is defined in (1). A head-to-tail alignment corresponds to $\theta=0$, and a parallel alignment corresponds to $\theta=90^{\circ}$. From the expression for $V$ it follows that $V<0$ (a favourable, energy-lowering interaction) for $0 \leq \theta<54.7^{\circ}, V=0$ when $\theta=54.7^{\circ}$ (at this angle $1-3 \cos ^{2} \theta=0$ ), and $V>0$ (an unfavourable, energy-raising interaction) for $54.7^{\circ}<\theta \leq 90^{\circ}$.\\
\includegraphics[max width=\textwidth, center]{2025_02_27_d4b95d9718f0b9e2e6b9g-710(1)}\\
\includegraphics[max width=\textwidth, center]{2025_02_27_d4b95d9718f0b9e2e6b9g-711(4)}

Figure 15G. 2 (a) The alignment of transition dipoles (the yellow arrows) shown here is energetically unfavourable, and the exciton absorption is shifted to higher energy (higher frequency). (b) The alignment shown here is energetically favourable for a transition in this orientation, and the exciton band occurs at lower frequency than in the isolated molecules.

In a head-to-tail arrangement, there is a favourable interaction between the region of partial positive charge in one molecule and the region of partial negative charge in the other molecule. In contrast, in a parallel arrangement, the molecular interaction is unfavourable because of the close approach of regions of partial charge with the same sign. It follows that an all-parallel arrangement of the transition dipoles (Fig. 15G.2a) is energetically unfavourable, so the absorption occurs at a higher frequency than in the isolated molecule. Conversely, a head-to-tail alignment of transition dipoles (Fig. 15G.2b) is energetically favourable, and the transition occurs at a lower frequency than in the isolated molecules.

If there are $N$ molecules per unit cell, there are $N$ exciton bands in the spectrum (if all of them are allowed). The splitting between the bands is the Davydov splitting. To understand the origin of the splitting, consider the case $N=2$ with the molecules arranged as in Fig. 15G. 3 and suppose that the\\
(a)\\
\includegraphics[max width=\textwidth, center]{2025_02_27_d4b95d9718f0b9e2e6b9g-711(3)}\\
\includegraphics[max width=\textwidth, center]{2025_02_27_d4b95d9718f0b9e2e6b9g-711}\\
(b)\\
\includegraphics[max width=\textwidth, center]{2025_02_27_d4b95d9718f0b9e2e6b9g-711(1)}\\
\includegraphics[max width=\textwidth, center]{2025_02_27_d4b95d9718f0b9e2e6b9g-711(2)}

Figure 15G.3 When the transition moments within a unit cell lie in different relative directions, as depicted in (a) and (b), the energies of the transitions are shifted and give rise to the two bands labelled (a) and (b) in the spectrum. The separation of the bands is the Davydov splitting.\\
transition dipoles are along the length of the molecules. The radiation stimulates the collective excitation of the transition dipoles that are in-phase between neighbouring unit cells. Within each unit cell the transition dipoles may be arrayed in the two different ways shown in the illustration. The two orientations correspond to different interaction energies, with interaction being unfavourable in one and favourable in the other, so the two transitions appear in the spectrum as two bands of different frequencies. The magnitude of the Davydov splitting is determined by the energy of interaction between the transition dipoles within the unit cell.

\subsection*{15G. 2 Metals and semiconductors}
Figure 15C. 8 shows the band structure in an idealized metallic conductor at $T=0$. The absorption of a photon can excite electrons from the occupied levels to the unoccupied levels. There is a near continuum of unoccupied energy levels above the Fermi level, so absorption occurs over a wide range of frequencies. In metals, the bands are sufficiently wide that radiation is absorbed from the radiofrequency to the ultraviolet region of the electromagnetic spectrum but not to very high-frequency electromagnetic radiation, such as X-rays and $\gamma$-rays, so metals are transparent at these frequencies. Because this range of absorbed frequencies includes the entire visible spectrum, it might be expected that all metals should be black. However, metals are in fact lustrous (that is, they reflect light) and some are coloured (that is, they absorb light of certain wavelengths), so the model clearly needs some improvement.

\subsection*{(a) Light absorption}
To explain the lustrous appearance of a smooth metal surface, it is important to realize that the absorbed energy can be reemitted very efficiently as light, with only a small fraction of the energy being released into the bulk as heat. Because the atoms near the surface of the material absorb most of the radiation, emission also occurs primarily from the surface. In essence, if the sample is excited with visible light, then electrons near the surface are driven into oscillation at the same frequency, and visible light is emitted from the surface, so accounting for the lustre of the material.

The perceived colour of a metal depends on the frequency range of reflected light. That in turn depends on the frequency range of light that can be absorbed and, by extension, on the band structure. Silver reflects light with nearly equal efficiency across the visible spectrum because its band structure has many unoccupied energy levels that can be populated by absorption of, and depopulated by emission of, visible light. On the other hand, copper has its characteristic colour because it has relatively fewer unoccupied energy levels that can\\
\includegraphics[max width=\textwidth, center]{2025_02_27_d4b95d9718f0b9e2e6b9g-712}

Figure 15G. 4 In some materials, the band gap $E_{\mathrm{g}}$ is very large and electron promotion can occur only by excitation with electromagnetic radiation.\\
be excited with violet, blue, and green light. The material reflects at all wavelengths, but more light is emitted at lower frequencies (corresponding to yellow, orange, and red). Similar arguments account for the colours of other metals, such as the yellow of gold. It is interesting to note that the colour of gold can be accounted for only by including relativistic effects in the calculation of its band structure.

Now consider semiconductors. If the band gap $E_{\mathrm{g}}$ is comparable to $k T$ the promotion of electrons from the conduction to the valence band of a semiconductor can be the result of thermal excitation. In some materials, the band gap is very large and electron promotion can occur only by excitation with electromagnetic radiation. However, as is seen from Fig. 15G.4, there is a frequency $v_{\text {min }}=E_{\mathrm{g}} / h$ below which light absorption cannot occur. Above this frequency threshold, a wide range of frequencies can be absorbed by the material, as in a metal.

\subsection*{Brief illustration 15G. 1}
The energy of the band gap in the semiconductor cadmium sulfide (CdS) is 2.4 eV (equivalent to $3.8 \times 10^{-19} \mathrm{~J}$ ). It follows that the minimum electronic absorption frequency is

$$
v_{\min }=\frac{3.8 \times 10^{-19} \mathrm{~J}}{6.626 \times 10^{-34} \mathrm{Js}}=5.8 \times 10^{14} \mathrm{~s}^{-1}
$$

This frequency, of 580 THz , corresponds to a wavelength of 520 nm (green light). Lower frequencies, corresponding to yellow, orange, and red, are not absorbed and consequently CdS appears yellow-orange.

\subsection*{(b) Light-emitting diodes and diode lasers}
The unique electrical properties of $\mathrm{p}-\mathrm{n}$ junctions between semiconductors (which are described in Topic 15E) can be put to good use in optical devices. In some materials, most notably gallium arsenide, GaAs, energy from electron-hole recombination is released not as heat but is carried away by photons as electrons move across the junction driven by the appropriate potential difference. Practical light-emitting diodes of this kind are widely used in electronic displays. The wavelength of\\
emitted light depends on the band gap of the semiconductor. Gallium arsenide itself emits infrared light, but its band gap is widened by incorporating phosphorus, and a material of composition approximately $\mathrm{GaAs}_{0.6} \mathrm{P}_{0.4}$ emits light in the red region of the spectrum.

A light-emitting diode is not a laser (Topic 11G) because stimulated emission is not involved. In diode lasers, light emission due to electron-hole recombination is employed as the basis of laser action, and the population inversion can be sustained by sweeping away the electrons that fall into the holes of the p-type semiconductor. One widely used material is $\mathrm{Ga}_{1-x} \mathrm{Al}_{x} \mathrm{As}$, which produces infrared laser radiation and is widely used in CD and DVD players. High-power diode lasers are also used to pump other lasers. One example is the pumping of $\mathrm{Nd}: Y A G$ lasers (Topic 11 G ) by $\mathrm{Ga}_{0.91} \mathrm{Al}_{0.09} \mathrm{As}^{2} / \mathrm{Ga}_{0.7} \mathrm{Al}_{0.3} \mathrm{As}$ diode lasers.

\subsection*{15G.3 Nonlinear optical phenomena}
Nonlinear optical phenomena arise from changes in the optical properties of a material in the presence of intense electromagnetic radiation. In frequency doubling (or 'second harmonic generation'), an intense laser beam is converted to radiation with twice (and in general a multiple) of its initial frequency as it passes through a suitable material. It follows that frequency doubling and tripling of an Nd:YAG laser, which emits radiation at 1064 nm (Topic 11G), produce green light at 532 nm and ultraviolet radiation at 355 nm , respectively. Common materials that can be used for frequency doubling in laser systems include crystals of potassium dihydrogenphosphate $\left(\mathrm{KH}_{2} \mathrm{PO}_{4}\right)$, lithium niobate $\left(\mathrm{LiNbO}_{3}\right)$, and $\beta$-barium borate $\left(\beta-\mathrm{BaB}_{2} \mathrm{O}_{4}\right)$.

Frequency doubling can be explained by examining how a substance responds nonlinearly to incident radiation of frequency $\omega=2 \pi v$. Radiation of a particular frequency arises from oscillations of an electric dipole at that frequency and the incident electric field $\mathcal{E}$ of the radiation induces an electric dipole moment of magnitude $\mu$, in the substance. At low light intensity, most materials respond linearly, in the sense that $\mu=\alpha E$, where $\alpha$ is the polarizability (Topic 14A). At high light intensity, the hyperpolarizability $\beta$ of the material becomes important (Topic 14A) and the induced dipole becomes

\[
\mu=\alpha E+\frac{1}{2} \beta E^{2}+\cdots \quad \begin{align*}
& \text { Induced dipole moment in }  \tag{15G.1}\\
& \text { terms of the hyperpolarizability }
\end{align*}
\]

The nonlinear term $\beta E^{2}$ can be expanded as follows if it is supposed that the incident electric field is $\mathcal{E}_{0} \cos \omega t$ :


\begin{equation*}
\beta E^{2}=\beta E_{0}^{2} \cos ^{2} \omega t=\frac{1}{2} \beta E_{0}^{2}(1+\cos 2 \omega t) \tag{15G.2}
\end{equation*}


Hence, the nonlinear term contributes an induced electric dipole that includes a component that oscillates at the frequency $2 \omega$ and that can act as a source of radiation of that frequency.

\subsection*{Checklist of concepts}
1. An exciton is an electron-hole pair caused by optical excitation in a solid; Frenkel excitons are localized on a single molecule, whereas Wannier excitons are spread over several molecules.2. If the unit cell contains $N$ molecules, there are $N$ exciton bands in the spectrum separated by the Davydov splitting.3. Nonlinear optical phenomena arise from changes in the optical properties of a material in the presence of intense electromagnetic radiation; they can give rise to frequency doubling.

\chapter{FOCUS 16 Molecules in motion}
This Focus is concerned with understanding how matter and other physical properties (such as energy and momentum) are transported from one place to another in both gases and liquids.

\subsection*{16A Transport properties of a perfect gas}
The transport of matter and physical properties can be described by a set of closely related empirical equations. For a gas, it is possible to understand the form of these equations by building a model based on the kinetic theory of gases discussed in Topic 1B. With this approach, the rate of diffusion, the rate of thermal conduction, viscosity, and effusion can all be related to quantities arising from the kinetic theory.\\
16A. 1 The phenomenological equations; 16A. 2 The transport parameters

\subsection*{16B Motion in liquids}
Molecular motion in liquids is different from that in gases on account of the presence of significant intermolecular interactions and the much higher density typical of a liquid. One way to monitor motion in such systems is to explore the electrical\\
resistance of electrolyte solutions and to analyse it in terms of the response of the ions to an applied electric field.\\
16B. 1 Experimental results; 16B. 2 The mobilities of ions

\subsection*{16C Diffusion}
The diffusion of solutes and various physical properties is an important process in liquids. It can be discussed by introducing the concept of a general 'thermodynamic force' that can be regarded as being responsible for the motion of molecules. This apparent force can be used to construct the important 'diffusion equation', which describes how solutes spread out in space with increasing time. An alternative model of diffusion as a random walk gives further insight.\\
16C. 1 The thermodynamic view; 16C. 2 The diffusion equation;\\
16C. 3 The statistical view

\subsection*{Web resource What is an application of this material?}
A great deal of chemistry, chemical engineering, and biology depends on the ability of molecules and ions to migrate through media of various kinds. Impact 25 explains how conductivity measurements are used to analyse the motion of ions through biological membranes.

\section*{TOPIC 16A Transport properties of a perfect gas }
\subsection*{Why do you need to know this material?}
Many physical processes take place by the transfer of a property from one region to another, and gas-phase chemical reactions depend on the rate at which molecules collide. The material presented here also includes general aspects of transport in fluid systems of any kind, and which are applicable to reactions taking place in solution.

\subsection*{What is the key idea?}
A molecule carries properties through space in steps of about the distance of its mean free path.

\subsection*{What do you need to know already?}
This Topic builds on and extends the kinetic theory of gases (Topic 1B). You need to be familiar with the concepts of the mean speed of molecules and the mean free path and their dependence on pressure and temperature.

A transport property is a process by which matter or an attribute of matter, such as momentum, is carried through a medium from one location to another. The rate of transport is commonly expressed in terms of an equation that is an empirical summary of experimental observations. These equations apply to all kinds of properties and media, and can be adapted to the discussion of transport properties of gases. In such cases, the kinetic theory of gases provides simple expressions that show how the rates of transport of these properties depend on the pressure and the temperature. The most important concept from the kinetic theory developed in Topic 1B, and used throughout this Topic, is the mean free path, $\lambda$, the average distance a molecule travels between collisions. According to eqn 1B.14, at a temperature $T$ and a pressure $p$

$$
\lambda=\frac{k T}{\sigma p}
$$

Mean free path\\[0pt]
[kinetic theory]\\
(16A.1a)

The parameter $\sigma$ is the collision cross-section of a molecule, a measure of the target area it presents in a collision. For the derivation and physical interpretation of this expression, see Topic 1B. Another important result from kinetic theory is the mean speed of molecules of molar mass $M$ at a temperature $T$, which is given by eqn 1B. 9 :

\[
v_{\text {mean }}=\left(\frac{8 R T}{\pi M}\right)^{1 / 2} \quad \begin{align*}
& \text { Mean speed }  \tag{16A.1b}\\
& {[\text { kinetic theory }]}
\end{align*}
\]

\subsection*{16A. 1 The phenomenological equations}
A 'phenomenological equation' is an equation that summarizes empirical observations on phenomena without, initially at least, being based on an understanding of the molecular processes responsible for the property. Such equations are encountered commonly in the study of fluids.

The rate of migration of a property is measured by its flux, $J$, the quantity of that property passing through a given area in a given time interval divided by the area and the duration of the interval. If matter is flowing (as in diffusion), the matter flux is reported as so many molecules per square metre per second (number or amount $\mathrm{m}^{-2} \mathrm{~s}^{-1}$ ). If the property migrating is energy (as in thermal conduction), then the energy flux is expressed in joules per square metre per second $\left(\mathrm{Jm}^{-2} \mathrm{~s}^{-1}\right)$, and so on. The total quantity of a property transferred through a given area $A$ in a given time interval $\Delta t$ is $|J| A \Delta t$. The flux $J$ may be positive or negative: the significance of its sign is discussed below.

Experimental observations on transport properties show that the flux of a property is usually proportional to the first derivative of a related quantity. For example, the flux of matter diffusing parallel to the $z$-axis of a container is found to be proportional to the gradient of the concentration along the same direction:

$$
J(\text { matter }) \propto \frac{\mathrm{d} \mathcal{N}}{\mathrm{~d} z}
$$

Fick's first law of diffusion\\
(16A.2)\\
where $\mathcal{N}$ is the number density of particles, with units number per metre cubed $\left(\mathrm{m}^{-3}\right)$. The proportionality of the flux of matter to the concentration gradient is sometimes called Fick's first law of diffusion: the law implies that diffusion is faster when the concentration varies steeply with position than when the concentration is nearly uniform. There is no net flux if the concentration is uniform $(\mathrm{d} \mathcal{N} / \mathrm{d} z=0)$. Similarly, the rate of thermal conduction (the flux of the energy associated with thermal motion) is found to be proportional to the temperature gradient:\\
$J$ (energy of thermal motion) $\propto \frac{\mathrm{d} T}{\mathrm{~d} z} \quad$ Flux of energy (16A.3)\\
A positive value of $J$ signifies a flux towards positive $z$; a negative value of $J$ signifies a flux towards negative $z$. Because matter flows down a concentration gradient, from high concentration to low concentration, $J$ is positive if $\mathrm{d} \mathcal{N} / \mathrm{d} z$ is negative (Fig. 16A.1). Therefore, the coefficient of proportionality in eqn 16A. 2 must be negative, and it is written as $-D$ :

$$
J(\text { matter })=-D \frac{\mathrm{~d} \mathcal{N}}{\mathrm{~d} z}
$$

Fick's first law in terms of the diffusion coefficient\\
(16A.4)

The constant $D$ is the called the diffusion coefficient; its SI units are metre squared per second $\left(\mathrm{m}^{2} \mathrm{~s}^{-1}\right)$. Energy migrates down a temperature gradient, and the same reasoning leads to


\begin{align*}
& J(\text { energy of thermal motion })=-\kappa \frac{\mathrm{d} T}{\mathrm{~d} z} \\
& \qquad \begin{array}{l}
\text { Flux of energy in terms of the } \\
\text { coefficient of thermal conductivity }
\end{array} \tag{16A.5}
\end{align*}


where $\kappa$ (kappa) is the coefficient of thermal conductivity. The units of $\kappa$ are joules per kelvin per metre per second $\left(\mathrm{J}^{-1} \mathrm{~m}^{-1} \mathrm{~s}^{-1}\right)$ or, because $1 \mathrm{Js}^{-1}=1 \mathrm{~W}$, watts per kelvin per metre ( $\mathrm{W} \mathrm{K}^{-1} \mathrm{~m}^{-1}$ ). Some experimental values are given in Table 16A.1.\\
\includegraphics[max width=\textwidth, center]{2025_02_27_d4b95d9718f0b9e2e6b9g-723}

Figure 16A. 1 The flux of particles down a concentration gradient. Fick's first law states that the flux of matter is proportional to the concentration gradient at that point.

Table 16A.1 Transport properties of gases at 1 atm*

\begin{center}
\begin{tabular}{lccc}
\hline
 & $\kappa /\left(\mathrm{mW} \mathrm{K}^{-1} \mathrm{~m}^{-1}\right)$ & \multicolumn{2}{c}{$\eta / \mu \mathrm{P}^{\ddagger}$} \\
\cline { 2 - 4 }
 & 273 K & 273 K & 293 K \\
\hline
Ar & 16.3 & 210 & 223 \\
$\mathrm{CO}_{2}$ & 14.5 & 136 & 147 \\
He & 144.2 & 187 & 196 \\
$\mathrm{~N}_{2}$ & 24.0 & 166 & 176 \\
\hline
\end{tabular}
\end{center}

\begin{itemize}
  \item More values are given in the Resource section.\\
${ }^{\ddagger} 1 \mu \mathrm{P}=10^{-7} \mathrm{~kg} \mathrm{~m}^{-1} \mathrm{~s}^{-1}$.
\end{itemize}

\subsection*{Brief illustration 16A. 1}
Suppose that two metal plates are placed perpendicular to the $z$-axis, with the first at $z=0$ and the second at $z=+1.0 \mathrm{~cm}$, and that the temperature of the second plate is 10 K higher than that of the first plate. The temperature gradient is

$$
\frac{\mathrm{d} T}{\mathrm{~d} z}=\frac{10 \mathrm{~K}}{1.0 \times 10^{-2} \mathrm{~m}}=+1.0 \times 10^{3} \mathrm{Km}^{-1}
$$

If the plates are separated by air, for which $\kappa=24.1 \mathrm{~mW} \mathrm{~K}^{-1} \mathrm{~m}^{-1}$, the energy flux is

$$
\begin{aligned}
J & =-\kappa \frac{\mathrm{d} T}{\mathrm{~d} z}=-\left(24.1 \mathrm{mWK}^{-1} \mathrm{~m}^{-1}\right) \times\left(1.0 \times 10^{3} \mathrm{Km}^{-1}\right) \\
& =-24 \mathrm{Wm}^{-2}
\end{aligned}
$$

The flux is negative because energy is transferred from the hotter plate at $z=+1.0 \mathrm{~cm}$ to the cooler plate at $z=0$, which is a flow in the $-z$ direction. The energy transferred through an area of $1.0 \mathrm{~cm}^{2}$ between the two plates in $1 \mathrm{~h}(3600 \mathrm{~s})$ is\\
Transfer $=|J| A \Delta t=\left(24 \mathrm{Wm}^{-2}\right) \times\left(1.0 \times 10^{-4} \mathrm{~m}^{2}\right) \times(3600 \mathrm{~s})=8.6 \mathrm{~J}$

Viscosity arises from the flux of linear momentum. To see the connection, consider a fluid in a state of Newtonian flow, in which a series of layers move past one another and, in this case, in the $x$-direction (Fig. 16A.2). The layer next to the wall of the vessel is stationary, and the velocity of successive layers varies linearly with distance, $z$, from the wall. Molecules ceaselessly move between the layers and bring with them the $x$-component of linear momentum they possessed in their original layer. A layer is retarded by molecules arriving from a more slowly moving layer because such molecules have a lower momentum in the $x$-direction. A layer is accelerated by molecules arriving from a more rapidly moving layer. The net retarding effect is interpreted as the viscosity of the fluid.

Because the retarding effect depends on the transfer of the $x$-component of linear momentum into the layer of interest,\\
\includegraphics[max width=\textwidth, center]{2025_02_27_d4b95d9718f0b9e2e6b9g-724}

Figure 16A.2 The viscosity of a fluid arises from the transport of linear momentum. In this illustration the fluid is undergoing Newtonian (laminar) flow in the $x$-direction, and particles bring their initial momentum when they enter a new layer.\\
the viscosity depends on the flux of this $x$-component in the $z$-direction. The flux of the $x$-component of momentum is proportional to $\mathrm{d} v_{x} / \mathrm{d} z$, where $v_{x}$ is the velocity in the $x$-direction; the flux can therefore be written

$$
J(x \text {-component of momentum })=-\eta \frac{\mathrm{d} v_{x}}{\mathrm{~d} z}
$$

\begin{displayquote}
Momentum flux in terms of the coefficient of viscosity\\
(16A.6)
\end{displayquote}

The constant of proportionality, $\eta$, is the coefficient of viscosity (or simply 'the viscosity'). Its units are kilograms per metre per second $\left(\mathrm{kg} \mathrm{m}^{-1} \mathrm{~s}^{-1}\right)$. Viscosities are often reported in the non-SI unit poise (P), with $1 \mathrm{P}=10^{-1} \mathrm{~kg} \mathrm{~m}^{-1} \mathrm{~s}^{-1}$. Some experimental values are given in Table 16A.1.

Although it is not strictly a transport property, closely related to diffusion is effusion, the escape of matter through a small hole. The essential empirical observations on effusion are summarized by Graham's law of effusion, which states that the rate of effusion is inversely proportional to the square root of the molar mass, $M$.

\subsection*{16A. 2 The transport parameters}
The kinetic theory of gases (Topic 1B) can be used to derive expressions for the diffusion characteristics of a perfect gas. All the expressions depend on knowing the collision flux, $Z_{\text {w }}$, which is the rate at which molecules strike a region in the gas (this region may be an imaginary window, a part of a wall, or a hole in a wall). Specifically, the collision flux is the number of collisions divided by the area of the region and the duration of the time interval. Its dependence on pressure and temperature can be derived from the kinetic theory.

How is that done? 16A. 1 Deriving an expression for the collision flux

Consider a wall of area $A$ perpendicular to the $x$-axis (Fig. 16A.3). In the following calculation, note that for a perfect gas the equation of state $p V=n R T$ can be used to relate the number density, $\mathcal{N}$, to the pressure by $\mathcal{N}=N / V=n N_{\mathrm{A}} / V$ $=n N_{\mathrm{A}} p / n R T=p / k T$. In the final equality $R=N_{\mathrm{A}} k$ was used, where $k$ is Boltzmann's constant.\\
\includegraphics[max width=\textwidth, center]{2025_02_27_d4b95d9718f0b9e2e6b9g-724(1)}

Figure 16A.3 A molecule will reach the wall on the right within an interval $\Delta t$ if it is within a distance $v_{x} \Delta t$ of the wall and travelling to the right.

Step 1 Identify the number of molecules that will strike an area If a molecule has $v_{x}>0$ (that is, it is travelling in the direction of positive $x$ ), then it will strike the wall within an interval $\Delta t$ if it lies within a distance $v_{x} \Delta t$ of the wall. Therefore, all molecules in the volume $A v_{x} \Delta t$, and with a positive $x$-component of velocity, will strike the wall in the interval $\Delta t$. The total number of collisions in this interval is therefore $\mathcal{N} A v_{x} \Delta t$, where $\mathcal{N}$ is the number density of molecules.

Step 2 Take into account the range of velocities\\
The velocity $v_{x}$ has a range of values described by the probability distribution $f\left(v_{x}\right)$ given in eqn 1B.3:

$$
f\left(v_{x}\right)=\left(\frac{m}{2 \pi k T}\right)^{1 / 2} \mathrm{e}^{-m v_{x}^{2} / 2 k T}
$$

with $f\left(v_{x}\right) \mathrm{d} v_{x}$ the probability of finding a molecule with a component of velocity between $v_{x}$ and $v_{x}+\mathrm{d} v_{x}$. The total number of collisions is found by summing $\mathcal{N} A v_{x} \Delta t$ over all positive values of $v_{x}$ (because only molecules with a positive component of velocity are moving towards the area of interest) with each value of $v_{x}$ being weighted by the probability of it occurring:

$$
\text { Number of collisions }=\mathcal{N} A \Delta t \int_{0}^{\infty} v_{x} f\left(v_{x}\right) \mathrm{d} v_{x}
$$

The collision flux is the number of collisions divided by $A$ and $\Delta t$, so

$$
Z_{\mathrm{W}}=\mathcal{N} \int_{0}^{\infty} v_{x} f\left(v_{x}\right) \mathrm{d} v_{x}
$$

Step 3 Evaluate the integral\\
Because

$$
\int_{0}^{\infty} v_{x} f\left(v_{x}\right) \mathrm{d} v_{x}=\left(\frac{m}{2 \pi k T}\right)^{1 / 2} \overbrace{\int_{0}^{\infty} v_{x} \mathrm{e}^{-m v_{x}^{2} / 2 k T} \mathrm{~d} v_{x}}^{\text {Integral G.2 }}=\left(\frac{k T}{2 \pi m}\right)^{1 / 2}
$$

it follows that

$$
Z_{\mathrm{W}}=\mathcal{N}\left(\frac{k T}{2 \pi m}\right)^{1 / 2} \stackrel{\frac{p}{k T}\left(\frac{k T}{2 \pi m}\right)^{1 / 2}}{=\frac{\mathcal{V}=p / k T}{}}
$$

and therefore\\
\includegraphics[max width=\textwidth, center]{2025_02_27_d4b95d9718f0b9e2e6b9g-725}

Step 4 Develop an alternative expression in terms of the mean speed\\
The mean speed is given by eqn 16A.1b as

$$
v_{\text {mean }}=\left(\frac{8 R T}{\pi M}\right)^{1 / 2}=\left(\frac{8 k T}{\pi m}\right)^{1 / 2}
$$

It follows that

$$
\left(\frac{k T}{2 \pi m}\right)^{1 / 2}=\frac{1}{4} v_{\text {mean }}
$$

and therefore $Z_{\mathrm{w}}=\mathcal{N}(\mathrm{kT} / 2 \pi \mathrm{~m})^{1 / 2}$ can be expressed as

$$
-\left\lvert\, Z_{\mathrm{w}}=\frac{1}{4} \mathcal{N} v_{\text {mean }} \nmid \begin{gathered}
\text { Collision flux } \\
\text { [perfect gas] }
\end{gathered}\right.
$$

According to eqn 16A.7a, the collision flux increases with pressure simply because increasing the pressure increases the number density and hence the number of collisions. The flux decreases with increasing mass of the molecules because heavy molecules move more slowly than light molecules. Caution, however, is needed with the interpretation of the role of temperature: it is wrong to conclude that because $T^{1 / 2}$ appears in the denominator that the collision flux decreases with increasing temperature. If the system has constant volume, the pressure increases with temperature $(p \propto T)$, so the collision flux is in fact proportional to $T / T^{1 / 2}=T^{1 / 2}$, and increases with temperature (because the molecules are moving faster).

\subsection*{Brief illustration 16A. 2}
The collision flux of $\mathrm{O}_{2}$ molecules, with $m=M / N_{\mathrm{A}}$ and $M=$ $32.00 \mathrm{~g} \mathrm{~mol}^{-1}$, at $25^{\circ} \mathrm{C}$ and 1.00 bar is

$$
\begin{aligned}
Z_{\mathrm{w}} & =\frac{\mathrm{kg} \mathrm{~m}^{-1} \mathrm{~s}^{-2}}{\left\{2 \pi \times \frac{32.00 \times 10^{-3} \mathrm{~kg} \mathrm{~mol}^{-1}}{6.022 \times 10^{23} \mathrm{~mol}^{-1}} \times\left(1.381 \times 10^{-23} \mathrm{JK}^{-1}\right) \times(298 \mathrm{~K})\right\}^{1 / 2}} \\
& =2.70 \times 10^{27} \mathrm{~m}^{-2} \mathrm{~s}^{-1}
\end{aligned}
$$

This flux corresponds to $0.45 \mathrm{~mol} \mathrm{~cm}^{-2} \mathrm{~s}^{-1}$.

\subsection*{(a) The diffusion coefficient}
The first application of the result in eqn 16A.7b is to use it to find an expression for the net flux of molecules arising from a concentration gradient.

How is that done? 16A. 2\\
Deriving an equation for the net flux of matter

Consider the arrangement depicted in Fig. 16A.4. The molecules passing through the area $A$ at $z=0$ have travelled an average of about one mean free path, $\lambda$, since their last collision.\\
\includegraphics[max width=\textwidth, center]{2025_02_27_d4b95d9718f0b9e2e6b9g-725(1)}

Figure 16A.4 The calculation of the rate of diffusion of a gas considers the net flux of molecules through a plane of area $A$ as a result of molecules arriving from an average distance $\lambda$ away in each direction.

Step 1 Set up expressions for the flux in each direction\\
If the number density at $z$ is $\mathcal{N}(z)$, then the number density at $z=\lambda$ can be estimated by using a Taylor expansion of the form $f(x)=f(0)+(\mathrm{d} f / \mathrm{d} x)_{0} x+\cdots$, truncated after the second term (see The chemist's toolkit 12 in Topic 5B):

$$
\mathcal{N}(+\lambda)=\mathcal{N}(0)+\lambda\left(\frac{\mathrm{d} \mathcal{N}}{\mathrm{~d} z}\right)_{0}
$$

Similarly, the number density at $z=-\lambda$ is

$$
\mathcal{N}(-\lambda)=\mathcal{N}(0)-\lambda\left(\frac{\mathrm{d} \mathcal{N}}{\mathrm{~d} z}\right)_{0}
$$

The average number of impacts on the imaginary region of area $A$ during an interval $\Delta t$ is $Z_{\mathrm{w}} A \Delta t$, where $Z_{\mathrm{w}}$ is the collision flux. Therefore, the matter flux from left to right, $J(\mathrm{~L} \rightarrow \mathrm{R})$, arising from the supply of molecules on the left, is this number of collisions divided by the time interval and the area:

$$
J(\mathrm{~L} \rightarrow \mathrm{R})=\frac{Z_{\mathrm{W}} A \Delta t}{A \Delta t}=Z_{\mathrm{W}}=\frac{1}{4} \mathcal{N}(-\lambda) v_{\text {mean }}
$$

The number density is that at $z=-\lambda$, where the molecules originate before striking the area. There is also a flux of molecules from right to left. The molecules making this journey have originated from $z=+\lambda$ where the number density is $\mathcal{N}(\lambda)$. Therefore,

$$
J(\mathrm{~L} \leftarrow \mathrm{R})=\frac{1}{4} \mathcal{N}(\lambda) v_{\text {mean }}
$$

\subsection*{Step 2 Evaluate the net flux}
The net flux from left to right is

$$
\begin{aligned}
J_{z} & =J(\mathrm{~L} \rightarrow \mathrm{R})-J(\mathrm{~L} \leftarrow \mathrm{R}) \\
& =\frac{1}{4} v_{\text {mean }}\{\mathcal{N}(-\lambda)-\mathcal{N}(\lambda)\} \\
& =\frac{1}{4} v_{\text {mean }}\left\{\left[\mathcal{N}(0)-\lambda\left(\frac{\mathrm{d} \mathcal{N}}{\mathrm{~d} z}\right)_{0}\right]-\left[\mathcal{N}(0)+\lambda\left(\frac{\mathrm{d} \mathcal{N}}{\mathrm{~d} z}\right)_{0}\right]\right\}
\end{aligned}
$$

That is,


\begin{equation*}
-J_{z}=-\frac{1}{2} v_{\operatorname{mean}} \lambda\left(\frac{\mathrm{d} \mathcal{N}}{\mathrm{~d} z}\right)_{0} \tag{16A.8}
\end{equation*}


This equation shows that the flux is proportional to the gradient of the concentration, in agreement with the empirical observation expressed by Fick's law, eqn 16A.2.

At this stage it looks as though a value of the diffusion coefficient can be picked out by comparing eqns 16A. 8 and 16A.4, so obtaining $D=\frac{1}{2} \lambda v_{\text {mean }}$. It must be remembered, however, that the calculation is quite crude, and is little more than an assessment of the order of magnitude of $D$. One aspect that has not been taken into account is illustrated in Fig. 16A.5, which shows that although a molecule may have begun its journey very close to the window, it could have a long flight before it gets there. Because the path is long, the molecule is likely to collide before reaching the window, so it ought not to be counted as passing through the window. Taking this effect into account results in the appearance of a factor of $\frac{2}{3}$ representing the lower flux. The modification results in

$$
D=\frac{1}{3} \lambda v_{\text {mean }}
$$

Diffusion coefficient\\
(16A.9)\\
\includegraphics[max width=\textwidth, center]{2025_02_27_d4b95d9718f0b9e2e6b9g-726}

Figure 16A. 5 One issue ignored in the simple treatment is that some molecules might make a long flight to the plane even though they are only a short perpendicular distance away from it. A molecule taking the longer flight has a higher chance of colliding during its journey.

\subsection*{Brief illustration 16A. 3}
In Brief illustration 1B. 3 of Topic 1B it is established that the mean free path of $\mathrm{N}_{2}$ molecules in a gas at 1.0 atm and $25^{\circ} \mathrm{C}$ is 91 nm ; in Example 1B. 1 of the same Topic it is calculated that under the same conditions the mean speed of $\mathrm{N}_{2}$ molecules is $475 \mathrm{~m} \mathrm{~s}^{-1}$. Therefore, the diffusion coefficient for $\mathrm{N}_{2}$ molecules under these conditions is

$$
D=\frac{1}{3} \times\left(91 \times 10^{-9} \mathrm{~m}\right) \times\left(475 \mathrm{~m} \mathrm{~s}^{-1}\right)=1.4 \times 10^{-5} \mathrm{~m}^{2} \mathrm{~s}^{-1}
$$

The experimental value is $1.5 \times 10^{-5} \mathrm{~m}^{2} \mathrm{~s}^{-1}$.

There are three points to note about eqn 16A.9:

\begin{itemize}
  \item The mean free path, $\lambda$, decreases as the pressure is increased (eqn 16A.1a), so $D$ decreases with increasing pressure and, as a result, the gas molecules diffuse more slowly.
  \item The mean speed, $v_{\text {mean }}$, increases with the temperature (eqn 16A.1b), so $D$ also increases with temperature. As a result, molecules in a hot gas diffuse more quickly than those when the gas is cool (for a given concentration gradient).
  \item Because the mean free path increases when the collision cross-section $\sigma$ of the molecules decreases, the diffusion coefficient is greater for small molecules than for large molecules.
\end{itemize}

\subsection*{(b) Thermal conductivity}
According to the equipartition theorem (The chemist's toolkit 7 in Topic 2A), each molecule carries an average energy $\varepsilon=v k T$, where $v$ depends on the number of quadratic contributions to the energy of the molecule and is a number of the order of 1. For atoms, which have three translational degrees of freedom,\\
$v=\frac{3}{2}$. When one molecule passes through the imaginary window, it transports that average energy. An argument similar to that used for diffusion can be used to discuss the transport of energy through this window.

\subsection*{How is that done? 16A. 3 Deriving an expression for the thermal conductivity}
Assume that the number density is uniform but that the temperature, and hence the average energy of the molecules, is not. Molecules arrive from the left after travelling a mean free path from their last collision in a hotter region, and therefore arrive with a higher energy. Molecules also arrive from the right after travelling a mean free path from a cooler region, and hence arrive with lower energy.

Step 1 Write expressions for the forward, reverse, and net energy flux\\
If the average energy of the molecules at $z$ is $\varepsilon(z)$, the two opposing energy fluxes are

$$
J(\mathrm{~L} \rightarrow \mathrm{R})=\overbrace{\frac{1}{4} \mathcal{N} \nu_{\text {mean }}}^{Z_{\mathrm{W}}} \varepsilon(-\lambda) \quad J(\mathrm{~L} \leftarrow \mathrm{R})=\overbrace{\frac{1}{4} \mathcal{N} \nu_{\text {mean }}}^{Z_{\mathrm{W}}} \varepsilon(\lambda)
$$

and the net flux is

$$
\begin{aligned}
J_{z} & =J(\mathrm{~L} \rightarrow \mathrm{R})-J(\mathrm{~L} \leftarrow \mathrm{R}) \\
& =\frac{1}{4} v_{\text {mean }} \mathcal{N}\{\varepsilon(-\lambda)-\varepsilon(\lambda)\} \\
& =\varepsilon(z)=\varepsilon(0)+z(\mathrm{~d} \varepsilon / \mathrm{dz})_{0_{0}, z}= \pm \lambda \\
& =\frac{1}{4} v_{\text {mean }} \mathcal{N}\left\{\left[\varepsilon(0)-\lambda\left(\frac{\mathrm{d} \varepsilon}{\mathrm{~d} z}\right)_{0}\right]-\left[\varepsilon(0)+\lambda\left(\frac{\mathrm{d} \varepsilon}{\mathrm{~d} z}\right)_{0}\right]\right\} \\
& =-\frac{1}{2} v_{\text {mean }} \lambda \mathcal{N}\left(\frac{\mathrm{d} \varepsilon}{\mathrm{~d} z}\right)_{0}
\end{aligned}
$$

Step 2 Express the energy gradient as a temperature gradient\\
Because $\varepsilon=v k T$, a gradient of energy can be expressed as a gradient of temperature: $\mathrm{d} \varepsilon / \mathrm{d} z=v k(\mathrm{~d} T / \mathrm{d} z)$; therefore

$$
J_{z}=-\frac{1}{2} v_{\operatorname{mean}} \lambda \mathcal{N}\left(\frac{\mathrm{d} \varepsilon}{\mathrm{~d} z}\right)_{0}=-\frac{1}{2} v k v_{\operatorname{mean}} \lambda \mathcal{N}\left(\frac{\mathrm{d} T}{\mathrm{~d} z}\right)_{0}
$$

This equation shows that the energy flux is proportional to the temperature gradient, which is the desired result. As before, the constant is multiplied by $\frac{2}{3}$ to take long flight paths into account, and comparison of this equation with eqn 16A.5 shows that\\
$-\left\lvert\, \kappa=\frac{1}{3} \nu k v_{\text {mean }} \lambda \mathcal{N} \longmapsto\right.$ (16A.10a)

The number density can be written in terms of the molar concentration [J] of the carrier particles J: $\mathcal{N}=n N_{\mathrm{A}} / V=$ [J] $N_{\mathrm{A}}$. Next, the quantity $v k N_{\mathrm{A}}$ can be identified as the molar constant-volume heat capacity of the gas: the molar energy is $N_{\mathrm{A}} v k T$, so it follows from $C_{V, \mathrm{~m}}=\left(\partial U_{\mathrm{m}} / \partial T\right)_{V}$ that $C_{V, \mathrm{~m}}=N_{\mathrm{A}} v k$. With these substitutions, eqn 16A.10a becomes

$$
\kappa=\frac{1}{3} v_{\text {mean }} \lambda[\mathrm{J}] C_{V, \mathrm{~m}}
$$

Thermal conductivity\\
(16A.10b)

Yet another form is found by starting with eqn 16A.10a, recognizing that $\mathcal{N}=p / k T$ and then using the expression for $D$ in eqn 16A.9:\\
\includegraphics[max width=\textwidth, center]{2025_02_27_d4b95d9718f0b9e2e6b9g-727}

\begin{displayquote}
Thermal conductivity (16A.10c)
\end{displayquote}

\subsection*{Brief illustration 16A.4}
In Brief illustration 16A. 3 the value $D=1.4 \times 10^{-5} \mathrm{~m}^{2} \mathrm{~s}^{-1}$ was calculated for $\mathrm{N}_{2}$ molecules at $25^{\circ} \mathrm{C}$ and 1.0 atm . The thermal conductivity can be calculated by using eqn 16A.10c and noting that for $\mathrm{N}_{2}$ molecules $v=\frac{5}{2}$ (three translational modes and two rotational modes; the vibrational mode is inactive for this 'stiff' molecule).

$$
\begin{aligned}
\kappa & =\frac{\frac{5}{2} \times(1.01 \times 10^{5} \overbrace{\mathrm{~Pa}}^{\mathrm{Jm}}) \times\left(1.4 \times 10^{-5} \mathrm{~m}^{2} \mathrm{~s}^{-1}\right)}{298 \mathrm{~K}} \\
& =1.2 \times 10^{-2} \mathrm{JK}^{-1} \mathrm{~m}^{-1} \mathrm{~s}^{-1}
\end{aligned}
$$

or $12 \mathrm{~mW} \mathrm{~K}^{-1} \mathrm{~m}^{-1}$. The experimental value is $26 \mathrm{~mW} \mathrm{~K}^{-1} \mathrm{~m}^{-1}$.

\subsection*{To interpret eqn 16A.10, note that:}
\begin{itemize}
  \item The mean free path is $\lambda$ is inversely proportional to the pressure, but the number density $\mathcal{N}$ is proportional to the pressure $(\mathcal{N}=p / k T)$. It follows that the product $\lambda \mathcal{N}$, which appears eqn 16A.10a, is independent of pressure, and so therefore is the thermal conductivity.
  \item The thermal conductivity is greater for gases with a high heat capacity (eqn 16A.10b) because a given temperature gradient then corresponds to a steeper energy gradient.
\end{itemize}

The physical reason for the pressure independence of the thermal conductivity is that the conductivity can be expected to be large when many molecules are available to transport the energy, but the presence of so many molecules limits their mean free path and they cannot carry the energy over a great distance.\\
\includegraphics[max width=\textwidth, center]{2025_02_27_d4b95d9718f0b9e2e6b9g-728}

Figure 16A. 6 The calculation of the viscosity of a gas examines the net $x$-component of momentum brought to a plane from faster and slower layers a mean free path away in each direction.

These two effects cancel. The thermal conductivity is indeed found experimentally to be independent of the pressure, except when the pressure is very low, and then $\kappa \propto p$. At very low pressures it is possible for $\lambda$ to exceed the dimensions of the apparatus, and the distance over which the energy is transported is then determined by the size of the container and not by collisions with the other molecules present. The flux is still proportional to the number of carriers, but the length of the journey no longer depends on $\lambda$, so $\kappa \propto[J]$, which implies that $\kappa \propto p$.

\subsection*{(c) Viscosity}
If the momentum in the $x$-direction depends on $z$ as $m v_{x}(z)$, then molecules moving from the right in Fig. 16A. 6 (from a fast layer to a slower one) transport a momentum $m v_{x}(\lambda)$ to their new layer at $z=0$; those travelling from the left transport $m v_{x}(-\lambda)$ to it. This picture can be used to build an expression for the coefficient of viscosity.

\subsection*{How is that done? 16A. 4}
Deriving an expression for\\
the viscosity\\
The strategy is the same as in the previous derivations, but now the property being transported is the linear momentum of the layers.

Step 1 Set up expressions for the flux of momentum in each direction and the net flux\\
Molecules arriving from the right bring a momentum

$$
m v_{x}(\lambda)=m v_{x}(0)+m \lambda\left(\frac{\mathrm{~d} v_{x}}{\mathrm{~d} z}\right)_{0}
$$

Those arriving from the left bring a momentum

$$
m v_{x}(-\lambda)=m v_{x}(0)-m \lambda\left(\frac{\mathrm{~d} v_{x}}{\mathrm{~d} z}\right)_{0}
$$

The net flux of $x$-momentum in the $z$-direction is therefore

$$
\begin{aligned}
J_{z} & =\frac{1}{4} v_{\text {mean }} \mathcal{N}\left\{\left[m v_{x}(0)-m \lambda\left(\frac{\mathrm{~d} v_{x}}{\mathrm{~d} z}\right)_{0}\right]-\left[m v_{x}(0)+m \lambda\left(\frac{\mathrm{~d} v_{x}}{\mathrm{~d} z}\right)_{0}\right]\right\} \\
& =-\frac{1}{2} v_{\text {mean }} \lambda m \mathcal{N}\left(\frac{\mathrm{~d} v_{x}}{\mathrm{~d} z}\right)_{0}
\end{aligned}
$$

Step 2 Identify the coefficient of viscosity\\
The flux is proportional to the velocity gradient, in line with the phenomenological equation. Comparison of this expression with eqn 16A.6, and multiplication by $\frac{2}{3}$ in the normal way, leads to\\
\includegraphics[max width=\textwidth, center]{2025_02_27_d4b95d9718f0b9e2e6b9g-728(1)}

Two alternative forms of this expression (after using $m N_{\mathrm{A}}=M$ and eqn 16A.9, $\left.D=\frac{1}{3} \lambda v_{\text {mean }}\right)$ are


\begin{align*}
& \eta=M D[\mathrm{~J}]  \tag{16A.11b}\\
& \eta=\frac{p M D}{R T} \tag{16A.11c}
\end{align*}


where [J] is the molar concentration of the gas molecules J and $M$ is their molar mass.

\subsection*{Brief illustration 16A. 5}
From Brief illustration 16A. 3 the value of $D$ for $\mathrm{N}_{2}$ molecules at $25^{\circ} \mathrm{C}$ and 1.0 atm is $1.4 \times 10^{-5} \mathrm{~m}^{2} \mathrm{~s}^{-1}$. Because $M=28.02 \mathrm{~g} \mathrm{~mol}^{-1}$, eqn 16A.11c gives

$$
\begin{aligned}
\eta & =\frac{\left(1.01 \times 10^{5} \mathrm{Ja}^{\mathrm{Pa}}\right) \times\left(28.02 \times 10^{-3} \mathrm{~kg} \mathrm{~mol}^{-1}\right) \times\left(1.4 \times 10^{-5} \mathrm{~m}^{2} \mathrm{~s}^{-1}\right)}{\left(8.3145 \mathrm{JK}^{-1} \mathrm{~mol}^{-1}\right) \times(298 \mathrm{~K})} \\
& =1.6 \times 10^{-5} \mathrm{~kg} \mathrm{~m}^{-1} \mathrm{~s}^{-1}
\end{aligned}
$$

or $160 \mu \mathrm{P}$. The experimental value is $177 \mu \mathrm{P}$.

The physical interpretation of eqn 16A.11a is as follows:

\begin{itemize}
  \item As has already been noted for thermal conductivity, the product $\lambda \mathcal{N}$ is independent of $p$. Therefore, like the thermal conductivity, the viscosity is independent of the pressure.
  \item Because $v_{\text {mean }} \propto T^{1 / 2}, \eta \propto T^{1 / 2}$ at constant volume (and $\eta \propto T^{3 / 2}$ at constant pressure). That is, the viscosity of a gas increases with temperature.
\end{itemize}

The physical reason for the pressure-independence of the viscosity is the same as for the thermal conductivity: more molecules are available to transport the momentum, but they carry it less far on account of the decrease in mean free path.

The increase of viscosity with temperature is explained when it is recalled that at high temperatures the molecules travel more quickly, so the flux of momentum is greater. In contrast, as discussed in Topic 16B, the viscosity of a liquid decreases with increase in temperature because intermolecular interactions must be overcome.

\subsection*{(d) Effusion}
Because the mean speed of molecules is inversely proportional to $M^{1 / 2}$, the rate at which they strike the area of the hole through which they are effusing is also inversely proportional to $M^{1 / 2}$, in accord with Graham's law. However, by using the expression for the collision flux, a more detailed expression for the rate of effusion can be obtained and used to interpret effusion data in a more sophisticated way.

When a gas at a pressure $p$ and temperature $T$ is separated from a vacuum by a small hole, the rate of escape of its molecules is equal to the rate at which they strike the area of the hole, which is the product of the collision flux and the area of the hole, $A_{0}$ :

$$
\text { Rate of effusion }=Z_{\mathrm{W}} A_{0}=\frac{p A_{0}}{(2 \pi m k T)^{1 / 2}}=\frac{p A_{0} N_{\mathrm{A}}}{(2 \pi M R T)^{1 / 2}}
$$

This rate is inversely proportional to $M^{1 / 2}$, in accord with Graham's law. Do not conclude, however, that because the expression includes a factor of $T^{-1 / 2}$ the rate of effusion decreases as the temperature increases. Because $p \propto T$, the rate is in fact proportional to $T^{1 / 2}$ and increases with temperature.

Equation 16A. 12 is the basis of the Knudsen method for the determination of the vapour pressures of liquids and solids, particularly of substances with very low vapour pressures and which cannot be measured directly. Thus, if the vapour pressure of a sample is $p$, and it is enclosed in a cavity with a small hole, then the rate of loss of mass from the container is proportional to $p$.

Example 16A. 1 Calculating the vapour pressure from a mass loss

Caesium (m.p. $29^{\circ} \mathrm{C}$, b.p. $686^{\circ} \mathrm{C}$ ) was introduced into a container and heated to $500^{\circ} \mathrm{C}$; the container is pierced by a hole of diameter 1.0 mm . It is found that in $1.00 \mathrm{~h}(3600 \mathrm{~s})$ the mass of the container decreased by 84.4 mg . Calculate the vapour pressure of liquid caesium at $500^{\circ} \mathrm{C}$.

Collect your thoughts The pressure of vapour is constant inside the container despite the effusion of atoms because the hot liquid metal replenishes the vapour. The rate of effusion is therefore constant, and given by eqn 16A.12. To express the rate in terms of mass, multiply the number of atoms that escape by the mass of each atom.

The solution The mass loss $\Delta m$ in an interval $\Delta t$ is equal to the number of molecules that strike the area of the hole in this interval multiplied by the mass of each molecule:

$$
\Delta m=\text { rate of effusion } \times \Delta t \times m=\overbrace{\frac{p A_{0} N_{\mathrm{A}}}{(2 \pi M R T)^{1 / 2}}}^{\text {eqn } 16 \mathrm{~A} .12} \times \Delta t \times m
$$

where $A_{0}$ is the area of the hole and $M$ is the molar mass. Rearrangement of this equation gives an expression for $p$ :

$$
p=\left(\frac{2 \pi R T}{M}\right)^{1 / 2} \frac{\Delta m}{A_{0} \Delta t}
$$

Substitution of the data and $M=132.9 \mathrm{~g} \mathrm{~mol}^{-1}$ gives

$$
\begin{aligned}
p= & \left(\frac{2 \pi \times\left(8.3145 \mathrm{JK}^{-1} \mathrm{~mol}^{-1}\right) \times(773 \mathrm{~K})}{0.1329 \mathrm{~kg} \mathrm{~mol}^{-1}}\right)^{1 / 2} \times \\
& \underbrace{\frac{84.4 \times 10^{-3} \mathrm{~kg}}{\pi \times\left(0.50 \times 10^{-3} \mathrm{~m}\right)^{2}} \times(3600 \mathrm{~s})}_{A_{0}}=1.6 \times 10^{4} \mathrm{~Pa} \text { or } 16 \mathrm{kPa}
\end{aligned}
$$

Self-test 16A.1 How long would it take for 200 mg of Cs atoms to effuse out of the oven under the same conditions?\\
\includegraphics[max width=\textwidth, center]{2025_02_27_d4b95d9718f0b9e2e6b9g-729}

\subsection*{Checklist of concepts}
\begin{enumerate}
  \item Flux is the quantity of a property passing through a given area in a given time interval divided by the area and the duration of the interval.
  \item Diffusion is the migration of matter down a concentration gradient.
  \item Fick's first law of diffusion states that the flux of matter is proportional to the concentration gradient.
  \item Thermal conduction is the migration of energy down a temperature gradient; the flux of energy is proportional to the temperature gradient.5. Viscosity is the migration of linear momentum down a velocity gradient; the flux of momentum is proportional to the velocity gradient.7. Graham's law of effusion states that the rate of effusion is inversely proportional to the square root of the molar mass.6. Effusion is the emergence of a gas from a container through a small hole.
\end{enumerate}

\subsection*{Checklist of equations}
\begin{center}
\begin{tabular}{llll}
\hline
Property & Equation & Comment & Equation number \\
\hline
Fick's first law of diffusion & $J=-D \mathrm{~d} \mathcal{N} / \mathrm{d} z$ &  & 16 A .4 \\
Flux of energy of thermal motion & $J=-\kappa \mathrm{d} T / \mathrm{d} z$ & 16 A .5 &  \\
Flux of momentum along $x$ & $J=-\eta \mathrm{d} v_{x} / \mathrm{d} z$ & 16 A .6 &  \\
Diffusion coefficient of a perfect gas & $D=\frac{1}{3} \lambda v_{\text {mean }}$ & $\mathrm{KMT}^{*}$ & 16 A .9 \\
Coefficient of thermal conductivity of a perfect gas & $\kappa=\frac{1}{3} v_{\text {mean }} \lambda[\mathrm{J}] C_{V, \mathrm{~m}}$ & KMT and equipartition & 16 A .10 b \\
Coefficient of viscosity of a perfect gas & $\eta=\frac{1}{3} v_{\text {mean }} \lambda m \mathcal{N}$ & KMT & 16 A .11 a \\
Rate of effusion & Rate $\propto 1 / M^{1 / 2}$ & Graham's law & 16 A .12 \\
\hline
\end{tabular}
\end{center}

\begin{itemize}
  \item KMT indicates that the equation is based on the kinetic theory of gases.
\end{itemize}

\section*{TOPIC 16B Motion in liquids}
\subsection*{Why do you need to know this material?}
Many chemical reactions take place in liquids, so for a full understanding of them it is important to know how solute molecules and ions move through such environments.

\subsection*{What is the key idea?}
lons reach a terminal velocity when the electrical force on them is balanced by the drag due to the viscosity of the solvent.

\subsection*{What do you need to know already?}
The discussion of viscosity starts with the definition of the coefficient of viscosity introduced in Topic 16A. Some calculations make use of the information about electrostatics set out in The chemist's toolkit 29.

The motion of ions and molecules is an important aspect of the properties of liquids and of the reactions taking place in them. It can be studied experimentally by a variety of methods. For example, bulk measurements of viscosity and its temperature dependence can be used to build models of the motion. At a more detailed level, relaxation time measurements in NMR (Topic 12C) and EPR can be used to show how molecules rotate. For instance, these observations show that large molecules in viscous fluids typically rotate in a series of small (about $5^{\circ}$ ) steps, whereas small molecules in non-viscous fluids typically jump through about 1 radian $\left(57^{\circ}\right)$ in each step. Another important technique is inelastic neutron scattering, in which the energy neutrons collect or discard as they pass through a sample is interpreted in terms of the motion of its molecules.

\subsection*{16B. 1 Experimental results}
There are two 'classical' methods of investigating molecular motion in liquids. One is the measurement of viscosity and its temperature dependence. Another involves inferring details about molecular motion by dragging ions through a solvent under the influence of an electric field.

\subsection*{(a) Liquid viscosity}
The coefficient of viscosity, $\eta$ (eta), is introduced in Topic 16A as the constant of proportionality in the relation between the flux of linear momentum and the velocity gradient in a fluid:


\begin{equation*}
J_{z}(x \text {-component of momentum })=-\eta \frac{\mathrm{d} v_{x}}{\mathrm{~d} z} \tag{16B.1}
\end{equation*}


The units of viscosity are kilograms per metre per second ( $\mathrm{kg} \mathrm{m}^{-1} \mathrm{~s}^{-1}$ ), but they may also be reported in the equivalent units of pascal seconds ( Pa s). The non-SI units poise ( P ) and centipoise ( cP ) are still widely encountered: $1 \mathrm{P}=10^{-1} \mathrm{Pas}$ and so $1 \mathrm{cP}=1 \mathrm{mPa}$. Table 16B. 1 lists some values of $\eta$ for liquids.

Unlike in a gas, for a molecule to move in a liquid it must acquire at least a minimum energy (an 'activation energy', $E_{\mathrm{a}}$, in the language of Topic 17D) to escape from its neighbours. From the Boltzmann distribution it follows that the probability that a molecule has at least an energy $E_{\mathrm{a}}$ is proportional to $\mathrm{e}^{-E_{a} / R T}$, so the mobility of the molecules in the liquid should follow this type of temperature dependence. As the temperature increases, the molecules become more mobile and so the viscosity is reduced; the expected temperature dependence of the viscosity, which decreases as the mobility of the molecules increases, is therefore of the form


\begin{equation*}
\eta=\eta_{0} \mathrm{e}^{E_{\mathrm{a}} / R T} \quad \text { Temperature dependence of viscosity (liquid) } \tag{16B.2}
\end{equation*}


(Note the positive sign of the exponent.) The activation energy typical of viscosity is comparable to the mean potential energy of intermolecular interactions. Equation 16B. 2 implies that the viscosity should decrease sharply with increasing temperature. Such a variation is found experimentally, at

Table 16B. 1 Viscosities of liquids at $298 \mathrm{~K}^{*}$

\begin{center}
\begin{tabular}{ll}
\hline
 & $\eta /\left(10^{-3} \mathrm{~kg} \mathrm{~m}^{-1} \mathrm{~s}^{-1}\right)$ \\
\hline
Benzene & 0.601 \\
Mercury & 1.55 \\
Pentane & 0.224 \\
Water $^{\ddagger}$ & 0.891 \\
\end{tabular}
\end{center}

\footnotetext{\begin{itemize}
  \item More values are given in the Resource section.\\
${ }^{\ddagger}$ Note that $1 \mathrm{cP}=10^{-3} \mathrm{~kg} \mathrm{~m}^{-1} \mathrm{~s}^{-1}$; the viscosity of water corresponds to 0.891 cP .
\end{itemize}
}
\includegraphics[max width=\textwidth, center]{2025_02_27_d4b95d9718f0b9e2e6b9g-732}

Figure 16B. 1 The temperature dependence of the viscosity of water. As the temperature is increased, more molecules are able to escape from the potential wells provided by their neighbours, and so the liquid becomes less viscous.\\
least over reasonably small temperature ranges (Fig. 16B.1). Intermolecular interactions govern the magnitude of $E_{\mathrm{a}}$, but the calculation of its value from an assumed form of the potential energy due to the interactions between the molecules is an immensely difficult and largely unsolved problem.

\subsection*{Brief illustration 16B. 1}
The viscosity of water at $25^{\circ} \mathrm{C}$ and $50^{\circ} \mathrm{C}$ is $0.890 \mathrm{mPa} s$ and 0.547 mPa s , respectively. It follows from eqn 16B. 2 that the activation energy for molecular migration is the solution of

$$
\frac{\eta\left(T_{2}\right)}{\eta\left(T_{1}\right)}=\mathrm{e}^{\left(E_{\mathrm{a}} / R\right)\left(1 / T_{2}-1 / T_{1}\right)}
$$

which, after taking logarithms of both sides, is

$$
\begin{aligned}
E_{\mathrm{a}} & =\frac{R \ln \left\{\eta\left(T_{2}\right) / \eta\left(T_{1}\right)\right\}}{1 / T_{2}-1 / T_{1}} \\
& =\frac{\left(8.3145 \mathrm{JK}^{-1} \mathrm{~mol}^{-1}\right) \ln \{(0.547 \mathrm{mPas}) /(0.890 \mathrm{mPas})\}}{1 /(323 \mathrm{~K})-1 /(298 \mathrm{~K})} \\
& =1.56 \times 10^{4} \mathrm{Jmol}^{-1}
\end{aligned}
$$

or $15.6 \mathrm{~kJ} \mathrm{~mol}^{-1}$. This value is comparable to the strength of a hydrogen bond.

One problem with the interpretation of viscosity measurements is that the change in density of the liquid as it is heated makes a pronounced contribution to the temperature variation of the viscosity. Thus, the temperature dependence of viscosity at constant volume, when the density is constant, is much less than that at constant pressure. At low temperatures, the viscosity of water decreases as the pressure is increased. This behaviour is consistent with the need to rupture hydrogen bonds for migration to occur.

\subsection*{(b) Electrolyte solutions}
When a potential difference is applied between two electrodes immersed in a solution containing ions there is a flow of current due to the migration of ions through the solution. The key electrical property of the solution is its resistance, $R$, which is expressed in ohms, $\Omega\left(1 \Omega=1 \mathrm{C}^{-1} \mathrm{~V}\right.$ s). It is often more convenient to work in terms of the conductance, $G$, which is the inverse of the resistance: $G=1 / R$, and therefore expressed in $\Omega^{-1}$. The reciprocal ohm used to be called the mho, but its SI designation is now the siemens, S , and $1 \mathrm{~S}=1 \Omega^{-1}=1 \mathrm{CV}^{-1} \mathrm{~s}^{-1}$. Electric current is expressed in amperes, A , with $1 \mathrm{~A}=1 \mathrm{C} \mathrm{s}^{-1}$, so a more physically revealing relation is $1 \mathrm{~S}=1 \mathrm{AV}^{-1}$.

The conductance of a sample is proportional to its crosssectional area, $A$, and inversely proportional to its length, $l$. Therefore

\[
G=\kappa \frac{A}{l} \quad \begin{array}{ll}
\text { Conductivity }  \tag{16B.3}\\
\text { [definition] }
\end{array}
\]

where the constant of proportionality $\kappa$ (kappa) is the electrical conductivity of the sample. With the conductance in siemens and the dimensions in metres, it follows that the units of $\kappa$ are siemens per metre $\left(\mathrm{Sm}^{-1}\right)$. The conductivity is a property of the material, whereas the conductance depends both on the material and its dimensions. The conductivity depends on the concentration of charge carriers in the sample, which suggests that it is sensible to introduce the molar conductivity, $\Lambda_{\mathrm{m}}$, defined as

\[
\Lambda_{\mathrm{m}}=\frac{\kappa}{c} \quad \begin{align*}
& \text { Molar conductivity }  \tag{16B.4}\\
& \text { [definition] }
\end{align*}
\]

where $c$ is the molar concentration of the added electrolyte. The units of molar conductivity are siemens metre-squared per mole ( $\mathrm{Sm}^{2} \mathrm{~mol}^{-1}$ ), and typical values are about $10 \mathrm{mS} \mathrm{m}^{2} \mathrm{~mol}^{-1}$ (where $1 \mathrm{mS}=10^{-3} \mathrm{~S}$ ).

Experimentally, it is found that the value of the molar conductivity varies with the concentration. One reason for this variation is that the number of ions in the solution might not be proportional to the nominal concentration of the electrolyte. For instance, the concentration of ions in a solution of a weak electrolyte depends on the degree of dissociation, which is a complicated function of the total amount of solute: doubling the total concentration of the solute does not double the number of ions. Secondly, even for fully dissociated strong electrolytes, because ions interact strongly with one another, the conductivity of a solution is not exactly proportional to the number of ions present.

In an extensive series of measurements during the nineteenth century, Friedrich Kohlrausch established the Kohlrausch law, that at low concentrations the molar conductivities of strong electrolytes depend on the square root of the concentration:

$$
\Lambda_{\mathrm{m}}=\Lambda_{\mathrm{m}}^{\circ}-K c^{1 / 2}
$$

Kohlrausch law\\
(16B.5)\\
where $\Lambda_{\mathrm{m}}^{\circ}$ is the limiting molar conductivity, the molar conductivity in the limit of zero concentration when the ions are so far apart that they move independently of each other. Kohlrausch also established that this limiting molar conductivity is the sum of contributions from the individual ions present. If the limiting molar conductivity of the cations is denoted $\lambda_{+}$and that of the anions $\lambda_{-}$, then his law of the independent migration of ions states that

$$
\Lambda_{\mathrm{m}}^{\circ}=v_{+} \lambda_{+}+v_{-} \lambda_{-} \quad \begin{aligned}
& \text { Law of the independent } \\
& \text { migration of ions }
\end{aligned}
$$

(16B.6)\\
where $v_{+}$and $v_{-}$are the numbers of cations and anions provided by each formula unit of electrolyte. For example, for $\mathrm{HCl}, \mathrm{NaCl}$, and $\mathrm{CuSO}_{4}, v_{+}=v_{-}=1$, but for $\mathrm{MgCl}_{2} v_{+}=1, v_{-}=2$.

\subsection*{Example 16B.1 Determining the limiting molar conductivity}
The conductivity of $\mathrm{KCl}(\mathrm{aq})$ at $25^{\circ} \mathrm{C}$ is $14.688 \mathrm{mS} \mathrm{m}^{-1}$ when $c=$ $1.0000 \mathrm{mmol} \mathrm{dm}^{-3}$, and $71.740 \mathrm{mS} \mathrm{m}^{-1}$ when $c=5.0000 \mathrm{mmol}$ $\mathrm{dm}^{-3}$. Determine the values of the limiting molar conductivity $\Lambda_{\mathrm{m}}^{\circ}$ and the Kohlrausch constant $\mathcal{K}$.

Collect your thoughts You need to use eqn 16B. 4 to determine the molar conductivities at the two concentrations. Then, by using eqn 16B.5, you can express the difference between these two values as $\Lambda_{\mathrm{m}}\left(c_{2}\right)-\Lambda_{\mathrm{m}}\left(c_{1}\right)=K\left(c_{1}^{1 / 2}-c_{2}^{1 / 2}\right)$. From this relation you can determine $\mathcal{K}$ and then go on to find $\Lambda_{\mathrm{m}}^{\circ}$ by using one of the values of the molar conductivity in eqn 16B.5, rearranged into $\Lambda_{\mathrm{m}}^{\circ}=\Lambda_{\mathrm{m}}+K c^{1 / 2}$.\\
The solution The molar conductivity of $\operatorname{KCl}(\mathrm{aq})$ when $c=1.0000 \mathrm{mmol} \mathrm{dm}^{-3}$ (which is the same as $1.0000 \mathrm{~mol} \mathrm{~m}^{-3}$ ) is

$$
\Lambda_{\mathrm{m}}=\frac{14.688 \mathrm{mSm}^{-1}}{1.0000 \mathrm{molm}^{-3}}=14.688 \mathrm{mSm}^{2} \mathrm{~mol}^{-1}
$$

Similarly, when $c=5.0000 \mathrm{~mol} \mathrm{dm}^{-3}$ its molar conductivity is $14.348 \mathrm{mS} \mathrm{m}^{2} \mathrm{~mol}^{-1}$. It then follows that

$$
\begin{aligned}
K & =\frac{\Lambda_{\mathrm{m}}\left(c_{2}\right)-\Lambda_{\mathrm{m}}\left(c_{1}\right)}{c_{1}^{1 / 2}-c_{2}^{1 / 2}}=\frac{(14.348-14.688) \mathrm{mSm}^{2} \mathrm{~mol}^{-1}}{\left(0.0010000^{1 / 2}-0.0050000^{1 / 2}\right)\left(\mathrm{mol} \mathrm{dm}^{-3}\right)^{1 / 2}} \\
& =8.698 \mathrm{mSm}^{2} \mathrm{~mol}^{-1} /\left(\mathrm{mol} \mathrm{dm}^{-3}\right)^{1 / 2}
\end{aligned}
$$

(For reasons that will become clear immediately below, it is best to keep this awkward but convenient array of units rather than converting them to the equivalent $10^{-3 / 2} \mathrm{Sm}^{7 / 2} \mathrm{~mol}^{-3 / 2}$.) The limiting molar conductivity is then found by using the data for $c=1.0000 \mathrm{mmol} \mathrm{dm}^{-3}$ :

$$
\begin{aligned}
\Lambda_{\mathrm{m}}^{\circ}= & 14.688 \mathrm{mSm}^{2} \mathrm{~mol}^{-1}+8.698 \frac{\mathrm{mSm}^{2} \mathrm{~mol}^{-1}}{\left(\mathrm{~mol} \mathrm{dm}^{-3}\right)^{1 / 2}} \\
& \times\left(1.0000 \times 10^{-3} \mathrm{moldm}^{-3}\right)^{1 / 2}=14.963 \mathrm{mSm}^{2} \mathrm{~mol}^{-1}
\end{aligned}
$$

Comment. Although the value of $\mathcal{K}$ has been given to four significant figures in conformity with the data, that degree of precision is probably over-optimistic in practice.

Self-test 16 B .1 The conductivity of $\mathrm{KClO}_{4}(\mathrm{aq})$ at $25^{\circ} \mathrm{C}$ is $13.780 \mathrm{mS} \mathrm{m}^{-1}$ when $c=1.000 \mathrm{mmol} \mathrm{dm}^{-3}$ and $67.045 \mathrm{mS} \mathrm{m}^{-1}$ when $c=5.000 \mathrm{mmol} \mathrm{dm}^{-3}$. Determine the values of the limiting molar conductivity $\Lambda_{\mathrm{m}}^{\circ}$ and the Kohlrausch constant $\mathcal{K}$ for this system.

\subsection*{16B.2 The mobilities of ions}
The reason why different ions have different molar conductivities in solution can be understood by analysing the motion of an ion subject to an electric field and at the same time surrounded by a viscous medium.

\subsection*{(a) The drift speed}
An ion in a vacuum is accelerated by an electric field, but in a viscous liquid the motion of the ion is impeded by the need for it to push its way through the tightly packed solvent molecules. The latter effect is called viscous drag. As the ion accelerates under the influence of the field, the viscous drag increases and the ion quickly reaches a steady terminal speed, called the drift speed, which can be found by balancing the two forces.

\subsection*{How is that done? 16B. 1}
Deriving an expression for\\
the drift speed\\
The starting point for this calculation is the result from electrostatics that when the potential difference between two planar electrodes a distance $l$ apart is $\Delta \phi$, the ions in the solution between them experience a uniform electric field of magnitude $E=\Delta \phi / l$. Here, and throughout this section, the sign of the charge number is disregarded so as to avoid notational complications.

\subsection*{Step 1 Find the force on the ion due to the field}
In an electric field $\mathcal{E}$, an ion of charge ze experiences a force of magnitude $z e E$ (see The chemist's toolkit 29). Therefore,

$$
F_{\text {electric }}=\frac{z e \Delta \phi}{l}
$$

Step 2 Find the force on the ion due to viscous drag\\
As the ion moves through the solvent it experiences a frictional retarding force proportional to its speed. For a spherical particle of radius $a$ travelling at a speed $s$, this force is given by Stokes' law, which Stokes derived by considering the hydrodynamics of the passage of a sphere through a continuous fluid:

$$
F_{\text {viscous }}=f_{s} \quad f=6 \pi \eta a \quad \text { Stokes'law }
$$

(16B.7)\\
where $\eta$ is the coefficient of viscosity. In this calculation it is assumed that Stokes' law applies on a molecular scale;\\
experimental evidence suggests that it often gives at least the right order of magnitude for the viscous force.

\subsection*{Step 3 Find the drift speed by balancing the two forces}
The two forces act in opposite directions and the ions quickly reach a terminal speed, the drift speed, $s$, when they are in balance. This balance occurs when $f s=z e \mathcal{E}$, and therefore


\begin{equation*}
\left.\eta s=\frac{z e E}{f} \right\rvert\, \tag{16B.8a}
\end{equation*}


Equation 16B.8a shows that the drift speed is proportional to the electric field strength. The constant of proportionality is called the mobility of the ion, $u$ :

$$
\begin{array}{ll}
s=u \mathcal{E} & \begin{array}{l}
\text { Mobility } \\
\text { [definition] }
\end{array}
\end{array}
$$

(16B.8b)

With the electric field strength in volts per metre $\left(\mathrm{Vm}^{-1}\right)$ and the drift speed in metres per second ( $\mathrm{m} \mathrm{s}^{-1}$ ) the slightly awkward units of $u$ are metres-squared per volt per second $\left(\mathrm{m}^{2} \mathrm{~V}^{-1} \mathrm{~s}^{-1}\right.$; note that $\mathrm{m}^{2} \mathrm{~V}^{-1} \mathrm{~s}^{-1} \times \mathrm{Vm}^{-1}=\mathrm{ms}^{-1}$ ); a selection of values is given in Table 16B.2. Comparison of the last two equations shows that


\begin{equation*}
u=\frac{z e}{f}=\frac{z e}{6 \pi \eta a} \tag{16B.9}
\end{equation*}


where the Stokes' law value for the frictional coefficient $f$ has been used.

\subsection*{Brief illustration 16B.2}
An order-of-magnitude estimate of the mobility can be found using eqn 16B. 9 with $z=1$ and $a=130 \mathrm{pm}$, which is typical of the radius of a hydrated ion; the viscosity of water at $25^{\circ} \mathrm{C}$ is 0.9 cP , or 0.9 mPa . Then

$$
u=\frac{\mathrm{JV}^{-1}}{6 \pi \times(0.9 \times 10^{-3} \underbrace{}_{\mathrm{Jm}} \mathrm{~m}^{-3}) \times\left(10^{-19} \stackrel{-}{\mathrm{C}}\right.} \mathrm{Pas}^{\mathrm{Pas}} \times 10^{-12} \mathrm{~m}) \quad 7.3 \times 10^{-8} \mathrm{~m}^{2} \mathrm{~V}^{-1} \mathrm{~s}^{-1}
$$

This value means that when there is a potential difference of 1.0 V across a solution of length 1.0 cm (so $E=100 \mathrm{~V} \mathrm{~m}^{-1}$ ), the drift speed is $7.3 \mu \mathrm{~m} \mathrm{~s}^{-1}$. That speed might seem slow, but not when expressed on a molecular scale, because it corresponds to an ion passing about $10^{4}$ solvent molecules per second.

Equation 16B. 9 implies that the mobility of an ion decreases with increasing solution viscosity and ion size. Experiments

\subsection*{The chemist's toolkit 29}
\subsection*{Electrostatics}
A charge $Q_{1}$ (units: coulomb, C) gives rise to a Coulomb potential $\phi$ (units: volt, V ). The potential energy (units: joule, J, with $1 \mathrm{~J}=1 \mathrm{VC}$ ) of a second charge $Q$ in that potential is

$$
E_{\mathrm{P}}=-Q \phi
$$

In one dimension, the electric field strength (units: volt per metre, $\mathrm{V} \mathrm{m}^{-1}$ ), $\mathcal{E}$, is the negative of the gradient of the electric potential $\phi$ :

$$
\mathcal{E}=-\frac{\mathrm{d} \phi}{\mathrm{~d} x}
$$

Electric field strength\\
In three dimensions the electric field is a vector, and

$$
\mathcal{E}=-\nabla \phi
$$

The electric field between two plane parallel plates separated by a distance $l$, and between which there is a potential difference $\Delta \phi$, is uniform and given by

$$
\mathcal{E}=-\frac{\Delta \phi}{l}
$$

A charge $Q$ experiences a force proportional to the electric field strength at its location:

$$
F_{\text {electric }}=Q \mathcal{E}
$$

A potential gives rise to a force only if it varies with distance.\\
confirm these predictions for bulky ions (such as $\mathrm{R}_{4} \mathrm{~N}^{+}$and $\mathrm{RCO}_{2}^{-}$) but not for small ions. For example, the mobilities of the alkali metal ions in water increase from $\mathrm{Li}^{+}$to $\mathrm{Rb}^{+}$ (Table 16B.2) even though the ionic radii increase. The paradox is resolved when it is realized that the radius $a$ in the Stokes formula is the hydrodynamic radius (or 'Stokes radius') of the ion, its effective radius in the solution taking into account all the $\mathrm{H}_{2} \mathrm{O}$ molecules it carries in its hydration shell. Small ions

Table 16B.2 Ionic mobilities in water at $298 \mathrm{~K}^{*}$

\begin{center}
\begin{tabular}{llll}
\hline
 & $\boldsymbol{u} /\left(\mathbf{1 0 ^ { - 8 }} \mathrm{m}^{2} \mathbf{V}^{-1} \mathbf{s}^{-1}\right)$ &  & $\boldsymbol{u} /\left(10^{-8} \mathrm{~m}^{2} \mathrm{~V}^{-1} \mathrm{~s}^{-1}\right)$ \\
\hline
$\mathrm{H}^{+}$ & 36.23 & $\mathrm{OH}^{-}$ & 20.64 \\
$\mathrm{Li}^{+}$ & 4.01 & $\mathrm{~F}^{-}$ & 5.70 \\
$\mathrm{Na}^{+}$ & 5.19 & $\mathrm{Cl}^{-}$ & 7.91 \\
$\mathrm{~K}^{+}$ & 7.62 & $\mathrm{Br}^{-}$ & 8.09 \\
$\mathrm{Rb}^{+}$ & 7.92 & $\mathrm{SO}_{4}^{2-}$ & 8.29 \\
\hline
\end{tabular}
\end{center}

\footnotetext{\begin{itemize}
  \item More values are given in the Resource section.
\end{itemize}
}
\includegraphics[max width=\textwidth, center]{2025_02_27_d4b95d9718f0b9e2e6b9g-735}

Figure 16B.2 A highly schematic diagram showing the effective motion of a proton in water.\\
give rise to stronger electric fields than large ones (the electric field at the surface of a sphere of radius $r$ is proportional to $z / r^{2}$, so the smaller the radius the stronger the field). Consequently, small ions are more extensively solvated than big ions and a small ion may have a large hydrodynamic radius because it drags many solvent molecules through the solution as it migrates. The hydrating $\mathrm{H}_{2} \mathrm{O}$ molecules are often very labile, however, and NMR and isotope studies have shown that the exchange between the coordination sphere of the ion and the bulk solvent is very rapid for ions of low charge but slow for ions of high charge.

The proton, although it is very small, has a very high mobility (Table 16B.2). Proton and ${ }^{17} \mathrm{O}-\mathrm{NMR}$ show that the characteristic lifetime of protons hopping from one molecule to the next is about 1.5 ps , which is comparable to the time that inelastic neutron scattering shows it takes a water molecule to turn through about 1 radian ( 1 to 2 ps ). According to the Grotthuss mechanism, there is an effective motion of a proton that involves the rearrangement of bonds in a group of water molecules (Fig. 16B.2). However, the actual mechanism is still highly contentious. The mobility of protons in liquid ammonia is also anomalous and presumably occurs by an analogous mechanism.

\subsection*{(b) Mobility and conductivity}
The limiting molar conductivity of an ion is a measurable quantity, and the mobility of an ion can, in principle, be calculated on the basis of a model of its motion through the solvent. It should be possible to find a relation between these two quantities.

\subsection*{How is that done? 16B. 2 Establishing the relation between}
 ionic mobility and limiting molar conductivityTo keep things notationally simple, ignore the signs of quantities in what follows and focus on their magnitudes. Consider a solution of an electrolyte at a molar concentration $c$. Let each formula unit give rise to $v_{+}$cations of charge $z_{+} e$ and $v_{-}$anions of charge $z_{-} e$. The molar concentration of each type of ion is therefore $v c$ (with $v=v_{+}$or $v_{-}$), and the number density of each type is $\mathcal{N}=v c N_{\mathrm{A}}$.

Step 1 Calculate the number of ions passing through an imaginary window\\
By referring to Fig. 16B. 3 you can see that the number of ions of one kind, moving at speed $s$, that pass through an imaginary window of area $A$ during an interval $\Delta t$ is equal to the number within the distance $s \Delta t$, and therefore to the number in the volume $s \Delta t A$. It follows that the number of ions passing through the window in that period is $\mathcal{N} s \Delta t A=s \Delta t A \nu c N_{\mathrm{A}}$.\\
\includegraphics[max width=\textwidth, center]{2025_02_27_d4b95d9718f0b9e2e6b9g-735(1)}

Figure 16B. 3 In the calculation of the current, all the ions, moving at speed $s$, within a distance $s \Delta t$ (i.e. those in the volume $s A \Delta t$ ) will pass through the area $A$.

Step 2 Calculate the charge passing through the window, and hence the current\\
Each ion carries a charge $z e$, so the charge passing through the window is $z e s \Delta t A v c N_{\mathrm{A}}$ which, by introducing Faraday's constant $F=e N_{\mathrm{A}}$, can be written $z s \Delta t A v c F$. The electrical current, $I$, is the rate of passage of charge, which in this case is the charge divided by the time interval $\Delta t$ : so $I=z s \Delta t A v c F / \Delta t=z s A v c F$.

Step 3 Set up expressions for the conductance, the conductivity, and the molar conductivity\\
The conductance is given by $G=I / \Delta \phi$, where $\Delta \phi$ is the potential difference across the solution. It follows that

$$
G=\frac{I}{\Delta \phi}=\frac{z s A v c F}{\Delta \phi}
$$

The conductivity is

$$
\kappa=\frac{G l}{A}=\frac{z s A v c F l}{\Delta \phi A}=\frac{z s v c F l}{\Delta \phi}
$$

The limiting molar ionic conductivity is

$$
\lambda=\frac{\kappa}{v c}=\frac{z s v c F l}{v c \Delta \phi}=\frac{z s F l}{\Delta \phi}
$$

\subsection*{Step 4 Introduce the ionic mobility}
At this point you can identify $\Delta \phi / l$ as the electric field strength $\mathcal{E}$, and $s / E$ as the mobility, $u$ :

\[
\lambda=\frac{z s F l}{\Delta \phi}=\frac{z s F}{\mathcal{E}}=z u F \quad \begin{array}{ll}
\frac{\Delta \phi}{T}=\mathcal{E} & \text { lon molar conductivity }  \tag{16B.10}\\
\text { in terms of mobility }
\end{array}
\]

Equation 16B. 10 applies to the cations and to the anions. For an electrolyte where there are $v_{+}$and $v_{-}$cations and anions, respectively, from each formula unit, it follows from eqn 16B. 6 that $\Lambda_{\mathrm{m}}^{\circ}=v_{+} \lambda_{+}+v_{-} \lambda_{\_}$. Therefore\\
\includegraphics[max width=\textwidth, center]{2025_02_27_d4b95d9718f0b9e2e6b9g-736}

For a symmetrical $z: z$ electrolyte (e.g. $\mathrm{CuSO}_{4}$ with $z_{+}=z_{-}=2$ ), this equation simplifies to


\begin{equation*}
\Lambda_{\mathrm{m}}^{\circ}=z\left(u_{+}+u_{-}\right) F \tag{16B.11b}
\end{equation*}


\subsection*{Brief illustration 16B.3}
In Brief illustration 16B. 2 the mobility of a typical ion is estimated as $7.3 \times 10^{-8} \mathrm{~m}^{2} \mathrm{~V}^{-1} \mathrm{~s}^{-1}$. For $z=1$, this value can be used to estimate a typical limiting molar conductivity of the ion as

$$
\begin{aligned}
\lambda & =1 \times\left(7.3 \times 10^{-8} \mathrm{~m}^{2} \mathrm{~V}^{-1} \mathrm{~s}^{-1}\right) \times\left(9.649 \times 10^{4} \mathrm{C} \mathrm{~mol}^{-1}\right) \\
& =7.0 \times 10^{-3} \mathrm{~m}^{2} \mathrm{~V}^{-1} \mathrm{~s}^{-1} \mathrm{Cmol}^{-1}
\end{aligned}
$$

Because $1 \mathrm{~V}^{-1} \mathrm{Cs}^{-1}=1 \mathrm{~S}$, the value can be expressed as $7.0 \mathrm{mS} \mathrm{m}^{2} \mathrm{~mol}^{-1}$. The experimental value for $\mathrm{K}^{+}(\mathrm{aq})$ is $7.4 \mathrm{mS} \mathrm{m}^{2} \mathrm{~mol}^{-1}$.

\subsection*{(c) The Einstein relations}
The relation between drift speed and the electric field strength in eqn 16B.8a $(s=z e E / f)$ is a special case of a more general relation derived in Topic 16C (eqn 16C.5):

$$
s=\frac{D F}{R T} \quad \begin{aligned}
& \text { Drift speed in terms of } \\
& \text { diffusion coefficient }
\end{aligned}
$$

(16B.12)\\
where $F$ is the force (per mole of ions) driving the ions through the viscous medium and $D$ is the diffusion coefficient for the species (Table 16B.3). For an ion in solution the drift speed is $s=u \mathcal{E}$ (eqn 16B.8b), and the force on each ion in an electric field of strength $\mathcal{E}$ is $e z \mathcal{E}$. It follows that the force per mole of ions is $N_{\mathrm{A}} e z \mathcal{E}$ which, by using $N_{\mathrm{A}} e=F$, can be written $z F E$. Substitution of these expressions for $s$ and $F$ into eqn 16B. 12 gives, on cancelling the $\mathcal{E}$, the Einstein relation:

$$
u=\frac{z D F}{R T}
$$

Einstein relation\\
(16B.13)

The mobility is high when the diffusion coefficient (which is inversely proportional to the viscosity) is high, indicating that

Table 16B. 3 Diffusion coefficients at $298 \mathrm{~K}, \mathrm{D} /\left(10^{-9} \mathrm{~m}^{2} \mathrm{~s}^{-1}\right)^{\star}$

\begin{center}
\begin{tabular}{llllll}
\hline
Molecules in liquids & \multicolumn{5}{l}{Ions in water} \\
\hline
$\mathrm{I}_{2}$ in hexane & 4.05 & $\mathrm{~K}^{+}$ & 1.96 & $\mathrm{Br}^{-}$ & 2.08 \\
in benzene & 2.13 & $\mathrm{H}^{+}$ & 9.31 & $\mathrm{Cl}^{-}$ & 2.03 \\
Glycine in water & 1.055 & $\mathrm{Na}^{+}$ & 1.33 & $\mathrm{I}^{-}$ & 2.05 \\
$\mathrm{H}_{2} \mathrm{O}$ in water & 2.26 &  &  & $\mathrm{OH}^{-}$ & 5.03 \\
Sucrose in water & 0.5216 &  &  &  &  \\
\hline
\end{tabular}
\end{center}

\begin{itemize}
  \item More values are given in the Resource section.\\
the solute molecules are very mobile. Don't be misled by the presence of temperature in the denominator into thinking that the ion mobility decreases with increasing temperature: the diffusion coefficient increases more rapidly with temperature than $T$ itself, so $u$ increases with increasing temperature.
\end{itemize}

\subsection*{Brief illustration 16B.4}
From Table 16B.2, the mobility of $\mathrm{SO}_{4}^{2-}$ is $8.29 \times 10^{-8} \mathrm{~m}^{2} \mathrm{~V}^{-1} \mathrm{~s}^{-1}$. It follows from eqn 16B. 13 in the form $D=u R T / z F$ that the diffusion coefficient for the ion in water at $25^{\circ} \mathrm{C}$ is

$$
\begin{aligned}
D & =\frac{\left(8.29 \times 10^{-8} \mathrm{~m}^{2} \mathrm{~V}^{-1} \mathrm{~s}^{-1}\right) \times\left(8.3145 \mathrm{JK}^{-1} \mathrm{~mol}^{-1}\right) \times(298 \mathrm{~K})}{2 \times(9.649 \times 10^{4} \underbrace{}_{\mathrm{JV}} \mathrm{~mol}^{-1})} \\
& =1.06 \times 10^{-9} \mathrm{~m}^{2} \mathrm{~s}^{-1}
\end{aligned}
$$

The Einstein relation can be developed to provide a link between the limiting molar conductivity of an electrolyte and the diffusion coefficients of its ions. First, by using eqns 16B. 10 and 16B. 13 the limiting molar conductivity of an ion can be written


\begin{equation*}
\lambda=z u F=\frac{z^{u=z D F / R T}}{z^{2} D F^{2}} \tag{16B.14}
\end{equation*}


Then, by using $\Lambda_{\mathrm{m}}^{\circ}=v_{+} \lambda_{+}+v_{-} \lambda_{-}$(eqn 16B.6), the limiting molar conductivity is

$$
\Lambda_{\mathrm{m}}^{\circ}=\left(v_{+} z_{+}^{2} D_{+}+v_{-} z_{-}^{2} D_{-}\right) \frac{F^{2}}{R T} \quad \begin{aligned}
& \text { Nernst-Einstein } \\
& \text { equation }
\end{aligned}
$$

(16B.15)\\
which is the Nernst-Einstein equation. One application of this equation is to the determination of ionic diffusion coefficients from conductivity measurements; another is to the prediction of conductivities based on models of ionic diffusion.

\subsection*{Checklist of concepts}
$\square$ 1. The viscosity of a liquid decreases with increasing temperature.2. Kohlrausch's law states that, at low concentrations, the molar conductivities of strong electrolytes vary as the square root of the concentration.3. The law of the independent migration of ions states that the molar conductivity, in the limit of zero concentration, is the sum of contributions from its individual ions.4. An ion reaches a drift speed when the acceleration due to the electrical force is balanced by the viscous drag.5. The hydrodynamic radius of an ion may be greater than its ionic radius.6. The high mobility of a proton in water is explained by the Grotthuss mechanism.7. The mobility of an ion can be related to its limiting molar conductivity and, via the Einstein relation, to its diffusion coefficient.

\subsection*{Checklist of equations}
\begin{center}
\begin{tabular}{|c|c|c|c|}
\hline
Property & Equation & Comment & Equation number \\
\hline
Viscosity of a liquid & $\eta=\eta_{0} \mathrm{e}^{E_{\mathrm{a}} / R T}$ & Over a narrow temperature range & 16B.2 \\
\hline
Conductivity & $\kappa=G l / A, G=1 / R$ & Definition & 16B.3 \\
\hline
Molar conductivity & $\Lambda_{\mathrm{m}}=\kappa / \mathrm{c}$ & Definition & 16B. 4 \\
\hline
Kohlrausch's law & $\Lambda_{\mathrm{m}}=\Lambda_{\mathrm{m}}^{\circ}-\mathcal{K} c^{1 / 2}$ & Empirical observation & 16B.5 \\
\hline
Law of independent migration of ions & $\Lambda_{\mathrm{m}}^{\circ}=v_{+} \lambda_{+}+v_{-} \lambda_{-}$ & Limiting law & 16B. 6 \\
\hline
Stokes' law & $F_{\text {viscous }}=f s, f=6 \pi \eta a$ &  & 16B.7 \\
\hline
Drift speed & $s=u \mathcal{E}$ & Defines $u$ & 16B.8b \\
\hline
Ion mobility & $u=z e / 6 \pi \eta a$ & Assumes Stokes' law & 16B. 9 \\
\hline
Conductivity and mobility & $\lambda=z u F$ &  & 16B.10 \\
\hline
Molar conductivity and mobility & $\Lambda_{\mathrm{m}}^{\circ}=\left(z_{+} u_{+} v_{+}+z_{-} u_{-} v_{-}\right) F$ &  & 16B.11a \\
\hline
Drift speed & $s=D F / R T$ & $F$ is a general (molar) force & 16B.12 \\
\hline
Einstein relation & $u=z D F / R T$ &  & 16B. 13 \\
\hline
Nernst-Einstein equation & $\Lambda_{\mathrm{m}}^{\circ}=\left(v_{+} z_{+}^{2} D_{+}+v_{-} z_{-}^{2} D_{-}\right)\left(F^{2} / R T\right)$ &  & 16B.15 \\
\hline
\end{tabular}
\end{center}

\section*{TOPIC 16C Diffusion}
\subsection*{Why do you need to know this material?}
The diffusion of chemical species through space determines the rates of many chemical reactions in chemical reactors, living cells, and the atmosphere.

\subsection*{$>$ What is the key idea?}
Molecules and ions tend to spread into a uniform distribution.

\subsection*{- What do you need to know already?}
This Topic draws on arguments relating to the calculation of flux (Topic 16A) and the notion of drift speed introduced in Topic 16B. It also uses the concept of chemical potential to discuss the direction of spontaneous change (Topic 5A). The final section uses a statistical argument like that used to discuss the properties of a random coil in Topic 14D.

That solutes in gases, liquids, and solids have a tendency to spread can be discussed from three points of view. The thermodynamic viewpoint makes use of the Second Law of thermodynamics and the tendency for entropy to increase and, if the temperature and pressure are constant, for the Gibbs energy to decrease. The second approach is to set up a differential equation for the change in concentration in a region by considering the flux of material through its boundaries. The third approach is based on a model in which diffusion is imagined as taking place in a series of random small steps.

Several derivations in this Topic use Fick's first law of diffusion, which is discussed in Topic 16A and repeated here for convenience:

$$
J(\text { number })=-D \frac{\mathrm{~d} \mathcal{N}}{\mathrm{~d} x}
$$ [number]

(16C.1a)\\
where $\mathcal{N}$ is the number density and $D$ is the diffusion coefficient. In a number of cases it is more convenient to discuss the flux in terms of the amount of molecules and the molar concentration, c. Division by Avogadro's constant turns eqn 16C.1a into

$$
J(\text { amount })=-D \frac{\mathrm{~d} c}{\mathrm{~d} x}
$$

Fick's first law\\
(16C.1b)

\subsection*{16c. 1 The thermodynamic view}
At constant temperature and pressure, the maximum nonexpansion work that can be done by a spontaneous process is equal to the change in the Gibbs energy (Topic 3D). In this case the spontaneous process is the spreading of a solute, and the work it could achieve per mole of solute molecules can be identified with the change in the chemical potential of the solute: $\mathrm{d} w_{\mathrm{m}}=\mathrm{d} \mu$. The difference in chemical potential between the locations $x+\mathrm{d} x$ and $x$ is

$$
\mathrm{d} \mu=\mu(x+\mathrm{d} x)-\mu(x)=\left(\frac{\partial \mu}{\partial x}\right)_{T, p} \mathrm{~d} x
$$

so the molar work associated with migration through $\mathrm{d} x$ is

$$
\mathrm{d} w_{\mathrm{m}}=\left(\frac{\partial \mu}{\partial x}\right)_{T, p} \mathrm{~d} x
$$

The work done in moving a distance $\mathrm{d} x$ against an opposing force $F$ (in this context, a molar quantity) is $\mathrm{d} w_{\mathrm{m}}=-F \mathrm{~d} x$. By comparing the two expressions for $\mathrm{d} w_{\mathrm{m}}$ it is seen that the slope of the chemical potential with respect to position can be interpreted as an effective force per mole of molecules. This thermodynamic force is written as

$$
F=-\left(\frac{\partial \mu}{\partial x}\right)_{T, p} \quad \begin{aligned}
& \text { Thermodynamic force } \\
& \text { [definition] }
\end{aligned}
$$

(16C.2)

There is not a real force pushing the molecules down the slope of the chemical potential: the apparent force represents the spontaneous tendency of the molecules to disperse as a consequence of the Second Law and the tendency towards greater entropy.

In a solution in which the activity of the solute is $a$, the chemical potential is $\mu=\mu^{\ominus}+R T \ln a$. The thermodynamic force can therefore be written in terms of the gradient of the logarithm of the activity:

$$
F=-\left(\frac{\partial \mu}{\partial x}\right)_{T, p} \stackrel{(\partial \mu / \partial x)_{T, p}=\partial\left(\mu^{\ominus}+R T \ln a\right) / \partial x}{=}-R T\left(\frac{\partial \ln a}{\partial x}\right)_{T, p}
$$

If the solution is ideal, $a$ may be replaced by $c / c^{\ominus}$, where $c$ is the molar concentration and $c^{\ominus}$ is its standard value ( $1 \mathrm{~mol} \mathrm{dm}^{-3}$ ):


\begin{equation*}
F=-R T\left(\frac{\partial \ln \left(c / c^{\ominus}\right)}{\partial x}\right)_{T, p}=-\frac{R T}{c}\left(\frac{\partial c}{\partial x}\right)_{T, p} \tag{16C.3b}
\end{equation*}


\subsection*{Example 16C. 1 Calculating the thermodynamic force}
Suppose that the concentration of a solute varies linearly along $x$ according to $c=c_{0}+\alpha x$, where $c_{0}$ is the concentration at $x=0$. Find an expression for the thermodynamic force at position $x$, and evaluate this force at $x=0$ and $x=1.0 \mathrm{~cm}$ for the case where $c_{0}=1.0 \mathrm{moldm}^{-3}$ and $\alpha=10 \mathrm{~mol} \mathrm{dm}^{-3} \mathrm{~m}^{-1}$. Take $T=298 \mathrm{~K}$.

Collect your thoughts You will need to use eqn 16C.3b to find the force, so begin by evaluating $\partial c / \partial x$.

The solution The gradient of the molar concentration is

$$
\left(\frac{\partial c}{\partial x}\right)_{T, p}=\left(\frac{\partial\left(c_{0}+\alpha x\right)}{\partial x}\right)_{T, p}=\alpha
$$

Then, from eqn 16C.3b, the thermodynamic force is

$$
F=-\frac{R T}{c}\left(\frac{\partial c}{\partial x}\right)_{T, p} \stackrel{(\partial c / \partial x)_{T, p}=\alpha_{i} c=c_{0}+\alpha x}{=-\frac{\alpha R T}{c_{0}+\alpha x}}
$$

At $x=0$,

$$
\begin{aligned}
F & =-\frac{\left(10 \mathrm{moldm}^{-3} \mathrm{~m}^{-1}\right) \times\left(8.3145 \mathrm{JK}^{-1} \mathrm{~mol}^{-1}\right) \times(298 \mathrm{~K})}{1.0 \mathrm{moldm}^{-3}} \\
& =-2.5 \times 10^{4} \mathrm{Jm}^{-1} \mathrm{~mol}^{-1}
\end{aligned}
$$

or $-25 \mathrm{kN} \mathrm{mol}^{-1}$. A similar calculation at $x=1.0 \mathrm{~cm}=1.0 \times$ $10^{-2} \mathrm{~m}$ gives the force as $-23 \mathrm{kN} \mathrm{mol}^{-1}$.

Comment. The negative sign indicates that the force is towards the left (towards negative $x$ ), because the concentration increases towards the right (as $c_{0}+\alpha x$ ) and there is therefore a tendency for the solute to migrate to the left under the influence of that apparent force. The magnitude of the thermodynamic force decreases as $x$ increases because the gradient of $\ln \left(c / c_{0}\right)$, which is $\alpha /\left(c_{0}+\alpha x\right)$, becomes smaller on going to the right (Fig. 16C.1).\\
\includegraphics[max width=\textwidth, center]{2025_02_27_d4b95d9718f0b9e2e6b9g-739}

Figure 16C. 1 The thermodynamic force is proportional to $-\partial \ln c / \partial x=-(1 / c) \partial c / \partial x$. The force thus drives the molecules from a region with higher concentration to one with lower concentration, and becomes smaller in magnitude as the concentration increases.

Self-test 16C. 1 Suppose that the concentration of a solute decreases exponentially to the right as $c(x)=c_{0} \mathrm{e}^{-x / l}$. Derive an expression for the thermodynamic force at any position.\\
$1 / L 甘=屯: \iota \partial м s u_{V}$

The thermodynamic force acts in many respects like a real physical force. In particular it is responsible for accelerating solute molecules until the viscous drag they experience balances the apparent driving force and they settle down to a steady 'drift speed' through the medium. By considering the balance of apparent driving force and the retarding viscous force it is possible to derive Fick's first law of diffusion and relate the diffusion coefficient to the properties of the medium.

\subsection*{How is that done? 16C. 1 Deriving Fick's first law of diffusion}
The flux due to a concentration gradient is the amount (in moles) of molecules passing through an area $A$ in an interval $\Delta t$ divided by the area and the interval. This flux can be related to the drift speed, $s$, by using an approach like that used in Topic 16A for diffusion in a gas.

Step 1 Find an expression for the flux due to molecules moving at the drift speed\\
In a time interval $\Delta t$ all the particles within a distance $s \Delta t$ can pass through the window, which means that all of the particles in a volume $s \Delta t A$ pass through the window (there is no reverse flux because, unlike in a gas, in this model all the solute molecules are moving down the concentration gradient). Hence, the amount (in moles) of solute molecules that can pass through the window is $s \Delta t A c$. The flux $J$ is this number divided by the area $A$ and by the time interval $\Delta t$ :

$$
J(\text { amount })=s c
$$

\subsection*{Step 2 Find an expression for the drift speed}
The apparent driving force (per mole) acting on the solute molecules is $F=-(R T / c) \mathrm{d} c / \mathrm{d} x$, eqn 16 C .3 b . The molecules also experience a viscous drag which is assumed to be proportional to the speed: expressed as a molar quantity this force is written $N_{A} f s$, where $f$ is a constant, the 'frictional constant', depending on the medium. When these two forces are in balance the molecules will be moving at the drift speed. It therefore follows that, with $R / N_{\mathrm{A}}=k$,

$$
N_{\mathrm{A}} f s=-\frac{R T}{c} \frac{\mathrm{~d} c}{\mathrm{~d} x} \text { hence } s=-\frac{R T}{N_{\mathrm{A}} f c} \frac{\mathrm{~d} c}{\mathrm{~d} x}=-\frac{k T}{f c} \frac{\mathrm{~d} c}{\mathrm{~d} x}
$$

The negative sign arises because the molecules are moving opposite to the direction of increasing concentration.

\subsection*{Step 3 Combine the two expressions}
Now substitute the expression for the drift speed into that for the flux to give

$$
J(\text { amount })=s c=-\frac{k T}{f c} \frac{\mathrm{~d} c}{\mathrm{~d} x} c=-\frac{k T}{f} \frac{\mathrm{~d} c}{\mathrm{~d} x}
$$

This expression has the same form as Fick's first law, eqn 16 C .1 b , that the flux is proportional to the concentration gradient. In addition, the diffusion constant $D$ can be identified as $k T / f$, which is the Stokes-Einstein relation:\\
\includegraphics[max width=\textwidth, center]{2025_02_27_d4b95d9718f0b9e2e6b9g-740}

The constant $f$ that appears in the Stokes-Einstein relation can be inferred from the hydrodynamic result known as Stokes' law for the viscous drag (this law is used in Topic 16B). According to this law the magnitude of the viscous force is $6 \pi \eta$ as for a spherical particle of radius $a$. It therefore follows that $f=6 \pi \eta a$, and substituting this into eqn 16C.4a gives the Stokes-Einstein equation

$$
D=\frac{k T}{6 \pi \eta a}
$$

Stokes-Einstein equation\\
(16C.4b)

This equation is an explicit relation between the diffusion coefficient and the viscosity for a species of hydrodynamic radius $a$ and confirms that $D$ is inversely proportional to $\eta$.

The drift speed can be related to the diffusion constant and the thermodynamic force by equating two expressions for the flux, $J=s c$ and $J=-D \mathrm{~d} c / \mathrm{d} x$, to obtain $s=-(D / c) \mathrm{d} c / \mathrm{d} x$. The concentration gradient can be expressed in terms of the thermodynamic force, that is, eqn 16C.3b rearranged in the form $\mathrm{d} c / \mathrm{d} x=-c F / R T$ to give


\begin{equation*}
s=\frac{D F}{R T} \tag{16C.5}
\end{equation*}


This relation provides a way to estimate the thermodynamic force from measurements of the drift speed and the diffusion coefficient.

\subsection*{Brief illustration 16C. 1}
Laser measurements show that a particular molecule has a drift speed of $1.0 \mu \mathrm{~m} \mathrm{~s}^{-1}$ in water at $25^{\circ} \mathrm{C}$, at which temperature the diffusion coefficient is $5.0 \times 10^{-9} \mathrm{~m}^{2} \mathrm{~s}^{-1}$. The corresponding thermodynamic force calculated using eqn 16C.5, rearranged into the form $F=s R T / D$, is

$$
\begin{aligned}
F & =\frac{\left(1.0 \times 10^{-6} \mathrm{~m} \mathrm{~s}^{-1}\right) \times\left(8.3145 \mathrm{JK}^{-1} \mathrm{~mol}^{-1}\right) \times(298 \mathrm{~K})}{5.0 \times 10^{-9} \mathrm{~m}^{2} \mathrm{~s}^{-1}} \\
& =5.0 \times 10^{5} \overbrace{\mathrm{Jm}^{-1}}^{\mathrm{N}} \mathrm{~mol}^{-1}
\end{aligned}
$$

or about $500 \mathrm{kN} \mathrm{mol}^{-1}$. This thermodynamic force is many times that of gravity, which explains why solutions do not sediment.

\subsection*{16C. 2 The diffusion equation}
Diffusion results in the modification of the distribution of concentration of the solute (or of a physical property) as inhomogeneities disappear. The discussion is expressed in terms of the diffusion of molecules, but similar arguments apply to the diffusion of other entities, such as ions, and of various physical properties, such as temperature.

\subsection*{(a) Simple diffusion}
The diffusion equation, one of the most important equations for discussing the properties of fluids, is an equation that expresses the rate of change of concentration of a species in terms of the inhomogeneity of its concentration. It is also called 'Fick's second law of diffusion', but that name is now rarely used. The diffusion equation can be derived on the basis of Fick's first law.

\subsection*{How is that done? 16C. 2 Deriving the diffusion equation}
The diffusion equation is developed by considering the net flux of particles entering a thin slab of cross-sectional area $A$ that extends from $x$ to $x+l$ (Fig. 16C.2) and therefore has volume Al.\\
\includegraphics[max width=\textwidth, center]{2025_02_27_d4b95d9718f0b9e2e6b9g-740(1)}

Figure 16C. 2 The net flux in a thin slab is the difference between the flux entering from the region of high concentration (on the left) and the flux leaving to the region of low concentration (on the right).

Step 1 Find an expression for the net rate of change of the concentration in the slab due to particles entering from each side

If the flux of molecules from the left is $J_{\mathrm{L}}$, then the rate at which the particles enter the slab is $J_{\mathrm{L}} A$. The rate of increase of the molar concentration inside the slab due to the flux in from the left is

$$
\left(\frac{\partial c}{\partial t}\right)_{\mathrm{L}}=\frac{J_{\mathrm{L}} A}{A l}=\frac{J_{\mathrm{L}}}{l}
$$

Molecules also flow out of the right face of the slab. If this flux is $J_{\mathrm{R}}$, then by a similar argument

$$
\left(\frac{\partial c}{\partial t}\right)_{\mathrm{R}}=-\frac{J_{\mathrm{R}}}{l}
$$

Note the minus sign: this flux reduces the concentration. The net rate of change of concentration in the slab is

$$
\frac{\partial c}{\partial t}=\left(\frac{\partial c}{\partial t}\right)_{\mathrm{L}}+\left(\frac{\partial c}{\partial t}\right)_{\mathrm{R}}=\frac{J_{\mathrm{L}}-J_{\mathrm{R}}}{l}
$$

Step 2 Relate the fluxes to the concentration gradients\\
From Fick's first law (eqn 16C.1b), each flux can be expressed in terms of the diffusion coefficient and the concentration gradient at each face:

$$
J_{\mathrm{L}}-J_{\mathrm{R}}=-D\left(\frac{\partial c}{\partial x}\right)_{\mathrm{L}}+D\left(\frac{\partial c}{\partial x}\right)_{\mathrm{R}}=D\left\{\left(\frac{\partial c}{\partial x}\right)_{\mathrm{R}}-\left(\frac{\partial c}{\partial x}\right)_{\mathrm{L}}\right\}
$$

where $(\partial c / \partial x)_{\mathrm{L}}$ is the concentration gradient at the left face of the slab, and similarly $(\partial c / \partial x)_{\mathrm{R}}$ is that at the right face. The concentration gradients at the two faces can be expressed in terms of the gradient (the first derivative of the concentration) at the centre of the slab, $(\partial c / \partial x)_{0}$, and the first derivative of that gradient (which is the second derivative of the concentration), $\partial^{2} c / \partial x^{2}$. The distances of the faces from the centre are $\frac{1}{2} l$ in each direction, so it follows that

$$
\left(\frac{\partial c}{\partial x}\right)_{\mathrm{R}}-\left(\frac{\partial c}{\partial x}\right)_{\mathrm{L}}=\left\{\left(\frac{\partial c}{\partial x}\right)_{0}+\frac{l}{2} \frac{\partial^{2} c}{\partial x^{2}}\right\}-\left\{\left(\frac{\partial c}{\partial x}\right)_{0}-\frac{l}{2} \frac{\partial^{2} c}{\partial x^{2}}\right\}=l \frac{\partial^{2} c}{\partial x^{2}}
$$

This expression is then substituted into that for the difference of the fluxes to give

$$
J_{\mathrm{L}}-J_{\mathrm{R}} \stackrel{\nabla}{=}=\frac{\partial^{2} c}{\partial c / \partial x)_{\mathrm{R}}-(\partial c \partial x)_{\mathrm{L}}=l\left(\partial^{2} c \partial x^{2}\right)}{ }^{2 x^{2}}
$$

Step 3 Combine the expression for the net flux with that for the time dependence of concentration\\
Substitute the last expression into the equation for the rate of change of concentration, cancel the $l$, and obtain the diffusion equation:

$$
-\left|\frac{\partial c}{\partial t}=D \frac{\partial^{2} c}{\partial x^{2}}\right| \quad \text { (16C.6) } \quad \text { Diffusion equation }
$$

The diffusion equation shows that the rate of change of concentration in a region is proportional to the curvature (more precisely, to the second derivative) of the concentration with respect to distance in that region. If the concentration changes sharply from point to point (if the distribution is highly wrinkled), then the concentration changes rapidly with time. Specifically:\\
\includegraphics[max width=\textwidth, center]{2025_02_27_d4b95d9718f0b9e2e6b9g-741}

Figure 16C. 3 The diffusion equation implies that, over time, peaks in a distribution (regions of negative curvature) spread and troughs (regions of positive curvature) fill in.

\begin{itemize}
  \item Where the curvature is positive (a dip, Fig. 16C.3), the change in concentration with time is positive; the dip tends to fill.
  \item Where the curvature is negative (a heap), the change in concentration with time is negative; the heap tends to spread.
  \item If the curvature is zero, then the concentration is constant in time.
\end{itemize}

The diffusion equation can be regarded as a mathematical formulation of the intuitive notion that there is a natural tendency for the wrinkles in a distribution to disappear.

\subsection*{Brief illustration 16C. 2}
If a concentration across a small region of space varies linearly as $c=c_{0}-\alpha x$ then it follows that $\partial^{2} c / \partial x^{2}=0$ and so from eqn $16 \mathrm{C} .6 \partial c / \partial t=0$. The concentration does not vary with time because the flow into one face of a slab is exactly matched by the flow out from the opposite face (Fig. 16C.4a). If the concentration varies as $c=c_{0}-\frac{1}{2} \beta x^{2}$ then $\partial^{2} c / \partial x^{2}=-\beta$ and consequently $\partial c / \partial t=-D \beta$. The concentration decreases with time because the flow out of the slab is greater than the flow into it (Fig. 16C.4b).\\
\includegraphics[max width=\textwidth, center]{2025_02_27_d4b95d9718f0b9e2e6b9g-741(1)}

Figure 16C. 4 The two instances treated in Brief illustration 16C.2: (a) linear concentration gradient, (b) parabolic concentration gradient.

\subsection*{(b) Diffusion with convection}
Convection is the bulk motion of regions of a fluid. This process contrasts with diffusion in which molecules move individually through the fluid. The flux due to convection can be analysed in a similar way to diffusion.

\subsection*{How is that done? 16C. 3}
Evaluating the change in\\
concentration due to convection\\
As in previous calculations, imagine the flux of molecules through an area $A$ in an interval $\Delta t$ but now due to convective flow in which the fluid moves at a speed $v$.

Step 1 Express the rate of change of concentration in terms of the net flux\\
As in the derivation of the diffusion equation,

$$
\frac{\partial c}{\partial t}=\frac{J_{\mathrm{L}, \text { conv }}-J_{\mathrm{R}, \mathrm{conv}}}{l}
$$

where $J_{\mathrm{L}, \text { conv }}$ and $J_{\mathrm{R}, \text { conv }}$ are, respectively, the fluxes from the left into and on the right out of the slab, but here due to convection.

Step 2 Evaluate the net flux\\
In an interval $\Delta t$ all the particles within a distance $v \Delta t$ and therefore in the volume $A v \Delta t$ pass through a face of the slab. If the molar concentration at the relevant face is $c$, then the amount passing through that face is $c A v \Delta t$. The 'convective flux' is this amount divided by the area of the face and the time interval:

$$
J_{\text {conv }}=\frac{c A v \Delta t}{A \Delta t}=c v
$$

Convective flux (16C.7)\\
The concentrations on the left $\left(c_{\mathrm{L}}\right)$ and right $\left(c_{\mathrm{R}}\right)$ faces of the slab are related to the concentration at its centre, $c_{0}$, by

$$
c_{\mathrm{R}}=c_{0}+\frac{1}{2} l\left(\frac{\partial c}{\partial x}\right) \quad c_{\mathrm{L}}=c_{0}-\frac{1}{2} l\left(\frac{\partial c}{\partial x}\right)
$$

where the first derivatives are evaluated at the centre of the slab. It follows from eqn 16C. 7 that

$$
\begin{aligned}
J_{\mathrm{L}, \text { conv }}-J_{\mathrm{R}, \text { conv }}=\left(c_{\mathrm{L}}-c_{\mathrm{R}}\right) v & =\left\{c_{0}-\frac{1}{2} l\left(\frac{\partial c}{\partial x}\right)\right\} v-\left\{c_{0}+\frac{1}{2} l\left(\frac{\partial c}{\partial x}\right)\right\} v \\
& =-\left(\frac{\partial c}{\partial x}\right) l v
\end{aligned}
$$

\subsection*{Step 3 Evaluate the net rate of change of concentration}
Now substitute this result into the expression for the rate of change of concentration derived in Step 1, and obtain

$$
-\left|\frac{\partial c}{\partial t}=-\left(\frac{\partial c}{\partial x}\right) v\right| \begin{array}{r}
\text { (16C.8) } \\
\end{array}
$$

\subsection*{Brief illustration 16C. 3}
If it is assumed that the concentration across a small region of space varies linearly as $c=c_{0}-\alpha x$, then $\partial c / \partial x=-\alpha$. If there is convective flow at velocity $v$, it follows from eqn 16C. 8 that $\partial c / \partial t=\alpha v$. The concentration in the slab increases because the convective flow into the left face outweighs the flow out from the right face; with this linear concentration dependence, there is no diffusion. If $\alpha=0.010 \mathrm{moldm}^{-3} \mathrm{~m}^{-1}$ and $v=$ $+1.0 \mathrm{~mm} \mathrm{~s}^{-1}$,

$$
\begin{aligned}
\frac{\partial c}{\partial t} & =\left(0.010 \mathrm{moldm}^{-3} \mathrm{~m}^{-1}\right) \times\left(1.0 \times 10^{-3} \mathrm{~m} \mathrm{~s}^{-1}\right) \\
& =1.0 \times 10^{-5} \mathrm{moldm}^{-3} \mathrm{~s}^{-1}
\end{aligned}
$$

The concentration increases at the rate of $10 \mu \mathrm{~mol} \mathrm{dm}^{-3} \mathrm{~s}^{-1}$.

When diffusion and convection occur together, the total rate of change of concentration in a region is the sum of the two effects, which is described by the generalized diffusion equation:

$$
\frac{\partial c}{\partial t}=D \frac{\partial^{2} c}{\partial x^{2}}-v \frac{\partial c}{\partial x} \quad \text { Generalized diffusion equation }
$$

(16C.9)\\
A further refinement, which is important in chemistry, is the possibility that the concentrations of molecules may change as a result of reaction. When reactions are included in eqn 16C. 9 (in Topic 18B) a differential equation is obtained that can be used to discuss the properties of reacting, diffusing, convecting systems. This equation is the basis for modelling reactors in the chemical industry and the utilization of resources in living cells.

\subsection*{(c) Solutions of the diffusion equation}
The diffusion equation (eqn 16C.6) is a second-order differential equation with respect to space and a first-order differential equation with respect to time. To find solutions it is necessary to know two boundary conditions for the spatial dependence and a single initial condition for the time dependence.

As an illustration, consider an arrangement in which there is a layer of a solute (such as sugar) at the bottom of a tall beaker of water (which for simplicity may be taken to be infinitely tall), with base area $A ; x$ is the distance measured up from the base. At $t=0$ it is assumed that all $N_{0}$ particles are concentrated on the $y z$-plane at $x=0$ : this is the initial condition. The two boundary conditions are derived from the requirements that the concentration must everywhere be finite and the total amount of particles present is $n_{0}$ (with $n_{0}=N_{0} / N_{A}$ ) at all times. With these conditions,

$$
c(x, t)=\frac{n_{0}}{A(\pi D t)^{1 / 2}} \mathrm{e}^{-x^{2} / 4 D t} \quad \begin{aligned}
& \text { One-dimensional } \\
& \text { diffusion }
\end{aligned}
$$

(16C.10)\\
\includegraphics[max width=\textwidth, center]{2025_02_27_d4b95d9718f0b9e2e6b9g-743}

Figure 16C. 5 The concentration profiles above a plane from which a solute is diffusing into pure solvent. The curves are labelled with the corresponding value of $D t$, and $x_{0}=\{4 D t\}^{1 / 2}$.\\
as may be verified by direct substitution. Figure 16C. 5 shows the shape of the concentration distribution at various times and illustrates how the concentration spreads.

Another useful result is for the diffusion in three dimensions arising from an initially localized concentration of solute (a sugar lump suspended in an infinitely large flask of water). The concentration of diffused solute is spherically symmetrical, and at a radius $r$ is

$$
c(r, t)=\frac{n_{0}}{8(\pi D t)^{3 / 2}} \mathrm{e}^{-r^{2} / 4 D t}
$$

Three-dimensional diffusion\\
(16C.11)

Other chemically (and physically) interesting arrangements, such as transport of substances across biological membranes can be treated, but the mathematical forms of the solutions are more cumbersome.

The diffusion equation is useful for the experimental determination of diffusion coefficients. In the capillary technique, a capillary tube, open at one end and containing a solution, is immersed in a well-stirred larger quantity of solvent, and the change of concentration in the tube is monitored. The solute diffuses from the open end of the capillary at a rate that can be calculated by solving the diffusion equation with the appropriate boundary and initial conditions, so $D$ may be determined. In the diaphragm technique, the diffusion occurs through the capillary pores of a sintered glass diaphragm separating the well-stirred solution and solvent. The concentrations are monitored and then related to the diffusion equation that has been solved for this arrangement. Diffusion coefficients may also be measured by a number of other techniques, including NMR spectroscopy (Topic 12C).

The diffusion equation can be used to predict the concentration of particles (or the value of some other physical quantity, such as the temperature in a non-uniform system) at any location. It can also be used to calculate the average displacement of the particles in a given time.

How is that done? 16C.4 Evaluating the average displacement in a one-dimensional system

To calculate the average value of the displacement $x$, denoted $\langle x\rangle$, you need to use eqn 16C. 10 to find an expression for the probability density of the particles $P(x)$, defined such that $P(x) \mathrm{d} x$ is the probability of finding a particle between $x$ and $x+\mathrm{d} x$. Then

$$
\langle x\rangle=\int_{0}^{\infty} x P(x) \mathrm{d} x
$$

Step 1 Set up an expression for the probability density\\
The number of particles in a slab of thickness $\mathrm{d} x$ at $x$ and at time $t$ is the volume of the slab, $A \mathrm{~d} x$, multiplied by the molar concentration at that location (and time) and Avogadro's constant: $N(x, t)=c(x, t) N_{\mathrm{A}} A \mathrm{~d} x$. The molar concentration at the position and time of interest is given by eqn 16C.10. The total number of particles is $N_{\mathrm{A}} n_{0}$, where $n_{0}$ is the total amount, so the probability of finding a molecule in the slab is

$$
\begin{aligned}
P(x) \mathrm{d} x & =\frac{N_{\mathrm{A}} c(x, t) A \mathrm{~d} x}{N_{\mathrm{A}} n_{0}}=\frac{A}{n_{0}} \frac{n_{0}}{A(\pi D t)^{1 / 2}} \mathrm{e}^{-x^{2} / 4 D t} \mathrm{~d} x \\
& =\frac{1}{(\pi D t)^{1 / 2}} \mathrm{e}^{-x^{2} / 4 D t} \mathrm{~d} x
\end{aligned}
$$

Step 2 Evaluate the integral\\
The integration required is

$$
\int_{0}^{\infty} x \frac{1}{(\pi D t)^{1 / 2}} \mathrm{e}^{-x^{2} / 4 D t} \mathrm{~d} x=\frac{1}{(\pi D t)^{1 / 2}} \overbrace{\int_{0}^{\infty} x \mathrm{e}^{-x^{2} / 4 D t} \mathrm{~d} x}^{\text {Integral G.2 }}=2\left(\frac{D t}{\pi}\right)^{1 / 2}
$$

That is, the average displacement of a diffusing particle in a time $t$ in a one-dimensional system is\\
$-\langle x\rangle=2\left(\frac{D t}{\pi}\right)^{1 / 2} \left\lvert\, \begin{aligned} & \text { Mean displacement } \\ & \text { [one dimension] }\end{aligned}\right.$\\
A similar calculation shows that the root-mean-square displacement in the same time is

\[
\left\langle x^{2}\right\rangle^{1 / 2}=(2 D t)^{1 / 2} \quad \begin{align*}
& \text { Root-mean-square displacement }  \tag{16C.13a}\\
& \text { [one dimension] }
\end{align*}
\]

This result is a useful measure of the spread of particles when they can diffuse in both directions from the origin, because in that case $\langle x\rangle=0$ at all times. The time-dependence of the root-mean-square displacement for particles with a typical diffusion coefficient in a liquid ( $D=5 \times 10^{-10} \mathrm{~m}^{2} \mathrm{~s}^{-1}$ ) is illustrated in Fig. 16C.6. The graph shows that diffusion is a very slow process (which is why solutions are stirred, to encourage mixing by convection).\\
\includegraphics[max width=\textwidth, center]{2025_02_27_d4b95d9718f0b9e2e6b9g-744}

Figure 16C. 6 The root-mean-square distance covered by particles with $D=5 \times 10^{-10} \mathrm{~m}^{2} \mathrm{~s}^{-1}$. Note the great slowness of diffusion.

In three dimensions the root-mean-square displacement is given by a similar expression:

$$
\left\langle r^{2}\right\rangle^{1 / 2}=(6 D t)^{1 / 2} \quad \begin{aligned}
& \text { Root-mean-square displacement } \\
& {[\text { three dimensions] }}
\end{aligned}
$$

(16C.13b)

\subsection*{16C. 3 The statistical view}
An intuitive picture of diffusion is of the particles moving in a series of small steps and gradually migrating from their original positions. This picture suggests a model in which each particle jumps through a distance $d$ after a time $\tau$. The total distance travelled by a particle in time $t$ is therefore $t d / \tau$. However, it is most unlikely that a particle will end up at this distance from the origin because each jump might be in a different direction.

The discussion is simplified by allowing the particles to travel only along a straight line (the $x$-axis) with each step a jump through distance $d$ to the left or to the right. This model is called the one-dimensional random walk. The probability that the walk will end up at a specific distance from the origin can be calculated by considering the statistics of the process. ${ }^{1}$

\subsection*{How is that done? 16C. 5 Evaluating the probability}
 distribution for a one-dimensional random walkImagine that a molecule has made $N$ steps in total, $N_{\mathrm{R}}$ of which are to the right and $N_{\mathrm{L}}$ to the left. The displacement from the origin is therefore $\left(N_{\mathrm{R}}-N_{\mathrm{L}}\right) d$, which is written $n d$ with $n=N_{\mathrm{R}}-N_{\mathrm{L}}$.

\footnotetext{${ }^{1}$ The calculation is essentially the same as in the discussion of the random coil structures of denatured polymers (Topic 14D).
}Step 1 Set up an expression for the probability of achieving a given final displacement\\
Many sequences of individual steps can arrive at a given final displacement. The number of these sequences is equal to the number of ways of choosing $N_{\mathrm{R}}$ steps to the right and $N_{\mathrm{L}}=$ $N-N_{\mathrm{R}}$ steps to the left:

$$
W=\frac{N!}{N_{\mathrm{L}}!N_{\mathrm{R}}!}=\frac{N!}{\left(N-N_{\mathrm{R}}\right)!N_{\mathrm{R}}!}
$$

At each step, the molecule can step to the left or right, so the total number of possible sequences of steps is $2^{N}$. The probability of achieving a final displacement $n d, P(n d)$, is therefore

$$
P(n d)=\frac{W}{2^{N}}=\frac{N!}{\left(N-N_{\mathrm{R}}\right)!N_{\mathrm{R}}!2^{N}}
$$

Step 2 Simplify this expression by using Stirling's approximation\\
This expression can be simplified by taking its logarithm to give

$$
\ln P=\ln N!-\left\{\ln \left(N-N_{\mathrm{R}}\right)!+\ln N_{\mathrm{R}}!+\ln 2^{N}\right\}
$$

and then using Stirling's approximation in the form

$$
\ln x!\approx \ln (2 \pi)^{1 / 2}+\left(x+\frac{1}{2}\right) \ln x-x
$$

to obtain

$$
\begin{aligned}
\ln P= & -\ln (2 \pi)^{1 / 2} 2^{N}+\left(N+\frac{1}{2}\right) \ln \frac{1}{1-N_{\mathrm{R}} / N} \\
& +N_{\mathrm{R}} \ln \frac{1-N_{\mathrm{R}} / N}{N_{\mathrm{R}} / N}-\frac{1}{2} \ln N_{\mathrm{R}}
\end{aligned}
$$

For reasons that will become clear shortly, it is convenient to introduce the new variable $\mu$ :

$$
\mu=\frac{N_{\mathrm{R}}}{N}-\frac{1}{2}
$$

from which definition it follows that $1-N_{\mathrm{R}} / N=\frac{1}{2}-\mu$ and $N_{\mathrm{R}} / N=\frac{1}{2}+\mu$. With these substitutions the expression for $\ln P$ can be written in terms of $N$ and $\mu$ alone:

$$
\begin{aligned}
\ln P= & -\ln (2 \pi)^{1 / 2} 2^{N}-\left(N+\frac{1}{2}\right) \ln \left(\frac{1}{2}-\mu\right)+N\left(\mu+\frac{1}{2}\right) \ln \left(\frac{1}{2}-\mu\right) \\
& -N\left(\mu+\frac{1}{2}\right) \ln \left(\mu+\frac{1}{2}\right)-\frac{1}{2} \ln N\left(\mu+\frac{1}{2}\right)
\end{aligned}
$$

\subsection*{Step 3 Expand the logarithms}
Because the probability of taking a step to the right or to the left is the same, it is expected that after many steps the total number to the right will be very close to half the number of steps. That is $N_{\mathrm{R}} / N \approx \frac{1}{2}$. It follows that $\mu \ll 1$, so you can use the series expansion

$$
\ln \left(\frac{1}{2} \pm \mu\right)=-\ln 2 \pm 2 \mu-2 \mu^{2}+\cdots
$$

and retain terms up to that in $\mu^{2}$. After rather of lot of algebra (see $A$ deeper look 11 on the website, where the details are given) you will obtain

$$
\ln P=-\ln (2 \pi N)^{1 / 2} 2^{N}+\ln 2^{N+1}-2(N-1) \mu^{2}
$$

At this point take antilogarithms of both sides and use $N \gg 1$ :

$$
P=\frac{2^{N+1} \mathrm{e}^{-2(N-1) \mu^{2}}}{2^{N}(2 \pi N)^{1 / 2}}=\frac{2 \mathrm{e}^{-2(N-1) \mu^{2}}}{(2 \pi N)^{1 / 2}} \approx \frac{2 \mathrm{e}^{-2 N \mu^{2}}}{(2 \pi N)^{1 / 2}}
$$

Step 4 Recast the expression for the probability in terms of the time for each step\\
The exponent $N \mu^{2}$ can be rewritten

$$
N \mu^{2}=\frac{\left(2 N_{\mathrm{R}}-N\right)^{2}}{4 N}=\frac{\left(N_{\mathrm{R}}-N_{\mathrm{L}}\right)^{2}}{4 N}=\frac{n^{2}}{4 N}
$$

The final distance from the origin, $x$, is equal to $n d$, and the number of steps taken in a time $t$ is $N=t / \tau$. It follows that $N \mu^{2}=n^{2} / 4 N=\tau x^{2} / 4 t d^{2}$. Substitution of these expressions for $N$ and $N \mu^{2}$ into the expression for $P$ gives

$$
-P(x, t)=\left(\frac{2 \tau}{\pi t}\right)^{1 / 2} \mathrm{e}^{-x^{2} \tau / 2 t d^{2}} \quad \begin{aligned}
& \text { One-dimensional random walk }
\end{aligned}
$$

The differences of detail between eqns 16C. 10 (for one-dimensional diffusion) and 16C. 14 arises from the fact that in the present calculation the particles can migrate in either direction from the origin. Moreover, they can be found only at discrete points separated by $d$ instead of being anywhere on a continuous line. The fact that the two expressions are so similar suggests that diffusion can indeed be interpreted as the outcome of a large number of steps in random directions.

By comparing the two exponents in eqn 16C. 10 and eqn 16 C .14 it is possible to relate the diffusion coefficient $D$ to the step length $d$ and the time between jumps, $\tau$. The result is the Einstein-Smoluchowski equation:

$$
D=\frac{d^{2}}{2 \tau}
$$

Einstein-Smoluchowski equation\\
(16C.15)

\subsection*{Brief illustration 16C. 4}
Suppose that in aqueous solution an $\mathrm{SO}_{4}^{2-}$ ion jumps through its own diameter of 500 pm each time it makes a move, then because $D=1.1 \times 10^{-9} \mathrm{~m}^{2} \mathrm{~s}^{-1}$ (as deduced from mobility measurements, Topic 16B), it follows from eqn 16C. 15 that

$$
\tau=\frac{d^{2}}{2 D}=\frac{\left(500 \times 10^{-12} \mathrm{~m}\right)^{2}}{2 \times\left(1.1 \times 10^{-9} \mathrm{~m}^{2} \mathrm{~s}^{-1}\right)}=1.1 \times 10^{-10} \mathrm{~s}
$$

or $\tau=110 \mathrm{ps}$. Because $\tau$ is the time for one jump, the ion makes about $1 \times 10^{10}$ jumps per second.

The Einstein-Smoluchowski equation makes the connection between the microscopic details of particle motion and the macroscopic parameters relating to diffusion. It also brings the discussion back full circle to the properties of the perfect gas treated in Topic 16A. If $d / \tau$ as is interpreted as $v_{\text {mean }}$, the mean speed of the molecules, and $d$ is interpreted as a mean free path $\lambda$, then the Einstein-Smoluchowski equation becomes $D=$ $\frac{1}{2} d(d / \tau)=\frac{1}{2} \lambda v_{\text {mean }}$, which is essentially the same expression as obtained from the kinetic model of gases (eqn 16A. 9 of Topic $16 \mathrm{~A}, D=\frac{1}{3} \lambda v_{\text {mean }}$ ). That is, the diffusion of a perfect gas is a random walk with an average step size equal to the mean free path.

\subsection*{Checklist of concepts}
1. A thermodynamic force is an apparent force that mirrors the spontaneous tendency of the molecules to disperse as a consequence of the Second Law and the tendency towards greater entropy.2. The drift speed is achieved when the viscous retarding force matches the thermodynamic force.3. The diffusion equation (Fick's second law) can be regarded as a mathematical formulation of the notionthat there is a natural tendency for concentration to become uniform.4. Convection is the bulk motion of regions of a fluid.\\
$\square$ 5. A model of diffusion is of the particles moving in a series of small steps, a random walk, and gradually migrating from their original positions.

\subsection*{Checklist of equations}
\begin{center}
\begin{tabular}{|c|c|c|c|}
\hline
Property & Equation & Comment & Equation number \\
\hline
Fick's first law & $J($ amount $)=-D \mathrm{~d} c / \mathrm{d} x$ &  & 16C.1b \\
\hline
Thermodynamic force & $F=-(\partial \mu / \partial x)_{T, p}$ & Definition & 16C. 2 \\
\hline
Stokes-Einstein relation & $D=k T / f$ & $f s$ is the frictional drag & 16C.4a \\
\hline
Drift speed & $s=D F / R T$ &  & 16C. 5 \\
\hline
Diffusion equation & $\partial c / \partial t=D \partial^{2} c / \partial x^{2}$ & One dimension & 16C. 6 \\
\hline
Generalized diffusion equation & $\partial c / \partial t=D \partial^{2} c / \partial x^{2}-v \partial c / \partial x$ & One dimension & 16C. 9 \\
\hline
Mean displacement & $x=2(D t / \pi)^{1 / 2}$ & One-dimensional diffusion & 16C. 12 \\
\hline
\multirow[t]{2}{*}{Root-mean-square displacement} & $\left\langle x^{2}\right\rangle^{1 / 2}=(2 D t)^{1 / 2}$ & One-dimensional diffusion & 16C.13a \\
\hline
 & $\left\langle r^{2}\right\rangle^{1 / 2}=(6 D t)^{1 / 2}$ & Three-dimensional diffusion & 16C.13b \\
\hline
Probability of displacement & $P(x, t)=(2 \tau / \pi t)^{1 / 2} \mathrm{e}^{-x^{2} \tau / 2 t d^{2}}$ & One-dimensional random walk & 16C. 14 \\
\hline
Einstein-Smoluchowski equation & $D=d^{2} / 2 \tau$ & One-dimensional random walk & 16C. 15 \\
\hline
\end{tabular}
\end{center}

\chapter{FOCUS 17 Chemical kinetics}
This Focus introduces the principles of 'chemical kinetics', the study of reaction rates. The rate of a chemical reaction might depend on variables that can be controlled, such as the pressure, the temperature, and the presence of a catalyst, and it is possible to optimize the rate by the appropriate choice of conditions.

\subsection*{17A The rates of chemical reactions}
This Topic discusses the definition of reaction rate and outlines the techniques for its measurement. The results of such measurements show that reaction rates depend on the concentration of reactants (and sometimes products) and 'rate constants' that are characteristic of the reaction. This dependence can be expressed in terms of differential equations known as 'rate laws'.\\
17A. 1 Monitoring the progress of a reaction; 17A. 2 The rates of reactions

\subsection*{17B Integrated rate laws}
'Integrated rate laws' are the solutions of the differential equations that describe rate laws. They are used to predict the concentrations of species at any time after the start of the reaction and to provide procedures for measuring rate constants. This Topic explores some simple yet useful integrated rate laws that appear throughout the Focus.\\
17B. 1 Zeroth-order reactions; 17B. 2 First-order reactions;\\
17B. 3 Second-order reactions

\subsection*{17C Reactions approaching equilibrium}
In general, rate laws must take into account both the forward and reverse reactions and describe the approach to equilibrium, when the forward and reverse rates are equal. The\\
results of the analysis are relations, which can be explored experimentally, between the equilibrium constant of the overall process and the rate constants of the forward and reverse reactions in the proposed mechanism.

\begin{verbatim}
17C.1 First-order reactions approaching equilibrium;
17C. 2 Relaxation methods
\end{verbatim}

\subsection*{17D The Arrhenius equation}
The rate constants of most reactions increase with increasing temperature. This Topic introduces the 'Arrhenius equation', which captures this temperature dependence by using only two parameters that can be determined experimentally.\\
17D. 1 The temperature dependence of reaction rates;\\
17D. 2 The interpretation of the Arrhenius parameters

\subsection*{17E Reaction mechanisms}
The study of reaction rates also leads to an understanding of the 'mechanisms' of reactions, their analysis into a sequence of elementary steps. This Topic shows how to construct rate laws from a proposed mechanism. The elementary steps themselves have simple rate laws which can be combined into an overall rate law by invoking the concept of the 'rate-determining step' of a reaction, by making the 'steady-state approximation', or by supposing that a 'pre-equilibrium' exists.\\
17E. 1 Elementary reactions; 17E. 2 Consecutive elementary reactions;\\
17E. 3 The steady-state approximation; 17E. 4 The rate-determining\\
step; 17E. 5 Pre-equilibria; 17E. 6 Kinetic and thermodynamic control\\
of reactions

\subsection*{17F Examples of reaction mechanisms}
This Topic develops three examples of reaction mechanisms. The first describes a special class of reactions in the gas phase\\
that depend on the collisions between reactants. The second gives insight into the formation of polymers and shows how the kinetics of their formation affects their properties. The third examines the general mechanism of action of 'enzymes', which are biological catalysts.\\
17F. 1 Unimolecular reactions; 17A. 2 Polymerization kinetics;\\
17F. 3 Enzyme-catalysed reactions

\subsection*{17G Photochemistry}
'Photochemistry' is the study of reactions that are initiated by light. This Topic explores the fate of the electronically excited molecules formed by the absorption of photons. One possible fate is energy transfer to another molecule; this process is par-\\
ticularly interesting as it is the basis for a method of estimating the distance between certain groups in large molecules.\\
17G. 1 Photochemical processes; 17G. 2 The primary quantum yield; 17G. 3 Mechanism of decay of excited singlet states; 17G. 4 Quenching; 17G. 5 Resonance energy transfer

\subsection*{Web resource What is an application of this material?}
Plants, algae, and some species of bacteria have evolved apparatus that performs 'photosynthesis', the capture of visible and near-infrared radiation for the purpose of synthesizing complex molecules in the cell. Impact 26 introduces the reaction steps involved.

\section*{TOPIC 17A The rates of chemical reactions}
\textbackslash begin\{abstract\}\\
> Why do you need to know this material?\\
Studies of the rates of consumption of reactants and formation of products make it possible to predict how quickly a reaction mixture approaches equilibrium. They also lead to detailed descriptions of the molecular events that transform reactants into products.\\
\textbackslash end\{abstract\}

\subsection*{- What is the key idea?}
Reaction rates are expressed as rate laws, which are empirical summaries of the rates in terms of the concentrations of reactants and, in some cases, products.

\subsection*{> What do you need to know already?}
This introductory Topic is the foundation of a sequence: all you need to be aware of initially is the significance of stoichiometric numbers (Topic 2C). For more background on the spectroscopic determination of concentration, refer to Topic 11A.

Chemical kinetics is the study of reaction rates. Experiments show that reaction rates depend on the concentration of reactants (and in some cases products) in characteristic ways that can be expressed in terms of differential equations known as 'rate laws'.

\subsection*{17A. 1 Monitoring the progress of a reaction}
The first step in the kinetic analysis of reactions is to establish the stoichiometry of the reaction and identify any side reactions. The basic data of chemical kinetics are then the concentrations of the reactants and products at different times after a reaction has been initiated.

\subsection*{(a) General considerations}
The rates of most chemical reactions are sensitive to the temperature (as described in Topic 17D), so in conventional experiments the temperature of the reaction mixture must be held constant throughout the course of the reaction. This requirement puts severe demands on the design of an experi-\\
ment. Gas-phase reactions, for instance, are often carried out in a vessel held in contact with a substantial block of metal. Liquid-phase reactions must be carried out in an efficient thermostat. Special efforts have to be made to study reactions at low temperatures, as in the study of the kinds of reactions that take place in interstellar clouds. Thus, supersonic expansion of the reacting gas can be used to attain temperatures as low as 10 K . For work in the liquid phase and the solid phase, very low temperatures are often reached by flowing cold liquid or cold gas around the reaction vessel. Alternatively, the entire reaction vessel is immersed in a thermally insulated container filled with a cryogenic liquid, such as liquid helium (for work at around 4 K ) or liquid nitrogen (for work at around 77 K ). Non-isothermal conditions are sometimes employed. For instance, the shelf life of an expensive pharmaceutical may be explored by slowly raising the temperature of a single sample.

Spectroscopy is widely applicable to the study of reaction kinetics, and is especially useful when one substance in the reaction mixture has a strong characteristic absorption in a conveniently accessible region of the electromagnetic spectrum. For example, the progress of the reaction $\mathrm{H}_{2}(\mathrm{~g})+$ $\mathrm{Br}_{2}(\mathrm{~g}) \rightarrow 2 \mathrm{HBr}(\mathrm{g})$ can be followed by measuring the absorption of visible radiation by bromine. A reaction that changes the number or type of ions present in a solution may be followed by monitoring the electrical conductivity of the solution. The replacement of neutral molecules by ionic products can result in dramatic changes in the conductivity, as in the reaction $\left(\mathrm{CH}_{3}\right)_{3} \mathrm{CCl}(\mathrm{aq})+\mathrm{H}_{2} \mathrm{O}(\mathrm{l}) \rightarrow\left(\mathrm{CH}_{3}\right)_{3} \mathrm{COH}(\mathrm{aq})+$ $\mathrm{H}^{+}(\mathrm{aq})+\mathrm{Cl}^{-}(\mathrm{aq})$. If hydrogen ions are produced or consumed, the reaction may be followed by monitoring the pH of the solution.

Other methods of determining composition include emission spectroscopy (Topic 11F), mass spectrometry, gas chromatography, nuclear magnetic resonance (Topics 12B and 12C), and electron paramagnetic resonance (for reactions involving radicals or paramagnetic d-metal ions; Topic 12D). A reaction in which at least one component is a gas might result in an overall change in pressure in a system of constant volume, so its progress may be followed by recording the variation of pressure with time.

\subsection*{Example 17A. 1 Relating the variation in the total pressure to the partial pressures of the species present}
The decomposition $\mathrm{N}_{2} \mathrm{O}_{5}(\mathrm{~g}) \rightarrow 2 \mathrm{NO}_{2}(\mathrm{~g})+\frac{1}{2} \mathrm{O}_{2}(\mathrm{~g})$ is monitored by measuring the total pressure in a constant-volume\\
reaction vessel held at constant temperature. If the initial pressure is $p_{0}$, and is due solely to $\mathrm{N}_{2} \mathrm{O}_{5}$, and the total pressure at any later time is $p$, derive expressions for the partial pressures of all three species in terms of $p_{0}$ and $p$.

Collect your thoughts You can assume perfect-gas behaviour, in which case the partial pressure of a gas is proportional to its amount. Then imagine that a certain amount of $\mathrm{N}_{2} \mathrm{O}_{5}$ decomposes such that its partial pressure falls from $p_{0}$ to $p_{0}-\Delta p$. Because $1 \mathrm{~mol}_{2} \mathrm{O}_{5}$ decomposes to give $2 \mathrm{~mol} \mathrm{NO}_{2}$, the partial pressure of $\mathrm{NO}_{2}$ increases from 0 to $2 \Delta p$. Likewise, the partial pressure of $\mathrm{O}_{2}$ increases from 0 to $\frac{1}{2} \Delta p$. The total pressure $p$ is the sum of the partial pressures of the three components.

The solution Draw up the following table:

\begin{center}
\begin{tabular}{lllll}
\hline
 & $p\left(\mathrm{~N}_{2} \mathrm{O}_{5}\right)$ & $p\left(\mathrm{NO}_{2}\right)$ & $p\left(\mathrm{O}_{2}\right)$ & $p$ \\
\hline
Initially & $p_{0}$ & 0 & 0 & $p_{0}$ \\
Later & $p_{0}-\Delta p$ & $2 \Delta p$ & $\frac{1}{2} \Delta p$ & $p_{0}+\frac{3}{2} \Delta p$ \\
Equivalent to & $\frac{5}{3} p_{0}-\frac{2}{3} p$ & $\frac{4}{3}\left(p-p_{0}\right)$ & $\frac{1}{3}\left(p-p_{0}\right)$ &  \\
\hline
\end{tabular}
\end{center}

The bottom line follows from $p=p_{0}+\frac{3}{2} \Delta p$, rearranged into $\Delta p=\frac{2}{3}\left(p-p_{0}\right)$, and then substituted into the line above.

Comment. A check of the calculation is to note that when all the $\mathrm{N}_{2} \mathrm{O}_{5}$ has been consumed its partial pressure is zero, which implies that $\frac{5}{3} p_{0}-\frac{2}{3} p_{\text {final }}=0$ and hence $p_{\text {final }}=\frac{5}{2} p_{0}$. This result is expected because $1 \mathrm{~mol} \mathrm{~N}_{2} \mathrm{O}_{5}$ is replaced by $2 \mathrm{~mol} \mathrm{NO}{ }_{2}$ and $\frac{1}{2} \mathrm{~mol} \mathrm{O}_{2}$ : the total amount of molecules therefore goes from 1 mol to $2 \frac{1}{2} \mathrm{~mol}$, resulting in the pressure increasing by a factor of $2 \frac{1}{2}=\frac{5}{2}$.

Self-test 17A. 1 Repeat the calculation of the partial pressures for the species in the reaction $2 \operatorname{NOBr}(\mathrm{~g}) \rightarrow 2 \mathrm{NO}(\mathrm{g})+\mathrm{Br}_{2}(\mathrm{~g})$, assuming that the initial pressure is $p_{0}$ and is due to NOBr alone.

$$
\left({ }^{0} d-d\right) \frac{\tau}{\mathrm{I}}={ }^{\tau_{\mathrm{IA}}} d{ }^{\circ} d-d={ }^{\mathrm{oN}} d \cdot d \frac{\tau}{\mathrm{I}}-{ }^{0} d \frac{\tau}{\varepsilon}={ }^{\mathrm{xq} \mathrm{\circ N}} d: \iota \partial \mathrm{Ms} u \mathrm{~V}
$$

\subsection*{(b) Special techniques}
The method used to monitor concentrations depends on the species involved and the rapidity with which their concentrations change. Many reactions reach equilibrium over periods of minutes or hours, and several techniques may then be used to follow the changing concentrations. In a real-time analysis the composition of the system is analysed while the reaction is in progress. Either a small sample is withdrawn or the bulk solution is monitored.\\
In the flow method the reactants are mixed as they flow together into a chamber (Fig. 17A.1). The reaction continues as the thoroughly mixed solutions flow through the outlet tube, and observation of the composition at different positions along the tube is equivalent to the observation of the reaction mixture at different times after mixing. The disadvantage of\\
\includegraphics[max width=\textwidth, center]{2025_02_27_d4b95d9718f0b9e2e6b9g-756(1)}

Figure 17A. 1 The arrangement used in the flow technique for studying reaction rates. The reactants are injected into the mixing chamber at a steady rate. The location of the spectrometer corresponds to different times after initiation of the reaction.\\
\includegraphics[max width=\textwidth, center]{2025_02_27_d4b95d9718f0b9e2e6b9g-756}

Figure 17A. 2 In the stopped-flow technique the reagents are driven quickly into the mixing chamber by the driving pistons and then the time dependence of the concentrations is monitored.\\
conventional flow techniques is that a large volume of reactant solution is necessary. This requirement makes the study of fast reactions particularly difficult because to spread the reaction over a length of tube the flow must be rapid. This disadvantage is avoided by the stopped-flow technique, in which the reagents are mixed very quickly in a small chamber fitted with a movable piston instead of an outlet tube (Fig. 17A.2). The flow pushes the piston back and ceases when it reaches a stop; the reaction continues in the mixed solutions. Observations, commonly using spectroscopic techniques such as ultravioletvisible absorption and fluorescence emission, are made on the sample as a function of time. The technique permits the study of reactions that occur on the millisecond to second timescale. The suitability of the stopped-flow method for the study of small samples means that it is appropriate for many biochemical reactions; it has been widely used to study the kinetics of protein folding and enzyme action.

Very fast reactions can be studied by flash photolysis, in which the sample is exposed to a brief flash of light, which initiates the reaction, and then the contents of the reaction chamber are monitored by electronic absorption or emission, infrared absorption, or Raman scattering. In the arrangement shown in Fig. 17A. 3 a strong and short laser pulse, the pump, promotes a molecule A to an excited electronic state $\mathrm{A}^{*}$ which can either emit a photon (as fluorescence or phosphorescence)\\
\includegraphics[max width=\textwidth, center]{2025_02_27_d4b95d9718f0b9e2e6b9g-757}

Figure 17A. 3 A configuration used for flash photolysis, in which the same pulsed laser is used to generate a monochromatic pump pulse and, after continuum generation, a 'white' light probe pulse. The time delay between the pump and probe pulses may be varied.\\
or react with another species $B$ to form first the intermediate AB and then the product C :

$$
\begin{array}{ll}
\mathrm{A}+h v \rightarrow \mathrm{~A}^{*} & \text { (absorption) } \\
\mathrm{A}^{*} \rightarrow \mathrm{~A} & \text { (emission) } \\
\mathrm{A}^{*}+\mathrm{B} \rightarrow \mathrm{AB} \rightarrow \mathrm{C} & \text { (reaction) }
\end{array}
$$

The rates of appearance and disappearance of the various species are determined by observing time-dependent changes in the absorption spectrum of the sample during the course of the reaction. This monitoring is done by passing a weak pulse of white light, the probe, through the sample at different times after the laser pulse. Pulsed 'white' light can be generated directly from the laser pulse by the phenomenon of continuum generation, in which focusing a short laser pulse on sapphire or a vessel containing water or carbon tetrachloride results in an outgoing beam with a wide distribution of frequencies. A time delay between the strong laser pulse and the 'white' light pulse can be introduced by allowing one of the beams to travel a longer distance before reaching the sample. For example, a difference in travel distance of $\Delta d=3 \mathrm{~mm}$ corresponds to a time delay $\Delta t=\Delta d / c \approx 10 \mathrm{ps}$ between two beams, where $c$ is the speed of light. The relative distances travelled by the two beams in Fig. 17A. 3 are controlled by directing the 'white' light beam to a motorized stage carrying a pair of mirrors.

In contrast to real-time analysis, quenching methods are based on 'quenching', or stopping, the reaction after it has been allowed to proceed for a certain time. In this way the composition is analysed at leisure and reaction intermediates may be trapped. These methods are suitable only for reactions slow enough for there to be little reaction during the time it takes to quench the mixture. In the chemical quench flow method, the reactants are mixed in much the same way as in the flow method but the reaction is quenched by another reagent, such as a solution of acid or base, after the mixture has\\
travelled along a fixed length of the outlet tube. Different reaction times can be selected by varying the flow rate along the outlet tube. An advantage of the chemical quench flow method over the stopped-flow method is that rapid spectroscopic measurements are not needed in order to measure the concentration of reactants and products. Once the reaction has been quenched, the solution may be examined by 'slow' techniques, such as mass spectrometry and chromatography. In the freeze quench method, the reaction is quenched by cooling the mixture within milliseconds and the concentrations of reactants, intermediates, and products are measured spectroscopically.

\subsection*{17A.2 The rates of reactions}
Reaction rates depend on the composition and the temperature of the reaction mixture. The next few sections look at these observations in more detail.

\subsection*{(a) The definition of rate}
Consider a reaction of the form $\mathrm{A}+2 \mathrm{~B} \rightarrow 3 \mathrm{C}+\mathrm{D}$, in which at some instant the molar concentration of a participant J is [J] and the volume of the system is constant. The instantaneous rate of consumption of a reactant or formation of a product is the slope of the tangent to the graph of concentration against time (expressed as a positive quantity). It follows that the instantaneous rate of consumption of one of the reactants at a given time is $-\mathrm{d}[\mathrm{R}] / \mathrm{d} t$, where R is A or B . This rate is a positive quantity (Fig. 17A.4). The rate of formation of one of the products (C or D , denoted P ) is $\mathrm{d}[\mathrm{P}] / \mathrm{d} t$ (note the difference in sign). This rate is also positive.\\
\includegraphics[max width=\textwidth, center]{2025_02_27_d4b95d9718f0b9e2e6b9g-757(1)}

Figure 17A. 4 The definition of (instantaneous) rate as the slope of the tangent drawn to the curve showing the variation of concentration of (a) products, (b) reactants with time. For negative slopes, the sign is changed when reporting the rate, so all reaction rates are positive.

It follows from the stoichiometry of the reaction $\mathrm{A}+2 \mathrm{~B} \rightarrow$ $3 \mathrm{C}+\mathrm{D}$ that

$$
\frac{\mathrm{d}[\mathrm{D}]}{\mathrm{d} t}=\frac{1}{3} \frac{\mathrm{~d}[\mathrm{C}]}{\mathrm{d} t}=-\frac{\mathrm{d}[\mathrm{~A}]}{\mathrm{d} t}=-\frac{1}{2} \frac{\mathrm{~d}[\mathrm{~B}]}{\mathrm{d} t}
$$

so there are several rates connected with the reaction. The undesirability of having different rates to describe the same reaction is avoided by introducing the extent of reaction, $\xi$ (xi), which is defined so that for each species J in the reaction, the change in amount of $\mathrm{J}, \mathrm{d} n_{\mathrm{J}}$, is


\begin{equation*}
\mathrm{d} n_{\mathrm{J}}=v_{\mathrm{J}} \mathrm{~d} \xi \tag{17A.1}
\end{equation*}


Extent of reaction [definition]\\
where $v_{\mathrm{J}}$ is the stoichiometric number of the species (Topic 2 C ; remember that $v_{\mathrm{J}}$ is negative for reactants and positive for products). The unique rate of reaction, $v$, is then defined as


\begin{equation*}
v=\frac{1}{V} \frac{\mathrm{~d} \xi}{\mathrm{~d} t} \tag{17A.2}
\end{equation*}


Rate of reaction [definition]\\
where $V$ is the volume of the system. For any species $\mathrm{J}, \mathrm{d} \xi=$ $\mathrm{d} n_{\mathrm{J}} / v_{\mathrm{J}}$, so


\begin{equation*}
v=\frac{1}{v_{\mathrm{J}}} \times \frac{1}{V} \frac{\mathrm{~d} n_{\mathrm{J}}}{\mathrm{~d} t} \tag{17A.3a}
\end{equation*}


For a homogeneous reaction in a constant-volume system the volume $V$ can be taken inside the differential and $n_{\mathrm{I}} / V$ is written as the molar concentration [J] to give


\begin{equation*}
v=\frac{1}{v_{\mathrm{J}}} \frac{\mathrm{~d}[\mathrm{~J}]}{\mathrm{d} t} \tag{17A.3b}
\end{equation*}


For a heterogeneous reaction the (constant) surface area, $A$, occupied by the species is used in place of $V$. Then, because the surface concentration is $\sigma_{\mathrm{J}}=n_{\mathrm{J}} / A$ it follows that


\begin{equation*}
v=\frac{1}{v_{\mathrm{J}}} \frac{\mathrm{~d} \sigma_{\mathrm{J}}}{\mathrm{~d} t} \tag{17A.3c}
\end{equation*}


In each case there is now a single rate for the reaction (for the chemical equation as written). With molar concentrations in moles per cubic decimetre and time in seconds, reaction rates of homogeneous reactions are reported in moles per cubic decimetre per second (moldm $\mathrm{s}^{-3}$ ) or related units. For gas-phase reactions, such as those taking place in the atmosphere, concentrations are often expressed in molecules per cubic centimetre (molecules $\mathrm{cm}^{-3}$ ) and rates in molecules per cubic centimetre per second (molecules $\mathrm{cm}^{-3} \mathrm{~s}^{-1}$ ). For heterogeneous reactions, rates are expressed in moles per square metre per second $\left(\mathrm{mol} \mathrm{m}^{-2} \mathrm{~s}^{-1}\right)$ or related units.

\subsection*{Brief illustration 17A. 1}
The rate of formation of $\mathrm{NO}, \mathrm{d}[\mathrm{NO}] / \mathrm{d} t$, in the reaction $2 \mathrm{NOBr}(\mathrm{g}) \rightarrow 2 \mathrm{NO}(\mathrm{g})+\mathrm{Br}_{2}(\mathrm{~g})$ is reported as $0.16 \mathrm{mmol} \mathrm{dm}^{-3} \mathrm{~s}^{-1}$. Because $v_{\mathrm{NO}}=+2$, the rate of the reaction is reported as\\
$v=\frac{1}{2} \mathrm{~d}[\mathrm{NO}] / \mathrm{d} t=0.080 \mathrm{mmoldm}^{-3} \mathrm{~s}^{-1}$. Because $v_{\mathrm{NOBr}}=-2$ it follows that the rate of reaction can be written in terms of $[\mathrm{NOBr}]$ as $v=-\frac{1}{2} \mathrm{~d}[\mathrm{NOBr}] / \mathrm{d} t$; hence $\mathrm{d}[\mathrm{NOBr}] / \mathrm{d} t=-2 v=$ $-0.16 \mathrm{mmoldm}^{-3} \mathrm{~s}^{-1}$. The rate of consumption of NOBr is therefore $0.16 \mathrm{mmoldm}^{-3} \mathrm{~s}^{-1}$, or $9.6 \times 10^{16}$ molecules $\mathrm{cm}^{-3} \mathrm{~s}^{-1}$.

\subsection*{(b) Rate laws and rate constants}
The rate of reaction is often found to be proportional to the concentrations of the reactants raised to a power. For example, the rate of a reaction might be found to be proportional to the molar concentrations of two reactants $A$ and $B$, so


\begin{equation*}
v=k_{\mathrm{r}}[\mathrm{~A}][\mathrm{B}] \tag{17.4}
\end{equation*}


The constant of proportionality $k_{\mathrm{r}}$ is called the rate constant for the reaction; it is independent of the concentrations but depends on the temperature. An experimentally determined equation of this kind is called the rate law of the reaction. More formally, a rate law is an equation that expresses the rate of reaction in terms of the concentrations of all the species present in the overall chemical equation for the reaction at the time of interest:

\[
v=f([\mathrm{~A}],[\mathrm{B}], \ldots) \quad \begin{align*}
& \text { Rate law in terms of concentrations }  \tag{17A.5a}\\
& \text { [general form] }
\end{align*}
\]

For homogeneous gas-phase reactions, it is often more convenient to express the rate law in terms of partial pressures, which for perfect gases are related to molar concentrations by $p_{\mathrm{J}}=R T[\mathrm{~J}]$. In this case,

\[
v=f\left(p_{\mathrm{A}}, p_{\mathrm{B}}, \ldots\right) \quad \begin{align*}
& \text { Rate law in terms of partial pressures }  \tag{17A.5b}\\
& {[\text { general form }]}
\end{align*}
\]

The rate law of a reaction is determined experimentally, and cannot in general be inferred from the chemical equation for the reaction. The reaction of hydrogen and bromine, for example, has a very simple stoichiometry, $\mathrm{H}_{2}(\mathrm{~g})+\mathrm{Br}_{2}(\mathrm{~g}) \rightarrow$ $2 \mathrm{HBr}(\mathrm{g})$, but its rate law is complicated:


\begin{equation*}
v=\frac{k_{\mathrm{a}}\left[\mathrm{H}_{2}\right]\left[\mathrm{Br}_{2}\right]^{3 / 2}}{\left[\mathrm{Br}_{2}\right]+k_{\mathrm{b}}[\mathrm{HBr}]} \tag{17A.6}
\end{equation*}


In certain cases the rate law does reflect the stoichiometry of the reaction; but that is either a coincidence or reflects a feature of the underlying reaction mechanism (Topic 17E).

A note on good (or, at least, our) practice A general rate constant is denoted $k_{\mathrm{r}}$ to distinguish it from the Boltzmann constant $k$. In some texts $k$ is used for the former and $k_{\mathrm{B}}$ for the latter. When expressing the rate constants in a more complicated rate law, such as that in eqn 17A. 6 , we use $k_{\mathrm{a}}, k_{\mathrm{b}}$, and so on.

The units of $k_{\mathrm{r}}$ are always such as to convert the product of concentrations, each raised to the appropriate power, into a\\
rate expressed as a change in concentration divided by time. For example, if the rate law is the one shown in eqn 17A.4, with concentrations expressed in $\mathrm{mol} \mathrm{dm}{ }^{-3}$, then the units of $k_{\mathrm{r}}$ will be $\mathrm{dm}^{3} \mathrm{~mol}^{-1} \mathrm{~s}^{-1}$ because

$$
\overbrace{\mathrm{dm}^{3} \mathrm{~mol}^{-1} \mathrm{~s}^{-1}}^{k_{i}} \times \overbrace{\mathrm{moldm}^{-3}}^{[A]} \times \overbrace{\mathrm{moldm}^{-3}}^{[B]}=\operatorname{moldm}^{-3} \mathrm{~s}^{-1}
$$

which are the units of $v$. If the concentrations are expressed in molecules $\mathrm{cm}^{-3}$, and the rate in molecules $\mathrm{cm}^{-3} \mathrm{~s}^{-1}$, then the rate constant is expressed in $\mathrm{cm}^{3}$ molecule ${ }^{-1} \mathrm{~s}^{-1}$. The approach just developed can be used to determine the units of the rate constant from rate laws of any form.

\subsection*{Brief illustration 17A. 2}
The rate constant for the reaction $\mathrm{O}(\mathrm{g})+\mathrm{O}_{3}(\mathrm{~g}) \rightarrow 2 \mathrm{O}_{2}(\mathrm{~g})$ is $8.0 \times 10^{-15} \mathrm{~cm}^{3}$ molecule ${ }^{-1} \mathrm{~s}^{-1}$ at 298 K . To express this rate constant in $\mathrm{dm}^{3} \mathrm{~mol}^{-1} \mathrm{~s}^{-1}$, make use of the relation $1 \mathrm{~cm}=10^{-1} \mathrm{dm}$ to convert the volume:

$$
\begin{aligned}
k_{\mathrm{r}} & =8.0 \times 10^{-15} \overbrace{\mathrm{~cm}^{3}}^{\left(10^{-1} \mathrm{dm}\right)^{3}} \text { molecule }{ }^{-1} \mathrm{~s}^{-1} \\
& =8.0 \times 10^{-18} \mathrm{dm}^{3} \text { molecule }{ }^{-1} \mathrm{~s}^{-1}
\end{aligned}
$$

Now note that the number of molecules can be expressed as an amount in moles by division by Avogadro's constant expressed as molecules per mole:

$$
\begin{aligned}
k_{\mathrm{r}} & =8.0 \times 10^{-18} \mathrm{dm}^{3} \text { molecule } \mathrm{e}^{-1} \mathrm{~s}^{-1} \\
& =8.0 \times 10^{-18} \mathrm{dm}^{3} \times\left(\frac{1 \text { molecule }}{6.022 \times 10^{23} \text { molecules } \mathrm{mol}^{-1}}\right)^{-1} \mathrm{~s}^{-1} \\
& =8.0 \times 10^{-18} \times 6.022 \times 10^{23} \mathrm{dm}^{3} \mathrm{~mol}^{-1} \mathrm{~s}^{-1} \\
& =4.8 \times 10^{6} \mathrm{dm}^{3} \mathrm{~mol}^{-1} \mathrm{~s}^{-1}
\end{aligned}
$$

A practical application of a rate law is that once the law and the value of the rate constant are known, it is possible to predict the rate of reaction from the composition of the mixture. Moreover, as demonstrated in Topic 17B, by knowing the rate law, it is also possible to predict the composition of the reaction mixture at a later stage of the reaction. A rate law also provides evidence used to assess the plausibility of a proposed mechanism of the reaction. This application is developed in Topic 17E.

\subsection*{(c) Reaction order}
Many reactions are found to have rate laws of the form


\begin{equation*}
v=k_{\mathrm{r}}[\mathrm{~A}]^{a}[\mathrm{~B}]^{b} \ldots \tag{17A.7}
\end{equation*}


The power to which the concentration of a species (a product or a reactant) is raised in a rate law of this kind is the order of the reaction with respect to that species. A reaction with the rate law in eqn 17A. 4 is first order in A and first order in B . The overall order of a reaction with a rate law like that in eqn 17 A .7 is the sum of the individual orders, $a+b+\cdots$. The overall order of the rate law in eqn 17A. 4 is $1+1=2$; the rate law is therefore said to be second-order overall.

A reaction need not have an integral order, and many gas-phase reactions do not. For example, a reaction with the rate law


\begin{equation*}
v=k_{\mathrm{r}}[\mathrm{~A}]^{1 / 2}[\mathrm{~B}] \tag{17A.8}
\end{equation*}


is half order in A , first order in B , and three-halves order overall.

\subsection*{Brief illustration 17A. 3}
The experimentally determined rate law for the gas-phase reaction $\mathrm{H}_{2}(\mathrm{~g})+\mathrm{Br}_{2}(\mathrm{~g}) \rightarrow 2 \mathrm{HBr}(\mathrm{g})$ is given by eqn 17A.6. In the rate law the concentration of $\mathrm{H}_{2}$ appears raised to the power +1 , so the reaction is first order in $\mathrm{H}_{2}$. However, the concentrations of $\mathrm{Br}_{2}$ and HBr do not appear as a single term raised to a power, so the reaction has an indefinite order with respect to both $\mathrm{Br}_{2}$ and HBr , and an indefinite order overall.

Some reactions obey a zeroth-order rate law, and therefore have a rate that is independent of the concentration of the reactant (so long as some is present). Thus, the catalytic decomposition of phosphine $\left(\mathrm{PH}_{3}\right)$ on hot tungsten at high pressures has the rate law


\begin{equation*}
v=k_{\mathrm{r}} \tag{17A.9}
\end{equation*}


This law means that $\mathrm{PH}_{3}$ decomposes at a constant rate until it has entirely disappeared.

As seen in Brief illustration 17A.3, when a rate law is not of the form in eqn 17A.7, the reaction does not have an overall order and might not even have definite orders with respect to each participant.

These remarks point to three important tasks:

\begin{itemize}
  \item To identify the rate law and obtain the rate constant from the experimental data. This aspect is discussed in this Topic.
  \item To account for the values of the rate constants and explain their temperature dependence. This task is undertaken in Topic 17D.
  \item To construct reaction mechanisms consistent with the rate law. The techniques for doing so are introduced in Topic 17E.
\end{itemize}

\subsection*{(d) The determination of the rate law}
The determination of a rate law is simplified by the isolation method, in which all the reactants except one are present in large excess. The dependence of the rate on each of the reactants can be found by isolating each of them in turn-by having all the other substances present in large excess-and piecing together a picture of the overall rate law.

If a reactant $B$ is in large excess its concentration is nearly constant throughout the reaction. Then, although the true rate law might be $v=k_{\mathrm{r}}[\mathrm{A}][\mathrm{B}]^{2}$, the current value of $[\mathrm{B}]$ can be approximated by its initial value $[\mathrm{B}]_{0}$ (from which it hardly changes in the course of the reaction) to give

$$
v=k_{\mathrm{r}, \text { eff }}[\mathrm{A}] \text {, with } k_{\mathrm{r}, \mathrm{eff}}=k_{\mathrm{r}}[\mathrm{~B}]_{0}^{2} \quad \begin{aligned}
& \text { A pseudofirst- } \\
& \text { order reaction, } \\
& \text { B in excess }
\end{aligned}
$$

(17A.10a)

Because the true rate law has been forced into first-order form by assuming a constant B concentration, the effective rate law is classified as a pseudofirst-order rate law and $k_{\text {reff }}$ is called the effective rate constant for a given, fixed concentration of B . If, instead, the concentration of $A$ is in large excess, and hence effectively constant, then the original rate law simplifies to


\begin{align*}
& v=k_{\mathrm{r}, \text { eff }}^{\prime}[\mathrm{B}]^{2} \text {, now with } k_{\mathrm{r}, \text { eff }}^{\prime}=k_{\mathrm{r}}[\mathrm{~A}]_{0} \\
& \text { A pseudosecond-order reaction, A in excess } \tag{17A.10b}
\end{align*}


This pseudosecond-order rate law is also much easier to analyse and identify than the complete law. Note that the order of the reaction and the form of the effective rate constant change according to whether A or B is in excess. In a similar manner, a reaction may even appear to be zeroth order. Many reactions in aqueous solution that are reported as first or second order are actually pseudofirst or pseudosecond order: the solvent water, for instance, might participate in a reaction but it is in such large excess that its concentration remains constant.

In the method of initial rates, which is often used in conjunction with the isolation method, the instantaneous rate is measured at the beginning of the reaction for several different initial concentrations of the isolated reactant. If the initial rate is doubled when the concentration of an isolated reactant A is doubled, then the reaction is first order in A ; if the initial rate is quadrupled, then the reaction is second order in A. More formally, with an eye on developing a graphical method for determining the order, suppose the rate law for a reaction with A isolated is

$$
v=k_{\mathrm{r}}[\mathrm{~A}]^{a}
$$

Then the initial rate of the reaction, $v_{0}$, is given by the initial concentration of A:

$$
v_{0}=k_{\mathrm{r}, \text { eff }}[\mathrm{A}]_{0}^{a} \quad \text { Initial rate of an ath-order reaction } \quad \text { (17A.11a) }
$$

Taking (common) logarithms gives


\begin{equation*}
\log v_{0}=\log \left(k_{\mathrm{r}, \mathrm{eff}}[\mathrm{~A}]_{0}^{a}\right) \stackrel{\log (x y)=\log x+\log y}{=} \log k_{\mathrm{r}, \mathrm{eff}}+\log [\mathrm{A}]_{0}^{a} \tag{17A.11b}
\end{equation*}


$$
\begin{aligned}
& \log x^{a}=a \log x \\
& =\log k_{\mathrm{r}, \mathrm{eff}}+a \log [\mathrm{~A}]_{0}
\end{aligned}
$$

This equation has the form of the equation for a straight line:


\begin{equation*}
\overbrace{\log v_{0}}^{y}=\overbrace{\log k_{\mathrm{r}, \text { eff }}}^{\text {intercept }}+\overbrace{a \log [\mathrm{~A}]_{0}}^{\text {slope } \times x} \tag{17A.11c}
\end{equation*}


It follows that, for a series of initial concentrations, a plot of the logarithms of the initial rates against the logarithms of the initial concentrations of A should be a straight line, and that the slope of the graph is $a$, the order of the reaction with respect to A .

\subsection*{Example 17A. 2}
\subsection*{Using the method of initial rates}
The recombination of I atoms in the gas phase in the presence of argon was investigated and the order of the reaction was determined by the method of initial rates. The initial rates of reaction of $2 \mathrm{I}(\mathrm{g})+\operatorname{Ar}(\mathrm{g}) \rightarrow \mathrm{I}_{2}(\mathrm{~g})+\operatorname{Ar}(\mathrm{g})$ were as follows:

$$
\begin{array}{lllll}
{[\mathrm{I}]_{0} /\left(10^{-5} \mathrm{~mol} \mathrm{dm}^{-3}\right)} & 1.0 & 2.0 & 4.0 & 6.0 \\
v_{0} /\left(\mathrm{mol} \mathrm{dm}^{-3} \mathrm{~s}^{-1}\right) & \text { (a) } 8.70 \times 10^{-4} & 3.48 \times 10^{-3} & 1.39 \times 10^{-2} & 3.13 \times 10^{-2} \\
& \text { (b) } 4.35 \times 10^{-3} & 1.74 \times 10^{-2} & 6.96 \times 10^{-2} & 1.57 \times 10^{-1} \\
& \text { (c) } 8.69 \times 10^{-3} & 3.47 \times 10^{-2} & 1.38 \times 10^{-1} & 3.13 \times 10^{-1}
\end{array}
$$

The Ar concentrations are (a) $1.0 \times 10^{-3} \mathrm{moldm}^{-3}$, (b) $5.0 \times$ $10^{-3} \mathrm{moldm}^{-3}$, and (c) $1.0 \times 10^{-2} \mathrm{moldm}^{-3}$. Find the orders of reaction with respect to I and Ar, and the rate constant.\\[0pt]
Collect your thoughts You need to identify sets of data in which only one reactant is changing (such as each row of data for constant [Ar]). The identification of order from such data involves the application of eqn 17A.11c, plotting the logarithm of the rate against the logarithm of a concentration of one of the reactants (in this case, arbitrarily chosen to be I). So, first tabulate the logarithms of the concentrations of I and the rates for the values at constant $[\mathrm{Ar}]_{0}$ (i.e. for each row of data). The slope gives the order with respect to [I] and the intercept at $\log [\mathrm{I}]_{0}=0$ gives $\log k_{\mathrm{r} \text {,eff }}$, with a different value for each $[\mathrm{Ar}]_{0}$. The effective rate constant obtained in this way is $k_{\mathrm{r} \text { eff }}=k_{\mathrm{r}}[\mathrm{Ar}]_{0}^{b}$, so to extract $k_{\mathrm{r}}$ and $b$, take logarithms, as in the text, to obtain

$$
\log k_{\mathrm{r}, \mathrm{eff}}=\log k_{\mathrm{r}}+b \log [\mathrm{Ar}]_{0}
$$

Now realize that you need to plot the $\log k_{\mathrm{r}, \text { eff }}$ found in the first part of the solution against $\log [A r]_{0}$. Then the slope gives $b$ and the intercept at $\log [\mathrm{Ar}]_{0}=0$ gives $\log k_{\mathrm{r}}$.\\
\includegraphics[max width=\textwidth, center]{2025_02_27_d4b95d9718f0b9e2e6b9g-761}

Figure 17A. 5 Analysis of the data in Example 17A.2. (a) Plots for finding the order with respect to $I$. The intercepts at $\log [I]_{0}=0$ are far to the right, and are shown in the inset. (b) The plots for finding the order with respect to Ar and the rate constant $k_{r}$. The intercept at $\log \left[\mathrm{Ar}_{0}=0\right.$ is far to the right, and is shown in the inset.

The solution The data give the following points for the graph:

\begin{center}
\begin{tabular}{llllll}
$\log \left([\mathrm{II}]_{0} / \mathrm{moldm}^{-3}\right)$ &  & -5.00 & -4.70 & -4.40 & -4.22 \\
$\log \left(v_{0} / \mathrm{moldm}^{-3} \mathrm{~s}^{-1}\right)$ & (a) -3.060 & -2.458 & -1.857 & -1.504 &  \\
 & (b) -2.362 & -1.759 & -1.157 & -0.804 &  \\
 & (c) -2.061 & -1.460 & -0.860 & -0.504 &  \\
\end{tabular}
\end{center}

The graph of the data for varying [I] but constant [Ar] is shown in Fig. 17A.5a. The slopes of the lines are 2, so the reaction is second order with respect to I. The effective rate constants $k_{\mathrm{r}, \mathrm{eff}}$ are as follows:

\begin{center}
\begin{tabular}{llll}
$[\mathrm{Ar}]_{0} /\left(\mathrm{mol} \mathrm{dm}^{-3}\right)$ & $1.0 \times 10^{-3}$ & $5.0 \times 10^{-3}$ & $1.0 \times 10^{-2}$ \\
$\log \left([\mathrm{Ar}]_{0} / \mathrm{mol} \mathrm{dm}^{-3}\right)$ & -3.00 & -2.30 & -2.00 \\
$\log \left(k_{\mathrm{r}, \mathrm{eff}} / \mathrm{dm}^{3} \mathrm{~mol}^{-1} \mathrm{~s}^{-1}\right)$ & 6.94 & 7.64 & 7.93 \\
\end{tabular}
\end{center}

Figure 17A.5b shows the plot of $\log \left\{k_{\text {reff }} /\left(\mathrm{dm}^{3} \mathrm{~mol}^{-1} \mathrm{~s}^{-1}\right)\right\}$ against $\log \left\{[\operatorname{Ar}]_{0} /\left(\mathrm{mol} \mathrm{dm}^{-3}\right)\right\}$. The slope is 1 , so $b=1$ and the reaction is first order with respect to Ar. The intercept at $\log \left\{[\operatorname{Ar}]_{0} /\left(\operatorname{moldm}{ }^{-3}\right)\right\}=0$ is $\log \left\{k_{\mathrm{r}, \text { eff }} /\left(\mathrm{dm}^{3} \mathrm{~mol}^{-1} \mathrm{~s}^{-1}\right)\right\}=9.94$, so $k_{\mathrm{r}}=8.7 \times 10^{9} \mathrm{dm}^{6} \mathrm{~mol}^{-2} \mathrm{~s}^{-1}$. The overall (initial) rate law is therefore $v=k_{\mathrm{r}}[\mathrm{I}]_{0}^{2}[\mathrm{Ar}]_{0}$.

A note on good practice When taking the common logarithm of a number of the form $x . x x \times 10^{n}$ (with $n<10$ ) there are four significant figures in the answer (for instance, $\log 1.23 \times 10^{4}=$ 4.090): the figure before the decimal point is simply the power of 10. Conversely, when taking the common antilogarithm of $y . y y y$, there are three significant figures in the answer (for instance, $10^{5.678}=4.76 \times 10^{5}$ ).

Self-test 17A.2 The initial rate of a certain reaction depended on the concentration of a substance J as follows:

\begin{center}
\begin{tabular}{lllll}
$[\mathrm{J}]_{0} /\left(10^{-3} \mathrm{~mol} \mathrm{dm}^{-3}\right)$ & 5.0 & 10.2 & 17 & 30 \\
$v_{0} /\left(10^{-7} \mathrm{moldm}^{-3} \mathrm{~s}^{-1}\right)$ & 3.6 & 9.6 & 4 & 130 \\
\end{tabular}
\end{center}

Find the order of the reaction with respect to $J$ and the rate constant.

$$
{ }_{\mathrm{I}-} \mathrm{s}_{\mathrm{I}} \mathrm{Iow}{ }_{\varepsilon} \operatorname{urp}_{\tau-} 0 \mathrm{I} \times 9 \cdot \mathrm{I} \text { ' } \tau: \jmath \partial m s u_{V}
$$

The method of initial rates might not reveal the full rate law because once the products have started to be generated they might participate in the reaction and affect its rate. For example, in the reaction between $\mathrm{H}_{2}$ and $\mathrm{Br}_{2}$, the rate law in eqn 17A. 6 shows that the rate depends on the concentration of the product HBr . To avoid this difficulty, the rate law should be fitted to the data throughout the reaction. The fitting may be done, in simple cases at least, by using a proposed rate law to predict the concentration of any component at any time, and comparing it with the data; methods based on this procedure are described in Topic 17B. A rate law should also be tested by observing whether the addition of products or, for gas-phase reactions, a change in the surface-to-volume ratio in the reaction chamber affects the rate.

\subsection*{Checklist of concepts}
\begin{enumerate}
  \item The rates of chemical reactions are measured by using techniques that monitor the concentrations of species present in the reaction mixture. Examples include realtime and quenching procedures, flow and stoppedflow techniques, and flash photolysis.\\
$\square$ 2. The instantaneous rate of consumption of a reactant or formation of a product is the slope of the tangent to the graph of concentration against time (expressed as a positive quantity).3. The rate of reaction is defined in terms of the extent of reaction in such a way that it is independent of the species being considered.4. A rate law is an expression for the reaction rate in terms of the concentrations of the species that occur in the overall chemical reaction.
  \item The order of a reaction is the power to which a participant is raised in the rate law; the overall order is the sum of these powers.
\end{enumerate}

\subsection*{Checklist of equations}
\begin{center}
\begin{tabular}{llll}
\hline
Property & Equation & Comment & Equation number \\
\hline
Rate of a reaction & $v=(1 / V)(\mathrm{d} \xi / \mathrm{d} t)$ & Definition & 17 A .2 \\
 & $v=\left(1 / v_{\mathrm{J}}\right)(\mathrm{d}[\mathrm{J}] / \mathrm{d} t)$ & Constant-volume system & 17 A .3 b \\
Rate law (in some cases) & $v=k_{\mathrm{r}}[\mathrm{A}]^{a}[\mathrm{~B}]^{b} \cdots$ & $a, b, \ldots:$ orders; $a+b+\cdots:$ overall order & 17 A .7 \\
Method of initial rates & $\log v_{0}=\log k_{\mathrm{r}, \text { eff }}+a \log [\mathrm{~A}]_{0}$ & Reactant A isolated & 17 A .11 c \\
\hline
\end{tabular}
\end{center}

\section*{TOPIC 17B Integrated rate laws}
\subsection*{> Why do you need to know this material?}
You need the integrated rate law if you want to predict the composition of a reaction mixture as it approaches equilibrium. The integrated rate law is also the basis of determining the order and rate constants of a reaction, which is a necessary step in the formulation of the mechanism of the reaction.

\subsection*{- What is the key idea?}
A rate law is a differential equation that can be integrated to find how the concentrations of reactants and products change with time.

\subsection*{> What do you need to know already?}
You need to be familiar with the concepts of rate law, reaction order, and rate constant (Topic 17A). The manipulation of simple rate laws requires only elementary techniques of integration (see the Resource section for standard integrals).

Rate laws (Topic 17A) are differential equations, which can be integrated to predict how the concentrations of the reactants and products change with time. Even the most complex rate laws may be integrated numerically. However, in a number of simple cases analytical solutions, known as integrated rate laws, are easily obtained and prove to be very useful.

\subsection*{17B.1 Zeroth-order reactions}
The rate of a zeroth-order reaction of the type $\mathrm{A} \rightarrow \mathrm{P}$ is constant (so long as reactant remains), so

$$
\frac{\mathrm{d}[\mathrm{~A}]}{\mathrm{d} t}=-k_{\mathrm{r}}
$$

It follows that the change in concentration of A is simply its rate of consumption (which is $-k_{\mathrm{r}}$ ) multiplied by the time $t$ for which the reaction has been in progress:

$$
[\mathrm{A}]-[\mathrm{A}]_{0}=-k_{\mathrm{r}} t
$$

\begin{center}
\includegraphics[max width=\textwidth]{2025_02_27_d4b95d9718f0b9e2e6b9g-763}
\end{center}

Figure 17B. 1 The linear decay of the reactant in a zeroth-order reaction.\\
where $[\mathrm{A}]$ is the concentration of A at $t$ and $[\mathrm{A}]_{0}$ is the initial concentration of A . This expression rearranges to


\begin{equation*}
[\mathrm{A}]=[\mathrm{A}]_{0}-k_{\mathrm{r}} t \quad \text { Integrated zeroth-order rate law } \tag{17B.1}
\end{equation*}


This expression applies until all the reactant has been used up at $t=[\mathrm{A}]_{0} / k_{\mathrm{r}}$, after which [A] remains zero (Fig. 17B.1).

\subsection*{17B.2 First-order reactions}
Consider the first-order rate law


\begin{equation*}
\frac{\mathrm{d}[\mathrm{~A}]}{\mathrm{d} t}=-k_{\mathrm{r}}[\mathrm{~A}] \tag{17B.2a}
\end{equation*}


This equation can be integrated to show how the concentration of A changes with time.

How is that done? 17B. 1\\
Deriving the first-order\\
integrated rate law\\
First, rearrange eqn 17B.2a into

$$
\frac{\mathrm{d}[\mathrm{~A}]}{[\mathrm{A}]}=-k_{\mathrm{r}} \mathrm{~d} t
$$

and recognize that $k_{\mathrm{r}}$ is a constant independent of $t$. Initially (at $t=0$ ) the concentration of A is $[\mathrm{A}]_{0}$, and at a later time $t$ it is\\[0pt]
[A], so make these values the matching limits of the integrals and write

$$
\int_{\left[A_{0}\right]_{0}}^{[\mathrm{A}]} \frac{\mathrm{d}[\mathrm{~A}]}{[\mathrm{A}]}=-k_{\mathrm{r}} \int_{0}^{t} \mathrm{~d} t
$$

Because the integral of $1 / x$ is $\ln x+$ constant, the left-hand side of this expression is

$$
\overbrace{\int_{[A]_{0}}^{[\mathrm{A}]} \frac{\mathrm{d}[\mathrm{~A}]}{\text { Integral A. } 2}}^{[\mathrm{A}]}=\ln [\mathrm{A}]+\text { constant }{ }_{[\mathrm{A}]_{0}}^{[\mathrm{A}]}=\ln [\mathrm{A}]-\ln [\mathrm{A}]_{0}=\ln \frac{[\mathrm{A}]}{[\mathrm{A}]_{0}}
$$

The right-hand side integrates to $-k_{\mathrm{r}} t$, and so

$$
-\ln \frac{[\mathrm{A}]}{[\mathrm{A}]_{0}}=-k_{\mathrm{r}} t, \quad[\mathrm{~A}]=[\mathrm{A}]_{0} \mathrm{e}^{-k_{t} t} \left\lvert\, \begin{gathered}
\text { Integrated first- } \\
\text { order rate law }
\end{gathered}\right.
$$

Equation 17 B .2 b shows that if $\ln \left([\mathrm{A}] /[\mathrm{A}]_{0}\right)$ is plotted against $t$, then a first-order reaction will give a straight line of slope $-k_{\mathrm{r}}$. Some rate constants determined in this way are given in Table 17B.1. The second expression in eqn 17B.2b shows that in a first-order reaction the reactant concentration decreases exponentially with time with a rate determined by $k_{\mathrm{r}}$ (Fig. 17B.2).

Table 17B. 1 Kinetic data for first-order reactions*

\begin{center}
\begin{tabular}{lllll}
\hline
Reaction & Phase & $\theta /{ }^{\circ} \mathrm{C}$ & $\boldsymbol{k}_{\mathrm{r}} / \mathrm{s}^{-1}$ & $\boldsymbol{t}_{1 / 2}$ \\
\hline
$2 \mathrm{~N}_{2} \mathrm{O}_{5} \rightarrow 4 \mathrm{NO}_{2}+\mathrm{O}_{2}$ & g & 25 & $3.38 \times 10^{-5}$ & 5.70 h \\
 & $\mathrm{Br}_{2}(\mathrm{l})$ & 25 & $4.27 \times 10^{-5}$ & 4.51 h \\
$\mathrm{C}_{2} \mathrm{H}_{6} \rightarrow 2 \mathrm{CH}_{3}$ & g & 700 & $5.36 \times 10^{-4}$ & 21.6 min \\
\hline
\end{tabular}
\end{center}

\begin{itemize}
  \item More values are given in the Resource section.\\
\includegraphics[max width=\textwidth, center]{2025_02_27_d4b95d9718f0b9e2e6b9g-764}
\end{itemize}

Figure 17B. 2 The exponential decay of the reactant in a firstorder reaction. The larger the rate constant, the more rapid is the decay: here $k_{\text {r,large }}=3 k_{\mathrm{r}, \text { small }}$.

The integrated rate law in eqn 17B.2b can be expressed in terms of the concentration of the product P by noting that for a reaction $\mathrm{A} \rightarrow \mathrm{P}$ the increase in the concentration of P matches the decrease in concentration of $A$. If it is assumed that there is no $P$ present at the start of the reaction, then $[\mathrm{P}]=[\mathrm{A}]_{0}-[\mathrm{A}]$, and hence $[\mathrm{A}]=[\mathrm{A}]_{0}-[\mathrm{P}]$. This expression for $[\mathrm{A}]$ can be substituted into eqn 17B. 2 b to give


\begin{equation*}
\ln \frac{[\mathrm{A}]_{0}-[\mathrm{P}]}{[\mathrm{A}]_{0}}=-k_{\mathrm{r}} t, \quad[\mathrm{P}]=[\mathrm{A}]_{0}\left(1-\mathrm{e}^{-k_{\mathrm{t}} t}\right) \tag{17B.2c}
\end{equation*}


A useful indication of the rate of a first-order chemical reaction is the half-life, $t_{1 / 2}$, of a substance, the time taken for the concentration of a reactant to fall to half its initial value. This quantity is readily obtained from the integrated rate law. Thus, the time for the concentration of $A$ to decrease from $[A]_{0}$ to $\frac{1}{2}[\mathrm{~A}]_{0}$ in a first-order reaction is given by eqn 17B. 2 b as

$$
k_{\mathrm{r}} t_{1 / 2}=-\ln \frac{\frac{1}{2}[\mathrm{~A}]_{0}}{[\mathrm{~A}]_{0}}=-\ln \frac{1}{2}=\ln 2
$$

Hence

\[
t_{1 / 2}=\frac{\ln 2}{k_{\mathrm{r}}} \quad \begin{align*}
& \text { Half-life }  \tag{17B.3}\\
& \text { [first-order reaction] }
\end{align*}
\]

(Note that $\ln 2=0.693$.) The main point to note about this result is that for a first-order reaction, the half-life of a reactant is independent of its initial concentration. Therefore, if the concentration of A at some arbitrary stage of the reaction is [A], then it will have fallen to $\frac{1}{2}[\mathrm{~A}]$ after a further interval of $(\ln 2) / k_{r}$. Some half-lives are given in Table 17B.1.

\subsection*{Example 17B. 1 Analysing a first-order reaction}
The variation in the partial pressure of azomethane with time was followed at 600 K , with the results given below. Confirm that the decomposition $\mathrm{CH}_{3} \mathrm{~N}_{2} \mathrm{CH}_{3}(\mathrm{~g}) \rightarrow \mathrm{CH}_{3} \mathrm{CH}_{3}(\mathrm{~g})+\mathrm{N}_{2}(\mathrm{~g})$ is first-order in azomethane, and find the rate constant and half-life at 600 K .

\begin{center}
\begin{tabular}{lcrrrr}
$t / \mathrm{s}$ & 0 & 1000 & \multicolumn{1}{c}{2000} & \multicolumn{1}{c}{3000} & \multicolumn{1}{c}{4000} \\
$p / \mathrm{Pa}$ & 10.9 & 7.63 & 5.32 & 3.71 & 2.59 \\
\end{tabular}
\end{center}

Collect your thoughts To confirm that a reaction is first order, plot $\ln \left([\mathrm{A}] /[\mathrm{A}]_{0}\right)$ against time and expect a straight line. Because the partial pressure of a gas is proportional to its concentration, an equivalent procedure is to plot $\ln \left(p / p_{0}\right)$ against $t$. If a straight line is obtained, its slope can be identified with $-k_{\mathrm{r}}$. The half-life is then calculated from $k_{\mathrm{r}}$ by using eqn 17B.3.

The solution Draw up the following table by using $p_{0}=10.9 \mathrm{~Pa}$ :

\begin{center}
\begin{tabular}{llrrrc}
$t / \mathrm{s}$ & 0 & \multicolumn{1}{l}{1000} & \multicolumn{1}{l}{2000} & \multicolumn{1}{l}{3000} & \multicolumn{1}{l}{4000} \\
$p / p_{0}$ & 1 & 0.700 & 0.488 & 0.340 & 0.0238 \\
$\ln \left(p / p_{0}\right)$ & 0 & -0.357 & -0.717 & -1.078 & -1.437 \\
\end{tabular}
\end{center}

Figure 17B. 3 shows the plot of $\ln \left(p / p_{0}\right)$ against $t / \mathrm{s}$. The plot is straight, confirming a first-order reaction, and its slope is $-3.6 \times 10^{-4}$. Therefore, $k_{\mathrm{r}}=3.6 \times 10^{-4} \mathrm{~s}^{-1}$. It follows from eqn 17B. 3 that the half-life is

$$
t_{1 / 2}=\frac{\ln 2}{3.6 \times 10^{-4} \mathrm{~s}^{-1}}=1.9 \times 10^{3} \mathrm{~s}
$$

\begin{center}
\includegraphics[max width=\textwidth]{2025_02_27_d4b95d9718f0b9e2e6b9g-765}
\end{center}

Figure 17B. 3 The determination of the rate constant of a first-order reaction: a straight line is obtained when $\ln [A]$ (or, as here, $\ln p / p_{0}$ ) is plotted against $t$; the slope is $-k_{r}$. The data plotted are from Example 17B.

Self-test 17B. 1 In a particular experiment, it was found that the concentration of $\mathrm{N}_{2} \mathrm{O}_{5}$ in liquid bromine varied with time as follows:

\begin{center}
\begin{tabular}{llcccc}
$t / \mathrm{s}$ & 0 & 200 & 400 & \multicolumn{1}{l}{600} & \multicolumn{1}{l}{1000} \\
$\left[\mathrm{~N}_{2} \mathrm{O}_{5}\right] /\left(\mathrm{mol} \mathrm{dm}^{-3}\right)$ & 0.110 & 0.073 & 0.048 & 0.032 & 0.014 \\
\end{tabular}
\end{center}

Confirm that the reaction is first order in $\mathrm{N}_{2} \mathrm{O}_{5}$ and determine the rate constant.

$$
{ }_{I_{-}^{-}} s_{\varepsilon-} 0 I \times I^{\prime} Z={ }^{1} y: \iota \partial s u_{V}
$$

\subsection*{17B. 3 Second-order reactions}
The integrated form of the second-order rate law,


\begin{equation*}
\frac{\mathrm{d}[\mathrm{~A}]}{\mathrm{d} t}=-k_{\mathrm{r}}[\mathrm{~A}]^{2} \tag{17B.4a}
\end{equation*}


can be found by much the same method as for first-order reactions.

\subsection*{How is that done? 17B. 2}
Deriving a second-order\\
integrated rate law\\
To integrate eqn 17B.4a, first rearrange it into

$$
\frac{\mathrm{d}[\mathrm{~A}]}{[\mathrm{A}]^{2}}=-k_{\mathrm{r}} \mathrm{~d} t
$$

The concentration is $[\mathrm{A}]_{0}$ at $t=0$ and at a later time $t$ it is $[\mathrm{A}]$. Therefore,

$$
-\overbrace{\int_{[A]_{0}}^{[A]} \frac{\mathrm{d}[\mathrm{~A}]}{[\mathrm{A}]^{2}}}^{\text {Integral A. } 1}=k_{\mathrm{r}} \int_{0}^{t} \mathrm{~d} t
$$

The integral on the left-hand side (including the minus sign) is

$$
\frac{1}{[\mathrm{~A}]}+\text { constant }\left.\right|_{[\mathrm{A}]_{0}} ^{[\mathrm{A}]}=\frac{1}{[\mathrm{~A}]}-\frac{1}{[\mathrm{~A}]_{0}}
$$

and that of the right-hand side is $k_{\mathrm{r}} t$. It follows that

$$
-\left|\frac{1}{[\mathrm{~A}]}-\frac{1}{[\mathrm{~A}]_{0}}=k_{\mathrm{r}} t, \quad[\mathrm{~A}]=\frac{[\mathrm{A}]_{0}}{1+k_{\mathrm{r}}[\mathrm{~A}]_{0}}\right| \begin{aligned}
& \text { Integrated second- } \\
& \text { order rate law }
\end{aligned}
$$

Equation 17B.4b shows that for a second-order reaction a plot of $1 /[\mathrm{A}]$ against $t$ is expected to be a straight line. The slope of the graph is $k_{\mathrm{r}}$. Some rate constants determined in this way are given in Table 17B.2. The alternative form of the equation can be used to predict the concentration of A at any time after the start of the reaction. It shows that the concentration of A approaches zero more slowly than in a first-order reaction with the same initial rate (Fig. 17B.4).

Table 17B. 2 Kinetic data for second-order reactions*

\begin{center}
\begin{tabular}{llll}
\hline
Reaction & Phase & $\theta /{ }^{\circ} \mathrm{C}$ & $\boldsymbol{k}_{\mathbf{r}} /\left(\mathrm{dm}^{3} \mathrm{~mol}^{-1} \mathrm{~s}^{-1}\right)$ \\
\hline
$2 \mathrm{NOBr} \rightarrow 2 \mathrm{NO}+\mathrm{Br}_{2}$ & g & 10 & 0.80 \\
$2 \mathrm{I} \rightarrow \mathrm{I}_{2}$ & g & 23 & $7 \times 10^{9}$ \\
\hline
\end{tabular}
\end{center}

\begin{itemize}
  \item More values are given in the Resource section.\\
\includegraphics[max width=\textwidth, center]{2025_02_27_d4b95d9718f0b9e2e6b9g-765(1)}
\end{itemize}

Figure 17B. 4 The variation with time of the concentration of a reactant in a second-order reaction. The grey lines are the corresponding decays in a first-order reaction with the same initial rate. For this illustration, $k_{\mathrm{r}, \text { arge }}=3 k_{\mathrm{r}, \text { small }}$.

As in the case of the first-order reactions, eqn 17B.4b can be rewritten in terms of the product P given the stoichiometry $\mathrm{A} \rightarrow \mathrm{P}$ by noting that $[\mathrm{A}]=[\mathrm{A}]_{0}-[\mathrm{P}]$. With this substitution, and after some rearrangement,


\begin{equation*}
\frac{[\mathrm{P}]}{\left([\mathrm{A}]_{0}-[\mathrm{P}]\right)[\mathrm{A}]_{0}}=k_{\mathrm{r}} t, \quad[\mathrm{P}]=\frac{k_{\mathrm{r}} \mathrm{t}[\mathrm{~A}]_{0}^{2}}{1+[\mathrm{A}]_{0} k_{\mathrm{r}} t} \tag{17B.4c}
\end{equation*}


It follows from eqn 17B.4b by substituting $t=t_{1 / 2}$ and $[\mathrm{A}]=\frac{1}{2}[\mathrm{~A}]_{0}$ that the half-life of a species A that is consumed in a second-order reaction is

\[
t_{1 / 2}=\frac{1}{k_{\mathrm{r}}[\mathrm{~A}]_{0}} \quad \begin{align*}
& \text { Half-life }  \tag{17B.5}\\
& {[\text { second-order reaction }]}
\end{align*}
\]

Therefore, unlike a first-order reaction, the half-life of a substance in a second-order reaction depends on the initial concentration. A practical consequence of this dependence is that species that decay by second-order reactions (which includes some environmentally harmful substances) may persist in low concentrations for long periods because their half-lives are long when their concentrations are low. In general, for an $n$ thorder reaction (with $n>1$ ) of the form $\mathrm{A} \rightarrow \mathrm{P}$, the half-life is related to the rate constant and the initial concentration of A by (see Problem P17B.15)

\[
t_{1 / 2}=\frac{2^{n-1}-1}{(n-1) k_{\mathrm{r}}[\mathrm{~A}]_{0}^{n-1}} \quad \begin{align*}
& \text { Half-life }  \tag{17B.6}\\
& {[n \text { th-order reaction, } n>1]}
\end{align*}
\]

Another type of second-order reaction is one that is first order in each of two reactants A and B:


\begin{equation*}
\frac{\mathrm{d}[\mathrm{~A}]}{\mathrm{d} t}=-k_{\mathrm{r}}[\mathrm{~A}][\mathrm{B}] \tag{17B.7a}
\end{equation*}


An example of a reaction that may have this rate law is $\mathrm{A}+$ $B \rightarrow P$. This rate law can be integrated to find the variation of the concentrations $[\mathrm{A}]$ and $[\mathrm{B}]$ with time.

\subsection*{How is that done? 17B. 3}
Deriving a second-order\\
integrated rate law for $A+B \rightarrow P$\\
Before integrating eqn 17B.7a, it is necessary to know how the concentration of B is related to that of A, which can be found from the reaction stoichiometry and the initial concentrations $[\mathrm{A}]_{0}$ and $[\mathrm{B}]_{0}$ which in this derivation are taken to be unequal. Follow these steps.

Step 1 Rewrite the rate law by considering the reaction stoichiometry\\
It follows from the reaction stoichiometry that when the concentration of A has fallen to $[\mathrm{A}]_{0}-x$, the concentration of B will have fallen to $[\mathrm{B}]_{0}-x$ (because each A that disap-\\
pears entails the disappearance of one B). Then eqn 17B.7a becomes

$$
\frac{\mathrm{d}[\mathrm{~A}]}{\mathrm{d} t}=-k_{\mathrm{r}}\left([\mathrm{~A}]_{0}-x\right)\left([\mathrm{B}]_{0}-x\right)
$$

Because $[\mathrm{A}]=[\mathrm{A}]_{0}-x$, it follows that $\mathrm{d}[\mathrm{A}] / \mathrm{d} t=-\mathrm{d} x / \mathrm{d} t$ and the rate law may be written

$$
\frac{\mathrm{d} x}{\mathrm{~d} t}=k_{\mathrm{r}}\left([\mathrm{~A}]_{0}-x\right)\left([\mathrm{B}]_{0}-x\right)
$$

\subsection*{Step 2 Integrate the rate law}
The initial condition is that $x=0$ when $t=0$; so the integrations required are

$$
\int_{0}^{x} \frac{\mathrm{~d} x}{\left([\mathrm{~A}]_{0}-x\right)\left([\mathrm{B}]_{0}-x\right)}=k_{\mathrm{r}} \int_{0}^{t} \mathrm{~d} t
$$

The right-hand side evaluates to $k_{\mathrm{r}} t$. The integral on the left is evaluated by using the method of partial fractions (see The chemist's toolkit 30 and the list of integrals in the Resource section):

$$
\overbrace{\int_{0}^{x} \frac{\mathrm{~d} x}{\left([\mathrm{~A}]_{0}-x\right)\left([\mathrm{B}]_{0}-x\right)}}^{\text {Integral A. } 3}=\frac{1}{[\mathrm{~B}]_{0}-[\mathrm{A}]_{0}}\left\{\ln \frac{[\mathrm{~A}]_{0}}{[\mathrm{~A}]_{0}-x}-\ln \frac{[\mathrm{B}]_{0}}{[\mathrm{~B}]_{0}-x}\right\}
$$

The two logarithms can be combined as follows:

$$
\begin{aligned}
\underbrace{\frac{[\mathrm{A}]_{0}}{[\mathrm{~A}]_{0}-x}}_{[\mathrm{A}]}-\ln \underbrace{\frac{[\mathrm{B}]]_{0}}{[\mathrm{~B}]_{0}-x}}_{[\mathrm{B}]} & =\ln \frac{[\mathrm{A}]_{0}}{[\mathrm{~A}]}-\ln \frac{[\mathrm{B}]_{0}}{[\mathrm{~B}]} \\
& =\ln \frac{1}{[\mathrm{~A}] /[\mathrm{A}]_{0}}-\ln \frac{1}{[\mathrm{~B}] /[\mathrm{B}]_{0}} \\
& =\ln \frac{[\mathrm{B}] /[\mathrm{B}]_{0}}{[\mathrm{~A}] /[\mathrm{A}]_{0}}
\end{aligned}
$$

Step 3 Finalize the expression\\
Combining all the results so far gives

$$
-\ln \frac{[\mathrm{B}] /[\mathrm{B}]_{0}}{[\mathrm{~A}] /[\mathrm{A}]_{0}}=\left([\mathrm{B}]_{0}-[\mathrm{A}]_{0}\right) k_{\mathrm{r}} t \left\lvert\, \begin{aligned}
& \text { Integrated rate law } \\
& \text { [second-order reaction of } \\
& \text { the type } \mathrm{A}+\mathrm{B} \rightarrow \mathrm{P}, \text { with } \\
& \left.[\mathrm{A}]_{0} \neq[\mathrm{B}]_{0}\right]
\end{aligned}\right.
$$

Therefore, a plot of the expression on the left against $t$ should be a straight line from which $k_{\mathrm{r}}$ can be obtained. As shown in the following Brief illustration, the rate constant may be estimated quickly by using data from only two measurements.

Similar calculations may be carried out to find the integrated rate laws for other orders, and some are listed in Table 17B.3.

\subsection*{The chemist's toolkit 30}
Integration by the method of partial fractions

To solve an integral of the form

$$
I=\int \frac{1}{(a-x)(b-x)} \mathrm{d} x
$$

where $a$ and $b$ are constants with $a \neq b$, use the method of partial fractions in which a fraction that is the product of terms (as in the denominator of this integrand) is written as a sum of fractions. To implement this procedure write the integrand as

$$
\frac{1}{(a-x)(b-x)}=\frac{1}{b-a}\left(\frac{1}{a-x}-\frac{1}{b-x}\right)
$$

Then integrate each term on the right. It follows that

$$
\begin{aligned}
I & =\frac{1}{b-a}(\overbrace{\int \frac{\mathrm{~d} x}{a-x}}^{\text {A.2 }}-\overbrace{\int \frac{\mathrm{d} x}{b-x}}^{\text {A.2 }}) \\
& =\frac{1}{b-a}\left(\ln \frac{1}{a-x}-\ln \frac{1}{b-x}\right)+\text { constant }
\end{aligned}
$$

\subsection*{Brief illustration 17B. 1}
Consider a second-order reaction of the type $\mathrm{A}+\mathrm{B} \rightarrow \mathrm{P}$ carried out in a solution. Initially, the concentrations of reactants are $[\mathrm{A}]_{0}=0.075 \mathrm{~mol} \mathrm{dm}^{-3}$ and $[\mathrm{B}]_{0}=0.050 \mathrm{~mol} \mathrm{dm}^{-3}$. After 1.0 h the concentration of $B$ has fallen to $[B]=0.020 \mathrm{~mol} \mathrm{dm}^{-3}$. Because the change in the concentration of $B$ is the same as that of A (and equal to $x$ ), it follows that during this time interval

$$
x=(0.050-0.020) \mathrm{mol} \mathrm{dm}^{-3}=0.030 \mathrm{~mol} \mathrm{dm}^{-3}
$$

Therefore, the concentration of A after 1.0 h is

$$
[\mathrm{A}]=[\mathrm{A}]_{0}-x=(0.075-0.030) \mathrm{mol} \mathrm{dm}^{-3}=0.045 \mathrm{~mol} \mathrm{dm}^{-3}
$$

and you are given that $[B]=0.020 \mathrm{moldm}^{-3}$. It follows from eqn 17B. 7 b that

$$
\begin{aligned}
k_{\mathrm{r}} & =\frac{1}{\left((0.050-0.075) \mathrm{mol} \mathrm{dm}^{-3}\right) \times(3600 \mathrm{~s})} \ln \frac{0.020 / 0.050}{0.045 / 0.075} \\
& =4.5 \times 10^{-3} \mathrm{dm}^{3} \mathrm{~mol}^{-1} \mathrm{~s}^{-1}
\end{aligned}
$$

Table 17B. 3 Integrated rate laws

\begin{center}
\begin{tabular}{|c|c|c|c|}
\hline
Order & Reaction & Rate law and its integrated form* & $t_{1 / 2}$ \\
\hline
\multirow[t]{3}{*}{0} & \multirow[t]{3}{*}{$\mathrm{A} \rightarrow \mathrm{P}$} & $v=k_{\mathrm{r}}$ & \multirow[t]{3}{*}{$[\mathrm{A}]_{0} / 2 k_{\mathrm{r}}$} \\
\hline
 &  & $k_{\mathrm{r}} t=[\mathrm{P}]$ for $0 \leq[\mathrm{P}] \leq[\mathrm{A}]_{0}$, &  \\
\hline
 &  & $[\mathrm{A}]=[\mathrm{A}]_{0}-k_{\mathrm{r}} \mathrm{t}$ for $0 \leq[\mathrm{A}] \leq[\mathrm{A}]_{0}$ &  \\
\hline
\multirow[t]{2}{*}{1} & \multirow[t]{2}{*}{$\mathrm{A} \rightarrow \mathrm{P}$} & \( v=k_{\mathrm{r}}[\mathrm{~A}] \) & \multirow[t]{2}{*}{$(\ln 2) / k_{\mathrm{r}}$} \\
\hline
 &  & \( k_{\mathrm{r}} t=\ln \frac{[\mathrm{A}]_{0}}{[\mathrm{~A}]}, \quad[\mathrm{A}]=[\mathrm{A}]_{0} \mathrm{e}^{-k_{\mathrm{r}} t}, \quad[\mathrm{P}]=[\mathrm{A}]_{0}\left(1-\mathrm{e}^{-k_{\mathrm{r}} t}\right) \) &  \\
\hline
\multirow[t]{7}{*}{2} & \multirow[t]{2}{*}{$\mathrm{A} \rightarrow \mathrm{P}$} & $v=k_{\mathrm{r}}[\mathrm{A}]^{2}$ & \multirow[t]{7}{*}{$1 / k_{\mathrm{r}}[\mathrm{A}]_{0}$} \\
\hline
 &  & \( k_{\mathrm{r}} t=\frac{[\mathrm{P}]}{[\mathrm{A}]_{0}\left([\mathrm{~A}]_{0}-[\mathrm{P}]\right)}, \quad[\mathrm{A}]=\frac{[\mathrm{A}]_{0}}{1+k_{\mathrm{r}} t[\mathrm{~A}]_{0}}, \quad[\mathrm{P}]=\frac{k_{\mathrm{r}} \mathrm{t}[\mathrm{~A}]_{0}^{2}}{1+[\mathrm{A}]_{0} k_{\mathrm{r}} t} \) &  \\
\hline
 & \multirow[t]{3}{*}{$\mathrm{A}+\mathrm{B} \rightarrow \mathrm{P}$} & $v=k_{\mathrm{r}}[\mathrm{A}][\mathrm{B}]$ &  \\
\hline
 &  & \( k_{\mathrm{r}} t=\frac{1}{[\mathrm{~B}]_{0}-[\mathrm{A}]_{0}} \ln \frac{[\mathrm{~A}]_{0}\left([\mathrm{~B}]_{0}-[\mathrm{P}]\right)}{\left([\mathrm{A}]_{0}-[\mathrm{P}]\right)[\mathrm{B}]_{0}}, \) &  \\
\hline
 &  & \( \ln \frac{[\mathrm{B}] /[\mathrm{B}]_{0}}{[\mathrm{~A}] /[\mathrm{A}]_{0}}=\left([\mathrm{B}]_{0}-[\mathrm{A}]_{0}\right) k_{\mathrm{r}} t, \quad[\mathrm{P}]=\frac{[\mathrm{A}]_{0}[\mathrm{~B}]_{0}\left(1-\mathrm{e}^{\left.(\mathrm{(B)}]_{0}-[\mathrm{A}]_{0}\right) k_{\mathrm{r}} t}\right)}{[\mathrm{A}]_{0}-[\mathrm{B}]_{0} \mathrm{e}^{\left.(\mathrm{(B)}]_{0}-[\mathrm{A}]_{0}\right) k_{\mathrm{r}} t}} \) &  \\
\hline
 & \multirow[t]{2}{*}{$\mathrm{A}+2 \mathrm{~B} \rightarrow \mathrm{P}$} & $v=k_{\mathrm{r}}[\mathrm{A}][\mathrm{B}]$ &  \\
\hline
 &  & \( k_{\mathrm{r}} t=\frac{1}{[\mathrm{~B}]_{0}-2[\mathrm{~A}]_{0}} \ln \frac{[\mathrm{~A}]_{0}\left([\mathrm{~B}]_{0}-2[\mathrm{P}]\right)}{\left([\mathrm{A}]_{0}-[\mathrm{P}]\right)[\mathrm{B}]_{0}}, \quad[\mathrm{P}]=\frac{[\mathrm{A}]_{0}[\mathrm{~B}]_{0}\left(1-\mathrm{e}^{\left(\left[\mathrm{B} \mathrm{~B}_{0}-2[\mathrm{~A}]_{0} k_{\mathrm{r}} t\right.\right.}\right)}{2[\mathrm{~A}]_{0}-[\mathrm{B}]_{0} \mathrm{e}^{\left.(\mathrm{B}]_{0}-2[\mathrm{~A}]_{0}\right) k_{\mathrm{r}} t}} \) &  \\
\hline
\multirow[t]{3}{*}{3} & \multirow[t]{3}{*}{$A+2 \mathrm{~B} \rightarrow \mathrm{P}$} & $v=k_{\mathrm{r}}[\mathrm{A}][\mathrm{B}]^{2}$ &  \\
\hline
 &  & \( k_{\mathrm{r}} t=\frac{2[\mathrm{P}]}{\left(2[\mathrm{~A}]_{0}-[\mathrm{B}]_{0}\right)\left([\mathrm{B}]_{0}-2[\mathrm{P}]\right)[\mathrm{B}]_{0}}+\frac{1}{\left(2[\mathrm{~A}]_{0}-[\mathrm{B}]_{0}\right)^{2}} \ln \frac{[\mathrm{~A}]_{0}\left([\mathrm{~B}]_{0}-2[\mathrm{P}]\right)}{\left([\mathrm{A}]_{0}-[\mathrm{P}]\right)[\mathrm{B}]_{0}} \) &  \\
\hline
 &  & [P] must be determined graphically or numerically &  \\
\hline
\multirow[t]{3}{*}{\(
n \geq 2
\)} & \multirow[t]{3}{*}{$\mathrm{A} \rightarrow \mathrm{P}$} & \( v=k_{\mathrm{r}}[\mathrm{~A}]^{n} \) & \multirow[t]{3}{*}{\(
\frac{2^{n}-1}{(n-1) k_{\mathrm{r}}[\mathrm{~A}]_{0}^{n-1}}
\)} \\
\hline
 &  & \( k_{\mathrm{r}} t=\frac{1}{n-1}\left\{\frac{1}{\left([\mathrm{~A}]_{0}-[\mathrm{P}]\right)^{n-1}}-\frac{1}{[\mathrm{~A}]_{0}^{n-1}}\right\} \) &  \\
\hline
 &  & No simple general solution for [P] for $n>3$ &  \\
\hline
\end{tabular}
\end{center}

\footnotetext{\begin{itemize}
  \item $v=\mathrm{d}[\mathrm{P}] / \mathrm{d} t$
\end{itemize}
}\subsection*{Checklist of concepts}
1. An integrated rate law is an expression for the concentration of a reactant or product as a function of time (Table 17B.3).2. The half-life of a reactant is the time it takes for its concentration to fall to half its initial value.$\square$ 3. Analysis of experimental data using integrated rate laws allow for the prediction of the composition of a reaction system at any stage, the verification of the rate law, and the determination of the rate constant.

\subsection*{Checklist of equations}
\begin{center}
\begin{tabular}{llll}
\hline
Property & Equation & Comment & Equation number \\
\hline
Integrated rate law & $[\mathrm{A}]=[\mathrm{A}]_{0}-k_{\mathrm{r}} t$ & Zeroth order, $\mathrm{A} \rightarrow \mathrm{P}$ & 17 B .1 \\
Integrated rate law & $\ln \left([\mathrm{A}] /[\mathrm{A}]_{0}\right)=-k_{\mathrm{r}} t$ or $[\mathrm{A}]=[\mathrm{A}]_{0} \mathrm{e}^{-k_{\mathrm{r}} t}$ & First order, $\mathrm{A} \rightarrow \mathrm{P}$ & 17 B .2 b \\
Half-life & $t_{1 / 2}=(\ln 2) / k_{\mathrm{r}}$ & First order, $\mathrm{A} \rightarrow \mathrm{P}$ & 17 B .3 \\
Integrated rate law & $1 /[\mathrm{A}]-1 /[\mathrm{A}]_{0}=k_{\mathrm{r}} t$ or $[\mathrm{A}]=[\mathrm{A}]_{0} /\left(1+k_{\mathrm{r}} t[\mathrm{~A}]_{0}\right)$ & Second order, $\mathrm{A} \rightarrow \mathrm{P}$ & 17 B .4 b \\
Half-life & $t_{1 / 2}=1 / k_{\mathrm{r}}[\mathrm{A}]_{0}$ & Second order, $\mathrm{A} \rightarrow \mathrm{P}$ & 17 B .5 \\
 & $t_{1 / 2}=\left(2^{n-1}-1\right) /(n-1) k_{\mathrm{r}}[\mathrm{A}]_{0}^{n-1}$ & $n$th order, $n>1$ & $17 \mathrm{~B} \cdot 6$ \\
Integrated rate law & $\ln \left\{\left([\mathrm{B}] /[\mathrm{B}]_{0}\right) /\left([\mathrm{A}] /[\mathrm{A}]_{0}\right)\right\}=\left([\mathrm{B}]_{0}-[\mathrm{A}]_{0}\right) k_{\mathrm{r}} t$ & Second order, $\mathrm{A}+\mathrm{B} \rightarrow \mathrm{P}$ & 17 B .7 b \\
\hline
\end{tabular}
\end{center}

\section*{TOPIC 17C Reactions approaching equilibrium }
\subsection*{> Why do you need to know this material?}
All reactions tend towards equilibrium and the rate laws can be used to describe the changing concentrations as they approach that composition. The analysis of the timedependence also reveals the connection between rate constants and equilibrium constants.

\subsection*{> What is the key idea?}
Both forward and reverse reactions must be incorporated into a reaction scheme in order to account for the approach to equilibrium.

\subsection*{> What do you need to know already?}
You need to be familiar with the concepts of rate law, reaction order, and rate constant (Topic 17A), integrated rate laws (Topic 17B), and equilibrium constants (Topic 6A). As in Topic 17B, the manipulation of simple rate laws requires only elementary techniques of integration.

In practice, most kinetic studies are made on reactions that are far from equilibrium and if products are in low concentration the reverse reactions are unimportant. Close to equilibrium, however, the products might be so abundant that the reverse reaction must be taken into account.

\subsection*{17C. 1 First-order reactions approaching equilibrium}
Considering a reaction in which A forms B and both forward and reverse reactions are first order (as in some isomerizations):

\[
\begin{array}{ll}
\mathrm{A} \rightarrow \mathrm{~B} & \frac{\mathrm{~d}[\mathrm{~A}]}{\mathrm{d} t}=-k_{\mathrm{r}}[\mathrm{~A}] \\
\mathrm{B} \rightarrow \mathrm{~A} & \frac{\mathrm{~d}[\mathrm{~A}]}{\mathrm{d} t}=k_{\mathrm{r}}^{\prime}[\mathrm{B}] \tag{17C.1}
\end{array}
\]

The concentration of $A$ is reduced by the forward reaction (at a rate $k_{\mathrm{r}}[\mathrm{A}]$ ) but it is increased by the reverse reaction (at a rate $\left.k_{\mathrm{r}}^{\prime}[\mathrm{B}]\right)$. The net rate of change at any stage is therefore


\begin{equation*}
\frac{\mathrm{d}[\mathrm{~A}]}{\mathrm{d} t}=-k_{\mathrm{r}}[\mathrm{~A}]+k_{\mathrm{r}}^{\prime}[\mathrm{B}] \tag{17C.2}
\end{equation*}


If the initial concentration of $A$ is $[A]_{0}$, and no $B$ is present initially, then at all times $[\mathrm{A}]+[\mathrm{B}]=[\mathrm{A}]_{0}$. Therefore,


\begin{align*}
\frac{\mathrm{d}[\mathrm{~A}]}{\mathrm{d} t} & =-k_{\mathrm{r}}[\mathrm{~A}]+k_{\mathrm{r}}^{\prime}\left([\mathrm{A}]_{0}-[\mathrm{A}]\right)  \tag{17C.3}\\
& =-\left(k_{\mathrm{r}}+k_{\mathrm{r}}^{\prime}\right)[\mathrm{A}]+k_{\mathrm{r}}^{\prime}[\mathrm{A}]_{0}
\end{align*}


The solution of this first-order differential equation (as may be checked by differentiation, Problem P17C.1) is


\begin{equation*}
[\mathrm{A}]=\frac{k_{\mathrm{r}}^{\prime}+k_{\mathrm{r}} \mathrm{e}^{-\left(k_{\mathrm{t}}+k_{\mathrm{r}}^{\prime}\right) t}}{k_{\mathrm{r}}+k_{\mathrm{r}}^{\prime}}[\mathrm{A}]_{0}, \quad[\mathrm{~B}]=[\mathrm{A}]_{0}-[\mathrm{A}] \tag{17C.4}
\end{equation*}


Figure 17C. 1 shows the time dependence predicted by this equation.

As $t \rightarrow \infty$, the exponential term in eqn 17C. 4 decreases to zero and the concentrations reach their equilibrium values, which are therefore


\begin{equation*}
[\mathrm{A}]_{\mathrm{eq}}=\frac{k_{\mathrm{r}}^{\prime}[\mathrm{A}]_{0}}{k_{\mathrm{r}}+k_{\mathrm{r}}^{\prime}}, \quad[\mathrm{B}]_{\mathrm{eq}}=[\mathrm{A}]_{0}-[\mathrm{A}]_{\mathrm{eq}}=\frac{k_{\mathrm{r}}[\mathrm{~A}]_{0}}{k_{\mathrm{r}}+k_{\mathrm{r}}^{\prime}} \tag{17C.5}
\end{equation*}


It follows that the equilibrium constant of the reaction is


\begin{equation*}
K=\frac{[\mathrm{B}]_{\mathrm{eq}}}{[\mathrm{~A}]_{\mathrm{eq}}}=\frac{k_{\mathrm{r}}}{k_{\mathrm{r}}^{\prime}} \tag{17C.6}
\end{equation*}


\begin{center}
\includegraphics[max width=\textwidth]{2025_02_27_d4b95d9718f0b9e2e6b9g-769}
\end{center}

Figure 17C. 1 The approach of concentrations to their equilibrium values as predicted by eqn 17C. 4 for a reaction $\mathrm{A} \rightleftharpoons \mathrm{B}$ that is first order in each direction, and for which $k_{\mathrm{r}}=2 k_{r}^{\prime}$.\\
(As explained in Topic 6A, the replacement of activities by the numerical values of molar concentrations is justified if the system is treated as ideal.) Exactly the same conclusion can be reached-more simply, in fact-by noting that, at equilibrium, the forward and reverse rates must be the same, so


\begin{equation*}
k_{\mathrm{r}}[\mathrm{~A}]_{\mathrm{eq}}=k_{\mathrm{r}}^{\prime}[\mathrm{B}]_{\mathrm{eq}} \tag{17C.7}
\end{equation*}


This relation rearranges into eqn 17C.6. The theoretical importance of eqn 17C. 6 is that it relates a thermodynamic quantity, the equilibrium constant, to quantities relating to rates. Its practical importance is that if one of the rate constants can be measured, then the other may be obtained if the equilibrium constant is known.

Equation 17C. 6 is valid even if the forward and reverse reactions have different orders, but in that case more care needs to be taken with units. For instance, if the reaction A + $\mathrm{B} \rightarrow \mathrm{C}$ is second-order in the forward direction and first-order in the reverse direction, then the condition for equilibrium is $k_{\mathrm{r}}[\mathrm{A}]_{\mathrm{eq}}[\mathrm{B}]_{\mathrm{eq}}=k_{\mathrm{r}}^{\prime}[\mathrm{C}]_{\mathrm{eq}}$ and the dimensionless equilibrium constant in full dress is

$$
K=\frac{[\mathrm{C}]_{\mathrm{eq}} / c^{\ominus}}{\left([\mathrm{A}]_{\mathrm{eq}} / c^{\ominus}\right)\left([\mathrm{B}]_{\mathrm{eq}} / c^{\ominus}\right)}=\left(\frac{[\mathrm{C}]}{[\mathrm{A}][\mathrm{B}]}\right)_{\mathrm{eq}} c^{\ominus}=\frac{k_{r}}{k_{r}^{\prime}} \times c^{\ominus}
$$

The presence of $c^{\ominus}=1 \mathrm{~mol} \mathrm{dm}^{-3}$ in the last term ensures that the ratio of second-order to first-order rate constants, with their different units, is turned into a dimensionless quantity.

\subsection*{Brief illustration 17C. 1}
The rate constants of the forward and reverse reactions for a dimerization reaction were found to be $8.0 \times 10^{8} \mathrm{dm}^{3} \mathrm{~mol}^{-1} \mathrm{~s}^{-1}$ (second order) and $2.0 \times 10^{6} \mathrm{~s}^{-1}$ (first order). The equilibrium constant for the dimerization is therefore

$$
K=\frac{8.0 \times 10^{8} \mathrm{dm}^{3} \mathrm{~mol}^{-1} \mathrm{~s}^{-1}}{2.0 \times 10^{6} \mathrm{~s}^{-1}} \times 1 \mathrm{moldm}^{-3}=4.0 \times 10^{2}
$$

For a more general reaction, the overall equilibrium constant can be expressed in terms of the rate constants for all the intermediate stages of the reaction mechanism (see Problem P17C.4):

\[
K=\frac{k_{\mathrm{a}}}{k_{\mathrm{a}}^{\prime}} \times \frac{k_{\mathrm{b}}}{k_{\mathrm{b}}^{\prime}} \times \cdots \quad \begin{align*}
& \text { The equilibrium constant in }  \tag{17C.8}\\
& \text { terms of the rate constants }
\end{align*}
\]

where the $k_{\mathrm{r}}$ are the rate constants for the individual steps and the $k_{\mathrm{r}}^{\prime}$ are those for the corresponding reverse steps. The appropriate powers of $c^{\ominus}$ should be included in each factor if the orders of the forward and reverse reactions are different.\\
\includegraphics[max width=\textwidth, center]{2025_02_27_d4b95d9718f0b9e2e6b9g-770}

Figure 17C. 2 The relaxation to the new equilibrium composition when a reaction initially at equilibrium at a temperature $T_{1}$ is subjected to a sudden change of temperature, which takes it to $T_{2}$.

\subsection*{17C. 2 Relaxation methods}
The term relaxation denotes the return of a system to equilibrium. It is used in chemical kinetics to indicate that an externally applied influence has shifted the equilibrium position of a reaction, often suddenly, and that the concentrations of the species involved then adjust towards the equilibrium values characteristic of the new conditions (Fig. 17C.2).

Consider the response of reaction rates to a temperature jump, a sudden change in temperature. As explained in Topic 6B, provided $\Delta_{\mathrm{r}} H^{\ominus}$ is non-zero the equilibrium composition of a reaction depends on the temperature, so a sudden shift in temperature acts as a perturbation on the system. One way of achieving a temperature jump is to subject the sample that has been made conducting by the addition of ions to an electric discharge; bursts of microwave radiation or intense electromagnetic pulses from lasers can also be used. Electrical discharges can achieve temperature jumps of between 5 and 10 K in about $1 \mu \mathrm{~s}$. The high energy output of a pulsed laser is sufficient to generate temperature jumps of between 10 and 30 K within nanoseconds in aqueous samples.

The response of a system to a sudden temperature increase can be analysed by considering the rate laws for the forward and reverse reactions and the temperature dependence of the rate constants.

\subsection*{How is that done? 17C. 1 Exploring the response to a temperature jump}
First consider a simple $A \rightleftharpoons B$ equilibrium that is first order in each direction and then an equilibrium $A \rightleftharpoons B+C$ that is first order forward and second order reverse. In each case, when the temperature is increased suddenly, the rate constants change from their original values to the new values $k_{\mathrm{r}}$ and $k_{\mathrm{r}}^{\prime}$ characteristic of the new temperature, but the concentrations of A and $B$ remain for an instant at their old equilibrium values.\\
(a) The equilibrium $A \rightleftharpoons B$ (first order forward and reverse)

As the system immediately after the temperature jump is no longer at equilibrium, it readjusts to the new equilibrium concentrations, which are now given by $k_{\mathrm{r}}[\mathrm{A}]_{\mathrm{eq}}=k_{\mathrm{r}}^{\prime}[\mathrm{B}]_{\mathrm{eq}}$, and it does so at a rate that depends on the new rate constants. Let the deviation of $[\mathrm{A}]$ from its new equilibrium value be $x$, so $[\mathrm{A}]=[\mathrm{A}]_{\mathrm{eq}}+x$; then the reaction stoichiometry implies that $[\mathrm{B}]=[\mathrm{B}]_{\text {eq }}-x$. At the new temperature the concentration of A changes as follows:

$$
\begin{aligned}
\frac{\mathrm{d}[\mathrm{~A}]}{\mathrm{d} t} & =-k_{\mathrm{r}}[\mathrm{~A}]+k_{\mathrm{r}}^{\prime}[\mathrm{B}] \\
& =-k_{\mathrm{r}}\left([\mathrm{~A}]_{\mathrm{eq}}+x\right)+k_{\mathrm{r}}^{\prime}\left([\mathrm{B}]_{\mathrm{eq}}-x\right) \\
& =\overbrace{k_{\mathrm{r}}[\mathrm{~A}]_{\mathrm{eq}}+k_{\mathrm{r}}^{\prime}[\mathrm{B}]_{\mathrm{eq}}}^{\text {Cancel }}-\left(k_{\mathrm{r}}+k_{\mathrm{r}}^{\prime}\right) x \\
& =-\left(k_{\mathrm{r}}+k_{\mathrm{r}}^{\prime}\right) x
\end{aligned}
$$

Because $\mathrm{d}[\mathrm{A}] / \mathrm{d} t=\mathrm{d} x / \mathrm{d} t$, this equation is a first-order differential equation with a solution that resembles eqn 17B.2b. If $x_{0}$ is the deviation from equilibrium immediately after the temperature jump, the time-dependence of $x$ is

$$
\left.\longrightarrow x=x_{0} \mathrm{e}^{-t / \tau} \tau=\frac{1}{k_{\mathrm{r}}+k_{\mathrm{r}}^{\prime}} \right\rvert\, \begin{aligned}
& \text { Relaxation after a } \\
& \text { temperature jump } \\
& \text { [first-order reactions] }
\end{aligned}
$$

(b) The equilibrium $A \rightleftharpoons B+C$ (first order forward and second order reverse)

As in the previous derivation, the system immediately after the temperature jump is no longer at equilibrium, so it readjusts to the new equilibrium concentrations, which are now given by $k_{\mathrm{r}}[\mathrm{A}]_{\text {eq }}=k_{\mathrm{r}}^{\prime}[\mathrm{B}]_{\mathrm{eq}}[\mathrm{C}]_{\mathrm{eq}}$, and it does so at a rate that depends on the new rate constants. Let the deviation of $[\mathrm{A}]$ from its new equilibrium value be $x$, so $[\mathrm{A}]=[\mathrm{A}]_{\text {eq }}+x$; then the reaction stoichiometry implies that $[\mathrm{B}]=[\mathrm{B}]_{\mathrm{eq}}-x$ and $[C]=[C]_{\mathrm{eq}}-x$.

Step 1 Set up and solve the rate equations\\
At the new temperature the concentration of A changes as follows:

$$
\begin{aligned}
\frac{\mathrm{d}[\mathrm{~A}]}{\mathrm{d} t} & =-k_{\mathrm{r}}[\mathrm{~A}]+k_{\mathrm{r}}^{\prime}[\mathrm{B}][\mathrm{C}] \\
& =-k_{\mathrm{r}}\left([\mathrm{~A}]_{\mathrm{eq}}+x\right)+k_{\mathrm{r}}^{\prime}\left([\mathrm{B}]_{\mathrm{eq}}-x\right)\left([\mathrm{C}]_{\mathrm{eq}}-x\right) \\
& =-\left\{k_{\mathrm{r}}+k_{\mathrm{r}}^{\prime}\left([\mathrm{B}]_{\mathrm{eq}}+[\mathrm{C}]_{\mathrm{eq}}\right)\right\} x \overbrace{-k_{\mathrm{r}}[\mathrm{~A}]_{\mathrm{eq}}+k_{\mathrm{r}}^{\prime}[\mathrm{B}]_{\mathrm{eq}}[\mathrm{C}]_{\mathrm{eq}}}^{0}+\overbrace{k_{\mathrm{r}}^{\prime} x^{2}}^{\text {Neglect }} \\
& =-\left\{k_{\mathrm{r}}+k_{\mathrm{r}}^{\prime}\left([\mathrm{B}]_{\mathrm{eq}}+[\mathrm{C}]_{\mathrm{eq}}\right)\right\} x
\end{aligned}
$$

As before, $\mathrm{d}[\mathrm{A}] / \mathrm{d} t=\mathrm{d} x / \mathrm{d} t$; the solution of this differential equation is an exponential decay proportional to $\mathrm{e}^{-t / \tau}$ with $\tau$ given by

$$
\frac{1}{\tau}=k_{\mathrm{r}}+k_{\mathrm{r}}^{\prime}\left([\mathrm{B}]_{\mathrm{eq}}+[\mathrm{C}]_{\mathrm{eq}}\right)
$$

Step 2 Relate the equilibrium concentrations by introducing an equilibrium constant\\
The equilibrium constant for the reaction (assuming ideal solutions) is

$$
K=\frac{\left([\mathrm{B}]_{\mathrm{eq}} / c^{\ominus}\right)\left([\mathrm{C}]_{\mathrm{eq}} / c^{\ominus}\right)}{\left([\mathrm{A}]_{\mathrm{eq}} / c^{\ominus}\right)}=\frac{[\mathrm{B}]_{\mathrm{eq}}[\mathrm{C}]_{\mathrm{eq}}}{[\mathrm{~A}]_{\mathrm{eq}} c^{\ominus}}
$$

The reaction stoichiometry implies that the concentrations of $B$ and C are the same, so

$$
[\mathrm{B}]_{\mathrm{eq}}=[\mathrm{C}]_{\mathrm{eq}}=\overbrace{\left(K[\mathrm{~A}]_{\mathrm{eq}} c^{\ominus}\right)^{1 / 2}}^{a}
$$

and the time constant becomes

$$
\frac{1}{\tau}=k_{\mathrm{r}}+2 a k_{\mathrm{r}}^{\prime}=k_{\mathrm{r}}^{\prime}\left(\frac{k_{\mathrm{r}}}{k_{\mathrm{r}}^{\prime}}+2 a\right)
$$

Step 3 Identify the equilibrium constant of the reaction\\
You should now recognize that the ratio of the rate constants is a form of the equilibrium constant at the new temperature. Specifically:

$$
k_{\mathrm{r}}[\mathrm{~A}]_{\mathrm{eq}}=k_{\mathrm{r}}^{\prime}[\mathrm{B}]_{\mathrm{eq}}[\mathrm{C}]_{\mathrm{eq}} \text { and } K=\frac{[\mathrm{B}]_{\mathrm{eq}}[\mathrm{C}]_{\mathrm{eq}}}{[\mathrm{~A}]_{\mathrm{eq}} c^{\ominus}}
$$

implying that

$$
\frac{k_{\mathrm{r}}}{k_{\mathrm{r}}^{\prime}}=\frac{[\mathrm{B}]_{\mathrm{eq}}[\mathrm{C}]_{\mathrm{eq}}}{[\mathrm{~A}]_{\mathrm{eq}}}=K c^{\ominus}
$$

and therefore

$$
\frac{1}{\tau}=k_{\mathrm{r}}^{\prime}\left(K c^{\ominus}+2 a\right)
$$

The time dependence of $x$ is therefore

$-$\begin{tabular}{ll|l}
$x=x_{0} \mathrm{e}^{-t / \tau}$ & $\tau=\frac{1}{k_{\mathrm{r}}^{\prime}\left(K c^{\ominus}+2 a\right)}$ & (17C.9b) \\
 & $=\left(K[\mathrm{~A}]_{\mathrm{eq}} c^{\ominus}\right)^{1 / 2}$ &  \\
\end{tabular}\$\textbackslash quad \textbackslash begin\{aligned\} \& \text{Relaxation after a} \\


\& \text{temperature jump} \\
\\
\& \text{[mixed-order reaction]}\textbackslash end\{aligned\}\$

\subsection*{Checklist of concepts}
\begin{enumerate}
  \item There is a relation between the equilibrium constant, a thermodynamic quantity, and the rate constants of the forward and reverse reactions (see below).\\
$\square$ 2. In relaxation methods of kinetic analysis, the equilibrium position of a reaction is shifted suddenly and then the time-dependence of the concentration of the species involved is followed.
\end{enumerate}

\subsection*{Checklist of equations}
\begin{center}
\begin{tabular}{llll}
\hline
Property & Equation & Comment & Equation number \\
\hline
Equilibrium constant in terms of rate constants & $K=k_{\mathrm{a}} / k_{\mathrm{a}}^{\prime} \times k_{\mathrm{b}} / k_{\mathrm{b}}^{\prime} \times \cdots$ & Include $c^{\ominus}$ as appropriate & 17 C .8 \\
Relaxation of an equilibrium $\mathrm{A} \rightleftharpoons \mathrm{B}$ after a temperature jump & $x=x_{0} \mathrm{e}^{-t / \tau}$ & First order in each direction & 17C.9a \\
 & $\tau=1 /\left(k_{\mathrm{r}}+k_{\mathrm{r}}^{\prime}\right)$ &  &  \\
\hline
\end{tabular}
\end{center}

\section*{TOPIC 17D The Arrhenius equation}
\subsection*{> Why do you need to know this material?}
Exploration of the dependence of reaction rates on temperature leads to the formulation of theories that reveal the details of the processes that occur when reactant molecules meet and undergo reaction.

\subsection*{$>$ What is the key idea?}
The temperature dependence of the rate of a reaction depends on the activation energy, the minimum energy needed for reaction to occur in an encounter between reactants.

\begin{displayquote}
What do you need to know already?\\
You need to know that the rate of a chemical reaction is expressed by a rate constant (Topic 17A).
\end{displayquote}

Chemical reactions usually go faster as the temperature is raised. It is found experimentally for many reactions that a plot of $\ln k_{\mathrm{r}}$ against $1 / T$ gives a straight line with a negative slope, indicating that an increase in $\ln k_{\mathrm{r}}$ (and therefore an increase in $k_{\mathrm{r}}$ ) results from a decrease in $1 / T$ (i.e. an increase in $T$ ).

\subsection*{17D. 1 The temperature dependence of reaction rates}
The temperature dependence characteristic of a reaction is normally expressed mathematically by introducing two parameters, one representing the intercept and the other the slope of the straight line of a so-called 'Arrhenius plot' of $\ln k_{\mathrm{r}}$ against $1 / T$ and writing the Arrhenius equation


\begin{equation*}
\ln k_{\mathrm{r}}=\ln A-\frac{E_{\mathrm{a}}}{R T} \quad \text { Arrhenius equation } \tag{17D.1}
\end{equation*}


The parameter $A$, which is obtained from the intercept of the line at $1 / T=0$ (at infinite temperature, Fig. 17D.1), is called the frequency factor (and still commonly the pre-exponential factor). The parameter $E_{\mathrm{a}}$, which is obtained from the slope of the line (which is equal to $-E_{a} / R$ ), is called the activation energy. Collectively the two quantities are called the Arrhenius parameters (Table 17D.1).\\
\includegraphics[max width=\textwidth, center]{2025_02_27_d4b95d9718f0b9e2e6b9g-773}

Figure 17D. 1 An Arrhenius plot, a plot of $\ln k_{r}$ against $1 / T$, is a straight line when the reaction follows the behaviour described by the Arrhenius equation (eqn 17D.1). The slope is $-E_{a} / R$ and the intercept at $1 / T=0$ is $\ln A$.

Table 17D. 1 Arrhenius parameters*

\begin{center}
\begin{tabular}{llll}
\hline
(1) First-order reactions & Phase & $A / \mathrm{s}^{-1}$ & $E_{\mathrm{a}} /\left(\mathrm{kJ} \mathrm{mol}^{-1}\right)$ \\
\hline
$\mathrm{CH}_{3} \mathrm{NC} \rightarrow \mathrm{CH}_{3} \mathrm{CN}$ & gas & $3.98 \times 10^{13}$ & 160 \\
$2 \mathrm{~N}_{2} \mathrm{O}_{5} \rightarrow 4 \mathrm{NO}_{2}+\mathrm{O}_{2}$ & gas & $4.94 \times 10^{13}$ & 103.4 \\
\hline
$(2)$ Second-order reactions & Phase & $A /\left(\mathrm{dm}^{3} \mathrm{~mol}^{-1} \mathrm{~s}^{-1}\right)$ & $E_{\mathrm{a}} /\left(\mathrm{kJ} \mathrm{mol}^{-1}\right)$ \\
\hline
$\mathrm{OH}+\mathrm{H}_{2} \rightarrow \mathrm{H}_{2} \mathrm{O}+\mathrm{H}$ & gas & $8.0 \times 10^{10}$ & 42 \\
$\mathrm{NaC}_{2} \mathrm{H}_{5} \mathrm{O}+\mathrm{CH}_{3} \mathrm{I}$ & in ethanol & $2.42 \times 10^{11}$ & 81.6 \\
\hline
\end{tabular}
\end{center}

\begin{itemize}
  \item More values are given in the Resource section.
\end{itemize}

\subsection*{Example 17D. 1}
\subsection*{Determining the Arrhenius parameters}
The rate of the second-order decomposition of ethanal (acetaldehyde, $\mathrm{CH}_{3} \mathrm{CHO}$ ) was measured over the temperature range $700-1000 \mathrm{~K}$, and the rate constants are reported below. Find $E_{\mathrm{a}}$ and $A$.

\begin{center}
\begin{tabular}{lcccccccc}
$T / \mathrm{K}$ & 700 & 730 & 760 & 790 & 810 & 840 & 910 & 1000 \\
$k_{\mathrm{r}} /$ &  &  &  &  &  &  &  &  \\
$\left(\mathrm{dm}^{3} \mathrm{~mol}^{-1} \mathrm{~s}^{-1}\right)$ & 0.011 & 0.035 & 0.105 & 0.343 & 0.789 & 2.17 & 20.0 & 145 \\
\end{tabular}
\end{center}

Collect your thoughts According to eqn 17D.1, the data can be analysed by plotting $\ln \left(k_{\mathrm{r}} / \mathrm{dm}^{3} \mathrm{~mol}^{-1} \mathrm{~s}^{-1}\right)$ against $1 /(T / \mathrm{K})$, or more conveniently $\left(10^{3} \mathrm{~K}\right) / T$, expecting to get a straight line. The activation energy is found from the dimensionless slope by writing $-E_{\mathrm{a}} / R=$ slope/units, where in this case 'units' $=$ $1 /\left(10^{3} \mathrm{~K}\right)$, so $E_{\mathrm{a}}=-$ slope $\times R \times 10^{3} \mathrm{~K}$. The intercept at $1 / T=0$\\
is $\ln \left(A / \mathrm{dm}^{3} \mathrm{~mol}^{-1} \mathrm{~s}^{-1}\right)$. Use a least-squares procedure to determine the slope and the intercept.

The solution Draw up the following table:

\begin{center}
\begin{tabular}{lrrrrrrrr}
$\left(10^{3} \mathrm{~K}\right) / T$ & 1.43 & 1.37 & 1.32 & 1.27 & 1.23 & 1.19 & 1.10 & 1.00 \\
$\ln \left(k_{\mathrm{r}} / \mathrm{dm}^{3} \mathrm{~mol}^{-1} \mathrm{~s}^{-1}\right)$ & -4.51 & -3.35 & -2.25 & -1.07 & -0.24 & 0.77 & 3.00 & 4.98 \\
\end{tabular}
\end{center}

Now plot $\ln k_{\mathrm{r}}$ against ( $10^{3} \mathrm{~K}$ )/T (Fig. 17D.2). The least-squares fit results in a line with slope -22.7 and intercept 27.7. Therefore,

$$
\begin{aligned}
& E_{\mathrm{a}}=-(-22.7) \times\left(8.3145 \mathrm{~J} \mathrm{~K}^{-1} \mathrm{~mol}^{-1}\right) \times\left(10^{3} \mathrm{~K}\right)=189 \mathrm{~kJ} \mathrm{~mol}^{-1} \\
& A=\mathrm{e}^{27.7} \mathrm{dm}^{3} \mathrm{~mol}^{-1} \mathrm{~s}^{-1}=1.1 \times 10^{12} \mathrm{dm}^{3} \mathrm{~mol}^{-1} \mathrm{~s}^{-1}
\end{aligned}
$$

Note that $A$ has the same units as $k_{r}$.\\
\includegraphics[max width=\textwidth, center]{2025_02_27_d4b95d9718f0b9e2e6b9g-774(1)}

Figure 17D. 2 The Arrhenius plot using the data in Example 17D.1.

Self-test 17D. 1 Determine $A$ and $E_{\mathrm{a}}$ from the following data:

\begin{center}
\begin{tabular}{llllll}
$T / \mathrm{K}$ & 300 & 350 & 400 & 450 & 500 \\
$k_{\mathrm{r}} /\left(\mathrm{dm}^{3} \mathrm{~mol}^{-1} \mathrm{~s}^{-1}\right)$ & $7.9 \times 10^{6}$ & $3.0 \times 10^{7}$ & $7.9 \times 10^{7}$ & $1.7 \times 10^{8}$ & $3.2 \times 10^{8}$ \\
\end{tabular}
\end{center}

\begin{center}
\includegraphics[max width=\textwidth]{2025_02_27_d4b95d9718f0b9e2e6b9g-774}
\end{center}

Once the activation energy of a reaction is known, the value of a rate constant $k_{\mathrm{r}, 2}$ at a temperature $T_{2}$ can be predicted from its value $k_{\mathrm{r}, 1}$ at another temperature $T_{1}$. To do so, write

$$
\ln k_{\mathrm{r}, 1}=\ln A-\frac{E_{\mathrm{a}}}{R T_{1}} \quad \ln k_{\mathrm{r}, 2}=\ln A-\frac{E_{\mathrm{a}}}{R T_{2}}
$$

and then subtract the first from the second to obtain

$$
\ln k_{\mathrm{r}, 2}-\ln k_{\mathrm{r}, 1}=-\frac{E_{\mathrm{a}}}{R T_{2}}+\frac{E_{\mathrm{a}}}{R T_{1}}
$$

which can be rearranged into


\begin{equation*}
\ln \frac{k_{\mathrm{r}, 2}}{k_{\mathrm{r}, 1}}=\frac{E_{\mathrm{a}}}{R}\left(\frac{1}{T_{1}}-\frac{1}{T_{2}}\right) \tag{17D.2}
\end{equation*}


\subsection*{Brief illustration 17D. 1}
For a reaction with an activation energy of $50 \mathrm{~kJ} \mathrm{~mol}^{-1}$, an increase in the temperature from $25^{\circ} \mathrm{C}$ to $37^{\circ} \mathrm{C}$ (body temperature) corresponds to

$$
\begin{aligned}
\ln \frac{k_{\mathrm{r}, 2}}{k_{\mathrm{r}, 1}} & =\frac{50 \times 10^{3} \mathrm{Jmol}^{-1}}{8.3145 \mathrm{JK}^{-1} \mathrm{~mol}^{-1}}\left(\frac{1}{298 \mathrm{~K}}-\frac{1}{310 \mathrm{~K}}\right) \\
& =\frac{50 \times 10^{3}}{8.3145}\left(\frac{1}{298}-\frac{1}{310}\right)=0.781 \ldots
\end{aligned}
$$

By taking natural antilogarithms (that is, by forming $\mathrm{e}^{x}$ ), $k_{\mathrm{r}, 2}=2.18 k_{\mathrm{r}, 1}$. This result corresponds to slightly more than a doubling of the rate constant as the temperature is increased from 298 K to 310 K .

The fact that $E_{\mathrm{a}}$ is given by the slope of the plot of $\ln k_{\mathrm{r}}$ against $1 / T$ leads to the following conclusions:

\begin{itemize}
  \item A high activation energy signifies that the rate constant depends strongly on temperature.
  \item If a reaction has zero activation energy, its rate is independent of temperature.
  \item A negative activation energy indicates that the rate decreases as the temperature is raised.
\end{itemize}

For some reactions it is found that a plot of $\ln k_{\mathrm{r}}$ against $1 / T$ does not give a straight line. It is still possible to define an activation energy for these 'non-Arrhenius reactions' at a particular temperature by writing

$$
-\frac{E_{\mathrm{a}}}{R}=\frac{\mathrm{d} \ln k_{\mathrm{r}}}{\mathrm{~d}(1 / T)}=-T^{2} \frac{\mathrm{~d} \ln k_{\mathrm{r}}}{\mathrm{~d} T}
$$

and therefore

\[
E_{\mathrm{a}}=R T^{2}\left(\frac{\mathrm{~d} \ln k_{\mathrm{r}}}{\mathrm{~d} T}\right) \quad \begin{align*}
& \text { Activation energy }  \tag{17D.3}\\
& {[\text { definition }]}
\end{align*}
\]

This expression is the formal definition of activation energy. It reduces to the earlier one (as the slope of a straight line) for a tem-perature-independent activation energy (see Problem P17D.1). Non-Arrhenius behaviour is sometimes a sign that quantum mechanical tunnelling is playing a significant role in the reaction (Topic 7D). In biological reactions it might signal that an enzyme has undergone a structural change and has become less efficient.

\subsection*{17D. 2 The interpretation of the Arrhenius parameters}
For the present Topic the Arrhenius parameters are regarded as purely empirical quantities which summarize the variation of rate constants with temperature. Focus 18 provides a more elaborate interpretation.

\subsection*{(a) A first look at the energy requirements of reactions}
To interpret $E_{\mathrm{a}}$, consider how the molecular potential energy changes in the course of a chemical reaction that begins with a collision between A and B molecules (Fig. 17D.3). In the gas phase that step is an actual collision; in solution it is best regarded as a close encounter, possibly with excess energy, which might involve the solvent too. As the reaction event proceeds, A and B come into contact, distort, and begin to exchange or discard atoms. The reaction coordinate summarizes the collection of motions, such as changes in interatomic distances and bond angles, that are directly involved in the formation of products from reactants. The reaction coordinate is essentially a geometrical concept and quite distinct from the extent of reaction. The potential energy rises to a maximum and the cluster of atoms that corresponds to the region close to the maximum is called the activated complex.

After the maximum, the potential energy falls as the atoms in the cluster rearrange and eventually it reaches a value characteristic of the products. The climax of the reaction is at the peak of the potential energy curve, which corresponds to the activation energy $E_{a}$. Here two reactant molecules have come to such a degree of closeness and distortion that a small further distortion will send them in the direction of products. This crucial configuration is called the transition state of the reaction. Although some molecules entering the transition state might revert to reactants, if they pass through this configuration then it is inevitable that products will emerge from the encounter.

A note on good practice The terms 'activated complex' and 'transition state' are often used as synonyms; however, there is a distinction, which is best kept in mind. An activated complex is a cluster of atoms that corresponds to the region close to the\\
\includegraphics[max width=\textwidth, center]{2025_02_27_d4b95d9718f0b9e2e6b9g-775(1)}

Figure 17D. 3 A potential energy profile for an exothermic reaction. The height of the barrier between the reactants and products is the activation energy of the forward reaction.\\
maximum; a transition state is a conformation of the atoms in the activated complex that, after a small further distortion, leads inevitably to products.

The conclusion from the preceding discussion is that\\
The activation energy is the minimum energy reactants must have in order to form products.

For example, in a reaction mixture there are numerous molecular encounters each second, but only very few are sufficiently energetic to lead to reaction. The fraction of close encounters between reactants with energy in excess of $E_{\mathrm{a}}$ is given by the Boltzmann distribution (Prologue and Topic 13A) as $\mathrm{e}^{-E_{\mathrm{a}} / R T}$. This interpretation is confirmed by comparing this expression with the Arrhenius equation written in the form

\[
\begin{array}{ll}
k_{\mathrm{r}}=A \mathrm{e}^{-E_{\mathrm{a}} / R T} & \begin{array}{l}
\text { Arrhenius equation } \\
\text { [alternative form] }
\end{array} \tag{17D.4}
\end{array}
\]

which is obtained by taking antilogarithms of both sides of eqn 17D.1. Insight into this expression can be obtained by considering the role of the Boltzmann distribution for a simple model system.

How is that done? 17D. 1 Interpreting the exponential factor in the Arrhenius equation

Suppose the energy levels available to the system form a uniform array of separation $\varepsilon$ such that the energy levels are $i \varepsilon$, with $i=0,1,2, \ldots$ (Fig. 17D.4).\\
\includegraphics[max width=\textwidth, center]{2025_02_27_d4b95d9718f0b9e2e6b9g-775}

Figure 17D. 4 Equally spaced energy levels of a model system. As shown in the text, the fraction of molecules with energy of at least $\varepsilon_{\text {min }}$ is $\mathrm{e}^{-\varepsilon_{\text {min }} / k T}$.

The Boltzmann distribution for this system is

$$
\frac{N_{i}}{N}=\frac{\mathrm{e}^{-i \varepsilon \beta}}{q}=\left(1-\mathrm{e}^{-\varepsilon \beta}\right) \mathrm{e}^{-i \varepsilon \beta}
$$

where $N_{i}$ is the number of molecules in state $i, N$ is the total number of molecules, $\beta=1 / k T$, and the partition function $q$ comes from the result in eqn 13B.2a. The total number\\
of molecules in states with energy greater than or equal to $i_{\text {min }} \varepsilon$ is

$$
\begin{aligned}
\sum_{i=i_{\min }}^{\infty} N_{i} & =\sum_{i=0}^{\infty} N_{i}-\sum_{i=0}^{i_{\min }-1} N_{i}=N-\frac{N}{G} \sum_{i=0}^{\left.N_{i}=(N / g) e^{-i \varepsilon \beta}\right)} \mathrm{e}^{-i \varepsilon \beta} \\
& =N-\frac{N}{G} \sum_{i=0}^{i_{\min }-1}\left(\mathrm{e}^{-\varepsilon \beta}\right)^{i}
\end{aligned}
$$

The sum in blue is a finite geometrical series of the form $1+$ $r+r^{2}+\cdots$, with $r=\mathrm{e}^{-\varepsilon \beta}$. The sum of $n-1$ terms of such a series is $\left(1-r^{n}\right) /(1-r)$. The term in blue can therefore be written

$$
\begin{aligned}
\sum_{i=0}^{i_{\min }-1}\left(\mathrm{e}^{-\varepsilon \beta}\right)^{i} & =\frac{1-\mathrm{e}^{-i_{\min } \varepsilon \beta}}{1-\mathrm{e}^{-\varepsilon \beta}} \\
& =\mathfrak{q}\left(1-\mathrm{e}^{-i_{\min } \varepsilon \beta},\right.
\end{aligned}
$$

Therefore, the fraction of molecules in states with energy of at least $\varepsilon_{\text {min }}=i_{\text {min }} \varepsilon$ is

$$
\begin{aligned}
\frac{1}{N} \sum_{i=i_{\min }}^{\infty} N_{i} & =\frac{1}{N}\left\{N-\frac{N}{q} \times \mathfrak{g}\left(1-\mathrm{e}^{-i_{\min } \varepsilon \beta}\right)\right\} \\
& =1-\left(1-\mathrm{e}^{-i_{\min } \varepsilon \beta}\right)=\mathrm{e}^{-i_{\min } \varepsilon \beta}=\mathrm{e}^{-\varepsilon_{\min } / k T}
\end{aligned}
$$

which has the form of eqn 17D.4.

\subsection*{Brief illustration 17D. 2}
The fraction of molecules with energy at least $\varepsilon_{\min }$ is $\mathrm{e}^{-\varepsilon_{\min } / k T}$. By multiplying $\varepsilon_{\text {min }}$ and $k$ by $N_{\mathrm{A}}$, Avogadro's constant, and identifying $N_{A} \varepsilon_{\min }$ with $E_{\mathrm{a}}$, then the fraction $f$ of molecular collisions that occur with at least a molar kinetic energy $E_{\mathrm{a}}$ becomes $f=\mathrm{e}^{-E_{\mathrm{a}} / R T}$. With $E_{\mathrm{a}}=50 \mathrm{~kJ} \mathrm{~mol}^{-1}=5.0 \times 10^{4} \mathrm{~J} \mathrm{~mol}^{-1}$ and $T=298 \mathrm{~K}$ :

$$
\left.f=\mathrm{e}^{-\left(5.0 \times 10^{4} \mathrm{Jmol}\right.}{ }^{-1}\right)\left(8.3145 \mathrm{~J} \mathrm{~K}^{-1} \mathrm{~mol}^{-1} \times 298 \mathrm{~K}\right)=1.7 \times 10^{-9}
$$

or about one or two in a billion.

If the activation energy is zero, each collision leads to reaction and, according to the Arrhenius equation, the rate constant is equal to the frequency factor $A$. This factor can therefore be identified as the rate constant in the limit that each collision is successful. The exponential factor, $\mathrm{e}^{-E_{\mathrm{a}} / R T}$, gives the fraction of collisions that are sufficiently energetic to be successful, so the rate constant is reduced from $A$ to $A e^{-E_{\mathrm{a}} / R T}$.\\
\includegraphics[max width=\textwidth, center]{2025_02_27_d4b95d9718f0b9e2e6b9g-776}

Figure 17D. 5 A catalyst provides a different path with a lower activation energy. The result is an increase in the rate of the reaction (in both directions).

\subsection*{(b) The effect of a catalyst on the activation energy}
The Arrhenius equation predicts that the rate constant of a reaction can be increased by increasing the temperature or by decreasing the activation energy. Changing the temperature of a reaction mixture is easy to do. Reducing the activation energy is more challenging, but is possible if a reaction takes place in the presence of a suitable catalyst, a substance that accelerates a reaction but undergoes no net chemical change. The catalyst lowers the activation energy of the reaction by providing an alternative path (Fig. 17D.5).

\subsection*{Brief illustration 17D. 3}
Enzymes are biological catalysts. Suppose that an enzyme reduces both the activation energy and the frequency factor of a reaction by a factor of ten. Letting the activation energy change from $80 \mathrm{~kJ} \mathrm{~mol}^{-1}$ to $8 \mathrm{~kJ} \mathrm{~mol}^{-1}$, and using eqn 17D.4, the ratio of rate constants at 298 K is

$$
\begin{aligned}
& \left.=\frac{1}{10} \times \mathrm{e}^{-\left(8 \times 10^{3} \mathrm{Jmol}\right.}{ }^{-1}-80 \times 10^{3} \mathrm{Jmol}{ }^{-1}\right) /\left(8.3145 \mathrm{~J} \mathrm{~K}^{-1} \mathrm{~mol}^{-1}\right) \times(298 \mathrm{~K}) \\
& =4.2 \times 10^{11}
\end{aligned}
$$

The calculation shows that a decrease in the activation energy by an order of magnitude has a much greater impact on the rate constant than a decrease by the same order of magnitude in the frequency factor.

\subsection*{Checklist of concepts}
1. The activation energy, the parameter $E_{\mathrm{a}}$ in the Arrhenius equation, is the minimum energy that a molecular encounter needs in order to result in reaction.2. The higher the activation energy, the more sensitive the rate constant is to the temperature.3. The frequency factor is the rate constant in the limit that all encounters, irrespective of their energy, lead to reaction.4. A catalyst lowers the activation energy of a reaction.\subsection*{Checklist of equations}
\begin{center}
\begin{tabular}{llll}
\hline
Property or process & Equation & Comment & Equation number \\
\hline
Arrhenius equation & $\ln k_{\mathrm{r}}=\ln A-E_{\mathrm{a}} / R T$ &  & 17 D .1 \\
Activation energy & $E_{\mathrm{a}}=R T^{2}\left(\mathrm{~d} \ln k_{\mathrm{r}} / \mathrm{d} T\right)$ & General definition & 17 D .3 \\
Arrhenius equation & $k_{\mathrm{r}}=A \mathrm{e}^{-E_{\mathrm{a}} / R T}$ & Alternative form & 17 D .4 \\
\hline
\end{tabular}
\end{center}

\section*{TOPIC 17E Reaction mechanisms}
\subsection*{- Why do you need to know this material?}
The ability to construct the rate law for a reaction that takes place by a sequence of steps provides insight into chemical reactions at the molecular level and also suggests how the yield of desired products can be optimized.

\subsection*{- What is the key idea?}
Many chemical reactions occur as a sequence of simpler steps, with corresponding rate laws that can be combined into an overall rate law by applying a variety of approximations.

\subsection*{> What do you need to know already?}
You need to be familiar with the concept of rate laws (Topic 17A), and how to integrate them (Topics 17B and 17C). You also need to be familiar with the Arrhenius equation for the effect of temperature on the rate constant (Topic 17D).

The study of reaction rates leads to an understanding of the mechanism of a reaction, its analysis into a sequence of elementary steps. Simple elementary steps have simple rate laws, which can be combined into an overall rate law by invoking one or more approximations.

\subsection*{17E. 1 Elementary reactions}
Many reactions occur in a sequence of steps called elementary reactions, each of which involves only a small number of molecules or ions. A typical elementary reaction is

$$
\mathrm{H}+\mathrm{Br}_{2} \rightarrow \mathrm{HBr}+\mathrm{Br}
$$

Note that the phase of the species is not specified in the chemical equation for an elementary reaction and the equation represents the specific process occurring to individual molecules. This equation, for instance, signifies that an H atom attacks a $\mathrm{Br}_{2}$ molecule to produce an HBr molecule and a Br atom.

The molecularity of an elementary reaction is the number of molecules coming together to react in an elementary reaction. In a unimolecular reaction, a single molecule shakes itself apart or its atoms into a new arrangement, as in the isomeriza-\\
tion of cyclopropane to propene. In a bimolecular reaction, a pair of molecules collide and exchange energy, atoms, or groups of atoms, or undergo some other kind of change. It is important to distinguish molecularity from order:

\begin{itemize}
  \item reaction order is an empirical quantity, and obtained from the experimentally determined rate law;
  \item molecularity refers to an elementary reaction proposed as an individual step in a mechanism.
\end{itemize}

The rate law of a unimolecular elementary reaction is firstorder in the reactant:

\[
\mathrm{A} \rightarrow \mathrm{P} \quad \frac{\mathrm{~d}[\mathrm{~A}]}{\mathrm{d} t}=-k_{\mathrm{r}}[\mathrm{~A}] \quad \begin{align*}
& \text { Unimolecular }  \tag{17E.1}\\
& \text { elementary reaction }
\end{align*}
\]

where P denotes products (several different species may be formed). A unimolecular reaction is first order because the number of A molecules that decay in a short interval is proportional to the number available to decay. For instance, ten times as many decay in the same interval when there are initially 1000 A molecules as when there are only 100 present. Therefore, the rate of decomposition of A is proportional to its concentration at any moment during the reaction.

An elementary bimolecular reaction has a second-order rate law:

\[
\mathrm{A}+\mathrm{B} \rightarrow \mathrm{P} \quad \frac{\mathrm{~d}[\mathrm{~A}]}{\mathrm{d} t}=-k_{\mathrm{r}}[\mathrm{~A}][\mathrm{B}] \quad \begin{align*}
& \text { Bimolecular }  \tag{17E.2}\\
& \text { elementary reaction }
\end{align*}
\]

A bimolecular reaction is second order because its rate is proportional to the rate at which the reactant species meet, which in turn is proportional to both their concentrations. Therefore, if there is evidence that a reaction is a single-step, bimolecular process, the rate law can simply be written down as in eqn 17E. 2 (and then tested against experimental data).

\subsection*{Brief illustration 17E. 1}
Bimolecular elementary reactions are believed to account for many homogeneous reactions, such as the dimerizations of alkenes and dienes and reactions such as

$$
\mathrm{CH}_{3} \mathrm{I}(\mathrm{alc})+\mathrm{CH}_{3} \mathrm{CH}_{2} \mathrm{O}^{-}(\text {alc }) \rightarrow \mathrm{CH}_{3} \mathrm{OCH}_{2} \mathrm{CH}_{3}(\text { alc })+\mathrm{I}^{-}(\mathrm{alc})
$$

(where 'alc' signifies alcohol solution). There is evidence that the mechanism of this reaction is a single elementary step:

$$
\mathrm{CH}_{3} \mathrm{I}+\mathrm{CH}_{3} \mathrm{CH}_{2} \mathrm{O}^{-} \rightarrow \mathrm{CH}_{3} \mathrm{OCH}_{2} \mathrm{CH}_{3}+\mathrm{I}^{-}
$$

This mechanism is consistent with the observed rate law

$$
v=k_{\mathrm{r}}\left[\mathrm{CH}_{3} \mathrm{I}\right]\left[\mathrm{CH}_{3} \mathrm{CH}_{2} \mathrm{O}^{-}\right]
$$

The following sections describe how a series of simple steps can be combined into a mechanism and how the corresponding overall rate law can be derived. For the present it is important to note that, if the reaction is an elementary bimolecular process, then it has second-order kinetics, but if the kinetics is second order, then the reaction might be complex. The postulated mechanism can be explored only by detailed detective work on the system and by investigating whether side products or intermediates appear during the course of the reaction. Detailed analysis of this kind was one of the ways, for example, in which the reaction $\mathrm{H}_{2}(\mathrm{~g})+\mathrm{I}_{2}(\mathrm{~g}) \rightarrow 2 \mathrm{HI}(\mathrm{g})$ was shown to proceed by a complex mechanism. For many years the reaction had been accepted on good but insufficiently meticulous evidence as a fine example of a simple bimolecular reaction, $\mathrm{H}_{2}+\mathrm{I}_{2} \rightarrow \mathrm{HI}+\mathrm{HI}$, in which atoms exchanged partners during a collision.

\subsection*{17E.2 Consecutive elementary reactions}
Some reactions proceed through the formation of an intermediate (denoted I), as in the consecutive unimolecular reactions

$$
\mathrm{A} \xrightarrow{k_{a}} \mathrm{I} \xrightarrow{k_{\mathrm{b}}} \mathrm{P}
$$

Note that the intermediate occurs in the reaction steps but does not appear in the overall reaction, which in this case is $A \rightarrow P$. Any reverse reactions are ignored here, so the reaction proceeds from all A to all P , not to an equilibrium mixture of the two. An example of this type of mechanism is the decay of a radioactive family, such as

$$
{ }^{239} \mathrm{U} \xrightarrow{23.5 \min }{ }^{239} \mathrm{~Np} \xrightarrow{2.35 \text { days }}{ }^{239} \mathrm{Pu}
$$

(The times are half-lives.) The characteristics of this type of reaction are discovered by setting up the rate laws for the net rate of change of the concentration of each substance and then combining them in the appropriate manner.

The rate of unimolecular decomposition of A is


\begin{equation*}
\frac{\mathrm{d}[\mathrm{~A}]}{\mathrm{d} t}=-k_{\mathrm{a}}[\mathrm{~A}] \tag{17E.3a}
\end{equation*}


The intermediate I is formed from A (at a rate $k_{\mathrm{a}}[\mathrm{A}]$ ) but decays to P (at a rate $k_{\mathrm{b}}[\mathrm{I}]$ ). The net rate of formation of I is therefore


\begin{equation*}
\frac{\mathrm{d}[\mathrm{I}]}{\mathrm{d} t}=k_{\mathrm{a}}[\mathrm{~A}]-k_{\mathrm{b}}[\mathrm{I}] \tag{17E.3b}
\end{equation*}


The product P is formed by the unimolecular decay of I :


\begin{equation*}
\frac{\mathrm{d}[\mathrm{P}]}{\mathrm{d} t}=k_{\mathrm{b}}[\mathrm{I}] \tag{17E.3c}
\end{equation*}


If it is assumed that initially the molar concentration of $A$ is $[\mathrm{A}]_{0}$, the first-order rate law of eqn 17E.3a can be integrated (as in Topic 17B) to give


\begin{equation*}
[\mathrm{A}]=[\mathrm{A}]_{0} \mathrm{e}^{-k_{\mathrm{a}} t} \tag{17E.4a}
\end{equation*}


When this equation is substituted into eqn 17E.3b, the result is, after rearrangement,

$$
\frac{\mathrm{d}[\mathrm{I}]}{\mathrm{d} t}+k_{\mathrm{b}}[\mathrm{I}]=k_{\mathrm{a}}[\mathrm{~A}]_{0} \mathrm{e}^{-k_{\mathrm{a}} t}
$$

This differential equation has a standard form in the sense that it has been studied and its solution listed. With the initial condition $[\mathrm{I}]_{0}=0$, because no intermediate is present initially, the solution of the differential equation (provided $k_{\mathrm{a}} \neq k_{\mathrm{b}}$ ) is


\begin{equation*}
[\mathrm{I}]=\frac{k_{\mathrm{a}}}{k_{\mathrm{b}}-k_{\mathrm{a}}}\left(\mathrm{e}^{-k_{\mathrm{a}} t}-\mathrm{e}^{-k_{\mathrm{b}} t}\right)[\mathrm{A}]_{0} \tag{17E.4b}
\end{equation*}


At all times $[\mathrm{A}]+[\mathrm{I}]+[\mathrm{P}]=[\mathrm{A}]_{0}$, so it follows that


\begin{equation*}
[\mathrm{P}]=\left\{1+\frac{k_{\mathrm{e}} \mathrm{e}^{-k_{\mathrm{b}} t}-k_{\mathrm{b}} \mathrm{e}^{-\mathrm{k}_{\mathrm{a}} t}}{k_{\mathrm{b}}-k_{\mathrm{a}}}\right\}[\mathrm{A}]_{0} \tag{17E.4c}
\end{equation*}


The concentration of the intermediate I rises to a maximum and then falls to zero (Fig. 17E.1). The concentration of the product $P$ rises from zero towards $[\mathrm{A}]_{0}$, when all A has been converted to $P$.\\
\includegraphics[max width=\textwidth, center]{2025_02_27_d4b95d9718f0b9e2e6b9g-779}

Figure 17E. 1 The concentrations of $A, I$, and $P$ in the consecutive reaction scheme $\mathrm{A} \rightarrow \mathrm{I} \rightarrow \mathrm{P}$. The curves are plots of eqns 17E.4a-C with $k_{\mathrm{a}}=10 k_{\mathrm{b}}$. If the intermediate I is in fact the desired product, it is important to be able to predict when its concentration is greatest; see Example 17E.1.

\subsection*{Example 17E. 1 Analysing consecutive reactions}
Suppose that in an industrial batch process a substance A produces the desired compound I which goes on to decay to a worthless product $P$, each step of the reaction being first order. At what time will I be present in greatest concentration?

Collect your thoughts The time dependence of the concentration of I is given by eqn 17E.4b. Find the time at which [I] passes through a maximum, $t_{\max }$, by calculating the derivative $\mathrm{d}[\mathrm{I}] / \mathrm{d} t$ and then setting it equal to zero.

The solution It follows from eqn 17E.4b that

$$
\frac{\mathrm{d}[\mathrm{I}]}{\mathrm{d} t}=-\frac{k_{\mathrm{a}}\left(k_{\mathrm{a}} \mathrm{e}^{-k_{\mathrm{a} t} t}-k_{\mathrm{b}} \mathrm{e}^{-k_{\mathrm{b}} t}\right)[\mathrm{A}]_{0}}{k_{\mathrm{b}}-k_{\mathrm{a}}}
$$

This derivative is equal to zero when $t=t_{\max }$ and $k_{\mathrm{a}} \mathrm{e}^{-k_{\mathrm{t}} t_{\text {max }}}=$ $k_{\mathrm{b}} \mathrm{e}^{-k_{\mathrm{b}} t_{\text {max }}}$. Therefore, taking natural logarithms of both sides gives

$$
\ln k_{\mathrm{a}}-k_{\mathrm{a}} t_{\max }=\ln k_{\mathrm{b}}-k_{\mathrm{b}} t_{\text {max }}
$$

which rearranges to

$$
\overbrace{\ln k_{\mathrm{a}}-\ln k_{\mathrm{b}}}^{\ln \left(k_{\mathrm{a}^{\prime}} / k_{\mathrm{b}}\right)}=k_{\mathrm{a}} t_{\max }-k_{\mathrm{b}} t_{\max }=\left(k_{\mathrm{a}}-k_{\mathrm{b}}\right) t_{\max }
$$

It then follows that

$$
t_{\max }=\frac{1}{k_{\mathrm{a}}-k_{\mathrm{b}}} \ln \frac{k_{\mathrm{a}}}{k_{\mathrm{b}}}
$$

Comment. For a given value of $k_{\mathrm{a}}$, as $k_{\mathrm{b}}$ increases both the time at which $[\mathrm{I}]$ is a maximum and the yield of I decrease.

Self-test 17E. 1 Calculate the maximum concentration of I and justify the last remark.

$$
\left.\left({ }^{p} y-{ }^{q} y\right)\right)^{q} y==^{c}\left({ }^{9} y /\left.\right|^{p} y\right)={ }^{0}[\mathrm{~V}] /^{\mathrm{xew}}[I]: \text { :zassu }
$$

\subsection*{17E. 3 The steady-state approximation}
One feature of the calculation so far has probably not gone unnoticed: there is a considerable increase in mathematical complexity as soon as the reaction mechanism has more than a couple of steps or reverse reactions are taken into account. A reaction scheme involving many steps is nearly always unsolvable analytically, and alternative methods of solution are necessary. One approach is to integrate the rate laws numerically. An alternative approach, which continues to be widely used because it leads to convenient expressions and more readily digestible results, is to make an approximation.

The steady-state approximation (which is also widely called the quasi-steady-state approximation to distinguish it from a true steady state) assumes that the intermediate, I , is in a low,\\
\includegraphics[max width=\textwidth, center]{2025_02_27_d4b95d9718f0b9e2e6b9g-780}

Figure 17E. 2 The basis of the steady-state approximation. It is supposed that the concentrations of intermediates remain small and hardly change during most of the course of the reaction.\\
constant concentration. More specifically, after an initial induction period, an interval during which the concentrations of intermediates rise from zero, the rates of change of the concentrations of all reaction intermediates are negligibly small during the major part of the reaction (Fig. 17E.2):


\begin{equation*}
\frac{\mathrm{d}[\mathrm{I}]}{\mathrm{d} t} \approx 0 \tag{17E.5}
\end{equation*}


Steady-state approximation

This approximation greatly simplifies the discussion of reaction schemes. For example, to apply the approximation to the consecutive first-order mechanism, $\mathrm{d}[\mathrm{I}] / \mathrm{d} t$ is set equal to 0 in eqn 17 E .3 b , which then becomes $k_{\mathrm{a}}[\mathrm{A}]-k_{\mathrm{b}}[\mathrm{I}]=0$. It then follows that


\begin{equation*}
[\mathrm{I}]=\frac{k_{\mathrm{a}}}{k_{\mathrm{b}}}[\mathrm{~A}] \tag{17E.6}
\end{equation*}


The steady-state approximation requires the concentration of the intermediate to be low relative to that of the reactants, which is the case when $k_{\mathrm{a}} \ll k_{\mathrm{b}}$. Equation 17E. 6 implies that the concentration of the intermediate changes as the concentration of A changes, but if $k_{\mathrm{a}} / k_{\mathrm{b}} \ll 1$ it changes very little. Both requirements of the steady-state approximation, the low concentration of intermediate and its slow change, are therefore satisfied.

On substituting the value of [I] in eqn 17E. 6 into eqn 17E.3c, that equation becomes


\begin{equation*}
\frac{\mathrm{d}[\mathrm{P}]}{\mathrm{d} t}=k_{\mathrm{b}}[\mathrm{I}] \stackrel{(1]=\left(k_{\mathrm{a}} / k_{\mathrm{b}}\right)[\mathrm{A}]}{\approx} k_{\mathrm{a}}[\mathrm{~A}] \tag{17E.7}
\end{equation*}


It follows that the rate of formation of P is the same as the rate of loss of A (as given by eqn 17E.3a), and that both processes are governed by the rate constant $k_{\mathrm{a}}$. In effect, the reactant A flows directly through to becoming P without any I accumulating on the way. The solution of eqn 17E. 7 is found by\\
\includegraphics[max width=\textwidth, center]{2025_02_27_d4b95d9718f0b9e2e6b9g-781(1)}

Figure 17E. 3 A comparison of the exact result for the concentrations of a consecutive reaction and the concentrations obtained by using the steady-state approximation (dotted lines) for $k_{\mathrm{b}}=20 k_{\mathrm{a}}$. (The curve for $[\mathrm{A}]$ is unchanged.)\\[0pt]
substituting the solution for [A], eqn 17E.4a, and integrating the resulting expression:

$$
\int_{[\mathrm{A}]=[\mathrm{A}]_{0} \mathrm{e}^{-k_{\mathrm{a}} t}}^{[\mathrm{P}]} \mathrm{d}[\mathrm{P}]=\int_{0}^{t} k_{\mathrm{a}}[\mathrm{~A}] \mathrm{d} t=k_{\mathrm{a}}[\mathrm{~A}]_{0} \overbrace{\int_{0}^{t} \mathrm{e}^{-k_{\mathrm{a}} t} \mathrm{~d} t}^{\text {Integral E. } 1}
$$

Hence


\begin{equation*}
[\mathrm{P}]=[\mathrm{A}]_{0}\left(1-\mathrm{e}^{-k_{\mathrm{a}} t}\right) \tag{17E.8}
\end{equation*}


This result is the same as in eqn 17E. 4 c when $k_{\mathrm{a}} \ll k_{\mathrm{b}}$; however, the use of the steady-state approximation is much simpler. Figure 17E. 3 compares the approximate solutions found here with the exact solutions found earlier: $k_{\mathrm{a}}$ does not have to be very much smaller than $k_{\mathrm{b}}$ for the approach to be reasonably accurate.

\subsection*{Example 17E. 2}
\subsection*{Using the steady-state approximation}
Devise the rate law for the decomposition of $\mathrm{N}_{2} \mathrm{O}_{5}, 2 \mathrm{~N}_{2} \mathrm{O}_{5}(\mathrm{~g})$ $\rightarrow 4 \mathrm{NO}_{2}(\mathrm{~g})+\mathrm{O}_{2}(\mathrm{~g})$ on the basis of the following mechanism:

$$
\begin{array}{ll}
\mathrm{N}_{2} \mathrm{O}_{5} \rightarrow \mathrm{NO}_{2}+\mathrm{NO}_{3} & k_{\mathrm{a}} \\
\mathrm{NO}_{2}+\mathrm{NO}_{3} \rightarrow \mathrm{~N}_{2} \mathrm{O}_{5} & k_{\mathrm{a}}^{\prime} \\
\mathrm{NO}_{2}+\mathrm{NO}_{3} \rightarrow \mathrm{NO}_{2}+\mathrm{O}_{2}+\mathrm{NO} & k_{\mathrm{b}} \\
\mathrm{NO}+\mathrm{N}_{2} \mathrm{O}_{5} \rightarrow \mathrm{NO}_{2}+\mathrm{NO}_{2}+\mathrm{NO}_{2} & k_{\mathrm{c}}
\end{array}
$$

A note on good practice Note that when writing the equation for an elementary reaction all the species are displayed individually; so write $\mathrm{A} \rightarrow \mathrm{B}+\mathrm{B}$, for instance, not $\mathrm{A} \rightarrow 2 \mathrm{~B}$.

Collect your thoughts First identify the intermediates and for each of them write an expression for the net rate of formation. Then apply the steady-state approximation and set these\\
net rates to zero. You can then solve the resulting equations algebraically to obtain expressions for the concentrations of the intermediates. Finally, use these solutions to obtain an expression for the overall rate of consumption of $\mathrm{N}_{2} \mathrm{O}_{5}$.\\
The solution The intermediates are NO and $\mathrm{NO}_{3}$; the net rates of change of their concentrations are

$$
\begin{aligned}
& \frac{\mathrm{d}[\mathrm{NO}]}{\mathrm{d} t}=k_{\mathrm{b}}\left[\mathrm{NO}_{2}\right]\left[\mathrm{NO}_{3}\right]-k_{\mathrm{c}}[\mathrm{NO}]\left[\mathrm{N}_{2} \mathrm{O}_{5}\right] \approx 0 \\
& \frac{\mathrm{~d}\left[\mathrm{NO}_{3}\right]}{\mathrm{d} t}=k_{\mathrm{a}}\left[\mathrm{~N}_{2} \mathrm{O}_{5}\right]-k_{\mathrm{a}}^{\prime}\left[\mathrm{NO}_{2}\right]\left[\mathrm{NO}_{3}\right]-k_{\mathrm{b}}\left[\mathrm{NO}_{2}\right]\left[\mathrm{NO}_{3}\right] \approx 0
\end{aligned}
$$

The solutions of these two simultaneous equations (in blue) are

$$
\left[\mathrm{NO}_{3}\right]=\frac{k_{\mathrm{a}}\left[\mathrm{~N}_{2} \mathrm{O}_{5}\right]}{\left(k_{\mathrm{a}}^{\prime}+k_{\mathrm{b}}\right)\left[\mathrm{NO}_{2}\right]} \quad[\mathrm{NO}]=\frac{k_{\mathrm{b}}\left[\mathrm{NO}_{2}\right]\left[\mathrm{NO}_{3}\right]}{k_{\mathrm{c}}\left[\mathrm{~N}_{2} \mathrm{O}_{5}\right]}=\frac{k_{\mathrm{a}} k_{\mathrm{b}}}{\left(k_{\mathrm{a}}^{\prime}+k_{\mathrm{b}}\right) k_{\mathrm{c}}}
$$

The net rate of change of concentration of $\mathrm{N}_{2} \mathrm{O}_{5}$ is then

$$
\begin{aligned}
\frac{\mathrm{d}\left[\mathrm{~N}_{2} \mathrm{O}_{5}\right]}{\mathrm{d} t} & =-k_{\mathrm{a}}\left[\mathrm{~N}_{2} \mathrm{O}_{5}\right]+k_{\mathrm{a}}^{\prime}\left[\mathrm{NO}_{2}\right]\left[\mathrm{NO}_{3}\right]-k_{\mathrm{c}}[\mathrm{NO}]\left[\mathrm{N}_{2} \mathrm{O}_{5}\right] \\
& =-k_{\mathrm{a}}\left[\mathrm{~N}_{2} \mathrm{O}_{5}\right]+\frac{k_{\mathrm{a}} k_{\mathrm{a}}^{\prime}\left[\mathrm{N}_{2} \mathrm{O}_{5}\right]}{k_{\mathrm{a}}^{\prime}+k_{\mathrm{b}}}-\frac{k_{\mathrm{a}} k_{\mathrm{b}}}{k_{\mathrm{a}}^{\prime}+k_{\mathrm{b}}}\left[\mathrm{~N}_{2} \mathrm{O}_{5}\right] \\
& =-\frac{2 k_{\mathrm{a}} k_{\mathrm{b}}\left[\mathrm{~N}_{2} \mathrm{O}_{5}\right]}{k_{\mathrm{a}}^{\prime}+k_{\mathrm{b}}}
\end{aligned}
$$

That is, $\mathrm{N}_{2} \mathrm{O}_{5}$ decays with a first-order rate law with a rate constant that depends on $k_{\mathrm{a}}, k_{\mathrm{a}}^{\prime}$, and $k_{\mathrm{b}}$ but not on $k_{c}$.\\
Self-test 17E. 2 Derive the rate law for the decomposition of ozone in the reaction $2 \mathrm{O}_{3}(\mathrm{~g}) \rightarrow 3 \mathrm{O}_{2}(\mathrm{~g})$ on the basis of the (incomplete) mechanism

$$
\begin{array}{ll}
\mathrm{O}_{3} \rightarrow \mathrm{O}_{2}+\mathrm{O} & k_{\mathrm{a}} \\
\mathrm{O}_{2}+\mathrm{O} \rightarrow \mathrm{O}_{3} & k_{\mathrm{a}}^{\prime} \\
\mathrm{O}+\mathrm{O}_{3} \rightarrow \mathrm{O}_{2}+\mathrm{O}_{2} & k_{\mathrm{b}}
\end{array}
$$

\begin{center}
\includegraphics[max width=\textwidth]{2025_02_27_d4b95d9718f0b9e2e6b9g-781}
\end{center}

\subsection*{17E. 4 The rate-determining step}
When the steady-state approximation is valid, which is when $k_{\mathrm{a}} \ll k_{\mathrm{b}}$ in the reaction $\mathrm{A} \rightarrow \mathrm{I} \rightarrow \mathrm{P}$, the decrease in the concentration of $A$ is matched by an increase in the concentration of P. It is important to realize that the rates of the steps $\mathrm{A} \rightarrow \mathrm{I}$ and $\mathrm{I} \rightarrow \mathrm{P}$ are the same: the concentration of I is so low compared to the concentration of A that even though $k_{\mathrm{a}} \ll k_{\mathrm{b}}$ (and therefore $k_{\mathrm{b}} \gg k_{\mathrm{a}}$ ) the rate of the second step, $k_{\mathrm{b}}[\mathrm{I}]$, matches that of the first, $k_{\mathrm{a}}[\mathrm{A}]$.\\
\includegraphics[max width=\textwidth, center]{2025_02_27_d4b95d9718f0b9e2e6b9g-782}

Figure 17E. 4 In these diagrams of reaction schemes, heavy arrows represent steps with large rate constants and light arrows represent steps with small rate constants. (a) The first step is rate-determining; (b) the second step is rate-determining; (c) although one step has a small rate constant, it is not ratedetermining because there is a route with a large rate constant that circumvents it.

A note on good practice It is commonly said that 'the first step is slow and the second is fast, so the first step is rate-determining'. Such a statement is incorrect: the two rates are equal; it is the rate constants that are different.

In general, the rate-determining step (RDS) is the step in a mechanism that controls the overall rate of the reaction (in the present example, the first step governed by $k_{\mathrm{a}}$, with $k_{\mathrm{a}} \ll k_{\mathrm{b}}$ ). The rate-determining step must be a crucial gateway for the formation of products, and not just a reaction with a small rate constant. If another reaction with a larger rate constant can also lead to products, then the step with the small rate constant is irrelevant because it can be sidestepped (Fig. 17E.4). In some cases, when a first-order reaction is in competition with a second-order reaction, the criterion has to be expressed in terms of the relative sizes of the first-order (for one step) and a pseudofirst-order (for the second step) rate constants, for only then can their magnitudes be compared. This point is illustrated in the next Brief illustration.\\
\includegraphics[max width=\textwidth, center]{2025_02_27_d4b95d9718f0b9e2e6b9g-782(1)}

Figure 17E. 5 The reaction profile for a mechanism in which the first step (RDS) is rate-determining.

The rate law of a reaction that has a rate-determining step can often-but certainly not always-be written down almost by inspection. If the first step in a mechanism is ratedetermining, then the rate of the overall reaction is equal to the rate of that step because the rate constants of the subsequent steps are such that the intermediates immediately flow through these steps to give products. Moreover, because the rate-determining step is the one with the smallest rate constant, then it follows that the rate-determining step is the one with the highest activation energy. Once over the initial barrier, the intermediates cascade into products (Fig. 17E.5).

\subsection*{Brief illustration 17E. 2}
The oxidation of NO to $\mathrm{NO}_{2}, 2 \mathrm{NO}(\mathrm{g})+\mathrm{O}_{2}(\mathrm{~g}) \rightarrow 2 \mathrm{NO}_{2}(\mathrm{~g})$, proceeds by the following mechanism:

$$
\begin{array}{ll}
\mathrm{NO}+\mathrm{NO} \rightarrow \mathrm{~N}_{2} \mathrm{O}_{2} & k_{\mathrm{a}} \\
\mathrm{~N}_{2} \mathrm{O}_{2} \rightarrow \mathrm{NO}+\mathrm{NO} & k_{\mathrm{a}}^{\prime} \\
\mathrm{N}_{2} \mathrm{O}_{2}+\mathrm{O}_{2} \rightarrow \mathrm{NO}_{2}+\mathrm{NO}_{2} & k_{\mathrm{b}}
\end{array}
$$

with rate law (see Problem P17E.6)

$$
\frac{\mathrm{d}\left[\mathrm{NO}_{2}\right]}{\mathrm{d} t}=\frac{2 k_{\mathrm{a}} k_{\mathrm{b}}[\mathrm{NO}]^{2}\left[\mathrm{O}_{2}\right]}{k_{\mathrm{a}}^{\prime}+k_{\mathrm{b}}\left[\mathrm{O}_{2}\right]}
$$

When the concentration of $\mathrm{O}_{2}$ in the reaction mixture is so high that $k_{\mathrm{a}}^{\prime} \ll k_{\mathrm{b}}\left[\mathrm{O}_{2}\right]$, the rate law then simplifies to

$$
\frac{\mathrm{d}\left[\mathrm{NO}_{2}\right]}{\mathrm{d} t}=2 k_{\mathrm{a}}[\mathrm{NO}]^{2}
$$

which shows that the formation of $\mathrm{N}_{2} \mathrm{O}_{2}$ in the first step is rate-determining. The rate law could also have been written by inspection of the mechanism, because the rate law for the overall reaction is simply the rate law of that rate-determining step.

\subsection*{17e. 5 Pre-equilibria}
Now consider a slightly more complicated mechanism in which an intermediate I reaches an equilibrium with the reactants A and B:

$$
\mathrm{A}+\mathrm{B} \underset{k_{\mathrm{d}}^{\prime}}{\stackrel{k_{a}}{\rightleftharpoons}} \mathrm{I} \xrightarrow{k_{b}} \mathrm{P}
$$

Pre-equilibrium\\
(17E.9)

This scheme involves a pre-equilibrium, in which an intermediate is in equilibrium with the reactants. A pre-equilibrium can arise when the rate of decay of the intermediate back into reactants is much faster than the rate at which it forms products; this condition is satisfied if $k_{\mathrm{a}}^{\prime} \gg k_{\mathrm{b}}$. Because A, B,\\
and I are assumed to be in equilibrium, it follows that (see Topic 17C)

$$
K=\frac{[\mathrm{I}] c^{\ominus}}{[\mathrm{A}][\mathrm{B}]} \text { and }[\mathrm{I}]=\frac{K}{c^{\ominus}}[\mathrm{A}][\mathrm{B}]=\frac{k_{\mathrm{a}}}{k_{\mathrm{a}}^{\prime}}[\mathrm{A}][\mathrm{B}]
$$

In writing these equations, the rate of reaction of I to form P is presumed to be too slow to affect the maintenance of the preequilibrium (see the following Example). The rate of formation of P may now be written

$$
\left.\frac{\mathrm{d}[\mathrm{P}]}{\mathrm{d} t}=k_{\mathrm{b}}[\mathrm{I}]=k_{\mathrm{b}} \frac{k_{\mathrm{a}}}{k_{\mathrm{a}}^{\prime}}[\mathrm{A}]=\left(k_{\mathrm{a}} / k_{\mathrm{a}}^{\prime}\right][\mathrm{A}][\mathrm{B}]\right]
$$

This rate law has the form of a second-order rate law with a composite rate constant:


\begin{equation*}
\frac{\mathrm{d}[\mathrm{P}]}{\mathrm{d} t}=k_{\mathrm{r}}[\mathrm{~A}][\mathrm{B}] \quad \text { with } \quad k_{\mathrm{r}}=\frac{k_{\mathrm{b}} k_{\mathrm{a}}}{k_{\mathrm{a}}^{\prime}} \tag{17E.12}
\end{equation*}


In this pre-equilibrium mechanism the final step, $I \rightarrow P$, is rate-determining. The preceding steps control the steady concentration of the intermediate.

\subsection*{Example 17E. 3 \\
 Analysing a pre-equilibrium using the}
 steady-state assumptionAnalyse the scheme shown in eqn 17E. 9 using the steady-state approximation.

Collect your thoughts Begin by writing the net rates of change of the concentrations of P and I, and then invoke the steadystate approximation for the intermediate I. Use the resulting expression to obtain the rate of change of the concentration of $P$.

The solution The net rates of change of P and I are

$$
\begin{aligned}
& \frac{\mathrm{d}[\mathrm{P}]}{\mathrm{d} t}=k_{\mathrm{b}}[\mathrm{I}] \\
& \frac{\mathrm{d}[\mathrm{I}]}{\mathrm{d} t}=k_{\mathrm{a}}[\mathrm{~A}][\mathrm{B}]-k_{\mathrm{a}}^{\prime}[\mathrm{I}]-k_{\mathrm{b}}[\mathrm{I}] \approx 0
\end{aligned}
$$

The second equation implies that

$$
[\mathrm{I}] \approx \frac{k_{\mathrm{a}}[\mathrm{~A}][\mathrm{B}]}{k_{\mathrm{a}}^{\prime}+k_{\mathrm{b}}}
$$

Now substitute this result into the expression for the rate of formation of P :

$$
\frac{\mathrm{d}[\mathrm{P}]}{\mathrm{d} t}=k_{\mathrm{b}}[\mathrm{I}] \approx k_{\mathrm{b}} \frac{k_{\mathrm{a}}[\mathrm{~A}][\mathrm{B}]}{k_{\mathrm{a}}^{\prime}+k_{\mathrm{b}}}=k_{\mathrm{r}}[\mathrm{~A}][\mathrm{B}] \quad \text { with } \quad k_{\mathrm{r}}=\frac{k_{\mathrm{a}} k_{\mathrm{b}}}{k_{\mathrm{a}}^{\prime}+k_{\mathrm{b}}}
$$

This expression reduces to that in eqn 17E. 12 when the rate constant for the decay of I into products is much smaller than that for its decay into reactants, $k_{\mathrm{b}} \ll k_{\mathrm{a}}^{\prime}$.

Self-test 17E. 3 Show that the pre-equilibrium mechanism in which $2 \mathrm{~A} \rightleftharpoons \mathrm{I}(\mathrm{K})$ followed by $\mathrm{I}+\mathrm{B} \rightarrow \mathrm{P}\left(k_{\mathrm{b}}\right)$ results in an overall third-order reaction.\\
\includegraphics[max width=\textwidth, center]{2025_02_27_d4b95d9718f0b9e2e6b9g-783}

One feature to note is that although each of the rate constants in eqn 17E. 12 increases with temperature, this might not be true of $k_{\mathrm{r}}$ itself. Thus, if the rate constant $k_{\mathrm{a}}^{\prime}$ increases more rapidly than the product $k_{\mathrm{a}} k_{\mathrm{b}}$ increases, then $k_{\mathrm{r}}=k_{\mathrm{a}} k_{\mathrm{b}} / k_{\mathrm{a}}^{\prime}$ decreases with increasing temperature and the reaction goes more slowly as the temperature is raised. Mathematically, the overall reaction is said to have a 'negative activation energy'. For example, suppose that each rate constant in eqn 17E. 12 has an Arrhenius temperature dependence (Topic 17D). It follows from the Arrhenius equation (eqn 17D.4, $k_{\mathrm{r}}=A \mathrm{e}^{-E_{\mathrm{a}} / R T}$ ) that

$$
k_{\mathrm{r}}=\frac{\left(A_{\mathrm{a}} \mathrm{e}^{-E_{\mathrm{a}, \mathrm{a}} / R T}\right)\left(A_{\mathrm{b}} \mathrm{e}^{-E_{\mathrm{a}, \mathrm{~b}} / R T}\right)}{A_{\mathrm{a}^{\prime}} \mathrm{e}^{-E_{\mathrm{a} \mathrm{a}^{\prime}} / R T}}=\frac{\begin{array}{l}
\mathrm{e}^{x+y}=\mathrm{e}^{x} \mathrm{e}^{y} \\
\mathrm{e}^{x-y}=\mathrm{e}^{\mathrm{x}} / \mathrm{e}^{y}
\end{array}}{A_{\mathrm{b}}} \mathrm{e}^{-\left(E_{\mathrm{a}, \mathrm{a}}+E_{\mathrm{a}, \mathrm{~b}}-E_{\mathrm{a}, \mathrm{a}^{\prime}}\right) / R T}
$$

The effective activation energy of the reaction is therefore


\begin{equation*}
E_{\mathrm{a}}=E_{\mathrm{a}, \mathrm{a}}+E_{\mathrm{a}, \mathrm{~b}}-E_{\mathrm{a}, \mathrm{a}^{\prime}} \tag{17E.13}
\end{equation*}


This activation energy is positive if $E_{\mathrm{a}, \mathrm{a}}+E_{\mathrm{a}, \mathrm{b}}>E_{\mathrm{a}, \mathrm{a}^{\prime}}$ (Fig. 17E.6a) but negative if $E_{\mathrm{a}, \mathrm{a}^{\prime}}>E_{\mathrm{a}, \mathrm{a}}+E_{\mathrm{a}, \mathrm{b}}$ (Fig. 17E.6b). An important consequence of this discussion is that it is necessary to be cautious when making predictions about the effect of temperature on reactions that are the outcome of several steps.\\
\includegraphics[max width=\textwidth, center]{2025_02_27_d4b95d9718f0b9e2e6b9g-783(1)}

Figure 17E. 6 For a reaction with a pre-equilibrium, there are three activation energies to take into account: two referring to the reversible steps of the pre-equilibrium and one for the final step. The relative magnitudes of the activation energies determine whether the overall activation energy is (a) positive or (b) negative.

\subsection*{17E. 6 Kinetic and thermodynamic control of reactions}
In some cases reactants can give rise to a variety of products, as in nitrations of mono-substituted benzene, when various proportions of the ortho-, meta-, and para-substituted products are obtained, depending on the directing power of the original substituent. Suppose two products, $\mathrm{P}_{1}$ and $\mathrm{P}_{2}$, are produced by the following competing reactions:

$$
\begin{array}{ll}
\mathrm{A}+\mathrm{B} \rightarrow \mathrm{P}_{1} & v\left(\mathrm{P}_{1}\right)=k_{\mathrm{r}, 1}[\mathrm{~A}][\mathrm{B}] \\
\mathrm{A}+\mathrm{B} \rightarrow \mathrm{P}_{2} & v\left(\mathrm{P}_{2}\right)=k_{\mathrm{r}, 2}[\mathrm{~A}][\mathrm{B}]
\end{array}
$$

The relative proportion in which the two products have been produced at a given stage of the reaction (before it has reached\\
equilibrium) is given by the ratio of the two rates, and therefore by the ratio of the two rate constants:


\begin{equation*}
\frac{\left[\mathrm{P}_{2}\right]}{\left[\mathrm{P}_{1}\right]}=\frac{k_{\mathrm{r}, 2}}{k_{\mathrm{r}, 1}} \tag{17E.14}
\end{equation*}


Kinetic control

This ratio represents the kinetic control over the proportions of products, control that stems from relative rates rather than thermodynamic considerations about equilibrium. Kinetic control is a common feature of the reactions encountered in organic chemistry, where reactants are chosen that facilitate pathways favouring the formation of a desired product. If a reaction is allowed to reach equilibrium, then the proportion of products is determined by thermodynamic rather than kinetic considerations, and the ratio of concentrations is controlled by considerations of the standard reaction Gibbs energy.

\subsection*{Checklist of concepts}
\begin{enumerate}
  \item The mechanism of reaction is the sequence of elementary steps that leads from reactants to products.
  \item The molecularity of an elementary reaction is the number of molecules coming together to react.
  \item An elementary unimolecular reaction has first-order kinetics; an elementary bimolecular reaction has sec-ond-order kinetics.
  \item The rate-determining step is the step in a reaction mechanism that controls the rate of the overall reaction.\\
$\square$ 5. In the steady-state approximation, it is assumed that the concentrations of all reaction intermediates remain constant and small throughout the reaction.6. Pre-equilibrium is a state in which an intermediate is in equilibrium with the reactants and which arises when the rate of decay of the intermediate back to reactants is much faster than the rate at which products are formed from the intermediate.
  \item Kinetic control over the proportions of products stems from relative rates rather than thermodynamic considerations about equilibrium.
\end{enumerate}

\subsection*{Checklist of equations}
\begin{center}
\begin{tabular}{|c|c|c|c|}
\hline
Property & Equation & Comment & Equation number \\
\hline
Unimolecular reaction & $\mathrm{d}[\mathrm{A}] / \mathrm{d} t=-k_{\mathrm{r}}[\mathrm{A}]$ & $\mathrm{A} \rightarrow \mathrm{P}$ & 17E. 1 \\
\hline
Bimolecular reaction & $\mathrm{d}[\mathrm{A}] / \mathrm{d} t=-k_{\mathrm{r}}[\mathrm{A}][\mathrm{B}]$ & $\mathrm{A}+\mathrm{B} \rightarrow \mathrm{P}$ & 17E. 2 \\
\hline
Consecutive reactions & \( \begin{aligned} & {[\mathrm{A}]=[\mathrm{A}]_{0} \mathrm{e}^{-k_{\mathrm{a}} t}} \\ & {[\mathrm{I}]=\left(k_{\mathrm{a}} /\left(k_{\mathrm{b}}-k_{\mathrm{a}}\right)\right)\left(\mathrm{e}^{-k_{\mathrm{a}} t}-\mathrm{e}^{-k_{\mathrm{b}} t}\right)[\mathrm{A}]_{0}} \\ & {[\mathrm{P}]=\left\{1+\left(k_{\mathrm{a}} \mathrm{e}^{-k_{\mathrm{b}} t}-k_{\mathrm{b}} \mathrm{e}^{-k_{\mathrm{a}} t}\right) /\left(k_{\mathrm{b}}-k_{\mathrm{a}}\right)\right\}[\mathrm{A}]_{0}} \end{aligned} \) & \( \mathrm{A} \xrightarrow{k_{a}} \mathrm{I} \xrightarrow{k_{\mathrm{b}}} \mathrm{P} \) & 17E. 4 \\
\hline
Steady-state approximation & $\mathrm{d}[\mathrm{I}] / \mathrm{d} t \approx 0$ & I is an intermediate & 17E. 5 \\
\hline
\end{tabular}
\end{center}

\section*{TOPIC 17F Examples of reaction mechanisms}
\subsection*{- Why do you need to know this material?}
Some important reactions have complex mechanisms and need special treatment, so you need to see how to make and implement assumptions about the relative rates of the steps in a mechanism.

\subsection*{> What is the key idea?}
The steady-state approximation can often be used to derive rate laws for proposed mechanisms.

\subsection*{- What do you need to know already?}
You need to be familiar with the concept of rate laws (Topic 17A) and the formulation of an overall rate law from a mechanism by using the steady-state approximation (Topic 17E).

Many reactions take place by mechanisms that involve several elementary steps. In each case it is possible to approach the setting up (and testing) of a rate law by proposing a mechanism and, when appropriate, applying the steady-state approximation.

\subsection*{17F. 1 Unimolecular reactions}
A number of gas-phase reactions follow first-order kinetics, as in the isomerization of cyclopropane to propene:

$$
\text { cyclo }-\mathrm{C}_{3} \mathrm{H}_{6}(\mathrm{~g}) \rightarrow \mathrm{CH}_{3} \mathrm{CH}=\mathrm{CH}_{2}(\mathrm{~g}) \quad v=k_{\mathrm{r}}\left[\text { cyclo }-\mathrm{C}_{3} \mathrm{H}_{6}\right]
$$

The problem with the interpretation of first-order rate laws is that presumably a molecule acquires enough energy to react as a result of collisions with other molecules. However, collisions are simple bimolecular events, so how can they result in a first-order rate law? First-order gas-phase reactions are widely called 'unimolecular reactions' because they also involve an elementary unimolecular step in which the reactant molecule changes into the product. This term must be used with cau-\\
\includegraphics[max width=\textwidth, center]{2025_02_27_d4b95d9718f0b9e2e6b9g-785}

Figure 17F. 1 A representation of the Lindemann-Hinshelwood mechanism of unimolecular reactions. The species $A$ is excited by collision with $A$, and the energized $A$ molecule ( $A^{*}$ ) may either be deactivated by a collision with A or go on to decay by a unimolecular process to form products.\\
tion, however, because the overall mechanism has bimolecular as well as unimolecular steps.

The first successful explanation of unimolecular reactions was provided by Frederick Lindemann in 1921 and then elaborated by Cyril Hinshelwood. In the Lindemann-Hinshelwood mechanism it is supposed that a reactant molecule A becomes energetically excited by collision with another A molecule in a bimolecular step (Fig. 17F.1):


\begin{equation*}
\mathrm{A}+\mathrm{A} \rightarrow \mathrm{~A}^{\star}+\mathrm{A} \quad \frac{\mathrm{~d}\left[\mathrm{~A}^{\star}\right]}{\mathrm{d} t}=k_{\mathrm{a}}[\mathrm{~A}]^{2} \tag{17F.1a}
\end{equation*}


The energized molecule ( $\mathrm{A}^{*}$ ) might lose its excess energy by collision with another molecule:


\begin{equation*}
\mathrm{A}+\mathrm{A}^{*} \rightarrow \mathrm{~A}+\mathrm{A} \quad \frac{\mathrm{~d}\left[\mathrm{~A}^{*}\right]}{\mathrm{d} t}=-k_{\mathrm{a}}^{\prime}[\mathrm{A}]\left[\mathrm{A}^{*}\right] \tag{17F.1b}
\end{equation*}


Alternatively, the excited molecule might shake itself apart or into a different arrangement of its atoms and form products $P$. That is, it might undergo the unimolecular decay


\begin{equation*}
\mathrm{A}^{*} \rightarrow \mathrm{P} \quad \frac{\mathrm{~d}\left[\mathrm{~A}^{*}\right]}{\mathrm{d} t}=-k_{\mathrm{b}}\left[\mathrm{~A}^{*}\right] \tag{17F.1c}
\end{equation*}


If the unimolecular step is the rate-determining step, the overall reaction will have first-order kinetics, as observed. This conclusion can be demonstrated explicitly by applying\\
the steady-state approximation to the net rate of formation of $A^{*}$ :


\begin{equation*}
\frac{\mathrm{d}\left[\mathrm{~A}^{*}\right]}{\mathrm{d} t}=k_{\mathrm{a}}[\mathrm{~A}]^{2}-k_{\mathrm{a}}^{\prime}[\mathrm{A}]\left[\mathrm{A}^{*}\right]-k_{\mathrm{b}}\left[\mathrm{~A}^{*}\right] \approx 0 \tag{17F.2}
\end{equation*}


Rearrangement of this equation gives


\begin{equation*}
\left[\mathrm{A}^{*}\right]=\frac{k_{\mathrm{a}}[\mathrm{~A}]^{2}}{k_{\mathrm{b}}+k_{\mathrm{a}}^{\prime}[\mathrm{A}]} \tag{17F.3}
\end{equation*}


so the rate law for the formation of P is


\begin{equation*}
\frac{\mathrm{d}[\mathrm{P}]}{\mathrm{d} t}=k_{\mathrm{b}}\left[\mathrm{~A}^{*}\right]=\frac{k_{\mathrm{a}} k_{\mathrm{b}}[\mathrm{~A}]^{2}}{k_{\mathrm{b}}+k_{\mathrm{a}}^{\prime}[\mathrm{A}]} \tag{17F.4}
\end{equation*}


At this stage the rate law is not first-order. However, if the rate of deactivation by $\left(\mathrm{A}^{*}, \mathrm{~A}\right)$ collisions is much greater than the rate of unimolecular decay, in the sense that $k_{\mathrm{a}}^{\prime}[\mathrm{A}]\left[\mathrm{A}^{*}\right] \gg k_{\mathrm{b}}\left[\mathrm{A}^{*}\right]$, or (after cancelling the $\left.\left[\mathrm{A}^{*}\right]\right), k_{\mathrm{a}}^{\prime}[\mathrm{A}] \gg k_{\mathrm{b}}$, then $k_{\mathrm{b}}$ can be neglected in the denominator of eqn 17F. 4 to obtain

$$
\frac{\mathrm{d}[\mathrm{P}]}{\mathrm{d} t}=k_{\mathrm{r}}[\mathrm{~A}] \quad \text { with } \quad k_{\mathrm{r}}=\frac{k_{\mathrm{a}} k_{\mathrm{b}}}{k_{\mathrm{a}}^{\prime}}
$$

Lindemann-Hinshelwood rate law\\
(17F.5)

Equation 17F. 5 is a first-order rate law, as required.\\
The Lindemann-Hinshelwood mechanism can be tested because it predicts that, as the concentration (and therefore the partial pressure) of A is reduced, the reaction should switch to overall second-order kinetics. Thus, when $k_{\mathrm{a}}^{\prime}[\mathrm{A}]\left[\mathrm{A}^{*}\right] \ll k_{\mathrm{b}}\left[\mathrm{A}^{*}\right]$ or (after cancelling the $\left.\left[\mathrm{A}^{*}\right]\right) k_{\mathrm{a}}^{\prime}[\mathrm{A}] \ll k_{\mathrm{b}}$, the rate law in eqn 17F. 4 becomes


\begin{equation*}
\frac{\mathrm{d}[\mathrm{P}]}{\mathrm{d} t}=k_{\mathrm{a}}[\mathrm{~A}]^{2} \tag{17F.6}
\end{equation*}


The physical reason for the change of order is that as the pressure is reduced the rate of the bimolecular process in which $\mathrm{A}^{*}$ loses its excess energy becomes negligible compared to the rate at which $\mathrm{A}^{\star}$ goes on to form products. The reaction mechanism is then a sequence of two steps, with the first step (which is bimolecular) being rate limiting. If the full rate law in eqn 17 F .4 is written as


\begin{equation*}
\frac{\mathrm{d}[\mathrm{P}]}{\mathrm{d} t}=k_{\mathrm{r}}[\mathrm{~A}] \quad \text { with } \quad k_{\mathrm{r}}=\frac{k_{\mathrm{a}} k_{\mathrm{b}}[\mathrm{~A}]}{k_{\mathrm{b}}+k_{\mathrm{a}}^{\prime}[\mathrm{A}]} \tag{17F.7}
\end{equation*}


then the expression for the effective rate constant, $k_{\mathrm{r}}$, can be rearranged (by inverting each side) to

\[
\frac{1}{k_{\mathrm{r}}}=\frac{k_{\mathrm{a}}^{\prime}}{k_{\mathrm{a}} k_{\mathrm{b}}}+\frac{1}{k_{\mathrm{a}}[\mathrm{~A}]} \quad \begin{align*}
& \text { Effective rate constant }  \tag{17F.8}\\
& \text { [Lindemann-Hinshelwood } \\
& \text { mechanism] }
\end{align*}
\]

Hence, a test of the theory is to plot $1 / k_{\mathrm{r}}$ against $1 /[\mathrm{A}]$ and expect a straight line. This behaviour is observed often at low concentrations but deviations are common at high concentrations. Topic 18A develops the description of the mechanism\\
further to take into account experimental results over a range of concentrations and pressures.

\subsection*{Example 17F. 1 Analysing data in terms of the Lindemann-Hinshelwood mechanism}
At 300 K the effective rate constant for a gaseous reaction $\mathrm{A} \rightarrow \mathrm{P}$, which has a Lindemann-Hinshelwood mechanism, is $k_{\mathrm{r}, 1}=2.50 \times 10^{-4} \mathrm{~s}^{-1}$ at $[\mathrm{A}]_{1}=5.21 \times 10^{-4} \mathrm{moldm}^{-3}$ and $k_{\mathrm{r}, 2}=$ $2.10 \times 10^{-5} \mathrm{~s}^{-1}$ at $[\mathrm{A}]_{2}=4.81 \times 10^{-6} \mathrm{moldm}^{-3}$. Calculate the rate constant $k_{\mathrm{a}}$ for the activation step in the mechanism.

Collect your thoughts Use eqn 17F. 8 to write an expression for the difference $1 / k_{\mathrm{r}, 2}-1 / k_{\mathrm{r}, 1}$, rearrange the expression for $k_{\mathrm{a}}$, and then insert the data.

The solution It follows from eqn 17F. 8 that

$$
\frac{1}{k_{\mathrm{r}, 2}}-\frac{1}{k_{\mathrm{r}, 1}}=\frac{1}{k_{\mathrm{a}}}\left(\frac{1}{[\mathrm{~A}]_{2}}-\frac{1}{[\mathrm{~A}]_{1}}\right)
$$

and so

$$
\begin{aligned}
k_{\mathrm{a}} & =\frac{1 /[\mathrm{A}]_{2}-1 /[\mathrm{A}]_{1}}{1 / k_{\mathrm{r}, 2}-1 / k_{\mathrm{r}, 1}} \\
& =\frac{1 /\left(4.81 \times 10^{-6} \mathrm{moldm}^{-3}\right)-1 /\left(5.21 \times 10^{-4} \mathrm{moldm}^{-3}\right)}{1 /\left(2.10 \times 10^{-5} \mathrm{~s}^{-1}\right)-1 /\left(2.50 \times 10^{-4} \mathrm{~s}^{-1}\right)} \\
& =4.72 \mathrm{dm}^{3} \mathrm{~mol}^{-1} \mathrm{~s}^{-1}
\end{aligned}
$$

Self-test 17F. 1 The effective rate constants for a gaseous reaction $\mathrm{A} \rightarrow \mathrm{P}$, which has a Lindemann-Hinshelwood mechanism, are $1.70 \times 10^{-3} \mathrm{~s}^{-1}$ and $2.20 \times 10^{-4} \mathrm{~s}^{-1}$ at $[\mathrm{A}]=4.37 \times 10^{-4}$ $\mathrm{mol} \mathrm{dm}{ }^{-3}$ and $1.00 \times 10^{-5} \mathrm{~mol} \mathrm{dm}^{-3}$, respectively. Calculate the rate constant for the activation step in the mechanism.\\
\includegraphics[max width=\textwidth, center]{2025_02_27_d4b95d9718f0b9e2e6b9g-786}

\subsection*{17F. 2 Polymerization kinetics}
There are two major classes of polymerization processes and in each one the average molar mass of the product varies with time in a distinctive way. In stepwise polymerization any two monomers present in the reaction mixture can link together at any time and growth of the polymer is not confined to chains that are already forming (Fig. 17F.2). As a result, monomers are consumed early in the reaction and, as will be seen, the average molar mass of the product grows linearly with time. In chain polymerization a monomer, $M$, attacks another monomer, links to it, then that unit attacks another monomer, and so on. The monomer is used up as it becomes linked to the growing chains (Fig. 17F.3). Polymers built from numerous monomers are formed rapidly and only the yield, not the average molar mass, of the polymer is increased by allowing long reaction times.\\
\includegraphics[max width=\textwidth, center]{2025_02_27_d4b95d9718f0b9e2e6b9g-787(1)}

Figure 17F. 2 In stepwise polymerization, growth can start at any pair of monomers (in green), and so new chains (in red) begin to form throughout the reaction.\\
\includegraphics[max width=\textwidth, center]{2025_02_27_d4b95d9718f0b9e2e6b9g-787}

Figure 17F.3 The process of chain polymerization. Chains (in red) grow as each chain acquires additional monomers (in green).

\subsection*{(a) Stepwise polymerization}
Stepwise polymerization commonly proceeds by a condensation reaction, in which a small molecule (typically $\mathrm{H}_{2} \mathrm{O}$ ) is eliminated in each step. Stepwise polymerization is the mechanism of production of polyamides, as in the formation of nylon-66:

$$
\begin{aligned}
& \mathrm{H}_{2} \mathrm{~N}\left(\mathrm{CH}_{2}\right)_{6} \mathrm{NH}_{2}+\mathrm{HOOC}\left(\mathrm{CH}_{2}\right)_{4} \mathrm{COOH} \rightarrow \\
& \mathrm{H}_{2} \mathrm{~N}\left(\mathrm{CH}_{2}\right)_{6} \mathrm{NHCO}\left(\mathrm{CH}_{2}\right)_{4} \mathrm{COOH}+\mathrm{H}_{2} \mathrm{O}
\end{aligned}
$$

continuing to

$$
\rightarrow \mathrm{H}-\left[\mathrm{HN}\left(\mathrm{CH}_{2}\right)_{6} \mathrm{NHCO}\left(\mathrm{CH}_{2}\right)_{4} \mathrm{CO}\right]_{n}-\mathrm{OH}
$$

Polyesters and polyurethanes are formed similarly (the latter without elimination). A polyester, for example, can be regarded as the outcome of the stepwise condensation of a hydroxyacid $\mathrm{HO}-\mathrm{R}-\mathrm{COOH}$. Consider the formation of a polyester from such a monomer. Its progress can be measured in terms of the concentration of the -COOH groups in the sample (denoted A), because these groups gradually disappear as the condensation proceeds. Because the condensation reaction can occur between molecules containing any number of monomer units, chains of many different lengths can grow in the reaction mixture.

In the absence of a catalyst, condensation is expected to be overall second-order in the concentration of the -OH and -COOH (or A) groups, so


\begin{equation*}
\frac{\mathrm{d}[\mathrm{~A}]}{\mathrm{d} t}=-k_{\mathrm{r}}[\mathrm{OH}][\mathrm{A}] \tag{17F.9a}
\end{equation*}


However, because there is one -OH group for each -COOH group, this equation is the same as


\begin{equation*}
\frac{\mathrm{d}[\mathrm{~A}]}{\mathrm{d} t}=-k_{\mathrm{r}}[\mathrm{~A}]^{2} \tag{17F.9b}
\end{equation*}


If the rate constant for the condensation is independent of the chain length, $k_{\mathrm{r}}$ remains constant throughout the reaction. The solution of this rate law is then given by eqn 17B. 4 b , and is


\begin{equation*}
[\mathrm{A}]=\frac{[\mathrm{A}]_{0}}{\left.1+k_{\mathrm{r}} t \mathrm{~A}\right]_{0}} \tag{17F.10}
\end{equation*}


The fraction, $p$, of -COOH groups that have condensed at time $t$ is therefore

$$
p=\frac{[\mathrm{A}]_{0}-[\mathrm{A}]}{[\mathrm{A}]_{0}}=\frac{\begin{array}{ll}
{[\mathrm{A}]=[\mathrm{A}]_{0} /\left(1+k_{\mathrm{r}} t[\mathrm{~A}]_{0}\right)} \\
\frac{k_{\mathrm{r}} t[\mathrm{~A}]_{0}}{1+k_{\mathrm{r}} t[\mathrm{~A}]_{0}} & \begin{array}{l}
\text { Fraction of condensed groups } \\
\text { [stepwise polymerization }]
\end{array}
\end{array} \text { later }}{}
$$

(17F.11)

The degree of polymerization, the average number of monomer residues per polymer molecule, can now be calculated. This quantity is the ratio of the initial concentration of $\mathrm{A},[\mathrm{A}]_{0}$, to the concentration of end groups, [A], at the time of interest, because there is one A group per polymer molecule. For example, if there were initially 1000 A groups and there are now only 10 , the average length of each polymer must be 100 units. Because [A] can be expressed in terms of $p$ (the first part of eqn 17F.11), the average number of monomers per polymer molecule, $\langle N\rangle$, is

$$
\langle N\rangle=\frac{[\mathrm{A}]_{0}}{[\mathrm{~A}]}=\frac{1}{1-p} \quad \begin{aligned}
& \text { Degree of polymerization } \\
& \text { [stepwise polymerization] }
\end{aligned}
$$

(17F.12a)\\
\includegraphics[max width=\textwidth, center]{2025_02_27_d4b95d9718f0b9e2e6b9g-788}

Figure 17F. 4 The average chain length of a polymer as a function of the fraction of reacted monomers, $p$. Note that $p$ must be very close to 1 for the chains to be long.

This result is illustrated in Fig. 17F.4. When $p$ is expressed in terms of the rate constant $k_{\mathrm{r}}$ (the second part of eqn 17F.11), the result is

$$
\langle N\rangle=1+k_{\mathrm{r}} t[\mathrm{~A}]_{0}
$$

Degree of polymerization in terms of the rate constant [stepwise polymerization]

The average length grows linearly with time. Therefore, the longer a stepwise polymerization proceeds, the higher the average molar mass of the product.

\subsection*{Brief illustration 17F. 1}
Consider a polymer formed by a stepwise process with $k_{\mathrm{r}}=1.00 \mathrm{dm}^{3} \mathrm{~mol}^{-1} \mathrm{~s}^{-1}$ and an initial monomer concentration of $[\mathrm{A}]_{0}=4.00 \times 10^{-3} \mathrm{moldm}^{-3}$. From eqn 17F.12b, the degree of polymerization at $t=1.5 \times 10^{4} \mathrm{~s}$ is

$$
\begin{aligned}
\langle N\rangle= & 1+\left(1.00 \mathrm{dm}^{3} \mathrm{~mol}^{-1} \mathrm{~s}^{-1}\right) \times\left(1.5 \times 10^{4} \mathrm{~s}\right) \\
& \times\left(4.00 \times 10^{-3} \mathrm{moldm}^{-3}\right)=61
\end{aligned}
$$

From eqn 17F.12a, the fraction condensed, $p$, is

$$
p=\frac{\langle N\rangle-1}{\langle N\rangle}=\frac{61-1}{61}=0.98
$$

\subsection*{(b) Chain polymerization}
Many gas-phase reactions and liquid-phase polymerization reactions are chain reactions. In a chain reaction, a reaction intermediate produced in one step generates an intermediate in a subsequent step, then that intermediate generates another intermediate, and so on. The intermediates in a chain reaction are called chain carriers. In a radical chain reaction the chain carriers are radicals (species with unpaired electrons).

Chain polymerization occurs by addition of monomers to a growing polymer, often by a radical chain process. It results in the rapid growth of an individual polymer chain for each monomer available to react. Examples include the addition polymerizations of ethene, methyl methacrylate, and styrene, as in

$$
-\mathrm{CH}_{2} \dot{\mathrm{C}} \mathrm{HX}+\mathrm{CH}_{2}=\mathrm{CHX} \rightarrow-\mathrm{CH}_{2} \mathrm{CHXCH}_{2} \dot{\mathrm{C}} \mathrm{HX}
$$

and subsequent reactions.\\
There are several characteristic steps in a chain reaction. Initiation, the formation of active radicals, may be written as

$$
\begin{aligned}
& \mathrm{In} \rightarrow \mathrm{R} \cdot+\mathrm{R} \cdot \\
& \mathrm{M}+\mathrm{R} \cdot \rightarrow \cdot \mathrm{M}_{1}
\end{aligned}
$$

where In is the initiator, $\mathrm{R} \cdot$ the radical that In forms, and $\cdot \mathrm{M}_{1}$ is a monomer radical. In this reaction a radical is produced, but in some polymerizations the initiation step leads to the formation of an ionic chain carrier. Initiation is followed by propagation, the continuation of the chain reaction:

$$
\begin{gathered}
\mathrm{M}+\cdot \mathrm{M}_{1} \rightarrow \cdot \mathrm{M}_{2} \\
\mathrm{M}+\cdot \mathrm{M}_{2} \rightarrow \cdot \mathrm{M}_{3} \\
\vdots \\
\mathrm{M}+\cdot \mathrm{M}_{n-1} \rightarrow \cdot \mathrm{M}_{n}
\end{gathered}
$$

where $M_{n}$ is a polymer consisting of $n$ monomer units. Polymerization may terminate in a number of ways. For example,

$$
\begin{array}{ll}
\text { mutual termination: } & \cdot \mathrm{M}_{n}+\cdot \mathrm{M}_{m} \rightarrow \mathrm{M}_{n+m} \\
\text { disproportionation: } & \cdot \mathrm{HM}_{n}+\cdot \mathrm{M}_{m} \rightarrow \mathrm{M}_{n}+\mathrm{HM}_{m} \\
\text { chain transfer: } & \mathrm{M}+\cdot \mathrm{M}_{n} \rightarrow \cdot \mathrm{M}+\mathrm{M}_{n}
\end{array}
$$

In mutual termination two growing radical chains combine. In termination by disproportionation a hydrogen atom transfers from one chain to another, corresponding to the oxidation of the donor and the reduction of the acceptor. In chain transfer, a new chain initiates at the expense of the one currently growing. As can be suspected, the mechanism is complicated, but can be explored by using the steady-state approximation.

\subsection*{How is that done? 17F. 1 Deriving an expression for the}
 rate of chain polymerizationThe kinetic analysis of chain polymerization must take into account initiation, propagation, and termination.

Step 1 Write an expression for the rate of initiation of the process If the initiation step is

$$
\begin{aligned}
& \text { In } \rightarrow \mathrm{R} \cdot+\mathrm{R} \cdot \quad v_{\mathrm{i}}=k_{\mathrm{i}}[\mathrm{In}] \\
& \mathrm{M}+\mathrm{R} \cdot \rightarrow \cdot \mathrm{M}_{1}
\end{aligned}
$$

If the rate constants of the chain propagation steps are large enough, the first of these two steps is rate-determining for the overall polymerization process and the rate of initiation is equal to $v_{\mathrm{i}}$.

\subsection*{Step 2 Write an expression for propagation}
If the rate of propagation is independent of chain size for sufficiently large chains, then the rate of propagation, $v_{\mathrm{p}}$, may be written

$$
v_{\mathrm{p}}=k_{\mathrm{p}}[\mathrm{M}][\cdot \mathrm{M}]
$$

where $\cdot \mathrm{M}$ stands for a polymer of any length. It follows from the remark in Step 1 that

$$
\left(\frac{\mathrm{d}[\cdot \mathrm{M}]}{\mathrm{d} t}\right)_{\text {production }}=2 f k_{\mathrm{i}}[\mathrm{In}]
$$

where $f$ is the fraction of radicals R. that successfully initiate a chain. The factor 2 recognizes that two radicals are formed in each initiation step.

\subsection*{Step 3 Consider the termination of the process}
For the present analysis, suppose that only mutual termination occurs. If the rate of termination is assumed to be independent of the length of the chain, the rate law for termination is

$$
v_{\mathrm{t}}=k_{\mathrm{t}}[\cdot \mathrm{M}]^{2}
$$

and the rate of change of radical concentration by this depletion process is

$$
\left(\frac{\mathrm{d}[\cdot \mathrm{M}]}{\mathrm{d} t}\right)_{\text {depletion }}=-2 k_{\mathrm{t}}[\cdot \mathrm{M}]^{2}
$$

In this case, the factor 2 recognizes that two radicals are removed in each depletion step.

Step 4 Apply the steady-state approximation\\
The net rate of formation of $\cdot \mathrm{M}$ is

$$
\frac{\mathrm{d}[\cdot \mathrm{M}]}{\mathrm{d} t}=\overbrace{2 f k_{\mathrm{i}}[\mathrm{In}]}^{\text {production }}-\overbrace{2 k_{\mathrm{t}}[\cdot \mathrm{M}]^{2}}^{\text {depletion }} \approx 0
$$

The steady-state concentration of radical chains is therefore

$$
[\cdot \mathrm{M}]=\left(\frac{f k_{\mathrm{i}}}{k_{\mathrm{t}}}\right)^{1 / 2}[\mathrm{In}]^{1 / 2}
$$

In Step 2 it is established that the overall rate of polymerization is equal to the rate of propagation, which is given by $v_{\mathrm{p}}=k_{\mathrm{p}}[\mathrm{M}][\cdot \mathrm{M}]$. The steady-state expression for $[\cdot \mathrm{M}]$ can now be inserted into this expression to give

$$
v_{\mathrm{p}}=k_{\mathrm{p}}[\cdot \mathrm{M}][\mathrm{M}]=k_{\mathrm{p}}\left(\frac{f k_{\mathrm{i}}}{k_{\mathrm{t}}}\right)^{1 / 2}[\mathrm{In}]^{1 / 2}[\mathrm{M}]
$$

The overall rate of polymerization is therefore proportional to the square root of the concentration of the initiator, In, and is given by

$$
-v=k_{\mathrm{r}}[\mathrm{In}]^{1 / 2}[\mathrm{M}] \text { with } \quad k_{\mathrm{r}}=k_{\mathrm{p}}\left(\frac{f k_{\mathrm{i}}}{k_{\mathrm{t}}}\right)^{1 / 2} \left\lvert\, \begin{aligned}
& \text { Rate of polymeri- } \\
& \begin{array}{l}
\text { zation } \\
\text { [chain polymeri- } \\
\text { zation] }
\end{array}
\end{aligned}\right.
$$

The kinetic chain length, $\lambda$, is the ratio of the number of monomer units consumed to the number of radicals produced in the initiation step:

$$
\lambda=\frac{\text { number of monomer units consumed }}{\text { number of radicals produced }}
$$

Kinetic chain length\\[0pt]
[definition]\\
(17F.14a)

The kinetic chain length can be imagined as the average number of molecules in a chain produced by one initiating radical. The kinetic chain length can be expressed in terms of the rate expressions above. To do so, recognize that monomers are consumed at the rate that chains propagate. Then,

$$
\lambda=\frac{\text { rate of propagation of chains }}{\text { rate of production of radicals }}
$$

Kinetic chain length in terms of reaction rates\\
(17F.14b)

In applying the steady-state approximation, the rate of production of radicals is set equal to the termination rate (Step 4 in the discussion above). Therefore, the kinetic chain length may be written as

$$
\lambda=\frac{k_{\mathrm{p}}[\cdot \mathrm{M}][\mathrm{M}]}{2 k_{\mathrm{t}}[\cdot \mathrm{M}]^{2}}=\frac{k_{\mathrm{p}}[\mathrm{M}]}{2 k_{\mathrm{t}}[\cdot \mathrm{M}]}
$$

The steady-state expression for $[\cdot \mathrm{M}],[\cdot \mathrm{M}]=\left(f k_{\mathrm{i}} / k_{\mathrm{t}}\right)^{1 / 2}[\mathrm{In}]^{1 / 2}$, is substituted for the radical concentration, to obtain

$$
\begin{aligned}
& \lambda=k_{\mathrm{r}}[\mathrm{M}][\mathrm{In}]^{-1 / 2} \text { with } k_{\mathrm{r}}=k_{\mathrm{p}}\left(4 f k_{\mathrm{i}} k_{\mathrm{t}}\right)^{-1 / 2} \\
& \begin{array}{l}
\text { Kinetic chain length } \\
\text { [chain polymerization] }
\end{array}
\end{aligned}
$$

(17F.14c)

In mutual termination, the average number of monomers in a polymer molecule, $\langle N\rangle$, produced by the reaction is the sum of the numbers of monomers in the two combining polymer chains. The average number of units in each chain is $\lambda$. Therefore,

$$
\langle N\rangle=2 \lambda=2 k_{\mathrm{r}}[\mathrm{M}][\operatorname{In}]^{-1 / 2} \quad \begin{aligned}
& \text { Degree of polymerization } \\
& \text { [chain polymerization] }
\end{aligned}
$$

(17F.15)\\
with $k_{\mathrm{r}}$ given in eqn 17F.14c. That is, the slower the initiation of the chain (the smaller the initiator concentration and the smaller the initiation rate constant), the greater is the kinetic chain length, and therefore the higher is the average molar mass of the polymer.

\subsection*{17F. 3 Enzyme-catalysed reactions}
A catalyst is a substance that accelerates a reaction but undergoes no net chemical change (Topic 17D): the catalyst lowers the activation energy of the reaction by providing an alternative path to that of the uncatalysed reaction (Fig. 17F.5). Enzymes, which are homogeneous biological catalysts, are very specific and can have a dramatic effect on the reactions they control. For example, the enzyme catalase accelerates the reaction it catalyses by a factor of $10^{12}$ at 298 K .

Enzymes contain an active site, which is responsible for binding the substrates, the reactants, and processing them into products. As is true of any catalyst, the active site returns to its original state after the products are released. Many enzymes consist primarily of proteins, some featuring organic or inorganic co-factors in their active sites. However, certain RNA molecules can also be biological catalysts, forming ribozymes.

The structure of the active site is specific to the reaction that it catalyses, with groups in the substrate attached to groups in the active site primarily by hydrogen bonding, electrostatic forces, and van der Waals interactions. Figure 17F. 6 shows two models that explain the binding of a substrate to the active site of an enzyme. In the lock-and-key model, the active site and substrate have complementary three-dimensional structures and dock without the need for major structural change. Experimental evidence however favours the induced\\
\includegraphics[max width=\textwidth, center]{2025_02_27_d4b95d9718f0b9e2e6b9g-790}

Figure 17F. 5 A catalyst provides a different path with a lower activation energy. The result is an increase in the rates of the forward (and reverse) reaction.\\
\includegraphics[max width=\textwidth, center]{2025_02_27_d4b95d9718f0b9e2e6b9g-790(1)}

Figure 17F. 6 Two models that explain the binding of a substrate to the active site of an enzyme. In the lock-and-key model, the active site and substrate have complementary threedimensional structures and dock without the need for major atomic rearrangements. In the induced fit model, binding of the substrate induces a conformational change in the active site. The substrate fits well in the active site after the conformational change has taken place.\\
fit model, in which binding of the substrate induces a conformational change in the active site. Only after the change does the substrate fit snugly in the active site.

Experimental studies of enzyme kinetics are typically conducted by monitoring the initial rate of product formation in a solution in which the enzyme is present at very low concentration. Indeed, enzymes are such efficient catalysts that significant accelerations may be observed even when their concentration is more than three orders of magnitude smaller than that of the substrate.

The principal features of many enzyme-catalysed reactions are as follows:

\begin{itemize}
  \item For a given initial concentration of substrate, $[\mathrm{S}]_{0}$, the initial rate of product formation is proportional to the total concentration of enzyme, $[E]_{0}$.
  \item For a given $[\mathrm{E}]_{0}$ and low values of $[\mathrm{S}]_{0}$, the rate of product formation is proportional to $[\mathrm{S}]_{0}$.
  \item For a given $[\mathrm{E}]_{0}$ and high values of $[\mathrm{S}]_{0}$, the rate of product formation becomes independent of $[\mathrm{S}]_{0}$, reaching a maximum value known as the maximum velocity, $v_{\text {max }}$.
\end{itemize}

The Michaelis-Menten mechanism accounts for these features. According to this mechanism, an enzyme-substrate complex is formed in the first step and then either the substrate is released unchanged or, after modification, released as the product:

$$
\begin{array}{ll}
\mathrm{E}+\mathrm{S} \underset{k_{\mathrm{a}}^{\prime}}{\stackrel{k_{\mathrm{a}}}{\rightleftarrows}} \mathrm{ES} & \begin{array}{l}
\text { Michaelis-Menten } \\
\text { mechanism }
\end{array} \\
\mathrm{ES} \xrightarrow{k_{\mathrm{b}}} \mathrm{P}+\mathrm{E} &
\end{array}
$$

Again, the mechanism may be analysed by using the steadystate approximation.

\subsection*{How is that done? 17F. 2 Deriving the Michaelis-Menten equation}
The rate of product formation according to the MichaelisMenten mechanism is

$$
v=k_{\mathrm{b}}[\mathrm{ES}]
$$

so the strategy is centred on finding an expression for the concentration of the intermediate, ES.

Step 1 Apply the steady-state approximation\\
The concentration of the enzyme-substrate complex is found by invoking the steady-state approximation and writing

$$
\frac{\mathrm{d}[\mathrm{ES}]}{\mathrm{d} t}=k_{\mathrm{a}}[\mathrm{E}][\mathrm{S}]-k_{\mathrm{a}}^{\prime}[\mathrm{ES}]-k_{\mathrm{b}}[\mathrm{ES}] \approx 0
$$

It follows that

$$
[\mathrm{ES}]=\frac{k_{\mathrm{a}}[\mathrm{E}][\mathrm{S}]}{k_{\mathrm{a}}^{\prime}+k_{\mathrm{b}}}
$$

where $[\mathrm{E}]$ and $[\mathrm{S}]$ are the concentrations of free enzyme and substrate, respectively.

Step 2: Simplify the expression for [ES]\\
Now define the Michaelis constant as

$$
K_{\mathrm{M}}=\frac{k_{\mathrm{a}}^{\prime}+k_{\mathrm{b}}}{k_{\mathrm{a}}}
$$

(Note that this constant has the dimensions of molar concentration.) To express the rate law in terms of the total concentrations of enzyme and the initial concentration of substrate first added, note that the total concentration of the enzyme is $[\mathrm{E}]_{0}=[\mathrm{E}]+[\mathrm{ES}]$. It follows that $[\mathrm{E}]=[\mathrm{E}]_{0}-[\mathrm{ES}]$. This expression for [ E$]$ is inserted into the steady-state expression for [ES] above to give

$$
[\mathrm{ES}]=\frac{\left([\mathrm{E}]_{0}-[\mathrm{ES}]\right)[\mathrm{S}]}{K_{\mathrm{M}}}
$$

Because the substrate is typically in large excess relative to the enzyme, the free substrate concentration is approximately equal to the initial substrate concentration: $[\mathrm{S}] \approx[\mathrm{S}]_{0}$. The solution of the resulting expression for [ES] is

$$
[\mathrm{ES}]=\frac{[\mathrm{S}]_{0}[\mathrm{E}]_{0}}{K_{\mathrm{M}}+[\mathrm{S}]_{0}}=\frac{[\mathrm{E}]_{0}}{1+K_{\mathrm{M}} /[\mathrm{S}]_{0}}
$$

Step 3 Write an expression for the rate law\\[0pt]
The expression for [ES] can now be substituted into $v=k_{\mathrm{b}}[\mathrm{ES}]$, to give the Michaelis-Menten equation\\
$\left.-v=\frac{k_{\mathrm{b}}[\mathrm{E}]_{0}}{1+K_{\mathrm{M}} /[\mathrm{S}]_{0}} \right\rvert\, \quad$ (17F.16)\\
\includegraphics[max width=\textwidth, center]{2025_02_27_d4b95d9718f0b9e2e6b9g-791}

Figure 17F. 7 The variation of the rate of an enzyme-catalysed reaction with substrate concentration. The approach to a maximum rate, $v_{\text {max }}$ for large $[\mathrm{S}]_{0}$ is explained by the MichaelisMenten mechanism.

Equation 17F. 16 predicts that, in accord with experimental observations (Fig. 17F.7):

\begin{itemize}
  \item When $[\mathrm{S}]_{0} \ll K_{\mathrm{M}}$, the rate is proportional to $[\mathrm{S}]_{0}$ :
\end{itemize}


\begin{equation*}
v=\frac{k_{\mathrm{b}}}{K_{\mathrm{M}}}[\mathrm{~S}]_{0}[\mathrm{E}]_{0} \tag{17F.17a}
\end{equation*}


\begin{itemize}
  \item When $[\mathrm{S}]_{0} \gg K_{\mathrm{M}}$, the rate reaches its maximum value and is independent of $[\mathrm{S}]_{0}$ :
\end{itemize}


\begin{equation*}
v_{\text {max }}=k_{\mathrm{b}}[\mathrm{E}]_{0} \tag{17F.17b}
\end{equation*}


Substitution of this definition of $v_{\text {max }}$ into eqn 17F. 16 gives


\begin{equation*}
v=\frac{v_{\max }}{1+K_{\mathrm{M}} /[\mathrm{S}]_{0}} \tag{17F.18a}
\end{equation*}


which can be rearranged into a form suitable for data analysis by linear regression by taking reciprocals of both sides:

$$
\frac{1}{v}=\frac{1}{v_{\max }}+\left(\frac{K_{\mathrm{M}}}{v_{\max }}\right) \frac{1}{[\mathrm{~S}]_{0}} \quad \quad \text { Lineweaver-Burk plot }
$$

(17F.18b)

A Lineweaver-Burk plot is a plot of $1 / v$ against $1 /[\mathrm{S}]_{0}$. According to eqn 17F.18b, it should yield a straight line with slope of $K_{\mathrm{M}} / v_{\text {max }}$, a $y$-intercept at $1 / v_{\text {max }}$, and an $x$-intercept at $-1 / K_{M}$ (Fig. 17F.8). The value of $K_{\mathrm{M}}$ can also be obtained from the ratio of the slope to the $y$-intercept. The value of $k_{\mathrm{b}}$ is then calculated from the $y$-intercept and eqn 17F.17b. However, the plot cannot give the individual rate constants $k_{\mathrm{a}}$ and $k_{\mathrm{a}}^{\prime}$ that appear in the definition of $K_{\mathrm{M}}$. The stopped-flow technique described in Topic 17A can give the additional data needed, because the rate of formation of the enzyme-substrate complex can be found by monitoring the concentration after mixing the enzyme and substrate. This procedure gives a value for $k_{\mathrm{a}}$, and $k_{\mathrm{a}}^{\prime}$ is then found by combining this result with the values of $k_{\mathrm{b}}$ and $K_{\mathrm{M}}$.\\
\includegraphics[max width=\textwidth, center]{2025_02_27_d4b95d9718f0b9e2e6b9g-792}

Figure 17F. 8 The structure of a Lineweaver-Burk plot for the analysis of an enzyme-catalysed reaction that proceeds by a Michaelis-Menten mechanism and the significance of the intercepts and the slope.

\subsection*{Example 17F.2 Analysing data by using a LineweaverBurk plot}
The enzyme carbonic anhydrase catalyses the hydration of $\mathrm{CO}_{2}$ in red blood cells to give hydrogencarbonate (bicarbonate) ion: $\mathrm{CO}_{2}(\mathrm{~g})+\mathrm{H}_{2} \mathrm{O}(\mathrm{l}) \rightarrow \mathrm{HCO}_{3}^{-}(\mathrm{aq})+\mathrm{H}^{+}(\mathrm{aq})$. The following data were obtained for the reaction at $\mathrm{pH}=7.1,273.5 \mathrm{~K}$, and an enzyme concentration of $2.3 \mathrm{nmol} \mathrm{dm}^{-3}$ :

\begin{center}
\begin{tabular}{lllll}
$\left[\mathrm{CO}_{2}\right] /\left(\mathrm{mmol} \mathrm{dm}^{-3}\right)$ & 1.25 & 2.5 & 5 & 20 \\
$\nu /\left(\mathrm{mmol} \mathrm{dm}^{-3} \mathrm{~s}^{-1}\right)$ & $2.78 \times 10^{-2}$ & $5.00 \times 10^{-2}$ & $8.33 \times 10^{-2}$ & $1.67 \times 10^{-1}$ \\
\end{tabular}
\end{center}

Determine the maximum velocity and the Michaelis constant for the reaction.

Collect your thoughts Prepare a Lineweaver-Burk plot as explained in the text and determine the values of $K_{\mathrm{M}}$ and $v_{\text {max }}$ by linear regression analysis.

The solution Draw up the following table:

\begin{center}
\begin{tabular}{lcccl}
$1 /\left(\left[\mathrm{CO}_{2}\right] /\left(\mathrm{mmol} \mathrm{dm}^{-3}\right)\right)$ & 0.800 & 0.400 & 0.200 & 0.0500 \\
$1 /\left(v /\left(\mathrm{mmol} \mathrm{dm}^{-3} \mathrm{~s}^{-1}\right)\right)$ & 36.0 & 20.0 & 12.0 & 6.0 \\
\end{tabular}
\end{center}

Figure 17F. 9 shows the Lineweaver-Burk plot for the data. The slope is 40.0 and the $y$-intercept is 4.00 . Hence,

$$
v_{\max } /\left(\mathrm{mmoldm}^{-3} \mathrm{~s}^{-1}\right)=\frac{1}{\text { intercept }}=\frac{1}{4.00}=0.250
$$

$$
\begin{aligned}
& \text { so } v_{\max }=0.250 \mathrm{mmol} \mathrm{dm}^{-3} \mathrm{~s}^{-1} . \\
& \qquad K_{\mathrm{M}} /\left(\mathrm{mmoldm}^{-3}\right)=\frac{\text { slope }}{\text { intercept }}=\frac{40.00}{4.00}=10.0
\end{aligned}
$$

and so $K_{M}=10.0 \mathrm{mmoldm}^{-3}$.\\
\includegraphics[max width=\textwidth, center]{2025_02_27_d4b95d9718f0b9e2e6b9g-792(1)}

Figure 17F.9 The Lineweaver-Burk plot of the data for Example 17F.2.

A note on good practice The slope and the intercept are unit-less: all graphs should be plotted as pure numbers.

Self-test 17F. 2 The enzyme $\alpha$-chymotrypsin is secreted in the pancreas of mammals and cleaves peptide bonds between certain amino acids. Several solutions containing the small peptide $N$-glutaryl-L-phenylalanine- $p$-nitroanilide at different concentrations were prepared and the same small amount of $\alpha$-chymotrypsin was added to each one. The following data were obtained on the initial rates of the formation of product:

\begin{center}
\begin{tabular}{lllllll}
$[\mathrm{S}] /\left(\mathrm{mmoldm}^{-3}\right)$ & 0.334 & 0.450 & 0.667 & 1.00 & 1.33 & 1.67 \\
$\nu /\left(\mathrm{mmol} \mathrm{dm}^{-3} \mathrm{~s}^{-1}\right)$ & 0.150 & 0.199 & 0.285 & 0.406 & 0.516 & 0.619 \\
\end{tabular}
\end{center}

Determine the maximum velocity and the Michaelis constant for the reaction.

The action of an enzyme may be partially suppressed by the presence of a foreign substance, which is called an inhibitor. An inhibitor may be a poison that has been administered to the organism, or it may be a substance that is naturally present in a cell and involved in its regulatory mechanism. An inhibitor typically works by blocking the active site or attaching elsewhere to the enzyme and forcing a change in geometry at the site so that it can no longer accommodate the substrate. ${ }^{1}$

\footnotetext{${ }^{1}$ The use of kinetic criteria to distinguish different types of inhibition is described in our Physical chemistry for the life sciences (2012).
}\subsection*{Checklist of concepts}
$\square$ 1. The Lindemann-Hinshelwood mechanism of 'unimolecular' reactions accounts for the first-order kinetics of some gas-phase reactions.2. In stepwise polymerization any two monomers in the reaction mixture can link together at any time.3. The longer a stepwise polymerization proceeds, the higher the average molar mass of the product.4. In chain polymerization an activated monomer attacks another monomer and links to it; the slower the initiation of the chain, the higher the average molar mass of the polymer.\\
$\square$ 5. The kinetic chain length is the ratio of the number of monomer units consumed to the number of radicals produced in the initiation step.6. Enzymes are homogeneous biological catalysts.\\
$\square$ 7. The Michaelis-Menten mechanism of enzyme kinetics accounts for the dependence of rate on the concentrations of the substrate and the enzyme.\\
$\square$ 8. A Lineweaver-Burk plot is used to determine the parameters that occur in the Michaelis-Menten mechanism.

\subsection*{Checklist of equations}
\begin{center}
\begin{tabular}{llll}
\hline
Property & Equation & Comment & Equation number \\
\hline
Lindemann-Hinshelwood rate law & $\mathrm{d}[\mathrm{P}] / \mathrm{d} t=k_{\mathrm{r}}[\mathrm{A}] \quad$ with $\quad k_{\mathrm{r}}=k_{\mathrm{a}} k_{\mathrm{b}} / k_{\mathrm{a}}^{\prime}$ & $k_{\mathrm{a}}^{\prime}[\mathrm{A}] \gg k_{\mathrm{b}}$ & 17 F .7 \\
Effective rate constant & $1 / k_{\mathrm{r}}=k_{\mathrm{a}}^{\prime} / k_{\mathrm{a}} k_{\mathrm{b}}+1 / k_{\mathrm{a}}[\mathrm{A}]$ & Lindemann-Hinshelwood mechanism & 17 F .8 \\
Fraction of condensed groups & $p=k_{\mathrm{r}} t[\mathrm{~A}]_{0} /\left(1+k_{\mathrm{r}}[\mathrm{A}]_{0}\right)$ & Stepwise polymerization & 17 F .11 \\
Degree of polymerization & $\langle N\rangle=1 /(1-p)=1+k_{\mathrm{r}} t[\mathrm{~A}]_{0}$ & Stepwise polymerization & 17 F .12 \\
Rate of polymerization & $v=k_{\mathrm{r}}[\mathrm{In}]^{1 / 2}[\mathrm{M}]$ & Chain polymerization & 17 F .13 \\
Kinetic chain length & $\lambda=k_{\mathrm{r}}[\mathrm{M}][\mathrm{In}]^{-1 / 2}, k_{\mathrm{r}}=k_{\mathrm{p}}\left(4 f k_{\mathrm{i}} k_{\mathrm{t}}\right)^{-1 / 2}$ & Chain polymerization & 17 F .14 c \\
Degree of polymerization & $\langle N\rangle=2 k_{\mathrm{r}}[\mathrm{M}][\mathrm{In}]^{-1 / 2}$ & Chain polymerization & 17 F .15 \\
Michaelis-Menten equation & $v=v_{\max } /\left(1+K_{\mathrm{M}} /[\mathrm{S}]_{0}\right)$ & 17 F .18 a & 17 F .18 b \\
Lineweaver-Burk plot & $1 / v=1 / v_{\max }+\left(K_{\mathrm{M}} / v_{\max }\right)\left(1 /[\mathrm{S}]_{0}\right)$ &  &  \\
\hline
\end{tabular}
\end{center}

\section*{TOPIC 17G Photochemistry}
\subsection*{> Why do you need to know this material?}
Many chemical and biological processes, including photosynthesis and vision, are initiated by the absorption of electromagnetic radiation, so you need to know how to include its effect in rate laws. The quantitative analysis of these processes provides insight into their mechanisms.

\subsection*{- What is the key idea?}
The mechanisms of many photochemical reactions lead to relatively simple rate laws that yield rate constants and quantitative measures of the efficiency with which radiant energy induces reactions.

\subsection*{> What do you need to know already?}
You need to be familiar with the concepts of singlet and triplet states (Topic 11F), modes of radiative decay (fluorescence and phosphorescence, Topic 11G), concepts of electronic spectroscopy (Topic 11F), and the formulation of a rate law from a proposed mechanism (Topic 17E).

Photochemical processes are initiated by the absorption of electromagnetic radiation. Among the most important of these processes are the ones that capture the radiant energy of the Sun. Some of these reactions lead to the heating of the atmosphere during the daytime by absorption of ultraviolet radiation. Others include the absorption of visible radiation during photosynthesis. Without photochemical processes, the Earth would probably be simply a warm, sterile, rock.

\subsection*{17G. 1 Photochemical processes}
Table 17G. 1 summarizes common photochemical reactions. Photochemical processes are initiated by the absorption of radiation by at least one component of a reaction mixture. In a primary process, products are formed directly from the excited state of a reactant. Examples include fluorescence (Topic 11G) and cis-trans photoisomerizations. Products of a secondary process originate from intermediates that are formed directly from the excited state of a reactant, such as oxidative processes initiated by the oxygen atom formed by the photodissociation of ozone.

Table 17G. 1 Examples of photochemical processes

\begin{center}
\begin{tabular}{|c|c|c|}
\hline
Process & General form & Example \\
\hline
Ionization & $\mathrm{A}^{*} \rightarrow \mathrm{~A}^{+}+\mathrm{e}^{-}$ & $\mathrm{NO}^{*} \rightarrow \mathrm{NO}^{+}+\mathrm{e}^{-}$ \\
\hline
Electron transfer & \( \begin{aligned} & \mathrm{A}^{*}+\mathrm{B} \rightarrow \mathrm{~A}^{+}+\mathrm{B}^{-} \\ & \quad \text { or } \mathrm{A}^{-}+\mathrm{B}^{+} \end{aligned} \) & $\mathrm{Ru}(\mathrm{bpy})_{3}^{2+\star}+\mathrm{Fe}^{3+} \rightarrow \mathrm{Ru}(\mathrm{bpy})_{3}^{3+}+\mathrm{Fe}^{2+}$ \\
\hline
\multirow[t]{2}{*}{Dissociation} & $\mathrm{A}^{*} \rightarrow \mathrm{~B}+\mathrm{C}$ & $\mathrm{O}_{3}^{*} \rightarrow \mathrm{O}_{2}+\mathrm{O}$ \\
\hline
 & \( \begin{array}{r} \mathrm{A}^{*}+\mathrm{B}-\mathrm{C} \rightarrow \\ \mathrm{~A}+\mathrm{B}+\mathrm{C} \end{array} \) & $\mathrm{Hg}^{*}+\mathrm{CH}_{4} \rightarrow \mathrm{Hg}+\mathrm{CH}_{3}+\mathrm{H}$ \\
\hline
\multirow[t]{2}{*}{Addition} & $\mathrm{A}^{*}+\mathrm{A}^{*} \rightarrow \mathrm{~B}$ & \includegraphics[max width=\textwidth]{2025_02_27_d4b95d9718f0b9e2e6b9g-794}
 \\
\hline
 & $\mathrm{A}^{*}+\mathrm{B} \rightarrow \mathrm{AB}$ & $\mathrm{Hg}^{*}+\mathrm{H}_{2} \rightarrow \mathrm{HgH}+\mathrm{H}$ \\
\hline
Abstraction & \( \begin{array}{r} \mathrm{A}^{*}+\mathrm{B}-\mathrm{C} \rightarrow \\ \mathrm{~A}-\mathrm{B}+\mathrm{C} \end{array} \) & $\mathrm{Hg}^{*}+\mathrm{CH}_{3}-\mathrm{H} \rightarrow \mathrm{Hg}-\mathrm{CH}_{3}+\mathrm{H}$ \\
\hline
\begin{tabular}{l}
Isomerization \\
or \\
rearrangement \\
\end{tabular} & $\mathrm{A}^{*} \rightarrow \mathrm{~A}^{\prime}$ & \includegraphics{smile-3884672e9ccbbaaa07ffe744f8c72c371a6074b6} \\
\hline
\end{tabular}
\end{center}

\begin{itemize}
  \item Excited state.
\end{itemize}

Competing with the formation of photochemical products are numerous primary photophysical processes that can deactivate the excited state (Table 17G.2). Therefore, it is important to consider the timescales of the formation and decay of excited states before describing the mechanisms of photochemical reactions.

Electronic transitions caused by absorption of ultraviolet and visible radiation occur within $10^{-16}-10^{-15} \mathrm{~s}$. The upper limit for the rate constant of a first-order photochemical reaction is then expected to be about $10^{16} \mathrm{~s}^{-1}$. Fluorescence is slower than absorption, with typical lifetimes of $10^{-12}-10^{-6} \mathrm{~s}$. Therefore, the excited singlet state can initiate very fast photochemical reactions in the femtosecond $\left(10^{-15} \mathrm{~s}\right)$ to picosecond $\left(10^{-12} \mathrm{~s}\right)$ range. Examples of such ultrafast reactions are the initial events of vision and of photosynthesis. Internal conversion (IC) occurs on a timescale similar to that for the release of vibrational energy in molecules, so can occur in less than $10^{-12}$ s. Typical intersystem crossing (ISC, Topic 11G) and phosphorescence lifetimes for large organic molecules are $10^{-12}-10^{-4} \mathrm{~s}$ and $10^{-6}-10^{-1} \mathrm{~s}$, respectively. As a consequence of their long lifetimes, excited triplet states are photochemically important. Indeed, because phosphorescence decay is several orders of magnitude slower than most typical reactions,

Table 17G. 2 Common photophysical processes

\begin{center}
\begin{tabular}{ll}
\hline
Primary absorption & $\mathrm{S}+h v \rightarrow \mathrm{~S}^{*}$ \\
Excited-state absorption & $\mathrm{S}^{*}+h v \rightarrow \mathrm{~S}^{* *}$ \\
 & $\mathrm{~T}^{*}+h v \rightarrow \mathrm{~T}^{* *}$ \\
Fluorescence & $\mathrm{S}^{*} \rightarrow \mathrm{~S}+h v$ \\
Stimulated emission & $\mathrm{S}^{*}+h v \rightarrow \mathrm{~S}+2 h v$ \\
Intersystem crossing (ISC) & $\mathrm{S}^{*} \rightarrow \mathrm{~T}^{*}$ \\
Phosphorescence & $\mathrm{T}^{*} \rightarrow \mathrm{~S}+h v$ \\
Internal conversion (IC) & $\mathrm{S}^{*} \rightarrow \mathrm{~S}$ \\
Collision-induced emission & $\mathrm{S}^{*}+\mathrm{M} \rightarrow \mathrm{S}+\mathrm{M}+h v$ \\
Collisional deactivation & $\mathrm{S}^{*}+\mathrm{M} \rightarrow \mathrm{S}+\mathrm{M}$ \\
 & $\mathrm{T}^{*}+\mathrm{M} \rightarrow \mathrm{S}+\mathrm{M}$ \\
Electronic energy transfer: &  \\
$\quad$ Singlet-singlet & $\mathrm{S}^{*}+\mathrm{S} \rightarrow \mathrm{S}+\mathrm{S}^{*}$ \\
$\quad$ Triplet-triplet & $\mathrm{T}^{*}+\mathrm{T} \rightarrow \mathrm{T}+\mathrm{T}^{*}$ \\
Excimer formation & $\mathrm{S}^{*}+\mathrm{S} \rightarrow(\mathrm{SS})^{*}$ \\
Energy pooling &  \\
$\quad$ Singlet-singlet & $\mathrm{S}^{*}+\mathrm{S}^{*} \rightarrow \mathrm{~S}^{* *}+\mathrm{S}$ \\
Triplet-triplet & $\mathrm{T}^{*}+\mathrm{T}^{*} \rightarrow \mathrm{~S}^{* *}+\mathrm{S}$ \\
\hline
\end{tabular}
\end{center}

\begin{itemize}
  \item Denotes an excited state; S is a singlet state, T a triplet state, and M a 'third body'.\\
species in excited triplet states can undergo a very large number of collisions with other reactants before they lose their energy radiatively.
\end{itemize}

\subsection*{Brief illustration 17G. 1}
To judge whether the excited singlet or triplet state of the reactant is a suitable product precursor, the emission lifetimes are compared with the half-life of the relevant chemical reaction (Topic 17B). Consider a unimolecular photochemical reaction with rate constant $k_{\mathrm{r}}=1.7 \times 10^{4} \mathrm{~s}^{-1}$, and therefore a half-life of $41 \mu \mathrm{~s}$. The observed fluorescence lifetime of the reactant is 1.0 ns and the observed phosphorescence lifetime is 1.0 ms . The excited singlet state is therefore too short-lived to be a major source of product in this reaction. On the other hand, the relatively long-lived excited triplet state is a good candidate for an intermediate.

\subsection*{17G. 2 The primary quantum yield}
The rates of deactivation of the excited state by radiative, non-radiative, and chemical processes determine the yield of product in a photochemical reaction. The primary quantum yield, $\phi$, is defined as the number of photophysical or photochemical events that lead to primary products divided\\
by the number of photons absorbed by the molecule in the same interval:

$$
\phi=\frac{\text { number of events }}{\text { number of photons absorbed }}=\frac{N_{\text {events }}}{N_{\mathrm{abs}}}
$$


\begin{align*}
& \text { Primary quantum yield }  \tag{17G.1a}\\
& \text { [definition] }
\end{align*}


When both the numerator and denominator of this expression are divided by the time interval over which the events occur, the primary quantum yield is also seen to be the rate of radiation-induced primary events divided by the rate of photon absorption, $I_{\text {abs }}$ :

$$
\phi=\frac{\text { rate of process }}{\text { rate of photon absorption }}=\frac{v}{I_{\mathrm{abs}}}
$$

Primary quantum yield in terms of rates of processes\\
(17G.1b)

\subsection*{Example 17G. 1 Calculating a primary quantum yield}
In an experiment to determine the quantum yield of a photochemical reaction, the absorbing substance was exposed to light of wavelength 490 nm from a 1.00 W laser source for 2700 s , with 60 per cent of the incident light being absorbed. As a result of irradiation, 3.44 mmol of the absorbing substance decomposed. What is the primary quantum yield?

Collect your thoughts You need to calculate the quantities in eqn \href{http://17G.la}{17G.la}. The number of photochemical events is simply the number of decomposed molecules, $N_{\text {events }}=$ $N_{\text {decomposed }}$. To calculate the number of absorbed photons $N_{\text {abs }}$, note that:

\begin{itemize}
  \item The energy absorbed by the substance is $E_{\text {abs }}=f P t$, where $P$ is the incident power, $t$ is the time of exposure, and the factor $f$ (in this case $f=0.60$ ) is the fraction of incident light that is absorbed.
  \item $E_{\text {abs }}$ is also related to the number $N_{\text {abs }}$ of absorbed photons through $E_{\text {abs }}=N_{\text {abs }} h \nu=N_{\text {abs }} h c / \lambda$, where $h c / \lambda$ is the energy of a single photon of wavelength $\lambda$.
\end{itemize}

You can combine these two expressions for the absorbed energy to obtain $N_{\text {abs }}$. The primary quantum yield follows from $\phi=N_{\text {decomposed }} / N_{\text {abs }}$.\\
The solution From the two expressions for the absorbed energy, it follows that

$$
f P t=N_{\mathrm{abs}}\left(\frac{h c}{\lambda}\right)
$$

and therefore that $N_{\text {abs }}=f P t \lambda / h c$. Now use eqn 17G.1a to write

$$
\phi=\frac{N_{\text {decomposed }}}{N_{\text {abs }}}=\frac{N_{\text {decomposes }} h c}{f P t \lambda}
$$

With $N_{\text {decomposed }}=\left(3.44 \times 10^{-3} \mathrm{~mol}\right) \times\left(6.022 \times 10^{23} \mathrm{~mol}^{-1}\right)=$ $2.07 \ldots \times 10^{21}, P=1.00 \mathrm{~W}=1.00 \mathrm{~J} \mathrm{~s}^{-1}, t=2700 \mathrm{~s}, \lambda=490 \mathrm{~nm}=$ $4.90 \times 10^{-7} \mathrm{~m}$, and $f=0.60$ it follows that

$$
\begin{aligned}
\phi & =\frac{\left(2.07 \ldots \times 10^{21}\right) \times\left(6.626 \times 10^{-34} \mathrm{Js}\right) \times\left(2.998 \times 10^{8} \mathrm{~m} \mathrm{~s}^{-1}\right)}{0.60 \times\left(1.00 \mathrm{Js} \mathrm{~s}^{-1}\right) \times(2700 \mathrm{~s}) \times\left(4.90 \times 10^{-7} \mathrm{~m}\right)} \\
& =0.52
\end{aligned}
$$

That is, about half the photons that are absorbed bring about photodissociation.

Self-test 17G. 1 In an experiment to measure the quantum yield of a photochemical reaction, the absorbing substance was exposed to 320 nm radiation from an 87.5 mW laser source for 38 min . The intensity of the transmitted light was 0.35 that of the incident light. As a result of irradiation, 0.324 mmol of the absorbing substance decomposed. Determine the primary quantum yield.\\
$\mathcal{E} 6^{\circ} 0=\phi: \iota \partial s u_{V}$

A molecule in an excited state must either decay to the ground state or form a photochemical product. Therefore, the total number of molecules deactivated by radiative processes, non-radiative processes, and photochemical reactions must be equal to the number of excited species produced by absorption of the incident radiation. It follows that the sum of primary quantum yields $\phi_{i}$ for all photophysical and photochemical events $i$ must be equal to 1 , regardless of the number of reactions involving the excited state:


\begin{equation*}
\sum_{i} \phi_{i}=\sum_{i} \frac{v_{i}}{I_{\mathrm{abs}}}=\frac{1}{I_{\mathrm{abs}}} \sum_{i} v_{i}=1 \tag{17G.2a}
\end{equation*}


Then, from eqn 17G.1b in the form $\phi_{i}=v_{i} / I_{\text {abs }}$ it follows that


\begin{equation*}
\phi_{i}=\frac{v_{i}}{\sum_{i} v_{i}} \tag{17G.2b}
\end{equation*}


Therefore, the primary quantum yield of a particular process may be determined directly from the experimental rates of all photophysical and photochemical processes that deactivate the excited state.

If it is assumed that the only photophysical processes for the excited singlet state are fluorescence, internal conversion, and phosphorescence, then it follows that

$$
\phi_{\mathrm{F}}+\phi_{\mathrm{IC}}+\phi_{\mathrm{P}}=1
$$

where $\phi_{\mathrm{F}}, \phi_{\mathrm{IC}}$, and $\phi_{\mathrm{P}}$ are the quantum yields of fluorescence, internal conversion, and phosphorescence, respectively (intersystem crossing from the singlet to the triplet state is taken into account by the presence of $\phi_{\mathrm{p}}$ ). The quantum yield of photon emission by fluorescence and phosphorescence is $\phi_{\text {emission }}=\phi_{\mathrm{F}}+$ $\phi_{\mathrm{p}}$, which is less than 1. If the excited singlet state also participates in a primary photochemical reaction with quantum yield $\phi_{r}$, then

$$
\phi_{\mathrm{F}}+\phi_{\mathrm{IC}}+\phi_{\mathrm{P}}+\phi_{\mathrm{r}}=1
$$

\subsection*{17G. 3 Mechanism of decay of excited singlet states}
Consider the formation and decay of an excited singlet state in the absence of a chemical reaction:\\
Absorption:\\
$\mathrm{S}+h v_{\mathrm{i}} \rightarrow \mathrm{S}^{*}$\\
$v_{\mathrm{abs}}=I_{\mathrm{abs}}$\\
Fluorescence:\\
Internal conversion:\\
Intersystem crossing:

$$
\begin{array}{ll}
\mathrm{S}^{*} \rightarrow \mathrm{~S}+h v_{\mathrm{f}} & v_{\mathrm{F}}=k_{\mathrm{F}}\left[\mathrm{~S}^{*}\right] \\
\mathrm{S}^{*} \rightarrow \mathrm{~S} & v_{\mathrm{IC}}=k_{\mathrm{IC}}\left[\mathrm{~S}^{*}\right]
\end{array}
$$

in which S is an absorbing singlet-state species, $\mathrm{S}^{*}$ an excited singlet state, $\mathrm{T}^{*}$ an excited triplet state, and $h v_{\mathrm{i}}$ and $h v_{\mathrm{f}}$ are the energies of the incident and fluorescent photons, respectively. From the methods presented in Topic 17E, the rate of formation of $\mathrm{S}^{*}$ and its net rate of disappearance may be written as:

$$
\begin{aligned}
\text { Rate of formation of } \mathrm{S}^{*} & =I_{\mathrm{abs}} \\
\text { Rate of disappearance of } \mathrm{S}^{*} & =k_{\mathrm{F}}\left[\mathrm{~S}^{*}\right]+k_{\mathrm{ISC}}\left[\mathrm{~S}^{*}\right]+k_{\mathrm{IC}}\left[\mathrm{~S}^{*}\right] \\
& =\left(k_{\mathrm{F}}+k_{\mathrm{ISC}}+k_{\mathrm{IC}}\right)\left[\mathrm{S}^{*}\right]
\end{aligned}
$$

It follows that the excited state decays by a first-order process so, when the light is turned off, the concentration of $\mathrm{S}^{\star}$ varies with time $t$ as


\begin{equation*}
\left[\mathrm{S}^{*}\right](t)=\left[\mathrm{S}^{*}\right]_{0} \mathrm{e}^{-t / \tau_{0}} \tag{17G.3a}
\end{equation*}


where the observed lifetime, $\tau_{0}$, of the excited singlet state is defined as


\begin{equation*}
\tau_{0}=\frac{1}{k_{\mathrm{F}}+k_{\mathrm{ISC}}+k_{\mathrm{IC}}} \tag{17G.3b}
\end{equation*}


Observed lifetime of the excited singlet state [definition]

This expression can be used in a kinetic analysis of the decay of $S^{*}$ to find an expression for the quantum yield of fluorescence.

\subsection*{How is that done? 17G. 1 Deriving an expression for the quantum yield of fluorescence}
Most fluorescence measurements are conducted by illuminating a dilute sample with a continuous and intense beam of visible or ultraviolet radiation. It follows that $\left[\mathrm{S}^{*}\right]$ is small and constant, so the steady-state approximation (Topic 17E) may be used for $\left[\mathrm{S}^{*}\right]$ :

$$
\begin{aligned}
\frac{\mathrm{d}\left[\mathrm{~S}^{*}\right]}{\mathrm{d} t} & =I_{\mathrm{abs}}-k_{\mathrm{F}}\left[\mathrm{~S}^{*}\right]-k_{\mathrm{ISC}}\left[\mathrm{~S}^{*}\right]-k_{\mathrm{IC}}\left[\mathrm{~S}^{*}\right] \\
& =I_{\mathrm{abs}}-\left(k_{\mathrm{F}}+k_{\mathrm{ISC}}+k_{\mathrm{lC}}\right)\left[\mathrm{S}^{*}\right] \approx 0
\end{aligned}
$$

Consequently,

$$
I_{\mathrm{abs}}=\left(k_{\mathrm{F}}+k_{\mathrm{ISC}}+k_{\mathrm{IC}}\right)\left[\mathrm{S}^{\star}\right]
$$

The rate of fluorescence, $v_{\mathrm{F}}$, is $k_{\mathrm{F}}\left[\mathrm{S}^{*}\right]$, so it follows from eqn 17 G .1 b that the quantum yield of fluorescence is

$$
\phi_{\mathrm{F}, 0}=\frac{v_{\mathrm{F}}}{I_{\mathrm{abs}}}=\frac{k_{\mathrm{F}}\left[\mathrm{~S}^{*}\right]}{\left(k_{\mathrm{F}}+k_{\mathrm{ISC}}+k_{\mathrm{IC}}\right)\left[\mathrm{S}^{*}\right]}
$$

which, by cancelling the [ $\mathrm{S}^{*}$ ], simplifies to

$$
\phi_{\mathrm{F}, 0}=\frac{k_{\mathrm{F}}}{k_{\mathrm{F}}+k_{\mathrm{ISC}}+k_{\mathrm{IC}}}
$$

Then, by using the result for the lifetime in eqn 17G.3b,

$$
-\left|\phi_{\mathrm{F}, 0}=k_{\mathrm{F}} \tau_{0}\right| \quad \text { Quantum yield of fluorescence }
$$

The observed fluorescence lifetime can be measured by using a pulsed laser technique. First, the sample is excited with a short light pulse from a laser using a wavelength at which S absorbs strongly. Then, the exponential decay of the fluorescence intensity after the pulse is monitored.

\subsection*{Brief illustration 17G. 2}
At a certain wavelength, the fluorescence quantum yield and observed fluorescence lifetime of tryptophan in water are $\phi_{\mathrm{F}, 0}=0.20$ and $\tau_{0}=2.6 \mathrm{~ns}$, respectively. It follows from eqn 17G. 4 that the fluorescence rate constant $k_{\mathrm{F}}$ is

$$
k_{\mathrm{F}}=\frac{\phi_{\mathrm{F}, 0}}{\tau_{0}}=\frac{0.20}{2.6 \times 10^{-9} \mathrm{~s}}=7.7 \times 10^{7} \mathrm{~s}^{-1}
$$

\subsection*{17G.4 Quenching}
The shortening of the lifetime of the excited state by the presence of another species is called quenching. Quenching may be either a desired process, such as in energy or electron transfer, or an undesired side reaction that can decrease the quantum yield of a desired photochemical process. Quenching effects are studied by monitoring the emission from the excited state that is involved in the photochemical reaction.

The addition of a quencher, Q , opens an additional channel for deactivation of $S^{*}$ :

$$
\text { Quenching: } \mathrm{S}^{*}+\mathrm{Q} \rightarrow \mathrm{~S}+\mathrm{Q} \quad v_{\mathrm{Q}}=k_{\mathrm{Q}}[\mathrm{Q}]\left[\mathrm{S}^{*}\right]
$$

The fluorescence quantum yields $\phi_{\mathrm{F}, 0}$ and $\phi_{\mathrm{F}}$ measured in the absence and presence of Q , respectively, can be expressed in terms of the molar concentration of the quencher, $[\mathrm{Q}]$.

How is that done? 17G. 2 Assessing the effect of a quencher on the fluorescence quantum yield

In the presence of quenching, the steady-state approximation for [ $\mathrm{S}^{\star}$ ] becomes

$$
\frac{\mathrm{d}\left[\mathrm{~S}^{\star}\right]}{\mathrm{d} t}=I_{\mathrm{abs}}-\left(k_{\mathrm{F}}+k_{\mathrm{ISC}}+k_{\mathrm{IC}}+k_{\mathrm{Q}}[\mathrm{Q}]\right)\left[\mathrm{S}^{\star}\right] \approx 0
$$

and the fluorescence quantum yield is

$$
\phi_{\mathrm{F}}=\frac{k_{\mathrm{F}}}{k_{\mathrm{F}}+k_{\mathrm{ISC}}+k_{\mathrm{IC}}+k_{\mathrm{Q}}[\mathrm{Q}]}
$$

The ratio of the quantum yields without and with a quencher present is

$$
\begin{aligned}
\frac{\phi_{\mathrm{F}, 0}}{\phi_{\mathrm{F}}} & =\frac{k_{\mathrm{F}}}{k_{\mathrm{F}}+k_{\mathrm{ISC}}+k_{\mathrm{IC}}} \times \frac{k_{\mathrm{F}}+k_{\mathrm{ISC}}+k_{\mathrm{IC}}+k_{\mathrm{Q}}[\mathrm{Q}]}{k_{\mathrm{F}}} \\
& =\frac{k_{\mathrm{F}}+k_{\mathrm{ISC}}+k_{\mathrm{IC}}+k_{\mathrm{Q}}[\mathrm{Q}]}{k_{\mathrm{F}}+k_{\mathrm{ISC}}+k_{\mathrm{IC}}} \\
& =1+\frac{k_{\mathrm{Q}}}{k_{\mathrm{F}}+k_{\mathrm{ISC}}+k_{\mathrm{IC}}}[\mathrm{Q}]
\end{aligned}
$$

By recognizing from eqn 17G.3b that $1 /\left(k_{\mathrm{F}}+k_{\mathrm{ISC}}+k_{\mathrm{IC}}\right)=\tau_{0}$, this expression becomes the Stern-Volmer equation:\\
$-\left|\frac{\phi_{\mathrm{F}, 0}}{\phi_{\mathrm{F}}}=1+\tau_{0} k_{\mathrm{Q}}[\mathrm{Q}]\right|$ Stern-Volmer equation

The Stern-Volmer equation implies that a plot of $\phi_{\mathrm{F}, 0} / \phi_{\mathrm{F}}$ against [Q] should be a straight line with slope $\tau_{0} k_{\mathrm{Q}}$. Such a plot is called a Stern-Volmer plot (Fig. 17G.1). The method may also be applied to the quenching of phosphorescence.

Equation 17G. 4 in the form $k_{\mathrm{F}}=\phi_{\mathrm{F}, 0} / \tau_{0}$ shows that the rate constant for fluorescence, and hence the rate of fluorescence (which determines the intensity of fluorescence), is proportional to the quantum yield. The ratio $\phi_{\mathrm{F}, 0} / \phi_{\mathrm{F}}$ is therefore equal to the ratio $I_{\mathrm{F}, 0} / I_{\mathrm{F}}$, where $I_{\mathrm{F}, 0}$ is the intensity of fluorescence in the absence of quencher and $I_{\mathrm{F}}$ the intensity when quencher is present. Similarly, from the same equation in the form $\tau_{0}=\phi_{\mathrm{F}, 0} / k_{\mathrm{F}}$, the fluorescence lifetime is also proportional to the quantum yield, so the ratio $\tau_{0} / \tau$ (where $\tau$ is the lifetime in the presence of the quencher) is also equal to $\phi_{\mathrm{F}, 0} / \phi_{\mathrm{F}}$.\\
\includegraphics[max width=\textwidth, center]{2025_02_27_d4b95d9718f0b9e2e6b9g-797}

Figure 17G. 1 The form of a Stern-Volmer plot and the interpretation of the slope in terms of the rate constant for quenching and the observed fluorescence lifetime in the absence of quenching.

Stern-Volmer plots can therefore be made by plotting either $I_{\mathrm{F}, 0} / I_{\mathrm{F}}$ or $\tau_{0} / \tau$ against the quencher concentration. The slope and intercept are the same as those shown for eqn 17G.5.

\subsection*{Example 17G. 2}
\subsection*{Determining the quenching rate constant}
The molecule $2,2^{\prime}$-bipyridine ( $1, \mathrm{bpy}$ ) forms a complex with the $\mathrm{Ru}^{2+}$ ion. Tris-(2, $2^{\prime}$-bipyridyl)ruthenium(II), $\left[\mathrm{Ru}(\mathrm{bpy})_{3}\right]^{2+}$ (2), has a strong metal-to-ligand charge transfer (MLCT) transition (Topic 11F) at 450 nm .\\


1 2,2'-Bipyridine (bpy)\\
\includegraphics[max width=\textwidth, center]{2025_02_27_d4b95d9718f0b9e2e6b9g-798}\\
$2\left[\mathrm{Ru}(\mathrm{bpy})_{3}\right]^{2+}$

The quenching of the ${ }^{*}\left[\mathrm{Ru}(\mathrm{bpy})_{3}\right]^{2+}$ excited state by $\mathrm{Fe}^{3+}$ (present as the complex ion $\left[\mathrm{Fe}\left(\mathrm{OH}_{2}\right)_{6}\right]^{3+}$ ) in acidic solution was monitored by measuring emission lifetimes at 600 nm . Determine the quenching rate constant for this reaction from the following data:

\begin{center}
\begin{tabular}{llllll}
$\left[\left[\mathrm{Fe}\left(\mathrm{OH}_{2}\right)_{6}\right]^{3+}\right] /\left(10^{-2} \mathrm{moldm}^{-3}\right)$ & 0 & 1.6 & 4.7 & 7.0 & 9.4 \\
$\tau /\left(10^{-7} \mathrm{~s}\right)$ & 6.00 & 4.05 & 3.37 & 2.96 & 2.17 \\
\end{tabular}
\end{center}

Collect your thoughts Rewrite the Stern-Volmer equation (eqn 17G.5) for use with lifetime data; then fit the data to a straight line.

The solution Substitute $\tau_{0} / \tau$ for $\phi_{\mathrm{F}, 0} / \phi_{\mathrm{F}}$ in eqn 17G. 5 and, after rearrangement, obtain

$$
\frac{1}{\tau}=\frac{1}{\tau_{0}}+k_{\mathrm{Q}}[\mathrm{Q}]
$$

Because the axes of plots should be labelled with pure numbers, it is necessary to introduce and handle units before using this equation for the analysis of the data. To bring the expression into a form suitable for plotting, it needs to be expressed in terms of $\tau /\left(10^{-7} \mathrm{~s}\right)$ and $[\mathrm{Q}] /\left(10^{-2} \mathrm{moldm}^{-3}\right)$ to match the data, and therefore (with these dimensionless terms in blue) to write it as

$$
\begin{aligned}
\frac{1}{\left(10^{-7} \mathrm{~s}\right) \tau /\left(10^{-7} \mathrm{~s}\right)}= & \frac{1}{\tau_{0}}+\left\{k_{\mathrm{Q}}[\mathrm{Q}] /\left(10^{-2} \mathrm{~mol} \mathrm{dm}^{-3}\right)\right\} \\
& \times\left(10^{-2} \mathrm{~mol} \mathrm{dm}^{-3}\right)
\end{aligned}
$$

Now multiply through by $10^{-7}$ s to obtain

$$
\begin{aligned}
\frac{1}{\tau /\left(10^{-7} \mathrm{~s}\right)}= & \frac{10^{-7} \mathrm{~s}}{\tau_{0}}+k_{\mathrm{Q}} \times\left(10^{-2} \mathrm{~mol} \mathrm{dm}^{-3}\right) \times\left(10^{-7} \mathrm{~s}\right) \\
& \times[\mathrm{Q}] /\left(10^{-2} \mathrm{~mol} \mathrm{dm}^{-3}\right)
\end{aligned}
$$

and collect terms:

$$
\begin{aligned}
\overbrace{\frac{1}{\tau /\left(10^{-7} \mathrm{~s}\right)}}^{y}= & \overbrace{\frac{1}{\tau_{0} /\left(10^{-7} \mathrm{~s}\right)}}^{y \text {-intercept }}+\overbrace{\left(k_{\mathrm{Q}} \times 10^{-9} \mathrm{~mol} \mathrm{dm}^{-3} \mathrm{~s}\right)}^{\text {slope }} \\
& \times \overbrace{[\mathrm{Q}] /\left(10^{-2} \mathrm{~mol} \mathrm{dm}\right.}{ }^{-3})
\end{aligned}
$$

Note that because slope $=k_{\mathrm{Q}} \times 10^{-9} \mathrm{~mol} \mathrm{dm}^{-3} \mathrm{~s}$, then $k_{\mathrm{Q}}=$ slope $\times$ $10^{9} \mathrm{dm}^{3} \mathrm{~mol}^{-1} \mathrm{~s}^{-1}$. Draw up the following table with $\mathrm{Q}=$ $\left[\mathrm{Fe}\left(\mathrm{OH}_{2}\right)_{6}\right]^{3+}$ :

$$
\begin{array}{llllll}
{\left[\left[\mathrm{Fe}\left(\mathrm{OH}_{2}\right)_{6}\right]^{3+}\right] /\left(10^{-2} \mathrm{~mol} \mathrm{dm}^{-3}\right)} & 0 & 1.6 & 4.7 & 7.0 & 9.4 \\
1 /\left(\tau / 10^{-7} \mathrm{~s}\right) & 0.167 & 0.247 & 0.297 & 0.338 & 0.461
\end{array}
$$

Figure 17G. 2 shows a plot of $1 /\left(\tau / 10^{-7} \mathrm{~s}\right)$ against $\left[\left[\mathrm{Fe}\left(\mathrm{OH}_{2}\right)_{6}\right]^{3+}\right] /$ $\left(10^{-2} \mathrm{moldm}^{-3}\right)$ and the results of a fit to this expression. The slope of the line is 0.029 , so $k_{\mathrm{Q}}=2.9 \times 10^{7} \mathrm{dm}^{3} \mathrm{~mol}^{-1} \mathrm{~s}^{-1}$.\\
\includegraphics[max width=\textwidth, center]{2025_02_27_d4b95d9718f0b9e2e6b9g-798(1)}

Figure 17G. 2 The Stern-Volmer plot of the data for Example 17G.2.

Comment. Measurements of emission lifetimes are preferred because they yield the value of $k_{\mathrm{Q}}$ directly. To determine the value of $k_{\mathrm{Q}}$ from intensity or quantum yield measurements, it is necessary to make an independent measurement of $\tau_{0}$.

Self-test 17G. 2 The quenching of tryptophan fluorescence by dissolved $\mathrm{O}_{2}$ gas was monitored by measuring emission lifetimes at 348 nm in aqueous solutions. Determine the quenching rate constant for this process from the following data:

\begin{center}
\begin{tabular}{|c|c|c|c|c|c|}
\hline
$\left[\mathrm{O}_{2}\right] /\left(10^{-2} \mathrm{moldm}^{-3}\right)$ & 0 & 2.3 & 5.5 & 8 & 10.8 \\
\hline
$\tau / \mathrm{ns}$ & 2.6 & 1.5 & 0.92 & 0.71 & 0.57 \\
\hline
\end{tabular}
\end{center}

Three common mechanisms for bimolecular quenching of an excited singlet (or triplet) state are:

$$
\begin{array}{ll}
\text { Collisional deactivation: } & \mathrm{S}^{*}+\mathrm{Q} \rightarrow \mathrm{~S}+\mathrm{Q} \\
\text { Resonance energy transfer: } & \mathrm{S}^{*}+\mathrm{Q} \rightarrow \mathrm{~S}+\mathrm{Q}^{*} \\
\text { Electron transfer: } & \mathrm{S}^{*}+\mathrm{Q} \rightarrow \mathrm{~S}^{+/-}+\mathrm{Q}^{-/+}
\end{array}
$$

The quenching rate constant itself does not give much insight into the mechanism of quenching. Collisional quenching is particularly efficient when Q is a species, such as iodide ion, which receives energy from $\mathrm{S}^{*}$ and then decays to the ground state primarily by releasing energy as heat. For the system of Example 17G.2, it is known that the quenching of the excited state of $\left[\mathrm{Ru}(\mathrm{bpy})_{3}\right]^{2+}$ is a result of electron transfer to $\mathrm{Fe}^{3+}$, but the quenching data do not prove the mechanism.

\subsection*{17G. 5 Resonance energy transfer}
The energy transfer process $\mathrm{S}^{*}+\mathrm{Q} \rightarrow \mathrm{S}+\mathrm{Q}^{*}$ can be regarded as taking place as follows. The oscillating electric field of the incoming electromagnetic radiation induces an oscillating electric dipole moment (a transition dipole moment) in S. Energy is absorbed by S if the frequency of the incident radiation, $v$, is such that $v=\Delta E_{\mathrm{S}} / h$, where $\Delta E_{\mathrm{S}}$ is the energy separation of the ground and excited electronic states of S and $h$ is Planck's constant. This is the 'resonance condition' for absorption of radiation (essentially the Bohr frequency condition, eqn 7A.9). The oscillating dipole on $S$ can now affect electrons bound to a nearby Q molecule by inducing an oscillating dipole moment (another transition dipole moment) in them. If the frequency of oscillation of the electric dipole moment in $S$ is such that $v=\Delta E_{\mathrm{Q}} / h$, where $\Delta E_{\mathrm{Q}}$ is the energy separation of the ground and excited electronic states of Q , then Q will absorb energy from $S$. The coupling of the two transition moments can be regarded as an exchange of a photon, in which a photon generated by $S$ is absorbed by Q .

The efficiency, $\eta_{\mathrm{T}}$, of resonance energy transfer is defined as

\[
\eta_{\mathrm{T}}=1-\frac{\phi_{\mathrm{F}}}{\phi_{\mathrm{F}, 0}} \quad \begin{align*}
& \text { Efficiency of resonance energy transfer }  \tag{17G.6}\\
& \text { [definition] }
\end{align*}
\]

According to the Förster theory of resonance energy transfer, energy transfer is efficient when:

\begin{itemize}
  \item The energy donor and acceptor are separated by a short distance (of the order of nanometres).
  \item The photon is regarded as emitted by the excited state of the donor and then absorbed directly by the acceptor.
\end{itemize}

For donor-acceptor systems held rigidly either by covalent bonds or by a protein 'scaffold', $\eta_{\mathrm{T}}$ increases with decreasing distance, $R$, according to

Table 17G. 3 Values of $R_{0}$ for some donor-acceptor pairs*

\begin{center}
\begin{tabular}{lll}
\hline
Donor $^{\ddagger}$ & Acceptor & $R_{0} / \mathrm{nm}$ \\
\hline
Naphthalene & Dansyl & 2.2 \\
Dansyl & ODR & 4.3 \\
Pyrene & Coumarin & 3.9 \\
1.5-I AEDANS & FITC & 4.9 \\
Tryptophan & 1.5-I AEDANS & 2.2 \\
Tryptophan & Haem (heme) & 2.9 \\
\hline
\end{tabular}
\end{center}

*Additional values may be found in J.R. Lacowicz, Principles of fluorescence spectroscopy, Kluwer Academic/Plenum, New York (1999).\\
${ }^{\ddagger}$ Abbreviations:\\
Dansyl: 5-dimethylamino-1-naphthalenesulfonic acid\\
FITC: fluorescein 5-isothiocyanate\\
1.5-I AEDANS: 5-(((2-iodoacetyl)amino)ethyl)amino)naphthalene-1-sulfonic acid (3) ODR: octadecyl-rhodamine

\[
\eta_{\mathrm{T}}=\frac{R_{0}^{6}}{R_{0}^{6}+R^{6}} \quad \begin{align*}
& \text { Efficiency of energy transfer in terms }  \tag{17G.7}\\
& \text { of the donor-acceptor distance }
\end{align*}
\]

where $R_{0}$ is a parameter (with dimensions of distance) that is characteristic of each donor-acceptor pair. It can be regarded as the distance at which energy transfer is 50 per cent efficient for a given donor-acceptor pair. (This assertion can be confirmed by using $R=R_{0}$ in eqn 17G.7.) Equation 17G. 7 has been verified experimentally and values of $R_{0}$ are available for a number of donor-acceptor pairs (Table 17G.3).

The emission and absorption spectra of molecules span a range of wavelengths, so the second requirement of the Förster theory is met when the emission spectrum of the donor molecule overlaps significantly with the absorption spectrum of the acceptor. In the overlap region, a photon emitted by the donor has the appropriate energy to be absorbed by the acceptor (Fig. 17G.3).\\
\includegraphics[max width=\textwidth, center]{2025_02_27_d4b95d9718f0b9e2e6b9g-799}

Figure 17G. 3 According to the Förster theory, the rate of energy transfer from a molecule $S^{*}$ in an excited state to a quencher molecule Q is optimized at radiation frequencies for which the emission spectrum of $S^{*}$ overlaps the absorption spectrum of Q , as shown in the (dark green) shaded region.

Equation 17G. 7 forms the basis of fluorescence resonance energy transfer (FRET), in which the dependence of the energy transfer efficiency, $\eta_{\mathrm{T}}$, on the distance, $R$, between energy donor and acceptor is used to measure distances in biological systems. In a typical FRET experiment, a site on a biopolymer or membrane is labelled covalently with an energy donor and another site is labelled covalently with an energy acceptor. In certain cases, the donor or acceptor may be natural constituents of the system, such as amino acid groups, cofactors, or enzyme substrates. The distance between the labels is then calculated from the known value of $R_{0}$ and eqn 17G.7. Several tests have shown that the FRET technique is useful for measuring distances ranging from 1 to 9 nm .

\subsection*{Brief illustration 17G. 3}
As an illustration of the FRET technique, consider a study of the protein rhodopsin. When an amino acid on the surface of rhodopsin was labelled covalently with the energy donor 1.5-I AEDANS (3), the fluorescence quantum yield of the label decreased from 0.75 to 0.68 due to quenching by the visual pigment 11-cis-retinal (4), which is attached elsewhere in the\\
protein. From eqn 17G. 6 it follows that $\eta_{\mathrm{T}}=1-0.68 / 0.75=$ 0.093 and from eqn 17G. 7 and the known value of $R_{0}=5.4 \mathrm{~nm}$ for the 1.5-I AEDANS/11-cis-retinal, $R=7.9 \mathrm{~nm}$. Therefore, take 7.9 nm to be the distance between the surface of the protein and 11-cis-retinal.\\
\includegraphics{smile-51d0e5e2c5d441e1d8c1a115dc88aa9e5b1388d3}

3 1.5-I AEDANS\\
\includegraphics{smile-8a8aff8beb66273dd743ef5635ad3a96d87d90f5}

4 11-cis-Retinal

If donor and acceptor molecules diffuse in solution or in the gas phase, Förster theory predicts that the efficiency of quenching by energy transfer increases as the average distance travelled between collisions of donor and acceptor decreases. That is, the quenching efficiency increases with concentration of quencher, as predicted by the Stern-Volmer equation.

\subsection*{Checklist of concepts}
\begin{enumerate}
  \item The primary quantum yield of a photochemical reaction is the number of reactant molecules producing specified primary products for each photon absorbed.\\
$\square$ 2. The observed lifetime of an excited state is related to the quantum yield and rate constant of emission.
  \item A Stern-Volmer plot is used to analyse the kinetics of fluorescence quenching in solution.
  \item Collisional deactivation, electron transfer, and resonance energy transfer are common fluorescence quenching processes.\\
$\square$ 5. The efficiency of resonance energy transfer decreases with increasing separation between donor and acceptor molecules.
\end{enumerate}

\subsection*{Checklist of equations}
\begin{center}
\begin{tabular}{llll}
\hline
Property & Equation & Comment & Equation number \\
\hline
Primary quantum yield & $\phi=\nu / I_{\mathrm{abs}}$ &  & 17 G .1 b \\
Excited state lifetime & $\tau_{0}=1 /\left(k_{\mathrm{F}}+k_{\mathrm{ISC}}+k_{\mathrm{IC}}\right)$ & No quencher present & 17 G .3 b \\
Quantum yield of fluorescence & $\phi_{\mathrm{F}, 0}=k_{\mathrm{F}} /\left(k_{\mathrm{F}}+k_{\mathrm{ISC}}+k_{\mathrm{IC}}\right)=k_{\mathrm{F}} \tau_{0}$ & Without quencher present & 17 G .4 \\
Stern-Volmer equation & $\phi_{\mathrm{F}, 0} / \phi_{\mathrm{F}}=1+\tau_{0} k_{\mathrm{Q}}[\mathrm{Q}]$ &  & 17 G .5 \\
Efficiency of resonance energy transfer & $\eta_{\mathrm{T}}=1-\phi_{\mathrm{F}} / \phi_{\mathrm{F}, 0}$ & Definition & $17 \mathrm{G.6}$ \\
 & $\eta_{\mathrm{T}}=R_{0}^{6} /\left(R_{0}^{6}+R^{6}\right)$ & Förster theory & 17 G .7 \\
\hline
\end{tabular}
\end{center}

\chapter{FOCUS 18 Reaction dynamics}
This Focus examines the details of what happens to molecules at the climax of reactions. Extensive changes of structure are taking place and energies the size of dissociation energies are being redistributed among bonds: old bonds are being ripped apart and new bonds are being formed. This is the heart of chemistry.

The calculation of the rates of such processes from first principles is very difficult. Nevertheless, like so many intricate problems, the broad features can be established quite simply. Only upon deeper inquiry do the complications emerge. Several approaches to the calculation of a rate constant for elementary bimolecular processes are explored here, ranging from electron transfer to chemical reactions involving bond breakage and formation. Although a great deal of information can be obtained from gas-phase reactions, many reactions of interest take place in condensed phases, and it is useful to attempt to predict their rates.

\subsection*{18A Collision theory}
This Topic explores 'collision theory', the simplest quantitative account of reaction rates. The treatment can be used only for the discussion of reactions between simple species in the gas phase. Basic collision theory considers only the impact of one molecule on another. An elaboration considered in this Topic takes into account how the resulting excitation energy accumulates in the bond where it is needed.\\
18A. 1 Reactive encounters; 18A. 2 The RRK model

\subsection*{18B Diffusion-controlled reactions}
Reactions in solution are classified into two types: 'diffusioncontrolled' where the rate is controlled by the frequency with which reactants meet, and 'activation-controlled', where the accumulation of sufficient energy in a pair that have met is ratedetermining. The rate constants for the former can be expressed quantitatively in terms of the diffusional characteristics of species in liquids. A more detailed account of the space- and time-development of products is obtained by using the diffusion equation.

\subsection*{18C Transition-state theory}
This Topic discusses 'transition-state theory', in which it is assumed that the reactant molecules form a complex that can be discussed in terms of the population of its energy levels. The theory inspires a thermodynamic approach to reaction rates, in which the rate constant is expressed in terms of thermodynamic parameters. This approach is useful for parametrizing the rates of reactions in solution.\\
18C. 1 The Eyring equation; 18C. 2 Thermodynamic aspects;\\
18C. 3 The kinetic isotope effect

\subsection*{18D The dynamics of molecular collisions}
The highest level of sophistication in the theoretical study of chemical reactions is in terms of potential energy surfaces and the motion of molecules on these surfaces. As explained in this Topic, such an approach gives an intimate picture of the events that occur when molecules collide, and provides a basis for studying them by using molecular beams.\\
18D. 1 Molecular beams; 18D. 2 Reactive collisions; 18D. 3 Potential energy surfaces; 18D. 4 Some results from experiments and calculations

\subsection*{18E Electron transfer in homogeneous systems}
In this Topic transition-state theory is used to examine the transfer of electrons in homogeneous systems, which include oxidation-reduction reactions in solution. One widely used theory, Marcus theory, establishes a relation between the activation parameters and the rate constant of electron transfer, and can be expressed in terms of structural parameters of the species involved.\\
18E. 1 The rate law; 18E. 2 The role of electron tunnelling; 18E. 3 The rate constant; 18E. 4 Experimental tests of the theory

\section*{TOPIC 18A Collision theory}
\subsection*{Why do you need to know this material?}
A major component of chemistry is the study of the detailed molecular mechanisms of chemical reactions. One of the earliest approaches, which continues to give insight into the details of mechanisms of gas-phase reactions, is collision theory.

\subsection*{What is the key idea?}
According to collision theory a bimolecular gas-phase reaction takes place when reactants collide, provided their relative kinetic energy exceeds a threshold value and certain steric requirements are fulfilled.

\subsection*{What do you need to know already?}
This Topic draws on the kinetic theory of gases, especially the expression for the mean speed of molecules (Topic 1B), and extends the account of the Lindemann-Hinshelwood mechanism of gas-phase reactions (Topic 17F). One argument draws on the Maxwell-Boltzmann distribution of molecular speeds (Topic 1B).

The rate constant of the bimolecular elementary reaction


\begin{equation*}
\mathrm{A}+\mathrm{B} \rightarrow \mathrm{P} \quad v=k_{\mathrm{r}}[\mathrm{~A}][\mathrm{B}] \tag{18A.1a}
\end{equation*}


depends on the temperature according to the Arrhenius expression (Topic 17D):


\begin{equation*}
k_{\mathrm{r}}=A \mathrm{e}^{-E_{\mathrm{a}} / R T} \tag{18A.1b}
\end{equation*}


Arrhenius expression\\
where $A$ is the 'frequency factor' and $E_{\mathrm{a}}$ is the 'activation energy'. This form of the Arrhenius expression can be explained by a model in which molecules in the gas collide and in the process may acquire sufficient energy to undergo reaction. Like all models, this one can be improved, but it is a good starting point for the discussion of gas-phase reactions.

\subsection*{18A. 1 Reactive encounters}
The general form of the expression for $k_{\mathrm{r}}$ in eqn 18A.1a can be anticipated by considering the physical requirements for reaction. The rate $v$ can be expected to be proportional to the fre-\\
quency of collisions, and therefore to the mean speed of the molecules, $v_{\text {mean }} \propto(T / M)^{1 / 2}$ where $M$ is some combination of the molar masses of A and B. The rate can also be expected to be proportional to the target area the molecules present, which is their collision cross-section, $\sigma$ (Topic 1B), and to the number densities $\mathcal{N}_{\mathrm{A}}$ and $\mathcal{N}_{\mathrm{B}}$ of A and B:

$$
v \propto \sigma(T / M)^{1 / 2} \mathcal{N}_{\mathrm{A}} \mathcal{N}_{\mathrm{B}} \propto \mathcal{N}_{\mathrm{N}} \propto[J]
$$

However, a collision is likely to be successful only if the kinetic energy of the molecules exceeds a minimum value, denoted $E^{\prime}$. This requirement suggests that the rate should also be proportional to a Boltzmann factor of the form $\mathrm{e}^{-E^{\prime} / R T}$ representing the fraction of collisions with at least the minimum required energy (Topic 17D). Therefore,

$$
v \propto \sigma(T / M)^{1 / 2} \mathrm{e}^{-E^{\prime} / R T}[\mathrm{~A}][\mathrm{B}]
$$

and, by writing the reaction rate in the form given in eqn 18A.1a, it follows that

$$
k_{\mathrm{r}} \propto \sigma(T / M)^{1 / 2} \mathrm{e}^{-E^{\prime} / R T}
$$

At this point, the form of the Arrhenius equation, eqn 18A.1b, begins to emerge, with the minimum kinetic energy $E^{\prime}$ identified as the activation energy $E_{\mathrm{a}}$ of the reaction. This identification, however, should not be regarded as precise, because collision theory is only a rudimentary model of chemical reactivity.

Not every collision will lead to reaction even if the energy requirement is satisfied, because the reactants might need to collide in a certain relative orientation. This 'steric requirement' suggests that a further factor, $P$, should be introduced, and that


\begin{equation*}
k_{\mathrm{r}} \propto P \sigma(T / M)^{1 / 2} \mathrm{e}^{-E^{\prime} / R T} \tag{18A.2}
\end{equation*}


As seen in detail below, this expression has the form predicted by collision theory. It reflects three aspects of a successful collision:\\
\includegraphics[max width=\textwidth, center]{2025_02_27_d4b95d9718f0b9e2e6b9g-812}

\subsection*{(a) Collision rates in gases}
As remarked, the reaction rate, and hence $k_{r}$, is expected to depend on the frequency with which molecules collide. The collision density, $Z_{A B}$, is the number of ( $\mathrm{A}, \mathrm{B}$ ) collisions in a region of the sample in an interval of time divided by the volume of the region and the duration of the interval. The frequency of collisions of a single molecule in a gas is calculated in Topic 1B (eqn 1B.12a, $z=\sigma v_{\text {rel }} \mathcal{V}_{A}$ ). That result can be adapted to derive an expression for $Z_{\mathrm{AB}}$.

\subsection*{How is that done? 18A. 1 \\
 Deriving an expression for the collision density}
The parameter $v_{\text {rel }}$ in the expression $z=\sigma v_{\text {rel }} \mathcal{N}$ is the mean relative speed of the colliding molecules and $\sigma$ is the collision cross-section: $\sigma=\pi d^{2}$, with $d=\frac{1}{2}\left(d_{\mathrm{A}}+d_{\mathrm{B}}\right)$, as shown in Fig. 18A.1. For collisions between A molecules of mass $m_{\mathrm{A}}$ and B molecules of mass $m_{\mathrm{B}}$, the mean relative speed is eqn 1B.11b $\left(v_{\text {rel }}=(8 k T / \pi \mu)^{1 / 2}\right.$, where $\left.\mu=m_{\mathrm{A}} m_{\mathrm{B}} /\left(m_{\mathrm{A}}+m_{\mathrm{B}}\right)\right)$. It follows that the collision rate of one A molecule with B molecules present at number density $\mathcal{N}_{\mathrm{B}}$ is $\sigma v_{\text {rel }} \mathcal{N}_{\mathrm{B}}$. The collision density is therefore this rate multiplied by the number density of A molecules, $\mathcal{N}_{\mathrm{A}}$ :


\begin{equation*}
Z_{\mathrm{AB}}=\sigma v_{\mathrm{rel}} \mathcal{N}_{\mathrm{A}} \mathcal{N}_{\mathrm{B}} \tag{18A.3}
\end{equation*}


The number density of a species J is $\mathcal{N}_{\mathrm{J}}=N_{\mathrm{A}}[\mathrm{J}]$, where [J] is its molar concentration and $N_{\mathrm{A}}$ is Avogadro's constant. It follows that


\begin{equation*}
\Longrightarrow Z_{\mathrm{AB}}=\sigma\left(\frac{8 k T}{\pi \mu}\right)^{1 / 2} N_{\mathrm{A}}^{2}[\mathrm{~A}][\mathrm{B}] \tag{18A.4a}
\end{equation*}


Collision density [KMT]\\
\includegraphics[max width=\textwidth, center]{2025_02_27_d4b95d9718f0b9e2e6b9g-813}

Figure 18A. 1 The collision cross-section for two molecules can be regarded to be the area within which the projectile molecule (A) must enter around the target molecule (B) in order for a collision to occur. If the diameters of the two molecules are $d_{\mathrm{A}}$ and $d_{\mathrm{B}}$, the radius of the target area is $d=\frac{1}{2}\left(d_{\mathrm{A}}+d_{\mathrm{B}}\right)$ and the cross-section is $\pi d^{2}$.

If the collision density is required in terms of the partial pressure of each gas J, then the molar concentrations in eqn

18A.4a are replaced by [J] $=n_{\mathrm{J}} / V=p_{\mathrm{J}} / R T$. For collisions between like molecules $\mu=\frac{1}{2} m_{\mathrm{A}}$ and eqn 18A.4a becomes

$$
\begin{aligned}
Z_{\mathrm{AA}} & =\frac{1}{2} \sigma\left(\frac{16 k T}{\pi m_{\mathrm{A}}}\right)^{1 / 2} N_{\mathrm{A}}^{2}[\mathrm{~A}]^{2} \\
& =\sigma\left(\frac{4 k T}{\pi m_{\mathrm{A}}}\right)^{1 / 2} N_{\mathrm{A}}^{2}[\mathrm{~A}]^{2}
\end{aligned}
$$

Collision density [identical molecules]\\
(18A.4b)\\
where the (blue) factor of $\frac{1}{2}$ has been introduced to avoid double counting of collisions (Smith with Jones and Jones with Smith, for instance).

\subsection*{Brief illustration 18A. 1}
In nitrogen at $25^{\circ} \mathrm{C}$ and 1.0 bar, when $\left[\mathrm{N}_{2}\right] \approx 40 \mathrm{~mol} \mathrm{~m}^{-3}$, with $\sigma=0.43 \mathrm{~nm}^{2}$ and $m_{\mathrm{N}_{2}}=28.02 m_{\mathrm{u}}$ the collision density is

$$
\begin{aligned}
Z_{\mathrm{N}_{2} \mathrm{~N}_{2}}= & \left(4.3 \times 10^{-19} \mathrm{~m}^{2}\right) \times\left(\frac{4 \times\left(1.381 \times 10^{-23} \mathrm{JK}^{-1}\right) \times(298 \mathrm{~K})}{\pi \times 28.02 \times\left(1.661 \times 10^{-27} \mathrm{~kg}\right)}\right)^{1 / 2} \\
& \times\left(6.022 \times 10^{23} \mathrm{~mol}^{-1}\right)^{2} \times\left(40 \mathrm{molm}^{-3}\right)^{2} \\
= & 8.4 \times 10^{34} \mathrm{~m}^{-3} \mathrm{~s}^{-1}
\end{aligned}
$$

This result shows that collision densities may be very large: even in $1 \mathrm{~cm}^{3}$, there are over $8 \times 10^{16}$ collisions in each picosecond.

\subsection*{(b) The energy requirement}
According to collision theory, the rate of change of $\mathcal{N}_{\mathrm{A}}$ due to reaction is the product of the collision density and the probability that a collision occurs with sufficient energy. The latter condition can be incorporated by writing the collision crosssection $\sigma$ as a function of the kinetic energy $\varepsilon$ of approach of the two colliding species, and setting the cross-section, $\sigma(\varepsilon)$, equal to zero if the kinetic energy of approach is below a certain threshold value, $\varepsilon_{\mathrm{a}}$. Later, $N_{\mathrm{A}} \varepsilon_{\mathrm{a}}$ will be identified as $E_{\mathrm{a}}$, the (molar) activation energy of the reaction. For a collision between A and B with a specific relative speed of approach $v_{\text {rel }}$ (not, at this stage, a mean value) it follows from eqn 18A. 3 that the rate of change of $\mathcal{N}_{\mathrm{A}}$ is


\begin{equation*}
\frac{\mathrm{d} \mathcal{N}_{\mathrm{A}}}{\mathrm{~d} t}=-\sigma(\varepsilon) v_{\mathrm{rel}} \mathcal{N}_{\mathrm{A}} \mathcal{N}_{\mathrm{B}} \tag{18A.5a}
\end{equation*}


or, in terms of molar concentrations,


\begin{equation*}
\frac{\mathrm{d}[\mathrm{~A}]}{\mathrm{d} t}=-\sigma(\varepsilon) v_{\mathrm{rel}} N_{\mathrm{A}}[\mathrm{~A}][\mathrm{B}] \tag{18A.5b}
\end{equation*}


The kinetic energy associated with the relative motion of the two particles is $\varepsilon=\frac{1}{2} \mu v_{\text {rel }}$. Therefore the relative speed can also be expressed in terms of the relative kinetic energy as $v_{\text {rel }}=$ $(2 \varepsilon / \mu)^{1 / 2}$. Because there is a wide range of approach energies\\
$\varepsilon$ in a sample, eqn 18A.5b must be averaged over a Boltzmann distribution of energies $f(\varepsilon)$ to give


\begin{equation*}
\frac{\mathrm{d}[\mathrm{~A}]}{\mathrm{d} t}=-\left\{\int_{0}^{\infty} \sigma(\varepsilon) v_{\mathrm{rel}} f(\varepsilon) \mathrm{d} \varepsilon\right\} N_{\mathrm{A}}[\mathrm{~A}][\mathrm{B}] \tag{18A.6}
\end{equation*}


where $f(\varepsilon) \mathrm{d} \varepsilon$ is the probability that the approach energy is between $\varepsilon$ and $\varepsilon+\mathrm{d} \varepsilon$. By comparison with eqn 18A.1a it follows that


\begin{equation*}
k_{\mathrm{r}}=N_{\mathrm{A}} \int_{0}^{\infty} \sigma(\varepsilon) \nu_{\mathrm{rel}} f(\varepsilon) \mathrm{d} \varepsilon \quad \text { Rate constant } \tag{18A.7}
\end{equation*}


To evaluate this integral it is necessary to establish the energy dependence of the collision cross-section, $\sigma(\varepsilon)$.

\subsection*{How is that done? 18A. 2 Deriving an expression for the energy dependence of the collision cross-section}
The key aspect of this model is that in a collision only the kinetic energy associated with a head-on collision is effective at bringing about reaction.

Step 1 Consider how the geometry of the collision affects the energy available for reaction\\
Consider two molecules A and B colliding with relative speed $v_{\text {rel }}$ and therefore relative kinetic energy $\varepsilon=\frac{1}{2} \mu \nu_{\text {rel }}^{2}$. Although a collision is counted when the centres of the molecules come within a distance $d$ of each other, that might be more a glancing blow than a head-on collision. Intuitively a head-on collision between A and B will be most effective in bringing about a chemical reaction, so from now on suppose that only the kinetic energy associated with the head-on component of the collision leads to reaction. This contribution to the kinetic energy depends on $v_{\text {rel, } A-B}$, the magnitude of the relative velocity component parallel to an axis connecting the centres of A and $B$.

Step 2 Find an expression for the head-on component of the velocity\\
As shown in the arrangement in Fig. 18A.2, the distance $a$ is the closest approach of the centres of the two molecules and $d$ is the distance between the centres. From trigonometry and the definition of the angle $\theta$ given in the diagram, it follows that

$$
v_{\text {rel }, \mathrm{A}-\mathrm{B}}=v_{\mathrm{rel}} \cos \theta=v_{\mathrm{rel}}\left(\frac{d^{2}-a^{2}}{d^{2}}\right)^{1 / 2}
$$

\subsection*{Step 3 Relate the velocities to energies}
The kinetic energy associated with the head-on collision is $\varepsilon_{\mathrm{A}-\mathrm{B}}=\frac{1}{2} \mu \nu_{\mathrm{rel}, \mathrm{A}-\mathrm{B}}^{2}$, and the total kinetic energy of the collision\\
\includegraphics[max width=\textwidth, center]{2025_02_27_d4b95d9718f0b9e2e6b9g-814}

Figure 18A. 2 The parameters used in the calculation of the relative kinetic energy associated with the head-on component of the collision of two molecules A and B.\\
is $\varepsilon=\frac{1}{2} \mu v_{\mathrm{rel}}^{2}$. You can relate these two quantities by using the result from Step 2:\\
$\boldsymbol{\varepsilon}_{\mathrm{A}-\mathrm{B}}=\frac{1}{2} \mu v_{\mathrm{rel}, \mathrm{A}-\mathrm{B}}^{2}=\frac{1}{2} \mu\left(v_{\text {rel }} \cos \theta\right)^{2}=\overbrace{\frac{1}{2} \mu v_{\text {rel } \mathrm{A}-\mathrm{B}}^{2}=v_{\text {rel }} \cos \theta}^{\varepsilon} \cos ^{2} \theta=\varepsilon \cos ^{\varepsilon}=\left(\frac{d^{2}-a^{2}}{d^{2}}\right)$

\subsection*{Step 4 Introduce an energy threshold}
As $a$ increases, the kinetic energy associated with the head-on collision decreases. The existence of an energy threshold, $\varepsilon_{\mathrm{a}}$, for the formation of products implies that there is a maximum value of $a, a_{\text {max }}$, above which reaction does not occur. Therefore, set $a=a_{\text {max }}$ and $\varepsilon_{\mathrm{A}-\mathrm{B}}=\varepsilon_{\mathrm{a}}$ and obtain

$$
\varepsilon_{\mathrm{a}}=\varepsilon\left(\frac{d^{2}-a_{\max }^{2}}{d^{2}}\right) \text { which rearranges to } a_{\max }^{2}=\left(1-\frac{\varepsilon_{\mathrm{a}}}{\varepsilon}\right) d^{2}
$$

Step 5 Rewrite the expression in terms of the collision crosssection\\
The energy-dependent collision cross-section is given in terms of $a_{\text {max }}$ as $\sigma(\varepsilon)=\pi a_{\text {max }}^{2}$, and $\pi d^{2}$ is identified as the (simple) collision cross-section $\sigma$ introduced in Fig. 18A.1. It follows that

$$
-\left|\sigma(\varepsilon)=\left(1-\frac{\varepsilon_{\mathrm{a}}}{\varepsilon}\right) \sigma\right| \begin{aligned}
& \begin{array}{l}
\text { Energy dependence of } \sigma(\varepsilon) \\
{\left[\varepsilon>\varepsilon_{\mathrm{a}}\right]}
\end{array}
\end{aligned}
$$

This form of the energy-dependence of $\sigma(\varepsilon)$ is broadly consistent with experimental determinations of the reaction between H and $\mathrm{D}_{2}$ as determined by molecular beam measurements of the kind described in Topic 18D (Fig. 18A.3).\\
With the energy dependence of the collision cross-section established, the rate constant can now be calculated.\\
\includegraphics[max width=\textwidth, center]{2025_02_27_d4b95d9718f0b9e2e6b9g-815}

Figure 18A. 3 The variation of the reactive cross-section with energy as expressed by eqn 18A.8. The data points are from experiments on the reaction $H+D_{2} \rightarrow H D+D$ (K. Tsukiyama et al., J. Chem. Phys. 84, 1934 (1986)).

\subsection*{How is that done? 18A. 3}
Deriving an expression for the rate constant

The calculation involves evaluating the integral in eqn 18A. 7 with the energy-dependent collision cross section given in eqn 18A.8.\\
Step 1 Use the Maxwell-Boltzmann distribution to write an expression for $f(\varepsilon) \mathrm{d} \varepsilon$\\
Adapt eqn 1B. 4 in Topic 1B by replacing $M / R$ with $\mu / k$ and so writing the distribution of relative molecular speeds as

$$
f\left(v_{\text {rel }}\right) \mathrm{d} v_{\text {rel }}=4 \pi\left(\frac{\mu}{2 \pi k T}\right)^{3 / 2} v_{\text {rel }}^{2} \mathrm{e}^{-\mu u_{\mathrm{rel}}^{2} / 2 k T} \mathrm{~d} v_{\text {rel }}
$$

You can write this expression in terms of the relative kinetic energy, $\varepsilon$, by noting that $\varepsilon=\frac{1}{2} \mu v_{\text {rel }}^{2}$. It follows that $v_{\text {rel }}=(2 \varepsilon / \mu)^{1 / 2}$ and so $\mathrm{d} v_{\text {rel }}=\frac{1}{2}(2 / \mu)^{1 / 2} \varepsilon^{-1 / 2} \mathrm{~d} \varepsilon=\mathrm{d} \varepsilon /(2 \mu \varepsilon)^{1 / 2}$. With this substitution the distribution becomes

$$
\begin{aligned}
f\left(v_{\text {rel }}\right) \mathrm{d} v_{\text {rel }} & =4 \pi\left(\frac{\mu}{2 \pi k T}\right)^{3 / 2}\left(\frac{2 \varepsilon}{\mu}\right) \mathrm{e}^{-\varepsilon / k T} \frac{\mathrm{~d} \varepsilon}{(2 \mu \varepsilon)^{1 / 2}} \\
& =\underbrace{2 \pi\left(\frac{1}{\pi k T}\right)^{3 / 2} \varepsilon^{1 / 2} \mathrm{e}^{-\varepsilon / k T} \mathrm{~d} \varepsilon}_{f(\varepsilon) \mathrm{d} \varepsilon}
\end{aligned}
$$

Step 2 Evaluate the integral\\
Now evaluate the integral

$$
\begin{aligned}
\int_{0}^{\infty} \sigma(\varepsilon) \overbrace{v_{\mathrm{rel}}}^{(2 \varepsilon / \mu)^{1 / 2}} f(\varepsilon) \mathrm{d} \varepsilon & =2 \pi\left(\frac{1}{\pi k T}\right)^{3 / 2} \int_{0}^{\infty} \sigma(\varepsilon)\left(\frac{2 \varepsilon}{\mu}\right)^{1 / 2} \varepsilon^{1 / 2} \mathrm{e}^{-\varepsilon / k T} \mathrm{~d} \varepsilon \\
& =\left(\frac{8}{\pi \mu k T}\right)^{1 / 2}\left(\frac{1}{k T}\right) \int_{0}^{\infty} \varepsilon \sigma(\varepsilon) \mathrm{e}^{-\varepsilon / k T} \mathrm{~d} \varepsilon
\end{aligned}
$$

With $\sigma(\varepsilon)$ from eqn 18A.8, write

$$
\left.\begin{array}{rl}
\int_{0}^{\infty} \varepsilon \sigma(\varepsilon) \mathrm{e}^{-\varepsilon / k T} \mathrm{~d} \varepsilon & =\sigma \int_{\varepsilon_{\mathrm{a}}}^{\infty} \varepsilon\left(1-\frac{\varepsilon_{\mathrm{a}}}{\varepsilon}\right) \mathrm{e}^{-\varepsilon / k T} \mathrm{~d} \varepsilon \\
& =\sigma\{\overbrace{\int_{\varepsilon_{\mathrm{a}}}^{\infty} \varepsilon \mathrm{e}^{-\varepsilon / k T} \mathrm{~d} \varepsilon} \varepsilon
\end{array}\right)
$$

It follows that

$$
\int_{0}^{\infty} \sigma(\varepsilon) v_{\text {rel }} f(\varepsilon) \mathrm{d} \varepsilon=\sigma\left(\frac{8 k T}{\pi \mu}\right)^{1 / 2} \mathrm{e}^{-\varepsilon_{\mathrm{a}} / k T}
$$

Step 3 Finalize the expression for the rate constant\\
Because $E_{\mathrm{a}}=N_{\mathrm{A}} \varepsilon_{\mathrm{a}}$ it follows that $\varepsilon_{\mathrm{a}} / k T=E_{\mathrm{a}} / R T$. With these substitutions it follows from the integral just evaluated and eqn 18A. 7 that

$$
-k_{\mathrm{r}}=\sigma N_{\mathrm{A}}\left(\frac{8 k T}{\pi \mu}\right)^{1 / 2} \mathrm{e}^{-E_{\mathrm{a}} / R T} \quad \begin{aligned}
& \text { Rate constant } \\
& \text { [collision theory] }
\end{aligned}
$$

This equation has the Arrhenius form $k_{\mathrm{r}}=A \mathrm{e}^{-E_{\mathrm{a}} / R T}$ provided the exponential temperature dependence dominates the weak square-root temperature dependence of the frequency factor. It follows that, within the constraints of collision theory, the activation energy, $E_{\mathrm{a}}$, can be identified with the minimum kinetic energy along the line of approach that is needed for reaction, and that the frequency factor (after multiplication by $[\mathrm{A}][\mathrm{B}]$ ) determines the rate at which collisions occur.

The simplest procedure for calculating $k_{\mathrm{r}}$ is to use for $\sigma$ the values obtained for non-reactive collisions (e.g. typically those obtained from viscosity measurements) or from tables of molecular radii. If the collision cross-sections of A and B are $\sigma_{\mathrm{A}}$ and $\sigma_{\mathrm{B}}$, then an approximate value of the AB cross-section is estimated from $\sigma=\pi d^{2}$, with $d=\frac{1}{2}\left(d_{\mathrm{A}}+d_{\mathrm{B}}\right)$. That is,

$$
\sigma \approx \frac{1}{4}\left(\sigma_{\mathrm{A}}^{1 / 2}+\sigma_{\mathrm{B}}^{1 / 2}\right)^{2}
$$

\subsection*{Brief illustration 18A. 2}
To estimate the rate constant for the reaction $\mathrm{H}_{2}+\mathrm{C}_{2} \mathrm{H}_{4} \rightarrow$ $\mathrm{C}_{2} \mathrm{H}_{6}$ at 628 K , first calculate $\mu$ by setting $m\left(\mathrm{H}_{2}\right)=2.016 m_{\mathrm{u}}$ and $m\left(\mathrm{C}_{2} \mathrm{H}_{4}\right)=28.05 m_{\mathrm{u}}$. A straightforward calculation gives $\mu=$ $3.123 \times 10^{-27} \mathrm{~kg}$. It then follows that

$$
\left(\frac{8 k T}{\pi \mu}\right)^{1 / 2}=\left(\frac{8 \times\left(1.381 \times 10^{-23} \mathrm{JK}^{-1}\right) \times(628 \mathrm{~K})}{\pi \times\left(3.123 \times 10^{-27} \mathrm{~kg}\right)}\right)^{1 / 2}=2.65 \ldots \mathrm{~km} \mathrm{~s}^{-1}
$$

From Table 1B.1, $\sigma\left(\mathrm{H}_{2}\right)=0.27 \mathrm{~nm}^{2}$ and $\sigma\left(\mathrm{C}_{2} \mathrm{H}_{4}\right)=0.64 \mathrm{~nm}^{2}$, giving $\sigma\left(\mathrm{H}_{2}, \mathrm{C}_{2} \mathrm{H}_{4}\right) \approx 0.44 \mathrm{~nm}^{2}$. The activation energy for this reaction is $180 \mathrm{~kJ} \mathrm{~mol}^{-1}$; therefore,

$$
\begin{aligned}
k_{\mathrm{r}}= & \left(4.4 \times 10^{-19} \mathrm{~m}^{2}\right) \times\left(6.022 \times 10^{23} \mathrm{~mol}^{-1}\right) \times\left(2.65 \ldots \times 10^{3} \mathrm{~m} \mathrm{~s}^{-1}\right) \\
& \times \mathrm{e}^{-\left(1.80 \times 10^{5} \mathrm{Jmol}^{-1}\right) /\left(8.3145 \mathrm{JK} \mathrm{~K}^{-1} \mathrm{~mol}^{-1}\right) \times(628 \mathrm{~K})} \\
= & 7.05 \ldots \times 10^{8} \mathrm{~m}^{3} \mathrm{~mol}^{-1} \mathrm{~s}^{-1} \times \mathrm{e}^{-34.5 \ldots}=7.5 \times 10^{-7} \mathrm{~m}^{3} \mathrm{~mol}^{-1} \mathrm{~s}^{-1}
\end{aligned}
$$

or $7.5 \times 10^{-4} \mathrm{dm}^{3} \mathrm{~mol}^{-1} \mathrm{~s}^{-1}$.

\subsection*{(c) The steric requirement}
Table 18A. 1 compares some values of the frequency factor calculated from collision cross-sections determined in other measurements with values obtained from Arrhenius plots. One of the reactions shows fair agreement between theory and experiment, but for others there are major discrepancies. In some cases the experimental values are orders of magnitude smaller than those calculated, which suggests that the collision energy is not the only criterion for reaction and that some other feature, such as the relative orientation of the colliding species, is important. Moreover, one reaction in the table has a pre-exponential factor larger than theory, which seems to indicate that the reaction occurs more quickly than the particles collide!\\
The disagreement between experiment and theory can be eliminated by introducing a steric factor, $P$, and expressing the reactive cross-section, $\sigma^{*}$, the actual cross-section for reactive collisions, as a multiple of the collision cross-section, $\sigma^{\star}=P \sigma$ (Fig. 18A.4). Then the rate constant becomes


\begin{equation*}
k_{\mathrm{r}}=P \sigma N_{\mathrm{A}}\left(\frac{8 k T}{\pi \mu}\right)^{1 / 2} \mathrm{e}^{-E_{\mathrm{a}} / R T} \tag{18A.10}
\end{equation*}


This expression has the form anticipated in eqn 18A.2. The steric factor is normally found to be several orders of magnitude smaller than 1.

Table 18A. 1 Arrhenius parameters for gas-phase reactions*

\begin{center}
\begin{tabular}{lllll}
\hline
 & \multicolumn{2}{c}{$A /\left(\mathrm{dm}^{3} \mathrm{~mol}^{-1} \mathrm{~s}^{-1}\right)$} & \multirow{2}{*}{$E_{\mathrm{a}} /\left(\mathbf{k J ~ m o l}^{-1}\right)$} & $P$ \\
\cline { 2 - 3 }
 & Experiment & Theory &  &  \\
\hline
$2 \mathrm{NOCl} \rightarrow$ & $9.4 \times 10^{9}$ & $5.9 \times 10^{10}$ & 102 & 0.16 \\
$2 \mathrm{NO}+2 \mathrm{Cl}$ &  &  &  &  \\
$2 \mathrm{ClO} \rightarrow \mathrm{Cl}_{2}+\mathrm{O}_{2}$ & $6.3 \times 10^{7}$ & $2.5 \times 10^{10}$ & 0 & $2.5 \times 10^{-3}$ \\
$\mathrm{H}_{2}+\mathrm{C}_{2} \mathrm{H}_{4} \rightarrow \mathrm{C}_{2} \mathrm{H}_{6}$ & $1.24 \times 10^{6}$ & $7.4 \times 10^{11}$ & 180 & $1.7 \times 10^{-6}$ \\
$\mathrm{~K}+\mathrm{Br}_{2} \rightarrow \mathrm{KBr}+\mathrm{Br}$ & $1.0 \times 10^{12}$ & $2.1 \times 10^{11}$ & 0 & 4.8 \\
\hline
\end{tabular}
\end{center}

\footnotetext{${ }^{*}$ More values are given in the Resource section.
}
\includegraphics[max width=\textwidth, center]{2025_02_27_d4b95d9718f0b9e2e6b9g-816}

Figure 18A. 4 The collision cross-section is the target area that results in simple deflection of the projectile molecule; the reactive cross-section is the corresponding area for chemical change to occur on collision.

\subsection*{Brief illustration 18A. 3}
It is found experimentally that the frequency factor for the reaction $\mathrm{H}_{2}+\mathrm{C}_{2} \mathrm{H}_{4} \rightarrow \mathrm{C}_{2} \mathrm{H}_{6}$ at 628 K is $1.24 \times 10^{6} \mathrm{dm}^{3} \mathrm{~mol}^{-1} \mathrm{~s}^{-1}$. The result in Brief illustration 18A. 2 can be expressed as $A=$ $7.05 \ldots \times 10^{11} \mathrm{dm}^{3} \mathrm{~mol}^{-1} \mathrm{~s}^{-1}$. It follows that the steric factor for this reaction is

$$
P=\frac{A_{\text {experimental }}}{A_{\text {calculated }}}=\frac{1.24 \times 10^{6} \mathrm{dm}^{3} \mathrm{~mol}^{-1} \mathrm{~s}^{-1}}{7.05 \ldots \times 10^{11} \mathrm{dm}^{3} \mathrm{~mol}^{-1} \mathrm{~s}^{-1}} \approx 1.8 \times 10^{-6}
$$

The very small value of $P$ is one reason, the other being the high activation energy, why catalysts are needed to bring this reaction about at a reasonable rate. As a general guide, the more complex the reactant molecules, the smaller is the value of $P$.

An example of a reaction for which it is possible to estimate the steric factor is $\mathrm{K}+\mathrm{Br}_{2} \rightarrow \mathrm{KBr}+\mathrm{Br}$, for which $P=$ 4.8. In this reaction, the distance of approach at which reaction occurs appears to be considerably larger than the distance needed for deflection of the path of the approaching molecules in a non-reactive collision. It has been proposed that the reaction proceeds by a harpoon mechanism. This brilliant name is based on a model of the reaction in which the K atom is pictured as approaching a $\mathrm{Br}_{2}$ molecule, and when the two are close enough an electron (the harpoon) flips across from K to $\mathrm{Br}_{2}$. In place of two neutral particles there are now two ions, so there is a Coulombic attraction between them: this attraction is the line on the harpoon. Under its influence the ions move together (the line is wound in), the reaction takes place, and $\mathrm{KBr}+\mathrm{Br}$ emerge. The harpoon extends the cross-section for the reactive encounter, and the reaction rate is significantly underestimated by taking for the collision cross-section the value for simple mechanical contact between K and $\mathrm{Br}_{2}$.

\subsection*{Example 18A． 1 Estimating a steric factor}
Estimate the value of $P$ for the harpoon mechanism by calcu－ lating the distance at which it becomes energetically favour－ able for the electron to leap from K to $\mathrm{Br}_{2}$ ．Take the sum of the radii of the reactants（treating them as spherical）to be 400 pm ．

Collect your thoughts Begin by identifying all the energy terms involved in the electron transfer process $\mathrm{K}+\mathrm{Br}_{2} \rightarrow \mathrm{~K}^{+}+$ $\mathrm{Br}_{2}^{-}$．There are three terms：the first is the ionization energy， $I$ ，of K ；the second is the electron affinity，$E_{\text {ea }}$ ，of $\mathrm{Br}_{2}$ ；and the third is the Coulombic interaction energy between the ions when they have been formed．When the separation of the ions is $R$ ，the Coulombic attraction energy is $-e^{2} / 4 \pi \varepsilon_{0} R$ ．The electron flips across when the sum of these three contributions changes from positive to negative（that is，when the sum becomes zero）so making the process energetically favourable．

The solution The net change in energy when the transfer occurs at a separation $R$ is

$$
E=I-E_{\mathrm{ea}}-\frac{e^{2}}{4 \pi \varepsilon_{0} R}
$$

This energy is zero when $R$ is equal to some critical value $R^{*}$ （and is negative for smaller values of $R$ ）

$$
0=I-E_{\mathrm{ea}}-\frac{e^{2}}{4 \pi \varepsilon_{0} R^{\star}} \text { rearranges to } R^{\star}=\frac{e^{2}}{4 \pi \varepsilon_{0}\left(I-E_{\mathrm{ea}}\right)}
$$

When the particles are at this separation，the harpoon shoots across from K to $\mathrm{Br}_{2}$ ．The reactive cross－section can therefore be identified as $\sigma^{*}=\pi R^{* 2}$ ．The non－reactive collision cross－ section is $\sigma=\pi d^{2}$ ，where $d=R(\mathrm{~K})+R\left(\mathrm{Br}_{2}\right)$ is the sum of the radii of the（assumed）spherical reactants．These values of $\sigma$ and $\sigma^{*}$ imply that the steric factor is

$$
P=\frac{\sigma^{\star}}{\sigma}=\frac{\pi R^{\star 2}}{\pi d^{2}}=\left\{\frac{e^{2}}{4 \pi \varepsilon_{0} d\left(I-E_{\text {ea }}\right)}\right\}^{2}
$$

With $I=420 \mathrm{~kJ} \mathrm{~mol}^{-1}$（corresponding to 0.70 aJ ），$E_{\text {ea }} \approx 250 \mathrm{~kJ}$ $\mathrm{mol}^{-1}$（corresponding to 0.42 aJ ），and $d=400 \mathrm{pm}$ ，the value of $P$ is 4．2，in good agreement with the experimental value（4．8）．

Self－test 18A．1 Estimate the value of $P$ for the harpoon reac－ tion between Na and $\mathrm{Cl}_{2}$ for which $d \approx 350 \mathrm{pm}$ ；take $E_{\text {ea }} \approx$ $230 \mathrm{~kJ} \mathrm{~mol}^{-1}$ ．

Example 18A． 1 illustrates two points about steric factors．First， the concept of a steric factor is not wholly useless because in some cases its numerical value can be estimated．Second，and more pessimistically，most reactions are much more complex than $\mathrm{K}+\mathrm{Br}_{2}$ ，and $P$ cannot be obtained so easily．

\subsection*{18A． 2 The RRK model}
The rate constants of＇unimolecular＇gas－phase reactions like those treated by the Lindemann－Hinshelwood mechanism （Topic 17F）can be estimated with a calculation based on the Rice－Ramsperger－Kassel model（RRK model）．That model was proposed in 1926 by O．K．Rice and H．C．Ramsperger and almost simultaneously by L．S．Kassel．It has been elaborated，largely by R．A．Marcus，into the＇RRKM model＇．The essential feature of the model is that although a molecule might have enough energy to react，that energy is distributed over all the modes of motion of the molecule，and reaction will occur only when enough of that energy has migrated into a particular location（such as a particu－ lar bond）in the molecule．The details are given in A deeper look 12 on the website of this book．Assuming that a molecule consists of $s$ identical harmonic oscillators，the principal conclusion is that the Kassel form of the unimolecular rate constant for the decay of the energized molecule $A^{*}$ to products is

\[
k_{\mathrm{b}}(E)=\left(1-\frac{E^{\star}}{E}\right)^{s-1} k_{\mathrm{b}} \text { for } E \geq E^{\star} \quad \begin{align*}
& \text { Unimolecular }  \tag{18A.11}\\
& \text { rate constant } \\
& \text { [Kassel form] }
\end{align*}
\]

where $k_{\mathrm{b}}$ is the rate constant used in the original Lindemann－ Hinshelwood theory for the decomposition of the energized molecule（Topic 17F），and $E^{\star}$ is the minimum energy that must be accumulated in a bond in order for it to break．

The energy dependence of the rate constant given by eqn 18A． 11 is shown in Fig．18A． 5 for various values of $s$ ．The equa－ tion can be interpreted as follows：\\
－The rate constant is smaller at a given excitation energy if $s$ is large，as it takes longer for the excitation energy to migrate through all the oscillators of a large molecule and accumulate in the location needed for reaction．\\
－As $E$ becomes very large，however，the term in paren－ theses approaches 1 ，and $k_{\mathrm{b}}(E)$ becomes independent of the energy and the number of oscillators in the molecule，as there is now enough energy to accumu－ late immediately in the critical mode regardless of the size of the molecule．\\
\includegraphics[max width=\textwidth, center]{2025_02_27_d4b95d9718f0b9e2e6b9g-817}

Figure 18A． 5 The energy dependence of the rate constant given by eqn 18A． 11 for three values of $s$ ．

\subsection*{Checklist of concepts}
1. In collision theory, it is supposed that the rate is proportional to the collision frequency, a steric factor, and the fraction of collisions that occur with at least the kinetic energy $E_{\mathrm{a}}$ along their lines of centres.2. The collision density is the number of collisions in a region of the sample in an interval of time divided by the volume of the region and the duration of the interval.3. The activation energy is the minimum kinetic energy along the line of approach of reactant molecules that is required for reaction.4. The steric factor is an adjustment that takes into account the orientational requirements for a successful collision.5. The rate constant for decomposition of an energized molecule can be estimated by using the RRK model.\subsection*{Checklist of equations}
\begin{center}
\begin{tabular}{llll}
\hline
Property & Equation & Comment & Equation number \\
\hline
Collision density & $Z_{\mathrm{AB}}=\sigma(8 k T / \pi \mu)^{1 / 2} N_{\mathrm{A}}^{2}[\mathrm{~A}][\mathrm{B}]$ & Unlike molecules, KMT (kinetic molecular theory) & 18 A .4 a \\
Energy dependence of $\sigma$ & $\sigma(\varepsilon)=\left(1-\varepsilon_{\mathrm{a}} / \varepsilon\right) \sigma$ & $\varepsilon \geq \varepsilon_{\mathrm{a}}, \sigma=0$ otherwise & 18 A .8 \\
Rate constant & $k_{\mathrm{r}}=P \sigma N_{\mathrm{A}}(8 k T / \pi \mu)^{1 / 2} \mathrm{e}^{-E_{\mathrm{a}} / R T}$ & KMT, collision theory & 18 A .10 \\
Unimolecular rate constant & $k_{\mathrm{b}}(E)=\left(1-E^{\star} / E\right)^{s-1} k_{\mathrm{b}}$ & RRK theory $\left(E \geq E^{*}\right)$ & 18 A .11 \\
\hline
\end{tabular}
\end{center}

\section*{TOPIC 18B Diffusion-controlled reactions }
\subsection*{Why do you need to know this material?}
Most chemical reactions take place in solution and for a thorough grasp of chemistry it is important to understand what controls their rates and how those rates can be modified.

\subsection*{- What is the key idea?}
The rate of a chemical reaction in solution is controlled either by the rate of diffusion of the reactants or by the activation energy of the step that leads to products.

\subsection*{What do you need to know already?}
This Topic makes use of the steady-state approximation (Topic 17E) and draws on Fick's first law of diffusion (Topic 16C). At one point it uses the Stokes-Einstein relation (Topic 16C).

Reactions in solution are entirely different from those in gases. No longer are there collisions of molecules hurtling through space; now there is the jostling of one molecule through a dense but mobile collection of molecules making up the fluid environment.

\subsection*{18B. 1 Reactions in solution}
Encounters between reactants in solution occur in a very different manner from encounters in gases. The encounters of reactant molecules dissolved in a solvent are considerably less frequent than in a gas. However, because a molecule also migrates only slowly away from a location, two reactant molecules that encounter each other stay near each other for much longer than in a gas. This lingering of one molecule near another on account of the hindering presence of solvent molecules is called the cage effect. Such an encounter pair may accumulate enough energy to react even though it does not have enough energy to do so when it first forms. The activation energy of a reaction is a much more complicated quantity\\
in solution than in a gas because the encounter pair is surrounded by solvent and the energy of the entire local assembly of reactant and solvent molecules must be considered.

\subsection*{(a) Classes of reaction}
The complicated overall process can be divided into simpler parts by setting up a simple kinetic scheme. Suppose the rate of formation of an encounter pair AB is first order in each of the reactants A and B :

$$
\mathrm{A}+\mathrm{B} \rightarrow \mathrm{AB} \quad v=k_{\mathrm{d}}[\mathrm{~A}][\mathrm{B}]
$$

As will be seen, $k_{\mathrm{d}}$ (where the d signifies diffusion) is determined by the diffusional characteristics of A and B. The encounter pair can break up without reaction or it can go on to form products P. If it is supposed that both processes are pseudofirst-order reactions (with the solvent perhaps playing a role), then the mechanism may be written

$$
\begin{aligned}
& \mathrm{AB} \rightarrow \mathrm{~A}+\mathrm{B} \quad v=k_{\mathrm{d}}^{\prime}[\mathrm{AB}] \\
& \mathrm{AB} \rightarrow \mathrm{P} \quad v=k_{\mathrm{a}}[\mathrm{AB}]
\end{aligned}
$$

The concentration of AB can now be found by applying the steady-state approximation (Topic 17E) to the equation for the net rate of change of concentration of AB :

$$
\frac{\mathrm{d}[\mathrm{AB}]}{\mathrm{d} t}=k_{\mathrm{d}}[\mathrm{~A}][\mathrm{B}]-k_{\mathrm{d}}^{\prime}[\mathrm{AB}]-k_{\mathrm{a}}[\mathrm{AB}]=0
$$

his expression solves to

$$
[\mathrm{AB}]=\frac{k_{\mathrm{d}}[\mathrm{~A}][\mathrm{B}]}{k_{\mathrm{a}}+k_{\mathrm{d}}^{\prime}}
$$

The rate of formation of products is therefore


\begin{equation*}
\frac{\mathrm{d}[\mathrm{P}]}{\mathrm{d} t}=k_{\mathrm{a}}[\mathrm{AB}]=k_{\mathrm{r}}[\mathrm{~A}][\mathrm{B}] \quad k_{\mathrm{r}}=\frac{k_{\mathrm{a}} k_{\mathrm{d}}}{k_{\mathrm{a}}+k_{\mathrm{d}}^{\prime}} \tag{18B.1}
\end{equation*}


Two limits can now be distinguished. If the rate of separation of the unreacted encounter pair is much slower than the rate at which it forms products, then $k_{\mathrm{d}}^{\prime}[\mathrm{AB}] \ll k_{\mathrm{a}}[\mathrm{AB}]$ (or, after cancelling the [AB], $k_{\mathrm{d}}^{\prime} \ll k_{\mathrm{a}}$ ), and the effective rate constant is


\begin{equation*}
k_{\mathrm{r}} \approx \frac{k_{\mathrm{a}} k_{\mathrm{d}}}{k_{\mathrm{a}}}=k_{\mathrm{d}} \tag{18B.2a}
\end{equation*}


Diffusion-controlled limit

Table 18B. 1 Arrhenius parameters for solvolysis reactions in solution

\begin{center}
\begin{tabular}{llll}
\hline
 & Solvent & $A /\left(\mathrm{dm}^{3} \mathrm{~mol}^{-1} \mathbf{s}^{-1}\right)$ & $E_{\mathbf{a}} /\left(\mathrm{kJ} \mathrm{mol}^{-1}\right)$ \\
\hline
$\left(\mathrm{CH}_{3}\right)_{3} \mathrm{CCl}$ & Water & $7.1 \times 10^{16}$ & 100 \\
 & Ethanol & $3.0 \times 10^{13}$ & 112 \\
 & Chloroform & $1.4 \times 10^{4}$ & 45 \\
$\mathrm{CH}_{3} \mathrm{CH}_{2} \mathrm{Br}$ & Ethanol & $4.3 \times 10^{11}$ & 90 \\
\hline
\end{tabular}
\end{center}

In this diffusion-controlled limit, the rate of reaction is governed by the rate at which the reactant molecules diffuse through the solvent. Because the combination of radicals involves very little activation energy, radical and atom recombination reactions are often diffusion-controlled. An activation-controlled reaction arises when a substantial activation energy is involved in the reaction $\mathrm{AB} \rightarrow \mathrm{P}$. Then $k_{\mathrm{a}}[\mathrm{AB}]$ $\ll k_{\mathrm{d}}^{\prime}[\mathrm{AB}]$ (implying $k_{\mathrm{a}} \ll k_{\mathrm{d}}^{\prime}$ ) and


\begin{equation*}
k_{\mathrm{r}} \approx \frac{k_{\mathrm{a}} k_{\mathrm{d}}}{k_{\mathrm{d}}^{\prime}}=k_{\mathrm{a}} \frac{K}{c^{\ominus}} \quad \text { Activation-controlled limit } \tag{18B.2b}
\end{equation*}


where $k_{\mathrm{d}} / k_{\mathrm{d}}^{\prime}=K / c^{\ominus}$ (see Topic 17C) and $K$ is the equilibrium constant for $\mathrm{A}+\mathrm{B} \rightleftharpoons \mathrm{AB}$. In this limit, the reaction proceeds at a rate that depends on the equilibrium concentration of encounter pairs and the rate at which energy accumulates in these pairs from the surrounding solvent. Some experimental data are given in Table 18B.1.

\subsection*{(b) Diffusion and reaction}
The rate of a diffusion-controlled reaction is calculated by considering the rate at which the reactants diffuse together.

\subsection*{How is that done? 18B. 1 Finding an expression for the}
 rate constant of a diffusion-controlled reactionSuppose that molecules of A and B in solution react immediately when they come within some critical distance $R^{*}$ of one another and the rate of reaction is controlled by the rate of encounters between A and B molecules as they diffuse together. As a result of the reaction, the concentration of $B$ molecules near $A$ is decreased and a concentration gradient of B molecules is established. There is a diffusive flux of B towards A as a result of that gradient, and the flux is constant while the reaction is in progress.

Step 1 Consider the rate at which B molecules cross the surface of a sphere centred on $A$

If the (molar) flux of B molecules towards A is $J_{\mathrm{B}}$, the rate (expressed as amount divided by time) at which B molecules\\
pass through a shell of radius $r$ and surface area $4 \pi r^{2}$ centred on A is

$$
v_{\mathrm{B}}=4 \pi r^{2} J_{\mathrm{B}}=4 \pi r^{2} D_{\mathrm{B}} \frac{\mathrm{~d}[\mathrm{~B}](r)}{\mathrm{d} r}
$$

An important point to recognize is that $v_{\mathrm{B}}$ is the same for a shell of any radius greater than or equal to $R^{\star}$, because no B molecules are lost until they have reached $R^{*}$. Also keep in mind that $v_{\mathrm{B}}$ is a rate expressed as amount/time, not concentration/time.

Step 2 Use the known values of [B] for the bulk to establish an expression for the variation of the concentration with distance For $v_{\mathrm{B}}$ to be independent of $r, r^{2} \mathrm{~d}[\mathrm{~B}](r) / \mathrm{d} r$ must be a constant, which implies that, provided $r>R^{*},[\mathrm{~B}](r)=a+b / r$. Thus, $\mathrm{d}[\mathrm{B}](r) / \mathrm{d} r=-b / r^{2}$, and $r^{2} \mathrm{~d}[\mathrm{~B}](r) / \mathrm{d} r=-b$, a constant, as required. You can find the values of the constants $a$ and $b$ by noting that as $r \rightarrow \infty,[\mathrm{~B}](r)$ tends to its bulk value, $[\mathrm{B}]$. Therefore $a=[\mathrm{B}]$ and hence $[\mathrm{B}](r)=[\mathrm{B}]+b / r$. When $r=R^{*},[\mathrm{~B}](r)=0$, which implies that $b=-[\mathrm{B}] R^{*}$. It follows that

$$
[\mathrm{B}](r)=[\mathrm{B}]\left(1-\frac{R^{*}}{r}\right)
$$

Figure 18B.1 illustrates the distance dependence of $[\mathrm{B}](r)$ according to this equation. The first derivative of $[\mathrm{B}](r)$ with respect to distance is $[\mathrm{B}] R^{\star} / r^{2}$, so

$$
v_{\mathrm{B}}=4 \pi r^{2} D_{\mathrm{B}} \frac{\mathrm{~d}[\mathrm{~B}](r)}{\mathrm{d} r}=4 \pi R^{\star} D_{\mathrm{B}}[\mathrm{~B}]
$$

\subsection*{Step 3 Write an expression for the overall rate of reaction}
To express the rate of reaction, $v_{\mathrm{B}}$ must be multiplied by the number of A molecules in the solution. If the bulk concentration of A is $[\mathrm{A}]$, then the number of A molecules in a solution of volume $V$ is $N_{\mathrm{A}}[\mathrm{A}] V$. Therefore, the rate of reaction (still as amount/time) is\\
\includegraphics[max width=\textwidth, center]{2025_02_27_d4b95d9718f0b9e2e6b9g-820}

Figure 18B. 1 The concentration profile for reaction in solution when a molecule B diffuses towards another reactant molecule and reacts if it reaches $R^{*}$.

It is unrealistic to suppose that all A molecules are stationary, so the diffusion coefficient $D_{\mathrm{B}}$ is now replaced by the sum of the diffusion coefficients of the two species, $D=D_{\mathrm{A}}+$ $D_{\mathrm{B}}$. Because it is more convenient to express rates as concentration/time, both sides of the equation are divided by the volume $V$, in which case

$$
v=\overbrace{4 \pi R^{*} D N_{\mathrm{A}}}^{k_{\mathrm{d}}}[\mathrm{~A}][\mathrm{B}]
$$

from which it follows that the diffusion-controlled rate constant is\\
$-1 k_{\mathrm{d}}=4 \pi R^{\star} D N_{\mathrm{A}} \xlongequal{\text { Rate constant of a diffusion-controlled reaction }}$

\subsection*{Brief illustration 18B. 1}
The order of magnitude of $R^{*}$ is $10^{-10} \mathrm{~m}(100 \mathrm{pm})$ and that of $D$ for a species in water is $10^{-9} \mathrm{~m}^{2} \mathrm{~s}^{-1}$. It follows from eqn 18B. 3 that

$$
\begin{aligned}
k_{\mathrm{d}} & \approx 4 \pi \times\left(10^{-10} \mathrm{~m}\right) \times\left(10^{-9} \mathrm{~m}^{2} \mathrm{~s}^{-1}\right) \times\left(6.022 \times 10^{23} \mathrm{~mol}^{-1}\right) \\
& \approx 8 \times 10^{5} \mathrm{~m}^{3} \mathrm{~mol}^{-1} \mathrm{~s}^{-1}
\end{aligned}
$$

which corresponds to about $10^{9} \mathrm{dm}^{3} \mathrm{~mol}^{-1} \mathrm{~s}^{-1}$. An indication that a reaction is diffusion-controlled is therefore that its rate constant is of that order of magnitude.

Equation 18B. 3 can be taken further by incorporating the Stokes-Einstein equation (eqn 16C.4b of Topic 16C, $D_{\mathrm{J}}=$ $k T / 6 \pi \eta R_{\mathrm{J}}$ ) for the relation between the diffusion constant and the hydrodynamic radius $R_{\mathrm{A}}$ and $R_{\mathrm{B}}$ of each molecule in a medium of viscosity $\eta$. As this relation is approximate, little extra error is introduced by writing $R_{\mathrm{A}}=R_{\mathrm{B}}=\frac{1}{2} R^{*}$, which leads to

$$
k_{\mathrm{d}}=\frac{8 R T}{3 \eta}
$$

Diffusion-controlled rate constant\\
(18B.4)\\
(The $R$ in this equation is the gas constant.) The radii have cancelled because, although the diffusion constants are smaller when the radii are large, the reactive collision radius is larger and the particles need travel a shorter distance to meet. In this approximation, the rate constant is independent of the identities of the reactants, and depends only on the temperature and the viscosity of the solvent.

\subsection*{Brief illustration 18B. 2}
The rate constant for the recombination of I atoms in hexane at 298 K , when the viscosity of the solvent is 0.326 cP (with $1 \mathrm{P}=$ $10^{-1} \mathrm{~kg} \mathrm{~m}^{-1} \mathrm{~s}^{-1}$ ) is

$$
k_{\mathrm{d}}=\frac{8 \times\left(8.3145 \mathrm{JK}^{-1} \mathrm{~mol}^{-1}\right) \times(298 \mathrm{~K})}{3 \times\left(3.26 \times 10^{-4} \mathrm{~kg} \mathrm{~m}^{-1} \mathrm{~s}^{-1}\right)}=2.0 \times 10^{7} \mathrm{~m}^{3} \mathrm{~mol}^{-1} \mathrm{~s}^{-1}
$$

where $1 \mathrm{~J}=1 \mathrm{~kg} \mathrm{~m}^{2} \mathrm{~s}^{-2}$. This result corresponds to $2.0 \times$ $10^{10} \mathrm{dm}^{3} \mathrm{~mol}^{-1} \mathrm{~s}^{-1}$. The experimental value is $1.3 \times 10^{10} \mathrm{dm}^{3} \mathrm{~mol}^{-1} \mathrm{~s}^{-1}$, so the agreement is very good considering the approximations involved.

\subsection*{18B.2 The material-balance equation}
The diffusion of reactants plays an important role in many chemical processes, such as the diffusion of $\mathrm{O}_{2}$ molecules into red blood cells and the diffusion of a gas towards a catalyst. To catch a glimpse of the kinds of calculations involved consider the diffusion equation (Topic 16C) generalized to take into account the possibility that the diffusing, convecting molecules are also reacting.

\subsection*{(a) The formulation of the equation}
Consider a small volume element in a chemical reactor (or a biological cell). The net rate at which J molecules enter the region by diffusion and convection is given by eqn 16C. 9 of Topic 16C:


\begin{equation*}
\frac{\partial[\mathrm{J}]}{\partial t}=D \frac{\partial^{2}[\mathrm{~J}]}{\partial x^{2}}-v \frac{\partial[\mathrm{~J}]}{\partial x} \quad \quad \text { Diffusion equation } \tag{18B.5}
\end{equation*}


where $v$ is the velocity of the convective flow of J and [J] in general depends on both position and time. If J disappears by a pseudofirst-order reaction, the net rate of change of molar concentration due to chemical reaction is

$$
\frac{\partial[\mathrm{J}]}{\partial t}=-k_{\mathrm{r}}[\mathrm{~J}]
$$

Therefore, the overall rate of change of the concentration of $J$ is\\
\includegraphics[max width=\textwidth, center]{2025_02_27_d4b95d9718f0b9e2e6b9g-821}

Equation 18B. 6 is called the material-balance equation. If the rate constant is large, then [J] will decline rapidly. However, if the diffusion constant is large, then the decline can be replenished as J diffuses rapidly into the region. The convection term, which may represent the effects of stirring, can sweep material either into or out of the region according to the signs of $v$ and the concentration gradient $\partial[\mathrm{J}] / \partial x$.

\subsection*{(b) Solutions of the equation}
The material-balance equation is a second-order partial differential equation and is far from easy to solve in general. Some idea of how it is solved can be obtained by considering the special case in which there is no convective motion (as in an unstirred reaction vessel):


\begin{equation*}
\frac{\partial[\mathrm{J}]}{\partial t}=D \frac{\partial^{2}[\mathrm{~J}]}{\partial x^{2}}-k_{\mathrm{r}}[\mathrm{~J}] \tag{18B.7}
\end{equation*}


As may be verified by substitution (Problem 18B.1), if the solution of this equation in the absence of reaction (that is, for $k_{\mathrm{r}}=$ $0)$ is $[J](x, t)$, then the solution $[J]^{*}(x, t)$ in the presence of reaction $\left(k_{r}>0\right)$ is

$$
[\mathrm{J}]^{*}(x, t)=[\mathrm{J}](x, t) \mathrm{e}^{-k_{\mathrm{r} t} t} \quad \text { Diffusion with reaction }
$$

(18B.8)\\
An example of a solution of the diffusion equation in the absence of reaction is that given in Topic 16C for a system in which initially a layer of $n_{0} N_{\mathrm{A}}$ molecules is spread over a plane of area $A$ :


\begin{equation*}
[J](x, t)=\frac{n_{0} \mathrm{e}^{-x^{2} / 4 D t}}{A(\pi D t)^{1 / 2}} \tag{18B.9}
\end{equation*}


When this expression is substituted into eqn 18B.8, the result is an expression for the concentration of J as it diffuses away from its initial surface layer and undergoes reaction in the overlying solution (Fig. 18B.2).

\subsection*{Brief illustration 18B. 3}
Suppose 1.0 g of iodine $\left(3.9 \mathrm{mmol}_{2}\right)$ is spread over a surface of area $5.0 \mathrm{~cm}^{2}$ under a column of hexane ( $D=4.1 \times 10^{-9} \mathrm{~m}^{2} \mathrm{~s}^{-1}$ ). As it diffuses upwards it reacts with a pseudofirst-order rate constant $k_{\mathrm{r}}=4.0 \times 10^{-5} \mathrm{~s}^{-1}$. By substituting these values into

$$
[\mathrm{J}]^{*}(x, t)=\frac{n_{0} \mathrm{e}^{-x^{2} / 4 D t-k_{t} t}}{A(\pi D t)^{1 / 2}}
$$

the following table of values can be constructed:

\begin{center}
\begin{tabular}{llll}
\hline
 & \multicolumn{3}{c}{$[\mathrm{J}]^{*} /\left(\mathrm{moldm}^{-3}\right)$ at $x$} \\
\cline { 2 - 4 }
$t$ & $\mathbf{1 ~ m m}$ & $\mathbf{5 ~ m m}$ & $\mathbf{1 ~ c m ~}$ \\
\hline
100 s & 3.72 & 0 & 0 \\
1000 s & 1.96 & 0.45 & 0.005 \\
10000 s & 0.46 & 0.40 & 0.25 \\
\hline
\end{tabular}
\end{center}

Even this relatively simple example has led to an equation that is difficult to solve, and only in some special cases can the full material-balance equation be solved analytically. Most modern work on reactor design and cell kinetics uses numerical methods to solve the equation, and detailed solutions for realistic environments, such as vessels of different shapes (which influence the boundary conditions on the solutions) and with a variety of inhomogeneously distributed reactants, can be obtained reasonably easily.\\
\includegraphics[max width=\textwidth, center]{2025_02_27_d4b95d9718f0b9e2e6b9g-822}

Figure 18B. 2 The concentration profiles for a diffusing, reacting system (e.g. a column of solution) in which one reactant is initially in a layer at $x=0$. In the absence of reaction (grey lines) the concentration profiles are similar to those in Fig. 16C.5.

\subsection*{Checklist of concepts}
1. The cage effect, the lingering of one reactant molecule near another due to the hindering presence of solvent molecules, results in the formation of an encounter pair of reactant molecules.2. A reaction in solution may be diffusion controlled if its rate is controlled by the rate at which reactant molecules encounter each other in solution.\begin{enumerate}
  \setcounter{enumi}{2}
  \item The rate of an activation-controlled reaction is controlled by the rate at which the encounter pair accumulates sufficient energy.
  \item The material-balance equation relates the overall rate of change of the concentration of a species to its rates of diffusion, convection, and reaction.
\end{enumerate}

\subsection*{Checklist of equations}
\begin{center}
\begin{tabular}{llll}
\hline
Property & Equation & Comment & Equation number \\
\hline
Diffusion-controlled limit & $k_{\mathrm{r}}=k_{\mathrm{d}}$ & $v=k_{\mathrm{d}}[\mathrm{A}][\mathrm{B}]$ for the encounter rate & 18B.2a \\
Activation-controlled limit & $k_{\mathrm{r}}=k_{\mathrm{a}}\left(K / c^{\ominus}\right)$ & $K$ for $\mathrm{A}+\mathrm{B} \rightleftharpoons \mathrm{AB}, k_{\mathrm{a}}$ for the decomposition of AB & 18 B .2 b \\
Diffusion-controlled rate constant & $k_{\mathrm{d}}=4 \pi R^{*} D N_{\mathrm{A}}$ & $D=D_{\mathrm{A}}+D_{\mathrm{B}}$ & 18 B .3 \\
 & $k_{\mathrm{d}}=8 R T / 3 \eta$ & Assumes Stokes-Einstein relation & 18 B .4 \\
Material-balance equation & $\partial[\mathrm{J}] / \partial t=D \partial^{2}[\mathrm{~J}] / \partial x^{2}$ & Diffusion and convection with first-order reaction & 18B.6 \\
 & $-v \partial[J] / \partial x-k_{\mathrm{r}}[J]$ &  &  \\
\hline
\end{tabular}
\end{center}

\section*{TOPIC 18C Transition-state theory }
\subsection*{Why do you need to know this material?}
Transition-state theory provides a way to relate the rate constant of reactions to models of the cluster of atoms supposed to form when reactants come together. It provides a link between information about the structures of reactants and the rate constant for their reaction.

\subsection*{What is the key idea?}
Reactants come together to form an activated complex, which decays into products.

\subsection*{What do you need to know already?}
This Topic makes use of two strands: one is the relation between equilibrium constants and partition functions (Topic 13F); the other is the relation between equilibrium constants and thermodynamic functions, such as the Gibbs energy, enthalpy, and entropy of reaction (Topic 6A). You need to be aware of the Arrhenius equation for the temperature dependence of the rate constant (Topic 17D).

In transition-state theory (which is also widely referred to as activated complex theory), the notion of the transition state is used in conjunction with concepts of statistical thermodynamics to provide a more detailed calculation of rate constants than collision theory provides (Topic 18A). Transition-state theory has the advantage that a quantity corresponding to the steric factor appears automatically and does not need to be grafted on to an equation as an afterthought; it is an attempt to identify the principal features governing the size of a rate constant in terms of a model of the events that take place during the reaction.

\subsection*{18c. 1 The Eyring equation}
In the course of a chemical reaction that begins with an encounter between molecules of A and molecules of B, the potential energy of the system typically changes in a manner shown in Fig. 18C.1. Although the illustration displays an exothermic reaction, a potential barrier is also common for endothermic reactions. As the reaction event proceeds,\\
\includegraphics[max width=\textwidth, center]{2025_02_27_d4b95d9718f0b9e2e6b9g-824}

Figure 18C. 1 A potential energy profile for an exothermic reaction. The height of the barrier between the reactants and products is the activation energy of the reaction.

A and B come into contact, distort, and begin to exchange or discard atoms.

\subsection*{(a) The formulation of the equation}
The reaction coordinate is a representation of the atomic displacements, such as changes in interatomic distances and bond angles, that are directly involved in the formation of products from reactants. The potential energy rises to a maximum and the cluster of atoms that corresponds to the region close to the maximum is called the activated complex. After the maximum, the potential energy falls as the atoms rearrange in the cluster and reaches a value characteristic of the products. The climax of the reaction is at the peak of the potential energy, which can be identified with the activation energy $E_{\mathrm{a}}$. However, as in collision theory, this identification should be regarded as approximate and is clarified later. At this peak, two reactant molecules have come to such a degree of closeness and distortion that a small further distortion will send them in the direction of products. This crucial configuration is called the transition state of the reaction. Although some molecules entering the transition state might revert to reactants, if they pass through this configuration then it is inevitable that products will emerge from the encounter.

A note on good practice The terms activated complex and transition state are often used as synonyms; however, it is best to preserve the distinction, with the former referring to the cluster of atoms in the vicinity of the peak of the potential energy curve, and the latter to their critical configuration.\\
\includegraphics[max width=\textwidth, center]{2025_02_27_d4b95d9718f0b9e2e6b9g-825}

Figure 18C. 2 A reaction profile (for an exothermic reaction). The horizontal axis is the reaction coordinate, and the vertical axis is potential energy. The activated complex is the region near the potential maximum, and the transition state corresponds to the maximum itself.

Transition-state theory pictures a reaction between $A$ and $B$ as proceeding through the formation of an activated complex, $\mathrm{C}^{\ddagger}$, in a rapid pre-equilibrium (Fig. 18C.2):


\begin{equation*}
\mathrm{A}+\mathrm{B} \rightleftharpoons \mathrm{C}^{\ddagger} \quad K^{\ddagger}=\frac{p_{\mathrm{C}^{\ddagger}} p^{\ominus}}{p_{\mathrm{A}} p_{\mathrm{B}}} \tag{18C.1}
\end{equation*}


where for this gas-phase reaction the activity of each species has been replaced by $p / p^{\ominus}$. The development of femtosecond (and even attosecond) pulsed lasers has made it possible to make observations on species that have such short lifetimes that in a number of respects they resemble an activated complex, which often survive for only a few picoseconds.

When the partial pressures, $p_{\mathrm{J}}$, are expressed in terms of the molar concentrations, $[J]$, by using $p_{\mathrm{J}}=R T[\mathrm{~J}]$, the concentration of activated complex is related to the (dimensionless) equilibrium constant by


\begin{equation*}
\left[\mathrm{C}^{\ddagger}\right]=\frac{R T}{p^{\ominus}} K^{\ddagger}[\mathrm{A}][\mathrm{B}] \tag{18C.2}
\end{equation*}


The activated complex falls apart by unimolecular decay into products, P , with a rate constant $k^{\ddagger}$ :


\begin{equation*}
\mathrm{C}^{\ddagger} \rightarrow \mathrm{P} \quad v=k^{\ddagger}\left[\mathrm{C}^{\ddagger}\right] \tag{18C.3}
\end{equation*}


It follows that


\begin{equation*}
v=k_{\mathrm{r}}[\mathrm{~A}][\mathrm{B}] \quad k_{\mathrm{r}}=\frac{R T}{p^{\ominus}} k^{\ddagger} K^{\ddagger} \tag{18C.4}
\end{equation*}


The next task is the calculation of the unimolecular rate con$\operatorname{stant} k^{\ddagger}$ and the equilibrium constant $K^{\ddagger}$.

\subsection*{(b) The rate of decay of the activated complex}
An activated complex forms products only if it passes through the transition state. As the reactant molecules approach the activated complex region, some bonds are forming and shorten-\\
ing while others are lengthening and breaking; therefore, along the reaction coordinate, there is a vibration-like motion of the atoms in the activated complex. If this motion occurs with a frequency $v^{\ddagger}$, then the frequency with which the cluster of atoms forming the complex approaches the transition state is also $v^{\ddagger}$. However, it is possible that not every oscillation along the reaction coordinate takes the complex through the transition state. For instance, the centrifugal effect of rotations might also be an important contribution to the break-up of the complex, and in some cases the complex might be rotating too slowly or rotating rapidly but about the wrong axis. Therefore, it is more appropriate to suppose that the rate of passage of the complex through the transition state is only proportional, not equal, to the vibrational frequency along the reaction coordinate, and to write


\begin{equation*}
k^{\ddagger}=\kappa v^{\ddagger} \tag{18C.5}
\end{equation*}


where $\kappa$ (kappa) is the transmission coefficient. In the absence of information to the contrary, $\kappa$ is assumed to be about 1 .

\subsection*{Brief illustration 18C. 1}
Typical vibrations of small molecules occur at wavenumbers of the order of $10^{3} \mathrm{~cm}^{-1}$ ( $\mathrm{C}-\mathrm{H}$ bends, for example, occur in the range $1340-1465 \mathrm{~cm}^{-1}$ ) and therefore occur at frequencies of the order of $10^{13} \mathrm{~Hz}$. Suppose that the loosely bound cluster vibrates at one or two orders of magnitude lower frequency, then $v^{\ddagger} \approx 10^{11}-10^{12} \mathrm{~Hz}$. These figures suggest that $k^{\ddagger} \approx 10^{11}-$ $10^{12} \mathrm{~s}^{-1}$, with $\kappa$ perhaps reducing that value further.

\subsection*{(c) The concentration of the activated complex}
Topic 13F explains how to calculate equilibrium constants from structural data. Equation 13F.10b of that Topic (which expresses $K$ in terms of the standard molar partition functions $q_{j}^{\ominus}$ ) can be used directly, which in this case gives


\begin{equation*}
K^{\ddagger}=\frac{N_{\mathrm{A}} \mathscr{G}_{\mathrm{C}^{\ddagger}}^{\ominus}}{\mathfrak{q}_{\mathrm{A}}^{\ominus} \mathfrak{G}_{\mathrm{B}}^{\ominus}} \mathrm{e}^{-\Delta E_{0} / R T} \tag{18C.6}
\end{equation*}


with


\begin{equation*}
\Delta E_{0}=E_{0}\left(\mathrm{C}^{\ddagger}\right)-E_{0}(\mathrm{~A})-E_{0}(\mathrm{~B}) \tag{18C.7}
\end{equation*}


Note that the units of $N_{\mathrm{A}}$ and the $\mathscr{q}_{\mathrm{J}}^{\ominus}$ are $\mathrm{mol}^{-1}$, so $K^{\ddagger}$ is dimensionless (as is appropriate for an equilibrium constant).

The focus of the final step of this part of the calculation is the partition function of the activated complex. For the special vibration of the activated complex $C^{\ddagger}$ that tips it through the transition state and has frequency $v^{\ddagger}$ the partition function may be written from eqn 13B. 15 of Topic 13B as

$$
q=\frac{1}{1-\mathrm{e}^{-h \nu^{\ddagger} / k T}}
$$

\begin{center}
\includegraphics[max width=\textwidth]{2025_02_27_d4b95d9718f0b9e2e6b9g-826}
\end{center}

Figure 18C. 3 In an elementary depiction of the activated complex close to the transition state, there is a broad, shallow dip in the potential energy surface along the reaction coordinate. The complex vibrates harmonically and almost classically in this well.

This frequency $v^{\ddagger}$ is much lower than for an ordinary molecular vibration because the oscillation corresponds to the complex falling apart (Fig. 18C.3), so the force constant is very low. Therefore, provided that $h v^{\ddagger} / k T \ll 1$, the exponential may be expanded and the partition function reduces to

$$
q=\frac{1}{1-\left(1-h v^{\ddagger} / k T+\cdots\right)} \approx \frac{k T}{h v^{\ddagger}}
$$

It follows that the partition function for the activated complex is


\begin{equation*}
\mathscr{G}_{\mathrm{C}^{\ddagger}}^{\ominus}=\frac{k T}{h v^{\ddagger}} \bar{G}_{\mathrm{C}^{\ddagger}}^{\ominus} \tag{18C.8}
\end{equation*}


where the bar in $\overline{\mathcal{G}}_{\mathrm{C}^{\ddagger}}^{\ominus}$ denotes that the partition function is for all the other modes of the complex. The constant $K^{\ddagger}$ is therefore


\begin{equation*}
K^{\ddagger}=\frac{k T}{h v^{\ddagger}} \bar{K}^{\ddagger} \quad \bar{K}^{\ddagger}=\frac{N_{\mathrm{A}} \overline{\mathscr{G}}_{\mathrm{C}^{\ddagger}}^{\ominus}}{\mathscr{q}_{\mathrm{A}}^{\ominus} \mathrm{G}_{\mathrm{B}}^{\ominus}} \mathrm{e}^{-\Delta E_{0} / R T} \tag{18C.9}
\end{equation*}


with $\bar{K}^{\ddagger}$ a kind of equilibrium constant, but with one vibrational mode of $\mathrm{C}^{\ddagger}$ discarded.

\subsection*{Brief illustration 18C. 2}
Consider the case of two structureless particles A and B colliding to give an activated complex that resembles a diatomic molecule. The activated complex is a diatomic cluster. It has one vibrational mode, but that mode corresponds to motion along the reaction coordinate and therefore does not appear in $\overline{\mathcal{G}}_{\mathrm{C}^{+}}^{\ominus}$. It follows that the standard molar partition function of the activated complex has only rotational and translational contributions.

\subsection*{(d) The rate constant}
All the parts of the calculation can now be combined into

$$
k_{\mathrm{r}}=\frac{R T}{p^{\ominus}} k^{\ddagger} K^{\ddagger}=\kappa v^{\ddagger} \frac{k T}{h v^{\ddagger}} \frac{R T}{p^{\ominus}} \bar{K}^{\ddagger}
$$

At this stage the unknown frequency $v^{\ddagger}$ (in blue) cancels to give one version of the Eyring equation:

$$
k_{\mathrm{r}}=\kappa \frac{k T}{h} \frac{R T}{p^{\ominus}} \bar{K}^{\ddagger}
$$

Eyring equation (18C.10)

The equilibrium constant $\bar{K}^{\ddagger}$, which here is expressed in terms of partial pressures, can be rewritten in terms of concentrations by using $[J]=p_{J} / R T$, and then with some rearrangement the equilibrium constant in terms of concentrations, $\bar{K}_{c}^{\ddagger}$ (the terms in blue), can be identified:\\
$\bar{K}^{\ddagger}=\frac{p_{\mathrm{C}^{\ddagger}} p^{\ominus}}{p_{\mathrm{A}} p_{\mathrm{B}}}=\frac{\left[\mathrm{C}^{\ddagger}\right]}{[\mathrm{A}][\mathrm{B}]} \frac{p^{\ominus}}{R T}=\frac{\overbrace{\left[\mathrm{C}^{\ddagger}\right] c^{\ominus}}^{\bar{K}_{\mathrm{c}}^{\ddagger}}}{[\mathrm{A}][\mathrm{B}]} \times \frac{1}{c^{\ominus}} \times \frac{p^{\ominus}}{R T}=\bar{K}_{\mathrm{c}}^{\ddagger} \frac{p^{\ominus}}{R T c^{\ominus}}(18 \mathrm{C} .11)$

Substitution of this relation into eqn 18C.10, gives and an alternative version of the Eyring equation:

$$
\begin{array}{ll}
k_{\mathrm{r}}=\kappa \frac{k T}{h c^{\ominus}} \bar{K}_{\mathrm{c}}^{\ddagger} & \begin{array}{l}
\text { Eyring equation } \\
\text { [alternative version] }
\end{array}
\end{array}
$$

(18C.12)

The equilibrium constant $\bar{K}^{\ddagger}$ can be computed from the partition functions of $A, B$, and $C^{\ddagger}$, so in principle the Eyring equation is an explicit expression for calculating the second-order rate constant for a bimolecular reaction in terms of the molecular parameters for the reactants and the activated complex, and the quantity $\kappa$.

The partition functions for the reactants can normally be calculated quite readily by using either spectroscopic information about their energy levels or the approximate expressions set out in the Checklist at the end of Topic 13B. The difficulty with the Eyring equation, however, lies in the calculation of the partition function of the activated complex: $\mathrm{C}^{\ddagger}$ is difficult to investigate spectroscopically, and in general it is necessary to make assumptions about its size, shape, and structure.

Example 18C. 1\\
Analysing the collision of structureless particles

Consider the case of two structureless (and different) particles $A$ and $B$ colliding to give an activated complex that resembles a diatomic molecule. Deduce an expression for the rate constant of the reaction $\mathrm{A}+\mathrm{B} \rightarrow \mathrm{P}$.

Collect your thoughts Because the reactants are structureless 'atoms', the only contribution to their partition functions is from translation. The activated complex is a diatomic cluster of mass $m_{\mathrm{C}^{\ddagger}}=m_{\mathrm{A}}+m_{\mathrm{B}}$ and moment of inertia $I$. It has one vibrational mode but, as explained in Brief illustration 18C.2, that mode corresponds to motion along the reaction coordinate. It follows that the standard molar partition function of the activated complex has only rotational and translational contributions. Expressions for the relevant partition functions are given at the end of Topic 13B.

The solution The translational partition functions are

$$
\mathfrak{q}_{\mathrm{J}}^{\ominus}=\frac{V_{\mathrm{m}}^{\ominus}}{\Lambda_{\mathrm{J}}^{3}} \quad \Lambda_{\mathrm{J}}=\frac{h}{\left(2 \pi m_{\mathrm{J}} k T\right)^{1 / 2}} \quad V_{\mathrm{m}}^{\ominus}=\frac{R T}{p^{\ominus}}
$$

with $\mathrm{J}=\mathrm{A}, \mathrm{B}$, and $\mathrm{C}^{\ddagger}$, and with $m_{\mathrm{C}^{\ddagger}}=m_{\mathrm{A}}+m_{\mathrm{B}}$. The expression for the partition function of the activated complex is

$$
\bar{G}_{\mathrm{C}^{\ddagger}}^{\ominus}=\frac{\overbrace{2 I k T}^{\text {rotation }}}{\hbar^{2}} \times \frac{\overbrace{V_{\mathrm{m}}^{\ominus}}^{\text {translation }}}{\Lambda_{\mathrm{C}^{\ddagger}}^{3}}
$$

where the high-temperature form of the rotational partition function has been used (Topic 13B). From eqn 18C.9, the constant $\bar{K}^{\ddagger}$ is expressed in terms of the partition functions as\\
$\bar{K}^{\ddagger}=\underbrace{N_{\mathrm{A}} \overbrace{\left(2 I k T / \hbar^{2}\right) V_{\mathrm{m}}^{\ominus} / \Lambda_{\mathrm{C}^{\ddagger}}^{3}}^{\left(V_{\mathrm{m}}^{\ominus} / \Lambda_{\mathrm{A}}^{3}\right)} \underbrace{\left(V_{\mathrm{m}}^{\ominus} / \Lambda_{\mathrm{B}}^{3}\right)}_{\bar{q}_{\mathrm{B}}^{\ominus}}}_{\bar{q}_{\mathrm{A}}^{\ominus}} \mathrm{C}^{\ominus} \mathrm{e}^{-\Delta E_{0} / R T}=\left(\frac{N_{\mathrm{A}} \Lambda_{\mathrm{A}}^{3} \Lambda_{\mathrm{B}}^{3}}{\Lambda_{\mathrm{C}^{\ddagger}}^{3} V_{\mathrm{m}}^{\ominus}}\right) \frac{2 I k T}{\hbar^{2}} \mathrm{e}^{-\Delta E_{0} / R T}$

It follows from eqn 18C. 10 that

$$
\begin{aligned}
k_{\mathrm{r}} & =\kappa \frac{k T}{h} \frac{R T}{p^{\ominus}} \overbrace{\left(\frac{N_{\mathrm{A}} \Lambda_{\mathrm{A}}^{3} \Lambda_{\mathrm{B}}^{3}}{\Lambda_{\mathrm{C}^{\ddagger}}^{3} V_{\mathrm{m}}^{\ominus}}\right) \frac{2 I k T}{\hbar^{2}} \mathrm{e}^{-\Delta E_{0} / R T}}^{\bar{K}^{\ddagger}} \\
& =\kappa \frac{k T}{h} N_{\mathrm{A}}\left(\frac{\Lambda_{\mathrm{A}} \Lambda_{\mathrm{B}}}{\Lambda_{\mathrm{C}^{\ddagger}}^{3}}\right)^{3} \frac{2 I k T}{\hbar^{2}} \mathrm{e}^{-\Delta E_{0} / R T}
\end{aligned}
$$

The moment of inertia of a diatomic molecule of bond length $r$ is $\mu r^{2}$, where $\mu=m_{\mathrm{A}} m_{\mathrm{B}} /\left(m_{\mathrm{A}}+m_{\mathrm{B}}\right)$, so after introducing the expressions for the thermal wavelengths $\Lambda$ and cancelling common terms, the result is

$$
k_{\mathrm{r}}=\kappa N_{\mathrm{A}}\left(\frac{8 k T}{\pi \mu}\right)^{1 / 2} \pi r^{2} \mathrm{e}^{-\Delta E_{0} / R T}
$$

Finally, by identifying $\kappa \pi r^{2}$ as the reactive cross-section $\sigma^{\star}$, the resulting expression is the same as that obtained from simple collision theory (eqn 18A.9):

$$
k_{\mathrm{r}}=N_{\mathrm{A}}\left(\frac{8 k T}{\pi \mu}\right)^{1 / 2} \sigma^{\star} \mathrm{e}^{-\Delta E_{0} / R T}
$$

Self-test 18C. 1 What additional contributions would there be to the partition functions of the reactants and of the activated complex if the reaction were $\mathrm{AB}+\mathrm{C} \rightarrow \mathrm{P}$, with a linear activated complex?\\
\includegraphics[max width=\textwidth, center]{2025_02_27_d4b95d9718f0b9e2e6b9g-827}\\
\includegraphics[max width=\textwidth, center]{2025_02_27_d4b95d9718f0b9e2e6b9g-827(1)}

\subsection*{18C.2 Thermodynamic aspects}
The statistical thermodynamic version of transition-state theory rapidly runs into difficulties because only in some cases is anything known about the structure of the activated complex. However, the concepts it introduces, principally that of an equilibrium between the reactants and the activated complex, have motivated a more general, empirical approach in which the activation process is expressed in terms of thermodynamic functions.

\subsection*{(a) Activation parameters}
If $\bar{K}^{\ddagger}$ is taken as an equilibrium constant (despite one mode of $\mathrm{C}^{\ddagger}$ having been discarded), then it can be expressed in terms of a Gibbs energy of activation, $\Delta^{\ddagger} G$, through the definition

\[
\Delta^{\ddagger} G=-R T \ln \bar{K}^{\ddagger} \quad \begin{align*}
& \text { Gibbs energy of activation }  \tag{18C.13}\\
& \text { [definition] }
\end{align*}
\]

All the $\Delta^{\ddagger} X$ in this section are standard thermodynamic quantities, $\Delta^{\ddagger} X^{\ominus}$, but the standard state sign will be omitted to avoid overburdening the notation. Then from eqn 18C. 10 the expression for the rate constant becomes


\begin{equation*}
k_{\mathrm{r}}=\kappa \frac{k T}{h} \frac{R T}{p^{\ominus}} \mathrm{e}^{-\Delta^{\ddagger} G / R T} \tag{18C.14}
\end{equation*}


Because $\Delta G=\Delta H-T \Delta S$, the Gibbs energy of activation can be divided into an entropy of activation, $\Delta^{\ddagger} S$, and an enthalpy of activation, $\Delta^{\ddagger} H$, by writing

\[
\Delta^{\ddagger} G=\Delta^{\ddagger} H-T \Delta^{\ddagger} S \quad \begin{align*}
& \text { Entropy and enthalpy of activation }  \tag{18C.15}\\
& \text { [definition] }
\end{align*}
\]

When eqn 18C. 15 is used in eqn 18C. 14 and $\kappa$ is absorbed into the entropy term, the result is


\begin{equation*}
k_{\mathrm{r}}=B \mathrm{e}^{\Delta^{\ddagger} S / R} \mathrm{e}^{-\Delta^{\ddagger} H / R T} \quad B=\frac{k T}{h} \frac{R T}{p^{\ominus}} \tag{18C.16}
\end{equation*}


To develop this expression further it is necessary to find a relation between the enthalpy of activation and the activation energy. The two are not the same, for two main reasons. One is that although it might be tempting to identify $E_{\mathrm{a}}$ with $\Delta^{\ddagger} U$,\\
that is valid only at $T=0$; at higher temperatures upper levels of all the species are occupied and contribute additional terms of the order of $R T$ (a value suggested by the equipartition principle). Secondly, for gas-phase processes (but not for those in solution), $\Delta^{\ddagger} H$ differs from $\Delta^{\ddagger} U$ by another contribution $R T$. These additional contributions need to be identified.

\subsection*{How is that done? 18C. 1 Relating the enthalpy of activation to the activation energy}
The relation between the enthalpy of activation and the activation energy depends on two equations. One is the expression for the temperature dependence of the 'equilibrium constant' $K^{\ddagger}$, which is eqn 6 B .2 of Topic 6 B in the form $\mathrm{d} \ln K^{\ddagger} / \mathrm{d} T=\Delta^{\ddagger} H / R T^{2}$, and the second is the definition of the activation energy, which is eqn 17D. 3 of Topic 17D in the form $\mathrm{d} \ln k_{\mathrm{r}} / \mathrm{d} T=E_{\mathrm{a}} / R T^{2}$. The link between the two expressions is the alternative version of the Eyring equation, eqn 18C.12, $k_{\mathrm{r}}=\left(\kappa k T / h c^{\ominus}\right) \bar{K}_{\mathrm{c}}^{\ddagger}$. Differentiation of the corresponding expression for $\ln k_{\mathrm{r}}$ with respect to $T$ gives

$$
\frac{\mathrm{d} \ln k_{\mathrm{r}}}{\mathrm{~d} T}=\frac{1}{T}+\frac{\mathrm{d} \ln \bar{K}_{\mathrm{c}}^{\ddagger}}{\mathrm{d} T}
$$

Then, by using $\mathrm{d} \ln \bar{k}_{\mathrm{r}} / \mathrm{d} T=E_{\mathrm{a}} / R T^{2}$, it follows that

$$
E_{\mathrm{a}}=R T+R T^{2} \frac{\mathrm{~d} \ln \bar{K}_{\mathrm{c}}^{\ddagger}}{\mathrm{d} T}
$$

At this point it is necessary to distinguish between gas-phase and solution-phase reactions of the form $A+B \rightleftharpoons C^{\ddagger}$. For the latter, the equilibrium constant is expressed in terms of concentrations and the second term in the preceding equation can be identified with $\Delta^{\ddagger} H$ without further calculation. It follows that

For a solution-phase reaction: $E_{\mathrm{a}}=R T+\Delta^{\ddagger} H$\\
One further step is needed for a gas-phase reaction because $\bar{K}_{c}^{\ddagger}$ is not the same as $\bar{K}^{\ddagger}$ (which is expressed in terms of partial pressures). According to eqn 18C. 11 the two are related by $\bar{K}^{\ddagger}=\left(p^{\ominus} / R T c^{\ominus}\right) K_{c}^{\ddagger}$, which implies that

$$
\frac{\mathrm{d} \ln \bar{K}_{\mathrm{c}}^{\ddagger}}{\mathrm{d} T}=\overbrace{\frac{\mathrm{d}}{\mathrm{~d} T} \ln \frac{R T c^{\ominus}}{p^{\ominus}}}^{1 / T}+\overbrace{\frac{\mathrm{d} \ln \bar{K}^{\ddagger}}{\mathrm{d} T}}^{\Delta^{\ddagger} H / R T^{2}}=\frac{1}{T}+\frac{\Delta^{\ddagger} H}{R T^{2}}
$$

Substitution of this expression into the previous expression for $E_{\mathrm{a}}$ in terms of $\mathrm{d} \ln \bar{K}_{\mathrm{c}}^{\ddagger} / \mathrm{d} T$ gives

$$
E_{\mathrm{a}}=R T+R T^{2} \frac{\mathrm{~d} \ln \bar{K}_{\mathrm{c}}^{\ddagger}}{\mathrm{d} T}=R T+R T^{2}\left(\frac{1}{T}+\frac{\Delta^{\ddagger} H}{R T^{2}}\right)=R T+R T+\Delta^{\ddagger} H
$$

It therefore follows that\\
For a gas-phase reaction: $E_{\mathrm{a}}=2 R T+\Delta^{\ddagger} H$

In summary:\\
(a) $\Delta^{\ddagger} H=E_{\mathrm{a}}-2 R T$\\
(bimolecular gas-phase reaction)\\
(18C.17)\\
(b) $\Delta^{\ddagger} H=E_{\mathrm{a}}-R T$

Relations between\\
(bimolecular reaction in solution)\\
$\Delta^{\ddagger} H$ and $E$\\
It now follows that

$$
k_{\mathrm{r}}=\mathrm{e}^{2} B \mathrm{e}^{\Delta^{\ddagger} S / R} \mathrm{e}^{-E_{\mathrm{a}} / R T} \quad \begin{aligned}
& \text { Rate constant } \\
& \text { [transition-state theory, bimolecular } \\
& \text { gas-phase reaction }]
\end{aligned}
$$

(18C.18a)\\
and

$$
k_{\mathrm{r}}=\mathrm{e} B \mathrm{e}^{\Delta^{\ddagger} S / R} \mathrm{e}^{-E_{\mathrm{a}} / R T} \quad \begin{aligned}
& \text { Rate constant } \\
& {[\text { [tansition-state theory, }} \\
& \text { bimolecular reaction in solution] }
\end{aligned}
$$

(18C.18b)\\
where, from eqn 18C.16, $B=(k T / h)\left(R T / p^{\ominus}\right)$. The Arrhenius frequency factors can be identified as

$$
A=\mathrm{e}^{2} B \mathrm{e}^{\Delta^{\ddagger} S / R} \quad \begin{aligned}
& \text { Frequency factor } \\
& \text { [transition-state theory, bimolecular } \\
& \text { gas-phase reaction] }
\end{aligned}
$$

(18C.19a)\\
and

$$
A=\mathrm{eB} \mathrm{e}^{\Delta^{\ddagger} S / R} \quad \begin{aligned}
& \text { Frequency factor } \\
& \text { [transition-state theory, bimolecular } \\
& \text { reaction in solution] }
\end{aligned}
$$

(18C.19b)\\
The entropy of activation is negative because throughout the system reactant species are combining to form reactive pairs. However, if there is a reduction in entropy below what would be expected for the simple encounter of $A$ and $B$, then the frequency factor $A$ will be reduced further. Indeed, that additional reduction in entropy, $\Delta^{\ddagger} S_{\text {steric }}$, can be identified as the origin of the steric factor $P$ of collision theory (Topic 18A), so that

$$
\begin{array}{ll}
P=\mathrm{e}^{\Delta^{\ddagger} S_{\text {steric }} / R} & \begin{array}{l}
\text { P-factor } \\
\\
\text { [transition-state theory] }
\end{array}
\end{array}
$$

(18C.20)\\
Thus, the more complex the steric requirements of the encounter, the more negative the value of $\Delta^{\ddagger} S_{\text {steric }}$, and the smaller the value of $P$.

\subsection*{Brief illustration 18C. 3}
The reaction of propylxanthate ion in ethanoic acid buffer solutions can be represented by the equation $\mathrm{A}^{-}+\mathrm{H}^{+} \rightarrow \mathrm{P}$. Near $30^{\circ} \mathrm{C}, A=2.05 \times 10^{13} \mathrm{dm}^{3} \mathrm{~mol}^{-1} \mathrm{~s}^{-1}$. To evaluate the entropy of activation at $30^{\circ} \mathrm{C}$, use eqn 18 C .19 b , rearranged as

$$
\Delta^{\ddagger} S=R \ln \frac{A}{\mathrm{e} B} \text { with } B=\frac{k T}{h} \frac{R T}{p^{\ominus}}=1.592 \times 10^{14} \mathrm{dm}^{3} \mathrm{~mol}^{-1} \mathrm{~s}^{-1}
$$

Therefore,

$$
\begin{aligned}
\Delta^{\ddagger} S & =R \ln \frac{2.05 \times 10^{13} \mathrm{dm}^{3} \mathrm{~mol}^{-1} \mathrm{~s}^{-1}}{\mathrm{e} \times 1.592 \times 10^{14} \mathrm{dm}^{3} \mathrm{~mol}^{-1} \mathrm{~s}^{-1}}=R \ln 0.0473 \ldots \\
& =-25.4 \mathrm{JK}^{-1} \mathrm{~mol}^{-1}
\end{aligned}
$$

\begin{center}
\includegraphics[max width=\textwidth]{2025_02_27_d4b95d9718f0b9e2e6b9g-829}
\end{center}

Figure 18C. 4 For a related series of reactions here denoted $a$ and $b$, as the standard reaction Gibbs energy becomes more negative on going from $a$ to $b$, the activation Gibbs energy decreases and the rate constant increases. The approximate linear correlation between $\Delta^{\ddagger} G$ and $\Delta_{\mathrm{r}} G^{\ominus}$ is the origin of 'linear free energy relations'.

Gibbs energies, enthalpies, and entropies of activation (and volumes and heat capacities of activation) are widely used to report experimental reaction rates, especially for organic reactions in solution. They are encountered when relationships between equilibrium constants and rates of reaction are explored by using correlation analysis, in which $\ln K$ (which is equal to $-\Delta_{\mathrm{r}} G^{\ominus} / R T$ ) is plotted against $\ln k_{\mathrm{r}}$ (which is proportional to $-\Delta^{\ddagger} G / R T$ ). In many cases the correlation is linear, signifying that as the reaction becomes thermodynamically more favourable, its rate constant increases (Fig. 18C.4). This linear correlation is the origin of the alternative name linear free energy relation (LFER).

\subsection*{(b) Reactions between ions}
The full statistical thermodynamic theory is very complicated for reactions involving ions in solution because the solvent plays a role in the activated complex. The thermodynamic version of transition-state theory simplifies the discussion and is applicable to non-ideal systems. In the thermodynamic approach, the rate law

$$
\frac{\mathrm{d}[\mathrm{P}]}{\mathrm{d} t}=k^{\ddagger}\left[\mathrm{C}^{\ddagger}\right]
$$

is combined with the thermodynamic equilibrium constant (Topic 6A)

$$
K=\frac{a_{\mathrm{C}^{\ddagger}}}{a_{\mathrm{A}} a_{\mathrm{B}}}=K_{\gamma} \frac{a_{\mathrm{j}}=\gamma_{\mathrm{J}}[\mathrm{~J}] / c^{\ominus}}{\left[\mathrm{C}^{\ddagger}\right] c^{\ominus}}[\mathrm{A}][\mathrm{B}] \quad K_{\gamma}=\frac{\gamma_{\mathrm{C}^{\ddagger}}}{\gamma_{\mathrm{A}} \gamma_{\mathrm{B}}}
$$

Then


\begin{equation*}
\frac{\mathrm{d}[\mathrm{P}]}{\mathrm{d} t}=k_{\mathrm{r}}[\mathrm{~A}][\mathrm{B}] \quad k_{\mathrm{r}}=\frac{k^{\ddagger} K}{K_{\gamma} c^{\ominus}} \tag{18C.21a}
\end{equation*}


If $k_{\mathrm{r}}^{\circ}$ is the rate constant when the activity coefficients are 1 (i.e. $\left.k_{\mathrm{r}}^{\circ}=k^{\ddagger} K / c^{\ominus}\right)$, then


\begin{equation*}
k_{\mathrm{r}}=\frac{k_{\mathrm{r}}^{\circ}}{K_{\gamma}} \quad \log k_{\mathrm{r}}=\log k_{\mathrm{r}}^{\circ}-\log K_{\gamma} \tag{18C.21b}
\end{equation*}


At low concentrations the activity coefficients can be expressed in terms of the ionic strength, $I$, of the solution by using the Debye-Hückel limiting law (Topic 5F, particularly eqn 5F.27, $\log \gamma_{ \pm}=-\mathcal{A}\left|z_{+} z_{-}\right| I^{1 / 2}$ ). However, the expressions needed are those for the individual ions rather than the mean value, and so it is more appropriate to write $\log \gamma_{J}=$ $-\mathcal{A} z_{J}^{2} I^{1 / 2}$ and


\begin{equation*}
\log \gamma_{\mathrm{A}}=-\mathcal{A} z_{\mathrm{A}}^{2} I^{1 / 2} \quad \log \gamma_{\mathrm{B}}=-\mathcal{A} z_{\mathrm{B}}^{2} I^{1 / 2} \tag{18C.22a}
\end{equation*}


with $\mathcal{A}=0.509$ in aqueous solution at 298 K and $z_{\mathrm{A}}$ and $z_{\mathrm{B}}$ the (signed) charge numbers of A and B, respectively. Because the activated complex forms from reaction of one of the ions of $A$ with one of the ions of $B$, the charge number of the activated complex is $z_{\mathrm{A}}+z_{\mathrm{B}}$ where $z_{\mathrm{J}}$ is positive for cations and negative for anions. Therefore


\begin{equation*}
\log \gamma_{\mathrm{C}^{\ddagger}}=-\mathcal{A}\left(z_{\mathrm{A}}+z_{\mathrm{B}}\right)^{2} I^{1 / 2} \tag{18C.22b}
\end{equation*}


When these expressions are inserted into eqn 18C.21b the result is


\begin{align*}
\log k_{\mathrm{r}} & =\log k_{\mathrm{r}}^{0}-\mathcal{A}\left\{z_{\mathrm{A}}^{2}+z_{\mathrm{B}}^{2}-\left(z_{\mathrm{A}}+z_{\mathrm{B}}\right)^{2}\right\} I^{1 / 2} \\
& =\log k_{\mathrm{r}}^{0}+2 \mathcal{A} z_{\mathrm{A}} z_{\mathrm{B}} I^{1 / 2} \tag{18C.23}
\end{align*}


Equation 18C. 23 expresses the kinetic salt effect, the variation of the rate constant of a reaction between ions with the ionic strength of the solution (Fig. 18C.5). The equation is interpreted as follows:

\begin{itemize}
  \item If the reactant ions have the same sign (as in a reaction between cations or between anions), then increasing the ionic strength by the addition of inert ions increases the rate constant.
  \item The formation of a single, highly charged ionic complex from two less highly charged ions is favoured by a high ionic strength because the new ion has a denser ionic atmosphere and interacts with that atmosphere more strongly.
  \item Conversely, ions of opposite charge react more slowly in solutions of high ionic strength. Now the charges cancel and the complex has a less favourable interaction with its atmosphere than the separated ions.\\
\includegraphics[max width=\textwidth, center]{2025_02_27_d4b95d9718f0b9e2e6b9g-830(1)}
\end{itemize}

Figure 18C. 5 Experimental tests of the kinetic salt effect for reactions in water at 298 K . The ion types are shown as spheres, and the slopes of the lines are those given by the Debye-Hückel limiting law and eqn 18C.23.

\subsection*{Example 18C. 2 \\
 Analysing data in terms of the kinetic salt effect}
The rate constant (at 298 K ) for the hydrolysis of $\left[\mathrm{CoBr}\left(\mathrm{NH}_{3}\right)_{5}\right]^{2+}$ under basic conditions in aqueous solution varies with ionic strength as in the following table. What can be deduced about the charge of the activated complex in the rate-determining step? What might this deduction imply about the mechanism?

\begin{center}
\begin{tabular}{lllllll}
$I$ & 0.0050 & 0.0100 & 0.0150 & 0.0200 & 0.0250 & 0.0300 \\
$k_{\mathrm{r}} / k_{\mathrm{r}}^{\circ}$ & 0.718 & 0.631 & 0.562 & 0.515 & 0.475 & 0.447 \\
\end{tabular}
\end{center}

Collect your thoughts According to eqn 18C.23, a plot of $\log \left(k_{\mathrm{r}} k_{\mathrm{r}}^{\circ}\right)$ against $I^{1 / 2}$ should have a slope of $2 \mathcal{A} z_{\mathrm{A}} z_{\mathrm{B}}$. Because $\mathcal{A}=0.509$ for aqueous solutions at 298 K , the slope will be $1.02 z_{\mathrm{A}} z_{\mathrm{B}}$. From the slope you can infer the charges of the ions involved in the formation of the activated complex.

The solution Form the following table:

\begin{center}
\begin{tabular}{lllllll}
$I$ & 0.0050 & 0.0100 & 0.0150 & 0.0200 & 0.0250 & 0.0300 \\
$I^{1 / 2}$ & 0.071 & 0.100 & 0.122 & 0.141 & 0.158 & 0.173 \\
$\log \left(k_{\mathrm{r}} / k_{\mathrm{r}}^{\circ}\right)$ & -0.14 & -0.20 & -0.25 & -0.29 & -0.32 & -0.35 \\
\end{tabular}
\end{center}

These values are plotted in Fig. 18C.6. The slope of the (least squares) straight line is -2.04 , indicating that $z_{\mathrm{A}} z_{\mathrm{B}}=-2$. A possible explanation for this conclusion is that the two species involved in the formation of the activated complex are $\left[\mathrm{CoBr}\left(\mathrm{NH}_{3}\right)_{5}\right]^{2+}$, which has $z=+2$, and $\mathrm{OH}^{-}$, which has $z=-1$. The product of the charges is therefore -2 .

Comment. Although the point is not pursued here, you should be aware that the rate constant is also influenced by the relative permittivity of the medium.

Self-test 18C. 2 An ion of charge number +1 is known to be involved in the activated complex of a reaction. Deduce the charge number of the other ion from the following data, recorded at 298 K in aqueous solution:\\
\includegraphics[max width=\textwidth, center]{2025_02_27_d4b95d9718f0b9e2e6b9g-830}

Figure 18C. 6 The experimental ionic strength dependence of the rate constant of a hydrolysis reaction: the slope gives information about the charge types involved in the activated complex of the rate-determining step. The data plotted are from Example 18C.2.

\begin{center}
\begin{tabular}{lllllll}
$I$ & 0.0050 & 0.0100 & 0.0150 & 0.0200 & 0.0250 & 0.0300 \\
$k_{\mathrm{r}} / k_{\mathrm{r}}^{\circ}$ & 0.847 & 0.791 & 0.750 & 0.717 & 0.690 & 0.666 \\
\end{tabular}
\end{center}

\subsection*{18C. 3 The kinetic isotope effect}
The postulation of a plausible reaction mechanism requires careful analysis of many experiments designed to determine the fate of atoms during the formation of products. Observation of the kinetic isotope effect, a decrease in the rate of a chemical reaction upon replacement of one atom in a reactant by a heavier isotope, facilitates the identification of bond-breaking events in the rate-determining step. A primary kinetic isotope effect is observed when the rate-determining step requires the scission of a bond involving the isotope. A secondary kinetic isotope effect is the reduction in reaction rate even though the bond involving the isotope is not broken to form product. In both cases, the effect arises from the change in activation energy that accompanies the replacement of an atom by a heavier isotope on account of changes in the zero-point vibrational energies. What follows is a description of the primary kinetic isotope effect.

Consider a reaction in which a $\mathrm{C}-\mathrm{H}$ bond is cleaved. If scission of this bond is the rate-determining step (Topic 17E), then the reaction coordinate corresponds to the stretching of the $\mathrm{C}-\mathrm{H}$ bond and the potential energy profile is shown in Fig. 18C.7. On deuteration, the dominant change is the reduction of the zero-point energy of the bond (because the deuterium atom is heavier). The whole reaction profile is not lowered, however, because the relevant vibration in the activated complex has a very low force constant, so there is little zero-point energy associated with the reaction coordinate in\\
\includegraphics[max width=\textwidth, center]{2025_02_27_d4b95d9718f0b9e2e6b9g-831}

Figure 18C.7 Changes in the reaction profile when a C-H bond undergoing cleavage is deuterated. In this illustration the C-H and C-D bonds are modelled as harmonic oscillators. The only significant change is in the zero-point energy of the reactants, which is lower for C-D than for C-H. As a result, the activation energy is greater for C-D cleavage than for C-H cleavage.\\
either form of the activated complex. From these considerations it is possible to investigate the effect of deuteration on the activation energy.

How is that done? 18C. 2\\
Exploring the primary kinetic isotope effect

Consider the cleavage of a C-H bond in a larger molecule. For such a reaction it is reasonable to assume that motion along the reaction coordinate is dominated by stretching and compression of the $\mathrm{C}-\mathrm{H}$ fragment. Therefore, to a good approximation, a change in the activation energy arises only from the change in zero-point energy of the stretching vibration, $\frac{1}{2} \hbar \omega$. It follows from Fig. 18C. 7 that

$$
\begin{aligned}
E_{\mathrm{a}}(\mathrm{C}-\mathrm{D})-E_{\mathrm{a}}(\mathrm{C}-\mathrm{H}) & =N_{\mathrm{A}}\left\{\frac{1}{2} \hbar \omega(\mathrm{C}-\mathrm{H})-\frac{1}{2} \hbar \omega(\mathrm{C}-\mathrm{D})\right\} \\
& =\frac{1}{2} N_{\mathrm{A}} \hbar\{\omega(\mathrm{C}-\mathrm{H})-\omega(\mathrm{C}-\mathrm{D})\}
\end{aligned}
$$

Then from Topic 11C $\omega(\mathrm{C}-\mathrm{D})=\left(\mu_{\mathrm{CH}} / \mu_{\mathrm{CD}}\right)^{1 / 2} \omega(\mathrm{C}-\mathrm{H})$, where $\mu$ is the relevant effective mass and therefore that

$$
E_{\mathrm{a}}(\mathrm{C}-\mathrm{D})-E_{\mathrm{a}}(\mathrm{C}-\mathrm{H})=\frac{1}{2} N_{\mathrm{A}} \hbar \omega(\mathrm{C}-\mathrm{H})\left\{1-\left(\frac{\mu_{\mathrm{CH}}}{\mu_{\mathrm{CD}}}\right)^{1 / 2}\right\}
$$

Effect of deuteration on the activation energy\\
(18C.24)\\
If the Arrhenius frequency factor does not change upon deuteration, the rate constants for the two species should be in the ratio

$$
\frac{k_{\mathrm{r}}(\mathrm{C}-\mathrm{D})}{k_{\mathrm{r}}(\mathrm{C}-\mathrm{H})}=\mathrm{e}^{-\left\{E_{\mathrm{a}}(\mathrm{C}-\mathrm{D})-E_{\mathrm{a}}(\mathrm{C}-\mathrm{H})\right\} / R T}=\mathrm{e}^{-\left\{E_{\mathrm{a}}(\mathrm{C}-\mathrm{D})-E_{\mathrm{a}}(\mathrm{C}-\mathrm{H})\right\} / N_{\mathrm{A}} k T}
$$

where $R=N_{\mathrm{A}} k$. After using eqn 18 C .24 for $E_{\mathrm{a}}(\mathrm{C}-\mathrm{D})-E_{\mathrm{a}}(\mathrm{C}-\mathrm{H})$ in this expression, the result is\\
$-\left\lvert\, \frac{k_{\mathrm{r}}(\mathrm{C}-\mathrm{D})}{k_{\mathrm{r}}(\mathrm{C}-\mathrm{H})}=\mathrm{e}^{-\zeta}\right.$ with $\left.\zeta=\frac{\hbar \omega(\mathrm{C}-\mathrm{H})}{2 k T}\left\{1-\left(\frac{\mu_{\mathrm{CH}}}{\mu_{\mathrm{CD}}}\right)^{1 / 2}\right\} \right\rvert\, \square$\\
Effect of deuteration on the rate constant\\
(18C.25)

Note that $\zeta>0$ ( $\zeta$ is zeta) because $\mu_{\mathrm{CD}}>\mu_{\mathrm{CH}}$ and so it follows that $k_{\mathrm{r}}(\mathrm{C}-\mathrm{D}) / k_{\mathrm{r}}(\mathrm{C}-\mathrm{H})<1$. As expected from Fig. 18C.7, the rate constant decreases upon deuteration.

\subsection*{Brief illustration 18C. 4}
From infrared spectra, the fundamental vibrational wavenumber $\tilde{v}$ for stretching of a $\mathrm{C}-\mathrm{H}$ bond is about $3000 \mathrm{~cm}^{-1}$. To convert this wavenumber to an angular frequency, $\omega=2 \pi v$, use $\omega=2 \pi c \tilde{v}$, so that

$$
\begin{aligned}
\omega & =2 \pi \times\left(2.998 \times 10^{10} \mathrm{~cm} \mathrm{~s}^{-1}\right) \times\left(3000 \mathrm{~cm}^{-1}\right) \\
& =5.65 \ldots \times 10^{14} \mathrm{~s}^{-1}
\end{aligned}
$$

The ratio of effective masses is

$$
\begin{aligned}
\frac{\mu_{\mathrm{CH}}}{\mu_{\mathrm{CD}}} & =\left(\frac{m_{\mathrm{C}} m_{\mathrm{H}}}{m_{\mathrm{C}}+m_{\mathrm{H}}}\right) \times\left(\frac{m_{\mathrm{C}}+m_{\mathrm{D}}}{m_{\mathrm{C}} m_{\mathrm{D}}}\right) \\
& =\left(\frac{12.01 \times 1.0078}{12.01+1.0078}\right) \times\left(\frac{12.01+2.0140}{12.01 \times 2.0140}\right) \\
& =0.539 \ldots
\end{aligned}
$$

Now use eqn 18C. 25 to calculate

$$
\begin{aligned}
\zeta & \left.=\frac{\left(1.055 \times 10^{-34} \mathrm{~J} \mathrm{~s}\right) \times\left(5.65 \ldots \times 10^{14} \mathrm{~s}^{-1}\right)}{2 \times\left(1.381 \times 10^{-23} \mathrm{~J} \mathrm{~K}\right.}{ }^{-1}\right) \times(298 \mathrm{~K})
\end{aligned}\left(1-0.539 \ldots{ }^{1 / 2}\right)
$$

and

$$
\frac{k_{\mathrm{r}}(\mathrm{C}-\mathrm{D})}{k_{\mathrm{r}}(\mathrm{C}-\mathrm{H})}=\mathrm{e}^{-1.92 \ldots}=0.146
$$

Therefore at room temperature cleavage of the $\mathrm{C}-\mathrm{H}$ bond should be about seven times faster than cleavage of the C-D bond, other conditions being equal. Experimental values of $k_{\mathrm{r}}(\mathrm{C}-\mathrm{D}) / k_{\mathrm{r}}(\mathrm{C}-\mathrm{H})$ can differ significantly from those predicted by eqn 18C. 25 on account of the severity of the assumptions in the model.

In some cases, substitution of deuterium for hydrogen results in values of $k_{\mathrm{r}}(\mathrm{C}-\mathrm{D}) / k_{\mathrm{r}}(\mathrm{C}-\mathrm{H})$ that are too low to be accounted for by eqn 18C.25, even when more complete models are used to predict ratios of rate constants. Such abnormal kinetic isotope effects are evidence for a path in which quantum mechanical tunnelling of hydrogen atoms takes place through the activation barrier (Fig. 18C.8). The probability\\
\includegraphics[max width=\textwidth, center]{2025_02_27_d4b95d9718f0b9e2e6b9g-832}

Figure 18C.8 A proton can tunnel through the activation energy barrier that separates reactants from products, so the effective height of the barrier is reduced and the rate of the proton transfer reaction increases. The effect is represented by drawing the wavefunction of the proton near the barrier. Proton tunnelling is important only at low temperatures, when most of the reactants are trapped on the left of the barrier.\\
of tunnelling through a barrier decreases as the mass of the particle increases (Topic 7D), so deuterium tunnels less efficiently through a barrier than hydrogen and its reactions are correspondingly slower. Quantum mechanical tunnelling can be the dominant process in reactions involving hydrogen atom or proton transfer when the temperature is so low that very few reactant molecules can overcome the activation energy barrier.

\subsection*{Checklist of concepts}
$\square$ 1. In transition-state theory, it is supposed that an activated complex is in equilibrium with the reactants.2. The rate at which the activated complex forms products depends on the rate at which it passes through a transition state.3. The rate constant may be parametrized in terms of the Gibbs energy, entropy, and enthalpy of activation.\\
4. The kinetic salt effect is the effect of an added inert salt on the rate constant of a reaction between ions.\\
5. The kinetic isotope effect is the decrease in the rate constant of a chemical reaction upon replacement of one atom in a reactant by a heavier isotope.

\subsection*{Checklist of equations}
\begin{center}
\begin{tabular}{|c|c|c|c|}
\hline
Property & Equation & Comment & Equation number \\
\hline
'Equilibrium constant' for activated complex formation & $\bar{K}^{\ddagger}=\left(N_{\mathrm{A}} \overline{\mathcal{G}}_{\mathrm{C}^{\ddagger}}^{\ominus} / \mathcal{G}_{\mathrm{A}}^{\ominus} \mathcal{G}_{\mathrm{B}}^{\ominus}\right) \mathrm{e}^{-\Delta E_{0} / R T}$ & Assume equilibrium; one vibrational mode of $\mathrm{C}^{\ddagger}$ discarded & 18C. 9 \\
\hline
\multirow[t]{2}{*}{Eyring equation} & $k_{\mathrm{r}}=\kappa(k T / h)\left(R T / p^{\ominus}\right) \bar{K}^{\ddagger}$ & \multirow[t]{2}{*}{Transition-state theory} & 18C. 10 \\
\hline
 & $k_{\mathrm{r}}=\kappa\left(k T / h c^{\ominus}\right) \bar{K}_{\mathrm{c}}^{*}$ &  & 18C. 12 \\
\hline
Gibbs energy of activation & $\Delta^{\ddagger} G=-R T \ln \bar{K}^{\ddagger}$ & Definition & 18C. 13 \\
\hline
Enthalpy and entropy of activation & $\Delta^{\ddagger} G=\Delta^{\ddagger} H-T \Delta^{\ddagger} S$ & Definition & 18C. 15 \\
\hline
Parametrization & $k_{\mathrm{r}}=\mathrm{e}^{n} B \mathrm{e}^{\Delta^{\ddagger} S / R} \mathrm{e}^{-E_{\mathrm{a}} / R T}$ & \multirow[t]{3}{*}{$n=2$ for bimolecular gas-phase reactions; $n=1$ for solution} & 18C. 18 \\
\hline
$A$-factor & $A=\mathrm{e}^{n} B \mathrm{e}^{ \pm^{\ddagger} S / R}$ &  & 18C. 19 \\
\hline
$P$-factor & $P=\mathrm{e}^{\mathrm{L}^{\ddagger} S_{\text {steric }} / R}$ &  & 18C. 20 \\
\hline
Kinetic salt effect & $\log k_{\mathrm{r}}=\log k_{\mathrm{r}}^{\circ}+2 \mathcal{A} z_{\mathrm{A}} z_{\mathrm{B}} I^{1 / 2}$ & Assumes Debye-Hückel limiting law valid & 18C. 23 \\
\hline
Primary kinetic isotope effect & \( \begin{aligned} & k_{\mathrm{r}}(\mathrm{C}-\mathrm{D}) / k_{\mathrm{r}}(\mathrm{C}-\mathrm{H})=\mathrm{e}^{-\zeta} \\ & \zeta=(\hbar \omega(\mathrm{C}-\mathrm{H}) / 2 k T) \times \\ & \left\{1-\left(\mu_{\mathrm{CH}} / \mu_{\mathrm{CD}}\right)^{1 / 2}\right\} \end{aligned} \) & Cleavage of a $\mathrm{C}-\mathrm{H} / \mathrm{D}$ bond in the ratedetermining step & 18C. 25 \\
\hline
\end{tabular}
\end{center}

\section*{TOPIC 18D The dynamics of molecular collisions }
\subsection*{Why do you need to know this material?}
Chemists are interested in the details of chemical reactions, and there is no more detailed approach than that involved in the study of the dynamics of reactive encounters, when one molecule collides with another and atoms exchange partners.

\subsection*{What is the key idea?}
The rates of reactions in the gas phase can be investigated by exploring the trajectories of molecules on potential energy surfaces.

\subsection*{> What do you need to know already?}
This Topic builds on the concept of rate constant (Topic 17A) and in one part of the discussion uses the concept of partition function (Topic 13B). The discussion of potential energy surfaces is qualitative, but the underlying calculations are those of self-consistent field theory (Topic 9E).

The investigation of the dynamics of the collisions between reactant molecules is the most detailed level of the examination of the factors that govern the rates of reactions. There are two approaches: an experimental one that uses molecular beams and a theoretical one that uses the results of computations.

\subsection*{18D.1 Molecular beams}
Molecular beams, which consist of collimated, narrow streams of molecules travelling through an evacuated vessel, allow collisions between molecules in preselected states (e.g. specific rotational and vibrational states) to be studied, and can be used to identify the states of the products of a reactive collision. Information of this kind is essential if a full picture of the reaction is to be built, because the rate constant is an average over events in which reactants in different initial states evolve into products in their final states.

\subsection*{(a) Techniques}
The basic arrangement of a molecular beam experiment is shown in Fig. 18D.1. Atoms or molecules emerge from a source chamber (which may be heated if the species is not already a gas) through a pinhole and out into a vacuum chamber. If the pressure of vapour in the source is increased so that the mean free path of the molecules in the emerging beam is much shorter than the diameter of the pinhole, many collisions take place even outside the source. The net effect of these collisions, which give rise to hydrodynamic flow, is to transfer momentum into the direction of the beam. The molecules in the beam then travel with very similar speeds, so further downstream few collisions take place between them. This condition is called molecular flow. If necessary, atoms or molecules moving with a given speed can be selected by using a velocity selector, as depicted in Fig. 18D.1.

In a molecular beam of this kind the spread of speeds is much smaller than that predicted by the Maxwell-Boltzmann distribution. The unexpected narrowness of the distribution is interpreted by assigning a low translational temperature to\\
\includegraphics[max width=\textwidth, center]{2025_02_27_d4b95d9718f0b9e2e6b9g-833}

Figure 18D. 1 The basic arrangement of a molecular beam apparatus. The atoms or molecules emerge from a source, and pass through a velocity selector, such as that discussed in Topic 1 B . The scattering occurs from the target gas (which might take the form of another beam), and the flux of particles entering the detector set at some angle is recorded. In a crossed-beam experiment, state-selected molecules are generated in two separate sources, and are directed perpendicular to one another. The detector responds to molecules (which may be product molecules if a chemical reaction occurs) scattered into a chosen direction.\\
\includegraphics[max width=\textwidth, center]{2025_02_27_d4b95d9718f0b9e2e6b9g-834(2)}

Figure 18D. 2 The shift in the mean speed and the width of the distribution brought about by use of a supersonic nozzle.\\
the molecules in the beam (Fig. 18D.2), which may be as low as 1 K . Such jets are called supersonic because the mean speed of the molecules in the jet is much greater than the speed of sound in the jet.

A supersonic jet can be converted into a more parallel supersonic beam if it is 'skimmed' in the region of hydrodynamic flow and the excess gas pumped away. A skimmer consists of a conical nozzle shaped to avoid any supersonic shock waves spreading back into the gas and so increasing the translational temperature (Fig. 18D.3). A jet or beam may also be formed by using helium or neon as the principal gas, and injecting molecules of interest into it in the hydrodynamic region of flow.

As well as having a low translational temperature, the molecules in the beam also have low rotational and vibrational temperatures. In this context, a rotational or vibrational temperature means the temperature that should be used in the Boltzmann distribution to reproduce the observed populations of the states. However, as rotational states equilibrate more slowly than translational states, and vibrational states equilibrate even more slowly, the rotational and vibrational populations of the species correspond to somewhat higher temperatures, of the order of 10 K for rotation and 100 K for vibrations.

The target gas may be either a bulk sample or another molecular beam. The detectors may consist of a chamber fitted\\
\includegraphics[max width=\textwidth, center]{2025_02_27_d4b95d9718f0b9e2e6b9g-834(3)}

Figure 18D. 3 A supersonic beam is generated by using a skimmer to remove some of the molecules from the beam, so leading to a greater degree of collimation.\\
with a sensitive pressure gauge, a 'bolometer' (a detector that responds to the incident energy by making use of the temper-ature-dependence of resistance), or an ionization detector, in which the incoming molecule is first ionized and then detected electronically. The rotational and vibrational state of the scattered molecules may also be determined spectroscopically.

\subsection*{(b) Experimental results}
The primary experimental information from a molecular beam experiment is the fraction of the molecules in the incident beam that are scattered into a particular direction. The fraction is normally expressed in terms of $\mathrm{d} I$, the number of molecules in a given time divided by the length of the interval, scattered into a cone (described by a solid angle $\mathrm{d} \Omega$ ) that represents the area covered by the 'eye' of the detector (Fig. 18D.4). This rate is reported as the differential scattering cross-section, $\sigma$, the constant of proportionality between the value of $d I$ and the intensity, $I$, of the incident beam, the number density of target molecules, $\mathcal{N}$, and the infinitesimal path length $\mathrm{d} x$ through the sample:

\[
\begin{array}{ll}
\mathrm{d} I=\sigma I \mathcal{N d} x & \begin{array}{l}
\text { Differential scattering cross-section } \\
\text { [definition] }
\end{array} \tag{18D.1}
\end{array}
\]

The value of $\sigma$ (which has the dimensions of area) depends on the impact parameter, $b$, the initial perpendicular separation of the paths of the colliding molecules (Fig. 18D.5), and the details of the intermolecular potential.

The role of the impact parameter is most easily seen by considering the impact of two hard spheres (Fig. 18D.6). If $b=0$, the projectile is on a trajectory that leads to a head-on collision, so the only scattering intensity is detected when the detector is at $\theta=\pi$. When the impact parameter is so great that the spheres do not make contact ( $b>R_{\mathrm{A}}+R_{\mathrm{B}}$ ), there is no\\
\includegraphics[max width=\textwidth, center]{2025_02_27_d4b95d9718f0b9e2e6b9g-834}

Figure 18D. 4 The definition of the solid angle, $\mathrm{d} \Omega$, for scattering.\\
\includegraphics[max width=\textwidth, center]{2025_02_27_d4b95d9718f0b9e2e6b9g-834(1)}

Figure 18D. 5 The definition of the impact parameter, $b$, as the perpendicular separation of the initial paths of the particles.\\
\includegraphics[max width=\textwidth, center]{2025_02_27_d4b95d9718f0b9e2e6b9g-835(3)}

Figure 18D. 6 Three typical cases for the collisions of two hard spheres: (a) $b=0$, giving backward scattering; (b) $b>R_{A}+R_{B}$, giving forward scattering; (c) $0<b<R_{A}+R_{B}$, leading to scattering into one direction on a ring of possibilities. (The target molecule is taken to be so heavy that it remains virtually stationary.)\\
scattering and the scattering cross-section is zero at all angles except $\theta=0$. Glancing blows, with $0<b \leq R_{\mathrm{A}}+R_{\mathrm{B}}$, lead to scattering intensity in cones around the forward direction.

The scattering pattern of real molecules, which are not hard spheres, depends on the details of the intermolecular potential energy and molecular shape. The scattering also depends on the relative speed of approach of the two molecules: a very fast molecule might pass through the interaction region without much deflection, whereas a slower one on the same path might be temporarily captured and undergo considerable deflection (Fig. 18D.7). The variation of the scattering cross-section with the relative speed of approach gives information about the strength and range of the intermolecular potential.

A further point is that the outcome of collisions is determined by quantum, not classical, mechanics. The wave nature of the molecules can be taken into account, at least to some extent, by drawing all classical trajectories that take the pro-\\
\includegraphics[max width=\textwidth, center]{2025_02_27_d4b95d9718f0b9e2e6b9g-835(1)}

Figure 18D. 7 The extent of scattering may depend on the relative speed of approach as well as the impact parameter. The dark central zone represents the repulsive core; the fuzzy outer zone represents the long-range attractive potential.\\
\includegraphics[max width=\textwidth, center]{2025_02_27_d4b95d9718f0b9e2e6b9g-835}

Figure 18D. 8 Two paths leading to the same destination will interfere quantum mechanically; in this case they give rise to quantum oscillations in the forward direction.\\
jectile molecule from source to detector, and then considering the effects of interference between them.

Two quantum mechanical effects are of great importance. A molecule with a certain impact parameter might approach the attractive region of the potential in such a way that it is deflected towards the repulsive core (Fig. 18D.8), which then repels it out through the attractive region to continue its flight in the forward direction. Some molecules, however, also travel in the forward direction because they have impact parameters so large that they are undeflected. The wavefunctions of the molecules that take the two types of path interfere, and the intensity in the forward direction is modified. The effect is called quantum oscillation. The same phenomenon accounts for the optical 'glory effect', in which a bright halo can sometimes be seen surrounding an illuminated object. (The coloured rings around the shadow of an aircraft cast on clouds by the Sun, and often seen in flight, are an example of an optical glory.)

The second quantum effect is the observation of a strongly enhanced scattering in a non-forward direction. This effect is called rainbow scattering because the same mechanism accounts for the appearance of an optical rainbow. The origin of the phenomenon is illustrated in Fig. 18D.9. As the impact parameter decreases, there comes a stage at which the scattering angle passes through a maximum and the interference between the paths results in a strongly scattered beam. The\\
\includegraphics[max width=\textwidth, center]{2025_02_27_d4b95d9718f0b9e2e6b9g-835(2)}

Figure 18D. 9 The interference of paths leading to rainbow scattering. The rainbow angle, $\theta_{r}$ is the maximum scattering angle reached as $b$ is decreased. Interference between the numerous paths at that angle modifies the scattering intensity markedly.\\
rainbow angle, $\theta_{\mathrm{r}}$, is the angle for which $\mathrm{d} \theta / \mathrm{d} b=0$ and the scattering is strong.

Another phenomenon that can occur in certain beams is the capturing of one species by another. The vibrational temperature in supersonic beams is so low that van der Waals molecules may be formed, which are complexes of the form AB in which $A$ and $B$ are held together by van der Waals forces or hydrogen bonds. Large numbers of such molecules have been studied spectroscopically, including $\mathrm{ArHCl},(\mathrm{HCl})_{2}, \mathrm{ArCO}_{2}$, and $\left(\mathrm{H}_{2} \mathrm{O}\right)_{2}$. The study of their spectroscopic properties gives detailed information about the intermolecular interactions involved.

\subsection*{18D. 2 Reactive collisions}
Detailed experimental information about the intimate processes that occur during reactive encounters comes from molecular beams, especially crossed molecular beams (see Fig. 18D.1). The detector for the products of the collision of molecules in the two beams can be moved to different angles to observe the angular distribution of the products. Because the molecules in the incoming beams can be prepared with different energies (e.g. with different translational energies by using rotating sectors and supersonic nozzles, with different vibrational energies by using selective excitation with lasers, and with different orientations by using electric fields), it is possible to study the dependence of the success of collisions on these variables and to study how they affect the properties of the product molecules.

\subsection*{(a) Probes of reactive collisions}
One method for examining the energy distribution in the products is infrared chemiluminescence, in which vibrationally excited molecules emit infrared radiation as they return to their ground states. By studying the intensities of the infrared emission spectrum, the populations of the vibrational states of the products may be determined (Fig. 18D.10). Another method makes use of laser-induced fluorescence. In this technique, a laser is used to excite a product molecule from a specific vibration-rotation level; the intensity of the fluorescence from the upper state is monitored and interpreted in terms of the population of the initial vibration-rotation state. When the molecules being studied do not fluoresce efficiently, versions of Raman spectroscopy (Topic 11A) can be used to monitor the progress of reaction.

Multiphoton ionization (MPI) techniques are also good alternatives for the study of weakly fluorescing molecules. In MPI, the absorption by a molecule of several photons from one or more pulsed lasers results in ionization if the total photon energy is greater than the ionization energy of the molecule. An important variant of MPI is resonant multiphoton\\
\includegraphics[max width=\textwidth, center]{2025_02_27_d4b95d9718f0b9e2e6b9g-836}

Figure 18D. 10 Infrared chemiluminescence from CO produced in the reaction $\mathrm{O}+\mathrm{CS} \rightarrow \mathrm{CO}+\mathrm{S}$ arises from the non-equilibrium populations of the vibrational states of CO. The horizontal bars indicate the relative populations of the vibrational states and the blue arrows indicate the observed transitions.\\
ionization (REMPI), in which one or more photons promote a molecule to an electronically excited state and then additional photons are used to generate ions from the excited state. The power of REMPI lies in the fact that the experimenter can choose which reactant or product to study by tuning the laser frequency to the electronic absorption band of a specific molecule.

The angular distribution of ions can be determined by reaction product imaging. In this technique, product ions are accelerated by an electric field towards a phosphorescent screen and the light emitted from specific spots where the ions struck the screen is imaged by a charge-coupled device (CCD).

\subsection*{(b) State-to-state reaction dynamics}
The concept of collision cross-section is introduced in connection with collision theory in Topic 18A, where it is shown that the second-order rate constant, $k_{\mathrm{r}}$, can be expressed as a Boltzmann-weighted average of the reactive cross-section and the relative speed of approach of the colliding reactant molecules. Equation 18A. 7 of that Topic $\left(k_{\mathrm{r}}=N_{\mathrm{A}} \int_{0}^{\infty} \sigma(\varepsilon) v_{\text {rel }} f(\varepsilon) \mathrm{d} \varepsilon\right)$ may be written as


\begin{equation*}
k_{\mathrm{r}}=\left\langle\sigma v_{\mathrm{rel}}\right\rangle N_{\mathrm{A}} \tag{18D.2}
\end{equation*}


where the angle brackets denote a Boltzmann average. Molecular beam studies provide a more sophisticated version of this quantity, for they provide the state-to-state cross-section, $\sigma_{n n^{\prime}}$, and hence the state-to-state rate constant, $k_{n n^{\prime}}$, for the reactive transition from initial state $n$ of the reactants to final state $n^{\prime}$ of the products:


\begin{equation*}
k_{n n^{\prime}}=\left\langle\sigma_{n n^{\prime}} v_{\text {rel }}\right\rangle N_{\mathrm{A}} \quad \text { State-to-state rate constant } \tag{18D.3}
\end{equation*}


The rate constant $k_{\mathrm{r}}$ is the sum of the state-to-state rate constants over all final states (because a reaction is successful whatever the final state of the products) and over a Boltzmann-weighted sum of initial states (because the reactants are initially present\\
with a characteristic distribution of populations at a temperature $T$ ):


\begin{equation*}
k_{\mathrm{r}}=\sum_{n, n^{\prime}} k_{n n^{\prime}}(T) f_{n}(T) \tag{18D.4}
\end{equation*}


where $f_{n}(T)$ is the Boltzmann factor at a temperature $T$. It follows that if the state-to-state cross-sections can be determined or calculated for a wide range of approach speeds and initial and final states, then there is a route to the calculation of the rate constant for the reaction.

\subsection*{Brief illustration 18D. 1}
Suppose a harmonic oscillator collides with another oscillator of the same effective mass and force constant. If the state-tostate rate constant for the excitation of the latter's vibration is $k_{v v^{\prime}}=k_{\mathrm{r}}^{\circ} \mathrm{\delta}_{v v^{\prime}}$ for all the states $v$ and $v^{\prime}$, implying that an excitation can flow only from any level to the same level of the second oscillator, then at a temperature $T$, when $f_{v}(T)=\mathrm{e}^{-v h \nu / k T} / q$, where $q$ is the molecular vibrational partition function (Topic 13B), the overall rate constant is

$$
k_{\mathrm{r}}=\frac{k_{\mathrm{r}}^{\circ}}{\mathrm{G}} \sum_{v, v^{\prime}} \delta_{v v^{\prime}} \mathrm{e}^{-v h \nu / k T}=\frac{k_{\mathrm{r}}^{\circ}}{\mathrm{G}} \overbrace{\sum_{v^{\prime}} \mathrm{e}^{\mathrm{v}^{\prime} h \nu / k T}}^{\mathcal{G}}=k_{\mathrm{r}}^{\circ}
$$

\subsection*{18D. 3 Potential energy surfaces}
One of the most important concepts for discussing beam results and calculating the state-to-state collision cross-section is the potential energy surface of a reaction, the potential energy as a function of the relative positions of all the atoms taking part in the reaction. Potential energy surfaces may be constructed from experimental data and from results of quantum chemical calculations (Topic 9E). The theoretical method requires the systematic calculation of the energies of the system in a large number of geometrical arrangements. Special computational techniques, such as those described in Topic 9E, are used to take into account electron correlation, which arises from interactions between electrons as they move closer to and farther from each other in a molecule or molecular cluster. Techniques that incorporate electron correlation accurately are very time consuming and, consequently, the most reliable results are for reactions between relatively simple particles, such as the reactions $\mathrm{H}+\mathrm{H}_{2} \rightarrow \mathrm{H}_{2}+\mathrm{H}$ and $\mathrm{H}+\mathrm{H}_{2} \mathrm{O} \rightarrow$ $\mathrm{OH}+\mathrm{H}_{2}$. An alternative is to use semi-empirical methods, in which results of calculations and experimental parameters are used to construct the potential energy surface.

To illustrate the features of a potential energy surface, consider the collision between an H atom and an $\mathrm{H}_{2}$ molecule. Detailed calculations show that the approach of an atom $\mathrm{H}_{\mathrm{A}}$\\
\includegraphics[max width=\textwidth, center]{2025_02_27_d4b95d9718f0b9e2e6b9g-837(1)}

Figure 18D. 11 The potential energy surface for the $H+H_{2} \rightarrow H_{2}+$ H reaction when the atoms are constrained to be collinear.\\
along the $H_{B}-H_{C}$ axis requires less energy for reaction than any other approach, so initially it is convenient to confine attention to that collinear approach. Two parameters are required to define the nuclear separations: the $\mathrm{H}_{\mathrm{A}}-\mathrm{H}_{\mathrm{B}}$ separation $R_{\mathrm{AB}}$ and the $\mathrm{H}_{\mathrm{B}}-\mathrm{H}_{\mathrm{C}}$ separation $R_{\mathrm{BC}}$.

At the start of the encounter $R_{\mathrm{AB}}$ is effectively infinite and $R_{\mathrm{BC}}$ is the $\mathrm{H}_{2}$ equilibrium bond length. At the end of a successful reactive encounter $R_{\mathrm{AB}}$ is equal to the equilibrium bond length and $R_{B C}$ is infinite. The total energy of the three-atom system depends on their relative separations, and can be found by doing an electronic structure calculation. The plot of the total energy of the system against $R_{\mathrm{AB}}$ and $R_{\mathrm{BC}}$ gives the potential energy surface of this collinear reaction (Fig. 18D.11). This surface is normally depicted as a contour diagram (Fig. 18D.12).

When $R_{\mathrm{AB}}$ is very large, the variation in potential energy represented by the surface as $R_{\mathrm{BC}}$ changes is that of an isolated $\mathrm{H}_{2}$ molecule as its bond length is altered. A section through the surface at $R_{\mathrm{AB}}=\infty$, for example, is the same as the $\mathrm{H}_{2}$ bonding potential energy curve. At the edge of the diagram where\\
\includegraphics[max width=\textwidth, center]{2025_02_27_d4b95d9718f0b9e2e6b9g-837}

Figure 18D. 12 The contour diagram (with contours of equal potential energy) corresponding to the surface in Fig. 18D.11. $R_{\mathrm{e}}$ marks the equilibrium bond length of an $\mathrm{H}_{2}$ molecule (strictly, it relates to the arrangement when the third atom is at infinity).\\
$R_{B C}$ is very large, a section through the surface is the molecular potential energy curve of an isolated $\mathrm{H}_{\mathrm{A}} \mathrm{H}_{\mathrm{B}}$ molecule.

\subsection*{Brief illustration 18D. 2}
The bimolecular reaction $\mathrm{H}+\mathrm{O}_{2} \rightarrow \mathrm{OH}+\mathrm{O}$ plays an important role in combustion processes. The reaction can be characterized in terms of the $\mathrm{HO}_{2}$ potential energy surface and the two distances for collinear approach $R_{\mathrm{HO}_{\mathrm{A}}}$ and $R_{\mathrm{O}_{\mathrm{A}} \mathrm{O}_{\mathrm{B}}}$. When $R_{\mathrm{HO}_{\mathrm{A}}}$ is very large, the variation of the $\mathrm{HO}_{2}$ potential energy with $R_{\mathrm{O}_{\mathrm{A}} \mathrm{O}_{\mathrm{B}}}$ is that of an isolated dioxygen molecule as its bond length is changed. Similarly, when $R_{\mathrm{O}_{\mathrm{A}} \mathrm{O}_{\mathrm{B}}}$ is very large, a section through the potential energy surface is the molecular potential energy curve of an isolated OH radical.

The actual path of the atoms in the course of the encounter depends on their total energy, the sum of their kinetic and potential energies. However, an initial idea of the paths available to the system can be obtained by identifying paths that correspond to least potential energy. For example, consider the changes in potential energy as $\mathrm{H}_{\mathrm{A}}$ approaches $\mathrm{H}_{\mathrm{B}} \mathrm{H}_{\mathrm{C}}$. If the $\mathrm{H}_{\mathrm{B}}-\mathrm{H}_{\mathrm{C}}$ bond length is constant during the initial approach of $H_{A}$, then the potential energy of the $H_{3}$ cluster rises along the path marked A in Fig. 18D.13. The potential energy reaches a high value as $\mathrm{H}_{\mathrm{A}}$ is pushed into the molecule and then decreases sharply as $H_{C}$ breaks off and separates to a great distance. An alternative reaction path can be imagined (B) in which the $H_{B}-H_{C}$ bond length increases while $H_{A}$ is still far away. Both paths, although feasible if the molecules have sufficient initial kinetic energy, take the three atoms to regions of high potential energy in the course of the encounter.

The path of least potential energy is the one marked C, corresponding to $R_{\mathrm{BC}}$ lengthening as $\mathrm{H}_{\mathrm{A}}$ approaches and begins\\
\includegraphics[max width=\textwidth, center]{2025_02_27_d4b95d9718f0b9e2e6b9g-838(1)}

Figure 18D. 13 Various trajectories through the potential energy surface shown in Fig. 18D.12. Path A corresponds to a route in which $R_{\mathrm{BC}}$ is held nearly constant as $\mathrm{H}_{\mathrm{A}}$ approaches; path B corresponds to a route in which $R_{\mathrm{BC}}$ lengthens at an early stage during the approach of $\mathrm{H}_{A^{\prime}}$ path C is the route along the floor of the potential valley.\\
\includegraphics[max width=\textwidth, center]{2025_02_27_d4b95d9718f0b9e2e6b9g-838}

Figure 18D. 14 The transition state is a set of configurations (here, marked by the red line across the saddle point) through which successful reactive trajectories must pass.\\
to form a bond with $\mathrm{H}_{\mathrm{B}}$. The $\mathrm{H}_{\mathrm{B}}-\mathrm{H}_{\mathrm{C}}$ bond relaxes at the demand of the incoming atom, and the potential energy climbs only as far as the saddle-shaped region of the surface, to the saddle point marked $C^{\ddagger}$. The encounter that requires the least increase in potential energy is the one in which the atoms take route C up the floor of the valley, through the saddle point, and down the floor of the other valley as $\mathrm{H}_{\mathrm{C}}$ recedes and the new $H_{A}-H_{B}$ bond achieves its equilibrium length. This path is the reaction coordinate.

It is now possible to make contact with the transition-state theory of reaction rates (Topic 18C). In terms of trajectories on potential surfaces with a total energy close to the saddle point energy, the transition state can be identified with a critical range of configurations such that every trajectory that goes through this configuration goes on to react (Fig. 18D.14). Most trajectories on potential energy surfaces do not go directly over the saddle point and therefore, to result in a reaction, they require a total energy significantly higher than the saddlepoint energy. As a result, the experimentally determined activation energy is often significantly higher than the calculated saddle-point energy.

\subsection*{18D. 4 Some results from experiments and calculations}
Although quantum mechanical tunnelling can play an important role in reactivity, particularly in hydrogen atom and electron transfer reactions, this discussion begins with consideration of the classical trajectories of particles over surfaces. From this viewpoint, to travel successfully from reactants to products, the incoming molecules must possess enough kinetic energy to be able to climb to the saddle point of the potential surface. Therefore, the shape of the surface can be explored experimentally by changing the relative speed of approach (by selecting the beam velocity) and the degree of\\
\includegraphics[max width=\textwidth, center]{2025_02_27_d4b95d9718f0b9e2e6b9g-839}

Figure 18D. 15 Some successful (*) and unsuccessful encounters. (a) $\mathrm{C}_{1}^{*}$ corresponds to the path along the foot of the valley; (b) $C_{2}^{\star}$ corresponds to an approach of $A$ to a vibrating $B C$ molecule, and the formation of a vibrating $A B$ molecule as $C$ departs. (c) $C_{3}$ corresponds to $A$ approaching a non-vibrating $B C$ molecule, but with insufficient translational kinetic energy; (d) $C_{4}$ corresponds to A approaching a vibrating BC molecule, but still the energy, and the phase of the vibration, is insufficient for reaction.\\
vibrational excitation and observing whether reaction occurs and whether the products emerge in a vibrationally excited state (Fig. 18D.15). For example, one question that can be answered is whether it is better to smash the reactants together with a lot of translational kinetic energy or to ensure instead that they approach in highly excited vibrational states. Thus, is trajectory $\mathrm{C}_{2}^{*}$, where the $\mathrm{H}_{\mathrm{B}} \mathrm{H}_{\mathrm{C}}$ molecule is initially vibrationally excited, more efficient at leading to reaction than the trajectory $C_{1}^{*}$, in which the total energy is the same but reactants have a high translational kinetic energy?

\subsection*{(a) The direction of attack and separation}
Figure 18D. 16 shows the results of a calculation of the potential energy as an H atom approaches an $\mathrm{H}_{2}$ molecule from different angles, the $\mathrm{H}_{2}$ bond being allowed to relax to the optimum length in each case. The potential barrier is least for collinear attack, as assumed earlier. (But be aware that other lines of attack are feasible and contribute to the overall rate.) In contrast, Fig. 18D. 17 shows the potential energy changes that occur as a Cl atom approaches an HI molecule. The lowest barrier occurs for approaches within a cone of half-angle $30^{\circ}$ surrounding the H atom. The relevance of this result to the calculation of the steric factor of collision theory should be noted: not every collision is successful, because they do not all lie within the reactive cone.\\
\includegraphics[max width=\textwidth, center]{2025_02_27_d4b95d9718f0b9e2e6b9g-839(1)}

Figure 18D. 16 An indication of how the anisotropy of the potential energy changes as H approaches $\mathrm{H}_{2}$ with different angles of attack. The collinear attack has the lowest potential barrier to reaction. The surface indicates the potential energy profile along the reaction coordinate for each configuration.

If the collision is sticky, so that when the reactants collide they orbit around each other, the products can be expected to emerge in random directions because all memory of the approach direction has been lost. A rotation takes about 1 ps , so if the collision is over in less than that time the complex will not have had time to rotate and the products will be thrown off in a specific direction. In the collision of $\mathrm{K}^{2}$ and $\mathrm{I}_{2}$, for example, most of the products are thrown off in the forward direction ('forward' and 'backward' refer to directions in a centre-ofmass coordinate system with the origin at the centre of mass of the colliding reactants and collision occurring when molecules are at the origin). This product distribution is consistent with the harpoon mechanism (Topic 18A) because the transition takes place at long range. In contrast, the collision of K with $\mathrm{CH}_{3} \mathrm{I}$ leads to reaction only if the molecules approach each other very closely. In this mechanism, $K$ effectively bumps into a brick wall, and the KI product bounces out in the backward direction. The detection of this anisotropy in the angular distribution of products gives an indication of the distance and orientation of approach needed for reaction, as well as showing that the event is complete in less than about 1 ps .\\
\includegraphics[max width=\textwidth, center]{2025_02_27_d4b95d9718f0b9e2e6b9g-839(2)}

Figure 18D. 17 The potential energy barrier for the approach of Cl to HI . In this case, successful encounters occur only when Cl approaches almost directly towards the H atom.\\
\includegraphics[max width=\textwidth, center]{2025_02_27_d4b95d9718f0b9e2e6b9g-840(1)}

Figure 18D.18 An attractive potential energy surface. A successful encounter ( $C^{*}$ ) involves high translational kinetic energy and results in a vibrationally excited product.

\subsection*{(b) Attractive and repulsive surfaces}
Some reactions are very sensitive to whether the energy has been deposited into a vibrational mode or left as the relative translational kinetic energy of the colliding molecules. For example, if two HI molecules are hurled together with more than twice the activation energy of the reaction, then no reaction occurs if all the energy is solely translational. For $\mathrm{F}+\mathrm{HCl} \rightarrow$ $\mathrm{Cl}+\mathrm{HF}$, for example, the reaction is about five times more efficient when the HCl is in its first vibrational excited state than when it is in its vibrational ground state, although HCl has the same total energy.

The origin of these requirements can be found by examining the potential energy surface. Figure 18D. 18 shows an attractive surface in which the saddle point occurs early in the reaction coordinate. Figure 18D. 19 shows a repulsive surface in which the saddle point occurs late. A surface that is attractive in one direction is repulsive in the reverse direction.

Consider first the attractive surface. If the original molecule is vibrationally excited, then a collision with an incoming molecule takes the system along C. This path is bottled up in the region of the reactants, and does not take the system to the saddle point. If,\\
\includegraphics[max width=\textwidth, center]{2025_02_27_d4b95d9718f0b9e2e6b9g-840}

Figure 18D. 19 A repulsive potential energy surface. A successful encounter $\left(C^{*}\right)$ involves initial vibrational excitation and the products have high translational kinetic energy. A reaction that is attractive in one direction is repulsive in the reverse direction.\\
however, the same amount of energy is present solely as translational kinetic energy, then the system moves along $\mathrm{C}^{\star}$ and travels smoothly over the saddle point into products. Therefore, reactions with attractive potential energy surfaces proceed more efficiently if the energy is in relative translational motion. Moreover, the potential energy surface shows that once past the saddle point the trajectory runs up the steep wall of the product valley, and then rolls from side to side as it falls to the foot of the valley as the products separate. In other words, the products emerge in a vibrationally excited state.

Now consider the repulsive surface. On trajectory $C$ the collisional energy is largely in translation. As the reactants approach, the potential energy rises. Their path takes them up the opposing face of the valley, and they are reflected back into the reactant region. This path corresponds to an unsuccessful encounter, even though the energy is sufficient for reaction. On $\mathrm{C}^{*}$ some of the energy is in the vibration of the reactant molecule and the motion causes the trajectory to weave from side to side up the valley as it approaches the saddle point. This motion may be sufficient to tip the system round the corner to the saddle point and then on to products. In this case, the product molecule is expected to be in an unexcited vibrational state. Reactions with repulsive potential surfaces can therefore be expected to proceed more efficiently if the excess energy is present as vibrations. This is the case with the $\mathrm{H}+\mathrm{Cl}_{2} \rightarrow \mathrm{HCl}+$ Cl reaction, for instance.

\subsection*{Brief illustration 18D. 3}
The reaction $\mathrm{H}+\mathrm{Cl}_{2} \rightarrow \mathrm{HCl}+\mathrm{Cl}$ has a repulsive potential surface. Of the following four reactive processes, $\mathrm{H}+\mathrm{Cl}_{2}(v) \rightarrow$ $\mathrm{HCl}\left(v^{\prime}\right)+\mathrm{Cl}$, denoted $\left(v, v^{\prime}\right)$, all at the same total energy, (a) $(0,0)$, (b) $(2,0)$, (c) $(0,2),(d)(2,2)$, reaction (b) is most probable with reactants vibrationally excited and products vibrationally unexcited.

\subsection*{(c) Quantum mechanical scattering theory}
A picture of the reaction event can be obtained by using classical mechanics to calculate the trajectories of the atoms taking place in a reaction from a set of initial conditions, such as velocities, relative orientations, and internal energies of the reacting particles. However, classical trajectory calculations do not recognize the fact that the motion of atoms, electrons, and nuclei is governed by quantum mechanics. The concept of trajectory then fades and is replaced by the unfolding of a wavefunction that represents initially the reactants and finally products.

Complete quantum mechanical calculations of trajectories and rate constants are very onerous because it is necessary to take into account all the allowed electronic, vibrational, and rotational states populated by each atom and molecule in the sys-\\
tem at a given temperature. It is common to define a 'channel' as a group of molecules in well-defined quantum mechanically allowed states. Then, at a given temperature, there are many channels that represent the reactants and many channels that represent possible products, with some transitions between channels being allowed but others not allowed. Furthermore, not every transition leads to a chemical reaction. For example, the process $\mathrm{H}_{2}^{*}+\mathrm{OH} \rightarrow \mathrm{H}_{2}+\mathrm{OH}^{*}$, where the asterisk denotes an excited state, amounts to energy transfer between $\mathrm{H}_{2}$ and OH , whereas the process $\mathrm{H}_{2}^{*}+\mathrm{OH} \rightarrow \mathrm{H}_{2} \mathrm{O}+\mathrm{H}$ represents a chemical reaction. What complicates a quantum mechanical calculation of rate constants even in this simple four-atom system is that many reacting channels present at a given temperature can lead to the desired products $\mathrm{H}_{2} \mathrm{O}+\mathrm{H}$, which themselves may be formed as many distinct channels. The cumulative reaction probability, $\bar{P}(E)$, at a fixed total energy $E$ is then written as

$$
\bar{P}(E)=\sum_{i, j} P_{i j}(E) \quad \text { Cumulative reaction probability }
$$

(18D.5)\\
where $P_{i j}(E)$ is the probability for a transition between a reacting channel $i$ and a product channel $j$ and the summation is over all possible transitions that lead to product. It is then possible to show that the rate constant is given by

$$
k_{\mathrm{r}}(T)=\frac{\int_{0}^{\infty} \bar{P}(E) \mathrm{e}^{-E / k T} \mathrm{~d} E}{h Q_{\mathrm{R}}(T)}
$$

Rate constant (18D.6)\\
where $Q_{R}(T)$ is the partition function density (the partition function divided by the volume) of the reactants at the temperature $T$. The significance of eqn 18D. 6 is that it provides a direct connection between an experimental quantity, the rate constant, and a theoretical quantity, $\bar{P}(E)$.

\subsection*{Checklist of concepts}
\begin{enumerate}
  \item A molecular beam is a collimated, narrow stream of molecules travelling through an evacuated vessel.2. In a molecular beam, the scattering pattern of molecules depends on quantum mechanical effects and the details of the intermolecular potential.3. A van der Waals molecule is a complex of the form $A B$ in which $A$ and $B$ are held together by van der Waals forces or hydrogen bonds.4. Techniques for the study of reactive collisions include infrared chemiluminescence, laser-induced fluorescence, multiphoton ionization (MPI), resonant mul-\\
tiphoton ionization (REMPI), and reaction product imaging.5. A potential energy surface maps the potential energy as a function of the relative positions of all the atoms taking part in a reaction.
  \item In an attractive surface, the saddle point (the highest point on the valley between reactants and products) occurs early on the reaction coordinate.7. In a repulsive surface, the saddle point occurs late on the reaction coordinate.
\end{enumerate}

\subsection*{Checklist of equations}
\begin{center}
\begin{tabular}{llll}
Property & Equation & Comment & Equation number \\
Rate of molecular scattering & $\mathrm{d} I=\sigma I \mathcal{N d} x$ & $\sigma$ is the differential scattering cross-section & 18 D .1 \\
Rate constant & $k_{\mathrm{r}}=\left\langle\sigma v_{\mathrm{rel}}\right\rangle N_{\mathrm{A}}$ &  & 18 D .2 \\
State-to-state rate constant & $k_{n n^{\prime}}=\left\langle\sigma_{n n^{\prime}} v_{\mathrm{rel}}\right\rangle N_{\mathrm{A}}$ & 18 D .3 &  \\
Overall rate constant & $k_{\mathrm{r}}=\sum_{n, n^{\prime}} k_{n n^{\prime}}(T) f_{n}(T)$ & 18 D .4 &  \\
Cumulative reaction probability & $\bar{P}(E)=\sum_{i, j} P_{i j}(E)$ & $Q_{\mathrm{R}}(T)$ is the partition function density & 18 D .5 \\
Rate constant & $k_{\mathrm{r}}(T)=\int_{0}^{\infty} \bar{P}(E) \mathrm{e}^{-E / k T} \mathrm{~d} E / h Q_{\mathrm{R}}(T)$ & 18 D .6 &  \\
\hline
\end{tabular}
\end{center}

\section*{TOPIC 18E Electron transfer in homogeneous systems}
\subsection*{Why do you need to know this material?}
Electron transfer reactions between protein-bound cofactors or between proteins play an important role in a variety of biological processes. Electron transfer is also important in homogeneous, non-biological catalysis.

\subsection*{What is the key idea?}
The rate constant of electron transfer in a donor-acceptor complex depends on the distance between electron donor and acceptor, the standard reaction Gibbs energy, and the energy needed to reach a particular arrangement of atoms.

What do you need to know already?\\
This Topic makes use of transition-state theory (Topic 18C). It also uses the concept of tunnelling (Topic 7D), the steady-state approximation (Topic 17E), and the FranckCondon principle (Topic 11F).

Concepts of transition-state theory and quantum theory can be applied to the study of a deceptively simple process, electron transfer between molecules in homogeneous systems. The approach developed here can be used to predict the rates of electron transfer of this kind with reasonable accuracy.

\subsection*{18E. 1 The rate law}
Consider electron transfer from a donor species $D$ to an acceptor species A in solution. The overall reaction is


\begin{equation*}
\mathrm{D}+\mathrm{A} \rightarrow \mathrm{D}^{+}+\mathrm{A}^{-} \quad v=k_{\mathrm{r}}[\mathrm{D}][\mathrm{A}] \tag{18E.1}
\end{equation*}


In the first step of the mechanism, D and A must diffuse through the solution and on meeting form a complex DA:


\begin{equation*}
\mathrm{D}+\mathrm{A} \underset{k_{\mathrm{a}}^{\prime}}{\stackrel{k_{\mathrm{a}}}{\rightleftarrows}} \mathrm{DA} \tag{18E.2a}
\end{equation*}


Suppose that in the complex D and A are separated by $d$, the distance between their outer surfaces. Then electron transfer occurs within the DA complex to yield $\mathrm{D}^{+} \mathrm{A}^{-}$:


\begin{equation*}
\mathrm{DA} \underset{k_{\mathrm{ct}}^{\prime}}{\stackrel{k_{\mathrm{e}}}{\rightleftarrows}} \mathrm{D}^{+} \mathrm{A}^{-} \tag{18E.2b}
\end{equation*}


The complex $\mathrm{D}^{+} \mathrm{A}^{-}$can also break apart and the ions diffuse through the solution:


\begin{equation*}
\mathrm{D}^{+} \mathrm{A}^{-} \xrightarrow{k_{\mathrm{d}}} \mathrm{D}^{+}+\mathrm{A}^{-} \tag{18E.2c}
\end{equation*}


Now the techniques described in Topic 17E can be used to write an expression for the rate constant for the overall reaction $\mathrm{D}+\mathrm{A} \rightarrow \mathrm{D}^{+}+\mathrm{A}^{-}$.

\subsection*{How is that done? 18E. 1}
Deriving an expression for the rate constant for electron transfer in solution

Identify the rate of the overall reaction (eqn 18E.1) with the rate of the step described by eqn 18E. 2 c because the products of the reaction are the separated ions:

$$
v=k_{\mathrm{d}}\left[\mathrm{D}^{+} \mathrm{A}^{-}\right]
$$

Then follow these steps.\\
Step 1 Apply the steady-state approximation to both intermediates\\
There are two reaction intermediates, DA and $\mathrm{D}^{+} \mathrm{A}^{-}$, so apply the steady-state approximation (Topic 17E) to both. From

$$
\frac{\mathrm{d}\left[\mathrm{D}^{+} \mathrm{A}^{-}\right]}{\mathrm{d} t}=k_{\mathrm{et}}[\mathrm{DA}]-k_{\mathrm{et}}^{\prime}\left[\mathrm{D}^{+} \mathrm{A}^{-}\right]-k_{\mathrm{d}}\left[\mathrm{D}^{+} \mathrm{A}^{-}\right]=0
$$

it follows that

$$
[\mathrm{DA}]=\frac{k_{\mathrm{et}}^{\prime}+k_{\mathrm{d}}}{k_{\mathrm{et}}}\left[\mathrm{D}^{+} \mathrm{A}^{-}\right]
$$

The steady-state expression for DA is

$$
\frac{\mathrm{d}[\mathrm{DA}]}{\mathrm{d} t}=k_{\mathrm{a}}[\mathrm{D}][\mathrm{A}]-k_{\mathrm{a}}^{\prime}[\mathrm{DA}]-k_{\mathrm{et}}[\mathrm{DA}]+k_{\mathrm{et}}^{\prime}\left[\mathrm{D}^{+} \mathrm{A}^{-}\right]=0
$$

Next, replace the terms in blue by the expression for [DA] from the preceding equation to give

$$
\begin{aligned}
k_{\mathrm{a}}[\mathrm{D}][\mathrm{A}] & -\frac{\left(k_{\mathrm{et}}^{\prime}+k_{\mathrm{d}}\right)\left(k_{\mathrm{a}}^{\prime}+k_{\mathrm{et}}\right)}{k_{\mathrm{et}}}\left[\mathrm{D}^{+} \mathrm{A}^{-}\right]+k_{\mathrm{et}}^{\prime}\left[\mathrm{D}^{+} \mathrm{A}^{-}\right] \\
& =k_{\mathrm{a}}[\mathrm{D}][\mathrm{A}]-\left[\mathrm{D}^{+} \mathrm{A}^{-}\right]\left[\frac{\left(k_{\mathrm{et}}^{\prime}+k_{\mathrm{d}}\right)\left(k_{\mathrm{a}}^{\prime}+k_{\mathrm{et}}\right)}{k_{\mathrm{et}}}-k_{\mathrm{et}}^{\prime}\right\} \\
& =k_{\mathrm{a}}[\mathrm{D}][\mathrm{A}]-\left[\mathrm{D}^{+} \mathrm{A}^{-}\right] \frac{\left(k_{\mathrm{et}}^{\prime}+k_{\mathrm{d}}\right)\left(k_{\mathrm{a}}^{\prime}+k_{\mathrm{et}}\right)-k_{\mathrm{et}}^{\prime} k_{\mathrm{et}}}{k_{\mathrm{et}}}=0
\end{aligned}
$$

It follows that

$$
\begin{aligned}
{\left[\mathrm{D}^{+} \mathrm{A}^{-}\right] } & =\frac{k_{\mathrm{a}} k_{\mathrm{et}}}{\left(k_{\mathrm{et}}^{\prime}+k_{\mathrm{d}}\right)\left(k_{\mathrm{a}}^{\prime}+k_{\mathrm{et}}\right)-k_{\mathrm{et}}^{\prime} k_{\mathrm{et}}}[\mathrm{D}][\mathrm{A}] \\
& =\frac{k_{\mathrm{a}} k_{\mathrm{et}}}{k_{\mathrm{et}}^{\prime} k_{\mathrm{a}}^{\prime}+k_{\mathrm{d}} k_{\mathrm{a}}^{\prime}+k_{\mathrm{d}} k_{\mathrm{et}}}[\mathrm{D}][\mathrm{A}]
\end{aligned}
$$

Step 2 Write an expression for the rate constant\\
The overall rate of reaction is $v=k_{\mathrm{d}}\left[\mathrm{D}^{+} \mathrm{A}^{-}\right]$, which now becomes

$$
v=k_{\mathrm{d}} \frac{k_{\mathrm{a}} k_{\mathrm{et}}}{k_{\mathrm{et}}^{\prime} k_{\mathrm{a}}^{\prime}+k_{\mathrm{d}} k_{\mathrm{a}}^{\prime}+k_{\mathrm{d}} k_{\mathrm{et}}}[\mathrm{D}][\mathrm{A}]
$$

This expression is of the form $v=k_{\mathrm{r}}[\mathrm{D}][\mathrm{A}]$, with $k_{\mathrm{r}}$ given by

$$
k_{\mathrm{r}}=\frac{k_{\mathrm{a}} k_{\mathrm{ee}} k_{\mathrm{d}}}{k_{\mathrm{a}}^{\prime} k_{\mathrm{et}}^{\prime}+k_{\mathrm{d}} k_{\mathrm{a}}^{\prime}+k_{\mathrm{d}} k_{\mathrm{et}}}=\frac{k_{\mathrm{a}} k_{\mathrm{et}} k_{\mathrm{d}}}{k_{\mathrm{a}}^{\prime}\left(k_{\mathrm{et}}^{\prime}+k_{\mathrm{d}}\right)+k_{\mathrm{d}} k_{\mathrm{et}}}
$$

Step 3 Rearrange the preceding expression\\
You can obtain a more convenient form of the preceding expression by dividing the numerator and denominator on the right-hand side by $k_{\mathrm{d}} k_{\mathrm{et}}$ to obtain

$$
k_{\mathrm{r}}=\frac{k_{\mathrm{a}}}{k_{\mathrm{a}}^{\prime}\left(k_{\mathrm{et}}^{\prime}+k_{\mathrm{d}}\right) / k_{\mathrm{et}} k_{\mathrm{d}}+1}
$$

The reciprocal of each side then gives

$$
\frac{1}{k_{\mathrm{r}}}=\frac{1}{k_{\mathrm{a}}}+\frac{k_{\mathrm{a}}^{\prime}}{k_{\mathrm{a}} k_{\mathrm{et}} k_{\mathrm{d}}}\left(k_{\mathrm{et}}^{\prime}+k_{\mathrm{d}}\right)
$$

and therefore

$$
\left.-\frac{1}{k_{\mathrm{r}}}=\frac{1}{k_{\mathrm{a}}}+\frac{k_{\mathrm{a}}^{\prime}}{k_{\mathrm{a}} k_{\mathrm{et}}}\left(1+\frac{k_{\mathrm{et}}^{\prime}}{k_{\mathrm{d}}}\right) \right\rvert\, \quad \text { (18E.3) } \quad \text { Electron transfer rate constant }
$$

To gain insight into this equation and the factors that determine the rate of electron transfer reactions in solution, assume that the main decay route for $\mathrm{D}^{+} \mathrm{A}^{-}$is dissociation of the complex into separated ions, and therefore that $k_{\mathrm{d}}\left[\mathrm{D}^{+} \mathrm{A}^{-}\right] \gg$ $k_{\text {et }}^{\prime}\left[\mathrm{D}^{+} \mathrm{A}^{-}\right]$, which implies that $k_{\mathrm{d}} \gg k_{\mathrm{et}}^{\prime}$. It follows that

$$
\frac{1}{k_{\mathrm{r}}} \approx \frac{1}{k_{\mathrm{a}}}\left(1+\frac{k_{\mathrm{a}}^{\prime}}{k_{\mathrm{et}}}\right)
$$

\begin{itemize}
  \item When $k_{\text {et }}[\mathrm{DA}] \gg k_{\mathrm{a}}^{\prime}[\mathrm{DA}]$, which implies that $k_{\mathrm{r}} \approx k_{\mathrm{a}}$, the rate of product formation is controlled by diffusion of D and A in solution.
  \item When $k_{\mathrm{et}}[\mathrm{DA}] \ll k_{\mathrm{a}}^{\prime}[\mathrm{DA}]$, it follows that $k_{\mathrm{r}} \approx\left(k_{\mathrm{a}} / k_{\mathrm{a}}^{\prime}\right) k_{\mathrm{et}}$ $=K k_{\text {et, }}$ where $K$ is the equilibrium constant for the diffusive encounter. The process is controlled by $k_{\text {et }}$ and therefore the activation energy of electron transfer in the DA complex.
\end{itemize}

\subsection*{18E.2 The role of electron tunnelling}
This analysis can be taken further by introducing the implication from transition-state theory that, at a given temperature, $k_{\text {et }} \propto \mathrm{e}^{-\Delta^{\ddagger} G / R T}$, where $\Delta^{\ddagger} G$ is the Gibbs energy of activation. The remaining task, therefore, is to write expressions for the proportionality constant and $\Delta^{\ddagger} G$. The discussion concentrates on the following two key aspects of the theory of electron transfer processes, which was developed independently by R.A. Marcus, N.S. Hush, V.G. Levich, and R.R. Dogonadze:

\begin{itemize}
  \item Electrons are transferred by tunnelling through a potential energy barrier, the height of which is partly determined by the ionization energies of the DA and $\mathrm{D}^{+} \mathrm{A}^{-}$ complexes. Electron tunnelling influences the magnitude of the proportionality constant in the expression for $k_{\mathrm{et}}$.
  \item The complex DA and the solvent molecules surrounding it undergo structural rearrangements prior to electron transfer. The energy associated with these rearrangements and the standard reaction Gibbs energy determine $\Delta^{\ddagger} G$.
\end{itemize}

According to the Franck-Condon principle (Topic 11F), electronic transitions are so fast that they can be regarded as taking place in a stationary nuclear framework. This principle also applies to an electron transfer process in which an electron migrates from one energy surface, representing the dependence of the energy of DA on its geometry, to another representing the energy of $\mathrm{D}^{+} \mathrm{A}^{-}$. The potential energy (and the Gibbs energy) surfaces of the two complexes (the reactant complex, DA , and the product complex, $\mathrm{D}^{+} \mathrm{A}^{-}$) can be represented by the parabolas characteristic of harmonic oscillators, with the displacement coordinate corresponding to the changing geometries (Fig. 18E.1). This coordinate represents a collective mode of the donor, acceptor, and solvent.

According to the Franck-Condon principle, the nuclei do not have time to move when the system passes from the reactant to the product surface as a result of the transfer of an electron. Therefore, electron transfer can occur only after thermal fluctuations bring the geometry of DA to $q^{\ddagger}$ in Fig. 18E.1, the value of the nuclear coordinate at which the two parabolas intersect.\\
\includegraphics[max width=\textwidth, center]{2025_02_27_d4b95d9718f0b9e2e6b9g-844}

Figure 18E. 1 The Gibbs energy surfaces of the complexes DA and $D^{+} A^{-}$involved in an electron transfer process are represented by parabolas characteristic of harmonic oscillators, with the displacement coordinate $q$ corresponding to the changing geometries of the system.

The proportionality constant in the expression for $k_{\mathrm{et}}$ is a measure of the rate at which the system converts from reactants (DA) to products ( $\mathrm{D}^{+} \mathrm{A}^{-}$) at $q^{\ddagger}$ by electron transfer within the thermally excited DA complex. To understand the process, consider the effect that the rearrangement of nuclear coordinates has on the electronic energy levels of DA and $\mathrm{D}^{+} \mathrm{A}^{-}$for a given distance $d$ between D and A (Fig. 18E.2). Initially, the HOMO of DA is lower than the LUMO of $\mathrm{D}^{+} \mathrm{A}^{-}$ (Fig. 18E.2a). As the nuclei rearrange into a configuration represented by $q^{\ddagger}$ in Fig. 18E.2b, the HOMO of DA and the LUMO of $\mathrm{D}^{+} \mathrm{A}^{-}$become similar in energy and electron transfer becomes feasible. Over reasonably short distances $d$, the main mechanism of electron transfer is tunnelling through the potential energy barrier depicted in Fig. 18E.2b. After an electron moves from the HOMO of DA to the LUMO of $\mathrm{D}^{+} \mathrm{A}^{-}$, the system relaxes to the configuration represented by $q_{0}^{p}$ in Fig. 18E.2c. As shown in the illustration, now the energy of $\mathrm{D}^{+} \mathrm{A}^{-}$is lower than that of DA , reflecting the thermodynamic tendency for A to remain reduced (as A') and for D to remain oxidized (as $\mathrm{D}^{+}$).\\
The tunnelling event responsible for electron transfer is similar to that described in Topic 7D, except that in this case the electron tunnels from an electronic level of DA , with wavefunction $\psi_{\mathrm{DA}}$, to an electronic level of $\mathrm{D}^{+} \mathrm{A}^{-}$, with wavefunction $\psi_{\mathrm{D}^{+} \mathrm{A}^{-}}$. The rate of an electronic transition from a level described by the wavefunction $\psi_{D A}$ to a level described by $\psi_{D^{+} A^{-}}$is proportional to the square of the integral

$$
H_{\mathrm{et}}=\int \psi_{\mathrm{DA}} \hat{h} \psi_{\mathrm{D}^{+} \mathrm{A}^{-}} \mathrm{d} \tau
$$

where $\hat{h}$ is a hamiltonian that describes the coupling of the electronic wavefunctions. The quantity $H_{\mathrm{et}}$ is often referred to as the 'electronic coupling matrix element'. The probability of tunnelling through a potential barrier typically has an ex-\\
\includegraphics[max width=\textwidth, center]{2025_02_27_d4b95d9718f0b9e2e6b9g-844(1)}

Figure 18E. 2 (a) At the nuclear configuration denoted by $q_{0}^{R}$, the electron to be transferred in DA is in the HOMO; the LUMO of $\mathrm{D}^{+} \mathrm{A}^{-}$ is too high in energy for efficient electron transfer. (b) As the nuclei rearrange to a configuration represented by $q^{\ddagger}$, the HOMO of DA and LUMO of $\mathrm{D}^{+} \mathrm{A}^{-}$become similar in energy and electron transfer occurs by tunnelling. (c) The system relaxes to the equilibrium nuclear configuration of $\mathrm{D}^{+} \mathrm{A}^{-}$denoted by $q_{0^{\prime}}^{\mathrm{p}}$ in which the LUMO of DA is higher in energy than the HOMO of $\mathrm{D}^{+} \mathrm{A}^{-}$. Adapted from R.A. Marcus and N. Sutin, Biochim. Biophys. Acta 811, 265 (1985).\\
ponential dependence on the width of the barrier (Topic 7D), suggesting that


\begin{equation*}
H_{\mathrm{et}}(d)^{2}=H_{\mathrm{et}}^{\circ 2} \mathrm{e}^{-\beta d} \tag{18E.4}
\end{equation*}


where $d$ is the edge-to-edge distance between D and $\mathrm{A}, \beta$ is a parameter that measures the sensitivity of the electronic coupling matrix element to distance, and $H_{\mathrm{et}}^{\circ}$ is the value of the electronic coupling matrix element at $d=0$. The value of $\beta$ depends on the medium through which the electron must travel from donor to acceptor. In a vacuum, $28 \mathrm{~nm}^{-1}<\beta<35 \mathrm{~nm}^{-1}$, whereas $\beta \approx 9 \mathrm{~nm}^{-1}$ when the intervening medium is a molecular link between donor and acceptor.

\subsection*{18E.3 The rate constant}
The proportionality constant in $k_{\text {et }} \propto \mathrm{e}^{-\Delta^{\ddagger} G / R T}$ is proportional to $H_{\mathrm{et}}(d)^{2}$, as expressed by eqn 18E.4. A detailed calculation (not reproduced here) shows that the full expression for $k_{\text {et }}$ is


\begin{equation*}
k_{\mathrm{et}}=\frac{1}{h}\left(\frac{\pi^{3}}{R T \Delta E_{\mathrm{R}}}\right)^{1 / 2} H_{\mathrm{et}}(d)^{2} \mathrm{e}^{-\Delta^{\ddagger} G / R T} \tag{18E.5}
\end{equation*}


where $\Delta E_{\mathrm{R}}$ is the reorganization energy, the energy change associated with the molecular rearrangement that must take place so that DA can take on the equilibrium geometry of\\
$\mathrm{D}^{+} \mathrm{A}^{-}$. These molecular rearrangements include the relative reorientation of the D and A molecules in DA and the relative reorientation of the solvent molecules surrounding DA. To use eqn 18E. 5 it is necessary to find an expression for the Gibbs energy of activation in terms of a simple model of the reaction.

How is that done? 18E. 2 Establishing an expression for the Gibbs energy of activation

The simplest way to derive an expression for the Gibbs energy of activation of electron transfer is to construct a model in which the energy surfaces of DA (the 'reactant complex', denoted R ) and $\mathrm{D}^{+} \mathrm{A}^{-}$(the 'product complex', denoted P ) are plotted against the reaction coordinate $q$ and assumed to be identical parabolic curves with displaced minima (Fig. 18E.3). For simplicity, $q$ can be set equal to 0 at the minimum of the reactant parabola, and Gibbs energies measured from that minimum. Then

$$
G_{\mathrm{m}, \mathrm{R}}(q)=k_{e} q^{2} \text { and } G_{\mathrm{m}, \mathrm{P}}(q)=k_{\mathrm{e}}\left(q-q_{0}^{\mathrm{P}}\right)^{2}+\Delta_{\mathrm{r}} G^{\ominus}
$$

where $q_{0}^{\mathrm{P}}$ is the location of the minimum of the product parabola, $k_{\mathrm{e}}$ is a constant describing the curvature of the parabola, and $\Delta_{\mathrm{r}} G^{\ominus}$ is the standard reaction Gibbs energy for the electron transfer process $\mathrm{DA} \rightarrow \mathrm{D}^{+} \mathrm{A}^{-}$.

Write $\Delta E_{\mathrm{R}}=G_{\mathrm{m}, \mathrm{R}}\left(q_{0}^{\mathrm{P}}\right)=k_{e} q_{0}^{\mathrm{P} 2}$ as the difference in the Gibbs energy of R when the coordinate changes from $q=0$ to the equilibrium value for P, $q_{0}^{\mathrm{P}}$. The Gibbs energy of activation, $\Delta^{\ddagger} G=$ $G_{\mathrm{m}, \mathrm{R}}\left(q^{\ddagger}\right)=k_{\mathrm{e}} q^{\ddagger 2}$, is the change in Gibbs energy of R when the coordinate changes from $q=0$ to $q^{\ddagger}$. Then follow these steps.

\subsection*{Step 1 Identify the location of the activated complex}
The activated complex occurs at the point where the two parabolas intersect, which is where

$$
\begin{aligned}
k_{e} q^{\ddagger 2} & =k_{\mathrm{e}}\left(q^{\ddagger}-q_{0}^{\mathrm{P}}\right)^{2}+\Delta_{\mathrm{r}} G^{\ominus} \\
& =k_{\mathrm{e}} q^{\ddagger 2}-2 k_{\mathrm{e}} q_{0}^{\mathrm{p}} q^{\ddagger}+\overbrace{k_{\mathrm{e}} q_{0}^{\mathrm{p} 2}}^{\Delta E_{\mathrm{R}}}+\Delta_{\mathrm{r}} G^{\ominus}
\end{aligned}
$$

\begin{center}
\includegraphics[max width=\textwidth]{2025_02_27_d4b95d9718f0b9e2e6b9g-845}
\end{center}

Figure 18E. 3 The model system used in the calculation of the Gibbs energy of activation for an electron transfer process.

After rearrangement, the expression for $q^{\ddagger}$ is

$$
q^{\ddagger}=\frac{\Delta E_{\mathrm{R}}+\Delta_{\mathrm{r}} G^{\ominus}}{2 k_{\mathrm{e}} q_{0}^{\mathrm{p}}}
$$

Step 2 Use the expression for $q^{\ddagger}$ to find an expression for $\Delta^{\ddagger} G$\\
Because Gibbs energies are measured from the minimum the reactant parabola, where $q=0$, it follows from the expression for $q^{\ddagger}$ that

$$
\Delta^{\ddagger} G=k_{\mathrm{e}} q^{\ddagger 2}=k_{\mathrm{e}}\left(\frac{\Delta E_{\mathrm{R}}+\Delta_{\mathrm{r}} G^{\ominus}}{2 k_{\mathrm{e}} q_{0}^{\mathrm{P}}}\right)^{2}=\frac{\left(\Delta E_{\mathrm{R}}+\Delta_{\mathrm{r}} G^{\ominus}\right)^{2}}{4 \underbrace{k_{e^{P} q_{0}^{P 2}}}_{\Delta E_{\mathrm{R}}}}
$$

which simplifies to


\begin{equation*}
-\Delta^{\ddagger} G=\frac{\left(\Delta_{\mathrm{r}} G^{\ominus}+\Delta E_{\mathrm{R}}\right)^{2}}{4 \Delta E_{\mathrm{R}}} \tag{18E.6}
\end{equation*}


This equation shows that $\Delta^{\ddagger} G=0$ when $\Delta_{\mathrm{r}} G^{\ominus}=-\Delta E_{\mathrm{R}}$, and the reorganization energy is cancelled by the standard reaction Gibbs energy. If this condition holds, the reaction is not slowed down by an activation barrier.

Equation 18E. 6 has some limitations. For instance, it describes processes with weak electronic coupling between donor and acceptor. Weak coupling is observed when the electroactive species are sufficiently far apart ( $d>1 \mathrm{~nm}$ ) that the tunnelling is an exponential function of distance. The weak coupling limit applies to a large number of electron transfer reactions, including those between proteins during metabolism. Strong coupling is observed when the wavefunctions $\psi_{\mathrm{A}}$ and $\psi_{\mathrm{D}}$ overlap very extensively and, as well as other complications, the tunnelling rate is no longer a simple exponential function of distance. Examples of strongly coupled systems are mixed-valence, binuclear d-metal complexes with the general structure $\mathrm{L}_{m} \mathrm{M}^{n+}-\mathrm{B}-\mathrm{M}^{p+} \mathrm{L}_{m}$, in which the electroactive metal ions are separated by a bridging ligand B. In these systems, $d<1.0 \mathrm{~nm}$.

\subsection*{18E. 4 Experimental tests of the theory}
The most meaningful experimental tests of the dependence of $k_{\text {et }}$ on $d$ are those in which the same donor and acceptor are positioned at a variety of distances, perhaps by covalent attachment to molecular linkers (see 1 for an example, in which the biphenyl group is the donor and A are various acceptors). Under these conditions, the term $\mathrm{e}^{-\Delta^{\ddagger} G / R T}$ becomes a constant and, after taking the natural logarithm of eqn 18E. 5 and using eqn 18E.4, the result is


\begin{equation*}
\ln k_{\mathrm{et}}=-\beta d+\mathrm{constant} \tag{18E.7}
\end{equation*}


which implies that a plot of $\ln k_{\text {et }}$ against $d$ should be a straight line of slope $-\beta$.\\
\includegraphics{smile-baedf1d3205e7878b0ace9f480bfe92c179896c7}

The dependence of $k_{\text {et }}$ on the standard reaction Gibbs energy has been investigated in systems where the edge-to-edge distance and the reorganization energy are constant for a series of reactions. Then, because $k_{\text {et }} \propto \mathrm{e}^{-\Delta^{\ddagger} G / R T}$ implies that $\ln k_{\mathrm{et}}=-\Delta^{\ddagger} G / R T+$ constant, it follows from eqn 18E. 6 that

$$
\begin{aligned}
\ln k_{\mathrm{et}} & =-\frac{\left(\Delta_{\mathrm{r}} G^{\ominus}+\Delta E_{\mathrm{R}}\right)^{2}}{4 R T \Delta E_{\mathrm{R}}}+\text { constant } \\
& =-\frac{\left(\Delta_{\mathrm{r}} G^{\ominus}\right)^{2}}{4 R T \Delta E_{\mathrm{R}}}-\frac{\Delta_{\mathrm{r}} G^{\ominus}}{2 R T} \overbrace{-\frac{\Delta E_{\mathrm{R}}}{4 R T}+\text { constant }}^{\text {A constant }}
\end{aligned}
$$

and therefore


\begin{equation*}
\ln k_{\mathrm{et}}=-\frac{R T}{4 \Delta E_{\mathrm{R}}}\left(\frac{\Delta_{\mathrm{r}} G^{\ominus}}{R T}\right)^{2}-\frac{1}{2}\left(\frac{\Delta_{\mathrm{r}} G^{\ominus}}{R T}\right)+\text { constant } \tag{18E.8}
\end{equation*}


A plot of $\ln k_{\mathrm{et}}\left(\right.$ or $\left.\log k_{\mathrm{et}}=\ln k_{\mathrm{et}} / \ln 10\right)$ against $\Delta_{\mathrm{r}} G^{\ominus}\left(\right.$ or $\left.-\Delta_{\mathrm{r}} G^{\ominus}\right)$ should therefore be a downward parabola (a curve of the form $y=a x^{2}+b x+c$; Fig. 18E.4). Equation 18E. 8 implies that the rate constant increases as $\Delta_{\mathrm{r}} G^{\ominus}$ decreases but only up to $-\Delta_{\mathrm{r}} G^{\ominus}=\Delta E_{\mathrm{R}}$. Beyond that, the reaction enters the inverted region, in which the rate constant decreases as the reaction becomes more exergonic ( $\Delta_{\mathrm{r}} G^{\ominus}$ becomes more negative). The inverted region has been observed in a series of special compounds in which the electron donor and acceptor are\\
\includegraphics[max width=\textwidth, center]{2025_02_27_d4b95d9718f0b9e2e6b9g-846(1)}

Figure 18E. 4 The parabolic dependence of $\ln k_{\text {et }}$ on $-\Delta_{\mathrm{r}} G^{\ominus}$ predicted by eqn 18E.8.\\
\includegraphics[max width=\textwidth, center]{2025_02_27_d4b95d9718f0b9e2e6b9g-846(2)}

Figure 18E.5 Variation of $\log k_{\text {et }}$ with $-\Delta_{\mathrm{r}} G^{\ominus}$ for a series of compounds with the structures given in (1) and as described in Brief illustration 18E.1. Based on J.R. Miller et al., J. Am. Chem. Soc. 106, 3047 (1984).\\
linked covalently to a molecular spacer of known and fixed size (Fig. 18E.5).

\subsection*{Brief illustration 18E. 1}
Kinetic measurements were conducted in 2-methyltetrahydrofuran and at 296 K for a series of compounds with the structures given in (1) with A the following groups:\\
\includegraphics[max width=\textwidth, center]{2025_02_27_d4b95d9718f0b9e2e6b9g-846}

The distance between donor (the reduced biphenyl group) and the acceptor is constant for all compounds in the series because the molecular linker remains the same. Each acceptor has a characteristic standard potential, so it follows that the standard Gibbs energy for the electron transfer process is different for each compound in the series. The line in Fig. 18E. 5 is a fit to a version of eqn 18E. 8 and the maximum of the parabola occurs at $-\Delta_{\mathrm{r}} G^{\ominus}=\Delta E_{\mathrm{R}}=1.4 \mathrm{eV}=1.4 \times 10^{2} \mathrm{~kJ} \mathrm{~mol}^{-1}$.

\subsection*{Checklist of concepts}
1. Electron transfer can occur through tunnelling only after thermal fluctuations bring the nuclear coordinate to the point at which the donor and acceptor have the same configuration.2. The reorganization energy is the energy change associated with molecular rearrangements that must take place so that DA can acquire the equilibrium geometry of $\mathrm{D}^{+} \mathrm{A}^{-}$.\subsection*{Checklist of equations}
\begin{center}
\begin{tabular}{llll}
\hline
Property & Equation & Comment & Equation number \\
\hline
Electron transfer rate constant & $1 / k_{\mathrm{r}}=1 / k_{\mathrm{a}}+\left(k_{\mathrm{a}}^{\prime} / k_{\mathrm{a}} k_{\mathrm{et}}\right)\left(1+k_{\mathrm{et}}^{\prime} / k_{\mathrm{d}}\right)$ & Steady-state assumption & 18 E .3 \\
Tunnelling probability & $H_{\mathrm{et}}(d)^{2}=H_{\mathrm{et}}^{\circ 2} \mathrm{e}^{-\beta d}$ & Assumed & 18 E .4 \\
Rate constant & $k_{\mathrm{et}}=(1 / h)\left(\pi^{3} / R T \Delta E_{\mathrm{R}}\right)^{1 / 2} H_{\mathrm{et}}(d)^{2} \mathrm{e}^{-\Delta^{\ddagger} G / R T}$ & Transition-state theory & 18 E .5 \\
Gibbs energy of activation & $\Delta^{\ddagger} G=\left(\Delta_{\mathrm{r}} G^{\ominus}+\Delta E_{\mathrm{R}}\right)^{2} / 4 \Delta E_{\mathrm{R}}$ & Assumes parabolic potential energy & 18 E .6 \\
\hline
\end{tabular}
\end{center}

\chapter{FOCUS 19 Processes at solid surfaces}
A great deal of chemistry occurs at solid surfaces. For example, the surfaces of solid catalysts provide sites where reactants can attach and then undergo chemical transformations. Even as simple an act as dissolving is intrinsically a surface phenomenon, with molecules or ions gradually escaping into the solvent from sites on its surface. Surface deposition, in which atoms are laid down on a surface to create layers, is crucial to the semiconductor industry as it is used in the fabrication of integrated circuits. Finally, in electrochemical processes, electron transfer takes place at the surface of the electrodes.

\subsection*{19A An introduction to solid surfaces}
In many cases the surface of a solid has a different structure from a slice through a bulk solid. It can have various types of imperfections, which turn out to have a profound effect on how atoms and molecules interact with it. This Topic describes the structure of surfaces, the attachment of molecules to them, and some of the techniques used to study them.\\
19A. 1 Surface growth; 19A. 2 Physisorption and chemisorption;\\
19A. 3 Experimental techniques

\subsection*{19B Adsorption and desorption}
A knowledge of the extent to which molecules attach themselves to a surface is crucial to understanding the way in which a surface influences chemical processes. The extent of adsorption can be explored with the aid of some simple models that allow quantitative predictions to be made about how the extent of surface coverage varies with both pressure and temperature.\\
19B. 1 Adsorption isotherms; 19B. 2 The rates of adsorption and desorption

\subsection*{19C Heterogeneous catalysis}
The chemical industry relies on the use of efficient catalysts to facilitate a wide variety of transformations, and the majority of these catalysts involve reactions at surfaces. This Topic describes how the concepts introduced in Topic 19B can be extended to provide a way of modelling surface reactions.\\
19C. 1 Mechanisms of heterogeneous catalysis; 19C. 2 Catalytic activity at surfaces

\subsection*{19D Processes at electrodes}
The key process at an electrode is the transfer of electrons, which takes place at the interface between a solid surface and an electrolyte. The process can be modelled by using a version of transition-state theory and leads to an understanding of electron transfer as an activated process and how it is affected by factors that can be controlled.\\
19D. 1 The electrode-solution interface; 19D. 2 The current density at an electrode; 19D. 3 Voltammetry; 19D. 4 Electrolysis; 19D. 5 Working galvanic cells

\subsection*{Web resources What is an application of this material?}
Almost the whole of modern chemical industry depends on the development, selection, and application of catalysts, with heterogeneous catalysts being particularly important. Impact 27 gives a brief overview of the types of catalysts used and the way in which they are thought to act. The search for efficient, portable or small-scale methods of generating electrical power has led to the development of fuel cells which convert hydrogen or hydrocarbon fuels directly into electrical power. Impact 28 reviews some of the developments in this area.

\section*{TOPIC 19A An introduction to solid surfaces }
\subsection*{Why do you need to know this material?}
To understand the processes occurring on a surface you need to know about its structure, how molecules are attached to it, and how it is studied experimentally.

\subsection*{What is the key idea?}
The attachment of molecules to a surface is influenced by structural features of the surface, including defects.

\subsection*{What do you need to know already?}
You need to be aware of the basic facts about the structures of solids (Topic 15A) and diffraction techniques (Topic 15B). This Topic draws on results from the kinetic theory of gases (Topic 1B).

Many events take place on surfaces. They include the growth of the surface itself as atoms or molecules condense on to it. They also include the attachment of other species, such as molecules from a gas. Adsorption is the attachment of particles to a solid surface; desorption is the reverse process. The substance that adsorbs is the adsorbate and the material to which it adsorbs is the adsorbent or substrate. Adsorption should be distinguished from absorption, the penetration of molecules into the interior of the solid; absorption is often preceded by adsorption.

\subsection*{19A. 1 Surface growth}
A simple picture of a perfect crystal surface is as a tray of oranges in a grocery store (Fig. 19A.1). A gas molecule that collides with the surface can be imagined as a ping-pong ball bouncing erratically over the oranges. The molecule loses energy as it bounces, but it is likely to escape from the surface before it has lost enough kinetic energy to be trapped. The same is true, to some extent, of an ionic crystal in contact with a solution. There is little energy advantage for an ion in solution to discard some of its solvating molecules and stick at an exposed position on the surface.\\
\includegraphics[max width=\textwidth, center]{2025_02_27_d4b95d9718f0b9e2e6b9g-856}

Figure 19A. 1 A schematic diagram of the flat surface of a solid. This primitive model is largely supported by scanning tunnelling microscope images.

The surface of a solid is, however, rarely a flat plane: it is a more rugged landscape, featuring different kinds of defects arising from incomplete layers of atoms or ions. A common type of surface defect is a step between two otherwise flat layers of atoms called terraces (Fig. 19A.2). A step defect might itself have defects, such as kinks. The presence of defects can have a strong influence on the adsorption process. For example, when a molecule strikes a terrace it bounces across it under the influence of the intermolecular potential, and might then come to a step or a corner formed by a kink. Instead of interacting with a single terrace atom, the molecule now interacts with several, and the interaction may be strong enough to trap it. Likewise, when ions deposit from solution, the loss of the solvation interaction is offset by a strong Coulombic\\
\includegraphics[max width=\textwidth, center]{2025_02_27_d4b95d9718f0b9e2e6b9g-856(1)}

Figure 19A. 2 Steps and kinks are two kinds of defects that may occur on otherwise perfect terraces. Defects play an important role in the adsorption of molecules and catalysis.\\
\includegraphics[max width=\textwidth, center]{2025_02_27_d4b95d9718f0b9e2e6b9g-857}

Figure 19A. 3 The slower-growing faces of a crystal dominate its final external appearance. Three successive stages of the growth are shown.\\
interaction between the arriving ions and several ions at the surface defect.

A crystal grows as molecules or ions accumulate on its surface. Different crystal planes grow at different rates, and the slowest growing faces dominate the appearance of the crystal. This feature is explained in Fig. 19A.3, where it can be seen that, although the horizontal face grows forward most rapidly, it grows itself out of existence, and the slower-growing faces survive.

Under normal conditions, a surface exposed to a gas is constantly bombarded with molecules so a freshly prepared surface is covered very quickly. Just how quickly can be estimated by using the kinetic theory of gases to derive an expression for the collision flux $Z_{\mathrm{w}}$ (eqn 16A.7a), which is the number of collisions that occur on a region of surface in a given interval divided by the area of the region and the duration of the interval:


\begin{equation*}
Z_{\mathrm{W}}=\frac{p}{(2 \pi m k T)^{1 / 2}}=\frac{p}{\left(2 \pi M k T / N_{\mathrm{A}}\right)^{1 / 2}} \quad \text { Collision flux } \tag{19A.1}
\end{equation*}


In this expression, $m$ is the mass of the molecule and $M$ is its molar mass.

\subsection*{Brief illustration 19A. 1}
For $\mathrm{N}_{2}$ gas, for which $M=28 \mathrm{~g} \mathrm{~mol}^{-1}$, at $1.0 \mathrm{bar}\left(1.0 \times 10^{5} \mathrm{~Pa}\right)$ and 298 K the collision flux is

$$
\begin{aligned}
Z_{\mathrm{W}} & =\frac{1.0 \times 10^{5} \overbrace{\mathrm{~Pa}}^{\mathrm{kg} \mathrm{~m}^{-1} \mathrm{~s}^{-2}}}{\left(2 \pi \times\left(0.028 \mathrm{~kg} \mathrm{~mol}^{-1} / 6.022 \times 10^{23} \mathrm{~mol}^{-1}\right) \times\left(1.381 \times 10^{-23} \mathrm{JK}^{-1}\right) \times(298 \mathrm{~K})\right)^{1 / 2}} \\
& =2.9 \times 10^{27} \mathrm{~m}^{-2} \mathrm{~s}^{-1}
\end{aligned}
$$

where $1 \mathrm{~J}=1 \mathrm{~kg} \mathrm{~m}^{2} \mathrm{~s}^{-2}$ has been used. Because $1 \mathrm{~m}^{2}$ of metal surface consists of about $10^{19}$ atoms, each atom is struck about $10^{8}$ times each second. Even if only a few collisions result in successful adsorption, the time for which a freshly prepared surface remains clean is very short.

\subsection*{19A. 2 Physisorption and chemisorption}
Molecules and atoms can attach to surfaces in two ways, by 'physisorption' and by 'chemisorption'.

In physisorption (a contraction of 'physical adsorption'), there is a van der Waals interaction (e.g. a dispersion or a dipolar interaction, Topic 14B) between the adsorbate and the substrate. Such interactions have a long range but are weak, and the energy released when a particle is physisorbed is of the same order of magnitude as the enthalpy of condensation. For a molecule to remain bound to the surface, the energy released on binding needs to be dissipated in some way; if it is not lost, the molecule will leave the surface again. The energy involved in physisorption is sufficiently small that it can be dispersed into vibrations of the lattice and dissipated as thermal motion: this process is called accommodation.

The enthalpy of physisorption can be measured by monitoring the rise in temperature of a sample of known heat capacity; typical values are in the region of $-20 \mathrm{~kJ} \mathrm{~mol}^{-1}$ (Table 19A.1). This small enthalpy change is insufficient to lead to bond breaking, so a physisorbed molecule retains its identity, although it might be distorted by the presence of the surface.

In chemisorption (a contraction of 'chemical adsorption'), the molecules (or atoms) attach to the surface by forming a chemical (usually covalent) bond, and tend to find sites that maximize their coordination number with the substrate. The enthalpy of chemisorption is very much greater than that for physisorption, and typical values are in the region of $-200 \mathrm{~kJ} \mathrm{~mol}^{-1}$ (Table 19A.2). The distance between the surface

Table 19A. 1 Maximum observed standard enthalpies of physisorption at $298 \mathrm{~K}^{*}$

\begin{center}
\begin{tabular}{ll}
\hline
Adsorbate & $\Delta_{\mathrm{ad}} \mathrm{H}^{\ominus} /\left(\mathrm{kJ} \mathrm{mol}^{-1}\right)$ \\
\hline
$\mathrm{CH}_{4}$ & -21 \\
$\mathrm{H}_{2}$ & -84 \\
$\mathrm{H}_{2} \mathrm{O}$ & -59 \\
$\mathrm{~N}_{2}$ & -21 \\
\hline
\end{tabular}
\end{center}

\begin{itemize}
  \item More values are given in the Resource section.
\end{itemize}

Table 19A. 2 Standard enthalpies of chemisorption, $\Delta_{\mathrm{ad}} H^{\ominus} /\left(\mathrm{kJ} \mathrm{mol}^{-1}\right)$, at $298 \mathrm{~K}^{*}$

\begin{center}
\begin{tabular}{lccc}
\hline
Adsorbate & \multicolumn{3}{c}{Adsorbent (substrate)} \\
\cline { 2 - 4 }
 & Cr & Fe & Ni \\
\hline
$\mathrm{C}_{2} \mathrm{H}_{4}$ & -427 & -285 & -243 \\
CO &  & -192 &  \\
$\mathrm{H}_{2}$ & -188 & -134 &  \\
$\mathrm{NH}_{3}$ &  & -188 & -155 \\
\hline
\end{tabular}
\end{center}

\footnotetext{\begin{itemize}
  \item More values are given in the Resource section.
\end{itemize}
}
and the closest adsorbate atom is also typically shorter for chemisorption than for physisorption. Bonds within a chemisorbed molecule may be broken due to strong interactions with the surface atoms, resulting in molecular fragments bound to the surface. Such species often play a key role in the catalytic properties of solid surfaces (Topic 19C).\\
Except in special cases, chemisorption must be exothermic. A spontaneous process requires $\Delta G<0$ at constant pressure and temperature. Because the translational freedom of the adsorbate is reduced when it is adsorbed, $\Delta S$ is negative. Therefore, in order for $\Delta G=\Delta H-T \Delta S$ to be negative, $\Delta H$ must be negative (i.e. the process must be exothermic). Exceptions may occur if the adsorbate dissociates and has high translational mobility on the surface. An example is the adsorption of $\mathrm{H}_{2}$ on glass, which is an endothermic process and in which adsorption is believed to be accompanied by dissociation to mobile H atoms.

The enthalpy of adsorption depends on the extent of surface coverage, mainly because the adsorbed species interact with each other. If they repel each other (as for CO on palladium) the adsorption becomes less exothermic (the enthalpy of adsorption less negative) as coverage increases. Moreover, studies show that such species settle on the surface in a disordered way until packing requirements demand order. If the adsorbate species attract one another (as for $\mathrm{O}_{2}$ on tungsten), then they tend to cluster together in islands, and growth occurs at the edges. These adsorbates also show order-disorder transitions when they are heated enough for thermal motion to overcome the interactions between the absorbed species, but not so much that they are desorbed.

The extent of surface coverage (by either physisorption or chemisorption) is normally expressed as the fractional coverage, $\theta$ :

$$
\theta=\frac{\text { number of adsorption sites occupied }}{\text { number of adsorption sites available }}=\frac{N_{\text {occupied }}}{N_{\text {available }}}, \begin{aligned}
& \text { Fractional coverage } \\
& \text { [definition] }
\end{aligned}
$$

(19A.2)\\
The fractional coverage is often calculated from the relation $\theta=$ $V / V_{\infty}$, where $V$ is the volume of adsorbate adsorbed and $V_{\infty}$ is the volume of adsorbate corresponding to complete monolayer coverage. The two volumes are of the free gas measured under the same conditions of temperature and pressure, not the volume that the adsorbed gas occupies when attached to the surface.

\subsection*{Brief illustration 19A. 2}
For the adsorption of CO on charcoal at $273 \mathrm{~K}, V_{\infty}=111 \mathrm{~cm}^{3}$, a value corrected to 1 atm . When the charcoal is exposed to a mixture of gases in which the partial pressure of CO is 80.0 kPa , the value of $V$ (also corrected to 1 atm ) is $41.6 \mathrm{~cm}^{3}$, so it follows that $\theta=\left(41.6 \mathrm{~cm}^{3}\right) /\left(111 \mathrm{~cm}^{3}\right)=0.375$ under these conditions.\\
\includegraphics[max width=\textwidth, center]{2025_02_27_d4b95d9718f0b9e2e6b9g-858}

Figure 19A. 4 Self-assembled monolayers of alkyl thiols formed onto a gold surface by reaction of the thiol groups with the surface and aggregation of the alkyl chains.

Chemisorption can be used as the basis for manipulation of surfaces on the nanometre scale. Of current interest are selfassembled monolayers (SAMs), ordered molecular aggregates that form a single layer of material on a surface. To understand the formation of SAMs, consider exposing molecules such as alkyl thiols, RSH, where R represents an alkyl chain, to an $\mathrm{Au}(0)$ surface. The thiols react with the surface, forming $\mathrm{RS}^{-}$ $\mathrm{Au}(\mathrm{I})$ adducts:

$$
\mathrm{RSH}+\mathrm{Au}(0)_{n} \rightarrow \mathrm{RS}^{-} \mathrm{Au}(\mathrm{I}) \cdot \mathrm{Au}(0)_{n-1}+\frac{1}{2} \mathrm{H}_{2}
$$

If $R$ is a sufficiently long chain, van der Waals interactions between the adsorbed RS units lead to the formation of a highly ordered monolayer on the surface (Fig. 19A.4). A selfassembled monolayer alters the properties of the surface. For example, a hydrophilic surface may be rendered hydrophobic once covered with a SAM. Furthermore, the attachment of functional groups to the exposed ends of the alkyl groups may impart specific chemical reactivity or ligand-binding properties to the surface, leading to applications in chemical (or biochemical) sensors and reactors.

\subsection*{19A.3 Experimental techniques}
Many experimental techniques are available for the study of the structure of the solid surface as well as the properties and arrangement of any adsorbed molecules. Some of these techniques offer atomic level resolution and allow the direct visualization of changes to the surface as adsorption and chemical reactions take place.

Experimental procedures must begin with a clean surface. The obvious way to retain cleanliness of a surface is to reduce the pressure and thereby reduce the number of impacts on the surface. When the pressure is reduced to 0.1 mPa (as in a simple vacuum system) the collision flux falls to about $10^{18} \mathrm{~m}^{-2} \mathrm{~s}^{-1}$, corresponding to one hit on a surface atom in each 0.1 s . Even\\
that is too frequent in most experiments, and in ultrahigh vacuum (UHV) techniques pressures as low as $0.1 \mu \mathrm{~Pa}$ (when $Z_{\mathrm{W}} \approx 10^{15} \mathrm{~m}^{-2} \mathrm{~s}^{-1}$ ) are reached on a routine basis and as low as 1 nPa (when $Z_{\mathrm{w}} \approx 10^{13} \mathrm{~m}^{-2} \mathrm{~s}^{-1}$ ) are reached with special care. These collision fluxes correspond to each surface atom being hit once every $10^{5}-10^{6} \mathrm{~s}$, or about once a day.

\subsection*{(a) Microscopy}
The basic approach of illuminating a small area of a sample and collecting light with a microscope has been used for many years to image small specimens. However, the resolution of a microscope, the minimum distance between two objects that leads to two distinct images, is on the order of the wavelength of the light being used. Therefore, conventional microscopes employing visible light have resolutions of the order of micrometres and are blind to features on a scale of nanometres.

One technique often used to image nanometre-sized objects is electron microscopy, in which a beam of electrons with a well-defined de Broglie wavelength (Topic 7A) replaces the light source found in traditional microscopes. Instead of glass or quartz lenses, magnetic fields are used to focus the beam. In transmission electron microscopy (TEM), the electron beam passes through the specimen and the image is collected on a screen. In scanning electron microscopy (SEM), electrons scattered back from a small area of the sample are detected and an image of the surface is then obtained by scanning the electron beam across the sample.

As in traditional light microscopy, the wavelength of the incident beam and the ability to focus it governs the resolution. It is now possible to achieve atomic resolution with TEM instruments, and SEM instruments can achieve resolution on the order of a few nanometres.

Scanning probe microscopy (SPM) is a collection of techniques that can be used to image and manipulate objects as small as atoms on surfaces. One version is scanning tunnelling microscopy (STM), in which a platinum-rhodium or tungsten needle is scanned across the surface of a conducting solid. When the tip of the needle is brought very close to the surface, electrons tunnel across the intervening space (Fig. 19A.5). In the 'constant-current mode' of operation, the tip moves up and down according to the form of the surface, and the topography of the surface, including any adsorbates, can be mapped on an atomic scale. The vertical motion of the tip is achieved by fixing it to a piezoelectric cylinder, which contracts or expands according to the potential difference it experiences. In the 'constant- $z$ mode', the vertical position of the tip is held constant and the current is monitored. Because the tunnelling probability is very sensitive to the size of the gap, the microscope can detect tiny, atom-scale variations in the height of the surface.

Figure 19A. 6 shows an example of the kind of image obtained from a surface, in this case that of copper atoms\\
\includegraphics[max width=\textwidth, center]{2025_02_27_d4b95d9718f0b9e2e6b9g-859}

Figure 19A. 5 A scanning tunnelling microscope makes use of the current due to electrons that tunnel between the surface and the tip. That current is very sensitive to the distance of the tip above the surface.\\
forming 'runways' on a surface. Each 'bump' on the surface corresponds to a single copper atom. In a further variation of the STM technique, the tip may be used to nudge single atoms around on the surface, making possible the fabrication of complex nanometre-sized materials and devices.

The diffusion characteristics of an adsorbate can be examined by using STM to follow the change in surface characteristics. Provided its path is not influenced by defects, an adsorbed atom makes a random walk across the surface, and the diffusion coefficient, $D$, is inferred from the mean distance, $d$, travelled in an interval $\tau$ by using the two-dimensional random walk expression $d=(D \tau)^{1 / 2}$ (which is derived like the onedimensional random walk expression in Topic 16C). Values of $D$ at different temperatures are interpreted by using an Arrhenius-like expression

\[
D=D_{0} \mathrm{e}^{-E_{\mathrm{a}, \mathrm{difif}} / R T} \quad \begin{align*}
& \text { Temperature dependence }  \tag{19A.3}\\
& \text { of the diffusion coefficient }
\end{align*}
\]

where $E_{\mathrm{a}, \mathrm{diff}}$ is the activation energy for diffusion and $D_{0}$ is the diffusion coefficient in the limit of infinite temperature. The variation of $D$ from one crystal plane to another can also be studied.\\
\includegraphics[max width=\textwidth, center]{2025_02_27_d4b95d9718f0b9e2e6b9g-859(1)}

Figure 19A.6 An STM image of copper atoms forming 'runways' on the surface of metallic copper when a layer of CuO is formed on the surface. (Image provided by Stephen Driver and Stephen Jenkins, University of Cambridge.)

\subsection*{Brief illustration 19A. 3}
For tungsten atoms on a tungsten surface it is found that $E_{\text {a,diff }}$ is in the range $57-87 \mathrm{~kJ} \mathrm{~mol}^{-1}$ and $D_{0} \approx 3.8 \times 10^{-11} \mathrm{~m}^{2} \mathrm{~s}^{-1}$. It follows from eqn 19A. 3 that at 800 K this range of activation energies corresponds to a value of the diffusion coefficient between

$$
\begin{aligned}
D & =\left(3.8 \times 10^{-11} \mathrm{~m}^{2} \mathrm{~s}^{-1}\right) \times \mathrm{e}^{-5.7 \times 10^{4} \mathrm{Jmol}^{-1} /\left(\left(8.3145 \mathrm{JK} \mathrm{~K}^{-1} \mathrm{~mol}^{-1}\right) \times(80 \mathrm{~K})\right)} \\
& =7.2 \times 10^{-15} \mathrm{~m}^{2} \mathrm{~s}^{-1}
\end{aligned}
$$

and

$$
\begin{aligned}
D & =\left(3.8 \times 10^{-11} \mathrm{~m}^{2} \mathrm{~s}^{-1}\right) \times \mathrm{e}^{-8.7 \times 10^{4} \mathrm{Jmol}}{ }^{-1 /\left(\left(8.3145 \mathrm{~K} \mathrm{~K}^{-1} \mathrm{~mol}^{-1}\right) \times(800 \mathrm{~K})\right)} \\
& =7.9 \times 10^{-17} \mathrm{~m}^{2} \mathrm{~s}^{-1}
\end{aligned}
$$

In atomic force microscopy (AFM), a sharpened tip attached to a cantilever is scanned across the surface. The force exerted by the surface and any molecules attached to it pushes or pulls on the tip and deflects the cantilever (Fig. 19A.7). The deflection is monitored by using a laser beam. Because no current needs to pass between the sample and the probe, the technique can be applied to non-conducting surfaces and to the study of solid-liquid interfaces.\\
Two modes of operation of AFM are common. In 'contact mode', or 'constant-force mode', the force between the tip and surface is held constant and the tip makes contact with the surface. This mode of operation can damage fragile samples on the surface. In 'non-contact', or 'tapping' mode, the tip bounces up and down with a specified frequency and never quite touches the surface. The amplitude of the oscillation of the tip changes when it passes over a species adsorbed on the surface.

\subsection*{(b) Ionization techniques}
The chemical composition of a surface can be determined by a variety of ionization techniques. The same techniques can be used to detect any remaining contamination after cleaning and to detect layers of material adsorbed later in the experiment.\\
\includegraphics[max width=\textwidth, center]{2025_02_27_d4b95d9718f0b9e2e6b9g-860}

Figure 19A. 7 In atomic force microscopy, a laser beam is used to monitor the tiny changes in position of a probe as it is attracted to or repelled by atoms on a surface.\\
\includegraphics[max width=\textwidth, center]{2025_02_27_d4b95d9718f0b9e2e6b9g-860(1)}

Figure 19A. 8 The X-ray photoelectron emission spectrum of a sample of gold contaminated with a surface layer of mercury. (M.W. Roberts and C.S. McKee, Chemistry of the metal-gas interface, Oxford (1978).)

One technique is photoemission spectroscopy, the origins of which lie in the photoelectric effect (Topic 7A). In this technique the sample is irradiated with photons of sufficient energy to eject electrons from the absorbed species, and the energies of these electrons are measured. If the ionizing radiation is in the ultraviolet (when the technique is denoted UPS), the ejected electrons are from the valence shells and so the technique can be used to infer details of the bonding between the adsorbate and the substrate. If X-rays are used (the technique is then denoted XPS), core electrons are ionized and their energies are characteristic of the atoms present (Fig. 19A.8), thus providing a fingerprint for the material present.

\subsection*{Brief illustration 19A. 4}
The principal difference between the UPS of free benzene and benzene adsorbed on palladium is in the energies of the $\pi$ electrons. This difference is interpreted as meaning that the benzene molecules interact with the surface through their $\pi$ orbitals, implying that the molecules lie parallel to the surface.

A very important technique, which is widely used in the microelectronics industry, is Auger electron spectroscopy (AES). The Auger effect (pronounced oh-zhey) is the emission of a second electron after high energy radiation has expelled another. The first electron to depart leaves a hole in a low-lying orbital, and then an electron in a higher energy orbital falls into it. The energy this transition releases may result either in the generation of radiation, which is called X-ray fluorescence (Fig. 19A.9a) or in the ejection of another electron (Fig. 19A.9b). The latter is the 'secondary electron' of the Auger effect. The energies of the secondary electrons are characteristic of the material present, so AES effectively provides a fingerprint of the sample. In practice, the Auger spectrum is normally obtained by using an electron beam (with energy in the range\\
\includegraphics[max width=\textwidth, center]{2025_02_27_d4b95d9718f0b9e2e6b9g-861(2)}

Figure 19A. 9 When an electron is expelled from a solid (a) an electron of higher energy may fall into the vacated orbital and emit an X-ray photon (X-ray fluorescence). Alternatively (b) the electron falling into the orbital may give up its energy to another electron, which is ejected as the secondary electron in the Auger effect.\\
$1-5 \mathrm{keV}$ ), rather than electromagnetic radiation, to eject the primary electron. In scanning Auger electron microscopy (SAM), the finely focused electron beam is scanned over the surface and a map of composition is compiled; the resolution can be less than about 50 nm .

\subsection*{(c) Diffraction techniques}
A technique for determining the arrangement of the atoms close to the surface is low energy electron diffraction (LEED). This technique is similar to X-ray diffraction (Topic 15B), but makes use of the wavelike properties of electrons. The use of low energy electrons (with energies in the range $10-200 \mathrm{eV}$, corresponding to wavelengths in the range $400-100 \mathrm{pm}$ ) ensures that the diffraction is caused only by atoms at or close to the surface. The experimental arrangement is shown in Fig. 19A.10, and typical LEED patterns, obtained by photographing the fluorescent screen through the viewing port, are shown in Fig. 19A.11.

Observations using LEED show that the surface of a crystal rarely has exactly the same form as a slice through the bulk\\
\includegraphics[max width=\textwidth, center]{2025_02_27_d4b95d9718f0b9e2e6b9g-861(1)}

Figure 19A.10 A schematic diagram of the apparatus used for a LEED experiment. The electrons diffracted by the surface layers are detected by the fluorescence they cause on the phosphor screen.\\
\includegraphics[max width=\textwidth, center]{2025_02_27_d4b95d9718f0b9e2e6b9g-861(3)}

Figure 19A. 11 LEED photographs of (a) a clean surface of $\mathrm{FeS}_{2}$ and (b) after its exposure to Mo atoms; it is thought that $\mathrm{MoS}_{2}$ forms on the surface. The black area is the shadow from the electron gun. (Photographs provided by Tao Liu, Israel Temprano, David King, Stephen Driver, and Stephen Jenkins, University of Cambridge.)\\
because surface and bulk atoms experience different forces. Reconstruction refers to processes by which atoms on the surface achieve their equilibrium structures. As a general rule, it is found that metal surfaces are simply truncations of the bulk lattice, but the distance between the top layer of atoms and the one below is contracted by around 5 per cent. Semiconductors generally have surfaces reconstructed to a depth of several layers. Reconstruction also occurs in ionic solids. For example, in lithium fluoride the $\mathrm{Li}^{+}$and $\mathrm{F}^{-}$ions close to the surface apparently lie on slightly different planes. An actual example of the detail that can now be obtained from refined LEED techniques is shown in Fig. 19A. 12 for $\mathrm{CH}_{3} \mathrm{C}$ - adsorbed on a (110) plane of rhodium.\\
\includegraphics[max width=\textwidth, center]{2025_02_27_d4b95d9718f0b9e2e6b9g-861}

Figure 19A. 12 The structure of a surface close to the point of attachment of $\mathrm{CH}_{3} \mathrm{C}$ - to the (110) surface of rhodium at 300 K and the changes in positions of the metal atoms that accompany chemisorption.

\subsection*{Example 19A. 1 Interpreting a LEED pattern}
The LEED pattern from a clean (110) face of palladium is shown in (a) below. The reconstructed surface gives a LEED pattern shown as (b). What can be inferred about the structure of the reconstructed surface?\\
\includegraphics[max width=\textwidth, center]{2025_02_27_d4b95d9718f0b9e2e6b9g-862}

Collect your thoughts Recall from Bragg's law (Topic 15B, $\lambda=2 d \sin \theta$ ), that for a given wavelength, the greater the separation $d$ of the layers, the smaller is the scattering angle (so that $2 d \sin \theta$ remains constant). It follows that, in terms of the LEED pattern, the farther apart the atoms responsible for the pattern, the closer the spots appear in the pattern. Twice the separation between the atoms corresponds to half the separation between the spots, and vice versa. Therefore, by comparing the two patterns you can decide if there has been any change in the spacing between the atoms.

The solution The horizontal separation between spots is unchanged, which indicates that the atoms remain in the same position in that dimension when reconstruction occurs. However, the vertical spacing is halved, which suggests that the atoms are twice as far apart in that direction as they are in the unreconstructed surface.

Self-test 19A. 1 Sketch the LEED pattern for a surface that differs from that shown in (a) above by tripling the vertical separation.\\
\includegraphics[max width=\textwidth, center]{2025_02_27_d4b95d9718f0b9e2e6b9g-862(1)}

The presence of terraces, steps, and kinks in a surface shows up in LEED patterns, and from such experiments the surface defect density (the number of defects in a region divided by the area of the region) can be estimated. The importance of this type of measurement will emerge later.

\subsection*{Brief illustration 19A. 5}
Three examples of how steps and kinks affect LEED patterns are shown in Fig. 19A.13. The samples were obtained by cleaving a crystal at different angles to a plane of atoms. Only terraces are produced when the cut is parallel to the plane, and the density of steps increases as the angle of the cut increases. The observation of additional structure in the LEED patterns, rather than blurring, shows that the steps are arrayed regularly.\\
\includegraphics[max width=\textwidth, center]{2025_02_27_d4b95d9718f0b9e2e6b9g-862(2)}

Figure 19A. 13 LEED patterns may be used to assess the defect density of a surface. The photographs correspond to a platinum surface with (a) low defect density, (b) regular steps separated by about six atoms, and (c) regular steps with kinks. (Photographs provided by Professor G.A. Samorjai.)

\subsection*{(d) Determination of the extent and rates of adsorption and desorption}
A common technique for measuring rates of processes on surfaces is to monitor the rates of flow of gas into and out of the system: the difference is the rate of gas uptake by the sample. Integration of this rate then gives the fractional coverage at any stage. Three other techniques for the determination of fractional coverage and the rate of adsorption are:

\begin{itemize}
  \item Gravimetry, in which the sample is weighed on a microbalance during the experiment.
\end{itemize}

The technique commonly uses a quartz crystal microbalance (QCM), in which the mass of a sample adsorbed on the surface of a quartz crystal is related to changes in the characteristic vibrational frequency of the crystal. Masses as small as a few nanograms can be measured reliably in this way.

\begin{itemize}
  \item Second harmonic generation (SHG), the conversion of an intense, pulsed laser beam to radiation with twice its initial frequency.
\end{itemize}

For example, adsorption of gas molecules on to a surface alters the intensity of the SHG signal. Because pulsed lasers are the excitation sources, time-resolved measurements of surface processes are possible on timescales as short as femtoseconds.

\begin{itemize}
  \item Surface plasmon resonance (SPR), the absorption of energy from an incident beam of electromagnetic radiation by surface 'plasmons'.
\end{itemize}

This technique is very sensitive and is now used routinely in the study of adsorption and desorption. To understand it, it is necessary to understand what is meant by a 'surface plasmon' and what kind of 'resonance' is involved.

The mobile delocalized valence electrons of metals form a plasma, a dense gas-like collection of charged particles. Bombardment of this plasma by light or an electron beam can cause transient changes in the distribution of electrons, with some regions becoming slightly denser than others. Coulombic repulsion in the regions of high density causes electrons to move away from each other, so lowering their density. The resulting oscillations in electron density, the plasmons, can be excited both in the bulk and on the surface of a metal. A surface plasmon propagates away from the surface, but the amplitude of the wave, also called an evanescent wave, decreases sharply with distance from the surface. The term 'resonance' in this context refers to the absorption that can be observed with appropriate choice of the wavelength and angle of incidence of the excitation beam.

To detect surface plasmon resonance it is common practice to use a monochromatic beam and to vary the angle of incidence (the $\phi$ in Fig. 19A.14). The beam passes into a prism and reflects off a face which has been coated with a thin film of gold or silver. Because the evanescent wave interacts with material a short distance away from the surface, the angle at which resonant absorption occurs depends on the refractive index of the medium on the opposite side of the metallic film. Thus, changing the identity and quantity of material on the surface changes the resonance angle.\\
\includegraphics[max width=\textwidth, center]{2025_02_27_d4b95d9718f0b9e2e6b9g-863}

Figure 19A. 14 The experimental arrangement for the observation of surface plasmon resonance (SPR), as explained in the text. Here receptors have been attached to the metallic film; the binding of ligands (red spheres) alters the angle at which SPR is detected.

The SPR technique can be used in the study of the binding of molecules to a surface or the binding of ligands to a biopolymer attached to the surface; this interaction mimics the biological recognition processes that occur in cells. Examples of complexes amenable to analysis include antibody-antigen and protein-DNA interactions. The most important advantage of SPR is its sensitivity: it is possible to measure the deposition of nanograms of material on to a surface. The main disadvantage of the technique is its requirement for immobilization of at least one of the components of the system under study.

\subsection*{Checklist of concepts}
$\square$ 1. Adsorption is the attachment of molecules to a surface; the substance that adsorbs is the adsorbate and the underlying material is the adsorbent or substrate. The reverse of adsorption is desorption.\\
$\square$ 2. Surface defects play an important role in the process of adsorption.\\
3. In physisorption molecules are attached to the surface by relatively weak forces, such as van der Waals interactions; in chemisorption the interactions are much stronger, and chemical bonds are formed between adsorbate and substrate.4. Reconstruction refers to processes by which atoms on the surface achieve their equilibrium structures.\\
5. Techniques for studying surfaces and their acronyms are:

\begin{center}
\begin{tabular}{ll}
AES & Auger electron spectroscopy \\
AFM & atomic force microscopy \\
LEED & low energy electron diffraction \\
QCM & quartz crystal microbalance \\
SEM & scanning electron microscopy \\
SHG & second harmonic generation \\
SPM & scanning probe microscopy \\
SPR & surface plasmon resonance \\
STM & tranning tunnelling microscopy \\
TEM & ultraviolet photoelectron spectroscopy \\
UPS & X-ray photoelectron spectroscopy \\
XPS &  \\
\end{tabular}
\end{center}

\subsection*{Checklist of equations}
\begin{center}
\begin{tabular}{llll}
\hline
Property & Equation & Comment & Equation number \\
\hline
Collision flux & $Z_{\mathrm{w}}=p /(2 \pi m k T)^{1 / 2}$ & Kinetic theory & 19 A .1 \\
Fractional coverage & $\theta=N_{\text {occupied }} / N_{\text {available }}$ & Definition & 19 A .2 \\
\hline
\end{tabular}
\end{center}

\section*{TOPIC 19B Adsorption and desorption}
\subsection*{Why do you need to know this material?}
To understand how surfaces can affect the rates of chemical reactions, you need to know how to assess the extent of surface coverage and the factors that determine the rates at which molecules attach to and detach from solid surfaces.

\subsection*{What is the key idea?}
The extent of surface coverage can be expressed in terms of isotherms derived on the basis of dynamic equilibria between adsorbed and free molecules.

\subsection*{What do you need to know already?}
This Topic extends the discussion of adsorption in Topic 19A. You need to be familiar with the basic ideas of chemical kinetics (Topics 17A-17C) and the Arrhenius equation (Topic 17D). One argument makes use of the relation between an equilibrium constant and the standard Gibbs energy of reaction (Topic 6A) and also of the GibbsHelmholtz equation (Topic 3E).

When a gas is adsorbed on a surface (Topic 19A) there is a dynamic equilibrium between the free and the adsorbed molecules. The fractional coverage, $\theta$, of the surface (eqn 19A.2) depends on the pressure of the overlying gas and the temperature; the expression describing its variation with pressure at a chosen temperature is called the adsorption isotherm.

\subsection*{19B. 1 Adsorption isotherms}
Many of the techniques discussed in Topic 19A can be used to measure $\theta$. Another is flash desorption, in which the sample is suddenly heated electrically and the resulting rise of pressure is interpreted in terms of the amount of adsorbate originally on the sample.

\subsection*{(a) The Langmuir isotherm}
An isotherm devised by Irving Langmuir is the simplest that is physically plausible. It is based on four assumptions:

\begin{itemize}
  \item Adsorption cannot proceed beyond monolayer coverage.
  \item All sites on the surface are equivalent.
  \item A molecule can be adsorbed only at a vacant site.
  \item The probability of adsorption is independent of the occupation of neighbouring sites (that is, there are no interactions between adsorbed molecules).
\end{itemize}

From these assumptions it is possible to develop an expression for the dependence of the fractional coverage on the pressure.

\subsection*{How is that done? 19B. 1}
Deriving the Langmuir isotherm\\
You need to consider the dynamic equilibrium between the molecules (A) in the gas phase and those on the surface (denoted AM):

$$
\mathrm{A}(\mathrm{~g})+\mathrm{M}(\text { surface }) \underset{k_{\mathrm{d}}}{\stackrel{k_{\mathrm{s}}}{\leftrightarrows}} \mathrm{AM}(\text { surface })
$$

\subsection*{Step 1 Write an expression for the rate of adsorption}
The rate of adsorption is proportional to the rate of collisions with the surface and therefore to the partial pressure $p$ of A. The rate is also proportional to the number of vacant sites, because molecules can be adsorbed only at these sites. If the total number of sites is $N$ and the fractional coverage is $\theta$, the number of vacant sites is $N(1-\theta)$. The rate of change of surface coverage, $\mathrm{d} \theta / \mathrm{d} t$, due to adsorption is therefore


\begin{equation*}
\frac{\mathrm{d} \theta}{\mathrm{~d} t}=k_{\mathrm{a}} p N(1-\theta) \tag{19B.1a}
\end{equation*}


Rate of adsorption

Step 2 Write an expression for the rate of desorption\\
The rate of change of fractional coverage due to desorption is proportional to the number of adsorbed species already present, which is equal to the number of occupied sites, $N \theta$ :


\begin{equation*}
\frac{\mathrm{d} \theta}{\mathrm{~d} t}=-k_{\mathrm{d}} N \theta \quad \text { Rate of desorption } \tag{19B.1b}
\end{equation*}


The term is negative because $\theta$ decreases as the molecules desorb.

\subsection*{Step 3 Equate the rates and construct the isotherm}
At equilibrium there is no net change in $\theta$, implying that the sum of these two rates must be zero: $k_{a} p N(1-\theta)-k_{d} N \theta=0$. Rearranging this equation gives the following expression, the Langmuir isotherm, which relates the surface coverage to the\\
pressure, and in which the parameter $\alpha$ has the dimensions of $1 /$ pressure:\\
$\jmath \theta=\frac{\alpha p}{1+\alpha p}$\\
$\alpha=\frac{k_{\mathrm{a}}}{k_{\mathrm{d}}}$

The Langmuir isotherm is tested by measuring the surface coverage as a function of the pressure, and then plotting these data in a form expected to give a straight line, as illustrated in the following Example.

\subsection*{Example 19B. 1 Using the Langmuir isotherm}
The following data are for the adsorption of CO on charcoal at 273 K. Confirm that they conform to the Langmuir isotherm, find the value of the parameter $\alpha$ and the volume corresponding to complete coverage. In each case $V$ has been corrected to $1 \mathrm{~atm}(101.325 \mathrm{kPa})$.

\begin{center}
\begin{tabular}{llllllll}
$p / \mathrm{kPa}$ & 13.3 & 26.7 & 40.0 & 53.3 & 66.7 & 80.0 & 93.3 \\
$V / \mathrm{cm}^{3}$ & 10.2 & 18.6 & 25.5 & 31.5 & 36.9 & 41.6 & 46.1 \\
\end{tabular}
\end{center}

Collect your thoughts The fractional coverage is given by $\theta=V / V_{\infty}$, where $V_{\infty}$ is the volume corresponding to complete coverage (eqn 19A.2). You need to manipulate the Langmuir isotherm so that you can plot a straight-line graph and extract the required parameters from its slope and intercept.

The solution Multiply both sides of eqn 19B. 2 by $(1+\alpha p)$ to give $\theta(1+\alpha p)=\alpha p$, and then substitute $\theta=V / V_{\infty}$ to give

$$
\frac{V}{V_{\infty}}+\frac{V \alpha p}{V_{\infty}}=\alpha p
$$

Division of both sides by $V \alpha$ gives

$$
\frac{1}{\alpha V_{\infty}}+\frac{p}{V_{\infty}}=\frac{p}{V} \text { which rearranges to } \quad \frac{p}{V}=\overbrace{\frac{1}{\alpha V_{\infty}}}^{\text {Intercept }}+\overbrace{\frac{1}{V_{\infty}} p}^{\text {Slope }}
$$

It follows that you should plot $p / V$ against $p$ and expect a straight line of slope $1 / V_{\infty}$ and intercept $1 / \alpha V_{\infty}$ at $p=0$; note that slope/intercept $=\left(1 / V_{\infty}\right) /\left(1 / \alpha V_{\infty}\right)=\alpha$.

The data for the plot are as follows:

\begin{center}
\begin{tabular}{lccccccc}
$p / \mathrm{kPa}$ & 13.3 & 26.7 & 40.0 & 53.3 & 66.7 & 80.0 & 93.3 \\
$(p / \mathrm{kPa}) /\left(\mathrm{V} / \mathrm{cm}^{3}\right)$ & 1.30 & 1.44 & 1.57 & 1.69 & 1.81 & 1.92 & 2.02 \\
\end{tabular}
\end{center}

The points are plotted in Fig. 19B.1. The (least squares) slope is $9.04 \times 10^{-3}$, so $V_{\infty}=1 /\left(9.04 \times 10^{-3} \mathrm{~cm}^{-3}\right)=111 \mathrm{~cm}^{3}$. The intercept at $p=0$ is $(p / \mathrm{kPa}) /\left(V / \mathrm{cm}^{3}\right)=1.20$, or $p / V=1.20 \mathrm{kPa} \mathrm{cm}^{-3}$ hence $1 / \alpha V_{\infty}=1.20 \mathrm{kPa} \mathrm{cm}^{-3}$. Therefore

$$
\alpha=\frac{1 / V_{\infty}}{1 / \alpha V_{\infty}}=\frac{1 / 111 \mathrm{~cm}^{3}}{1.20 \mathrm{kPacm}^{-3}}=7.51 \times 10^{-3} \mathrm{kPa}^{-1}
$$

\begin{center}
\includegraphics[max width=\textwidth]{2025_02_27_d4b95d9718f0b9e2e6b9g-865(1)}
\end{center}

Figure 19B. 1 The plot of the data in Example 19B.1. As illustrated here, the Langmuir isotherm predicts that a straight line should be obtained when $p / V$ is plotted against $p$.

Self-test 19B. 1 Repeat the calculation for the following data:

\begin{center}
\begin{tabular}{|c|c|c|c|c|c|c|c|}
\hline
$p / \mathrm{kPa}$ & 13.3 & 26.7 & 40.0 & 53.3 & 66.7 & 80.0 & 93.3 \\
\hline
\multirow[t]{2}{*}{$\mathrm{V} / \mathrm{cm}^{3}$} & 10.3 & 19.3 & 27.3 & 34.1 & 40.0 & 45.5 & 48.0 \\
\hline
 &  &  &  & \multicolumn{4}{|l|}{\includegraphics[max width=\textwidth]{2025_02_27_d4b95d9718f0b9e2e6b9g-865}
} \\
\hline
\end{tabular}
\end{center}

If a molecule $A_{2}$ adsorbs with dissociation to give two $A$ fragments on the surface, the rate of adsorption is proportional to the pressure and to the square of the number of vacant sites: two sites are needed to accommodate the two A. The rate of change of the fractional coverage due to adsorption is then


\begin{equation*}
\frac{\mathrm{d} \theta}{\mathrm{~d} t}=k_{\mathrm{a}} p\{N(1-\theta)\}^{2} \tag{19B.3a}
\end{equation*}


Desorption requires that two of the species A encounter one another so they can leave as $\mathrm{A}_{2}$. The rate of change of the fractional coverage is therefore second-order in the number of sites occupied:


\begin{equation*}
\frac{\mathrm{d} \theta}{\mathrm{~d} t}=-k_{\mathrm{d}}(N \theta)^{2} \tag{19B.3b}
\end{equation*}


The condition for no net change, which means setting the rates in eqns 19B.3a and 19B.3b equal to each other, leads to the isotherm

\[
\theta=\frac{(\alpha p)^{1 / 2}}{1+(\alpha p)^{1 / 2}} \quad \begin{align*}
& \text { Langmuir isotherm for }  \tag{19B.4}\\
& \text { adsorption with dissociation }
\end{align*}
\]

The surface coverage depends more weakly on pressure than it does for non-dissociative adsorption.

The shapes of the Langmuir isotherms with and without dissociation are shown in Figs. 19B. 2 and 19B.3. The fractional coverage increases with increasing pressure, and approaches 1 only at very high pressure when the gas is forced on to every available site of the surface.\\
\includegraphics[max width=\textwidth, center]{2025_02_27_d4b95d9718f0b9e2e6b9g-866}

Figure 19B. 2 The Langmuir isotherm for non-dissociative adsorption for different values of $\alpha$.\\
\includegraphics[max width=\textwidth, center]{2025_02_27_d4b95d9718f0b9e2e6b9g-866(1)}

Figure 19B. 3 The Langmuir isotherm for dissociative adsorption, $\mathrm{A}_{2}(\mathrm{~g}) \rightarrow 2 \mathrm{~A}$ (surface), for different values of $\alpha$.

\subsection*{(b) The isosteric enthalpy of adsorption}
The Langmuir isotherm depends on the value of $\alpha=k_{a} / k_{d}$, which in turn depends on the temperature. It is possible to relate this temperature dependence to the isosteric enthalpy of adsorption, $\Delta_{\mathrm{ad}} H^{\ominus}$, which is the standard enthalpy of adsorption at a fixed surface coverage.

\subsection*{How is that done? 19B. 2 \\
 Relating the temperature dependence of $\alpha$ to the isosteric enthalpy of adsorption}
The quantity $\alpha=k_{\mathrm{a}} / k_{\mathrm{d}}$ is the ratio of the rate constants for the forward and reverse reactions in the equilibrium $\mathrm{A}(\mathrm{g})+$ M (surface) $\rightleftharpoons \mathrm{AM}$ (surface). It follows from the discussion in Topic 17C that $\alpha$ is related to the equilibrium constant for this reaction, and so its temperature dependence can be developed in the same way as for any other equilibrium constant (Topic 6B).

Step 1 Relate $\alpha$ to the equilibrium constant\\
Because the dimensions of $\alpha$ are those of $1 /$ pressure, its relation to the dimensionless equilibrium constant is $K=\left(k_{\mathrm{a}} / k_{\mathrm{d}}\right) \times p^{\ominus}$ $=\alpha p^{\ominus}$.

Step 2 Relate the equilibrium constant to the standard Gibbs energy of adsorption\\
From eqn 6A.15 ( $\left.\Delta_{\mathrm{r}} G^{\ominus}=-R T \ln K\right)$ it follows that $\Delta_{\mathrm{ad}} G^{\ominus}=$ $-R T \ln \left(\alpha p^{\ominus}\right)$, where $\Delta_{\mathrm{ad}} G^{\ominus}$ is the standard Gibbs energy of adsorption. This expression can be rearranged to

$$
-R \ln \left(\alpha p^{\ominus}\right)=\frac{\Delta_{\mathrm{ad}} G^{\ominus}}{T}
$$

Step 3 Use the Gibbs-Helmholtz equation to relate the temperature dependence of $\Delta G^{\ominus} / T$ to the enthalpy of adsorption The derivative with respect to $T$ of the last expression is

$$
\frac{\mathrm{d}\left(-R \ln \left(\alpha p^{\ominus}\right)\right)}{\mathrm{d} T}=\frac{\mathrm{d}}{\mathrm{~d} T} \frac{\Delta_{\mathrm{ad}} G^{\ominus}}{T}
$$

Now use the Gibbs-Helmholtz equation (eqn 3E.11, $\mathrm{d}(\Delta G / T) / \mathrm{d} T$ $=-\Delta H / T^{2}$ ) to write the right-hand side of this equation as $-\Delta_{\mathrm{ad}} H^{\ominus} / T^{2}$, and therefore obtain

$$
\frac{\mathrm{d}\left(-R \ln \left(\alpha p^{\ominus}\right)\right)}{\mathrm{d} T}=\frac{-\Delta_{\mathrm{ad}} H^{\ominus}}{T^{2}} \text { hence } \frac{\mathrm{d} \ln \left(\alpha p^{\ominus}\right)}{\mathrm{d} T}=\frac{\Delta_{\mathrm{ad}} H^{\ominus}}{R T^{2}}
$$

There is a possibility that the standard enthalpy of adsorption depends on the fractional coverage, so this expression is restricted to constant $\theta$. The derivative is therefore a partial derivative evaluated at constant $\theta$ and $\Delta_{\mathrm{ad}} H^{\ominus}$ must be interpreted as the isosteric enthalpy of adsorption. The final result, therefore, is an expression for obtaining this quantity from the temperature dependence of $\alpha$ :


\begin{equation*}
\left(\frac{\partial \ln \left(\alpha p^{\ominus}\right)}{\partial T}\right)_{\theta}=\frac{\Delta_{\mathrm{ad}} H^{\ominus}}{R T^{2}} \tag{19B.5a}
\end{equation*}


This expression can be cast into a more useful form by using $\mathrm{d}(1 / T) / \mathrm{d} T=-1 / T^{2}$ to rewrite is as\\
$\left.-\left(\frac{\partial \ln \left(\alpha p^{\ominus}\right)}{\partial(1 / T)}\right)_{\theta}=-\frac{\Delta_{\text {ad }} H^{\ominus}}{R} \right\rvert\, \begin{aligned} & \text { Isosteric enthalpy of adsorption }\end{aligned}$

The following Example shows how eqn 19B.5b leads to a graphical method for determining a value of the isosteric enthalpy of adsorption.

\subsection*{Example 19B. 2 \\
 Measuring the isosteric enthalpy of adsorption}
The following data show the pressures of CO needed for the volume of adsorbed gas (corrected to 1 atm and $0^{\circ} \mathrm{C}$ ) to be $10.0 \mathrm{~cm}^{3}$ using the same sample as in Example 19B.1. In this case, there is no dissociation. Calculate the isosteric enthalpy of adsorption at this fractional coverage.

\begin{center}
\begin{tabular}{lrrrrrr}
$T / \mathrm{K}$ & \multicolumn{2}{l}{200} & \multicolumn{1}{l}{210} & \multicolumn{1}{l}{220} & \multicolumn{1}{l}{230} & \multicolumn{1}{l}{240} \\
$p / \mathrm{kPa}$ & 4.00 & 4.95 & 6.03 & 7.20 & 8.47 & 9.85 \\
\end{tabular}
\end{center}

Collect your thoughts The same volume is adsorbed at each temperature, so the surface coverage is the same at all temperatures; that is, the data are for isosteric conditions. You first need to relate the given pressures to a value of $\alpha$ by using the Langmuir isotherm (eqn 19B.2) arranged into $\alpha=\theta / p(1-\theta)$. However, because $\theta$ is constant, this expression reduces to $\alpha=C / p$, where $C$ is a dimensionless constant. It follows that $\ln \left(\alpha p^{\ominus}\right)=\ln \left(C p^{\ominus} / p\right)=-\ln \left(p / p^{\ominus}\right)+\ln C$ and therefore, from eqn 19B.5b, that a plot of $\ln \left(p / p^{\ominus}\right)$ against $1 / T$ should therefore be a straight line of slope $\Delta_{\mathrm{ad}} H^{\ominus} / R$.

The solution With $p^{\ominus}=1 \mathrm{bar}=10^{2} \mathrm{kPa}$, draw up the following table:

\begin{center}
\begin{tabular}{lrrrrrr}
$T / K$ & \multicolumn{1}{l}{200} & \multicolumn{1}{l}{210} & \multicolumn{1}{l}{220} & \multicolumn{1}{l}{230} & \multicolumn{1}{l}{240} & \multicolumn{1}{l}{250} \\
$10^{3} /(T / \mathrm{K})$ & 5.00 & 4.76 & 4.55 & 4.35 & 4.17 & 4.00 \\
$\left(p / p^{\ominus}\right) \times 10^{2}$ & 4.00 & 4.95 & 6.03 & 7.20 & 8.47 & 9.85 \\
$\ln \left(p / p^{\ominus}\right)$ & -3.22 & -3.01 & -2.81 & -2.63 & -2.47 & -2.32 \\
\end{tabular}
\end{center}

The points are plotted in Fig. 19B.4. The slope (of the least squares fitted line) is -0.901 , so $\left(-\Delta_{\mathrm{ad}} H^{\ominus} / R\right) / 10^{3}=0.901 \mathrm{~K}$ and hence

$$
\Delta_{\mathrm{ad}} H^{\ominus}=-\left(0.901 \times 10^{3} \mathrm{~K}\right) \times\left(8.3145 \mathrm{~J} \mathrm{~K}^{-1} \mathrm{~mol}^{-1}\right)=-7.5 \mathrm{~kJ} \mathrm{~mol}^{-1}
$$

\begin{center}
\includegraphics[max width=\textwidth]{2025_02_27_d4b95d9718f0b9e2e6b9g-867(1)}
\end{center}

Figure 19B. 4 The isosteric enthalpy of adsorption can be obtained from the slope of the plot of $\ln \left(p / p^{\ominus}\right)$ against $1 / T$, where $p$ is the pressure needed to achieve the specified surface coverage. The data used are from Example 19B.2.

Self-test 19B.2 Repeat the calculation using the following data, which are isosteric:

\begin{center}
\begin{tabular}{|c|c|c|c|c|c|c|}
\hline
T/K & 200 & 210 & 220 & 230 & 240 & 250 \\
\hline
$p / \mathrm{kPa}$ & 4.32 & 5.59 & 7.07 & 8.80 & 10.67 & 12.80 \\
\hline
 &  &  &  &  & \multicolumn{2}{|l|}{\includegraphics[max width=\textwidth]{2025_02_27_d4b95d9718f0b9e2e6b9g-867}
} \\
\hline
\end{tabular}
\end{center}

Two assumptions of the Langmuir isotherm are the independence and equivalence of the adsorption sites. Deviations from the isotherm can often be traced to the failure of these assumptions. For example, the enthalpy of adsorption often becomes less negative as $\theta$ increases, which suggests that the\\
energetically most favourable sites are occupied first. Also interactions between the molecules already adsorbed on the surface can be important.

\subsection*{(c) The BET isotherm}
A number of isotherms have been developed to deal with cases where deviations from the Langmuir isotherm are important. If the initial adsorbed layer can act as a substrate for further (e.g. physical) adsorption, then instead of the volume of gas adsorbed levelling off at high pressures to a value corresponding to a complete monolayer, it can be expected to rise indefinitely. The most widely used isotherm dealing with multilayer adsorption was derived by Stephen Brunauer, Paul Emmett, and Edward Teller and is called the BET isotherm:

\[
\frac{V}{V_{\operatorname{mon}}}=\frac{c z}{(1-z)\{1-(1-c) z\}} \quad \text { with } z=\frac{p}{p^{*}} \quad \begin{align*}
& \text { BET }  \tag{19B.6}\\
& \text { isotherm }
\end{align*}
\]

where $p^{*}$ is the vapour pressure of the pure liquid substrate, $V$ is the volume of gas adsorbed, and $V_{\operatorname{mon}}$ is the volume of gas corresponding to a complete monolayer. The constant $c$ is characteristic of the system: $c=\alpha_{0} / \alpha_{1}$ where $\alpha_{0}=k_{\mathrm{a}, 0} / k_{\mathrm{d}, 0}$ is the ratio of the rate constants for adsorption and desorption from the substrate, and $\alpha_{1}=k_{\mathrm{a}, 1} / k_{\mathrm{d}, 1}$ is the similar ratio for the subsequent layers. The rather fiddly derivation of this isotherm is given in A deeper look 13 on the website of this book.

Figure 19B. 5 illustrates the shapes of BET isotherms. They rise indefinitely as the pressure is increased because there is no limit to the amount of material that may condense when multilayer coverage is possible. A BET isotherm is not accurate at all pressures, but it is widely used in industry to determine the surface areas of solids.\\
\includegraphics[max width=\textwidth, center]{2025_02_27_d4b95d9718f0b9e2e6b9g-867(2)}

Figure 19B. 5 Plots of the BET isotherm for different values of $c$. The value of $V / V_{\text {mon }}$ rises indefinitely because the model permits the formation of multiple layers on the surface.

The form in which the BET isotherm is commonly used is obtained by inverting both sides of eqn 19B. 6 to obtain

$$
\frac{V_{\text {mon }}}{V}=\frac{(1-z)\{1-(1-c) z\}}{c z}
$$

and then multiplying both sides by $z /(1-z) V_{\text {mon }}$ to obtain

$$
\frac{z}{(1-z) V}=\frac{\{1-(1-c) z\}}{c V_{\operatorname{mon}}}
$$

The right-hand side separates into two terms to give


\begin{equation*}
\frac{z}{(1-z) V}=\overbrace{\frac{1}{c V_{\text {mon }}}}^{\text {Intercept }}+\overbrace{\frac{(c-1)}{c V_{\text {mon }}}}^{\text {Slope }} z \tag{19B.7}
\end{equation*}


Therefore, a plot of $z /(1-z) V$ against $z$ is expected to be a straight line with slope $(c-1) / c V_{\text {mon }}$ and intercept $1 / c V_{\text {mon }}$ at $z=0$. Note that slope/intercept $=\left\{(c-1) / c V_{\text {mon }}\right\} /\left(1 / c V_{\text {mon }}\right)=c-1$.

\subsection*{Example 19B.3 Using the BET isotherm}
The data below relate to the adsorption of $\mathrm{N}_{2}(\mathrm{~g})$ on rutile $\left(\mathrm{TiO}_{2}\right)$ at 75 K .

$$
\begin{array}{lccccccc}
p / \mathrm{kPa} & 0.160 & 1.87 & 6.11 & 11.67 & 17.02 & 21.92 & 27.29 \\
V / \mathrm{mm}^{3} & 601 & 720 & 822 & 935 & 1046 & 1146 & 1254
\end{array}
$$

The volumes have been corrected to 1.00 atm and 273 K and refer to 1.00 g of substrate. At 75 K , the vapour pressure of liquid nitrogen is $p^{*}=76.0 \mathrm{kPa}$. Confirm that these data fit a BET isotherm, and determine the values of $V_{\operatorname{mon}}$ and $c$.

Collect your thoughts Equation 19B. 7 indicates that a plot of $z /(1-z) V$ against $z$, with $z=p / p^{*}$, gives a straight line of slope $(c-1) / c V_{\text {mon }}$ and intercept $1 / c V_{\text {mon }}$ at $z=0$. As remarked in the text, the ratio of the slope to the intercept gives $c-1$. Make sure that the coordinates, slope, and intercept are all dimensionless and interpret them appropriately.

The solution Draw up the following table:

\begin{center}
\begin{tabular}{llcccccc}
$p / \mathrm{kPa}$ & 0.160 & 1.87 & 6.11 & 11.67 & 17.02 & 21.92 & 27.29 \\
$10^{3} z$ & 2.11 & 24.6 & 80.4 & 154 & 224 & 288 & 359 \\
\begin{tabular}{rl}
$10^{4} z /\{(1-z)$ \\
$\left.\left(V / \mathrm{mm}^{3}\right)\right\}$ \\
\end{tabular} & 0.035 & 0.350 & 1.06 & 1.94 & 2.76 & 3.54 & 4.47 \\
\end{tabular}
\end{center}

These points are plotted in Fig. 19B.6. The least squares best line has an intercept at $z=0$ of $10^{4} z /\left\{(1-z)\left(V / \mathrm{mm}^{3}\right)\right\}=0.0411$, or $z /\left\{(1-z)\left(V / \mathrm{mm}^{3}\right)\right\}=4.11 \times 10^{-6}$, so\\
$\frac{1}{c\left(V_{\text {mon }} / \mathrm{mm}^{3}\right)}=4.11 \times 10^{-6}$ and therefore $\frac{1}{c V_{\text {mon }}}=4.11 \times 10^{-6} \mathrm{~mm}^{-3}$\\
The slope of the plot of $10^{4} z /\left\{(1-z)\left(V / \mathrm{mm}^{3}\right)\right\}$ against $10^{3} z$ is $1.22 \times 10^{-2}$, so the slope of $z /\left\{(1-z)\left(V / \mathrm{mm}^{3}\right)\right\}$ against $z$ is $1.22 \times 10^{-2} \times 10^{-4} \times 10^{3}=1.22 \times 10^{-3}$. Therefore

$$
\frac{c-1}{c\left(V_{\text {mon }} / \mathrm{mm}^{3}\right)}=1.22 \times 10^{-3} \text { and so } \frac{c-1}{c V_{\text {mon }}}=1.22 \times 10^{-3} \mathrm{~mm}^{-3}
$$

The ratio of $(c-1) / c V_{\text {mon }}$ and $\left(1 / c V_{\text {mon }}\right)$, from the previous two expressions, is

$$
c-1=\frac{1.22 \times 10^{-3} \mathrm{~mm}^{-3}}{4.11 \times 10^{-6} \mathrm{~mm}^{-3}}=297
$$

so $c=298$. Then

$$
V_{\mathrm{mon}}=\frac{1}{298 \times\left(4.11 \times 10^{-6} \mathrm{~mm}^{-3}\right)}=816 \mathrm{~mm}^{3}
$$

\begin{center}
\includegraphics[max width=\textwidth]{2025_02_27_d4b95d9718f0b9e2e6b9g-868}
\end{center}

Figure 19B. 6 The BET isotherm can be tested, and the parameters determined, by plotting $z /(1-z) V$ against $z$. The data are from Example 19B.3.

Comment. At 1.00 atm and $273 \mathrm{~K}, 816 \mathrm{~mm}^{3}$ corresponds to $3.6 \times 10^{-5} \mathrm{~mol}$, or $2.2 \times 10^{19}$ atoms. Because each atom occupies about $0.16 \mathrm{~nm}^{2}$, the surface area of the sample is about $3.5 \mathrm{~m}^{2}$.

Self-test 19B. 3 Repeat the calculation for the following data which refer to the adsorption of $\mathrm{N}_{2}(\mathrm{~g})$ at 75 K . The volumes have been corrected to 1.00 atm and 273 K .

\begin{center}
\begin{tabular}{|c|c|c|c|c|c|c|c|}
\hline
$p / \mathrm{kPa}$ & 0.160 & 1.87 & 6.11 & 11.67 & 17.02 & 21.92 & 27.29 \\
\hline
$V / \mathrm{cm}^{3}$ & 235 & 559 & 649 & 719 & 790 & 860 & 950 \\
\hline
\end{tabular}
\end{center}

The constant $c$ depends on the temperature and can be related to the enthalpy changes associated with the formation of the first and subsequent monolayers.

How is that done? 19B. 3\\
Relating the constant $c$ in the\\
BET isotherm to relevant enthalpy changes\\
Just as the Gibbs-Helmholtz equation can be used to express the temperature dependence of the 'equilibrium constant' $\alpha$ that appears in the Langmuir isotherm, it can also be used to express the temperature dependence of $\alpha_{0}$ and $\alpha_{1}$, and therefore of their ratio, $c$.

Step 1 Write $\alpha_{0}$ and $\alpha_{1}$ in terms of the relevant Gibbs energy changes\\
The parameter $\alpha_{0}=k_{\mathrm{a}, 0} / k_{\mathrm{d}, 0}$ refers to the formation of the first monolayer (the one attached to the surface) which occurs in the Langmuir isotherm. It follows that $\alpha_{0}$ is related to the standard Gibbs energy of adsorption, $\Delta_{\mathrm{ad}} G^{\ominus}: \Delta_{\mathrm{ad}} G^{\ominus}=$ $-R T \ln \left(\alpha_{0} p^{\ominus}\right)$. It will turn out to be convenient to replace the Gibbs energy of adsorption by the Gibbs energy of desorption, with $\Delta_{\text {des }} G^{\ominus}=-\Delta_{\text {ad }} G^{\ominus}$. It follows that $\Delta_{\text {des }} G^{\ominus}=R T \ln \left(\alpha_{0} p^{\ominus}\right)$, or $\alpha_{0} p^{\ominus}=\mathrm{e}^{\Delta_{\mathrm{des}} G^{\ominus} / R T}$.

The parameter $\alpha_{1}=k_{\mathrm{a}, 1} / k_{\mathrm{d}, 1}$ refers to the formation of the second and subsequent monolayers, which is analogous to the condensation of a gas into a liquid. It follows that $\alpha_{1}$ is related to the standard Gibbs energy of condensation, $\Delta_{\text {con }} G^{\ominus}=-R T \ln \left(\alpha_{1} p^{\ominus}\right)$. In terms of the standard Gibbs energy of vaporization, $\Delta_{\text {vap }} G^{\ominus}=-\Delta_{\text {con }} G^{\ominus}$, it follows that $\Delta_{\text {vap }} G^{\ominus}=$ $R T \ln \left(\alpha_{1} p^{\ominus}\right)$ or $\alpha_{1} p^{\ominus}=\mathrm{e}^{\mathrm{smap}^{\ominus} / R T}$.\\
Step 2 Write c in terms of the relevant Gibbs energy changes\\
You can now use the results from Step 1 to write $c=\alpha_{0} / \alpha_{1}$ in terms of the Gibbs energy changes

$$
c=\frac{\alpha_{0}}{\alpha_{1}}=\frac{\mathrm{e}^{\Delta_{\mathrm{des}} G^{\ominus} / R T}}{\mathrm{e}^{\Delta_{\mathrm{vap}} G^{\ominus} / R T}}=\frac{\mathrm{e}^{\Delta_{\mathrm{des}} H^{\ominus} / R T} \mathrm{e}^{-\Delta_{\mathrm{des}} \mathrm{~s}^{\ominus} / R}}{\mathrm{e}_{\mathrm{mpq}} H^{\ominus} / R T} \mathrm{e}^{-\Delta_{\mathrm{rap}} \mathrm{~s}^{\ominus} / R}
$$

where the Gibbs energies have been written in terms of the corresponding enthalpies and entropies.

\subsection*{Step 3 Simplify the expression}
The entropies of desorption and vaporization can be assumed to be the same because they correspond to similar processes involving the escape of the condensed adsorbate to the gas phase. They cancel to give an expression for $c$ in terms of the standard enthalpies of desorption and vaporization:

$$
-\left|c=\mathrm{e}^{\left(\Delta_{\text {des }} H^{\ominus}-\Delta_{\text {vap }} H^{\ominus}\right) / R T}\right| \begin{aligned}
& \text { The } c \text { constant from the BET isotherm } \\
& \text { in terms of enthalpy changes }
\end{aligned}
$$

From eqn 19B. 8 it follows that the constant $c$ is large when the enthalpy of desorption for the first monolayer is large compared with the enthalpy of vaporization of the liquid adsorbate. In Fig. 19B. 5 it is seen that full monolayer coverage (when $V / V_{\text {mon }}=1$ ) is reached at lower pressures when $c$ is large. This behaviour is consistent with the formation of the first layer becoming more favourable as $\Delta_{\text {vap }} H^{\ominus}$ becomes more negative (and $\Delta_{\text {des }} H^{\ominus}$ more positive).

Example 19B. 3 indicates that $c$ is of the order of $10^{2}$. When $c \gg 1$, the BET isotherm takes the simpler form


\begin{equation*}
\frac{V}{V_{\operatorname{mon}}}=\frac{1}{1-z} \tag{19B.9}
\end{equation*}


BET isotherm when $c \gg 1$\\
This expression is applicable to unreactive gases on polar surfaces, because $\Delta_{\text {des }} H^{\ominus}$ is then significantly greater than $\Delta_{\text {vap }} H^{\ominus}$.

The BET isotherm fits experimental observations moderately well over restricted pressure ranges, but it errs by underestimating the extent of adsorption at low pressures and by overestimating it at high pressures.

\subsection*{(d) The Temkin and Freundlich isotherms}
The Langmuir isotherm assumes that all sites are equivalent and independent, which implies that the enthalpy of adsorption is independent of the surface coverage. Experimentally, it is often found that the enthalpy of adsorption becomes less negative as $\theta$ increases, which suggests that the energetically most favourable sites are occupied first. Various attempts have been made to take these variations into account. The Temkin isotherm,

$$
\theta=c_{1} \ln \left(c_{2} p\right) \quad \text { Temkin isotherm }
$$

(19B.10)\\
where $c_{1}$ and $c_{2}$ are constants, corresponds to supposing that the adsorption enthalpy changes linearly with pressure. The Freundlich isotherm

$$
\theta=c_{1} p^{1 / c_{2}}
$$

Freundlich isotherm (19B.11)\\
corresponds to a logarithmic change. This isotherm attempts to incorporate the role of interactions between the adsorbed molecules on the surface.

Different isotherms agree with experiment more or less well over restricted ranges of pressure, but they remain largely empirical. Empirical, however, does not mean useless for, if the parameters of a reliable isotherm are known, reasonably reliable results can be obtained for the extent of surface coverage under various conditions. This kind of information is essential for any discussion of heterogeneous catalysis (Topic 19C).

\subsection*{19B.2 The rates of adsorption and desorption}
This section takes a more detailed look at adsorption and desorption at a molecular level, focusing in particular on the energetics of chemisorption.

\subsection*{(a) The precursor state}
Figure 19B. 7 shows how the potential energy of a molecule varies with its distance from the surface of the substrate. As the molecule approaches the surface its potential energy falls as it becomes physisorbed into the precursor state for chemisorption. Dissociation into fragments often takes place as a molecule moves into its chemisorbed state, and after an initial increase of energy as bonds in the molecule are distorted there\\
\includegraphics[max width=\textwidth, center]{2025_02_27_d4b95d9718f0b9e2e6b9g-870(1)}

Figure 19B. 7 The potential energy profiles for the dissociative chemisorption of an $A_{2}$ molecule. In each case, $P$ is the enthalpy of (non-dissociative) physisorption and $C$ that for chemisorption (at $T=0$ ). The height of the intermediate peak determines whether the chemisorption is (a) not activated or (b) activated.\\
is a sharp decrease as the adsorbate-substrate bonds reach their full strength. Even if the molecule does not fragment, there is likely to be an initial increase of potential energy as the molecule approaches the surface and its bonds adjust.\\
In most cases, therefore, a potential energy barrier separating the precursor and chemisorbed states is expected. This barrier, though, might be low, and might not rise above the energy zero, which is the energy when the adsorbate is far away (Fig. 19B.7a). In this case, chemisorption is not an activated process and can be expected to be rapid, which is the case for many gas adsorptions on clean metals. In some cases, however, the barrier rises above zero (as in Fig. 19B.7b); such chemisorptions are activated and slower than the non-activated kind. An example is $\mathrm{H}_{2}$ on copper, which has an activation energy in the region of $20-40 \mathrm{~kJ} \mathrm{~mol}^{-1}$.

One point that emerges from this discussion is that rates are not good criteria for distinguishing between physisorption and chemisorption. Chemisorption can be fast if the activation energy is small or zero, but it may be slow if the activation energy is large. Physisorption is usually fast, but it can appear to be slow if adsorption is taking place on a porous medium.

\subsection*{Brief illustration 19B. 1}
Consider two adsorption experiments for hydrogen on different faces of a copper crystal. For adsorption on face 1 the activation energy is $28 \mathrm{~kJ} \mathrm{~mol}^{-1}$ and on face 2 the activation energy is $33 \mathrm{~kJ} \mathrm{~mol}^{-1}$. If Arrhenius behaviour is assumed, and the frequency factors are the same, the ratio of the rates of adsorption on equal areas of the two faces at 250 K is

$$
\begin{aligned}
\frac{\operatorname{Rate}(1)}{\operatorname{Rate}(2)} & =\frac{A \mathrm{e}^{-E_{\mathrm{asad}}(1) / R T}}{A \mathrm{e}^{-E_{\mathrm{adad}}(2) / R T}}=\mathrm{e}^{-\left\{E_{\mathrm{aad}}(1)-E_{\mathrm{atad}}(2)\right\} / R T} \\
& =\mathrm{e}^{5 \times 10^{3} \mathrm{~J} \mathrm{~mol}^{-1} /\left(8.3145 \mathrm{~J} \mathrm{~K}^{-1} \mathrm{~mol}^{-1}\right) \times(250 \mathrm{~K})}=11
\end{aligned}
$$

\subsection*{(b) Adsorption and desorption at the molecular level}
The rate at which a surface is covered by adsorbate depends on the ability of the substrate to dissipate the energy of the incoming particle as it collides with the surface, the process of 'accommodation'. If the energy is not dissipated quickly, the particle migrates over the surface until it reaches an edge or until a vibration expels it into the overlying gas. The proportion of collisions with the surface that successfully lead to adsorption is called the sticking probability, $s$ :\\
$s=\frac{\text { rate of adsorption of particles by the surface }}{\text { rate of collision of particles with the surface }}$\\[0pt]
Sticking probability [definition]\\
(19B.12)\\
The denominator can be calculated from the kinetic model (from $Z_{\mathrm{w}}$, Topic 19A), and the numerator can be measured by observing the rate of change of pressure.

Values of $s$ vary widely. For example, at room temperature CO has $s$ in the range $0.1-1.0$ for several d-metal surfaces, but for $\mathrm{N}_{2}$ on rhenium $s<10^{-2}$, indicating that more than a hundred collisions are needed before one molecule becomes stuck to the surface. Studies on specific crystal planes show a pronounced specificity: for $\mathrm{N}_{2}$ on tungsten at room temperature, $s$ ranges from 0.74 on the (320) faces down to less than 0.01 on the (110) faces. The sticking probability decreases as the surface coverage increases (Fig. 19B.8). A simple assumption is that $s$ is proportional to $1-\theta$, the fraction uncovered, and it is common to write

$$
s=(1-\theta) s_{0}
$$

Commonly used form of the sticking probability\\
(19B.13)\\
where $s_{0}$ is the sticking probability on a perfectly clean surface. The results in the illustration do not fit this expression because\\
\includegraphics[max width=\textwidth, center]{2025_02_27_d4b95d9718f0b9e2e6b9g-870}

Figure 19B.8 The sticking probability of $\mathrm{N}_{2}$ on various faces of a tungsten crystal and its dependence on surface coverage. Note the very low sticking probability for the (110) and (111) faces. (Data provided by Professor D.A. King.)\\
they show that $s$ remains close to $s_{0}$ until the coverage has risen to about $6 \times 10^{13}$ molecules $\mathrm{cm}^{-2}$, and then falls steeply. The explanation is probably that the colliding molecule does not enter the chemisorbed state at once, but moves over the surface until it encounters an empty site.

Desorption is always activated because the particles have to be lifted from the bottom of a potential well. A physisorbed particle vibrates in its shallow potential well, and might shake itself off the surface after a short time. If the temperature dependence of the first-order rate constant for desorption follows Arrhenius behaviour, then $k_{\mathrm{d}}=A \mathrm{e}^{-E_{\mathrm{adse}} / R T}$, where $E_{\mathrm{a}, \text { des }}$ is the activation energy for desorption. Therefore, the temperature dependence of the residence half-life, the half-life for remaining on the surface, is

$$
t_{1 / 2}=\frac{\ln 2}{k_{\mathrm{d}}}=\tau_{0} \mathrm{e}^{E_{\mathrm{ades}} / R T} \quad \tau_{0}=\frac{\ln 2}{A} \quad \text { Residence half-life }
$$

(19B.14)

The time $\tau_{0}$ is the residence half-life in the limit of very high temperature, which is when the activation barrier has negligible effect; $\tau_{0}$ is the lower limit of the residence half-life. Note the positive sign in the exponent: the greater the activation energy for desorption, the larger is the residence half-life.

\subsection*{Brief illustration 19B. 2}
If it is supposed that $1 / \tau_{0}$ is approximately the same as the vibrational frequency of the weak adsorbate-surface bond (about $10^{12} \mathrm{~Hz}$ ) and $E_{\text {a,des }} \approx 25 \mathrm{~kJ} \mathrm{~mol}^{-1}$, then residence half-lives of around 25 ns are predicted at room temperature. Lifetimes close to 1 s are obtained by lowering the temperature to about 100 K . For chemisorption, with $E_{\mathrm{a}, \text { des }}=100 \mathrm{~kJ} \mathrm{~mol}^{-1}$ and guessing that $\tau_{0}=10^{-14} \mathrm{~s}$ (because the adsorbate-substrate bond is quite stiff), a residence half-life of about $3 \times 10^{3} \mathrm{~s}$ (about an hour) at room temperature is expected, decreasing to 1 s at about 370 K .

The desorption activation energy can be measured in several ways. However, such values must be interpreted with caution because they often depend on the fractional coverage, and so might change as desorption proceeds. Moreover, the transfer of concepts such as 'reaction order' and 'rate constant' from bulk studies to surfaces is hazardous, and there are few examples of strictly first-order or second-order desorption kinetics (just as there are few integral-order reactions in the gas phase too).

If the complications are disregarded, one way of measuring the desorption activation energy is to monitor the rate of increase in pressure when the sample is maintained at a series of temperatures, and then to attempt to make an Arrhenius plot. A more sophisticated technique is temperature-programmed desorption (TPD) or thermal desorption spectroscopy (TDS). In these experiments the temperature of the sample is raised\\
\includegraphics[max width=\textwidth, center]{2025_02_27_d4b95d9718f0b9e2e6b9g-871}

Figure 19B.9 The TPD spectrum of $\mathrm{H}_{2}$ desorbing from a PtW layer deposited on the surface of $\gamma$-alumina $\left(\gamma-\mathrm{Al}_{2} \mathrm{O}_{3}\right)$. The profiles correspond to different fractional surface coverages of $\mathrm{H}_{2}$. (Based on F. Lai, D-W. Kim, O.S. Alexeev, G.W. Graham, M. Shelef, and B.C. Gates, Phys. Chem. Chem. Phys. 2, 1997 (2000).)\\
linearly and a surge in the desorption rate (as monitored by a mass spectrometer) is observed when the temperature reaches the point at which desorption occurs rapidly. However, once the desorption is complete (in the sense that there is no more adsorbate to escape from the surface), the desorption rate falls away as the temperature continues to rise. The TPD spectrum, the plot of desorption rate against temperature, therefore shows a peak, the location of which depends on the desorption activation energy (Fig. 19B.9).

In many cases only a single activation energy (and a single peak in the TPD spectrum) is observed. When several peaks are observed they might correspond to adsorption on different crystal planes or to multilayer adsorption. For instance, Cd atoms on tungsten show two activation energies, one of $18 \mathrm{~kJ} \mathrm{~mol}^{-1}$ and the other of $90 \mathrm{~kJ} \mathrm{~mol}^{-1}$. The explanation is that the more tightly bound Cd atoms are attached directly to the substrate, and the less strongly bound are in a layer (or layers) above the first layer. Another example of a system showing two desorption activation energies is CO on tungsten, the values being $120 \mathrm{~kJ} \mathrm{~mol}^{-1}$ and $300 \mathrm{~kJ} \mathrm{~mol}^{-1}$. The explanation is believed to be the existence of two types of metal-adsorbate binding site, one involving a simple $\mathrm{M}-\mathrm{CO}$ bond, the other adsorption with dissociation into individually adsorbed C and O atoms.

\subsection*{(c) Mobility on surfaces}
A further aspect of the strength of the interactions between adsorbate and substrate is the mobility of the adsorbate. Mobility is often a vital feature of a catalyst's activity, because a catalyst might be ineffective if the reactant molecules adsorb so strongly that they cannot migrate.

The activation energy for diffusion over a surface need not be the same as for desorption because the particles may be\\
able to move through valleys between potential peaks without leaving the surface completely. In general, the activation energy for migration is about 10-20 per cent of the energy of the surface-adsorbate bond, but the actual value depends on the extent of coverage. The defect structure of the sample (which depends on the temperature) may also play a dominant role\\
because the adsorbed molecules might find it easier to skip across a terrace than to roll along the foot of a step, and these molecules might become trapped in vacancies in an otherwise flat terrace. Diffusion may also be easier across one crystal face than another, and so the surface mobility depends on which lattice planes are exposed.

\subsection*{Checklist of concepts}
\begin{enumerate}
  \item An adsorption isotherm expresses the variation of the fractional coverage $\theta$ with pressure at constant temperature.2. Flash desorption is a technique in which the sample is suddenly heated and the resulting rise of pressure is interpreted in terms of the amount of adsorbate originally on the substrate.\\
$\square$ 3. Examples of adsorption isotherms include the Langmuir, BET, Temkin, and Freundlich isotherms.4. The sticking probability is the proportion of collisions with the surface that successfully lead to adsorption.5. Desorption is an activated process; the desorption activation energy is measured by temperature-programmed desorption or thermal desorption spectroscopy.
\end{enumerate}

\subsection*{Checklist of equations}
\begin{center}
\begin{tabular}{|c|c|c|c|}
\hline
Property & Equation & Comment & Equation number \\
\hline
\multicolumn{4}{|l|}{Langmuir isotherm:} \\
\hline
(a) without dissociation & $\theta=\alpha p /(1+\alpha p)$ & Independent and equivalent sites, & 19B. 2 \\
\hline
(b) with dissociation & $\theta=(\alpha p)^{1 / 2} /\left\{1+(\alpha p)^{1 / 2}\right\}$ & monolayer coverage & 19B. 4 \\
\hline
Isosteric enthalpy of adsorption & $\left(\partial \ln \left(\alpha p^{\ominus}\right) / \partial(1 / T)\right)_{\theta}=-\Delta_{\text {ad }} H^{\ominus} / R$ &  & 19B.5b \\
\hline
BET isotherm & \( \begin{aligned} & V / V_{\text {mon }}=c z /(1-z)\{1-(1-c) z\}, \\ & z=p / p^{*}, c=\mathrm{e}^{\left(\Delta_{\text {des }} H^{*}-\Delta_{\text {vap }} H^{\ominus}\right) / R T} \end{aligned} \) & Multilayer adsorption & 19B. 6 and 19B. 8 \\
\hline
Temkin isotherm & $\theta=c_{1} \ln \left(c_{2} p\right)$ & Enthalpy of adsorption varies with $\theta$ & 19B. 10 \\
\hline
Freundlich isotherm & $\theta=c_{1} p^{1 / c_{2}}$ & Adsorbate-adsorbate interactions & 19B.11 \\
\hline
Sticking probability & $s=(1-\theta) s_{0}$ & Approximate form & 19B. 13 \\
\hline
\end{tabular}
\end{center}

\section*{TOPIC 19C Heterogeneous catalysis}
\subsection*{> Why do you need to know this material?}
Because the chemical industry relies on heterogeneous catalysis for many of its most important large-scale processes, to see how they might be improved it is necessary to understand their mechanisms.

\subsection*{- What is the key idea?}
Heterogeneous catalysis commonly involves chemisorption of one or more reactants and a consequent lowering of the activation energy.

\subsection*{What do you need to know already?}
Catalysis is introduced in Topic 17F. This Topic builds on the discussion of reaction mechanisms (Topic 17E), and uses the Arrhenius equation (Topic 17D) and adsorption isotherms (Topic 19B).

A heterogeneous catalyst is a catalyst in a different phase from that in which reactants and products are found. An example is the iron-based solid catalyst for the reaction of hydrogen and nitrogen to form ammonia. The metal provides a surface to which the reactants bind, so preparing them for reaction and facilitating their encounters.

\subsection*{19C. 1 Mechanisms of heterogeneous catalysis}
Many catalysts depend on co-adsorption, the adsorption of two or more species. One consequence of the presence of a second species may be the modification of the electronic structure at the surface of a metal. For instance, partial coverage of d-metal surfaces by alkali metals has a pronounced effect on the electron distribution at the surface and reduces the work function of the metal (the energy needed to remove an electron). Such modifiers can act as 'promoters' (to enhance the action of catalysts) or as 'poisons' (to inhibit catalytic action).

Figure 19C. 1 shows the potential energy curve for a reaction in the presence of a heterogeneous catalyst. Differences between Fig. 19C. 1 and Fig. 17F. 5 arise from the fact that heterogeneous catalysis normally depends on at least one reactant being\\
\includegraphics[max width=\textwidth, center]{2025_02_27_d4b95d9718f0b9e2e6b9g-873}

Figure 19C. 1 The reaction profile for catalysed and uncatalysed reactions. The catalysed reaction path includes activation energies for adsorption and desorption as well as an overall lower activation energy for the process.\\
adsorbed (usually chemisorbed) and modified into a form in which it readily undergoes reaction, followed by desorption of products. Modification of the reactant often takes the form of a fragmentation of the reactant molecules. In practice, the catalyst is dispersed as very small particles of linear dimension less than 2 nm on a porous oxide support. Shape-selective catalysts, such as the zeolites, which have a pore size that can distinguish shapes and sizes at a molecular scale, have high internal specific surface areas, in the range of $100-500 \mathrm{~m}^{2} \mathrm{~g}^{-1}$.

\subsection*{(a) Unimolecular reactions}
A surface-catalysed unimolecular reaction is one in which an adsorbed molecule undergoes decomposition on a surface. Its rate law can be written in terms of an adsorption isotherm if it is assumed that the rate is proportional to the extent of surface coverage. For example, if the fractional coverage $\theta$ is given by the Langmuir isotherm (eqn 19B.2, $\theta=\alpha p /(1+\alpha p)$ ), the rate is


\begin{equation*}
v=k_{\mathrm{r}} \theta=\frac{k_{\mathrm{r}} \alpha p}{1+\alpha p} \tag{19C.1}
\end{equation*}


where $p$ is the pressure of the adsorbing substance.

\subsection*{Brief illustration 19C. 1}
The decomposition of phosphine $\left(\mathrm{PH}_{3}\right)$ on tungsten is found to be first-order at low pressures. This order can be understood by using eqn 19C. 1 and noting that when $\alpha p \ll 1, v=k_{\mathrm{r}} \alpha p$,\\
a first-order rate law. On the other hand, when $\alpha p \gg 1$, the 1 in the denominator can be ignored and $v=k_{\mathrm{r}}$, a zeroth-order law. A zeroth-order rate law is expected when the pressure is so high that the entire surface is covered; in this limit, the coverage, and hence the rate, is unaffected by a further rise in the pressure.

\subsection*{(b) The Langmuir-Hinshelwood mechanism}
In the Langmuir-Hinshelwood mechanism (LH mechanism) of surface-catalysed reactions, it is proposed that the reaction takes place by encounters between molecules adsorbed on the surface. For a reaction between species $A$ and $B$, the rate law is expected to be first order in the fractional coverage of $A\left(\theta_{A}\right)$, and of $\mathrm{B}\left(\theta_{\mathrm{B}}\right)$, and second-order overall:

\[
\mathrm{A}+\mathrm{B} \rightarrow \mathrm{P} \quad v=k_{\mathrm{r}} \theta_{\mathrm{A}} \theta_{\mathrm{B}} \quad \begin{align*}
& \text { Langmuir-Hinshelwood }  \tag{19С.2a}\\
& \text { (LH) rate law }
\end{align*}
\]

An example of a reaction thought to proceed by this mechanism is the catalytic oxidation of CO to $\mathrm{CO}_{2}$. The LH rate law can be developed by using an isotherm to relate the fractional coverage of each species to its partial pressure.

How is that done? 19C. 1 Developing the rate law of the LH mechanism

You can derive expressions for the fractional coverage of A and $B$ by analysing the dynamic equilibrium between free and adsorbed molecules in much the same way as in the derivation of the Langmuir isotherm itself (Topic 19B), the difference being that two different species are now competing for the same adsorption sites.

Step 1 Write expressions for the rates of adsorption and desorption of $A$ and $B$\\
The rate of adsorption of $A$ is proportional to the partial pressure of $A, p_{A}$, and the number of vacant sites. If the number of surface sites is $N$, then the number of vacant sites is $N\left(1-\theta_{\mathrm{A}}-\theta_{\mathrm{B}}\right)$. Therefore

Rate of adsorption of $\mathrm{A}=k_{\mathrm{a}, \mathrm{A}} p_{\mathrm{A}} N\left(1-\theta_{\mathrm{A}}-\theta_{\mathrm{B}}\right)$\\
where $k_{\mathrm{a}, \mathrm{A}}$ is the rate constant for adsorption of A . The rate of desorption of A is proportional to the number of sites occupied by A molecules, $N \theta_{\mathrm{A}}$ :

Rate of desorption of $\mathrm{A}=k_{\mathrm{d}, \mathrm{A}} N \theta_{\mathrm{A}}$\\
where $k_{\mathrm{d}, \mathrm{A}}$ is the rate constant for desorption of A . The analogous expressions for B are

\begin{displayquote}
Rate of adsorption of $\mathrm{B}=k_{\mathrm{a}, \mathrm{B}} p_{\mathrm{B}} N\left(1-\theta_{\mathrm{A}}-\theta_{\mathrm{B}}\right)$\\
Rate of desorption of $\mathrm{B}=k_{\mathrm{d}, \mathrm{B}} N \theta_{\mathrm{B}}$
\end{displayquote}

Step 2 Set the rates of adsorption and desorption to be equal At equilibrium, the rates of adsorption and desorption for each species are equal. For the species A

$$
k_{\mathrm{a}, \mathrm{~A}} P_{\mathrm{A}} N\left(1-\theta_{\mathrm{A}}-\theta_{\mathrm{B}}\right)=k_{\mathrm{d}, \mathrm{~A}} N \theta_{\mathrm{A}}
$$

This expression is simplified by introducing $\alpha_{\mathrm{A}}=k_{\mathrm{a}, \mathrm{A}} / k_{\mathrm{d}, \mathrm{A}}$ to give

$$
\alpha_{A} p_{\mathrm{A}}\left(1-\theta_{\mathrm{A}}-\theta_{\mathrm{B}}\right)=\theta_{\mathrm{A}}
$$

and therefore

$$
\left(\alpha_{\mathrm{A}} p_{\mathrm{A}}+1\right) \theta_{\mathrm{A}}+\alpha_{\mathrm{A}} p_{\mathrm{A}} \theta_{\mathrm{B}}=\alpha_{\mathrm{A}} p_{\mathrm{A}}
$$

Similarly for B , and with $\alpha_{\mathrm{B}}=k_{\mathrm{a}, \mathrm{B}} / k_{\mathrm{d}, \mathrm{B}}$,

$$
\alpha_{\mathrm{B}} p_{\mathrm{B}}\left(1-\theta_{\mathrm{A}}-\theta_{\mathrm{B}}\right)=\theta_{\mathrm{B}}
$$

and therefore

$$
\alpha_{\mathrm{B}} p_{\mathrm{B}} \theta_{\mathrm{A}}+\left(\alpha_{\mathrm{B}} p_{\mathrm{B}}+1\right) \theta_{\mathrm{B}}=\alpha_{\mathrm{B}} p_{\mathrm{B}}
$$

The solutions of these two simultaneous equations for $\theta_{\mathrm{A}}$ and $\theta_{\mathrm{B}}$ are

$$
\theta_{\mathrm{A}}=\frac{\alpha_{\mathrm{A}} p_{\mathrm{A}}}{1+\alpha_{\mathrm{A}} p_{\mathrm{A}}+\alpha_{\mathrm{B}} p_{\mathrm{B}}} \quad \theta_{\mathrm{B}}=\frac{\alpha_{\mathrm{B}} p_{\mathrm{B}}}{1+\alpha_{\mathrm{A}} p_{\mathrm{A}}+\alpha_{\mathrm{B}} p_{\mathrm{B}}}
$$

Step 3 Use the expressions for $\theta_{A}$ and $\theta_{B}$ in the rate law\\
Now substitute the expressions for the fractional surface coverage into the rate law, eqn 19C.2a, to give

$$
v=k_{\mathrm{r}} \theta_{\mathrm{A}} \theta_{\mathrm{B}}=k_{\mathrm{r}} \frac{\alpha_{\mathrm{A}} p_{\mathrm{A}}}{1+\alpha_{\mathrm{A}} p_{\mathrm{A}}+\alpha_{\mathrm{B}} p_{\mathrm{B}}} \frac{\alpha_{\mathrm{B}} p_{\mathrm{B}}}{1+\alpha_{\mathrm{A}} p_{\mathrm{A}}+\alpha_{\mathrm{B}} p_{\mathrm{B}}}
$$

which, after minor rearrangement, gives the rate law in terms of the partial pressures and the parameters $\alpha_{\mathrm{A}}$ and $\alpha_{\mathrm{B}}$ :\\
\includegraphics[max width=\textwidth, center]{2025_02_27_d4b95d9718f0b9e2e6b9g-874}

The parameters $\alpha$ and the rate constant $k_{\mathrm{r}}$ are all temperaturedependent, so the overall temperature dependence of the rate may be strongly non-Arrhenius (in the sense that the reaction rate constant is unlikely to be proportional to $\left.\mathrm{e}^{-E_{a} / R T}\right)$. The LH mechanism has been analysed in terms of the adsorption and subsequent reaction between species A and B, but it may be that on adsorption either or both species dissociates to give fragments which then react.

\subsection*{(c) The Eley-Rideal mechanism}
In the Eley-Rideal mechanism (ER mechanism) of a surfacecatalysed reaction, it is proposed that a gas-phase molecule collides with another molecule already adsorbed on the surface. The rate of formation of product is proportional to the partial pressure, $p_{\mathrm{B}}$, of the non-adsorbed gas B and the fractional surface coverage, $\theta_{A}$, of the adsorbed gas A , to give the rate law

$$
v=k_{\mathrm{r}} p_{\mathrm{B}} \theta_{\mathrm{A}} \quad \text { Eley-Rideal rate law } \quad \text { (19C.3) }
$$

The rate of the catalysed reaction might be much larger than for the uncatalysed gas-phase reaction because the reaction on the surface has a low activation energy and the adsorption itself is often not activated. If it is assumed that the Langmuir isotherm applies to species $A$, the fractional coverage is $\theta_{\mathrm{A}}=\alpha p_{\mathrm{A}} /\left(1+\alpha p_{\mathrm{A}}\right)$ so the EL rate law becomes


\begin{equation*}
v=\frac{k_{\mathrm{r}} \alpha p_{\mathrm{B}} p_{\mathrm{A}}}{1+\alpha p_{\mathrm{A}}} \tag{19C.4}
\end{equation*}


\subsection*{Brief illustration 19C. 2}
According to eqn 19C.4, when the partial pressure of A is high (in the sense $\alpha p_{\mathrm{A}} \gg 1$ ), the denominator is simply $\alpha p_{\mathrm{A}}$, and the rate is $k_{\mathrm{r}} p_{\mathrm{B}}$. At such high pressures the surface is completely covered with A , so increasing the pressure of A has no effect. The rate of reaction is then limited by the rate at which B reacts with adsorbed A. When the pressure of A is low ( $\alpha p_{\mathrm{A}} \ll 1$ ), the term $\alpha p_{\mathrm{A}}$ in the denominator can be ignored and the rate of reaction becomes $k_{\mathrm{r}} \alpha p_{\mathrm{A}} p_{\mathrm{B}}$. Now the coverage of the surface is low and so increasing the pressure of A increases the surface coverage, and hence the rate.

Almost all surface-catalysed reactions are thought to take place by the LH mechanism, but a number of reactions with an ER mechanism have also been identified from molecular beam investigations. For example, the reaction between gaseous H atoms and adsorbed D atoms to form gaseous HD is thought to proceed by the ER mechanism in which an H atom collides directly with an adsorbed D atom, picking it up to form HD. However, the two mechanisms should really be thought of as ideal limits with all reactions lying somewhere between the two and showing features of each one.

\subsection*{19C. 2 Catalytic activity at surfaces}
It has become possible to investigate how the catalytic activity of a surface depends on its structure as well as its composition. For instance, the cleavage of $\mathrm{C}-\mathrm{H}$ and $\mathrm{H}-\mathrm{H}$ bonds appears to\\
depend on the presence of steps and kinks, and a terrace often has only minimal catalytic activity.

The reaction $\mathrm{H}_{2}+\mathrm{D}_{2} \rightarrow 2 \mathrm{HD}$ has been studied in detail. For this reaction, terrace sites are inactive but one molecule in ten reacts when it strikes a step. Although the step itself might be the important feature, it may be that the presence of the step merely exposes a more reactive crystal face (the step face itself). Likewise, the dehydrogenation of hexane to hexene depends strongly on the kink density, and it appears that kinks are needed to cleave $\mathrm{C}-\mathrm{C}$ bonds. These observations suggest a reason why even small amounts of impurities may poison a catalyst: they are likely to attach to step and kink sites, and so impair the activity of the catalyst entirely. A constructive outcome is that the extent of dehydrogenation may be controlled relative to other types of reactions by seeking impurities that adsorb at kinks and act as specific poisons.

The activity of a catalyst depends on the strength of chemisorption as indicated by the 'volcano' curve in Fig. 19C. 2 (which is so-called on account of its general shape). To be active, the catalyst should be extensively covered by adsorbate, which is the case if chemisorption is strong. On the other hand, if the strength of the substrate-adsorbate interaction becomes too great, the activity declines either because the other reactant molecules cannot react with the adsorbate or because the adsorbate molecules are immobilized on the surface. This pattern of behaviour suggests that the activity of a catalyst should initially increase with strength of adsorption (as measured, for instance, by the enthalpy of adsorption) and then decline, and that the most active catalysts should be those lying near the summit of the volcano. Most active metals are those that lie close to the middle of the d block. Many metals are suitable for adsorbing gases, and some trends are summarized in Table 19C.1.\\
\includegraphics[max width=\textwidth, center]{2025_02_27_d4b95d9718f0b9e2e6b9g-875}

Figure 19C. 2 A 'volcano curve' of catalytic activity arises because although the reactants must adsorb reasonably strongly, they must not adsorb so strongly that they are immobilized. The lower curve refers to the first series of d-block metals, the upper curve to the second and third series of d-block metals.

Table 19C. 1 Chemisorption abilities*

\begin{center}
\begin{tabular}{llllllll}
\hline
 & $\mathrm{O}_{2}$ & $\mathrm{C}_{2} \mathrm{H}_{2}$ & $\mathrm{C}_{2} \mathrm{H}_{4}$ & CO & $\mathrm{H}_{2}$ & $\mathrm{CO}_{2}$ & $\mathrm{~N}_{2}$ \\
\hline
$\mathrm{Ti}, \mathrm{Cr}, \mathrm{Mo}, \mathrm{Fe}$ & + & + & + & + & + & + & + \\
$\mathrm{Ni}, \mathrm{Co}$ & + & + & + & + & + & + & - \\
$\mathrm{Pd}, \mathrm{Pt}$ & + & + & + & + & + & - & - \\
$\mathrm{Mn}, \mathrm{Cu}$ & + & + & + & + & $\pm$ & - & - \\
$\mathrm{Al}, \mathrm{Au}$ & + & + & + & - & - & - & - \\
$\mathrm{Li}, \mathrm{Na}, \mathrm{K}$ & + & + & - & - & - & - & - \\
$\mathrm{Mg}, \mathrm{Ag}, \mathrm{Zn}, \mathrm{Pb}$ & + & - & - & - & - & - & - \\
\hline
\end{tabular}
\end{center}

${ }^{\star}+$, Strong chemisorption; $\pm$, chemisorption; - , no chemisorption.

\subsection*{Brief illustration 19C. 3}
The data in Table 19C. 1 show that for a number of metals the general order of chemisorption ability decreases along the series $\mathrm{O}_{2}, \mathrm{C}_{2} \mathrm{H}_{2}, \mathrm{C}_{2} \mathrm{H}_{4}, \mathrm{CO}, \mathrm{H}_{2}, \mathrm{CO}_{2}, \mathrm{~N}_{2}$. Some of these molecules adsorb dissociatively (e.g. $\mathrm{H}_{2}$ ). Elements from the d block, such as iron, titanium, and chromium, adsorb all these gases, but manganese and copper are unable to adsorb $\mathrm{N}_{2}$ and $\mathrm{CO}_{2}$. Metals towards the left of the periodic table (e.g. magnesium) can adsorb (and, in fact, react with) only the most active gas $\left(\mathrm{O}_{2}\right)$.

\subsection*{Checklist of concepts}
1. A heterogeneous catalyst is a catalyst in a different phase from the reaction mixture.2. In the Langmuir-Hinshelwood mechanism of sur-face-catalysed reactions, the reaction takes place by encounters between molecules adsorbed on the surface.\begin{enumerate}
  \setcounter{enumi}{2}
  \item In the Eley-Rideal mechanism of a surface-catalysed reaction, a gas-phase molecule collides with another molecule already adsorbed on the surface.4. The activity of a catalyst depends on the strength of chemisorption.
\end{enumerate}

\subsection*{Checklist of equations}
\begin{center}
\begin{tabular}{llll}
\hline
Property & Equation & Comment & Equation number \\
\hline
Langmuir-Hinshelwood mechanism & $v=k_{\mathrm{r}} \theta_{\mathrm{A}} \theta_{\mathrm{B}}$ & A and B both adsorbed & 19C.2a \\
Eley-Rideal mechanism & $v=k_{\mathrm{r}} p_{\mathrm{B}} \theta_{\mathrm{A}}$ & Only A adsorbed & 19C.3 \\
\hline
\end{tabular}
\end{center}

\section*{TOPIC 19D Processes at electrodes}
\subsection*{Why do you need to know this material?}
A knowledge of the factors that determine the rate of electron transfer at electrodes leads to a better understanding of the charging and discharging of batteries, the production of power using solar cells, and manufacturing using electrolysis, all of which are important technologies with wide impact.

\subsection*{What is the key idea?}
The rate of oxidation and reduction at an electrode depends on the height of the activation barrier, which can be modified by applying a potential difference across the solution/electrode interface.

\subsection*{What do you need to know already?}
You need to be familiar with electrochemical cells (Topic 6 C ), electrode potentials (Topic 6D), and the thermodynamic version of transition-state theory (Topic 18C), particularly the Gibbs energy of activation.

The surface of a solid electrode is in contact with the ions in an electrolyte solution. The rates of oxidation and reduction at this interface depend on how rapidly electrons can be transferred through it.

\subsection*{19D. 1 The electrode-solution interface}
An electrode in contact with a solution of an electrolyte acquires a charge as a result either of the escape of atoms into the solution as cations, leaving behind a negative charge, or as a result of ions becoming attached to the surface. As the electrode becomes charged, an electrical potential difference develops across the interface and makes that process more difficult. For example, if the charge arises from the escape of atoms as cations, the increasing negative charge on the electrode makes it more unfavourable for the cations to leave. Eventually equilibrium is reached with a characteristic potential difference between the electrode and the solution.

The charge on the electrode affects the composition of the surrounding electrolyte solution because it is energetically\\
favourable for ions with the opposite charge to cluster nearby. This tendency, however, is disrupted by thermal motion and various models have been developed to describe the outcome of this competition, some simply by ignoring it. The modification of the local concentrations near an electrode implies that it might be misleading to use activity coefficients characteristic of the bulk to discuss the thermodynamic properties of ions near the interface. This is one of the reasons why measurements of the dynamics of electrode processes are almost always done by using a large excess of supporting electrolyte (e.g. a $1 \mathrm{moldm}^{-3}$ solution of a salt, an acid, or a base). Under such conditions, the activity coefficients are almost constant because the inert ions dominate the effects of local changes caused by any reactions taking place. The use of a concentrated solution also minimizes ion migration effects.

The most primitive model of the boundary between the electrode and the electrolyte solution is as an electrical double layer, in which it is supposed that there is a sheet of positive charge at the surface of the electrode and a sheet of negative charge next to it in the solution (or vice versa).

More sophisticated models for the interface introduce a more gradual change in the structure of the solution. In the Helmholtz layer model solvated ions lie along the surface of the electrode but are held away from it by their hydration spheres (Fig. 19D.1). The location of the sheet of ionic charge, which is called the outer Helmholtz plane (OHP), is identified\\
\includegraphics[max width=\textwidth, center]{2025_02_27_d4b95d9718f0b9e2e6b9g-877}

Figure 19D. 1 A simple model of the electrode-solution interface treats it as two rigid planes of charge. One plane, the outer Helmholtz plane (OHP), is due to the ions with their solvating molecules and the other plane is that of the electrode itself. The plot shows the dependence of the electric potential with distance from the electrode surface according to this model. Between the electrode surface and the OHP, the potential varies linearly from $\phi_{M^{\prime}}$ the value in the metal, to $\phi_{5}$, the value in the bulk of the solution.\\
\includegraphics[max width=\textwidth, center]{2025_02_27_d4b95d9718f0b9e2e6b9g-878(1)}

Figure 19D. 2 The Gouy-Chapman model of the electrical double layer treats the outer region as an atmosphere of counter-charge, similar to the Debye-Hückel model of an ionic atmosphere.\\
as the plane running through the solvated ions. In this simple model, the electrical potential changes linearly within the layer from $\phi_{\mathrm{M}}$ at the metal to $\phi_{\mathrm{S}}$, the value characteristic of the solution, at the OHP. In a refinement of this model, ions that have discarded their solvating molecules and have become attached to the electrode surface by chemical bonds are regarded as forming the inner Helmholtz plane (IHP).

The Helmholtz layer model ignores the disrupting effect of thermal motion, which tends to break up and disperse the rigid outer plane of charge. In the Gouy-Chapman model of the diffuse double layer, the disordering effect of thermal motion is taken into account in much the same way as the Debye-Hückel model describes the ionic atmosphere of an ion (Topic 5F). The difference is that the central ion is replaced by an infinite plane electrode. Figure 19D. 2 shows how, in the Gouy-Chapman model, the local concentrations of cations and anions differ from their bulk concentrations. Ions of opposite charge cluster close to the electrode and ions of the same charge are repelled from it. As a result, the potential changes smoothly from $\phi_{M}$ to $\phi_{S}$.\\
Neither the Helmholtz nor the Gouy-Chapman model is a very good representation of the structure of the double layer. The former overemphasizes the rigidity of the local solution; the latter underemphasizes its structure. The two are combined in the Stern model, in which the ions closest to the electrode are constrained into a rigid Helmholtz plane while beyond that plane the ions are dispersed as in the GouyChapman model (Fig. 19D.3). Yet another level of sophistication is found in the Grahame model, which adds an inner Helmholtz plane to the Stern model.\\
The potential difference between the bulk metal and the bulk solution is the Galvani potential difference, $\Delta \phi=$ $\phi_{\mathrm{M}}-\phi_{\mathrm{S}}$. If the electrode is part of a cell from which no current is being drawn the Galvani potential difference can be identified with the electrode potential discussed in Topic 6D.\\
\includegraphics[max width=\textwidth, center]{2025_02_27_d4b95d9718f0b9e2e6b9g-878}

Figure 19D. 3 A representation of the Stern model of the electrode-solution interface. The model incorporates the concepts of an outer Helmholtz plane near the electrode surface and of a diffuse double layer further away from the surface.

However, the value of $\Delta \phi$ can be altered at will by the application of an external electrical potential difference to the cell, and when the cell is producing current the potential difference at the electrode/electrolyte interface also changes from its zero-current value.

\subsection*{19D.2 The current density at an electrode}
The current density, $j$, is the electric current (in amperes, $A$; $1 \mathrm{~A}=1 \mathrm{Cs}^{-1}$ ) flowing through a region of an electrode divided by the area of the region (in square metres or a submultiple, such as square centimetres). Current is the rate of flow of charge, so a current density of $1 \mathrm{Acm}^{-2}$ represents a flow of about $10 \mu \mathrm{~mol}$ of electrons per second per square centimetre. The current density is a measure of the rate of the electrontransfer process occurring at the electrode.

\subsection*{(a) The Butler-Volmer equation}
As explained in Topic 6C, a cathode is the site of reduction and an anode is the site of oxidation. This nomenclature is carried over into the classification of the current density. A flow of electrons from the electrode to bring about reduction of the electroactive species in the solution is called the cathodic current density, $j_{c}$. The opposite flow, from solution into the electrode due to oxidation of the electroactive species, is called the anodic current density, $j_{2}$. The net current density, $j$, is the difference of these two current densities, $j=j_{c}-j_{\mathrm{a}}$. Reduction is dominant if $j_{\mathrm{c}}>j_{\mathrm{a}}$ and therefore $j>0$ : the current density is then said to be 'cathodic'. Oxidation is dominant if $j_{c}<j_{\mathrm{a}}$, corresponding to $j<0$ : the current density is then 'anodic' (Fig. 19D.4).\\
\includegraphics[max width=\textwidth, center]{2025_02_27_d4b95d9718f0b9e2e6b9g-879}

Figure 19D. 4 The net current density is defined as the difference between the cathodic and anodic current densities. (a) When $j_{\mathrm{a}}>j_{c^{\prime}}$ the net current is anodic, and there is a net oxidation of the species in solution. (b) When $j_{c}>j_{a^{\prime}}$ the net current is cathodic, and the net process is reduction.

When the electrode is at equilibrium there is no net current flow, and $\Delta \phi$ can be identified with the electrode potential $E$. This equilibrium is a dynamic one in which both the anodic and cathodic processes are taking place, but in such a way that their current densities are equal. The anodic (or cathodic) current density at equilibrium is called the exchange-current density, $j_{0}$.

If the potential difference at the interface, $\Delta \phi$, differs from $E$, net current flows through the electrode. The overpotential $\eta$ is defined as

$$
\eta=E^{\prime}-E
$$

Overpotential [definition]\\
(19D.1)\\
where $E^{\prime}$ is the potential difference applied to the cell or its potential difference under working conditions. It follows that $\Delta \phi=E+\eta$.

The energies of the charged species involved in the electron transfer process depend on the electrical potential on each side of the interface and are therefore affected by the potential difference across it. The rate of electron transfer is similarly affected, so the task is to relate the resulting current density to the overpotential.

\subsection*{How is that done? 19D. 1 Deriving the relation between current density and the overpotential}
The current density is determined by the rate of the electrontransfer process taking place at the electrode between an oxidized species $\mathrm{Ox}^{+}$and a reduced species Red. Both Red and $\mathrm{Ox}^{+}$are in solution; the electrons involved in the redox process are in the electrode. The rate constant for the reduction step (the cathodic process) is $k_{c}$, and for the reverse oxidation step (the anodic process) is $k_{\mathrm{a}}$.

$$
\mathrm{Ox}^{+}(\text {solution })+\mathrm{e}^{-}(\text {electrode }) \underset{k_{\mathrm{s}}}{\stackrel{k_{\mathrm{s}}}{\leftrightarrows}} \operatorname{Red}(\text { solution })
$$

The model does not depend on this choice of the charges; they have been chosen for convenience.

Step 1 Write expressions for the rate of oxidation and reduction\\
An electrode reaction is heterogeneous, so its rate is specified by the flux of material. This flux is the amount of material produced over a region of the electrode surface in an interval of time, divided by the area of the region and the duration of the interval. A first-order heterogeneous rate law has the form

$$
\text { Product flux }=k_{\mathrm{r}}[\mathrm{X}]
$$

where $[\mathrm{X}]$ is the molar concentration of the relevant electroactive species in the solution. The rate constant has dimensions of length/time (with units, for example, of centimetres per second, $\mathrm{cm} \mathrm{s}^{-1}$ ). If the molar concentrations of the oxidized and reduced species are $\left[\mathrm{Ox}^{+}\right]$and [Red], respectively, then the rate of reduction of $\mathrm{Ox}^{+}$is $k_{c}\left[\mathrm{Ox}^{+}\right]$and the rate of oxidation of Red is $k_{\mathrm{a}}$ [Red].

Step 2 Write expressions for the current density in terms of the rate\\
The cathodic current density, $j_{c}$, is equal to the flux multiplied by Faraday's constant, $F=N_{A} e$, the magnitude of the charge per mole of electrons:

$$
j_{\mathrm{c}}=F k_{c}\left[\mathrm{Ox}^{+}\right] \text {for } \mathrm{Ox}^{+}+\mathrm{e}^{-} \rightarrow \mathrm{Red}
$$

Similarly, the anodic current density, $j_{\mathrm{a}}$, is

$$
j_{\mathrm{a}}=F k_{\mathrm{a}}[\operatorname{Red}] \text { for Red } \rightarrow \mathrm{Ox}^{+}+\mathrm{e}^{-}
$$

The net current density at the electrode is the difference

$$
j=j_{\mathrm{a}}-j_{\mathrm{c}}=F k_{\mathrm{a}}[\operatorname{Red}]-F k_{\mathrm{c}}\left[\mathrm{Ox}^{+}\right]
$$

Step 3 Write the rate constants in terms of the Gibbs energies of activation\\
Now write the two rate constants in a form suggested by transition-state theory (Topic 18C) as

$$
k_{\mathrm{r}}=B \mathrm{e}^{-\Delta^{\ddagger} G / R T}
$$

where $\Delta^{\ddagger} G$ is the Gibbs energy of activation and $B$ is a constant with the same dimensions as $k_{r}$. Then

$$
j_{\mathrm{a}}=F B_{\mathrm{a}}[\text { Red }] \mathrm{e}^{-\mathrm{-}^{ \pm} G_{\mathrm{a}} / R T} \quad j_{\mathrm{c}}=F B_{\mathrm{c}}\left[\mathrm{Ox}^{+}\right] \mathrm{e}^{-\Delta^{\star} G_{c} / R T}
$$

Step 4 Relate the Gibbs energies of activation to the electrical potential difference\\
If a species of charge number $z$ (for $\mathrm{Ox}^{+}, z=+1$; for Red, $z=0$; for $\mathrm{e}^{-}, z=-1$ ) is present in a region of electrical potential $\phi$ its standard chemical potential is

$$
\bar{\mu}^{\ominus}=\mu_{0}^{\ominus}+z F \phi
$$

where $\mu_{0}^{\ominus}$ is the standard chemical potential in the absence of an electrical potential. The quantity $\bar{\mu}$ is called the electrochemical potential. The species $\mathrm{Ox}^{+}$and Red are in the solution, and so experience the potential $\phi_{s}$, whereas the electrons\\
are in the metallic electrode and so experience the potential $\phi_{\mathrm{M}}$. The reduced species is electrically neutral (in this formalism). The standard electrochemical potentials of the three species are therefore

$$
\begin{aligned}
& \bar{\mu}^{\ominus}\left(\mathrm{Ox}^{+}\right)=\mu_{0}^{\ominus}\left(\mathrm{Ox}^{+}\right)+F \phi_{\mathrm{S}} \quad \bar{\mu}^{\ominus}(\mathrm{Red})=\mu_{0}^{\ominus}(\mathrm{Red}) \\
& \bar{\mu}^{\ominus}\left(\mathrm{e}^{-}\right)=\mu_{0}^{\ominus}\left(\mathrm{e}^{-}\right)-F \phi_{\mathrm{M}}
\end{aligned}
$$

In the reaction $\mathrm{Ox}^{+}+\mathrm{e}^{-} \rightarrow$ Red the standard Gibbs energy of the reactants is therefore

$$
\begin{aligned}
G_{\mathrm{m}}^{\ominus}(\text { reactants }) & =\bar{\mu}^{\ominus}\left(\mathrm{Ox}^{+}\right)+\bar{\mu}^{\ominus}\left(\mathrm{e}^{-}\right)=\mu_{0}^{\ominus}\left(\mathrm{Ox}^{+}\right)+\mu_{0}^{\ominus}\left(\mathrm{e}^{-}\right)+\overbrace{F \phi_{\mathrm{s}}-F \phi_{\mathrm{M}}}^{-F \Delta \phi} \\
& =G_{\mathrm{m}}^{\ominus}(0, \text { reactants })-F \Delta \phi
\end{aligned}
$$

where $G_{\mathrm{m}}^{\ominus}(0$, reactants) is the standard molar Gibbs energy when no electrical potential is applied, and $\Delta \phi=\phi_{\mathrm{M}}-\phi_{\mathrm{S}}$.

If the activated complex appears early along the reaction pathway, meaning that it has a structure not too dissimilar from the reactants $\left(\mathrm{Ox}^{+}+\mathrm{e}^{-}\right)$, then its Gibbs energy is affected by the applied electrical potential in a similar way. Therefore, because the potential energy difference has the same effect on the Gibbs energies of both the reactants and activated complex, the Gibbs energy of activation is unaffected by the value of $\Delta \phi$ (Fig. 19D.5a). In contrast, if the activated complex appears late in the reaction pathway, and so resembles the electrically neutral product Red, then its Gibbs energy is unchanged by the potential difference. As $\Delta \phi$ increases, the standard molar Gibbs energy of the reactants $\left(\mathrm{Ox}^{+}+\mathrm{e}^{-}\right)$is lowered by $F \Delta \phi$, so in this case the Gibbs energy of activation is increased by $F \Delta \phi$ (Fig. 19D.5b).\\
\includegraphics[max width=\textwidth, center]{2025_02_27_d4b95d9718f0b9e2e6b9g-880}

Figure 19D. 5 Profiles of how the Gibbs energy varies between the oxidized species ( $\mathrm{Ox}^{+}+\mathrm{e}^{-}$) and the reduced species (Red) at an electrode. The purple line shows the profile, and the blue line shows how it is modified when a potential difference $\Delta \phi$ is applied across the electrode. In\\
(a) the transition state resembles the oxidized species; in\\
(b) the transition state resembles the reduced species.

These two special cases can be brought together if the Gibbs energy of activation for the cathodic (reduction) process is written

$$
\Delta^{\ddagger} G_{c}=\Delta^{\ddagger} G_{c}(0)+\alpha F \Delta \phi
$$

where $\Delta^{\ddagger} G_{\mathrm{c}}(0)$ is the Gibbs energy of activation when $\Delta \phi=0$. The parameter $\alpha$, the transfer coefficient, lies in the range 0 to 1 : it is 0 if the activated complex closely resembles the reactants (Fig. 19D.5a), and 1 if the complex closely resembles the products (Fig. 19D.5b). Experimentally, $\alpha$ is often found to be about 0.5.\\
A similar argument applies to the anodic process, the oxidation Red $\rightarrow \mathrm{Ox}^{+}+\mathrm{e}^{-}$, which is the reverse of the cathodic process. As is evident from Fig. 19D.5a, if the activated complex resembles $\mathrm{Ox}^{+}+\mathrm{e}^{-}(\alpha=0)$, the Gibbs energy of activation for the anodic step is decreased by $-F \Delta \phi$. On the other hand, if the activated complex resembles Red (Fig. 19D.5b, $\alpha=1$ ), the Gibbs energy of activation for the anodic step is unaffected by a change in $\Delta \phi$. The overall effect on the Gibbs energy of activation for the anodic process can therefore be written

$$
\Delta^{\ddagger} G_{\mathrm{a}}=\Delta^{\ddagger} G_{\mathrm{a}}(0)-(1-\alpha) F \Delta \phi
$$

Step 5 Write the rate constants using the expressions for the Gibbs energy of activation

Now insert the Gibbs energies of activation into the expressions for $j_{\mathrm{a}}$ and $j_{\mathrm{c}}$ to give

$$
j_{\mathrm{a}}=F B_{\mathrm{a}}[\text { Red }] \mathrm{e}^{-\Delta^{ \pm} G_{\mathrm{a}}(0) / R T} \mathrm{e}^{(1-\alpha) f \Delta \phi} \quad j_{\mathrm{c}}=F B_{\mathrm{c}}[\mathrm{Ox}] \mathrm{e}^{-\Delta^{\ddagger} G_{\mathrm{c}}(0) / R T} \mathrm{e}^{-\alpha f \Delta \phi}
$$

where the appearance of the expressions has been simplified by writing $F / R T=f$.

\subsection*{Step 6 Consider the effect of the overpotential}
If a potential difference is applied such that the net current density is zero, $\Delta \phi$ can be identified as the electrode potential, $E$. The current densities are then both equal to the exchangecurrent density, $j_{0}$ :

$$
j_{0}=F B_{\mathrm{a}}[\text { Red }] \mathrm{e}^{-\Delta^{\ddagger} G_{\mathrm{a}}(0) / R T} \mathrm{e}^{(1-\alpha) f E}=F B_{\mathrm{c}}[\mathrm{Ox}] \mathrm{e}^{-\Delta^{\ddagger} G_{c}(0) / R T} \mathrm{e}^{-\alpha f E}
$$

The role of the overpotential can now be identified by substituting $\Delta \phi=E+\eta$ :

$$
\begin{aligned}
j_{\mathrm{c}} & =F B_{\mathrm{c}}[\mathrm{Ox}] \mathrm{e}^{-\Delta^{\ddagger} G_{\mathrm{c}}(0) / R T} \mathrm{e}^{-\alpha f \Delta \phi}=F B_{\mathrm{c}}[\mathrm{Ox}] \mathrm{e}^{-\Delta^{\ddagger} G_{c}(0) / R T} \mathrm{e}^{-\alpha f(E+\eta)} \\
& =\overbrace{F B_{\mathrm{c}}[\mathrm{Ox}] \mathrm{e}^{-\Delta^{\ddagger} G_{c}(0) / R T} \mathrm{e}^{-\alpha f E}}^{j_{0}} \mathrm{e}^{-\alpha f \eta}=j_{0} \mathrm{e}^{-\alpha f \eta}
\end{aligned}
$$

A similar argument gives

$$
j_{\mathrm{a}}=j_{0} \mathrm{e}^{(1-\alpha) f \eta}
$$

The net current density is $j=j_{a}-j_{c}$; therefore\\
$-j=j_{0}\left\{\mathrm{e}^{(1-\alpha) f \eta}-\mathrm{e}^{-\alpha f \eta}\right\}$\\
(19D.2)\\
Butler-Volmer equation

Equation 19D. 2 is the Butler-Volmer equation. It can be interpreted as follows:

\begin{itemize}
  \item When the overpotential $\eta$ is zero, there is no net current density (there are equal and opposite flows).
  \item If $\alpha=0$, the cathodic current density is equal to the exchange-current density and is independent of the overpotential.
  \item If $\alpha=1$, the anodic current density is equal to the exchange-current density and is independent of the overpotential.
  \item Provided $0<\alpha<1$, as $\eta$ becomes increasingly positive the anodic current density dominates the cathodic current density and the dominant process is the oxidation Red $\rightarrow \mathrm{Ox}^{+}+\mathrm{e}^{-}$. As $\eta$ becomes increasingly negative, the cathodic current density dominates the anodic current density and the dominant process is the reduction $\mathrm{Ox}^{+}+\mathrm{e}^{-} \rightarrow$ Red.
\end{itemize}

Figure 19D. 6 shows how eqn 19D. 2 predicts the dependence of the net current density on the overpotential for different values of the transfer coefficient. When the overpotential is so small that $f \eta \ll 1$ (in practice, $\eta$ less than about 10 mV ) the exponentials in eqn 19D. 2 can be expanded by using $\mathrm{e}^{x}=1+$ $x+\cdots$ to give


\begin{equation*}
j=j_{0} \overbrace{\{1+(1-\alpha) f \eta+\cdots}^{\mathrm{e}^{(1-\alpha) f \eta}}-\overbrace{(1-\alpha f \eta+\cdots)\}}^{\mathrm{e}^{-\alpha f \eta}} \approx j_{0} f \eta \tag{19D.3}
\end{equation*}


This equation shows that the net current density is proportional to the overpotential, so at low overpotentials the interface obeys Ohm's law. The relation can also be reversed to calculate the overpotential that must exist if a current density $j$ has been established by some external circuit:


\begin{equation*}
\eta=\frac{R T j}{F j_{0}} \tag{19D.4}
\end{equation*}


\begin{center}
\includegraphics[max width=\textwidth]{2025_02_27_d4b95d9718f0b9e2e6b9g-881}
\end{center}

Figure 19D. 6 The dependence of the current density on the overpotential for different values of the transfer coefficient.

\subsection*{Brief illustration 19D. 1}
The exchange-current density of a $\operatorname{Pt}(\mathrm{s})\left|\mathrm{H}_{2}(\mathrm{~g})\right| \mathrm{H}^{+}(\mathrm{aq})$ electrode at 298 K is $0.79 \mathrm{~mA} \mathrm{~cm}^{-2}$. The current density when the overpotential is +5.0 mV is obtained by using eqn 19D. 3 and $f=F / R T=1 /(25.69 \mathrm{mV})$ :

$$
j=j_{0} f \eta=\frac{\left(0.79 \mathrm{mAcm}^{-2}\right) \times(5.0 \mathrm{mV})}{25.69 \mathrm{mV}}=0.15 \mathrm{mAcm}^{-2}
$$

The current through an electrode of total area $5.0 \mathrm{~cm}^{2}$ is therefore 0.75 mA .

Some experimental values for the Butler-Volmer parameters are given in Table 19D.1. From them it is seen that ex-change-current densities vary over a very wide range. Their values are generally large when the redox process involves no bond breaking (as in the $\left[\mathrm{Fe}(\mathrm{CN})_{6}\right]^{3-},\left[\mathrm{Fe}(\mathrm{CN})_{6}\right]^{4-}$ couple) or if only weak bonds are broken (as in $\mathrm{Cl}_{2}, \mathrm{Cl}^{-}$). They are generally small when more than one electron needs to be transferred, or when multiple or strong bonds are broken, such as in the $\mathrm{N}_{2}, \mathrm{~N}_{3}^{-}$couple and in redox reactions of organic compounds.

A further consequence of the Butler-Volmer equation is illustrated by the curves in Fig. 19D.7, in which the cathodic and anodic current densities are plotted separately against the overpotential. When the overpotential is zero the two currents are equal. As the overpotential increases, the cathodic current decreases and the anodic current increases. Note, however, that for modest values of the overpotential $(|\eta f| \leq 3)$ both currents are significant.

If the exchange-current density is decreased then, although the curves have the same general shape, a much greater overpotential is needed in order to achieve the same current density. This dependence leads to a distinction between different kinds of electrodes. If the exchange-current density is 'large', then modest values of the overpotential lead to significant net current flow. Such an electrode is described as reversible in the sense that both the cathodic and anodic processes are taking place to a significant extent. In contrast, if the exchangecurrent density is 'small' a much larger overpotential is needed to achieve the same current. With such a value for the overpotential either the anodic or the cathodic current dominates. Such an electrode is described as irreversible.

Table 19D. 1 Exchange-current densities and transfer coefficients at $298 \mathrm{~K}^{*}$

\begin{center}
\begin{tabular}{llll}
\hline
Reaction & Electrode & $j_{0} /\left(\mathrm{Acm}^{-2}\right)$ & $\alpha$ \\
\hline
$2 \mathrm{H}^{+}+2 \mathrm{e}^{-} \rightarrow \mathrm{H}_{2}$ & Pt & $7.9 \times 10^{-4}$ &  \\
 & Ni & $6.3 \times 10^{-6}$ & 0.58 \\
 & Pb & $5.0 \times 10^{-12}$ &  \\
$\mathrm{Fe}^{3+}+\mathrm{e}^{-} \rightarrow \mathrm{Fe}^{2+}$ & Pt & $2.5 \times 10^{-3}$ & 0.58 \\
\hline
\end{tabular}
\end{center}

\footnotetext{\begin{itemize}
  \item More values are given in the Resource section.
\end{itemize}
}
\includegraphics[max width=\textwidth, center]{2025_02_27_d4b95d9718f0b9e2e6b9g-882(1)}

Figure 19D. 7 The dependence of the anodic (purple) and cathodic (blue) current density on the overpotential (for $\alpha=0.5$ ). The dotted lines are the corresponding current densities when the exchange-current density is one tenth of the value for the solid lines.

\subsection*{(b) Tafel plots}
When the overpotential is large and positive (in practice, $\eta \geq$ 0.12 V ), the anodic process is dominant. The current density is then given by the first term in eqn 19D.2:


\begin{equation*}
j=j_{0} \mathrm{e}^{(1-\alpha) f \eta} \text { hence } \ln j=\ln j_{0}+(1-\alpha) f \eta \tag{19D.5a}
\end{equation*}


A plot of the logarithm of the current density against the overpotential is called a Tafel plot. The slope, which is equal to $(1-\alpha) f$, gives the value of $\alpha$ and the intercept at $\eta=0$ gives the exchange-current density. If the overpotential is large and negative (in practice, $\eta \leq-0.12 \mathrm{~V}$ ), the cathodic process is dominant. The current density is then given by the second term in eqn 19D.2:


\begin{equation*}
j=j_{0} \mathrm{e}^{-\alpha f \eta} \text { hence } \ln j=\ln j_{0}-\alpha f \eta \tag{19D.5b}
\end{equation*}


In this case the slope of the Tafel plot is $-\alpha f$.

\subsection*{Example 19D. 1 Analysing data using a Tafel plot}
The following data refer to the anodic current through a platinum electrode of area $2.0 \mathrm{~cm}^{2}$ in contact with an $\mathrm{Fe}^{3+}, \mathrm{Fe}^{2+}$ aqueous solution at 298 K . Determine the exchange-current density and the transfer coefficient for the electrode process.

\begin{center}
\begin{tabular}{lccccc}
$\eta / \mathrm{mV}$ & 50 & 100 & 150 & 200 & 250 \\
$I / \mathrm{mA}$ & 8.8 & 25.0 & 58.0 & 131 & 298 \\
\end{tabular}
\end{center}

Collect your thoughts Because the current is anodic, the appropriate plot is of $\ln j$ against $\eta$. The intercept at $\eta=0$ is $\ln j_{0}$ and the slope is $(1-\alpha) f$. The current density is obtained by dividing the current by the area of the electrode.

The solution Draw up the following table:

\begin{center}
\begin{tabular}{lccccc}
$\eta / \mathrm{mV}$ & 50 & 100 & 150 & 200 & 250 \\
$j /\left(\mathrm{mA} \mathrm{cm}^{-2}\right)$ & 4.4 & 12.5 & 29.0 & 65.5 & 149 \\
$\ln \left(j /\left(\mathrm{mA} \mathrm{cm}^{-2}\right)\right)$ & 1.48 & 2.53 & 3.37 & 4.18 & 5.00 \\
\end{tabular}
\end{center}

The points are plotted in Fig. 19D.8. The data points for $\eta \geq 100 \mathrm{mV}$ give a straight line of extrapolated intercept 0.88 and slope 0.0165 . From the intercept it follows that $\ln \left(j_{0} /\left(\mathrm{mAcm}^{-2}\right)\right)=0.88$, so $j_{0}=2.4 \mathrm{~mA} \mathrm{~cm}^{-2}$. From the slope it follows that $(1-\alpha) f=0.0165 \mathrm{mV}^{-1}$, and because $f=F / R T=$ $38.9 \mathrm{~V}^{-1}, \alpha=0.58$.\\
\includegraphics[max width=\textwidth, center]{2025_02_27_d4b95d9718f0b9e2e6b9g-882(2)}

Figure 19D. 8 A Tafel plot is used to measure the exchangecurrent density (given by the extrapolated intercept at $\eta=0$ ) and the transfer coefficient (from the slope). The data are from Example 19D.1.

Comment. Note that the Tafel plot is nonlinear for $\eta<$ 100 mV ; in this region $\alpha f \eta=2.3$ and the condition $\alpha f \eta \gg 1$ is not satisfied.

Self-test 19D. 1 Repeat the analysis using the following cathodic current data recorded at 298 K and for an electrode of area $2.0 \mathrm{~cm}^{2}$ :\\
\includegraphics[max width=\textwidth, center]{2025_02_27_d4b95d9718f0b9e2e6b9g-882}

\subsection*{19D. 3 Voltammetry}
In the derivation of the Butler-Volmer equation it is assumed that the concentrations of the electroactive species are those of the bulk solution. Provided the current density is low, this approximation is likely to be valid because there will be only a small amount of the electroactive species converted from one form to another. However, this assumption fails at high current densities because the consumption of electroactive species close to the electrode results in a concentration gradient. The\\
diffusion of the species towards the electrode from the bulk is slow and may become rate determining; if this is the case, increasing the overpotential leads to no further increase in the current. This effect is called concentration polarization. Concentration polarization is important in the interpretation of voltammetry, the measurement of the current through an electrode as the applied potential difference is changed.

In linear-sweep voltammetry the current is measured as the applied potential difference is increased linearly with time (Fig. 19D.9a); Fig. 19D.9b shows typical data obtained in this way. If the applied potential difference becomes more negative as the sweep proceeds, the cathodic current due to reduction increases and the anodic current decreases. As the applied potential difference becomes more negative than the electrode potential, the overpotential becomes negative and the ButlerVolmer equation predicts that the cathodic current will increase exponentially. This accounts for the rapid increase in the current seen in Fig. 19D.9b.

The Butler-Volmer equation predicts that the current will go on rising as the potential difference becomes more negative, but in practice the current reaches a maximum and then declines. The explanation for this decrease is that the concentration of the electroactive species ( $\mathrm{Ox}^{+}$in this case) near the electrode is being depleted by the reduction process, thereby resulting in a decrease in the current. Diffusion of $\mathrm{Ox}^{+}$from the bulk solution towards the electrode replenishes its concentration there, and the balance between the rate of this process and the rate of the reduction determines the way in which the current declines at more negative potential differences.

If the bulk concentration of $\mathrm{Ox}^{+}$is increased, the peak current is increased, as shown in Fig. 19D.9b. The maximum current is proportional to the molar concentration of $\mathrm{Ox}^{+}$, so its concentration can be determined from the peak height after subtraction of an extrapolated baseline. It is also found that increasing the sweep rate increases the peak current. This\\
\includegraphics[max width=\textwidth, center]{2025_02_27_d4b95d9718f0b9e2e6b9g-883}\\
effect is a consequence of the balance between the rates of reduction and diffusion. The molecules close to the electrode become reduced, leading to a concentration gradient between the electrode surface and the bulk solution. This gradient drives the diffusion process: the greater the gradient, the faster is the diffusion and hence the greater is the current that can be sustained. A fast sweep leads to more rapid depletion of $\mathrm{Ox}^{+}$at the electrode, hence to a larger concentration gradient, faster diffusion, and therefore a larger current. According to the diffusion equation (Topic 16C), the net distance a molecule migrates is proportional to the square root of the time, and this dependence results in the peak current being proportional to the square root of the sweep rate.

In cyclic voltammetry the potential difference is applied with a triangular waveform (linearly up, then linearly down, Fig. 19D.10a) and the current is monitored. Cyclic voltammetry data are obtained at scan rates of about $50 \mathrm{mV} \mathrm{s}^{-1}$, so a scan over a range of 2 V takes about 80 s . A typical cyclic voltammogram is shown in Fig. 19D.10b; note that only the oxidized species is present at the start of the experiment. For the first part of the sweep, up to the time $t_{3}$ when the direction of the sweep is reversed, the shape of the curve and its interpretation is just as for a linear-sweep experiment (Fig. 19D.9b). After $t_{3}$ the potential difference becomes less negative and consequently the rate of the cathodic process (reduction) decreases; however, the rate of the anodic process (oxidation) increases. The current therefore decreases, reflecting the reduced rate of the cathodic process and the growing rate of the anodic process in which the layer of the reduced species formed on the electrode during the first part of the sweep is progressively oxidized. Eventually, the anodic process dominates and the direction of the current reverses. There then follows a maximum in the current followed by a decrease, which are explained in the same way as for the first part of the sweep.

The voltammogram shown in Fig. 19D.10b is for a reversible process in which only a small overpotential is required to\\
\includegraphics[max width=\textwidth, center]{2025_02_27_d4b95d9718f0b9e2e6b9g-883(1)}

Figure 19D. 10 (a) The change of potential difference with time used to record a cyclic voltammogram. The resulting current/ potential curve for: (b) a reversible electrode process, and (c) an irreversible process.\\
give a significant current. Such a voltammogram is broadly symmetrical about the standard potential of the couple. The potential differences at which the peak current is found for forward and reverse sweeps are symmetrical about the standard potential of the couple, allowing it to be estimated. An example of such a system is the $\left[\mathrm{Fe}(\mathrm{CN})_{6}\right]^{3-},\left[\mathrm{Fe}(\mathrm{CN})_{6}\right]^{4-}$ couple.\\
If the electrode process is irreversible (Fig. 19D.10c), a large overpotential is required to give a significant current, and the form of the cyclic voltammogram is affected. The first part of the voltammogram up to time $t_{3}$ is much the same as for the reversible case, except that it is shifted to a higher overpotential. When the sweep is reversed, the current decreases, reflecting the slower cathodic process. However, because a significant (positive) overpotential is needed for the anodic process to become important, the current falls back to zero as the cathodic process slows. The anodic process is not significant, and the current does not therefore change sign. The shape of the voltammogram therefore is very different from that for a reversible electrode.

The overall shape of a voltammogram gives details of the kinetics of the electrode process. Furthermore, the appearance of the curve may depend on the timescale of the sweep, because if the sweep is too fast some processes might not have time to occur.

\subsection*{Example 19D. 2 Analysing a cyclic voltammetry experiment}
The electroreduction of $p$-bromonitrobenzene in liquid ammonia is believed to occur by the following mechanism:

$$
\begin{array}{ll}
a & \mathrm{BrC}_{6} \mathrm{H}_{4} \mathrm{NO}_{2}+\mathrm{e}^{-} \rightarrow \mathrm{BrC}_{6} \mathrm{H}_{4} \mathrm{NO}_{2}^{-} \\
b & \mathrm{BrC}_{6} \mathrm{H}_{4} \mathrm{NO}_{2}^{-} \rightarrow \mathrm{C}_{6} \mathrm{H}_{4} \mathrm{NO}_{2}+\mathrm{Br}^{-} \\
c & \mathrm{C}_{6} \mathrm{H}_{4} \mathrm{NO}_{2}+\mathrm{e}^{-} \rightarrow \mathrm{C}_{6} \mathrm{H}_{4} \mathrm{NO}_{2}^{-} \\
d & \mathrm{C}_{6} \mathrm{H}_{4} \mathrm{NO}_{2}^{-}+\mathrm{H}^{+} \rightarrow \mathrm{C}_{6} \mathrm{H}_{5} \mathrm{NO}_{2}
\end{array}
$$

Figure 19D. 11 shows cyclic voltammograms recorded for this system at two different sweep rates. Interpret these in terms of the above mechanism.\\
\includegraphics[max width=\textwidth, center]{2025_02_27_d4b95d9718f0b9e2e6b9g-884(3)}

Figure 19D. 11 Cyclic voltammograms referred to in Example 19D.2: (a) for a slow sweep rate, and (b) for a fast sweep rate.

Collect your thoughts As described in the text, the shape of the voltammogram is influenced by whether or not the process is reversible on the timescale of the sweep. That distinction depends on the relative rates of the steps involved, both the redox steps and the other reactions.

The solution The voltammogram for the slow sweep, Fig. 19D.11a, is reminiscent of the one already described for an irreversible electrode process. However, in this case the process is irreversible because the reduced species formed in step $a$ goes on to react further, so $\mathrm{BrC}_{6} \mathrm{H}_{4} \mathrm{NO}_{2}^{-}$is not available to be oxidized in the second part of the sweep. The current therefore does not change sign. If step $b$ is fast, then $\cdot \mathrm{C}_{6} \mathrm{H}_{4} \mathrm{NO}_{2}$ may be reduced further in step $c$, but the process is made irreversible by the reduced species being removed in step $d$. At faster sweep rates, Fig. 19D.11b, the voltammogram is reminiscent of the one described for a reversible electrode process. In this case step $b$ is not fast enough to remove $\mathrm{BrC}_{6} \mathrm{H}_{4} \mathrm{NO}_{2}^{-}$, which remains available for oxidation during the second half of the sweep, leading to a change in the direction of the current.

Self-test 19D. 2 Suggest an interpretation of the cyclic voltammogram shown in Fig. 19D. 12 for the reduction of $\mathrm{ClC}_{6} \mathrm{H}_{4} \mathrm{CN}$ in acid solution in terms of the following reaction scheme:\\
a $\mathrm{ClC}_{6} \mathrm{H}_{4} \mathrm{CN}+\mathrm{e}^{-} \rightleftharpoons \mathrm{ClC}_{6} \mathrm{H}_{4} \mathrm{CN}^{-}$\\
$b \mathrm{ClC}_{6} \mathrm{H}_{4} \mathrm{CN}^{-}+\mathrm{H}^{+}+\mathrm{e}^{-} \rightarrow \mathrm{C}_{6} \mathrm{H}_{5} \mathrm{CN}+\mathrm{Cl}^{-}$(irreversible)\\
c $\mathrm{C}_{6} \mathrm{H}_{5} \mathrm{CN}+\mathrm{e}^{-} \rightleftharpoons \mathrm{C}_{6} \mathrm{H}_{5} \mathrm{CN}^{-}$\\
\includegraphics[max width=\textwidth, center]{2025_02_27_d4b95d9718f0b9e2e6b9g-884(2)}

Potential difference\\
Figure 19D. 12 The cyclic voltammogram referred to in Selftest 19D.2.\\
\includegraphics[max width=\textwidth, center]{2025_02_27_d4b95d9718f0b9e2e6b9g-884(4)}\\
\includegraphics[max width=\textwidth, center]{2025_02_27_d4b95d9718f0b9e2e6b9g-884}\\
\includegraphics[max width=\textwidth, center]{2025_02_27_d4b95d9718f0b9e2e6b9g-884(1)}\\
\includegraphics[max width=\textwidth, center]{2025_02_27_d4b95d9718f0b9e2e6b9g-884(5)}

\subsection*{19D. 4 Electrolysis}
To induce current to flow through an electrolytic cell and bring about a nonspontaneous cell reaction, the applied potential difference must exceed the zero-current potential by at least the cell overpotential, the sum of the overpotentials at the two\\
electrodes and the ohmic drop ( $I R_{\mathrm{s}}$, where $R_{\mathrm{s}}$ is the internal resistance of the cell) due to the current through the electrolyte. The additional potential difference needed to achieve a detectable rate of reaction might need to be large when the ex-change-current density at the electrodes is small. For similar reasons, a working galvanic cell generates a smaller potential difference than under zero-current conditions.

The relative rates of gas evolution or metal deposition during electrolysis can be estimated from the Butler-Volmer equation and tables of exchange-current densities. From eqn 19D.5b and assuming equal transfer coefficients, the ratio of the cathodic currents is


\begin{equation*}
\frac{j^{\prime}}{j}=\frac{j_{0}^{\prime}}{j_{0}} \mathrm{e}^{\left(\eta-\eta^{\prime}\right) \alpha f} \tag{19D.6}
\end{equation*}


where $j^{\prime}$ is the current density for electrodeposition, $j$ is that for gas evolution, and $j_{0}^{\prime}$ and $j_{0}$ are the corresponding exchangecurrent densities. This equation shows that metal deposition is favoured by a large exchange-current density and relatively high gas evolution overpotential (so $\eta-\eta^{\prime}$ is positive and large). Note that $\eta<0$ for a cathodic process, so $-\eta^{\prime}>0$. The exchange-current density depends strongly on the nature of the electrode surface, and changes in the course of the electrodeposition of one metal on another. A very crude criterion is that significant evolution or deposition occurs only if the overpotential exceeds about 0.6 V .

A glance at Table 19D. 1 shows the wide range of exchangecurrent densities for a metal/hydrogen electrode. The smallest exchange currents occur for lead and mercury: $1 \mathrm{pAcm}^{-2}$ corresponds to a monolayer of atoms being replaced in about 5 years. For such systems, a high overpotential is needed to induce significant hydrogen evolution. In contrast, the value for platinum $\left(1 \mathrm{mAcm}^{-2}\right)$ corresponds to a monolayer being replaced in 0.1 s , so significant gas evolution occurs for a much lower overpotential.

The exchange-current density also depends on the crystal face exposed. For the deposition of copper on copper, the (100) face has $j_{0}=1 \mathrm{~mA} \mathrm{~cm}^{-2}$, so for the same overpotential the (100) face grows at 2.5 times the rate of the (111) face, for which $j_{0}=$ $0.4 \mathrm{~mA} \mathrm{~cm}^{-2}$.

\subsection*{Quantitative laws of electrolysis}

The amount (in moles) of a substance deposited or dissolved at an electrode is determined by the total charge passed and the stoichiometry of the electrode half-reaction. If a constant current $I$ is applied for time $t$, the total charge is $Q = It$. The relationship between charge and moles of reaction is

\begin{equation*}
n_{\mathrm{substance}} = \frac{Q}{nF} = \frac{It}{nF} \tag{19D.9}
\end{equation*}

where $n$ is the stoichiometric number of electrons transferred per formula unit (e.g.\ $n = 2$ for $\mathrm{Cu}^{2+} + 2e^{-} \rightarrow \mathrm{Cu}$) and $F = 96\,485\ \mathrm{C\,mol^{-1}}$ is Faraday's constant. The mass of substance deposited is therefore

\begin{equation*}
m = \frac{M I t}{nF} \tag{19D.10}
\end{equation*}

where $M$ is the molar mass. The time required to deposit a mass $m$ is $t = mnF/(MI)$. These relations hold for each electrode independently; the same current $I$ flows through the whole circuit, so the amounts deposited at cathode and anode are related by their respective $n$ values. This principle underlies the use of coulometers (e.g.\ silver or copper coulometers) to measure total charge from deposited mass.

When two or more species could be deposited, the one with the less negative (more positive) electrode potential is preferentially reduced. In practice, the order of deposition is modified by overpotentials: the species with the smallest magnitude of cathodic overpotential at a given current density deposits first.

\subsection*{19D. 5 Working galvanic cells}
In working galvanic cells (those not balanced against an external potential), the overpotential leads to a smaller cell potential than under zero-current conditions. Furthermore, the cell potential is expected to decrease as current is generated because it is then no longer working reversibly and can therefore do less than maximum work.

Consider the cell $\mathrm{M}\left|\mathrm{M}^{+}(\mathrm{aq}) \| \mathrm{M}^{\prime+}(\mathrm{aq})\right| \mathrm{M}^{\prime}$, and ignore all the complications arising from liquid junctions. The potential\\
difference generated by the cell is $E^{\prime}=\Delta \phi_{\mathrm{R}}-\Delta \phi_{\mathrm{L}}$. Because the electrode potentials differ from their zero-current values by overpotentials, they can be written $\Delta \phi_{\mathrm{X}}=E_{\mathrm{X}}+\eta_{\mathrm{x}}$ where X is L or R for the left or right electrode, respectively. The cell potential is therefore


\begin{equation*}
E^{\prime}=E+\eta_{\mathrm{R}}-\eta_{\mathrm{L}} \tag{19D.7a}
\end{equation*}


To avoid confusion about signs ( $\eta_{\mathrm{R}}$ is negative, $\eta_{\mathrm{L}}$ is positive), and to emphasize that a working cell generates a lower potential difference than a zero-current cell, this expression is written as


\begin{equation*}
E^{\prime}=E-\left|\eta_{\mathrm{R}}\right|-\left|\eta_{\mathrm{L}}\right| \tag{19D.7b}
\end{equation*}


with $E$ the cell potential. The ohmic potential difference $I R_{s}$, where $R_{\mathrm{s}}$ is the cell's internal resistance, should also be subtracted


\begin{equation*}
E^{\prime}=E-\left|\eta_{\mathrm{R}}\right|-\left|\eta_{\mathrm{L}}\right|-I R_{\mathrm{s}} \tag{19D.7c}
\end{equation*}


The ohmic term is a contribution to the cell's irreversibility-it is a thermal dissipation term-so the sign of $I R_{\mathrm{s}}$ is always such as to reduce the potential difference in the direction of zero.

The overpotentials in eqn 19D. 7 can be calculated from the Butler-Volmer equation for a given current, $I$, being drawn. The equations are simplified by supposing that the areas, $A$, of the electrodes are the same, that only one electron is transferred in the rate-determining steps at the electrodes, that the transfer coefficients are both $\frac{1}{2}$, and that the highoverpotential limit of the Butler-Volmer equation may be used. Then from eqns 19D.5a, 19D.5b, and 19D.7c it follows that


\begin{equation*}
E^{\prime}=E-I R_{\mathrm{s}}-\frac{4 R T}{F} \ln \left(\frac{I}{A \bar{j}}\right) \quad \bar{j}=\left(j_{0 \mathrm{~L}} j_{0 \mathrm{R}}\right)^{1 / 2} \tag{19D.8}
\end{equation*}


where $j_{0 L}$ and $j_{O R}$ are the exchange-current densities for the two electrodes.

\subsection*{Brief illustration 19D. 2}
Suppose that a cell consists of two electrodes each of area $10 \mathrm{~cm}^{2}$ with exchange-current densities $5 \mu \mathrm{Acm}^{-2}$ and has internal resistance $10 \Omega$. At $298 \mathrm{~K}, R T / F=25.7 \mathrm{mV}$. The zerocurrent cell potential is 1.5 V . If the cell is producing a current of 10 mA , its working potential will be

$$
\begin{aligned}
E^{\prime}= & 1.5 \mathrm{~V}-\overbrace{(0.010 \mathrm{~A} \times 10 \Omega)}^{0.10 \mathrm{~V}} \\
& \overbrace{-4(0.0257 \mathrm{~V}) \ln \left(\frac{1000 \times 10 \mu \mathrm{~A}}{\left(10 \mathrm{~cm}^{2}\right) \times\left(5 \mu \mathrm{Acm}^{-2}\right)}\right)}^{0.54 \ldots \mathrm{~V}}=0.9 \mathrm{~V}
\end{aligned}
$$

where $1 \mathrm{~A} \Omega=1 \mathrm{~V}$ has been used. Various other factors that reduce the cell potential, such as the inability of reactants to diffuse rapidly enough to the electrodes, have been ignored.

Electric storage cells operate as galvanic cells while they are producing electricity but as electrolytic cells while they are being charged by an external supply. The lead-acid battery is an old device, but one well suited to the job of starting cars\\
(and the only one available). During charging the cathode reaction is the reduction of $\mathrm{Pb}^{2+}$ and its deposition as lead on the lead electrode. Deposition occurs instead of the reduction of the acid to hydrogen because the latter has a low exchangecurrent density on lead. The anode reaction during charging is the oxidation of $\mathrm{Pb}(\mathrm{II})$ to $\mathrm{Pb}(\mathrm{IV})$, which is deposited as the oxide $\mathrm{PbO}_{2}$. On discharge, the two reactions run in reverse. Because they have such high exchange-current densities the discharge can occur rapidly, which is why the lead battery can produce large currents on demand.

\subsection*{Checklist of concepts}
$\square$ 1. An electrical double layer consists of sheets of opposite charge at the surface of the electrode and next to it in the solution.2. Descriptions of the double layer include the Helmholtz layer model and the Gouy-Chapman model.3. The Galvani potential difference is the potential difference between the bulk of the metal electrode and the bulk of the solution.4. The current density at an electrode is expressed by the Butler-Volmer equation.\\
$\square$ 5. A Tafel plot is the plot of the logarithm of the current density against the overpotential (see below).6. Voltammetry is the study of the current through an electrode as a function of the applied potential difference.\\
7. To induce current to flow through an electrolytic cell and bring about a nonspontaneous cell reaction, the applied potential difference must exceed the cell potential by at least the cell overpotential.\\
$\square$ 8. In working galvanic cells the overpotential leads to a smaller potential difference.

\subsection*{Checklist of equations}
\begin{center}
\begin{tabular}{llll}
\hline
Property & Equation & Comment & Equation number \\
\hline
Butler-Volmer equation & $j=j_{0}\left\{\mathrm{e}^{(1-\alpha) f \eta}-\mathrm{e}^{-\alpha f \eta}\right\}$ &  & 19D.2 \\
Tafel plots & $\ln j=\ln j_{0}+(1-\alpha) f \eta$ & Anodic current density & 19D.5a \\
 & $\ln j=\ln j_{0}-\alpha f \eta$ & Cathodic current density & 19D.5b \\
Potential of a working galvanic cell & $E^{\prime}=E-I R_{\mathrm{s}}-(4 R T / F) \ln (I / A \bar{j})$ & $\bar{j}=\left(j_{0 \mathrm{~L}} j_{0 \mathrm{R}}\right)^{1 / 2}$ & 19 D .8 \\
\hline
\end{tabular}
\end{center}

\end{document}